% =====================================================================
%  Lezione zero 3 — Agile e Scrum nella pratica
%  Fonte: trascrizione della terza "lezione zero" (GD0_3): una lezione
%  introdotta da Diego Ricchiuti e tenuta da due ospiti, Simone Cicchetti
%  (ex studente, oggi level designer e producer) e Cristian Liberatore
%  (studente del terzo anno, game designer e project manager), che
%  raccontano come hanno applicato davvero Agile e Scrum nei loro progetti.
%  Pacchetti richiesti nel preambolo del libro:
%    \usepackage[italian]{babel}, \usepackage[T1]{fontenc},
%    \usepackage[utf8]{inputenc} (solo pdfLaTeX), \usepackage{booktabs}
%
%  Collocazione nel libro: parte "Lezioni zero", subito dopo
%  GD0_2_tempo_idee_e_leadership, con la stessa numerazione propria (0.3).
%  Nel file principale del libro, ad esempio:
%
%    \part{Lezioni zero: il game designer nell'industria}
%    \renewcommand{\thechapter}{0.\arabic{chapter}}
%    \setcounter{chapter}{0}
%    % =====================================================================
%  Lezione zero 1 — Lavorare in squadra: la cultura del feedback
%  Fonte: trascrizione della prima "lezione zero" (GD0), prima giornata di
%  un workshop di due giorni tenuto prima dell'inizio del corso di Game
%  Design, dedicato alla figura professionale e al lavoro in team.
%  Pacchetti richiesti nel preambolo del libro:
%    \usepackage[italian]{babel}, \usepackage[T1]{fontenc},
%    \usepackage[utf8]{inputenc} (solo pdfLaTeX), \usepackage{booktabs}
%
%  Collocazione consigliata nel libro: le lezioni zero precedono il corso,
%  quindi vanno in una parte separata PRIMA di GD1, con numerazione propria
%  (0.1, 0.2, ...) così che i capitoli del corso conservino i numeri 1, 2, 3.
%  Nel file principale del libro, ad esempio:
%
%    \part{Lezioni zero: il game designer nell'industria}
%    \renewcommand{\thechapter}{0.\arabic{chapter}}
%    \setcounter{chapter}{0}
%    % =====================================================================
%  Lezione zero 1 — Lavorare in squadra: la cultura del feedback
%  Fonte: trascrizione della prima "lezione zero" (GD0), prima giornata di
%  un workshop di due giorni tenuto prima dell'inizio del corso di Game
%  Design, dedicato alla figura professionale e al lavoro in team.
%  Pacchetti richiesti nel preambolo del libro:
%    \usepackage[italian]{babel}, \usepackage[T1]{fontenc},
%    \usepackage[utf8]{inputenc} (solo pdfLaTeX), \usepackage{booktabs}
%
%  Collocazione consigliata nel libro: le lezioni zero precedono il corso,
%  quindi vanno in una parte separata PRIMA di GD1, con numerazione propria
%  (0.1, 0.2, ...) così che i capitoli del corso conservino i numeri 1, 2, 3.
%  Nel file principale del libro, ad esempio:
%
%    \part{Lezioni zero: il game designer nell'industria}
%    \renewcommand{\thechapter}{0.\arabic{chapter}}
%    \setcounter{chapter}{0}
%    % =====================================================================
%  Lezione zero 1 — Lavorare in squadra: la cultura del feedback
%  Fonte: trascrizione della prima "lezione zero" (GD0), prima giornata di
%  un workshop di due giorni tenuto prima dell'inizio del corso di Game
%  Design, dedicato alla figura professionale e al lavoro in team.
%  Pacchetti richiesti nel preambolo del libro:
%    \usepackage[italian]{babel}, \usepackage[T1]{fontenc},
%    \usepackage[utf8]{inputenc} (solo pdfLaTeX), \usepackage{booktabs}
%
%  Collocazione consigliata nel libro: le lezioni zero precedono il corso,
%  quindi vanno in una parte separata PRIMA di GD1, con numerazione propria
%  (0.1, 0.2, ...) così che i capitoli del corso conservino i numeri 1, 2, 3.
%  Nel file principale del libro, ad esempio:
%
%    \part{Lezioni zero: il game designer nell'industria}
%    \renewcommand{\thechapter}{0.\arabic{chapter}}
%    \setcounter{chapter}{0}
%    \input{Seminari/GD0_1_cultura_del_feedback}
%    % ... altre lezioni zero ...
%    \part{Il corso di Game Design}
%    \renewcommand{\thechapter}{\arabic{chapter}}
%    \setcounter{chapter}{0}
%    \input{GD1/GD1_introduzione_al_game_design}
%    % ...
%
%  (Con hyperref aggiungere anche \renewcommand{\theHchapter}{z.\arabic{chapter}}
%  nella parte zero, per evitare ancore duplicate.)
%  Tutte le etichette di questo capitolo hanno il prefisso "z1:" per non
%  collidere con quelle dei capitoli del corso.
% =====================================================================

\chapter[La cultura del feedback]{Lavorare in squadra: la cultura del feedback}
\label{cap:z1-feedback}

Questo capitolo apre le \emph{lezioni zero}: incontri che precedono l'inizio
vero e proprio del corso di Game Design e che non riguardano una tecnica di
design in senso stretto, ma \textbf{la figura professionale} di chi lavora
nell'industria dei videogiochi. Prima di imparare a progettare un gioco
conviene infatti capire in che tipo di ambiente lo si progetterà: un ambiente
fatto di team, di gerarchie, di scadenze e, soprattutto, di persone che devono
riuscire a lavorare insieme.

La lezione è la prima giornata di un \textbf{workshop di due giorni} (circa
dieci ore in totale, pause pranzo comprese) tenuto da Stefano
Cozzi\footnote{Il cognome è riportato come compare nella trascrizione
automatica, che in più punti storpia i nomi propri; andrebbe verificato.
Nella stessa trascrizione lo studio per cui lavora appare come
``Stormine'', ``Stormline'' e ``Stormind'': si tratta di Stormind Games,
studio con sede a Catania, coerentemente con quanto detto in lezione.}, un
professionista con circa quindici anni di esperienza nel settore, tra
Blizzard Entertainment, Riot Games e, oggi, Stormind Games. Il workshop è
nato da un confronto con i coordinatori dell'accademia\footnote{Nella
trascrizione l'accademia è resa in modi diversi (``Hive'', ``a dive'',
``D.I.V.''); i coordinatori citati sono Enrico e Andrea, insieme al
direttore Diego Ricchiuti. Nel resto del capitolo si usa semplicemente
``l'accademia'', come negli altri capitoli.}: serviva un momento dedicato a
\emph{come ci si comporta in un ambiente professionale}, cosa che riguarda sia
il lavoro futuro sia i progetti di gruppo svolti durante il percorso di
studi, dove ci si deve già abituare a comportarsi ``come se si fosse in
un'azienda''.

La prima giornata ha un unico grande tema, sezionato in più parti:
\textbf{la cultura del feedback}. La seconda giornata (che sarà oggetto della
prossima lezione zero) sarà più pratica: \emph{time management},
\emph{brainstorming} e il ``guidare da dentro'' un team anche senza esserne
il leader (sezione~\ref{sec:z1-domani}).

Il percorso del capitolo segue quello della lezione:
\begin{enumerate}
  \item il formato del workshop e il percorso professionale del relatore
        (sezione~\ref{sec:z1-relatore});
  \item perché le aziende tengono tanto al lavoro di squadra e al
        \emph{fit} culturale (sezione~\ref{sec:z1-squadra});
  \item che cos'è il feedback e da che cosa dipende il modo di darlo
        (sezione~\ref{sec:z1-cose-feedback});
  \item le cinque fasi di vita di un team e il valore del
        \emph{post-mortem} (sezioni~\ref{sec:z1-fasi-team}
        e~\ref{sec:z1-qa-fasi});
  \item il confronto con i colleghi e alcune domande sul lavoro nei grandi
        studi (sezioni~\ref{sec:z1-rivalita} e~\ref{sec:z1-industria});
  \item il circolo virtuoso tra feedback, miglioramento e fiducia, e i suoi
        effetti sulla \emph{retention} e sulla cultura aziendale
        (sezioni~\ref{sec:z1-triangolo} e~\ref{sec:z1-cultura});
  \item il ``gioco del traditore'' e il principio del
        \textbf{default to trust} (sezioni~\ref{sec:z1-among-us}
        e~\ref{sec:z1-default-trust});
  \item le traiettorie di carriera e la valutazione delle performance
        (sezione~\ref{sec:z1-traiettorie});
  \item come si dà, come si riceve e quando si dà un feedback
        (sezioni~\ref{sec:z1-dare}, \ref{sec:z1-ricevere}
        e~\ref{sec:z1-quando});
  \item il feedback come regalo, le domande finali e le anticipazioni
        della seconda giornata (sezioni~\ref{sec:z1-regalo},
        \ref{sec:z1-qa-finale} e~\ref{sec:z1-domani}).
\end{enumerate}

Una premessa di metodo: il relatore ha alternato continuamente teoria,
domande alla classe e racconti della propria esperienza. I racconti non sono
digressioni: sono gli esempi concreti con cui i principi vengono resi
tangibili, e per questo sono riportati con cura.

% ---------------------------------------------------------------------
\section{Il workshop e il relatore}
\label{sec:z1-relatore}

\subsection{Formato e regole di partecipazione}

Il relatore ha cercato di ``riassumere quindici anni di esperienza in dieci
ore''. Ha materiale per circa il doppio delle ore disponibili: le cose
strettamente necessarie vengono trattate per prime, il resto è tenuto ``in
riserva'' come \emph{nice to have}. Le regole del workshop:
\begin{itemize}
  \item le domande si fanno \textbf{alla fine di ogni slide}, così da chiudere
        prima il discorso; non c'è un'unica sessione finale di domande e
        risposte (anche se ci sarà), perché il relatore preferisce fermarsi
        spesso e non fare ``uno spiegone di dieci ore'';
  \item le domande non devono per forza riguardare la slide appena vista: se
        una slide dà lo spunto per altro, se ne può parlare; al massimo sarà
        il relatore a rispondere in breve o a rimandare a una slide
        successiva;
  \item si può chiedere anche di esperienze passate, opinioni e sensazioni;
        se per NDA o per serietà professionale non può rispondere, lo dirà
        esplicitamente (``non lasciate che sia un \emph{blocker}'');
  \item su Zoom si usa l'alzata di mano o la chat, che il relatore tiene
        aperta sul secondo monitor;
  \item non ci sono pause ufficiali: la lezione è registrata e ognuno è
        libero di assentarsi quando serve, cercando però di non disturbare il
        flusso della lezione;
  \item le slide verranno fornite integralmente tramite la segreteria
        dell'accademia (Vanessa), insieme alle registrazioni: le slide senza
        il parlato servono poco, ma sono un buon promemoria.
\end{itemize}

\subsection{Il percorso professionale}

Il relatore si presenta (scherzosamente) come la ``guida escursionistica''
che accompagna la classe nella ``foresta incantata'' della cultura del
feedback, perché per molti è un territorio sconosciuto. Il suo percorso:
\begin{itemize}
  \item in Italia ha lavorato per aziende che fornivano servizi a terzi, tra
        cui il lavoro su \emph{Warhammer Online}, e per un canale televisivo
        sul digitale terrestre dedicato ai videogiochi\footnote{Nella
        trascrizione: ``Digital Game Network Channel''. Il nome esatto del
        canale andrebbe verificato.};
  \item a \textbf{Blizzard Entertainment} (circa sei anni, a Parigi) ha fatto
        il \emph{community manager} per l'Italia (\emph{StarCraft~II},
        \emph{World of Warcraft} da \emph{Mists of Pandaria} in poi,
        \emph{Hearthstone}, \emph{Diablo~III}, le espansioni di
        \emph{StarCraft}, il lancio di Battle.net), per poi diventare il
        manager dei community manager europei;
  \item a \textbf{Riot Games} (circa cinque anni, a Dublino) è stato
        brevemente community manager italiano e poi \emph{publishing
        manager} di un team di circa trenta persone per l'Europa. Il
        \emph{publishing} non è solo community: è l'unione di
        \textbf{community, marketing e PR};
  \item a \textbf{PG Esports} si è occupato di eventi dal vivo e online
        (ad esempio la Milano Games Week e tornei di \emph{Valorant} e
        \emph{League of Legends}), restando in contatto con Riot e Blizzard;
  \item da circa un anno è tornato nello sviluppo, a \textbf{Stormind
        Games}, dove oggi si occupa di risorse umane e organizzazione del
        lavoro; lo studio verrà presentato più ampiamente al \emph{Career
        Day} dell'accademia insieme alla sua HR manager;
  \item fa inoltre consulenze aziendali anche al di fuori dei videogiochi e
        insegna in master universitari (IULM di Milano).
\end{itemize}

Il suo ruolo ricorrente, sia a Blizzard sia a Riot, è stato quello di
\emph{Mr.~Wolf} (``per chi ha visto \emph{Pulp Fiction}''): quello che
risolve i problemi. Veniva inserito nei team più in difficoltà, ad esempio
dal punto di vista delle performance, per sistemarli e passare al successivo.
Era presente in alcuni dei momenti più critici, come il lancio di
\emph{Diablo~III} o il lancio di espansioni di \emph{World of
Warcraft}\footnote{La trascrizione riporta ``il lancio con i Gary's'' e
``World of Trenor'': si tratta verosimilmente di \emph{Mists of Pandaria} e
di \emph{Warlords of Draenor}.}, a fare \emph{damage control}.

La scoperta fondamentale di questi anni è che i problemi da risolvere
\textbf{quasi mai riguardavano il prodotto} (non era certo lui a sistemare i
server di \emph{Diablo~III}), ma i team che c'erano dietro: persone che
litigavano, non collaboravano, non riuscivano a stare nella stessa stanza
senza urlarsi contro. E nel 90\% dei casi tutto questo si riduce a
\textbf{problemi di comunicazione}: è raro che perfino le incompatibilità
caratteriali non siano risolvibili parlando.

\subsection{Struttura non vuol dire burocrazia}

A Stormind il relatore sta accompagnando un team passato da 30 a 95 persone
in meno di un anno, attraverso tutti i cambiamenti di cui parla il workshop.
Darsi una struttura viene spesso percepito male, come ``più processi, più
procedure'', ma bisogna distinguere:
\begin{itemize}
  \item un processo come ``timbrare il cartellino'' può essere davvero
        invasivo e noioso;
  \item definire \emph{come funziona un documento di game design, chi lo
        approva, chi lo visiona e chi no} è invece una struttura necessaria.
\end{itemize}
Uno studio di 30 persone nella stessa sede può farne a meno: quando c'è un
dubbio ci si alza e si va alla scrivania del collega. Uno studio di 95
persone sparse per l'Italia (e in parte per l'Europa) non può più farlo:
servono una struttura diversa, un modo di lavorare diverso e una
pianificazione del futuro diversa.

% ---------------------------------------------------------------------
\section{Perché il lavoro è sempre lavoro di squadra}
\label{sec:z1-squadra}

\subsection{Non si va in ufficio per fare amicizia}

Il primo messaggio è netto: \textbf{non si va in ufficio per fare amici} o
per trovare compagni di vita. Si va in ufficio perché c'è un lavoro da fare, e
lo scopo è riuscire a farlo \emph{anche con chi sta antipatico}. Con le
persone simpatiche è facile: diventeranno amici, e non sarà mai un problema.
È con chi sta antipatico che le cose si complicano, ed è lì che emergono i
\textbf{veri professionisti}: quelli che riescono a lavorare con chiunque
senza che l'antipatia pesi sul modo in cui lavorano e collaborano, mettendo il
prodotto e il team davanti al proprio ego e alle proprie preferenze personali.

La frase che fa da filo conduttore a tutto il workshop (e che il relatore
promette di riproporre ``fino alla nausea'') è:

\begin{quote}
  \emph{La capacità di un gruppo di persone di fare cose straordinarie
  dipende dal modo in cui queste persone operano insieme.}
\end{quote}

Le aziende raramente pensano che l'innesto di una singola persona possa
cambiare il destino di un progetto. Questo è tanto più vero quanto più grande
è il team, ma vale anche per i piccoli studi, anche per un gruppo di cinque
persone. Come negli sport di squadra: l'individuo conta fino a un certo punto,
e anche se si ha Dončić in squadra, se si gioca uno contro cinque si perde
comunque\footnote{Il riferimento è alle finali NBA in corso nei giorni del
workshop tra Dallas Mavericks e Boston Celtics (giugno 2024), il che
permette di datare approssimativamente la lezione.}.

\subsection{Lo \emph{useless genius} e il costo delle persone tossiche}

Molti dipartimenti HR usano la definizione di \textbf{useless
genius}\footnote{Il concetto è più noto in letteratura manageriale come
\emph{brilliant jerk}; in lezione è usata l'espressione \emph{useless
genius}.}: si può essere bravissimi, a volte i migliori nel proprio lavoro,
ma se non ci si sa comportare, relazionare e operare in un gruppo si è
fondamentalmente inutili ai fini del gruppo. Nel migliore dei casi inutili,
nel peggiore \emph{dannosi}.

In passato le aziende tendevano a ``chiudere un occhio'': la persona è
fondamentale, la teniamo. Oggi la domanda è diventata: \emph{quanto costa
all'azienda, oltre al suo stipendio, tenere una persona che crea un ambiente
tossico e peggiora chi le sta intorno?} Ci sono costi sommersi che le aziende
cominciano a vedere, e la tendenza è liberarsi di queste persone il prima
possibile. Il ragionamento è semplice:
\begin{itemize}
  \item le \textbf{competenze si possono trasmettere} o sostituire: si trova
        qualcuno che fa lo stesso lavoro magari nel doppio del tempo, oppure
        più persone che se lo dividono;
  \item il \textbf{carattere delle persone} invece cambia molto raramente.
\end{itemize}

\subsection{Il fit culturale}

Quando qualcuno \textbf{non è un \emph{fit} culturale} per un'azienda è
meglio separarsi presto. Ma non essere adatti a \emph{quell'}azienda non
significa non essere adatti a lavorare: è perfettamente legittimo che una
persona inadatta a un'azienda sia adattissima a un'altra.

Il relatore l'ha vissuto in prima persona. Dopo sei anni a Blizzard viene
contattato da Riot e affronta un processo di selezione lunghissimo (``quasi
undici colloqui''), al termine del quale viene scartato. Scrive allora una
mail che gli ha ``cambiato la traiettoria di carriera'': chiede
semplicemente perché non è stato ritenuto adatto. La risposta: in nessun
colloquio si era esposto abbastanza da far capire \emph{perché} volesse
lasciare Blizzard, e quindi pensavano che in realtà stesse bene dov'era. Il
motivo era che, per educazione professionale, non gli piace parlare male del
proprio posto di lavoro; portata all'eccesso, però, questa riservatezza era
diventata un ``è tutto bellissimo, ma vorrei provare Riot'', senza un perché.
Lo invitano comunque a un incontro dal vivo (``dal vivo è diverso''), e alla
fine viene assunto.

Una volta arrivato, i sei mesi di prova (\emph{probation}) sono serviti tutti
per smettere di ragionare e lavorare ``come a Blizzard''. La differenza
culturale era enorme, tanto che in Riot esisteva un gruppo di supporto
interno tra ex dipendenti Blizzard: un collega gliel'ha descritto come
``uscire da una relazione tossica e doversi riabituare a una relazione
normale''\footnote{Il relatore precisa che parla della Blizzard di allora,
che non esiste più: oggi è parte di Microsoft e nessuno dei suoi ex colleghi
vi lavora ancora, anche a causa dei licenziamenti di massa.}.
\begin{description}
  \item[Blizzard] era \emph{developer first}: il focus era chiaramente sugli
        sviluppatori, prima di tutti gli altri. Dopo sei anni c'erano ancora
        meeting a cui non poteva partecipare perché non abbastanza
        \emph{senior}, informazioni che non poteva sapere, piani che non
        poteva vedere.
  \item[Riot] era \emph{player first}: prima i giocatori, e i dipendenti come
        mezzo per trasferire il feedback dei giocatori dentro il gioco. Il
        secondo giorno lo portano in una stanza e gli mostrano il piano dei
        cinque anni successivi. Lui, sorpreso, pensa a un errore (``sono qui
        da oggi, non vorrei mettervi nei guai''); la risposta è: ``sei un
        dipendente Riot o no? Allora non ci sono livelli di accesso: devi
        sapere queste cose per fare meglio il tuo lavoro''.
\end{description}

Da questa esperienza ha imparato molto sulla \textbf{trasparenza}: condividere
le informazioni aiuta il lavoro. Non sempre, però: a volte non si condivide
qualcosa non per segretezza, ma perché non è ancora pronta per essere
presentata, esattamente come non si mostra un gioco quando è ancora solo
teoria (``te lo faccio vedere quando ho qualcosa di tangibile'').

La lezione per gli studenti è doppia: le aziende valuteranno il \emph{fit}
culturale sempre di più, ma capire se si è un \emph{fit} richiede anche
\textbf{lavoro da parte del candidato}. Bisogna avere la presenza di spirito
di dire: ``conosco i valori di questa azienda, non mando nemmeno il CV,
perché so che durerei un mese e mezzo''.

% ---------------------------------------------------------------------
\section{Che cos'è il feedback}
\label{sec:z1-cose-feedback}

\subsection{Le prime risposte della classe}

La prima domanda alla classe è: \emph{quando sentite la parola
``feedback'', a cosa pensate?} Le risposte raccolte, a voce e in chat:
\begin{itemize}
  \item un \textbf{riscontro} su ciò che si fa e su come lo si fa, sul lavoro
        o sul modo in cui ci si interfaccia con il gruppo; può essere
        lavorativo oppure comunicativo e comportamentale. Chi ha risposto lo
        vede come qualcosa che, se l'azienda non dà, va attivamente
        richiesto, perché è importante sapere come sta andando per
        migliorarsi;
  \item né positivo né negativo, ma \textbf{costruttivo}: in entrambi i casi
        serve a migliorarsi;
  \item uno \emph{scambio di opinioni}, un \emph{riscontro produttivo};
  \item un feedback ``di reazione'', non verbale: si fa una cosa e si
        osserva come cambia l'ambiente o il comportamento delle persone
        intorno (anche perché a voce la gente spesso si tiene le cose
        negative);
  \item per un designer, anche il \textbf{linguaggio del corpo}: osservando
        un giocatore giocare si può dedurre l'emozione che prova (frustrazione,
        incomprensione).
\end{itemize}

Sulla forma (scritto o orale?) le opinioni divergono: c'è chi lo preferisce
\emph{mille volte} scritto; chi pensa che dipenda dalle dimensioni
dell'azienda (orale nelle piccole, scritto in quelle con team di centinaia di
persone); chi pensa che dipenda dall'importanza del feedback (scritto se
tecnico o pieno di dati, da tenere sotto gli occhi; orale se è un consiglio
veloce).

Il relatore usa due esempi per mostrare che \emph{dipende}:
\begin{itemize}
  \item un feedback sulla programmazione ha senso darlo scritto, magari
        commentando direttamente il codice, più che ``andando a prendere un
        caffè'';
  \item ma se il feedback è ``ieri, durante il meeting, ho percepito male il
        fatto che mi interrompessi di continuo'', ce lo si aspetta
        \textbf{orale}.
\end{itemize}

\subsection{Da che cosa dipende il modo di dare feedback}

Lo \emph{spoiler} della giornata è che \textbf{non esiste una risposta univoca}
su come dare feedback. Il workshop non insegna ``il'' modo, ma fornisce
\emph{tutte le opzioni}, perché il modo giusto cambia in base a molti fattori:
\begin{itemize}
  \item a chi lo si dà: un collega, un superiore, il proprio capo, un proprio
        \emph{report} (un sottoposto);
  \item se è personale o professionale;
  \item se è la prima volta o la terza volta che si dà lo stesso feedback;
  \item se riguarda un progetto importante o secondario.
\end{itemize}

Manca però l'elemento fondamentale, che la classe individua subito:
\textbf{l'altro}. A parità di tutti gli altri fattori, il \emph{carattere
della persona che riceve il feedback} dovrebbe far cambiare il modo in cui lo
si dà. Per questo, prima della lezione, agli studenti era stato chiesto di
fare un test di personalità: non un test di psicologia avanzata, ma uno
strumento per rendersi conto di quanta varietà ci sia anche solo in una
classe di 53 persone (il che non vuol dire trovare 52 modi diversi di dare
feedback, ma lavorarci un po' sì).

Il relatore porta se stesso come esempio. Preferisce:
\begin{itemize}
  \item i feedback \textbf{positivi scritti}, perché è a disagio a ricevere
        complimenti in pubblico;
  \item i feedback sul proprio lavoro \textbf{diretti}: ``vai subito al
        punto e non girarmi intorno''.
\end{itemize}
È molto orgoglioso del suo lavoro e, quando qualcuno lo critica, la sua
primissima reazione istintiva è ``come ti permetti''. Negli anni ha imparato
a farla durare un secondo, ma chi gli dà feedback di persona dice di
vederla ancora negli occhi. Per mitigarla, dice esplicitamente a chi gli dà
feedback: \emph{``ignora il modo in cui reagisco nei primi cinque secondi,
vai avanti, poi mi passa''}. È considerato da tutti molto aperto al feedback
e non ha problemi a dire ``hai ragione, facciamo diversamente'': ma ci sono
reazioni istintive e caratteriali che è difficile controllare al cento per
cento.

\subsection{Con chi si lavora meglio}

Guardando le persone con cui si lavora abitualmente (anche nei progetti
dell'accademia), si noterà che si lavora meglio con due profili:
\begin{enumerate}
  \item chi è \textbf{esattamente come noi}, ma qui è un 50 e 50: alcuni
        caratteri vanno d'accordo con i propri simili, altri non vogliono
        nessuno con lo stesso carattere nel gruppo (e si potrebbe rientrare
        in questa categoria);
  \item chi è \textbf{affine}: ha qualcosa in comune con noi, ma non tutto,
        e proprio per questo può dare punti di vista differenti.
\end{enumerate}
Il lavoro della giornata è capire come valorizzare il secondo caso: ottenere
feedback da persone con caratteri, visioni e background diversi, per
completare una visione del proprio lavoro che è necessariamente limitata. È
come esplorare una mappa con la nebbia di guerra: un'altra persona, al di là
del muro, vede cose che noi non vediamo.

\subsection{Il feedback come somma di piccoli miglioramenti}

Per il relatore il feedback è questo: \textbf{applicare tanti piccoli
miglioramenti periodici}, che si sommano nel tempo. Migliorare dello 0,1\%
ogni giorno porta a fine anno a un miglioramento enorme, e questo vale per i
prodotti, per i giochi, per le persone e per i team.

\paragraph{L'esempio del ciclismo britannico.} La squadra di ciclismo
britannica, fino a una quindicina d'anni fa, non aveva vinto quasi nulla.
Cambiato il team manager, tutti si aspettavano grandi cambiamenti sportivi
(convocazioni, strategie di gara). Arrivarono invece \emph{minuscoli
cambiamenti non sportivi}: un altro chef in trasferta, materassi migliori nei
camper per riposare meglio, un sellino più comodo a fine tappa. Cambiamenti
marginali che da soli non avrebbero cambiato niente; eppure, dopo due o tre
anni di questo lavoro quotidiano, la squadra ha cominciato a vincere a
ripetizione, Olimpiadi comprese, cambiando il modo in cui il ciclismo è
visto in Gran Bretagna\footnote{È la nota strategia delle \emph{marginal
gains} associata a Dave Brailsford (British Cycling e Team Sky); il nome non
è citato in lezione.}. Certo, si è investito anche sulle persone giuste, ma
tutto è partito dai piccoli miglioramenti periodici.

Per gli studenti vale lo stesso: chiarire un conflitto con una persona con
cui si ha ancora poco a che fare, o una situazione che ha lasciato
``l'amaro in bocca'', sembra inutile oggi, ma darà i suoi frutti tra tre, sei,
nove mesi. Sono \emph{investimenti sul futuro}.

\paragraph{Peggiorare ogni giorno.} Esiste anche il percorso inverso:
lasciar correre, ``buttare sotto il tappeto'' le cose perché ``tanto è
inutile parlarne''. Queste cose tornano sempre fuori dopo otto mesi, come il
\emph{libro dei rancori} dei nani di Warhammer: ``ti ricordi quando, un anno
e mezzo fa, all'inizio del progetto, mi hai detto quella cosa? Mi sono
offeso''. A quel punto è quasi impossibile risolvere, perché nessuno ricorda
più con precisione cosa sia successo (``ma quale cena?'', ``forse era in
pizzeria'', ``no, eravamo da un'altra parte''). \textbf{Tenersi le cose nel
cassetto per tirarle fuori quando servono è sempre sconsigliato}: non servono
ed è impossibile usarle in modo costruttivo.

% ---------------------------------------------------------------------
\section{Le cinque fasi di vita di un team}
\label{sec:z1-fasi-team}

\subsection{Perché conoscerle}

Quasi tutti gli studenti lavorano già in gruppo (al primo anno su un gioco da
tavolo, negli anni successivi su un videogioco) e si saranno accorti che un
team attraversa delle fasi: il modo di lavorare insieme cambia nel tempo.
Spesso ci si fa prendere di sorpresa: ``eravamo così affiatati, cosa è
successo?''.

Un esempio da una consulenza fuori dai videogiochi: un'azienda farmaceutica
gli affida un team commerciale (rappresentanti che presentano medicinali a
ospedali, medici e privati). Il problema: ``prima andava tutto bene; ora che
siamo in diciotto dovremmo lavorare meno e meglio, e invece va peggio''. Va
peggio perché le persone nuove sono appena arrivate, mentre il management si
aspettava già di vedere i risultati dell'investimento.

Fino a pochi anni fa le fasi di un team venivano disegnate come una
\textbf{linea dritta verso l'alto}: le fasi erano le stesse, ma viste come un
miglioramento costante. Non si vedevano le parti \emph{fisiologiche} del
lavoro in team, cioè anche i peggioramenti. Accettarle è il primo passo per
accettare i fallimenti: nessun percorso personale o di team è lineare.

Il relatore ammette che il suo modo di lavorare è nato per il 70\% dal
proprio impegno a imparare, fare bene e rispettare le persone, e per il 30\%
dall'aver lavorato con ``incapaci colossali'' che gli hanno insegnato come
\emph{non} fare le cose. A Blizzard un manager gli fece capire che, pur non
sapendo ancora che tipo di \emph{people manager} sarebbe diventato, sapeva
benissimo quale tipo di manager \emph{non} voleva essere.

\subsection{Forming, storming, norming, performing, adjourning}

Le fasi sono cinque, un po' come le fasi del lutto. Molte aziende ne
considerano solo quattro (``e sono quelle che solitamente chiudono''); le
aziende al passo con i tempi ne considerano cinque\footnote{Il modello è
quello di Bruce Tuckman (1965), a cui la quinta fase, \emph{adjourning}, fu
aggiunta nel 1977. Il nome di Tuckman non è citato in lezione; la
trascrizione riporta la quinta fase come ``journing''.}.

\begin{description}
  \item[Forming] Si viene messi insieme e si comincia a lavorare. Non si sa
        nulla degli altri; spesso non si conoscono nemmeno i nomi, ed è
        normale.
  \item[Storming] Si cominciano a ``dare le testate come le capre in
        montagna per determinare il territorio'': questa cosa la faccio io,
        questa la vorrei fare così, chi decide? Si stabilisce chi fa cosa e
        come si divide il lavoro. Nel grafico è una \textbf{fase di
        peggioramento}: che il team litighi o discuta civilmente non cambia
        nulla, il team deve investire parte di tempo ed energie per risolvere
        queste questioni. È normale e \textbf{non è evitabile}.
  \item[Norming] Si comincia a capire come lavorano gli altri, cosa sanno
        fare e come far funzionare il gruppo; iniziano i miglioramenti. È una
        forma di \emph{accettazione} degli altri, che deriva dall'aver capito
        cosa devo fare io e cosa devono fare loro.
  \item[Performing] C'è uno sforzo \textbf{proattivo}: ``questa cosa
        toccherebbe a me, ma so che Alessio la fa molto meglio, quindi gli
        chiedo aiuto''. Si sfruttano i punti di forza degli altri, perché li
        si conosce e perché si sa che passare il testimone non avrà effetti
        negativi (il capo non dirà ``perché non l'hai fatta tu?'').
  \item[Adjourning] Il progetto finisce, si chiude o viene riorganizzato, il
        team viene sciolto. Quello che andrebbe fatto è analizzare
        \textbf{cosa ha funzionato, cosa no e cosa migliorare} per il
        progetto successivo.
\end{description}

\subsection{La fase che tutti saltano}

``Spoiler incredibile'': l'\emph{adjourning} è la fase che molte aziende
saltano, \textbf{perché c'è sempre il progetto dopo}.

A PG Esports, per la Milano Games Week (come per qualsiasi evento dal vivo a
cadenza annuale: Lucca Comics, Etna Comics, Comicon\dots), finita l'edizione
di un anno si cominciava dal giorno dopo a lavorare a quella successiva,
perché un anno sembra tanto ma non basta mai. Nei videogiochi è lo stesso:
finito un gioco, a parte le patch di mantenimento e del \emph{day one}, si
passa subito al successivo.

Il cambiamento introdotto dal relatore (in parte già praticato, ma da lui
strutturato) è stato: finita la Games Week, \textbf{tutti in ferie per una
settimana}; poi ci si ritrova per definire cosa ha funzionato e cosa no, con
un \textbf{post-mortem} (o \emph{debrief}); solo dopo si inizia a lavorare
all'edizione successiva. Senza questa fase manca la parte di
\emph{learning}, cioè di ciò che si è imparato dal progetto.

Il consiglio vale per gli studenti: quando si finisce il gioco da tavolo, e
poi il gioco digitale, bisogna trovare il tempo per il post-mortem,
altrimenti si butta via \textbf{il 90\% dell'utilità del progetto}. Una
struttura semplice per farlo è descritta nella sezione~\ref{sec:z1-qa-finale}.

% ---------------------------------------------------------------------
\section{Domande sulle fasi del team}
\label{sec:z1-qa-fasi}

\subsection{Lo storming che ritorna}

\paragraph{Capita che lo storming si ritagli spazio anche nelle altre fasi?}
Lo \emph{storming} è una fase precisa, che però molti manager, coordinatori e
team cercano di \textbf{saltare}: ``se io sono bravo come manager e loro sono
bravi come team, passeremo da forming a norming senza problemi''. È un errore,
perché senza una fase in cui si discute davvero non si sviluppa una
\emph{skill} fondamentale: \textbf{saper discutere}. Tutti conoscono persone
non abituate a farlo, per cui ogni minimo disaccordo diventa un affronto
personale (``avevamo deciso la pizza, perché andiamo al sushi? Perché mi fate
questo affronto?''): sono persone che non sono abituate ad avere conflitti
costruttivi.

Se si pensa di essere in \emph{norming} e continuano a emergere conflitti, il
punto è che \textbf{non si era davvero in norming}: lo storming non ``si
trascina'' nelle fasi successive, semplicemente non se n'è mai usciti.
Allo stesso modo, a chi chiede se ogni mini-fase di storming durante la vita
del gruppo sia frutto di uno storming fatto male o solo di carattere, la
risposta è già contenuta qui.

Un partecipante osserva che spesso, ancor prima di finire un progetto, si è
già messi su quello successivo: è il modo migliore per essere in costante
\emph{damage control}. Era l'esperienza del relatore a Blizzard: veniva
spostato fisicamente da un team all'altro, e finito il lancio di
\emph{Diablo~III} era già su \emph{World of Warcraft}, senza vivere il
``dopo''. Si è dovuto ritagliare da solo lo spazio per un debrief personale,
ed è lì che ha deciso che tipo di manager non voleva essere. Il consiglio:
\textbf{se non lo fa il team, fatelo voi}.

\subsection{Dove si passa la maggior parte del tempo}

\emph{In quale fase si passa più tempo nella propria carriera?} Gli studenti
rispondono \emph{storming} e \emph{norming}. Il relatore, in quindici anni,
calcola di aver passato in \emph{performing} circa \textbf{dieci mesi}. Non è
pessimismo: arrivare a performing deve essere sempre l'obiettivo, sia da
manager sia da membro del team; ma è difficile arrivarci, ed è normale che non
sempre ci si arrivi.

Lo storming può durare a lungo, ma con team molto esperti dura pochissimo,
perché c'è chi lo sa gestire. In realtà \textbf{la maggior parte del tempo si
passa in forming}, perché ogni equilibrio raggiunto vale per \emph{quelle}
persone, \emph{quel} progetto, \emph{quelle} tempistiche. Basta che cambi un
qualsiasi fattore (una deadline spostata, una persona nuova, qualcuno che si
dimette, un progetto cancellato) e si ricomincia, perché i bilanciamenti
precedenti non valgono più.

Esempio: un team di cinque persone affiatatissimo, in pieno performing. La
dirigenza, entusiasta, ``raddoppia le risorse'' e lo porta a dieci. O si
lascia tutto com'è (e i cinque nuovi non fanno niente), o bisogna ridiscutere
chi fa cosa e come i nuovi possono aiutare: che piaccia o no, si ricomincia
da forming.

Quindi la bravura di un manager, ma anche di un membro del team,
\textbf{non sta nel passare tanto tempo in performing, ma nell'accelerare il
passaggio da una fase all'altra}, e nel fare bene la fase di
\emph{adjourning}.

\subsection{Analizzare anche i successi}

C'è la tendenza a fare post-mortem solo dei progetti andati male: ``questa
cosa non ha funzionato, vediamo perché''. Va benissimo. Ma quando un progetto
va bene si dà per scontato che sia andato bene ``perché era buono''. Bisogna
invece chiedersi \emph{perché} è andato bene: era il prodotto? La
tempistica? Il modo in cui abbiamo fatto una cosa piuttosto che un'altra?
\textbf{Se non si capisce davvero il motivo del proprio successo, si
tenderanno a replicare le cose sbagliate.}

\paragraph{Blockbuster.} Quando arrivò la notizia che Netflix avrebbe
lanciato un servizio di streaming, Blockbuster (la catena di noleggio di
videocassette e DVD) aveva tutto il tempo per preparare un'offerta
concorrente. In una riunione del management raccontata in un libro (di cui il
relatore non ricorda il titolo), la direzione sembrava orientata a investire
nello streaming, finché un dirigente fece notare che il 12\% del fatturato
derivava dalle \emph{multe per riconsegna tardiva}, impossibili con lo
streaming, mentre secondo le stime lo streaming avrebbe prodotto il 10\% del
fatturato. ``Non possiamo rinunciare a questo 2\%'': niente streaming. Oggi
non c'è più un Blockbuster ``manco a cercarlo col lanternino''. Non avevano
capito qual era il vero motivo del loro successo, cioè il loro modello di
business\footnote{Le percentuali sono quelle riferite in lezione e non sono
state verificate.}.

\paragraph{The Founder.} Nel film sulla nascita di McDonald's c'è una scena in
cui il futuro socio del protagonista gli dice: ``tu pensi di essere nella
\emph{food industry}, e invece sei nella \emph{real estate industry}''. Non
dovrebbe investire sugli hamburger, ma sui terreni su cui costruisce i
ristoranti. Anche per un team vale lo stesso: bisogna capire che cosa
l'ha reso davvero di successo, oppure no.

\subsection{L'autoanalisi personale}

\paragraph{E quando si cambia azienda?} Uno studente chiede se il debrief sia
possibile anche quando si lascia un'azienda. Non si può tornare in ufficio a
discutere della propria esperienza, ma lo si può fare \textbf{per se stessi}.
Il relatore l'ha fatto, anche se all'epoca in modo poco consapevole, ed è
così che è passato dal \emph{community management} al \emph{people
management}. Analizzando cosa aveva fatto bene a Blizzard si era reso conto
che:
\begin{itemize}
  \item aveva costruito un ottimo rapporto di fiducia con i giocatori (con
        alcuni proprietari di fansite parla ancora oggi);
  \item pur senza la seniority, era diventato un punto di riferimento per il
        team, visto come imparziale ed equilibrato anche su temi che non lo
        riguardavano direttamente;
  \item sapeva \textbf{risolvere problemi}: durante il lancio di
        \emph{Diablo~III} era il più tranquillo di tutti, non perché
        sminuisse il problema, ma perché sapeva che si sarebbe risolto con i
        suoi tempi; il lavoro non era spegnere l'incendio ma ``far sfollare
        la gente'' e aspettare che tutto tornasse in ordine.
\end{itemize}
E aveva capito cosa sapeva fare meno bene, ad esempio \textbf{gestire il
proprio manager}. Avere un manager non significa che il manager gestisca te:
il manager gestisce il tuo lavoro, ma tu devi saper gestire i tuoi
\emph{stakeholder}. Spesso lo stakeholder è il proprio manager; da manager,
gli stakeholder saranno CEO, CFO e così via. Gestirli significa conoscerli,
sapere come portare loro le informazioni, come far presente che qualcosa non
va. Fa parte del lavoro, anche se a lui non piace, e ci ha lavorato.

Nel mese e mezzo di pausa prima di entrare in Riot si è concentrato su
questo: ``cosa porto come mio punto di forza?''. L'autoanalisi è culminata
nel chiedere espressamente alla persona che l'aveva assunto: ``prima mi avete
detto di no, ora mi avete preso: perché?''.

\paragraph{Come si capisce se un progetto è andato bene?} Dipende
dall'\textbf{obiettivo}. Se l'obiettivo era vendere milioni di copie e il
gioco le ha vendute, bisogna capire perché. Se, più ragionevolmente,
l'obiettivo era un gioco da tavolo che permettesse di mettere in pratica le
conoscenze del corso, va valutato per quello: ha permesso di farlo? Se sì,
perché? Se no, perché no? L'obiettivo va fissato \textbf{all'inizio del
progetto, non alla fine}, proprio per poterlo valutare dopo.

\subsection{Quando il team cambia di continuo}

\paragraph{Tutti questi cambiamenti non sono dannosi per il progetto?}
Assolutamente sì, ma quasi mai sono volontari. I progetti andati male sono
spesso progetti grandi.

Il relatore racconta di \emph{Cyberpunk 2077}. Conosceva persone in CD
Projekt Red (``il mondo è piccolissimo, le cose si sanno''). Alla Gamescom va
con Diego Ricchiuti, direttore dell'accademia, a vedere la presentazione del
primo gameplay, quella ``in cui regalavano il giubbottino da samurai''.
Finito il video, tra gli applausi, si guardano e capiscono entrambi che
\textbf{il gioco non è pronto}, a soli quattro mesi dall'uscita prevista:
Diego aveva notato il combattimento e la guida. Gli amici di CD Projekt Red
(uno aveva lavorato con lui a Blizzard per sei anni) gli confermano tra le
righe che ``ce la faremo, ma ci sono ancora cose da sistemare''. Il gioco fu
poi rimandato più volte\footnote{In lezione il relatore dice ``due, forse
tre volte''; l'uscita di \emph{Cyberpunk 2077} è stata effettivamente
rinviata tre volte prima del dicembre 2020.}.

Per deformazione professionale, chi ha sviluppato l'occhio critico non
riesce più a godersi un trailer: in un trailer di \emph{The Day Before} Diego
gli aveva fatto notare che l'effetto della fiammata di un fucile era stato
chiaramente aggiunto in postproduzione e non faceva parte del gameplay. ``È
il triste epilogo di questa storia'': non riescono più a godersi un trailer.

Il punto è che più il team è grande, più queste variabili vengono imposte, e
più è facile che il fallimento \textbf{non dipenda dal fatto che il prodotto
sia brutto} (\emph{Cyberpunk} è un bellissimo gioco sotto molti punti di
vista, e il relatore lo sta rigiocando ora che è stato corretto ed espanso).
Nei grandi gruppi qualcuno arriva, qualcuno si dimette, cambia il budget,
arriva il rinvio; nel caso di \emph{Anthem}, il team ha fatto così tante volte
forming--storming--forming--storming che ne è uscito ``un prodotto che
dovevamo espellere: l'abbiamo espulso, buttato fuori dal nostro organismo pur
di liberarcene''.

Tra i cambiamenti c'è anche la crescita verticale: se il capo viene promosso
perché lavora benissimo, bisogna scegliere un nuovo capo, ed è sempre un
caos (chi ci resta male, chi pensava di meritarlo, chi voleva il suo amico,
chi non sopporta il prescelto). Per questo il relatore è \textbf{contrario
alle promozioni durante i progetti}: si rischia di rovinare tutto il lavoro di
storming.

\subsection{Obiettivi chiari nei piccoli gruppi}

Nei gruppi piccoli, come quelli degli studenti, il problema deriva più
spesso dal \textbf{non avere chiaro che cosa si vuole ottenere}: un prodotto
puramente educativo e formativo, o un prodotto richiesto dall'accademia che,
se viene bene, si potrebbe anche vendere? Un gioco commerciale o un gioco
artistico? Va chiarito subito, perché ci saranno sempre momenti in cui
scegliere tra A e B, e a guidare la scelta sarà la domanda: \emph{quale delle
due opzioni ci porta più vicino all'obiettivo?} Se l'obiettivo non è chiaro,
non si sa cosa scegliere.

Questo porta molte aziende a fare \textbf{scelte di design che in realtà sono
scelte commerciali}. Le microtransazioni sono scelte commerciali: nessuno ha
mai pensato ``sai cosa piacerebbe tanto ai giocatori? Pagare le bustine''.
Prima le bustine si guadagnavano giocando; ora la tendenza è ``potresti
giocare, ma se paghi è meglio''. Si pensi alle polemiche su Ubisoft per
l'opzione di ``saltare il gioco'' pagando, o ai famosi \emph{boost} di livello
di \emph{World of Warcraft}: ``devo pagarti per saltare il gioco?''.

Una studentessa (Silvia) cita \emph{Final Fantasy XIV}, che a un certo punto fu
``resettato'' perché non funzionava: un esempio positivo. Il relatore aggiunge
quello negativo di \emph{Age of Conan}, che non è mai riuscito a fare quel
reset iniziale (``è una schifezza, ve lo ripresentiamo quando è pronto''). Ci
vuole consapevolezza, ma anche i mezzi: un conto è Square Enix, che può
permetterselo, un altro è un'azienda che non può chiudere tutto per un anno
continuando a pagare gli stipendi.

Riassumendo: nei piccoli gruppi i problemi sono più di
\textbf{comunicazione e chiarezza} (obiettivi non chiari o che si spostano);
nelle grandi aziende di più di \textbf{persone che vanno e vengono, obiettivi
e tempistiche che cambiano}. In accademia le scadenze sono fisse (il gioco da
tavolo si consegna a fine anno scolastico, a giugno, e giugno resta); in
azienda può succedere che una feature prevista per la \emph{vertical slice}
di luglio venga anticipata a maggio per portarla alla Gamescom, e un sistema
perfetto per consegnare a giugno può non esserlo più per consegnare a maggio.

\subsection{L'ambiente determina la qualità del prodotto}

\paragraph{È vero che chi passa per Electronic Arts finisce per fare brutti
giochi?} Uno studente cita BioWare, che faceva ottimi giochi finché era
autonoma. Il problema non è dei developer, che non sono ``improvvisamente
diventati incapaci'': è evidente che il problema è altrove, e che
\textbf{l'ambiente in cui si fa un gioco determina molto spesso la qualità del
prodotto}.

Un esempio locale positivo: Tommaso Bonanni, conosciuto dal relatore insieme a
Diego e già docente dell'accademia, ha fondato a Roma Caracal Games Studio,
applicando fin dall'inizio questi concetti di cultura aziendale. L'azienda è
rimasta piccola, ma ha molto chiaro cosa vuole fare e come. Quando si
predispone fin dall'inizio l'azienda e il team a lavorare in un certo modo è
fattibilissimo; è molto più difficile quando si è EA, si acquisisce
un'azienda e le si dice ``da oggi si fa così''.

Il relatore ha vissuto la transizione da Blizzard ad Activision Blizzard (uno
dei motivi per cui se n'è andato, volontariamente). Ha visto la cultura
interna cambiare per gradi:
\begin{enumerate}
  \item ``facciamo giochi fighissimi, perché i giochi devono essere
        fighissimi'';
  \item ``facciamo giochi fighi così vendono di più'';
  \item ``vendiamo di più facendo giochi fighi'' (sembra uguale, ma è molto
        diverso);
  \item ``facciamo più soldi''.
\end{enumerate}
Lo si capiva dalle scelte di design, osservando quale fosse il peso più
grande sulla bilancia. Lo stesso varrà per gli studenti, nei loro giochi e nei
colloqui di lavoro: capire subito che cosa vuole l'altra parte.

\paragraph{Leader e visione.} Nei gruppi dell'accademia, dal secondo anno, c'è
una persona che assume il ruolo di \emph{lead} o \emph{scrum master}. È
diverso dall'essere un \emph{game director} (``il gioco si fa come dico io,
aiutatemi a farlo con la mia visione''), come avviene nei grandi studi: da lì
nascono i nomi nei titoli, come \emph{Sid Meier's Civilization}, che indicano
la persona che ha fatto quel gioco, oppure un puro branding, come \emph{Tom
Clancy's Rainbow Six}.

\paragraph{Inclusion come jolly.} A margine, uno studente osserva che in molte
aziende ``la risposta giusta è sempre \emph{inclusion}'', usata come un
jolly. È comunissimo, conferma il relatore: \emph{diversity and inclusion}
viene spesso usata per dire ``così includiamo più clienti'', ma non è quello
il suo punto (``altro discorso'').

% ---------------------------------------------------------------------
\section{Confrontarsi con chi è più bravo}
\label{sec:z1-rivalita}

Arrivato a Riot dopo sei anni a Blizzard, il relatore pensava di avere un buon
bagaglio di esperienza. Si è trovato invece ``come un fante buttato in mezzo
ai Navy SEAL'': persone che avevano superato 11, 12, 13 colloqui, con un
livello molto alto.

Individua subito un collega francese che, nella sua testa, diventa ``la
versione migliore di me'': stesso ruolo (community manager), esperienza
simile (arrivava da Ubisoft, come quasi tutti i francesi di Riot), ma molto
più spigliato, sempre con la battuta pronta, ``l'anima della festa''. Il
relatore invece è introverso: un workshop come questo per lui è
\emph{emotivamente drenante}, non ama parlare in pubblico pur amando
condividere le proprie esperienze, e non riesce a sostenere conversazioni
come ``cosa hai fatto questo weekend?''. Lavorando raramente insieme, aveva
sviluppato una rivalità ``da allenatore di Pokémon'' (lui era Ash, il
relatore l'antagonista).

Dopo sei mesi, a un evento di \emph{team building} aziendale (a Tenerife), il
suo capo, molto attento a queste dinamiche, li mette apposta nello stesso
gruppo. A fine serata, complice l'open bar, il relatore gli confessa di averlo
mal sopportato dal primo giorno per invidia della sua estroversione. Il
collega ride e risponde che pochi giorni prima aveva detto al capo la stessa
cosa al contrario: invidiava a Stefano il suo essere sempre organizzato, sul
pezzo, presentarsi nella sala riunioni giusta (mentre lui finiva in quella
sbagliata), arrivare con gli appunti, con i dati a supporto e le analisi di
mercato, mentre lui spesso non aveva nemmeno letto la mail del meeting.
Un'invidia in due direzioni, che li ha portati a lavorare molto insieme e a
migliorarsi a vicenda.

Le lezioni:
\begin{itemize}
  \item non bisogna demoralizzarsi vedendo qualcuno che fa lo stesso lavoro
        ed è più avanti sotto alcuni aspetti: spesso c'è un bilanciamento. I
        ``punti caratteristica'' sono limitati: se li ha messi in carisma,
        non li avrà messi in saggezza;
  \item conviene \textbf{chiedere a queste persone come mai} sono così brave
        in qualcosa, come si fa con un progetto andato bene: magari non lo si
        farà mai bene quanto loro, ma si comincerà ad avere il
        \emph{mindset} giusto; e viceversa, loro saranno spinti a chiedere a
        noi;
  \item spesso \textbf{le proprie doti le scopriamo dagli altri}. Il relatore
        non si era mai accorto di essere ``quello che resta calmo durante le
        emergenze'' (il lancio di \emph{Diablo~III}, quello di
        \emph{Draenor}, i problemi dal vivo alla Games Week) finché non
        gliel'hanno fatto notare tutti: per lui era la normalità. A chi pensa
        ``non sono capace, forse non è il mio percorso'', può darsi che
        semplicemente nessuno gli abbia ancora detto ``questa cosa la sai
        fare davvero bene, concentrati su questo''.
\end{itemize}

% ---------------------------------------------------------------------
\section{Domande sul lavoro nell'industria}
\label{sec:z1-industria}

Diverse domande degli studenti hanno riguardato le condizioni di lavoro nei
grandi studi internazionali. Le risposte sono raccolte qui per argomento.

\subsection{Presenza, remoto e trasferimenti}

\paragraph{Grandi studi: quasi solo in presenza.} Sia a Blizzard sia a Riot si
lavorava \textbf{solo in presenza}, senza smart working: sei anni a Parigi e
cinque a Dublino. Si parlava solo inglese: anche nei meeting tra soli
italiani si tendeva a comunicare e a scrivere in inglese, per correttezza
verso chi si fosse aggiunto e per lasciare una traccia comprensibile a tutti.
Il francese serviva per fare la spesa e sopravvivere, non in ufficio.
L'ambiente era estremamente eterogeneo (francesi, polacchi, spagnoli e
molti altri): durante il colloquio dal vivo a Dublino, mentre si preoccupava
di essersi vestito in modo adeguato, vide passare dietro la parete a vetro un
collega svedese con i dreadlock, a piedi scalzi, con una pianta sotto un
braccio e un aquilone sotto l'altro, e capì di essere ``a posto così''.
Un'esposizione a molte culture e modi di pensare che il relatore considera
bellissima.

Anche oggi, tranne alcuni studi particolari, la maggior parte del lavoro si
svolge di persona: amici a Paradox lavorano a Stoccolma, i pochi rimasti a
Blizzard si sono dovuti trasferire a Irvine. Per posizioni \emph{junior} o
\emph{entry level} i grandi studi raramente concedono il lavoro da casa; per
profili senior può esserci qualche eccezione (ad esempio nove mesi da remoto
e tre in sede).

\paragraph{Relocation.} Spesso tocca spostarsi, ma il relatore non ha mai
pagato un trasloco di tasca propria: queste aziende offrono la
\emph{relocation} (aiuto a trovare casa, ad aprire il conto in banca e
supporto pratico), perché è un investimento sulla persona.

\paragraph{Colloqui mentre si lavora altrove.} Per i colloqui in presenza
prendeva ferie, senza dire dove andava. Il consiglio generale: per fare
colloqui \textbf{prendete ferie o permessi}, non datevi malati, ``che non è un
bel modo di comportarsi''.

\paragraph{Remoto ``doppio''.} Un collega lavora per l'azienda coreana che
gestisce \emph{PUBG}, con ufficio ad Amsterdam, ma lavora da Milano: un
``doppio remoto''. Si trova, ma in genere un ambiente anglofono con sedi in
Europa chiederà di trasferirsi in Europa.

\paragraph{Famiglia e trasferimenti.} Il relatore ha avuto la fortuna di
sposare una persona che fa lo stesso lavoro (conosciuta lavorando su
\emph{Warhammer Online}, poi con lui a Blizzard e a Riot, nella
localizzazione e nel doppiaggio, ad esempio di \emph{World of Warcraft} in
italiano): avendo competenze rare, spesso venivano assunti entrambi, al
massimo a pochi mesi di distanza. Per i developer la soluzione frequente è
che la famiglia segua la persona assunta, perché lo stipendio basta per
tutti. Ma non sempre c'è equilibrio, e a volte uno dei due deve sacrificarsi:
sua moglie ha lasciato Blizzard per seguirlo a Dublino (per poi essere
assunta da Riot); lui ha lasciato Riot per tornare in Italia quando lei ha
trovato lavoro qui. Ha dovuto rinunciare più volte a lavorare in Islanda per
CCP (\emph{EVE Online}), perché non era una scelta solo sua. Si fanno
compromessi, ma si riesce a far quadrare.

\paragraph{Ci si sposta per ogni progetto?} No: si sono spostati
esclusivamente per cambiare azienda. Cambiando progetto nella stessa azienda
è raro, quasi mai (``non l'ho mai visto succedere''), che venga chiesto di
trasferirsi. In Francia stavano benissimo (vivevano fuori Parigi; ``i
problemi con i francesi sono in realtà problemi con i parigini''), parlavano
francese ed erano vicini alla doppia cittadinanza: se ne sono andati perché
non vedevano un futuro a Blizzard, per il modo in cui l'azienda trattava i
colleghi.

\paragraph{Il mercato oggi.} Non esiste un \emph{silver bullet} per trovare
lavoro domani dove si vuole. Con tutti i \emph{layoff} recenti ci sono
migliaia di esperti del settore che cercano lavoro: serve pazienza, ma è
fattibile.

\paragraph{Studi indie e remoto.} Negli studi più piccoli il remoto è più
diffuso: a Stormind circa 60 persone su 95 lavorano da remoto, ma lo studio
si è strutturato fin dall'inizio per farlo. Serve consapevolezza. Bending
Spoons (che sviluppa app per mobile, non giochi, ed è l'azienda dell'app
Immuni), pur essendo grandissima e con sede a Milano, lavora da remoto e
lascia l'ufficio a chi vuole: lo smart working è posizionato quasi come un
\emph{benefit}, ma permette anche di avere un ufficio da 300 persone invece
che da 700. Le sedi fisiche costano decine di milioni all'anno (Blizzard ha
chiuso le sedi di Amsterdam e di Parigi). Tuttavia, per lavori molto
verticali e specifici gli studi preferiscono avere le persone dal vivo; i
ruoli più spesso remoti sono quelli che non fanno sviluppo (customer service,
PR\dots).

Anche gli artisti lavorano spesso in presenza: l'accademia può dare solo una
\emph{simulazione} di quanto stretto sarà il contatto tra programmazione,
design e art. ``Sarà ancora di più.''

\paragraph{La comunicazione da remoto.} È più difficile. Pur essendo
estremamente introverso, il relatore \textbf{obbliga tutti a usare la webcam}
nei meeting in cui conta: non in un aggiornamento con novanta persone
collegate, ma certamente in un colloquio a due sulla performance dell'ultimo
anno. La componente fisica è una parte enorme della
comunicazione\footnote{In lezione si parla di ``quasi l'ottanta per cento'';
il dato va preso come indicativo, poiché le percentuali di questo tipo
circolano in forme diverse e sono spesso citate fuori contesto.}: anche se si
è timidi, accendere la webcam è l'approssimazione più vicina all'essere dal
vivo. Negli ambienti dove si lavorava senza webcam la situazione era
``esponenzialmente peggiore''.

\subsection{Differenze culturali nel team}

\paragraph{Gestire sensibilità diverse.} A Riot i team erano divisi per
\emph{lingua}, non per nazione: il team di lingua spagnola includeva anche un
messicano, quello di lingua russa includeva russi e ucraini, proprio nel
periodo della guerra in Crimea. Tra i colleghi non c'era alcun problema (erano
amici), ma ci si è posti la domanda: \emph{come comunichiamo con una community
che sarà divisa su questo tema?} Nei giochi online la divisione era evidente:
in \emph{Rocket League} c'era chi, con un compagno di squadra con la bandiera
avversaria, si faceva otto autogol pur di farlo perdere. Allo stesso modo,
con molti americani mandati a Dublino a coprire posizioni scoperte, c'erano le
elezioni (Clinton--Trump, Biden--Trump) e varie guerre: serviva sempre una
certa sensibilità.

Imparare a navigare queste situazioni lo ha reso un professionista migliore:
gli ha insegnato a fermarsi e chiedersi ``io la percepisco così; gli altri
dieci, con background, storie e lingue diverse, come la penseranno?'', a
ragionare più \emph{long term} e più \emph{big picture}, con più empatia.

\paragraph{Stili di comunicazione.} Un partecipante che ha lavorato nel Regno
Unito chiede dello stile comunicativo italiano, a volte frainteso. Uno dei più
grandi complimenti ricevuti dal relatore è stato ``non sembri italiano''
(``gesticoli poco per essere italiano'', ``parli bene inglese per essere
italiano''). Ha imparato ad \textbf{adattare la comunicazione alla persona}:
con spagnoli e francesi la battuta sulla nazione funziona, perché c'è un
``odio--amore'' condiviso; con russi e giapponesi no. Doveva adattare anche
l'umorismo. Un collega metà russo e metà messicano diceva di aver preso ``il
lato peggiore di entrambe le nazionalità'', ed è stato un ottimo maestro per
imparare a leggere culture diverse.

Non bisogna generalizzare, ma gli stereotipi possono essere un punto di
partenza ragionevole, \emph{sapendo che troverete sempre le eccezioni}. E
vanno letti correttamente: gli americani tendono a essere più estroversi e ad
avere sempre qualcosa da dire non perché ``sono fatti così'', ma perché
vengono educati così (all'università tra i corsi c'è il \emph{debate}, saper
argomentare in pubblico); un collega giapponese può essere più introverso per
ragioni culturali. Nel dubbio si parte da lì, poi si conosce la persona.

\subsection{Negoziare, discriminazione e stipendi}

\paragraph{La famiglia come leva.} Avere un partner riduce le opzioni, ma le
aziende lo sanno: uno spostamento è per due. Si può usare come leva per
negoziare (un compenso diverso, un supporto per il trasloco): se un'azienda ha
deciso di volere quella persona, 2, 3, 4 o 10.000 euro in più all'anno, per
un'azienda che fattura centinaia di milioni o miliardi, non cambiano nulla.

\paragraph{Essere sposati è un motivo per non essere assunti?} No. In alcuni
paesi è proprio illegale discriminare su questo: in Francia era illegale
chiedere CV con data di nascita, foto, sesso o religione, e chiedere se si
vogliono figli. Purtroppo, dai racconti dei suoi studenti di master (in
settori come \emph{beauty and wellness} e \emph{food and wine}), in piccole
realtà italiane queste domande (``hai 25 anni, hai intenzione di
sposarti?'') vengono ancora fatte, e servono a filtrare i candidati. Nelle
grandi aziende ci sono sistemi per evitarlo; anzi, molte danno \emph{più}
benefit a chi ha famiglia (``porta pure la famiglia, è tutto pagato''),
perché le persone competenti hanno spesso una certa età e quindi,
ragionevolmente, una famiglia.

\paragraph{Gli stipendi.} Dipendono molto da dove si vive. Lo stipendio a Riot
non era paragonabile a quello di Blizzard perché in Francia era tassato al
52\% e in Irlanda al 17\%\footnote{Aliquote come riportate in lezione,
riferite alla situazione personale del relatore.}: nessuno stipendio
competitivo può compensare il 35\% di differenza di tasse. Negli Stati Uniti,
con pochissima previdenza pubblica, gli stipendi sono molto più alti perché
poi ci si paga tutto (``se hai un raffreddore paghi tu''). Bisogna valutare il
potere d'acquisto di quei soldi nel paese in cui si andrà a vivere.

\subsection{Freelance e \emph{work for hire}}

All'inizio della sessione pomeridiana uno studente che lavora nell'audio
chiede quali siano le condizioni per chi lavora come \emph{freelance}.
\begin{itemize}
  \item Dipende molto dal progetto: spesso il contratto è a progetto, per
        portare a termine una parte o un intero progetto.
  \item Non tutte le aziende usano freelance: Riot e Blizzard lo facevano
        pochissimo e preferivano assumere (magari trattando poi la persona
        ``come un freelance'', lasciandola a casa a fine progetto).
  \item Spesso i lavori offerti come freelance sono in realtà
        \textbf{sostitutivi di un full time} (``freelance, ma sei ore al
        giorno'').
  \item Si può fare, ma bisogna essere molto chiari nel definire l'impegno
        (quante ore, come vengono tracciate): non tutti specificano cosa
        significa davvero ``freelance''. Lavorando così per due o tre aziende
        si può arrivare a un tempo pieno.
  \item È più comune il \emph{work for hire}: aziende che forniscono personale
        ad altre aziende; si è assunti da X ma si lavora ``embeddati'' presso
        Y (una sorta di subappalto). Anche lo studio del relatore lo fa.
  \item Il puro freelance di quattro mesi esiste, ma di solito per quei
        quattro mesi diventa un lavoro a tempo pieno; è raro ``quattro mesi,
        due ore al giorno''. Nelle grandi aziende è improbabile lavorare nello
        stesso mese per Ubisoft, Riot e Blizzard tre ore ciascuna; è più
        realistico lavorare per più startup indie italiane, che non possono
        permettersi un'assunzione ma possono pagare due ore al giorno.
  \item Molto passa dai contatti: trovare le aziende giuste e rendere esplicite
        le condizioni.
\end{itemize}

% ---------------------------------------------------------------------
\section{Il circolo virtuoso: feedback, miglioramento, fiducia}
\label{sec:z1-triangolo}

\subsection{Il triangolo}

L'obiettivo a cui il workshop vuole arrivare è un \textbf{triangolo che si
autoalimenta}:
\begin{enumerate}
  \item si dà un \textbf{feedback} con l'obiettivo di migliorare le
        performance (inclusi elementi caratteriali);
  \item la persona \textbf{migliora} e vede di averne beneficiato;
  \item questo genera \textbf{fiducia}, che rende più facile dare e ricevere
        il feedback successivo.
\end{enumerate}
Togliendo uno qualsiasi dei tre elementi, il circolo non è altrettanto forte.

In qualsiasi fase si parta, il rispetto per i colleghi è inizialmente dettato
dall'\emph{educazione}, non dalla fiducia. La fiducia arriva quando i colleghi
cominciano a dire cose giuste e sensate e a \textbf{mantenere la parola}.
Dare e ricevere feedback è il modo in cui, in un certo senso,
``inganniamo il cervello'' per creare una \emph{case history} di fiducia: si
ha motivo di fidarsi di più, di ascoltare di più, di essere più ricettivi. E
si diventa sempre più bravi ad accettare il feedback e ad agire di
conseguenza.

\subsection{Accettare un feedback non vuol dire fare come dice l'altro}

Dieci anni fa il relatore non era ostile ai feedback, ma molto
\emph{difensivo}: ``sì, sì, va bene, ora ci guardo''; oppure, a chi gli diceva
``avresti dovuto farla così'', rispondeva spiegando perché l'aveva fatta in
quel modo. Oggi ha imparato a dire ``ok, grazie, ora ci ragiono'': non deve
difendere la propria posizione da chi viene solo a dire ``hai pensato a
questa cosa?''. Prima lo viveva come un attacco (``mi viene a dire che non so
fare il mio lavoro''), ora come un ``ottimo punto di vista, lo prendo in
considerazione''.

Con una postilla fondamentale: \textbf{accettare il feedback non significa
fare come dice l'altro}. ``Mi hai detto che preferisci questo mob verde. Ho
pensato anche a farlo verde? Sì. Penso comunque che giallo sia meglio? Sì, lo
faccio giallo.'' Il relatore accetta \emph{sempre} il feedback, valido o no:
garantisce di tenerlo a mente quando lavora. Che poi faccia come gli viene
detto è un altro discorso.

\subsection{Fiducia e scadenze}

La fiducia funziona allo stesso modo con gli impegni. Se si chiede a qualcuno
una cosa per venerdì e la consegna venerdì, la volta dopo gliela si chiederà
per venerdì. Se invece arriva il giovedì della settimana dopo, la volta dopo
gli si darà una \textbf{deadline anticipata}, come con chi arriva sempre in
ritardo agli appuntamenti (``se ti dico alle 7 arrivi alle 7 e mezza; la
prossima volta ti dico alle 6 e mezza'').

Al lavoro questo è dannosissimo, perché crea \textbf{deadline fittizie} che
non sono mai quelle vere, e nascono discussioni del tipo ``perché devo farlo
in due giorni se si fa in cinque?'' (``sì, se tu lo facessi in cinque''). La
conversazione giusta è un'altra: \emph{``che problemi hai avuto nel
consegnare venerdì?'' ``Questo, questo e questo.'' ``Cosa posso fare per
evitare che succeda?''}

Il relatore non vuole mai sentirsi dire ``non sono riuscito a farlo''.
Preferisce di gran lunga sentirsi dire, tre giorni prima, \textbf{``potrei non
riuscire a farlo, per questi motivi''}. Chi ha la visione d'insieme può dare
più budget, più persone, spostare la deadline, cambiare lo scopo del progetto:
può fare qualsiasi cosa per aiutare. Ma se lo si scopre il venerdì, cinque
minuti prima della deadline, non può fare più niente. Il messaggio è:
\textbf{``aiutami a poterti aiutare''}. Non penserà mai male di chi dice ``non
ce la faccio per questi motivi'', perché non sta dicendo ``non ho voglia'': sta
indicando i problemi, e lavorare su quei problemi è compito del manager. Non
può invece fidarsi di chi fino al giorno prima dice ``tranquillo, ce la
faccio'' e il giorno dopo ``non ce l'ho fatta''.

Questo è fondamentale per chi lavora in programmazione, design o art, perché
ci saranno sempre deadline: bisogna sempre \textbf{lavorare sulle aspettative},
dire al proprio manager (o al proprio report) cosa è lecito aspettarsi. Non
bisogna fare come i pinguini di \emph{Madagascar} (``sorridete e salutate, va
tutto bene''): se c'è un problema, va detto, altrimenti si peggiora solo la
propria situazione.

Da italiani si tende a non voler mai alzare la mano per chiedere aiuto
(``no, no, ce la facciamo, lo faccio io''). A volte non ce la si può fare, non
perché non si è capaci, ma perché vari fattori lo impediscono. Accettarlo è
importante: \textbf{è molto meglio avvisare e chiedere una mano che dire
``non ce l'ho fatta''}.

\subsection{Domande sul dare e ricevere feedback}

\paragraph{``Ne parliamo a fare, se poi fai come vuoi?''} Uno studente (Cristian)
riceve feedback che, dopo averli analizzati, non ritiene congrui; quando non
li segue gli viene detto ``ne parliamo a fare?''. C'è un problema di
\textbf{aspettative}: l'idea che accettare un feedback voglia dire fare come
dice l'altro. Accettarlo significa prendere in considerazione ciò che è stato
detto. Chi decide diversamente ha però un \textbf{onere}: argomentare perché
si fa comunque a modo proprio (``mi hai detto verde, la facciamo gialla
perché\dots''). Poi bisogna solo aspettare i risultati per vedere chi aveva
ragione.

La decisione spetta sempre \textbf{a chi ha la responsabilità di quella
decisione}. Se il game director deve scegliere tra A e B, anche se il
feedback dice B, la scelta è sua perché la responsabilità è sua. Non sarebbe
corretto essere ritenuti responsabili di una scelta delegata a chi dà
feedback. La risposta può essere semplicemente: ``l'ho tenuto in
considerazione, ma visto che ne sono responsabile io e le conseguenze sono
mie, ho scelto di fare X''.

\paragraph{E se l'altro fa ripicca?} Un partecipante chiede di chi, vedendo
il proprio feedback non accolto, si vendica ignorando a sua volta feedback
buoni. Per il relatore è un problema caratteriale: chi si comporta così ``è
già sulla strada per uscire''. Non c'è spazio per l'ego in team di questo
tipo. Il fenomeno nasce dall'idea che \emph{il feedback negativo faccia fare
brutta figura}, che il fatto che si potesse fare meglio tolga qualcosa alle
proprie capacità.

Non è così: \textbf{il feedback non è un giudizio}. Si può passare un anno a
Riot ricevendo solo feedback negativi e avere una valutazione annuale
stellare, perché ricevere molti feedback negativi significa essere una persona
approcciabile, che agisce sul feedback e invita gli altri a darlo. Il relatore
ha visto valutazioni molto peggiori per chi riceveva solo feedback positivi,
perché la gente aveva paura di dargliene di negativi: ma solo ``pacche sulle
spalle'' non fanno migliorare. Il problema va affrontato parlando con la
persona; a volte le conversazioni iniziano con \textbf{``ti do un feedback sul
feedback''}, cioè sul modo in cui l'altro ha dato il suo.

Esempi dalla sua esperienza:
\begin{itemize}
  \item lui chiede spesso di \textbf{non ricevere feedback senza esempi
        tangibili}, perché fa fatica ad astrarre: ``sei stato maleducato'' non
        significa nulla, ``ieri al meeting mi hai risposto male su questa
        cosa'' sì;
  \item ad altri serve l'opposto: ha ricevuto il feedback ``quando devi dirmi
        che ho sbagliato, puoi girarci un po' intorno e darmi del contesto?'',
        perché il suo stile diretto (``questa cosa non ha funzionato, dovremmo
        fare così'') feriva. Ha adattato il suo modo di dare feedback: con
        alcuni è diretto, con altri gira intorno. L'importante è saperlo.
\end{itemize}

\subsection{Non portare la soluzione}

Un consiglio centrale: \textbf{quando si dà un feedback, non si dà una
soluzione}. Non si dice ``io farei così'', ma ``secondo me questa cosa si
potrebbe fare diversamente, per questo e questo motivo: come potremmo farla
diversamente?''. È come chi smette di fumare perché vuole e chi smette perché
glielo dicono moglie e figli: se non si è convinti, spesso non si smette.

Quando il relatore diceva ``c'è questo problema, io farei così'', il meeting
successivo iniziava inevitabilmente con ``ci ho provato come hai detto tu, con
me non ha funzionato''. Sempre. Perché la soluzione era sua. La
responsabilità della soluzione deve essere invece \textbf{congiunta}: ogni
feedback deve creare un piccolo team che ha come obiettivo risolvere quel
problema. Questo toglie l'antagonismo (non ``ti dico che non sai fare il tuo
lavoro'', ma ``ti voglio aiutare a farlo meglio, lavoriamoci insieme'') e
sposta il focus: il problema non sei tu, è ciò che sta accadendo.

\paragraph{E se sono i piani alti a tenere all'oscuro dei problemi?} È una
questione di cultura aziendale. Il relatore fa sempre presente al proprio capo
il concetto di \textbf{service level agreement} (che verrà ripreso nella
seconda giornata): ``se mi dai questi input, ti garantisco questo output in
questi tempi''. Se il capo chiede delle cose senza dire budget, tempi e
perché, non può aspettarsi lo stesso risultato. È una comunicazione a due
vie: si dice al manager cosa serve per fare il lavoro, e lui deve fornirlo o
accettare che il risultato non sarà quello voluto.

\subsection{Amici, parenti e parole proibite}

\paragraph{``Non si lavora con amici e parenti''?} Per sua definizione
personale il relatore ha pochissimi amici (``stanno sulle dita di una mano'')
e una soglia molto alta per considerare amico un collega. Allo stesso tempo,
in quindici anni e con centinaia di colleghi, ha avuto solo \textbf{due}
esperienze in cui non è riuscito a lavorare con qualcuno: non serve essere
amici per far funzionare le cose. Con gli amici, anzi, gli viene più facile:
se dice a un amico che una cosa ``fa schifo'', l'amico non pensa ``come si
permette'', ma ``se perfino Stefano, che è mio amico, me lo dice, allora è
proprio così''. Diventa un rafforzativo, grazie alla fiducia costruita prima.
Il problema del lavorare con amici e parenti nasce dal \emph{non voler} dire
all'amico che sta sbagliando; se invece la confidenza permette di dirlo, si
può saltare la ``confezione'' del messaggio e passare direttamente al
contenuto.

\paragraph{Le frasi peggiori.} Alla domanda se esistano parole peggiori di
``ottimo lavoro'', il relatore indica la sua classifica: ``ottimo lavoro''
(detto in modo vuoto)\footnote{Il passaggio della trascrizione è
frammentario: si intende verosimilmente il complimento generico e di
circostanza, che non dice nulla di utile.}; poi la frase che lo ``fa uscire di
testa'', \textbf{``abbiamo sempre fatto così''} (``e infatti avete chiuso'',
come Blockbuster); infine \textbf{``siamo una famiglia''}, riferito al posto di
lavoro. Un'azienda non è mai una famiglia, perché non ci si può aspettare lo
stesso impegno da chi la possiede e da chi ci lavora: non c'è un
corrispettivo proporzionale tra impegno e ritorno. Per quanto inclusiva,
flessibile e razionale, non sarà mai una famiglia.

\subsection{Ego, muri e il peggior modo di scegliere un leader}

\paragraph{Chi mette davanti l'ego e alza muri.} Esistono, ma di solito durano
pochissimo in azienda. Non si tratta di avere un ego forte, ma di \textbf{non
saper ricevere feedback}. Spesso sono persone bravissime nel loro lavoro che
poi si chiedono perché tutti intorno diventino più bravi di loro: perché gli
altri accettano feedback e hanno più fonti da cui attingere, mentre loro si
sono rintanati nella loro torre.

Molti manager tendono ad \textbf{assumere persone un filo peggiori di loro},
per non fare brutta figura. Il relatore assume solo persone più brave di lui,
``perché non è una gara'': ha lavorato con molti community manager più bravi
di lui come community manager, ma lui era più bravo a fare il manager. Non
viene valutato su quanto sa fare il community manager, ma su quanto il suo
team sa farlo. Un partecipante chiede se la paura di avere qualcuno di
migliore esista anche in ambito accademico. Assolutamente sì: al relatore è
capitato di non essere preso come docente in un corso di formazione sui
videogiochi perché il responsabile aveva percepito che ne sapesse più di
lui\footnote{La trascrizione dice ``che non sapevo più di lui'', ma il senso
del passaggio richiede il contrario.}. ``Ma perché te ne importa? Tu sei quello che ha portato chi
sa fare il suo lavoro: ben venga.'' O si ha questa cultura, o si ha paura per
il proprio posto.

\paragraph{Promuovere il migliore.} Quando bisogna scegliere un leader per un
team (programmazione, design, art), molto spesso si sceglie \textbf{il migliore
tecnico}: il miglior programmatore, il miglior designer. È spesso la scelta
sbagliata: è come prendere un idraulico per fare l'elettricista. Fare il
\emph{people manager} non è come fare il designer o il programmatore; si può
essere un ottimo programmatore e un pessimo manager, e viceversa. Le aziende
l'hanno capito e lo fanno meno di prima. A Blizzard c'era un grave problema di
\emph{middle management}: molte persone in ruoli importanti erano state
promosse cercando le skill sbagliate, erano frustrate dal non saper fare il
proprio lavoro e schiacciavano chi stava sotto per paura di fare brutta
figura.

Un partecipante cita una teoria simile: una persona viene promossa finché non
raggiunge \emph{il proprio livello di incompetenza}, dove si blocca\footnote{È
il \emph{principio di Peter} (Laurence J.\ Peter, 1969), citato in lezione
come ``la teoria di Peter''.}. Un altro cita Michael Scott di \emph{The
Office}, su cui il relatore tornerà (sezione~\ref{sec:z1-quando}).

% ---------------------------------------------------------------------
\section{Perché la cultura conta: retention, costi e fiducia}
\label{sec:z1-cultura}

\subsection{L'ambiente di lavoro e la retention}

Perché tutto questo è così importante? Perché l'ambiente di lavoro conta
moltissimo per la \textbf{retention}, cioè per trattenere le persone e farle
lavorare in modo efficiente. Se contasse solo lo stipendio, molti non
lascerebbero mai Riot o Blizzard; e più si va avanti con l'età, più si impara
a dare valore anche al resto.

Inoltre, secondo gli studi citati nelle slide, chi vive il proprio lavoro in
modo negativo è \textbf{meno efficiente nel problem solving}: l'energia
mentale va verso ``come faccio a non farmi dare la colpa'' o ``come faccio a
fare bella figura'', non verso la soluzione del problema. In ambienti negativi,
dove ci si deve preoccupare della propria ``sopravvivenza'', si lavora
peggio.

\subsection{Il costo di chi se ne va}

Ogni volta che qualcuno si dimette o viene licenziato, il costo di trovare
una nuova persona, fare i colloqui, formarla e farle attraversare forming,
storming, norming\dots\ è circa \textbf{un terzo del suo stipendio}. Per questo
le aziende tendono a fare più \emph{retention} che \emph{acquisition}, almeno
con le figure considerate utili alla produzione.

Allo stesso tempo si vedono sempre più spesso \emph{layoff} e \emph{hiring
wave}: si licenziano 3.000 persone e sei mesi dopo se ne riassumono 2.500 o
3.500, ma tutte più junior. Il relatore ne ha scritto in un articolo su
Eurogamer.it, in cui spiega anche alcuni ``trucchi del mestiere'': perché i
licenziamenti avvengono sempre in certi periodi e sempre ``a gruppi'' (prima
Bethesda, poi Blizzard, poi un'altra\dots).

\paragraph{L'offerta di Riot.} A Riot esisteva una pratica (di cui il relatore
non ricorda il nome) per cui, firmato il contratto, si aveva \textbf{un anno
di tempo per dire ``non mi piace, non me la sento, me ne vado''} e ricevere
subito come buonuscita un anno di stipendio. Avevano stimato che fosse molto
meglio pagare subito un anno a qualcuno e lasciarlo andare che tenerlo due,
tre, quattro anni come \emph{underperformer}: la responsabilità della scelta
di assumere era dell'azienda\footnote{La pratica, nota come ``The Offer'', è
descritta in fonti pubbliche con condizioni diverse (un periodo più breve e
una somma pari a una frazione dello stipendio annuo). Si riportano le
condizioni come raccontate in lezione; andrebbero verificate.}. L'hanno usata
pochissimi, ma il solo fatto che esistesse faceva lavorare meglio anche lui da
manager, perché sapeva che chi non si trovava bene poteva andarsene con un
anno di tempo per trovare un altro lavoro.

\paragraph{Il ricircolo.} C'è un forte ricircolo nell'industria: molti ex
colleghi passano da Ubisoft a Paradox, da Paradox a Funcom, da Funcom a
Blizzard. Il relatore ha scelto di uscirne lavorando solo per aziende
\textbf{indipendenti}: Stormind non è un indie di cinque persone, ma è
indipendente, e non deve accontentare azionisti. Quando vede un'azienda con
utili miliardari licenziare migliaia di persone si chiede ``dov'è la linea''
(\emph{when enough is enough}): con azionisti che vogliono un ritorno, se
quello è l'obiettivo stanno facendo la cosa giusta, ma lui non condivide
l'obiettivo\footnote{In lezione si parla di ``6,5 miliardi di utile'' e di
``3.000 persone licenziate''; le cifre sono riportate come dette e non sono
state verificate.}.

\paragraph{L'open space vuoto.} A Riot, durante un'ondata di licenziamenti,
del suo team di trenta persone è rimasto solo lui: ``teniamo te come
coordinatore, poi ricostruirai il team''. Ha rifiutato di ricostruire un team
cancellato ``con un colpo di spugna'' e ha deciso che non avrebbe più lavorato
per una grande azienda. Nessuno aveva pensato di spostarlo, e ha lavorato per
tre mesi in un open space di cento metri quadri con la sua scrivania e altre
ventinove vuote. Su una lavagna aveva scritto i nomi dei suoi ragazzi e li
cancellava man mano che, anche grazie ai suoi \emph{referral}, trovavano
lavoro. Cancellato l'ultimo nome, si è dimesso: per lui il lavoro era finito.

\paragraph{Team che si dissolvono.} Uno studente cita un articolo su Larian
Studios, il cui management è contrario a ricostruire ogni volta il team a fine
progetto. È lo stesso problema: si pensa ai ``giochi Blizzard'' o ai ``giochi
BioWare'', ma le persone che li hanno fatti non ci sono più; è rimasto solo il
brand, e non ha senso avere le stesse aspettative da un gioco fatto da persone
diverse. E anche chi sopravvive ai layoff vive pensando che il prossimo
potrebbe essere lui: alcune persone, licenziate nell'ultima ondata a
Blizzard, gli hanno detto ``erano due anni che lo sapevo''. Vivevano a Cork,
``un paesotto irlandese dove o vai per lavorare a Blizzard o non ci vai'',
avevano comprato casa: non ci si può aspettare chissà che morale.

\subsection{La cultura è forte quando si lavora per gli altri}

\begin{quote}
  \emph{La cultura è forte quando le persone lavorano per gli altri e a favore
  degli altri, e non per se stesse.}
\end{quote}

Se si dà feedback per far migliorare l'altro e non per imporre il proprio
modo, se si aiuta per il bene del progetto e non per fare bella figura, la
cultura si rafforza. È dai \textbf{piccoli gesti} che emerge la cultura
aziendale: ad esempio, dal non dover chiedere ``perché mi stai aiutando?''.

\paragraph{L'incognita.} \textbf{L'incognita è spesso dannosa quanto le
cattive intenzioni.} In un esperimento (``mi pare a Stoccolma, cinque o sei
anni fa'') si distribuivano per strada banconote da 500 euro, e la maggior
parte delle persone le rifiutava chiedendosi ``dov'è la fregatura?''
\footnote{L'esperimento è riportato in modo vago in lezione e non è stato
possibile identificarlo con certezza.}. Lo stesso accadeva al relatore,
arrivato a PG e a Stormind da Riot: a ``vedo che sei in difficoltà, posso
aiutarti?'' la risposta era sempre ``perché?''. Alla fine rispondeva: ``perché
quando sarò in difficoltà io vorrei che tu mi aiutassi''.

Quando si comincia a lavorare per gli altri e si smette di pensare ``così avrò
la promozione, così avrò l'aumento'', promozioni, aumenti e rispetto arrivano
da soli. A Blizzard aveva fallito tre volte i colloqui per essere promosso a
manager (inteso come coordinatore del team); quando ha deciso di smettere di
pensarci e di fare solo il bene del proprio lavoro, la volta dopo è stato
promosso. Funziona così: chi aiuta i colleghi e fa andare bene il progetto
viene notato (``questa persona fa migliorare cinque persone, non solo se
stessa: vale di più'').

\subsection{Il lavoro è valutato come opera d'insieme}

Il relatore riconosce che buona parte del suo successo è venuto dal fatto che
il suo team lavorava bene, non dalla sua bravura personale. Anche i progetti
degli studenti, presentati a Diego, a Enrico, ad Andrea, verranno valutati
\textbf{come opera d'insieme}. Nessuno dice ``il \emph{combat system} di questo
gioco l'ha fatto X, che bravo'': si dice ``BioWare ha fatto un bel gioco'',
``Blizzard ha fatto un bel gioco''.

Lo ha replicato anche lato HR. In molte aziende ci sono \textbf{bonus
personali} legati a obiettivi (se il gioco vende X, se la feature è pronta in
tempo). Il relatore li ha dimezzati: metà restano personali, metà diventano
\textbf{bonus di team}, che tutti ricevono se il team raggiunge gli obiettivi.
Un incentivo concreto a collaborare: non solo si dice che il lavoro di gruppo è
importante, lo si dimostra dandogli più valore economico. Aveva proposto il
100\% di bonus di team, ma gli è stato concesso il 50 e 50; ha funzionato
comunque, anche se sarebbe stato ancora più potente. Ovviamente, alla
domanda ``perché mi aiuti?'' si può anche rispondere ``perché voglio il
bonus'': un aiuto dato ``per il motivo sbagliato'' era un compromesso che era
disposto ad accettare.

\subsection{Il coraggio di chiedere aiuto}

Un partecipante osserva che aiutarsi dovrebbe essere la norma. Il modello
diffuso, invece, è ``devi essere forte e indipendente, non chiedere mai una
mano'', e vale in tutte le direzioni (dal maschilismo alla ``donna forte e
indipendente che non ha bisogno di nessuno''): \textbf{non ammettere di avere
bisogno di aiuto è sempre uno dei motivi per cui le aziende vanno male}, che si
tratti di uomini, donne o persone non binarie. Il problema è non avere la
sicurezza di poter alzare la mano.

Il relatore ha cercato di togliere ogni conseguenza dall'aver alzato la mano:
se si \emph{depenalizza} qualcosa, le persone non sono più disincentivate a
farla. Ha detto espressamente che non avrebbe mai valutato nessuno per il
numero di volte in cui chiedeva aiuto o per il numero di feedback negativi
ricevuti, e questo ha funzionato.

\subsection{Perché si dà un feedback}

\emph{Perché avete dato il vostro ultimo feedback? Che cosa volevate ottenere?}
Le risposte della classe:
\begin{itemize}
  \item per \textbf{riconoscere il valore} di un collega, anche verso il
        manager, perché qualcuno che non aveva visibilità vedesse quanto
        quella persona fosse efficace;
  \item per \textbf{spingere una persona a superarsi};
  \item per aiutare un collega a rendere il prodotto all'altezza di ciò che
        voleva rappresentare;
  \item per \textbf{genuinità}: c'era qualcosa che poteva essere fatto meglio.
\end{itemize}
Il relatore aggiunge che la forma cambia col tempo: all'inizio era ``devo
darti un feedback, hai dieci minuti?''; ora è un messaggio informale su
Discord, perché ha costruito un rapporto di fiducia e sa che non verrà letto
come ``l'HR manager mi scrive, sono nei guai?'', ma come ``Stefano è così, mi
dà feedback, resterà tra noi e lo risolveremo insieme''.

Un altro motivo è \textbf{risolvere incompatibilità caratteriali}: le
conversazioni più difficili, del tipo ``io e te non stiamo funzionando, cosa
possiamo fare diversamente?''. Sono importanti perché, se due responsabili non
vanno d'accordo, il team se ne accorge, comincia a raccontarsi storie (``si
odiano, ci saranno problemi'') e, per difendere il proprio manager, comincia a
giudicare male il lavoro dell'altro gruppo. È la dinamica da fermare:
\textbf{valutare gli altri solo per difendere la propria fazione}.

Dare e ricevere feedback sono \textbf{due muscoli da allenare} (il relatore è
sempre stato più bravo a darlo che a riceverlo), e si impara col tempo. È come
fare colloqui: fare colloqui è una skill a parte, che non indica quanto si è
bravi nel proprio lavoro. Essere scartati non significa non saper fare quel
lavoro: spesso significa aver fatto male il colloquio o che è stato preferito
qualcun altro.

\paragraph{Feedback pubblici e privati.} Uno studente racconta di aver dato un
feedback negativo a un superiore, senza un motivo particolare, per dimostrare
a sé e ai colleghi che non si viene licenziati per questo. Il relatore coglie
l'occasione per precisare che \textbf{i feedback negativi si danno
individualmente}: lui dà moltissimi feedback di cui il resto dell'azienda non
sa nulla. Reagirebbe molto male se, in un meeting, qualcuno alzasse la mano per
dire ``comunque questa cosa non l'hai fatta bene'': quello non è feedback, è
discussione, accettabile solo se è un brainstorming e se l'opinione è stata
chiesta. Alcune aziende hanno piattaforme pubbliche di feedback, ma spesso
servono solo ad applaudire; lo stesso accade nelle aziende che
\emph{obbligano} a lasciare un certo numero di feedback a settimana, che
diventano ``lui è bravo, lei è competente, anche questa settimana l'abbiamo
fatto''. Bisognerebbe chiedere a chi le adotta cosa pensa di ottenere (forse
trasparenza), ma difficilmente una persona ``attaccata'' pubblicamente reagisce
bene.

% ---------------------------------------------------------------------
\section{Among Us: la caccia al traditore}
\label{sec:z1-among-us}

Prima della pausa pranzo il relatore mostra tre film di fantascienza di
qualche decennio fa e chiede che cosa li accomuni\footnote{I titoli non sono
leggibili nella trascrizione: viene citato \emph{Event Horizon} da uno
studente, e il film ``in mezzo'' è indicato come ``l'originale'',
verosimilmente \emph{La cosa} (\emph{The Thing}).}. In ognuno c'è un gruppo di
persone con un obiettivo e la presenza, reale o percepita, di qualcuno che
rema contro: \textbf{l'energia del gruppo viene spesa per trovare il
responsabile}.

\paragraph{L'esperimento degli alieni.} Nei master in presenza il relatore
propone un esercizio. Distribuisce segretamente cartellini plastificati con
scritto ``essere umano'' o ``alieno''; le carte alieno contengono istruzioni di
comportamento con cui gli alieni, che cercano di fingersi umani, si
tradirebbero (dare sempre ragione a chi ha parlato prima, finire ogni frase
con una domanda, parlare in un certo modo). Lo scopo: in dieci minuti formare
gruppi, e \textbf{vince il gruppo più numeroso senza alieni}; basta un alieno
per perdere. Si formano sempre gruppi piccoli (due, cinque, sei persone),
perché la tendenza è ``mi fido al massimo di voi due''.

Ma vincono i gruppi più grandi, per un motivo semplice: il relatore
\textbf{non distribuisce nessuna carta alieno}. Mostra a tutti entrambe le
carte, ma consegna solo carte ``umano''. Vuole simulare ciò che accade in
ufficio: la tendenza a percepire sempre qualcuno che rema contro (i manager,
l'altro team, quel collega), quando in realtà non c'è nessuno, ``è tutto nella
nostra testa''. Quando chiede le motivazioni, emergono sospetti
\emph{fabbricati}: ``faceva sempre domande'', ``parlava solo all'infinito''
(``no, parla così e basta''). Ci si fabbrica da soli i motivi per non fidarsi
degli altri, ed è comunissimo.

\paragraph{Among Us.} L'autosabotaggio è, per tornare ad \emph{Among Us},
\textbf{concentrarsi sul cercare cadaveri e traditori invece di fare le
task}: non si vince perché non si sono completate le task, mentre si cercava il
colpevole. Succede tantissimo negli uffici, ed è il motivo per cui l'80\% dei
problemi che il relatore risolveva erano di comunicazione. Entrando nei team
sentiva immancabilmente: ``facciamo il possibile, ma questa lotta continua con
l'altro gruppo è estenuante''. ``Perché c'è questo problema?'' ``Non lo so,
c'era già quando sono arrivato.'' Un paradosso: c'è un problema, nessuno sa
dire quale sia, ma c'è sempre stato e quindi lo si accetta.

\paragraph{L'esperimento delle scimmie.} È l'esempio, ``finto o reale'', delle
scimmie in gabbia\footnote{È il noto racconto delle ``cinque scimmie'', spesso
presentato come esperimento scientifico ma probabilmente apocrifo (lo stesso
relatore lo definisce ``finto barra reale''); nella trascrizione il numero
delle scimmie varia.}. In una gabbia con alcune scimmie pende una banana
elettrificata: chi la tocca fa prendere la scossa a tutte. Le scimmie imparano
e, quando una si avvicina alla banana, le altre la picchiano. Poi si
sostituisce una scimmia alla volta con una nuova: la nuova va verso la banana,
le altre la picchiano, e smette. Si continua finché nella gabbia ci sono solo
scimmie nuove, poi si toglie la corrente. Nessuna scimmia tocca la banana:
nessuna ha mai preso la scossa, ma tutte hanno imparato che chi tocca la
banana prende le botte.

Nei team è lo stesso: il motivo per cui si ``odia'' l'altro team si
trasforma, ma un vero motivo non c'è. Ne derivano la sfiducia, il non alzare
mai la mano per chiedere aiuto, il non dire mai ``bel lavoro'' agli altri
perché non sia mai che qualcuno ti veda fare un complimento ``a quelli''.
Risolvere queste dinamiche è lungo e tedioso, e sempre pieno di ``non mi
ricordo'', ``io non c'ero'' e ``abbiamo sempre fatto così'', la cosa più
frustrante per chi cerca di risolverle.

% ---------------------------------------------------------------------
\section{Il principio del default to trust}
\label{sec:z1-default-trust}

\subsection{La definizione}

Il relatore indica questa come la slide più importante della presentazione.
\textbf{Default to trust} è uno dei valori aziendali di Riot Games, che lui ha
tradotto, adattato e ``esportato'':

\begin{quote}
  \emph{Approccio ogni situazione di conflitto partendo dal presupposto che una
  persona abbia agito in un certo modo perché, con le informazioni a sua
  disposizione in quel momento, ha pensato che fosse la soluzione migliore, e
  non perché sia un'incapace o volesse danneggiare il progetto
  intenzionalmente.}
\end{quote}

In altre parole: \textbf{di default, mi fido}.

\subsection{L'esempio del ritardo al daily stand-up}

A Riot il relatore aveva un \emph{daily stand-up} dal vivo alle 15 con il suo
team di venti persone. Una ragazza arrivava regolarmente alle 15:15, per una,
due, tre settimane. Alla quarta ci ha parlato. \emph{Come iniziereste la
conversazione?} Le proposte della classe:
\begin{itemize}
  \item sottolineare l'importanza di esserci dall'inizio, perché potrebbe
        perdersi qualcosa che le interessa, e per rispetto dei colleghi
        presenti;
  \item ``devi arrivare in orario'', spiegando perché è importante;
  \item \textbf{chiedere perché}.
\end{itemize}

Il relatore osserva, anche dalle sue consulenze fuori dai videogiochi, che
\textbf{più si sale con la seniority}, più la risposta è ``devi smettere di
arrivare in ritardo'', spesso con un \emph{or else} finale (``altrimenti
succede quest'altra cosa''). Più si scende con l'età, e quindi con meno
abitudini consolidate, più la prima domanda diventa ``dimmi perché arrivi in
ritardo''. Già spiegare perché è importante è ``miglia avanti'' rispetto al
semplice ``smetti''.

Partire dal perché gli ha permesso di scoprire la causa: il meeting precedente
della ragazza, in un'altra zona del campus a un quarto d'ora di distanza,
finiva sempre tardi, e lei doveva attraversare il campus di corsa. Non aveva
la confidenza per dire all'altro manager, che insisteva perché si restasse
fino alla fine, che doveva andare; sacrificava quindi il meeting del relatore,
sapendo che era il più comprensivo dei due.

Il relatore si è offerto di \textbf{mediare}: ha detto all'altro manager che
o anticipava l'inizio del suo meeting, così da finire comunque in tempo,
oppure lasciava uscire la ragazza, perché quel meeting gli serviva. C'è anche
una questione di \textbf{economia di scala}, spesso sottovalutata: se venti
persone aspettano qualcuno per dieci minuti, si sono sprecati dieci minuti per
venti. Per lo stesso motivo il workshop è iniziato alle 10:05: iniziare alle
10:10 avrebbe significato far aspettare cinquanta persone cinque minuti in
più. A Riot c'era una scrivania vuota con un pupazzo gonfiabile, \textbf{Mr.\
Late}: sommando tutti i minuti di ritardo a meeting, call e conversazioni, in
un anno veniva fuori praticamente \emph{una persona a tempo pieno}. Serviva a
ricordare che anche cinque minuti contano, e che bisogna rispettare gli altri
arrivando in orario o almeno avvisando.

\paragraph{Gli effetti.} La conversazione era avvenuta nel primissimo periodo a
Riot, quando il team non aveva ancora motivo di fidarsi di lui. Risolvendo il
problema ha eliminato anche tutto lo strato di micro-dinamiche spiacevoli che
un manager impara a vedere (due colleghi che si guardano con la risatina
quando lei entra in ritardo). La ragazza ha cominciato a entrare tranquilla,
con una certa leggerezza: non le faceva piacere arrivare in ritardo (``ci sono
anche persone che lo fanno apposta, e quello è un problema da affrontare
diversamente''). Soprattutto, con i colleghi non parlava di lui come del capo
che minacciava, ma come di chi le aveva risolto un problema: era diventato
un \textbf{alleato}. La settimana dopo altre tre persone sono andate da lui
con un problema da risolvere. Ripetuto cinque, sei, sette, dodici volte, in un
paio di mesi poteva dire quasi qualsiasi cosa nei meeting sapendo che
richieste e critiche sarebbero state percepite come ``me lo chiede perché vuole
darmi una mano'', non per fare bella figura o per presentare i successi del
team ai suoi superiori.

\subsection{Condividere i successi, schermare i fallimenti}

Un altro elemento fondamentale (ripreso nella seconda giornata): il manager
\textbf{condivide i successi e protegge il team dai fallimenti}. ``Se va male
qualcosa è colpa mia; se va bene qualcosa, voglio che ne parliate voi.'' Il
lavoro del manager non dovrebbe essere valutato su ciò che presenta del lavoro
del team: tutti sanno comunque chi è il manager. Tanto vale lasciar presentare
il team, che potrà dire con orgoglio di aver presentato a Mike
Morhaime\footnote{Cofondatore ed ex presidente di Blizzard; nella
trascrizione ``Mike Morayme''.} la propria campagna europea per
\emph{StarCraft~II}. Ci sono manager che dicono ``come mi prendo le colpe, mi
prendo i meriti'', ma i meriti li hanno già: è inutile ostentarli. Con i
fallimenti invece il manager è la \textbf{rete di sicurezza} del team. Il
relatore ne ha fatto una questione di principio: \textbf{non scaricare mai la
colpa su qualcun altro}.

\paragraph{Colpevole o causa.} Un esempio che il relatore attribuisce a Pirate
Software, creatore di contenuti ed ex dipendente Blizzard\footnote{Il relatore
aggiunge che si tratterebbe del figlio dello sviluppatore a cui si sarebbe
ispirato il personaggio del giocatore ``senza vita'' nell'episodio di
\emph{South Park} dedicato a \emph{World of Warcraft}. L'informazione è
riportata come detta in lezione e non è verificata.}. Contattando il servizio
clienti Blizzard per motivi personali, aveva scoperto che i colleghi gli
avevano scritto nei ticket informazioni confidenziali che non avrebbero dovuto
condividere in quel modo, nemmeno con un collega. Non ha segnalato le persone
da punire: ha detto ``c'è un problema, evidentemente le policy non sono
chiare''. Quando gli hanno chiesto i nomi, ha rifiutato: ``se vi dico chi è,
punite quelli, e gli altri cinquecento continueranno a farlo; non avremo
cambiato la causa alla fonte''.

\begin{quote}
  \emph{È inutile trovare il colpevole. Voglio trovare la causa. Perché il
  colpevole cambia.}
\end{quote}

Il colpevole a volte è vittima della situazione: per questo la definizione
dice \emph{``con le informazioni a sua disposizione in quel momento''}. Se il
relatore avesse rimproverato la ragazza in ritardo, non solo non avrebbe
risolto il problema, ma avrebbe creato un'\textbf{ingiustizia}: lei, che
attraversava il campus di corsa per una colpa non sua e di cui non poteva
parlare, sarebbe stata anche rimproverata.

Nessuno fa un brutto lavoro apposta: nessuno è contento quando sbaglia o
perde. Perché al lavoro diamo per scontato che gli altri sbaglino
intenzionalmente? Se qualcuno ha fatto una cosa male, è perché in quel momento
credeva fosse la scelta migliore. La domanda giusta non è ``perché hai fatto
questa cosa che è stata un fallimento?'', ma:

\begin{quote}
  \emph{Che cosa ti ha fatto pensare che questa fosse la scelta giusta?}
\end{quote}

Magari erano limiti di budget, di tempo, di scopo: si agisce su quelli, non
sulla persona. Se si cambia la persona ma le si danno gli stessi fattori, è
molto probabile che arrivi alle stesse conclusioni, e ci si troverà a dover
``disciplinare'' anche lei.

\paragraph{La reputazione costruita sui rumor.} Durante il giro di
presentazione in un ufficio (non in nessuna delle aziende citate), al relatore
è stato detto: ``questo è il dipartimento di design; se ti serve qualcosa
chiedi a lei, non a lui, perché lui consegna sempre in ritardo''. Oltre a
trovarlo poco professionale, una settimana dopo ha chiesto: ``cosa ti ha fatto
questa persona?'' ``A me niente, però ho sentito che\dots''. Si torna alle
scimmie: magari quella persona non consegna in ritardo da mesi, ma ormai lo
dicono tutti. È un circolo vizioso da fermare subito.

\paragraph{Niente escalation inutili.} Lo stesso vale per il feedback: quando
qualcuno va dal relatore dicendo ``ho un problema con questa persona'', lui
chiede: \textbf{``perché ne stai parlando a me e non a lei?''}. ``Perché sei il
mio capo'' non basta. Passare da lui significa dare per scontato che l'altra
persona non abbia interesse a risolvere, mentre magari basta parlarle. È la
strada più facile (``lo dico a Stefano, ci pensa Stefano''), e proprio per
questo bisogna mettere dei paletti: il manager interviene quando serve, ma le
cose che si possono gestire da soli vanno gestite da soli. (Quando invece il
confronto diretto fallisce, l'escalation è legittima: si veda la
sezione~\ref{sec:z1-qa-finale}.)

% ---------------------------------------------------------------------
\section{Traiettorie di carriera e valutazione delle performance}
\label{sec:z1-traiettorie}

\subsection{Quale traiettoria vorreste?}

Il relatore mostra un grafico con quattro curve colorate, contrassegnate
anche da lettere (A, B, C, D) per eventuali persone daltoniche (``a me è
capitato che qualcuno scoprisse di esserlo durante una lezione''). Le curve
rappresentano la performance lavorativa nel tempo. Due domande: \emph{quale
vorreste essere? E quale pensate realisticamente di essere?} Un partecipante
aspira a B e pensa di essere C; quasi tutti scelgono B.

Perché non A, che sale più rapidamente? Un partecipante risponde che una
crescita così repentina può essere più dannosa che utile, come una ``rush'' di
studio molto faticosa che poi non si riesce a sostenere. Il relatore propone
una lettura suggerita da un amico psicologo, che spiega la forma del grafico:
di B, C e D \textbf{vediamo dove finiscono}; di A no. Potrebbe continuare a
salire oltre la scala del grafico, o crollare sotto D. Non saperlo ci spinge a
preferire B. (A chi obietta che anche le altre proseguono oltre il bordo, il
relatore precisa: il bordo destro è ``oggi'', e di A non sappiamo come sia
arrivata fin lì.)

\subsection{La carriera non è una linea retta}

Tutti immaginano la carriera come una linea: si parte da X e, attraverso
stage e posizioni, si arriva a Y. In realtà carriera e performance avranno un
andamento irregolare, con alti e bassi. Come non si accetta che un team possa
temporaneamente \emph{underperformare}, tanto meno lo si accetta da se stessi.
Bisogna accettare che ci saranno momenti bassi.

Per questo il relatore \textbf{detesta le valutazioni annuali delle
performance}: prendono un punto casuale della traiettoria, che se si è
fortunati è il picco più alto e se si è sfortunati il più basso, e lo usano
per valutare un intero anno. Insegna ai manager a fare più punti di controllo
e a guardare altre cose:
\begin{enumerate}
  \item \textbf{quanto tempo passa tra un picco e l'altro}, cioè quanto
        \emph{cooldown} serve a una persona. È il motivo per cui, dopo la
        Games Week, tutti staccavano una o due settimane. Il relatore stesso,
        che si impegna moltissimo e tende al burnout, non potrebbe rifare lo
        stesso workshop il lunedì e il martedì successivi con la stessa
        qualità;
  \item \textbf{la differenza tra la miglior performance e la peggiore}: che
        cosa ha portato a lavorare così bene, e così male? Con chi si
        lavorava, su cosa ci si concentrava, quanto tempo e budget c'erano?
        L'obiettivo è capire come \emph{replicare} le condizioni che fanno
        lavorare bene, invece del solito ``hai lavorato male, è colpa tua''.
        Il relatore ha avuto molti manager per cui quando lavorava male era
        colpa sua e quando lavorava bene ``che bravo'', ma sapeva che raramente
        era l'unico fattore della propria performance. È la fallacia del ``se
        sai fare 100, perché non fai sempre 100?'', come nella scena di
        \emph{Schindler's List} sull'operaio che produce una vite in pochi
        secondi: non sempre un ritmo è sostenibile;
  \item \textbf{il fallimento di oggi confrontato con quello di sei mesi fa}.
        Il relatore non è contento della presentazione di oggi (l'avrebbe
        migliorata per altri due anni), ma rispetto a quelle di cinque anni fa
        vede un grande miglioramento di chiarezza: il suo peggior fallimento di
        oggi equivale al suo primo successo di inizio carriera.
\end{enumerate}
Quello che conta è la \textbf{traiettoria} dall'inizio dell'anno a oggi, non
i singoli punti. E se si riesce a capire insieme quando e perché si ha
successo e quando e perché si fallisce, si può lavorare per replicare solo le
situazioni che fanno lavorare bene.

\paragraph{Lo stesso vale per i giochi.} È il motivo principale per cui i
giochi delle grandi case vanno bene o male. Quando si riescono a replicare le
condizioni che hanno prodotto \emph{Warcraft~II}, esce \emph{Warcraft~III};
quando non ci si riesce, esce ``il nuovo \emph{Dragon Age}'': la gente non è
arrabbiata perché il gioco non è bello, ma perché non è come gli altri
\emph{Dragon Age}; non ne replica il \emph{feeling}, anche se potrebbe
essere un gioco favoloso. L'eccezione positiva è il passaggio da \emph{Dune} a
\emph{Dune~II}: il primo era di fatto un'avventura grafica, il secondo ha
generato il genere degli strategici in tempo reale. Gli sviluppatori avevano
capito che il successo non stava nel tipo di gioco, ma nell'ambientazione che
attirava le persone: un bel gioco in quell'ambientazione sarebbe piaciuto
comunque. Il relatore comprò \emph{Dune~II} aspettandosi un seguito
dell'avventura e, pur sorpreso, lo finì.

\subsection{Lying, hiding, faking}

Quando si è valutati su un singolo punto nel tempo, e da quel punto dipendono
bonus, stipendio e il rischio di licenziamento, si comincia a fare
\textbf{lying, hiding e faking}:
\begin{description}
  \item[Lying] dire che è sempre andato tutto bene, che si finirà sempre in
        orario, che non c'è nessun problema;
  \item[Hiding] nascondere: ``non era un mio task'', ``non c'ero'', ``non ho
        visto niente'';
  \item[Faking] fingere, anche in cose apparentemente innocue come ``sì, sì,
        ho capito tutto'' quando si ha solo paura di dire ``me lo rispieghi?
        Non ho capito niente''.
\end{description}

\subsection{``Non ho capito'' e la conferma delle istruzioni}

Dire ``non ho capito, ripetimelo'' è ancora visto come una figura da sciocco.
Il relatore non si arrabbierà mai se qualcuno gli chiede di rispiegare
un'intera slide: le cose ovvie per lui possono non esserlo per altri. Più si
conosce qualcosa, \textbf{più è difficile spiegarla}, perché si danno per
scontate molte cose. Ne parlava proprio con Diego, che gli diceva che fare
l'insegnante diventa sempre più difficile man mano che si diventa esperti. È
come spiegare \emph{EVE Online}, uno dei giochi più complessi in circolazione:
chi lo conosce bene dà per scontato che ``ovviamente non puoi fare questa
cosa'', mentre per un nuovo giocatore non è affatto ovvio. A scuola chi dice
``non ho capito'' viene visto come lo ``scemo''; nessuno si chiede ``non è che
l'ho spiegato male?''.

Il racconto più ``traumatico'' dei ragazzi che fanno stage nelle grandi
aziende è sempre lo stesso: ``Cinzia, devi fare questa cosa.'' Cinzia la fa.
``Non te l'avevo chiesta così.'' Il problema è che nessuno ha
\textbf{confermato che Cinzia avesse capito}. Il consiglio, valido sempre:
\begin{itemize}
  \item quando si \emph{ricevono} istruzioni, ripeterle: ``ho capito bene?
        Vuoi che faccia questa cosa in questo modo?'';
  \item quando si \emph{danno} istruzioni, chiedere a chi le riceve di
        ripeterle, non perché lo si ritenga incapace, ma per capire che cosa
        ha capito dalla propria spiegazione (``quindi vuoi un quadrato?''
        ``No, ho detto un triangolo\dots'' ``Però mi hai detto una forma con
        quattro angoli.'' ``Hai ragione: voglio un quadrato'').
\end{itemize}
Frasi come ``non te lo ripeto due volte'' o ``stammi dietro'' portano solo a
fare le cose come si pensa che vadano fatte e a sentirsi dire a fine progetto
``non la volevo così'': che si tratti di una feature, di un design o di un
modello 3D. (Un partecipante osserva che lo stesso accade insegnando le lingue:
chi insegna dà per scontato, ad esempio, che si sappia dell'esistenza del
genere neutro in tedesco.)

\paragraph{Come convincere il team a chiedere?} Il relatore mette sempre come
premessa che \textbf{l'obiettivo non è trovare la persona più intelligente
nella stanza, ma fare un buon lavoro}. Che Gianni capisca alla prima e un altro
alla quinta non gli interessa, se alla fine entrambi lavorano bene: ripetere
cinque volte è comunque meno lavoro che far rifare le cose sette volte. È un
investimento a lungo termine, e non dissimula alcun fastidio: davvero non gli
dà fastidio. Ma ripetere le stesse parole non aiuta: ``non ho sentito'' si
risolve ripetendo; ``non ho capito'' si risolve cambiando spiegazione. Meglio
ancora: ``dimmi che cosa hai capito tu, e vediamo cosa non torna''. Se
intravedesse nel team uno sguardo del tipo ``è la quinta volta che lo
chiede'', interverrebbe: spesso chi giudica è proprio chi dice ``ho capito'' e
poi sbaglia. Preferisce di gran lunga chi chiede cinque volte.

\subsection{Il feedback più utile che abbiate ricevuto}

\emph{Qual è il feedback più utile che abbiate mai ricevuto?} Il relatore ne ha
almeno dieci che gli hanno cambiato il modo di lavorare. Le risposte:
\begin{itemize}
  \item un semplice \textbf{``grazie''} per un lavoro fatto, più utile di un
        ``questa cosa fa schifo'';
  \item un feedback sulla propria \textbf{impulsività} nel rispondere: fermarsi
        un attimo e capire di cosa si sta parlando prima di intervenire;
  \item \textbf{essere se stessi}, anche al lavoro, e non scusarsi
        continuamente per ogni cosa.
\end{itemize}
I feedback che toccano di più sono quelli \textbf{veri}. Il ``grazie''
significa ``ho visto quello che hai fatto e l'ho apprezzato''. Quelli negativi
che si ricordano sono quelli che hanno fatto realizzare ``questa cosa che
pensavo di saper fare, non la so fare'', oppure ``la so fare, ma questa volta
non mi è venuta'' (non è che se una cosa una volta non riesce non la si sa più
fare: può capitare di sbagliare una porta). Arrivano sempre da un
\textbf{punto cieco}: qualcosa che non si sapeva fosse percepito dagli altri.
Conta più ciò che le persone ti dicono di ciò che pensi che pensino di te; e
ciò che il relatore pensa che gli altri pensino di lui è sempre terribile,
mentre il loro giudizio si rivela molto più positivo del previsto.

\paragraph{Il denominatore comune siete voi.} Si potrebbe pensare che questi
concetti non siano applicabili da chi arriva in un'azienda senza poterne
cambiare la cultura. Ma c'è sempre un denominatore comune a cui applicarli:
\emph{se stessi}. La cultura aziendale può non spingere a dare feedback, può
vedere male chi chiede aiuto; ma il relatore ha visto che quando lui lo fa, le
persone intorno tendono a farlo, e un gruppo che lo fa lavora meglio. Perfino
i CEO più ostici a un certo punto gli chiedevano: ``perché il tuo team ora
performa meglio degli altri?''. La differenza è il default to trust: una
pratica che col tempo fa operare il team \emph{come un team} e non come un
insieme di individui. Un videogioco è un progetto corale: c'è chi fa giochi da
solo, ma non si fa da soli \emph{Red Dead Redemption~2}; e se il gruppo non è
affiatato e allineato non esce \emph{Red Dead Redemption~2}, esce
\emph{Anthem} (``oggi la mia fissazione è \emph{Anthem}, domani sarà
un'altra'').

% ---------------------------------------------------------------------
\section{Come dare un feedback}
\label{sec:z1-dare}

\subsection{Concentrarsi sull'azione e sui suoi effetti}

Il primo principio: \textbf{il feedback non riguarda mai la persona}. Non si
dice ``tu hai fatto questa cosa che non va bene'', ma ``questa cosa non va
bene: come facciamo a non farla più?''. Che sia un comportamento, un'azione,
una decisione o una strategia, ci si concentra su quello, e si
\textbf{esplicita l'effetto} che ha. Nell'esempio del ritardo l'effetto era
``quindici minuti di attesa moltiplicati per venti persone''. In un progetto
può essere una catena:
\begin{quote}
  se non rispetti le deadline, l'animazione che viene dopo resta ferma; se
  loro sono in ritardo, la \emph{vertical slice} è in ritardo; se la vertical
  slice non è pronta per la Gamescom, non andiamo alla Gamescom e non troviamo
  un publisher; se non troviamo un publisher, chiudiamo.
\end{quote}
Più la conseguenza è grave, più dovrebbe sembrare piccolo il gesto di dire
``datemi una mano'': è il \emph{male minore}. Poi, a cose fatte, ci si dà
feedback: cosa ha funzionato, cosa no, cosa fare di più, di meno, cosa smettere
di fare.

\subsection{Il pericolo degli ambienti solo positivi}

Gli ambienti dove ci si fa sempre grandi complimenti e nessuno parla dei
problemi creano difficoltà vere alle aziende. Il problema non è mai la
mancanza di complimenti, ma un \textbf{eccesso di ottimismo}: si comincia a
pensare che le cose andranno come si spera, perché tutti intorno danno solo
feedback positivi.

A Blizzard i feedback determinavano gli \textbf{aumenti annuali}: ognuno
doveva indicare cinque persone che avrebbero dato feedback su di lui, e
ovviamente tutti indicavano gli amici, o colleghi di progetti andati bene. Un
sistema nato con le migliori intenzioni diventava uno \emph{scambio di voti}
(``se mi metti follow ti metto follow''), senza alcun miglioramento.

Al contrario, \textbf{le conversazioni più brutte sono state le più utili}:
affrontare i licenziamenti a Riot, o la cancellazione di \emph{Project Titan}
a Blizzard, il progetto che doveva essere il nuovo MMO dopo \emph{World of
Warcraft} e da cui poi è nato \emph{Overwatch}\footnote{Nella trascrizione si
legge anche ``doveva essere Warcraft 2'', probabilmente un lapsus o un errore
di trascrizione.}. Andando a vedere perché lo si cancellava, c'erano solo
motivazioni razionali e pragmatiche, di business, chiare a tutti, non un ``Mike
Morhaime ci odia''. Sono anche le conversazioni in cui ha lavorato su se
stesso: è molto \textbf{avverso ai conflitti}, e questo a volte gli ha
procurato conseguenze molto più gravi di quanto sarebbe costato un piccolo
disagio. Dieci anni fa faceva fatica ad accettare questo difetto; ha dovuto
subire le conseguenze di conflitti evitati e poi cresciuti.

\subsection{Che cosa fare}

Le linee guida della slide, che il relatore legge rapidamente perché saranno
disponibili nel materiale:
\begin{description}
  \item[Attenzione al linguaggio del corpo] dare un feedback a braccia
        incrociate o guardando altrove è diverso dal guardare l'altro negli
        occhi.
  \item[Sii specifico] con esempi tangibili.
  \item[Esplicita gli effetti] come visto sopra.
  \item[Discuti il prima possibile] non aspettare un mese per parlare di
        qualcosa successo ieri, altrimenti nessuno ricorderà più bene cosa sia
        accaduto.
  \item[Concentrati sull'azione e sul problema, non sulla persona] il nemico è
        l'azione, non la persona che si ha davanti. Non bisogna far cambiare
        l'altro: bisogna evitare che il problema si ripresenti.
\end{description}

\paragraph{Gli insulti sono rivolti al ruolo.} Un esempio solo apparentemente
scollegato. Il relatore ha lavorato su \emph{League of Legends}, ``non
esattamente noto per avere una community da orsetti del cuore'', in varie
\emph{task force} sulla tossicità. A chi smetteva di giocare veniva inviato un
sondaggio sul perché; il 33\% indicava la tossicità della community. Dopo varie
campagne di sensibilizzazione e cambiamenti si è scesi al 30\%: sembra solo
un 3\% in meno, ma valeva circa \textbf{120 milioni di revenue in più
all'anno}, perché quei giocatori restavano esposti a skin e contenuti e
continuavano a comprare. Anche qui, piccoli guadagni che si sommano.

Lavorando con degli psicologi, aveva notato che gli insulti \textbf{non sono
quasi mai rivolti alla persona, ma al ruolo}: ``il mio jungler fa schifo'',
``il mio DPS non è capace'', ``il nostro portiere è una pippa''. Chiamare
qualcuno per ruolo lo \emph{depersonalizza} e rende più facile insultarlo. Lo
stesso nell'esempio del giro dell'ufficio: non ``non chiedere a Giulio'', ma
``non chiedere a quel grafico, consegna sempre in ritardo''. Concentrarsi
sull'azione, senza attaccare la persona, fa sì che il feedback non sia (e non
venga percepito come) un attacco personale.

\paragraph{Considera il carattere.} Qui tornano utili i test di personalità.
Servono più agli studenti che al relatore (che sa già che usciranno 53 profili
diversi, senza un pattern): bisogna conoscere il proprio profilo e
\textbf{condividerlo con chi lavora con noi}, chiedendo anche agli altri di
fare lo stesso. In ufficio il relatore ha usato una versione più avanzata, a
pagamento, con un report di circa trenta pagine: il \textbf{profilo
DISC}\footnote{Nella trascrizione ``disk''. Si tratta del modello DISC
(Dominance, Influence, Steadiness, Conscientiousness).}. Ha mappato tutta
l'azienda e ha fatto stampare il profilo di ciascuno sulla sua scrivania, così
che tutti sapessero con chi stavano parlando. Capire come ragionano gli altri
(``è introverso, non gliene parlo davanti a tutti'') è una delle cose che gli
sono tornate più utili: \emph{adattare ciò che si vuole dire a chi si ha
davanti}.

\paragraph{Assicurati di non disturbare.} Sotto il profilo di ognuno c'era un
piccolo \textbf{semaforo}: verde, ``parlami pure, interrompimi''; giallo,
``interrompimi solo se è urgente''; rosso, ``mi sto concentrando, non mi
parlare''.

\paragraph{Lascia finire il lavoro.} Durante gli eventi dal vivo il relatore si
era reso conto di quanto fosse grave dare feedback mentre qualcuno sta ancora
lavorando. Se una cosa sta andando male, dirlo in quel momento non aiuta: tra
la pressione di finire e quella del feedback, peggiorano entrambe le cose, e
anche il rapporto. Meglio lasciarla finire, anche male, e parlarne dopo.

\subsection{Che cosa non fare}

\begin{description}
  \item[Non basarti su supposizioni] ``mi hanno detto che ieri in un meeting
        hai fatto\dots'': o c'eri e l'hai visto, o non esiste.
  \item[Non giudicare la persona.]
  \item[Non fornire feedback vago.]
  \item[Non accettare e non fornire feedback di terze parti] Se qualcuno vi
        dicesse ``mi hanno detto che hai fatto un brutto lavoro'', il primo
        pensiero non sarebbe ``cosa ho fatto male?'', ma \emph{``chi te l'ha
        detto? Chi è la vipera che ho in seno? Chi è l'impostore?''}. Siamo
        fatti così. Ogni conversazione che inizia con ``mi hanno detto che'' è
        inutile: l'altro non ascolterà, penserà solo a chi dare la caccia, e
        spesso penserà alla persona sbagliata, creando antipatie tra persone
        che non c'entravano nulla. Quando al relatore dicono ``mi hanno detto
        che lui ha fatto\dots'', risponde ``vai a parlarne con lui''; se gli
        dicono ``mi hanno detto che tu hai fatto\dots'', risponde ``mandami la
        persona che te l'ha detto, altrimenti non ne parliamo''. Non per
        cattiveria: se chiedesse ``che cosa ho detto di preciso a Marco?'' e la
        risposta fosse ``non lo so, non c'ero'', di cosa si starebbe parlando?
        \textbf{Parlate con la fonte}, altrimenti fate come se quel feedback non
        esistesse.
  \item[Non dare consigli non richiesti] come visto, non portare la soluzione:
        la soluzione deve venire da chi riceve il feedback.
  \item[``Non mi piace'' non è un feedback] ``Penso che questo design non
        funzioni perché, visto che il nostro obiettivo è\dots'' ha argomenti;
        ``non mi piace il combat system'' è un'opinione. Anche ``io l'avrei
        fatto diversamente'' non è un feedback.
  \item[Accettare non vuol dire fare come proponi] deve essere chiaro a chi dà
        il feedback: l'aspettativa è essere ascoltati e presi in
        considerazione, ma se la responsabilità è dell'altro, deciderà
        l'altro.
  \item[Non offenderti] se il feedback non viene seguito. Restarci male
        significa pensare di saperne più dell'altro, cioè \emph{venire meno al
        default to trust}. ``Decidi tu'' significa ``mi fido di te, perché sei
        lo specialista, e mi fido che mi ascolterai e farai la cosa migliore
        per tutti'', che potrebbe anche non essere ciò che ho detto io. Se
        voglio che l'altro faccia sempre come dico io, mi sto fidando solo di
        me stesso.
\end{description}

\subsection{Il primo feedback e il debriefing}

La cosa più difficile è dare il \textbf{primo} feedback: è un muscolo da
allenare. Chi non l'ha mai fatto (soprattutto con questa struttura) dovrebbe
provarci, anche su minuzie (``abbiamo concordato di usare WhatsApp nel gruppo:
non potremmo usare Telegram?''), anche solo per vedere come va. Altre
indicazioni:
\begin{itemize}
  \item feedback \textbf{solo uno a uno}, mai davanti ad altri, soprattutto se
        negativo;
  \item \textbf{prepararsi} prima;
  \item fare il \textbf{debriefing del feedback}: dopo, fermarsi a pensare a
        com'è andato. Il relatore, quando media un feedback tra due persone,
        si sente spesso rispondere ``benissimo'' da una e ``malissimo''
        dall'altra: la percezione è molto diversa. Non si fa problemi a tornare
        dalla stessa persona: ``ieri ti ho dato un feedback: com'è andata? Sei
        soddisfatto? Come poteva andare meglio?''. Con le persone con cui si ha
        meno a che fare, questo aiuta molto la volta successiva.
\end{itemize}

% ---------------------------------------------------------------------
\section{Come ricevere un feedback}
\label{sec:z1-ricevere}

Per il relatore è la parte più difficile. Si riassume in quattro passi:
\begin{enumerate}
  \item \textbf{Ascolta}: si parte sempre ascoltando.
  \item \textbf{Chiedi dettagli}, se serve, ma per capire meglio, non per
        difenderti.
  \item \textbf{Accetta, a prescindere}: non si sta decidendo chi ha ragione;
        si sta ascoltando un'opinione. Non c'è nulla da cui difendersi.
  \item \textbf{Agisci}: fatti un piano. Può essere banale (``mi hanno detto
        che interrompo sempre: cerco di non interrompere''), oppure ``devo
        migliorare questa cosa: come?''. Pensa anche a \emph{chi} potrebbe
        aiutarti e a \emph{che supporto} ti serve: il piano non deve essere
        per forza ``devo risolvere tutto da solo''.
\end{enumerate}
Una domanda che molti sottovalutano: \emph{``sono junior in un'azienda: sarò
mica il primo ad avere questo problema?''}. Probabilmente qualcun altro l'ha
già avuto: basta chiedere ``tu come hai fatto?''.

\subsection{Il feedback positivo}

Tutto questo vale per i feedback costruttivi negativi, a cui si pensa di
default; ma esistono anche i \textbf{feedback positivi}, e vanno usati. Quando
si riceve un ``grazie'' o un ``bravo'', la domanda è: \emph{come faccio a
ricreare le stesse condizioni?}

Il relatore ha lavorato anche in aziende dove il feedback era solo negativo.
Il rischio? Si immagini l'impiegato ideale: vuole sempre migliorare, è sempre
pronto a dire ``magari sbaglio io''. Se nessuno gli dice cosa sta facendo
\emph{bene}, spinto dal desiderio di migliorare rischia di cambiare cose che
erano già buone, e di farle peggio. Il rinforzo positivo \textbf{non è solo una
gentilezza}: è un atto pratico, che evita che le persone peggiorino ciò che
sanno già fare bene.

\subsection{Il feedback accidentale}

Esiste infine il \textbf{feedback accidentale}, molto frequente negli uffici:
\begin{itemize}
  \item il ``chiedi a lei, non a quel grafico'' del giro dell'ufficio è un
        feedback negativo su una persona, dato accidentalmente, né richiesto
        né voluto;
  \item a Riot li chiamavano \textbf{drive-by}, come nelle scene dei film in
        cui passa un'auto, si abbassa il finestrino, si spara e si riparte. In
        molte aziende con una ``cultura del feedback'' in cui tutti devono
        dire la loro, si riceve una valanga di commenti del tipo ``comunque il
        sito a me non piace, ciao''. Non si sa perché, non si sa cosa fare di
        diverso, non si sa che rischi ci siano: \emph{a che cosa è servita
        questa conversazione?} A niente.
\end{itemize}

% ---------------------------------------------------------------------
\section{Quando, quanto e con chi}
\label{sec:z1-quando}

\subsection{Frequentemente e recentemente}

\emph{Quando avete ricevuto il vostro ultimo feedback?} Per gli studenti,
qualche giorno fa (nel corso) o qualche mese fa (al lavoro, dopo una
consegna). Come aveva osservato Silvia, il feedback andrebbe dato
spesso.

Il relatore usa l'esempio di una nave da crociera (``la Costa Concordia'') che
deve andare da A a B. La rotta ideale (azzurra) è quasi impossibile: nessun
progetto va esattamente come previsto. C'è la rotta ``panoramica'' (verde), che
gira intorno all'isola d'Elba; ci sono poi rotte (arancione e rosa) che a un
certo punto si allontanano molto e poi correggono bruscamente: qualcuno si è
accorto che si stava sbagliando, ``ma non è quello che avevo chiesto'',
``non è dove dovevamo andare''. Sono i famosi \textbf{sanity check} o
\emph{checkpoint}.

L'ideale è avere \textbf{tanti piccoli check per fare tante micro-correzioni}:
correggere la rotta di un paio di gradi è molto più rapido che farlo dopo
essersi allontanati. Bastano trenta secondi di ``ehi, vado bene?''. Spesso non
lo si fa per paura di sembrare insicuri o poco autonomi, ma per il relatore è
l'opposto: quando qualcuno gli chiede ``vado bene così?'' e lui risponde
``sì'', si alleggerisce anche del proprio carico mentale. Un capo che assegna
un compito e lo vede solo alla fine vive sempre col dubbio (``avranno capito?
Lo staranno facendo giusto? Arriverà in tempo?''); con un check di dieci
secondi al giorno (``tutto ok, Stefano'') è tranquillo, e la persona va avanti
sapendo che quanto fatto finora va bene. Non serve fare ogni giorno il giro del
dipartimento: anche ogni uno o due giorni basta a evitare la rotta verde.

Per il feedback è lo stesso: \textbf{va dato subito}. Dire oggi ``ieri hai
fatto questa cosa'' è molto diverso da dire ``negli ultimi sei mesi hai fatto
questa cosa''. Il relatore si è sentito dire più volte, alla valutazione di
dicembre, ``da maggio hai fatto questa cosa che non andava''. La sua risposta:
``e aspetti dicembre per dirmelo? Mi hai lasciato sbagliare per sei mesi?''. È
una mancanza di rispetto, e la conversazione non può più essere costruttiva.
Per questo non ama le valutazioni annuali: i manager crescono abituandosi a
prendere appunti delle cose negative per tirarle fuori alla valutazione
(``a marzo hai fatto questa cosa'' ``sì, però poi ho smesso'' ``sì, però a
marzo l'hai fatta'').

In sintesi:
\begin{description}
  \item[Frequentemente] riguarda l'intervallo tra un check e l'altro.
  \item[Recentemente] significa non lasciar passare un anno prima di dare il
        feedback, altrimenti non serve più a niente.
\end{description}

\subsection{Quanto dura un feedback}

Un feedback dura \textbf{da trenta secondi a quindici minuti}; raramente di
più. Chi non ne ha mai dato uno dovrebbe sforzarsi di darne almeno uno, magari
positivo per cominciare. Ogni feedback dato o ricevuto è un
\textbf{deposito nel proprio ``fondo fiduciario''} (\emph{trust fund}), che
costruisce la reputazione futura. Il relatore si è costruito la reputazione di
chi dà feedback il più possibile oggettivi e analitici: giusti o sbagliati che
siano, costringono a pensare. Non impone mai il suo punto di vista (``si fa
come dico io''), ma chiede ``tu come faresti? Come posso aiutarti?''. Il
feedback, per lui, serve a \emph{far crescere qualcun altro}.

\subsection{I manager peggiori: ``a che numero sto pensando?''}

I peggiori manager che il relatore abbia avuto erano quelli che volevano che si
facesse come volevano loro, anche quando chiedevano ``tu come la faresti?''.
In realtà la domanda era ``\emph{a che numero sto pensando?}'': se dicevi 5,
``no, dovremmo fare 9''; se dicevi 9, ``no, 5''. Lo stress costante di dover
indovinare cosa vuole il capo nasce dalla mancanza di fiducia (``se mi chiedi
come fare, fammelo fare'').

Lo stesso gli succede con i CEO che lo chiamano per ``sistemare la cultura
aziendale'' e poi dicono ``non sono sicuro, io farei diversamente''. La sua
risposta: ``lo so che faresti diversamente, per questo mi hai chiamato''. Se si
chiama un esperto, bisogna lasciargli fare il suo lavoro: non si va dal
meccanico a cambiare le gomme per poi dirgli ``non dovresti girare un po' di
più quel bullone?''. Chiedere un parere a un esperto significa dargli anche il
potere di scegliere. Il problema è il \emph{micromanagement}, o in generale il
non accettare punti di vista diversi dal proprio (magari perché si è il
fondatore e si pensa di essere l'unico a saper fare le cose).

\subsection{Dare feedback a un superiore egocentrico}

\paragraph{Come si dà un feedback negativo a un superiore con un grande ego?}
Il trucco è \textbf{non puntare sull'ego, ma sull'azione e sui suoi effetti}.
Non ``tu fai così e io mi sento così'': anzi, si alimenta il suo ego
rendendolo il \emph{main character} dell'avventura, l'eroe che risolve il
problema (``tu, grande manager, aiutaci a risolvere questo problema che
affligge noi che facciamo l'operatività''). Non è l'ideale (``io cambierei
manager''), ma pochissimi possono scegliersi i colleghi: il relatore stesso ha
quasi sempre ricevuto team preassemblati (``ecco le chiavi, questa è la tua
macchina''). Bisogna giocare con le carte che si hanno in mano.

\paragraph{Un esempio pratico.} Con un manager che fa micromanagement e mette
sempre ``il becco'' nel proprio lavoro, si può dire: ``se ti fidi di me per
questo lavoro, con le tue capacità potresti dedicarti a questa cosa più
strategica e di lungo termine: pensa quanto sarebbe più utile per l'azienda''.
In sostanza: ``guarda cosa potresti raggiungere se mi lasciassi fare il mio
lavoro''. Non come un ``tu fai il tuo e io faccio il mio'', ma come un
``ti sollevo da questa parte del carico, che so che ritieni importante: lascia
che ci pensi io, così puoi dedicarti ad altro'', che magari gli dà anche più
visibilità.

Si gioca su ciò che quella persona pensa di saper fare, e che probabilmente sa
anche fare. Anche Michael Scott di \emph{The Office}, per quanto preso in giro,
è raccontato come un venditore bravissimo, che ha chiuso i contratti più
importanti per l'azienda, e rispetto ai manager che gli succedono è forse
``uno dei meno peggiori''. Tutti, anche i manager peggiori, hanno punti di
forza: \textbf{giocate su quelli}.

% ---------------------------------------------------------------------
\section{Il feedback come regalo}
\label{sec:z1-regalo}

Per chiudere, il relatore risponde alla domanda d'apertura, che aveva lasciato
apposta in sospeso: che cos'è il feedback per lui?

\begin{quote}
  \emph{Il feedback è un regalo che mi viene fatto da qualcuno che so che vuole
  il mio interesse.}
\end{quote}

Ha costruito un team di persone che si fidano di lui e che non hanno motivo di
sabotarlo o di ferirlo: per questo interpreta anche le cose peggiori con il
default to trust. È un regalo, e cerca di prenderlo in considerazione il più
possibile.

Ci è arrivato partendo da se stesso. Oltre a essere avverso ai conflitti, si
definisce ``avverso alla comunicazione umana'': paradossalmente, una persona
introversa e timida finita a fare comunicazione, marketing e ora HR (``succede
spesso''). Prima di ordinare una pizza al telefono ha dieci minuti di
preparazione mentale, e se l'ordine delle domande cambia va in confusione. Per
lui un feedback di dieci minuti sono spesso \textbf{tre ore di preparazione}:
costruisce un \textbf{albero dei dialoghi}, come in un'avventura grafica
personale, con tutte le opzioni. Poi spesso finisce con ``hai fatto questa
cosa, andrebbe migliorata'' ``hai ragione, grazie'', e basta.

Così ha pensato: chissà quanti si fanno gli stessi problemi con lui, e passano
ore a prepararsi per dirgli ``questa cosa l'hai fatta male'', sapendo che
riceveranno comunque quell'occhiata fugace del ``come ti permetti''. Eppure
vengono a dirglielo. Non può far altro che rispondere: ``grazie, hai ragione,
cercherò di fare meglio''. È così che ha metabolizzato il ricevere feedback:
capendo che non c'è nessuna cattiva intenzione e che, se non proprio felice
(``ci ho messo una pietra sopra''), può almeno essere \emph{grato} che ci sia
ancora qualcuno che ha così a cuore il suo interesse da dargli feedback.

\paragraph{Il silenzio fa più paura.} Non bisogna spaventarsi dei feedback
negativi che emergono dai sondaggi aziendali: bisogna spaventarsi del
\textbf{silenzio}, quando nessuno dice più nulla. Significa che nessuno si
fida più che si sappiano risolvere i problemi.

\paragraph{La sindrome dell'impostore.} Il suo motivatore personale, all'inizio,
era soprattutto ``non voglio fallire'': per questo cercava di migliorare, e
migliorare è diventato accettare feedback, perché più feedback riceve, più
migliora. Il feedback è un \textbf{acceleratore}, e lui cerca di incentivare
gli altri a darglielo. Resta però sempre la vocina del ``e se si accorgono che
non sei capace?''. Con la \emph{impostor syndrome} convive da anni e
continuerà a conviverci; se ce l'avete, conviveteci anche voi. Avrete sempre
la sensazione che ``se ne accorgeranno'': forse sì, ma sapete fare anche
altre cose. Non concentratevi solo su ciò che non sapete fare. Gli errori
capitano: la perfezione non è di questo mondo.

\paragraph{Esercitarsi in accademia.} Per fare pratica con il feedback,
l'ambiente migliore è quello dei colleghi di corso: persone nella stessa
situazione, con esperienza simile e lo stesso obiettivo, in un contesto in cui
non si può essere licenziati e l'unico effetto a lungo termine è ``ho parlato
con questa persona, ora le cose vanno meglio''. E bisogna avere il coraggio di
parlare anche con i docenti (``dire a Diego: secondo me questa lezione potevi
farla in questo modo''; ``so che è una persona che lo accetta''), senza farsi
bloccare dalla gerarchia. Nell'azienda giusta (come quella in cui lavora ora) il
relatore va sempre dal CEO a dare feedback, sapendo che non verrà licenziato:
il CEO gli ha detto ``ti ho preso per darmi feedback''.

La frase di chiusura:

\begin{quote}
  \emph{Per superare le vostre difficoltà, tutto quello che vi viene chiesto è
  il coraggio di chiedere aiuto.}
\end{quote}

Chiedere aiuto non è una debolezza: è un segno di consapevolezza di sé
(\emph{self-awareness}) e di forza personale. Se tutti lo facessero, i
progetti andrebbero molto meglio; spesso invece ci si fa frenare dal ``non
voglio fare brutta figura'', ``non voglio passare per quello scemo'', e sono
proprio queste paure a rallentarci e a impedirci di raggiungere gli obiettivi.

% ---------------------------------------------------------------------
\section{Domande e risposte finali}
\label{sec:z1-qa-finale}

L'ultima mezz'ora è dedicata a domande libere.

\paragraph{E se una persona peggiora per uno o due anni di fila?} Se nessuno
gliel'ha mai detto, c'è un grosso problema, non suo ma dell'azienda. Se invece
c'è consapevolezza, il percorso è: ``cosa vuoi fare per risolvere?''; ``ho
provato questo, questo e questo''; ``visto che non ha funzionato, vuoi un mio
consiglio? Proviamo con quello?''; ``che supporto ti serve?''. Se, con tutto il
supporto, il tempo e le soluzioni concordate per mesi, il miglioramento non
arriva, per il relatore la soluzione è: \textbf{``dimmi un prezzo per andare
via''}. Preferisce pagare una buonuscita (\emph{severance package}) che
avviare la strada dei richiami formali sulla performance fino al
licenziamento per giusta causa. A un certo punto andare avanti non è un favore
per nessuno.

\paragraph{È vero che nei videogiochi si viene sfruttati? Meglio un piccolo
studio italiano?} Per il relatore si è quasi più a rischio a Blizzard che in un
piccolo studio di Roma o Torino. Nel piccolo studio si sa dove ci si va a
cacciare: non ci sarà il team building a Santo Domingo, ma ci si concentra su
un gioco alla volta e, se va bene, si fa il successivo. A Blizzard si può
lavorare su \emph{Overwatch}, che vende milioni, e perdere comunque il posto,
in modo totalmente scollegato dal proprio lavoro. Gli stipendi più alti
all'estero compensano anche questo rischio, che spesso non dipende da come va
il gioco (anzi, ``il gioco non ha fatto come volevamo'' è spesso una scusa per
licenziamenti che non avevano motivo).

\paragraph{Cambiare settore a trent'anni, da web developer full remote?} Lo
stipendio indicato dallo studente non è fuori dalla portata di vari studi
italiani, anzi è superabile. Non va esclusa la possibilità di lavorare da
remoto per aziende piccole ma estere. Alcuni amici del relatore, pagati con
stipendi italiani, si sono trasferiti a Bali o in paesi con un costo della vita
più basso; il collega di \emph{PUBG} percepisce uno stipendio olandese vivendo
a Milano, che per quanto cara non costa come i Paesi Bassi.

\paragraph{Come si diventa manager?} Dipende dall'azienda: alcune promuovono i
migliori (processo tendenzialmente sbagliato), altre propongono un percorso di
\emph{people management}. Diventare manager significa \textbf{abbandonare ciò
che si sa fare per imparare a fare altro}. Il relatore ha dovuto accettare che
non sarebbe più stato ``il miglior community manager italiano del mondo''
(``eravamo in pochi, quindi era facile''): dedicandosi alle persone non poteva
più fare il suo lavoro bene come prima. È stata una negoziazione con se
stesso: fare il manager significa non essere più l'esperto di qualcosa, ma
saper gestire e far crescere persone. Significa anche \textbf{cedere il
controllo della propria carriera}: se il team decidesse di non fare più nulla,
il manager ne pagherebbe le conseguenze. Il suo lavoro diventa scegliere le
persone giuste e farle lavorare bene. Lui lo esplicita sempre: sa di affidare
al team parte della propria carriera, e il team sa che lui combatterà per
loro, anche prendendosi le colpe.

\paragraph{L'inglese.} \textbf{Va assolutamente imparato}, anche lavorando in
Italia. Non saperlo preclude circa il 95\% del mercato (tutto quello fuori
dall'Italia) e soprattutto il \emph{networking} con persone che potrebbero dare
occasioni future: l'inglese è spesso l'unica lingua in comune con chi lavora
all'estero ed è oggi la lingua universale dei videogiochi.
\begin{itemize}
  \item \emph{Che livello serve?} Non importa B2 o C2: bisogna saper sostenere
        una conversazione, guardare almeno Netflix con sottotitoli in inglese,
        comunicare senza dodici pause e senza ``come si dice\dots''. Il
        certificato non verrà mai chiesto: verrà chiesto se si sa parlare
        inglese. Il relatore ha sentito persone con un C2 che non sapevano
        parlarlo.
  \item \emph{Come allenarlo senza occasioni di conversazione?} Netflix, Prime
        Video, Now TV, Udemy, Coursera, YouTube, Duolingo, siti di scambio
        linguistico con madrelingua, Discord: non sono mai esistiti così tanti
        mezzi gratuiti. Il relatore ha imparato l'inglese giocando (con una
        pronuncia sbagliata, perché leggeva senza sentire, su cui ha poi
        dovuto lavorare); guardare un film in inglese con i sottotitoli in
        italiano insegna anche la traduzione.
  \item Oltre all'ascolto passivo è fondamentale \textbf{parlare}: in una
        flotta di \emph{EVE Online} con 150 persone quasi tutte madrelingua
        inglese si è esposti a decine di accenti e modi di dire (mille modi di
        dire ``come stai'' oltre a \emph{how are you}); lavorando all'estero si
        impara a distinguere gli accenti (Manchester, Londra) e ad apprezzarli
        anche nei media.
\end{itemize}

\paragraph{Quanto deve sapere un junior?} Dipende da quanto è tecnico il
ruolo. Un game designer junior che non sa qualcosa può essere affiancato da un
senior; un programmatore che non sa programmare, in quel momento, non è utile
all'azienda.

\paragraph{Annunci ``junior con cinque anni di esperienza''.} ``Cercasi
ventenne con trent'anni di esperienza.'' Non è l'annuncio per un junior: è un
errore dell'azienda, che in realtà cerca un \emph{mid}. Riceverà profili
junior, si stupirà che non abbiano cinque anni di esperienza, rifarà l'annuncio
e avrà perso sei mesi. Il consiglio: tenere d'occhio quella pagina, perché
dopo qualche mese l'annuncio verrà corretto.

\paragraph{Una struttura per il post-mortem.} A chi chiede un formato di base,
il relatore anticipa quanto vedrà nella seconda giornata a proposito di
brainstorming: una lavagna \textbf{divisa in quattro aree}, su cui ognuno
attacca tutti i post-it che vuole:
\begin{center}
  \begin{tabular}{@{}p{0.44\textwidth}p{0.44\textwidth}@{}}
    \toprule
    \textbf{Smettere di fare} & \textbf{Iniziare a fare} \\
    ciò che non ha funzionato & ciò che non abbiamo fatto e andava fatto \\
    \midrule
    \textbf{Fare di meno} & \textbf{Fare di più} \\
    ciò che ha fatto perdere tempo rispetto all'impatto ottenuto &
    ciò che ha funzionato ma, fatto con troppa leggerezza, non ha dato i
    frutti sperati \\
    \bottomrule
  \end{tabular}
\end{center}
Con i post-it ognuno ha voce e tutto risulta catalogato; per il report basta
una foto della lavagna (o uno screenshot, se la lavagna è digitale, come Miro).

\paragraph{Come sbloccare la paralisi del feedback verso il management?} Uno
studente ha visto manager dire ``non abbiate paura, la mia porta è sempre
aperta'' senza che cambiasse nulla. C'è una dissonanza tra ciò che dicono e
ciò che fanno: è ``la mamma che dice vieni qua che non ti faccio niente''.
Qualcosa deve aver fatto pensare alle persone che non è sicuro. Prima va
sistemato quello; poi si può provare con i \textbf{feedback anonimi}.
L'obiettivo del relatore resta il feedback aperto e firmato, ma a Stormind ha
iniziato con feedback anonimi verso il management per verificare se
l'anonimato fosse un fattore. Non lo era (evidentemente non c'erano problemi a
parlare con il CEO), e la volta successiva intende abbandonarli. Ma è un buon
\emph{check}: se senza anonimato non arriva nulla e con l'anonimato arrivano
cento feedback, il problema è chiaro.

\paragraph{Quando è giusto rivolgersi a un superiore?} La strada giusta resta
parlare direttamente con l'interessato. Ma se si prova e non si trova una
soluzione, la cosa più ovvia è andare dal proprio manager insieme: ``ci
aiuti?''. Quando il relatore parlava di mediare i dialoghi, era proprio questo:
fare da moderatore, aiutare a sbloccare i punti fermi o, se non si riesce a
scegliere tra A e B, scegliere lui. L'escalation, quando necessaria, va fatta.

\paragraph{Il tempo libero veniva rispettato?} No, e se l'è dovuto imporre
lui. Il lavoro gli piaceva così tanto che continuare a casa non gli sembrava
lavorare. Gli si è ritorto contro con il lancio di \emph{Diablo~III}, fatto in
orario americano (le due o l'una di notte in Europa), con l'ufficio chiuso
perché a nessuno era stato chiesto di restare. Il lancio va malissimo, i
giocatori protestano, e lui, sveglio, risponde sui forum fino alle tre, quattro,
cinque del mattino: ``non posso risolvere il problema, ma resto qui ad
aggiornarvi con quello che riesco a sapere dal team di sviluppo''. Il giorno
dopo trova sulla scrivania un \textbf{richiamo formale dell'HR}: aver lavorato
fuori orario senza autorizzazione, condotta non professionale. Per lui è stato
un interruttore: ``ora me ne vado''. Due settimane dopo faceva colloqui con
Riot. Ancora oggi è contrario a quel richiamo, ma gli ha fatto capire come
ragiona un'azienda: ``hai un contratto 9--18, non mi interessa se l'hai fatto
con le migliori intenzioni''. Eppure per anni ha goduto di un'ottima
reputazione in quelle community come ``quello che è rimasto con noi fino alle
cinque del mattino'', e anche l'azienda ne ha beneficiato.

Ci sono poi aziende in cui, soprattutto per i developer, viene richiesto il
famigerato \textbf{crunch}: lavoro straordinario continuativo, senza pause.
Dipende molto dall'azienda; sta venendo (``grazie al cielo'') deprecato, perché
molte aziende capiscono che non è sano ed è visto male anche dal pubblico, ma
esiste ancora.

\paragraph{E se il tuo capo demotiva qualcuno del tuo team?} Qui ``diventa
personale'': il team del relatore si fida di lui perché sa che in questi casi
``diventa una bestia''. ``Te la puoi prendere con me quanto vuoi, ma il mio
team non lo tocchi.'' Lo gestisce direttamente con l'altro manager, per quanto
brutta possa diventare la conversazione. Lo sforzo più grande è \textbf{non
andarne poi a parlare con il team} (``ragazzi, ho fatto questa cosa''), perché
minerebbe la fiducia verso l'alto: non vuole un rapporto ottimo fino a lui e
pessimo oltre. L'obiettivo è che l'altro manager ``si faccia perdonare''
cambiando il proprio comportamento, e che il team percepisca da solo il
cambiamento.

\paragraph{Chiusura della giornata.} Il relatore chiede un feedback collettivo
in chat (``se domani ci dai \emph{more of the same} siamo contenti, oppure
ditemi cosa limare''), si scusa per i ripetuti crash della connessione e
ringrazia per la proattività nelle domande, che gli servono anche come
\emph{check}: se nessuno chiede nulla su un argomento, significa che può
andare avanti. Racconta di aver passato una settimana a decidere che cosa
mettere nelle slide e di averle poi preparate in poco tempo (rifacendone
alcune fino alle tre di notte). Ricorda infine di usare il test di
personalità per sé e per il proprio gruppo, facendolo fare a chi non l'ha
ancora fatto: tornerà molto utile per cooperare meglio.

% ---------------------------------------------------------------------
\section{Anticipazioni sulla seconda giornata}
\label{sec:z1-domani}

La seconda giornata del workshop sarà più pratica: ``oggi vi ho spiegato
perché, domani vi insegno come lo fanno queste aziende''. Gli argomenti
previsti:
\begin{description}
  \item[Time management] come organizzarsi una giornata professionale in
        azienda: come strutturare una \emph{task list}, cosa fare e cosa non
        fare, come e cosa prioritizzare.
  \item[Brainstorming] i gruppi dovranno fare moltissimi meeting (come facciamo
        il gioco, come lo chiamiamo, che scopo gli diamo). Se non si sanno fare
        bene, i brainstorming sono \textbf{dannosi ed eterni}: un brainstorming
        dovrebbe durare \textbf{al massimo due ore}. Verranno mostrati metodi
        che accorciano la durata, massimizzano l'effetto e, soprattutto,
        \textbf{fanno partecipare tutti}: non c'è niente di peggio di un
        brainstorming in cui parlano in due per il 90\% del tempo. Il classico
        ``ragazzi, non esistono idee sbagliate, sparate'' è il modo peggiore di
        farlo; non è colpa di nessuno, ma si può fare diversamente.
  \item[Guidare da dentro] come essere leader nel proprio team anche senza
        esserne il leader: sviluppare responsabilità personale e verso gli
        altri. Non essere il manager di qualcuno non significa non esserne
        responsabili: la persona accanto è responsabilità di tutti, perché si
        è un team e non un insieme di individui. ``Se qualcuno cade, tutti lo
        facciamo rialzare: non c'è l'addetto al rialzamento.'' È la versione
        \emph{advanced} della giornata: come comportarsi da manager senza
        averne il titolo.
\end{description}
Torneranno inoltre il \emph{service level agreement} e il tema della
condivisione dei successi e della protezione del team dai fallimenti.

% ---------------------------------------------------------------------
\section*{Punti chiave del capitolo}
\addcontentsline{toc}{section}{Punti chiave del capitolo}

\begin{itemize}
  \item ``La capacità di un gruppo di persone di fare cose straordinarie
        dipende dal modo in cui queste persone operano insieme.'' Il 90\% dei
        problemi dei team sono \textbf{problemi di comunicazione}.
  \item Le aziende si liberano sempre più presto dello \textbf{useless
        genius}: le competenze si trasmettono, il carattere no. Il
        \textbf{fit culturale} va valutato anche dal candidato.
  \item Non esiste un unico modo di dare feedback: dipende dal contesto e
        soprattutto da \textbf{chi lo riceve}. Il feedback è una somma di
        piccoli miglioramenti periodici; tenersi le cose nel cassetto non
        serve mai.
  \item I team attraversano \textbf{forming, storming, norming, performing e
        adjourning}. Lo storming non si salta; si passa la maggior parte del
        tempo in forming; la bravura sta nell'accelerare i passaggi e nel
        fare bene il \textbf{post-mortem}, anche dei successi.
  \item Gli obiettivi di un progetto vanno fissati \textbf{all'inizio}: solo
        così si può valutare se è andato bene e scegliere tra A e B.
  \item Il circolo virtuoso è \textbf{feedback $\rightarrow$ miglioramento
        $\rightarrow$ fiducia}. Accettare un feedback non significa fare come
        dice l'altro: decide chi ha la responsabilità, ma deve argomentare.
  \item ``Aiutami a poterti aiutare'': meglio dire in anticipo ``potrei non
        farcela'' che ``non ce l'ho fatta''. Chiedere aiuto non è una
        debolezza.
  \item Quando si dà feedback \textbf{non si porta la soluzione}: la si
        costruisce insieme a chi lo riceve.
  \item La cultura è forte quando si lavora \textbf{per gli altri}; il lavoro
        è valutato come opera d'insieme (anche con bonus di team).
  \item Come in \emph{Among Us}, i team sprecano energie a cercare il
        traditore invece di fare le task. \textbf{Default to trust}: chi ha
        sbagliato ha agito con le informazioni che aveva. Cercare
        \textbf{la causa, non il colpevole}: ``che cosa ti ha fatto pensare
        che fosse la scelta giusta?''.
  \item Carriere e performance non sono lineari: vanno valutate le
        \textbf{traiettorie}, non un punto dell'anno. Le valutazioni puntuali
        generano \emph{lying, hiding, faking}.
  \item Confermare sempre le istruzioni facendosele ripetere; ``non ho
        capito'' non è una colpa.
  \item Dare feedback: sull'\textbf{azione} e sui suoi \textbf{effetti}, specifico,
        presto, uno a uno, adattato al carattere, senza terze parti, senza
        ``non mi piace'', senza offendersi se non viene seguito; poi fare il
        debriefing.
  \item Ricevere feedback: \textbf{ascolta, chiedi dettagli, accetta, agisci}.
        Il feedback positivo serve a non peggiorare ciò che si fa bene; il
        feedback accidentale e il \emph{drive-by} non servono a nulla.
  \item Feedback \textbf{frequenti e recenti}: tanti piccoli \emph{sanity
        check} costano meno di una grande correzione di rotta.
  \item Il feedback è \textbf{un regalo} di chi ha a cuore il nostro
        interesse; fa paura il silenzio, non le critiche.
\end{itemize}

%    % ... altre lezioni zero ...
%    \part{Il corso di Game Design}
%    \renewcommand{\thechapter}{\arabic{chapter}}
%    \setcounter{chapter}{0}
%    % =====================================================================
%  Capitolo 1 — Introduzione al corso e al mestiere del game designer
%  Fonte: trascrizione della lezione GD1 (Game Design, primo anno)
%  Pacchetti richiesti nel preambolo del libro:
%    \usepackage[italian]{babel}, \usepackage[T1]{fontenc},
%    \usepackage[utf8]{inputenc} (solo pdfLaTeX), \usepackage{booktabs}
% =====================================================================

\chapter[Il mestiere del game designer]{Il mestiere del game designer: introduzione al corso}
\label{cap:introduzione-game-design}

Questo primo capitolo fa da soglia. Prima di affrontare metodi, framework e
documenti, è necessario capire \emph{in che modo} si impara il game design,
\emph{con chi} lo si fa, \emph{con quali strumenti} e, soprattutto,
\emph{che cosa fa concretamente} un game designer nel suo lavoro quotidiano.
Il capitolo segue l'impostazione della lezione inaugurale del corso: si parte
dal contesto formativo (il corso come simulazione di lavoro), si passa al
metodo di apprendimento (l'\emph{inquiry-based learning}), si elencano
letture e risorse di riferimento, e si arriva infine a una prima definizione
operativa del ruolo del designer, dei limiti entro cui lavora e delle sue
possibili specializzazioni.

Una premessa terminologica: il game design è una disciplina giovane e quasi
tutta la sua letteratura è scritta in inglese. Per questo motivo molti termini
tecnici verranno lasciati in lingua originale (\emph{pipeline},
\emph{pitching}, \emph{tuning}, \emph{feature team}, \dots): sono i termini che
si useranno nel lavoro, ed è importante imparare a riconoscerli e a usarli
correttamente fin da subito.

% ---------------------------------------------------------------------
\section{Il contesto: un corso come simulazione di lavoro}
\label{sec:simulazione-lavoro}

\subsection{Il principio fondamentale}

Il primo concetto da fissare è che il corso \textbf{si presenta come una
simulazione del lavoro}. Dal momento in cui si entra in aula si viene trattati
come \emph{potenziali professionisti dell'industria}: non soltanto studenti,
ma lavoratori effettivi di uno studio (l'accademia), con una gerarchia, dei
referenti, dei compiti e delle scadenze.

Ne derivano alcune conseguenze pratiche:
\begin{itemize}
  \item il docente non assegna un ``compitino'' da risolvere per accumulare
        punti o presenze: ogni \emph{task} è pensato per far dedicare tempo
        all'apprendimento di qualcosa di concreto;
  \item non consegnare un task non danneggia il docente (che ``vive ancora
        domani e viene pagato ugualmente''), ma lo studente, che perde
        l'occasione di mettere le mani su un argomento e diventare più abile;
  \item docenti e studenti devono trattarsi \textbf{da colleghi}, perché è
        molto probabile che un domani lo diventino davvero.
\end{itemize}

\subsection{Perché l'impegno conta più che mai}

Il lavoro del game designer è descritto come \emph{molto romantico, ma anche
molto duro}, e sempre più duro col passare degli anni. Durante il periodo del
Covid l'industria ha visto grandi investimenti; quei ``rubinetti'' sono stati
chiusi da tempo, il rapporto con i publisher è diventato più complesso e le
opportunità lavorative non sono più quelle di tre o quattro anni fa. Questo
significa:
\begin{itemize}
  \item più competizione per ogni posizione;
  \item aspettative più alte: ci si aspetta che i candidati siano ancora più
        preparati e sicuri delle proprie competenze;
  \item un unico modo per diventarlo: \textbf{lavorare}. Più tempo si dedica
        oggi ad apprendere il mestiere, più facilmente si sarà assumibili
        domani e più agevole sarà il lavoro di tutti i giorni.
\end{itemize}

\subsection{Sbagliare qui per non sbagliare fuori}

Il mantra del corso è: \textbf{quanto più provate qui, anche in modo goffo e
sbagliando, tanto meno sbaglierete là fuori}. In accademia, per quanto la
simulazione sia realistica, nessuno toglie lo stipendio e nessuno licenzia da
un momento all'altro. Nell'industria sì, e un errore può costare molto denaro
all'azienda, ancora di più se l'azienda è la propria.

\subsection{Il valore del fallimento}

Il fallimento, se affrontato con la mentalità corretta, è probabilmente
\emph{la cosa più vantaggiosa} che si possa fare durante la formazione:
\begin{itemize}
  \item un successo (``avete superato questa deadline'') conferma le proprie
        abilità, ma insegna relativamente poco;
  \item dopo $n$ fallimenti si impara moltissimo, e i successi arrivano;
  \item è fondamentale non demoralizzarsi: il feedback negativo non è mai un
        giudizio sulla persona, ma su ciò che è stato consegnato.
\end{itemize}

Il docente riconosce apertamente che dedicare tempo ed energia a un progetto e
vederlo rifiutato, o vederlo non uscire mai, è demotivante: è successo a lui e
lo ha visto accadere anche su progetti molto grandi, con team numerosi. Il
punto però non è ciò che si è ottenuto materialmente, ma \textbf{quanto si è
cresciuti}. Da ogni esperienza ci si portano a casa due cose:
\begin{enumerate}
  \item \textbf{l'esperienza}: quanto si è diventati bravi, veloci,
        praticamente abili nell'applicare un framework, quanto si è cresciuti
        come designer;
  \item \textbf{il rapporto umano}: nel lavoro si ``suda sangue'' fianco a
        fianco con altre persone, e questo crea legami molto forti che
        portano avanti l'industria nel senso più positivo del termine: nascono
        progetti privati, ci si porta da un'azienda all'altra, si formano team
        con basi solide.
\end{enumerate}

Va evitato anche l'eccesso opposto: la cultura del \emph{grinding} si è
appropriata di questi discorsi in modo distorto. Non bisogna diventare robot.

% ---------------------------------------------------------------------
\section{Organizzazione del corso e gerarchia}
\label{sec:gerarchia-corso}

\subsection{La struttura gerarchica}

La struttura dei referenti del corso ricalca volutamente una
\textbf{gerarchia tipica dell'industria}:
\begin{itemize}
  \item al vertice c'è il \textbf{direttore del corso} (Diego Ricciuti), che
        corrisponde al \emph{creative director}: supervisiona la visione del
        progetto (il concetto di ``visione'' verrà approfondito più avanti);
  \item sotto di lui ci sono i \textbf{referenti di dipartimento}, uno per
        disciplina, con un corrispettivo per ciascuna sede (Milano e Roma);
  \item in fondo ci sono gli studenti, che occupano il primo gradino di
        un'azienda: l'\emph{intern} (stagista) o il \emph{junior}, il passo
        che precede il lavoratore ``regolare''.
\end{itemize}

\begin{table}[htbp]
  \centering
  \begin{tabular}{@{}p{0.34\textwidth}p{0.29\textwidth}p{0.29\textwidth}@{}}
    \toprule
    \textbf{Disciplina} & \textbf{Milano} & \textbf{Roma} \\
    \midrule
    Direzione (\emph{creative director}) & \multicolumn{2}{p{0.58\textwidth}}{Diego Ricciuti} \\
    Game design      & Mattia Podini        & Giacomo ``Jack'' Danna \\
    Level design     & Stefano Zocchi       & Andrea Dispoli \\
    Narrative design & \multicolumn{2}{p{0.58\textwidth}}{Daniele Fusetto (lezioni condivise), Giacomo Masi} \\
    Programmazione (Unreal) & Alessandro Pregnolato & Matteo Velardi \\
    \bottomrule
  \end{tabular}
  \caption{Referenti del corso per disciplina e sede (nomi come riportati in lezione).}
  \label{tab:referenti}
\end{table}

Per il primo anno, il referente diretto degli studenti di game design è il
docente di game design: è a lui che ci si rivolge per dubbi e domande. Altre
figure utili sono il direttore didattico (severo, ``giustamente severo'', da
contattare solo per questioni particolarmente importanti), l'amministrazione
(contratti e documentazione) e la segreteria.

\subsection{I canali di comunicazione}

\begin{description}
  \item[Discord] è il canale per le comunicazioni ufficiali del corso e per
        le domande che non richiedono urgenza immediata; può avere qualche
        ritardo.
  \item[E-mail] è il canale per tutto ciò che è formale o ``ancora più
        ufficiale'' e per le comunicazioni rapide e immediate.
  \item[Telegram] esistono residui di vecchi canali precedenti a Discord, da
        evitare.
\end{description}

Le domande che valgono per tutti vanno fatte in classe (anche da un portavoce
a nome del gruppo); quelle singole e personali via Discord o, se formali, via
e-mail. Se il docente non risponde a una domanda c'è un motivo, ad esempio
ragioni legali, e lo dirà esplicitamente.

\subsection{I tre moduli e la settimana tipo}

Le tre discipline principali del primo anno (e non solo) sono:
\begin{itemize}
  \item \textbf{Game Design} (tendenzialmente il lunedì);
  \item \textbf{Level Design} (tendenzialmente il mercoledì);
  \item \textbf{Programmazione} (tendenzialmente il venerdì).
\end{itemize}
Nel primo mese le lezioni sono tutte di game design (lunedì, mercoledì e
venerdì). Sono previste anche lezioni di narrativa, ma il piano formativo si
aggiorna ogni anno: le ore di narrativa sono state ridotte a favore di quelle
di programmazione. Le poche lezioni di narrativa daranno quindi dei
\emph{cenni}, non la preparazione di un narrative designer. Ciò che serve
davvero è imparare \textbf{come comunicano gli altri dipartimenti} e come
lavorare con persone che fanno un lavoro diverso ma complementare al proprio.

\subsection{Le due anime delle lezioni di game design}

\begin{description}
  \item[Design methods (metodologie)] lezioni prevalentemente teoriche, ma con
        una buona componente pratica, in cui si studiano metodologie di design
        e i cosiddetti \emph{framework} dell'industria.
  \item[Documentazione e pitching] il cuore del lavoro quotidiano del
        designer.
\end{description}

\paragraph{La documentazione.}
Il docente si definisce ``il più grande nerd della documentazione'' che si
possa incontrare. È una parte dura, perché a pochi piace documentare, ma il
lavoro del designer \emph{è principalmente questo}: scrivere documenti e, da
quei documenti, passare alla pratica e all'implementazione. Chi non sa
documentare bene parte male. Quando non si è presenti, il documento è il
proprio \textbf{biglietto da visita}: qualcun altro lo leggerà, quindi va
scritto nel modo più chiaro e comprensibile possibile.

Un esempio concreto del perché la documentazione conta: CD Projekt Red ha
sviluppato un proprio engine proprietario, ma buona parte della sua
documentazione è stata ``buttata all'acqua''. Chi deve usarlo o lo sa già
usare, o ha competenze di un altro engine da trasferire, oppure è in
difficoltà. Si documenta \textbf{perché altri possano leggere}.

\paragraph{Il pitching.}
Il \emph{pitching} è l'esposizione, la presentazione orale. La parte creativa
del team presenta spesso:
\begin{itemize}
  \item le proprie idee ai colleghi;
  \item il gioco a convention ed eventi;
  \item il pitch del gioco agli investitori, per ottenere i fondi necessari a
        realizzarlo.
\end{itemize}
Saper parlare e comunicare chiaramente le idee è fondamentale. In sintesi:
\emph{il game designer parla}, e bisogna farsene una ragione.

\subsection{Altre componenti del corso}

\begin{description}
  \item[Gioco critico e playtest] la classe viene divisa in due gruppi: uno
        gioca, l'altro osserva. Chi osserva non applaude ``come le
        scimmiette'': annota, segna e analizza il \emph{playtest}, esattamente
        come si fa nell'industria. Il playtest fa parte delle mansioni del
        designer, ma esiste anche come specializzazione autonoma.
  \item[Pipeline] termine che indica tutto ciò che riguarda i processi
        formalizzati dell'industria; di solito si intende la \emph{pipeline di
        produzione}: come si divide il team, qual è la sua gerarchia, come il
        \emph{producer} gestisce le deadline. In breve: \emph{come} e
        \emph{quando} si fanno le cose. Al primo anno è meno rilevante (si
        dovrà ``nuotare'' un po' da soli); dal secondo anno si vedranno le
        metodologie effettivamente applicate ai progetti digitali. Si
        scoprirà quanto sia difficile lavorare a un progetto digitale, e
        ancor più fisico, senza un piano di sviluppo comprovato, come quelli
        usati per qualsiasi software in qualsiasi parte del mondo.
  \item[Engine] il motore di riferimento è \textbf{Unreal Engine}
        (sezione~\ref{sec:engine}).
  \item[Programmazione] principalmente tramite \emph{Blueprint} in Unreal,
        più eventuali sistemi di versionamento; il corso parte da zero.
  \item[Progetti pratici] probabilmente la cosa più importante. Al primo
        anno i progetti di gruppo restano circoscritti al corso di design,
        perché mancano ancora alcune nozioni (come comunicare, affrontare
        problemi, dare e ricevere feedback). In seguito si allargheranno a
        progetti digitali, molto più difficili di quelli fisici, con membri
        di altri team: un rapporto non unilaterale, ma ``covalente''.
\end{description}

% ---------------------------------------------------------------------
\section{Il metodo: l'inquiry-based learning}
\label{sec:inquiry}

\subsection{Che cos'è}

Il corso è costruito secondo l'\textbf{inquiry-based learning}
(apprendimento basato sull'indagine), un approccio pedagogico nato intorno
agli anni Sessanta, che si richiama a pensatori come Dewey e Piaget. In
generale, le discipline da cui il design attinge sono quasi sempre la
sociologia, la psicologia e, talvolta, l'economia.

Il funzionamento è un ciclo:
\begin{enumerate}
  \item il docente fornisce metodologie e framework dell'industria durante
        la lezione;
  \item lo studente, fuori dalle ore di lezione, lavora sui \emph{task}:
        mette a terra le nozioni, le maneggia, le assorbe meglio;
  \item da questo lavoro nascono nuovi dubbi;
  \item i dubbi tornano in aula: lo studente chiede, il docente risponde
        (ammesso che dare la risposta sia utile, e ammesso che una risposta
        esista).
\end{enumerate}

\subsection{Non sempre esiste una risposta giusta}

Non tutto ciò che si studia ha ``la risposta migliore possibile in assoluto
nel vuoto''. Esistono metodologie, limiti e framework che aiutano il lavoro,
ma non sono fissi: cambiano da azienda ad azienda, da lavoro a lavoro, da
progetto a progetto, e a volte vengono stravolti. Sta alla
\textbf{sensibilità del designer} capire quale strumento o framework usare in
un dato momento, e questa sensibilità si impara solo con l'esperienza. Nel
lavoro nessuno dirà ``usa questa cosa'' o ``è meglio quest'altra'': sarà il
designer a scegliere.

Anche l'esperienza del docente ha confini: sei-sette anni di lavoro, un modo
personale di approcciare il design e un certo insieme di progetti (ad esempio
nessun MMO). Oltre quei confini si cercano risposte nella ricerca, studiando
come aziende con esperienza specifica (ad esempio Riot) hanno affrontato
problemi simili e hanno condiviso le loro soluzioni, sempre valutando quali
metodi siano adatti alla dimensione del team, al tipo di progetto e al tempo
disponibile.

\subsection{Sensibilità, non memoria}

Non interessa imparare a memoria un framework o la teoria di un autore. Certe
terminologie vanno conosciute, ma per ragioni di comunicazione e chiarezza.
Ciò che conta è sviluppare la \textbf{capacità di farsi domande}, di farle
agli altri, di fare ricerca e di trovare una propria risposta. È esattamente
ciò che si fa nel lavoro: davanti a qualcosa che non si sa fare, l'abilità
del designer sta nel cercare come qualcun altro ha risolto un problema simile.

\subsection{Le lezioni come corsi di aggiornamento}

Le lezioni frontali simulano i \emph{corsi di aggiornamento} aziendali: le
sessioni che un lead di design tiene al suo team di design (o il lead di
level design al suo team, e così via). Alla fine dell'aggiornamento c'è quasi
sempre un lavoro da fare, un obiettivo da portare a terra:
\textbf{l'obiettivo della settimana}. Si lavora di settimana in settimana, e
il perché verrà chiarito in seguito.

Struttura tipica di una lezione:
\begin{enumerate}
  \item breve sessione di domande e risposte sulla lezione precedente;
  \item analisi di eventuali esercizi particolarmente critici;
  \item lezione vera e propria;
  \item discussione di gruppo, se c'è tempo;
  \item assegnazione del task della settimana.
\end{enumerate}
Il docente tende a ``parlare e straparlare'', ad aprire parentesi e a saltare
da un argomento all'altro: l'unico modo per fermarlo è che siano gli studenti
a interromperlo con domande.

\subsection{La gerarchia della competenza}

L'inquiry-based learning può essere rappresentato come una scala di
competenza:
\begin{enumerate}
  \item la sola lezione frontale porta a \textbf{conoscenze incomplete},
        perché manca la parte pratica e di ricerca;
  \item il lavoro a casa (task, ricerca) porta a una competenza \textbf{più
        completa}, non sbagliata in sé, ma ancora parziale e complementare ad
        altre;
  \item la \textbf{condivisione} dei task e dei feedback con gli altri porta
        idealmente alla competenza più completa possibile;
  \item il picco della competenza è la sua \textbf{applicazione nei
        progetti} (fisici o digitali, di gruppo): ``con le cose che avete
        imparato e ricercato, riuscite a fare un gioco?''. Quando la risposta
        diventa positiva, o comunque si inizia a navigare nei processi di
        design per realizzarne uno, quella è la vittoria.
\end{enumerate}
Non si parte subito creando giochi: si parte dall'analisi, dal design e dalle
regole generali di design, prima ``nel vuoto'', poi contestualizzando i
framework, infine unendo i puntini nei progetti.

Rimarranno comunque delle lacune, per tutti, docente compreso, perché il game
design \textbf{non è una scienza esatta}, anzi, a rigore non è affatto una
scienza. Prende in prestito metodologie dalle discipline scientifiche, ma
contiene una componente imprevedibile e incalcolabile: \textbf{la
produzione}. Si possono applicare perfettamente le regole di un framework, e
quel framework può comunque non essere adatto al gioco, al proprio modo di
sviluppare o al tempo di produzione disponibile (alcuni framework richiedono
una ricerca e un'attenzione che il tempo non concede).

C'è poi la componente creativa, la parte ``romantica'' del lavoro: esiste, ma
è \textbf{metà del lavoro}. L'altra metà è ricerca, metodo scientifico al
100\%. Dagli stessi dati $n$ designer diversi, con sensibilità diverse,
possono ottenere $n$ risposte diverse. Non tutte saranno valide: vanno
applicate, messe a terra, e questa scrematura fa parte del lavoro.

\subsection{I livelli dell'apprendimento}

Un secondo modello, a gradini, descrive i livelli di apprendimento (è la
cosiddetta tassonomia SOLO):
\begin{description}
  \item[Pre-strutturale] il docente dà il compito e lo studente lo risolve,
        possibilmente il prima possibile. È esattamente ciò che \emph{non} si
        vuole.
  \item[Uni-strutturale] si affronta una sola disciplina: è il primo passo
        necessario, ed è il motivo per cui nel primo mese ci sono solo lezioni
        di game design.
  \item[Multi-strutturale] entrano in gioco altri dipartimenti: level design,
        game design, programmazione e così via.
  \item[Relazionale] si capisce che alcune cose valgono per più
        dipartimenti: alcuni framework di narrativa derivano dal design o
        possono essere usati per fare design; le regole di documentazione si
        applicano anche ad altri dipartimenti; le regole di level design si
        mescolano con quelle di game design. Dopotutto, per fare un gioco
        servono molti dipartimenti.
  \item[Astratto esteso] il punto finale: muoversi autonomamente e capire
        cosa usare e quando, sulla base della propria esperienza. L'unico modo
        per arrivarci è fare pratica e fare progetti.
\end{description}

% ---------------------------------------------------------------------
\section{Task, feedback e regole del corso}
\label{sec:task-feedback}

\subsection{I task come opportunità}

I task non sono strettamente obbligatori (esiste un numero minimo di
consegne), ma sono \textbf{opportunità}. Quanto e come ci si lavora dipende
dallo studente. Un task fatto in cinque minuti, o in cui si è stati ``furbi'',
non verrà letto né riceverà feedback.

Domande da porsi davanti a ogni task:
\begin{enumerate}
  \item \textbf{Perché} mi è stato assegnato? Qual è l'obiettivo?
  \item Come ci lavoro per ottenere un risultato \textbf{che abbia senso}? Se
        la ricerca e i dati indicano una direzione, è inutile andare in
        un'altra.
\end{enumerate}

\subsection{Quanto lavorare: un calcolo indicativo}

Se il corso è una simulazione di lavoro, vale il carico di un lavoro normale:
\begin{itemize}
  \item un lavoro a tempo pieno comprende da 30 a 45 ore settimanali; si
        considera una media di \textbf{40 ore};
  \item si sottraggono circa 9 ore di lezione: restano $\approx 31$, diciamo
        \textbf{30 ore};
  \item con circa 3 task a settimana: $30 / 3 = $ \textbf{circa 10 ore per
        task}.
\end{itemize}
Il numero può variare, ma rende l'idea: dedicare un'ora a un task e poi fare
altro non fa guadagnare nulla a nessuno. Lo studente ``butta soldi'', il
docente risparmia tempo leggendo il task in un minuto (o non leggendolo).
Dedicando invece attenzione, ricerca, cura nella formattazione e
sperimentazione, e facendolo con cognizione di causa, si impara.

\subsection{Niente voti, solo feedback}

Ad eccezione del corso di programmazione, non ci sono voti (né ``ottimo'', né
10, né 100). Ci sono \textbf{feedback}:
\begin{itemize}
  \item in game design, generalmente sulla documentazione;
  \item in level design, sui progetti pratici.
\end{itemize}
Il docente indica cosa va bene e cosa va modificato, concentrandosi in
particolare su ciò che va cambiato. Il feedback \textbf{non è mai un giudizio
sulla persona}, ma su ciò che è stato consegnato nel concreto.

\subsection{La condivisione dei task}

Tutti i task consegnati vengono \textbf{condivisi} con la classe, e possono
essere letti solo dopo aver ricevuto i feedback del docente. Si cresce così
non solo con i propri feedback, ma anche con quelli dati agli altri. Regola
importante: ciò che il docente dice a uno studente \textbf{vale per tutti}. Se
a una persona viene detto ``questa cosa non va mai più fatta'', da quel
momento vale per l'intera classe.

\subsection{Non copiare}

Copiare i task degli altri è una delle cose peggiori che si possano fare, non
per una questione morale, ma perché annulla il senso del percorso. Si può
dedicare a ogni task il tempo e l'attenzione che si ritengono opportuni, ed è
normale che col tempo ci si specializzi (preferendo, ad esempio, il level
design o la programmazione tecnica) e si investa meno dove si sente meno
necessità. In questi casi l'alternativa corretta è:
\begin{itemize}
  \item non consegnare;
  \item studiare ciò che hanno fatto gli altri;
  \item eventualmente riconsegnare in seguito.
\end{itemize}
Ci saranno momenti in cui si verrà sommersi dalle cose da fare: bisogna
imparare a scegliere.

\subsection{Il ``non leggo''}

Il docente ha piena facoltà di non dare alcun feedback a un task (il
cosiddetto \textbf{``non leggo''}) se:
\begin{itemize}
  \item non si è lavorato con attenzione;
  \item non si rispettano le indicazioni date;
  \item è stata violata una regola estremamente importante.
\end{itemize}
Un task ``non letto'' equivale a un task non consegnato e non entra nel
conteggio. Attenzione: un task \emph{perfetto} e un task che \emph{fallisce}
nel comprendere o applicare un framework sono \textbf{ugualmente validi},
purché ci si sia impegnati e si sia dedicato tempo, pazienza e ``sudore''.

\paragraph{Riconsegna.} Un task non letto, in ritardo o che si vuole
migliorare può essere riconsegnato e verrà conteggiato, a patto che:
\begin{itemize}
  \item i problemi segnalati siano stati risolti;
  \item sia stato fatto ulteriore lavoro, integrando anche tutto ciò che nel
        frattempo è stato detto agli altri colleghi nei loro feedback.
\end{itemize}
Il task va sviluppato con le competenze che si hanno nel momento della
riconsegna, non con quelle del momento della prima consegna.

\subsection{Requisiti per superare l'anno}

\begin{itemize}
  \item consegnare almeno l'\textbf{80\% dei task} assegnati (soglia
        malleabile);
  \item partecipare almeno all'\textbf{80\% delle lezioni}, in presenza o
        da remoto, preferibilmente in presenza; i ritardi vengono valutati
        caso per caso;
  \item comportarsi da persone corrette e rispettose (è stata firmata una
        carta dello studente): non urlare, non alzare le mani, non
        vandalizzare, non compiere atti illegali.
\end{itemize}
Impegnarsi di più porta inoltre a essere selezionati per varie opportunità.

\subsection{Scadenze}

\begin{itemize}
  \item la scadenza di un task è \textbf{sempre prima dell'inizio della
        lezione} successiva indicata;
  \item se la classe, come team, mantiene comportamenti scorretti, il docente
        può \textbf{anticipare le deadline} per tutti.
\end{itemize}

\subsection{Accessibilità e DSA}

Gli studenti con disturbi specifici dell'apprendimento o neurodivergenze
(dislessia, ADHD, autismo, ecc.) sono invitati a segnalarlo. Quando si
affronteranno le regole di formattazione e documentazione (errori
grammaticali, mancanze, incoerenze) se ne terrà conto e si potranno trovare
dei compromessi, pur ricordando che nel mondo del lavoro nessuno farà
``sconti''. La documentazione può essere inviata all'amministrazione o alla
segreteria, che la comunicherà a tutti i docenti, ma il consiglio è parlarne
anche direttamente con ogni singolo docente.

% ---------------------------------------------------------------------
\section{Le competenze trasversali}
\label{sec:soft-skill}

\subsection{I giochi si fanno in gruppo}

Si possono fare giochi da soli, ed è un'impresa ammirevole; ma lavorando da
soli \textbf{tutto ciò che si guadagna in libertà si perde in produttività}.

\paragraph{Il caso Stardew Valley.} \emph{Stardew Valley} è stato realizzato
da una sola persona, Eric Barone, in anni di sviluppo e di sacrifici: faceva
un altro lavoro per mantenersi (in un cinema) e di notte programmava, faceva
l'arte e componeva la musica. Il risultato è stato straordinario e lo ha reso
ricco, ma si tratta di un caso estremamente raro. Quanti giochi simili, nati
nelle stesse condizioni, sono falliti? Molti di più.

\paragraph{Il rischio d'impresa.} Lavorare da soli non è sbagliato di per sé
e ha aspetti molto positivi: si ha il privilegio di creare esattamente ciò che
si vuole. Ma ci si fa carico anche di tutto il \textbf{rischio d'impresa}: se
si lavora a un gioco per cinque anni e poi il gioco va male, resta il
vantaggio di averlo pubblicato, ma il dopo è sperare di vivere di quella
release o di finanziare con essa il gioco successivo. Le grandi aziende hanno
gli stessi problemi, ma li distribuiscono su più persone e più responsabilità.

L'obiettivo di un professionista è \textbf{saper lavorare in diverse
condizioni}, e la più comune è quella del dipendente (o del freelance /
libero professionista con partita IVA): si viene assunti per le proprie
capacità da un'azienda, piccola o grande che sia, e si lavora con altre
persone.

\subsection{Essere persone con cui si vuole lavorare}

Se le persone non vogliono lavorare con voi, quello è \textbf{il primo motivo
di licenziamento}. E se non si viene licenziati, il clima nell'azienda
diventa terribile: la produttività crolla, la motivazione crolla. Le
\emph{soft skill}, cioè l'insieme di abilità che rendono una persona
piacevole con cui lavorare, vanno coltivate e diventano sempre più importanti
col passare degli anni (relativamente meno ora, molto di più quando i team
saranno più grandi). Anche chi è un ``lupo solitario'' deve imparare a:
\begin{itemize}
  \item lavorare con altre persone;
  \item comunicare con altre persone;
  \item trovare un punto in comune da cui partire per fare le cose.
\end{itemize}

\subsection{Indipendenza e proattività}

Sono le altre due parole chiave del corso:
\begin{description}
  \item[Indipendenza] spesso si verrà lasciati da soli, non per cattiveria ma
        per sfida: quando si cerca una risposta da soli, ciò che conta non è
        tanto la risposta quanto \textbf{il processo}. Una risposta data dal
        docente non vale niente.
  \item[Proattività] chiedere, domandare, agire da soli senza aspettare
        direttive, perché spesso quelle direttive non ci sono, a maggior
        ragione nei team piccoli.
\end{description}
Tutto questo si concretizza nel parlare, nel presentare e nel fare
\emph{post-mortem} (l'analisi a posteriori di un progetto concluso, che verrà
approfondita più avanti), e in generale nell'attenzione alla crescita
personale.

\subsection{Fare domande}

Il corso insiste sulle domande: alzare la mano, intervenire anche ``a
bomba'', fermare il docente ogni volta che qualcosa non è chiaro. Se una
persona non conosce qualcosa, molto probabilmente c'è almeno un'altra persona
che non la conosce. Il docente può dare per scontata la conoscenza di termini
che non lo sono, ed è del tutto normale, soprattutto all'inizio.

% ---------------------------------------------------------------------
\section{Strumenti: engine, Unreal e materiali}
\label{sec:engine}

\subsection{Unreal Engine e gli altri motori}

Il motore di riferimento del corso, usato in level design, è \textbf{Unreal
Engine}: ``un gran mattone'', ma che dà a un designer molte possibilità. Il
vantaggio principale è che permette a un designer non pratico di
programmazione di \textbf{prototipare tramite i Blueprint} (il sistema di
scripting visuale di Unreal). Questo non impedisce di passare a Unity o ad
altri engine, e viceversa: il docente stesso ha studiato e lavorato a lungo
con Unreal, ma la sua ultima esperienza professionale (come level designer) è
stata su Unity, un engine che non conosceva, a dimostrazione che il salto da
un engine all'altro è possibile.

Panorama dei motori citati:
\begin{itemize}
  \item \textbf{Unreal Engine} e \textbf{Unity}: i due più diffusi, con
        modelli di sorgente ``più o meno'' aperto;
  \item \textbf{Godot}: engine open source;
  \item \textbf{CryEngine}: proprietario di Crytek;
  \item \textbf{Snowdrop}: engine di Ubisoft;
  \item Rockstar usa un proprio engine proprietario;
  \item alcune aziende usano versioni \textbf{personalizzate} di engine
        esistenti: ad esempio Milestone, nei giochi di corse, usa un Unreal
        customizzato;
  \item CD Projekt Red ha un engine proprietario (con i problemi di
        documentazione già citati).
\end{itemize}
In sintesi: ogni azienda può avere un engine proprietario oppure una versione
modificata di un engine commerciale, ma si parte sempre da una base.

\subsection{Programmazione: cosa ci si aspetta da un designer}

Non è richiesto alcun bagaglio pregresso di programmazione: il corso parte da
zero. Gli studenti sono stati selezionati (il corso è a numero chiuso) per
competenze analitiche e critiche sopra la media. Anche nel mercato del lavoro,
per una posizione di game designer generalmente \textbf{non ci si aspetta
competenza di scripting} (C++, Java, C\#, ecc.), a meno di voler fare il
\emph{technical designer}, e anche lì varia da azienda ad azienda.

\subsection{Versione e preparazione}

Indicazioni pratiche emerse in aula:
\begin{itemize}
  \item scaricare Unreal Engine; la versione consigliata per il corso è
        stata indicata nella \textbf{5.6} (con conferma successiva). La
        stabilità dipende da cosa si deve fare, da quali tool e plugin si
        usano: i problemi nascono spesso dai plugin. Nel passaggio da Unreal 4
        a Unreal 5, ad esempio, il sistema \emph{Landscape} era inizialmente
        inutilizzabile, mentre in Unreal 4 funzionava molto bene;
  \item è consigliato installare tramite \textbf{Epic Games Launcher}, più
        comodo: permette anche di riscattare giochi gratuiti periodici e di
        accedere ai pacchetti di \textbf{Fab} (il marketplace di asset
        acquisito da Epic), che ogni mese offre alcuni asset gratuiti utili
        per esercitarsi nel level design (ad esempio per fare \emph{speed
        level design});
  \item molti pacchetti di UI o di sistemi di interazione sono utili per
        prototipare al volo, ma spesso inutilmente complicati;
  \item un plugin di \textbf{materiali di prototipazione a griglia} (adattivi,
        con metriche, allineabili alla griglia, con materiali già pronti e
        colori modificabili, dal costo di circa 5 euro) è indicato come
        ``salvavita'' per il level designer; se necessario, i pacchetti
        indispensabili vengono forniti;
  \item portare \textbf{sempre il portatile} alle lezioni di level design
        (da novembre, a partire dall'interfaccia dell'engine); alle lezioni
        di game design basta che lo porti almeno metà della classe, per i
        lavori a coppie. Serve comunque qualcosa per prendere appunti
        (``tavolette d'argilla e cannucce, se volete fare i babilonesi'',
        oppure carta e penna).
\end{itemize}
Chi ha un portatile poco potente dovrebbe comunque provare a installare
l'engine e aprire un template per vedere come gira: all'inizio non si faranno
cose pesanti. Nelle prime lezioni di level design è importante mettere subito
le mani sull'interfaccia: ascoltare dove si trova uno strumento (ad esempio lo
\emph{snapping}) senza provarlo porta a dimenticarlo.

% ---------------------------------------------------------------------
\section{Letture e risorse di riferimento}
\label{sec:risorse}

\subsection{I libri}

\begin{description}
  \item[\emph{The Art of Game Design: A Book of Lenses} (Jesse Schell)] il
        testo \textbf{davvero importante}. Alcune lezioni finali dell'anno si
        baseranno sulle sue terminologie. È utile perché apre la mente sulla
        possibilità di applicare metodologie diverse a giochi diversi, ma
        soprattutto per la \textbf{terminologia} comune. Esiste anche in
        audiolibro (circa 25 ore).
  \item[\emph{Challenges for Game Designers} (Brenda Romero e Ian Schreiber)]
        libro \emph{entry level} con moltissimi esercizi, molto caro al
        direttore del corso.
  \item[\emph{21st Century Game Design} (Chris Bateman)] utile come
        retrospettiva: è un libro visionario che ha previsto metodi e
        soprattutto \emph{archetipizzazioni} dei giocatori tuttora valide nel
        design e nel marketing. Servirà per una delle lezioni più importanti
        dell'anno. È utile leggerlo per vedere quanto aveva previsto, ma è
        utile ``fino a un certo punto''.
  \item[\emph{An Architectural Approach to Level Design} (Christopher
        Totten)] non strettamente necessario, ma di lettura semplice. Ha due
        anime: parte dall'architettura (babilonese, greca, ecc.) e la trasla
        su esempi di level design (come \emph{Doom} e \emph{Team Fortress}).
        Un tempo ne esisteva un riassunto su Gamasutra (oggi
        \emph{Game Developer}).
  \item[\emph{A Theory of Fun for Game Design} (Raph Koster)] utile fino a un
        certo punto: contiene diverse idee superate, che lo stesso Koster ha
        affrontato in vari talk reperibili nel GDC Vault. Verrà ripreso nelle
        prime lezioni, poi messo da parte.
\end{description}

\subsection{Canali e conferenze}

\begin{description}
  \item[GDC (Game Developers Conference)] il suo canale raccoglie talk e
        panel estremamente importanti: è ``uno strumento misterioso che ci
        servirà più tardi''. \textbf{Attenzione}: per molti talk mancano ancora
        le competenze per interpretarli, e alcuni potrebbero fuorviare o
        mettere in testa idee strane. La GDC si concentra spesso su
        produttività e mobile.
  \item[Il \emph{Dragon Speech} di Chris Crawford] il talk del fondatore della
        GDC: appartiene a un passato lontano, ma è molto motivante. Da vedere
        subito.
  \item[Game Maker's Toolkit] canale particolarmente utile per nozioni
        pratiche, con video brevi e alcune interviste.
  \item[Nordic Game] altra conferenza importante, con talk e panel di esperti
        del settore.
\end{description}

\subsection{Associazioni e premi: come leggerli}

Le associazioni di settore (ad esempio IIDEA in Italia, o la IGDA) non sono
il male, anzi: sono un canale importante per affacciarsi all'estero o per
accedere a fondi. Va però ricordato che \textbf{tutelano i propri associati}:
non esiste solo il loro palinsesto, e bisogna valutare quanto spazio occupano
nel panorama della visibilità.

Lo stesso vale per i \textbf{premi} (Game Awards, ``Best Mobile Game'',
``Games for Change'', ecc.): al designer \emph{non interessano}. Che un gioco
abbia vinto ``più di 90 premi'' non dice nulla di utile. Servono a dare
visibilità e lustro ai giochi ed è giusto che esistano, perché ampliano il
panorama e il numero di giocatori. Ma ciò che interessa al designer sono i
\textbf{dati}:
\begin{itemize}
  \item come sono state implementate certe cose;
  \item come è stato fatto il marketing e dove il gioco è stato presentato;
  \item non quante wishlist ha fatto, ma \emph{come} le ha fatte;
  \item in che periodo storico è uscito e quali erano i competitor in quel
        momento.
\end{itemize}

% ---------------------------------------------------------------------
\section{Che cosa fa un game designer}
\label{sec:cosa-fa}

Per arrivare a una prima definizione, la classe ha condotto un brainstorming:
ognuno doveva indicare una cosa che, secondo lui, fa il game designer (senza
ripetere quelle già dette). Ne è emersa la lista seguente, commentata punto
per punto.

\subsection{Progettare meccaniche e struttura del gioco}

Il ruolo più generico è quello del \textbf{core mechanic designer}, detto
anche \emph{gameplay designer} e, un tempo, \emph{black box designer}: si
occupa delle \textbf{meccaniche di gioco}. Si va dal piccolo al grande (in
funzione anche della dimensione dell'azienda):
\begin{itemize}
  \item nel \textbf{piccolo} ci sono le singole meccaniche;
  \item nel \textbf{grande} ci sono i \textbf{sistemi}. Nei giochi
        \emph{sistemici}, con molti sistemi interconnessi (ad esempio
        \emph{GTA}), esistono specializzazioni dedicate a bilanciare e
        supervisionare i sistemi: è il \emph{system design}.
\end{itemize}

\paragraph{Quando un gioco è ``sistemico''?}
Non dipende solo da come le meccaniche interagiscono con il giocatore (le
\emph{dinamiche}, secondo i framework che passano dalle meccaniche alle
dinamiche, che verranno studiati più avanti), ma anche da \emph{quante}
meccaniche ci sono e da quanto interagiscono tra loro. Non diremmo che
\emph{Super Mario} sia un gioco sistemico, eppure contiene dei sistemi: la
linea è sottile. Quando i sistemi aumentano molto, sono complessi e si
influenzano a vicenda, il gioco diventa sistemico. Esempio: in \emph{The
Legend of Zelda: Breath of the Wild} il vento modifica quanto il personaggio
sale o scende; il fuoco può incendiare l'erba; l'erba prende fuoco se colpita
da una freccia infuocata; e così via. Nel mercato moderno i giochi sistemici
sono spesso, ma non sempre, legati a mappe \emph{open world}.

\subsection{Le 3C: camera, character, control}

La base del gioco, da cui generalmente si inizia a lavorare (anche nei
prototipi), è la ``sacra trinità videoludica'' delle \textbf{3C}:
\begin{description}
  \item[Camera] che tipo di inquadratura ha il gioco: prima persona, terza
        persona, isometrica, ecc.
  \item[Character] quanto è grande il personaggio, che tipo di personaggio è
        e, di conseguenza, quali azioni può compiere.
  \item[Control] i controlli: che cosa il personaggio può fare, quali sono i
        \emph{verbi} di gioco.
\end{description}
Le 3C sono definite dal designer, ed esiste anche una specializzazione per
questo lavoro. Esistono varianti: quella del direttore del corso è
\textbf{4C + J}, cioè \emph{camera, character, control, conflict} (conflitto)
e \emph{juice} (il ``succo'', il feedback che rende piacevole l'interazione).
Verranno approfondite più avanti.

\subsection{Il personaggio}

Il designer definisce il personaggio dal punto di vista delle meccaniche
(dimensioni, azioni, capacità), all'interno delle 3C. L'aspetto estetico è
invece competenza di altri dipartimenti.

\subsection{Documentare le informazioni per gli altri dipartimenti}

Assolutamente vero: la documentazione è una delle prime attività del
designer. Più l'azienda è grande, più questo lavoro diventa prominente; più è
piccola, meno è rilevante. In \textbf{pre-produzione}, però, è sempre
fondamentale.

\subsection{Mediatore e leader ``al servizio''}

Il designer è un \textbf{mediatore} nei conflitti interni al team, una sorta
di collante che si relaziona con dipartimenti diversi (non sempre gli stessi,
a seconda della fase di produzione). A volte fa anche da leader, a seconda di
come sono gestite la pipeline e la gerarchia.

Questa leadership è difficile da definire perché \textbf{non è verticale}: è
una leadership \emph{al servizio} (\emph{servant leadership}). Il designer
``è asservito al programmatore'': il suo compito non è arrivare con un'idea e
scaricarla addosso agli altri, né portare la soluzione e lasciare il problema
a qualcun altro, ma \textbf{seguire il processo a basso livello}, sporcandosi
le mani.

\paragraph{I feature team.} Nei team medio-piccoli o medio-grandi possono
esistere sotto-team (\emph{feature team}) dedicati a un obiettivo specifico:
ad esempio un designer incaricato del character o della camera, affiancato da
uno o due programmatori e da un artista. In questi casi il designer
supervisiona il lavoro e gestisce la produzione del sotto-team.

\subsection{Gestire l'avanzamento del prodotto}

Esiste una specializzazione che si occupa solo di questo, il
\textbf{producer}, a un livello gerarchico completamente diverso (e che
potrebbe essere una strada percorribile al termine del percorso di studi).
All'interno di un feature team anche il designer può farlo, ma non è detto.
Può certamente \emph{dare la direzione}: nella sua documentazione può
indicare che per arrivare a uno stato giocabile del prototipo serve un certo
lavoro di arte e di programmazione, per lo stato successivo un altro, e così
via. In ambito produttivo esistono metodologie strutturate come
\textbf{Agile}, che verranno viste in seguito.

\subsection{Presentare le idee}

Il pitching (sezione~\ref{sec:gerarchia-corso}) è richiesto al designer sia
in pre-produzione sia in produzione, in misura che dipende dall'azienda. Ad
esempio Ubisoft lascia ai designer, e in realtà a qualsiasi membro, la
possibilità di \emph{pitchare} un'idea prima di una nuova produzione, perché
possa iniziare a prendere vita.

\subsection{Bozze e prototipi}

Le bozze vanno bene, ma devono essere \textbf{chiare}. C'è un'eterna lotta tra
documentazione e gioco giocabile: se si fa solo documentazione il gioco non
avanza; se si fa solo il gioco manca la documentazione e si perde la visione.
Sono due cose \textbf{complementari}.

Il modo per comunicare un'idea in forma grezza ma giocabile è il
\textbf{prototipo}. Un prototipo può essere qualsiasi cosa: di carta, in
Unreal con Blueprint ``brutti''. \emph{Non deve essere bello, deve essere
giocabile}.

\paragraph{Esempio: il prototipo di carta di Mario + Rabbids.} Durante lo
sviluppo di \emph{Mario + Rabbids}, i programmatori di Ubisoft erano impegnati
su altri progetti (come \emph{Ghost Recon}). I designer realizzarono allora
un \textbf{gioco da tavolo} che riproduceva le meccaniche del gioco, quelle
di uno strategico a turni sul modello di \emph{XCOM}. E funzionava.
L'obiettivo era \textbf{validare un'idea}, validare un gameplay.

Altri modi per rendere chiara un'idea sono i \emph{mock-up} nella
documentazione e le immagini modificate o i \emph{mash-up} di qualsiasi
natura. Lato documentazione l'importante è che sia chiaro; lato prototipo,
che sia giocabile.

\subsection{Fare ricerca per l'idea, e poi smontarla}

Il designer ha una componente creativa importante e fa ricerca per arrivare a
un'idea. Ma ciò che deve poi fare necessariamente è \textbf{smontarla}. Il
suo lavoro è scientifico, esattamente come il \textbf{metodo scientifico}, il
cui obiettivo è \emph{confutare un'ipotesi}: si arriva a un'ipotesi tramite
sperimentazione e poi si cerca $n$ volte di confutarla; se resiste, la si
assume come valida \emph{in un certo contesto} (le regole della fisica
classica, ad esempio, non valgono nella meccanica quantistica).

Allo stesso modo il designer arriva a un'idea (tramite ricerca o chiedendo
agli altri: le idee le hanno tutti, non sono un privilegio del designer),
cerca di smontarla e poi cerca di \textbf{metterla a terra}. Se l'idea non è
applicabile, non è implementabile, finché non si trova un modo per
realizzarla va scartata. Può essere l'idea più bella dell'universo, ma
\textbf{se non può essere messa a terra, non esiste}.

\paragraph{L'originalità.} Serve un'idea originale? Dipende dall'obiettivo e
dalla visione del progetto:
\begin{itemize}
  \item se l'obiettivo è vendere molto, un'idea originale non serve
        necessariamente;
  \item se si vuole creare un nuovo filone, serve un'idea, se non originale,
        almeno interessante: può essere un discorso di marketing o una
        meccanica innovativa, con il rischio di non vendere.
\end{itemize}
Il ruolo del \emph{director} e il concetto di \emph{visione} verranno
approfonditi in seguito.

\subsection{Analizzare i competitor}

Arrivati a un'idea (``vogliamo un certo tipo di combattimento''), si va a
vedere se qualcuno l'ha già fatta e \emph{come}. Le meccaniche non sono
privilegio di nessuno, salvo quando vengono brevettate, come il
\textbf{Nemesis System} di \emph{La Terra di Mezzo: L'Ombra di Mordor}.

\paragraph{Esempio: il combattimento da ridimensionare.} In un progetto
(presso Stormind Games) il team si trovò con un sistema di combattimento
troppo costoso: mancavano le risorse e gli animatori per realizzarlo. Poco
output lato gameplay significava non poter fare qualcosa di complesso e
profondo come in \emph{Elden Ring}; al contrario, il team aveva molti
artisti e poteva puntare su qualcosa di più artistico. Si iniziò così a
orbitare intorno al corpo a corpo di \emph{The Last of Us} (parte I e II).
Sul canale della GDC esistono diversi talk che raccontano cosa è stato fatto
e come, con successi, insuccessi e costi (ad esempio del \emph{motion
capture}): da quei talk si ricavarono moltissime informazioni utili.

I competitor possono essere:
\begin{itemize}
  \item \textbf{diretti}: un gioco molto simile al proprio;
  \item \textbf{di mercato}: un gioco che domina il mercato in cui ci si vuole
        affacciare;
  \item \textbf{indiretti}: un gioco con un sistema simile al proprio, ma
        tutto il resto diverso.
\end{itemize}
Regola d'oro: \textbf{là fuori, nascosto nell'ombra, c'è sicuramente un gioco
molto simile al vostro}. Va solo trovato. E se per qualche ragione non ha
venduto, bisogna capire perché.

\subsection{Lavorare in engine e fare tuning}

Il level designer lavora in engine, così come il technical designer; ma
qualsiasi game designer, a seconda di ciò che deve fare, può finire in
engine, soprattutto per il \textbf{bilanciamento} e il \textbf{tuning}.

Il \emph{tuning} è l'azione di bilanciare con precisione i dati. Esempio: in
\emph{Super Mario} si modifica quanto in alto salta il personaggio, con che
velocità, in quanto tempo raggiunge l'apice, e così via. Questo lavoro rende
il gioco:
\begin{itemize}
  \item ad alto livello, dotato della \textbf{progressione} giusta;
  \item a basso livello, \textbf{piacevole da giocare}.
\end{itemize}
Il \emph{fine tuning} è un lavoro di dettaglio, paziente e minuzioso.

\subsection{Bilanciare}

Il bilanciamento, emerso tardi nel brainstorming, è una delle attività
centrali del designer e ha anche una specializzazione dedicata
(l'\emph{economy designer}, sezione~\ref{sec:tipi-designer}).

\subsection{Analizzare il mercato}

In pre-produzione, di solito, sì: capire le tendenze, i generi che cambiano
nel tempo, la compatibilità del gioco con il momento storico, in vista ad
esempio di un possibile publisher. Va capito quanto, come e con quanto tempo
farlo, perché si è sempre limitati da molte cose.

\subsection{Sintesi}

\begin{quote}
Il designer parte dall'idea ma non si ferma lì: la confuta, lavora sui
dettagli, bilancia, fa \emph{fine tuning}; lavora con il team e fa da
\emph{manager} per diversi dipartimenti; e naviga costantemente tra i limiti.
\end{quote}

% ---------------------------------------------------------------------
\section{I limiti del design}
\label{sec:limiti}

I limiti ci sono sempre, e nel mondo del design sono spesso numerosi. I più
importanti entro cui navigare sono:
\begin{description}
  \item[Tempo] ad alto livello, l'idea deve essere realizzabile nel tempo di
        produzione disponibile; a basso livello, il proprio lavoro
        settimanale deve essere fattibile nel tempo che si ha. Non ci si può
        dividere tra mille feature team se realisticamente se ne può seguire
        uno solo.
  \item[Tecnologia] si può o non si può fare qualcosa? Ci sono i server per un
        multiplayer con un certo numero di giocatori? Ci si può permettere,
        ad esempio, una fauna che reagisce alla fisica come quella di
        \emph{GTA VI}, che richiede il meglio dei \emph{technical artist}?
  \item[Background e competenze] che cosa sa fare il team? Realisticamente,
        chi ha lavorato dieci anni su giochi horror viene assunto per un
        horror, perché ha le competenze per risolvere problemi già visti.
  \item[Persone e motivazione] il team ha la motivazione e le competenze per
        fare un certo tipo di gioco, e per lavorare insieme?
  \item[Denaro] tutto ruota intorno ai soldi. Che si lavori da soli o in mille,
        le persone vanno pagate e il gioco deve avere un ritorno economico.
        Bisogna decidere dove tracciare la linea: raggiungere il
        \textbf{break-even} (guadagnare esattamente quanto speso) o puntare a
        grandi guadagni. Con un \textbf{publisher}, inoltre, una parte
        consistente dei ricavi va a lui (ripartizioni dell'ordine di 80/20 o
        70/30, che dipendono dal contratto firmato).
\end{description}

\subsection{Non siete ``gli uomini delle idee''}

Per qualcuno può essere un problema scoprirlo, ma il designer non è
``l'uomo (o la donna) delle idee'': accanto alla parte creativa c'è una parte
tecnica, e si è fortemente limitati da molti fattori. Questi limiti, però,
sono un aiuto:
\begin{itemize}
  \item riducono il ventaglio di possibilità potenzialmente distruttive per il
        team (cose che non si riescono a fare, che non piacciono o per cui
        mancano le risorse);
  \item aiutano a mirare bene, trovando soluzioni creative in uno spazio più
        piccolo e più attuabile.
\end{itemize}

\begin{quote}
\centering
\textbf{Le idee sono nulla, l'implementazione è tutto.}
\end{quote}

% ---------------------------------------------------------------------
\section{Ruoli e specializzazioni}
\label{sec:tipi-designer}

\subsection{La scala gerarchica}

\begin{description}
  \item[Creative Director (CD)] mantiene la \textbf{visione} e dà l'idea, in
        modo che tutto segua una certa direzione. In teoria può diventarlo
        chiunque, ma con un'esperienza notevole, nell'ordine dei decenni: si
        assume che in quel tempo abbia visto come lavorano molti dipartimenti.
        Naturalmente un CD con un background artistico ha meno competenze di
        design, e viceversa.
  \item[Director di disciplina] uno per dipartimento: \emph{game director},
        \emph{level director}, \emph{art director}, ecc.
  \item[Lead] il capo del dipartimento, che gestisce i gruppi di persone. Chi
        lavora in un dipartimento risponde al proprio lead (ad esempio un game
        designer risponde al \emph{lead designer}, non al game director).
  \item[Senior / Principal] i profili più esperti (a seconda dell'azienda).
  \item[Junior] il profilo di inizio carriera.
  \item[Intern] lo stagista.
\end{description}
Lo schema si replica uno a uno per ogni dipartimento: arte, suono, animazione
(a volte arte e animazione insieme, o \emph{rigging} e animazione insieme),
programmazione, UX/UI e così via. A fare la differenza, a parità di ruolo, è
il \textbf{livello di seniority}.

\subsection{UX e UI}

UX e UI stanno a metà tra game design e arte: a seconda dell'azienda il ruolo
si chiama \emph{UX/UI Artist} o \emph{UX/UI Designer}.
\begin{description}
  \item[UX (User Experience)] come una persona interagisce con le cose. Esempio:
        su uno smartphone, un pulsante accende e spegne lo schermo; per
        digitare la password si scorre verso l'alto e compaiono i punti da
        toccare.
  \item[UI (User Interface)] come sono ``vestite'' queste cose: ad esempio
        l'aspetto della barra dei punti vita (HP), realizzata con asset 2D o
        3D.
\end{description}

\subsection{Le specializzazioni del game design}

Più l'azienda è grande, più la specializzazione è puntuale.
\begin{description}
  \item[Gameplay / Core mechanics designer] (sono la stessa cosa) si occupa
        delle meccaniche e delle 3C.
  \item[System designer] si occupa dei sistemi e delle loro interazioni.
  \item[Economy designer] si occupa del bilanciamento; lavora soprattutto con
        fogli di calcolo (Excel, Google Sheets).
  \item[Combat designer] si occupa del combattimento; negli annunci di lavoro
        indica quasi sempre un combattimento \emph{action}, con gameplay
        reattivo (non strategico), anche se non è una regola.
  \item[Multiplayer designer] si occupa delle interazioni multigiocatore.
  \item[Technical designer] fa da ponte tra game design e programmazione.
\end{description}

\subsection{Le specializzazioni del level design}

\begin{description}
  \item[Encounter designer] posiziona i nemici nella mappa in modo che gli
        scontri siano bilanciati e la progressione equilibrata come
        desiderato.
  \item[World / map designer] popola il mondo di gioco.
  \item[Quest designer] popola il mondo con le missioni e pensa alla
        progressione al loro interno, cioè a dove il giocatore deve andare (si
        vedano ad esempio gli annunci di lavoro di CD Projekt Red o di
        Larian).
  \item[Technical level designer] fa da ponte tra level design e
        programmazione.
\end{description}

\subsection{Discipline vicine: suono, UI, produzione}

Il monte ore del corso non permette di uscire come \emph{narrative designer}
(o \emph{sound designer}) completi. Ciò non toglie che si possano coltivare
queste discipline:
\begin{itemize}
  \item nei progetti di gruppo si può occuparsi anche del suono, della UI/UX,
        o provare ad applicare metodologie come Agile per fare
        \emph{producing} e gestire il team;
  \item il tempo dedicato a queste attività va però sottratto al game design,
        e dipende da quante persone ci sono nel team, da quanto lavoro c'è e
        da quanto è grande il progetto;
  \item avere sviluppato e rilasciato un progetto, anche piccolo (anche una
        \emph{game jam}), documentando i propri processi, è già un ottimo
        argomento da presentare in futuro;
  \item la via migliore resta lavorare nell'industria e, da lì, chiedere di
        imparare anche altre discipline.
\end{itemize}
Un gioco si può fare su qualsiasi cosa: sugli scarafaggi, sui piedini delle
sedie, sui pennarelli. Si parte avvantaggiati da ciò che è vicino alla propria
sensibilità: chi ama la musica può pensare a un gioco con meccaniche legate
alla musica, con un gameplay innovativo ma contenuto e realizzabile. Il design
è un cappello che comprende molte discipline.

% ---------------------------------------------------------------------
\section{Approfondimenti emersi in lezione}
\label{sec:approfondimenti}

\subsection{Roguelike e roguelite}

Un esempio di come il marketing conia generi: Supergiant Games, con
\emph{Hades}, rese celebre il termine \textbf{roguelite} al posto di
\emph{roguelike} (pur non essendo il primo roguelite in assoluto). La
differenza:
\begin{itemize}
  \item il \textbf{roguelike} si fonda su una serie di punti fissati in una
        sorta di \emph{manifesto} che definisce che cosa rende un gioco
        roguelike; uno su tutti: a ogni morte si ricomincia la \emph{run} da
        capo;
  \item il \textbf{roguelite} rompe alcuni di questi punti, in particolare
        introducendo una \textbf{progressione esterna alla run}: il giocatore
        può potenziarsi fuori dalla singola partita per vincere più
        facilmente in quelle successive. \emph{Hades} ne è il caso principe.
\end{itemize}
Analogamente, un mobile game sviluppato dal direttore del corso fu
promosso come ``il primo \emph{donation game} al mondo'': un gioco che
convertiva il tempo di gioco in donazioni reali a enti benefici per animali.
Anche qui un neologismo creato lato marketing per definire un genere.

\subsection{Applicazioni fuori dall'industria videoludica}

Le competenze sviluppate non servono solo ai videogiochi:
\begin{itemize}
  \item i \textbf{render architettonici} oggi possono essere realizzati con
        Unreal, e diversi professionisti formati nel game design lavorano in
        questo campo;
  \item la \textbf{gamification}, cioè l'applicazione di regole videoludiche
        a contesti non ludici (mostre interattive, progetti culturali,
        consulenze come quella per applicare questi concetti alle macchine da
        allenamento di un'azienda di fitness), è in Italia uno dei modi più
        redditizi di spendere queste competenze, spesso con finanziamenti
        pubblici.
\end{itemize}

% ---------------------------------------------------------------------
\section{Il primo task}
\label{sec:primo-task}

\begin{enumerate}
  \item Creare un indirizzo e-mail \textbf{professionale} (consigliato: un indirizzo Gmail nella forma
        \texttt{nome.cognome@gmail.com}; l'anno di nascita è accettabile).
        Va bene anche un indirizzo professionale già esistente usato per
        lavoro. Da evitare indirizzi poco seri (``non \texttt{blackfire67}'').
  \item Una persona della classe raccoglie tutti gli indirizzi in un unico
        documento e lo invia via e-mail al docente, che aggiungerà tutti al
        Drive condiviso del corso (dove verranno caricate lezioni, slide e
        materiali, per avere i documenti tutti nello stesso posto).
  \item Installare Discord e iscriversi al server del corso, dove passano le
        comunicazioni ufficiali. È consuetudine, e utile, creare anche un
        gruppo tra studenti.
\end{enumerate}
\textbf{Scadenza}: prima della lezione del venerdì.

% ---------------------------------------------------------------------
\section*{Punti chiave del capitolo}
\addcontentsline{toc}{section}{Punti chiave del capitolo}

\begin{itemize}
  \item Il corso è una \textbf{simulazione di lavoro}: ci si tratta da
        colleghi, con una gerarchia da studio di sviluppo.
  \item Si impara con l'\textbf{inquiry-based learning}: lezione, task,
        dubbi, domande, condivisione, progetti.
  \item Il fallimento affrontato con la mentalità giusta è la principale fonte
        di crescita; i feedback riguardano il lavoro, mai la persona.
  \item Il lavoro del designer è per metà creativo e per metà
        \textbf{ricerca e metodo scientifico}; il game design non è una
        scienza esatta.
  \item \textbf{Documentazione} e \textbf{pitching} sono il cuore del
        mestiere.
  \item Il designer progetta meccaniche e sistemi (partendo dalle
        \textbf{3C}), documenta, media, prototipa, confuta le idee, studia i
        competitor, bilancia e fa tuning.
  \item La leadership del designer è \textbf{al servizio} del team, non
        verticale.
  \item I limiti (tempo, tecnologia, competenze, persone, denaro) restringono
        e indirizzano la creatività.
  \item \textbf{Le idee sono nulla, l'implementazione è tutto.}
\end{itemize}

%    % ...
%
%  (Con hyperref aggiungere anche \renewcommand{\theHchapter}{z.\arabic{chapter}}
%  nella parte zero, per evitare ancore duplicate.)
%  Tutte le etichette di questo capitolo hanno il prefisso "z1:" per non
%  collidere con quelle dei capitoli del corso.
% =====================================================================

\chapter[La cultura del feedback]{Lavorare in squadra: la cultura del feedback}
\label{cap:z1-feedback}

Questo capitolo apre le \emph{lezioni zero}: incontri che precedono l'inizio
vero e proprio del corso di Game Design e che non riguardano una tecnica di
design in senso stretto, ma \textbf{la figura professionale} di chi lavora
nell'industria dei videogiochi. Prima di imparare a progettare un gioco
conviene infatti capire in che tipo di ambiente lo si progetterà: un ambiente
fatto di team, di gerarchie, di scadenze e, soprattutto, di persone che devono
riuscire a lavorare insieme.

La lezione è la prima giornata di un \textbf{workshop di due giorni} (circa
dieci ore in totale, pause pranzo comprese) tenuto da Stefano
Cozzi\footnote{Il cognome è riportato come compare nella trascrizione
automatica, che in più punti storpia i nomi propri; andrebbe verificato.
Nella stessa trascrizione lo studio per cui lavora appare come
``Stormine'', ``Stormline'' e ``Stormind'': si tratta di Stormind Games,
studio con sede a Catania, coerentemente con quanto detto in lezione.}, un
professionista con circa quindici anni di esperienza nel settore, tra
Blizzard Entertainment, Riot Games e, oggi, Stormind Games. Il workshop è
nato da un confronto con i coordinatori dell'accademia\footnote{Nella
trascrizione l'accademia è resa in modi diversi (``Hive'', ``a dive'',
``D.I.V.''); i coordinatori citati sono Enrico e Andrea, insieme al
direttore Diego Ricchiuti. Nel resto del capitolo si usa semplicemente
``l'accademia'', come negli altri capitoli.}: serviva un momento dedicato a
\emph{come ci si comporta in un ambiente professionale}, cosa che riguarda sia
il lavoro futuro sia i progetti di gruppo svolti durante il percorso di
studi, dove ci si deve già abituare a comportarsi ``come se si fosse in
un'azienda''.

La prima giornata ha un unico grande tema, sezionato in più parti:
\textbf{la cultura del feedback}. La seconda giornata (che sarà oggetto della
prossima lezione zero) sarà più pratica: \emph{time management},
\emph{brainstorming} e il ``guidare da dentro'' un team anche senza esserne
il leader (sezione~\ref{sec:z1-domani}).

Il percorso del capitolo segue quello della lezione:
\begin{enumerate}
  \item il formato del workshop e il percorso professionale del relatore
        (sezione~\ref{sec:z1-relatore});
  \item perché le aziende tengono tanto al lavoro di squadra e al
        \emph{fit} culturale (sezione~\ref{sec:z1-squadra});
  \item che cos'è il feedback e da che cosa dipende il modo di darlo
        (sezione~\ref{sec:z1-cose-feedback});
  \item le cinque fasi di vita di un team e il valore del
        \emph{post-mortem} (sezioni~\ref{sec:z1-fasi-team}
        e~\ref{sec:z1-qa-fasi});
  \item il confronto con i colleghi e alcune domande sul lavoro nei grandi
        studi (sezioni~\ref{sec:z1-rivalita} e~\ref{sec:z1-industria});
  \item il circolo virtuoso tra feedback, miglioramento e fiducia, e i suoi
        effetti sulla \emph{retention} e sulla cultura aziendale
        (sezioni~\ref{sec:z1-triangolo} e~\ref{sec:z1-cultura});
  \item il ``gioco del traditore'' e il principio del
        \textbf{default to trust} (sezioni~\ref{sec:z1-among-us}
        e~\ref{sec:z1-default-trust});
  \item le traiettorie di carriera e la valutazione delle performance
        (sezione~\ref{sec:z1-traiettorie});
  \item come si dà, come si riceve e quando si dà un feedback
        (sezioni~\ref{sec:z1-dare}, \ref{sec:z1-ricevere}
        e~\ref{sec:z1-quando});
  \item il feedback come regalo, le domande finali e le anticipazioni
        della seconda giornata (sezioni~\ref{sec:z1-regalo},
        \ref{sec:z1-qa-finale} e~\ref{sec:z1-domani}).
\end{enumerate}

Una premessa di metodo: il relatore ha alternato continuamente teoria,
domande alla classe e racconti della propria esperienza. I racconti non sono
digressioni: sono gli esempi concreti con cui i principi vengono resi
tangibili, e per questo sono riportati con cura.

% ---------------------------------------------------------------------
\section{Il workshop e il relatore}
\label{sec:z1-relatore}

\subsection{Formato e regole di partecipazione}

Il relatore ha cercato di ``riassumere quindici anni di esperienza in dieci
ore''. Ha materiale per circa il doppio delle ore disponibili: le cose
strettamente necessarie vengono trattate per prime, il resto è tenuto ``in
riserva'' come \emph{nice to have}. Le regole del workshop:
\begin{itemize}
  \item le domande si fanno \textbf{alla fine di ogni slide}, così da chiudere
        prima il discorso; non c'è un'unica sessione finale di domande e
        risposte (anche se ci sarà), perché il relatore preferisce fermarsi
        spesso e non fare ``uno spiegone di dieci ore'';
  \item le domande non devono per forza riguardare la slide appena vista: se
        una slide dà lo spunto per altro, se ne può parlare; al massimo sarà
        il relatore a rispondere in breve o a rimandare a una slide
        successiva;
  \item si può chiedere anche di esperienze passate, opinioni e sensazioni;
        se per NDA o per serietà professionale non può rispondere, lo dirà
        esplicitamente (``non lasciate che sia un \emph{blocker}'');
  \item su Zoom si usa l'alzata di mano o la chat, che il relatore tiene
        aperta sul secondo monitor;
  \item non ci sono pause ufficiali: la lezione è registrata e ognuno è
        libero di assentarsi quando serve, cercando però di non disturbare il
        flusso della lezione;
  \item le slide verranno fornite integralmente tramite la segreteria
        dell'accademia (Vanessa), insieme alle registrazioni: le slide senza
        il parlato servono poco, ma sono un buon promemoria.
\end{itemize}

\subsection{Il percorso professionale}

Il relatore si presenta (scherzosamente) come la ``guida escursionistica''
che accompagna la classe nella ``foresta incantata'' della cultura del
feedback, perché per molti è un territorio sconosciuto. Il suo percorso:
\begin{itemize}
  \item in Italia ha lavorato per aziende che fornivano servizi a terzi, tra
        cui il lavoro su \emph{Warhammer Online}, e per un canale televisivo
        sul digitale terrestre dedicato ai videogiochi\footnote{Nella
        trascrizione: ``Digital Game Network Channel''. Il nome esatto del
        canale andrebbe verificato.};
  \item a \textbf{Blizzard Entertainment} (circa sei anni, a Parigi) ha fatto
        il \emph{community manager} per l'Italia (\emph{StarCraft~II},
        \emph{World of Warcraft} da \emph{Mists of Pandaria} in poi,
        \emph{Hearthstone}, \emph{Diablo~III}, le espansioni di
        \emph{StarCraft}, il lancio di Battle.net), per poi diventare il
        manager dei community manager europei;
  \item a \textbf{Riot Games} (circa cinque anni, a Dublino) è stato
        brevemente community manager italiano e poi \emph{publishing
        manager} di un team di circa trenta persone per l'Europa. Il
        \emph{publishing} non è solo community: è l'unione di
        \textbf{community, marketing e PR};
  \item a \textbf{PG Esports} si è occupato di eventi dal vivo e online
        (ad esempio la Milano Games Week e tornei di \emph{Valorant} e
        \emph{League of Legends}), restando in contatto con Riot e Blizzard;
  \item da circa un anno è tornato nello sviluppo, a \textbf{Stormind
        Games}, dove oggi si occupa di risorse umane e organizzazione del
        lavoro; lo studio verrà presentato più ampiamente al \emph{Career
        Day} dell'accademia insieme alla sua HR manager;
  \item fa inoltre consulenze aziendali anche al di fuori dei videogiochi e
        insegna in master universitari (IULM di Milano).
\end{itemize}

Il suo ruolo ricorrente, sia a Blizzard sia a Riot, è stato quello di
\emph{Mr.~Wolf} (``per chi ha visto \emph{Pulp Fiction}''): quello che
risolve i problemi. Veniva inserito nei team più in difficoltà, ad esempio
dal punto di vista delle performance, per sistemarli e passare al successivo.
Era presente in alcuni dei momenti più critici, come il lancio di
\emph{Diablo~III} o il lancio di espansioni di \emph{World of
Warcraft}\footnote{La trascrizione riporta ``il lancio con i Gary's'' e
``World of Trenor'': si tratta verosimilmente di \emph{Mists of Pandaria} e
di \emph{Warlords of Draenor}.}, a fare \emph{damage control}.

La scoperta fondamentale di questi anni è che i problemi da risolvere
\textbf{quasi mai riguardavano il prodotto} (non era certo lui a sistemare i
server di \emph{Diablo~III}), ma i team che c'erano dietro: persone che
litigavano, non collaboravano, non riuscivano a stare nella stessa stanza
senza urlarsi contro. E nel 90\% dei casi tutto questo si riduce a
\textbf{problemi di comunicazione}: è raro che perfino le incompatibilità
caratteriali non siano risolvibili parlando.

\subsection{Struttura non vuol dire burocrazia}

A Stormind il relatore sta accompagnando un team passato da 30 a 95 persone
in meno di un anno, attraverso tutti i cambiamenti di cui parla il workshop.
Darsi una struttura viene spesso percepito male, come ``più processi, più
procedure'', ma bisogna distinguere:
\begin{itemize}
  \item un processo come ``timbrare il cartellino'' può essere davvero
        invasivo e noioso;
  \item definire \emph{come funziona un documento di game design, chi lo
        approva, chi lo visiona e chi no} è invece una struttura necessaria.
\end{itemize}
Uno studio di 30 persone nella stessa sede può farne a meno: quando c'è un
dubbio ci si alza e si va alla scrivania del collega. Uno studio di 95
persone sparse per l'Italia (e in parte per l'Europa) non può più farlo:
servono una struttura diversa, un modo di lavorare diverso e una
pianificazione del futuro diversa.

% ---------------------------------------------------------------------
\section{Perché il lavoro è sempre lavoro di squadra}
\label{sec:z1-squadra}

\subsection{Non si va in ufficio per fare amicizia}

Il primo messaggio è netto: \textbf{non si va in ufficio per fare amici} o
per trovare compagni di vita. Si va in ufficio perché c'è un lavoro da fare, e
lo scopo è riuscire a farlo \emph{anche con chi sta antipatico}. Con le
persone simpatiche è facile: diventeranno amici, e non sarà mai un problema.
È con chi sta antipatico che le cose si complicano, ed è lì che emergono i
\textbf{veri professionisti}: quelli che riescono a lavorare con chiunque
senza che l'antipatia pesi sul modo in cui lavorano e collaborano, mettendo il
prodotto e il team davanti al proprio ego e alle proprie preferenze personali.

La frase che fa da filo conduttore a tutto il workshop (e che il relatore
promette di riproporre ``fino alla nausea'') è:

\begin{quote}
  \emph{La capacità di un gruppo di persone di fare cose straordinarie
  dipende dal modo in cui queste persone operano insieme.}
\end{quote}

Le aziende raramente pensano che l'innesto di una singola persona possa
cambiare il destino di un progetto. Questo è tanto più vero quanto più grande
è il team, ma vale anche per i piccoli studi, anche per un gruppo di cinque
persone. Come negli sport di squadra: l'individuo conta fino a un certo punto,
e anche se si ha Dončić in squadra, se si gioca uno contro cinque si perde
comunque\footnote{Il riferimento è alle finali NBA in corso nei giorni del
workshop tra Dallas Mavericks e Boston Celtics (giugno 2024), il che
permette di datare approssimativamente la lezione.}.

\subsection{Lo \emph{useless genius} e il costo delle persone tossiche}

Molti dipartimenti HR usano la definizione di \textbf{useless
genius}\footnote{Il concetto è più noto in letteratura manageriale come
\emph{brilliant jerk}; in lezione è usata l'espressione \emph{useless
genius}.}: si può essere bravissimi, a volte i migliori nel proprio lavoro,
ma se non ci si sa comportare, relazionare e operare in un gruppo si è
fondamentalmente inutili ai fini del gruppo. Nel migliore dei casi inutili,
nel peggiore \emph{dannosi}.

In passato le aziende tendevano a ``chiudere un occhio'': la persona è
fondamentale, la teniamo. Oggi la domanda è diventata: \emph{quanto costa
all'azienda, oltre al suo stipendio, tenere una persona che crea un ambiente
tossico e peggiora chi le sta intorno?} Ci sono costi sommersi che le aziende
cominciano a vedere, e la tendenza è liberarsi di queste persone il prima
possibile. Il ragionamento è semplice:
\begin{itemize}
  \item le \textbf{competenze si possono trasmettere} o sostituire: si trova
        qualcuno che fa lo stesso lavoro magari nel doppio del tempo, oppure
        più persone che se lo dividono;
  \item il \textbf{carattere delle persone} invece cambia molto raramente.
\end{itemize}

\subsection{Il fit culturale}

Quando qualcuno \textbf{non è un \emph{fit} culturale} per un'azienda è
meglio separarsi presto. Ma non essere adatti a \emph{quell'}azienda non
significa non essere adatti a lavorare: è perfettamente legittimo che una
persona inadatta a un'azienda sia adattissima a un'altra.

Il relatore l'ha vissuto in prima persona. Dopo sei anni a Blizzard viene
contattato da Riot e affronta un processo di selezione lunghissimo (``quasi
undici colloqui''), al termine del quale viene scartato. Scrive allora una
mail che gli ha ``cambiato la traiettoria di carriera'': chiede
semplicemente perché non è stato ritenuto adatto. La risposta: in nessun
colloquio si era esposto abbastanza da far capire \emph{perché} volesse
lasciare Blizzard, e quindi pensavano che in realtà stesse bene dov'era. Il
motivo era che, per educazione professionale, non gli piace parlare male del
proprio posto di lavoro; portata all'eccesso, però, questa riservatezza era
diventata un ``è tutto bellissimo, ma vorrei provare Riot'', senza un perché.
Lo invitano comunque a un incontro dal vivo (``dal vivo è diverso''), e alla
fine viene assunto.

Una volta arrivato, i sei mesi di prova (\emph{probation}) sono serviti tutti
per smettere di ragionare e lavorare ``come a Blizzard''. La differenza
culturale era enorme, tanto che in Riot esisteva un gruppo di supporto
interno tra ex dipendenti Blizzard: un collega gliel'ha descritto come
``uscire da una relazione tossica e doversi riabituare a una relazione
normale''\footnote{Il relatore precisa che parla della Blizzard di allora,
che non esiste più: oggi è parte di Microsoft e nessuno dei suoi ex colleghi
vi lavora ancora, anche a causa dei licenziamenti di massa.}.
\begin{description}
  \item[Blizzard] era \emph{developer first}: il focus era chiaramente sugli
        sviluppatori, prima di tutti gli altri. Dopo sei anni c'erano ancora
        meeting a cui non poteva partecipare perché non abbastanza
        \emph{senior}, informazioni che non poteva sapere, piani che non
        poteva vedere.
  \item[Riot] era \emph{player first}: prima i giocatori, e i dipendenti come
        mezzo per trasferire il feedback dei giocatori dentro il gioco. Il
        secondo giorno lo portano in una stanza e gli mostrano il piano dei
        cinque anni successivi. Lui, sorpreso, pensa a un errore (``sono qui
        da oggi, non vorrei mettervi nei guai''); la risposta è: ``sei un
        dipendente Riot o no? Allora non ci sono livelli di accesso: devi
        sapere queste cose per fare meglio il tuo lavoro''.
\end{description}

Da questa esperienza ha imparato molto sulla \textbf{trasparenza}: condividere
le informazioni aiuta il lavoro. Non sempre, però: a volte non si condivide
qualcosa non per segretezza, ma perché non è ancora pronta per essere
presentata, esattamente come non si mostra un gioco quando è ancora solo
teoria (``te lo faccio vedere quando ho qualcosa di tangibile'').

La lezione per gli studenti è doppia: le aziende valuteranno il \emph{fit}
culturale sempre di più, ma capire se si è un \emph{fit} richiede anche
\textbf{lavoro da parte del candidato}. Bisogna avere la presenza di spirito
di dire: ``conosco i valori di questa azienda, non mando nemmeno il CV,
perché so che durerei un mese e mezzo''.

% ---------------------------------------------------------------------
\section{Che cos'è il feedback}
\label{sec:z1-cose-feedback}

\subsection{Le prime risposte della classe}

La prima domanda alla classe è: \emph{quando sentite la parola
``feedback'', a cosa pensate?} Le risposte raccolte, a voce e in chat:
\begin{itemize}
  \item un \textbf{riscontro} su ciò che si fa e su come lo si fa, sul lavoro
        o sul modo in cui ci si interfaccia con il gruppo; può essere
        lavorativo oppure comunicativo e comportamentale. Chi ha risposto lo
        vede come qualcosa che, se l'azienda non dà, va attivamente
        richiesto, perché è importante sapere come sta andando per
        migliorarsi;
  \item né positivo né negativo, ma \textbf{costruttivo}: in entrambi i casi
        serve a migliorarsi;
  \item uno \emph{scambio di opinioni}, un \emph{riscontro produttivo};
  \item un feedback ``di reazione'', non verbale: si fa una cosa e si
        osserva come cambia l'ambiente o il comportamento delle persone
        intorno (anche perché a voce la gente spesso si tiene le cose
        negative);
  \item per un designer, anche il \textbf{linguaggio del corpo}: osservando
        un giocatore giocare si può dedurre l'emozione che prova (frustrazione,
        incomprensione).
\end{itemize}

Sulla forma (scritto o orale?) le opinioni divergono: c'è chi lo preferisce
\emph{mille volte} scritto; chi pensa che dipenda dalle dimensioni
dell'azienda (orale nelle piccole, scritto in quelle con team di centinaia di
persone); chi pensa che dipenda dall'importanza del feedback (scritto se
tecnico o pieno di dati, da tenere sotto gli occhi; orale se è un consiglio
veloce).

Il relatore usa due esempi per mostrare che \emph{dipende}:
\begin{itemize}
  \item un feedback sulla programmazione ha senso darlo scritto, magari
        commentando direttamente il codice, più che ``andando a prendere un
        caffè'';
  \item ma se il feedback è ``ieri, durante il meeting, ho percepito male il
        fatto che mi interrompessi di continuo'', ce lo si aspetta
        \textbf{orale}.
\end{itemize}

\subsection{Da che cosa dipende il modo di dare feedback}

Lo \emph{spoiler} della giornata è che \textbf{non esiste una risposta univoca}
su come dare feedback. Il workshop non insegna ``il'' modo, ma fornisce
\emph{tutte le opzioni}, perché il modo giusto cambia in base a molti fattori:
\begin{itemize}
  \item a chi lo si dà: un collega, un superiore, il proprio capo, un proprio
        \emph{report} (un sottoposto);
  \item se è personale o professionale;
  \item se è la prima volta o la terza volta che si dà lo stesso feedback;
  \item se riguarda un progetto importante o secondario.
\end{itemize}

Manca però l'elemento fondamentale, che la classe individua subito:
\textbf{l'altro}. A parità di tutti gli altri fattori, il \emph{carattere
della persona che riceve il feedback} dovrebbe far cambiare il modo in cui lo
si dà. Per questo, prima della lezione, agli studenti era stato chiesto di
fare un test di personalità: non un test di psicologia avanzata, ma uno
strumento per rendersi conto di quanta varietà ci sia anche solo in una
classe di 53 persone (il che non vuol dire trovare 52 modi diversi di dare
feedback, ma lavorarci un po' sì).

Il relatore porta se stesso come esempio. Preferisce:
\begin{itemize}
  \item i feedback \textbf{positivi scritti}, perché è a disagio a ricevere
        complimenti in pubblico;
  \item i feedback sul proprio lavoro \textbf{diretti}: ``vai subito al
        punto e non girarmi intorno''.
\end{itemize}
È molto orgoglioso del suo lavoro e, quando qualcuno lo critica, la sua
primissima reazione istintiva è ``come ti permetti''. Negli anni ha imparato
a farla durare un secondo, ma chi gli dà feedback di persona dice di
vederla ancora negli occhi. Per mitigarla, dice esplicitamente a chi gli dà
feedback: \emph{``ignora il modo in cui reagisco nei primi cinque secondi,
vai avanti, poi mi passa''}. È considerato da tutti molto aperto al feedback
e non ha problemi a dire ``hai ragione, facciamo diversamente'': ma ci sono
reazioni istintive e caratteriali che è difficile controllare al cento per
cento.

\subsection{Con chi si lavora meglio}

Guardando le persone con cui si lavora abitualmente (anche nei progetti
dell'accademia), si noterà che si lavora meglio con due profili:
\begin{enumerate}
  \item chi è \textbf{esattamente come noi}, ma qui è un 50 e 50: alcuni
        caratteri vanno d'accordo con i propri simili, altri non vogliono
        nessuno con lo stesso carattere nel gruppo (e si potrebbe rientrare
        in questa categoria);
  \item chi è \textbf{affine}: ha qualcosa in comune con noi, ma non tutto,
        e proprio per questo può dare punti di vista differenti.
\end{enumerate}
Il lavoro della giornata è capire come valorizzare il secondo caso: ottenere
feedback da persone con caratteri, visioni e background diversi, per
completare una visione del proprio lavoro che è necessariamente limitata. È
come esplorare una mappa con la nebbia di guerra: un'altra persona, al di là
del muro, vede cose che noi non vediamo.

\subsection{Il feedback come somma di piccoli miglioramenti}

Per il relatore il feedback è questo: \textbf{applicare tanti piccoli
miglioramenti periodici}, che si sommano nel tempo. Migliorare dello 0,1\%
ogni giorno porta a fine anno a un miglioramento enorme, e questo vale per i
prodotti, per i giochi, per le persone e per i team.

\paragraph{L'esempio del ciclismo britannico.} La squadra di ciclismo
britannica, fino a una quindicina d'anni fa, non aveva vinto quasi nulla.
Cambiato il team manager, tutti si aspettavano grandi cambiamenti sportivi
(convocazioni, strategie di gara). Arrivarono invece \emph{minuscoli
cambiamenti non sportivi}: un altro chef in trasferta, materassi migliori nei
camper per riposare meglio, un sellino più comodo a fine tappa. Cambiamenti
marginali che da soli non avrebbero cambiato niente; eppure, dopo due o tre
anni di questo lavoro quotidiano, la squadra ha cominciato a vincere a
ripetizione, Olimpiadi comprese, cambiando il modo in cui il ciclismo è
visto in Gran Bretagna\footnote{È la nota strategia delle \emph{marginal
gains} associata a Dave Brailsford (British Cycling e Team Sky); il nome non
è citato in lezione.}. Certo, si è investito anche sulle persone giuste, ma
tutto è partito dai piccoli miglioramenti periodici.

Per gli studenti vale lo stesso: chiarire un conflitto con una persona con
cui si ha ancora poco a che fare, o una situazione che ha lasciato
``l'amaro in bocca'', sembra inutile oggi, ma darà i suoi frutti tra tre, sei,
nove mesi. Sono \emph{investimenti sul futuro}.

\paragraph{Peggiorare ogni giorno.} Esiste anche il percorso inverso:
lasciar correre, ``buttare sotto il tappeto'' le cose perché ``tanto è
inutile parlarne''. Queste cose tornano sempre fuori dopo otto mesi, come il
\emph{libro dei rancori} dei nani di Warhammer: ``ti ricordi quando, un anno
e mezzo fa, all'inizio del progetto, mi hai detto quella cosa? Mi sono
offeso''. A quel punto è quasi impossibile risolvere, perché nessuno ricorda
più con precisione cosa sia successo (``ma quale cena?'', ``forse era in
pizzeria'', ``no, eravamo da un'altra parte''). \textbf{Tenersi le cose nel
cassetto per tirarle fuori quando servono è sempre sconsigliato}: non servono
ed è impossibile usarle in modo costruttivo.

% ---------------------------------------------------------------------
\section{Le cinque fasi di vita di un team}
\label{sec:z1-fasi-team}

\subsection{Perché conoscerle}

Quasi tutti gli studenti lavorano già in gruppo (al primo anno su un gioco da
tavolo, negli anni successivi su un videogioco) e si saranno accorti che un
team attraversa delle fasi: il modo di lavorare insieme cambia nel tempo.
Spesso ci si fa prendere di sorpresa: ``eravamo così affiatati, cosa è
successo?''.

Un esempio da una consulenza fuori dai videogiochi: un'azienda farmaceutica
gli affida un team commerciale (rappresentanti che presentano medicinali a
ospedali, medici e privati). Il problema: ``prima andava tutto bene; ora che
siamo in diciotto dovremmo lavorare meno e meglio, e invece va peggio''. Va
peggio perché le persone nuove sono appena arrivate, mentre il management si
aspettava già di vedere i risultati dell'investimento.

Fino a pochi anni fa le fasi di un team venivano disegnate come una
\textbf{linea dritta verso l'alto}: le fasi erano le stesse, ma viste come un
miglioramento costante. Non si vedevano le parti \emph{fisiologiche} del
lavoro in team, cioè anche i peggioramenti. Accettarle è il primo passo per
accettare i fallimenti: nessun percorso personale o di team è lineare.

Il relatore ammette che il suo modo di lavorare è nato per il 70\% dal
proprio impegno a imparare, fare bene e rispettare le persone, e per il 30\%
dall'aver lavorato con ``incapaci colossali'' che gli hanno insegnato come
\emph{non} fare le cose. A Blizzard un manager gli fece capire che, pur non
sapendo ancora che tipo di \emph{people manager} sarebbe diventato, sapeva
benissimo quale tipo di manager \emph{non} voleva essere.

\subsection{Forming, storming, norming, performing, adjourning}

Le fasi sono cinque, un po' come le fasi del lutto. Molte aziende ne
considerano solo quattro (``e sono quelle che solitamente chiudono''); le
aziende al passo con i tempi ne considerano cinque\footnote{Il modello è
quello di Bruce Tuckman (1965), a cui la quinta fase, \emph{adjourning}, fu
aggiunta nel 1977. Il nome di Tuckman non è citato in lezione; la
trascrizione riporta la quinta fase come ``journing''.}.

\begin{description}
  \item[Forming] Si viene messi insieme e si comincia a lavorare. Non si sa
        nulla degli altri; spesso non si conoscono nemmeno i nomi, ed è
        normale.
  \item[Storming] Si cominciano a ``dare le testate come le capre in
        montagna per determinare il territorio'': questa cosa la faccio io,
        questa la vorrei fare così, chi decide? Si stabilisce chi fa cosa e
        come si divide il lavoro. Nel grafico è una \textbf{fase di
        peggioramento}: che il team litighi o discuta civilmente non cambia
        nulla, il team deve investire parte di tempo ed energie per risolvere
        queste questioni. È normale e \textbf{non è evitabile}.
  \item[Norming] Si comincia a capire come lavorano gli altri, cosa sanno
        fare e come far funzionare il gruppo; iniziano i miglioramenti. È una
        forma di \emph{accettazione} degli altri, che deriva dall'aver capito
        cosa devo fare io e cosa devono fare loro.
  \item[Performing] C'è uno sforzo \textbf{proattivo}: ``questa cosa
        toccherebbe a me, ma so che Alessio la fa molto meglio, quindi gli
        chiedo aiuto''. Si sfruttano i punti di forza degli altri, perché li
        si conosce e perché si sa che passare il testimone non avrà effetti
        negativi (il capo non dirà ``perché non l'hai fatta tu?'').
  \item[Adjourning] Il progetto finisce, si chiude o viene riorganizzato, il
        team viene sciolto. Quello che andrebbe fatto è analizzare
        \textbf{cosa ha funzionato, cosa no e cosa migliorare} per il
        progetto successivo.
\end{description}

\subsection{La fase che tutti saltano}

``Spoiler incredibile'': l'\emph{adjourning} è la fase che molte aziende
saltano, \textbf{perché c'è sempre il progetto dopo}.

A PG Esports, per la Milano Games Week (come per qualsiasi evento dal vivo a
cadenza annuale: Lucca Comics, Etna Comics, Comicon\dots), finita l'edizione
di un anno si cominciava dal giorno dopo a lavorare a quella successiva,
perché un anno sembra tanto ma non basta mai. Nei videogiochi è lo stesso:
finito un gioco, a parte le patch di mantenimento e del \emph{day one}, si
passa subito al successivo.

Il cambiamento introdotto dal relatore (in parte già praticato, ma da lui
strutturato) è stato: finita la Games Week, \textbf{tutti in ferie per una
settimana}; poi ci si ritrova per definire cosa ha funzionato e cosa no, con
un \textbf{post-mortem} (o \emph{debrief}); solo dopo si inizia a lavorare
all'edizione successiva. Senza questa fase manca la parte di
\emph{learning}, cioè di ciò che si è imparato dal progetto.

Il consiglio vale per gli studenti: quando si finisce il gioco da tavolo, e
poi il gioco digitale, bisogna trovare il tempo per il post-mortem,
altrimenti si butta via \textbf{il 90\% dell'utilità del progetto}. Una
struttura semplice per farlo è descritta nella sezione~\ref{sec:z1-qa-finale}.

% ---------------------------------------------------------------------
\section{Domande sulle fasi del team}
\label{sec:z1-qa-fasi}

\subsection{Lo storming che ritorna}

\paragraph{Capita che lo storming si ritagli spazio anche nelle altre fasi?}
Lo \emph{storming} è una fase precisa, che però molti manager, coordinatori e
team cercano di \textbf{saltare}: ``se io sono bravo come manager e loro sono
bravi come team, passeremo da forming a norming senza problemi''. È un errore,
perché senza una fase in cui si discute davvero non si sviluppa una
\emph{skill} fondamentale: \textbf{saper discutere}. Tutti conoscono persone
non abituate a farlo, per cui ogni minimo disaccordo diventa un affronto
personale (``avevamo deciso la pizza, perché andiamo al sushi? Perché mi fate
questo affronto?''): sono persone che non sono abituate ad avere conflitti
costruttivi.

Se si pensa di essere in \emph{norming} e continuano a emergere conflitti, il
punto è che \textbf{non si era davvero in norming}: lo storming non ``si
trascina'' nelle fasi successive, semplicemente non se n'è mai usciti.
Allo stesso modo, a chi chiede se ogni mini-fase di storming durante la vita
del gruppo sia frutto di uno storming fatto male o solo di carattere, la
risposta è già contenuta qui.

Un partecipante osserva che spesso, ancor prima di finire un progetto, si è
già messi su quello successivo: è il modo migliore per essere in costante
\emph{damage control}. Era l'esperienza del relatore a Blizzard: veniva
spostato fisicamente da un team all'altro, e finito il lancio di
\emph{Diablo~III} era già su \emph{World of Warcraft}, senza vivere il
``dopo''. Si è dovuto ritagliare da solo lo spazio per un debrief personale,
ed è lì che ha deciso che tipo di manager non voleva essere. Il consiglio:
\textbf{se non lo fa il team, fatelo voi}.

\subsection{Dove si passa la maggior parte del tempo}

\emph{In quale fase si passa più tempo nella propria carriera?} Gli studenti
rispondono \emph{storming} e \emph{norming}. Il relatore, in quindici anni,
calcola di aver passato in \emph{performing} circa \textbf{dieci mesi}. Non è
pessimismo: arrivare a performing deve essere sempre l'obiettivo, sia da
manager sia da membro del team; ma è difficile arrivarci, ed è normale che non
sempre ci si arrivi.

Lo storming può durare a lungo, ma con team molto esperti dura pochissimo,
perché c'è chi lo sa gestire. In realtà \textbf{la maggior parte del tempo si
passa in forming}, perché ogni equilibrio raggiunto vale per \emph{quelle}
persone, \emph{quel} progetto, \emph{quelle} tempistiche. Basta che cambi un
qualsiasi fattore (una deadline spostata, una persona nuova, qualcuno che si
dimette, un progetto cancellato) e si ricomincia, perché i bilanciamenti
precedenti non valgono più.

Esempio: un team di cinque persone affiatatissimo, in pieno performing. La
dirigenza, entusiasta, ``raddoppia le risorse'' e lo porta a dieci. O si
lascia tutto com'è (e i cinque nuovi non fanno niente), o bisogna ridiscutere
chi fa cosa e come i nuovi possono aiutare: che piaccia o no, si ricomincia
da forming.

Quindi la bravura di un manager, ma anche di un membro del team,
\textbf{non sta nel passare tanto tempo in performing, ma nell'accelerare il
passaggio da una fase all'altra}, e nel fare bene la fase di
\emph{adjourning}.

\subsection{Analizzare anche i successi}

C'è la tendenza a fare post-mortem solo dei progetti andati male: ``questa
cosa non ha funzionato, vediamo perché''. Va benissimo. Ma quando un progetto
va bene si dà per scontato che sia andato bene ``perché era buono''. Bisogna
invece chiedersi \emph{perché} è andato bene: era il prodotto? La
tempistica? Il modo in cui abbiamo fatto una cosa piuttosto che un'altra?
\textbf{Se non si capisce davvero il motivo del proprio successo, si
tenderanno a replicare le cose sbagliate.}

\paragraph{Blockbuster.} Quando arrivò la notizia che Netflix avrebbe
lanciato un servizio di streaming, Blockbuster (la catena di noleggio di
videocassette e DVD) aveva tutto il tempo per preparare un'offerta
concorrente. In una riunione del management raccontata in un libro (di cui il
relatore non ricorda il titolo), la direzione sembrava orientata a investire
nello streaming, finché un dirigente fece notare che il 12\% del fatturato
derivava dalle \emph{multe per riconsegna tardiva}, impossibili con lo
streaming, mentre secondo le stime lo streaming avrebbe prodotto il 10\% del
fatturato. ``Non possiamo rinunciare a questo 2\%'': niente streaming. Oggi
non c'è più un Blockbuster ``manco a cercarlo col lanternino''. Non avevano
capito qual era il vero motivo del loro successo, cioè il loro modello di
business\footnote{Le percentuali sono quelle riferite in lezione e non sono
state verificate.}.

\paragraph{The Founder.} Nel film sulla nascita di McDonald's c'è una scena in
cui il futuro socio del protagonista gli dice: ``tu pensi di essere nella
\emph{food industry}, e invece sei nella \emph{real estate industry}''. Non
dovrebbe investire sugli hamburger, ma sui terreni su cui costruisce i
ristoranti. Anche per un team vale lo stesso: bisogna capire che cosa
l'ha reso davvero di successo, oppure no.

\subsection{L'autoanalisi personale}

\paragraph{E quando si cambia azienda?} Uno studente chiede se il debrief sia
possibile anche quando si lascia un'azienda. Non si può tornare in ufficio a
discutere della propria esperienza, ma lo si può fare \textbf{per se stessi}.
Il relatore l'ha fatto, anche se all'epoca in modo poco consapevole, ed è
così che è passato dal \emph{community management} al \emph{people
management}. Analizzando cosa aveva fatto bene a Blizzard si era reso conto
che:
\begin{itemize}
  \item aveva costruito un ottimo rapporto di fiducia con i giocatori (con
        alcuni proprietari di fansite parla ancora oggi);
  \item pur senza la seniority, era diventato un punto di riferimento per il
        team, visto come imparziale ed equilibrato anche su temi che non lo
        riguardavano direttamente;
  \item sapeva \textbf{risolvere problemi}: durante il lancio di
        \emph{Diablo~III} era il più tranquillo di tutti, non perché
        sminuisse il problema, ma perché sapeva che si sarebbe risolto con i
        suoi tempi; il lavoro non era spegnere l'incendio ma ``far sfollare
        la gente'' e aspettare che tutto tornasse in ordine.
\end{itemize}
E aveva capito cosa sapeva fare meno bene, ad esempio \textbf{gestire il
proprio manager}. Avere un manager non significa che il manager gestisca te:
il manager gestisce il tuo lavoro, ma tu devi saper gestire i tuoi
\emph{stakeholder}. Spesso lo stakeholder è il proprio manager; da manager,
gli stakeholder saranno CEO, CFO e così via. Gestirli significa conoscerli,
sapere come portare loro le informazioni, come far presente che qualcosa non
va. Fa parte del lavoro, anche se a lui non piace, e ci ha lavorato.

Nel mese e mezzo di pausa prima di entrare in Riot si è concentrato su
questo: ``cosa porto come mio punto di forza?''. L'autoanalisi è culminata
nel chiedere espressamente alla persona che l'aveva assunto: ``prima mi avete
detto di no, ora mi avete preso: perché?''.

\paragraph{Come si capisce se un progetto è andato bene?} Dipende
dall'\textbf{obiettivo}. Se l'obiettivo era vendere milioni di copie e il
gioco le ha vendute, bisogna capire perché. Se, più ragionevolmente,
l'obiettivo era un gioco da tavolo che permettesse di mettere in pratica le
conoscenze del corso, va valutato per quello: ha permesso di farlo? Se sì,
perché? Se no, perché no? L'obiettivo va fissato \textbf{all'inizio del
progetto, non alla fine}, proprio per poterlo valutare dopo.

\subsection{Quando il team cambia di continuo}

\paragraph{Tutti questi cambiamenti non sono dannosi per il progetto?}
Assolutamente sì, ma quasi mai sono volontari. I progetti andati male sono
spesso progetti grandi.

Il relatore racconta di \emph{Cyberpunk 2077}. Conosceva persone in CD
Projekt Red (``il mondo è piccolissimo, le cose si sanno''). Alla Gamescom va
con Diego Ricchiuti, direttore dell'accademia, a vedere la presentazione del
primo gameplay, quella ``in cui regalavano il giubbottino da samurai''.
Finito il video, tra gli applausi, si guardano e capiscono entrambi che
\textbf{il gioco non è pronto}, a soli quattro mesi dall'uscita prevista:
Diego aveva notato il combattimento e la guida. Gli amici di CD Projekt Red
(uno aveva lavorato con lui a Blizzard per sei anni) gli confermano tra le
righe che ``ce la faremo, ma ci sono ancora cose da sistemare''. Il gioco fu
poi rimandato più volte\footnote{In lezione il relatore dice ``due, forse
tre volte''; l'uscita di \emph{Cyberpunk 2077} è stata effettivamente
rinviata tre volte prima del dicembre 2020.}.

Per deformazione professionale, chi ha sviluppato l'occhio critico non
riesce più a godersi un trailer: in un trailer di \emph{The Day Before} Diego
gli aveva fatto notare che l'effetto della fiammata di un fucile era stato
chiaramente aggiunto in postproduzione e non faceva parte del gameplay. ``È
il triste epilogo di questa storia'': non riescono più a godersi un trailer.

Il punto è che più il team è grande, più queste variabili vengono imposte, e
più è facile che il fallimento \textbf{non dipenda dal fatto che il prodotto
sia brutto} (\emph{Cyberpunk} è un bellissimo gioco sotto molti punti di
vista, e il relatore lo sta rigiocando ora che è stato corretto ed espanso).
Nei grandi gruppi qualcuno arriva, qualcuno si dimette, cambia il budget,
arriva il rinvio; nel caso di \emph{Anthem}, il team ha fatto così tante volte
forming--storming--forming--storming che ne è uscito ``un prodotto che
dovevamo espellere: l'abbiamo espulso, buttato fuori dal nostro organismo pur
di liberarcene''.

Tra i cambiamenti c'è anche la crescita verticale: se il capo viene promosso
perché lavora benissimo, bisogna scegliere un nuovo capo, ed è sempre un
caos (chi ci resta male, chi pensava di meritarlo, chi voleva il suo amico,
chi non sopporta il prescelto). Per questo il relatore è \textbf{contrario
alle promozioni durante i progetti}: si rischia di rovinare tutto il lavoro di
storming.

\subsection{Obiettivi chiari nei piccoli gruppi}

Nei gruppi piccoli, come quelli degli studenti, il problema deriva più
spesso dal \textbf{non avere chiaro che cosa si vuole ottenere}: un prodotto
puramente educativo e formativo, o un prodotto richiesto dall'accademia che,
se viene bene, si potrebbe anche vendere? Un gioco commerciale o un gioco
artistico? Va chiarito subito, perché ci saranno sempre momenti in cui
scegliere tra A e B, e a guidare la scelta sarà la domanda: \emph{quale delle
due opzioni ci porta più vicino all'obiettivo?} Se l'obiettivo non è chiaro,
non si sa cosa scegliere.

Questo porta molte aziende a fare \textbf{scelte di design che in realtà sono
scelte commerciali}. Le microtransazioni sono scelte commerciali: nessuno ha
mai pensato ``sai cosa piacerebbe tanto ai giocatori? Pagare le bustine''.
Prima le bustine si guadagnavano giocando; ora la tendenza è ``potresti
giocare, ma se paghi è meglio''. Si pensi alle polemiche su Ubisoft per
l'opzione di ``saltare il gioco'' pagando, o ai famosi \emph{boost} di livello
di \emph{World of Warcraft}: ``devo pagarti per saltare il gioco?''.

Una studentessa (Silvia) cita \emph{Final Fantasy XIV}, che a un certo punto fu
``resettato'' perché non funzionava: un esempio positivo. Il relatore aggiunge
quello negativo di \emph{Age of Conan}, che non è mai riuscito a fare quel
reset iniziale (``è una schifezza, ve lo ripresentiamo quando è pronto''). Ci
vuole consapevolezza, ma anche i mezzi: un conto è Square Enix, che può
permetterselo, un altro è un'azienda che non può chiudere tutto per un anno
continuando a pagare gli stipendi.

Riassumendo: nei piccoli gruppi i problemi sono più di
\textbf{comunicazione e chiarezza} (obiettivi non chiari o che si spostano);
nelle grandi aziende di più di \textbf{persone che vanno e vengono, obiettivi
e tempistiche che cambiano}. In accademia le scadenze sono fisse (il gioco da
tavolo si consegna a fine anno scolastico, a giugno, e giugno resta); in
azienda può succedere che una feature prevista per la \emph{vertical slice}
di luglio venga anticipata a maggio per portarla alla Gamescom, e un sistema
perfetto per consegnare a giugno può non esserlo più per consegnare a maggio.

\subsection{L'ambiente determina la qualità del prodotto}

\paragraph{È vero che chi passa per Electronic Arts finisce per fare brutti
giochi?} Uno studente cita BioWare, che faceva ottimi giochi finché era
autonoma. Il problema non è dei developer, che non sono ``improvvisamente
diventati incapaci'': è evidente che il problema è altrove, e che
\textbf{l'ambiente in cui si fa un gioco determina molto spesso la qualità del
prodotto}.

Un esempio locale positivo: Tommaso Bonanni, conosciuto dal relatore insieme a
Diego e già docente dell'accademia, ha fondato a Roma Caracal Games Studio,
applicando fin dall'inizio questi concetti di cultura aziendale. L'azienda è
rimasta piccola, ma ha molto chiaro cosa vuole fare e come. Quando si
predispone fin dall'inizio l'azienda e il team a lavorare in un certo modo è
fattibilissimo; è molto più difficile quando si è EA, si acquisisce
un'azienda e le si dice ``da oggi si fa così''.

Il relatore ha vissuto la transizione da Blizzard ad Activision Blizzard (uno
dei motivi per cui se n'è andato, volontariamente). Ha visto la cultura
interna cambiare per gradi:
\begin{enumerate}
  \item ``facciamo giochi fighissimi, perché i giochi devono essere
        fighissimi'';
  \item ``facciamo giochi fighi così vendono di più'';
  \item ``vendiamo di più facendo giochi fighi'' (sembra uguale, ma è molto
        diverso);
  \item ``facciamo più soldi''.
\end{enumerate}
Lo si capiva dalle scelte di design, osservando quale fosse il peso più
grande sulla bilancia. Lo stesso varrà per gli studenti, nei loro giochi e nei
colloqui di lavoro: capire subito che cosa vuole l'altra parte.

\paragraph{Leader e visione.} Nei gruppi dell'accademia, dal secondo anno, c'è
una persona che assume il ruolo di \emph{lead} o \emph{scrum master}. È
diverso dall'essere un \emph{game director} (``il gioco si fa come dico io,
aiutatemi a farlo con la mia visione''), come avviene nei grandi studi: da lì
nascono i nomi nei titoli, come \emph{Sid Meier's Civilization}, che indicano
la persona che ha fatto quel gioco, oppure un puro branding, come \emph{Tom
Clancy's Rainbow Six}.

\paragraph{Inclusion come jolly.} A margine, uno studente osserva che in molte
aziende ``la risposta giusta è sempre \emph{inclusion}'', usata come un
jolly. È comunissimo, conferma il relatore: \emph{diversity and inclusion}
viene spesso usata per dire ``così includiamo più clienti'', ma non è quello
il suo punto (``altro discorso'').

% ---------------------------------------------------------------------
\section{Confrontarsi con chi è più bravo}
\label{sec:z1-rivalita}

Arrivato a Riot dopo sei anni a Blizzard, il relatore pensava di avere un buon
bagaglio di esperienza. Si è trovato invece ``come un fante buttato in mezzo
ai Navy SEAL'': persone che avevano superato 11, 12, 13 colloqui, con un
livello molto alto.

Individua subito un collega francese che, nella sua testa, diventa ``la
versione migliore di me'': stesso ruolo (community manager), esperienza
simile (arrivava da Ubisoft, come quasi tutti i francesi di Riot), ma molto
più spigliato, sempre con la battuta pronta, ``l'anima della festa''. Il
relatore invece è introverso: un workshop come questo per lui è
\emph{emotivamente drenante}, non ama parlare in pubblico pur amando
condividere le proprie esperienze, e non riesce a sostenere conversazioni
come ``cosa hai fatto questo weekend?''. Lavorando raramente insieme, aveva
sviluppato una rivalità ``da allenatore di Pokémon'' (lui era Ash, il
relatore l'antagonista).

Dopo sei mesi, a un evento di \emph{team building} aziendale (a Tenerife), il
suo capo, molto attento a queste dinamiche, li mette apposta nello stesso
gruppo. A fine serata, complice l'open bar, il relatore gli confessa di averlo
mal sopportato dal primo giorno per invidia della sua estroversione. Il
collega ride e risponde che pochi giorni prima aveva detto al capo la stessa
cosa al contrario: invidiava a Stefano il suo essere sempre organizzato, sul
pezzo, presentarsi nella sala riunioni giusta (mentre lui finiva in quella
sbagliata), arrivare con gli appunti, con i dati a supporto e le analisi di
mercato, mentre lui spesso non aveva nemmeno letto la mail del meeting.
Un'invidia in due direzioni, che li ha portati a lavorare molto insieme e a
migliorarsi a vicenda.

Le lezioni:
\begin{itemize}
  \item non bisogna demoralizzarsi vedendo qualcuno che fa lo stesso lavoro
        ed è più avanti sotto alcuni aspetti: spesso c'è un bilanciamento. I
        ``punti caratteristica'' sono limitati: se li ha messi in carisma,
        non li avrà messi in saggezza;
  \item conviene \textbf{chiedere a queste persone come mai} sono così brave
        in qualcosa, come si fa con un progetto andato bene: magari non lo si
        farà mai bene quanto loro, ma si comincerà ad avere il
        \emph{mindset} giusto; e viceversa, loro saranno spinti a chiedere a
        noi;
  \item spesso \textbf{le proprie doti le scopriamo dagli altri}. Il relatore
        non si era mai accorto di essere ``quello che resta calmo durante le
        emergenze'' (il lancio di \emph{Diablo~III}, quello di
        \emph{Draenor}, i problemi dal vivo alla Games Week) finché non
        gliel'hanno fatto notare tutti: per lui era la normalità. A chi pensa
        ``non sono capace, forse non è il mio percorso'', può darsi che
        semplicemente nessuno gli abbia ancora detto ``questa cosa la sai
        fare davvero bene, concentrati su questo''.
\end{itemize}

% ---------------------------------------------------------------------
\section{Domande sul lavoro nell'industria}
\label{sec:z1-industria}

Diverse domande degli studenti hanno riguardato le condizioni di lavoro nei
grandi studi internazionali. Le risposte sono raccolte qui per argomento.

\subsection{Presenza, remoto e trasferimenti}

\paragraph{Grandi studi: quasi solo in presenza.} Sia a Blizzard sia a Riot si
lavorava \textbf{solo in presenza}, senza smart working: sei anni a Parigi e
cinque a Dublino. Si parlava solo inglese: anche nei meeting tra soli
italiani si tendeva a comunicare e a scrivere in inglese, per correttezza
verso chi si fosse aggiunto e per lasciare una traccia comprensibile a tutti.
Il francese serviva per fare la spesa e sopravvivere, non in ufficio.
L'ambiente era estremamente eterogeneo (francesi, polacchi, spagnoli e
molti altri): durante il colloquio dal vivo a Dublino, mentre si preoccupava
di essersi vestito in modo adeguato, vide passare dietro la parete a vetro un
collega svedese con i dreadlock, a piedi scalzi, con una pianta sotto un
braccio e un aquilone sotto l'altro, e capì di essere ``a posto così''.
Un'esposizione a molte culture e modi di pensare che il relatore considera
bellissima.

Anche oggi, tranne alcuni studi particolari, la maggior parte del lavoro si
svolge di persona: amici a Paradox lavorano a Stoccolma, i pochi rimasti a
Blizzard si sono dovuti trasferire a Irvine. Per posizioni \emph{junior} o
\emph{entry level} i grandi studi raramente concedono il lavoro da casa; per
profili senior può esserci qualche eccezione (ad esempio nove mesi da remoto
e tre in sede).

\paragraph{Relocation.} Spesso tocca spostarsi, ma il relatore non ha mai
pagato un trasloco di tasca propria: queste aziende offrono la
\emph{relocation} (aiuto a trovare casa, ad aprire il conto in banca e
supporto pratico), perché è un investimento sulla persona.

\paragraph{Colloqui mentre si lavora altrove.} Per i colloqui in presenza
prendeva ferie, senza dire dove andava. Il consiglio generale: per fare
colloqui \textbf{prendete ferie o permessi}, non datevi malati, ``che non è un
bel modo di comportarsi''.

\paragraph{Remoto ``doppio''.} Un collega lavora per l'azienda coreana che
gestisce \emph{PUBG}, con ufficio ad Amsterdam, ma lavora da Milano: un
``doppio remoto''. Si trova, ma in genere un ambiente anglofono con sedi in
Europa chiederà di trasferirsi in Europa.

\paragraph{Famiglia e trasferimenti.} Il relatore ha avuto la fortuna di
sposare una persona che fa lo stesso lavoro (conosciuta lavorando su
\emph{Warhammer Online}, poi con lui a Blizzard e a Riot, nella
localizzazione e nel doppiaggio, ad esempio di \emph{World of Warcraft} in
italiano): avendo competenze rare, spesso venivano assunti entrambi, al
massimo a pochi mesi di distanza. Per i developer la soluzione frequente è
che la famiglia segua la persona assunta, perché lo stipendio basta per
tutti. Ma non sempre c'è equilibrio, e a volte uno dei due deve sacrificarsi:
sua moglie ha lasciato Blizzard per seguirlo a Dublino (per poi essere
assunta da Riot); lui ha lasciato Riot per tornare in Italia quando lei ha
trovato lavoro qui. Ha dovuto rinunciare più volte a lavorare in Islanda per
CCP (\emph{EVE Online}), perché non era una scelta solo sua. Si fanno
compromessi, ma si riesce a far quadrare.

\paragraph{Ci si sposta per ogni progetto?} No: si sono spostati
esclusivamente per cambiare azienda. Cambiando progetto nella stessa azienda
è raro, quasi mai (``non l'ho mai visto succedere''), che venga chiesto di
trasferirsi. In Francia stavano benissimo (vivevano fuori Parigi; ``i
problemi con i francesi sono in realtà problemi con i parigini''), parlavano
francese ed erano vicini alla doppia cittadinanza: se ne sono andati perché
non vedevano un futuro a Blizzard, per il modo in cui l'azienda trattava i
colleghi.

\paragraph{Il mercato oggi.} Non esiste un \emph{silver bullet} per trovare
lavoro domani dove si vuole. Con tutti i \emph{layoff} recenti ci sono
migliaia di esperti del settore che cercano lavoro: serve pazienza, ma è
fattibile.

\paragraph{Studi indie e remoto.} Negli studi più piccoli il remoto è più
diffuso: a Stormind circa 60 persone su 95 lavorano da remoto, ma lo studio
si è strutturato fin dall'inizio per farlo. Serve consapevolezza. Bending
Spoons (che sviluppa app per mobile, non giochi, ed è l'azienda dell'app
Immuni), pur essendo grandissima e con sede a Milano, lavora da remoto e
lascia l'ufficio a chi vuole: lo smart working è posizionato quasi come un
\emph{benefit}, ma permette anche di avere un ufficio da 300 persone invece
che da 700. Le sedi fisiche costano decine di milioni all'anno (Blizzard ha
chiuso le sedi di Amsterdam e di Parigi). Tuttavia, per lavori molto
verticali e specifici gli studi preferiscono avere le persone dal vivo; i
ruoli più spesso remoti sono quelli che non fanno sviluppo (customer service,
PR\dots).

Anche gli artisti lavorano spesso in presenza: l'accademia può dare solo una
\emph{simulazione} di quanto stretto sarà il contatto tra programmazione,
design e art. ``Sarà ancora di più.''

\paragraph{La comunicazione da remoto.} È più difficile. Pur essendo
estremamente introverso, il relatore \textbf{obbliga tutti a usare la webcam}
nei meeting in cui conta: non in un aggiornamento con novanta persone
collegate, ma certamente in un colloquio a due sulla performance dell'ultimo
anno. La componente fisica è una parte enorme della
comunicazione\footnote{In lezione si parla di ``quasi l'ottanta per cento'';
il dato va preso come indicativo, poiché le percentuali di questo tipo
circolano in forme diverse e sono spesso citate fuori contesto.}: anche se si
è timidi, accendere la webcam è l'approssimazione più vicina all'essere dal
vivo. Negli ambienti dove si lavorava senza webcam la situazione era
``esponenzialmente peggiore''.

\subsection{Differenze culturali nel team}

\paragraph{Gestire sensibilità diverse.} A Riot i team erano divisi per
\emph{lingua}, non per nazione: il team di lingua spagnola includeva anche un
messicano, quello di lingua russa includeva russi e ucraini, proprio nel
periodo della guerra in Crimea. Tra i colleghi non c'era alcun problema (erano
amici), ma ci si è posti la domanda: \emph{come comunichiamo con una community
che sarà divisa su questo tema?} Nei giochi online la divisione era evidente:
in \emph{Rocket League} c'era chi, con un compagno di squadra con la bandiera
avversaria, si faceva otto autogol pur di farlo perdere. Allo stesso modo,
con molti americani mandati a Dublino a coprire posizioni scoperte, c'erano le
elezioni (Clinton--Trump, Biden--Trump) e varie guerre: serviva sempre una
certa sensibilità.

Imparare a navigare queste situazioni lo ha reso un professionista migliore:
gli ha insegnato a fermarsi e chiedersi ``io la percepisco così; gli altri
dieci, con background, storie e lingue diverse, come la penseranno?'', a
ragionare più \emph{long term} e più \emph{big picture}, con più empatia.

\paragraph{Stili di comunicazione.} Un partecipante che ha lavorato nel Regno
Unito chiede dello stile comunicativo italiano, a volte frainteso. Uno dei più
grandi complimenti ricevuti dal relatore è stato ``non sembri italiano''
(``gesticoli poco per essere italiano'', ``parli bene inglese per essere
italiano''). Ha imparato ad \textbf{adattare la comunicazione alla persona}:
con spagnoli e francesi la battuta sulla nazione funziona, perché c'è un
``odio--amore'' condiviso; con russi e giapponesi no. Doveva adattare anche
l'umorismo. Un collega metà russo e metà messicano diceva di aver preso ``il
lato peggiore di entrambe le nazionalità'', ed è stato un ottimo maestro per
imparare a leggere culture diverse.

Non bisogna generalizzare, ma gli stereotipi possono essere un punto di
partenza ragionevole, \emph{sapendo che troverete sempre le eccezioni}. E
vanno letti correttamente: gli americani tendono a essere più estroversi e ad
avere sempre qualcosa da dire non perché ``sono fatti così'', ma perché
vengono educati così (all'università tra i corsi c'è il \emph{debate}, saper
argomentare in pubblico); un collega giapponese può essere più introverso per
ragioni culturali. Nel dubbio si parte da lì, poi si conosce la persona.

\subsection{Negoziare, discriminazione e stipendi}

\paragraph{La famiglia come leva.} Avere un partner riduce le opzioni, ma le
aziende lo sanno: uno spostamento è per due. Si può usare come leva per
negoziare (un compenso diverso, un supporto per il trasloco): se un'azienda ha
deciso di volere quella persona, 2, 3, 4 o 10.000 euro in più all'anno, per
un'azienda che fattura centinaia di milioni o miliardi, non cambiano nulla.

\paragraph{Essere sposati è un motivo per non essere assunti?} No. In alcuni
paesi è proprio illegale discriminare su questo: in Francia era illegale
chiedere CV con data di nascita, foto, sesso o religione, e chiedere se si
vogliono figli. Purtroppo, dai racconti dei suoi studenti di master (in
settori come \emph{beauty and wellness} e \emph{food and wine}), in piccole
realtà italiane queste domande (``hai 25 anni, hai intenzione di
sposarti?'') vengono ancora fatte, e servono a filtrare i candidati. Nelle
grandi aziende ci sono sistemi per evitarlo; anzi, molte danno \emph{più}
benefit a chi ha famiglia (``porta pure la famiglia, è tutto pagato''),
perché le persone competenti hanno spesso una certa età e quindi,
ragionevolmente, una famiglia.

\paragraph{Gli stipendi.} Dipendono molto da dove si vive. Lo stipendio a Riot
non era paragonabile a quello di Blizzard perché in Francia era tassato al
52\% e in Irlanda al 17\%\footnote{Aliquote come riportate in lezione,
riferite alla situazione personale del relatore.}: nessuno stipendio
competitivo può compensare il 35\% di differenza di tasse. Negli Stati Uniti,
con pochissima previdenza pubblica, gli stipendi sono molto più alti perché
poi ci si paga tutto (``se hai un raffreddore paghi tu''). Bisogna valutare il
potere d'acquisto di quei soldi nel paese in cui si andrà a vivere.

\subsection{Freelance e \emph{work for hire}}

All'inizio della sessione pomeridiana uno studente che lavora nell'audio
chiede quali siano le condizioni per chi lavora come \emph{freelance}.
\begin{itemize}
  \item Dipende molto dal progetto: spesso il contratto è a progetto, per
        portare a termine una parte o un intero progetto.
  \item Non tutte le aziende usano freelance: Riot e Blizzard lo facevano
        pochissimo e preferivano assumere (magari trattando poi la persona
        ``come un freelance'', lasciandola a casa a fine progetto).
  \item Spesso i lavori offerti come freelance sono in realtà
        \textbf{sostitutivi di un full time} (``freelance, ma sei ore al
        giorno'').
  \item Si può fare, ma bisogna essere molto chiari nel definire l'impegno
        (quante ore, come vengono tracciate): non tutti specificano cosa
        significa davvero ``freelance''. Lavorando così per due o tre aziende
        si può arrivare a un tempo pieno.
  \item È più comune il \emph{work for hire}: aziende che forniscono personale
        ad altre aziende; si è assunti da X ma si lavora ``embeddati'' presso
        Y (una sorta di subappalto). Anche lo studio del relatore lo fa.
  \item Il puro freelance di quattro mesi esiste, ma di solito per quei
        quattro mesi diventa un lavoro a tempo pieno; è raro ``quattro mesi,
        due ore al giorno''. Nelle grandi aziende è improbabile lavorare nello
        stesso mese per Ubisoft, Riot e Blizzard tre ore ciascuna; è più
        realistico lavorare per più startup indie italiane, che non possono
        permettersi un'assunzione ma possono pagare due ore al giorno.
  \item Molto passa dai contatti: trovare le aziende giuste e rendere esplicite
        le condizioni.
\end{itemize}

% ---------------------------------------------------------------------
\section{Il circolo virtuoso: feedback, miglioramento, fiducia}
\label{sec:z1-triangolo}

\subsection{Il triangolo}

L'obiettivo a cui il workshop vuole arrivare è un \textbf{triangolo che si
autoalimenta}:
\begin{enumerate}
  \item si dà un \textbf{feedback} con l'obiettivo di migliorare le
        performance (inclusi elementi caratteriali);
  \item la persona \textbf{migliora} e vede di averne beneficiato;
  \item questo genera \textbf{fiducia}, che rende più facile dare e ricevere
        il feedback successivo.
\end{enumerate}
Togliendo uno qualsiasi dei tre elementi, il circolo non è altrettanto forte.

In qualsiasi fase si parta, il rispetto per i colleghi è inizialmente dettato
dall'\emph{educazione}, non dalla fiducia. La fiducia arriva quando i colleghi
cominciano a dire cose giuste e sensate e a \textbf{mantenere la parola}.
Dare e ricevere feedback è il modo in cui, in un certo senso,
``inganniamo il cervello'' per creare una \emph{case history} di fiducia: si
ha motivo di fidarsi di più, di ascoltare di più, di essere più ricettivi. E
si diventa sempre più bravi ad accettare il feedback e ad agire di
conseguenza.

\subsection{Accettare un feedback non vuol dire fare come dice l'altro}

Dieci anni fa il relatore non era ostile ai feedback, ma molto
\emph{difensivo}: ``sì, sì, va bene, ora ci guardo''; oppure, a chi gli diceva
``avresti dovuto farla così'', rispondeva spiegando perché l'aveva fatta in
quel modo. Oggi ha imparato a dire ``ok, grazie, ora ci ragiono'': non deve
difendere la propria posizione da chi viene solo a dire ``hai pensato a
questa cosa?''. Prima lo viveva come un attacco (``mi viene a dire che non so
fare il mio lavoro''), ora come un ``ottimo punto di vista, lo prendo in
considerazione''.

Con una postilla fondamentale: \textbf{accettare il feedback non significa
fare come dice l'altro}. ``Mi hai detto che preferisci questo mob verde. Ho
pensato anche a farlo verde? Sì. Penso comunque che giallo sia meglio? Sì, lo
faccio giallo.'' Il relatore accetta \emph{sempre} il feedback, valido o no:
garantisce di tenerlo a mente quando lavora. Che poi faccia come gli viene
detto è un altro discorso.

\subsection{Fiducia e scadenze}

La fiducia funziona allo stesso modo con gli impegni. Se si chiede a qualcuno
una cosa per venerdì e la consegna venerdì, la volta dopo gliela si chiederà
per venerdì. Se invece arriva il giovedì della settimana dopo, la volta dopo
gli si darà una \textbf{deadline anticipata}, come con chi arriva sempre in
ritardo agli appuntamenti (``se ti dico alle 7 arrivi alle 7 e mezza; la
prossima volta ti dico alle 6 e mezza'').

Al lavoro questo è dannosissimo, perché crea \textbf{deadline fittizie} che
non sono mai quelle vere, e nascono discussioni del tipo ``perché devo farlo
in due giorni se si fa in cinque?'' (``sì, se tu lo facessi in cinque''). La
conversazione giusta è un'altra: \emph{``che problemi hai avuto nel
consegnare venerdì?'' ``Questo, questo e questo.'' ``Cosa posso fare per
evitare che succeda?''}

Il relatore non vuole mai sentirsi dire ``non sono riuscito a farlo''.
Preferisce di gran lunga sentirsi dire, tre giorni prima, \textbf{``potrei non
riuscire a farlo, per questi motivi''}. Chi ha la visione d'insieme può dare
più budget, più persone, spostare la deadline, cambiare lo scopo del progetto:
può fare qualsiasi cosa per aiutare. Ma se lo si scopre il venerdì, cinque
minuti prima della deadline, non può fare più niente. Il messaggio è:
\textbf{``aiutami a poterti aiutare''}. Non penserà mai male di chi dice ``non
ce la faccio per questi motivi'', perché non sta dicendo ``non ho voglia'': sta
indicando i problemi, e lavorare su quei problemi è compito del manager. Non
può invece fidarsi di chi fino al giorno prima dice ``tranquillo, ce la
faccio'' e il giorno dopo ``non ce l'ho fatta''.

Questo è fondamentale per chi lavora in programmazione, design o art, perché
ci saranno sempre deadline: bisogna sempre \textbf{lavorare sulle aspettative},
dire al proprio manager (o al proprio report) cosa è lecito aspettarsi. Non
bisogna fare come i pinguini di \emph{Madagascar} (``sorridete e salutate, va
tutto bene''): se c'è un problema, va detto, altrimenti si peggiora solo la
propria situazione.

Da italiani si tende a non voler mai alzare la mano per chiedere aiuto
(``no, no, ce la facciamo, lo faccio io''). A volte non ce la si può fare, non
perché non si è capaci, ma perché vari fattori lo impediscono. Accettarlo è
importante: \textbf{è molto meglio avvisare e chiedere una mano che dire
``non ce l'ho fatta''}.

\subsection{Domande sul dare e ricevere feedback}

\paragraph{``Ne parliamo a fare, se poi fai come vuoi?''} Uno studente (Cristian)
riceve feedback che, dopo averli analizzati, non ritiene congrui; quando non
li segue gli viene detto ``ne parliamo a fare?''. C'è un problema di
\textbf{aspettative}: l'idea che accettare un feedback voglia dire fare come
dice l'altro. Accettarlo significa prendere in considerazione ciò che è stato
detto. Chi decide diversamente ha però un \textbf{onere}: argomentare perché
si fa comunque a modo proprio (``mi hai detto verde, la facciamo gialla
perché\dots''). Poi bisogna solo aspettare i risultati per vedere chi aveva
ragione.

La decisione spetta sempre \textbf{a chi ha la responsabilità di quella
decisione}. Se il game director deve scegliere tra A e B, anche se il
feedback dice B, la scelta è sua perché la responsabilità è sua. Non sarebbe
corretto essere ritenuti responsabili di una scelta delegata a chi dà
feedback. La risposta può essere semplicemente: ``l'ho tenuto in
considerazione, ma visto che ne sono responsabile io e le conseguenze sono
mie, ho scelto di fare X''.

\paragraph{E se l'altro fa ripicca?} Un partecipante chiede di chi, vedendo
il proprio feedback non accolto, si vendica ignorando a sua volta feedback
buoni. Per il relatore è un problema caratteriale: chi si comporta così ``è
già sulla strada per uscire''. Non c'è spazio per l'ego in team di questo
tipo. Il fenomeno nasce dall'idea che \emph{il feedback negativo faccia fare
brutta figura}, che il fatto che si potesse fare meglio tolga qualcosa alle
proprie capacità.

Non è così: \textbf{il feedback non è un giudizio}. Si può passare un anno a
Riot ricevendo solo feedback negativi e avere una valutazione annuale
stellare, perché ricevere molti feedback negativi significa essere una persona
approcciabile, che agisce sul feedback e invita gli altri a darlo. Il relatore
ha visto valutazioni molto peggiori per chi riceveva solo feedback positivi,
perché la gente aveva paura di dargliene di negativi: ma solo ``pacche sulle
spalle'' non fanno migliorare. Il problema va affrontato parlando con la
persona; a volte le conversazioni iniziano con \textbf{``ti do un feedback sul
feedback''}, cioè sul modo in cui l'altro ha dato il suo.

Esempi dalla sua esperienza:
\begin{itemize}
  \item lui chiede spesso di \textbf{non ricevere feedback senza esempi
        tangibili}, perché fa fatica ad astrarre: ``sei stato maleducato'' non
        significa nulla, ``ieri al meeting mi hai risposto male su questa
        cosa'' sì;
  \item ad altri serve l'opposto: ha ricevuto il feedback ``quando devi dirmi
        che ho sbagliato, puoi girarci un po' intorno e darmi del contesto?'',
        perché il suo stile diretto (``questa cosa non ha funzionato, dovremmo
        fare così'') feriva. Ha adattato il suo modo di dare feedback: con
        alcuni è diretto, con altri gira intorno. L'importante è saperlo.
\end{itemize}

\subsection{Non portare la soluzione}

Un consiglio centrale: \textbf{quando si dà un feedback, non si dà una
soluzione}. Non si dice ``io farei così'', ma ``secondo me questa cosa si
potrebbe fare diversamente, per questo e questo motivo: come potremmo farla
diversamente?''. È come chi smette di fumare perché vuole e chi smette perché
glielo dicono moglie e figli: se non si è convinti, spesso non si smette.

Quando il relatore diceva ``c'è questo problema, io farei così'', il meeting
successivo iniziava inevitabilmente con ``ci ho provato come hai detto tu, con
me non ha funzionato''. Sempre. Perché la soluzione era sua. La
responsabilità della soluzione deve essere invece \textbf{congiunta}: ogni
feedback deve creare un piccolo team che ha come obiettivo risolvere quel
problema. Questo toglie l'antagonismo (non ``ti dico che non sai fare il tuo
lavoro'', ma ``ti voglio aiutare a farlo meglio, lavoriamoci insieme'') e
sposta il focus: il problema non sei tu, è ciò che sta accadendo.

\paragraph{E se sono i piani alti a tenere all'oscuro dei problemi?} È una
questione di cultura aziendale. Il relatore fa sempre presente al proprio capo
il concetto di \textbf{service level agreement} (che verrà ripreso nella
seconda giornata): ``se mi dai questi input, ti garantisco questo output in
questi tempi''. Se il capo chiede delle cose senza dire budget, tempi e
perché, non può aspettarsi lo stesso risultato. È una comunicazione a due
vie: si dice al manager cosa serve per fare il lavoro, e lui deve fornirlo o
accettare che il risultato non sarà quello voluto.

\subsection{Amici, parenti e parole proibite}

\paragraph{``Non si lavora con amici e parenti''?} Per sua definizione
personale il relatore ha pochissimi amici (``stanno sulle dita di una mano'')
e una soglia molto alta per considerare amico un collega. Allo stesso tempo,
in quindici anni e con centinaia di colleghi, ha avuto solo \textbf{due}
esperienze in cui non è riuscito a lavorare con qualcuno: non serve essere
amici per far funzionare le cose. Con gli amici, anzi, gli viene più facile:
se dice a un amico che una cosa ``fa schifo'', l'amico non pensa ``come si
permette'', ma ``se perfino Stefano, che è mio amico, me lo dice, allora è
proprio così''. Diventa un rafforzativo, grazie alla fiducia costruita prima.
Il problema del lavorare con amici e parenti nasce dal \emph{non voler} dire
all'amico che sta sbagliando; se invece la confidenza permette di dirlo, si
può saltare la ``confezione'' del messaggio e passare direttamente al
contenuto.

\paragraph{Le frasi peggiori.} Alla domanda se esistano parole peggiori di
``ottimo lavoro'', il relatore indica la sua classifica: ``ottimo lavoro''
(detto in modo vuoto)\footnote{Il passaggio della trascrizione è
frammentario: si intende verosimilmente il complimento generico e di
circostanza, che non dice nulla di utile.}; poi la frase che lo ``fa uscire di
testa'', \textbf{``abbiamo sempre fatto così''} (``e infatti avete chiuso'',
come Blockbuster); infine \textbf{``siamo una famiglia''}, riferito al posto di
lavoro. Un'azienda non è mai una famiglia, perché non ci si può aspettare lo
stesso impegno da chi la possiede e da chi ci lavora: non c'è un
corrispettivo proporzionale tra impegno e ritorno. Per quanto inclusiva,
flessibile e razionale, non sarà mai una famiglia.

\subsection{Ego, muri e il peggior modo di scegliere un leader}

\paragraph{Chi mette davanti l'ego e alza muri.} Esistono, ma di solito durano
pochissimo in azienda. Non si tratta di avere un ego forte, ma di \textbf{non
saper ricevere feedback}. Spesso sono persone bravissime nel loro lavoro che
poi si chiedono perché tutti intorno diventino più bravi di loro: perché gli
altri accettano feedback e hanno più fonti da cui attingere, mentre loro si
sono rintanati nella loro torre.

Molti manager tendono ad \textbf{assumere persone un filo peggiori di loro},
per non fare brutta figura. Il relatore assume solo persone più brave di lui,
``perché non è una gara'': ha lavorato con molti community manager più bravi
di lui come community manager, ma lui era più bravo a fare il manager. Non
viene valutato su quanto sa fare il community manager, ma su quanto il suo
team sa farlo. Un partecipante chiede se la paura di avere qualcuno di
migliore esista anche in ambito accademico. Assolutamente sì: al relatore è
capitato di non essere preso come docente in un corso di formazione sui
videogiochi perché il responsabile aveva percepito che ne sapesse più di
lui\footnote{La trascrizione dice ``che non sapevo più di lui'', ma il senso
del passaggio richiede il contrario.}. ``Ma perché te ne importa? Tu sei quello che ha portato chi
sa fare il suo lavoro: ben venga.'' O si ha questa cultura, o si ha paura per
il proprio posto.

\paragraph{Promuovere il migliore.} Quando bisogna scegliere un leader per un
team (programmazione, design, art), molto spesso si sceglie \textbf{il migliore
tecnico}: il miglior programmatore, il miglior designer. È spesso la scelta
sbagliata: è come prendere un idraulico per fare l'elettricista. Fare il
\emph{people manager} non è come fare il designer o il programmatore; si può
essere un ottimo programmatore e un pessimo manager, e viceversa. Le aziende
l'hanno capito e lo fanno meno di prima. A Blizzard c'era un grave problema di
\emph{middle management}: molte persone in ruoli importanti erano state
promosse cercando le skill sbagliate, erano frustrate dal non saper fare il
proprio lavoro e schiacciavano chi stava sotto per paura di fare brutta
figura.

Un partecipante cita una teoria simile: una persona viene promossa finché non
raggiunge \emph{il proprio livello di incompetenza}, dove si blocca\footnote{È
il \emph{principio di Peter} (Laurence J.\ Peter, 1969), citato in lezione
come ``la teoria di Peter''.}. Un altro cita Michael Scott di \emph{The
Office}, su cui il relatore tornerà (sezione~\ref{sec:z1-quando}).

% ---------------------------------------------------------------------
\section{Perché la cultura conta: retention, costi e fiducia}
\label{sec:z1-cultura}

\subsection{L'ambiente di lavoro e la retention}

Perché tutto questo è così importante? Perché l'ambiente di lavoro conta
moltissimo per la \textbf{retention}, cioè per trattenere le persone e farle
lavorare in modo efficiente. Se contasse solo lo stipendio, molti non
lascerebbero mai Riot o Blizzard; e più si va avanti con l'età, più si impara
a dare valore anche al resto.

Inoltre, secondo gli studi citati nelle slide, chi vive il proprio lavoro in
modo negativo è \textbf{meno efficiente nel problem solving}: l'energia
mentale va verso ``come faccio a non farmi dare la colpa'' o ``come faccio a
fare bella figura'', non verso la soluzione del problema. In ambienti negativi,
dove ci si deve preoccupare della propria ``sopravvivenza'', si lavora
peggio.

\subsection{Il costo di chi se ne va}

Ogni volta che qualcuno si dimette o viene licenziato, il costo di trovare
una nuova persona, fare i colloqui, formarla e farle attraversare forming,
storming, norming\dots\ è circa \textbf{un terzo del suo stipendio}. Per questo
le aziende tendono a fare più \emph{retention} che \emph{acquisition}, almeno
con le figure considerate utili alla produzione.

Allo stesso tempo si vedono sempre più spesso \emph{layoff} e \emph{hiring
wave}: si licenziano 3.000 persone e sei mesi dopo se ne riassumono 2.500 o
3.500, ma tutte più junior. Il relatore ne ha scritto in un articolo su
Eurogamer.it, in cui spiega anche alcuni ``trucchi del mestiere'': perché i
licenziamenti avvengono sempre in certi periodi e sempre ``a gruppi'' (prima
Bethesda, poi Blizzard, poi un'altra\dots).

\paragraph{L'offerta di Riot.} A Riot esisteva una pratica (di cui il relatore
non ricorda il nome) per cui, firmato il contratto, si aveva \textbf{un anno
di tempo per dire ``non mi piace, non me la sento, me ne vado''} e ricevere
subito come buonuscita un anno di stipendio. Avevano stimato che fosse molto
meglio pagare subito un anno a qualcuno e lasciarlo andare che tenerlo due,
tre, quattro anni come \emph{underperformer}: la responsabilità della scelta
di assumere era dell'azienda\footnote{La pratica, nota come ``The Offer'', è
descritta in fonti pubbliche con condizioni diverse (un periodo più breve e
una somma pari a una frazione dello stipendio annuo). Si riportano le
condizioni come raccontate in lezione; andrebbero verificate.}. L'hanno usata
pochissimi, ma il solo fatto che esistesse faceva lavorare meglio anche lui da
manager, perché sapeva che chi non si trovava bene poteva andarsene con un
anno di tempo per trovare un altro lavoro.

\paragraph{Il ricircolo.} C'è un forte ricircolo nell'industria: molti ex
colleghi passano da Ubisoft a Paradox, da Paradox a Funcom, da Funcom a
Blizzard. Il relatore ha scelto di uscirne lavorando solo per aziende
\textbf{indipendenti}: Stormind non è un indie di cinque persone, ma è
indipendente, e non deve accontentare azionisti. Quando vede un'azienda con
utili miliardari licenziare migliaia di persone si chiede ``dov'è la linea''
(\emph{when enough is enough}): con azionisti che vogliono un ritorno, se
quello è l'obiettivo stanno facendo la cosa giusta, ma lui non condivide
l'obiettivo\footnote{In lezione si parla di ``6,5 miliardi di utile'' e di
``3.000 persone licenziate''; le cifre sono riportate come dette e non sono
state verificate.}.

\paragraph{L'open space vuoto.} A Riot, durante un'ondata di licenziamenti,
del suo team di trenta persone è rimasto solo lui: ``teniamo te come
coordinatore, poi ricostruirai il team''. Ha rifiutato di ricostruire un team
cancellato ``con un colpo di spugna'' e ha deciso che non avrebbe più lavorato
per una grande azienda. Nessuno aveva pensato di spostarlo, e ha lavorato per
tre mesi in un open space di cento metri quadri con la sua scrivania e altre
ventinove vuote. Su una lavagna aveva scritto i nomi dei suoi ragazzi e li
cancellava man mano che, anche grazie ai suoi \emph{referral}, trovavano
lavoro. Cancellato l'ultimo nome, si è dimesso: per lui il lavoro era finito.

\paragraph{Team che si dissolvono.} Uno studente cita un articolo su Larian
Studios, il cui management è contrario a ricostruire ogni volta il team a fine
progetto. È lo stesso problema: si pensa ai ``giochi Blizzard'' o ai ``giochi
BioWare'', ma le persone che li hanno fatti non ci sono più; è rimasto solo il
brand, e non ha senso avere le stesse aspettative da un gioco fatto da persone
diverse. E anche chi sopravvive ai layoff vive pensando che il prossimo
potrebbe essere lui: alcune persone, licenziate nell'ultima ondata a
Blizzard, gli hanno detto ``erano due anni che lo sapevo''. Vivevano a Cork,
``un paesotto irlandese dove o vai per lavorare a Blizzard o non ci vai'',
avevano comprato casa: non ci si può aspettare chissà che morale.

\subsection{La cultura è forte quando si lavora per gli altri}

\begin{quote}
  \emph{La cultura è forte quando le persone lavorano per gli altri e a favore
  degli altri, e non per se stesse.}
\end{quote}

Se si dà feedback per far migliorare l'altro e non per imporre il proprio
modo, se si aiuta per il bene del progetto e non per fare bella figura, la
cultura si rafforza. È dai \textbf{piccoli gesti} che emerge la cultura
aziendale: ad esempio, dal non dover chiedere ``perché mi stai aiutando?''.

\paragraph{L'incognita.} \textbf{L'incognita è spesso dannosa quanto le
cattive intenzioni.} In un esperimento (``mi pare a Stoccolma, cinque o sei
anni fa'') si distribuivano per strada banconote da 500 euro, e la maggior
parte delle persone le rifiutava chiedendosi ``dov'è la fregatura?''
\footnote{L'esperimento è riportato in modo vago in lezione e non è stato
possibile identificarlo con certezza.}. Lo stesso accadeva al relatore,
arrivato a PG e a Stormind da Riot: a ``vedo che sei in difficoltà, posso
aiutarti?'' la risposta era sempre ``perché?''. Alla fine rispondeva: ``perché
quando sarò in difficoltà io vorrei che tu mi aiutassi''.

Quando si comincia a lavorare per gli altri e si smette di pensare ``così avrò
la promozione, così avrò l'aumento'', promozioni, aumenti e rispetto arrivano
da soli. A Blizzard aveva fallito tre volte i colloqui per essere promosso a
manager (inteso come coordinatore del team); quando ha deciso di smettere di
pensarci e di fare solo il bene del proprio lavoro, la volta dopo è stato
promosso. Funziona così: chi aiuta i colleghi e fa andare bene il progetto
viene notato (``questa persona fa migliorare cinque persone, non solo se
stessa: vale di più'').

\subsection{Il lavoro è valutato come opera d'insieme}

Il relatore riconosce che buona parte del suo successo è venuto dal fatto che
il suo team lavorava bene, non dalla sua bravura personale. Anche i progetti
degli studenti, presentati a Diego, a Enrico, ad Andrea, verranno valutati
\textbf{come opera d'insieme}. Nessuno dice ``il \emph{combat system} di questo
gioco l'ha fatto X, che bravo'': si dice ``BioWare ha fatto un bel gioco'',
``Blizzard ha fatto un bel gioco''.

Lo ha replicato anche lato HR. In molte aziende ci sono \textbf{bonus
personali} legati a obiettivi (se il gioco vende X, se la feature è pronta in
tempo). Il relatore li ha dimezzati: metà restano personali, metà diventano
\textbf{bonus di team}, che tutti ricevono se il team raggiunge gli obiettivi.
Un incentivo concreto a collaborare: non solo si dice che il lavoro di gruppo è
importante, lo si dimostra dandogli più valore economico. Aveva proposto il
100\% di bonus di team, ma gli è stato concesso il 50 e 50; ha funzionato
comunque, anche se sarebbe stato ancora più potente. Ovviamente, alla
domanda ``perché mi aiuti?'' si può anche rispondere ``perché voglio il
bonus'': un aiuto dato ``per il motivo sbagliato'' era un compromesso che era
disposto ad accettare.

\subsection{Il coraggio di chiedere aiuto}

Un partecipante osserva che aiutarsi dovrebbe essere la norma. Il modello
diffuso, invece, è ``devi essere forte e indipendente, non chiedere mai una
mano'', e vale in tutte le direzioni (dal maschilismo alla ``donna forte e
indipendente che non ha bisogno di nessuno''): \textbf{non ammettere di avere
bisogno di aiuto è sempre uno dei motivi per cui le aziende vanno male}, che si
tratti di uomini, donne o persone non binarie. Il problema è non avere la
sicurezza di poter alzare la mano.

Il relatore ha cercato di togliere ogni conseguenza dall'aver alzato la mano:
se si \emph{depenalizza} qualcosa, le persone non sono più disincentivate a
farla. Ha detto espressamente che non avrebbe mai valutato nessuno per il
numero di volte in cui chiedeva aiuto o per il numero di feedback negativi
ricevuti, e questo ha funzionato.

\subsection{Perché si dà un feedback}

\emph{Perché avete dato il vostro ultimo feedback? Che cosa volevate ottenere?}
Le risposte della classe:
\begin{itemize}
  \item per \textbf{riconoscere il valore} di un collega, anche verso il
        manager, perché qualcuno che non aveva visibilità vedesse quanto
        quella persona fosse efficace;
  \item per \textbf{spingere una persona a superarsi};
  \item per aiutare un collega a rendere il prodotto all'altezza di ciò che
        voleva rappresentare;
  \item per \textbf{genuinità}: c'era qualcosa che poteva essere fatto meglio.
\end{itemize}
Il relatore aggiunge che la forma cambia col tempo: all'inizio era ``devo
darti un feedback, hai dieci minuti?''; ora è un messaggio informale su
Discord, perché ha costruito un rapporto di fiducia e sa che non verrà letto
come ``l'HR manager mi scrive, sono nei guai?'', ma come ``Stefano è così, mi
dà feedback, resterà tra noi e lo risolveremo insieme''.

Un altro motivo è \textbf{risolvere incompatibilità caratteriali}: le
conversazioni più difficili, del tipo ``io e te non stiamo funzionando, cosa
possiamo fare diversamente?''. Sono importanti perché, se due responsabili non
vanno d'accordo, il team se ne accorge, comincia a raccontarsi storie (``si
odiano, ci saranno problemi'') e, per difendere il proprio manager, comincia a
giudicare male il lavoro dell'altro gruppo. È la dinamica da fermare:
\textbf{valutare gli altri solo per difendere la propria fazione}.

Dare e ricevere feedback sono \textbf{due muscoli da allenare} (il relatore è
sempre stato più bravo a darlo che a riceverlo), e si impara col tempo. È come
fare colloqui: fare colloqui è una skill a parte, che non indica quanto si è
bravi nel proprio lavoro. Essere scartati non significa non saper fare quel
lavoro: spesso significa aver fatto male il colloquio o che è stato preferito
qualcun altro.

\paragraph{Feedback pubblici e privati.} Uno studente racconta di aver dato un
feedback negativo a un superiore, senza un motivo particolare, per dimostrare
a sé e ai colleghi che non si viene licenziati per questo. Il relatore coglie
l'occasione per precisare che \textbf{i feedback negativi si danno
individualmente}: lui dà moltissimi feedback di cui il resto dell'azienda non
sa nulla. Reagirebbe molto male se, in un meeting, qualcuno alzasse la mano per
dire ``comunque questa cosa non l'hai fatta bene'': quello non è feedback, è
discussione, accettabile solo se è un brainstorming e se l'opinione è stata
chiesta. Alcune aziende hanno piattaforme pubbliche di feedback, ma spesso
servono solo ad applaudire; lo stesso accade nelle aziende che
\emph{obbligano} a lasciare un certo numero di feedback a settimana, che
diventano ``lui è bravo, lei è competente, anche questa settimana l'abbiamo
fatto''. Bisognerebbe chiedere a chi le adotta cosa pensa di ottenere (forse
trasparenza), ma difficilmente una persona ``attaccata'' pubblicamente reagisce
bene.

% ---------------------------------------------------------------------
\section{Among Us: la caccia al traditore}
\label{sec:z1-among-us}

Prima della pausa pranzo il relatore mostra tre film di fantascienza di
qualche decennio fa e chiede che cosa li accomuni\footnote{I titoli non sono
leggibili nella trascrizione: viene citato \emph{Event Horizon} da uno
studente, e il film ``in mezzo'' è indicato come ``l'originale'',
verosimilmente \emph{La cosa} (\emph{The Thing}).}. In ognuno c'è un gruppo di
persone con un obiettivo e la presenza, reale o percepita, di qualcuno che
rema contro: \textbf{l'energia del gruppo viene spesa per trovare il
responsabile}.

\paragraph{L'esperimento degli alieni.} Nei master in presenza il relatore
propone un esercizio. Distribuisce segretamente cartellini plastificati con
scritto ``essere umano'' o ``alieno''; le carte alieno contengono istruzioni di
comportamento con cui gli alieni, che cercano di fingersi umani, si
tradirebbero (dare sempre ragione a chi ha parlato prima, finire ogni frase
con una domanda, parlare in un certo modo). Lo scopo: in dieci minuti formare
gruppi, e \textbf{vince il gruppo più numeroso senza alieni}; basta un alieno
per perdere. Si formano sempre gruppi piccoli (due, cinque, sei persone),
perché la tendenza è ``mi fido al massimo di voi due''.

Ma vincono i gruppi più grandi, per un motivo semplice: il relatore
\textbf{non distribuisce nessuna carta alieno}. Mostra a tutti entrambe le
carte, ma consegna solo carte ``umano''. Vuole simulare ciò che accade in
ufficio: la tendenza a percepire sempre qualcuno che rema contro (i manager,
l'altro team, quel collega), quando in realtà non c'è nessuno, ``è tutto nella
nostra testa''. Quando chiede le motivazioni, emergono sospetti
\emph{fabbricati}: ``faceva sempre domande'', ``parlava solo all'infinito''
(``no, parla così e basta''). Ci si fabbrica da soli i motivi per non fidarsi
degli altri, ed è comunissimo.

\paragraph{Among Us.} L'autosabotaggio è, per tornare ad \emph{Among Us},
\textbf{concentrarsi sul cercare cadaveri e traditori invece di fare le
task}: non si vince perché non si sono completate le task, mentre si cercava il
colpevole. Succede tantissimo negli uffici, ed è il motivo per cui l'80\% dei
problemi che il relatore risolveva erano di comunicazione. Entrando nei team
sentiva immancabilmente: ``facciamo il possibile, ma questa lotta continua con
l'altro gruppo è estenuante''. ``Perché c'è questo problema?'' ``Non lo so,
c'era già quando sono arrivato.'' Un paradosso: c'è un problema, nessuno sa
dire quale sia, ma c'è sempre stato e quindi lo si accetta.

\paragraph{L'esperimento delle scimmie.} È l'esempio, ``finto o reale'', delle
scimmie in gabbia\footnote{È il noto racconto delle ``cinque scimmie'', spesso
presentato come esperimento scientifico ma probabilmente apocrifo (lo stesso
relatore lo definisce ``finto barra reale''); nella trascrizione il numero
delle scimmie varia.}. In una gabbia con alcune scimmie pende una banana
elettrificata: chi la tocca fa prendere la scossa a tutte. Le scimmie imparano
e, quando una si avvicina alla banana, le altre la picchiano. Poi si
sostituisce una scimmia alla volta con una nuova: la nuova va verso la banana,
le altre la picchiano, e smette. Si continua finché nella gabbia ci sono solo
scimmie nuove, poi si toglie la corrente. Nessuna scimmia tocca la banana:
nessuna ha mai preso la scossa, ma tutte hanno imparato che chi tocca la
banana prende le botte.

Nei team è lo stesso: il motivo per cui si ``odia'' l'altro team si
trasforma, ma un vero motivo non c'è. Ne derivano la sfiducia, il non alzare
mai la mano per chiedere aiuto, il non dire mai ``bel lavoro'' agli altri
perché non sia mai che qualcuno ti veda fare un complimento ``a quelli''.
Risolvere queste dinamiche è lungo e tedioso, e sempre pieno di ``non mi
ricordo'', ``io non c'ero'' e ``abbiamo sempre fatto così'', la cosa più
frustrante per chi cerca di risolverle.

% ---------------------------------------------------------------------
\section{Il principio del default to trust}
\label{sec:z1-default-trust}

\subsection{La definizione}

Il relatore indica questa come la slide più importante della presentazione.
\textbf{Default to trust} è uno dei valori aziendali di Riot Games, che lui ha
tradotto, adattato e ``esportato'':

\begin{quote}
  \emph{Approccio ogni situazione di conflitto partendo dal presupposto che una
  persona abbia agito in un certo modo perché, con le informazioni a sua
  disposizione in quel momento, ha pensato che fosse la soluzione migliore, e
  non perché sia un'incapace o volesse danneggiare il progetto
  intenzionalmente.}
\end{quote}

In altre parole: \textbf{di default, mi fido}.

\subsection{L'esempio del ritardo al daily stand-up}

A Riot il relatore aveva un \emph{daily stand-up} dal vivo alle 15 con il suo
team di venti persone. Una ragazza arrivava regolarmente alle 15:15, per una,
due, tre settimane. Alla quarta ci ha parlato. \emph{Come iniziereste la
conversazione?} Le proposte della classe:
\begin{itemize}
  \item sottolineare l'importanza di esserci dall'inizio, perché potrebbe
        perdersi qualcosa che le interessa, e per rispetto dei colleghi
        presenti;
  \item ``devi arrivare in orario'', spiegando perché è importante;
  \item \textbf{chiedere perché}.
\end{itemize}

Il relatore osserva, anche dalle sue consulenze fuori dai videogiochi, che
\textbf{più si sale con la seniority}, più la risposta è ``devi smettere di
arrivare in ritardo'', spesso con un \emph{or else} finale (``altrimenti
succede quest'altra cosa''). Più si scende con l'età, e quindi con meno
abitudini consolidate, più la prima domanda diventa ``dimmi perché arrivi in
ritardo''. Già spiegare perché è importante è ``miglia avanti'' rispetto al
semplice ``smetti''.

Partire dal perché gli ha permesso di scoprire la causa: il meeting precedente
della ragazza, in un'altra zona del campus a un quarto d'ora di distanza,
finiva sempre tardi, e lei doveva attraversare il campus di corsa. Non aveva
la confidenza per dire all'altro manager, che insisteva perché si restasse
fino alla fine, che doveva andare; sacrificava quindi il meeting del relatore,
sapendo che era il più comprensivo dei due.

Il relatore si è offerto di \textbf{mediare}: ha detto all'altro manager che
o anticipava l'inizio del suo meeting, così da finire comunque in tempo,
oppure lasciava uscire la ragazza, perché quel meeting gli serviva. C'è anche
una questione di \textbf{economia di scala}, spesso sottovalutata: se venti
persone aspettano qualcuno per dieci minuti, si sono sprecati dieci minuti per
venti. Per lo stesso motivo il workshop è iniziato alle 10:05: iniziare alle
10:10 avrebbe significato far aspettare cinquanta persone cinque minuti in
più. A Riot c'era una scrivania vuota con un pupazzo gonfiabile, \textbf{Mr.\
Late}: sommando tutti i minuti di ritardo a meeting, call e conversazioni, in
un anno veniva fuori praticamente \emph{una persona a tempo pieno}. Serviva a
ricordare che anche cinque minuti contano, e che bisogna rispettare gli altri
arrivando in orario o almeno avvisando.

\paragraph{Gli effetti.} La conversazione era avvenuta nel primissimo periodo a
Riot, quando il team non aveva ancora motivo di fidarsi di lui. Risolvendo il
problema ha eliminato anche tutto lo strato di micro-dinamiche spiacevoli che
un manager impara a vedere (due colleghi che si guardano con la risatina
quando lei entra in ritardo). La ragazza ha cominciato a entrare tranquilla,
con una certa leggerezza: non le faceva piacere arrivare in ritardo (``ci sono
anche persone che lo fanno apposta, e quello è un problema da affrontare
diversamente''). Soprattutto, con i colleghi non parlava di lui come del capo
che minacciava, ma come di chi le aveva risolto un problema: era diventato
un \textbf{alleato}. La settimana dopo altre tre persone sono andate da lui
con un problema da risolvere. Ripetuto cinque, sei, sette, dodici volte, in un
paio di mesi poteva dire quasi qualsiasi cosa nei meeting sapendo che
richieste e critiche sarebbero state percepite come ``me lo chiede perché vuole
darmi una mano'', non per fare bella figura o per presentare i successi del
team ai suoi superiori.

\subsection{Condividere i successi, schermare i fallimenti}

Un altro elemento fondamentale (ripreso nella seconda giornata): il manager
\textbf{condivide i successi e protegge il team dai fallimenti}. ``Se va male
qualcosa è colpa mia; se va bene qualcosa, voglio che ne parliate voi.'' Il
lavoro del manager non dovrebbe essere valutato su ciò che presenta del lavoro
del team: tutti sanno comunque chi è il manager. Tanto vale lasciar presentare
il team, che potrà dire con orgoglio di aver presentato a Mike
Morhaime\footnote{Cofondatore ed ex presidente di Blizzard; nella
trascrizione ``Mike Morayme''.} la propria campagna europea per
\emph{StarCraft~II}. Ci sono manager che dicono ``come mi prendo le colpe, mi
prendo i meriti'', ma i meriti li hanno già: è inutile ostentarli. Con i
fallimenti invece il manager è la \textbf{rete di sicurezza} del team. Il
relatore ne ha fatto una questione di principio: \textbf{non scaricare mai la
colpa su qualcun altro}.

\paragraph{Colpevole o causa.} Un esempio che il relatore attribuisce a Pirate
Software, creatore di contenuti ed ex dipendente Blizzard\footnote{Il relatore
aggiunge che si tratterebbe del figlio dello sviluppatore a cui si sarebbe
ispirato il personaggio del giocatore ``senza vita'' nell'episodio di
\emph{South Park} dedicato a \emph{World of Warcraft}. L'informazione è
riportata come detta in lezione e non è verificata.}. Contattando il servizio
clienti Blizzard per motivi personali, aveva scoperto che i colleghi gli
avevano scritto nei ticket informazioni confidenziali che non avrebbero dovuto
condividere in quel modo, nemmeno con un collega. Non ha segnalato le persone
da punire: ha detto ``c'è un problema, evidentemente le policy non sono
chiare''. Quando gli hanno chiesto i nomi, ha rifiutato: ``se vi dico chi è,
punite quelli, e gli altri cinquecento continueranno a farlo; non avremo
cambiato la causa alla fonte''.

\begin{quote}
  \emph{È inutile trovare il colpevole. Voglio trovare la causa. Perché il
  colpevole cambia.}
\end{quote}

Il colpevole a volte è vittima della situazione: per questo la definizione
dice \emph{``con le informazioni a sua disposizione in quel momento''}. Se il
relatore avesse rimproverato la ragazza in ritardo, non solo non avrebbe
risolto il problema, ma avrebbe creato un'\textbf{ingiustizia}: lei, che
attraversava il campus di corsa per una colpa non sua e di cui non poteva
parlare, sarebbe stata anche rimproverata.

Nessuno fa un brutto lavoro apposta: nessuno è contento quando sbaglia o
perde. Perché al lavoro diamo per scontato che gli altri sbaglino
intenzionalmente? Se qualcuno ha fatto una cosa male, è perché in quel momento
credeva fosse la scelta migliore. La domanda giusta non è ``perché hai fatto
questa cosa che è stata un fallimento?'', ma:

\begin{quote}
  \emph{Che cosa ti ha fatto pensare che questa fosse la scelta giusta?}
\end{quote}

Magari erano limiti di budget, di tempo, di scopo: si agisce su quelli, non
sulla persona. Se si cambia la persona ma le si danno gli stessi fattori, è
molto probabile che arrivi alle stesse conclusioni, e ci si troverà a dover
``disciplinare'' anche lei.

\paragraph{La reputazione costruita sui rumor.} Durante il giro di
presentazione in un ufficio (non in nessuna delle aziende citate), al relatore
è stato detto: ``questo è il dipartimento di design; se ti serve qualcosa
chiedi a lei, non a lui, perché lui consegna sempre in ritardo''. Oltre a
trovarlo poco professionale, una settimana dopo ha chiesto: ``cosa ti ha fatto
questa persona?'' ``A me niente, però ho sentito che\dots''. Si torna alle
scimmie: magari quella persona non consegna in ritardo da mesi, ma ormai lo
dicono tutti. È un circolo vizioso da fermare subito.

\paragraph{Niente escalation inutili.} Lo stesso vale per il feedback: quando
qualcuno va dal relatore dicendo ``ho un problema con questa persona'', lui
chiede: \textbf{``perché ne stai parlando a me e non a lei?''}. ``Perché sei il
mio capo'' non basta. Passare da lui significa dare per scontato che l'altra
persona non abbia interesse a risolvere, mentre magari basta parlarle. È la
strada più facile (``lo dico a Stefano, ci pensa Stefano''), e proprio per
questo bisogna mettere dei paletti: il manager interviene quando serve, ma le
cose che si possono gestire da soli vanno gestite da soli. (Quando invece il
confronto diretto fallisce, l'escalation è legittima: si veda la
sezione~\ref{sec:z1-qa-finale}.)

% ---------------------------------------------------------------------
\section{Traiettorie di carriera e valutazione delle performance}
\label{sec:z1-traiettorie}

\subsection{Quale traiettoria vorreste?}

Il relatore mostra un grafico con quattro curve colorate, contrassegnate
anche da lettere (A, B, C, D) per eventuali persone daltoniche (``a me è
capitato che qualcuno scoprisse di esserlo durante una lezione''). Le curve
rappresentano la performance lavorativa nel tempo. Due domande: \emph{quale
vorreste essere? E quale pensate realisticamente di essere?} Un partecipante
aspira a B e pensa di essere C; quasi tutti scelgono B.

Perché non A, che sale più rapidamente? Un partecipante risponde che una
crescita così repentina può essere più dannosa che utile, come una ``rush'' di
studio molto faticosa che poi non si riesce a sostenere. Il relatore propone
una lettura suggerita da un amico psicologo, che spiega la forma del grafico:
di B, C e D \textbf{vediamo dove finiscono}; di A no. Potrebbe continuare a
salire oltre la scala del grafico, o crollare sotto D. Non saperlo ci spinge a
preferire B. (A chi obietta che anche le altre proseguono oltre il bordo, il
relatore precisa: il bordo destro è ``oggi'', e di A non sappiamo come sia
arrivata fin lì.)

\subsection{La carriera non è una linea retta}

Tutti immaginano la carriera come una linea: si parte da X e, attraverso
stage e posizioni, si arriva a Y. In realtà carriera e performance avranno un
andamento irregolare, con alti e bassi. Come non si accetta che un team possa
temporaneamente \emph{underperformare}, tanto meno lo si accetta da se stessi.
Bisogna accettare che ci saranno momenti bassi.

Per questo il relatore \textbf{detesta le valutazioni annuali delle
performance}: prendono un punto casuale della traiettoria, che se si è
fortunati è il picco più alto e se si è sfortunati il più basso, e lo usano
per valutare un intero anno. Insegna ai manager a fare più punti di controllo
e a guardare altre cose:
\begin{enumerate}
  \item \textbf{quanto tempo passa tra un picco e l'altro}, cioè quanto
        \emph{cooldown} serve a una persona. È il motivo per cui, dopo la
        Games Week, tutti staccavano una o due settimane. Il relatore stesso,
        che si impegna moltissimo e tende al burnout, non potrebbe rifare lo
        stesso workshop il lunedì e il martedì successivi con la stessa
        qualità;
  \item \textbf{la differenza tra la miglior performance e la peggiore}: che
        cosa ha portato a lavorare così bene, e così male? Con chi si
        lavorava, su cosa ci si concentrava, quanto tempo e budget c'erano?
        L'obiettivo è capire come \emph{replicare} le condizioni che fanno
        lavorare bene, invece del solito ``hai lavorato male, è colpa tua''.
        Il relatore ha avuto molti manager per cui quando lavorava male era
        colpa sua e quando lavorava bene ``che bravo'', ma sapeva che raramente
        era l'unico fattore della propria performance. È la fallacia del ``se
        sai fare 100, perché non fai sempre 100?'', come nella scena di
        \emph{Schindler's List} sull'operaio che produce una vite in pochi
        secondi: non sempre un ritmo è sostenibile;
  \item \textbf{il fallimento di oggi confrontato con quello di sei mesi fa}.
        Il relatore non è contento della presentazione di oggi (l'avrebbe
        migliorata per altri due anni), ma rispetto a quelle di cinque anni fa
        vede un grande miglioramento di chiarezza: il suo peggior fallimento di
        oggi equivale al suo primo successo di inizio carriera.
\end{enumerate}
Quello che conta è la \textbf{traiettoria} dall'inizio dell'anno a oggi, non
i singoli punti. E se si riesce a capire insieme quando e perché si ha
successo e quando e perché si fallisce, si può lavorare per replicare solo le
situazioni che fanno lavorare bene.

\paragraph{Lo stesso vale per i giochi.} È il motivo principale per cui i
giochi delle grandi case vanno bene o male. Quando si riescono a replicare le
condizioni che hanno prodotto \emph{Warcraft~II}, esce \emph{Warcraft~III};
quando non ci si riesce, esce ``il nuovo \emph{Dragon Age}'': la gente non è
arrabbiata perché il gioco non è bello, ma perché non è come gli altri
\emph{Dragon Age}; non ne replica il \emph{feeling}, anche se potrebbe
essere un gioco favoloso. L'eccezione positiva è il passaggio da \emph{Dune} a
\emph{Dune~II}: il primo era di fatto un'avventura grafica, il secondo ha
generato il genere degli strategici in tempo reale. Gli sviluppatori avevano
capito che il successo non stava nel tipo di gioco, ma nell'ambientazione che
attirava le persone: un bel gioco in quell'ambientazione sarebbe piaciuto
comunque. Il relatore comprò \emph{Dune~II} aspettandosi un seguito
dell'avventura e, pur sorpreso, lo finì.

\subsection{Lying, hiding, faking}

Quando si è valutati su un singolo punto nel tempo, e da quel punto dipendono
bonus, stipendio e il rischio di licenziamento, si comincia a fare
\textbf{lying, hiding e faking}:
\begin{description}
  \item[Lying] dire che è sempre andato tutto bene, che si finirà sempre in
        orario, che non c'è nessun problema;
  \item[Hiding] nascondere: ``non era un mio task'', ``non c'ero'', ``non ho
        visto niente'';
  \item[Faking] fingere, anche in cose apparentemente innocue come ``sì, sì,
        ho capito tutto'' quando si ha solo paura di dire ``me lo rispieghi?
        Non ho capito niente''.
\end{description}

\subsection{``Non ho capito'' e la conferma delle istruzioni}

Dire ``non ho capito, ripetimelo'' è ancora visto come una figura da sciocco.
Il relatore non si arrabbierà mai se qualcuno gli chiede di rispiegare
un'intera slide: le cose ovvie per lui possono non esserlo per altri. Più si
conosce qualcosa, \textbf{più è difficile spiegarla}, perché si danno per
scontate molte cose. Ne parlava proprio con Diego, che gli diceva che fare
l'insegnante diventa sempre più difficile man mano che si diventa esperti. È
come spiegare \emph{EVE Online}, uno dei giochi più complessi in circolazione:
chi lo conosce bene dà per scontato che ``ovviamente non puoi fare questa
cosa'', mentre per un nuovo giocatore non è affatto ovvio. A scuola chi dice
``non ho capito'' viene visto come lo ``scemo''; nessuno si chiede ``non è che
l'ho spiegato male?''.

Il racconto più ``traumatico'' dei ragazzi che fanno stage nelle grandi
aziende è sempre lo stesso: ``Cinzia, devi fare questa cosa.'' Cinzia la fa.
``Non te l'avevo chiesta così.'' Il problema è che nessuno ha
\textbf{confermato che Cinzia avesse capito}. Il consiglio, valido sempre:
\begin{itemize}
  \item quando si \emph{ricevono} istruzioni, ripeterle: ``ho capito bene?
        Vuoi che faccia questa cosa in questo modo?'';
  \item quando si \emph{danno} istruzioni, chiedere a chi le riceve di
        ripeterle, non perché lo si ritenga incapace, ma per capire che cosa
        ha capito dalla propria spiegazione (``quindi vuoi un quadrato?''
        ``No, ho detto un triangolo\dots'' ``Però mi hai detto una forma con
        quattro angoli.'' ``Hai ragione: voglio un quadrato'').
\end{itemize}
Frasi come ``non te lo ripeto due volte'' o ``stammi dietro'' portano solo a
fare le cose come si pensa che vadano fatte e a sentirsi dire a fine progetto
``non la volevo così'': che si tratti di una feature, di un design o di un
modello 3D. (Un partecipante osserva che lo stesso accade insegnando le lingue:
chi insegna dà per scontato, ad esempio, che si sappia dell'esistenza del
genere neutro in tedesco.)

\paragraph{Come convincere il team a chiedere?} Il relatore mette sempre come
premessa che \textbf{l'obiettivo non è trovare la persona più intelligente
nella stanza, ma fare un buon lavoro}. Che Gianni capisca alla prima e un altro
alla quinta non gli interessa, se alla fine entrambi lavorano bene: ripetere
cinque volte è comunque meno lavoro che far rifare le cose sette volte. È un
investimento a lungo termine, e non dissimula alcun fastidio: davvero non gli
dà fastidio. Ma ripetere le stesse parole non aiuta: ``non ho sentito'' si
risolve ripetendo; ``non ho capito'' si risolve cambiando spiegazione. Meglio
ancora: ``dimmi che cosa hai capito tu, e vediamo cosa non torna''. Se
intravedesse nel team uno sguardo del tipo ``è la quinta volta che lo
chiede'', interverrebbe: spesso chi giudica è proprio chi dice ``ho capito'' e
poi sbaglia. Preferisce di gran lunga chi chiede cinque volte.

\subsection{Il feedback più utile che abbiate ricevuto}

\emph{Qual è il feedback più utile che abbiate mai ricevuto?} Il relatore ne ha
almeno dieci che gli hanno cambiato il modo di lavorare. Le risposte:
\begin{itemize}
  \item un semplice \textbf{``grazie''} per un lavoro fatto, più utile di un
        ``questa cosa fa schifo'';
  \item un feedback sulla propria \textbf{impulsività} nel rispondere: fermarsi
        un attimo e capire di cosa si sta parlando prima di intervenire;
  \item \textbf{essere se stessi}, anche al lavoro, e non scusarsi
        continuamente per ogni cosa.
\end{itemize}
I feedback che toccano di più sono quelli \textbf{veri}. Il ``grazie''
significa ``ho visto quello che hai fatto e l'ho apprezzato''. Quelli negativi
che si ricordano sono quelli che hanno fatto realizzare ``questa cosa che
pensavo di saper fare, non la so fare'', oppure ``la so fare, ma questa volta
non mi è venuta'' (non è che se una cosa una volta non riesce non la si sa più
fare: può capitare di sbagliare una porta). Arrivano sempre da un
\textbf{punto cieco}: qualcosa che non si sapeva fosse percepito dagli altri.
Conta più ciò che le persone ti dicono di ciò che pensi che pensino di te; e
ciò che il relatore pensa che gli altri pensino di lui è sempre terribile,
mentre il loro giudizio si rivela molto più positivo del previsto.

\paragraph{Il denominatore comune siete voi.} Si potrebbe pensare che questi
concetti non siano applicabili da chi arriva in un'azienda senza poterne
cambiare la cultura. Ma c'è sempre un denominatore comune a cui applicarli:
\emph{se stessi}. La cultura aziendale può non spingere a dare feedback, può
vedere male chi chiede aiuto; ma il relatore ha visto che quando lui lo fa, le
persone intorno tendono a farlo, e un gruppo che lo fa lavora meglio. Perfino
i CEO più ostici a un certo punto gli chiedevano: ``perché il tuo team ora
performa meglio degli altri?''. La differenza è il default to trust: una
pratica che col tempo fa operare il team \emph{come un team} e non come un
insieme di individui. Un videogioco è un progetto corale: c'è chi fa giochi da
solo, ma non si fa da soli \emph{Red Dead Redemption~2}; e se il gruppo non è
affiatato e allineato non esce \emph{Red Dead Redemption~2}, esce
\emph{Anthem} (``oggi la mia fissazione è \emph{Anthem}, domani sarà
un'altra'').

% ---------------------------------------------------------------------
\section{Come dare un feedback}
\label{sec:z1-dare}

\subsection{Concentrarsi sull'azione e sui suoi effetti}

Il primo principio: \textbf{il feedback non riguarda mai la persona}. Non si
dice ``tu hai fatto questa cosa che non va bene'', ma ``questa cosa non va
bene: come facciamo a non farla più?''. Che sia un comportamento, un'azione,
una decisione o una strategia, ci si concentra su quello, e si
\textbf{esplicita l'effetto} che ha. Nell'esempio del ritardo l'effetto era
``quindici minuti di attesa moltiplicati per venti persone''. In un progetto
può essere una catena:
\begin{quote}
  se non rispetti le deadline, l'animazione che viene dopo resta ferma; se
  loro sono in ritardo, la \emph{vertical slice} è in ritardo; se la vertical
  slice non è pronta per la Gamescom, non andiamo alla Gamescom e non troviamo
  un publisher; se non troviamo un publisher, chiudiamo.
\end{quote}
Più la conseguenza è grave, più dovrebbe sembrare piccolo il gesto di dire
``datemi una mano'': è il \emph{male minore}. Poi, a cose fatte, ci si dà
feedback: cosa ha funzionato, cosa no, cosa fare di più, di meno, cosa smettere
di fare.

\subsection{Il pericolo degli ambienti solo positivi}

Gli ambienti dove ci si fa sempre grandi complimenti e nessuno parla dei
problemi creano difficoltà vere alle aziende. Il problema non è mai la
mancanza di complimenti, ma un \textbf{eccesso di ottimismo}: si comincia a
pensare che le cose andranno come si spera, perché tutti intorno danno solo
feedback positivi.

A Blizzard i feedback determinavano gli \textbf{aumenti annuali}: ognuno
doveva indicare cinque persone che avrebbero dato feedback su di lui, e
ovviamente tutti indicavano gli amici, o colleghi di progetti andati bene. Un
sistema nato con le migliori intenzioni diventava uno \emph{scambio di voti}
(``se mi metti follow ti metto follow''), senza alcun miglioramento.

Al contrario, \textbf{le conversazioni più brutte sono state le più utili}:
affrontare i licenziamenti a Riot, o la cancellazione di \emph{Project Titan}
a Blizzard, il progetto che doveva essere il nuovo MMO dopo \emph{World of
Warcraft} e da cui poi è nato \emph{Overwatch}\footnote{Nella trascrizione si
legge anche ``doveva essere Warcraft 2'', probabilmente un lapsus o un errore
di trascrizione.}. Andando a vedere perché lo si cancellava, c'erano solo
motivazioni razionali e pragmatiche, di business, chiare a tutti, non un ``Mike
Morhaime ci odia''. Sono anche le conversazioni in cui ha lavorato su se
stesso: è molto \textbf{avverso ai conflitti}, e questo a volte gli ha
procurato conseguenze molto più gravi di quanto sarebbe costato un piccolo
disagio. Dieci anni fa faceva fatica ad accettare questo difetto; ha dovuto
subire le conseguenze di conflitti evitati e poi cresciuti.

\subsection{Che cosa fare}

Le linee guida della slide, che il relatore legge rapidamente perché saranno
disponibili nel materiale:
\begin{description}
  \item[Attenzione al linguaggio del corpo] dare un feedback a braccia
        incrociate o guardando altrove è diverso dal guardare l'altro negli
        occhi.
  \item[Sii specifico] con esempi tangibili.
  \item[Esplicita gli effetti] come visto sopra.
  \item[Discuti il prima possibile] non aspettare un mese per parlare di
        qualcosa successo ieri, altrimenti nessuno ricorderà più bene cosa sia
        accaduto.
  \item[Concentrati sull'azione e sul problema, non sulla persona] il nemico è
        l'azione, non la persona che si ha davanti. Non bisogna far cambiare
        l'altro: bisogna evitare che il problema si ripresenti.
\end{description}

\paragraph{Gli insulti sono rivolti al ruolo.} Un esempio solo apparentemente
scollegato. Il relatore ha lavorato su \emph{League of Legends}, ``non
esattamente noto per avere una community da orsetti del cuore'', in varie
\emph{task force} sulla tossicità. A chi smetteva di giocare veniva inviato un
sondaggio sul perché; il 33\% indicava la tossicità della community. Dopo varie
campagne di sensibilizzazione e cambiamenti si è scesi al 30\%: sembra solo
un 3\% in meno, ma valeva circa \textbf{120 milioni di revenue in più
all'anno}, perché quei giocatori restavano esposti a skin e contenuti e
continuavano a comprare. Anche qui, piccoli guadagni che si sommano.

Lavorando con degli psicologi, aveva notato che gli insulti \textbf{non sono
quasi mai rivolti alla persona, ma al ruolo}: ``il mio jungler fa schifo'',
``il mio DPS non è capace'', ``il nostro portiere è una pippa''. Chiamare
qualcuno per ruolo lo \emph{depersonalizza} e rende più facile insultarlo. Lo
stesso nell'esempio del giro dell'ufficio: non ``non chiedere a Giulio'', ma
``non chiedere a quel grafico, consegna sempre in ritardo''. Concentrarsi
sull'azione, senza attaccare la persona, fa sì che il feedback non sia (e non
venga percepito come) un attacco personale.

\paragraph{Considera il carattere.} Qui tornano utili i test di personalità.
Servono più agli studenti che al relatore (che sa già che usciranno 53 profili
diversi, senza un pattern): bisogna conoscere il proprio profilo e
\textbf{condividerlo con chi lavora con noi}, chiedendo anche agli altri di
fare lo stesso. In ufficio il relatore ha usato una versione più avanzata, a
pagamento, con un report di circa trenta pagine: il \textbf{profilo
DISC}\footnote{Nella trascrizione ``disk''. Si tratta del modello DISC
(Dominance, Influence, Steadiness, Conscientiousness).}. Ha mappato tutta
l'azienda e ha fatto stampare il profilo di ciascuno sulla sua scrivania, così
che tutti sapessero con chi stavano parlando. Capire come ragionano gli altri
(``è introverso, non gliene parlo davanti a tutti'') è una delle cose che gli
sono tornate più utili: \emph{adattare ciò che si vuole dire a chi si ha
davanti}.

\paragraph{Assicurati di non disturbare.} Sotto il profilo di ognuno c'era un
piccolo \textbf{semaforo}: verde, ``parlami pure, interrompimi''; giallo,
``interrompimi solo se è urgente''; rosso, ``mi sto concentrando, non mi
parlare''.

\paragraph{Lascia finire il lavoro.} Durante gli eventi dal vivo il relatore si
era reso conto di quanto fosse grave dare feedback mentre qualcuno sta ancora
lavorando. Se una cosa sta andando male, dirlo in quel momento non aiuta: tra
la pressione di finire e quella del feedback, peggiorano entrambe le cose, e
anche il rapporto. Meglio lasciarla finire, anche male, e parlarne dopo.

\subsection{Che cosa non fare}

\begin{description}
  \item[Non basarti su supposizioni] ``mi hanno detto che ieri in un meeting
        hai fatto\dots'': o c'eri e l'hai visto, o non esiste.
  \item[Non giudicare la persona.]
  \item[Non fornire feedback vago.]
  \item[Non accettare e non fornire feedback di terze parti] Se qualcuno vi
        dicesse ``mi hanno detto che hai fatto un brutto lavoro'', il primo
        pensiero non sarebbe ``cosa ho fatto male?'', ma \emph{``chi te l'ha
        detto? Chi è la vipera che ho in seno? Chi è l'impostore?''}. Siamo
        fatti così. Ogni conversazione che inizia con ``mi hanno detto che'' è
        inutile: l'altro non ascolterà, penserà solo a chi dare la caccia, e
        spesso penserà alla persona sbagliata, creando antipatie tra persone
        che non c'entravano nulla. Quando al relatore dicono ``mi hanno detto
        che lui ha fatto\dots'', risponde ``vai a parlarne con lui''; se gli
        dicono ``mi hanno detto che tu hai fatto\dots'', risponde ``mandami la
        persona che te l'ha detto, altrimenti non ne parliamo''. Non per
        cattiveria: se chiedesse ``che cosa ho detto di preciso a Marco?'' e la
        risposta fosse ``non lo so, non c'ero'', di cosa si starebbe parlando?
        \textbf{Parlate con la fonte}, altrimenti fate come se quel feedback non
        esistesse.
  \item[Non dare consigli non richiesti] come visto, non portare la soluzione:
        la soluzione deve venire da chi riceve il feedback.
  \item[``Non mi piace'' non è un feedback] ``Penso che questo design non
        funzioni perché, visto che il nostro obiettivo è\dots'' ha argomenti;
        ``non mi piace il combat system'' è un'opinione. Anche ``io l'avrei
        fatto diversamente'' non è un feedback.
  \item[Accettare non vuol dire fare come proponi] deve essere chiaro a chi dà
        il feedback: l'aspettativa è essere ascoltati e presi in
        considerazione, ma se la responsabilità è dell'altro, deciderà
        l'altro.
  \item[Non offenderti] se il feedback non viene seguito. Restarci male
        significa pensare di saperne più dell'altro, cioè \emph{venire meno al
        default to trust}. ``Decidi tu'' significa ``mi fido di te, perché sei
        lo specialista, e mi fido che mi ascolterai e farai la cosa migliore
        per tutti'', che potrebbe anche non essere ciò che ho detto io. Se
        voglio che l'altro faccia sempre come dico io, mi sto fidando solo di
        me stesso.
\end{description}

\subsection{Il primo feedback e il debriefing}

La cosa più difficile è dare il \textbf{primo} feedback: è un muscolo da
allenare. Chi non l'ha mai fatto (soprattutto con questa struttura) dovrebbe
provarci, anche su minuzie (``abbiamo concordato di usare WhatsApp nel gruppo:
non potremmo usare Telegram?''), anche solo per vedere come va. Altre
indicazioni:
\begin{itemize}
  \item feedback \textbf{solo uno a uno}, mai davanti ad altri, soprattutto se
        negativo;
  \item \textbf{prepararsi} prima;
  \item fare il \textbf{debriefing del feedback}: dopo, fermarsi a pensare a
        com'è andato. Il relatore, quando media un feedback tra due persone,
        si sente spesso rispondere ``benissimo'' da una e ``malissimo''
        dall'altra: la percezione è molto diversa. Non si fa problemi a tornare
        dalla stessa persona: ``ieri ti ho dato un feedback: com'è andata? Sei
        soddisfatto? Come poteva andare meglio?''. Con le persone con cui si ha
        meno a che fare, questo aiuta molto la volta successiva.
\end{itemize}

% ---------------------------------------------------------------------
\section{Come ricevere un feedback}
\label{sec:z1-ricevere}

Per il relatore è la parte più difficile. Si riassume in quattro passi:
\begin{enumerate}
  \item \textbf{Ascolta}: si parte sempre ascoltando.
  \item \textbf{Chiedi dettagli}, se serve, ma per capire meglio, non per
        difenderti.
  \item \textbf{Accetta, a prescindere}: non si sta decidendo chi ha ragione;
        si sta ascoltando un'opinione. Non c'è nulla da cui difendersi.
  \item \textbf{Agisci}: fatti un piano. Può essere banale (``mi hanno detto
        che interrompo sempre: cerco di non interrompere''), oppure ``devo
        migliorare questa cosa: come?''. Pensa anche a \emph{chi} potrebbe
        aiutarti e a \emph{che supporto} ti serve: il piano non deve essere
        per forza ``devo risolvere tutto da solo''.
\end{enumerate}
Una domanda che molti sottovalutano: \emph{``sono junior in un'azienda: sarò
mica il primo ad avere questo problema?''}. Probabilmente qualcun altro l'ha
già avuto: basta chiedere ``tu come hai fatto?''.

\subsection{Il feedback positivo}

Tutto questo vale per i feedback costruttivi negativi, a cui si pensa di
default; ma esistono anche i \textbf{feedback positivi}, e vanno usati. Quando
si riceve un ``grazie'' o un ``bravo'', la domanda è: \emph{come faccio a
ricreare le stesse condizioni?}

Il relatore ha lavorato anche in aziende dove il feedback era solo negativo.
Il rischio? Si immagini l'impiegato ideale: vuole sempre migliorare, è sempre
pronto a dire ``magari sbaglio io''. Se nessuno gli dice cosa sta facendo
\emph{bene}, spinto dal desiderio di migliorare rischia di cambiare cose che
erano già buone, e di farle peggio. Il rinforzo positivo \textbf{non è solo una
gentilezza}: è un atto pratico, che evita che le persone peggiorino ciò che
sanno già fare bene.

\subsection{Il feedback accidentale}

Esiste infine il \textbf{feedback accidentale}, molto frequente negli uffici:
\begin{itemize}
  \item il ``chiedi a lei, non a quel grafico'' del giro dell'ufficio è un
        feedback negativo su una persona, dato accidentalmente, né richiesto
        né voluto;
  \item a Riot li chiamavano \textbf{drive-by}, come nelle scene dei film in
        cui passa un'auto, si abbassa il finestrino, si spara e si riparte. In
        molte aziende con una ``cultura del feedback'' in cui tutti devono
        dire la loro, si riceve una valanga di commenti del tipo ``comunque il
        sito a me non piace, ciao''. Non si sa perché, non si sa cosa fare di
        diverso, non si sa che rischi ci siano: \emph{a che cosa è servita
        questa conversazione?} A niente.
\end{itemize}

% ---------------------------------------------------------------------
\section{Quando, quanto e con chi}
\label{sec:z1-quando}

\subsection{Frequentemente e recentemente}

\emph{Quando avete ricevuto il vostro ultimo feedback?} Per gli studenti,
qualche giorno fa (nel corso) o qualche mese fa (al lavoro, dopo una
consegna). Come aveva osservato Silvia, il feedback andrebbe dato
spesso.

Il relatore usa l'esempio di una nave da crociera (``la Costa Concordia'') che
deve andare da A a B. La rotta ideale (azzurra) è quasi impossibile: nessun
progetto va esattamente come previsto. C'è la rotta ``panoramica'' (verde), che
gira intorno all'isola d'Elba; ci sono poi rotte (arancione e rosa) che a un
certo punto si allontanano molto e poi correggono bruscamente: qualcuno si è
accorto che si stava sbagliando, ``ma non è quello che avevo chiesto'',
``non è dove dovevamo andare''. Sono i famosi \textbf{sanity check} o
\emph{checkpoint}.

L'ideale è avere \textbf{tanti piccoli check per fare tante micro-correzioni}:
correggere la rotta di un paio di gradi è molto più rapido che farlo dopo
essersi allontanati. Bastano trenta secondi di ``ehi, vado bene?''. Spesso non
lo si fa per paura di sembrare insicuri o poco autonomi, ma per il relatore è
l'opposto: quando qualcuno gli chiede ``vado bene così?'' e lui risponde
``sì'', si alleggerisce anche del proprio carico mentale. Un capo che assegna
un compito e lo vede solo alla fine vive sempre col dubbio (``avranno capito?
Lo staranno facendo giusto? Arriverà in tempo?''); con un check di dieci
secondi al giorno (``tutto ok, Stefano'') è tranquillo, e la persona va avanti
sapendo che quanto fatto finora va bene. Non serve fare ogni giorno il giro del
dipartimento: anche ogni uno o due giorni basta a evitare la rotta verde.

Per il feedback è lo stesso: \textbf{va dato subito}. Dire oggi ``ieri hai
fatto questa cosa'' è molto diverso da dire ``negli ultimi sei mesi hai fatto
questa cosa''. Il relatore si è sentito dire più volte, alla valutazione di
dicembre, ``da maggio hai fatto questa cosa che non andava''. La sua risposta:
``e aspetti dicembre per dirmelo? Mi hai lasciato sbagliare per sei mesi?''. È
una mancanza di rispetto, e la conversazione non può più essere costruttiva.
Per questo non ama le valutazioni annuali: i manager crescono abituandosi a
prendere appunti delle cose negative per tirarle fuori alla valutazione
(``a marzo hai fatto questa cosa'' ``sì, però poi ho smesso'' ``sì, però a
marzo l'hai fatta'').

In sintesi:
\begin{description}
  \item[Frequentemente] riguarda l'intervallo tra un check e l'altro.
  \item[Recentemente] significa non lasciar passare un anno prima di dare il
        feedback, altrimenti non serve più a niente.
\end{description}

\subsection{Quanto dura un feedback}

Un feedback dura \textbf{da trenta secondi a quindici minuti}; raramente di
più. Chi non ne ha mai dato uno dovrebbe sforzarsi di darne almeno uno, magari
positivo per cominciare. Ogni feedback dato o ricevuto è un
\textbf{deposito nel proprio ``fondo fiduciario''} (\emph{trust fund}), che
costruisce la reputazione futura. Il relatore si è costruito la reputazione di
chi dà feedback il più possibile oggettivi e analitici: giusti o sbagliati che
siano, costringono a pensare. Non impone mai il suo punto di vista (``si fa
come dico io''), ma chiede ``tu come faresti? Come posso aiutarti?''. Il
feedback, per lui, serve a \emph{far crescere qualcun altro}.

\subsection{I manager peggiori: ``a che numero sto pensando?''}

I peggiori manager che il relatore abbia avuto erano quelli che volevano che si
facesse come volevano loro, anche quando chiedevano ``tu come la faresti?''.
In realtà la domanda era ``\emph{a che numero sto pensando?}'': se dicevi 5,
``no, dovremmo fare 9''; se dicevi 9, ``no, 5''. Lo stress costante di dover
indovinare cosa vuole il capo nasce dalla mancanza di fiducia (``se mi chiedi
come fare, fammelo fare'').

Lo stesso gli succede con i CEO che lo chiamano per ``sistemare la cultura
aziendale'' e poi dicono ``non sono sicuro, io farei diversamente''. La sua
risposta: ``lo so che faresti diversamente, per questo mi hai chiamato''. Se si
chiama un esperto, bisogna lasciargli fare il suo lavoro: non si va dal
meccanico a cambiare le gomme per poi dirgli ``non dovresti girare un po' di
più quel bullone?''. Chiedere un parere a un esperto significa dargli anche il
potere di scegliere. Il problema è il \emph{micromanagement}, o in generale il
non accettare punti di vista diversi dal proprio (magari perché si è il
fondatore e si pensa di essere l'unico a saper fare le cose).

\subsection{Dare feedback a un superiore egocentrico}

\paragraph{Come si dà un feedback negativo a un superiore con un grande ego?}
Il trucco è \textbf{non puntare sull'ego, ma sull'azione e sui suoi effetti}.
Non ``tu fai così e io mi sento così'': anzi, si alimenta il suo ego
rendendolo il \emph{main character} dell'avventura, l'eroe che risolve il
problema (``tu, grande manager, aiutaci a risolvere questo problema che
affligge noi che facciamo l'operatività''). Non è l'ideale (``io cambierei
manager''), ma pochissimi possono scegliersi i colleghi: il relatore stesso ha
quasi sempre ricevuto team preassemblati (``ecco le chiavi, questa è la tua
macchina''). Bisogna giocare con le carte che si hanno in mano.

\paragraph{Un esempio pratico.} Con un manager che fa micromanagement e mette
sempre ``il becco'' nel proprio lavoro, si può dire: ``se ti fidi di me per
questo lavoro, con le tue capacità potresti dedicarti a questa cosa più
strategica e di lungo termine: pensa quanto sarebbe più utile per l'azienda''.
In sostanza: ``guarda cosa potresti raggiungere se mi lasciassi fare il mio
lavoro''. Non come un ``tu fai il tuo e io faccio il mio'', ma come un
``ti sollevo da questa parte del carico, che so che ritieni importante: lascia
che ci pensi io, così puoi dedicarti ad altro'', che magari gli dà anche più
visibilità.

Si gioca su ciò che quella persona pensa di saper fare, e che probabilmente sa
anche fare. Anche Michael Scott di \emph{The Office}, per quanto preso in giro,
è raccontato come un venditore bravissimo, che ha chiuso i contratti più
importanti per l'azienda, e rispetto ai manager che gli succedono è forse
``uno dei meno peggiori''. Tutti, anche i manager peggiori, hanno punti di
forza: \textbf{giocate su quelli}.

% ---------------------------------------------------------------------
\section{Il feedback come regalo}
\label{sec:z1-regalo}

Per chiudere, il relatore risponde alla domanda d'apertura, che aveva lasciato
apposta in sospeso: che cos'è il feedback per lui?

\begin{quote}
  \emph{Il feedback è un regalo che mi viene fatto da qualcuno che so che vuole
  il mio interesse.}
\end{quote}

Ha costruito un team di persone che si fidano di lui e che non hanno motivo di
sabotarlo o di ferirlo: per questo interpreta anche le cose peggiori con il
default to trust. È un regalo, e cerca di prenderlo in considerazione il più
possibile.

Ci è arrivato partendo da se stesso. Oltre a essere avverso ai conflitti, si
definisce ``avverso alla comunicazione umana'': paradossalmente, una persona
introversa e timida finita a fare comunicazione, marketing e ora HR (``succede
spesso''). Prima di ordinare una pizza al telefono ha dieci minuti di
preparazione mentale, e se l'ordine delle domande cambia va in confusione. Per
lui un feedback di dieci minuti sono spesso \textbf{tre ore di preparazione}:
costruisce un \textbf{albero dei dialoghi}, come in un'avventura grafica
personale, con tutte le opzioni. Poi spesso finisce con ``hai fatto questa
cosa, andrebbe migliorata'' ``hai ragione, grazie'', e basta.

Così ha pensato: chissà quanti si fanno gli stessi problemi con lui, e passano
ore a prepararsi per dirgli ``questa cosa l'hai fatta male'', sapendo che
riceveranno comunque quell'occhiata fugace del ``come ti permetti''. Eppure
vengono a dirglielo. Non può far altro che rispondere: ``grazie, hai ragione,
cercherò di fare meglio''. È così che ha metabolizzato il ricevere feedback:
capendo che non c'è nessuna cattiva intenzione e che, se non proprio felice
(``ci ho messo una pietra sopra''), può almeno essere \emph{grato} che ci sia
ancora qualcuno che ha così a cuore il suo interesse da dargli feedback.

\paragraph{Il silenzio fa più paura.} Non bisogna spaventarsi dei feedback
negativi che emergono dai sondaggi aziendali: bisogna spaventarsi del
\textbf{silenzio}, quando nessuno dice più nulla. Significa che nessuno si
fida più che si sappiano risolvere i problemi.

\paragraph{La sindrome dell'impostore.} Il suo motivatore personale, all'inizio,
era soprattutto ``non voglio fallire'': per questo cercava di migliorare, e
migliorare è diventato accettare feedback, perché più feedback riceve, più
migliora. Il feedback è un \textbf{acceleratore}, e lui cerca di incentivare
gli altri a darglielo. Resta però sempre la vocina del ``e se si accorgono che
non sei capace?''. Con la \emph{impostor syndrome} convive da anni e
continuerà a conviverci; se ce l'avete, conviveteci anche voi. Avrete sempre
la sensazione che ``se ne accorgeranno'': forse sì, ma sapete fare anche
altre cose. Non concentratevi solo su ciò che non sapete fare. Gli errori
capitano: la perfezione non è di questo mondo.

\paragraph{Esercitarsi in accademia.} Per fare pratica con il feedback,
l'ambiente migliore è quello dei colleghi di corso: persone nella stessa
situazione, con esperienza simile e lo stesso obiettivo, in un contesto in cui
non si può essere licenziati e l'unico effetto a lungo termine è ``ho parlato
con questa persona, ora le cose vanno meglio''. E bisogna avere il coraggio di
parlare anche con i docenti (``dire a Diego: secondo me questa lezione potevi
farla in questo modo''; ``so che è una persona che lo accetta''), senza farsi
bloccare dalla gerarchia. Nell'azienda giusta (come quella in cui lavora ora) il
relatore va sempre dal CEO a dare feedback, sapendo che non verrà licenziato:
il CEO gli ha detto ``ti ho preso per darmi feedback''.

La frase di chiusura:

\begin{quote}
  \emph{Per superare le vostre difficoltà, tutto quello che vi viene chiesto è
  il coraggio di chiedere aiuto.}
\end{quote}

Chiedere aiuto non è una debolezza: è un segno di consapevolezza di sé
(\emph{self-awareness}) e di forza personale. Se tutti lo facessero, i
progetti andrebbero molto meglio; spesso invece ci si fa frenare dal ``non
voglio fare brutta figura'', ``non voglio passare per quello scemo'', e sono
proprio queste paure a rallentarci e a impedirci di raggiungere gli obiettivi.

% ---------------------------------------------------------------------
\section{Domande e risposte finali}
\label{sec:z1-qa-finale}

L'ultima mezz'ora è dedicata a domande libere.

\paragraph{E se una persona peggiora per uno o due anni di fila?} Se nessuno
gliel'ha mai detto, c'è un grosso problema, non suo ma dell'azienda. Se invece
c'è consapevolezza, il percorso è: ``cosa vuoi fare per risolvere?''; ``ho
provato questo, questo e questo''; ``visto che non ha funzionato, vuoi un mio
consiglio? Proviamo con quello?''; ``che supporto ti serve?''. Se, con tutto il
supporto, il tempo e le soluzioni concordate per mesi, il miglioramento non
arriva, per il relatore la soluzione è: \textbf{``dimmi un prezzo per andare
via''}. Preferisce pagare una buonuscita (\emph{severance package}) che
avviare la strada dei richiami formali sulla performance fino al
licenziamento per giusta causa. A un certo punto andare avanti non è un favore
per nessuno.

\paragraph{È vero che nei videogiochi si viene sfruttati? Meglio un piccolo
studio italiano?} Per il relatore si è quasi più a rischio a Blizzard che in un
piccolo studio di Roma o Torino. Nel piccolo studio si sa dove ci si va a
cacciare: non ci sarà il team building a Santo Domingo, ma ci si concentra su
un gioco alla volta e, se va bene, si fa il successivo. A Blizzard si può
lavorare su \emph{Overwatch}, che vende milioni, e perdere comunque il posto,
in modo totalmente scollegato dal proprio lavoro. Gli stipendi più alti
all'estero compensano anche questo rischio, che spesso non dipende da come va
il gioco (anzi, ``il gioco non ha fatto come volevamo'' è spesso una scusa per
licenziamenti che non avevano motivo).

\paragraph{Cambiare settore a trent'anni, da web developer full remote?} Lo
stipendio indicato dallo studente non è fuori dalla portata di vari studi
italiani, anzi è superabile. Non va esclusa la possibilità di lavorare da
remoto per aziende piccole ma estere. Alcuni amici del relatore, pagati con
stipendi italiani, si sono trasferiti a Bali o in paesi con un costo della vita
più basso; il collega di \emph{PUBG} percepisce uno stipendio olandese vivendo
a Milano, che per quanto cara non costa come i Paesi Bassi.

\paragraph{Come si diventa manager?} Dipende dall'azienda: alcune promuovono i
migliori (processo tendenzialmente sbagliato), altre propongono un percorso di
\emph{people management}. Diventare manager significa \textbf{abbandonare ciò
che si sa fare per imparare a fare altro}. Il relatore ha dovuto accettare che
non sarebbe più stato ``il miglior community manager italiano del mondo''
(``eravamo in pochi, quindi era facile''): dedicandosi alle persone non poteva
più fare il suo lavoro bene come prima. È stata una negoziazione con se
stesso: fare il manager significa non essere più l'esperto di qualcosa, ma
saper gestire e far crescere persone. Significa anche \textbf{cedere il
controllo della propria carriera}: se il team decidesse di non fare più nulla,
il manager ne pagherebbe le conseguenze. Il suo lavoro diventa scegliere le
persone giuste e farle lavorare bene. Lui lo esplicita sempre: sa di affidare
al team parte della propria carriera, e il team sa che lui combatterà per
loro, anche prendendosi le colpe.

\paragraph{L'inglese.} \textbf{Va assolutamente imparato}, anche lavorando in
Italia. Non saperlo preclude circa il 95\% del mercato (tutto quello fuori
dall'Italia) e soprattutto il \emph{networking} con persone che potrebbero dare
occasioni future: l'inglese è spesso l'unica lingua in comune con chi lavora
all'estero ed è oggi la lingua universale dei videogiochi.
\begin{itemize}
  \item \emph{Che livello serve?} Non importa B2 o C2: bisogna saper sostenere
        una conversazione, guardare almeno Netflix con sottotitoli in inglese,
        comunicare senza dodici pause e senza ``come si dice\dots''. Il
        certificato non verrà mai chiesto: verrà chiesto se si sa parlare
        inglese. Il relatore ha sentito persone con un C2 che non sapevano
        parlarlo.
  \item \emph{Come allenarlo senza occasioni di conversazione?} Netflix, Prime
        Video, Now TV, Udemy, Coursera, YouTube, Duolingo, siti di scambio
        linguistico con madrelingua, Discord: non sono mai esistiti così tanti
        mezzi gratuiti. Il relatore ha imparato l'inglese giocando (con una
        pronuncia sbagliata, perché leggeva senza sentire, su cui ha poi
        dovuto lavorare); guardare un film in inglese con i sottotitoli in
        italiano insegna anche la traduzione.
  \item Oltre all'ascolto passivo è fondamentale \textbf{parlare}: in una
        flotta di \emph{EVE Online} con 150 persone quasi tutte madrelingua
        inglese si è esposti a decine di accenti e modi di dire (mille modi di
        dire ``come stai'' oltre a \emph{how are you}); lavorando all'estero si
        impara a distinguere gli accenti (Manchester, Londra) e ad apprezzarli
        anche nei media.
\end{itemize}

\paragraph{Quanto deve sapere un junior?} Dipende da quanto è tecnico il
ruolo. Un game designer junior che non sa qualcosa può essere affiancato da un
senior; un programmatore che non sa programmare, in quel momento, non è utile
all'azienda.

\paragraph{Annunci ``junior con cinque anni di esperienza''.} ``Cercasi
ventenne con trent'anni di esperienza.'' Non è l'annuncio per un junior: è un
errore dell'azienda, che in realtà cerca un \emph{mid}. Riceverà profili
junior, si stupirà che non abbiano cinque anni di esperienza, rifarà l'annuncio
e avrà perso sei mesi. Il consiglio: tenere d'occhio quella pagina, perché
dopo qualche mese l'annuncio verrà corretto.

\paragraph{Una struttura per il post-mortem.} A chi chiede un formato di base,
il relatore anticipa quanto vedrà nella seconda giornata a proposito di
brainstorming: una lavagna \textbf{divisa in quattro aree}, su cui ognuno
attacca tutti i post-it che vuole:
\begin{center}
  \begin{tabular}{@{}p{0.44\textwidth}p{0.44\textwidth}@{}}
    \toprule
    \textbf{Smettere di fare} & \textbf{Iniziare a fare} \\
    ciò che non ha funzionato & ciò che non abbiamo fatto e andava fatto \\
    \midrule
    \textbf{Fare di meno} & \textbf{Fare di più} \\
    ciò che ha fatto perdere tempo rispetto all'impatto ottenuto &
    ciò che ha funzionato ma, fatto con troppa leggerezza, non ha dato i
    frutti sperati \\
    \bottomrule
  \end{tabular}
\end{center}
Con i post-it ognuno ha voce e tutto risulta catalogato; per il report basta
una foto della lavagna (o uno screenshot, se la lavagna è digitale, come Miro).

\paragraph{Come sbloccare la paralisi del feedback verso il management?} Uno
studente ha visto manager dire ``non abbiate paura, la mia porta è sempre
aperta'' senza che cambiasse nulla. C'è una dissonanza tra ciò che dicono e
ciò che fanno: è ``la mamma che dice vieni qua che non ti faccio niente''.
Qualcosa deve aver fatto pensare alle persone che non è sicuro. Prima va
sistemato quello; poi si può provare con i \textbf{feedback anonimi}.
L'obiettivo del relatore resta il feedback aperto e firmato, ma a Stormind ha
iniziato con feedback anonimi verso il management per verificare se
l'anonimato fosse un fattore. Non lo era (evidentemente non c'erano problemi a
parlare con il CEO), e la volta successiva intende abbandonarli. Ma è un buon
\emph{check}: se senza anonimato non arriva nulla e con l'anonimato arrivano
cento feedback, il problema è chiaro.

\paragraph{Quando è giusto rivolgersi a un superiore?} La strada giusta resta
parlare direttamente con l'interessato. Ma se si prova e non si trova una
soluzione, la cosa più ovvia è andare dal proprio manager insieme: ``ci
aiuti?''. Quando il relatore parlava di mediare i dialoghi, era proprio questo:
fare da moderatore, aiutare a sbloccare i punti fermi o, se non si riesce a
scegliere tra A e B, scegliere lui. L'escalation, quando necessaria, va fatta.

\paragraph{Il tempo libero veniva rispettato?} No, e se l'è dovuto imporre
lui. Il lavoro gli piaceva così tanto che continuare a casa non gli sembrava
lavorare. Gli si è ritorto contro con il lancio di \emph{Diablo~III}, fatto in
orario americano (le due o l'una di notte in Europa), con l'ufficio chiuso
perché a nessuno era stato chiesto di restare. Il lancio va malissimo, i
giocatori protestano, e lui, sveglio, risponde sui forum fino alle tre, quattro,
cinque del mattino: ``non posso risolvere il problema, ma resto qui ad
aggiornarvi con quello che riesco a sapere dal team di sviluppo''. Il giorno
dopo trova sulla scrivania un \textbf{richiamo formale dell'HR}: aver lavorato
fuori orario senza autorizzazione, condotta non professionale. Per lui è stato
un interruttore: ``ora me ne vado''. Due settimane dopo faceva colloqui con
Riot. Ancora oggi è contrario a quel richiamo, ma gli ha fatto capire come
ragiona un'azienda: ``hai un contratto 9--18, non mi interessa se l'hai fatto
con le migliori intenzioni''. Eppure per anni ha goduto di un'ottima
reputazione in quelle community come ``quello che è rimasto con noi fino alle
cinque del mattino'', e anche l'azienda ne ha beneficiato.

Ci sono poi aziende in cui, soprattutto per i developer, viene richiesto il
famigerato \textbf{crunch}: lavoro straordinario continuativo, senza pause.
Dipende molto dall'azienda; sta venendo (``grazie al cielo'') deprecato, perché
molte aziende capiscono che non è sano ed è visto male anche dal pubblico, ma
esiste ancora.

\paragraph{E se il tuo capo demotiva qualcuno del tuo team?} Qui ``diventa
personale'': il team del relatore si fida di lui perché sa che in questi casi
``diventa una bestia''. ``Te la puoi prendere con me quanto vuoi, ma il mio
team non lo tocchi.'' Lo gestisce direttamente con l'altro manager, per quanto
brutta possa diventare la conversazione. Lo sforzo più grande è \textbf{non
andarne poi a parlare con il team} (``ragazzi, ho fatto questa cosa''), perché
minerebbe la fiducia verso l'alto: non vuole un rapporto ottimo fino a lui e
pessimo oltre. L'obiettivo è che l'altro manager ``si faccia perdonare''
cambiando il proprio comportamento, e che il team percepisca da solo il
cambiamento.

\paragraph{Chiusura della giornata.} Il relatore chiede un feedback collettivo
in chat (``se domani ci dai \emph{more of the same} siamo contenti, oppure
ditemi cosa limare''), si scusa per i ripetuti crash della connessione e
ringrazia per la proattività nelle domande, che gli servono anche come
\emph{check}: se nessuno chiede nulla su un argomento, significa che può
andare avanti. Racconta di aver passato una settimana a decidere che cosa
mettere nelle slide e di averle poi preparate in poco tempo (rifacendone
alcune fino alle tre di notte). Ricorda infine di usare il test di
personalità per sé e per il proprio gruppo, facendolo fare a chi non l'ha
ancora fatto: tornerà molto utile per cooperare meglio.

% ---------------------------------------------------------------------
\section{Anticipazioni sulla seconda giornata}
\label{sec:z1-domani}

La seconda giornata del workshop sarà più pratica: ``oggi vi ho spiegato
perché, domani vi insegno come lo fanno queste aziende''. Gli argomenti
previsti:
\begin{description}
  \item[Time management] come organizzarsi una giornata professionale in
        azienda: come strutturare una \emph{task list}, cosa fare e cosa non
        fare, come e cosa prioritizzare.
  \item[Brainstorming] i gruppi dovranno fare moltissimi meeting (come facciamo
        il gioco, come lo chiamiamo, che scopo gli diamo). Se non si sanno fare
        bene, i brainstorming sono \textbf{dannosi ed eterni}: un brainstorming
        dovrebbe durare \textbf{al massimo due ore}. Verranno mostrati metodi
        che accorciano la durata, massimizzano l'effetto e, soprattutto,
        \textbf{fanno partecipare tutti}: non c'è niente di peggio di un
        brainstorming in cui parlano in due per il 90\% del tempo. Il classico
        ``ragazzi, non esistono idee sbagliate, sparate'' è il modo peggiore di
        farlo; non è colpa di nessuno, ma si può fare diversamente.
  \item[Guidare da dentro] come essere leader nel proprio team anche senza
        esserne il leader: sviluppare responsabilità personale e verso gli
        altri. Non essere il manager di qualcuno non significa non esserne
        responsabili: la persona accanto è responsabilità di tutti, perché si
        è un team e non un insieme di individui. ``Se qualcuno cade, tutti lo
        facciamo rialzare: non c'è l'addetto al rialzamento.'' È la versione
        \emph{advanced} della giornata: come comportarsi da manager senza
        averne il titolo.
\end{description}
Torneranno inoltre il \emph{service level agreement} e il tema della
condivisione dei successi e della protezione del team dai fallimenti.

% ---------------------------------------------------------------------
\section*{Punti chiave del capitolo}
\addcontentsline{toc}{section}{Punti chiave del capitolo}

\begin{itemize}
  \item ``La capacità di un gruppo di persone di fare cose straordinarie
        dipende dal modo in cui queste persone operano insieme.'' Il 90\% dei
        problemi dei team sono \textbf{problemi di comunicazione}.
  \item Le aziende si liberano sempre più presto dello \textbf{useless
        genius}: le competenze si trasmettono, il carattere no. Il
        \textbf{fit culturale} va valutato anche dal candidato.
  \item Non esiste un unico modo di dare feedback: dipende dal contesto e
        soprattutto da \textbf{chi lo riceve}. Il feedback è una somma di
        piccoli miglioramenti periodici; tenersi le cose nel cassetto non
        serve mai.
  \item I team attraversano \textbf{forming, storming, norming, performing e
        adjourning}. Lo storming non si salta; si passa la maggior parte del
        tempo in forming; la bravura sta nell'accelerare i passaggi e nel
        fare bene il \textbf{post-mortem}, anche dei successi.
  \item Gli obiettivi di un progetto vanno fissati \textbf{all'inizio}: solo
        così si può valutare se è andato bene e scegliere tra A e B.
  \item Il circolo virtuoso è \textbf{feedback $\rightarrow$ miglioramento
        $\rightarrow$ fiducia}. Accettare un feedback non significa fare come
        dice l'altro: decide chi ha la responsabilità, ma deve argomentare.
  \item ``Aiutami a poterti aiutare'': meglio dire in anticipo ``potrei non
        farcela'' che ``non ce l'ho fatta''. Chiedere aiuto non è una
        debolezza.
  \item Quando si dà feedback \textbf{non si porta la soluzione}: la si
        costruisce insieme a chi lo riceve.
  \item La cultura è forte quando si lavora \textbf{per gli altri}; il lavoro
        è valutato come opera d'insieme (anche con bonus di team).
  \item Come in \emph{Among Us}, i team sprecano energie a cercare il
        traditore invece di fare le task. \textbf{Default to trust}: chi ha
        sbagliato ha agito con le informazioni che aveva. Cercare
        \textbf{la causa, non il colpevole}: ``che cosa ti ha fatto pensare
        che fosse la scelta giusta?''.
  \item Carriere e performance non sono lineari: vanno valutate le
        \textbf{traiettorie}, non un punto dell'anno. Le valutazioni puntuali
        generano \emph{lying, hiding, faking}.
  \item Confermare sempre le istruzioni facendosele ripetere; ``non ho
        capito'' non è una colpa.
  \item Dare feedback: sull'\textbf{azione} e sui suoi \textbf{effetti}, specifico,
        presto, uno a uno, adattato al carattere, senza terze parti, senza
        ``non mi piace'', senza offendersi se non viene seguito; poi fare il
        debriefing.
  \item Ricevere feedback: \textbf{ascolta, chiedi dettagli, accetta, agisci}.
        Il feedback positivo serve a non peggiorare ciò che si fa bene; il
        feedback accidentale e il \emph{drive-by} non servono a nulla.
  \item Feedback \textbf{frequenti e recenti}: tanti piccoli \emph{sanity
        check} costano meno di una grande correzione di rotta.
  \item Il feedback è \textbf{un regalo} di chi ha a cuore il nostro
        interesse; fa paura il silenzio, non le critiche.
\end{itemize}

%    % ... altre lezioni zero ...
%    \part{Il corso di Game Design}
%    \renewcommand{\thechapter}{\arabic{chapter}}
%    \setcounter{chapter}{0}
%    % =====================================================================
%  Capitolo 1 — Introduzione al corso e al mestiere del game designer
%  Fonte: trascrizione della lezione GD1 (Game Design, primo anno)
%  Pacchetti richiesti nel preambolo del libro:
%    \usepackage[italian]{babel}, \usepackage[T1]{fontenc},
%    \usepackage[utf8]{inputenc} (solo pdfLaTeX), \usepackage{booktabs}
% =====================================================================

\chapter[Il mestiere del game designer]{Il mestiere del game designer: introduzione al corso}
\label{cap:introduzione-game-design}

Questo primo capitolo fa da soglia. Prima di affrontare metodi, framework e
documenti, è necessario capire \emph{in che modo} si impara il game design,
\emph{con chi} lo si fa, \emph{con quali strumenti} e, soprattutto,
\emph{che cosa fa concretamente} un game designer nel suo lavoro quotidiano.
Il capitolo segue l'impostazione della lezione inaugurale del corso: si parte
dal contesto formativo (il corso come simulazione di lavoro), si passa al
metodo di apprendimento (l'\emph{inquiry-based learning}), si elencano
letture e risorse di riferimento, e si arriva infine a una prima definizione
operativa del ruolo del designer, dei limiti entro cui lavora e delle sue
possibili specializzazioni.

Una premessa terminologica: il game design è una disciplina giovane e quasi
tutta la sua letteratura è scritta in inglese. Per questo motivo molti termini
tecnici verranno lasciati in lingua originale (\emph{pipeline},
\emph{pitching}, \emph{tuning}, \emph{feature team}, \dots): sono i termini che
si useranno nel lavoro, ed è importante imparare a riconoscerli e a usarli
correttamente fin da subito.

% ---------------------------------------------------------------------
\section{Il contesto: un corso come simulazione di lavoro}
\label{sec:simulazione-lavoro}

\subsection{Il principio fondamentale}

Il primo concetto da fissare è che il corso \textbf{si presenta come una
simulazione del lavoro}. Dal momento in cui si entra in aula si viene trattati
come \emph{potenziali professionisti dell'industria}: non soltanto studenti,
ma lavoratori effettivi di uno studio (l'accademia), con una gerarchia, dei
referenti, dei compiti e delle scadenze.

Ne derivano alcune conseguenze pratiche:
\begin{itemize}
  \item il docente non assegna un ``compitino'' da risolvere per accumulare
        punti o presenze: ogni \emph{task} è pensato per far dedicare tempo
        all'apprendimento di qualcosa di concreto;
  \item non consegnare un task non danneggia il docente (che ``vive ancora
        domani e viene pagato ugualmente''), ma lo studente, che perde
        l'occasione di mettere le mani su un argomento e diventare più abile;
  \item docenti e studenti devono trattarsi \textbf{da colleghi}, perché è
        molto probabile che un domani lo diventino davvero.
\end{itemize}

\subsection{Perché l'impegno conta più che mai}

Il lavoro del game designer è descritto come \emph{molto romantico, ma anche
molto duro}, e sempre più duro col passare degli anni. Durante il periodo del
Covid l'industria ha visto grandi investimenti; quei ``rubinetti'' sono stati
chiusi da tempo, il rapporto con i publisher è diventato più complesso e le
opportunità lavorative non sono più quelle di tre o quattro anni fa. Questo
significa:
\begin{itemize}
  \item più competizione per ogni posizione;
  \item aspettative più alte: ci si aspetta che i candidati siano ancora più
        preparati e sicuri delle proprie competenze;
  \item un unico modo per diventarlo: \textbf{lavorare}. Più tempo si dedica
        oggi ad apprendere il mestiere, più facilmente si sarà assumibili
        domani e più agevole sarà il lavoro di tutti i giorni.
\end{itemize}

\subsection{Sbagliare qui per non sbagliare fuori}

Il mantra del corso è: \textbf{quanto più provate qui, anche in modo goffo e
sbagliando, tanto meno sbaglierete là fuori}. In accademia, per quanto la
simulazione sia realistica, nessuno toglie lo stipendio e nessuno licenzia da
un momento all'altro. Nell'industria sì, e un errore può costare molto denaro
all'azienda, ancora di più se l'azienda è la propria.

\subsection{Il valore del fallimento}

Il fallimento, se affrontato con la mentalità corretta, è probabilmente
\emph{la cosa più vantaggiosa} che si possa fare durante la formazione:
\begin{itemize}
  \item un successo (``avete superato questa deadline'') conferma le proprie
        abilità, ma insegna relativamente poco;
  \item dopo $n$ fallimenti si impara moltissimo, e i successi arrivano;
  \item è fondamentale non demoralizzarsi: il feedback negativo non è mai un
        giudizio sulla persona, ma su ciò che è stato consegnato.
\end{itemize}

Il docente riconosce apertamente che dedicare tempo ed energia a un progetto e
vederlo rifiutato, o vederlo non uscire mai, è demotivante: è successo a lui e
lo ha visto accadere anche su progetti molto grandi, con team numerosi. Il
punto però non è ciò che si è ottenuto materialmente, ma \textbf{quanto si è
cresciuti}. Da ogni esperienza ci si portano a casa due cose:
\begin{enumerate}
  \item \textbf{l'esperienza}: quanto si è diventati bravi, veloci,
        praticamente abili nell'applicare un framework, quanto si è cresciuti
        come designer;
  \item \textbf{il rapporto umano}: nel lavoro si ``suda sangue'' fianco a
        fianco con altre persone, e questo crea legami molto forti che
        portano avanti l'industria nel senso più positivo del termine: nascono
        progetti privati, ci si porta da un'azienda all'altra, si formano team
        con basi solide.
\end{enumerate}

Va evitato anche l'eccesso opposto: la cultura del \emph{grinding} si è
appropriata di questi discorsi in modo distorto. Non bisogna diventare robot.

% ---------------------------------------------------------------------
\section{Organizzazione del corso e gerarchia}
\label{sec:gerarchia-corso}

\subsection{La struttura gerarchica}

La struttura dei referenti del corso ricalca volutamente una
\textbf{gerarchia tipica dell'industria}:
\begin{itemize}
  \item al vertice c'è il \textbf{direttore del corso} (Diego Ricciuti), che
        corrisponde al \emph{creative director}: supervisiona la visione del
        progetto (il concetto di ``visione'' verrà approfondito più avanti);
  \item sotto di lui ci sono i \textbf{referenti di dipartimento}, uno per
        disciplina, con un corrispettivo per ciascuna sede (Milano e Roma);
  \item in fondo ci sono gli studenti, che occupano il primo gradino di
        un'azienda: l'\emph{intern} (stagista) o il \emph{junior}, il passo
        che precede il lavoratore ``regolare''.
\end{itemize}

\begin{table}[htbp]
  \centering
  \begin{tabular}{@{}p{0.34\textwidth}p{0.29\textwidth}p{0.29\textwidth}@{}}
    \toprule
    \textbf{Disciplina} & \textbf{Milano} & \textbf{Roma} \\
    \midrule
    Direzione (\emph{creative director}) & \multicolumn{2}{p{0.58\textwidth}}{Diego Ricciuti} \\
    Game design      & Mattia Podini        & Giacomo ``Jack'' Danna \\
    Level design     & Stefano Zocchi       & Andrea Dispoli \\
    Narrative design & \multicolumn{2}{p{0.58\textwidth}}{Daniele Fusetto (lezioni condivise), Giacomo Masi} \\
    Programmazione (Unreal) & Alessandro Pregnolato & Matteo Velardi \\
    \bottomrule
  \end{tabular}
  \caption{Referenti del corso per disciplina e sede (nomi come riportati in lezione).}
  \label{tab:referenti}
\end{table}

Per il primo anno, il referente diretto degli studenti di game design è il
docente di game design: è a lui che ci si rivolge per dubbi e domande. Altre
figure utili sono il direttore didattico (severo, ``giustamente severo'', da
contattare solo per questioni particolarmente importanti), l'amministrazione
(contratti e documentazione) e la segreteria.

\subsection{I canali di comunicazione}

\begin{description}
  \item[Discord] è il canale per le comunicazioni ufficiali del corso e per
        le domande che non richiedono urgenza immediata; può avere qualche
        ritardo.
  \item[E-mail] è il canale per tutto ciò che è formale o ``ancora più
        ufficiale'' e per le comunicazioni rapide e immediate.
  \item[Telegram] esistono residui di vecchi canali precedenti a Discord, da
        evitare.
\end{description}

Le domande che valgono per tutti vanno fatte in classe (anche da un portavoce
a nome del gruppo); quelle singole e personali via Discord o, se formali, via
e-mail. Se il docente non risponde a una domanda c'è un motivo, ad esempio
ragioni legali, e lo dirà esplicitamente.

\subsection{I tre moduli e la settimana tipo}

Le tre discipline principali del primo anno (e non solo) sono:
\begin{itemize}
  \item \textbf{Game Design} (tendenzialmente il lunedì);
  \item \textbf{Level Design} (tendenzialmente il mercoledì);
  \item \textbf{Programmazione} (tendenzialmente il venerdì).
\end{itemize}
Nel primo mese le lezioni sono tutte di game design (lunedì, mercoledì e
venerdì). Sono previste anche lezioni di narrativa, ma il piano formativo si
aggiorna ogni anno: le ore di narrativa sono state ridotte a favore di quelle
di programmazione. Le poche lezioni di narrativa daranno quindi dei
\emph{cenni}, non la preparazione di un narrative designer. Ciò che serve
davvero è imparare \textbf{come comunicano gli altri dipartimenti} e come
lavorare con persone che fanno un lavoro diverso ma complementare al proprio.

\subsection{Le due anime delle lezioni di game design}

\begin{description}
  \item[Design methods (metodologie)] lezioni prevalentemente teoriche, ma con
        una buona componente pratica, in cui si studiano metodologie di design
        e i cosiddetti \emph{framework} dell'industria.
  \item[Documentazione e pitching] il cuore del lavoro quotidiano del
        designer.
\end{description}

\paragraph{La documentazione.}
Il docente si definisce ``il più grande nerd della documentazione'' che si
possa incontrare. È una parte dura, perché a pochi piace documentare, ma il
lavoro del designer \emph{è principalmente questo}: scrivere documenti e, da
quei documenti, passare alla pratica e all'implementazione. Chi non sa
documentare bene parte male. Quando non si è presenti, il documento è il
proprio \textbf{biglietto da visita}: qualcun altro lo leggerà, quindi va
scritto nel modo più chiaro e comprensibile possibile.

Un esempio concreto del perché la documentazione conta: CD Projekt Red ha
sviluppato un proprio engine proprietario, ma buona parte della sua
documentazione è stata ``buttata all'acqua''. Chi deve usarlo o lo sa già
usare, o ha competenze di un altro engine da trasferire, oppure è in
difficoltà. Si documenta \textbf{perché altri possano leggere}.

\paragraph{Il pitching.}
Il \emph{pitching} è l'esposizione, la presentazione orale. La parte creativa
del team presenta spesso:
\begin{itemize}
  \item le proprie idee ai colleghi;
  \item il gioco a convention ed eventi;
  \item il pitch del gioco agli investitori, per ottenere i fondi necessari a
        realizzarlo.
\end{itemize}
Saper parlare e comunicare chiaramente le idee è fondamentale. In sintesi:
\emph{il game designer parla}, e bisogna farsene una ragione.

\subsection{Altre componenti del corso}

\begin{description}
  \item[Gioco critico e playtest] la classe viene divisa in due gruppi: uno
        gioca, l'altro osserva. Chi osserva non applaude ``come le
        scimmiette'': annota, segna e analizza il \emph{playtest}, esattamente
        come si fa nell'industria. Il playtest fa parte delle mansioni del
        designer, ma esiste anche come specializzazione autonoma.
  \item[Pipeline] termine che indica tutto ciò che riguarda i processi
        formalizzati dell'industria; di solito si intende la \emph{pipeline di
        produzione}: come si divide il team, qual è la sua gerarchia, come il
        \emph{producer} gestisce le deadline. In breve: \emph{come} e
        \emph{quando} si fanno le cose. Al primo anno è meno rilevante (si
        dovrà ``nuotare'' un po' da soli); dal secondo anno si vedranno le
        metodologie effettivamente applicate ai progetti digitali. Si
        scoprirà quanto sia difficile lavorare a un progetto digitale, e
        ancor più fisico, senza un piano di sviluppo comprovato, come quelli
        usati per qualsiasi software in qualsiasi parte del mondo.
  \item[Engine] il motore di riferimento è \textbf{Unreal Engine}
        (sezione~\ref{sec:engine}).
  \item[Programmazione] principalmente tramite \emph{Blueprint} in Unreal,
        più eventuali sistemi di versionamento; il corso parte da zero.
  \item[Progetti pratici] probabilmente la cosa più importante. Al primo
        anno i progetti di gruppo restano circoscritti al corso di design,
        perché mancano ancora alcune nozioni (come comunicare, affrontare
        problemi, dare e ricevere feedback). In seguito si allargheranno a
        progetti digitali, molto più difficili di quelli fisici, con membri
        di altri team: un rapporto non unilaterale, ma ``covalente''.
\end{description}

% ---------------------------------------------------------------------
\section{Il metodo: l'inquiry-based learning}
\label{sec:inquiry}

\subsection{Che cos'è}

Il corso è costruito secondo l'\textbf{inquiry-based learning}
(apprendimento basato sull'indagine), un approccio pedagogico nato intorno
agli anni Sessanta, che si richiama a pensatori come Dewey e Piaget. In
generale, le discipline da cui il design attinge sono quasi sempre la
sociologia, la psicologia e, talvolta, l'economia.

Il funzionamento è un ciclo:
\begin{enumerate}
  \item il docente fornisce metodologie e framework dell'industria durante
        la lezione;
  \item lo studente, fuori dalle ore di lezione, lavora sui \emph{task}:
        mette a terra le nozioni, le maneggia, le assorbe meglio;
  \item da questo lavoro nascono nuovi dubbi;
  \item i dubbi tornano in aula: lo studente chiede, il docente risponde
        (ammesso che dare la risposta sia utile, e ammesso che una risposta
        esista).
\end{enumerate}

\subsection{Non sempre esiste una risposta giusta}

Non tutto ciò che si studia ha ``la risposta migliore possibile in assoluto
nel vuoto''. Esistono metodologie, limiti e framework che aiutano il lavoro,
ma non sono fissi: cambiano da azienda ad azienda, da lavoro a lavoro, da
progetto a progetto, e a volte vengono stravolti. Sta alla
\textbf{sensibilità del designer} capire quale strumento o framework usare in
un dato momento, e questa sensibilità si impara solo con l'esperienza. Nel
lavoro nessuno dirà ``usa questa cosa'' o ``è meglio quest'altra'': sarà il
designer a scegliere.

Anche l'esperienza del docente ha confini: sei-sette anni di lavoro, un modo
personale di approcciare il design e un certo insieme di progetti (ad esempio
nessun MMO). Oltre quei confini si cercano risposte nella ricerca, studiando
come aziende con esperienza specifica (ad esempio Riot) hanno affrontato
problemi simili e hanno condiviso le loro soluzioni, sempre valutando quali
metodi siano adatti alla dimensione del team, al tipo di progetto e al tempo
disponibile.

\subsection{Sensibilità, non memoria}

Non interessa imparare a memoria un framework o la teoria di un autore. Certe
terminologie vanno conosciute, ma per ragioni di comunicazione e chiarezza.
Ciò che conta è sviluppare la \textbf{capacità di farsi domande}, di farle
agli altri, di fare ricerca e di trovare una propria risposta. È esattamente
ciò che si fa nel lavoro: davanti a qualcosa che non si sa fare, l'abilità
del designer sta nel cercare come qualcun altro ha risolto un problema simile.

\subsection{Le lezioni come corsi di aggiornamento}

Le lezioni frontali simulano i \emph{corsi di aggiornamento} aziendali: le
sessioni che un lead di design tiene al suo team di design (o il lead di
level design al suo team, e così via). Alla fine dell'aggiornamento c'è quasi
sempre un lavoro da fare, un obiettivo da portare a terra:
\textbf{l'obiettivo della settimana}. Si lavora di settimana in settimana, e
il perché verrà chiarito in seguito.

Struttura tipica di una lezione:
\begin{enumerate}
  \item breve sessione di domande e risposte sulla lezione precedente;
  \item analisi di eventuali esercizi particolarmente critici;
  \item lezione vera e propria;
  \item discussione di gruppo, se c'è tempo;
  \item assegnazione del task della settimana.
\end{enumerate}
Il docente tende a ``parlare e straparlare'', ad aprire parentesi e a saltare
da un argomento all'altro: l'unico modo per fermarlo è che siano gli studenti
a interromperlo con domande.

\subsection{La gerarchia della competenza}

L'inquiry-based learning può essere rappresentato come una scala di
competenza:
\begin{enumerate}
  \item la sola lezione frontale porta a \textbf{conoscenze incomplete},
        perché manca la parte pratica e di ricerca;
  \item il lavoro a casa (task, ricerca) porta a una competenza \textbf{più
        completa}, non sbagliata in sé, ma ancora parziale e complementare ad
        altre;
  \item la \textbf{condivisione} dei task e dei feedback con gli altri porta
        idealmente alla competenza più completa possibile;
  \item il picco della competenza è la sua \textbf{applicazione nei
        progetti} (fisici o digitali, di gruppo): ``con le cose che avete
        imparato e ricercato, riuscite a fare un gioco?''. Quando la risposta
        diventa positiva, o comunque si inizia a navigare nei processi di
        design per realizzarne uno, quella è la vittoria.
\end{enumerate}
Non si parte subito creando giochi: si parte dall'analisi, dal design e dalle
regole generali di design, prima ``nel vuoto'', poi contestualizzando i
framework, infine unendo i puntini nei progetti.

Rimarranno comunque delle lacune, per tutti, docente compreso, perché il game
design \textbf{non è una scienza esatta}, anzi, a rigore non è affatto una
scienza. Prende in prestito metodologie dalle discipline scientifiche, ma
contiene una componente imprevedibile e incalcolabile: \textbf{la
produzione}. Si possono applicare perfettamente le regole di un framework, e
quel framework può comunque non essere adatto al gioco, al proprio modo di
sviluppare o al tempo di produzione disponibile (alcuni framework richiedono
una ricerca e un'attenzione che il tempo non concede).

C'è poi la componente creativa, la parte ``romantica'' del lavoro: esiste, ma
è \textbf{metà del lavoro}. L'altra metà è ricerca, metodo scientifico al
100\%. Dagli stessi dati $n$ designer diversi, con sensibilità diverse,
possono ottenere $n$ risposte diverse. Non tutte saranno valide: vanno
applicate, messe a terra, e questa scrematura fa parte del lavoro.

\subsection{I livelli dell'apprendimento}

Un secondo modello, a gradini, descrive i livelli di apprendimento (è la
cosiddetta tassonomia SOLO):
\begin{description}
  \item[Pre-strutturale] il docente dà il compito e lo studente lo risolve,
        possibilmente il prima possibile. È esattamente ciò che \emph{non} si
        vuole.
  \item[Uni-strutturale] si affronta una sola disciplina: è il primo passo
        necessario, ed è il motivo per cui nel primo mese ci sono solo lezioni
        di game design.
  \item[Multi-strutturale] entrano in gioco altri dipartimenti: level design,
        game design, programmazione e così via.
  \item[Relazionale] si capisce che alcune cose valgono per più
        dipartimenti: alcuni framework di narrativa derivano dal design o
        possono essere usati per fare design; le regole di documentazione si
        applicano anche ad altri dipartimenti; le regole di level design si
        mescolano con quelle di game design. Dopotutto, per fare un gioco
        servono molti dipartimenti.
  \item[Astratto esteso] il punto finale: muoversi autonomamente e capire
        cosa usare e quando, sulla base della propria esperienza. L'unico modo
        per arrivarci è fare pratica e fare progetti.
\end{description}

% ---------------------------------------------------------------------
\section{Task, feedback e regole del corso}
\label{sec:task-feedback}

\subsection{I task come opportunità}

I task non sono strettamente obbligatori (esiste un numero minimo di
consegne), ma sono \textbf{opportunità}. Quanto e come ci si lavora dipende
dallo studente. Un task fatto in cinque minuti, o in cui si è stati ``furbi'',
non verrà letto né riceverà feedback.

Domande da porsi davanti a ogni task:
\begin{enumerate}
  \item \textbf{Perché} mi è stato assegnato? Qual è l'obiettivo?
  \item Come ci lavoro per ottenere un risultato \textbf{che abbia senso}? Se
        la ricerca e i dati indicano una direzione, è inutile andare in
        un'altra.
\end{enumerate}

\subsection{Quanto lavorare: un calcolo indicativo}

Se il corso è una simulazione di lavoro, vale il carico di un lavoro normale:
\begin{itemize}
  \item un lavoro a tempo pieno comprende da 30 a 45 ore settimanali; si
        considera una media di \textbf{40 ore};
  \item si sottraggono circa 9 ore di lezione: restano $\approx 31$, diciamo
        \textbf{30 ore};
  \item con circa 3 task a settimana: $30 / 3 = $ \textbf{circa 10 ore per
        task}.
\end{itemize}
Il numero può variare, ma rende l'idea: dedicare un'ora a un task e poi fare
altro non fa guadagnare nulla a nessuno. Lo studente ``butta soldi'', il
docente risparmia tempo leggendo il task in un minuto (o non leggendolo).
Dedicando invece attenzione, ricerca, cura nella formattazione e
sperimentazione, e facendolo con cognizione di causa, si impara.

\subsection{Niente voti, solo feedback}

Ad eccezione del corso di programmazione, non ci sono voti (né ``ottimo'', né
10, né 100). Ci sono \textbf{feedback}:
\begin{itemize}
  \item in game design, generalmente sulla documentazione;
  \item in level design, sui progetti pratici.
\end{itemize}
Il docente indica cosa va bene e cosa va modificato, concentrandosi in
particolare su ciò che va cambiato. Il feedback \textbf{non è mai un giudizio
sulla persona}, ma su ciò che è stato consegnato nel concreto.

\subsection{La condivisione dei task}

Tutti i task consegnati vengono \textbf{condivisi} con la classe, e possono
essere letti solo dopo aver ricevuto i feedback del docente. Si cresce così
non solo con i propri feedback, ma anche con quelli dati agli altri. Regola
importante: ciò che il docente dice a uno studente \textbf{vale per tutti}. Se
a una persona viene detto ``questa cosa non va mai più fatta'', da quel
momento vale per l'intera classe.

\subsection{Non copiare}

Copiare i task degli altri è una delle cose peggiori che si possano fare, non
per una questione morale, ma perché annulla il senso del percorso. Si può
dedicare a ogni task il tempo e l'attenzione che si ritengono opportuni, ed è
normale che col tempo ci si specializzi (preferendo, ad esempio, il level
design o la programmazione tecnica) e si investa meno dove si sente meno
necessità. In questi casi l'alternativa corretta è:
\begin{itemize}
  \item non consegnare;
  \item studiare ciò che hanno fatto gli altri;
  \item eventualmente riconsegnare in seguito.
\end{itemize}
Ci saranno momenti in cui si verrà sommersi dalle cose da fare: bisogna
imparare a scegliere.

\subsection{Il ``non leggo''}

Il docente ha piena facoltà di non dare alcun feedback a un task (il
cosiddetto \textbf{``non leggo''}) se:
\begin{itemize}
  \item non si è lavorato con attenzione;
  \item non si rispettano le indicazioni date;
  \item è stata violata una regola estremamente importante.
\end{itemize}
Un task ``non letto'' equivale a un task non consegnato e non entra nel
conteggio. Attenzione: un task \emph{perfetto} e un task che \emph{fallisce}
nel comprendere o applicare un framework sono \textbf{ugualmente validi},
purché ci si sia impegnati e si sia dedicato tempo, pazienza e ``sudore''.

\paragraph{Riconsegna.} Un task non letto, in ritardo o che si vuole
migliorare può essere riconsegnato e verrà conteggiato, a patto che:
\begin{itemize}
  \item i problemi segnalati siano stati risolti;
  \item sia stato fatto ulteriore lavoro, integrando anche tutto ciò che nel
        frattempo è stato detto agli altri colleghi nei loro feedback.
\end{itemize}
Il task va sviluppato con le competenze che si hanno nel momento della
riconsegna, non con quelle del momento della prima consegna.

\subsection{Requisiti per superare l'anno}

\begin{itemize}
  \item consegnare almeno l'\textbf{80\% dei task} assegnati (soglia
        malleabile);
  \item partecipare almeno all'\textbf{80\% delle lezioni}, in presenza o
        da remoto, preferibilmente in presenza; i ritardi vengono valutati
        caso per caso;
  \item comportarsi da persone corrette e rispettose (è stata firmata una
        carta dello studente): non urlare, non alzare le mani, non
        vandalizzare, non compiere atti illegali.
\end{itemize}
Impegnarsi di più porta inoltre a essere selezionati per varie opportunità.

\subsection{Scadenze}

\begin{itemize}
  \item la scadenza di un task è \textbf{sempre prima dell'inizio della
        lezione} successiva indicata;
  \item se la classe, come team, mantiene comportamenti scorretti, il docente
        può \textbf{anticipare le deadline} per tutti.
\end{itemize}

\subsection{Accessibilità e DSA}

Gli studenti con disturbi specifici dell'apprendimento o neurodivergenze
(dislessia, ADHD, autismo, ecc.) sono invitati a segnalarlo. Quando si
affronteranno le regole di formattazione e documentazione (errori
grammaticali, mancanze, incoerenze) se ne terrà conto e si potranno trovare
dei compromessi, pur ricordando che nel mondo del lavoro nessuno farà
``sconti''. La documentazione può essere inviata all'amministrazione o alla
segreteria, che la comunicherà a tutti i docenti, ma il consiglio è parlarne
anche direttamente con ogni singolo docente.

% ---------------------------------------------------------------------
\section{Le competenze trasversali}
\label{sec:soft-skill}

\subsection{I giochi si fanno in gruppo}

Si possono fare giochi da soli, ed è un'impresa ammirevole; ma lavorando da
soli \textbf{tutto ciò che si guadagna in libertà si perde in produttività}.

\paragraph{Il caso Stardew Valley.} \emph{Stardew Valley} è stato realizzato
da una sola persona, Eric Barone, in anni di sviluppo e di sacrifici: faceva
un altro lavoro per mantenersi (in un cinema) e di notte programmava, faceva
l'arte e componeva la musica. Il risultato è stato straordinario e lo ha reso
ricco, ma si tratta di un caso estremamente raro. Quanti giochi simili, nati
nelle stesse condizioni, sono falliti? Molti di più.

\paragraph{Il rischio d'impresa.} Lavorare da soli non è sbagliato di per sé
e ha aspetti molto positivi: si ha il privilegio di creare esattamente ciò che
si vuole. Ma ci si fa carico anche di tutto il \textbf{rischio d'impresa}: se
si lavora a un gioco per cinque anni e poi il gioco va male, resta il
vantaggio di averlo pubblicato, ma il dopo è sperare di vivere di quella
release o di finanziare con essa il gioco successivo. Le grandi aziende hanno
gli stessi problemi, ma li distribuiscono su più persone e più responsabilità.

L'obiettivo di un professionista è \textbf{saper lavorare in diverse
condizioni}, e la più comune è quella del dipendente (o del freelance /
libero professionista con partita IVA): si viene assunti per le proprie
capacità da un'azienda, piccola o grande che sia, e si lavora con altre
persone.

\subsection{Essere persone con cui si vuole lavorare}

Se le persone non vogliono lavorare con voi, quello è \textbf{il primo motivo
di licenziamento}. E se non si viene licenziati, il clima nell'azienda
diventa terribile: la produttività crolla, la motivazione crolla. Le
\emph{soft skill}, cioè l'insieme di abilità che rendono una persona
piacevole con cui lavorare, vanno coltivate e diventano sempre più importanti
col passare degli anni (relativamente meno ora, molto di più quando i team
saranno più grandi). Anche chi è un ``lupo solitario'' deve imparare a:
\begin{itemize}
  \item lavorare con altre persone;
  \item comunicare con altre persone;
  \item trovare un punto in comune da cui partire per fare le cose.
\end{itemize}

\subsection{Indipendenza e proattività}

Sono le altre due parole chiave del corso:
\begin{description}
  \item[Indipendenza] spesso si verrà lasciati da soli, non per cattiveria ma
        per sfida: quando si cerca una risposta da soli, ciò che conta non è
        tanto la risposta quanto \textbf{il processo}. Una risposta data dal
        docente non vale niente.
  \item[Proattività] chiedere, domandare, agire da soli senza aspettare
        direttive, perché spesso quelle direttive non ci sono, a maggior
        ragione nei team piccoli.
\end{description}
Tutto questo si concretizza nel parlare, nel presentare e nel fare
\emph{post-mortem} (l'analisi a posteriori di un progetto concluso, che verrà
approfondita più avanti), e in generale nell'attenzione alla crescita
personale.

\subsection{Fare domande}

Il corso insiste sulle domande: alzare la mano, intervenire anche ``a
bomba'', fermare il docente ogni volta che qualcosa non è chiaro. Se una
persona non conosce qualcosa, molto probabilmente c'è almeno un'altra persona
che non la conosce. Il docente può dare per scontata la conoscenza di termini
che non lo sono, ed è del tutto normale, soprattutto all'inizio.

% ---------------------------------------------------------------------
\section{Strumenti: engine, Unreal e materiali}
\label{sec:engine}

\subsection{Unreal Engine e gli altri motori}

Il motore di riferimento del corso, usato in level design, è \textbf{Unreal
Engine}: ``un gran mattone'', ma che dà a un designer molte possibilità. Il
vantaggio principale è che permette a un designer non pratico di
programmazione di \textbf{prototipare tramite i Blueprint} (il sistema di
scripting visuale di Unreal). Questo non impedisce di passare a Unity o ad
altri engine, e viceversa: il docente stesso ha studiato e lavorato a lungo
con Unreal, ma la sua ultima esperienza professionale (come level designer) è
stata su Unity, un engine che non conosceva, a dimostrazione che il salto da
un engine all'altro è possibile.

Panorama dei motori citati:
\begin{itemize}
  \item \textbf{Unreal Engine} e \textbf{Unity}: i due più diffusi, con
        modelli di sorgente ``più o meno'' aperto;
  \item \textbf{Godot}: engine open source;
  \item \textbf{CryEngine}: proprietario di Crytek;
  \item \textbf{Snowdrop}: engine di Ubisoft;
  \item Rockstar usa un proprio engine proprietario;
  \item alcune aziende usano versioni \textbf{personalizzate} di engine
        esistenti: ad esempio Milestone, nei giochi di corse, usa un Unreal
        customizzato;
  \item CD Projekt Red ha un engine proprietario (con i problemi di
        documentazione già citati).
\end{itemize}
In sintesi: ogni azienda può avere un engine proprietario oppure una versione
modificata di un engine commerciale, ma si parte sempre da una base.

\subsection{Programmazione: cosa ci si aspetta da un designer}

Non è richiesto alcun bagaglio pregresso di programmazione: il corso parte da
zero. Gli studenti sono stati selezionati (il corso è a numero chiuso) per
competenze analitiche e critiche sopra la media. Anche nel mercato del lavoro,
per una posizione di game designer generalmente \textbf{non ci si aspetta
competenza di scripting} (C++, Java, C\#, ecc.), a meno di voler fare il
\emph{technical designer}, e anche lì varia da azienda ad azienda.

\subsection{Versione e preparazione}

Indicazioni pratiche emerse in aula:
\begin{itemize}
  \item scaricare Unreal Engine; la versione consigliata per il corso è
        stata indicata nella \textbf{5.6} (con conferma successiva). La
        stabilità dipende da cosa si deve fare, da quali tool e plugin si
        usano: i problemi nascono spesso dai plugin. Nel passaggio da Unreal 4
        a Unreal 5, ad esempio, il sistema \emph{Landscape} era inizialmente
        inutilizzabile, mentre in Unreal 4 funzionava molto bene;
  \item è consigliato installare tramite \textbf{Epic Games Launcher}, più
        comodo: permette anche di riscattare giochi gratuiti periodici e di
        accedere ai pacchetti di \textbf{Fab} (il marketplace di asset
        acquisito da Epic), che ogni mese offre alcuni asset gratuiti utili
        per esercitarsi nel level design (ad esempio per fare \emph{speed
        level design});
  \item molti pacchetti di UI o di sistemi di interazione sono utili per
        prototipare al volo, ma spesso inutilmente complicati;
  \item un plugin di \textbf{materiali di prototipazione a griglia} (adattivi,
        con metriche, allineabili alla griglia, con materiali già pronti e
        colori modificabili, dal costo di circa 5 euro) è indicato come
        ``salvavita'' per il level designer; se necessario, i pacchetti
        indispensabili vengono forniti;
  \item portare \textbf{sempre il portatile} alle lezioni di level design
        (da novembre, a partire dall'interfaccia dell'engine); alle lezioni
        di game design basta che lo porti almeno metà della classe, per i
        lavori a coppie. Serve comunque qualcosa per prendere appunti
        (``tavolette d'argilla e cannucce, se volete fare i babilonesi'',
        oppure carta e penna).
\end{itemize}
Chi ha un portatile poco potente dovrebbe comunque provare a installare
l'engine e aprire un template per vedere come gira: all'inizio non si faranno
cose pesanti. Nelle prime lezioni di level design è importante mettere subito
le mani sull'interfaccia: ascoltare dove si trova uno strumento (ad esempio lo
\emph{snapping}) senza provarlo porta a dimenticarlo.

% ---------------------------------------------------------------------
\section{Letture e risorse di riferimento}
\label{sec:risorse}

\subsection{I libri}

\begin{description}
  \item[\emph{The Art of Game Design: A Book of Lenses} (Jesse Schell)] il
        testo \textbf{davvero importante}. Alcune lezioni finali dell'anno si
        baseranno sulle sue terminologie. È utile perché apre la mente sulla
        possibilità di applicare metodologie diverse a giochi diversi, ma
        soprattutto per la \textbf{terminologia} comune. Esiste anche in
        audiolibro (circa 25 ore).
  \item[\emph{Challenges for Game Designers} (Brenda Romero e Ian Schreiber)]
        libro \emph{entry level} con moltissimi esercizi, molto caro al
        direttore del corso.
  \item[\emph{21st Century Game Design} (Chris Bateman)] utile come
        retrospettiva: è un libro visionario che ha previsto metodi e
        soprattutto \emph{archetipizzazioni} dei giocatori tuttora valide nel
        design e nel marketing. Servirà per una delle lezioni più importanti
        dell'anno. È utile leggerlo per vedere quanto aveva previsto, ma è
        utile ``fino a un certo punto''.
  \item[\emph{An Architectural Approach to Level Design} (Christopher
        Totten)] non strettamente necessario, ma di lettura semplice. Ha due
        anime: parte dall'architettura (babilonese, greca, ecc.) e la trasla
        su esempi di level design (come \emph{Doom} e \emph{Team Fortress}).
        Un tempo ne esisteva un riassunto su Gamasutra (oggi
        \emph{Game Developer}).
  \item[\emph{A Theory of Fun for Game Design} (Raph Koster)] utile fino a un
        certo punto: contiene diverse idee superate, che lo stesso Koster ha
        affrontato in vari talk reperibili nel GDC Vault. Verrà ripreso nelle
        prime lezioni, poi messo da parte.
\end{description}

\subsection{Canali e conferenze}

\begin{description}
  \item[GDC (Game Developers Conference)] il suo canale raccoglie talk e
        panel estremamente importanti: è ``uno strumento misterioso che ci
        servirà più tardi''. \textbf{Attenzione}: per molti talk mancano ancora
        le competenze per interpretarli, e alcuni potrebbero fuorviare o
        mettere in testa idee strane. La GDC si concentra spesso su
        produttività e mobile.
  \item[Il \emph{Dragon Speech} di Chris Crawford] il talk del fondatore della
        GDC: appartiene a un passato lontano, ma è molto motivante. Da vedere
        subito.
  \item[Game Maker's Toolkit] canale particolarmente utile per nozioni
        pratiche, con video brevi e alcune interviste.
  \item[Nordic Game] altra conferenza importante, con talk e panel di esperti
        del settore.
\end{description}

\subsection{Associazioni e premi: come leggerli}

Le associazioni di settore (ad esempio IIDEA in Italia, o la IGDA) non sono
il male, anzi: sono un canale importante per affacciarsi all'estero o per
accedere a fondi. Va però ricordato che \textbf{tutelano i propri associati}:
non esiste solo il loro palinsesto, e bisogna valutare quanto spazio occupano
nel panorama della visibilità.

Lo stesso vale per i \textbf{premi} (Game Awards, ``Best Mobile Game'',
``Games for Change'', ecc.): al designer \emph{non interessano}. Che un gioco
abbia vinto ``più di 90 premi'' non dice nulla di utile. Servono a dare
visibilità e lustro ai giochi ed è giusto che esistano, perché ampliano il
panorama e il numero di giocatori. Ma ciò che interessa al designer sono i
\textbf{dati}:
\begin{itemize}
  \item come sono state implementate certe cose;
  \item come è stato fatto il marketing e dove il gioco è stato presentato;
  \item non quante wishlist ha fatto, ma \emph{come} le ha fatte;
  \item in che periodo storico è uscito e quali erano i competitor in quel
        momento.
\end{itemize}

% ---------------------------------------------------------------------
\section{Che cosa fa un game designer}
\label{sec:cosa-fa}

Per arrivare a una prima definizione, la classe ha condotto un brainstorming:
ognuno doveva indicare una cosa che, secondo lui, fa il game designer (senza
ripetere quelle già dette). Ne è emersa la lista seguente, commentata punto
per punto.

\subsection{Progettare meccaniche e struttura del gioco}

Il ruolo più generico è quello del \textbf{core mechanic designer}, detto
anche \emph{gameplay designer} e, un tempo, \emph{black box designer}: si
occupa delle \textbf{meccaniche di gioco}. Si va dal piccolo al grande (in
funzione anche della dimensione dell'azienda):
\begin{itemize}
  \item nel \textbf{piccolo} ci sono le singole meccaniche;
  \item nel \textbf{grande} ci sono i \textbf{sistemi}. Nei giochi
        \emph{sistemici}, con molti sistemi interconnessi (ad esempio
        \emph{GTA}), esistono specializzazioni dedicate a bilanciare e
        supervisionare i sistemi: è il \emph{system design}.
\end{itemize}

\paragraph{Quando un gioco è ``sistemico''?}
Non dipende solo da come le meccaniche interagiscono con il giocatore (le
\emph{dinamiche}, secondo i framework che passano dalle meccaniche alle
dinamiche, che verranno studiati più avanti), ma anche da \emph{quante}
meccaniche ci sono e da quanto interagiscono tra loro. Non diremmo che
\emph{Super Mario} sia un gioco sistemico, eppure contiene dei sistemi: la
linea è sottile. Quando i sistemi aumentano molto, sono complessi e si
influenzano a vicenda, il gioco diventa sistemico. Esempio: in \emph{The
Legend of Zelda: Breath of the Wild} il vento modifica quanto il personaggio
sale o scende; il fuoco può incendiare l'erba; l'erba prende fuoco se colpita
da una freccia infuocata; e così via. Nel mercato moderno i giochi sistemici
sono spesso, ma non sempre, legati a mappe \emph{open world}.

\subsection{Le 3C: camera, character, control}

La base del gioco, da cui generalmente si inizia a lavorare (anche nei
prototipi), è la ``sacra trinità videoludica'' delle \textbf{3C}:
\begin{description}
  \item[Camera] che tipo di inquadratura ha il gioco: prima persona, terza
        persona, isometrica, ecc.
  \item[Character] quanto è grande il personaggio, che tipo di personaggio è
        e, di conseguenza, quali azioni può compiere.
  \item[Control] i controlli: che cosa il personaggio può fare, quali sono i
        \emph{verbi} di gioco.
\end{description}
Le 3C sono definite dal designer, ed esiste anche una specializzazione per
questo lavoro. Esistono varianti: quella del direttore del corso è
\textbf{4C + J}, cioè \emph{camera, character, control, conflict} (conflitto)
e \emph{juice} (il ``succo'', il feedback che rende piacevole l'interazione).
Verranno approfondite più avanti.

\subsection{Il personaggio}

Il designer definisce il personaggio dal punto di vista delle meccaniche
(dimensioni, azioni, capacità), all'interno delle 3C. L'aspetto estetico è
invece competenza di altri dipartimenti.

\subsection{Documentare le informazioni per gli altri dipartimenti}

Assolutamente vero: la documentazione è una delle prime attività del
designer. Più l'azienda è grande, più questo lavoro diventa prominente; più è
piccola, meno è rilevante. In \textbf{pre-produzione}, però, è sempre
fondamentale.

\subsection{Mediatore e leader ``al servizio''}

Il designer è un \textbf{mediatore} nei conflitti interni al team, una sorta
di collante che si relaziona con dipartimenti diversi (non sempre gli stessi,
a seconda della fase di produzione). A volte fa anche da leader, a seconda di
come sono gestite la pipeline e la gerarchia.

Questa leadership è difficile da definire perché \textbf{non è verticale}: è
una leadership \emph{al servizio} (\emph{servant leadership}). Il designer
``è asservito al programmatore'': il suo compito non è arrivare con un'idea e
scaricarla addosso agli altri, né portare la soluzione e lasciare il problema
a qualcun altro, ma \textbf{seguire il processo a basso livello}, sporcandosi
le mani.

\paragraph{I feature team.} Nei team medio-piccoli o medio-grandi possono
esistere sotto-team (\emph{feature team}) dedicati a un obiettivo specifico:
ad esempio un designer incaricato del character o della camera, affiancato da
uno o due programmatori e da un artista. In questi casi il designer
supervisiona il lavoro e gestisce la produzione del sotto-team.

\subsection{Gestire l'avanzamento del prodotto}

Esiste una specializzazione che si occupa solo di questo, il
\textbf{producer}, a un livello gerarchico completamente diverso (e che
potrebbe essere una strada percorribile al termine del percorso di studi).
All'interno di un feature team anche il designer può farlo, ma non è detto.
Può certamente \emph{dare la direzione}: nella sua documentazione può
indicare che per arrivare a uno stato giocabile del prototipo serve un certo
lavoro di arte e di programmazione, per lo stato successivo un altro, e così
via. In ambito produttivo esistono metodologie strutturate come
\textbf{Agile}, che verranno viste in seguito.

\subsection{Presentare le idee}

Il pitching (sezione~\ref{sec:gerarchia-corso}) è richiesto al designer sia
in pre-produzione sia in produzione, in misura che dipende dall'azienda. Ad
esempio Ubisoft lascia ai designer, e in realtà a qualsiasi membro, la
possibilità di \emph{pitchare} un'idea prima di una nuova produzione, perché
possa iniziare a prendere vita.

\subsection{Bozze e prototipi}

Le bozze vanno bene, ma devono essere \textbf{chiare}. C'è un'eterna lotta tra
documentazione e gioco giocabile: se si fa solo documentazione il gioco non
avanza; se si fa solo il gioco manca la documentazione e si perde la visione.
Sono due cose \textbf{complementari}.

Il modo per comunicare un'idea in forma grezza ma giocabile è il
\textbf{prototipo}. Un prototipo può essere qualsiasi cosa: di carta, in
Unreal con Blueprint ``brutti''. \emph{Non deve essere bello, deve essere
giocabile}.

\paragraph{Esempio: il prototipo di carta di Mario + Rabbids.} Durante lo
sviluppo di \emph{Mario + Rabbids}, i programmatori di Ubisoft erano impegnati
su altri progetti (come \emph{Ghost Recon}). I designer realizzarono allora
un \textbf{gioco da tavolo} che riproduceva le meccaniche del gioco, quelle
di uno strategico a turni sul modello di \emph{XCOM}. E funzionava.
L'obiettivo era \textbf{validare un'idea}, validare un gameplay.

Altri modi per rendere chiara un'idea sono i \emph{mock-up} nella
documentazione e le immagini modificate o i \emph{mash-up} di qualsiasi
natura. Lato documentazione l'importante è che sia chiaro; lato prototipo,
che sia giocabile.

\subsection{Fare ricerca per l'idea, e poi smontarla}

Il designer ha una componente creativa importante e fa ricerca per arrivare a
un'idea. Ma ciò che deve poi fare necessariamente è \textbf{smontarla}. Il
suo lavoro è scientifico, esattamente come il \textbf{metodo scientifico}, il
cui obiettivo è \emph{confutare un'ipotesi}: si arriva a un'ipotesi tramite
sperimentazione e poi si cerca $n$ volte di confutarla; se resiste, la si
assume come valida \emph{in un certo contesto} (le regole della fisica
classica, ad esempio, non valgono nella meccanica quantistica).

Allo stesso modo il designer arriva a un'idea (tramite ricerca o chiedendo
agli altri: le idee le hanno tutti, non sono un privilegio del designer),
cerca di smontarla e poi cerca di \textbf{metterla a terra}. Se l'idea non è
applicabile, non è implementabile, finché non si trova un modo per
realizzarla va scartata. Può essere l'idea più bella dell'universo, ma
\textbf{se non può essere messa a terra, non esiste}.

\paragraph{L'originalità.} Serve un'idea originale? Dipende dall'obiettivo e
dalla visione del progetto:
\begin{itemize}
  \item se l'obiettivo è vendere molto, un'idea originale non serve
        necessariamente;
  \item se si vuole creare un nuovo filone, serve un'idea, se non originale,
        almeno interessante: può essere un discorso di marketing o una
        meccanica innovativa, con il rischio di non vendere.
\end{itemize}
Il ruolo del \emph{director} e il concetto di \emph{visione} verranno
approfonditi in seguito.

\subsection{Analizzare i competitor}

Arrivati a un'idea (``vogliamo un certo tipo di combattimento''), si va a
vedere se qualcuno l'ha già fatta e \emph{come}. Le meccaniche non sono
privilegio di nessuno, salvo quando vengono brevettate, come il
\textbf{Nemesis System} di \emph{La Terra di Mezzo: L'Ombra di Mordor}.

\paragraph{Esempio: il combattimento da ridimensionare.} In un progetto
(presso Stormind Games) il team si trovò con un sistema di combattimento
troppo costoso: mancavano le risorse e gli animatori per realizzarlo. Poco
output lato gameplay significava non poter fare qualcosa di complesso e
profondo come in \emph{Elden Ring}; al contrario, il team aveva molti
artisti e poteva puntare su qualcosa di più artistico. Si iniziò così a
orbitare intorno al corpo a corpo di \emph{The Last of Us} (parte I e II).
Sul canale della GDC esistono diversi talk che raccontano cosa è stato fatto
e come, con successi, insuccessi e costi (ad esempio del \emph{motion
capture}): da quei talk si ricavarono moltissime informazioni utili.

I competitor possono essere:
\begin{itemize}
  \item \textbf{diretti}: un gioco molto simile al proprio;
  \item \textbf{di mercato}: un gioco che domina il mercato in cui ci si vuole
        affacciare;
  \item \textbf{indiretti}: un gioco con un sistema simile al proprio, ma
        tutto il resto diverso.
\end{itemize}
Regola d'oro: \textbf{là fuori, nascosto nell'ombra, c'è sicuramente un gioco
molto simile al vostro}. Va solo trovato. E se per qualche ragione non ha
venduto, bisogna capire perché.

\subsection{Lavorare in engine e fare tuning}

Il level designer lavora in engine, così come il technical designer; ma
qualsiasi game designer, a seconda di ciò che deve fare, può finire in
engine, soprattutto per il \textbf{bilanciamento} e il \textbf{tuning}.

Il \emph{tuning} è l'azione di bilanciare con precisione i dati. Esempio: in
\emph{Super Mario} si modifica quanto in alto salta il personaggio, con che
velocità, in quanto tempo raggiunge l'apice, e così via. Questo lavoro rende
il gioco:
\begin{itemize}
  \item ad alto livello, dotato della \textbf{progressione} giusta;
  \item a basso livello, \textbf{piacevole da giocare}.
\end{itemize}
Il \emph{fine tuning} è un lavoro di dettaglio, paziente e minuzioso.

\subsection{Bilanciare}

Il bilanciamento, emerso tardi nel brainstorming, è una delle attività
centrali del designer e ha anche una specializzazione dedicata
(l'\emph{economy designer}, sezione~\ref{sec:tipi-designer}).

\subsection{Analizzare il mercato}

In pre-produzione, di solito, sì: capire le tendenze, i generi che cambiano
nel tempo, la compatibilità del gioco con il momento storico, in vista ad
esempio di un possibile publisher. Va capito quanto, come e con quanto tempo
farlo, perché si è sempre limitati da molte cose.

\subsection{Sintesi}

\begin{quote}
Il designer parte dall'idea ma non si ferma lì: la confuta, lavora sui
dettagli, bilancia, fa \emph{fine tuning}; lavora con il team e fa da
\emph{manager} per diversi dipartimenti; e naviga costantemente tra i limiti.
\end{quote}

% ---------------------------------------------------------------------
\section{I limiti del design}
\label{sec:limiti}

I limiti ci sono sempre, e nel mondo del design sono spesso numerosi. I più
importanti entro cui navigare sono:
\begin{description}
  \item[Tempo] ad alto livello, l'idea deve essere realizzabile nel tempo di
        produzione disponibile; a basso livello, il proprio lavoro
        settimanale deve essere fattibile nel tempo che si ha. Non ci si può
        dividere tra mille feature team se realisticamente se ne può seguire
        uno solo.
  \item[Tecnologia] si può o non si può fare qualcosa? Ci sono i server per un
        multiplayer con un certo numero di giocatori? Ci si può permettere,
        ad esempio, una fauna che reagisce alla fisica come quella di
        \emph{GTA VI}, che richiede il meglio dei \emph{technical artist}?
  \item[Background e competenze] che cosa sa fare il team? Realisticamente,
        chi ha lavorato dieci anni su giochi horror viene assunto per un
        horror, perché ha le competenze per risolvere problemi già visti.
  \item[Persone e motivazione] il team ha la motivazione e le competenze per
        fare un certo tipo di gioco, e per lavorare insieme?
  \item[Denaro] tutto ruota intorno ai soldi. Che si lavori da soli o in mille,
        le persone vanno pagate e il gioco deve avere un ritorno economico.
        Bisogna decidere dove tracciare la linea: raggiungere il
        \textbf{break-even} (guadagnare esattamente quanto speso) o puntare a
        grandi guadagni. Con un \textbf{publisher}, inoltre, una parte
        consistente dei ricavi va a lui (ripartizioni dell'ordine di 80/20 o
        70/30, che dipendono dal contratto firmato).
\end{description}

\subsection{Non siete ``gli uomini delle idee''}

Per qualcuno può essere un problema scoprirlo, ma il designer non è
``l'uomo (o la donna) delle idee'': accanto alla parte creativa c'è una parte
tecnica, e si è fortemente limitati da molti fattori. Questi limiti, però,
sono un aiuto:
\begin{itemize}
  \item riducono il ventaglio di possibilità potenzialmente distruttive per il
        team (cose che non si riescono a fare, che non piacciono o per cui
        mancano le risorse);
  \item aiutano a mirare bene, trovando soluzioni creative in uno spazio più
        piccolo e più attuabile.
\end{itemize}

\begin{quote}
\centering
\textbf{Le idee sono nulla, l'implementazione è tutto.}
\end{quote}

% ---------------------------------------------------------------------
\section{Ruoli e specializzazioni}
\label{sec:tipi-designer}

\subsection{La scala gerarchica}

\begin{description}
  \item[Creative Director (CD)] mantiene la \textbf{visione} e dà l'idea, in
        modo che tutto segua una certa direzione. In teoria può diventarlo
        chiunque, ma con un'esperienza notevole, nell'ordine dei decenni: si
        assume che in quel tempo abbia visto come lavorano molti dipartimenti.
        Naturalmente un CD con un background artistico ha meno competenze di
        design, e viceversa.
  \item[Director di disciplina] uno per dipartimento: \emph{game director},
        \emph{level director}, \emph{art director}, ecc.
  \item[Lead] il capo del dipartimento, che gestisce i gruppi di persone. Chi
        lavora in un dipartimento risponde al proprio lead (ad esempio un game
        designer risponde al \emph{lead designer}, non al game director).
  \item[Senior / Principal] i profili più esperti (a seconda dell'azienda).
  \item[Junior] il profilo di inizio carriera.
  \item[Intern] lo stagista.
\end{description}
Lo schema si replica uno a uno per ogni dipartimento: arte, suono, animazione
(a volte arte e animazione insieme, o \emph{rigging} e animazione insieme),
programmazione, UX/UI e così via. A fare la differenza, a parità di ruolo, è
il \textbf{livello di seniority}.

\subsection{UX e UI}

UX e UI stanno a metà tra game design e arte: a seconda dell'azienda il ruolo
si chiama \emph{UX/UI Artist} o \emph{UX/UI Designer}.
\begin{description}
  \item[UX (User Experience)] come una persona interagisce con le cose. Esempio:
        su uno smartphone, un pulsante accende e spegne lo schermo; per
        digitare la password si scorre verso l'alto e compaiono i punti da
        toccare.
  \item[UI (User Interface)] come sono ``vestite'' queste cose: ad esempio
        l'aspetto della barra dei punti vita (HP), realizzata con asset 2D o
        3D.
\end{description}

\subsection{Le specializzazioni del game design}

Più l'azienda è grande, più la specializzazione è puntuale.
\begin{description}
  \item[Gameplay / Core mechanics designer] (sono la stessa cosa) si occupa
        delle meccaniche e delle 3C.
  \item[System designer] si occupa dei sistemi e delle loro interazioni.
  \item[Economy designer] si occupa del bilanciamento; lavora soprattutto con
        fogli di calcolo (Excel, Google Sheets).
  \item[Combat designer] si occupa del combattimento; negli annunci di lavoro
        indica quasi sempre un combattimento \emph{action}, con gameplay
        reattivo (non strategico), anche se non è una regola.
  \item[Multiplayer designer] si occupa delle interazioni multigiocatore.
  \item[Technical designer] fa da ponte tra game design e programmazione.
\end{description}

\subsection{Le specializzazioni del level design}

\begin{description}
  \item[Encounter designer] posiziona i nemici nella mappa in modo che gli
        scontri siano bilanciati e la progressione equilibrata come
        desiderato.
  \item[World / map designer] popola il mondo di gioco.
  \item[Quest designer] popola il mondo con le missioni e pensa alla
        progressione al loro interno, cioè a dove il giocatore deve andare (si
        vedano ad esempio gli annunci di lavoro di CD Projekt Red o di
        Larian).
  \item[Technical level designer] fa da ponte tra level design e
        programmazione.
\end{description}

\subsection{Discipline vicine: suono, UI, produzione}

Il monte ore del corso non permette di uscire come \emph{narrative designer}
(o \emph{sound designer}) completi. Ciò non toglie che si possano coltivare
queste discipline:
\begin{itemize}
  \item nei progetti di gruppo si può occuparsi anche del suono, della UI/UX,
        o provare ad applicare metodologie come Agile per fare
        \emph{producing} e gestire il team;
  \item il tempo dedicato a queste attività va però sottratto al game design,
        e dipende da quante persone ci sono nel team, da quanto lavoro c'è e
        da quanto è grande il progetto;
  \item avere sviluppato e rilasciato un progetto, anche piccolo (anche una
        \emph{game jam}), documentando i propri processi, è già un ottimo
        argomento da presentare in futuro;
  \item la via migliore resta lavorare nell'industria e, da lì, chiedere di
        imparare anche altre discipline.
\end{itemize}
Un gioco si può fare su qualsiasi cosa: sugli scarafaggi, sui piedini delle
sedie, sui pennarelli. Si parte avvantaggiati da ciò che è vicino alla propria
sensibilità: chi ama la musica può pensare a un gioco con meccaniche legate
alla musica, con un gameplay innovativo ma contenuto e realizzabile. Il design
è un cappello che comprende molte discipline.

% ---------------------------------------------------------------------
\section{Approfondimenti emersi in lezione}
\label{sec:approfondimenti}

\subsection{Roguelike e roguelite}

Un esempio di come il marketing conia generi: Supergiant Games, con
\emph{Hades}, rese celebre il termine \textbf{roguelite} al posto di
\emph{roguelike} (pur non essendo il primo roguelite in assoluto). La
differenza:
\begin{itemize}
  \item il \textbf{roguelike} si fonda su una serie di punti fissati in una
        sorta di \emph{manifesto} che definisce che cosa rende un gioco
        roguelike; uno su tutti: a ogni morte si ricomincia la \emph{run} da
        capo;
  \item il \textbf{roguelite} rompe alcuni di questi punti, in particolare
        introducendo una \textbf{progressione esterna alla run}: il giocatore
        può potenziarsi fuori dalla singola partita per vincere più
        facilmente in quelle successive. \emph{Hades} ne è il caso principe.
\end{itemize}
Analogamente, un mobile game sviluppato dal direttore del corso fu
promosso come ``il primo \emph{donation game} al mondo'': un gioco che
convertiva il tempo di gioco in donazioni reali a enti benefici per animali.
Anche qui un neologismo creato lato marketing per definire un genere.

\subsection{Applicazioni fuori dall'industria videoludica}

Le competenze sviluppate non servono solo ai videogiochi:
\begin{itemize}
  \item i \textbf{render architettonici} oggi possono essere realizzati con
        Unreal, e diversi professionisti formati nel game design lavorano in
        questo campo;
  \item la \textbf{gamification}, cioè l'applicazione di regole videoludiche
        a contesti non ludici (mostre interattive, progetti culturali,
        consulenze come quella per applicare questi concetti alle macchine da
        allenamento di un'azienda di fitness), è in Italia uno dei modi più
        redditizi di spendere queste competenze, spesso con finanziamenti
        pubblici.
\end{itemize}

% ---------------------------------------------------------------------
\section{Il primo task}
\label{sec:primo-task}

\begin{enumerate}
  \item Creare un indirizzo e-mail \textbf{professionale} (consigliato: un indirizzo Gmail nella forma
        \texttt{nome.cognome@gmail.com}; l'anno di nascita è accettabile).
        Va bene anche un indirizzo professionale già esistente usato per
        lavoro. Da evitare indirizzi poco seri (``non \texttt{blackfire67}'').
  \item Una persona della classe raccoglie tutti gli indirizzi in un unico
        documento e lo invia via e-mail al docente, che aggiungerà tutti al
        Drive condiviso del corso (dove verranno caricate lezioni, slide e
        materiali, per avere i documenti tutti nello stesso posto).
  \item Installare Discord e iscriversi al server del corso, dove passano le
        comunicazioni ufficiali. È consuetudine, e utile, creare anche un
        gruppo tra studenti.
\end{enumerate}
\textbf{Scadenza}: prima della lezione del venerdì.

% ---------------------------------------------------------------------
\section*{Punti chiave del capitolo}
\addcontentsline{toc}{section}{Punti chiave del capitolo}

\begin{itemize}
  \item Il corso è una \textbf{simulazione di lavoro}: ci si tratta da
        colleghi, con una gerarchia da studio di sviluppo.
  \item Si impara con l'\textbf{inquiry-based learning}: lezione, task,
        dubbi, domande, condivisione, progetti.
  \item Il fallimento affrontato con la mentalità giusta è la principale fonte
        di crescita; i feedback riguardano il lavoro, mai la persona.
  \item Il lavoro del designer è per metà creativo e per metà
        \textbf{ricerca e metodo scientifico}; il game design non è una
        scienza esatta.
  \item \textbf{Documentazione} e \textbf{pitching} sono il cuore del
        mestiere.
  \item Il designer progetta meccaniche e sistemi (partendo dalle
        \textbf{3C}), documenta, media, prototipa, confuta le idee, studia i
        competitor, bilancia e fa tuning.
  \item La leadership del designer è \textbf{al servizio} del team, non
        verticale.
  \item I limiti (tempo, tecnologia, competenze, persone, denaro) restringono
        e indirizzano la creatività.
  \item \textbf{Le idee sono nulla, l'implementazione è tutto.}
\end{itemize}

%    % ...
%
%  (Con hyperref aggiungere anche \renewcommand{\theHchapter}{z.\arabic{chapter}}
%  nella parte zero, per evitare ancore duplicate.)
%  Tutte le etichette di questo capitolo hanno il prefisso "z1:" per non
%  collidere con quelle dei capitoli del corso.
% =====================================================================

\chapter[La cultura del feedback]{Lavorare in squadra: la cultura del feedback}
\label{cap:z1-feedback}

Questo capitolo apre le \emph{lezioni zero}: incontri che precedono l'inizio
vero e proprio del corso di Game Design e che non riguardano una tecnica di
design in senso stretto, ma \textbf{la figura professionale} di chi lavora
nell'industria dei videogiochi. Prima di imparare a progettare un gioco
conviene infatti capire in che tipo di ambiente lo si progetterà: un ambiente
fatto di team, di gerarchie, di scadenze e, soprattutto, di persone che devono
riuscire a lavorare insieme.

La lezione è la prima giornata di un \textbf{workshop di due giorni} (circa
dieci ore in totale, pause pranzo comprese) tenuto da Stefano
Cozzi\footnote{Il cognome è riportato come compare nella trascrizione
automatica, che in più punti storpia i nomi propri; andrebbe verificato.
Nella stessa trascrizione lo studio per cui lavora appare come
``Stormine'', ``Stormline'' e ``Stormind'': si tratta di Stormind Games,
studio con sede a Catania, coerentemente con quanto detto in lezione.}, un
professionista con circa quindici anni di esperienza nel settore, tra
Blizzard Entertainment, Riot Games e, oggi, Stormind Games. Il workshop è
nato da un confronto con i coordinatori dell'accademia\footnote{Nella
trascrizione l'accademia è resa in modi diversi (``Hive'', ``a dive'',
``D.I.V.''); i coordinatori citati sono Enrico e Andrea, insieme al
direttore Diego Ricchiuti. Nel resto del capitolo si usa semplicemente
``l'accademia'', come negli altri capitoli.}: serviva un momento dedicato a
\emph{come ci si comporta in un ambiente professionale}, cosa che riguarda sia
il lavoro futuro sia i progetti di gruppo svolti durante il percorso di
studi, dove ci si deve già abituare a comportarsi ``come se si fosse in
un'azienda''.

La prima giornata ha un unico grande tema, sezionato in più parti:
\textbf{la cultura del feedback}. La seconda giornata (che sarà oggetto della
prossima lezione zero) sarà più pratica: \emph{time management},
\emph{brainstorming} e il ``guidare da dentro'' un team anche senza esserne
il leader (sezione~\ref{sec:z1-domani}).

Il percorso del capitolo segue quello della lezione:
\begin{enumerate}
  \item il formato del workshop e il percorso professionale del relatore
        (sezione~\ref{sec:z1-relatore});
  \item perché le aziende tengono tanto al lavoro di squadra e al
        \emph{fit} culturale (sezione~\ref{sec:z1-squadra});
  \item che cos'è il feedback e da che cosa dipende il modo di darlo
        (sezione~\ref{sec:z1-cose-feedback});
  \item le cinque fasi di vita di un team e il valore del
        \emph{post-mortem} (sezioni~\ref{sec:z1-fasi-team}
        e~\ref{sec:z1-qa-fasi});
  \item il confronto con i colleghi e alcune domande sul lavoro nei grandi
        studi (sezioni~\ref{sec:z1-rivalita} e~\ref{sec:z1-industria});
  \item il circolo virtuoso tra feedback, miglioramento e fiducia, e i suoi
        effetti sulla \emph{retention} e sulla cultura aziendale
        (sezioni~\ref{sec:z1-triangolo} e~\ref{sec:z1-cultura});
  \item il ``gioco del traditore'' e il principio del
        \textbf{default to trust} (sezioni~\ref{sec:z1-among-us}
        e~\ref{sec:z1-default-trust});
  \item le traiettorie di carriera e la valutazione delle performance
        (sezione~\ref{sec:z1-traiettorie});
  \item come si dà, come si riceve e quando si dà un feedback
        (sezioni~\ref{sec:z1-dare}, \ref{sec:z1-ricevere}
        e~\ref{sec:z1-quando});
  \item il feedback come regalo, le domande finali e le anticipazioni
        della seconda giornata (sezioni~\ref{sec:z1-regalo},
        \ref{sec:z1-qa-finale} e~\ref{sec:z1-domani}).
\end{enumerate}

Una premessa di metodo: il relatore ha alternato continuamente teoria,
domande alla classe e racconti della propria esperienza. I racconti non sono
digressioni: sono gli esempi concreti con cui i principi vengono resi
tangibili, e per questo sono riportati con cura.

% ---------------------------------------------------------------------
\section{Il workshop e il relatore}
\label{sec:z1-relatore}

\subsection{Formato e regole di partecipazione}

Il relatore ha cercato di ``riassumere quindici anni di esperienza in dieci
ore''. Ha materiale per circa il doppio delle ore disponibili: le cose
strettamente necessarie vengono trattate per prime, il resto è tenuto ``in
riserva'' come \emph{nice to have}. Le regole del workshop:
\begin{itemize}
  \item le domande si fanno \textbf{alla fine di ogni slide}, così da chiudere
        prima il discorso; non c'è un'unica sessione finale di domande e
        risposte (anche se ci sarà), perché il relatore preferisce fermarsi
        spesso e non fare ``uno spiegone di dieci ore'';
  \item le domande non devono per forza riguardare la slide appena vista: se
        una slide dà lo spunto per altro, se ne può parlare; al massimo sarà
        il relatore a rispondere in breve o a rimandare a una slide
        successiva;
  \item si può chiedere anche di esperienze passate, opinioni e sensazioni;
        se per NDA o per serietà professionale non può rispondere, lo dirà
        esplicitamente (``non lasciate che sia un \emph{blocker}'');
  \item su Zoom si usa l'alzata di mano o la chat, che il relatore tiene
        aperta sul secondo monitor;
  \item non ci sono pause ufficiali: la lezione è registrata e ognuno è
        libero di assentarsi quando serve, cercando però di non disturbare il
        flusso della lezione;
  \item le slide verranno fornite integralmente tramite la segreteria
        dell'accademia (Vanessa), insieme alle registrazioni: le slide senza
        il parlato servono poco, ma sono un buon promemoria.
\end{itemize}

\subsection{Il percorso professionale}

Il relatore si presenta (scherzosamente) come la ``guida escursionistica''
che accompagna la classe nella ``foresta incantata'' della cultura del
feedback, perché per molti è un territorio sconosciuto. Il suo percorso:
\begin{itemize}
  \item in Italia ha lavorato per aziende che fornivano servizi a terzi, tra
        cui il lavoro su \emph{Warhammer Online}, e per un canale televisivo
        sul digitale terrestre dedicato ai videogiochi\footnote{Nella
        trascrizione: ``Digital Game Network Channel''. Il nome esatto del
        canale andrebbe verificato.};
  \item a \textbf{Blizzard Entertainment} (circa sei anni, a Parigi) ha fatto
        il \emph{community manager} per l'Italia (\emph{StarCraft~II},
        \emph{World of Warcraft} da \emph{Mists of Pandaria} in poi,
        \emph{Hearthstone}, \emph{Diablo~III}, le espansioni di
        \emph{StarCraft}, il lancio di Battle.net), per poi diventare il
        manager dei community manager europei;
  \item a \textbf{Riot Games} (circa cinque anni, a Dublino) è stato
        brevemente community manager italiano e poi \emph{publishing
        manager} di un team di circa trenta persone per l'Europa. Il
        \emph{publishing} non è solo community: è l'unione di
        \textbf{community, marketing e PR};
  \item a \textbf{PG Esports} si è occupato di eventi dal vivo e online
        (ad esempio la Milano Games Week e tornei di \emph{Valorant} e
        \emph{League of Legends}), restando in contatto con Riot e Blizzard;
  \item da circa un anno è tornato nello sviluppo, a \textbf{Stormind
        Games}, dove oggi si occupa di risorse umane e organizzazione del
        lavoro; lo studio verrà presentato più ampiamente al \emph{Career
        Day} dell'accademia insieme alla sua HR manager;
  \item fa inoltre consulenze aziendali anche al di fuori dei videogiochi e
        insegna in master universitari (IULM di Milano).
\end{itemize}

Il suo ruolo ricorrente, sia a Blizzard sia a Riot, è stato quello di
\emph{Mr.~Wolf} (``per chi ha visto \emph{Pulp Fiction}''): quello che
risolve i problemi. Veniva inserito nei team più in difficoltà, ad esempio
dal punto di vista delle performance, per sistemarli e passare al successivo.
Era presente in alcuni dei momenti più critici, come il lancio di
\emph{Diablo~III} o il lancio di espansioni di \emph{World of
Warcraft}\footnote{La trascrizione riporta ``il lancio con i Gary's'' e
``World of Trenor'': si tratta verosimilmente di \emph{Mists of Pandaria} e
di \emph{Warlords of Draenor}.}, a fare \emph{damage control}.

La scoperta fondamentale di questi anni è che i problemi da risolvere
\textbf{quasi mai riguardavano il prodotto} (non era certo lui a sistemare i
server di \emph{Diablo~III}), ma i team che c'erano dietro: persone che
litigavano, non collaboravano, non riuscivano a stare nella stessa stanza
senza urlarsi contro. E nel 90\% dei casi tutto questo si riduce a
\textbf{problemi di comunicazione}: è raro che perfino le incompatibilità
caratteriali non siano risolvibili parlando.

\subsection{Struttura non vuol dire burocrazia}

A Stormind il relatore sta accompagnando un team passato da 30 a 95 persone
in meno di un anno, attraverso tutti i cambiamenti di cui parla il workshop.
Darsi una struttura viene spesso percepito male, come ``più processi, più
procedure'', ma bisogna distinguere:
\begin{itemize}
  \item un processo come ``timbrare il cartellino'' può essere davvero
        invasivo e noioso;
  \item definire \emph{come funziona un documento di game design, chi lo
        approva, chi lo visiona e chi no} è invece una struttura necessaria.
\end{itemize}
Uno studio di 30 persone nella stessa sede può farne a meno: quando c'è un
dubbio ci si alza e si va alla scrivania del collega. Uno studio di 95
persone sparse per l'Italia (e in parte per l'Europa) non può più farlo:
servono una struttura diversa, un modo di lavorare diverso e una
pianificazione del futuro diversa.

% ---------------------------------------------------------------------
\section{Perché il lavoro è sempre lavoro di squadra}
\label{sec:z1-squadra}

\subsection{Non si va in ufficio per fare amicizia}

Il primo messaggio è netto: \textbf{non si va in ufficio per fare amici} o
per trovare compagni di vita. Si va in ufficio perché c'è un lavoro da fare, e
lo scopo è riuscire a farlo \emph{anche con chi sta antipatico}. Con le
persone simpatiche è facile: diventeranno amici, e non sarà mai un problema.
È con chi sta antipatico che le cose si complicano, ed è lì che emergono i
\textbf{veri professionisti}: quelli che riescono a lavorare con chiunque
senza che l'antipatia pesi sul modo in cui lavorano e collaborano, mettendo il
prodotto e il team davanti al proprio ego e alle proprie preferenze personali.

La frase che fa da filo conduttore a tutto il workshop (e che il relatore
promette di riproporre ``fino alla nausea'') è:

\begin{quote}
  \emph{La capacità di un gruppo di persone di fare cose straordinarie
  dipende dal modo in cui queste persone operano insieme.}
\end{quote}

Le aziende raramente pensano che l'innesto di una singola persona possa
cambiare il destino di un progetto. Questo è tanto più vero quanto più grande
è il team, ma vale anche per i piccoli studi, anche per un gruppo di cinque
persone. Come negli sport di squadra: l'individuo conta fino a un certo punto,
e anche se si ha Dončić in squadra, se si gioca uno contro cinque si perde
comunque\footnote{Il riferimento è alle finali NBA in corso nei giorni del
workshop tra Dallas Mavericks e Boston Celtics (giugno 2024), il che
permette di datare approssimativamente la lezione.}.

\subsection{Lo \emph{useless genius} e il costo delle persone tossiche}

Molti dipartimenti HR usano la definizione di \textbf{useless
genius}\footnote{Il concetto è più noto in letteratura manageriale come
\emph{brilliant jerk}; in lezione è usata l'espressione \emph{useless
genius}.}: si può essere bravissimi, a volte i migliori nel proprio lavoro,
ma se non ci si sa comportare, relazionare e operare in un gruppo si è
fondamentalmente inutili ai fini del gruppo. Nel migliore dei casi inutili,
nel peggiore \emph{dannosi}.

In passato le aziende tendevano a ``chiudere un occhio'': la persona è
fondamentale, la teniamo. Oggi la domanda è diventata: \emph{quanto costa
all'azienda, oltre al suo stipendio, tenere una persona che crea un ambiente
tossico e peggiora chi le sta intorno?} Ci sono costi sommersi che le aziende
cominciano a vedere, e la tendenza è liberarsi di queste persone il prima
possibile. Il ragionamento è semplice:
\begin{itemize}
  \item le \textbf{competenze si possono trasmettere} o sostituire: si trova
        qualcuno che fa lo stesso lavoro magari nel doppio del tempo, oppure
        più persone che se lo dividono;
  \item il \textbf{carattere delle persone} invece cambia molto raramente.
\end{itemize}

\subsection{Il fit culturale}

Quando qualcuno \textbf{non è un \emph{fit} culturale} per un'azienda è
meglio separarsi presto. Ma non essere adatti a \emph{quell'}azienda non
significa non essere adatti a lavorare: è perfettamente legittimo che una
persona inadatta a un'azienda sia adattissima a un'altra.

Il relatore l'ha vissuto in prima persona. Dopo sei anni a Blizzard viene
contattato da Riot e affronta un processo di selezione lunghissimo (``quasi
undici colloqui''), al termine del quale viene scartato. Scrive allora una
mail che gli ha ``cambiato la traiettoria di carriera'': chiede
semplicemente perché non è stato ritenuto adatto. La risposta: in nessun
colloquio si era esposto abbastanza da far capire \emph{perché} volesse
lasciare Blizzard, e quindi pensavano che in realtà stesse bene dov'era. Il
motivo era che, per educazione professionale, non gli piace parlare male del
proprio posto di lavoro; portata all'eccesso, però, questa riservatezza era
diventata un ``è tutto bellissimo, ma vorrei provare Riot'', senza un perché.
Lo invitano comunque a un incontro dal vivo (``dal vivo è diverso''), e alla
fine viene assunto.

Una volta arrivato, i sei mesi di prova (\emph{probation}) sono serviti tutti
per smettere di ragionare e lavorare ``come a Blizzard''. La differenza
culturale era enorme, tanto che in Riot esisteva un gruppo di supporto
interno tra ex dipendenti Blizzard: un collega gliel'ha descritto come
``uscire da una relazione tossica e doversi riabituare a una relazione
normale''\footnote{Il relatore precisa che parla della Blizzard di allora,
che non esiste più: oggi è parte di Microsoft e nessuno dei suoi ex colleghi
vi lavora ancora, anche a causa dei licenziamenti di massa.}.
\begin{description}
  \item[Blizzard] era \emph{developer first}: il focus era chiaramente sugli
        sviluppatori, prima di tutti gli altri. Dopo sei anni c'erano ancora
        meeting a cui non poteva partecipare perché non abbastanza
        \emph{senior}, informazioni che non poteva sapere, piani che non
        poteva vedere.
  \item[Riot] era \emph{player first}: prima i giocatori, e i dipendenti come
        mezzo per trasferire il feedback dei giocatori dentro il gioco. Il
        secondo giorno lo portano in una stanza e gli mostrano il piano dei
        cinque anni successivi. Lui, sorpreso, pensa a un errore (``sono qui
        da oggi, non vorrei mettervi nei guai''); la risposta è: ``sei un
        dipendente Riot o no? Allora non ci sono livelli di accesso: devi
        sapere queste cose per fare meglio il tuo lavoro''.
\end{description}

Da questa esperienza ha imparato molto sulla \textbf{trasparenza}: condividere
le informazioni aiuta il lavoro. Non sempre, però: a volte non si condivide
qualcosa non per segretezza, ma perché non è ancora pronta per essere
presentata, esattamente come non si mostra un gioco quando è ancora solo
teoria (``te lo faccio vedere quando ho qualcosa di tangibile'').

La lezione per gli studenti è doppia: le aziende valuteranno il \emph{fit}
culturale sempre di più, ma capire se si è un \emph{fit} richiede anche
\textbf{lavoro da parte del candidato}. Bisogna avere la presenza di spirito
di dire: ``conosco i valori di questa azienda, non mando nemmeno il CV,
perché so che durerei un mese e mezzo''.

% ---------------------------------------------------------------------
\section{Che cos'è il feedback}
\label{sec:z1-cose-feedback}

\subsection{Le prime risposte della classe}

La prima domanda alla classe è: \emph{quando sentite la parola
``feedback'', a cosa pensate?} Le risposte raccolte, a voce e in chat:
\begin{itemize}
  \item un \textbf{riscontro} su ciò che si fa e su come lo si fa, sul lavoro
        o sul modo in cui ci si interfaccia con il gruppo; può essere
        lavorativo oppure comunicativo e comportamentale. Chi ha risposto lo
        vede come qualcosa che, se l'azienda non dà, va attivamente
        richiesto, perché è importante sapere come sta andando per
        migliorarsi;
  \item né positivo né negativo, ma \textbf{costruttivo}: in entrambi i casi
        serve a migliorarsi;
  \item uno \emph{scambio di opinioni}, un \emph{riscontro produttivo};
  \item un feedback ``di reazione'', non verbale: si fa una cosa e si
        osserva come cambia l'ambiente o il comportamento delle persone
        intorno (anche perché a voce la gente spesso si tiene le cose
        negative);
  \item per un designer, anche il \textbf{linguaggio del corpo}: osservando
        un giocatore giocare si può dedurre l'emozione che prova (frustrazione,
        incomprensione).
\end{itemize}

Sulla forma (scritto o orale?) le opinioni divergono: c'è chi lo preferisce
\emph{mille volte} scritto; chi pensa che dipenda dalle dimensioni
dell'azienda (orale nelle piccole, scritto in quelle con team di centinaia di
persone); chi pensa che dipenda dall'importanza del feedback (scritto se
tecnico o pieno di dati, da tenere sotto gli occhi; orale se è un consiglio
veloce).

Il relatore usa due esempi per mostrare che \emph{dipende}:
\begin{itemize}
  \item un feedback sulla programmazione ha senso darlo scritto, magari
        commentando direttamente il codice, più che ``andando a prendere un
        caffè'';
  \item ma se il feedback è ``ieri, durante il meeting, ho percepito male il
        fatto che mi interrompessi di continuo'', ce lo si aspetta
        \textbf{orale}.
\end{itemize}

\subsection{Da che cosa dipende il modo di dare feedback}

Lo \emph{spoiler} della giornata è che \textbf{non esiste una risposta univoca}
su come dare feedback. Il workshop non insegna ``il'' modo, ma fornisce
\emph{tutte le opzioni}, perché il modo giusto cambia in base a molti fattori:
\begin{itemize}
  \item a chi lo si dà: un collega, un superiore, il proprio capo, un proprio
        \emph{report} (un sottoposto);
  \item se è personale o professionale;
  \item se è la prima volta o la terza volta che si dà lo stesso feedback;
  \item se riguarda un progetto importante o secondario.
\end{itemize}

Manca però l'elemento fondamentale, che la classe individua subito:
\textbf{l'altro}. A parità di tutti gli altri fattori, il \emph{carattere
della persona che riceve il feedback} dovrebbe far cambiare il modo in cui lo
si dà. Per questo, prima della lezione, agli studenti era stato chiesto di
fare un test di personalità: non un test di psicologia avanzata, ma uno
strumento per rendersi conto di quanta varietà ci sia anche solo in una
classe di 53 persone (il che non vuol dire trovare 52 modi diversi di dare
feedback, ma lavorarci un po' sì).

Il relatore porta se stesso come esempio. Preferisce:
\begin{itemize}
  \item i feedback \textbf{positivi scritti}, perché è a disagio a ricevere
        complimenti in pubblico;
  \item i feedback sul proprio lavoro \textbf{diretti}: ``vai subito al
        punto e non girarmi intorno''.
\end{itemize}
È molto orgoglioso del suo lavoro e, quando qualcuno lo critica, la sua
primissima reazione istintiva è ``come ti permetti''. Negli anni ha imparato
a farla durare un secondo, ma chi gli dà feedback di persona dice di
vederla ancora negli occhi. Per mitigarla, dice esplicitamente a chi gli dà
feedback: \emph{``ignora il modo in cui reagisco nei primi cinque secondi,
vai avanti, poi mi passa''}. È considerato da tutti molto aperto al feedback
e non ha problemi a dire ``hai ragione, facciamo diversamente'': ma ci sono
reazioni istintive e caratteriali che è difficile controllare al cento per
cento.

\subsection{Con chi si lavora meglio}

Guardando le persone con cui si lavora abitualmente (anche nei progetti
dell'accademia), si noterà che si lavora meglio con due profili:
\begin{enumerate}
  \item chi è \textbf{esattamente come noi}, ma qui è un 50 e 50: alcuni
        caratteri vanno d'accordo con i propri simili, altri non vogliono
        nessuno con lo stesso carattere nel gruppo (e si potrebbe rientrare
        in questa categoria);
  \item chi è \textbf{affine}: ha qualcosa in comune con noi, ma non tutto,
        e proprio per questo può dare punti di vista differenti.
\end{enumerate}
Il lavoro della giornata è capire come valorizzare il secondo caso: ottenere
feedback da persone con caratteri, visioni e background diversi, per
completare una visione del proprio lavoro che è necessariamente limitata. È
come esplorare una mappa con la nebbia di guerra: un'altra persona, al di là
del muro, vede cose che noi non vediamo.

\subsection{Il feedback come somma di piccoli miglioramenti}

Per il relatore il feedback è questo: \textbf{applicare tanti piccoli
miglioramenti periodici}, che si sommano nel tempo. Migliorare dello 0,1\%
ogni giorno porta a fine anno a un miglioramento enorme, e questo vale per i
prodotti, per i giochi, per le persone e per i team.

\paragraph{L'esempio del ciclismo britannico.} La squadra di ciclismo
britannica, fino a una quindicina d'anni fa, non aveva vinto quasi nulla.
Cambiato il team manager, tutti si aspettavano grandi cambiamenti sportivi
(convocazioni, strategie di gara). Arrivarono invece \emph{minuscoli
cambiamenti non sportivi}: un altro chef in trasferta, materassi migliori nei
camper per riposare meglio, un sellino più comodo a fine tappa. Cambiamenti
marginali che da soli non avrebbero cambiato niente; eppure, dopo due o tre
anni di questo lavoro quotidiano, la squadra ha cominciato a vincere a
ripetizione, Olimpiadi comprese, cambiando il modo in cui il ciclismo è
visto in Gran Bretagna\footnote{È la nota strategia delle \emph{marginal
gains} associata a Dave Brailsford (British Cycling e Team Sky); il nome non
è citato in lezione.}. Certo, si è investito anche sulle persone giuste, ma
tutto è partito dai piccoli miglioramenti periodici.

Per gli studenti vale lo stesso: chiarire un conflitto con una persona con
cui si ha ancora poco a che fare, o una situazione che ha lasciato
``l'amaro in bocca'', sembra inutile oggi, ma darà i suoi frutti tra tre, sei,
nove mesi. Sono \emph{investimenti sul futuro}.

\paragraph{Peggiorare ogni giorno.} Esiste anche il percorso inverso:
lasciar correre, ``buttare sotto il tappeto'' le cose perché ``tanto è
inutile parlarne''. Queste cose tornano sempre fuori dopo otto mesi, come il
\emph{libro dei rancori} dei nani di Warhammer: ``ti ricordi quando, un anno
e mezzo fa, all'inizio del progetto, mi hai detto quella cosa? Mi sono
offeso''. A quel punto è quasi impossibile risolvere, perché nessuno ricorda
più con precisione cosa sia successo (``ma quale cena?'', ``forse era in
pizzeria'', ``no, eravamo da un'altra parte''). \textbf{Tenersi le cose nel
cassetto per tirarle fuori quando servono è sempre sconsigliato}: non servono
ed è impossibile usarle in modo costruttivo.

% ---------------------------------------------------------------------
\section{Le cinque fasi di vita di un team}
\label{sec:z1-fasi-team}

\subsection{Perché conoscerle}

Quasi tutti gli studenti lavorano già in gruppo (al primo anno su un gioco da
tavolo, negli anni successivi su un videogioco) e si saranno accorti che un
team attraversa delle fasi: il modo di lavorare insieme cambia nel tempo.
Spesso ci si fa prendere di sorpresa: ``eravamo così affiatati, cosa è
successo?''.

Un esempio da una consulenza fuori dai videogiochi: un'azienda farmaceutica
gli affida un team commerciale (rappresentanti che presentano medicinali a
ospedali, medici e privati). Il problema: ``prima andava tutto bene; ora che
siamo in diciotto dovremmo lavorare meno e meglio, e invece va peggio''. Va
peggio perché le persone nuove sono appena arrivate, mentre il management si
aspettava già di vedere i risultati dell'investimento.

Fino a pochi anni fa le fasi di un team venivano disegnate come una
\textbf{linea dritta verso l'alto}: le fasi erano le stesse, ma viste come un
miglioramento costante. Non si vedevano le parti \emph{fisiologiche} del
lavoro in team, cioè anche i peggioramenti. Accettarle è il primo passo per
accettare i fallimenti: nessun percorso personale o di team è lineare.

Il relatore ammette che il suo modo di lavorare è nato per il 70\% dal
proprio impegno a imparare, fare bene e rispettare le persone, e per il 30\%
dall'aver lavorato con ``incapaci colossali'' che gli hanno insegnato come
\emph{non} fare le cose. A Blizzard un manager gli fece capire che, pur non
sapendo ancora che tipo di \emph{people manager} sarebbe diventato, sapeva
benissimo quale tipo di manager \emph{non} voleva essere.

\subsection{Forming, storming, norming, performing, adjourning}

Le fasi sono cinque, un po' come le fasi del lutto. Molte aziende ne
considerano solo quattro (``e sono quelle che solitamente chiudono''); le
aziende al passo con i tempi ne considerano cinque\footnote{Il modello è
quello di Bruce Tuckman (1965), a cui la quinta fase, \emph{adjourning}, fu
aggiunta nel 1977. Il nome di Tuckman non è citato in lezione; la
trascrizione riporta la quinta fase come ``journing''.}.

\begin{description}
  \item[Forming] Si viene messi insieme e si comincia a lavorare. Non si sa
        nulla degli altri; spesso non si conoscono nemmeno i nomi, ed è
        normale.
  \item[Storming] Si cominciano a ``dare le testate come le capre in
        montagna per determinare il territorio'': questa cosa la faccio io,
        questa la vorrei fare così, chi decide? Si stabilisce chi fa cosa e
        come si divide il lavoro. Nel grafico è una \textbf{fase di
        peggioramento}: che il team litighi o discuta civilmente non cambia
        nulla, il team deve investire parte di tempo ed energie per risolvere
        queste questioni. È normale e \textbf{non è evitabile}.
  \item[Norming] Si comincia a capire come lavorano gli altri, cosa sanno
        fare e come far funzionare il gruppo; iniziano i miglioramenti. È una
        forma di \emph{accettazione} degli altri, che deriva dall'aver capito
        cosa devo fare io e cosa devono fare loro.
  \item[Performing] C'è uno sforzo \textbf{proattivo}: ``questa cosa
        toccherebbe a me, ma so che Alessio la fa molto meglio, quindi gli
        chiedo aiuto''. Si sfruttano i punti di forza degli altri, perché li
        si conosce e perché si sa che passare il testimone non avrà effetti
        negativi (il capo non dirà ``perché non l'hai fatta tu?'').
  \item[Adjourning] Il progetto finisce, si chiude o viene riorganizzato, il
        team viene sciolto. Quello che andrebbe fatto è analizzare
        \textbf{cosa ha funzionato, cosa no e cosa migliorare} per il
        progetto successivo.
\end{description}

\subsection{La fase che tutti saltano}

``Spoiler incredibile'': l'\emph{adjourning} è la fase che molte aziende
saltano, \textbf{perché c'è sempre il progetto dopo}.

A PG Esports, per la Milano Games Week (come per qualsiasi evento dal vivo a
cadenza annuale: Lucca Comics, Etna Comics, Comicon\dots), finita l'edizione
di un anno si cominciava dal giorno dopo a lavorare a quella successiva,
perché un anno sembra tanto ma non basta mai. Nei videogiochi è lo stesso:
finito un gioco, a parte le patch di mantenimento e del \emph{day one}, si
passa subito al successivo.

Il cambiamento introdotto dal relatore (in parte già praticato, ma da lui
strutturato) è stato: finita la Games Week, \textbf{tutti in ferie per una
settimana}; poi ci si ritrova per definire cosa ha funzionato e cosa no, con
un \textbf{post-mortem} (o \emph{debrief}); solo dopo si inizia a lavorare
all'edizione successiva. Senza questa fase manca la parte di
\emph{learning}, cioè di ciò che si è imparato dal progetto.

Il consiglio vale per gli studenti: quando si finisce il gioco da tavolo, e
poi il gioco digitale, bisogna trovare il tempo per il post-mortem,
altrimenti si butta via \textbf{il 90\% dell'utilità del progetto}. Una
struttura semplice per farlo è descritta nella sezione~\ref{sec:z1-qa-finale}.

% ---------------------------------------------------------------------
\section{Domande sulle fasi del team}
\label{sec:z1-qa-fasi}

\subsection{Lo storming che ritorna}

\paragraph{Capita che lo storming si ritagli spazio anche nelle altre fasi?}
Lo \emph{storming} è una fase precisa, che però molti manager, coordinatori e
team cercano di \textbf{saltare}: ``se io sono bravo come manager e loro sono
bravi come team, passeremo da forming a norming senza problemi''. È un errore,
perché senza una fase in cui si discute davvero non si sviluppa una
\emph{skill} fondamentale: \textbf{saper discutere}. Tutti conoscono persone
non abituate a farlo, per cui ogni minimo disaccordo diventa un affronto
personale (``avevamo deciso la pizza, perché andiamo al sushi? Perché mi fate
questo affronto?''): sono persone che non sono abituate ad avere conflitti
costruttivi.

Se si pensa di essere in \emph{norming} e continuano a emergere conflitti, il
punto è che \textbf{non si era davvero in norming}: lo storming non ``si
trascina'' nelle fasi successive, semplicemente non se n'è mai usciti.
Allo stesso modo, a chi chiede se ogni mini-fase di storming durante la vita
del gruppo sia frutto di uno storming fatto male o solo di carattere, la
risposta è già contenuta qui.

Un partecipante osserva che spesso, ancor prima di finire un progetto, si è
già messi su quello successivo: è il modo migliore per essere in costante
\emph{damage control}. Era l'esperienza del relatore a Blizzard: veniva
spostato fisicamente da un team all'altro, e finito il lancio di
\emph{Diablo~III} era già su \emph{World of Warcraft}, senza vivere il
``dopo''. Si è dovuto ritagliare da solo lo spazio per un debrief personale,
ed è lì che ha deciso che tipo di manager non voleva essere. Il consiglio:
\textbf{se non lo fa il team, fatelo voi}.

\subsection{Dove si passa la maggior parte del tempo}

\emph{In quale fase si passa più tempo nella propria carriera?} Gli studenti
rispondono \emph{storming} e \emph{norming}. Il relatore, in quindici anni,
calcola di aver passato in \emph{performing} circa \textbf{dieci mesi}. Non è
pessimismo: arrivare a performing deve essere sempre l'obiettivo, sia da
manager sia da membro del team; ma è difficile arrivarci, ed è normale che non
sempre ci si arrivi.

Lo storming può durare a lungo, ma con team molto esperti dura pochissimo,
perché c'è chi lo sa gestire. In realtà \textbf{la maggior parte del tempo si
passa in forming}, perché ogni equilibrio raggiunto vale per \emph{quelle}
persone, \emph{quel} progetto, \emph{quelle} tempistiche. Basta che cambi un
qualsiasi fattore (una deadline spostata, una persona nuova, qualcuno che si
dimette, un progetto cancellato) e si ricomincia, perché i bilanciamenti
precedenti non valgono più.

Esempio: un team di cinque persone affiatatissimo, in pieno performing. La
dirigenza, entusiasta, ``raddoppia le risorse'' e lo porta a dieci. O si
lascia tutto com'è (e i cinque nuovi non fanno niente), o bisogna ridiscutere
chi fa cosa e come i nuovi possono aiutare: che piaccia o no, si ricomincia
da forming.

Quindi la bravura di un manager, ma anche di un membro del team,
\textbf{non sta nel passare tanto tempo in performing, ma nell'accelerare il
passaggio da una fase all'altra}, e nel fare bene la fase di
\emph{adjourning}.

\subsection{Analizzare anche i successi}

C'è la tendenza a fare post-mortem solo dei progetti andati male: ``questa
cosa non ha funzionato, vediamo perché''. Va benissimo. Ma quando un progetto
va bene si dà per scontato che sia andato bene ``perché era buono''. Bisogna
invece chiedersi \emph{perché} è andato bene: era il prodotto? La
tempistica? Il modo in cui abbiamo fatto una cosa piuttosto che un'altra?
\textbf{Se non si capisce davvero il motivo del proprio successo, si
tenderanno a replicare le cose sbagliate.}

\paragraph{Blockbuster.} Quando arrivò la notizia che Netflix avrebbe
lanciato un servizio di streaming, Blockbuster (la catena di noleggio di
videocassette e DVD) aveva tutto il tempo per preparare un'offerta
concorrente. In una riunione del management raccontata in un libro (di cui il
relatore non ricorda il titolo), la direzione sembrava orientata a investire
nello streaming, finché un dirigente fece notare che il 12\% del fatturato
derivava dalle \emph{multe per riconsegna tardiva}, impossibili con lo
streaming, mentre secondo le stime lo streaming avrebbe prodotto il 10\% del
fatturato. ``Non possiamo rinunciare a questo 2\%'': niente streaming. Oggi
non c'è più un Blockbuster ``manco a cercarlo col lanternino''. Non avevano
capito qual era il vero motivo del loro successo, cioè il loro modello di
business\footnote{Le percentuali sono quelle riferite in lezione e non sono
state verificate.}.

\paragraph{The Founder.} Nel film sulla nascita di McDonald's c'è una scena in
cui il futuro socio del protagonista gli dice: ``tu pensi di essere nella
\emph{food industry}, e invece sei nella \emph{real estate industry}''. Non
dovrebbe investire sugli hamburger, ma sui terreni su cui costruisce i
ristoranti. Anche per un team vale lo stesso: bisogna capire che cosa
l'ha reso davvero di successo, oppure no.

\subsection{L'autoanalisi personale}

\paragraph{E quando si cambia azienda?} Uno studente chiede se il debrief sia
possibile anche quando si lascia un'azienda. Non si può tornare in ufficio a
discutere della propria esperienza, ma lo si può fare \textbf{per se stessi}.
Il relatore l'ha fatto, anche se all'epoca in modo poco consapevole, ed è
così che è passato dal \emph{community management} al \emph{people
management}. Analizzando cosa aveva fatto bene a Blizzard si era reso conto
che:
\begin{itemize}
  \item aveva costruito un ottimo rapporto di fiducia con i giocatori (con
        alcuni proprietari di fansite parla ancora oggi);
  \item pur senza la seniority, era diventato un punto di riferimento per il
        team, visto come imparziale ed equilibrato anche su temi che non lo
        riguardavano direttamente;
  \item sapeva \textbf{risolvere problemi}: durante il lancio di
        \emph{Diablo~III} era il più tranquillo di tutti, non perché
        sminuisse il problema, ma perché sapeva che si sarebbe risolto con i
        suoi tempi; il lavoro non era spegnere l'incendio ma ``far sfollare
        la gente'' e aspettare che tutto tornasse in ordine.
\end{itemize}
E aveva capito cosa sapeva fare meno bene, ad esempio \textbf{gestire il
proprio manager}. Avere un manager non significa che il manager gestisca te:
il manager gestisce il tuo lavoro, ma tu devi saper gestire i tuoi
\emph{stakeholder}. Spesso lo stakeholder è il proprio manager; da manager,
gli stakeholder saranno CEO, CFO e così via. Gestirli significa conoscerli,
sapere come portare loro le informazioni, come far presente che qualcosa non
va. Fa parte del lavoro, anche se a lui non piace, e ci ha lavorato.

Nel mese e mezzo di pausa prima di entrare in Riot si è concentrato su
questo: ``cosa porto come mio punto di forza?''. L'autoanalisi è culminata
nel chiedere espressamente alla persona che l'aveva assunto: ``prima mi avete
detto di no, ora mi avete preso: perché?''.

\paragraph{Come si capisce se un progetto è andato bene?} Dipende
dall'\textbf{obiettivo}. Se l'obiettivo era vendere milioni di copie e il
gioco le ha vendute, bisogna capire perché. Se, più ragionevolmente,
l'obiettivo era un gioco da tavolo che permettesse di mettere in pratica le
conoscenze del corso, va valutato per quello: ha permesso di farlo? Se sì,
perché? Se no, perché no? L'obiettivo va fissato \textbf{all'inizio del
progetto, non alla fine}, proprio per poterlo valutare dopo.

\subsection{Quando il team cambia di continuo}

\paragraph{Tutti questi cambiamenti non sono dannosi per il progetto?}
Assolutamente sì, ma quasi mai sono volontari. I progetti andati male sono
spesso progetti grandi.

Il relatore racconta di \emph{Cyberpunk 2077}. Conosceva persone in CD
Projekt Red (``il mondo è piccolissimo, le cose si sanno''). Alla Gamescom va
con Diego Ricchiuti, direttore dell'accademia, a vedere la presentazione del
primo gameplay, quella ``in cui regalavano il giubbottino da samurai''.
Finito il video, tra gli applausi, si guardano e capiscono entrambi che
\textbf{il gioco non è pronto}, a soli quattro mesi dall'uscita prevista:
Diego aveva notato il combattimento e la guida. Gli amici di CD Projekt Red
(uno aveva lavorato con lui a Blizzard per sei anni) gli confermano tra le
righe che ``ce la faremo, ma ci sono ancora cose da sistemare''. Il gioco fu
poi rimandato più volte\footnote{In lezione il relatore dice ``due, forse
tre volte''; l'uscita di \emph{Cyberpunk 2077} è stata effettivamente
rinviata tre volte prima del dicembre 2020.}.

Per deformazione professionale, chi ha sviluppato l'occhio critico non
riesce più a godersi un trailer: in un trailer di \emph{The Day Before} Diego
gli aveva fatto notare che l'effetto della fiammata di un fucile era stato
chiaramente aggiunto in postproduzione e non faceva parte del gameplay. ``È
il triste epilogo di questa storia'': non riescono più a godersi un trailer.

Il punto è che più il team è grande, più queste variabili vengono imposte, e
più è facile che il fallimento \textbf{non dipenda dal fatto che il prodotto
sia brutto} (\emph{Cyberpunk} è un bellissimo gioco sotto molti punti di
vista, e il relatore lo sta rigiocando ora che è stato corretto ed espanso).
Nei grandi gruppi qualcuno arriva, qualcuno si dimette, cambia il budget,
arriva il rinvio; nel caso di \emph{Anthem}, il team ha fatto così tante volte
forming--storming--forming--storming che ne è uscito ``un prodotto che
dovevamo espellere: l'abbiamo espulso, buttato fuori dal nostro organismo pur
di liberarcene''.

Tra i cambiamenti c'è anche la crescita verticale: se il capo viene promosso
perché lavora benissimo, bisogna scegliere un nuovo capo, ed è sempre un
caos (chi ci resta male, chi pensava di meritarlo, chi voleva il suo amico,
chi non sopporta il prescelto). Per questo il relatore è \textbf{contrario
alle promozioni durante i progetti}: si rischia di rovinare tutto il lavoro di
storming.

\subsection{Obiettivi chiari nei piccoli gruppi}

Nei gruppi piccoli, come quelli degli studenti, il problema deriva più
spesso dal \textbf{non avere chiaro che cosa si vuole ottenere}: un prodotto
puramente educativo e formativo, o un prodotto richiesto dall'accademia che,
se viene bene, si potrebbe anche vendere? Un gioco commerciale o un gioco
artistico? Va chiarito subito, perché ci saranno sempre momenti in cui
scegliere tra A e B, e a guidare la scelta sarà la domanda: \emph{quale delle
due opzioni ci porta più vicino all'obiettivo?} Se l'obiettivo non è chiaro,
non si sa cosa scegliere.

Questo porta molte aziende a fare \textbf{scelte di design che in realtà sono
scelte commerciali}. Le microtransazioni sono scelte commerciali: nessuno ha
mai pensato ``sai cosa piacerebbe tanto ai giocatori? Pagare le bustine''.
Prima le bustine si guadagnavano giocando; ora la tendenza è ``potresti
giocare, ma se paghi è meglio''. Si pensi alle polemiche su Ubisoft per
l'opzione di ``saltare il gioco'' pagando, o ai famosi \emph{boost} di livello
di \emph{World of Warcraft}: ``devo pagarti per saltare il gioco?''.

Una studentessa (Silvia) cita \emph{Final Fantasy XIV}, che a un certo punto fu
``resettato'' perché non funzionava: un esempio positivo. Il relatore aggiunge
quello negativo di \emph{Age of Conan}, che non è mai riuscito a fare quel
reset iniziale (``è una schifezza, ve lo ripresentiamo quando è pronto''). Ci
vuole consapevolezza, ma anche i mezzi: un conto è Square Enix, che può
permetterselo, un altro è un'azienda che non può chiudere tutto per un anno
continuando a pagare gli stipendi.

Riassumendo: nei piccoli gruppi i problemi sono più di
\textbf{comunicazione e chiarezza} (obiettivi non chiari o che si spostano);
nelle grandi aziende di più di \textbf{persone che vanno e vengono, obiettivi
e tempistiche che cambiano}. In accademia le scadenze sono fisse (il gioco da
tavolo si consegna a fine anno scolastico, a giugno, e giugno resta); in
azienda può succedere che una feature prevista per la \emph{vertical slice}
di luglio venga anticipata a maggio per portarla alla Gamescom, e un sistema
perfetto per consegnare a giugno può non esserlo più per consegnare a maggio.

\subsection{L'ambiente determina la qualità del prodotto}

\paragraph{È vero che chi passa per Electronic Arts finisce per fare brutti
giochi?} Uno studente cita BioWare, che faceva ottimi giochi finché era
autonoma. Il problema non è dei developer, che non sono ``improvvisamente
diventati incapaci'': è evidente che il problema è altrove, e che
\textbf{l'ambiente in cui si fa un gioco determina molto spesso la qualità del
prodotto}.

Un esempio locale positivo: Tommaso Bonanni, conosciuto dal relatore insieme a
Diego e già docente dell'accademia, ha fondato a Roma Caracal Games Studio,
applicando fin dall'inizio questi concetti di cultura aziendale. L'azienda è
rimasta piccola, ma ha molto chiaro cosa vuole fare e come. Quando si
predispone fin dall'inizio l'azienda e il team a lavorare in un certo modo è
fattibilissimo; è molto più difficile quando si è EA, si acquisisce
un'azienda e le si dice ``da oggi si fa così''.

Il relatore ha vissuto la transizione da Blizzard ad Activision Blizzard (uno
dei motivi per cui se n'è andato, volontariamente). Ha visto la cultura
interna cambiare per gradi:
\begin{enumerate}
  \item ``facciamo giochi fighissimi, perché i giochi devono essere
        fighissimi'';
  \item ``facciamo giochi fighi così vendono di più'';
  \item ``vendiamo di più facendo giochi fighi'' (sembra uguale, ma è molto
        diverso);
  \item ``facciamo più soldi''.
\end{enumerate}
Lo si capiva dalle scelte di design, osservando quale fosse il peso più
grande sulla bilancia. Lo stesso varrà per gli studenti, nei loro giochi e nei
colloqui di lavoro: capire subito che cosa vuole l'altra parte.

\paragraph{Leader e visione.} Nei gruppi dell'accademia, dal secondo anno, c'è
una persona che assume il ruolo di \emph{lead} o \emph{scrum master}. È
diverso dall'essere un \emph{game director} (``il gioco si fa come dico io,
aiutatemi a farlo con la mia visione''), come avviene nei grandi studi: da lì
nascono i nomi nei titoli, come \emph{Sid Meier's Civilization}, che indicano
la persona che ha fatto quel gioco, oppure un puro branding, come \emph{Tom
Clancy's Rainbow Six}.

\paragraph{Inclusion come jolly.} A margine, uno studente osserva che in molte
aziende ``la risposta giusta è sempre \emph{inclusion}'', usata come un
jolly. È comunissimo, conferma il relatore: \emph{diversity and inclusion}
viene spesso usata per dire ``così includiamo più clienti'', ma non è quello
il suo punto (``altro discorso'').

% ---------------------------------------------------------------------
\section{Confrontarsi con chi è più bravo}
\label{sec:z1-rivalita}

Arrivato a Riot dopo sei anni a Blizzard, il relatore pensava di avere un buon
bagaglio di esperienza. Si è trovato invece ``come un fante buttato in mezzo
ai Navy SEAL'': persone che avevano superato 11, 12, 13 colloqui, con un
livello molto alto.

Individua subito un collega francese che, nella sua testa, diventa ``la
versione migliore di me'': stesso ruolo (community manager), esperienza
simile (arrivava da Ubisoft, come quasi tutti i francesi di Riot), ma molto
più spigliato, sempre con la battuta pronta, ``l'anima della festa''. Il
relatore invece è introverso: un workshop come questo per lui è
\emph{emotivamente drenante}, non ama parlare in pubblico pur amando
condividere le proprie esperienze, e non riesce a sostenere conversazioni
come ``cosa hai fatto questo weekend?''. Lavorando raramente insieme, aveva
sviluppato una rivalità ``da allenatore di Pokémon'' (lui era Ash, il
relatore l'antagonista).

Dopo sei mesi, a un evento di \emph{team building} aziendale (a Tenerife), il
suo capo, molto attento a queste dinamiche, li mette apposta nello stesso
gruppo. A fine serata, complice l'open bar, il relatore gli confessa di averlo
mal sopportato dal primo giorno per invidia della sua estroversione. Il
collega ride e risponde che pochi giorni prima aveva detto al capo la stessa
cosa al contrario: invidiava a Stefano il suo essere sempre organizzato, sul
pezzo, presentarsi nella sala riunioni giusta (mentre lui finiva in quella
sbagliata), arrivare con gli appunti, con i dati a supporto e le analisi di
mercato, mentre lui spesso non aveva nemmeno letto la mail del meeting.
Un'invidia in due direzioni, che li ha portati a lavorare molto insieme e a
migliorarsi a vicenda.

Le lezioni:
\begin{itemize}
  \item non bisogna demoralizzarsi vedendo qualcuno che fa lo stesso lavoro
        ed è più avanti sotto alcuni aspetti: spesso c'è un bilanciamento. I
        ``punti caratteristica'' sono limitati: se li ha messi in carisma,
        non li avrà messi in saggezza;
  \item conviene \textbf{chiedere a queste persone come mai} sono così brave
        in qualcosa, come si fa con un progetto andato bene: magari non lo si
        farà mai bene quanto loro, ma si comincerà ad avere il
        \emph{mindset} giusto; e viceversa, loro saranno spinti a chiedere a
        noi;
  \item spesso \textbf{le proprie doti le scopriamo dagli altri}. Il relatore
        non si era mai accorto di essere ``quello che resta calmo durante le
        emergenze'' (il lancio di \emph{Diablo~III}, quello di
        \emph{Draenor}, i problemi dal vivo alla Games Week) finché non
        gliel'hanno fatto notare tutti: per lui era la normalità. A chi pensa
        ``non sono capace, forse non è il mio percorso'', può darsi che
        semplicemente nessuno gli abbia ancora detto ``questa cosa la sai
        fare davvero bene, concentrati su questo''.
\end{itemize}

% ---------------------------------------------------------------------
\section{Domande sul lavoro nell'industria}
\label{sec:z1-industria}

Diverse domande degli studenti hanno riguardato le condizioni di lavoro nei
grandi studi internazionali. Le risposte sono raccolte qui per argomento.

\subsection{Presenza, remoto e trasferimenti}

\paragraph{Grandi studi: quasi solo in presenza.} Sia a Blizzard sia a Riot si
lavorava \textbf{solo in presenza}, senza smart working: sei anni a Parigi e
cinque a Dublino. Si parlava solo inglese: anche nei meeting tra soli
italiani si tendeva a comunicare e a scrivere in inglese, per correttezza
verso chi si fosse aggiunto e per lasciare una traccia comprensibile a tutti.
Il francese serviva per fare la spesa e sopravvivere, non in ufficio.
L'ambiente era estremamente eterogeneo (francesi, polacchi, spagnoli e
molti altri): durante il colloquio dal vivo a Dublino, mentre si preoccupava
di essersi vestito in modo adeguato, vide passare dietro la parete a vetro un
collega svedese con i dreadlock, a piedi scalzi, con una pianta sotto un
braccio e un aquilone sotto l'altro, e capì di essere ``a posto così''.
Un'esposizione a molte culture e modi di pensare che il relatore considera
bellissima.

Anche oggi, tranne alcuni studi particolari, la maggior parte del lavoro si
svolge di persona: amici a Paradox lavorano a Stoccolma, i pochi rimasti a
Blizzard si sono dovuti trasferire a Irvine. Per posizioni \emph{junior} o
\emph{entry level} i grandi studi raramente concedono il lavoro da casa; per
profili senior può esserci qualche eccezione (ad esempio nove mesi da remoto
e tre in sede).

\paragraph{Relocation.} Spesso tocca spostarsi, ma il relatore non ha mai
pagato un trasloco di tasca propria: queste aziende offrono la
\emph{relocation} (aiuto a trovare casa, ad aprire il conto in banca e
supporto pratico), perché è un investimento sulla persona.

\paragraph{Colloqui mentre si lavora altrove.} Per i colloqui in presenza
prendeva ferie, senza dire dove andava. Il consiglio generale: per fare
colloqui \textbf{prendete ferie o permessi}, non datevi malati, ``che non è un
bel modo di comportarsi''.

\paragraph{Remoto ``doppio''.} Un collega lavora per l'azienda coreana che
gestisce \emph{PUBG}, con ufficio ad Amsterdam, ma lavora da Milano: un
``doppio remoto''. Si trova, ma in genere un ambiente anglofono con sedi in
Europa chiederà di trasferirsi in Europa.

\paragraph{Famiglia e trasferimenti.} Il relatore ha avuto la fortuna di
sposare una persona che fa lo stesso lavoro (conosciuta lavorando su
\emph{Warhammer Online}, poi con lui a Blizzard e a Riot, nella
localizzazione e nel doppiaggio, ad esempio di \emph{World of Warcraft} in
italiano): avendo competenze rare, spesso venivano assunti entrambi, al
massimo a pochi mesi di distanza. Per i developer la soluzione frequente è
che la famiglia segua la persona assunta, perché lo stipendio basta per
tutti. Ma non sempre c'è equilibrio, e a volte uno dei due deve sacrificarsi:
sua moglie ha lasciato Blizzard per seguirlo a Dublino (per poi essere
assunta da Riot); lui ha lasciato Riot per tornare in Italia quando lei ha
trovato lavoro qui. Ha dovuto rinunciare più volte a lavorare in Islanda per
CCP (\emph{EVE Online}), perché non era una scelta solo sua. Si fanno
compromessi, ma si riesce a far quadrare.

\paragraph{Ci si sposta per ogni progetto?} No: si sono spostati
esclusivamente per cambiare azienda. Cambiando progetto nella stessa azienda
è raro, quasi mai (``non l'ho mai visto succedere''), che venga chiesto di
trasferirsi. In Francia stavano benissimo (vivevano fuori Parigi; ``i
problemi con i francesi sono in realtà problemi con i parigini''), parlavano
francese ed erano vicini alla doppia cittadinanza: se ne sono andati perché
non vedevano un futuro a Blizzard, per il modo in cui l'azienda trattava i
colleghi.

\paragraph{Il mercato oggi.} Non esiste un \emph{silver bullet} per trovare
lavoro domani dove si vuole. Con tutti i \emph{layoff} recenti ci sono
migliaia di esperti del settore che cercano lavoro: serve pazienza, ma è
fattibile.

\paragraph{Studi indie e remoto.} Negli studi più piccoli il remoto è più
diffuso: a Stormind circa 60 persone su 95 lavorano da remoto, ma lo studio
si è strutturato fin dall'inizio per farlo. Serve consapevolezza. Bending
Spoons (che sviluppa app per mobile, non giochi, ed è l'azienda dell'app
Immuni), pur essendo grandissima e con sede a Milano, lavora da remoto e
lascia l'ufficio a chi vuole: lo smart working è posizionato quasi come un
\emph{benefit}, ma permette anche di avere un ufficio da 300 persone invece
che da 700. Le sedi fisiche costano decine di milioni all'anno (Blizzard ha
chiuso le sedi di Amsterdam e di Parigi). Tuttavia, per lavori molto
verticali e specifici gli studi preferiscono avere le persone dal vivo; i
ruoli più spesso remoti sono quelli che non fanno sviluppo (customer service,
PR\dots).

Anche gli artisti lavorano spesso in presenza: l'accademia può dare solo una
\emph{simulazione} di quanto stretto sarà il contatto tra programmazione,
design e art. ``Sarà ancora di più.''

\paragraph{La comunicazione da remoto.} È più difficile. Pur essendo
estremamente introverso, il relatore \textbf{obbliga tutti a usare la webcam}
nei meeting in cui conta: non in un aggiornamento con novanta persone
collegate, ma certamente in un colloquio a due sulla performance dell'ultimo
anno. La componente fisica è una parte enorme della
comunicazione\footnote{In lezione si parla di ``quasi l'ottanta per cento'';
il dato va preso come indicativo, poiché le percentuali di questo tipo
circolano in forme diverse e sono spesso citate fuori contesto.}: anche se si
è timidi, accendere la webcam è l'approssimazione più vicina all'essere dal
vivo. Negli ambienti dove si lavorava senza webcam la situazione era
``esponenzialmente peggiore''.

\subsection{Differenze culturali nel team}

\paragraph{Gestire sensibilità diverse.} A Riot i team erano divisi per
\emph{lingua}, non per nazione: il team di lingua spagnola includeva anche un
messicano, quello di lingua russa includeva russi e ucraini, proprio nel
periodo della guerra in Crimea. Tra i colleghi non c'era alcun problema (erano
amici), ma ci si è posti la domanda: \emph{come comunichiamo con una community
che sarà divisa su questo tema?} Nei giochi online la divisione era evidente:
in \emph{Rocket League} c'era chi, con un compagno di squadra con la bandiera
avversaria, si faceva otto autogol pur di farlo perdere. Allo stesso modo,
con molti americani mandati a Dublino a coprire posizioni scoperte, c'erano le
elezioni (Clinton--Trump, Biden--Trump) e varie guerre: serviva sempre una
certa sensibilità.

Imparare a navigare queste situazioni lo ha reso un professionista migliore:
gli ha insegnato a fermarsi e chiedersi ``io la percepisco così; gli altri
dieci, con background, storie e lingue diverse, come la penseranno?'', a
ragionare più \emph{long term} e più \emph{big picture}, con più empatia.

\paragraph{Stili di comunicazione.} Un partecipante che ha lavorato nel Regno
Unito chiede dello stile comunicativo italiano, a volte frainteso. Uno dei più
grandi complimenti ricevuti dal relatore è stato ``non sembri italiano''
(``gesticoli poco per essere italiano'', ``parli bene inglese per essere
italiano''). Ha imparato ad \textbf{adattare la comunicazione alla persona}:
con spagnoli e francesi la battuta sulla nazione funziona, perché c'è un
``odio--amore'' condiviso; con russi e giapponesi no. Doveva adattare anche
l'umorismo. Un collega metà russo e metà messicano diceva di aver preso ``il
lato peggiore di entrambe le nazionalità'', ed è stato un ottimo maestro per
imparare a leggere culture diverse.

Non bisogna generalizzare, ma gli stereotipi possono essere un punto di
partenza ragionevole, \emph{sapendo che troverete sempre le eccezioni}. E
vanno letti correttamente: gli americani tendono a essere più estroversi e ad
avere sempre qualcosa da dire non perché ``sono fatti così'', ma perché
vengono educati così (all'università tra i corsi c'è il \emph{debate}, saper
argomentare in pubblico); un collega giapponese può essere più introverso per
ragioni culturali. Nel dubbio si parte da lì, poi si conosce la persona.

\subsection{Negoziare, discriminazione e stipendi}

\paragraph{La famiglia come leva.} Avere un partner riduce le opzioni, ma le
aziende lo sanno: uno spostamento è per due. Si può usare come leva per
negoziare (un compenso diverso, un supporto per il trasloco): se un'azienda ha
deciso di volere quella persona, 2, 3, 4 o 10.000 euro in più all'anno, per
un'azienda che fattura centinaia di milioni o miliardi, non cambiano nulla.

\paragraph{Essere sposati è un motivo per non essere assunti?} No. In alcuni
paesi è proprio illegale discriminare su questo: in Francia era illegale
chiedere CV con data di nascita, foto, sesso o religione, e chiedere se si
vogliono figli. Purtroppo, dai racconti dei suoi studenti di master (in
settori come \emph{beauty and wellness} e \emph{food and wine}), in piccole
realtà italiane queste domande (``hai 25 anni, hai intenzione di
sposarti?'') vengono ancora fatte, e servono a filtrare i candidati. Nelle
grandi aziende ci sono sistemi per evitarlo; anzi, molte danno \emph{più}
benefit a chi ha famiglia (``porta pure la famiglia, è tutto pagato''),
perché le persone competenti hanno spesso una certa età e quindi,
ragionevolmente, una famiglia.

\paragraph{Gli stipendi.} Dipendono molto da dove si vive. Lo stipendio a Riot
non era paragonabile a quello di Blizzard perché in Francia era tassato al
52\% e in Irlanda al 17\%\footnote{Aliquote come riportate in lezione,
riferite alla situazione personale del relatore.}: nessuno stipendio
competitivo può compensare il 35\% di differenza di tasse. Negli Stati Uniti,
con pochissima previdenza pubblica, gli stipendi sono molto più alti perché
poi ci si paga tutto (``se hai un raffreddore paghi tu''). Bisogna valutare il
potere d'acquisto di quei soldi nel paese in cui si andrà a vivere.

\subsection{Freelance e \emph{work for hire}}

All'inizio della sessione pomeridiana uno studente che lavora nell'audio
chiede quali siano le condizioni per chi lavora come \emph{freelance}.
\begin{itemize}
  \item Dipende molto dal progetto: spesso il contratto è a progetto, per
        portare a termine una parte o un intero progetto.
  \item Non tutte le aziende usano freelance: Riot e Blizzard lo facevano
        pochissimo e preferivano assumere (magari trattando poi la persona
        ``come un freelance'', lasciandola a casa a fine progetto).
  \item Spesso i lavori offerti come freelance sono in realtà
        \textbf{sostitutivi di un full time} (``freelance, ma sei ore al
        giorno'').
  \item Si può fare, ma bisogna essere molto chiari nel definire l'impegno
        (quante ore, come vengono tracciate): non tutti specificano cosa
        significa davvero ``freelance''. Lavorando così per due o tre aziende
        si può arrivare a un tempo pieno.
  \item È più comune il \emph{work for hire}: aziende che forniscono personale
        ad altre aziende; si è assunti da X ma si lavora ``embeddati'' presso
        Y (una sorta di subappalto). Anche lo studio del relatore lo fa.
  \item Il puro freelance di quattro mesi esiste, ma di solito per quei
        quattro mesi diventa un lavoro a tempo pieno; è raro ``quattro mesi,
        due ore al giorno''. Nelle grandi aziende è improbabile lavorare nello
        stesso mese per Ubisoft, Riot e Blizzard tre ore ciascuna; è più
        realistico lavorare per più startup indie italiane, che non possono
        permettersi un'assunzione ma possono pagare due ore al giorno.
  \item Molto passa dai contatti: trovare le aziende giuste e rendere esplicite
        le condizioni.
\end{itemize}

% ---------------------------------------------------------------------
\section{Il circolo virtuoso: feedback, miglioramento, fiducia}
\label{sec:z1-triangolo}

\subsection{Il triangolo}

L'obiettivo a cui il workshop vuole arrivare è un \textbf{triangolo che si
autoalimenta}:
\begin{enumerate}
  \item si dà un \textbf{feedback} con l'obiettivo di migliorare le
        performance (inclusi elementi caratteriali);
  \item la persona \textbf{migliora} e vede di averne beneficiato;
  \item questo genera \textbf{fiducia}, che rende più facile dare e ricevere
        il feedback successivo.
\end{enumerate}
Togliendo uno qualsiasi dei tre elementi, il circolo non è altrettanto forte.

In qualsiasi fase si parta, il rispetto per i colleghi è inizialmente dettato
dall'\emph{educazione}, non dalla fiducia. La fiducia arriva quando i colleghi
cominciano a dire cose giuste e sensate e a \textbf{mantenere la parola}.
Dare e ricevere feedback è il modo in cui, in un certo senso,
``inganniamo il cervello'' per creare una \emph{case history} di fiducia: si
ha motivo di fidarsi di più, di ascoltare di più, di essere più ricettivi. E
si diventa sempre più bravi ad accettare il feedback e ad agire di
conseguenza.

\subsection{Accettare un feedback non vuol dire fare come dice l'altro}

Dieci anni fa il relatore non era ostile ai feedback, ma molto
\emph{difensivo}: ``sì, sì, va bene, ora ci guardo''; oppure, a chi gli diceva
``avresti dovuto farla così'', rispondeva spiegando perché l'aveva fatta in
quel modo. Oggi ha imparato a dire ``ok, grazie, ora ci ragiono'': non deve
difendere la propria posizione da chi viene solo a dire ``hai pensato a
questa cosa?''. Prima lo viveva come un attacco (``mi viene a dire che non so
fare il mio lavoro''), ora come un ``ottimo punto di vista, lo prendo in
considerazione''.

Con una postilla fondamentale: \textbf{accettare il feedback non significa
fare come dice l'altro}. ``Mi hai detto che preferisci questo mob verde. Ho
pensato anche a farlo verde? Sì. Penso comunque che giallo sia meglio? Sì, lo
faccio giallo.'' Il relatore accetta \emph{sempre} il feedback, valido o no:
garantisce di tenerlo a mente quando lavora. Che poi faccia come gli viene
detto è un altro discorso.

\subsection{Fiducia e scadenze}

La fiducia funziona allo stesso modo con gli impegni. Se si chiede a qualcuno
una cosa per venerdì e la consegna venerdì, la volta dopo gliela si chiederà
per venerdì. Se invece arriva il giovedì della settimana dopo, la volta dopo
gli si darà una \textbf{deadline anticipata}, come con chi arriva sempre in
ritardo agli appuntamenti (``se ti dico alle 7 arrivi alle 7 e mezza; la
prossima volta ti dico alle 6 e mezza'').

Al lavoro questo è dannosissimo, perché crea \textbf{deadline fittizie} che
non sono mai quelle vere, e nascono discussioni del tipo ``perché devo farlo
in due giorni se si fa in cinque?'' (``sì, se tu lo facessi in cinque''). La
conversazione giusta è un'altra: \emph{``che problemi hai avuto nel
consegnare venerdì?'' ``Questo, questo e questo.'' ``Cosa posso fare per
evitare che succeda?''}

Il relatore non vuole mai sentirsi dire ``non sono riuscito a farlo''.
Preferisce di gran lunga sentirsi dire, tre giorni prima, \textbf{``potrei non
riuscire a farlo, per questi motivi''}. Chi ha la visione d'insieme può dare
più budget, più persone, spostare la deadline, cambiare lo scopo del progetto:
può fare qualsiasi cosa per aiutare. Ma se lo si scopre il venerdì, cinque
minuti prima della deadline, non può fare più niente. Il messaggio è:
\textbf{``aiutami a poterti aiutare''}. Non penserà mai male di chi dice ``non
ce la faccio per questi motivi'', perché non sta dicendo ``non ho voglia'': sta
indicando i problemi, e lavorare su quei problemi è compito del manager. Non
può invece fidarsi di chi fino al giorno prima dice ``tranquillo, ce la
faccio'' e il giorno dopo ``non ce l'ho fatta''.

Questo è fondamentale per chi lavora in programmazione, design o art, perché
ci saranno sempre deadline: bisogna sempre \textbf{lavorare sulle aspettative},
dire al proprio manager (o al proprio report) cosa è lecito aspettarsi. Non
bisogna fare come i pinguini di \emph{Madagascar} (``sorridete e salutate, va
tutto bene''): se c'è un problema, va detto, altrimenti si peggiora solo la
propria situazione.

Da italiani si tende a non voler mai alzare la mano per chiedere aiuto
(``no, no, ce la facciamo, lo faccio io''). A volte non ce la si può fare, non
perché non si è capaci, ma perché vari fattori lo impediscono. Accettarlo è
importante: \textbf{è molto meglio avvisare e chiedere una mano che dire
``non ce l'ho fatta''}.

\subsection{Domande sul dare e ricevere feedback}

\paragraph{``Ne parliamo a fare, se poi fai come vuoi?''} Uno studente (Cristian)
riceve feedback che, dopo averli analizzati, non ritiene congrui; quando non
li segue gli viene detto ``ne parliamo a fare?''. C'è un problema di
\textbf{aspettative}: l'idea che accettare un feedback voglia dire fare come
dice l'altro. Accettarlo significa prendere in considerazione ciò che è stato
detto. Chi decide diversamente ha però un \textbf{onere}: argomentare perché
si fa comunque a modo proprio (``mi hai detto verde, la facciamo gialla
perché\dots''). Poi bisogna solo aspettare i risultati per vedere chi aveva
ragione.

La decisione spetta sempre \textbf{a chi ha la responsabilità di quella
decisione}. Se il game director deve scegliere tra A e B, anche se il
feedback dice B, la scelta è sua perché la responsabilità è sua. Non sarebbe
corretto essere ritenuti responsabili di una scelta delegata a chi dà
feedback. La risposta può essere semplicemente: ``l'ho tenuto in
considerazione, ma visto che ne sono responsabile io e le conseguenze sono
mie, ho scelto di fare X''.

\paragraph{E se l'altro fa ripicca?} Un partecipante chiede di chi, vedendo
il proprio feedback non accolto, si vendica ignorando a sua volta feedback
buoni. Per il relatore è un problema caratteriale: chi si comporta così ``è
già sulla strada per uscire''. Non c'è spazio per l'ego in team di questo
tipo. Il fenomeno nasce dall'idea che \emph{il feedback negativo faccia fare
brutta figura}, che il fatto che si potesse fare meglio tolga qualcosa alle
proprie capacità.

Non è così: \textbf{il feedback non è un giudizio}. Si può passare un anno a
Riot ricevendo solo feedback negativi e avere una valutazione annuale
stellare, perché ricevere molti feedback negativi significa essere una persona
approcciabile, che agisce sul feedback e invita gli altri a darlo. Il relatore
ha visto valutazioni molto peggiori per chi riceveva solo feedback positivi,
perché la gente aveva paura di dargliene di negativi: ma solo ``pacche sulle
spalle'' non fanno migliorare. Il problema va affrontato parlando con la
persona; a volte le conversazioni iniziano con \textbf{``ti do un feedback sul
feedback''}, cioè sul modo in cui l'altro ha dato il suo.

Esempi dalla sua esperienza:
\begin{itemize}
  \item lui chiede spesso di \textbf{non ricevere feedback senza esempi
        tangibili}, perché fa fatica ad astrarre: ``sei stato maleducato'' non
        significa nulla, ``ieri al meeting mi hai risposto male su questa
        cosa'' sì;
  \item ad altri serve l'opposto: ha ricevuto il feedback ``quando devi dirmi
        che ho sbagliato, puoi girarci un po' intorno e darmi del contesto?'',
        perché il suo stile diretto (``questa cosa non ha funzionato, dovremmo
        fare così'') feriva. Ha adattato il suo modo di dare feedback: con
        alcuni è diretto, con altri gira intorno. L'importante è saperlo.
\end{itemize}

\subsection{Non portare la soluzione}

Un consiglio centrale: \textbf{quando si dà un feedback, non si dà una
soluzione}. Non si dice ``io farei così'', ma ``secondo me questa cosa si
potrebbe fare diversamente, per questo e questo motivo: come potremmo farla
diversamente?''. È come chi smette di fumare perché vuole e chi smette perché
glielo dicono moglie e figli: se non si è convinti, spesso non si smette.

Quando il relatore diceva ``c'è questo problema, io farei così'', il meeting
successivo iniziava inevitabilmente con ``ci ho provato come hai detto tu, con
me non ha funzionato''. Sempre. Perché la soluzione era sua. La
responsabilità della soluzione deve essere invece \textbf{congiunta}: ogni
feedback deve creare un piccolo team che ha come obiettivo risolvere quel
problema. Questo toglie l'antagonismo (non ``ti dico che non sai fare il tuo
lavoro'', ma ``ti voglio aiutare a farlo meglio, lavoriamoci insieme'') e
sposta il focus: il problema non sei tu, è ciò che sta accadendo.

\paragraph{E se sono i piani alti a tenere all'oscuro dei problemi?} È una
questione di cultura aziendale. Il relatore fa sempre presente al proprio capo
il concetto di \textbf{service level agreement} (che verrà ripreso nella
seconda giornata): ``se mi dai questi input, ti garantisco questo output in
questi tempi''. Se il capo chiede delle cose senza dire budget, tempi e
perché, non può aspettarsi lo stesso risultato. È una comunicazione a due
vie: si dice al manager cosa serve per fare il lavoro, e lui deve fornirlo o
accettare che il risultato non sarà quello voluto.

\subsection{Amici, parenti e parole proibite}

\paragraph{``Non si lavora con amici e parenti''?} Per sua definizione
personale il relatore ha pochissimi amici (``stanno sulle dita di una mano'')
e una soglia molto alta per considerare amico un collega. Allo stesso tempo,
in quindici anni e con centinaia di colleghi, ha avuto solo \textbf{due}
esperienze in cui non è riuscito a lavorare con qualcuno: non serve essere
amici per far funzionare le cose. Con gli amici, anzi, gli viene più facile:
se dice a un amico che una cosa ``fa schifo'', l'amico non pensa ``come si
permette'', ma ``se perfino Stefano, che è mio amico, me lo dice, allora è
proprio così''. Diventa un rafforzativo, grazie alla fiducia costruita prima.
Il problema del lavorare con amici e parenti nasce dal \emph{non voler} dire
all'amico che sta sbagliando; se invece la confidenza permette di dirlo, si
può saltare la ``confezione'' del messaggio e passare direttamente al
contenuto.

\paragraph{Le frasi peggiori.} Alla domanda se esistano parole peggiori di
``ottimo lavoro'', il relatore indica la sua classifica: ``ottimo lavoro''
(detto in modo vuoto)\footnote{Il passaggio della trascrizione è
frammentario: si intende verosimilmente il complimento generico e di
circostanza, che non dice nulla di utile.}; poi la frase che lo ``fa uscire di
testa'', \textbf{``abbiamo sempre fatto così''} (``e infatti avete chiuso'',
come Blockbuster); infine \textbf{``siamo una famiglia''}, riferito al posto di
lavoro. Un'azienda non è mai una famiglia, perché non ci si può aspettare lo
stesso impegno da chi la possiede e da chi ci lavora: non c'è un
corrispettivo proporzionale tra impegno e ritorno. Per quanto inclusiva,
flessibile e razionale, non sarà mai una famiglia.

\subsection{Ego, muri e il peggior modo di scegliere un leader}

\paragraph{Chi mette davanti l'ego e alza muri.} Esistono, ma di solito durano
pochissimo in azienda. Non si tratta di avere un ego forte, ma di \textbf{non
saper ricevere feedback}. Spesso sono persone bravissime nel loro lavoro che
poi si chiedono perché tutti intorno diventino più bravi di loro: perché gli
altri accettano feedback e hanno più fonti da cui attingere, mentre loro si
sono rintanati nella loro torre.

Molti manager tendono ad \textbf{assumere persone un filo peggiori di loro},
per non fare brutta figura. Il relatore assume solo persone più brave di lui,
``perché non è una gara'': ha lavorato con molti community manager più bravi
di lui come community manager, ma lui era più bravo a fare il manager. Non
viene valutato su quanto sa fare il community manager, ma su quanto il suo
team sa farlo. Un partecipante chiede se la paura di avere qualcuno di
migliore esista anche in ambito accademico. Assolutamente sì: al relatore è
capitato di non essere preso come docente in un corso di formazione sui
videogiochi perché il responsabile aveva percepito che ne sapesse più di
lui\footnote{La trascrizione dice ``che non sapevo più di lui'', ma il senso
del passaggio richiede il contrario.}. ``Ma perché te ne importa? Tu sei quello che ha portato chi
sa fare il suo lavoro: ben venga.'' O si ha questa cultura, o si ha paura per
il proprio posto.

\paragraph{Promuovere il migliore.} Quando bisogna scegliere un leader per un
team (programmazione, design, art), molto spesso si sceglie \textbf{il migliore
tecnico}: il miglior programmatore, il miglior designer. È spesso la scelta
sbagliata: è come prendere un idraulico per fare l'elettricista. Fare il
\emph{people manager} non è come fare il designer o il programmatore; si può
essere un ottimo programmatore e un pessimo manager, e viceversa. Le aziende
l'hanno capito e lo fanno meno di prima. A Blizzard c'era un grave problema di
\emph{middle management}: molte persone in ruoli importanti erano state
promosse cercando le skill sbagliate, erano frustrate dal non saper fare il
proprio lavoro e schiacciavano chi stava sotto per paura di fare brutta
figura.

Un partecipante cita una teoria simile: una persona viene promossa finché non
raggiunge \emph{il proprio livello di incompetenza}, dove si blocca\footnote{È
il \emph{principio di Peter} (Laurence J.\ Peter, 1969), citato in lezione
come ``la teoria di Peter''.}. Un altro cita Michael Scott di \emph{The
Office}, su cui il relatore tornerà (sezione~\ref{sec:z1-quando}).

% ---------------------------------------------------------------------
\section{Perché la cultura conta: retention, costi e fiducia}
\label{sec:z1-cultura}

\subsection{L'ambiente di lavoro e la retention}

Perché tutto questo è così importante? Perché l'ambiente di lavoro conta
moltissimo per la \textbf{retention}, cioè per trattenere le persone e farle
lavorare in modo efficiente. Se contasse solo lo stipendio, molti non
lascerebbero mai Riot o Blizzard; e più si va avanti con l'età, più si impara
a dare valore anche al resto.

Inoltre, secondo gli studi citati nelle slide, chi vive il proprio lavoro in
modo negativo è \textbf{meno efficiente nel problem solving}: l'energia
mentale va verso ``come faccio a non farmi dare la colpa'' o ``come faccio a
fare bella figura'', non verso la soluzione del problema. In ambienti negativi,
dove ci si deve preoccupare della propria ``sopravvivenza'', si lavora
peggio.

\subsection{Il costo di chi se ne va}

Ogni volta che qualcuno si dimette o viene licenziato, il costo di trovare
una nuova persona, fare i colloqui, formarla e farle attraversare forming,
storming, norming\dots\ è circa \textbf{un terzo del suo stipendio}. Per questo
le aziende tendono a fare più \emph{retention} che \emph{acquisition}, almeno
con le figure considerate utili alla produzione.

Allo stesso tempo si vedono sempre più spesso \emph{layoff} e \emph{hiring
wave}: si licenziano 3.000 persone e sei mesi dopo se ne riassumono 2.500 o
3.500, ma tutte più junior. Il relatore ne ha scritto in un articolo su
Eurogamer.it, in cui spiega anche alcuni ``trucchi del mestiere'': perché i
licenziamenti avvengono sempre in certi periodi e sempre ``a gruppi'' (prima
Bethesda, poi Blizzard, poi un'altra\dots).

\paragraph{L'offerta di Riot.} A Riot esisteva una pratica (di cui il relatore
non ricorda il nome) per cui, firmato il contratto, si aveva \textbf{un anno
di tempo per dire ``non mi piace, non me la sento, me ne vado''} e ricevere
subito come buonuscita un anno di stipendio. Avevano stimato che fosse molto
meglio pagare subito un anno a qualcuno e lasciarlo andare che tenerlo due,
tre, quattro anni come \emph{underperformer}: la responsabilità della scelta
di assumere era dell'azienda\footnote{La pratica, nota come ``The Offer'', è
descritta in fonti pubbliche con condizioni diverse (un periodo più breve e
una somma pari a una frazione dello stipendio annuo). Si riportano le
condizioni come raccontate in lezione; andrebbero verificate.}. L'hanno usata
pochissimi, ma il solo fatto che esistesse faceva lavorare meglio anche lui da
manager, perché sapeva che chi non si trovava bene poteva andarsene con un
anno di tempo per trovare un altro lavoro.

\paragraph{Il ricircolo.} C'è un forte ricircolo nell'industria: molti ex
colleghi passano da Ubisoft a Paradox, da Paradox a Funcom, da Funcom a
Blizzard. Il relatore ha scelto di uscirne lavorando solo per aziende
\textbf{indipendenti}: Stormind non è un indie di cinque persone, ma è
indipendente, e non deve accontentare azionisti. Quando vede un'azienda con
utili miliardari licenziare migliaia di persone si chiede ``dov'è la linea''
(\emph{when enough is enough}): con azionisti che vogliono un ritorno, se
quello è l'obiettivo stanno facendo la cosa giusta, ma lui non condivide
l'obiettivo\footnote{In lezione si parla di ``6,5 miliardi di utile'' e di
``3.000 persone licenziate''; le cifre sono riportate come dette e non sono
state verificate.}.

\paragraph{L'open space vuoto.} A Riot, durante un'ondata di licenziamenti,
del suo team di trenta persone è rimasto solo lui: ``teniamo te come
coordinatore, poi ricostruirai il team''. Ha rifiutato di ricostruire un team
cancellato ``con un colpo di spugna'' e ha deciso che non avrebbe più lavorato
per una grande azienda. Nessuno aveva pensato di spostarlo, e ha lavorato per
tre mesi in un open space di cento metri quadri con la sua scrivania e altre
ventinove vuote. Su una lavagna aveva scritto i nomi dei suoi ragazzi e li
cancellava man mano che, anche grazie ai suoi \emph{referral}, trovavano
lavoro. Cancellato l'ultimo nome, si è dimesso: per lui il lavoro era finito.

\paragraph{Team che si dissolvono.} Uno studente cita un articolo su Larian
Studios, il cui management è contrario a ricostruire ogni volta il team a fine
progetto. È lo stesso problema: si pensa ai ``giochi Blizzard'' o ai ``giochi
BioWare'', ma le persone che li hanno fatti non ci sono più; è rimasto solo il
brand, e non ha senso avere le stesse aspettative da un gioco fatto da persone
diverse. E anche chi sopravvive ai layoff vive pensando che il prossimo
potrebbe essere lui: alcune persone, licenziate nell'ultima ondata a
Blizzard, gli hanno detto ``erano due anni che lo sapevo''. Vivevano a Cork,
``un paesotto irlandese dove o vai per lavorare a Blizzard o non ci vai'',
avevano comprato casa: non ci si può aspettare chissà che morale.

\subsection{La cultura è forte quando si lavora per gli altri}

\begin{quote}
  \emph{La cultura è forte quando le persone lavorano per gli altri e a favore
  degli altri, e non per se stesse.}
\end{quote}

Se si dà feedback per far migliorare l'altro e non per imporre il proprio
modo, se si aiuta per il bene del progetto e non per fare bella figura, la
cultura si rafforza. È dai \textbf{piccoli gesti} che emerge la cultura
aziendale: ad esempio, dal non dover chiedere ``perché mi stai aiutando?''.

\paragraph{L'incognita.} \textbf{L'incognita è spesso dannosa quanto le
cattive intenzioni.} In un esperimento (``mi pare a Stoccolma, cinque o sei
anni fa'') si distribuivano per strada banconote da 500 euro, e la maggior
parte delle persone le rifiutava chiedendosi ``dov'è la fregatura?''
\footnote{L'esperimento è riportato in modo vago in lezione e non è stato
possibile identificarlo con certezza.}. Lo stesso accadeva al relatore,
arrivato a PG e a Stormind da Riot: a ``vedo che sei in difficoltà, posso
aiutarti?'' la risposta era sempre ``perché?''. Alla fine rispondeva: ``perché
quando sarò in difficoltà io vorrei che tu mi aiutassi''.

Quando si comincia a lavorare per gli altri e si smette di pensare ``così avrò
la promozione, così avrò l'aumento'', promozioni, aumenti e rispetto arrivano
da soli. A Blizzard aveva fallito tre volte i colloqui per essere promosso a
manager (inteso come coordinatore del team); quando ha deciso di smettere di
pensarci e di fare solo il bene del proprio lavoro, la volta dopo è stato
promosso. Funziona così: chi aiuta i colleghi e fa andare bene il progetto
viene notato (``questa persona fa migliorare cinque persone, non solo se
stessa: vale di più'').

\subsection{Il lavoro è valutato come opera d'insieme}

Il relatore riconosce che buona parte del suo successo è venuto dal fatto che
il suo team lavorava bene, non dalla sua bravura personale. Anche i progetti
degli studenti, presentati a Diego, a Enrico, ad Andrea, verranno valutati
\textbf{come opera d'insieme}. Nessuno dice ``il \emph{combat system} di questo
gioco l'ha fatto X, che bravo'': si dice ``BioWare ha fatto un bel gioco'',
``Blizzard ha fatto un bel gioco''.

Lo ha replicato anche lato HR. In molte aziende ci sono \textbf{bonus
personali} legati a obiettivi (se il gioco vende X, se la feature è pronta in
tempo). Il relatore li ha dimezzati: metà restano personali, metà diventano
\textbf{bonus di team}, che tutti ricevono se il team raggiunge gli obiettivi.
Un incentivo concreto a collaborare: non solo si dice che il lavoro di gruppo è
importante, lo si dimostra dandogli più valore economico. Aveva proposto il
100\% di bonus di team, ma gli è stato concesso il 50 e 50; ha funzionato
comunque, anche se sarebbe stato ancora più potente. Ovviamente, alla
domanda ``perché mi aiuti?'' si può anche rispondere ``perché voglio il
bonus'': un aiuto dato ``per il motivo sbagliato'' era un compromesso che era
disposto ad accettare.

\subsection{Il coraggio di chiedere aiuto}

Un partecipante osserva che aiutarsi dovrebbe essere la norma. Il modello
diffuso, invece, è ``devi essere forte e indipendente, non chiedere mai una
mano'', e vale in tutte le direzioni (dal maschilismo alla ``donna forte e
indipendente che non ha bisogno di nessuno''): \textbf{non ammettere di avere
bisogno di aiuto è sempre uno dei motivi per cui le aziende vanno male}, che si
tratti di uomini, donne o persone non binarie. Il problema è non avere la
sicurezza di poter alzare la mano.

Il relatore ha cercato di togliere ogni conseguenza dall'aver alzato la mano:
se si \emph{depenalizza} qualcosa, le persone non sono più disincentivate a
farla. Ha detto espressamente che non avrebbe mai valutato nessuno per il
numero di volte in cui chiedeva aiuto o per il numero di feedback negativi
ricevuti, e questo ha funzionato.

\subsection{Perché si dà un feedback}

\emph{Perché avete dato il vostro ultimo feedback? Che cosa volevate ottenere?}
Le risposte della classe:
\begin{itemize}
  \item per \textbf{riconoscere il valore} di un collega, anche verso il
        manager, perché qualcuno che non aveva visibilità vedesse quanto
        quella persona fosse efficace;
  \item per \textbf{spingere una persona a superarsi};
  \item per aiutare un collega a rendere il prodotto all'altezza di ciò che
        voleva rappresentare;
  \item per \textbf{genuinità}: c'era qualcosa che poteva essere fatto meglio.
\end{itemize}
Il relatore aggiunge che la forma cambia col tempo: all'inizio era ``devo
darti un feedback, hai dieci minuti?''; ora è un messaggio informale su
Discord, perché ha costruito un rapporto di fiducia e sa che non verrà letto
come ``l'HR manager mi scrive, sono nei guai?'', ma come ``Stefano è così, mi
dà feedback, resterà tra noi e lo risolveremo insieme''.

Un altro motivo è \textbf{risolvere incompatibilità caratteriali}: le
conversazioni più difficili, del tipo ``io e te non stiamo funzionando, cosa
possiamo fare diversamente?''. Sono importanti perché, se due responsabili non
vanno d'accordo, il team se ne accorge, comincia a raccontarsi storie (``si
odiano, ci saranno problemi'') e, per difendere il proprio manager, comincia a
giudicare male il lavoro dell'altro gruppo. È la dinamica da fermare:
\textbf{valutare gli altri solo per difendere la propria fazione}.

Dare e ricevere feedback sono \textbf{due muscoli da allenare} (il relatore è
sempre stato più bravo a darlo che a riceverlo), e si impara col tempo. È come
fare colloqui: fare colloqui è una skill a parte, che non indica quanto si è
bravi nel proprio lavoro. Essere scartati non significa non saper fare quel
lavoro: spesso significa aver fatto male il colloquio o che è stato preferito
qualcun altro.

\paragraph{Feedback pubblici e privati.} Uno studente racconta di aver dato un
feedback negativo a un superiore, senza un motivo particolare, per dimostrare
a sé e ai colleghi che non si viene licenziati per questo. Il relatore coglie
l'occasione per precisare che \textbf{i feedback negativi si danno
individualmente}: lui dà moltissimi feedback di cui il resto dell'azienda non
sa nulla. Reagirebbe molto male se, in un meeting, qualcuno alzasse la mano per
dire ``comunque questa cosa non l'hai fatta bene'': quello non è feedback, è
discussione, accettabile solo se è un brainstorming e se l'opinione è stata
chiesta. Alcune aziende hanno piattaforme pubbliche di feedback, ma spesso
servono solo ad applaudire; lo stesso accade nelle aziende che
\emph{obbligano} a lasciare un certo numero di feedback a settimana, che
diventano ``lui è bravo, lei è competente, anche questa settimana l'abbiamo
fatto''. Bisognerebbe chiedere a chi le adotta cosa pensa di ottenere (forse
trasparenza), ma difficilmente una persona ``attaccata'' pubblicamente reagisce
bene.

% ---------------------------------------------------------------------
\section{Among Us: la caccia al traditore}
\label{sec:z1-among-us}

Prima della pausa pranzo il relatore mostra tre film di fantascienza di
qualche decennio fa e chiede che cosa li accomuni\footnote{I titoli non sono
leggibili nella trascrizione: viene citato \emph{Event Horizon} da uno
studente, e il film ``in mezzo'' è indicato come ``l'originale'',
verosimilmente \emph{La cosa} (\emph{The Thing}).}. In ognuno c'è un gruppo di
persone con un obiettivo e la presenza, reale o percepita, di qualcuno che
rema contro: \textbf{l'energia del gruppo viene spesa per trovare il
responsabile}.

\paragraph{L'esperimento degli alieni.} Nei master in presenza il relatore
propone un esercizio. Distribuisce segretamente cartellini plastificati con
scritto ``essere umano'' o ``alieno''; le carte alieno contengono istruzioni di
comportamento con cui gli alieni, che cercano di fingersi umani, si
tradirebbero (dare sempre ragione a chi ha parlato prima, finire ogni frase
con una domanda, parlare in un certo modo). Lo scopo: in dieci minuti formare
gruppi, e \textbf{vince il gruppo più numeroso senza alieni}; basta un alieno
per perdere. Si formano sempre gruppi piccoli (due, cinque, sei persone),
perché la tendenza è ``mi fido al massimo di voi due''.

Ma vincono i gruppi più grandi, per un motivo semplice: il relatore
\textbf{non distribuisce nessuna carta alieno}. Mostra a tutti entrambe le
carte, ma consegna solo carte ``umano''. Vuole simulare ciò che accade in
ufficio: la tendenza a percepire sempre qualcuno che rema contro (i manager,
l'altro team, quel collega), quando in realtà non c'è nessuno, ``è tutto nella
nostra testa''. Quando chiede le motivazioni, emergono sospetti
\emph{fabbricati}: ``faceva sempre domande'', ``parlava solo all'infinito''
(``no, parla così e basta''). Ci si fabbrica da soli i motivi per non fidarsi
degli altri, ed è comunissimo.

\paragraph{Among Us.} L'autosabotaggio è, per tornare ad \emph{Among Us},
\textbf{concentrarsi sul cercare cadaveri e traditori invece di fare le
task}: non si vince perché non si sono completate le task, mentre si cercava il
colpevole. Succede tantissimo negli uffici, ed è il motivo per cui l'80\% dei
problemi che il relatore risolveva erano di comunicazione. Entrando nei team
sentiva immancabilmente: ``facciamo il possibile, ma questa lotta continua con
l'altro gruppo è estenuante''. ``Perché c'è questo problema?'' ``Non lo so,
c'era già quando sono arrivato.'' Un paradosso: c'è un problema, nessuno sa
dire quale sia, ma c'è sempre stato e quindi lo si accetta.

\paragraph{L'esperimento delle scimmie.} È l'esempio, ``finto o reale'', delle
scimmie in gabbia\footnote{È il noto racconto delle ``cinque scimmie'', spesso
presentato come esperimento scientifico ma probabilmente apocrifo (lo stesso
relatore lo definisce ``finto barra reale''); nella trascrizione il numero
delle scimmie varia.}. In una gabbia con alcune scimmie pende una banana
elettrificata: chi la tocca fa prendere la scossa a tutte. Le scimmie imparano
e, quando una si avvicina alla banana, le altre la picchiano. Poi si
sostituisce una scimmia alla volta con una nuova: la nuova va verso la banana,
le altre la picchiano, e smette. Si continua finché nella gabbia ci sono solo
scimmie nuove, poi si toglie la corrente. Nessuna scimmia tocca la banana:
nessuna ha mai preso la scossa, ma tutte hanno imparato che chi tocca la
banana prende le botte.

Nei team è lo stesso: il motivo per cui si ``odia'' l'altro team si
trasforma, ma un vero motivo non c'è. Ne derivano la sfiducia, il non alzare
mai la mano per chiedere aiuto, il non dire mai ``bel lavoro'' agli altri
perché non sia mai che qualcuno ti veda fare un complimento ``a quelli''.
Risolvere queste dinamiche è lungo e tedioso, e sempre pieno di ``non mi
ricordo'', ``io non c'ero'' e ``abbiamo sempre fatto così'', la cosa più
frustrante per chi cerca di risolverle.

% ---------------------------------------------------------------------
\section{Il principio del default to trust}
\label{sec:z1-default-trust}

\subsection{La definizione}

Il relatore indica questa come la slide più importante della presentazione.
\textbf{Default to trust} è uno dei valori aziendali di Riot Games, che lui ha
tradotto, adattato e ``esportato'':

\begin{quote}
  \emph{Approccio ogni situazione di conflitto partendo dal presupposto che una
  persona abbia agito in un certo modo perché, con le informazioni a sua
  disposizione in quel momento, ha pensato che fosse la soluzione migliore, e
  non perché sia un'incapace o volesse danneggiare il progetto
  intenzionalmente.}
\end{quote}

In altre parole: \textbf{di default, mi fido}.

\subsection{L'esempio del ritardo al daily stand-up}

A Riot il relatore aveva un \emph{daily stand-up} dal vivo alle 15 con il suo
team di venti persone. Una ragazza arrivava regolarmente alle 15:15, per una,
due, tre settimane. Alla quarta ci ha parlato. \emph{Come iniziereste la
conversazione?} Le proposte della classe:
\begin{itemize}
  \item sottolineare l'importanza di esserci dall'inizio, perché potrebbe
        perdersi qualcosa che le interessa, e per rispetto dei colleghi
        presenti;
  \item ``devi arrivare in orario'', spiegando perché è importante;
  \item \textbf{chiedere perché}.
\end{itemize}

Il relatore osserva, anche dalle sue consulenze fuori dai videogiochi, che
\textbf{più si sale con la seniority}, più la risposta è ``devi smettere di
arrivare in ritardo'', spesso con un \emph{or else} finale (``altrimenti
succede quest'altra cosa''). Più si scende con l'età, e quindi con meno
abitudini consolidate, più la prima domanda diventa ``dimmi perché arrivi in
ritardo''. Già spiegare perché è importante è ``miglia avanti'' rispetto al
semplice ``smetti''.

Partire dal perché gli ha permesso di scoprire la causa: il meeting precedente
della ragazza, in un'altra zona del campus a un quarto d'ora di distanza,
finiva sempre tardi, e lei doveva attraversare il campus di corsa. Non aveva
la confidenza per dire all'altro manager, che insisteva perché si restasse
fino alla fine, che doveva andare; sacrificava quindi il meeting del relatore,
sapendo che era il più comprensivo dei due.

Il relatore si è offerto di \textbf{mediare}: ha detto all'altro manager che
o anticipava l'inizio del suo meeting, così da finire comunque in tempo,
oppure lasciava uscire la ragazza, perché quel meeting gli serviva. C'è anche
una questione di \textbf{economia di scala}, spesso sottovalutata: se venti
persone aspettano qualcuno per dieci minuti, si sono sprecati dieci minuti per
venti. Per lo stesso motivo il workshop è iniziato alle 10:05: iniziare alle
10:10 avrebbe significato far aspettare cinquanta persone cinque minuti in
più. A Riot c'era una scrivania vuota con un pupazzo gonfiabile, \textbf{Mr.\
Late}: sommando tutti i minuti di ritardo a meeting, call e conversazioni, in
un anno veniva fuori praticamente \emph{una persona a tempo pieno}. Serviva a
ricordare che anche cinque minuti contano, e che bisogna rispettare gli altri
arrivando in orario o almeno avvisando.

\paragraph{Gli effetti.} La conversazione era avvenuta nel primissimo periodo a
Riot, quando il team non aveva ancora motivo di fidarsi di lui. Risolvendo il
problema ha eliminato anche tutto lo strato di micro-dinamiche spiacevoli che
un manager impara a vedere (due colleghi che si guardano con la risatina
quando lei entra in ritardo). La ragazza ha cominciato a entrare tranquilla,
con una certa leggerezza: non le faceva piacere arrivare in ritardo (``ci sono
anche persone che lo fanno apposta, e quello è un problema da affrontare
diversamente''). Soprattutto, con i colleghi non parlava di lui come del capo
che minacciava, ma come di chi le aveva risolto un problema: era diventato
un \textbf{alleato}. La settimana dopo altre tre persone sono andate da lui
con un problema da risolvere. Ripetuto cinque, sei, sette, dodici volte, in un
paio di mesi poteva dire quasi qualsiasi cosa nei meeting sapendo che
richieste e critiche sarebbero state percepite come ``me lo chiede perché vuole
darmi una mano'', non per fare bella figura o per presentare i successi del
team ai suoi superiori.

\subsection{Condividere i successi, schermare i fallimenti}

Un altro elemento fondamentale (ripreso nella seconda giornata): il manager
\textbf{condivide i successi e protegge il team dai fallimenti}. ``Se va male
qualcosa è colpa mia; se va bene qualcosa, voglio che ne parliate voi.'' Il
lavoro del manager non dovrebbe essere valutato su ciò che presenta del lavoro
del team: tutti sanno comunque chi è il manager. Tanto vale lasciar presentare
il team, che potrà dire con orgoglio di aver presentato a Mike
Morhaime\footnote{Cofondatore ed ex presidente di Blizzard; nella
trascrizione ``Mike Morayme''.} la propria campagna europea per
\emph{StarCraft~II}. Ci sono manager che dicono ``come mi prendo le colpe, mi
prendo i meriti'', ma i meriti li hanno già: è inutile ostentarli. Con i
fallimenti invece il manager è la \textbf{rete di sicurezza} del team. Il
relatore ne ha fatto una questione di principio: \textbf{non scaricare mai la
colpa su qualcun altro}.

\paragraph{Colpevole o causa.} Un esempio che il relatore attribuisce a Pirate
Software, creatore di contenuti ed ex dipendente Blizzard\footnote{Il relatore
aggiunge che si tratterebbe del figlio dello sviluppatore a cui si sarebbe
ispirato il personaggio del giocatore ``senza vita'' nell'episodio di
\emph{South Park} dedicato a \emph{World of Warcraft}. L'informazione è
riportata come detta in lezione e non è verificata.}. Contattando il servizio
clienti Blizzard per motivi personali, aveva scoperto che i colleghi gli
avevano scritto nei ticket informazioni confidenziali che non avrebbero dovuto
condividere in quel modo, nemmeno con un collega. Non ha segnalato le persone
da punire: ha detto ``c'è un problema, evidentemente le policy non sono
chiare''. Quando gli hanno chiesto i nomi, ha rifiutato: ``se vi dico chi è,
punite quelli, e gli altri cinquecento continueranno a farlo; non avremo
cambiato la causa alla fonte''.

\begin{quote}
  \emph{È inutile trovare il colpevole. Voglio trovare la causa. Perché il
  colpevole cambia.}
\end{quote}

Il colpevole a volte è vittima della situazione: per questo la definizione
dice \emph{``con le informazioni a sua disposizione in quel momento''}. Se il
relatore avesse rimproverato la ragazza in ritardo, non solo non avrebbe
risolto il problema, ma avrebbe creato un'\textbf{ingiustizia}: lei, che
attraversava il campus di corsa per una colpa non sua e di cui non poteva
parlare, sarebbe stata anche rimproverata.

Nessuno fa un brutto lavoro apposta: nessuno è contento quando sbaglia o
perde. Perché al lavoro diamo per scontato che gli altri sbaglino
intenzionalmente? Se qualcuno ha fatto una cosa male, è perché in quel momento
credeva fosse la scelta migliore. La domanda giusta non è ``perché hai fatto
questa cosa che è stata un fallimento?'', ma:

\begin{quote}
  \emph{Che cosa ti ha fatto pensare che questa fosse la scelta giusta?}
\end{quote}

Magari erano limiti di budget, di tempo, di scopo: si agisce su quelli, non
sulla persona. Se si cambia la persona ma le si danno gli stessi fattori, è
molto probabile che arrivi alle stesse conclusioni, e ci si troverà a dover
``disciplinare'' anche lei.

\paragraph{La reputazione costruita sui rumor.} Durante il giro di
presentazione in un ufficio (non in nessuna delle aziende citate), al relatore
è stato detto: ``questo è il dipartimento di design; se ti serve qualcosa
chiedi a lei, non a lui, perché lui consegna sempre in ritardo''. Oltre a
trovarlo poco professionale, una settimana dopo ha chiesto: ``cosa ti ha fatto
questa persona?'' ``A me niente, però ho sentito che\dots''. Si torna alle
scimmie: magari quella persona non consegna in ritardo da mesi, ma ormai lo
dicono tutti. È un circolo vizioso da fermare subito.

\paragraph{Niente escalation inutili.} Lo stesso vale per il feedback: quando
qualcuno va dal relatore dicendo ``ho un problema con questa persona'', lui
chiede: \textbf{``perché ne stai parlando a me e non a lei?''}. ``Perché sei il
mio capo'' non basta. Passare da lui significa dare per scontato che l'altra
persona non abbia interesse a risolvere, mentre magari basta parlarle. È la
strada più facile (``lo dico a Stefano, ci pensa Stefano''), e proprio per
questo bisogna mettere dei paletti: il manager interviene quando serve, ma le
cose che si possono gestire da soli vanno gestite da soli. (Quando invece il
confronto diretto fallisce, l'escalation è legittima: si veda la
sezione~\ref{sec:z1-qa-finale}.)

% ---------------------------------------------------------------------
\section{Traiettorie di carriera e valutazione delle performance}
\label{sec:z1-traiettorie}

\subsection{Quale traiettoria vorreste?}

Il relatore mostra un grafico con quattro curve colorate, contrassegnate
anche da lettere (A, B, C, D) per eventuali persone daltoniche (``a me è
capitato che qualcuno scoprisse di esserlo durante una lezione''). Le curve
rappresentano la performance lavorativa nel tempo. Due domande: \emph{quale
vorreste essere? E quale pensate realisticamente di essere?} Un partecipante
aspira a B e pensa di essere C; quasi tutti scelgono B.

Perché non A, che sale più rapidamente? Un partecipante risponde che una
crescita così repentina può essere più dannosa che utile, come una ``rush'' di
studio molto faticosa che poi non si riesce a sostenere. Il relatore propone
una lettura suggerita da un amico psicologo, che spiega la forma del grafico:
di B, C e D \textbf{vediamo dove finiscono}; di A no. Potrebbe continuare a
salire oltre la scala del grafico, o crollare sotto D. Non saperlo ci spinge a
preferire B. (A chi obietta che anche le altre proseguono oltre il bordo, il
relatore precisa: il bordo destro è ``oggi'', e di A non sappiamo come sia
arrivata fin lì.)

\subsection{La carriera non è una linea retta}

Tutti immaginano la carriera come una linea: si parte da X e, attraverso
stage e posizioni, si arriva a Y. In realtà carriera e performance avranno un
andamento irregolare, con alti e bassi. Come non si accetta che un team possa
temporaneamente \emph{underperformare}, tanto meno lo si accetta da se stessi.
Bisogna accettare che ci saranno momenti bassi.

Per questo il relatore \textbf{detesta le valutazioni annuali delle
performance}: prendono un punto casuale della traiettoria, che se si è
fortunati è il picco più alto e se si è sfortunati il più basso, e lo usano
per valutare un intero anno. Insegna ai manager a fare più punti di controllo
e a guardare altre cose:
\begin{enumerate}
  \item \textbf{quanto tempo passa tra un picco e l'altro}, cioè quanto
        \emph{cooldown} serve a una persona. È il motivo per cui, dopo la
        Games Week, tutti staccavano una o due settimane. Il relatore stesso,
        che si impegna moltissimo e tende al burnout, non potrebbe rifare lo
        stesso workshop il lunedì e il martedì successivi con la stessa
        qualità;
  \item \textbf{la differenza tra la miglior performance e la peggiore}: che
        cosa ha portato a lavorare così bene, e così male? Con chi si
        lavorava, su cosa ci si concentrava, quanto tempo e budget c'erano?
        L'obiettivo è capire come \emph{replicare} le condizioni che fanno
        lavorare bene, invece del solito ``hai lavorato male, è colpa tua''.
        Il relatore ha avuto molti manager per cui quando lavorava male era
        colpa sua e quando lavorava bene ``che bravo'', ma sapeva che raramente
        era l'unico fattore della propria performance. È la fallacia del ``se
        sai fare 100, perché non fai sempre 100?'', come nella scena di
        \emph{Schindler's List} sull'operaio che produce una vite in pochi
        secondi: non sempre un ritmo è sostenibile;
  \item \textbf{il fallimento di oggi confrontato con quello di sei mesi fa}.
        Il relatore non è contento della presentazione di oggi (l'avrebbe
        migliorata per altri due anni), ma rispetto a quelle di cinque anni fa
        vede un grande miglioramento di chiarezza: il suo peggior fallimento di
        oggi equivale al suo primo successo di inizio carriera.
\end{enumerate}
Quello che conta è la \textbf{traiettoria} dall'inizio dell'anno a oggi, non
i singoli punti. E se si riesce a capire insieme quando e perché si ha
successo e quando e perché si fallisce, si può lavorare per replicare solo le
situazioni che fanno lavorare bene.

\paragraph{Lo stesso vale per i giochi.} È il motivo principale per cui i
giochi delle grandi case vanno bene o male. Quando si riescono a replicare le
condizioni che hanno prodotto \emph{Warcraft~II}, esce \emph{Warcraft~III};
quando non ci si riesce, esce ``il nuovo \emph{Dragon Age}'': la gente non è
arrabbiata perché il gioco non è bello, ma perché non è come gli altri
\emph{Dragon Age}; non ne replica il \emph{feeling}, anche se potrebbe
essere un gioco favoloso. L'eccezione positiva è il passaggio da \emph{Dune} a
\emph{Dune~II}: il primo era di fatto un'avventura grafica, il secondo ha
generato il genere degli strategici in tempo reale. Gli sviluppatori avevano
capito che il successo non stava nel tipo di gioco, ma nell'ambientazione che
attirava le persone: un bel gioco in quell'ambientazione sarebbe piaciuto
comunque. Il relatore comprò \emph{Dune~II} aspettandosi un seguito
dell'avventura e, pur sorpreso, lo finì.

\subsection{Lying, hiding, faking}

Quando si è valutati su un singolo punto nel tempo, e da quel punto dipendono
bonus, stipendio e il rischio di licenziamento, si comincia a fare
\textbf{lying, hiding e faking}:
\begin{description}
  \item[Lying] dire che è sempre andato tutto bene, che si finirà sempre in
        orario, che non c'è nessun problema;
  \item[Hiding] nascondere: ``non era un mio task'', ``non c'ero'', ``non ho
        visto niente'';
  \item[Faking] fingere, anche in cose apparentemente innocue come ``sì, sì,
        ho capito tutto'' quando si ha solo paura di dire ``me lo rispieghi?
        Non ho capito niente''.
\end{description}

\subsection{``Non ho capito'' e la conferma delle istruzioni}

Dire ``non ho capito, ripetimelo'' è ancora visto come una figura da sciocco.
Il relatore non si arrabbierà mai se qualcuno gli chiede di rispiegare
un'intera slide: le cose ovvie per lui possono non esserlo per altri. Più si
conosce qualcosa, \textbf{più è difficile spiegarla}, perché si danno per
scontate molte cose. Ne parlava proprio con Diego, che gli diceva che fare
l'insegnante diventa sempre più difficile man mano che si diventa esperti. È
come spiegare \emph{EVE Online}, uno dei giochi più complessi in circolazione:
chi lo conosce bene dà per scontato che ``ovviamente non puoi fare questa
cosa'', mentre per un nuovo giocatore non è affatto ovvio. A scuola chi dice
``non ho capito'' viene visto come lo ``scemo''; nessuno si chiede ``non è che
l'ho spiegato male?''.

Il racconto più ``traumatico'' dei ragazzi che fanno stage nelle grandi
aziende è sempre lo stesso: ``Cinzia, devi fare questa cosa.'' Cinzia la fa.
``Non te l'avevo chiesta così.'' Il problema è che nessuno ha
\textbf{confermato che Cinzia avesse capito}. Il consiglio, valido sempre:
\begin{itemize}
  \item quando si \emph{ricevono} istruzioni, ripeterle: ``ho capito bene?
        Vuoi che faccia questa cosa in questo modo?'';
  \item quando si \emph{danno} istruzioni, chiedere a chi le riceve di
        ripeterle, non perché lo si ritenga incapace, ma per capire che cosa
        ha capito dalla propria spiegazione (``quindi vuoi un quadrato?''
        ``No, ho detto un triangolo\dots'' ``Però mi hai detto una forma con
        quattro angoli.'' ``Hai ragione: voglio un quadrato'').
\end{itemize}
Frasi come ``non te lo ripeto due volte'' o ``stammi dietro'' portano solo a
fare le cose come si pensa che vadano fatte e a sentirsi dire a fine progetto
``non la volevo così'': che si tratti di una feature, di un design o di un
modello 3D. (Un partecipante osserva che lo stesso accade insegnando le lingue:
chi insegna dà per scontato, ad esempio, che si sappia dell'esistenza del
genere neutro in tedesco.)

\paragraph{Come convincere il team a chiedere?} Il relatore mette sempre come
premessa che \textbf{l'obiettivo non è trovare la persona più intelligente
nella stanza, ma fare un buon lavoro}. Che Gianni capisca alla prima e un altro
alla quinta non gli interessa, se alla fine entrambi lavorano bene: ripetere
cinque volte è comunque meno lavoro che far rifare le cose sette volte. È un
investimento a lungo termine, e non dissimula alcun fastidio: davvero non gli
dà fastidio. Ma ripetere le stesse parole non aiuta: ``non ho sentito'' si
risolve ripetendo; ``non ho capito'' si risolve cambiando spiegazione. Meglio
ancora: ``dimmi che cosa hai capito tu, e vediamo cosa non torna''. Se
intravedesse nel team uno sguardo del tipo ``è la quinta volta che lo
chiede'', interverrebbe: spesso chi giudica è proprio chi dice ``ho capito'' e
poi sbaglia. Preferisce di gran lunga chi chiede cinque volte.

\subsection{Il feedback più utile che abbiate ricevuto}

\emph{Qual è il feedback più utile che abbiate mai ricevuto?} Il relatore ne ha
almeno dieci che gli hanno cambiato il modo di lavorare. Le risposte:
\begin{itemize}
  \item un semplice \textbf{``grazie''} per un lavoro fatto, più utile di un
        ``questa cosa fa schifo'';
  \item un feedback sulla propria \textbf{impulsività} nel rispondere: fermarsi
        un attimo e capire di cosa si sta parlando prima di intervenire;
  \item \textbf{essere se stessi}, anche al lavoro, e non scusarsi
        continuamente per ogni cosa.
\end{itemize}
I feedback che toccano di più sono quelli \textbf{veri}. Il ``grazie''
significa ``ho visto quello che hai fatto e l'ho apprezzato''. Quelli negativi
che si ricordano sono quelli che hanno fatto realizzare ``questa cosa che
pensavo di saper fare, non la so fare'', oppure ``la so fare, ma questa volta
non mi è venuta'' (non è che se una cosa una volta non riesce non la si sa più
fare: può capitare di sbagliare una porta). Arrivano sempre da un
\textbf{punto cieco}: qualcosa che non si sapeva fosse percepito dagli altri.
Conta più ciò che le persone ti dicono di ciò che pensi che pensino di te; e
ciò che il relatore pensa che gli altri pensino di lui è sempre terribile,
mentre il loro giudizio si rivela molto più positivo del previsto.

\paragraph{Il denominatore comune siete voi.} Si potrebbe pensare che questi
concetti non siano applicabili da chi arriva in un'azienda senza poterne
cambiare la cultura. Ma c'è sempre un denominatore comune a cui applicarli:
\emph{se stessi}. La cultura aziendale può non spingere a dare feedback, può
vedere male chi chiede aiuto; ma il relatore ha visto che quando lui lo fa, le
persone intorno tendono a farlo, e un gruppo che lo fa lavora meglio. Perfino
i CEO più ostici a un certo punto gli chiedevano: ``perché il tuo team ora
performa meglio degli altri?''. La differenza è il default to trust: una
pratica che col tempo fa operare il team \emph{come un team} e non come un
insieme di individui. Un videogioco è un progetto corale: c'è chi fa giochi da
solo, ma non si fa da soli \emph{Red Dead Redemption~2}; e se il gruppo non è
affiatato e allineato non esce \emph{Red Dead Redemption~2}, esce
\emph{Anthem} (``oggi la mia fissazione è \emph{Anthem}, domani sarà
un'altra'').

% ---------------------------------------------------------------------
\section{Come dare un feedback}
\label{sec:z1-dare}

\subsection{Concentrarsi sull'azione e sui suoi effetti}

Il primo principio: \textbf{il feedback non riguarda mai la persona}. Non si
dice ``tu hai fatto questa cosa che non va bene'', ma ``questa cosa non va
bene: come facciamo a non farla più?''. Che sia un comportamento, un'azione,
una decisione o una strategia, ci si concentra su quello, e si
\textbf{esplicita l'effetto} che ha. Nell'esempio del ritardo l'effetto era
``quindici minuti di attesa moltiplicati per venti persone''. In un progetto
può essere una catena:
\begin{quote}
  se non rispetti le deadline, l'animazione che viene dopo resta ferma; se
  loro sono in ritardo, la \emph{vertical slice} è in ritardo; se la vertical
  slice non è pronta per la Gamescom, non andiamo alla Gamescom e non troviamo
  un publisher; se non troviamo un publisher, chiudiamo.
\end{quote}
Più la conseguenza è grave, più dovrebbe sembrare piccolo il gesto di dire
``datemi una mano'': è il \emph{male minore}. Poi, a cose fatte, ci si dà
feedback: cosa ha funzionato, cosa no, cosa fare di più, di meno, cosa smettere
di fare.

\subsection{Il pericolo degli ambienti solo positivi}

Gli ambienti dove ci si fa sempre grandi complimenti e nessuno parla dei
problemi creano difficoltà vere alle aziende. Il problema non è mai la
mancanza di complimenti, ma un \textbf{eccesso di ottimismo}: si comincia a
pensare che le cose andranno come si spera, perché tutti intorno danno solo
feedback positivi.

A Blizzard i feedback determinavano gli \textbf{aumenti annuali}: ognuno
doveva indicare cinque persone che avrebbero dato feedback su di lui, e
ovviamente tutti indicavano gli amici, o colleghi di progetti andati bene. Un
sistema nato con le migliori intenzioni diventava uno \emph{scambio di voti}
(``se mi metti follow ti metto follow''), senza alcun miglioramento.

Al contrario, \textbf{le conversazioni più brutte sono state le più utili}:
affrontare i licenziamenti a Riot, o la cancellazione di \emph{Project Titan}
a Blizzard, il progetto che doveva essere il nuovo MMO dopo \emph{World of
Warcraft} e da cui poi è nato \emph{Overwatch}\footnote{Nella trascrizione si
legge anche ``doveva essere Warcraft 2'', probabilmente un lapsus o un errore
di trascrizione.}. Andando a vedere perché lo si cancellava, c'erano solo
motivazioni razionali e pragmatiche, di business, chiare a tutti, non un ``Mike
Morhaime ci odia''. Sono anche le conversazioni in cui ha lavorato su se
stesso: è molto \textbf{avverso ai conflitti}, e questo a volte gli ha
procurato conseguenze molto più gravi di quanto sarebbe costato un piccolo
disagio. Dieci anni fa faceva fatica ad accettare questo difetto; ha dovuto
subire le conseguenze di conflitti evitati e poi cresciuti.

\subsection{Che cosa fare}

Le linee guida della slide, che il relatore legge rapidamente perché saranno
disponibili nel materiale:
\begin{description}
  \item[Attenzione al linguaggio del corpo] dare un feedback a braccia
        incrociate o guardando altrove è diverso dal guardare l'altro negli
        occhi.
  \item[Sii specifico] con esempi tangibili.
  \item[Esplicita gli effetti] come visto sopra.
  \item[Discuti il prima possibile] non aspettare un mese per parlare di
        qualcosa successo ieri, altrimenti nessuno ricorderà più bene cosa sia
        accaduto.
  \item[Concentrati sull'azione e sul problema, non sulla persona] il nemico è
        l'azione, non la persona che si ha davanti. Non bisogna far cambiare
        l'altro: bisogna evitare che il problema si ripresenti.
\end{description}

\paragraph{Gli insulti sono rivolti al ruolo.} Un esempio solo apparentemente
scollegato. Il relatore ha lavorato su \emph{League of Legends}, ``non
esattamente noto per avere una community da orsetti del cuore'', in varie
\emph{task force} sulla tossicità. A chi smetteva di giocare veniva inviato un
sondaggio sul perché; il 33\% indicava la tossicità della community. Dopo varie
campagne di sensibilizzazione e cambiamenti si è scesi al 30\%: sembra solo
un 3\% in meno, ma valeva circa \textbf{120 milioni di revenue in più
all'anno}, perché quei giocatori restavano esposti a skin e contenuti e
continuavano a comprare. Anche qui, piccoli guadagni che si sommano.

Lavorando con degli psicologi, aveva notato che gli insulti \textbf{non sono
quasi mai rivolti alla persona, ma al ruolo}: ``il mio jungler fa schifo'',
``il mio DPS non è capace'', ``il nostro portiere è una pippa''. Chiamare
qualcuno per ruolo lo \emph{depersonalizza} e rende più facile insultarlo. Lo
stesso nell'esempio del giro dell'ufficio: non ``non chiedere a Giulio'', ma
``non chiedere a quel grafico, consegna sempre in ritardo''. Concentrarsi
sull'azione, senza attaccare la persona, fa sì che il feedback non sia (e non
venga percepito come) un attacco personale.

\paragraph{Considera il carattere.} Qui tornano utili i test di personalità.
Servono più agli studenti che al relatore (che sa già che usciranno 53 profili
diversi, senza un pattern): bisogna conoscere il proprio profilo e
\textbf{condividerlo con chi lavora con noi}, chiedendo anche agli altri di
fare lo stesso. In ufficio il relatore ha usato una versione più avanzata, a
pagamento, con un report di circa trenta pagine: il \textbf{profilo
DISC}\footnote{Nella trascrizione ``disk''. Si tratta del modello DISC
(Dominance, Influence, Steadiness, Conscientiousness).}. Ha mappato tutta
l'azienda e ha fatto stampare il profilo di ciascuno sulla sua scrivania, così
che tutti sapessero con chi stavano parlando. Capire come ragionano gli altri
(``è introverso, non gliene parlo davanti a tutti'') è una delle cose che gli
sono tornate più utili: \emph{adattare ciò che si vuole dire a chi si ha
davanti}.

\paragraph{Assicurati di non disturbare.} Sotto il profilo di ognuno c'era un
piccolo \textbf{semaforo}: verde, ``parlami pure, interrompimi''; giallo,
``interrompimi solo se è urgente''; rosso, ``mi sto concentrando, non mi
parlare''.

\paragraph{Lascia finire il lavoro.} Durante gli eventi dal vivo il relatore si
era reso conto di quanto fosse grave dare feedback mentre qualcuno sta ancora
lavorando. Se una cosa sta andando male, dirlo in quel momento non aiuta: tra
la pressione di finire e quella del feedback, peggiorano entrambe le cose, e
anche il rapporto. Meglio lasciarla finire, anche male, e parlarne dopo.

\subsection{Che cosa non fare}

\begin{description}
  \item[Non basarti su supposizioni] ``mi hanno detto che ieri in un meeting
        hai fatto\dots'': o c'eri e l'hai visto, o non esiste.
  \item[Non giudicare la persona.]
  \item[Non fornire feedback vago.]
  \item[Non accettare e non fornire feedback di terze parti] Se qualcuno vi
        dicesse ``mi hanno detto che hai fatto un brutto lavoro'', il primo
        pensiero non sarebbe ``cosa ho fatto male?'', ma \emph{``chi te l'ha
        detto? Chi è la vipera che ho in seno? Chi è l'impostore?''}. Siamo
        fatti così. Ogni conversazione che inizia con ``mi hanno detto che'' è
        inutile: l'altro non ascolterà, penserà solo a chi dare la caccia, e
        spesso penserà alla persona sbagliata, creando antipatie tra persone
        che non c'entravano nulla. Quando al relatore dicono ``mi hanno detto
        che lui ha fatto\dots'', risponde ``vai a parlarne con lui''; se gli
        dicono ``mi hanno detto che tu hai fatto\dots'', risponde ``mandami la
        persona che te l'ha detto, altrimenti non ne parliamo''. Non per
        cattiveria: se chiedesse ``che cosa ho detto di preciso a Marco?'' e la
        risposta fosse ``non lo so, non c'ero'', di cosa si starebbe parlando?
        \textbf{Parlate con la fonte}, altrimenti fate come se quel feedback non
        esistesse.
  \item[Non dare consigli non richiesti] come visto, non portare la soluzione:
        la soluzione deve venire da chi riceve il feedback.
  \item[``Non mi piace'' non è un feedback] ``Penso che questo design non
        funzioni perché, visto che il nostro obiettivo è\dots'' ha argomenti;
        ``non mi piace il combat system'' è un'opinione. Anche ``io l'avrei
        fatto diversamente'' non è un feedback.
  \item[Accettare non vuol dire fare come proponi] deve essere chiaro a chi dà
        il feedback: l'aspettativa è essere ascoltati e presi in
        considerazione, ma se la responsabilità è dell'altro, deciderà
        l'altro.
  \item[Non offenderti] se il feedback non viene seguito. Restarci male
        significa pensare di saperne più dell'altro, cioè \emph{venire meno al
        default to trust}. ``Decidi tu'' significa ``mi fido di te, perché sei
        lo specialista, e mi fido che mi ascolterai e farai la cosa migliore
        per tutti'', che potrebbe anche non essere ciò che ho detto io. Se
        voglio che l'altro faccia sempre come dico io, mi sto fidando solo di
        me stesso.
\end{description}

\subsection{Il primo feedback e il debriefing}

La cosa più difficile è dare il \textbf{primo} feedback: è un muscolo da
allenare. Chi non l'ha mai fatto (soprattutto con questa struttura) dovrebbe
provarci, anche su minuzie (``abbiamo concordato di usare WhatsApp nel gruppo:
non potremmo usare Telegram?''), anche solo per vedere come va. Altre
indicazioni:
\begin{itemize}
  \item feedback \textbf{solo uno a uno}, mai davanti ad altri, soprattutto se
        negativo;
  \item \textbf{prepararsi} prima;
  \item fare il \textbf{debriefing del feedback}: dopo, fermarsi a pensare a
        com'è andato. Il relatore, quando media un feedback tra due persone,
        si sente spesso rispondere ``benissimo'' da una e ``malissimo''
        dall'altra: la percezione è molto diversa. Non si fa problemi a tornare
        dalla stessa persona: ``ieri ti ho dato un feedback: com'è andata? Sei
        soddisfatto? Come poteva andare meglio?''. Con le persone con cui si ha
        meno a che fare, questo aiuta molto la volta successiva.
\end{itemize}

% ---------------------------------------------------------------------
\section{Come ricevere un feedback}
\label{sec:z1-ricevere}

Per il relatore è la parte più difficile. Si riassume in quattro passi:
\begin{enumerate}
  \item \textbf{Ascolta}: si parte sempre ascoltando.
  \item \textbf{Chiedi dettagli}, se serve, ma per capire meglio, non per
        difenderti.
  \item \textbf{Accetta, a prescindere}: non si sta decidendo chi ha ragione;
        si sta ascoltando un'opinione. Non c'è nulla da cui difendersi.
  \item \textbf{Agisci}: fatti un piano. Può essere banale (``mi hanno detto
        che interrompo sempre: cerco di non interrompere''), oppure ``devo
        migliorare questa cosa: come?''. Pensa anche a \emph{chi} potrebbe
        aiutarti e a \emph{che supporto} ti serve: il piano non deve essere
        per forza ``devo risolvere tutto da solo''.
\end{enumerate}
Una domanda che molti sottovalutano: \emph{``sono junior in un'azienda: sarò
mica il primo ad avere questo problema?''}. Probabilmente qualcun altro l'ha
già avuto: basta chiedere ``tu come hai fatto?''.

\subsection{Il feedback positivo}

Tutto questo vale per i feedback costruttivi negativi, a cui si pensa di
default; ma esistono anche i \textbf{feedback positivi}, e vanno usati. Quando
si riceve un ``grazie'' o un ``bravo'', la domanda è: \emph{come faccio a
ricreare le stesse condizioni?}

Il relatore ha lavorato anche in aziende dove il feedback era solo negativo.
Il rischio? Si immagini l'impiegato ideale: vuole sempre migliorare, è sempre
pronto a dire ``magari sbaglio io''. Se nessuno gli dice cosa sta facendo
\emph{bene}, spinto dal desiderio di migliorare rischia di cambiare cose che
erano già buone, e di farle peggio. Il rinforzo positivo \textbf{non è solo una
gentilezza}: è un atto pratico, che evita che le persone peggiorino ciò che
sanno già fare bene.

\subsection{Il feedback accidentale}

Esiste infine il \textbf{feedback accidentale}, molto frequente negli uffici:
\begin{itemize}
  \item il ``chiedi a lei, non a quel grafico'' del giro dell'ufficio è un
        feedback negativo su una persona, dato accidentalmente, né richiesto
        né voluto;
  \item a Riot li chiamavano \textbf{drive-by}, come nelle scene dei film in
        cui passa un'auto, si abbassa il finestrino, si spara e si riparte. In
        molte aziende con una ``cultura del feedback'' in cui tutti devono
        dire la loro, si riceve una valanga di commenti del tipo ``comunque il
        sito a me non piace, ciao''. Non si sa perché, non si sa cosa fare di
        diverso, non si sa che rischi ci siano: \emph{a che cosa è servita
        questa conversazione?} A niente.
\end{itemize}

% ---------------------------------------------------------------------
\section{Quando, quanto e con chi}
\label{sec:z1-quando}

\subsection{Frequentemente e recentemente}

\emph{Quando avete ricevuto il vostro ultimo feedback?} Per gli studenti,
qualche giorno fa (nel corso) o qualche mese fa (al lavoro, dopo una
consegna). Come aveva osservato Silvia, il feedback andrebbe dato
spesso.

Il relatore usa l'esempio di una nave da crociera (``la Costa Concordia'') che
deve andare da A a B. La rotta ideale (azzurra) è quasi impossibile: nessun
progetto va esattamente come previsto. C'è la rotta ``panoramica'' (verde), che
gira intorno all'isola d'Elba; ci sono poi rotte (arancione e rosa) che a un
certo punto si allontanano molto e poi correggono bruscamente: qualcuno si è
accorto che si stava sbagliando, ``ma non è quello che avevo chiesto'',
``non è dove dovevamo andare''. Sono i famosi \textbf{sanity check} o
\emph{checkpoint}.

L'ideale è avere \textbf{tanti piccoli check per fare tante micro-correzioni}:
correggere la rotta di un paio di gradi è molto più rapido che farlo dopo
essersi allontanati. Bastano trenta secondi di ``ehi, vado bene?''. Spesso non
lo si fa per paura di sembrare insicuri o poco autonomi, ma per il relatore è
l'opposto: quando qualcuno gli chiede ``vado bene così?'' e lui risponde
``sì'', si alleggerisce anche del proprio carico mentale. Un capo che assegna
un compito e lo vede solo alla fine vive sempre col dubbio (``avranno capito?
Lo staranno facendo giusto? Arriverà in tempo?''); con un check di dieci
secondi al giorno (``tutto ok, Stefano'') è tranquillo, e la persona va avanti
sapendo che quanto fatto finora va bene. Non serve fare ogni giorno il giro del
dipartimento: anche ogni uno o due giorni basta a evitare la rotta verde.

Per il feedback è lo stesso: \textbf{va dato subito}. Dire oggi ``ieri hai
fatto questa cosa'' è molto diverso da dire ``negli ultimi sei mesi hai fatto
questa cosa''. Il relatore si è sentito dire più volte, alla valutazione di
dicembre, ``da maggio hai fatto questa cosa che non andava''. La sua risposta:
``e aspetti dicembre per dirmelo? Mi hai lasciato sbagliare per sei mesi?''. È
una mancanza di rispetto, e la conversazione non può più essere costruttiva.
Per questo non ama le valutazioni annuali: i manager crescono abituandosi a
prendere appunti delle cose negative per tirarle fuori alla valutazione
(``a marzo hai fatto questa cosa'' ``sì, però poi ho smesso'' ``sì, però a
marzo l'hai fatta'').

In sintesi:
\begin{description}
  \item[Frequentemente] riguarda l'intervallo tra un check e l'altro.
  \item[Recentemente] significa non lasciar passare un anno prima di dare il
        feedback, altrimenti non serve più a niente.
\end{description}

\subsection{Quanto dura un feedback}

Un feedback dura \textbf{da trenta secondi a quindici minuti}; raramente di
più. Chi non ne ha mai dato uno dovrebbe sforzarsi di darne almeno uno, magari
positivo per cominciare. Ogni feedback dato o ricevuto è un
\textbf{deposito nel proprio ``fondo fiduciario''} (\emph{trust fund}), che
costruisce la reputazione futura. Il relatore si è costruito la reputazione di
chi dà feedback il più possibile oggettivi e analitici: giusti o sbagliati che
siano, costringono a pensare. Non impone mai il suo punto di vista (``si fa
come dico io''), ma chiede ``tu come faresti? Come posso aiutarti?''. Il
feedback, per lui, serve a \emph{far crescere qualcun altro}.

\subsection{I manager peggiori: ``a che numero sto pensando?''}

I peggiori manager che il relatore abbia avuto erano quelli che volevano che si
facesse come volevano loro, anche quando chiedevano ``tu come la faresti?''.
In realtà la domanda era ``\emph{a che numero sto pensando?}'': se dicevi 5,
``no, dovremmo fare 9''; se dicevi 9, ``no, 5''. Lo stress costante di dover
indovinare cosa vuole il capo nasce dalla mancanza di fiducia (``se mi chiedi
come fare, fammelo fare'').

Lo stesso gli succede con i CEO che lo chiamano per ``sistemare la cultura
aziendale'' e poi dicono ``non sono sicuro, io farei diversamente''. La sua
risposta: ``lo so che faresti diversamente, per questo mi hai chiamato''. Se si
chiama un esperto, bisogna lasciargli fare il suo lavoro: non si va dal
meccanico a cambiare le gomme per poi dirgli ``non dovresti girare un po' di
più quel bullone?''. Chiedere un parere a un esperto significa dargli anche il
potere di scegliere. Il problema è il \emph{micromanagement}, o in generale il
non accettare punti di vista diversi dal proprio (magari perché si è il
fondatore e si pensa di essere l'unico a saper fare le cose).

\subsection{Dare feedback a un superiore egocentrico}

\paragraph{Come si dà un feedback negativo a un superiore con un grande ego?}
Il trucco è \textbf{non puntare sull'ego, ma sull'azione e sui suoi effetti}.
Non ``tu fai così e io mi sento così'': anzi, si alimenta il suo ego
rendendolo il \emph{main character} dell'avventura, l'eroe che risolve il
problema (``tu, grande manager, aiutaci a risolvere questo problema che
affligge noi che facciamo l'operatività''). Non è l'ideale (``io cambierei
manager''), ma pochissimi possono scegliersi i colleghi: il relatore stesso ha
quasi sempre ricevuto team preassemblati (``ecco le chiavi, questa è la tua
macchina''). Bisogna giocare con le carte che si hanno in mano.

\paragraph{Un esempio pratico.} Con un manager che fa micromanagement e mette
sempre ``il becco'' nel proprio lavoro, si può dire: ``se ti fidi di me per
questo lavoro, con le tue capacità potresti dedicarti a questa cosa più
strategica e di lungo termine: pensa quanto sarebbe più utile per l'azienda''.
In sostanza: ``guarda cosa potresti raggiungere se mi lasciassi fare il mio
lavoro''. Non come un ``tu fai il tuo e io faccio il mio'', ma come un
``ti sollevo da questa parte del carico, che so che ritieni importante: lascia
che ci pensi io, così puoi dedicarti ad altro'', che magari gli dà anche più
visibilità.

Si gioca su ciò che quella persona pensa di saper fare, e che probabilmente sa
anche fare. Anche Michael Scott di \emph{The Office}, per quanto preso in giro,
è raccontato come un venditore bravissimo, che ha chiuso i contratti più
importanti per l'azienda, e rispetto ai manager che gli succedono è forse
``uno dei meno peggiori''. Tutti, anche i manager peggiori, hanno punti di
forza: \textbf{giocate su quelli}.

% ---------------------------------------------------------------------
\section{Il feedback come regalo}
\label{sec:z1-regalo}

Per chiudere, il relatore risponde alla domanda d'apertura, che aveva lasciato
apposta in sospeso: che cos'è il feedback per lui?

\begin{quote}
  \emph{Il feedback è un regalo che mi viene fatto da qualcuno che so che vuole
  il mio interesse.}
\end{quote}

Ha costruito un team di persone che si fidano di lui e che non hanno motivo di
sabotarlo o di ferirlo: per questo interpreta anche le cose peggiori con il
default to trust. È un regalo, e cerca di prenderlo in considerazione il più
possibile.

Ci è arrivato partendo da se stesso. Oltre a essere avverso ai conflitti, si
definisce ``avverso alla comunicazione umana'': paradossalmente, una persona
introversa e timida finita a fare comunicazione, marketing e ora HR (``succede
spesso''). Prima di ordinare una pizza al telefono ha dieci minuti di
preparazione mentale, e se l'ordine delle domande cambia va in confusione. Per
lui un feedback di dieci minuti sono spesso \textbf{tre ore di preparazione}:
costruisce un \textbf{albero dei dialoghi}, come in un'avventura grafica
personale, con tutte le opzioni. Poi spesso finisce con ``hai fatto questa
cosa, andrebbe migliorata'' ``hai ragione, grazie'', e basta.

Così ha pensato: chissà quanti si fanno gli stessi problemi con lui, e passano
ore a prepararsi per dirgli ``questa cosa l'hai fatta male'', sapendo che
riceveranno comunque quell'occhiata fugace del ``come ti permetti''. Eppure
vengono a dirglielo. Non può far altro che rispondere: ``grazie, hai ragione,
cercherò di fare meglio''. È così che ha metabolizzato il ricevere feedback:
capendo che non c'è nessuna cattiva intenzione e che, se non proprio felice
(``ci ho messo una pietra sopra''), può almeno essere \emph{grato} che ci sia
ancora qualcuno che ha così a cuore il suo interesse da dargli feedback.

\paragraph{Il silenzio fa più paura.} Non bisogna spaventarsi dei feedback
negativi che emergono dai sondaggi aziendali: bisogna spaventarsi del
\textbf{silenzio}, quando nessuno dice più nulla. Significa che nessuno si
fida più che si sappiano risolvere i problemi.

\paragraph{La sindrome dell'impostore.} Il suo motivatore personale, all'inizio,
era soprattutto ``non voglio fallire'': per questo cercava di migliorare, e
migliorare è diventato accettare feedback, perché più feedback riceve, più
migliora. Il feedback è un \textbf{acceleratore}, e lui cerca di incentivare
gli altri a darglielo. Resta però sempre la vocina del ``e se si accorgono che
non sei capace?''. Con la \emph{impostor syndrome} convive da anni e
continuerà a conviverci; se ce l'avete, conviveteci anche voi. Avrete sempre
la sensazione che ``se ne accorgeranno'': forse sì, ma sapete fare anche
altre cose. Non concentratevi solo su ciò che non sapete fare. Gli errori
capitano: la perfezione non è di questo mondo.

\paragraph{Esercitarsi in accademia.} Per fare pratica con il feedback,
l'ambiente migliore è quello dei colleghi di corso: persone nella stessa
situazione, con esperienza simile e lo stesso obiettivo, in un contesto in cui
non si può essere licenziati e l'unico effetto a lungo termine è ``ho parlato
con questa persona, ora le cose vanno meglio''. E bisogna avere il coraggio di
parlare anche con i docenti (``dire a Diego: secondo me questa lezione potevi
farla in questo modo''; ``so che è una persona che lo accetta''), senza farsi
bloccare dalla gerarchia. Nell'azienda giusta (come quella in cui lavora ora) il
relatore va sempre dal CEO a dare feedback, sapendo che non verrà licenziato:
il CEO gli ha detto ``ti ho preso per darmi feedback''.

La frase di chiusura:

\begin{quote}
  \emph{Per superare le vostre difficoltà, tutto quello che vi viene chiesto è
  il coraggio di chiedere aiuto.}
\end{quote}

Chiedere aiuto non è una debolezza: è un segno di consapevolezza di sé
(\emph{self-awareness}) e di forza personale. Se tutti lo facessero, i
progetti andrebbero molto meglio; spesso invece ci si fa frenare dal ``non
voglio fare brutta figura'', ``non voglio passare per quello scemo'', e sono
proprio queste paure a rallentarci e a impedirci di raggiungere gli obiettivi.

% ---------------------------------------------------------------------
\section{Domande e risposte finali}
\label{sec:z1-qa-finale}

L'ultima mezz'ora è dedicata a domande libere.

\paragraph{E se una persona peggiora per uno o due anni di fila?} Se nessuno
gliel'ha mai detto, c'è un grosso problema, non suo ma dell'azienda. Se invece
c'è consapevolezza, il percorso è: ``cosa vuoi fare per risolvere?''; ``ho
provato questo, questo e questo''; ``visto che non ha funzionato, vuoi un mio
consiglio? Proviamo con quello?''; ``che supporto ti serve?''. Se, con tutto il
supporto, il tempo e le soluzioni concordate per mesi, il miglioramento non
arriva, per il relatore la soluzione è: \textbf{``dimmi un prezzo per andare
via''}. Preferisce pagare una buonuscita (\emph{severance package}) che
avviare la strada dei richiami formali sulla performance fino al
licenziamento per giusta causa. A un certo punto andare avanti non è un favore
per nessuno.

\paragraph{È vero che nei videogiochi si viene sfruttati? Meglio un piccolo
studio italiano?} Per il relatore si è quasi più a rischio a Blizzard che in un
piccolo studio di Roma o Torino. Nel piccolo studio si sa dove ci si va a
cacciare: non ci sarà il team building a Santo Domingo, ma ci si concentra su
un gioco alla volta e, se va bene, si fa il successivo. A Blizzard si può
lavorare su \emph{Overwatch}, che vende milioni, e perdere comunque il posto,
in modo totalmente scollegato dal proprio lavoro. Gli stipendi più alti
all'estero compensano anche questo rischio, che spesso non dipende da come va
il gioco (anzi, ``il gioco non ha fatto come volevamo'' è spesso una scusa per
licenziamenti che non avevano motivo).

\paragraph{Cambiare settore a trent'anni, da web developer full remote?} Lo
stipendio indicato dallo studente non è fuori dalla portata di vari studi
italiani, anzi è superabile. Non va esclusa la possibilità di lavorare da
remoto per aziende piccole ma estere. Alcuni amici del relatore, pagati con
stipendi italiani, si sono trasferiti a Bali o in paesi con un costo della vita
più basso; il collega di \emph{PUBG} percepisce uno stipendio olandese vivendo
a Milano, che per quanto cara non costa come i Paesi Bassi.

\paragraph{Come si diventa manager?} Dipende dall'azienda: alcune promuovono i
migliori (processo tendenzialmente sbagliato), altre propongono un percorso di
\emph{people management}. Diventare manager significa \textbf{abbandonare ciò
che si sa fare per imparare a fare altro}. Il relatore ha dovuto accettare che
non sarebbe più stato ``il miglior community manager italiano del mondo''
(``eravamo in pochi, quindi era facile''): dedicandosi alle persone non poteva
più fare il suo lavoro bene come prima. È stata una negoziazione con se
stesso: fare il manager significa non essere più l'esperto di qualcosa, ma
saper gestire e far crescere persone. Significa anche \textbf{cedere il
controllo della propria carriera}: se il team decidesse di non fare più nulla,
il manager ne pagherebbe le conseguenze. Il suo lavoro diventa scegliere le
persone giuste e farle lavorare bene. Lui lo esplicita sempre: sa di affidare
al team parte della propria carriera, e il team sa che lui combatterà per
loro, anche prendendosi le colpe.

\paragraph{L'inglese.} \textbf{Va assolutamente imparato}, anche lavorando in
Italia. Non saperlo preclude circa il 95\% del mercato (tutto quello fuori
dall'Italia) e soprattutto il \emph{networking} con persone che potrebbero dare
occasioni future: l'inglese è spesso l'unica lingua in comune con chi lavora
all'estero ed è oggi la lingua universale dei videogiochi.
\begin{itemize}
  \item \emph{Che livello serve?} Non importa B2 o C2: bisogna saper sostenere
        una conversazione, guardare almeno Netflix con sottotitoli in inglese,
        comunicare senza dodici pause e senza ``come si dice\dots''. Il
        certificato non verrà mai chiesto: verrà chiesto se si sa parlare
        inglese. Il relatore ha sentito persone con un C2 che non sapevano
        parlarlo.
  \item \emph{Come allenarlo senza occasioni di conversazione?} Netflix, Prime
        Video, Now TV, Udemy, Coursera, YouTube, Duolingo, siti di scambio
        linguistico con madrelingua, Discord: non sono mai esistiti così tanti
        mezzi gratuiti. Il relatore ha imparato l'inglese giocando (con una
        pronuncia sbagliata, perché leggeva senza sentire, su cui ha poi
        dovuto lavorare); guardare un film in inglese con i sottotitoli in
        italiano insegna anche la traduzione.
  \item Oltre all'ascolto passivo è fondamentale \textbf{parlare}: in una
        flotta di \emph{EVE Online} con 150 persone quasi tutte madrelingua
        inglese si è esposti a decine di accenti e modi di dire (mille modi di
        dire ``come stai'' oltre a \emph{how are you}); lavorando all'estero si
        impara a distinguere gli accenti (Manchester, Londra) e ad apprezzarli
        anche nei media.
\end{itemize}

\paragraph{Quanto deve sapere un junior?} Dipende da quanto è tecnico il
ruolo. Un game designer junior che non sa qualcosa può essere affiancato da un
senior; un programmatore che non sa programmare, in quel momento, non è utile
all'azienda.

\paragraph{Annunci ``junior con cinque anni di esperienza''.} ``Cercasi
ventenne con trent'anni di esperienza.'' Non è l'annuncio per un junior: è un
errore dell'azienda, che in realtà cerca un \emph{mid}. Riceverà profili
junior, si stupirà che non abbiano cinque anni di esperienza, rifarà l'annuncio
e avrà perso sei mesi. Il consiglio: tenere d'occhio quella pagina, perché
dopo qualche mese l'annuncio verrà corretto.

\paragraph{Una struttura per il post-mortem.} A chi chiede un formato di base,
il relatore anticipa quanto vedrà nella seconda giornata a proposito di
brainstorming: una lavagna \textbf{divisa in quattro aree}, su cui ognuno
attacca tutti i post-it che vuole:
\begin{center}
  \begin{tabular}{@{}p{0.44\textwidth}p{0.44\textwidth}@{}}
    \toprule
    \textbf{Smettere di fare} & \textbf{Iniziare a fare} \\
    ciò che non ha funzionato & ciò che non abbiamo fatto e andava fatto \\
    \midrule
    \textbf{Fare di meno} & \textbf{Fare di più} \\
    ciò che ha fatto perdere tempo rispetto all'impatto ottenuto &
    ciò che ha funzionato ma, fatto con troppa leggerezza, non ha dato i
    frutti sperati \\
    \bottomrule
  \end{tabular}
\end{center}
Con i post-it ognuno ha voce e tutto risulta catalogato; per il report basta
una foto della lavagna (o uno screenshot, se la lavagna è digitale, come Miro).

\paragraph{Come sbloccare la paralisi del feedback verso il management?} Uno
studente ha visto manager dire ``non abbiate paura, la mia porta è sempre
aperta'' senza che cambiasse nulla. C'è una dissonanza tra ciò che dicono e
ciò che fanno: è ``la mamma che dice vieni qua che non ti faccio niente''.
Qualcosa deve aver fatto pensare alle persone che non è sicuro. Prima va
sistemato quello; poi si può provare con i \textbf{feedback anonimi}.
L'obiettivo del relatore resta il feedback aperto e firmato, ma a Stormind ha
iniziato con feedback anonimi verso il management per verificare se
l'anonimato fosse un fattore. Non lo era (evidentemente non c'erano problemi a
parlare con il CEO), e la volta successiva intende abbandonarli. Ma è un buon
\emph{check}: se senza anonimato non arriva nulla e con l'anonimato arrivano
cento feedback, il problema è chiaro.

\paragraph{Quando è giusto rivolgersi a un superiore?} La strada giusta resta
parlare direttamente con l'interessato. Ma se si prova e non si trova una
soluzione, la cosa più ovvia è andare dal proprio manager insieme: ``ci
aiuti?''. Quando il relatore parlava di mediare i dialoghi, era proprio questo:
fare da moderatore, aiutare a sbloccare i punti fermi o, se non si riesce a
scegliere tra A e B, scegliere lui. L'escalation, quando necessaria, va fatta.

\paragraph{Il tempo libero veniva rispettato?} No, e se l'è dovuto imporre
lui. Il lavoro gli piaceva così tanto che continuare a casa non gli sembrava
lavorare. Gli si è ritorto contro con il lancio di \emph{Diablo~III}, fatto in
orario americano (le due o l'una di notte in Europa), con l'ufficio chiuso
perché a nessuno era stato chiesto di restare. Il lancio va malissimo, i
giocatori protestano, e lui, sveglio, risponde sui forum fino alle tre, quattro,
cinque del mattino: ``non posso risolvere il problema, ma resto qui ad
aggiornarvi con quello che riesco a sapere dal team di sviluppo''. Il giorno
dopo trova sulla scrivania un \textbf{richiamo formale dell'HR}: aver lavorato
fuori orario senza autorizzazione, condotta non professionale. Per lui è stato
un interruttore: ``ora me ne vado''. Due settimane dopo faceva colloqui con
Riot. Ancora oggi è contrario a quel richiamo, ma gli ha fatto capire come
ragiona un'azienda: ``hai un contratto 9--18, non mi interessa se l'hai fatto
con le migliori intenzioni''. Eppure per anni ha goduto di un'ottima
reputazione in quelle community come ``quello che è rimasto con noi fino alle
cinque del mattino'', e anche l'azienda ne ha beneficiato.

Ci sono poi aziende in cui, soprattutto per i developer, viene richiesto il
famigerato \textbf{crunch}: lavoro straordinario continuativo, senza pause.
Dipende molto dall'azienda; sta venendo (``grazie al cielo'') deprecato, perché
molte aziende capiscono che non è sano ed è visto male anche dal pubblico, ma
esiste ancora.

\paragraph{E se il tuo capo demotiva qualcuno del tuo team?} Qui ``diventa
personale'': il team del relatore si fida di lui perché sa che in questi casi
``diventa una bestia''. ``Te la puoi prendere con me quanto vuoi, ma il mio
team non lo tocchi.'' Lo gestisce direttamente con l'altro manager, per quanto
brutta possa diventare la conversazione. Lo sforzo più grande è \textbf{non
andarne poi a parlare con il team} (``ragazzi, ho fatto questa cosa''), perché
minerebbe la fiducia verso l'alto: non vuole un rapporto ottimo fino a lui e
pessimo oltre. L'obiettivo è che l'altro manager ``si faccia perdonare''
cambiando il proprio comportamento, e che il team percepisca da solo il
cambiamento.

\paragraph{Chiusura della giornata.} Il relatore chiede un feedback collettivo
in chat (``se domani ci dai \emph{more of the same} siamo contenti, oppure
ditemi cosa limare''), si scusa per i ripetuti crash della connessione e
ringrazia per la proattività nelle domande, che gli servono anche come
\emph{check}: se nessuno chiede nulla su un argomento, significa che può
andare avanti. Racconta di aver passato una settimana a decidere che cosa
mettere nelle slide e di averle poi preparate in poco tempo (rifacendone
alcune fino alle tre di notte). Ricorda infine di usare il test di
personalità per sé e per il proprio gruppo, facendolo fare a chi non l'ha
ancora fatto: tornerà molto utile per cooperare meglio.

% ---------------------------------------------------------------------
\section{Anticipazioni sulla seconda giornata}
\label{sec:z1-domani}

La seconda giornata del workshop sarà più pratica: ``oggi vi ho spiegato
perché, domani vi insegno come lo fanno queste aziende''. Gli argomenti
previsti:
\begin{description}
  \item[Time management] come organizzarsi una giornata professionale in
        azienda: come strutturare una \emph{task list}, cosa fare e cosa non
        fare, come e cosa prioritizzare.
  \item[Brainstorming] i gruppi dovranno fare moltissimi meeting (come facciamo
        il gioco, come lo chiamiamo, che scopo gli diamo). Se non si sanno fare
        bene, i brainstorming sono \textbf{dannosi ed eterni}: un brainstorming
        dovrebbe durare \textbf{al massimo due ore}. Verranno mostrati metodi
        che accorciano la durata, massimizzano l'effetto e, soprattutto,
        \textbf{fanno partecipare tutti}: non c'è niente di peggio di un
        brainstorming in cui parlano in due per il 90\% del tempo. Il classico
        ``ragazzi, non esistono idee sbagliate, sparate'' è il modo peggiore di
        farlo; non è colpa di nessuno, ma si può fare diversamente.
  \item[Guidare da dentro] come essere leader nel proprio team anche senza
        esserne il leader: sviluppare responsabilità personale e verso gli
        altri. Non essere il manager di qualcuno non significa non esserne
        responsabili: la persona accanto è responsabilità di tutti, perché si
        è un team e non un insieme di individui. ``Se qualcuno cade, tutti lo
        facciamo rialzare: non c'è l'addetto al rialzamento.'' È la versione
        \emph{advanced} della giornata: come comportarsi da manager senza
        averne il titolo.
\end{description}
Torneranno inoltre il \emph{service level agreement} e il tema della
condivisione dei successi e della protezione del team dai fallimenti.

% ---------------------------------------------------------------------
\section*{Punti chiave del capitolo}
\addcontentsline{toc}{section}{Punti chiave del capitolo}

\begin{itemize}
  \item ``La capacità di un gruppo di persone di fare cose straordinarie
        dipende dal modo in cui queste persone operano insieme.'' Il 90\% dei
        problemi dei team sono \textbf{problemi di comunicazione}.
  \item Le aziende si liberano sempre più presto dello \textbf{useless
        genius}: le competenze si trasmettono, il carattere no. Il
        \textbf{fit culturale} va valutato anche dal candidato.
  \item Non esiste un unico modo di dare feedback: dipende dal contesto e
        soprattutto da \textbf{chi lo riceve}. Il feedback è una somma di
        piccoli miglioramenti periodici; tenersi le cose nel cassetto non
        serve mai.
  \item I team attraversano \textbf{forming, storming, norming, performing e
        adjourning}. Lo storming non si salta; si passa la maggior parte del
        tempo in forming; la bravura sta nell'accelerare i passaggi e nel
        fare bene il \textbf{post-mortem}, anche dei successi.
  \item Gli obiettivi di un progetto vanno fissati \textbf{all'inizio}: solo
        così si può valutare se è andato bene e scegliere tra A e B.
  \item Il circolo virtuoso è \textbf{feedback $\rightarrow$ miglioramento
        $\rightarrow$ fiducia}. Accettare un feedback non significa fare come
        dice l'altro: decide chi ha la responsabilità, ma deve argomentare.
  \item ``Aiutami a poterti aiutare'': meglio dire in anticipo ``potrei non
        farcela'' che ``non ce l'ho fatta''. Chiedere aiuto non è una
        debolezza.
  \item Quando si dà feedback \textbf{non si porta la soluzione}: la si
        costruisce insieme a chi lo riceve.
  \item La cultura è forte quando si lavora \textbf{per gli altri}; il lavoro
        è valutato come opera d'insieme (anche con bonus di team).
  \item Come in \emph{Among Us}, i team sprecano energie a cercare il
        traditore invece di fare le task. \textbf{Default to trust}: chi ha
        sbagliato ha agito con le informazioni che aveva. Cercare
        \textbf{la causa, non il colpevole}: ``che cosa ti ha fatto pensare
        che fosse la scelta giusta?''.
  \item Carriere e performance non sono lineari: vanno valutate le
        \textbf{traiettorie}, non un punto dell'anno. Le valutazioni puntuali
        generano \emph{lying, hiding, faking}.
  \item Confermare sempre le istruzioni facendosele ripetere; ``non ho
        capito'' non è una colpa.
  \item Dare feedback: sull'\textbf{azione} e sui suoi \textbf{effetti}, specifico,
        presto, uno a uno, adattato al carattere, senza terze parti, senza
        ``non mi piace'', senza offendersi se non viene seguito; poi fare il
        debriefing.
  \item Ricevere feedback: \textbf{ascolta, chiedi dettagli, accetta, agisci}.
        Il feedback positivo serve a non peggiorare ciò che si fa bene; il
        feedback accidentale e il \emph{drive-by} non servono a nulla.
  \item Feedback \textbf{frequenti e recenti}: tanti piccoli \emph{sanity
        check} costano meno di una grande correzione di rotta.
  \item Il feedback è \textbf{un regalo} di chi ha a cuore il nostro
        interesse; fa paura il silenzio, non le critiche.
\end{itemize}

%    % =====================================================================
%  Lezione zero 2 — Gestire il tempo, generare idee, guidare da dentro
%  Fonte: trascrizione della seconda "lezione zero" (GD0_2), seconda e
%  ultima giornata del workshop di due giorni tenuto prima dell'inizio del
%  corso di Game Design, dedicato alla figura professionale e al lavoro in
%  team. Prosegue direttamente il capitolo 0.1 (cultura del feedback).
%  Pacchetti richiesti nel preambolo del libro:
%    \usepackage[italian]{babel}, \usepackage[T1]{fontenc},
%    \usepackage[utf8]{inputenc} (solo pdfLaTeX), \usepackage{booktabs}
%
%  Collocazione nel libro: parte "Lezioni zero", subito dopo
%  GD0_1_cultura_del_feedback, con la stessa numerazione propria (0.2).
%  Nel file principale del libro, ad esempio:
%
%    \part{Lezioni zero: il game designer nell'industria}
%    \renewcommand{\thechapter}{0.\arabic{chapter}}
%    \setcounter{chapter}{0}
%    % =====================================================================
%  Lezione zero 1 — Lavorare in squadra: la cultura del feedback
%  Fonte: trascrizione della prima "lezione zero" (GD0), prima giornata di
%  un workshop di due giorni tenuto prima dell'inizio del corso di Game
%  Design, dedicato alla figura professionale e al lavoro in team.
%  Pacchetti richiesti nel preambolo del libro:
%    \usepackage[italian]{babel}, \usepackage[T1]{fontenc},
%    \usepackage[utf8]{inputenc} (solo pdfLaTeX), \usepackage{booktabs}
%
%  Collocazione consigliata nel libro: le lezioni zero precedono il corso,
%  quindi vanno in una parte separata PRIMA di GD1, con numerazione propria
%  (0.1, 0.2, ...) così che i capitoli del corso conservino i numeri 1, 2, 3.
%  Nel file principale del libro, ad esempio:
%
%    \part{Lezioni zero: il game designer nell'industria}
%    \renewcommand{\thechapter}{0.\arabic{chapter}}
%    \setcounter{chapter}{0}
%    % =====================================================================
%  Lezione zero 1 — Lavorare in squadra: la cultura del feedback
%  Fonte: trascrizione della prima "lezione zero" (GD0), prima giornata di
%  un workshop di due giorni tenuto prima dell'inizio del corso di Game
%  Design, dedicato alla figura professionale e al lavoro in team.
%  Pacchetti richiesti nel preambolo del libro:
%    \usepackage[italian]{babel}, \usepackage[T1]{fontenc},
%    \usepackage[utf8]{inputenc} (solo pdfLaTeX), \usepackage{booktabs}
%
%  Collocazione consigliata nel libro: le lezioni zero precedono il corso,
%  quindi vanno in una parte separata PRIMA di GD1, con numerazione propria
%  (0.1, 0.2, ...) così che i capitoli del corso conservino i numeri 1, 2, 3.
%  Nel file principale del libro, ad esempio:
%
%    \part{Lezioni zero: il game designer nell'industria}
%    \renewcommand{\thechapter}{0.\arabic{chapter}}
%    \setcounter{chapter}{0}
%    \input{Seminari/GD0_1_cultura_del_feedback}
%    % ... altre lezioni zero ...
%    \part{Il corso di Game Design}
%    \renewcommand{\thechapter}{\arabic{chapter}}
%    \setcounter{chapter}{0}
%    \input{GD1/GD1_introduzione_al_game_design}
%    % ...
%
%  (Con hyperref aggiungere anche \renewcommand{\theHchapter}{z.\arabic{chapter}}
%  nella parte zero, per evitare ancore duplicate.)
%  Tutte le etichette di questo capitolo hanno il prefisso "z1:" per non
%  collidere con quelle dei capitoli del corso.
% =====================================================================

\chapter[La cultura del feedback]{Lavorare in squadra: la cultura del feedback}
\label{cap:z1-feedback}

Questo capitolo apre le \emph{lezioni zero}: incontri che precedono l'inizio
vero e proprio del corso di Game Design e che non riguardano una tecnica di
design in senso stretto, ma \textbf{la figura professionale} di chi lavora
nell'industria dei videogiochi. Prima di imparare a progettare un gioco
conviene infatti capire in che tipo di ambiente lo si progetterà: un ambiente
fatto di team, di gerarchie, di scadenze e, soprattutto, di persone che devono
riuscire a lavorare insieme.

La lezione è la prima giornata di un \textbf{workshop di due giorni} (circa
dieci ore in totale, pause pranzo comprese) tenuto da Stefano
Cozzi\footnote{Il cognome è riportato come compare nella trascrizione
automatica, che in più punti storpia i nomi propri; andrebbe verificato.
Nella stessa trascrizione lo studio per cui lavora appare come
``Stormine'', ``Stormline'' e ``Stormind'': si tratta di Stormind Games,
studio con sede a Catania, coerentemente con quanto detto in lezione.}, un
professionista con circa quindici anni di esperienza nel settore, tra
Blizzard Entertainment, Riot Games e, oggi, Stormind Games. Il workshop è
nato da un confronto con i coordinatori dell'accademia\footnote{Nella
trascrizione l'accademia è resa in modi diversi (``Hive'', ``a dive'',
``D.I.V.''); i coordinatori citati sono Enrico e Andrea, insieme al
direttore Diego Ricchiuti. Nel resto del capitolo si usa semplicemente
``l'accademia'', come negli altri capitoli.}: serviva un momento dedicato a
\emph{come ci si comporta in un ambiente professionale}, cosa che riguarda sia
il lavoro futuro sia i progetti di gruppo svolti durante il percorso di
studi, dove ci si deve già abituare a comportarsi ``come se si fosse in
un'azienda''.

La prima giornata ha un unico grande tema, sezionato in più parti:
\textbf{la cultura del feedback}. La seconda giornata (che sarà oggetto della
prossima lezione zero) sarà più pratica: \emph{time management},
\emph{brainstorming} e il ``guidare da dentro'' un team anche senza esserne
il leader (sezione~\ref{sec:z1-domani}).

Il percorso del capitolo segue quello della lezione:
\begin{enumerate}
  \item il formato del workshop e il percorso professionale del relatore
        (sezione~\ref{sec:z1-relatore});
  \item perché le aziende tengono tanto al lavoro di squadra e al
        \emph{fit} culturale (sezione~\ref{sec:z1-squadra});
  \item che cos'è il feedback e da che cosa dipende il modo di darlo
        (sezione~\ref{sec:z1-cose-feedback});
  \item le cinque fasi di vita di un team e il valore del
        \emph{post-mortem} (sezioni~\ref{sec:z1-fasi-team}
        e~\ref{sec:z1-qa-fasi});
  \item il confronto con i colleghi e alcune domande sul lavoro nei grandi
        studi (sezioni~\ref{sec:z1-rivalita} e~\ref{sec:z1-industria});
  \item il circolo virtuoso tra feedback, miglioramento e fiducia, e i suoi
        effetti sulla \emph{retention} e sulla cultura aziendale
        (sezioni~\ref{sec:z1-triangolo} e~\ref{sec:z1-cultura});
  \item il ``gioco del traditore'' e il principio del
        \textbf{default to trust} (sezioni~\ref{sec:z1-among-us}
        e~\ref{sec:z1-default-trust});
  \item le traiettorie di carriera e la valutazione delle performance
        (sezione~\ref{sec:z1-traiettorie});
  \item come si dà, come si riceve e quando si dà un feedback
        (sezioni~\ref{sec:z1-dare}, \ref{sec:z1-ricevere}
        e~\ref{sec:z1-quando});
  \item il feedback come regalo, le domande finali e le anticipazioni
        della seconda giornata (sezioni~\ref{sec:z1-regalo},
        \ref{sec:z1-qa-finale} e~\ref{sec:z1-domani}).
\end{enumerate}

Una premessa di metodo: il relatore ha alternato continuamente teoria,
domande alla classe e racconti della propria esperienza. I racconti non sono
digressioni: sono gli esempi concreti con cui i principi vengono resi
tangibili, e per questo sono riportati con cura.

% ---------------------------------------------------------------------
\section{Il workshop e il relatore}
\label{sec:z1-relatore}

\subsection{Formato e regole di partecipazione}

Il relatore ha cercato di ``riassumere quindici anni di esperienza in dieci
ore''. Ha materiale per circa il doppio delle ore disponibili: le cose
strettamente necessarie vengono trattate per prime, il resto è tenuto ``in
riserva'' come \emph{nice to have}. Le regole del workshop:
\begin{itemize}
  \item le domande si fanno \textbf{alla fine di ogni slide}, così da chiudere
        prima il discorso; non c'è un'unica sessione finale di domande e
        risposte (anche se ci sarà), perché il relatore preferisce fermarsi
        spesso e non fare ``uno spiegone di dieci ore'';
  \item le domande non devono per forza riguardare la slide appena vista: se
        una slide dà lo spunto per altro, se ne può parlare; al massimo sarà
        il relatore a rispondere in breve o a rimandare a una slide
        successiva;
  \item si può chiedere anche di esperienze passate, opinioni e sensazioni;
        se per NDA o per serietà professionale non può rispondere, lo dirà
        esplicitamente (``non lasciate che sia un \emph{blocker}'');
  \item su Zoom si usa l'alzata di mano o la chat, che il relatore tiene
        aperta sul secondo monitor;
  \item non ci sono pause ufficiali: la lezione è registrata e ognuno è
        libero di assentarsi quando serve, cercando però di non disturbare il
        flusso della lezione;
  \item le slide verranno fornite integralmente tramite la segreteria
        dell'accademia (Vanessa), insieme alle registrazioni: le slide senza
        il parlato servono poco, ma sono un buon promemoria.
\end{itemize}

\subsection{Il percorso professionale}

Il relatore si presenta (scherzosamente) come la ``guida escursionistica''
che accompagna la classe nella ``foresta incantata'' della cultura del
feedback, perché per molti è un territorio sconosciuto. Il suo percorso:
\begin{itemize}
  \item in Italia ha lavorato per aziende che fornivano servizi a terzi, tra
        cui il lavoro su \emph{Warhammer Online}, e per un canale televisivo
        sul digitale terrestre dedicato ai videogiochi\footnote{Nella
        trascrizione: ``Digital Game Network Channel''. Il nome esatto del
        canale andrebbe verificato.};
  \item a \textbf{Blizzard Entertainment} (circa sei anni, a Parigi) ha fatto
        il \emph{community manager} per l'Italia (\emph{StarCraft~II},
        \emph{World of Warcraft} da \emph{Mists of Pandaria} in poi,
        \emph{Hearthstone}, \emph{Diablo~III}, le espansioni di
        \emph{StarCraft}, il lancio di Battle.net), per poi diventare il
        manager dei community manager europei;
  \item a \textbf{Riot Games} (circa cinque anni, a Dublino) è stato
        brevemente community manager italiano e poi \emph{publishing
        manager} di un team di circa trenta persone per l'Europa. Il
        \emph{publishing} non è solo community: è l'unione di
        \textbf{community, marketing e PR};
  \item a \textbf{PG Esports} si è occupato di eventi dal vivo e online
        (ad esempio la Milano Games Week e tornei di \emph{Valorant} e
        \emph{League of Legends}), restando in contatto con Riot e Blizzard;
  \item da circa un anno è tornato nello sviluppo, a \textbf{Stormind
        Games}, dove oggi si occupa di risorse umane e organizzazione del
        lavoro; lo studio verrà presentato più ampiamente al \emph{Career
        Day} dell'accademia insieme alla sua HR manager;
  \item fa inoltre consulenze aziendali anche al di fuori dei videogiochi e
        insegna in master universitari (IULM di Milano).
\end{itemize}

Il suo ruolo ricorrente, sia a Blizzard sia a Riot, è stato quello di
\emph{Mr.~Wolf} (``per chi ha visto \emph{Pulp Fiction}''): quello che
risolve i problemi. Veniva inserito nei team più in difficoltà, ad esempio
dal punto di vista delle performance, per sistemarli e passare al successivo.
Era presente in alcuni dei momenti più critici, come il lancio di
\emph{Diablo~III} o il lancio di espansioni di \emph{World of
Warcraft}\footnote{La trascrizione riporta ``il lancio con i Gary's'' e
``World of Trenor'': si tratta verosimilmente di \emph{Mists of Pandaria} e
di \emph{Warlords of Draenor}.}, a fare \emph{damage control}.

La scoperta fondamentale di questi anni è che i problemi da risolvere
\textbf{quasi mai riguardavano il prodotto} (non era certo lui a sistemare i
server di \emph{Diablo~III}), ma i team che c'erano dietro: persone che
litigavano, non collaboravano, non riuscivano a stare nella stessa stanza
senza urlarsi contro. E nel 90\% dei casi tutto questo si riduce a
\textbf{problemi di comunicazione}: è raro che perfino le incompatibilità
caratteriali non siano risolvibili parlando.

\subsection{Struttura non vuol dire burocrazia}

A Stormind il relatore sta accompagnando un team passato da 30 a 95 persone
in meno di un anno, attraverso tutti i cambiamenti di cui parla il workshop.
Darsi una struttura viene spesso percepito male, come ``più processi, più
procedure'', ma bisogna distinguere:
\begin{itemize}
  \item un processo come ``timbrare il cartellino'' può essere davvero
        invasivo e noioso;
  \item definire \emph{come funziona un documento di game design, chi lo
        approva, chi lo visiona e chi no} è invece una struttura necessaria.
\end{itemize}
Uno studio di 30 persone nella stessa sede può farne a meno: quando c'è un
dubbio ci si alza e si va alla scrivania del collega. Uno studio di 95
persone sparse per l'Italia (e in parte per l'Europa) non può più farlo:
servono una struttura diversa, un modo di lavorare diverso e una
pianificazione del futuro diversa.

% ---------------------------------------------------------------------
\section{Perché il lavoro è sempre lavoro di squadra}
\label{sec:z1-squadra}

\subsection{Non si va in ufficio per fare amicizia}

Il primo messaggio è netto: \textbf{non si va in ufficio per fare amici} o
per trovare compagni di vita. Si va in ufficio perché c'è un lavoro da fare, e
lo scopo è riuscire a farlo \emph{anche con chi sta antipatico}. Con le
persone simpatiche è facile: diventeranno amici, e non sarà mai un problema.
È con chi sta antipatico che le cose si complicano, ed è lì che emergono i
\textbf{veri professionisti}: quelli che riescono a lavorare con chiunque
senza che l'antipatia pesi sul modo in cui lavorano e collaborano, mettendo il
prodotto e il team davanti al proprio ego e alle proprie preferenze personali.

La frase che fa da filo conduttore a tutto il workshop (e che il relatore
promette di riproporre ``fino alla nausea'') è:

\begin{quote}
  \emph{La capacità di un gruppo di persone di fare cose straordinarie
  dipende dal modo in cui queste persone operano insieme.}
\end{quote}

Le aziende raramente pensano che l'innesto di una singola persona possa
cambiare il destino di un progetto. Questo è tanto più vero quanto più grande
è il team, ma vale anche per i piccoli studi, anche per un gruppo di cinque
persone. Come negli sport di squadra: l'individuo conta fino a un certo punto,
e anche se si ha Dončić in squadra, se si gioca uno contro cinque si perde
comunque\footnote{Il riferimento è alle finali NBA in corso nei giorni del
workshop tra Dallas Mavericks e Boston Celtics (giugno 2024), il che
permette di datare approssimativamente la lezione.}.

\subsection{Lo \emph{useless genius} e il costo delle persone tossiche}

Molti dipartimenti HR usano la definizione di \textbf{useless
genius}\footnote{Il concetto è più noto in letteratura manageriale come
\emph{brilliant jerk}; in lezione è usata l'espressione \emph{useless
genius}.}: si può essere bravissimi, a volte i migliori nel proprio lavoro,
ma se non ci si sa comportare, relazionare e operare in un gruppo si è
fondamentalmente inutili ai fini del gruppo. Nel migliore dei casi inutili,
nel peggiore \emph{dannosi}.

In passato le aziende tendevano a ``chiudere un occhio'': la persona è
fondamentale, la teniamo. Oggi la domanda è diventata: \emph{quanto costa
all'azienda, oltre al suo stipendio, tenere una persona che crea un ambiente
tossico e peggiora chi le sta intorno?} Ci sono costi sommersi che le aziende
cominciano a vedere, e la tendenza è liberarsi di queste persone il prima
possibile. Il ragionamento è semplice:
\begin{itemize}
  \item le \textbf{competenze si possono trasmettere} o sostituire: si trova
        qualcuno che fa lo stesso lavoro magari nel doppio del tempo, oppure
        più persone che se lo dividono;
  \item il \textbf{carattere delle persone} invece cambia molto raramente.
\end{itemize}

\subsection{Il fit culturale}

Quando qualcuno \textbf{non è un \emph{fit} culturale} per un'azienda è
meglio separarsi presto. Ma non essere adatti a \emph{quell'}azienda non
significa non essere adatti a lavorare: è perfettamente legittimo che una
persona inadatta a un'azienda sia adattissima a un'altra.

Il relatore l'ha vissuto in prima persona. Dopo sei anni a Blizzard viene
contattato da Riot e affronta un processo di selezione lunghissimo (``quasi
undici colloqui''), al termine del quale viene scartato. Scrive allora una
mail che gli ha ``cambiato la traiettoria di carriera'': chiede
semplicemente perché non è stato ritenuto adatto. La risposta: in nessun
colloquio si era esposto abbastanza da far capire \emph{perché} volesse
lasciare Blizzard, e quindi pensavano che in realtà stesse bene dov'era. Il
motivo era che, per educazione professionale, non gli piace parlare male del
proprio posto di lavoro; portata all'eccesso, però, questa riservatezza era
diventata un ``è tutto bellissimo, ma vorrei provare Riot'', senza un perché.
Lo invitano comunque a un incontro dal vivo (``dal vivo è diverso''), e alla
fine viene assunto.

Una volta arrivato, i sei mesi di prova (\emph{probation}) sono serviti tutti
per smettere di ragionare e lavorare ``come a Blizzard''. La differenza
culturale era enorme, tanto che in Riot esisteva un gruppo di supporto
interno tra ex dipendenti Blizzard: un collega gliel'ha descritto come
``uscire da una relazione tossica e doversi riabituare a una relazione
normale''\footnote{Il relatore precisa che parla della Blizzard di allora,
che non esiste più: oggi è parte di Microsoft e nessuno dei suoi ex colleghi
vi lavora ancora, anche a causa dei licenziamenti di massa.}.
\begin{description}
  \item[Blizzard] era \emph{developer first}: il focus era chiaramente sugli
        sviluppatori, prima di tutti gli altri. Dopo sei anni c'erano ancora
        meeting a cui non poteva partecipare perché non abbastanza
        \emph{senior}, informazioni che non poteva sapere, piani che non
        poteva vedere.
  \item[Riot] era \emph{player first}: prima i giocatori, e i dipendenti come
        mezzo per trasferire il feedback dei giocatori dentro il gioco. Il
        secondo giorno lo portano in una stanza e gli mostrano il piano dei
        cinque anni successivi. Lui, sorpreso, pensa a un errore (``sono qui
        da oggi, non vorrei mettervi nei guai''); la risposta è: ``sei un
        dipendente Riot o no? Allora non ci sono livelli di accesso: devi
        sapere queste cose per fare meglio il tuo lavoro''.
\end{description}

Da questa esperienza ha imparato molto sulla \textbf{trasparenza}: condividere
le informazioni aiuta il lavoro. Non sempre, però: a volte non si condivide
qualcosa non per segretezza, ma perché non è ancora pronta per essere
presentata, esattamente come non si mostra un gioco quando è ancora solo
teoria (``te lo faccio vedere quando ho qualcosa di tangibile'').

La lezione per gli studenti è doppia: le aziende valuteranno il \emph{fit}
culturale sempre di più, ma capire se si è un \emph{fit} richiede anche
\textbf{lavoro da parte del candidato}. Bisogna avere la presenza di spirito
di dire: ``conosco i valori di questa azienda, non mando nemmeno il CV,
perché so che durerei un mese e mezzo''.

% ---------------------------------------------------------------------
\section{Che cos'è il feedback}
\label{sec:z1-cose-feedback}

\subsection{Le prime risposte della classe}

La prima domanda alla classe è: \emph{quando sentite la parola
``feedback'', a cosa pensate?} Le risposte raccolte, a voce e in chat:
\begin{itemize}
  \item un \textbf{riscontro} su ciò che si fa e su come lo si fa, sul lavoro
        o sul modo in cui ci si interfaccia con il gruppo; può essere
        lavorativo oppure comunicativo e comportamentale. Chi ha risposto lo
        vede come qualcosa che, se l'azienda non dà, va attivamente
        richiesto, perché è importante sapere come sta andando per
        migliorarsi;
  \item né positivo né negativo, ma \textbf{costruttivo}: in entrambi i casi
        serve a migliorarsi;
  \item uno \emph{scambio di opinioni}, un \emph{riscontro produttivo};
  \item un feedback ``di reazione'', non verbale: si fa una cosa e si
        osserva come cambia l'ambiente o il comportamento delle persone
        intorno (anche perché a voce la gente spesso si tiene le cose
        negative);
  \item per un designer, anche il \textbf{linguaggio del corpo}: osservando
        un giocatore giocare si può dedurre l'emozione che prova (frustrazione,
        incomprensione).
\end{itemize}

Sulla forma (scritto o orale?) le opinioni divergono: c'è chi lo preferisce
\emph{mille volte} scritto; chi pensa che dipenda dalle dimensioni
dell'azienda (orale nelle piccole, scritto in quelle con team di centinaia di
persone); chi pensa che dipenda dall'importanza del feedback (scritto se
tecnico o pieno di dati, da tenere sotto gli occhi; orale se è un consiglio
veloce).

Il relatore usa due esempi per mostrare che \emph{dipende}:
\begin{itemize}
  \item un feedback sulla programmazione ha senso darlo scritto, magari
        commentando direttamente il codice, più che ``andando a prendere un
        caffè'';
  \item ma se il feedback è ``ieri, durante il meeting, ho percepito male il
        fatto che mi interrompessi di continuo'', ce lo si aspetta
        \textbf{orale}.
\end{itemize}

\subsection{Da che cosa dipende il modo di dare feedback}

Lo \emph{spoiler} della giornata è che \textbf{non esiste una risposta univoca}
su come dare feedback. Il workshop non insegna ``il'' modo, ma fornisce
\emph{tutte le opzioni}, perché il modo giusto cambia in base a molti fattori:
\begin{itemize}
  \item a chi lo si dà: un collega, un superiore, il proprio capo, un proprio
        \emph{report} (un sottoposto);
  \item se è personale o professionale;
  \item se è la prima volta o la terza volta che si dà lo stesso feedback;
  \item se riguarda un progetto importante o secondario.
\end{itemize}

Manca però l'elemento fondamentale, che la classe individua subito:
\textbf{l'altro}. A parità di tutti gli altri fattori, il \emph{carattere
della persona che riceve il feedback} dovrebbe far cambiare il modo in cui lo
si dà. Per questo, prima della lezione, agli studenti era stato chiesto di
fare un test di personalità: non un test di psicologia avanzata, ma uno
strumento per rendersi conto di quanta varietà ci sia anche solo in una
classe di 53 persone (il che non vuol dire trovare 52 modi diversi di dare
feedback, ma lavorarci un po' sì).

Il relatore porta se stesso come esempio. Preferisce:
\begin{itemize}
  \item i feedback \textbf{positivi scritti}, perché è a disagio a ricevere
        complimenti in pubblico;
  \item i feedback sul proprio lavoro \textbf{diretti}: ``vai subito al
        punto e non girarmi intorno''.
\end{itemize}
È molto orgoglioso del suo lavoro e, quando qualcuno lo critica, la sua
primissima reazione istintiva è ``come ti permetti''. Negli anni ha imparato
a farla durare un secondo, ma chi gli dà feedback di persona dice di
vederla ancora negli occhi. Per mitigarla, dice esplicitamente a chi gli dà
feedback: \emph{``ignora il modo in cui reagisco nei primi cinque secondi,
vai avanti, poi mi passa''}. È considerato da tutti molto aperto al feedback
e non ha problemi a dire ``hai ragione, facciamo diversamente'': ma ci sono
reazioni istintive e caratteriali che è difficile controllare al cento per
cento.

\subsection{Con chi si lavora meglio}

Guardando le persone con cui si lavora abitualmente (anche nei progetti
dell'accademia), si noterà che si lavora meglio con due profili:
\begin{enumerate}
  \item chi è \textbf{esattamente come noi}, ma qui è un 50 e 50: alcuni
        caratteri vanno d'accordo con i propri simili, altri non vogliono
        nessuno con lo stesso carattere nel gruppo (e si potrebbe rientrare
        in questa categoria);
  \item chi è \textbf{affine}: ha qualcosa in comune con noi, ma non tutto,
        e proprio per questo può dare punti di vista differenti.
\end{enumerate}
Il lavoro della giornata è capire come valorizzare il secondo caso: ottenere
feedback da persone con caratteri, visioni e background diversi, per
completare una visione del proprio lavoro che è necessariamente limitata. È
come esplorare una mappa con la nebbia di guerra: un'altra persona, al di là
del muro, vede cose che noi non vediamo.

\subsection{Il feedback come somma di piccoli miglioramenti}

Per il relatore il feedback è questo: \textbf{applicare tanti piccoli
miglioramenti periodici}, che si sommano nel tempo. Migliorare dello 0,1\%
ogni giorno porta a fine anno a un miglioramento enorme, e questo vale per i
prodotti, per i giochi, per le persone e per i team.

\paragraph{L'esempio del ciclismo britannico.} La squadra di ciclismo
britannica, fino a una quindicina d'anni fa, non aveva vinto quasi nulla.
Cambiato il team manager, tutti si aspettavano grandi cambiamenti sportivi
(convocazioni, strategie di gara). Arrivarono invece \emph{minuscoli
cambiamenti non sportivi}: un altro chef in trasferta, materassi migliori nei
camper per riposare meglio, un sellino più comodo a fine tappa. Cambiamenti
marginali che da soli non avrebbero cambiato niente; eppure, dopo due o tre
anni di questo lavoro quotidiano, la squadra ha cominciato a vincere a
ripetizione, Olimpiadi comprese, cambiando il modo in cui il ciclismo è
visto in Gran Bretagna\footnote{È la nota strategia delle \emph{marginal
gains} associata a Dave Brailsford (British Cycling e Team Sky); il nome non
è citato in lezione.}. Certo, si è investito anche sulle persone giuste, ma
tutto è partito dai piccoli miglioramenti periodici.

Per gli studenti vale lo stesso: chiarire un conflitto con una persona con
cui si ha ancora poco a che fare, o una situazione che ha lasciato
``l'amaro in bocca'', sembra inutile oggi, ma darà i suoi frutti tra tre, sei,
nove mesi. Sono \emph{investimenti sul futuro}.

\paragraph{Peggiorare ogni giorno.} Esiste anche il percorso inverso:
lasciar correre, ``buttare sotto il tappeto'' le cose perché ``tanto è
inutile parlarne''. Queste cose tornano sempre fuori dopo otto mesi, come il
\emph{libro dei rancori} dei nani di Warhammer: ``ti ricordi quando, un anno
e mezzo fa, all'inizio del progetto, mi hai detto quella cosa? Mi sono
offeso''. A quel punto è quasi impossibile risolvere, perché nessuno ricorda
più con precisione cosa sia successo (``ma quale cena?'', ``forse era in
pizzeria'', ``no, eravamo da un'altra parte''). \textbf{Tenersi le cose nel
cassetto per tirarle fuori quando servono è sempre sconsigliato}: non servono
ed è impossibile usarle in modo costruttivo.

% ---------------------------------------------------------------------
\section{Le cinque fasi di vita di un team}
\label{sec:z1-fasi-team}

\subsection{Perché conoscerle}

Quasi tutti gli studenti lavorano già in gruppo (al primo anno su un gioco da
tavolo, negli anni successivi su un videogioco) e si saranno accorti che un
team attraversa delle fasi: il modo di lavorare insieme cambia nel tempo.
Spesso ci si fa prendere di sorpresa: ``eravamo così affiatati, cosa è
successo?''.

Un esempio da una consulenza fuori dai videogiochi: un'azienda farmaceutica
gli affida un team commerciale (rappresentanti che presentano medicinali a
ospedali, medici e privati). Il problema: ``prima andava tutto bene; ora che
siamo in diciotto dovremmo lavorare meno e meglio, e invece va peggio''. Va
peggio perché le persone nuove sono appena arrivate, mentre il management si
aspettava già di vedere i risultati dell'investimento.

Fino a pochi anni fa le fasi di un team venivano disegnate come una
\textbf{linea dritta verso l'alto}: le fasi erano le stesse, ma viste come un
miglioramento costante. Non si vedevano le parti \emph{fisiologiche} del
lavoro in team, cioè anche i peggioramenti. Accettarle è il primo passo per
accettare i fallimenti: nessun percorso personale o di team è lineare.

Il relatore ammette che il suo modo di lavorare è nato per il 70\% dal
proprio impegno a imparare, fare bene e rispettare le persone, e per il 30\%
dall'aver lavorato con ``incapaci colossali'' che gli hanno insegnato come
\emph{non} fare le cose. A Blizzard un manager gli fece capire che, pur non
sapendo ancora che tipo di \emph{people manager} sarebbe diventato, sapeva
benissimo quale tipo di manager \emph{non} voleva essere.

\subsection{Forming, storming, norming, performing, adjourning}

Le fasi sono cinque, un po' come le fasi del lutto. Molte aziende ne
considerano solo quattro (``e sono quelle che solitamente chiudono''); le
aziende al passo con i tempi ne considerano cinque\footnote{Il modello è
quello di Bruce Tuckman (1965), a cui la quinta fase, \emph{adjourning}, fu
aggiunta nel 1977. Il nome di Tuckman non è citato in lezione; la
trascrizione riporta la quinta fase come ``journing''.}.

\begin{description}
  \item[Forming] Si viene messi insieme e si comincia a lavorare. Non si sa
        nulla degli altri; spesso non si conoscono nemmeno i nomi, ed è
        normale.
  \item[Storming] Si cominciano a ``dare le testate come le capre in
        montagna per determinare il territorio'': questa cosa la faccio io,
        questa la vorrei fare così, chi decide? Si stabilisce chi fa cosa e
        come si divide il lavoro. Nel grafico è una \textbf{fase di
        peggioramento}: che il team litighi o discuta civilmente non cambia
        nulla, il team deve investire parte di tempo ed energie per risolvere
        queste questioni. È normale e \textbf{non è evitabile}.
  \item[Norming] Si comincia a capire come lavorano gli altri, cosa sanno
        fare e come far funzionare il gruppo; iniziano i miglioramenti. È una
        forma di \emph{accettazione} degli altri, che deriva dall'aver capito
        cosa devo fare io e cosa devono fare loro.
  \item[Performing] C'è uno sforzo \textbf{proattivo}: ``questa cosa
        toccherebbe a me, ma so che Alessio la fa molto meglio, quindi gli
        chiedo aiuto''. Si sfruttano i punti di forza degli altri, perché li
        si conosce e perché si sa che passare il testimone non avrà effetti
        negativi (il capo non dirà ``perché non l'hai fatta tu?'').
  \item[Adjourning] Il progetto finisce, si chiude o viene riorganizzato, il
        team viene sciolto. Quello che andrebbe fatto è analizzare
        \textbf{cosa ha funzionato, cosa no e cosa migliorare} per il
        progetto successivo.
\end{description}

\subsection{La fase che tutti saltano}

``Spoiler incredibile'': l'\emph{adjourning} è la fase che molte aziende
saltano, \textbf{perché c'è sempre il progetto dopo}.

A PG Esports, per la Milano Games Week (come per qualsiasi evento dal vivo a
cadenza annuale: Lucca Comics, Etna Comics, Comicon\dots), finita l'edizione
di un anno si cominciava dal giorno dopo a lavorare a quella successiva,
perché un anno sembra tanto ma non basta mai. Nei videogiochi è lo stesso:
finito un gioco, a parte le patch di mantenimento e del \emph{day one}, si
passa subito al successivo.

Il cambiamento introdotto dal relatore (in parte già praticato, ma da lui
strutturato) è stato: finita la Games Week, \textbf{tutti in ferie per una
settimana}; poi ci si ritrova per definire cosa ha funzionato e cosa no, con
un \textbf{post-mortem} (o \emph{debrief}); solo dopo si inizia a lavorare
all'edizione successiva. Senza questa fase manca la parte di
\emph{learning}, cioè di ciò che si è imparato dal progetto.

Il consiglio vale per gli studenti: quando si finisce il gioco da tavolo, e
poi il gioco digitale, bisogna trovare il tempo per il post-mortem,
altrimenti si butta via \textbf{il 90\% dell'utilità del progetto}. Una
struttura semplice per farlo è descritta nella sezione~\ref{sec:z1-qa-finale}.

% ---------------------------------------------------------------------
\section{Domande sulle fasi del team}
\label{sec:z1-qa-fasi}

\subsection{Lo storming che ritorna}

\paragraph{Capita che lo storming si ritagli spazio anche nelle altre fasi?}
Lo \emph{storming} è una fase precisa, che però molti manager, coordinatori e
team cercano di \textbf{saltare}: ``se io sono bravo come manager e loro sono
bravi come team, passeremo da forming a norming senza problemi''. È un errore,
perché senza una fase in cui si discute davvero non si sviluppa una
\emph{skill} fondamentale: \textbf{saper discutere}. Tutti conoscono persone
non abituate a farlo, per cui ogni minimo disaccordo diventa un affronto
personale (``avevamo deciso la pizza, perché andiamo al sushi? Perché mi fate
questo affronto?''): sono persone che non sono abituate ad avere conflitti
costruttivi.

Se si pensa di essere in \emph{norming} e continuano a emergere conflitti, il
punto è che \textbf{non si era davvero in norming}: lo storming non ``si
trascina'' nelle fasi successive, semplicemente non se n'è mai usciti.
Allo stesso modo, a chi chiede se ogni mini-fase di storming durante la vita
del gruppo sia frutto di uno storming fatto male o solo di carattere, la
risposta è già contenuta qui.

Un partecipante osserva che spesso, ancor prima di finire un progetto, si è
già messi su quello successivo: è il modo migliore per essere in costante
\emph{damage control}. Era l'esperienza del relatore a Blizzard: veniva
spostato fisicamente da un team all'altro, e finito il lancio di
\emph{Diablo~III} era già su \emph{World of Warcraft}, senza vivere il
``dopo''. Si è dovuto ritagliare da solo lo spazio per un debrief personale,
ed è lì che ha deciso che tipo di manager non voleva essere. Il consiglio:
\textbf{se non lo fa il team, fatelo voi}.

\subsection{Dove si passa la maggior parte del tempo}

\emph{In quale fase si passa più tempo nella propria carriera?} Gli studenti
rispondono \emph{storming} e \emph{norming}. Il relatore, in quindici anni,
calcola di aver passato in \emph{performing} circa \textbf{dieci mesi}. Non è
pessimismo: arrivare a performing deve essere sempre l'obiettivo, sia da
manager sia da membro del team; ma è difficile arrivarci, ed è normale che non
sempre ci si arrivi.

Lo storming può durare a lungo, ma con team molto esperti dura pochissimo,
perché c'è chi lo sa gestire. In realtà \textbf{la maggior parte del tempo si
passa in forming}, perché ogni equilibrio raggiunto vale per \emph{quelle}
persone, \emph{quel} progetto, \emph{quelle} tempistiche. Basta che cambi un
qualsiasi fattore (una deadline spostata, una persona nuova, qualcuno che si
dimette, un progetto cancellato) e si ricomincia, perché i bilanciamenti
precedenti non valgono più.

Esempio: un team di cinque persone affiatatissimo, in pieno performing. La
dirigenza, entusiasta, ``raddoppia le risorse'' e lo porta a dieci. O si
lascia tutto com'è (e i cinque nuovi non fanno niente), o bisogna ridiscutere
chi fa cosa e come i nuovi possono aiutare: che piaccia o no, si ricomincia
da forming.

Quindi la bravura di un manager, ma anche di un membro del team,
\textbf{non sta nel passare tanto tempo in performing, ma nell'accelerare il
passaggio da una fase all'altra}, e nel fare bene la fase di
\emph{adjourning}.

\subsection{Analizzare anche i successi}

C'è la tendenza a fare post-mortem solo dei progetti andati male: ``questa
cosa non ha funzionato, vediamo perché''. Va benissimo. Ma quando un progetto
va bene si dà per scontato che sia andato bene ``perché era buono''. Bisogna
invece chiedersi \emph{perché} è andato bene: era il prodotto? La
tempistica? Il modo in cui abbiamo fatto una cosa piuttosto che un'altra?
\textbf{Se non si capisce davvero il motivo del proprio successo, si
tenderanno a replicare le cose sbagliate.}

\paragraph{Blockbuster.} Quando arrivò la notizia che Netflix avrebbe
lanciato un servizio di streaming, Blockbuster (la catena di noleggio di
videocassette e DVD) aveva tutto il tempo per preparare un'offerta
concorrente. In una riunione del management raccontata in un libro (di cui il
relatore non ricorda il titolo), la direzione sembrava orientata a investire
nello streaming, finché un dirigente fece notare che il 12\% del fatturato
derivava dalle \emph{multe per riconsegna tardiva}, impossibili con lo
streaming, mentre secondo le stime lo streaming avrebbe prodotto il 10\% del
fatturato. ``Non possiamo rinunciare a questo 2\%'': niente streaming. Oggi
non c'è più un Blockbuster ``manco a cercarlo col lanternino''. Non avevano
capito qual era il vero motivo del loro successo, cioè il loro modello di
business\footnote{Le percentuali sono quelle riferite in lezione e non sono
state verificate.}.

\paragraph{The Founder.} Nel film sulla nascita di McDonald's c'è una scena in
cui il futuro socio del protagonista gli dice: ``tu pensi di essere nella
\emph{food industry}, e invece sei nella \emph{real estate industry}''. Non
dovrebbe investire sugli hamburger, ma sui terreni su cui costruisce i
ristoranti. Anche per un team vale lo stesso: bisogna capire che cosa
l'ha reso davvero di successo, oppure no.

\subsection{L'autoanalisi personale}

\paragraph{E quando si cambia azienda?} Uno studente chiede se il debrief sia
possibile anche quando si lascia un'azienda. Non si può tornare in ufficio a
discutere della propria esperienza, ma lo si può fare \textbf{per se stessi}.
Il relatore l'ha fatto, anche se all'epoca in modo poco consapevole, ed è
così che è passato dal \emph{community management} al \emph{people
management}. Analizzando cosa aveva fatto bene a Blizzard si era reso conto
che:
\begin{itemize}
  \item aveva costruito un ottimo rapporto di fiducia con i giocatori (con
        alcuni proprietari di fansite parla ancora oggi);
  \item pur senza la seniority, era diventato un punto di riferimento per il
        team, visto come imparziale ed equilibrato anche su temi che non lo
        riguardavano direttamente;
  \item sapeva \textbf{risolvere problemi}: durante il lancio di
        \emph{Diablo~III} era il più tranquillo di tutti, non perché
        sminuisse il problema, ma perché sapeva che si sarebbe risolto con i
        suoi tempi; il lavoro non era spegnere l'incendio ma ``far sfollare
        la gente'' e aspettare che tutto tornasse in ordine.
\end{itemize}
E aveva capito cosa sapeva fare meno bene, ad esempio \textbf{gestire il
proprio manager}. Avere un manager non significa che il manager gestisca te:
il manager gestisce il tuo lavoro, ma tu devi saper gestire i tuoi
\emph{stakeholder}. Spesso lo stakeholder è il proprio manager; da manager,
gli stakeholder saranno CEO, CFO e così via. Gestirli significa conoscerli,
sapere come portare loro le informazioni, come far presente che qualcosa non
va. Fa parte del lavoro, anche se a lui non piace, e ci ha lavorato.

Nel mese e mezzo di pausa prima di entrare in Riot si è concentrato su
questo: ``cosa porto come mio punto di forza?''. L'autoanalisi è culminata
nel chiedere espressamente alla persona che l'aveva assunto: ``prima mi avete
detto di no, ora mi avete preso: perché?''.

\paragraph{Come si capisce se un progetto è andato bene?} Dipende
dall'\textbf{obiettivo}. Se l'obiettivo era vendere milioni di copie e il
gioco le ha vendute, bisogna capire perché. Se, più ragionevolmente,
l'obiettivo era un gioco da tavolo che permettesse di mettere in pratica le
conoscenze del corso, va valutato per quello: ha permesso di farlo? Se sì,
perché? Se no, perché no? L'obiettivo va fissato \textbf{all'inizio del
progetto, non alla fine}, proprio per poterlo valutare dopo.

\subsection{Quando il team cambia di continuo}

\paragraph{Tutti questi cambiamenti non sono dannosi per il progetto?}
Assolutamente sì, ma quasi mai sono volontari. I progetti andati male sono
spesso progetti grandi.

Il relatore racconta di \emph{Cyberpunk 2077}. Conosceva persone in CD
Projekt Red (``il mondo è piccolissimo, le cose si sanno''). Alla Gamescom va
con Diego Ricchiuti, direttore dell'accademia, a vedere la presentazione del
primo gameplay, quella ``in cui regalavano il giubbottino da samurai''.
Finito il video, tra gli applausi, si guardano e capiscono entrambi che
\textbf{il gioco non è pronto}, a soli quattro mesi dall'uscita prevista:
Diego aveva notato il combattimento e la guida. Gli amici di CD Projekt Red
(uno aveva lavorato con lui a Blizzard per sei anni) gli confermano tra le
righe che ``ce la faremo, ma ci sono ancora cose da sistemare''. Il gioco fu
poi rimandato più volte\footnote{In lezione il relatore dice ``due, forse
tre volte''; l'uscita di \emph{Cyberpunk 2077} è stata effettivamente
rinviata tre volte prima del dicembre 2020.}.

Per deformazione professionale, chi ha sviluppato l'occhio critico non
riesce più a godersi un trailer: in un trailer di \emph{The Day Before} Diego
gli aveva fatto notare che l'effetto della fiammata di un fucile era stato
chiaramente aggiunto in postproduzione e non faceva parte del gameplay. ``È
il triste epilogo di questa storia'': non riescono più a godersi un trailer.

Il punto è che più il team è grande, più queste variabili vengono imposte, e
più è facile che il fallimento \textbf{non dipenda dal fatto che il prodotto
sia brutto} (\emph{Cyberpunk} è un bellissimo gioco sotto molti punti di
vista, e il relatore lo sta rigiocando ora che è stato corretto ed espanso).
Nei grandi gruppi qualcuno arriva, qualcuno si dimette, cambia il budget,
arriva il rinvio; nel caso di \emph{Anthem}, il team ha fatto così tante volte
forming--storming--forming--storming che ne è uscito ``un prodotto che
dovevamo espellere: l'abbiamo espulso, buttato fuori dal nostro organismo pur
di liberarcene''.

Tra i cambiamenti c'è anche la crescita verticale: se il capo viene promosso
perché lavora benissimo, bisogna scegliere un nuovo capo, ed è sempre un
caos (chi ci resta male, chi pensava di meritarlo, chi voleva il suo amico,
chi non sopporta il prescelto). Per questo il relatore è \textbf{contrario
alle promozioni durante i progetti}: si rischia di rovinare tutto il lavoro di
storming.

\subsection{Obiettivi chiari nei piccoli gruppi}

Nei gruppi piccoli, come quelli degli studenti, il problema deriva più
spesso dal \textbf{non avere chiaro che cosa si vuole ottenere}: un prodotto
puramente educativo e formativo, o un prodotto richiesto dall'accademia che,
se viene bene, si potrebbe anche vendere? Un gioco commerciale o un gioco
artistico? Va chiarito subito, perché ci saranno sempre momenti in cui
scegliere tra A e B, e a guidare la scelta sarà la domanda: \emph{quale delle
due opzioni ci porta più vicino all'obiettivo?} Se l'obiettivo non è chiaro,
non si sa cosa scegliere.

Questo porta molte aziende a fare \textbf{scelte di design che in realtà sono
scelte commerciali}. Le microtransazioni sono scelte commerciali: nessuno ha
mai pensato ``sai cosa piacerebbe tanto ai giocatori? Pagare le bustine''.
Prima le bustine si guadagnavano giocando; ora la tendenza è ``potresti
giocare, ma se paghi è meglio''. Si pensi alle polemiche su Ubisoft per
l'opzione di ``saltare il gioco'' pagando, o ai famosi \emph{boost} di livello
di \emph{World of Warcraft}: ``devo pagarti per saltare il gioco?''.

Una studentessa (Silvia) cita \emph{Final Fantasy XIV}, che a un certo punto fu
``resettato'' perché non funzionava: un esempio positivo. Il relatore aggiunge
quello negativo di \emph{Age of Conan}, che non è mai riuscito a fare quel
reset iniziale (``è una schifezza, ve lo ripresentiamo quando è pronto''). Ci
vuole consapevolezza, ma anche i mezzi: un conto è Square Enix, che può
permetterselo, un altro è un'azienda che non può chiudere tutto per un anno
continuando a pagare gli stipendi.

Riassumendo: nei piccoli gruppi i problemi sono più di
\textbf{comunicazione e chiarezza} (obiettivi non chiari o che si spostano);
nelle grandi aziende di più di \textbf{persone che vanno e vengono, obiettivi
e tempistiche che cambiano}. In accademia le scadenze sono fisse (il gioco da
tavolo si consegna a fine anno scolastico, a giugno, e giugno resta); in
azienda può succedere che una feature prevista per la \emph{vertical slice}
di luglio venga anticipata a maggio per portarla alla Gamescom, e un sistema
perfetto per consegnare a giugno può non esserlo più per consegnare a maggio.

\subsection{L'ambiente determina la qualità del prodotto}

\paragraph{È vero che chi passa per Electronic Arts finisce per fare brutti
giochi?} Uno studente cita BioWare, che faceva ottimi giochi finché era
autonoma. Il problema non è dei developer, che non sono ``improvvisamente
diventati incapaci'': è evidente che il problema è altrove, e che
\textbf{l'ambiente in cui si fa un gioco determina molto spesso la qualità del
prodotto}.

Un esempio locale positivo: Tommaso Bonanni, conosciuto dal relatore insieme a
Diego e già docente dell'accademia, ha fondato a Roma Caracal Games Studio,
applicando fin dall'inizio questi concetti di cultura aziendale. L'azienda è
rimasta piccola, ma ha molto chiaro cosa vuole fare e come. Quando si
predispone fin dall'inizio l'azienda e il team a lavorare in un certo modo è
fattibilissimo; è molto più difficile quando si è EA, si acquisisce
un'azienda e le si dice ``da oggi si fa così''.

Il relatore ha vissuto la transizione da Blizzard ad Activision Blizzard (uno
dei motivi per cui se n'è andato, volontariamente). Ha visto la cultura
interna cambiare per gradi:
\begin{enumerate}
  \item ``facciamo giochi fighissimi, perché i giochi devono essere
        fighissimi'';
  \item ``facciamo giochi fighi così vendono di più'';
  \item ``vendiamo di più facendo giochi fighi'' (sembra uguale, ma è molto
        diverso);
  \item ``facciamo più soldi''.
\end{enumerate}
Lo si capiva dalle scelte di design, osservando quale fosse il peso più
grande sulla bilancia. Lo stesso varrà per gli studenti, nei loro giochi e nei
colloqui di lavoro: capire subito che cosa vuole l'altra parte.

\paragraph{Leader e visione.} Nei gruppi dell'accademia, dal secondo anno, c'è
una persona che assume il ruolo di \emph{lead} o \emph{scrum master}. È
diverso dall'essere un \emph{game director} (``il gioco si fa come dico io,
aiutatemi a farlo con la mia visione''), come avviene nei grandi studi: da lì
nascono i nomi nei titoli, come \emph{Sid Meier's Civilization}, che indicano
la persona che ha fatto quel gioco, oppure un puro branding, come \emph{Tom
Clancy's Rainbow Six}.

\paragraph{Inclusion come jolly.} A margine, uno studente osserva che in molte
aziende ``la risposta giusta è sempre \emph{inclusion}'', usata come un
jolly. È comunissimo, conferma il relatore: \emph{diversity and inclusion}
viene spesso usata per dire ``così includiamo più clienti'', ma non è quello
il suo punto (``altro discorso'').

% ---------------------------------------------------------------------
\section{Confrontarsi con chi è più bravo}
\label{sec:z1-rivalita}

Arrivato a Riot dopo sei anni a Blizzard, il relatore pensava di avere un buon
bagaglio di esperienza. Si è trovato invece ``come un fante buttato in mezzo
ai Navy SEAL'': persone che avevano superato 11, 12, 13 colloqui, con un
livello molto alto.

Individua subito un collega francese che, nella sua testa, diventa ``la
versione migliore di me'': stesso ruolo (community manager), esperienza
simile (arrivava da Ubisoft, come quasi tutti i francesi di Riot), ma molto
più spigliato, sempre con la battuta pronta, ``l'anima della festa''. Il
relatore invece è introverso: un workshop come questo per lui è
\emph{emotivamente drenante}, non ama parlare in pubblico pur amando
condividere le proprie esperienze, e non riesce a sostenere conversazioni
come ``cosa hai fatto questo weekend?''. Lavorando raramente insieme, aveva
sviluppato una rivalità ``da allenatore di Pokémon'' (lui era Ash, il
relatore l'antagonista).

Dopo sei mesi, a un evento di \emph{team building} aziendale (a Tenerife), il
suo capo, molto attento a queste dinamiche, li mette apposta nello stesso
gruppo. A fine serata, complice l'open bar, il relatore gli confessa di averlo
mal sopportato dal primo giorno per invidia della sua estroversione. Il
collega ride e risponde che pochi giorni prima aveva detto al capo la stessa
cosa al contrario: invidiava a Stefano il suo essere sempre organizzato, sul
pezzo, presentarsi nella sala riunioni giusta (mentre lui finiva in quella
sbagliata), arrivare con gli appunti, con i dati a supporto e le analisi di
mercato, mentre lui spesso non aveva nemmeno letto la mail del meeting.
Un'invidia in due direzioni, che li ha portati a lavorare molto insieme e a
migliorarsi a vicenda.

Le lezioni:
\begin{itemize}
  \item non bisogna demoralizzarsi vedendo qualcuno che fa lo stesso lavoro
        ed è più avanti sotto alcuni aspetti: spesso c'è un bilanciamento. I
        ``punti caratteristica'' sono limitati: se li ha messi in carisma,
        non li avrà messi in saggezza;
  \item conviene \textbf{chiedere a queste persone come mai} sono così brave
        in qualcosa, come si fa con un progetto andato bene: magari non lo si
        farà mai bene quanto loro, ma si comincerà ad avere il
        \emph{mindset} giusto; e viceversa, loro saranno spinti a chiedere a
        noi;
  \item spesso \textbf{le proprie doti le scopriamo dagli altri}. Il relatore
        non si era mai accorto di essere ``quello che resta calmo durante le
        emergenze'' (il lancio di \emph{Diablo~III}, quello di
        \emph{Draenor}, i problemi dal vivo alla Games Week) finché non
        gliel'hanno fatto notare tutti: per lui era la normalità. A chi pensa
        ``non sono capace, forse non è il mio percorso'', può darsi che
        semplicemente nessuno gli abbia ancora detto ``questa cosa la sai
        fare davvero bene, concentrati su questo''.
\end{itemize}

% ---------------------------------------------------------------------
\section{Domande sul lavoro nell'industria}
\label{sec:z1-industria}

Diverse domande degli studenti hanno riguardato le condizioni di lavoro nei
grandi studi internazionali. Le risposte sono raccolte qui per argomento.

\subsection{Presenza, remoto e trasferimenti}

\paragraph{Grandi studi: quasi solo in presenza.} Sia a Blizzard sia a Riot si
lavorava \textbf{solo in presenza}, senza smart working: sei anni a Parigi e
cinque a Dublino. Si parlava solo inglese: anche nei meeting tra soli
italiani si tendeva a comunicare e a scrivere in inglese, per correttezza
verso chi si fosse aggiunto e per lasciare una traccia comprensibile a tutti.
Il francese serviva per fare la spesa e sopravvivere, non in ufficio.
L'ambiente era estremamente eterogeneo (francesi, polacchi, spagnoli e
molti altri): durante il colloquio dal vivo a Dublino, mentre si preoccupava
di essersi vestito in modo adeguato, vide passare dietro la parete a vetro un
collega svedese con i dreadlock, a piedi scalzi, con una pianta sotto un
braccio e un aquilone sotto l'altro, e capì di essere ``a posto così''.
Un'esposizione a molte culture e modi di pensare che il relatore considera
bellissima.

Anche oggi, tranne alcuni studi particolari, la maggior parte del lavoro si
svolge di persona: amici a Paradox lavorano a Stoccolma, i pochi rimasti a
Blizzard si sono dovuti trasferire a Irvine. Per posizioni \emph{junior} o
\emph{entry level} i grandi studi raramente concedono il lavoro da casa; per
profili senior può esserci qualche eccezione (ad esempio nove mesi da remoto
e tre in sede).

\paragraph{Relocation.} Spesso tocca spostarsi, ma il relatore non ha mai
pagato un trasloco di tasca propria: queste aziende offrono la
\emph{relocation} (aiuto a trovare casa, ad aprire il conto in banca e
supporto pratico), perché è un investimento sulla persona.

\paragraph{Colloqui mentre si lavora altrove.} Per i colloqui in presenza
prendeva ferie, senza dire dove andava. Il consiglio generale: per fare
colloqui \textbf{prendete ferie o permessi}, non datevi malati, ``che non è un
bel modo di comportarsi''.

\paragraph{Remoto ``doppio''.} Un collega lavora per l'azienda coreana che
gestisce \emph{PUBG}, con ufficio ad Amsterdam, ma lavora da Milano: un
``doppio remoto''. Si trova, ma in genere un ambiente anglofono con sedi in
Europa chiederà di trasferirsi in Europa.

\paragraph{Famiglia e trasferimenti.} Il relatore ha avuto la fortuna di
sposare una persona che fa lo stesso lavoro (conosciuta lavorando su
\emph{Warhammer Online}, poi con lui a Blizzard e a Riot, nella
localizzazione e nel doppiaggio, ad esempio di \emph{World of Warcraft} in
italiano): avendo competenze rare, spesso venivano assunti entrambi, al
massimo a pochi mesi di distanza. Per i developer la soluzione frequente è
che la famiglia segua la persona assunta, perché lo stipendio basta per
tutti. Ma non sempre c'è equilibrio, e a volte uno dei due deve sacrificarsi:
sua moglie ha lasciato Blizzard per seguirlo a Dublino (per poi essere
assunta da Riot); lui ha lasciato Riot per tornare in Italia quando lei ha
trovato lavoro qui. Ha dovuto rinunciare più volte a lavorare in Islanda per
CCP (\emph{EVE Online}), perché non era una scelta solo sua. Si fanno
compromessi, ma si riesce a far quadrare.

\paragraph{Ci si sposta per ogni progetto?} No: si sono spostati
esclusivamente per cambiare azienda. Cambiando progetto nella stessa azienda
è raro, quasi mai (``non l'ho mai visto succedere''), che venga chiesto di
trasferirsi. In Francia stavano benissimo (vivevano fuori Parigi; ``i
problemi con i francesi sono in realtà problemi con i parigini''), parlavano
francese ed erano vicini alla doppia cittadinanza: se ne sono andati perché
non vedevano un futuro a Blizzard, per il modo in cui l'azienda trattava i
colleghi.

\paragraph{Il mercato oggi.} Non esiste un \emph{silver bullet} per trovare
lavoro domani dove si vuole. Con tutti i \emph{layoff} recenti ci sono
migliaia di esperti del settore che cercano lavoro: serve pazienza, ma è
fattibile.

\paragraph{Studi indie e remoto.} Negli studi più piccoli il remoto è più
diffuso: a Stormind circa 60 persone su 95 lavorano da remoto, ma lo studio
si è strutturato fin dall'inizio per farlo. Serve consapevolezza. Bending
Spoons (che sviluppa app per mobile, non giochi, ed è l'azienda dell'app
Immuni), pur essendo grandissima e con sede a Milano, lavora da remoto e
lascia l'ufficio a chi vuole: lo smart working è posizionato quasi come un
\emph{benefit}, ma permette anche di avere un ufficio da 300 persone invece
che da 700. Le sedi fisiche costano decine di milioni all'anno (Blizzard ha
chiuso le sedi di Amsterdam e di Parigi). Tuttavia, per lavori molto
verticali e specifici gli studi preferiscono avere le persone dal vivo; i
ruoli più spesso remoti sono quelli che non fanno sviluppo (customer service,
PR\dots).

Anche gli artisti lavorano spesso in presenza: l'accademia può dare solo una
\emph{simulazione} di quanto stretto sarà il contatto tra programmazione,
design e art. ``Sarà ancora di più.''

\paragraph{La comunicazione da remoto.} È più difficile. Pur essendo
estremamente introverso, il relatore \textbf{obbliga tutti a usare la webcam}
nei meeting in cui conta: non in un aggiornamento con novanta persone
collegate, ma certamente in un colloquio a due sulla performance dell'ultimo
anno. La componente fisica è una parte enorme della
comunicazione\footnote{In lezione si parla di ``quasi l'ottanta per cento'';
il dato va preso come indicativo, poiché le percentuali di questo tipo
circolano in forme diverse e sono spesso citate fuori contesto.}: anche se si
è timidi, accendere la webcam è l'approssimazione più vicina all'essere dal
vivo. Negli ambienti dove si lavorava senza webcam la situazione era
``esponenzialmente peggiore''.

\subsection{Differenze culturali nel team}

\paragraph{Gestire sensibilità diverse.} A Riot i team erano divisi per
\emph{lingua}, non per nazione: il team di lingua spagnola includeva anche un
messicano, quello di lingua russa includeva russi e ucraini, proprio nel
periodo della guerra in Crimea. Tra i colleghi non c'era alcun problema (erano
amici), ma ci si è posti la domanda: \emph{come comunichiamo con una community
che sarà divisa su questo tema?} Nei giochi online la divisione era evidente:
in \emph{Rocket League} c'era chi, con un compagno di squadra con la bandiera
avversaria, si faceva otto autogol pur di farlo perdere. Allo stesso modo,
con molti americani mandati a Dublino a coprire posizioni scoperte, c'erano le
elezioni (Clinton--Trump, Biden--Trump) e varie guerre: serviva sempre una
certa sensibilità.

Imparare a navigare queste situazioni lo ha reso un professionista migliore:
gli ha insegnato a fermarsi e chiedersi ``io la percepisco così; gli altri
dieci, con background, storie e lingue diverse, come la penseranno?'', a
ragionare più \emph{long term} e più \emph{big picture}, con più empatia.

\paragraph{Stili di comunicazione.} Un partecipante che ha lavorato nel Regno
Unito chiede dello stile comunicativo italiano, a volte frainteso. Uno dei più
grandi complimenti ricevuti dal relatore è stato ``non sembri italiano''
(``gesticoli poco per essere italiano'', ``parli bene inglese per essere
italiano''). Ha imparato ad \textbf{adattare la comunicazione alla persona}:
con spagnoli e francesi la battuta sulla nazione funziona, perché c'è un
``odio--amore'' condiviso; con russi e giapponesi no. Doveva adattare anche
l'umorismo. Un collega metà russo e metà messicano diceva di aver preso ``il
lato peggiore di entrambe le nazionalità'', ed è stato un ottimo maestro per
imparare a leggere culture diverse.

Non bisogna generalizzare, ma gli stereotipi possono essere un punto di
partenza ragionevole, \emph{sapendo che troverete sempre le eccezioni}. E
vanno letti correttamente: gli americani tendono a essere più estroversi e ad
avere sempre qualcosa da dire non perché ``sono fatti così'', ma perché
vengono educati così (all'università tra i corsi c'è il \emph{debate}, saper
argomentare in pubblico); un collega giapponese può essere più introverso per
ragioni culturali. Nel dubbio si parte da lì, poi si conosce la persona.

\subsection{Negoziare, discriminazione e stipendi}

\paragraph{La famiglia come leva.} Avere un partner riduce le opzioni, ma le
aziende lo sanno: uno spostamento è per due. Si può usare come leva per
negoziare (un compenso diverso, un supporto per il trasloco): se un'azienda ha
deciso di volere quella persona, 2, 3, 4 o 10.000 euro in più all'anno, per
un'azienda che fattura centinaia di milioni o miliardi, non cambiano nulla.

\paragraph{Essere sposati è un motivo per non essere assunti?} No. In alcuni
paesi è proprio illegale discriminare su questo: in Francia era illegale
chiedere CV con data di nascita, foto, sesso o religione, e chiedere se si
vogliono figli. Purtroppo, dai racconti dei suoi studenti di master (in
settori come \emph{beauty and wellness} e \emph{food and wine}), in piccole
realtà italiane queste domande (``hai 25 anni, hai intenzione di
sposarti?'') vengono ancora fatte, e servono a filtrare i candidati. Nelle
grandi aziende ci sono sistemi per evitarlo; anzi, molte danno \emph{più}
benefit a chi ha famiglia (``porta pure la famiglia, è tutto pagato''),
perché le persone competenti hanno spesso una certa età e quindi,
ragionevolmente, una famiglia.

\paragraph{Gli stipendi.} Dipendono molto da dove si vive. Lo stipendio a Riot
non era paragonabile a quello di Blizzard perché in Francia era tassato al
52\% e in Irlanda al 17\%\footnote{Aliquote come riportate in lezione,
riferite alla situazione personale del relatore.}: nessuno stipendio
competitivo può compensare il 35\% di differenza di tasse. Negli Stati Uniti,
con pochissima previdenza pubblica, gli stipendi sono molto più alti perché
poi ci si paga tutto (``se hai un raffreddore paghi tu''). Bisogna valutare il
potere d'acquisto di quei soldi nel paese in cui si andrà a vivere.

\subsection{Freelance e \emph{work for hire}}

All'inizio della sessione pomeridiana uno studente che lavora nell'audio
chiede quali siano le condizioni per chi lavora come \emph{freelance}.
\begin{itemize}
  \item Dipende molto dal progetto: spesso il contratto è a progetto, per
        portare a termine una parte o un intero progetto.
  \item Non tutte le aziende usano freelance: Riot e Blizzard lo facevano
        pochissimo e preferivano assumere (magari trattando poi la persona
        ``come un freelance'', lasciandola a casa a fine progetto).
  \item Spesso i lavori offerti come freelance sono in realtà
        \textbf{sostitutivi di un full time} (``freelance, ma sei ore al
        giorno'').
  \item Si può fare, ma bisogna essere molto chiari nel definire l'impegno
        (quante ore, come vengono tracciate): non tutti specificano cosa
        significa davvero ``freelance''. Lavorando così per due o tre aziende
        si può arrivare a un tempo pieno.
  \item È più comune il \emph{work for hire}: aziende che forniscono personale
        ad altre aziende; si è assunti da X ma si lavora ``embeddati'' presso
        Y (una sorta di subappalto). Anche lo studio del relatore lo fa.
  \item Il puro freelance di quattro mesi esiste, ma di solito per quei
        quattro mesi diventa un lavoro a tempo pieno; è raro ``quattro mesi,
        due ore al giorno''. Nelle grandi aziende è improbabile lavorare nello
        stesso mese per Ubisoft, Riot e Blizzard tre ore ciascuna; è più
        realistico lavorare per più startup indie italiane, che non possono
        permettersi un'assunzione ma possono pagare due ore al giorno.
  \item Molto passa dai contatti: trovare le aziende giuste e rendere esplicite
        le condizioni.
\end{itemize}

% ---------------------------------------------------------------------
\section{Il circolo virtuoso: feedback, miglioramento, fiducia}
\label{sec:z1-triangolo}

\subsection{Il triangolo}

L'obiettivo a cui il workshop vuole arrivare è un \textbf{triangolo che si
autoalimenta}:
\begin{enumerate}
  \item si dà un \textbf{feedback} con l'obiettivo di migliorare le
        performance (inclusi elementi caratteriali);
  \item la persona \textbf{migliora} e vede di averne beneficiato;
  \item questo genera \textbf{fiducia}, che rende più facile dare e ricevere
        il feedback successivo.
\end{enumerate}
Togliendo uno qualsiasi dei tre elementi, il circolo non è altrettanto forte.

In qualsiasi fase si parta, il rispetto per i colleghi è inizialmente dettato
dall'\emph{educazione}, non dalla fiducia. La fiducia arriva quando i colleghi
cominciano a dire cose giuste e sensate e a \textbf{mantenere la parola}.
Dare e ricevere feedback è il modo in cui, in un certo senso,
``inganniamo il cervello'' per creare una \emph{case history} di fiducia: si
ha motivo di fidarsi di più, di ascoltare di più, di essere più ricettivi. E
si diventa sempre più bravi ad accettare il feedback e ad agire di
conseguenza.

\subsection{Accettare un feedback non vuol dire fare come dice l'altro}

Dieci anni fa il relatore non era ostile ai feedback, ma molto
\emph{difensivo}: ``sì, sì, va bene, ora ci guardo''; oppure, a chi gli diceva
``avresti dovuto farla così'', rispondeva spiegando perché l'aveva fatta in
quel modo. Oggi ha imparato a dire ``ok, grazie, ora ci ragiono'': non deve
difendere la propria posizione da chi viene solo a dire ``hai pensato a
questa cosa?''. Prima lo viveva come un attacco (``mi viene a dire che non so
fare il mio lavoro''), ora come un ``ottimo punto di vista, lo prendo in
considerazione''.

Con una postilla fondamentale: \textbf{accettare il feedback non significa
fare come dice l'altro}. ``Mi hai detto che preferisci questo mob verde. Ho
pensato anche a farlo verde? Sì. Penso comunque che giallo sia meglio? Sì, lo
faccio giallo.'' Il relatore accetta \emph{sempre} il feedback, valido o no:
garantisce di tenerlo a mente quando lavora. Che poi faccia come gli viene
detto è un altro discorso.

\subsection{Fiducia e scadenze}

La fiducia funziona allo stesso modo con gli impegni. Se si chiede a qualcuno
una cosa per venerdì e la consegna venerdì, la volta dopo gliela si chiederà
per venerdì. Se invece arriva il giovedì della settimana dopo, la volta dopo
gli si darà una \textbf{deadline anticipata}, come con chi arriva sempre in
ritardo agli appuntamenti (``se ti dico alle 7 arrivi alle 7 e mezza; la
prossima volta ti dico alle 6 e mezza'').

Al lavoro questo è dannosissimo, perché crea \textbf{deadline fittizie} che
non sono mai quelle vere, e nascono discussioni del tipo ``perché devo farlo
in due giorni se si fa in cinque?'' (``sì, se tu lo facessi in cinque''). La
conversazione giusta è un'altra: \emph{``che problemi hai avuto nel
consegnare venerdì?'' ``Questo, questo e questo.'' ``Cosa posso fare per
evitare che succeda?''}

Il relatore non vuole mai sentirsi dire ``non sono riuscito a farlo''.
Preferisce di gran lunga sentirsi dire, tre giorni prima, \textbf{``potrei non
riuscire a farlo, per questi motivi''}. Chi ha la visione d'insieme può dare
più budget, più persone, spostare la deadline, cambiare lo scopo del progetto:
può fare qualsiasi cosa per aiutare. Ma se lo si scopre il venerdì, cinque
minuti prima della deadline, non può fare più niente. Il messaggio è:
\textbf{``aiutami a poterti aiutare''}. Non penserà mai male di chi dice ``non
ce la faccio per questi motivi'', perché non sta dicendo ``non ho voglia'': sta
indicando i problemi, e lavorare su quei problemi è compito del manager. Non
può invece fidarsi di chi fino al giorno prima dice ``tranquillo, ce la
faccio'' e il giorno dopo ``non ce l'ho fatta''.

Questo è fondamentale per chi lavora in programmazione, design o art, perché
ci saranno sempre deadline: bisogna sempre \textbf{lavorare sulle aspettative},
dire al proprio manager (o al proprio report) cosa è lecito aspettarsi. Non
bisogna fare come i pinguini di \emph{Madagascar} (``sorridete e salutate, va
tutto bene''): se c'è un problema, va detto, altrimenti si peggiora solo la
propria situazione.

Da italiani si tende a non voler mai alzare la mano per chiedere aiuto
(``no, no, ce la facciamo, lo faccio io''). A volte non ce la si può fare, non
perché non si è capaci, ma perché vari fattori lo impediscono. Accettarlo è
importante: \textbf{è molto meglio avvisare e chiedere una mano che dire
``non ce l'ho fatta''}.

\subsection{Domande sul dare e ricevere feedback}

\paragraph{``Ne parliamo a fare, se poi fai come vuoi?''} Uno studente (Cristian)
riceve feedback che, dopo averli analizzati, non ritiene congrui; quando non
li segue gli viene detto ``ne parliamo a fare?''. C'è un problema di
\textbf{aspettative}: l'idea che accettare un feedback voglia dire fare come
dice l'altro. Accettarlo significa prendere in considerazione ciò che è stato
detto. Chi decide diversamente ha però un \textbf{onere}: argomentare perché
si fa comunque a modo proprio (``mi hai detto verde, la facciamo gialla
perché\dots''). Poi bisogna solo aspettare i risultati per vedere chi aveva
ragione.

La decisione spetta sempre \textbf{a chi ha la responsabilità di quella
decisione}. Se il game director deve scegliere tra A e B, anche se il
feedback dice B, la scelta è sua perché la responsabilità è sua. Non sarebbe
corretto essere ritenuti responsabili di una scelta delegata a chi dà
feedback. La risposta può essere semplicemente: ``l'ho tenuto in
considerazione, ma visto che ne sono responsabile io e le conseguenze sono
mie, ho scelto di fare X''.

\paragraph{E se l'altro fa ripicca?} Un partecipante chiede di chi, vedendo
il proprio feedback non accolto, si vendica ignorando a sua volta feedback
buoni. Per il relatore è un problema caratteriale: chi si comporta così ``è
già sulla strada per uscire''. Non c'è spazio per l'ego in team di questo
tipo. Il fenomeno nasce dall'idea che \emph{il feedback negativo faccia fare
brutta figura}, che il fatto che si potesse fare meglio tolga qualcosa alle
proprie capacità.

Non è così: \textbf{il feedback non è un giudizio}. Si può passare un anno a
Riot ricevendo solo feedback negativi e avere una valutazione annuale
stellare, perché ricevere molti feedback negativi significa essere una persona
approcciabile, che agisce sul feedback e invita gli altri a darlo. Il relatore
ha visto valutazioni molto peggiori per chi riceveva solo feedback positivi,
perché la gente aveva paura di dargliene di negativi: ma solo ``pacche sulle
spalle'' non fanno migliorare. Il problema va affrontato parlando con la
persona; a volte le conversazioni iniziano con \textbf{``ti do un feedback sul
feedback''}, cioè sul modo in cui l'altro ha dato il suo.

Esempi dalla sua esperienza:
\begin{itemize}
  \item lui chiede spesso di \textbf{non ricevere feedback senza esempi
        tangibili}, perché fa fatica ad astrarre: ``sei stato maleducato'' non
        significa nulla, ``ieri al meeting mi hai risposto male su questa
        cosa'' sì;
  \item ad altri serve l'opposto: ha ricevuto il feedback ``quando devi dirmi
        che ho sbagliato, puoi girarci un po' intorno e darmi del contesto?'',
        perché il suo stile diretto (``questa cosa non ha funzionato, dovremmo
        fare così'') feriva. Ha adattato il suo modo di dare feedback: con
        alcuni è diretto, con altri gira intorno. L'importante è saperlo.
\end{itemize}

\subsection{Non portare la soluzione}

Un consiglio centrale: \textbf{quando si dà un feedback, non si dà una
soluzione}. Non si dice ``io farei così'', ma ``secondo me questa cosa si
potrebbe fare diversamente, per questo e questo motivo: come potremmo farla
diversamente?''. È come chi smette di fumare perché vuole e chi smette perché
glielo dicono moglie e figli: se non si è convinti, spesso non si smette.

Quando il relatore diceva ``c'è questo problema, io farei così'', il meeting
successivo iniziava inevitabilmente con ``ci ho provato come hai detto tu, con
me non ha funzionato''. Sempre. Perché la soluzione era sua. La
responsabilità della soluzione deve essere invece \textbf{congiunta}: ogni
feedback deve creare un piccolo team che ha come obiettivo risolvere quel
problema. Questo toglie l'antagonismo (non ``ti dico che non sai fare il tuo
lavoro'', ma ``ti voglio aiutare a farlo meglio, lavoriamoci insieme'') e
sposta il focus: il problema non sei tu, è ciò che sta accadendo.

\paragraph{E se sono i piani alti a tenere all'oscuro dei problemi?} È una
questione di cultura aziendale. Il relatore fa sempre presente al proprio capo
il concetto di \textbf{service level agreement} (che verrà ripreso nella
seconda giornata): ``se mi dai questi input, ti garantisco questo output in
questi tempi''. Se il capo chiede delle cose senza dire budget, tempi e
perché, non può aspettarsi lo stesso risultato. È una comunicazione a due
vie: si dice al manager cosa serve per fare il lavoro, e lui deve fornirlo o
accettare che il risultato non sarà quello voluto.

\subsection{Amici, parenti e parole proibite}

\paragraph{``Non si lavora con amici e parenti''?} Per sua definizione
personale il relatore ha pochissimi amici (``stanno sulle dita di una mano'')
e una soglia molto alta per considerare amico un collega. Allo stesso tempo,
in quindici anni e con centinaia di colleghi, ha avuto solo \textbf{due}
esperienze in cui non è riuscito a lavorare con qualcuno: non serve essere
amici per far funzionare le cose. Con gli amici, anzi, gli viene più facile:
se dice a un amico che una cosa ``fa schifo'', l'amico non pensa ``come si
permette'', ma ``se perfino Stefano, che è mio amico, me lo dice, allora è
proprio così''. Diventa un rafforzativo, grazie alla fiducia costruita prima.
Il problema del lavorare con amici e parenti nasce dal \emph{non voler} dire
all'amico che sta sbagliando; se invece la confidenza permette di dirlo, si
può saltare la ``confezione'' del messaggio e passare direttamente al
contenuto.

\paragraph{Le frasi peggiori.} Alla domanda se esistano parole peggiori di
``ottimo lavoro'', il relatore indica la sua classifica: ``ottimo lavoro''
(detto in modo vuoto)\footnote{Il passaggio della trascrizione è
frammentario: si intende verosimilmente il complimento generico e di
circostanza, che non dice nulla di utile.}; poi la frase che lo ``fa uscire di
testa'', \textbf{``abbiamo sempre fatto così''} (``e infatti avete chiuso'',
come Blockbuster); infine \textbf{``siamo una famiglia''}, riferito al posto di
lavoro. Un'azienda non è mai una famiglia, perché non ci si può aspettare lo
stesso impegno da chi la possiede e da chi ci lavora: non c'è un
corrispettivo proporzionale tra impegno e ritorno. Per quanto inclusiva,
flessibile e razionale, non sarà mai una famiglia.

\subsection{Ego, muri e il peggior modo di scegliere un leader}

\paragraph{Chi mette davanti l'ego e alza muri.} Esistono, ma di solito durano
pochissimo in azienda. Non si tratta di avere un ego forte, ma di \textbf{non
saper ricevere feedback}. Spesso sono persone bravissime nel loro lavoro che
poi si chiedono perché tutti intorno diventino più bravi di loro: perché gli
altri accettano feedback e hanno più fonti da cui attingere, mentre loro si
sono rintanati nella loro torre.

Molti manager tendono ad \textbf{assumere persone un filo peggiori di loro},
per non fare brutta figura. Il relatore assume solo persone più brave di lui,
``perché non è una gara'': ha lavorato con molti community manager più bravi
di lui come community manager, ma lui era più bravo a fare il manager. Non
viene valutato su quanto sa fare il community manager, ma su quanto il suo
team sa farlo. Un partecipante chiede se la paura di avere qualcuno di
migliore esista anche in ambito accademico. Assolutamente sì: al relatore è
capitato di non essere preso come docente in un corso di formazione sui
videogiochi perché il responsabile aveva percepito che ne sapesse più di
lui\footnote{La trascrizione dice ``che non sapevo più di lui'', ma il senso
del passaggio richiede il contrario.}. ``Ma perché te ne importa? Tu sei quello che ha portato chi
sa fare il suo lavoro: ben venga.'' O si ha questa cultura, o si ha paura per
il proprio posto.

\paragraph{Promuovere il migliore.} Quando bisogna scegliere un leader per un
team (programmazione, design, art), molto spesso si sceglie \textbf{il migliore
tecnico}: il miglior programmatore, il miglior designer. È spesso la scelta
sbagliata: è come prendere un idraulico per fare l'elettricista. Fare il
\emph{people manager} non è come fare il designer o il programmatore; si può
essere un ottimo programmatore e un pessimo manager, e viceversa. Le aziende
l'hanno capito e lo fanno meno di prima. A Blizzard c'era un grave problema di
\emph{middle management}: molte persone in ruoli importanti erano state
promosse cercando le skill sbagliate, erano frustrate dal non saper fare il
proprio lavoro e schiacciavano chi stava sotto per paura di fare brutta
figura.

Un partecipante cita una teoria simile: una persona viene promossa finché non
raggiunge \emph{il proprio livello di incompetenza}, dove si blocca\footnote{È
il \emph{principio di Peter} (Laurence J.\ Peter, 1969), citato in lezione
come ``la teoria di Peter''.}. Un altro cita Michael Scott di \emph{The
Office}, su cui il relatore tornerà (sezione~\ref{sec:z1-quando}).

% ---------------------------------------------------------------------
\section{Perché la cultura conta: retention, costi e fiducia}
\label{sec:z1-cultura}

\subsection{L'ambiente di lavoro e la retention}

Perché tutto questo è così importante? Perché l'ambiente di lavoro conta
moltissimo per la \textbf{retention}, cioè per trattenere le persone e farle
lavorare in modo efficiente. Se contasse solo lo stipendio, molti non
lascerebbero mai Riot o Blizzard; e più si va avanti con l'età, più si impara
a dare valore anche al resto.

Inoltre, secondo gli studi citati nelle slide, chi vive il proprio lavoro in
modo negativo è \textbf{meno efficiente nel problem solving}: l'energia
mentale va verso ``come faccio a non farmi dare la colpa'' o ``come faccio a
fare bella figura'', non verso la soluzione del problema. In ambienti negativi,
dove ci si deve preoccupare della propria ``sopravvivenza'', si lavora
peggio.

\subsection{Il costo di chi se ne va}

Ogni volta che qualcuno si dimette o viene licenziato, il costo di trovare
una nuova persona, fare i colloqui, formarla e farle attraversare forming,
storming, norming\dots\ è circa \textbf{un terzo del suo stipendio}. Per questo
le aziende tendono a fare più \emph{retention} che \emph{acquisition}, almeno
con le figure considerate utili alla produzione.

Allo stesso tempo si vedono sempre più spesso \emph{layoff} e \emph{hiring
wave}: si licenziano 3.000 persone e sei mesi dopo se ne riassumono 2.500 o
3.500, ma tutte più junior. Il relatore ne ha scritto in un articolo su
Eurogamer.it, in cui spiega anche alcuni ``trucchi del mestiere'': perché i
licenziamenti avvengono sempre in certi periodi e sempre ``a gruppi'' (prima
Bethesda, poi Blizzard, poi un'altra\dots).

\paragraph{L'offerta di Riot.} A Riot esisteva una pratica (di cui il relatore
non ricorda il nome) per cui, firmato il contratto, si aveva \textbf{un anno
di tempo per dire ``non mi piace, non me la sento, me ne vado''} e ricevere
subito come buonuscita un anno di stipendio. Avevano stimato che fosse molto
meglio pagare subito un anno a qualcuno e lasciarlo andare che tenerlo due,
tre, quattro anni come \emph{underperformer}: la responsabilità della scelta
di assumere era dell'azienda\footnote{La pratica, nota come ``The Offer'', è
descritta in fonti pubbliche con condizioni diverse (un periodo più breve e
una somma pari a una frazione dello stipendio annuo). Si riportano le
condizioni come raccontate in lezione; andrebbero verificate.}. L'hanno usata
pochissimi, ma il solo fatto che esistesse faceva lavorare meglio anche lui da
manager, perché sapeva che chi non si trovava bene poteva andarsene con un
anno di tempo per trovare un altro lavoro.

\paragraph{Il ricircolo.} C'è un forte ricircolo nell'industria: molti ex
colleghi passano da Ubisoft a Paradox, da Paradox a Funcom, da Funcom a
Blizzard. Il relatore ha scelto di uscirne lavorando solo per aziende
\textbf{indipendenti}: Stormind non è un indie di cinque persone, ma è
indipendente, e non deve accontentare azionisti. Quando vede un'azienda con
utili miliardari licenziare migliaia di persone si chiede ``dov'è la linea''
(\emph{when enough is enough}): con azionisti che vogliono un ritorno, se
quello è l'obiettivo stanno facendo la cosa giusta, ma lui non condivide
l'obiettivo\footnote{In lezione si parla di ``6,5 miliardi di utile'' e di
``3.000 persone licenziate''; le cifre sono riportate come dette e non sono
state verificate.}.

\paragraph{L'open space vuoto.} A Riot, durante un'ondata di licenziamenti,
del suo team di trenta persone è rimasto solo lui: ``teniamo te come
coordinatore, poi ricostruirai il team''. Ha rifiutato di ricostruire un team
cancellato ``con un colpo di spugna'' e ha deciso che non avrebbe più lavorato
per una grande azienda. Nessuno aveva pensato di spostarlo, e ha lavorato per
tre mesi in un open space di cento metri quadri con la sua scrivania e altre
ventinove vuote. Su una lavagna aveva scritto i nomi dei suoi ragazzi e li
cancellava man mano che, anche grazie ai suoi \emph{referral}, trovavano
lavoro. Cancellato l'ultimo nome, si è dimesso: per lui il lavoro era finito.

\paragraph{Team che si dissolvono.} Uno studente cita un articolo su Larian
Studios, il cui management è contrario a ricostruire ogni volta il team a fine
progetto. È lo stesso problema: si pensa ai ``giochi Blizzard'' o ai ``giochi
BioWare'', ma le persone che li hanno fatti non ci sono più; è rimasto solo il
brand, e non ha senso avere le stesse aspettative da un gioco fatto da persone
diverse. E anche chi sopravvive ai layoff vive pensando che il prossimo
potrebbe essere lui: alcune persone, licenziate nell'ultima ondata a
Blizzard, gli hanno detto ``erano due anni che lo sapevo''. Vivevano a Cork,
``un paesotto irlandese dove o vai per lavorare a Blizzard o non ci vai'',
avevano comprato casa: non ci si può aspettare chissà che morale.

\subsection{La cultura è forte quando si lavora per gli altri}

\begin{quote}
  \emph{La cultura è forte quando le persone lavorano per gli altri e a favore
  degli altri, e non per se stesse.}
\end{quote}

Se si dà feedback per far migliorare l'altro e non per imporre il proprio
modo, se si aiuta per il bene del progetto e non per fare bella figura, la
cultura si rafforza. È dai \textbf{piccoli gesti} che emerge la cultura
aziendale: ad esempio, dal non dover chiedere ``perché mi stai aiutando?''.

\paragraph{L'incognita.} \textbf{L'incognita è spesso dannosa quanto le
cattive intenzioni.} In un esperimento (``mi pare a Stoccolma, cinque o sei
anni fa'') si distribuivano per strada banconote da 500 euro, e la maggior
parte delle persone le rifiutava chiedendosi ``dov'è la fregatura?''
\footnote{L'esperimento è riportato in modo vago in lezione e non è stato
possibile identificarlo con certezza.}. Lo stesso accadeva al relatore,
arrivato a PG e a Stormind da Riot: a ``vedo che sei in difficoltà, posso
aiutarti?'' la risposta era sempre ``perché?''. Alla fine rispondeva: ``perché
quando sarò in difficoltà io vorrei che tu mi aiutassi''.

Quando si comincia a lavorare per gli altri e si smette di pensare ``così avrò
la promozione, così avrò l'aumento'', promozioni, aumenti e rispetto arrivano
da soli. A Blizzard aveva fallito tre volte i colloqui per essere promosso a
manager (inteso come coordinatore del team); quando ha deciso di smettere di
pensarci e di fare solo il bene del proprio lavoro, la volta dopo è stato
promosso. Funziona così: chi aiuta i colleghi e fa andare bene il progetto
viene notato (``questa persona fa migliorare cinque persone, non solo se
stessa: vale di più'').

\subsection{Il lavoro è valutato come opera d'insieme}

Il relatore riconosce che buona parte del suo successo è venuto dal fatto che
il suo team lavorava bene, non dalla sua bravura personale. Anche i progetti
degli studenti, presentati a Diego, a Enrico, ad Andrea, verranno valutati
\textbf{come opera d'insieme}. Nessuno dice ``il \emph{combat system} di questo
gioco l'ha fatto X, che bravo'': si dice ``BioWare ha fatto un bel gioco'',
``Blizzard ha fatto un bel gioco''.

Lo ha replicato anche lato HR. In molte aziende ci sono \textbf{bonus
personali} legati a obiettivi (se il gioco vende X, se la feature è pronta in
tempo). Il relatore li ha dimezzati: metà restano personali, metà diventano
\textbf{bonus di team}, che tutti ricevono se il team raggiunge gli obiettivi.
Un incentivo concreto a collaborare: non solo si dice che il lavoro di gruppo è
importante, lo si dimostra dandogli più valore economico. Aveva proposto il
100\% di bonus di team, ma gli è stato concesso il 50 e 50; ha funzionato
comunque, anche se sarebbe stato ancora più potente. Ovviamente, alla
domanda ``perché mi aiuti?'' si può anche rispondere ``perché voglio il
bonus'': un aiuto dato ``per il motivo sbagliato'' era un compromesso che era
disposto ad accettare.

\subsection{Il coraggio di chiedere aiuto}

Un partecipante osserva che aiutarsi dovrebbe essere la norma. Il modello
diffuso, invece, è ``devi essere forte e indipendente, non chiedere mai una
mano'', e vale in tutte le direzioni (dal maschilismo alla ``donna forte e
indipendente che non ha bisogno di nessuno''): \textbf{non ammettere di avere
bisogno di aiuto è sempre uno dei motivi per cui le aziende vanno male}, che si
tratti di uomini, donne o persone non binarie. Il problema è non avere la
sicurezza di poter alzare la mano.

Il relatore ha cercato di togliere ogni conseguenza dall'aver alzato la mano:
se si \emph{depenalizza} qualcosa, le persone non sono più disincentivate a
farla. Ha detto espressamente che non avrebbe mai valutato nessuno per il
numero di volte in cui chiedeva aiuto o per il numero di feedback negativi
ricevuti, e questo ha funzionato.

\subsection{Perché si dà un feedback}

\emph{Perché avete dato il vostro ultimo feedback? Che cosa volevate ottenere?}
Le risposte della classe:
\begin{itemize}
  \item per \textbf{riconoscere il valore} di un collega, anche verso il
        manager, perché qualcuno che non aveva visibilità vedesse quanto
        quella persona fosse efficace;
  \item per \textbf{spingere una persona a superarsi};
  \item per aiutare un collega a rendere il prodotto all'altezza di ciò che
        voleva rappresentare;
  \item per \textbf{genuinità}: c'era qualcosa che poteva essere fatto meglio.
\end{itemize}
Il relatore aggiunge che la forma cambia col tempo: all'inizio era ``devo
darti un feedback, hai dieci minuti?''; ora è un messaggio informale su
Discord, perché ha costruito un rapporto di fiducia e sa che non verrà letto
come ``l'HR manager mi scrive, sono nei guai?'', ma come ``Stefano è così, mi
dà feedback, resterà tra noi e lo risolveremo insieme''.

Un altro motivo è \textbf{risolvere incompatibilità caratteriali}: le
conversazioni più difficili, del tipo ``io e te non stiamo funzionando, cosa
possiamo fare diversamente?''. Sono importanti perché, se due responsabili non
vanno d'accordo, il team se ne accorge, comincia a raccontarsi storie (``si
odiano, ci saranno problemi'') e, per difendere il proprio manager, comincia a
giudicare male il lavoro dell'altro gruppo. È la dinamica da fermare:
\textbf{valutare gli altri solo per difendere la propria fazione}.

Dare e ricevere feedback sono \textbf{due muscoli da allenare} (il relatore è
sempre stato più bravo a darlo che a riceverlo), e si impara col tempo. È come
fare colloqui: fare colloqui è una skill a parte, che non indica quanto si è
bravi nel proprio lavoro. Essere scartati non significa non saper fare quel
lavoro: spesso significa aver fatto male il colloquio o che è stato preferito
qualcun altro.

\paragraph{Feedback pubblici e privati.} Uno studente racconta di aver dato un
feedback negativo a un superiore, senza un motivo particolare, per dimostrare
a sé e ai colleghi che non si viene licenziati per questo. Il relatore coglie
l'occasione per precisare che \textbf{i feedback negativi si danno
individualmente}: lui dà moltissimi feedback di cui il resto dell'azienda non
sa nulla. Reagirebbe molto male se, in un meeting, qualcuno alzasse la mano per
dire ``comunque questa cosa non l'hai fatta bene'': quello non è feedback, è
discussione, accettabile solo se è un brainstorming e se l'opinione è stata
chiesta. Alcune aziende hanno piattaforme pubbliche di feedback, ma spesso
servono solo ad applaudire; lo stesso accade nelle aziende che
\emph{obbligano} a lasciare un certo numero di feedback a settimana, che
diventano ``lui è bravo, lei è competente, anche questa settimana l'abbiamo
fatto''. Bisognerebbe chiedere a chi le adotta cosa pensa di ottenere (forse
trasparenza), ma difficilmente una persona ``attaccata'' pubblicamente reagisce
bene.

% ---------------------------------------------------------------------
\section{Among Us: la caccia al traditore}
\label{sec:z1-among-us}

Prima della pausa pranzo il relatore mostra tre film di fantascienza di
qualche decennio fa e chiede che cosa li accomuni\footnote{I titoli non sono
leggibili nella trascrizione: viene citato \emph{Event Horizon} da uno
studente, e il film ``in mezzo'' è indicato come ``l'originale'',
verosimilmente \emph{La cosa} (\emph{The Thing}).}. In ognuno c'è un gruppo di
persone con un obiettivo e la presenza, reale o percepita, di qualcuno che
rema contro: \textbf{l'energia del gruppo viene spesa per trovare il
responsabile}.

\paragraph{L'esperimento degli alieni.} Nei master in presenza il relatore
propone un esercizio. Distribuisce segretamente cartellini plastificati con
scritto ``essere umano'' o ``alieno''; le carte alieno contengono istruzioni di
comportamento con cui gli alieni, che cercano di fingersi umani, si
tradirebbero (dare sempre ragione a chi ha parlato prima, finire ogni frase
con una domanda, parlare in un certo modo). Lo scopo: in dieci minuti formare
gruppi, e \textbf{vince il gruppo più numeroso senza alieni}; basta un alieno
per perdere. Si formano sempre gruppi piccoli (due, cinque, sei persone),
perché la tendenza è ``mi fido al massimo di voi due''.

Ma vincono i gruppi più grandi, per un motivo semplice: il relatore
\textbf{non distribuisce nessuna carta alieno}. Mostra a tutti entrambe le
carte, ma consegna solo carte ``umano''. Vuole simulare ciò che accade in
ufficio: la tendenza a percepire sempre qualcuno che rema contro (i manager,
l'altro team, quel collega), quando in realtà non c'è nessuno, ``è tutto nella
nostra testa''. Quando chiede le motivazioni, emergono sospetti
\emph{fabbricati}: ``faceva sempre domande'', ``parlava solo all'infinito''
(``no, parla così e basta''). Ci si fabbrica da soli i motivi per non fidarsi
degli altri, ed è comunissimo.

\paragraph{Among Us.} L'autosabotaggio è, per tornare ad \emph{Among Us},
\textbf{concentrarsi sul cercare cadaveri e traditori invece di fare le
task}: non si vince perché non si sono completate le task, mentre si cercava il
colpevole. Succede tantissimo negli uffici, ed è il motivo per cui l'80\% dei
problemi che il relatore risolveva erano di comunicazione. Entrando nei team
sentiva immancabilmente: ``facciamo il possibile, ma questa lotta continua con
l'altro gruppo è estenuante''. ``Perché c'è questo problema?'' ``Non lo so,
c'era già quando sono arrivato.'' Un paradosso: c'è un problema, nessuno sa
dire quale sia, ma c'è sempre stato e quindi lo si accetta.

\paragraph{L'esperimento delle scimmie.} È l'esempio, ``finto o reale'', delle
scimmie in gabbia\footnote{È il noto racconto delle ``cinque scimmie'', spesso
presentato come esperimento scientifico ma probabilmente apocrifo (lo stesso
relatore lo definisce ``finto barra reale''); nella trascrizione il numero
delle scimmie varia.}. In una gabbia con alcune scimmie pende una banana
elettrificata: chi la tocca fa prendere la scossa a tutte. Le scimmie imparano
e, quando una si avvicina alla banana, le altre la picchiano. Poi si
sostituisce una scimmia alla volta con una nuova: la nuova va verso la banana,
le altre la picchiano, e smette. Si continua finché nella gabbia ci sono solo
scimmie nuove, poi si toglie la corrente. Nessuna scimmia tocca la banana:
nessuna ha mai preso la scossa, ma tutte hanno imparato che chi tocca la
banana prende le botte.

Nei team è lo stesso: il motivo per cui si ``odia'' l'altro team si
trasforma, ma un vero motivo non c'è. Ne derivano la sfiducia, il non alzare
mai la mano per chiedere aiuto, il non dire mai ``bel lavoro'' agli altri
perché non sia mai che qualcuno ti veda fare un complimento ``a quelli''.
Risolvere queste dinamiche è lungo e tedioso, e sempre pieno di ``non mi
ricordo'', ``io non c'ero'' e ``abbiamo sempre fatto così'', la cosa più
frustrante per chi cerca di risolverle.

% ---------------------------------------------------------------------
\section{Il principio del default to trust}
\label{sec:z1-default-trust}

\subsection{La definizione}

Il relatore indica questa come la slide più importante della presentazione.
\textbf{Default to trust} è uno dei valori aziendali di Riot Games, che lui ha
tradotto, adattato e ``esportato'':

\begin{quote}
  \emph{Approccio ogni situazione di conflitto partendo dal presupposto che una
  persona abbia agito in un certo modo perché, con le informazioni a sua
  disposizione in quel momento, ha pensato che fosse la soluzione migliore, e
  non perché sia un'incapace o volesse danneggiare il progetto
  intenzionalmente.}
\end{quote}

In altre parole: \textbf{di default, mi fido}.

\subsection{L'esempio del ritardo al daily stand-up}

A Riot il relatore aveva un \emph{daily stand-up} dal vivo alle 15 con il suo
team di venti persone. Una ragazza arrivava regolarmente alle 15:15, per una,
due, tre settimane. Alla quarta ci ha parlato. \emph{Come iniziereste la
conversazione?} Le proposte della classe:
\begin{itemize}
  \item sottolineare l'importanza di esserci dall'inizio, perché potrebbe
        perdersi qualcosa che le interessa, e per rispetto dei colleghi
        presenti;
  \item ``devi arrivare in orario'', spiegando perché è importante;
  \item \textbf{chiedere perché}.
\end{itemize}

Il relatore osserva, anche dalle sue consulenze fuori dai videogiochi, che
\textbf{più si sale con la seniority}, più la risposta è ``devi smettere di
arrivare in ritardo'', spesso con un \emph{or else} finale (``altrimenti
succede quest'altra cosa''). Più si scende con l'età, e quindi con meno
abitudini consolidate, più la prima domanda diventa ``dimmi perché arrivi in
ritardo''. Già spiegare perché è importante è ``miglia avanti'' rispetto al
semplice ``smetti''.

Partire dal perché gli ha permesso di scoprire la causa: il meeting precedente
della ragazza, in un'altra zona del campus a un quarto d'ora di distanza,
finiva sempre tardi, e lei doveva attraversare il campus di corsa. Non aveva
la confidenza per dire all'altro manager, che insisteva perché si restasse
fino alla fine, che doveva andare; sacrificava quindi il meeting del relatore,
sapendo che era il più comprensivo dei due.

Il relatore si è offerto di \textbf{mediare}: ha detto all'altro manager che
o anticipava l'inizio del suo meeting, così da finire comunque in tempo,
oppure lasciava uscire la ragazza, perché quel meeting gli serviva. C'è anche
una questione di \textbf{economia di scala}, spesso sottovalutata: se venti
persone aspettano qualcuno per dieci minuti, si sono sprecati dieci minuti per
venti. Per lo stesso motivo il workshop è iniziato alle 10:05: iniziare alle
10:10 avrebbe significato far aspettare cinquanta persone cinque minuti in
più. A Riot c'era una scrivania vuota con un pupazzo gonfiabile, \textbf{Mr.\
Late}: sommando tutti i minuti di ritardo a meeting, call e conversazioni, in
un anno veniva fuori praticamente \emph{una persona a tempo pieno}. Serviva a
ricordare che anche cinque minuti contano, e che bisogna rispettare gli altri
arrivando in orario o almeno avvisando.

\paragraph{Gli effetti.} La conversazione era avvenuta nel primissimo periodo a
Riot, quando il team non aveva ancora motivo di fidarsi di lui. Risolvendo il
problema ha eliminato anche tutto lo strato di micro-dinamiche spiacevoli che
un manager impara a vedere (due colleghi che si guardano con la risatina
quando lei entra in ritardo). La ragazza ha cominciato a entrare tranquilla,
con una certa leggerezza: non le faceva piacere arrivare in ritardo (``ci sono
anche persone che lo fanno apposta, e quello è un problema da affrontare
diversamente''). Soprattutto, con i colleghi non parlava di lui come del capo
che minacciava, ma come di chi le aveva risolto un problema: era diventato
un \textbf{alleato}. La settimana dopo altre tre persone sono andate da lui
con un problema da risolvere. Ripetuto cinque, sei, sette, dodici volte, in un
paio di mesi poteva dire quasi qualsiasi cosa nei meeting sapendo che
richieste e critiche sarebbero state percepite come ``me lo chiede perché vuole
darmi una mano'', non per fare bella figura o per presentare i successi del
team ai suoi superiori.

\subsection{Condividere i successi, schermare i fallimenti}

Un altro elemento fondamentale (ripreso nella seconda giornata): il manager
\textbf{condivide i successi e protegge il team dai fallimenti}. ``Se va male
qualcosa è colpa mia; se va bene qualcosa, voglio che ne parliate voi.'' Il
lavoro del manager non dovrebbe essere valutato su ciò che presenta del lavoro
del team: tutti sanno comunque chi è il manager. Tanto vale lasciar presentare
il team, che potrà dire con orgoglio di aver presentato a Mike
Morhaime\footnote{Cofondatore ed ex presidente di Blizzard; nella
trascrizione ``Mike Morayme''.} la propria campagna europea per
\emph{StarCraft~II}. Ci sono manager che dicono ``come mi prendo le colpe, mi
prendo i meriti'', ma i meriti li hanno già: è inutile ostentarli. Con i
fallimenti invece il manager è la \textbf{rete di sicurezza} del team. Il
relatore ne ha fatto una questione di principio: \textbf{non scaricare mai la
colpa su qualcun altro}.

\paragraph{Colpevole o causa.} Un esempio che il relatore attribuisce a Pirate
Software, creatore di contenuti ed ex dipendente Blizzard\footnote{Il relatore
aggiunge che si tratterebbe del figlio dello sviluppatore a cui si sarebbe
ispirato il personaggio del giocatore ``senza vita'' nell'episodio di
\emph{South Park} dedicato a \emph{World of Warcraft}. L'informazione è
riportata come detta in lezione e non è verificata.}. Contattando il servizio
clienti Blizzard per motivi personali, aveva scoperto che i colleghi gli
avevano scritto nei ticket informazioni confidenziali che non avrebbero dovuto
condividere in quel modo, nemmeno con un collega. Non ha segnalato le persone
da punire: ha detto ``c'è un problema, evidentemente le policy non sono
chiare''. Quando gli hanno chiesto i nomi, ha rifiutato: ``se vi dico chi è,
punite quelli, e gli altri cinquecento continueranno a farlo; non avremo
cambiato la causa alla fonte''.

\begin{quote}
  \emph{È inutile trovare il colpevole. Voglio trovare la causa. Perché il
  colpevole cambia.}
\end{quote}

Il colpevole a volte è vittima della situazione: per questo la definizione
dice \emph{``con le informazioni a sua disposizione in quel momento''}. Se il
relatore avesse rimproverato la ragazza in ritardo, non solo non avrebbe
risolto il problema, ma avrebbe creato un'\textbf{ingiustizia}: lei, che
attraversava il campus di corsa per una colpa non sua e di cui non poteva
parlare, sarebbe stata anche rimproverata.

Nessuno fa un brutto lavoro apposta: nessuno è contento quando sbaglia o
perde. Perché al lavoro diamo per scontato che gli altri sbaglino
intenzionalmente? Se qualcuno ha fatto una cosa male, è perché in quel momento
credeva fosse la scelta migliore. La domanda giusta non è ``perché hai fatto
questa cosa che è stata un fallimento?'', ma:

\begin{quote}
  \emph{Che cosa ti ha fatto pensare che questa fosse la scelta giusta?}
\end{quote}

Magari erano limiti di budget, di tempo, di scopo: si agisce su quelli, non
sulla persona. Se si cambia la persona ma le si danno gli stessi fattori, è
molto probabile che arrivi alle stesse conclusioni, e ci si troverà a dover
``disciplinare'' anche lei.

\paragraph{La reputazione costruita sui rumor.} Durante il giro di
presentazione in un ufficio (non in nessuna delle aziende citate), al relatore
è stato detto: ``questo è il dipartimento di design; se ti serve qualcosa
chiedi a lei, non a lui, perché lui consegna sempre in ritardo''. Oltre a
trovarlo poco professionale, una settimana dopo ha chiesto: ``cosa ti ha fatto
questa persona?'' ``A me niente, però ho sentito che\dots''. Si torna alle
scimmie: magari quella persona non consegna in ritardo da mesi, ma ormai lo
dicono tutti. È un circolo vizioso da fermare subito.

\paragraph{Niente escalation inutili.} Lo stesso vale per il feedback: quando
qualcuno va dal relatore dicendo ``ho un problema con questa persona'', lui
chiede: \textbf{``perché ne stai parlando a me e non a lei?''}. ``Perché sei il
mio capo'' non basta. Passare da lui significa dare per scontato che l'altra
persona non abbia interesse a risolvere, mentre magari basta parlarle. È la
strada più facile (``lo dico a Stefano, ci pensa Stefano''), e proprio per
questo bisogna mettere dei paletti: il manager interviene quando serve, ma le
cose che si possono gestire da soli vanno gestite da soli. (Quando invece il
confronto diretto fallisce, l'escalation è legittima: si veda la
sezione~\ref{sec:z1-qa-finale}.)

% ---------------------------------------------------------------------
\section{Traiettorie di carriera e valutazione delle performance}
\label{sec:z1-traiettorie}

\subsection{Quale traiettoria vorreste?}

Il relatore mostra un grafico con quattro curve colorate, contrassegnate
anche da lettere (A, B, C, D) per eventuali persone daltoniche (``a me è
capitato che qualcuno scoprisse di esserlo durante una lezione''). Le curve
rappresentano la performance lavorativa nel tempo. Due domande: \emph{quale
vorreste essere? E quale pensate realisticamente di essere?} Un partecipante
aspira a B e pensa di essere C; quasi tutti scelgono B.

Perché non A, che sale più rapidamente? Un partecipante risponde che una
crescita così repentina può essere più dannosa che utile, come una ``rush'' di
studio molto faticosa che poi non si riesce a sostenere. Il relatore propone
una lettura suggerita da un amico psicologo, che spiega la forma del grafico:
di B, C e D \textbf{vediamo dove finiscono}; di A no. Potrebbe continuare a
salire oltre la scala del grafico, o crollare sotto D. Non saperlo ci spinge a
preferire B. (A chi obietta che anche le altre proseguono oltre il bordo, il
relatore precisa: il bordo destro è ``oggi'', e di A non sappiamo come sia
arrivata fin lì.)

\subsection{La carriera non è una linea retta}

Tutti immaginano la carriera come una linea: si parte da X e, attraverso
stage e posizioni, si arriva a Y. In realtà carriera e performance avranno un
andamento irregolare, con alti e bassi. Come non si accetta che un team possa
temporaneamente \emph{underperformare}, tanto meno lo si accetta da se stessi.
Bisogna accettare che ci saranno momenti bassi.

Per questo il relatore \textbf{detesta le valutazioni annuali delle
performance}: prendono un punto casuale della traiettoria, che se si è
fortunati è il picco più alto e se si è sfortunati il più basso, e lo usano
per valutare un intero anno. Insegna ai manager a fare più punti di controllo
e a guardare altre cose:
\begin{enumerate}
  \item \textbf{quanto tempo passa tra un picco e l'altro}, cioè quanto
        \emph{cooldown} serve a una persona. È il motivo per cui, dopo la
        Games Week, tutti staccavano una o due settimane. Il relatore stesso,
        che si impegna moltissimo e tende al burnout, non potrebbe rifare lo
        stesso workshop il lunedì e il martedì successivi con la stessa
        qualità;
  \item \textbf{la differenza tra la miglior performance e la peggiore}: che
        cosa ha portato a lavorare così bene, e così male? Con chi si
        lavorava, su cosa ci si concentrava, quanto tempo e budget c'erano?
        L'obiettivo è capire come \emph{replicare} le condizioni che fanno
        lavorare bene, invece del solito ``hai lavorato male, è colpa tua''.
        Il relatore ha avuto molti manager per cui quando lavorava male era
        colpa sua e quando lavorava bene ``che bravo'', ma sapeva che raramente
        era l'unico fattore della propria performance. È la fallacia del ``se
        sai fare 100, perché non fai sempre 100?'', come nella scena di
        \emph{Schindler's List} sull'operaio che produce una vite in pochi
        secondi: non sempre un ritmo è sostenibile;
  \item \textbf{il fallimento di oggi confrontato con quello di sei mesi fa}.
        Il relatore non è contento della presentazione di oggi (l'avrebbe
        migliorata per altri due anni), ma rispetto a quelle di cinque anni fa
        vede un grande miglioramento di chiarezza: il suo peggior fallimento di
        oggi equivale al suo primo successo di inizio carriera.
\end{enumerate}
Quello che conta è la \textbf{traiettoria} dall'inizio dell'anno a oggi, non
i singoli punti. E se si riesce a capire insieme quando e perché si ha
successo e quando e perché si fallisce, si può lavorare per replicare solo le
situazioni che fanno lavorare bene.

\paragraph{Lo stesso vale per i giochi.} È il motivo principale per cui i
giochi delle grandi case vanno bene o male. Quando si riescono a replicare le
condizioni che hanno prodotto \emph{Warcraft~II}, esce \emph{Warcraft~III};
quando non ci si riesce, esce ``il nuovo \emph{Dragon Age}'': la gente non è
arrabbiata perché il gioco non è bello, ma perché non è come gli altri
\emph{Dragon Age}; non ne replica il \emph{feeling}, anche se potrebbe
essere un gioco favoloso. L'eccezione positiva è il passaggio da \emph{Dune} a
\emph{Dune~II}: il primo era di fatto un'avventura grafica, il secondo ha
generato il genere degli strategici in tempo reale. Gli sviluppatori avevano
capito che il successo non stava nel tipo di gioco, ma nell'ambientazione che
attirava le persone: un bel gioco in quell'ambientazione sarebbe piaciuto
comunque. Il relatore comprò \emph{Dune~II} aspettandosi un seguito
dell'avventura e, pur sorpreso, lo finì.

\subsection{Lying, hiding, faking}

Quando si è valutati su un singolo punto nel tempo, e da quel punto dipendono
bonus, stipendio e il rischio di licenziamento, si comincia a fare
\textbf{lying, hiding e faking}:
\begin{description}
  \item[Lying] dire che è sempre andato tutto bene, che si finirà sempre in
        orario, che non c'è nessun problema;
  \item[Hiding] nascondere: ``non era un mio task'', ``non c'ero'', ``non ho
        visto niente'';
  \item[Faking] fingere, anche in cose apparentemente innocue come ``sì, sì,
        ho capito tutto'' quando si ha solo paura di dire ``me lo rispieghi?
        Non ho capito niente''.
\end{description}

\subsection{``Non ho capito'' e la conferma delle istruzioni}

Dire ``non ho capito, ripetimelo'' è ancora visto come una figura da sciocco.
Il relatore non si arrabbierà mai se qualcuno gli chiede di rispiegare
un'intera slide: le cose ovvie per lui possono non esserlo per altri. Più si
conosce qualcosa, \textbf{più è difficile spiegarla}, perché si danno per
scontate molte cose. Ne parlava proprio con Diego, che gli diceva che fare
l'insegnante diventa sempre più difficile man mano che si diventa esperti. È
come spiegare \emph{EVE Online}, uno dei giochi più complessi in circolazione:
chi lo conosce bene dà per scontato che ``ovviamente non puoi fare questa
cosa'', mentre per un nuovo giocatore non è affatto ovvio. A scuola chi dice
``non ho capito'' viene visto come lo ``scemo''; nessuno si chiede ``non è che
l'ho spiegato male?''.

Il racconto più ``traumatico'' dei ragazzi che fanno stage nelle grandi
aziende è sempre lo stesso: ``Cinzia, devi fare questa cosa.'' Cinzia la fa.
``Non te l'avevo chiesta così.'' Il problema è che nessuno ha
\textbf{confermato che Cinzia avesse capito}. Il consiglio, valido sempre:
\begin{itemize}
  \item quando si \emph{ricevono} istruzioni, ripeterle: ``ho capito bene?
        Vuoi che faccia questa cosa in questo modo?'';
  \item quando si \emph{danno} istruzioni, chiedere a chi le riceve di
        ripeterle, non perché lo si ritenga incapace, ma per capire che cosa
        ha capito dalla propria spiegazione (``quindi vuoi un quadrato?''
        ``No, ho detto un triangolo\dots'' ``Però mi hai detto una forma con
        quattro angoli.'' ``Hai ragione: voglio un quadrato'').
\end{itemize}
Frasi come ``non te lo ripeto due volte'' o ``stammi dietro'' portano solo a
fare le cose come si pensa che vadano fatte e a sentirsi dire a fine progetto
``non la volevo così'': che si tratti di una feature, di un design o di un
modello 3D. (Un partecipante osserva che lo stesso accade insegnando le lingue:
chi insegna dà per scontato, ad esempio, che si sappia dell'esistenza del
genere neutro in tedesco.)

\paragraph{Come convincere il team a chiedere?} Il relatore mette sempre come
premessa che \textbf{l'obiettivo non è trovare la persona più intelligente
nella stanza, ma fare un buon lavoro}. Che Gianni capisca alla prima e un altro
alla quinta non gli interessa, se alla fine entrambi lavorano bene: ripetere
cinque volte è comunque meno lavoro che far rifare le cose sette volte. È un
investimento a lungo termine, e non dissimula alcun fastidio: davvero non gli
dà fastidio. Ma ripetere le stesse parole non aiuta: ``non ho sentito'' si
risolve ripetendo; ``non ho capito'' si risolve cambiando spiegazione. Meglio
ancora: ``dimmi che cosa hai capito tu, e vediamo cosa non torna''. Se
intravedesse nel team uno sguardo del tipo ``è la quinta volta che lo
chiede'', interverrebbe: spesso chi giudica è proprio chi dice ``ho capito'' e
poi sbaglia. Preferisce di gran lunga chi chiede cinque volte.

\subsection{Il feedback più utile che abbiate ricevuto}

\emph{Qual è il feedback più utile che abbiate mai ricevuto?} Il relatore ne ha
almeno dieci che gli hanno cambiato il modo di lavorare. Le risposte:
\begin{itemize}
  \item un semplice \textbf{``grazie''} per un lavoro fatto, più utile di un
        ``questa cosa fa schifo'';
  \item un feedback sulla propria \textbf{impulsività} nel rispondere: fermarsi
        un attimo e capire di cosa si sta parlando prima di intervenire;
  \item \textbf{essere se stessi}, anche al lavoro, e non scusarsi
        continuamente per ogni cosa.
\end{itemize}
I feedback che toccano di più sono quelli \textbf{veri}. Il ``grazie''
significa ``ho visto quello che hai fatto e l'ho apprezzato''. Quelli negativi
che si ricordano sono quelli che hanno fatto realizzare ``questa cosa che
pensavo di saper fare, non la so fare'', oppure ``la so fare, ma questa volta
non mi è venuta'' (non è che se una cosa una volta non riesce non la si sa più
fare: può capitare di sbagliare una porta). Arrivano sempre da un
\textbf{punto cieco}: qualcosa che non si sapeva fosse percepito dagli altri.
Conta più ciò che le persone ti dicono di ciò che pensi che pensino di te; e
ciò che il relatore pensa che gli altri pensino di lui è sempre terribile,
mentre il loro giudizio si rivela molto più positivo del previsto.

\paragraph{Il denominatore comune siete voi.} Si potrebbe pensare che questi
concetti non siano applicabili da chi arriva in un'azienda senza poterne
cambiare la cultura. Ma c'è sempre un denominatore comune a cui applicarli:
\emph{se stessi}. La cultura aziendale può non spingere a dare feedback, può
vedere male chi chiede aiuto; ma il relatore ha visto che quando lui lo fa, le
persone intorno tendono a farlo, e un gruppo che lo fa lavora meglio. Perfino
i CEO più ostici a un certo punto gli chiedevano: ``perché il tuo team ora
performa meglio degli altri?''. La differenza è il default to trust: una
pratica che col tempo fa operare il team \emph{come un team} e non come un
insieme di individui. Un videogioco è un progetto corale: c'è chi fa giochi da
solo, ma non si fa da soli \emph{Red Dead Redemption~2}; e se il gruppo non è
affiatato e allineato non esce \emph{Red Dead Redemption~2}, esce
\emph{Anthem} (``oggi la mia fissazione è \emph{Anthem}, domani sarà
un'altra'').

% ---------------------------------------------------------------------
\section{Come dare un feedback}
\label{sec:z1-dare}

\subsection{Concentrarsi sull'azione e sui suoi effetti}

Il primo principio: \textbf{il feedback non riguarda mai la persona}. Non si
dice ``tu hai fatto questa cosa che non va bene'', ma ``questa cosa non va
bene: come facciamo a non farla più?''. Che sia un comportamento, un'azione,
una decisione o una strategia, ci si concentra su quello, e si
\textbf{esplicita l'effetto} che ha. Nell'esempio del ritardo l'effetto era
``quindici minuti di attesa moltiplicati per venti persone''. In un progetto
può essere una catena:
\begin{quote}
  se non rispetti le deadline, l'animazione che viene dopo resta ferma; se
  loro sono in ritardo, la \emph{vertical slice} è in ritardo; se la vertical
  slice non è pronta per la Gamescom, non andiamo alla Gamescom e non troviamo
  un publisher; se non troviamo un publisher, chiudiamo.
\end{quote}
Più la conseguenza è grave, più dovrebbe sembrare piccolo il gesto di dire
``datemi una mano'': è il \emph{male minore}. Poi, a cose fatte, ci si dà
feedback: cosa ha funzionato, cosa no, cosa fare di più, di meno, cosa smettere
di fare.

\subsection{Il pericolo degli ambienti solo positivi}

Gli ambienti dove ci si fa sempre grandi complimenti e nessuno parla dei
problemi creano difficoltà vere alle aziende. Il problema non è mai la
mancanza di complimenti, ma un \textbf{eccesso di ottimismo}: si comincia a
pensare che le cose andranno come si spera, perché tutti intorno danno solo
feedback positivi.

A Blizzard i feedback determinavano gli \textbf{aumenti annuali}: ognuno
doveva indicare cinque persone che avrebbero dato feedback su di lui, e
ovviamente tutti indicavano gli amici, o colleghi di progetti andati bene. Un
sistema nato con le migliori intenzioni diventava uno \emph{scambio di voti}
(``se mi metti follow ti metto follow''), senza alcun miglioramento.

Al contrario, \textbf{le conversazioni più brutte sono state le più utili}:
affrontare i licenziamenti a Riot, o la cancellazione di \emph{Project Titan}
a Blizzard, il progetto che doveva essere il nuovo MMO dopo \emph{World of
Warcraft} e da cui poi è nato \emph{Overwatch}\footnote{Nella trascrizione si
legge anche ``doveva essere Warcraft 2'', probabilmente un lapsus o un errore
di trascrizione.}. Andando a vedere perché lo si cancellava, c'erano solo
motivazioni razionali e pragmatiche, di business, chiare a tutti, non un ``Mike
Morhaime ci odia''. Sono anche le conversazioni in cui ha lavorato su se
stesso: è molto \textbf{avverso ai conflitti}, e questo a volte gli ha
procurato conseguenze molto più gravi di quanto sarebbe costato un piccolo
disagio. Dieci anni fa faceva fatica ad accettare questo difetto; ha dovuto
subire le conseguenze di conflitti evitati e poi cresciuti.

\subsection{Che cosa fare}

Le linee guida della slide, che il relatore legge rapidamente perché saranno
disponibili nel materiale:
\begin{description}
  \item[Attenzione al linguaggio del corpo] dare un feedback a braccia
        incrociate o guardando altrove è diverso dal guardare l'altro negli
        occhi.
  \item[Sii specifico] con esempi tangibili.
  \item[Esplicita gli effetti] come visto sopra.
  \item[Discuti il prima possibile] non aspettare un mese per parlare di
        qualcosa successo ieri, altrimenti nessuno ricorderà più bene cosa sia
        accaduto.
  \item[Concentrati sull'azione e sul problema, non sulla persona] il nemico è
        l'azione, non la persona che si ha davanti. Non bisogna far cambiare
        l'altro: bisogna evitare che il problema si ripresenti.
\end{description}

\paragraph{Gli insulti sono rivolti al ruolo.} Un esempio solo apparentemente
scollegato. Il relatore ha lavorato su \emph{League of Legends}, ``non
esattamente noto per avere una community da orsetti del cuore'', in varie
\emph{task force} sulla tossicità. A chi smetteva di giocare veniva inviato un
sondaggio sul perché; il 33\% indicava la tossicità della community. Dopo varie
campagne di sensibilizzazione e cambiamenti si è scesi al 30\%: sembra solo
un 3\% in meno, ma valeva circa \textbf{120 milioni di revenue in più
all'anno}, perché quei giocatori restavano esposti a skin e contenuti e
continuavano a comprare. Anche qui, piccoli guadagni che si sommano.

Lavorando con degli psicologi, aveva notato che gli insulti \textbf{non sono
quasi mai rivolti alla persona, ma al ruolo}: ``il mio jungler fa schifo'',
``il mio DPS non è capace'', ``il nostro portiere è una pippa''. Chiamare
qualcuno per ruolo lo \emph{depersonalizza} e rende più facile insultarlo. Lo
stesso nell'esempio del giro dell'ufficio: non ``non chiedere a Giulio'', ma
``non chiedere a quel grafico, consegna sempre in ritardo''. Concentrarsi
sull'azione, senza attaccare la persona, fa sì che il feedback non sia (e non
venga percepito come) un attacco personale.

\paragraph{Considera il carattere.} Qui tornano utili i test di personalità.
Servono più agli studenti che al relatore (che sa già che usciranno 53 profili
diversi, senza un pattern): bisogna conoscere il proprio profilo e
\textbf{condividerlo con chi lavora con noi}, chiedendo anche agli altri di
fare lo stesso. In ufficio il relatore ha usato una versione più avanzata, a
pagamento, con un report di circa trenta pagine: il \textbf{profilo
DISC}\footnote{Nella trascrizione ``disk''. Si tratta del modello DISC
(Dominance, Influence, Steadiness, Conscientiousness).}. Ha mappato tutta
l'azienda e ha fatto stampare il profilo di ciascuno sulla sua scrivania, così
che tutti sapessero con chi stavano parlando. Capire come ragionano gli altri
(``è introverso, non gliene parlo davanti a tutti'') è una delle cose che gli
sono tornate più utili: \emph{adattare ciò che si vuole dire a chi si ha
davanti}.

\paragraph{Assicurati di non disturbare.} Sotto il profilo di ognuno c'era un
piccolo \textbf{semaforo}: verde, ``parlami pure, interrompimi''; giallo,
``interrompimi solo se è urgente''; rosso, ``mi sto concentrando, non mi
parlare''.

\paragraph{Lascia finire il lavoro.} Durante gli eventi dal vivo il relatore si
era reso conto di quanto fosse grave dare feedback mentre qualcuno sta ancora
lavorando. Se una cosa sta andando male, dirlo in quel momento non aiuta: tra
la pressione di finire e quella del feedback, peggiorano entrambe le cose, e
anche il rapporto. Meglio lasciarla finire, anche male, e parlarne dopo.

\subsection{Che cosa non fare}

\begin{description}
  \item[Non basarti su supposizioni] ``mi hanno detto che ieri in un meeting
        hai fatto\dots'': o c'eri e l'hai visto, o non esiste.
  \item[Non giudicare la persona.]
  \item[Non fornire feedback vago.]
  \item[Non accettare e non fornire feedback di terze parti] Se qualcuno vi
        dicesse ``mi hanno detto che hai fatto un brutto lavoro'', il primo
        pensiero non sarebbe ``cosa ho fatto male?'', ma \emph{``chi te l'ha
        detto? Chi è la vipera che ho in seno? Chi è l'impostore?''}. Siamo
        fatti così. Ogni conversazione che inizia con ``mi hanno detto che'' è
        inutile: l'altro non ascolterà, penserà solo a chi dare la caccia, e
        spesso penserà alla persona sbagliata, creando antipatie tra persone
        che non c'entravano nulla. Quando al relatore dicono ``mi hanno detto
        che lui ha fatto\dots'', risponde ``vai a parlarne con lui''; se gli
        dicono ``mi hanno detto che tu hai fatto\dots'', risponde ``mandami la
        persona che te l'ha detto, altrimenti non ne parliamo''. Non per
        cattiveria: se chiedesse ``che cosa ho detto di preciso a Marco?'' e la
        risposta fosse ``non lo so, non c'ero'', di cosa si starebbe parlando?
        \textbf{Parlate con la fonte}, altrimenti fate come se quel feedback non
        esistesse.
  \item[Non dare consigli non richiesti] come visto, non portare la soluzione:
        la soluzione deve venire da chi riceve il feedback.
  \item[``Non mi piace'' non è un feedback] ``Penso che questo design non
        funzioni perché, visto che il nostro obiettivo è\dots'' ha argomenti;
        ``non mi piace il combat system'' è un'opinione. Anche ``io l'avrei
        fatto diversamente'' non è un feedback.
  \item[Accettare non vuol dire fare come proponi] deve essere chiaro a chi dà
        il feedback: l'aspettativa è essere ascoltati e presi in
        considerazione, ma se la responsabilità è dell'altro, deciderà
        l'altro.
  \item[Non offenderti] se il feedback non viene seguito. Restarci male
        significa pensare di saperne più dell'altro, cioè \emph{venire meno al
        default to trust}. ``Decidi tu'' significa ``mi fido di te, perché sei
        lo specialista, e mi fido che mi ascolterai e farai la cosa migliore
        per tutti'', che potrebbe anche non essere ciò che ho detto io. Se
        voglio che l'altro faccia sempre come dico io, mi sto fidando solo di
        me stesso.
\end{description}

\subsection{Il primo feedback e il debriefing}

La cosa più difficile è dare il \textbf{primo} feedback: è un muscolo da
allenare. Chi non l'ha mai fatto (soprattutto con questa struttura) dovrebbe
provarci, anche su minuzie (``abbiamo concordato di usare WhatsApp nel gruppo:
non potremmo usare Telegram?''), anche solo per vedere come va. Altre
indicazioni:
\begin{itemize}
  \item feedback \textbf{solo uno a uno}, mai davanti ad altri, soprattutto se
        negativo;
  \item \textbf{prepararsi} prima;
  \item fare il \textbf{debriefing del feedback}: dopo, fermarsi a pensare a
        com'è andato. Il relatore, quando media un feedback tra due persone,
        si sente spesso rispondere ``benissimo'' da una e ``malissimo''
        dall'altra: la percezione è molto diversa. Non si fa problemi a tornare
        dalla stessa persona: ``ieri ti ho dato un feedback: com'è andata? Sei
        soddisfatto? Come poteva andare meglio?''. Con le persone con cui si ha
        meno a che fare, questo aiuta molto la volta successiva.
\end{itemize}

% ---------------------------------------------------------------------
\section{Come ricevere un feedback}
\label{sec:z1-ricevere}

Per il relatore è la parte più difficile. Si riassume in quattro passi:
\begin{enumerate}
  \item \textbf{Ascolta}: si parte sempre ascoltando.
  \item \textbf{Chiedi dettagli}, se serve, ma per capire meglio, non per
        difenderti.
  \item \textbf{Accetta, a prescindere}: non si sta decidendo chi ha ragione;
        si sta ascoltando un'opinione. Non c'è nulla da cui difendersi.
  \item \textbf{Agisci}: fatti un piano. Può essere banale (``mi hanno detto
        che interrompo sempre: cerco di non interrompere''), oppure ``devo
        migliorare questa cosa: come?''. Pensa anche a \emph{chi} potrebbe
        aiutarti e a \emph{che supporto} ti serve: il piano non deve essere
        per forza ``devo risolvere tutto da solo''.
\end{enumerate}
Una domanda che molti sottovalutano: \emph{``sono junior in un'azienda: sarò
mica il primo ad avere questo problema?''}. Probabilmente qualcun altro l'ha
già avuto: basta chiedere ``tu come hai fatto?''.

\subsection{Il feedback positivo}

Tutto questo vale per i feedback costruttivi negativi, a cui si pensa di
default; ma esistono anche i \textbf{feedback positivi}, e vanno usati. Quando
si riceve un ``grazie'' o un ``bravo'', la domanda è: \emph{come faccio a
ricreare le stesse condizioni?}

Il relatore ha lavorato anche in aziende dove il feedback era solo negativo.
Il rischio? Si immagini l'impiegato ideale: vuole sempre migliorare, è sempre
pronto a dire ``magari sbaglio io''. Se nessuno gli dice cosa sta facendo
\emph{bene}, spinto dal desiderio di migliorare rischia di cambiare cose che
erano già buone, e di farle peggio. Il rinforzo positivo \textbf{non è solo una
gentilezza}: è un atto pratico, che evita che le persone peggiorino ciò che
sanno già fare bene.

\subsection{Il feedback accidentale}

Esiste infine il \textbf{feedback accidentale}, molto frequente negli uffici:
\begin{itemize}
  \item il ``chiedi a lei, non a quel grafico'' del giro dell'ufficio è un
        feedback negativo su una persona, dato accidentalmente, né richiesto
        né voluto;
  \item a Riot li chiamavano \textbf{drive-by}, come nelle scene dei film in
        cui passa un'auto, si abbassa il finestrino, si spara e si riparte. In
        molte aziende con una ``cultura del feedback'' in cui tutti devono
        dire la loro, si riceve una valanga di commenti del tipo ``comunque il
        sito a me non piace, ciao''. Non si sa perché, non si sa cosa fare di
        diverso, non si sa che rischi ci siano: \emph{a che cosa è servita
        questa conversazione?} A niente.
\end{itemize}

% ---------------------------------------------------------------------
\section{Quando, quanto e con chi}
\label{sec:z1-quando}

\subsection{Frequentemente e recentemente}

\emph{Quando avete ricevuto il vostro ultimo feedback?} Per gli studenti,
qualche giorno fa (nel corso) o qualche mese fa (al lavoro, dopo una
consegna). Come aveva osservato Silvia, il feedback andrebbe dato
spesso.

Il relatore usa l'esempio di una nave da crociera (``la Costa Concordia'') che
deve andare da A a B. La rotta ideale (azzurra) è quasi impossibile: nessun
progetto va esattamente come previsto. C'è la rotta ``panoramica'' (verde), che
gira intorno all'isola d'Elba; ci sono poi rotte (arancione e rosa) che a un
certo punto si allontanano molto e poi correggono bruscamente: qualcuno si è
accorto che si stava sbagliando, ``ma non è quello che avevo chiesto'',
``non è dove dovevamo andare''. Sono i famosi \textbf{sanity check} o
\emph{checkpoint}.

L'ideale è avere \textbf{tanti piccoli check per fare tante micro-correzioni}:
correggere la rotta di un paio di gradi è molto più rapido che farlo dopo
essersi allontanati. Bastano trenta secondi di ``ehi, vado bene?''. Spesso non
lo si fa per paura di sembrare insicuri o poco autonomi, ma per il relatore è
l'opposto: quando qualcuno gli chiede ``vado bene così?'' e lui risponde
``sì'', si alleggerisce anche del proprio carico mentale. Un capo che assegna
un compito e lo vede solo alla fine vive sempre col dubbio (``avranno capito?
Lo staranno facendo giusto? Arriverà in tempo?''); con un check di dieci
secondi al giorno (``tutto ok, Stefano'') è tranquillo, e la persona va avanti
sapendo che quanto fatto finora va bene. Non serve fare ogni giorno il giro del
dipartimento: anche ogni uno o due giorni basta a evitare la rotta verde.

Per il feedback è lo stesso: \textbf{va dato subito}. Dire oggi ``ieri hai
fatto questa cosa'' è molto diverso da dire ``negli ultimi sei mesi hai fatto
questa cosa''. Il relatore si è sentito dire più volte, alla valutazione di
dicembre, ``da maggio hai fatto questa cosa che non andava''. La sua risposta:
``e aspetti dicembre per dirmelo? Mi hai lasciato sbagliare per sei mesi?''. È
una mancanza di rispetto, e la conversazione non può più essere costruttiva.
Per questo non ama le valutazioni annuali: i manager crescono abituandosi a
prendere appunti delle cose negative per tirarle fuori alla valutazione
(``a marzo hai fatto questa cosa'' ``sì, però poi ho smesso'' ``sì, però a
marzo l'hai fatta'').

In sintesi:
\begin{description}
  \item[Frequentemente] riguarda l'intervallo tra un check e l'altro.
  \item[Recentemente] significa non lasciar passare un anno prima di dare il
        feedback, altrimenti non serve più a niente.
\end{description}

\subsection{Quanto dura un feedback}

Un feedback dura \textbf{da trenta secondi a quindici minuti}; raramente di
più. Chi non ne ha mai dato uno dovrebbe sforzarsi di darne almeno uno, magari
positivo per cominciare. Ogni feedback dato o ricevuto è un
\textbf{deposito nel proprio ``fondo fiduciario''} (\emph{trust fund}), che
costruisce la reputazione futura. Il relatore si è costruito la reputazione di
chi dà feedback il più possibile oggettivi e analitici: giusti o sbagliati che
siano, costringono a pensare. Non impone mai il suo punto di vista (``si fa
come dico io''), ma chiede ``tu come faresti? Come posso aiutarti?''. Il
feedback, per lui, serve a \emph{far crescere qualcun altro}.

\subsection{I manager peggiori: ``a che numero sto pensando?''}

I peggiori manager che il relatore abbia avuto erano quelli che volevano che si
facesse come volevano loro, anche quando chiedevano ``tu come la faresti?''.
In realtà la domanda era ``\emph{a che numero sto pensando?}'': se dicevi 5,
``no, dovremmo fare 9''; se dicevi 9, ``no, 5''. Lo stress costante di dover
indovinare cosa vuole il capo nasce dalla mancanza di fiducia (``se mi chiedi
come fare, fammelo fare'').

Lo stesso gli succede con i CEO che lo chiamano per ``sistemare la cultura
aziendale'' e poi dicono ``non sono sicuro, io farei diversamente''. La sua
risposta: ``lo so che faresti diversamente, per questo mi hai chiamato''. Se si
chiama un esperto, bisogna lasciargli fare il suo lavoro: non si va dal
meccanico a cambiare le gomme per poi dirgli ``non dovresti girare un po' di
più quel bullone?''. Chiedere un parere a un esperto significa dargli anche il
potere di scegliere. Il problema è il \emph{micromanagement}, o in generale il
non accettare punti di vista diversi dal proprio (magari perché si è il
fondatore e si pensa di essere l'unico a saper fare le cose).

\subsection{Dare feedback a un superiore egocentrico}

\paragraph{Come si dà un feedback negativo a un superiore con un grande ego?}
Il trucco è \textbf{non puntare sull'ego, ma sull'azione e sui suoi effetti}.
Non ``tu fai così e io mi sento così'': anzi, si alimenta il suo ego
rendendolo il \emph{main character} dell'avventura, l'eroe che risolve il
problema (``tu, grande manager, aiutaci a risolvere questo problema che
affligge noi che facciamo l'operatività''). Non è l'ideale (``io cambierei
manager''), ma pochissimi possono scegliersi i colleghi: il relatore stesso ha
quasi sempre ricevuto team preassemblati (``ecco le chiavi, questa è la tua
macchina''). Bisogna giocare con le carte che si hanno in mano.

\paragraph{Un esempio pratico.} Con un manager che fa micromanagement e mette
sempre ``il becco'' nel proprio lavoro, si può dire: ``se ti fidi di me per
questo lavoro, con le tue capacità potresti dedicarti a questa cosa più
strategica e di lungo termine: pensa quanto sarebbe più utile per l'azienda''.
In sostanza: ``guarda cosa potresti raggiungere se mi lasciassi fare il mio
lavoro''. Non come un ``tu fai il tuo e io faccio il mio'', ma come un
``ti sollevo da questa parte del carico, che so che ritieni importante: lascia
che ci pensi io, così puoi dedicarti ad altro'', che magari gli dà anche più
visibilità.

Si gioca su ciò che quella persona pensa di saper fare, e che probabilmente sa
anche fare. Anche Michael Scott di \emph{The Office}, per quanto preso in giro,
è raccontato come un venditore bravissimo, che ha chiuso i contratti più
importanti per l'azienda, e rispetto ai manager che gli succedono è forse
``uno dei meno peggiori''. Tutti, anche i manager peggiori, hanno punti di
forza: \textbf{giocate su quelli}.

% ---------------------------------------------------------------------
\section{Il feedback come regalo}
\label{sec:z1-regalo}

Per chiudere, il relatore risponde alla domanda d'apertura, che aveva lasciato
apposta in sospeso: che cos'è il feedback per lui?

\begin{quote}
  \emph{Il feedback è un regalo che mi viene fatto da qualcuno che so che vuole
  il mio interesse.}
\end{quote}

Ha costruito un team di persone che si fidano di lui e che non hanno motivo di
sabotarlo o di ferirlo: per questo interpreta anche le cose peggiori con il
default to trust. È un regalo, e cerca di prenderlo in considerazione il più
possibile.

Ci è arrivato partendo da se stesso. Oltre a essere avverso ai conflitti, si
definisce ``avverso alla comunicazione umana'': paradossalmente, una persona
introversa e timida finita a fare comunicazione, marketing e ora HR (``succede
spesso''). Prima di ordinare una pizza al telefono ha dieci minuti di
preparazione mentale, e se l'ordine delle domande cambia va in confusione. Per
lui un feedback di dieci minuti sono spesso \textbf{tre ore di preparazione}:
costruisce un \textbf{albero dei dialoghi}, come in un'avventura grafica
personale, con tutte le opzioni. Poi spesso finisce con ``hai fatto questa
cosa, andrebbe migliorata'' ``hai ragione, grazie'', e basta.

Così ha pensato: chissà quanti si fanno gli stessi problemi con lui, e passano
ore a prepararsi per dirgli ``questa cosa l'hai fatta male'', sapendo che
riceveranno comunque quell'occhiata fugace del ``come ti permetti''. Eppure
vengono a dirglielo. Non può far altro che rispondere: ``grazie, hai ragione,
cercherò di fare meglio''. È così che ha metabolizzato il ricevere feedback:
capendo che non c'è nessuna cattiva intenzione e che, se non proprio felice
(``ci ho messo una pietra sopra''), può almeno essere \emph{grato} che ci sia
ancora qualcuno che ha così a cuore il suo interesse da dargli feedback.

\paragraph{Il silenzio fa più paura.} Non bisogna spaventarsi dei feedback
negativi che emergono dai sondaggi aziendali: bisogna spaventarsi del
\textbf{silenzio}, quando nessuno dice più nulla. Significa che nessuno si
fida più che si sappiano risolvere i problemi.

\paragraph{La sindrome dell'impostore.} Il suo motivatore personale, all'inizio,
era soprattutto ``non voglio fallire'': per questo cercava di migliorare, e
migliorare è diventato accettare feedback, perché più feedback riceve, più
migliora. Il feedback è un \textbf{acceleratore}, e lui cerca di incentivare
gli altri a darglielo. Resta però sempre la vocina del ``e se si accorgono che
non sei capace?''. Con la \emph{impostor syndrome} convive da anni e
continuerà a conviverci; se ce l'avete, conviveteci anche voi. Avrete sempre
la sensazione che ``se ne accorgeranno'': forse sì, ma sapete fare anche
altre cose. Non concentratevi solo su ciò che non sapete fare. Gli errori
capitano: la perfezione non è di questo mondo.

\paragraph{Esercitarsi in accademia.} Per fare pratica con il feedback,
l'ambiente migliore è quello dei colleghi di corso: persone nella stessa
situazione, con esperienza simile e lo stesso obiettivo, in un contesto in cui
non si può essere licenziati e l'unico effetto a lungo termine è ``ho parlato
con questa persona, ora le cose vanno meglio''. E bisogna avere il coraggio di
parlare anche con i docenti (``dire a Diego: secondo me questa lezione potevi
farla in questo modo''; ``so che è una persona che lo accetta''), senza farsi
bloccare dalla gerarchia. Nell'azienda giusta (come quella in cui lavora ora) il
relatore va sempre dal CEO a dare feedback, sapendo che non verrà licenziato:
il CEO gli ha detto ``ti ho preso per darmi feedback''.

La frase di chiusura:

\begin{quote}
  \emph{Per superare le vostre difficoltà, tutto quello che vi viene chiesto è
  il coraggio di chiedere aiuto.}
\end{quote}

Chiedere aiuto non è una debolezza: è un segno di consapevolezza di sé
(\emph{self-awareness}) e di forza personale. Se tutti lo facessero, i
progetti andrebbero molto meglio; spesso invece ci si fa frenare dal ``non
voglio fare brutta figura'', ``non voglio passare per quello scemo'', e sono
proprio queste paure a rallentarci e a impedirci di raggiungere gli obiettivi.

% ---------------------------------------------------------------------
\section{Domande e risposte finali}
\label{sec:z1-qa-finale}

L'ultima mezz'ora è dedicata a domande libere.

\paragraph{E se una persona peggiora per uno o due anni di fila?} Se nessuno
gliel'ha mai detto, c'è un grosso problema, non suo ma dell'azienda. Se invece
c'è consapevolezza, il percorso è: ``cosa vuoi fare per risolvere?''; ``ho
provato questo, questo e questo''; ``visto che non ha funzionato, vuoi un mio
consiglio? Proviamo con quello?''; ``che supporto ti serve?''. Se, con tutto il
supporto, il tempo e le soluzioni concordate per mesi, il miglioramento non
arriva, per il relatore la soluzione è: \textbf{``dimmi un prezzo per andare
via''}. Preferisce pagare una buonuscita (\emph{severance package}) che
avviare la strada dei richiami formali sulla performance fino al
licenziamento per giusta causa. A un certo punto andare avanti non è un favore
per nessuno.

\paragraph{È vero che nei videogiochi si viene sfruttati? Meglio un piccolo
studio italiano?} Per il relatore si è quasi più a rischio a Blizzard che in un
piccolo studio di Roma o Torino. Nel piccolo studio si sa dove ci si va a
cacciare: non ci sarà il team building a Santo Domingo, ma ci si concentra su
un gioco alla volta e, se va bene, si fa il successivo. A Blizzard si può
lavorare su \emph{Overwatch}, che vende milioni, e perdere comunque il posto,
in modo totalmente scollegato dal proprio lavoro. Gli stipendi più alti
all'estero compensano anche questo rischio, che spesso non dipende da come va
il gioco (anzi, ``il gioco non ha fatto come volevamo'' è spesso una scusa per
licenziamenti che non avevano motivo).

\paragraph{Cambiare settore a trent'anni, da web developer full remote?} Lo
stipendio indicato dallo studente non è fuori dalla portata di vari studi
italiani, anzi è superabile. Non va esclusa la possibilità di lavorare da
remoto per aziende piccole ma estere. Alcuni amici del relatore, pagati con
stipendi italiani, si sono trasferiti a Bali o in paesi con un costo della vita
più basso; il collega di \emph{PUBG} percepisce uno stipendio olandese vivendo
a Milano, che per quanto cara non costa come i Paesi Bassi.

\paragraph{Come si diventa manager?} Dipende dall'azienda: alcune promuovono i
migliori (processo tendenzialmente sbagliato), altre propongono un percorso di
\emph{people management}. Diventare manager significa \textbf{abbandonare ciò
che si sa fare per imparare a fare altro}. Il relatore ha dovuto accettare che
non sarebbe più stato ``il miglior community manager italiano del mondo''
(``eravamo in pochi, quindi era facile''): dedicandosi alle persone non poteva
più fare il suo lavoro bene come prima. È stata una negoziazione con se
stesso: fare il manager significa non essere più l'esperto di qualcosa, ma
saper gestire e far crescere persone. Significa anche \textbf{cedere il
controllo della propria carriera}: se il team decidesse di non fare più nulla,
il manager ne pagherebbe le conseguenze. Il suo lavoro diventa scegliere le
persone giuste e farle lavorare bene. Lui lo esplicita sempre: sa di affidare
al team parte della propria carriera, e il team sa che lui combatterà per
loro, anche prendendosi le colpe.

\paragraph{L'inglese.} \textbf{Va assolutamente imparato}, anche lavorando in
Italia. Non saperlo preclude circa il 95\% del mercato (tutto quello fuori
dall'Italia) e soprattutto il \emph{networking} con persone che potrebbero dare
occasioni future: l'inglese è spesso l'unica lingua in comune con chi lavora
all'estero ed è oggi la lingua universale dei videogiochi.
\begin{itemize}
  \item \emph{Che livello serve?} Non importa B2 o C2: bisogna saper sostenere
        una conversazione, guardare almeno Netflix con sottotitoli in inglese,
        comunicare senza dodici pause e senza ``come si dice\dots''. Il
        certificato non verrà mai chiesto: verrà chiesto se si sa parlare
        inglese. Il relatore ha sentito persone con un C2 che non sapevano
        parlarlo.
  \item \emph{Come allenarlo senza occasioni di conversazione?} Netflix, Prime
        Video, Now TV, Udemy, Coursera, YouTube, Duolingo, siti di scambio
        linguistico con madrelingua, Discord: non sono mai esistiti così tanti
        mezzi gratuiti. Il relatore ha imparato l'inglese giocando (con una
        pronuncia sbagliata, perché leggeva senza sentire, su cui ha poi
        dovuto lavorare); guardare un film in inglese con i sottotitoli in
        italiano insegna anche la traduzione.
  \item Oltre all'ascolto passivo è fondamentale \textbf{parlare}: in una
        flotta di \emph{EVE Online} con 150 persone quasi tutte madrelingua
        inglese si è esposti a decine di accenti e modi di dire (mille modi di
        dire ``come stai'' oltre a \emph{how are you}); lavorando all'estero si
        impara a distinguere gli accenti (Manchester, Londra) e ad apprezzarli
        anche nei media.
\end{itemize}

\paragraph{Quanto deve sapere un junior?} Dipende da quanto è tecnico il
ruolo. Un game designer junior che non sa qualcosa può essere affiancato da un
senior; un programmatore che non sa programmare, in quel momento, non è utile
all'azienda.

\paragraph{Annunci ``junior con cinque anni di esperienza''.} ``Cercasi
ventenne con trent'anni di esperienza.'' Non è l'annuncio per un junior: è un
errore dell'azienda, che in realtà cerca un \emph{mid}. Riceverà profili
junior, si stupirà che non abbiano cinque anni di esperienza, rifarà l'annuncio
e avrà perso sei mesi. Il consiglio: tenere d'occhio quella pagina, perché
dopo qualche mese l'annuncio verrà corretto.

\paragraph{Una struttura per il post-mortem.} A chi chiede un formato di base,
il relatore anticipa quanto vedrà nella seconda giornata a proposito di
brainstorming: una lavagna \textbf{divisa in quattro aree}, su cui ognuno
attacca tutti i post-it che vuole:
\begin{center}
  \begin{tabular}{@{}p{0.44\textwidth}p{0.44\textwidth}@{}}
    \toprule
    \textbf{Smettere di fare} & \textbf{Iniziare a fare} \\
    ciò che non ha funzionato & ciò che non abbiamo fatto e andava fatto \\
    \midrule
    \textbf{Fare di meno} & \textbf{Fare di più} \\
    ciò che ha fatto perdere tempo rispetto all'impatto ottenuto &
    ciò che ha funzionato ma, fatto con troppa leggerezza, non ha dato i
    frutti sperati \\
    \bottomrule
  \end{tabular}
\end{center}
Con i post-it ognuno ha voce e tutto risulta catalogato; per il report basta
una foto della lavagna (o uno screenshot, se la lavagna è digitale, come Miro).

\paragraph{Come sbloccare la paralisi del feedback verso il management?} Uno
studente ha visto manager dire ``non abbiate paura, la mia porta è sempre
aperta'' senza che cambiasse nulla. C'è una dissonanza tra ciò che dicono e
ciò che fanno: è ``la mamma che dice vieni qua che non ti faccio niente''.
Qualcosa deve aver fatto pensare alle persone che non è sicuro. Prima va
sistemato quello; poi si può provare con i \textbf{feedback anonimi}.
L'obiettivo del relatore resta il feedback aperto e firmato, ma a Stormind ha
iniziato con feedback anonimi verso il management per verificare se
l'anonimato fosse un fattore. Non lo era (evidentemente non c'erano problemi a
parlare con il CEO), e la volta successiva intende abbandonarli. Ma è un buon
\emph{check}: se senza anonimato non arriva nulla e con l'anonimato arrivano
cento feedback, il problema è chiaro.

\paragraph{Quando è giusto rivolgersi a un superiore?} La strada giusta resta
parlare direttamente con l'interessato. Ma se si prova e non si trova una
soluzione, la cosa più ovvia è andare dal proprio manager insieme: ``ci
aiuti?''. Quando il relatore parlava di mediare i dialoghi, era proprio questo:
fare da moderatore, aiutare a sbloccare i punti fermi o, se non si riesce a
scegliere tra A e B, scegliere lui. L'escalation, quando necessaria, va fatta.

\paragraph{Il tempo libero veniva rispettato?} No, e se l'è dovuto imporre
lui. Il lavoro gli piaceva così tanto che continuare a casa non gli sembrava
lavorare. Gli si è ritorto contro con il lancio di \emph{Diablo~III}, fatto in
orario americano (le due o l'una di notte in Europa), con l'ufficio chiuso
perché a nessuno era stato chiesto di restare. Il lancio va malissimo, i
giocatori protestano, e lui, sveglio, risponde sui forum fino alle tre, quattro,
cinque del mattino: ``non posso risolvere il problema, ma resto qui ad
aggiornarvi con quello che riesco a sapere dal team di sviluppo''. Il giorno
dopo trova sulla scrivania un \textbf{richiamo formale dell'HR}: aver lavorato
fuori orario senza autorizzazione, condotta non professionale. Per lui è stato
un interruttore: ``ora me ne vado''. Due settimane dopo faceva colloqui con
Riot. Ancora oggi è contrario a quel richiamo, ma gli ha fatto capire come
ragiona un'azienda: ``hai un contratto 9--18, non mi interessa se l'hai fatto
con le migliori intenzioni''. Eppure per anni ha goduto di un'ottima
reputazione in quelle community come ``quello che è rimasto con noi fino alle
cinque del mattino'', e anche l'azienda ne ha beneficiato.

Ci sono poi aziende in cui, soprattutto per i developer, viene richiesto il
famigerato \textbf{crunch}: lavoro straordinario continuativo, senza pause.
Dipende molto dall'azienda; sta venendo (``grazie al cielo'') deprecato, perché
molte aziende capiscono che non è sano ed è visto male anche dal pubblico, ma
esiste ancora.

\paragraph{E se il tuo capo demotiva qualcuno del tuo team?} Qui ``diventa
personale'': il team del relatore si fida di lui perché sa che in questi casi
``diventa una bestia''. ``Te la puoi prendere con me quanto vuoi, ma il mio
team non lo tocchi.'' Lo gestisce direttamente con l'altro manager, per quanto
brutta possa diventare la conversazione. Lo sforzo più grande è \textbf{non
andarne poi a parlare con il team} (``ragazzi, ho fatto questa cosa''), perché
minerebbe la fiducia verso l'alto: non vuole un rapporto ottimo fino a lui e
pessimo oltre. L'obiettivo è che l'altro manager ``si faccia perdonare''
cambiando il proprio comportamento, e che il team percepisca da solo il
cambiamento.

\paragraph{Chiusura della giornata.} Il relatore chiede un feedback collettivo
in chat (``se domani ci dai \emph{more of the same} siamo contenti, oppure
ditemi cosa limare''), si scusa per i ripetuti crash della connessione e
ringrazia per la proattività nelle domande, che gli servono anche come
\emph{check}: se nessuno chiede nulla su un argomento, significa che può
andare avanti. Racconta di aver passato una settimana a decidere che cosa
mettere nelle slide e di averle poi preparate in poco tempo (rifacendone
alcune fino alle tre di notte). Ricorda infine di usare il test di
personalità per sé e per il proprio gruppo, facendolo fare a chi non l'ha
ancora fatto: tornerà molto utile per cooperare meglio.

% ---------------------------------------------------------------------
\section{Anticipazioni sulla seconda giornata}
\label{sec:z1-domani}

La seconda giornata del workshop sarà più pratica: ``oggi vi ho spiegato
perché, domani vi insegno come lo fanno queste aziende''. Gli argomenti
previsti:
\begin{description}
  \item[Time management] come organizzarsi una giornata professionale in
        azienda: come strutturare una \emph{task list}, cosa fare e cosa non
        fare, come e cosa prioritizzare.
  \item[Brainstorming] i gruppi dovranno fare moltissimi meeting (come facciamo
        il gioco, come lo chiamiamo, che scopo gli diamo). Se non si sanno fare
        bene, i brainstorming sono \textbf{dannosi ed eterni}: un brainstorming
        dovrebbe durare \textbf{al massimo due ore}. Verranno mostrati metodi
        che accorciano la durata, massimizzano l'effetto e, soprattutto,
        \textbf{fanno partecipare tutti}: non c'è niente di peggio di un
        brainstorming in cui parlano in due per il 90\% del tempo. Il classico
        ``ragazzi, non esistono idee sbagliate, sparate'' è il modo peggiore di
        farlo; non è colpa di nessuno, ma si può fare diversamente.
  \item[Guidare da dentro] come essere leader nel proprio team anche senza
        esserne il leader: sviluppare responsabilità personale e verso gli
        altri. Non essere il manager di qualcuno non significa non esserne
        responsabili: la persona accanto è responsabilità di tutti, perché si
        è un team e non un insieme di individui. ``Se qualcuno cade, tutti lo
        facciamo rialzare: non c'è l'addetto al rialzamento.'' È la versione
        \emph{advanced} della giornata: come comportarsi da manager senza
        averne il titolo.
\end{description}
Torneranno inoltre il \emph{service level agreement} e il tema della
condivisione dei successi e della protezione del team dai fallimenti.

% ---------------------------------------------------------------------
\section*{Punti chiave del capitolo}
\addcontentsline{toc}{section}{Punti chiave del capitolo}

\begin{itemize}
  \item ``La capacità di un gruppo di persone di fare cose straordinarie
        dipende dal modo in cui queste persone operano insieme.'' Il 90\% dei
        problemi dei team sono \textbf{problemi di comunicazione}.
  \item Le aziende si liberano sempre più presto dello \textbf{useless
        genius}: le competenze si trasmettono, il carattere no. Il
        \textbf{fit culturale} va valutato anche dal candidato.
  \item Non esiste un unico modo di dare feedback: dipende dal contesto e
        soprattutto da \textbf{chi lo riceve}. Il feedback è una somma di
        piccoli miglioramenti periodici; tenersi le cose nel cassetto non
        serve mai.
  \item I team attraversano \textbf{forming, storming, norming, performing e
        adjourning}. Lo storming non si salta; si passa la maggior parte del
        tempo in forming; la bravura sta nell'accelerare i passaggi e nel
        fare bene il \textbf{post-mortem}, anche dei successi.
  \item Gli obiettivi di un progetto vanno fissati \textbf{all'inizio}: solo
        così si può valutare se è andato bene e scegliere tra A e B.
  \item Il circolo virtuoso è \textbf{feedback $\rightarrow$ miglioramento
        $\rightarrow$ fiducia}. Accettare un feedback non significa fare come
        dice l'altro: decide chi ha la responsabilità, ma deve argomentare.
  \item ``Aiutami a poterti aiutare'': meglio dire in anticipo ``potrei non
        farcela'' che ``non ce l'ho fatta''. Chiedere aiuto non è una
        debolezza.
  \item Quando si dà feedback \textbf{non si porta la soluzione}: la si
        costruisce insieme a chi lo riceve.
  \item La cultura è forte quando si lavora \textbf{per gli altri}; il lavoro
        è valutato come opera d'insieme (anche con bonus di team).
  \item Come in \emph{Among Us}, i team sprecano energie a cercare il
        traditore invece di fare le task. \textbf{Default to trust}: chi ha
        sbagliato ha agito con le informazioni che aveva. Cercare
        \textbf{la causa, non il colpevole}: ``che cosa ti ha fatto pensare
        che fosse la scelta giusta?''.
  \item Carriere e performance non sono lineari: vanno valutate le
        \textbf{traiettorie}, non un punto dell'anno. Le valutazioni puntuali
        generano \emph{lying, hiding, faking}.
  \item Confermare sempre le istruzioni facendosele ripetere; ``non ho
        capito'' non è una colpa.
  \item Dare feedback: sull'\textbf{azione} e sui suoi \textbf{effetti}, specifico,
        presto, uno a uno, adattato al carattere, senza terze parti, senza
        ``non mi piace'', senza offendersi se non viene seguito; poi fare il
        debriefing.
  \item Ricevere feedback: \textbf{ascolta, chiedi dettagli, accetta, agisci}.
        Il feedback positivo serve a non peggiorare ciò che si fa bene; il
        feedback accidentale e il \emph{drive-by} non servono a nulla.
  \item Feedback \textbf{frequenti e recenti}: tanti piccoli \emph{sanity
        check} costano meno di una grande correzione di rotta.
  \item Il feedback è \textbf{un regalo} di chi ha a cuore il nostro
        interesse; fa paura il silenzio, non le critiche.
\end{itemize}

%    % ... altre lezioni zero ...
%    \part{Il corso di Game Design}
%    \renewcommand{\thechapter}{\arabic{chapter}}
%    \setcounter{chapter}{0}
%    % =====================================================================
%  Capitolo 1 — Introduzione al corso e al mestiere del game designer
%  Fonte: trascrizione della lezione GD1 (Game Design, primo anno)
%  Pacchetti richiesti nel preambolo del libro:
%    \usepackage[italian]{babel}, \usepackage[T1]{fontenc},
%    \usepackage[utf8]{inputenc} (solo pdfLaTeX), \usepackage{booktabs}
% =====================================================================

\chapter[Il mestiere del game designer]{Il mestiere del game designer: introduzione al corso}
\label{cap:introduzione-game-design}

Questo primo capitolo fa da soglia. Prima di affrontare metodi, framework e
documenti, è necessario capire \emph{in che modo} si impara il game design,
\emph{con chi} lo si fa, \emph{con quali strumenti} e, soprattutto,
\emph{che cosa fa concretamente} un game designer nel suo lavoro quotidiano.
Il capitolo segue l'impostazione della lezione inaugurale del corso: si parte
dal contesto formativo (il corso come simulazione di lavoro), si passa al
metodo di apprendimento (l'\emph{inquiry-based learning}), si elencano
letture e risorse di riferimento, e si arriva infine a una prima definizione
operativa del ruolo del designer, dei limiti entro cui lavora e delle sue
possibili specializzazioni.

Una premessa terminologica: il game design è una disciplina giovane e quasi
tutta la sua letteratura è scritta in inglese. Per questo motivo molti termini
tecnici verranno lasciati in lingua originale (\emph{pipeline},
\emph{pitching}, \emph{tuning}, \emph{feature team}, \dots): sono i termini che
si useranno nel lavoro, ed è importante imparare a riconoscerli e a usarli
correttamente fin da subito.

% ---------------------------------------------------------------------
\section{Il contesto: un corso come simulazione di lavoro}
\label{sec:simulazione-lavoro}

\subsection{Il principio fondamentale}

Il primo concetto da fissare è che il corso \textbf{si presenta come una
simulazione del lavoro}. Dal momento in cui si entra in aula si viene trattati
come \emph{potenziali professionisti dell'industria}: non soltanto studenti,
ma lavoratori effettivi di uno studio (l'accademia), con una gerarchia, dei
referenti, dei compiti e delle scadenze.

Ne derivano alcune conseguenze pratiche:
\begin{itemize}
  \item il docente non assegna un ``compitino'' da risolvere per accumulare
        punti o presenze: ogni \emph{task} è pensato per far dedicare tempo
        all'apprendimento di qualcosa di concreto;
  \item non consegnare un task non danneggia il docente (che ``vive ancora
        domani e viene pagato ugualmente''), ma lo studente, che perde
        l'occasione di mettere le mani su un argomento e diventare più abile;
  \item docenti e studenti devono trattarsi \textbf{da colleghi}, perché è
        molto probabile che un domani lo diventino davvero.
\end{itemize}

\subsection{Perché l'impegno conta più che mai}

Il lavoro del game designer è descritto come \emph{molto romantico, ma anche
molto duro}, e sempre più duro col passare degli anni. Durante il periodo del
Covid l'industria ha visto grandi investimenti; quei ``rubinetti'' sono stati
chiusi da tempo, il rapporto con i publisher è diventato più complesso e le
opportunità lavorative non sono più quelle di tre o quattro anni fa. Questo
significa:
\begin{itemize}
  \item più competizione per ogni posizione;
  \item aspettative più alte: ci si aspetta che i candidati siano ancora più
        preparati e sicuri delle proprie competenze;
  \item un unico modo per diventarlo: \textbf{lavorare}. Più tempo si dedica
        oggi ad apprendere il mestiere, più facilmente si sarà assumibili
        domani e più agevole sarà il lavoro di tutti i giorni.
\end{itemize}

\subsection{Sbagliare qui per non sbagliare fuori}

Il mantra del corso è: \textbf{quanto più provate qui, anche in modo goffo e
sbagliando, tanto meno sbaglierete là fuori}. In accademia, per quanto la
simulazione sia realistica, nessuno toglie lo stipendio e nessuno licenzia da
un momento all'altro. Nell'industria sì, e un errore può costare molto denaro
all'azienda, ancora di più se l'azienda è la propria.

\subsection{Il valore del fallimento}

Il fallimento, se affrontato con la mentalità corretta, è probabilmente
\emph{la cosa più vantaggiosa} che si possa fare durante la formazione:
\begin{itemize}
  \item un successo (``avete superato questa deadline'') conferma le proprie
        abilità, ma insegna relativamente poco;
  \item dopo $n$ fallimenti si impara moltissimo, e i successi arrivano;
  \item è fondamentale non demoralizzarsi: il feedback negativo non è mai un
        giudizio sulla persona, ma su ciò che è stato consegnato.
\end{itemize}

Il docente riconosce apertamente che dedicare tempo ed energia a un progetto e
vederlo rifiutato, o vederlo non uscire mai, è demotivante: è successo a lui e
lo ha visto accadere anche su progetti molto grandi, con team numerosi. Il
punto però non è ciò che si è ottenuto materialmente, ma \textbf{quanto si è
cresciuti}. Da ogni esperienza ci si portano a casa due cose:
\begin{enumerate}
  \item \textbf{l'esperienza}: quanto si è diventati bravi, veloci,
        praticamente abili nell'applicare un framework, quanto si è cresciuti
        come designer;
  \item \textbf{il rapporto umano}: nel lavoro si ``suda sangue'' fianco a
        fianco con altre persone, e questo crea legami molto forti che
        portano avanti l'industria nel senso più positivo del termine: nascono
        progetti privati, ci si porta da un'azienda all'altra, si formano team
        con basi solide.
\end{enumerate}

Va evitato anche l'eccesso opposto: la cultura del \emph{grinding} si è
appropriata di questi discorsi in modo distorto. Non bisogna diventare robot.

% ---------------------------------------------------------------------
\section{Organizzazione del corso e gerarchia}
\label{sec:gerarchia-corso}

\subsection{La struttura gerarchica}

La struttura dei referenti del corso ricalca volutamente una
\textbf{gerarchia tipica dell'industria}:
\begin{itemize}
  \item al vertice c'è il \textbf{direttore del corso} (Diego Ricciuti), che
        corrisponde al \emph{creative director}: supervisiona la visione del
        progetto (il concetto di ``visione'' verrà approfondito più avanti);
  \item sotto di lui ci sono i \textbf{referenti di dipartimento}, uno per
        disciplina, con un corrispettivo per ciascuna sede (Milano e Roma);
  \item in fondo ci sono gli studenti, che occupano il primo gradino di
        un'azienda: l'\emph{intern} (stagista) o il \emph{junior}, il passo
        che precede il lavoratore ``regolare''.
\end{itemize}

\begin{table}[htbp]
  \centering
  \begin{tabular}{@{}p{0.34\textwidth}p{0.29\textwidth}p{0.29\textwidth}@{}}
    \toprule
    \textbf{Disciplina} & \textbf{Milano} & \textbf{Roma} \\
    \midrule
    Direzione (\emph{creative director}) & \multicolumn{2}{p{0.58\textwidth}}{Diego Ricciuti} \\
    Game design      & Mattia Podini        & Giacomo ``Jack'' Danna \\
    Level design     & Stefano Zocchi       & Andrea Dispoli \\
    Narrative design & \multicolumn{2}{p{0.58\textwidth}}{Daniele Fusetto (lezioni condivise), Giacomo Masi} \\
    Programmazione (Unreal) & Alessandro Pregnolato & Matteo Velardi \\
    \bottomrule
  \end{tabular}
  \caption{Referenti del corso per disciplina e sede (nomi come riportati in lezione).}
  \label{tab:referenti}
\end{table}

Per il primo anno, il referente diretto degli studenti di game design è il
docente di game design: è a lui che ci si rivolge per dubbi e domande. Altre
figure utili sono il direttore didattico (severo, ``giustamente severo'', da
contattare solo per questioni particolarmente importanti), l'amministrazione
(contratti e documentazione) e la segreteria.

\subsection{I canali di comunicazione}

\begin{description}
  \item[Discord] è il canale per le comunicazioni ufficiali del corso e per
        le domande che non richiedono urgenza immediata; può avere qualche
        ritardo.
  \item[E-mail] è il canale per tutto ciò che è formale o ``ancora più
        ufficiale'' e per le comunicazioni rapide e immediate.
  \item[Telegram] esistono residui di vecchi canali precedenti a Discord, da
        evitare.
\end{description}

Le domande che valgono per tutti vanno fatte in classe (anche da un portavoce
a nome del gruppo); quelle singole e personali via Discord o, se formali, via
e-mail. Se il docente non risponde a una domanda c'è un motivo, ad esempio
ragioni legali, e lo dirà esplicitamente.

\subsection{I tre moduli e la settimana tipo}

Le tre discipline principali del primo anno (e non solo) sono:
\begin{itemize}
  \item \textbf{Game Design} (tendenzialmente il lunedì);
  \item \textbf{Level Design} (tendenzialmente il mercoledì);
  \item \textbf{Programmazione} (tendenzialmente il venerdì).
\end{itemize}
Nel primo mese le lezioni sono tutte di game design (lunedì, mercoledì e
venerdì). Sono previste anche lezioni di narrativa, ma il piano formativo si
aggiorna ogni anno: le ore di narrativa sono state ridotte a favore di quelle
di programmazione. Le poche lezioni di narrativa daranno quindi dei
\emph{cenni}, non la preparazione di un narrative designer. Ciò che serve
davvero è imparare \textbf{come comunicano gli altri dipartimenti} e come
lavorare con persone che fanno un lavoro diverso ma complementare al proprio.

\subsection{Le due anime delle lezioni di game design}

\begin{description}
  \item[Design methods (metodologie)] lezioni prevalentemente teoriche, ma con
        una buona componente pratica, in cui si studiano metodologie di design
        e i cosiddetti \emph{framework} dell'industria.
  \item[Documentazione e pitching] il cuore del lavoro quotidiano del
        designer.
\end{description}

\paragraph{La documentazione.}
Il docente si definisce ``il più grande nerd della documentazione'' che si
possa incontrare. È una parte dura, perché a pochi piace documentare, ma il
lavoro del designer \emph{è principalmente questo}: scrivere documenti e, da
quei documenti, passare alla pratica e all'implementazione. Chi non sa
documentare bene parte male. Quando non si è presenti, il documento è il
proprio \textbf{biglietto da visita}: qualcun altro lo leggerà, quindi va
scritto nel modo più chiaro e comprensibile possibile.

Un esempio concreto del perché la documentazione conta: CD Projekt Red ha
sviluppato un proprio engine proprietario, ma buona parte della sua
documentazione è stata ``buttata all'acqua''. Chi deve usarlo o lo sa già
usare, o ha competenze di un altro engine da trasferire, oppure è in
difficoltà. Si documenta \textbf{perché altri possano leggere}.

\paragraph{Il pitching.}
Il \emph{pitching} è l'esposizione, la presentazione orale. La parte creativa
del team presenta spesso:
\begin{itemize}
  \item le proprie idee ai colleghi;
  \item il gioco a convention ed eventi;
  \item il pitch del gioco agli investitori, per ottenere i fondi necessari a
        realizzarlo.
\end{itemize}
Saper parlare e comunicare chiaramente le idee è fondamentale. In sintesi:
\emph{il game designer parla}, e bisogna farsene una ragione.

\subsection{Altre componenti del corso}

\begin{description}
  \item[Gioco critico e playtest] la classe viene divisa in due gruppi: uno
        gioca, l'altro osserva. Chi osserva non applaude ``come le
        scimmiette'': annota, segna e analizza il \emph{playtest}, esattamente
        come si fa nell'industria. Il playtest fa parte delle mansioni del
        designer, ma esiste anche come specializzazione autonoma.
  \item[Pipeline] termine che indica tutto ciò che riguarda i processi
        formalizzati dell'industria; di solito si intende la \emph{pipeline di
        produzione}: come si divide il team, qual è la sua gerarchia, come il
        \emph{producer} gestisce le deadline. In breve: \emph{come} e
        \emph{quando} si fanno le cose. Al primo anno è meno rilevante (si
        dovrà ``nuotare'' un po' da soli); dal secondo anno si vedranno le
        metodologie effettivamente applicate ai progetti digitali. Si
        scoprirà quanto sia difficile lavorare a un progetto digitale, e
        ancor più fisico, senza un piano di sviluppo comprovato, come quelli
        usati per qualsiasi software in qualsiasi parte del mondo.
  \item[Engine] il motore di riferimento è \textbf{Unreal Engine}
        (sezione~\ref{sec:engine}).
  \item[Programmazione] principalmente tramite \emph{Blueprint} in Unreal,
        più eventuali sistemi di versionamento; il corso parte da zero.
  \item[Progetti pratici] probabilmente la cosa più importante. Al primo
        anno i progetti di gruppo restano circoscritti al corso di design,
        perché mancano ancora alcune nozioni (come comunicare, affrontare
        problemi, dare e ricevere feedback). In seguito si allargheranno a
        progetti digitali, molto più difficili di quelli fisici, con membri
        di altri team: un rapporto non unilaterale, ma ``covalente''.
\end{description}

% ---------------------------------------------------------------------
\section{Il metodo: l'inquiry-based learning}
\label{sec:inquiry}

\subsection{Che cos'è}

Il corso è costruito secondo l'\textbf{inquiry-based learning}
(apprendimento basato sull'indagine), un approccio pedagogico nato intorno
agli anni Sessanta, che si richiama a pensatori come Dewey e Piaget. In
generale, le discipline da cui il design attinge sono quasi sempre la
sociologia, la psicologia e, talvolta, l'economia.

Il funzionamento è un ciclo:
\begin{enumerate}
  \item il docente fornisce metodologie e framework dell'industria durante
        la lezione;
  \item lo studente, fuori dalle ore di lezione, lavora sui \emph{task}:
        mette a terra le nozioni, le maneggia, le assorbe meglio;
  \item da questo lavoro nascono nuovi dubbi;
  \item i dubbi tornano in aula: lo studente chiede, il docente risponde
        (ammesso che dare la risposta sia utile, e ammesso che una risposta
        esista).
\end{enumerate}

\subsection{Non sempre esiste una risposta giusta}

Non tutto ciò che si studia ha ``la risposta migliore possibile in assoluto
nel vuoto''. Esistono metodologie, limiti e framework che aiutano il lavoro,
ma non sono fissi: cambiano da azienda ad azienda, da lavoro a lavoro, da
progetto a progetto, e a volte vengono stravolti. Sta alla
\textbf{sensibilità del designer} capire quale strumento o framework usare in
un dato momento, e questa sensibilità si impara solo con l'esperienza. Nel
lavoro nessuno dirà ``usa questa cosa'' o ``è meglio quest'altra'': sarà il
designer a scegliere.

Anche l'esperienza del docente ha confini: sei-sette anni di lavoro, un modo
personale di approcciare il design e un certo insieme di progetti (ad esempio
nessun MMO). Oltre quei confini si cercano risposte nella ricerca, studiando
come aziende con esperienza specifica (ad esempio Riot) hanno affrontato
problemi simili e hanno condiviso le loro soluzioni, sempre valutando quali
metodi siano adatti alla dimensione del team, al tipo di progetto e al tempo
disponibile.

\subsection{Sensibilità, non memoria}

Non interessa imparare a memoria un framework o la teoria di un autore. Certe
terminologie vanno conosciute, ma per ragioni di comunicazione e chiarezza.
Ciò che conta è sviluppare la \textbf{capacità di farsi domande}, di farle
agli altri, di fare ricerca e di trovare una propria risposta. È esattamente
ciò che si fa nel lavoro: davanti a qualcosa che non si sa fare, l'abilità
del designer sta nel cercare come qualcun altro ha risolto un problema simile.

\subsection{Le lezioni come corsi di aggiornamento}

Le lezioni frontali simulano i \emph{corsi di aggiornamento} aziendali: le
sessioni che un lead di design tiene al suo team di design (o il lead di
level design al suo team, e così via). Alla fine dell'aggiornamento c'è quasi
sempre un lavoro da fare, un obiettivo da portare a terra:
\textbf{l'obiettivo della settimana}. Si lavora di settimana in settimana, e
il perché verrà chiarito in seguito.

Struttura tipica di una lezione:
\begin{enumerate}
  \item breve sessione di domande e risposte sulla lezione precedente;
  \item analisi di eventuali esercizi particolarmente critici;
  \item lezione vera e propria;
  \item discussione di gruppo, se c'è tempo;
  \item assegnazione del task della settimana.
\end{enumerate}
Il docente tende a ``parlare e straparlare'', ad aprire parentesi e a saltare
da un argomento all'altro: l'unico modo per fermarlo è che siano gli studenti
a interromperlo con domande.

\subsection{La gerarchia della competenza}

L'inquiry-based learning può essere rappresentato come una scala di
competenza:
\begin{enumerate}
  \item la sola lezione frontale porta a \textbf{conoscenze incomplete},
        perché manca la parte pratica e di ricerca;
  \item il lavoro a casa (task, ricerca) porta a una competenza \textbf{più
        completa}, non sbagliata in sé, ma ancora parziale e complementare ad
        altre;
  \item la \textbf{condivisione} dei task e dei feedback con gli altri porta
        idealmente alla competenza più completa possibile;
  \item il picco della competenza è la sua \textbf{applicazione nei
        progetti} (fisici o digitali, di gruppo): ``con le cose che avete
        imparato e ricercato, riuscite a fare un gioco?''. Quando la risposta
        diventa positiva, o comunque si inizia a navigare nei processi di
        design per realizzarne uno, quella è la vittoria.
\end{enumerate}
Non si parte subito creando giochi: si parte dall'analisi, dal design e dalle
regole generali di design, prima ``nel vuoto'', poi contestualizzando i
framework, infine unendo i puntini nei progetti.

Rimarranno comunque delle lacune, per tutti, docente compreso, perché il game
design \textbf{non è una scienza esatta}, anzi, a rigore non è affatto una
scienza. Prende in prestito metodologie dalle discipline scientifiche, ma
contiene una componente imprevedibile e incalcolabile: \textbf{la
produzione}. Si possono applicare perfettamente le regole di un framework, e
quel framework può comunque non essere adatto al gioco, al proprio modo di
sviluppare o al tempo di produzione disponibile (alcuni framework richiedono
una ricerca e un'attenzione che il tempo non concede).

C'è poi la componente creativa, la parte ``romantica'' del lavoro: esiste, ma
è \textbf{metà del lavoro}. L'altra metà è ricerca, metodo scientifico al
100\%. Dagli stessi dati $n$ designer diversi, con sensibilità diverse,
possono ottenere $n$ risposte diverse. Non tutte saranno valide: vanno
applicate, messe a terra, e questa scrematura fa parte del lavoro.

\subsection{I livelli dell'apprendimento}

Un secondo modello, a gradini, descrive i livelli di apprendimento (è la
cosiddetta tassonomia SOLO):
\begin{description}
  \item[Pre-strutturale] il docente dà il compito e lo studente lo risolve,
        possibilmente il prima possibile. È esattamente ciò che \emph{non} si
        vuole.
  \item[Uni-strutturale] si affronta una sola disciplina: è il primo passo
        necessario, ed è il motivo per cui nel primo mese ci sono solo lezioni
        di game design.
  \item[Multi-strutturale] entrano in gioco altri dipartimenti: level design,
        game design, programmazione e così via.
  \item[Relazionale] si capisce che alcune cose valgono per più
        dipartimenti: alcuni framework di narrativa derivano dal design o
        possono essere usati per fare design; le regole di documentazione si
        applicano anche ad altri dipartimenti; le regole di level design si
        mescolano con quelle di game design. Dopotutto, per fare un gioco
        servono molti dipartimenti.
  \item[Astratto esteso] il punto finale: muoversi autonomamente e capire
        cosa usare e quando, sulla base della propria esperienza. L'unico modo
        per arrivarci è fare pratica e fare progetti.
\end{description}

% ---------------------------------------------------------------------
\section{Task, feedback e regole del corso}
\label{sec:task-feedback}

\subsection{I task come opportunità}

I task non sono strettamente obbligatori (esiste un numero minimo di
consegne), ma sono \textbf{opportunità}. Quanto e come ci si lavora dipende
dallo studente. Un task fatto in cinque minuti, o in cui si è stati ``furbi'',
non verrà letto né riceverà feedback.

Domande da porsi davanti a ogni task:
\begin{enumerate}
  \item \textbf{Perché} mi è stato assegnato? Qual è l'obiettivo?
  \item Come ci lavoro per ottenere un risultato \textbf{che abbia senso}? Se
        la ricerca e i dati indicano una direzione, è inutile andare in
        un'altra.
\end{enumerate}

\subsection{Quanto lavorare: un calcolo indicativo}

Se il corso è una simulazione di lavoro, vale il carico di un lavoro normale:
\begin{itemize}
  \item un lavoro a tempo pieno comprende da 30 a 45 ore settimanali; si
        considera una media di \textbf{40 ore};
  \item si sottraggono circa 9 ore di lezione: restano $\approx 31$, diciamo
        \textbf{30 ore};
  \item con circa 3 task a settimana: $30 / 3 = $ \textbf{circa 10 ore per
        task}.
\end{itemize}
Il numero può variare, ma rende l'idea: dedicare un'ora a un task e poi fare
altro non fa guadagnare nulla a nessuno. Lo studente ``butta soldi'', il
docente risparmia tempo leggendo il task in un minuto (o non leggendolo).
Dedicando invece attenzione, ricerca, cura nella formattazione e
sperimentazione, e facendolo con cognizione di causa, si impara.

\subsection{Niente voti, solo feedback}

Ad eccezione del corso di programmazione, non ci sono voti (né ``ottimo'', né
10, né 100). Ci sono \textbf{feedback}:
\begin{itemize}
  \item in game design, generalmente sulla documentazione;
  \item in level design, sui progetti pratici.
\end{itemize}
Il docente indica cosa va bene e cosa va modificato, concentrandosi in
particolare su ciò che va cambiato. Il feedback \textbf{non è mai un giudizio
sulla persona}, ma su ciò che è stato consegnato nel concreto.

\subsection{La condivisione dei task}

Tutti i task consegnati vengono \textbf{condivisi} con la classe, e possono
essere letti solo dopo aver ricevuto i feedback del docente. Si cresce così
non solo con i propri feedback, ma anche con quelli dati agli altri. Regola
importante: ciò che il docente dice a uno studente \textbf{vale per tutti}. Se
a una persona viene detto ``questa cosa non va mai più fatta'', da quel
momento vale per l'intera classe.

\subsection{Non copiare}

Copiare i task degli altri è una delle cose peggiori che si possano fare, non
per una questione morale, ma perché annulla il senso del percorso. Si può
dedicare a ogni task il tempo e l'attenzione che si ritengono opportuni, ed è
normale che col tempo ci si specializzi (preferendo, ad esempio, il level
design o la programmazione tecnica) e si investa meno dove si sente meno
necessità. In questi casi l'alternativa corretta è:
\begin{itemize}
  \item non consegnare;
  \item studiare ciò che hanno fatto gli altri;
  \item eventualmente riconsegnare in seguito.
\end{itemize}
Ci saranno momenti in cui si verrà sommersi dalle cose da fare: bisogna
imparare a scegliere.

\subsection{Il ``non leggo''}

Il docente ha piena facoltà di non dare alcun feedback a un task (il
cosiddetto \textbf{``non leggo''}) se:
\begin{itemize}
  \item non si è lavorato con attenzione;
  \item non si rispettano le indicazioni date;
  \item è stata violata una regola estremamente importante.
\end{itemize}
Un task ``non letto'' equivale a un task non consegnato e non entra nel
conteggio. Attenzione: un task \emph{perfetto} e un task che \emph{fallisce}
nel comprendere o applicare un framework sono \textbf{ugualmente validi},
purché ci si sia impegnati e si sia dedicato tempo, pazienza e ``sudore''.

\paragraph{Riconsegna.} Un task non letto, in ritardo o che si vuole
migliorare può essere riconsegnato e verrà conteggiato, a patto che:
\begin{itemize}
  \item i problemi segnalati siano stati risolti;
  \item sia stato fatto ulteriore lavoro, integrando anche tutto ciò che nel
        frattempo è stato detto agli altri colleghi nei loro feedback.
\end{itemize}
Il task va sviluppato con le competenze che si hanno nel momento della
riconsegna, non con quelle del momento della prima consegna.

\subsection{Requisiti per superare l'anno}

\begin{itemize}
  \item consegnare almeno l'\textbf{80\% dei task} assegnati (soglia
        malleabile);
  \item partecipare almeno all'\textbf{80\% delle lezioni}, in presenza o
        da remoto, preferibilmente in presenza; i ritardi vengono valutati
        caso per caso;
  \item comportarsi da persone corrette e rispettose (è stata firmata una
        carta dello studente): non urlare, non alzare le mani, non
        vandalizzare, non compiere atti illegali.
\end{itemize}
Impegnarsi di più porta inoltre a essere selezionati per varie opportunità.

\subsection{Scadenze}

\begin{itemize}
  \item la scadenza di un task è \textbf{sempre prima dell'inizio della
        lezione} successiva indicata;
  \item se la classe, come team, mantiene comportamenti scorretti, il docente
        può \textbf{anticipare le deadline} per tutti.
\end{itemize}

\subsection{Accessibilità e DSA}

Gli studenti con disturbi specifici dell'apprendimento o neurodivergenze
(dislessia, ADHD, autismo, ecc.) sono invitati a segnalarlo. Quando si
affronteranno le regole di formattazione e documentazione (errori
grammaticali, mancanze, incoerenze) se ne terrà conto e si potranno trovare
dei compromessi, pur ricordando che nel mondo del lavoro nessuno farà
``sconti''. La documentazione può essere inviata all'amministrazione o alla
segreteria, che la comunicherà a tutti i docenti, ma il consiglio è parlarne
anche direttamente con ogni singolo docente.

% ---------------------------------------------------------------------
\section{Le competenze trasversali}
\label{sec:soft-skill}

\subsection{I giochi si fanno in gruppo}

Si possono fare giochi da soli, ed è un'impresa ammirevole; ma lavorando da
soli \textbf{tutto ciò che si guadagna in libertà si perde in produttività}.

\paragraph{Il caso Stardew Valley.} \emph{Stardew Valley} è stato realizzato
da una sola persona, Eric Barone, in anni di sviluppo e di sacrifici: faceva
un altro lavoro per mantenersi (in un cinema) e di notte programmava, faceva
l'arte e componeva la musica. Il risultato è stato straordinario e lo ha reso
ricco, ma si tratta di un caso estremamente raro. Quanti giochi simili, nati
nelle stesse condizioni, sono falliti? Molti di più.

\paragraph{Il rischio d'impresa.} Lavorare da soli non è sbagliato di per sé
e ha aspetti molto positivi: si ha il privilegio di creare esattamente ciò che
si vuole. Ma ci si fa carico anche di tutto il \textbf{rischio d'impresa}: se
si lavora a un gioco per cinque anni e poi il gioco va male, resta il
vantaggio di averlo pubblicato, ma il dopo è sperare di vivere di quella
release o di finanziare con essa il gioco successivo. Le grandi aziende hanno
gli stessi problemi, ma li distribuiscono su più persone e più responsabilità.

L'obiettivo di un professionista è \textbf{saper lavorare in diverse
condizioni}, e la più comune è quella del dipendente (o del freelance /
libero professionista con partita IVA): si viene assunti per le proprie
capacità da un'azienda, piccola o grande che sia, e si lavora con altre
persone.

\subsection{Essere persone con cui si vuole lavorare}

Se le persone non vogliono lavorare con voi, quello è \textbf{il primo motivo
di licenziamento}. E se non si viene licenziati, il clima nell'azienda
diventa terribile: la produttività crolla, la motivazione crolla. Le
\emph{soft skill}, cioè l'insieme di abilità che rendono una persona
piacevole con cui lavorare, vanno coltivate e diventano sempre più importanti
col passare degli anni (relativamente meno ora, molto di più quando i team
saranno più grandi). Anche chi è un ``lupo solitario'' deve imparare a:
\begin{itemize}
  \item lavorare con altre persone;
  \item comunicare con altre persone;
  \item trovare un punto in comune da cui partire per fare le cose.
\end{itemize}

\subsection{Indipendenza e proattività}

Sono le altre due parole chiave del corso:
\begin{description}
  \item[Indipendenza] spesso si verrà lasciati da soli, non per cattiveria ma
        per sfida: quando si cerca una risposta da soli, ciò che conta non è
        tanto la risposta quanto \textbf{il processo}. Una risposta data dal
        docente non vale niente.
  \item[Proattività] chiedere, domandare, agire da soli senza aspettare
        direttive, perché spesso quelle direttive non ci sono, a maggior
        ragione nei team piccoli.
\end{description}
Tutto questo si concretizza nel parlare, nel presentare e nel fare
\emph{post-mortem} (l'analisi a posteriori di un progetto concluso, che verrà
approfondita più avanti), e in generale nell'attenzione alla crescita
personale.

\subsection{Fare domande}

Il corso insiste sulle domande: alzare la mano, intervenire anche ``a
bomba'', fermare il docente ogni volta che qualcosa non è chiaro. Se una
persona non conosce qualcosa, molto probabilmente c'è almeno un'altra persona
che non la conosce. Il docente può dare per scontata la conoscenza di termini
che non lo sono, ed è del tutto normale, soprattutto all'inizio.

% ---------------------------------------------------------------------
\section{Strumenti: engine, Unreal e materiali}
\label{sec:engine}

\subsection{Unreal Engine e gli altri motori}

Il motore di riferimento del corso, usato in level design, è \textbf{Unreal
Engine}: ``un gran mattone'', ma che dà a un designer molte possibilità. Il
vantaggio principale è che permette a un designer non pratico di
programmazione di \textbf{prototipare tramite i Blueprint} (il sistema di
scripting visuale di Unreal). Questo non impedisce di passare a Unity o ad
altri engine, e viceversa: il docente stesso ha studiato e lavorato a lungo
con Unreal, ma la sua ultima esperienza professionale (come level designer) è
stata su Unity, un engine che non conosceva, a dimostrazione che il salto da
un engine all'altro è possibile.

Panorama dei motori citati:
\begin{itemize}
  \item \textbf{Unreal Engine} e \textbf{Unity}: i due più diffusi, con
        modelli di sorgente ``più o meno'' aperto;
  \item \textbf{Godot}: engine open source;
  \item \textbf{CryEngine}: proprietario di Crytek;
  \item \textbf{Snowdrop}: engine di Ubisoft;
  \item Rockstar usa un proprio engine proprietario;
  \item alcune aziende usano versioni \textbf{personalizzate} di engine
        esistenti: ad esempio Milestone, nei giochi di corse, usa un Unreal
        customizzato;
  \item CD Projekt Red ha un engine proprietario (con i problemi di
        documentazione già citati).
\end{itemize}
In sintesi: ogni azienda può avere un engine proprietario oppure una versione
modificata di un engine commerciale, ma si parte sempre da una base.

\subsection{Programmazione: cosa ci si aspetta da un designer}

Non è richiesto alcun bagaglio pregresso di programmazione: il corso parte da
zero. Gli studenti sono stati selezionati (il corso è a numero chiuso) per
competenze analitiche e critiche sopra la media. Anche nel mercato del lavoro,
per una posizione di game designer generalmente \textbf{non ci si aspetta
competenza di scripting} (C++, Java, C\#, ecc.), a meno di voler fare il
\emph{technical designer}, e anche lì varia da azienda ad azienda.

\subsection{Versione e preparazione}

Indicazioni pratiche emerse in aula:
\begin{itemize}
  \item scaricare Unreal Engine; la versione consigliata per il corso è
        stata indicata nella \textbf{5.6} (con conferma successiva). La
        stabilità dipende da cosa si deve fare, da quali tool e plugin si
        usano: i problemi nascono spesso dai plugin. Nel passaggio da Unreal 4
        a Unreal 5, ad esempio, il sistema \emph{Landscape} era inizialmente
        inutilizzabile, mentre in Unreal 4 funzionava molto bene;
  \item è consigliato installare tramite \textbf{Epic Games Launcher}, più
        comodo: permette anche di riscattare giochi gratuiti periodici e di
        accedere ai pacchetti di \textbf{Fab} (il marketplace di asset
        acquisito da Epic), che ogni mese offre alcuni asset gratuiti utili
        per esercitarsi nel level design (ad esempio per fare \emph{speed
        level design});
  \item molti pacchetti di UI o di sistemi di interazione sono utili per
        prototipare al volo, ma spesso inutilmente complicati;
  \item un plugin di \textbf{materiali di prototipazione a griglia} (adattivi,
        con metriche, allineabili alla griglia, con materiali già pronti e
        colori modificabili, dal costo di circa 5 euro) è indicato come
        ``salvavita'' per il level designer; se necessario, i pacchetti
        indispensabili vengono forniti;
  \item portare \textbf{sempre il portatile} alle lezioni di level design
        (da novembre, a partire dall'interfaccia dell'engine); alle lezioni
        di game design basta che lo porti almeno metà della classe, per i
        lavori a coppie. Serve comunque qualcosa per prendere appunti
        (``tavolette d'argilla e cannucce, se volete fare i babilonesi'',
        oppure carta e penna).
\end{itemize}
Chi ha un portatile poco potente dovrebbe comunque provare a installare
l'engine e aprire un template per vedere come gira: all'inizio non si faranno
cose pesanti. Nelle prime lezioni di level design è importante mettere subito
le mani sull'interfaccia: ascoltare dove si trova uno strumento (ad esempio lo
\emph{snapping}) senza provarlo porta a dimenticarlo.

% ---------------------------------------------------------------------
\section{Letture e risorse di riferimento}
\label{sec:risorse}

\subsection{I libri}

\begin{description}
  \item[\emph{The Art of Game Design: A Book of Lenses} (Jesse Schell)] il
        testo \textbf{davvero importante}. Alcune lezioni finali dell'anno si
        baseranno sulle sue terminologie. È utile perché apre la mente sulla
        possibilità di applicare metodologie diverse a giochi diversi, ma
        soprattutto per la \textbf{terminologia} comune. Esiste anche in
        audiolibro (circa 25 ore).
  \item[\emph{Challenges for Game Designers} (Brenda Romero e Ian Schreiber)]
        libro \emph{entry level} con moltissimi esercizi, molto caro al
        direttore del corso.
  \item[\emph{21st Century Game Design} (Chris Bateman)] utile come
        retrospettiva: è un libro visionario che ha previsto metodi e
        soprattutto \emph{archetipizzazioni} dei giocatori tuttora valide nel
        design e nel marketing. Servirà per una delle lezioni più importanti
        dell'anno. È utile leggerlo per vedere quanto aveva previsto, ma è
        utile ``fino a un certo punto''.
  \item[\emph{An Architectural Approach to Level Design} (Christopher
        Totten)] non strettamente necessario, ma di lettura semplice. Ha due
        anime: parte dall'architettura (babilonese, greca, ecc.) e la trasla
        su esempi di level design (come \emph{Doom} e \emph{Team Fortress}).
        Un tempo ne esisteva un riassunto su Gamasutra (oggi
        \emph{Game Developer}).
  \item[\emph{A Theory of Fun for Game Design} (Raph Koster)] utile fino a un
        certo punto: contiene diverse idee superate, che lo stesso Koster ha
        affrontato in vari talk reperibili nel GDC Vault. Verrà ripreso nelle
        prime lezioni, poi messo da parte.
\end{description}

\subsection{Canali e conferenze}

\begin{description}
  \item[GDC (Game Developers Conference)] il suo canale raccoglie talk e
        panel estremamente importanti: è ``uno strumento misterioso che ci
        servirà più tardi''. \textbf{Attenzione}: per molti talk mancano ancora
        le competenze per interpretarli, e alcuni potrebbero fuorviare o
        mettere in testa idee strane. La GDC si concentra spesso su
        produttività e mobile.
  \item[Il \emph{Dragon Speech} di Chris Crawford] il talk del fondatore della
        GDC: appartiene a un passato lontano, ma è molto motivante. Da vedere
        subito.
  \item[Game Maker's Toolkit] canale particolarmente utile per nozioni
        pratiche, con video brevi e alcune interviste.
  \item[Nordic Game] altra conferenza importante, con talk e panel di esperti
        del settore.
\end{description}

\subsection{Associazioni e premi: come leggerli}

Le associazioni di settore (ad esempio IIDEA in Italia, o la IGDA) non sono
il male, anzi: sono un canale importante per affacciarsi all'estero o per
accedere a fondi. Va però ricordato che \textbf{tutelano i propri associati}:
non esiste solo il loro palinsesto, e bisogna valutare quanto spazio occupano
nel panorama della visibilità.

Lo stesso vale per i \textbf{premi} (Game Awards, ``Best Mobile Game'',
``Games for Change'', ecc.): al designer \emph{non interessano}. Che un gioco
abbia vinto ``più di 90 premi'' non dice nulla di utile. Servono a dare
visibilità e lustro ai giochi ed è giusto che esistano, perché ampliano il
panorama e il numero di giocatori. Ma ciò che interessa al designer sono i
\textbf{dati}:
\begin{itemize}
  \item come sono state implementate certe cose;
  \item come è stato fatto il marketing e dove il gioco è stato presentato;
  \item non quante wishlist ha fatto, ma \emph{come} le ha fatte;
  \item in che periodo storico è uscito e quali erano i competitor in quel
        momento.
\end{itemize}

% ---------------------------------------------------------------------
\section{Che cosa fa un game designer}
\label{sec:cosa-fa}

Per arrivare a una prima definizione, la classe ha condotto un brainstorming:
ognuno doveva indicare una cosa che, secondo lui, fa il game designer (senza
ripetere quelle già dette). Ne è emersa la lista seguente, commentata punto
per punto.

\subsection{Progettare meccaniche e struttura del gioco}

Il ruolo più generico è quello del \textbf{core mechanic designer}, detto
anche \emph{gameplay designer} e, un tempo, \emph{black box designer}: si
occupa delle \textbf{meccaniche di gioco}. Si va dal piccolo al grande (in
funzione anche della dimensione dell'azienda):
\begin{itemize}
  \item nel \textbf{piccolo} ci sono le singole meccaniche;
  \item nel \textbf{grande} ci sono i \textbf{sistemi}. Nei giochi
        \emph{sistemici}, con molti sistemi interconnessi (ad esempio
        \emph{GTA}), esistono specializzazioni dedicate a bilanciare e
        supervisionare i sistemi: è il \emph{system design}.
\end{itemize}

\paragraph{Quando un gioco è ``sistemico''?}
Non dipende solo da come le meccaniche interagiscono con il giocatore (le
\emph{dinamiche}, secondo i framework che passano dalle meccaniche alle
dinamiche, che verranno studiati più avanti), ma anche da \emph{quante}
meccaniche ci sono e da quanto interagiscono tra loro. Non diremmo che
\emph{Super Mario} sia un gioco sistemico, eppure contiene dei sistemi: la
linea è sottile. Quando i sistemi aumentano molto, sono complessi e si
influenzano a vicenda, il gioco diventa sistemico. Esempio: in \emph{The
Legend of Zelda: Breath of the Wild} il vento modifica quanto il personaggio
sale o scende; il fuoco può incendiare l'erba; l'erba prende fuoco se colpita
da una freccia infuocata; e così via. Nel mercato moderno i giochi sistemici
sono spesso, ma non sempre, legati a mappe \emph{open world}.

\subsection{Le 3C: camera, character, control}

La base del gioco, da cui generalmente si inizia a lavorare (anche nei
prototipi), è la ``sacra trinità videoludica'' delle \textbf{3C}:
\begin{description}
  \item[Camera] che tipo di inquadratura ha il gioco: prima persona, terza
        persona, isometrica, ecc.
  \item[Character] quanto è grande il personaggio, che tipo di personaggio è
        e, di conseguenza, quali azioni può compiere.
  \item[Control] i controlli: che cosa il personaggio può fare, quali sono i
        \emph{verbi} di gioco.
\end{description}
Le 3C sono definite dal designer, ed esiste anche una specializzazione per
questo lavoro. Esistono varianti: quella del direttore del corso è
\textbf{4C + J}, cioè \emph{camera, character, control, conflict} (conflitto)
e \emph{juice} (il ``succo'', il feedback che rende piacevole l'interazione).
Verranno approfondite più avanti.

\subsection{Il personaggio}

Il designer definisce il personaggio dal punto di vista delle meccaniche
(dimensioni, azioni, capacità), all'interno delle 3C. L'aspetto estetico è
invece competenza di altri dipartimenti.

\subsection{Documentare le informazioni per gli altri dipartimenti}

Assolutamente vero: la documentazione è una delle prime attività del
designer. Più l'azienda è grande, più questo lavoro diventa prominente; più è
piccola, meno è rilevante. In \textbf{pre-produzione}, però, è sempre
fondamentale.

\subsection{Mediatore e leader ``al servizio''}

Il designer è un \textbf{mediatore} nei conflitti interni al team, una sorta
di collante che si relaziona con dipartimenti diversi (non sempre gli stessi,
a seconda della fase di produzione). A volte fa anche da leader, a seconda di
come sono gestite la pipeline e la gerarchia.

Questa leadership è difficile da definire perché \textbf{non è verticale}: è
una leadership \emph{al servizio} (\emph{servant leadership}). Il designer
``è asservito al programmatore'': il suo compito non è arrivare con un'idea e
scaricarla addosso agli altri, né portare la soluzione e lasciare il problema
a qualcun altro, ma \textbf{seguire il processo a basso livello}, sporcandosi
le mani.

\paragraph{I feature team.} Nei team medio-piccoli o medio-grandi possono
esistere sotto-team (\emph{feature team}) dedicati a un obiettivo specifico:
ad esempio un designer incaricato del character o della camera, affiancato da
uno o due programmatori e da un artista. In questi casi il designer
supervisiona il lavoro e gestisce la produzione del sotto-team.

\subsection{Gestire l'avanzamento del prodotto}

Esiste una specializzazione che si occupa solo di questo, il
\textbf{producer}, a un livello gerarchico completamente diverso (e che
potrebbe essere una strada percorribile al termine del percorso di studi).
All'interno di un feature team anche il designer può farlo, ma non è detto.
Può certamente \emph{dare la direzione}: nella sua documentazione può
indicare che per arrivare a uno stato giocabile del prototipo serve un certo
lavoro di arte e di programmazione, per lo stato successivo un altro, e così
via. In ambito produttivo esistono metodologie strutturate come
\textbf{Agile}, che verranno viste in seguito.

\subsection{Presentare le idee}

Il pitching (sezione~\ref{sec:gerarchia-corso}) è richiesto al designer sia
in pre-produzione sia in produzione, in misura che dipende dall'azienda. Ad
esempio Ubisoft lascia ai designer, e in realtà a qualsiasi membro, la
possibilità di \emph{pitchare} un'idea prima di una nuova produzione, perché
possa iniziare a prendere vita.

\subsection{Bozze e prototipi}

Le bozze vanno bene, ma devono essere \textbf{chiare}. C'è un'eterna lotta tra
documentazione e gioco giocabile: se si fa solo documentazione il gioco non
avanza; se si fa solo il gioco manca la documentazione e si perde la visione.
Sono due cose \textbf{complementari}.

Il modo per comunicare un'idea in forma grezza ma giocabile è il
\textbf{prototipo}. Un prototipo può essere qualsiasi cosa: di carta, in
Unreal con Blueprint ``brutti''. \emph{Non deve essere bello, deve essere
giocabile}.

\paragraph{Esempio: il prototipo di carta di Mario + Rabbids.} Durante lo
sviluppo di \emph{Mario + Rabbids}, i programmatori di Ubisoft erano impegnati
su altri progetti (come \emph{Ghost Recon}). I designer realizzarono allora
un \textbf{gioco da tavolo} che riproduceva le meccaniche del gioco, quelle
di uno strategico a turni sul modello di \emph{XCOM}. E funzionava.
L'obiettivo era \textbf{validare un'idea}, validare un gameplay.

Altri modi per rendere chiara un'idea sono i \emph{mock-up} nella
documentazione e le immagini modificate o i \emph{mash-up} di qualsiasi
natura. Lato documentazione l'importante è che sia chiaro; lato prototipo,
che sia giocabile.

\subsection{Fare ricerca per l'idea, e poi smontarla}

Il designer ha una componente creativa importante e fa ricerca per arrivare a
un'idea. Ma ciò che deve poi fare necessariamente è \textbf{smontarla}. Il
suo lavoro è scientifico, esattamente come il \textbf{metodo scientifico}, il
cui obiettivo è \emph{confutare un'ipotesi}: si arriva a un'ipotesi tramite
sperimentazione e poi si cerca $n$ volte di confutarla; se resiste, la si
assume come valida \emph{in un certo contesto} (le regole della fisica
classica, ad esempio, non valgono nella meccanica quantistica).

Allo stesso modo il designer arriva a un'idea (tramite ricerca o chiedendo
agli altri: le idee le hanno tutti, non sono un privilegio del designer),
cerca di smontarla e poi cerca di \textbf{metterla a terra}. Se l'idea non è
applicabile, non è implementabile, finché non si trova un modo per
realizzarla va scartata. Può essere l'idea più bella dell'universo, ma
\textbf{se non può essere messa a terra, non esiste}.

\paragraph{L'originalità.} Serve un'idea originale? Dipende dall'obiettivo e
dalla visione del progetto:
\begin{itemize}
  \item se l'obiettivo è vendere molto, un'idea originale non serve
        necessariamente;
  \item se si vuole creare un nuovo filone, serve un'idea, se non originale,
        almeno interessante: può essere un discorso di marketing o una
        meccanica innovativa, con il rischio di non vendere.
\end{itemize}
Il ruolo del \emph{director} e il concetto di \emph{visione} verranno
approfonditi in seguito.

\subsection{Analizzare i competitor}

Arrivati a un'idea (``vogliamo un certo tipo di combattimento''), si va a
vedere se qualcuno l'ha già fatta e \emph{come}. Le meccaniche non sono
privilegio di nessuno, salvo quando vengono brevettate, come il
\textbf{Nemesis System} di \emph{La Terra di Mezzo: L'Ombra di Mordor}.

\paragraph{Esempio: il combattimento da ridimensionare.} In un progetto
(presso Stormind Games) il team si trovò con un sistema di combattimento
troppo costoso: mancavano le risorse e gli animatori per realizzarlo. Poco
output lato gameplay significava non poter fare qualcosa di complesso e
profondo come in \emph{Elden Ring}; al contrario, il team aveva molti
artisti e poteva puntare su qualcosa di più artistico. Si iniziò così a
orbitare intorno al corpo a corpo di \emph{The Last of Us} (parte I e II).
Sul canale della GDC esistono diversi talk che raccontano cosa è stato fatto
e come, con successi, insuccessi e costi (ad esempio del \emph{motion
capture}): da quei talk si ricavarono moltissime informazioni utili.

I competitor possono essere:
\begin{itemize}
  \item \textbf{diretti}: un gioco molto simile al proprio;
  \item \textbf{di mercato}: un gioco che domina il mercato in cui ci si vuole
        affacciare;
  \item \textbf{indiretti}: un gioco con un sistema simile al proprio, ma
        tutto il resto diverso.
\end{itemize}
Regola d'oro: \textbf{là fuori, nascosto nell'ombra, c'è sicuramente un gioco
molto simile al vostro}. Va solo trovato. E se per qualche ragione non ha
venduto, bisogna capire perché.

\subsection{Lavorare in engine e fare tuning}

Il level designer lavora in engine, così come il technical designer; ma
qualsiasi game designer, a seconda di ciò che deve fare, può finire in
engine, soprattutto per il \textbf{bilanciamento} e il \textbf{tuning}.

Il \emph{tuning} è l'azione di bilanciare con precisione i dati. Esempio: in
\emph{Super Mario} si modifica quanto in alto salta il personaggio, con che
velocità, in quanto tempo raggiunge l'apice, e così via. Questo lavoro rende
il gioco:
\begin{itemize}
  \item ad alto livello, dotato della \textbf{progressione} giusta;
  \item a basso livello, \textbf{piacevole da giocare}.
\end{itemize}
Il \emph{fine tuning} è un lavoro di dettaglio, paziente e minuzioso.

\subsection{Bilanciare}

Il bilanciamento, emerso tardi nel brainstorming, è una delle attività
centrali del designer e ha anche una specializzazione dedicata
(l'\emph{economy designer}, sezione~\ref{sec:tipi-designer}).

\subsection{Analizzare il mercato}

In pre-produzione, di solito, sì: capire le tendenze, i generi che cambiano
nel tempo, la compatibilità del gioco con il momento storico, in vista ad
esempio di un possibile publisher. Va capito quanto, come e con quanto tempo
farlo, perché si è sempre limitati da molte cose.

\subsection{Sintesi}

\begin{quote}
Il designer parte dall'idea ma non si ferma lì: la confuta, lavora sui
dettagli, bilancia, fa \emph{fine tuning}; lavora con il team e fa da
\emph{manager} per diversi dipartimenti; e naviga costantemente tra i limiti.
\end{quote}

% ---------------------------------------------------------------------
\section{I limiti del design}
\label{sec:limiti}

I limiti ci sono sempre, e nel mondo del design sono spesso numerosi. I più
importanti entro cui navigare sono:
\begin{description}
  \item[Tempo] ad alto livello, l'idea deve essere realizzabile nel tempo di
        produzione disponibile; a basso livello, il proprio lavoro
        settimanale deve essere fattibile nel tempo che si ha. Non ci si può
        dividere tra mille feature team se realisticamente se ne può seguire
        uno solo.
  \item[Tecnologia] si può o non si può fare qualcosa? Ci sono i server per un
        multiplayer con un certo numero di giocatori? Ci si può permettere,
        ad esempio, una fauna che reagisce alla fisica come quella di
        \emph{GTA VI}, che richiede il meglio dei \emph{technical artist}?
  \item[Background e competenze] che cosa sa fare il team? Realisticamente,
        chi ha lavorato dieci anni su giochi horror viene assunto per un
        horror, perché ha le competenze per risolvere problemi già visti.
  \item[Persone e motivazione] il team ha la motivazione e le competenze per
        fare un certo tipo di gioco, e per lavorare insieme?
  \item[Denaro] tutto ruota intorno ai soldi. Che si lavori da soli o in mille,
        le persone vanno pagate e il gioco deve avere un ritorno economico.
        Bisogna decidere dove tracciare la linea: raggiungere il
        \textbf{break-even} (guadagnare esattamente quanto speso) o puntare a
        grandi guadagni. Con un \textbf{publisher}, inoltre, una parte
        consistente dei ricavi va a lui (ripartizioni dell'ordine di 80/20 o
        70/30, che dipendono dal contratto firmato).
\end{description}

\subsection{Non siete ``gli uomini delle idee''}

Per qualcuno può essere un problema scoprirlo, ma il designer non è
``l'uomo (o la donna) delle idee'': accanto alla parte creativa c'è una parte
tecnica, e si è fortemente limitati da molti fattori. Questi limiti, però,
sono un aiuto:
\begin{itemize}
  \item riducono il ventaglio di possibilità potenzialmente distruttive per il
        team (cose che non si riescono a fare, che non piacciono o per cui
        mancano le risorse);
  \item aiutano a mirare bene, trovando soluzioni creative in uno spazio più
        piccolo e più attuabile.
\end{itemize}

\begin{quote}
\centering
\textbf{Le idee sono nulla, l'implementazione è tutto.}
\end{quote}

% ---------------------------------------------------------------------
\section{Ruoli e specializzazioni}
\label{sec:tipi-designer}

\subsection{La scala gerarchica}

\begin{description}
  \item[Creative Director (CD)] mantiene la \textbf{visione} e dà l'idea, in
        modo che tutto segua una certa direzione. In teoria può diventarlo
        chiunque, ma con un'esperienza notevole, nell'ordine dei decenni: si
        assume che in quel tempo abbia visto come lavorano molti dipartimenti.
        Naturalmente un CD con un background artistico ha meno competenze di
        design, e viceversa.
  \item[Director di disciplina] uno per dipartimento: \emph{game director},
        \emph{level director}, \emph{art director}, ecc.
  \item[Lead] il capo del dipartimento, che gestisce i gruppi di persone. Chi
        lavora in un dipartimento risponde al proprio lead (ad esempio un game
        designer risponde al \emph{lead designer}, non al game director).
  \item[Senior / Principal] i profili più esperti (a seconda dell'azienda).
  \item[Junior] il profilo di inizio carriera.
  \item[Intern] lo stagista.
\end{description}
Lo schema si replica uno a uno per ogni dipartimento: arte, suono, animazione
(a volte arte e animazione insieme, o \emph{rigging} e animazione insieme),
programmazione, UX/UI e così via. A fare la differenza, a parità di ruolo, è
il \textbf{livello di seniority}.

\subsection{UX e UI}

UX e UI stanno a metà tra game design e arte: a seconda dell'azienda il ruolo
si chiama \emph{UX/UI Artist} o \emph{UX/UI Designer}.
\begin{description}
  \item[UX (User Experience)] come una persona interagisce con le cose. Esempio:
        su uno smartphone, un pulsante accende e spegne lo schermo; per
        digitare la password si scorre verso l'alto e compaiono i punti da
        toccare.
  \item[UI (User Interface)] come sono ``vestite'' queste cose: ad esempio
        l'aspetto della barra dei punti vita (HP), realizzata con asset 2D o
        3D.
\end{description}

\subsection{Le specializzazioni del game design}

Più l'azienda è grande, più la specializzazione è puntuale.
\begin{description}
  \item[Gameplay / Core mechanics designer] (sono la stessa cosa) si occupa
        delle meccaniche e delle 3C.
  \item[System designer] si occupa dei sistemi e delle loro interazioni.
  \item[Economy designer] si occupa del bilanciamento; lavora soprattutto con
        fogli di calcolo (Excel, Google Sheets).
  \item[Combat designer] si occupa del combattimento; negli annunci di lavoro
        indica quasi sempre un combattimento \emph{action}, con gameplay
        reattivo (non strategico), anche se non è una regola.
  \item[Multiplayer designer] si occupa delle interazioni multigiocatore.
  \item[Technical designer] fa da ponte tra game design e programmazione.
\end{description}

\subsection{Le specializzazioni del level design}

\begin{description}
  \item[Encounter designer] posiziona i nemici nella mappa in modo che gli
        scontri siano bilanciati e la progressione equilibrata come
        desiderato.
  \item[World / map designer] popola il mondo di gioco.
  \item[Quest designer] popola il mondo con le missioni e pensa alla
        progressione al loro interno, cioè a dove il giocatore deve andare (si
        vedano ad esempio gli annunci di lavoro di CD Projekt Red o di
        Larian).
  \item[Technical level designer] fa da ponte tra level design e
        programmazione.
\end{description}

\subsection{Discipline vicine: suono, UI, produzione}

Il monte ore del corso non permette di uscire come \emph{narrative designer}
(o \emph{sound designer}) completi. Ciò non toglie che si possano coltivare
queste discipline:
\begin{itemize}
  \item nei progetti di gruppo si può occuparsi anche del suono, della UI/UX,
        o provare ad applicare metodologie come Agile per fare
        \emph{producing} e gestire il team;
  \item il tempo dedicato a queste attività va però sottratto al game design,
        e dipende da quante persone ci sono nel team, da quanto lavoro c'è e
        da quanto è grande il progetto;
  \item avere sviluppato e rilasciato un progetto, anche piccolo (anche una
        \emph{game jam}), documentando i propri processi, è già un ottimo
        argomento da presentare in futuro;
  \item la via migliore resta lavorare nell'industria e, da lì, chiedere di
        imparare anche altre discipline.
\end{itemize}
Un gioco si può fare su qualsiasi cosa: sugli scarafaggi, sui piedini delle
sedie, sui pennarelli. Si parte avvantaggiati da ciò che è vicino alla propria
sensibilità: chi ama la musica può pensare a un gioco con meccaniche legate
alla musica, con un gameplay innovativo ma contenuto e realizzabile. Il design
è un cappello che comprende molte discipline.

% ---------------------------------------------------------------------
\section{Approfondimenti emersi in lezione}
\label{sec:approfondimenti}

\subsection{Roguelike e roguelite}

Un esempio di come il marketing conia generi: Supergiant Games, con
\emph{Hades}, rese celebre il termine \textbf{roguelite} al posto di
\emph{roguelike} (pur non essendo il primo roguelite in assoluto). La
differenza:
\begin{itemize}
  \item il \textbf{roguelike} si fonda su una serie di punti fissati in una
        sorta di \emph{manifesto} che definisce che cosa rende un gioco
        roguelike; uno su tutti: a ogni morte si ricomincia la \emph{run} da
        capo;
  \item il \textbf{roguelite} rompe alcuni di questi punti, in particolare
        introducendo una \textbf{progressione esterna alla run}: il giocatore
        può potenziarsi fuori dalla singola partita per vincere più
        facilmente in quelle successive. \emph{Hades} ne è il caso principe.
\end{itemize}
Analogamente, un mobile game sviluppato dal direttore del corso fu
promosso come ``il primo \emph{donation game} al mondo'': un gioco che
convertiva il tempo di gioco in donazioni reali a enti benefici per animali.
Anche qui un neologismo creato lato marketing per definire un genere.

\subsection{Applicazioni fuori dall'industria videoludica}

Le competenze sviluppate non servono solo ai videogiochi:
\begin{itemize}
  \item i \textbf{render architettonici} oggi possono essere realizzati con
        Unreal, e diversi professionisti formati nel game design lavorano in
        questo campo;
  \item la \textbf{gamification}, cioè l'applicazione di regole videoludiche
        a contesti non ludici (mostre interattive, progetti culturali,
        consulenze come quella per applicare questi concetti alle macchine da
        allenamento di un'azienda di fitness), è in Italia uno dei modi più
        redditizi di spendere queste competenze, spesso con finanziamenti
        pubblici.
\end{itemize}

% ---------------------------------------------------------------------
\section{Il primo task}
\label{sec:primo-task}

\begin{enumerate}
  \item Creare un indirizzo e-mail \textbf{professionale} (consigliato: un indirizzo Gmail nella forma
        \texttt{nome.cognome@gmail.com}; l'anno di nascita è accettabile).
        Va bene anche un indirizzo professionale già esistente usato per
        lavoro. Da evitare indirizzi poco seri (``non \texttt{blackfire67}'').
  \item Una persona della classe raccoglie tutti gli indirizzi in un unico
        documento e lo invia via e-mail al docente, che aggiungerà tutti al
        Drive condiviso del corso (dove verranno caricate lezioni, slide e
        materiali, per avere i documenti tutti nello stesso posto).
  \item Installare Discord e iscriversi al server del corso, dove passano le
        comunicazioni ufficiali. È consuetudine, e utile, creare anche un
        gruppo tra studenti.
\end{enumerate}
\textbf{Scadenza}: prima della lezione del venerdì.

% ---------------------------------------------------------------------
\section*{Punti chiave del capitolo}
\addcontentsline{toc}{section}{Punti chiave del capitolo}

\begin{itemize}
  \item Il corso è una \textbf{simulazione di lavoro}: ci si tratta da
        colleghi, con una gerarchia da studio di sviluppo.
  \item Si impara con l'\textbf{inquiry-based learning}: lezione, task,
        dubbi, domande, condivisione, progetti.
  \item Il fallimento affrontato con la mentalità giusta è la principale fonte
        di crescita; i feedback riguardano il lavoro, mai la persona.
  \item Il lavoro del designer è per metà creativo e per metà
        \textbf{ricerca e metodo scientifico}; il game design non è una
        scienza esatta.
  \item \textbf{Documentazione} e \textbf{pitching} sono il cuore del
        mestiere.
  \item Il designer progetta meccaniche e sistemi (partendo dalle
        \textbf{3C}), documenta, media, prototipa, confuta le idee, studia i
        competitor, bilancia e fa tuning.
  \item La leadership del designer è \textbf{al servizio} del team, non
        verticale.
  \item I limiti (tempo, tecnologia, competenze, persone, denaro) restringono
        e indirizzano la creatività.
  \item \textbf{Le idee sono nulla, l'implementazione è tutto.}
\end{itemize}

%    % ...
%
%  (Con hyperref aggiungere anche \renewcommand{\theHchapter}{z.\arabic{chapter}}
%  nella parte zero, per evitare ancore duplicate.)
%  Tutte le etichette di questo capitolo hanno il prefisso "z1:" per non
%  collidere con quelle dei capitoli del corso.
% =====================================================================

\chapter[La cultura del feedback]{Lavorare in squadra: la cultura del feedback}
\label{cap:z1-feedback}

Questo capitolo apre le \emph{lezioni zero}: incontri che precedono l'inizio
vero e proprio del corso di Game Design e che non riguardano una tecnica di
design in senso stretto, ma \textbf{la figura professionale} di chi lavora
nell'industria dei videogiochi. Prima di imparare a progettare un gioco
conviene infatti capire in che tipo di ambiente lo si progetterà: un ambiente
fatto di team, di gerarchie, di scadenze e, soprattutto, di persone che devono
riuscire a lavorare insieme.

La lezione è la prima giornata di un \textbf{workshop di due giorni} (circa
dieci ore in totale, pause pranzo comprese) tenuto da Stefano
Cozzi\footnote{Il cognome è riportato come compare nella trascrizione
automatica, che in più punti storpia i nomi propri; andrebbe verificato.
Nella stessa trascrizione lo studio per cui lavora appare come
``Stormine'', ``Stormline'' e ``Stormind'': si tratta di Stormind Games,
studio con sede a Catania, coerentemente con quanto detto in lezione.}, un
professionista con circa quindici anni di esperienza nel settore, tra
Blizzard Entertainment, Riot Games e, oggi, Stormind Games. Il workshop è
nato da un confronto con i coordinatori dell'accademia\footnote{Nella
trascrizione l'accademia è resa in modi diversi (``Hive'', ``a dive'',
``D.I.V.''); i coordinatori citati sono Enrico e Andrea, insieme al
direttore Diego Ricchiuti. Nel resto del capitolo si usa semplicemente
``l'accademia'', come negli altri capitoli.}: serviva un momento dedicato a
\emph{come ci si comporta in un ambiente professionale}, cosa che riguarda sia
il lavoro futuro sia i progetti di gruppo svolti durante il percorso di
studi, dove ci si deve già abituare a comportarsi ``come se si fosse in
un'azienda''.

La prima giornata ha un unico grande tema, sezionato in più parti:
\textbf{la cultura del feedback}. La seconda giornata (che sarà oggetto della
prossima lezione zero) sarà più pratica: \emph{time management},
\emph{brainstorming} e il ``guidare da dentro'' un team anche senza esserne
il leader (sezione~\ref{sec:z1-domani}).

Il percorso del capitolo segue quello della lezione:
\begin{enumerate}
  \item il formato del workshop e il percorso professionale del relatore
        (sezione~\ref{sec:z1-relatore});
  \item perché le aziende tengono tanto al lavoro di squadra e al
        \emph{fit} culturale (sezione~\ref{sec:z1-squadra});
  \item che cos'è il feedback e da che cosa dipende il modo di darlo
        (sezione~\ref{sec:z1-cose-feedback});
  \item le cinque fasi di vita di un team e il valore del
        \emph{post-mortem} (sezioni~\ref{sec:z1-fasi-team}
        e~\ref{sec:z1-qa-fasi});
  \item il confronto con i colleghi e alcune domande sul lavoro nei grandi
        studi (sezioni~\ref{sec:z1-rivalita} e~\ref{sec:z1-industria});
  \item il circolo virtuoso tra feedback, miglioramento e fiducia, e i suoi
        effetti sulla \emph{retention} e sulla cultura aziendale
        (sezioni~\ref{sec:z1-triangolo} e~\ref{sec:z1-cultura});
  \item il ``gioco del traditore'' e il principio del
        \textbf{default to trust} (sezioni~\ref{sec:z1-among-us}
        e~\ref{sec:z1-default-trust});
  \item le traiettorie di carriera e la valutazione delle performance
        (sezione~\ref{sec:z1-traiettorie});
  \item come si dà, come si riceve e quando si dà un feedback
        (sezioni~\ref{sec:z1-dare}, \ref{sec:z1-ricevere}
        e~\ref{sec:z1-quando});
  \item il feedback come regalo, le domande finali e le anticipazioni
        della seconda giornata (sezioni~\ref{sec:z1-regalo},
        \ref{sec:z1-qa-finale} e~\ref{sec:z1-domani}).
\end{enumerate}

Una premessa di metodo: il relatore ha alternato continuamente teoria,
domande alla classe e racconti della propria esperienza. I racconti non sono
digressioni: sono gli esempi concreti con cui i principi vengono resi
tangibili, e per questo sono riportati con cura.

% ---------------------------------------------------------------------
\section{Il workshop e il relatore}
\label{sec:z1-relatore}

\subsection{Formato e regole di partecipazione}

Il relatore ha cercato di ``riassumere quindici anni di esperienza in dieci
ore''. Ha materiale per circa il doppio delle ore disponibili: le cose
strettamente necessarie vengono trattate per prime, il resto è tenuto ``in
riserva'' come \emph{nice to have}. Le regole del workshop:
\begin{itemize}
  \item le domande si fanno \textbf{alla fine di ogni slide}, così da chiudere
        prima il discorso; non c'è un'unica sessione finale di domande e
        risposte (anche se ci sarà), perché il relatore preferisce fermarsi
        spesso e non fare ``uno spiegone di dieci ore'';
  \item le domande non devono per forza riguardare la slide appena vista: se
        una slide dà lo spunto per altro, se ne può parlare; al massimo sarà
        il relatore a rispondere in breve o a rimandare a una slide
        successiva;
  \item si può chiedere anche di esperienze passate, opinioni e sensazioni;
        se per NDA o per serietà professionale non può rispondere, lo dirà
        esplicitamente (``non lasciate che sia un \emph{blocker}'');
  \item su Zoom si usa l'alzata di mano o la chat, che il relatore tiene
        aperta sul secondo monitor;
  \item non ci sono pause ufficiali: la lezione è registrata e ognuno è
        libero di assentarsi quando serve, cercando però di non disturbare il
        flusso della lezione;
  \item le slide verranno fornite integralmente tramite la segreteria
        dell'accademia (Vanessa), insieme alle registrazioni: le slide senza
        il parlato servono poco, ma sono un buon promemoria.
\end{itemize}

\subsection{Il percorso professionale}

Il relatore si presenta (scherzosamente) come la ``guida escursionistica''
che accompagna la classe nella ``foresta incantata'' della cultura del
feedback, perché per molti è un territorio sconosciuto. Il suo percorso:
\begin{itemize}
  \item in Italia ha lavorato per aziende che fornivano servizi a terzi, tra
        cui il lavoro su \emph{Warhammer Online}, e per un canale televisivo
        sul digitale terrestre dedicato ai videogiochi\footnote{Nella
        trascrizione: ``Digital Game Network Channel''. Il nome esatto del
        canale andrebbe verificato.};
  \item a \textbf{Blizzard Entertainment} (circa sei anni, a Parigi) ha fatto
        il \emph{community manager} per l'Italia (\emph{StarCraft~II},
        \emph{World of Warcraft} da \emph{Mists of Pandaria} in poi,
        \emph{Hearthstone}, \emph{Diablo~III}, le espansioni di
        \emph{StarCraft}, il lancio di Battle.net), per poi diventare il
        manager dei community manager europei;
  \item a \textbf{Riot Games} (circa cinque anni, a Dublino) è stato
        brevemente community manager italiano e poi \emph{publishing
        manager} di un team di circa trenta persone per l'Europa. Il
        \emph{publishing} non è solo community: è l'unione di
        \textbf{community, marketing e PR};
  \item a \textbf{PG Esports} si è occupato di eventi dal vivo e online
        (ad esempio la Milano Games Week e tornei di \emph{Valorant} e
        \emph{League of Legends}), restando in contatto con Riot e Blizzard;
  \item da circa un anno è tornato nello sviluppo, a \textbf{Stormind
        Games}, dove oggi si occupa di risorse umane e organizzazione del
        lavoro; lo studio verrà presentato più ampiamente al \emph{Career
        Day} dell'accademia insieme alla sua HR manager;
  \item fa inoltre consulenze aziendali anche al di fuori dei videogiochi e
        insegna in master universitari (IULM di Milano).
\end{itemize}

Il suo ruolo ricorrente, sia a Blizzard sia a Riot, è stato quello di
\emph{Mr.~Wolf} (``per chi ha visto \emph{Pulp Fiction}''): quello che
risolve i problemi. Veniva inserito nei team più in difficoltà, ad esempio
dal punto di vista delle performance, per sistemarli e passare al successivo.
Era presente in alcuni dei momenti più critici, come il lancio di
\emph{Diablo~III} o il lancio di espansioni di \emph{World of
Warcraft}\footnote{La trascrizione riporta ``il lancio con i Gary's'' e
``World of Trenor'': si tratta verosimilmente di \emph{Mists of Pandaria} e
di \emph{Warlords of Draenor}.}, a fare \emph{damage control}.

La scoperta fondamentale di questi anni è che i problemi da risolvere
\textbf{quasi mai riguardavano il prodotto} (non era certo lui a sistemare i
server di \emph{Diablo~III}), ma i team che c'erano dietro: persone che
litigavano, non collaboravano, non riuscivano a stare nella stessa stanza
senza urlarsi contro. E nel 90\% dei casi tutto questo si riduce a
\textbf{problemi di comunicazione}: è raro che perfino le incompatibilità
caratteriali non siano risolvibili parlando.

\subsection{Struttura non vuol dire burocrazia}

A Stormind il relatore sta accompagnando un team passato da 30 a 95 persone
in meno di un anno, attraverso tutti i cambiamenti di cui parla il workshop.
Darsi una struttura viene spesso percepito male, come ``più processi, più
procedure'', ma bisogna distinguere:
\begin{itemize}
  \item un processo come ``timbrare il cartellino'' può essere davvero
        invasivo e noioso;
  \item definire \emph{come funziona un documento di game design, chi lo
        approva, chi lo visiona e chi no} è invece una struttura necessaria.
\end{itemize}
Uno studio di 30 persone nella stessa sede può farne a meno: quando c'è un
dubbio ci si alza e si va alla scrivania del collega. Uno studio di 95
persone sparse per l'Italia (e in parte per l'Europa) non può più farlo:
servono una struttura diversa, un modo di lavorare diverso e una
pianificazione del futuro diversa.

% ---------------------------------------------------------------------
\section{Perché il lavoro è sempre lavoro di squadra}
\label{sec:z1-squadra}

\subsection{Non si va in ufficio per fare amicizia}

Il primo messaggio è netto: \textbf{non si va in ufficio per fare amici} o
per trovare compagni di vita. Si va in ufficio perché c'è un lavoro da fare, e
lo scopo è riuscire a farlo \emph{anche con chi sta antipatico}. Con le
persone simpatiche è facile: diventeranno amici, e non sarà mai un problema.
È con chi sta antipatico che le cose si complicano, ed è lì che emergono i
\textbf{veri professionisti}: quelli che riescono a lavorare con chiunque
senza che l'antipatia pesi sul modo in cui lavorano e collaborano, mettendo il
prodotto e il team davanti al proprio ego e alle proprie preferenze personali.

La frase che fa da filo conduttore a tutto il workshop (e che il relatore
promette di riproporre ``fino alla nausea'') è:

\begin{quote}
  \emph{La capacità di un gruppo di persone di fare cose straordinarie
  dipende dal modo in cui queste persone operano insieme.}
\end{quote}

Le aziende raramente pensano che l'innesto di una singola persona possa
cambiare il destino di un progetto. Questo è tanto più vero quanto più grande
è il team, ma vale anche per i piccoli studi, anche per un gruppo di cinque
persone. Come negli sport di squadra: l'individuo conta fino a un certo punto,
e anche se si ha Dončić in squadra, se si gioca uno contro cinque si perde
comunque\footnote{Il riferimento è alle finali NBA in corso nei giorni del
workshop tra Dallas Mavericks e Boston Celtics (giugno 2024), il che
permette di datare approssimativamente la lezione.}.

\subsection{Lo \emph{useless genius} e il costo delle persone tossiche}

Molti dipartimenti HR usano la definizione di \textbf{useless
genius}\footnote{Il concetto è più noto in letteratura manageriale come
\emph{brilliant jerk}; in lezione è usata l'espressione \emph{useless
genius}.}: si può essere bravissimi, a volte i migliori nel proprio lavoro,
ma se non ci si sa comportare, relazionare e operare in un gruppo si è
fondamentalmente inutili ai fini del gruppo. Nel migliore dei casi inutili,
nel peggiore \emph{dannosi}.

In passato le aziende tendevano a ``chiudere un occhio'': la persona è
fondamentale, la teniamo. Oggi la domanda è diventata: \emph{quanto costa
all'azienda, oltre al suo stipendio, tenere una persona che crea un ambiente
tossico e peggiora chi le sta intorno?} Ci sono costi sommersi che le aziende
cominciano a vedere, e la tendenza è liberarsi di queste persone il prima
possibile. Il ragionamento è semplice:
\begin{itemize}
  \item le \textbf{competenze si possono trasmettere} o sostituire: si trova
        qualcuno che fa lo stesso lavoro magari nel doppio del tempo, oppure
        più persone che se lo dividono;
  \item il \textbf{carattere delle persone} invece cambia molto raramente.
\end{itemize}

\subsection{Il fit culturale}

Quando qualcuno \textbf{non è un \emph{fit} culturale} per un'azienda è
meglio separarsi presto. Ma non essere adatti a \emph{quell'}azienda non
significa non essere adatti a lavorare: è perfettamente legittimo che una
persona inadatta a un'azienda sia adattissima a un'altra.

Il relatore l'ha vissuto in prima persona. Dopo sei anni a Blizzard viene
contattato da Riot e affronta un processo di selezione lunghissimo (``quasi
undici colloqui''), al termine del quale viene scartato. Scrive allora una
mail che gli ha ``cambiato la traiettoria di carriera'': chiede
semplicemente perché non è stato ritenuto adatto. La risposta: in nessun
colloquio si era esposto abbastanza da far capire \emph{perché} volesse
lasciare Blizzard, e quindi pensavano che in realtà stesse bene dov'era. Il
motivo era che, per educazione professionale, non gli piace parlare male del
proprio posto di lavoro; portata all'eccesso, però, questa riservatezza era
diventata un ``è tutto bellissimo, ma vorrei provare Riot'', senza un perché.
Lo invitano comunque a un incontro dal vivo (``dal vivo è diverso''), e alla
fine viene assunto.

Una volta arrivato, i sei mesi di prova (\emph{probation}) sono serviti tutti
per smettere di ragionare e lavorare ``come a Blizzard''. La differenza
culturale era enorme, tanto che in Riot esisteva un gruppo di supporto
interno tra ex dipendenti Blizzard: un collega gliel'ha descritto come
``uscire da una relazione tossica e doversi riabituare a una relazione
normale''\footnote{Il relatore precisa che parla della Blizzard di allora,
che non esiste più: oggi è parte di Microsoft e nessuno dei suoi ex colleghi
vi lavora ancora, anche a causa dei licenziamenti di massa.}.
\begin{description}
  \item[Blizzard] era \emph{developer first}: il focus era chiaramente sugli
        sviluppatori, prima di tutti gli altri. Dopo sei anni c'erano ancora
        meeting a cui non poteva partecipare perché non abbastanza
        \emph{senior}, informazioni che non poteva sapere, piani che non
        poteva vedere.
  \item[Riot] era \emph{player first}: prima i giocatori, e i dipendenti come
        mezzo per trasferire il feedback dei giocatori dentro il gioco. Il
        secondo giorno lo portano in una stanza e gli mostrano il piano dei
        cinque anni successivi. Lui, sorpreso, pensa a un errore (``sono qui
        da oggi, non vorrei mettervi nei guai''); la risposta è: ``sei un
        dipendente Riot o no? Allora non ci sono livelli di accesso: devi
        sapere queste cose per fare meglio il tuo lavoro''.
\end{description}

Da questa esperienza ha imparato molto sulla \textbf{trasparenza}: condividere
le informazioni aiuta il lavoro. Non sempre, però: a volte non si condivide
qualcosa non per segretezza, ma perché non è ancora pronta per essere
presentata, esattamente come non si mostra un gioco quando è ancora solo
teoria (``te lo faccio vedere quando ho qualcosa di tangibile'').

La lezione per gli studenti è doppia: le aziende valuteranno il \emph{fit}
culturale sempre di più, ma capire se si è un \emph{fit} richiede anche
\textbf{lavoro da parte del candidato}. Bisogna avere la presenza di spirito
di dire: ``conosco i valori di questa azienda, non mando nemmeno il CV,
perché so che durerei un mese e mezzo''.

% ---------------------------------------------------------------------
\section{Che cos'è il feedback}
\label{sec:z1-cose-feedback}

\subsection{Le prime risposte della classe}

La prima domanda alla classe è: \emph{quando sentite la parola
``feedback'', a cosa pensate?} Le risposte raccolte, a voce e in chat:
\begin{itemize}
  \item un \textbf{riscontro} su ciò che si fa e su come lo si fa, sul lavoro
        o sul modo in cui ci si interfaccia con il gruppo; può essere
        lavorativo oppure comunicativo e comportamentale. Chi ha risposto lo
        vede come qualcosa che, se l'azienda non dà, va attivamente
        richiesto, perché è importante sapere come sta andando per
        migliorarsi;
  \item né positivo né negativo, ma \textbf{costruttivo}: in entrambi i casi
        serve a migliorarsi;
  \item uno \emph{scambio di opinioni}, un \emph{riscontro produttivo};
  \item un feedback ``di reazione'', non verbale: si fa una cosa e si
        osserva come cambia l'ambiente o il comportamento delle persone
        intorno (anche perché a voce la gente spesso si tiene le cose
        negative);
  \item per un designer, anche il \textbf{linguaggio del corpo}: osservando
        un giocatore giocare si può dedurre l'emozione che prova (frustrazione,
        incomprensione).
\end{itemize}

Sulla forma (scritto o orale?) le opinioni divergono: c'è chi lo preferisce
\emph{mille volte} scritto; chi pensa che dipenda dalle dimensioni
dell'azienda (orale nelle piccole, scritto in quelle con team di centinaia di
persone); chi pensa che dipenda dall'importanza del feedback (scritto se
tecnico o pieno di dati, da tenere sotto gli occhi; orale se è un consiglio
veloce).

Il relatore usa due esempi per mostrare che \emph{dipende}:
\begin{itemize}
  \item un feedback sulla programmazione ha senso darlo scritto, magari
        commentando direttamente il codice, più che ``andando a prendere un
        caffè'';
  \item ma se il feedback è ``ieri, durante il meeting, ho percepito male il
        fatto che mi interrompessi di continuo'', ce lo si aspetta
        \textbf{orale}.
\end{itemize}

\subsection{Da che cosa dipende il modo di dare feedback}

Lo \emph{spoiler} della giornata è che \textbf{non esiste una risposta univoca}
su come dare feedback. Il workshop non insegna ``il'' modo, ma fornisce
\emph{tutte le opzioni}, perché il modo giusto cambia in base a molti fattori:
\begin{itemize}
  \item a chi lo si dà: un collega, un superiore, il proprio capo, un proprio
        \emph{report} (un sottoposto);
  \item se è personale o professionale;
  \item se è la prima volta o la terza volta che si dà lo stesso feedback;
  \item se riguarda un progetto importante o secondario.
\end{itemize}

Manca però l'elemento fondamentale, che la classe individua subito:
\textbf{l'altro}. A parità di tutti gli altri fattori, il \emph{carattere
della persona che riceve il feedback} dovrebbe far cambiare il modo in cui lo
si dà. Per questo, prima della lezione, agli studenti era stato chiesto di
fare un test di personalità: non un test di psicologia avanzata, ma uno
strumento per rendersi conto di quanta varietà ci sia anche solo in una
classe di 53 persone (il che non vuol dire trovare 52 modi diversi di dare
feedback, ma lavorarci un po' sì).

Il relatore porta se stesso come esempio. Preferisce:
\begin{itemize}
  \item i feedback \textbf{positivi scritti}, perché è a disagio a ricevere
        complimenti in pubblico;
  \item i feedback sul proprio lavoro \textbf{diretti}: ``vai subito al
        punto e non girarmi intorno''.
\end{itemize}
È molto orgoglioso del suo lavoro e, quando qualcuno lo critica, la sua
primissima reazione istintiva è ``come ti permetti''. Negli anni ha imparato
a farla durare un secondo, ma chi gli dà feedback di persona dice di
vederla ancora negli occhi. Per mitigarla, dice esplicitamente a chi gli dà
feedback: \emph{``ignora il modo in cui reagisco nei primi cinque secondi,
vai avanti, poi mi passa''}. È considerato da tutti molto aperto al feedback
e non ha problemi a dire ``hai ragione, facciamo diversamente'': ma ci sono
reazioni istintive e caratteriali che è difficile controllare al cento per
cento.

\subsection{Con chi si lavora meglio}

Guardando le persone con cui si lavora abitualmente (anche nei progetti
dell'accademia), si noterà che si lavora meglio con due profili:
\begin{enumerate}
  \item chi è \textbf{esattamente come noi}, ma qui è un 50 e 50: alcuni
        caratteri vanno d'accordo con i propri simili, altri non vogliono
        nessuno con lo stesso carattere nel gruppo (e si potrebbe rientrare
        in questa categoria);
  \item chi è \textbf{affine}: ha qualcosa in comune con noi, ma non tutto,
        e proprio per questo può dare punti di vista differenti.
\end{enumerate}
Il lavoro della giornata è capire come valorizzare il secondo caso: ottenere
feedback da persone con caratteri, visioni e background diversi, per
completare una visione del proprio lavoro che è necessariamente limitata. È
come esplorare una mappa con la nebbia di guerra: un'altra persona, al di là
del muro, vede cose che noi non vediamo.

\subsection{Il feedback come somma di piccoli miglioramenti}

Per il relatore il feedback è questo: \textbf{applicare tanti piccoli
miglioramenti periodici}, che si sommano nel tempo. Migliorare dello 0,1\%
ogni giorno porta a fine anno a un miglioramento enorme, e questo vale per i
prodotti, per i giochi, per le persone e per i team.

\paragraph{L'esempio del ciclismo britannico.} La squadra di ciclismo
britannica, fino a una quindicina d'anni fa, non aveva vinto quasi nulla.
Cambiato il team manager, tutti si aspettavano grandi cambiamenti sportivi
(convocazioni, strategie di gara). Arrivarono invece \emph{minuscoli
cambiamenti non sportivi}: un altro chef in trasferta, materassi migliori nei
camper per riposare meglio, un sellino più comodo a fine tappa. Cambiamenti
marginali che da soli non avrebbero cambiato niente; eppure, dopo due o tre
anni di questo lavoro quotidiano, la squadra ha cominciato a vincere a
ripetizione, Olimpiadi comprese, cambiando il modo in cui il ciclismo è
visto in Gran Bretagna\footnote{È la nota strategia delle \emph{marginal
gains} associata a Dave Brailsford (British Cycling e Team Sky); il nome non
è citato in lezione.}. Certo, si è investito anche sulle persone giuste, ma
tutto è partito dai piccoli miglioramenti periodici.

Per gli studenti vale lo stesso: chiarire un conflitto con una persona con
cui si ha ancora poco a che fare, o una situazione che ha lasciato
``l'amaro in bocca'', sembra inutile oggi, ma darà i suoi frutti tra tre, sei,
nove mesi. Sono \emph{investimenti sul futuro}.

\paragraph{Peggiorare ogni giorno.} Esiste anche il percorso inverso:
lasciar correre, ``buttare sotto il tappeto'' le cose perché ``tanto è
inutile parlarne''. Queste cose tornano sempre fuori dopo otto mesi, come il
\emph{libro dei rancori} dei nani di Warhammer: ``ti ricordi quando, un anno
e mezzo fa, all'inizio del progetto, mi hai detto quella cosa? Mi sono
offeso''. A quel punto è quasi impossibile risolvere, perché nessuno ricorda
più con precisione cosa sia successo (``ma quale cena?'', ``forse era in
pizzeria'', ``no, eravamo da un'altra parte''). \textbf{Tenersi le cose nel
cassetto per tirarle fuori quando servono è sempre sconsigliato}: non servono
ed è impossibile usarle in modo costruttivo.

% ---------------------------------------------------------------------
\section{Le cinque fasi di vita di un team}
\label{sec:z1-fasi-team}

\subsection{Perché conoscerle}

Quasi tutti gli studenti lavorano già in gruppo (al primo anno su un gioco da
tavolo, negli anni successivi su un videogioco) e si saranno accorti che un
team attraversa delle fasi: il modo di lavorare insieme cambia nel tempo.
Spesso ci si fa prendere di sorpresa: ``eravamo così affiatati, cosa è
successo?''.

Un esempio da una consulenza fuori dai videogiochi: un'azienda farmaceutica
gli affida un team commerciale (rappresentanti che presentano medicinali a
ospedali, medici e privati). Il problema: ``prima andava tutto bene; ora che
siamo in diciotto dovremmo lavorare meno e meglio, e invece va peggio''. Va
peggio perché le persone nuove sono appena arrivate, mentre il management si
aspettava già di vedere i risultati dell'investimento.

Fino a pochi anni fa le fasi di un team venivano disegnate come una
\textbf{linea dritta verso l'alto}: le fasi erano le stesse, ma viste come un
miglioramento costante. Non si vedevano le parti \emph{fisiologiche} del
lavoro in team, cioè anche i peggioramenti. Accettarle è il primo passo per
accettare i fallimenti: nessun percorso personale o di team è lineare.

Il relatore ammette che il suo modo di lavorare è nato per il 70\% dal
proprio impegno a imparare, fare bene e rispettare le persone, e per il 30\%
dall'aver lavorato con ``incapaci colossali'' che gli hanno insegnato come
\emph{non} fare le cose. A Blizzard un manager gli fece capire che, pur non
sapendo ancora che tipo di \emph{people manager} sarebbe diventato, sapeva
benissimo quale tipo di manager \emph{non} voleva essere.

\subsection{Forming, storming, norming, performing, adjourning}

Le fasi sono cinque, un po' come le fasi del lutto. Molte aziende ne
considerano solo quattro (``e sono quelle che solitamente chiudono''); le
aziende al passo con i tempi ne considerano cinque\footnote{Il modello è
quello di Bruce Tuckman (1965), a cui la quinta fase, \emph{adjourning}, fu
aggiunta nel 1977. Il nome di Tuckman non è citato in lezione; la
trascrizione riporta la quinta fase come ``journing''.}.

\begin{description}
  \item[Forming] Si viene messi insieme e si comincia a lavorare. Non si sa
        nulla degli altri; spesso non si conoscono nemmeno i nomi, ed è
        normale.
  \item[Storming] Si cominciano a ``dare le testate come le capre in
        montagna per determinare il territorio'': questa cosa la faccio io,
        questa la vorrei fare così, chi decide? Si stabilisce chi fa cosa e
        come si divide il lavoro. Nel grafico è una \textbf{fase di
        peggioramento}: che il team litighi o discuta civilmente non cambia
        nulla, il team deve investire parte di tempo ed energie per risolvere
        queste questioni. È normale e \textbf{non è evitabile}.
  \item[Norming] Si comincia a capire come lavorano gli altri, cosa sanno
        fare e come far funzionare il gruppo; iniziano i miglioramenti. È una
        forma di \emph{accettazione} degli altri, che deriva dall'aver capito
        cosa devo fare io e cosa devono fare loro.
  \item[Performing] C'è uno sforzo \textbf{proattivo}: ``questa cosa
        toccherebbe a me, ma so che Alessio la fa molto meglio, quindi gli
        chiedo aiuto''. Si sfruttano i punti di forza degli altri, perché li
        si conosce e perché si sa che passare il testimone non avrà effetti
        negativi (il capo non dirà ``perché non l'hai fatta tu?'').
  \item[Adjourning] Il progetto finisce, si chiude o viene riorganizzato, il
        team viene sciolto. Quello che andrebbe fatto è analizzare
        \textbf{cosa ha funzionato, cosa no e cosa migliorare} per il
        progetto successivo.
\end{description}

\subsection{La fase che tutti saltano}

``Spoiler incredibile'': l'\emph{adjourning} è la fase che molte aziende
saltano, \textbf{perché c'è sempre il progetto dopo}.

A PG Esports, per la Milano Games Week (come per qualsiasi evento dal vivo a
cadenza annuale: Lucca Comics, Etna Comics, Comicon\dots), finita l'edizione
di un anno si cominciava dal giorno dopo a lavorare a quella successiva,
perché un anno sembra tanto ma non basta mai. Nei videogiochi è lo stesso:
finito un gioco, a parte le patch di mantenimento e del \emph{day one}, si
passa subito al successivo.

Il cambiamento introdotto dal relatore (in parte già praticato, ma da lui
strutturato) è stato: finita la Games Week, \textbf{tutti in ferie per una
settimana}; poi ci si ritrova per definire cosa ha funzionato e cosa no, con
un \textbf{post-mortem} (o \emph{debrief}); solo dopo si inizia a lavorare
all'edizione successiva. Senza questa fase manca la parte di
\emph{learning}, cioè di ciò che si è imparato dal progetto.

Il consiglio vale per gli studenti: quando si finisce il gioco da tavolo, e
poi il gioco digitale, bisogna trovare il tempo per il post-mortem,
altrimenti si butta via \textbf{il 90\% dell'utilità del progetto}. Una
struttura semplice per farlo è descritta nella sezione~\ref{sec:z1-qa-finale}.

% ---------------------------------------------------------------------
\section{Domande sulle fasi del team}
\label{sec:z1-qa-fasi}

\subsection{Lo storming che ritorna}

\paragraph{Capita che lo storming si ritagli spazio anche nelle altre fasi?}
Lo \emph{storming} è una fase precisa, che però molti manager, coordinatori e
team cercano di \textbf{saltare}: ``se io sono bravo come manager e loro sono
bravi come team, passeremo da forming a norming senza problemi''. È un errore,
perché senza una fase in cui si discute davvero non si sviluppa una
\emph{skill} fondamentale: \textbf{saper discutere}. Tutti conoscono persone
non abituate a farlo, per cui ogni minimo disaccordo diventa un affronto
personale (``avevamo deciso la pizza, perché andiamo al sushi? Perché mi fate
questo affronto?''): sono persone che non sono abituate ad avere conflitti
costruttivi.

Se si pensa di essere in \emph{norming} e continuano a emergere conflitti, il
punto è che \textbf{non si era davvero in norming}: lo storming non ``si
trascina'' nelle fasi successive, semplicemente non se n'è mai usciti.
Allo stesso modo, a chi chiede se ogni mini-fase di storming durante la vita
del gruppo sia frutto di uno storming fatto male o solo di carattere, la
risposta è già contenuta qui.

Un partecipante osserva che spesso, ancor prima di finire un progetto, si è
già messi su quello successivo: è il modo migliore per essere in costante
\emph{damage control}. Era l'esperienza del relatore a Blizzard: veniva
spostato fisicamente da un team all'altro, e finito il lancio di
\emph{Diablo~III} era già su \emph{World of Warcraft}, senza vivere il
``dopo''. Si è dovuto ritagliare da solo lo spazio per un debrief personale,
ed è lì che ha deciso che tipo di manager non voleva essere. Il consiglio:
\textbf{se non lo fa il team, fatelo voi}.

\subsection{Dove si passa la maggior parte del tempo}

\emph{In quale fase si passa più tempo nella propria carriera?} Gli studenti
rispondono \emph{storming} e \emph{norming}. Il relatore, in quindici anni,
calcola di aver passato in \emph{performing} circa \textbf{dieci mesi}. Non è
pessimismo: arrivare a performing deve essere sempre l'obiettivo, sia da
manager sia da membro del team; ma è difficile arrivarci, ed è normale che non
sempre ci si arrivi.

Lo storming può durare a lungo, ma con team molto esperti dura pochissimo,
perché c'è chi lo sa gestire. In realtà \textbf{la maggior parte del tempo si
passa in forming}, perché ogni equilibrio raggiunto vale per \emph{quelle}
persone, \emph{quel} progetto, \emph{quelle} tempistiche. Basta che cambi un
qualsiasi fattore (una deadline spostata, una persona nuova, qualcuno che si
dimette, un progetto cancellato) e si ricomincia, perché i bilanciamenti
precedenti non valgono più.

Esempio: un team di cinque persone affiatatissimo, in pieno performing. La
dirigenza, entusiasta, ``raddoppia le risorse'' e lo porta a dieci. O si
lascia tutto com'è (e i cinque nuovi non fanno niente), o bisogna ridiscutere
chi fa cosa e come i nuovi possono aiutare: che piaccia o no, si ricomincia
da forming.

Quindi la bravura di un manager, ma anche di un membro del team,
\textbf{non sta nel passare tanto tempo in performing, ma nell'accelerare il
passaggio da una fase all'altra}, e nel fare bene la fase di
\emph{adjourning}.

\subsection{Analizzare anche i successi}

C'è la tendenza a fare post-mortem solo dei progetti andati male: ``questa
cosa non ha funzionato, vediamo perché''. Va benissimo. Ma quando un progetto
va bene si dà per scontato che sia andato bene ``perché era buono''. Bisogna
invece chiedersi \emph{perché} è andato bene: era il prodotto? La
tempistica? Il modo in cui abbiamo fatto una cosa piuttosto che un'altra?
\textbf{Se non si capisce davvero il motivo del proprio successo, si
tenderanno a replicare le cose sbagliate.}

\paragraph{Blockbuster.} Quando arrivò la notizia che Netflix avrebbe
lanciato un servizio di streaming, Blockbuster (la catena di noleggio di
videocassette e DVD) aveva tutto il tempo per preparare un'offerta
concorrente. In una riunione del management raccontata in un libro (di cui il
relatore non ricorda il titolo), la direzione sembrava orientata a investire
nello streaming, finché un dirigente fece notare che il 12\% del fatturato
derivava dalle \emph{multe per riconsegna tardiva}, impossibili con lo
streaming, mentre secondo le stime lo streaming avrebbe prodotto il 10\% del
fatturato. ``Non possiamo rinunciare a questo 2\%'': niente streaming. Oggi
non c'è più un Blockbuster ``manco a cercarlo col lanternino''. Non avevano
capito qual era il vero motivo del loro successo, cioè il loro modello di
business\footnote{Le percentuali sono quelle riferite in lezione e non sono
state verificate.}.

\paragraph{The Founder.} Nel film sulla nascita di McDonald's c'è una scena in
cui il futuro socio del protagonista gli dice: ``tu pensi di essere nella
\emph{food industry}, e invece sei nella \emph{real estate industry}''. Non
dovrebbe investire sugli hamburger, ma sui terreni su cui costruisce i
ristoranti. Anche per un team vale lo stesso: bisogna capire che cosa
l'ha reso davvero di successo, oppure no.

\subsection{L'autoanalisi personale}

\paragraph{E quando si cambia azienda?} Uno studente chiede se il debrief sia
possibile anche quando si lascia un'azienda. Non si può tornare in ufficio a
discutere della propria esperienza, ma lo si può fare \textbf{per se stessi}.
Il relatore l'ha fatto, anche se all'epoca in modo poco consapevole, ed è
così che è passato dal \emph{community management} al \emph{people
management}. Analizzando cosa aveva fatto bene a Blizzard si era reso conto
che:
\begin{itemize}
  \item aveva costruito un ottimo rapporto di fiducia con i giocatori (con
        alcuni proprietari di fansite parla ancora oggi);
  \item pur senza la seniority, era diventato un punto di riferimento per il
        team, visto come imparziale ed equilibrato anche su temi che non lo
        riguardavano direttamente;
  \item sapeva \textbf{risolvere problemi}: durante il lancio di
        \emph{Diablo~III} era il più tranquillo di tutti, non perché
        sminuisse il problema, ma perché sapeva che si sarebbe risolto con i
        suoi tempi; il lavoro non era spegnere l'incendio ma ``far sfollare
        la gente'' e aspettare che tutto tornasse in ordine.
\end{itemize}
E aveva capito cosa sapeva fare meno bene, ad esempio \textbf{gestire il
proprio manager}. Avere un manager non significa che il manager gestisca te:
il manager gestisce il tuo lavoro, ma tu devi saper gestire i tuoi
\emph{stakeholder}. Spesso lo stakeholder è il proprio manager; da manager,
gli stakeholder saranno CEO, CFO e così via. Gestirli significa conoscerli,
sapere come portare loro le informazioni, come far presente che qualcosa non
va. Fa parte del lavoro, anche se a lui non piace, e ci ha lavorato.

Nel mese e mezzo di pausa prima di entrare in Riot si è concentrato su
questo: ``cosa porto come mio punto di forza?''. L'autoanalisi è culminata
nel chiedere espressamente alla persona che l'aveva assunto: ``prima mi avete
detto di no, ora mi avete preso: perché?''.

\paragraph{Come si capisce se un progetto è andato bene?} Dipende
dall'\textbf{obiettivo}. Se l'obiettivo era vendere milioni di copie e il
gioco le ha vendute, bisogna capire perché. Se, più ragionevolmente,
l'obiettivo era un gioco da tavolo che permettesse di mettere in pratica le
conoscenze del corso, va valutato per quello: ha permesso di farlo? Se sì,
perché? Se no, perché no? L'obiettivo va fissato \textbf{all'inizio del
progetto, non alla fine}, proprio per poterlo valutare dopo.

\subsection{Quando il team cambia di continuo}

\paragraph{Tutti questi cambiamenti non sono dannosi per il progetto?}
Assolutamente sì, ma quasi mai sono volontari. I progetti andati male sono
spesso progetti grandi.

Il relatore racconta di \emph{Cyberpunk 2077}. Conosceva persone in CD
Projekt Red (``il mondo è piccolissimo, le cose si sanno''). Alla Gamescom va
con Diego Ricchiuti, direttore dell'accademia, a vedere la presentazione del
primo gameplay, quella ``in cui regalavano il giubbottino da samurai''.
Finito il video, tra gli applausi, si guardano e capiscono entrambi che
\textbf{il gioco non è pronto}, a soli quattro mesi dall'uscita prevista:
Diego aveva notato il combattimento e la guida. Gli amici di CD Projekt Red
(uno aveva lavorato con lui a Blizzard per sei anni) gli confermano tra le
righe che ``ce la faremo, ma ci sono ancora cose da sistemare''. Il gioco fu
poi rimandato più volte\footnote{In lezione il relatore dice ``due, forse
tre volte''; l'uscita di \emph{Cyberpunk 2077} è stata effettivamente
rinviata tre volte prima del dicembre 2020.}.

Per deformazione professionale, chi ha sviluppato l'occhio critico non
riesce più a godersi un trailer: in un trailer di \emph{The Day Before} Diego
gli aveva fatto notare che l'effetto della fiammata di un fucile era stato
chiaramente aggiunto in postproduzione e non faceva parte del gameplay. ``È
il triste epilogo di questa storia'': non riescono più a godersi un trailer.

Il punto è che più il team è grande, più queste variabili vengono imposte, e
più è facile che il fallimento \textbf{non dipenda dal fatto che il prodotto
sia brutto} (\emph{Cyberpunk} è un bellissimo gioco sotto molti punti di
vista, e il relatore lo sta rigiocando ora che è stato corretto ed espanso).
Nei grandi gruppi qualcuno arriva, qualcuno si dimette, cambia il budget,
arriva il rinvio; nel caso di \emph{Anthem}, il team ha fatto così tante volte
forming--storming--forming--storming che ne è uscito ``un prodotto che
dovevamo espellere: l'abbiamo espulso, buttato fuori dal nostro organismo pur
di liberarcene''.

Tra i cambiamenti c'è anche la crescita verticale: se il capo viene promosso
perché lavora benissimo, bisogna scegliere un nuovo capo, ed è sempre un
caos (chi ci resta male, chi pensava di meritarlo, chi voleva il suo amico,
chi non sopporta il prescelto). Per questo il relatore è \textbf{contrario
alle promozioni durante i progetti}: si rischia di rovinare tutto il lavoro di
storming.

\subsection{Obiettivi chiari nei piccoli gruppi}

Nei gruppi piccoli, come quelli degli studenti, il problema deriva più
spesso dal \textbf{non avere chiaro che cosa si vuole ottenere}: un prodotto
puramente educativo e formativo, o un prodotto richiesto dall'accademia che,
se viene bene, si potrebbe anche vendere? Un gioco commerciale o un gioco
artistico? Va chiarito subito, perché ci saranno sempre momenti in cui
scegliere tra A e B, e a guidare la scelta sarà la domanda: \emph{quale delle
due opzioni ci porta più vicino all'obiettivo?} Se l'obiettivo non è chiaro,
non si sa cosa scegliere.

Questo porta molte aziende a fare \textbf{scelte di design che in realtà sono
scelte commerciali}. Le microtransazioni sono scelte commerciali: nessuno ha
mai pensato ``sai cosa piacerebbe tanto ai giocatori? Pagare le bustine''.
Prima le bustine si guadagnavano giocando; ora la tendenza è ``potresti
giocare, ma se paghi è meglio''. Si pensi alle polemiche su Ubisoft per
l'opzione di ``saltare il gioco'' pagando, o ai famosi \emph{boost} di livello
di \emph{World of Warcraft}: ``devo pagarti per saltare il gioco?''.

Una studentessa (Silvia) cita \emph{Final Fantasy XIV}, che a un certo punto fu
``resettato'' perché non funzionava: un esempio positivo. Il relatore aggiunge
quello negativo di \emph{Age of Conan}, che non è mai riuscito a fare quel
reset iniziale (``è una schifezza, ve lo ripresentiamo quando è pronto''). Ci
vuole consapevolezza, ma anche i mezzi: un conto è Square Enix, che può
permetterselo, un altro è un'azienda che non può chiudere tutto per un anno
continuando a pagare gli stipendi.

Riassumendo: nei piccoli gruppi i problemi sono più di
\textbf{comunicazione e chiarezza} (obiettivi non chiari o che si spostano);
nelle grandi aziende di più di \textbf{persone che vanno e vengono, obiettivi
e tempistiche che cambiano}. In accademia le scadenze sono fisse (il gioco da
tavolo si consegna a fine anno scolastico, a giugno, e giugno resta); in
azienda può succedere che una feature prevista per la \emph{vertical slice}
di luglio venga anticipata a maggio per portarla alla Gamescom, e un sistema
perfetto per consegnare a giugno può non esserlo più per consegnare a maggio.

\subsection{L'ambiente determina la qualità del prodotto}

\paragraph{È vero che chi passa per Electronic Arts finisce per fare brutti
giochi?} Uno studente cita BioWare, che faceva ottimi giochi finché era
autonoma. Il problema non è dei developer, che non sono ``improvvisamente
diventati incapaci'': è evidente che il problema è altrove, e che
\textbf{l'ambiente in cui si fa un gioco determina molto spesso la qualità del
prodotto}.

Un esempio locale positivo: Tommaso Bonanni, conosciuto dal relatore insieme a
Diego e già docente dell'accademia, ha fondato a Roma Caracal Games Studio,
applicando fin dall'inizio questi concetti di cultura aziendale. L'azienda è
rimasta piccola, ma ha molto chiaro cosa vuole fare e come. Quando si
predispone fin dall'inizio l'azienda e il team a lavorare in un certo modo è
fattibilissimo; è molto più difficile quando si è EA, si acquisisce
un'azienda e le si dice ``da oggi si fa così''.

Il relatore ha vissuto la transizione da Blizzard ad Activision Blizzard (uno
dei motivi per cui se n'è andato, volontariamente). Ha visto la cultura
interna cambiare per gradi:
\begin{enumerate}
  \item ``facciamo giochi fighissimi, perché i giochi devono essere
        fighissimi'';
  \item ``facciamo giochi fighi così vendono di più'';
  \item ``vendiamo di più facendo giochi fighi'' (sembra uguale, ma è molto
        diverso);
  \item ``facciamo più soldi''.
\end{enumerate}
Lo si capiva dalle scelte di design, osservando quale fosse il peso più
grande sulla bilancia. Lo stesso varrà per gli studenti, nei loro giochi e nei
colloqui di lavoro: capire subito che cosa vuole l'altra parte.

\paragraph{Leader e visione.} Nei gruppi dell'accademia, dal secondo anno, c'è
una persona che assume il ruolo di \emph{lead} o \emph{scrum master}. È
diverso dall'essere un \emph{game director} (``il gioco si fa come dico io,
aiutatemi a farlo con la mia visione''), come avviene nei grandi studi: da lì
nascono i nomi nei titoli, come \emph{Sid Meier's Civilization}, che indicano
la persona che ha fatto quel gioco, oppure un puro branding, come \emph{Tom
Clancy's Rainbow Six}.

\paragraph{Inclusion come jolly.} A margine, uno studente osserva che in molte
aziende ``la risposta giusta è sempre \emph{inclusion}'', usata come un
jolly. È comunissimo, conferma il relatore: \emph{diversity and inclusion}
viene spesso usata per dire ``così includiamo più clienti'', ma non è quello
il suo punto (``altro discorso'').

% ---------------------------------------------------------------------
\section{Confrontarsi con chi è più bravo}
\label{sec:z1-rivalita}

Arrivato a Riot dopo sei anni a Blizzard, il relatore pensava di avere un buon
bagaglio di esperienza. Si è trovato invece ``come un fante buttato in mezzo
ai Navy SEAL'': persone che avevano superato 11, 12, 13 colloqui, con un
livello molto alto.

Individua subito un collega francese che, nella sua testa, diventa ``la
versione migliore di me'': stesso ruolo (community manager), esperienza
simile (arrivava da Ubisoft, come quasi tutti i francesi di Riot), ma molto
più spigliato, sempre con la battuta pronta, ``l'anima della festa''. Il
relatore invece è introverso: un workshop come questo per lui è
\emph{emotivamente drenante}, non ama parlare in pubblico pur amando
condividere le proprie esperienze, e non riesce a sostenere conversazioni
come ``cosa hai fatto questo weekend?''. Lavorando raramente insieme, aveva
sviluppato una rivalità ``da allenatore di Pokémon'' (lui era Ash, il
relatore l'antagonista).

Dopo sei mesi, a un evento di \emph{team building} aziendale (a Tenerife), il
suo capo, molto attento a queste dinamiche, li mette apposta nello stesso
gruppo. A fine serata, complice l'open bar, il relatore gli confessa di averlo
mal sopportato dal primo giorno per invidia della sua estroversione. Il
collega ride e risponde che pochi giorni prima aveva detto al capo la stessa
cosa al contrario: invidiava a Stefano il suo essere sempre organizzato, sul
pezzo, presentarsi nella sala riunioni giusta (mentre lui finiva in quella
sbagliata), arrivare con gli appunti, con i dati a supporto e le analisi di
mercato, mentre lui spesso non aveva nemmeno letto la mail del meeting.
Un'invidia in due direzioni, che li ha portati a lavorare molto insieme e a
migliorarsi a vicenda.

Le lezioni:
\begin{itemize}
  \item non bisogna demoralizzarsi vedendo qualcuno che fa lo stesso lavoro
        ed è più avanti sotto alcuni aspetti: spesso c'è un bilanciamento. I
        ``punti caratteristica'' sono limitati: se li ha messi in carisma,
        non li avrà messi in saggezza;
  \item conviene \textbf{chiedere a queste persone come mai} sono così brave
        in qualcosa, come si fa con un progetto andato bene: magari non lo si
        farà mai bene quanto loro, ma si comincerà ad avere il
        \emph{mindset} giusto; e viceversa, loro saranno spinti a chiedere a
        noi;
  \item spesso \textbf{le proprie doti le scopriamo dagli altri}. Il relatore
        non si era mai accorto di essere ``quello che resta calmo durante le
        emergenze'' (il lancio di \emph{Diablo~III}, quello di
        \emph{Draenor}, i problemi dal vivo alla Games Week) finché non
        gliel'hanno fatto notare tutti: per lui era la normalità. A chi pensa
        ``non sono capace, forse non è il mio percorso'', può darsi che
        semplicemente nessuno gli abbia ancora detto ``questa cosa la sai
        fare davvero bene, concentrati su questo''.
\end{itemize}

% ---------------------------------------------------------------------
\section{Domande sul lavoro nell'industria}
\label{sec:z1-industria}

Diverse domande degli studenti hanno riguardato le condizioni di lavoro nei
grandi studi internazionali. Le risposte sono raccolte qui per argomento.

\subsection{Presenza, remoto e trasferimenti}

\paragraph{Grandi studi: quasi solo in presenza.} Sia a Blizzard sia a Riot si
lavorava \textbf{solo in presenza}, senza smart working: sei anni a Parigi e
cinque a Dublino. Si parlava solo inglese: anche nei meeting tra soli
italiani si tendeva a comunicare e a scrivere in inglese, per correttezza
verso chi si fosse aggiunto e per lasciare una traccia comprensibile a tutti.
Il francese serviva per fare la spesa e sopravvivere, non in ufficio.
L'ambiente era estremamente eterogeneo (francesi, polacchi, spagnoli e
molti altri): durante il colloquio dal vivo a Dublino, mentre si preoccupava
di essersi vestito in modo adeguato, vide passare dietro la parete a vetro un
collega svedese con i dreadlock, a piedi scalzi, con una pianta sotto un
braccio e un aquilone sotto l'altro, e capì di essere ``a posto così''.
Un'esposizione a molte culture e modi di pensare che il relatore considera
bellissima.

Anche oggi, tranne alcuni studi particolari, la maggior parte del lavoro si
svolge di persona: amici a Paradox lavorano a Stoccolma, i pochi rimasti a
Blizzard si sono dovuti trasferire a Irvine. Per posizioni \emph{junior} o
\emph{entry level} i grandi studi raramente concedono il lavoro da casa; per
profili senior può esserci qualche eccezione (ad esempio nove mesi da remoto
e tre in sede).

\paragraph{Relocation.} Spesso tocca spostarsi, ma il relatore non ha mai
pagato un trasloco di tasca propria: queste aziende offrono la
\emph{relocation} (aiuto a trovare casa, ad aprire il conto in banca e
supporto pratico), perché è un investimento sulla persona.

\paragraph{Colloqui mentre si lavora altrove.} Per i colloqui in presenza
prendeva ferie, senza dire dove andava. Il consiglio generale: per fare
colloqui \textbf{prendete ferie o permessi}, non datevi malati, ``che non è un
bel modo di comportarsi''.

\paragraph{Remoto ``doppio''.} Un collega lavora per l'azienda coreana che
gestisce \emph{PUBG}, con ufficio ad Amsterdam, ma lavora da Milano: un
``doppio remoto''. Si trova, ma in genere un ambiente anglofono con sedi in
Europa chiederà di trasferirsi in Europa.

\paragraph{Famiglia e trasferimenti.} Il relatore ha avuto la fortuna di
sposare una persona che fa lo stesso lavoro (conosciuta lavorando su
\emph{Warhammer Online}, poi con lui a Blizzard e a Riot, nella
localizzazione e nel doppiaggio, ad esempio di \emph{World of Warcraft} in
italiano): avendo competenze rare, spesso venivano assunti entrambi, al
massimo a pochi mesi di distanza. Per i developer la soluzione frequente è
che la famiglia segua la persona assunta, perché lo stipendio basta per
tutti. Ma non sempre c'è equilibrio, e a volte uno dei due deve sacrificarsi:
sua moglie ha lasciato Blizzard per seguirlo a Dublino (per poi essere
assunta da Riot); lui ha lasciato Riot per tornare in Italia quando lei ha
trovato lavoro qui. Ha dovuto rinunciare più volte a lavorare in Islanda per
CCP (\emph{EVE Online}), perché non era una scelta solo sua. Si fanno
compromessi, ma si riesce a far quadrare.

\paragraph{Ci si sposta per ogni progetto?} No: si sono spostati
esclusivamente per cambiare azienda. Cambiando progetto nella stessa azienda
è raro, quasi mai (``non l'ho mai visto succedere''), che venga chiesto di
trasferirsi. In Francia stavano benissimo (vivevano fuori Parigi; ``i
problemi con i francesi sono in realtà problemi con i parigini''), parlavano
francese ed erano vicini alla doppia cittadinanza: se ne sono andati perché
non vedevano un futuro a Blizzard, per il modo in cui l'azienda trattava i
colleghi.

\paragraph{Il mercato oggi.} Non esiste un \emph{silver bullet} per trovare
lavoro domani dove si vuole. Con tutti i \emph{layoff} recenti ci sono
migliaia di esperti del settore che cercano lavoro: serve pazienza, ma è
fattibile.

\paragraph{Studi indie e remoto.} Negli studi più piccoli il remoto è più
diffuso: a Stormind circa 60 persone su 95 lavorano da remoto, ma lo studio
si è strutturato fin dall'inizio per farlo. Serve consapevolezza. Bending
Spoons (che sviluppa app per mobile, non giochi, ed è l'azienda dell'app
Immuni), pur essendo grandissima e con sede a Milano, lavora da remoto e
lascia l'ufficio a chi vuole: lo smart working è posizionato quasi come un
\emph{benefit}, ma permette anche di avere un ufficio da 300 persone invece
che da 700. Le sedi fisiche costano decine di milioni all'anno (Blizzard ha
chiuso le sedi di Amsterdam e di Parigi). Tuttavia, per lavori molto
verticali e specifici gli studi preferiscono avere le persone dal vivo; i
ruoli più spesso remoti sono quelli che non fanno sviluppo (customer service,
PR\dots).

Anche gli artisti lavorano spesso in presenza: l'accademia può dare solo una
\emph{simulazione} di quanto stretto sarà il contatto tra programmazione,
design e art. ``Sarà ancora di più.''

\paragraph{La comunicazione da remoto.} È più difficile. Pur essendo
estremamente introverso, il relatore \textbf{obbliga tutti a usare la webcam}
nei meeting in cui conta: non in un aggiornamento con novanta persone
collegate, ma certamente in un colloquio a due sulla performance dell'ultimo
anno. La componente fisica è una parte enorme della
comunicazione\footnote{In lezione si parla di ``quasi l'ottanta per cento'';
il dato va preso come indicativo, poiché le percentuali di questo tipo
circolano in forme diverse e sono spesso citate fuori contesto.}: anche se si
è timidi, accendere la webcam è l'approssimazione più vicina all'essere dal
vivo. Negli ambienti dove si lavorava senza webcam la situazione era
``esponenzialmente peggiore''.

\subsection{Differenze culturali nel team}

\paragraph{Gestire sensibilità diverse.} A Riot i team erano divisi per
\emph{lingua}, non per nazione: il team di lingua spagnola includeva anche un
messicano, quello di lingua russa includeva russi e ucraini, proprio nel
periodo della guerra in Crimea. Tra i colleghi non c'era alcun problema (erano
amici), ma ci si è posti la domanda: \emph{come comunichiamo con una community
che sarà divisa su questo tema?} Nei giochi online la divisione era evidente:
in \emph{Rocket League} c'era chi, con un compagno di squadra con la bandiera
avversaria, si faceva otto autogol pur di farlo perdere. Allo stesso modo,
con molti americani mandati a Dublino a coprire posizioni scoperte, c'erano le
elezioni (Clinton--Trump, Biden--Trump) e varie guerre: serviva sempre una
certa sensibilità.

Imparare a navigare queste situazioni lo ha reso un professionista migliore:
gli ha insegnato a fermarsi e chiedersi ``io la percepisco così; gli altri
dieci, con background, storie e lingue diverse, come la penseranno?'', a
ragionare più \emph{long term} e più \emph{big picture}, con più empatia.

\paragraph{Stili di comunicazione.} Un partecipante che ha lavorato nel Regno
Unito chiede dello stile comunicativo italiano, a volte frainteso. Uno dei più
grandi complimenti ricevuti dal relatore è stato ``non sembri italiano''
(``gesticoli poco per essere italiano'', ``parli bene inglese per essere
italiano''). Ha imparato ad \textbf{adattare la comunicazione alla persona}:
con spagnoli e francesi la battuta sulla nazione funziona, perché c'è un
``odio--amore'' condiviso; con russi e giapponesi no. Doveva adattare anche
l'umorismo. Un collega metà russo e metà messicano diceva di aver preso ``il
lato peggiore di entrambe le nazionalità'', ed è stato un ottimo maestro per
imparare a leggere culture diverse.

Non bisogna generalizzare, ma gli stereotipi possono essere un punto di
partenza ragionevole, \emph{sapendo che troverete sempre le eccezioni}. E
vanno letti correttamente: gli americani tendono a essere più estroversi e ad
avere sempre qualcosa da dire non perché ``sono fatti così'', ma perché
vengono educati così (all'università tra i corsi c'è il \emph{debate}, saper
argomentare in pubblico); un collega giapponese può essere più introverso per
ragioni culturali. Nel dubbio si parte da lì, poi si conosce la persona.

\subsection{Negoziare, discriminazione e stipendi}

\paragraph{La famiglia come leva.} Avere un partner riduce le opzioni, ma le
aziende lo sanno: uno spostamento è per due. Si può usare come leva per
negoziare (un compenso diverso, un supporto per il trasloco): se un'azienda ha
deciso di volere quella persona, 2, 3, 4 o 10.000 euro in più all'anno, per
un'azienda che fattura centinaia di milioni o miliardi, non cambiano nulla.

\paragraph{Essere sposati è un motivo per non essere assunti?} No. In alcuni
paesi è proprio illegale discriminare su questo: in Francia era illegale
chiedere CV con data di nascita, foto, sesso o religione, e chiedere se si
vogliono figli. Purtroppo, dai racconti dei suoi studenti di master (in
settori come \emph{beauty and wellness} e \emph{food and wine}), in piccole
realtà italiane queste domande (``hai 25 anni, hai intenzione di
sposarti?'') vengono ancora fatte, e servono a filtrare i candidati. Nelle
grandi aziende ci sono sistemi per evitarlo; anzi, molte danno \emph{più}
benefit a chi ha famiglia (``porta pure la famiglia, è tutto pagato''),
perché le persone competenti hanno spesso una certa età e quindi,
ragionevolmente, una famiglia.

\paragraph{Gli stipendi.} Dipendono molto da dove si vive. Lo stipendio a Riot
non era paragonabile a quello di Blizzard perché in Francia era tassato al
52\% e in Irlanda al 17\%\footnote{Aliquote come riportate in lezione,
riferite alla situazione personale del relatore.}: nessuno stipendio
competitivo può compensare il 35\% di differenza di tasse. Negli Stati Uniti,
con pochissima previdenza pubblica, gli stipendi sono molto più alti perché
poi ci si paga tutto (``se hai un raffreddore paghi tu''). Bisogna valutare il
potere d'acquisto di quei soldi nel paese in cui si andrà a vivere.

\subsection{Freelance e \emph{work for hire}}

All'inizio della sessione pomeridiana uno studente che lavora nell'audio
chiede quali siano le condizioni per chi lavora come \emph{freelance}.
\begin{itemize}
  \item Dipende molto dal progetto: spesso il contratto è a progetto, per
        portare a termine una parte o un intero progetto.
  \item Non tutte le aziende usano freelance: Riot e Blizzard lo facevano
        pochissimo e preferivano assumere (magari trattando poi la persona
        ``come un freelance'', lasciandola a casa a fine progetto).
  \item Spesso i lavori offerti come freelance sono in realtà
        \textbf{sostitutivi di un full time} (``freelance, ma sei ore al
        giorno'').
  \item Si può fare, ma bisogna essere molto chiari nel definire l'impegno
        (quante ore, come vengono tracciate): non tutti specificano cosa
        significa davvero ``freelance''. Lavorando così per due o tre aziende
        si può arrivare a un tempo pieno.
  \item È più comune il \emph{work for hire}: aziende che forniscono personale
        ad altre aziende; si è assunti da X ma si lavora ``embeddati'' presso
        Y (una sorta di subappalto). Anche lo studio del relatore lo fa.
  \item Il puro freelance di quattro mesi esiste, ma di solito per quei
        quattro mesi diventa un lavoro a tempo pieno; è raro ``quattro mesi,
        due ore al giorno''. Nelle grandi aziende è improbabile lavorare nello
        stesso mese per Ubisoft, Riot e Blizzard tre ore ciascuna; è più
        realistico lavorare per più startup indie italiane, che non possono
        permettersi un'assunzione ma possono pagare due ore al giorno.
  \item Molto passa dai contatti: trovare le aziende giuste e rendere esplicite
        le condizioni.
\end{itemize}

% ---------------------------------------------------------------------
\section{Il circolo virtuoso: feedback, miglioramento, fiducia}
\label{sec:z1-triangolo}

\subsection{Il triangolo}

L'obiettivo a cui il workshop vuole arrivare è un \textbf{triangolo che si
autoalimenta}:
\begin{enumerate}
  \item si dà un \textbf{feedback} con l'obiettivo di migliorare le
        performance (inclusi elementi caratteriali);
  \item la persona \textbf{migliora} e vede di averne beneficiato;
  \item questo genera \textbf{fiducia}, che rende più facile dare e ricevere
        il feedback successivo.
\end{enumerate}
Togliendo uno qualsiasi dei tre elementi, il circolo non è altrettanto forte.

In qualsiasi fase si parta, il rispetto per i colleghi è inizialmente dettato
dall'\emph{educazione}, non dalla fiducia. La fiducia arriva quando i colleghi
cominciano a dire cose giuste e sensate e a \textbf{mantenere la parola}.
Dare e ricevere feedback è il modo in cui, in un certo senso,
``inganniamo il cervello'' per creare una \emph{case history} di fiducia: si
ha motivo di fidarsi di più, di ascoltare di più, di essere più ricettivi. E
si diventa sempre più bravi ad accettare il feedback e ad agire di
conseguenza.

\subsection{Accettare un feedback non vuol dire fare come dice l'altro}

Dieci anni fa il relatore non era ostile ai feedback, ma molto
\emph{difensivo}: ``sì, sì, va bene, ora ci guardo''; oppure, a chi gli diceva
``avresti dovuto farla così'', rispondeva spiegando perché l'aveva fatta in
quel modo. Oggi ha imparato a dire ``ok, grazie, ora ci ragiono'': non deve
difendere la propria posizione da chi viene solo a dire ``hai pensato a
questa cosa?''. Prima lo viveva come un attacco (``mi viene a dire che non so
fare il mio lavoro''), ora come un ``ottimo punto di vista, lo prendo in
considerazione''.

Con una postilla fondamentale: \textbf{accettare il feedback non significa
fare come dice l'altro}. ``Mi hai detto che preferisci questo mob verde. Ho
pensato anche a farlo verde? Sì. Penso comunque che giallo sia meglio? Sì, lo
faccio giallo.'' Il relatore accetta \emph{sempre} il feedback, valido o no:
garantisce di tenerlo a mente quando lavora. Che poi faccia come gli viene
detto è un altro discorso.

\subsection{Fiducia e scadenze}

La fiducia funziona allo stesso modo con gli impegni. Se si chiede a qualcuno
una cosa per venerdì e la consegna venerdì, la volta dopo gliela si chiederà
per venerdì. Se invece arriva il giovedì della settimana dopo, la volta dopo
gli si darà una \textbf{deadline anticipata}, come con chi arriva sempre in
ritardo agli appuntamenti (``se ti dico alle 7 arrivi alle 7 e mezza; la
prossima volta ti dico alle 6 e mezza'').

Al lavoro questo è dannosissimo, perché crea \textbf{deadline fittizie} che
non sono mai quelle vere, e nascono discussioni del tipo ``perché devo farlo
in due giorni se si fa in cinque?'' (``sì, se tu lo facessi in cinque''). La
conversazione giusta è un'altra: \emph{``che problemi hai avuto nel
consegnare venerdì?'' ``Questo, questo e questo.'' ``Cosa posso fare per
evitare che succeda?''}

Il relatore non vuole mai sentirsi dire ``non sono riuscito a farlo''.
Preferisce di gran lunga sentirsi dire, tre giorni prima, \textbf{``potrei non
riuscire a farlo, per questi motivi''}. Chi ha la visione d'insieme può dare
più budget, più persone, spostare la deadline, cambiare lo scopo del progetto:
può fare qualsiasi cosa per aiutare. Ma se lo si scopre il venerdì, cinque
minuti prima della deadline, non può fare più niente. Il messaggio è:
\textbf{``aiutami a poterti aiutare''}. Non penserà mai male di chi dice ``non
ce la faccio per questi motivi'', perché non sta dicendo ``non ho voglia'': sta
indicando i problemi, e lavorare su quei problemi è compito del manager. Non
può invece fidarsi di chi fino al giorno prima dice ``tranquillo, ce la
faccio'' e il giorno dopo ``non ce l'ho fatta''.

Questo è fondamentale per chi lavora in programmazione, design o art, perché
ci saranno sempre deadline: bisogna sempre \textbf{lavorare sulle aspettative},
dire al proprio manager (o al proprio report) cosa è lecito aspettarsi. Non
bisogna fare come i pinguini di \emph{Madagascar} (``sorridete e salutate, va
tutto bene''): se c'è un problema, va detto, altrimenti si peggiora solo la
propria situazione.

Da italiani si tende a non voler mai alzare la mano per chiedere aiuto
(``no, no, ce la facciamo, lo faccio io''). A volte non ce la si può fare, non
perché non si è capaci, ma perché vari fattori lo impediscono. Accettarlo è
importante: \textbf{è molto meglio avvisare e chiedere una mano che dire
``non ce l'ho fatta''}.

\subsection{Domande sul dare e ricevere feedback}

\paragraph{``Ne parliamo a fare, se poi fai come vuoi?''} Uno studente (Cristian)
riceve feedback che, dopo averli analizzati, non ritiene congrui; quando non
li segue gli viene detto ``ne parliamo a fare?''. C'è un problema di
\textbf{aspettative}: l'idea che accettare un feedback voglia dire fare come
dice l'altro. Accettarlo significa prendere in considerazione ciò che è stato
detto. Chi decide diversamente ha però un \textbf{onere}: argomentare perché
si fa comunque a modo proprio (``mi hai detto verde, la facciamo gialla
perché\dots''). Poi bisogna solo aspettare i risultati per vedere chi aveva
ragione.

La decisione spetta sempre \textbf{a chi ha la responsabilità di quella
decisione}. Se il game director deve scegliere tra A e B, anche se il
feedback dice B, la scelta è sua perché la responsabilità è sua. Non sarebbe
corretto essere ritenuti responsabili di una scelta delegata a chi dà
feedback. La risposta può essere semplicemente: ``l'ho tenuto in
considerazione, ma visto che ne sono responsabile io e le conseguenze sono
mie, ho scelto di fare X''.

\paragraph{E se l'altro fa ripicca?} Un partecipante chiede di chi, vedendo
il proprio feedback non accolto, si vendica ignorando a sua volta feedback
buoni. Per il relatore è un problema caratteriale: chi si comporta così ``è
già sulla strada per uscire''. Non c'è spazio per l'ego in team di questo
tipo. Il fenomeno nasce dall'idea che \emph{il feedback negativo faccia fare
brutta figura}, che il fatto che si potesse fare meglio tolga qualcosa alle
proprie capacità.

Non è così: \textbf{il feedback non è un giudizio}. Si può passare un anno a
Riot ricevendo solo feedback negativi e avere una valutazione annuale
stellare, perché ricevere molti feedback negativi significa essere una persona
approcciabile, che agisce sul feedback e invita gli altri a darlo. Il relatore
ha visto valutazioni molto peggiori per chi riceveva solo feedback positivi,
perché la gente aveva paura di dargliene di negativi: ma solo ``pacche sulle
spalle'' non fanno migliorare. Il problema va affrontato parlando con la
persona; a volte le conversazioni iniziano con \textbf{``ti do un feedback sul
feedback''}, cioè sul modo in cui l'altro ha dato il suo.

Esempi dalla sua esperienza:
\begin{itemize}
  \item lui chiede spesso di \textbf{non ricevere feedback senza esempi
        tangibili}, perché fa fatica ad astrarre: ``sei stato maleducato'' non
        significa nulla, ``ieri al meeting mi hai risposto male su questa
        cosa'' sì;
  \item ad altri serve l'opposto: ha ricevuto il feedback ``quando devi dirmi
        che ho sbagliato, puoi girarci un po' intorno e darmi del contesto?'',
        perché il suo stile diretto (``questa cosa non ha funzionato, dovremmo
        fare così'') feriva. Ha adattato il suo modo di dare feedback: con
        alcuni è diretto, con altri gira intorno. L'importante è saperlo.
\end{itemize}

\subsection{Non portare la soluzione}

Un consiglio centrale: \textbf{quando si dà un feedback, non si dà una
soluzione}. Non si dice ``io farei così'', ma ``secondo me questa cosa si
potrebbe fare diversamente, per questo e questo motivo: come potremmo farla
diversamente?''. È come chi smette di fumare perché vuole e chi smette perché
glielo dicono moglie e figli: se non si è convinti, spesso non si smette.

Quando il relatore diceva ``c'è questo problema, io farei così'', il meeting
successivo iniziava inevitabilmente con ``ci ho provato come hai detto tu, con
me non ha funzionato''. Sempre. Perché la soluzione era sua. La
responsabilità della soluzione deve essere invece \textbf{congiunta}: ogni
feedback deve creare un piccolo team che ha come obiettivo risolvere quel
problema. Questo toglie l'antagonismo (non ``ti dico che non sai fare il tuo
lavoro'', ma ``ti voglio aiutare a farlo meglio, lavoriamoci insieme'') e
sposta il focus: il problema non sei tu, è ciò che sta accadendo.

\paragraph{E se sono i piani alti a tenere all'oscuro dei problemi?} È una
questione di cultura aziendale. Il relatore fa sempre presente al proprio capo
il concetto di \textbf{service level agreement} (che verrà ripreso nella
seconda giornata): ``se mi dai questi input, ti garantisco questo output in
questi tempi''. Se il capo chiede delle cose senza dire budget, tempi e
perché, non può aspettarsi lo stesso risultato. È una comunicazione a due
vie: si dice al manager cosa serve per fare il lavoro, e lui deve fornirlo o
accettare che il risultato non sarà quello voluto.

\subsection{Amici, parenti e parole proibite}

\paragraph{``Non si lavora con amici e parenti''?} Per sua definizione
personale il relatore ha pochissimi amici (``stanno sulle dita di una mano'')
e una soglia molto alta per considerare amico un collega. Allo stesso tempo,
in quindici anni e con centinaia di colleghi, ha avuto solo \textbf{due}
esperienze in cui non è riuscito a lavorare con qualcuno: non serve essere
amici per far funzionare le cose. Con gli amici, anzi, gli viene più facile:
se dice a un amico che una cosa ``fa schifo'', l'amico non pensa ``come si
permette'', ma ``se perfino Stefano, che è mio amico, me lo dice, allora è
proprio così''. Diventa un rafforzativo, grazie alla fiducia costruita prima.
Il problema del lavorare con amici e parenti nasce dal \emph{non voler} dire
all'amico che sta sbagliando; se invece la confidenza permette di dirlo, si
può saltare la ``confezione'' del messaggio e passare direttamente al
contenuto.

\paragraph{Le frasi peggiori.} Alla domanda se esistano parole peggiori di
``ottimo lavoro'', il relatore indica la sua classifica: ``ottimo lavoro''
(detto in modo vuoto)\footnote{Il passaggio della trascrizione è
frammentario: si intende verosimilmente il complimento generico e di
circostanza, che non dice nulla di utile.}; poi la frase che lo ``fa uscire di
testa'', \textbf{``abbiamo sempre fatto così''} (``e infatti avete chiuso'',
come Blockbuster); infine \textbf{``siamo una famiglia''}, riferito al posto di
lavoro. Un'azienda non è mai una famiglia, perché non ci si può aspettare lo
stesso impegno da chi la possiede e da chi ci lavora: non c'è un
corrispettivo proporzionale tra impegno e ritorno. Per quanto inclusiva,
flessibile e razionale, non sarà mai una famiglia.

\subsection{Ego, muri e il peggior modo di scegliere un leader}

\paragraph{Chi mette davanti l'ego e alza muri.} Esistono, ma di solito durano
pochissimo in azienda. Non si tratta di avere un ego forte, ma di \textbf{non
saper ricevere feedback}. Spesso sono persone bravissime nel loro lavoro che
poi si chiedono perché tutti intorno diventino più bravi di loro: perché gli
altri accettano feedback e hanno più fonti da cui attingere, mentre loro si
sono rintanati nella loro torre.

Molti manager tendono ad \textbf{assumere persone un filo peggiori di loro},
per non fare brutta figura. Il relatore assume solo persone più brave di lui,
``perché non è una gara'': ha lavorato con molti community manager più bravi
di lui come community manager, ma lui era più bravo a fare il manager. Non
viene valutato su quanto sa fare il community manager, ma su quanto il suo
team sa farlo. Un partecipante chiede se la paura di avere qualcuno di
migliore esista anche in ambito accademico. Assolutamente sì: al relatore è
capitato di non essere preso come docente in un corso di formazione sui
videogiochi perché il responsabile aveva percepito che ne sapesse più di
lui\footnote{La trascrizione dice ``che non sapevo più di lui'', ma il senso
del passaggio richiede il contrario.}. ``Ma perché te ne importa? Tu sei quello che ha portato chi
sa fare il suo lavoro: ben venga.'' O si ha questa cultura, o si ha paura per
il proprio posto.

\paragraph{Promuovere il migliore.} Quando bisogna scegliere un leader per un
team (programmazione, design, art), molto spesso si sceglie \textbf{il migliore
tecnico}: il miglior programmatore, il miglior designer. È spesso la scelta
sbagliata: è come prendere un idraulico per fare l'elettricista. Fare il
\emph{people manager} non è come fare il designer o il programmatore; si può
essere un ottimo programmatore e un pessimo manager, e viceversa. Le aziende
l'hanno capito e lo fanno meno di prima. A Blizzard c'era un grave problema di
\emph{middle management}: molte persone in ruoli importanti erano state
promosse cercando le skill sbagliate, erano frustrate dal non saper fare il
proprio lavoro e schiacciavano chi stava sotto per paura di fare brutta
figura.

Un partecipante cita una teoria simile: una persona viene promossa finché non
raggiunge \emph{il proprio livello di incompetenza}, dove si blocca\footnote{È
il \emph{principio di Peter} (Laurence J.\ Peter, 1969), citato in lezione
come ``la teoria di Peter''.}. Un altro cita Michael Scott di \emph{The
Office}, su cui il relatore tornerà (sezione~\ref{sec:z1-quando}).

% ---------------------------------------------------------------------
\section{Perché la cultura conta: retention, costi e fiducia}
\label{sec:z1-cultura}

\subsection{L'ambiente di lavoro e la retention}

Perché tutto questo è così importante? Perché l'ambiente di lavoro conta
moltissimo per la \textbf{retention}, cioè per trattenere le persone e farle
lavorare in modo efficiente. Se contasse solo lo stipendio, molti non
lascerebbero mai Riot o Blizzard; e più si va avanti con l'età, più si impara
a dare valore anche al resto.

Inoltre, secondo gli studi citati nelle slide, chi vive il proprio lavoro in
modo negativo è \textbf{meno efficiente nel problem solving}: l'energia
mentale va verso ``come faccio a non farmi dare la colpa'' o ``come faccio a
fare bella figura'', non verso la soluzione del problema. In ambienti negativi,
dove ci si deve preoccupare della propria ``sopravvivenza'', si lavora
peggio.

\subsection{Il costo di chi se ne va}

Ogni volta che qualcuno si dimette o viene licenziato, il costo di trovare
una nuova persona, fare i colloqui, formarla e farle attraversare forming,
storming, norming\dots\ è circa \textbf{un terzo del suo stipendio}. Per questo
le aziende tendono a fare più \emph{retention} che \emph{acquisition}, almeno
con le figure considerate utili alla produzione.

Allo stesso tempo si vedono sempre più spesso \emph{layoff} e \emph{hiring
wave}: si licenziano 3.000 persone e sei mesi dopo se ne riassumono 2.500 o
3.500, ma tutte più junior. Il relatore ne ha scritto in un articolo su
Eurogamer.it, in cui spiega anche alcuni ``trucchi del mestiere'': perché i
licenziamenti avvengono sempre in certi periodi e sempre ``a gruppi'' (prima
Bethesda, poi Blizzard, poi un'altra\dots).

\paragraph{L'offerta di Riot.} A Riot esisteva una pratica (di cui il relatore
non ricorda il nome) per cui, firmato il contratto, si aveva \textbf{un anno
di tempo per dire ``non mi piace, non me la sento, me ne vado''} e ricevere
subito come buonuscita un anno di stipendio. Avevano stimato che fosse molto
meglio pagare subito un anno a qualcuno e lasciarlo andare che tenerlo due,
tre, quattro anni come \emph{underperformer}: la responsabilità della scelta
di assumere era dell'azienda\footnote{La pratica, nota come ``The Offer'', è
descritta in fonti pubbliche con condizioni diverse (un periodo più breve e
una somma pari a una frazione dello stipendio annuo). Si riportano le
condizioni come raccontate in lezione; andrebbero verificate.}. L'hanno usata
pochissimi, ma il solo fatto che esistesse faceva lavorare meglio anche lui da
manager, perché sapeva che chi non si trovava bene poteva andarsene con un
anno di tempo per trovare un altro lavoro.

\paragraph{Il ricircolo.} C'è un forte ricircolo nell'industria: molti ex
colleghi passano da Ubisoft a Paradox, da Paradox a Funcom, da Funcom a
Blizzard. Il relatore ha scelto di uscirne lavorando solo per aziende
\textbf{indipendenti}: Stormind non è un indie di cinque persone, ma è
indipendente, e non deve accontentare azionisti. Quando vede un'azienda con
utili miliardari licenziare migliaia di persone si chiede ``dov'è la linea''
(\emph{when enough is enough}): con azionisti che vogliono un ritorno, se
quello è l'obiettivo stanno facendo la cosa giusta, ma lui non condivide
l'obiettivo\footnote{In lezione si parla di ``6,5 miliardi di utile'' e di
``3.000 persone licenziate''; le cifre sono riportate come dette e non sono
state verificate.}.

\paragraph{L'open space vuoto.} A Riot, durante un'ondata di licenziamenti,
del suo team di trenta persone è rimasto solo lui: ``teniamo te come
coordinatore, poi ricostruirai il team''. Ha rifiutato di ricostruire un team
cancellato ``con un colpo di spugna'' e ha deciso che non avrebbe più lavorato
per una grande azienda. Nessuno aveva pensato di spostarlo, e ha lavorato per
tre mesi in un open space di cento metri quadri con la sua scrivania e altre
ventinove vuote. Su una lavagna aveva scritto i nomi dei suoi ragazzi e li
cancellava man mano che, anche grazie ai suoi \emph{referral}, trovavano
lavoro. Cancellato l'ultimo nome, si è dimesso: per lui il lavoro era finito.

\paragraph{Team che si dissolvono.} Uno studente cita un articolo su Larian
Studios, il cui management è contrario a ricostruire ogni volta il team a fine
progetto. È lo stesso problema: si pensa ai ``giochi Blizzard'' o ai ``giochi
BioWare'', ma le persone che li hanno fatti non ci sono più; è rimasto solo il
brand, e non ha senso avere le stesse aspettative da un gioco fatto da persone
diverse. E anche chi sopravvive ai layoff vive pensando che il prossimo
potrebbe essere lui: alcune persone, licenziate nell'ultima ondata a
Blizzard, gli hanno detto ``erano due anni che lo sapevo''. Vivevano a Cork,
``un paesotto irlandese dove o vai per lavorare a Blizzard o non ci vai'',
avevano comprato casa: non ci si può aspettare chissà che morale.

\subsection{La cultura è forte quando si lavora per gli altri}

\begin{quote}
  \emph{La cultura è forte quando le persone lavorano per gli altri e a favore
  degli altri, e non per se stesse.}
\end{quote}

Se si dà feedback per far migliorare l'altro e non per imporre il proprio
modo, se si aiuta per il bene del progetto e non per fare bella figura, la
cultura si rafforza. È dai \textbf{piccoli gesti} che emerge la cultura
aziendale: ad esempio, dal non dover chiedere ``perché mi stai aiutando?''.

\paragraph{L'incognita.} \textbf{L'incognita è spesso dannosa quanto le
cattive intenzioni.} In un esperimento (``mi pare a Stoccolma, cinque o sei
anni fa'') si distribuivano per strada banconote da 500 euro, e la maggior
parte delle persone le rifiutava chiedendosi ``dov'è la fregatura?''
\footnote{L'esperimento è riportato in modo vago in lezione e non è stato
possibile identificarlo con certezza.}. Lo stesso accadeva al relatore,
arrivato a PG e a Stormind da Riot: a ``vedo che sei in difficoltà, posso
aiutarti?'' la risposta era sempre ``perché?''. Alla fine rispondeva: ``perché
quando sarò in difficoltà io vorrei che tu mi aiutassi''.

Quando si comincia a lavorare per gli altri e si smette di pensare ``così avrò
la promozione, così avrò l'aumento'', promozioni, aumenti e rispetto arrivano
da soli. A Blizzard aveva fallito tre volte i colloqui per essere promosso a
manager (inteso come coordinatore del team); quando ha deciso di smettere di
pensarci e di fare solo il bene del proprio lavoro, la volta dopo è stato
promosso. Funziona così: chi aiuta i colleghi e fa andare bene il progetto
viene notato (``questa persona fa migliorare cinque persone, non solo se
stessa: vale di più'').

\subsection{Il lavoro è valutato come opera d'insieme}

Il relatore riconosce che buona parte del suo successo è venuto dal fatto che
il suo team lavorava bene, non dalla sua bravura personale. Anche i progetti
degli studenti, presentati a Diego, a Enrico, ad Andrea, verranno valutati
\textbf{come opera d'insieme}. Nessuno dice ``il \emph{combat system} di questo
gioco l'ha fatto X, che bravo'': si dice ``BioWare ha fatto un bel gioco'',
``Blizzard ha fatto un bel gioco''.

Lo ha replicato anche lato HR. In molte aziende ci sono \textbf{bonus
personali} legati a obiettivi (se il gioco vende X, se la feature è pronta in
tempo). Il relatore li ha dimezzati: metà restano personali, metà diventano
\textbf{bonus di team}, che tutti ricevono se il team raggiunge gli obiettivi.
Un incentivo concreto a collaborare: non solo si dice che il lavoro di gruppo è
importante, lo si dimostra dandogli più valore economico. Aveva proposto il
100\% di bonus di team, ma gli è stato concesso il 50 e 50; ha funzionato
comunque, anche se sarebbe stato ancora più potente. Ovviamente, alla
domanda ``perché mi aiuti?'' si può anche rispondere ``perché voglio il
bonus'': un aiuto dato ``per il motivo sbagliato'' era un compromesso che era
disposto ad accettare.

\subsection{Il coraggio di chiedere aiuto}

Un partecipante osserva che aiutarsi dovrebbe essere la norma. Il modello
diffuso, invece, è ``devi essere forte e indipendente, non chiedere mai una
mano'', e vale in tutte le direzioni (dal maschilismo alla ``donna forte e
indipendente che non ha bisogno di nessuno''): \textbf{non ammettere di avere
bisogno di aiuto è sempre uno dei motivi per cui le aziende vanno male}, che si
tratti di uomini, donne o persone non binarie. Il problema è non avere la
sicurezza di poter alzare la mano.

Il relatore ha cercato di togliere ogni conseguenza dall'aver alzato la mano:
se si \emph{depenalizza} qualcosa, le persone non sono più disincentivate a
farla. Ha detto espressamente che non avrebbe mai valutato nessuno per il
numero di volte in cui chiedeva aiuto o per il numero di feedback negativi
ricevuti, e questo ha funzionato.

\subsection{Perché si dà un feedback}

\emph{Perché avete dato il vostro ultimo feedback? Che cosa volevate ottenere?}
Le risposte della classe:
\begin{itemize}
  \item per \textbf{riconoscere il valore} di un collega, anche verso il
        manager, perché qualcuno che non aveva visibilità vedesse quanto
        quella persona fosse efficace;
  \item per \textbf{spingere una persona a superarsi};
  \item per aiutare un collega a rendere il prodotto all'altezza di ciò che
        voleva rappresentare;
  \item per \textbf{genuinità}: c'era qualcosa che poteva essere fatto meglio.
\end{itemize}
Il relatore aggiunge che la forma cambia col tempo: all'inizio era ``devo
darti un feedback, hai dieci minuti?''; ora è un messaggio informale su
Discord, perché ha costruito un rapporto di fiducia e sa che non verrà letto
come ``l'HR manager mi scrive, sono nei guai?'', ma come ``Stefano è così, mi
dà feedback, resterà tra noi e lo risolveremo insieme''.

Un altro motivo è \textbf{risolvere incompatibilità caratteriali}: le
conversazioni più difficili, del tipo ``io e te non stiamo funzionando, cosa
possiamo fare diversamente?''. Sono importanti perché, se due responsabili non
vanno d'accordo, il team se ne accorge, comincia a raccontarsi storie (``si
odiano, ci saranno problemi'') e, per difendere il proprio manager, comincia a
giudicare male il lavoro dell'altro gruppo. È la dinamica da fermare:
\textbf{valutare gli altri solo per difendere la propria fazione}.

Dare e ricevere feedback sono \textbf{due muscoli da allenare} (il relatore è
sempre stato più bravo a darlo che a riceverlo), e si impara col tempo. È come
fare colloqui: fare colloqui è una skill a parte, che non indica quanto si è
bravi nel proprio lavoro. Essere scartati non significa non saper fare quel
lavoro: spesso significa aver fatto male il colloquio o che è stato preferito
qualcun altro.

\paragraph{Feedback pubblici e privati.} Uno studente racconta di aver dato un
feedback negativo a un superiore, senza un motivo particolare, per dimostrare
a sé e ai colleghi che non si viene licenziati per questo. Il relatore coglie
l'occasione per precisare che \textbf{i feedback negativi si danno
individualmente}: lui dà moltissimi feedback di cui il resto dell'azienda non
sa nulla. Reagirebbe molto male se, in un meeting, qualcuno alzasse la mano per
dire ``comunque questa cosa non l'hai fatta bene'': quello non è feedback, è
discussione, accettabile solo se è un brainstorming e se l'opinione è stata
chiesta. Alcune aziende hanno piattaforme pubbliche di feedback, ma spesso
servono solo ad applaudire; lo stesso accade nelle aziende che
\emph{obbligano} a lasciare un certo numero di feedback a settimana, che
diventano ``lui è bravo, lei è competente, anche questa settimana l'abbiamo
fatto''. Bisognerebbe chiedere a chi le adotta cosa pensa di ottenere (forse
trasparenza), ma difficilmente una persona ``attaccata'' pubblicamente reagisce
bene.

% ---------------------------------------------------------------------
\section{Among Us: la caccia al traditore}
\label{sec:z1-among-us}

Prima della pausa pranzo il relatore mostra tre film di fantascienza di
qualche decennio fa e chiede che cosa li accomuni\footnote{I titoli non sono
leggibili nella trascrizione: viene citato \emph{Event Horizon} da uno
studente, e il film ``in mezzo'' è indicato come ``l'originale'',
verosimilmente \emph{La cosa} (\emph{The Thing}).}. In ognuno c'è un gruppo di
persone con un obiettivo e la presenza, reale o percepita, di qualcuno che
rema contro: \textbf{l'energia del gruppo viene spesa per trovare il
responsabile}.

\paragraph{L'esperimento degli alieni.} Nei master in presenza il relatore
propone un esercizio. Distribuisce segretamente cartellini plastificati con
scritto ``essere umano'' o ``alieno''; le carte alieno contengono istruzioni di
comportamento con cui gli alieni, che cercano di fingersi umani, si
tradirebbero (dare sempre ragione a chi ha parlato prima, finire ogni frase
con una domanda, parlare in un certo modo). Lo scopo: in dieci minuti formare
gruppi, e \textbf{vince il gruppo più numeroso senza alieni}; basta un alieno
per perdere. Si formano sempre gruppi piccoli (due, cinque, sei persone),
perché la tendenza è ``mi fido al massimo di voi due''.

Ma vincono i gruppi più grandi, per un motivo semplice: il relatore
\textbf{non distribuisce nessuna carta alieno}. Mostra a tutti entrambe le
carte, ma consegna solo carte ``umano''. Vuole simulare ciò che accade in
ufficio: la tendenza a percepire sempre qualcuno che rema contro (i manager,
l'altro team, quel collega), quando in realtà non c'è nessuno, ``è tutto nella
nostra testa''. Quando chiede le motivazioni, emergono sospetti
\emph{fabbricati}: ``faceva sempre domande'', ``parlava solo all'infinito''
(``no, parla così e basta''). Ci si fabbrica da soli i motivi per non fidarsi
degli altri, ed è comunissimo.

\paragraph{Among Us.} L'autosabotaggio è, per tornare ad \emph{Among Us},
\textbf{concentrarsi sul cercare cadaveri e traditori invece di fare le
task}: non si vince perché non si sono completate le task, mentre si cercava il
colpevole. Succede tantissimo negli uffici, ed è il motivo per cui l'80\% dei
problemi che il relatore risolveva erano di comunicazione. Entrando nei team
sentiva immancabilmente: ``facciamo il possibile, ma questa lotta continua con
l'altro gruppo è estenuante''. ``Perché c'è questo problema?'' ``Non lo so,
c'era già quando sono arrivato.'' Un paradosso: c'è un problema, nessuno sa
dire quale sia, ma c'è sempre stato e quindi lo si accetta.

\paragraph{L'esperimento delle scimmie.} È l'esempio, ``finto o reale'', delle
scimmie in gabbia\footnote{È il noto racconto delle ``cinque scimmie'', spesso
presentato come esperimento scientifico ma probabilmente apocrifo (lo stesso
relatore lo definisce ``finto barra reale''); nella trascrizione il numero
delle scimmie varia.}. In una gabbia con alcune scimmie pende una banana
elettrificata: chi la tocca fa prendere la scossa a tutte. Le scimmie imparano
e, quando una si avvicina alla banana, le altre la picchiano. Poi si
sostituisce una scimmia alla volta con una nuova: la nuova va verso la banana,
le altre la picchiano, e smette. Si continua finché nella gabbia ci sono solo
scimmie nuove, poi si toglie la corrente. Nessuna scimmia tocca la banana:
nessuna ha mai preso la scossa, ma tutte hanno imparato che chi tocca la
banana prende le botte.

Nei team è lo stesso: il motivo per cui si ``odia'' l'altro team si
trasforma, ma un vero motivo non c'è. Ne derivano la sfiducia, il non alzare
mai la mano per chiedere aiuto, il non dire mai ``bel lavoro'' agli altri
perché non sia mai che qualcuno ti veda fare un complimento ``a quelli''.
Risolvere queste dinamiche è lungo e tedioso, e sempre pieno di ``non mi
ricordo'', ``io non c'ero'' e ``abbiamo sempre fatto così'', la cosa più
frustrante per chi cerca di risolverle.

% ---------------------------------------------------------------------
\section{Il principio del default to trust}
\label{sec:z1-default-trust}

\subsection{La definizione}

Il relatore indica questa come la slide più importante della presentazione.
\textbf{Default to trust} è uno dei valori aziendali di Riot Games, che lui ha
tradotto, adattato e ``esportato'':

\begin{quote}
  \emph{Approccio ogni situazione di conflitto partendo dal presupposto che una
  persona abbia agito in un certo modo perché, con le informazioni a sua
  disposizione in quel momento, ha pensato che fosse la soluzione migliore, e
  non perché sia un'incapace o volesse danneggiare il progetto
  intenzionalmente.}
\end{quote}

In altre parole: \textbf{di default, mi fido}.

\subsection{L'esempio del ritardo al daily stand-up}

A Riot il relatore aveva un \emph{daily stand-up} dal vivo alle 15 con il suo
team di venti persone. Una ragazza arrivava regolarmente alle 15:15, per una,
due, tre settimane. Alla quarta ci ha parlato. \emph{Come iniziereste la
conversazione?} Le proposte della classe:
\begin{itemize}
  \item sottolineare l'importanza di esserci dall'inizio, perché potrebbe
        perdersi qualcosa che le interessa, e per rispetto dei colleghi
        presenti;
  \item ``devi arrivare in orario'', spiegando perché è importante;
  \item \textbf{chiedere perché}.
\end{itemize}

Il relatore osserva, anche dalle sue consulenze fuori dai videogiochi, che
\textbf{più si sale con la seniority}, più la risposta è ``devi smettere di
arrivare in ritardo'', spesso con un \emph{or else} finale (``altrimenti
succede quest'altra cosa''). Più si scende con l'età, e quindi con meno
abitudini consolidate, più la prima domanda diventa ``dimmi perché arrivi in
ritardo''. Già spiegare perché è importante è ``miglia avanti'' rispetto al
semplice ``smetti''.

Partire dal perché gli ha permesso di scoprire la causa: il meeting precedente
della ragazza, in un'altra zona del campus a un quarto d'ora di distanza,
finiva sempre tardi, e lei doveva attraversare il campus di corsa. Non aveva
la confidenza per dire all'altro manager, che insisteva perché si restasse
fino alla fine, che doveva andare; sacrificava quindi il meeting del relatore,
sapendo che era il più comprensivo dei due.

Il relatore si è offerto di \textbf{mediare}: ha detto all'altro manager che
o anticipava l'inizio del suo meeting, così da finire comunque in tempo,
oppure lasciava uscire la ragazza, perché quel meeting gli serviva. C'è anche
una questione di \textbf{economia di scala}, spesso sottovalutata: se venti
persone aspettano qualcuno per dieci minuti, si sono sprecati dieci minuti per
venti. Per lo stesso motivo il workshop è iniziato alle 10:05: iniziare alle
10:10 avrebbe significato far aspettare cinquanta persone cinque minuti in
più. A Riot c'era una scrivania vuota con un pupazzo gonfiabile, \textbf{Mr.\
Late}: sommando tutti i minuti di ritardo a meeting, call e conversazioni, in
un anno veniva fuori praticamente \emph{una persona a tempo pieno}. Serviva a
ricordare che anche cinque minuti contano, e che bisogna rispettare gli altri
arrivando in orario o almeno avvisando.

\paragraph{Gli effetti.} La conversazione era avvenuta nel primissimo periodo a
Riot, quando il team non aveva ancora motivo di fidarsi di lui. Risolvendo il
problema ha eliminato anche tutto lo strato di micro-dinamiche spiacevoli che
un manager impara a vedere (due colleghi che si guardano con la risatina
quando lei entra in ritardo). La ragazza ha cominciato a entrare tranquilla,
con una certa leggerezza: non le faceva piacere arrivare in ritardo (``ci sono
anche persone che lo fanno apposta, e quello è un problema da affrontare
diversamente''). Soprattutto, con i colleghi non parlava di lui come del capo
che minacciava, ma come di chi le aveva risolto un problema: era diventato
un \textbf{alleato}. La settimana dopo altre tre persone sono andate da lui
con un problema da risolvere. Ripetuto cinque, sei, sette, dodici volte, in un
paio di mesi poteva dire quasi qualsiasi cosa nei meeting sapendo che
richieste e critiche sarebbero state percepite come ``me lo chiede perché vuole
darmi una mano'', non per fare bella figura o per presentare i successi del
team ai suoi superiori.

\subsection{Condividere i successi, schermare i fallimenti}

Un altro elemento fondamentale (ripreso nella seconda giornata): il manager
\textbf{condivide i successi e protegge il team dai fallimenti}. ``Se va male
qualcosa è colpa mia; se va bene qualcosa, voglio che ne parliate voi.'' Il
lavoro del manager non dovrebbe essere valutato su ciò che presenta del lavoro
del team: tutti sanno comunque chi è il manager. Tanto vale lasciar presentare
il team, che potrà dire con orgoglio di aver presentato a Mike
Morhaime\footnote{Cofondatore ed ex presidente di Blizzard; nella
trascrizione ``Mike Morayme''.} la propria campagna europea per
\emph{StarCraft~II}. Ci sono manager che dicono ``come mi prendo le colpe, mi
prendo i meriti'', ma i meriti li hanno già: è inutile ostentarli. Con i
fallimenti invece il manager è la \textbf{rete di sicurezza} del team. Il
relatore ne ha fatto una questione di principio: \textbf{non scaricare mai la
colpa su qualcun altro}.

\paragraph{Colpevole o causa.} Un esempio che il relatore attribuisce a Pirate
Software, creatore di contenuti ed ex dipendente Blizzard\footnote{Il relatore
aggiunge che si tratterebbe del figlio dello sviluppatore a cui si sarebbe
ispirato il personaggio del giocatore ``senza vita'' nell'episodio di
\emph{South Park} dedicato a \emph{World of Warcraft}. L'informazione è
riportata come detta in lezione e non è verificata.}. Contattando il servizio
clienti Blizzard per motivi personali, aveva scoperto che i colleghi gli
avevano scritto nei ticket informazioni confidenziali che non avrebbero dovuto
condividere in quel modo, nemmeno con un collega. Non ha segnalato le persone
da punire: ha detto ``c'è un problema, evidentemente le policy non sono
chiare''. Quando gli hanno chiesto i nomi, ha rifiutato: ``se vi dico chi è,
punite quelli, e gli altri cinquecento continueranno a farlo; non avremo
cambiato la causa alla fonte''.

\begin{quote}
  \emph{È inutile trovare il colpevole. Voglio trovare la causa. Perché il
  colpevole cambia.}
\end{quote}

Il colpevole a volte è vittima della situazione: per questo la definizione
dice \emph{``con le informazioni a sua disposizione in quel momento''}. Se il
relatore avesse rimproverato la ragazza in ritardo, non solo non avrebbe
risolto il problema, ma avrebbe creato un'\textbf{ingiustizia}: lei, che
attraversava il campus di corsa per una colpa non sua e di cui non poteva
parlare, sarebbe stata anche rimproverata.

Nessuno fa un brutto lavoro apposta: nessuno è contento quando sbaglia o
perde. Perché al lavoro diamo per scontato che gli altri sbaglino
intenzionalmente? Se qualcuno ha fatto una cosa male, è perché in quel momento
credeva fosse la scelta migliore. La domanda giusta non è ``perché hai fatto
questa cosa che è stata un fallimento?'', ma:

\begin{quote}
  \emph{Che cosa ti ha fatto pensare che questa fosse la scelta giusta?}
\end{quote}

Magari erano limiti di budget, di tempo, di scopo: si agisce su quelli, non
sulla persona. Se si cambia la persona ma le si danno gli stessi fattori, è
molto probabile che arrivi alle stesse conclusioni, e ci si troverà a dover
``disciplinare'' anche lei.

\paragraph{La reputazione costruita sui rumor.} Durante il giro di
presentazione in un ufficio (non in nessuna delle aziende citate), al relatore
è stato detto: ``questo è il dipartimento di design; se ti serve qualcosa
chiedi a lei, non a lui, perché lui consegna sempre in ritardo''. Oltre a
trovarlo poco professionale, una settimana dopo ha chiesto: ``cosa ti ha fatto
questa persona?'' ``A me niente, però ho sentito che\dots''. Si torna alle
scimmie: magari quella persona non consegna in ritardo da mesi, ma ormai lo
dicono tutti. È un circolo vizioso da fermare subito.

\paragraph{Niente escalation inutili.} Lo stesso vale per il feedback: quando
qualcuno va dal relatore dicendo ``ho un problema con questa persona'', lui
chiede: \textbf{``perché ne stai parlando a me e non a lei?''}. ``Perché sei il
mio capo'' non basta. Passare da lui significa dare per scontato che l'altra
persona non abbia interesse a risolvere, mentre magari basta parlarle. È la
strada più facile (``lo dico a Stefano, ci pensa Stefano''), e proprio per
questo bisogna mettere dei paletti: il manager interviene quando serve, ma le
cose che si possono gestire da soli vanno gestite da soli. (Quando invece il
confronto diretto fallisce, l'escalation è legittima: si veda la
sezione~\ref{sec:z1-qa-finale}.)

% ---------------------------------------------------------------------
\section{Traiettorie di carriera e valutazione delle performance}
\label{sec:z1-traiettorie}

\subsection{Quale traiettoria vorreste?}

Il relatore mostra un grafico con quattro curve colorate, contrassegnate
anche da lettere (A, B, C, D) per eventuali persone daltoniche (``a me è
capitato che qualcuno scoprisse di esserlo durante una lezione''). Le curve
rappresentano la performance lavorativa nel tempo. Due domande: \emph{quale
vorreste essere? E quale pensate realisticamente di essere?} Un partecipante
aspira a B e pensa di essere C; quasi tutti scelgono B.

Perché non A, che sale più rapidamente? Un partecipante risponde che una
crescita così repentina può essere più dannosa che utile, come una ``rush'' di
studio molto faticosa che poi non si riesce a sostenere. Il relatore propone
una lettura suggerita da un amico psicologo, che spiega la forma del grafico:
di B, C e D \textbf{vediamo dove finiscono}; di A no. Potrebbe continuare a
salire oltre la scala del grafico, o crollare sotto D. Non saperlo ci spinge a
preferire B. (A chi obietta che anche le altre proseguono oltre il bordo, il
relatore precisa: il bordo destro è ``oggi'', e di A non sappiamo come sia
arrivata fin lì.)

\subsection{La carriera non è una linea retta}

Tutti immaginano la carriera come una linea: si parte da X e, attraverso
stage e posizioni, si arriva a Y. In realtà carriera e performance avranno un
andamento irregolare, con alti e bassi. Come non si accetta che un team possa
temporaneamente \emph{underperformare}, tanto meno lo si accetta da se stessi.
Bisogna accettare che ci saranno momenti bassi.

Per questo il relatore \textbf{detesta le valutazioni annuali delle
performance}: prendono un punto casuale della traiettoria, che se si è
fortunati è il picco più alto e se si è sfortunati il più basso, e lo usano
per valutare un intero anno. Insegna ai manager a fare più punti di controllo
e a guardare altre cose:
\begin{enumerate}
  \item \textbf{quanto tempo passa tra un picco e l'altro}, cioè quanto
        \emph{cooldown} serve a una persona. È il motivo per cui, dopo la
        Games Week, tutti staccavano una o due settimane. Il relatore stesso,
        che si impegna moltissimo e tende al burnout, non potrebbe rifare lo
        stesso workshop il lunedì e il martedì successivi con la stessa
        qualità;
  \item \textbf{la differenza tra la miglior performance e la peggiore}: che
        cosa ha portato a lavorare così bene, e così male? Con chi si
        lavorava, su cosa ci si concentrava, quanto tempo e budget c'erano?
        L'obiettivo è capire come \emph{replicare} le condizioni che fanno
        lavorare bene, invece del solito ``hai lavorato male, è colpa tua''.
        Il relatore ha avuto molti manager per cui quando lavorava male era
        colpa sua e quando lavorava bene ``che bravo'', ma sapeva che raramente
        era l'unico fattore della propria performance. È la fallacia del ``se
        sai fare 100, perché non fai sempre 100?'', come nella scena di
        \emph{Schindler's List} sull'operaio che produce una vite in pochi
        secondi: non sempre un ritmo è sostenibile;
  \item \textbf{il fallimento di oggi confrontato con quello di sei mesi fa}.
        Il relatore non è contento della presentazione di oggi (l'avrebbe
        migliorata per altri due anni), ma rispetto a quelle di cinque anni fa
        vede un grande miglioramento di chiarezza: il suo peggior fallimento di
        oggi equivale al suo primo successo di inizio carriera.
\end{enumerate}
Quello che conta è la \textbf{traiettoria} dall'inizio dell'anno a oggi, non
i singoli punti. E se si riesce a capire insieme quando e perché si ha
successo e quando e perché si fallisce, si può lavorare per replicare solo le
situazioni che fanno lavorare bene.

\paragraph{Lo stesso vale per i giochi.} È il motivo principale per cui i
giochi delle grandi case vanno bene o male. Quando si riescono a replicare le
condizioni che hanno prodotto \emph{Warcraft~II}, esce \emph{Warcraft~III};
quando non ci si riesce, esce ``il nuovo \emph{Dragon Age}'': la gente non è
arrabbiata perché il gioco non è bello, ma perché non è come gli altri
\emph{Dragon Age}; non ne replica il \emph{feeling}, anche se potrebbe
essere un gioco favoloso. L'eccezione positiva è il passaggio da \emph{Dune} a
\emph{Dune~II}: il primo era di fatto un'avventura grafica, il secondo ha
generato il genere degli strategici in tempo reale. Gli sviluppatori avevano
capito che il successo non stava nel tipo di gioco, ma nell'ambientazione che
attirava le persone: un bel gioco in quell'ambientazione sarebbe piaciuto
comunque. Il relatore comprò \emph{Dune~II} aspettandosi un seguito
dell'avventura e, pur sorpreso, lo finì.

\subsection{Lying, hiding, faking}

Quando si è valutati su un singolo punto nel tempo, e da quel punto dipendono
bonus, stipendio e il rischio di licenziamento, si comincia a fare
\textbf{lying, hiding e faking}:
\begin{description}
  \item[Lying] dire che è sempre andato tutto bene, che si finirà sempre in
        orario, che non c'è nessun problema;
  \item[Hiding] nascondere: ``non era un mio task'', ``non c'ero'', ``non ho
        visto niente'';
  \item[Faking] fingere, anche in cose apparentemente innocue come ``sì, sì,
        ho capito tutto'' quando si ha solo paura di dire ``me lo rispieghi?
        Non ho capito niente''.
\end{description}

\subsection{``Non ho capito'' e la conferma delle istruzioni}

Dire ``non ho capito, ripetimelo'' è ancora visto come una figura da sciocco.
Il relatore non si arrabbierà mai se qualcuno gli chiede di rispiegare
un'intera slide: le cose ovvie per lui possono non esserlo per altri. Più si
conosce qualcosa, \textbf{più è difficile spiegarla}, perché si danno per
scontate molte cose. Ne parlava proprio con Diego, che gli diceva che fare
l'insegnante diventa sempre più difficile man mano che si diventa esperti. È
come spiegare \emph{EVE Online}, uno dei giochi più complessi in circolazione:
chi lo conosce bene dà per scontato che ``ovviamente non puoi fare questa
cosa'', mentre per un nuovo giocatore non è affatto ovvio. A scuola chi dice
``non ho capito'' viene visto come lo ``scemo''; nessuno si chiede ``non è che
l'ho spiegato male?''.

Il racconto più ``traumatico'' dei ragazzi che fanno stage nelle grandi
aziende è sempre lo stesso: ``Cinzia, devi fare questa cosa.'' Cinzia la fa.
``Non te l'avevo chiesta così.'' Il problema è che nessuno ha
\textbf{confermato che Cinzia avesse capito}. Il consiglio, valido sempre:
\begin{itemize}
  \item quando si \emph{ricevono} istruzioni, ripeterle: ``ho capito bene?
        Vuoi che faccia questa cosa in questo modo?'';
  \item quando si \emph{danno} istruzioni, chiedere a chi le riceve di
        ripeterle, non perché lo si ritenga incapace, ma per capire che cosa
        ha capito dalla propria spiegazione (``quindi vuoi un quadrato?''
        ``No, ho detto un triangolo\dots'' ``Però mi hai detto una forma con
        quattro angoli.'' ``Hai ragione: voglio un quadrato'').
\end{itemize}
Frasi come ``non te lo ripeto due volte'' o ``stammi dietro'' portano solo a
fare le cose come si pensa che vadano fatte e a sentirsi dire a fine progetto
``non la volevo così'': che si tratti di una feature, di un design o di un
modello 3D. (Un partecipante osserva che lo stesso accade insegnando le lingue:
chi insegna dà per scontato, ad esempio, che si sappia dell'esistenza del
genere neutro in tedesco.)

\paragraph{Come convincere il team a chiedere?} Il relatore mette sempre come
premessa che \textbf{l'obiettivo non è trovare la persona più intelligente
nella stanza, ma fare un buon lavoro}. Che Gianni capisca alla prima e un altro
alla quinta non gli interessa, se alla fine entrambi lavorano bene: ripetere
cinque volte è comunque meno lavoro che far rifare le cose sette volte. È un
investimento a lungo termine, e non dissimula alcun fastidio: davvero non gli
dà fastidio. Ma ripetere le stesse parole non aiuta: ``non ho sentito'' si
risolve ripetendo; ``non ho capito'' si risolve cambiando spiegazione. Meglio
ancora: ``dimmi che cosa hai capito tu, e vediamo cosa non torna''. Se
intravedesse nel team uno sguardo del tipo ``è la quinta volta che lo
chiede'', interverrebbe: spesso chi giudica è proprio chi dice ``ho capito'' e
poi sbaglia. Preferisce di gran lunga chi chiede cinque volte.

\subsection{Il feedback più utile che abbiate ricevuto}

\emph{Qual è il feedback più utile che abbiate mai ricevuto?} Il relatore ne ha
almeno dieci che gli hanno cambiato il modo di lavorare. Le risposte:
\begin{itemize}
  \item un semplice \textbf{``grazie''} per un lavoro fatto, più utile di un
        ``questa cosa fa schifo'';
  \item un feedback sulla propria \textbf{impulsività} nel rispondere: fermarsi
        un attimo e capire di cosa si sta parlando prima di intervenire;
  \item \textbf{essere se stessi}, anche al lavoro, e non scusarsi
        continuamente per ogni cosa.
\end{itemize}
I feedback che toccano di più sono quelli \textbf{veri}. Il ``grazie''
significa ``ho visto quello che hai fatto e l'ho apprezzato''. Quelli negativi
che si ricordano sono quelli che hanno fatto realizzare ``questa cosa che
pensavo di saper fare, non la so fare'', oppure ``la so fare, ma questa volta
non mi è venuta'' (non è che se una cosa una volta non riesce non la si sa più
fare: può capitare di sbagliare una porta). Arrivano sempre da un
\textbf{punto cieco}: qualcosa che non si sapeva fosse percepito dagli altri.
Conta più ciò che le persone ti dicono di ciò che pensi che pensino di te; e
ciò che il relatore pensa che gli altri pensino di lui è sempre terribile,
mentre il loro giudizio si rivela molto più positivo del previsto.

\paragraph{Il denominatore comune siete voi.} Si potrebbe pensare che questi
concetti non siano applicabili da chi arriva in un'azienda senza poterne
cambiare la cultura. Ma c'è sempre un denominatore comune a cui applicarli:
\emph{se stessi}. La cultura aziendale può non spingere a dare feedback, può
vedere male chi chiede aiuto; ma il relatore ha visto che quando lui lo fa, le
persone intorno tendono a farlo, e un gruppo che lo fa lavora meglio. Perfino
i CEO più ostici a un certo punto gli chiedevano: ``perché il tuo team ora
performa meglio degli altri?''. La differenza è il default to trust: una
pratica che col tempo fa operare il team \emph{come un team} e non come un
insieme di individui. Un videogioco è un progetto corale: c'è chi fa giochi da
solo, ma non si fa da soli \emph{Red Dead Redemption~2}; e se il gruppo non è
affiatato e allineato non esce \emph{Red Dead Redemption~2}, esce
\emph{Anthem} (``oggi la mia fissazione è \emph{Anthem}, domani sarà
un'altra'').

% ---------------------------------------------------------------------
\section{Come dare un feedback}
\label{sec:z1-dare}

\subsection{Concentrarsi sull'azione e sui suoi effetti}

Il primo principio: \textbf{il feedback non riguarda mai la persona}. Non si
dice ``tu hai fatto questa cosa che non va bene'', ma ``questa cosa non va
bene: come facciamo a non farla più?''. Che sia un comportamento, un'azione,
una decisione o una strategia, ci si concentra su quello, e si
\textbf{esplicita l'effetto} che ha. Nell'esempio del ritardo l'effetto era
``quindici minuti di attesa moltiplicati per venti persone''. In un progetto
può essere una catena:
\begin{quote}
  se non rispetti le deadline, l'animazione che viene dopo resta ferma; se
  loro sono in ritardo, la \emph{vertical slice} è in ritardo; se la vertical
  slice non è pronta per la Gamescom, non andiamo alla Gamescom e non troviamo
  un publisher; se non troviamo un publisher, chiudiamo.
\end{quote}
Più la conseguenza è grave, più dovrebbe sembrare piccolo il gesto di dire
``datemi una mano'': è il \emph{male minore}. Poi, a cose fatte, ci si dà
feedback: cosa ha funzionato, cosa no, cosa fare di più, di meno, cosa smettere
di fare.

\subsection{Il pericolo degli ambienti solo positivi}

Gli ambienti dove ci si fa sempre grandi complimenti e nessuno parla dei
problemi creano difficoltà vere alle aziende. Il problema non è mai la
mancanza di complimenti, ma un \textbf{eccesso di ottimismo}: si comincia a
pensare che le cose andranno come si spera, perché tutti intorno danno solo
feedback positivi.

A Blizzard i feedback determinavano gli \textbf{aumenti annuali}: ognuno
doveva indicare cinque persone che avrebbero dato feedback su di lui, e
ovviamente tutti indicavano gli amici, o colleghi di progetti andati bene. Un
sistema nato con le migliori intenzioni diventava uno \emph{scambio di voti}
(``se mi metti follow ti metto follow''), senza alcun miglioramento.

Al contrario, \textbf{le conversazioni più brutte sono state le più utili}:
affrontare i licenziamenti a Riot, o la cancellazione di \emph{Project Titan}
a Blizzard, il progetto che doveva essere il nuovo MMO dopo \emph{World of
Warcraft} e da cui poi è nato \emph{Overwatch}\footnote{Nella trascrizione si
legge anche ``doveva essere Warcraft 2'', probabilmente un lapsus o un errore
di trascrizione.}. Andando a vedere perché lo si cancellava, c'erano solo
motivazioni razionali e pragmatiche, di business, chiare a tutti, non un ``Mike
Morhaime ci odia''. Sono anche le conversazioni in cui ha lavorato su se
stesso: è molto \textbf{avverso ai conflitti}, e questo a volte gli ha
procurato conseguenze molto più gravi di quanto sarebbe costato un piccolo
disagio. Dieci anni fa faceva fatica ad accettare questo difetto; ha dovuto
subire le conseguenze di conflitti evitati e poi cresciuti.

\subsection{Che cosa fare}

Le linee guida della slide, che il relatore legge rapidamente perché saranno
disponibili nel materiale:
\begin{description}
  \item[Attenzione al linguaggio del corpo] dare un feedback a braccia
        incrociate o guardando altrove è diverso dal guardare l'altro negli
        occhi.
  \item[Sii specifico] con esempi tangibili.
  \item[Esplicita gli effetti] come visto sopra.
  \item[Discuti il prima possibile] non aspettare un mese per parlare di
        qualcosa successo ieri, altrimenti nessuno ricorderà più bene cosa sia
        accaduto.
  \item[Concentrati sull'azione e sul problema, non sulla persona] il nemico è
        l'azione, non la persona che si ha davanti. Non bisogna far cambiare
        l'altro: bisogna evitare che il problema si ripresenti.
\end{description}

\paragraph{Gli insulti sono rivolti al ruolo.} Un esempio solo apparentemente
scollegato. Il relatore ha lavorato su \emph{League of Legends}, ``non
esattamente noto per avere una community da orsetti del cuore'', in varie
\emph{task force} sulla tossicità. A chi smetteva di giocare veniva inviato un
sondaggio sul perché; il 33\% indicava la tossicità della community. Dopo varie
campagne di sensibilizzazione e cambiamenti si è scesi al 30\%: sembra solo
un 3\% in meno, ma valeva circa \textbf{120 milioni di revenue in più
all'anno}, perché quei giocatori restavano esposti a skin e contenuti e
continuavano a comprare. Anche qui, piccoli guadagni che si sommano.

Lavorando con degli psicologi, aveva notato che gli insulti \textbf{non sono
quasi mai rivolti alla persona, ma al ruolo}: ``il mio jungler fa schifo'',
``il mio DPS non è capace'', ``il nostro portiere è una pippa''. Chiamare
qualcuno per ruolo lo \emph{depersonalizza} e rende più facile insultarlo. Lo
stesso nell'esempio del giro dell'ufficio: non ``non chiedere a Giulio'', ma
``non chiedere a quel grafico, consegna sempre in ritardo''. Concentrarsi
sull'azione, senza attaccare la persona, fa sì che il feedback non sia (e non
venga percepito come) un attacco personale.

\paragraph{Considera il carattere.} Qui tornano utili i test di personalità.
Servono più agli studenti che al relatore (che sa già che usciranno 53 profili
diversi, senza un pattern): bisogna conoscere il proprio profilo e
\textbf{condividerlo con chi lavora con noi}, chiedendo anche agli altri di
fare lo stesso. In ufficio il relatore ha usato una versione più avanzata, a
pagamento, con un report di circa trenta pagine: il \textbf{profilo
DISC}\footnote{Nella trascrizione ``disk''. Si tratta del modello DISC
(Dominance, Influence, Steadiness, Conscientiousness).}. Ha mappato tutta
l'azienda e ha fatto stampare il profilo di ciascuno sulla sua scrivania, così
che tutti sapessero con chi stavano parlando. Capire come ragionano gli altri
(``è introverso, non gliene parlo davanti a tutti'') è una delle cose che gli
sono tornate più utili: \emph{adattare ciò che si vuole dire a chi si ha
davanti}.

\paragraph{Assicurati di non disturbare.} Sotto il profilo di ognuno c'era un
piccolo \textbf{semaforo}: verde, ``parlami pure, interrompimi''; giallo,
``interrompimi solo se è urgente''; rosso, ``mi sto concentrando, non mi
parlare''.

\paragraph{Lascia finire il lavoro.} Durante gli eventi dal vivo il relatore si
era reso conto di quanto fosse grave dare feedback mentre qualcuno sta ancora
lavorando. Se una cosa sta andando male, dirlo in quel momento non aiuta: tra
la pressione di finire e quella del feedback, peggiorano entrambe le cose, e
anche il rapporto. Meglio lasciarla finire, anche male, e parlarne dopo.

\subsection{Che cosa non fare}

\begin{description}
  \item[Non basarti su supposizioni] ``mi hanno detto che ieri in un meeting
        hai fatto\dots'': o c'eri e l'hai visto, o non esiste.
  \item[Non giudicare la persona.]
  \item[Non fornire feedback vago.]
  \item[Non accettare e non fornire feedback di terze parti] Se qualcuno vi
        dicesse ``mi hanno detto che hai fatto un brutto lavoro'', il primo
        pensiero non sarebbe ``cosa ho fatto male?'', ma \emph{``chi te l'ha
        detto? Chi è la vipera che ho in seno? Chi è l'impostore?''}. Siamo
        fatti così. Ogni conversazione che inizia con ``mi hanno detto che'' è
        inutile: l'altro non ascolterà, penserà solo a chi dare la caccia, e
        spesso penserà alla persona sbagliata, creando antipatie tra persone
        che non c'entravano nulla. Quando al relatore dicono ``mi hanno detto
        che lui ha fatto\dots'', risponde ``vai a parlarne con lui''; se gli
        dicono ``mi hanno detto che tu hai fatto\dots'', risponde ``mandami la
        persona che te l'ha detto, altrimenti non ne parliamo''. Non per
        cattiveria: se chiedesse ``che cosa ho detto di preciso a Marco?'' e la
        risposta fosse ``non lo so, non c'ero'', di cosa si starebbe parlando?
        \textbf{Parlate con la fonte}, altrimenti fate come se quel feedback non
        esistesse.
  \item[Non dare consigli non richiesti] come visto, non portare la soluzione:
        la soluzione deve venire da chi riceve il feedback.
  \item[``Non mi piace'' non è un feedback] ``Penso che questo design non
        funzioni perché, visto che il nostro obiettivo è\dots'' ha argomenti;
        ``non mi piace il combat system'' è un'opinione. Anche ``io l'avrei
        fatto diversamente'' non è un feedback.
  \item[Accettare non vuol dire fare come proponi] deve essere chiaro a chi dà
        il feedback: l'aspettativa è essere ascoltati e presi in
        considerazione, ma se la responsabilità è dell'altro, deciderà
        l'altro.
  \item[Non offenderti] se il feedback non viene seguito. Restarci male
        significa pensare di saperne più dell'altro, cioè \emph{venire meno al
        default to trust}. ``Decidi tu'' significa ``mi fido di te, perché sei
        lo specialista, e mi fido che mi ascolterai e farai la cosa migliore
        per tutti'', che potrebbe anche non essere ciò che ho detto io. Se
        voglio che l'altro faccia sempre come dico io, mi sto fidando solo di
        me stesso.
\end{description}

\subsection{Il primo feedback e il debriefing}

La cosa più difficile è dare il \textbf{primo} feedback: è un muscolo da
allenare. Chi non l'ha mai fatto (soprattutto con questa struttura) dovrebbe
provarci, anche su minuzie (``abbiamo concordato di usare WhatsApp nel gruppo:
non potremmo usare Telegram?''), anche solo per vedere come va. Altre
indicazioni:
\begin{itemize}
  \item feedback \textbf{solo uno a uno}, mai davanti ad altri, soprattutto se
        negativo;
  \item \textbf{prepararsi} prima;
  \item fare il \textbf{debriefing del feedback}: dopo, fermarsi a pensare a
        com'è andato. Il relatore, quando media un feedback tra due persone,
        si sente spesso rispondere ``benissimo'' da una e ``malissimo''
        dall'altra: la percezione è molto diversa. Non si fa problemi a tornare
        dalla stessa persona: ``ieri ti ho dato un feedback: com'è andata? Sei
        soddisfatto? Come poteva andare meglio?''. Con le persone con cui si ha
        meno a che fare, questo aiuta molto la volta successiva.
\end{itemize}

% ---------------------------------------------------------------------
\section{Come ricevere un feedback}
\label{sec:z1-ricevere}

Per il relatore è la parte più difficile. Si riassume in quattro passi:
\begin{enumerate}
  \item \textbf{Ascolta}: si parte sempre ascoltando.
  \item \textbf{Chiedi dettagli}, se serve, ma per capire meglio, non per
        difenderti.
  \item \textbf{Accetta, a prescindere}: non si sta decidendo chi ha ragione;
        si sta ascoltando un'opinione. Non c'è nulla da cui difendersi.
  \item \textbf{Agisci}: fatti un piano. Può essere banale (``mi hanno detto
        che interrompo sempre: cerco di non interrompere''), oppure ``devo
        migliorare questa cosa: come?''. Pensa anche a \emph{chi} potrebbe
        aiutarti e a \emph{che supporto} ti serve: il piano non deve essere
        per forza ``devo risolvere tutto da solo''.
\end{enumerate}
Una domanda che molti sottovalutano: \emph{``sono junior in un'azienda: sarò
mica il primo ad avere questo problema?''}. Probabilmente qualcun altro l'ha
già avuto: basta chiedere ``tu come hai fatto?''.

\subsection{Il feedback positivo}

Tutto questo vale per i feedback costruttivi negativi, a cui si pensa di
default; ma esistono anche i \textbf{feedback positivi}, e vanno usati. Quando
si riceve un ``grazie'' o un ``bravo'', la domanda è: \emph{come faccio a
ricreare le stesse condizioni?}

Il relatore ha lavorato anche in aziende dove il feedback era solo negativo.
Il rischio? Si immagini l'impiegato ideale: vuole sempre migliorare, è sempre
pronto a dire ``magari sbaglio io''. Se nessuno gli dice cosa sta facendo
\emph{bene}, spinto dal desiderio di migliorare rischia di cambiare cose che
erano già buone, e di farle peggio. Il rinforzo positivo \textbf{non è solo una
gentilezza}: è un atto pratico, che evita che le persone peggiorino ciò che
sanno già fare bene.

\subsection{Il feedback accidentale}

Esiste infine il \textbf{feedback accidentale}, molto frequente negli uffici:
\begin{itemize}
  \item il ``chiedi a lei, non a quel grafico'' del giro dell'ufficio è un
        feedback negativo su una persona, dato accidentalmente, né richiesto
        né voluto;
  \item a Riot li chiamavano \textbf{drive-by}, come nelle scene dei film in
        cui passa un'auto, si abbassa il finestrino, si spara e si riparte. In
        molte aziende con una ``cultura del feedback'' in cui tutti devono
        dire la loro, si riceve una valanga di commenti del tipo ``comunque il
        sito a me non piace, ciao''. Non si sa perché, non si sa cosa fare di
        diverso, non si sa che rischi ci siano: \emph{a che cosa è servita
        questa conversazione?} A niente.
\end{itemize}

% ---------------------------------------------------------------------
\section{Quando, quanto e con chi}
\label{sec:z1-quando}

\subsection{Frequentemente e recentemente}

\emph{Quando avete ricevuto il vostro ultimo feedback?} Per gli studenti,
qualche giorno fa (nel corso) o qualche mese fa (al lavoro, dopo una
consegna). Come aveva osservato Silvia, il feedback andrebbe dato
spesso.

Il relatore usa l'esempio di una nave da crociera (``la Costa Concordia'') che
deve andare da A a B. La rotta ideale (azzurra) è quasi impossibile: nessun
progetto va esattamente come previsto. C'è la rotta ``panoramica'' (verde), che
gira intorno all'isola d'Elba; ci sono poi rotte (arancione e rosa) che a un
certo punto si allontanano molto e poi correggono bruscamente: qualcuno si è
accorto che si stava sbagliando, ``ma non è quello che avevo chiesto'',
``non è dove dovevamo andare''. Sono i famosi \textbf{sanity check} o
\emph{checkpoint}.

L'ideale è avere \textbf{tanti piccoli check per fare tante micro-correzioni}:
correggere la rotta di un paio di gradi è molto più rapido che farlo dopo
essersi allontanati. Bastano trenta secondi di ``ehi, vado bene?''. Spesso non
lo si fa per paura di sembrare insicuri o poco autonomi, ma per il relatore è
l'opposto: quando qualcuno gli chiede ``vado bene così?'' e lui risponde
``sì'', si alleggerisce anche del proprio carico mentale. Un capo che assegna
un compito e lo vede solo alla fine vive sempre col dubbio (``avranno capito?
Lo staranno facendo giusto? Arriverà in tempo?''); con un check di dieci
secondi al giorno (``tutto ok, Stefano'') è tranquillo, e la persona va avanti
sapendo che quanto fatto finora va bene. Non serve fare ogni giorno il giro del
dipartimento: anche ogni uno o due giorni basta a evitare la rotta verde.

Per il feedback è lo stesso: \textbf{va dato subito}. Dire oggi ``ieri hai
fatto questa cosa'' è molto diverso da dire ``negli ultimi sei mesi hai fatto
questa cosa''. Il relatore si è sentito dire più volte, alla valutazione di
dicembre, ``da maggio hai fatto questa cosa che non andava''. La sua risposta:
``e aspetti dicembre per dirmelo? Mi hai lasciato sbagliare per sei mesi?''. È
una mancanza di rispetto, e la conversazione non può più essere costruttiva.
Per questo non ama le valutazioni annuali: i manager crescono abituandosi a
prendere appunti delle cose negative per tirarle fuori alla valutazione
(``a marzo hai fatto questa cosa'' ``sì, però poi ho smesso'' ``sì, però a
marzo l'hai fatta'').

In sintesi:
\begin{description}
  \item[Frequentemente] riguarda l'intervallo tra un check e l'altro.
  \item[Recentemente] significa non lasciar passare un anno prima di dare il
        feedback, altrimenti non serve più a niente.
\end{description}

\subsection{Quanto dura un feedback}

Un feedback dura \textbf{da trenta secondi a quindici minuti}; raramente di
più. Chi non ne ha mai dato uno dovrebbe sforzarsi di darne almeno uno, magari
positivo per cominciare. Ogni feedback dato o ricevuto è un
\textbf{deposito nel proprio ``fondo fiduciario''} (\emph{trust fund}), che
costruisce la reputazione futura. Il relatore si è costruito la reputazione di
chi dà feedback il più possibile oggettivi e analitici: giusti o sbagliati che
siano, costringono a pensare. Non impone mai il suo punto di vista (``si fa
come dico io''), ma chiede ``tu come faresti? Come posso aiutarti?''. Il
feedback, per lui, serve a \emph{far crescere qualcun altro}.

\subsection{I manager peggiori: ``a che numero sto pensando?''}

I peggiori manager che il relatore abbia avuto erano quelli che volevano che si
facesse come volevano loro, anche quando chiedevano ``tu come la faresti?''.
In realtà la domanda era ``\emph{a che numero sto pensando?}'': se dicevi 5,
``no, dovremmo fare 9''; se dicevi 9, ``no, 5''. Lo stress costante di dover
indovinare cosa vuole il capo nasce dalla mancanza di fiducia (``se mi chiedi
come fare, fammelo fare'').

Lo stesso gli succede con i CEO che lo chiamano per ``sistemare la cultura
aziendale'' e poi dicono ``non sono sicuro, io farei diversamente''. La sua
risposta: ``lo so che faresti diversamente, per questo mi hai chiamato''. Se si
chiama un esperto, bisogna lasciargli fare il suo lavoro: non si va dal
meccanico a cambiare le gomme per poi dirgli ``non dovresti girare un po' di
più quel bullone?''. Chiedere un parere a un esperto significa dargli anche il
potere di scegliere. Il problema è il \emph{micromanagement}, o in generale il
non accettare punti di vista diversi dal proprio (magari perché si è il
fondatore e si pensa di essere l'unico a saper fare le cose).

\subsection{Dare feedback a un superiore egocentrico}

\paragraph{Come si dà un feedback negativo a un superiore con un grande ego?}
Il trucco è \textbf{non puntare sull'ego, ma sull'azione e sui suoi effetti}.
Non ``tu fai così e io mi sento così'': anzi, si alimenta il suo ego
rendendolo il \emph{main character} dell'avventura, l'eroe che risolve il
problema (``tu, grande manager, aiutaci a risolvere questo problema che
affligge noi che facciamo l'operatività''). Non è l'ideale (``io cambierei
manager''), ma pochissimi possono scegliersi i colleghi: il relatore stesso ha
quasi sempre ricevuto team preassemblati (``ecco le chiavi, questa è la tua
macchina''). Bisogna giocare con le carte che si hanno in mano.

\paragraph{Un esempio pratico.} Con un manager che fa micromanagement e mette
sempre ``il becco'' nel proprio lavoro, si può dire: ``se ti fidi di me per
questo lavoro, con le tue capacità potresti dedicarti a questa cosa più
strategica e di lungo termine: pensa quanto sarebbe più utile per l'azienda''.
In sostanza: ``guarda cosa potresti raggiungere se mi lasciassi fare il mio
lavoro''. Non come un ``tu fai il tuo e io faccio il mio'', ma come un
``ti sollevo da questa parte del carico, che so che ritieni importante: lascia
che ci pensi io, così puoi dedicarti ad altro'', che magari gli dà anche più
visibilità.

Si gioca su ciò che quella persona pensa di saper fare, e che probabilmente sa
anche fare. Anche Michael Scott di \emph{The Office}, per quanto preso in giro,
è raccontato come un venditore bravissimo, che ha chiuso i contratti più
importanti per l'azienda, e rispetto ai manager che gli succedono è forse
``uno dei meno peggiori''. Tutti, anche i manager peggiori, hanno punti di
forza: \textbf{giocate su quelli}.

% ---------------------------------------------------------------------
\section{Il feedback come regalo}
\label{sec:z1-regalo}

Per chiudere, il relatore risponde alla domanda d'apertura, che aveva lasciato
apposta in sospeso: che cos'è il feedback per lui?

\begin{quote}
  \emph{Il feedback è un regalo che mi viene fatto da qualcuno che so che vuole
  il mio interesse.}
\end{quote}

Ha costruito un team di persone che si fidano di lui e che non hanno motivo di
sabotarlo o di ferirlo: per questo interpreta anche le cose peggiori con il
default to trust. È un regalo, e cerca di prenderlo in considerazione il più
possibile.

Ci è arrivato partendo da se stesso. Oltre a essere avverso ai conflitti, si
definisce ``avverso alla comunicazione umana'': paradossalmente, una persona
introversa e timida finita a fare comunicazione, marketing e ora HR (``succede
spesso''). Prima di ordinare una pizza al telefono ha dieci minuti di
preparazione mentale, e se l'ordine delle domande cambia va in confusione. Per
lui un feedback di dieci minuti sono spesso \textbf{tre ore di preparazione}:
costruisce un \textbf{albero dei dialoghi}, come in un'avventura grafica
personale, con tutte le opzioni. Poi spesso finisce con ``hai fatto questa
cosa, andrebbe migliorata'' ``hai ragione, grazie'', e basta.

Così ha pensato: chissà quanti si fanno gli stessi problemi con lui, e passano
ore a prepararsi per dirgli ``questa cosa l'hai fatta male'', sapendo che
riceveranno comunque quell'occhiata fugace del ``come ti permetti''. Eppure
vengono a dirglielo. Non può far altro che rispondere: ``grazie, hai ragione,
cercherò di fare meglio''. È così che ha metabolizzato il ricevere feedback:
capendo che non c'è nessuna cattiva intenzione e che, se non proprio felice
(``ci ho messo una pietra sopra''), può almeno essere \emph{grato} che ci sia
ancora qualcuno che ha così a cuore il suo interesse da dargli feedback.

\paragraph{Il silenzio fa più paura.} Non bisogna spaventarsi dei feedback
negativi che emergono dai sondaggi aziendali: bisogna spaventarsi del
\textbf{silenzio}, quando nessuno dice più nulla. Significa che nessuno si
fida più che si sappiano risolvere i problemi.

\paragraph{La sindrome dell'impostore.} Il suo motivatore personale, all'inizio,
era soprattutto ``non voglio fallire'': per questo cercava di migliorare, e
migliorare è diventato accettare feedback, perché più feedback riceve, più
migliora. Il feedback è un \textbf{acceleratore}, e lui cerca di incentivare
gli altri a darglielo. Resta però sempre la vocina del ``e se si accorgono che
non sei capace?''. Con la \emph{impostor syndrome} convive da anni e
continuerà a conviverci; se ce l'avete, conviveteci anche voi. Avrete sempre
la sensazione che ``se ne accorgeranno'': forse sì, ma sapete fare anche
altre cose. Non concentratevi solo su ciò che non sapete fare. Gli errori
capitano: la perfezione non è di questo mondo.

\paragraph{Esercitarsi in accademia.} Per fare pratica con il feedback,
l'ambiente migliore è quello dei colleghi di corso: persone nella stessa
situazione, con esperienza simile e lo stesso obiettivo, in un contesto in cui
non si può essere licenziati e l'unico effetto a lungo termine è ``ho parlato
con questa persona, ora le cose vanno meglio''. E bisogna avere il coraggio di
parlare anche con i docenti (``dire a Diego: secondo me questa lezione potevi
farla in questo modo''; ``so che è una persona che lo accetta''), senza farsi
bloccare dalla gerarchia. Nell'azienda giusta (come quella in cui lavora ora) il
relatore va sempre dal CEO a dare feedback, sapendo che non verrà licenziato:
il CEO gli ha detto ``ti ho preso per darmi feedback''.

La frase di chiusura:

\begin{quote}
  \emph{Per superare le vostre difficoltà, tutto quello che vi viene chiesto è
  il coraggio di chiedere aiuto.}
\end{quote}

Chiedere aiuto non è una debolezza: è un segno di consapevolezza di sé
(\emph{self-awareness}) e di forza personale. Se tutti lo facessero, i
progetti andrebbero molto meglio; spesso invece ci si fa frenare dal ``non
voglio fare brutta figura'', ``non voglio passare per quello scemo'', e sono
proprio queste paure a rallentarci e a impedirci di raggiungere gli obiettivi.

% ---------------------------------------------------------------------
\section{Domande e risposte finali}
\label{sec:z1-qa-finale}

L'ultima mezz'ora è dedicata a domande libere.

\paragraph{E se una persona peggiora per uno o due anni di fila?} Se nessuno
gliel'ha mai detto, c'è un grosso problema, non suo ma dell'azienda. Se invece
c'è consapevolezza, il percorso è: ``cosa vuoi fare per risolvere?''; ``ho
provato questo, questo e questo''; ``visto che non ha funzionato, vuoi un mio
consiglio? Proviamo con quello?''; ``che supporto ti serve?''. Se, con tutto il
supporto, il tempo e le soluzioni concordate per mesi, il miglioramento non
arriva, per il relatore la soluzione è: \textbf{``dimmi un prezzo per andare
via''}. Preferisce pagare una buonuscita (\emph{severance package}) che
avviare la strada dei richiami formali sulla performance fino al
licenziamento per giusta causa. A un certo punto andare avanti non è un favore
per nessuno.

\paragraph{È vero che nei videogiochi si viene sfruttati? Meglio un piccolo
studio italiano?} Per il relatore si è quasi più a rischio a Blizzard che in un
piccolo studio di Roma o Torino. Nel piccolo studio si sa dove ci si va a
cacciare: non ci sarà il team building a Santo Domingo, ma ci si concentra su
un gioco alla volta e, se va bene, si fa il successivo. A Blizzard si può
lavorare su \emph{Overwatch}, che vende milioni, e perdere comunque il posto,
in modo totalmente scollegato dal proprio lavoro. Gli stipendi più alti
all'estero compensano anche questo rischio, che spesso non dipende da come va
il gioco (anzi, ``il gioco non ha fatto come volevamo'' è spesso una scusa per
licenziamenti che non avevano motivo).

\paragraph{Cambiare settore a trent'anni, da web developer full remote?} Lo
stipendio indicato dallo studente non è fuori dalla portata di vari studi
italiani, anzi è superabile. Non va esclusa la possibilità di lavorare da
remoto per aziende piccole ma estere. Alcuni amici del relatore, pagati con
stipendi italiani, si sono trasferiti a Bali o in paesi con un costo della vita
più basso; il collega di \emph{PUBG} percepisce uno stipendio olandese vivendo
a Milano, che per quanto cara non costa come i Paesi Bassi.

\paragraph{Come si diventa manager?} Dipende dall'azienda: alcune promuovono i
migliori (processo tendenzialmente sbagliato), altre propongono un percorso di
\emph{people management}. Diventare manager significa \textbf{abbandonare ciò
che si sa fare per imparare a fare altro}. Il relatore ha dovuto accettare che
non sarebbe più stato ``il miglior community manager italiano del mondo''
(``eravamo in pochi, quindi era facile''): dedicandosi alle persone non poteva
più fare il suo lavoro bene come prima. È stata una negoziazione con se
stesso: fare il manager significa non essere più l'esperto di qualcosa, ma
saper gestire e far crescere persone. Significa anche \textbf{cedere il
controllo della propria carriera}: se il team decidesse di non fare più nulla,
il manager ne pagherebbe le conseguenze. Il suo lavoro diventa scegliere le
persone giuste e farle lavorare bene. Lui lo esplicita sempre: sa di affidare
al team parte della propria carriera, e il team sa che lui combatterà per
loro, anche prendendosi le colpe.

\paragraph{L'inglese.} \textbf{Va assolutamente imparato}, anche lavorando in
Italia. Non saperlo preclude circa il 95\% del mercato (tutto quello fuori
dall'Italia) e soprattutto il \emph{networking} con persone che potrebbero dare
occasioni future: l'inglese è spesso l'unica lingua in comune con chi lavora
all'estero ed è oggi la lingua universale dei videogiochi.
\begin{itemize}
  \item \emph{Che livello serve?} Non importa B2 o C2: bisogna saper sostenere
        una conversazione, guardare almeno Netflix con sottotitoli in inglese,
        comunicare senza dodici pause e senza ``come si dice\dots''. Il
        certificato non verrà mai chiesto: verrà chiesto se si sa parlare
        inglese. Il relatore ha sentito persone con un C2 che non sapevano
        parlarlo.
  \item \emph{Come allenarlo senza occasioni di conversazione?} Netflix, Prime
        Video, Now TV, Udemy, Coursera, YouTube, Duolingo, siti di scambio
        linguistico con madrelingua, Discord: non sono mai esistiti così tanti
        mezzi gratuiti. Il relatore ha imparato l'inglese giocando (con una
        pronuncia sbagliata, perché leggeva senza sentire, su cui ha poi
        dovuto lavorare); guardare un film in inglese con i sottotitoli in
        italiano insegna anche la traduzione.
  \item Oltre all'ascolto passivo è fondamentale \textbf{parlare}: in una
        flotta di \emph{EVE Online} con 150 persone quasi tutte madrelingua
        inglese si è esposti a decine di accenti e modi di dire (mille modi di
        dire ``come stai'' oltre a \emph{how are you}); lavorando all'estero si
        impara a distinguere gli accenti (Manchester, Londra) e ad apprezzarli
        anche nei media.
\end{itemize}

\paragraph{Quanto deve sapere un junior?} Dipende da quanto è tecnico il
ruolo. Un game designer junior che non sa qualcosa può essere affiancato da un
senior; un programmatore che non sa programmare, in quel momento, non è utile
all'azienda.

\paragraph{Annunci ``junior con cinque anni di esperienza''.} ``Cercasi
ventenne con trent'anni di esperienza.'' Non è l'annuncio per un junior: è un
errore dell'azienda, che in realtà cerca un \emph{mid}. Riceverà profili
junior, si stupirà che non abbiano cinque anni di esperienza, rifarà l'annuncio
e avrà perso sei mesi. Il consiglio: tenere d'occhio quella pagina, perché
dopo qualche mese l'annuncio verrà corretto.

\paragraph{Una struttura per il post-mortem.} A chi chiede un formato di base,
il relatore anticipa quanto vedrà nella seconda giornata a proposito di
brainstorming: una lavagna \textbf{divisa in quattro aree}, su cui ognuno
attacca tutti i post-it che vuole:
\begin{center}
  \begin{tabular}{@{}p{0.44\textwidth}p{0.44\textwidth}@{}}
    \toprule
    \textbf{Smettere di fare} & \textbf{Iniziare a fare} \\
    ciò che non ha funzionato & ciò che non abbiamo fatto e andava fatto \\
    \midrule
    \textbf{Fare di meno} & \textbf{Fare di più} \\
    ciò che ha fatto perdere tempo rispetto all'impatto ottenuto &
    ciò che ha funzionato ma, fatto con troppa leggerezza, non ha dato i
    frutti sperati \\
    \bottomrule
  \end{tabular}
\end{center}
Con i post-it ognuno ha voce e tutto risulta catalogato; per il report basta
una foto della lavagna (o uno screenshot, se la lavagna è digitale, come Miro).

\paragraph{Come sbloccare la paralisi del feedback verso il management?} Uno
studente ha visto manager dire ``non abbiate paura, la mia porta è sempre
aperta'' senza che cambiasse nulla. C'è una dissonanza tra ciò che dicono e
ciò che fanno: è ``la mamma che dice vieni qua che non ti faccio niente''.
Qualcosa deve aver fatto pensare alle persone che non è sicuro. Prima va
sistemato quello; poi si può provare con i \textbf{feedback anonimi}.
L'obiettivo del relatore resta il feedback aperto e firmato, ma a Stormind ha
iniziato con feedback anonimi verso il management per verificare se
l'anonimato fosse un fattore. Non lo era (evidentemente non c'erano problemi a
parlare con il CEO), e la volta successiva intende abbandonarli. Ma è un buon
\emph{check}: se senza anonimato non arriva nulla e con l'anonimato arrivano
cento feedback, il problema è chiaro.

\paragraph{Quando è giusto rivolgersi a un superiore?} La strada giusta resta
parlare direttamente con l'interessato. Ma se si prova e non si trova una
soluzione, la cosa più ovvia è andare dal proprio manager insieme: ``ci
aiuti?''. Quando il relatore parlava di mediare i dialoghi, era proprio questo:
fare da moderatore, aiutare a sbloccare i punti fermi o, se non si riesce a
scegliere tra A e B, scegliere lui. L'escalation, quando necessaria, va fatta.

\paragraph{Il tempo libero veniva rispettato?} No, e se l'è dovuto imporre
lui. Il lavoro gli piaceva così tanto che continuare a casa non gli sembrava
lavorare. Gli si è ritorto contro con il lancio di \emph{Diablo~III}, fatto in
orario americano (le due o l'una di notte in Europa), con l'ufficio chiuso
perché a nessuno era stato chiesto di restare. Il lancio va malissimo, i
giocatori protestano, e lui, sveglio, risponde sui forum fino alle tre, quattro,
cinque del mattino: ``non posso risolvere il problema, ma resto qui ad
aggiornarvi con quello che riesco a sapere dal team di sviluppo''. Il giorno
dopo trova sulla scrivania un \textbf{richiamo formale dell'HR}: aver lavorato
fuori orario senza autorizzazione, condotta non professionale. Per lui è stato
un interruttore: ``ora me ne vado''. Due settimane dopo faceva colloqui con
Riot. Ancora oggi è contrario a quel richiamo, ma gli ha fatto capire come
ragiona un'azienda: ``hai un contratto 9--18, non mi interessa se l'hai fatto
con le migliori intenzioni''. Eppure per anni ha goduto di un'ottima
reputazione in quelle community come ``quello che è rimasto con noi fino alle
cinque del mattino'', e anche l'azienda ne ha beneficiato.

Ci sono poi aziende in cui, soprattutto per i developer, viene richiesto il
famigerato \textbf{crunch}: lavoro straordinario continuativo, senza pause.
Dipende molto dall'azienda; sta venendo (``grazie al cielo'') deprecato, perché
molte aziende capiscono che non è sano ed è visto male anche dal pubblico, ma
esiste ancora.

\paragraph{E se il tuo capo demotiva qualcuno del tuo team?} Qui ``diventa
personale'': il team del relatore si fida di lui perché sa che in questi casi
``diventa una bestia''. ``Te la puoi prendere con me quanto vuoi, ma il mio
team non lo tocchi.'' Lo gestisce direttamente con l'altro manager, per quanto
brutta possa diventare la conversazione. Lo sforzo più grande è \textbf{non
andarne poi a parlare con il team} (``ragazzi, ho fatto questa cosa''), perché
minerebbe la fiducia verso l'alto: non vuole un rapporto ottimo fino a lui e
pessimo oltre. L'obiettivo è che l'altro manager ``si faccia perdonare''
cambiando il proprio comportamento, e che il team percepisca da solo il
cambiamento.

\paragraph{Chiusura della giornata.} Il relatore chiede un feedback collettivo
in chat (``se domani ci dai \emph{more of the same} siamo contenti, oppure
ditemi cosa limare''), si scusa per i ripetuti crash della connessione e
ringrazia per la proattività nelle domande, che gli servono anche come
\emph{check}: se nessuno chiede nulla su un argomento, significa che può
andare avanti. Racconta di aver passato una settimana a decidere che cosa
mettere nelle slide e di averle poi preparate in poco tempo (rifacendone
alcune fino alle tre di notte). Ricorda infine di usare il test di
personalità per sé e per il proprio gruppo, facendolo fare a chi non l'ha
ancora fatto: tornerà molto utile per cooperare meglio.

% ---------------------------------------------------------------------
\section{Anticipazioni sulla seconda giornata}
\label{sec:z1-domani}

La seconda giornata del workshop sarà più pratica: ``oggi vi ho spiegato
perché, domani vi insegno come lo fanno queste aziende''. Gli argomenti
previsti:
\begin{description}
  \item[Time management] come organizzarsi una giornata professionale in
        azienda: come strutturare una \emph{task list}, cosa fare e cosa non
        fare, come e cosa prioritizzare.
  \item[Brainstorming] i gruppi dovranno fare moltissimi meeting (come facciamo
        il gioco, come lo chiamiamo, che scopo gli diamo). Se non si sanno fare
        bene, i brainstorming sono \textbf{dannosi ed eterni}: un brainstorming
        dovrebbe durare \textbf{al massimo due ore}. Verranno mostrati metodi
        che accorciano la durata, massimizzano l'effetto e, soprattutto,
        \textbf{fanno partecipare tutti}: non c'è niente di peggio di un
        brainstorming in cui parlano in due per il 90\% del tempo. Il classico
        ``ragazzi, non esistono idee sbagliate, sparate'' è il modo peggiore di
        farlo; non è colpa di nessuno, ma si può fare diversamente.
  \item[Guidare da dentro] come essere leader nel proprio team anche senza
        esserne il leader: sviluppare responsabilità personale e verso gli
        altri. Non essere il manager di qualcuno non significa non esserne
        responsabili: la persona accanto è responsabilità di tutti, perché si
        è un team e non un insieme di individui. ``Se qualcuno cade, tutti lo
        facciamo rialzare: non c'è l'addetto al rialzamento.'' È la versione
        \emph{advanced} della giornata: come comportarsi da manager senza
        averne il titolo.
\end{description}
Torneranno inoltre il \emph{service level agreement} e il tema della
condivisione dei successi e della protezione del team dai fallimenti.

% ---------------------------------------------------------------------
\section*{Punti chiave del capitolo}
\addcontentsline{toc}{section}{Punti chiave del capitolo}

\begin{itemize}
  \item ``La capacità di un gruppo di persone di fare cose straordinarie
        dipende dal modo in cui queste persone operano insieme.'' Il 90\% dei
        problemi dei team sono \textbf{problemi di comunicazione}.
  \item Le aziende si liberano sempre più presto dello \textbf{useless
        genius}: le competenze si trasmettono, il carattere no. Il
        \textbf{fit culturale} va valutato anche dal candidato.
  \item Non esiste un unico modo di dare feedback: dipende dal contesto e
        soprattutto da \textbf{chi lo riceve}. Il feedback è una somma di
        piccoli miglioramenti periodici; tenersi le cose nel cassetto non
        serve mai.
  \item I team attraversano \textbf{forming, storming, norming, performing e
        adjourning}. Lo storming non si salta; si passa la maggior parte del
        tempo in forming; la bravura sta nell'accelerare i passaggi e nel
        fare bene il \textbf{post-mortem}, anche dei successi.
  \item Gli obiettivi di un progetto vanno fissati \textbf{all'inizio}: solo
        così si può valutare se è andato bene e scegliere tra A e B.
  \item Il circolo virtuoso è \textbf{feedback $\rightarrow$ miglioramento
        $\rightarrow$ fiducia}. Accettare un feedback non significa fare come
        dice l'altro: decide chi ha la responsabilità, ma deve argomentare.
  \item ``Aiutami a poterti aiutare'': meglio dire in anticipo ``potrei non
        farcela'' che ``non ce l'ho fatta''. Chiedere aiuto non è una
        debolezza.
  \item Quando si dà feedback \textbf{non si porta la soluzione}: la si
        costruisce insieme a chi lo riceve.
  \item La cultura è forte quando si lavora \textbf{per gli altri}; il lavoro
        è valutato come opera d'insieme (anche con bonus di team).
  \item Come in \emph{Among Us}, i team sprecano energie a cercare il
        traditore invece di fare le task. \textbf{Default to trust}: chi ha
        sbagliato ha agito con le informazioni che aveva. Cercare
        \textbf{la causa, non il colpevole}: ``che cosa ti ha fatto pensare
        che fosse la scelta giusta?''.
  \item Carriere e performance non sono lineari: vanno valutate le
        \textbf{traiettorie}, non un punto dell'anno. Le valutazioni puntuali
        generano \emph{lying, hiding, faking}.
  \item Confermare sempre le istruzioni facendosele ripetere; ``non ho
        capito'' non è una colpa.
  \item Dare feedback: sull'\textbf{azione} e sui suoi \textbf{effetti}, specifico,
        presto, uno a uno, adattato al carattere, senza terze parti, senza
        ``non mi piace'', senza offendersi se non viene seguito; poi fare il
        debriefing.
  \item Ricevere feedback: \textbf{ascolta, chiedi dettagli, accetta, agisci}.
        Il feedback positivo serve a non peggiorare ciò che si fa bene; il
        feedback accidentale e il \emph{drive-by} non servono a nulla.
  \item Feedback \textbf{frequenti e recenti}: tanti piccoli \emph{sanity
        check} costano meno di una grande correzione di rotta.
  \item Il feedback è \textbf{un regalo} di chi ha a cuore il nostro
        interesse; fa paura il silenzio, non le critiche.
\end{itemize}

%    % =====================================================================
%  Lezione zero 2 — Gestire il tempo, generare idee, guidare da dentro
%  Fonte: trascrizione della seconda "lezione zero" (GD0_2), seconda e
%  ultima giornata del workshop di due giorni tenuto prima dell'inizio del
%  corso di Game Design, dedicato alla figura professionale e al lavoro in
%  team. Prosegue direttamente il capitolo 0.1 (cultura del feedback).
%  Pacchetti richiesti nel preambolo del libro:
%    \usepackage[italian]{babel}, \usepackage[T1]{fontenc},
%    \usepackage[utf8]{inputenc} (solo pdfLaTeX), \usepackage{booktabs}
%
%  Collocazione nel libro: parte "Lezioni zero", subito dopo
%  GD0_1_cultura_del_feedback, con la stessa numerazione propria (0.2).
%  Nel file principale del libro, ad esempio:
%
%    \part{Lezioni zero: il game designer nell'industria}
%    \renewcommand{\thechapter}{0.\arabic{chapter}}
%    \setcounter{chapter}{0}
%    % =====================================================================
%  Lezione zero 1 — Lavorare in squadra: la cultura del feedback
%  Fonte: trascrizione della prima "lezione zero" (GD0), prima giornata di
%  un workshop di due giorni tenuto prima dell'inizio del corso di Game
%  Design, dedicato alla figura professionale e al lavoro in team.
%  Pacchetti richiesti nel preambolo del libro:
%    \usepackage[italian]{babel}, \usepackage[T1]{fontenc},
%    \usepackage[utf8]{inputenc} (solo pdfLaTeX), \usepackage{booktabs}
%
%  Collocazione consigliata nel libro: le lezioni zero precedono il corso,
%  quindi vanno in una parte separata PRIMA di GD1, con numerazione propria
%  (0.1, 0.2, ...) così che i capitoli del corso conservino i numeri 1, 2, 3.
%  Nel file principale del libro, ad esempio:
%
%    \part{Lezioni zero: il game designer nell'industria}
%    \renewcommand{\thechapter}{0.\arabic{chapter}}
%    \setcounter{chapter}{0}
%    \input{Seminari/GD0_1_cultura_del_feedback}
%    % ... altre lezioni zero ...
%    \part{Il corso di Game Design}
%    \renewcommand{\thechapter}{\arabic{chapter}}
%    \setcounter{chapter}{0}
%    \input{GD1/GD1_introduzione_al_game_design}
%    % ...
%
%  (Con hyperref aggiungere anche \renewcommand{\theHchapter}{z.\arabic{chapter}}
%  nella parte zero, per evitare ancore duplicate.)
%  Tutte le etichette di questo capitolo hanno il prefisso "z1:" per non
%  collidere con quelle dei capitoli del corso.
% =====================================================================

\chapter[La cultura del feedback]{Lavorare in squadra: la cultura del feedback}
\label{cap:z1-feedback}

Questo capitolo apre le \emph{lezioni zero}: incontri che precedono l'inizio
vero e proprio del corso di Game Design e che non riguardano una tecnica di
design in senso stretto, ma \textbf{la figura professionale} di chi lavora
nell'industria dei videogiochi. Prima di imparare a progettare un gioco
conviene infatti capire in che tipo di ambiente lo si progetterà: un ambiente
fatto di team, di gerarchie, di scadenze e, soprattutto, di persone che devono
riuscire a lavorare insieme.

La lezione è la prima giornata di un \textbf{workshop di due giorni} (circa
dieci ore in totale, pause pranzo comprese) tenuto da Stefano
Cozzi\footnote{Il cognome è riportato come compare nella trascrizione
automatica, che in più punti storpia i nomi propri; andrebbe verificato.
Nella stessa trascrizione lo studio per cui lavora appare come
``Stormine'', ``Stormline'' e ``Stormind'': si tratta di Stormind Games,
studio con sede a Catania, coerentemente con quanto detto in lezione.}, un
professionista con circa quindici anni di esperienza nel settore, tra
Blizzard Entertainment, Riot Games e, oggi, Stormind Games. Il workshop è
nato da un confronto con i coordinatori dell'accademia\footnote{Nella
trascrizione l'accademia è resa in modi diversi (``Hive'', ``a dive'',
``D.I.V.''); i coordinatori citati sono Enrico e Andrea, insieme al
direttore Diego Ricchiuti. Nel resto del capitolo si usa semplicemente
``l'accademia'', come negli altri capitoli.}: serviva un momento dedicato a
\emph{come ci si comporta in un ambiente professionale}, cosa che riguarda sia
il lavoro futuro sia i progetti di gruppo svolti durante il percorso di
studi, dove ci si deve già abituare a comportarsi ``come se si fosse in
un'azienda''.

La prima giornata ha un unico grande tema, sezionato in più parti:
\textbf{la cultura del feedback}. La seconda giornata (che sarà oggetto della
prossima lezione zero) sarà più pratica: \emph{time management},
\emph{brainstorming} e il ``guidare da dentro'' un team anche senza esserne
il leader (sezione~\ref{sec:z1-domani}).

Il percorso del capitolo segue quello della lezione:
\begin{enumerate}
  \item il formato del workshop e il percorso professionale del relatore
        (sezione~\ref{sec:z1-relatore});
  \item perché le aziende tengono tanto al lavoro di squadra e al
        \emph{fit} culturale (sezione~\ref{sec:z1-squadra});
  \item che cos'è il feedback e da che cosa dipende il modo di darlo
        (sezione~\ref{sec:z1-cose-feedback});
  \item le cinque fasi di vita di un team e il valore del
        \emph{post-mortem} (sezioni~\ref{sec:z1-fasi-team}
        e~\ref{sec:z1-qa-fasi});
  \item il confronto con i colleghi e alcune domande sul lavoro nei grandi
        studi (sezioni~\ref{sec:z1-rivalita} e~\ref{sec:z1-industria});
  \item il circolo virtuoso tra feedback, miglioramento e fiducia, e i suoi
        effetti sulla \emph{retention} e sulla cultura aziendale
        (sezioni~\ref{sec:z1-triangolo} e~\ref{sec:z1-cultura});
  \item il ``gioco del traditore'' e il principio del
        \textbf{default to trust} (sezioni~\ref{sec:z1-among-us}
        e~\ref{sec:z1-default-trust});
  \item le traiettorie di carriera e la valutazione delle performance
        (sezione~\ref{sec:z1-traiettorie});
  \item come si dà, come si riceve e quando si dà un feedback
        (sezioni~\ref{sec:z1-dare}, \ref{sec:z1-ricevere}
        e~\ref{sec:z1-quando});
  \item il feedback come regalo, le domande finali e le anticipazioni
        della seconda giornata (sezioni~\ref{sec:z1-regalo},
        \ref{sec:z1-qa-finale} e~\ref{sec:z1-domani}).
\end{enumerate}

Una premessa di metodo: il relatore ha alternato continuamente teoria,
domande alla classe e racconti della propria esperienza. I racconti non sono
digressioni: sono gli esempi concreti con cui i principi vengono resi
tangibili, e per questo sono riportati con cura.

% ---------------------------------------------------------------------
\section{Il workshop e il relatore}
\label{sec:z1-relatore}

\subsection{Formato e regole di partecipazione}

Il relatore ha cercato di ``riassumere quindici anni di esperienza in dieci
ore''. Ha materiale per circa il doppio delle ore disponibili: le cose
strettamente necessarie vengono trattate per prime, il resto è tenuto ``in
riserva'' come \emph{nice to have}. Le regole del workshop:
\begin{itemize}
  \item le domande si fanno \textbf{alla fine di ogni slide}, così da chiudere
        prima il discorso; non c'è un'unica sessione finale di domande e
        risposte (anche se ci sarà), perché il relatore preferisce fermarsi
        spesso e non fare ``uno spiegone di dieci ore'';
  \item le domande non devono per forza riguardare la slide appena vista: se
        una slide dà lo spunto per altro, se ne può parlare; al massimo sarà
        il relatore a rispondere in breve o a rimandare a una slide
        successiva;
  \item si può chiedere anche di esperienze passate, opinioni e sensazioni;
        se per NDA o per serietà professionale non può rispondere, lo dirà
        esplicitamente (``non lasciate che sia un \emph{blocker}'');
  \item su Zoom si usa l'alzata di mano o la chat, che il relatore tiene
        aperta sul secondo monitor;
  \item non ci sono pause ufficiali: la lezione è registrata e ognuno è
        libero di assentarsi quando serve, cercando però di non disturbare il
        flusso della lezione;
  \item le slide verranno fornite integralmente tramite la segreteria
        dell'accademia (Vanessa), insieme alle registrazioni: le slide senza
        il parlato servono poco, ma sono un buon promemoria.
\end{itemize}

\subsection{Il percorso professionale}

Il relatore si presenta (scherzosamente) come la ``guida escursionistica''
che accompagna la classe nella ``foresta incantata'' della cultura del
feedback, perché per molti è un territorio sconosciuto. Il suo percorso:
\begin{itemize}
  \item in Italia ha lavorato per aziende che fornivano servizi a terzi, tra
        cui il lavoro su \emph{Warhammer Online}, e per un canale televisivo
        sul digitale terrestre dedicato ai videogiochi\footnote{Nella
        trascrizione: ``Digital Game Network Channel''. Il nome esatto del
        canale andrebbe verificato.};
  \item a \textbf{Blizzard Entertainment} (circa sei anni, a Parigi) ha fatto
        il \emph{community manager} per l'Italia (\emph{StarCraft~II},
        \emph{World of Warcraft} da \emph{Mists of Pandaria} in poi,
        \emph{Hearthstone}, \emph{Diablo~III}, le espansioni di
        \emph{StarCraft}, il lancio di Battle.net), per poi diventare il
        manager dei community manager europei;
  \item a \textbf{Riot Games} (circa cinque anni, a Dublino) è stato
        brevemente community manager italiano e poi \emph{publishing
        manager} di un team di circa trenta persone per l'Europa. Il
        \emph{publishing} non è solo community: è l'unione di
        \textbf{community, marketing e PR};
  \item a \textbf{PG Esports} si è occupato di eventi dal vivo e online
        (ad esempio la Milano Games Week e tornei di \emph{Valorant} e
        \emph{League of Legends}), restando in contatto con Riot e Blizzard;
  \item da circa un anno è tornato nello sviluppo, a \textbf{Stormind
        Games}, dove oggi si occupa di risorse umane e organizzazione del
        lavoro; lo studio verrà presentato più ampiamente al \emph{Career
        Day} dell'accademia insieme alla sua HR manager;
  \item fa inoltre consulenze aziendali anche al di fuori dei videogiochi e
        insegna in master universitari (IULM di Milano).
\end{itemize}

Il suo ruolo ricorrente, sia a Blizzard sia a Riot, è stato quello di
\emph{Mr.~Wolf} (``per chi ha visto \emph{Pulp Fiction}''): quello che
risolve i problemi. Veniva inserito nei team più in difficoltà, ad esempio
dal punto di vista delle performance, per sistemarli e passare al successivo.
Era presente in alcuni dei momenti più critici, come il lancio di
\emph{Diablo~III} o il lancio di espansioni di \emph{World of
Warcraft}\footnote{La trascrizione riporta ``il lancio con i Gary's'' e
``World of Trenor'': si tratta verosimilmente di \emph{Mists of Pandaria} e
di \emph{Warlords of Draenor}.}, a fare \emph{damage control}.

La scoperta fondamentale di questi anni è che i problemi da risolvere
\textbf{quasi mai riguardavano il prodotto} (non era certo lui a sistemare i
server di \emph{Diablo~III}), ma i team che c'erano dietro: persone che
litigavano, non collaboravano, non riuscivano a stare nella stessa stanza
senza urlarsi contro. E nel 90\% dei casi tutto questo si riduce a
\textbf{problemi di comunicazione}: è raro che perfino le incompatibilità
caratteriali non siano risolvibili parlando.

\subsection{Struttura non vuol dire burocrazia}

A Stormind il relatore sta accompagnando un team passato da 30 a 95 persone
in meno di un anno, attraverso tutti i cambiamenti di cui parla il workshop.
Darsi una struttura viene spesso percepito male, come ``più processi, più
procedure'', ma bisogna distinguere:
\begin{itemize}
  \item un processo come ``timbrare il cartellino'' può essere davvero
        invasivo e noioso;
  \item definire \emph{come funziona un documento di game design, chi lo
        approva, chi lo visiona e chi no} è invece una struttura necessaria.
\end{itemize}
Uno studio di 30 persone nella stessa sede può farne a meno: quando c'è un
dubbio ci si alza e si va alla scrivania del collega. Uno studio di 95
persone sparse per l'Italia (e in parte per l'Europa) non può più farlo:
servono una struttura diversa, un modo di lavorare diverso e una
pianificazione del futuro diversa.

% ---------------------------------------------------------------------
\section{Perché il lavoro è sempre lavoro di squadra}
\label{sec:z1-squadra}

\subsection{Non si va in ufficio per fare amicizia}

Il primo messaggio è netto: \textbf{non si va in ufficio per fare amici} o
per trovare compagni di vita. Si va in ufficio perché c'è un lavoro da fare, e
lo scopo è riuscire a farlo \emph{anche con chi sta antipatico}. Con le
persone simpatiche è facile: diventeranno amici, e non sarà mai un problema.
È con chi sta antipatico che le cose si complicano, ed è lì che emergono i
\textbf{veri professionisti}: quelli che riescono a lavorare con chiunque
senza che l'antipatia pesi sul modo in cui lavorano e collaborano, mettendo il
prodotto e il team davanti al proprio ego e alle proprie preferenze personali.

La frase che fa da filo conduttore a tutto il workshop (e che il relatore
promette di riproporre ``fino alla nausea'') è:

\begin{quote}
  \emph{La capacità di un gruppo di persone di fare cose straordinarie
  dipende dal modo in cui queste persone operano insieme.}
\end{quote}

Le aziende raramente pensano che l'innesto di una singola persona possa
cambiare il destino di un progetto. Questo è tanto più vero quanto più grande
è il team, ma vale anche per i piccoli studi, anche per un gruppo di cinque
persone. Come negli sport di squadra: l'individuo conta fino a un certo punto,
e anche se si ha Dončić in squadra, se si gioca uno contro cinque si perde
comunque\footnote{Il riferimento è alle finali NBA in corso nei giorni del
workshop tra Dallas Mavericks e Boston Celtics (giugno 2024), il che
permette di datare approssimativamente la lezione.}.

\subsection{Lo \emph{useless genius} e il costo delle persone tossiche}

Molti dipartimenti HR usano la definizione di \textbf{useless
genius}\footnote{Il concetto è più noto in letteratura manageriale come
\emph{brilliant jerk}; in lezione è usata l'espressione \emph{useless
genius}.}: si può essere bravissimi, a volte i migliori nel proprio lavoro,
ma se non ci si sa comportare, relazionare e operare in un gruppo si è
fondamentalmente inutili ai fini del gruppo. Nel migliore dei casi inutili,
nel peggiore \emph{dannosi}.

In passato le aziende tendevano a ``chiudere un occhio'': la persona è
fondamentale, la teniamo. Oggi la domanda è diventata: \emph{quanto costa
all'azienda, oltre al suo stipendio, tenere una persona che crea un ambiente
tossico e peggiora chi le sta intorno?} Ci sono costi sommersi che le aziende
cominciano a vedere, e la tendenza è liberarsi di queste persone il prima
possibile. Il ragionamento è semplice:
\begin{itemize}
  \item le \textbf{competenze si possono trasmettere} o sostituire: si trova
        qualcuno che fa lo stesso lavoro magari nel doppio del tempo, oppure
        più persone che se lo dividono;
  \item il \textbf{carattere delle persone} invece cambia molto raramente.
\end{itemize}

\subsection{Il fit culturale}

Quando qualcuno \textbf{non è un \emph{fit} culturale} per un'azienda è
meglio separarsi presto. Ma non essere adatti a \emph{quell'}azienda non
significa non essere adatti a lavorare: è perfettamente legittimo che una
persona inadatta a un'azienda sia adattissima a un'altra.

Il relatore l'ha vissuto in prima persona. Dopo sei anni a Blizzard viene
contattato da Riot e affronta un processo di selezione lunghissimo (``quasi
undici colloqui''), al termine del quale viene scartato. Scrive allora una
mail che gli ha ``cambiato la traiettoria di carriera'': chiede
semplicemente perché non è stato ritenuto adatto. La risposta: in nessun
colloquio si era esposto abbastanza da far capire \emph{perché} volesse
lasciare Blizzard, e quindi pensavano che in realtà stesse bene dov'era. Il
motivo era che, per educazione professionale, non gli piace parlare male del
proprio posto di lavoro; portata all'eccesso, però, questa riservatezza era
diventata un ``è tutto bellissimo, ma vorrei provare Riot'', senza un perché.
Lo invitano comunque a un incontro dal vivo (``dal vivo è diverso''), e alla
fine viene assunto.

Una volta arrivato, i sei mesi di prova (\emph{probation}) sono serviti tutti
per smettere di ragionare e lavorare ``come a Blizzard''. La differenza
culturale era enorme, tanto che in Riot esisteva un gruppo di supporto
interno tra ex dipendenti Blizzard: un collega gliel'ha descritto come
``uscire da una relazione tossica e doversi riabituare a una relazione
normale''\footnote{Il relatore precisa che parla della Blizzard di allora,
che non esiste più: oggi è parte di Microsoft e nessuno dei suoi ex colleghi
vi lavora ancora, anche a causa dei licenziamenti di massa.}.
\begin{description}
  \item[Blizzard] era \emph{developer first}: il focus era chiaramente sugli
        sviluppatori, prima di tutti gli altri. Dopo sei anni c'erano ancora
        meeting a cui non poteva partecipare perché non abbastanza
        \emph{senior}, informazioni che non poteva sapere, piani che non
        poteva vedere.
  \item[Riot] era \emph{player first}: prima i giocatori, e i dipendenti come
        mezzo per trasferire il feedback dei giocatori dentro il gioco. Il
        secondo giorno lo portano in una stanza e gli mostrano il piano dei
        cinque anni successivi. Lui, sorpreso, pensa a un errore (``sono qui
        da oggi, non vorrei mettervi nei guai''); la risposta è: ``sei un
        dipendente Riot o no? Allora non ci sono livelli di accesso: devi
        sapere queste cose per fare meglio il tuo lavoro''.
\end{description}

Da questa esperienza ha imparato molto sulla \textbf{trasparenza}: condividere
le informazioni aiuta il lavoro. Non sempre, però: a volte non si condivide
qualcosa non per segretezza, ma perché non è ancora pronta per essere
presentata, esattamente come non si mostra un gioco quando è ancora solo
teoria (``te lo faccio vedere quando ho qualcosa di tangibile'').

La lezione per gli studenti è doppia: le aziende valuteranno il \emph{fit}
culturale sempre di più, ma capire se si è un \emph{fit} richiede anche
\textbf{lavoro da parte del candidato}. Bisogna avere la presenza di spirito
di dire: ``conosco i valori di questa azienda, non mando nemmeno il CV,
perché so che durerei un mese e mezzo''.

% ---------------------------------------------------------------------
\section{Che cos'è il feedback}
\label{sec:z1-cose-feedback}

\subsection{Le prime risposte della classe}

La prima domanda alla classe è: \emph{quando sentite la parola
``feedback'', a cosa pensate?} Le risposte raccolte, a voce e in chat:
\begin{itemize}
  \item un \textbf{riscontro} su ciò che si fa e su come lo si fa, sul lavoro
        o sul modo in cui ci si interfaccia con il gruppo; può essere
        lavorativo oppure comunicativo e comportamentale. Chi ha risposto lo
        vede come qualcosa che, se l'azienda non dà, va attivamente
        richiesto, perché è importante sapere come sta andando per
        migliorarsi;
  \item né positivo né negativo, ma \textbf{costruttivo}: in entrambi i casi
        serve a migliorarsi;
  \item uno \emph{scambio di opinioni}, un \emph{riscontro produttivo};
  \item un feedback ``di reazione'', non verbale: si fa una cosa e si
        osserva come cambia l'ambiente o il comportamento delle persone
        intorno (anche perché a voce la gente spesso si tiene le cose
        negative);
  \item per un designer, anche il \textbf{linguaggio del corpo}: osservando
        un giocatore giocare si può dedurre l'emozione che prova (frustrazione,
        incomprensione).
\end{itemize}

Sulla forma (scritto o orale?) le opinioni divergono: c'è chi lo preferisce
\emph{mille volte} scritto; chi pensa che dipenda dalle dimensioni
dell'azienda (orale nelle piccole, scritto in quelle con team di centinaia di
persone); chi pensa che dipenda dall'importanza del feedback (scritto se
tecnico o pieno di dati, da tenere sotto gli occhi; orale se è un consiglio
veloce).

Il relatore usa due esempi per mostrare che \emph{dipende}:
\begin{itemize}
  \item un feedback sulla programmazione ha senso darlo scritto, magari
        commentando direttamente il codice, più che ``andando a prendere un
        caffè'';
  \item ma se il feedback è ``ieri, durante il meeting, ho percepito male il
        fatto che mi interrompessi di continuo'', ce lo si aspetta
        \textbf{orale}.
\end{itemize}

\subsection{Da che cosa dipende il modo di dare feedback}

Lo \emph{spoiler} della giornata è che \textbf{non esiste una risposta univoca}
su come dare feedback. Il workshop non insegna ``il'' modo, ma fornisce
\emph{tutte le opzioni}, perché il modo giusto cambia in base a molti fattori:
\begin{itemize}
  \item a chi lo si dà: un collega, un superiore, il proprio capo, un proprio
        \emph{report} (un sottoposto);
  \item se è personale o professionale;
  \item se è la prima volta o la terza volta che si dà lo stesso feedback;
  \item se riguarda un progetto importante o secondario.
\end{itemize}

Manca però l'elemento fondamentale, che la classe individua subito:
\textbf{l'altro}. A parità di tutti gli altri fattori, il \emph{carattere
della persona che riceve il feedback} dovrebbe far cambiare il modo in cui lo
si dà. Per questo, prima della lezione, agli studenti era stato chiesto di
fare un test di personalità: non un test di psicologia avanzata, ma uno
strumento per rendersi conto di quanta varietà ci sia anche solo in una
classe di 53 persone (il che non vuol dire trovare 52 modi diversi di dare
feedback, ma lavorarci un po' sì).

Il relatore porta se stesso come esempio. Preferisce:
\begin{itemize}
  \item i feedback \textbf{positivi scritti}, perché è a disagio a ricevere
        complimenti in pubblico;
  \item i feedback sul proprio lavoro \textbf{diretti}: ``vai subito al
        punto e non girarmi intorno''.
\end{itemize}
È molto orgoglioso del suo lavoro e, quando qualcuno lo critica, la sua
primissima reazione istintiva è ``come ti permetti''. Negli anni ha imparato
a farla durare un secondo, ma chi gli dà feedback di persona dice di
vederla ancora negli occhi. Per mitigarla, dice esplicitamente a chi gli dà
feedback: \emph{``ignora il modo in cui reagisco nei primi cinque secondi,
vai avanti, poi mi passa''}. È considerato da tutti molto aperto al feedback
e non ha problemi a dire ``hai ragione, facciamo diversamente'': ma ci sono
reazioni istintive e caratteriali che è difficile controllare al cento per
cento.

\subsection{Con chi si lavora meglio}

Guardando le persone con cui si lavora abitualmente (anche nei progetti
dell'accademia), si noterà che si lavora meglio con due profili:
\begin{enumerate}
  \item chi è \textbf{esattamente come noi}, ma qui è un 50 e 50: alcuni
        caratteri vanno d'accordo con i propri simili, altri non vogliono
        nessuno con lo stesso carattere nel gruppo (e si potrebbe rientrare
        in questa categoria);
  \item chi è \textbf{affine}: ha qualcosa in comune con noi, ma non tutto,
        e proprio per questo può dare punti di vista differenti.
\end{enumerate}
Il lavoro della giornata è capire come valorizzare il secondo caso: ottenere
feedback da persone con caratteri, visioni e background diversi, per
completare una visione del proprio lavoro che è necessariamente limitata. È
come esplorare una mappa con la nebbia di guerra: un'altra persona, al di là
del muro, vede cose che noi non vediamo.

\subsection{Il feedback come somma di piccoli miglioramenti}

Per il relatore il feedback è questo: \textbf{applicare tanti piccoli
miglioramenti periodici}, che si sommano nel tempo. Migliorare dello 0,1\%
ogni giorno porta a fine anno a un miglioramento enorme, e questo vale per i
prodotti, per i giochi, per le persone e per i team.

\paragraph{L'esempio del ciclismo britannico.} La squadra di ciclismo
britannica, fino a una quindicina d'anni fa, non aveva vinto quasi nulla.
Cambiato il team manager, tutti si aspettavano grandi cambiamenti sportivi
(convocazioni, strategie di gara). Arrivarono invece \emph{minuscoli
cambiamenti non sportivi}: un altro chef in trasferta, materassi migliori nei
camper per riposare meglio, un sellino più comodo a fine tappa. Cambiamenti
marginali che da soli non avrebbero cambiato niente; eppure, dopo due o tre
anni di questo lavoro quotidiano, la squadra ha cominciato a vincere a
ripetizione, Olimpiadi comprese, cambiando il modo in cui il ciclismo è
visto in Gran Bretagna\footnote{È la nota strategia delle \emph{marginal
gains} associata a Dave Brailsford (British Cycling e Team Sky); il nome non
è citato in lezione.}. Certo, si è investito anche sulle persone giuste, ma
tutto è partito dai piccoli miglioramenti periodici.

Per gli studenti vale lo stesso: chiarire un conflitto con una persona con
cui si ha ancora poco a che fare, o una situazione che ha lasciato
``l'amaro in bocca'', sembra inutile oggi, ma darà i suoi frutti tra tre, sei,
nove mesi. Sono \emph{investimenti sul futuro}.

\paragraph{Peggiorare ogni giorno.} Esiste anche il percorso inverso:
lasciar correre, ``buttare sotto il tappeto'' le cose perché ``tanto è
inutile parlarne''. Queste cose tornano sempre fuori dopo otto mesi, come il
\emph{libro dei rancori} dei nani di Warhammer: ``ti ricordi quando, un anno
e mezzo fa, all'inizio del progetto, mi hai detto quella cosa? Mi sono
offeso''. A quel punto è quasi impossibile risolvere, perché nessuno ricorda
più con precisione cosa sia successo (``ma quale cena?'', ``forse era in
pizzeria'', ``no, eravamo da un'altra parte''). \textbf{Tenersi le cose nel
cassetto per tirarle fuori quando servono è sempre sconsigliato}: non servono
ed è impossibile usarle in modo costruttivo.

% ---------------------------------------------------------------------
\section{Le cinque fasi di vita di un team}
\label{sec:z1-fasi-team}

\subsection{Perché conoscerle}

Quasi tutti gli studenti lavorano già in gruppo (al primo anno su un gioco da
tavolo, negli anni successivi su un videogioco) e si saranno accorti che un
team attraversa delle fasi: il modo di lavorare insieme cambia nel tempo.
Spesso ci si fa prendere di sorpresa: ``eravamo così affiatati, cosa è
successo?''.

Un esempio da una consulenza fuori dai videogiochi: un'azienda farmaceutica
gli affida un team commerciale (rappresentanti che presentano medicinali a
ospedali, medici e privati). Il problema: ``prima andava tutto bene; ora che
siamo in diciotto dovremmo lavorare meno e meglio, e invece va peggio''. Va
peggio perché le persone nuove sono appena arrivate, mentre il management si
aspettava già di vedere i risultati dell'investimento.

Fino a pochi anni fa le fasi di un team venivano disegnate come una
\textbf{linea dritta verso l'alto}: le fasi erano le stesse, ma viste come un
miglioramento costante. Non si vedevano le parti \emph{fisiologiche} del
lavoro in team, cioè anche i peggioramenti. Accettarle è il primo passo per
accettare i fallimenti: nessun percorso personale o di team è lineare.

Il relatore ammette che il suo modo di lavorare è nato per il 70\% dal
proprio impegno a imparare, fare bene e rispettare le persone, e per il 30\%
dall'aver lavorato con ``incapaci colossali'' che gli hanno insegnato come
\emph{non} fare le cose. A Blizzard un manager gli fece capire che, pur non
sapendo ancora che tipo di \emph{people manager} sarebbe diventato, sapeva
benissimo quale tipo di manager \emph{non} voleva essere.

\subsection{Forming, storming, norming, performing, adjourning}

Le fasi sono cinque, un po' come le fasi del lutto. Molte aziende ne
considerano solo quattro (``e sono quelle che solitamente chiudono''); le
aziende al passo con i tempi ne considerano cinque\footnote{Il modello è
quello di Bruce Tuckman (1965), a cui la quinta fase, \emph{adjourning}, fu
aggiunta nel 1977. Il nome di Tuckman non è citato in lezione; la
trascrizione riporta la quinta fase come ``journing''.}.

\begin{description}
  \item[Forming] Si viene messi insieme e si comincia a lavorare. Non si sa
        nulla degli altri; spesso non si conoscono nemmeno i nomi, ed è
        normale.
  \item[Storming] Si cominciano a ``dare le testate come le capre in
        montagna per determinare il territorio'': questa cosa la faccio io,
        questa la vorrei fare così, chi decide? Si stabilisce chi fa cosa e
        come si divide il lavoro. Nel grafico è una \textbf{fase di
        peggioramento}: che il team litighi o discuta civilmente non cambia
        nulla, il team deve investire parte di tempo ed energie per risolvere
        queste questioni. È normale e \textbf{non è evitabile}.
  \item[Norming] Si comincia a capire come lavorano gli altri, cosa sanno
        fare e come far funzionare il gruppo; iniziano i miglioramenti. È una
        forma di \emph{accettazione} degli altri, che deriva dall'aver capito
        cosa devo fare io e cosa devono fare loro.
  \item[Performing] C'è uno sforzo \textbf{proattivo}: ``questa cosa
        toccherebbe a me, ma so che Alessio la fa molto meglio, quindi gli
        chiedo aiuto''. Si sfruttano i punti di forza degli altri, perché li
        si conosce e perché si sa che passare il testimone non avrà effetti
        negativi (il capo non dirà ``perché non l'hai fatta tu?'').
  \item[Adjourning] Il progetto finisce, si chiude o viene riorganizzato, il
        team viene sciolto. Quello che andrebbe fatto è analizzare
        \textbf{cosa ha funzionato, cosa no e cosa migliorare} per il
        progetto successivo.
\end{description}

\subsection{La fase che tutti saltano}

``Spoiler incredibile'': l'\emph{adjourning} è la fase che molte aziende
saltano, \textbf{perché c'è sempre il progetto dopo}.

A PG Esports, per la Milano Games Week (come per qualsiasi evento dal vivo a
cadenza annuale: Lucca Comics, Etna Comics, Comicon\dots), finita l'edizione
di un anno si cominciava dal giorno dopo a lavorare a quella successiva,
perché un anno sembra tanto ma non basta mai. Nei videogiochi è lo stesso:
finito un gioco, a parte le patch di mantenimento e del \emph{day one}, si
passa subito al successivo.

Il cambiamento introdotto dal relatore (in parte già praticato, ma da lui
strutturato) è stato: finita la Games Week, \textbf{tutti in ferie per una
settimana}; poi ci si ritrova per definire cosa ha funzionato e cosa no, con
un \textbf{post-mortem} (o \emph{debrief}); solo dopo si inizia a lavorare
all'edizione successiva. Senza questa fase manca la parte di
\emph{learning}, cioè di ciò che si è imparato dal progetto.

Il consiglio vale per gli studenti: quando si finisce il gioco da tavolo, e
poi il gioco digitale, bisogna trovare il tempo per il post-mortem,
altrimenti si butta via \textbf{il 90\% dell'utilità del progetto}. Una
struttura semplice per farlo è descritta nella sezione~\ref{sec:z1-qa-finale}.

% ---------------------------------------------------------------------
\section{Domande sulle fasi del team}
\label{sec:z1-qa-fasi}

\subsection{Lo storming che ritorna}

\paragraph{Capita che lo storming si ritagli spazio anche nelle altre fasi?}
Lo \emph{storming} è una fase precisa, che però molti manager, coordinatori e
team cercano di \textbf{saltare}: ``se io sono bravo come manager e loro sono
bravi come team, passeremo da forming a norming senza problemi''. È un errore,
perché senza una fase in cui si discute davvero non si sviluppa una
\emph{skill} fondamentale: \textbf{saper discutere}. Tutti conoscono persone
non abituate a farlo, per cui ogni minimo disaccordo diventa un affronto
personale (``avevamo deciso la pizza, perché andiamo al sushi? Perché mi fate
questo affronto?''): sono persone che non sono abituate ad avere conflitti
costruttivi.

Se si pensa di essere in \emph{norming} e continuano a emergere conflitti, il
punto è che \textbf{non si era davvero in norming}: lo storming non ``si
trascina'' nelle fasi successive, semplicemente non se n'è mai usciti.
Allo stesso modo, a chi chiede se ogni mini-fase di storming durante la vita
del gruppo sia frutto di uno storming fatto male o solo di carattere, la
risposta è già contenuta qui.

Un partecipante osserva che spesso, ancor prima di finire un progetto, si è
già messi su quello successivo: è il modo migliore per essere in costante
\emph{damage control}. Era l'esperienza del relatore a Blizzard: veniva
spostato fisicamente da un team all'altro, e finito il lancio di
\emph{Diablo~III} era già su \emph{World of Warcraft}, senza vivere il
``dopo''. Si è dovuto ritagliare da solo lo spazio per un debrief personale,
ed è lì che ha deciso che tipo di manager non voleva essere. Il consiglio:
\textbf{se non lo fa il team, fatelo voi}.

\subsection{Dove si passa la maggior parte del tempo}

\emph{In quale fase si passa più tempo nella propria carriera?} Gli studenti
rispondono \emph{storming} e \emph{norming}. Il relatore, in quindici anni,
calcola di aver passato in \emph{performing} circa \textbf{dieci mesi}. Non è
pessimismo: arrivare a performing deve essere sempre l'obiettivo, sia da
manager sia da membro del team; ma è difficile arrivarci, ed è normale che non
sempre ci si arrivi.

Lo storming può durare a lungo, ma con team molto esperti dura pochissimo,
perché c'è chi lo sa gestire. In realtà \textbf{la maggior parte del tempo si
passa in forming}, perché ogni equilibrio raggiunto vale per \emph{quelle}
persone, \emph{quel} progetto, \emph{quelle} tempistiche. Basta che cambi un
qualsiasi fattore (una deadline spostata, una persona nuova, qualcuno che si
dimette, un progetto cancellato) e si ricomincia, perché i bilanciamenti
precedenti non valgono più.

Esempio: un team di cinque persone affiatatissimo, in pieno performing. La
dirigenza, entusiasta, ``raddoppia le risorse'' e lo porta a dieci. O si
lascia tutto com'è (e i cinque nuovi non fanno niente), o bisogna ridiscutere
chi fa cosa e come i nuovi possono aiutare: che piaccia o no, si ricomincia
da forming.

Quindi la bravura di un manager, ma anche di un membro del team,
\textbf{non sta nel passare tanto tempo in performing, ma nell'accelerare il
passaggio da una fase all'altra}, e nel fare bene la fase di
\emph{adjourning}.

\subsection{Analizzare anche i successi}

C'è la tendenza a fare post-mortem solo dei progetti andati male: ``questa
cosa non ha funzionato, vediamo perché''. Va benissimo. Ma quando un progetto
va bene si dà per scontato che sia andato bene ``perché era buono''. Bisogna
invece chiedersi \emph{perché} è andato bene: era il prodotto? La
tempistica? Il modo in cui abbiamo fatto una cosa piuttosto che un'altra?
\textbf{Se non si capisce davvero il motivo del proprio successo, si
tenderanno a replicare le cose sbagliate.}

\paragraph{Blockbuster.} Quando arrivò la notizia che Netflix avrebbe
lanciato un servizio di streaming, Blockbuster (la catena di noleggio di
videocassette e DVD) aveva tutto il tempo per preparare un'offerta
concorrente. In una riunione del management raccontata in un libro (di cui il
relatore non ricorda il titolo), la direzione sembrava orientata a investire
nello streaming, finché un dirigente fece notare che il 12\% del fatturato
derivava dalle \emph{multe per riconsegna tardiva}, impossibili con lo
streaming, mentre secondo le stime lo streaming avrebbe prodotto il 10\% del
fatturato. ``Non possiamo rinunciare a questo 2\%'': niente streaming. Oggi
non c'è più un Blockbuster ``manco a cercarlo col lanternino''. Non avevano
capito qual era il vero motivo del loro successo, cioè il loro modello di
business\footnote{Le percentuali sono quelle riferite in lezione e non sono
state verificate.}.

\paragraph{The Founder.} Nel film sulla nascita di McDonald's c'è una scena in
cui il futuro socio del protagonista gli dice: ``tu pensi di essere nella
\emph{food industry}, e invece sei nella \emph{real estate industry}''. Non
dovrebbe investire sugli hamburger, ma sui terreni su cui costruisce i
ristoranti. Anche per un team vale lo stesso: bisogna capire che cosa
l'ha reso davvero di successo, oppure no.

\subsection{L'autoanalisi personale}

\paragraph{E quando si cambia azienda?} Uno studente chiede se il debrief sia
possibile anche quando si lascia un'azienda. Non si può tornare in ufficio a
discutere della propria esperienza, ma lo si può fare \textbf{per se stessi}.
Il relatore l'ha fatto, anche se all'epoca in modo poco consapevole, ed è
così che è passato dal \emph{community management} al \emph{people
management}. Analizzando cosa aveva fatto bene a Blizzard si era reso conto
che:
\begin{itemize}
  \item aveva costruito un ottimo rapporto di fiducia con i giocatori (con
        alcuni proprietari di fansite parla ancora oggi);
  \item pur senza la seniority, era diventato un punto di riferimento per il
        team, visto come imparziale ed equilibrato anche su temi che non lo
        riguardavano direttamente;
  \item sapeva \textbf{risolvere problemi}: durante il lancio di
        \emph{Diablo~III} era il più tranquillo di tutti, non perché
        sminuisse il problema, ma perché sapeva che si sarebbe risolto con i
        suoi tempi; il lavoro non era spegnere l'incendio ma ``far sfollare
        la gente'' e aspettare che tutto tornasse in ordine.
\end{itemize}
E aveva capito cosa sapeva fare meno bene, ad esempio \textbf{gestire il
proprio manager}. Avere un manager non significa che il manager gestisca te:
il manager gestisce il tuo lavoro, ma tu devi saper gestire i tuoi
\emph{stakeholder}. Spesso lo stakeholder è il proprio manager; da manager,
gli stakeholder saranno CEO, CFO e così via. Gestirli significa conoscerli,
sapere come portare loro le informazioni, come far presente che qualcosa non
va. Fa parte del lavoro, anche se a lui non piace, e ci ha lavorato.

Nel mese e mezzo di pausa prima di entrare in Riot si è concentrato su
questo: ``cosa porto come mio punto di forza?''. L'autoanalisi è culminata
nel chiedere espressamente alla persona che l'aveva assunto: ``prima mi avete
detto di no, ora mi avete preso: perché?''.

\paragraph{Come si capisce se un progetto è andato bene?} Dipende
dall'\textbf{obiettivo}. Se l'obiettivo era vendere milioni di copie e il
gioco le ha vendute, bisogna capire perché. Se, più ragionevolmente,
l'obiettivo era un gioco da tavolo che permettesse di mettere in pratica le
conoscenze del corso, va valutato per quello: ha permesso di farlo? Se sì,
perché? Se no, perché no? L'obiettivo va fissato \textbf{all'inizio del
progetto, non alla fine}, proprio per poterlo valutare dopo.

\subsection{Quando il team cambia di continuo}

\paragraph{Tutti questi cambiamenti non sono dannosi per il progetto?}
Assolutamente sì, ma quasi mai sono volontari. I progetti andati male sono
spesso progetti grandi.

Il relatore racconta di \emph{Cyberpunk 2077}. Conosceva persone in CD
Projekt Red (``il mondo è piccolissimo, le cose si sanno''). Alla Gamescom va
con Diego Ricchiuti, direttore dell'accademia, a vedere la presentazione del
primo gameplay, quella ``in cui regalavano il giubbottino da samurai''.
Finito il video, tra gli applausi, si guardano e capiscono entrambi che
\textbf{il gioco non è pronto}, a soli quattro mesi dall'uscita prevista:
Diego aveva notato il combattimento e la guida. Gli amici di CD Projekt Red
(uno aveva lavorato con lui a Blizzard per sei anni) gli confermano tra le
righe che ``ce la faremo, ma ci sono ancora cose da sistemare''. Il gioco fu
poi rimandato più volte\footnote{In lezione il relatore dice ``due, forse
tre volte''; l'uscita di \emph{Cyberpunk 2077} è stata effettivamente
rinviata tre volte prima del dicembre 2020.}.

Per deformazione professionale, chi ha sviluppato l'occhio critico non
riesce più a godersi un trailer: in un trailer di \emph{The Day Before} Diego
gli aveva fatto notare che l'effetto della fiammata di un fucile era stato
chiaramente aggiunto in postproduzione e non faceva parte del gameplay. ``È
il triste epilogo di questa storia'': non riescono più a godersi un trailer.

Il punto è che più il team è grande, più queste variabili vengono imposte, e
più è facile che il fallimento \textbf{non dipenda dal fatto che il prodotto
sia brutto} (\emph{Cyberpunk} è un bellissimo gioco sotto molti punti di
vista, e il relatore lo sta rigiocando ora che è stato corretto ed espanso).
Nei grandi gruppi qualcuno arriva, qualcuno si dimette, cambia il budget,
arriva il rinvio; nel caso di \emph{Anthem}, il team ha fatto così tante volte
forming--storming--forming--storming che ne è uscito ``un prodotto che
dovevamo espellere: l'abbiamo espulso, buttato fuori dal nostro organismo pur
di liberarcene''.

Tra i cambiamenti c'è anche la crescita verticale: se il capo viene promosso
perché lavora benissimo, bisogna scegliere un nuovo capo, ed è sempre un
caos (chi ci resta male, chi pensava di meritarlo, chi voleva il suo amico,
chi non sopporta il prescelto). Per questo il relatore è \textbf{contrario
alle promozioni durante i progetti}: si rischia di rovinare tutto il lavoro di
storming.

\subsection{Obiettivi chiari nei piccoli gruppi}

Nei gruppi piccoli, come quelli degli studenti, il problema deriva più
spesso dal \textbf{non avere chiaro che cosa si vuole ottenere}: un prodotto
puramente educativo e formativo, o un prodotto richiesto dall'accademia che,
se viene bene, si potrebbe anche vendere? Un gioco commerciale o un gioco
artistico? Va chiarito subito, perché ci saranno sempre momenti in cui
scegliere tra A e B, e a guidare la scelta sarà la domanda: \emph{quale delle
due opzioni ci porta più vicino all'obiettivo?} Se l'obiettivo non è chiaro,
non si sa cosa scegliere.

Questo porta molte aziende a fare \textbf{scelte di design che in realtà sono
scelte commerciali}. Le microtransazioni sono scelte commerciali: nessuno ha
mai pensato ``sai cosa piacerebbe tanto ai giocatori? Pagare le bustine''.
Prima le bustine si guadagnavano giocando; ora la tendenza è ``potresti
giocare, ma se paghi è meglio''. Si pensi alle polemiche su Ubisoft per
l'opzione di ``saltare il gioco'' pagando, o ai famosi \emph{boost} di livello
di \emph{World of Warcraft}: ``devo pagarti per saltare il gioco?''.

Una studentessa (Silvia) cita \emph{Final Fantasy XIV}, che a un certo punto fu
``resettato'' perché non funzionava: un esempio positivo. Il relatore aggiunge
quello negativo di \emph{Age of Conan}, che non è mai riuscito a fare quel
reset iniziale (``è una schifezza, ve lo ripresentiamo quando è pronto''). Ci
vuole consapevolezza, ma anche i mezzi: un conto è Square Enix, che può
permetterselo, un altro è un'azienda che non può chiudere tutto per un anno
continuando a pagare gli stipendi.

Riassumendo: nei piccoli gruppi i problemi sono più di
\textbf{comunicazione e chiarezza} (obiettivi non chiari o che si spostano);
nelle grandi aziende di più di \textbf{persone che vanno e vengono, obiettivi
e tempistiche che cambiano}. In accademia le scadenze sono fisse (il gioco da
tavolo si consegna a fine anno scolastico, a giugno, e giugno resta); in
azienda può succedere che una feature prevista per la \emph{vertical slice}
di luglio venga anticipata a maggio per portarla alla Gamescom, e un sistema
perfetto per consegnare a giugno può non esserlo più per consegnare a maggio.

\subsection{L'ambiente determina la qualità del prodotto}

\paragraph{È vero che chi passa per Electronic Arts finisce per fare brutti
giochi?} Uno studente cita BioWare, che faceva ottimi giochi finché era
autonoma. Il problema non è dei developer, che non sono ``improvvisamente
diventati incapaci'': è evidente che il problema è altrove, e che
\textbf{l'ambiente in cui si fa un gioco determina molto spesso la qualità del
prodotto}.

Un esempio locale positivo: Tommaso Bonanni, conosciuto dal relatore insieme a
Diego e già docente dell'accademia, ha fondato a Roma Caracal Games Studio,
applicando fin dall'inizio questi concetti di cultura aziendale. L'azienda è
rimasta piccola, ma ha molto chiaro cosa vuole fare e come. Quando si
predispone fin dall'inizio l'azienda e il team a lavorare in un certo modo è
fattibilissimo; è molto più difficile quando si è EA, si acquisisce
un'azienda e le si dice ``da oggi si fa così''.

Il relatore ha vissuto la transizione da Blizzard ad Activision Blizzard (uno
dei motivi per cui se n'è andato, volontariamente). Ha visto la cultura
interna cambiare per gradi:
\begin{enumerate}
  \item ``facciamo giochi fighissimi, perché i giochi devono essere
        fighissimi'';
  \item ``facciamo giochi fighi così vendono di più'';
  \item ``vendiamo di più facendo giochi fighi'' (sembra uguale, ma è molto
        diverso);
  \item ``facciamo più soldi''.
\end{enumerate}
Lo si capiva dalle scelte di design, osservando quale fosse il peso più
grande sulla bilancia. Lo stesso varrà per gli studenti, nei loro giochi e nei
colloqui di lavoro: capire subito che cosa vuole l'altra parte.

\paragraph{Leader e visione.} Nei gruppi dell'accademia, dal secondo anno, c'è
una persona che assume il ruolo di \emph{lead} o \emph{scrum master}. È
diverso dall'essere un \emph{game director} (``il gioco si fa come dico io,
aiutatemi a farlo con la mia visione''), come avviene nei grandi studi: da lì
nascono i nomi nei titoli, come \emph{Sid Meier's Civilization}, che indicano
la persona che ha fatto quel gioco, oppure un puro branding, come \emph{Tom
Clancy's Rainbow Six}.

\paragraph{Inclusion come jolly.} A margine, uno studente osserva che in molte
aziende ``la risposta giusta è sempre \emph{inclusion}'', usata come un
jolly. È comunissimo, conferma il relatore: \emph{diversity and inclusion}
viene spesso usata per dire ``così includiamo più clienti'', ma non è quello
il suo punto (``altro discorso'').

% ---------------------------------------------------------------------
\section{Confrontarsi con chi è più bravo}
\label{sec:z1-rivalita}

Arrivato a Riot dopo sei anni a Blizzard, il relatore pensava di avere un buon
bagaglio di esperienza. Si è trovato invece ``come un fante buttato in mezzo
ai Navy SEAL'': persone che avevano superato 11, 12, 13 colloqui, con un
livello molto alto.

Individua subito un collega francese che, nella sua testa, diventa ``la
versione migliore di me'': stesso ruolo (community manager), esperienza
simile (arrivava da Ubisoft, come quasi tutti i francesi di Riot), ma molto
più spigliato, sempre con la battuta pronta, ``l'anima della festa''. Il
relatore invece è introverso: un workshop come questo per lui è
\emph{emotivamente drenante}, non ama parlare in pubblico pur amando
condividere le proprie esperienze, e non riesce a sostenere conversazioni
come ``cosa hai fatto questo weekend?''. Lavorando raramente insieme, aveva
sviluppato una rivalità ``da allenatore di Pokémon'' (lui era Ash, il
relatore l'antagonista).

Dopo sei mesi, a un evento di \emph{team building} aziendale (a Tenerife), il
suo capo, molto attento a queste dinamiche, li mette apposta nello stesso
gruppo. A fine serata, complice l'open bar, il relatore gli confessa di averlo
mal sopportato dal primo giorno per invidia della sua estroversione. Il
collega ride e risponde che pochi giorni prima aveva detto al capo la stessa
cosa al contrario: invidiava a Stefano il suo essere sempre organizzato, sul
pezzo, presentarsi nella sala riunioni giusta (mentre lui finiva in quella
sbagliata), arrivare con gli appunti, con i dati a supporto e le analisi di
mercato, mentre lui spesso non aveva nemmeno letto la mail del meeting.
Un'invidia in due direzioni, che li ha portati a lavorare molto insieme e a
migliorarsi a vicenda.

Le lezioni:
\begin{itemize}
  \item non bisogna demoralizzarsi vedendo qualcuno che fa lo stesso lavoro
        ed è più avanti sotto alcuni aspetti: spesso c'è un bilanciamento. I
        ``punti caratteristica'' sono limitati: se li ha messi in carisma,
        non li avrà messi in saggezza;
  \item conviene \textbf{chiedere a queste persone come mai} sono così brave
        in qualcosa, come si fa con un progetto andato bene: magari non lo si
        farà mai bene quanto loro, ma si comincerà ad avere il
        \emph{mindset} giusto; e viceversa, loro saranno spinti a chiedere a
        noi;
  \item spesso \textbf{le proprie doti le scopriamo dagli altri}. Il relatore
        non si era mai accorto di essere ``quello che resta calmo durante le
        emergenze'' (il lancio di \emph{Diablo~III}, quello di
        \emph{Draenor}, i problemi dal vivo alla Games Week) finché non
        gliel'hanno fatto notare tutti: per lui era la normalità. A chi pensa
        ``non sono capace, forse non è il mio percorso'', può darsi che
        semplicemente nessuno gli abbia ancora detto ``questa cosa la sai
        fare davvero bene, concentrati su questo''.
\end{itemize}

% ---------------------------------------------------------------------
\section{Domande sul lavoro nell'industria}
\label{sec:z1-industria}

Diverse domande degli studenti hanno riguardato le condizioni di lavoro nei
grandi studi internazionali. Le risposte sono raccolte qui per argomento.

\subsection{Presenza, remoto e trasferimenti}

\paragraph{Grandi studi: quasi solo in presenza.} Sia a Blizzard sia a Riot si
lavorava \textbf{solo in presenza}, senza smart working: sei anni a Parigi e
cinque a Dublino. Si parlava solo inglese: anche nei meeting tra soli
italiani si tendeva a comunicare e a scrivere in inglese, per correttezza
verso chi si fosse aggiunto e per lasciare una traccia comprensibile a tutti.
Il francese serviva per fare la spesa e sopravvivere, non in ufficio.
L'ambiente era estremamente eterogeneo (francesi, polacchi, spagnoli e
molti altri): durante il colloquio dal vivo a Dublino, mentre si preoccupava
di essersi vestito in modo adeguato, vide passare dietro la parete a vetro un
collega svedese con i dreadlock, a piedi scalzi, con una pianta sotto un
braccio e un aquilone sotto l'altro, e capì di essere ``a posto così''.
Un'esposizione a molte culture e modi di pensare che il relatore considera
bellissima.

Anche oggi, tranne alcuni studi particolari, la maggior parte del lavoro si
svolge di persona: amici a Paradox lavorano a Stoccolma, i pochi rimasti a
Blizzard si sono dovuti trasferire a Irvine. Per posizioni \emph{junior} o
\emph{entry level} i grandi studi raramente concedono il lavoro da casa; per
profili senior può esserci qualche eccezione (ad esempio nove mesi da remoto
e tre in sede).

\paragraph{Relocation.} Spesso tocca spostarsi, ma il relatore non ha mai
pagato un trasloco di tasca propria: queste aziende offrono la
\emph{relocation} (aiuto a trovare casa, ad aprire il conto in banca e
supporto pratico), perché è un investimento sulla persona.

\paragraph{Colloqui mentre si lavora altrove.} Per i colloqui in presenza
prendeva ferie, senza dire dove andava. Il consiglio generale: per fare
colloqui \textbf{prendete ferie o permessi}, non datevi malati, ``che non è un
bel modo di comportarsi''.

\paragraph{Remoto ``doppio''.} Un collega lavora per l'azienda coreana che
gestisce \emph{PUBG}, con ufficio ad Amsterdam, ma lavora da Milano: un
``doppio remoto''. Si trova, ma in genere un ambiente anglofono con sedi in
Europa chiederà di trasferirsi in Europa.

\paragraph{Famiglia e trasferimenti.} Il relatore ha avuto la fortuna di
sposare una persona che fa lo stesso lavoro (conosciuta lavorando su
\emph{Warhammer Online}, poi con lui a Blizzard e a Riot, nella
localizzazione e nel doppiaggio, ad esempio di \emph{World of Warcraft} in
italiano): avendo competenze rare, spesso venivano assunti entrambi, al
massimo a pochi mesi di distanza. Per i developer la soluzione frequente è
che la famiglia segua la persona assunta, perché lo stipendio basta per
tutti. Ma non sempre c'è equilibrio, e a volte uno dei due deve sacrificarsi:
sua moglie ha lasciato Blizzard per seguirlo a Dublino (per poi essere
assunta da Riot); lui ha lasciato Riot per tornare in Italia quando lei ha
trovato lavoro qui. Ha dovuto rinunciare più volte a lavorare in Islanda per
CCP (\emph{EVE Online}), perché non era una scelta solo sua. Si fanno
compromessi, ma si riesce a far quadrare.

\paragraph{Ci si sposta per ogni progetto?} No: si sono spostati
esclusivamente per cambiare azienda. Cambiando progetto nella stessa azienda
è raro, quasi mai (``non l'ho mai visto succedere''), che venga chiesto di
trasferirsi. In Francia stavano benissimo (vivevano fuori Parigi; ``i
problemi con i francesi sono in realtà problemi con i parigini''), parlavano
francese ed erano vicini alla doppia cittadinanza: se ne sono andati perché
non vedevano un futuro a Blizzard, per il modo in cui l'azienda trattava i
colleghi.

\paragraph{Il mercato oggi.} Non esiste un \emph{silver bullet} per trovare
lavoro domani dove si vuole. Con tutti i \emph{layoff} recenti ci sono
migliaia di esperti del settore che cercano lavoro: serve pazienza, ma è
fattibile.

\paragraph{Studi indie e remoto.} Negli studi più piccoli il remoto è più
diffuso: a Stormind circa 60 persone su 95 lavorano da remoto, ma lo studio
si è strutturato fin dall'inizio per farlo. Serve consapevolezza. Bending
Spoons (che sviluppa app per mobile, non giochi, ed è l'azienda dell'app
Immuni), pur essendo grandissima e con sede a Milano, lavora da remoto e
lascia l'ufficio a chi vuole: lo smart working è posizionato quasi come un
\emph{benefit}, ma permette anche di avere un ufficio da 300 persone invece
che da 700. Le sedi fisiche costano decine di milioni all'anno (Blizzard ha
chiuso le sedi di Amsterdam e di Parigi). Tuttavia, per lavori molto
verticali e specifici gli studi preferiscono avere le persone dal vivo; i
ruoli più spesso remoti sono quelli che non fanno sviluppo (customer service,
PR\dots).

Anche gli artisti lavorano spesso in presenza: l'accademia può dare solo una
\emph{simulazione} di quanto stretto sarà il contatto tra programmazione,
design e art. ``Sarà ancora di più.''

\paragraph{La comunicazione da remoto.} È più difficile. Pur essendo
estremamente introverso, il relatore \textbf{obbliga tutti a usare la webcam}
nei meeting in cui conta: non in un aggiornamento con novanta persone
collegate, ma certamente in un colloquio a due sulla performance dell'ultimo
anno. La componente fisica è una parte enorme della
comunicazione\footnote{In lezione si parla di ``quasi l'ottanta per cento'';
il dato va preso come indicativo, poiché le percentuali di questo tipo
circolano in forme diverse e sono spesso citate fuori contesto.}: anche se si
è timidi, accendere la webcam è l'approssimazione più vicina all'essere dal
vivo. Negli ambienti dove si lavorava senza webcam la situazione era
``esponenzialmente peggiore''.

\subsection{Differenze culturali nel team}

\paragraph{Gestire sensibilità diverse.} A Riot i team erano divisi per
\emph{lingua}, non per nazione: il team di lingua spagnola includeva anche un
messicano, quello di lingua russa includeva russi e ucraini, proprio nel
periodo della guerra in Crimea. Tra i colleghi non c'era alcun problema (erano
amici), ma ci si è posti la domanda: \emph{come comunichiamo con una community
che sarà divisa su questo tema?} Nei giochi online la divisione era evidente:
in \emph{Rocket League} c'era chi, con un compagno di squadra con la bandiera
avversaria, si faceva otto autogol pur di farlo perdere. Allo stesso modo,
con molti americani mandati a Dublino a coprire posizioni scoperte, c'erano le
elezioni (Clinton--Trump, Biden--Trump) e varie guerre: serviva sempre una
certa sensibilità.

Imparare a navigare queste situazioni lo ha reso un professionista migliore:
gli ha insegnato a fermarsi e chiedersi ``io la percepisco così; gli altri
dieci, con background, storie e lingue diverse, come la penseranno?'', a
ragionare più \emph{long term} e più \emph{big picture}, con più empatia.

\paragraph{Stili di comunicazione.} Un partecipante che ha lavorato nel Regno
Unito chiede dello stile comunicativo italiano, a volte frainteso. Uno dei più
grandi complimenti ricevuti dal relatore è stato ``non sembri italiano''
(``gesticoli poco per essere italiano'', ``parli bene inglese per essere
italiano''). Ha imparato ad \textbf{adattare la comunicazione alla persona}:
con spagnoli e francesi la battuta sulla nazione funziona, perché c'è un
``odio--amore'' condiviso; con russi e giapponesi no. Doveva adattare anche
l'umorismo. Un collega metà russo e metà messicano diceva di aver preso ``il
lato peggiore di entrambe le nazionalità'', ed è stato un ottimo maestro per
imparare a leggere culture diverse.

Non bisogna generalizzare, ma gli stereotipi possono essere un punto di
partenza ragionevole, \emph{sapendo che troverete sempre le eccezioni}. E
vanno letti correttamente: gli americani tendono a essere più estroversi e ad
avere sempre qualcosa da dire non perché ``sono fatti così'', ma perché
vengono educati così (all'università tra i corsi c'è il \emph{debate}, saper
argomentare in pubblico); un collega giapponese può essere più introverso per
ragioni culturali. Nel dubbio si parte da lì, poi si conosce la persona.

\subsection{Negoziare, discriminazione e stipendi}

\paragraph{La famiglia come leva.} Avere un partner riduce le opzioni, ma le
aziende lo sanno: uno spostamento è per due. Si può usare come leva per
negoziare (un compenso diverso, un supporto per il trasloco): se un'azienda ha
deciso di volere quella persona, 2, 3, 4 o 10.000 euro in più all'anno, per
un'azienda che fattura centinaia di milioni o miliardi, non cambiano nulla.

\paragraph{Essere sposati è un motivo per non essere assunti?} No. In alcuni
paesi è proprio illegale discriminare su questo: in Francia era illegale
chiedere CV con data di nascita, foto, sesso o religione, e chiedere se si
vogliono figli. Purtroppo, dai racconti dei suoi studenti di master (in
settori come \emph{beauty and wellness} e \emph{food and wine}), in piccole
realtà italiane queste domande (``hai 25 anni, hai intenzione di
sposarti?'') vengono ancora fatte, e servono a filtrare i candidati. Nelle
grandi aziende ci sono sistemi per evitarlo; anzi, molte danno \emph{più}
benefit a chi ha famiglia (``porta pure la famiglia, è tutto pagato''),
perché le persone competenti hanno spesso una certa età e quindi,
ragionevolmente, una famiglia.

\paragraph{Gli stipendi.} Dipendono molto da dove si vive. Lo stipendio a Riot
non era paragonabile a quello di Blizzard perché in Francia era tassato al
52\% e in Irlanda al 17\%\footnote{Aliquote come riportate in lezione,
riferite alla situazione personale del relatore.}: nessuno stipendio
competitivo può compensare il 35\% di differenza di tasse. Negli Stati Uniti,
con pochissima previdenza pubblica, gli stipendi sono molto più alti perché
poi ci si paga tutto (``se hai un raffreddore paghi tu''). Bisogna valutare il
potere d'acquisto di quei soldi nel paese in cui si andrà a vivere.

\subsection{Freelance e \emph{work for hire}}

All'inizio della sessione pomeridiana uno studente che lavora nell'audio
chiede quali siano le condizioni per chi lavora come \emph{freelance}.
\begin{itemize}
  \item Dipende molto dal progetto: spesso il contratto è a progetto, per
        portare a termine una parte o un intero progetto.
  \item Non tutte le aziende usano freelance: Riot e Blizzard lo facevano
        pochissimo e preferivano assumere (magari trattando poi la persona
        ``come un freelance'', lasciandola a casa a fine progetto).
  \item Spesso i lavori offerti come freelance sono in realtà
        \textbf{sostitutivi di un full time} (``freelance, ma sei ore al
        giorno'').
  \item Si può fare, ma bisogna essere molto chiari nel definire l'impegno
        (quante ore, come vengono tracciate): non tutti specificano cosa
        significa davvero ``freelance''. Lavorando così per due o tre aziende
        si può arrivare a un tempo pieno.
  \item È più comune il \emph{work for hire}: aziende che forniscono personale
        ad altre aziende; si è assunti da X ma si lavora ``embeddati'' presso
        Y (una sorta di subappalto). Anche lo studio del relatore lo fa.
  \item Il puro freelance di quattro mesi esiste, ma di solito per quei
        quattro mesi diventa un lavoro a tempo pieno; è raro ``quattro mesi,
        due ore al giorno''. Nelle grandi aziende è improbabile lavorare nello
        stesso mese per Ubisoft, Riot e Blizzard tre ore ciascuna; è più
        realistico lavorare per più startup indie italiane, che non possono
        permettersi un'assunzione ma possono pagare due ore al giorno.
  \item Molto passa dai contatti: trovare le aziende giuste e rendere esplicite
        le condizioni.
\end{itemize}

% ---------------------------------------------------------------------
\section{Il circolo virtuoso: feedback, miglioramento, fiducia}
\label{sec:z1-triangolo}

\subsection{Il triangolo}

L'obiettivo a cui il workshop vuole arrivare è un \textbf{triangolo che si
autoalimenta}:
\begin{enumerate}
  \item si dà un \textbf{feedback} con l'obiettivo di migliorare le
        performance (inclusi elementi caratteriali);
  \item la persona \textbf{migliora} e vede di averne beneficiato;
  \item questo genera \textbf{fiducia}, che rende più facile dare e ricevere
        il feedback successivo.
\end{enumerate}
Togliendo uno qualsiasi dei tre elementi, il circolo non è altrettanto forte.

In qualsiasi fase si parta, il rispetto per i colleghi è inizialmente dettato
dall'\emph{educazione}, non dalla fiducia. La fiducia arriva quando i colleghi
cominciano a dire cose giuste e sensate e a \textbf{mantenere la parola}.
Dare e ricevere feedback è il modo in cui, in un certo senso,
``inganniamo il cervello'' per creare una \emph{case history} di fiducia: si
ha motivo di fidarsi di più, di ascoltare di più, di essere più ricettivi. E
si diventa sempre più bravi ad accettare il feedback e ad agire di
conseguenza.

\subsection{Accettare un feedback non vuol dire fare come dice l'altro}

Dieci anni fa il relatore non era ostile ai feedback, ma molto
\emph{difensivo}: ``sì, sì, va bene, ora ci guardo''; oppure, a chi gli diceva
``avresti dovuto farla così'', rispondeva spiegando perché l'aveva fatta in
quel modo. Oggi ha imparato a dire ``ok, grazie, ora ci ragiono'': non deve
difendere la propria posizione da chi viene solo a dire ``hai pensato a
questa cosa?''. Prima lo viveva come un attacco (``mi viene a dire che non so
fare il mio lavoro''), ora come un ``ottimo punto di vista, lo prendo in
considerazione''.

Con una postilla fondamentale: \textbf{accettare il feedback non significa
fare come dice l'altro}. ``Mi hai detto che preferisci questo mob verde. Ho
pensato anche a farlo verde? Sì. Penso comunque che giallo sia meglio? Sì, lo
faccio giallo.'' Il relatore accetta \emph{sempre} il feedback, valido o no:
garantisce di tenerlo a mente quando lavora. Che poi faccia come gli viene
detto è un altro discorso.

\subsection{Fiducia e scadenze}

La fiducia funziona allo stesso modo con gli impegni. Se si chiede a qualcuno
una cosa per venerdì e la consegna venerdì, la volta dopo gliela si chiederà
per venerdì. Se invece arriva il giovedì della settimana dopo, la volta dopo
gli si darà una \textbf{deadline anticipata}, come con chi arriva sempre in
ritardo agli appuntamenti (``se ti dico alle 7 arrivi alle 7 e mezza; la
prossima volta ti dico alle 6 e mezza'').

Al lavoro questo è dannosissimo, perché crea \textbf{deadline fittizie} che
non sono mai quelle vere, e nascono discussioni del tipo ``perché devo farlo
in due giorni se si fa in cinque?'' (``sì, se tu lo facessi in cinque''). La
conversazione giusta è un'altra: \emph{``che problemi hai avuto nel
consegnare venerdì?'' ``Questo, questo e questo.'' ``Cosa posso fare per
evitare che succeda?''}

Il relatore non vuole mai sentirsi dire ``non sono riuscito a farlo''.
Preferisce di gran lunga sentirsi dire, tre giorni prima, \textbf{``potrei non
riuscire a farlo, per questi motivi''}. Chi ha la visione d'insieme può dare
più budget, più persone, spostare la deadline, cambiare lo scopo del progetto:
può fare qualsiasi cosa per aiutare. Ma se lo si scopre il venerdì, cinque
minuti prima della deadline, non può fare più niente. Il messaggio è:
\textbf{``aiutami a poterti aiutare''}. Non penserà mai male di chi dice ``non
ce la faccio per questi motivi'', perché non sta dicendo ``non ho voglia'': sta
indicando i problemi, e lavorare su quei problemi è compito del manager. Non
può invece fidarsi di chi fino al giorno prima dice ``tranquillo, ce la
faccio'' e il giorno dopo ``non ce l'ho fatta''.

Questo è fondamentale per chi lavora in programmazione, design o art, perché
ci saranno sempre deadline: bisogna sempre \textbf{lavorare sulle aspettative},
dire al proprio manager (o al proprio report) cosa è lecito aspettarsi. Non
bisogna fare come i pinguini di \emph{Madagascar} (``sorridete e salutate, va
tutto bene''): se c'è un problema, va detto, altrimenti si peggiora solo la
propria situazione.

Da italiani si tende a non voler mai alzare la mano per chiedere aiuto
(``no, no, ce la facciamo, lo faccio io''). A volte non ce la si può fare, non
perché non si è capaci, ma perché vari fattori lo impediscono. Accettarlo è
importante: \textbf{è molto meglio avvisare e chiedere una mano che dire
``non ce l'ho fatta''}.

\subsection{Domande sul dare e ricevere feedback}

\paragraph{``Ne parliamo a fare, se poi fai come vuoi?''} Uno studente (Cristian)
riceve feedback che, dopo averli analizzati, non ritiene congrui; quando non
li segue gli viene detto ``ne parliamo a fare?''. C'è un problema di
\textbf{aspettative}: l'idea che accettare un feedback voglia dire fare come
dice l'altro. Accettarlo significa prendere in considerazione ciò che è stato
detto. Chi decide diversamente ha però un \textbf{onere}: argomentare perché
si fa comunque a modo proprio (``mi hai detto verde, la facciamo gialla
perché\dots''). Poi bisogna solo aspettare i risultati per vedere chi aveva
ragione.

La decisione spetta sempre \textbf{a chi ha la responsabilità di quella
decisione}. Se il game director deve scegliere tra A e B, anche se il
feedback dice B, la scelta è sua perché la responsabilità è sua. Non sarebbe
corretto essere ritenuti responsabili di una scelta delegata a chi dà
feedback. La risposta può essere semplicemente: ``l'ho tenuto in
considerazione, ma visto che ne sono responsabile io e le conseguenze sono
mie, ho scelto di fare X''.

\paragraph{E se l'altro fa ripicca?} Un partecipante chiede di chi, vedendo
il proprio feedback non accolto, si vendica ignorando a sua volta feedback
buoni. Per il relatore è un problema caratteriale: chi si comporta così ``è
già sulla strada per uscire''. Non c'è spazio per l'ego in team di questo
tipo. Il fenomeno nasce dall'idea che \emph{il feedback negativo faccia fare
brutta figura}, che il fatto che si potesse fare meglio tolga qualcosa alle
proprie capacità.

Non è così: \textbf{il feedback non è un giudizio}. Si può passare un anno a
Riot ricevendo solo feedback negativi e avere una valutazione annuale
stellare, perché ricevere molti feedback negativi significa essere una persona
approcciabile, che agisce sul feedback e invita gli altri a darlo. Il relatore
ha visto valutazioni molto peggiori per chi riceveva solo feedback positivi,
perché la gente aveva paura di dargliene di negativi: ma solo ``pacche sulle
spalle'' non fanno migliorare. Il problema va affrontato parlando con la
persona; a volte le conversazioni iniziano con \textbf{``ti do un feedback sul
feedback''}, cioè sul modo in cui l'altro ha dato il suo.

Esempi dalla sua esperienza:
\begin{itemize}
  \item lui chiede spesso di \textbf{non ricevere feedback senza esempi
        tangibili}, perché fa fatica ad astrarre: ``sei stato maleducato'' non
        significa nulla, ``ieri al meeting mi hai risposto male su questa
        cosa'' sì;
  \item ad altri serve l'opposto: ha ricevuto il feedback ``quando devi dirmi
        che ho sbagliato, puoi girarci un po' intorno e darmi del contesto?'',
        perché il suo stile diretto (``questa cosa non ha funzionato, dovremmo
        fare così'') feriva. Ha adattato il suo modo di dare feedback: con
        alcuni è diretto, con altri gira intorno. L'importante è saperlo.
\end{itemize}

\subsection{Non portare la soluzione}

Un consiglio centrale: \textbf{quando si dà un feedback, non si dà una
soluzione}. Non si dice ``io farei così'', ma ``secondo me questa cosa si
potrebbe fare diversamente, per questo e questo motivo: come potremmo farla
diversamente?''. È come chi smette di fumare perché vuole e chi smette perché
glielo dicono moglie e figli: se non si è convinti, spesso non si smette.

Quando il relatore diceva ``c'è questo problema, io farei così'', il meeting
successivo iniziava inevitabilmente con ``ci ho provato come hai detto tu, con
me non ha funzionato''. Sempre. Perché la soluzione era sua. La
responsabilità della soluzione deve essere invece \textbf{congiunta}: ogni
feedback deve creare un piccolo team che ha come obiettivo risolvere quel
problema. Questo toglie l'antagonismo (non ``ti dico che non sai fare il tuo
lavoro'', ma ``ti voglio aiutare a farlo meglio, lavoriamoci insieme'') e
sposta il focus: il problema non sei tu, è ciò che sta accadendo.

\paragraph{E se sono i piani alti a tenere all'oscuro dei problemi?} È una
questione di cultura aziendale. Il relatore fa sempre presente al proprio capo
il concetto di \textbf{service level agreement} (che verrà ripreso nella
seconda giornata): ``se mi dai questi input, ti garantisco questo output in
questi tempi''. Se il capo chiede delle cose senza dire budget, tempi e
perché, non può aspettarsi lo stesso risultato. È una comunicazione a due
vie: si dice al manager cosa serve per fare il lavoro, e lui deve fornirlo o
accettare che il risultato non sarà quello voluto.

\subsection{Amici, parenti e parole proibite}

\paragraph{``Non si lavora con amici e parenti''?} Per sua definizione
personale il relatore ha pochissimi amici (``stanno sulle dita di una mano'')
e una soglia molto alta per considerare amico un collega. Allo stesso tempo,
in quindici anni e con centinaia di colleghi, ha avuto solo \textbf{due}
esperienze in cui non è riuscito a lavorare con qualcuno: non serve essere
amici per far funzionare le cose. Con gli amici, anzi, gli viene più facile:
se dice a un amico che una cosa ``fa schifo'', l'amico non pensa ``come si
permette'', ma ``se perfino Stefano, che è mio amico, me lo dice, allora è
proprio così''. Diventa un rafforzativo, grazie alla fiducia costruita prima.
Il problema del lavorare con amici e parenti nasce dal \emph{non voler} dire
all'amico che sta sbagliando; se invece la confidenza permette di dirlo, si
può saltare la ``confezione'' del messaggio e passare direttamente al
contenuto.

\paragraph{Le frasi peggiori.} Alla domanda se esistano parole peggiori di
``ottimo lavoro'', il relatore indica la sua classifica: ``ottimo lavoro''
(detto in modo vuoto)\footnote{Il passaggio della trascrizione è
frammentario: si intende verosimilmente il complimento generico e di
circostanza, che non dice nulla di utile.}; poi la frase che lo ``fa uscire di
testa'', \textbf{``abbiamo sempre fatto così''} (``e infatti avete chiuso'',
come Blockbuster); infine \textbf{``siamo una famiglia''}, riferito al posto di
lavoro. Un'azienda non è mai una famiglia, perché non ci si può aspettare lo
stesso impegno da chi la possiede e da chi ci lavora: non c'è un
corrispettivo proporzionale tra impegno e ritorno. Per quanto inclusiva,
flessibile e razionale, non sarà mai una famiglia.

\subsection{Ego, muri e il peggior modo di scegliere un leader}

\paragraph{Chi mette davanti l'ego e alza muri.} Esistono, ma di solito durano
pochissimo in azienda. Non si tratta di avere un ego forte, ma di \textbf{non
saper ricevere feedback}. Spesso sono persone bravissime nel loro lavoro che
poi si chiedono perché tutti intorno diventino più bravi di loro: perché gli
altri accettano feedback e hanno più fonti da cui attingere, mentre loro si
sono rintanati nella loro torre.

Molti manager tendono ad \textbf{assumere persone un filo peggiori di loro},
per non fare brutta figura. Il relatore assume solo persone più brave di lui,
``perché non è una gara'': ha lavorato con molti community manager più bravi
di lui come community manager, ma lui era più bravo a fare il manager. Non
viene valutato su quanto sa fare il community manager, ma su quanto il suo
team sa farlo. Un partecipante chiede se la paura di avere qualcuno di
migliore esista anche in ambito accademico. Assolutamente sì: al relatore è
capitato di non essere preso come docente in un corso di formazione sui
videogiochi perché il responsabile aveva percepito che ne sapesse più di
lui\footnote{La trascrizione dice ``che non sapevo più di lui'', ma il senso
del passaggio richiede il contrario.}. ``Ma perché te ne importa? Tu sei quello che ha portato chi
sa fare il suo lavoro: ben venga.'' O si ha questa cultura, o si ha paura per
il proprio posto.

\paragraph{Promuovere il migliore.} Quando bisogna scegliere un leader per un
team (programmazione, design, art), molto spesso si sceglie \textbf{il migliore
tecnico}: il miglior programmatore, il miglior designer. È spesso la scelta
sbagliata: è come prendere un idraulico per fare l'elettricista. Fare il
\emph{people manager} non è come fare il designer o il programmatore; si può
essere un ottimo programmatore e un pessimo manager, e viceversa. Le aziende
l'hanno capito e lo fanno meno di prima. A Blizzard c'era un grave problema di
\emph{middle management}: molte persone in ruoli importanti erano state
promosse cercando le skill sbagliate, erano frustrate dal non saper fare il
proprio lavoro e schiacciavano chi stava sotto per paura di fare brutta
figura.

Un partecipante cita una teoria simile: una persona viene promossa finché non
raggiunge \emph{il proprio livello di incompetenza}, dove si blocca\footnote{È
il \emph{principio di Peter} (Laurence J.\ Peter, 1969), citato in lezione
come ``la teoria di Peter''.}. Un altro cita Michael Scott di \emph{The
Office}, su cui il relatore tornerà (sezione~\ref{sec:z1-quando}).

% ---------------------------------------------------------------------
\section{Perché la cultura conta: retention, costi e fiducia}
\label{sec:z1-cultura}

\subsection{L'ambiente di lavoro e la retention}

Perché tutto questo è così importante? Perché l'ambiente di lavoro conta
moltissimo per la \textbf{retention}, cioè per trattenere le persone e farle
lavorare in modo efficiente. Se contasse solo lo stipendio, molti non
lascerebbero mai Riot o Blizzard; e più si va avanti con l'età, più si impara
a dare valore anche al resto.

Inoltre, secondo gli studi citati nelle slide, chi vive il proprio lavoro in
modo negativo è \textbf{meno efficiente nel problem solving}: l'energia
mentale va verso ``come faccio a non farmi dare la colpa'' o ``come faccio a
fare bella figura'', non verso la soluzione del problema. In ambienti negativi,
dove ci si deve preoccupare della propria ``sopravvivenza'', si lavora
peggio.

\subsection{Il costo di chi se ne va}

Ogni volta che qualcuno si dimette o viene licenziato, il costo di trovare
una nuova persona, fare i colloqui, formarla e farle attraversare forming,
storming, norming\dots\ è circa \textbf{un terzo del suo stipendio}. Per questo
le aziende tendono a fare più \emph{retention} che \emph{acquisition}, almeno
con le figure considerate utili alla produzione.

Allo stesso tempo si vedono sempre più spesso \emph{layoff} e \emph{hiring
wave}: si licenziano 3.000 persone e sei mesi dopo se ne riassumono 2.500 o
3.500, ma tutte più junior. Il relatore ne ha scritto in un articolo su
Eurogamer.it, in cui spiega anche alcuni ``trucchi del mestiere'': perché i
licenziamenti avvengono sempre in certi periodi e sempre ``a gruppi'' (prima
Bethesda, poi Blizzard, poi un'altra\dots).

\paragraph{L'offerta di Riot.} A Riot esisteva una pratica (di cui il relatore
non ricorda il nome) per cui, firmato il contratto, si aveva \textbf{un anno
di tempo per dire ``non mi piace, non me la sento, me ne vado''} e ricevere
subito come buonuscita un anno di stipendio. Avevano stimato che fosse molto
meglio pagare subito un anno a qualcuno e lasciarlo andare che tenerlo due,
tre, quattro anni come \emph{underperformer}: la responsabilità della scelta
di assumere era dell'azienda\footnote{La pratica, nota come ``The Offer'', è
descritta in fonti pubbliche con condizioni diverse (un periodo più breve e
una somma pari a una frazione dello stipendio annuo). Si riportano le
condizioni come raccontate in lezione; andrebbero verificate.}. L'hanno usata
pochissimi, ma il solo fatto che esistesse faceva lavorare meglio anche lui da
manager, perché sapeva che chi non si trovava bene poteva andarsene con un
anno di tempo per trovare un altro lavoro.

\paragraph{Il ricircolo.} C'è un forte ricircolo nell'industria: molti ex
colleghi passano da Ubisoft a Paradox, da Paradox a Funcom, da Funcom a
Blizzard. Il relatore ha scelto di uscirne lavorando solo per aziende
\textbf{indipendenti}: Stormind non è un indie di cinque persone, ma è
indipendente, e non deve accontentare azionisti. Quando vede un'azienda con
utili miliardari licenziare migliaia di persone si chiede ``dov'è la linea''
(\emph{when enough is enough}): con azionisti che vogliono un ritorno, se
quello è l'obiettivo stanno facendo la cosa giusta, ma lui non condivide
l'obiettivo\footnote{In lezione si parla di ``6,5 miliardi di utile'' e di
``3.000 persone licenziate''; le cifre sono riportate come dette e non sono
state verificate.}.

\paragraph{L'open space vuoto.} A Riot, durante un'ondata di licenziamenti,
del suo team di trenta persone è rimasto solo lui: ``teniamo te come
coordinatore, poi ricostruirai il team''. Ha rifiutato di ricostruire un team
cancellato ``con un colpo di spugna'' e ha deciso che non avrebbe più lavorato
per una grande azienda. Nessuno aveva pensato di spostarlo, e ha lavorato per
tre mesi in un open space di cento metri quadri con la sua scrivania e altre
ventinove vuote. Su una lavagna aveva scritto i nomi dei suoi ragazzi e li
cancellava man mano che, anche grazie ai suoi \emph{referral}, trovavano
lavoro. Cancellato l'ultimo nome, si è dimesso: per lui il lavoro era finito.

\paragraph{Team che si dissolvono.} Uno studente cita un articolo su Larian
Studios, il cui management è contrario a ricostruire ogni volta il team a fine
progetto. È lo stesso problema: si pensa ai ``giochi Blizzard'' o ai ``giochi
BioWare'', ma le persone che li hanno fatti non ci sono più; è rimasto solo il
brand, e non ha senso avere le stesse aspettative da un gioco fatto da persone
diverse. E anche chi sopravvive ai layoff vive pensando che il prossimo
potrebbe essere lui: alcune persone, licenziate nell'ultima ondata a
Blizzard, gli hanno detto ``erano due anni che lo sapevo''. Vivevano a Cork,
``un paesotto irlandese dove o vai per lavorare a Blizzard o non ci vai'',
avevano comprato casa: non ci si può aspettare chissà che morale.

\subsection{La cultura è forte quando si lavora per gli altri}

\begin{quote}
  \emph{La cultura è forte quando le persone lavorano per gli altri e a favore
  degli altri, e non per se stesse.}
\end{quote}

Se si dà feedback per far migliorare l'altro e non per imporre il proprio
modo, se si aiuta per il bene del progetto e non per fare bella figura, la
cultura si rafforza. È dai \textbf{piccoli gesti} che emerge la cultura
aziendale: ad esempio, dal non dover chiedere ``perché mi stai aiutando?''.

\paragraph{L'incognita.} \textbf{L'incognita è spesso dannosa quanto le
cattive intenzioni.} In un esperimento (``mi pare a Stoccolma, cinque o sei
anni fa'') si distribuivano per strada banconote da 500 euro, e la maggior
parte delle persone le rifiutava chiedendosi ``dov'è la fregatura?''
\footnote{L'esperimento è riportato in modo vago in lezione e non è stato
possibile identificarlo con certezza.}. Lo stesso accadeva al relatore,
arrivato a PG e a Stormind da Riot: a ``vedo che sei in difficoltà, posso
aiutarti?'' la risposta era sempre ``perché?''. Alla fine rispondeva: ``perché
quando sarò in difficoltà io vorrei che tu mi aiutassi''.

Quando si comincia a lavorare per gli altri e si smette di pensare ``così avrò
la promozione, così avrò l'aumento'', promozioni, aumenti e rispetto arrivano
da soli. A Blizzard aveva fallito tre volte i colloqui per essere promosso a
manager (inteso come coordinatore del team); quando ha deciso di smettere di
pensarci e di fare solo il bene del proprio lavoro, la volta dopo è stato
promosso. Funziona così: chi aiuta i colleghi e fa andare bene il progetto
viene notato (``questa persona fa migliorare cinque persone, non solo se
stessa: vale di più'').

\subsection{Il lavoro è valutato come opera d'insieme}

Il relatore riconosce che buona parte del suo successo è venuto dal fatto che
il suo team lavorava bene, non dalla sua bravura personale. Anche i progetti
degli studenti, presentati a Diego, a Enrico, ad Andrea, verranno valutati
\textbf{come opera d'insieme}. Nessuno dice ``il \emph{combat system} di questo
gioco l'ha fatto X, che bravo'': si dice ``BioWare ha fatto un bel gioco'',
``Blizzard ha fatto un bel gioco''.

Lo ha replicato anche lato HR. In molte aziende ci sono \textbf{bonus
personali} legati a obiettivi (se il gioco vende X, se la feature è pronta in
tempo). Il relatore li ha dimezzati: metà restano personali, metà diventano
\textbf{bonus di team}, che tutti ricevono se il team raggiunge gli obiettivi.
Un incentivo concreto a collaborare: non solo si dice che il lavoro di gruppo è
importante, lo si dimostra dandogli più valore economico. Aveva proposto il
100\% di bonus di team, ma gli è stato concesso il 50 e 50; ha funzionato
comunque, anche se sarebbe stato ancora più potente. Ovviamente, alla
domanda ``perché mi aiuti?'' si può anche rispondere ``perché voglio il
bonus'': un aiuto dato ``per il motivo sbagliato'' era un compromesso che era
disposto ad accettare.

\subsection{Il coraggio di chiedere aiuto}

Un partecipante osserva che aiutarsi dovrebbe essere la norma. Il modello
diffuso, invece, è ``devi essere forte e indipendente, non chiedere mai una
mano'', e vale in tutte le direzioni (dal maschilismo alla ``donna forte e
indipendente che non ha bisogno di nessuno''): \textbf{non ammettere di avere
bisogno di aiuto è sempre uno dei motivi per cui le aziende vanno male}, che si
tratti di uomini, donne o persone non binarie. Il problema è non avere la
sicurezza di poter alzare la mano.

Il relatore ha cercato di togliere ogni conseguenza dall'aver alzato la mano:
se si \emph{depenalizza} qualcosa, le persone non sono più disincentivate a
farla. Ha detto espressamente che non avrebbe mai valutato nessuno per il
numero di volte in cui chiedeva aiuto o per il numero di feedback negativi
ricevuti, e questo ha funzionato.

\subsection{Perché si dà un feedback}

\emph{Perché avete dato il vostro ultimo feedback? Che cosa volevate ottenere?}
Le risposte della classe:
\begin{itemize}
  \item per \textbf{riconoscere il valore} di un collega, anche verso il
        manager, perché qualcuno che non aveva visibilità vedesse quanto
        quella persona fosse efficace;
  \item per \textbf{spingere una persona a superarsi};
  \item per aiutare un collega a rendere il prodotto all'altezza di ciò che
        voleva rappresentare;
  \item per \textbf{genuinità}: c'era qualcosa che poteva essere fatto meglio.
\end{itemize}
Il relatore aggiunge che la forma cambia col tempo: all'inizio era ``devo
darti un feedback, hai dieci minuti?''; ora è un messaggio informale su
Discord, perché ha costruito un rapporto di fiducia e sa che non verrà letto
come ``l'HR manager mi scrive, sono nei guai?'', ma come ``Stefano è così, mi
dà feedback, resterà tra noi e lo risolveremo insieme''.

Un altro motivo è \textbf{risolvere incompatibilità caratteriali}: le
conversazioni più difficili, del tipo ``io e te non stiamo funzionando, cosa
possiamo fare diversamente?''. Sono importanti perché, se due responsabili non
vanno d'accordo, il team se ne accorge, comincia a raccontarsi storie (``si
odiano, ci saranno problemi'') e, per difendere il proprio manager, comincia a
giudicare male il lavoro dell'altro gruppo. È la dinamica da fermare:
\textbf{valutare gli altri solo per difendere la propria fazione}.

Dare e ricevere feedback sono \textbf{due muscoli da allenare} (il relatore è
sempre stato più bravo a darlo che a riceverlo), e si impara col tempo. È come
fare colloqui: fare colloqui è una skill a parte, che non indica quanto si è
bravi nel proprio lavoro. Essere scartati non significa non saper fare quel
lavoro: spesso significa aver fatto male il colloquio o che è stato preferito
qualcun altro.

\paragraph{Feedback pubblici e privati.} Uno studente racconta di aver dato un
feedback negativo a un superiore, senza un motivo particolare, per dimostrare
a sé e ai colleghi che non si viene licenziati per questo. Il relatore coglie
l'occasione per precisare che \textbf{i feedback negativi si danno
individualmente}: lui dà moltissimi feedback di cui il resto dell'azienda non
sa nulla. Reagirebbe molto male se, in un meeting, qualcuno alzasse la mano per
dire ``comunque questa cosa non l'hai fatta bene'': quello non è feedback, è
discussione, accettabile solo se è un brainstorming e se l'opinione è stata
chiesta. Alcune aziende hanno piattaforme pubbliche di feedback, ma spesso
servono solo ad applaudire; lo stesso accade nelle aziende che
\emph{obbligano} a lasciare un certo numero di feedback a settimana, che
diventano ``lui è bravo, lei è competente, anche questa settimana l'abbiamo
fatto''. Bisognerebbe chiedere a chi le adotta cosa pensa di ottenere (forse
trasparenza), ma difficilmente una persona ``attaccata'' pubblicamente reagisce
bene.

% ---------------------------------------------------------------------
\section{Among Us: la caccia al traditore}
\label{sec:z1-among-us}

Prima della pausa pranzo il relatore mostra tre film di fantascienza di
qualche decennio fa e chiede che cosa li accomuni\footnote{I titoli non sono
leggibili nella trascrizione: viene citato \emph{Event Horizon} da uno
studente, e il film ``in mezzo'' è indicato come ``l'originale'',
verosimilmente \emph{La cosa} (\emph{The Thing}).}. In ognuno c'è un gruppo di
persone con un obiettivo e la presenza, reale o percepita, di qualcuno che
rema contro: \textbf{l'energia del gruppo viene spesa per trovare il
responsabile}.

\paragraph{L'esperimento degli alieni.} Nei master in presenza il relatore
propone un esercizio. Distribuisce segretamente cartellini plastificati con
scritto ``essere umano'' o ``alieno''; le carte alieno contengono istruzioni di
comportamento con cui gli alieni, che cercano di fingersi umani, si
tradirebbero (dare sempre ragione a chi ha parlato prima, finire ogni frase
con una domanda, parlare in un certo modo). Lo scopo: in dieci minuti formare
gruppi, e \textbf{vince il gruppo più numeroso senza alieni}; basta un alieno
per perdere. Si formano sempre gruppi piccoli (due, cinque, sei persone),
perché la tendenza è ``mi fido al massimo di voi due''.

Ma vincono i gruppi più grandi, per un motivo semplice: il relatore
\textbf{non distribuisce nessuna carta alieno}. Mostra a tutti entrambe le
carte, ma consegna solo carte ``umano''. Vuole simulare ciò che accade in
ufficio: la tendenza a percepire sempre qualcuno che rema contro (i manager,
l'altro team, quel collega), quando in realtà non c'è nessuno, ``è tutto nella
nostra testa''. Quando chiede le motivazioni, emergono sospetti
\emph{fabbricati}: ``faceva sempre domande'', ``parlava solo all'infinito''
(``no, parla così e basta''). Ci si fabbrica da soli i motivi per non fidarsi
degli altri, ed è comunissimo.

\paragraph{Among Us.} L'autosabotaggio è, per tornare ad \emph{Among Us},
\textbf{concentrarsi sul cercare cadaveri e traditori invece di fare le
task}: non si vince perché non si sono completate le task, mentre si cercava il
colpevole. Succede tantissimo negli uffici, ed è il motivo per cui l'80\% dei
problemi che il relatore risolveva erano di comunicazione. Entrando nei team
sentiva immancabilmente: ``facciamo il possibile, ma questa lotta continua con
l'altro gruppo è estenuante''. ``Perché c'è questo problema?'' ``Non lo so,
c'era già quando sono arrivato.'' Un paradosso: c'è un problema, nessuno sa
dire quale sia, ma c'è sempre stato e quindi lo si accetta.

\paragraph{L'esperimento delle scimmie.} È l'esempio, ``finto o reale'', delle
scimmie in gabbia\footnote{È il noto racconto delle ``cinque scimmie'', spesso
presentato come esperimento scientifico ma probabilmente apocrifo (lo stesso
relatore lo definisce ``finto barra reale''); nella trascrizione il numero
delle scimmie varia.}. In una gabbia con alcune scimmie pende una banana
elettrificata: chi la tocca fa prendere la scossa a tutte. Le scimmie imparano
e, quando una si avvicina alla banana, le altre la picchiano. Poi si
sostituisce una scimmia alla volta con una nuova: la nuova va verso la banana,
le altre la picchiano, e smette. Si continua finché nella gabbia ci sono solo
scimmie nuove, poi si toglie la corrente. Nessuna scimmia tocca la banana:
nessuna ha mai preso la scossa, ma tutte hanno imparato che chi tocca la
banana prende le botte.

Nei team è lo stesso: il motivo per cui si ``odia'' l'altro team si
trasforma, ma un vero motivo non c'è. Ne derivano la sfiducia, il non alzare
mai la mano per chiedere aiuto, il non dire mai ``bel lavoro'' agli altri
perché non sia mai che qualcuno ti veda fare un complimento ``a quelli''.
Risolvere queste dinamiche è lungo e tedioso, e sempre pieno di ``non mi
ricordo'', ``io non c'ero'' e ``abbiamo sempre fatto così'', la cosa più
frustrante per chi cerca di risolverle.

% ---------------------------------------------------------------------
\section{Il principio del default to trust}
\label{sec:z1-default-trust}

\subsection{La definizione}

Il relatore indica questa come la slide più importante della presentazione.
\textbf{Default to trust} è uno dei valori aziendali di Riot Games, che lui ha
tradotto, adattato e ``esportato'':

\begin{quote}
  \emph{Approccio ogni situazione di conflitto partendo dal presupposto che una
  persona abbia agito in un certo modo perché, con le informazioni a sua
  disposizione in quel momento, ha pensato che fosse la soluzione migliore, e
  non perché sia un'incapace o volesse danneggiare il progetto
  intenzionalmente.}
\end{quote}

In altre parole: \textbf{di default, mi fido}.

\subsection{L'esempio del ritardo al daily stand-up}

A Riot il relatore aveva un \emph{daily stand-up} dal vivo alle 15 con il suo
team di venti persone. Una ragazza arrivava regolarmente alle 15:15, per una,
due, tre settimane. Alla quarta ci ha parlato. \emph{Come iniziereste la
conversazione?} Le proposte della classe:
\begin{itemize}
  \item sottolineare l'importanza di esserci dall'inizio, perché potrebbe
        perdersi qualcosa che le interessa, e per rispetto dei colleghi
        presenti;
  \item ``devi arrivare in orario'', spiegando perché è importante;
  \item \textbf{chiedere perché}.
\end{itemize}

Il relatore osserva, anche dalle sue consulenze fuori dai videogiochi, che
\textbf{più si sale con la seniority}, più la risposta è ``devi smettere di
arrivare in ritardo'', spesso con un \emph{or else} finale (``altrimenti
succede quest'altra cosa''). Più si scende con l'età, e quindi con meno
abitudini consolidate, più la prima domanda diventa ``dimmi perché arrivi in
ritardo''. Già spiegare perché è importante è ``miglia avanti'' rispetto al
semplice ``smetti''.

Partire dal perché gli ha permesso di scoprire la causa: il meeting precedente
della ragazza, in un'altra zona del campus a un quarto d'ora di distanza,
finiva sempre tardi, e lei doveva attraversare il campus di corsa. Non aveva
la confidenza per dire all'altro manager, che insisteva perché si restasse
fino alla fine, che doveva andare; sacrificava quindi il meeting del relatore,
sapendo che era il più comprensivo dei due.

Il relatore si è offerto di \textbf{mediare}: ha detto all'altro manager che
o anticipava l'inizio del suo meeting, così da finire comunque in tempo,
oppure lasciava uscire la ragazza, perché quel meeting gli serviva. C'è anche
una questione di \textbf{economia di scala}, spesso sottovalutata: se venti
persone aspettano qualcuno per dieci minuti, si sono sprecati dieci minuti per
venti. Per lo stesso motivo il workshop è iniziato alle 10:05: iniziare alle
10:10 avrebbe significato far aspettare cinquanta persone cinque minuti in
più. A Riot c'era una scrivania vuota con un pupazzo gonfiabile, \textbf{Mr.\
Late}: sommando tutti i minuti di ritardo a meeting, call e conversazioni, in
un anno veniva fuori praticamente \emph{una persona a tempo pieno}. Serviva a
ricordare che anche cinque minuti contano, e che bisogna rispettare gli altri
arrivando in orario o almeno avvisando.

\paragraph{Gli effetti.} La conversazione era avvenuta nel primissimo periodo a
Riot, quando il team non aveva ancora motivo di fidarsi di lui. Risolvendo il
problema ha eliminato anche tutto lo strato di micro-dinamiche spiacevoli che
un manager impara a vedere (due colleghi che si guardano con la risatina
quando lei entra in ritardo). La ragazza ha cominciato a entrare tranquilla,
con una certa leggerezza: non le faceva piacere arrivare in ritardo (``ci sono
anche persone che lo fanno apposta, e quello è un problema da affrontare
diversamente''). Soprattutto, con i colleghi non parlava di lui come del capo
che minacciava, ma come di chi le aveva risolto un problema: era diventato
un \textbf{alleato}. La settimana dopo altre tre persone sono andate da lui
con un problema da risolvere. Ripetuto cinque, sei, sette, dodici volte, in un
paio di mesi poteva dire quasi qualsiasi cosa nei meeting sapendo che
richieste e critiche sarebbero state percepite come ``me lo chiede perché vuole
darmi una mano'', non per fare bella figura o per presentare i successi del
team ai suoi superiori.

\subsection{Condividere i successi, schermare i fallimenti}

Un altro elemento fondamentale (ripreso nella seconda giornata): il manager
\textbf{condivide i successi e protegge il team dai fallimenti}. ``Se va male
qualcosa è colpa mia; se va bene qualcosa, voglio che ne parliate voi.'' Il
lavoro del manager non dovrebbe essere valutato su ciò che presenta del lavoro
del team: tutti sanno comunque chi è il manager. Tanto vale lasciar presentare
il team, che potrà dire con orgoglio di aver presentato a Mike
Morhaime\footnote{Cofondatore ed ex presidente di Blizzard; nella
trascrizione ``Mike Morayme''.} la propria campagna europea per
\emph{StarCraft~II}. Ci sono manager che dicono ``come mi prendo le colpe, mi
prendo i meriti'', ma i meriti li hanno già: è inutile ostentarli. Con i
fallimenti invece il manager è la \textbf{rete di sicurezza} del team. Il
relatore ne ha fatto una questione di principio: \textbf{non scaricare mai la
colpa su qualcun altro}.

\paragraph{Colpevole o causa.} Un esempio che il relatore attribuisce a Pirate
Software, creatore di contenuti ed ex dipendente Blizzard\footnote{Il relatore
aggiunge che si tratterebbe del figlio dello sviluppatore a cui si sarebbe
ispirato il personaggio del giocatore ``senza vita'' nell'episodio di
\emph{South Park} dedicato a \emph{World of Warcraft}. L'informazione è
riportata come detta in lezione e non è verificata.}. Contattando il servizio
clienti Blizzard per motivi personali, aveva scoperto che i colleghi gli
avevano scritto nei ticket informazioni confidenziali che non avrebbero dovuto
condividere in quel modo, nemmeno con un collega. Non ha segnalato le persone
da punire: ha detto ``c'è un problema, evidentemente le policy non sono
chiare''. Quando gli hanno chiesto i nomi, ha rifiutato: ``se vi dico chi è,
punite quelli, e gli altri cinquecento continueranno a farlo; non avremo
cambiato la causa alla fonte''.

\begin{quote}
  \emph{È inutile trovare il colpevole. Voglio trovare la causa. Perché il
  colpevole cambia.}
\end{quote}

Il colpevole a volte è vittima della situazione: per questo la definizione
dice \emph{``con le informazioni a sua disposizione in quel momento''}. Se il
relatore avesse rimproverato la ragazza in ritardo, non solo non avrebbe
risolto il problema, ma avrebbe creato un'\textbf{ingiustizia}: lei, che
attraversava il campus di corsa per una colpa non sua e di cui non poteva
parlare, sarebbe stata anche rimproverata.

Nessuno fa un brutto lavoro apposta: nessuno è contento quando sbaglia o
perde. Perché al lavoro diamo per scontato che gli altri sbaglino
intenzionalmente? Se qualcuno ha fatto una cosa male, è perché in quel momento
credeva fosse la scelta migliore. La domanda giusta non è ``perché hai fatto
questa cosa che è stata un fallimento?'', ma:

\begin{quote}
  \emph{Che cosa ti ha fatto pensare che questa fosse la scelta giusta?}
\end{quote}

Magari erano limiti di budget, di tempo, di scopo: si agisce su quelli, non
sulla persona. Se si cambia la persona ma le si danno gli stessi fattori, è
molto probabile che arrivi alle stesse conclusioni, e ci si troverà a dover
``disciplinare'' anche lei.

\paragraph{La reputazione costruita sui rumor.} Durante il giro di
presentazione in un ufficio (non in nessuna delle aziende citate), al relatore
è stato detto: ``questo è il dipartimento di design; se ti serve qualcosa
chiedi a lei, non a lui, perché lui consegna sempre in ritardo''. Oltre a
trovarlo poco professionale, una settimana dopo ha chiesto: ``cosa ti ha fatto
questa persona?'' ``A me niente, però ho sentito che\dots''. Si torna alle
scimmie: magari quella persona non consegna in ritardo da mesi, ma ormai lo
dicono tutti. È un circolo vizioso da fermare subito.

\paragraph{Niente escalation inutili.} Lo stesso vale per il feedback: quando
qualcuno va dal relatore dicendo ``ho un problema con questa persona'', lui
chiede: \textbf{``perché ne stai parlando a me e non a lei?''}. ``Perché sei il
mio capo'' non basta. Passare da lui significa dare per scontato che l'altra
persona non abbia interesse a risolvere, mentre magari basta parlarle. È la
strada più facile (``lo dico a Stefano, ci pensa Stefano''), e proprio per
questo bisogna mettere dei paletti: il manager interviene quando serve, ma le
cose che si possono gestire da soli vanno gestite da soli. (Quando invece il
confronto diretto fallisce, l'escalation è legittima: si veda la
sezione~\ref{sec:z1-qa-finale}.)

% ---------------------------------------------------------------------
\section{Traiettorie di carriera e valutazione delle performance}
\label{sec:z1-traiettorie}

\subsection{Quale traiettoria vorreste?}

Il relatore mostra un grafico con quattro curve colorate, contrassegnate
anche da lettere (A, B, C, D) per eventuali persone daltoniche (``a me è
capitato che qualcuno scoprisse di esserlo durante una lezione''). Le curve
rappresentano la performance lavorativa nel tempo. Due domande: \emph{quale
vorreste essere? E quale pensate realisticamente di essere?} Un partecipante
aspira a B e pensa di essere C; quasi tutti scelgono B.

Perché non A, che sale più rapidamente? Un partecipante risponde che una
crescita così repentina può essere più dannosa che utile, come una ``rush'' di
studio molto faticosa che poi non si riesce a sostenere. Il relatore propone
una lettura suggerita da un amico psicologo, che spiega la forma del grafico:
di B, C e D \textbf{vediamo dove finiscono}; di A no. Potrebbe continuare a
salire oltre la scala del grafico, o crollare sotto D. Non saperlo ci spinge a
preferire B. (A chi obietta che anche le altre proseguono oltre il bordo, il
relatore precisa: il bordo destro è ``oggi'', e di A non sappiamo come sia
arrivata fin lì.)

\subsection{La carriera non è una linea retta}

Tutti immaginano la carriera come una linea: si parte da X e, attraverso
stage e posizioni, si arriva a Y. In realtà carriera e performance avranno un
andamento irregolare, con alti e bassi. Come non si accetta che un team possa
temporaneamente \emph{underperformare}, tanto meno lo si accetta da se stessi.
Bisogna accettare che ci saranno momenti bassi.

Per questo il relatore \textbf{detesta le valutazioni annuali delle
performance}: prendono un punto casuale della traiettoria, che se si è
fortunati è il picco più alto e se si è sfortunati il più basso, e lo usano
per valutare un intero anno. Insegna ai manager a fare più punti di controllo
e a guardare altre cose:
\begin{enumerate}
  \item \textbf{quanto tempo passa tra un picco e l'altro}, cioè quanto
        \emph{cooldown} serve a una persona. È il motivo per cui, dopo la
        Games Week, tutti staccavano una o due settimane. Il relatore stesso,
        che si impegna moltissimo e tende al burnout, non potrebbe rifare lo
        stesso workshop il lunedì e il martedì successivi con la stessa
        qualità;
  \item \textbf{la differenza tra la miglior performance e la peggiore}: che
        cosa ha portato a lavorare così bene, e così male? Con chi si
        lavorava, su cosa ci si concentrava, quanto tempo e budget c'erano?
        L'obiettivo è capire come \emph{replicare} le condizioni che fanno
        lavorare bene, invece del solito ``hai lavorato male, è colpa tua''.
        Il relatore ha avuto molti manager per cui quando lavorava male era
        colpa sua e quando lavorava bene ``che bravo'', ma sapeva che raramente
        era l'unico fattore della propria performance. È la fallacia del ``se
        sai fare 100, perché non fai sempre 100?'', come nella scena di
        \emph{Schindler's List} sull'operaio che produce una vite in pochi
        secondi: non sempre un ritmo è sostenibile;
  \item \textbf{il fallimento di oggi confrontato con quello di sei mesi fa}.
        Il relatore non è contento della presentazione di oggi (l'avrebbe
        migliorata per altri due anni), ma rispetto a quelle di cinque anni fa
        vede un grande miglioramento di chiarezza: il suo peggior fallimento di
        oggi equivale al suo primo successo di inizio carriera.
\end{enumerate}
Quello che conta è la \textbf{traiettoria} dall'inizio dell'anno a oggi, non
i singoli punti. E se si riesce a capire insieme quando e perché si ha
successo e quando e perché si fallisce, si può lavorare per replicare solo le
situazioni che fanno lavorare bene.

\paragraph{Lo stesso vale per i giochi.} È il motivo principale per cui i
giochi delle grandi case vanno bene o male. Quando si riescono a replicare le
condizioni che hanno prodotto \emph{Warcraft~II}, esce \emph{Warcraft~III};
quando non ci si riesce, esce ``il nuovo \emph{Dragon Age}'': la gente non è
arrabbiata perché il gioco non è bello, ma perché non è come gli altri
\emph{Dragon Age}; non ne replica il \emph{feeling}, anche se potrebbe
essere un gioco favoloso. L'eccezione positiva è il passaggio da \emph{Dune} a
\emph{Dune~II}: il primo era di fatto un'avventura grafica, il secondo ha
generato il genere degli strategici in tempo reale. Gli sviluppatori avevano
capito che il successo non stava nel tipo di gioco, ma nell'ambientazione che
attirava le persone: un bel gioco in quell'ambientazione sarebbe piaciuto
comunque. Il relatore comprò \emph{Dune~II} aspettandosi un seguito
dell'avventura e, pur sorpreso, lo finì.

\subsection{Lying, hiding, faking}

Quando si è valutati su un singolo punto nel tempo, e da quel punto dipendono
bonus, stipendio e il rischio di licenziamento, si comincia a fare
\textbf{lying, hiding e faking}:
\begin{description}
  \item[Lying] dire che è sempre andato tutto bene, che si finirà sempre in
        orario, che non c'è nessun problema;
  \item[Hiding] nascondere: ``non era un mio task'', ``non c'ero'', ``non ho
        visto niente'';
  \item[Faking] fingere, anche in cose apparentemente innocue come ``sì, sì,
        ho capito tutto'' quando si ha solo paura di dire ``me lo rispieghi?
        Non ho capito niente''.
\end{description}

\subsection{``Non ho capito'' e la conferma delle istruzioni}

Dire ``non ho capito, ripetimelo'' è ancora visto come una figura da sciocco.
Il relatore non si arrabbierà mai se qualcuno gli chiede di rispiegare
un'intera slide: le cose ovvie per lui possono non esserlo per altri. Più si
conosce qualcosa, \textbf{più è difficile spiegarla}, perché si danno per
scontate molte cose. Ne parlava proprio con Diego, che gli diceva che fare
l'insegnante diventa sempre più difficile man mano che si diventa esperti. È
come spiegare \emph{EVE Online}, uno dei giochi più complessi in circolazione:
chi lo conosce bene dà per scontato che ``ovviamente non puoi fare questa
cosa'', mentre per un nuovo giocatore non è affatto ovvio. A scuola chi dice
``non ho capito'' viene visto come lo ``scemo''; nessuno si chiede ``non è che
l'ho spiegato male?''.

Il racconto più ``traumatico'' dei ragazzi che fanno stage nelle grandi
aziende è sempre lo stesso: ``Cinzia, devi fare questa cosa.'' Cinzia la fa.
``Non te l'avevo chiesta così.'' Il problema è che nessuno ha
\textbf{confermato che Cinzia avesse capito}. Il consiglio, valido sempre:
\begin{itemize}
  \item quando si \emph{ricevono} istruzioni, ripeterle: ``ho capito bene?
        Vuoi che faccia questa cosa in questo modo?'';
  \item quando si \emph{danno} istruzioni, chiedere a chi le riceve di
        ripeterle, non perché lo si ritenga incapace, ma per capire che cosa
        ha capito dalla propria spiegazione (``quindi vuoi un quadrato?''
        ``No, ho detto un triangolo\dots'' ``Però mi hai detto una forma con
        quattro angoli.'' ``Hai ragione: voglio un quadrato'').
\end{itemize}
Frasi come ``non te lo ripeto due volte'' o ``stammi dietro'' portano solo a
fare le cose come si pensa che vadano fatte e a sentirsi dire a fine progetto
``non la volevo così'': che si tratti di una feature, di un design o di un
modello 3D. (Un partecipante osserva che lo stesso accade insegnando le lingue:
chi insegna dà per scontato, ad esempio, che si sappia dell'esistenza del
genere neutro in tedesco.)

\paragraph{Come convincere il team a chiedere?} Il relatore mette sempre come
premessa che \textbf{l'obiettivo non è trovare la persona più intelligente
nella stanza, ma fare un buon lavoro}. Che Gianni capisca alla prima e un altro
alla quinta non gli interessa, se alla fine entrambi lavorano bene: ripetere
cinque volte è comunque meno lavoro che far rifare le cose sette volte. È un
investimento a lungo termine, e non dissimula alcun fastidio: davvero non gli
dà fastidio. Ma ripetere le stesse parole non aiuta: ``non ho sentito'' si
risolve ripetendo; ``non ho capito'' si risolve cambiando spiegazione. Meglio
ancora: ``dimmi che cosa hai capito tu, e vediamo cosa non torna''. Se
intravedesse nel team uno sguardo del tipo ``è la quinta volta che lo
chiede'', interverrebbe: spesso chi giudica è proprio chi dice ``ho capito'' e
poi sbaglia. Preferisce di gran lunga chi chiede cinque volte.

\subsection{Il feedback più utile che abbiate ricevuto}

\emph{Qual è il feedback più utile che abbiate mai ricevuto?} Il relatore ne ha
almeno dieci che gli hanno cambiato il modo di lavorare. Le risposte:
\begin{itemize}
  \item un semplice \textbf{``grazie''} per un lavoro fatto, più utile di un
        ``questa cosa fa schifo'';
  \item un feedback sulla propria \textbf{impulsività} nel rispondere: fermarsi
        un attimo e capire di cosa si sta parlando prima di intervenire;
  \item \textbf{essere se stessi}, anche al lavoro, e non scusarsi
        continuamente per ogni cosa.
\end{itemize}
I feedback che toccano di più sono quelli \textbf{veri}. Il ``grazie''
significa ``ho visto quello che hai fatto e l'ho apprezzato''. Quelli negativi
che si ricordano sono quelli che hanno fatto realizzare ``questa cosa che
pensavo di saper fare, non la so fare'', oppure ``la so fare, ma questa volta
non mi è venuta'' (non è che se una cosa una volta non riesce non la si sa più
fare: può capitare di sbagliare una porta). Arrivano sempre da un
\textbf{punto cieco}: qualcosa che non si sapeva fosse percepito dagli altri.
Conta più ciò che le persone ti dicono di ciò che pensi che pensino di te; e
ciò che il relatore pensa che gli altri pensino di lui è sempre terribile,
mentre il loro giudizio si rivela molto più positivo del previsto.

\paragraph{Il denominatore comune siete voi.} Si potrebbe pensare che questi
concetti non siano applicabili da chi arriva in un'azienda senza poterne
cambiare la cultura. Ma c'è sempre un denominatore comune a cui applicarli:
\emph{se stessi}. La cultura aziendale può non spingere a dare feedback, può
vedere male chi chiede aiuto; ma il relatore ha visto che quando lui lo fa, le
persone intorno tendono a farlo, e un gruppo che lo fa lavora meglio. Perfino
i CEO più ostici a un certo punto gli chiedevano: ``perché il tuo team ora
performa meglio degli altri?''. La differenza è il default to trust: una
pratica che col tempo fa operare il team \emph{come un team} e non come un
insieme di individui. Un videogioco è un progetto corale: c'è chi fa giochi da
solo, ma non si fa da soli \emph{Red Dead Redemption~2}; e se il gruppo non è
affiatato e allineato non esce \emph{Red Dead Redemption~2}, esce
\emph{Anthem} (``oggi la mia fissazione è \emph{Anthem}, domani sarà
un'altra'').

% ---------------------------------------------------------------------
\section{Come dare un feedback}
\label{sec:z1-dare}

\subsection{Concentrarsi sull'azione e sui suoi effetti}

Il primo principio: \textbf{il feedback non riguarda mai la persona}. Non si
dice ``tu hai fatto questa cosa che non va bene'', ma ``questa cosa non va
bene: come facciamo a non farla più?''. Che sia un comportamento, un'azione,
una decisione o una strategia, ci si concentra su quello, e si
\textbf{esplicita l'effetto} che ha. Nell'esempio del ritardo l'effetto era
``quindici minuti di attesa moltiplicati per venti persone''. In un progetto
può essere una catena:
\begin{quote}
  se non rispetti le deadline, l'animazione che viene dopo resta ferma; se
  loro sono in ritardo, la \emph{vertical slice} è in ritardo; se la vertical
  slice non è pronta per la Gamescom, non andiamo alla Gamescom e non troviamo
  un publisher; se non troviamo un publisher, chiudiamo.
\end{quote}
Più la conseguenza è grave, più dovrebbe sembrare piccolo il gesto di dire
``datemi una mano'': è il \emph{male minore}. Poi, a cose fatte, ci si dà
feedback: cosa ha funzionato, cosa no, cosa fare di più, di meno, cosa smettere
di fare.

\subsection{Il pericolo degli ambienti solo positivi}

Gli ambienti dove ci si fa sempre grandi complimenti e nessuno parla dei
problemi creano difficoltà vere alle aziende. Il problema non è mai la
mancanza di complimenti, ma un \textbf{eccesso di ottimismo}: si comincia a
pensare che le cose andranno come si spera, perché tutti intorno danno solo
feedback positivi.

A Blizzard i feedback determinavano gli \textbf{aumenti annuali}: ognuno
doveva indicare cinque persone che avrebbero dato feedback su di lui, e
ovviamente tutti indicavano gli amici, o colleghi di progetti andati bene. Un
sistema nato con le migliori intenzioni diventava uno \emph{scambio di voti}
(``se mi metti follow ti metto follow''), senza alcun miglioramento.

Al contrario, \textbf{le conversazioni più brutte sono state le più utili}:
affrontare i licenziamenti a Riot, o la cancellazione di \emph{Project Titan}
a Blizzard, il progetto che doveva essere il nuovo MMO dopo \emph{World of
Warcraft} e da cui poi è nato \emph{Overwatch}\footnote{Nella trascrizione si
legge anche ``doveva essere Warcraft 2'', probabilmente un lapsus o un errore
di trascrizione.}. Andando a vedere perché lo si cancellava, c'erano solo
motivazioni razionali e pragmatiche, di business, chiare a tutti, non un ``Mike
Morhaime ci odia''. Sono anche le conversazioni in cui ha lavorato su se
stesso: è molto \textbf{avverso ai conflitti}, e questo a volte gli ha
procurato conseguenze molto più gravi di quanto sarebbe costato un piccolo
disagio. Dieci anni fa faceva fatica ad accettare questo difetto; ha dovuto
subire le conseguenze di conflitti evitati e poi cresciuti.

\subsection{Che cosa fare}

Le linee guida della slide, che il relatore legge rapidamente perché saranno
disponibili nel materiale:
\begin{description}
  \item[Attenzione al linguaggio del corpo] dare un feedback a braccia
        incrociate o guardando altrove è diverso dal guardare l'altro negli
        occhi.
  \item[Sii specifico] con esempi tangibili.
  \item[Esplicita gli effetti] come visto sopra.
  \item[Discuti il prima possibile] non aspettare un mese per parlare di
        qualcosa successo ieri, altrimenti nessuno ricorderà più bene cosa sia
        accaduto.
  \item[Concentrati sull'azione e sul problema, non sulla persona] il nemico è
        l'azione, non la persona che si ha davanti. Non bisogna far cambiare
        l'altro: bisogna evitare che il problema si ripresenti.
\end{description}

\paragraph{Gli insulti sono rivolti al ruolo.} Un esempio solo apparentemente
scollegato. Il relatore ha lavorato su \emph{League of Legends}, ``non
esattamente noto per avere una community da orsetti del cuore'', in varie
\emph{task force} sulla tossicità. A chi smetteva di giocare veniva inviato un
sondaggio sul perché; il 33\% indicava la tossicità della community. Dopo varie
campagne di sensibilizzazione e cambiamenti si è scesi al 30\%: sembra solo
un 3\% in meno, ma valeva circa \textbf{120 milioni di revenue in più
all'anno}, perché quei giocatori restavano esposti a skin e contenuti e
continuavano a comprare. Anche qui, piccoli guadagni che si sommano.

Lavorando con degli psicologi, aveva notato che gli insulti \textbf{non sono
quasi mai rivolti alla persona, ma al ruolo}: ``il mio jungler fa schifo'',
``il mio DPS non è capace'', ``il nostro portiere è una pippa''. Chiamare
qualcuno per ruolo lo \emph{depersonalizza} e rende più facile insultarlo. Lo
stesso nell'esempio del giro dell'ufficio: non ``non chiedere a Giulio'', ma
``non chiedere a quel grafico, consegna sempre in ritardo''. Concentrarsi
sull'azione, senza attaccare la persona, fa sì che il feedback non sia (e non
venga percepito come) un attacco personale.

\paragraph{Considera il carattere.} Qui tornano utili i test di personalità.
Servono più agli studenti che al relatore (che sa già che usciranno 53 profili
diversi, senza un pattern): bisogna conoscere il proprio profilo e
\textbf{condividerlo con chi lavora con noi}, chiedendo anche agli altri di
fare lo stesso. In ufficio il relatore ha usato una versione più avanzata, a
pagamento, con un report di circa trenta pagine: il \textbf{profilo
DISC}\footnote{Nella trascrizione ``disk''. Si tratta del modello DISC
(Dominance, Influence, Steadiness, Conscientiousness).}. Ha mappato tutta
l'azienda e ha fatto stampare il profilo di ciascuno sulla sua scrivania, così
che tutti sapessero con chi stavano parlando. Capire come ragionano gli altri
(``è introverso, non gliene parlo davanti a tutti'') è una delle cose che gli
sono tornate più utili: \emph{adattare ciò che si vuole dire a chi si ha
davanti}.

\paragraph{Assicurati di non disturbare.} Sotto il profilo di ognuno c'era un
piccolo \textbf{semaforo}: verde, ``parlami pure, interrompimi''; giallo,
``interrompimi solo se è urgente''; rosso, ``mi sto concentrando, non mi
parlare''.

\paragraph{Lascia finire il lavoro.} Durante gli eventi dal vivo il relatore si
era reso conto di quanto fosse grave dare feedback mentre qualcuno sta ancora
lavorando. Se una cosa sta andando male, dirlo in quel momento non aiuta: tra
la pressione di finire e quella del feedback, peggiorano entrambe le cose, e
anche il rapporto. Meglio lasciarla finire, anche male, e parlarne dopo.

\subsection{Che cosa non fare}

\begin{description}
  \item[Non basarti su supposizioni] ``mi hanno detto che ieri in un meeting
        hai fatto\dots'': o c'eri e l'hai visto, o non esiste.
  \item[Non giudicare la persona.]
  \item[Non fornire feedback vago.]
  \item[Non accettare e non fornire feedback di terze parti] Se qualcuno vi
        dicesse ``mi hanno detto che hai fatto un brutto lavoro'', il primo
        pensiero non sarebbe ``cosa ho fatto male?'', ma \emph{``chi te l'ha
        detto? Chi è la vipera che ho in seno? Chi è l'impostore?''}. Siamo
        fatti così. Ogni conversazione che inizia con ``mi hanno detto che'' è
        inutile: l'altro non ascolterà, penserà solo a chi dare la caccia, e
        spesso penserà alla persona sbagliata, creando antipatie tra persone
        che non c'entravano nulla. Quando al relatore dicono ``mi hanno detto
        che lui ha fatto\dots'', risponde ``vai a parlarne con lui''; se gli
        dicono ``mi hanno detto che tu hai fatto\dots'', risponde ``mandami la
        persona che te l'ha detto, altrimenti non ne parliamo''. Non per
        cattiveria: se chiedesse ``che cosa ho detto di preciso a Marco?'' e la
        risposta fosse ``non lo so, non c'ero'', di cosa si starebbe parlando?
        \textbf{Parlate con la fonte}, altrimenti fate come se quel feedback non
        esistesse.
  \item[Non dare consigli non richiesti] come visto, non portare la soluzione:
        la soluzione deve venire da chi riceve il feedback.
  \item[``Non mi piace'' non è un feedback] ``Penso che questo design non
        funzioni perché, visto che il nostro obiettivo è\dots'' ha argomenti;
        ``non mi piace il combat system'' è un'opinione. Anche ``io l'avrei
        fatto diversamente'' non è un feedback.
  \item[Accettare non vuol dire fare come proponi] deve essere chiaro a chi dà
        il feedback: l'aspettativa è essere ascoltati e presi in
        considerazione, ma se la responsabilità è dell'altro, deciderà
        l'altro.
  \item[Non offenderti] se il feedback non viene seguito. Restarci male
        significa pensare di saperne più dell'altro, cioè \emph{venire meno al
        default to trust}. ``Decidi tu'' significa ``mi fido di te, perché sei
        lo specialista, e mi fido che mi ascolterai e farai la cosa migliore
        per tutti'', che potrebbe anche non essere ciò che ho detto io. Se
        voglio che l'altro faccia sempre come dico io, mi sto fidando solo di
        me stesso.
\end{description}

\subsection{Il primo feedback e il debriefing}

La cosa più difficile è dare il \textbf{primo} feedback: è un muscolo da
allenare. Chi non l'ha mai fatto (soprattutto con questa struttura) dovrebbe
provarci, anche su minuzie (``abbiamo concordato di usare WhatsApp nel gruppo:
non potremmo usare Telegram?''), anche solo per vedere come va. Altre
indicazioni:
\begin{itemize}
  \item feedback \textbf{solo uno a uno}, mai davanti ad altri, soprattutto se
        negativo;
  \item \textbf{prepararsi} prima;
  \item fare il \textbf{debriefing del feedback}: dopo, fermarsi a pensare a
        com'è andato. Il relatore, quando media un feedback tra due persone,
        si sente spesso rispondere ``benissimo'' da una e ``malissimo''
        dall'altra: la percezione è molto diversa. Non si fa problemi a tornare
        dalla stessa persona: ``ieri ti ho dato un feedback: com'è andata? Sei
        soddisfatto? Come poteva andare meglio?''. Con le persone con cui si ha
        meno a che fare, questo aiuta molto la volta successiva.
\end{itemize}

% ---------------------------------------------------------------------
\section{Come ricevere un feedback}
\label{sec:z1-ricevere}

Per il relatore è la parte più difficile. Si riassume in quattro passi:
\begin{enumerate}
  \item \textbf{Ascolta}: si parte sempre ascoltando.
  \item \textbf{Chiedi dettagli}, se serve, ma per capire meglio, non per
        difenderti.
  \item \textbf{Accetta, a prescindere}: non si sta decidendo chi ha ragione;
        si sta ascoltando un'opinione. Non c'è nulla da cui difendersi.
  \item \textbf{Agisci}: fatti un piano. Può essere banale (``mi hanno detto
        che interrompo sempre: cerco di non interrompere''), oppure ``devo
        migliorare questa cosa: come?''. Pensa anche a \emph{chi} potrebbe
        aiutarti e a \emph{che supporto} ti serve: il piano non deve essere
        per forza ``devo risolvere tutto da solo''.
\end{enumerate}
Una domanda che molti sottovalutano: \emph{``sono junior in un'azienda: sarò
mica il primo ad avere questo problema?''}. Probabilmente qualcun altro l'ha
già avuto: basta chiedere ``tu come hai fatto?''.

\subsection{Il feedback positivo}

Tutto questo vale per i feedback costruttivi negativi, a cui si pensa di
default; ma esistono anche i \textbf{feedback positivi}, e vanno usati. Quando
si riceve un ``grazie'' o un ``bravo'', la domanda è: \emph{come faccio a
ricreare le stesse condizioni?}

Il relatore ha lavorato anche in aziende dove il feedback era solo negativo.
Il rischio? Si immagini l'impiegato ideale: vuole sempre migliorare, è sempre
pronto a dire ``magari sbaglio io''. Se nessuno gli dice cosa sta facendo
\emph{bene}, spinto dal desiderio di migliorare rischia di cambiare cose che
erano già buone, e di farle peggio. Il rinforzo positivo \textbf{non è solo una
gentilezza}: è un atto pratico, che evita che le persone peggiorino ciò che
sanno già fare bene.

\subsection{Il feedback accidentale}

Esiste infine il \textbf{feedback accidentale}, molto frequente negli uffici:
\begin{itemize}
  \item il ``chiedi a lei, non a quel grafico'' del giro dell'ufficio è un
        feedback negativo su una persona, dato accidentalmente, né richiesto
        né voluto;
  \item a Riot li chiamavano \textbf{drive-by}, come nelle scene dei film in
        cui passa un'auto, si abbassa il finestrino, si spara e si riparte. In
        molte aziende con una ``cultura del feedback'' in cui tutti devono
        dire la loro, si riceve una valanga di commenti del tipo ``comunque il
        sito a me non piace, ciao''. Non si sa perché, non si sa cosa fare di
        diverso, non si sa che rischi ci siano: \emph{a che cosa è servita
        questa conversazione?} A niente.
\end{itemize}

% ---------------------------------------------------------------------
\section{Quando, quanto e con chi}
\label{sec:z1-quando}

\subsection{Frequentemente e recentemente}

\emph{Quando avete ricevuto il vostro ultimo feedback?} Per gli studenti,
qualche giorno fa (nel corso) o qualche mese fa (al lavoro, dopo una
consegna). Come aveva osservato Silvia, il feedback andrebbe dato
spesso.

Il relatore usa l'esempio di una nave da crociera (``la Costa Concordia'') che
deve andare da A a B. La rotta ideale (azzurra) è quasi impossibile: nessun
progetto va esattamente come previsto. C'è la rotta ``panoramica'' (verde), che
gira intorno all'isola d'Elba; ci sono poi rotte (arancione e rosa) che a un
certo punto si allontanano molto e poi correggono bruscamente: qualcuno si è
accorto che si stava sbagliando, ``ma non è quello che avevo chiesto'',
``non è dove dovevamo andare''. Sono i famosi \textbf{sanity check} o
\emph{checkpoint}.

L'ideale è avere \textbf{tanti piccoli check per fare tante micro-correzioni}:
correggere la rotta di un paio di gradi è molto più rapido che farlo dopo
essersi allontanati. Bastano trenta secondi di ``ehi, vado bene?''. Spesso non
lo si fa per paura di sembrare insicuri o poco autonomi, ma per il relatore è
l'opposto: quando qualcuno gli chiede ``vado bene così?'' e lui risponde
``sì'', si alleggerisce anche del proprio carico mentale. Un capo che assegna
un compito e lo vede solo alla fine vive sempre col dubbio (``avranno capito?
Lo staranno facendo giusto? Arriverà in tempo?''); con un check di dieci
secondi al giorno (``tutto ok, Stefano'') è tranquillo, e la persona va avanti
sapendo che quanto fatto finora va bene. Non serve fare ogni giorno il giro del
dipartimento: anche ogni uno o due giorni basta a evitare la rotta verde.

Per il feedback è lo stesso: \textbf{va dato subito}. Dire oggi ``ieri hai
fatto questa cosa'' è molto diverso da dire ``negli ultimi sei mesi hai fatto
questa cosa''. Il relatore si è sentito dire più volte, alla valutazione di
dicembre, ``da maggio hai fatto questa cosa che non andava''. La sua risposta:
``e aspetti dicembre per dirmelo? Mi hai lasciato sbagliare per sei mesi?''. È
una mancanza di rispetto, e la conversazione non può più essere costruttiva.
Per questo non ama le valutazioni annuali: i manager crescono abituandosi a
prendere appunti delle cose negative per tirarle fuori alla valutazione
(``a marzo hai fatto questa cosa'' ``sì, però poi ho smesso'' ``sì, però a
marzo l'hai fatta'').

In sintesi:
\begin{description}
  \item[Frequentemente] riguarda l'intervallo tra un check e l'altro.
  \item[Recentemente] significa non lasciar passare un anno prima di dare il
        feedback, altrimenti non serve più a niente.
\end{description}

\subsection{Quanto dura un feedback}

Un feedback dura \textbf{da trenta secondi a quindici minuti}; raramente di
più. Chi non ne ha mai dato uno dovrebbe sforzarsi di darne almeno uno, magari
positivo per cominciare. Ogni feedback dato o ricevuto è un
\textbf{deposito nel proprio ``fondo fiduciario''} (\emph{trust fund}), che
costruisce la reputazione futura. Il relatore si è costruito la reputazione di
chi dà feedback il più possibile oggettivi e analitici: giusti o sbagliati che
siano, costringono a pensare. Non impone mai il suo punto di vista (``si fa
come dico io''), ma chiede ``tu come faresti? Come posso aiutarti?''. Il
feedback, per lui, serve a \emph{far crescere qualcun altro}.

\subsection{I manager peggiori: ``a che numero sto pensando?''}

I peggiori manager che il relatore abbia avuto erano quelli che volevano che si
facesse come volevano loro, anche quando chiedevano ``tu come la faresti?''.
In realtà la domanda era ``\emph{a che numero sto pensando?}'': se dicevi 5,
``no, dovremmo fare 9''; se dicevi 9, ``no, 5''. Lo stress costante di dover
indovinare cosa vuole il capo nasce dalla mancanza di fiducia (``se mi chiedi
come fare, fammelo fare'').

Lo stesso gli succede con i CEO che lo chiamano per ``sistemare la cultura
aziendale'' e poi dicono ``non sono sicuro, io farei diversamente''. La sua
risposta: ``lo so che faresti diversamente, per questo mi hai chiamato''. Se si
chiama un esperto, bisogna lasciargli fare il suo lavoro: non si va dal
meccanico a cambiare le gomme per poi dirgli ``non dovresti girare un po' di
più quel bullone?''. Chiedere un parere a un esperto significa dargli anche il
potere di scegliere. Il problema è il \emph{micromanagement}, o in generale il
non accettare punti di vista diversi dal proprio (magari perché si è il
fondatore e si pensa di essere l'unico a saper fare le cose).

\subsection{Dare feedback a un superiore egocentrico}

\paragraph{Come si dà un feedback negativo a un superiore con un grande ego?}
Il trucco è \textbf{non puntare sull'ego, ma sull'azione e sui suoi effetti}.
Non ``tu fai così e io mi sento così'': anzi, si alimenta il suo ego
rendendolo il \emph{main character} dell'avventura, l'eroe che risolve il
problema (``tu, grande manager, aiutaci a risolvere questo problema che
affligge noi che facciamo l'operatività''). Non è l'ideale (``io cambierei
manager''), ma pochissimi possono scegliersi i colleghi: il relatore stesso ha
quasi sempre ricevuto team preassemblati (``ecco le chiavi, questa è la tua
macchina''). Bisogna giocare con le carte che si hanno in mano.

\paragraph{Un esempio pratico.} Con un manager che fa micromanagement e mette
sempre ``il becco'' nel proprio lavoro, si può dire: ``se ti fidi di me per
questo lavoro, con le tue capacità potresti dedicarti a questa cosa più
strategica e di lungo termine: pensa quanto sarebbe più utile per l'azienda''.
In sostanza: ``guarda cosa potresti raggiungere se mi lasciassi fare il mio
lavoro''. Non come un ``tu fai il tuo e io faccio il mio'', ma come un
``ti sollevo da questa parte del carico, che so che ritieni importante: lascia
che ci pensi io, così puoi dedicarti ad altro'', che magari gli dà anche più
visibilità.

Si gioca su ciò che quella persona pensa di saper fare, e che probabilmente sa
anche fare. Anche Michael Scott di \emph{The Office}, per quanto preso in giro,
è raccontato come un venditore bravissimo, che ha chiuso i contratti più
importanti per l'azienda, e rispetto ai manager che gli succedono è forse
``uno dei meno peggiori''. Tutti, anche i manager peggiori, hanno punti di
forza: \textbf{giocate su quelli}.

% ---------------------------------------------------------------------
\section{Il feedback come regalo}
\label{sec:z1-regalo}

Per chiudere, il relatore risponde alla domanda d'apertura, che aveva lasciato
apposta in sospeso: che cos'è il feedback per lui?

\begin{quote}
  \emph{Il feedback è un regalo che mi viene fatto da qualcuno che so che vuole
  il mio interesse.}
\end{quote}

Ha costruito un team di persone che si fidano di lui e che non hanno motivo di
sabotarlo o di ferirlo: per questo interpreta anche le cose peggiori con il
default to trust. È un regalo, e cerca di prenderlo in considerazione il più
possibile.

Ci è arrivato partendo da se stesso. Oltre a essere avverso ai conflitti, si
definisce ``avverso alla comunicazione umana'': paradossalmente, una persona
introversa e timida finita a fare comunicazione, marketing e ora HR (``succede
spesso''). Prima di ordinare una pizza al telefono ha dieci minuti di
preparazione mentale, e se l'ordine delle domande cambia va in confusione. Per
lui un feedback di dieci minuti sono spesso \textbf{tre ore di preparazione}:
costruisce un \textbf{albero dei dialoghi}, come in un'avventura grafica
personale, con tutte le opzioni. Poi spesso finisce con ``hai fatto questa
cosa, andrebbe migliorata'' ``hai ragione, grazie'', e basta.

Così ha pensato: chissà quanti si fanno gli stessi problemi con lui, e passano
ore a prepararsi per dirgli ``questa cosa l'hai fatta male'', sapendo che
riceveranno comunque quell'occhiata fugace del ``come ti permetti''. Eppure
vengono a dirglielo. Non può far altro che rispondere: ``grazie, hai ragione,
cercherò di fare meglio''. È così che ha metabolizzato il ricevere feedback:
capendo che non c'è nessuna cattiva intenzione e che, se non proprio felice
(``ci ho messo una pietra sopra''), può almeno essere \emph{grato} che ci sia
ancora qualcuno che ha così a cuore il suo interesse da dargli feedback.

\paragraph{Il silenzio fa più paura.} Non bisogna spaventarsi dei feedback
negativi che emergono dai sondaggi aziendali: bisogna spaventarsi del
\textbf{silenzio}, quando nessuno dice più nulla. Significa che nessuno si
fida più che si sappiano risolvere i problemi.

\paragraph{La sindrome dell'impostore.} Il suo motivatore personale, all'inizio,
era soprattutto ``non voglio fallire'': per questo cercava di migliorare, e
migliorare è diventato accettare feedback, perché più feedback riceve, più
migliora. Il feedback è un \textbf{acceleratore}, e lui cerca di incentivare
gli altri a darglielo. Resta però sempre la vocina del ``e se si accorgono che
non sei capace?''. Con la \emph{impostor syndrome} convive da anni e
continuerà a conviverci; se ce l'avete, conviveteci anche voi. Avrete sempre
la sensazione che ``se ne accorgeranno'': forse sì, ma sapete fare anche
altre cose. Non concentratevi solo su ciò che non sapete fare. Gli errori
capitano: la perfezione non è di questo mondo.

\paragraph{Esercitarsi in accademia.} Per fare pratica con il feedback,
l'ambiente migliore è quello dei colleghi di corso: persone nella stessa
situazione, con esperienza simile e lo stesso obiettivo, in un contesto in cui
non si può essere licenziati e l'unico effetto a lungo termine è ``ho parlato
con questa persona, ora le cose vanno meglio''. E bisogna avere il coraggio di
parlare anche con i docenti (``dire a Diego: secondo me questa lezione potevi
farla in questo modo''; ``so che è una persona che lo accetta''), senza farsi
bloccare dalla gerarchia. Nell'azienda giusta (come quella in cui lavora ora) il
relatore va sempre dal CEO a dare feedback, sapendo che non verrà licenziato:
il CEO gli ha detto ``ti ho preso per darmi feedback''.

La frase di chiusura:

\begin{quote}
  \emph{Per superare le vostre difficoltà, tutto quello che vi viene chiesto è
  il coraggio di chiedere aiuto.}
\end{quote}

Chiedere aiuto non è una debolezza: è un segno di consapevolezza di sé
(\emph{self-awareness}) e di forza personale. Se tutti lo facessero, i
progetti andrebbero molto meglio; spesso invece ci si fa frenare dal ``non
voglio fare brutta figura'', ``non voglio passare per quello scemo'', e sono
proprio queste paure a rallentarci e a impedirci di raggiungere gli obiettivi.

% ---------------------------------------------------------------------
\section{Domande e risposte finali}
\label{sec:z1-qa-finale}

L'ultima mezz'ora è dedicata a domande libere.

\paragraph{E se una persona peggiora per uno o due anni di fila?} Se nessuno
gliel'ha mai detto, c'è un grosso problema, non suo ma dell'azienda. Se invece
c'è consapevolezza, il percorso è: ``cosa vuoi fare per risolvere?''; ``ho
provato questo, questo e questo''; ``visto che non ha funzionato, vuoi un mio
consiglio? Proviamo con quello?''; ``che supporto ti serve?''. Se, con tutto il
supporto, il tempo e le soluzioni concordate per mesi, il miglioramento non
arriva, per il relatore la soluzione è: \textbf{``dimmi un prezzo per andare
via''}. Preferisce pagare una buonuscita (\emph{severance package}) che
avviare la strada dei richiami formali sulla performance fino al
licenziamento per giusta causa. A un certo punto andare avanti non è un favore
per nessuno.

\paragraph{È vero che nei videogiochi si viene sfruttati? Meglio un piccolo
studio italiano?} Per il relatore si è quasi più a rischio a Blizzard che in un
piccolo studio di Roma o Torino. Nel piccolo studio si sa dove ci si va a
cacciare: non ci sarà il team building a Santo Domingo, ma ci si concentra su
un gioco alla volta e, se va bene, si fa il successivo. A Blizzard si può
lavorare su \emph{Overwatch}, che vende milioni, e perdere comunque il posto,
in modo totalmente scollegato dal proprio lavoro. Gli stipendi più alti
all'estero compensano anche questo rischio, che spesso non dipende da come va
il gioco (anzi, ``il gioco non ha fatto come volevamo'' è spesso una scusa per
licenziamenti che non avevano motivo).

\paragraph{Cambiare settore a trent'anni, da web developer full remote?} Lo
stipendio indicato dallo studente non è fuori dalla portata di vari studi
italiani, anzi è superabile. Non va esclusa la possibilità di lavorare da
remoto per aziende piccole ma estere. Alcuni amici del relatore, pagati con
stipendi italiani, si sono trasferiti a Bali o in paesi con un costo della vita
più basso; il collega di \emph{PUBG} percepisce uno stipendio olandese vivendo
a Milano, che per quanto cara non costa come i Paesi Bassi.

\paragraph{Come si diventa manager?} Dipende dall'azienda: alcune promuovono i
migliori (processo tendenzialmente sbagliato), altre propongono un percorso di
\emph{people management}. Diventare manager significa \textbf{abbandonare ciò
che si sa fare per imparare a fare altro}. Il relatore ha dovuto accettare che
non sarebbe più stato ``il miglior community manager italiano del mondo''
(``eravamo in pochi, quindi era facile''): dedicandosi alle persone non poteva
più fare il suo lavoro bene come prima. È stata una negoziazione con se
stesso: fare il manager significa non essere più l'esperto di qualcosa, ma
saper gestire e far crescere persone. Significa anche \textbf{cedere il
controllo della propria carriera}: se il team decidesse di non fare più nulla,
il manager ne pagherebbe le conseguenze. Il suo lavoro diventa scegliere le
persone giuste e farle lavorare bene. Lui lo esplicita sempre: sa di affidare
al team parte della propria carriera, e il team sa che lui combatterà per
loro, anche prendendosi le colpe.

\paragraph{L'inglese.} \textbf{Va assolutamente imparato}, anche lavorando in
Italia. Non saperlo preclude circa il 95\% del mercato (tutto quello fuori
dall'Italia) e soprattutto il \emph{networking} con persone che potrebbero dare
occasioni future: l'inglese è spesso l'unica lingua in comune con chi lavora
all'estero ed è oggi la lingua universale dei videogiochi.
\begin{itemize}
  \item \emph{Che livello serve?} Non importa B2 o C2: bisogna saper sostenere
        una conversazione, guardare almeno Netflix con sottotitoli in inglese,
        comunicare senza dodici pause e senza ``come si dice\dots''. Il
        certificato non verrà mai chiesto: verrà chiesto se si sa parlare
        inglese. Il relatore ha sentito persone con un C2 che non sapevano
        parlarlo.
  \item \emph{Come allenarlo senza occasioni di conversazione?} Netflix, Prime
        Video, Now TV, Udemy, Coursera, YouTube, Duolingo, siti di scambio
        linguistico con madrelingua, Discord: non sono mai esistiti così tanti
        mezzi gratuiti. Il relatore ha imparato l'inglese giocando (con una
        pronuncia sbagliata, perché leggeva senza sentire, su cui ha poi
        dovuto lavorare); guardare un film in inglese con i sottotitoli in
        italiano insegna anche la traduzione.
  \item Oltre all'ascolto passivo è fondamentale \textbf{parlare}: in una
        flotta di \emph{EVE Online} con 150 persone quasi tutte madrelingua
        inglese si è esposti a decine di accenti e modi di dire (mille modi di
        dire ``come stai'' oltre a \emph{how are you}); lavorando all'estero si
        impara a distinguere gli accenti (Manchester, Londra) e ad apprezzarli
        anche nei media.
\end{itemize}

\paragraph{Quanto deve sapere un junior?} Dipende da quanto è tecnico il
ruolo. Un game designer junior che non sa qualcosa può essere affiancato da un
senior; un programmatore che non sa programmare, in quel momento, non è utile
all'azienda.

\paragraph{Annunci ``junior con cinque anni di esperienza''.} ``Cercasi
ventenne con trent'anni di esperienza.'' Non è l'annuncio per un junior: è un
errore dell'azienda, che in realtà cerca un \emph{mid}. Riceverà profili
junior, si stupirà che non abbiano cinque anni di esperienza, rifarà l'annuncio
e avrà perso sei mesi. Il consiglio: tenere d'occhio quella pagina, perché
dopo qualche mese l'annuncio verrà corretto.

\paragraph{Una struttura per il post-mortem.} A chi chiede un formato di base,
il relatore anticipa quanto vedrà nella seconda giornata a proposito di
brainstorming: una lavagna \textbf{divisa in quattro aree}, su cui ognuno
attacca tutti i post-it che vuole:
\begin{center}
  \begin{tabular}{@{}p{0.44\textwidth}p{0.44\textwidth}@{}}
    \toprule
    \textbf{Smettere di fare} & \textbf{Iniziare a fare} \\
    ciò che non ha funzionato & ciò che non abbiamo fatto e andava fatto \\
    \midrule
    \textbf{Fare di meno} & \textbf{Fare di più} \\
    ciò che ha fatto perdere tempo rispetto all'impatto ottenuto &
    ciò che ha funzionato ma, fatto con troppa leggerezza, non ha dato i
    frutti sperati \\
    \bottomrule
  \end{tabular}
\end{center}
Con i post-it ognuno ha voce e tutto risulta catalogato; per il report basta
una foto della lavagna (o uno screenshot, se la lavagna è digitale, come Miro).

\paragraph{Come sbloccare la paralisi del feedback verso il management?} Uno
studente ha visto manager dire ``non abbiate paura, la mia porta è sempre
aperta'' senza che cambiasse nulla. C'è una dissonanza tra ciò che dicono e
ciò che fanno: è ``la mamma che dice vieni qua che non ti faccio niente''.
Qualcosa deve aver fatto pensare alle persone che non è sicuro. Prima va
sistemato quello; poi si può provare con i \textbf{feedback anonimi}.
L'obiettivo del relatore resta il feedback aperto e firmato, ma a Stormind ha
iniziato con feedback anonimi verso il management per verificare se
l'anonimato fosse un fattore. Non lo era (evidentemente non c'erano problemi a
parlare con il CEO), e la volta successiva intende abbandonarli. Ma è un buon
\emph{check}: se senza anonimato non arriva nulla e con l'anonimato arrivano
cento feedback, il problema è chiaro.

\paragraph{Quando è giusto rivolgersi a un superiore?} La strada giusta resta
parlare direttamente con l'interessato. Ma se si prova e non si trova una
soluzione, la cosa più ovvia è andare dal proprio manager insieme: ``ci
aiuti?''. Quando il relatore parlava di mediare i dialoghi, era proprio questo:
fare da moderatore, aiutare a sbloccare i punti fermi o, se non si riesce a
scegliere tra A e B, scegliere lui. L'escalation, quando necessaria, va fatta.

\paragraph{Il tempo libero veniva rispettato?} No, e se l'è dovuto imporre
lui. Il lavoro gli piaceva così tanto che continuare a casa non gli sembrava
lavorare. Gli si è ritorto contro con il lancio di \emph{Diablo~III}, fatto in
orario americano (le due o l'una di notte in Europa), con l'ufficio chiuso
perché a nessuno era stato chiesto di restare. Il lancio va malissimo, i
giocatori protestano, e lui, sveglio, risponde sui forum fino alle tre, quattro,
cinque del mattino: ``non posso risolvere il problema, ma resto qui ad
aggiornarvi con quello che riesco a sapere dal team di sviluppo''. Il giorno
dopo trova sulla scrivania un \textbf{richiamo formale dell'HR}: aver lavorato
fuori orario senza autorizzazione, condotta non professionale. Per lui è stato
un interruttore: ``ora me ne vado''. Due settimane dopo faceva colloqui con
Riot. Ancora oggi è contrario a quel richiamo, ma gli ha fatto capire come
ragiona un'azienda: ``hai un contratto 9--18, non mi interessa se l'hai fatto
con le migliori intenzioni''. Eppure per anni ha goduto di un'ottima
reputazione in quelle community come ``quello che è rimasto con noi fino alle
cinque del mattino'', e anche l'azienda ne ha beneficiato.

Ci sono poi aziende in cui, soprattutto per i developer, viene richiesto il
famigerato \textbf{crunch}: lavoro straordinario continuativo, senza pause.
Dipende molto dall'azienda; sta venendo (``grazie al cielo'') deprecato, perché
molte aziende capiscono che non è sano ed è visto male anche dal pubblico, ma
esiste ancora.

\paragraph{E se il tuo capo demotiva qualcuno del tuo team?} Qui ``diventa
personale'': il team del relatore si fida di lui perché sa che in questi casi
``diventa una bestia''. ``Te la puoi prendere con me quanto vuoi, ma il mio
team non lo tocchi.'' Lo gestisce direttamente con l'altro manager, per quanto
brutta possa diventare la conversazione. Lo sforzo più grande è \textbf{non
andarne poi a parlare con il team} (``ragazzi, ho fatto questa cosa''), perché
minerebbe la fiducia verso l'alto: non vuole un rapporto ottimo fino a lui e
pessimo oltre. L'obiettivo è che l'altro manager ``si faccia perdonare''
cambiando il proprio comportamento, e che il team percepisca da solo il
cambiamento.

\paragraph{Chiusura della giornata.} Il relatore chiede un feedback collettivo
in chat (``se domani ci dai \emph{more of the same} siamo contenti, oppure
ditemi cosa limare''), si scusa per i ripetuti crash della connessione e
ringrazia per la proattività nelle domande, che gli servono anche come
\emph{check}: se nessuno chiede nulla su un argomento, significa che può
andare avanti. Racconta di aver passato una settimana a decidere che cosa
mettere nelle slide e di averle poi preparate in poco tempo (rifacendone
alcune fino alle tre di notte). Ricorda infine di usare il test di
personalità per sé e per il proprio gruppo, facendolo fare a chi non l'ha
ancora fatto: tornerà molto utile per cooperare meglio.

% ---------------------------------------------------------------------
\section{Anticipazioni sulla seconda giornata}
\label{sec:z1-domani}

La seconda giornata del workshop sarà più pratica: ``oggi vi ho spiegato
perché, domani vi insegno come lo fanno queste aziende''. Gli argomenti
previsti:
\begin{description}
  \item[Time management] come organizzarsi una giornata professionale in
        azienda: come strutturare una \emph{task list}, cosa fare e cosa non
        fare, come e cosa prioritizzare.
  \item[Brainstorming] i gruppi dovranno fare moltissimi meeting (come facciamo
        il gioco, come lo chiamiamo, che scopo gli diamo). Se non si sanno fare
        bene, i brainstorming sono \textbf{dannosi ed eterni}: un brainstorming
        dovrebbe durare \textbf{al massimo due ore}. Verranno mostrati metodi
        che accorciano la durata, massimizzano l'effetto e, soprattutto,
        \textbf{fanno partecipare tutti}: non c'è niente di peggio di un
        brainstorming in cui parlano in due per il 90\% del tempo. Il classico
        ``ragazzi, non esistono idee sbagliate, sparate'' è il modo peggiore di
        farlo; non è colpa di nessuno, ma si può fare diversamente.
  \item[Guidare da dentro] come essere leader nel proprio team anche senza
        esserne il leader: sviluppare responsabilità personale e verso gli
        altri. Non essere il manager di qualcuno non significa non esserne
        responsabili: la persona accanto è responsabilità di tutti, perché si
        è un team e non un insieme di individui. ``Se qualcuno cade, tutti lo
        facciamo rialzare: non c'è l'addetto al rialzamento.'' È la versione
        \emph{advanced} della giornata: come comportarsi da manager senza
        averne il titolo.
\end{description}
Torneranno inoltre il \emph{service level agreement} e il tema della
condivisione dei successi e della protezione del team dai fallimenti.

% ---------------------------------------------------------------------
\section*{Punti chiave del capitolo}
\addcontentsline{toc}{section}{Punti chiave del capitolo}

\begin{itemize}
  \item ``La capacità di un gruppo di persone di fare cose straordinarie
        dipende dal modo in cui queste persone operano insieme.'' Il 90\% dei
        problemi dei team sono \textbf{problemi di comunicazione}.
  \item Le aziende si liberano sempre più presto dello \textbf{useless
        genius}: le competenze si trasmettono, il carattere no. Il
        \textbf{fit culturale} va valutato anche dal candidato.
  \item Non esiste un unico modo di dare feedback: dipende dal contesto e
        soprattutto da \textbf{chi lo riceve}. Il feedback è una somma di
        piccoli miglioramenti periodici; tenersi le cose nel cassetto non
        serve mai.
  \item I team attraversano \textbf{forming, storming, norming, performing e
        adjourning}. Lo storming non si salta; si passa la maggior parte del
        tempo in forming; la bravura sta nell'accelerare i passaggi e nel
        fare bene il \textbf{post-mortem}, anche dei successi.
  \item Gli obiettivi di un progetto vanno fissati \textbf{all'inizio}: solo
        così si può valutare se è andato bene e scegliere tra A e B.
  \item Il circolo virtuoso è \textbf{feedback $\rightarrow$ miglioramento
        $\rightarrow$ fiducia}. Accettare un feedback non significa fare come
        dice l'altro: decide chi ha la responsabilità, ma deve argomentare.
  \item ``Aiutami a poterti aiutare'': meglio dire in anticipo ``potrei non
        farcela'' che ``non ce l'ho fatta''. Chiedere aiuto non è una
        debolezza.
  \item Quando si dà feedback \textbf{non si porta la soluzione}: la si
        costruisce insieme a chi lo riceve.
  \item La cultura è forte quando si lavora \textbf{per gli altri}; il lavoro
        è valutato come opera d'insieme (anche con bonus di team).
  \item Come in \emph{Among Us}, i team sprecano energie a cercare il
        traditore invece di fare le task. \textbf{Default to trust}: chi ha
        sbagliato ha agito con le informazioni che aveva. Cercare
        \textbf{la causa, non il colpevole}: ``che cosa ti ha fatto pensare
        che fosse la scelta giusta?''.
  \item Carriere e performance non sono lineari: vanno valutate le
        \textbf{traiettorie}, non un punto dell'anno. Le valutazioni puntuali
        generano \emph{lying, hiding, faking}.
  \item Confermare sempre le istruzioni facendosele ripetere; ``non ho
        capito'' non è una colpa.
  \item Dare feedback: sull'\textbf{azione} e sui suoi \textbf{effetti}, specifico,
        presto, uno a uno, adattato al carattere, senza terze parti, senza
        ``non mi piace'', senza offendersi se non viene seguito; poi fare il
        debriefing.
  \item Ricevere feedback: \textbf{ascolta, chiedi dettagli, accetta, agisci}.
        Il feedback positivo serve a non peggiorare ciò che si fa bene; il
        feedback accidentale e il \emph{drive-by} non servono a nulla.
  \item Feedback \textbf{frequenti e recenti}: tanti piccoli \emph{sanity
        check} costano meno di una grande correzione di rotta.
  \item Il feedback è \textbf{un regalo} di chi ha a cuore il nostro
        interesse; fa paura il silenzio, non le critiche.
\end{itemize}

%    % =====================================================================
%  Lezione zero 2 — Gestire il tempo, generare idee, guidare da dentro
%  Fonte: trascrizione della seconda "lezione zero" (GD0_2), seconda e
%  ultima giornata del workshop di due giorni tenuto prima dell'inizio del
%  corso di Game Design, dedicato alla figura professionale e al lavoro in
%  team. Prosegue direttamente il capitolo 0.1 (cultura del feedback).
%  Pacchetti richiesti nel preambolo del libro:
%    \usepackage[italian]{babel}, \usepackage[T1]{fontenc},
%    \usepackage[utf8]{inputenc} (solo pdfLaTeX), \usepackage{booktabs}
%
%  Collocazione nel libro: parte "Lezioni zero", subito dopo
%  GD0_1_cultura_del_feedback, con la stessa numerazione propria (0.2).
%  Nel file principale del libro, ad esempio:
%
%    \part{Lezioni zero: il game designer nell'industria}
%    \renewcommand{\thechapter}{0.\arabic{chapter}}
%    \setcounter{chapter}{0}
%    \input{Seminari/GD0_1_cultura_del_feedback}
%    \input{Seminari/GD0_2_tempo_idee_e_leadership}
%    \part{Il corso di Game Design}
%    \renewcommand{\thechapter}{\arabic{chapter}}
%    \setcounter{chapter}{0}
%    \input{GD1/GD1_introduzione_al_game_design}
%    % ...
%
%  (Con hyperref aggiungere anche \renewcommand{\theHchapter}{z.\arabic{chapter}}
%  nella parte zero, per evitare ancore duplicate.)
%  Tutte le etichette di questo capitolo hanno il prefisso "z2:"; i
%  rimandi al capitolo precedente usano le etichette "z1:" già definite in
%  GD0_1_cultura_del_feedback.tex.
% =====================================================================

\chapter[Tempo, idee e leadership]{Gestire il tempo, generare idee, guidare da dentro}
\label{cap:z2-tempo-idee-leadership}

Questo capitolo raccoglie la \textbf{seconda giornata} del workshop di due
giorni con cui si aprono le lezioni zero, tenuto dallo stesso relatore del
capitolo~\ref{cap:z1-feedback}, Stefano Cozzi\footnote{Per il cognome e per
il nome dello studio (Stormind Games) valgono le avvertenze della nota
all'inizio del capitolo~\ref{cap:z1-feedback}. In questa giornata la
trascrizione riporta soltanto il nome di battesimo, ``Stefano'', con cui il
relatore si cita più volte nei suoi esempi.}. Se la prima giornata spiegava
\emph{perché} le aziende tengono tanto alla cultura del feedback, la seconda
è dichiaratamente più pratica: ``ieri vi ho spiegato perché, oggi vi insegno
come lo fanno queste aziende''. I concetti della prima giornata (il
\emph{default to trust}, il ``potrei non farcela'' detto in anticipo, la
condivisione dei successi, il feedback frequente) ritornano di continuo,
applicati alla gestione del tempo, alle riunioni creative e alla leadership.

La giornata si è svolta di domenica, a chiusura di un fine settimana di
workshop\footnote{Il relatore dice ``oggi chiaramente è domenica'' e saluta
la classe con ``buon resto del weekend''; la prima giornata si era quindi
svolta di sabato. Più avanti parla di ``dodici ore'' di workshop in totale,
mentre nel capitolo~\ref{cap:z1-feedback} si parlava di circa dieci ore: la
differenza dipende probabilmente dal conteggio delle pause.}, con il
relatore alle prese con la voce un po' calata dal giorno prima (``la sfida
principale di oggi sarà contro la mia voce'').

Il percorso del capitolo segue quello della lezione:
\begin{enumerate}
  \item il programma della giornata (sezione~\ref{sec:z2-programma});
  \item le domande sulla prima giornata, che si trasformano in una lunga
        sessione su \textbf{colloqui e selezione}
        (sezione~\ref{sec:z2-colloqui});
  \item il \textbf{time management}: pianificare la giornata, le liste, le
        checklist, le interruzioni (sezione~\ref{sec:z2-pianificare});
  \item la \textbf{matrice di Eisenhower} per prioritizzare
        (sezione~\ref{sec:z2-eisenhower});
  \item come proteggere il proprio tempo: sprechi, distrazioni, il saper dire
        di no (anche ai meeting), il mito del multitasking e la chiusura della
        giornata (sezione~\ref{sec:z2-proteggere});
  \item le domande sul time management e una parentesi su tempo libero, ferie
        e benessere in azienda (sezioni~\ref{sec:z2-qa-tempo}
        e~\ref{sec:z2-benessere});
  \item il \textbf{brainstorming}: obiettivi, problemi tipici, un metodo
        completo (\emph{rapid ideation} e \emph{dot voting}) e la
        \textbf{matrice impatto--sforzo} per decidere
        (sezioni~\ref{sec:z2-brainstorming}, \ref{sec:z2-metodo}
        e~\ref{sec:z2-impatto});
  \item i consigli pratici, le domande sul brainstorming e il principio del
        \textbf{disagree and commit} (sezioni~\ref{sec:z2-qa-brainstorming}
        e~\ref{sec:z2-disagree});
  \item il \textbf{leading from within}: perché la leadership non è un
        titolo, le carriere divise tra manager ed esperti
        (sezione~\ref{sec:z2-leadership});
  \item il ``momento Piero Angela'' sulla biologia del gruppo e le domande
        che ne sono nate (sezioni~\ref{sec:z2-biologia}
        e~\ref{sec:z2-qa-biologia});
  \item la \textbf{piramide della fiducia}, il prendersi cura del gruppo e
        che cosa cambia quando si diventa manager
        (sezioni~\ref{sec:z2-piramide}, \ref{sec:z2-gruppo}
        e~\ref{sec:z2-manager});
  \item le domande finali e la chiusura del workshop
        (sezioni~\ref{sec:z2-qa-finale} e~\ref{sec:z2-chiusura}).
\end{enumerate}

Come per il capitolo precedente, una premessa di metodo: buona parte della
giornata è fatta di domande della classe e di racconti dell'esperienza del
relatore. Sono riportati con cura perché sono il modo in cui i principi
vengono resi concreti, e perché spesso contengono le indicazioni più
direttamente utilizzabili (soprattutto sui colloqui di lavoro).

% ---------------------------------------------------------------------
\section{Il programma della giornata}
\label{sec:z2-programma}

Come annunciato alla fine della prima giornata
(sezione~\ref{sec:z1-domani}), gli argomenti sono due:
\begin{itemize}
  \item un ``doppio topic'', \textbf{time management e brainstorming};
  \item il \textbf{leading from within}, che il relatore definisce ``la
        versione avanzata di ieri'': come guidare un team stando dentro al
        team, verso i colleghi e verso il gruppo con cui si lavora.
\end{itemize}
Poiché l'ordine è indifferente, il relatore lascia scegliere alla classe: dalle
poche preferenze espresse, la leadership va al pomeriggio, così che le domande
finali si concentrino su quella, mentre la mattina è dedicata a time
management e brainstorming. Il materiale c'è, ma il relatore avverte che la
giornata potrebbe finire prima delle quattro: non vuole trattenere nessuno se
non c'è niente di specifico da fare. Per l'eventuale tempo avanzato prima della
pausa pranzo propone un \textbf{argomento bonus}: i colloqui di lavoro,
selezione e \emph{recruitment}, visto che il relatore ha fatto e fa tuttora
selezione del personale e che per molti studenti sarà la cosa più immediata da
affrontare finito il percorso di studi (domande comuni, come rispondere,
eccetera).

Prima di cominciare, però, il relatore apre a \textbf{domande sulla giornata
precedente}: la classe ha avuto modo di ``metabolizzare'' i contenuti e
potrebbero essere emersi dubbi a posteriori. È proprio da una di queste domande
che nasce, in anticipo, la parte sui colloqui.

% ---------------------------------------------------------------------
\section{Colloqui di lavoro: capire che cosa cerca l'azienda}
\label{sec:z2-colloqui}

\subsection{L'annuncio come prima fonte di informazione}

La domanda: il giorno prima si era detto che è importante capire che cosa
l'azienda cerca davvero; ma se il primo colloquio è con le risorse umane, cioè
con qualcuno che magari ``di tecnico non sa'', come si fa a capirlo? Il
relatore risponde subito, nel dubbio che più tardi manchi il tempo.

La prima fonte è \textbf{l'annuncio}. Non viene scritto a caso: è approvato
sia dall'HR sia dall'\emph{hiring manager}\footnote{Nella trascrizione
``Airing Manager'' e ``Iron Manager''; allo stesso modo ``e-char'' e ``char''
stanno per HR.}, cioè la persona che seguirà il candidato una volta assunto,
e spesso dall'intero dipartimento HR. È calcolato e scritto in un certo modo
per far capire che cosa serve. Molti annunci moderni, oltre alla parte su
\emph{che cosa si verrà a fare}, contengono una parte su \emph{che cosa ci si
aspetta da questa persona}, cioè il comportamento e il carattere che servono
in quella posizione: ed è quella che dice molto sulla \textbf{cultura
aziendale}.

L'esempio è quello di due annunci per la stessa posizione (programmatore
senior) descritta in modo diverso:
\begin{enumerate}
  \item ``cerchiamo un programmatore esperto che possa prendersi cura dei
        nostri programmatori junior, aiutarli a crescere, sviluppare un
        ambiente di lavoro inclusivo e stimolante, in modo che i nostri
        progetti futuri siano all'avanguardia'';
  \item ``cerchiamo un programmatore esperto che sappia rispettare le
        deadline, seguire le indicazioni dei dati, sia indipendente nel suo
        lavoro e sappia relazionarsi con gli \emph{stakeholder} per capire e
        proporre nuove direzioni per il progetto''.
\end{enumerate}
Nel primo caso ci si aspetta un team di junior, probabilmente con un'età media
più bassa, con cui ci si relazionerà soprattutto; per un senior è l'occasione
di sviluppare le doti da \emph{people manager}, e al colloquio arriveranno
domande come ``hai mai lavorato con dei team? Con che età media? Ti sei mai
trovato a guidare un gruppo?''. Nel secondo caso il senior deve eseguire un
piano deciso o proporre qualcosa, e si relazionerà soprattutto \emph{verso
l'alto}; al colloquio arriverà piuttosto ``fammi questo test di
programmazione''. Sono, di fatto, due colloqui completamente diversi.

\subsection{Quanti colloqui e con chi}

Ormai è molto raro essere assunti con un solo colloquio: ne servono almeno
due, realisticamente \textbf{tre}, che è diventata la media. L'iter tipico:
\begin{enumerate}
  \item il primo con l'HR;
  \item il secondo con l'HR o con il futuro manager;
  \item il terzo sicuramente con il futuro manager, perché l'azienda vuole far
        conoscere al candidato la persona con cui lavorerà, e viceversa.
\end{enumerate}
Già le domande diverse dei vari interlocutori aiutano a capire che cosa
cercano. Ma c'è una parte fondamentale, che secondo il relatore molti
sottovalutano: il momento finale delle \textbf{``hai domande?''}.

\subsection{Le domande da fare al selezionatore}

Un principio di base: \textbf{i colloqui non sono ``devo farmi selezionare
dall'azienda'', sono ``dobbiamo selezionarci a vicenda''}. Le domande vanno
quindi preparate: ci sarà sempre qualcosa che, dalla \emph{job description} o
dal colloquio, non è chiaro al cento per cento. Nel dubbio, il relatore
suggerisce tre domande ``che funzionano sempre'':

\paragraph{``Se iniziassi domani, quale sarebbe il mio \emph{day to day}?''}
Ottiene due effetti. Il primo è pratico: descrive la giornata tipo e fa capire
se piacerebbe o no. Il secondo è più psicologico, ma il relatore sa che
funziona: costringe l'interlocutore a pensare al candidato \emph{già dentro
l'azienda}, cioè a che cosa potrebbe fare da subito e a quale impatto potrebbe
avere con le sue competenze.

\paragraph{``Se aveste una bacchetta magica, che cosa cambiereste subito
dell'azienda?''} Sposta la conversazione su di loro e va bene sia con l'HR sia
con l'hiring manager. Li obbliga a indicare una sfida interna che vorrebbero
risolvere, ed è sempre un'informazione utile: se entrambi rispondono che
cambierebbero il modo in cui si prendono le decisioni, viene da pensare che non
ci sia un percorso chiaro e trasparente da parte di chi sta in alto.

\paragraph{Lo scambio: ``vi do un motivo per non prendermi, voi datemene uno
per non venire''.} È più avanzata e il relatore non la usa sempre, ma la trova
particolarmente interessante. Permette al candidato di esprimere i propri dubbi
in modo costruttivo (``mi dite che lavorerei a stretto contatto con un team di
junior: un motivo per non prendermi è che non ho anni di esperienza nella
gestione di giovani sviluppatori, sarebbe la mia prima volta; però penso di
poterlo fare perché ho questo interesse e queste competenze''). Dall'altra
parte, il selezionatore dirà che cosa secondo lui sarebbe più difficile per il
candidato: si scopre così che cosa sta già pensando, e l'informazione torna
utile anche per i colloqui successivi.

\paragraph{E se il selezionatore si innervosisce?} ``Se un recruiter si
innervosisce per le domande del candidato, ha sbagliato lavoro.'' Al massimo il
relatore si innervosisce se un senior gli chiede che cosa faccia di preciso un
programmatore senior (``lo fai da dieci anni, lo saprai tu''); ma se un junior
chiede che cosa dovrebbe fare ogni giorno, lo si spiega volentieri. E in realtà
vale anche per i senior, perché \textbf{in ogni azienda lo stesso ruolo fa cose
diverse}: conviene sempre farsi spiegare che cosa intendono per quel titolo.
L'esempio viene dal mondo delle community: \emph{community manager} a Riot
indicava la persona a capo del team di \emph{community representative} di una
nazionalità (``community manager italiano''); a Blizzard indicava la persona
sopra l'intero dipartimento che gestiva quei team. Candidarsi per lo stesso
titolo non garantisce di aver capito per che cosa ci si sta candidando, e i
recruiter lo sanno.

Se poi un recruiter si innervosisce comunque, sta dando più risposte di quelle
che pensa: vuol dire che l'azienda non vuole domande, o non vuole interazioni
del tipo ``ma perché facciamo questa cosa?''. Che abbia formato i suoi
recruiter a innervosirsi, o che accetti che lo facciano, è informativo.

\subsection{Le domande che dicono molto sull'azienda}

Allo stesso modo, come detto il giorno prima, è informativo se un recruiter
chiede a una donna se vuole avere figli: ``ah, quindi a voi interessa sapere se
voglio avere figli: questo già mi dice che probabilmente non voglio lavorare
qui''.

Un esempio \emph{gender neutral}: la richiesta di vedere la \textbf{busta paga
precedente}. ``Quanto è la tua RAL? Mandami una busta paga.'' La risposta del
relatore è no: ``ti ho detto qual è la mia RAL; se vogliamo instaurare un
rapporto in cui non ti fidi di quello che dico, abbiamo un altro problema''. È
di nuovo \emph{default to trust} (sezione~\ref{sec:z1-default-trust}). Il
relatore è quasi sicuro che legalmente non si possa chiedere la busta paga di
un candidato\footnote{È l'opinione del relatore, espressa come tale (``sono
quasi sicuro''); non si tratta di un parere legale e non è stata verificata.}:
a quel punto diventa lui antagonista del recruiter, e visto che la
conversazione è finita in una ``zona grigia'', fa lui le domande (``come mai è
legale questa cosa che mi chiedi?''). Ha già deciso che probabilmente non
lavorerà lì, ma così magari evita che i candidati successivi debbano subire le
stesse richieste. Cita l'esempio perché sa che in Italia è abbastanza comune,
mentre all'estero non gli è mai successo, anche se avrebbe potuto ``sparare
qualsiasi cifra''.

Il motivo è che \textbf{non c'è alcun onere di dimostrare quanto si
guadagna}: chi spara troppo alto si autoesclude. Recruiter e manager non
arrivano al colloquio con carta bianca né con l'indicazione di risparmiare il
più possibile, ma con un \emph{range}: ``per questa posizione possiamo
spendere dai 30 ai 40 mila euro l'anno; paga questa persona a seconda della
sua seniority''. Se un candidato che guadagna 30 dice di prenderne 55 per
``provare a fare i soldi'' e il budget arriva al massimo a 40, il relatore non
chiederebbe nessuna prova: direbbe ``grazie, siamo in un range diverso, possiamo
salutarci''. La domanda sulla busta paga è quindi \emph{formulata male}, perché
è poco furba. Quella giusta è: \textbf{``hai già una RAL in mente per questa
posizione?''}. Si ottengono le stesse informazioni senza passare per la zona
grigia.

\subsection{La cifra minima mentale}

E se la domanda la fanno a un junior che non ha mai lavorato? Il consiglio vale
per tutti: \textbf{preparatevi sempre la vostra cifra minima mentale}, per
qualsiasi offerta di lavoro, prima ancora di fare il colloquio. Non serve
conoscere la media italiana per quella posizione o quanto paga quell'azienda:
serve sapere \emph{quanto si è disposti ad accettare senza andare in negativo},
tenendo conto dei propri costi (la benzina per andare in ufficio ogni giorno, un
eventuale trasferimento, e così via). Quella cifra va sempre avuta in testa.

\subsection{La scrematura iniziale e il portfolio}

Uno studente chiede come funzioni la risposta a una candidatura: arriva un
rifiuto diretto, oppure c'è una scrematura iniziale e poi i colloqui? E vale
per tutti i ruoli? Sì: c'è una \textbf{scrematura iniziale} fatta dall'HR o, in
alcune aziende, direttamente dal software dell'HR. Si scarta chi è chiaramente
fuori da ciò che si cerca:
\begin{itemize}
  \item si cerca un \emph{graphic designer} e arriva il CV di un programmatore
        (``faccio questi esempi perché purtroppo succede'');
  \item si cerca un \emph{animator} tecnico bravo con MotionBuilder, lo si
        scrive nell'annuncio, e il candidato risponde ``so fare tutto tranne
        MotionBuilder'': c'è un motivo tecnico per non poterlo prendere;
  \item problemi di seniority: si cerca un senior e risponde un junior con un
        anno di esperienza.
\end{itemize}
Se la figura non è già chiaramente fuori target, si fa il primo colloquio.

\paragraph{Gli anni di esperienza negli annunci junior.} Spesso gli annunci
junior chiedono come minimo tre anni di esperienza. \textbf{Se l'annuncio dice
``junior'', ignorate i requisiti di esperienza}: o sono lì per sbaglio, o sono lì
per scremare. Il secondo caso è più \emph{borderline}, ma il relatore l'ha
visto fare in alcune aziende americane: si mettono requisiti molto alti,
chiaramente irraggiungibili, per vedere chi ha la ``faccia tosta'' di mandare
comunque il CV (``non me ne frega se chiedete un \emph{bachelor's degree}, non ce
l'ho, mando lo stesso''). Così però si filtra una grossa fetta di popolazione e
si seleziona un certo tipo di carattere: si stanno cercando ``i Top Gun''. Non è
quindi necessario iniziare con uno stage: si può iniziare direttamente da
junior.

\paragraph{Chi guarda il portfolio.} Il relatore non ha mai avuto un portfolio,
ma per le figure che lo richiedono parte comunque dal CV (per escludere, ad
esempio, il designer che si candida da programmatore) e guarda il portfolio
solo a parità con altri candidati: da HR può aprirlo e dire ``bello'', ma non
può dare un giudizio tecnico. Il portfolio lo apre chi sarà l'hiring manager,
che è sempre molto più tecnico. L'ordine quindi è: l'HR guarda il CV e si salva
il portfolio; se il candidato passa la prima scrematura, CV e portfolio vanno
a un tecnico che conferma (``sì, sa fare queste cose'') o boccia (``il
portfolio indica chiaramente che non sarebbe in grado, digli pure che non
facciamo il colloquio'').

\subsection{Test pratici, da casa e con qualsiasi strumento}

\paragraph{Si può chiedere un test pratico al posto del portfolio?} Uno
studente chiede se, con un portfolio che non rispecchia le proprie capacità, si
possa chiedere direttamente un test. No: se l'azienda ha bisogno di verificare
una competenza ha già pronto un test e lo fa fare a prescindere. Il relatore,
per esempio, fa fare un test ai programmatori e lo fa correggere dal suo senior
developer; per gli artisti si chiedono animazioni di prova, chiaramente veloci.
Non chiede mai test troppo lunghi, perché non vuole far lavorare nessuno gratis:
non chiede ``fammi un pitch di progetto'', chiede uno spunto per capire come si
ragiona e come si lavora. Per il resto, \textbf{se il portfolio non è
indicativo, va aggiornato il portfolio}.

\paragraph{Test dal vivo o a casa?} Il relatore li fa fare sempre a casa. Sa
benissimo che un programmatore a casa lavorerà con due monitor, il codice da
una parte e Google o ChatGPT dall'altra. Ne ha discusso con i suoi senior
programmer: interessa che qualcuno sappia le cose a memoria, o che sappia
trovarle e fare quel lavoro a prescindere? A lui interessa che il lavoro
\textbf{venga fatto}: che il candidato abbia imparato il codice a memoria o che
in venti minuti trovi le informazioni, le metta insieme, le faccia elaborare a
ChatGPT, le corregga e consegni il codice che serve, non ha preferenze. Quando
si lavora le indicazioni saranno ``fai questa cosa il prima possibile'': se la
si sa fare anche con strumenti aggiuntivi, ben venga, ``anzi, spiegalo anche
agli altri''. Farla a memoria va benissimo, ``ma perché mettere l'\emph{hard
mode}?''. E comunque ci sono cose che non si possono delegare a uno strumento:
se arriva un codice che non funziona e va fatto il \emph{debug}, bisogna
saperlo fare e saper progettare i test. La figura più completa, per il
relatore, è anche quella che sa trovarsi le informazioni e gli aiuti.

Le prove dal vivo sono soprattutto domande di \emph{problem solving}: il
relatore vuole sentire come si ragiona e chiede espressamente di \textbf{pensare
ad alta voce}, per seguire il percorso mentale che porta alla conclusione. Sono
più domande impegnative che veri test. I test veri e propri si fanno con i
propri tempi e i propri strumenti: ``scegli il momento, ti mandiamo il test, hai
due ore per rimandarlo indietro''. Questo è l'accordo; in quelle due ore si può
fare quello che si vuole.

\subsection{Junior, mid e senior}

\paragraph{Perché ci sono così pochi annunci per figure mid?} Uno studente ha
notato su LinkedIn e sulle piattaforme di recruiting un divario notevole tra le
richieste di junior e quelle di mid. Il motivo è che per molte aziende,
Stormind inclusa, \textbf{la differenza tra junior e mid è sparita}: o si è
senior o non lo si è. Ci si è resi conto che la differenza tra un junior e un
mid è che il mid si è abituato a lavorare in quell'azienda e il junior non
ancora; ma un mid che arriva con tre anni di esperienza in un altro studio (ad
esempio Milestone\footnote{Nella trascrizione ``se ha fatto tre anni in
milestone'': si tratta verosimilmente di Milestone, studio di sviluppo
milanese.}) va comunque ``ripreso come junior'', perché si lavora in modo
diverso e deve reimparare. Un senior, invece, ha molte più competenze e le sa
applicare in un certo modo, e quella distinzione resta netta.

\paragraph{Ci si può candidare anche se il livello non coincide?} Un'altra
studentessa dava per scontato che, se un annuncio non specifica ``junior'', sia
per mid. Se non è specificato ``senior'', secondo il relatore si può mandare il
CV chiedendo di valutare. Certo, se è specificato junior e si è junior le
possibilità di essere contattati sono maggiori. Ma vale lo stesso discorso del
livello di inglese fatto il giorno prima: si cerca qualcuno che \emph{sappia
fare quelle cose}, non una certificazione. Se serve un programmatore che sappia
usare Unreal Engine~5 e sistemare la cosa X o Y, che sia mid da due anni o
junior da uno non è un problema.

\subsection{Gli strumenti di intelligenza artificiale}

Una domanda in chat riguarda l'uso quotidiano dell'intelligenza artificiale in
azienda. Il relatore sta facendo fare corsi di IA al suo team ``perché voglio che
la gente lavori meno'': vuole che si sfrutti qualsiasi strumento disponibile,
perché gli altri lo stanno facendo e non vuole restare indietro. Se un lavoro
si può fare in tre ore di meno, non c'è alcun vantaggio a farlo in modo lungo e
tortuoso. Il rischio esiste: se si fa fare tutto a ChatGPT e non ci si tiene al
passo con il proprio \emph{know-how}, quando ci sarà da risolvere qualcosa ci si
potrebbe ``impallare''. Ma sta al singolo restare aggiornato e saper fare il
proprio lavoro: \textbf{gli strumenti devono essere strumenti, non un sostituto
delle competenze}. ``Non è una partita di calcio, con regole su come si devono
fare le cose o un \emph{dress code}'': il lavoro va fatto nel modo più efficiente
possibile.

A uno studente a cui era saltata la connessione proprio in questo momento il
relatore risponde che la lezione è registrata e si può recuperare dal video.
Poi chiude la parentesi: ``la parte difficile era finire ieri in tempo, ce
l'abbiamo fatta''.

% ---------------------------------------------------------------------
\section{Time management: pianificare la giornata}
\label{sec:z2-pianificare}

\subsection{Time management o prioritizzazione?}

La domanda di partenza alla classe è: \emph{chi di voi, in questo momento, fa
time management su se stesso}, cioè si organizza la giornata all'inizio e
organizza il proprio tempo? Qualcuno risponde di sì, qualcuno ha uno
\emph{schedule} settimanale con orari da rispettare, qualcuno ``ci prova''
(``hai fatto bene a specificare \emph{ci provo}: magari non la rispetto''). Il
relatore lo chiede perché è uno di quegli argomenti in cui molti, dopo aver
visto la spiegazione, si accorgono di non fare davvero time management: fanno
piuttosto \textbf{prioritizzazione}, ad esempio una lista di cose da fare, che
è un po' diverso.

Nessuna delle tecniche è complicatissima, ma è utile sapersi muovere. Il
relatore colloca volutamente questo modulo \emph{dopo} quello sul feedback: saper
gestire il proprio tempo permette anche di dare feedback agli altri sulla loro
gestione del tempo.

\subsection{La \emph{must do list}}

Il punto di partenza è il \textbf{planning iniziale}: non improvvisare la
giornata. Improvvisando si è in balia di una grande limitazione umana: ci si
dimentica le cose, oppure ci si lascia trascinare dal task più urgente o
dall'ultimo arrivato. Esistono modi di dire come ``se arriva una mail o rispondi
subito o diventa un task per domani'', perché il ``rispondo dopo'' non arriva
mai; più avanti si vedranno modi più raffinati di gestire le zone grigie, ma la
pianificazione resta fondamentale.

La prima cosa da fare \emph{non} è la lista delle cose da fare oggi, ma la lista
delle cose che \textbf{assolutamente} vanno fatte oggi. Le due liste non devono
coincidere: la seconda richiede un livello di ragionamento in più. Sono le cose
per cui, se alle sei si stacca e non sono fatte, si resta finché non sono fatte,
perché sono imprescindibili.

\paragraph{Il primo momento per dire ``potrei non farcela''.} La \emph{must do
list} è anche il primo indicatore della giornata. Se è troppo corposa (dipende
naturalmente dalla complessità dei task), quello è il momento, ``alle 9:01 se
iniziate alle 9'', per dire al proprio manager \emph{oggi sarò in difficoltà}.
Alle 9:01 il manager ha l'intera giornata per aiutare; se invece si pecca di
ottimismo o di presunzione (``ce la posso fare'') e glielo si dice alle 17:59,
non può fare nulla. È lo stesso discorso del giorno prima
(sezione~\ref{sec:z1-triangolo}): anche facendo la to-do list si applicano i
concetti della cultura del feedback, ed è il motivo per cui questo modulo viene
dopo. Ci saranno vari momenti di questo tipo nella giornata, ma questo è il
primo.

\paragraph{E se ci si accorge a metà giornata di non farcela?} Uno studente
chiede che cosa fare se, convinto di farcela, si accorge a metà giornata di non
poter completare bene la \emph{must do list}. Il relatore lo concederebbe
\emph{una volta}: se si vogliono testare i propri limiti, va bene, ma
comunque si dice al capo ``oggi potrei essere tirato, ci provo'', e sarà lui a
valutare se il rischio è accettabile. Dipende dal rischio: se il \emph{must do}
è consegnare una milestone a un cliente, ``nel dubbio facciamolo fare a
qualcun altro, non possiamo correre quel rischio''; se il rischio è che domani
un animator riceva il codice con tre ore di ritardo, lo si può correre. ``È un
po' questo il lavoro.''

\paragraph{Cinque minuti e una stima.} Tutto questo va fatto in circa cinque
minuti. Se non è chiaro che cosa si deve fare in una giornata, \textbf{si
chiede}: ogni giorno dovrebbe essere chiaro. Poi si cerca di \textbf{stimare il
tempo} di ogni cosa. Si diventa più bravi col tempo: alla quinta volta che si
fa il report mensile si sa che ci vuole un'ora e ventidue minuti; la prima volta
si può chiedere a un collega quanto ci mette lui, sapendo che per noi potrebbe
essere diverso.

\subsection{To-do list e checklist}

Esistono molti strumenti per tenere to-do list e checklist; il relatore non ne
consiglia uno in particolare e cita come esempi Any.do, i promemoria di Windows
o Samsung Notes\footnote{Nella trascrizione ``AnyDo, Windows Memo''. Any.do è
un'app di liste e promemoria; ``Windows Memo'' indica probabilmente le note o i
promemoria integrati in Windows (Sticky Notes, Microsoft To Do).}: ``potete usare
davvero la qualunque''. L'importante è tenere la lista \textbf{visibile} da
qualche parte sul computer. Alcuni accorgimenti:
\begin{itemize}
  \item i \textbf{must do} vanno sempre \textbf{in cima} alla lista;
  \item quando durante la giornata arriva un task nuovo, la cosa da fare subito
        non è \emph{eseguirlo}, ma \textbf{aggiungerlo alla lista}: ``se arriva
        una mail, aggiungete subito alla lista \emph{rispondi alla mail}'',
        decidendo se è diventato un \emph{must do} o se è comunque un task per
        oggi;
  \item non limitarsi alle cose operative: il relatore ha iniziato ad
        aggiungere anche task come ``parla con Tizio'' o ``parla con Caio'',
        perché è da un po' che non chiede loro un aggiornamento; se non li
        mette in lista si accorge di non farli;
  \item spuntare le voci completate: per quanto minimo, rilascia
        \textbf{dopamina} (se ne parlerà nel pomeriggio,
        sezione~\ref{sec:z2-biologia}), aiuta a motivarsi e ad andare avanti.
        Sono piccole cose, ma si sommano.
\end{itemize}

\paragraph{Priorità contro recenza.} Che cosa succede se si esegue sempre
l'ultimo task arrivato? Lo chiede il relatore alla classe. Si interrompe il
flusso di lavoro e si cambia continuamente attività, certo; ma soprattutto si
stanno \textbf{mescolando priorità e recenza}: si sta dicendo ``non importa come
ho prioritizzato i task, la mia priorità è sempre l'ultimo arrivato'', che non
vuol dire fare le cose più importanti. È il motivo per cui tanti dipartimenti IT,
quando si chiede loro qualcosa, rispondono ``mi apri un ticket?''. È diventata
quasi una battuta, ma il concetto è chiarissimo: devono poter prioritizzare
quella richiesta rispetto alle altre 112 arrivate in giornata, e senza ticket
non possono confrontarla con le altre. Il relatore è totalmente a favore: ``certo
che ti apro un ticket, altrimenti tu come fai a tracciare il tuo lavoro?''.

\subsection{La checklist: l'esempio dell'aviazione}

La \textbf{checklist} non è la task list. La task list contiene le cose da fare;
la checklist contiene \textbf{i passi da seguire quando si esegue un task
specifico}. L'esempio è quello delle \textbf{checklist pre-volo}, introdotte
nell'aviazione (secondo il relatore tra gli anni Sessanta e Settanta): chi vola
spesso nota che i piloti seguono sempre un ordine simile nei controlli, tanto
che il relatore capisce quanto manca al decollo da quando controllano i flap, i
motori, quando accendono e spengono le luci. L'introduzione delle checklist
avrebbe ridotto di circa il 75\% gli incidenti in fase di decollo\footnote{Date
e percentuale sono quelle riferite in lezione e non sono state verificate.
L'origine della checklist pre-volo viene di solito fatta risalire al 1935, dopo
l'incidente del prototipo del Boeing B-17, causato da un controllo dimenticato da
un equipaggio molto esperto; la diffusione sistematica delle checklist e dei
protocolli di cabina è invece successiva.}, e soprattutto per i \textbf{piloti
più senior}. Perché? Perché chi ha fatto un decollo centomila volte non va a
controllare l'alettone destro: ha funzionato centocinquantamila volte, funzionerà
anche stavolta. Quando ci si abitua a fare qualcosa, si sbaglia proprio quando
qualcosa non va. La checklist ha obbligato a fare quei controlli ogni volta e ha
ridotto in modo drastico gli incidenti.

Il relatore usa checklist anche per cose banali, ad esempio le \textbf{email}:
essendo una persona abbastanza ansiosa su queste cose, se prepara una mail alla
leggera finisce per passare troppo tempo a ricontrollare l'oggetto, il
destinatario, l'allegato, e comunque si dimentica qualcosa. Seguire una checklist
sembra più complesso, ma \textbf{fa risparmiare tempo}: se la checklist è
completa, si sa di aver fatto tutto e si può andare avanti senza continuare a
pensare ``ho chiuso il gas? ho spento la luce?''. È come la lista della spesa: se
tutto è segnato e preso, punto. È un investimento: trenta secondi in più per
fare la cosa, invece di cinque minuti a sperare di averla fatta giusta.

\paragraph{E se ci si dimentica qualcosa nella checklist?} Succede. A quel
punto \emph{preparare la checklist} è un task in sé: un giorno in cui si ha
tempo si fa quel task pensando a quali passi si seguono e a come assicurarsi che
funzioni. Nessuna checklist è perfetta: ci saranno sempre compromessi tra ``devo
davvero controllarlo ogni volta?'' e ``posso farne a meno?''. Però, in generale,
\textbf{se qualcosa anche solo una volta ha dato problemi, va in checklist},
perché vuol dire che potrebbe darne di nuovo. Nessun manager rimprovererà
``come hai fatto a non prevedere che questa cosa poteva andare male?'': si
diventa esperti anche nelle checklist di un task facendo quel task, ed è un
processo iterativo. Se qualcosa è stato dimenticato e ha creato problemi, lo si
aggiunge per la volta dopo. Sbagliare una o due volte non è un problema;
l'importante è che la situazione migliori la volta successiva.

\paragraph{L'inizio di un junior e l'\emph{onboarding}.} A uno studente che
osserva che l'inizio di un junior è quindi inevitabilmente un po' travagliato,
il relatore risponde che è giusto così: ``se volessi il dipendente perfetto da
subito, cercherei un senior, non un junior''. Allo stesso modo non capisce le
aziende che prendono un junior e non lo fanno seguire da nessuno: che cosa ci si
aspetta, che assorba passivamente tutto quello che ha intorno e impari ``per
osmosi''? Può succedere, ma imparerà a fare le cose come vuole lui, perché
nessuno gli sta dicendo ``meglio farle così''. Moltissime aziende, oltre al
manager (che magari ha quindici \emph{report} e non ha il tempo materiale di
seguire un junior come andrebbe seguito), assegnano un \textbf{referente}, di
solito un senior, con il compito di prendersi cura della persona nuova e
insegnarle a fare le cose. Altrimenti l'\textbf{onboarding}, cioè il tempo
necessario per arrivare al cento per cento nel proprio lavoro, dura un anno.
Più risorse l'azienda dà all'inizio, prima la persona diventa operativa: è
nell'interesse dell'azienda dare il massimo supporto a chi arriva.

Il relatore diffida delle aziende che ``vogliono vedere se uno se la cava''. È
il discorso dei cento metri e della corsa a ostacoli\footnote{Il relatore lo
introduce come un discorso già fatto (``si torna al discorso: vuoi fare i cento
metri o la corsa a ostacoli?''), ma il passaggio non compare nella trascrizione
della prima giornata riportata nel capitolo~\ref{cap:z1-feedback}.}: non si può
far correre uno sui cento metri e l'altro con gli ostacoli e aspettarsi tempi
simili, perché chi ha gli ostacoli ci metterà per forza di più. Il supporto ai junior è un
investimento sensato.

\subsection{\emph{First contact resolution}}

Un concetto che viene dai call center: \textbf{quante richieste vengono risolte
al primo contatto}. È importante quando un task richiede iterazioni o
comunicazione con altri: bisogna organizzarsi come se si dovesse risolvere la
cosa \emph{in una sola passata}. Se viene chiesto un modello 3D, lo si fa giusto
la prima volta, al meglio delle proprie possibilità. Se si deve parlare con
qualcuno, si raccolgono tutte le richieste in una.

L'esempio è quello dei ticket che il relatore mandava ai colleghi del supporto di
Blizzard: ``ciao, ho un problema, non riesco a loggarmi'' --- ``mi dici lo
username?'' --- ``è questo'' --- ``mi dici l'email?'' --- ``è questa'' --- e così
via, con un paio d'ore tra un messaggio e l'altro perché il ticket deve fare il
suo giro. Tutto poteva essere riassunto in ``ciao, mi servono username, email e
questo e quest'altro''. Unificare le richieste su una singola cosa aiuta a
risolverla il prima possibile; bisogna avere chiaro \emph{che cosa serve prima
di fare una richiesta}, in modo che ce ne sia una sola. In molti sistemi di
supporto, se si apre un secondo ticket per aggiungere informazioni al primo, il
primo, anche se era quasi in lavorazione, viene rimandato in fondo alla coda:
serve a disincentivare chi manda continuamente ticket del tipo ``mi sono
ricordato questa cosa''. O si manda tutto subito, o ci vorrà tempo.

\subsection{Non iniziare senza le informazioni}

Corollario: \textbf{se non si ha tutto quello che serve per completare un task,
non lo si inizia}. Prima ci si procura quello che manca. Al relatore capita con
la pianificazione dei due anni di lavoro successivi: se non ha tutte le
informazioni, prima lavora per averle (``mi dai questa informazione, mi dai
quell'altra''), poi pianifica; altrimenti pianificherebbe qualcosa di cui non ha
completa visibilità.

\subsection{Il costo delle interruzioni}

Alternare continuamente l'attenzione tra i task e, soprattutto, reagire di
continuo agli input esterni (la mail, il messaggio, il collega che parla) fa
\textbf{sparire il tempo dalla giornata}. Se mentre si prepara un report in
Excel si viene interrotti (dal gatto, da una mail, da un collega, da una
telefonata) e poi si torna al foglio, servono in media \textbf{un quarto d'ora}
per tornare mentalmente al cento per cento; chi è particolarmente bravo ci mette
dieci minuti, ma si parla comunque di dieci-quindici minuti\footnote{Ordini di
grandezza come riportati in lezione. Le ricerche più citate sul tema (ad esempio
quelle di Gloria Mark sulle interruzioni nel lavoro d'ufficio) indicano tempi
medi di ripresa anche superiori, intorno ai venti minuti.}. Se si viene
interrotti di continuo, si passa più tempo a cercare di tornare al cento per
cento che a lavorare. È una partita persa in partenza.

Per questo il giorno prima si era parlato dei \textbf{semafori} sulla scrivania
(sezione~\ref{sec:z1-dare}): il semaforo può essere anche semplicemente il
proprio stato su Slack, su Discord o su qualunque strumento si usi con i
colleghi. Se si è occupati, si è occupati: è una propria valutazione che,
accettando di essere disturbati di continuo, quella cosa non si riuscirà a
farla, o almeno non in tempo.

% ---------------------------------------------------------------------
\section{Prioritizzare: la matrice di Eisenhower}
\label{sec:z2-eisenhower}

Ora che si hanno task list, checklist, \emph{must do} e \emph{nice to do}, resta
la domanda di prima: se arriva un task nuovo e si decide di non fare per prime le
cose nuove, \emph{come si capisce se è davvero urgente}? Serve un metodo, ed è la
\textbf{matrice di Eisenhower}\footnote{Nella trascrizione anche ``Jason
Howard''. Dal nome del presidente statunitense Dwight D.~Eisenhower, a cui si
attribuisce la distinzione tra urgente e importante.} (la classe prova a
indovinarne il nome proponendo altre matrici, come una matrice dei rischi, ma è
un'altra cosa). Ogni task che arriva va posizionato su questa tabella:

\begin{center}
  \begin{tabular}{@{}p{0.2\textwidth}p{0.36\textwidth}p{0.36\textwidth}@{}}
    \toprule
    & \textbf{Urgente} & \textbf{Non urgente} \\
    \midrule
    \textbf{Importante} &
    \textbf{Fallo} (\emph{do}): diventa un task di oggi. &
    \textbf{Pianificalo} (\emph{schedule}): va fatto, ma non subito. \\
    \addlinespace
    \textbf{Non importante} &
    \textbf{Delegalo} (\emph{delegate}): va fatto, ma non serve che lo faccia
    tu. &
    \textbf{Eliminalo} (\emph{delete}): non lo fai. \\
    \bottomrule
  \end{tabular}
\end{center}

\paragraph{Urgente e importante: lo fai.} Sono di solito task con una
\textbf{deadline molto chiara} (entro oggi, entro un'ora, entro X) e con
\textbf{conseguenze} se non vengono fatti: da qui l'``importante''.

\paragraph{Importante ma non urgente: lo pianifichi.} Va fatto, ma non subito:
o lo si mette tra le cose di oggi ``ma dopo'', o anche alla settimana successiva.
L'esempio: un collaboratore scrive ``il mio rinnovo è tra tre mesi, ne possiamo
parlare?''. È importante, perché il rinnovo di una persona è importante; non è
urgente, se ne può parlare anche domani, così da potersi preparare e fare altre
cose nel frattempo.

\paragraph{Urgente ma non importante: lo deleghi.} Qui ``non importante''
\textbf{non} vuol dire che il task non sia importante, ma che \textbf{non è
importante che lo faccia tu}, che è molto diverso. Se arriva ``mi serve questo
cambio a questa animazione 3D'' e ciò che conta è il cambio, non chi lo fa, si può
dire ``sono d'accordo che è urgente, ma non serve che lo faccia io'' e chiedere a
qualcun altro. Sono task che vanno fatti ma non richiedono la propria particolare
\emph{expertise}.

Il relatore insiste molto su questo punto con manager e senior: quando arriva un
task urgente, i senior tendono a dire ``lo faccio io, perché è urgente''. Ma se un
junior può farlo, lo si fa fare al junior e si investono le competenze più alte
del senior in qualcos'altro: il senior può fare il lavoro del junior, ma il
junior non può fare il lavoro del senior mentre il senior fa quello del junior.
Le persone non sono intercambiabili. Si usa il \textbf{minimo sforzo
necessario}: se una cosa la può fare un mid, la fa il mid, non il senior.

\paragraph{Né urgente né importante: lo elimini.} È il quadrante che quasi
nessuno usa: tutti fanno, delegano o pianificano, ma non pensano mai a
\emph{delete}. Invece è fondamentale. A volte è corretto dire: ``questa cosa non
dovrei farla io, non è urgente, non la faccio; sappi che non la faccio, e se c'è
un problema ne parli con il manager''. Non si sta delegando a qualcun altro: si
sta dicendo che quella cosa così com'è \textbf{non andrebbe fatta}, perché non
ne sono chiari l'importanza e il perché. Nella propria lista di oggi non ha
spazio, e non è abbastanza importante per essere fatta da sé nemmeno nei giorni
successivi.

Un esempio di \emph{delete} che non è un vero \emph{delete}: arriva una richiesta
di meeting di tre ore sul futuro del progetto, con invitati tutti gli art
director, il game director e così via. Il relatore spesso risponde ``io non
vengo'': per lui è \emph{delete}, ma il meeting si fa lo stesso. Uno studente
osserva che a volte si tratta di una richiesta mal formulata o poco chiara: è
esattamente così. ``Nella sua forma per me è no; se poi mi spieghi meglio perché
dovrei farlo io o perché è urgente, allora cambia.''

% ---------------------------------------------------------------------
\section{Proteggere il proprio tempo}
\label{sec:z2-proteggere}

\subsection{Evitare gli sprechi: un limite che nessuno può insegnare}

Evitare gli sprechi di tempo è difficile, soprattutto da remoto: è una sfida
costante con se stessi. Ha molto a che fare con il carattere, il
\emph{self-control} e (``passatemi anche questa'') la professionalità. Il
relatore può dare tutti i consigli del mondo, ma se dentro di sé non si percepisce
l'urgenza di fare il proprio lavoro quando lo si deve fare, non lo si può
insegnare. Se si è a casa e invece di lavorare si gioca a qualcosa su Steam,
nessuno deve spiegare che non si dovrebbe. Ci sono limiti oltre i quali non si
può andare.

Quello che si può fare è sapere che \textbf{ogni interruzione ha un costo
sommerso}: ogni messaggio, domanda, telefonata, ogni cambio di focus. Sono i
quindici minuti di prima, che possono arrivare a trenta a seconda del tipo di
task.

\paragraph{È più difficile tornare su un task logico o creativo?} Il relatore lo
chiede alla classe: fa più fatica un artista a tornare sul suo modello 3D o un
\emph{data analyst} a tornare sul suo foglio Excel? Le risposte si dividono
(``creativo'', ``entrambi''), ed è proprio il punto: dipende a chi lo si chiede. I
creativi rispondono ``creativo'', gli analisti ``analitico'', perché ognuno ha una
chiara comprensione di che cosa serve per essere davvero concentrati su ciò che
fa. Il creativo sa che per essere pienamente creativo gli servono l'ambiente, la
concentrazione, gli spunti, gli esempi da tenere chiari in testa; l'analista deve
avere in testa i conti, i fogli di dati e così via. Sono \textbf{ugualmente
difficili}, ma solo chi fa quella cosa ne apprezza fino in fondo la difficoltà. Si
può avere un assaggio dell'altra parte, ma la piena comprensione arriva solo
facendo: è come guardare uno sport in tv e pensare ``potrei farlo anch'io'',
finché non ci si trova lì.

\subsection{Distrazioni volontarie e involontarie}

Uno studente, che sta leggendo manuali di \emph{design thinking}, chiede come
conciliare tutto questo con l'idea che distrarsi, cioè non restare focalizzati,
aiuti a trovare un nuovo punto di vista per risolvere un problema. Qual è il
giusto equilibrio? Il relatore è \textbf{totalmente d'accordo} con quel concetto:
sta parlando di \textbf{interruzioni non volontarie}. Sa benissimo che dopo cinque
ore su un foglio Excel ``esplode la testa''; ma se si pianifica un quarto d'ora di
pausa e lo si rispetta, si sta investendo sulla propria creatività ed
efficienza. Non è un investimento dover rispondere a una chiamata ogni tre
minuti.

Anzi, il consiglio è \textbf{crearsi interruzioni pianificate}, sotto il proprio
controllo: fare pause, fare altro, soprattutto nel \emph{problem solving}. Il
relatore stesso, se non trova una soluzione, si ferma, fa un altro task e poi ci
torna; ma è una sua scelta. La discriminante è che la pausa sia \emph{scelta} e
non imposta: anche se si fanno le pause giuste, se si viene interrotti proprio
nel momento di massima concentrazione (raggiunta magari grazie alla pausa) si
butta via tutto. Uno studente cita la \textbf{tecnica del pomodoro}: è
esattamente quello, piccole pause pianificate che aiutano a essere creativi,
purché restino pianificate. Se i quindici minuti diventano ``vabbè, facciamo
mezz'ora'' o ``lo faccio nel pomeriggio'', non è più la tecnica del pomodoro:
diventano interruzioni, o un pretesto per non lavorare.

\subsection{Il peggior nemico siamo noi}

E chi sono i peggiori nemici in questo? ``Noi stessi'': siamo noi ad
autoconvincerci che venti minuti vanno bene quanto quindici. Le
\textbf{distrazioni autoinflitte} sono spesso le peggiori. Il relatore porta
l'esempio di una sua giornata, che ha annotato perché lo aveva fatto ridere
quanto tempo era riuscito a buttare via senza che nessuno glielo imponesse:
\begin{itemize}
  \item gli piace avere qualcosa in sottofondo mentre lavora (``ha a che fare con
        l'ADHD sicuramente''\footnote{Nella trascrizione ``IDACD''. Più avanti il
        relatore dice di avere l'ADHD ``come dicevo ieri''; il capitolo 0.1 non
        lo riporta, ma riferisce il suo racconto sull'``overthinking'' anche
        nell'ordinare una pizza (sezione~\ref{sec:z1-regalo}).}),
        e in genere lo aiuta. Il problema è quando il sottofondo è troppo
        interessante: la musica va bene, ma se nella playlist di Spotify parte un
        podcast di Alessandro Barbero sulla guerra civile americana, finisce per
        ascoltare quello più del lavoro;
  \item un'abitudine che ha poi perso: se faceva straordinari alle due di notte,
        da solo, per qualcosa di urgentissimo, apriva comunque Slack a
        controllare i messaggi, e la mail. ``Come se alle due di notte qualcuno
        mi scrivesse qualcosa di urgente'';
  \item la pausa con un po' di TikTok, dopo la quale si accorgeva di aver
        dimenticato su quale slide della presentazione stava lavorando.
\end{itemize}
La regola: \textbf{inserire i break è meglio che subirli passivamente}. Un conto
è dire ``faccio una pausa intenzionalmente'', un conto è ``mi sono distratto e ho
perso due ore di lavoro''. Chi ne è particolarmente soggetto non deve
demoralizzarsi: è anche un allenamento.

\subsection{Spezzettare i task grandi}

Un altro concetto utile: se si fa molta fatica a iniziare i task grossi, perché
sembrano infiniti (``più tardi inizio, più tardi inizia questa sofferenza''),
\textbf{vanno spezzettati}. Un task grosso non è quasi mai un task grosso: è
spesso un insieme di task piccoli. Si comincia a farne dei pezzi e si va avanti
piano piano, spuntandoli. Aiuta tantissimo anche nel \emph{problem solving}: un
problema gigante va scomposto in problemi più piccoli, da risolvere per primi, per
poi ricomporre il puzzle.

\subsection{Accettare di dire di no}

Tornando alla matrice di Eisenhower: bisogna \textbf{accettare di dover dire di
no}, perché dire sempre sì porta a problemi ben peggiori di quanto si pensi.

\paragraph{Il collega che diceva sempre sì.} A Riot il relatore aveva un collega
estremamente disponibile, amichevole e collaborativo: qualsiasi cosa gli si
chiedesse, rispondeva ``certo, la faccio io'' (``mi porti fuori il cane?'' ---
``certo''). Il problema è che era umano: a un certo punto si supera la quantità di
cose umanamente possibili in una giornata. In più tendeva a offrirsi volontario
per far vedere che era intraprendente e aveva iniziativa, e nei progetti che
coordinava il relatore doveva spesso rifiutare la sua candidatura perché era
sicuro che non ce l'avrebbe fatta. Il suo dire sempre sì, a un certo punto, lo ha
danneggiato: era quello che dice sempre che lo fa lui, ma non è detto che lo
faccia. ``Ce la fai?'' --- ``Sì, sì, assolutamente'' --- e poi ``scusa, non ce l'ho
fatta''. ``Ok, però ora hai messo nei guai me e il resto del team.'' Quella
disponibilità era più dannosa di un ``no, non posso''. È lo stesso discorso del
giorno prima: dire all'ultimo momento che non ce la si fa è peggio che dirlo a
inizio settimana.

Uno studente lo paragona alla \emph{sindrome dello yes man}. Non proprio: lo
\emph{yes man} dà sempre ragione all'altro perché è avverso al conflitto e non vuole
la discussione del ``spiegami perché sei in disaccordo''. Il collega, invece,
pensava davvero di poter fare molto più di quanto potesse, e rifiutava l'idea di
non riuscirci.

\paragraph{``Ci hai pensato bene?''} Vale per i progetti e per le deadline. Con i
junior soprattutto, quando a ``vuoi fare questa cosa?'' arriva un sì troppo
convinto, il relatore chiede sempre: ``ci hai pensato bene? Me lo stai dicendo
per farmi contento? Sappi che se mi dici di no non penso peggio di te; penso
peggio di te se mi dici di sì e poi non ce la fai. Vuoi rivalutare la risposta, e
magari dirmi \emph{non lo so} invece di un sì categorico?''.

\paragraph{Che cos'è una sconfitta.} Uno studente chiede se davanti alle
``sconfitte'' non si dovrebbe rimediare. Il relatore chiarisce che per lui la
sconfitta non è il task non fatto: quella è una sconfitta \emph{a breve termine},
si troverà un modo di farlo, è un danno ``ma chi se ne frega''. La vera sconfitta,
un danno \emph{a lungo termine}, è aver pensato di riuscirci e non esserci
riusciti: se non si risolve il motivo per cui si è detto sì quando era no, lo si
ripeterà altre cinque, dieci, quindici volte. Nella testa del collega di Riot c'era
sempre una scusante (``c'erano questi imprevisti che mi hanno bloccato''), come
quella con cui ci si giustifica le interruzioni. Il relatore gli rispondeva: ``se
ogni volta c'è un imprevisto, forse c'è un \emph{pattern}: il problema non è che
non riesci a fare le cose, è che non prevedi mai gli imprevisti''. Dopo sette
ritardi su dieci dovrebbe scattare l'ingranaggio mentale (``forse sbaglio io a
dire sempre di sì''), ma non è sempre facile farlo scattare, e a volte serve una
spinta. È una delle situazioni più frequenti. E la cosa più difficile è che,
quando glielo si fa notare, queste persone lo prendono come un ``non sei capace''.
Il messaggio è invece: \textbf{``sei bravo, ma sei umano anche tu''}. Non si può
fare un numero infinito di cose in un'ora: bisogna accettare i propri limiti ``dati
dalla carne'' e non pensare di essere un cyborg.

\paragraph{Il cliente non ha sempre ragione.} Uno studente osserva che vale anche
per i project manager e per il principio ``il cliente ha sempre ragione''. Anzi:
il relatore ha lavorato con molti clienti e a volte \textbf{il cliente non sa che
cosa vuole}, che è molto diverso dall'avere sempre ragione. Soprattutto con i
clienti che chiedevano pitch creativi, la sua squadra presentava sì un'opzione
con la loro richiesta, ma anche le proprie (``se vuoi raggiungere questo, sarebbe
meglio fare così''), e molto spesso il cliente sceglieva quella proposta dal team.
Spesso il cliente ha in mente \emph{una soluzione}, ma bisogna capire \emph{quale
obiettivo} ha in mente e poi dargli una soluzione: non vuole la soluzione, vuole il
\emph{goal}, e se c'è un modo migliore per arrivarci, forse è meglio un'altra
soluzione.

\subsection{Dire di no ai meeting}

La prima cosa a cui il relatore ha imparato a dire di no sono i \textbf{meeting}.
Ha passato anni in ambienti in cui il suo calendario era letteralmente solo
meeting dalla mattina alla sera: estenuante per lui e inutile per il team. Tra i
fattori c'è il \textbf{FOMO}: tanti manager vogliono essere in tutti i meeting per
non perdersi niente ed essere sempre informati su quello che fa il team. In realtà
vanno in \emph{burnout}, passano i pomeriggi nei meeting e finiscono per essere i
più distratti, quelli che non ascoltano e che fanno domande inutili solo per far
sentire che fanno domande, senza aiutare davvero.

Prima di accettare un invito il relatore si fa una serie di domande, \textbf{come
una checklist} (e le fa anche agli altri):
\begin{enumerate}
  \item \textbf{Sono davvero necessario?} Servo davvero io, o se mancassi non
        succederebbe niente? Ai meeting ``di cortesia'' non va: se pensa di
        essere necessario lo dice lui, e il fatto che si faccia un meeting
        importante e urgente senza di lui non danneggia il suo ego. ``Non
        aggiungermi solo perché pensi che serva che io venga aggiunto.''
  \item \textbf{Uno di noi potrebbe partecipare e informare il resto del team?}
        C'è già qualcuno del suo team? Se sì, non può andare solo lui? Per un
        meeting informativo (una nuova policy, una nuova procedura da spiegare a
        tutti i dipartimenti), andarci tutti significa investire tre ore a
        testa; se ci va uno solo e poi spiega agli altri, si risparmia tempo.
  \item \textbf{Il riassunto sarà reso disponibile?} Per il relatore è
        obbligatorio. Per lo stesso motivo per cui la lezione viene
        registrata: se ci si perde un pezzo lo si può rivedere, se qualcosa non è
        chiaro lo si può risentire. Ma c'è anche un motivo di \emph{office
        politics}: se in un meeting si chiede ``mi fai questa cosa per venerdì?''
        e la risposta è sì, non si vuole sentire venerdì ``intendevo venerdì
        prossimo'' o ``non ero sicuro''. Nelle note del meeting c'è scritto che hai
        detto sì. ``Lo so che è un po' paraculo, ma in realtà è così che si
        lavora'': bisogna dare più peso alla parola data, e quindi tenerne
        traccia. Finché tutto resta orale, niente rimane scritto e tutto è
        interpretabile. Oggi è anche facile: Zoom, con l'assistente IA, produce
        in automatico trascrizione e sottotitoli, senza bisogno di qualcuno
        addetto alle \emph{meeting minutes}. In ogni caso, se in un meeting viene
        presa una decisione, il relatore vuole che sia \textbf{ripetuta per
        iscritto}: non deve ipotizzare, deve poter andare a rileggere.
  \item \textbf{Il meeting aggiunge valore rispetto al documento?} Se arriva un
        invito con allegato un documento di venti pagine, il relatore chiede: ``il
        meeting consiste nel leggermi il documento? Perché altrimenti non
        vengo''. Se non c'è un valore aggiunto, se lo legge da solo e manda
        conferma di averlo letto. È lo stesso motivo per cui si sconsigliano le
        presentazioni in cui si leggono le slide; e infatti, osserva, in quel
        momento lui stesso stava leggendo una slide, ma aggiungendo informazioni.
        ``Se tutta la lezione fosse io che vi leggo le slide, non servirei io: vi
        manderei le slide.''
  \item \textbf{Per che cos'altro potrei usare questo tempo?} È forse la più
        importante. Anche ammesso che serva, che serva proprio lui e che il
        meeting sia importante, in quel momento dovrebbe fare qualcosa di più
        importante? Se sì, fa quella. Non bisogna avere paura di dire ``no, non
        vengo'' anche se si serve, oppure ``ho un'altra cosa importante,
        spostiamo il meeting''. Il relatore si è sforzato di essere più
        \textbf{assertivo}: ``capisco la richiesta, ma ho qualcosa di più urgente
        da fare in questo momento''.
\end{enumerate}

\paragraph{E le richieste a voce?} Uno studente chiede se una richiesta fatta a
voce da un senior (``devi fare questo'') vada accettata così o fatta mandare per
mail o per messaggio. Accettarla a voce ci può stare. Ma se la richiesta è ``devi
farla in questo modo particolare e preciso'', è meglio farsi mandare il modo
preciso \textbf{per iscritto}: non si deve lasciare a interpretazione \emph{come}
fare quella cosa, per non ricadere nel ``ah, ma io non la volevo così''. Oppure si
conferma: ``sì, la faccio; va bene se la faccio così, così e così?''. Se dice sì,
si è a posto.

\paragraph{Non dovrebbero pensarci gli organizzatori?} Dalla chat: queste domande
non dovrebbero farsele già gli organizzatori dei meeting? Sì, sarebbe meglio che
tutti lo facessero, e il relatore stesso se le fa quando invita qualcuno: ``è
davvero necessario Marco? È davvero necessario Andrea?'', e solo se sì li invita.
Ma lavora in aziende anche di cinquecento dipendenti, dove non è detto che tutti
conoscano o applichino questa logica; per questo la applica per sé, sperando che
la applichino tutti. Ha sempre chiarito con tutti che da lui non arriveranno mai
inviti di cortesia o dettati dalla gerarchia (``deve esserci anche il mio capo''):
``devi esserci tu; se vuoi portarti il capo, portatelo, ma non lo invito solo
perché qui si vuole sempre il capo presente: sono sicuro che avrà di meglio da
fare''.

\paragraph{E i task?} Dire di no vale naturalmente anche per i task: si applica la
matrice di Eisenhower (sezione~\ref{sec:z2-eisenhower}) e si decide se fare,
delegare, pianificare o cancellare. Non c'è molto da aggiungere.

\subsection{Il multitasking non esiste}

``Questo potrà urtare tanti ego'', ma \textbf{il multitasking non esiste}. Si può
essere più bravi o più portati di altri a fare più cose contemporaneamente; la
realtà è che, concentrandosi su più cose, \textbf{la somma non farà mai cento}. Se
si lavora su due cose non si lavora al 50\% su una e al 50\% sull'altra, ma al 40\%
e al 40\%, e il 20\% va perso in ``costo di manutenzione'', cioè nella capacità di
seguirle entrambe. Il consiglio è dedicarsi al cento per cento a un task, poi
passare al successivo, poi al successivo, invece di farne cinque insieme: si può
avere la percezione di andare più veloci, ma non è così.

Uno studente cita una lettura di neuroscienze secondo cui il cervello è in grado
di fare una sola operazione alla volta. Il relatore conferma: smettiamola di
parlarne come di un pregio o di una \emph{skill}. È un modo che alcuni cervelli
trovano per svagarsi, l'equivalente di accendere un podcast: fare due cose
insieme vuol dire cercare di distrarsi alternandole, e farle entrambe peggio di
quanto si farebbero una alla volta. A chi preferisce fare tante cose insieme
suggerisce di provare a farne una alla volta: le farà più velocemente. Anche il
semplice passaggio da una cosa all'altra non è istantaneo: ci sono sempre una
decina di secondi in cui il cervello si chiede ``che cosa stavo facendo?''; per
cinquanta cambi, già questo spiega la perdita di velocità. È una questione di
\textbf{bandwidth mentale}, un concetto che tornerà nel pomeriggio: anche lo
\emph{switch} occupa banda.

\subsection{La fine della giornata}

A fine giornata si \textbf{rimette a fuoco la giornata}: si riguardano la task list
e la \emph{must do list} e si conferma di aver fatto tutto. È letteralmente il
controllo prima di uscire di casa: chiavi, porta chiusa, cellulare, portafoglio.
Al relatore è capitato di dimenticare un \emph{must do} rimasto in lista per tutto
il giorno, ed entra in gioco quella che chiama la \textbf{sindrome del
tirannosauro}\footnote{Il riferimento è alla celebre scena di \emph{Jurassic Park},
in cui si dice che il tirannosauro vede la preda solo se si muove.}: se una cosa è
lì, non è detto che la si veda finché non si muove. Averla in lista tutto il giorno
non vuol dire averla percepita come da fare.

Può anche capitare di non farcela perché quel giorno sono arrivate così tante cose
urgenti e importanti da rendere inutili le stime: la giornata no capita e non c'è
nessun problema. Ma se era davvero un \emph{must do}, \textbf{va fatto comunque}:
è lì che scattano la richiesta di straordinario o il ``devo farlo comunque''.
Attenzione: se si chiude la giornata con un \emph{must do} ancora aperto e si
decide ``vabbè, lo faccio domani'', può starci, ma allora \textbf{non doveva essere
in quella lista}, perché era rimandabile.

Al relatore questo controllo serve anche ``per quieto vivere'': se chiude ogni
giornata con la \emph{must do list} completa, sa che magari non ha fatto tutto
quello che c'era nella to-do list, ma i \emph{must do} sì. ``Il mio, oggi, l'ho
fatto'': può stare tranquillo senza continuare a chiedersi se ha fatto quello che
doveva.

% ---------------------------------------------------------------------
\section{Domande sul time management}
\label{sec:z2-qa-tempo}

\paragraph{C'è un numero medio ``sano'' di \emph{must do} quotidiani?} No, perché un
conto è avere dieci \emph{must do} come ``rispondere a dieci mail'', un conto è
averne uno solo come ``inventarmi un nuovo gioco''. Dipende dalla dimensione: se ne
gestiscono di più se sono minuscoli, di meno se sono giganteschi. Ci sono lavori
che per definizione hanno tantissimi \emph{must do} quotidiani e lavori che ne hanno
uno al giorno. Si imparerà da sé qual è il numero giusto.

\paragraph{C'è un modo per stimare quanto si riesce a fare in una giornata?} ``Lo so
che non è una bella risposta'', ma l'unico modo è la propria \textbf{storia}:
sapere quanto si riesce a fare nel quotidiano. È normale e fisiologico che le
prime stime siano sbagliate, e probabilmente eccessivamente ottimistiche;
l'importante è rendersene conto. È letteralmente la \emph{seniority}: abituarsi a
fare qualcosa insegna quanto ci si mette. Va benissimo metterci di più all'inizio,
ed è normale che il tempo cambi man mano che si migliora: ciò che oggi si stima in
un'ora, tra sei mesi lo si stimerà in mezz'ora. Chi ha lavorato con Scrum lo sa:
su un progetto nuovo tutti sbagliano la dimensione delle \emph{stories} all'inizio,
perché nessuno sa quanto ci vorrà, e al terzo \emph{sprint} consecutivo è più
chiaro. Le stime sono stime, altrimenti si chiamerebbero ``inserire quanto tempo ci
metterai''.

Se però si dà una stima e poi il lavoro si rivela il triplo, il problema non è la
stima: è che \textbf{bisognava chiedere più informazioni} prima di darla. La
risposta giusta non è la stima, ma ``per darti una stima accurata mi serve sapere
questo, questo e questo''. Il problema non è sbagliare la stima, è non darsi i mezzi
per azzeccarla.

\paragraph{Come si fa a far entrare questi metodi nella routine? È solo
determinazione?} È un misto. Il relatore risponde da persona con ADHD, con la
tendenza a distrarsi e a fare un mucchio di \emph{overthinking} (``come dicevo ieri,
anche ordinare la pizza è overthinking''). È un po' disciplina, ma soprattutto far
diventare qualcosa un'\textbf{abitudine}: lo sforzo per fare le checklist era molto
più alto un anno prima, oggi è un'abitudine. Di solito servono \textbf{circa tre
mesi} perché qualcosa diventi un'abitudine\footnote{Stima come riportata in lezione.
Gli studi più citati sulla formazione delle abitudini indicano tempi molto variabili
da persona a persona, con una mediana intorno ai due mesi.}: la teoria dice che chi
inizia oggi ad andare in palestra ogni giorno, dopo tre mesi non penserà più ``oggi
devo andare in palestra'', ma la palestra farà parte della sua giornata. Non è più
uno sforzo, ma ``è così che lavoro''.

\paragraph{Chi fa davvero time management?} Prima di passare al brainstorming il
relatore torna alla domanda iniziale: qualcuno, ora che sa che cosa intendeva,
cambierebbe risposta? Uno studente risponde che lo fa, ma senza checklist, perché
non ne sente il bisogno (non gli capita spesso di dimenticare le cose); la
\emph{must do list} invece la segue sempre, e le cose meno importanti capitano nelle
giornate in cui ne salta qualcuna. Va benissimo: non esiste un unico modo, ognuno
troverà il proprio. Il relatore porta il suo esempio perché sa che è servito a molte
persone con cui ha lavorato; con chi faceva fatica ad applicare un metodo (come
chiedeva prima uno studente) ne hanno trovati altri, perché ciò che funziona per uno
non è detto che funzioni per un altro. L'importante è che \textbf{ognuno trovi il
proprio metodo}, perché la soluzione non è non fare time management: ci si sentirà
più ``liberi'' improvvisando la giornata, ma non la si sta sfruttando a dovere.

E non si parla dei video sul \emph{mindset} in cui ci si alza alle cinque, si fa la
checklist e si fa colazione leggendo le mail: organizzarsi meglio la giornata
\textbf{dà più tempo libero}. E per il relatore il tempo libero delle persone è
fondamentale, perché è ciò che permette loro di lavorare meglio.

% ---------------------------------------------------------------------
\section{Una parentesi: tempo libero, ferie e benessere}
\label{sec:z2-benessere}

\subsection{Le tre condizioni del relatore}

Il relatore è un sostenitore della \textbf{settimana lavorativa di quattro giorni},
che ha applicato in varie aziende; dove non ci è riuscito, ha introdotto almeno
\textbf{venerdì liberi}, due volte al mese o una volta al mese, cioè giornate in
cui non si fa niente. Come le pause brevi sono importanti per essere creativi,
anche le pause lunghe, cioè giornate intere, lo sono per essere creativi ed
efficienti. È lo stesso motivo per cui, dopo la Games Week, dava una settimana
libera alle persone (sezione~\ref{sec:z1-fasi-team}): ha bisogno che siano
riposate, e quella settimana ``persa'' è in realtà un investimento per le
settimane successive. In chat qualcuno lo paragona al ``metodo finlandese'': è un
misto tra quello e il metodo del buonsenso.

Il relatore viene spesso chiamato a \textbf{cambiare le culture aziendali}, ed è un
po' il suo focus. Come un allenatore che accetta una squadra solo se può portare il
suo staff (il preparatore, il tecnico), lui chiede tre elementi:
\begin{enumerate}
  \item possibilmente \textbf{quattro giorni lavorativi}, o altrimenti giornate
        libere per dare spazio alle persone;
  \item \textbf{ferie illimitate}: nessun limite massimo, ma un \textbf{limite
        minimo garantito}. Ferie illimitate non devono voler dire poter scegliere
        di non prenderle: ``devi prendere almeno X giorni, fino a infinito'';
  \item la \textbf{feedback culture}.
\end{enumerate}
Senza questi tre elementi non garantisce il risultato: è un modo di gestire le
aspettative fin dall'inizio, come si diceva il giorno prima.

\paragraph{La feedback culture.} Il perché è stato spiegato nella prima giornata, ma
il relatore lo sottolinea: un ambiente in cui ci si può fidare del team vale
moltissimo. Quando va a fare lezione allo IULM, anche in settimana, non deve
preoccuparsi di che cosa stia facendo il team mentre lui non lavora (``oggi è
domenica e non ho questa preoccupazione a prescindere, ma anche se fosse mercoledì
non l'avrei''): sa che stanno facendo del loro meglio, e non gli interessa sapere
altro.

\paragraph{Ferie e permessi illimitati.} Nell'esperienza del relatore, sia a Riot sia
a Blizzard, avere ferie illimitate significava spesso \emph{prenderne meno} di quelle
obbligatorie: c'è una vocina che dice ``ti serve davvero un giorno di ferie oggi?
Tanto sono illimitate, te le prenderai quando ti serviranno'', e sapere di averne
illimitate spingeva a lavorare di più. Per lui era altrettanto dannoso, e per questo
ha introdotto un \textbf{minimo obbligatorio di 24 giorni} di ferie.

C'è stato però un effetto non previsto, e molto desiderato: in tutte le aziende in
cui ha introdotto ferie e permessi illimitati, \textbf{i giorni di malattia sono
calati del 70\%}\footnote{Dato riferito dal relatore sulla base della propria
esperienza, non verificabile.}. La gente si ammala di meno? O piuttosto il giorno in
cui ci si sveglia e si pensa ``oggi non ne ho voglia, mi do malato'' diventa il
giorno in cui si prende un permesso? Si smette di usare come malattia un giorno che
non era di malattia.

\paragraph{L'\emph{Hangover Day}.} Sia a Riot sia a PG Esports il relatore aveva
introdotto quello che a Riot, ``essendo Riot a Dublino'', si chiamava
\textbf{Hangover Day}\footnote{Nella trascrizione ``Angover Day''.}: una giornata
all'anno in cui si aveva il diritto di telefonare alle nove del mattino e dire ``oggi
non vengo perché non ne ho voglia'' (tradotto: perché ieri sera ho bevuto troppo).
Altrimenti quella giornata sarebbe diventata un giorno di malattia o di ferie. Il
fatto è che \textbf{quasi nessuno lo usava}: il relatore stesso a volte ci ha
pensato, ma poi la mattina si chiedeva ``voglio davvero giocarmelo oggi? Tutto
sommato potrei anche andare'', e ci andava. Ma anche solo sapere di averlo lo faceva
stare tranquillo. Allo stesso modo, con le ferie illimitate, chi voleva prendersi un
mese lo faceva (``buon per loro''), chi non lo faceva andava bene lo stesso:
l'importante era che staccassero. C'era un \textbf{benessere più diffuso} nel team,
perché sapevano che, volendo, avrebbero potuto.

\paragraph{Il \emph{Neko Day}.} Il primo punto, i giorni liberi, era chiamato
\textbf{Neko Day}, da \emph{neko}, ``gatto'' in giapponese: il giorno in cui si può
fare il gatto, rannicchiati davanti al camino a non fare niente. Perché
introdurlo? Chi lavora dal lunedì al venerdì passa di solito il sabato o la domenica
a fare le commissioni che non può fare nei giorni lavorativi: la spesa, le bollette,
il bollo dell'auto. In questo modo sfrutta meno il weekend. Con un venerdì libero al
mese (o due, o tre) le persone facevano quel giorno le commissioni del weekend, e
sabato e domenica si riposavano davvero (``anche non fare niente andava bene,
chiaramente''). Il lunedì arrivavano molto più cariche e riposate. Questo ha aiutato
tantissimo a mantenere un team sano, fisicamente e mentalmente.

\paragraph{Il supporto psicologico.} A PG Esports il relatore ha introdotto anche il
\textbf{supporto psicoterapeutico offerto dall'azienda}: sei sedute da un'ora con
uno psicoterapeuta, online. Era il periodo in cui si parlava del bonus psicologo
statale, ``che poi non è mai arrivato''\footnote{Così in lezione. In realtà il bonus
psicologo è stato istituito in Italia nel 2022, ma con fondi limitati rispetto alle
richieste e con forti ritardi nell'erogazione; il relatore si riferisce
probabilmente a questo.}, e lui ha pensato: l'idea è buona, perché non introdurla a
livello aziendale? Erano gli anni dei lockdown: l'azienda stava perdendo molte persone
che andavano in depressione o non avevano più voglia di lavorare, perché avevano
difficoltà a vivere quel periodo (``io stesso non è che l'abbia vissuto benissimo'').
Qualcuno ne ha usufruito e gli ha detto: ``mi stavo per dimettere, poi ho risolto le
cose e ho preferito restare''.

Secondo il relatore oggi c'è una consapevolezza maggiore da parte delle aziende
sulla \textbf{salute mentale}: vuol dire dare più ferie, dare questo tipo di
supporto, dare più \emph{benefit} (anche i buoni pasto). Il benessere non è solo
economico, è anche mentale, ed è una cosa positiva da sapere.

% ---------------------------------------------------------------------
\section{Brainstorming: obiettivo e problemi tipici}
\label{sec:z2-brainstorming}

\subsection{Perché il brainstorming riguarda tutti}

Secondo il relatore questa è la parte che servirà di più, sia nei gruppi e nei
progetti del percorso di studi (il gioco da tavolo, il gioco digitale) sia nelle
aziende. Rispetto al time management ha una pertinenza più immediata: tutti
hanno esempi di \textbf{brainstorming andati male}, da cui si è usciti esausti, con
la sensazione di non aver ottenuto niente, o di non aver parlato abbastanza (``ha
parlato solo quello''), o di non essere sicuri della scelta. Il brainstorming lascia
una sensazione molto chiara, a fine meeting, che qualcosa non torna; il time
management dà piuttosto la tristezza del ``non ce l'ho fatta a fare quello che
dovevo oggi, come mai?''.

\subsection{Generare idee, non decidere}

Il primo chiarimento riguarda \textbf{l'obiettivo}, che il relatore vede spesso
travisato. L'obiettivo del brainstorming è \textbf{generare idee}; invece diventa
spesso decisionale, con lo scopo di scegliere che cosa fare. Ma allora non sarebbe
un brainstorming, sarebbe un \emph{decision making meeting}. Generare molte idee e
scegliere la migliore nello stesso meeting non funziona né per generare idee né per
scegliere la migliore, perché sono \textbf{due metodi che si contrappongono}. Il
brainstorming va fatto quando si devono generare idee \emph{per poi} scegliere che
cosa fare: deve venire \textbf{prima} del \emph{decision making}, come primo passo
della fase di esplorazione di un nuovo progetto.

\subsection{``Non esistono cattive idee''?}

Bisogna evitare che i membri del team \textbf{autofiltrino} le proprie idee per paura
di essere rifiutati o presi alla leggera. Da qui nasce la frase ``non esistono
cattive idee, buttate fuori tutto'', che è vera solo fino a un certo punto. Le
cattive idee esistono, oggettivamente; quello che non si deve fare è
\emph{colpevolizzarle}. Né si deve per forza perseguire ogni opzione: parlare di
opzioni che sono effettivamente brutte è uno spreco di tempo, e il brainstorming ha
anche un problema di gestione del tempo. Il relatore vede spesso brainstorming
gonfiati a ore e ore in cui si parla anche di cose che non meritano nemmeno dieci
minuti, come i feedback buoni ma inutili del giorno prima. \textbf{Trattare tutte le
idee allo stesso modo è sbagliato}, perché non hanno la stessa qualità. Bisogna
trovare un modo di valutarle senza che chi le ha proposte si senta preso di mira,
bullizzato o sminuito.

Per questo a ogni brainstorming dovrebbe seguire un \textbf{follow-up}, il momento in
cui si decide che cosa fare: le due parti vanno separate, altrimenti ``si crea solo
casino''. Il vero lavoro arriva quando le idee vengono messe a terra: c'è la parte di
brainstorming e di \emph{theorycrafting}, e la parte esecutiva viene dopo.

\subsection{I problemi più comuni}

\begin{description}
  \item[Le prime idee monopolizzano la sessione] Le prime due idee ``sparate a
        caso'' sono quelle su cui ci si concentra per la maggior parte del tempo,
        sviscerandone i dettagli (``intendevi così o così?''), per poi chiudere con
        ``altre idee? Questa, questa e questa? Sì, no, ok, andiamo avanti''. Si
        parla troppo delle prime idee e poco delle ultime.
  \item[Poche persone parlano per la maggior parte del tempo] È un effetto
        collegato al precedente.
  \item[\emph{Anchoring}] Ci si lega alle prime idee, e spesso si votano o si
        scelgono non perché siano le migliori, ma perché sembra di conoscerle meglio
        di quelle menzionate in trenta secondi alla fine. Non c'è nessun motivo per
        pensare che un'idea detta alla fine sia peggiore: era solo arrivata dopo.
        C'è un problema di \emph{timing}. Anzi, un'idea venuta istantaneamente è
        forse quella meno ragionata, perché ci si è pensato meno tempo.
  \item[Silenzi imbarazzanti] Spesso le persone arrivano al meeting senza nemmeno
        sapere che è un brainstorming. Quando il relatore convoca un brainstorming,
        comunica prima a tutti il \emph{topic}, il tipo di idee cercate, i paletti e
        le restrizioni, così che anche chi è meno portato all'improvvisazione arrivi
        con qualche idea pronta. Se la gente arriva impreparata rischia di non dire
        niente per un'ora, e allora perché è stata invitata? Non è un difetto loro:
        tanti manager dicevano al relatore ``ci sono persone inadatte ai
        brainstorming perché non sanno improvvisare''. No: si sta facendo
        brainstorming con un metodo che premia solo chi sa improvvisare ed è
        abbastanza estroverso da parlare sul momento. Non significa che chi non
        parla non avrebbe potuto contribuire, ma che non gli è stato dato modo di
        farlo.
  \item[Estroversi e pensatori rapidi dominano] Le personalità estroverse e chi
        pensa velocemente tendono a dominare la conversazione, più o meno
        intenzionalmente, senza lasciare agli altri il tempo di contribuire. Non c'è
        uno spazio in cui un introverso possa dire ``veramente io volevo\dots'':
        ``no, no, stiamo parlando di questa bellissima idea, facciamo questa''.
\end{description}

% ---------------------------------------------------------------------
\section{Un metodo di brainstorming}
\label{sec:z2-metodo}

Il relatore propone una soluzione principale, declinata in vari modi, che tocca
tutti questi punti.

\subsection{Preparazione e ambiente inclusivo}

\paragraph{Arrivare con un obiettivo chiaro.} Bisogna arrivare al brainstorming con
chiaro e preciso che cosa si sta cercando, cioè \textbf{che cosa si sta cercando di
risolvere}, \textbf{perché} e soprattutto \textbf{con quali vincoli}. Si stanno solo
raccogliendo idee per un cliente? Si cercano idee fattibili, o idee innovative? Che
cosa si vuole produrre da questo meeting? Deve essere chiaro a tutti \emph{prima}
del meeting, non spiegato durante, perché altrimenti si ricade nel problema delle
persone non pronte a intervenire.

\paragraph{Non giudicare le idee degli altri.} Al relatore è capitato spessissimo di
sentire un'idea terribile da una persona e, mezz'ora dopo, l'idea perfetta dalla
stessa persona. Ma se all'idea terribile tutti rispondono ``che stupidaggine, per
favore stai zitto'', quella persona non dirà la successiva, perché non vuole
rivivere l'esperienza. Giudicando le idee degli altri, di fatto, li si zittisce.
Le idee migliori nascono quando si propone quella che si pensa sia l'idea perfetta,
magari non è buona e non viene scelta (``fa niente''), si ascoltano anche le altre e
poi la si \textbf{riformula} integrando i pareri sentiti nel frattempo. Sono queste
\textbf{idee combinate} a essere spesso le soluzioni migliori.

\subsection{Separare generazione e discussione: \emph{brainwriting} e
\emph{rapid ideation}}

Con il classico ``non ci sono cattive idee, sparate'', succede spesso che uno prenda
la parola, spieghi la sua idea e il suo concetto, e gli altri comincino subito a
fargli domande (``ma quindi la vedi così?''): si sta unendo la persona che
\emph{genera} l'idea con la persona che la \emph{discute}, nello stesso momento.
Invece bisogna \textbf{dividere la generazione dalla discussione}: un momento in cui
si generano idee e un momento in cui si discutono. È il principio del
\textbf{brainwriting}. Funziona molto quando nel team c'è una figura con una
posizione gerarchica più alta: se quella persona esprime un'idea e i due momenti
sono uniti, tutti tenderanno ad allinearsi. E si vuole che \emph{tutti} possano
avere idee.

Lo strumento è la \textbf{rapid ideation}. Il relatore di solito dice: ``per i primi
cinque minuti non parliamo; scrivete su dei post-it il pitch delle vostre idee''. Se
ne possono scrivere quanti se ne vuole, dieci o uno, e si mettono da parte. Ma il
tempo è di \textbf{cinque minuti}: non dieci, non due, e soprattutto non mezz'ora.
Il \textbf{limite di tempo breve} è importantissimo, perché costringe,
consapevolmente o meno, a filtrare le idee mentalmente e a scrivere solo quelle che
si ritengono migliori: magari non è chiaro perché si preferisca un'idea all'altra,
ma il cervello sta facendo il lavoro in background e fa venire in mente le cose più
importanti o più giuste. È ottimo anche per i team che tendono a distrarsi, perché
sono solo cinque minuti.

\subsection{\emph{Dot voting}}

Poi si passa al \textbf{dot voting}:
\begin{enumerate}
  \item si attaccano \textbf{tutte} le idee su una lavagna fisica, un muro o una
        lavagna digitale (come Miro\footnote{Nella trascrizione ``mirror board'':
        si tratta della lavagna digitale Miro, già citata nella prima
        giornata.});
  \item si \textbf{uniscono le idee che hanno lo stesso concetto di base}. Se si sta
        pensando a un evento dal vivo e la proposta del relatore è un parco di
        dinosauri, mentre quella di un collega è un parco di dinosauri ``del
        paleolitico'' con i tirannosauri, le si raggruppa sotto ``parco di
        dinosauri'': non sono due idee separate. Un altro gruppo potrebbe essere,
        ad esempio, un parco a tema Avengers;
  \item ognuno ha \textbf{da tre a cinque voti} e va alla lavagna a mettere i propri
        pallini sulle idee che gli piacciono e di cui vorrebbe sapere di più.
\end{enumerate}
Il risultato tipico: un'idea con molti voti, una con un paio, altre con un numero
discreto.

\paragraph{Che cosa si vede subito?} Il relatore lo chiede alla classe. Si vede una
\textbf{preferenza} per un tema o per un'idea. Non c'è \textbf{polarizzazione} su una
o due idee iniziali, perché c'è varietà. Il voto è \textbf{anonimo}: non ci sono nomi.
E soprattutto manca un'altra cosa importantissima: \textbf{il giudizio negativo
espresso}. Se c'è un'idea con sette voti, una con cinque e una con tre, è
ragionevole parlare prima di quelle, perché tutti hanno espresso interesse e vogliono
saperne di più; ma nessuno sta dicendo ``non voglio assolutamente parlare di questa
perché mi fa schifo''. È semplicemente oggettivo che quell'idea non ha ricevuto voti.

\paragraph{Autovotarsi.} Il sistema permette anche di votare la propria idea, e non c'è
nessun problema: se la propria idea ha un pallino e gli altri ne mettono sette su
un'altra, il proprio voto è ininfluente. Il relatore vedeva spesso che all'inizio la
gente tendeva ad autovotarsi, poi capiva che era inutile, perché non sarebbe mai stato
il voto decisivo, e tanto valeva votare le idee che interessavano davvero. La
tendenza spariva con il tempo.

\paragraph{Che cosa si è ottenuto.} Tre ore di meeting sono state filtrate in una
discussione concentrata solo sulle idee che tutti ritengono interessanti, con un
\emph{agreement} di fondo: non si parla dell'idea numero uno perché è uscita per prima,
non è una questione di \emph{timing} ma di qualità. Il relatore ha vissuto tanti
meeting con la frustrazione interna del ``stiamo sprecando tempo su questa idea,
quando vorrei dire la mia che secondo me è meglio''. Qui non si può sprecare tempo,
perché si parla di qualcosa che tutti vogliono sentire.

\paragraph{E se si formano sottogruppi che si votano a vicenda?} Uno studente chiede se
non possano esistere ``cordate''. Anche se esistessero, significherebbe che il gruppo di
quel meeting vuole parlare di quella cosa, e va bene: l'importante è non sprecare tempo
su idee che non interessano a nessuno. Se si investono due ore, quarantacinque minuti
su ciascuna di tre idee vanno benissimo; dieci minuti per ogni idea no. E comunque il
voto è anonimo: non si può sapere che un'idea è di Gianni per dare supporto a Gianni,
quindi un eventuale sostegno reciproco sarebbe involontario.

\paragraph{E chi è lento a elaborare?} Uno studente osserva che il problema degli
introversi potrebbe ripresentarsi con i post-it: con un tempo così ristretto, chi è più
lento potrebbe andare ancora più in panico. Per questo bisogna \textbf{arrivare
preparati}: si conosce il topic in anticipo, quindi almeno un paio di idee già
formalizzate si hanno in mente. Magari non se ne scriveranno quindici, ma un paio sì.

\subsection{Quanti voti e quali varianti}

Dare più pallini a testa sarebbe controproducente? Il relatore indicava da tre a
cinque voti. Chiede alla classe quanti siano i gruppi di lavoro: dai quattro agli otto,
alcuni gruppi del secondo anno sono nove. La sua prassi:
\begin{itemize}
  \item \textbf{dai cinque partecipanti in su}, \textbf{tre voti} a testa, più che
        sufficienti per popolare molto la lavagna;
  \item \textbf{fino a cinque partecipanti}, \textbf{cinque voti} a testa, che
        permettono anche di variare un po'.
\end{itemize}
Esistono varianti: ad esempio tre voti di colori diversi, il blu per l'idea più
innovativa, il verde per quella più bella, il giallo per quella più facile, così da
qualificare \emph{perché} un'idea piace. Se l'obiettivo è proporre a un cliente tre
pitch (uno facile, uno medio, uno avanzato), la votazione è già fatta. All'inizio,
però, meglio partire da tre o cinque voti generici; con il tempo si imparerà ad
accelerare.

\subsection{Il moderatore e gli invitati}

\paragraph{Serve un moderatore?} Può essere utile. Il relatore fa spessissimo il
moderatore: partecipa ai brainstorming senza fare altro che moderare, ``perché so che
le persone con le idee fighe sono loro e non io''. Ci si può anche autogestire. Nei team
piccoli non ha senso ``perdere'' una persona per la moderazione; nei team più grandi
(dai cinque in su, sette-otto persone), se qualcuno dice ``non penso di avere idee
particolari su questo, vi modero così vi aiuto'', è un ottimo modo di contribuire.

\paragraph{Chi invitare?} Dipende da che cosa si cerca. Per inventare un nuovo gioco il
relatore invita i game director, gli art director e magari il \emph{publishing manager}
per un input commerciale; per decidere il \emph{combat system} invita magari il senior
game designer e il game director. La domanda che si fa chi convoca il meeting è: dato
l'obiettivo, chi è più ragionevole che possa contribuire?

Un 3D artist della classe chiede se allora gli artisti 3D non vengano inclusi. Dipende:
se l'obiettivo è definire la linea creativa dello stile dei personaggi, il relatore
invita anche il 3D artist, e non solo i \emph{concept artist}. Il 3D artist può dare un
input diverso dal concept: vota un'idea perché per lui è \emph{più fattibile}. Un
concept può essere bellissimo ma ``un inferno'' da modellare, e chi sa di doverlo fare
entro giugno non vota un'idea che sa già infattibile.

In realtà si può invitare chiunque: per un pitch su un'area dedicata al basket alla
Games Week, il relatore ha invitato un tecnico di regia che gioca a basket, per avere
anche la prospettiva di chi ha familiarità con il tema. A uno studente che porta
l'esempio di un elemento di interfaccia sul personaggio, il relatore risponde che,
volendo, si trova un motivo per invitare chiunque: l'importante è avere chiaro
\textbf{chi è necessario}. Il gruppo di lavoro può essere di otto e il brainstorming di
dodici, con input esterni a chi poi farà il lavoro. Basta non esagerare: se la maggior
parte dei voti arriva da chi non farà quel lavoro, si vota su qualcosa che poi non
faranno loro, e non ha più senso.

% ---------------------------------------------------------------------
\section{Dalla votazione alla decisione: la matrice impatto--sforzo}
\label{sec:z2-impatto}

\subsection{Massimo impatto, minimo sforzo}

Fatto questo, non si è finito: ora si sa quali idee interessano di più, ma non si è
ancora deciso che cosa fare. Serve la parte decisionale. L'obiettivo di qualsiasi
azienda (ed è il motivo per cui al relatore non importa se il programmatore fa il test
da casa cercando su Google) è \textbf{massimizzare l'impatto e minimizzare lo sforzo}.
Vale per i progetti, i task, il recruitment, per qualsiasi cosa: ``non perché sono
pigro, ma perché sarebbe stupido fare diversamente''.

Le idee selezionate si mettono quindi in un'altra matrice, con due assi:
\begin{itemize}
  \item sull'asse verticale il \textbf{valore}, cioè quanto \textbf{impatto} avrebbe
        l'idea (in alto l'impatto alto);
  \item sull'asse orizzontale lo \textbf{sforzo} (\emph{effort}), da basso (a sinistra)
        ad alto (a destra).
\end{itemize}
In alto a sinistra ci sono quindi le idee ad alto impatto e basso sforzo, in alto a
destra quelle ad alto impatto e alto sforzo, e così via.

\paragraph{L'esempio Adidas.} Se il cliente è Adidas, con la linea di abbigliamento per
il basket, un'idea a basso valore potrebbe essere una campagna su TikTok in cui gli
utenti condividono i loro momenti con un pallone da basket; un'idea ad alto valore
potrebbe essere fondare un parco divertimenti dove i giocatori NBA vengono a sfidarti
uno contro uno. Questo non dice che una sia meglio dell'altra: dice che l'impatto è
diverso. Allo stesso modo si valuta lo sforzo: la campagna social è a basso sforzo, il
parco divertimenti ad alto sforzo per molti motivi (i permessi di costruzione, i
giocatori NBA).

\subsection{Posizionare le idee, insieme e con precisione}

Il gruppo che a quel punto è \textbf{il gruppo decisionale}, cioè chi farà il lavoro,
deve posizionare sulla matrice le idee ritenute migliori, e deve essere \textbf{d'accordo
tutto il gruppo}: ``siamo tutti d'accordo che il parco dei dinosauri di Stefano e Andrea è
ad alto valore ma anche ad alto sforzo? Sì? Allora lo mettiamo qui''. E bisogna indicare
\textbf{un punto preciso}, non solo il quadrante: piazzare un'idea in un punto non è come
piazzarla poco più in là. Il risultato è una sorta di \textbf{mosaico} delle idee sul
grafico.

\subsection{Leggere la matrice in diagonale}

Davanti a una situazione tipica, con varie idee sparse, il relatore chiede alla classe:
se foste i project manager, se doveste fare voi queste cose, \emph{quali} fareste? Al
plurale, perché spesso da un brainstorming si esce con più idee realizzabili e si deve
scegliere su quali concentrarsi. Le risposte variano: c'è chi sceglie le due con più
impatto, chi aggiunge che dipende da tempi, budget e risorse, chi preferisce quella più
innovativa ma anche più rischiosa.

Il metodo che secondo il relatore va applicato è \textbf{procedere in diagonale}, dalle
idee a basso sforzo e alto impatto verso quelle ad alto sforzo e basso
impatto\footnote{Nella slide le idee erano rappresentate da post-it colorati (rossa,
rosa o salmone, azzurra, verde acqua, viola, blu, gialla), e le indicazioni del
relatore si riferiscono a quei colori. Dalla sola trascrizione non è possibile
ricostruire con certezza la posizione di ciascun post-it, e alcune indicazioni sui
colori risultano contraddittorie; qui è riportato il principio generale, che il
relatore chiarisce nelle risposte alle domande.}:
\begin{enumerate}
  \item un'idea a \textbf{basso sforzo e alto impatto} (nella slide, quella rossa) va
        fatta \textbf{per prima, sempre}: ``non dovrebbe essere nemmeno in discussione'';
  \item poi, a seconda delle risorse e del tempo, si possono fare
        \textbf{contemporaneamente} le idee che stanno sulla stessa diagonale
        successiva, che hanno cioè un rapporto simile tra impatto e sforzo;
  \item poi le successive;
  \item le idee ad \textbf{alto sforzo e basso impatto} (nella slide, la gialla)
        \textbf{non vanno nemmeno considerate}. Nessuno, infatti, l'aveva scelta.
\end{enumerate}
Per le scelte intermedie valgono due criteri, chiariti rispondendo a uno studente che
non capiva perché alcune idee fossero preferibili ad altre: \textbf{a parità di sforzo si
sceglie l'idea con più impatto}; \textbf{a parità di impatto si sceglie quella con meno
sforzo}. Con risorse infinite si farebbero anche le idee più costose; ma di solito le
risorse sono limitate, e allora tra due idee con lo stesso sforzo e un impatto molto
diverso la scelta è evidente. Se ci piacerebbe fare un'idea più ambiziosa ma abbiamo
risorse solo per una meno costosa sulla stessa diagonale, va bene lo stesso: abbiamo
comunque stabilito che è meglio delle idee sulle diagonali successive. Quanto fare dipende
dalla \textbf{bandwidth} a disposizione.

\paragraph{``La gialla era la mia idea: devo offendermi?''} No. Qualcuno ha detto che
l'autore fa schifo? No. Si è concordato che la sua idea, per quanto votata (``qualcuno ha
votato per la mia idea''), è piuttosto difficile da realizzare rispetto al suo impatto.
Non si sta dicendo che non sa più fare il suo lavoro. Anzi, se lui stesso dovesse
scegliere tra la propria idea gialla e la rossa di un altro, sceglierebbe la rossa: non è
nemmeno una questione di discussione. La matrice rende \textbf{visivamente evidente} che
cosa fare, e permette di aggirare molte discussioni che non sono più soggettive:
oggettivamente, dovremmo fare così.

\subsection{Il ruolo del moderatore e dell'esperto}

È qui che lavora il moderatore: ``non è più una discussione; visto il posizionamento,
dobbiamo fare queste cose''. Non si decide a maggioranza dove mettere le idee: se uno
posiziona un'idea e un altro non è d'accordo, difficilmente dirà ``la rossa dovrebbe stare
al posto della gialla''; dirà ``secondo me è un po' più a destra, un po' più in basso'', e
si fanno gli aggiustamenti. Se invece nascono discussioni del tipo ``la rossa dovrebbe
stare al posto della gialla'', vuol dire che non si è definita bene l'idea e non è chiaro
quanto sia davvero difficile da realizzare.

Oppure interviene \textbf{l'esperto}: è l'esempio del concept artist e del 3D artist. Il
concept artist dice che un'idea è a basso sforzo e alto impatto; il 3D artist risponde ``ne
condivido l'impatto, ma dovendola fare io ti assicuro che lo sforzo non è basso, è molto
alto''. Allora è un altro discorso, perché l'esperto sta dicendo che lo sforzo è diverso
da quello stimato. Ma alla fine tutto deve essere concordato.

\subsection{Chi decide}

\paragraph{Il gruppo decisionale può essere diverso da quello del brainstorming?} Sì: a
questo punto escono quelli invitati come ``consulenti'', che hanno dato idee ma non devono
poi realizzarle. Farli votare anche qui significherebbe perdere il collegamento tra chi
deve fare le cose e che cosa si decide di fare. In quel momento non hanno responsabilità:
hanno dato un input, e se decidessero che cosa devono fare gli altri dovrebbero esserne
i manager. Non dovrebbero poter dire agli altri che cosa fare se non ne hanno la
responsabilità. Vale per tutti i team, piccoli e grandi; secondo il relatore chi decide
dovrebbe essere \textbf{al massimo una decina di persone}. In un team di dieci che fa il
brainstorming tutto insieme, va benissimo che chi vota coincida con chi ha partecipato:
l'importante è che \textbf{chi vota sia chi farà il lavoro}.

\paragraph{Lo sforzo è il costo o la fatica?} Uno studente chiede se l'alto impatto
includa anche il costo. No, il costo sta nello sforzo, che per il relatore va
\textbf{ragionato in risorse}: denaro, tempo e così via. Un'idea può avere un basso costo
monetario ma un alto costo in tempo, e allora va bilanciata e si troverà a metà. È lì che
si spostano le idee e che avviene la conversazione: i tecnici dicono che lo sforzo tecnico
è un certo, poi arriva l'\emph{accountant} che dice che a livello economico è un altro, e
la si sposta.

\paragraph{E il piacere di fare il progetto?} In un gruppo di due o tre persone, osserva
uno studente, si finisce per votare anche per preferenza personale: quando si fa un gioco
bisogna avere anche il piacere di farlo. Ma \textbf{anche quello è impatto}: l'impatto
sulla motivazione. Come lo sforzo sono le risorse, l'impatto è l'impatto \emph{su di te,
sul team, sul mercato}, su ciò che si vuole. Il relatore lo ha tenuto volutamente
generico: può essere quanti soldi si fanno o quanto è utile per il proprio CV, e
l'impatto cambia tra un progetto accademico e uno professionale.

% ---------------------------------------------------------------------
\section{Consigli pratici e domande sul brainstorming}
\label{sec:z2-qa-brainstorming}

\subsection{Musica e cambio di scenario}

``Sembrano tutte stupidaggini, ma vi assicuro che aiutano.''

\paragraph{Musica di sottofondo.} È dimostrato che aiuta\footnote{Le affermazioni sugli
effetti della musica (chiave maggiore, rumore ambientale) sono riportate come dette in
lezione e non sono state verificate.}. Meglio una musica con \textbf{tempo e volume
costanti}: niente \emph{bass drop}, niente variazioni di tema o di ritmo. Va bene la
musica fatta di rumore ambientale a basso volume; è stato rilevato anche che le canzoni
in \textbf{chiave maggiore} aiutano la parte creativa a produrre idee più originali.
Meglio la \textbf{musica strumentale}, senza testo: per quanto ci si impegni a non
ascoltarlo, il cervello cercherà di capire le parole. È un effetto marginale e non
succede niente se capita la canzone con il testo, ma si mette in conto che ``l'uno per
cento del cervello'' sarà impegnato a capirlo, come quando in treno si ascolta, che lo si
voglia o no, la telefonata di qualcun altro. È \emph{min-maxing}.

\paragraph{Cambiare scenario.} A Riot c'era una stanza dedicata al brainstorming, pensata
per essere creativa: le finestre sul fiume, le porte a vetri colorati, tutto fatto per
ispirare creatività. Ma fare brainstorming sempre nella stessa stanza generava sempre lo
stesso tipo di idee: anche solo cambiare ambiente aiutava ad averne di diverse. Il
relatore spesso spostava il meeting, ad esempio nello Starbucks sotto l'ufficio. Lo
spazio fisico gioca un ruolo importante, ma anche quello virtuale: se lo si fa sempre su
Discord, si può provare con ognuno uno sfondo diverso e simpatico. Oppure: ha sempre
moderato la stessa persona? Si prova con un altro moderatore. Si cambiano i fattori
perché i risultati non siano sempre gli stessi.

Uno studente chiede se il cambio di scenario possa essere anche qualcosa come passarsi una
palla per avere la parola, o un altro espediente per tenere viva l'attenzione. Con questo
modello, che non richiede di prendere la parola, l'esempio della palla perde senso, ma come
esempio va benissimo: anche il comportamento fisico e il \emph{modus operandi} sono un ottimo
cambio di scenario. Come lo sono cambiare il moderatore o \textbf{cambiare l'obiettivo}: se
il cliente chiede X, gli si darà un'idea X, ma si può provare anche a votare per un'idea Y,
facendo più votazioni con vari criteri di valutazione.

\subsection{Domande}

\paragraph{L'idea scelta deve essere approvata da qualcun altro?} Se un manager dice ``fate
questo brainstorming senza di me'', sta \textbf{delegando la scelta}: ``mi fido di quello che
sceglierete''. Avendo preparato prima limiti, paletti e \emph{constraint}, si aspetta che la
soluzione portata sia in linea con essi; ma se non partecipa deve fidarsi della scelta. Se
poi la scelta debba passare anche da un livello superiore dipende dall'azienda e dagli
accordi precedenti.

\paragraph{Si può fare un brainstorming ibrido, tra presenza e remoto?} Uno studente fa
l'esempio di un gruppo in cui uno è a New York e gli altri a Milano. Sì, senza problemi: si
fa una call, il portatile punta verso il muro con le idee, o le si mostra fisicamente con
il cellulare. Chi è da remoto scrive i suoi post-it, ad esempio mandandoli su WhatsApp al
moderatore o a qualcuno presente che li scrive per lui, e vota dicendo ``mettimi un pallino
qui, qui e qui'' (in quel caso il voto non è anonimo, ma si può fare tranquillamente); se
deve parlare, parla.

\paragraph{E se i colleghi non riconoscono o non condividono il metodo?} Ci deve essere un
\textbf{agreement}: non si può arrivare e dire ``facciamo così'' se gli altri non vogliono.
Questo metodo funziona solo se tutti sono d'accordo a usarlo. Il relatore ne ha proposto
uno tra i cento metodi di brainstorming esistenti, un buon compromesso tra semplicità ed
efficienza (ce ne sono di più semplici e di molto più efficienti ma più complessi). Se il
gruppo non ne vede il valore, al massimo insisterebbe per \textbf{fare almeno una prova},
altrimenti diventa chiusura mentale; se poi per quel team non funziona, se ne cercano
altri.

\paragraph{E nei gruppi molto piccoli, dove diventa un duello?} Uno studente ha lavorato in
gruppi di tre o anche di due persone, dove il brainstorming diventa quasi un duello in cui
ognuno difende la propria idea. Bisogna \textbf{aspettare la discussione}: invece di ``scegliamo
tra la mia idea e la tua'', si cerca di concordare quale idea abbia l'impatto maggiore e
quella lo sforzo minore. E se non si è d'accordo, bisogna \textbf{portare esempi}: se uno
dice ``questa idea costa di più'', l'altro chiede ``in che senso? Che cosa costa di più?''. Se
non sa rispondere, non si può fare qualcosa senza sapere come va a finire: meglio rimandare
la discussione a quando si avranno le informazioni. A volte la domanda giusta è \emph{perché
non abbiamo queste informazioni?}: forse è \textbf{prematuro} decidere, perché non si hanno
ancora tutti i fattori.

Nei team piccoli non dovrebbe esserci un ``si fa come dici tu o come dico io'', una scelta a
pelle: si dovrebbero conoscere così bene le proprie proposte da poterle posizionare con
assoluta certezza sulla matrice, e avere anche la presenza di spirito di dire ``la tua è
chiaramente migliore, si fa la tua''. Si fa la cosa migliore per entrambi, che dovrebbe
essere oggettiva e non soggettiva. Non è facile arrivarci, ma deve essere l'obiettivo del
team.

% ---------------------------------------------------------------------
\section{Disagree and commit}
\label{sec:z2-disagree}

\subsection{Quando la scelta non ci convince}

Uno studente teme che chi vede scartata la propria idea perda motivazione, in un gruppo
piccolo come in uno grande. Non è un deragliamento: è una conseguenza potenziale di
qualsiasi scelta. Il relatore ha visto gente demotivata perfino dalla settimana di quattro
giorni: ci sarà sempre qualcuno più o meno motivato dalle scelte. Ma qual era l'alternativa?
Un brainstorming classico in cui tutti parlano, chi più chi meno, e si sceglie un'idea che
tutti pensano sia la migliore senza esserne davvero sicuri, perché non si è votato, né
espresso una preferenza, né discusso davvero le idee migliori.

Qui entra il concetto di teamwork del \textbf{disagree and commit}: ``non sono d'accordo con
questa decisione, non mi piace, non penso sia la migliore, ma la faccio''. Per due motivi:
\begin{enumerate}
  \item perché \textbf{è il mio lavoro}, e dovrebbe essere il punto principale;
  \item perché \textbf{voi ritenete che sia la scelta migliore}. Se si è la voce fuori dal coro
        e tutti hanno scelto un'altra strada, non si rema contro: è \emph{default to trust}.
        Se gli altri vi hanno visto più potenzialità, e (grande premessa) \textbf{hanno
        ascoltato} il proprio disaccordo e i suoi motivi, e hanno comunque scelto così, ci si
        fida di loro.
\end{enumerate}
Con una postilla: se le cose vanno male, \textbf{non si dà la colpa agli altri} e non si dice
``ve l'avevo detto'', come se non si facesse parte del team. Disagree and commit vuol dire ``mi
fido del team e facciamo questa cosa''; se poi si rivela sbagliata, \emph{il team} ha
sbagliato, e la volta dopo si cambierà il modo di scegliere. Altrimenti non è disagree and
commit, è ``\emph{disagree and ostracize}'': disaccordo più ostilità. Non bisogna cercare il
colpevole, perché si è deciso tutti insieme (sezione~\ref{sec:z1-default-trust}).

Se si pensa davvero che si stia sbagliando clamorosamente il posizionamento di un'idea, si
dice piuttosto ``non andiamo avanti finché non definiamo bene lo sforzo di questa cosa''; ma
bisogna saperlo giustificare, altrimenti è chi dissente a dover fare un passo indietro.

\subsection{La demotivazione fa parte del lavoro}

La demotivazione ci sarà sempre, perché non si sceglierà mai solo ciò che piace: per quanto
piaccia il proprio lavoro, ci saranno task sgradevoli e scelte che non spettano a noi. E
ci saranno scelte (``ve lo dico per esperienza personale'') che si è convinti siano
sbagliate, al punto di pensare ``piuttosto lascio il team'', ma che si finisce per accettare,
per poi rendersi conto che gli altri avevano ragione. Come c'è un \emph{bias} verso l'amico e
il nemico, c'è anche un \textbf{bias verso noi stessi}: crediamo di conoscere bene una cosa, e
se altre dieci persone dicono ``secondo me è fattibile'', bisognerebbe aprirsi alla
possibilità che ne sappiano più di noi.

``Non lo accetto ma lo faccio'' dovrebbe essere alla base del lavoro insieme, perché è
\emph{default to trust}: rifiutarsi (``non mi fido di voi, non lo faccio'') significa non
fidarsi del giudizio del proprio team. E dire ad alta voce ``non sono d'accordo, ma ti darò
comunque il massimo'' è la cosa che \textbf{fa guadagnare di più il rispetto degli altri}.
Non c'è niente di meglio che vedere qualcuno che aveva un'opinione diversa impegnarsi per
far vedere di avere torto, facendo andare bene la cosa. Per il relatore dimostra solo che è
una persona matura e importante; ha avuto persone che, passate da questa situazione, gli
hanno fatto pensare ``è molto più matura di quanto credessi'', e che ha valorizzato molto di
più in seguito, perché sa di potersi fidare di chi fa disagree and commit: è stato aperto
nella comunicazione, ha detto che cosa secondo lui non andava, ma non ha fatto niente per
remare contro, anzi ha dato il suo supporto. ``Non posso chiedere di più a una persona.''

\subsection{Chi arriva a lavori iniziati}

\paragraph{Se il team cambia in corso d'opera, chi subentra come si rapporta alle decisioni
già prese?} (per dimissioni, promozioni e così via). Deve \textbf{accettare lo status quo}. Al
relatore capita spesso di entrare in team con progetti già in corso: si fa dire che cosa è
stato scelto finora e va \textbf{in commit di default}: ``non ho preso parte a questa decisione,
quindi non sono in \emph{agreement}, ma ora faccio parte del team, quindi vado in commit''.

\paragraph{E se l'esperienza dice che una scelta ha uno sforzo molto diverso?} Ad esempio
che un'idea ``rossa'' ha in realtà lo sforzo di una ``viola''. Lo si fa presente, ma dipende dal
\textbf{tempismo}: se la scelta è appena stata fatta e c'è ancora modo di cambiarla, la si
segnala; se si è al 90\% del progetto, è controproducente dire ``fermi tutti, passo
indietro'', e a un certo punto bisogna rispettare la scelta del team (``avete scelto così,
andiamo avanti e vediamo come va''). E se si sa che, calcolando ferie e periodo estivo,
un'animazione prevista per giugno arriverà a settembre? È molto \emph{case by case}: a
seconda dell'apporto si può modificare il risultato o no. È come riaprire un processo se
emergono nuove prove: se dimostrano che l'imputato non è l'assassino, è giusto riaprirlo; se
dimostrano che l'assassino ha sparato con la sinistra invece che con la destra, conta fino a
un certo punto. È giusto segnalarlo, ma non sempre quello che si porta giustifica uno
``\emph{scratch that}, rifacciamo tutto'': diventa \textbf{un feedback, non una risoluzione}
(``state facendo così, questo è il mio feedback, poi vediamo come va'').

Con cinque minuti di anticipo si va in pausa pranzo, con l'accordo per il pomeriggio: circa
un'ora di presentazione più ``standard'' sul leading from within, poi domande e risposte,
``e quando finiscono le domande abbiamo finito''.

% ---------------------------------------------------------------------
\section{Leading from within: perché parlare di leadership}
\label{sec:z2-leadership}

\subsection{La ripresa del pomeriggio}

Alla ripresa, mentre si aspettano gli ultimi arrivati, il relatore chiede se durante la
pausa siano venute in mente domande. L'unica è logistica: le slide verranno date a fine
workshop? Sì: il relatore le passerà all'accademia (a Vanessa, Andrea ed
Enrico\footnote{Nella trascrizione l'accademia è chiamata ``Ive'' e ``Hive'', come nella
prima giornata; per i nomi dei referenti vedi la nota all'inizio del
capitolo~\ref{cap:z1-feedback}.}), che le farà avere alla classe. Gli appunti vanno
comunque presi, meglio su quello che dice che su quello che mostra.

L'ultima parte, dice, dovrebbe essere abbastanza veloce da presentare, ma vale la pena
fermarsi a parlarne, perché è un discorso più avanzato. È la logica chiusura di quanto
detto tra la prima giornata e la mattina: tirare le fila e mostrare come tutto faccia parte
di un unico filone di ragionamento e progettualità, con un unico scopo, un unico obiettivo e
un unico modo di ragionare.

La voce sembra a posto, ma il relatore è stanco dopo due giorni, come immagina lo sia la
classe, nello \emph{stretch} finale di due giornate molto piene. Sono cose che in condizioni
normali spiega al suo team nell'arco di settimane, con il tempo di metterle in pratica,
discuterle e approfondirle. Qui è tutto concentrato: servirà tempo per metterlo in atto e
fare le proprie prove, quindi non bisogna spaventarsi se non verrà tutto giusto la prima
volta.

\subsection{Il leader non è il più bravo}

Perché parlare di leadership in un workshop non rivolto solo ai leader? La risposta
completa arriverà più avanti, ma il punto di partenza è che il relatore \textbf{non crede
nei leader di team o di progetto che detengono da soli il sapere e il modo di fare le
cose}. Come detto il giorno prima, essere un leader, un manager o un capo non dovrebbe
voler dire essere il migliore in quel lavoro, ma essere \textbf{la persona più adatta a
ricoprire quel ruolo}: guidare un gruppo significa \textbf{occuparsi delle persone} al suo
interno.

L'obiettivo che il relatore si dà, e che secondo lui ogni manager dovrebbe avere in mente,
è che quando si torna dal lavoro ci si senta \textbf{ispirati, sicuri e realizzati}, e che
questo sia \textbf{un diritto naturale di tutti}, non un lusso riservato a chi ricopre certe
posizioni, a pochi fortunati o ``a qualcuno sì e a qualcuno no''. Ha notato,
specialmente in Italia, la tendenza a vedere tutto ciò che viene prima del diventare
manager come la \emph{gavetta} necessaria per arrivarci. Invece chiunque faccia parte di
un'organizzazione, che sia un'azienda di videogiochi o altro, dovrebbe sentirsi soddisfatto,
ispirato e sicuro in quello che fa; e non è corretto che un'azienda tratti bene le persone
``da manager in su''. Anche quando non lo dicono, a volte i fatti sono quelli: da manager si
viene trattati in un certo modo e si ha accesso a certe cose, altrimenti no. Non si tratta di
\textbf{appiattire la gerarchia}: si tratta di assicurarsi che, pur dando cose diverse ai
diversi livelli, a tutti i livelli venga riconosciuto uno strato comune di
\textbf{rispetto, fiducia e professionalità}.

\subsection{Le carriere divise: manager o \emph{subject matter expert}}

Uno studente chiede conferma: quindi un \emph{lead artist} non è necessariamente più bravo
di un senior a fare l'artista? Esatto. E qui il relatore apre una parentesi. In molte
aziende, e soprattutto in quelle in cui tende a lavorare lui, la motivazione principale per
diventare manager è avere una promozione e quindi un aumento, più che il desiderio di
prendersi cura del team e delle persone. Come si ovvia al problema? \textbf{Dividendo le
carriere}. Arrivati al livello senior, il relatore fa con ognuno dei suoi senior
\textbf{``the Talk''}, per esporre le due opzioni che hanno davanti:
\begin{description}
  \item[Busta A: diventare manager] È il percorso standard: si valuta se si è adatti a fare
        il leader e a guidare un team, e si danno gli strumenti per farlo, perché
        ``nessuno nasce leader imparato'': prima di provarlo sulla propria pelle nessuno sa
        che cosa voglia dire. Per quanto si possa essere carismatici, fare il leader è
        un'altra cosa, e non vuol dire solo essere carismatici. Ci sono formazione,
        strumenti e un percorso di crescita.
  \item[Busta B: diventare \emph{subject matter expert}] Una persona che va oltre il senior
        nella stessa disciplina: diventa \textbf{l'esperto dell'intera azienda} su quella
        disciplina, perché la sua priorità non è aiutare gli altri ma essere il migliore.
        ``Non c'è assolutamente nulla di male.'' Il suo ruolo è essere il più specializzato
        in quel campo, il punto di riferimento da cui tutti vanno a farsi spiegare le cose.
\end{description}
Il concetto è abbastanza nuovo, e il punto chiave è che l'azienda \textbf{paga il
\emph{subject matter expert} e il leader esattamente allo stesso modo}: non c'è un guadagno
economico nello scegliere di fare il manager invece dello specialista, e il relatore valuta
quel tipo di conoscenza quanto il saper gestire persone. Così chi vuole essere il migliore,
o anche solo un aumento, può continuare a fare ciò che sa fare e avere l'aumento; chi vuole
l'aumento \emph{e} prendersi cura delle persone sceglie l'altra strada. Così si fa una scrematura sul
\emph{motivo} per cui si vuole fare il manager. Sembra una scelta piccola, ma è molto
importante capire se ci si vuole occupare di se stessi o degli altri. Non è un giudizio
(``non sto dicendo che sia meglio l'uno o l'altro''): le due cose \textbf{non convivono}, e
bisogna sapere che cosa si vuole fare, sapendo che l'azienda tratterà allo stesso modo le
due scelte.

\paragraph{Si può essere entrambi?} Si può scegliere ad esempio il percorso da lead artist
e voler comunque continuare a migliorare le proprie capacità. Ma, come nelle simulazioni di
conversazione del giorno prima, il relatore chiederebbe: ``se avessi le risorse per fare solo
una delle due cose, quale sceglieresti?''. Perché non le si avrà: non ci sarà il tempo e il
modo di fare bene entrambe. Il suo ``cinismo'' arriva a dire che bisogna scegliere, perché
tenersi buone entrambe le strade porta a farle male tutte e due. Preferisce un ottimo
\emph{people manager} o un ottimo specialista; le figure intermedie, in questo campo,
rischiano di diventare pericolose per sé e per il team.

\paragraph{Specializzazioni verticali.} La strada dello specialista permette
specializzazioni molto verticali: in chat qualcuno cita il \emph{mecha designer}, ma può
esserci anche lo specialista di animazioni delle espressioni facciali. Naturalmente l'azienda
deve avere bisogno di quella competenza (non ci si specializza in cose che non servono), e non
è detto che ne serva uno solo. È un percorso di carriera parallelo e alternativo a quello del
manager ``standard''.

% ---------------------------------------------------------------------
\section{Il ``momento Piero Angela'': la biologia del gruppo}
\label{sec:z2-biologia}

\subsection{Perché una digressione scientifica}

Il relatore annuncia un ``piccolo excursus biologico'', fondamentale per capire dove vuole
arrivare, e prova a mettere un sottofondo musicale adatto: la condivisione dell'audio non
funziona, quindi manda il link in chat perché ognuno lo faccia partire per conto suo. Un
studente lo battezza subito ``momento Piero Angela''. Il relatore cerca di far combaciare le
motivazioni aziendali del suo metodo di leadership e di progettualità con \textbf{basi
scientifiche}: in passato ha collaborato con varie università che si occupavano di
neuromarketing e neurobiologia, per avere una conferma o una smentita del fatto che ciò che
ipotizzava ``non fosse una cavolata colossale''. Semplificando moltissimo, si guardano solo
alcune sostanze che si traducono direttamente nel modo in cui si vive l'essere in un gruppo,
in un team, in un'azienda\footnote{L'impianto di questa parte (endorfine, dopamina,
serotonina, ossitocina e cortisolo, il villaggio protetto dai pericoli esterni, il ``cerchio
di sicurezza'') ricalca da vicino quello del libro di Simon Sinek \emph{Leaders Eat Last}
(2014), che il relatore cita poco più avanti. Le descrizioni biologiche sono molto
semplificate, come avverte lo stesso relatore, e vanno lette come tali.}.

\subsection{Endorfine}

Le \textbf{endorfine} sono rilasciate per \textbf{mascherare il dolore}. Lo sa chi corre, chi
fa pesi o ginnastica, ma anche chi ha riso così tanto per una battuta da avere male ai
fianchi: è quella sensazione di dolore in cui però si è contenti di provarlo. Il \emph{runner's
high} è l'esempio più calzante: essere distrutti dopo una corsa ma soddisfatti. È la risposta
biologica allo sforzarsi in cose che sono per il nostro bene, un modo con cui il corpo ci
``convince'' a farle.

\subsection{Dopamina}

La \textbf{dopamina} è più familiare: è la \textbf{ricompensa per un'azione}. Dalle cose più
banali (si mangia e si riceve dopamina) a quelle viste in mattinata: anche completare una
checklist rilascia dopamina (sezione~\ref{sec:z2-pianificare}). Viene rilasciata quando si
vede l'obiettivo, come il traguardo durante una corsa; è la stessa cosa che accadeva ai
cacciatori del villaggio in tempi antichi, che rilasciavano dopamina quando vedevano le tracce
e venivano così incentivati a seguirle. Sono concetti fondamentalmente ``animali'', ma si
applicano tuttora anche in azienda.

\paragraph{Una nota di colore: la dipendenza.} La dopamina \textbf{crea dipendenza}: molti casi
di ossessione per alcol, nicotina, gioco d'azzardo e social media sono problemi legati alla
dipendenza da dopamina. Secondo il relatore non è un caso che le diagnosi di ADD e ADHD siano
aumentate del 60\% negli ultimi dieci anni\footnote{Nella trascrizione ``EDD e HD''. Il dato
e il nesso causale con la maggiore disponibilità di alcol, nicotina, gioco d'azzardo e social
media sono riportati come detti in lezione e non sono stati verificati; l'aumento delle
diagnosi viene in genere attribuito anche ad altri fattori, come la maggiore consapevolezza e
il cambiamento dei criteri diagnostici.}, perché buona parte di quegli stimoli è diventata
ampiamente disponibile proprio in quel periodo. Anche la dipendenza dal lavoro, a volte, ha la
stessa origine: si vuole il rispetto che deriva dal lavoro, il \emph{rush} che deriva
dall'averlo.

\subsection{Il villaggio di diecimila anni fa}

Si fa poi un salto di diecimila anni, in un villaggio. Gli esseri umani, per quanto più o meno
introversi o estroversi, sono \textbf{animali sociali} che tendono a vivere esperienze di
gruppo; le aziende nascono da un gruppo di persone che si danno un obiettivo condiviso, come
in antichità un gruppo di persone che voleva sopravvivere insieme. Quello che non è cambiato è
che \textbf{fuori dal gruppo ci sono dei pericoli}: tigri, lupi, draghi (``do per scontato
che nell'antichità ci fossero i draghi, perché D\&D mi insegna così''). Per evitarli, i gruppi
tendono a riunirsi e a creare difese intorno a sé: da qui le prime difese dei primi villaggi,
per isolarsi dai pericoli esterni. ``Vi giuro che avrà tutto senso.''

\subsection{Serotonina}

La \textbf{serotonina} viene rilasciata anche \emph{indirettamente}: non deve per forza
riguardare noi. Quando si assiste a una cerimonia o celebrazione pubblica che include una
persona cara, si rilascia serotonina come se si fosse noi l'oggetto del riconoscimento. È uno
dei motivi per cui si fanno diplomi e lauree ``all'americana'', con la consegna sul palco
davanti a tutti: gli studi avrebbero mostrato che anche i genitori in platea rilasciano
serotonina, come chi riceve il premio. Lo stesso vale per parenti, amici, matrimoni: sono
esperienze condivise, e si tende a condividere in pubblico riconoscimenti e successi perché in
tutti i partecipanti viene rilasciata serotonina.

Purtroppo ha a che fare anche con i beni materiali: la tendenza all'\emph{hoarding}, cioè a
collezionare eccessivamente qualcosa, nasce dal fatto che quell'attività rilascia serotonina,
ingannando la mente nel credere di aver ottenuto un successo.

È chiamata anche \textbf{la sostanza della leadership}, perché aiuta a riconoscere in qualcun
altro chi ha valore per il gruppo. Nel villaggio, quando arriva un pericolo, c'è sempre un
gruppo di persone che dice ``ce ne occupiamo noi'': il gruppo sa che si stanno prendendo un
rischio, e la serotonina è il modo per riconoscere che meritano rispetto. Il relatore consiglia
il libro di Simon Sinek \emph{Leaders Eat Last}\footnote{Nella trascrizione ``Leadership
Lust''. Il titolo è \emph{Leaders Eat Last: Why Some Teams Pull Together and Others Don't}
(2014).}: i leader mangiano per ultimi,
all'opposto di quanto accadeva in antichità, quando al ritorno dalla caccia mangiavano per
primi i cacciatori, che dovevano essere in forze per la caccia successiva, poi donne e
bambini, poi gli altri. C'è un \textbf{riconoscimento sociale}: tu fai qualcosa per proteggere
il gruppo, noi ti riconosciamo qualcos'altro in cambio.

Le aziende sono strutturate in modo molto diverso? È una domanda retorica: quando si dice che
serve un capo, un manager, un rappresentante, si sta dicendo che quella persona dovrebbe avere
l'interesse di proteggere il gruppo e di prendersi cura di tutti. Per questo sapere che il
proprio capo guadagna di più non fa particolarmente arrabbiare: lo si vede come quasi naturale,
perché più responsabilità si ha verso il gruppo più si guadagna, dato che quello dovrebbe essere
il suo scopo. Ed è anche il motivo per cui il relatore si arrabbia molto quando legge di CEO o
\emph{executive} che, a fronte di licenziamenti di migliaia di persone, si assegnano dei bonus.
È una rabbia non del tutto razionale, che sente da dentro: è come se si stesse violando un
\textbf{contratto sociale}. Il contratto è ``tu ti prendi dei rischi per me e mi difendi, io ti
riconosco più cibo, più soldi''; quando si abusa di quella posizione per un beneficio personale
a scapito di chi si dovrebbe proteggere, tutti lo vedono come un'ingiustizia, perché si viene
meno al contratto stabilito diecimila anni fa con la nostra biologia.

\subsection{Ossitocina}

L'\textbf{ossitocina}\footnote{Nella trascrizione ``ossidocina''.} è la \textbf{sensazione di
sicurezza} che si genera quando si è vicini a qualcuno. Non è un caso che tra i legami di
affinità e collaborazione più forti ci siano quelli dei team che si uniscono contro un
\textbf{nemico comune}: i team che superano insieme avversità e difficoltà ne escono molto
legati. Tra le persone che il relatore sente più spesso in questo periodo ci sono quelle con cui
ha condiviso momenti difficili, un cattivo manager, una brutta esperienza lavorativa: sa che
erano lì con lui. Vale anche per i veterani di guerra o per i sopravvissuti a una stessa
tragedia, che sviluppano un rapporto molto forte (se ne è parlato molto negli Stati Uniti a
proposito dei veterani), perché hanno rilasciato immense dosi di ossitocina. È anche ciò che fa
accettare subito il figlio alla madre quando lo prende in braccio dopo il parto, ed è il motivo
per cui, a migliaia di anni di distanza, si stringono ancora accordi con una stretta di mano o
un abbraccio: c'è sempre qualcosa di fisico, generato dal bisogno di ossitocina.

Viene generata anche quando si compiono o si \textbf{assiste ad atti di generosità} senza
aspettarsi nulla in cambio. Se il relatore dicesse di aver donato cinquanta euro a
un'associazione di beneficenza, la reazione sarebbe ``ok, bene''; se dicesse che ogni sabato
pomeriggio fa volontariato, lo si vedrebbe come qualcosa di molto più prezioso, e la cifra
potrebbe crescere parecchio prima che qualcuno la consideri più preziosa dei suoi pomeriggi:
tutti sanno che la vera donazione è \textbf{il tempo} investito, che vale più del bene
materiale. E basta assistere: vedere qualcuno che aiuta un'anziana ad attraversare la strada fa
stare meglio anche noi. È percepita a livello di gruppo. Infine, l'ossitocina rafforza il
\textbf{sistema immunitario}: sentirsi sicuri e tranquilli fa stare meglio e, in un certo senso,
ammalare di meno.

\subsection{Cortisolo}

L'opposto è il \textbf{cortisolo}: la sensazione di stress, paranoia e paura, la risposta
\emph{fight or flight} (o scappo o combatto), con battito accelerato, muscoli in tensione e così
via. Tra i suoi effetti:
\begin{itemize}
  \item \textbf{blocca il sistema immunitario}: se si sta per morire, al corpo non importa di
        prendersi un raffreddore. È il motivo per cui le persone stressate si ammalano di più;
  \item \textbf{blocca l'ossitocina}: non ci si può sentire tranquilli quando si è in ansia;
  \item soprattutto, e questo è il punto più importante per le aziende, è
        \textbf{contagioso}. In qualsiasi documentario, quando una gazzella che bruca alza la
        testa e percepisce un movimento nella savana, scappano tutte, non solo lei. Succede,
        non sorprendentemente, anche in ufficio: quando qualcuno ha paura perché potrebbero
        esserci licenziamenti, o perché il manager è molto aggressivo, la paura si diffonde nel
        gruppo ed è molto dannosa.
\end{itemize}
È anche l'effetto alla base della \textbf{sindrome dell'impostore}: sentirsi costantemente
inadeguati perché si teme che succeda qualcosa significa vivere in uno stato di stress. Ed è il
motivo per cui chi è insoddisfatto del lavoro tende a stare peggio anche fisicamente e ad
ammalarsi di più.

\begin{center}
  \begin{tabular}{@{}p{0.17\textwidth}p{0.36\textwidth}p{0.37\textwidth}@{}}
    \toprule
    \textbf{Sostanza} & \textbf{Che cosa fa (in sintesi)} & \textbf{Esempi dalla lezione} \\
    \midrule
    Endorfine & mascherano il dolore di uno sforzo che ci fa bene &
      \emph{runner's high}, ridere fino ad avere male ai fianchi \\
    \addlinespace
    Dopamina & ricompensa un'azione, si attiva vedendo l'obiettivo; crea dipendenza &
      spuntare una checklist, le tracce per il cacciatore, gioco d'azzardo, social \\
    \addlinespace
    Serotonina & riconoscimento sociale, anche indiretto; ``sostanza della leadership'' &
      lauree e premi in pubblico, chi difende il villaggio, \emph{hoarding} \\
    \addlinespace
    Ossitocina & sicurezza e legame; rafforza il sistema immunitario &
      il nemico comune, i veterani, la stretta di mano, il volontariato \\
    \addlinespace
    Cortisolo & stress, \emph{fight or flight}; blocca difese immunitarie e ossitocina; è
      contagioso & la gazzella che scappa, la paura dei licenziamenti, la sindrome
      dell'impostore \\
    \bottomrule
  \end{tabular}
\end{center}

% ---------------------------------------------------------------------
\section{Domande sulla biologia del gruppo}
\label{sec:z2-qa-biologia}

\paragraph{Il gioco d'azzardo include i gacha e le loot box?} Assolutamente sì, anzi ne sono il
\textbf{driver principale}: i gacha, i FIFA Ultimate Team, lo ``sbustare'' da \emph{Magic} a
\emph{Hearthstone}, e in generale il ``vediamo se stavolta mi va bene'' sono tutti gioco
d'azzardo. Lo studente aggiunge che molti giochi mobile sono costruiti su questi meccanismi e che
c'è chi gioca più a quelli che ai giochi ``normali''. È il motivo per cui alcuni stati hanno
dichiarato illegali certi giochi\footnote{Il caso più noto è quello del Belgio, che dal 2018
considera gioco d'azzardo le \emph{loot box} a pagamento; anche i Paesi Bassi le hanno
contestate.}.

\paragraph{Come si genera l'ossitocina nei rapporti da remoto?} ``Bellissima domanda.'' È il motivo
per cui, come detto il giorno prima, il relatore insiste perché i meeting si facciano almeno in
\textbf{webcam}, perché i riconoscimenti siano dati e i successi celebrati \textbf{pubblicamente}
(quando non mette troppo a disagio la persona), e in generale cerca di creare \textbf{situazioni
condivise}. Anche lavorando da remoto, insiste per avere il \emph{team building} dal vivo e almeno
\textbf{un incontro all'anno} di persona in ufficio, perché crea situazioni di rapporto diretto.

\paragraph{Come si contrasta il cortisolo?} In chat qualcuno suggerisce la meditazione: sì,
assolutamente. Per il lavoro, la risposta arriva subito dopo nella presentazione.

\paragraph{Il nemico comune è una strategia che un manager può usare?} Spesso \textbf{il manager è
il nemico comune}, senza bisogno di sforzarsi. Ma il relatore chiede di capire bene il punto:
alcuni manager applicano consapevolmente il concetto ``mi ergo io a nemico del team, così sono
tutti uniti contro di me e lavorano bene'', e ne fanno uno stile di management. Decidono di essere
il \emph{bad cop}, il \emph{villain}, ed è così che ottengono risultati. È una scelta lecita e
funziona; il relatore preferisce fare diversamente, ma non per questo dice che non funzioni. Molte
volte si sottovaluta quanto sia consapevole: ``non si rende conto di quanto sia antipatico a
tutti''; forse se ne rende perfettamente conto e lo fa apposta per essere il bersaglio. Il relatore
l'ha visto fare spesso come soluzione disperata, quando il team era troppo impegnato a litigare al
suo interno: ``piuttosto che vederli litigare tra loro, gli do io un motivo per odiare me, così si
uniscono contro di me''. Non lo consiglia assolutamente, ma l'ha visto funzionare.

Uno studente racconta che nella sua azienda il manager fa vedere \emph{il lavoro stesso} come il
nemico (``prima finiamo, meglio è''). È un modello più basilare ma, come dice lo studente stesso,
funziona.

\paragraph{Quanto è etico puntare sulla dipendenza del giocatore per fare soldi?} ``Non è
assolutamente etico, ed è assolutamente profittevole, altrimenti le aziende non lo farebbero.'' Il
relatore racconta una cosa che lo ha sconvolto quando l'ha scoperta, a Blizzard. Lavorava su
\emph{Hearthstone}\footnote{Nella trascrizione ``Arsdol''.} e si stava discutendo di come
monetizzare il gioco; venne preso come consulente un esperto, israeliano se ricorda bene, che
lavorava a Malta per società di poker online e di \emph{gambling}. Raccontò che molte aziende che
monetizzano sulla dipendenza dal gioco d'azzardo avevano sviluppato metodi più o meno etici (``più
spesso meno''). L'esempio culminante: analizzavano gli utenti occasionali, quelli che giocavano
magari una settimana al mese, per capire \emph{quando arrivava loro lo stipendio} (di solito si
deposita lo stipendio e poi si gioca quello che resta), per colpirli il giorno dopo con email di
offerte per tornare a giocare. Li si stimolava proprio nel momento in cui avevano disponibilità, sapendo
che erano sensibili, per aumentare le probabilità di riaverli come giocatori. ``È orribile. È
terrificante e, come tutte le cose terrificanti, è anche altamente funzionale.''

Gli ha fatto riflettere molto su dove ci si possa spingere per fare soldi. La versione etica, per
lui, è ad esempio far uscire un gioco sotto Natale, quando si è più disposti a regalarlo ai figli:
si stanno ottimizzando le proprie possibilità, e lo accetta. Si scade nel \textbf{predatorio}
quando si sfrutta la debolezza di qualcuno in un momento preciso per estorcergli soldi che non può
permettersi di perdere. Poi ognuno tragga le proprie conclusioni: ``non è un workshop sull'etica del
lavoro''. Non pensa di dover imporre la propria opinione, né si arroga il diritto di dire che chi lo
fa è una bestia senza cuore, visto che è il loro lavoro: ``chiunque riesca poi a dormire in pace con
se stesso, buon percorso''.

% ---------------------------------------------------------------------
\section{La piramide della fiducia}
\label{sec:z2-piramide}

\subsection{La leadership non è un titolo}

Chiuso il ``momento Piero Angela'', il relatore arriva al punto. La frase della prima giornata (``la
capacità di un gruppo di persone di fare cose straordinarie dipende dal modo in cui queste persone
operano insieme'') porta a capire che \textbf{la leadership non è un titolo lavorativo}. Essere
manager da contratto, o avere ``manager'' nel titolo, rende un manager, non un leader. Il relatore
non ha mai visto un team rispettare qualcuno solo perché da contratto doveva farlo: si può avere
paura o timore del proprio manager, ma non lo si deve rispettare perché è il manager. \textbf{Il
rispetto si guadagna}, e si possono influenzare le persone intorno solo attraverso il rispetto
guadagnato, quando si sono accumulati abbastanza esempi e \emph{case history} da giustificare
rispetto e fiducia. Altrimenti, dice, sarebbe il primo a chiedere ``perché vi fidate di me, se non
vi ho ancora dato motivo di farlo?'': non sarebbe razionale.

Da qui la \textbf{piramide della fiducia}, che si costruisce partendo dalla base.

\subsection{Alla base: l'integrità morale}

Per essere un buon leader o un buon manager bisogna \textbf{voler genuinamente aiutare gli altri
membri del team}. È il motivo del discorso sulle carriere divise: è a questo livello che si
screma chi vuole la promozione per i soldi. Chi vuole essere leader non deve cercare il guadagno
personale, ma quello del gruppo. Come detto il giorno prima, il relatore è pronto a litigare per il
suo team e a prendersi le colpe per il suo team, perché vuole il bene del suo team e sa che poi loro
si prenderanno cura di lui. Chi ha in mente questo sta in questa categoria. Poi:
\begin{itemize}
  \item \textbf{essere onesti}. Il relatore può affermare di non aver mai mentito
        intenzionalmente a un collega, a un cliente o a un giocatore. Gli è capitato di dare
        informazioni errate perché in quel momento sapeva quelle: ad esempio annunciare che
        un'espansione sarebbe uscita in un certo mese, e poi l'uscita è slittata\footnote{Nella
        trascrizione l'esempio non è coerente: prima si parla di un'uscita a novembre rimandata
        di una settimana, poi di ``pensavo fosse settembre''. Il senso resta chiaro.}. Ha detto
        una cosa sbagliata, ma tutti hanno capito che non stava mentendo: in quel momento pensava
        che fosse vera;
  \item soprattutto, \textbf{accettare la responsabilità dei fallimenti}: quando si sbaglia,
        saper dire ``ho sbagliato''. Il relatore lo ha messo perché ``non avete idea di quanti
        facciano fatica'' ad ammettere di poter sbagliare: c'è sempre una scusante, una
        mitigazione, un ``sì, ma''. Lui vuole sentire se è colpa tua o no;
  \item \textbf{chiedere scusa} e \textbf{ammettere di non sapere}: anche dire ``non lo so'' viene
        spesso evitato, come accettare un fallimento.
\end{itemize}

\paragraph{``Raccontami il tuo più grande fallimento''.} Per lo stesso motivo, nei colloqui il
relatore chiede spesso, più che ``parlami di un tuo grande successo'' o di un progetto andato bene,
\textbf{``mi racconti il tuo più grande fallimento?''}, e osserva come lo raccontano. Chi non si
aspetta la domanda spesso racconta di quella volta in cui ha provato in tutti i modi a fare in un
certo modo, ma il capo gli ha detto di fare diversamente e il progetto è andato male. ``Non mi hai
risposto: mi hai raccontato un progetto andato male, ma il tuo messaggio è che volevi fare
diversamente e qualcuno ti ha detto di fare così. Quindi non è un tuo fallimento.'' È come
rispondere a ``qual è la tua debolezza?'' con ``sono così attento ai dettagli che a volte\dots'':
si fa passare un punto di forza per una debolezza. Il relatore vuole un esempio in cui \emph{tu},
senza altri fattori, hai sbagliato, e come ti sei comportato. Gli dice molto su come quella persona
reagirà ai suoi feedback: se ha davanti ``un muro'' che dirà sempre ``non è colpa mia'', non è adatta
a un ambiente in cui il feedback è alla base di ogni comunicazione. Chi non ammette di aver mai
sbagliato in vita sua, lo scarta.

\paragraph{È possibile colpevolizzarsi troppo per un fallimento?} Sì, assolutamente. Rientra molto
nella sindrome dell'impostore, o nell'aver sempre vissuto rapporti di lavoro tossici: ci si abitua ad
accettare che ci vengano urlate contro le cose e che sia sempre colpa propria, e ci si autoconvince
che lo sia.

\subsection{Al centro: l'affidabilità}

Superato questo livello si arriva all'\textbf{affidabilità}. Per essere un buon leader non bisogna
essere il migliore, ma affidabili: basta essere riconosciuti come bravi in qualcosa, non
necessariamente i migliori, e quell'affidabilità ce la si porterà anche nella leadership. Significa:
\begin{itemize}
  \item \textbf{rispettare la parola data}: se si promette ``parleremo di un aumento al prossimo
        slot disponibile'', se ne parlerà. Non ci si dimentica, e non si dice poi ``in realtà non si
        poteva fare, ma ho voluto provarci fino all'ultimo''. Una promessa va mantenuta;
  \item \textbf{impegnarsi per il tempo necessario} a raggiungere gli obiettivi;
  \item \textbf{saper rifiutare un lavoro} quando necessario: è il saper dire di no della mattina,
        ancora la matrice di Eisenhower (sezione~\ref{sec:z2-eisenhower}). ``Sono sempre gli
        stessi concetti.''
\end{itemize}

\subsection{In cima: la leadership, cioè dare l'esempio}

L'ultimo gradino è la \textbf{leadership}\footnote{Nella trascrizione ``dealership''.}, cioè
\textbf{dare l'esempio}. Non bisogna essere i più bravi; bisogna saper \textbf{offrirsi volontari
non solo quando il progetto è un successo sicuro, ma anche quando c'è il rischio di un
fallimento}. Sono tutti bravi a offrirsi per il gioco nuovo che andrà sicuramente bene.

\paragraph{Il rappresentante dei licenziati.} Al relatore è capitato di offrirsi per rappresentare
il team durante i \emph{layoff}, con il rischio di essere ostracizzato dall'azienda; con quelle
persone è ancora in contatto. Quando Blizzard fece i licenziamenti nella sede francese, vicino a
Versailles\footnote{Nella trascrizione ``i layoff di Versailles''. La sede europea di Blizzard si
trovava a Vélizy-Villacoublay, nei pressi di Versailles; il relatore non precisa l'anno.}, che
coinvolsero quasi tutto l'ufficio, lui non era tra i licenziati, ma venne \textbf{eletto
rappresentante sindacale} di chi lo era. Per lui è stato molto importante: pur non essendo parte in
causa, lo ritennero degno della loro fiducia. Per quanto ne sapevano, avrebbe potuto negoziare e
accettare tutto quello che proponeva Blizzard, perché non ne era toccato: che si negoziasse un mese
di liquidazione in più o in meno non gli portava un soldo. Avrebbe persino potuto usare la
situazione per negoziare un aumento per sé (``gli facciamo accettare tre mesi invece di sette e a te
diamo un aumento''): non era fuori discussione. Se si sono fidati è solo perché c'era una piramide di
questo tipo e lui aveva dato loro abbastanza motivi per farlo.

\paragraph{Autorevolezza, non autorità.} A uno studente che lo riassume così, il relatore risponde
che è proprio questo: un carisma diverso dall'essere simpatici, l'essere un punto di riferimento,
\textbf{più un pilastro di supporto che il motore} che fa andare avanti le cose. Non il motore davanti
a tutti, ma il \textbf{meccanico} che, quando si rompe qualcosa, sai che non ti frega: ti sistema le
cose senza ``farti la cresta'' sulle pastiglie dei freni.

\paragraph{Non abbandonare nessuno.} E soprattutto \textbf{non abbandonare chi è in difficoltà}.
Sembra una banalità, ma è alla base del motto dei Marines statunitensi, \emph{no one left behind}:
sapere che, se si viene feriti, verranno a prenderti anche sotto i missili è un forte elemento
motivazionale per correre quel rischio. In azienda è il motivo per cui il relatore ribadisce sempre
che, se qualcuno è in difficoltà, la sua prima opzione non è licenziarlo ma \textbf{dargli una mano},
e lo ha fatto a volte anche all'eccesso, per far vedere agli altri che quando saranno in difficoltà
non li lascerà in mezzo a una strada. A volte ce l'ha fatta: ha esempi di persone ``a tanto così''
dal licenziamento recuperate per i capelli e diventate ottimi \emph{performer}. Altre volte la
persona non era recuperabile e alla fine è stata licenziata. Ma chi attraversa un periodo difficile
(il \emph{burnout}, la depressione, una relazione finita) sa che la prima risposta non sarà ``ci sei
servito finora, ora puoi andartene'', ma ``che cosa possiamo fare?''. E questo fa lavorare tutti
tranquilli.

% ---------------------------------------------------------------------
\section{Prendersi cura del gruppo}
\label{sec:z2-gruppo}

\subsection{La soluzione: più leader possibile}

In un caso del genere manager e leader coincidono, ed è l'ideale: un manager dovrebbe essere anche
un leader. Ma la vera soluzione è \textbf{avere più leader possibile}. Se l'obiettivo del leader è
prendersi cura degli altri, e il manager dovrebbe essere ``il primo tra i leader'', ancora meglio è
che \textbf{tutti siano leader del proprio gruppo}, perché la leadership non deriva da un titolo né
dall'essere il manager. Il relatore ha visto tanti gruppi andare avanti perché qualcuno al loro
interno era il leader, non perché lo era il manager.

Vuol dire mettere in pratica tutto il discorso della prima giornata su reciprocità, gentilezza,
aiuto e collaborazione. Si sa che, nel peggiore dei casi, c'è il manager, ma in condizioni normali ci
sono i colleghi: e allora si va da chi è in difficoltà a dire ``posso darti una mano?''. Così, quando
saremo noi in difficoltà, magari senza nemmeno rendercene conto, sarà qualcun altro a dirci ``ti vedo
in difficoltà, come posso aiutarti?''.

\paragraph{L'esempio è contagioso, nel bene e nel male.} Quando si ha una posizione anche solo di
rispetto agli occhi degli altri, tutto ciò che si fa ha effetto sugli altri. In negativo vale la
\textbf{slippery slope}\footnote{L'esempio del cartone abbandonato ricorda la cosiddetta ``teoria
delle finestre rotte'' (\emph{broken windows theory}), secondo cui i piccoli segni di degrado
tollerati favoriscono comportamenti peggiori.}: se si accetta un cartone di pizza abbandonato per
strada, domani se ne accetteranno due, dopodomani tre. Ma vale anche in positivo: chi vede qualcuno
fare qualcosa di positivo pensa ``lo voglio fare anch'io''. ``Come fai una cosa, fai tutto'', dice un
proverbio che il relatore attribuisce alla tradizione cinese\footnote{La massima (``how you do
anything is how you do everything'') è attribuita a fonti diverse, spesso alla tradizione zen;
l'attribuzione è incerta.}.

\subsection{Il villaggio di oggi}

Si torna al villaggio del ``momento Piero Angela'', con i difensori e le bestie (reali o no) intorno.
Oggi fuori dal gruppo non ci sono più tigri, ma le \textbf{deadline}, l'\textbf{overwork}, i
\textbf{layoff}, il manager che urla, le aspettative, i \textbf{KPI}. Avere manager che si prendono
cura del gruppo è un buon inizio, ma è ancora meglio avere \textbf{il gruppo che si prende cura del
gruppo}. Non bisogna più limitarsi a pensare che sia responsabilità del manager prendersi cura di
noi: è responsabilità nostra prenderci cura del gruppo di cui facciamo parte.

Che cosa serve? La piramide vista prima. Serve essere percepiti come \textbf{autentici}: il
relatore ormai riconosce abbastanza facilmente chi si vuole ergere a leader perché gli piace
\emph{sentirsi} leader, non perché vuole \emph{fare} il leader, e quella mancanza di genuinità non
aiuta. In generale bisogna \textbf{generare sicurezza} nel resto del gruppo.

\subsection{Smettere di dire ``ci penserà qualcun altro''}

Un esempio concreto è smettere di dire \textbf{``ci penserà qualcun altro''}. Nello sviluppo è
evidente: se si trova un bug che non riguarda il proprio lavoro, tutti sono d'accordo che vada
comunque segnalato, perché altrimenti potrebbe arrivare in \emph{live}, nella \emph{release}.

\paragraph{Le lasagne di PG Esports.} Circa un anno e mezzo prima, quando lavorava ancora a PG
Esports (che organizza anche eventi dal vivo), l'azienda ospitava in ufficio dei potenziali sponsor
venuti a conoscere il team per decidere se sponsorizzare la Milano Games Week. C'era un catering, con
i classici buffet. In mezzo alla confusione (una ventina di persone di PG Esports e una trentina tra
manager, sponsor ed \emph{executive}), il relatore vede per terra, in mezzo alla stanza, una mezza
porzione di lasagne abbandonata. Senza pensarci la raccoglie per portarla via: non è sua, ma è un
piatto in mezzo alla stanza. Nel tragitto fino al bancone del buffet, almeno altri sei tra executive e
sponsor che non lo conoscevano gli consegnano il loro piatto non finito, come a un cameriere: nella
loro testa l'unico motivo per cui qualcuno potesse raccogliere un piatto da terra era che fosse il
suo lavoro. Sul momento non ci fa caso; arrivato al bancone si gira e vede tutto il suo team che ride
della scena.

Torna da loro e ne nasce una discussione: ``avete visto quanto è stato educativo?''. Avevano visto
esattamente come ragionano i gruppi che non sono abituati a pensare ``adesso ci penso io'', ma ``ci
penserà qualcun altro''. In quelle persone non c'era malizia: sicuramente era un cameriere, perché
altrimenti farlo? ``Non è il mio lavoro, lo avrà fatto perché è il suo.'' E diventa un problema quando
la stessa tendenza (``non c'entra con il mio lavoro, ci penserà qualcun altro'') si applica a
qualsiasi altra cosa. Il relatore si aspetta che \textbf{tutti raccolgano quel piattino}: che chi vede
un bug, se non può correggerlo lo segnali, e se può correggerlo lo corregga anche se non è il suo
lavoro, ``perché è un piattino in mezzo alla stanza''.

\paragraph{Da reattivi a proattivi.} Significa passare da un atteggiamento \textbf{reattivo}, cioè di
fatto passivo, a uno \textbf{proattivo}: prevedere i problemi e cercarli intorno a sé per risolverli,
invece di aspettare che esplodano (``finché non esplode non è un problema mio''). Ragionare nell'ottica
del ``e se succedesse, e io me ne fossi accorto senza fare niente?''. Uno studente lo chiama
\emph{mindset da tester}: il relatore si aspetta che, se si vede qualcosa, ci si assicuri di non essere
gli unici ad averlo visto. A volte basta farlo presente: il relatore non potrebbe mai correggere un
bug, quindi se lo trova lo segnala, e va bene così. Ma non bisogna dare per scontato che ``è un bug così
evidente che qualcuno l'avrà già segnalato''. È il paradosso dell'incidente stradale: \textbf{più
l'incidente è grosso, più è probabile che nessuno chiami l'ambulanza}, perché tutti pensano che
qualcuno l'abbia già fatto\footnote{È il cosiddetto \emph{effetto spettatore} (\emph{bystander
effect}), studiato in psicologia sociale a partire dagli anni Sessanta.}.

\subsection{Domande sul ``me ne occupo io''}

\paragraph{Se correggo un bug, devo avvisare il collega?} Dopo una breve pausa per problemi audio
(``facciamo \emph{double testing}''), uno studente chiede se, correggendo il lavoro di un collega, il
dovere sia poi andare a dargli il feedback. A prescindere bisogna informare qualcuno; e la
conversazione dovrebbe essere: ``ho trovato questa cosa, vuoi che te la sistemi io?''. Correggere
senza dirlo a nessuno potrebbe essere un danno: magari era \emph{working as intended}, o qualcuno ci
stava già lavorando.

\paragraph{E se divento ``quello che risolve quel problema''?} Uno studente porta la sua esperienza:
in un'azienda in cui ha lavorato tre anni, risolvendo un problema facile che non era di sua
competenza è diventato, di fatto, l'incaricato di quel problema, distolto dalle sue mansioni. È
un'arma a doppio taglio. Per il relatore lì il problema è del \textbf{management}. In alcune aziende si
applica il criterio ``chi trova il problema lo risolve'' (``se te ne sei accorto tu, pensaci tu''), ed
è un altro modo di scaricare la responsabilità sugli altri. Non è ciò che dice lui: deve essere chi
trova il problema a valutare se è giusto e positivo che lo risolva lui. Il cartone di pizza lo
raccoglie perché ci mette dieci secondi; ma se ha trovato un bug nel sistema di pagamento che richiede
tre settimane di lavoro, e chi l'ha introdotto fa finta di niente, non deve risolverlo lui: deve far
presente che c'è un problema, che è stato ignorato e che va risolto (``pensateci voi, ma io vi dico
che c'è''). Va fatto caso per caso: se è immediato, facile e nel proprio interesse, lo si risolve.
Ma il relatore sconsiglia alle aziende la regola ``se lo trovi lo risolvi'', perché \textbf{disincentiva
a trovare i problemi}: qualcuno potrebbe non segnalarli per non avere altro lavoro e aspettare che li
trovi qualcun altro, ed è negativo quanto la situazione descritta. Per chiudere, ne discuterebbe con il
management: ``possiamo rivedere questa policy? Incentiva l'azione sbagliata''.

\paragraph{E se ci dicono di non mettere mano?} Uno studente che lavora come tester racconta che spesso
potrebbero correggere certe cose, ma viene loro detto di non toccarle. Esatto: si comunica. ``Ho trovato
questa cosa, ne siete consapevoli? Sì? Ci state già lavorando? No? Posso farlo io? Sì? Allora lo
faccio.'' Non bisogna esagerare con l'iniziativa: anche qui vale il \emph{default to trust}. Si parte
dal presupposto che qualcuno potrebbe averlo già visto e starci già lavorando, quindi si chiede e si
offre il proprio supporto.

% ---------------------------------------------------------------------
\section{Se si diventa manager}
\label{sec:z2-manager}

\subsection{Meriti, colpe e \emph{big picture}}

Una breve digressione per chi, un giorno, diventerà manager: che cosa cambierebbe? Oltre a tutto
quanto detto, ci sono il discorso su \textbf{meriti e colpe} e quello su \textbf{progetto e task}:
\textbf{delegare i meriti e prendersi le colpe}, e pensare al progetto più che ai singoli task, cioè
alla \emph{big picture}, invece che al ``voglio tornare a fare le cose che sapevo fare''.

È un altro problema dei manager promossi dal gruppo di chi sapeva fare bene un certo lavoro: quando
hanno l'occasione di tornare a farlo, lo fanno molto volentieri. ``C'è un task da fare, lo faccio io'',
magari con la scusa di dare l'esempio. No: si sta andando nella propria \textbf{comfort zone}. Il
relatore lo capisce (anche lui tornerebbe volentieri a fare ciò che faceva prima), ma bisogna accettare
che \textbf{quelle cose non le si deve più fare}.

\subsection{Il ruolo è più grande della persona}

Anche da manager, il ruolo che si ricopre per il gruppo è \textbf{più importante di sé}. Non importa
chi sia il community manager di Blizzard, di Riot o di un'altra azienda: conta ciò che rappresenta per i
giocatori in quanto community manager, non il fatto che sia lui piuttosto che un altro.

Prendersi cura del gruppo vuol dire applicare concetti molto semplici:
\begin{center}
  \begin{tabular}{@{}p{0.42\textwidth}p{0.48\textwidth}@{}}
    \toprule
    \textbf{Invece di dire\dots} & \textbf{\dots dire} \\
    \midrule
    ``Qui non si accettano errori.'' &
      ``Spiegami perché abbiamo sbagliato'': fammi capire perché è successo
      (sezione~\ref{sec:z1-default-trust}). \\
    \addlinespace
    ``Qui si fa come dico io.'' & ``Tu come faresti?'' \\
    \bottomrule
  \end{tabular}
\end{center}

\subsection{Il tetto della bravura del capo}

Il secondo punto è cruciale soprattutto quando si promuove il più bravo in un lavoro: se il criterio
noto è che si promuove il più bravo, il promosso lo sa, e la prima cosa che penserà è ``faccio lavorare
tutti come me, visto che sono il più bravo''. Il relatore chiede alla classe: che cosa si nega l'azienda
con questa dinamica? Anche ammesso che sia davvero il più bravo, se tutti fanno come lui diventeranno
tutti più bravi, no?

Il problema è che \textbf{si fissa come tetto massimo la bravura di quella persona}: se tutti fanno solo
come dice lui, si arriverà al massimo al suo livello, mentre ognuno ha il proprio metodo. Da qui quanto
detto il giorno prima sulla tendenza dei manager ad assumere solo persone peggiori di sé: se il criterio
di promozione è la bravura, si prenderanno sempre persone un po' meno brave per evitare che diventino il
proprio capo. Se invece le carriere sono divise (``se sei bravo fai questo, se vuoi fare il \emph{people
manager} fai quest'altro''), il manager è incentivato a prendere persone \textbf{migliori di sé}; e se,
come detto il giorno prima, il bonus dipende anche dalle performance del gruppo, è nell'interesse di
tutti che il gruppo sia pieno di superstar, un ulteriore incentivo ad assumere i più bravi.

\paragraph{L'opzione C.} C'è poi un'altra cosa, forse minima, ma per il relatore molto importante: si
perdono le \textbf{occasioni di innovare}. Nelle aziende precedenti era visto come il punto di riferimento
per il \emph{problem solving}: in un'emergenza, se c'era da scegliere tra A e B, ``decide Stefano''. Andava
benissimo, ma la tendenza era che a ogni emergenza venivano da lui. Così ci si negava \textbf{l'opzione C}:
quella a cui lui non avrebbe pensato e qualcun altro sì. Magari la prima volta sbaglia, la seconda anche,
e la terza ha un'idea migliore della sua, proprio perché lui non è più l'esperto di quel settore. Fissare
tutto come ``al massimo decide Stefano'' significa non accettare altre idee: un limite molto grosso
all'innovazione, all'aggiornamento e al fare le cose in modo più moderno e flessibile.

\subsection{Celebrare i successi e i fallimenti}

\paragraph{Celebrare i successi del team.} Il relatore vuole assicurarsi che tutti capiscano quando si è
ottenuta una vittoria o un successo: come detto il giorno prima, se non si sa che cosa si fa bene si
rischia di smettere di farlo. ``Questa cosa l'abbiamo fatta bene, complimenti, bravi tutti.''

\paragraph{Celebrare i fallimenti.} Sui fallimenti si era spinto ben oltre il semplice accettarli: aveva
indetto \textbf{meeting in cui si celebravano i fallimenti}, con le persone che raccontavano i propri errori
per distribuire quella conoscenza al resto del team. ``Non mi fa arrabbiare l'errore fatto, mi fa arrabbiare
lo stesso errore fatto quindici volte.'' Condividere un fallimento con il gruppo non significa ``guardate
questa persona sciocca che ha fatto questo errore'', ma ``questa persona ci sta permettendo di evitare lo
stesso errore in quindici'': sta salvando l'azienda dal ripeterlo quindici volte, perché la gente ha paura di
dire agli altri che ha sbagliato. Mostrare dove si è sbagliato diventava quindi un incentivo, un servizio
all'azienda. E si torna al punto di prima: se il proprio ruolo per il gruppo è più importante di sé, non
importa dover alzare la mano e dire ``ho sbagliato, ho fatto così'', perché gli altri penseranno ``grazie,
così io non lo faccio''. Si dovrebbe essere lieti che qualcuno abbia fatto quell'errore prima di noi, e
andare a ringraziarlo per averlo condiviso.

\paragraph{Condividere i successi.} Infine bisogna \textbf{dare il merito a chi lo ha}. Al relatore è
capitato spesso di ricevere ringraziamenti come ``grazie a Stefano e al suo team per aver fatto questa
cosa'', e la sua risposta è stata ``grazie, ma ha fatto tutto Gianni''. Non bisogna avere paura di dire
``grazie del feedback, ma è stato lui'': è giusto ringraziare lui e il team (il relatore si sente parte del
team, chiaramente). Non si ottiene niente dal dire ``grazie, riferirò i complimenti al team'': il relatore
non ha mai creduto nel ``riferire'' che il management ha detto ``ottimo lavoro''. Chiede al management di
dirlo \textbf{direttamente} al team, perché sentirlo dire dal management ha un impatto diverso che sentirlo
riferire da lui. In questo senso, osserva uno studente, ogni gruppo è il gruppo di sostegno di se stesso:
``assolutamente sì''.

\subsection{Quanto ci vuole?}

Non c'è una scadenza. Il rispetto non si ottiene pagando per fare prima, ``con il \emph{boost} a livello
60''. Servono tempo ed energia, ed è difficilmente misurabile: il relatore non può dire quanto il team si
fidi di lui in questo momento. Lo percepisce da come lavorano, ma lo si vedrà davvero solo quando ci sarà un
problema grosso; fino ad allora è tutta una sua percezione. Il team lo sente ripetere queste cose ogni giorno
e lo vede metterle in pratica ogni giorno, e lui può solo sperare che il messaggio passi; quando ci saranno
davvero ``gatte da pelare'' capirà come è messo, perché non c'è momento migliore di una difficoltà per vedere
come reagisce un team.

\textbf{Essere leader è prendersi un impegno a lungo termine verso i propri colleghi}: non ci sono modi per
velocizzarlo. Bisogna fare il proprio lavoro ogni giorno ``a testa bassa'', aiutando gli altri, facendolo
vedere ed \textbf{esplicitando perché lo si fa}. Il relatore lo fa quotidianamente: ``ti sto aiutando; nel
caso avessi un dubbio, lo faccio perché il mio lavoro è aiutarti e io ne traggo vantaggio dal fatto che tu
lavori meglio. Sono qui per te, ma lo faccio anche per me: non pensare che sia solo gentilezza. È il mio
lavoro, ma sappi che lo faccio spontaneamente''.

\subsection{L'ultima slide}

L'obiettivo dovrebbe sempre essere lo stesso: tutti dovremmo poter finire la giornata, staccare e tornare a
casa sentendoci
\begin{itemize}
  \item \textbf{ispirati}, perché abbiamo fatto qualcosa che ha un effetto tangibile;
  \item \textbf{sicuri}, perché sappiamo di avere colleghi, un leader e dei manager che si prendono cura di
        noi;
  \item \textbf{realizzati}, per quello che facciamo e perché è ciò che ci piace fare.
\end{itemize}
E questo dovrebbe essere \textbf{un diritto naturale di tutti}, non un ``sì, ma solo il manager'', ``sì, ma
solo se sei un esperto'', ``sì, ma solo se sei un game director''. Ognuno ha le proprie capacità: bisogna
solo trovare ciò che piace e che si vuole fare.

% ---------------------------------------------------------------------
\section{Domande e risposte finali}
\label{sec:z2-qa-finale}

Finita la presentazione, si passa alle domande: ``chiedete tutto quello che volete, perché è
l'ultima possibilità''. Le domande sono riportate per temi.

\subsection{Leadership e organizzazione del lavoro}

\paragraph{La leadership riguarda tutti o solo chi vuole diventare leader?} Il discorso sui
manager vale solo per i manager; ma il consiglio è \textbf{essere tutti leader del proprio gruppo,
da dentro}. Essere un leader significa sapere che gli altri possono fare affidamento su di sé, e
di queste figure in un team ce ne possono essere infinite: non è una gara a chi è il leader.
``Siate tutti una di quelle figure a cui gli altri guardano quando sono in difficoltà.''

\paragraph{Da manager non si hanno più orari e si fanno straordinari di continuo?} No: il relatore
ha notato di essere molto più efficiente come manager quando si dà \textbf{paletti molto precisi}.
Se stacca alle sei e alle sei e cinque arriva un messaggio, risponde ``se non mi dici che è
urgentissimo, non lo guardo nemmeno, lo leggo domani''; se non è urgente, lo guarda il giorno dopo.
Ma gestisce le aspettative: ``se è urgente ci sono e ci sarò sempre, ma non aspettatevi che faccia
orari eccessivi solo perché sono manager, così come non me lo aspetto da voi''. Si aspetta che
ognuno faccia ciò che deve; se è urgente tutti faranno di più, ma se non lo è, no.

\paragraph{L'orario davvero flessibile.} Qui il relatore aggiunge una cosa che si era dimenticato di
dire parlando di ferie e permessi illimitati: introduce anche l'\textbf{orario davvero
flessibile}. Non il flessibile di molte aziende, in cui si può iniziare alle 9 e finire alle 18
oppure iniziare alle 10 e finire alle 19: se oggi si vuole iniziare alle 10 e finire alle 15, e si è
fatto tutto quello che si doveva, va bene così. Ci si prendono le ore libere in più per altro, per
staccare, per riposarsi, sapendo che se la settimana dopo bisognerà lavorare dalle 9 alle 21, lo si
farà (o si dirà ``ieri ho fatto dalle 9 alle 21, domani faccio dalle 10 alle 15''). L'obiettivo dato
è un lavoro e un progetto, non un orario: si lavora per completare i propri task e rispettare le
deadline. L'orario indicato è solo indicativo, ed è quello dei meeting: \textbf{nessun meeting prima
delle 10 o dopo le 17}. Se poi qualcuno vuole lavorare dalle 10 alle 11 e poi di notte, perché quel
giorno preferisce lavorare di notte, più rilassato, faccia pure: l'importante è che trovi il tempo per
riposare e staccare. Ci si autogestisce senza timbrare cartellini; l'unico accordo è ``hai un task, va
consegnato entro il giorno X e mi garantisci che venga fatto; come ci arrivi non mi interessa''.

\paragraph{Le è mai capitato di essere aiutato da qualcuno, in un momento di difficoltà?} Sì: come
detto, questo comportamento è \textbf{contagioso in senso positivo}. Se ogni volta che qualcuno è in
difficoltà qualcun altro lo aiuta, prima o poi capita anche a chi aiuta. Il relatore non lo dice
``tanto per dire'': ha avuto momenti in cui faceva davvero fatica ad alzarsi dal letto per lavorare,
perché esausto o demotivato per altri motivi, e il team gli ha detto ``stai a letto, ci pensiamo noi
oggi''. Questo lo aiuta ad aiutare, perché sa che la cosa è bidirezionale.

\paragraph{Tra le soft skill del workshop, da quale partire?} \textbf{Il feedback}: saperlo dare e
ricevere. È la più difficile, molto più del ``voglio aiutarti'', ma è anche quella su cui si costruisce
tutto il resto: attraverso il feedback si costruisce la \textbf{fiducia}, attraverso la fiducia il
\textbf{rispetto}, e il rispetto diventa la base portante di tutto il resto.

\paragraph{E se il leader se ne va?} A una domanda su che cosa fare quando viene meno la figura di
riferimento del gruppo\footnote{La domanda non è leggibile nella trascrizione; dalla risposta si
capisce che riguardava il caso in cui il leader che incarnava la visione del gruppo se ne va.}, il
relatore risponde che la domanda da farsi è un'altra: \textbf{è cambiato l'obiettivo?} La direzione è
sempre la stessa? Se sì, chi è la persona più adatta a portarci lì? Non conta tanto chi guida, ma
l'obiettivo che ci accomunava, di solito la visione di un leader in cui tutti si allineano: ci
crediamo ancora tutti, o è cambiato qualcosa? È il momento di fare autoanalisi: ``ero qui per quella
persona o per la visione?''. Si dice, ed è vero, che \textbf{le persone non si dimettono dai posti di
lavoro, ma da capi che non rispettano più}, perché spesso l'unica interfaccia con l'azienda è il
proprio capo: le dimissioni non sono un messaggio per l'azienda, ma per il capo (``non voglio più
lavorare con te''). Allo stesso modo, se non c'è più il leader della visione, bisogna prima capire se
si condivide ancora qualcosa.

\paragraph{Lavorare così è più faticoso?} Sì e no. ``Non voglio demotivarvi, perché non è il mio
lavoro'', ma non è tutto rose e fiori: lavorare in questo modo è anche più stancante
\textbf{emotivamente}, perché non si può più pensare ``faccio il mio, alle cinque ho finito e chi si è
visto si è visto''. Si sente un senso di responsabilità verso gli altri. Succede che si finisca il
proprio lavoro del pomeriggio, ci si giri e si veda un collega che finirà alle nove di sera, e
l'istinto sarà ``posso farti un pezzo del lavoro?'': così si finisce entrambi alle sei. In un certo
senso ci si accolla più lavoro, con la consapevolezza che quando si sarà in difficoltà ci sarà qualcuno
che si girerà a dire ``ti vedo stressato, posso darti una mano?''. È bellissimo, ma non dà solo benefici
diretti: richiede anche un investimento personale che non sempre le persone vogliono fare. Pro e
contro: è il modo che sceglierebbe il relatore, ma non è per tutti. Conosce persone che vogliono poter
staccare appena finito, e va benissimo: si cerca un'azienda in cui si possa lavorare così.

\subsection{Condividere fallimenti e successi in pratica}

\paragraph{Come si celebrano i fallimenti e si condividono i successi?} Condividere i successi sembra
più semplice (``bravo Mario che ha fatto questa cosa''). Il relatore dà uno spunto ancora migliore:
\textbf{``bravo Mario, e ora Mario ci racconta come l'ha fatto e com'è andata''}. Sono sottigliezze,
ma contano: raccontare in prima persona ciò che si è fatto dà un moto d'orgoglio che non dà sentirlo
raccontare da un altro.

Per i fallimenti, che sembrano dover partire dalla persona stessa, l'obiettivo era che tutti li
condividessero spontaneamente; ma all'inizio era difficile, perché non viene naturale alzare la mano e
dire ``ho sbagliato''. Per questo il relatore ha creato artificialmente dei \textbf{``festival del
fallimento''}, dal tono chiaramente ironico: durante una pausa pranzo, mentre tutti mangiavano insieme,
a turno si saliva ``sul palco'' a raccontare la stupidaggine fatta in settimana. Molto spesso iniziava
lui, raccontando i propri errori (``ho sbagliato questa cosa, sappiate che non si fa così''). Dopo aver
forzato l'interazione per una, due, tre settimane, alla quarta lo facevano da soli: era diventato
organico dire a tutti ``ragazzi, attenzione, ho fatto questa sciocchezza, non fatela anche voi'', e c'era
perfino un \textbf{tag apposito su Slack}. È stato il team a decidere come condividere i fallimenti: come
sempre, invece di imporre il modo, il relatore ha lasciato scegliere a loro, e funziona. È impegnativo
tanto quanto lo si vuole rendere tale.

\paragraph{Che tipo di fallimenti si condividono? Anche quelli personali?} Di solito si parte da quelli
\textbf{professionali}, cioè da ciò che si riconosce come utile anche agli altri: qualcosa che ha
rallentato, fatto perdere tempo o che rischia di ripetersi (``sappiate che, per evitare che lo facciate
anche voi\dots''). Il resto è più da conversazione privata, \emph{one to one}, con la persona con cui ci
si vuole sfogare, o con gli amici: se non ha un potenziale impatto sugli altri, non c'è un incentivo a
dirlo. E un errore che si è fatto per il proprio tipo di personalità? Qui tornano i test di personalità
della prima giornata: conoscendo il proprio profilo e quello dei colleghi, si sa che qualcuno ha lo
stesso carattere, e quindi lo si aiuta.

\subsection{Cultura aziendale e formazione}

\paragraph{Come si arriva a poter fare queste cose in azienda?} Conta molto la cultura aziendale. Il
relatore sa di lavorare bene perché oggi ha il \emph{background} e i successi alle spalle per dimostrare
che il suo metodo funziona, e quindi glielo lasciano fare; ma se dieci anni prima avesse proposto le
stesse cose, ``mi avrebbero mandato a quel paese''. C'è un momento e un posto per fare le cose, e non è
sempre possibile fare ciò che si ha in testa; ma ci si può lavorare. Per lui la domanda è diventata:
``sono convinto che così sarebbe meglio; come faccio ad avere carta bianca?''. \textbf{Portando
risultati}: si è concentrato su quello, e ci è arrivato progressivamente.

\paragraph{Come funziona la formazione con workshop come questo, in azienda?} Il relatore li fa in modo
un po' più rilassato ed esteso, con molte più interazioni ed esempi, perché in azienda si possono usare
esempi che si hanno davanti ogni giorno, e parlare di problemi e discussioni diventa molto più reale. Ma
in generale è sempre \textbf{formazione pratica}, ripetuta finché il messaggio non è recepito. Il segno
che qualcosa è stato recepito è quando arriva una persona nuova e \textbf{non è lui a spiegarle le cose}:
è il team stesso a spiegarle che cos'è la feedback culture, e lui non deve parlare. È la conferma che il
concetto viene perpetuato dal team.

\paragraph{Quanti manager dell'industria condividono questi principi?} Il relatore ha fatto molti
interventi all'IGDA di Roma e ne farà uno a Milano la settimana successiva (``il 24, mi pare''); interviene
spesso a workshop ed eventi, e il riscontro è sempre ``che bello, vorremmo farlo anche noi in Italia'', e ci
stanno lavorando. La moglie lavora in un'azienda che non ha niente a che fare con i videogiochi e lo stanno
già applicando. Magari in Italia è tra i primi a parlarne, perché il paese è ancora un po' indietro, ma non è
assolutamente l'unico: è più diffuso di quanto si pensi, non è una cosa di nicchia.

\paragraph{Questi spunti si applicano anche all'IT e alla consulenza?} ``Praticamente tutti.'' Si applicano
facilmente quando c'è un gruppo di persone su cui applicarli; se si cambia continuamente cliente è più
difficile portare questa mentalità da loro, ma questo non impedisce a chi lavora con noi di applicarla
direttamente con noi.

\paragraph{Ci si può tenere in contatto?} Su LinkedIn: il relatore accetta sempre tutti e risponde alle
domande quando riesce (ne riceve tante, ma cerca di stare al passo almeno su base settimanale). E, ``a costo
di sembrare quello che si fa pubblicità'', chi lavora in un'azienda che fa fatica ad applicare questi
concetti o non sa che esistono può metterla in contatto con lui: è spesso così che ha trovato nuovi modi di
applicarli. Vale anche per chi tra sei mesi lavorerà altrove. I suoi profili sono pubblici e non c'è nessun
problema a contattarlo.

\subsection{Performance e persone in difficoltà}

\paragraph{Se un'azienda è a corto di personale, conviene comunque eliminare chi performa poco?} Il giorno
prima si era detto che conviene separarsi presto da chi non performa (sezione~\ref{sec:z1-squadra}); ma non
si finisce per caricare di lavoro tutti gli altri? Per il relatore \textbf{eliminare non vuol dire non
sostituire}: vuol dire cambiare la persona con qualcuno magari più junior, ma più adatto come \emph{fit}
culturale. Il vecchio discorso del ``fare lo stesso lavoro con meno gente'' non è sostenibile: lo si può fare
fino a un certo punto, in caso di necessità, ma non è una soluzione. Non è ottimizzare, è far andare in
\emph{burnout} chi resta.

\paragraph{Ha più senso tenere una persona poco performante e non in sintonia con il team, piuttosto che
licenziarla senza poterla sostituire?} No: se una persona è poco in sintonia, performa poco ed è in generale
tossica e negativa per il team, \textbf{prima ci si separa, meglio è}. Questo non vuol dire non provare a
recuperarla, ma bisogna andare dritti al punto: ``abbiamo un problema: ti vedo poco performante, non ti vedo
allineato con il team, vedo che il team non si fida di te. Che cosa vogliamo fare?''. Se la risposta è ``non lo
so'', la strada porta dritta all'uscita. Bisogna \textbf{recidere i rami morti} il prima possibile, facendo
però attenzione a essere sicuri che si tratti di un ramo morto e non di un ramo recuperabile: tagliare un ramo
recuperabile mette negli altri la paura di essere i prossimi, ``alla prima foglia gialla''. Quello che si
vuole è che, quando un ramo è morto, la pianta smetta di alimentarlo, e ne crescerà un altro. Entrambi i
messaggi devono essere chiari: \textbf{non si subisce passivamente la situazione}, e allo stesso tempo
\textbf{non si abbandona nessuno} (se qualcuno è in difficoltà, la prima soluzione non è tagliare il ramo). È
un compromesso tra le due cose. (A questo punto arriva, puntuale, il crash della connessione: ``non era ancora
successo oggi''.)

\paragraph{Che cosa fare se qualcuno nel team sembra soffrire di depressione?} È un discorso complesso, perché
il relatore \textbf{non è uno psicologo}. Può dare la sua disponibilità a concedere permessi, coprire il lavoro
della persona e in generale darle modo di seguire un percorso con uno specialista: non può risolvere la
depressione, ma può alleviarne la parte lavorativa e permettere di ricevere supporto da chi è qualificato. È lo
stesso motivo per cui, a Blizzard e a Riot, quando arrivavano ticket con minacce di suicidio (``mi avete
bannato, quindi mi ammazzo'', ticket purtroppo abbastanza comuni), il \emph{customer service} era istruito a
non intervenire direttamente: non dovevano essere loro a fermare la persona, perché non qualificati, ma
indirizzarla a qualcuno di qualificato, prendendo tempo. Una persona non qualificata che prova a farlo rischia
di peggiorare la situazione; come davanti a un incidente si chiama l'ambulanza e non si fa una tracheotomia.
Per il relatore salute mentale e salute fisica sono la stessa cosa: come non si azzarderebbe a operare qualcuno
di appendicite, non si azzarda a fargli da terapeuta.

\paragraph{Come si danno feedback adatti a ogni tipo di personalità?} Bisognerebbe leggere i risultati del
profilo della persona sullo stesso sito su cui si è fatto il test di personalità della prima giornata, che
spiega anche come interagire con quel profilo. In generale: con chi è più \textbf{analitico}, portare dati a
supporto; con chi è più \textbf{emotivo}, metterlo in una condizione di comfort; con chi è molto
\textbf{pratico}, andare dritti al punto senza girarci intorno. Poi fare pratica. Oppure il modo ``perfetto''
per il relatore: provarci come viene e poi \textbf{chiedere} ``andava bene così o preferisci in un altro modo?
Potevo fare meglio?''. Bisogna creare anche quel rapporto: fidarsi di discutere anche di queste cose.

\paragraph{Se nel team fossimo quasi tutti leader, non ci sarebbero ``troppi galli''?} No, perché per leader si
intende chi si prende cura degli altri, e magari lo fossero tutti. Non significa che tutti scelgano che cosa
fare: spesso lo sceglierà comunque il manager, o guiderà la scelta. Ma avere tutte persone con iniziativa e
capaci di scegliere quando serve è qualcosa per cui il relatore ``metterebbe la firma''. Non ha mai corso il
rischio dei ``troppi galli'': se ci sono quei caratteri, il team è stato assemblato male, perché non tutti sono
compatibili con questo modo di lavorare e magari non si hanno le persone giuste per applicarlo.

\paragraph{Se la fiducia richiede tempo ed esperienza, quanto è difficile per un junior?} Meno di quanto si
pensi: proprio perché su un junior non ci sono aspettative, un junior che chiede ``posso darti una mano?'' dà un
segno chiaro che tiene più al team che a se stesso. Si può essere leader anche non dal punto di vista
professionale, ma come chi va a dare supporto agli altri: ``il medico del team, l'\emph{healer}'', e lo si può
fare anche da junior. È così che ha iniziato il relatore: andare ad aiutare a risolvere problemi non suoi gli ha
fatto accelerare molto il percorso per ottenere rispetto e stima, e lo ha esposto a cose a cui altrimenti non
sarebbe mai arrivato. È una questione di carattere: se ci si tiene agli altri, lo si fa vedere, e il resto
arriva come risultato. ``È quasi un reddito passivo'': un investimento nel far vedere che ci si è, e il resto
arriva (``Stefano si è preso cura di me, quindi ora mi prendo cura di lui'').

\paragraph{Aiutare gli altri solo dopo aver finito il proprio lavoro?} Dipende dall'urgenza. Il relatore spesso,
se qualcosa è davvero urgente, mette in pausa ciò che sta facendo, risolve e poi torna al suo lavoro; ma è troppo
\emph{case by case}. Se due colleghi litigano su chi debba avere le ferie di Natale, lo si può risolvere anche
domani; se litigano perché si è perso il \emph{repository} e domani c'è la consegna, lo si risolve subito.

\paragraph{Come mantenere motivato un team su progetti non retribuiti, come quelli accademici?} Essendo il più
\textbf{trasparenti} possibile nei processi decisionali: se si deve decidere se fare un action o un RPG, tutti
devono sapere perché si fa l'uno invece dell'altro (da qui il metodo di brainstorming della mattina). Deve essere
sempre chiaro a tutti perché si stanno facendo le cose in quel modo, perché \textbf{la demotivazione nasce dal
disallineamento} rispetto all'obiettivo del gruppo, e l'unico modo di essere allineati è essere sempre
consapevoli e partecipi del processo che decide dove sta andando il gruppo. E soprattutto: \textbf{non chiudere
mai un meeting senza aver concordato che cosa si vuole fare} e senza essere tutti d'accordo. Se non lo si è,
bisogna continuare a parlarne, per quanto scomodo, lungo e noioso: si deve sempre uscire da un meeting
sapendo che si è tutti d'accordo, o al massimo in \emph{disagree and commit} (sezione~\ref{sec:z2-disagree}).
``Quello è il massimo accettabile.''

\subsection{Carriera, colloqui e candidature}

\paragraph{Come ha iniziato a lavorare all'estero?} Negli annunci online si trovano spesso requisiti come un titolo
di studio conseguito nel paese o l'abilitazione a lavorarci. Questo vale per i paesi fuori dall'Unione europea: per
lavorare negli Stati Uniti o in Sud America serve un visto o una \emph{green card}. In molti casi, però, è proprio
ottenere il lavoro a permettere di lavorare lì: quando il relatore doveva andare a lavorare in California (a Irvine,
a Santa Monica, dove capitava), era l'azienda a occuparsi dei permessi necessari, perché ``per venire qui devi avere
un lavoro, e il lavoro te lo diamo noi''.

Il suo inizio è stato abbastanza semplice. In Italia lavorava per Digital Bros, occupandosi dei giochi online gestiti
dall'azienda e di quello che era il canale Game Network sul digitale terrestre\footnote{Questo passaggio chiarisce il
dubbio segnalato nel capitolo~\ref{cap:z1-feedback} sul ``Digital Game Network Channel'': si tratta di Game Network,
canale televisivo dedicato ai videogiochi legato al gruppo Digital Bros.}. Vide un annuncio a Dublino per
\emph{Warhammer Online}, mandò la candidatura, fece i colloqui e gli dissero ``ti paghiamo la \emph{relocation}, vieni
qui''. Come detto il giorno prima, sono le aziende a occuparsi del trasferimento: se devono creare a Dublino un'unità
che segue dodici lingue europee, sanno di dover far arrivare persone da fuori, ed è già a budget. Non è nemmeno una
richiesta da fare: è prevista. Per il relatore, da questo punto di vista, è indifferente lavorare per un'azienda
italiana o estera: se si è adatti, si è adatti. Per questo insisteva sull'inglese: se lo si parla, è indifferente di
che cosa ci si occupi e quale sia la propria lingua madre, a meno che la posizione non sia esplicitamente ``community
manager italiano''. Un programmatore polacco o italiano non cambia niente. Lavorare all'estero non è più difficile: c'è
più concorrenza solo se l'azienda cerca attivamente persone da tutto il mondo e non si pone problemi di \emph{location}.
Cambiano l'azienda, il fatto di lavorare da remoto o dal vivo, l'eventuale trasferimento; per il resto non cambia
molto.

\paragraph{Gli anni passati nell'IT contano se si vuole passare ai videogiochi?} Dipende da che cosa si vuole fare: se
da responsabile IT si vuole diventare programmatore, bisogna saper dimostrare di saper programmare. Vale anche per i
passaggi interni: se un game designer vuole diventare art director, al relatore va bene, deve solo dimostrare di saper
fare l'art director. Anzi, incentiva le \textbf{promozioni interne}, perché si conosce già il candidato; ma gli
spostamenti non sono liberi ``a sentimento'': bisogna saper fare quella cosa. I recruiter valuteranno non tanto gli anni
quanto ciò che si sa fare; gli anni servono comunque, perché dimostrano di saper lavorare in un ambiente professionale,
che è un altro insieme di competenze. Ma se si sa fare ciò che chiedono, aver lavorato nell'azienda X o Y non dovrebbe
cambiare molto.

\paragraph{Quanto vengono considerate le candidature spontanee?} Non tantissimo. Soprattutto per le aziende non grandi,
se non stanno cercando qualcuno è perché non se lo possono permettere: una candidatura spontanea non le convincerà ad
assumere. Più l'azienda è piccola, più conviene aspettare l'annuncio, perché se non c'è vuol dire che non c'è budget, e
potrebbero rispondere ``ti vorremmo, ma non possiamo prenderti''. C'è anche il rischio che finisca nel vuoto (l'HR magari
non controlla nemmeno la mail, perché non sta assumendo) o, peggio, che rispondano ``non ci servi'' e una settimana dopo
esca la posizione, con il candidato già nel cestino dei rifiutati. Si rischia di \textbf{bruciarsi} la possibilità.

\paragraph{Come ci si accorge di essere riconosciuti come leader, e come lo si ``vende'' a un colloquio?} Uno studente
osserva che la leadership, così descritta, non la si riconosce da sé ma la riconoscono gli altri. Internamente il
relatore non ha mai dovuto ``venderla'', perché parlavano i risultati. Finché non lo si vive è difficile capirlo, ma
quando si arriva alla posizione in cui il team rispetta e vede qualcuno come punto di riferimento, non ci sono grandi
dubbi: è evidente che vengono da lui. Si può riflettere su \textbf{quanto spesso gli altri vengono da noi}, cercano di
emulare il nostro modo di fare, vanno a rivedere le nostre cose per capire dove hanno sbagliato, chiedono un parere, ci
parlano di un problema: la somma di tutto questo significa molto spesso che si fidano. Per ``venderla'' all'esterno si
può dire di essersi occupati di un team in cui, attraverso il lavoro e il rispetto degli altri, si è arrivati a essere
visti come punto di riferimento; ma il suggerimento è \textbf{portare esempi tangibili} (``c'era questo progetto, non
sapevamo come risolverlo, il team mi ha chiesto una mano e l'ho risolto così''), invece di ``incensarsi'' dicendo ``sono
il punto di riferimento del team''. Se il relatore racconta di essere stato eletto rappresentante dei licenziati senza
essere tra loro, un'azienda capisce che si fidavano di lui, punto: non ci sono altre conclusioni da trarre.

\paragraph{Quali sono gli errori più gravi dei junior ai colloqui, a livello di soft skill?} In ordine di priorità:
\begin{enumerate}
  \item \textbf{non arrivare preparati}, cioè \textbf{senza esempi}. Il relatore ha sempre un esempio pronto per ogni
        cosa di cui vuole parlare: per la fiducia, per gli \emph{achievement}, per i risultati. Servono sempre esempi
        tangibili per far capire che cosa si intende;
  \item \textbf{non avere domande alla fine}, perché dimostra di non avere chiaro che cosa si andrebbe a fare. Anche solo
        chiedere che cosa si farebbe, che cosa si può migliorare, o ``che impressioni hai avuto finora dal mio
        colloquio?'': si possono fare domande molto aperte al recruiter per capire come sta andando;
  \item \textbf{non leggere bene la \emph{job description}}: presentarsi con vaghezza, senza aver capito che cosa cerca
        l'azienda. Lo si impara anche con la pratica; ma il designer che si candida per una posizione da programmatore,
        senza nulla nel CV che faccia pensare che sappia anche programmare, il relatore non solo lo scarta: lo mette in
        una \textbf{blacklist} di persone che forse non considererà mai. Se non si riesce a interpretare un testo come una
        \emph{job description}, è un problema troppo grosso da superare in una collaborazione;
  \item \textbf{errori di battitura e CV impaginati male}, che molti sottovalutano. Il relatore ha letto letteralmente,
        in una lettera di presentazione per una posizione da \emph{copywriter}, la frase ``sono una persona molto attenta
        ai dettaggi''. Come si può considerare quella persona per quel ruolo? Non ha nemmeno incollato la lettera in Word
        per il controllo automatico: non è un errore sfuggito al correttore perché la parola esiste con un altro senso,
        ``dettaggi'' viene segnato in rosso in qualsiasi lingua. Come affidarle la comunicazione aziendale? Vale per tutti
        gli errori: ciò che si manda è \textbf{il primo biglietto da visita} con l'azienda, e non si manda un biglietto da
        visita sgrammaticato. Un conto è non andare ai colloqui in giacca e cravatta (``non l'ho mai fatto e probabilmente
        non lo farò mai''), un altro è presentarsi con errori;
  \item un altro esempio comune: \textbf{l'email sbagliata nel CV}. Se la mail torna indietro, il relatore non telefona
        per farsi dare quella corretta: se non si è in grado di scrivere giusta la propria email, non vale la pena. Più
        l'errore è banale e facilmente evitabile, più lo fa arrabbiare (``scusate se mi sono inalberato''): evidentemente
        il candidato non teneva abbastanza a quel posto.
\end{enumerate}

\paragraph{Quanto è importante la lettera di presentazione?} Sempre di più: sempre più aziende la richiedono come
obbligatoria, quindi conviene prepararne una, sapendo di doverla adattare al destinatario. Deve essere \textbf{meno lunga
di una pagina}, in \textbf{tre paragrafi}:
\begin{enumerate}
  \item \textbf{chi siete}: un minimo di \emph{backstory}, che cosa avete fatto, che cosa volete fare, da dove arrivate;
  \item \textbf{perché avete mandato il CV}: che cosa vi ha interessato di questa azienda;
  \item \textbf{perché dovrebbero chiamare voi}: che cosa dovrebbe spingere l'azienda a scegliere proprio voi.
\end{enumerate}
Uno studente dice di riuscire bene con i primi due, ma di trovare difficile il terzo, che è poi la domanda che tutte le
aziende fanno. In realtà basta \textbf{collegarsi alla \emph{job description}}: se cercano un 3D artist in grado di fare
X, Y e Z, si conferma di averlo letto e si spiega che con la propria esperienza si è imparato a fare proprio quelle cose.
Non vogliono sentire niente di particolare, ``non devi cambiargli il mondo''.

\paragraph{Che cosa guarda un recruiter in un portfolio?} Il relatore li guarda poco: fa un \emph{sanity check} generale,
per vedere che non sia evidente perfino a lui, ``a digiuno di tante cose tecniche'', che c'è un problema (se se ne accorge
perfino lui, il problema c'è). Per il resto guarda che il portfolio sia \textbf{coerente con il CV}: se arriva un CV da
senior con dodici progetti completati e anni di lavoro, e il portfolio sono quattro immagini, non è coerente. Vuole una
concordanza tra ciò che si racconta e ciò che si vede.

\paragraph{Che strumenti per creare un CV?} \textbf{Canva}, che ha un catalogo di template gratuiti e a pagamento
amplissimo e personalizzabile: anche se non piace, lo si prende e lo si adatta, ma fa già il 90\% del lavoro. Insieme a
\textbf{ChatGPT}, perché ``è inutile far finta che non esista'': ci si fa scrivere la lettera di presentazione e poi la si
corregge e la si adatta, e lo stesso per il CV. L'importante è renderli adatti e coerenti, non lasciarli solo frutto di
ChatGPT; ma per fare il grosso del lavoro è un ottimo strumento.

\paragraph{Conviene lasciare nel CV esperienze non del settore?} Uno studente ha alle spalle il liceo, vari stage e
lavori che non c'entrano niente con la modellazione 3D. Si possono mettere finché c'è spazio: il relatore consiglia
\textbf{al massimo due pagine}, e quando ha cominciato a sforare nella terza ha iniziato a togliere ciò che non era più
rilevante. Ormai nel suo CV la formazione è un paragrafo: ha così tante esperienze lavorative che non interessa a nessuno
che abbia fatto il liceo scientifico o classico, e non lo mette nemmeno più. Si filtra in base a ciò che è più rilevante
per lo spazio disponibile.

\paragraph{Che cosa fare con le aziende che dicono ``ti faremo sapere'' e poi spariscono?} Non aspettare uno o due anni
per scrivere. Il suggerimento: tra le domande finali del colloquio, l'ultima del relatore è sempre \textbf{``tra quanto
mi fate sapere?''}. Se rispondono dieci giorni, all'undicesimo scrive: ``mi avete detto dieci''. Non si sente in difetto a
mettere pressione, perché sono stati loro a dare la scadenza. E se viene ignorato del tutto, quell'azienda per lui è
``morta'': non si ricandiderà mai, perché non vede perché il rispetto debba andare in una sola direzione.

\paragraph{Quanto essere trasparenti, nel CV e nella lettera, su problemi personali come l'ADHD?} Uno studente con
ADHD\footnote{Nella trascrizione ``essendo 10D''; nel resto della risposta ``di'' e ``dei 10D'' stanno per ADHD.} dice di
tendere ad avere performance altalenanti. Il relatore lo invita a \textbf{non usare l'ADHD come scusa} per essere
altalenante: anche lui ha l'ADHD, ma non è altalenante. Esserlo è senz'altro un \emph{malus}, ma esistono i mezzi per non
avere problemi di questo tipo: lui ha trovato i suoi metodi (la to-do list, la \emph{must do list}, il delegare il lavoro
al team, perché occuparsi delle cose per gli altri lo mantiene molto motivato, sveglio, all'erta, attento a come vanno le
cose). Sta anche a ciascuno trovare un lavoro in cui compensare ciò che viene richiesto: se domani lo rimettessero a fare
il programmatore, come all'inizio, probabilmente impazzirebbe, perché non riusciva a restare concentrato sul codice oltre
un certo tempo. Lì ha capito che non poteva fare il programmatore, ma ha trovato qualcos'altro che gli permette di essere
com'è e fare ciò che fa. Bisogna capire i propri punti di forza e di debolezza e giocare su quelli: si possono avere tutti
i difetti del mondo e voler comunque fare una cosa, ma bisogna accettare che sarà più difficile.

Quanto alla trasparenza, il consiglio è \textbf{non mentire mai nel CV}, ``perché tanto vi beccano''. E per il discorso della
\emph{bandwidth}: tutte le energie spese a mentire, e a tenere viva la bugia, sono energie tolte allo spiegare bene ciò che
si sa fare. È un investimento nel posto sbagliato. Bisogna concentrarsi su ciò che si sa fare, con esempi tangibili, e dire
``non lo so fare'' quando non lo si sa fare. Il peggio non è essere scoperti, ma \textbf{non esserlo}: se si finge di sapere
il C++, si viene assunti e messi a lavorare in C++, che cosa si fa il giorno dopo? Si rischia davvero di rovinarsi la
carriera. Se non si è un \emph{fit}, si sarà il \emph{fit} per la prossima. L'incubo del relatore è essere assunto in un
posto dove pensano che sappia fare qualcosa che non sa fare; anzi, tende a esplicitarlo: ``cercate questa cosa? Io non la so
fare, va bene lo stesso?'' --- ``Sì, te la insegniamo noi'' --- ``Voglio solo che sappiate che non la so fare''. Meglio
questo che vivere con l'ansia che prima o poi si accorgano che non sa programmare.

Gli capita spesso con chi mente sull'inglese. Al primo colloquio chiede ``sai parlare inglese? A che livello?'' --- ``C1,
\emph{fluent}, quasi madrelingua'' --- ``Perfetto'', e lascia cadere l'argomento. All'ultimo colloquio: ``ti va se lo facciamo
in inglese?'' --- ``No, ma io intendevo che\dots'' --- ``No, mi avevi detto \emph{fluent}: ora facciamo il colloquio in
inglese''. Prima o poi le bugie raggiungono chi le dice.

% ---------------------------------------------------------------------
\section{Chiusura del workshop}
\label{sec:z2-chiusura}

Il relatore chiede alla classe di far sapere com'è andata a Vanessa, il punto di contatto con l'accademia per il workshop:
sarà lei a riferirgli eventuali lamentele o suggerimenti, che gli permetteranno di fare meglio i prossimi. Ringrazia per aver
``retto per dodici ore'' e spera che sia stato utile, consapevole che tutti sono stanchi quanto lui.

Gli studenti rispondono che è stato chiarificatore su molte cose che non sapevano, soprattutto su come interfacciarsi nei
team, e che sarà materiale da rivedere per fissare bene i concetti, vista la mole in due giorni; qualcuno si aspettava che
sarebbe stato molto più pesante e l'ha trovato invece scorrevole. Il relatore ribadisce che \textbf{niente sostituirà
l'esperienza diretta}: il workshop dà un'infarinatura e un punto di vista, ma i concetti si metteranno in pratica solo nel
\emph{day to day}. Bisogna avere \textbf{pazienza con se stessi}, sapendo che non verrà tutto bene dal primo giorno.
Confessa di essere stato un po' preoccupato per la durata, ma il materiale c'era, e farlo in un giorno avrebbe rischiato di
trasformarlo in una corsa: ``rischiava di essere uno \emph{sprint}, e invece è una maratona''.

Infine ricorda che si può aggiungerlo su LinkedIn, che registrazione e slide (in forma integrale) verranno fatte avere
tramite Vanessa (``dite che vi ho autorizzato io a chiederle'') e che, per qualsiasi dubbio, resta a disposizione. ``Buon
resto del weekend.''

% ---------------------------------------------------------------------
\section*{Punti chiave del capitolo}
\addcontentsline{toc}{section}{Punti chiave del capitolo}

\begin{itemize}
  \item I colloqui sono una \textbf{selezione reciproca}. L'annuncio dice molto sulla cultura aziendale; alla fine del
        colloquio vanno sempre fatte domande (il \emph{day to day}, la ``bacchetta magica'', ``tra quanto mi fate
        sapere?''). Si arriva con una \textbf{cifra minima mentale} e con \textbf{esempi tangibili}; non si mente mai.
  \item Time management non è una lista di cose da fare: si parte dalla \textbf{must do list}, che è anche il primo
        momento per dire ``oggi potrei non farcela''; si stimano i tempi e si chiede quando qualcosa non è chiaro.
  \item Ogni nuovo task va \textbf{prima aggiunto alla lista}, non eseguito: priorità non è recenza. Le
        \textbf{checklist} fanno risparmiare tempo, anche (e soprattutto) ai più esperti.
  \item \textbf{First contact resolution}: raccogliere tutto ciò che serve in una sola richiesta, e non iniziare un task
        senza le informazioni necessarie. Ogni interruzione costa dieci-quindici minuti di concentrazione.
  \item \textbf{Matrice di Eisenhower}: fai, pianifica, delega, elimina. ``Non importante'' significa che non è importante
        che lo faccia tu; il quadrante dimenticato è \emph{delete}.
  \item Le pause vanno \textbf{scelte, non subite}; i task grossi vanno spezzettati; bisogna \textbf{saper dire di no},
        anche ai meeting (sono necessario? può andarci uno solo? ci sarà un resoconto scritto? per che cos'altro potrei usare
        questo tempo?). \textbf{Il multitasking non esiste}.
  \item A fine giornata si verifica la \emph{must do list}: un \emph{must do} rimandabile non era un \emph{must do}.
  \item Ferie illimitate con un minimo, giornate libere e feedback culture sono, per il relatore, le condizioni per un team
        sano; il benessere è anche mentale.
  \item Il \textbf{brainstorming genera idee, non decide}. Va preparato (topic, obiettivo, vincoli), separa la
        \textbf{generazione} dalla \textbf{discussione} (\emph{rapid ideation} di cinque minuti, \emph{dot voting}
        anonimo) e combatte l'\emph{anchoring} e il dominio degli estroversi.
  \item La decisione si prende con la \textbf{matrice impatto--sforzo}: prima le idee a basso sforzo e alto impatto, poi in
        diagonale. Lo sforzo sono le risorse, l'impatto comprende anche la motivazione; decide chi farà il lavoro.
  \item \textbf{Disagree and commit}: si dichiara il disaccordo, poi ci si impegna al massimo. Se va male ha sbagliato il
        team, non ``ve l'avevo detto''.
  \item \textbf{La leadership non è un titolo}: il rispetto si guadagna. Le carriere vanno divise tra \textbf{manager} e
        \textbf{subject matter expert}, pagati allo stesso modo, così che diventi manager solo chi vuole prendersi cura degli
        altri.
  \item Dopamina, serotonina e ossitocina legano il gruppo; il \textbf{cortisolo} è contagioso e lo disgrega. Il
        ``contratto sociale'' del villaggio vale ancora: chi protegge il gruppo riceve di più, chi ne abusa tradisce.
  \item La \textbf{piramide della fiducia}: integrità (aiutare davvero, onestà, ammettere gli errori), affidabilità
        (mantenere la parola, saper rifiutare), leadership (dare l'esempio, offrirsi anche quando si rischia, non lasciare
        indietro nessuno).
  \item Il gruppo deve prendersi cura del gruppo: \textbf{raccogliere il piattino}, smettere di dire ``ci penserà qualcun
        altro'', segnalare e chiedere prima di intervenire.
  \item Da manager: delegare i meriti e prendersi le colpe, lasciare la comfort zone, chiedere ``tu come faresti?'',
        celebrare successi e \textbf{fallimenti}, dare il merito a chi lo ha. Essere leader è un \textbf{impegno a lungo
        termine}.
  \item L'obiettivo: tornare a casa \textbf{ispirati, sicuri e realizzati}, come diritto di tutti e non privilegio di pochi.
\end{itemize}

%    \part{Il corso di Game Design}
%    \renewcommand{\thechapter}{\arabic{chapter}}
%    \setcounter{chapter}{0}
%    % =====================================================================
%  Capitolo 1 — Introduzione al corso e al mestiere del game designer
%  Fonte: trascrizione della lezione GD1 (Game Design, primo anno)
%  Pacchetti richiesti nel preambolo del libro:
%    \usepackage[italian]{babel}, \usepackage[T1]{fontenc},
%    \usepackage[utf8]{inputenc} (solo pdfLaTeX), \usepackage{booktabs}
% =====================================================================

\chapter[Il mestiere del game designer]{Il mestiere del game designer: introduzione al corso}
\label{cap:introduzione-game-design}

Questo primo capitolo fa da soglia. Prima di affrontare metodi, framework e
documenti, è necessario capire \emph{in che modo} si impara il game design,
\emph{con chi} lo si fa, \emph{con quali strumenti} e, soprattutto,
\emph{che cosa fa concretamente} un game designer nel suo lavoro quotidiano.
Il capitolo segue l'impostazione della lezione inaugurale del corso: si parte
dal contesto formativo (il corso come simulazione di lavoro), si passa al
metodo di apprendimento (l'\emph{inquiry-based learning}), si elencano
letture e risorse di riferimento, e si arriva infine a una prima definizione
operativa del ruolo del designer, dei limiti entro cui lavora e delle sue
possibili specializzazioni.

Una premessa terminologica: il game design è una disciplina giovane e quasi
tutta la sua letteratura è scritta in inglese. Per questo motivo molti termini
tecnici verranno lasciati in lingua originale (\emph{pipeline},
\emph{pitching}, \emph{tuning}, \emph{feature team}, \dots): sono i termini che
si useranno nel lavoro, ed è importante imparare a riconoscerli e a usarli
correttamente fin da subito.

% ---------------------------------------------------------------------
\section{Il contesto: un corso come simulazione di lavoro}
\label{sec:simulazione-lavoro}

\subsection{Il principio fondamentale}

Il primo concetto da fissare è che il corso \textbf{si presenta come una
simulazione del lavoro}. Dal momento in cui si entra in aula si viene trattati
come \emph{potenziali professionisti dell'industria}: non soltanto studenti,
ma lavoratori effettivi di uno studio (l'accademia), con una gerarchia, dei
referenti, dei compiti e delle scadenze.

Ne derivano alcune conseguenze pratiche:
\begin{itemize}
  \item il docente non assegna un ``compitino'' da risolvere per accumulare
        punti o presenze: ogni \emph{task} è pensato per far dedicare tempo
        all'apprendimento di qualcosa di concreto;
  \item non consegnare un task non danneggia il docente (che ``vive ancora
        domani e viene pagato ugualmente''), ma lo studente, che perde
        l'occasione di mettere le mani su un argomento e diventare più abile;
  \item docenti e studenti devono trattarsi \textbf{da colleghi}, perché è
        molto probabile che un domani lo diventino davvero.
\end{itemize}

\subsection{Perché l'impegno conta più che mai}

Il lavoro del game designer è descritto come \emph{molto romantico, ma anche
molto duro}, e sempre più duro col passare degli anni. Durante il periodo del
Covid l'industria ha visto grandi investimenti; quei ``rubinetti'' sono stati
chiusi da tempo, il rapporto con i publisher è diventato più complesso e le
opportunità lavorative non sono più quelle di tre o quattro anni fa. Questo
significa:
\begin{itemize}
  \item più competizione per ogni posizione;
  \item aspettative più alte: ci si aspetta che i candidati siano ancora più
        preparati e sicuri delle proprie competenze;
  \item un unico modo per diventarlo: \textbf{lavorare}. Più tempo si dedica
        oggi ad apprendere il mestiere, più facilmente si sarà assumibili
        domani e più agevole sarà il lavoro di tutti i giorni.
\end{itemize}

\subsection{Sbagliare qui per non sbagliare fuori}

Il mantra del corso è: \textbf{quanto più provate qui, anche in modo goffo e
sbagliando, tanto meno sbaglierete là fuori}. In accademia, per quanto la
simulazione sia realistica, nessuno toglie lo stipendio e nessuno licenzia da
un momento all'altro. Nell'industria sì, e un errore può costare molto denaro
all'azienda, ancora di più se l'azienda è la propria.

\subsection{Il valore del fallimento}

Il fallimento, se affrontato con la mentalità corretta, è probabilmente
\emph{la cosa più vantaggiosa} che si possa fare durante la formazione:
\begin{itemize}
  \item un successo (``avete superato questa deadline'') conferma le proprie
        abilità, ma insegna relativamente poco;
  \item dopo $n$ fallimenti si impara moltissimo, e i successi arrivano;
  \item è fondamentale non demoralizzarsi: il feedback negativo non è mai un
        giudizio sulla persona, ma su ciò che è stato consegnato.
\end{itemize}

Il docente riconosce apertamente che dedicare tempo ed energia a un progetto e
vederlo rifiutato, o vederlo non uscire mai, è demotivante: è successo a lui e
lo ha visto accadere anche su progetti molto grandi, con team numerosi. Il
punto però non è ciò che si è ottenuto materialmente, ma \textbf{quanto si è
cresciuti}. Da ogni esperienza ci si portano a casa due cose:
\begin{enumerate}
  \item \textbf{l'esperienza}: quanto si è diventati bravi, veloci,
        praticamente abili nell'applicare un framework, quanto si è cresciuti
        come designer;
  \item \textbf{il rapporto umano}: nel lavoro si ``suda sangue'' fianco a
        fianco con altre persone, e questo crea legami molto forti che
        portano avanti l'industria nel senso più positivo del termine: nascono
        progetti privati, ci si porta da un'azienda all'altra, si formano team
        con basi solide.
\end{enumerate}

Va evitato anche l'eccesso opposto: la cultura del \emph{grinding} si è
appropriata di questi discorsi in modo distorto. Non bisogna diventare robot.

% ---------------------------------------------------------------------
\section{Organizzazione del corso e gerarchia}
\label{sec:gerarchia-corso}

\subsection{La struttura gerarchica}

La struttura dei referenti del corso ricalca volutamente una
\textbf{gerarchia tipica dell'industria}:
\begin{itemize}
  \item al vertice c'è il \textbf{direttore del corso} (Diego Ricciuti), che
        corrisponde al \emph{creative director}: supervisiona la visione del
        progetto (il concetto di ``visione'' verrà approfondito più avanti);
  \item sotto di lui ci sono i \textbf{referenti di dipartimento}, uno per
        disciplina, con un corrispettivo per ciascuna sede (Milano e Roma);
  \item in fondo ci sono gli studenti, che occupano il primo gradino di
        un'azienda: l'\emph{intern} (stagista) o il \emph{junior}, il passo
        che precede il lavoratore ``regolare''.
\end{itemize}

\begin{table}[htbp]
  \centering
  \begin{tabular}{@{}p{0.34\textwidth}p{0.29\textwidth}p{0.29\textwidth}@{}}
    \toprule
    \textbf{Disciplina} & \textbf{Milano} & \textbf{Roma} \\
    \midrule
    Direzione (\emph{creative director}) & \multicolumn{2}{p{0.58\textwidth}}{Diego Ricciuti} \\
    Game design      & Mattia Podini        & Giacomo ``Jack'' Danna \\
    Level design     & Stefano Zocchi       & Andrea Dispoli \\
    Narrative design & \multicolumn{2}{p{0.58\textwidth}}{Daniele Fusetto (lezioni condivise), Giacomo Masi} \\
    Programmazione (Unreal) & Alessandro Pregnolato & Matteo Velardi \\
    \bottomrule
  \end{tabular}
  \caption{Referenti del corso per disciplina e sede (nomi come riportati in lezione).}
  \label{tab:referenti}
\end{table}

Per il primo anno, il referente diretto degli studenti di game design è il
docente di game design: è a lui che ci si rivolge per dubbi e domande. Altre
figure utili sono il direttore didattico (severo, ``giustamente severo'', da
contattare solo per questioni particolarmente importanti), l'amministrazione
(contratti e documentazione) e la segreteria.

\subsection{I canali di comunicazione}

\begin{description}
  \item[Discord] è il canale per le comunicazioni ufficiali del corso e per
        le domande che non richiedono urgenza immediata; può avere qualche
        ritardo.
  \item[E-mail] è il canale per tutto ciò che è formale o ``ancora più
        ufficiale'' e per le comunicazioni rapide e immediate.
  \item[Telegram] esistono residui di vecchi canali precedenti a Discord, da
        evitare.
\end{description}

Le domande che valgono per tutti vanno fatte in classe (anche da un portavoce
a nome del gruppo); quelle singole e personali via Discord o, se formali, via
e-mail. Se il docente non risponde a una domanda c'è un motivo, ad esempio
ragioni legali, e lo dirà esplicitamente.

\subsection{I tre moduli e la settimana tipo}

Le tre discipline principali del primo anno (e non solo) sono:
\begin{itemize}
  \item \textbf{Game Design} (tendenzialmente il lunedì);
  \item \textbf{Level Design} (tendenzialmente il mercoledì);
  \item \textbf{Programmazione} (tendenzialmente il venerdì).
\end{itemize}
Nel primo mese le lezioni sono tutte di game design (lunedì, mercoledì e
venerdì). Sono previste anche lezioni di narrativa, ma il piano formativo si
aggiorna ogni anno: le ore di narrativa sono state ridotte a favore di quelle
di programmazione. Le poche lezioni di narrativa daranno quindi dei
\emph{cenni}, non la preparazione di un narrative designer. Ciò che serve
davvero è imparare \textbf{come comunicano gli altri dipartimenti} e come
lavorare con persone che fanno un lavoro diverso ma complementare al proprio.

\subsection{Le due anime delle lezioni di game design}

\begin{description}
  \item[Design methods (metodologie)] lezioni prevalentemente teoriche, ma con
        una buona componente pratica, in cui si studiano metodologie di design
        e i cosiddetti \emph{framework} dell'industria.
  \item[Documentazione e pitching] il cuore del lavoro quotidiano del
        designer.
\end{description}

\paragraph{La documentazione.}
Il docente si definisce ``il più grande nerd della documentazione'' che si
possa incontrare. È una parte dura, perché a pochi piace documentare, ma il
lavoro del designer \emph{è principalmente questo}: scrivere documenti e, da
quei documenti, passare alla pratica e all'implementazione. Chi non sa
documentare bene parte male. Quando non si è presenti, il documento è il
proprio \textbf{biglietto da visita}: qualcun altro lo leggerà, quindi va
scritto nel modo più chiaro e comprensibile possibile.

Un esempio concreto del perché la documentazione conta: CD Projekt Red ha
sviluppato un proprio engine proprietario, ma buona parte della sua
documentazione è stata ``buttata all'acqua''. Chi deve usarlo o lo sa già
usare, o ha competenze di un altro engine da trasferire, oppure è in
difficoltà. Si documenta \textbf{perché altri possano leggere}.

\paragraph{Il pitching.}
Il \emph{pitching} è l'esposizione, la presentazione orale. La parte creativa
del team presenta spesso:
\begin{itemize}
  \item le proprie idee ai colleghi;
  \item il gioco a convention ed eventi;
  \item il pitch del gioco agli investitori, per ottenere i fondi necessari a
        realizzarlo.
\end{itemize}
Saper parlare e comunicare chiaramente le idee è fondamentale. In sintesi:
\emph{il game designer parla}, e bisogna farsene una ragione.

\subsection{Altre componenti del corso}

\begin{description}
  \item[Gioco critico e playtest] la classe viene divisa in due gruppi: uno
        gioca, l'altro osserva. Chi osserva non applaude ``come le
        scimmiette'': annota, segna e analizza il \emph{playtest}, esattamente
        come si fa nell'industria. Il playtest fa parte delle mansioni del
        designer, ma esiste anche come specializzazione autonoma.
  \item[Pipeline] termine che indica tutto ciò che riguarda i processi
        formalizzati dell'industria; di solito si intende la \emph{pipeline di
        produzione}: come si divide il team, qual è la sua gerarchia, come il
        \emph{producer} gestisce le deadline. In breve: \emph{come} e
        \emph{quando} si fanno le cose. Al primo anno è meno rilevante (si
        dovrà ``nuotare'' un po' da soli); dal secondo anno si vedranno le
        metodologie effettivamente applicate ai progetti digitali. Si
        scoprirà quanto sia difficile lavorare a un progetto digitale, e
        ancor più fisico, senza un piano di sviluppo comprovato, come quelli
        usati per qualsiasi software in qualsiasi parte del mondo.
  \item[Engine] il motore di riferimento è \textbf{Unreal Engine}
        (sezione~\ref{sec:engine}).
  \item[Programmazione] principalmente tramite \emph{Blueprint} in Unreal,
        più eventuali sistemi di versionamento; il corso parte da zero.
  \item[Progetti pratici] probabilmente la cosa più importante. Al primo
        anno i progetti di gruppo restano circoscritti al corso di design,
        perché mancano ancora alcune nozioni (come comunicare, affrontare
        problemi, dare e ricevere feedback). In seguito si allargheranno a
        progetti digitali, molto più difficili di quelli fisici, con membri
        di altri team: un rapporto non unilaterale, ma ``covalente''.
\end{description}

% ---------------------------------------------------------------------
\section{Il metodo: l'inquiry-based learning}
\label{sec:inquiry}

\subsection{Che cos'è}

Il corso è costruito secondo l'\textbf{inquiry-based learning}
(apprendimento basato sull'indagine), un approccio pedagogico nato intorno
agli anni Sessanta, che si richiama a pensatori come Dewey e Piaget. In
generale, le discipline da cui il design attinge sono quasi sempre la
sociologia, la psicologia e, talvolta, l'economia.

Il funzionamento è un ciclo:
\begin{enumerate}
  \item il docente fornisce metodologie e framework dell'industria durante
        la lezione;
  \item lo studente, fuori dalle ore di lezione, lavora sui \emph{task}:
        mette a terra le nozioni, le maneggia, le assorbe meglio;
  \item da questo lavoro nascono nuovi dubbi;
  \item i dubbi tornano in aula: lo studente chiede, il docente risponde
        (ammesso che dare la risposta sia utile, e ammesso che una risposta
        esista).
\end{enumerate}

\subsection{Non sempre esiste una risposta giusta}

Non tutto ciò che si studia ha ``la risposta migliore possibile in assoluto
nel vuoto''. Esistono metodologie, limiti e framework che aiutano il lavoro,
ma non sono fissi: cambiano da azienda ad azienda, da lavoro a lavoro, da
progetto a progetto, e a volte vengono stravolti. Sta alla
\textbf{sensibilità del designer} capire quale strumento o framework usare in
un dato momento, e questa sensibilità si impara solo con l'esperienza. Nel
lavoro nessuno dirà ``usa questa cosa'' o ``è meglio quest'altra'': sarà il
designer a scegliere.

Anche l'esperienza del docente ha confini: sei-sette anni di lavoro, un modo
personale di approcciare il design e un certo insieme di progetti (ad esempio
nessun MMO). Oltre quei confini si cercano risposte nella ricerca, studiando
come aziende con esperienza specifica (ad esempio Riot) hanno affrontato
problemi simili e hanno condiviso le loro soluzioni, sempre valutando quali
metodi siano adatti alla dimensione del team, al tipo di progetto e al tempo
disponibile.

\subsection{Sensibilità, non memoria}

Non interessa imparare a memoria un framework o la teoria di un autore. Certe
terminologie vanno conosciute, ma per ragioni di comunicazione e chiarezza.
Ciò che conta è sviluppare la \textbf{capacità di farsi domande}, di farle
agli altri, di fare ricerca e di trovare una propria risposta. È esattamente
ciò che si fa nel lavoro: davanti a qualcosa che non si sa fare, l'abilità
del designer sta nel cercare come qualcun altro ha risolto un problema simile.

\subsection{Le lezioni come corsi di aggiornamento}

Le lezioni frontali simulano i \emph{corsi di aggiornamento} aziendali: le
sessioni che un lead di design tiene al suo team di design (o il lead di
level design al suo team, e così via). Alla fine dell'aggiornamento c'è quasi
sempre un lavoro da fare, un obiettivo da portare a terra:
\textbf{l'obiettivo della settimana}. Si lavora di settimana in settimana, e
il perché verrà chiarito in seguito.

Struttura tipica di una lezione:
\begin{enumerate}
  \item breve sessione di domande e risposte sulla lezione precedente;
  \item analisi di eventuali esercizi particolarmente critici;
  \item lezione vera e propria;
  \item discussione di gruppo, se c'è tempo;
  \item assegnazione del task della settimana.
\end{enumerate}
Il docente tende a ``parlare e straparlare'', ad aprire parentesi e a saltare
da un argomento all'altro: l'unico modo per fermarlo è che siano gli studenti
a interromperlo con domande.

\subsection{La gerarchia della competenza}

L'inquiry-based learning può essere rappresentato come una scala di
competenza:
\begin{enumerate}
  \item la sola lezione frontale porta a \textbf{conoscenze incomplete},
        perché manca la parte pratica e di ricerca;
  \item il lavoro a casa (task, ricerca) porta a una competenza \textbf{più
        completa}, non sbagliata in sé, ma ancora parziale e complementare ad
        altre;
  \item la \textbf{condivisione} dei task e dei feedback con gli altri porta
        idealmente alla competenza più completa possibile;
  \item il picco della competenza è la sua \textbf{applicazione nei
        progetti} (fisici o digitali, di gruppo): ``con le cose che avete
        imparato e ricercato, riuscite a fare un gioco?''. Quando la risposta
        diventa positiva, o comunque si inizia a navigare nei processi di
        design per realizzarne uno, quella è la vittoria.
\end{enumerate}
Non si parte subito creando giochi: si parte dall'analisi, dal design e dalle
regole generali di design, prima ``nel vuoto'', poi contestualizzando i
framework, infine unendo i puntini nei progetti.

Rimarranno comunque delle lacune, per tutti, docente compreso, perché il game
design \textbf{non è una scienza esatta}, anzi, a rigore non è affatto una
scienza. Prende in prestito metodologie dalle discipline scientifiche, ma
contiene una componente imprevedibile e incalcolabile: \textbf{la
produzione}. Si possono applicare perfettamente le regole di un framework, e
quel framework può comunque non essere adatto al gioco, al proprio modo di
sviluppare o al tempo di produzione disponibile (alcuni framework richiedono
una ricerca e un'attenzione che il tempo non concede).

C'è poi la componente creativa, la parte ``romantica'' del lavoro: esiste, ma
è \textbf{metà del lavoro}. L'altra metà è ricerca, metodo scientifico al
100\%. Dagli stessi dati $n$ designer diversi, con sensibilità diverse,
possono ottenere $n$ risposte diverse. Non tutte saranno valide: vanno
applicate, messe a terra, e questa scrematura fa parte del lavoro.

\subsection{I livelli dell'apprendimento}

Un secondo modello, a gradini, descrive i livelli di apprendimento (è la
cosiddetta tassonomia SOLO):
\begin{description}
  \item[Pre-strutturale] il docente dà il compito e lo studente lo risolve,
        possibilmente il prima possibile. È esattamente ciò che \emph{non} si
        vuole.
  \item[Uni-strutturale] si affronta una sola disciplina: è il primo passo
        necessario, ed è il motivo per cui nel primo mese ci sono solo lezioni
        di game design.
  \item[Multi-strutturale] entrano in gioco altri dipartimenti: level design,
        game design, programmazione e così via.
  \item[Relazionale] si capisce che alcune cose valgono per più
        dipartimenti: alcuni framework di narrativa derivano dal design o
        possono essere usati per fare design; le regole di documentazione si
        applicano anche ad altri dipartimenti; le regole di level design si
        mescolano con quelle di game design. Dopotutto, per fare un gioco
        servono molti dipartimenti.
  \item[Astratto esteso] il punto finale: muoversi autonomamente e capire
        cosa usare e quando, sulla base della propria esperienza. L'unico modo
        per arrivarci è fare pratica e fare progetti.
\end{description}

% ---------------------------------------------------------------------
\section{Task, feedback e regole del corso}
\label{sec:task-feedback}

\subsection{I task come opportunità}

I task non sono strettamente obbligatori (esiste un numero minimo di
consegne), ma sono \textbf{opportunità}. Quanto e come ci si lavora dipende
dallo studente. Un task fatto in cinque minuti, o in cui si è stati ``furbi'',
non verrà letto né riceverà feedback.

Domande da porsi davanti a ogni task:
\begin{enumerate}
  \item \textbf{Perché} mi è stato assegnato? Qual è l'obiettivo?
  \item Come ci lavoro per ottenere un risultato \textbf{che abbia senso}? Se
        la ricerca e i dati indicano una direzione, è inutile andare in
        un'altra.
\end{enumerate}

\subsection{Quanto lavorare: un calcolo indicativo}

Se il corso è una simulazione di lavoro, vale il carico di un lavoro normale:
\begin{itemize}
  \item un lavoro a tempo pieno comprende da 30 a 45 ore settimanali; si
        considera una media di \textbf{40 ore};
  \item si sottraggono circa 9 ore di lezione: restano $\approx 31$, diciamo
        \textbf{30 ore};
  \item con circa 3 task a settimana: $30 / 3 = $ \textbf{circa 10 ore per
        task}.
\end{itemize}
Il numero può variare, ma rende l'idea: dedicare un'ora a un task e poi fare
altro non fa guadagnare nulla a nessuno. Lo studente ``butta soldi'', il
docente risparmia tempo leggendo il task in un minuto (o non leggendolo).
Dedicando invece attenzione, ricerca, cura nella formattazione e
sperimentazione, e facendolo con cognizione di causa, si impara.

\subsection{Niente voti, solo feedback}

Ad eccezione del corso di programmazione, non ci sono voti (né ``ottimo'', né
10, né 100). Ci sono \textbf{feedback}:
\begin{itemize}
  \item in game design, generalmente sulla documentazione;
  \item in level design, sui progetti pratici.
\end{itemize}
Il docente indica cosa va bene e cosa va modificato, concentrandosi in
particolare su ciò che va cambiato. Il feedback \textbf{non è mai un giudizio
sulla persona}, ma su ciò che è stato consegnato nel concreto.

\subsection{La condivisione dei task}

Tutti i task consegnati vengono \textbf{condivisi} con la classe, e possono
essere letti solo dopo aver ricevuto i feedback del docente. Si cresce così
non solo con i propri feedback, ma anche con quelli dati agli altri. Regola
importante: ciò che il docente dice a uno studente \textbf{vale per tutti}. Se
a una persona viene detto ``questa cosa non va mai più fatta'', da quel
momento vale per l'intera classe.

\subsection{Non copiare}

Copiare i task degli altri è una delle cose peggiori che si possano fare, non
per una questione morale, ma perché annulla il senso del percorso. Si può
dedicare a ogni task il tempo e l'attenzione che si ritengono opportuni, ed è
normale che col tempo ci si specializzi (preferendo, ad esempio, il level
design o la programmazione tecnica) e si investa meno dove si sente meno
necessità. In questi casi l'alternativa corretta è:
\begin{itemize}
  \item non consegnare;
  \item studiare ciò che hanno fatto gli altri;
  \item eventualmente riconsegnare in seguito.
\end{itemize}
Ci saranno momenti in cui si verrà sommersi dalle cose da fare: bisogna
imparare a scegliere.

\subsection{Il ``non leggo''}

Il docente ha piena facoltà di non dare alcun feedback a un task (il
cosiddetto \textbf{``non leggo''}) se:
\begin{itemize}
  \item non si è lavorato con attenzione;
  \item non si rispettano le indicazioni date;
  \item è stata violata una regola estremamente importante.
\end{itemize}
Un task ``non letto'' equivale a un task non consegnato e non entra nel
conteggio. Attenzione: un task \emph{perfetto} e un task che \emph{fallisce}
nel comprendere o applicare un framework sono \textbf{ugualmente validi},
purché ci si sia impegnati e si sia dedicato tempo, pazienza e ``sudore''.

\paragraph{Riconsegna.} Un task non letto, in ritardo o che si vuole
migliorare può essere riconsegnato e verrà conteggiato, a patto che:
\begin{itemize}
  \item i problemi segnalati siano stati risolti;
  \item sia stato fatto ulteriore lavoro, integrando anche tutto ciò che nel
        frattempo è stato detto agli altri colleghi nei loro feedback.
\end{itemize}
Il task va sviluppato con le competenze che si hanno nel momento della
riconsegna, non con quelle del momento della prima consegna.

\subsection{Requisiti per superare l'anno}

\begin{itemize}
  \item consegnare almeno l'\textbf{80\% dei task} assegnati (soglia
        malleabile);
  \item partecipare almeno all'\textbf{80\% delle lezioni}, in presenza o
        da remoto, preferibilmente in presenza; i ritardi vengono valutati
        caso per caso;
  \item comportarsi da persone corrette e rispettose (è stata firmata una
        carta dello studente): non urlare, non alzare le mani, non
        vandalizzare, non compiere atti illegali.
\end{itemize}
Impegnarsi di più porta inoltre a essere selezionati per varie opportunità.

\subsection{Scadenze}

\begin{itemize}
  \item la scadenza di un task è \textbf{sempre prima dell'inizio della
        lezione} successiva indicata;
  \item se la classe, come team, mantiene comportamenti scorretti, il docente
        può \textbf{anticipare le deadline} per tutti.
\end{itemize}

\subsection{Accessibilità e DSA}

Gli studenti con disturbi specifici dell'apprendimento o neurodivergenze
(dislessia, ADHD, autismo, ecc.) sono invitati a segnalarlo. Quando si
affronteranno le regole di formattazione e documentazione (errori
grammaticali, mancanze, incoerenze) se ne terrà conto e si potranno trovare
dei compromessi, pur ricordando che nel mondo del lavoro nessuno farà
``sconti''. La documentazione può essere inviata all'amministrazione o alla
segreteria, che la comunicherà a tutti i docenti, ma il consiglio è parlarne
anche direttamente con ogni singolo docente.

% ---------------------------------------------------------------------
\section{Le competenze trasversali}
\label{sec:soft-skill}

\subsection{I giochi si fanno in gruppo}

Si possono fare giochi da soli, ed è un'impresa ammirevole; ma lavorando da
soli \textbf{tutto ciò che si guadagna in libertà si perde in produttività}.

\paragraph{Il caso Stardew Valley.} \emph{Stardew Valley} è stato realizzato
da una sola persona, Eric Barone, in anni di sviluppo e di sacrifici: faceva
un altro lavoro per mantenersi (in un cinema) e di notte programmava, faceva
l'arte e componeva la musica. Il risultato è stato straordinario e lo ha reso
ricco, ma si tratta di un caso estremamente raro. Quanti giochi simili, nati
nelle stesse condizioni, sono falliti? Molti di più.

\paragraph{Il rischio d'impresa.} Lavorare da soli non è sbagliato di per sé
e ha aspetti molto positivi: si ha il privilegio di creare esattamente ciò che
si vuole. Ma ci si fa carico anche di tutto il \textbf{rischio d'impresa}: se
si lavora a un gioco per cinque anni e poi il gioco va male, resta il
vantaggio di averlo pubblicato, ma il dopo è sperare di vivere di quella
release o di finanziare con essa il gioco successivo. Le grandi aziende hanno
gli stessi problemi, ma li distribuiscono su più persone e più responsabilità.

L'obiettivo di un professionista è \textbf{saper lavorare in diverse
condizioni}, e la più comune è quella del dipendente (o del freelance /
libero professionista con partita IVA): si viene assunti per le proprie
capacità da un'azienda, piccola o grande che sia, e si lavora con altre
persone.

\subsection{Essere persone con cui si vuole lavorare}

Se le persone non vogliono lavorare con voi, quello è \textbf{il primo motivo
di licenziamento}. E se non si viene licenziati, il clima nell'azienda
diventa terribile: la produttività crolla, la motivazione crolla. Le
\emph{soft skill}, cioè l'insieme di abilità che rendono una persona
piacevole con cui lavorare, vanno coltivate e diventano sempre più importanti
col passare degli anni (relativamente meno ora, molto di più quando i team
saranno più grandi). Anche chi è un ``lupo solitario'' deve imparare a:
\begin{itemize}
  \item lavorare con altre persone;
  \item comunicare con altre persone;
  \item trovare un punto in comune da cui partire per fare le cose.
\end{itemize}

\subsection{Indipendenza e proattività}

Sono le altre due parole chiave del corso:
\begin{description}
  \item[Indipendenza] spesso si verrà lasciati da soli, non per cattiveria ma
        per sfida: quando si cerca una risposta da soli, ciò che conta non è
        tanto la risposta quanto \textbf{il processo}. Una risposta data dal
        docente non vale niente.
  \item[Proattività] chiedere, domandare, agire da soli senza aspettare
        direttive, perché spesso quelle direttive non ci sono, a maggior
        ragione nei team piccoli.
\end{description}
Tutto questo si concretizza nel parlare, nel presentare e nel fare
\emph{post-mortem} (l'analisi a posteriori di un progetto concluso, che verrà
approfondita più avanti), e in generale nell'attenzione alla crescita
personale.

\subsection{Fare domande}

Il corso insiste sulle domande: alzare la mano, intervenire anche ``a
bomba'', fermare il docente ogni volta che qualcosa non è chiaro. Se una
persona non conosce qualcosa, molto probabilmente c'è almeno un'altra persona
che non la conosce. Il docente può dare per scontata la conoscenza di termini
che non lo sono, ed è del tutto normale, soprattutto all'inizio.

% ---------------------------------------------------------------------
\section{Strumenti: engine, Unreal e materiali}
\label{sec:engine}

\subsection{Unreal Engine e gli altri motori}

Il motore di riferimento del corso, usato in level design, è \textbf{Unreal
Engine}: ``un gran mattone'', ma che dà a un designer molte possibilità. Il
vantaggio principale è che permette a un designer non pratico di
programmazione di \textbf{prototipare tramite i Blueprint} (il sistema di
scripting visuale di Unreal). Questo non impedisce di passare a Unity o ad
altri engine, e viceversa: il docente stesso ha studiato e lavorato a lungo
con Unreal, ma la sua ultima esperienza professionale (come level designer) è
stata su Unity, un engine che non conosceva, a dimostrazione che il salto da
un engine all'altro è possibile.

Panorama dei motori citati:
\begin{itemize}
  \item \textbf{Unreal Engine} e \textbf{Unity}: i due più diffusi, con
        modelli di sorgente ``più o meno'' aperto;
  \item \textbf{Godot}: engine open source;
  \item \textbf{CryEngine}: proprietario di Crytek;
  \item \textbf{Snowdrop}: engine di Ubisoft;
  \item Rockstar usa un proprio engine proprietario;
  \item alcune aziende usano versioni \textbf{personalizzate} di engine
        esistenti: ad esempio Milestone, nei giochi di corse, usa un Unreal
        customizzato;
  \item CD Projekt Red ha un engine proprietario (con i problemi di
        documentazione già citati).
\end{itemize}
In sintesi: ogni azienda può avere un engine proprietario oppure una versione
modificata di un engine commerciale, ma si parte sempre da una base.

\subsection{Programmazione: cosa ci si aspetta da un designer}

Non è richiesto alcun bagaglio pregresso di programmazione: il corso parte da
zero. Gli studenti sono stati selezionati (il corso è a numero chiuso) per
competenze analitiche e critiche sopra la media. Anche nel mercato del lavoro,
per una posizione di game designer generalmente \textbf{non ci si aspetta
competenza di scripting} (C++, Java, C\#, ecc.), a meno di voler fare il
\emph{technical designer}, e anche lì varia da azienda ad azienda.

\subsection{Versione e preparazione}

Indicazioni pratiche emerse in aula:
\begin{itemize}
  \item scaricare Unreal Engine; la versione consigliata per il corso è
        stata indicata nella \textbf{5.6} (con conferma successiva). La
        stabilità dipende da cosa si deve fare, da quali tool e plugin si
        usano: i problemi nascono spesso dai plugin. Nel passaggio da Unreal 4
        a Unreal 5, ad esempio, il sistema \emph{Landscape} era inizialmente
        inutilizzabile, mentre in Unreal 4 funzionava molto bene;
  \item è consigliato installare tramite \textbf{Epic Games Launcher}, più
        comodo: permette anche di riscattare giochi gratuiti periodici e di
        accedere ai pacchetti di \textbf{Fab} (il marketplace di asset
        acquisito da Epic), che ogni mese offre alcuni asset gratuiti utili
        per esercitarsi nel level design (ad esempio per fare \emph{speed
        level design});
  \item molti pacchetti di UI o di sistemi di interazione sono utili per
        prototipare al volo, ma spesso inutilmente complicati;
  \item un plugin di \textbf{materiali di prototipazione a griglia} (adattivi,
        con metriche, allineabili alla griglia, con materiali già pronti e
        colori modificabili, dal costo di circa 5 euro) è indicato come
        ``salvavita'' per il level designer; se necessario, i pacchetti
        indispensabili vengono forniti;
  \item portare \textbf{sempre il portatile} alle lezioni di level design
        (da novembre, a partire dall'interfaccia dell'engine); alle lezioni
        di game design basta che lo porti almeno metà della classe, per i
        lavori a coppie. Serve comunque qualcosa per prendere appunti
        (``tavolette d'argilla e cannucce, se volete fare i babilonesi'',
        oppure carta e penna).
\end{itemize}
Chi ha un portatile poco potente dovrebbe comunque provare a installare
l'engine e aprire un template per vedere come gira: all'inizio non si faranno
cose pesanti. Nelle prime lezioni di level design è importante mettere subito
le mani sull'interfaccia: ascoltare dove si trova uno strumento (ad esempio lo
\emph{snapping}) senza provarlo porta a dimenticarlo.

% ---------------------------------------------------------------------
\section{Letture e risorse di riferimento}
\label{sec:risorse}

\subsection{I libri}

\begin{description}
  \item[\emph{The Art of Game Design: A Book of Lenses} (Jesse Schell)] il
        testo \textbf{davvero importante}. Alcune lezioni finali dell'anno si
        baseranno sulle sue terminologie. È utile perché apre la mente sulla
        possibilità di applicare metodologie diverse a giochi diversi, ma
        soprattutto per la \textbf{terminologia} comune. Esiste anche in
        audiolibro (circa 25 ore).
  \item[\emph{Challenges for Game Designers} (Brenda Romero e Ian Schreiber)]
        libro \emph{entry level} con moltissimi esercizi, molto caro al
        direttore del corso.
  \item[\emph{21st Century Game Design} (Chris Bateman)] utile come
        retrospettiva: è un libro visionario che ha previsto metodi e
        soprattutto \emph{archetipizzazioni} dei giocatori tuttora valide nel
        design e nel marketing. Servirà per una delle lezioni più importanti
        dell'anno. È utile leggerlo per vedere quanto aveva previsto, ma è
        utile ``fino a un certo punto''.
  \item[\emph{An Architectural Approach to Level Design} (Christopher
        Totten)] non strettamente necessario, ma di lettura semplice. Ha due
        anime: parte dall'architettura (babilonese, greca, ecc.) e la trasla
        su esempi di level design (come \emph{Doom} e \emph{Team Fortress}).
        Un tempo ne esisteva un riassunto su Gamasutra (oggi
        \emph{Game Developer}).
  \item[\emph{A Theory of Fun for Game Design} (Raph Koster)] utile fino a un
        certo punto: contiene diverse idee superate, che lo stesso Koster ha
        affrontato in vari talk reperibili nel GDC Vault. Verrà ripreso nelle
        prime lezioni, poi messo da parte.
\end{description}

\subsection{Canali e conferenze}

\begin{description}
  \item[GDC (Game Developers Conference)] il suo canale raccoglie talk e
        panel estremamente importanti: è ``uno strumento misterioso che ci
        servirà più tardi''. \textbf{Attenzione}: per molti talk mancano ancora
        le competenze per interpretarli, e alcuni potrebbero fuorviare o
        mettere in testa idee strane. La GDC si concentra spesso su
        produttività e mobile.
  \item[Il \emph{Dragon Speech} di Chris Crawford] il talk del fondatore della
        GDC: appartiene a un passato lontano, ma è molto motivante. Da vedere
        subito.
  \item[Game Maker's Toolkit] canale particolarmente utile per nozioni
        pratiche, con video brevi e alcune interviste.
  \item[Nordic Game] altra conferenza importante, con talk e panel di esperti
        del settore.
\end{description}

\subsection{Associazioni e premi: come leggerli}

Le associazioni di settore (ad esempio IIDEA in Italia, o la IGDA) non sono
il male, anzi: sono un canale importante per affacciarsi all'estero o per
accedere a fondi. Va però ricordato che \textbf{tutelano i propri associati}:
non esiste solo il loro palinsesto, e bisogna valutare quanto spazio occupano
nel panorama della visibilità.

Lo stesso vale per i \textbf{premi} (Game Awards, ``Best Mobile Game'',
``Games for Change'', ecc.): al designer \emph{non interessano}. Che un gioco
abbia vinto ``più di 90 premi'' non dice nulla di utile. Servono a dare
visibilità e lustro ai giochi ed è giusto che esistano, perché ampliano il
panorama e il numero di giocatori. Ma ciò che interessa al designer sono i
\textbf{dati}:
\begin{itemize}
  \item come sono state implementate certe cose;
  \item come è stato fatto il marketing e dove il gioco è stato presentato;
  \item non quante wishlist ha fatto, ma \emph{come} le ha fatte;
  \item in che periodo storico è uscito e quali erano i competitor in quel
        momento.
\end{itemize}

% ---------------------------------------------------------------------
\section{Che cosa fa un game designer}
\label{sec:cosa-fa}

Per arrivare a una prima definizione, la classe ha condotto un brainstorming:
ognuno doveva indicare una cosa che, secondo lui, fa il game designer (senza
ripetere quelle già dette). Ne è emersa la lista seguente, commentata punto
per punto.

\subsection{Progettare meccaniche e struttura del gioco}

Il ruolo più generico è quello del \textbf{core mechanic designer}, detto
anche \emph{gameplay designer} e, un tempo, \emph{black box designer}: si
occupa delle \textbf{meccaniche di gioco}. Si va dal piccolo al grande (in
funzione anche della dimensione dell'azienda):
\begin{itemize}
  \item nel \textbf{piccolo} ci sono le singole meccaniche;
  \item nel \textbf{grande} ci sono i \textbf{sistemi}. Nei giochi
        \emph{sistemici}, con molti sistemi interconnessi (ad esempio
        \emph{GTA}), esistono specializzazioni dedicate a bilanciare e
        supervisionare i sistemi: è il \emph{system design}.
\end{itemize}

\paragraph{Quando un gioco è ``sistemico''?}
Non dipende solo da come le meccaniche interagiscono con il giocatore (le
\emph{dinamiche}, secondo i framework che passano dalle meccaniche alle
dinamiche, che verranno studiati più avanti), ma anche da \emph{quante}
meccaniche ci sono e da quanto interagiscono tra loro. Non diremmo che
\emph{Super Mario} sia un gioco sistemico, eppure contiene dei sistemi: la
linea è sottile. Quando i sistemi aumentano molto, sono complessi e si
influenzano a vicenda, il gioco diventa sistemico. Esempio: in \emph{The
Legend of Zelda: Breath of the Wild} il vento modifica quanto il personaggio
sale o scende; il fuoco può incendiare l'erba; l'erba prende fuoco se colpita
da una freccia infuocata; e così via. Nel mercato moderno i giochi sistemici
sono spesso, ma non sempre, legati a mappe \emph{open world}.

\subsection{Le 3C: camera, character, control}

La base del gioco, da cui generalmente si inizia a lavorare (anche nei
prototipi), è la ``sacra trinità videoludica'' delle \textbf{3C}:
\begin{description}
  \item[Camera] che tipo di inquadratura ha il gioco: prima persona, terza
        persona, isometrica, ecc.
  \item[Character] quanto è grande il personaggio, che tipo di personaggio è
        e, di conseguenza, quali azioni può compiere.
  \item[Control] i controlli: che cosa il personaggio può fare, quali sono i
        \emph{verbi} di gioco.
\end{description}
Le 3C sono definite dal designer, ed esiste anche una specializzazione per
questo lavoro. Esistono varianti: quella del direttore del corso è
\textbf{4C + J}, cioè \emph{camera, character, control, conflict} (conflitto)
e \emph{juice} (il ``succo'', il feedback che rende piacevole l'interazione).
Verranno approfondite più avanti.

\subsection{Il personaggio}

Il designer definisce il personaggio dal punto di vista delle meccaniche
(dimensioni, azioni, capacità), all'interno delle 3C. L'aspetto estetico è
invece competenza di altri dipartimenti.

\subsection{Documentare le informazioni per gli altri dipartimenti}

Assolutamente vero: la documentazione è una delle prime attività del
designer. Più l'azienda è grande, più questo lavoro diventa prominente; più è
piccola, meno è rilevante. In \textbf{pre-produzione}, però, è sempre
fondamentale.

\subsection{Mediatore e leader ``al servizio''}

Il designer è un \textbf{mediatore} nei conflitti interni al team, una sorta
di collante che si relaziona con dipartimenti diversi (non sempre gli stessi,
a seconda della fase di produzione). A volte fa anche da leader, a seconda di
come sono gestite la pipeline e la gerarchia.

Questa leadership è difficile da definire perché \textbf{non è verticale}: è
una leadership \emph{al servizio} (\emph{servant leadership}). Il designer
``è asservito al programmatore'': il suo compito non è arrivare con un'idea e
scaricarla addosso agli altri, né portare la soluzione e lasciare il problema
a qualcun altro, ma \textbf{seguire il processo a basso livello}, sporcandosi
le mani.

\paragraph{I feature team.} Nei team medio-piccoli o medio-grandi possono
esistere sotto-team (\emph{feature team}) dedicati a un obiettivo specifico:
ad esempio un designer incaricato del character o della camera, affiancato da
uno o due programmatori e da un artista. In questi casi il designer
supervisiona il lavoro e gestisce la produzione del sotto-team.

\subsection{Gestire l'avanzamento del prodotto}

Esiste una specializzazione che si occupa solo di questo, il
\textbf{producer}, a un livello gerarchico completamente diverso (e che
potrebbe essere una strada percorribile al termine del percorso di studi).
All'interno di un feature team anche il designer può farlo, ma non è detto.
Può certamente \emph{dare la direzione}: nella sua documentazione può
indicare che per arrivare a uno stato giocabile del prototipo serve un certo
lavoro di arte e di programmazione, per lo stato successivo un altro, e così
via. In ambito produttivo esistono metodologie strutturate come
\textbf{Agile}, che verranno viste in seguito.

\subsection{Presentare le idee}

Il pitching (sezione~\ref{sec:gerarchia-corso}) è richiesto al designer sia
in pre-produzione sia in produzione, in misura che dipende dall'azienda. Ad
esempio Ubisoft lascia ai designer, e in realtà a qualsiasi membro, la
possibilità di \emph{pitchare} un'idea prima di una nuova produzione, perché
possa iniziare a prendere vita.

\subsection{Bozze e prototipi}

Le bozze vanno bene, ma devono essere \textbf{chiare}. C'è un'eterna lotta tra
documentazione e gioco giocabile: se si fa solo documentazione il gioco non
avanza; se si fa solo il gioco manca la documentazione e si perde la visione.
Sono due cose \textbf{complementari}.

Il modo per comunicare un'idea in forma grezza ma giocabile è il
\textbf{prototipo}. Un prototipo può essere qualsiasi cosa: di carta, in
Unreal con Blueprint ``brutti''. \emph{Non deve essere bello, deve essere
giocabile}.

\paragraph{Esempio: il prototipo di carta di Mario + Rabbids.} Durante lo
sviluppo di \emph{Mario + Rabbids}, i programmatori di Ubisoft erano impegnati
su altri progetti (come \emph{Ghost Recon}). I designer realizzarono allora
un \textbf{gioco da tavolo} che riproduceva le meccaniche del gioco, quelle
di uno strategico a turni sul modello di \emph{XCOM}. E funzionava.
L'obiettivo era \textbf{validare un'idea}, validare un gameplay.

Altri modi per rendere chiara un'idea sono i \emph{mock-up} nella
documentazione e le immagini modificate o i \emph{mash-up} di qualsiasi
natura. Lato documentazione l'importante è che sia chiaro; lato prototipo,
che sia giocabile.

\subsection{Fare ricerca per l'idea, e poi smontarla}

Il designer ha una componente creativa importante e fa ricerca per arrivare a
un'idea. Ma ciò che deve poi fare necessariamente è \textbf{smontarla}. Il
suo lavoro è scientifico, esattamente come il \textbf{metodo scientifico}, il
cui obiettivo è \emph{confutare un'ipotesi}: si arriva a un'ipotesi tramite
sperimentazione e poi si cerca $n$ volte di confutarla; se resiste, la si
assume come valida \emph{in un certo contesto} (le regole della fisica
classica, ad esempio, non valgono nella meccanica quantistica).

Allo stesso modo il designer arriva a un'idea (tramite ricerca o chiedendo
agli altri: le idee le hanno tutti, non sono un privilegio del designer),
cerca di smontarla e poi cerca di \textbf{metterla a terra}. Se l'idea non è
applicabile, non è implementabile, finché non si trova un modo per
realizzarla va scartata. Può essere l'idea più bella dell'universo, ma
\textbf{se non può essere messa a terra, non esiste}.

\paragraph{L'originalità.} Serve un'idea originale? Dipende dall'obiettivo e
dalla visione del progetto:
\begin{itemize}
  \item se l'obiettivo è vendere molto, un'idea originale non serve
        necessariamente;
  \item se si vuole creare un nuovo filone, serve un'idea, se non originale,
        almeno interessante: può essere un discorso di marketing o una
        meccanica innovativa, con il rischio di non vendere.
\end{itemize}
Il ruolo del \emph{director} e il concetto di \emph{visione} verranno
approfonditi in seguito.

\subsection{Analizzare i competitor}

Arrivati a un'idea (``vogliamo un certo tipo di combattimento''), si va a
vedere se qualcuno l'ha già fatta e \emph{come}. Le meccaniche non sono
privilegio di nessuno, salvo quando vengono brevettate, come il
\textbf{Nemesis System} di \emph{La Terra di Mezzo: L'Ombra di Mordor}.

\paragraph{Esempio: il combattimento da ridimensionare.} In un progetto
(presso Stormind Games) il team si trovò con un sistema di combattimento
troppo costoso: mancavano le risorse e gli animatori per realizzarlo. Poco
output lato gameplay significava non poter fare qualcosa di complesso e
profondo come in \emph{Elden Ring}; al contrario, il team aveva molti
artisti e poteva puntare su qualcosa di più artistico. Si iniziò così a
orbitare intorno al corpo a corpo di \emph{The Last of Us} (parte I e II).
Sul canale della GDC esistono diversi talk che raccontano cosa è stato fatto
e come, con successi, insuccessi e costi (ad esempio del \emph{motion
capture}): da quei talk si ricavarono moltissime informazioni utili.

I competitor possono essere:
\begin{itemize}
  \item \textbf{diretti}: un gioco molto simile al proprio;
  \item \textbf{di mercato}: un gioco che domina il mercato in cui ci si vuole
        affacciare;
  \item \textbf{indiretti}: un gioco con un sistema simile al proprio, ma
        tutto il resto diverso.
\end{itemize}
Regola d'oro: \textbf{là fuori, nascosto nell'ombra, c'è sicuramente un gioco
molto simile al vostro}. Va solo trovato. E se per qualche ragione non ha
venduto, bisogna capire perché.

\subsection{Lavorare in engine e fare tuning}

Il level designer lavora in engine, così come il technical designer; ma
qualsiasi game designer, a seconda di ciò che deve fare, può finire in
engine, soprattutto per il \textbf{bilanciamento} e il \textbf{tuning}.

Il \emph{tuning} è l'azione di bilanciare con precisione i dati. Esempio: in
\emph{Super Mario} si modifica quanto in alto salta il personaggio, con che
velocità, in quanto tempo raggiunge l'apice, e così via. Questo lavoro rende
il gioco:
\begin{itemize}
  \item ad alto livello, dotato della \textbf{progressione} giusta;
  \item a basso livello, \textbf{piacevole da giocare}.
\end{itemize}
Il \emph{fine tuning} è un lavoro di dettaglio, paziente e minuzioso.

\subsection{Bilanciare}

Il bilanciamento, emerso tardi nel brainstorming, è una delle attività
centrali del designer e ha anche una specializzazione dedicata
(l'\emph{economy designer}, sezione~\ref{sec:tipi-designer}).

\subsection{Analizzare il mercato}

In pre-produzione, di solito, sì: capire le tendenze, i generi che cambiano
nel tempo, la compatibilità del gioco con il momento storico, in vista ad
esempio di un possibile publisher. Va capito quanto, come e con quanto tempo
farlo, perché si è sempre limitati da molte cose.

\subsection{Sintesi}

\begin{quote}
Il designer parte dall'idea ma non si ferma lì: la confuta, lavora sui
dettagli, bilancia, fa \emph{fine tuning}; lavora con il team e fa da
\emph{manager} per diversi dipartimenti; e naviga costantemente tra i limiti.
\end{quote}

% ---------------------------------------------------------------------
\section{I limiti del design}
\label{sec:limiti}

I limiti ci sono sempre, e nel mondo del design sono spesso numerosi. I più
importanti entro cui navigare sono:
\begin{description}
  \item[Tempo] ad alto livello, l'idea deve essere realizzabile nel tempo di
        produzione disponibile; a basso livello, il proprio lavoro
        settimanale deve essere fattibile nel tempo che si ha. Non ci si può
        dividere tra mille feature team se realisticamente se ne può seguire
        uno solo.
  \item[Tecnologia] si può o non si può fare qualcosa? Ci sono i server per un
        multiplayer con un certo numero di giocatori? Ci si può permettere,
        ad esempio, una fauna che reagisce alla fisica come quella di
        \emph{GTA VI}, che richiede il meglio dei \emph{technical artist}?
  \item[Background e competenze] che cosa sa fare il team? Realisticamente,
        chi ha lavorato dieci anni su giochi horror viene assunto per un
        horror, perché ha le competenze per risolvere problemi già visti.
  \item[Persone e motivazione] il team ha la motivazione e le competenze per
        fare un certo tipo di gioco, e per lavorare insieme?
  \item[Denaro] tutto ruota intorno ai soldi. Che si lavori da soli o in mille,
        le persone vanno pagate e il gioco deve avere un ritorno economico.
        Bisogna decidere dove tracciare la linea: raggiungere il
        \textbf{break-even} (guadagnare esattamente quanto speso) o puntare a
        grandi guadagni. Con un \textbf{publisher}, inoltre, una parte
        consistente dei ricavi va a lui (ripartizioni dell'ordine di 80/20 o
        70/30, che dipendono dal contratto firmato).
\end{description}

\subsection{Non siete ``gli uomini delle idee''}

Per qualcuno può essere un problema scoprirlo, ma il designer non è
``l'uomo (o la donna) delle idee'': accanto alla parte creativa c'è una parte
tecnica, e si è fortemente limitati da molti fattori. Questi limiti, però,
sono un aiuto:
\begin{itemize}
  \item riducono il ventaglio di possibilità potenzialmente distruttive per il
        team (cose che non si riescono a fare, che non piacciono o per cui
        mancano le risorse);
  \item aiutano a mirare bene, trovando soluzioni creative in uno spazio più
        piccolo e più attuabile.
\end{itemize}

\begin{quote}
\centering
\textbf{Le idee sono nulla, l'implementazione è tutto.}
\end{quote}

% ---------------------------------------------------------------------
\section{Ruoli e specializzazioni}
\label{sec:tipi-designer}

\subsection{La scala gerarchica}

\begin{description}
  \item[Creative Director (CD)] mantiene la \textbf{visione} e dà l'idea, in
        modo che tutto segua una certa direzione. In teoria può diventarlo
        chiunque, ma con un'esperienza notevole, nell'ordine dei decenni: si
        assume che in quel tempo abbia visto come lavorano molti dipartimenti.
        Naturalmente un CD con un background artistico ha meno competenze di
        design, e viceversa.
  \item[Director di disciplina] uno per dipartimento: \emph{game director},
        \emph{level director}, \emph{art director}, ecc.
  \item[Lead] il capo del dipartimento, che gestisce i gruppi di persone. Chi
        lavora in un dipartimento risponde al proprio lead (ad esempio un game
        designer risponde al \emph{lead designer}, non al game director).
  \item[Senior / Principal] i profili più esperti (a seconda dell'azienda).
  \item[Junior] il profilo di inizio carriera.
  \item[Intern] lo stagista.
\end{description}
Lo schema si replica uno a uno per ogni dipartimento: arte, suono, animazione
(a volte arte e animazione insieme, o \emph{rigging} e animazione insieme),
programmazione, UX/UI e così via. A fare la differenza, a parità di ruolo, è
il \textbf{livello di seniority}.

\subsection{UX e UI}

UX e UI stanno a metà tra game design e arte: a seconda dell'azienda il ruolo
si chiama \emph{UX/UI Artist} o \emph{UX/UI Designer}.
\begin{description}
  \item[UX (User Experience)] come una persona interagisce con le cose. Esempio:
        su uno smartphone, un pulsante accende e spegne lo schermo; per
        digitare la password si scorre verso l'alto e compaiono i punti da
        toccare.
  \item[UI (User Interface)] come sono ``vestite'' queste cose: ad esempio
        l'aspetto della barra dei punti vita (HP), realizzata con asset 2D o
        3D.
\end{description}

\subsection{Le specializzazioni del game design}

Più l'azienda è grande, più la specializzazione è puntuale.
\begin{description}
  \item[Gameplay / Core mechanics designer] (sono la stessa cosa) si occupa
        delle meccaniche e delle 3C.
  \item[System designer] si occupa dei sistemi e delle loro interazioni.
  \item[Economy designer] si occupa del bilanciamento; lavora soprattutto con
        fogli di calcolo (Excel, Google Sheets).
  \item[Combat designer] si occupa del combattimento; negli annunci di lavoro
        indica quasi sempre un combattimento \emph{action}, con gameplay
        reattivo (non strategico), anche se non è una regola.
  \item[Multiplayer designer] si occupa delle interazioni multigiocatore.
  \item[Technical designer] fa da ponte tra game design e programmazione.
\end{description}

\subsection{Le specializzazioni del level design}

\begin{description}
  \item[Encounter designer] posiziona i nemici nella mappa in modo che gli
        scontri siano bilanciati e la progressione equilibrata come
        desiderato.
  \item[World / map designer] popola il mondo di gioco.
  \item[Quest designer] popola il mondo con le missioni e pensa alla
        progressione al loro interno, cioè a dove il giocatore deve andare (si
        vedano ad esempio gli annunci di lavoro di CD Projekt Red o di
        Larian).
  \item[Technical level designer] fa da ponte tra level design e
        programmazione.
\end{description}

\subsection{Discipline vicine: suono, UI, produzione}

Il monte ore del corso non permette di uscire come \emph{narrative designer}
(o \emph{sound designer}) completi. Ciò non toglie che si possano coltivare
queste discipline:
\begin{itemize}
  \item nei progetti di gruppo si può occuparsi anche del suono, della UI/UX,
        o provare ad applicare metodologie come Agile per fare
        \emph{producing} e gestire il team;
  \item il tempo dedicato a queste attività va però sottratto al game design,
        e dipende da quante persone ci sono nel team, da quanto lavoro c'è e
        da quanto è grande il progetto;
  \item avere sviluppato e rilasciato un progetto, anche piccolo (anche una
        \emph{game jam}), documentando i propri processi, è già un ottimo
        argomento da presentare in futuro;
  \item la via migliore resta lavorare nell'industria e, da lì, chiedere di
        imparare anche altre discipline.
\end{itemize}
Un gioco si può fare su qualsiasi cosa: sugli scarafaggi, sui piedini delle
sedie, sui pennarelli. Si parte avvantaggiati da ciò che è vicino alla propria
sensibilità: chi ama la musica può pensare a un gioco con meccaniche legate
alla musica, con un gameplay innovativo ma contenuto e realizzabile. Il design
è un cappello che comprende molte discipline.

% ---------------------------------------------------------------------
\section{Approfondimenti emersi in lezione}
\label{sec:approfondimenti}

\subsection{Roguelike e roguelite}

Un esempio di come il marketing conia generi: Supergiant Games, con
\emph{Hades}, rese celebre il termine \textbf{roguelite} al posto di
\emph{roguelike} (pur non essendo il primo roguelite in assoluto). La
differenza:
\begin{itemize}
  \item il \textbf{roguelike} si fonda su una serie di punti fissati in una
        sorta di \emph{manifesto} che definisce che cosa rende un gioco
        roguelike; uno su tutti: a ogni morte si ricomincia la \emph{run} da
        capo;
  \item il \textbf{roguelite} rompe alcuni di questi punti, in particolare
        introducendo una \textbf{progressione esterna alla run}: il giocatore
        può potenziarsi fuori dalla singola partita per vincere più
        facilmente in quelle successive. \emph{Hades} ne è il caso principe.
\end{itemize}
Analogamente, un mobile game sviluppato dal direttore del corso fu
promosso come ``il primo \emph{donation game} al mondo'': un gioco che
convertiva il tempo di gioco in donazioni reali a enti benefici per animali.
Anche qui un neologismo creato lato marketing per definire un genere.

\subsection{Applicazioni fuori dall'industria videoludica}

Le competenze sviluppate non servono solo ai videogiochi:
\begin{itemize}
  \item i \textbf{render architettonici} oggi possono essere realizzati con
        Unreal, e diversi professionisti formati nel game design lavorano in
        questo campo;
  \item la \textbf{gamification}, cioè l'applicazione di regole videoludiche
        a contesti non ludici (mostre interattive, progetti culturali,
        consulenze come quella per applicare questi concetti alle macchine da
        allenamento di un'azienda di fitness), è in Italia uno dei modi più
        redditizi di spendere queste competenze, spesso con finanziamenti
        pubblici.
\end{itemize}

% ---------------------------------------------------------------------
\section{Il primo task}
\label{sec:primo-task}

\begin{enumerate}
  \item Creare un indirizzo e-mail \textbf{professionale} (consigliato: un indirizzo Gmail nella forma
        \texttt{nome.cognome@gmail.com}; l'anno di nascita è accettabile).
        Va bene anche un indirizzo professionale già esistente usato per
        lavoro. Da evitare indirizzi poco seri (``non \texttt{blackfire67}'').
  \item Una persona della classe raccoglie tutti gli indirizzi in un unico
        documento e lo invia via e-mail al docente, che aggiungerà tutti al
        Drive condiviso del corso (dove verranno caricate lezioni, slide e
        materiali, per avere i documenti tutti nello stesso posto).
  \item Installare Discord e iscriversi al server del corso, dove passano le
        comunicazioni ufficiali. È consuetudine, e utile, creare anche un
        gruppo tra studenti.
\end{enumerate}
\textbf{Scadenza}: prima della lezione del venerdì.

% ---------------------------------------------------------------------
\section*{Punti chiave del capitolo}
\addcontentsline{toc}{section}{Punti chiave del capitolo}

\begin{itemize}
  \item Il corso è una \textbf{simulazione di lavoro}: ci si tratta da
        colleghi, con una gerarchia da studio di sviluppo.
  \item Si impara con l'\textbf{inquiry-based learning}: lezione, task,
        dubbi, domande, condivisione, progetti.
  \item Il fallimento affrontato con la mentalità giusta è la principale fonte
        di crescita; i feedback riguardano il lavoro, mai la persona.
  \item Il lavoro del designer è per metà creativo e per metà
        \textbf{ricerca e metodo scientifico}; il game design non è una
        scienza esatta.
  \item \textbf{Documentazione} e \textbf{pitching} sono il cuore del
        mestiere.
  \item Il designer progetta meccaniche e sistemi (partendo dalle
        \textbf{3C}), documenta, media, prototipa, confuta le idee, studia i
        competitor, bilancia e fa tuning.
  \item La leadership del designer è \textbf{al servizio} del team, non
        verticale.
  \item I limiti (tempo, tecnologia, competenze, persone, denaro) restringono
        e indirizzano la creatività.
  \item \textbf{Le idee sono nulla, l'implementazione è tutto.}
\end{itemize}

%    % ...
%
%  (Con hyperref aggiungere anche \renewcommand{\theHchapter}{z.\arabic{chapter}}
%  nella parte zero, per evitare ancore duplicate.)
%  Tutte le etichette di questo capitolo hanno il prefisso "z2:"; i
%  rimandi al capitolo precedente usano le etichette "z1:" già definite in
%  GD0_1_cultura_del_feedback.tex.
% =====================================================================

\chapter[Tempo, idee e leadership]{Gestire il tempo, generare idee, guidare da dentro}
\label{cap:z2-tempo-idee-leadership}

Questo capitolo raccoglie la \textbf{seconda giornata} del workshop di due
giorni con cui si aprono le lezioni zero, tenuto dallo stesso relatore del
capitolo~\ref{cap:z1-feedback}, Stefano Cozzi\footnote{Per il cognome e per
il nome dello studio (Stormind Games) valgono le avvertenze della nota
all'inizio del capitolo~\ref{cap:z1-feedback}. In questa giornata la
trascrizione riporta soltanto il nome di battesimo, ``Stefano'', con cui il
relatore si cita più volte nei suoi esempi.}. Se la prima giornata spiegava
\emph{perché} le aziende tengono tanto alla cultura del feedback, la seconda
è dichiaratamente più pratica: ``ieri vi ho spiegato perché, oggi vi insegno
come lo fanno queste aziende''. I concetti della prima giornata (il
\emph{default to trust}, il ``potrei non farcela'' detto in anticipo, la
condivisione dei successi, il feedback frequente) ritornano di continuo,
applicati alla gestione del tempo, alle riunioni creative e alla leadership.

La giornata si è svolta di domenica, a chiusura di un fine settimana di
workshop\footnote{Il relatore dice ``oggi chiaramente è domenica'' e saluta
la classe con ``buon resto del weekend''; la prima giornata si era quindi
svolta di sabato. Più avanti parla di ``dodici ore'' di workshop in totale,
mentre nel capitolo~\ref{cap:z1-feedback} si parlava di circa dieci ore: la
differenza dipende probabilmente dal conteggio delle pause.}, con il
relatore alle prese con la voce un po' calata dal giorno prima (``la sfida
principale di oggi sarà contro la mia voce'').

Il percorso del capitolo segue quello della lezione:
\begin{enumerate}
  \item il programma della giornata (sezione~\ref{sec:z2-programma});
  \item le domande sulla prima giornata, che si trasformano in una lunga
        sessione su \textbf{colloqui e selezione}
        (sezione~\ref{sec:z2-colloqui});
  \item il \textbf{time management}: pianificare la giornata, le liste, le
        checklist, le interruzioni (sezione~\ref{sec:z2-pianificare});
  \item la \textbf{matrice di Eisenhower} per prioritizzare
        (sezione~\ref{sec:z2-eisenhower});
  \item come proteggere il proprio tempo: sprechi, distrazioni, il saper dire
        di no (anche ai meeting), il mito del multitasking e la chiusura della
        giornata (sezione~\ref{sec:z2-proteggere});
  \item le domande sul time management e una parentesi su tempo libero, ferie
        e benessere in azienda (sezioni~\ref{sec:z2-qa-tempo}
        e~\ref{sec:z2-benessere});
  \item il \textbf{brainstorming}: obiettivi, problemi tipici, un metodo
        completo (\emph{rapid ideation} e \emph{dot voting}) e la
        \textbf{matrice impatto--sforzo} per decidere
        (sezioni~\ref{sec:z2-brainstorming}, \ref{sec:z2-metodo}
        e~\ref{sec:z2-impatto});
  \item i consigli pratici, le domande sul brainstorming e il principio del
        \textbf{disagree and commit} (sezioni~\ref{sec:z2-qa-brainstorming}
        e~\ref{sec:z2-disagree});
  \item il \textbf{leading from within}: perché la leadership non è un
        titolo, le carriere divise tra manager ed esperti
        (sezione~\ref{sec:z2-leadership});
  \item il ``momento Piero Angela'' sulla biologia del gruppo e le domande
        che ne sono nate (sezioni~\ref{sec:z2-biologia}
        e~\ref{sec:z2-qa-biologia});
  \item la \textbf{piramide della fiducia}, il prendersi cura del gruppo e
        che cosa cambia quando si diventa manager
        (sezioni~\ref{sec:z2-piramide}, \ref{sec:z2-gruppo}
        e~\ref{sec:z2-manager});
  \item le domande finali e la chiusura del workshop
        (sezioni~\ref{sec:z2-qa-finale} e~\ref{sec:z2-chiusura}).
\end{enumerate}

Come per il capitolo precedente, una premessa di metodo: buona parte della
giornata è fatta di domande della classe e di racconti dell'esperienza del
relatore. Sono riportati con cura perché sono il modo in cui i principi
vengono resi concreti, e perché spesso contengono le indicazioni più
direttamente utilizzabili (soprattutto sui colloqui di lavoro).

% ---------------------------------------------------------------------
\section{Il programma della giornata}
\label{sec:z2-programma}

Come annunciato alla fine della prima giornata
(sezione~\ref{sec:z1-domani}), gli argomenti sono due:
\begin{itemize}
  \item un ``doppio topic'', \textbf{time management e brainstorming};
  \item il \textbf{leading from within}, che il relatore definisce ``la
        versione avanzata di ieri'': come guidare un team stando dentro al
        team, verso i colleghi e verso il gruppo con cui si lavora.
\end{itemize}
Poiché l'ordine è indifferente, il relatore lascia scegliere alla classe: dalle
poche preferenze espresse, la leadership va al pomeriggio, così che le domande
finali si concentrino su quella, mentre la mattina è dedicata a time
management e brainstorming. Il materiale c'è, ma il relatore avverte che la
giornata potrebbe finire prima delle quattro: non vuole trattenere nessuno se
non c'è niente di specifico da fare. Per l'eventuale tempo avanzato prima della
pausa pranzo propone un \textbf{argomento bonus}: i colloqui di lavoro,
selezione e \emph{recruitment}, visto che il relatore ha fatto e fa tuttora
selezione del personale e che per molti studenti sarà la cosa più immediata da
affrontare finito il percorso di studi (domande comuni, come rispondere,
eccetera).

Prima di cominciare, però, il relatore apre a \textbf{domande sulla giornata
precedente}: la classe ha avuto modo di ``metabolizzare'' i contenuti e
potrebbero essere emersi dubbi a posteriori. È proprio da una di queste domande
che nasce, in anticipo, la parte sui colloqui.

% ---------------------------------------------------------------------
\section{Colloqui di lavoro: capire che cosa cerca l'azienda}
\label{sec:z2-colloqui}

\subsection{L'annuncio come prima fonte di informazione}

La domanda: il giorno prima si era detto che è importante capire che cosa
l'azienda cerca davvero; ma se il primo colloquio è con le risorse umane, cioè
con qualcuno che magari ``di tecnico non sa'', come si fa a capirlo? Il
relatore risponde subito, nel dubbio che più tardi manchi il tempo.

La prima fonte è \textbf{l'annuncio}. Non viene scritto a caso: è approvato
sia dall'HR sia dall'\emph{hiring manager}\footnote{Nella trascrizione
``Airing Manager'' e ``Iron Manager''; allo stesso modo ``e-char'' e ``char''
stanno per HR.}, cioè la persona che seguirà il candidato una volta assunto,
e spesso dall'intero dipartimento HR. È calcolato e scritto in un certo modo
per far capire che cosa serve. Molti annunci moderni, oltre alla parte su
\emph{che cosa si verrà a fare}, contengono una parte su \emph{che cosa ci si
aspetta da questa persona}, cioè il comportamento e il carattere che servono
in quella posizione: ed è quella che dice molto sulla \textbf{cultura
aziendale}.

L'esempio è quello di due annunci per la stessa posizione (programmatore
senior) descritta in modo diverso:
\begin{enumerate}
  \item ``cerchiamo un programmatore esperto che possa prendersi cura dei
        nostri programmatori junior, aiutarli a crescere, sviluppare un
        ambiente di lavoro inclusivo e stimolante, in modo che i nostri
        progetti futuri siano all'avanguardia'';
  \item ``cerchiamo un programmatore esperto che sappia rispettare le
        deadline, seguire le indicazioni dei dati, sia indipendente nel suo
        lavoro e sappia relazionarsi con gli \emph{stakeholder} per capire e
        proporre nuove direzioni per il progetto''.
\end{enumerate}
Nel primo caso ci si aspetta un team di junior, probabilmente con un'età media
più bassa, con cui ci si relazionerà soprattutto; per un senior è l'occasione
di sviluppare le doti da \emph{people manager}, e al colloquio arriveranno
domande come ``hai mai lavorato con dei team? Con che età media? Ti sei mai
trovato a guidare un gruppo?''. Nel secondo caso il senior deve eseguire un
piano deciso o proporre qualcosa, e si relazionerà soprattutto \emph{verso
l'alto}; al colloquio arriverà piuttosto ``fammi questo test di
programmazione''. Sono, di fatto, due colloqui completamente diversi.

\subsection{Quanti colloqui e con chi}

Ormai è molto raro essere assunti con un solo colloquio: ne servono almeno
due, realisticamente \textbf{tre}, che è diventata la media. L'iter tipico:
\begin{enumerate}
  \item il primo con l'HR;
  \item il secondo con l'HR o con il futuro manager;
  \item il terzo sicuramente con il futuro manager, perché l'azienda vuole far
        conoscere al candidato la persona con cui lavorerà, e viceversa.
\end{enumerate}
Già le domande diverse dei vari interlocutori aiutano a capire che cosa
cercano. Ma c'è una parte fondamentale, che secondo il relatore molti
sottovalutano: il momento finale delle \textbf{``hai domande?''}.

\subsection{Le domande da fare al selezionatore}

Un principio di base: \textbf{i colloqui non sono ``devo farmi selezionare
dall'azienda'', sono ``dobbiamo selezionarci a vicenda''}. Le domande vanno
quindi preparate: ci sarà sempre qualcosa che, dalla \emph{job description} o
dal colloquio, non è chiaro al cento per cento. Nel dubbio, il relatore
suggerisce tre domande ``che funzionano sempre'':

\paragraph{``Se iniziassi domani, quale sarebbe il mio \emph{day to day}?''}
Ottiene due effetti. Il primo è pratico: descrive la giornata tipo e fa capire
se piacerebbe o no. Il secondo è più psicologico, ma il relatore sa che
funziona: costringe l'interlocutore a pensare al candidato \emph{già dentro
l'azienda}, cioè a che cosa potrebbe fare da subito e a quale impatto potrebbe
avere con le sue competenze.

\paragraph{``Se aveste una bacchetta magica, che cosa cambiereste subito
dell'azienda?''} Sposta la conversazione su di loro e va bene sia con l'HR sia
con l'hiring manager. Li obbliga a indicare una sfida interna che vorrebbero
risolvere, ed è sempre un'informazione utile: se entrambi rispondono che
cambierebbero il modo in cui si prendono le decisioni, viene da pensare che non
ci sia un percorso chiaro e trasparente da parte di chi sta in alto.

\paragraph{Lo scambio: ``vi do un motivo per non prendermi, voi datemene uno
per non venire''.} È più avanzata e il relatore non la usa sempre, ma la trova
particolarmente interessante. Permette al candidato di esprimere i propri dubbi
in modo costruttivo (``mi dite che lavorerei a stretto contatto con un team di
junior: un motivo per non prendermi è che non ho anni di esperienza nella
gestione di giovani sviluppatori, sarebbe la mia prima volta; però penso di
poterlo fare perché ho questo interesse e queste competenze''). Dall'altra
parte, il selezionatore dirà che cosa secondo lui sarebbe più difficile per il
candidato: si scopre così che cosa sta già pensando, e l'informazione torna
utile anche per i colloqui successivi.

\paragraph{E se il selezionatore si innervosisce?} ``Se un recruiter si
innervosisce per le domande del candidato, ha sbagliato lavoro.'' Al massimo il
relatore si innervosisce se un senior gli chiede che cosa faccia di preciso un
programmatore senior (``lo fai da dieci anni, lo saprai tu''); ma se un junior
chiede che cosa dovrebbe fare ogni giorno, lo si spiega volentieri. E in realtà
vale anche per i senior, perché \textbf{in ogni azienda lo stesso ruolo fa cose
diverse}: conviene sempre farsi spiegare che cosa intendono per quel titolo.
L'esempio viene dal mondo delle community: \emph{community manager} a Riot
indicava la persona a capo del team di \emph{community representative} di una
nazionalità (``community manager italiano''); a Blizzard indicava la persona
sopra l'intero dipartimento che gestiva quei team. Candidarsi per lo stesso
titolo non garantisce di aver capito per che cosa ci si sta candidando, e i
recruiter lo sanno.

Se poi un recruiter si innervosisce comunque, sta dando più risposte di quelle
che pensa: vuol dire che l'azienda non vuole domande, o non vuole interazioni
del tipo ``ma perché facciamo questa cosa?''. Che abbia formato i suoi
recruiter a innervosirsi, o che accetti che lo facciano, è informativo.

\subsection{Le domande che dicono molto sull'azienda}

Allo stesso modo, come detto il giorno prima, è informativo se un recruiter
chiede a una donna se vuole avere figli: ``ah, quindi a voi interessa sapere se
voglio avere figli: questo già mi dice che probabilmente non voglio lavorare
qui''.

Un esempio \emph{gender neutral}: la richiesta di vedere la \textbf{busta paga
precedente}. ``Quanto è la tua RAL? Mandami una busta paga.'' La risposta del
relatore è no: ``ti ho detto qual è la mia RAL; se vogliamo instaurare un
rapporto in cui non ti fidi di quello che dico, abbiamo un altro problema''. È
di nuovo \emph{default to trust} (sezione~\ref{sec:z1-default-trust}). Il
relatore è quasi sicuro che legalmente non si possa chiedere la busta paga di
un candidato\footnote{È l'opinione del relatore, espressa come tale (``sono
quasi sicuro''); non si tratta di un parere legale e non è stata verificata.}:
a quel punto diventa lui antagonista del recruiter, e visto che la
conversazione è finita in una ``zona grigia'', fa lui le domande (``come mai è
legale questa cosa che mi chiedi?''). Ha già deciso che probabilmente non
lavorerà lì, ma così magari evita che i candidati successivi debbano subire le
stesse richieste. Cita l'esempio perché sa che in Italia è abbastanza comune,
mentre all'estero non gli è mai successo, anche se avrebbe potuto ``sparare
qualsiasi cifra''.

Il motivo è che \textbf{non c'è alcun onere di dimostrare quanto si
guadagna}: chi spara troppo alto si autoesclude. Recruiter e manager non
arrivano al colloquio con carta bianca né con l'indicazione di risparmiare il
più possibile, ma con un \emph{range}: ``per questa posizione possiamo
spendere dai 30 ai 40 mila euro l'anno; paga questa persona a seconda della
sua seniority''. Se un candidato che guadagna 30 dice di prenderne 55 per
``provare a fare i soldi'' e il budget arriva al massimo a 40, il relatore non
chiederebbe nessuna prova: direbbe ``grazie, siamo in un range diverso, possiamo
salutarci''. La domanda sulla busta paga è quindi \emph{formulata male}, perché
è poco furba. Quella giusta è: \textbf{``hai già una RAL in mente per questa
posizione?''}. Si ottengono le stesse informazioni senza passare per la zona
grigia.

\subsection{La cifra minima mentale}

E se la domanda la fanno a un junior che non ha mai lavorato? Il consiglio vale
per tutti: \textbf{preparatevi sempre la vostra cifra minima mentale}, per
qualsiasi offerta di lavoro, prima ancora di fare il colloquio. Non serve
conoscere la media italiana per quella posizione o quanto paga quell'azienda:
serve sapere \emph{quanto si è disposti ad accettare senza andare in negativo},
tenendo conto dei propri costi (la benzina per andare in ufficio ogni giorno, un
eventuale trasferimento, e così via). Quella cifra va sempre avuta in testa.

\subsection{La scrematura iniziale e il portfolio}

Uno studente chiede come funzioni la risposta a una candidatura: arriva un
rifiuto diretto, oppure c'è una scrematura iniziale e poi i colloqui? E vale
per tutti i ruoli? Sì: c'è una \textbf{scrematura iniziale} fatta dall'HR o, in
alcune aziende, direttamente dal software dell'HR. Si scarta chi è chiaramente
fuori da ciò che si cerca:
\begin{itemize}
  \item si cerca un \emph{graphic designer} e arriva il CV di un programmatore
        (``faccio questi esempi perché purtroppo succede'');
  \item si cerca un \emph{animator} tecnico bravo con MotionBuilder, lo si
        scrive nell'annuncio, e il candidato risponde ``so fare tutto tranne
        MotionBuilder'': c'è un motivo tecnico per non poterlo prendere;
  \item problemi di seniority: si cerca un senior e risponde un junior con un
        anno di esperienza.
\end{itemize}
Se la figura non è già chiaramente fuori target, si fa il primo colloquio.

\paragraph{Gli anni di esperienza negli annunci junior.} Spesso gli annunci
junior chiedono come minimo tre anni di esperienza. \textbf{Se l'annuncio dice
``junior'', ignorate i requisiti di esperienza}: o sono lì per sbaglio, o sono lì
per scremare. Il secondo caso è più \emph{borderline}, ma il relatore l'ha
visto fare in alcune aziende americane: si mettono requisiti molto alti,
chiaramente irraggiungibili, per vedere chi ha la ``faccia tosta'' di mandare
comunque il CV (``non me ne frega se chiedete un \emph{bachelor's degree}, non ce
l'ho, mando lo stesso''). Così però si filtra una grossa fetta di popolazione e
si seleziona un certo tipo di carattere: si stanno cercando ``i Top Gun''. Non è
quindi necessario iniziare con uno stage: si può iniziare direttamente da
junior.

\paragraph{Chi guarda il portfolio.} Il relatore non ha mai avuto un portfolio,
ma per le figure che lo richiedono parte comunque dal CV (per escludere, ad
esempio, il designer che si candida da programmatore) e guarda il portfolio
solo a parità con altri candidati: da HR può aprirlo e dire ``bello'', ma non
può dare un giudizio tecnico. Il portfolio lo apre chi sarà l'hiring manager,
che è sempre molto più tecnico. L'ordine quindi è: l'HR guarda il CV e si salva
il portfolio; se il candidato passa la prima scrematura, CV e portfolio vanno
a un tecnico che conferma (``sì, sa fare queste cose'') o boccia (``il
portfolio indica chiaramente che non sarebbe in grado, digli pure che non
facciamo il colloquio'').

\subsection{Test pratici, da casa e con qualsiasi strumento}

\paragraph{Si può chiedere un test pratico al posto del portfolio?} Uno
studente chiede se, con un portfolio che non rispecchia le proprie capacità, si
possa chiedere direttamente un test. No: se l'azienda ha bisogno di verificare
una competenza ha già pronto un test e lo fa fare a prescindere. Il relatore,
per esempio, fa fare un test ai programmatori e lo fa correggere dal suo senior
developer; per gli artisti si chiedono animazioni di prova, chiaramente veloci.
Non chiede mai test troppo lunghi, perché non vuole far lavorare nessuno gratis:
non chiede ``fammi un pitch di progetto'', chiede uno spunto per capire come si
ragiona e come si lavora. Per il resto, \textbf{se il portfolio non è
indicativo, va aggiornato il portfolio}.

\paragraph{Test dal vivo o a casa?} Il relatore li fa fare sempre a casa. Sa
benissimo che un programmatore a casa lavorerà con due monitor, il codice da
una parte e Google o ChatGPT dall'altra. Ne ha discusso con i suoi senior
programmer: interessa che qualcuno sappia le cose a memoria, o che sappia
trovarle e fare quel lavoro a prescindere? A lui interessa che il lavoro
\textbf{venga fatto}: che il candidato abbia imparato il codice a memoria o che
in venti minuti trovi le informazioni, le metta insieme, le faccia elaborare a
ChatGPT, le corregga e consegni il codice che serve, non ha preferenze. Quando
si lavora le indicazioni saranno ``fai questa cosa il prima possibile'': se la
si sa fare anche con strumenti aggiuntivi, ben venga, ``anzi, spiegalo anche
agli altri''. Farla a memoria va benissimo, ``ma perché mettere l'\emph{hard
mode}?''. E comunque ci sono cose che non si possono delegare a uno strumento:
se arriva un codice che non funziona e va fatto il \emph{debug}, bisogna
saperlo fare e saper progettare i test. La figura più completa, per il
relatore, è anche quella che sa trovarsi le informazioni e gli aiuti.

Le prove dal vivo sono soprattutto domande di \emph{problem solving}: il
relatore vuole sentire come si ragiona e chiede espressamente di \textbf{pensare
ad alta voce}, per seguire il percorso mentale che porta alla conclusione. Sono
più domande impegnative che veri test. I test veri e propri si fanno con i
propri tempi e i propri strumenti: ``scegli il momento, ti mandiamo il test, hai
due ore per rimandarlo indietro''. Questo è l'accordo; in quelle due ore si può
fare quello che si vuole.

\subsection{Junior, mid e senior}

\paragraph{Perché ci sono così pochi annunci per figure mid?} Uno studente ha
notato su LinkedIn e sulle piattaforme di recruiting un divario notevole tra le
richieste di junior e quelle di mid. Il motivo è che per molte aziende,
Stormind inclusa, \textbf{la differenza tra junior e mid è sparita}: o si è
senior o non lo si è. Ci si è resi conto che la differenza tra un junior e un
mid è che il mid si è abituato a lavorare in quell'azienda e il junior non
ancora; ma un mid che arriva con tre anni di esperienza in un altro studio (ad
esempio Milestone\footnote{Nella trascrizione ``se ha fatto tre anni in
milestone'': si tratta verosimilmente di Milestone, studio di sviluppo
milanese.}) va comunque ``ripreso come junior'', perché si lavora in modo
diverso e deve reimparare. Un senior, invece, ha molte più competenze e le sa
applicare in un certo modo, e quella distinzione resta netta.

\paragraph{Ci si può candidare anche se il livello non coincide?} Un'altra
studentessa dava per scontato che, se un annuncio non specifica ``junior'', sia
per mid. Se non è specificato ``senior'', secondo il relatore si può mandare il
CV chiedendo di valutare. Certo, se è specificato junior e si è junior le
possibilità di essere contattati sono maggiori. Ma vale lo stesso discorso del
livello di inglese fatto il giorno prima: si cerca qualcuno che \emph{sappia
fare quelle cose}, non una certificazione. Se serve un programmatore che sappia
usare Unreal Engine~5 e sistemare la cosa X o Y, che sia mid da due anni o
junior da uno non è un problema.

\subsection{Gli strumenti di intelligenza artificiale}

Una domanda in chat riguarda l'uso quotidiano dell'intelligenza artificiale in
azienda. Il relatore sta facendo fare corsi di IA al suo team ``perché voglio che
la gente lavori meno'': vuole che si sfrutti qualsiasi strumento disponibile,
perché gli altri lo stanno facendo e non vuole restare indietro. Se un lavoro
si può fare in tre ore di meno, non c'è alcun vantaggio a farlo in modo lungo e
tortuoso. Il rischio esiste: se si fa fare tutto a ChatGPT e non ci si tiene al
passo con il proprio \emph{know-how}, quando ci sarà da risolvere qualcosa ci si
potrebbe ``impallare''. Ma sta al singolo restare aggiornato e saper fare il
proprio lavoro: \textbf{gli strumenti devono essere strumenti, non un sostituto
delle competenze}. ``Non è una partita di calcio, con regole su come si devono
fare le cose o un \emph{dress code}'': il lavoro va fatto nel modo più efficiente
possibile.

A uno studente a cui era saltata la connessione proprio in questo momento il
relatore risponde che la lezione è registrata e si può recuperare dal video.
Poi chiude la parentesi: ``la parte difficile era finire ieri in tempo, ce
l'abbiamo fatta''.

% ---------------------------------------------------------------------
\section{Time management: pianificare la giornata}
\label{sec:z2-pianificare}

\subsection{Time management o prioritizzazione?}

La domanda di partenza alla classe è: \emph{chi di voi, in questo momento, fa
time management su se stesso}, cioè si organizza la giornata all'inizio e
organizza il proprio tempo? Qualcuno risponde di sì, qualcuno ha uno
\emph{schedule} settimanale con orari da rispettare, qualcuno ``ci prova''
(``hai fatto bene a specificare \emph{ci provo}: magari non la rispetto''). Il
relatore lo chiede perché è uno di quegli argomenti in cui molti, dopo aver
visto la spiegazione, si accorgono di non fare davvero time management: fanno
piuttosto \textbf{prioritizzazione}, ad esempio una lista di cose da fare, che
è un po' diverso.

Nessuna delle tecniche è complicatissima, ma è utile sapersi muovere. Il
relatore colloca volutamente questo modulo \emph{dopo} quello sul feedback: saper
gestire il proprio tempo permette anche di dare feedback agli altri sulla loro
gestione del tempo.

\subsection{La \emph{must do list}}

Il punto di partenza è il \textbf{planning iniziale}: non improvvisare la
giornata. Improvvisando si è in balia di una grande limitazione umana: ci si
dimentica le cose, oppure ci si lascia trascinare dal task più urgente o
dall'ultimo arrivato. Esistono modi di dire come ``se arriva una mail o rispondi
subito o diventa un task per domani'', perché il ``rispondo dopo'' non arriva
mai; più avanti si vedranno modi più raffinati di gestire le zone grigie, ma la
pianificazione resta fondamentale.

La prima cosa da fare \emph{non} è la lista delle cose da fare oggi, ma la lista
delle cose che \textbf{assolutamente} vanno fatte oggi. Le due liste non devono
coincidere: la seconda richiede un livello di ragionamento in più. Sono le cose
per cui, se alle sei si stacca e non sono fatte, si resta finché non sono fatte,
perché sono imprescindibili.

\paragraph{Il primo momento per dire ``potrei non farcela''.} La \emph{must do
list} è anche il primo indicatore della giornata. Se è troppo corposa (dipende
naturalmente dalla complessità dei task), quello è il momento, ``alle 9:01 se
iniziate alle 9'', per dire al proprio manager \emph{oggi sarò in difficoltà}.
Alle 9:01 il manager ha l'intera giornata per aiutare; se invece si pecca di
ottimismo o di presunzione (``ce la posso fare'') e glielo si dice alle 17:59,
non può fare nulla. È lo stesso discorso del giorno prima
(sezione~\ref{sec:z1-triangolo}): anche facendo la to-do list si applicano i
concetti della cultura del feedback, ed è il motivo per cui questo modulo viene
dopo. Ci saranno vari momenti di questo tipo nella giornata, ma questo è il
primo.

\paragraph{E se ci si accorge a metà giornata di non farcela?} Uno studente
chiede che cosa fare se, convinto di farcela, si accorge a metà giornata di non
poter completare bene la \emph{must do list}. Il relatore lo concederebbe
\emph{una volta}: se si vogliono testare i propri limiti, va bene, ma
comunque si dice al capo ``oggi potrei essere tirato, ci provo'', e sarà lui a
valutare se il rischio è accettabile. Dipende dal rischio: se il \emph{must do}
è consegnare una milestone a un cliente, ``nel dubbio facciamolo fare a
qualcun altro, non possiamo correre quel rischio''; se il rischio è che domani
un animator riceva il codice con tre ore di ritardo, lo si può correre. ``È un
po' questo il lavoro.''

\paragraph{Cinque minuti e una stima.} Tutto questo va fatto in circa cinque
minuti. Se non è chiaro che cosa si deve fare in una giornata, \textbf{si
chiede}: ogni giorno dovrebbe essere chiaro. Poi si cerca di \textbf{stimare il
tempo} di ogni cosa. Si diventa più bravi col tempo: alla quinta volta che si
fa il report mensile si sa che ci vuole un'ora e ventidue minuti; la prima volta
si può chiedere a un collega quanto ci mette lui, sapendo che per noi potrebbe
essere diverso.

\subsection{To-do list e checklist}

Esistono molti strumenti per tenere to-do list e checklist; il relatore non ne
consiglia uno in particolare e cita come esempi Any.do, i promemoria di Windows
o Samsung Notes\footnote{Nella trascrizione ``AnyDo, Windows Memo''. Any.do è
un'app di liste e promemoria; ``Windows Memo'' indica probabilmente le note o i
promemoria integrati in Windows (Sticky Notes, Microsoft To Do).}: ``potete usare
davvero la qualunque''. L'importante è tenere la lista \textbf{visibile} da
qualche parte sul computer. Alcuni accorgimenti:
\begin{itemize}
  \item i \textbf{must do} vanno sempre \textbf{in cima} alla lista;
  \item quando durante la giornata arriva un task nuovo, la cosa da fare subito
        non è \emph{eseguirlo}, ma \textbf{aggiungerlo alla lista}: ``se arriva
        una mail, aggiungete subito alla lista \emph{rispondi alla mail}'',
        decidendo se è diventato un \emph{must do} o se è comunque un task per
        oggi;
  \item non limitarsi alle cose operative: il relatore ha iniziato ad
        aggiungere anche task come ``parla con Tizio'' o ``parla con Caio'',
        perché è da un po' che non chiede loro un aggiornamento; se non li
        mette in lista si accorge di non farli;
  \item spuntare le voci completate: per quanto minimo, rilascia
        \textbf{dopamina} (se ne parlerà nel pomeriggio,
        sezione~\ref{sec:z2-biologia}), aiuta a motivarsi e ad andare avanti.
        Sono piccole cose, ma si sommano.
\end{itemize}

\paragraph{Priorità contro recenza.} Che cosa succede se si esegue sempre
l'ultimo task arrivato? Lo chiede il relatore alla classe. Si interrompe il
flusso di lavoro e si cambia continuamente attività, certo; ma soprattutto si
stanno \textbf{mescolando priorità e recenza}: si sta dicendo ``non importa come
ho prioritizzato i task, la mia priorità è sempre l'ultimo arrivato'', che non
vuol dire fare le cose più importanti. È il motivo per cui tanti dipartimenti IT,
quando si chiede loro qualcosa, rispondono ``mi apri un ticket?''. È diventata
quasi una battuta, ma il concetto è chiarissimo: devono poter prioritizzare
quella richiesta rispetto alle altre 112 arrivate in giornata, e senza ticket
non possono confrontarla con le altre. Il relatore è totalmente a favore: ``certo
che ti apro un ticket, altrimenti tu come fai a tracciare il tuo lavoro?''.

\subsection{La checklist: l'esempio dell'aviazione}

La \textbf{checklist} non è la task list. La task list contiene le cose da fare;
la checklist contiene \textbf{i passi da seguire quando si esegue un task
specifico}. L'esempio è quello delle \textbf{checklist pre-volo}, introdotte
nell'aviazione (secondo il relatore tra gli anni Sessanta e Settanta): chi vola
spesso nota che i piloti seguono sempre un ordine simile nei controlli, tanto
che il relatore capisce quanto manca al decollo da quando controllano i flap, i
motori, quando accendono e spengono le luci. L'introduzione delle checklist
avrebbe ridotto di circa il 75\% gli incidenti in fase di decollo\footnote{Date
e percentuale sono quelle riferite in lezione e non sono state verificate.
L'origine della checklist pre-volo viene di solito fatta risalire al 1935, dopo
l'incidente del prototipo del Boeing B-17, causato da un controllo dimenticato da
un equipaggio molto esperto; la diffusione sistematica delle checklist e dei
protocolli di cabina è invece successiva.}, e soprattutto per i \textbf{piloti
più senior}. Perché? Perché chi ha fatto un decollo centomila volte non va a
controllare l'alettone destro: ha funzionato centocinquantamila volte, funzionerà
anche stavolta. Quando ci si abitua a fare qualcosa, si sbaglia proprio quando
qualcosa non va. La checklist ha obbligato a fare quei controlli ogni volta e ha
ridotto in modo drastico gli incidenti.

Il relatore usa checklist anche per cose banali, ad esempio le \textbf{email}:
essendo una persona abbastanza ansiosa su queste cose, se prepara una mail alla
leggera finisce per passare troppo tempo a ricontrollare l'oggetto, il
destinatario, l'allegato, e comunque si dimentica qualcosa. Seguire una checklist
sembra più complesso, ma \textbf{fa risparmiare tempo}: se la checklist è
completa, si sa di aver fatto tutto e si può andare avanti senza continuare a
pensare ``ho chiuso il gas? ho spento la luce?''. È come la lista della spesa: se
tutto è segnato e preso, punto. È un investimento: trenta secondi in più per
fare la cosa, invece di cinque minuti a sperare di averla fatta giusta.

\paragraph{E se ci si dimentica qualcosa nella checklist?} Succede. A quel
punto \emph{preparare la checklist} è un task in sé: un giorno in cui si ha
tempo si fa quel task pensando a quali passi si seguono e a come assicurarsi che
funzioni. Nessuna checklist è perfetta: ci saranno sempre compromessi tra ``devo
davvero controllarlo ogni volta?'' e ``posso farne a meno?''. Però, in generale,
\textbf{se qualcosa anche solo una volta ha dato problemi, va in checklist},
perché vuol dire che potrebbe darne di nuovo. Nessun manager rimprovererà
``come hai fatto a non prevedere che questa cosa poteva andare male?'': si
diventa esperti anche nelle checklist di un task facendo quel task, ed è un
processo iterativo. Se qualcosa è stato dimenticato e ha creato problemi, lo si
aggiunge per la volta dopo. Sbagliare una o due volte non è un problema;
l'importante è che la situazione migliori la volta successiva.

\paragraph{L'inizio di un junior e l'\emph{onboarding}.} A uno studente che
osserva che l'inizio di un junior è quindi inevitabilmente un po' travagliato,
il relatore risponde che è giusto così: ``se volessi il dipendente perfetto da
subito, cercherei un senior, non un junior''. Allo stesso modo non capisce le
aziende che prendono un junior e non lo fanno seguire da nessuno: che cosa ci si
aspetta, che assorba passivamente tutto quello che ha intorno e impari ``per
osmosi''? Può succedere, ma imparerà a fare le cose come vuole lui, perché
nessuno gli sta dicendo ``meglio farle così''. Moltissime aziende, oltre al
manager (che magari ha quindici \emph{report} e non ha il tempo materiale di
seguire un junior come andrebbe seguito), assegnano un \textbf{referente}, di
solito un senior, con il compito di prendersi cura della persona nuova e
insegnarle a fare le cose. Altrimenti l'\textbf{onboarding}, cioè il tempo
necessario per arrivare al cento per cento nel proprio lavoro, dura un anno.
Più risorse l'azienda dà all'inizio, prima la persona diventa operativa: è
nell'interesse dell'azienda dare il massimo supporto a chi arriva.

Il relatore diffida delle aziende che ``vogliono vedere se uno se la cava''. È
il discorso dei cento metri e della corsa a ostacoli\footnote{Il relatore lo
introduce come un discorso già fatto (``si torna al discorso: vuoi fare i cento
metri o la corsa a ostacoli?''), ma il passaggio non compare nella trascrizione
della prima giornata riportata nel capitolo~\ref{cap:z1-feedback}.}: non si può
far correre uno sui cento metri e l'altro con gli ostacoli e aspettarsi tempi
simili, perché chi ha gli ostacoli ci metterà per forza di più. Il supporto ai junior è un
investimento sensato.

\subsection{\emph{First contact resolution}}

Un concetto che viene dai call center: \textbf{quante richieste vengono risolte
al primo contatto}. È importante quando un task richiede iterazioni o
comunicazione con altri: bisogna organizzarsi come se si dovesse risolvere la
cosa \emph{in una sola passata}. Se viene chiesto un modello 3D, lo si fa giusto
la prima volta, al meglio delle proprie possibilità. Se si deve parlare con
qualcuno, si raccolgono tutte le richieste in una.

L'esempio è quello dei ticket che il relatore mandava ai colleghi del supporto di
Blizzard: ``ciao, ho un problema, non riesco a loggarmi'' --- ``mi dici lo
username?'' --- ``è questo'' --- ``mi dici l'email?'' --- ``è questa'' --- e così
via, con un paio d'ore tra un messaggio e l'altro perché il ticket deve fare il
suo giro. Tutto poteva essere riassunto in ``ciao, mi servono username, email e
questo e quest'altro''. Unificare le richieste su una singola cosa aiuta a
risolverla il prima possibile; bisogna avere chiaro \emph{che cosa serve prima
di fare una richiesta}, in modo che ce ne sia una sola. In molti sistemi di
supporto, se si apre un secondo ticket per aggiungere informazioni al primo, il
primo, anche se era quasi in lavorazione, viene rimandato in fondo alla coda:
serve a disincentivare chi manda continuamente ticket del tipo ``mi sono
ricordato questa cosa''. O si manda tutto subito, o ci vorrà tempo.

\subsection{Non iniziare senza le informazioni}

Corollario: \textbf{se non si ha tutto quello che serve per completare un task,
non lo si inizia}. Prima ci si procura quello che manca. Al relatore capita con
la pianificazione dei due anni di lavoro successivi: se non ha tutte le
informazioni, prima lavora per averle (``mi dai questa informazione, mi dai
quell'altra''), poi pianifica; altrimenti pianificherebbe qualcosa di cui non ha
completa visibilità.

\subsection{Il costo delle interruzioni}

Alternare continuamente l'attenzione tra i task e, soprattutto, reagire di
continuo agli input esterni (la mail, il messaggio, il collega che parla) fa
\textbf{sparire il tempo dalla giornata}. Se mentre si prepara un report in
Excel si viene interrotti (dal gatto, da una mail, da un collega, da una
telefonata) e poi si torna al foglio, servono in media \textbf{un quarto d'ora}
per tornare mentalmente al cento per cento; chi è particolarmente bravo ci mette
dieci minuti, ma si parla comunque di dieci-quindici minuti\footnote{Ordini di
grandezza come riportati in lezione. Le ricerche più citate sul tema (ad esempio
quelle di Gloria Mark sulle interruzioni nel lavoro d'ufficio) indicano tempi
medi di ripresa anche superiori, intorno ai venti minuti.}. Se si viene
interrotti di continuo, si passa più tempo a cercare di tornare al cento per
cento che a lavorare. È una partita persa in partenza.

Per questo il giorno prima si era parlato dei \textbf{semafori} sulla scrivania
(sezione~\ref{sec:z1-dare}): il semaforo può essere anche semplicemente il
proprio stato su Slack, su Discord o su qualunque strumento si usi con i
colleghi. Se si è occupati, si è occupati: è una propria valutazione che,
accettando di essere disturbati di continuo, quella cosa non si riuscirà a
farla, o almeno non in tempo.

% ---------------------------------------------------------------------
\section{Prioritizzare: la matrice di Eisenhower}
\label{sec:z2-eisenhower}

Ora che si hanno task list, checklist, \emph{must do} e \emph{nice to do}, resta
la domanda di prima: se arriva un task nuovo e si decide di non fare per prime le
cose nuove, \emph{come si capisce se è davvero urgente}? Serve un metodo, ed è la
\textbf{matrice di Eisenhower}\footnote{Nella trascrizione anche ``Jason
Howard''. Dal nome del presidente statunitense Dwight D.~Eisenhower, a cui si
attribuisce la distinzione tra urgente e importante.} (la classe prova a
indovinarne il nome proponendo altre matrici, come una matrice dei rischi, ma è
un'altra cosa). Ogni task che arriva va posizionato su questa tabella:

\begin{center}
  \begin{tabular}{@{}p{0.2\textwidth}p{0.36\textwidth}p{0.36\textwidth}@{}}
    \toprule
    & \textbf{Urgente} & \textbf{Non urgente} \\
    \midrule
    \textbf{Importante} &
    \textbf{Fallo} (\emph{do}): diventa un task di oggi. &
    \textbf{Pianificalo} (\emph{schedule}): va fatto, ma non subito. \\
    \addlinespace
    \textbf{Non importante} &
    \textbf{Delegalo} (\emph{delegate}): va fatto, ma non serve che lo faccia
    tu. &
    \textbf{Eliminalo} (\emph{delete}): non lo fai. \\
    \bottomrule
  \end{tabular}
\end{center}

\paragraph{Urgente e importante: lo fai.} Sono di solito task con una
\textbf{deadline molto chiara} (entro oggi, entro un'ora, entro X) e con
\textbf{conseguenze} se non vengono fatti: da qui l'``importante''.

\paragraph{Importante ma non urgente: lo pianifichi.} Va fatto, ma non subito:
o lo si mette tra le cose di oggi ``ma dopo'', o anche alla settimana successiva.
L'esempio: un collaboratore scrive ``il mio rinnovo è tra tre mesi, ne possiamo
parlare?''. È importante, perché il rinnovo di una persona è importante; non è
urgente, se ne può parlare anche domani, così da potersi preparare e fare altre
cose nel frattempo.

\paragraph{Urgente ma non importante: lo deleghi.} Qui ``non importante''
\textbf{non} vuol dire che il task non sia importante, ma che \textbf{non è
importante che lo faccia tu}, che è molto diverso. Se arriva ``mi serve questo
cambio a questa animazione 3D'' e ciò che conta è il cambio, non chi lo fa, si può
dire ``sono d'accordo che è urgente, ma non serve che lo faccia io'' e chiedere a
qualcun altro. Sono task che vanno fatti ma non richiedono la propria particolare
\emph{expertise}.

Il relatore insiste molto su questo punto con manager e senior: quando arriva un
task urgente, i senior tendono a dire ``lo faccio io, perché è urgente''. Ma se un
junior può farlo, lo si fa fare al junior e si investono le competenze più alte
del senior in qualcos'altro: il senior può fare il lavoro del junior, ma il
junior non può fare il lavoro del senior mentre il senior fa quello del junior.
Le persone non sono intercambiabili. Si usa il \textbf{minimo sforzo
necessario}: se una cosa la può fare un mid, la fa il mid, non il senior.

\paragraph{Né urgente né importante: lo elimini.} È il quadrante che quasi
nessuno usa: tutti fanno, delegano o pianificano, ma non pensano mai a
\emph{delete}. Invece è fondamentale. A volte è corretto dire: ``questa cosa non
dovrei farla io, non è urgente, non la faccio; sappi che non la faccio, e se c'è
un problema ne parli con il manager''. Non si sta delegando a qualcun altro: si
sta dicendo che quella cosa così com'è \textbf{non andrebbe fatta}, perché non
ne sono chiari l'importanza e il perché. Nella propria lista di oggi non ha
spazio, e non è abbastanza importante per essere fatta da sé nemmeno nei giorni
successivi.

Un esempio di \emph{delete} che non è un vero \emph{delete}: arriva una richiesta
di meeting di tre ore sul futuro del progetto, con invitati tutti gli art
director, il game director e così via. Il relatore spesso risponde ``io non
vengo'': per lui è \emph{delete}, ma il meeting si fa lo stesso. Uno studente
osserva che a volte si tratta di una richiesta mal formulata o poco chiara: è
esattamente così. ``Nella sua forma per me è no; se poi mi spieghi meglio perché
dovrei farlo io o perché è urgente, allora cambia.''

% ---------------------------------------------------------------------
\section{Proteggere il proprio tempo}
\label{sec:z2-proteggere}

\subsection{Evitare gli sprechi: un limite che nessuno può insegnare}

Evitare gli sprechi di tempo è difficile, soprattutto da remoto: è una sfida
costante con se stessi. Ha molto a che fare con il carattere, il
\emph{self-control} e (``passatemi anche questa'') la professionalità. Il
relatore può dare tutti i consigli del mondo, ma se dentro di sé non si percepisce
l'urgenza di fare il proprio lavoro quando lo si deve fare, non lo si può
insegnare. Se si è a casa e invece di lavorare si gioca a qualcosa su Steam,
nessuno deve spiegare che non si dovrebbe. Ci sono limiti oltre i quali non si
può andare.

Quello che si può fare è sapere che \textbf{ogni interruzione ha un costo
sommerso}: ogni messaggio, domanda, telefonata, ogni cambio di focus. Sono i
quindici minuti di prima, che possono arrivare a trenta a seconda del tipo di
task.

\paragraph{È più difficile tornare su un task logico o creativo?} Il relatore lo
chiede alla classe: fa più fatica un artista a tornare sul suo modello 3D o un
\emph{data analyst} a tornare sul suo foglio Excel? Le risposte si dividono
(``creativo'', ``entrambi''), ed è proprio il punto: dipende a chi lo si chiede. I
creativi rispondono ``creativo'', gli analisti ``analitico'', perché ognuno ha una
chiara comprensione di che cosa serve per essere davvero concentrati su ciò che
fa. Il creativo sa che per essere pienamente creativo gli servono l'ambiente, la
concentrazione, gli spunti, gli esempi da tenere chiari in testa; l'analista deve
avere in testa i conti, i fogli di dati e così via. Sono \textbf{ugualmente
difficili}, ma solo chi fa quella cosa ne apprezza fino in fondo la difficoltà. Si
può avere un assaggio dell'altra parte, ma la piena comprensione arriva solo
facendo: è come guardare uno sport in tv e pensare ``potrei farlo anch'io'',
finché non ci si trova lì.

\subsection{Distrazioni volontarie e involontarie}

Uno studente, che sta leggendo manuali di \emph{design thinking}, chiede come
conciliare tutto questo con l'idea che distrarsi, cioè non restare focalizzati,
aiuti a trovare un nuovo punto di vista per risolvere un problema. Qual è il
giusto equilibrio? Il relatore è \textbf{totalmente d'accordo} con quel concetto:
sta parlando di \textbf{interruzioni non volontarie}. Sa benissimo che dopo cinque
ore su un foglio Excel ``esplode la testa''; ma se si pianifica un quarto d'ora di
pausa e lo si rispetta, si sta investendo sulla propria creatività ed
efficienza. Non è un investimento dover rispondere a una chiamata ogni tre
minuti.

Anzi, il consiglio è \textbf{crearsi interruzioni pianificate}, sotto il proprio
controllo: fare pause, fare altro, soprattutto nel \emph{problem solving}. Il
relatore stesso, se non trova una soluzione, si ferma, fa un altro task e poi ci
torna; ma è una sua scelta. La discriminante è che la pausa sia \emph{scelta} e
non imposta: anche se si fanno le pause giuste, se si viene interrotti proprio
nel momento di massima concentrazione (raggiunta magari grazie alla pausa) si
butta via tutto. Uno studente cita la \textbf{tecnica del pomodoro}: è
esattamente quello, piccole pause pianificate che aiutano a essere creativi,
purché restino pianificate. Se i quindici minuti diventano ``vabbè, facciamo
mezz'ora'' o ``lo faccio nel pomeriggio'', non è più la tecnica del pomodoro:
diventano interruzioni, o un pretesto per non lavorare.

\subsection{Il peggior nemico siamo noi}

E chi sono i peggiori nemici in questo? ``Noi stessi'': siamo noi ad
autoconvincerci che venti minuti vanno bene quanto quindici. Le
\textbf{distrazioni autoinflitte} sono spesso le peggiori. Il relatore porta
l'esempio di una sua giornata, che ha annotato perché lo aveva fatto ridere
quanto tempo era riuscito a buttare via senza che nessuno glielo imponesse:
\begin{itemize}
  \item gli piace avere qualcosa in sottofondo mentre lavora (``ha a che fare con
        l'ADHD sicuramente''\footnote{Nella trascrizione ``IDACD''. Più avanti il
        relatore dice di avere l'ADHD ``come dicevo ieri''; il capitolo 0.1 non
        lo riporta, ma riferisce il suo racconto sull'``overthinking'' anche
        nell'ordinare una pizza (sezione~\ref{sec:z1-regalo}).}),
        e in genere lo aiuta. Il problema è quando il sottofondo è troppo
        interessante: la musica va bene, ma se nella playlist di Spotify parte un
        podcast di Alessandro Barbero sulla guerra civile americana, finisce per
        ascoltare quello più del lavoro;
  \item un'abitudine che ha poi perso: se faceva straordinari alle due di notte,
        da solo, per qualcosa di urgentissimo, apriva comunque Slack a
        controllare i messaggi, e la mail. ``Come se alle due di notte qualcuno
        mi scrivesse qualcosa di urgente'';
  \item la pausa con un po' di TikTok, dopo la quale si accorgeva di aver
        dimenticato su quale slide della presentazione stava lavorando.
\end{itemize}
La regola: \textbf{inserire i break è meglio che subirli passivamente}. Un conto
è dire ``faccio una pausa intenzionalmente'', un conto è ``mi sono distratto e ho
perso due ore di lavoro''. Chi ne è particolarmente soggetto non deve
demoralizzarsi: è anche un allenamento.

\subsection{Spezzettare i task grandi}

Un altro concetto utile: se si fa molta fatica a iniziare i task grossi, perché
sembrano infiniti (``più tardi inizio, più tardi inizia questa sofferenza''),
\textbf{vanno spezzettati}. Un task grosso non è quasi mai un task grosso: è
spesso un insieme di task piccoli. Si comincia a farne dei pezzi e si va avanti
piano piano, spuntandoli. Aiuta tantissimo anche nel \emph{problem solving}: un
problema gigante va scomposto in problemi più piccoli, da risolvere per primi, per
poi ricomporre il puzzle.

\subsection{Accettare di dire di no}

Tornando alla matrice di Eisenhower: bisogna \textbf{accettare di dover dire di
no}, perché dire sempre sì porta a problemi ben peggiori di quanto si pensi.

\paragraph{Il collega che diceva sempre sì.} A Riot il relatore aveva un collega
estremamente disponibile, amichevole e collaborativo: qualsiasi cosa gli si
chiedesse, rispondeva ``certo, la faccio io'' (``mi porti fuori il cane?'' ---
``certo''). Il problema è che era umano: a un certo punto si supera la quantità di
cose umanamente possibili in una giornata. In più tendeva a offrirsi volontario
per far vedere che era intraprendente e aveva iniziativa, e nei progetti che
coordinava il relatore doveva spesso rifiutare la sua candidatura perché era
sicuro che non ce l'avrebbe fatta. Il suo dire sempre sì, a un certo punto, lo ha
danneggiato: era quello che dice sempre che lo fa lui, ma non è detto che lo
faccia. ``Ce la fai?'' --- ``Sì, sì, assolutamente'' --- e poi ``scusa, non ce l'ho
fatta''. ``Ok, però ora hai messo nei guai me e il resto del team.'' Quella
disponibilità era più dannosa di un ``no, non posso''. È lo stesso discorso del
giorno prima: dire all'ultimo momento che non ce la si fa è peggio che dirlo a
inizio settimana.

Uno studente lo paragona alla \emph{sindrome dello yes man}. Non proprio: lo
\emph{yes man} dà sempre ragione all'altro perché è avverso al conflitto e non vuole
la discussione del ``spiegami perché sei in disaccordo''. Il collega, invece,
pensava davvero di poter fare molto più di quanto potesse, e rifiutava l'idea di
non riuscirci.

\paragraph{``Ci hai pensato bene?''} Vale per i progetti e per le deadline. Con i
junior soprattutto, quando a ``vuoi fare questa cosa?'' arriva un sì troppo
convinto, il relatore chiede sempre: ``ci hai pensato bene? Me lo stai dicendo
per farmi contento? Sappi che se mi dici di no non penso peggio di te; penso
peggio di te se mi dici di sì e poi non ce la fai. Vuoi rivalutare la risposta, e
magari dirmi \emph{non lo so} invece di un sì categorico?''.

\paragraph{Che cos'è una sconfitta.} Uno studente chiede se davanti alle
``sconfitte'' non si dovrebbe rimediare. Il relatore chiarisce che per lui la
sconfitta non è il task non fatto: quella è una sconfitta \emph{a breve termine},
si troverà un modo di farlo, è un danno ``ma chi se ne frega''. La vera sconfitta,
un danno \emph{a lungo termine}, è aver pensato di riuscirci e non esserci
riusciti: se non si risolve il motivo per cui si è detto sì quando era no, lo si
ripeterà altre cinque, dieci, quindici volte. Nella testa del collega di Riot c'era
sempre una scusante (``c'erano questi imprevisti che mi hanno bloccato''), come
quella con cui ci si giustifica le interruzioni. Il relatore gli rispondeva: ``se
ogni volta c'è un imprevisto, forse c'è un \emph{pattern}: il problema non è che
non riesci a fare le cose, è che non prevedi mai gli imprevisti''. Dopo sette
ritardi su dieci dovrebbe scattare l'ingranaggio mentale (``forse sbaglio io a
dire sempre di sì''), ma non è sempre facile farlo scattare, e a volte serve una
spinta. È una delle situazioni più frequenti. E la cosa più difficile è che,
quando glielo si fa notare, queste persone lo prendono come un ``non sei capace''.
Il messaggio è invece: \textbf{``sei bravo, ma sei umano anche tu''}. Non si può
fare un numero infinito di cose in un'ora: bisogna accettare i propri limiti ``dati
dalla carne'' e non pensare di essere un cyborg.

\paragraph{Il cliente non ha sempre ragione.} Uno studente osserva che vale anche
per i project manager e per il principio ``il cliente ha sempre ragione''. Anzi:
il relatore ha lavorato con molti clienti e a volte \textbf{il cliente non sa che
cosa vuole}, che è molto diverso dall'avere sempre ragione. Soprattutto con i
clienti che chiedevano pitch creativi, la sua squadra presentava sì un'opzione
con la loro richiesta, ma anche le proprie (``se vuoi raggiungere questo, sarebbe
meglio fare così''), e molto spesso il cliente sceglieva quella proposta dal team.
Spesso il cliente ha in mente \emph{una soluzione}, ma bisogna capire \emph{quale
obiettivo} ha in mente e poi dargli una soluzione: non vuole la soluzione, vuole il
\emph{goal}, e se c'è un modo migliore per arrivarci, forse è meglio un'altra
soluzione.

\subsection{Dire di no ai meeting}

La prima cosa a cui il relatore ha imparato a dire di no sono i \textbf{meeting}.
Ha passato anni in ambienti in cui il suo calendario era letteralmente solo
meeting dalla mattina alla sera: estenuante per lui e inutile per il team. Tra i
fattori c'è il \textbf{FOMO}: tanti manager vogliono essere in tutti i meeting per
non perdersi niente ed essere sempre informati su quello che fa il team. In realtà
vanno in \emph{burnout}, passano i pomeriggi nei meeting e finiscono per essere i
più distratti, quelli che non ascoltano e che fanno domande inutili solo per far
sentire che fanno domande, senza aiutare davvero.

Prima di accettare un invito il relatore si fa una serie di domande, \textbf{come
una checklist} (e le fa anche agli altri):
\begin{enumerate}
  \item \textbf{Sono davvero necessario?} Servo davvero io, o se mancassi non
        succederebbe niente? Ai meeting ``di cortesia'' non va: se pensa di
        essere necessario lo dice lui, e il fatto che si faccia un meeting
        importante e urgente senza di lui non danneggia il suo ego. ``Non
        aggiungermi solo perché pensi che serva che io venga aggiunto.''
  \item \textbf{Uno di noi potrebbe partecipare e informare il resto del team?}
        C'è già qualcuno del suo team? Se sì, non può andare solo lui? Per un
        meeting informativo (una nuova policy, una nuova procedura da spiegare a
        tutti i dipartimenti), andarci tutti significa investire tre ore a
        testa; se ci va uno solo e poi spiega agli altri, si risparmia tempo.
  \item \textbf{Il riassunto sarà reso disponibile?} Per il relatore è
        obbligatorio. Per lo stesso motivo per cui la lezione viene
        registrata: se ci si perde un pezzo lo si può rivedere, se qualcosa non è
        chiaro lo si può risentire. Ma c'è anche un motivo di \emph{office
        politics}: se in un meeting si chiede ``mi fai questa cosa per venerdì?''
        e la risposta è sì, non si vuole sentire venerdì ``intendevo venerdì
        prossimo'' o ``non ero sicuro''. Nelle note del meeting c'è scritto che hai
        detto sì. ``Lo so che è un po' paraculo, ma in realtà è così che si
        lavora'': bisogna dare più peso alla parola data, e quindi tenerne
        traccia. Finché tutto resta orale, niente rimane scritto e tutto è
        interpretabile. Oggi è anche facile: Zoom, con l'assistente IA, produce
        in automatico trascrizione e sottotitoli, senza bisogno di qualcuno
        addetto alle \emph{meeting minutes}. In ogni caso, se in un meeting viene
        presa una decisione, il relatore vuole che sia \textbf{ripetuta per
        iscritto}: non deve ipotizzare, deve poter andare a rileggere.
  \item \textbf{Il meeting aggiunge valore rispetto al documento?} Se arriva un
        invito con allegato un documento di venti pagine, il relatore chiede: ``il
        meeting consiste nel leggermi il documento? Perché altrimenti non
        vengo''. Se non c'è un valore aggiunto, se lo legge da solo e manda
        conferma di averlo letto. È lo stesso motivo per cui si sconsigliano le
        presentazioni in cui si leggono le slide; e infatti, osserva, in quel
        momento lui stesso stava leggendo una slide, ma aggiungendo informazioni.
        ``Se tutta la lezione fosse io che vi leggo le slide, non servirei io: vi
        manderei le slide.''
  \item \textbf{Per che cos'altro potrei usare questo tempo?} È forse la più
        importante. Anche ammesso che serva, che serva proprio lui e che il
        meeting sia importante, in quel momento dovrebbe fare qualcosa di più
        importante? Se sì, fa quella. Non bisogna avere paura di dire ``no, non
        vengo'' anche se si serve, oppure ``ho un'altra cosa importante,
        spostiamo il meeting''. Il relatore si è sforzato di essere più
        \textbf{assertivo}: ``capisco la richiesta, ma ho qualcosa di più urgente
        da fare in questo momento''.
\end{enumerate}

\paragraph{E le richieste a voce?} Uno studente chiede se una richiesta fatta a
voce da un senior (``devi fare questo'') vada accettata così o fatta mandare per
mail o per messaggio. Accettarla a voce ci può stare. Ma se la richiesta è ``devi
farla in questo modo particolare e preciso'', è meglio farsi mandare il modo
preciso \textbf{per iscritto}: non si deve lasciare a interpretazione \emph{come}
fare quella cosa, per non ricadere nel ``ah, ma io non la volevo così''. Oppure si
conferma: ``sì, la faccio; va bene se la faccio così, così e così?''. Se dice sì,
si è a posto.

\paragraph{Non dovrebbero pensarci gli organizzatori?} Dalla chat: queste domande
non dovrebbero farsele già gli organizzatori dei meeting? Sì, sarebbe meglio che
tutti lo facessero, e il relatore stesso se le fa quando invita qualcuno: ``è
davvero necessario Marco? È davvero necessario Andrea?'', e solo se sì li invita.
Ma lavora in aziende anche di cinquecento dipendenti, dove non è detto che tutti
conoscano o applichino questa logica; per questo la applica per sé, sperando che
la applichino tutti. Ha sempre chiarito con tutti che da lui non arriveranno mai
inviti di cortesia o dettati dalla gerarchia (``deve esserci anche il mio capo''):
``devi esserci tu; se vuoi portarti il capo, portatelo, ma non lo invito solo
perché qui si vuole sempre il capo presente: sono sicuro che avrà di meglio da
fare''.

\paragraph{E i task?} Dire di no vale naturalmente anche per i task: si applica la
matrice di Eisenhower (sezione~\ref{sec:z2-eisenhower}) e si decide se fare,
delegare, pianificare o cancellare. Non c'è molto da aggiungere.

\subsection{Il multitasking non esiste}

``Questo potrà urtare tanti ego'', ma \textbf{il multitasking non esiste}. Si può
essere più bravi o più portati di altri a fare più cose contemporaneamente; la
realtà è che, concentrandosi su più cose, \textbf{la somma non farà mai cento}. Se
si lavora su due cose non si lavora al 50\% su una e al 50\% sull'altra, ma al 40\%
e al 40\%, e il 20\% va perso in ``costo di manutenzione'', cioè nella capacità di
seguirle entrambe. Il consiglio è dedicarsi al cento per cento a un task, poi
passare al successivo, poi al successivo, invece di farne cinque insieme: si può
avere la percezione di andare più veloci, ma non è così.

Uno studente cita una lettura di neuroscienze secondo cui il cervello è in grado
di fare una sola operazione alla volta. Il relatore conferma: smettiamola di
parlarne come di un pregio o di una \emph{skill}. È un modo che alcuni cervelli
trovano per svagarsi, l'equivalente di accendere un podcast: fare due cose
insieme vuol dire cercare di distrarsi alternandole, e farle entrambe peggio di
quanto si farebbero una alla volta. A chi preferisce fare tante cose insieme
suggerisce di provare a farne una alla volta: le farà più velocemente. Anche il
semplice passaggio da una cosa all'altra non è istantaneo: ci sono sempre una
decina di secondi in cui il cervello si chiede ``che cosa stavo facendo?''; per
cinquanta cambi, già questo spiega la perdita di velocità. È una questione di
\textbf{bandwidth mentale}, un concetto che tornerà nel pomeriggio: anche lo
\emph{switch} occupa banda.

\subsection{La fine della giornata}

A fine giornata si \textbf{rimette a fuoco la giornata}: si riguardano la task list
e la \emph{must do list} e si conferma di aver fatto tutto. È letteralmente il
controllo prima di uscire di casa: chiavi, porta chiusa, cellulare, portafoglio.
Al relatore è capitato di dimenticare un \emph{must do} rimasto in lista per tutto
il giorno, ed entra in gioco quella che chiama la \textbf{sindrome del
tirannosauro}\footnote{Il riferimento è alla celebre scena di \emph{Jurassic Park},
in cui si dice che il tirannosauro vede la preda solo se si muove.}: se una cosa è
lì, non è detto che la si veda finché non si muove. Averla in lista tutto il giorno
non vuol dire averla percepita come da fare.

Può anche capitare di non farcela perché quel giorno sono arrivate così tante cose
urgenti e importanti da rendere inutili le stime: la giornata no capita e non c'è
nessun problema. Ma se era davvero un \emph{must do}, \textbf{va fatto comunque}:
è lì che scattano la richiesta di straordinario o il ``devo farlo comunque''.
Attenzione: se si chiude la giornata con un \emph{must do} ancora aperto e si
decide ``vabbè, lo faccio domani'', può starci, ma allora \textbf{non doveva essere
in quella lista}, perché era rimandabile.

Al relatore questo controllo serve anche ``per quieto vivere'': se chiude ogni
giornata con la \emph{must do list} completa, sa che magari non ha fatto tutto
quello che c'era nella to-do list, ma i \emph{must do} sì. ``Il mio, oggi, l'ho
fatto'': può stare tranquillo senza continuare a chiedersi se ha fatto quello che
doveva.

% ---------------------------------------------------------------------
\section{Domande sul time management}
\label{sec:z2-qa-tempo}

\paragraph{C'è un numero medio ``sano'' di \emph{must do} quotidiani?} No, perché un
conto è avere dieci \emph{must do} come ``rispondere a dieci mail'', un conto è
averne uno solo come ``inventarmi un nuovo gioco''. Dipende dalla dimensione: se ne
gestiscono di più se sono minuscoli, di meno se sono giganteschi. Ci sono lavori
che per definizione hanno tantissimi \emph{must do} quotidiani e lavori che ne hanno
uno al giorno. Si imparerà da sé qual è il numero giusto.

\paragraph{C'è un modo per stimare quanto si riesce a fare in una giornata?} ``Lo so
che non è una bella risposta'', ma l'unico modo è la propria \textbf{storia}:
sapere quanto si riesce a fare nel quotidiano. È normale e fisiologico che le
prime stime siano sbagliate, e probabilmente eccessivamente ottimistiche;
l'importante è rendersene conto. È letteralmente la \emph{seniority}: abituarsi a
fare qualcosa insegna quanto ci si mette. Va benissimo metterci di più all'inizio,
ed è normale che il tempo cambi man mano che si migliora: ciò che oggi si stima in
un'ora, tra sei mesi lo si stimerà in mezz'ora. Chi ha lavorato con Scrum lo sa:
su un progetto nuovo tutti sbagliano la dimensione delle \emph{stories} all'inizio,
perché nessuno sa quanto ci vorrà, e al terzo \emph{sprint} consecutivo è più
chiaro. Le stime sono stime, altrimenti si chiamerebbero ``inserire quanto tempo ci
metterai''.

Se però si dà una stima e poi il lavoro si rivela il triplo, il problema non è la
stima: è che \textbf{bisognava chiedere più informazioni} prima di darla. La
risposta giusta non è la stima, ma ``per darti una stima accurata mi serve sapere
questo, questo e questo''. Il problema non è sbagliare la stima, è non darsi i mezzi
per azzeccarla.

\paragraph{Come si fa a far entrare questi metodi nella routine? È solo
determinazione?} È un misto. Il relatore risponde da persona con ADHD, con la
tendenza a distrarsi e a fare un mucchio di \emph{overthinking} (``come dicevo ieri,
anche ordinare la pizza è overthinking''). È un po' disciplina, ma soprattutto far
diventare qualcosa un'\textbf{abitudine}: lo sforzo per fare le checklist era molto
più alto un anno prima, oggi è un'abitudine. Di solito servono \textbf{circa tre
mesi} perché qualcosa diventi un'abitudine\footnote{Stima come riportata in lezione.
Gli studi più citati sulla formazione delle abitudini indicano tempi molto variabili
da persona a persona, con una mediana intorno ai due mesi.}: la teoria dice che chi
inizia oggi ad andare in palestra ogni giorno, dopo tre mesi non penserà più ``oggi
devo andare in palestra'', ma la palestra farà parte della sua giornata. Non è più
uno sforzo, ma ``è così che lavoro''.

\paragraph{Chi fa davvero time management?} Prima di passare al brainstorming il
relatore torna alla domanda iniziale: qualcuno, ora che sa che cosa intendeva,
cambierebbe risposta? Uno studente risponde che lo fa, ma senza checklist, perché
non ne sente il bisogno (non gli capita spesso di dimenticare le cose); la
\emph{must do list} invece la segue sempre, e le cose meno importanti capitano nelle
giornate in cui ne salta qualcuna. Va benissimo: non esiste un unico modo, ognuno
troverà il proprio. Il relatore porta il suo esempio perché sa che è servito a molte
persone con cui ha lavorato; con chi faceva fatica ad applicare un metodo (come
chiedeva prima uno studente) ne hanno trovati altri, perché ciò che funziona per uno
non è detto che funzioni per un altro. L'importante è che \textbf{ognuno trovi il
proprio metodo}, perché la soluzione non è non fare time management: ci si sentirà
più ``liberi'' improvvisando la giornata, ma non la si sta sfruttando a dovere.

E non si parla dei video sul \emph{mindset} in cui ci si alza alle cinque, si fa la
checklist e si fa colazione leggendo le mail: organizzarsi meglio la giornata
\textbf{dà più tempo libero}. E per il relatore il tempo libero delle persone è
fondamentale, perché è ciò che permette loro di lavorare meglio.

% ---------------------------------------------------------------------
\section{Una parentesi: tempo libero, ferie e benessere}
\label{sec:z2-benessere}

\subsection{Le tre condizioni del relatore}

Il relatore è un sostenitore della \textbf{settimana lavorativa di quattro giorni},
che ha applicato in varie aziende; dove non ci è riuscito, ha introdotto almeno
\textbf{venerdì liberi}, due volte al mese o una volta al mese, cioè giornate in
cui non si fa niente. Come le pause brevi sono importanti per essere creativi,
anche le pause lunghe, cioè giornate intere, lo sono per essere creativi ed
efficienti. È lo stesso motivo per cui, dopo la Games Week, dava una settimana
libera alle persone (sezione~\ref{sec:z1-fasi-team}): ha bisogno che siano
riposate, e quella settimana ``persa'' è in realtà un investimento per le
settimane successive. In chat qualcuno lo paragona al ``metodo finlandese'': è un
misto tra quello e il metodo del buonsenso.

Il relatore viene spesso chiamato a \textbf{cambiare le culture aziendali}, ed è un
po' il suo focus. Come un allenatore che accetta una squadra solo se può portare il
suo staff (il preparatore, il tecnico), lui chiede tre elementi:
\begin{enumerate}
  \item possibilmente \textbf{quattro giorni lavorativi}, o altrimenti giornate
        libere per dare spazio alle persone;
  \item \textbf{ferie illimitate}: nessun limite massimo, ma un \textbf{limite
        minimo garantito}. Ferie illimitate non devono voler dire poter scegliere
        di non prenderle: ``devi prendere almeno X giorni, fino a infinito'';
  \item la \textbf{feedback culture}.
\end{enumerate}
Senza questi tre elementi non garantisce il risultato: è un modo di gestire le
aspettative fin dall'inizio, come si diceva il giorno prima.

\paragraph{La feedback culture.} Il perché è stato spiegato nella prima giornata, ma
il relatore lo sottolinea: un ambiente in cui ci si può fidare del team vale
moltissimo. Quando va a fare lezione allo IULM, anche in settimana, non deve
preoccuparsi di che cosa stia facendo il team mentre lui non lavora (``oggi è
domenica e non ho questa preoccupazione a prescindere, ma anche se fosse mercoledì
non l'avrei''): sa che stanno facendo del loro meglio, e non gli interessa sapere
altro.

\paragraph{Ferie e permessi illimitati.} Nell'esperienza del relatore, sia a Riot sia
a Blizzard, avere ferie illimitate significava spesso \emph{prenderne meno} di quelle
obbligatorie: c'è una vocina che dice ``ti serve davvero un giorno di ferie oggi?
Tanto sono illimitate, te le prenderai quando ti serviranno'', e sapere di averne
illimitate spingeva a lavorare di più. Per lui era altrettanto dannoso, e per questo
ha introdotto un \textbf{minimo obbligatorio di 24 giorni} di ferie.

C'è stato però un effetto non previsto, e molto desiderato: in tutte le aziende in
cui ha introdotto ferie e permessi illimitati, \textbf{i giorni di malattia sono
calati del 70\%}\footnote{Dato riferito dal relatore sulla base della propria
esperienza, non verificabile.}. La gente si ammala di meno? O piuttosto il giorno in
cui ci si sveglia e si pensa ``oggi non ne ho voglia, mi do malato'' diventa il
giorno in cui si prende un permesso? Si smette di usare come malattia un giorno che
non era di malattia.

\paragraph{L'\emph{Hangover Day}.} Sia a Riot sia a PG Esports il relatore aveva
introdotto quello che a Riot, ``essendo Riot a Dublino'', si chiamava
\textbf{Hangover Day}\footnote{Nella trascrizione ``Angover Day''.}: una giornata
all'anno in cui si aveva il diritto di telefonare alle nove del mattino e dire ``oggi
non vengo perché non ne ho voglia'' (tradotto: perché ieri sera ho bevuto troppo).
Altrimenti quella giornata sarebbe diventata un giorno di malattia o di ferie. Il
fatto è che \textbf{quasi nessuno lo usava}: il relatore stesso a volte ci ha
pensato, ma poi la mattina si chiedeva ``voglio davvero giocarmelo oggi? Tutto
sommato potrei anche andare'', e ci andava. Ma anche solo sapere di averlo lo faceva
stare tranquillo. Allo stesso modo, con le ferie illimitate, chi voleva prendersi un
mese lo faceva (``buon per loro''), chi non lo faceva andava bene lo stesso:
l'importante era che staccassero. C'era un \textbf{benessere più diffuso} nel team,
perché sapevano che, volendo, avrebbero potuto.

\paragraph{Il \emph{Neko Day}.} Il primo punto, i giorni liberi, era chiamato
\textbf{Neko Day}, da \emph{neko}, ``gatto'' in giapponese: il giorno in cui si può
fare il gatto, rannicchiati davanti al camino a non fare niente. Perché
introdurlo? Chi lavora dal lunedì al venerdì passa di solito il sabato o la domenica
a fare le commissioni che non può fare nei giorni lavorativi: la spesa, le bollette,
il bollo dell'auto. In questo modo sfrutta meno il weekend. Con un venerdì libero al
mese (o due, o tre) le persone facevano quel giorno le commissioni del weekend, e
sabato e domenica si riposavano davvero (``anche non fare niente andava bene,
chiaramente''). Il lunedì arrivavano molto più cariche e riposate. Questo ha aiutato
tantissimo a mantenere un team sano, fisicamente e mentalmente.

\paragraph{Il supporto psicologico.} A PG Esports il relatore ha introdotto anche il
\textbf{supporto psicoterapeutico offerto dall'azienda}: sei sedute da un'ora con
uno psicoterapeuta, online. Era il periodo in cui si parlava del bonus psicologo
statale, ``che poi non è mai arrivato''\footnote{Così in lezione. In realtà il bonus
psicologo è stato istituito in Italia nel 2022, ma con fondi limitati rispetto alle
richieste e con forti ritardi nell'erogazione; il relatore si riferisce
probabilmente a questo.}, e lui ha pensato: l'idea è buona, perché non introdurla a
livello aziendale? Erano gli anni dei lockdown: l'azienda stava perdendo molte persone
che andavano in depressione o non avevano più voglia di lavorare, perché avevano
difficoltà a vivere quel periodo (``io stesso non è che l'abbia vissuto benissimo'').
Qualcuno ne ha usufruito e gli ha detto: ``mi stavo per dimettere, poi ho risolto le
cose e ho preferito restare''.

Secondo il relatore oggi c'è una consapevolezza maggiore da parte delle aziende
sulla \textbf{salute mentale}: vuol dire dare più ferie, dare questo tipo di
supporto, dare più \emph{benefit} (anche i buoni pasto). Il benessere non è solo
economico, è anche mentale, ed è una cosa positiva da sapere.

% ---------------------------------------------------------------------
\section{Brainstorming: obiettivo e problemi tipici}
\label{sec:z2-brainstorming}

\subsection{Perché il brainstorming riguarda tutti}

Secondo il relatore questa è la parte che servirà di più, sia nei gruppi e nei
progetti del percorso di studi (il gioco da tavolo, il gioco digitale) sia nelle
aziende. Rispetto al time management ha una pertinenza più immediata: tutti
hanno esempi di \textbf{brainstorming andati male}, da cui si è usciti esausti, con
la sensazione di non aver ottenuto niente, o di non aver parlato abbastanza (``ha
parlato solo quello''), o di non essere sicuri della scelta. Il brainstorming lascia
una sensazione molto chiara, a fine meeting, che qualcosa non torna; il time
management dà piuttosto la tristezza del ``non ce l'ho fatta a fare quello che
dovevo oggi, come mai?''.

\subsection{Generare idee, non decidere}

Il primo chiarimento riguarda \textbf{l'obiettivo}, che il relatore vede spesso
travisato. L'obiettivo del brainstorming è \textbf{generare idee}; invece diventa
spesso decisionale, con lo scopo di scegliere che cosa fare. Ma allora non sarebbe
un brainstorming, sarebbe un \emph{decision making meeting}. Generare molte idee e
scegliere la migliore nello stesso meeting non funziona né per generare idee né per
scegliere la migliore, perché sono \textbf{due metodi che si contrappongono}. Il
brainstorming va fatto quando si devono generare idee \emph{per poi} scegliere che
cosa fare: deve venire \textbf{prima} del \emph{decision making}, come primo passo
della fase di esplorazione di un nuovo progetto.

\subsection{``Non esistono cattive idee''?}

Bisogna evitare che i membri del team \textbf{autofiltrino} le proprie idee per paura
di essere rifiutati o presi alla leggera. Da qui nasce la frase ``non esistono
cattive idee, buttate fuori tutto'', che è vera solo fino a un certo punto. Le
cattive idee esistono, oggettivamente; quello che non si deve fare è
\emph{colpevolizzarle}. Né si deve per forza perseguire ogni opzione: parlare di
opzioni che sono effettivamente brutte è uno spreco di tempo, e il brainstorming ha
anche un problema di gestione del tempo. Il relatore vede spesso brainstorming
gonfiati a ore e ore in cui si parla anche di cose che non meritano nemmeno dieci
minuti, come i feedback buoni ma inutili del giorno prima. \textbf{Trattare tutte le
idee allo stesso modo è sbagliato}, perché non hanno la stessa qualità. Bisogna
trovare un modo di valutarle senza che chi le ha proposte si senta preso di mira,
bullizzato o sminuito.

Per questo a ogni brainstorming dovrebbe seguire un \textbf{follow-up}, il momento in
cui si decide che cosa fare: le due parti vanno separate, altrimenti ``si crea solo
casino''. Il vero lavoro arriva quando le idee vengono messe a terra: c'è la parte di
brainstorming e di \emph{theorycrafting}, e la parte esecutiva viene dopo.

\subsection{I problemi più comuni}

\begin{description}
  \item[Le prime idee monopolizzano la sessione] Le prime due idee ``sparate a
        caso'' sono quelle su cui ci si concentra per la maggior parte del tempo,
        sviscerandone i dettagli (``intendevi così o così?''), per poi chiudere con
        ``altre idee? Questa, questa e questa? Sì, no, ok, andiamo avanti''. Si
        parla troppo delle prime idee e poco delle ultime.
  \item[Poche persone parlano per la maggior parte del tempo] È un effetto
        collegato al precedente.
  \item[\emph{Anchoring}] Ci si lega alle prime idee, e spesso si votano o si
        scelgono non perché siano le migliori, ma perché sembra di conoscerle meglio
        di quelle menzionate in trenta secondi alla fine. Non c'è nessun motivo per
        pensare che un'idea detta alla fine sia peggiore: era solo arrivata dopo.
        C'è un problema di \emph{timing}. Anzi, un'idea venuta istantaneamente è
        forse quella meno ragionata, perché ci si è pensato meno tempo.
  \item[Silenzi imbarazzanti] Spesso le persone arrivano al meeting senza nemmeno
        sapere che è un brainstorming. Quando il relatore convoca un brainstorming,
        comunica prima a tutti il \emph{topic}, il tipo di idee cercate, i paletti e
        le restrizioni, così che anche chi è meno portato all'improvvisazione arrivi
        con qualche idea pronta. Se la gente arriva impreparata rischia di non dire
        niente per un'ora, e allora perché è stata invitata? Non è un difetto loro:
        tanti manager dicevano al relatore ``ci sono persone inadatte ai
        brainstorming perché non sanno improvvisare''. No: si sta facendo
        brainstorming con un metodo che premia solo chi sa improvvisare ed è
        abbastanza estroverso da parlare sul momento. Non significa che chi non
        parla non avrebbe potuto contribuire, ma che non gli è stato dato modo di
        farlo.
  \item[Estroversi e pensatori rapidi dominano] Le personalità estroverse e chi
        pensa velocemente tendono a dominare la conversazione, più o meno
        intenzionalmente, senza lasciare agli altri il tempo di contribuire. Non c'è
        uno spazio in cui un introverso possa dire ``veramente io volevo\dots'':
        ``no, no, stiamo parlando di questa bellissima idea, facciamo questa''.
\end{description}

% ---------------------------------------------------------------------
\section{Un metodo di brainstorming}
\label{sec:z2-metodo}

Il relatore propone una soluzione principale, declinata in vari modi, che tocca
tutti questi punti.

\subsection{Preparazione e ambiente inclusivo}

\paragraph{Arrivare con un obiettivo chiaro.} Bisogna arrivare al brainstorming con
chiaro e preciso che cosa si sta cercando, cioè \textbf{che cosa si sta cercando di
risolvere}, \textbf{perché} e soprattutto \textbf{con quali vincoli}. Si stanno solo
raccogliendo idee per un cliente? Si cercano idee fattibili, o idee innovative? Che
cosa si vuole produrre da questo meeting? Deve essere chiaro a tutti \emph{prima}
del meeting, non spiegato durante, perché altrimenti si ricade nel problema delle
persone non pronte a intervenire.

\paragraph{Non giudicare le idee degli altri.} Al relatore è capitato spessissimo di
sentire un'idea terribile da una persona e, mezz'ora dopo, l'idea perfetta dalla
stessa persona. Ma se all'idea terribile tutti rispondono ``che stupidaggine, per
favore stai zitto'', quella persona non dirà la successiva, perché non vuole
rivivere l'esperienza. Giudicando le idee degli altri, di fatto, li si zittisce.
Le idee migliori nascono quando si propone quella che si pensa sia l'idea perfetta,
magari non è buona e non viene scelta (``fa niente''), si ascoltano anche le altre e
poi la si \textbf{riformula} integrando i pareri sentiti nel frattempo. Sono queste
\textbf{idee combinate} a essere spesso le soluzioni migliori.

\subsection{Separare generazione e discussione: \emph{brainwriting} e
\emph{rapid ideation}}

Con il classico ``non ci sono cattive idee, sparate'', succede spesso che uno prenda
la parola, spieghi la sua idea e il suo concetto, e gli altri comincino subito a
fargli domande (``ma quindi la vedi così?''): si sta unendo la persona che
\emph{genera} l'idea con la persona che la \emph{discute}, nello stesso momento.
Invece bisogna \textbf{dividere la generazione dalla discussione}: un momento in cui
si generano idee e un momento in cui si discutono. È il principio del
\textbf{brainwriting}. Funziona molto quando nel team c'è una figura con una
posizione gerarchica più alta: se quella persona esprime un'idea e i due momenti
sono uniti, tutti tenderanno ad allinearsi. E si vuole che \emph{tutti} possano
avere idee.

Lo strumento è la \textbf{rapid ideation}. Il relatore di solito dice: ``per i primi
cinque minuti non parliamo; scrivete su dei post-it il pitch delle vostre idee''. Se
ne possono scrivere quanti se ne vuole, dieci o uno, e si mettono da parte. Ma il
tempo è di \textbf{cinque minuti}: non dieci, non due, e soprattutto non mezz'ora.
Il \textbf{limite di tempo breve} è importantissimo, perché costringe,
consapevolmente o meno, a filtrare le idee mentalmente e a scrivere solo quelle che
si ritengono migliori: magari non è chiaro perché si preferisca un'idea all'altra,
ma il cervello sta facendo il lavoro in background e fa venire in mente le cose più
importanti o più giuste. È ottimo anche per i team che tendono a distrarsi, perché
sono solo cinque minuti.

\subsection{\emph{Dot voting}}

Poi si passa al \textbf{dot voting}:
\begin{enumerate}
  \item si attaccano \textbf{tutte} le idee su una lavagna fisica, un muro o una
        lavagna digitale (come Miro\footnote{Nella trascrizione ``mirror board'':
        si tratta della lavagna digitale Miro, già citata nella prima
        giornata.});
  \item si \textbf{uniscono le idee che hanno lo stesso concetto di base}. Se si sta
        pensando a un evento dal vivo e la proposta del relatore è un parco di
        dinosauri, mentre quella di un collega è un parco di dinosauri ``del
        paleolitico'' con i tirannosauri, le si raggruppa sotto ``parco di
        dinosauri'': non sono due idee separate. Un altro gruppo potrebbe essere,
        ad esempio, un parco a tema Avengers;
  \item ognuno ha \textbf{da tre a cinque voti} e va alla lavagna a mettere i propri
        pallini sulle idee che gli piacciono e di cui vorrebbe sapere di più.
\end{enumerate}
Il risultato tipico: un'idea con molti voti, una con un paio, altre con un numero
discreto.

\paragraph{Che cosa si vede subito?} Il relatore lo chiede alla classe. Si vede una
\textbf{preferenza} per un tema o per un'idea. Non c'è \textbf{polarizzazione} su una
o due idee iniziali, perché c'è varietà. Il voto è \textbf{anonimo}: non ci sono nomi.
E soprattutto manca un'altra cosa importantissima: \textbf{il giudizio negativo
espresso}. Se c'è un'idea con sette voti, una con cinque e una con tre, è
ragionevole parlare prima di quelle, perché tutti hanno espresso interesse e vogliono
saperne di più; ma nessuno sta dicendo ``non voglio assolutamente parlare di questa
perché mi fa schifo''. È semplicemente oggettivo che quell'idea non ha ricevuto voti.

\paragraph{Autovotarsi.} Il sistema permette anche di votare la propria idea, e non c'è
nessun problema: se la propria idea ha un pallino e gli altri ne mettono sette su
un'altra, il proprio voto è ininfluente. Il relatore vedeva spesso che all'inizio la
gente tendeva ad autovotarsi, poi capiva che era inutile, perché non sarebbe mai stato
il voto decisivo, e tanto valeva votare le idee che interessavano davvero. La
tendenza spariva con il tempo.

\paragraph{Che cosa si è ottenuto.} Tre ore di meeting sono state filtrate in una
discussione concentrata solo sulle idee che tutti ritengono interessanti, con un
\emph{agreement} di fondo: non si parla dell'idea numero uno perché è uscita per prima,
non è una questione di \emph{timing} ma di qualità. Il relatore ha vissuto tanti
meeting con la frustrazione interna del ``stiamo sprecando tempo su questa idea,
quando vorrei dire la mia che secondo me è meglio''. Qui non si può sprecare tempo,
perché si parla di qualcosa che tutti vogliono sentire.

\paragraph{E se si formano sottogruppi che si votano a vicenda?} Uno studente chiede se
non possano esistere ``cordate''. Anche se esistessero, significherebbe che il gruppo di
quel meeting vuole parlare di quella cosa, e va bene: l'importante è non sprecare tempo
su idee che non interessano a nessuno. Se si investono due ore, quarantacinque minuti
su ciascuna di tre idee vanno benissimo; dieci minuti per ogni idea no. E comunque il
voto è anonimo: non si può sapere che un'idea è di Gianni per dare supporto a Gianni,
quindi un eventuale sostegno reciproco sarebbe involontario.

\paragraph{E chi è lento a elaborare?} Uno studente osserva che il problema degli
introversi potrebbe ripresentarsi con i post-it: con un tempo così ristretto, chi è più
lento potrebbe andare ancora più in panico. Per questo bisogna \textbf{arrivare
preparati}: si conosce il topic in anticipo, quindi almeno un paio di idee già
formalizzate si hanno in mente. Magari non se ne scriveranno quindici, ma un paio sì.

\subsection{Quanti voti e quali varianti}

Dare più pallini a testa sarebbe controproducente? Il relatore indicava da tre a
cinque voti. Chiede alla classe quanti siano i gruppi di lavoro: dai quattro agli otto,
alcuni gruppi del secondo anno sono nove. La sua prassi:
\begin{itemize}
  \item \textbf{dai cinque partecipanti in su}, \textbf{tre voti} a testa, più che
        sufficienti per popolare molto la lavagna;
  \item \textbf{fino a cinque partecipanti}, \textbf{cinque voti} a testa, che
        permettono anche di variare un po'.
\end{itemize}
Esistono varianti: ad esempio tre voti di colori diversi, il blu per l'idea più
innovativa, il verde per quella più bella, il giallo per quella più facile, così da
qualificare \emph{perché} un'idea piace. Se l'obiettivo è proporre a un cliente tre
pitch (uno facile, uno medio, uno avanzato), la votazione è già fatta. All'inizio,
però, meglio partire da tre o cinque voti generici; con il tempo si imparerà ad
accelerare.

\subsection{Il moderatore e gli invitati}

\paragraph{Serve un moderatore?} Può essere utile. Il relatore fa spessissimo il
moderatore: partecipa ai brainstorming senza fare altro che moderare, ``perché so che
le persone con le idee fighe sono loro e non io''. Ci si può anche autogestire. Nei team
piccoli non ha senso ``perdere'' una persona per la moderazione; nei team più grandi
(dai cinque in su, sette-otto persone), se qualcuno dice ``non penso di avere idee
particolari su questo, vi modero così vi aiuto'', è un ottimo modo di contribuire.

\paragraph{Chi invitare?} Dipende da che cosa si cerca. Per inventare un nuovo gioco il
relatore invita i game director, gli art director e magari il \emph{publishing manager}
per un input commerciale; per decidere il \emph{combat system} invita magari il senior
game designer e il game director. La domanda che si fa chi convoca il meeting è: dato
l'obiettivo, chi è più ragionevole che possa contribuire?

Un 3D artist della classe chiede se allora gli artisti 3D non vengano inclusi. Dipende:
se l'obiettivo è definire la linea creativa dello stile dei personaggi, il relatore
invita anche il 3D artist, e non solo i \emph{concept artist}. Il 3D artist può dare un
input diverso dal concept: vota un'idea perché per lui è \emph{più fattibile}. Un
concept può essere bellissimo ma ``un inferno'' da modellare, e chi sa di doverlo fare
entro giugno non vota un'idea che sa già infattibile.

In realtà si può invitare chiunque: per un pitch su un'area dedicata al basket alla
Games Week, il relatore ha invitato un tecnico di regia che gioca a basket, per avere
anche la prospettiva di chi ha familiarità con il tema. A uno studente che porta
l'esempio di un elemento di interfaccia sul personaggio, il relatore risponde che,
volendo, si trova un motivo per invitare chiunque: l'importante è avere chiaro
\textbf{chi è necessario}. Il gruppo di lavoro può essere di otto e il brainstorming di
dodici, con input esterni a chi poi farà il lavoro. Basta non esagerare: se la maggior
parte dei voti arriva da chi non farà quel lavoro, si vota su qualcosa che poi non
faranno loro, e non ha più senso.

% ---------------------------------------------------------------------
\section{Dalla votazione alla decisione: la matrice impatto--sforzo}
\label{sec:z2-impatto}

\subsection{Massimo impatto, minimo sforzo}

Fatto questo, non si è finito: ora si sa quali idee interessano di più, ma non si è
ancora deciso che cosa fare. Serve la parte decisionale. L'obiettivo di qualsiasi
azienda (ed è il motivo per cui al relatore non importa se il programmatore fa il test
da casa cercando su Google) è \textbf{massimizzare l'impatto e minimizzare lo sforzo}.
Vale per i progetti, i task, il recruitment, per qualsiasi cosa: ``non perché sono
pigro, ma perché sarebbe stupido fare diversamente''.

Le idee selezionate si mettono quindi in un'altra matrice, con due assi:
\begin{itemize}
  \item sull'asse verticale il \textbf{valore}, cioè quanto \textbf{impatto} avrebbe
        l'idea (in alto l'impatto alto);
  \item sull'asse orizzontale lo \textbf{sforzo} (\emph{effort}), da basso (a sinistra)
        ad alto (a destra).
\end{itemize}
In alto a sinistra ci sono quindi le idee ad alto impatto e basso sforzo, in alto a
destra quelle ad alto impatto e alto sforzo, e così via.

\paragraph{L'esempio Adidas.} Se il cliente è Adidas, con la linea di abbigliamento per
il basket, un'idea a basso valore potrebbe essere una campagna su TikTok in cui gli
utenti condividono i loro momenti con un pallone da basket; un'idea ad alto valore
potrebbe essere fondare un parco divertimenti dove i giocatori NBA vengono a sfidarti
uno contro uno. Questo non dice che una sia meglio dell'altra: dice che l'impatto è
diverso. Allo stesso modo si valuta lo sforzo: la campagna social è a basso sforzo, il
parco divertimenti ad alto sforzo per molti motivi (i permessi di costruzione, i
giocatori NBA).

\subsection{Posizionare le idee, insieme e con precisione}

Il gruppo che a quel punto è \textbf{il gruppo decisionale}, cioè chi farà il lavoro,
deve posizionare sulla matrice le idee ritenute migliori, e deve essere \textbf{d'accordo
tutto il gruppo}: ``siamo tutti d'accordo che il parco dei dinosauri di Stefano e Andrea è
ad alto valore ma anche ad alto sforzo? Sì? Allora lo mettiamo qui''. E bisogna indicare
\textbf{un punto preciso}, non solo il quadrante: piazzare un'idea in un punto non è come
piazzarla poco più in là. Il risultato è una sorta di \textbf{mosaico} delle idee sul
grafico.

\subsection{Leggere la matrice in diagonale}

Davanti a una situazione tipica, con varie idee sparse, il relatore chiede alla classe:
se foste i project manager, se doveste fare voi queste cose, \emph{quali} fareste? Al
plurale, perché spesso da un brainstorming si esce con più idee realizzabili e si deve
scegliere su quali concentrarsi. Le risposte variano: c'è chi sceglie le due con più
impatto, chi aggiunge che dipende da tempi, budget e risorse, chi preferisce quella più
innovativa ma anche più rischiosa.

Il metodo che secondo il relatore va applicato è \textbf{procedere in diagonale}, dalle
idee a basso sforzo e alto impatto verso quelle ad alto sforzo e basso
impatto\footnote{Nella slide le idee erano rappresentate da post-it colorati (rossa,
rosa o salmone, azzurra, verde acqua, viola, blu, gialla), e le indicazioni del
relatore si riferiscono a quei colori. Dalla sola trascrizione non è possibile
ricostruire con certezza la posizione di ciascun post-it, e alcune indicazioni sui
colori risultano contraddittorie; qui è riportato il principio generale, che il
relatore chiarisce nelle risposte alle domande.}:
\begin{enumerate}
  \item un'idea a \textbf{basso sforzo e alto impatto} (nella slide, quella rossa) va
        fatta \textbf{per prima, sempre}: ``non dovrebbe essere nemmeno in discussione'';
  \item poi, a seconda delle risorse e del tempo, si possono fare
        \textbf{contemporaneamente} le idee che stanno sulla stessa diagonale
        successiva, che hanno cioè un rapporto simile tra impatto e sforzo;
  \item poi le successive;
  \item le idee ad \textbf{alto sforzo e basso impatto} (nella slide, la gialla)
        \textbf{non vanno nemmeno considerate}. Nessuno, infatti, l'aveva scelta.
\end{enumerate}
Per le scelte intermedie valgono due criteri, chiariti rispondendo a uno studente che
non capiva perché alcune idee fossero preferibili ad altre: \textbf{a parità di sforzo si
sceglie l'idea con più impatto}; \textbf{a parità di impatto si sceglie quella con meno
sforzo}. Con risorse infinite si farebbero anche le idee più costose; ma di solito le
risorse sono limitate, e allora tra due idee con lo stesso sforzo e un impatto molto
diverso la scelta è evidente. Se ci piacerebbe fare un'idea più ambiziosa ma abbiamo
risorse solo per una meno costosa sulla stessa diagonale, va bene lo stesso: abbiamo
comunque stabilito che è meglio delle idee sulle diagonali successive. Quanto fare dipende
dalla \textbf{bandwidth} a disposizione.

\paragraph{``La gialla era la mia idea: devo offendermi?''} No. Qualcuno ha detto che
l'autore fa schifo? No. Si è concordato che la sua idea, per quanto votata (``qualcuno ha
votato per la mia idea''), è piuttosto difficile da realizzare rispetto al suo impatto.
Non si sta dicendo che non sa più fare il suo lavoro. Anzi, se lui stesso dovesse
scegliere tra la propria idea gialla e la rossa di un altro, sceglierebbe la rossa: non è
nemmeno una questione di discussione. La matrice rende \textbf{visivamente evidente} che
cosa fare, e permette di aggirare molte discussioni che non sono più soggettive:
oggettivamente, dovremmo fare così.

\subsection{Il ruolo del moderatore e dell'esperto}

È qui che lavora il moderatore: ``non è più una discussione; visto il posizionamento,
dobbiamo fare queste cose''. Non si decide a maggioranza dove mettere le idee: se uno
posiziona un'idea e un altro non è d'accordo, difficilmente dirà ``la rossa dovrebbe stare
al posto della gialla''; dirà ``secondo me è un po' più a destra, un po' più in basso'', e
si fanno gli aggiustamenti. Se invece nascono discussioni del tipo ``la rossa dovrebbe
stare al posto della gialla'', vuol dire che non si è definita bene l'idea e non è chiaro
quanto sia davvero difficile da realizzare.

Oppure interviene \textbf{l'esperto}: è l'esempio del concept artist e del 3D artist. Il
concept artist dice che un'idea è a basso sforzo e alto impatto; il 3D artist risponde ``ne
condivido l'impatto, ma dovendola fare io ti assicuro che lo sforzo non è basso, è molto
alto''. Allora è un altro discorso, perché l'esperto sta dicendo che lo sforzo è diverso
da quello stimato. Ma alla fine tutto deve essere concordato.

\subsection{Chi decide}

\paragraph{Il gruppo decisionale può essere diverso da quello del brainstorming?} Sì: a
questo punto escono quelli invitati come ``consulenti'', che hanno dato idee ma non devono
poi realizzarle. Farli votare anche qui significherebbe perdere il collegamento tra chi
deve fare le cose e che cosa si decide di fare. In quel momento non hanno responsabilità:
hanno dato un input, e se decidessero che cosa devono fare gli altri dovrebbero esserne
i manager. Non dovrebbero poter dire agli altri che cosa fare se non ne hanno la
responsabilità. Vale per tutti i team, piccoli e grandi; secondo il relatore chi decide
dovrebbe essere \textbf{al massimo una decina di persone}. In un team di dieci che fa il
brainstorming tutto insieme, va benissimo che chi vota coincida con chi ha partecipato:
l'importante è che \textbf{chi vota sia chi farà il lavoro}.

\paragraph{Lo sforzo è il costo o la fatica?} Uno studente chiede se l'alto impatto
includa anche il costo. No, il costo sta nello sforzo, che per il relatore va
\textbf{ragionato in risorse}: denaro, tempo e così via. Un'idea può avere un basso costo
monetario ma un alto costo in tempo, e allora va bilanciata e si troverà a metà. È lì che
si spostano le idee e che avviene la conversazione: i tecnici dicono che lo sforzo tecnico
è un certo, poi arriva l'\emph{accountant} che dice che a livello economico è un altro, e
la si sposta.

\paragraph{E il piacere di fare il progetto?} In un gruppo di due o tre persone, osserva
uno studente, si finisce per votare anche per preferenza personale: quando si fa un gioco
bisogna avere anche il piacere di farlo. Ma \textbf{anche quello è impatto}: l'impatto
sulla motivazione. Come lo sforzo sono le risorse, l'impatto è l'impatto \emph{su di te,
sul team, sul mercato}, su ciò che si vuole. Il relatore lo ha tenuto volutamente
generico: può essere quanti soldi si fanno o quanto è utile per il proprio CV, e
l'impatto cambia tra un progetto accademico e uno professionale.

% ---------------------------------------------------------------------
\section{Consigli pratici e domande sul brainstorming}
\label{sec:z2-qa-brainstorming}

\subsection{Musica e cambio di scenario}

``Sembrano tutte stupidaggini, ma vi assicuro che aiutano.''

\paragraph{Musica di sottofondo.} È dimostrato che aiuta\footnote{Le affermazioni sugli
effetti della musica (chiave maggiore, rumore ambientale) sono riportate come dette in
lezione e non sono state verificate.}. Meglio una musica con \textbf{tempo e volume
costanti}: niente \emph{bass drop}, niente variazioni di tema o di ritmo. Va bene la
musica fatta di rumore ambientale a basso volume; è stato rilevato anche che le canzoni
in \textbf{chiave maggiore} aiutano la parte creativa a produrre idee più originali.
Meglio la \textbf{musica strumentale}, senza testo: per quanto ci si impegni a non
ascoltarlo, il cervello cercherà di capire le parole. È un effetto marginale e non
succede niente se capita la canzone con il testo, ma si mette in conto che ``l'uno per
cento del cervello'' sarà impegnato a capirlo, come quando in treno si ascolta, che lo si
voglia o no, la telefonata di qualcun altro. È \emph{min-maxing}.

\paragraph{Cambiare scenario.} A Riot c'era una stanza dedicata al brainstorming, pensata
per essere creativa: le finestre sul fiume, le porte a vetri colorati, tutto fatto per
ispirare creatività. Ma fare brainstorming sempre nella stessa stanza generava sempre lo
stesso tipo di idee: anche solo cambiare ambiente aiutava ad averne di diverse. Il
relatore spesso spostava il meeting, ad esempio nello Starbucks sotto l'ufficio. Lo
spazio fisico gioca un ruolo importante, ma anche quello virtuale: se lo si fa sempre su
Discord, si può provare con ognuno uno sfondo diverso e simpatico. Oppure: ha sempre
moderato la stessa persona? Si prova con un altro moderatore. Si cambiano i fattori
perché i risultati non siano sempre gli stessi.

Uno studente chiede se il cambio di scenario possa essere anche qualcosa come passarsi una
palla per avere la parola, o un altro espediente per tenere viva l'attenzione. Con questo
modello, che non richiede di prendere la parola, l'esempio della palla perde senso, ma come
esempio va benissimo: anche il comportamento fisico e il \emph{modus operandi} sono un ottimo
cambio di scenario. Come lo sono cambiare il moderatore o \textbf{cambiare l'obiettivo}: se
il cliente chiede X, gli si darà un'idea X, ma si può provare anche a votare per un'idea Y,
facendo più votazioni con vari criteri di valutazione.

\subsection{Domande}

\paragraph{L'idea scelta deve essere approvata da qualcun altro?} Se un manager dice ``fate
questo brainstorming senza di me'', sta \textbf{delegando la scelta}: ``mi fido di quello che
sceglierete''. Avendo preparato prima limiti, paletti e \emph{constraint}, si aspetta che la
soluzione portata sia in linea con essi; ma se non partecipa deve fidarsi della scelta. Se
poi la scelta debba passare anche da un livello superiore dipende dall'azienda e dagli
accordi precedenti.

\paragraph{Si può fare un brainstorming ibrido, tra presenza e remoto?} Uno studente fa
l'esempio di un gruppo in cui uno è a New York e gli altri a Milano. Sì, senza problemi: si
fa una call, il portatile punta verso il muro con le idee, o le si mostra fisicamente con
il cellulare. Chi è da remoto scrive i suoi post-it, ad esempio mandandoli su WhatsApp al
moderatore o a qualcuno presente che li scrive per lui, e vota dicendo ``mettimi un pallino
qui, qui e qui'' (in quel caso il voto non è anonimo, ma si può fare tranquillamente); se
deve parlare, parla.

\paragraph{E se i colleghi non riconoscono o non condividono il metodo?} Ci deve essere un
\textbf{agreement}: non si può arrivare e dire ``facciamo così'' se gli altri non vogliono.
Questo metodo funziona solo se tutti sono d'accordo a usarlo. Il relatore ne ha proposto
uno tra i cento metodi di brainstorming esistenti, un buon compromesso tra semplicità ed
efficienza (ce ne sono di più semplici e di molto più efficienti ma più complessi). Se il
gruppo non ne vede il valore, al massimo insisterebbe per \textbf{fare almeno una prova},
altrimenti diventa chiusura mentale; se poi per quel team non funziona, se ne cercano
altri.

\paragraph{E nei gruppi molto piccoli, dove diventa un duello?} Uno studente ha lavorato in
gruppi di tre o anche di due persone, dove il brainstorming diventa quasi un duello in cui
ognuno difende la propria idea. Bisogna \textbf{aspettare la discussione}: invece di ``scegliamo
tra la mia idea e la tua'', si cerca di concordare quale idea abbia l'impatto maggiore e
quella lo sforzo minore. E se non si è d'accordo, bisogna \textbf{portare esempi}: se uno
dice ``questa idea costa di più'', l'altro chiede ``in che senso? Che cosa costa di più?''. Se
non sa rispondere, non si può fare qualcosa senza sapere come va a finire: meglio rimandare
la discussione a quando si avranno le informazioni. A volte la domanda giusta è \emph{perché
non abbiamo queste informazioni?}: forse è \textbf{prematuro} decidere, perché non si hanno
ancora tutti i fattori.

Nei team piccoli non dovrebbe esserci un ``si fa come dici tu o come dico io'', una scelta a
pelle: si dovrebbero conoscere così bene le proprie proposte da poterle posizionare con
assoluta certezza sulla matrice, e avere anche la presenza di spirito di dire ``la tua è
chiaramente migliore, si fa la tua''. Si fa la cosa migliore per entrambi, che dovrebbe
essere oggettiva e non soggettiva. Non è facile arrivarci, ma deve essere l'obiettivo del
team.

% ---------------------------------------------------------------------
\section{Disagree and commit}
\label{sec:z2-disagree}

\subsection{Quando la scelta non ci convince}

Uno studente teme che chi vede scartata la propria idea perda motivazione, in un gruppo
piccolo come in uno grande. Non è un deragliamento: è una conseguenza potenziale di
qualsiasi scelta. Il relatore ha visto gente demotivata perfino dalla settimana di quattro
giorni: ci sarà sempre qualcuno più o meno motivato dalle scelte. Ma qual era l'alternativa?
Un brainstorming classico in cui tutti parlano, chi più chi meno, e si sceglie un'idea che
tutti pensano sia la migliore senza esserne davvero sicuri, perché non si è votato, né
espresso una preferenza, né discusso davvero le idee migliori.

Qui entra il concetto di teamwork del \textbf{disagree and commit}: ``non sono d'accordo con
questa decisione, non mi piace, non penso sia la migliore, ma la faccio''. Per due motivi:
\begin{enumerate}
  \item perché \textbf{è il mio lavoro}, e dovrebbe essere il punto principale;
  \item perché \textbf{voi ritenete che sia la scelta migliore}. Se si è la voce fuori dal coro
        e tutti hanno scelto un'altra strada, non si rema contro: è \emph{default to trust}.
        Se gli altri vi hanno visto più potenzialità, e (grande premessa) \textbf{hanno
        ascoltato} il proprio disaccordo e i suoi motivi, e hanno comunque scelto così, ci si
        fida di loro.
\end{enumerate}
Con una postilla: se le cose vanno male, \textbf{non si dà la colpa agli altri} e non si dice
``ve l'avevo detto'', come se non si facesse parte del team. Disagree and commit vuol dire ``mi
fido del team e facciamo questa cosa''; se poi si rivela sbagliata, \emph{il team} ha
sbagliato, e la volta dopo si cambierà il modo di scegliere. Altrimenti non è disagree and
commit, è ``\emph{disagree and ostracize}'': disaccordo più ostilità. Non bisogna cercare il
colpevole, perché si è deciso tutti insieme (sezione~\ref{sec:z1-default-trust}).

Se si pensa davvero che si stia sbagliando clamorosamente il posizionamento di un'idea, si
dice piuttosto ``non andiamo avanti finché non definiamo bene lo sforzo di questa cosa''; ma
bisogna saperlo giustificare, altrimenti è chi dissente a dover fare un passo indietro.

\subsection{La demotivazione fa parte del lavoro}

La demotivazione ci sarà sempre, perché non si sceglierà mai solo ciò che piace: per quanto
piaccia il proprio lavoro, ci saranno task sgradevoli e scelte che non spettano a noi. E
ci saranno scelte (``ve lo dico per esperienza personale'') che si è convinti siano
sbagliate, al punto di pensare ``piuttosto lascio il team'', ma che si finisce per accettare,
per poi rendersi conto che gli altri avevano ragione. Come c'è un \emph{bias} verso l'amico e
il nemico, c'è anche un \textbf{bias verso noi stessi}: crediamo di conoscere bene una cosa, e
se altre dieci persone dicono ``secondo me è fattibile'', bisognerebbe aprirsi alla
possibilità che ne sappiano più di noi.

``Non lo accetto ma lo faccio'' dovrebbe essere alla base del lavoro insieme, perché è
\emph{default to trust}: rifiutarsi (``non mi fido di voi, non lo faccio'') significa non
fidarsi del giudizio del proprio team. E dire ad alta voce ``non sono d'accordo, ma ti darò
comunque il massimo'' è la cosa che \textbf{fa guadagnare di più il rispetto degli altri}.
Non c'è niente di meglio che vedere qualcuno che aveva un'opinione diversa impegnarsi per
far vedere di avere torto, facendo andare bene la cosa. Per il relatore dimostra solo che è
una persona matura e importante; ha avuto persone che, passate da questa situazione, gli
hanno fatto pensare ``è molto più matura di quanto credessi'', e che ha valorizzato molto di
più in seguito, perché sa di potersi fidare di chi fa disagree and commit: è stato aperto
nella comunicazione, ha detto che cosa secondo lui non andava, ma non ha fatto niente per
remare contro, anzi ha dato il suo supporto. ``Non posso chiedere di più a una persona.''

\subsection{Chi arriva a lavori iniziati}

\paragraph{Se il team cambia in corso d'opera, chi subentra come si rapporta alle decisioni
già prese?} (per dimissioni, promozioni e così via). Deve \textbf{accettare lo status quo}. Al
relatore capita spesso di entrare in team con progetti già in corso: si fa dire che cosa è
stato scelto finora e va \textbf{in commit di default}: ``non ho preso parte a questa decisione,
quindi non sono in \emph{agreement}, ma ora faccio parte del team, quindi vado in commit''.

\paragraph{E se l'esperienza dice che una scelta ha uno sforzo molto diverso?} Ad esempio
che un'idea ``rossa'' ha in realtà lo sforzo di una ``viola''. Lo si fa presente, ma dipende dal
\textbf{tempismo}: se la scelta è appena stata fatta e c'è ancora modo di cambiarla, la si
segnala; se si è al 90\% del progetto, è controproducente dire ``fermi tutti, passo
indietro'', e a un certo punto bisogna rispettare la scelta del team (``avete scelto così,
andiamo avanti e vediamo come va''). E se si sa che, calcolando ferie e periodo estivo,
un'animazione prevista per giugno arriverà a settembre? È molto \emph{case by case}: a
seconda dell'apporto si può modificare il risultato o no. È come riaprire un processo se
emergono nuove prove: se dimostrano che l'imputato non è l'assassino, è giusto riaprirlo; se
dimostrano che l'assassino ha sparato con la sinistra invece che con la destra, conta fino a
un certo punto. È giusto segnalarlo, ma non sempre quello che si porta giustifica uno
``\emph{scratch that}, rifacciamo tutto'': diventa \textbf{un feedback, non una risoluzione}
(``state facendo così, questo è il mio feedback, poi vediamo come va'').

Con cinque minuti di anticipo si va in pausa pranzo, con l'accordo per il pomeriggio: circa
un'ora di presentazione più ``standard'' sul leading from within, poi domande e risposte,
``e quando finiscono le domande abbiamo finito''.

% ---------------------------------------------------------------------
\section{Leading from within: perché parlare di leadership}
\label{sec:z2-leadership}

\subsection{La ripresa del pomeriggio}

Alla ripresa, mentre si aspettano gli ultimi arrivati, il relatore chiede se durante la
pausa siano venute in mente domande. L'unica è logistica: le slide verranno date a fine
workshop? Sì: il relatore le passerà all'accademia (a Vanessa, Andrea ed
Enrico\footnote{Nella trascrizione l'accademia è chiamata ``Ive'' e ``Hive'', come nella
prima giornata; per i nomi dei referenti vedi la nota all'inizio del
capitolo~\ref{cap:z1-feedback}.}), che le farà avere alla classe. Gli appunti vanno
comunque presi, meglio su quello che dice che su quello che mostra.

L'ultima parte, dice, dovrebbe essere abbastanza veloce da presentare, ma vale la pena
fermarsi a parlarne, perché è un discorso più avanzato. È la logica chiusura di quanto
detto tra la prima giornata e la mattina: tirare le fila e mostrare come tutto faccia parte
di un unico filone di ragionamento e progettualità, con un unico scopo, un unico obiettivo e
un unico modo di ragionare.

La voce sembra a posto, ma il relatore è stanco dopo due giorni, come immagina lo sia la
classe, nello \emph{stretch} finale di due giornate molto piene. Sono cose che in condizioni
normali spiega al suo team nell'arco di settimane, con il tempo di metterle in pratica,
discuterle e approfondirle. Qui è tutto concentrato: servirà tempo per metterlo in atto e
fare le proprie prove, quindi non bisogna spaventarsi se non verrà tutto giusto la prima
volta.

\subsection{Il leader non è il più bravo}

Perché parlare di leadership in un workshop non rivolto solo ai leader? La risposta
completa arriverà più avanti, ma il punto di partenza è che il relatore \textbf{non crede
nei leader di team o di progetto che detengono da soli il sapere e il modo di fare le
cose}. Come detto il giorno prima, essere un leader, un manager o un capo non dovrebbe
voler dire essere il migliore in quel lavoro, ma essere \textbf{la persona più adatta a
ricoprire quel ruolo}: guidare un gruppo significa \textbf{occuparsi delle persone} al suo
interno.

L'obiettivo che il relatore si dà, e che secondo lui ogni manager dovrebbe avere in mente,
è che quando si torna dal lavoro ci si senta \textbf{ispirati, sicuri e realizzati}, e che
questo sia \textbf{un diritto naturale di tutti}, non un lusso riservato a chi ricopre certe
posizioni, a pochi fortunati o ``a qualcuno sì e a qualcuno no''. Ha notato,
specialmente in Italia, la tendenza a vedere tutto ciò che viene prima del diventare
manager come la \emph{gavetta} necessaria per arrivarci. Invece chiunque faccia parte di
un'organizzazione, che sia un'azienda di videogiochi o altro, dovrebbe sentirsi soddisfatto,
ispirato e sicuro in quello che fa; e non è corretto che un'azienda tratti bene le persone
``da manager in su''. Anche quando non lo dicono, a volte i fatti sono quelli: da manager si
viene trattati in un certo modo e si ha accesso a certe cose, altrimenti no. Non si tratta di
\textbf{appiattire la gerarchia}: si tratta di assicurarsi che, pur dando cose diverse ai
diversi livelli, a tutti i livelli venga riconosciuto uno strato comune di
\textbf{rispetto, fiducia e professionalità}.

\subsection{Le carriere divise: manager o \emph{subject matter expert}}

Uno studente chiede conferma: quindi un \emph{lead artist} non è necessariamente più bravo
di un senior a fare l'artista? Esatto. E qui il relatore apre una parentesi. In molte
aziende, e soprattutto in quelle in cui tende a lavorare lui, la motivazione principale per
diventare manager è avere una promozione e quindi un aumento, più che il desiderio di
prendersi cura del team e delle persone. Come si ovvia al problema? \textbf{Dividendo le
carriere}. Arrivati al livello senior, il relatore fa con ognuno dei suoi senior
\textbf{``the Talk''}, per esporre le due opzioni che hanno davanti:
\begin{description}
  \item[Busta A: diventare manager] È il percorso standard: si valuta se si è adatti a fare
        il leader e a guidare un team, e si danno gli strumenti per farlo, perché
        ``nessuno nasce leader imparato'': prima di provarlo sulla propria pelle nessuno sa
        che cosa voglia dire. Per quanto si possa essere carismatici, fare il leader è
        un'altra cosa, e non vuol dire solo essere carismatici. Ci sono formazione,
        strumenti e un percorso di crescita.
  \item[Busta B: diventare \emph{subject matter expert}] Una persona che va oltre il senior
        nella stessa disciplina: diventa \textbf{l'esperto dell'intera azienda} su quella
        disciplina, perché la sua priorità non è aiutare gli altri ma essere il migliore.
        ``Non c'è assolutamente nulla di male.'' Il suo ruolo è essere il più specializzato
        in quel campo, il punto di riferimento da cui tutti vanno a farsi spiegare le cose.
\end{description}
Il concetto è abbastanza nuovo, e il punto chiave è che l'azienda \textbf{paga il
\emph{subject matter expert} e il leader esattamente allo stesso modo}: non c'è un guadagno
economico nello scegliere di fare il manager invece dello specialista, e il relatore valuta
quel tipo di conoscenza quanto il saper gestire persone. Così chi vuole essere il migliore,
o anche solo un aumento, può continuare a fare ciò che sa fare e avere l'aumento; chi vuole
l'aumento \emph{e} prendersi cura delle persone sceglie l'altra strada. Così si fa una scrematura sul
\emph{motivo} per cui si vuole fare il manager. Sembra una scelta piccola, ma è molto
importante capire se ci si vuole occupare di se stessi o degli altri. Non è un giudizio
(``non sto dicendo che sia meglio l'uno o l'altro''): le due cose \textbf{non convivono}, e
bisogna sapere che cosa si vuole fare, sapendo che l'azienda tratterà allo stesso modo le
due scelte.

\paragraph{Si può essere entrambi?} Si può scegliere ad esempio il percorso da lead artist
e voler comunque continuare a migliorare le proprie capacità. Ma, come nelle simulazioni di
conversazione del giorno prima, il relatore chiederebbe: ``se avessi le risorse per fare solo
una delle due cose, quale sceglieresti?''. Perché non le si avrà: non ci sarà il tempo e il
modo di fare bene entrambe. Il suo ``cinismo'' arriva a dire che bisogna scegliere, perché
tenersi buone entrambe le strade porta a farle male tutte e due. Preferisce un ottimo
\emph{people manager} o un ottimo specialista; le figure intermedie, in questo campo,
rischiano di diventare pericolose per sé e per il team.

\paragraph{Specializzazioni verticali.} La strada dello specialista permette
specializzazioni molto verticali: in chat qualcuno cita il \emph{mecha designer}, ma può
esserci anche lo specialista di animazioni delle espressioni facciali. Naturalmente l'azienda
deve avere bisogno di quella competenza (non ci si specializza in cose che non servono), e non
è detto che ne serva uno solo. È un percorso di carriera parallelo e alternativo a quello del
manager ``standard''.

% ---------------------------------------------------------------------
\section{Il ``momento Piero Angela'': la biologia del gruppo}
\label{sec:z2-biologia}

\subsection{Perché una digressione scientifica}

Il relatore annuncia un ``piccolo excursus biologico'', fondamentale per capire dove vuole
arrivare, e prova a mettere un sottofondo musicale adatto: la condivisione dell'audio non
funziona, quindi manda il link in chat perché ognuno lo faccia partire per conto suo. Un
studente lo battezza subito ``momento Piero Angela''. Il relatore cerca di far combaciare le
motivazioni aziendali del suo metodo di leadership e di progettualità con \textbf{basi
scientifiche}: in passato ha collaborato con varie università che si occupavano di
neuromarketing e neurobiologia, per avere una conferma o una smentita del fatto che ciò che
ipotizzava ``non fosse una cavolata colossale''. Semplificando moltissimo, si guardano solo
alcune sostanze che si traducono direttamente nel modo in cui si vive l'essere in un gruppo,
in un team, in un'azienda\footnote{L'impianto di questa parte (endorfine, dopamina,
serotonina, ossitocina e cortisolo, il villaggio protetto dai pericoli esterni, il ``cerchio
di sicurezza'') ricalca da vicino quello del libro di Simon Sinek \emph{Leaders Eat Last}
(2014), che il relatore cita poco più avanti. Le descrizioni biologiche sono molto
semplificate, come avverte lo stesso relatore, e vanno lette come tali.}.

\subsection{Endorfine}

Le \textbf{endorfine} sono rilasciate per \textbf{mascherare il dolore}. Lo sa chi corre, chi
fa pesi o ginnastica, ma anche chi ha riso così tanto per una battuta da avere male ai
fianchi: è quella sensazione di dolore in cui però si è contenti di provarlo. Il \emph{runner's
high} è l'esempio più calzante: essere distrutti dopo una corsa ma soddisfatti. È la risposta
biologica allo sforzarsi in cose che sono per il nostro bene, un modo con cui il corpo ci
``convince'' a farle.

\subsection{Dopamina}

La \textbf{dopamina} è più familiare: è la \textbf{ricompensa per un'azione}. Dalle cose più
banali (si mangia e si riceve dopamina) a quelle viste in mattinata: anche completare una
checklist rilascia dopamina (sezione~\ref{sec:z2-pianificare}). Viene rilasciata quando si
vede l'obiettivo, come il traguardo durante una corsa; è la stessa cosa che accadeva ai
cacciatori del villaggio in tempi antichi, che rilasciavano dopamina quando vedevano le tracce
e venivano così incentivati a seguirle. Sono concetti fondamentalmente ``animali'', ma si
applicano tuttora anche in azienda.

\paragraph{Una nota di colore: la dipendenza.} La dopamina \textbf{crea dipendenza}: molti casi
di ossessione per alcol, nicotina, gioco d'azzardo e social media sono problemi legati alla
dipendenza da dopamina. Secondo il relatore non è un caso che le diagnosi di ADD e ADHD siano
aumentate del 60\% negli ultimi dieci anni\footnote{Nella trascrizione ``EDD e HD''. Il dato
e il nesso causale con la maggiore disponibilità di alcol, nicotina, gioco d'azzardo e social
media sono riportati come detti in lezione e non sono stati verificati; l'aumento delle
diagnosi viene in genere attribuito anche ad altri fattori, come la maggiore consapevolezza e
il cambiamento dei criteri diagnostici.}, perché buona parte di quegli stimoli è diventata
ampiamente disponibile proprio in quel periodo. Anche la dipendenza dal lavoro, a volte, ha la
stessa origine: si vuole il rispetto che deriva dal lavoro, il \emph{rush} che deriva
dall'averlo.

\subsection{Il villaggio di diecimila anni fa}

Si fa poi un salto di diecimila anni, in un villaggio. Gli esseri umani, per quanto più o meno
introversi o estroversi, sono \textbf{animali sociali} che tendono a vivere esperienze di
gruppo; le aziende nascono da un gruppo di persone che si danno un obiettivo condiviso, come
in antichità un gruppo di persone che voleva sopravvivere insieme. Quello che non è cambiato è
che \textbf{fuori dal gruppo ci sono dei pericoli}: tigri, lupi, draghi (``do per scontato
che nell'antichità ci fossero i draghi, perché D\&D mi insegna così''). Per evitarli, i gruppi
tendono a riunirsi e a creare difese intorno a sé: da qui le prime difese dei primi villaggi,
per isolarsi dai pericoli esterni. ``Vi giuro che avrà tutto senso.''

\subsection{Serotonina}

La \textbf{serotonina} viene rilasciata anche \emph{indirettamente}: non deve per forza
riguardare noi. Quando si assiste a una cerimonia o celebrazione pubblica che include una
persona cara, si rilascia serotonina come se si fosse noi l'oggetto del riconoscimento. È uno
dei motivi per cui si fanno diplomi e lauree ``all'americana'', con la consegna sul palco
davanti a tutti: gli studi avrebbero mostrato che anche i genitori in platea rilasciano
serotonina, come chi riceve il premio. Lo stesso vale per parenti, amici, matrimoni: sono
esperienze condivise, e si tende a condividere in pubblico riconoscimenti e successi perché in
tutti i partecipanti viene rilasciata serotonina.

Purtroppo ha a che fare anche con i beni materiali: la tendenza all'\emph{hoarding}, cioè a
collezionare eccessivamente qualcosa, nasce dal fatto che quell'attività rilascia serotonina,
ingannando la mente nel credere di aver ottenuto un successo.

È chiamata anche \textbf{la sostanza della leadership}, perché aiuta a riconoscere in qualcun
altro chi ha valore per il gruppo. Nel villaggio, quando arriva un pericolo, c'è sempre un
gruppo di persone che dice ``ce ne occupiamo noi'': il gruppo sa che si stanno prendendo un
rischio, e la serotonina è il modo per riconoscere che meritano rispetto. Il relatore consiglia
il libro di Simon Sinek \emph{Leaders Eat Last}\footnote{Nella trascrizione ``Leadership
Lust''. Il titolo è \emph{Leaders Eat Last: Why Some Teams Pull Together and Others Don't}
(2014).}: i leader mangiano per ultimi,
all'opposto di quanto accadeva in antichità, quando al ritorno dalla caccia mangiavano per
primi i cacciatori, che dovevano essere in forze per la caccia successiva, poi donne e
bambini, poi gli altri. C'è un \textbf{riconoscimento sociale}: tu fai qualcosa per proteggere
il gruppo, noi ti riconosciamo qualcos'altro in cambio.

Le aziende sono strutturate in modo molto diverso? È una domanda retorica: quando si dice che
serve un capo, un manager, un rappresentante, si sta dicendo che quella persona dovrebbe avere
l'interesse di proteggere il gruppo e di prendersi cura di tutti. Per questo sapere che il
proprio capo guadagna di più non fa particolarmente arrabbiare: lo si vede come quasi naturale,
perché più responsabilità si ha verso il gruppo più si guadagna, dato che quello dovrebbe essere
il suo scopo. Ed è anche il motivo per cui il relatore si arrabbia molto quando legge di CEO o
\emph{executive} che, a fronte di licenziamenti di migliaia di persone, si assegnano dei bonus.
È una rabbia non del tutto razionale, che sente da dentro: è come se si stesse violando un
\textbf{contratto sociale}. Il contratto è ``tu ti prendi dei rischi per me e mi difendi, io ti
riconosco più cibo, più soldi''; quando si abusa di quella posizione per un beneficio personale
a scapito di chi si dovrebbe proteggere, tutti lo vedono come un'ingiustizia, perché si viene
meno al contratto stabilito diecimila anni fa con la nostra biologia.

\subsection{Ossitocina}

L'\textbf{ossitocina}\footnote{Nella trascrizione ``ossidocina''.} è la \textbf{sensazione di
sicurezza} che si genera quando si è vicini a qualcuno. Non è un caso che tra i legami di
affinità e collaborazione più forti ci siano quelli dei team che si uniscono contro un
\textbf{nemico comune}: i team che superano insieme avversità e difficoltà ne escono molto
legati. Tra le persone che il relatore sente più spesso in questo periodo ci sono quelle con cui
ha condiviso momenti difficili, un cattivo manager, una brutta esperienza lavorativa: sa che
erano lì con lui. Vale anche per i veterani di guerra o per i sopravvissuti a una stessa
tragedia, che sviluppano un rapporto molto forte (se ne è parlato molto negli Stati Uniti a
proposito dei veterani), perché hanno rilasciato immense dosi di ossitocina. È anche ciò che fa
accettare subito il figlio alla madre quando lo prende in braccio dopo il parto, ed è il motivo
per cui, a migliaia di anni di distanza, si stringono ancora accordi con una stretta di mano o
un abbraccio: c'è sempre qualcosa di fisico, generato dal bisogno di ossitocina.

Viene generata anche quando si compiono o si \textbf{assiste ad atti di generosità} senza
aspettarsi nulla in cambio. Se il relatore dicesse di aver donato cinquanta euro a
un'associazione di beneficenza, la reazione sarebbe ``ok, bene''; se dicesse che ogni sabato
pomeriggio fa volontariato, lo si vedrebbe come qualcosa di molto più prezioso, e la cifra
potrebbe crescere parecchio prima che qualcuno la consideri più preziosa dei suoi pomeriggi:
tutti sanno che la vera donazione è \textbf{il tempo} investito, che vale più del bene
materiale. E basta assistere: vedere qualcuno che aiuta un'anziana ad attraversare la strada fa
stare meglio anche noi. È percepita a livello di gruppo. Infine, l'ossitocina rafforza il
\textbf{sistema immunitario}: sentirsi sicuri e tranquilli fa stare meglio e, in un certo senso,
ammalare di meno.

\subsection{Cortisolo}

L'opposto è il \textbf{cortisolo}: la sensazione di stress, paranoia e paura, la risposta
\emph{fight or flight} (o scappo o combatto), con battito accelerato, muscoli in tensione e così
via. Tra i suoi effetti:
\begin{itemize}
  \item \textbf{blocca il sistema immunitario}: se si sta per morire, al corpo non importa di
        prendersi un raffreddore. È il motivo per cui le persone stressate si ammalano di più;
  \item \textbf{blocca l'ossitocina}: non ci si può sentire tranquilli quando si è in ansia;
  \item soprattutto, e questo è il punto più importante per le aziende, è
        \textbf{contagioso}. In qualsiasi documentario, quando una gazzella che bruca alza la
        testa e percepisce un movimento nella savana, scappano tutte, non solo lei. Succede,
        non sorprendentemente, anche in ufficio: quando qualcuno ha paura perché potrebbero
        esserci licenziamenti, o perché il manager è molto aggressivo, la paura si diffonde nel
        gruppo ed è molto dannosa.
\end{itemize}
È anche l'effetto alla base della \textbf{sindrome dell'impostore}: sentirsi costantemente
inadeguati perché si teme che succeda qualcosa significa vivere in uno stato di stress. Ed è il
motivo per cui chi è insoddisfatto del lavoro tende a stare peggio anche fisicamente e ad
ammalarsi di più.

\begin{center}
  \begin{tabular}{@{}p{0.17\textwidth}p{0.36\textwidth}p{0.37\textwidth}@{}}
    \toprule
    \textbf{Sostanza} & \textbf{Che cosa fa (in sintesi)} & \textbf{Esempi dalla lezione} \\
    \midrule
    Endorfine & mascherano il dolore di uno sforzo che ci fa bene &
      \emph{runner's high}, ridere fino ad avere male ai fianchi \\
    \addlinespace
    Dopamina & ricompensa un'azione, si attiva vedendo l'obiettivo; crea dipendenza &
      spuntare una checklist, le tracce per il cacciatore, gioco d'azzardo, social \\
    \addlinespace
    Serotonina & riconoscimento sociale, anche indiretto; ``sostanza della leadership'' &
      lauree e premi in pubblico, chi difende il villaggio, \emph{hoarding} \\
    \addlinespace
    Ossitocina & sicurezza e legame; rafforza il sistema immunitario &
      il nemico comune, i veterani, la stretta di mano, il volontariato \\
    \addlinespace
    Cortisolo & stress, \emph{fight or flight}; blocca difese immunitarie e ossitocina; è
      contagioso & la gazzella che scappa, la paura dei licenziamenti, la sindrome
      dell'impostore \\
    \bottomrule
  \end{tabular}
\end{center}

% ---------------------------------------------------------------------
\section{Domande sulla biologia del gruppo}
\label{sec:z2-qa-biologia}

\paragraph{Il gioco d'azzardo include i gacha e le loot box?} Assolutamente sì, anzi ne sono il
\textbf{driver principale}: i gacha, i FIFA Ultimate Team, lo ``sbustare'' da \emph{Magic} a
\emph{Hearthstone}, e in generale il ``vediamo se stavolta mi va bene'' sono tutti gioco
d'azzardo. Lo studente aggiunge che molti giochi mobile sono costruiti su questi meccanismi e che
c'è chi gioca più a quelli che ai giochi ``normali''. È il motivo per cui alcuni stati hanno
dichiarato illegali certi giochi\footnote{Il caso più noto è quello del Belgio, che dal 2018
considera gioco d'azzardo le \emph{loot box} a pagamento; anche i Paesi Bassi le hanno
contestate.}.

\paragraph{Come si genera l'ossitocina nei rapporti da remoto?} ``Bellissima domanda.'' È il motivo
per cui, come detto il giorno prima, il relatore insiste perché i meeting si facciano almeno in
\textbf{webcam}, perché i riconoscimenti siano dati e i successi celebrati \textbf{pubblicamente}
(quando non mette troppo a disagio la persona), e in generale cerca di creare \textbf{situazioni
condivise}. Anche lavorando da remoto, insiste per avere il \emph{team building} dal vivo e almeno
\textbf{un incontro all'anno} di persona in ufficio, perché crea situazioni di rapporto diretto.

\paragraph{Come si contrasta il cortisolo?} In chat qualcuno suggerisce la meditazione: sì,
assolutamente. Per il lavoro, la risposta arriva subito dopo nella presentazione.

\paragraph{Il nemico comune è una strategia che un manager può usare?} Spesso \textbf{il manager è
il nemico comune}, senza bisogno di sforzarsi. Ma il relatore chiede di capire bene il punto:
alcuni manager applicano consapevolmente il concetto ``mi ergo io a nemico del team, così sono
tutti uniti contro di me e lavorano bene'', e ne fanno uno stile di management. Decidono di essere
il \emph{bad cop}, il \emph{villain}, ed è così che ottengono risultati. È una scelta lecita e
funziona; il relatore preferisce fare diversamente, ma non per questo dice che non funzioni. Molte
volte si sottovaluta quanto sia consapevole: ``non si rende conto di quanto sia antipatico a
tutti''; forse se ne rende perfettamente conto e lo fa apposta per essere il bersaglio. Il relatore
l'ha visto fare spesso come soluzione disperata, quando il team era troppo impegnato a litigare al
suo interno: ``piuttosto che vederli litigare tra loro, gli do io un motivo per odiare me, così si
uniscono contro di me''. Non lo consiglia assolutamente, ma l'ha visto funzionare.

Uno studente racconta che nella sua azienda il manager fa vedere \emph{il lavoro stesso} come il
nemico (``prima finiamo, meglio è''). È un modello più basilare ma, come dice lo studente stesso,
funziona.

\paragraph{Quanto è etico puntare sulla dipendenza del giocatore per fare soldi?} ``Non è
assolutamente etico, ed è assolutamente profittevole, altrimenti le aziende non lo farebbero.'' Il
relatore racconta una cosa che lo ha sconvolto quando l'ha scoperta, a Blizzard. Lavorava su
\emph{Hearthstone}\footnote{Nella trascrizione ``Arsdol''.} e si stava discutendo di come
monetizzare il gioco; venne preso come consulente un esperto, israeliano se ricorda bene, che
lavorava a Malta per società di poker online e di \emph{gambling}. Raccontò che molte aziende che
monetizzano sulla dipendenza dal gioco d'azzardo avevano sviluppato metodi più o meno etici (``più
spesso meno''). L'esempio culminante: analizzavano gli utenti occasionali, quelli che giocavano
magari una settimana al mese, per capire \emph{quando arrivava loro lo stipendio} (di solito si
deposita lo stipendio e poi si gioca quello che resta), per colpirli il giorno dopo con email di
offerte per tornare a giocare. Li si stimolava proprio nel momento in cui avevano disponibilità, sapendo
che erano sensibili, per aumentare le probabilità di riaverli come giocatori. ``È orribile. È
terrificante e, come tutte le cose terrificanti, è anche altamente funzionale.''

Gli ha fatto riflettere molto su dove ci si possa spingere per fare soldi. La versione etica, per
lui, è ad esempio far uscire un gioco sotto Natale, quando si è più disposti a regalarlo ai figli:
si stanno ottimizzando le proprie possibilità, e lo accetta. Si scade nel \textbf{predatorio}
quando si sfrutta la debolezza di qualcuno in un momento preciso per estorcergli soldi che non può
permettersi di perdere. Poi ognuno tragga le proprie conclusioni: ``non è un workshop sull'etica del
lavoro''. Non pensa di dover imporre la propria opinione, né si arroga il diritto di dire che chi lo
fa è una bestia senza cuore, visto che è il loro lavoro: ``chiunque riesca poi a dormire in pace con
se stesso, buon percorso''.

% ---------------------------------------------------------------------
\section{La piramide della fiducia}
\label{sec:z2-piramide}

\subsection{La leadership non è un titolo}

Chiuso il ``momento Piero Angela'', il relatore arriva al punto. La frase della prima giornata (``la
capacità di un gruppo di persone di fare cose straordinarie dipende dal modo in cui queste persone
operano insieme'') porta a capire che \textbf{la leadership non è un titolo lavorativo}. Essere
manager da contratto, o avere ``manager'' nel titolo, rende un manager, non un leader. Il relatore
non ha mai visto un team rispettare qualcuno solo perché da contratto doveva farlo: si può avere
paura o timore del proprio manager, ma non lo si deve rispettare perché è il manager. \textbf{Il
rispetto si guadagna}, e si possono influenzare le persone intorno solo attraverso il rispetto
guadagnato, quando si sono accumulati abbastanza esempi e \emph{case history} da giustificare
rispetto e fiducia. Altrimenti, dice, sarebbe il primo a chiedere ``perché vi fidate di me, se non
vi ho ancora dato motivo di farlo?'': non sarebbe razionale.

Da qui la \textbf{piramide della fiducia}, che si costruisce partendo dalla base.

\subsection{Alla base: l'integrità morale}

Per essere un buon leader o un buon manager bisogna \textbf{voler genuinamente aiutare gli altri
membri del team}. È il motivo del discorso sulle carriere divise: è a questo livello che si
screma chi vuole la promozione per i soldi. Chi vuole essere leader non deve cercare il guadagno
personale, ma quello del gruppo. Come detto il giorno prima, il relatore è pronto a litigare per il
suo team e a prendersi le colpe per il suo team, perché vuole il bene del suo team e sa che poi loro
si prenderanno cura di lui. Chi ha in mente questo sta in questa categoria. Poi:
\begin{itemize}
  \item \textbf{essere onesti}. Il relatore può affermare di non aver mai mentito
        intenzionalmente a un collega, a un cliente o a un giocatore. Gli è capitato di dare
        informazioni errate perché in quel momento sapeva quelle: ad esempio annunciare che
        un'espansione sarebbe uscita in un certo mese, e poi l'uscita è slittata\footnote{Nella
        trascrizione l'esempio non è coerente: prima si parla di un'uscita a novembre rimandata
        di una settimana, poi di ``pensavo fosse settembre''. Il senso resta chiaro.}. Ha detto
        una cosa sbagliata, ma tutti hanno capito che non stava mentendo: in quel momento pensava
        che fosse vera;
  \item soprattutto, \textbf{accettare la responsabilità dei fallimenti}: quando si sbaglia,
        saper dire ``ho sbagliato''. Il relatore lo ha messo perché ``non avete idea di quanti
        facciano fatica'' ad ammettere di poter sbagliare: c'è sempre una scusante, una
        mitigazione, un ``sì, ma''. Lui vuole sentire se è colpa tua o no;
  \item \textbf{chiedere scusa} e \textbf{ammettere di non sapere}: anche dire ``non lo so'' viene
        spesso evitato, come accettare un fallimento.
\end{itemize}

\paragraph{``Raccontami il tuo più grande fallimento''.} Per lo stesso motivo, nei colloqui il
relatore chiede spesso, più che ``parlami di un tuo grande successo'' o di un progetto andato bene,
\textbf{``mi racconti il tuo più grande fallimento?''}, e osserva come lo raccontano. Chi non si
aspetta la domanda spesso racconta di quella volta in cui ha provato in tutti i modi a fare in un
certo modo, ma il capo gli ha detto di fare diversamente e il progetto è andato male. ``Non mi hai
risposto: mi hai raccontato un progetto andato male, ma il tuo messaggio è che volevi fare
diversamente e qualcuno ti ha detto di fare così. Quindi non è un tuo fallimento.'' È come
rispondere a ``qual è la tua debolezza?'' con ``sono così attento ai dettagli che a volte\dots'':
si fa passare un punto di forza per una debolezza. Il relatore vuole un esempio in cui \emph{tu},
senza altri fattori, hai sbagliato, e come ti sei comportato. Gli dice molto su come quella persona
reagirà ai suoi feedback: se ha davanti ``un muro'' che dirà sempre ``non è colpa mia'', non è adatta
a un ambiente in cui il feedback è alla base di ogni comunicazione. Chi non ammette di aver mai
sbagliato in vita sua, lo scarta.

\paragraph{È possibile colpevolizzarsi troppo per un fallimento?} Sì, assolutamente. Rientra molto
nella sindrome dell'impostore, o nell'aver sempre vissuto rapporti di lavoro tossici: ci si abitua ad
accettare che ci vengano urlate contro le cose e che sia sempre colpa propria, e ci si autoconvince
che lo sia.

\subsection{Al centro: l'affidabilità}

Superato questo livello si arriva all'\textbf{affidabilità}. Per essere un buon leader non bisogna
essere il migliore, ma affidabili: basta essere riconosciuti come bravi in qualcosa, non
necessariamente i migliori, e quell'affidabilità ce la si porterà anche nella leadership. Significa:
\begin{itemize}
  \item \textbf{rispettare la parola data}: se si promette ``parleremo di un aumento al prossimo
        slot disponibile'', se ne parlerà. Non ci si dimentica, e non si dice poi ``in realtà non si
        poteva fare, ma ho voluto provarci fino all'ultimo''. Una promessa va mantenuta;
  \item \textbf{impegnarsi per il tempo necessario} a raggiungere gli obiettivi;
  \item \textbf{saper rifiutare un lavoro} quando necessario: è il saper dire di no della mattina,
        ancora la matrice di Eisenhower (sezione~\ref{sec:z2-eisenhower}). ``Sono sempre gli
        stessi concetti.''
\end{itemize}

\subsection{In cima: la leadership, cioè dare l'esempio}

L'ultimo gradino è la \textbf{leadership}\footnote{Nella trascrizione ``dealership''.}, cioè
\textbf{dare l'esempio}. Non bisogna essere i più bravi; bisogna saper \textbf{offrirsi volontari
non solo quando il progetto è un successo sicuro, ma anche quando c'è il rischio di un
fallimento}. Sono tutti bravi a offrirsi per il gioco nuovo che andrà sicuramente bene.

\paragraph{Il rappresentante dei licenziati.} Al relatore è capitato di offrirsi per rappresentare
il team durante i \emph{layoff}, con il rischio di essere ostracizzato dall'azienda; con quelle
persone è ancora in contatto. Quando Blizzard fece i licenziamenti nella sede francese, vicino a
Versailles\footnote{Nella trascrizione ``i layoff di Versailles''. La sede europea di Blizzard si
trovava a Vélizy-Villacoublay, nei pressi di Versailles; il relatore non precisa l'anno.}, che
coinvolsero quasi tutto l'ufficio, lui non era tra i licenziati, ma venne \textbf{eletto
rappresentante sindacale} di chi lo era. Per lui è stato molto importante: pur non essendo parte in
causa, lo ritennero degno della loro fiducia. Per quanto ne sapevano, avrebbe potuto negoziare e
accettare tutto quello che proponeva Blizzard, perché non ne era toccato: che si negoziasse un mese
di liquidazione in più o in meno non gli portava un soldo. Avrebbe persino potuto usare la
situazione per negoziare un aumento per sé (``gli facciamo accettare tre mesi invece di sette e a te
diamo un aumento''): non era fuori discussione. Se si sono fidati è solo perché c'era una piramide di
questo tipo e lui aveva dato loro abbastanza motivi per farlo.

\paragraph{Autorevolezza, non autorità.} A uno studente che lo riassume così, il relatore risponde
che è proprio questo: un carisma diverso dall'essere simpatici, l'essere un punto di riferimento,
\textbf{più un pilastro di supporto che il motore} che fa andare avanti le cose. Non il motore davanti
a tutti, ma il \textbf{meccanico} che, quando si rompe qualcosa, sai che non ti frega: ti sistema le
cose senza ``farti la cresta'' sulle pastiglie dei freni.

\paragraph{Non abbandonare nessuno.} E soprattutto \textbf{non abbandonare chi è in difficoltà}.
Sembra una banalità, ma è alla base del motto dei Marines statunitensi, \emph{no one left behind}:
sapere che, se si viene feriti, verranno a prenderti anche sotto i missili è un forte elemento
motivazionale per correre quel rischio. In azienda è il motivo per cui il relatore ribadisce sempre
che, se qualcuno è in difficoltà, la sua prima opzione non è licenziarlo ma \textbf{dargli una mano},
e lo ha fatto a volte anche all'eccesso, per far vedere agli altri che quando saranno in difficoltà
non li lascerà in mezzo a una strada. A volte ce l'ha fatta: ha esempi di persone ``a tanto così''
dal licenziamento recuperate per i capelli e diventate ottimi \emph{performer}. Altre volte la
persona non era recuperabile e alla fine è stata licenziata. Ma chi attraversa un periodo difficile
(il \emph{burnout}, la depressione, una relazione finita) sa che la prima risposta non sarà ``ci sei
servito finora, ora puoi andartene'', ma ``che cosa possiamo fare?''. E questo fa lavorare tutti
tranquilli.

% ---------------------------------------------------------------------
\section{Prendersi cura del gruppo}
\label{sec:z2-gruppo}

\subsection{La soluzione: più leader possibile}

In un caso del genere manager e leader coincidono, ed è l'ideale: un manager dovrebbe essere anche
un leader. Ma la vera soluzione è \textbf{avere più leader possibile}. Se l'obiettivo del leader è
prendersi cura degli altri, e il manager dovrebbe essere ``il primo tra i leader'', ancora meglio è
che \textbf{tutti siano leader del proprio gruppo}, perché la leadership non deriva da un titolo né
dall'essere il manager. Il relatore ha visto tanti gruppi andare avanti perché qualcuno al loro
interno era il leader, non perché lo era il manager.

Vuol dire mettere in pratica tutto il discorso della prima giornata su reciprocità, gentilezza,
aiuto e collaborazione. Si sa che, nel peggiore dei casi, c'è il manager, ma in condizioni normali ci
sono i colleghi: e allora si va da chi è in difficoltà a dire ``posso darti una mano?''. Così, quando
saremo noi in difficoltà, magari senza nemmeno rendercene conto, sarà qualcun altro a dirci ``ti vedo
in difficoltà, come posso aiutarti?''.

\paragraph{L'esempio è contagioso, nel bene e nel male.} Quando si ha una posizione anche solo di
rispetto agli occhi degli altri, tutto ciò che si fa ha effetto sugli altri. In negativo vale la
\textbf{slippery slope}\footnote{L'esempio del cartone abbandonato ricorda la cosiddetta ``teoria
delle finestre rotte'' (\emph{broken windows theory}), secondo cui i piccoli segni di degrado
tollerati favoriscono comportamenti peggiori.}: se si accetta un cartone di pizza abbandonato per
strada, domani se ne accetteranno due, dopodomani tre. Ma vale anche in positivo: chi vede qualcuno
fare qualcosa di positivo pensa ``lo voglio fare anch'io''. ``Come fai una cosa, fai tutto'', dice un
proverbio che il relatore attribuisce alla tradizione cinese\footnote{La massima (``how you do
anything is how you do everything'') è attribuita a fonti diverse, spesso alla tradizione zen;
l'attribuzione è incerta.}.

\subsection{Il villaggio di oggi}

Si torna al villaggio del ``momento Piero Angela'', con i difensori e le bestie (reali o no) intorno.
Oggi fuori dal gruppo non ci sono più tigri, ma le \textbf{deadline}, l'\textbf{overwork}, i
\textbf{layoff}, il manager che urla, le aspettative, i \textbf{KPI}. Avere manager che si prendono
cura del gruppo è un buon inizio, ma è ancora meglio avere \textbf{il gruppo che si prende cura del
gruppo}. Non bisogna più limitarsi a pensare che sia responsabilità del manager prendersi cura di
noi: è responsabilità nostra prenderci cura del gruppo di cui facciamo parte.

Che cosa serve? La piramide vista prima. Serve essere percepiti come \textbf{autentici}: il
relatore ormai riconosce abbastanza facilmente chi si vuole ergere a leader perché gli piace
\emph{sentirsi} leader, non perché vuole \emph{fare} il leader, e quella mancanza di genuinità non
aiuta. In generale bisogna \textbf{generare sicurezza} nel resto del gruppo.

\subsection{Smettere di dire ``ci penserà qualcun altro''}

Un esempio concreto è smettere di dire \textbf{``ci penserà qualcun altro''}. Nello sviluppo è
evidente: se si trova un bug che non riguarda il proprio lavoro, tutti sono d'accordo che vada
comunque segnalato, perché altrimenti potrebbe arrivare in \emph{live}, nella \emph{release}.

\paragraph{Le lasagne di PG Esports.} Circa un anno e mezzo prima, quando lavorava ancora a PG
Esports (che organizza anche eventi dal vivo), l'azienda ospitava in ufficio dei potenziali sponsor
venuti a conoscere il team per decidere se sponsorizzare la Milano Games Week. C'era un catering, con
i classici buffet. In mezzo alla confusione (una ventina di persone di PG Esports e una trentina tra
manager, sponsor ed \emph{executive}), il relatore vede per terra, in mezzo alla stanza, una mezza
porzione di lasagne abbandonata. Senza pensarci la raccoglie per portarla via: non è sua, ma è un
piatto in mezzo alla stanza. Nel tragitto fino al bancone del buffet, almeno altri sei tra executive e
sponsor che non lo conoscevano gli consegnano il loro piatto non finito, come a un cameriere: nella
loro testa l'unico motivo per cui qualcuno potesse raccogliere un piatto da terra era che fosse il
suo lavoro. Sul momento non ci fa caso; arrivato al bancone si gira e vede tutto il suo team che ride
della scena.

Torna da loro e ne nasce una discussione: ``avete visto quanto è stato educativo?''. Avevano visto
esattamente come ragionano i gruppi che non sono abituati a pensare ``adesso ci penso io'', ma ``ci
penserà qualcun altro''. In quelle persone non c'era malizia: sicuramente era un cameriere, perché
altrimenti farlo? ``Non è il mio lavoro, lo avrà fatto perché è il suo.'' E diventa un problema quando
la stessa tendenza (``non c'entra con il mio lavoro, ci penserà qualcun altro'') si applica a
qualsiasi altra cosa. Il relatore si aspetta che \textbf{tutti raccolgano quel piattino}: che chi vede
un bug, se non può correggerlo lo segnali, e se può correggerlo lo corregga anche se non è il suo
lavoro, ``perché è un piattino in mezzo alla stanza''.

\paragraph{Da reattivi a proattivi.} Significa passare da un atteggiamento \textbf{reattivo}, cioè di
fatto passivo, a uno \textbf{proattivo}: prevedere i problemi e cercarli intorno a sé per risolverli,
invece di aspettare che esplodano (``finché non esplode non è un problema mio''). Ragionare nell'ottica
del ``e se succedesse, e io me ne fossi accorto senza fare niente?''. Uno studente lo chiama
\emph{mindset da tester}: il relatore si aspetta che, se si vede qualcosa, ci si assicuri di non essere
gli unici ad averlo visto. A volte basta farlo presente: il relatore non potrebbe mai correggere un
bug, quindi se lo trova lo segnala, e va bene così. Ma non bisogna dare per scontato che ``è un bug così
evidente che qualcuno l'avrà già segnalato''. È il paradosso dell'incidente stradale: \textbf{più
l'incidente è grosso, più è probabile che nessuno chiami l'ambulanza}, perché tutti pensano che
qualcuno l'abbia già fatto\footnote{È il cosiddetto \emph{effetto spettatore} (\emph{bystander
effect}), studiato in psicologia sociale a partire dagli anni Sessanta.}.

\subsection{Domande sul ``me ne occupo io''}

\paragraph{Se correggo un bug, devo avvisare il collega?} Dopo una breve pausa per problemi audio
(``facciamo \emph{double testing}''), uno studente chiede se, correggendo il lavoro di un collega, il
dovere sia poi andare a dargli il feedback. A prescindere bisogna informare qualcuno; e la
conversazione dovrebbe essere: ``ho trovato questa cosa, vuoi che te la sistemi io?''. Correggere
senza dirlo a nessuno potrebbe essere un danno: magari era \emph{working as intended}, o qualcuno ci
stava già lavorando.

\paragraph{E se divento ``quello che risolve quel problema''?} Uno studente porta la sua esperienza:
in un'azienda in cui ha lavorato tre anni, risolvendo un problema facile che non era di sua
competenza è diventato, di fatto, l'incaricato di quel problema, distolto dalle sue mansioni. È
un'arma a doppio taglio. Per il relatore lì il problema è del \textbf{management}. In alcune aziende si
applica il criterio ``chi trova il problema lo risolve'' (``se te ne sei accorto tu, pensaci tu''), ed
è un altro modo di scaricare la responsabilità sugli altri. Non è ciò che dice lui: deve essere chi
trova il problema a valutare se è giusto e positivo che lo risolva lui. Il cartone di pizza lo
raccoglie perché ci mette dieci secondi; ma se ha trovato un bug nel sistema di pagamento che richiede
tre settimane di lavoro, e chi l'ha introdotto fa finta di niente, non deve risolverlo lui: deve far
presente che c'è un problema, che è stato ignorato e che va risolto (``pensateci voi, ma io vi dico
che c'è''). Va fatto caso per caso: se è immediato, facile e nel proprio interesse, lo si risolve.
Ma il relatore sconsiglia alle aziende la regola ``se lo trovi lo risolvi'', perché \textbf{disincentiva
a trovare i problemi}: qualcuno potrebbe non segnalarli per non avere altro lavoro e aspettare che li
trovi qualcun altro, ed è negativo quanto la situazione descritta. Per chiudere, ne discuterebbe con il
management: ``possiamo rivedere questa policy? Incentiva l'azione sbagliata''.

\paragraph{E se ci dicono di non mettere mano?} Uno studente che lavora come tester racconta che spesso
potrebbero correggere certe cose, ma viene loro detto di non toccarle. Esatto: si comunica. ``Ho trovato
questa cosa, ne siete consapevoli? Sì? Ci state già lavorando? No? Posso farlo io? Sì? Allora lo
faccio.'' Non bisogna esagerare con l'iniziativa: anche qui vale il \emph{default to trust}. Si parte
dal presupposto che qualcuno potrebbe averlo già visto e starci già lavorando, quindi si chiede e si
offre il proprio supporto.

% ---------------------------------------------------------------------
\section{Se si diventa manager}
\label{sec:z2-manager}

\subsection{Meriti, colpe e \emph{big picture}}

Una breve digressione per chi, un giorno, diventerà manager: che cosa cambierebbe? Oltre a tutto
quanto detto, ci sono il discorso su \textbf{meriti e colpe} e quello su \textbf{progetto e task}:
\textbf{delegare i meriti e prendersi le colpe}, e pensare al progetto più che ai singoli task, cioè
alla \emph{big picture}, invece che al ``voglio tornare a fare le cose che sapevo fare''.

È un altro problema dei manager promossi dal gruppo di chi sapeva fare bene un certo lavoro: quando
hanno l'occasione di tornare a farlo, lo fanno molto volentieri. ``C'è un task da fare, lo faccio io'',
magari con la scusa di dare l'esempio. No: si sta andando nella propria \textbf{comfort zone}. Il
relatore lo capisce (anche lui tornerebbe volentieri a fare ciò che faceva prima), ma bisogna accettare
che \textbf{quelle cose non le si deve più fare}.

\subsection{Il ruolo è più grande della persona}

Anche da manager, il ruolo che si ricopre per il gruppo è \textbf{più importante di sé}. Non importa
chi sia il community manager di Blizzard, di Riot o di un'altra azienda: conta ciò che rappresenta per i
giocatori in quanto community manager, non il fatto che sia lui piuttosto che un altro.

Prendersi cura del gruppo vuol dire applicare concetti molto semplici:
\begin{center}
  \begin{tabular}{@{}p{0.42\textwidth}p{0.48\textwidth}@{}}
    \toprule
    \textbf{Invece di dire\dots} & \textbf{\dots dire} \\
    \midrule
    ``Qui non si accettano errori.'' &
      ``Spiegami perché abbiamo sbagliato'': fammi capire perché è successo
      (sezione~\ref{sec:z1-default-trust}). \\
    \addlinespace
    ``Qui si fa come dico io.'' & ``Tu come faresti?'' \\
    \bottomrule
  \end{tabular}
\end{center}

\subsection{Il tetto della bravura del capo}

Il secondo punto è cruciale soprattutto quando si promuove il più bravo in un lavoro: se il criterio
noto è che si promuove il più bravo, il promosso lo sa, e la prima cosa che penserà è ``faccio lavorare
tutti come me, visto che sono il più bravo''. Il relatore chiede alla classe: che cosa si nega l'azienda
con questa dinamica? Anche ammesso che sia davvero il più bravo, se tutti fanno come lui diventeranno
tutti più bravi, no?

Il problema è che \textbf{si fissa come tetto massimo la bravura di quella persona}: se tutti fanno solo
come dice lui, si arriverà al massimo al suo livello, mentre ognuno ha il proprio metodo. Da qui quanto
detto il giorno prima sulla tendenza dei manager ad assumere solo persone peggiori di sé: se il criterio
di promozione è la bravura, si prenderanno sempre persone un po' meno brave per evitare che diventino il
proprio capo. Se invece le carriere sono divise (``se sei bravo fai questo, se vuoi fare il \emph{people
manager} fai quest'altro''), il manager è incentivato a prendere persone \textbf{migliori di sé}; e se,
come detto il giorno prima, il bonus dipende anche dalle performance del gruppo, è nell'interesse di
tutti che il gruppo sia pieno di superstar, un ulteriore incentivo ad assumere i più bravi.

\paragraph{L'opzione C.} C'è poi un'altra cosa, forse minima, ma per il relatore molto importante: si
perdono le \textbf{occasioni di innovare}. Nelle aziende precedenti era visto come il punto di riferimento
per il \emph{problem solving}: in un'emergenza, se c'era da scegliere tra A e B, ``decide Stefano''. Andava
benissimo, ma la tendenza era che a ogni emergenza venivano da lui. Così ci si negava \textbf{l'opzione C}:
quella a cui lui non avrebbe pensato e qualcun altro sì. Magari la prima volta sbaglia, la seconda anche,
e la terza ha un'idea migliore della sua, proprio perché lui non è più l'esperto di quel settore. Fissare
tutto come ``al massimo decide Stefano'' significa non accettare altre idee: un limite molto grosso
all'innovazione, all'aggiornamento e al fare le cose in modo più moderno e flessibile.

\subsection{Celebrare i successi e i fallimenti}

\paragraph{Celebrare i successi del team.} Il relatore vuole assicurarsi che tutti capiscano quando si è
ottenuta una vittoria o un successo: come detto il giorno prima, se non si sa che cosa si fa bene si
rischia di smettere di farlo. ``Questa cosa l'abbiamo fatta bene, complimenti, bravi tutti.''

\paragraph{Celebrare i fallimenti.} Sui fallimenti si era spinto ben oltre il semplice accettarli: aveva
indetto \textbf{meeting in cui si celebravano i fallimenti}, con le persone che raccontavano i propri errori
per distribuire quella conoscenza al resto del team. ``Non mi fa arrabbiare l'errore fatto, mi fa arrabbiare
lo stesso errore fatto quindici volte.'' Condividere un fallimento con il gruppo non significa ``guardate
questa persona sciocca che ha fatto questo errore'', ma ``questa persona ci sta permettendo di evitare lo
stesso errore in quindici'': sta salvando l'azienda dal ripeterlo quindici volte, perché la gente ha paura di
dire agli altri che ha sbagliato. Mostrare dove si è sbagliato diventava quindi un incentivo, un servizio
all'azienda. E si torna al punto di prima: se il proprio ruolo per il gruppo è più importante di sé, non
importa dover alzare la mano e dire ``ho sbagliato, ho fatto così'', perché gli altri penseranno ``grazie,
così io non lo faccio''. Si dovrebbe essere lieti che qualcuno abbia fatto quell'errore prima di noi, e
andare a ringraziarlo per averlo condiviso.

\paragraph{Condividere i successi.} Infine bisogna \textbf{dare il merito a chi lo ha}. Al relatore è
capitato spesso di ricevere ringraziamenti come ``grazie a Stefano e al suo team per aver fatto questa
cosa'', e la sua risposta è stata ``grazie, ma ha fatto tutto Gianni''. Non bisogna avere paura di dire
``grazie del feedback, ma è stato lui'': è giusto ringraziare lui e il team (il relatore si sente parte del
team, chiaramente). Non si ottiene niente dal dire ``grazie, riferirò i complimenti al team'': il relatore
non ha mai creduto nel ``riferire'' che il management ha detto ``ottimo lavoro''. Chiede al management di
dirlo \textbf{direttamente} al team, perché sentirlo dire dal management ha un impatto diverso che sentirlo
riferire da lui. In questo senso, osserva uno studente, ogni gruppo è il gruppo di sostegno di se stesso:
``assolutamente sì''.

\subsection{Quanto ci vuole?}

Non c'è una scadenza. Il rispetto non si ottiene pagando per fare prima, ``con il \emph{boost} a livello
60''. Servono tempo ed energia, ed è difficilmente misurabile: il relatore non può dire quanto il team si
fidi di lui in questo momento. Lo percepisce da come lavorano, ma lo si vedrà davvero solo quando ci sarà un
problema grosso; fino ad allora è tutta una sua percezione. Il team lo sente ripetere queste cose ogni giorno
e lo vede metterle in pratica ogni giorno, e lui può solo sperare che il messaggio passi; quando ci saranno
davvero ``gatte da pelare'' capirà come è messo, perché non c'è momento migliore di una difficoltà per vedere
come reagisce un team.

\textbf{Essere leader è prendersi un impegno a lungo termine verso i propri colleghi}: non ci sono modi per
velocizzarlo. Bisogna fare il proprio lavoro ogni giorno ``a testa bassa'', aiutando gli altri, facendolo
vedere ed \textbf{esplicitando perché lo si fa}. Il relatore lo fa quotidianamente: ``ti sto aiutando; nel
caso avessi un dubbio, lo faccio perché il mio lavoro è aiutarti e io ne traggo vantaggio dal fatto che tu
lavori meglio. Sono qui per te, ma lo faccio anche per me: non pensare che sia solo gentilezza. È il mio
lavoro, ma sappi che lo faccio spontaneamente''.

\subsection{L'ultima slide}

L'obiettivo dovrebbe sempre essere lo stesso: tutti dovremmo poter finire la giornata, staccare e tornare a
casa sentendoci
\begin{itemize}
  \item \textbf{ispirati}, perché abbiamo fatto qualcosa che ha un effetto tangibile;
  \item \textbf{sicuri}, perché sappiamo di avere colleghi, un leader e dei manager che si prendono cura di
        noi;
  \item \textbf{realizzati}, per quello che facciamo e perché è ciò che ci piace fare.
\end{itemize}
E questo dovrebbe essere \textbf{un diritto naturale di tutti}, non un ``sì, ma solo il manager'', ``sì, ma
solo se sei un esperto'', ``sì, ma solo se sei un game director''. Ognuno ha le proprie capacità: bisogna
solo trovare ciò che piace e che si vuole fare.

% ---------------------------------------------------------------------
\section{Domande e risposte finali}
\label{sec:z2-qa-finale}

Finita la presentazione, si passa alle domande: ``chiedete tutto quello che volete, perché è
l'ultima possibilità''. Le domande sono riportate per temi.

\subsection{Leadership e organizzazione del lavoro}

\paragraph{La leadership riguarda tutti o solo chi vuole diventare leader?} Il discorso sui
manager vale solo per i manager; ma il consiglio è \textbf{essere tutti leader del proprio gruppo,
da dentro}. Essere un leader significa sapere che gli altri possono fare affidamento su di sé, e
di queste figure in un team ce ne possono essere infinite: non è una gara a chi è il leader.
``Siate tutti una di quelle figure a cui gli altri guardano quando sono in difficoltà.''

\paragraph{Da manager non si hanno più orari e si fanno straordinari di continuo?} No: il relatore
ha notato di essere molto più efficiente come manager quando si dà \textbf{paletti molto precisi}.
Se stacca alle sei e alle sei e cinque arriva un messaggio, risponde ``se non mi dici che è
urgentissimo, non lo guardo nemmeno, lo leggo domani''; se non è urgente, lo guarda il giorno dopo.
Ma gestisce le aspettative: ``se è urgente ci sono e ci sarò sempre, ma non aspettatevi che faccia
orari eccessivi solo perché sono manager, così come non me lo aspetto da voi''. Si aspetta che
ognuno faccia ciò che deve; se è urgente tutti faranno di più, ma se non lo è, no.

\paragraph{L'orario davvero flessibile.} Qui il relatore aggiunge una cosa che si era dimenticato di
dire parlando di ferie e permessi illimitati: introduce anche l'\textbf{orario davvero
flessibile}. Non il flessibile di molte aziende, in cui si può iniziare alle 9 e finire alle 18
oppure iniziare alle 10 e finire alle 19: se oggi si vuole iniziare alle 10 e finire alle 15, e si è
fatto tutto quello che si doveva, va bene così. Ci si prendono le ore libere in più per altro, per
staccare, per riposarsi, sapendo che se la settimana dopo bisognerà lavorare dalle 9 alle 21, lo si
farà (o si dirà ``ieri ho fatto dalle 9 alle 21, domani faccio dalle 10 alle 15''). L'obiettivo dato
è un lavoro e un progetto, non un orario: si lavora per completare i propri task e rispettare le
deadline. L'orario indicato è solo indicativo, ed è quello dei meeting: \textbf{nessun meeting prima
delle 10 o dopo le 17}. Se poi qualcuno vuole lavorare dalle 10 alle 11 e poi di notte, perché quel
giorno preferisce lavorare di notte, più rilassato, faccia pure: l'importante è che trovi il tempo per
riposare e staccare. Ci si autogestisce senza timbrare cartellini; l'unico accordo è ``hai un task, va
consegnato entro il giorno X e mi garantisci che venga fatto; come ci arrivi non mi interessa''.

\paragraph{Le è mai capitato di essere aiutato da qualcuno, in un momento di difficoltà?} Sì: come
detto, questo comportamento è \textbf{contagioso in senso positivo}. Se ogni volta che qualcuno è in
difficoltà qualcun altro lo aiuta, prima o poi capita anche a chi aiuta. Il relatore non lo dice
``tanto per dire'': ha avuto momenti in cui faceva davvero fatica ad alzarsi dal letto per lavorare,
perché esausto o demotivato per altri motivi, e il team gli ha detto ``stai a letto, ci pensiamo noi
oggi''. Questo lo aiuta ad aiutare, perché sa che la cosa è bidirezionale.

\paragraph{Tra le soft skill del workshop, da quale partire?} \textbf{Il feedback}: saperlo dare e
ricevere. È la più difficile, molto più del ``voglio aiutarti'', ma è anche quella su cui si costruisce
tutto il resto: attraverso il feedback si costruisce la \textbf{fiducia}, attraverso la fiducia il
\textbf{rispetto}, e il rispetto diventa la base portante di tutto il resto.

\paragraph{E se il leader se ne va?} A una domanda su che cosa fare quando viene meno la figura di
riferimento del gruppo\footnote{La domanda non è leggibile nella trascrizione; dalla risposta si
capisce che riguardava il caso in cui il leader che incarnava la visione del gruppo se ne va.}, il
relatore risponde che la domanda da farsi è un'altra: \textbf{è cambiato l'obiettivo?} La direzione è
sempre la stessa? Se sì, chi è la persona più adatta a portarci lì? Non conta tanto chi guida, ma
l'obiettivo che ci accomunava, di solito la visione di un leader in cui tutti si allineano: ci
crediamo ancora tutti, o è cambiato qualcosa? È il momento di fare autoanalisi: ``ero qui per quella
persona o per la visione?''. Si dice, ed è vero, che \textbf{le persone non si dimettono dai posti di
lavoro, ma da capi che non rispettano più}, perché spesso l'unica interfaccia con l'azienda è il
proprio capo: le dimissioni non sono un messaggio per l'azienda, ma per il capo (``non voglio più
lavorare con te''). Allo stesso modo, se non c'è più il leader della visione, bisogna prima capire se
si condivide ancora qualcosa.

\paragraph{Lavorare così è più faticoso?} Sì e no. ``Non voglio demotivarvi, perché non è il mio
lavoro'', ma non è tutto rose e fiori: lavorare in questo modo è anche più stancante
\textbf{emotivamente}, perché non si può più pensare ``faccio il mio, alle cinque ho finito e chi si è
visto si è visto''. Si sente un senso di responsabilità verso gli altri. Succede che si finisca il
proprio lavoro del pomeriggio, ci si giri e si veda un collega che finirà alle nove di sera, e
l'istinto sarà ``posso farti un pezzo del lavoro?'': così si finisce entrambi alle sei. In un certo
senso ci si accolla più lavoro, con la consapevolezza che quando si sarà in difficoltà ci sarà qualcuno
che si girerà a dire ``ti vedo stressato, posso darti una mano?''. È bellissimo, ma non dà solo benefici
diretti: richiede anche un investimento personale che non sempre le persone vogliono fare. Pro e
contro: è il modo che sceglierebbe il relatore, ma non è per tutti. Conosce persone che vogliono poter
staccare appena finito, e va benissimo: si cerca un'azienda in cui si possa lavorare così.

\subsection{Condividere fallimenti e successi in pratica}

\paragraph{Come si celebrano i fallimenti e si condividono i successi?} Condividere i successi sembra
più semplice (``bravo Mario che ha fatto questa cosa''). Il relatore dà uno spunto ancora migliore:
\textbf{``bravo Mario, e ora Mario ci racconta come l'ha fatto e com'è andata''}. Sono sottigliezze,
ma contano: raccontare in prima persona ciò che si è fatto dà un moto d'orgoglio che non dà sentirlo
raccontare da un altro.

Per i fallimenti, che sembrano dover partire dalla persona stessa, l'obiettivo era che tutti li
condividessero spontaneamente; ma all'inizio era difficile, perché non viene naturale alzare la mano e
dire ``ho sbagliato''. Per questo il relatore ha creato artificialmente dei \textbf{``festival del
fallimento''}, dal tono chiaramente ironico: durante una pausa pranzo, mentre tutti mangiavano insieme,
a turno si saliva ``sul palco'' a raccontare la stupidaggine fatta in settimana. Molto spesso iniziava
lui, raccontando i propri errori (``ho sbagliato questa cosa, sappiate che non si fa così''). Dopo aver
forzato l'interazione per una, due, tre settimane, alla quarta lo facevano da soli: era diventato
organico dire a tutti ``ragazzi, attenzione, ho fatto questa sciocchezza, non fatela anche voi'', e c'era
perfino un \textbf{tag apposito su Slack}. È stato il team a decidere come condividere i fallimenti: come
sempre, invece di imporre il modo, il relatore ha lasciato scegliere a loro, e funziona. È impegnativo
tanto quanto lo si vuole rendere tale.

\paragraph{Che tipo di fallimenti si condividono? Anche quelli personali?} Di solito si parte da quelli
\textbf{professionali}, cioè da ciò che si riconosce come utile anche agli altri: qualcosa che ha
rallentato, fatto perdere tempo o che rischia di ripetersi (``sappiate che, per evitare che lo facciate
anche voi\dots''). Il resto è più da conversazione privata, \emph{one to one}, con la persona con cui ci
si vuole sfogare, o con gli amici: se non ha un potenziale impatto sugli altri, non c'è un incentivo a
dirlo. E un errore che si è fatto per il proprio tipo di personalità? Qui tornano i test di personalità
della prima giornata: conoscendo il proprio profilo e quello dei colleghi, si sa che qualcuno ha lo
stesso carattere, e quindi lo si aiuta.

\subsection{Cultura aziendale e formazione}

\paragraph{Come si arriva a poter fare queste cose in azienda?} Conta molto la cultura aziendale. Il
relatore sa di lavorare bene perché oggi ha il \emph{background} e i successi alle spalle per dimostrare
che il suo metodo funziona, e quindi glielo lasciano fare; ma se dieci anni prima avesse proposto le
stesse cose, ``mi avrebbero mandato a quel paese''. C'è un momento e un posto per fare le cose, e non è
sempre possibile fare ciò che si ha in testa; ma ci si può lavorare. Per lui la domanda è diventata:
``sono convinto che così sarebbe meglio; come faccio ad avere carta bianca?''. \textbf{Portando
risultati}: si è concentrato su quello, e ci è arrivato progressivamente.

\paragraph{Come funziona la formazione con workshop come questo, in azienda?} Il relatore li fa in modo
un po' più rilassato ed esteso, con molte più interazioni ed esempi, perché in azienda si possono usare
esempi che si hanno davanti ogni giorno, e parlare di problemi e discussioni diventa molto più reale. Ma
in generale è sempre \textbf{formazione pratica}, ripetuta finché il messaggio non è recepito. Il segno
che qualcosa è stato recepito è quando arriva una persona nuova e \textbf{non è lui a spiegarle le cose}:
è il team stesso a spiegarle che cos'è la feedback culture, e lui non deve parlare. È la conferma che il
concetto viene perpetuato dal team.

\paragraph{Quanti manager dell'industria condividono questi principi?} Il relatore ha fatto molti
interventi all'IGDA di Roma e ne farà uno a Milano la settimana successiva (``il 24, mi pare''); interviene
spesso a workshop ed eventi, e il riscontro è sempre ``che bello, vorremmo farlo anche noi in Italia'', e ci
stanno lavorando. La moglie lavora in un'azienda che non ha niente a che fare con i videogiochi e lo stanno
già applicando. Magari in Italia è tra i primi a parlarne, perché il paese è ancora un po' indietro, ma non è
assolutamente l'unico: è più diffuso di quanto si pensi, non è una cosa di nicchia.

\paragraph{Questi spunti si applicano anche all'IT e alla consulenza?} ``Praticamente tutti.'' Si applicano
facilmente quando c'è un gruppo di persone su cui applicarli; se si cambia continuamente cliente è più
difficile portare questa mentalità da loro, ma questo non impedisce a chi lavora con noi di applicarla
direttamente con noi.

\paragraph{Ci si può tenere in contatto?} Su LinkedIn: il relatore accetta sempre tutti e risponde alle
domande quando riesce (ne riceve tante, ma cerca di stare al passo almeno su base settimanale). E, ``a costo
di sembrare quello che si fa pubblicità'', chi lavora in un'azienda che fa fatica ad applicare questi
concetti o non sa che esistono può metterla in contatto con lui: è spesso così che ha trovato nuovi modi di
applicarli. Vale anche per chi tra sei mesi lavorerà altrove. I suoi profili sono pubblici e non c'è nessun
problema a contattarlo.

\subsection{Performance e persone in difficoltà}

\paragraph{Se un'azienda è a corto di personale, conviene comunque eliminare chi performa poco?} Il giorno
prima si era detto che conviene separarsi presto da chi non performa (sezione~\ref{sec:z1-squadra}); ma non
si finisce per caricare di lavoro tutti gli altri? Per il relatore \textbf{eliminare non vuol dire non
sostituire}: vuol dire cambiare la persona con qualcuno magari più junior, ma più adatto come \emph{fit}
culturale. Il vecchio discorso del ``fare lo stesso lavoro con meno gente'' non è sostenibile: lo si può fare
fino a un certo punto, in caso di necessità, ma non è una soluzione. Non è ottimizzare, è far andare in
\emph{burnout} chi resta.

\paragraph{Ha più senso tenere una persona poco performante e non in sintonia con il team, piuttosto che
licenziarla senza poterla sostituire?} No: se una persona è poco in sintonia, performa poco ed è in generale
tossica e negativa per il team, \textbf{prima ci si separa, meglio è}. Questo non vuol dire non provare a
recuperarla, ma bisogna andare dritti al punto: ``abbiamo un problema: ti vedo poco performante, non ti vedo
allineato con il team, vedo che il team non si fida di te. Che cosa vogliamo fare?''. Se la risposta è ``non lo
so'', la strada porta dritta all'uscita. Bisogna \textbf{recidere i rami morti} il prima possibile, facendo
però attenzione a essere sicuri che si tratti di un ramo morto e non di un ramo recuperabile: tagliare un ramo
recuperabile mette negli altri la paura di essere i prossimi, ``alla prima foglia gialla''. Quello che si
vuole è che, quando un ramo è morto, la pianta smetta di alimentarlo, e ne crescerà un altro. Entrambi i
messaggi devono essere chiari: \textbf{non si subisce passivamente la situazione}, e allo stesso tempo
\textbf{non si abbandona nessuno} (se qualcuno è in difficoltà, la prima soluzione non è tagliare il ramo). È
un compromesso tra le due cose. (A questo punto arriva, puntuale, il crash della connessione: ``non era ancora
successo oggi''.)

\paragraph{Che cosa fare se qualcuno nel team sembra soffrire di depressione?} È un discorso complesso, perché
il relatore \textbf{non è uno psicologo}. Può dare la sua disponibilità a concedere permessi, coprire il lavoro
della persona e in generale darle modo di seguire un percorso con uno specialista: non può risolvere la
depressione, ma può alleviarne la parte lavorativa e permettere di ricevere supporto da chi è qualificato. È lo
stesso motivo per cui, a Blizzard e a Riot, quando arrivavano ticket con minacce di suicidio (``mi avete
bannato, quindi mi ammazzo'', ticket purtroppo abbastanza comuni), il \emph{customer service} era istruito a
non intervenire direttamente: non dovevano essere loro a fermare la persona, perché non qualificati, ma
indirizzarla a qualcuno di qualificato, prendendo tempo. Una persona non qualificata che prova a farlo rischia
di peggiorare la situazione; come davanti a un incidente si chiama l'ambulanza e non si fa una tracheotomia.
Per il relatore salute mentale e salute fisica sono la stessa cosa: come non si azzarderebbe a operare qualcuno
di appendicite, non si azzarda a fargli da terapeuta.

\paragraph{Come si danno feedback adatti a ogni tipo di personalità?} Bisognerebbe leggere i risultati del
profilo della persona sullo stesso sito su cui si è fatto il test di personalità della prima giornata, che
spiega anche come interagire con quel profilo. In generale: con chi è più \textbf{analitico}, portare dati a
supporto; con chi è più \textbf{emotivo}, metterlo in una condizione di comfort; con chi è molto
\textbf{pratico}, andare dritti al punto senza girarci intorno. Poi fare pratica. Oppure il modo ``perfetto''
per il relatore: provarci come viene e poi \textbf{chiedere} ``andava bene così o preferisci in un altro modo?
Potevo fare meglio?''. Bisogna creare anche quel rapporto: fidarsi di discutere anche di queste cose.

\paragraph{Se nel team fossimo quasi tutti leader, non ci sarebbero ``troppi galli''?} No, perché per leader si
intende chi si prende cura degli altri, e magari lo fossero tutti. Non significa che tutti scelgano che cosa
fare: spesso lo sceglierà comunque il manager, o guiderà la scelta. Ma avere tutte persone con iniziativa e
capaci di scegliere quando serve è qualcosa per cui il relatore ``metterebbe la firma''. Non ha mai corso il
rischio dei ``troppi galli'': se ci sono quei caratteri, il team è stato assemblato male, perché non tutti sono
compatibili con questo modo di lavorare e magari non si hanno le persone giuste per applicarlo.

\paragraph{Se la fiducia richiede tempo ed esperienza, quanto è difficile per un junior?} Meno di quanto si
pensi: proprio perché su un junior non ci sono aspettative, un junior che chiede ``posso darti una mano?'' dà un
segno chiaro che tiene più al team che a se stesso. Si può essere leader anche non dal punto di vista
professionale, ma come chi va a dare supporto agli altri: ``il medico del team, l'\emph{healer}'', e lo si può
fare anche da junior. È così che ha iniziato il relatore: andare ad aiutare a risolvere problemi non suoi gli ha
fatto accelerare molto il percorso per ottenere rispetto e stima, e lo ha esposto a cose a cui altrimenti non
sarebbe mai arrivato. È una questione di carattere: se ci si tiene agli altri, lo si fa vedere, e il resto
arriva come risultato. ``È quasi un reddito passivo'': un investimento nel far vedere che ci si è, e il resto
arriva (``Stefano si è preso cura di me, quindi ora mi prendo cura di lui'').

\paragraph{Aiutare gli altri solo dopo aver finito il proprio lavoro?} Dipende dall'urgenza. Il relatore spesso,
se qualcosa è davvero urgente, mette in pausa ciò che sta facendo, risolve e poi torna al suo lavoro; ma è troppo
\emph{case by case}. Se due colleghi litigano su chi debba avere le ferie di Natale, lo si può risolvere anche
domani; se litigano perché si è perso il \emph{repository} e domani c'è la consegna, lo si risolve subito.

\paragraph{Come mantenere motivato un team su progetti non retribuiti, come quelli accademici?} Essendo il più
\textbf{trasparenti} possibile nei processi decisionali: se si deve decidere se fare un action o un RPG, tutti
devono sapere perché si fa l'uno invece dell'altro (da qui il metodo di brainstorming della mattina). Deve essere
sempre chiaro a tutti perché si stanno facendo le cose in quel modo, perché \textbf{la demotivazione nasce dal
disallineamento} rispetto all'obiettivo del gruppo, e l'unico modo di essere allineati è essere sempre
consapevoli e partecipi del processo che decide dove sta andando il gruppo. E soprattutto: \textbf{non chiudere
mai un meeting senza aver concordato che cosa si vuole fare} e senza essere tutti d'accordo. Se non lo si è,
bisogna continuare a parlarne, per quanto scomodo, lungo e noioso: si deve sempre uscire da un meeting
sapendo che si è tutti d'accordo, o al massimo in \emph{disagree and commit} (sezione~\ref{sec:z2-disagree}).
``Quello è il massimo accettabile.''

\subsection{Carriera, colloqui e candidature}

\paragraph{Come ha iniziato a lavorare all'estero?} Negli annunci online si trovano spesso requisiti come un titolo
di studio conseguito nel paese o l'abilitazione a lavorarci. Questo vale per i paesi fuori dall'Unione europea: per
lavorare negli Stati Uniti o in Sud America serve un visto o una \emph{green card}. In molti casi, però, è proprio
ottenere il lavoro a permettere di lavorare lì: quando il relatore doveva andare a lavorare in California (a Irvine,
a Santa Monica, dove capitava), era l'azienda a occuparsi dei permessi necessari, perché ``per venire qui devi avere
un lavoro, e il lavoro te lo diamo noi''.

Il suo inizio è stato abbastanza semplice. In Italia lavorava per Digital Bros, occupandosi dei giochi online gestiti
dall'azienda e di quello che era il canale Game Network sul digitale terrestre\footnote{Questo passaggio chiarisce il
dubbio segnalato nel capitolo~\ref{cap:z1-feedback} sul ``Digital Game Network Channel'': si tratta di Game Network,
canale televisivo dedicato ai videogiochi legato al gruppo Digital Bros.}. Vide un annuncio a Dublino per
\emph{Warhammer Online}, mandò la candidatura, fece i colloqui e gli dissero ``ti paghiamo la \emph{relocation}, vieni
qui''. Come detto il giorno prima, sono le aziende a occuparsi del trasferimento: se devono creare a Dublino un'unità
che segue dodici lingue europee, sanno di dover far arrivare persone da fuori, ed è già a budget. Non è nemmeno una
richiesta da fare: è prevista. Per il relatore, da questo punto di vista, è indifferente lavorare per un'azienda
italiana o estera: se si è adatti, si è adatti. Per questo insisteva sull'inglese: se lo si parla, è indifferente di
che cosa ci si occupi e quale sia la propria lingua madre, a meno che la posizione non sia esplicitamente ``community
manager italiano''. Un programmatore polacco o italiano non cambia niente. Lavorare all'estero non è più difficile: c'è
più concorrenza solo se l'azienda cerca attivamente persone da tutto il mondo e non si pone problemi di \emph{location}.
Cambiano l'azienda, il fatto di lavorare da remoto o dal vivo, l'eventuale trasferimento; per il resto non cambia
molto.

\paragraph{Gli anni passati nell'IT contano se si vuole passare ai videogiochi?} Dipende da che cosa si vuole fare: se
da responsabile IT si vuole diventare programmatore, bisogna saper dimostrare di saper programmare. Vale anche per i
passaggi interni: se un game designer vuole diventare art director, al relatore va bene, deve solo dimostrare di saper
fare l'art director. Anzi, incentiva le \textbf{promozioni interne}, perché si conosce già il candidato; ma gli
spostamenti non sono liberi ``a sentimento'': bisogna saper fare quella cosa. I recruiter valuteranno non tanto gli anni
quanto ciò che si sa fare; gli anni servono comunque, perché dimostrano di saper lavorare in un ambiente professionale,
che è un altro insieme di competenze. Ma se si sa fare ciò che chiedono, aver lavorato nell'azienda X o Y non dovrebbe
cambiare molto.

\paragraph{Quanto vengono considerate le candidature spontanee?} Non tantissimo. Soprattutto per le aziende non grandi,
se non stanno cercando qualcuno è perché non se lo possono permettere: una candidatura spontanea non le convincerà ad
assumere. Più l'azienda è piccola, più conviene aspettare l'annuncio, perché se non c'è vuol dire che non c'è budget, e
potrebbero rispondere ``ti vorremmo, ma non possiamo prenderti''. C'è anche il rischio che finisca nel vuoto (l'HR magari
non controlla nemmeno la mail, perché non sta assumendo) o, peggio, che rispondano ``non ci servi'' e una settimana dopo
esca la posizione, con il candidato già nel cestino dei rifiutati. Si rischia di \textbf{bruciarsi} la possibilità.

\paragraph{Come ci si accorge di essere riconosciuti come leader, e come lo si ``vende'' a un colloquio?} Uno studente
osserva che la leadership, così descritta, non la si riconosce da sé ma la riconoscono gli altri. Internamente il
relatore non ha mai dovuto ``venderla'', perché parlavano i risultati. Finché non lo si vive è difficile capirlo, ma
quando si arriva alla posizione in cui il team rispetta e vede qualcuno come punto di riferimento, non ci sono grandi
dubbi: è evidente che vengono da lui. Si può riflettere su \textbf{quanto spesso gli altri vengono da noi}, cercano di
emulare il nostro modo di fare, vanno a rivedere le nostre cose per capire dove hanno sbagliato, chiedono un parere, ci
parlano di un problema: la somma di tutto questo significa molto spesso che si fidano. Per ``venderla'' all'esterno si
può dire di essersi occupati di un team in cui, attraverso il lavoro e il rispetto degli altri, si è arrivati a essere
visti come punto di riferimento; ma il suggerimento è \textbf{portare esempi tangibili} (``c'era questo progetto, non
sapevamo come risolverlo, il team mi ha chiesto una mano e l'ho risolto così''), invece di ``incensarsi'' dicendo ``sono
il punto di riferimento del team''. Se il relatore racconta di essere stato eletto rappresentante dei licenziati senza
essere tra loro, un'azienda capisce che si fidavano di lui, punto: non ci sono altre conclusioni da trarre.

\paragraph{Quali sono gli errori più gravi dei junior ai colloqui, a livello di soft skill?} In ordine di priorità:
\begin{enumerate}
  \item \textbf{non arrivare preparati}, cioè \textbf{senza esempi}. Il relatore ha sempre un esempio pronto per ogni
        cosa di cui vuole parlare: per la fiducia, per gli \emph{achievement}, per i risultati. Servono sempre esempi
        tangibili per far capire che cosa si intende;
  \item \textbf{non avere domande alla fine}, perché dimostra di non avere chiaro che cosa si andrebbe a fare. Anche solo
        chiedere che cosa si farebbe, che cosa si può migliorare, o ``che impressioni hai avuto finora dal mio
        colloquio?'': si possono fare domande molto aperte al recruiter per capire come sta andando;
  \item \textbf{non leggere bene la \emph{job description}}: presentarsi con vaghezza, senza aver capito che cosa cerca
        l'azienda. Lo si impara anche con la pratica; ma il designer che si candida per una posizione da programmatore,
        senza nulla nel CV che faccia pensare che sappia anche programmare, il relatore non solo lo scarta: lo mette in
        una \textbf{blacklist} di persone che forse non considererà mai. Se non si riesce a interpretare un testo come una
        \emph{job description}, è un problema troppo grosso da superare in una collaborazione;
  \item \textbf{errori di battitura e CV impaginati male}, che molti sottovalutano. Il relatore ha letto letteralmente,
        in una lettera di presentazione per una posizione da \emph{copywriter}, la frase ``sono una persona molto attenta
        ai dettaggi''. Come si può considerare quella persona per quel ruolo? Non ha nemmeno incollato la lettera in Word
        per il controllo automatico: non è un errore sfuggito al correttore perché la parola esiste con un altro senso,
        ``dettaggi'' viene segnato in rosso in qualsiasi lingua. Come affidarle la comunicazione aziendale? Vale per tutti
        gli errori: ciò che si manda è \textbf{il primo biglietto da visita} con l'azienda, e non si manda un biglietto da
        visita sgrammaticato. Un conto è non andare ai colloqui in giacca e cravatta (``non l'ho mai fatto e probabilmente
        non lo farò mai''), un altro è presentarsi con errori;
  \item un altro esempio comune: \textbf{l'email sbagliata nel CV}. Se la mail torna indietro, il relatore non telefona
        per farsi dare quella corretta: se non si è in grado di scrivere giusta la propria email, non vale la pena. Più
        l'errore è banale e facilmente evitabile, più lo fa arrabbiare (``scusate se mi sono inalberato''): evidentemente
        il candidato non teneva abbastanza a quel posto.
\end{enumerate}

\paragraph{Quanto è importante la lettera di presentazione?} Sempre di più: sempre più aziende la richiedono come
obbligatoria, quindi conviene prepararne una, sapendo di doverla adattare al destinatario. Deve essere \textbf{meno lunga
di una pagina}, in \textbf{tre paragrafi}:
\begin{enumerate}
  \item \textbf{chi siete}: un minimo di \emph{backstory}, che cosa avete fatto, che cosa volete fare, da dove arrivate;
  \item \textbf{perché avete mandato il CV}: che cosa vi ha interessato di questa azienda;
  \item \textbf{perché dovrebbero chiamare voi}: che cosa dovrebbe spingere l'azienda a scegliere proprio voi.
\end{enumerate}
Uno studente dice di riuscire bene con i primi due, ma di trovare difficile il terzo, che è poi la domanda che tutte le
aziende fanno. In realtà basta \textbf{collegarsi alla \emph{job description}}: se cercano un 3D artist in grado di fare
X, Y e Z, si conferma di averlo letto e si spiega che con la propria esperienza si è imparato a fare proprio quelle cose.
Non vogliono sentire niente di particolare, ``non devi cambiargli il mondo''.

\paragraph{Che cosa guarda un recruiter in un portfolio?} Il relatore li guarda poco: fa un \emph{sanity check} generale,
per vedere che non sia evidente perfino a lui, ``a digiuno di tante cose tecniche'', che c'è un problema (se se ne accorge
perfino lui, il problema c'è). Per il resto guarda che il portfolio sia \textbf{coerente con il CV}: se arriva un CV da
senior con dodici progetti completati e anni di lavoro, e il portfolio sono quattro immagini, non è coerente. Vuole una
concordanza tra ciò che si racconta e ciò che si vede.

\paragraph{Che strumenti per creare un CV?} \textbf{Canva}, che ha un catalogo di template gratuiti e a pagamento
amplissimo e personalizzabile: anche se non piace, lo si prende e lo si adatta, ma fa già il 90\% del lavoro. Insieme a
\textbf{ChatGPT}, perché ``è inutile far finta che non esista'': ci si fa scrivere la lettera di presentazione e poi la si
corregge e la si adatta, e lo stesso per il CV. L'importante è renderli adatti e coerenti, non lasciarli solo frutto di
ChatGPT; ma per fare il grosso del lavoro è un ottimo strumento.

\paragraph{Conviene lasciare nel CV esperienze non del settore?} Uno studente ha alle spalle il liceo, vari stage e
lavori che non c'entrano niente con la modellazione 3D. Si possono mettere finché c'è spazio: il relatore consiglia
\textbf{al massimo due pagine}, e quando ha cominciato a sforare nella terza ha iniziato a togliere ciò che non era più
rilevante. Ormai nel suo CV la formazione è un paragrafo: ha così tante esperienze lavorative che non interessa a nessuno
che abbia fatto il liceo scientifico o classico, e non lo mette nemmeno più. Si filtra in base a ciò che è più rilevante
per lo spazio disponibile.

\paragraph{Che cosa fare con le aziende che dicono ``ti faremo sapere'' e poi spariscono?} Non aspettare uno o due anni
per scrivere. Il suggerimento: tra le domande finali del colloquio, l'ultima del relatore è sempre \textbf{``tra quanto
mi fate sapere?''}. Se rispondono dieci giorni, all'undicesimo scrive: ``mi avete detto dieci''. Non si sente in difetto a
mettere pressione, perché sono stati loro a dare la scadenza. E se viene ignorato del tutto, quell'azienda per lui è
``morta'': non si ricandiderà mai, perché non vede perché il rispetto debba andare in una sola direzione.

\paragraph{Quanto essere trasparenti, nel CV e nella lettera, su problemi personali come l'ADHD?} Uno studente con
ADHD\footnote{Nella trascrizione ``essendo 10D''; nel resto della risposta ``di'' e ``dei 10D'' stanno per ADHD.} dice di
tendere ad avere performance altalenanti. Il relatore lo invita a \textbf{non usare l'ADHD come scusa} per essere
altalenante: anche lui ha l'ADHD, ma non è altalenante. Esserlo è senz'altro un \emph{malus}, ma esistono i mezzi per non
avere problemi di questo tipo: lui ha trovato i suoi metodi (la to-do list, la \emph{must do list}, il delegare il lavoro
al team, perché occuparsi delle cose per gli altri lo mantiene molto motivato, sveglio, all'erta, attento a come vanno le
cose). Sta anche a ciascuno trovare un lavoro in cui compensare ciò che viene richiesto: se domani lo rimettessero a fare
il programmatore, come all'inizio, probabilmente impazzirebbe, perché non riusciva a restare concentrato sul codice oltre
un certo tempo. Lì ha capito che non poteva fare il programmatore, ma ha trovato qualcos'altro che gli permette di essere
com'è e fare ciò che fa. Bisogna capire i propri punti di forza e di debolezza e giocare su quelli: si possono avere tutti
i difetti del mondo e voler comunque fare una cosa, ma bisogna accettare che sarà più difficile.

Quanto alla trasparenza, il consiglio è \textbf{non mentire mai nel CV}, ``perché tanto vi beccano''. E per il discorso della
\emph{bandwidth}: tutte le energie spese a mentire, e a tenere viva la bugia, sono energie tolte allo spiegare bene ciò che
si sa fare. È un investimento nel posto sbagliato. Bisogna concentrarsi su ciò che si sa fare, con esempi tangibili, e dire
``non lo so fare'' quando non lo si sa fare. Il peggio non è essere scoperti, ma \textbf{non esserlo}: se si finge di sapere
il C++, si viene assunti e messi a lavorare in C++, che cosa si fa il giorno dopo? Si rischia davvero di rovinarsi la
carriera. Se non si è un \emph{fit}, si sarà il \emph{fit} per la prossima. L'incubo del relatore è essere assunto in un
posto dove pensano che sappia fare qualcosa che non sa fare; anzi, tende a esplicitarlo: ``cercate questa cosa? Io non la so
fare, va bene lo stesso?'' --- ``Sì, te la insegniamo noi'' --- ``Voglio solo che sappiate che non la so fare''. Meglio
questo che vivere con l'ansia che prima o poi si accorgano che non sa programmare.

Gli capita spesso con chi mente sull'inglese. Al primo colloquio chiede ``sai parlare inglese? A che livello?'' --- ``C1,
\emph{fluent}, quasi madrelingua'' --- ``Perfetto'', e lascia cadere l'argomento. All'ultimo colloquio: ``ti va se lo facciamo
in inglese?'' --- ``No, ma io intendevo che\dots'' --- ``No, mi avevi detto \emph{fluent}: ora facciamo il colloquio in
inglese''. Prima o poi le bugie raggiungono chi le dice.

% ---------------------------------------------------------------------
\section{Chiusura del workshop}
\label{sec:z2-chiusura}

Il relatore chiede alla classe di far sapere com'è andata a Vanessa, il punto di contatto con l'accademia per il workshop:
sarà lei a riferirgli eventuali lamentele o suggerimenti, che gli permetteranno di fare meglio i prossimi. Ringrazia per aver
``retto per dodici ore'' e spera che sia stato utile, consapevole che tutti sono stanchi quanto lui.

Gli studenti rispondono che è stato chiarificatore su molte cose che non sapevano, soprattutto su come interfacciarsi nei
team, e che sarà materiale da rivedere per fissare bene i concetti, vista la mole in due giorni; qualcuno si aspettava che
sarebbe stato molto più pesante e l'ha trovato invece scorrevole. Il relatore ribadisce che \textbf{niente sostituirà
l'esperienza diretta}: il workshop dà un'infarinatura e un punto di vista, ma i concetti si metteranno in pratica solo nel
\emph{day to day}. Bisogna avere \textbf{pazienza con se stessi}, sapendo che non verrà tutto bene dal primo giorno.
Confessa di essere stato un po' preoccupato per la durata, ma il materiale c'era, e farlo in un giorno avrebbe rischiato di
trasformarlo in una corsa: ``rischiava di essere uno \emph{sprint}, e invece è una maratona''.

Infine ricorda che si può aggiungerlo su LinkedIn, che registrazione e slide (in forma integrale) verranno fatte avere
tramite Vanessa (``dite che vi ho autorizzato io a chiederle'') e che, per qualsiasi dubbio, resta a disposizione. ``Buon
resto del weekend.''

% ---------------------------------------------------------------------
\section*{Punti chiave del capitolo}
\addcontentsline{toc}{section}{Punti chiave del capitolo}

\begin{itemize}
  \item I colloqui sono una \textbf{selezione reciproca}. L'annuncio dice molto sulla cultura aziendale; alla fine del
        colloquio vanno sempre fatte domande (il \emph{day to day}, la ``bacchetta magica'', ``tra quanto mi fate
        sapere?''). Si arriva con una \textbf{cifra minima mentale} e con \textbf{esempi tangibili}; non si mente mai.
  \item Time management non è una lista di cose da fare: si parte dalla \textbf{must do list}, che è anche il primo
        momento per dire ``oggi potrei non farcela''; si stimano i tempi e si chiede quando qualcosa non è chiaro.
  \item Ogni nuovo task va \textbf{prima aggiunto alla lista}, non eseguito: priorità non è recenza. Le
        \textbf{checklist} fanno risparmiare tempo, anche (e soprattutto) ai più esperti.
  \item \textbf{First contact resolution}: raccogliere tutto ciò che serve in una sola richiesta, e non iniziare un task
        senza le informazioni necessarie. Ogni interruzione costa dieci-quindici minuti di concentrazione.
  \item \textbf{Matrice di Eisenhower}: fai, pianifica, delega, elimina. ``Non importante'' significa che non è importante
        che lo faccia tu; il quadrante dimenticato è \emph{delete}.
  \item Le pause vanno \textbf{scelte, non subite}; i task grossi vanno spezzettati; bisogna \textbf{saper dire di no},
        anche ai meeting (sono necessario? può andarci uno solo? ci sarà un resoconto scritto? per che cos'altro potrei usare
        questo tempo?). \textbf{Il multitasking non esiste}.
  \item A fine giornata si verifica la \emph{must do list}: un \emph{must do} rimandabile non era un \emph{must do}.
  \item Ferie illimitate con un minimo, giornate libere e feedback culture sono, per il relatore, le condizioni per un team
        sano; il benessere è anche mentale.
  \item Il \textbf{brainstorming genera idee, non decide}. Va preparato (topic, obiettivo, vincoli), separa la
        \textbf{generazione} dalla \textbf{discussione} (\emph{rapid ideation} di cinque minuti, \emph{dot voting}
        anonimo) e combatte l'\emph{anchoring} e il dominio degli estroversi.
  \item La decisione si prende con la \textbf{matrice impatto--sforzo}: prima le idee a basso sforzo e alto impatto, poi in
        diagonale. Lo sforzo sono le risorse, l'impatto comprende anche la motivazione; decide chi farà il lavoro.
  \item \textbf{Disagree and commit}: si dichiara il disaccordo, poi ci si impegna al massimo. Se va male ha sbagliato il
        team, non ``ve l'avevo detto''.
  \item \textbf{La leadership non è un titolo}: il rispetto si guadagna. Le carriere vanno divise tra \textbf{manager} e
        \textbf{subject matter expert}, pagati allo stesso modo, così che diventi manager solo chi vuole prendersi cura degli
        altri.
  \item Dopamina, serotonina e ossitocina legano il gruppo; il \textbf{cortisolo} è contagioso e lo disgrega. Il
        ``contratto sociale'' del villaggio vale ancora: chi protegge il gruppo riceve di più, chi ne abusa tradisce.
  \item La \textbf{piramide della fiducia}: integrità (aiutare davvero, onestà, ammettere gli errori), affidabilità
        (mantenere la parola, saper rifiutare), leadership (dare l'esempio, offrirsi anche quando si rischia, non lasciare
        indietro nessuno).
  \item Il gruppo deve prendersi cura del gruppo: \textbf{raccogliere il piattino}, smettere di dire ``ci penserà qualcun
        altro'', segnalare e chiedere prima di intervenire.
  \item Da manager: delegare i meriti e prendersi le colpe, lasciare la comfort zone, chiedere ``tu come faresti?'',
        celebrare successi e \textbf{fallimenti}, dare il merito a chi lo ha. Essere leader è un \textbf{impegno a lungo
        termine}.
  \item L'obiettivo: tornare a casa \textbf{ispirati, sicuri e realizzati}, come diritto di tutti e non privilegio di pochi.
\end{itemize}

%    \part{Il corso di Game Design}
%    \renewcommand{\thechapter}{\arabic{chapter}}
%    \setcounter{chapter}{0}
%    % =====================================================================
%  Capitolo 1 — Introduzione al corso e al mestiere del game designer
%  Fonte: trascrizione della lezione GD1 (Game Design, primo anno)
%  Pacchetti richiesti nel preambolo del libro:
%    \usepackage[italian]{babel}, \usepackage[T1]{fontenc},
%    \usepackage[utf8]{inputenc} (solo pdfLaTeX), \usepackage{booktabs}
% =====================================================================

\chapter[Il mestiere del game designer]{Il mestiere del game designer: introduzione al corso}
\label{cap:introduzione-game-design}

Questo primo capitolo fa da soglia. Prima di affrontare metodi, framework e
documenti, è necessario capire \emph{in che modo} si impara il game design,
\emph{con chi} lo si fa, \emph{con quali strumenti} e, soprattutto,
\emph{che cosa fa concretamente} un game designer nel suo lavoro quotidiano.
Il capitolo segue l'impostazione della lezione inaugurale del corso: si parte
dal contesto formativo (il corso come simulazione di lavoro), si passa al
metodo di apprendimento (l'\emph{inquiry-based learning}), si elencano
letture e risorse di riferimento, e si arriva infine a una prima definizione
operativa del ruolo del designer, dei limiti entro cui lavora e delle sue
possibili specializzazioni.

Una premessa terminologica: il game design è una disciplina giovane e quasi
tutta la sua letteratura è scritta in inglese. Per questo motivo molti termini
tecnici verranno lasciati in lingua originale (\emph{pipeline},
\emph{pitching}, \emph{tuning}, \emph{feature team}, \dots): sono i termini che
si useranno nel lavoro, ed è importante imparare a riconoscerli e a usarli
correttamente fin da subito.

% ---------------------------------------------------------------------
\section{Il contesto: un corso come simulazione di lavoro}
\label{sec:simulazione-lavoro}

\subsection{Il principio fondamentale}

Il primo concetto da fissare è che il corso \textbf{si presenta come una
simulazione del lavoro}. Dal momento in cui si entra in aula si viene trattati
come \emph{potenziali professionisti dell'industria}: non soltanto studenti,
ma lavoratori effettivi di uno studio (l'accademia), con una gerarchia, dei
referenti, dei compiti e delle scadenze.

Ne derivano alcune conseguenze pratiche:
\begin{itemize}
  \item il docente non assegna un ``compitino'' da risolvere per accumulare
        punti o presenze: ogni \emph{task} è pensato per far dedicare tempo
        all'apprendimento di qualcosa di concreto;
  \item non consegnare un task non danneggia il docente (che ``vive ancora
        domani e viene pagato ugualmente''), ma lo studente, che perde
        l'occasione di mettere le mani su un argomento e diventare più abile;
  \item docenti e studenti devono trattarsi \textbf{da colleghi}, perché è
        molto probabile che un domani lo diventino davvero.
\end{itemize}

\subsection{Perché l'impegno conta più che mai}

Il lavoro del game designer è descritto come \emph{molto romantico, ma anche
molto duro}, e sempre più duro col passare degli anni. Durante il periodo del
Covid l'industria ha visto grandi investimenti; quei ``rubinetti'' sono stati
chiusi da tempo, il rapporto con i publisher è diventato più complesso e le
opportunità lavorative non sono più quelle di tre o quattro anni fa. Questo
significa:
\begin{itemize}
  \item più competizione per ogni posizione;
  \item aspettative più alte: ci si aspetta che i candidati siano ancora più
        preparati e sicuri delle proprie competenze;
  \item un unico modo per diventarlo: \textbf{lavorare}. Più tempo si dedica
        oggi ad apprendere il mestiere, più facilmente si sarà assumibili
        domani e più agevole sarà il lavoro di tutti i giorni.
\end{itemize}

\subsection{Sbagliare qui per non sbagliare fuori}

Il mantra del corso è: \textbf{quanto più provate qui, anche in modo goffo e
sbagliando, tanto meno sbaglierete là fuori}. In accademia, per quanto la
simulazione sia realistica, nessuno toglie lo stipendio e nessuno licenzia da
un momento all'altro. Nell'industria sì, e un errore può costare molto denaro
all'azienda, ancora di più se l'azienda è la propria.

\subsection{Il valore del fallimento}

Il fallimento, se affrontato con la mentalità corretta, è probabilmente
\emph{la cosa più vantaggiosa} che si possa fare durante la formazione:
\begin{itemize}
  \item un successo (``avete superato questa deadline'') conferma le proprie
        abilità, ma insegna relativamente poco;
  \item dopo $n$ fallimenti si impara moltissimo, e i successi arrivano;
  \item è fondamentale non demoralizzarsi: il feedback negativo non è mai un
        giudizio sulla persona, ma su ciò che è stato consegnato.
\end{itemize}

Il docente riconosce apertamente che dedicare tempo ed energia a un progetto e
vederlo rifiutato, o vederlo non uscire mai, è demotivante: è successo a lui e
lo ha visto accadere anche su progetti molto grandi, con team numerosi. Il
punto però non è ciò che si è ottenuto materialmente, ma \textbf{quanto si è
cresciuti}. Da ogni esperienza ci si portano a casa due cose:
\begin{enumerate}
  \item \textbf{l'esperienza}: quanto si è diventati bravi, veloci,
        praticamente abili nell'applicare un framework, quanto si è cresciuti
        come designer;
  \item \textbf{il rapporto umano}: nel lavoro si ``suda sangue'' fianco a
        fianco con altre persone, e questo crea legami molto forti che
        portano avanti l'industria nel senso più positivo del termine: nascono
        progetti privati, ci si porta da un'azienda all'altra, si formano team
        con basi solide.
\end{enumerate}

Va evitato anche l'eccesso opposto: la cultura del \emph{grinding} si è
appropriata di questi discorsi in modo distorto. Non bisogna diventare robot.

% ---------------------------------------------------------------------
\section{Organizzazione del corso e gerarchia}
\label{sec:gerarchia-corso}

\subsection{La struttura gerarchica}

La struttura dei referenti del corso ricalca volutamente una
\textbf{gerarchia tipica dell'industria}:
\begin{itemize}
  \item al vertice c'è il \textbf{direttore del corso} (Diego Ricciuti), che
        corrisponde al \emph{creative director}: supervisiona la visione del
        progetto (il concetto di ``visione'' verrà approfondito più avanti);
  \item sotto di lui ci sono i \textbf{referenti di dipartimento}, uno per
        disciplina, con un corrispettivo per ciascuna sede (Milano e Roma);
  \item in fondo ci sono gli studenti, che occupano il primo gradino di
        un'azienda: l'\emph{intern} (stagista) o il \emph{junior}, il passo
        che precede il lavoratore ``regolare''.
\end{itemize}

\begin{table}[htbp]
  \centering
  \begin{tabular}{@{}p{0.34\textwidth}p{0.29\textwidth}p{0.29\textwidth}@{}}
    \toprule
    \textbf{Disciplina} & \textbf{Milano} & \textbf{Roma} \\
    \midrule
    Direzione (\emph{creative director}) & \multicolumn{2}{p{0.58\textwidth}}{Diego Ricciuti} \\
    Game design      & Mattia Podini        & Giacomo ``Jack'' Danna \\
    Level design     & Stefano Zocchi       & Andrea Dispoli \\
    Narrative design & \multicolumn{2}{p{0.58\textwidth}}{Daniele Fusetto (lezioni condivise), Giacomo Masi} \\
    Programmazione (Unreal) & Alessandro Pregnolato & Matteo Velardi \\
    \bottomrule
  \end{tabular}
  \caption{Referenti del corso per disciplina e sede (nomi come riportati in lezione).}
  \label{tab:referenti}
\end{table}

Per il primo anno, il referente diretto degli studenti di game design è il
docente di game design: è a lui che ci si rivolge per dubbi e domande. Altre
figure utili sono il direttore didattico (severo, ``giustamente severo'', da
contattare solo per questioni particolarmente importanti), l'amministrazione
(contratti e documentazione) e la segreteria.

\subsection{I canali di comunicazione}

\begin{description}
  \item[Discord] è il canale per le comunicazioni ufficiali del corso e per
        le domande che non richiedono urgenza immediata; può avere qualche
        ritardo.
  \item[E-mail] è il canale per tutto ciò che è formale o ``ancora più
        ufficiale'' e per le comunicazioni rapide e immediate.
  \item[Telegram] esistono residui di vecchi canali precedenti a Discord, da
        evitare.
\end{description}

Le domande che valgono per tutti vanno fatte in classe (anche da un portavoce
a nome del gruppo); quelle singole e personali via Discord o, se formali, via
e-mail. Se il docente non risponde a una domanda c'è un motivo, ad esempio
ragioni legali, e lo dirà esplicitamente.

\subsection{I tre moduli e la settimana tipo}

Le tre discipline principali del primo anno (e non solo) sono:
\begin{itemize}
  \item \textbf{Game Design} (tendenzialmente il lunedì);
  \item \textbf{Level Design} (tendenzialmente il mercoledì);
  \item \textbf{Programmazione} (tendenzialmente il venerdì).
\end{itemize}
Nel primo mese le lezioni sono tutte di game design (lunedì, mercoledì e
venerdì). Sono previste anche lezioni di narrativa, ma il piano formativo si
aggiorna ogni anno: le ore di narrativa sono state ridotte a favore di quelle
di programmazione. Le poche lezioni di narrativa daranno quindi dei
\emph{cenni}, non la preparazione di un narrative designer. Ciò che serve
davvero è imparare \textbf{come comunicano gli altri dipartimenti} e come
lavorare con persone che fanno un lavoro diverso ma complementare al proprio.

\subsection{Le due anime delle lezioni di game design}

\begin{description}
  \item[Design methods (metodologie)] lezioni prevalentemente teoriche, ma con
        una buona componente pratica, in cui si studiano metodologie di design
        e i cosiddetti \emph{framework} dell'industria.
  \item[Documentazione e pitching] il cuore del lavoro quotidiano del
        designer.
\end{description}

\paragraph{La documentazione.}
Il docente si definisce ``il più grande nerd della documentazione'' che si
possa incontrare. È una parte dura, perché a pochi piace documentare, ma il
lavoro del designer \emph{è principalmente questo}: scrivere documenti e, da
quei documenti, passare alla pratica e all'implementazione. Chi non sa
documentare bene parte male. Quando non si è presenti, il documento è il
proprio \textbf{biglietto da visita}: qualcun altro lo leggerà, quindi va
scritto nel modo più chiaro e comprensibile possibile.

Un esempio concreto del perché la documentazione conta: CD Projekt Red ha
sviluppato un proprio engine proprietario, ma buona parte della sua
documentazione è stata ``buttata all'acqua''. Chi deve usarlo o lo sa già
usare, o ha competenze di un altro engine da trasferire, oppure è in
difficoltà. Si documenta \textbf{perché altri possano leggere}.

\paragraph{Il pitching.}
Il \emph{pitching} è l'esposizione, la presentazione orale. La parte creativa
del team presenta spesso:
\begin{itemize}
  \item le proprie idee ai colleghi;
  \item il gioco a convention ed eventi;
  \item il pitch del gioco agli investitori, per ottenere i fondi necessari a
        realizzarlo.
\end{itemize}
Saper parlare e comunicare chiaramente le idee è fondamentale. In sintesi:
\emph{il game designer parla}, e bisogna farsene una ragione.

\subsection{Altre componenti del corso}

\begin{description}
  \item[Gioco critico e playtest] la classe viene divisa in due gruppi: uno
        gioca, l'altro osserva. Chi osserva non applaude ``come le
        scimmiette'': annota, segna e analizza il \emph{playtest}, esattamente
        come si fa nell'industria. Il playtest fa parte delle mansioni del
        designer, ma esiste anche come specializzazione autonoma.
  \item[Pipeline] termine che indica tutto ciò che riguarda i processi
        formalizzati dell'industria; di solito si intende la \emph{pipeline di
        produzione}: come si divide il team, qual è la sua gerarchia, come il
        \emph{producer} gestisce le deadline. In breve: \emph{come} e
        \emph{quando} si fanno le cose. Al primo anno è meno rilevante (si
        dovrà ``nuotare'' un po' da soli); dal secondo anno si vedranno le
        metodologie effettivamente applicate ai progetti digitali. Si
        scoprirà quanto sia difficile lavorare a un progetto digitale, e
        ancor più fisico, senza un piano di sviluppo comprovato, come quelli
        usati per qualsiasi software in qualsiasi parte del mondo.
  \item[Engine] il motore di riferimento è \textbf{Unreal Engine}
        (sezione~\ref{sec:engine}).
  \item[Programmazione] principalmente tramite \emph{Blueprint} in Unreal,
        più eventuali sistemi di versionamento; il corso parte da zero.
  \item[Progetti pratici] probabilmente la cosa più importante. Al primo
        anno i progetti di gruppo restano circoscritti al corso di design,
        perché mancano ancora alcune nozioni (come comunicare, affrontare
        problemi, dare e ricevere feedback). In seguito si allargheranno a
        progetti digitali, molto più difficili di quelli fisici, con membri
        di altri team: un rapporto non unilaterale, ma ``covalente''.
\end{description}

% ---------------------------------------------------------------------
\section{Il metodo: l'inquiry-based learning}
\label{sec:inquiry}

\subsection{Che cos'è}

Il corso è costruito secondo l'\textbf{inquiry-based learning}
(apprendimento basato sull'indagine), un approccio pedagogico nato intorno
agli anni Sessanta, che si richiama a pensatori come Dewey e Piaget. In
generale, le discipline da cui il design attinge sono quasi sempre la
sociologia, la psicologia e, talvolta, l'economia.

Il funzionamento è un ciclo:
\begin{enumerate}
  \item il docente fornisce metodologie e framework dell'industria durante
        la lezione;
  \item lo studente, fuori dalle ore di lezione, lavora sui \emph{task}:
        mette a terra le nozioni, le maneggia, le assorbe meglio;
  \item da questo lavoro nascono nuovi dubbi;
  \item i dubbi tornano in aula: lo studente chiede, il docente risponde
        (ammesso che dare la risposta sia utile, e ammesso che una risposta
        esista).
\end{enumerate}

\subsection{Non sempre esiste una risposta giusta}

Non tutto ciò che si studia ha ``la risposta migliore possibile in assoluto
nel vuoto''. Esistono metodologie, limiti e framework che aiutano il lavoro,
ma non sono fissi: cambiano da azienda ad azienda, da lavoro a lavoro, da
progetto a progetto, e a volte vengono stravolti. Sta alla
\textbf{sensibilità del designer} capire quale strumento o framework usare in
un dato momento, e questa sensibilità si impara solo con l'esperienza. Nel
lavoro nessuno dirà ``usa questa cosa'' o ``è meglio quest'altra'': sarà il
designer a scegliere.

Anche l'esperienza del docente ha confini: sei-sette anni di lavoro, un modo
personale di approcciare il design e un certo insieme di progetti (ad esempio
nessun MMO). Oltre quei confini si cercano risposte nella ricerca, studiando
come aziende con esperienza specifica (ad esempio Riot) hanno affrontato
problemi simili e hanno condiviso le loro soluzioni, sempre valutando quali
metodi siano adatti alla dimensione del team, al tipo di progetto e al tempo
disponibile.

\subsection{Sensibilità, non memoria}

Non interessa imparare a memoria un framework o la teoria di un autore. Certe
terminologie vanno conosciute, ma per ragioni di comunicazione e chiarezza.
Ciò che conta è sviluppare la \textbf{capacità di farsi domande}, di farle
agli altri, di fare ricerca e di trovare una propria risposta. È esattamente
ciò che si fa nel lavoro: davanti a qualcosa che non si sa fare, l'abilità
del designer sta nel cercare come qualcun altro ha risolto un problema simile.

\subsection{Le lezioni come corsi di aggiornamento}

Le lezioni frontali simulano i \emph{corsi di aggiornamento} aziendali: le
sessioni che un lead di design tiene al suo team di design (o il lead di
level design al suo team, e così via). Alla fine dell'aggiornamento c'è quasi
sempre un lavoro da fare, un obiettivo da portare a terra:
\textbf{l'obiettivo della settimana}. Si lavora di settimana in settimana, e
il perché verrà chiarito in seguito.

Struttura tipica di una lezione:
\begin{enumerate}
  \item breve sessione di domande e risposte sulla lezione precedente;
  \item analisi di eventuali esercizi particolarmente critici;
  \item lezione vera e propria;
  \item discussione di gruppo, se c'è tempo;
  \item assegnazione del task della settimana.
\end{enumerate}
Il docente tende a ``parlare e straparlare'', ad aprire parentesi e a saltare
da un argomento all'altro: l'unico modo per fermarlo è che siano gli studenti
a interromperlo con domande.

\subsection{La gerarchia della competenza}

L'inquiry-based learning può essere rappresentato come una scala di
competenza:
\begin{enumerate}
  \item la sola lezione frontale porta a \textbf{conoscenze incomplete},
        perché manca la parte pratica e di ricerca;
  \item il lavoro a casa (task, ricerca) porta a una competenza \textbf{più
        completa}, non sbagliata in sé, ma ancora parziale e complementare ad
        altre;
  \item la \textbf{condivisione} dei task e dei feedback con gli altri porta
        idealmente alla competenza più completa possibile;
  \item il picco della competenza è la sua \textbf{applicazione nei
        progetti} (fisici o digitali, di gruppo): ``con le cose che avete
        imparato e ricercato, riuscite a fare un gioco?''. Quando la risposta
        diventa positiva, o comunque si inizia a navigare nei processi di
        design per realizzarne uno, quella è la vittoria.
\end{enumerate}
Non si parte subito creando giochi: si parte dall'analisi, dal design e dalle
regole generali di design, prima ``nel vuoto'', poi contestualizzando i
framework, infine unendo i puntini nei progetti.

Rimarranno comunque delle lacune, per tutti, docente compreso, perché il game
design \textbf{non è una scienza esatta}, anzi, a rigore non è affatto una
scienza. Prende in prestito metodologie dalle discipline scientifiche, ma
contiene una componente imprevedibile e incalcolabile: \textbf{la
produzione}. Si possono applicare perfettamente le regole di un framework, e
quel framework può comunque non essere adatto al gioco, al proprio modo di
sviluppare o al tempo di produzione disponibile (alcuni framework richiedono
una ricerca e un'attenzione che il tempo non concede).

C'è poi la componente creativa, la parte ``romantica'' del lavoro: esiste, ma
è \textbf{metà del lavoro}. L'altra metà è ricerca, metodo scientifico al
100\%. Dagli stessi dati $n$ designer diversi, con sensibilità diverse,
possono ottenere $n$ risposte diverse. Non tutte saranno valide: vanno
applicate, messe a terra, e questa scrematura fa parte del lavoro.

\subsection{I livelli dell'apprendimento}

Un secondo modello, a gradini, descrive i livelli di apprendimento (è la
cosiddetta tassonomia SOLO):
\begin{description}
  \item[Pre-strutturale] il docente dà il compito e lo studente lo risolve,
        possibilmente il prima possibile. È esattamente ciò che \emph{non} si
        vuole.
  \item[Uni-strutturale] si affronta una sola disciplina: è il primo passo
        necessario, ed è il motivo per cui nel primo mese ci sono solo lezioni
        di game design.
  \item[Multi-strutturale] entrano in gioco altri dipartimenti: level design,
        game design, programmazione e così via.
  \item[Relazionale] si capisce che alcune cose valgono per più
        dipartimenti: alcuni framework di narrativa derivano dal design o
        possono essere usati per fare design; le regole di documentazione si
        applicano anche ad altri dipartimenti; le regole di level design si
        mescolano con quelle di game design. Dopotutto, per fare un gioco
        servono molti dipartimenti.
  \item[Astratto esteso] il punto finale: muoversi autonomamente e capire
        cosa usare e quando, sulla base della propria esperienza. L'unico modo
        per arrivarci è fare pratica e fare progetti.
\end{description}

% ---------------------------------------------------------------------
\section{Task, feedback e regole del corso}
\label{sec:task-feedback}

\subsection{I task come opportunità}

I task non sono strettamente obbligatori (esiste un numero minimo di
consegne), ma sono \textbf{opportunità}. Quanto e come ci si lavora dipende
dallo studente. Un task fatto in cinque minuti, o in cui si è stati ``furbi'',
non verrà letto né riceverà feedback.

Domande da porsi davanti a ogni task:
\begin{enumerate}
  \item \textbf{Perché} mi è stato assegnato? Qual è l'obiettivo?
  \item Come ci lavoro per ottenere un risultato \textbf{che abbia senso}? Se
        la ricerca e i dati indicano una direzione, è inutile andare in
        un'altra.
\end{enumerate}

\subsection{Quanto lavorare: un calcolo indicativo}

Se il corso è una simulazione di lavoro, vale il carico di un lavoro normale:
\begin{itemize}
  \item un lavoro a tempo pieno comprende da 30 a 45 ore settimanali; si
        considera una media di \textbf{40 ore};
  \item si sottraggono circa 9 ore di lezione: restano $\approx 31$, diciamo
        \textbf{30 ore};
  \item con circa 3 task a settimana: $30 / 3 = $ \textbf{circa 10 ore per
        task}.
\end{itemize}
Il numero può variare, ma rende l'idea: dedicare un'ora a un task e poi fare
altro non fa guadagnare nulla a nessuno. Lo studente ``butta soldi'', il
docente risparmia tempo leggendo il task in un minuto (o non leggendolo).
Dedicando invece attenzione, ricerca, cura nella formattazione e
sperimentazione, e facendolo con cognizione di causa, si impara.

\subsection{Niente voti, solo feedback}

Ad eccezione del corso di programmazione, non ci sono voti (né ``ottimo'', né
10, né 100). Ci sono \textbf{feedback}:
\begin{itemize}
  \item in game design, generalmente sulla documentazione;
  \item in level design, sui progetti pratici.
\end{itemize}
Il docente indica cosa va bene e cosa va modificato, concentrandosi in
particolare su ciò che va cambiato. Il feedback \textbf{non è mai un giudizio
sulla persona}, ma su ciò che è stato consegnato nel concreto.

\subsection{La condivisione dei task}

Tutti i task consegnati vengono \textbf{condivisi} con la classe, e possono
essere letti solo dopo aver ricevuto i feedback del docente. Si cresce così
non solo con i propri feedback, ma anche con quelli dati agli altri. Regola
importante: ciò che il docente dice a uno studente \textbf{vale per tutti}. Se
a una persona viene detto ``questa cosa non va mai più fatta'', da quel
momento vale per l'intera classe.

\subsection{Non copiare}

Copiare i task degli altri è una delle cose peggiori che si possano fare, non
per una questione morale, ma perché annulla il senso del percorso. Si può
dedicare a ogni task il tempo e l'attenzione che si ritengono opportuni, ed è
normale che col tempo ci si specializzi (preferendo, ad esempio, il level
design o la programmazione tecnica) e si investa meno dove si sente meno
necessità. In questi casi l'alternativa corretta è:
\begin{itemize}
  \item non consegnare;
  \item studiare ciò che hanno fatto gli altri;
  \item eventualmente riconsegnare in seguito.
\end{itemize}
Ci saranno momenti in cui si verrà sommersi dalle cose da fare: bisogna
imparare a scegliere.

\subsection{Il ``non leggo''}

Il docente ha piena facoltà di non dare alcun feedback a un task (il
cosiddetto \textbf{``non leggo''}) se:
\begin{itemize}
  \item non si è lavorato con attenzione;
  \item non si rispettano le indicazioni date;
  \item è stata violata una regola estremamente importante.
\end{itemize}
Un task ``non letto'' equivale a un task non consegnato e non entra nel
conteggio. Attenzione: un task \emph{perfetto} e un task che \emph{fallisce}
nel comprendere o applicare un framework sono \textbf{ugualmente validi},
purché ci si sia impegnati e si sia dedicato tempo, pazienza e ``sudore''.

\paragraph{Riconsegna.} Un task non letto, in ritardo o che si vuole
migliorare può essere riconsegnato e verrà conteggiato, a patto che:
\begin{itemize}
  \item i problemi segnalati siano stati risolti;
  \item sia stato fatto ulteriore lavoro, integrando anche tutto ciò che nel
        frattempo è stato detto agli altri colleghi nei loro feedback.
\end{itemize}
Il task va sviluppato con le competenze che si hanno nel momento della
riconsegna, non con quelle del momento della prima consegna.

\subsection{Requisiti per superare l'anno}

\begin{itemize}
  \item consegnare almeno l'\textbf{80\% dei task} assegnati (soglia
        malleabile);
  \item partecipare almeno all'\textbf{80\% delle lezioni}, in presenza o
        da remoto, preferibilmente in presenza; i ritardi vengono valutati
        caso per caso;
  \item comportarsi da persone corrette e rispettose (è stata firmata una
        carta dello studente): non urlare, non alzare le mani, non
        vandalizzare, non compiere atti illegali.
\end{itemize}
Impegnarsi di più porta inoltre a essere selezionati per varie opportunità.

\subsection{Scadenze}

\begin{itemize}
  \item la scadenza di un task è \textbf{sempre prima dell'inizio della
        lezione} successiva indicata;
  \item se la classe, come team, mantiene comportamenti scorretti, il docente
        può \textbf{anticipare le deadline} per tutti.
\end{itemize}

\subsection{Accessibilità e DSA}

Gli studenti con disturbi specifici dell'apprendimento o neurodivergenze
(dislessia, ADHD, autismo, ecc.) sono invitati a segnalarlo. Quando si
affronteranno le regole di formattazione e documentazione (errori
grammaticali, mancanze, incoerenze) se ne terrà conto e si potranno trovare
dei compromessi, pur ricordando che nel mondo del lavoro nessuno farà
``sconti''. La documentazione può essere inviata all'amministrazione o alla
segreteria, che la comunicherà a tutti i docenti, ma il consiglio è parlarne
anche direttamente con ogni singolo docente.

% ---------------------------------------------------------------------
\section{Le competenze trasversali}
\label{sec:soft-skill}

\subsection{I giochi si fanno in gruppo}

Si possono fare giochi da soli, ed è un'impresa ammirevole; ma lavorando da
soli \textbf{tutto ciò che si guadagna in libertà si perde in produttività}.

\paragraph{Il caso Stardew Valley.} \emph{Stardew Valley} è stato realizzato
da una sola persona, Eric Barone, in anni di sviluppo e di sacrifici: faceva
un altro lavoro per mantenersi (in un cinema) e di notte programmava, faceva
l'arte e componeva la musica. Il risultato è stato straordinario e lo ha reso
ricco, ma si tratta di un caso estremamente raro. Quanti giochi simili, nati
nelle stesse condizioni, sono falliti? Molti di più.

\paragraph{Il rischio d'impresa.} Lavorare da soli non è sbagliato di per sé
e ha aspetti molto positivi: si ha il privilegio di creare esattamente ciò che
si vuole. Ma ci si fa carico anche di tutto il \textbf{rischio d'impresa}: se
si lavora a un gioco per cinque anni e poi il gioco va male, resta il
vantaggio di averlo pubblicato, ma il dopo è sperare di vivere di quella
release o di finanziare con essa il gioco successivo. Le grandi aziende hanno
gli stessi problemi, ma li distribuiscono su più persone e più responsabilità.

L'obiettivo di un professionista è \textbf{saper lavorare in diverse
condizioni}, e la più comune è quella del dipendente (o del freelance /
libero professionista con partita IVA): si viene assunti per le proprie
capacità da un'azienda, piccola o grande che sia, e si lavora con altre
persone.

\subsection{Essere persone con cui si vuole lavorare}

Se le persone non vogliono lavorare con voi, quello è \textbf{il primo motivo
di licenziamento}. E se non si viene licenziati, il clima nell'azienda
diventa terribile: la produttività crolla, la motivazione crolla. Le
\emph{soft skill}, cioè l'insieme di abilità che rendono una persona
piacevole con cui lavorare, vanno coltivate e diventano sempre più importanti
col passare degli anni (relativamente meno ora, molto di più quando i team
saranno più grandi). Anche chi è un ``lupo solitario'' deve imparare a:
\begin{itemize}
  \item lavorare con altre persone;
  \item comunicare con altre persone;
  \item trovare un punto in comune da cui partire per fare le cose.
\end{itemize}

\subsection{Indipendenza e proattività}

Sono le altre due parole chiave del corso:
\begin{description}
  \item[Indipendenza] spesso si verrà lasciati da soli, non per cattiveria ma
        per sfida: quando si cerca una risposta da soli, ciò che conta non è
        tanto la risposta quanto \textbf{il processo}. Una risposta data dal
        docente non vale niente.
  \item[Proattività] chiedere, domandare, agire da soli senza aspettare
        direttive, perché spesso quelle direttive non ci sono, a maggior
        ragione nei team piccoli.
\end{description}
Tutto questo si concretizza nel parlare, nel presentare e nel fare
\emph{post-mortem} (l'analisi a posteriori di un progetto concluso, che verrà
approfondita più avanti), e in generale nell'attenzione alla crescita
personale.

\subsection{Fare domande}

Il corso insiste sulle domande: alzare la mano, intervenire anche ``a
bomba'', fermare il docente ogni volta che qualcosa non è chiaro. Se una
persona non conosce qualcosa, molto probabilmente c'è almeno un'altra persona
che non la conosce. Il docente può dare per scontata la conoscenza di termini
che non lo sono, ed è del tutto normale, soprattutto all'inizio.

% ---------------------------------------------------------------------
\section{Strumenti: engine, Unreal e materiali}
\label{sec:engine}

\subsection{Unreal Engine e gli altri motori}

Il motore di riferimento del corso, usato in level design, è \textbf{Unreal
Engine}: ``un gran mattone'', ma che dà a un designer molte possibilità. Il
vantaggio principale è che permette a un designer non pratico di
programmazione di \textbf{prototipare tramite i Blueprint} (il sistema di
scripting visuale di Unreal). Questo non impedisce di passare a Unity o ad
altri engine, e viceversa: il docente stesso ha studiato e lavorato a lungo
con Unreal, ma la sua ultima esperienza professionale (come level designer) è
stata su Unity, un engine che non conosceva, a dimostrazione che il salto da
un engine all'altro è possibile.

Panorama dei motori citati:
\begin{itemize}
  \item \textbf{Unreal Engine} e \textbf{Unity}: i due più diffusi, con
        modelli di sorgente ``più o meno'' aperto;
  \item \textbf{Godot}: engine open source;
  \item \textbf{CryEngine}: proprietario di Crytek;
  \item \textbf{Snowdrop}: engine di Ubisoft;
  \item Rockstar usa un proprio engine proprietario;
  \item alcune aziende usano versioni \textbf{personalizzate} di engine
        esistenti: ad esempio Milestone, nei giochi di corse, usa un Unreal
        customizzato;
  \item CD Projekt Red ha un engine proprietario (con i problemi di
        documentazione già citati).
\end{itemize}
In sintesi: ogni azienda può avere un engine proprietario oppure una versione
modificata di un engine commerciale, ma si parte sempre da una base.

\subsection{Programmazione: cosa ci si aspetta da un designer}

Non è richiesto alcun bagaglio pregresso di programmazione: il corso parte da
zero. Gli studenti sono stati selezionati (il corso è a numero chiuso) per
competenze analitiche e critiche sopra la media. Anche nel mercato del lavoro,
per una posizione di game designer generalmente \textbf{non ci si aspetta
competenza di scripting} (C++, Java, C\#, ecc.), a meno di voler fare il
\emph{technical designer}, e anche lì varia da azienda ad azienda.

\subsection{Versione e preparazione}

Indicazioni pratiche emerse in aula:
\begin{itemize}
  \item scaricare Unreal Engine; la versione consigliata per il corso è
        stata indicata nella \textbf{5.6} (con conferma successiva). La
        stabilità dipende da cosa si deve fare, da quali tool e plugin si
        usano: i problemi nascono spesso dai plugin. Nel passaggio da Unreal 4
        a Unreal 5, ad esempio, il sistema \emph{Landscape} era inizialmente
        inutilizzabile, mentre in Unreal 4 funzionava molto bene;
  \item è consigliato installare tramite \textbf{Epic Games Launcher}, più
        comodo: permette anche di riscattare giochi gratuiti periodici e di
        accedere ai pacchetti di \textbf{Fab} (il marketplace di asset
        acquisito da Epic), che ogni mese offre alcuni asset gratuiti utili
        per esercitarsi nel level design (ad esempio per fare \emph{speed
        level design});
  \item molti pacchetti di UI o di sistemi di interazione sono utili per
        prototipare al volo, ma spesso inutilmente complicati;
  \item un plugin di \textbf{materiali di prototipazione a griglia} (adattivi,
        con metriche, allineabili alla griglia, con materiali già pronti e
        colori modificabili, dal costo di circa 5 euro) è indicato come
        ``salvavita'' per il level designer; se necessario, i pacchetti
        indispensabili vengono forniti;
  \item portare \textbf{sempre il portatile} alle lezioni di level design
        (da novembre, a partire dall'interfaccia dell'engine); alle lezioni
        di game design basta che lo porti almeno metà della classe, per i
        lavori a coppie. Serve comunque qualcosa per prendere appunti
        (``tavolette d'argilla e cannucce, se volete fare i babilonesi'',
        oppure carta e penna).
\end{itemize}
Chi ha un portatile poco potente dovrebbe comunque provare a installare
l'engine e aprire un template per vedere come gira: all'inizio non si faranno
cose pesanti. Nelle prime lezioni di level design è importante mettere subito
le mani sull'interfaccia: ascoltare dove si trova uno strumento (ad esempio lo
\emph{snapping}) senza provarlo porta a dimenticarlo.

% ---------------------------------------------------------------------
\section{Letture e risorse di riferimento}
\label{sec:risorse}

\subsection{I libri}

\begin{description}
  \item[\emph{The Art of Game Design: A Book of Lenses} (Jesse Schell)] il
        testo \textbf{davvero importante}. Alcune lezioni finali dell'anno si
        baseranno sulle sue terminologie. È utile perché apre la mente sulla
        possibilità di applicare metodologie diverse a giochi diversi, ma
        soprattutto per la \textbf{terminologia} comune. Esiste anche in
        audiolibro (circa 25 ore).
  \item[\emph{Challenges for Game Designers} (Brenda Romero e Ian Schreiber)]
        libro \emph{entry level} con moltissimi esercizi, molto caro al
        direttore del corso.
  \item[\emph{21st Century Game Design} (Chris Bateman)] utile come
        retrospettiva: è un libro visionario che ha previsto metodi e
        soprattutto \emph{archetipizzazioni} dei giocatori tuttora valide nel
        design e nel marketing. Servirà per una delle lezioni più importanti
        dell'anno. È utile leggerlo per vedere quanto aveva previsto, ma è
        utile ``fino a un certo punto''.
  \item[\emph{An Architectural Approach to Level Design} (Christopher
        Totten)] non strettamente necessario, ma di lettura semplice. Ha due
        anime: parte dall'architettura (babilonese, greca, ecc.) e la trasla
        su esempi di level design (come \emph{Doom} e \emph{Team Fortress}).
        Un tempo ne esisteva un riassunto su Gamasutra (oggi
        \emph{Game Developer}).
  \item[\emph{A Theory of Fun for Game Design} (Raph Koster)] utile fino a un
        certo punto: contiene diverse idee superate, che lo stesso Koster ha
        affrontato in vari talk reperibili nel GDC Vault. Verrà ripreso nelle
        prime lezioni, poi messo da parte.
\end{description}

\subsection{Canali e conferenze}

\begin{description}
  \item[GDC (Game Developers Conference)] il suo canale raccoglie talk e
        panel estremamente importanti: è ``uno strumento misterioso che ci
        servirà più tardi''. \textbf{Attenzione}: per molti talk mancano ancora
        le competenze per interpretarli, e alcuni potrebbero fuorviare o
        mettere in testa idee strane. La GDC si concentra spesso su
        produttività e mobile.
  \item[Il \emph{Dragon Speech} di Chris Crawford] il talk del fondatore della
        GDC: appartiene a un passato lontano, ma è molto motivante. Da vedere
        subito.
  \item[Game Maker's Toolkit] canale particolarmente utile per nozioni
        pratiche, con video brevi e alcune interviste.
  \item[Nordic Game] altra conferenza importante, con talk e panel di esperti
        del settore.
\end{description}

\subsection{Associazioni e premi: come leggerli}

Le associazioni di settore (ad esempio IIDEA in Italia, o la IGDA) non sono
il male, anzi: sono un canale importante per affacciarsi all'estero o per
accedere a fondi. Va però ricordato che \textbf{tutelano i propri associati}:
non esiste solo il loro palinsesto, e bisogna valutare quanto spazio occupano
nel panorama della visibilità.

Lo stesso vale per i \textbf{premi} (Game Awards, ``Best Mobile Game'',
``Games for Change'', ecc.): al designer \emph{non interessano}. Che un gioco
abbia vinto ``più di 90 premi'' non dice nulla di utile. Servono a dare
visibilità e lustro ai giochi ed è giusto che esistano, perché ampliano il
panorama e il numero di giocatori. Ma ciò che interessa al designer sono i
\textbf{dati}:
\begin{itemize}
  \item come sono state implementate certe cose;
  \item come è stato fatto il marketing e dove il gioco è stato presentato;
  \item non quante wishlist ha fatto, ma \emph{come} le ha fatte;
  \item in che periodo storico è uscito e quali erano i competitor in quel
        momento.
\end{itemize}

% ---------------------------------------------------------------------
\section{Che cosa fa un game designer}
\label{sec:cosa-fa}

Per arrivare a una prima definizione, la classe ha condotto un brainstorming:
ognuno doveva indicare una cosa che, secondo lui, fa il game designer (senza
ripetere quelle già dette). Ne è emersa la lista seguente, commentata punto
per punto.

\subsection{Progettare meccaniche e struttura del gioco}

Il ruolo più generico è quello del \textbf{core mechanic designer}, detto
anche \emph{gameplay designer} e, un tempo, \emph{black box designer}: si
occupa delle \textbf{meccaniche di gioco}. Si va dal piccolo al grande (in
funzione anche della dimensione dell'azienda):
\begin{itemize}
  \item nel \textbf{piccolo} ci sono le singole meccaniche;
  \item nel \textbf{grande} ci sono i \textbf{sistemi}. Nei giochi
        \emph{sistemici}, con molti sistemi interconnessi (ad esempio
        \emph{GTA}), esistono specializzazioni dedicate a bilanciare e
        supervisionare i sistemi: è il \emph{system design}.
\end{itemize}

\paragraph{Quando un gioco è ``sistemico''?}
Non dipende solo da come le meccaniche interagiscono con il giocatore (le
\emph{dinamiche}, secondo i framework che passano dalle meccaniche alle
dinamiche, che verranno studiati più avanti), ma anche da \emph{quante}
meccaniche ci sono e da quanto interagiscono tra loro. Non diremmo che
\emph{Super Mario} sia un gioco sistemico, eppure contiene dei sistemi: la
linea è sottile. Quando i sistemi aumentano molto, sono complessi e si
influenzano a vicenda, il gioco diventa sistemico. Esempio: in \emph{The
Legend of Zelda: Breath of the Wild} il vento modifica quanto il personaggio
sale o scende; il fuoco può incendiare l'erba; l'erba prende fuoco se colpita
da una freccia infuocata; e così via. Nel mercato moderno i giochi sistemici
sono spesso, ma non sempre, legati a mappe \emph{open world}.

\subsection{Le 3C: camera, character, control}

La base del gioco, da cui generalmente si inizia a lavorare (anche nei
prototipi), è la ``sacra trinità videoludica'' delle \textbf{3C}:
\begin{description}
  \item[Camera] che tipo di inquadratura ha il gioco: prima persona, terza
        persona, isometrica, ecc.
  \item[Character] quanto è grande il personaggio, che tipo di personaggio è
        e, di conseguenza, quali azioni può compiere.
  \item[Control] i controlli: che cosa il personaggio può fare, quali sono i
        \emph{verbi} di gioco.
\end{description}
Le 3C sono definite dal designer, ed esiste anche una specializzazione per
questo lavoro. Esistono varianti: quella del direttore del corso è
\textbf{4C + J}, cioè \emph{camera, character, control, conflict} (conflitto)
e \emph{juice} (il ``succo'', il feedback che rende piacevole l'interazione).
Verranno approfondite più avanti.

\subsection{Il personaggio}

Il designer definisce il personaggio dal punto di vista delle meccaniche
(dimensioni, azioni, capacità), all'interno delle 3C. L'aspetto estetico è
invece competenza di altri dipartimenti.

\subsection{Documentare le informazioni per gli altri dipartimenti}

Assolutamente vero: la documentazione è una delle prime attività del
designer. Più l'azienda è grande, più questo lavoro diventa prominente; più è
piccola, meno è rilevante. In \textbf{pre-produzione}, però, è sempre
fondamentale.

\subsection{Mediatore e leader ``al servizio''}

Il designer è un \textbf{mediatore} nei conflitti interni al team, una sorta
di collante che si relaziona con dipartimenti diversi (non sempre gli stessi,
a seconda della fase di produzione). A volte fa anche da leader, a seconda di
come sono gestite la pipeline e la gerarchia.

Questa leadership è difficile da definire perché \textbf{non è verticale}: è
una leadership \emph{al servizio} (\emph{servant leadership}). Il designer
``è asservito al programmatore'': il suo compito non è arrivare con un'idea e
scaricarla addosso agli altri, né portare la soluzione e lasciare il problema
a qualcun altro, ma \textbf{seguire il processo a basso livello}, sporcandosi
le mani.

\paragraph{I feature team.} Nei team medio-piccoli o medio-grandi possono
esistere sotto-team (\emph{feature team}) dedicati a un obiettivo specifico:
ad esempio un designer incaricato del character o della camera, affiancato da
uno o due programmatori e da un artista. In questi casi il designer
supervisiona il lavoro e gestisce la produzione del sotto-team.

\subsection{Gestire l'avanzamento del prodotto}

Esiste una specializzazione che si occupa solo di questo, il
\textbf{producer}, a un livello gerarchico completamente diverso (e che
potrebbe essere una strada percorribile al termine del percorso di studi).
All'interno di un feature team anche il designer può farlo, ma non è detto.
Può certamente \emph{dare la direzione}: nella sua documentazione può
indicare che per arrivare a uno stato giocabile del prototipo serve un certo
lavoro di arte e di programmazione, per lo stato successivo un altro, e così
via. In ambito produttivo esistono metodologie strutturate come
\textbf{Agile}, che verranno viste in seguito.

\subsection{Presentare le idee}

Il pitching (sezione~\ref{sec:gerarchia-corso}) è richiesto al designer sia
in pre-produzione sia in produzione, in misura che dipende dall'azienda. Ad
esempio Ubisoft lascia ai designer, e in realtà a qualsiasi membro, la
possibilità di \emph{pitchare} un'idea prima di una nuova produzione, perché
possa iniziare a prendere vita.

\subsection{Bozze e prototipi}

Le bozze vanno bene, ma devono essere \textbf{chiare}. C'è un'eterna lotta tra
documentazione e gioco giocabile: se si fa solo documentazione il gioco non
avanza; se si fa solo il gioco manca la documentazione e si perde la visione.
Sono due cose \textbf{complementari}.

Il modo per comunicare un'idea in forma grezza ma giocabile è il
\textbf{prototipo}. Un prototipo può essere qualsiasi cosa: di carta, in
Unreal con Blueprint ``brutti''. \emph{Non deve essere bello, deve essere
giocabile}.

\paragraph{Esempio: il prototipo di carta di Mario + Rabbids.} Durante lo
sviluppo di \emph{Mario + Rabbids}, i programmatori di Ubisoft erano impegnati
su altri progetti (come \emph{Ghost Recon}). I designer realizzarono allora
un \textbf{gioco da tavolo} che riproduceva le meccaniche del gioco, quelle
di uno strategico a turni sul modello di \emph{XCOM}. E funzionava.
L'obiettivo era \textbf{validare un'idea}, validare un gameplay.

Altri modi per rendere chiara un'idea sono i \emph{mock-up} nella
documentazione e le immagini modificate o i \emph{mash-up} di qualsiasi
natura. Lato documentazione l'importante è che sia chiaro; lato prototipo,
che sia giocabile.

\subsection{Fare ricerca per l'idea, e poi smontarla}

Il designer ha una componente creativa importante e fa ricerca per arrivare a
un'idea. Ma ciò che deve poi fare necessariamente è \textbf{smontarla}. Il
suo lavoro è scientifico, esattamente come il \textbf{metodo scientifico}, il
cui obiettivo è \emph{confutare un'ipotesi}: si arriva a un'ipotesi tramite
sperimentazione e poi si cerca $n$ volte di confutarla; se resiste, la si
assume come valida \emph{in un certo contesto} (le regole della fisica
classica, ad esempio, non valgono nella meccanica quantistica).

Allo stesso modo il designer arriva a un'idea (tramite ricerca o chiedendo
agli altri: le idee le hanno tutti, non sono un privilegio del designer),
cerca di smontarla e poi cerca di \textbf{metterla a terra}. Se l'idea non è
applicabile, non è implementabile, finché non si trova un modo per
realizzarla va scartata. Può essere l'idea più bella dell'universo, ma
\textbf{se non può essere messa a terra, non esiste}.

\paragraph{L'originalità.} Serve un'idea originale? Dipende dall'obiettivo e
dalla visione del progetto:
\begin{itemize}
  \item se l'obiettivo è vendere molto, un'idea originale non serve
        necessariamente;
  \item se si vuole creare un nuovo filone, serve un'idea, se non originale,
        almeno interessante: può essere un discorso di marketing o una
        meccanica innovativa, con il rischio di non vendere.
\end{itemize}
Il ruolo del \emph{director} e il concetto di \emph{visione} verranno
approfonditi in seguito.

\subsection{Analizzare i competitor}

Arrivati a un'idea (``vogliamo un certo tipo di combattimento''), si va a
vedere se qualcuno l'ha già fatta e \emph{come}. Le meccaniche non sono
privilegio di nessuno, salvo quando vengono brevettate, come il
\textbf{Nemesis System} di \emph{La Terra di Mezzo: L'Ombra di Mordor}.

\paragraph{Esempio: il combattimento da ridimensionare.} In un progetto
(presso Stormind Games) il team si trovò con un sistema di combattimento
troppo costoso: mancavano le risorse e gli animatori per realizzarlo. Poco
output lato gameplay significava non poter fare qualcosa di complesso e
profondo come in \emph{Elden Ring}; al contrario, il team aveva molti
artisti e poteva puntare su qualcosa di più artistico. Si iniziò così a
orbitare intorno al corpo a corpo di \emph{The Last of Us} (parte I e II).
Sul canale della GDC esistono diversi talk che raccontano cosa è stato fatto
e come, con successi, insuccessi e costi (ad esempio del \emph{motion
capture}): da quei talk si ricavarono moltissime informazioni utili.

I competitor possono essere:
\begin{itemize}
  \item \textbf{diretti}: un gioco molto simile al proprio;
  \item \textbf{di mercato}: un gioco che domina il mercato in cui ci si vuole
        affacciare;
  \item \textbf{indiretti}: un gioco con un sistema simile al proprio, ma
        tutto il resto diverso.
\end{itemize}
Regola d'oro: \textbf{là fuori, nascosto nell'ombra, c'è sicuramente un gioco
molto simile al vostro}. Va solo trovato. E se per qualche ragione non ha
venduto, bisogna capire perché.

\subsection{Lavorare in engine e fare tuning}

Il level designer lavora in engine, così come il technical designer; ma
qualsiasi game designer, a seconda di ciò che deve fare, può finire in
engine, soprattutto per il \textbf{bilanciamento} e il \textbf{tuning}.

Il \emph{tuning} è l'azione di bilanciare con precisione i dati. Esempio: in
\emph{Super Mario} si modifica quanto in alto salta il personaggio, con che
velocità, in quanto tempo raggiunge l'apice, e così via. Questo lavoro rende
il gioco:
\begin{itemize}
  \item ad alto livello, dotato della \textbf{progressione} giusta;
  \item a basso livello, \textbf{piacevole da giocare}.
\end{itemize}
Il \emph{fine tuning} è un lavoro di dettaglio, paziente e minuzioso.

\subsection{Bilanciare}

Il bilanciamento, emerso tardi nel brainstorming, è una delle attività
centrali del designer e ha anche una specializzazione dedicata
(l'\emph{economy designer}, sezione~\ref{sec:tipi-designer}).

\subsection{Analizzare il mercato}

In pre-produzione, di solito, sì: capire le tendenze, i generi che cambiano
nel tempo, la compatibilità del gioco con il momento storico, in vista ad
esempio di un possibile publisher. Va capito quanto, come e con quanto tempo
farlo, perché si è sempre limitati da molte cose.

\subsection{Sintesi}

\begin{quote}
Il designer parte dall'idea ma non si ferma lì: la confuta, lavora sui
dettagli, bilancia, fa \emph{fine tuning}; lavora con il team e fa da
\emph{manager} per diversi dipartimenti; e naviga costantemente tra i limiti.
\end{quote}

% ---------------------------------------------------------------------
\section{I limiti del design}
\label{sec:limiti}

I limiti ci sono sempre, e nel mondo del design sono spesso numerosi. I più
importanti entro cui navigare sono:
\begin{description}
  \item[Tempo] ad alto livello, l'idea deve essere realizzabile nel tempo di
        produzione disponibile; a basso livello, il proprio lavoro
        settimanale deve essere fattibile nel tempo che si ha. Non ci si può
        dividere tra mille feature team se realisticamente se ne può seguire
        uno solo.
  \item[Tecnologia] si può o non si può fare qualcosa? Ci sono i server per un
        multiplayer con un certo numero di giocatori? Ci si può permettere,
        ad esempio, una fauna che reagisce alla fisica come quella di
        \emph{GTA VI}, che richiede il meglio dei \emph{technical artist}?
  \item[Background e competenze] che cosa sa fare il team? Realisticamente,
        chi ha lavorato dieci anni su giochi horror viene assunto per un
        horror, perché ha le competenze per risolvere problemi già visti.
  \item[Persone e motivazione] il team ha la motivazione e le competenze per
        fare un certo tipo di gioco, e per lavorare insieme?
  \item[Denaro] tutto ruota intorno ai soldi. Che si lavori da soli o in mille,
        le persone vanno pagate e il gioco deve avere un ritorno economico.
        Bisogna decidere dove tracciare la linea: raggiungere il
        \textbf{break-even} (guadagnare esattamente quanto speso) o puntare a
        grandi guadagni. Con un \textbf{publisher}, inoltre, una parte
        consistente dei ricavi va a lui (ripartizioni dell'ordine di 80/20 o
        70/30, che dipendono dal contratto firmato).
\end{description}

\subsection{Non siete ``gli uomini delle idee''}

Per qualcuno può essere un problema scoprirlo, ma il designer non è
``l'uomo (o la donna) delle idee'': accanto alla parte creativa c'è una parte
tecnica, e si è fortemente limitati da molti fattori. Questi limiti, però,
sono un aiuto:
\begin{itemize}
  \item riducono il ventaglio di possibilità potenzialmente distruttive per il
        team (cose che non si riescono a fare, che non piacciono o per cui
        mancano le risorse);
  \item aiutano a mirare bene, trovando soluzioni creative in uno spazio più
        piccolo e più attuabile.
\end{itemize}

\begin{quote}
\centering
\textbf{Le idee sono nulla, l'implementazione è tutto.}
\end{quote}

% ---------------------------------------------------------------------
\section{Ruoli e specializzazioni}
\label{sec:tipi-designer}

\subsection{La scala gerarchica}

\begin{description}
  \item[Creative Director (CD)] mantiene la \textbf{visione} e dà l'idea, in
        modo che tutto segua una certa direzione. In teoria può diventarlo
        chiunque, ma con un'esperienza notevole, nell'ordine dei decenni: si
        assume che in quel tempo abbia visto come lavorano molti dipartimenti.
        Naturalmente un CD con un background artistico ha meno competenze di
        design, e viceversa.
  \item[Director di disciplina] uno per dipartimento: \emph{game director},
        \emph{level director}, \emph{art director}, ecc.
  \item[Lead] il capo del dipartimento, che gestisce i gruppi di persone. Chi
        lavora in un dipartimento risponde al proprio lead (ad esempio un game
        designer risponde al \emph{lead designer}, non al game director).
  \item[Senior / Principal] i profili più esperti (a seconda dell'azienda).
  \item[Junior] il profilo di inizio carriera.
  \item[Intern] lo stagista.
\end{description}
Lo schema si replica uno a uno per ogni dipartimento: arte, suono, animazione
(a volte arte e animazione insieme, o \emph{rigging} e animazione insieme),
programmazione, UX/UI e così via. A fare la differenza, a parità di ruolo, è
il \textbf{livello di seniority}.

\subsection{UX e UI}

UX e UI stanno a metà tra game design e arte: a seconda dell'azienda il ruolo
si chiama \emph{UX/UI Artist} o \emph{UX/UI Designer}.
\begin{description}
  \item[UX (User Experience)] come una persona interagisce con le cose. Esempio:
        su uno smartphone, un pulsante accende e spegne lo schermo; per
        digitare la password si scorre verso l'alto e compaiono i punti da
        toccare.
  \item[UI (User Interface)] come sono ``vestite'' queste cose: ad esempio
        l'aspetto della barra dei punti vita (HP), realizzata con asset 2D o
        3D.
\end{description}

\subsection{Le specializzazioni del game design}

Più l'azienda è grande, più la specializzazione è puntuale.
\begin{description}
  \item[Gameplay / Core mechanics designer] (sono la stessa cosa) si occupa
        delle meccaniche e delle 3C.
  \item[System designer] si occupa dei sistemi e delle loro interazioni.
  \item[Economy designer] si occupa del bilanciamento; lavora soprattutto con
        fogli di calcolo (Excel, Google Sheets).
  \item[Combat designer] si occupa del combattimento; negli annunci di lavoro
        indica quasi sempre un combattimento \emph{action}, con gameplay
        reattivo (non strategico), anche se non è una regola.
  \item[Multiplayer designer] si occupa delle interazioni multigiocatore.
  \item[Technical designer] fa da ponte tra game design e programmazione.
\end{description}

\subsection{Le specializzazioni del level design}

\begin{description}
  \item[Encounter designer] posiziona i nemici nella mappa in modo che gli
        scontri siano bilanciati e la progressione equilibrata come
        desiderato.
  \item[World / map designer] popola il mondo di gioco.
  \item[Quest designer] popola il mondo con le missioni e pensa alla
        progressione al loro interno, cioè a dove il giocatore deve andare (si
        vedano ad esempio gli annunci di lavoro di CD Projekt Red o di
        Larian).
  \item[Technical level designer] fa da ponte tra level design e
        programmazione.
\end{description}

\subsection{Discipline vicine: suono, UI, produzione}

Il monte ore del corso non permette di uscire come \emph{narrative designer}
(o \emph{sound designer}) completi. Ciò non toglie che si possano coltivare
queste discipline:
\begin{itemize}
  \item nei progetti di gruppo si può occuparsi anche del suono, della UI/UX,
        o provare ad applicare metodologie come Agile per fare
        \emph{producing} e gestire il team;
  \item il tempo dedicato a queste attività va però sottratto al game design,
        e dipende da quante persone ci sono nel team, da quanto lavoro c'è e
        da quanto è grande il progetto;
  \item avere sviluppato e rilasciato un progetto, anche piccolo (anche una
        \emph{game jam}), documentando i propri processi, è già un ottimo
        argomento da presentare in futuro;
  \item la via migliore resta lavorare nell'industria e, da lì, chiedere di
        imparare anche altre discipline.
\end{itemize}
Un gioco si può fare su qualsiasi cosa: sugli scarafaggi, sui piedini delle
sedie, sui pennarelli. Si parte avvantaggiati da ciò che è vicino alla propria
sensibilità: chi ama la musica può pensare a un gioco con meccaniche legate
alla musica, con un gameplay innovativo ma contenuto e realizzabile. Il design
è un cappello che comprende molte discipline.

% ---------------------------------------------------------------------
\section{Approfondimenti emersi in lezione}
\label{sec:approfondimenti}

\subsection{Roguelike e roguelite}

Un esempio di come il marketing conia generi: Supergiant Games, con
\emph{Hades}, rese celebre il termine \textbf{roguelite} al posto di
\emph{roguelike} (pur non essendo il primo roguelite in assoluto). La
differenza:
\begin{itemize}
  \item il \textbf{roguelike} si fonda su una serie di punti fissati in una
        sorta di \emph{manifesto} che definisce che cosa rende un gioco
        roguelike; uno su tutti: a ogni morte si ricomincia la \emph{run} da
        capo;
  \item il \textbf{roguelite} rompe alcuni di questi punti, in particolare
        introducendo una \textbf{progressione esterna alla run}: il giocatore
        può potenziarsi fuori dalla singola partita per vincere più
        facilmente in quelle successive. \emph{Hades} ne è il caso principe.
\end{itemize}
Analogamente, un mobile game sviluppato dal direttore del corso fu
promosso come ``il primo \emph{donation game} al mondo'': un gioco che
convertiva il tempo di gioco in donazioni reali a enti benefici per animali.
Anche qui un neologismo creato lato marketing per definire un genere.

\subsection{Applicazioni fuori dall'industria videoludica}

Le competenze sviluppate non servono solo ai videogiochi:
\begin{itemize}
  \item i \textbf{render architettonici} oggi possono essere realizzati con
        Unreal, e diversi professionisti formati nel game design lavorano in
        questo campo;
  \item la \textbf{gamification}, cioè l'applicazione di regole videoludiche
        a contesti non ludici (mostre interattive, progetti culturali,
        consulenze come quella per applicare questi concetti alle macchine da
        allenamento di un'azienda di fitness), è in Italia uno dei modi più
        redditizi di spendere queste competenze, spesso con finanziamenti
        pubblici.
\end{itemize}

% ---------------------------------------------------------------------
\section{Il primo task}
\label{sec:primo-task}

\begin{enumerate}
  \item Creare un indirizzo e-mail \textbf{professionale} (consigliato: un indirizzo Gmail nella forma
        \texttt{nome.cognome@gmail.com}; l'anno di nascita è accettabile).
        Va bene anche un indirizzo professionale già esistente usato per
        lavoro. Da evitare indirizzi poco seri (``non \texttt{blackfire67}'').
  \item Una persona della classe raccoglie tutti gli indirizzi in un unico
        documento e lo invia via e-mail al docente, che aggiungerà tutti al
        Drive condiviso del corso (dove verranno caricate lezioni, slide e
        materiali, per avere i documenti tutti nello stesso posto).
  \item Installare Discord e iscriversi al server del corso, dove passano le
        comunicazioni ufficiali. È consuetudine, e utile, creare anche un
        gruppo tra studenti.
\end{enumerate}
\textbf{Scadenza}: prima della lezione del venerdì.

% ---------------------------------------------------------------------
\section*{Punti chiave del capitolo}
\addcontentsline{toc}{section}{Punti chiave del capitolo}

\begin{itemize}
  \item Il corso è una \textbf{simulazione di lavoro}: ci si tratta da
        colleghi, con una gerarchia da studio di sviluppo.
  \item Si impara con l'\textbf{inquiry-based learning}: lezione, task,
        dubbi, domande, condivisione, progetti.
  \item Il fallimento affrontato con la mentalità giusta è la principale fonte
        di crescita; i feedback riguardano il lavoro, mai la persona.
  \item Il lavoro del designer è per metà creativo e per metà
        \textbf{ricerca e metodo scientifico}; il game design non è una
        scienza esatta.
  \item \textbf{Documentazione} e \textbf{pitching} sono il cuore del
        mestiere.
  \item Il designer progetta meccaniche e sistemi (partendo dalle
        \textbf{3C}), documenta, media, prototipa, confuta le idee, studia i
        competitor, bilancia e fa tuning.
  \item La leadership del designer è \textbf{al servizio} del team, non
        verticale.
  \item I limiti (tempo, tecnologia, competenze, persone, denaro) restringono
        e indirizzano la creatività.
  \item \textbf{Le idee sono nulla, l'implementazione è tutto.}
\end{itemize}

%    % ...
%
%  (Con hyperref aggiungere anche \renewcommand{\theHchapter}{z.\arabic{chapter}}
%  nella parte zero, per evitare ancore duplicate.)
%  Tutte le etichette di questo capitolo hanno il prefisso "z2:"; i
%  rimandi al capitolo precedente usano le etichette "z1:" già definite in
%  GD0_1_cultura_del_feedback.tex.
% =====================================================================

\chapter[Tempo, idee e leadership]{Gestire il tempo, generare idee, guidare da dentro}
\label{cap:z2-tempo-idee-leadership}

Questo capitolo raccoglie la \textbf{seconda giornata} del workshop di due
giorni con cui si aprono le lezioni zero, tenuto dallo stesso relatore del
capitolo~\ref{cap:z1-feedback}, Stefano Cozzi\footnote{Per il cognome e per
il nome dello studio (Stormind Games) valgono le avvertenze della nota
all'inizio del capitolo~\ref{cap:z1-feedback}. In questa giornata la
trascrizione riporta soltanto il nome di battesimo, ``Stefano'', con cui il
relatore si cita più volte nei suoi esempi.}. Se la prima giornata spiegava
\emph{perché} le aziende tengono tanto alla cultura del feedback, la seconda
è dichiaratamente più pratica: ``ieri vi ho spiegato perché, oggi vi insegno
come lo fanno queste aziende''. I concetti della prima giornata (il
\emph{default to trust}, il ``potrei non farcela'' detto in anticipo, la
condivisione dei successi, il feedback frequente) ritornano di continuo,
applicati alla gestione del tempo, alle riunioni creative e alla leadership.

La giornata si è svolta di domenica, a chiusura di un fine settimana di
workshop\footnote{Il relatore dice ``oggi chiaramente è domenica'' e saluta
la classe con ``buon resto del weekend''; la prima giornata si era quindi
svolta di sabato. Più avanti parla di ``dodici ore'' di workshop in totale,
mentre nel capitolo~\ref{cap:z1-feedback} si parlava di circa dieci ore: la
differenza dipende probabilmente dal conteggio delle pause.}, con il
relatore alle prese con la voce un po' calata dal giorno prima (``la sfida
principale di oggi sarà contro la mia voce'').

Il percorso del capitolo segue quello della lezione:
\begin{enumerate}
  \item il programma della giornata (sezione~\ref{sec:z2-programma});
  \item le domande sulla prima giornata, che si trasformano in una lunga
        sessione su \textbf{colloqui e selezione}
        (sezione~\ref{sec:z2-colloqui});
  \item il \textbf{time management}: pianificare la giornata, le liste, le
        checklist, le interruzioni (sezione~\ref{sec:z2-pianificare});
  \item la \textbf{matrice di Eisenhower} per prioritizzare
        (sezione~\ref{sec:z2-eisenhower});
  \item come proteggere il proprio tempo: sprechi, distrazioni, il saper dire
        di no (anche ai meeting), il mito del multitasking e la chiusura della
        giornata (sezione~\ref{sec:z2-proteggere});
  \item le domande sul time management e una parentesi su tempo libero, ferie
        e benessere in azienda (sezioni~\ref{sec:z2-qa-tempo}
        e~\ref{sec:z2-benessere});
  \item il \textbf{brainstorming}: obiettivi, problemi tipici, un metodo
        completo (\emph{rapid ideation} e \emph{dot voting}) e la
        \textbf{matrice impatto--sforzo} per decidere
        (sezioni~\ref{sec:z2-brainstorming}, \ref{sec:z2-metodo}
        e~\ref{sec:z2-impatto});
  \item i consigli pratici, le domande sul brainstorming e il principio del
        \textbf{disagree and commit} (sezioni~\ref{sec:z2-qa-brainstorming}
        e~\ref{sec:z2-disagree});
  \item il \textbf{leading from within}: perché la leadership non è un
        titolo, le carriere divise tra manager ed esperti
        (sezione~\ref{sec:z2-leadership});
  \item il ``momento Piero Angela'' sulla biologia del gruppo e le domande
        che ne sono nate (sezioni~\ref{sec:z2-biologia}
        e~\ref{sec:z2-qa-biologia});
  \item la \textbf{piramide della fiducia}, il prendersi cura del gruppo e
        che cosa cambia quando si diventa manager
        (sezioni~\ref{sec:z2-piramide}, \ref{sec:z2-gruppo}
        e~\ref{sec:z2-manager});
  \item le domande finali e la chiusura del workshop
        (sezioni~\ref{sec:z2-qa-finale} e~\ref{sec:z2-chiusura}).
\end{enumerate}

Come per il capitolo precedente, una premessa di metodo: buona parte della
giornata è fatta di domande della classe e di racconti dell'esperienza del
relatore. Sono riportati con cura perché sono il modo in cui i principi
vengono resi concreti, e perché spesso contengono le indicazioni più
direttamente utilizzabili (soprattutto sui colloqui di lavoro).

% ---------------------------------------------------------------------
\section{Il programma della giornata}
\label{sec:z2-programma}

Come annunciato alla fine della prima giornata
(sezione~\ref{sec:z1-domani}), gli argomenti sono due:
\begin{itemize}
  \item un ``doppio topic'', \textbf{time management e brainstorming};
  \item il \textbf{leading from within}, che il relatore definisce ``la
        versione avanzata di ieri'': come guidare un team stando dentro al
        team, verso i colleghi e verso il gruppo con cui si lavora.
\end{itemize}
Poiché l'ordine è indifferente, il relatore lascia scegliere alla classe: dalle
poche preferenze espresse, la leadership va al pomeriggio, così che le domande
finali si concentrino su quella, mentre la mattina è dedicata a time
management e brainstorming. Il materiale c'è, ma il relatore avverte che la
giornata potrebbe finire prima delle quattro: non vuole trattenere nessuno se
non c'è niente di specifico da fare. Per l'eventuale tempo avanzato prima della
pausa pranzo propone un \textbf{argomento bonus}: i colloqui di lavoro,
selezione e \emph{recruitment}, visto che il relatore ha fatto e fa tuttora
selezione del personale e che per molti studenti sarà la cosa più immediata da
affrontare finito il percorso di studi (domande comuni, come rispondere,
eccetera).

Prima di cominciare, però, il relatore apre a \textbf{domande sulla giornata
precedente}: la classe ha avuto modo di ``metabolizzare'' i contenuti e
potrebbero essere emersi dubbi a posteriori. È proprio da una di queste domande
che nasce, in anticipo, la parte sui colloqui.

% ---------------------------------------------------------------------
\section{Colloqui di lavoro: capire che cosa cerca l'azienda}
\label{sec:z2-colloqui}

\subsection{L'annuncio come prima fonte di informazione}

La domanda: il giorno prima si era detto che è importante capire che cosa
l'azienda cerca davvero; ma se il primo colloquio è con le risorse umane, cioè
con qualcuno che magari ``di tecnico non sa'', come si fa a capirlo? Il
relatore risponde subito, nel dubbio che più tardi manchi il tempo.

La prima fonte è \textbf{l'annuncio}. Non viene scritto a caso: è approvato
sia dall'HR sia dall'\emph{hiring manager}\footnote{Nella trascrizione
``Airing Manager'' e ``Iron Manager''; allo stesso modo ``e-char'' e ``char''
stanno per HR.}, cioè la persona che seguirà il candidato una volta assunto,
e spesso dall'intero dipartimento HR. È calcolato e scritto in un certo modo
per far capire che cosa serve. Molti annunci moderni, oltre alla parte su
\emph{che cosa si verrà a fare}, contengono una parte su \emph{che cosa ci si
aspetta da questa persona}, cioè il comportamento e il carattere che servono
in quella posizione: ed è quella che dice molto sulla \textbf{cultura
aziendale}.

L'esempio è quello di due annunci per la stessa posizione (programmatore
senior) descritta in modo diverso:
\begin{enumerate}
  \item ``cerchiamo un programmatore esperto che possa prendersi cura dei
        nostri programmatori junior, aiutarli a crescere, sviluppare un
        ambiente di lavoro inclusivo e stimolante, in modo che i nostri
        progetti futuri siano all'avanguardia'';
  \item ``cerchiamo un programmatore esperto che sappia rispettare le
        deadline, seguire le indicazioni dei dati, sia indipendente nel suo
        lavoro e sappia relazionarsi con gli \emph{stakeholder} per capire e
        proporre nuove direzioni per il progetto''.
\end{enumerate}
Nel primo caso ci si aspetta un team di junior, probabilmente con un'età media
più bassa, con cui ci si relazionerà soprattutto; per un senior è l'occasione
di sviluppare le doti da \emph{people manager}, e al colloquio arriveranno
domande come ``hai mai lavorato con dei team? Con che età media? Ti sei mai
trovato a guidare un gruppo?''. Nel secondo caso il senior deve eseguire un
piano deciso o proporre qualcosa, e si relazionerà soprattutto \emph{verso
l'alto}; al colloquio arriverà piuttosto ``fammi questo test di
programmazione''. Sono, di fatto, due colloqui completamente diversi.

\subsection{Quanti colloqui e con chi}

Ormai è molto raro essere assunti con un solo colloquio: ne servono almeno
due, realisticamente \textbf{tre}, che è diventata la media. L'iter tipico:
\begin{enumerate}
  \item il primo con l'HR;
  \item il secondo con l'HR o con il futuro manager;
  \item il terzo sicuramente con il futuro manager, perché l'azienda vuole far
        conoscere al candidato la persona con cui lavorerà, e viceversa.
\end{enumerate}
Già le domande diverse dei vari interlocutori aiutano a capire che cosa
cercano. Ma c'è una parte fondamentale, che secondo il relatore molti
sottovalutano: il momento finale delle \textbf{``hai domande?''}.

\subsection{Le domande da fare al selezionatore}

Un principio di base: \textbf{i colloqui non sono ``devo farmi selezionare
dall'azienda'', sono ``dobbiamo selezionarci a vicenda''}. Le domande vanno
quindi preparate: ci sarà sempre qualcosa che, dalla \emph{job description} o
dal colloquio, non è chiaro al cento per cento. Nel dubbio, il relatore
suggerisce tre domande ``che funzionano sempre'':

\paragraph{``Se iniziassi domani, quale sarebbe il mio \emph{day to day}?''}
Ottiene due effetti. Il primo è pratico: descrive la giornata tipo e fa capire
se piacerebbe o no. Il secondo è più psicologico, ma il relatore sa che
funziona: costringe l'interlocutore a pensare al candidato \emph{già dentro
l'azienda}, cioè a che cosa potrebbe fare da subito e a quale impatto potrebbe
avere con le sue competenze.

\paragraph{``Se aveste una bacchetta magica, che cosa cambiereste subito
dell'azienda?''} Sposta la conversazione su di loro e va bene sia con l'HR sia
con l'hiring manager. Li obbliga a indicare una sfida interna che vorrebbero
risolvere, ed è sempre un'informazione utile: se entrambi rispondono che
cambierebbero il modo in cui si prendono le decisioni, viene da pensare che non
ci sia un percorso chiaro e trasparente da parte di chi sta in alto.

\paragraph{Lo scambio: ``vi do un motivo per non prendermi, voi datemene uno
per non venire''.} È più avanzata e il relatore non la usa sempre, ma la trova
particolarmente interessante. Permette al candidato di esprimere i propri dubbi
in modo costruttivo (``mi dite che lavorerei a stretto contatto con un team di
junior: un motivo per non prendermi è che non ho anni di esperienza nella
gestione di giovani sviluppatori, sarebbe la mia prima volta; però penso di
poterlo fare perché ho questo interesse e queste competenze''). Dall'altra
parte, il selezionatore dirà che cosa secondo lui sarebbe più difficile per il
candidato: si scopre così che cosa sta già pensando, e l'informazione torna
utile anche per i colloqui successivi.

\paragraph{E se il selezionatore si innervosisce?} ``Se un recruiter si
innervosisce per le domande del candidato, ha sbagliato lavoro.'' Al massimo il
relatore si innervosisce se un senior gli chiede che cosa faccia di preciso un
programmatore senior (``lo fai da dieci anni, lo saprai tu''); ma se un junior
chiede che cosa dovrebbe fare ogni giorno, lo si spiega volentieri. E in realtà
vale anche per i senior, perché \textbf{in ogni azienda lo stesso ruolo fa cose
diverse}: conviene sempre farsi spiegare che cosa intendono per quel titolo.
L'esempio viene dal mondo delle community: \emph{community manager} a Riot
indicava la persona a capo del team di \emph{community representative} di una
nazionalità (``community manager italiano''); a Blizzard indicava la persona
sopra l'intero dipartimento che gestiva quei team. Candidarsi per lo stesso
titolo non garantisce di aver capito per che cosa ci si sta candidando, e i
recruiter lo sanno.

Se poi un recruiter si innervosisce comunque, sta dando più risposte di quelle
che pensa: vuol dire che l'azienda non vuole domande, o non vuole interazioni
del tipo ``ma perché facciamo questa cosa?''. Che abbia formato i suoi
recruiter a innervosirsi, o che accetti che lo facciano, è informativo.

\subsection{Le domande che dicono molto sull'azienda}

Allo stesso modo, come detto il giorno prima, è informativo se un recruiter
chiede a una donna se vuole avere figli: ``ah, quindi a voi interessa sapere se
voglio avere figli: questo già mi dice che probabilmente non voglio lavorare
qui''.

Un esempio \emph{gender neutral}: la richiesta di vedere la \textbf{busta paga
precedente}. ``Quanto è la tua RAL? Mandami una busta paga.'' La risposta del
relatore è no: ``ti ho detto qual è la mia RAL; se vogliamo instaurare un
rapporto in cui non ti fidi di quello che dico, abbiamo un altro problema''. È
di nuovo \emph{default to trust} (sezione~\ref{sec:z1-default-trust}). Il
relatore è quasi sicuro che legalmente non si possa chiedere la busta paga di
un candidato\footnote{È l'opinione del relatore, espressa come tale (``sono
quasi sicuro''); non si tratta di un parere legale e non è stata verificata.}:
a quel punto diventa lui antagonista del recruiter, e visto che la
conversazione è finita in una ``zona grigia'', fa lui le domande (``come mai è
legale questa cosa che mi chiedi?''). Ha già deciso che probabilmente non
lavorerà lì, ma così magari evita che i candidati successivi debbano subire le
stesse richieste. Cita l'esempio perché sa che in Italia è abbastanza comune,
mentre all'estero non gli è mai successo, anche se avrebbe potuto ``sparare
qualsiasi cifra''.

Il motivo è che \textbf{non c'è alcun onere di dimostrare quanto si
guadagna}: chi spara troppo alto si autoesclude. Recruiter e manager non
arrivano al colloquio con carta bianca né con l'indicazione di risparmiare il
più possibile, ma con un \emph{range}: ``per questa posizione possiamo
spendere dai 30 ai 40 mila euro l'anno; paga questa persona a seconda della
sua seniority''. Se un candidato che guadagna 30 dice di prenderne 55 per
``provare a fare i soldi'' e il budget arriva al massimo a 40, il relatore non
chiederebbe nessuna prova: direbbe ``grazie, siamo in un range diverso, possiamo
salutarci''. La domanda sulla busta paga è quindi \emph{formulata male}, perché
è poco furba. Quella giusta è: \textbf{``hai già una RAL in mente per questa
posizione?''}. Si ottengono le stesse informazioni senza passare per la zona
grigia.

\subsection{La cifra minima mentale}

E se la domanda la fanno a un junior che non ha mai lavorato? Il consiglio vale
per tutti: \textbf{preparatevi sempre la vostra cifra minima mentale}, per
qualsiasi offerta di lavoro, prima ancora di fare il colloquio. Non serve
conoscere la media italiana per quella posizione o quanto paga quell'azienda:
serve sapere \emph{quanto si è disposti ad accettare senza andare in negativo},
tenendo conto dei propri costi (la benzina per andare in ufficio ogni giorno, un
eventuale trasferimento, e così via). Quella cifra va sempre avuta in testa.

\subsection{La scrematura iniziale e il portfolio}

Uno studente chiede come funzioni la risposta a una candidatura: arriva un
rifiuto diretto, oppure c'è una scrematura iniziale e poi i colloqui? E vale
per tutti i ruoli? Sì: c'è una \textbf{scrematura iniziale} fatta dall'HR o, in
alcune aziende, direttamente dal software dell'HR. Si scarta chi è chiaramente
fuori da ciò che si cerca:
\begin{itemize}
  \item si cerca un \emph{graphic designer} e arriva il CV di un programmatore
        (``faccio questi esempi perché purtroppo succede'');
  \item si cerca un \emph{animator} tecnico bravo con MotionBuilder, lo si
        scrive nell'annuncio, e il candidato risponde ``so fare tutto tranne
        MotionBuilder'': c'è un motivo tecnico per non poterlo prendere;
  \item problemi di seniority: si cerca un senior e risponde un junior con un
        anno di esperienza.
\end{itemize}
Se la figura non è già chiaramente fuori target, si fa il primo colloquio.

\paragraph{Gli anni di esperienza negli annunci junior.} Spesso gli annunci
junior chiedono come minimo tre anni di esperienza. \textbf{Se l'annuncio dice
``junior'', ignorate i requisiti di esperienza}: o sono lì per sbaglio, o sono lì
per scremare. Il secondo caso è più \emph{borderline}, ma il relatore l'ha
visto fare in alcune aziende americane: si mettono requisiti molto alti,
chiaramente irraggiungibili, per vedere chi ha la ``faccia tosta'' di mandare
comunque il CV (``non me ne frega se chiedete un \emph{bachelor's degree}, non ce
l'ho, mando lo stesso''). Così però si filtra una grossa fetta di popolazione e
si seleziona un certo tipo di carattere: si stanno cercando ``i Top Gun''. Non è
quindi necessario iniziare con uno stage: si può iniziare direttamente da
junior.

\paragraph{Chi guarda il portfolio.} Il relatore non ha mai avuto un portfolio,
ma per le figure che lo richiedono parte comunque dal CV (per escludere, ad
esempio, il designer che si candida da programmatore) e guarda il portfolio
solo a parità con altri candidati: da HR può aprirlo e dire ``bello'', ma non
può dare un giudizio tecnico. Il portfolio lo apre chi sarà l'hiring manager,
che è sempre molto più tecnico. L'ordine quindi è: l'HR guarda il CV e si salva
il portfolio; se il candidato passa la prima scrematura, CV e portfolio vanno
a un tecnico che conferma (``sì, sa fare queste cose'') o boccia (``il
portfolio indica chiaramente che non sarebbe in grado, digli pure che non
facciamo il colloquio'').

\subsection{Test pratici, da casa e con qualsiasi strumento}

\paragraph{Si può chiedere un test pratico al posto del portfolio?} Uno
studente chiede se, con un portfolio che non rispecchia le proprie capacità, si
possa chiedere direttamente un test. No: se l'azienda ha bisogno di verificare
una competenza ha già pronto un test e lo fa fare a prescindere. Il relatore,
per esempio, fa fare un test ai programmatori e lo fa correggere dal suo senior
developer; per gli artisti si chiedono animazioni di prova, chiaramente veloci.
Non chiede mai test troppo lunghi, perché non vuole far lavorare nessuno gratis:
non chiede ``fammi un pitch di progetto'', chiede uno spunto per capire come si
ragiona e come si lavora. Per il resto, \textbf{se il portfolio non è
indicativo, va aggiornato il portfolio}.

\paragraph{Test dal vivo o a casa?} Il relatore li fa fare sempre a casa. Sa
benissimo che un programmatore a casa lavorerà con due monitor, il codice da
una parte e Google o ChatGPT dall'altra. Ne ha discusso con i suoi senior
programmer: interessa che qualcuno sappia le cose a memoria, o che sappia
trovarle e fare quel lavoro a prescindere? A lui interessa che il lavoro
\textbf{venga fatto}: che il candidato abbia imparato il codice a memoria o che
in venti minuti trovi le informazioni, le metta insieme, le faccia elaborare a
ChatGPT, le corregga e consegni il codice che serve, non ha preferenze. Quando
si lavora le indicazioni saranno ``fai questa cosa il prima possibile'': se la
si sa fare anche con strumenti aggiuntivi, ben venga, ``anzi, spiegalo anche
agli altri''. Farla a memoria va benissimo, ``ma perché mettere l'\emph{hard
mode}?''. E comunque ci sono cose che non si possono delegare a uno strumento:
se arriva un codice che non funziona e va fatto il \emph{debug}, bisogna
saperlo fare e saper progettare i test. La figura più completa, per il
relatore, è anche quella che sa trovarsi le informazioni e gli aiuti.

Le prove dal vivo sono soprattutto domande di \emph{problem solving}: il
relatore vuole sentire come si ragiona e chiede espressamente di \textbf{pensare
ad alta voce}, per seguire il percorso mentale che porta alla conclusione. Sono
più domande impegnative che veri test. I test veri e propri si fanno con i
propri tempi e i propri strumenti: ``scegli il momento, ti mandiamo il test, hai
due ore per rimandarlo indietro''. Questo è l'accordo; in quelle due ore si può
fare quello che si vuole.

\subsection{Junior, mid e senior}

\paragraph{Perché ci sono così pochi annunci per figure mid?} Uno studente ha
notato su LinkedIn e sulle piattaforme di recruiting un divario notevole tra le
richieste di junior e quelle di mid. Il motivo è che per molte aziende,
Stormind inclusa, \textbf{la differenza tra junior e mid è sparita}: o si è
senior o non lo si è. Ci si è resi conto che la differenza tra un junior e un
mid è che il mid si è abituato a lavorare in quell'azienda e il junior non
ancora; ma un mid che arriva con tre anni di esperienza in un altro studio (ad
esempio Milestone\footnote{Nella trascrizione ``se ha fatto tre anni in
milestone'': si tratta verosimilmente di Milestone, studio di sviluppo
milanese.}) va comunque ``ripreso come junior'', perché si lavora in modo
diverso e deve reimparare. Un senior, invece, ha molte più competenze e le sa
applicare in un certo modo, e quella distinzione resta netta.

\paragraph{Ci si può candidare anche se il livello non coincide?} Un'altra
studentessa dava per scontato che, se un annuncio non specifica ``junior'', sia
per mid. Se non è specificato ``senior'', secondo il relatore si può mandare il
CV chiedendo di valutare. Certo, se è specificato junior e si è junior le
possibilità di essere contattati sono maggiori. Ma vale lo stesso discorso del
livello di inglese fatto il giorno prima: si cerca qualcuno che \emph{sappia
fare quelle cose}, non una certificazione. Se serve un programmatore che sappia
usare Unreal Engine~5 e sistemare la cosa X o Y, che sia mid da due anni o
junior da uno non è un problema.

\subsection{Gli strumenti di intelligenza artificiale}

Una domanda in chat riguarda l'uso quotidiano dell'intelligenza artificiale in
azienda. Il relatore sta facendo fare corsi di IA al suo team ``perché voglio che
la gente lavori meno'': vuole che si sfrutti qualsiasi strumento disponibile,
perché gli altri lo stanno facendo e non vuole restare indietro. Se un lavoro
si può fare in tre ore di meno, non c'è alcun vantaggio a farlo in modo lungo e
tortuoso. Il rischio esiste: se si fa fare tutto a ChatGPT e non ci si tiene al
passo con il proprio \emph{know-how}, quando ci sarà da risolvere qualcosa ci si
potrebbe ``impallare''. Ma sta al singolo restare aggiornato e saper fare il
proprio lavoro: \textbf{gli strumenti devono essere strumenti, non un sostituto
delle competenze}. ``Non è una partita di calcio, con regole su come si devono
fare le cose o un \emph{dress code}'': il lavoro va fatto nel modo più efficiente
possibile.

A uno studente a cui era saltata la connessione proprio in questo momento il
relatore risponde che la lezione è registrata e si può recuperare dal video.
Poi chiude la parentesi: ``la parte difficile era finire ieri in tempo, ce
l'abbiamo fatta''.

% ---------------------------------------------------------------------
\section{Time management: pianificare la giornata}
\label{sec:z2-pianificare}

\subsection{Time management o prioritizzazione?}

La domanda di partenza alla classe è: \emph{chi di voi, in questo momento, fa
time management su se stesso}, cioè si organizza la giornata all'inizio e
organizza il proprio tempo? Qualcuno risponde di sì, qualcuno ha uno
\emph{schedule} settimanale con orari da rispettare, qualcuno ``ci prova''
(``hai fatto bene a specificare \emph{ci provo}: magari non la rispetto''). Il
relatore lo chiede perché è uno di quegli argomenti in cui molti, dopo aver
visto la spiegazione, si accorgono di non fare davvero time management: fanno
piuttosto \textbf{prioritizzazione}, ad esempio una lista di cose da fare, che
è un po' diverso.

Nessuna delle tecniche è complicatissima, ma è utile sapersi muovere. Il
relatore colloca volutamente questo modulo \emph{dopo} quello sul feedback: saper
gestire il proprio tempo permette anche di dare feedback agli altri sulla loro
gestione del tempo.

\subsection{La \emph{must do list}}

Il punto di partenza è il \textbf{planning iniziale}: non improvvisare la
giornata. Improvvisando si è in balia di una grande limitazione umana: ci si
dimentica le cose, oppure ci si lascia trascinare dal task più urgente o
dall'ultimo arrivato. Esistono modi di dire come ``se arriva una mail o rispondi
subito o diventa un task per domani'', perché il ``rispondo dopo'' non arriva
mai; più avanti si vedranno modi più raffinati di gestire le zone grigie, ma la
pianificazione resta fondamentale.

La prima cosa da fare \emph{non} è la lista delle cose da fare oggi, ma la lista
delle cose che \textbf{assolutamente} vanno fatte oggi. Le due liste non devono
coincidere: la seconda richiede un livello di ragionamento in più. Sono le cose
per cui, se alle sei si stacca e non sono fatte, si resta finché non sono fatte,
perché sono imprescindibili.

\paragraph{Il primo momento per dire ``potrei non farcela''.} La \emph{must do
list} è anche il primo indicatore della giornata. Se è troppo corposa (dipende
naturalmente dalla complessità dei task), quello è il momento, ``alle 9:01 se
iniziate alle 9'', per dire al proprio manager \emph{oggi sarò in difficoltà}.
Alle 9:01 il manager ha l'intera giornata per aiutare; se invece si pecca di
ottimismo o di presunzione (``ce la posso fare'') e glielo si dice alle 17:59,
non può fare nulla. È lo stesso discorso del giorno prima
(sezione~\ref{sec:z1-triangolo}): anche facendo la to-do list si applicano i
concetti della cultura del feedback, ed è il motivo per cui questo modulo viene
dopo. Ci saranno vari momenti di questo tipo nella giornata, ma questo è il
primo.

\paragraph{E se ci si accorge a metà giornata di non farcela?} Uno studente
chiede che cosa fare se, convinto di farcela, si accorge a metà giornata di non
poter completare bene la \emph{must do list}. Il relatore lo concederebbe
\emph{una volta}: se si vogliono testare i propri limiti, va bene, ma
comunque si dice al capo ``oggi potrei essere tirato, ci provo'', e sarà lui a
valutare se il rischio è accettabile. Dipende dal rischio: se il \emph{must do}
è consegnare una milestone a un cliente, ``nel dubbio facciamolo fare a
qualcun altro, non possiamo correre quel rischio''; se il rischio è che domani
un animator riceva il codice con tre ore di ritardo, lo si può correre. ``È un
po' questo il lavoro.''

\paragraph{Cinque minuti e una stima.} Tutto questo va fatto in circa cinque
minuti. Se non è chiaro che cosa si deve fare in una giornata, \textbf{si
chiede}: ogni giorno dovrebbe essere chiaro. Poi si cerca di \textbf{stimare il
tempo} di ogni cosa. Si diventa più bravi col tempo: alla quinta volta che si
fa il report mensile si sa che ci vuole un'ora e ventidue minuti; la prima volta
si può chiedere a un collega quanto ci mette lui, sapendo che per noi potrebbe
essere diverso.

\subsection{To-do list e checklist}

Esistono molti strumenti per tenere to-do list e checklist; il relatore non ne
consiglia uno in particolare e cita come esempi Any.do, i promemoria di Windows
o Samsung Notes\footnote{Nella trascrizione ``AnyDo, Windows Memo''. Any.do è
un'app di liste e promemoria; ``Windows Memo'' indica probabilmente le note o i
promemoria integrati in Windows (Sticky Notes, Microsoft To Do).}: ``potete usare
davvero la qualunque''. L'importante è tenere la lista \textbf{visibile} da
qualche parte sul computer. Alcuni accorgimenti:
\begin{itemize}
  \item i \textbf{must do} vanno sempre \textbf{in cima} alla lista;
  \item quando durante la giornata arriva un task nuovo, la cosa da fare subito
        non è \emph{eseguirlo}, ma \textbf{aggiungerlo alla lista}: ``se arriva
        una mail, aggiungete subito alla lista \emph{rispondi alla mail}'',
        decidendo se è diventato un \emph{must do} o se è comunque un task per
        oggi;
  \item non limitarsi alle cose operative: il relatore ha iniziato ad
        aggiungere anche task come ``parla con Tizio'' o ``parla con Caio'',
        perché è da un po' che non chiede loro un aggiornamento; se non li
        mette in lista si accorge di non farli;
  \item spuntare le voci completate: per quanto minimo, rilascia
        \textbf{dopamina} (se ne parlerà nel pomeriggio,
        sezione~\ref{sec:z2-biologia}), aiuta a motivarsi e ad andare avanti.
        Sono piccole cose, ma si sommano.
\end{itemize}

\paragraph{Priorità contro recenza.} Che cosa succede se si esegue sempre
l'ultimo task arrivato? Lo chiede il relatore alla classe. Si interrompe il
flusso di lavoro e si cambia continuamente attività, certo; ma soprattutto si
stanno \textbf{mescolando priorità e recenza}: si sta dicendo ``non importa come
ho prioritizzato i task, la mia priorità è sempre l'ultimo arrivato'', che non
vuol dire fare le cose più importanti. È il motivo per cui tanti dipartimenti IT,
quando si chiede loro qualcosa, rispondono ``mi apri un ticket?''. È diventata
quasi una battuta, ma il concetto è chiarissimo: devono poter prioritizzare
quella richiesta rispetto alle altre 112 arrivate in giornata, e senza ticket
non possono confrontarla con le altre. Il relatore è totalmente a favore: ``certo
che ti apro un ticket, altrimenti tu come fai a tracciare il tuo lavoro?''.

\subsection{La checklist: l'esempio dell'aviazione}

La \textbf{checklist} non è la task list. La task list contiene le cose da fare;
la checklist contiene \textbf{i passi da seguire quando si esegue un task
specifico}. L'esempio è quello delle \textbf{checklist pre-volo}, introdotte
nell'aviazione (secondo il relatore tra gli anni Sessanta e Settanta): chi vola
spesso nota che i piloti seguono sempre un ordine simile nei controlli, tanto
che il relatore capisce quanto manca al decollo da quando controllano i flap, i
motori, quando accendono e spengono le luci. L'introduzione delle checklist
avrebbe ridotto di circa il 75\% gli incidenti in fase di decollo\footnote{Date
e percentuale sono quelle riferite in lezione e non sono state verificate.
L'origine della checklist pre-volo viene di solito fatta risalire al 1935, dopo
l'incidente del prototipo del Boeing B-17, causato da un controllo dimenticato da
un equipaggio molto esperto; la diffusione sistematica delle checklist e dei
protocolli di cabina è invece successiva.}, e soprattutto per i \textbf{piloti
più senior}. Perché? Perché chi ha fatto un decollo centomila volte non va a
controllare l'alettone destro: ha funzionato centocinquantamila volte, funzionerà
anche stavolta. Quando ci si abitua a fare qualcosa, si sbaglia proprio quando
qualcosa non va. La checklist ha obbligato a fare quei controlli ogni volta e ha
ridotto in modo drastico gli incidenti.

Il relatore usa checklist anche per cose banali, ad esempio le \textbf{email}:
essendo una persona abbastanza ansiosa su queste cose, se prepara una mail alla
leggera finisce per passare troppo tempo a ricontrollare l'oggetto, il
destinatario, l'allegato, e comunque si dimentica qualcosa. Seguire una checklist
sembra più complesso, ma \textbf{fa risparmiare tempo}: se la checklist è
completa, si sa di aver fatto tutto e si può andare avanti senza continuare a
pensare ``ho chiuso il gas? ho spento la luce?''. È come la lista della spesa: se
tutto è segnato e preso, punto. È un investimento: trenta secondi in più per
fare la cosa, invece di cinque minuti a sperare di averla fatta giusta.

\paragraph{E se ci si dimentica qualcosa nella checklist?} Succede. A quel
punto \emph{preparare la checklist} è un task in sé: un giorno in cui si ha
tempo si fa quel task pensando a quali passi si seguono e a come assicurarsi che
funzioni. Nessuna checklist è perfetta: ci saranno sempre compromessi tra ``devo
davvero controllarlo ogni volta?'' e ``posso farne a meno?''. Però, in generale,
\textbf{se qualcosa anche solo una volta ha dato problemi, va in checklist},
perché vuol dire che potrebbe darne di nuovo. Nessun manager rimprovererà
``come hai fatto a non prevedere che questa cosa poteva andare male?'': si
diventa esperti anche nelle checklist di un task facendo quel task, ed è un
processo iterativo. Se qualcosa è stato dimenticato e ha creato problemi, lo si
aggiunge per la volta dopo. Sbagliare una o due volte non è un problema;
l'importante è che la situazione migliori la volta successiva.

\paragraph{L'inizio di un junior e l'\emph{onboarding}.} A uno studente che
osserva che l'inizio di un junior è quindi inevitabilmente un po' travagliato,
il relatore risponde che è giusto così: ``se volessi il dipendente perfetto da
subito, cercherei un senior, non un junior''. Allo stesso modo non capisce le
aziende che prendono un junior e non lo fanno seguire da nessuno: che cosa ci si
aspetta, che assorba passivamente tutto quello che ha intorno e impari ``per
osmosi''? Può succedere, ma imparerà a fare le cose come vuole lui, perché
nessuno gli sta dicendo ``meglio farle così''. Moltissime aziende, oltre al
manager (che magari ha quindici \emph{report} e non ha il tempo materiale di
seguire un junior come andrebbe seguito), assegnano un \textbf{referente}, di
solito un senior, con il compito di prendersi cura della persona nuova e
insegnarle a fare le cose. Altrimenti l'\textbf{onboarding}, cioè il tempo
necessario per arrivare al cento per cento nel proprio lavoro, dura un anno.
Più risorse l'azienda dà all'inizio, prima la persona diventa operativa: è
nell'interesse dell'azienda dare il massimo supporto a chi arriva.

Il relatore diffida delle aziende che ``vogliono vedere se uno se la cava''. È
il discorso dei cento metri e della corsa a ostacoli\footnote{Il relatore lo
introduce come un discorso già fatto (``si torna al discorso: vuoi fare i cento
metri o la corsa a ostacoli?''), ma il passaggio non compare nella trascrizione
della prima giornata riportata nel capitolo~\ref{cap:z1-feedback}.}: non si può
far correre uno sui cento metri e l'altro con gli ostacoli e aspettarsi tempi
simili, perché chi ha gli ostacoli ci metterà per forza di più. Il supporto ai junior è un
investimento sensato.

\subsection{\emph{First contact resolution}}

Un concetto che viene dai call center: \textbf{quante richieste vengono risolte
al primo contatto}. È importante quando un task richiede iterazioni o
comunicazione con altri: bisogna organizzarsi come se si dovesse risolvere la
cosa \emph{in una sola passata}. Se viene chiesto un modello 3D, lo si fa giusto
la prima volta, al meglio delle proprie possibilità. Se si deve parlare con
qualcuno, si raccolgono tutte le richieste in una.

L'esempio è quello dei ticket che il relatore mandava ai colleghi del supporto di
Blizzard: ``ciao, ho un problema, non riesco a loggarmi'' --- ``mi dici lo
username?'' --- ``è questo'' --- ``mi dici l'email?'' --- ``è questa'' --- e così
via, con un paio d'ore tra un messaggio e l'altro perché il ticket deve fare il
suo giro. Tutto poteva essere riassunto in ``ciao, mi servono username, email e
questo e quest'altro''. Unificare le richieste su una singola cosa aiuta a
risolverla il prima possibile; bisogna avere chiaro \emph{che cosa serve prima
di fare una richiesta}, in modo che ce ne sia una sola. In molti sistemi di
supporto, se si apre un secondo ticket per aggiungere informazioni al primo, il
primo, anche se era quasi in lavorazione, viene rimandato in fondo alla coda:
serve a disincentivare chi manda continuamente ticket del tipo ``mi sono
ricordato questa cosa''. O si manda tutto subito, o ci vorrà tempo.

\subsection{Non iniziare senza le informazioni}

Corollario: \textbf{se non si ha tutto quello che serve per completare un task,
non lo si inizia}. Prima ci si procura quello che manca. Al relatore capita con
la pianificazione dei due anni di lavoro successivi: se non ha tutte le
informazioni, prima lavora per averle (``mi dai questa informazione, mi dai
quell'altra''), poi pianifica; altrimenti pianificherebbe qualcosa di cui non ha
completa visibilità.

\subsection{Il costo delle interruzioni}

Alternare continuamente l'attenzione tra i task e, soprattutto, reagire di
continuo agli input esterni (la mail, il messaggio, il collega che parla) fa
\textbf{sparire il tempo dalla giornata}. Se mentre si prepara un report in
Excel si viene interrotti (dal gatto, da una mail, da un collega, da una
telefonata) e poi si torna al foglio, servono in media \textbf{un quarto d'ora}
per tornare mentalmente al cento per cento; chi è particolarmente bravo ci mette
dieci minuti, ma si parla comunque di dieci-quindici minuti\footnote{Ordini di
grandezza come riportati in lezione. Le ricerche più citate sul tema (ad esempio
quelle di Gloria Mark sulle interruzioni nel lavoro d'ufficio) indicano tempi
medi di ripresa anche superiori, intorno ai venti minuti.}. Se si viene
interrotti di continuo, si passa più tempo a cercare di tornare al cento per
cento che a lavorare. È una partita persa in partenza.

Per questo il giorno prima si era parlato dei \textbf{semafori} sulla scrivania
(sezione~\ref{sec:z1-dare}): il semaforo può essere anche semplicemente il
proprio stato su Slack, su Discord o su qualunque strumento si usi con i
colleghi. Se si è occupati, si è occupati: è una propria valutazione che,
accettando di essere disturbati di continuo, quella cosa non si riuscirà a
farla, o almeno non in tempo.

% ---------------------------------------------------------------------
\section{Prioritizzare: la matrice di Eisenhower}
\label{sec:z2-eisenhower}

Ora che si hanno task list, checklist, \emph{must do} e \emph{nice to do}, resta
la domanda di prima: se arriva un task nuovo e si decide di non fare per prime le
cose nuove, \emph{come si capisce se è davvero urgente}? Serve un metodo, ed è la
\textbf{matrice di Eisenhower}\footnote{Nella trascrizione anche ``Jason
Howard''. Dal nome del presidente statunitense Dwight D.~Eisenhower, a cui si
attribuisce la distinzione tra urgente e importante.} (la classe prova a
indovinarne il nome proponendo altre matrici, come una matrice dei rischi, ma è
un'altra cosa). Ogni task che arriva va posizionato su questa tabella:

\begin{center}
  \begin{tabular}{@{}p{0.2\textwidth}p{0.36\textwidth}p{0.36\textwidth}@{}}
    \toprule
    & \textbf{Urgente} & \textbf{Non urgente} \\
    \midrule
    \textbf{Importante} &
    \textbf{Fallo} (\emph{do}): diventa un task di oggi. &
    \textbf{Pianificalo} (\emph{schedule}): va fatto, ma non subito. \\
    \addlinespace
    \textbf{Non importante} &
    \textbf{Delegalo} (\emph{delegate}): va fatto, ma non serve che lo faccia
    tu. &
    \textbf{Eliminalo} (\emph{delete}): non lo fai. \\
    \bottomrule
  \end{tabular}
\end{center}

\paragraph{Urgente e importante: lo fai.} Sono di solito task con una
\textbf{deadline molto chiara} (entro oggi, entro un'ora, entro X) e con
\textbf{conseguenze} se non vengono fatti: da qui l'``importante''.

\paragraph{Importante ma non urgente: lo pianifichi.} Va fatto, ma non subito:
o lo si mette tra le cose di oggi ``ma dopo'', o anche alla settimana successiva.
L'esempio: un collaboratore scrive ``il mio rinnovo è tra tre mesi, ne possiamo
parlare?''. È importante, perché il rinnovo di una persona è importante; non è
urgente, se ne può parlare anche domani, così da potersi preparare e fare altre
cose nel frattempo.

\paragraph{Urgente ma non importante: lo deleghi.} Qui ``non importante''
\textbf{non} vuol dire che il task non sia importante, ma che \textbf{non è
importante che lo faccia tu}, che è molto diverso. Se arriva ``mi serve questo
cambio a questa animazione 3D'' e ciò che conta è il cambio, non chi lo fa, si può
dire ``sono d'accordo che è urgente, ma non serve che lo faccia io'' e chiedere a
qualcun altro. Sono task che vanno fatti ma non richiedono la propria particolare
\emph{expertise}.

Il relatore insiste molto su questo punto con manager e senior: quando arriva un
task urgente, i senior tendono a dire ``lo faccio io, perché è urgente''. Ma se un
junior può farlo, lo si fa fare al junior e si investono le competenze più alte
del senior in qualcos'altro: il senior può fare il lavoro del junior, ma il
junior non può fare il lavoro del senior mentre il senior fa quello del junior.
Le persone non sono intercambiabili. Si usa il \textbf{minimo sforzo
necessario}: se una cosa la può fare un mid, la fa il mid, non il senior.

\paragraph{Né urgente né importante: lo elimini.} È il quadrante che quasi
nessuno usa: tutti fanno, delegano o pianificano, ma non pensano mai a
\emph{delete}. Invece è fondamentale. A volte è corretto dire: ``questa cosa non
dovrei farla io, non è urgente, non la faccio; sappi che non la faccio, e se c'è
un problema ne parli con il manager''. Non si sta delegando a qualcun altro: si
sta dicendo che quella cosa così com'è \textbf{non andrebbe fatta}, perché non
ne sono chiari l'importanza e il perché. Nella propria lista di oggi non ha
spazio, e non è abbastanza importante per essere fatta da sé nemmeno nei giorni
successivi.

Un esempio di \emph{delete} che non è un vero \emph{delete}: arriva una richiesta
di meeting di tre ore sul futuro del progetto, con invitati tutti gli art
director, il game director e così via. Il relatore spesso risponde ``io non
vengo'': per lui è \emph{delete}, ma il meeting si fa lo stesso. Uno studente
osserva che a volte si tratta di una richiesta mal formulata o poco chiara: è
esattamente così. ``Nella sua forma per me è no; se poi mi spieghi meglio perché
dovrei farlo io o perché è urgente, allora cambia.''

% ---------------------------------------------------------------------
\section{Proteggere il proprio tempo}
\label{sec:z2-proteggere}

\subsection{Evitare gli sprechi: un limite che nessuno può insegnare}

Evitare gli sprechi di tempo è difficile, soprattutto da remoto: è una sfida
costante con se stessi. Ha molto a che fare con il carattere, il
\emph{self-control} e (``passatemi anche questa'') la professionalità. Il
relatore può dare tutti i consigli del mondo, ma se dentro di sé non si percepisce
l'urgenza di fare il proprio lavoro quando lo si deve fare, non lo si può
insegnare. Se si è a casa e invece di lavorare si gioca a qualcosa su Steam,
nessuno deve spiegare che non si dovrebbe. Ci sono limiti oltre i quali non si
può andare.

Quello che si può fare è sapere che \textbf{ogni interruzione ha un costo
sommerso}: ogni messaggio, domanda, telefonata, ogni cambio di focus. Sono i
quindici minuti di prima, che possono arrivare a trenta a seconda del tipo di
task.

\paragraph{È più difficile tornare su un task logico o creativo?} Il relatore lo
chiede alla classe: fa più fatica un artista a tornare sul suo modello 3D o un
\emph{data analyst} a tornare sul suo foglio Excel? Le risposte si dividono
(``creativo'', ``entrambi''), ed è proprio il punto: dipende a chi lo si chiede. I
creativi rispondono ``creativo'', gli analisti ``analitico'', perché ognuno ha una
chiara comprensione di che cosa serve per essere davvero concentrati su ciò che
fa. Il creativo sa che per essere pienamente creativo gli servono l'ambiente, la
concentrazione, gli spunti, gli esempi da tenere chiari in testa; l'analista deve
avere in testa i conti, i fogli di dati e così via. Sono \textbf{ugualmente
difficili}, ma solo chi fa quella cosa ne apprezza fino in fondo la difficoltà. Si
può avere un assaggio dell'altra parte, ma la piena comprensione arriva solo
facendo: è come guardare uno sport in tv e pensare ``potrei farlo anch'io'',
finché non ci si trova lì.

\subsection{Distrazioni volontarie e involontarie}

Uno studente, che sta leggendo manuali di \emph{design thinking}, chiede come
conciliare tutto questo con l'idea che distrarsi, cioè non restare focalizzati,
aiuti a trovare un nuovo punto di vista per risolvere un problema. Qual è il
giusto equilibrio? Il relatore è \textbf{totalmente d'accordo} con quel concetto:
sta parlando di \textbf{interruzioni non volontarie}. Sa benissimo che dopo cinque
ore su un foglio Excel ``esplode la testa''; ma se si pianifica un quarto d'ora di
pausa e lo si rispetta, si sta investendo sulla propria creatività ed
efficienza. Non è un investimento dover rispondere a una chiamata ogni tre
minuti.

Anzi, il consiglio è \textbf{crearsi interruzioni pianificate}, sotto il proprio
controllo: fare pause, fare altro, soprattutto nel \emph{problem solving}. Il
relatore stesso, se non trova una soluzione, si ferma, fa un altro task e poi ci
torna; ma è una sua scelta. La discriminante è che la pausa sia \emph{scelta} e
non imposta: anche se si fanno le pause giuste, se si viene interrotti proprio
nel momento di massima concentrazione (raggiunta magari grazie alla pausa) si
butta via tutto. Uno studente cita la \textbf{tecnica del pomodoro}: è
esattamente quello, piccole pause pianificate che aiutano a essere creativi,
purché restino pianificate. Se i quindici minuti diventano ``vabbè, facciamo
mezz'ora'' o ``lo faccio nel pomeriggio'', non è più la tecnica del pomodoro:
diventano interruzioni, o un pretesto per non lavorare.

\subsection{Il peggior nemico siamo noi}

E chi sono i peggiori nemici in questo? ``Noi stessi'': siamo noi ad
autoconvincerci che venti minuti vanno bene quanto quindici. Le
\textbf{distrazioni autoinflitte} sono spesso le peggiori. Il relatore porta
l'esempio di una sua giornata, che ha annotato perché lo aveva fatto ridere
quanto tempo era riuscito a buttare via senza che nessuno glielo imponesse:
\begin{itemize}
  \item gli piace avere qualcosa in sottofondo mentre lavora (``ha a che fare con
        l'ADHD sicuramente''\footnote{Nella trascrizione ``IDACD''. Più avanti il
        relatore dice di avere l'ADHD ``come dicevo ieri''; il capitolo 0.1 non
        lo riporta, ma riferisce il suo racconto sull'``overthinking'' anche
        nell'ordinare una pizza (sezione~\ref{sec:z1-regalo}).}),
        e in genere lo aiuta. Il problema è quando il sottofondo è troppo
        interessante: la musica va bene, ma se nella playlist di Spotify parte un
        podcast di Alessandro Barbero sulla guerra civile americana, finisce per
        ascoltare quello più del lavoro;
  \item un'abitudine che ha poi perso: se faceva straordinari alle due di notte,
        da solo, per qualcosa di urgentissimo, apriva comunque Slack a
        controllare i messaggi, e la mail. ``Come se alle due di notte qualcuno
        mi scrivesse qualcosa di urgente'';
  \item la pausa con un po' di TikTok, dopo la quale si accorgeva di aver
        dimenticato su quale slide della presentazione stava lavorando.
\end{itemize}
La regola: \textbf{inserire i break è meglio che subirli passivamente}. Un conto
è dire ``faccio una pausa intenzionalmente'', un conto è ``mi sono distratto e ho
perso due ore di lavoro''. Chi ne è particolarmente soggetto non deve
demoralizzarsi: è anche un allenamento.

\subsection{Spezzettare i task grandi}

Un altro concetto utile: se si fa molta fatica a iniziare i task grossi, perché
sembrano infiniti (``più tardi inizio, più tardi inizia questa sofferenza''),
\textbf{vanno spezzettati}. Un task grosso non è quasi mai un task grosso: è
spesso un insieme di task piccoli. Si comincia a farne dei pezzi e si va avanti
piano piano, spuntandoli. Aiuta tantissimo anche nel \emph{problem solving}: un
problema gigante va scomposto in problemi più piccoli, da risolvere per primi, per
poi ricomporre il puzzle.

\subsection{Accettare di dire di no}

Tornando alla matrice di Eisenhower: bisogna \textbf{accettare di dover dire di
no}, perché dire sempre sì porta a problemi ben peggiori di quanto si pensi.

\paragraph{Il collega che diceva sempre sì.} A Riot il relatore aveva un collega
estremamente disponibile, amichevole e collaborativo: qualsiasi cosa gli si
chiedesse, rispondeva ``certo, la faccio io'' (``mi porti fuori il cane?'' ---
``certo''). Il problema è che era umano: a un certo punto si supera la quantità di
cose umanamente possibili in una giornata. In più tendeva a offrirsi volontario
per far vedere che era intraprendente e aveva iniziativa, e nei progetti che
coordinava il relatore doveva spesso rifiutare la sua candidatura perché era
sicuro che non ce l'avrebbe fatta. Il suo dire sempre sì, a un certo punto, lo ha
danneggiato: era quello che dice sempre che lo fa lui, ma non è detto che lo
faccia. ``Ce la fai?'' --- ``Sì, sì, assolutamente'' --- e poi ``scusa, non ce l'ho
fatta''. ``Ok, però ora hai messo nei guai me e il resto del team.'' Quella
disponibilità era più dannosa di un ``no, non posso''. È lo stesso discorso del
giorno prima: dire all'ultimo momento che non ce la si fa è peggio che dirlo a
inizio settimana.

Uno studente lo paragona alla \emph{sindrome dello yes man}. Non proprio: lo
\emph{yes man} dà sempre ragione all'altro perché è avverso al conflitto e non vuole
la discussione del ``spiegami perché sei in disaccordo''. Il collega, invece,
pensava davvero di poter fare molto più di quanto potesse, e rifiutava l'idea di
non riuscirci.

\paragraph{``Ci hai pensato bene?''} Vale per i progetti e per le deadline. Con i
junior soprattutto, quando a ``vuoi fare questa cosa?'' arriva un sì troppo
convinto, il relatore chiede sempre: ``ci hai pensato bene? Me lo stai dicendo
per farmi contento? Sappi che se mi dici di no non penso peggio di te; penso
peggio di te se mi dici di sì e poi non ce la fai. Vuoi rivalutare la risposta, e
magari dirmi \emph{non lo so} invece di un sì categorico?''.

\paragraph{Che cos'è una sconfitta.} Uno studente chiede se davanti alle
``sconfitte'' non si dovrebbe rimediare. Il relatore chiarisce che per lui la
sconfitta non è il task non fatto: quella è una sconfitta \emph{a breve termine},
si troverà un modo di farlo, è un danno ``ma chi se ne frega''. La vera sconfitta,
un danno \emph{a lungo termine}, è aver pensato di riuscirci e non esserci
riusciti: se non si risolve il motivo per cui si è detto sì quando era no, lo si
ripeterà altre cinque, dieci, quindici volte. Nella testa del collega di Riot c'era
sempre una scusante (``c'erano questi imprevisti che mi hanno bloccato''), come
quella con cui ci si giustifica le interruzioni. Il relatore gli rispondeva: ``se
ogni volta c'è un imprevisto, forse c'è un \emph{pattern}: il problema non è che
non riesci a fare le cose, è che non prevedi mai gli imprevisti''. Dopo sette
ritardi su dieci dovrebbe scattare l'ingranaggio mentale (``forse sbaglio io a
dire sempre di sì''), ma non è sempre facile farlo scattare, e a volte serve una
spinta. È una delle situazioni più frequenti. E la cosa più difficile è che,
quando glielo si fa notare, queste persone lo prendono come un ``non sei capace''.
Il messaggio è invece: \textbf{``sei bravo, ma sei umano anche tu''}. Non si può
fare un numero infinito di cose in un'ora: bisogna accettare i propri limiti ``dati
dalla carne'' e non pensare di essere un cyborg.

\paragraph{Il cliente non ha sempre ragione.} Uno studente osserva che vale anche
per i project manager e per il principio ``il cliente ha sempre ragione''. Anzi:
il relatore ha lavorato con molti clienti e a volte \textbf{il cliente non sa che
cosa vuole}, che è molto diverso dall'avere sempre ragione. Soprattutto con i
clienti che chiedevano pitch creativi, la sua squadra presentava sì un'opzione
con la loro richiesta, ma anche le proprie (``se vuoi raggiungere questo, sarebbe
meglio fare così''), e molto spesso il cliente sceglieva quella proposta dal team.
Spesso il cliente ha in mente \emph{una soluzione}, ma bisogna capire \emph{quale
obiettivo} ha in mente e poi dargli una soluzione: non vuole la soluzione, vuole il
\emph{goal}, e se c'è un modo migliore per arrivarci, forse è meglio un'altra
soluzione.

\subsection{Dire di no ai meeting}

La prima cosa a cui il relatore ha imparato a dire di no sono i \textbf{meeting}.
Ha passato anni in ambienti in cui il suo calendario era letteralmente solo
meeting dalla mattina alla sera: estenuante per lui e inutile per il team. Tra i
fattori c'è il \textbf{FOMO}: tanti manager vogliono essere in tutti i meeting per
non perdersi niente ed essere sempre informati su quello che fa il team. In realtà
vanno in \emph{burnout}, passano i pomeriggi nei meeting e finiscono per essere i
più distratti, quelli che non ascoltano e che fanno domande inutili solo per far
sentire che fanno domande, senza aiutare davvero.

Prima di accettare un invito il relatore si fa una serie di domande, \textbf{come
una checklist} (e le fa anche agli altri):
\begin{enumerate}
  \item \textbf{Sono davvero necessario?} Servo davvero io, o se mancassi non
        succederebbe niente? Ai meeting ``di cortesia'' non va: se pensa di
        essere necessario lo dice lui, e il fatto che si faccia un meeting
        importante e urgente senza di lui non danneggia il suo ego. ``Non
        aggiungermi solo perché pensi che serva che io venga aggiunto.''
  \item \textbf{Uno di noi potrebbe partecipare e informare il resto del team?}
        C'è già qualcuno del suo team? Se sì, non può andare solo lui? Per un
        meeting informativo (una nuova policy, una nuova procedura da spiegare a
        tutti i dipartimenti), andarci tutti significa investire tre ore a
        testa; se ci va uno solo e poi spiega agli altri, si risparmia tempo.
  \item \textbf{Il riassunto sarà reso disponibile?} Per il relatore è
        obbligatorio. Per lo stesso motivo per cui la lezione viene
        registrata: se ci si perde un pezzo lo si può rivedere, se qualcosa non è
        chiaro lo si può risentire. Ma c'è anche un motivo di \emph{office
        politics}: se in un meeting si chiede ``mi fai questa cosa per venerdì?''
        e la risposta è sì, non si vuole sentire venerdì ``intendevo venerdì
        prossimo'' o ``non ero sicuro''. Nelle note del meeting c'è scritto che hai
        detto sì. ``Lo so che è un po' paraculo, ma in realtà è così che si
        lavora'': bisogna dare più peso alla parola data, e quindi tenerne
        traccia. Finché tutto resta orale, niente rimane scritto e tutto è
        interpretabile. Oggi è anche facile: Zoom, con l'assistente IA, produce
        in automatico trascrizione e sottotitoli, senza bisogno di qualcuno
        addetto alle \emph{meeting minutes}. In ogni caso, se in un meeting viene
        presa una decisione, il relatore vuole che sia \textbf{ripetuta per
        iscritto}: non deve ipotizzare, deve poter andare a rileggere.
  \item \textbf{Il meeting aggiunge valore rispetto al documento?} Se arriva un
        invito con allegato un documento di venti pagine, il relatore chiede: ``il
        meeting consiste nel leggermi il documento? Perché altrimenti non
        vengo''. Se non c'è un valore aggiunto, se lo legge da solo e manda
        conferma di averlo letto. È lo stesso motivo per cui si sconsigliano le
        presentazioni in cui si leggono le slide; e infatti, osserva, in quel
        momento lui stesso stava leggendo una slide, ma aggiungendo informazioni.
        ``Se tutta la lezione fosse io che vi leggo le slide, non servirei io: vi
        manderei le slide.''
  \item \textbf{Per che cos'altro potrei usare questo tempo?} È forse la più
        importante. Anche ammesso che serva, che serva proprio lui e che il
        meeting sia importante, in quel momento dovrebbe fare qualcosa di più
        importante? Se sì, fa quella. Non bisogna avere paura di dire ``no, non
        vengo'' anche se si serve, oppure ``ho un'altra cosa importante,
        spostiamo il meeting''. Il relatore si è sforzato di essere più
        \textbf{assertivo}: ``capisco la richiesta, ma ho qualcosa di più urgente
        da fare in questo momento''.
\end{enumerate}

\paragraph{E le richieste a voce?} Uno studente chiede se una richiesta fatta a
voce da un senior (``devi fare questo'') vada accettata così o fatta mandare per
mail o per messaggio. Accettarla a voce ci può stare. Ma se la richiesta è ``devi
farla in questo modo particolare e preciso'', è meglio farsi mandare il modo
preciso \textbf{per iscritto}: non si deve lasciare a interpretazione \emph{come}
fare quella cosa, per non ricadere nel ``ah, ma io non la volevo così''. Oppure si
conferma: ``sì, la faccio; va bene se la faccio così, così e così?''. Se dice sì,
si è a posto.

\paragraph{Non dovrebbero pensarci gli organizzatori?} Dalla chat: queste domande
non dovrebbero farsele già gli organizzatori dei meeting? Sì, sarebbe meglio che
tutti lo facessero, e il relatore stesso se le fa quando invita qualcuno: ``è
davvero necessario Marco? È davvero necessario Andrea?'', e solo se sì li invita.
Ma lavora in aziende anche di cinquecento dipendenti, dove non è detto che tutti
conoscano o applichino questa logica; per questo la applica per sé, sperando che
la applichino tutti. Ha sempre chiarito con tutti che da lui non arriveranno mai
inviti di cortesia o dettati dalla gerarchia (``deve esserci anche il mio capo''):
``devi esserci tu; se vuoi portarti il capo, portatelo, ma non lo invito solo
perché qui si vuole sempre il capo presente: sono sicuro che avrà di meglio da
fare''.

\paragraph{E i task?} Dire di no vale naturalmente anche per i task: si applica la
matrice di Eisenhower (sezione~\ref{sec:z2-eisenhower}) e si decide se fare,
delegare, pianificare o cancellare. Non c'è molto da aggiungere.

\subsection{Il multitasking non esiste}

``Questo potrà urtare tanti ego'', ma \textbf{il multitasking non esiste}. Si può
essere più bravi o più portati di altri a fare più cose contemporaneamente; la
realtà è che, concentrandosi su più cose, \textbf{la somma non farà mai cento}. Se
si lavora su due cose non si lavora al 50\% su una e al 50\% sull'altra, ma al 40\%
e al 40\%, e il 20\% va perso in ``costo di manutenzione'', cioè nella capacità di
seguirle entrambe. Il consiglio è dedicarsi al cento per cento a un task, poi
passare al successivo, poi al successivo, invece di farne cinque insieme: si può
avere la percezione di andare più veloci, ma non è così.

Uno studente cita una lettura di neuroscienze secondo cui il cervello è in grado
di fare una sola operazione alla volta. Il relatore conferma: smettiamola di
parlarne come di un pregio o di una \emph{skill}. È un modo che alcuni cervelli
trovano per svagarsi, l'equivalente di accendere un podcast: fare due cose
insieme vuol dire cercare di distrarsi alternandole, e farle entrambe peggio di
quanto si farebbero una alla volta. A chi preferisce fare tante cose insieme
suggerisce di provare a farne una alla volta: le farà più velocemente. Anche il
semplice passaggio da una cosa all'altra non è istantaneo: ci sono sempre una
decina di secondi in cui il cervello si chiede ``che cosa stavo facendo?''; per
cinquanta cambi, già questo spiega la perdita di velocità. È una questione di
\textbf{bandwidth mentale}, un concetto che tornerà nel pomeriggio: anche lo
\emph{switch} occupa banda.

\subsection{La fine della giornata}

A fine giornata si \textbf{rimette a fuoco la giornata}: si riguardano la task list
e la \emph{must do list} e si conferma di aver fatto tutto. È letteralmente il
controllo prima di uscire di casa: chiavi, porta chiusa, cellulare, portafoglio.
Al relatore è capitato di dimenticare un \emph{must do} rimasto in lista per tutto
il giorno, ed entra in gioco quella che chiama la \textbf{sindrome del
tirannosauro}\footnote{Il riferimento è alla celebre scena di \emph{Jurassic Park},
in cui si dice che il tirannosauro vede la preda solo se si muove.}: se una cosa è
lì, non è detto che la si veda finché non si muove. Averla in lista tutto il giorno
non vuol dire averla percepita come da fare.

Può anche capitare di non farcela perché quel giorno sono arrivate così tante cose
urgenti e importanti da rendere inutili le stime: la giornata no capita e non c'è
nessun problema. Ma se era davvero un \emph{must do}, \textbf{va fatto comunque}:
è lì che scattano la richiesta di straordinario o il ``devo farlo comunque''.
Attenzione: se si chiude la giornata con un \emph{must do} ancora aperto e si
decide ``vabbè, lo faccio domani'', può starci, ma allora \textbf{non doveva essere
in quella lista}, perché era rimandabile.

Al relatore questo controllo serve anche ``per quieto vivere'': se chiude ogni
giornata con la \emph{must do list} completa, sa che magari non ha fatto tutto
quello che c'era nella to-do list, ma i \emph{must do} sì. ``Il mio, oggi, l'ho
fatto'': può stare tranquillo senza continuare a chiedersi se ha fatto quello che
doveva.

% ---------------------------------------------------------------------
\section{Domande sul time management}
\label{sec:z2-qa-tempo}

\paragraph{C'è un numero medio ``sano'' di \emph{must do} quotidiani?} No, perché un
conto è avere dieci \emph{must do} come ``rispondere a dieci mail'', un conto è
averne uno solo come ``inventarmi un nuovo gioco''. Dipende dalla dimensione: se ne
gestiscono di più se sono minuscoli, di meno se sono giganteschi. Ci sono lavori
che per definizione hanno tantissimi \emph{must do} quotidiani e lavori che ne hanno
uno al giorno. Si imparerà da sé qual è il numero giusto.

\paragraph{C'è un modo per stimare quanto si riesce a fare in una giornata?} ``Lo so
che non è una bella risposta'', ma l'unico modo è la propria \textbf{storia}:
sapere quanto si riesce a fare nel quotidiano. È normale e fisiologico che le
prime stime siano sbagliate, e probabilmente eccessivamente ottimistiche;
l'importante è rendersene conto. È letteralmente la \emph{seniority}: abituarsi a
fare qualcosa insegna quanto ci si mette. Va benissimo metterci di più all'inizio,
ed è normale che il tempo cambi man mano che si migliora: ciò che oggi si stima in
un'ora, tra sei mesi lo si stimerà in mezz'ora. Chi ha lavorato con Scrum lo sa:
su un progetto nuovo tutti sbagliano la dimensione delle \emph{stories} all'inizio,
perché nessuno sa quanto ci vorrà, e al terzo \emph{sprint} consecutivo è più
chiaro. Le stime sono stime, altrimenti si chiamerebbero ``inserire quanto tempo ci
metterai''.

Se però si dà una stima e poi il lavoro si rivela il triplo, il problema non è la
stima: è che \textbf{bisognava chiedere più informazioni} prima di darla. La
risposta giusta non è la stima, ma ``per darti una stima accurata mi serve sapere
questo, questo e questo''. Il problema non è sbagliare la stima, è non darsi i mezzi
per azzeccarla.

\paragraph{Come si fa a far entrare questi metodi nella routine? È solo
determinazione?} È un misto. Il relatore risponde da persona con ADHD, con la
tendenza a distrarsi e a fare un mucchio di \emph{overthinking} (``come dicevo ieri,
anche ordinare la pizza è overthinking''). È un po' disciplina, ma soprattutto far
diventare qualcosa un'\textbf{abitudine}: lo sforzo per fare le checklist era molto
più alto un anno prima, oggi è un'abitudine. Di solito servono \textbf{circa tre
mesi} perché qualcosa diventi un'abitudine\footnote{Stima come riportata in lezione.
Gli studi più citati sulla formazione delle abitudini indicano tempi molto variabili
da persona a persona, con una mediana intorno ai due mesi.}: la teoria dice che chi
inizia oggi ad andare in palestra ogni giorno, dopo tre mesi non penserà più ``oggi
devo andare in palestra'', ma la palestra farà parte della sua giornata. Non è più
uno sforzo, ma ``è così che lavoro''.

\paragraph{Chi fa davvero time management?} Prima di passare al brainstorming il
relatore torna alla domanda iniziale: qualcuno, ora che sa che cosa intendeva,
cambierebbe risposta? Uno studente risponde che lo fa, ma senza checklist, perché
non ne sente il bisogno (non gli capita spesso di dimenticare le cose); la
\emph{must do list} invece la segue sempre, e le cose meno importanti capitano nelle
giornate in cui ne salta qualcuna. Va benissimo: non esiste un unico modo, ognuno
troverà il proprio. Il relatore porta il suo esempio perché sa che è servito a molte
persone con cui ha lavorato; con chi faceva fatica ad applicare un metodo (come
chiedeva prima uno studente) ne hanno trovati altri, perché ciò che funziona per uno
non è detto che funzioni per un altro. L'importante è che \textbf{ognuno trovi il
proprio metodo}, perché la soluzione non è non fare time management: ci si sentirà
più ``liberi'' improvvisando la giornata, ma non la si sta sfruttando a dovere.

E non si parla dei video sul \emph{mindset} in cui ci si alza alle cinque, si fa la
checklist e si fa colazione leggendo le mail: organizzarsi meglio la giornata
\textbf{dà più tempo libero}. E per il relatore il tempo libero delle persone è
fondamentale, perché è ciò che permette loro di lavorare meglio.

% ---------------------------------------------------------------------
\section{Una parentesi: tempo libero, ferie e benessere}
\label{sec:z2-benessere}

\subsection{Le tre condizioni del relatore}

Il relatore è un sostenitore della \textbf{settimana lavorativa di quattro giorni},
che ha applicato in varie aziende; dove non ci è riuscito, ha introdotto almeno
\textbf{venerdì liberi}, due volte al mese o una volta al mese, cioè giornate in
cui non si fa niente. Come le pause brevi sono importanti per essere creativi,
anche le pause lunghe, cioè giornate intere, lo sono per essere creativi ed
efficienti. È lo stesso motivo per cui, dopo la Games Week, dava una settimana
libera alle persone (sezione~\ref{sec:z1-fasi-team}): ha bisogno che siano
riposate, e quella settimana ``persa'' è in realtà un investimento per le
settimane successive. In chat qualcuno lo paragona al ``metodo finlandese'': è un
misto tra quello e il metodo del buonsenso.

Il relatore viene spesso chiamato a \textbf{cambiare le culture aziendali}, ed è un
po' il suo focus. Come un allenatore che accetta una squadra solo se può portare il
suo staff (il preparatore, il tecnico), lui chiede tre elementi:
\begin{enumerate}
  \item possibilmente \textbf{quattro giorni lavorativi}, o altrimenti giornate
        libere per dare spazio alle persone;
  \item \textbf{ferie illimitate}: nessun limite massimo, ma un \textbf{limite
        minimo garantito}. Ferie illimitate non devono voler dire poter scegliere
        di non prenderle: ``devi prendere almeno X giorni, fino a infinito'';
  \item la \textbf{feedback culture}.
\end{enumerate}
Senza questi tre elementi non garantisce il risultato: è un modo di gestire le
aspettative fin dall'inizio, come si diceva il giorno prima.

\paragraph{La feedback culture.} Il perché è stato spiegato nella prima giornata, ma
il relatore lo sottolinea: un ambiente in cui ci si può fidare del team vale
moltissimo. Quando va a fare lezione allo IULM, anche in settimana, non deve
preoccuparsi di che cosa stia facendo il team mentre lui non lavora (``oggi è
domenica e non ho questa preoccupazione a prescindere, ma anche se fosse mercoledì
non l'avrei''): sa che stanno facendo del loro meglio, e non gli interessa sapere
altro.

\paragraph{Ferie e permessi illimitati.} Nell'esperienza del relatore, sia a Riot sia
a Blizzard, avere ferie illimitate significava spesso \emph{prenderne meno} di quelle
obbligatorie: c'è una vocina che dice ``ti serve davvero un giorno di ferie oggi?
Tanto sono illimitate, te le prenderai quando ti serviranno'', e sapere di averne
illimitate spingeva a lavorare di più. Per lui era altrettanto dannoso, e per questo
ha introdotto un \textbf{minimo obbligatorio di 24 giorni} di ferie.

C'è stato però un effetto non previsto, e molto desiderato: in tutte le aziende in
cui ha introdotto ferie e permessi illimitati, \textbf{i giorni di malattia sono
calati del 70\%}\footnote{Dato riferito dal relatore sulla base della propria
esperienza, non verificabile.}. La gente si ammala di meno? O piuttosto il giorno in
cui ci si sveglia e si pensa ``oggi non ne ho voglia, mi do malato'' diventa il
giorno in cui si prende un permesso? Si smette di usare come malattia un giorno che
non era di malattia.

\paragraph{L'\emph{Hangover Day}.} Sia a Riot sia a PG Esports il relatore aveva
introdotto quello che a Riot, ``essendo Riot a Dublino'', si chiamava
\textbf{Hangover Day}\footnote{Nella trascrizione ``Angover Day''.}: una giornata
all'anno in cui si aveva il diritto di telefonare alle nove del mattino e dire ``oggi
non vengo perché non ne ho voglia'' (tradotto: perché ieri sera ho bevuto troppo).
Altrimenti quella giornata sarebbe diventata un giorno di malattia o di ferie. Il
fatto è che \textbf{quasi nessuno lo usava}: il relatore stesso a volte ci ha
pensato, ma poi la mattina si chiedeva ``voglio davvero giocarmelo oggi? Tutto
sommato potrei anche andare'', e ci andava. Ma anche solo sapere di averlo lo faceva
stare tranquillo. Allo stesso modo, con le ferie illimitate, chi voleva prendersi un
mese lo faceva (``buon per loro''), chi non lo faceva andava bene lo stesso:
l'importante era che staccassero. C'era un \textbf{benessere più diffuso} nel team,
perché sapevano che, volendo, avrebbero potuto.

\paragraph{Il \emph{Neko Day}.} Il primo punto, i giorni liberi, era chiamato
\textbf{Neko Day}, da \emph{neko}, ``gatto'' in giapponese: il giorno in cui si può
fare il gatto, rannicchiati davanti al camino a non fare niente. Perché
introdurlo? Chi lavora dal lunedì al venerdì passa di solito il sabato o la domenica
a fare le commissioni che non può fare nei giorni lavorativi: la spesa, le bollette,
il bollo dell'auto. In questo modo sfrutta meno il weekend. Con un venerdì libero al
mese (o due, o tre) le persone facevano quel giorno le commissioni del weekend, e
sabato e domenica si riposavano davvero (``anche non fare niente andava bene,
chiaramente''). Il lunedì arrivavano molto più cariche e riposate. Questo ha aiutato
tantissimo a mantenere un team sano, fisicamente e mentalmente.

\paragraph{Il supporto psicologico.} A PG Esports il relatore ha introdotto anche il
\textbf{supporto psicoterapeutico offerto dall'azienda}: sei sedute da un'ora con
uno psicoterapeuta, online. Era il periodo in cui si parlava del bonus psicologo
statale, ``che poi non è mai arrivato''\footnote{Così in lezione. In realtà il bonus
psicologo è stato istituito in Italia nel 2022, ma con fondi limitati rispetto alle
richieste e con forti ritardi nell'erogazione; il relatore si riferisce
probabilmente a questo.}, e lui ha pensato: l'idea è buona, perché non introdurla a
livello aziendale? Erano gli anni dei lockdown: l'azienda stava perdendo molte persone
che andavano in depressione o non avevano più voglia di lavorare, perché avevano
difficoltà a vivere quel periodo (``io stesso non è che l'abbia vissuto benissimo'').
Qualcuno ne ha usufruito e gli ha detto: ``mi stavo per dimettere, poi ho risolto le
cose e ho preferito restare''.

Secondo il relatore oggi c'è una consapevolezza maggiore da parte delle aziende
sulla \textbf{salute mentale}: vuol dire dare più ferie, dare questo tipo di
supporto, dare più \emph{benefit} (anche i buoni pasto). Il benessere non è solo
economico, è anche mentale, ed è una cosa positiva da sapere.

% ---------------------------------------------------------------------
\section{Brainstorming: obiettivo e problemi tipici}
\label{sec:z2-brainstorming}

\subsection{Perché il brainstorming riguarda tutti}

Secondo il relatore questa è la parte che servirà di più, sia nei gruppi e nei
progetti del percorso di studi (il gioco da tavolo, il gioco digitale) sia nelle
aziende. Rispetto al time management ha una pertinenza più immediata: tutti
hanno esempi di \textbf{brainstorming andati male}, da cui si è usciti esausti, con
la sensazione di non aver ottenuto niente, o di non aver parlato abbastanza (``ha
parlato solo quello''), o di non essere sicuri della scelta. Il brainstorming lascia
una sensazione molto chiara, a fine meeting, che qualcosa non torna; il time
management dà piuttosto la tristezza del ``non ce l'ho fatta a fare quello che
dovevo oggi, come mai?''.

\subsection{Generare idee, non decidere}

Il primo chiarimento riguarda \textbf{l'obiettivo}, che il relatore vede spesso
travisato. L'obiettivo del brainstorming è \textbf{generare idee}; invece diventa
spesso decisionale, con lo scopo di scegliere che cosa fare. Ma allora non sarebbe
un brainstorming, sarebbe un \emph{decision making meeting}. Generare molte idee e
scegliere la migliore nello stesso meeting non funziona né per generare idee né per
scegliere la migliore, perché sono \textbf{due metodi che si contrappongono}. Il
brainstorming va fatto quando si devono generare idee \emph{per poi} scegliere che
cosa fare: deve venire \textbf{prima} del \emph{decision making}, come primo passo
della fase di esplorazione di un nuovo progetto.

\subsection{``Non esistono cattive idee''?}

Bisogna evitare che i membri del team \textbf{autofiltrino} le proprie idee per paura
di essere rifiutati o presi alla leggera. Da qui nasce la frase ``non esistono
cattive idee, buttate fuori tutto'', che è vera solo fino a un certo punto. Le
cattive idee esistono, oggettivamente; quello che non si deve fare è
\emph{colpevolizzarle}. Né si deve per forza perseguire ogni opzione: parlare di
opzioni che sono effettivamente brutte è uno spreco di tempo, e il brainstorming ha
anche un problema di gestione del tempo. Il relatore vede spesso brainstorming
gonfiati a ore e ore in cui si parla anche di cose che non meritano nemmeno dieci
minuti, come i feedback buoni ma inutili del giorno prima. \textbf{Trattare tutte le
idee allo stesso modo è sbagliato}, perché non hanno la stessa qualità. Bisogna
trovare un modo di valutarle senza che chi le ha proposte si senta preso di mira,
bullizzato o sminuito.

Per questo a ogni brainstorming dovrebbe seguire un \textbf{follow-up}, il momento in
cui si decide che cosa fare: le due parti vanno separate, altrimenti ``si crea solo
casino''. Il vero lavoro arriva quando le idee vengono messe a terra: c'è la parte di
brainstorming e di \emph{theorycrafting}, e la parte esecutiva viene dopo.

\subsection{I problemi più comuni}

\begin{description}
  \item[Le prime idee monopolizzano la sessione] Le prime due idee ``sparate a
        caso'' sono quelle su cui ci si concentra per la maggior parte del tempo,
        sviscerandone i dettagli (``intendevi così o così?''), per poi chiudere con
        ``altre idee? Questa, questa e questa? Sì, no, ok, andiamo avanti''. Si
        parla troppo delle prime idee e poco delle ultime.
  \item[Poche persone parlano per la maggior parte del tempo] È un effetto
        collegato al precedente.
  \item[\emph{Anchoring}] Ci si lega alle prime idee, e spesso si votano o si
        scelgono non perché siano le migliori, ma perché sembra di conoscerle meglio
        di quelle menzionate in trenta secondi alla fine. Non c'è nessun motivo per
        pensare che un'idea detta alla fine sia peggiore: era solo arrivata dopo.
        C'è un problema di \emph{timing}. Anzi, un'idea venuta istantaneamente è
        forse quella meno ragionata, perché ci si è pensato meno tempo.
  \item[Silenzi imbarazzanti] Spesso le persone arrivano al meeting senza nemmeno
        sapere che è un brainstorming. Quando il relatore convoca un brainstorming,
        comunica prima a tutti il \emph{topic}, il tipo di idee cercate, i paletti e
        le restrizioni, così che anche chi è meno portato all'improvvisazione arrivi
        con qualche idea pronta. Se la gente arriva impreparata rischia di non dire
        niente per un'ora, e allora perché è stata invitata? Non è un difetto loro:
        tanti manager dicevano al relatore ``ci sono persone inadatte ai
        brainstorming perché non sanno improvvisare''. No: si sta facendo
        brainstorming con un metodo che premia solo chi sa improvvisare ed è
        abbastanza estroverso da parlare sul momento. Non significa che chi non
        parla non avrebbe potuto contribuire, ma che non gli è stato dato modo di
        farlo.
  \item[Estroversi e pensatori rapidi dominano] Le personalità estroverse e chi
        pensa velocemente tendono a dominare la conversazione, più o meno
        intenzionalmente, senza lasciare agli altri il tempo di contribuire. Non c'è
        uno spazio in cui un introverso possa dire ``veramente io volevo\dots'':
        ``no, no, stiamo parlando di questa bellissima idea, facciamo questa''.
\end{description}

% ---------------------------------------------------------------------
\section{Un metodo di brainstorming}
\label{sec:z2-metodo}

Il relatore propone una soluzione principale, declinata in vari modi, che tocca
tutti questi punti.

\subsection{Preparazione e ambiente inclusivo}

\paragraph{Arrivare con un obiettivo chiaro.} Bisogna arrivare al brainstorming con
chiaro e preciso che cosa si sta cercando, cioè \textbf{che cosa si sta cercando di
risolvere}, \textbf{perché} e soprattutto \textbf{con quali vincoli}. Si stanno solo
raccogliendo idee per un cliente? Si cercano idee fattibili, o idee innovative? Che
cosa si vuole produrre da questo meeting? Deve essere chiaro a tutti \emph{prima}
del meeting, non spiegato durante, perché altrimenti si ricade nel problema delle
persone non pronte a intervenire.

\paragraph{Non giudicare le idee degli altri.} Al relatore è capitato spessissimo di
sentire un'idea terribile da una persona e, mezz'ora dopo, l'idea perfetta dalla
stessa persona. Ma se all'idea terribile tutti rispondono ``che stupidaggine, per
favore stai zitto'', quella persona non dirà la successiva, perché non vuole
rivivere l'esperienza. Giudicando le idee degli altri, di fatto, li si zittisce.
Le idee migliori nascono quando si propone quella che si pensa sia l'idea perfetta,
magari non è buona e non viene scelta (``fa niente''), si ascoltano anche le altre e
poi la si \textbf{riformula} integrando i pareri sentiti nel frattempo. Sono queste
\textbf{idee combinate} a essere spesso le soluzioni migliori.

\subsection{Separare generazione e discussione: \emph{brainwriting} e
\emph{rapid ideation}}

Con il classico ``non ci sono cattive idee, sparate'', succede spesso che uno prenda
la parola, spieghi la sua idea e il suo concetto, e gli altri comincino subito a
fargli domande (``ma quindi la vedi così?''): si sta unendo la persona che
\emph{genera} l'idea con la persona che la \emph{discute}, nello stesso momento.
Invece bisogna \textbf{dividere la generazione dalla discussione}: un momento in cui
si generano idee e un momento in cui si discutono. È il principio del
\textbf{brainwriting}. Funziona molto quando nel team c'è una figura con una
posizione gerarchica più alta: se quella persona esprime un'idea e i due momenti
sono uniti, tutti tenderanno ad allinearsi. E si vuole che \emph{tutti} possano
avere idee.

Lo strumento è la \textbf{rapid ideation}. Il relatore di solito dice: ``per i primi
cinque minuti non parliamo; scrivete su dei post-it il pitch delle vostre idee''. Se
ne possono scrivere quanti se ne vuole, dieci o uno, e si mettono da parte. Ma il
tempo è di \textbf{cinque minuti}: non dieci, non due, e soprattutto non mezz'ora.
Il \textbf{limite di tempo breve} è importantissimo, perché costringe,
consapevolmente o meno, a filtrare le idee mentalmente e a scrivere solo quelle che
si ritengono migliori: magari non è chiaro perché si preferisca un'idea all'altra,
ma il cervello sta facendo il lavoro in background e fa venire in mente le cose più
importanti o più giuste. È ottimo anche per i team che tendono a distrarsi, perché
sono solo cinque minuti.

\subsection{\emph{Dot voting}}

Poi si passa al \textbf{dot voting}:
\begin{enumerate}
  \item si attaccano \textbf{tutte} le idee su una lavagna fisica, un muro o una
        lavagna digitale (come Miro\footnote{Nella trascrizione ``mirror board'':
        si tratta della lavagna digitale Miro, già citata nella prima
        giornata.});
  \item si \textbf{uniscono le idee che hanno lo stesso concetto di base}. Se si sta
        pensando a un evento dal vivo e la proposta del relatore è un parco di
        dinosauri, mentre quella di un collega è un parco di dinosauri ``del
        paleolitico'' con i tirannosauri, le si raggruppa sotto ``parco di
        dinosauri'': non sono due idee separate. Un altro gruppo potrebbe essere,
        ad esempio, un parco a tema Avengers;
  \item ognuno ha \textbf{da tre a cinque voti} e va alla lavagna a mettere i propri
        pallini sulle idee che gli piacciono e di cui vorrebbe sapere di più.
\end{enumerate}
Il risultato tipico: un'idea con molti voti, una con un paio, altre con un numero
discreto.

\paragraph{Che cosa si vede subito?} Il relatore lo chiede alla classe. Si vede una
\textbf{preferenza} per un tema o per un'idea. Non c'è \textbf{polarizzazione} su una
o due idee iniziali, perché c'è varietà. Il voto è \textbf{anonimo}: non ci sono nomi.
E soprattutto manca un'altra cosa importantissima: \textbf{il giudizio negativo
espresso}. Se c'è un'idea con sette voti, una con cinque e una con tre, è
ragionevole parlare prima di quelle, perché tutti hanno espresso interesse e vogliono
saperne di più; ma nessuno sta dicendo ``non voglio assolutamente parlare di questa
perché mi fa schifo''. È semplicemente oggettivo che quell'idea non ha ricevuto voti.

\paragraph{Autovotarsi.} Il sistema permette anche di votare la propria idea, e non c'è
nessun problema: se la propria idea ha un pallino e gli altri ne mettono sette su
un'altra, il proprio voto è ininfluente. Il relatore vedeva spesso che all'inizio la
gente tendeva ad autovotarsi, poi capiva che era inutile, perché non sarebbe mai stato
il voto decisivo, e tanto valeva votare le idee che interessavano davvero. La
tendenza spariva con il tempo.

\paragraph{Che cosa si è ottenuto.} Tre ore di meeting sono state filtrate in una
discussione concentrata solo sulle idee che tutti ritengono interessanti, con un
\emph{agreement} di fondo: non si parla dell'idea numero uno perché è uscita per prima,
non è una questione di \emph{timing} ma di qualità. Il relatore ha vissuto tanti
meeting con la frustrazione interna del ``stiamo sprecando tempo su questa idea,
quando vorrei dire la mia che secondo me è meglio''. Qui non si può sprecare tempo,
perché si parla di qualcosa che tutti vogliono sentire.

\paragraph{E se si formano sottogruppi che si votano a vicenda?} Uno studente chiede se
non possano esistere ``cordate''. Anche se esistessero, significherebbe che il gruppo di
quel meeting vuole parlare di quella cosa, e va bene: l'importante è non sprecare tempo
su idee che non interessano a nessuno. Se si investono due ore, quarantacinque minuti
su ciascuna di tre idee vanno benissimo; dieci minuti per ogni idea no. E comunque il
voto è anonimo: non si può sapere che un'idea è di Gianni per dare supporto a Gianni,
quindi un eventuale sostegno reciproco sarebbe involontario.

\paragraph{E chi è lento a elaborare?} Uno studente osserva che il problema degli
introversi potrebbe ripresentarsi con i post-it: con un tempo così ristretto, chi è più
lento potrebbe andare ancora più in panico. Per questo bisogna \textbf{arrivare
preparati}: si conosce il topic in anticipo, quindi almeno un paio di idee già
formalizzate si hanno in mente. Magari non se ne scriveranno quindici, ma un paio sì.

\subsection{Quanti voti e quali varianti}

Dare più pallini a testa sarebbe controproducente? Il relatore indicava da tre a
cinque voti. Chiede alla classe quanti siano i gruppi di lavoro: dai quattro agli otto,
alcuni gruppi del secondo anno sono nove. La sua prassi:
\begin{itemize}
  \item \textbf{dai cinque partecipanti in su}, \textbf{tre voti} a testa, più che
        sufficienti per popolare molto la lavagna;
  \item \textbf{fino a cinque partecipanti}, \textbf{cinque voti} a testa, che
        permettono anche di variare un po'.
\end{itemize}
Esistono varianti: ad esempio tre voti di colori diversi, il blu per l'idea più
innovativa, il verde per quella più bella, il giallo per quella più facile, così da
qualificare \emph{perché} un'idea piace. Se l'obiettivo è proporre a un cliente tre
pitch (uno facile, uno medio, uno avanzato), la votazione è già fatta. All'inizio,
però, meglio partire da tre o cinque voti generici; con il tempo si imparerà ad
accelerare.

\subsection{Il moderatore e gli invitati}

\paragraph{Serve un moderatore?} Può essere utile. Il relatore fa spessissimo il
moderatore: partecipa ai brainstorming senza fare altro che moderare, ``perché so che
le persone con le idee fighe sono loro e non io''. Ci si può anche autogestire. Nei team
piccoli non ha senso ``perdere'' una persona per la moderazione; nei team più grandi
(dai cinque in su, sette-otto persone), se qualcuno dice ``non penso di avere idee
particolari su questo, vi modero così vi aiuto'', è un ottimo modo di contribuire.

\paragraph{Chi invitare?} Dipende da che cosa si cerca. Per inventare un nuovo gioco il
relatore invita i game director, gli art director e magari il \emph{publishing manager}
per un input commerciale; per decidere il \emph{combat system} invita magari il senior
game designer e il game director. La domanda che si fa chi convoca il meeting è: dato
l'obiettivo, chi è più ragionevole che possa contribuire?

Un 3D artist della classe chiede se allora gli artisti 3D non vengano inclusi. Dipende:
se l'obiettivo è definire la linea creativa dello stile dei personaggi, il relatore
invita anche il 3D artist, e non solo i \emph{concept artist}. Il 3D artist può dare un
input diverso dal concept: vota un'idea perché per lui è \emph{più fattibile}. Un
concept può essere bellissimo ma ``un inferno'' da modellare, e chi sa di doverlo fare
entro giugno non vota un'idea che sa già infattibile.

In realtà si può invitare chiunque: per un pitch su un'area dedicata al basket alla
Games Week, il relatore ha invitato un tecnico di regia che gioca a basket, per avere
anche la prospettiva di chi ha familiarità con il tema. A uno studente che porta
l'esempio di un elemento di interfaccia sul personaggio, il relatore risponde che,
volendo, si trova un motivo per invitare chiunque: l'importante è avere chiaro
\textbf{chi è necessario}. Il gruppo di lavoro può essere di otto e il brainstorming di
dodici, con input esterni a chi poi farà il lavoro. Basta non esagerare: se la maggior
parte dei voti arriva da chi non farà quel lavoro, si vota su qualcosa che poi non
faranno loro, e non ha più senso.

% ---------------------------------------------------------------------
\section{Dalla votazione alla decisione: la matrice impatto--sforzo}
\label{sec:z2-impatto}

\subsection{Massimo impatto, minimo sforzo}

Fatto questo, non si è finito: ora si sa quali idee interessano di più, ma non si è
ancora deciso che cosa fare. Serve la parte decisionale. L'obiettivo di qualsiasi
azienda (ed è il motivo per cui al relatore non importa se il programmatore fa il test
da casa cercando su Google) è \textbf{massimizzare l'impatto e minimizzare lo sforzo}.
Vale per i progetti, i task, il recruitment, per qualsiasi cosa: ``non perché sono
pigro, ma perché sarebbe stupido fare diversamente''.

Le idee selezionate si mettono quindi in un'altra matrice, con due assi:
\begin{itemize}
  \item sull'asse verticale il \textbf{valore}, cioè quanto \textbf{impatto} avrebbe
        l'idea (in alto l'impatto alto);
  \item sull'asse orizzontale lo \textbf{sforzo} (\emph{effort}), da basso (a sinistra)
        ad alto (a destra).
\end{itemize}
In alto a sinistra ci sono quindi le idee ad alto impatto e basso sforzo, in alto a
destra quelle ad alto impatto e alto sforzo, e così via.

\paragraph{L'esempio Adidas.} Se il cliente è Adidas, con la linea di abbigliamento per
il basket, un'idea a basso valore potrebbe essere una campagna su TikTok in cui gli
utenti condividono i loro momenti con un pallone da basket; un'idea ad alto valore
potrebbe essere fondare un parco divertimenti dove i giocatori NBA vengono a sfidarti
uno contro uno. Questo non dice che una sia meglio dell'altra: dice che l'impatto è
diverso. Allo stesso modo si valuta lo sforzo: la campagna social è a basso sforzo, il
parco divertimenti ad alto sforzo per molti motivi (i permessi di costruzione, i
giocatori NBA).

\subsection{Posizionare le idee, insieme e con precisione}

Il gruppo che a quel punto è \textbf{il gruppo decisionale}, cioè chi farà il lavoro,
deve posizionare sulla matrice le idee ritenute migliori, e deve essere \textbf{d'accordo
tutto il gruppo}: ``siamo tutti d'accordo che il parco dei dinosauri di Stefano e Andrea è
ad alto valore ma anche ad alto sforzo? Sì? Allora lo mettiamo qui''. E bisogna indicare
\textbf{un punto preciso}, non solo il quadrante: piazzare un'idea in un punto non è come
piazzarla poco più in là. Il risultato è una sorta di \textbf{mosaico} delle idee sul
grafico.

\subsection{Leggere la matrice in diagonale}

Davanti a una situazione tipica, con varie idee sparse, il relatore chiede alla classe:
se foste i project manager, se doveste fare voi queste cose, \emph{quali} fareste? Al
plurale, perché spesso da un brainstorming si esce con più idee realizzabili e si deve
scegliere su quali concentrarsi. Le risposte variano: c'è chi sceglie le due con più
impatto, chi aggiunge che dipende da tempi, budget e risorse, chi preferisce quella più
innovativa ma anche più rischiosa.

Il metodo che secondo il relatore va applicato è \textbf{procedere in diagonale}, dalle
idee a basso sforzo e alto impatto verso quelle ad alto sforzo e basso
impatto\footnote{Nella slide le idee erano rappresentate da post-it colorati (rossa,
rosa o salmone, azzurra, verde acqua, viola, blu, gialla), e le indicazioni del
relatore si riferiscono a quei colori. Dalla sola trascrizione non è possibile
ricostruire con certezza la posizione di ciascun post-it, e alcune indicazioni sui
colori risultano contraddittorie; qui è riportato il principio generale, che il
relatore chiarisce nelle risposte alle domande.}:
\begin{enumerate}
  \item un'idea a \textbf{basso sforzo e alto impatto} (nella slide, quella rossa) va
        fatta \textbf{per prima, sempre}: ``non dovrebbe essere nemmeno in discussione'';
  \item poi, a seconda delle risorse e del tempo, si possono fare
        \textbf{contemporaneamente} le idee che stanno sulla stessa diagonale
        successiva, che hanno cioè un rapporto simile tra impatto e sforzo;
  \item poi le successive;
  \item le idee ad \textbf{alto sforzo e basso impatto} (nella slide, la gialla)
        \textbf{non vanno nemmeno considerate}. Nessuno, infatti, l'aveva scelta.
\end{enumerate}
Per le scelte intermedie valgono due criteri, chiariti rispondendo a uno studente che
non capiva perché alcune idee fossero preferibili ad altre: \textbf{a parità di sforzo si
sceglie l'idea con più impatto}; \textbf{a parità di impatto si sceglie quella con meno
sforzo}. Con risorse infinite si farebbero anche le idee più costose; ma di solito le
risorse sono limitate, e allora tra due idee con lo stesso sforzo e un impatto molto
diverso la scelta è evidente. Se ci piacerebbe fare un'idea più ambiziosa ma abbiamo
risorse solo per una meno costosa sulla stessa diagonale, va bene lo stesso: abbiamo
comunque stabilito che è meglio delle idee sulle diagonali successive. Quanto fare dipende
dalla \textbf{bandwidth} a disposizione.

\paragraph{``La gialla era la mia idea: devo offendermi?''} No. Qualcuno ha detto che
l'autore fa schifo? No. Si è concordato che la sua idea, per quanto votata (``qualcuno ha
votato per la mia idea''), è piuttosto difficile da realizzare rispetto al suo impatto.
Non si sta dicendo che non sa più fare il suo lavoro. Anzi, se lui stesso dovesse
scegliere tra la propria idea gialla e la rossa di un altro, sceglierebbe la rossa: non è
nemmeno una questione di discussione. La matrice rende \textbf{visivamente evidente} che
cosa fare, e permette di aggirare molte discussioni che non sono più soggettive:
oggettivamente, dovremmo fare così.

\subsection{Il ruolo del moderatore e dell'esperto}

È qui che lavora il moderatore: ``non è più una discussione; visto il posizionamento,
dobbiamo fare queste cose''. Non si decide a maggioranza dove mettere le idee: se uno
posiziona un'idea e un altro non è d'accordo, difficilmente dirà ``la rossa dovrebbe stare
al posto della gialla''; dirà ``secondo me è un po' più a destra, un po' più in basso'', e
si fanno gli aggiustamenti. Se invece nascono discussioni del tipo ``la rossa dovrebbe
stare al posto della gialla'', vuol dire che non si è definita bene l'idea e non è chiaro
quanto sia davvero difficile da realizzare.

Oppure interviene \textbf{l'esperto}: è l'esempio del concept artist e del 3D artist. Il
concept artist dice che un'idea è a basso sforzo e alto impatto; il 3D artist risponde ``ne
condivido l'impatto, ma dovendola fare io ti assicuro che lo sforzo non è basso, è molto
alto''. Allora è un altro discorso, perché l'esperto sta dicendo che lo sforzo è diverso
da quello stimato. Ma alla fine tutto deve essere concordato.

\subsection{Chi decide}

\paragraph{Il gruppo decisionale può essere diverso da quello del brainstorming?} Sì: a
questo punto escono quelli invitati come ``consulenti'', che hanno dato idee ma non devono
poi realizzarle. Farli votare anche qui significherebbe perdere il collegamento tra chi
deve fare le cose e che cosa si decide di fare. In quel momento non hanno responsabilità:
hanno dato un input, e se decidessero che cosa devono fare gli altri dovrebbero esserne
i manager. Non dovrebbero poter dire agli altri che cosa fare se non ne hanno la
responsabilità. Vale per tutti i team, piccoli e grandi; secondo il relatore chi decide
dovrebbe essere \textbf{al massimo una decina di persone}. In un team di dieci che fa il
brainstorming tutto insieme, va benissimo che chi vota coincida con chi ha partecipato:
l'importante è che \textbf{chi vota sia chi farà il lavoro}.

\paragraph{Lo sforzo è il costo o la fatica?} Uno studente chiede se l'alto impatto
includa anche il costo. No, il costo sta nello sforzo, che per il relatore va
\textbf{ragionato in risorse}: denaro, tempo e così via. Un'idea può avere un basso costo
monetario ma un alto costo in tempo, e allora va bilanciata e si troverà a metà. È lì che
si spostano le idee e che avviene la conversazione: i tecnici dicono che lo sforzo tecnico
è un certo, poi arriva l'\emph{accountant} che dice che a livello economico è un altro, e
la si sposta.

\paragraph{E il piacere di fare il progetto?} In un gruppo di due o tre persone, osserva
uno studente, si finisce per votare anche per preferenza personale: quando si fa un gioco
bisogna avere anche il piacere di farlo. Ma \textbf{anche quello è impatto}: l'impatto
sulla motivazione. Come lo sforzo sono le risorse, l'impatto è l'impatto \emph{su di te,
sul team, sul mercato}, su ciò che si vuole. Il relatore lo ha tenuto volutamente
generico: può essere quanti soldi si fanno o quanto è utile per il proprio CV, e
l'impatto cambia tra un progetto accademico e uno professionale.

% ---------------------------------------------------------------------
\section{Consigli pratici e domande sul brainstorming}
\label{sec:z2-qa-brainstorming}

\subsection{Musica e cambio di scenario}

``Sembrano tutte stupidaggini, ma vi assicuro che aiutano.''

\paragraph{Musica di sottofondo.} È dimostrato che aiuta\footnote{Le affermazioni sugli
effetti della musica (chiave maggiore, rumore ambientale) sono riportate come dette in
lezione e non sono state verificate.}. Meglio una musica con \textbf{tempo e volume
costanti}: niente \emph{bass drop}, niente variazioni di tema o di ritmo. Va bene la
musica fatta di rumore ambientale a basso volume; è stato rilevato anche che le canzoni
in \textbf{chiave maggiore} aiutano la parte creativa a produrre idee più originali.
Meglio la \textbf{musica strumentale}, senza testo: per quanto ci si impegni a non
ascoltarlo, il cervello cercherà di capire le parole. È un effetto marginale e non
succede niente se capita la canzone con il testo, ma si mette in conto che ``l'uno per
cento del cervello'' sarà impegnato a capirlo, come quando in treno si ascolta, che lo si
voglia o no, la telefonata di qualcun altro. È \emph{min-maxing}.

\paragraph{Cambiare scenario.} A Riot c'era una stanza dedicata al brainstorming, pensata
per essere creativa: le finestre sul fiume, le porte a vetri colorati, tutto fatto per
ispirare creatività. Ma fare brainstorming sempre nella stessa stanza generava sempre lo
stesso tipo di idee: anche solo cambiare ambiente aiutava ad averne di diverse. Il
relatore spesso spostava il meeting, ad esempio nello Starbucks sotto l'ufficio. Lo
spazio fisico gioca un ruolo importante, ma anche quello virtuale: se lo si fa sempre su
Discord, si può provare con ognuno uno sfondo diverso e simpatico. Oppure: ha sempre
moderato la stessa persona? Si prova con un altro moderatore. Si cambiano i fattori
perché i risultati non siano sempre gli stessi.

Uno studente chiede se il cambio di scenario possa essere anche qualcosa come passarsi una
palla per avere la parola, o un altro espediente per tenere viva l'attenzione. Con questo
modello, che non richiede di prendere la parola, l'esempio della palla perde senso, ma come
esempio va benissimo: anche il comportamento fisico e il \emph{modus operandi} sono un ottimo
cambio di scenario. Come lo sono cambiare il moderatore o \textbf{cambiare l'obiettivo}: se
il cliente chiede X, gli si darà un'idea X, ma si può provare anche a votare per un'idea Y,
facendo più votazioni con vari criteri di valutazione.

\subsection{Domande}

\paragraph{L'idea scelta deve essere approvata da qualcun altro?} Se un manager dice ``fate
questo brainstorming senza di me'', sta \textbf{delegando la scelta}: ``mi fido di quello che
sceglierete''. Avendo preparato prima limiti, paletti e \emph{constraint}, si aspetta che la
soluzione portata sia in linea con essi; ma se non partecipa deve fidarsi della scelta. Se
poi la scelta debba passare anche da un livello superiore dipende dall'azienda e dagli
accordi precedenti.

\paragraph{Si può fare un brainstorming ibrido, tra presenza e remoto?} Uno studente fa
l'esempio di un gruppo in cui uno è a New York e gli altri a Milano. Sì, senza problemi: si
fa una call, il portatile punta verso il muro con le idee, o le si mostra fisicamente con
il cellulare. Chi è da remoto scrive i suoi post-it, ad esempio mandandoli su WhatsApp al
moderatore o a qualcuno presente che li scrive per lui, e vota dicendo ``mettimi un pallino
qui, qui e qui'' (in quel caso il voto non è anonimo, ma si può fare tranquillamente); se
deve parlare, parla.

\paragraph{E se i colleghi non riconoscono o non condividono il metodo?} Ci deve essere un
\textbf{agreement}: non si può arrivare e dire ``facciamo così'' se gli altri non vogliono.
Questo metodo funziona solo se tutti sono d'accordo a usarlo. Il relatore ne ha proposto
uno tra i cento metodi di brainstorming esistenti, un buon compromesso tra semplicità ed
efficienza (ce ne sono di più semplici e di molto più efficienti ma più complessi). Se il
gruppo non ne vede il valore, al massimo insisterebbe per \textbf{fare almeno una prova},
altrimenti diventa chiusura mentale; se poi per quel team non funziona, se ne cercano
altri.

\paragraph{E nei gruppi molto piccoli, dove diventa un duello?} Uno studente ha lavorato in
gruppi di tre o anche di due persone, dove il brainstorming diventa quasi un duello in cui
ognuno difende la propria idea. Bisogna \textbf{aspettare la discussione}: invece di ``scegliamo
tra la mia idea e la tua'', si cerca di concordare quale idea abbia l'impatto maggiore e
quella lo sforzo minore. E se non si è d'accordo, bisogna \textbf{portare esempi}: se uno
dice ``questa idea costa di più'', l'altro chiede ``in che senso? Che cosa costa di più?''. Se
non sa rispondere, non si può fare qualcosa senza sapere come va a finire: meglio rimandare
la discussione a quando si avranno le informazioni. A volte la domanda giusta è \emph{perché
non abbiamo queste informazioni?}: forse è \textbf{prematuro} decidere, perché non si hanno
ancora tutti i fattori.

Nei team piccoli non dovrebbe esserci un ``si fa come dici tu o come dico io'', una scelta a
pelle: si dovrebbero conoscere così bene le proprie proposte da poterle posizionare con
assoluta certezza sulla matrice, e avere anche la presenza di spirito di dire ``la tua è
chiaramente migliore, si fa la tua''. Si fa la cosa migliore per entrambi, che dovrebbe
essere oggettiva e non soggettiva. Non è facile arrivarci, ma deve essere l'obiettivo del
team.

% ---------------------------------------------------------------------
\section{Disagree and commit}
\label{sec:z2-disagree}

\subsection{Quando la scelta non ci convince}

Uno studente teme che chi vede scartata la propria idea perda motivazione, in un gruppo
piccolo come in uno grande. Non è un deragliamento: è una conseguenza potenziale di
qualsiasi scelta. Il relatore ha visto gente demotivata perfino dalla settimana di quattro
giorni: ci sarà sempre qualcuno più o meno motivato dalle scelte. Ma qual era l'alternativa?
Un brainstorming classico in cui tutti parlano, chi più chi meno, e si sceglie un'idea che
tutti pensano sia la migliore senza esserne davvero sicuri, perché non si è votato, né
espresso una preferenza, né discusso davvero le idee migliori.

Qui entra il concetto di teamwork del \textbf{disagree and commit}: ``non sono d'accordo con
questa decisione, non mi piace, non penso sia la migliore, ma la faccio''. Per due motivi:
\begin{enumerate}
  \item perché \textbf{è il mio lavoro}, e dovrebbe essere il punto principale;
  \item perché \textbf{voi ritenete che sia la scelta migliore}. Se si è la voce fuori dal coro
        e tutti hanno scelto un'altra strada, non si rema contro: è \emph{default to trust}.
        Se gli altri vi hanno visto più potenzialità, e (grande premessa) \textbf{hanno
        ascoltato} il proprio disaccordo e i suoi motivi, e hanno comunque scelto così, ci si
        fida di loro.
\end{enumerate}
Con una postilla: se le cose vanno male, \textbf{non si dà la colpa agli altri} e non si dice
``ve l'avevo detto'', come se non si facesse parte del team. Disagree and commit vuol dire ``mi
fido del team e facciamo questa cosa''; se poi si rivela sbagliata, \emph{il team} ha
sbagliato, e la volta dopo si cambierà il modo di scegliere. Altrimenti non è disagree and
commit, è ``\emph{disagree and ostracize}'': disaccordo più ostilità. Non bisogna cercare il
colpevole, perché si è deciso tutti insieme (sezione~\ref{sec:z1-default-trust}).

Se si pensa davvero che si stia sbagliando clamorosamente il posizionamento di un'idea, si
dice piuttosto ``non andiamo avanti finché non definiamo bene lo sforzo di questa cosa''; ma
bisogna saperlo giustificare, altrimenti è chi dissente a dover fare un passo indietro.

\subsection{La demotivazione fa parte del lavoro}

La demotivazione ci sarà sempre, perché non si sceglierà mai solo ciò che piace: per quanto
piaccia il proprio lavoro, ci saranno task sgradevoli e scelte che non spettano a noi. E
ci saranno scelte (``ve lo dico per esperienza personale'') che si è convinti siano
sbagliate, al punto di pensare ``piuttosto lascio il team'', ma che si finisce per accettare,
per poi rendersi conto che gli altri avevano ragione. Come c'è un \emph{bias} verso l'amico e
il nemico, c'è anche un \textbf{bias verso noi stessi}: crediamo di conoscere bene una cosa, e
se altre dieci persone dicono ``secondo me è fattibile'', bisognerebbe aprirsi alla
possibilità che ne sappiano più di noi.

``Non lo accetto ma lo faccio'' dovrebbe essere alla base del lavoro insieme, perché è
\emph{default to trust}: rifiutarsi (``non mi fido di voi, non lo faccio'') significa non
fidarsi del giudizio del proprio team. E dire ad alta voce ``non sono d'accordo, ma ti darò
comunque il massimo'' è la cosa che \textbf{fa guadagnare di più il rispetto degli altri}.
Non c'è niente di meglio che vedere qualcuno che aveva un'opinione diversa impegnarsi per
far vedere di avere torto, facendo andare bene la cosa. Per il relatore dimostra solo che è
una persona matura e importante; ha avuto persone che, passate da questa situazione, gli
hanno fatto pensare ``è molto più matura di quanto credessi'', e che ha valorizzato molto di
più in seguito, perché sa di potersi fidare di chi fa disagree and commit: è stato aperto
nella comunicazione, ha detto che cosa secondo lui non andava, ma non ha fatto niente per
remare contro, anzi ha dato il suo supporto. ``Non posso chiedere di più a una persona.''

\subsection{Chi arriva a lavori iniziati}

\paragraph{Se il team cambia in corso d'opera, chi subentra come si rapporta alle decisioni
già prese?} (per dimissioni, promozioni e così via). Deve \textbf{accettare lo status quo}. Al
relatore capita spesso di entrare in team con progetti già in corso: si fa dire che cosa è
stato scelto finora e va \textbf{in commit di default}: ``non ho preso parte a questa decisione,
quindi non sono in \emph{agreement}, ma ora faccio parte del team, quindi vado in commit''.

\paragraph{E se l'esperienza dice che una scelta ha uno sforzo molto diverso?} Ad esempio
che un'idea ``rossa'' ha in realtà lo sforzo di una ``viola''. Lo si fa presente, ma dipende dal
\textbf{tempismo}: se la scelta è appena stata fatta e c'è ancora modo di cambiarla, la si
segnala; se si è al 90\% del progetto, è controproducente dire ``fermi tutti, passo
indietro'', e a un certo punto bisogna rispettare la scelta del team (``avete scelto così,
andiamo avanti e vediamo come va''). E se si sa che, calcolando ferie e periodo estivo,
un'animazione prevista per giugno arriverà a settembre? È molto \emph{case by case}: a
seconda dell'apporto si può modificare il risultato o no. È come riaprire un processo se
emergono nuove prove: se dimostrano che l'imputato non è l'assassino, è giusto riaprirlo; se
dimostrano che l'assassino ha sparato con la sinistra invece che con la destra, conta fino a
un certo punto. È giusto segnalarlo, ma non sempre quello che si porta giustifica uno
``\emph{scratch that}, rifacciamo tutto'': diventa \textbf{un feedback, non una risoluzione}
(``state facendo così, questo è il mio feedback, poi vediamo come va'').

Con cinque minuti di anticipo si va in pausa pranzo, con l'accordo per il pomeriggio: circa
un'ora di presentazione più ``standard'' sul leading from within, poi domande e risposte,
``e quando finiscono le domande abbiamo finito''.

% ---------------------------------------------------------------------
\section{Leading from within: perché parlare di leadership}
\label{sec:z2-leadership}

\subsection{La ripresa del pomeriggio}

Alla ripresa, mentre si aspettano gli ultimi arrivati, il relatore chiede se durante la
pausa siano venute in mente domande. L'unica è logistica: le slide verranno date a fine
workshop? Sì: il relatore le passerà all'accademia (a Vanessa, Andrea ed
Enrico\footnote{Nella trascrizione l'accademia è chiamata ``Ive'' e ``Hive'', come nella
prima giornata; per i nomi dei referenti vedi la nota all'inizio del
capitolo~\ref{cap:z1-feedback}.}), che le farà avere alla classe. Gli appunti vanno
comunque presi, meglio su quello che dice che su quello che mostra.

L'ultima parte, dice, dovrebbe essere abbastanza veloce da presentare, ma vale la pena
fermarsi a parlarne, perché è un discorso più avanzato. È la logica chiusura di quanto
detto tra la prima giornata e la mattina: tirare le fila e mostrare come tutto faccia parte
di un unico filone di ragionamento e progettualità, con un unico scopo, un unico obiettivo e
un unico modo di ragionare.

La voce sembra a posto, ma il relatore è stanco dopo due giorni, come immagina lo sia la
classe, nello \emph{stretch} finale di due giornate molto piene. Sono cose che in condizioni
normali spiega al suo team nell'arco di settimane, con il tempo di metterle in pratica,
discuterle e approfondirle. Qui è tutto concentrato: servirà tempo per metterlo in atto e
fare le proprie prove, quindi non bisogna spaventarsi se non verrà tutto giusto la prima
volta.

\subsection{Il leader non è il più bravo}

Perché parlare di leadership in un workshop non rivolto solo ai leader? La risposta
completa arriverà più avanti, ma il punto di partenza è che il relatore \textbf{non crede
nei leader di team o di progetto che detengono da soli il sapere e il modo di fare le
cose}. Come detto il giorno prima, essere un leader, un manager o un capo non dovrebbe
voler dire essere il migliore in quel lavoro, ma essere \textbf{la persona più adatta a
ricoprire quel ruolo}: guidare un gruppo significa \textbf{occuparsi delle persone} al suo
interno.

L'obiettivo che il relatore si dà, e che secondo lui ogni manager dovrebbe avere in mente,
è che quando si torna dal lavoro ci si senta \textbf{ispirati, sicuri e realizzati}, e che
questo sia \textbf{un diritto naturale di tutti}, non un lusso riservato a chi ricopre certe
posizioni, a pochi fortunati o ``a qualcuno sì e a qualcuno no''. Ha notato,
specialmente in Italia, la tendenza a vedere tutto ciò che viene prima del diventare
manager come la \emph{gavetta} necessaria per arrivarci. Invece chiunque faccia parte di
un'organizzazione, che sia un'azienda di videogiochi o altro, dovrebbe sentirsi soddisfatto,
ispirato e sicuro in quello che fa; e non è corretto che un'azienda tratti bene le persone
``da manager in su''. Anche quando non lo dicono, a volte i fatti sono quelli: da manager si
viene trattati in un certo modo e si ha accesso a certe cose, altrimenti no. Non si tratta di
\textbf{appiattire la gerarchia}: si tratta di assicurarsi che, pur dando cose diverse ai
diversi livelli, a tutti i livelli venga riconosciuto uno strato comune di
\textbf{rispetto, fiducia e professionalità}.

\subsection{Le carriere divise: manager o \emph{subject matter expert}}

Uno studente chiede conferma: quindi un \emph{lead artist} non è necessariamente più bravo
di un senior a fare l'artista? Esatto. E qui il relatore apre una parentesi. In molte
aziende, e soprattutto in quelle in cui tende a lavorare lui, la motivazione principale per
diventare manager è avere una promozione e quindi un aumento, più che il desiderio di
prendersi cura del team e delle persone. Come si ovvia al problema? \textbf{Dividendo le
carriere}. Arrivati al livello senior, il relatore fa con ognuno dei suoi senior
\textbf{``the Talk''}, per esporre le due opzioni che hanno davanti:
\begin{description}
  \item[Busta A: diventare manager] È il percorso standard: si valuta se si è adatti a fare
        il leader e a guidare un team, e si danno gli strumenti per farlo, perché
        ``nessuno nasce leader imparato'': prima di provarlo sulla propria pelle nessuno sa
        che cosa voglia dire. Per quanto si possa essere carismatici, fare il leader è
        un'altra cosa, e non vuol dire solo essere carismatici. Ci sono formazione,
        strumenti e un percorso di crescita.
  \item[Busta B: diventare \emph{subject matter expert}] Una persona che va oltre il senior
        nella stessa disciplina: diventa \textbf{l'esperto dell'intera azienda} su quella
        disciplina, perché la sua priorità non è aiutare gli altri ma essere il migliore.
        ``Non c'è assolutamente nulla di male.'' Il suo ruolo è essere il più specializzato
        in quel campo, il punto di riferimento da cui tutti vanno a farsi spiegare le cose.
\end{description}
Il concetto è abbastanza nuovo, e il punto chiave è che l'azienda \textbf{paga il
\emph{subject matter expert} e il leader esattamente allo stesso modo}: non c'è un guadagno
economico nello scegliere di fare il manager invece dello specialista, e il relatore valuta
quel tipo di conoscenza quanto il saper gestire persone. Così chi vuole essere il migliore,
o anche solo un aumento, può continuare a fare ciò che sa fare e avere l'aumento; chi vuole
l'aumento \emph{e} prendersi cura delle persone sceglie l'altra strada. Così si fa una scrematura sul
\emph{motivo} per cui si vuole fare il manager. Sembra una scelta piccola, ma è molto
importante capire se ci si vuole occupare di se stessi o degli altri. Non è un giudizio
(``non sto dicendo che sia meglio l'uno o l'altro''): le due cose \textbf{non convivono}, e
bisogna sapere che cosa si vuole fare, sapendo che l'azienda tratterà allo stesso modo le
due scelte.

\paragraph{Si può essere entrambi?} Si può scegliere ad esempio il percorso da lead artist
e voler comunque continuare a migliorare le proprie capacità. Ma, come nelle simulazioni di
conversazione del giorno prima, il relatore chiederebbe: ``se avessi le risorse per fare solo
una delle due cose, quale sceglieresti?''. Perché non le si avrà: non ci sarà il tempo e il
modo di fare bene entrambe. Il suo ``cinismo'' arriva a dire che bisogna scegliere, perché
tenersi buone entrambe le strade porta a farle male tutte e due. Preferisce un ottimo
\emph{people manager} o un ottimo specialista; le figure intermedie, in questo campo,
rischiano di diventare pericolose per sé e per il team.

\paragraph{Specializzazioni verticali.} La strada dello specialista permette
specializzazioni molto verticali: in chat qualcuno cita il \emph{mecha designer}, ma può
esserci anche lo specialista di animazioni delle espressioni facciali. Naturalmente l'azienda
deve avere bisogno di quella competenza (non ci si specializza in cose che non servono), e non
è detto che ne serva uno solo. È un percorso di carriera parallelo e alternativo a quello del
manager ``standard''.

% ---------------------------------------------------------------------
\section{Il ``momento Piero Angela'': la biologia del gruppo}
\label{sec:z2-biologia}

\subsection{Perché una digressione scientifica}

Il relatore annuncia un ``piccolo excursus biologico'', fondamentale per capire dove vuole
arrivare, e prova a mettere un sottofondo musicale adatto: la condivisione dell'audio non
funziona, quindi manda il link in chat perché ognuno lo faccia partire per conto suo. Un
studente lo battezza subito ``momento Piero Angela''. Il relatore cerca di far combaciare le
motivazioni aziendali del suo metodo di leadership e di progettualità con \textbf{basi
scientifiche}: in passato ha collaborato con varie università che si occupavano di
neuromarketing e neurobiologia, per avere una conferma o una smentita del fatto che ciò che
ipotizzava ``non fosse una cavolata colossale''. Semplificando moltissimo, si guardano solo
alcune sostanze che si traducono direttamente nel modo in cui si vive l'essere in un gruppo,
in un team, in un'azienda\footnote{L'impianto di questa parte (endorfine, dopamina,
serotonina, ossitocina e cortisolo, il villaggio protetto dai pericoli esterni, il ``cerchio
di sicurezza'') ricalca da vicino quello del libro di Simon Sinek \emph{Leaders Eat Last}
(2014), che il relatore cita poco più avanti. Le descrizioni biologiche sono molto
semplificate, come avverte lo stesso relatore, e vanno lette come tali.}.

\subsection{Endorfine}

Le \textbf{endorfine} sono rilasciate per \textbf{mascherare il dolore}. Lo sa chi corre, chi
fa pesi o ginnastica, ma anche chi ha riso così tanto per una battuta da avere male ai
fianchi: è quella sensazione di dolore in cui però si è contenti di provarlo. Il \emph{runner's
high} è l'esempio più calzante: essere distrutti dopo una corsa ma soddisfatti. È la risposta
biologica allo sforzarsi in cose che sono per il nostro bene, un modo con cui il corpo ci
``convince'' a farle.

\subsection{Dopamina}

La \textbf{dopamina} è più familiare: è la \textbf{ricompensa per un'azione}. Dalle cose più
banali (si mangia e si riceve dopamina) a quelle viste in mattinata: anche completare una
checklist rilascia dopamina (sezione~\ref{sec:z2-pianificare}). Viene rilasciata quando si
vede l'obiettivo, come il traguardo durante una corsa; è la stessa cosa che accadeva ai
cacciatori del villaggio in tempi antichi, che rilasciavano dopamina quando vedevano le tracce
e venivano così incentivati a seguirle. Sono concetti fondamentalmente ``animali'', ma si
applicano tuttora anche in azienda.

\paragraph{Una nota di colore: la dipendenza.} La dopamina \textbf{crea dipendenza}: molti casi
di ossessione per alcol, nicotina, gioco d'azzardo e social media sono problemi legati alla
dipendenza da dopamina. Secondo il relatore non è un caso che le diagnosi di ADD e ADHD siano
aumentate del 60\% negli ultimi dieci anni\footnote{Nella trascrizione ``EDD e HD''. Il dato
e il nesso causale con la maggiore disponibilità di alcol, nicotina, gioco d'azzardo e social
media sono riportati come detti in lezione e non sono stati verificati; l'aumento delle
diagnosi viene in genere attribuito anche ad altri fattori, come la maggiore consapevolezza e
il cambiamento dei criteri diagnostici.}, perché buona parte di quegli stimoli è diventata
ampiamente disponibile proprio in quel periodo. Anche la dipendenza dal lavoro, a volte, ha la
stessa origine: si vuole il rispetto che deriva dal lavoro, il \emph{rush} che deriva
dall'averlo.

\subsection{Il villaggio di diecimila anni fa}

Si fa poi un salto di diecimila anni, in un villaggio. Gli esseri umani, per quanto più o meno
introversi o estroversi, sono \textbf{animali sociali} che tendono a vivere esperienze di
gruppo; le aziende nascono da un gruppo di persone che si danno un obiettivo condiviso, come
in antichità un gruppo di persone che voleva sopravvivere insieme. Quello che non è cambiato è
che \textbf{fuori dal gruppo ci sono dei pericoli}: tigri, lupi, draghi (``do per scontato
che nell'antichità ci fossero i draghi, perché D\&D mi insegna così''). Per evitarli, i gruppi
tendono a riunirsi e a creare difese intorno a sé: da qui le prime difese dei primi villaggi,
per isolarsi dai pericoli esterni. ``Vi giuro che avrà tutto senso.''

\subsection{Serotonina}

La \textbf{serotonina} viene rilasciata anche \emph{indirettamente}: non deve per forza
riguardare noi. Quando si assiste a una cerimonia o celebrazione pubblica che include una
persona cara, si rilascia serotonina come se si fosse noi l'oggetto del riconoscimento. È uno
dei motivi per cui si fanno diplomi e lauree ``all'americana'', con la consegna sul palco
davanti a tutti: gli studi avrebbero mostrato che anche i genitori in platea rilasciano
serotonina, come chi riceve il premio. Lo stesso vale per parenti, amici, matrimoni: sono
esperienze condivise, e si tende a condividere in pubblico riconoscimenti e successi perché in
tutti i partecipanti viene rilasciata serotonina.

Purtroppo ha a che fare anche con i beni materiali: la tendenza all'\emph{hoarding}, cioè a
collezionare eccessivamente qualcosa, nasce dal fatto che quell'attività rilascia serotonina,
ingannando la mente nel credere di aver ottenuto un successo.

È chiamata anche \textbf{la sostanza della leadership}, perché aiuta a riconoscere in qualcun
altro chi ha valore per il gruppo. Nel villaggio, quando arriva un pericolo, c'è sempre un
gruppo di persone che dice ``ce ne occupiamo noi'': il gruppo sa che si stanno prendendo un
rischio, e la serotonina è il modo per riconoscere che meritano rispetto. Il relatore consiglia
il libro di Simon Sinek \emph{Leaders Eat Last}\footnote{Nella trascrizione ``Leadership
Lust''. Il titolo è \emph{Leaders Eat Last: Why Some Teams Pull Together and Others Don't}
(2014).}: i leader mangiano per ultimi,
all'opposto di quanto accadeva in antichità, quando al ritorno dalla caccia mangiavano per
primi i cacciatori, che dovevano essere in forze per la caccia successiva, poi donne e
bambini, poi gli altri. C'è un \textbf{riconoscimento sociale}: tu fai qualcosa per proteggere
il gruppo, noi ti riconosciamo qualcos'altro in cambio.

Le aziende sono strutturate in modo molto diverso? È una domanda retorica: quando si dice che
serve un capo, un manager, un rappresentante, si sta dicendo che quella persona dovrebbe avere
l'interesse di proteggere il gruppo e di prendersi cura di tutti. Per questo sapere che il
proprio capo guadagna di più non fa particolarmente arrabbiare: lo si vede come quasi naturale,
perché più responsabilità si ha verso il gruppo più si guadagna, dato che quello dovrebbe essere
il suo scopo. Ed è anche il motivo per cui il relatore si arrabbia molto quando legge di CEO o
\emph{executive} che, a fronte di licenziamenti di migliaia di persone, si assegnano dei bonus.
È una rabbia non del tutto razionale, che sente da dentro: è come se si stesse violando un
\textbf{contratto sociale}. Il contratto è ``tu ti prendi dei rischi per me e mi difendi, io ti
riconosco più cibo, più soldi''; quando si abusa di quella posizione per un beneficio personale
a scapito di chi si dovrebbe proteggere, tutti lo vedono come un'ingiustizia, perché si viene
meno al contratto stabilito diecimila anni fa con la nostra biologia.

\subsection{Ossitocina}

L'\textbf{ossitocina}\footnote{Nella trascrizione ``ossidocina''.} è la \textbf{sensazione di
sicurezza} che si genera quando si è vicini a qualcuno. Non è un caso che tra i legami di
affinità e collaborazione più forti ci siano quelli dei team che si uniscono contro un
\textbf{nemico comune}: i team che superano insieme avversità e difficoltà ne escono molto
legati. Tra le persone che il relatore sente più spesso in questo periodo ci sono quelle con cui
ha condiviso momenti difficili, un cattivo manager, una brutta esperienza lavorativa: sa che
erano lì con lui. Vale anche per i veterani di guerra o per i sopravvissuti a una stessa
tragedia, che sviluppano un rapporto molto forte (se ne è parlato molto negli Stati Uniti a
proposito dei veterani), perché hanno rilasciato immense dosi di ossitocina. È anche ciò che fa
accettare subito il figlio alla madre quando lo prende in braccio dopo il parto, ed è il motivo
per cui, a migliaia di anni di distanza, si stringono ancora accordi con una stretta di mano o
un abbraccio: c'è sempre qualcosa di fisico, generato dal bisogno di ossitocina.

Viene generata anche quando si compiono o si \textbf{assiste ad atti di generosità} senza
aspettarsi nulla in cambio. Se il relatore dicesse di aver donato cinquanta euro a
un'associazione di beneficenza, la reazione sarebbe ``ok, bene''; se dicesse che ogni sabato
pomeriggio fa volontariato, lo si vedrebbe come qualcosa di molto più prezioso, e la cifra
potrebbe crescere parecchio prima che qualcuno la consideri più preziosa dei suoi pomeriggi:
tutti sanno che la vera donazione è \textbf{il tempo} investito, che vale più del bene
materiale. E basta assistere: vedere qualcuno che aiuta un'anziana ad attraversare la strada fa
stare meglio anche noi. È percepita a livello di gruppo. Infine, l'ossitocina rafforza il
\textbf{sistema immunitario}: sentirsi sicuri e tranquilli fa stare meglio e, in un certo senso,
ammalare di meno.

\subsection{Cortisolo}

L'opposto è il \textbf{cortisolo}: la sensazione di stress, paranoia e paura, la risposta
\emph{fight or flight} (o scappo o combatto), con battito accelerato, muscoli in tensione e così
via. Tra i suoi effetti:
\begin{itemize}
  \item \textbf{blocca il sistema immunitario}: se si sta per morire, al corpo non importa di
        prendersi un raffreddore. È il motivo per cui le persone stressate si ammalano di più;
  \item \textbf{blocca l'ossitocina}: non ci si può sentire tranquilli quando si è in ansia;
  \item soprattutto, e questo è il punto più importante per le aziende, è
        \textbf{contagioso}. In qualsiasi documentario, quando una gazzella che bruca alza la
        testa e percepisce un movimento nella savana, scappano tutte, non solo lei. Succede,
        non sorprendentemente, anche in ufficio: quando qualcuno ha paura perché potrebbero
        esserci licenziamenti, o perché il manager è molto aggressivo, la paura si diffonde nel
        gruppo ed è molto dannosa.
\end{itemize}
È anche l'effetto alla base della \textbf{sindrome dell'impostore}: sentirsi costantemente
inadeguati perché si teme che succeda qualcosa significa vivere in uno stato di stress. Ed è il
motivo per cui chi è insoddisfatto del lavoro tende a stare peggio anche fisicamente e ad
ammalarsi di più.

\begin{center}
  \begin{tabular}{@{}p{0.17\textwidth}p{0.36\textwidth}p{0.37\textwidth}@{}}
    \toprule
    \textbf{Sostanza} & \textbf{Che cosa fa (in sintesi)} & \textbf{Esempi dalla lezione} \\
    \midrule
    Endorfine & mascherano il dolore di uno sforzo che ci fa bene &
      \emph{runner's high}, ridere fino ad avere male ai fianchi \\
    \addlinespace
    Dopamina & ricompensa un'azione, si attiva vedendo l'obiettivo; crea dipendenza &
      spuntare una checklist, le tracce per il cacciatore, gioco d'azzardo, social \\
    \addlinespace
    Serotonina & riconoscimento sociale, anche indiretto; ``sostanza della leadership'' &
      lauree e premi in pubblico, chi difende il villaggio, \emph{hoarding} \\
    \addlinespace
    Ossitocina & sicurezza e legame; rafforza il sistema immunitario &
      il nemico comune, i veterani, la stretta di mano, il volontariato \\
    \addlinespace
    Cortisolo & stress, \emph{fight or flight}; blocca difese immunitarie e ossitocina; è
      contagioso & la gazzella che scappa, la paura dei licenziamenti, la sindrome
      dell'impostore \\
    \bottomrule
  \end{tabular}
\end{center}

% ---------------------------------------------------------------------
\section{Domande sulla biologia del gruppo}
\label{sec:z2-qa-biologia}

\paragraph{Il gioco d'azzardo include i gacha e le loot box?} Assolutamente sì, anzi ne sono il
\textbf{driver principale}: i gacha, i FIFA Ultimate Team, lo ``sbustare'' da \emph{Magic} a
\emph{Hearthstone}, e in generale il ``vediamo se stavolta mi va bene'' sono tutti gioco
d'azzardo. Lo studente aggiunge che molti giochi mobile sono costruiti su questi meccanismi e che
c'è chi gioca più a quelli che ai giochi ``normali''. È il motivo per cui alcuni stati hanno
dichiarato illegali certi giochi\footnote{Il caso più noto è quello del Belgio, che dal 2018
considera gioco d'azzardo le \emph{loot box} a pagamento; anche i Paesi Bassi le hanno
contestate.}.

\paragraph{Come si genera l'ossitocina nei rapporti da remoto?} ``Bellissima domanda.'' È il motivo
per cui, come detto il giorno prima, il relatore insiste perché i meeting si facciano almeno in
\textbf{webcam}, perché i riconoscimenti siano dati e i successi celebrati \textbf{pubblicamente}
(quando non mette troppo a disagio la persona), e in generale cerca di creare \textbf{situazioni
condivise}. Anche lavorando da remoto, insiste per avere il \emph{team building} dal vivo e almeno
\textbf{un incontro all'anno} di persona in ufficio, perché crea situazioni di rapporto diretto.

\paragraph{Come si contrasta il cortisolo?} In chat qualcuno suggerisce la meditazione: sì,
assolutamente. Per il lavoro, la risposta arriva subito dopo nella presentazione.

\paragraph{Il nemico comune è una strategia che un manager può usare?} Spesso \textbf{il manager è
il nemico comune}, senza bisogno di sforzarsi. Ma il relatore chiede di capire bene il punto:
alcuni manager applicano consapevolmente il concetto ``mi ergo io a nemico del team, così sono
tutti uniti contro di me e lavorano bene'', e ne fanno uno stile di management. Decidono di essere
il \emph{bad cop}, il \emph{villain}, ed è così che ottengono risultati. È una scelta lecita e
funziona; il relatore preferisce fare diversamente, ma non per questo dice che non funzioni. Molte
volte si sottovaluta quanto sia consapevole: ``non si rende conto di quanto sia antipatico a
tutti''; forse se ne rende perfettamente conto e lo fa apposta per essere il bersaglio. Il relatore
l'ha visto fare spesso come soluzione disperata, quando il team era troppo impegnato a litigare al
suo interno: ``piuttosto che vederli litigare tra loro, gli do io un motivo per odiare me, così si
uniscono contro di me''. Non lo consiglia assolutamente, ma l'ha visto funzionare.

Uno studente racconta che nella sua azienda il manager fa vedere \emph{il lavoro stesso} come il
nemico (``prima finiamo, meglio è''). È un modello più basilare ma, come dice lo studente stesso,
funziona.

\paragraph{Quanto è etico puntare sulla dipendenza del giocatore per fare soldi?} ``Non è
assolutamente etico, ed è assolutamente profittevole, altrimenti le aziende non lo farebbero.'' Il
relatore racconta una cosa che lo ha sconvolto quando l'ha scoperta, a Blizzard. Lavorava su
\emph{Hearthstone}\footnote{Nella trascrizione ``Arsdol''.} e si stava discutendo di come
monetizzare il gioco; venne preso come consulente un esperto, israeliano se ricorda bene, che
lavorava a Malta per società di poker online e di \emph{gambling}. Raccontò che molte aziende che
monetizzano sulla dipendenza dal gioco d'azzardo avevano sviluppato metodi più o meno etici (``più
spesso meno''). L'esempio culminante: analizzavano gli utenti occasionali, quelli che giocavano
magari una settimana al mese, per capire \emph{quando arrivava loro lo stipendio} (di solito si
deposita lo stipendio e poi si gioca quello che resta), per colpirli il giorno dopo con email di
offerte per tornare a giocare. Li si stimolava proprio nel momento in cui avevano disponibilità, sapendo
che erano sensibili, per aumentare le probabilità di riaverli come giocatori. ``È orribile. È
terrificante e, come tutte le cose terrificanti, è anche altamente funzionale.''

Gli ha fatto riflettere molto su dove ci si possa spingere per fare soldi. La versione etica, per
lui, è ad esempio far uscire un gioco sotto Natale, quando si è più disposti a regalarlo ai figli:
si stanno ottimizzando le proprie possibilità, e lo accetta. Si scade nel \textbf{predatorio}
quando si sfrutta la debolezza di qualcuno in un momento preciso per estorcergli soldi che non può
permettersi di perdere. Poi ognuno tragga le proprie conclusioni: ``non è un workshop sull'etica del
lavoro''. Non pensa di dover imporre la propria opinione, né si arroga il diritto di dire che chi lo
fa è una bestia senza cuore, visto che è il loro lavoro: ``chiunque riesca poi a dormire in pace con
se stesso, buon percorso''.

% ---------------------------------------------------------------------
\section{La piramide della fiducia}
\label{sec:z2-piramide}

\subsection{La leadership non è un titolo}

Chiuso il ``momento Piero Angela'', il relatore arriva al punto. La frase della prima giornata (``la
capacità di un gruppo di persone di fare cose straordinarie dipende dal modo in cui queste persone
operano insieme'') porta a capire che \textbf{la leadership non è un titolo lavorativo}. Essere
manager da contratto, o avere ``manager'' nel titolo, rende un manager, non un leader. Il relatore
non ha mai visto un team rispettare qualcuno solo perché da contratto doveva farlo: si può avere
paura o timore del proprio manager, ma non lo si deve rispettare perché è il manager. \textbf{Il
rispetto si guadagna}, e si possono influenzare le persone intorno solo attraverso il rispetto
guadagnato, quando si sono accumulati abbastanza esempi e \emph{case history} da giustificare
rispetto e fiducia. Altrimenti, dice, sarebbe il primo a chiedere ``perché vi fidate di me, se non
vi ho ancora dato motivo di farlo?'': non sarebbe razionale.

Da qui la \textbf{piramide della fiducia}, che si costruisce partendo dalla base.

\subsection{Alla base: l'integrità morale}

Per essere un buon leader o un buon manager bisogna \textbf{voler genuinamente aiutare gli altri
membri del team}. È il motivo del discorso sulle carriere divise: è a questo livello che si
screma chi vuole la promozione per i soldi. Chi vuole essere leader non deve cercare il guadagno
personale, ma quello del gruppo. Come detto il giorno prima, il relatore è pronto a litigare per il
suo team e a prendersi le colpe per il suo team, perché vuole il bene del suo team e sa che poi loro
si prenderanno cura di lui. Chi ha in mente questo sta in questa categoria. Poi:
\begin{itemize}
  \item \textbf{essere onesti}. Il relatore può affermare di non aver mai mentito
        intenzionalmente a un collega, a un cliente o a un giocatore. Gli è capitato di dare
        informazioni errate perché in quel momento sapeva quelle: ad esempio annunciare che
        un'espansione sarebbe uscita in un certo mese, e poi l'uscita è slittata\footnote{Nella
        trascrizione l'esempio non è coerente: prima si parla di un'uscita a novembre rimandata
        di una settimana, poi di ``pensavo fosse settembre''. Il senso resta chiaro.}. Ha detto
        una cosa sbagliata, ma tutti hanno capito che non stava mentendo: in quel momento pensava
        che fosse vera;
  \item soprattutto, \textbf{accettare la responsabilità dei fallimenti}: quando si sbaglia,
        saper dire ``ho sbagliato''. Il relatore lo ha messo perché ``non avete idea di quanti
        facciano fatica'' ad ammettere di poter sbagliare: c'è sempre una scusante, una
        mitigazione, un ``sì, ma''. Lui vuole sentire se è colpa tua o no;
  \item \textbf{chiedere scusa} e \textbf{ammettere di non sapere}: anche dire ``non lo so'' viene
        spesso evitato, come accettare un fallimento.
\end{itemize}

\paragraph{``Raccontami il tuo più grande fallimento''.} Per lo stesso motivo, nei colloqui il
relatore chiede spesso, più che ``parlami di un tuo grande successo'' o di un progetto andato bene,
\textbf{``mi racconti il tuo più grande fallimento?''}, e osserva come lo raccontano. Chi non si
aspetta la domanda spesso racconta di quella volta in cui ha provato in tutti i modi a fare in un
certo modo, ma il capo gli ha detto di fare diversamente e il progetto è andato male. ``Non mi hai
risposto: mi hai raccontato un progetto andato male, ma il tuo messaggio è che volevi fare
diversamente e qualcuno ti ha detto di fare così. Quindi non è un tuo fallimento.'' È come
rispondere a ``qual è la tua debolezza?'' con ``sono così attento ai dettagli che a volte\dots'':
si fa passare un punto di forza per una debolezza. Il relatore vuole un esempio in cui \emph{tu},
senza altri fattori, hai sbagliato, e come ti sei comportato. Gli dice molto su come quella persona
reagirà ai suoi feedback: se ha davanti ``un muro'' che dirà sempre ``non è colpa mia'', non è adatta
a un ambiente in cui il feedback è alla base di ogni comunicazione. Chi non ammette di aver mai
sbagliato in vita sua, lo scarta.

\paragraph{È possibile colpevolizzarsi troppo per un fallimento?} Sì, assolutamente. Rientra molto
nella sindrome dell'impostore, o nell'aver sempre vissuto rapporti di lavoro tossici: ci si abitua ad
accettare che ci vengano urlate contro le cose e che sia sempre colpa propria, e ci si autoconvince
che lo sia.

\subsection{Al centro: l'affidabilità}

Superato questo livello si arriva all'\textbf{affidabilità}. Per essere un buon leader non bisogna
essere il migliore, ma affidabili: basta essere riconosciuti come bravi in qualcosa, non
necessariamente i migliori, e quell'affidabilità ce la si porterà anche nella leadership. Significa:
\begin{itemize}
  \item \textbf{rispettare la parola data}: se si promette ``parleremo di un aumento al prossimo
        slot disponibile'', se ne parlerà. Non ci si dimentica, e non si dice poi ``in realtà non si
        poteva fare, ma ho voluto provarci fino all'ultimo''. Una promessa va mantenuta;
  \item \textbf{impegnarsi per il tempo necessario} a raggiungere gli obiettivi;
  \item \textbf{saper rifiutare un lavoro} quando necessario: è il saper dire di no della mattina,
        ancora la matrice di Eisenhower (sezione~\ref{sec:z2-eisenhower}). ``Sono sempre gli
        stessi concetti.''
\end{itemize}

\subsection{In cima: la leadership, cioè dare l'esempio}

L'ultimo gradino è la \textbf{leadership}\footnote{Nella trascrizione ``dealership''.}, cioè
\textbf{dare l'esempio}. Non bisogna essere i più bravi; bisogna saper \textbf{offrirsi volontari
non solo quando il progetto è un successo sicuro, ma anche quando c'è il rischio di un
fallimento}. Sono tutti bravi a offrirsi per il gioco nuovo che andrà sicuramente bene.

\paragraph{Il rappresentante dei licenziati.} Al relatore è capitato di offrirsi per rappresentare
il team durante i \emph{layoff}, con il rischio di essere ostracizzato dall'azienda; con quelle
persone è ancora in contatto. Quando Blizzard fece i licenziamenti nella sede francese, vicino a
Versailles\footnote{Nella trascrizione ``i layoff di Versailles''. La sede europea di Blizzard si
trovava a Vélizy-Villacoublay, nei pressi di Versailles; il relatore non precisa l'anno.}, che
coinvolsero quasi tutto l'ufficio, lui non era tra i licenziati, ma venne \textbf{eletto
rappresentante sindacale} di chi lo era. Per lui è stato molto importante: pur non essendo parte in
causa, lo ritennero degno della loro fiducia. Per quanto ne sapevano, avrebbe potuto negoziare e
accettare tutto quello che proponeva Blizzard, perché non ne era toccato: che si negoziasse un mese
di liquidazione in più o in meno non gli portava un soldo. Avrebbe persino potuto usare la
situazione per negoziare un aumento per sé (``gli facciamo accettare tre mesi invece di sette e a te
diamo un aumento''): non era fuori discussione. Se si sono fidati è solo perché c'era una piramide di
questo tipo e lui aveva dato loro abbastanza motivi per farlo.

\paragraph{Autorevolezza, non autorità.} A uno studente che lo riassume così, il relatore risponde
che è proprio questo: un carisma diverso dall'essere simpatici, l'essere un punto di riferimento,
\textbf{più un pilastro di supporto che il motore} che fa andare avanti le cose. Non il motore davanti
a tutti, ma il \textbf{meccanico} che, quando si rompe qualcosa, sai che non ti frega: ti sistema le
cose senza ``farti la cresta'' sulle pastiglie dei freni.

\paragraph{Non abbandonare nessuno.} E soprattutto \textbf{non abbandonare chi è in difficoltà}.
Sembra una banalità, ma è alla base del motto dei Marines statunitensi, \emph{no one left behind}:
sapere che, se si viene feriti, verranno a prenderti anche sotto i missili è un forte elemento
motivazionale per correre quel rischio. In azienda è il motivo per cui il relatore ribadisce sempre
che, se qualcuno è in difficoltà, la sua prima opzione non è licenziarlo ma \textbf{dargli una mano},
e lo ha fatto a volte anche all'eccesso, per far vedere agli altri che quando saranno in difficoltà
non li lascerà in mezzo a una strada. A volte ce l'ha fatta: ha esempi di persone ``a tanto così''
dal licenziamento recuperate per i capelli e diventate ottimi \emph{performer}. Altre volte la
persona non era recuperabile e alla fine è stata licenziata. Ma chi attraversa un periodo difficile
(il \emph{burnout}, la depressione, una relazione finita) sa che la prima risposta non sarà ``ci sei
servito finora, ora puoi andartene'', ma ``che cosa possiamo fare?''. E questo fa lavorare tutti
tranquilli.

% ---------------------------------------------------------------------
\section{Prendersi cura del gruppo}
\label{sec:z2-gruppo}

\subsection{La soluzione: più leader possibile}

In un caso del genere manager e leader coincidono, ed è l'ideale: un manager dovrebbe essere anche
un leader. Ma la vera soluzione è \textbf{avere più leader possibile}. Se l'obiettivo del leader è
prendersi cura degli altri, e il manager dovrebbe essere ``il primo tra i leader'', ancora meglio è
che \textbf{tutti siano leader del proprio gruppo}, perché la leadership non deriva da un titolo né
dall'essere il manager. Il relatore ha visto tanti gruppi andare avanti perché qualcuno al loro
interno era il leader, non perché lo era il manager.

Vuol dire mettere in pratica tutto il discorso della prima giornata su reciprocità, gentilezza,
aiuto e collaborazione. Si sa che, nel peggiore dei casi, c'è il manager, ma in condizioni normali ci
sono i colleghi: e allora si va da chi è in difficoltà a dire ``posso darti una mano?''. Così, quando
saremo noi in difficoltà, magari senza nemmeno rendercene conto, sarà qualcun altro a dirci ``ti vedo
in difficoltà, come posso aiutarti?''.

\paragraph{L'esempio è contagioso, nel bene e nel male.} Quando si ha una posizione anche solo di
rispetto agli occhi degli altri, tutto ciò che si fa ha effetto sugli altri. In negativo vale la
\textbf{slippery slope}\footnote{L'esempio del cartone abbandonato ricorda la cosiddetta ``teoria
delle finestre rotte'' (\emph{broken windows theory}), secondo cui i piccoli segni di degrado
tollerati favoriscono comportamenti peggiori.}: se si accetta un cartone di pizza abbandonato per
strada, domani se ne accetteranno due, dopodomani tre. Ma vale anche in positivo: chi vede qualcuno
fare qualcosa di positivo pensa ``lo voglio fare anch'io''. ``Come fai una cosa, fai tutto'', dice un
proverbio che il relatore attribuisce alla tradizione cinese\footnote{La massima (``how you do
anything is how you do everything'') è attribuita a fonti diverse, spesso alla tradizione zen;
l'attribuzione è incerta.}.

\subsection{Il villaggio di oggi}

Si torna al villaggio del ``momento Piero Angela'', con i difensori e le bestie (reali o no) intorno.
Oggi fuori dal gruppo non ci sono più tigri, ma le \textbf{deadline}, l'\textbf{overwork}, i
\textbf{layoff}, il manager che urla, le aspettative, i \textbf{KPI}. Avere manager che si prendono
cura del gruppo è un buon inizio, ma è ancora meglio avere \textbf{il gruppo che si prende cura del
gruppo}. Non bisogna più limitarsi a pensare che sia responsabilità del manager prendersi cura di
noi: è responsabilità nostra prenderci cura del gruppo di cui facciamo parte.

Che cosa serve? La piramide vista prima. Serve essere percepiti come \textbf{autentici}: il
relatore ormai riconosce abbastanza facilmente chi si vuole ergere a leader perché gli piace
\emph{sentirsi} leader, non perché vuole \emph{fare} il leader, e quella mancanza di genuinità non
aiuta. In generale bisogna \textbf{generare sicurezza} nel resto del gruppo.

\subsection{Smettere di dire ``ci penserà qualcun altro''}

Un esempio concreto è smettere di dire \textbf{``ci penserà qualcun altro''}. Nello sviluppo è
evidente: se si trova un bug che non riguarda il proprio lavoro, tutti sono d'accordo che vada
comunque segnalato, perché altrimenti potrebbe arrivare in \emph{live}, nella \emph{release}.

\paragraph{Le lasagne di PG Esports.} Circa un anno e mezzo prima, quando lavorava ancora a PG
Esports (che organizza anche eventi dal vivo), l'azienda ospitava in ufficio dei potenziali sponsor
venuti a conoscere il team per decidere se sponsorizzare la Milano Games Week. C'era un catering, con
i classici buffet. In mezzo alla confusione (una ventina di persone di PG Esports e una trentina tra
manager, sponsor ed \emph{executive}), il relatore vede per terra, in mezzo alla stanza, una mezza
porzione di lasagne abbandonata. Senza pensarci la raccoglie per portarla via: non è sua, ma è un
piatto in mezzo alla stanza. Nel tragitto fino al bancone del buffet, almeno altri sei tra executive e
sponsor che non lo conoscevano gli consegnano il loro piatto non finito, come a un cameriere: nella
loro testa l'unico motivo per cui qualcuno potesse raccogliere un piatto da terra era che fosse il
suo lavoro. Sul momento non ci fa caso; arrivato al bancone si gira e vede tutto il suo team che ride
della scena.

Torna da loro e ne nasce una discussione: ``avete visto quanto è stato educativo?''. Avevano visto
esattamente come ragionano i gruppi che non sono abituati a pensare ``adesso ci penso io'', ma ``ci
penserà qualcun altro''. In quelle persone non c'era malizia: sicuramente era un cameriere, perché
altrimenti farlo? ``Non è il mio lavoro, lo avrà fatto perché è il suo.'' E diventa un problema quando
la stessa tendenza (``non c'entra con il mio lavoro, ci penserà qualcun altro'') si applica a
qualsiasi altra cosa. Il relatore si aspetta che \textbf{tutti raccolgano quel piattino}: che chi vede
un bug, se non può correggerlo lo segnali, e se può correggerlo lo corregga anche se non è il suo
lavoro, ``perché è un piattino in mezzo alla stanza''.

\paragraph{Da reattivi a proattivi.} Significa passare da un atteggiamento \textbf{reattivo}, cioè di
fatto passivo, a uno \textbf{proattivo}: prevedere i problemi e cercarli intorno a sé per risolverli,
invece di aspettare che esplodano (``finché non esplode non è un problema mio''). Ragionare nell'ottica
del ``e se succedesse, e io me ne fossi accorto senza fare niente?''. Uno studente lo chiama
\emph{mindset da tester}: il relatore si aspetta che, se si vede qualcosa, ci si assicuri di non essere
gli unici ad averlo visto. A volte basta farlo presente: il relatore non potrebbe mai correggere un
bug, quindi se lo trova lo segnala, e va bene così. Ma non bisogna dare per scontato che ``è un bug così
evidente che qualcuno l'avrà già segnalato''. È il paradosso dell'incidente stradale: \textbf{più
l'incidente è grosso, più è probabile che nessuno chiami l'ambulanza}, perché tutti pensano che
qualcuno l'abbia già fatto\footnote{È il cosiddetto \emph{effetto spettatore} (\emph{bystander
effect}), studiato in psicologia sociale a partire dagli anni Sessanta.}.

\subsection{Domande sul ``me ne occupo io''}

\paragraph{Se correggo un bug, devo avvisare il collega?} Dopo una breve pausa per problemi audio
(``facciamo \emph{double testing}''), uno studente chiede se, correggendo il lavoro di un collega, il
dovere sia poi andare a dargli il feedback. A prescindere bisogna informare qualcuno; e la
conversazione dovrebbe essere: ``ho trovato questa cosa, vuoi che te la sistemi io?''. Correggere
senza dirlo a nessuno potrebbe essere un danno: magari era \emph{working as intended}, o qualcuno ci
stava già lavorando.

\paragraph{E se divento ``quello che risolve quel problema''?} Uno studente porta la sua esperienza:
in un'azienda in cui ha lavorato tre anni, risolvendo un problema facile che non era di sua
competenza è diventato, di fatto, l'incaricato di quel problema, distolto dalle sue mansioni. È
un'arma a doppio taglio. Per il relatore lì il problema è del \textbf{management}. In alcune aziende si
applica il criterio ``chi trova il problema lo risolve'' (``se te ne sei accorto tu, pensaci tu''), ed
è un altro modo di scaricare la responsabilità sugli altri. Non è ciò che dice lui: deve essere chi
trova il problema a valutare se è giusto e positivo che lo risolva lui. Il cartone di pizza lo
raccoglie perché ci mette dieci secondi; ma se ha trovato un bug nel sistema di pagamento che richiede
tre settimane di lavoro, e chi l'ha introdotto fa finta di niente, non deve risolverlo lui: deve far
presente che c'è un problema, che è stato ignorato e che va risolto (``pensateci voi, ma io vi dico
che c'è''). Va fatto caso per caso: se è immediato, facile e nel proprio interesse, lo si risolve.
Ma il relatore sconsiglia alle aziende la regola ``se lo trovi lo risolvi'', perché \textbf{disincentiva
a trovare i problemi}: qualcuno potrebbe non segnalarli per non avere altro lavoro e aspettare che li
trovi qualcun altro, ed è negativo quanto la situazione descritta. Per chiudere, ne discuterebbe con il
management: ``possiamo rivedere questa policy? Incentiva l'azione sbagliata''.

\paragraph{E se ci dicono di non mettere mano?} Uno studente che lavora come tester racconta che spesso
potrebbero correggere certe cose, ma viene loro detto di non toccarle. Esatto: si comunica. ``Ho trovato
questa cosa, ne siete consapevoli? Sì? Ci state già lavorando? No? Posso farlo io? Sì? Allora lo
faccio.'' Non bisogna esagerare con l'iniziativa: anche qui vale il \emph{default to trust}. Si parte
dal presupposto che qualcuno potrebbe averlo già visto e starci già lavorando, quindi si chiede e si
offre il proprio supporto.

% ---------------------------------------------------------------------
\section{Se si diventa manager}
\label{sec:z2-manager}

\subsection{Meriti, colpe e \emph{big picture}}

Una breve digressione per chi, un giorno, diventerà manager: che cosa cambierebbe? Oltre a tutto
quanto detto, ci sono il discorso su \textbf{meriti e colpe} e quello su \textbf{progetto e task}:
\textbf{delegare i meriti e prendersi le colpe}, e pensare al progetto più che ai singoli task, cioè
alla \emph{big picture}, invece che al ``voglio tornare a fare le cose che sapevo fare''.

È un altro problema dei manager promossi dal gruppo di chi sapeva fare bene un certo lavoro: quando
hanno l'occasione di tornare a farlo, lo fanno molto volentieri. ``C'è un task da fare, lo faccio io'',
magari con la scusa di dare l'esempio. No: si sta andando nella propria \textbf{comfort zone}. Il
relatore lo capisce (anche lui tornerebbe volentieri a fare ciò che faceva prima), ma bisogna accettare
che \textbf{quelle cose non le si deve più fare}.

\subsection{Il ruolo è più grande della persona}

Anche da manager, il ruolo che si ricopre per il gruppo è \textbf{più importante di sé}. Non importa
chi sia il community manager di Blizzard, di Riot o di un'altra azienda: conta ciò che rappresenta per i
giocatori in quanto community manager, non il fatto che sia lui piuttosto che un altro.

Prendersi cura del gruppo vuol dire applicare concetti molto semplici:
\begin{center}
  \begin{tabular}{@{}p{0.42\textwidth}p{0.48\textwidth}@{}}
    \toprule
    \textbf{Invece di dire\dots} & \textbf{\dots dire} \\
    \midrule
    ``Qui non si accettano errori.'' &
      ``Spiegami perché abbiamo sbagliato'': fammi capire perché è successo
      (sezione~\ref{sec:z1-default-trust}). \\
    \addlinespace
    ``Qui si fa come dico io.'' & ``Tu come faresti?'' \\
    \bottomrule
  \end{tabular}
\end{center}

\subsection{Il tetto della bravura del capo}

Il secondo punto è cruciale soprattutto quando si promuove il più bravo in un lavoro: se il criterio
noto è che si promuove il più bravo, il promosso lo sa, e la prima cosa che penserà è ``faccio lavorare
tutti come me, visto che sono il più bravo''. Il relatore chiede alla classe: che cosa si nega l'azienda
con questa dinamica? Anche ammesso che sia davvero il più bravo, se tutti fanno come lui diventeranno
tutti più bravi, no?

Il problema è che \textbf{si fissa come tetto massimo la bravura di quella persona}: se tutti fanno solo
come dice lui, si arriverà al massimo al suo livello, mentre ognuno ha il proprio metodo. Da qui quanto
detto il giorno prima sulla tendenza dei manager ad assumere solo persone peggiori di sé: se il criterio
di promozione è la bravura, si prenderanno sempre persone un po' meno brave per evitare che diventino il
proprio capo. Se invece le carriere sono divise (``se sei bravo fai questo, se vuoi fare il \emph{people
manager} fai quest'altro''), il manager è incentivato a prendere persone \textbf{migliori di sé}; e se,
come detto il giorno prima, il bonus dipende anche dalle performance del gruppo, è nell'interesse di
tutti che il gruppo sia pieno di superstar, un ulteriore incentivo ad assumere i più bravi.

\paragraph{L'opzione C.} C'è poi un'altra cosa, forse minima, ma per il relatore molto importante: si
perdono le \textbf{occasioni di innovare}. Nelle aziende precedenti era visto come il punto di riferimento
per il \emph{problem solving}: in un'emergenza, se c'era da scegliere tra A e B, ``decide Stefano''. Andava
benissimo, ma la tendenza era che a ogni emergenza venivano da lui. Così ci si negava \textbf{l'opzione C}:
quella a cui lui non avrebbe pensato e qualcun altro sì. Magari la prima volta sbaglia, la seconda anche,
e la terza ha un'idea migliore della sua, proprio perché lui non è più l'esperto di quel settore. Fissare
tutto come ``al massimo decide Stefano'' significa non accettare altre idee: un limite molto grosso
all'innovazione, all'aggiornamento e al fare le cose in modo più moderno e flessibile.

\subsection{Celebrare i successi e i fallimenti}

\paragraph{Celebrare i successi del team.} Il relatore vuole assicurarsi che tutti capiscano quando si è
ottenuta una vittoria o un successo: come detto il giorno prima, se non si sa che cosa si fa bene si
rischia di smettere di farlo. ``Questa cosa l'abbiamo fatta bene, complimenti, bravi tutti.''

\paragraph{Celebrare i fallimenti.} Sui fallimenti si era spinto ben oltre il semplice accettarli: aveva
indetto \textbf{meeting in cui si celebravano i fallimenti}, con le persone che raccontavano i propri errori
per distribuire quella conoscenza al resto del team. ``Non mi fa arrabbiare l'errore fatto, mi fa arrabbiare
lo stesso errore fatto quindici volte.'' Condividere un fallimento con il gruppo non significa ``guardate
questa persona sciocca che ha fatto questo errore'', ma ``questa persona ci sta permettendo di evitare lo
stesso errore in quindici'': sta salvando l'azienda dal ripeterlo quindici volte, perché la gente ha paura di
dire agli altri che ha sbagliato. Mostrare dove si è sbagliato diventava quindi un incentivo, un servizio
all'azienda. E si torna al punto di prima: se il proprio ruolo per il gruppo è più importante di sé, non
importa dover alzare la mano e dire ``ho sbagliato, ho fatto così'', perché gli altri penseranno ``grazie,
così io non lo faccio''. Si dovrebbe essere lieti che qualcuno abbia fatto quell'errore prima di noi, e
andare a ringraziarlo per averlo condiviso.

\paragraph{Condividere i successi.} Infine bisogna \textbf{dare il merito a chi lo ha}. Al relatore è
capitato spesso di ricevere ringraziamenti come ``grazie a Stefano e al suo team per aver fatto questa
cosa'', e la sua risposta è stata ``grazie, ma ha fatto tutto Gianni''. Non bisogna avere paura di dire
``grazie del feedback, ma è stato lui'': è giusto ringraziare lui e il team (il relatore si sente parte del
team, chiaramente). Non si ottiene niente dal dire ``grazie, riferirò i complimenti al team'': il relatore
non ha mai creduto nel ``riferire'' che il management ha detto ``ottimo lavoro''. Chiede al management di
dirlo \textbf{direttamente} al team, perché sentirlo dire dal management ha un impatto diverso che sentirlo
riferire da lui. In questo senso, osserva uno studente, ogni gruppo è il gruppo di sostegno di se stesso:
``assolutamente sì''.

\subsection{Quanto ci vuole?}

Non c'è una scadenza. Il rispetto non si ottiene pagando per fare prima, ``con il \emph{boost} a livello
60''. Servono tempo ed energia, ed è difficilmente misurabile: il relatore non può dire quanto il team si
fidi di lui in questo momento. Lo percepisce da come lavorano, ma lo si vedrà davvero solo quando ci sarà un
problema grosso; fino ad allora è tutta una sua percezione. Il team lo sente ripetere queste cose ogni giorno
e lo vede metterle in pratica ogni giorno, e lui può solo sperare che il messaggio passi; quando ci saranno
davvero ``gatte da pelare'' capirà come è messo, perché non c'è momento migliore di una difficoltà per vedere
come reagisce un team.

\textbf{Essere leader è prendersi un impegno a lungo termine verso i propri colleghi}: non ci sono modi per
velocizzarlo. Bisogna fare il proprio lavoro ogni giorno ``a testa bassa'', aiutando gli altri, facendolo
vedere ed \textbf{esplicitando perché lo si fa}. Il relatore lo fa quotidianamente: ``ti sto aiutando; nel
caso avessi un dubbio, lo faccio perché il mio lavoro è aiutarti e io ne traggo vantaggio dal fatto che tu
lavori meglio. Sono qui per te, ma lo faccio anche per me: non pensare che sia solo gentilezza. È il mio
lavoro, ma sappi che lo faccio spontaneamente''.

\subsection{L'ultima slide}

L'obiettivo dovrebbe sempre essere lo stesso: tutti dovremmo poter finire la giornata, staccare e tornare a
casa sentendoci
\begin{itemize}
  \item \textbf{ispirati}, perché abbiamo fatto qualcosa che ha un effetto tangibile;
  \item \textbf{sicuri}, perché sappiamo di avere colleghi, un leader e dei manager che si prendono cura di
        noi;
  \item \textbf{realizzati}, per quello che facciamo e perché è ciò che ci piace fare.
\end{itemize}
E questo dovrebbe essere \textbf{un diritto naturale di tutti}, non un ``sì, ma solo il manager'', ``sì, ma
solo se sei un esperto'', ``sì, ma solo se sei un game director''. Ognuno ha le proprie capacità: bisogna
solo trovare ciò che piace e che si vuole fare.

% ---------------------------------------------------------------------
\section{Domande e risposte finali}
\label{sec:z2-qa-finale}

Finita la presentazione, si passa alle domande: ``chiedete tutto quello che volete, perché è
l'ultima possibilità''. Le domande sono riportate per temi.

\subsection{Leadership e organizzazione del lavoro}

\paragraph{La leadership riguarda tutti o solo chi vuole diventare leader?} Il discorso sui
manager vale solo per i manager; ma il consiglio è \textbf{essere tutti leader del proprio gruppo,
da dentro}. Essere un leader significa sapere che gli altri possono fare affidamento su di sé, e
di queste figure in un team ce ne possono essere infinite: non è una gara a chi è il leader.
``Siate tutti una di quelle figure a cui gli altri guardano quando sono in difficoltà.''

\paragraph{Da manager non si hanno più orari e si fanno straordinari di continuo?} No: il relatore
ha notato di essere molto più efficiente come manager quando si dà \textbf{paletti molto precisi}.
Se stacca alle sei e alle sei e cinque arriva un messaggio, risponde ``se non mi dici che è
urgentissimo, non lo guardo nemmeno, lo leggo domani''; se non è urgente, lo guarda il giorno dopo.
Ma gestisce le aspettative: ``se è urgente ci sono e ci sarò sempre, ma non aspettatevi che faccia
orari eccessivi solo perché sono manager, così come non me lo aspetto da voi''. Si aspetta che
ognuno faccia ciò che deve; se è urgente tutti faranno di più, ma se non lo è, no.

\paragraph{L'orario davvero flessibile.} Qui il relatore aggiunge una cosa che si era dimenticato di
dire parlando di ferie e permessi illimitati: introduce anche l'\textbf{orario davvero
flessibile}. Non il flessibile di molte aziende, in cui si può iniziare alle 9 e finire alle 18
oppure iniziare alle 10 e finire alle 19: se oggi si vuole iniziare alle 10 e finire alle 15, e si è
fatto tutto quello che si doveva, va bene così. Ci si prendono le ore libere in più per altro, per
staccare, per riposarsi, sapendo che se la settimana dopo bisognerà lavorare dalle 9 alle 21, lo si
farà (o si dirà ``ieri ho fatto dalle 9 alle 21, domani faccio dalle 10 alle 15''). L'obiettivo dato
è un lavoro e un progetto, non un orario: si lavora per completare i propri task e rispettare le
deadline. L'orario indicato è solo indicativo, ed è quello dei meeting: \textbf{nessun meeting prima
delle 10 o dopo le 17}. Se poi qualcuno vuole lavorare dalle 10 alle 11 e poi di notte, perché quel
giorno preferisce lavorare di notte, più rilassato, faccia pure: l'importante è che trovi il tempo per
riposare e staccare. Ci si autogestisce senza timbrare cartellini; l'unico accordo è ``hai un task, va
consegnato entro il giorno X e mi garantisci che venga fatto; come ci arrivi non mi interessa''.

\paragraph{Le è mai capitato di essere aiutato da qualcuno, in un momento di difficoltà?} Sì: come
detto, questo comportamento è \textbf{contagioso in senso positivo}. Se ogni volta che qualcuno è in
difficoltà qualcun altro lo aiuta, prima o poi capita anche a chi aiuta. Il relatore non lo dice
``tanto per dire'': ha avuto momenti in cui faceva davvero fatica ad alzarsi dal letto per lavorare,
perché esausto o demotivato per altri motivi, e il team gli ha detto ``stai a letto, ci pensiamo noi
oggi''. Questo lo aiuta ad aiutare, perché sa che la cosa è bidirezionale.

\paragraph{Tra le soft skill del workshop, da quale partire?} \textbf{Il feedback}: saperlo dare e
ricevere. È la più difficile, molto più del ``voglio aiutarti'', ma è anche quella su cui si costruisce
tutto il resto: attraverso il feedback si costruisce la \textbf{fiducia}, attraverso la fiducia il
\textbf{rispetto}, e il rispetto diventa la base portante di tutto il resto.

\paragraph{E se il leader se ne va?} A una domanda su che cosa fare quando viene meno la figura di
riferimento del gruppo\footnote{La domanda non è leggibile nella trascrizione; dalla risposta si
capisce che riguardava il caso in cui il leader che incarnava la visione del gruppo se ne va.}, il
relatore risponde che la domanda da farsi è un'altra: \textbf{è cambiato l'obiettivo?} La direzione è
sempre la stessa? Se sì, chi è la persona più adatta a portarci lì? Non conta tanto chi guida, ma
l'obiettivo che ci accomunava, di solito la visione di un leader in cui tutti si allineano: ci
crediamo ancora tutti, o è cambiato qualcosa? È il momento di fare autoanalisi: ``ero qui per quella
persona o per la visione?''. Si dice, ed è vero, che \textbf{le persone non si dimettono dai posti di
lavoro, ma da capi che non rispettano più}, perché spesso l'unica interfaccia con l'azienda è il
proprio capo: le dimissioni non sono un messaggio per l'azienda, ma per il capo (``non voglio più
lavorare con te''). Allo stesso modo, se non c'è più il leader della visione, bisogna prima capire se
si condivide ancora qualcosa.

\paragraph{Lavorare così è più faticoso?} Sì e no. ``Non voglio demotivarvi, perché non è il mio
lavoro'', ma non è tutto rose e fiori: lavorare in questo modo è anche più stancante
\textbf{emotivamente}, perché non si può più pensare ``faccio il mio, alle cinque ho finito e chi si è
visto si è visto''. Si sente un senso di responsabilità verso gli altri. Succede che si finisca il
proprio lavoro del pomeriggio, ci si giri e si veda un collega che finirà alle nove di sera, e
l'istinto sarà ``posso farti un pezzo del lavoro?'': così si finisce entrambi alle sei. In un certo
senso ci si accolla più lavoro, con la consapevolezza che quando si sarà in difficoltà ci sarà qualcuno
che si girerà a dire ``ti vedo stressato, posso darti una mano?''. È bellissimo, ma non dà solo benefici
diretti: richiede anche un investimento personale che non sempre le persone vogliono fare. Pro e
contro: è il modo che sceglierebbe il relatore, ma non è per tutti. Conosce persone che vogliono poter
staccare appena finito, e va benissimo: si cerca un'azienda in cui si possa lavorare così.

\subsection{Condividere fallimenti e successi in pratica}

\paragraph{Come si celebrano i fallimenti e si condividono i successi?} Condividere i successi sembra
più semplice (``bravo Mario che ha fatto questa cosa''). Il relatore dà uno spunto ancora migliore:
\textbf{``bravo Mario, e ora Mario ci racconta come l'ha fatto e com'è andata''}. Sono sottigliezze,
ma contano: raccontare in prima persona ciò che si è fatto dà un moto d'orgoglio che non dà sentirlo
raccontare da un altro.

Per i fallimenti, che sembrano dover partire dalla persona stessa, l'obiettivo era che tutti li
condividessero spontaneamente; ma all'inizio era difficile, perché non viene naturale alzare la mano e
dire ``ho sbagliato''. Per questo il relatore ha creato artificialmente dei \textbf{``festival del
fallimento''}, dal tono chiaramente ironico: durante una pausa pranzo, mentre tutti mangiavano insieme,
a turno si saliva ``sul palco'' a raccontare la stupidaggine fatta in settimana. Molto spesso iniziava
lui, raccontando i propri errori (``ho sbagliato questa cosa, sappiate che non si fa così''). Dopo aver
forzato l'interazione per una, due, tre settimane, alla quarta lo facevano da soli: era diventato
organico dire a tutti ``ragazzi, attenzione, ho fatto questa sciocchezza, non fatela anche voi'', e c'era
perfino un \textbf{tag apposito su Slack}. È stato il team a decidere come condividere i fallimenti: come
sempre, invece di imporre il modo, il relatore ha lasciato scegliere a loro, e funziona. È impegnativo
tanto quanto lo si vuole rendere tale.

\paragraph{Che tipo di fallimenti si condividono? Anche quelli personali?} Di solito si parte da quelli
\textbf{professionali}, cioè da ciò che si riconosce come utile anche agli altri: qualcosa che ha
rallentato, fatto perdere tempo o che rischia di ripetersi (``sappiate che, per evitare che lo facciate
anche voi\dots''). Il resto è più da conversazione privata, \emph{one to one}, con la persona con cui ci
si vuole sfogare, o con gli amici: se non ha un potenziale impatto sugli altri, non c'è un incentivo a
dirlo. E un errore che si è fatto per il proprio tipo di personalità? Qui tornano i test di personalità
della prima giornata: conoscendo il proprio profilo e quello dei colleghi, si sa che qualcuno ha lo
stesso carattere, e quindi lo si aiuta.

\subsection{Cultura aziendale e formazione}

\paragraph{Come si arriva a poter fare queste cose in azienda?} Conta molto la cultura aziendale. Il
relatore sa di lavorare bene perché oggi ha il \emph{background} e i successi alle spalle per dimostrare
che il suo metodo funziona, e quindi glielo lasciano fare; ma se dieci anni prima avesse proposto le
stesse cose, ``mi avrebbero mandato a quel paese''. C'è un momento e un posto per fare le cose, e non è
sempre possibile fare ciò che si ha in testa; ma ci si può lavorare. Per lui la domanda è diventata:
``sono convinto che così sarebbe meglio; come faccio ad avere carta bianca?''. \textbf{Portando
risultati}: si è concentrato su quello, e ci è arrivato progressivamente.

\paragraph{Come funziona la formazione con workshop come questo, in azienda?} Il relatore li fa in modo
un po' più rilassato ed esteso, con molte più interazioni ed esempi, perché in azienda si possono usare
esempi che si hanno davanti ogni giorno, e parlare di problemi e discussioni diventa molto più reale. Ma
in generale è sempre \textbf{formazione pratica}, ripetuta finché il messaggio non è recepito. Il segno
che qualcosa è stato recepito è quando arriva una persona nuova e \textbf{non è lui a spiegarle le cose}:
è il team stesso a spiegarle che cos'è la feedback culture, e lui non deve parlare. È la conferma che il
concetto viene perpetuato dal team.

\paragraph{Quanti manager dell'industria condividono questi principi?} Il relatore ha fatto molti
interventi all'IGDA di Roma e ne farà uno a Milano la settimana successiva (``il 24, mi pare''); interviene
spesso a workshop ed eventi, e il riscontro è sempre ``che bello, vorremmo farlo anche noi in Italia'', e ci
stanno lavorando. La moglie lavora in un'azienda che non ha niente a che fare con i videogiochi e lo stanno
già applicando. Magari in Italia è tra i primi a parlarne, perché il paese è ancora un po' indietro, ma non è
assolutamente l'unico: è più diffuso di quanto si pensi, non è una cosa di nicchia.

\paragraph{Questi spunti si applicano anche all'IT e alla consulenza?} ``Praticamente tutti.'' Si applicano
facilmente quando c'è un gruppo di persone su cui applicarli; se si cambia continuamente cliente è più
difficile portare questa mentalità da loro, ma questo non impedisce a chi lavora con noi di applicarla
direttamente con noi.

\paragraph{Ci si può tenere in contatto?} Su LinkedIn: il relatore accetta sempre tutti e risponde alle
domande quando riesce (ne riceve tante, ma cerca di stare al passo almeno su base settimanale). E, ``a costo
di sembrare quello che si fa pubblicità'', chi lavora in un'azienda che fa fatica ad applicare questi
concetti o non sa che esistono può metterla in contatto con lui: è spesso così che ha trovato nuovi modi di
applicarli. Vale anche per chi tra sei mesi lavorerà altrove. I suoi profili sono pubblici e non c'è nessun
problema a contattarlo.

\subsection{Performance e persone in difficoltà}

\paragraph{Se un'azienda è a corto di personale, conviene comunque eliminare chi performa poco?} Il giorno
prima si era detto che conviene separarsi presto da chi non performa (sezione~\ref{sec:z1-squadra}); ma non
si finisce per caricare di lavoro tutti gli altri? Per il relatore \textbf{eliminare non vuol dire non
sostituire}: vuol dire cambiare la persona con qualcuno magari più junior, ma più adatto come \emph{fit}
culturale. Il vecchio discorso del ``fare lo stesso lavoro con meno gente'' non è sostenibile: lo si può fare
fino a un certo punto, in caso di necessità, ma non è una soluzione. Non è ottimizzare, è far andare in
\emph{burnout} chi resta.

\paragraph{Ha più senso tenere una persona poco performante e non in sintonia con il team, piuttosto che
licenziarla senza poterla sostituire?} No: se una persona è poco in sintonia, performa poco ed è in generale
tossica e negativa per il team, \textbf{prima ci si separa, meglio è}. Questo non vuol dire non provare a
recuperarla, ma bisogna andare dritti al punto: ``abbiamo un problema: ti vedo poco performante, non ti vedo
allineato con il team, vedo che il team non si fida di te. Che cosa vogliamo fare?''. Se la risposta è ``non lo
so'', la strada porta dritta all'uscita. Bisogna \textbf{recidere i rami morti} il prima possibile, facendo
però attenzione a essere sicuri che si tratti di un ramo morto e non di un ramo recuperabile: tagliare un ramo
recuperabile mette negli altri la paura di essere i prossimi, ``alla prima foglia gialla''. Quello che si
vuole è che, quando un ramo è morto, la pianta smetta di alimentarlo, e ne crescerà un altro. Entrambi i
messaggi devono essere chiari: \textbf{non si subisce passivamente la situazione}, e allo stesso tempo
\textbf{non si abbandona nessuno} (se qualcuno è in difficoltà, la prima soluzione non è tagliare il ramo). È
un compromesso tra le due cose. (A questo punto arriva, puntuale, il crash della connessione: ``non era ancora
successo oggi''.)

\paragraph{Che cosa fare se qualcuno nel team sembra soffrire di depressione?} È un discorso complesso, perché
il relatore \textbf{non è uno psicologo}. Può dare la sua disponibilità a concedere permessi, coprire il lavoro
della persona e in generale darle modo di seguire un percorso con uno specialista: non può risolvere la
depressione, ma può alleviarne la parte lavorativa e permettere di ricevere supporto da chi è qualificato. È lo
stesso motivo per cui, a Blizzard e a Riot, quando arrivavano ticket con minacce di suicidio (``mi avete
bannato, quindi mi ammazzo'', ticket purtroppo abbastanza comuni), il \emph{customer service} era istruito a
non intervenire direttamente: non dovevano essere loro a fermare la persona, perché non qualificati, ma
indirizzarla a qualcuno di qualificato, prendendo tempo. Una persona non qualificata che prova a farlo rischia
di peggiorare la situazione; come davanti a un incidente si chiama l'ambulanza e non si fa una tracheotomia.
Per il relatore salute mentale e salute fisica sono la stessa cosa: come non si azzarderebbe a operare qualcuno
di appendicite, non si azzarda a fargli da terapeuta.

\paragraph{Come si danno feedback adatti a ogni tipo di personalità?} Bisognerebbe leggere i risultati del
profilo della persona sullo stesso sito su cui si è fatto il test di personalità della prima giornata, che
spiega anche come interagire con quel profilo. In generale: con chi è più \textbf{analitico}, portare dati a
supporto; con chi è più \textbf{emotivo}, metterlo in una condizione di comfort; con chi è molto
\textbf{pratico}, andare dritti al punto senza girarci intorno. Poi fare pratica. Oppure il modo ``perfetto''
per il relatore: provarci come viene e poi \textbf{chiedere} ``andava bene così o preferisci in un altro modo?
Potevo fare meglio?''. Bisogna creare anche quel rapporto: fidarsi di discutere anche di queste cose.

\paragraph{Se nel team fossimo quasi tutti leader, non ci sarebbero ``troppi galli''?} No, perché per leader si
intende chi si prende cura degli altri, e magari lo fossero tutti. Non significa che tutti scelgano che cosa
fare: spesso lo sceglierà comunque il manager, o guiderà la scelta. Ma avere tutte persone con iniziativa e
capaci di scegliere quando serve è qualcosa per cui il relatore ``metterebbe la firma''. Non ha mai corso il
rischio dei ``troppi galli'': se ci sono quei caratteri, il team è stato assemblato male, perché non tutti sono
compatibili con questo modo di lavorare e magari non si hanno le persone giuste per applicarlo.

\paragraph{Se la fiducia richiede tempo ed esperienza, quanto è difficile per un junior?} Meno di quanto si
pensi: proprio perché su un junior non ci sono aspettative, un junior che chiede ``posso darti una mano?'' dà un
segno chiaro che tiene più al team che a se stesso. Si può essere leader anche non dal punto di vista
professionale, ma come chi va a dare supporto agli altri: ``il medico del team, l'\emph{healer}'', e lo si può
fare anche da junior. È così che ha iniziato il relatore: andare ad aiutare a risolvere problemi non suoi gli ha
fatto accelerare molto il percorso per ottenere rispetto e stima, e lo ha esposto a cose a cui altrimenti non
sarebbe mai arrivato. È una questione di carattere: se ci si tiene agli altri, lo si fa vedere, e il resto
arriva come risultato. ``È quasi un reddito passivo'': un investimento nel far vedere che ci si è, e il resto
arriva (``Stefano si è preso cura di me, quindi ora mi prendo cura di lui'').

\paragraph{Aiutare gli altri solo dopo aver finito il proprio lavoro?} Dipende dall'urgenza. Il relatore spesso,
se qualcosa è davvero urgente, mette in pausa ciò che sta facendo, risolve e poi torna al suo lavoro; ma è troppo
\emph{case by case}. Se due colleghi litigano su chi debba avere le ferie di Natale, lo si può risolvere anche
domani; se litigano perché si è perso il \emph{repository} e domani c'è la consegna, lo si risolve subito.

\paragraph{Come mantenere motivato un team su progetti non retribuiti, come quelli accademici?} Essendo il più
\textbf{trasparenti} possibile nei processi decisionali: se si deve decidere se fare un action o un RPG, tutti
devono sapere perché si fa l'uno invece dell'altro (da qui il metodo di brainstorming della mattina). Deve essere
sempre chiaro a tutti perché si stanno facendo le cose in quel modo, perché \textbf{la demotivazione nasce dal
disallineamento} rispetto all'obiettivo del gruppo, e l'unico modo di essere allineati è essere sempre
consapevoli e partecipi del processo che decide dove sta andando il gruppo. E soprattutto: \textbf{non chiudere
mai un meeting senza aver concordato che cosa si vuole fare} e senza essere tutti d'accordo. Se non lo si è,
bisogna continuare a parlarne, per quanto scomodo, lungo e noioso: si deve sempre uscire da un meeting
sapendo che si è tutti d'accordo, o al massimo in \emph{disagree and commit} (sezione~\ref{sec:z2-disagree}).
``Quello è il massimo accettabile.''

\subsection{Carriera, colloqui e candidature}

\paragraph{Come ha iniziato a lavorare all'estero?} Negli annunci online si trovano spesso requisiti come un titolo
di studio conseguito nel paese o l'abilitazione a lavorarci. Questo vale per i paesi fuori dall'Unione europea: per
lavorare negli Stati Uniti o in Sud America serve un visto o una \emph{green card}. In molti casi, però, è proprio
ottenere il lavoro a permettere di lavorare lì: quando il relatore doveva andare a lavorare in California (a Irvine,
a Santa Monica, dove capitava), era l'azienda a occuparsi dei permessi necessari, perché ``per venire qui devi avere
un lavoro, e il lavoro te lo diamo noi''.

Il suo inizio è stato abbastanza semplice. In Italia lavorava per Digital Bros, occupandosi dei giochi online gestiti
dall'azienda e di quello che era il canale Game Network sul digitale terrestre\footnote{Questo passaggio chiarisce il
dubbio segnalato nel capitolo~\ref{cap:z1-feedback} sul ``Digital Game Network Channel'': si tratta di Game Network,
canale televisivo dedicato ai videogiochi legato al gruppo Digital Bros.}. Vide un annuncio a Dublino per
\emph{Warhammer Online}, mandò la candidatura, fece i colloqui e gli dissero ``ti paghiamo la \emph{relocation}, vieni
qui''. Come detto il giorno prima, sono le aziende a occuparsi del trasferimento: se devono creare a Dublino un'unità
che segue dodici lingue europee, sanno di dover far arrivare persone da fuori, ed è già a budget. Non è nemmeno una
richiesta da fare: è prevista. Per il relatore, da questo punto di vista, è indifferente lavorare per un'azienda
italiana o estera: se si è adatti, si è adatti. Per questo insisteva sull'inglese: se lo si parla, è indifferente di
che cosa ci si occupi e quale sia la propria lingua madre, a meno che la posizione non sia esplicitamente ``community
manager italiano''. Un programmatore polacco o italiano non cambia niente. Lavorare all'estero non è più difficile: c'è
più concorrenza solo se l'azienda cerca attivamente persone da tutto il mondo e non si pone problemi di \emph{location}.
Cambiano l'azienda, il fatto di lavorare da remoto o dal vivo, l'eventuale trasferimento; per il resto non cambia
molto.

\paragraph{Gli anni passati nell'IT contano se si vuole passare ai videogiochi?} Dipende da che cosa si vuole fare: se
da responsabile IT si vuole diventare programmatore, bisogna saper dimostrare di saper programmare. Vale anche per i
passaggi interni: se un game designer vuole diventare art director, al relatore va bene, deve solo dimostrare di saper
fare l'art director. Anzi, incentiva le \textbf{promozioni interne}, perché si conosce già il candidato; ma gli
spostamenti non sono liberi ``a sentimento'': bisogna saper fare quella cosa. I recruiter valuteranno non tanto gli anni
quanto ciò che si sa fare; gli anni servono comunque, perché dimostrano di saper lavorare in un ambiente professionale,
che è un altro insieme di competenze. Ma se si sa fare ciò che chiedono, aver lavorato nell'azienda X o Y non dovrebbe
cambiare molto.

\paragraph{Quanto vengono considerate le candidature spontanee?} Non tantissimo. Soprattutto per le aziende non grandi,
se non stanno cercando qualcuno è perché non se lo possono permettere: una candidatura spontanea non le convincerà ad
assumere. Più l'azienda è piccola, più conviene aspettare l'annuncio, perché se non c'è vuol dire che non c'è budget, e
potrebbero rispondere ``ti vorremmo, ma non possiamo prenderti''. C'è anche il rischio che finisca nel vuoto (l'HR magari
non controlla nemmeno la mail, perché non sta assumendo) o, peggio, che rispondano ``non ci servi'' e una settimana dopo
esca la posizione, con il candidato già nel cestino dei rifiutati. Si rischia di \textbf{bruciarsi} la possibilità.

\paragraph{Come ci si accorge di essere riconosciuti come leader, e come lo si ``vende'' a un colloquio?} Uno studente
osserva che la leadership, così descritta, non la si riconosce da sé ma la riconoscono gli altri. Internamente il
relatore non ha mai dovuto ``venderla'', perché parlavano i risultati. Finché non lo si vive è difficile capirlo, ma
quando si arriva alla posizione in cui il team rispetta e vede qualcuno come punto di riferimento, non ci sono grandi
dubbi: è evidente che vengono da lui. Si può riflettere su \textbf{quanto spesso gli altri vengono da noi}, cercano di
emulare il nostro modo di fare, vanno a rivedere le nostre cose per capire dove hanno sbagliato, chiedono un parere, ci
parlano di un problema: la somma di tutto questo significa molto spesso che si fidano. Per ``venderla'' all'esterno si
può dire di essersi occupati di un team in cui, attraverso il lavoro e il rispetto degli altri, si è arrivati a essere
visti come punto di riferimento; ma il suggerimento è \textbf{portare esempi tangibili} (``c'era questo progetto, non
sapevamo come risolverlo, il team mi ha chiesto una mano e l'ho risolto così''), invece di ``incensarsi'' dicendo ``sono
il punto di riferimento del team''. Se il relatore racconta di essere stato eletto rappresentante dei licenziati senza
essere tra loro, un'azienda capisce che si fidavano di lui, punto: non ci sono altre conclusioni da trarre.

\paragraph{Quali sono gli errori più gravi dei junior ai colloqui, a livello di soft skill?} In ordine di priorità:
\begin{enumerate}
  \item \textbf{non arrivare preparati}, cioè \textbf{senza esempi}. Il relatore ha sempre un esempio pronto per ogni
        cosa di cui vuole parlare: per la fiducia, per gli \emph{achievement}, per i risultati. Servono sempre esempi
        tangibili per far capire che cosa si intende;
  \item \textbf{non avere domande alla fine}, perché dimostra di non avere chiaro che cosa si andrebbe a fare. Anche solo
        chiedere che cosa si farebbe, che cosa si può migliorare, o ``che impressioni hai avuto finora dal mio
        colloquio?'': si possono fare domande molto aperte al recruiter per capire come sta andando;
  \item \textbf{non leggere bene la \emph{job description}}: presentarsi con vaghezza, senza aver capito che cosa cerca
        l'azienda. Lo si impara anche con la pratica; ma il designer che si candida per una posizione da programmatore,
        senza nulla nel CV che faccia pensare che sappia anche programmare, il relatore non solo lo scarta: lo mette in
        una \textbf{blacklist} di persone che forse non considererà mai. Se non si riesce a interpretare un testo come una
        \emph{job description}, è un problema troppo grosso da superare in una collaborazione;
  \item \textbf{errori di battitura e CV impaginati male}, che molti sottovalutano. Il relatore ha letto letteralmente,
        in una lettera di presentazione per una posizione da \emph{copywriter}, la frase ``sono una persona molto attenta
        ai dettaggi''. Come si può considerare quella persona per quel ruolo? Non ha nemmeno incollato la lettera in Word
        per il controllo automatico: non è un errore sfuggito al correttore perché la parola esiste con un altro senso,
        ``dettaggi'' viene segnato in rosso in qualsiasi lingua. Come affidarle la comunicazione aziendale? Vale per tutti
        gli errori: ciò che si manda è \textbf{il primo biglietto da visita} con l'azienda, e non si manda un biglietto da
        visita sgrammaticato. Un conto è non andare ai colloqui in giacca e cravatta (``non l'ho mai fatto e probabilmente
        non lo farò mai''), un altro è presentarsi con errori;
  \item un altro esempio comune: \textbf{l'email sbagliata nel CV}. Se la mail torna indietro, il relatore non telefona
        per farsi dare quella corretta: se non si è in grado di scrivere giusta la propria email, non vale la pena. Più
        l'errore è banale e facilmente evitabile, più lo fa arrabbiare (``scusate se mi sono inalberato''): evidentemente
        il candidato non teneva abbastanza a quel posto.
\end{enumerate}

\paragraph{Quanto è importante la lettera di presentazione?} Sempre di più: sempre più aziende la richiedono come
obbligatoria, quindi conviene prepararne una, sapendo di doverla adattare al destinatario. Deve essere \textbf{meno lunga
di una pagina}, in \textbf{tre paragrafi}:
\begin{enumerate}
  \item \textbf{chi siete}: un minimo di \emph{backstory}, che cosa avete fatto, che cosa volete fare, da dove arrivate;
  \item \textbf{perché avete mandato il CV}: che cosa vi ha interessato di questa azienda;
  \item \textbf{perché dovrebbero chiamare voi}: che cosa dovrebbe spingere l'azienda a scegliere proprio voi.
\end{enumerate}
Uno studente dice di riuscire bene con i primi due, ma di trovare difficile il terzo, che è poi la domanda che tutte le
aziende fanno. In realtà basta \textbf{collegarsi alla \emph{job description}}: se cercano un 3D artist in grado di fare
X, Y e Z, si conferma di averlo letto e si spiega che con la propria esperienza si è imparato a fare proprio quelle cose.
Non vogliono sentire niente di particolare, ``non devi cambiargli il mondo''.

\paragraph{Che cosa guarda un recruiter in un portfolio?} Il relatore li guarda poco: fa un \emph{sanity check} generale,
per vedere che non sia evidente perfino a lui, ``a digiuno di tante cose tecniche'', che c'è un problema (se se ne accorge
perfino lui, il problema c'è). Per il resto guarda che il portfolio sia \textbf{coerente con il CV}: se arriva un CV da
senior con dodici progetti completati e anni di lavoro, e il portfolio sono quattro immagini, non è coerente. Vuole una
concordanza tra ciò che si racconta e ciò che si vede.

\paragraph{Che strumenti per creare un CV?} \textbf{Canva}, che ha un catalogo di template gratuiti e a pagamento
amplissimo e personalizzabile: anche se non piace, lo si prende e lo si adatta, ma fa già il 90\% del lavoro. Insieme a
\textbf{ChatGPT}, perché ``è inutile far finta che non esista'': ci si fa scrivere la lettera di presentazione e poi la si
corregge e la si adatta, e lo stesso per il CV. L'importante è renderli adatti e coerenti, non lasciarli solo frutto di
ChatGPT; ma per fare il grosso del lavoro è un ottimo strumento.

\paragraph{Conviene lasciare nel CV esperienze non del settore?} Uno studente ha alle spalle il liceo, vari stage e
lavori che non c'entrano niente con la modellazione 3D. Si possono mettere finché c'è spazio: il relatore consiglia
\textbf{al massimo due pagine}, e quando ha cominciato a sforare nella terza ha iniziato a togliere ciò che non era più
rilevante. Ormai nel suo CV la formazione è un paragrafo: ha così tante esperienze lavorative che non interessa a nessuno
che abbia fatto il liceo scientifico o classico, e non lo mette nemmeno più. Si filtra in base a ciò che è più rilevante
per lo spazio disponibile.

\paragraph{Che cosa fare con le aziende che dicono ``ti faremo sapere'' e poi spariscono?} Non aspettare uno o due anni
per scrivere. Il suggerimento: tra le domande finali del colloquio, l'ultima del relatore è sempre \textbf{``tra quanto
mi fate sapere?''}. Se rispondono dieci giorni, all'undicesimo scrive: ``mi avete detto dieci''. Non si sente in difetto a
mettere pressione, perché sono stati loro a dare la scadenza. E se viene ignorato del tutto, quell'azienda per lui è
``morta'': non si ricandiderà mai, perché non vede perché il rispetto debba andare in una sola direzione.

\paragraph{Quanto essere trasparenti, nel CV e nella lettera, su problemi personali come l'ADHD?} Uno studente con
ADHD\footnote{Nella trascrizione ``essendo 10D''; nel resto della risposta ``di'' e ``dei 10D'' stanno per ADHD.} dice di
tendere ad avere performance altalenanti. Il relatore lo invita a \textbf{non usare l'ADHD come scusa} per essere
altalenante: anche lui ha l'ADHD, ma non è altalenante. Esserlo è senz'altro un \emph{malus}, ma esistono i mezzi per non
avere problemi di questo tipo: lui ha trovato i suoi metodi (la to-do list, la \emph{must do list}, il delegare il lavoro
al team, perché occuparsi delle cose per gli altri lo mantiene molto motivato, sveglio, all'erta, attento a come vanno le
cose). Sta anche a ciascuno trovare un lavoro in cui compensare ciò che viene richiesto: se domani lo rimettessero a fare
il programmatore, come all'inizio, probabilmente impazzirebbe, perché non riusciva a restare concentrato sul codice oltre
un certo tempo. Lì ha capito che non poteva fare il programmatore, ma ha trovato qualcos'altro che gli permette di essere
com'è e fare ciò che fa. Bisogna capire i propri punti di forza e di debolezza e giocare su quelli: si possono avere tutti
i difetti del mondo e voler comunque fare una cosa, ma bisogna accettare che sarà più difficile.

Quanto alla trasparenza, il consiglio è \textbf{non mentire mai nel CV}, ``perché tanto vi beccano''. E per il discorso della
\emph{bandwidth}: tutte le energie spese a mentire, e a tenere viva la bugia, sono energie tolte allo spiegare bene ciò che
si sa fare. È un investimento nel posto sbagliato. Bisogna concentrarsi su ciò che si sa fare, con esempi tangibili, e dire
``non lo so fare'' quando non lo si sa fare. Il peggio non è essere scoperti, ma \textbf{non esserlo}: se si finge di sapere
il C++, si viene assunti e messi a lavorare in C++, che cosa si fa il giorno dopo? Si rischia davvero di rovinarsi la
carriera. Se non si è un \emph{fit}, si sarà il \emph{fit} per la prossima. L'incubo del relatore è essere assunto in un
posto dove pensano che sappia fare qualcosa che non sa fare; anzi, tende a esplicitarlo: ``cercate questa cosa? Io non la so
fare, va bene lo stesso?'' --- ``Sì, te la insegniamo noi'' --- ``Voglio solo che sappiate che non la so fare''. Meglio
questo che vivere con l'ansia che prima o poi si accorgano che non sa programmare.

Gli capita spesso con chi mente sull'inglese. Al primo colloquio chiede ``sai parlare inglese? A che livello?'' --- ``C1,
\emph{fluent}, quasi madrelingua'' --- ``Perfetto'', e lascia cadere l'argomento. All'ultimo colloquio: ``ti va se lo facciamo
in inglese?'' --- ``No, ma io intendevo che\dots'' --- ``No, mi avevi detto \emph{fluent}: ora facciamo il colloquio in
inglese''. Prima o poi le bugie raggiungono chi le dice.

% ---------------------------------------------------------------------
\section{Chiusura del workshop}
\label{sec:z2-chiusura}

Il relatore chiede alla classe di far sapere com'è andata a Vanessa, il punto di contatto con l'accademia per il workshop:
sarà lei a riferirgli eventuali lamentele o suggerimenti, che gli permetteranno di fare meglio i prossimi. Ringrazia per aver
``retto per dodici ore'' e spera che sia stato utile, consapevole che tutti sono stanchi quanto lui.

Gli studenti rispondono che è stato chiarificatore su molte cose che non sapevano, soprattutto su come interfacciarsi nei
team, e che sarà materiale da rivedere per fissare bene i concetti, vista la mole in due giorni; qualcuno si aspettava che
sarebbe stato molto più pesante e l'ha trovato invece scorrevole. Il relatore ribadisce che \textbf{niente sostituirà
l'esperienza diretta}: il workshop dà un'infarinatura e un punto di vista, ma i concetti si metteranno in pratica solo nel
\emph{day to day}. Bisogna avere \textbf{pazienza con se stessi}, sapendo che non verrà tutto bene dal primo giorno.
Confessa di essere stato un po' preoccupato per la durata, ma il materiale c'era, e farlo in un giorno avrebbe rischiato di
trasformarlo in una corsa: ``rischiava di essere uno \emph{sprint}, e invece è una maratona''.

Infine ricorda che si può aggiungerlo su LinkedIn, che registrazione e slide (in forma integrale) verranno fatte avere
tramite Vanessa (``dite che vi ho autorizzato io a chiederle'') e che, per qualsiasi dubbio, resta a disposizione. ``Buon
resto del weekend.''

% ---------------------------------------------------------------------
\section*{Punti chiave del capitolo}
\addcontentsline{toc}{section}{Punti chiave del capitolo}

\begin{itemize}
  \item I colloqui sono una \textbf{selezione reciproca}. L'annuncio dice molto sulla cultura aziendale; alla fine del
        colloquio vanno sempre fatte domande (il \emph{day to day}, la ``bacchetta magica'', ``tra quanto mi fate
        sapere?''). Si arriva con una \textbf{cifra minima mentale} e con \textbf{esempi tangibili}; non si mente mai.
  \item Time management non è una lista di cose da fare: si parte dalla \textbf{must do list}, che è anche il primo
        momento per dire ``oggi potrei non farcela''; si stimano i tempi e si chiede quando qualcosa non è chiaro.
  \item Ogni nuovo task va \textbf{prima aggiunto alla lista}, non eseguito: priorità non è recenza. Le
        \textbf{checklist} fanno risparmiare tempo, anche (e soprattutto) ai più esperti.
  \item \textbf{First contact resolution}: raccogliere tutto ciò che serve in una sola richiesta, e non iniziare un task
        senza le informazioni necessarie. Ogni interruzione costa dieci-quindici minuti di concentrazione.
  \item \textbf{Matrice di Eisenhower}: fai, pianifica, delega, elimina. ``Non importante'' significa che non è importante
        che lo faccia tu; il quadrante dimenticato è \emph{delete}.
  \item Le pause vanno \textbf{scelte, non subite}; i task grossi vanno spezzettati; bisogna \textbf{saper dire di no},
        anche ai meeting (sono necessario? può andarci uno solo? ci sarà un resoconto scritto? per che cos'altro potrei usare
        questo tempo?). \textbf{Il multitasking non esiste}.
  \item A fine giornata si verifica la \emph{must do list}: un \emph{must do} rimandabile non era un \emph{must do}.
  \item Ferie illimitate con un minimo, giornate libere e feedback culture sono, per il relatore, le condizioni per un team
        sano; il benessere è anche mentale.
  \item Il \textbf{brainstorming genera idee, non decide}. Va preparato (topic, obiettivo, vincoli), separa la
        \textbf{generazione} dalla \textbf{discussione} (\emph{rapid ideation} di cinque minuti, \emph{dot voting}
        anonimo) e combatte l'\emph{anchoring} e il dominio degli estroversi.
  \item La decisione si prende con la \textbf{matrice impatto--sforzo}: prima le idee a basso sforzo e alto impatto, poi in
        diagonale. Lo sforzo sono le risorse, l'impatto comprende anche la motivazione; decide chi farà il lavoro.
  \item \textbf{Disagree and commit}: si dichiara il disaccordo, poi ci si impegna al massimo. Se va male ha sbagliato il
        team, non ``ve l'avevo detto''.
  \item \textbf{La leadership non è un titolo}: il rispetto si guadagna. Le carriere vanno divise tra \textbf{manager} e
        \textbf{subject matter expert}, pagati allo stesso modo, così che diventi manager solo chi vuole prendersi cura degli
        altri.
  \item Dopamina, serotonina e ossitocina legano il gruppo; il \textbf{cortisolo} è contagioso e lo disgrega. Il
        ``contratto sociale'' del villaggio vale ancora: chi protegge il gruppo riceve di più, chi ne abusa tradisce.
  \item La \textbf{piramide della fiducia}: integrità (aiutare davvero, onestà, ammettere gli errori), affidabilità
        (mantenere la parola, saper rifiutare), leadership (dare l'esempio, offrirsi anche quando si rischia, non lasciare
        indietro nessuno).
  \item Il gruppo deve prendersi cura del gruppo: \textbf{raccogliere il piattino}, smettere di dire ``ci penserà qualcun
        altro'', segnalare e chiedere prima di intervenire.
  \item Da manager: delegare i meriti e prendersi le colpe, lasciare la comfort zone, chiedere ``tu come faresti?'',
        celebrare successi e \textbf{fallimenti}, dare il merito a chi lo ha. Essere leader è un \textbf{impegno a lungo
        termine}.
  \item L'obiettivo: tornare a casa \textbf{ispirati, sicuri e realizzati}, come diritto di tutti e non privilegio di pochi.
\end{itemize}

%    % =====================================================================
%  Lezione zero 3 — Agile e Scrum nella pratica
%  Fonte: trascrizione della terza "lezione zero" (GD0_3): una lezione
%  introdotta da Diego Ricchiuti e tenuta da due ospiti, Simone Cicchetti
%  (ex studente, oggi level designer e producer) e Cristian Liberatore
%  (studente del terzo anno, game designer e project manager), che
%  raccontano come hanno applicato davvero Agile e Scrum nei loro progetti.
%  Pacchetti richiesti nel preambolo del libro:
%    \usepackage[italian]{babel}, \usepackage[T1]{fontenc},
%    \usepackage[utf8]{inputenc} (solo pdfLaTeX), \usepackage{booktabs}
%
%  Collocazione nel libro: parte "Lezioni zero", subito dopo
%  GD0_2_tempo_idee_e_leadership, con la stessa numerazione propria (0.3).
%  Nel file principale del libro, ad esempio:
%
%    \part{Lezioni zero: il game designer nell'industria}
%    \renewcommand{\thechapter}{0.\arabic{chapter}}
%    \setcounter{chapter}{0}
%    % =====================================================================
%  Lezione zero 1 — Lavorare in squadra: la cultura del feedback
%  Fonte: trascrizione della prima "lezione zero" (GD0), prima giornata di
%  un workshop di due giorni tenuto prima dell'inizio del corso di Game
%  Design, dedicato alla figura professionale e al lavoro in team.
%  Pacchetti richiesti nel preambolo del libro:
%    \usepackage[italian]{babel}, \usepackage[T1]{fontenc},
%    \usepackage[utf8]{inputenc} (solo pdfLaTeX), \usepackage{booktabs}
%
%  Collocazione consigliata nel libro: le lezioni zero precedono il corso,
%  quindi vanno in una parte separata PRIMA di GD1, con numerazione propria
%  (0.1, 0.2, ...) così che i capitoli del corso conservino i numeri 1, 2, 3.
%  Nel file principale del libro, ad esempio:
%
%    \part{Lezioni zero: il game designer nell'industria}
%    \renewcommand{\thechapter}{0.\arabic{chapter}}
%    \setcounter{chapter}{0}
%    % =====================================================================
%  Lezione zero 1 — Lavorare in squadra: la cultura del feedback
%  Fonte: trascrizione della prima "lezione zero" (GD0), prima giornata di
%  un workshop di due giorni tenuto prima dell'inizio del corso di Game
%  Design, dedicato alla figura professionale e al lavoro in team.
%  Pacchetti richiesti nel preambolo del libro:
%    \usepackage[italian]{babel}, \usepackage[T1]{fontenc},
%    \usepackage[utf8]{inputenc} (solo pdfLaTeX), \usepackage{booktabs}
%
%  Collocazione consigliata nel libro: le lezioni zero precedono il corso,
%  quindi vanno in una parte separata PRIMA di GD1, con numerazione propria
%  (0.1, 0.2, ...) così che i capitoli del corso conservino i numeri 1, 2, 3.
%  Nel file principale del libro, ad esempio:
%
%    \part{Lezioni zero: il game designer nell'industria}
%    \renewcommand{\thechapter}{0.\arabic{chapter}}
%    \setcounter{chapter}{0}
%    \input{Seminari/GD0_1_cultura_del_feedback}
%    % ... altre lezioni zero ...
%    \part{Il corso di Game Design}
%    \renewcommand{\thechapter}{\arabic{chapter}}
%    \setcounter{chapter}{0}
%    \input{GD1/GD1_introduzione_al_game_design}
%    % ...
%
%  (Con hyperref aggiungere anche \renewcommand{\theHchapter}{z.\arabic{chapter}}
%  nella parte zero, per evitare ancore duplicate.)
%  Tutte le etichette di questo capitolo hanno il prefisso "z1:" per non
%  collidere con quelle dei capitoli del corso.
% =====================================================================

\chapter[La cultura del feedback]{Lavorare in squadra: la cultura del feedback}
\label{cap:z1-feedback}

Questo capitolo apre le \emph{lezioni zero}: incontri che precedono l'inizio
vero e proprio del corso di Game Design e che non riguardano una tecnica di
design in senso stretto, ma \textbf{la figura professionale} di chi lavora
nell'industria dei videogiochi. Prima di imparare a progettare un gioco
conviene infatti capire in che tipo di ambiente lo si progetterà: un ambiente
fatto di team, di gerarchie, di scadenze e, soprattutto, di persone che devono
riuscire a lavorare insieme.

La lezione è la prima giornata di un \textbf{workshop di due giorni} (circa
dieci ore in totale, pause pranzo comprese) tenuto da Stefano
Cozzi\footnote{Il cognome è riportato come compare nella trascrizione
automatica, che in più punti storpia i nomi propri; andrebbe verificato.
Nella stessa trascrizione lo studio per cui lavora appare come
``Stormine'', ``Stormline'' e ``Stormind'': si tratta di Stormind Games,
studio con sede a Catania, coerentemente con quanto detto in lezione.}, un
professionista con circa quindici anni di esperienza nel settore, tra
Blizzard Entertainment, Riot Games e, oggi, Stormind Games. Il workshop è
nato da un confronto con i coordinatori dell'accademia\footnote{Nella
trascrizione l'accademia è resa in modi diversi (``Hive'', ``a dive'',
``D.I.V.''); i coordinatori citati sono Enrico e Andrea, insieme al
direttore Diego Ricchiuti. Nel resto del capitolo si usa semplicemente
``l'accademia'', come negli altri capitoli.}: serviva un momento dedicato a
\emph{come ci si comporta in un ambiente professionale}, cosa che riguarda sia
il lavoro futuro sia i progetti di gruppo svolti durante il percorso di
studi, dove ci si deve già abituare a comportarsi ``come se si fosse in
un'azienda''.

La prima giornata ha un unico grande tema, sezionato in più parti:
\textbf{la cultura del feedback}. La seconda giornata (che sarà oggetto della
prossima lezione zero) sarà più pratica: \emph{time management},
\emph{brainstorming} e il ``guidare da dentro'' un team anche senza esserne
il leader (sezione~\ref{sec:z1-domani}).

Il percorso del capitolo segue quello della lezione:
\begin{enumerate}
  \item il formato del workshop e il percorso professionale del relatore
        (sezione~\ref{sec:z1-relatore});
  \item perché le aziende tengono tanto al lavoro di squadra e al
        \emph{fit} culturale (sezione~\ref{sec:z1-squadra});
  \item che cos'è il feedback e da che cosa dipende il modo di darlo
        (sezione~\ref{sec:z1-cose-feedback});
  \item le cinque fasi di vita di un team e il valore del
        \emph{post-mortem} (sezioni~\ref{sec:z1-fasi-team}
        e~\ref{sec:z1-qa-fasi});
  \item il confronto con i colleghi e alcune domande sul lavoro nei grandi
        studi (sezioni~\ref{sec:z1-rivalita} e~\ref{sec:z1-industria});
  \item il circolo virtuoso tra feedback, miglioramento e fiducia, e i suoi
        effetti sulla \emph{retention} e sulla cultura aziendale
        (sezioni~\ref{sec:z1-triangolo} e~\ref{sec:z1-cultura});
  \item il ``gioco del traditore'' e il principio del
        \textbf{default to trust} (sezioni~\ref{sec:z1-among-us}
        e~\ref{sec:z1-default-trust});
  \item le traiettorie di carriera e la valutazione delle performance
        (sezione~\ref{sec:z1-traiettorie});
  \item come si dà, come si riceve e quando si dà un feedback
        (sezioni~\ref{sec:z1-dare}, \ref{sec:z1-ricevere}
        e~\ref{sec:z1-quando});
  \item il feedback come regalo, le domande finali e le anticipazioni
        della seconda giornata (sezioni~\ref{sec:z1-regalo},
        \ref{sec:z1-qa-finale} e~\ref{sec:z1-domani}).
\end{enumerate}

Una premessa di metodo: il relatore ha alternato continuamente teoria,
domande alla classe e racconti della propria esperienza. I racconti non sono
digressioni: sono gli esempi concreti con cui i principi vengono resi
tangibili, e per questo sono riportati con cura.

% ---------------------------------------------------------------------
\section{Il workshop e il relatore}
\label{sec:z1-relatore}

\subsection{Formato e regole di partecipazione}

Il relatore ha cercato di ``riassumere quindici anni di esperienza in dieci
ore''. Ha materiale per circa il doppio delle ore disponibili: le cose
strettamente necessarie vengono trattate per prime, il resto è tenuto ``in
riserva'' come \emph{nice to have}. Le regole del workshop:
\begin{itemize}
  \item le domande si fanno \textbf{alla fine di ogni slide}, così da chiudere
        prima il discorso; non c'è un'unica sessione finale di domande e
        risposte (anche se ci sarà), perché il relatore preferisce fermarsi
        spesso e non fare ``uno spiegone di dieci ore'';
  \item le domande non devono per forza riguardare la slide appena vista: se
        una slide dà lo spunto per altro, se ne può parlare; al massimo sarà
        il relatore a rispondere in breve o a rimandare a una slide
        successiva;
  \item si può chiedere anche di esperienze passate, opinioni e sensazioni;
        se per NDA o per serietà professionale non può rispondere, lo dirà
        esplicitamente (``non lasciate che sia un \emph{blocker}'');
  \item su Zoom si usa l'alzata di mano o la chat, che il relatore tiene
        aperta sul secondo monitor;
  \item non ci sono pause ufficiali: la lezione è registrata e ognuno è
        libero di assentarsi quando serve, cercando però di non disturbare il
        flusso della lezione;
  \item le slide verranno fornite integralmente tramite la segreteria
        dell'accademia (Vanessa), insieme alle registrazioni: le slide senza
        il parlato servono poco, ma sono un buon promemoria.
\end{itemize}

\subsection{Il percorso professionale}

Il relatore si presenta (scherzosamente) come la ``guida escursionistica''
che accompagna la classe nella ``foresta incantata'' della cultura del
feedback, perché per molti è un territorio sconosciuto. Il suo percorso:
\begin{itemize}
  \item in Italia ha lavorato per aziende che fornivano servizi a terzi, tra
        cui il lavoro su \emph{Warhammer Online}, e per un canale televisivo
        sul digitale terrestre dedicato ai videogiochi\footnote{Nella
        trascrizione: ``Digital Game Network Channel''. Il nome esatto del
        canale andrebbe verificato.};
  \item a \textbf{Blizzard Entertainment} (circa sei anni, a Parigi) ha fatto
        il \emph{community manager} per l'Italia (\emph{StarCraft~II},
        \emph{World of Warcraft} da \emph{Mists of Pandaria} in poi,
        \emph{Hearthstone}, \emph{Diablo~III}, le espansioni di
        \emph{StarCraft}, il lancio di Battle.net), per poi diventare il
        manager dei community manager europei;
  \item a \textbf{Riot Games} (circa cinque anni, a Dublino) è stato
        brevemente community manager italiano e poi \emph{publishing
        manager} di un team di circa trenta persone per l'Europa. Il
        \emph{publishing} non è solo community: è l'unione di
        \textbf{community, marketing e PR};
  \item a \textbf{PG Esports} si è occupato di eventi dal vivo e online
        (ad esempio la Milano Games Week e tornei di \emph{Valorant} e
        \emph{League of Legends}), restando in contatto con Riot e Blizzard;
  \item da circa un anno è tornato nello sviluppo, a \textbf{Stormind
        Games}, dove oggi si occupa di risorse umane e organizzazione del
        lavoro; lo studio verrà presentato più ampiamente al \emph{Career
        Day} dell'accademia insieme alla sua HR manager;
  \item fa inoltre consulenze aziendali anche al di fuori dei videogiochi e
        insegna in master universitari (IULM di Milano).
\end{itemize}

Il suo ruolo ricorrente, sia a Blizzard sia a Riot, è stato quello di
\emph{Mr.~Wolf} (``per chi ha visto \emph{Pulp Fiction}''): quello che
risolve i problemi. Veniva inserito nei team più in difficoltà, ad esempio
dal punto di vista delle performance, per sistemarli e passare al successivo.
Era presente in alcuni dei momenti più critici, come il lancio di
\emph{Diablo~III} o il lancio di espansioni di \emph{World of
Warcraft}\footnote{La trascrizione riporta ``il lancio con i Gary's'' e
``World of Trenor'': si tratta verosimilmente di \emph{Mists of Pandaria} e
di \emph{Warlords of Draenor}.}, a fare \emph{damage control}.

La scoperta fondamentale di questi anni è che i problemi da risolvere
\textbf{quasi mai riguardavano il prodotto} (non era certo lui a sistemare i
server di \emph{Diablo~III}), ma i team che c'erano dietro: persone che
litigavano, non collaboravano, non riuscivano a stare nella stessa stanza
senza urlarsi contro. E nel 90\% dei casi tutto questo si riduce a
\textbf{problemi di comunicazione}: è raro che perfino le incompatibilità
caratteriali non siano risolvibili parlando.

\subsection{Struttura non vuol dire burocrazia}

A Stormind il relatore sta accompagnando un team passato da 30 a 95 persone
in meno di un anno, attraverso tutti i cambiamenti di cui parla il workshop.
Darsi una struttura viene spesso percepito male, come ``più processi, più
procedure'', ma bisogna distinguere:
\begin{itemize}
  \item un processo come ``timbrare il cartellino'' può essere davvero
        invasivo e noioso;
  \item definire \emph{come funziona un documento di game design, chi lo
        approva, chi lo visiona e chi no} è invece una struttura necessaria.
\end{itemize}
Uno studio di 30 persone nella stessa sede può farne a meno: quando c'è un
dubbio ci si alza e si va alla scrivania del collega. Uno studio di 95
persone sparse per l'Italia (e in parte per l'Europa) non può più farlo:
servono una struttura diversa, un modo di lavorare diverso e una
pianificazione del futuro diversa.

% ---------------------------------------------------------------------
\section{Perché il lavoro è sempre lavoro di squadra}
\label{sec:z1-squadra}

\subsection{Non si va in ufficio per fare amicizia}

Il primo messaggio è netto: \textbf{non si va in ufficio per fare amici} o
per trovare compagni di vita. Si va in ufficio perché c'è un lavoro da fare, e
lo scopo è riuscire a farlo \emph{anche con chi sta antipatico}. Con le
persone simpatiche è facile: diventeranno amici, e non sarà mai un problema.
È con chi sta antipatico che le cose si complicano, ed è lì che emergono i
\textbf{veri professionisti}: quelli che riescono a lavorare con chiunque
senza che l'antipatia pesi sul modo in cui lavorano e collaborano, mettendo il
prodotto e il team davanti al proprio ego e alle proprie preferenze personali.

La frase che fa da filo conduttore a tutto il workshop (e che il relatore
promette di riproporre ``fino alla nausea'') è:

\begin{quote}
  \emph{La capacità di un gruppo di persone di fare cose straordinarie
  dipende dal modo in cui queste persone operano insieme.}
\end{quote}

Le aziende raramente pensano che l'innesto di una singola persona possa
cambiare il destino di un progetto. Questo è tanto più vero quanto più grande
è il team, ma vale anche per i piccoli studi, anche per un gruppo di cinque
persone. Come negli sport di squadra: l'individuo conta fino a un certo punto,
e anche se si ha Dončić in squadra, se si gioca uno contro cinque si perde
comunque\footnote{Il riferimento è alle finali NBA in corso nei giorni del
workshop tra Dallas Mavericks e Boston Celtics (giugno 2024), il che
permette di datare approssimativamente la lezione.}.

\subsection{Lo \emph{useless genius} e il costo delle persone tossiche}

Molti dipartimenti HR usano la definizione di \textbf{useless
genius}\footnote{Il concetto è più noto in letteratura manageriale come
\emph{brilliant jerk}; in lezione è usata l'espressione \emph{useless
genius}.}: si può essere bravissimi, a volte i migliori nel proprio lavoro,
ma se non ci si sa comportare, relazionare e operare in un gruppo si è
fondamentalmente inutili ai fini del gruppo. Nel migliore dei casi inutili,
nel peggiore \emph{dannosi}.

In passato le aziende tendevano a ``chiudere un occhio'': la persona è
fondamentale, la teniamo. Oggi la domanda è diventata: \emph{quanto costa
all'azienda, oltre al suo stipendio, tenere una persona che crea un ambiente
tossico e peggiora chi le sta intorno?} Ci sono costi sommersi che le aziende
cominciano a vedere, e la tendenza è liberarsi di queste persone il prima
possibile. Il ragionamento è semplice:
\begin{itemize}
  \item le \textbf{competenze si possono trasmettere} o sostituire: si trova
        qualcuno che fa lo stesso lavoro magari nel doppio del tempo, oppure
        più persone che se lo dividono;
  \item il \textbf{carattere delle persone} invece cambia molto raramente.
\end{itemize}

\subsection{Il fit culturale}

Quando qualcuno \textbf{non è un \emph{fit} culturale} per un'azienda è
meglio separarsi presto. Ma non essere adatti a \emph{quell'}azienda non
significa non essere adatti a lavorare: è perfettamente legittimo che una
persona inadatta a un'azienda sia adattissima a un'altra.

Il relatore l'ha vissuto in prima persona. Dopo sei anni a Blizzard viene
contattato da Riot e affronta un processo di selezione lunghissimo (``quasi
undici colloqui''), al termine del quale viene scartato. Scrive allora una
mail che gli ha ``cambiato la traiettoria di carriera'': chiede
semplicemente perché non è stato ritenuto adatto. La risposta: in nessun
colloquio si era esposto abbastanza da far capire \emph{perché} volesse
lasciare Blizzard, e quindi pensavano che in realtà stesse bene dov'era. Il
motivo era che, per educazione professionale, non gli piace parlare male del
proprio posto di lavoro; portata all'eccesso, però, questa riservatezza era
diventata un ``è tutto bellissimo, ma vorrei provare Riot'', senza un perché.
Lo invitano comunque a un incontro dal vivo (``dal vivo è diverso''), e alla
fine viene assunto.

Una volta arrivato, i sei mesi di prova (\emph{probation}) sono serviti tutti
per smettere di ragionare e lavorare ``come a Blizzard''. La differenza
culturale era enorme, tanto che in Riot esisteva un gruppo di supporto
interno tra ex dipendenti Blizzard: un collega gliel'ha descritto come
``uscire da una relazione tossica e doversi riabituare a una relazione
normale''\footnote{Il relatore precisa che parla della Blizzard di allora,
che non esiste più: oggi è parte di Microsoft e nessuno dei suoi ex colleghi
vi lavora ancora, anche a causa dei licenziamenti di massa.}.
\begin{description}
  \item[Blizzard] era \emph{developer first}: il focus era chiaramente sugli
        sviluppatori, prima di tutti gli altri. Dopo sei anni c'erano ancora
        meeting a cui non poteva partecipare perché non abbastanza
        \emph{senior}, informazioni che non poteva sapere, piani che non
        poteva vedere.
  \item[Riot] era \emph{player first}: prima i giocatori, e i dipendenti come
        mezzo per trasferire il feedback dei giocatori dentro il gioco. Il
        secondo giorno lo portano in una stanza e gli mostrano il piano dei
        cinque anni successivi. Lui, sorpreso, pensa a un errore (``sono qui
        da oggi, non vorrei mettervi nei guai''); la risposta è: ``sei un
        dipendente Riot o no? Allora non ci sono livelli di accesso: devi
        sapere queste cose per fare meglio il tuo lavoro''.
\end{description}

Da questa esperienza ha imparato molto sulla \textbf{trasparenza}: condividere
le informazioni aiuta il lavoro. Non sempre, però: a volte non si condivide
qualcosa non per segretezza, ma perché non è ancora pronta per essere
presentata, esattamente come non si mostra un gioco quando è ancora solo
teoria (``te lo faccio vedere quando ho qualcosa di tangibile'').

La lezione per gli studenti è doppia: le aziende valuteranno il \emph{fit}
culturale sempre di più, ma capire se si è un \emph{fit} richiede anche
\textbf{lavoro da parte del candidato}. Bisogna avere la presenza di spirito
di dire: ``conosco i valori di questa azienda, non mando nemmeno il CV,
perché so che durerei un mese e mezzo''.

% ---------------------------------------------------------------------
\section{Che cos'è il feedback}
\label{sec:z1-cose-feedback}

\subsection{Le prime risposte della classe}

La prima domanda alla classe è: \emph{quando sentite la parola
``feedback'', a cosa pensate?} Le risposte raccolte, a voce e in chat:
\begin{itemize}
  \item un \textbf{riscontro} su ciò che si fa e su come lo si fa, sul lavoro
        o sul modo in cui ci si interfaccia con il gruppo; può essere
        lavorativo oppure comunicativo e comportamentale. Chi ha risposto lo
        vede come qualcosa che, se l'azienda non dà, va attivamente
        richiesto, perché è importante sapere come sta andando per
        migliorarsi;
  \item né positivo né negativo, ma \textbf{costruttivo}: in entrambi i casi
        serve a migliorarsi;
  \item uno \emph{scambio di opinioni}, un \emph{riscontro produttivo};
  \item un feedback ``di reazione'', non verbale: si fa una cosa e si
        osserva come cambia l'ambiente o il comportamento delle persone
        intorno (anche perché a voce la gente spesso si tiene le cose
        negative);
  \item per un designer, anche il \textbf{linguaggio del corpo}: osservando
        un giocatore giocare si può dedurre l'emozione che prova (frustrazione,
        incomprensione).
\end{itemize}

Sulla forma (scritto o orale?) le opinioni divergono: c'è chi lo preferisce
\emph{mille volte} scritto; chi pensa che dipenda dalle dimensioni
dell'azienda (orale nelle piccole, scritto in quelle con team di centinaia di
persone); chi pensa che dipenda dall'importanza del feedback (scritto se
tecnico o pieno di dati, da tenere sotto gli occhi; orale se è un consiglio
veloce).

Il relatore usa due esempi per mostrare che \emph{dipende}:
\begin{itemize}
  \item un feedback sulla programmazione ha senso darlo scritto, magari
        commentando direttamente il codice, più che ``andando a prendere un
        caffè'';
  \item ma se il feedback è ``ieri, durante il meeting, ho percepito male il
        fatto che mi interrompessi di continuo'', ce lo si aspetta
        \textbf{orale}.
\end{itemize}

\subsection{Da che cosa dipende il modo di dare feedback}

Lo \emph{spoiler} della giornata è che \textbf{non esiste una risposta univoca}
su come dare feedback. Il workshop non insegna ``il'' modo, ma fornisce
\emph{tutte le opzioni}, perché il modo giusto cambia in base a molti fattori:
\begin{itemize}
  \item a chi lo si dà: un collega, un superiore, il proprio capo, un proprio
        \emph{report} (un sottoposto);
  \item se è personale o professionale;
  \item se è la prima volta o la terza volta che si dà lo stesso feedback;
  \item se riguarda un progetto importante o secondario.
\end{itemize}

Manca però l'elemento fondamentale, che la classe individua subito:
\textbf{l'altro}. A parità di tutti gli altri fattori, il \emph{carattere
della persona che riceve il feedback} dovrebbe far cambiare il modo in cui lo
si dà. Per questo, prima della lezione, agli studenti era stato chiesto di
fare un test di personalità: non un test di psicologia avanzata, ma uno
strumento per rendersi conto di quanta varietà ci sia anche solo in una
classe di 53 persone (il che non vuol dire trovare 52 modi diversi di dare
feedback, ma lavorarci un po' sì).

Il relatore porta se stesso come esempio. Preferisce:
\begin{itemize}
  \item i feedback \textbf{positivi scritti}, perché è a disagio a ricevere
        complimenti in pubblico;
  \item i feedback sul proprio lavoro \textbf{diretti}: ``vai subito al
        punto e non girarmi intorno''.
\end{itemize}
È molto orgoglioso del suo lavoro e, quando qualcuno lo critica, la sua
primissima reazione istintiva è ``come ti permetti''. Negli anni ha imparato
a farla durare un secondo, ma chi gli dà feedback di persona dice di
vederla ancora negli occhi. Per mitigarla, dice esplicitamente a chi gli dà
feedback: \emph{``ignora il modo in cui reagisco nei primi cinque secondi,
vai avanti, poi mi passa''}. È considerato da tutti molto aperto al feedback
e non ha problemi a dire ``hai ragione, facciamo diversamente'': ma ci sono
reazioni istintive e caratteriali che è difficile controllare al cento per
cento.

\subsection{Con chi si lavora meglio}

Guardando le persone con cui si lavora abitualmente (anche nei progetti
dell'accademia), si noterà che si lavora meglio con due profili:
\begin{enumerate}
  \item chi è \textbf{esattamente come noi}, ma qui è un 50 e 50: alcuni
        caratteri vanno d'accordo con i propri simili, altri non vogliono
        nessuno con lo stesso carattere nel gruppo (e si potrebbe rientrare
        in questa categoria);
  \item chi è \textbf{affine}: ha qualcosa in comune con noi, ma non tutto,
        e proprio per questo può dare punti di vista differenti.
\end{enumerate}
Il lavoro della giornata è capire come valorizzare il secondo caso: ottenere
feedback da persone con caratteri, visioni e background diversi, per
completare una visione del proprio lavoro che è necessariamente limitata. È
come esplorare una mappa con la nebbia di guerra: un'altra persona, al di là
del muro, vede cose che noi non vediamo.

\subsection{Il feedback come somma di piccoli miglioramenti}

Per il relatore il feedback è questo: \textbf{applicare tanti piccoli
miglioramenti periodici}, che si sommano nel tempo. Migliorare dello 0,1\%
ogni giorno porta a fine anno a un miglioramento enorme, e questo vale per i
prodotti, per i giochi, per le persone e per i team.

\paragraph{L'esempio del ciclismo britannico.} La squadra di ciclismo
britannica, fino a una quindicina d'anni fa, non aveva vinto quasi nulla.
Cambiato il team manager, tutti si aspettavano grandi cambiamenti sportivi
(convocazioni, strategie di gara). Arrivarono invece \emph{minuscoli
cambiamenti non sportivi}: un altro chef in trasferta, materassi migliori nei
camper per riposare meglio, un sellino più comodo a fine tappa. Cambiamenti
marginali che da soli non avrebbero cambiato niente; eppure, dopo due o tre
anni di questo lavoro quotidiano, la squadra ha cominciato a vincere a
ripetizione, Olimpiadi comprese, cambiando il modo in cui il ciclismo è
visto in Gran Bretagna\footnote{È la nota strategia delle \emph{marginal
gains} associata a Dave Brailsford (British Cycling e Team Sky); il nome non
è citato in lezione.}. Certo, si è investito anche sulle persone giuste, ma
tutto è partito dai piccoli miglioramenti periodici.

Per gli studenti vale lo stesso: chiarire un conflitto con una persona con
cui si ha ancora poco a che fare, o una situazione che ha lasciato
``l'amaro in bocca'', sembra inutile oggi, ma darà i suoi frutti tra tre, sei,
nove mesi. Sono \emph{investimenti sul futuro}.

\paragraph{Peggiorare ogni giorno.} Esiste anche il percorso inverso:
lasciar correre, ``buttare sotto il tappeto'' le cose perché ``tanto è
inutile parlarne''. Queste cose tornano sempre fuori dopo otto mesi, come il
\emph{libro dei rancori} dei nani di Warhammer: ``ti ricordi quando, un anno
e mezzo fa, all'inizio del progetto, mi hai detto quella cosa? Mi sono
offeso''. A quel punto è quasi impossibile risolvere, perché nessuno ricorda
più con precisione cosa sia successo (``ma quale cena?'', ``forse era in
pizzeria'', ``no, eravamo da un'altra parte''). \textbf{Tenersi le cose nel
cassetto per tirarle fuori quando servono è sempre sconsigliato}: non servono
ed è impossibile usarle in modo costruttivo.

% ---------------------------------------------------------------------
\section{Le cinque fasi di vita di un team}
\label{sec:z1-fasi-team}

\subsection{Perché conoscerle}

Quasi tutti gli studenti lavorano già in gruppo (al primo anno su un gioco da
tavolo, negli anni successivi su un videogioco) e si saranno accorti che un
team attraversa delle fasi: il modo di lavorare insieme cambia nel tempo.
Spesso ci si fa prendere di sorpresa: ``eravamo così affiatati, cosa è
successo?''.

Un esempio da una consulenza fuori dai videogiochi: un'azienda farmaceutica
gli affida un team commerciale (rappresentanti che presentano medicinali a
ospedali, medici e privati). Il problema: ``prima andava tutto bene; ora che
siamo in diciotto dovremmo lavorare meno e meglio, e invece va peggio''. Va
peggio perché le persone nuove sono appena arrivate, mentre il management si
aspettava già di vedere i risultati dell'investimento.

Fino a pochi anni fa le fasi di un team venivano disegnate come una
\textbf{linea dritta verso l'alto}: le fasi erano le stesse, ma viste come un
miglioramento costante. Non si vedevano le parti \emph{fisiologiche} del
lavoro in team, cioè anche i peggioramenti. Accettarle è il primo passo per
accettare i fallimenti: nessun percorso personale o di team è lineare.

Il relatore ammette che il suo modo di lavorare è nato per il 70\% dal
proprio impegno a imparare, fare bene e rispettare le persone, e per il 30\%
dall'aver lavorato con ``incapaci colossali'' che gli hanno insegnato come
\emph{non} fare le cose. A Blizzard un manager gli fece capire che, pur non
sapendo ancora che tipo di \emph{people manager} sarebbe diventato, sapeva
benissimo quale tipo di manager \emph{non} voleva essere.

\subsection{Forming, storming, norming, performing, adjourning}

Le fasi sono cinque, un po' come le fasi del lutto. Molte aziende ne
considerano solo quattro (``e sono quelle che solitamente chiudono''); le
aziende al passo con i tempi ne considerano cinque\footnote{Il modello è
quello di Bruce Tuckman (1965), a cui la quinta fase, \emph{adjourning}, fu
aggiunta nel 1977. Il nome di Tuckman non è citato in lezione; la
trascrizione riporta la quinta fase come ``journing''.}.

\begin{description}
  \item[Forming] Si viene messi insieme e si comincia a lavorare. Non si sa
        nulla degli altri; spesso non si conoscono nemmeno i nomi, ed è
        normale.
  \item[Storming] Si cominciano a ``dare le testate come le capre in
        montagna per determinare il territorio'': questa cosa la faccio io,
        questa la vorrei fare così, chi decide? Si stabilisce chi fa cosa e
        come si divide il lavoro. Nel grafico è una \textbf{fase di
        peggioramento}: che il team litighi o discuta civilmente non cambia
        nulla, il team deve investire parte di tempo ed energie per risolvere
        queste questioni. È normale e \textbf{non è evitabile}.
  \item[Norming] Si comincia a capire come lavorano gli altri, cosa sanno
        fare e come far funzionare il gruppo; iniziano i miglioramenti. È una
        forma di \emph{accettazione} degli altri, che deriva dall'aver capito
        cosa devo fare io e cosa devono fare loro.
  \item[Performing] C'è uno sforzo \textbf{proattivo}: ``questa cosa
        toccherebbe a me, ma so che Alessio la fa molto meglio, quindi gli
        chiedo aiuto''. Si sfruttano i punti di forza degli altri, perché li
        si conosce e perché si sa che passare il testimone non avrà effetti
        negativi (il capo non dirà ``perché non l'hai fatta tu?'').
  \item[Adjourning] Il progetto finisce, si chiude o viene riorganizzato, il
        team viene sciolto. Quello che andrebbe fatto è analizzare
        \textbf{cosa ha funzionato, cosa no e cosa migliorare} per il
        progetto successivo.
\end{description}

\subsection{La fase che tutti saltano}

``Spoiler incredibile'': l'\emph{adjourning} è la fase che molte aziende
saltano, \textbf{perché c'è sempre il progetto dopo}.

A PG Esports, per la Milano Games Week (come per qualsiasi evento dal vivo a
cadenza annuale: Lucca Comics, Etna Comics, Comicon\dots), finita l'edizione
di un anno si cominciava dal giorno dopo a lavorare a quella successiva,
perché un anno sembra tanto ma non basta mai. Nei videogiochi è lo stesso:
finito un gioco, a parte le patch di mantenimento e del \emph{day one}, si
passa subito al successivo.

Il cambiamento introdotto dal relatore (in parte già praticato, ma da lui
strutturato) è stato: finita la Games Week, \textbf{tutti in ferie per una
settimana}; poi ci si ritrova per definire cosa ha funzionato e cosa no, con
un \textbf{post-mortem} (o \emph{debrief}); solo dopo si inizia a lavorare
all'edizione successiva. Senza questa fase manca la parte di
\emph{learning}, cioè di ciò che si è imparato dal progetto.

Il consiglio vale per gli studenti: quando si finisce il gioco da tavolo, e
poi il gioco digitale, bisogna trovare il tempo per il post-mortem,
altrimenti si butta via \textbf{il 90\% dell'utilità del progetto}. Una
struttura semplice per farlo è descritta nella sezione~\ref{sec:z1-qa-finale}.

% ---------------------------------------------------------------------
\section{Domande sulle fasi del team}
\label{sec:z1-qa-fasi}

\subsection{Lo storming che ritorna}

\paragraph{Capita che lo storming si ritagli spazio anche nelle altre fasi?}
Lo \emph{storming} è una fase precisa, che però molti manager, coordinatori e
team cercano di \textbf{saltare}: ``se io sono bravo come manager e loro sono
bravi come team, passeremo da forming a norming senza problemi''. È un errore,
perché senza una fase in cui si discute davvero non si sviluppa una
\emph{skill} fondamentale: \textbf{saper discutere}. Tutti conoscono persone
non abituate a farlo, per cui ogni minimo disaccordo diventa un affronto
personale (``avevamo deciso la pizza, perché andiamo al sushi? Perché mi fate
questo affronto?''): sono persone che non sono abituate ad avere conflitti
costruttivi.

Se si pensa di essere in \emph{norming} e continuano a emergere conflitti, il
punto è che \textbf{non si era davvero in norming}: lo storming non ``si
trascina'' nelle fasi successive, semplicemente non se n'è mai usciti.
Allo stesso modo, a chi chiede se ogni mini-fase di storming durante la vita
del gruppo sia frutto di uno storming fatto male o solo di carattere, la
risposta è già contenuta qui.

Un partecipante osserva che spesso, ancor prima di finire un progetto, si è
già messi su quello successivo: è il modo migliore per essere in costante
\emph{damage control}. Era l'esperienza del relatore a Blizzard: veniva
spostato fisicamente da un team all'altro, e finito il lancio di
\emph{Diablo~III} era già su \emph{World of Warcraft}, senza vivere il
``dopo''. Si è dovuto ritagliare da solo lo spazio per un debrief personale,
ed è lì che ha deciso che tipo di manager non voleva essere. Il consiglio:
\textbf{se non lo fa il team, fatelo voi}.

\subsection{Dove si passa la maggior parte del tempo}

\emph{In quale fase si passa più tempo nella propria carriera?} Gli studenti
rispondono \emph{storming} e \emph{norming}. Il relatore, in quindici anni,
calcola di aver passato in \emph{performing} circa \textbf{dieci mesi}. Non è
pessimismo: arrivare a performing deve essere sempre l'obiettivo, sia da
manager sia da membro del team; ma è difficile arrivarci, ed è normale che non
sempre ci si arrivi.

Lo storming può durare a lungo, ma con team molto esperti dura pochissimo,
perché c'è chi lo sa gestire. In realtà \textbf{la maggior parte del tempo si
passa in forming}, perché ogni equilibrio raggiunto vale per \emph{quelle}
persone, \emph{quel} progetto, \emph{quelle} tempistiche. Basta che cambi un
qualsiasi fattore (una deadline spostata, una persona nuova, qualcuno che si
dimette, un progetto cancellato) e si ricomincia, perché i bilanciamenti
precedenti non valgono più.

Esempio: un team di cinque persone affiatatissimo, in pieno performing. La
dirigenza, entusiasta, ``raddoppia le risorse'' e lo porta a dieci. O si
lascia tutto com'è (e i cinque nuovi non fanno niente), o bisogna ridiscutere
chi fa cosa e come i nuovi possono aiutare: che piaccia o no, si ricomincia
da forming.

Quindi la bravura di un manager, ma anche di un membro del team,
\textbf{non sta nel passare tanto tempo in performing, ma nell'accelerare il
passaggio da una fase all'altra}, e nel fare bene la fase di
\emph{adjourning}.

\subsection{Analizzare anche i successi}

C'è la tendenza a fare post-mortem solo dei progetti andati male: ``questa
cosa non ha funzionato, vediamo perché''. Va benissimo. Ma quando un progetto
va bene si dà per scontato che sia andato bene ``perché era buono''. Bisogna
invece chiedersi \emph{perché} è andato bene: era il prodotto? La
tempistica? Il modo in cui abbiamo fatto una cosa piuttosto che un'altra?
\textbf{Se non si capisce davvero il motivo del proprio successo, si
tenderanno a replicare le cose sbagliate.}

\paragraph{Blockbuster.} Quando arrivò la notizia che Netflix avrebbe
lanciato un servizio di streaming, Blockbuster (la catena di noleggio di
videocassette e DVD) aveva tutto il tempo per preparare un'offerta
concorrente. In una riunione del management raccontata in un libro (di cui il
relatore non ricorda il titolo), la direzione sembrava orientata a investire
nello streaming, finché un dirigente fece notare che il 12\% del fatturato
derivava dalle \emph{multe per riconsegna tardiva}, impossibili con lo
streaming, mentre secondo le stime lo streaming avrebbe prodotto il 10\% del
fatturato. ``Non possiamo rinunciare a questo 2\%'': niente streaming. Oggi
non c'è più un Blockbuster ``manco a cercarlo col lanternino''. Non avevano
capito qual era il vero motivo del loro successo, cioè il loro modello di
business\footnote{Le percentuali sono quelle riferite in lezione e non sono
state verificate.}.

\paragraph{The Founder.} Nel film sulla nascita di McDonald's c'è una scena in
cui il futuro socio del protagonista gli dice: ``tu pensi di essere nella
\emph{food industry}, e invece sei nella \emph{real estate industry}''. Non
dovrebbe investire sugli hamburger, ma sui terreni su cui costruisce i
ristoranti. Anche per un team vale lo stesso: bisogna capire che cosa
l'ha reso davvero di successo, oppure no.

\subsection{L'autoanalisi personale}

\paragraph{E quando si cambia azienda?} Uno studente chiede se il debrief sia
possibile anche quando si lascia un'azienda. Non si può tornare in ufficio a
discutere della propria esperienza, ma lo si può fare \textbf{per se stessi}.
Il relatore l'ha fatto, anche se all'epoca in modo poco consapevole, ed è
così che è passato dal \emph{community management} al \emph{people
management}. Analizzando cosa aveva fatto bene a Blizzard si era reso conto
che:
\begin{itemize}
  \item aveva costruito un ottimo rapporto di fiducia con i giocatori (con
        alcuni proprietari di fansite parla ancora oggi);
  \item pur senza la seniority, era diventato un punto di riferimento per il
        team, visto come imparziale ed equilibrato anche su temi che non lo
        riguardavano direttamente;
  \item sapeva \textbf{risolvere problemi}: durante il lancio di
        \emph{Diablo~III} era il più tranquillo di tutti, non perché
        sminuisse il problema, ma perché sapeva che si sarebbe risolto con i
        suoi tempi; il lavoro non era spegnere l'incendio ma ``far sfollare
        la gente'' e aspettare che tutto tornasse in ordine.
\end{itemize}
E aveva capito cosa sapeva fare meno bene, ad esempio \textbf{gestire il
proprio manager}. Avere un manager non significa che il manager gestisca te:
il manager gestisce il tuo lavoro, ma tu devi saper gestire i tuoi
\emph{stakeholder}. Spesso lo stakeholder è il proprio manager; da manager,
gli stakeholder saranno CEO, CFO e così via. Gestirli significa conoscerli,
sapere come portare loro le informazioni, come far presente che qualcosa non
va. Fa parte del lavoro, anche se a lui non piace, e ci ha lavorato.

Nel mese e mezzo di pausa prima di entrare in Riot si è concentrato su
questo: ``cosa porto come mio punto di forza?''. L'autoanalisi è culminata
nel chiedere espressamente alla persona che l'aveva assunto: ``prima mi avete
detto di no, ora mi avete preso: perché?''.

\paragraph{Come si capisce se un progetto è andato bene?} Dipende
dall'\textbf{obiettivo}. Se l'obiettivo era vendere milioni di copie e il
gioco le ha vendute, bisogna capire perché. Se, più ragionevolmente,
l'obiettivo era un gioco da tavolo che permettesse di mettere in pratica le
conoscenze del corso, va valutato per quello: ha permesso di farlo? Se sì,
perché? Se no, perché no? L'obiettivo va fissato \textbf{all'inizio del
progetto, non alla fine}, proprio per poterlo valutare dopo.

\subsection{Quando il team cambia di continuo}

\paragraph{Tutti questi cambiamenti non sono dannosi per il progetto?}
Assolutamente sì, ma quasi mai sono volontari. I progetti andati male sono
spesso progetti grandi.

Il relatore racconta di \emph{Cyberpunk 2077}. Conosceva persone in CD
Projekt Red (``il mondo è piccolissimo, le cose si sanno''). Alla Gamescom va
con Diego Ricchiuti, direttore dell'accademia, a vedere la presentazione del
primo gameplay, quella ``in cui regalavano il giubbottino da samurai''.
Finito il video, tra gli applausi, si guardano e capiscono entrambi che
\textbf{il gioco non è pronto}, a soli quattro mesi dall'uscita prevista:
Diego aveva notato il combattimento e la guida. Gli amici di CD Projekt Red
(uno aveva lavorato con lui a Blizzard per sei anni) gli confermano tra le
righe che ``ce la faremo, ma ci sono ancora cose da sistemare''. Il gioco fu
poi rimandato più volte\footnote{In lezione il relatore dice ``due, forse
tre volte''; l'uscita di \emph{Cyberpunk 2077} è stata effettivamente
rinviata tre volte prima del dicembre 2020.}.

Per deformazione professionale, chi ha sviluppato l'occhio critico non
riesce più a godersi un trailer: in un trailer di \emph{The Day Before} Diego
gli aveva fatto notare che l'effetto della fiammata di un fucile era stato
chiaramente aggiunto in postproduzione e non faceva parte del gameplay. ``È
il triste epilogo di questa storia'': non riescono più a godersi un trailer.

Il punto è che più il team è grande, più queste variabili vengono imposte, e
più è facile che il fallimento \textbf{non dipenda dal fatto che il prodotto
sia brutto} (\emph{Cyberpunk} è un bellissimo gioco sotto molti punti di
vista, e il relatore lo sta rigiocando ora che è stato corretto ed espanso).
Nei grandi gruppi qualcuno arriva, qualcuno si dimette, cambia il budget,
arriva il rinvio; nel caso di \emph{Anthem}, il team ha fatto così tante volte
forming--storming--forming--storming che ne è uscito ``un prodotto che
dovevamo espellere: l'abbiamo espulso, buttato fuori dal nostro organismo pur
di liberarcene''.

Tra i cambiamenti c'è anche la crescita verticale: se il capo viene promosso
perché lavora benissimo, bisogna scegliere un nuovo capo, ed è sempre un
caos (chi ci resta male, chi pensava di meritarlo, chi voleva il suo amico,
chi non sopporta il prescelto). Per questo il relatore è \textbf{contrario
alle promozioni durante i progetti}: si rischia di rovinare tutto il lavoro di
storming.

\subsection{Obiettivi chiari nei piccoli gruppi}

Nei gruppi piccoli, come quelli degli studenti, il problema deriva più
spesso dal \textbf{non avere chiaro che cosa si vuole ottenere}: un prodotto
puramente educativo e formativo, o un prodotto richiesto dall'accademia che,
se viene bene, si potrebbe anche vendere? Un gioco commerciale o un gioco
artistico? Va chiarito subito, perché ci saranno sempre momenti in cui
scegliere tra A e B, e a guidare la scelta sarà la domanda: \emph{quale delle
due opzioni ci porta più vicino all'obiettivo?} Se l'obiettivo non è chiaro,
non si sa cosa scegliere.

Questo porta molte aziende a fare \textbf{scelte di design che in realtà sono
scelte commerciali}. Le microtransazioni sono scelte commerciali: nessuno ha
mai pensato ``sai cosa piacerebbe tanto ai giocatori? Pagare le bustine''.
Prima le bustine si guadagnavano giocando; ora la tendenza è ``potresti
giocare, ma se paghi è meglio''. Si pensi alle polemiche su Ubisoft per
l'opzione di ``saltare il gioco'' pagando, o ai famosi \emph{boost} di livello
di \emph{World of Warcraft}: ``devo pagarti per saltare il gioco?''.

Una studentessa (Silvia) cita \emph{Final Fantasy XIV}, che a un certo punto fu
``resettato'' perché non funzionava: un esempio positivo. Il relatore aggiunge
quello negativo di \emph{Age of Conan}, che non è mai riuscito a fare quel
reset iniziale (``è una schifezza, ve lo ripresentiamo quando è pronto''). Ci
vuole consapevolezza, ma anche i mezzi: un conto è Square Enix, che può
permetterselo, un altro è un'azienda che non può chiudere tutto per un anno
continuando a pagare gli stipendi.

Riassumendo: nei piccoli gruppi i problemi sono più di
\textbf{comunicazione e chiarezza} (obiettivi non chiari o che si spostano);
nelle grandi aziende di più di \textbf{persone che vanno e vengono, obiettivi
e tempistiche che cambiano}. In accademia le scadenze sono fisse (il gioco da
tavolo si consegna a fine anno scolastico, a giugno, e giugno resta); in
azienda può succedere che una feature prevista per la \emph{vertical slice}
di luglio venga anticipata a maggio per portarla alla Gamescom, e un sistema
perfetto per consegnare a giugno può non esserlo più per consegnare a maggio.

\subsection{L'ambiente determina la qualità del prodotto}

\paragraph{È vero che chi passa per Electronic Arts finisce per fare brutti
giochi?} Uno studente cita BioWare, che faceva ottimi giochi finché era
autonoma. Il problema non è dei developer, che non sono ``improvvisamente
diventati incapaci'': è evidente che il problema è altrove, e che
\textbf{l'ambiente in cui si fa un gioco determina molto spesso la qualità del
prodotto}.

Un esempio locale positivo: Tommaso Bonanni, conosciuto dal relatore insieme a
Diego e già docente dell'accademia, ha fondato a Roma Caracal Games Studio,
applicando fin dall'inizio questi concetti di cultura aziendale. L'azienda è
rimasta piccola, ma ha molto chiaro cosa vuole fare e come. Quando si
predispone fin dall'inizio l'azienda e il team a lavorare in un certo modo è
fattibilissimo; è molto più difficile quando si è EA, si acquisisce
un'azienda e le si dice ``da oggi si fa così''.

Il relatore ha vissuto la transizione da Blizzard ad Activision Blizzard (uno
dei motivi per cui se n'è andato, volontariamente). Ha visto la cultura
interna cambiare per gradi:
\begin{enumerate}
  \item ``facciamo giochi fighissimi, perché i giochi devono essere
        fighissimi'';
  \item ``facciamo giochi fighi così vendono di più'';
  \item ``vendiamo di più facendo giochi fighi'' (sembra uguale, ma è molto
        diverso);
  \item ``facciamo più soldi''.
\end{enumerate}
Lo si capiva dalle scelte di design, osservando quale fosse il peso più
grande sulla bilancia. Lo stesso varrà per gli studenti, nei loro giochi e nei
colloqui di lavoro: capire subito che cosa vuole l'altra parte.

\paragraph{Leader e visione.} Nei gruppi dell'accademia, dal secondo anno, c'è
una persona che assume il ruolo di \emph{lead} o \emph{scrum master}. È
diverso dall'essere un \emph{game director} (``il gioco si fa come dico io,
aiutatemi a farlo con la mia visione''), come avviene nei grandi studi: da lì
nascono i nomi nei titoli, come \emph{Sid Meier's Civilization}, che indicano
la persona che ha fatto quel gioco, oppure un puro branding, come \emph{Tom
Clancy's Rainbow Six}.

\paragraph{Inclusion come jolly.} A margine, uno studente osserva che in molte
aziende ``la risposta giusta è sempre \emph{inclusion}'', usata come un
jolly. È comunissimo, conferma il relatore: \emph{diversity and inclusion}
viene spesso usata per dire ``così includiamo più clienti'', ma non è quello
il suo punto (``altro discorso'').

% ---------------------------------------------------------------------
\section{Confrontarsi con chi è più bravo}
\label{sec:z1-rivalita}

Arrivato a Riot dopo sei anni a Blizzard, il relatore pensava di avere un buon
bagaglio di esperienza. Si è trovato invece ``come un fante buttato in mezzo
ai Navy SEAL'': persone che avevano superato 11, 12, 13 colloqui, con un
livello molto alto.

Individua subito un collega francese che, nella sua testa, diventa ``la
versione migliore di me'': stesso ruolo (community manager), esperienza
simile (arrivava da Ubisoft, come quasi tutti i francesi di Riot), ma molto
più spigliato, sempre con la battuta pronta, ``l'anima della festa''. Il
relatore invece è introverso: un workshop come questo per lui è
\emph{emotivamente drenante}, non ama parlare in pubblico pur amando
condividere le proprie esperienze, e non riesce a sostenere conversazioni
come ``cosa hai fatto questo weekend?''. Lavorando raramente insieme, aveva
sviluppato una rivalità ``da allenatore di Pokémon'' (lui era Ash, il
relatore l'antagonista).

Dopo sei mesi, a un evento di \emph{team building} aziendale (a Tenerife), il
suo capo, molto attento a queste dinamiche, li mette apposta nello stesso
gruppo. A fine serata, complice l'open bar, il relatore gli confessa di averlo
mal sopportato dal primo giorno per invidia della sua estroversione. Il
collega ride e risponde che pochi giorni prima aveva detto al capo la stessa
cosa al contrario: invidiava a Stefano il suo essere sempre organizzato, sul
pezzo, presentarsi nella sala riunioni giusta (mentre lui finiva in quella
sbagliata), arrivare con gli appunti, con i dati a supporto e le analisi di
mercato, mentre lui spesso non aveva nemmeno letto la mail del meeting.
Un'invidia in due direzioni, che li ha portati a lavorare molto insieme e a
migliorarsi a vicenda.

Le lezioni:
\begin{itemize}
  \item non bisogna demoralizzarsi vedendo qualcuno che fa lo stesso lavoro
        ed è più avanti sotto alcuni aspetti: spesso c'è un bilanciamento. I
        ``punti caratteristica'' sono limitati: se li ha messi in carisma,
        non li avrà messi in saggezza;
  \item conviene \textbf{chiedere a queste persone come mai} sono così brave
        in qualcosa, come si fa con un progetto andato bene: magari non lo si
        farà mai bene quanto loro, ma si comincerà ad avere il
        \emph{mindset} giusto; e viceversa, loro saranno spinti a chiedere a
        noi;
  \item spesso \textbf{le proprie doti le scopriamo dagli altri}. Il relatore
        non si era mai accorto di essere ``quello che resta calmo durante le
        emergenze'' (il lancio di \emph{Diablo~III}, quello di
        \emph{Draenor}, i problemi dal vivo alla Games Week) finché non
        gliel'hanno fatto notare tutti: per lui era la normalità. A chi pensa
        ``non sono capace, forse non è il mio percorso'', può darsi che
        semplicemente nessuno gli abbia ancora detto ``questa cosa la sai
        fare davvero bene, concentrati su questo''.
\end{itemize}

% ---------------------------------------------------------------------
\section{Domande sul lavoro nell'industria}
\label{sec:z1-industria}

Diverse domande degli studenti hanno riguardato le condizioni di lavoro nei
grandi studi internazionali. Le risposte sono raccolte qui per argomento.

\subsection{Presenza, remoto e trasferimenti}

\paragraph{Grandi studi: quasi solo in presenza.} Sia a Blizzard sia a Riot si
lavorava \textbf{solo in presenza}, senza smart working: sei anni a Parigi e
cinque a Dublino. Si parlava solo inglese: anche nei meeting tra soli
italiani si tendeva a comunicare e a scrivere in inglese, per correttezza
verso chi si fosse aggiunto e per lasciare una traccia comprensibile a tutti.
Il francese serviva per fare la spesa e sopravvivere, non in ufficio.
L'ambiente era estremamente eterogeneo (francesi, polacchi, spagnoli e
molti altri): durante il colloquio dal vivo a Dublino, mentre si preoccupava
di essersi vestito in modo adeguato, vide passare dietro la parete a vetro un
collega svedese con i dreadlock, a piedi scalzi, con una pianta sotto un
braccio e un aquilone sotto l'altro, e capì di essere ``a posto così''.
Un'esposizione a molte culture e modi di pensare che il relatore considera
bellissima.

Anche oggi, tranne alcuni studi particolari, la maggior parte del lavoro si
svolge di persona: amici a Paradox lavorano a Stoccolma, i pochi rimasti a
Blizzard si sono dovuti trasferire a Irvine. Per posizioni \emph{junior} o
\emph{entry level} i grandi studi raramente concedono il lavoro da casa; per
profili senior può esserci qualche eccezione (ad esempio nove mesi da remoto
e tre in sede).

\paragraph{Relocation.} Spesso tocca spostarsi, ma il relatore non ha mai
pagato un trasloco di tasca propria: queste aziende offrono la
\emph{relocation} (aiuto a trovare casa, ad aprire il conto in banca e
supporto pratico), perché è un investimento sulla persona.

\paragraph{Colloqui mentre si lavora altrove.} Per i colloqui in presenza
prendeva ferie, senza dire dove andava. Il consiglio generale: per fare
colloqui \textbf{prendete ferie o permessi}, non datevi malati, ``che non è un
bel modo di comportarsi''.

\paragraph{Remoto ``doppio''.} Un collega lavora per l'azienda coreana che
gestisce \emph{PUBG}, con ufficio ad Amsterdam, ma lavora da Milano: un
``doppio remoto''. Si trova, ma in genere un ambiente anglofono con sedi in
Europa chiederà di trasferirsi in Europa.

\paragraph{Famiglia e trasferimenti.} Il relatore ha avuto la fortuna di
sposare una persona che fa lo stesso lavoro (conosciuta lavorando su
\emph{Warhammer Online}, poi con lui a Blizzard e a Riot, nella
localizzazione e nel doppiaggio, ad esempio di \emph{World of Warcraft} in
italiano): avendo competenze rare, spesso venivano assunti entrambi, al
massimo a pochi mesi di distanza. Per i developer la soluzione frequente è
che la famiglia segua la persona assunta, perché lo stipendio basta per
tutti. Ma non sempre c'è equilibrio, e a volte uno dei due deve sacrificarsi:
sua moglie ha lasciato Blizzard per seguirlo a Dublino (per poi essere
assunta da Riot); lui ha lasciato Riot per tornare in Italia quando lei ha
trovato lavoro qui. Ha dovuto rinunciare più volte a lavorare in Islanda per
CCP (\emph{EVE Online}), perché non era una scelta solo sua. Si fanno
compromessi, ma si riesce a far quadrare.

\paragraph{Ci si sposta per ogni progetto?} No: si sono spostati
esclusivamente per cambiare azienda. Cambiando progetto nella stessa azienda
è raro, quasi mai (``non l'ho mai visto succedere''), che venga chiesto di
trasferirsi. In Francia stavano benissimo (vivevano fuori Parigi; ``i
problemi con i francesi sono in realtà problemi con i parigini''), parlavano
francese ed erano vicini alla doppia cittadinanza: se ne sono andati perché
non vedevano un futuro a Blizzard, per il modo in cui l'azienda trattava i
colleghi.

\paragraph{Il mercato oggi.} Non esiste un \emph{silver bullet} per trovare
lavoro domani dove si vuole. Con tutti i \emph{layoff} recenti ci sono
migliaia di esperti del settore che cercano lavoro: serve pazienza, ma è
fattibile.

\paragraph{Studi indie e remoto.} Negli studi più piccoli il remoto è più
diffuso: a Stormind circa 60 persone su 95 lavorano da remoto, ma lo studio
si è strutturato fin dall'inizio per farlo. Serve consapevolezza. Bending
Spoons (che sviluppa app per mobile, non giochi, ed è l'azienda dell'app
Immuni), pur essendo grandissima e con sede a Milano, lavora da remoto e
lascia l'ufficio a chi vuole: lo smart working è posizionato quasi come un
\emph{benefit}, ma permette anche di avere un ufficio da 300 persone invece
che da 700. Le sedi fisiche costano decine di milioni all'anno (Blizzard ha
chiuso le sedi di Amsterdam e di Parigi). Tuttavia, per lavori molto
verticali e specifici gli studi preferiscono avere le persone dal vivo; i
ruoli più spesso remoti sono quelli che non fanno sviluppo (customer service,
PR\dots).

Anche gli artisti lavorano spesso in presenza: l'accademia può dare solo una
\emph{simulazione} di quanto stretto sarà il contatto tra programmazione,
design e art. ``Sarà ancora di più.''

\paragraph{La comunicazione da remoto.} È più difficile. Pur essendo
estremamente introverso, il relatore \textbf{obbliga tutti a usare la webcam}
nei meeting in cui conta: non in un aggiornamento con novanta persone
collegate, ma certamente in un colloquio a due sulla performance dell'ultimo
anno. La componente fisica è una parte enorme della
comunicazione\footnote{In lezione si parla di ``quasi l'ottanta per cento'';
il dato va preso come indicativo, poiché le percentuali di questo tipo
circolano in forme diverse e sono spesso citate fuori contesto.}: anche se si
è timidi, accendere la webcam è l'approssimazione più vicina all'essere dal
vivo. Negli ambienti dove si lavorava senza webcam la situazione era
``esponenzialmente peggiore''.

\subsection{Differenze culturali nel team}

\paragraph{Gestire sensibilità diverse.} A Riot i team erano divisi per
\emph{lingua}, non per nazione: il team di lingua spagnola includeva anche un
messicano, quello di lingua russa includeva russi e ucraini, proprio nel
periodo della guerra in Crimea. Tra i colleghi non c'era alcun problema (erano
amici), ma ci si è posti la domanda: \emph{come comunichiamo con una community
che sarà divisa su questo tema?} Nei giochi online la divisione era evidente:
in \emph{Rocket League} c'era chi, con un compagno di squadra con la bandiera
avversaria, si faceva otto autogol pur di farlo perdere. Allo stesso modo,
con molti americani mandati a Dublino a coprire posizioni scoperte, c'erano le
elezioni (Clinton--Trump, Biden--Trump) e varie guerre: serviva sempre una
certa sensibilità.

Imparare a navigare queste situazioni lo ha reso un professionista migliore:
gli ha insegnato a fermarsi e chiedersi ``io la percepisco così; gli altri
dieci, con background, storie e lingue diverse, come la penseranno?'', a
ragionare più \emph{long term} e più \emph{big picture}, con più empatia.

\paragraph{Stili di comunicazione.} Un partecipante che ha lavorato nel Regno
Unito chiede dello stile comunicativo italiano, a volte frainteso. Uno dei più
grandi complimenti ricevuti dal relatore è stato ``non sembri italiano''
(``gesticoli poco per essere italiano'', ``parli bene inglese per essere
italiano''). Ha imparato ad \textbf{adattare la comunicazione alla persona}:
con spagnoli e francesi la battuta sulla nazione funziona, perché c'è un
``odio--amore'' condiviso; con russi e giapponesi no. Doveva adattare anche
l'umorismo. Un collega metà russo e metà messicano diceva di aver preso ``il
lato peggiore di entrambe le nazionalità'', ed è stato un ottimo maestro per
imparare a leggere culture diverse.

Non bisogna generalizzare, ma gli stereotipi possono essere un punto di
partenza ragionevole, \emph{sapendo che troverete sempre le eccezioni}. E
vanno letti correttamente: gli americani tendono a essere più estroversi e ad
avere sempre qualcosa da dire non perché ``sono fatti così'', ma perché
vengono educati così (all'università tra i corsi c'è il \emph{debate}, saper
argomentare in pubblico); un collega giapponese può essere più introverso per
ragioni culturali. Nel dubbio si parte da lì, poi si conosce la persona.

\subsection{Negoziare, discriminazione e stipendi}

\paragraph{La famiglia come leva.} Avere un partner riduce le opzioni, ma le
aziende lo sanno: uno spostamento è per due. Si può usare come leva per
negoziare (un compenso diverso, un supporto per il trasloco): se un'azienda ha
deciso di volere quella persona, 2, 3, 4 o 10.000 euro in più all'anno, per
un'azienda che fattura centinaia di milioni o miliardi, non cambiano nulla.

\paragraph{Essere sposati è un motivo per non essere assunti?} No. In alcuni
paesi è proprio illegale discriminare su questo: in Francia era illegale
chiedere CV con data di nascita, foto, sesso o religione, e chiedere se si
vogliono figli. Purtroppo, dai racconti dei suoi studenti di master (in
settori come \emph{beauty and wellness} e \emph{food and wine}), in piccole
realtà italiane queste domande (``hai 25 anni, hai intenzione di
sposarti?'') vengono ancora fatte, e servono a filtrare i candidati. Nelle
grandi aziende ci sono sistemi per evitarlo; anzi, molte danno \emph{più}
benefit a chi ha famiglia (``porta pure la famiglia, è tutto pagato''),
perché le persone competenti hanno spesso una certa età e quindi,
ragionevolmente, una famiglia.

\paragraph{Gli stipendi.} Dipendono molto da dove si vive. Lo stipendio a Riot
non era paragonabile a quello di Blizzard perché in Francia era tassato al
52\% e in Irlanda al 17\%\footnote{Aliquote come riportate in lezione,
riferite alla situazione personale del relatore.}: nessuno stipendio
competitivo può compensare il 35\% di differenza di tasse. Negli Stati Uniti,
con pochissima previdenza pubblica, gli stipendi sono molto più alti perché
poi ci si paga tutto (``se hai un raffreddore paghi tu''). Bisogna valutare il
potere d'acquisto di quei soldi nel paese in cui si andrà a vivere.

\subsection{Freelance e \emph{work for hire}}

All'inizio della sessione pomeridiana uno studente che lavora nell'audio
chiede quali siano le condizioni per chi lavora come \emph{freelance}.
\begin{itemize}
  \item Dipende molto dal progetto: spesso il contratto è a progetto, per
        portare a termine una parte o un intero progetto.
  \item Non tutte le aziende usano freelance: Riot e Blizzard lo facevano
        pochissimo e preferivano assumere (magari trattando poi la persona
        ``come un freelance'', lasciandola a casa a fine progetto).
  \item Spesso i lavori offerti come freelance sono in realtà
        \textbf{sostitutivi di un full time} (``freelance, ma sei ore al
        giorno'').
  \item Si può fare, ma bisogna essere molto chiari nel definire l'impegno
        (quante ore, come vengono tracciate): non tutti specificano cosa
        significa davvero ``freelance''. Lavorando così per due o tre aziende
        si può arrivare a un tempo pieno.
  \item È più comune il \emph{work for hire}: aziende che forniscono personale
        ad altre aziende; si è assunti da X ma si lavora ``embeddati'' presso
        Y (una sorta di subappalto). Anche lo studio del relatore lo fa.
  \item Il puro freelance di quattro mesi esiste, ma di solito per quei
        quattro mesi diventa un lavoro a tempo pieno; è raro ``quattro mesi,
        due ore al giorno''. Nelle grandi aziende è improbabile lavorare nello
        stesso mese per Ubisoft, Riot e Blizzard tre ore ciascuna; è più
        realistico lavorare per più startup indie italiane, che non possono
        permettersi un'assunzione ma possono pagare due ore al giorno.
  \item Molto passa dai contatti: trovare le aziende giuste e rendere esplicite
        le condizioni.
\end{itemize}

% ---------------------------------------------------------------------
\section{Il circolo virtuoso: feedback, miglioramento, fiducia}
\label{sec:z1-triangolo}

\subsection{Il triangolo}

L'obiettivo a cui il workshop vuole arrivare è un \textbf{triangolo che si
autoalimenta}:
\begin{enumerate}
  \item si dà un \textbf{feedback} con l'obiettivo di migliorare le
        performance (inclusi elementi caratteriali);
  \item la persona \textbf{migliora} e vede di averne beneficiato;
  \item questo genera \textbf{fiducia}, che rende più facile dare e ricevere
        il feedback successivo.
\end{enumerate}
Togliendo uno qualsiasi dei tre elementi, il circolo non è altrettanto forte.

In qualsiasi fase si parta, il rispetto per i colleghi è inizialmente dettato
dall'\emph{educazione}, non dalla fiducia. La fiducia arriva quando i colleghi
cominciano a dire cose giuste e sensate e a \textbf{mantenere la parola}.
Dare e ricevere feedback è il modo in cui, in un certo senso,
``inganniamo il cervello'' per creare una \emph{case history} di fiducia: si
ha motivo di fidarsi di più, di ascoltare di più, di essere più ricettivi. E
si diventa sempre più bravi ad accettare il feedback e ad agire di
conseguenza.

\subsection{Accettare un feedback non vuol dire fare come dice l'altro}

Dieci anni fa il relatore non era ostile ai feedback, ma molto
\emph{difensivo}: ``sì, sì, va bene, ora ci guardo''; oppure, a chi gli diceva
``avresti dovuto farla così'', rispondeva spiegando perché l'aveva fatta in
quel modo. Oggi ha imparato a dire ``ok, grazie, ora ci ragiono'': non deve
difendere la propria posizione da chi viene solo a dire ``hai pensato a
questa cosa?''. Prima lo viveva come un attacco (``mi viene a dire che non so
fare il mio lavoro''), ora come un ``ottimo punto di vista, lo prendo in
considerazione''.

Con una postilla fondamentale: \textbf{accettare il feedback non significa
fare come dice l'altro}. ``Mi hai detto che preferisci questo mob verde. Ho
pensato anche a farlo verde? Sì. Penso comunque che giallo sia meglio? Sì, lo
faccio giallo.'' Il relatore accetta \emph{sempre} il feedback, valido o no:
garantisce di tenerlo a mente quando lavora. Che poi faccia come gli viene
detto è un altro discorso.

\subsection{Fiducia e scadenze}

La fiducia funziona allo stesso modo con gli impegni. Se si chiede a qualcuno
una cosa per venerdì e la consegna venerdì, la volta dopo gliela si chiederà
per venerdì. Se invece arriva il giovedì della settimana dopo, la volta dopo
gli si darà una \textbf{deadline anticipata}, come con chi arriva sempre in
ritardo agli appuntamenti (``se ti dico alle 7 arrivi alle 7 e mezza; la
prossima volta ti dico alle 6 e mezza'').

Al lavoro questo è dannosissimo, perché crea \textbf{deadline fittizie} che
non sono mai quelle vere, e nascono discussioni del tipo ``perché devo farlo
in due giorni se si fa in cinque?'' (``sì, se tu lo facessi in cinque''). La
conversazione giusta è un'altra: \emph{``che problemi hai avuto nel
consegnare venerdì?'' ``Questo, questo e questo.'' ``Cosa posso fare per
evitare che succeda?''}

Il relatore non vuole mai sentirsi dire ``non sono riuscito a farlo''.
Preferisce di gran lunga sentirsi dire, tre giorni prima, \textbf{``potrei non
riuscire a farlo, per questi motivi''}. Chi ha la visione d'insieme può dare
più budget, più persone, spostare la deadline, cambiare lo scopo del progetto:
può fare qualsiasi cosa per aiutare. Ma se lo si scopre il venerdì, cinque
minuti prima della deadline, non può fare più niente. Il messaggio è:
\textbf{``aiutami a poterti aiutare''}. Non penserà mai male di chi dice ``non
ce la faccio per questi motivi'', perché non sta dicendo ``non ho voglia'': sta
indicando i problemi, e lavorare su quei problemi è compito del manager. Non
può invece fidarsi di chi fino al giorno prima dice ``tranquillo, ce la
faccio'' e il giorno dopo ``non ce l'ho fatta''.

Questo è fondamentale per chi lavora in programmazione, design o art, perché
ci saranno sempre deadline: bisogna sempre \textbf{lavorare sulle aspettative},
dire al proprio manager (o al proprio report) cosa è lecito aspettarsi. Non
bisogna fare come i pinguini di \emph{Madagascar} (``sorridete e salutate, va
tutto bene''): se c'è un problema, va detto, altrimenti si peggiora solo la
propria situazione.

Da italiani si tende a non voler mai alzare la mano per chiedere aiuto
(``no, no, ce la facciamo, lo faccio io''). A volte non ce la si può fare, non
perché non si è capaci, ma perché vari fattori lo impediscono. Accettarlo è
importante: \textbf{è molto meglio avvisare e chiedere una mano che dire
``non ce l'ho fatta''}.

\subsection{Domande sul dare e ricevere feedback}

\paragraph{``Ne parliamo a fare, se poi fai come vuoi?''} Uno studente (Cristian)
riceve feedback che, dopo averli analizzati, non ritiene congrui; quando non
li segue gli viene detto ``ne parliamo a fare?''. C'è un problema di
\textbf{aspettative}: l'idea che accettare un feedback voglia dire fare come
dice l'altro. Accettarlo significa prendere in considerazione ciò che è stato
detto. Chi decide diversamente ha però un \textbf{onere}: argomentare perché
si fa comunque a modo proprio (``mi hai detto verde, la facciamo gialla
perché\dots''). Poi bisogna solo aspettare i risultati per vedere chi aveva
ragione.

La decisione spetta sempre \textbf{a chi ha la responsabilità di quella
decisione}. Se il game director deve scegliere tra A e B, anche se il
feedback dice B, la scelta è sua perché la responsabilità è sua. Non sarebbe
corretto essere ritenuti responsabili di una scelta delegata a chi dà
feedback. La risposta può essere semplicemente: ``l'ho tenuto in
considerazione, ma visto che ne sono responsabile io e le conseguenze sono
mie, ho scelto di fare X''.

\paragraph{E se l'altro fa ripicca?} Un partecipante chiede di chi, vedendo
il proprio feedback non accolto, si vendica ignorando a sua volta feedback
buoni. Per il relatore è un problema caratteriale: chi si comporta così ``è
già sulla strada per uscire''. Non c'è spazio per l'ego in team di questo
tipo. Il fenomeno nasce dall'idea che \emph{il feedback negativo faccia fare
brutta figura}, che il fatto che si potesse fare meglio tolga qualcosa alle
proprie capacità.

Non è così: \textbf{il feedback non è un giudizio}. Si può passare un anno a
Riot ricevendo solo feedback negativi e avere una valutazione annuale
stellare, perché ricevere molti feedback negativi significa essere una persona
approcciabile, che agisce sul feedback e invita gli altri a darlo. Il relatore
ha visto valutazioni molto peggiori per chi riceveva solo feedback positivi,
perché la gente aveva paura di dargliene di negativi: ma solo ``pacche sulle
spalle'' non fanno migliorare. Il problema va affrontato parlando con la
persona; a volte le conversazioni iniziano con \textbf{``ti do un feedback sul
feedback''}, cioè sul modo in cui l'altro ha dato il suo.

Esempi dalla sua esperienza:
\begin{itemize}
  \item lui chiede spesso di \textbf{non ricevere feedback senza esempi
        tangibili}, perché fa fatica ad astrarre: ``sei stato maleducato'' non
        significa nulla, ``ieri al meeting mi hai risposto male su questa
        cosa'' sì;
  \item ad altri serve l'opposto: ha ricevuto il feedback ``quando devi dirmi
        che ho sbagliato, puoi girarci un po' intorno e darmi del contesto?'',
        perché il suo stile diretto (``questa cosa non ha funzionato, dovremmo
        fare così'') feriva. Ha adattato il suo modo di dare feedback: con
        alcuni è diretto, con altri gira intorno. L'importante è saperlo.
\end{itemize}

\subsection{Non portare la soluzione}

Un consiglio centrale: \textbf{quando si dà un feedback, non si dà una
soluzione}. Non si dice ``io farei così'', ma ``secondo me questa cosa si
potrebbe fare diversamente, per questo e questo motivo: come potremmo farla
diversamente?''. È come chi smette di fumare perché vuole e chi smette perché
glielo dicono moglie e figli: se non si è convinti, spesso non si smette.

Quando il relatore diceva ``c'è questo problema, io farei così'', il meeting
successivo iniziava inevitabilmente con ``ci ho provato come hai detto tu, con
me non ha funzionato''. Sempre. Perché la soluzione era sua. La
responsabilità della soluzione deve essere invece \textbf{congiunta}: ogni
feedback deve creare un piccolo team che ha come obiettivo risolvere quel
problema. Questo toglie l'antagonismo (non ``ti dico che non sai fare il tuo
lavoro'', ma ``ti voglio aiutare a farlo meglio, lavoriamoci insieme'') e
sposta il focus: il problema non sei tu, è ciò che sta accadendo.

\paragraph{E se sono i piani alti a tenere all'oscuro dei problemi?} È una
questione di cultura aziendale. Il relatore fa sempre presente al proprio capo
il concetto di \textbf{service level agreement} (che verrà ripreso nella
seconda giornata): ``se mi dai questi input, ti garantisco questo output in
questi tempi''. Se il capo chiede delle cose senza dire budget, tempi e
perché, non può aspettarsi lo stesso risultato. È una comunicazione a due
vie: si dice al manager cosa serve per fare il lavoro, e lui deve fornirlo o
accettare che il risultato non sarà quello voluto.

\subsection{Amici, parenti e parole proibite}

\paragraph{``Non si lavora con amici e parenti''?} Per sua definizione
personale il relatore ha pochissimi amici (``stanno sulle dita di una mano'')
e una soglia molto alta per considerare amico un collega. Allo stesso tempo,
in quindici anni e con centinaia di colleghi, ha avuto solo \textbf{due}
esperienze in cui non è riuscito a lavorare con qualcuno: non serve essere
amici per far funzionare le cose. Con gli amici, anzi, gli viene più facile:
se dice a un amico che una cosa ``fa schifo'', l'amico non pensa ``come si
permette'', ma ``se perfino Stefano, che è mio amico, me lo dice, allora è
proprio così''. Diventa un rafforzativo, grazie alla fiducia costruita prima.
Il problema del lavorare con amici e parenti nasce dal \emph{non voler} dire
all'amico che sta sbagliando; se invece la confidenza permette di dirlo, si
può saltare la ``confezione'' del messaggio e passare direttamente al
contenuto.

\paragraph{Le frasi peggiori.} Alla domanda se esistano parole peggiori di
``ottimo lavoro'', il relatore indica la sua classifica: ``ottimo lavoro''
(detto in modo vuoto)\footnote{Il passaggio della trascrizione è
frammentario: si intende verosimilmente il complimento generico e di
circostanza, che non dice nulla di utile.}; poi la frase che lo ``fa uscire di
testa'', \textbf{``abbiamo sempre fatto così''} (``e infatti avete chiuso'',
come Blockbuster); infine \textbf{``siamo una famiglia''}, riferito al posto di
lavoro. Un'azienda non è mai una famiglia, perché non ci si può aspettare lo
stesso impegno da chi la possiede e da chi ci lavora: non c'è un
corrispettivo proporzionale tra impegno e ritorno. Per quanto inclusiva,
flessibile e razionale, non sarà mai una famiglia.

\subsection{Ego, muri e il peggior modo di scegliere un leader}

\paragraph{Chi mette davanti l'ego e alza muri.} Esistono, ma di solito durano
pochissimo in azienda. Non si tratta di avere un ego forte, ma di \textbf{non
saper ricevere feedback}. Spesso sono persone bravissime nel loro lavoro che
poi si chiedono perché tutti intorno diventino più bravi di loro: perché gli
altri accettano feedback e hanno più fonti da cui attingere, mentre loro si
sono rintanati nella loro torre.

Molti manager tendono ad \textbf{assumere persone un filo peggiori di loro},
per non fare brutta figura. Il relatore assume solo persone più brave di lui,
``perché non è una gara'': ha lavorato con molti community manager più bravi
di lui come community manager, ma lui era più bravo a fare il manager. Non
viene valutato su quanto sa fare il community manager, ma su quanto il suo
team sa farlo. Un partecipante chiede se la paura di avere qualcuno di
migliore esista anche in ambito accademico. Assolutamente sì: al relatore è
capitato di non essere preso come docente in un corso di formazione sui
videogiochi perché il responsabile aveva percepito che ne sapesse più di
lui\footnote{La trascrizione dice ``che non sapevo più di lui'', ma il senso
del passaggio richiede il contrario.}. ``Ma perché te ne importa? Tu sei quello che ha portato chi
sa fare il suo lavoro: ben venga.'' O si ha questa cultura, o si ha paura per
il proprio posto.

\paragraph{Promuovere il migliore.} Quando bisogna scegliere un leader per un
team (programmazione, design, art), molto spesso si sceglie \textbf{il migliore
tecnico}: il miglior programmatore, il miglior designer. È spesso la scelta
sbagliata: è come prendere un idraulico per fare l'elettricista. Fare il
\emph{people manager} non è come fare il designer o il programmatore; si può
essere un ottimo programmatore e un pessimo manager, e viceversa. Le aziende
l'hanno capito e lo fanno meno di prima. A Blizzard c'era un grave problema di
\emph{middle management}: molte persone in ruoli importanti erano state
promosse cercando le skill sbagliate, erano frustrate dal non saper fare il
proprio lavoro e schiacciavano chi stava sotto per paura di fare brutta
figura.

Un partecipante cita una teoria simile: una persona viene promossa finché non
raggiunge \emph{il proprio livello di incompetenza}, dove si blocca\footnote{È
il \emph{principio di Peter} (Laurence J.\ Peter, 1969), citato in lezione
come ``la teoria di Peter''.}. Un altro cita Michael Scott di \emph{The
Office}, su cui il relatore tornerà (sezione~\ref{sec:z1-quando}).

% ---------------------------------------------------------------------
\section{Perché la cultura conta: retention, costi e fiducia}
\label{sec:z1-cultura}

\subsection{L'ambiente di lavoro e la retention}

Perché tutto questo è così importante? Perché l'ambiente di lavoro conta
moltissimo per la \textbf{retention}, cioè per trattenere le persone e farle
lavorare in modo efficiente. Se contasse solo lo stipendio, molti non
lascerebbero mai Riot o Blizzard; e più si va avanti con l'età, più si impara
a dare valore anche al resto.

Inoltre, secondo gli studi citati nelle slide, chi vive il proprio lavoro in
modo negativo è \textbf{meno efficiente nel problem solving}: l'energia
mentale va verso ``come faccio a non farmi dare la colpa'' o ``come faccio a
fare bella figura'', non verso la soluzione del problema. In ambienti negativi,
dove ci si deve preoccupare della propria ``sopravvivenza'', si lavora
peggio.

\subsection{Il costo di chi se ne va}

Ogni volta che qualcuno si dimette o viene licenziato, il costo di trovare
una nuova persona, fare i colloqui, formarla e farle attraversare forming,
storming, norming\dots\ è circa \textbf{un terzo del suo stipendio}. Per questo
le aziende tendono a fare più \emph{retention} che \emph{acquisition}, almeno
con le figure considerate utili alla produzione.

Allo stesso tempo si vedono sempre più spesso \emph{layoff} e \emph{hiring
wave}: si licenziano 3.000 persone e sei mesi dopo se ne riassumono 2.500 o
3.500, ma tutte più junior. Il relatore ne ha scritto in un articolo su
Eurogamer.it, in cui spiega anche alcuni ``trucchi del mestiere'': perché i
licenziamenti avvengono sempre in certi periodi e sempre ``a gruppi'' (prima
Bethesda, poi Blizzard, poi un'altra\dots).

\paragraph{L'offerta di Riot.} A Riot esisteva una pratica (di cui il relatore
non ricorda il nome) per cui, firmato il contratto, si aveva \textbf{un anno
di tempo per dire ``non mi piace, non me la sento, me ne vado''} e ricevere
subito come buonuscita un anno di stipendio. Avevano stimato che fosse molto
meglio pagare subito un anno a qualcuno e lasciarlo andare che tenerlo due,
tre, quattro anni come \emph{underperformer}: la responsabilità della scelta
di assumere era dell'azienda\footnote{La pratica, nota come ``The Offer'', è
descritta in fonti pubbliche con condizioni diverse (un periodo più breve e
una somma pari a una frazione dello stipendio annuo). Si riportano le
condizioni come raccontate in lezione; andrebbero verificate.}. L'hanno usata
pochissimi, ma il solo fatto che esistesse faceva lavorare meglio anche lui da
manager, perché sapeva che chi non si trovava bene poteva andarsene con un
anno di tempo per trovare un altro lavoro.

\paragraph{Il ricircolo.} C'è un forte ricircolo nell'industria: molti ex
colleghi passano da Ubisoft a Paradox, da Paradox a Funcom, da Funcom a
Blizzard. Il relatore ha scelto di uscirne lavorando solo per aziende
\textbf{indipendenti}: Stormind non è un indie di cinque persone, ma è
indipendente, e non deve accontentare azionisti. Quando vede un'azienda con
utili miliardari licenziare migliaia di persone si chiede ``dov'è la linea''
(\emph{when enough is enough}): con azionisti che vogliono un ritorno, se
quello è l'obiettivo stanno facendo la cosa giusta, ma lui non condivide
l'obiettivo\footnote{In lezione si parla di ``6,5 miliardi di utile'' e di
``3.000 persone licenziate''; le cifre sono riportate come dette e non sono
state verificate.}.

\paragraph{L'open space vuoto.} A Riot, durante un'ondata di licenziamenti,
del suo team di trenta persone è rimasto solo lui: ``teniamo te come
coordinatore, poi ricostruirai il team''. Ha rifiutato di ricostruire un team
cancellato ``con un colpo di spugna'' e ha deciso che non avrebbe più lavorato
per una grande azienda. Nessuno aveva pensato di spostarlo, e ha lavorato per
tre mesi in un open space di cento metri quadri con la sua scrivania e altre
ventinove vuote. Su una lavagna aveva scritto i nomi dei suoi ragazzi e li
cancellava man mano che, anche grazie ai suoi \emph{referral}, trovavano
lavoro. Cancellato l'ultimo nome, si è dimesso: per lui il lavoro era finito.

\paragraph{Team che si dissolvono.} Uno studente cita un articolo su Larian
Studios, il cui management è contrario a ricostruire ogni volta il team a fine
progetto. È lo stesso problema: si pensa ai ``giochi Blizzard'' o ai ``giochi
BioWare'', ma le persone che li hanno fatti non ci sono più; è rimasto solo il
brand, e non ha senso avere le stesse aspettative da un gioco fatto da persone
diverse. E anche chi sopravvive ai layoff vive pensando che il prossimo
potrebbe essere lui: alcune persone, licenziate nell'ultima ondata a
Blizzard, gli hanno detto ``erano due anni che lo sapevo''. Vivevano a Cork,
``un paesotto irlandese dove o vai per lavorare a Blizzard o non ci vai'',
avevano comprato casa: non ci si può aspettare chissà che morale.

\subsection{La cultura è forte quando si lavora per gli altri}

\begin{quote}
  \emph{La cultura è forte quando le persone lavorano per gli altri e a favore
  degli altri, e non per se stesse.}
\end{quote}

Se si dà feedback per far migliorare l'altro e non per imporre il proprio
modo, se si aiuta per il bene del progetto e non per fare bella figura, la
cultura si rafforza. È dai \textbf{piccoli gesti} che emerge la cultura
aziendale: ad esempio, dal non dover chiedere ``perché mi stai aiutando?''.

\paragraph{L'incognita.} \textbf{L'incognita è spesso dannosa quanto le
cattive intenzioni.} In un esperimento (``mi pare a Stoccolma, cinque o sei
anni fa'') si distribuivano per strada banconote da 500 euro, e la maggior
parte delle persone le rifiutava chiedendosi ``dov'è la fregatura?''
\footnote{L'esperimento è riportato in modo vago in lezione e non è stato
possibile identificarlo con certezza.}. Lo stesso accadeva al relatore,
arrivato a PG e a Stormind da Riot: a ``vedo che sei in difficoltà, posso
aiutarti?'' la risposta era sempre ``perché?''. Alla fine rispondeva: ``perché
quando sarò in difficoltà io vorrei che tu mi aiutassi''.

Quando si comincia a lavorare per gli altri e si smette di pensare ``così avrò
la promozione, così avrò l'aumento'', promozioni, aumenti e rispetto arrivano
da soli. A Blizzard aveva fallito tre volte i colloqui per essere promosso a
manager (inteso come coordinatore del team); quando ha deciso di smettere di
pensarci e di fare solo il bene del proprio lavoro, la volta dopo è stato
promosso. Funziona così: chi aiuta i colleghi e fa andare bene il progetto
viene notato (``questa persona fa migliorare cinque persone, non solo se
stessa: vale di più'').

\subsection{Il lavoro è valutato come opera d'insieme}

Il relatore riconosce che buona parte del suo successo è venuto dal fatto che
il suo team lavorava bene, non dalla sua bravura personale. Anche i progetti
degli studenti, presentati a Diego, a Enrico, ad Andrea, verranno valutati
\textbf{come opera d'insieme}. Nessuno dice ``il \emph{combat system} di questo
gioco l'ha fatto X, che bravo'': si dice ``BioWare ha fatto un bel gioco'',
``Blizzard ha fatto un bel gioco''.

Lo ha replicato anche lato HR. In molte aziende ci sono \textbf{bonus
personali} legati a obiettivi (se il gioco vende X, se la feature è pronta in
tempo). Il relatore li ha dimezzati: metà restano personali, metà diventano
\textbf{bonus di team}, che tutti ricevono se il team raggiunge gli obiettivi.
Un incentivo concreto a collaborare: non solo si dice che il lavoro di gruppo è
importante, lo si dimostra dandogli più valore economico. Aveva proposto il
100\% di bonus di team, ma gli è stato concesso il 50 e 50; ha funzionato
comunque, anche se sarebbe stato ancora più potente. Ovviamente, alla
domanda ``perché mi aiuti?'' si può anche rispondere ``perché voglio il
bonus'': un aiuto dato ``per il motivo sbagliato'' era un compromesso che era
disposto ad accettare.

\subsection{Il coraggio di chiedere aiuto}

Un partecipante osserva che aiutarsi dovrebbe essere la norma. Il modello
diffuso, invece, è ``devi essere forte e indipendente, non chiedere mai una
mano'', e vale in tutte le direzioni (dal maschilismo alla ``donna forte e
indipendente che non ha bisogno di nessuno''): \textbf{non ammettere di avere
bisogno di aiuto è sempre uno dei motivi per cui le aziende vanno male}, che si
tratti di uomini, donne o persone non binarie. Il problema è non avere la
sicurezza di poter alzare la mano.

Il relatore ha cercato di togliere ogni conseguenza dall'aver alzato la mano:
se si \emph{depenalizza} qualcosa, le persone non sono più disincentivate a
farla. Ha detto espressamente che non avrebbe mai valutato nessuno per il
numero di volte in cui chiedeva aiuto o per il numero di feedback negativi
ricevuti, e questo ha funzionato.

\subsection{Perché si dà un feedback}

\emph{Perché avete dato il vostro ultimo feedback? Che cosa volevate ottenere?}
Le risposte della classe:
\begin{itemize}
  \item per \textbf{riconoscere il valore} di un collega, anche verso il
        manager, perché qualcuno che non aveva visibilità vedesse quanto
        quella persona fosse efficace;
  \item per \textbf{spingere una persona a superarsi};
  \item per aiutare un collega a rendere il prodotto all'altezza di ciò che
        voleva rappresentare;
  \item per \textbf{genuinità}: c'era qualcosa che poteva essere fatto meglio.
\end{itemize}
Il relatore aggiunge che la forma cambia col tempo: all'inizio era ``devo
darti un feedback, hai dieci minuti?''; ora è un messaggio informale su
Discord, perché ha costruito un rapporto di fiducia e sa che non verrà letto
come ``l'HR manager mi scrive, sono nei guai?'', ma come ``Stefano è così, mi
dà feedback, resterà tra noi e lo risolveremo insieme''.

Un altro motivo è \textbf{risolvere incompatibilità caratteriali}: le
conversazioni più difficili, del tipo ``io e te non stiamo funzionando, cosa
possiamo fare diversamente?''. Sono importanti perché, se due responsabili non
vanno d'accordo, il team se ne accorge, comincia a raccontarsi storie (``si
odiano, ci saranno problemi'') e, per difendere il proprio manager, comincia a
giudicare male il lavoro dell'altro gruppo. È la dinamica da fermare:
\textbf{valutare gli altri solo per difendere la propria fazione}.

Dare e ricevere feedback sono \textbf{due muscoli da allenare} (il relatore è
sempre stato più bravo a darlo che a riceverlo), e si impara col tempo. È come
fare colloqui: fare colloqui è una skill a parte, che non indica quanto si è
bravi nel proprio lavoro. Essere scartati non significa non saper fare quel
lavoro: spesso significa aver fatto male il colloquio o che è stato preferito
qualcun altro.

\paragraph{Feedback pubblici e privati.} Uno studente racconta di aver dato un
feedback negativo a un superiore, senza un motivo particolare, per dimostrare
a sé e ai colleghi che non si viene licenziati per questo. Il relatore coglie
l'occasione per precisare che \textbf{i feedback negativi si danno
individualmente}: lui dà moltissimi feedback di cui il resto dell'azienda non
sa nulla. Reagirebbe molto male se, in un meeting, qualcuno alzasse la mano per
dire ``comunque questa cosa non l'hai fatta bene'': quello non è feedback, è
discussione, accettabile solo se è un brainstorming e se l'opinione è stata
chiesta. Alcune aziende hanno piattaforme pubbliche di feedback, ma spesso
servono solo ad applaudire; lo stesso accade nelle aziende che
\emph{obbligano} a lasciare un certo numero di feedback a settimana, che
diventano ``lui è bravo, lei è competente, anche questa settimana l'abbiamo
fatto''. Bisognerebbe chiedere a chi le adotta cosa pensa di ottenere (forse
trasparenza), ma difficilmente una persona ``attaccata'' pubblicamente reagisce
bene.

% ---------------------------------------------------------------------
\section{Among Us: la caccia al traditore}
\label{sec:z1-among-us}

Prima della pausa pranzo il relatore mostra tre film di fantascienza di
qualche decennio fa e chiede che cosa li accomuni\footnote{I titoli non sono
leggibili nella trascrizione: viene citato \emph{Event Horizon} da uno
studente, e il film ``in mezzo'' è indicato come ``l'originale'',
verosimilmente \emph{La cosa} (\emph{The Thing}).}. In ognuno c'è un gruppo di
persone con un obiettivo e la presenza, reale o percepita, di qualcuno che
rema contro: \textbf{l'energia del gruppo viene spesa per trovare il
responsabile}.

\paragraph{L'esperimento degli alieni.} Nei master in presenza il relatore
propone un esercizio. Distribuisce segretamente cartellini plastificati con
scritto ``essere umano'' o ``alieno''; le carte alieno contengono istruzioni di
comportamento con cui gli alieni, che cercano di fingersi umani, si
tradirebbero (dare sempre ragione a chi ha parlato prima, finire ogni frase
con una domanda, parlare in un certo modo). Lo scopo: in dieci minuti formare
gruppi, e \textbf{vince il gruppo più numeroso senza alieni}; basta un alieno
per perdere. Si formano sempre gruppi piccoli (due, cinque, sei persone),
perché la tendenza è ``mi fido al massimo di voi due''.

Ma vincono i gruppi più grandi, per un motivo semplice: il relatore
\textbf{non distribuisce nessuna carta alieno}. Mostra a tutti entrambe le
carte, ma consegna solo carte ``umano''. Vuole simulare ciò che accade in
ufficio: la tendenza a percepire sempre qualcuno che rema contro (i manager,
l'altro team, quel collega), quando in realtà non c'è nessuno, ``è tutto nella
nostra testa''. Quando chiede le motivazioni, emergono sospetti
\emph{fabbricati}: ``faceva sempre domande'', ``parlava solo all'infinito''
(``no, parla così e basta''). Ci si fabbrica da soli i motivi per non fidarsi
degli altri, ed è comunissimo.

\paragraph{Among Us.} L'autosabotaggio è, per tornare ad \emph{Among Us},
\textbf{concentrarsi sul cercare cadaveri e traditori invece di fare le
task}: non si vince perché non si sono completate le task, mentre si cercava il
colpevole. Succede tantissimo negli uffici, ed è il motivo per cui l'80\% dei
problemi che il relatore risolveva erano di comunicazione. Entrando nei team
sentiva immancabilmente: ``facciamo il possibile, ma questa lotta continua con
l'altro gruppo è estenuante''. ``Perché c'è questo problema?'' ``Non lo so,
c'era già quando sono arrivato.'' Un paradosso: c'è un problema, nessuno sa
dire quale sia, ma c'è sempre stato e quindi lo si accetta.

\paragraph{L'esperimento delle scimmie.} È l'esempio, ``finto o reale'', delle
scimmie in gabbia\footnote{È il noto racconto delle ``cinque scimmie'', spesso
presentato come esperimento scientifico ma probabilmente apocrifo (lo stesso
relatore lo definisce ``finto barra reale''); nella trascrizione il numero
delle scimmie varia.}. In una gabbia con alcune scimmie pende una banana
elettrificata: chi la tocca fa prendere la scossa a tutte. Le scimmie imparano
e, quando una si avvicina alla banana, le altre la picchiano. Poi si
sostituisce una scimmia alla volta con una nuova: la nuova va verso la banana,
le altre la picchiano, e smette. Si continua finché nella gabbia ci sono solo
scimmie nuove, poi si toglie la corrente. Nessuna scimmia tocca la banana:
nessuna ha mai preso la scossa, ma tutte hanno imparato che chi tocca la
banana prende le botte.

Nei team è lo stesso: il motivo per cui si ``odia'' l'altro team si
trasforma, ma un vero motivo non c'è. Ne derivano la sfiducia, il non alzare
mai la mano per chiedere aiuto, il non dire mai ``bel lavoro'' agli altri
perché non sia mai che qualcuno ti veda fare un complimento ``a quelli''.
Risolvere queste dinamiche è lungo e tedioso, e sempre pieno di ``non mi
ricordo'', ``io non c'ero'' e ``abbiamo sempre fatto così'', la cosa più
frustrante per chi cerca di risolverle.

% ---------------------------------------------------------------------
\section{Il principio del default to trust}
\label{sec:z1-default-trust}

\subsection{La definizione}

Il relatore indica questa come la slide più importante della presentazione.
\textbf{Default to trust} è uno dei valori aziendali di Riot Games, che lui ha
tradotto, adattato e ``esportato'':

\begin{quote}
  \emph{Approccio ogni situazione di conflitto partendo dal presupposto che una
  persona abbia agito in un certo modo perché, con le informazioni a sua
  disposizione in quel momento, ha pensato che fosse la soluzione migliore, e
  non perché sia un'incapace o volesse danneggiare il progetto
  intenzionalmente.}
\end{quote}

In altre parole: \textbf{di default, mi fido}.

\subsection{L'esempio del ritardo al daily stand-up}

A Riot il relatore aveva un \emph{daily stand-up} dal vivo alle 15 con il suo
team di venti persone. Una ragazza arrivava regolarmente alle 15:15, per una,
due, tre settimane. Alla quarta ci ha parlato. \emph{Come iniziereste la
conversazione?} Le proposte della classe:
\begin{itemize}
  \item sottolineare l'importanza di esserci dall'inizio, perché potrebbe
        perdersi qualcosa che le interessa, e per rispetto dei colleghi
        presenti;
  \item ``devi arrivare in orario'', spiegando perché è importante;
  \item \textbf{chiedere perché}.
\end{itemize}

Il relatore osserva, anche dalle sue consulenze fuori dai videogiochi, che
\textbf{più si sale con la seniority}, più la risposta è ``devi smettere di
arrivare in ritardo'', spesso con un \emph{or else} finale (``altrimenti
succede quest'altra cosa''). Più si scende con l'età, e quindi con meno
abitudini consolidate, più la prima domanda diventa ``dimmi perché arrivi in
ritardo''. Già spiegare perché è importante è ``miglia avanti'' rispetto al
semplice ``smetti''.

Partire dal perché gli ha permesso di scoprire la causa: il meeting precedente
della ragazza, in un'altra zona del campus a un quarto d'ora di distanza,
finiva sempre tardi, e lei doveva attraversare il campus di corsa. Non aveva
la confidenza per dire all'altro manager, che insisteva perché si restasse
fino alla fine, che doveva andare; sacrificava quindi il meeting del relatore,
sapendo che era il più comprensivo dei due.

Il relatore si è offerto di \textbf{mediare}: ha detto all'altro manager che
o anticipava l'inizio del suo meeting, così da finire comunque in tempo,
oppure lasciava uscire la ragazza, perché quel meeting gli serviva. C'è anche
una questione di \textbf{economia di scala}, spesso sottovalutata: se venti
persone aspettano qualcuno per dieci minuti, si sono sprecati dieci minuti per
venti. Per lo stesso motivo il workshop è iniziato alle 10:05: iniziare alle
10:10 avrebbe significato far aspettare cinquanta persone cinque minuti in
più. A Riot c'era una scrivania vuota con un pupazzo gonfiabile, \textbf{Mr.\
Late}: sommando tutti i minuti di ritardo a meeting, call e conversazioni, in
un anno veniva fuori praticamente \emph{una persona a tempo pieno}. Serviva a
ricordare che anche cinque minuti contano, e che bisogna rispettare gli altri
arrivando in orario o almeno avvisando.

\paragraph{Gli effetti.} La conversazione era avvenuta nel primissimo periodo a
Riot, quando il team non aveva ancora motivo di fidarsi di lui. Risolvendo il
problema ha eliminato anche tutto lo strato di micro-dinamiche spiacevoli che
un manager impara a vedere (due colleghi che si guardano con la risatina
quando lei entra in ritardo). La ragazza ha cominciato a entrare tranquilla,
con una certa leggerezza: non le faceva piacere arrivare in ritardo (``ci sono
anche persone che lo fanno apposta, e quello è un problema da affrontare
diversamente''). Soprattutto, con i colleghi non parlava di lui come del capo
che minacciava, ma come di chi le aveva risolto un problema: era diventato
un \textbf{alleato}. La settimana dopo altre tre persone sono andate da lui
con un problema da risolvere. Ripetuto cinque, sei, sette, dodici volte, in un
paio di mesi poteva dire quasi qualsiasi cosa nei meeting sapendo che
richieste e critiche sarebbero state percepite come ``me lo chiede perché vuole
darmi una mano'', non per fare bella figura o per presentare i successi del
team ai suoi superiori.

\subsection{Condividere i successi, schermare i fallimenti}

Un altro elemento fondamentale (ripreso nella seconda giornata): il manager
\textbf{condivide i successi e protegge il team dai fallimenti}. ``Se va male
qualcosa è colpa mia; se va bene qualcosa, voglio che ne parliate voi.'' Il
lavoro del manager non dovrebbe essere valutato su ciò che presenta del lavoro
del team: tutti sanno comunque chi è il manager. Tanto vale lasciar presentare
il team, che potrà dire con orgoglio di aver presentato a Mike
Morhaime\footnote{Cofondatore ed ex presidente di Blizzard; nella
trascrizione ``Mike Morayme''.} la propria campagna europea per
\emph{StarCraft~II}. Ci sono manager che dicono ``come mi prendo le colpe, mi
prendo i meriti'', ma i meriti li hanno già: è inutile ostentarli. Con i
fallimenti invece il manager è la \textbf{rete di sicurezza} del team. Il
relatore ne ha fatto una questione di principio: \textbf{non scaricare mai la
colpa su qualcun altro}.

\paragraph{Colpevole o causa.} Un esempio che il relatore attribuisce a Pirate
Software, creatore di contenuti ed ex dipendente Blizzard\footnote{Il relatore
aggiunge che si tratterebbe del figlio dello sviluppatore a cui si sarebbe
ispirato il personaggio del giocatore ``senza vita'' nell'episodio di
\emph{South Park} dedicato a \emph{World of Warcraft}. L'informazione è
riportata come detta in lezione e non è verificata.}. Contattando il servizio
clienti Blizzard per motivi personali, aveva scoperto che i colleghi gli
avevano scritto nei ticket informazioni confidenziali che non avrebbero dovuto
condividere in quel modo, nemmeno con un collega. Non ha segnalato le persone
da punire: ha detto ``c'è un problema, evidentemente le policy non sono
chiare''. Quando gli hanno chiesto i nomi, ha rifiutato: ``se vi dico chi è,
punite quelli, e gli altri cinquecento continueranno a farlo; non avremo
cambiato la causa alla fonte''.

\begin{quote}
  \emph{È inutile trovare il colpevole. Voglio trovare la causa. Perché il
  colpevole cambia.}
\end{quote}

Il colpevole a volte è vittima della situazione: per questo la definizione
dice \emph{``con le informazioni a sua disposizione in quel momento''}. Se il
relatore avesse rimproverato la ragazza in ritardo, non solo non avrebbe
risolto il problema, ma avrebbe creato un'\textbf{ingiustizia}: lei, che
attraversava il campus di corsa per una colpa non sua e di cui non poteva
parlare, sarebbe stata anche rimproverata.

Nessuno fa un brutto lavoro apposta: nessuno è contento quando sbaglia o
perde. Perché al lavoro diamo per scontato che gli altri sbaglino
intenzionalmente? Se qualcuno ha fatto una cosa male, è perché in quel momento
credeva fosse la scelta migliore. La domanda giusta non è ``perché hai fatto
questa cosa che è stata un fallimento?'', ma:

\begin{quote}
  \emph{Che cosa ti ha fatto pensare che questa fosse la scelta giusta?}
\end{quote}

Magari erano limiti di budget, di tempo, di scopo: si agisce su quelli, non
sulla persona. Se si cambia la persona ma le si danno gli stessi fattori, è
molto probabile che arrivi alle stesse conclusioni, e ci si troverà a dover
``disciplinare'' anche lei.

\paragraph{La reputazione costruita sui rumor.} Durante il giro di
presentazione in un ufficio (non in nessuna delle aziende citate), al relatore
è stato detto: ``questo è il dipartimento di design; se ti serve qualcosa
chiedi a lei, non a lui, perché lui consegna sempre in ritardo''. Oltre a
trovarlo poco professionale, una settimana dopo ha chiesto: ``cosa ti ha fatto
questa persona?'' ``A me niente, però ho sentito che\dots''. Si torna alle
scimmie: magari quella persona non consegna in ritardo da mesi, ma ormai lo
dicono tutti. È un circolo vizioso da fermare subito.

\paragraph{Niente escalation inutili.} Lo stesso vale per il feedback: quando
qualcuno va dal relatore dicendo ``ho un problema con questa persona'', lui
chiede: \textbf{``perché ne stai parlando a me e non a lei?''}. ``Perché sei il
mio capo'' non basta. Passare da lui significa dare per scontato che l'altra
persona non abbia interesse a risolvere, mentre magari basta parlarle. È la
strada più facile (``lo dico a Stefano, ci pensa Stefano''), e proprio per
questo bisogna mettere dei paletti: il manager interviene quando serve, ma le
cose che si possono gestire da soli vanno gestite da soli. (Quando invece il
confronto diretto fallisce, l'escalation è legittima: si veda la
sezione~\ref{sec:z1-qa-finale}.)

% ---------------------------------------------------------------------
\section{Traiettorie di carriera e valutazione delle performance}
\label{sec:z1-traiettorie}

\subsection{Quale traiettoria vorreste?}

Il relatore mostra un grafico con quattro curve colorate, contrassegnate
anche da lettere (A, B, C, D) per eventuali persone daltoniche (``a me è
capitato che qualcuno scoprisse di esserlo durante una lezione''). Le curve
rappresentano la performance lavorativa nel tempo. Due domande: \emph{quale
vorreste essere? E quale pensate realisticamente di essere?} Un partecipante
aspira a B e pensa di essere C; quasi tutti scelgono B.

Perché non A, che sale più rapidamente? Un partecipante risponde che una
crescita così repentina può essere più dannosa che utile, come una ``rush'' di
studio molto faticosa che poi non si riesce a sostenere. Il relatore propone
una lettura suggerita da un amico psicologo, che spiega la forma del grafico:
di B, C e D \textbf{vediamo dove finiscono}; di A no. Potrebbe continuare a
salire oltre la scala del grafico, o crollare sotto D. Non saperlo ci spinge a
preferire B. (A chi obietta che anche le altre proseguono oltre il bordo, il
relatore precisa: il bordo destro è ``oggi'', e di A non sappiamo come sia
arrivata fin lì.)

\subsection{La carriera non è una linea retta}

Tutti immaginano la carriera come una linea: si parte da X e, attraverso
stage e posizioni, si arriva a Y. In realtà carriera e performance avranno un
andamento irregolare, con alti e bassi. Come non si accetta che un team possa
temporaneamente \emph{underperformare}, tanto meno lo si accetta da se stessi.
Bisogna accettare che ci saranno momenti bassi.

Per questo il relatore \textbf{detesta le valutazioni annuali delle
performance}: prendono un punto casuale della traiettoria, che se si è
fortunati è il picco più alto e se si è sfortunati il più basso, e lo usano
per valutare un intero anno. Insegna ai manager a fare più punti di controllo
e a guardare altre cose:
\begin{enumerate}
  \item \textbf{quanto tempo passa tra un picco e l'altro}, cioè quanto
        \emph{cooldown} serve a una persona. È il motivo per cui, dopo la
        Games Week, tutti staccavano una o due settimane. Il relatore stesso,
        che si impegna moltissimo e tende al burnout, non potrebbe rifare lo
        stesso workshop il lunedì e il martedì successivi con la stessa
        qualità;
  \item \textbf{la differenza tra la miglior performance e la peggiore}: che
        cosa ha portato a lavorare così bene, e così male? Con chi si
        lavorava, su cosa ci si concentrava, quanto tempo e budget c'erano?
        L'obiettivo è capire come \emph{replicare} le condizioni che fanno
        lavorare bene, invece del solito ``hai lavorato male, è colpa tua''.
        Il relatore ha avuto molti manager per cui quando lavorava male era
        colpa sua e quando lavorava bene ``che bravo'', ma sapeva che raramente
        era l'unico fattore della propria performance. È la fallacia del ``se
        sai fare 100, perché non fai sempre 100?'', come nella scena di
        \emph{Schindler's List} sull'operaio che produce una vite in pochi
        secondi: non sempre un ritmo è sostenibile;
  \item \textbf{il fallimento di oggi confrontato con quello di sei mesi fa}.
        Il relatore non è contento della presentazione di oggi (l'avrebbe
        migliorata per altri due anni), ma rispetto a quelle di cinque anni fa
        vede un grande miglioramento di chiarezza: il suo peggior fallimento di
        oggi equivale al suo primo successo di inizio carriera.
\end{enumerate}
Quello che conta è la \textbf{traiettoria} dall'inizio dell'anno a oggi, non
i singoli punti. E se si riesce a capire insieme quando e perché si ha
successo e quando e perché si fallisce, si può lavorare per replicare solo le
situazioni che fanno lavorare bene.

\paragraph{Lo stesso vale per i giochi.} È il motivo principale per cui i
giochi delle grandi case vanno bene o male. Quando si riescono a replicare le
condizioni che hanno prodotto \emph{Warcraft~II}, esce \emph{Warcraft~III};
quando non ci si riesce, esce ``il nuovo \emph{Dragon Age}'': la gente non è
arrabbiata perché il gioco non è bello, ma perché non è come gli altri
\emph{Dragon Age}; non ne replica il \emph{feeling}, anche se potrebbe
essere un gioco favoloso. L'eccezione positiva è il passaggio da \emph{Dune} a
\emph{Dune~II}: il primo era di fatto un'avventura grafica, il secondo ha
generato il genere degli strategici in tempo reale. Gli sviluppatori avevano
capito che il successo non stava nel tipo di gioco, ma nell'ambientazione che
attirava le persone: un bel gioco in quell'ambientazione sarebbe piaciuto
comunque. Il relatore comprò \emph{Dune~II} aspettandosi un seguito
dell'avventura e, pur sorpreso, lo finì.

\subsection{Lying, hiding, faking}

Quando si è valutati su un singolo punto nel tempo, e da quel punto dipendono
bonus, stipendio e il rischio di licenziamento, si comincia a fare
\textbf{lying, hiding e faking}:
\begin{description}
  \item[Lying] dire che è sempre andato tutto bene, che si finirà sempre in
        orario, che non c'è nessun problema;
  \item[Hiding] nascondere: ``non era un mio task'', ``non c'ero'', ``non ho
        visto niente'';
  \item[Faking] fingere, anche in cose apparentemente innocue come ``sì, sì,
        ho capito tutto'' quando si ha solo paura di dire ``me lo rispieghi?
        Non ho capito niente''.
\end{description}

\subsection{``Non ho capito'' e la conferma delle istruzioni}

Dire ``non ho capito, ripetimelo'' è ancora visto come una figura da sciocco.
Il relatore non si arrabbierà mai se qualcuno gli chiede di rispiegare
un'intera slide: le cose ovvie per lui possono non esserlo per altri. Più si
conosce qualcosa, \textbf{più è difficile spiegarla}, perché si danno per
scontate molte cose. Ne parlava proprio con Diego, che gli diceva che fare
l'insegnante diventa sempre più difficile man mano che si diventa esperti. È
come spiegare \emph{EVE Online}, uno dei giochi più complessi in circolazione:
chi lo conosce bene dà per scontato che ``ovviamente non puoi fare questa
cosa'', mentre per un nuovo giocatore non è affatto ovvio. A scuola chi dice
``non ho capito'' viene visto come lo ``scemo''; nessuno si chiede ``non è che
l'ho spiegato male?''.

Il racconto più ``traumatico'' dei ragazzi che fanno stage nelle grandi
aziende è sempre lo stesso: ``Cinzia, devi fare questa cosa.'' Cinzia la fa.
``Non te l'avevo chiesta così.'' Il problema è che nessuno ha
\textbf{confermato che Cinzia avesse capito}. Il consiglio, valido sempre:
\begin{itemize}
  \item quando si \emph{ricevono} istruzioni, ripeterle: ``ho capito bene?
        Vuoi che faccia questa cosa in questo modo?'';
  \item quando si \emph{danno} istruzioni, chiedere a chi le riceve di
        ripeterle, non perché lo si ritenga incapace, ma per capire che cosa
        ha capito dalla propria spiegazione (``quindi vuoi un quadrato?''
        ``No, ho detto un triangolo\dots'' ``Però mi hai detto una forma con
        quattro angoli.'' ``Hai ragione: voglio un quadrato'').
\end{itemize}
Frasi come ``non te lo ripeto due volte'' o ``stammi dietro'' portano solo a
fare le cose come si pensa che vadano fatte e a sentirsi dire a fine progetto
``non la volevo così'': che si tratti di una feature, di un design o di un
modello 3D. (Un partecipante osserva che lo stesso accade insegnando le lingue:
chi insegna dà per scontato, ad esempio, che si sappia dell'esistenza del
genere neutro in tedesco.)

\paragraph{Come convincere il team a chiedere?} Il relatore mette sempre come
premessa che \textbf{l'obiettivo non è trovare la persona più intelligente
nella stanza, ma fare un buon lavoro}. Che Gianni capisca alla prima e un altro
alla quinta non gli interessa, se alla fine entrambi lavorano bene: ripetere
cinque volte è comunque meno lavoro che far rifare le cose sette volte. È un
investimento a lungo termine, e non dissimula alcun fastidio: davvero non gli
dà fastidio. Ma ripetere le stesse parole non aiuta: ``non ho sentito'' si
risolve ripetendo; ``non ho capito'' si risolve cambiando spiegazione. Meglio
ancora: ``dimmi che cosa hai capito tu, e vediamo cosa non torna''. Se
intravedesse nel team uno sguardo del tipo ``è la quinta volta che lo
chiede'', interverrebbe: spesso chi giudica è proprio chi dice ``ho capito'' e
poi sbaglia. Preferisce di gran lunga chi chiede cinque volte.

\subsection{Il feedback più utile che abbiate ricevuto}

\emph{Qual è il feedback più utile che abbiate mai ricevuto?} Il relatore ne ha
almeno dieci che gli hanno cambiato il modo di lavorare. Le risposte:
\begin{itemize}
  \item un semplice \textbf{``grazie''} per un lavoro fatto, più utile di un
        ``questa cosa fa schifo'';
  \item un feedback sulla propria \textbf{impulsività} nel rispondere: fermarsi
        un attimo e capire di cosa si sta parlando prima di intervenire;
  \item \textbf{essere se stessi}, anche al lavoro, e non scusarsi
        continuamente per ogni cosa.
\end{itemize}
I feedback che toccano di più sono quelli \textbf{veri}. Il ``grazie''
significa ``ho visto quello che hai fatto e l'ho apprezzato''. Quelli negativi
che si ricordano sono quelli che hanno fatto realizzare ``questa cosa che
pensavo di saper fare, non la so fare'', oppure ``la so fare, ma questa volta
non mi è venuta'' (non è che se una cosa una volta non riesce non la si sa più
fare: può capitare di sbagliare una porta). Arrivano sempre da un
\textbf{punto cieco}: qualcosa che non si sapeva fosse percepito dagli altri.
Conta più ciò che le persone ti dicono di ciò che pensi che pensino di te; e
ciò che il relatore pensa che gli altri pensino di lui è sempre terribile,
mentre il loro giudizio si rivela molto più positivo del previsto.

\paragraph{Il denominatore comune siete voi.} Si potrebbe pensare che questi
concetti non siano applicabili da chi arriva in un'azienda senza poterne
cambiare la cultura. Ma c'è sempre un denominatore comune a cui applicarli:
\emph{se stessi}. La cultura aziendale può non spingere a dare feedback, può
vedere male chi chiede aiuto; ma il relatore ha visto che quando lui lo fa, le
persone intorno tendono a farlo, e un gruppo che lo fa lavora meglio. Perfino
i CEO più ostici a un certo punto gli chiedevano: ``perché il tuo team ora
performa meglio degli altri?''. La differenza è il default to trust: una
pratica che col tempo fa operare il team \emph{come un team} e non come un
insieme di individui. Un videogioco è un progetto corale: c'è chi fa giochi da
solo, ma non si fa da soli \emph{Red Dead Redemption~2}; e se il gruppo non è
affiatato e allineato non esce \emph{Red Dead Redemption~2}, esce
\emph{Anthem} (``oggi la mia fissazione è \emph{Anthem}, domani sarà
un'altra'').

% ---------------------------------------------------------------------
\section{Come dare un feedback}
\label{sec:z1-dare}

\subsection{Concentrarsi sull'azione e sui suoi effetti}

Il primo principio: \textbf{il feedback non riguarda mai la persona}. Non si
dice ``tu hai fatto questa cosa che non va bene'', ma ``questa cosa non va
bene: come facciamo a non farla più?''. Che sia un comportamento, un'azione,
una decisione o una strategia, ci si concentra su quello, e si
\textbf{esplicita l'effetto} che ha. Nell'esempio del ritardo l'effetto era
``quindici minuti di attesa moltiplicati per venti persone''. In un progetto
può essere una catena:
\begin{quote}
  se non rispetti le deadline, l'animazione che viene dopo resta ferma; se
  loro sono in ritardo, la \emph{vertical slice} è in ritardo; se la vertical
  slice non è pronta per la Gamescom, non andiamo alla Gamescom e non troviamo
  un publisher; se non troviamo un publisher, chiudiamo.
\end{quote}
Più la conseguenza è grave, più dovrebbe sembrare piccolo il gesto di dire
``datemi una mano'': è il \emph{male minore}. Poi, a cose fatte, ci si dà
feedback: cosa ha funzionato, cosa no, cosa fare di più, di meno, cosa smettere
di fare.

\subsection{Il pericolo degli ambienti solo positivi}

Gli ambienti dove ci si fa sempre grandi complimenti e nessuno parla dei
problemi creano difficoltà vere alle aziende. Il problema non è mai la
mancanza di complimenti, ma un \textbf{eccesso di ottimismo}: si comincia a
pensare che le cose andranno come si spera, perché tutti intorno danno solo
feedback positivi.

A Blizzard i feedback determinavano gli \textbf{aumenti annuali}: ognuno
doveva indicare cinque persone che avrebbero dato feedback su di lui, e
ovviamente tutti indicavano gli amici, o colleghi di progetti andati bene. Un
sistema nato con le migliori intenzioni diventava uno \emph{scambio di voti}
(``se mi metti follow ti metto follow''), senza alcun miglioramento.

Al contrario, \textbf{le conversazioni più brutte sono state le più utili}:
affrontare i licenziamenti a Riot, o la cancellazione di \emph{Project Titan}
a Blizzard, il progetto che doveva essere il nuovo MMO dopo \emph{World of
Warcraft} e da cui poi è nato \emph{Overwatch}\footnote{Nella trascrizione si
legge anche ``doveva essere Warcraft 2'', probabilmente un lapsus o un errore
di trascrizione.}. Andando a vedere perché lo si cancellava, c'erano solo
motivazioni razionali e pragmatiche, di business, chiare a tutti, non un ``Mike
Morhaime ci odia''. Sono anche le conversazioni in cui ha lavorato su se
stesso: è molto \textbf{avverso ai conflitti}, e questo a volte gli ha
procurato conseguenze molto più gravi di quanto sarebbe costato un piccolo
disagio. Dieci anni fa faceva fatica ad accettare questo difetto; ha dovuto
subire le conseguenze di conflitti evitati e poi cresciuti.

\subsection{Che cosa fare}

Le linee guida della slide, che il relatore legge rapidamente perché saranno
disponibili nel materiale:
\begin{description}
  \item[Attenzione al linguaggio del corpo] dare un feedback a braccia
        incrociate o guardando altrove è diverso dal guardare l'altro negli
        occhi.
  \item[Sii specifico] con esempi tangibili.
  \item[Esplicita gli effetti] come visto sopra.
  \item[Discuti il prima possibile] non aspettare un mese per parlare di
        qualcosa successo ieri, altrimenti nessuno ricorderà più bene cosa sia
        accaduto.
  \item[Concentrati sull'azione e sul problema, non sulla persona] il nemico è
        l'azione, non la persona che si ha davanti. Non bisogna far cambiare
        l'altro: bisogna evitare che il problema si ripresenti.
\end{description}

\paragraph{Gli insulti sono rivolti al ruolo.} Un esempio solo apparentemente
scollegato. Il relatore ha lavorato su \emph{League of Legends}, ``non
esattamente noto per avere una community da orsetti del cuore'', in varie
\emph{task force} sulla tossicità. A chi smetteva di giocare veniva inviato un
sondaggio sul perché; il 33\% indicava la tossicità della community. Dopo varie
campagne di sensibilizzazione e cambiamenti si è scesi al 30\%: sembra solo
un 3\% in meno, ma valeva circa \textbf{120 milioni di revenue in più
all'anno}, perché quei giocatori restavano esposti a skin e contenuti e
continuavano a comprare. Anche qui, piccoli guadagni che si sommano.

Lavorando con degli psicologi, aveva notato che gli insulti \textbf{non sono
quasi mai rivolti alla persona, ma al ruolo}: ``il mio jungler fa schifo'',
``il mio DPS non è capace'', ``il nostro portiere è una pippa''. Chiamare
qualcuno per ruolo lo \emph{depersonalizza} e rende più facile insultarlo. Lo
stesso nell'esempio del giro dell'ufficio: non ``non chiedere a Giulio'', ma
``non chiedere a quel grafico, consegna sempre in ritardo''. Concentrarsi
sull'azione, senza attaccare la persona, fa sì che il feedback non sia (e non
venga percepito come) un attacco personale.

\paragraph{Considera il carattere.} Qui tornano utili i test di personalità.
Servono più agli studenti che al relatore (che sa già che usciranno 53 profili
diversi, senza un pattern): bisogna conoscere il proprio profilo e
\textbf{condividerlo con chi lavora con noi}, chiedendo anche agli altri di
fare lo stesso. In ufficio il relatore ha usato una versione più avanzata, a
pagamento, con un report di circa trenta pagine: il \textbf{profilo
DISC}\footnote{Nella trascrizione ``disk''. Si tratta del modello DISC
(Dominance, Influence, Steadiness, Conscientiousness).}. Ha mappato tutta
l'azienda e ha fatto stampare il profilo di ciascuno sulla sua scrivania, così
che tutti sapessero con chi stavano parlando. Capire come ragionano gli altri
(``è introverso, non gliene parlo davanti a tutti'') è una delle cose che gli
sono tornate più utili: \emph{adattare ciò che si vuole dire a chi si ha
davanti}.

\paragraph{Assicurati di non disturbare.} Sotto il profilo di ognuno c'era un
piccolo \textbf{semaforo}: verde, ``parlami pure, interrompimi''; giallo,
``interrompimi solo se è urgente''; rosso, ``mi sto concentrando, non mi
parlare''.

\paragraph{Lascia finire il lavoro.} Durante gli eventi dal vivo il relatore si
era reso conto di quanto fosse grave dare feedback mentre qualcuno sta ancora
lavorando. Se una cosa sta andando male, dirlo in quel momento non aiuta: tra
la pressione di finire e quella del feedback, peggiorano entrambe le cose, e
anche il rapporto. Meglio lasciarla finire, anche male, e parlarne dopo.

\subsection{Che cosa non fare}

\begin{description}
  \item[Non basarti su supposizioni] ``mi hanno detto che ieri in un meeting
        hai fatto\dots'': o c'eri e l'hai visto, o non esiste.
  \item[Non giudicare la persona.]
  \item[Non fornire feedback vago.]
  \item[Non accettare e non fornire feedback di terze parti] Se qualcuno vi
        dicesse ``mi hanno detto che hai fatto un brutto lavoro'', il primo
        pensiero non sarebbe ``cosa ho fatto male?'', ma \emph{``chi te l'ha
        detto? Chi è la vipera che ho in seno? Chi è l'impostore?''}. Siamo
        fatti così. Ogni conversazione che inizia con ``mi hanno detto che'' è
        inutile: l'altro non ascolterà, penserà solo a chi dare la caccia, e
        spesso penserà alla persona sbagliata, creando antipatie tra persone
        che non c'entravano nulla. Quando al relatore dicono ``mi hanno detto
        che lui ha fatto\dots'', risponde ``vai a parlarne con lui''; se gli
        dicono ``mi hanno detto che tu hai fatto\dots'', risponde ``mandami la
        persona che te l'ha detto, altrimenti non ne parliamo''. Non per
        cattiveria: se chiedesse ``che cosa ho detto di preciso a Marco?'' e la
        risposta fosse ``non lo so, non c'ero'', di cosa si starebbe parlando?
        \textbf{Parlate con la fonte}, altrimenti fate come se quel feedback non
        esistesse.
  \item[Non dare consigli non richiesti] come visto, non portare la soluzione:
        la soluzione deve venire da chi riceve il feedback.
  \item[``Non mi piace'' non è un feedback] ``Penso che questo design non
        funzioni perché, visto che il nostro obiettivo è\dots'' ha argomenti;
        ``non mi piace il combat system'' è un'opinione. Anche ``io l'avrei
        fatto diversamente'' non è un feedback.
  \item[Accettare non vuol dire fare come proponi] deve essere chiaro a chi dà
        il feedback: l'aspettativa è essere ascoltati e presi in
        considerazione, ma se la responsabilità è dell'altro, deciderà
        l'altro.
  \item[Non offenderti] se il feedback non viene seguito. Restarci male
        significa pensare di saperne più dell'altro, cioè \emph{venire meno al
        default to trust}. ``Decidi tu'' significa ``mi fido di te, perché sei
        lo specialista, e mi fido che mi ascolterai e farai la cosa migliore
        per tutti'', che potrebbe anche non essere ciò che ho detto io. Se
        voglio che l'altro faccia sempre come dico io, mi sto fidando solo di
        me stesso.
\end{description}

\subsection{Il primo feedback e il debriefing}

La cosa più difficile è dare il \textbf{primo} feedback: è un muscolo da
allenare. Chi non l'ha mai fatto (soprattutto con questa struttura) dovrebbe
provarci, anche su minuzie (``abbiamo concordato di usare WhatsApp nel gruppo:
non potremmo usare Telegram?''), anche solo per vedere come va. Altre
indicazioni:
\begin{itemize}
  \item feedback \textbf{solo uno a uno}, mai davanti ad altri, soprattutto se
        negativo;
  \item \textbf{prepararsi} prima;
  \item fare il \textbf{debriefing del feedback}: dopo, fermarsi a pensare a
        com'è andato. Il relatore, quando media un feedback tra due persone,
        si sente spesso rispondere ``benissimo'' da una e ``malissimo''
        dall'altra: la percezione è molto diversa. Non si fa problemi a tornare
        dalla stessa persona: ``ieri ti ho dato un feedback: com'è andata? Sei
        soddisfatto? Come poteva andare meglio?''. Con le persone con cui si ha
        meno a che fare, questo aiuta molto la volta successiva.
\end{itemize}

% ---------------------------------------------------------------------
\section{Come ricevere un feedback}
\label{sec:z1-ricevere}

Per il relatore è la parte più difficile. Si riassume in quattro passi:
\begin{enumerate}
  \item \textbf{Ascolta}: si parte sempre ascoltando.
  \item \textbf{Chiedi dettagli}, se serve, ma per capire meglio, non per
        difenderti.
  \item \textbf{Accetta, a prescindere}: non si sta decidendo chi ha ragione;
        si sta ascoltando un'opinione. Non c'è nulla da cui difendersi.
  \item \textbf{Agisci}: fatti un piano. Può essere banale (``mi hanno detto
        che interrompo sempre: cerco di non interrompere''), oppure ``devo
        migliorare questa cosa: come?''. Pensa anche a \emph{chi} potrebbe
        aiutarti e a \emph{che supporto} ti serve: il piano non deve essere
        per forza ``devo risolvere tutto da solo''.
\end{enumerate}
Una domanda che molti sottovalutano: \emph{``sono junior in un'azienda: sarò
mica il primo ad avere questo problema?''}. Probabilmente qualcun altro l'ha
già avuto: basta chiedere ``tu come hai fatto?''.

\subsection{Il feedback positivo}

Tutto questo vale per i feedback costruttivi negativi, a cui si pensa di
default; ma esistono anche i \textbf{feedback positivi}, e vanno usati. Quando
si riceve un ``grazie'' o un ``bravo'', la domanda è: \emph{come faccio a
ricreare le stesse condizioni?}

Il relatore ha lavorato anche in aziende dove il feedback era solo negativo.
Il rischio? Si immagini l'impiegato ideale: vuole sempre migliorare, è sempre
pronto a dire ``magari sbaglio io''. Se nessuno gli dice cosa sta facendo
\emph{bene}, spinto dal desiderio di migliorare rischia di cambiare cose che
erano già buone, e di farle peggio. Il rinforzo positivo \textbf{non è solo una
gentilezza}: è un atto pratico, che evita che le persone peggiorino ciò che
sanno già fare bene.

\subsection{Il feedback accidentale}

Esiste infine il \textbf{feedback accidentale}, molto frequente negli uffici:
\begin{itemize}
  \item il ``chiedi a lei, non a quel grafico'' del giro dell'ufficio è un
        feedback negativo su una persona, dato accidentalmente, né richiesto
        né voluto;
  \item a Riot li chiamavano \textbf{drive-by}, come nelle scene dei film in
        cui passa un'auto, si abbassa il finestrino, si spara e si riparte. In
        molte aziende con una ``cultura del feedback'' in cui tutti devono
        dire la loro, si riceve una valanga di commenti del tipo ``comunque il
        sito a me non piace, ciao''. Non si sa perché, non si sa cosa fare di
        diverso, non si sa che rischi ci siano: \emph{a che cosa è servita
        questa conversazione?} A niente.
\end{itemize}

% ---------------------------------------------------------------------
\section{Quando, quanto e con chi}
\label{sec:z1-quando}

\subsection{Frequentemente e recentemente}

\emph{Quando avete ricevuto il vostro ultimo feedback?} Per gli studenti,
qualche giorno fa (nel corso) o qualche mese fa (al lavoro, dopo una
consegna). Come aveva osservato Silvia, il feedback andrebbe dato
spesso.

Il relatore usa l'esempio di una nave da crociera (``la Costa Concordia'') che
deve andare da A a B. La rotta ideale (azzurra) è quasi impossibile: nessun
progetto va esattamente come previsto. C'è la rotta ``panoramica'' (verde), che
gira intorno all'isola d'Elba; ci sono poi rotte (arancione e rosa) che a un
certo punto si allontanano molto e poi correggono bruscamente: qualcuno si è
accorto che si stava sbagliando, ``ma non è quello che avevo chiesto'',
``non è dove dovevamo andare''. Sono i famosi \textbf{sanity check} o
\emph{checkpoint}.

L'ideale è avere \textbf{tanti piccoli check per fare tante micro-correzioni}:
correggere la rotta di un paio di gradi è molto più rapido che farlo dopo
essersi allontanati. Bastano trenta secondi di ``ehi, vado bene?''. Spesso non
lo si fa per paura di sembrare insicuri o poco autonomi, ma per il relatore è
l'opposto: quando qualcuno gli chiede ``vado bene così?'' e lui risponde
``sì'', si alleggerisce anche del proprio carico mentale. Un capo che assegna
un compito e lo vede solo alla fine vive sempre col dubbio (``avranno capito?
Lo staranno facendo giusto? Arriverà in tempo?''); con un check di dieci
secondi al giorno (``tutto ok, Stefano'') è tranquillo, e la persona va avanti
sapendo che quanto fatto finora va bene. Non serve fare ogni giorno il giro del
dipartimento: anche ogni uno o due giorni basta a evitare la rotta verde.

Per il feedback è lo stesso: \textbf{va dato subito}. Dire oggi ``ieri hai
fatto questa cosa'' è molto diverso da dire ``negli ultimi sei mesi hai fatto
questa cosa''. Il relatore si è sentito dire più volte, alla valutazione di
dicembre, ``da maggio hai fatto questa cosa che non andava''. La sua risposta:
``e aspetti dicembre per dirmelo? Mi hai lasciato sbagliare per sei mesi?''. È
una mancanza di rispetto, e la conversazione non può più essere costruttiva.
Per questo non ama le valutazioni annuali: i manager crescono abituandosi a
prendere appunti delle cose negative per tirarle fuori alla valutazione
(``a marzo hai fatto questa cosa'' ``sì, però poi ho smesso'' ``sì, però a
marzo l'hai fatta'').

In sintesi:
\begin{description}
  \item[Frequentemente] riguarda l'intervallo tra un check e l'altro.
  \item[Recentemente] significa non lasciar passare un anno prima di dare il
        feedback, altrimenti non serve più a niente.
\end{description}

\subsection{Quanto dura un feedback}

Un feedback dura \textbf{da trenta secondi a quindici minuti}; raramente di
più. Chi non ne ha mai dato uno dovrebbe sforzarsi di darne almeno uno, magari
positivo per cominciare. Ogni feedback dato o ricevuto è un
\textbf{deposito nel proprio ``fondo fiduciario''} (\emph{trust fund}), che
costruisce la reputazione futura. Il relatore si è costruito la reputazione di
chi dà feedback il più possibile oggettivi e analitici: giusti o sbagliati che
siano, costringono a pensare. Non impone mai il suo punto di vista (``si fa
come dico io''), ma chiede ``tu come faresti? Come posso aiutarti?''. Il
feedback, per lui, serve a \emph{far crescere qualcun altro}.

\subsection{I manager peggiori: ``a che numero sto pensando?''}

I peggiori manager che il relatore abbia avuto erano quelli che volevano che si
facesse come volevano loro, anche quando chiedevano ``tu come la faresti?''.
In realtà la domanda era ``\emph{a che numero sto pensando?}'': se dicevi 5,
``no, dovremmo fare 9''; se dicevi 9, ``no, 5''. Lo stress costante di dover
indovinare cosa vuole il capo nasce dalla mancanza di fiducia (``se mi chiedi
come fare, fammelo fare'').

Lo stesso gli succede con i CEO che lo chiamano per ``sistemare la cultura
aziendale'' e poi dicono ``non sono sicuro, io farei diversamente''. La sua
risposta: ``lo so che faresti diversamente, per questo mi hai chiamato''. Se si
chiama un esperto, bisogna lasciargli fare il suo lavoro: non si va dal
meccanico a cambiare le gomme per poi dirgli ``non dovresti girare un po' di
più quel bullone?''. Chiedere un parere a un esperto significa dargli anche il
potere di scegliere. Il problema è il \emph{micromanagement}, o in generale il
non accettare punti di vista diversi dal proprio (magari perché si è il
fondatore e si pensa di essere l'unico a saper fare le cose).

\subsection{Dare feedback a un superiore egocentrico}

\paragraph{Come si dà un feedback negativo a un superiore con un grande ego?}
Il trucco è \textbf{non puntare sull'ego, ma sull'azione e sui suoi effetti}.
Non ``tu fai così e io mi sento così'': anzi, si alimenta il suo ego
rendendolo il \emph{main character} dell'avventura, l'eroe che risolve il
problema (``tu, grande manager, aiutaci a risolvere questo problema che
affligge noi che facciamo l'operatività''). Non è l'ideale (``io cambierei
manager''), ma pochissimi possono scegliersi i colleghi: il relatore stesso ha
quasi sempre ricevuto team preassemblati (``ecco le chiavi, questa è la tua
macchina''). Bisogna giocare con le carte che si hanno in mano.

\paragraph{Un esempio pratico.} Con un manager che fa micromanagement e mette
sempre ``il becco'' nel proprio lavoro, si può dire: ``se ti fidi di me per
questo lavoro, con le tue capacità potresti dedicarti a questa cosa più
strategica e di lungo termine: pensa quanto sarebbe più utile per l'azienda''.
In sostanza: ``guarda cosa potresti raggiungere se mi lasciassi fare il mio
lavoro''. Non come un ``tu fai il tuo e io faccio il mio'', ma come un
``ti sollevo da questa parte del carico, che so che ritieni importante: lascia
che ci pensi io, così puoi dedicarti ad altro'', che magari gli dà anche più
visibilità.

Si gioca su ciò che quella persona pensa di saper fare, e che probabilmente sa
anche fare. Anche Michael Scott di \emph{The Office}, per quanto preso in giro,
è raccontato come un venditore bravissimo, che ha chiuso i contratti più
importanti per l'azienda, e rispetto ai manager che gli succedono è forse
``uno dei meno peggiori''. Tutti, anche i manager peggiori, hanno punti di
forza: \textbf{giocate su quelli}.

% ---------------------------------------------------------------------
\section{Il feedback come regalo}
\label{sec:z1-regalo}

Per chiudere, il relatore risponde alla domanda d'apertura, che aveva lasciato
apposta in sospeso: che cos'è il feedback per lui?

\begin{quote}
  \emph{Il feedback è un regalo che mi viene fatto da qualcuno che so che vuole
  il mio interesse.}
\end{quote}

Ha costruito un team di persone che si fidano di lui e che non hanno motivo di
sabotarlo o di ferirlo: per questo interpreta anche le cose peggiori con il
default to trust. È un regalo, e cerca di prenderlo in considerazione il più
possibile.

Ci è arrivato partendo da se stesso. Oltre a essere avverso ai conflitti, si
definisce ``avverso alla comunicazione umana'': paradossalmente, una persona
introversa e timida finita a fare comunicazione, marketing e ora HR (``succede
spesso''). Prima di ordinare una pizza al telefono ha dieci minuti di
preparazione mentale, e se l'ordine delle domande cambia va in confusione. Per
lui un feedback di dieci minuti sono spesso \textbf{tre ore di preparazione}:
costruisce un \textbf{albero dei dialoghi}, come in un'avventura grafica
personale, con tutte le opzioni. Poi spesso finisce con ``hai fatto questa
cosa, andrebbe migliorata'' ``hai ragione, grazie'', e basta.

Così ha pensato: chissà quanti si fanno gli stessi problemi con lui, e passano
ore a prepararsi per dirgli ``questa cosa l'hai fatta male'', sapendo che
riceveranno comunque quell'occhiata fugace del ``come ti permetti''. Eppure
vengono a dirglielo. Non può far altro che rispondere: ``grazie, hai ragione,
cercherò di fare meglio''. È così che ha metabolizzato il ricevere feedback:
capendo che non c'è nessuna cattiva intenzione e che, se non proprio felice
(``ci ho messo una pietra sopra''), può almeno essere \emph{grato} che ci sia
ancora qualcuno che ha così a cuore il suo interesse da dargli feedback.

\paragraph{Il silenzio fa più paura.} Non bisogna spaventarsi dei feedback
negativi che emergono dai sondaggi aziendali: bisogna spaventarsi del
\textbf{silenzio}, quando nessuno dice più nulla. Significa che nessuno si
fida più che si sappiano risolvere i problemi.

\paragraph{La sindrome dell'impostore.} Il suo motivatore personale, all'inizio,
era soprattutto ``non voglio fallire'': per questo cercava di migliorare, e
migliorare è diventato accettare feedback, perché più feedback riceve, più
migliora. Il feedback è un \textbf{acceleratore}, e lui cerca di incentivare
gli altri a darglielo. Resta però sempre la vocina del ``e se si accorgono che
non sei capace?''. Con la \emph{impostor syndrome} convive da anni e
continuerà a conviverci; se ce l'avete, conviveteci anche voi. Avrete sempre
la sensazione che ``se ne accorgeranno'': forse sì, ma sapete fare anche
altre cose. Non concentratevi solo su ciò che non sapete fare. Gli errori
capitano: la perfezione non è di questo mondo.

\paragraph{Esercitarsi in accademia.} Per fare pratica con il feedback,
l'ambiente migliore è quello dei colleghi di corso: persone nella stessa
situazione, con esperienza simile e lo stesso obiettivo, in un contesto in cui
non si può essere licenziati e l'unico effetto a lungo termine è ``ho parlato
con questa persona, ora le cose vanno meglio''. E bisogna avere il coraggio di
parlare anche con i docenti (``dire a Diego: secondo me questa lezione potevi
farla in questo modo''; ``so che è una persona che lo accetta''), senza farsi
bloccare dalla gerarchia. Nell'azienda giusta (come quella in cui lavora ora) il
relatore va sempre dal CEO a dare feedback, sapendo che non verrà licenziato:
il CEO gli ha detto ``ti ho preso per darmi feedback''.

La frase di chiusura:

\begin{quote}
  \emph{Per superare le vostre difficoltà, tutto quello che vi viene chiesto è
  il coraggio di chiedere aiuto.}
\end{quote}

Chiedere aiuto non è una debolezza: è un segno di consapevolezza di sé
(\emph{self-awareness}) e di forza personale. Se tutti lo facessero, i
progetti andrebbero molto meglio; spesso invece ci si fa frenare dal ``non
voglio fare brutta figura'', ``non voglio passare per quello scemo'', e sono
proprio queste paure a rallentarci e a impedirci di raggiungere gli obiettivi.

% ---------------------------------------------------------------------
\section{Domande e risposte finali}
\label{sec:z1-qa-finale}

L'ultima mezz'ora è dedicata a domande libere.

\paragraph{E se una persona peggiora per uno o due anni di fila?} Se nessuno
gliel'ha mai detto, c'è un grosso problema, non suo ma dell'azienda. Se invece
c'è consapevolezza, il percorso è: ``cosa vuoi fare per risolvere?''; ``ho
provato questo, questo e questo''; ``visto che non ha funzionato, vuoi un mio
consiglio? Proviamo con quello?''; ``che supporto ti serve?''. Se, con tutto il
supporto, il tempo e le soluzioni concordate per mesi, il miglioramento non
arriva, per il relatore la soluzione è: \textbf{``dimmi un prezzo per andare
via''}. Preferisce pagare una buonuscita (\emph{severance package}) che
avviare la strada dei richiami formali sulla performance fino al
licenziamento per giusta causa. A un certo punto andare avanti non è un favore
per nessuno.

\paragraph{È vero che nei videogiochi si viene sfruttati? Meglio un piccolo
studio italiano?} Per il relatore si è quasi più a rischio a Blizzard che in un
piccolo studio di Roma o Torino. Nel piccolo studio si sa dove ci si va a
cacciare: non ci sarà il team building a Santo Domingo, ma ci si concentra su
un gioco alla volta e, se va bene, si fa il successivo. A Blizzard si può
lavorare su \emph{Overwatch}, che vende milioni, e perdere comunque il posto,
in modo totalmente scollegato dal proprio lavoro. Gli stipendi più alti
all'estero compensano anche questo rischio, che spesso non dipende da come va
il gioco (anzi, ``il gioco non ha fatto come volevamo'' è spesso una scusa per
licenziamenti che non avevano motivo).

\paragraph{Cambiare settore a trent'anni, da web developer full remote?} Lo
stipendio indicato dallo studente non è fuori dalla portata di vari studi
italiani, anzi è superabile. Non va esclusa la possibilità di lavorare da
remoto per aziende piccole ma estere. Alcuni amici del relatore, pagati con
stipendi italiani, si sono trasferiti a Bali o in paesi con un costo della vita
più basso; il collega di \emph{PUBG} percepisce uno stipendio olandese vivendo
a Milano, che per quanto cara non costa come i Paesi Bassi.

\paragraph{Come si diventa manager?} Dipende dall'azienda: alcune promuovono i
migliori (processo tendenzialmente sbagliato), altre propongono un percorso di
\emph{people management}. Diventare manager significa \textbf{abbandonare ciò
che si sa fare per imparare a fare altro}. Il relatore ha dovuto accettare che
non sarebbe più stato ``il miglior community manager italiano del mondo''
(``eravamo in pochi, quindi era facile''): dedicandosi alle persone non poteva
più fare il suo lavoro bene come prima. È stata una negoziazione con se
stesso: fare il manager significa non essere più l'esperto di qualcosa, ma
saper gestire e far crescere persone. Significa anche \textbf{cedere il
controllo della propria carriera}: se il team decidesse di non fare più nulla,
il manager ne pagherebbe le conseguenze. Il suo lavoro diventa scegliere le
persone giuste e farle lavorare bene. Lui lo esplicita sempre: sa di affidare
al team parte della propria carriera, e il team sa che lui combatterà per
loro, anche prendendosi le colpe.

\paragraph{L'inglese.} \textbf{Va assolutamente imparato}, anche lavorando in
Italia. Non saperlo preclude circa il 95\% del mercato (tutto quello fuori
dall'Italia) e soprattutto il \emph{networking} con persone che potrebbero dare
occasioni future: l'inglese è spesso l'unica lingua in comune con chi lavora
all'estero ed è oggi la lingua universale dei videogiochi.
\begin{itemize}
  \item \emph{Che livello serve?} Non importa B2 o C2: bisogna saper sostenere
        una conversazione, guardare almeno Netflix con sottotitoli in inglese,
        comunicare senza dodici pause e senza ``come si dice\dots''. Il
        certificato non verrà mai chiesto: verrà chiesto se si sa parlare
        inglese. Il relatore ha sentito persone con un C2 che non sapevano
        parlarlo.
  \item \emph{Come allenarlo senza occasioni di conversazione?} Netflix, Prime
        Video, Now TV, Udemy, Coursera, YouTube, Duolingo, siti di scambio
        linguistico con madrelingua, Discord: non sono mai esistiti così tanti
        mezzi gratuiti. Il relatore ha imparato l'inglese giocando (con una
        pronuncia sbagliata, perché leggeva senza sentire, su cui ha poi
        dovuto lavorare); guardare un film in inglese con i sottotitoli in
        italiano insegna anche la traduzione.
  \item Oltre all'ascolto passivo è fondamentale \textbf{parlare}: in una
        flotta di \emph{EVE Online} con 150 persone quasi tutte madrelingua
        inglese si è esposti a decine di accenti e modi di dire (mille modi di
        dire ``come stai'' oltre a \emph{how are you}); lavorando all'estero si
        impara a distinguere gli accenti (Manchester, Londra) e ad apprezzarli
        anche nei media.
\end{itemize}

\paragraph{Quanto deve sapere un junior?} Dipende da quanto è tecnico il
ruolo. Un game designer junior che non sa qualcosa può essere affiancato da un
senior; un programmatore che non sa programmare, in quel momento, non è utile
all'azienda.

\paragraph{Annunci ``junior con cinque anni di esperienza''.} ``Cercasi
ventenne con trent'anni di esperienza.'' Non è l'annuncio per un junior: è un
errore dell'azienda, che in realtà cerca un \emph{mid}. Riceverà profili
junior, si stupirà che non abbiano cinque anni di esperienza, rifarà l'annuncio
e avrà perso sei mesi. Il consiglio: tenere d'occhio quella pagina, perché
dopo qualche mese l'annuncio verrà corretto.

\paragraph{Una struttura per il post-mortem.} A chi chiede un formato di base,
il relatore anticipa quanto vedrà nella seconda giornata a proposito di
brainstorming: una lavagna \textbf{divisa in quattro aree}, su cui ognuno
attacca tutti i post-it che vuole:
\begin{center}
  \begin{tabular}{@{}p{0.44\textwidth}p{0.44\textwidth}@{}}
    \toprule
    \textbf{Smettere di fare} & \textbf{Iniziare a fare} \\
    ciò che non ha funzionato & ciò che non abbiamo fatto e andava fatto \\
    \midrule
    \textbf{Fare di meno} & \textbf{Fare di più} \\
    ciò che ha fatto perdere tempo rispetto all'impatto ottenuto &
    ciò che ha funzionato ma, fatto con troppa leggerezza, non ha dato i
    frutti sperati \\
    \bottomrule
  \end{tabular}
\end{center}
Con i post-it ognuno ha voce e tutto risulta catalogato; per il report basta
una foto della lavagna (o uno screenshot, se la lavagna è digitale, come Miro).

\paragraph{Come sbloccare la paralisi del feedback verso il management?} Uno
studente ha visto manager dire ``non abbiate paura, la mia porta è sempre
aperta'' senza che cambiasse nulla. C'è una dissonanza tra ciò che dicono e
ciò che fanno: è ``la mamma che dice vieni qua che non ti faccio niente''.
Qualcosa deve aver fatto pensare alle persone che non è sicuro. Prima va
sistemato quello; poi si può provare con i \textbf{feedback anonimi}.
L'obiettivo del relatore resta il feedback aperto e firmato, ma a Stormind ha
iniziato con feedback anonimi verso il management per verificare se
l'anonimato fosse un fattore. Non lo era (evidentemente non c'erano problemi a
parlare con il CEO), e la volta successiva intende abbandonarli. Ma è un buon
\emph{check}: se senza anonimato non arriva nulla e con l'anonimato arrivano
cento feedback, il problema è chiaro.

\paragraph{Quando è giusto rivolgersi a un superiore?} La strada giusta resta
parlare direttamente con l'interessato. Ma se si prova e non si trova una
soluzione, la cosa più ovvia è andare dal proprio manager insieme: ``ci
aiuti?''. Quando il relatore parlava di mediare i dialoghi, era proprio questo:
fare da moderatore, aiutare a sbloccare i punti fermi o, se non si riesce a
scegliere tra A e B, scegliere lui. L'escalation, quando necessaria, va fatta.

\paragraph{Il tempo libero veniva rispettato?} No, e se l'è dovuto imporre
lui. Il lavoro gli piaceva così tanto che continuare a casa non gli sembrava
lavorare. Gli si è ritorto contro con il lancio di \emph{Diablo~III}, fatto in
orario americano (le due o l'una di notte in Europa), con l'ufficio chiuso
perché a nessuno era stato chiesto di restare. Il lancio va malissimo, i
giocatori protestano, e lui, sveglio, risponde sui forum fino alle tre, quattro,
cinque del mattino: ``non posso risolvere il problema, ma resto qui ad
aggiornarvi con quello che riesco a sapere dal team di sviluppo''. Il giorno
dopo trova sulla scrivania un \textbf{richiamo formale dell'HR}: aver lavorato
fuori orario senza autorizzazione, condotta non professionale. Per lui è stato
un interruttore: ``ora me ne vado''. Due settimane dopo faceva colloqui con
Riot. Ancora oggi è contrario a quel richiamo, ma gli ha fatto capire come
ragiona un'azienda: ``hai un contratto 9--18, non mi interessa se l'hai fatto
con le migliori intenzioni''. Eppure per anni ha goduto di un'ottima
reputazione in quelle community come ``quello che è rimasto con noi fino alle
cinque del mattino'', e anche l'azienda ne ha beneficiato.

Ci sono poi aziende in cui, soprattutto per i developer, viene richiesto il
famigerato \textbf{crunch}: lavoro straordinario continuativo, senza pause.
Dipende molto dall'azienda; sta venendo (``grazie al cielo'') deprecato, perché
molte aziende capiscono che non è sano ed è visto male anche dal pubblico, ma
esiste ancora.

\paragraph{E se il tuo capo demotiva qualcuno del tuo team?} Qui ``diventa
personale'': il team del relatore si fida di lui perché sa che in questi casi
``diventa una bestia''. ``Te la puoi prendere con me quanto vuoi, ma il mio
team non lo tocchi.'' Lo gestisce direttamente con l'altro manager, per quanto
brutta possa diventare la conversazione. Lo sforzo più grande è \textbf{non
andarne poi a parlare con il team} (``ragazzi, ho fatto questa cosa''), perché
minerebbe la fiducia verso l'alto: non vuole un rapporto ottimo fino a lui e
pessimo oltre. L'obiettivo è che l'altro manager ``si faccia perdonare''
cambiando il proprio comportamento, e che il team percepisca da solo il
cambiamento.

\paragraph{Chiusura della giornata.} Il relatore chiede un feedback collettivo
in chat (``se domani ci dai \emph{more of the same} siamo contenti, oppure
ditemi cosa limare''), si scusa per i ripetuti crash della connessione e
ringrazia per la proattività nelle domande, che gli servono anche come
\emph{check}: se nessuno chiede nulla su un argomento, significa che può
andare avanti. Racconta di aver passato una settimana a decidere che cosa
mettere nelle slide e di averle poi preparate in poco tempo (rifacendone
alcune fino alle tre di notte). Ricorda infine di usare il test di
personalità per sé e per il proprio gruppo, facendolo fare a chi non l'ha
ancora fatto: tornerà molto utile per cooperare meglio.

% ---------------------------------------------------------------------
\section{Anticipazioni sulla seconda giornata}
\label{sec:z1-domani}

La seconda giornata del workshop sarà più pratica: ``oggi vi ho spiegato
perché, domani vi insegno come lo fanno queste aziende''. Gli argomenti
previsti:
\begin{description}
  \item[Time management] come organizzarsi una giornata professionale in
        azienda: come strutturare una \emph{task list}, cosa fare e cosa non
        fare, come e cosa prioritizzare.
  \item[Brainstorming] i gruppi dovranno fare moltissimi meeting (come facciamo
        il gioco, come lo chiamiamo, che scopo gli diamo). Se non si sanno fare
        bene, i brainstorming sono \textbf{dannosi ed eterni}: un brainstorming
        dovrebbe durare \textbf{al massimo due ore}. Verranno mostrati metodi
        che accorciano la durata, massimizzano l'effetto e, soprattutto,
        \textbf{fanno partecipare tutti}: non c'è niente di peggio di un
        brainstorming in cui parlano in due per il 90\% del tempo. Il classico
        ``ragazzi, non esistono idee sbagliate, sparate'' è il modo peggiore di
        farlo; non è colpa di nessuno, ma si può fare diversamente.
  \item[Guidare da dentro] come essere leader nel proprio team anche senza
        esserne il leader: sviluppare responsabilità personale e verso gli
        altri. Non essere il manager di qualcuno non significa non esserne
        responsabili: la persona accanto è responsabilità di tutti, perché si
        è un team e non un insieme di individui. ``Se qualcuno cade, tutti lo
        facciamo rialzare: non c'è l'addetto al rialzamento.'' È la versione
        \emph{advanced} della giornata: come comportarsi da manager senza
        averne il titolo.
\end{description}
Torneranno inoltre il \emph{service level agreement} e il tema della
condivisione dei successi e della protezione del team dai fallimenti.

% ---------------------------------------------------------------------
\section*{Punti chiave del capitolo}
\addcontentsline{toc}{section}{Punti chiave del capitolo}

\begin{itemize}
  \item ``La capacità di un gruppo di persone di fare cose straordinarie
        dipende dal modo in cui queste persone operano insieme.'' Il 90\% dei
        problemi dei team sono \textbf{problemi di comunicazione}.
  \item Le aziende si liberano sempre più presto dello \textbf{useless
        genius}: le competenze si trasmettono, il carattere no. Il
        \textbf{fit culturale} va valutato anche dal candidato.
  \item Non esiste un unico modo di dare feedback: dipende dal contesto e
        soprattutto da \textbf{chi lo riceve}. Il feedback è una somma di
        piccoli miglioramenti periodici; tenersi le cose nel cassetto non
        serve mai.
  \item I team attraversano \textbf{forming, storming, norming, performing e
        adjourning}. Lo storming non si salta; si passa la maggior parte del
        tempo in forming; la bravura sta nell'accelerare i passaggi e nel
        fare bene il \textbf{post-mortem}, anche dei successi.
  \item Gli obiettivi di un progetto vanno fissati \textbf{all'inizio}: solo
        così si può valutare se è andato bene e scegliere tra A e B.
  \item Il circolo virtuoso è \textbf{feedback $\rightarrow$ miglioramento
        $\rightarrow$ fiducia}. Accettare un feedback non significa fare come
        dice l'altro: decide chi ha la responsabilità, ma deve argomentare.
  \item ``Aiutami a poterti aiutare'': meglio dire in anticipo ``potrei non
        farcela'' che ``non ce l'ho fatta''. Chiedere aiuto non è una
        debolezza.
  \item Quando si dà feedback \textbf{non si porta la soluzione}: la si
        costruisce insieme a chi lo riceve.
  \item La cultura è forte quando si lavora \textbf{per gli altri}; il lavoro
        è valutato come opera d'insieme (anche con bonus di team).
  \item Come in \emph{Among Us}, i team sprecano energie a cercare il
        traditore invece di fare le task. \textbf{Default to trust}: chi ha
        sbagliato ha agito con le informazioni che aveva. Cercare
        \textbf{la causa, non il colpevole}: ``che cosa ti ha fatto pensare
        che fosse la scelta giusta?''.
  \item Carriere e performance non sono lineari: vanno valutate le
        \textbf{traiettorie}, non un punto dell'anno. Le valutazioni puntuali
        generano \emph{lying, hiding, faking}.
  \item Confermare sempre le istruzioni facendosele ripetere; ``non ho
        capito'' non è una colpa.
  \item Dare feedback: sull'\textbf{azione} e sui suoi \textbf{effetti}, specifico,
        presto, uno a uno, adattato al carattere, senza terze parti, senza
        ``non mi piace'', senza offendersi se non viene seguito; poi fare il
        debriefing.
  \item Ricevere feedback: \textbf{ascolta, chiedi dettagli, accetta, agisci}.
        Il feedback positivo serve a non peggiorare ciò che si fa bene; il
        feedback accidentale e il \emph{drive-by} non servono a nulla.
  \item Feedback \textbf{frequenti e recenti}: tanti piccoli \emph{sanity
        check} costano meno di una grande correzione di rotta.
  \item Il feedback è \textbf{un regalo} di chi ha a cuore il nostro
        interesse; fa paura il silenzio, non le critiche.
\end{itemize}

%    % ... altre lezioni zero ...
%    \part{Il corso di Game Design}
%    \renewcommand{\thechapter}{\arabic{chapter}}
%    \setcounter{chapter}{0}
%    % =====================================================================
%  Capitolo 1 — Introduzione al corso e al mestiere del game designer
%  Fonte: trascrizione della lezione GD1 (Game Design, primo anno)
%  Pacchetti richiesti nel preambolo del libro:
%    \usepackage[italian]{babel}, \usepackage[T1]{fontenc},
%    \usepackage[utf8]{inputenc} (solo pdfLaTeX), \usepackage{booktabs}
% =====================================================================

\chapter[Il mestiere del game designer]{Il mestiere del game designer: introduzione al corso}
\label{cap:introduzione-game-design}

Questo primo capitolo fa da soglia. Prima di affrontare metodi, framework e
documenti, è necessario capire \emph{in che modo} si impara il game design,
\emph{con chi} lo si fa, \emph{con quali strumenti} e, soprattutto,
\emph{che cosa fa concretamente} un game designer nel suo lavoro quotidiano.
Il capitolo segue l'impostazione della lezione inaugurale del corso: si parte
dal contesto formativo (il corso come simulazione di lavoro), si passa al
metodo di apprendimento (l'\emph{inquiry-based learning}), si elencano
letture e risorse di riferimento, e si arriva infine a una prima definizione
operativa del ruolo del designer, dei limiti entro cui lavora e delle sue
possibili specializzazioni.

Una premessa terminologica: il game design è una disciplina giovane e quasi
tutta la sua letteratura è scritta in inglese. Per questo motivo molti termini
tecnici verranno lasciati in lingua originale (\emph{pipeline},
\emph{pitching}, \emph{tuning}, \emph{feature team}, \dots): sono i termini che
si useranno nel lavoro, ed è importante imparare a riconoscerli e a usarli
correttamente fin da subito.

% ---------------------------------------------------------------------
\section{Il contesto: un corso come simulazione di lavoro}
\label{sec:simulazione-lavoro}

\subsection{Il principio fondamentale}

Il primo concetto da fissare è che il corso \textbf{si presenta come una
simulazione del lavoro}. Dal momento in cui si entra in aula si viene trattati
come \emph{potenziali professionisti dell'industria}: non soltanto studenti,
ma lavoratori effettivi di uno studio (l'accademia), con una gerarchia, dei
referenti, dei compiti e delle scadenze.

Ne derivano alcune conseguenze pratiche:
\begin{itemize}
  \item il docente non assegna un ``compitino'' da risolvere per accumulare
        punti o presenze: ogni \emph{task} è pensato per far dedicare tempo
        all'apprendimento di qualcosa di concreto;
  \item non consegnare un task non danneggia il docente (che ``vive ancora
        domani e viene pagato ugualmente''), ma lo studente, che perde
        l'occasione di mettere le mani su un argomento e diventare più abile;
  \item docenti e studenti devono trattarsi \textbf{da colleghi}, perché è
        molto probabile che un domani lo diventino davvero.
\end{itemize}

\subsection{Perché l'impegno conta più che mai}

Il lavoro del game designer è descritto come \emph{molto romantico, ma anche
molto duro}, e sempre più duro col passare degli anni. Durante il periodo del
Covid l'industria ha visto grandi investimenti; quei ``rubinetti'' sono stati
chiusi da tempo, il rapporto con i publisher è diventato più complesso e le
opportunità lavorative non sono più quelle di tre o quattro anni fa. Questo
significa:
\begin{itemize}
  \item più competizione per ogni posizione;
  \item aspettative più alte: ci si aspetta che i candidati siano ancora più
        preparati e sicuri delle proprie competenze;
  \item un unico modo per diventarlo: \textbf{lavorare}. Più tempo si dedica
        oggi ad apprendere il mestiere, più facilmente si sarà assumibili
        domani e più agevole sarà il lavoro di tutti i giorni.
\end{itemize}

\subsection{Sbagliare qui per non sbagliare fuori}

Il mantra del corso è: \textbf{quanto più provate qui, anche in modo goffo e
sbagliando, tanto meno sbaglierete là fuori}. In accademia, per quanto la
simulazione sia realistica, nessuno toglie lo stipendio e nessuno licenzia da
un momento all'altro. Nell'industria sì, e un errore può costare molto denaro
all'azienda, ancora di più se l'azienda è la propria.

\subsection{Il valore del fallimento}

Il fallimento, se affrontato con la mentalità corretta, è probabilmente
\emph{la cosa più vantaggiosa} che si possa fare durante la formazione:
\begin{itemize}
  \item un successo (``avete superato questa deadline'') conferma le proprie
        abilità, ma insegna relativamente poco;
  \item dopo $n$ fallimenti si impara moltissimo, e i successi arrivano;
  \item è fondamentale non demoralizzarsi: il feedback negativo non è mai un
        giudizio sulla persona, ma su ciò che è stato consegnato.
\end{itemize}

Il docente riconosce apertamente che dedicare tempo ed energia a un progetto e
vederlo rifiutato, o vederlo non uscire mai, è demotivante: è successo a lui e
lo ha visto accadere anche su progetti molto grandi, con team numerosi. Il
punto però non è ciò che si è ottenuto materialmente, ma \textbf{quanto si è
cresciuti}. Da ogni esperienza ci si portano a casa due cose:
\begin{enumerate}
  \item \textbf{l'esperienza}: quanto si è diventati bravi, veloci,
        praticamente abili nell'applicare un framework, quanto si è cresciuti
        come designer;
  \item \textbf{il rapporto umano}: nel lavoro si ``suda sangue'' fianco a
        fianco con altre persone, e questo crea legami molto forti che
        portano avanti l'industria nel senso più positivo del termine: nascono
        progetti privati, ci si porta da un'azienda all'altra, si formano team
        con basi solide.
\end{enumerate}

Va evitato anche l'eccesso opposto: la cultura del \emph{grinding} si è
appropriata di questi discorsi in modo distorto. Non bisogna diventare robot.

% ---------------------------------------------------------------------
\section{Organizzazione del corso e gerarchia}
\label{sec:gerarchia-corso}

\subsection{La struttura gerarchica}

La struttura dei referenti del corso ricalca volutamente una
\textbf{gerarchia tipica dell'industria}:
\begin{itemize}
  \item al vertice c'è il \textbf{direttore del corso} (Diego Ricciuti), che
        corrisponde al \emph{creative director}: supervisiona la visione del
        progetto (il concetto di ``visione'' verrà approfondito più avanti);
  \item sotto di lui ci sono i \textbf{referenti di dipartimento}, uno per
        disciplina, con un corrispettivo per ciascuna sede (Milano e Roma);
  \item in fondo ci sono gli studenti, che occupano il primo gradino di
        un'azienda: l'\emph{intern} (stagista) o il \emph{junior}, il passo
        che precede il lavoratore ``regolare''.
\end{itemize}

\begin{table}[htbp]
  \centering
  \begin{tabular}{@{}p{0.34\textwidth}p{0.29\textwidth}p{0.29\textwidth}@{}}
    \toprule
    \textbf{Disciplina} & \textbf{Milano} & \textbf{Roma} \\
    \midrule
    Direzione (\emph{creative director}) & \multicolumn{2}{p{0.58\textwidth}}{Diego Ricciuti} \\
    Game design      & Mattia Podini        & Giacomo ``Jack'' Danna \\
    Level design     & Stefano Zocchi       & Andrea Dispoli \\
    Narrative design & \multicolumn{2}{p{0.58\textwidth}}{Daniele Fusetto (lezioni condivise), Giacomo Masi} \\
    Programmazione (Unreal) & Alessandro Pregnolato & Matteo Velardi \\
    \bottomrule
  \end{tabular}
  \caption{Referenti del corso per disciplina e sede (nomi come riportati in lezione).}
  \label{tab:referenti}
\end{table}

Per il primo anno, il referente diretto degli studenti di game design è il
docente di game design: è a lui che ci si rivolge per dubbi e domande. Altre
figure utili sono il direttore didattico (severo, ``giustamente severo'', da
contattare solo per questioni particolarmente importanti), l'amministrazione
(contratti e documentazione) e la segreteria.

\subsection{I canali di comunicazione}

\begin{description}
  \item[Discord] è il canale per le comunicazioni ufficiali del corso e per
        le domande che non richiedono urgenza immediata; può avere qualche
        ritardo.
  \item[E-mail] è il canale per tutto ciò che è formale o ``ancora più
        ufficiale'' e per le comunicazioni rapide e immediate.
  \item[Telegram] esistono residui di vecchi canali precedenti a Discord, da
        evitare.
\end{description}

Le domande che valgono per tutti vanno fatte in classe (anche da un portavoce
a nome del gruppo); quelle singole e personali via Discord o, se formali, via
e-mail. Se il docente non risponde a una domanda c'è un motivo, ad esempio
ragioni legali, e lo dirà esplicitamente.

\subsection{I tre moduli e la settimana tipo}

Le tre discipline principali del primo anno (e non solo) sono:
\begin{itemize}
  \item \textbf{Game Design} (tendenzialmente il lunedì);
  \item \textbf{Level Design} (tendenzialmente il mercoledì);
  \item \textbf{Programmazione} (tendenzialmente il venerdì).
\end{itemize}
Nel primo mese le lezioni sono tutte di game design (lunedì, mercoledì e
venerdì). Sono previste anche lezioni di narrativa, ma il piano formativo si
aggiorna ogni anno: le ore di narrativa sono state ridotte a favore di quelle
di programmazione. Le poche lezioni di narrativa daranno quindi dei
\emph{cenni}, non la preparazione di un narrative designer. Ciò che serve
davvero è imparare \textbf{come comunicano gli altri dipartimenti} e come
lavorare con persone che fanno un lavoro diverso ma complementare al proprio.

\subsection{Le due anime delle lezioni di game design}

\begin{description}
  \item[Design methods (metodologie)] lezioni prevalentemente teoriche, ma con
        una buona componente pratica, in cui si studiano metodologie di design
        e i cosiddetti \emph{framework} dell'industria.
  \item[Documentazione e pitching] il cuore del lavoro quotidiano del
        designer.
\end{description}

\paragraph{La documentazione.}
Il docente si definisce ``il più grande nerd della documentazione'' che si
possa incontrare. È una parte dura, perché a pochi piace documentare, ma il
lavoro del designer \emph{è principalmente questo}: scrivere documenti e, da
quei documenti, passare alla pratica e all'implementazione. Chi non sa
documentare bene parte male. Quando non si è presenti, il documento è il
proprio \textbf{biglietto da visita}: qualcun altro lo leggerà, quindi va
scritto nel modo più chiaro e comprensibile possibile.

Un esempio concreto del perché la documentazione conta: CD Projekt Red ha
sviluppato un proprio engine proprietario, ma buona parte della sua
documentazione è stata ``buttata all'acqua''. Chi deve usarlo o lo sa già
usare, o ha competenze di un altro engine da trasferire, oppure è in
difficoltà. Si documenta \textbf{perché altri possano leggere}.

\paragraph{Il pitching.}
Il \emph{pitching} è l'esposizione, la presentazione orale. La parte creativa
del team presenta spesso:
\begin{itemize}
  \item le proprie idee ai colleghi;
  \item il gioco a convention ed eventi;
  \item il pitch del gioco agli investitori, per ottenere i fondi necessari a
        realizzarlo.
\end{itemize}
Saper parlare e comunicare chiaramente le idee è fondamentale. In sintesi:
\emph{il game designer parla}, e bisogna farsene una ragione.

\subsection{Altre componenti del corso}

\begin{description}
  \item[Gioco critico e playtest] la classe viene divisa in due gruppi: uno
        gioca, l'altro osserva. Chi osserva non applaude ``come le
        scimmiette'': annota, segna e analizza il \emph{playtest}, esattamente
        come si fa nell'industria. Il playtest fa parte delle mansioni del
        designer, ma esiste anche come specializzazione autonoma.
  \item[Pipeline] termine che indica tutto ciò che riguarda i processi
        formalizzati dell'industria; di solito si intende la \emph{pipeline di
        produzione}: come si divide il team, qual è la sua gerarchia, come il
        \emph{producer} gestisce le deadline. In breve: \emph{come} e
        \emph{quando} si fanno le cose. Al primo anno è meno rilevante (si
        dovrà ``nuotare'' un po' da soli); dal secondo anno si vedranno le
        metodologie effettivamente applicate ai progetti digitali. Si
        scoprirà quanto sia difficile lavorare a un progetto digitale, e
        ancor più fisico, senza un piano di sviluppo comprovato, come quelli
        usati per qualsiasi software in qualsiasi parte del mondo.
  \item[Engine] il motore di riferimento è \textbf{Unreal Engine}
        (sezione~\ref{sec:engine}).
  \item[Programmazione] principalmente tramite \emph{Blueprint} in Unreal,
        più eventuali sistemi di versionamento; il corso parte da zero.
  \item[Progetti pratici] probabilmente la cosa più importante. Al primo
        anno i progetti di gruppo restano circoscritti al corso di design,
        perché mancano ancora alcune nozioni (come comunicare, affrontare
        problemi, dare e ricevere feedback). In seguito si allargheranno a
        progetti digitali, molto più difficili di quelli fisici, con membri
        di altri team: un rapporto non unilaterale, ma ``covalente''.
\end{description}

% ---------------------------------------------------------------------
\section{Il metodo: l'inquiry-based learning}
\label{sec:inquiry}

\subsection{Che cos'è}

Il corso è costruito secondo l'\textbf{inquiry-based learning}
(apprendimento basato sull'indagine), un approccio pedagogico nato intorno
agli anni Sessanta, che si richiama a pensatori come Dewey e Piaget. In
generale, le discipline da cui il design attinge sono quasi sempre la
sociologia, la psicologia e, talvolta, l'economia.

Il funzionamento è un ciclo:
\begin{enumerate}
  \item il docente fornisce metodologie e framework dell'industria durante
        la lezione;
  \item lo studente, fuori dalle ore di lezione, lavora sui \emph{task}:
        mette a terra le nozioni, le maneggia, le assorbe meglio;
  \item da questo lavoro nascono nuovi dubbi;
  \item i dubbi tornano in aula: lo studente chiede, il docente risponde
        (ammesso che dare la risposta sia utile, e ammesso che una risposta
        esista).
\end{enumerate}

\subsection{Non sempre esiste una risposta giusta}

Non tutto ciò che si studia ha ``la risposta migliore possibile in assoluto
nel vuoto''. Esistono metodologie, limiti e framework che aiutano il lavoro,
ma non sono fissi: cambiano da azienda ad azienda, da lavoro a lavoro, da
progetto a progetto, e a volte vengono stravolti. Sta alla
\textbf{sensibilità del designer} capire quale strumento o framework usare in
un dato momento, e questa sensibilità si impara solo con l'esperienza. Nel
lavoro nessuno dirà ``usa questa cosa'' o ``è meglio quest'altra'': sarà il
designer a scegliere.

Anche l'esperienza del docente ha confini: sei-sette anni di lavoro, un modo
personale di approcciare il design e un certo insieme di progetti (ad esempio
nessun MMO). Oltre quei confini si cercano risposte nella ricerca, studiando
come aziende con esperienza specifica (ad esempio Riot) hanno affrontato
problemi simili e hanno condiviso le loro soluzioni, sempre valutando quali
metodi siano adatti alla dimensione del team, al tipo di progetto e al tempo
disponibile.

\subsection{Sensibilità, non memoria}

Non interessa imparare a memoria un framework o la teoria di un autore. Certe
terminologie vanno conosciute, ma per ragioni di comunicazione e chiarezza.
Ciò che conta è sviluppare la \textbf{capacità di farsi domande}, di farle
agli altri, di fare ricerca e di trovare una propria risposta. È esattamente
ciò che si fa nel lavoro: davanti a qualcosa che non si sa fare, l'abilità
del designer sta nel cercare come qualcun altro ha risolto un problema simile.

\subsection{Le lezioni come corsi di aggiornamento}

Le lezioni frontali simulano i \emph{corsi di aggiornamento} aziendali: le
sessioni che un lead di design tiene al suo team di design (o il lead di
level design al suo team, e così via). Alla fine dell'aggiornamento c'è quasi
sempre un lavoro da fare, un obiettivo da portare a terra:
\textbf{l'obiettivo della settimana}. Si lavora di settimana in settimana, e
il perché verrà chiarito in seguito.

Struttura tipica di una lezione:
\begin{enumerate}
  \item breve sessione di domande e risposte sulla lezione precedente;
  \item analisi di eventuali esercizi particolarmente critici;
  \item lezione vera e propria;
  \item discussione di gruppo, se c'è tempo;
  \item assegnazione del task della settimana.
\end{enumerate}
Il docente tende a ``parlare e straparlare'', ad aprire parentesi e a saltare
da un argomento all'altro: l'unico modo per fermarlo è che siano gli studenti
a interromperlo con domande.

\subsection{La gerarchia della competenza}

L'inquiry-based learning può essere rappresentato come una scala di
competenza:
\begin{enumerate}
  \item la sola lezione frontale porta a \textbf{conoscenze incomplete},
        perché manca la parte pratica e di ricerca;
  \item il lavoro a casa (task, ricerca) porta a una competenza \textbf{più
        completa}, non sbagliata in sé, ma ancora parziale e complementare ad
        altre;
  \item la \textbf{condivisione} dei task e dei feedback con gli altri porta
        idealmente alla competenza più completa possibile;
  \item il picco della competenza è la sua \textbf{applicazione nei
        progetti} (fisici o digitali, di gruppo): ``con le cose che avete
        imparato e ricercato, riuscite a fare un gioco?''. Quando la risposta
        diventa positiva, o comunque si inizia a navigare nei processi di
        design per realizzarne uno, quella è la vittoria.
\end{enumerate}
Non si parte subito creando giochi: si parte dall'analisi, dal design e dalle
regole generali di design, prima ``nel vuoto'', poi contestualizzando i
framework, infine unendo i puntini nei progetti.

Rimarranno comunque delle lacune, per tutti, docente compreso, perché il game
design \textbf{non è una scienza esatta}, anzi, a rigore non è affatto una
scienza. Prende in prestito metodologie dalle discipline scientifiche, ma
contiene una componente imprevedibile e incalcolabile: \textbf{la
produzione}. Si possono applicare perfettamente le regole di un framework, e
quel framework può comunque non essere adatto al gioco, al proprio modo di
sviluppare o al tempo di produzione disponibile (alcuni framework richiedono
una ricerca e un'attenzione che il tempo non concede).

C'è poi la componente creativa, la parte ``romantica'' del lavoro: esiste, ma
è \textbf{metà del lavoro}. L'altra metà è ricerca, metodo scientifico al
100\%. Dagli stessi dati $n$ designer diversi, con sensibilità diverse,
possono ottenere $n$ risposte diverse. Non tutte saranno valide: vanno
applicate, messe a terra, e questa scrematura fa parte del lavoro.

\subsection{I livelli dell'apprendimento}

Un secondo modello, a gradini, descrive i livelli di apprendimento (è la
cosiddetta tassonomia SOLO):
\begin{description}
  \item[Pre-strutturale] il docente dà il compito e lo studente lo risolve,
        possibilmente il prima possibile. È esattamente ciò che \emph{non} si
        vuole.
  \item[Uni-strutturale] si affronta una sola disciplina: è il primo passo
        necessario, ed è il motivo per cui nel primo mese ci sono solo lezioni
        di game design.
  \item[Multi-strutturale] entrano in gioco altri dipartimenti: level design,
        game design, programmazione e così via.
  \item[Relazionale] si capisce che alcune cose valgono per più
        dipartimenti: alcuni framework di narrativa derivano dal design o
        possono essere usati per fare design; le regole di documentazione si
        applicano anche ad altri dipartimenti; le regole di level design si
        mescolano con quelle di game design. Dopotutto, per fare un gioco
        servono molti dipartimenti.
  \item[Astratto esteso] il punto finale: muoversi autonomamente e capire
        cosa usare e quando, sulla base della propria esperienza. L'unico modo
        per arrivarci è fare pratica e fare progetti.
\end{description}

% ---------------------------------------------------------------------
\section{Task, feedback e regole del corso}
\label{sec:task-feedback}

\subsection{I task come opportunità}

I task non sono strettamente obbligatori (esiste un numero minimo di
consegne), ma sono \textbf{opportunità}. Quanto e come ci si lavora dipende
dallo studente. Un task fatto in cinque minuti, o in cui si è stati ``furbi'',
non verrà letto né riceverà feedback.

Domande da porsi davanti a ogni task:
\begin{enumerate}
  \item \textbf{Perché} mi è stato assegnato? Qual è l'obiettivo?
  \item Come ci lavoro per ottenere un risultato \textbf{che abbia senso}? Se
        la ricerca e i dati indicano una direzione, è inutile andare in
        un'altra.
\end{enumerate}

\subsection{Quanto lavorare: un calcolo indicativo}

Se il corso è una simulazione di lavoro, vale il carico di un lavoro normale:
\begin{itemize}
  \item un lavoro a tempo pieno comprende da 30 a 45 ore settimanali; si
        considera una media di \textbf{40 ore};
  \item si sottraggono circa 9 ore di lezione: restano $\approx 31$, diciamo
        \textbf{30 ore};
  \item con circa 3 task a settimana: $30 / 3 = $ \textbf{circa 10 ore per
        task}.
\end{itemize}
Il numero può variare, ma rende l'idea: dedicare un'ora a un task e poi fare
altro non fa guadagnare nulla a nessuno. Lo studente ``butta soldi'', il
docente risparmia tempo leggendo il task in un minuto (o non leggendolo).
Dedicando invece attenzione, ricerca, cura nella formattazione e
sperimentazione, e facendolo con cognizione di causa, si impara.

\subsection{Niente voti, solo feedback}

Ad eccezione del corso di programmazione, non ci sono voti (né ``ottimo'', né
10, né 100). Ci sono \textbf{feedback}:
\begin{itemize}
  \item in game design, generalmente sulla documentazione;
  \item in level design, sui progetti pratici.
\end{itemize}
Il docente indica cosa va bene e cosa va modificato, concentrandosi in
particolare su ciò che va cambiato. Il feedback \textbf{non è mai un giudizio
sulla persona}, ma su ciò che è stato consegnato nel concreto.

\subsection{La condivisione dei task}

Tutti i task consegnati vengono \textbf{condivisi} con la classe, e possono
essere letti solo dopo aver ricevuto i feedback del docente. Si cresce così
non solo con i propri feedback, ma anche con quelli dati agli altri. Regola
importante: ciò che il docente dice a uno studente \textbf{vale per tutti}. Se
a una persona viene detto ``questa cosa non va mai più fatta'', da quel
momento vale per l'intera classe.

\subsection{Non copiare}

Copiare i task degli altri è una delle cose peggiori che si possano fare, non
per una questione morale, ma perché annulla il senso del percorso. Si può
dedicare a ogni task il tempo e l'attenzione che si ritengono opportuni, ed è
normale che col tempo ci si specializzi (preferendo, ad esempio, il level
design o la programmazione tecnica) e si investa meno dove si sente meno
necessità. In questi casi l'alternativa corretta è:
\begin{itemize}
  \item non consegnare;
  \item studiare ciò che hanno fatto gli altri;
  \item eventualmente riconsegnare in seguito.
\end{itemize}
Ci saranno momenti in cui si verrà sommersi dalle cose da fare: bisogna
imparare a scegliere.

\subsection{Il ``non leggo''}

Il docente ha piena facoltà di non dare alcun feedback a un task (il
cosiddetto \textbf{``non leggo''}) se:
\begin{itemize}
  \item non si è lavorato con attenzione;
  \item non si rispettano le indicazioni date;
  \item è stata violata una regola estremamente importante.
\end{itemize}
Un task ``non letto'' equivale a un task non consegnato e non entra nel
conteggio. Attenzione: un task \emph{perfetto} e un task che \emph{fallisce}
nel comprendere o applicare un framework sono \textbf{ugualmente validi},
purché ci si sia impegnati e si sia dedicato tempo, pazienza e ``sudore''.

\paragraph{Riconsegna.} Un task non letto, in ritardo o che si vuole
migliorare può essere riconsegnato e verrà conteggiato, a patto che:
\begin{itemize}
  \item i problemi segnalati siano stati risolti;
  \item sia stato fatto ulteriore lavoro, integrando anche tutto ciò che nel
        frattempo è stato detto agli altri colleghi nei loro feedback.
\end{itemize}
Il task va sviluppato con le competenze che si hanno nel momento della
riconsegna, non con quelle del momento della prima consegna.

\subsection{Requisiti per superare l'anno}

\begin{itemize}
  \item consegnare almeno l'\textbf{80\% dei task} assegnati (soglia
        malleabile);
  \item partecipare almeno all'\textbf{80\% delle lezioni}, in presenza o
        da remoto, preferibilmente in presenza; i ritardi vengono valutati
        caso per caso;
  \item comportarsi da persone corrette e rispettose (è stata firmata una
        carta dello studente): non urlare, non alzare le mani, non
        vandalizzare, non compiere atti illegali.
\end{itemize}
Impegnarsi di più porta inoltre a essere selezionati per varie opportunità.

\subsection{Scadenze}

\begin{itemize}
  \item la scadenza di un task è \textbf{sempre prima dell'inizio della
        lezione} successiva indicata;
  \item se la classe, come team, mantiene comportamenti scorretti, il docente
        può \textbf{anticipare le deadline} per tutti.
\end{itemize}

\subsection{Accessibilità e DSA}

Gli studenti con disturbi specifici dell'apprendimento o neurodivergenze
(dislessia, ADHD, autismo, ecc.) sono invitati a segnalarlo. Quando si
affronteranno le regole di formattazione e documentazione (errori
grammaticali, mancanze, incoerenze) se ne terrà conto e si potranno trovare
dei compromessi, pur ricordando che nel mondo del lavoro nessuno farà
``sconti''. La documentazione può essere inviata all'amministrazione o alla
segreteria, che la comunicherà a tutti i docenti, ma il consiglio è parlarne
anche direttamente con ogni singolo docente.

% ---------------------------------------------------------------------
\section{Le competenze trasversali}
\label{sec:soft-skill}

\subsection{I giochi si fanno in gruppo}

Si possono fare giochi da soli, ed è un'impresa ammirevole; ma lavorando da
soli \textbf{tutto ciò che si guadagna in libertà si perde in produttività}.

\paragraph{Il caso Stardew Valley.} \emph{Stardew Valley} è stato realizzato
da una sola persona, Eric Barone, in anni di sviluppo e di sacrifici: faceva
un altro lavoro per mantenersi (in un cinema) e di notte programmava, faceva
l'arte e componeva la musica. Il risultato è stato straordinario e lo ha reso
ricco, ma si tratta di un caso estremamente raro. Quanti giochi simili, nati
nelle stesse condizioni, sono falliti? Molti di più.

\paragraph{Il rischio d'impresa.} Lavorare da soli non è sbagliato di per sé
e ha aspetti molto positivi: si ha il privilegio di creare esattamente ciò che
si vuole. Ma ci si fa carico anche di tutto il \textbf{rischio d'impresa}: se
si lavora a un gioco per cinque anni e poi il gioco va male, resta il
vantaggio di averlo pubblicato, ma il dopo è sperare di vivere di quella
release o di finanziare con essa il gioco successivo. Le grandi aziende hanno
gli stessi problemi, ma li distribuiscono su più persone e più responsabilità.

L'obiettivo di un professionista è \textbf{saper lavorare in diverse
condizioni}, e la più comune è quella del dipendente (o del freelance /
libero professionista con partita IVA): si viene assunti per le proprie
capacità da un'azienda, piccola o grande che sia, e si lavora con altre
persone.

\subsection{Essere persone con cui si vuole lavorare}

Se le persone non vogliono lavorare con voi, quello è \textbf{il primo motivo
di licenziamento}. E se non si viene licenziati, il clima nell'azienda
diventa terribile: la produttività crolla, la motivazione crolla. Le
\emph{soft skill}, cioè l'insieme di abilità che rendono una persona
piacevole con cui lavorare, vanno coltivate e diventano sempre più importanti
col passare degli anni (relativamente meno ora, molto di più quando i team
saranno più grandi). Anche chi è un ``lupo solitario'' deve imparare a:
\begin{itemize}
  \item lavorare con altre persone;
  \item comunicare con altre persone;
  \item trovare un punto in comune da cui partire per fare le cose.
\end{itemize}

\subsection{Indipendenza e proattività}

Sono le altre due parole chiave del corso:
\begin{description}
  \item[Indipendenza] spesso si verrà lasciati da soli, non per cattiveria ma
        per sfida: quando si cerca una risposta da soli, ciò che conta non è
        tanto la risposta quanto \textbf{il processo}. Una risposta data dal
        docente non vale niente.
  \item[Proattività] chiedere, domandare, agire da soli senza aspettare
        direttive, perché spesso quelle direttive non ci sono, a maggior
        ragione nei team piccoli.
\end{description}
Tutto questo si concretizza nel parlare, nel presentare e nel fare
\emph{post-mortem} (l'analisi a posteriori di un progetto concluso, che verrà
approfondita più avanti), e in generale nell'attenzione alla crescita
personale.

\subsection{Fare domande}

Il corso insiste sulle domande: alzare la mano, intervenire anche ``a
bomba'', fermare il docente ogni volta che qualcosa non è chiaro. Se una
persona non conosce qualcosa, molto probabilmente c'è almeno un'altra persona
che non la conosce. Il docente può dare per scontata la conoscenza di termini
che non lo sono, ed è del tutto normale, soprattutto all'inizio.

% ---------------------------------------------------------------------
\section{Strumenti: engine, Unreal e materiali}
\label{sec:engine}

\subsection{Unreal Engine e gli altri motori}

Il motore di riferimento del corso, usato in level design, è \textbf{Unreal
Engine}: ``un gran mattone'', ma che dà a un designer molte possibilità. Il
vantaggio principale è che permette a un designer non pratico di
programmazione di \textbf{prototipare tramite i Blueprint} (il sistema di
scripting visuale di Unreal). Questo non impedisce di passare a Unity o ad
altri engine, e viceversa: il docente stesso ha studiato e lavorato a lungo
con Unreal, ma la sua ultima esperienza professionale (come level designer) è
stata su Unity, un engine che non conosceva, a dimostrazione che il salto da
un engine all'altro è possibile.

Panorama dei motori citati:
\begin{itemize}
  \item \textbf{Unreal Engine} e \textbf{Unity}: i due più diffusi, con
        modelli di sorgente ``più o meno'' aperto;
  \item \textbf{Godot}: engine open source;
  \item \textbf{CryEngine}: proprietario di Crytek;
  \item \textbf{Snowdrop}: engine di Ubisoft;
  \item Rockstar usa un proprio engine proprietario;
  \item alcune aziende usano versioni \textbf{personalizzate} di engine
        esistenti: ad esempio Milestone, nei giochi di corse, usa un Unreal
        customizzato;
  \item CD Projekt Red ha un engine proprietario (con i problemi di
        documentazione già citati).
\end{itemize}
In sintesi: ogni azienda può avere un engine proprietario oppure una versione
modificata di un engine commerciale, ma si parte sempre da una base.

\subsection{Programmazione: cosa ci si aspetta da un designer}

Non è richiesto alcun bagaglio pregresso di programmazione: il corso parte da
zero. Gli studenti sono stati selezionati (il corso è a numero chiuso) per
competenze analitiche e critiche sopra la media. Anche nel mercato del lavoro,
per una posizione di game designer generalmente \textbf{non ci si aspetta
competenza di scripting} (C++, Java, C\#, ecc.), a meno di voler fare il
\emph{technical designer}, e anche lì varia da azienda ad azienda.

\subsection{Versione e preparazione}

Indicazioni pratiche emerse in aula:
\begin{itemize}
  \item scaricare Unreal Engine; la versione consigliata per il corso è
        stata indicata nella \textbf{5.6} (con conferma successiva). La
        stabilità dipende da cosa si deve fare, da quali tool e plugin si
        usano: i problemi nascono spesso dai plugin. Nel passaggio da Unreal 4
        a Unreal 5, ad esempio, il sistema \emph{Landscape} era inizialmente
        inutilizzabile, mentre in Unreal 4 funzionava molto bene;
  \item è consigliato installare tramite \textbf{Epic Games Launcher}, più
        comodo: permette anche di riscattare giochi gratuiti periodici e di
        accedere ai pacchetti di \textbf{Fab} (il marketplace di asset
        acquisito da Epic), che ogni mese offre alcuni asset gratuiti utili
        per esercitarsi nel level design (ad esempio per fare \emph{speed
        level design});
  \item molti pacchetti di UI o di sistemi di interazione sono utili per
        prototipare al volo, ma spesso inutilmente complicati;
  \item un plugin di \textbf{materiali di prototipazione a griglia} (adattivi,
        con metriche, allineabili alla griglia, con materiali già pronti e
        colori modificabili, dal costo di circa 5 euro) è indicato come
        ``salvavita'' per il level designer; se necessario, i pacchetti
        indispensabili vengono forniti;
  \item portare \textbf{sempre il portatile} alle lezioni di level design
        (da novembre, a partire dall'interfaccia dell'engine); alle lezioni
        di game design basta che lo porti almeno metà della classe, per i
        lavori a coppie. Serve comunque qualcosa per prendere appunti
        (``tavolette d'argilla e cannucce, se volete fare i babilonesi'',
        oppure carta e penna).
\end{itemize}
Chi ha un portatile poco potente dovrebbe comunque provare a installare
l'engine e aprire un template per vedere come gira: all'inizio non si faranno
cose pesanti. Nelle prime lezioni di level design è importante mettere subito
le mani sull'interfaccia: ascoltare dove si trova uno strumento (ad esempio lo
\emph{snapping}) senza provarlo porta a dimenticarlo.

% ---------------------------------------------------------------------
\section{Letture e risorse di riferimento}
\label{sec:risorse}

\subsection{I libri}

\begin{description}
  \item[\emph{The Art of Game Design: A Book of Lenses} (Jesse Schell)] il
        testo \textbf{davvero importante}. Alcune lezioni finali dell'anno si
        baseranno sulle sue terminologie. È utile perché apre la mente sulla
        possibilità di applicare metodologie diverse a giochi diversi, ma
        soprattutto per la \textbf{terminologia} comune. Esiste anche in
        audiolibro (circa 25 ore).
  \item[\emph{Challenges for Game Designers} (Brenda Romero e Ian Schreiber)]
        libro \emph{entry level} con moltissimi esercizi, molto caro al
        direttore del corso.
  \item[\emph{21st Century Game Design} (Chris Bateman)] utile come
        retrospettiva: è un libro visionario che ha previsto metodi e
        soprattutto \emph{archetipizzazioni} dei giocatori tuttora valide nel
        design e nel marketing. Servirà per una delle lezioni più importanti
        dell'anno. È utile leggerlo per vedere quanto aveva previsto, ma è
        utile ``fino a un certo punto''.
  \item[\emph{An Architectural Approach to Level Design} (Christopher
        Totten)] non strettamente necessario, ma di lettura semplice. Ha due
        anime: parte dall'architettura (babilonese, greca, ecc.) e la trasla
        su esempi di level design (come \emph{Doom} e \emph{Team Fortress}).
        Un tempo ne esisteva un riassunto su Gamasutra (oggi
        \emph{Game Developer}).
  \item[\emph{A Theory of Fun for Game Design} (Raph Koster)] utile fino a un
        certo punto: contiene diverse idee superate, che lo stesso Koster ha
        affrontato in vari talk reperibili nel GDC Vault. Verrà ripreso nelle
        prime lezioni, poi messo da parte.
\end{description}

\subsection{Canali e conferenze}

\begin{description}
  \item[GDC (Game Developers Conference)] il suo canale raccoglie talk e
        panel estremamente importanti: è ``uno strumento misterioso che ci
        servirà più tardi''. \textbf{Attenzione}: per molti talk mancano ancora
        le competenze per interpretarli, e alcuni potrebbero fuorviare o
        mettere in testa idee strane. La GDC si concentra spesso su
        produttività e mobile.
  \item[Il \emph{Dragon Speech} di Chris Crawford] il talk del fondatore della
        GDC: appartiene a un passato lontano, ma è molto motivante. Da vedere
        subito.
  \item[Game Maker's Toolkit] canale particolarmente utile per nozioni
        pratiche, con video brevi e alcune interviste.
  \item[Nordic Game] altra conferenza importante, con talk e panel di esperti
        del settore.
\end{description}

\subsection{Associazioni e premi: come leggerli}

Le associazioni di settore (ad esempio IIDEA in Italia, o la IGDA) non sono
il male, anzi: sono un canale importante per affacciarsi all'estero o per
accedere a fondi. Va però ricordato che \textbf{tutelano i propri associati}:
non esiste solo il loro palinsesto, e bisogna valutare quanto spazio occupano
nel panorama della visibilità.

Lo stesso vale per i \textbf{premi} (Game Awards, ``Best Mobile Game'',
``Games for Change'', ecc.): al designer \emph{non interessano}. Che un gioco
abbia vinto ``più di 90 premi'' non dice nulla di utile. Servono a dare
visibilità e lustro ai giochi ed è giusto che esistano, perché ampliano il
panorama e il numero di giocatori. Ma ciò che interessa al designer sono i
\textbf{dati}:
\begin{itemize}
  \item come sono state implementate certe cose;
  \item come è stato fatto il marketing e dove il gioco è stato presentato;
  \item non quante wishlist ha fatto, ma \emph{come} le ha fatte;
  \item in che periodo storico è uscito e quali erano i competitor in quel
        momento.
\end{itemize}

% ---------------------------------------------------------------------
\section{Che cosa fa un game designer}
\label{sec:cosa-fa}

Per arrivare a una prima definizione, la classe ha condotto un brainstorming:
ognuno doveva indicare una cosa che, secondo lui, fa il game designer (senza
ripetere quelle già dette). Ne è emersa la lista seguente, commentata punto
per punto.

\subsection{Progettare meccaniche e struttura del gioco}

Il ruolo più generico è quello del \textbf{core mechanic designer}, detto
anche \emph{gameplay designer} e, un tempo, \emph{black box designer}: si
occupa delle \textbf{meccaniche di gioco}. Si va dal piccolo al grande (in
funzione anche della dimensione dell'azienda):
\begin{itemize}
  \item nel \textbf{piccolo} ci sono le singole meccaniche;
  \item nel \textbf{grande} ci sono i \textbf{sistemi}. Nei giochi
        \emph{sistemici}, con molti sistemi interconnessi (ad esempio
        \emph{GTA}), esistono specializzazioni dedicate a bilanciare e
        supervisionare i sistemi: è il \emph{system design}.
\end{itemize}

\paragraph{Quando un gioco è ``sistemico''?}
Non dipende solo da come le meccaniche interagiscono con il giocatore (le
\emph{dinamiche}, secondo i framework che passano dalle meccaniche alle
dinamiche, che verranno studiati più avanti), ma anche da \emph{quante}
meccaniche ci sono e da quanto interagiscono tra loro. Non diremmo che
\emph{Super Mario} sia un gioco sistemico, eppure contiene dei sistemi: la
linea è sottile. Quando i sistemi aumentano molto, sono complessi e si
influenzano a vicenda, il gioco diventa sistemico. Esempio: in \emph{The
Legend of Zelda: Breath of the Wild} il vento modifica quanto il personaggio
sale o scende; il fuoco può incendiare l'erba; l'erba prende fuoco se colpita
da una freccia infuocata; e così via. Nel mercato moderno i giochi sistemici
sono spesso, ma non sempre, legati a mappe \emph{open world}.

\subsection{Le 3C: camera, character, control}

La base del gioco, da cui generalmente si inizia a lavorare (anche nei
prototipi), è la ``sacra trinità videoludica'' delle \textbf{3C}:
\begin{description}
  \item[Camera] che tipo di inquadratura ha il gioco: prima persona, terza
        persona, isometrica, ecc.
  \item[Character] quanto è grande il personaggio, che tipo di personaggio è
        e, di conseguenza, quali azioni può compiere.
  \item[Control] i controlli: che cosa il personaggio può fare, quali sono i
        \emph{verbi} di gioco.
\end{description}
Le 3C sono definite dal designer, ed esiste anche una specializzazione per
questo lavoro. Esistono varianti: quella del direttore del corso è
\textbf{4C + J}, cioè \emph{camera, character, control, conflict} (conflitto)
e \emph{juice} (il ``succo'', il feedback che rende piacevole l'interazione).
Verranno approfondite più avanti.

\subsection{Il personaggio}

Il designer definisce il personaggio dal punto di vista delle meccaniche
(dimensioni, azioni, capacità), all'interno delle 3C. L'aspetto estetico è
invece competenza di altri dipartimenti.

\subsection{Documentare le informazioni per gli altri dipartimenti}

Assolutamente vero: la documentazione è una delle prime attività del
designer. Più l'azienda è grande, più questo lavoro diventa prominente; più è
piccola, meno è rilevante. In \textbf{pre-produzione}, però, è sempre
fondamentale.

\subsection{Mediatore e leader ``al servizio''}

Il designer è un \textbf{mediatore} nei conflitti interni al team, una sorta
di collante che si relaziona con dipartimenti diversi (non sempre gli stessi,
a seconda della fase di produzione). A volte fa anche da leader, a seconda di
come sono gestite la pipeline e la gerarchia.

Questa leadership è difficile da definire perché \textbf{non è verticale}: è
una leadership \emph{al servizio} (\emph{servant leadership}). Il designer
``è asservito al programmatore'': il suo compito non è arrivare con un'idea e
scaricarla addosso agli altri, né portare la soluzione e lasciare il problema
a qualcun altro, ma \textbf{seguire il processo a basso livello}, sporcandosi
le mani.

\paragraph{I feature team.} Nei team medio-piccoli o medio-grandi possono
esistere sotto-team (\emph{feature team}) dedicati a un obiettivo specifico:
ad esempio un designer incaricato del character o della camera, affiancato da
uno o due programmatori e da un artista. In questi casi il designer
supervisiona il lavoro e gestisce la produzione del sotto-team.

\subsection{Gestire l'avanzamento del prodotto}

Esiste una specializzazione che si occupa solo di questo, il
\textbf{producer}, a un livello gerarchico completamente diverso (e che
potrebbe essere una strada percorribile al termine del percorso di studi).
All'interno di un feature team anche il designer può farlo, ma non è detto.
Può certamente \emph{dare la direzione}: nella sua documentazione può
indicare che per arrivare a uno stato giocabile del prototipo serve un certo
lavoro di arte e di programmazione, per lo stato successivo un altro, e così
via. In ambito produttivo esistono metodologie strutturate come
\textbf{Agile}, che verranno viste in seguito.

\subsection{Presentare le idee}

Il pitching (sezione~\ref{sec:gerarchia-corso}) è richiesto al designer sia
in pre-produzione sia in produzione, in misura che dipende dall'azienda. Ad
esempio Ubisoft lascia ai designer, e in realtà a qualsiasi membro, la
possibilità di \emph{pitchare} un'idea prima di una nuova produzione, perché
possa iniziare a prendere vita.

\subsection{Bozze e prototipi}

Le bozze vanno bene, ma devono essere \textbf{chiare}. C'è un'eterna lotta tra
documentazione e gioco giocabile: se si fa solo documentazione il gioco non
avanza; se si fa solo il gioco manca la documentazione e si perde la visione.
Sono due cose \textbf{complementari}.

Il modo per comunicare un'idea in forma grezza ma giocabile è il
\textbf{prototipo}. Un prototipo può essere qualsiasi cosa: di carta, in
Unreal con Blueprint ``brutti''. \emph{Non deve essere bello, deve essere
giocabile}.

\paragraph{Esempio: il prototipo di carta di Mario + Rabbids.} Durante lo
sviluppo di \emph{Mario + Rabbids}, i programmatori di Ubisoft erano impegnati
su altri progetti (come \emph{Ghost Recon}). I designer realizzarono allora
un \textbf{gioco da tavolo} che riproduceva le meccaniche del gioco, quelle
di uno strategico a turni sul modello di \emph{XCOM}. E funzionava.
L'obiettivo era \textbf{validare un'idea}, validare un gameplay.

Altri modi per rendere chiara un'idea sono i \emph{mock-up} nella
documentazione e le immagini modificate o i \emph{mash-up} di qualsiasi
natura. Lato documentazione l'importante è che sia chiaro; lato prototipo,
che sia giocabile.

\subsection{Fare ricerca per l'idea, e poi smontarla}

Il designer ha una componente creativa importante e fa ricerca per arrivare a
un'idea. Ma ciò che deve poi fare necessariamente è \textbf{smontarla}. Il
suo lavoro è scientifico, esattamente come il \textbf{metodo scientifico}, il
cui obiettivo è \emph{confutare un'ipotesi}: si arriva a un'ipotesi tramite
sperimentazione e poi si cerca $n$ volte di confutarla; se resiste, la si
assume come valida \emph{in un certo contesto} (le regole della fisica
classica, ad esempio, non valgono nella meccanica quantistica).

Allo stesso modo il designer arriva a un'idea (tramite ricerca o chiedendo
agli altri: le idee le hanno tutti, non sono un privilegio del designer),
cerca di smontarla e poi cerca di \textbf{metterla a terra}. Se l'idea non è
applicabile, non è implementabile, finché non si trova un modo per
realizzarla va scartata. Può essere l'idea più bella dell'universo, ma
\textbf{se non può essere messa a terra, non esiste}.

\paragraph{L'originalità.} Serve un'idea originale? Dipende dall'obiettivo e
dalla visione del progetto:
\begin{itemize}
  \item se l'obiettivo è vendere molto, un'idea originale non serve
        necessariamente;
  \item se si vuole creare un nuovo filone, serve un'idea, se non originale,
        almeno interessante: può essere un discorso di marketing o una
        meccanica innovativa, con il rischio di non vendere.
\end{itemize}
Il ruolo del \emph{director} e il concetto di \emph{visione} verranno
approfonditi in seguito.

\subsection{Analizzare i competitor}

Arrivati a un'idea (``vogliamo un certo tipo di combattimento''), si va a
vedere se qualcuno l'ha già fatta e \emph{come}. Le meccaniche non sono
privilegio di nessuno, salvo quando vengono brevettate, come il
\textbf{Nemesis System} di \emph{La Terra di Mezzo: L'Ombra di Mordor}.

\paragraph{Esempio: il combattimento da ridimensionare.} In un progetto
(presso Stormind Games) il team si trovò con un sistema di combattimento
troppo costoso: mancavano le risorse e gli animatori per realizzarlo. Poco
output lato gameplay significava non poter fare qualcosa di complesso e
profondo come in \emph{Elden Ring}; al contrario, il team aveva molti
artisti e poteva puntare su qualcosa di più artistico. Si iniziò così a
orbitare intorno al corpo a corpo di \emph{The Last of Us} (parte I e II).
Sul canale della GDC esistono diversi talk che raccontano cosa è stato fatto
e come, con successi, insuccessi e costi (ad esempio del \emph{motion
capture}): da quei talk si ricavarono moltissime informazioni utili.

I competitor possono essere:
\begin{itemize}
  \item \textbf{diretti}: un gioco molto simile al proprio;
  \item \textbf{di mercato}: un gioco che domina il mercato in cui ci si vuole
        affacciare;
  \item \textbf{indiretti}: un gioco con un sistema simile al proprio, ma
        tutto il resto diverso.
\end{itemize}
Regola d'oro: \textbf{là fuori, nascosto nell'ombra, c'è sicuramente un gioco
molto simile al vostro}. Va solo trovato. E se per qualche ragione non ha
venduto, bisogna capire perché.

\subsection{Lavorare in engine e fare tuning}

Il level designer lavora in engine, così come il technical designer; ma
qualsiasi game designer, a seconda di ciò che deve fare, può finire in
engine, soprattutto per il \textbf{bilanciamento} e il \textbf{tuning}.

Il \emph{tuning} è l'azione di bilanciare con precisione i dati. Esempio: in
\emph{Super Mario} si modifica quanto in alto salta il personaggio, con che
velocità, in quanto tempo raggiunge l'apice, e così via. Questo lavoro rende
il gioco:
\begin{itemize}
  \item ad alto livello, dotato della \textbf{progressione} giusta;
  \item a basso livello, \textbf{piacevole da giocare}.
\end{itemize}
Il \emph{fine tuning} è un lavoro di dettaglio, paziente e minuzioso.

\subsection{Bilanciare}

Il bilanciamento, emerso tardi nel brainstorming, è una delle attività
centrali del designer e ha anche una specializzazione dedicata
(l'\emph{economy designer}, sezione~\ref{sec:tipi-designer}).

\subsection{Analizzare il mercato}

In pre-produzione, di solito, sì: capire le tendenze, i generi che cambiano
nel tempo, la compatibilità del gioco con il momento storico, in vista ad
esempio di un possibile publisher. Va capito quanto, come e con quanto tempo
farlo, perché si è sempre limitati da molte cose.

\subsection{Sintesi}

\begin{quote}
Il designer parte dall'idea ma non si ferma lì: la confuta, lavora sui
dettagli, bilancia, fa \emph{fine tuning}; lavora con il team e fa da
\emph{manager} per diversi dipartimenti; e naviga costantemente tra i limiti.
\end{quote}

% ---------------------------------------------------------------------
\section{I limiti del design}
\label{sec:limiti}

I limiti ci sono sempre, e nel mondo del design sono spesso numerosi. I più
importanti entro cui navigare sono:
\begin{description}
  \item[Tempo] ad alto livello, l'idea deve essere realizzabile nel tempo di
        produzione disponibile; a basso livello, il proprio lavoro
        settimanale deve essere fattibile nel tempo che si ha. Non ci si può
        dividere tra mille feature team se realisticamente se ne può seguire
        uno solo.
  \item[Tecnologia] si può o non si può fare qualcosa? Ci sono i server per un
        multiplayer con un certo numero di giocatori? Ci si può permettere,
        ad esempio, una fauna che reagisce alla fisica come quella di
        \emph{GTA VI}, che richiede il meglio dei \emph{technical artist}?
  \item[Background e competenze] che cosa sa fare il team? Realisticamente,
        chi ha lavorato dieci anni su giochi horror viene assunto per un
        horror, perché ha le competenze per risolvere problemi già visti.
  \item[Persone e motivazione] il team ha la motivazione e le competenze per
        fare un certo tipo di gioco, e per lavorare insieme?
  \item[Denaro] tutto ruota intorno ai soldi. Che si lavori da soli o in mille,
        le persone vanno pagate e il gioco deve avere un ritorno economico.
        Bisogna decidere dove tracciare la linea: raggiungere il
        \textbf{break-even} (guadagnare esattamente quanto speso) o puntare a
        grandi guadagni. Con un \textbf{publisher}, inoltre, una parte
        consistente dei ricavi va a lui (ripartizioni dell'ordine di 80/20 o
        70/30, che dipendono dal contratto firmato).
\end{description}

\subsection{Non siete ``gli uomini delle idee''}

Per qualcuno può essere un problema scoprirlo, ma il designer non è
``l'uomo (o la donna) delle idee'': accanto alla parte creativa c'è una parte
tecnica, e si è fortemente limitati da molti fattori. Questi limiti, però,
sono un aiuto:
\begin{itemize}
  \item riducono il ventaglio di possibilità potenzialmente distruttive per il
        team (cose che non si riescono a fare, che non piacciono o per cui
        mancano le risorse);
  \item aiutano a mirare bene, trovando soluzioni creative in uno spazio più
        piccolo e più attuabile.
\end{itemize}

\begin{quote}
\centering
\textbf{Le idee sono nulla, l'implementazione è tutto.}
\end{quote}

% ---------------------------------------------------------------------
\section{Ruoli e specializzazioni}
\label{sec:tipi-designer}

\subsection{La scala gerarchica}

\begin{description}
  \item[Creative Director (CD)] mantiene la \textbf{visione} e dà l'idea, in
        modo che tutto segua una certa direzione. In teoria può diventarlo
        chiunque, ma con un'esperienza notevole, nell'ordine dei decenni: si
        assume che in quel tempo abbia visto come lavorano molti dipartimenti.
        Naturalmente un CD con un background artistico ha meno competenze di
        design, e viceversa.
  \item[Director di disciplina] uno per dipartimento: \emph{game director},
        \emph{level director}, \emph{art director}, ecc.
  \item[Lead] il capo del dipartimento, che gestisce i gruppi di persone. Chi
        lavora in un dipartimento risponde al proprio lead (ad esempio un game
        designer risponde al \emph{lead designer}, non al game director).
  \item[Senior / Principal] i profili più esperti (a seconda dell'azienda).
  \item[Junior] il profilo di inizio carriera.
  \item[Intern] lo stagista.
\end{description}
Lo schema si replica uno a uno per ogni dipartimento: arte, suono, animazione
(a volte arte e animazione insieme, o \emph{rigging} e animazione insieme),
programmazione, UX/UI e così via. A fare la differenza, a parità di ruolo, è
il \textbf{livello di seniority}.

\subsection{UX e UI}

UX e UI stanno a metà tra game design e arte: a seconda dell'azienda il ruolo
si chiama \emph{UX/UI Artist} o \emph{UX/UI Designer}.
\begin{description}
  \item[UX (User Experience)] come una persona interagisce con le cose. Esempio:
        su uno smartphone, un pulsante accende e spegne lo schermo; per
        digitare la password si scorre verso l'alto e compaiono i punti da
        toccare.
  \item[UI (User Interface)] come sono ``vestite'' queste cose: ad esempio
        l'aspetto della barra dei punti vita (HP), realizzata con asset 2D o
        3D.
\end{description}

\subsection{Le specializzazioni del game design}

Più l'azienda è grande, più la specializzazione è puntuale.
\begin{description}
  \item[Gameplay / Core mechanics designer] (sono la stessa cosa) si occupa
        delle meccaniche e delle 3C.
  \item[System designer] si occupa dei sistemi e delle loro interazioni.
  \item[Economy designer] si occupa del bilanciamento; lavora soprattutto con
        fogli di calcolo (Excel, Google Sheets).
  \item[Combat designer] si occupa del combattimento; negli annunci di lavoro
        indica quasi sempre un combattimento \emph{action}, con gameplay
        reattivo (non strategico), anche se non è una regola.
  \item[Multiplayer designer] si occupa delle interazioni multigiocatore.
  \item[Technical designer] fa da ponte tra game design e programmazione.
\end{description}

\subsection{Le specializzazioni del level design}

\begin{description}
  \item[Encounter designer] posiziona i nemici nella mappa in modo che gli
        scontri siano bilanciati e la progressione equilibrata come
        desiderato.
  \item[World / map designer] popola il mondo di gioco.
  \item[Quest designer] popola il mondo con le missioni e pensa alla
        progressione al loro interno, cioè a dove il giocatore deve andare (si
        vedano ad esempio gli annunci di lavoro di CD Projekt Red o di
        Larian).
  \item[Technical level designer] fa da ponte tra level design e
        programmazione.
\end{description}

\subsection{Discipline vicine: suono, UI, produzione}

Il monte ore del corso non permette di uscire come \emph{narrative designer}
(o \emph{sound designer}) completi. Ciò non toglie che si possano coltivare
queste discipline:
\begin{itemize}
  \item nei progetti di gruppo si può occuparsi anche del suono, della UI/UX,
        o provare ad applicare metodologie come Agile per fare
        \emph{producing} e gestire il team;
  \item il tempo dedicato a queste attività va però sottratto al game design,
        e dipende da quante persone ci sono nel team, da quanto lavoro c'è e
        da quanto è grande il progetto;
  \item avere sviluppato e rilasciato un progetto, anche piccolo (anche una
        \emph{game jam}), documentando i propri processi, è già un ottimo
        argomento da presentare in futuro;
  \item la via migliore resta lavorare nell'industria e, da lì, chiedere di
        imparare anche altre discipline.
\end{itemize}
Un gioco si può fare su qualsiasi cosa: sugli scarafaggi, sui piedini delle
sedie, sui pennarelli. Si parte avvantaggiati da ciò che è vicino alla propria
sensibilità: chi ama la musica può pensare a un gioco con meccaniche legate
alla musica, con un gameplay innovativo ma contenuto e realizzabile. Il design
è un cappello che comprende molte discipline.

% ---------------------------------------------------------------------
\section{Approfondimenti emersi in lezione}
\label{sec:approfondimenti}

\subsection{Roguelike e roguelite}

Un esempio di come il marketing conia generi: Supergiant Games, con
\emph{Hades}, rese celebre il termine \textbf{roguelite} al posto di
\emph{roguelike} (pur non essendo il primo roguelite in assoluto). La
differenza:
\begin{itemize}
  \item il \textbf{roguelike} si fonda su una serie di punti fissati in una
        sorta di \emph{manifesto} che definisce che cosa rende un gioco
        roguelike; uno su tutti: a ogni morte si ricomincia la \emph{run} da
        capo;
  \item il \textbf{roguelite} rompe alcuni di questi punti, in particolare
        introducendo una \textbf{progressione esterna alla run}: il giocatore
        può potenziarsi fuori dalla singola partita per vincere più
        facilmente in quelle successive. \emph{Hades} ne è il caso principe.
\end{itemize}
Analogamente, un mobile game sviluppato dal direttore del corso fu
promosso come ``il primo \emph{donation game} al mondo'': un gioco che
convertiva il tempo di gioco in donazioni reali a enti benefici per animali.
Anche qui un neologismo creato lato marketing per definire un genere.

\subsection{Applicazioni fuori dall'industria videoludica}

Le competenze sviluppate non servono solo ai videogiochi:
\begin{itemize}
  \item i \textbf{render architettonici} oggi possono essere realizzati con
        Unreal, e diversi professionisti formati nel game design lavorano in
        questo campo;
  \item la \textbf{gamification}, cioè l'applicazione di regole videoludiche
        a contesti non ludici (mostre interattive, progetti culturali,
        consulenze come quella per applicare questi concetti alle macchine da
        allenamento di un'azienda di fitness), è in Italia uno dei modi più
        redditizi di spendere queste competenze, spesso con finanziamenti
        pubblici.
\end{itemize}

% ---------------------------------------------------------------------
\section{Il primo task}
\label{sec:primo-task}

\begin{enumerate}
  \item Creare un indirizzo e-mail \textbf{professionale} (consigliato: un indirizzo Gmail nella forma
        \texttt{nome.cognome@gmail.com}; l'anno di nascita è accettabile).
        Va bene anche un indirizzo professionale già esistente usato per
        lavoro. Da evitare indirizzi poco seri (``non \texttt{blackfire67}'').
  \item Una persona della classe raccoglie tutti gli indirizzi in un unico
        documento e lo invia via e-mail al docente, che aggiungerà tutti al
        Drive condiviso del corso (dove verranno caricate lezioni, slide e
        materiali, per avere i documenti tutti nello stesso posto).
  \item Installare Discord e iscriversi al server del corso, dove passano le
        comunicazioni ufficiali. È consuetudine, e utile, creare anche un
        gruppo tra studenti.
\end{enumerate}
\textbf{Scadenza}: prima della lezione del venerdì.

% ---------------------------------------------------------------------
\section*{Punti chiave del capitolo}
\addcontentsline{toc}{section}{Punti chiave del capitolo}

\begin{itemize}
  \item Il corso è una \textbf{simulazione di lavoro}: ci si tratta da
        colleghi, con una gerarchia da studio di sviluppo.
  \item Si impara con l'\textbf{inquiry-based learning}: lezione, task,
        dubbi, domande, condivisione, progetti.
  \item Il fallimento affrontato con la mentalità giusta è la principale fonte
        di crescita; i feedback riguardano il lavoro, mai la persona.
  \item Il lavoro del designer è per metà creativo e per metà
        \textbf{ricerca e metodo scientifico}; il game design non è una
        scienza esatta.
  \item \textbf{Documentazione} e \textbf{pitching} sono il cuore del
        mestiere.
  \item Il designer progetta meccaniche e sistemi (partendo dalle
        \textbf{3C}), documenta, media, prototipa, confuta le idee, studia i
        competitor, bilancia e fa tuning.
  \item La leadership del designer è \textbf{al servizio} del team, non
        verticale.
  \item I limiti (tempo, tecnologia, competenze, persone, denaro) restringono
        e indirizzano la creatività.
  \item \textbf{Le idee sono nulla, l'implementazione è tutto.}
\end{itemize}

%    % ...
%
%  (Con hyperref aggiungere anche \renewcommand{\theHchapter}{z.\arabic{chapter}}
%  nella parte zero, per evitare ancore duplicate.)
%  Tutte le etichette di questo capitolo hanno il prefisso "z1:" per non
%  collidere con quelle dei capitoli del corso.
% =====================================================================

\chapter[La cultura del feedback]{Lavorare in squadra: la cultura del feedback}
\label{cap:z1-feedback}

Questo capitolo apre le \emph{lezioni zero}: incontri che precedono l'inizio
vero e proprio del corso di Game Design e che non riguardano una tecnica di
design in senso stretto, ma \textbf{la figura professionale} di chi lavora
nell'industria dei videogiochi. Prima di imparare a progettare un gioco
conviene infatti capire in che tipo di ambiente lo si progetterà: un ambiente
fatto di team, di gerarchie, di scadenze e, soprattutto, di persone che devono
riuscire a lavorare insieme.

La lezione è la prima giornata di un \textbf{workshop di due giorni} (circa
dieci ore in totale, pause pranzo comprese) tenuto da Stefano
Cozzi\footnote{Il cognome è riportato come compare nella trascrizione
automatica, che in più punti storpia i nomi propri; andrebbe verificato.
Nella stessa trascrizione lo studio per cui lavora appare come
``Stormine'', ``Stormline'' e ``Stormind'': si tratta di Stormind Games,
studio con sede a Catania, coerentemente con quanto detto in lezione.}, un
professionista con circa quindici anni di esperienza nel settore, tra
Blizzard Entertainment, Riot Games e, oggi, Stormind Games. Il workshop è
nato da un confronto con i coordinatori dell'accademia\footnote{Nella
trascrizione l'accademia è resa in modi diversi (``Hive'', ``a dive'',
``D.I.V.''); i coordinatori citati sono Enrico e Andrea, insieme al
direttore Diego Ricchiuti. Nel resto del capitolo si usa semplicemente
``l'accademia'', come negli altri capitoli.}: serviva un momento dedicato a
\emph{come ci si comporta in un ambiente professionale}, cosa che riguarda sia
il lavoro futuro sia i progetti di gruppo svolti durante il percorso di
studi, dove ci si deve già abituare a comportarsi ``come se si fosse in
un'azienda''.

La prima giornata ha un unico grande tema, sezionato in più parti:
\textbf{la cultura del feedback}. La seconda giornata (che sarà oggetto della
prossima lezione zero) sarà più pratica: \emph{time management},
\emph{brainstorming} e il ``guidare da dentro'' un team anche senza esserne
il leader (sezione~\ref{sec:z1-domani}).

Il percorso del capitolo segue quello della lezione:
\begin{enumerate}
  \item il formato del workshop e il percorso professionale del relatore
        (sezione~\ref{sec:z1-relatore});
  \item perché le aziende tengono tanto al lavoro di squadra e al
        \emph{fit} culturale (sezione~\ref{sec:z1-squadra});
  \item che cos'è il feedback e da che cosa dipende il modo di darlo
        (sezione~\ref{sec:z1-cose-feedback});
  \item le cinque fasi di vita di un team e il valore del
        \emph{post-mortem} (sezioni~\ref{sec:z1-fasi-team}
        e~\ref{sec:z1-qa-fasi});
  \item il confronto con i colleghi e alcune domande sul lavoro nei grandi
        studi (sezioni~\ref{sec:z1-rivalita} e~\ref{sec:z1-industria});
  \item il circolo virtuoso tra feedback, miglioramento e fiducia, e i suoi
        effetti sulla \emph{retention} e sulla cultura aziendale
        (sezioni~\ref{sec:z1-triangolo} e~\ref{sec:z1-cultura});
  \item il ``gioco del traditore'' e il principio del
        \textbf{default to trust} (sezioni~\ref{sec:z1-among-us}
        e~\ref{sec:z1-default-trust});
  \item le traiettorie di carriera e la valutazione delle performance
        (sezione~\ref{sec:z1-traiettorie});
  \item come si dà, come si riceve e quando si dà un feedback
        (sezioni~\ref{sec:z1-dare}, \ref{sec:z1-ricevere}
        e~\ref{sec:z1-quando});
  \item il feedback come regalo, le domande finali e le anticipazioni
        della seconda giornata (sezioni~\ref{sec:z1-regalo},
        \ref{sec:z1-qa-finale} e~\ref{sec:z1-domani}).
\end{enumerate}

Una premessa di metodo: il relatore ha alternato continuamente teoria,
domande alla classe e racconti della propria esperienza. I racconti non sono
digressioni: sono gli esempi concreti con cui i principi vengono resi
tangibili, e per questo sono riportati con cura.

% ---------------------------------------------------------------------
\section{Il workshop e il relatore}
\label{sec:z1-relatore}

\subsection{Formato e regole di partecipazione}

Il relatore ha cercato di ``riassumere quindici anni di esperienza in dieci
ore''. Ha materiale per circa il doppio delle ore disponibili: le cose
strettamente necessarie vengono trattate per prime, il resto è tenuto ``in
riserva'' come \emph{nice to have}. Le regole del workshop:
\begin{itemize}
  \item le domande si fanno \textbf{alla fine di ogni slide}, così da chiudere
        prima il discorso; non c'è un'unica sessione finale di domande e
        risposte (anche se ci sarà), perché il relatore preferisce fermarsi
        spesso e non fare ``uno spiegone di dieci ore'';
  \item le domande non devono per forza riguardare la slide appena vista: se
        una slide dà lo spunto per altro, se ne può parlare; al massimo sarà
        il relatore a rispondere in breve o a rimandare a una slide
        successiva;
  \item si può chiedere anche di esperienze passate, opinioni e sensazioni;
        se per NDA o per serietà professionale non può rispondere, lo dirà
        esplicitamente (``non lasciate che sia un \emph{blocker}'');
  \item su Zoom si usa l'alzata di mano o la chat, che il relatore tiene
        aperta sul secondo monitor;
  \item non ci sono pause ufficiali: la lezione è registrata e ognuno è
        libero di assentarsi quando serve, cercando però di non disturbare il
        flusso della lezione;
  \item le slide verranno fornite integralmente tramite la segreteria
        dell'accademia (Vanessa), insieme alle registrazioni: le slide senza
        il parlato servono poco, ma sono un buon promemoria.
\end{itemize}

\subsection{Il percorso professionale}

Il relatore si presenta (scherzosamente) come la ``guida escursionistica''
che accompagna la classe nella ``foresta incantata'' della cultura del
feedback, perché per molti è un territorio sconosciuto. Il suo percorso:
\begin{itemize}
  \item in Italia ha lavorato per aziende che fornivano servizi a terzi, tra
        cui il lavoro su \emph{Warhammer Online}, e per un canale televisivo
        sul digitale terrestre dedicato ai videogiochi\footnote{Nella
        trascrizione: ``Digital Game Network Channel''. Il nome esatto del
        canale andrebbe verificato.};
  \item a \textbf{Blizzard Entertainment} (circa sei anni, a Parigi) ha fatto
        il \emph{community manager} per l'Italia (\emph{StarCraft~II},
        \emph{World of Warcraft} da \emph{Mists of Pandaria} in poi,
        \emph{Hearthstone}, \emph{Diablo~III}, le espansioni di
        \emph{StarCraft}, il lancio di Battle.net), per poi diventare il
        manager dei community manager europei;
  \item a \textbf{Riot Games} (circa cinque anni, a Dublino) è stato
        brevemente community manager italiano e poi \emph{publishing
        manager} di un team di circa trenta persone per l'Europa. Il
        \emph{publishing} non è solo community: è l'unione di
        \textbf{community, marketing e PR};
  \item a \textbf{PG Esports} si è occupato di eventi dal vivo e online
        (ad esempio la Milano Games Week e tornei di \emph{Valorant} e
        \emph{League of Legends}), restando in contatto con Riot e Blizzard;
  \item da circa un anno è tornato nello sviluppo, a \textbf{Stormind
        Games}, dove oggi si occupa di risorse umane e organizzazione del
        lavoro; lo studio verrà presentato più ampiamente al \emph{Career
        Day} dell'accademia insieme alla sua HR manager;
  \item fa inoltre consulenze aziendali anche al di fuori dei videogiochi e
        insegna in master universitari (IULM di Milano).
\end{itemize}

Il suo ruolo ricorrente, sia a Blizzard sia a Riot, è stato quello di
\emph{Mr.~Wolf} (``per chi ha visto \emph{Pulp Fiction}''): quello che
risolve i problemi. Veniva inserito nei team più in difficoltà, ad esempio
dal punto di vista delle performance, per sistemarli e passare al successivo.
Era presente in alcuni dei momenti più critici, come il lancio di
\emph{Diablo~III} o il lancio di espansioni di \emph{World of
Warcraft}\footnote{La trascrizione riporta ``il lancio con i Gary's'' e
``World of Trenor'': si tratta verosimilmente di \emph{Mists of Pandaria} e
di \emph{Warlords of Draenor}.}, a fare \emph{damage control}.

La scoperta fondamentale di questi anni è che i problemi da risolvere
\textbf{quasi mai riguardavano il prodotto} (non era certo lui a sistemare i
server di \emph{Diablo~III}), ma i team che c'erano dietro: persone che
litigavano, non collaboravano, non riuscivano a stare nella stessa stanza
senza urlarsi contro. E nel 90\% dei casi tutto questo si riduce a
\textbf{problemi di comunicazione}: è raro che perfino le incompatibilità
caratteriali non siano risolvibili parlando.

\subsection{Struttura non vuol dire burocrazia}

A Stormind il relatore sta accompagnando un team passato da 30 a 95 persone
in meno di un anno, attraverso tutti i cambiamenti di cui parla il workshop.
Darsi una struttura viene spesso percepito male, come ``più processi, più
procedure'', ma bisogna distinguere:
\begin{itemize}
  \item un processo come ``timbrare il cartellino'' può essere davvero
        invasivo e noioso;
  \item definire \emph{come funziona un documento di game design, chi lo
        approva, chi lo visiona e chi no} è invece una struttura necessaria.
\end{itemize}
Uno studio di 30 persone nella stessa sede può farne a meno: quando c'è un
dubbio ci si alza e si va alla scrivania del collega. Uno studio di 95
persone sparse per l'Italia (e in parte per l'Europa) non può più farlo:
servono una struttura diversa, un modo di lavorare diverso e una
pianificazione del futuro diversa.

% ---------------------------------------------------------------------
\section{Perché il lavoro è sempre lavoro di squadra}
\label{sec:z1-squadra}

\subsection{Non si va in ufficio per fare amicizia}

Il primo messaggio è netto: \textbf{non si va in ufficio per fare amici} o
per trovare compagni di vita. Si va in ufficio perché c'è un lavoro da fare, e
lo scopo è riuscire a farlo \emph{anche con chi sta antipatico}. Con le
persone simpatiche è facile: diventeranno amici, e non sarà mai un problema.
È con chi sta antipatico che le cose si complicano, ed è lì che emergono i
\textbf{veri professionisti}: quelli che riescono a lavorare con chiunque
senza che l'antipatia pesi sul modo in cui lavorano e collaborano, mettendo il
prodotto e il team davanti al proprio ego e alle proprie preferenze personali.

La frase che fa da filo conduttore a tutto il workshop (e che il relatore
promette di riproporre ``fino alla nausea'') è:

\begin{quote}
  \emph{La capacità di un gruppo di persone di fare cose straordinarie
  dipende dal modo in cui queste persone operano insieme.}
\end{quote}

Le aziende raramente pensano che l'innesto di una singola persona possa
cambiare il destino di un progetto. Questo è tanto più vero quanto più grande
è il team, ma vale anche per i piccoli studi, anche per un gruppo di cinque
persone. Come negli sport di squadra: l'individuo conta fino a un certo punto,
e anche se si ha Dončić in squadra, se si gioca uno contro cinque si perde
comunque\footnote{Il riferimento è alle finali NBA in corso nei giorni del
workshop tra Dallas Mavericks e Boston Celtics (giugno 2024), il che
permette di datare approssimativamente la lezione.}.

\subsection{Lo \emph{useless genius} e il costo delle persone tossiche}

Molti dipartimenti HR usano la definizione di \textbf{useless
genius}\footnote{Il concetto è più noto in letteratura manageriale come
\emph{brilliant jerk}; in lezione è usata l'espressione \emph{useless
genius}.}: si può essere bravissimi, a volte i migliori nel proprio lavoro,
ma se non ci si sa comportare, relazionare e operare in un gruppo si è
fondamentalmente inutili ai fini del gruppo. Nel migliore dei casi inutili,
nel peggiore \emph{dannosi}.

In passato le aziende tendevano a ``chiudere un occhio'': la persona è
fondamentale, la teniamo. Oggi la domanda è diventata: \emph{quanto costa
all'azienda, oltre al suo stipendio, tenere una persona che crea un ambiente
tossico e peggiora chi le sta intorno?} Ci sono costi sommersi che le aziende
cominciano a vedere, e la tendenza è liberarsi di queste persone il prima
possibile. Il ragionamento è semplice:
\begin{itemize}
  \item le \textbf{competenze si possono trasmettere} o sostituire: si trova
        qualcuno che fa lo stesso lavoro magari nel doppio del tempo, oppure
        più persone che se lo dividono;
  \item il \textbf{carattere delle persone} invece cambia molto raramente.
\end{itemize}

\subsection{Il fit culturale}

Quando qualcuno \textbf{non è un \emph{fit} culturale} per un'azienda è
meglio separarsi presto. Ma non essere adatti a \emph{quell'}azienda non
significa non essere adatti a lavorare: è perfettamente legittimo che una
persona inadatta a un'azienda sia adattissima a un'altra.

Il relatore l'ha vissuto in prima persona. Dopo sei anni a Blizzard viene
contattato da Riot e affronta un processo di selezione lunghissimo (``quasi
undici colloqui''), al termine del quale viene scartato. Scrive allora una
mail che gli ha ``cambiato la traiettoria di carriera'': chiede
semplicemente perché non è stato ritenuto adatto. La risposta: in nessun
colloquio si era esposto abbastanza da far capire \emph{perché} volesse
lasciare Blizzard, e quindi pensavano che in realtà stesse bene dov'era. Il
motivo era che, per educazione professionale, non gli piace parlare male del
proprio posto di lavoro; portata all'eccesso, però, questa riservatezza era
diventata un ``è tutto bellissimo, ma vorrei provare Riot'', senza un perché.
Lo invitano comunque a un incontro dal vivo (``dal vivo è diverso''), e alla
fine viene assunto.

Una volta arrivato, i sei mesi di prova (\emph{probation}) sono serviti tutti
per smettere di ragionare e lavorare ``come a Blizzard''. La differenza
culturale era enorme, tanto che in Riot esisteva un gruppo di supporto
interno tra ex dipendenti Blizzard: un collega gliel'ha descritto come
``uscire da una relazione tossica e doversi riabituare a una relazione
normale''\footnote{Il relatore precisa che parla della Blizzard di allora,
che non esiste più: oggi è parte di Microsoft e nessuno dei suoi ex colleghi
vi lavora ancora, anche a causa dei licenziamenti di massa.}.
\begin{description}
  \item[Blizzard] era \emph{developer first}: il focus era chiaramente sugli
        sviluppatori, prima di tutti gli altri. Dopo sei anni c'erano ancora
        meeting a cui non poteva partecipare perché non abbastanza
        \emph{senior}, informazioni che non poteva sapere, piani che non
        poteva vedere.
  \item[Riot] era \emph{player first}: prima i giocatori, e i dipendenti come
        mezzo per trasferire il feedback dei giocatori dentro il gioco. Il
        secondo giorno lo portano in una stanza e gli mostrano il piano dei
        cinque anni successivi. Lui, sorpreso, pensa a un errore (``sono qui
        da oggi, non vorrei mettervi nei guai''); la risposta è: ``sei un
        dipendente Riot o no? Allora non ci sono livelli di accesso: devi
        sapere queste cose per fare meglio il tuo lavoro''.
\end{description}

Da questa esperienza ha imparato molto sulla \textbf{trasparenza}: condividere
le informazioni aiuta il lavoro. Non sempre, però: a volte non si condivide
qualcosa non per segretezza, ma perché non è ancora pronta per essere
presentata, esattamente come non si mostra un gioco quando è ancora solo
teoria (``te lo faccio vedere quando ho qualcosa di tangibile'').

La lezione per gli studenti è doppia: le aziende valuteranno il \emph{fit}
culturale sempre di più, ma capire se si è un \emph{fit} richiede anche
\textbf{lavoro da parte del candidato}. Bisogna avere la presenza di spirito
di dire: ``conosco i valori di questa azienda, non mando nemmeno il CV,
perché so che durerei un mese e mezzo''.

% ---------------------------------------------------------------------
\section{Che cos'è il feedback}
\label{sec:z1-cose-feedback}

\subsection{Le prime risposte della classe}

La prima domanda alla classe è: \emph{quando sentite la parola
``feedback'', a cosa pensate?} Le risposte raccolte, a voce e in chat:
\begin{itemize}
  \item un \textbf{riscontro} su ciò che si fa e su come lo si fa, sul lavoro
        o sul modo in cui ci si interfaccia con il gruppo; può essere
        lavorativo oppure comunicativo e comportamentale. Chi ha risposto lo
        vede come qualcosa che, se l'azienda non dà, va attivamente
        richiesto, perché è importante sapere come sta andando per
        migliorarsi;
  \item né positivo né negativo, ma \textbf{costruttivo}: in entrambi i casi
        serve a migliorarsi;
  \item uno \emph{scambio di opinioni}, un \emph{riscontro produttivo};
  \item un feedback ``di reazione'', non verbale: si fa una cosa e si
        osserva come cambia l'ambiente o il comportamento delle persone
        intorno (anche perché a voce la gente spesso si tiene le cose
        negative);
  \item per un designer, anche il \textbf{linguaggio del corpo}: osservando
        un giocatore giocare si può dedurre l'emozione che prova (frustrazione,
        incomprensione).
\end{itemize}

Sulla forma (scritto o orale?) le opinioni divergono: c'è chi lo preferisce
\emph{mille volte} scritto; chi pensa che dipenda dalle dimensioni
dell'azienda (orale nelle piccole, scritto in quelle con team di centinaia di
persone); chi pensa che dipenda dall'importanza del feedback (scritto se
tecnico o pieno di dati, da tenere sotto gli occhi; orale se è un consiglio
veloce).

Il relatore usa due esempi per mostrare che \emph{dipende}:
\begin{itemize}
  \item un feedback sulla programmazione ha senso darlo scritto, magari
        commentando direttamente il codice, più che ``andando a prendere un
        caffè'';
  \item ma se il feedback è ``ieri, durante il meeting, ho percepito male il
        fatto che mi interrompessi di continuo'', ce lo si aspetta
        \textbf{orale}.
\end{itemize}

\subsection{Da che cosa dipende il modo di dare feedback}

Lo \emph{spoiler} della giornata è che \textbf{non esiste una risposta univoca}
su come dare feedback. Il workshop non insegna ``il'' modo, ma fornisce
\emph{tutte le opzioni}, perché il modo giusto cambia in base a molti fattori:
\begin{itemize}
  \item a chi lo si dà: un collega, un superiore, il proprio capo, un proprio
        \emph{report} (un sottoposto);
  \item se è personale o professionale;
  \item se è la prima volta o la terza volta che si dà lo stesso feedback;
  \item se riguarda un progetto importante o secondario.
\end{itemize}

Manca però l'elemento fondamentale, che la classe individua subito:
\textbf{l'altro}. A parità di tutti gli altri fattori, il \emph{carattere
della persona che riceve il feedback} dovrebbe far cambiare il modo in cui lo
si dà. Per questo, prima della lezione, agli studenti era stato chiesto di
fare un test di personalità: non un test di psicologia avanzata, ma uno
strumento per rendersi conto di quanta varietà ci sia anche solo in una
classe di 53 persone (il che non vuol dire trovare 52 modi diversi di dare
feedback, ma lavorarci un po' sì).

Il relatore porta se stesso come esempio. Preferisce:
\begin{itemize}
  \item i feedback \textbf{positivi scritti}, perché è a disagio a ricevere
        complimenti in pubblico;
  \item i feedback sul proprio lavoro \textbf{diretti}: ``vai subito al
        punto e non girarmi intorno''.
\end{itemize}
È molto orgoglioso del suo lavoro e, quando qualcuno lo critica, la sua
primissima reazione istintiva è ``come ti permetti''. Negli anni ha imparato
a farla durare un secondo, ma chi gli dà feedback di persona dice di
vederla ancora negli occhi. Per mitigarla, dice esplicitamente a chi gli dà
feedback: \emph{``ignora il modo in cui reagisco nei primi cinque secondi,
vai avanti, poi mi passa''}. È considerato da tutti molto aperto al feedback
e non ha problemi a dire ``hai ragione, facciamo diversamente'': ma ci sono
reazioni istintive e caratteriali che è difficile controllare al cento per
cento.

\subsection{Con chi si lavora meglio}

Guardando le persone con cui si lavora abitualmente (anche nei progetti
dell'accademia), si noterà che si lavora meglio con due profili:
\begin{enumerate}
  \item chi è \textbf{esattamente come noi}, ma qui è un 50 e 50: alcuni
        caratteri vanno d'accordo con i propri simili, altri non vogliono
        nessuno con lo stesso carattere nel gruppo (e si potrebbe rientrare
        in questa categoria);
  \item chi è \textbf{affine}: ha qualcosa in comune con noi, ma non tutto,
        e proprio per questo può dare punti di vista differenti.
\end{enumerate}
Il lavoro della giornata è capire come valorizzare il secondo caso: ottenere
feedback da persone con caratteri, visioni e background diversi, per
completare una visione del proprio lavoro che è necessariamente limitata. È
come esplorare una mappa con la nebbia di guerra: un'altra persona, al di là
del muro, vede cose che noi non vediamo.

\subsection{Il feedback come somma di piccoli miglioramenti}

Per il relatore il feedback è questo: \textbf{applicare tanti piccoli
miglioramenti periodici}, che si sommano nel tempo. Migliorare dello 0,1\%
ogni giorno porta a fine anno a un miglioramento enorme, e questo vale per i
prodotti, per i giochi, per le persone e per i team.

\paragraph{L'esempio del ciclismo britannico.} La squadra di ciclismo
britannica, fino a una quindicina d'anni fa, non aveva vinto quasi nulla.
Cambiato il team manager, tutti si aspettavano grandi cambiamenti sportivi
(convocazioni, strategie di gara). Arrivarono invece \emph{minuscoli
cambiamenti non sportivi}: un altro chef in trasferta, materassi migliori nei
camper per riposare meglio, un sellino più comodo a fine tappa. Cambiamenti
marginali che da soli non avrebbero cambiato niente; eppure, dopo due o tre
anni di questo lavoro quotidiano, la squadra ha cominciato a vincere a
ripetizione, Olimpiadi comprese, cambiando il modo in cui il ciclismo è
visto in Gran Bretagna\footnote{È la nota strategia delle \emph{marginal
gains} associata a Dave Brailsford (British Cycling e Team Sky); il nome non
è citato in lezione.}. Certo, si è investito anche sulle persone giuste, ma
tutto è partito dai piccoli miglioramenti periodici.

Per gli studenti vale lo stesso: chiarire un conflitto con una persona con
cui si ha ancora poco a che fare, o una situazione che ha lasciato
``l'amaro in bocca'', sembra inutile oggi, ma darà i suoi frutti tra tre, sei,
nove mesi. Sono \emph{investimenti sul futuro}.

\paragraph{Peggiorare ogni giorno.} Esiste anche il percorso inverso:
lasciar correre, ``buttare sotto il tappeto'' le cose perché ``tanto è
inutile parlarne''. Queste cose tornano sempre fuori dopo otto mesi, come il
\emph{libro dei rancori} dei nani di Warhammer: ``ti ricordi quando, un anno
e mezzo fa, all'inizio del progetto, mi hai detto quella cosa? Mi sono
offeso''. A quel punto è quasi impossibile risolvere, perché nessuno ricorda
più con precisione cosa sia successo (``ma quale cena?'', ``forse era in
pizzeria'', ``no, eravamo da un'altra parte''). \textbf{Tenersi le cose nel
cassetto per tirarle fuori quando servono è sempre sconsigliato}: non servono
ed è impossibile usarle in modo costruttivo.

% ---------------------------------------------------------------------
\section{Le cinque fasi di vita di un team}
\label{sec:z1-fasi-team}

\subsection{Perché conoscerle}

Quasi tutti gli studenti lavorano già in gruppo (al primo anno su un gioco da
tavolo, negli anni successivi su un videogioco) e si saranno accorti che un
team attraversa delle fasi: il modo di lavorare insieme cambia nel tempo.
Spesso ci si fa prendere di sorpresa: ``eravamo così affiatati, cosa è
successo?''.

Un esempio da una consulenza fuori dai videogiochi: un'azienda farmaceutica
gli affida un team commerciale (rappresentanti che presentano medicinali a
ospedali, medici e privati). Il problema: ``prima andava tutto bene; ora che
siamo in diciotto dovremmo lavorare meno e meglio, e invece va peggio''. Va
peggio perché le persone nuove sono appena arrivate, mentre il management si
aspettava già di vedere i risultati dell'investimento.

Fino a pochi anni fa le fasi di un team venivano disegnate come una
\textbf{linea dritta verso l'alto}: le fasi erano le stesse, ma viste come un
miglioramento costante. Non si vedevano le parti \emph{fisiologiche} del
lavoro in team, cioè anche i peggioramenti. Accettarle è il primo passo per
accettare i fallimenti: nessun percorso personale o di team è lineare.

Il relatore ammette che il suo modo di lavorare è nato per il 70\% dal
proprio impegno a imparare, fare bene e rispettare le persone, e per il 30\%
dall'aver lavorato con ``incapaci colossali'' che gli hanno insegnato come
\emph{non} fare le cose. A Blizzard un manager gli fece capire che, pur non
sapendo ancora che tipo di \emph{people manager} sarebbe diventato, sapeva
benissimo quale tipo di manager \emph{non} voleva essere.

\subsection{Forming, storming, norming, performing, adjourning}

Le fasi sono cinque, un po' come le fasi del lutto. Molte aziende ne
considerano solo quattro (``e sono quelle che solitamente chiudono''); le
aziende al passo con i tempi ne considerano cinque\footnote{Il modello è
quello di Bruce Tuckman (1965), a cui la quinta fase, \emph{adjourning}, fu
aggiunta nel 1977. Il nome di Tuckman non è citato in lezione; la
trascrizione riporta la quinta fase come ``journing''.}.

\begin{description}
  \item[Forming] Si viene messi insieme e si comincia a lavorare. Non si sa
        nulla degli altri; spesso non si conoscono nemmeno i nomi, ed è
        normale.
  \item[Storming] Si cominciano a ``dare le testate come le capre in
        montagna per determinare il territorio'': questa cosa la faccio io,
        questa la vorrei fare così, chi decide? Si stabilisce chi fa cosa e
        come si divide il lavoro. Nel grafico è una \textbf{fase di
        peggioramento}: che il team litighi o discuta civilmente non cambia
        nulla, il team deve investire parte di tempo ed energie per risolvere
        queste questioni. È normale e \textbf{non è evitabile}.
  \item[Norming] Si comincia a capire come lavorano gli altri, cosa sanno
        fare e come far funzionare il gruppo; iniziano i miglioramenti. È una
        forma di \emph{accettazione} degli altri, che deriva dall'aver capito
        cosa devo fare io e cosa devono fare loro.
  \item[Performing] C'è uno sforzo \textbf{proattivo}: ``questa cosa
        toccherebbe a me, ma so che Alessio la fa molto meglio, quindi gli
        chiedo aiuto''. Si sfruttano i punti di forza degli altri, perché li
        si conosce e perché si sa che passare il testimone non avrà effetti
        negativi (il capo non dirà ``perché non l'hai fatta tu?'').
  \item[Adjourning] Il progetto finisce, si chiude o viene riorganizzato, il
        team viene sciolto. Quello che andrebbe fatto è analizzare
        \textbf{cosa ha funzionato, cosa no e cosa migliorare} per il
        progetto successivo.
\end{description}

\subsection{La fase che tutti saltano}

``Spoiler incredibile'': l'\emph{adjourning} è la fase che molte aziende
saltano, \textbf{perché c'è sempre il progetto dopo}.

A PG Esports, per la Milano Games Week (come per qualsiasi evento dal vivo a
cadenza annuale: Lucca Comics, Etna Comics, Comicon\dots), finita l'edizione
di un anno si cominciava dal giorno dopo a lavorare a quella successiva,
perché un anno sembra tanto ma non basta mai. Nei videogiochi è lo stesso:
finito un gioco, a parte le patch di mantenimento e del \emph{day one}, si
passa subito al successivo.

Il cambiamento introdotto dal relatore (in parte già praticato, ma da lui
strutturato) è stato: finita la Games Week, \textbf{tutti in ferie per una
settimana}; poi ci si ritrova per definire cosa ha funzionato e cosa no, con
un \textbf{post-mortem} (o \emph{debrief}); solo dopo si inizia a lavorare
all'edizione successiva. Senza questa fase manca la parte di
\emph{learning}, cioè di ciò che si è imparato dal progetto.

Il consiglio vale per gli studenti: quando si finisce il gioco da tavolo, e
poi il gioco digitale, bisogna trovare il tempo per il post-mortem,
altrimenti si butta via \textbf{il 90\% dell'utilità del progetto}. Una
struttura semplice per farlo è descritta nella sezione~\ref{sec:z1-qa-finale}.

% ---------------------------------------------------------------------
\section{Domande sulle fasi del team}
\label{sec:z1-qa-fasi}

\subsection{Lo storming che ritorna}

\paragraph{Capita che lo storming si ritagli spazio anche nelle altre fasi?}
Lo \emph{storming} è una fase precisa, che però molti manager, coordinatori e
team cercano di \textbf{saltare}: ``se io sono bravo come manager e loro sono
bravi come team, passeremo da forming a norming senza problemi''. È un errore,
perché senza una fase in cui si discute davvero non si sviluppa una
\emph{skill} fondamentale: \textbf{saper discutere}. Tutti conoscono persone
non abituate a farlo, per cui ogni minimo disaccordo diventa un affronto
personale (``avevamo deciso la pizza, perché andiamo al sushi? Perché mi fate
questo affronto?''): sono persone che non sono abituate ad avere conflitti
costruttivi.

Se si pensa di essere in \emph{norming} e continuano a emergere conflitti, il
punto è che \textbf{non si era davvero in norming}: lo storming non ``si
trascina'' nelle fasi successive, semplicemente non se n'è mai usciti.
Allo stesso modo, a chi chiede se ogni mini-fase di storming durante la vita
del gruppo sia frutto di uno storming fatto male o solo di carattere, la
risposta è già contenuta qui.

Un partecipante osserva che spesso, ancor prima di finire un progetto, si è
già messi su quello successivo: è il modo migliore per essere in costante
\emph{damage control}. Era l'esperienza del relatore a Blizzard: veniva
spostato fisicamente da un team all'altro, e finito il lancio di
\emph{Diablo~III} era già su \emph{World of Warcraft}, senza vivere il
``dopo''. Si è dovuto ritagliare da solo lo spazio per un debrief personale,
ed è lì che ha deciso che tipo di manager non voleva essere. Il consiglio:
\textbf{se non lo fa il team, fatelo voi}.

\subsection{Dove si passa la maggior parte del tempo}

\emph{In quale fase si passa più tempo nella propria carriera?} Gli studenti
rispondono \emph{storming} e \emph{norming}. Il relatore, in quindici anni,
calcola di aver passato in \emph{performing} circa \textbf{dieci mesi}. Non è
pessimismo: arrivare a performing deve essere sempre l'obiettivo, sia da
manager sia da membro del team; ma è difficile arrivarci, ed è normale che non
sempre ci si arrivi.

Lo storming può durare a lungo, ma con team molto esperti dura pochissimo,
perché c'è chi lo sa gestire. In realtà \textbf{la maggior parte del tempo si
passa in forming}, perché ogni equilibrio raggiunto vale per \emph{quelle}
persone, \emph{quel} progetto, \emph{quelle} tempistiche. Basta che cambi un
qualsiasi fattore (una deadline spostata, una persona nuova, qualcuno che si
dimette, un progetto cancellato) e si ricomincia, perché i bilanciamenti
precedenti non valgono più.

Esempio: un team di cinque persone affiatatissimo, in pieno performing. La
dirigenza, entusiasta, ``raddoppia le risorse'' e lo porta a dieci. O si
lascia tutto com'è (e i cinque nuovi non fanno niente), o bisogna ridiscutere
chi fa cosa e come i nuovi possono aiutare: che piaccia o no, si ricomincia
da forming.

Quindi la bravura di un manager, ma anche di un membro del team,
\textbf{non sta nel passare tanto tempo in performing, ma nell'accelerare il
passaggio da una fase all'altra}, e nel fare bene la fase di
\emph{adjourning}.

\subsection{Analizzare anche i successi}

C'è la tendenza a fare post-mortem solo dei progetti andati male: ``questa
cosa non ha funzionato, vediamo perché''. Va benissimo. Ma quando un progetto
va bene si dà per scontato che sia andato bene ``perché era buono''. Bisogna
invece chiedersi \emph{perché} è andato bene: era il prodotto? La
tempistica? Il modo in cui abbiamo fatto una cosa piuttosto che un'altra?
\textbf{Se non si capisce davvero il motivo del proprio successo, si
tenderanno a replicare le cose sbagliate.}

\paragraph{Blockbuster.} Quando arrivò la notizia che Netflix avrebbe
lanciato un servizio di streaming, Blockbuster (la catena di noleggio di
videocassette e DVD) aveva tutto il tempo per preparare un'offerta
concorrente. In una riunione del management raccontata in un libro (di cui il
relatore non ricorda il titolo), la direzione sembrava orientata a investire
nello streaming, finché un dirigente fece notare che il 12\% del fatturato
derivava dalle \emph{multe per riconsegna tardiva}, impossibili con lo
streaming, mentre secondo le stime lo streaming avrebbe prodotto il 10\% del
fatturato. ``Non possiamo rinunciare a questo 2\%'': niente streaming. Oggi
non c'è più un Blockbuster ``manco a cercarlo col lanternino''. Non avevano
capito qual era il vero motivo del loro successo, cioè il loro modello di
business\footnote{Le percentuali sono quelle riferite in lezione e non sono
state verificate.}.

\paragraph{The Founder.} Nel film sulla nascita di McDonald's c'è una scena in
cui il futuro socio del protagonista gli dice: ``tu pensi di essere nella
\emph{food industry}, e invece sei nella \emph{real estate industry}''. Non
dovrebbe investire sugli hamburger, ma sui terreni su cui costruisce i
ristoranti. Anche per un team vale lo stesso: bisogna capire che cosa
l'ha reso davvero di successo, oppure no.

\subsection{L'autoanalisi personale}

\paragraph{E quando si cambia azienda?} Uno studente chiede se il debrief sia
possibile anche quando si lascia un'azienda. Non si può tornare in ufficio a
discutere della propria esperienza, ma lo si può fare \textbf{per se stessi}.
Il relatore l'ha fatto, anche se all'epoca in modo poco consapevole, ed è
così che è passato dal \emph{community management} al \emph{people
management}. Analizzando cosa aveva fatto bene a Blizzard si era reso conto
che:
\begin{itemize}
  \item aveva costruito un ottimo rapporto di fiducia con i giocatori (con
        alcuni proprietari di fansite parla ancora oggi);
  \item pur senza la seniority, era diventato un punto di riferimento per il
        team, visto come imparziale ed equilibrato anche su temi che non lo
        riguardavano direttamente;
  \item sapeva \textbf{risolvere problemi}: durante il lancio di
        \emph{Diablo~III} era il più tranquillo di tutti, non perché
        sminuisse il problema, ma perché sapeva che si sarebbe risolto con i
        suoi tempi; il lavoro non era spegnere l'incendio ma ``far sfollare
        la gente'' e aspettare che tutto tornasse in ordine.
\end{itemize}
E aveva capito cosa sapeva fare meno bene, ad esempio \textbf{gestire il
proprio manager}. Avere un manager non significa che il manager gestisca te:
il manager gestisce il tuo lavoro, ma tu devi saper gestire i tuoi
\emph{stakeholder}. Spesso lo stakeholder è il proprio manager; da manager,
gli stakeholder saranno CEO, CFO e così via. Gestirli significa conoscerli,
sapere come portare loro le informazioni, come far presente che qualcosa non
va. Fa parte del lavoro, anche se a lui non piace, e ci ha lavorato.

Nel mese e mezzo di pausa prima di entrare in Riot si è concentrato su
questo: ``cosa porto come mio punto di forza?''. L'autoanalisi è culminata
nel chiedere espressamente alla persona che l'aveva assunto: ``prima mi avete
detto di no, ora mi avete preso: perché?''.

\paragraph{Come si capisce se un progetto è andato bene?} Dipende
dall'\textbf{obiettivo}. Se l'obiettivo era vendere milioni di copie e il
gioco le ha vendute, bisogna capire perché. Se, più ragionevolmente,
l'obiettivo era un gioco da tavolo che permettesse di mettere in pratica le
conoscenze del corso, va valutato per quello: ha permesso di farlo? Se sì,
perché? Se no, perché no? L'obiettivo va fissato \textbf{all'inizio del
progetto, non alla fine}, proprio per poterlo valutare dopo.

\subsection{Quando il team cambia di continuo}

\paragraph{Tutti questi cambiamenti non sono dannosi per il progetto?}
Assolutamente sì, ma quasi mai sono volontari. I progetti andati male sono
spesso progetti grandi.

Il relatore racconta di \emph{Cyberpunk 2077}. Conosceva persone in CD
Projekt Red (``il mondo è piccolissimo, le cose si sanno''). Alla Gamescom va
con Diego Ricchiuti, direttore dell'accademia, a vedere la presentazione del
primo gameplay, quella ``in cui regalavano il giubbottino da samurai''.
Finito il video, tra gli applausi, si guardano e capiscono entrambi che
\textbf{il gioco non è pronto}, a soli quattro mesi dall'uscita prevista:
Diego aveva notato il combattimento e la guida. Gli amici di CD Projekt Red
(uno aveva lavorato con lui a Blizzard per sei anni) gli confermano tra le
righe che ``ce la faremo, ma ci sono ancora cose da sistemare''. Il gioco fu
poi rimandato più volte\footnote{In lezione il relatore dice ``due, forse
tre volte''; l'uscita di \emph{Cyberpunk 2077} è stata effettivamente
rinviata tre volte prima del dicembre 2020.}.

Per deformazione professionale, chi ha sviluppato l'occhio critico non
riesce più a godersi un trailer: in un trailer di \emph{The Day Before} Diego
gli aveva fatto notare che l'effetto della fiammata di un fucile era stato
chiaramente aggiunto in postproduzione e non faceva parte del gameplay. ``È
il triste epilogo di questa storia'': non riescono più a godersi un trailer.

Il punto è che più il team è grande, più queste variabili vengono imposte, e
più è facile che il fallimento \textbf{non dipenda dal fatto che il prodotto
sia brutto} (\emph{Cyberpunk} è un bellissimo gioco sotto molti punti di
vista, e il relatore lo sta rigiocando ora che è stato corretto ed espanso).
Nei grandi gruppi qualcuno arriva, qualcuno si dimette, cambia il budget,
arriva il rinvio; nel caso di \emph{Anthem}, il team ha fatto così tante volte
forming--storming--forming--storming che ne è uscito ``un prodotto che
dovevamo espellere: l'abbiamo espulso, buttato fuori dal nostro organismo pur
di liberarcene''.

Tra i cambiamenti c'è anche la crescita verticale: se il capo viene promosso
perché lavora benissimo, bisogna scegliere un nuovo capo, ed è sempre un
caos (chi ci resta male, chi pensava di meritarlo, chi voleva il suo amico,
chi non sopporta il prescelto). Per questo il relatore è \textbf{contrario
alle promozioni durante i progetti}: si rischia di rovinare tutto il lavoro di
storming.

\subsection{Obiettivi chiari nei piccoli gruppi}

Nei gruppi piccoli, come quelli degli studenti, il problema deriva più
spesso dal \textbf{non avere chiaro che cosa si vuole ottenere}: un prodotto
puramente educativo e formativo, o un prodotto richiesto dall'accademia che,
se viene bene, si potrebbe anche vendere? Un gioco commerciale o un gioco
artistico? Va chiarito subito, perché ci saranno sempre momenti in cui
scegliere tra A e B, e a guidare la scelta sarà la domanda: \emph{quale delle
due opzioni ci porta più vicino all'obiettivo?} Se l'obiettivo non è chiaro,
non si sa cosa scegliere.

Questo porta molte aziende a fare \textbf{scelte di design che in realtà sono
scelte commerciali}. Le microtransazioni sono scelte commerciali: nessuno ha
mai pensato ``sai cosa piacerebbe tanto ai giocatori? Pagare le bustine''.
Prima le bustine si guadagnavano giocando; ora la tendenza è ``potresti
giocare, ma se paghi è meglio''. Si pensi alle polemiche su Ubisoft per
l'opzione di ``saltare il gioco'' pagando, o ai famosi \emph{boost} di livello
di \emph{World of Warcraft}: ``devo pagarti per saltare il gioco?''.

Una studentessa (Silvia) cita \emph{Final Fantasy XIV}, che a un certo punto fu
``resettato'' perché non funzionava: un esempio positivo. Il relatore aggiunge
quello negativo di \emph{Age of Conan}, che non è mai riuscito a fare quel
reset iniziale (``è una schifezza, ve lo ripresentiamo quando è pronto''). Ci
vuole consapevolezza, ma anche i mezzi: un conto è Square Enix, che può
permetterselo, un altro è un'azienda che non può chiudere tutto per un anno
continuando a pagare gli stipendi.

Riassumendo: nei piccoli gruppi i problemi sono più di
\textbf{comunicazione e chiarezza} (obiettivi non chiari o che si spostano);
nelle grandi aziende di più di \textbf{persone che vanno e vengono, obiettivi
e tempistiche che cambiano}. In accademia le scadenze sono fisse (il gioco da
tavolo si consegna a fine anno scolastico, a giugno, e giugno resta); in
azienda può succedere che una feature prevista per la \emph{vertical slice}
di luglio venga anticipata a maggio per portarla alla Gamescom, e un sistema
perfetto per consegnare a giugno può non esserlo più per consegnare a maggio.

\subsection{L'ambiente determina la qualità del prodotto}

\paragraph{È vero che chi passa per Electronic Arts finisce per fare brutti
giochi?} Uno studente cita BioWare, che faceva ottimi giochi finché era
autonoma. Il problema non è dei developer, che non sono ``improvvisamente
diventati incapaci'': è evidente che il problema è altrove, e che
\textbf{l'ambiente in cui si fa un gioco determina molto spesso la qualità del
prodotto}.

Un esempio locale positivo: Tommaso Bonanni, conosciuto dal relatore insieme a
Diego e già docente dell'accademia, ha fondato a Roma Caracal Games Studio,
applicando fin dall'inizio questi concetti di cultura aziendale. L'azienda è
rimasta piccola, ma ha molto chiaro cosa vuole fare e come. Quando si
predispone fin dall'inizio l'azienda e il team a lavorare in un certo modo è
fattibilissimo; è molto più difficile quando si è EA, si acquisisce
un'azienda e le si dice ``da oggi si fa così''.

Il relatore ha vissuto la transizione da Blizzard ad Activision Blizzard (uno
dei motivi per cui se n'è andato, volontariamente). Ha visto la cultura
interna cambiare per gradi:
\begin{enumerate}
  \item ``facciamo giochi fighissimi, perché i giochi devono essere
        fighissimi'';
  \item ``facciamo giochi fighi così vendono di più'';
  \item ``vendiamo di più facendo giochi fighi'' (sembra uguale, ma è molto
        diverso);
  \item ``facciamo più soldi''.
\end{enumerate}
Lo si capiva dalle scelte di design, osservando quale fosse il peso più
grande sulla bilancia. Lo stesso varrà per gli studenti, nei loro giochi e nei
colloqui di lavoro: capire subito che cosa vuole l'altra parte.

\paragraph{Leader e visione.} Nei gruppi dell'accademia, dal secondo anno, c'è
una persona che assume il ruolo di \emph{lead} o \emph{scrum master}. È
diverso dall'essere un \emph{game director} (``il gioco si fa come dico io,
aiutatemi a farlo con la mia visione''), come avviene nei grandi studi: da lì
nascono i nomi nei titoli, come \emph{Sid Meier's Civilization}, che indicano
la persona che ha fatto quel gioco, oppure un puro branding, come \emph{Tom
Clancy's Rainbow Six}.

\paragraph{Inclusion come jolly.} A margine, uno studente osserva che in molte
aziende ``la risposta giusta è sempre \emph{inclusion}'', usata come un
jolly. È comunissimo, conferma il relatore: \emph{diversity and inclusion}
viene spesso usata per dire ``così includiamo più clienti'', ma non è quello
il suo punto (``altro discorso'').

% ---------------------------------------------------------------------
\section{Confrontarsi con chi è più bravo}
\label{sec:z1-rivalita}

Arrivato a Riot dopo sei anni a Blizzard, il relatore pensava di avere un buon
bagaglio di esperienza. Si è trovato invece ``come un fante buttato in mezzo
ai Navy SEAL'': persone che avevano superato 11, 12, 13 colloqui, con un
livello molto alto.

Individua subito un collega francese che, nella sua testa, diventa ``la
versione migliore di me'': stesso ruolo (community manager), esperienza
simile (arrivava da Ubisoft, come quasi tutti i francesi di Riot), ma molto
più spigliato, sempre con la battuta pronta, ``l'anima della festa''. Il
relatore invece è introverso: un workshop come questo per lui è
\emph{emotivamente drenante}, non ama parlare in pubblico pur amando
condividere le proprie esperienze, e non riesce a sostenere conversazioni
come ``cosa hai fatto questo weekend?''. Lavorando raramente insieme, aveva
sviluppato una rivalità ``da allenatore di Pokémon'' (lui era Ash, il
relatore l'antagonista).

Dopo sei mesi, a un evento di \emph{team building} aziendale (a Tenerife), il
suo capo, molto attento a queste dinamiche, li mette apposta nello stesso
gruppo. A fine serata, complice l'open bar, il relatore gli confessa di averlo
mal sopportato dal primo giorno per invidia della sua estroversione. Il
collega ride e risponde che pochi giorni prima aveva detto al capo la stessa
cosa al contrario: invidiava a Stefano il suo essere sempre organizzato, sul
pezzo, presentarsi nella sala riunioni giusta (mentre lui finiva in quella
sbagliata), arrivare con gli appunti, con i dati a supporto e le analisi di
mercato, mentre lui spesso non aveva nemmeno letto la mail del meeting.
Un'invidia in due direzioni, che li ha portati a lavorare molto insieme e a
migliorarsi a vicenda.

Le lezioni:
\begin{itemize}
  \item non bisogna demoralizzarsi vedendo qualcuno che fa lo stesso lavoro
        ed è più avanti sotto alcuni aspetti: spesso c'è un bilanciamento. I
        ``punti caratteristica'' sono limitati: se li ha messi in carisma,
        non li avrà messi in saggezza;
  \item conviene \textbf{chiedere a queste persone come mai} sono così brave
        in qualcosa, come si fa con un progetto andato bene: magari non lo si
        farà mai bene quanto loro, ma si comincerà ad avere il
        \emph{mindset} giusto; e viceversa, loro saranno spinti a chiedere a
        noi;
  \item spesso \textbf{le proprie doti le scopriamo dagli altri}. Il relatore
        non si era mai accorto di essere ``quello che resta calmo durante le
        emergenze'' (il lancio di \emph{Diablo~III}, quello di
        \emph{Draenor}, i problemi dal vivo alla Games Week) finché non
        gliel'hanno fatto notare tutti: per lui era la normalità. A chi pensa
        ``non sono capace, forse non è il mio percorso'', può darsi che
        semplicemente nessuno gli abbia ancora detto ``questa cosa la sai
        fare davvero bene, concentrati su questo''.
\end{itemize}

% ---------------------------------------------------------------------
\section{Domande sul lavoro nell'industria}
\label{sec:z1-industria}

Diverse domande degli studenti hanno riguardato le condizioni di lavoro nei
grandi studi internazionali. Le risposte sono raccolte qui per argomento.

\subsection{Presenza, remoto e trasferimenti}

\paragraph{Grandi studi: quasi solo in presenza.} Sia a Blizzard sia a Riot si
lavorava \textbf{solo in presenza}, senza smart working: sei anni a Parigi e
cinque a Dublino. Si parlava solo inglese: anche nei meeting tra soli
italiani si tendeva a comunicare e a scrivere in inglese, per correttezza
verso chi si fosse aggiunto e per lasciare una traccia comprensibile a tutti.
Il francese serviva per fare la spesa e sopravvivere, non in ufficio.
L'ambiente era estremamente eterogeneo (francesi, polacchi, spagnoli e
molti altri): durante il colloquio dal vivo a Dublino, mentre si preoccupava
di essersi vestito in modo adeguato, vide passare dietro la parete a vetro un
collega svedese con i dreadlock, a piedi scalzi, con una pianta sotto un
braccio e un aquilone sotto l'altro, e capì di essere ``a posto così''.
Un'esposizione a molte culture e modi di pensare che il relatore considera
bellissima.

Anche oggi, tranne alcuni studi particolari, la maggior parte del lavoro si
svolge di persona: amici a Paradox lavorano a Stoccolma, i pochi rimasti a
Blizzard si sono dovuti trasferire a Irvine. Per posizioni \emph{junior} o
\emph{entry level} i grandi studi raramente concedono il lavoro da casa; per
profili senior può esserci qualche eccezione (ad esempio nove mesi da remoto
e tre in sede).

\paragraph{Relocation.} Spesso tocca spostarsi, ma il relatore non ha mai
pagato un trasloco di tasca propria: queste aziende offrono la
\emph{relocation} (aiuto a trovare casa, ad aprire il conto in banca e
supporto pratico), perché è un investimento sulla persona.

\paragraph{Colloqui mentre si lavora altrove.} Per i colloqui in presenza
prendeva ferie, senza dire dove andava. Il consiglio generale: per fare
colloqui \textbf{prendete ferie o permessi}, non datevi malati, ``che non è un
bel modo di comportarsi''.

\paragraph{Remoto ``doppio''.} Un collega lavora per l'azienda coreana che
gestisce \emph{PUBG}, con ufficio ad Amsterdam, ma lavora da Milano: un
``doppio remoto''. Si trova, ma in genere un ambiente anglofono con sedi in
Europa chiederà di trasferirsi in Europa.

\paragraph{Famiglia e trasferimenti.} Il relatore ha avuto la fortuna di
sposare una persona che fa lo stesso lavoro (conosciuta lavorando su
\emph{Warhammer Online}, poi con lui a Blizzard e a Riot, nella
localizzazione e nel doppiaggio, ad esempio di \emph{World of Warcraft} in
italiano): avendo competenze rare, spesso venivano assunti entrambi, al
massimo a pochi mesi di distanza. Per i developer la soluzione frequente è
che la famiglia segua la persona assunta, perché lo stipendio basta per
tutti. Ma non sempre c'è equilibrio, e a volte uno dei due deve sacrificarsi:
sua moglie ha lasciato Blizzard per seguirlo a Dublino (per poi essere
assunta da Riot); lui ha lasciato Riot per tornare in Italia quando lei ha
trovato lavoro qui. Ha dovuto rinunciare più volte a lavorare in Islanda per
CCP (\emph{EVE Online}), perché non era una scelta solo sua. Si fanno
compromessi, ma si riesce a far quadrare.

\paragraph{Ci si sposta per ogni progetto?} No: si sono spostati
esclusivamente per cambiare azienda. Cambiando progetto nella stessa azienda
è raro, quasi mai (``non l'ho mai visto succedere''), che venga chiesto di
trasferirsi. In Francia stavano benissimo (vivevano fuori Parigi; ``i
problemi con i francesi sono in realtà problemi con i parigini''), parlavano
francese ed erano vicini alla doppia cittadinanza: se ne sono andati perché
non vedevano un futuro a Blizzard, per il modo in cui l'azienda trattava i
colleghi.

\paragraph{Il mercato oggi.} Non esiste un \emph{silver bullet} per trovare
lavoro domani dove si vuole. Con tutti i \emph{layoff} recenti ci sono
migliaia di esperti del settore che cercano lavoro: serve pazienza, ma è
fattibile.

\paragraph{Studi indie e remoto.} Negli studi più piccoli il remoto è più
diffuso: a Stormind circa 60 persone su 95 lavorano da remoto, ma lo studio
si è strutturato fin dall'inizio per farlo. Serve consapevolezza. Bending
Spoons (che sviluppa app per mobile, non giochi, ed è l'azienda dell'app
Immuni), pur essendo grandissima e con sede a Milano, lavora da remoto e
lascia l'ufficio a chi vuole: lo smart working è posizionato quasi come un
\emph{benefit}, ma permette anche di avere un ufficio da 300 persone invece
che da 700. Le sedi fisiche costano decine di milioni all'anno (Blizzard ha
chiuso le sedi di Amsterdam e di Parigi). Tuttavia, per lavori molto
verticali e specifici gli studi preferiscono avere le persone dal vivo; i
ruoli più spesso remoti sono quelli che non fanno sviluppo (customer service,
PR\dots).

Anche gli artisti lavorano spesso in presenza: l'accademia può dare solo una
\emph{simulazione} di quanto stretto sarà il contatto tra programmazione,
design e art. ``Sarà ancora di più.''

\paragraph{La comunicazione da remoto.} È più difficile. Pur essendo
estremamente introverso, il relatore \textbf{obbliga tutti a usare la webcam}
nei meeting in cui conta: non in un aggiornamento con novanta persone
collegate, ma certamente in un colloquio a due sulla performance dell'ultimo
anno. La componente fisica è una parte enorme della
comunicazione\footnote{In lezione si parla di ``quasi l'ottanta per cento'';
il dato va preso come indicativo, poiché le percentuali di questo tipo
circolano in forme diverse e sono spesso citate fuori contesto.}: anche se si
è timidi, accendere la webcam è l'approssimazione più vicina all'essere dal
vivo. Negli ambienti dove si lavorava senza webcam la situazione era
``esponenzialmente peggiore''.

\subsection{Differenze culturali nel team}

\paragraph{Gestire sensibilità diverse.} A Riot i team erano divisi per
\emph{lingua}, non per nazione: il team di lingua spagnola includeva anche un
messicano, quello di lingua russa includeva russi e ucraini, proprio nel
periodo della guerra in Crimea. Tra i colleghi non c'era alcun problema (erano
amici), ma ci si è posti la domanda: \emph{come comunichiamo con una community
che sarà divisa su questo tema?} Nei giochi online la divisione era evidente:
in \emph{Rocket League} c'era chi, con un compagno di squadra con la bandiera
avversaria, si faceva otto autogol pur di farlo perdere. Allo stesso modo,
con molti americani mandati a Dublino a coprire posizioni scoperte, c'erano le
elezioni (Clinton--Trump, Biden--Trump) e varie guerre: serviva sempre una
certa sensibilità.

Imparare a navigare queste situazioni lo ha reso un professionista migliore:
gli ha insegnato a fermarsi e chiedersi ``io la percepisco così; gli altri
dieci, con background, storie e lingue diverse, come la penseranno?'', a
ragionare più \emph{long term} e più \emph{big picture}, con più empatia.

\paragraph{Stili di comunicazione.} Un partecipante che ha lavorato nel Regno
Unito chiede dello stile comunicativo italiano, a volte frainteso. Uno dei più
grandi complimenti ricevuti dal relatore è stato ``non sembri italiano''
(``gesticoli poco per essere italiano'', ``parli bene inglese per essere
italiano''). Ha imparato ad \textbf{adattare la comunicazione alla persona}:
con spagnoli e francesi la battuta sulla nazione funziona, perché c'è un
``odio--amore'' condiviso; con russi e giapponesi no. Doveva adattare anche
l'umorismo. Un collega metà russo e metà messicano diceva di aver preso ``il
lato peggiore di entrambe le nazionalità'', ed è stato un ottimo maestro per
imparare a leggere culture diverse.

Non bisogna generalizzare, ma gli stereotipi possono essere un punto di
partenza ragionevole, \emph{sapendo che troverete sempre le eccezioni}. E
vanno letti correttamente: gli americani tendono a essere più estroversi e ad
avere sempre qualcosa da dire non perché ``sono fatti così'', ma perché
vengono educati così (all'università tra i corsi c'è il \emph{debate}, saper
argomentare in pubblico); un collega giapponese può essere più introverso per
ragioni culturali. Nel dubbio si parte da lì, poi si conosce la persona.

\subsection{Negoziare, discriminazione e stipendi}

\paragraph{La famiglia come leva.} Avere un partner riduce le opzioni, ma le
aziende lo sanno: uno spostamento è per due. Si può usare come leva per
negoziare (un compenso diverso, un supporto per il trasloco): se un'azienda ha
deciso di volere quella persona, 2, 3, 4 o 10.000 euro in più all'anno, per
un'azienda che fattura centinaia di milioni o miliardi, non cambiano nulla.

\paragraph{Essere sposati è un motivo per non essere assunti?} No. In alcuni
paesi è proprio illegale discriminare su questo: in Francia era illegale
chiedere CV con data di nascita, foto, sesso o religione, e chiedere se si
vogliono figli. Purtroppo, dai racconti dei suoi studenti di master (in
settori come \emph{beauty and wellness} e \emph{food and wine}), in piccole
realtà italiane queste domande (``hai 25 anni, hai intenzione di
sposarti?'') vengono ancora fatte, e servono a filtrare i candidati. Nelle
grandi aziende ci sono sistemi per evitarlo; anzi, molte danno \emph{più}
benefit a chi ha famiglia (``porta pure la famiglia, è tutto pagato''),
perché le persone competenti hanno spesso una certa età e quindi,
ragionevolmente, una famiglia.

\paragraph{Gli stipendi.} Dipendono molto da dove si vive. Lo stipendio a Riot
non era paragonabile a quello di Blizzard perché in Francia era tassato al
52\% e in Irlanda al 17\%\footnote{Aliquote come riportate in lezione,
riferite alla situazione personale del relatore.}: nessuno stipendio
competitivo può compensare il 35\% di differenza di tasse. Negli Stati Uniti,
con pochissima previdenza pubblica, gli stipendi sono molto più alti perché
poi ci si paga tutto (``se hai un raffreddore paghi tu''). Bisogna valutare il
potere d'acquisto di quei soldi nel paese in cui si andrà a vivere.

\subsection{Freelance e \emph{work for hire}}

All'inizio della sessione pomeridiana uno studente che lavora nell'audio
chiede quali siano le condizioni per chi lavora come \emph{freelance}.
\begin{itemize}
  \item Dipende molto dal progetto: spesso il contratto è a progetto, per
        portare a termine una parte o un intero progetto.
  \item Non tutte le aziende usano freelance: Riot e Blizzard lo facevano
        pochissimo e preferivano assumere (magari trattando poi la persona
        ``come un freelance'', lasciandola a casa a fine progetto).
  \item Spesso i lavori offerti come freelance sono in realtà
        \textbf{sostitutivi di un full time} (``freelance, ma sei ore al
        giorno'').
  \item Si può fare, ma bisogna essere molto chiari nel definire l'impegno
        (quante ore, come vengono tracciate): non tutti specificano cosa
        significa davvero ``freelance''. Lavorando così per due o tre aziende
        si può arrivare a un tempo pieno.
  \item È più comune il \emph{work for hire}: aziende che forniscono personale
        ad altre aziende; si è assunti da X ma si lavora ``embeddati'' presso
        Y (una sorta di subappalto). Anche lo studio del relatore lo fa.
  \item Il puro freelance di quattro mesi esiste, ma di solito per quei
        quattro mesi diventa un lavoro a tempo pieno; è raro ``quattro mesi,
        due ore al giorno''. Nelle grandi aziende è improbabile lavorare nello
        stesso mese per Ubisoft, Riot e Blizzard tre ore ciascuna; è più
        realistico lavorare per più startup indie italiane, che non possono
        permettersi un'assunzione ma possono pagare due ore al giorno.
  \item Molto passa dai contatti: trovare le aziende giuste e rendere esplicite
        le condizioni.
\end{itemize}

% ---------------------------------------------------------------------
\section{Il circolo virtuoso: feedback, miglioramento, fiducia}
\label{sec:z1-triangolo}

\subsection{Il triangolo}

L'obiettivo a cui il workshop vuole arrivare è un \textbf{triangolo che si
autoalimenta}:
\begin{enumerate}
  \item si dà un \textbf{feedback} con l'obiettivo di migliorare le
        performance (inclusi elementi caratteriali);
  \item la persona \textbf{migliora} e vede di averne beneficiato;
  \item questo genera \textbf{fiducia}, che rende più facile dare e ricevere
        il feedback successivo.
\end{enumerate}
Togliendo uno qualsiasi dei tre elementi, il circolo non è altrettanto forte.

In qualsiasi fase si parta, il rispetto per i colleghi è inizialmente dettato
dall'\emph{educazione}, non dalla fiducia. La fiducia arriva quando i colleghi
cominciano a dire cose giuste e sensate e a \textbf{mantenere la parola}.
Dare e ricevere feedback è il modo in cui, in un certo senso,
``inganniamo il cervello'' per creare una \emph{case history} di fiducia: si
ha motivo di fidarsi di più, di ascoltare di più, di essere più ricettivi. E
si diventa sempre più bravi ad accettare il feedback e ad agire di
conseguenza.

\subsection{Accettare un feedback non vuol dire fare come dice l'altro}

Dieci anni fa il relatore non era ostile ai feedback, ma molto
\emph{difensivo}: ``sì, sì, va bene, ora ci guardo''; oppure, a chi gli diceva
``avresti dovuto farla così'', rispondeva spiegando perché l'aveva fatta in
quel modo. Oggi ha imparato a dire ``ok, grazie, ora ci ragiono'': non deve
difendere la propria posizione da chi viene solo a dire ``hai pensato a
questa cosa?''. Prima lo viveva come un attacco (``mi viene a dire che non so
fare il mio lavoro''), ora come un ``ottimo punto di vista, lo prendo in
considerazione''.

Con una postilla fondamentale: \textbf{accettare il feedback non significa
fare come dice l'altro}. ``Mi hai detto che preferisci questo mob verde. Ho
pensato anche a farlo verde? Sì. Penso comunque che giallo sia meglio? Sì, lo
faccio giallo.'' Il relatore accetta \emph{sempre} il feedback, valido o no:
garantisce di tenerlo a mente quando lavora. Che poi faccia come gli viene
detto è un altro discorso.

\subsection{Fiducia e scadenze}

La fiducia funziona allo stesso modo con gli impegni. Se si chiede a qualcuno
una cosa per venerdì e la consegna venerdì, la volta dopo gliela si chiederà
per venerdì. Se invece arriva il giovedì della settimana dopo, la volta dopo
gli si darà una \textbf{deadline anticipata}, come con chi arriva sempre in
ritardo agli appuntamenti (``se ti dico alle 7 arrivi alle 7 e mezza; la
prossima volta ti dico alle 6 e mezza'').

Al lavoro questo è dannosissimo, perché crea \textbf{deadline fittizie} che
non sono mai quelle vere, e nascono discussioni del tipo ``perché devo farlo
in due giorni se si fa in cinque?'' (``sì, se tu lo facessi in cinque''). La
conversazione giusta è un'altra: \emph{``che problemi hai avuto nel
consegnare venerdì?'' ``Questo, questo e questo.'' ``Cosa posso fare per
evitare che succeda?''}

Il relatore non vuole mai sentirsi dire ``non sono riuscito a farlo''.
Preferisce di gran lunga sentirsi dire, tre giorni prima, \textbf{``potrei non
riuscire a farlo, per questi motivi''}. Chi ha la visione d'insieme può dare
più budget, più persone, spostare la deadline, cambiare lo scopo del progetto:
può fare qualsiasi cosa per aiutare. Ma se lo si scopre il venerdì, cinque
minuti prima della deadline, non può fare più niente. Il messaggio è:
\textbf{``aiutami a poterti aiutare''}. Non penserà mai male di chi dice ``non
ce la faccio per questi motivi'', perché non sta dicendo ``non ho voglia'': sta
indicando i problemi, e lavorare su quei problemi è compito del manager. Non
può invece fidarsi di chi fino al giorno prima dice ``tranquillo, ce la
faccio'' e il giorno dopo ``non ce l'ho fatta''.

Questo è fondamentale per chi lavora in programmazione, design o art, perché
ci saranno sempre deadline: bisogna sempre \textbf{lavorare sulle aspettative},
dire al proprio manager (o al proprio report) cosa è lecito aspettarsi. Non
bisogna fare come i pinguini di \emph{Madagascar} (``sorridete e salutate, va
tutto bene''): se c'è un problema, va detto, altrimenti si peggiora solo la
propria situazione.

Da italiani si tende a non voler mai alzare la mano per chiedere aiuto
(``no, no, ce la facciamo, lo faccio io''). A volte non ce la si può fare, non
perché non si è capaci, ma perché vari fattori lo impediscono. Accettarlo è
importante: \textbf{è molto meglio avvisare e chiedere una mano che dire
``non ce l'ho fatta''}.

\subsection{Domande sul dare e ricevere feedback}

\paragraph{``Ne parliamo a fare, se poi fai come vuoi?''} Uno studente (Cristian)
riceve feedback che, dopo averli analizzati, non ritiene congrui; quando non
li segue gli viene detto ``ne parliamo a fare?''. C'è un problema di
\textbf{aspettative}: l'idea che accettare un feedback voglia dire fare come
dice l'altro. Accettarlo significa prendere in considerazione ciò che è stato
detto. Chi decide diversamente ha però un \textbf{onere}: argomentare perché
si fa comunque a modo proprio (``mi hai detto verde, la facciamo gialla
perché\dots''). Poi bisogna solo aspettare i risultati per vedere chi aveva
ragione.

La decisione spetta sempre \textbf{a chi ha la responsabilità di quella
decisione}. Se il game director deve scegliere tra A e B, anche se il
feedback dice B, la scelta è sua perché la responsabilità è sua. Non sarebbe
corretto essere ritenuti responsabili di una scelta delegata a chi dà
feedback. La risposta può essere semplicemente: ``l'ho tenuto in
considerazione, ma visto che ne sono responsabile io e le conseguenze sono
mie, ho scelto di fare X''.

\paragraph{E se l'altro fa ripicca?} Un partecipante chiede di chi, vedendo
il proprio feedback non accolto, si vendica ignorando a sua volta feedback
buoni. Per il relatore è un problema caratteriale: chi si comporta così ``è
già sulla strada per uscire''. Non c'è spazio per l'ego in team di questo
tipo. Il fenomeno nasce dall'idea che \emph{il feedback negativo faccia fare
brutta figura}, che il fatto che si potesse fare meglio tolga qualcosa alle
proprie capacità.

Non è così: \textbf{il feedback non è un giudizio}. Si può passare un anno a
Riot ricevendo solo feedback negativi e avere una valutazione annuale
stellare, perché ricevere molti feedback negativi significa essere una persona
approcciabile, che agisce sul feedback e invita gli altri a darlo. Il relatore
ha visto valutazioni molto peggiori per chi riceveva solo feedback positivi,
perché la gente aveva paura di dargliene di negativi: ma solo ``pacche sulle
spalle'' non fanno migliorare. Il problema va affrontato parlando con la
persona; a volte le conversazioni iniziano con \textbf{``ti do un feedback sul
feedback''}, cioè sul modo in cui l'altro ha dato il suo.

Esempi dalla sua esperienza:
\begin{itemize}
  \item lui chiede spesso di \textbf{non ricevere feedback senza esempi
        tangibili}, perché fa fatica ad astrarre: ``sei stato maleducato'' non
        significa nulla, ``ieri al meeting mi hai risposto male su questa
        cosa'' sì;
  \item ad altri serve l'opposto: ha ricevuto il feedback ``quando devi dirmi
        che ho sbagliato, puoi girarci un po' intorno e darmi del contesto?'',
        perché il suo stile diretto (``questa cosa non ha funzionato, dovremmo
        fare così'') feriva. Ha adattato il suo modo di dare feedback: con
        alcuni è diretto, con altri gira intorno. L'importante è saperlo.
\end{itemize}

\subsection{Non portare la soluzione}

Un consiglio centrale: \textbf{quando si dà un feedback, non si dà una
soluzione}. Non si dice ``io farei così'', ma ``secondo me questa cosa si
potrebbe fare diversamente, per questo e questo motivo: come potremmo farla
diversamente?''. È come chi smette di fumare perché vuole e chi smette perché
glielo dicono moglie e figli: se non si è convinti, spesso non si smette.

Quando il relatore diceva ``c'è questo problema, io farei così'', il meeting
successivo iniziava inevitabilmente con ``ci ho provato come hai detto tu, con
me non ha funzionato''. Sempre. Perché la soluzione era sua. La
responsabilità della soluzione deve essere invece \textbf{congiunta}: ogni
feedback deve creare un piccolo team che ha come obiettivo risolvere quel
problema. Questo toglie l'antagonismo (non ``ti dico che non sai fare il tuo
lavoro'', ma ``ti voglio aiutare a farlo meglio, lavoriamoci insieme'') e
sposta il focus: il problema non sei tu, è ciò che sta accadendo.

\paragraph{E se sono i piani alti a tenere all'oscuro dei problemi?} È una
questione di cultura aziendale. Il relatore fa sempre presente al proprio capo
il concetto di \textbf{service level agreement} (che verrà ripreso nella
seconda giornata): ``se mi dai questi input, ti garantisco questo output in
questi tempi''. Se il capo chiede delle cose senza dire budget, tempi e
perché, non può aspettarsi lo stesso risultato. È una comunicazione a due
vie: si dice al manager cosa serve per fare il lavoro, e lui deve fornirlo o
accettare che il risultato non sarà quello voluto.

\subsection{Amici, parenti e parole proibite}

\paragraph{``Non si lavora con amici e parenti''?} Per sua definizione
personale il relatore ha pochissimi amici (``stanno sulle dita di una mano'')
e una soglia molto alta per considerare amico un collega. Allo stesso tempo,
in quindici anni e con centinaia di colleghi, ha avuto solo \textbf{due}
esperienze in cui non è riuscito a lavorare con qualcuno: non serve essere
amici per far funzionare le cose. Con gli amici, anzi, gli viene più facile:
se dice a un amico che una cosa ``fa schifo'', l'amico non pensa ``come si
permette'', ma ``se perfino Stefano, che è mio amico, me lo dice, allora è
proprio così''. Diventa un rafforzativo, grazie alla fiducia costruita prima.
Il problema del lavorare con amici e parenti nasce dal \emph{non voler} dire
all'amico che sta sbagliando; se invece la confidenza permette di dirlo, si
può saltare la ``confezione'' del messaggio e passare direttamente al
contenuto.

\paragraph{Le frasi peggiori.} Alla domanda se esistano parole peggiori di
``ottimo lavoro'', il relatore indica la sua classifica: ``ottimo lavoro''
(detto in modo vuoto)\footnote{Il passaggio della trascrizione è
frammentario: si intende verosimilmente il complimento generico e di
circostanza, che non dice nulla di utile.}; poi la frase che lo ``fa uscire di
testa'', \textbf{``abbiamo sempre fatto così''} (``e infatti avete chiuso'',
come Blockbuster); infine \textbf{``siamo una famiglia''}, riferito al posto di
lavoro. Un'azienda non è mai una famiglia, perché non ci si può aspettare lo
stesso impegno da chi la possiede e da chi ci lavora: non c'è un
corrispettivo proporzionale tra impegno e ritorno. Per quanto inclusiva,
flessibile e razionale, non sarà mai una famiglia.

\subsection{Ego, muri e il peggior modo di scegliere un leader}

\paragraph{Chi mette davanti l'ego e alza muri.} Esistono, ma di solito durano
pochissimo in azienda. Non si tratta di avere un ego forte, ma di \textbf{non
saper ricevere feedback}. Spesso sono persone bravissime nel loro lavoro che
poi si chiedono perché tutti intorno diventino più bravi di loro: perché gli
altri accettano feedback e hanno più fonti da cui attingere, mentre loro si
sono rintanati nella loro torre.

Molti manager tendono ad \textbf{assumere persone un filo peggiori di loro},
per non fare brutta figura. Il relatore assume solo persone più brave di lui,
``perché non è una gara'': ha lavorato con molti community manager più bravi
di lui come community manager, ma lui era più bravo a fare il manager. Non
viene valutato su quanto sa fare il community manager, ma su quanto il suo
team sa farlo. Un partecipante chiede se la paura di avere qualcuno di
migliore esista anche in ambito accademico. Assolutamente sì: al relatore è
capitato di non essere preso come docente in un corso di formazione sui
videogiochi perché il responsabile aveva percepito che ne sapesse più di
lui\footnote{La trascrizione dice ``che non sapevo più di lui'', ma il senso
del passaggio richiede il contrario.}. ``Ma perché te ne importa? Tu sei quello che ha portato chi
sa fare il suo lavoro: ben venga.'' O si ha questa cultura, o si ha paura per
il proprio posto.

\paragraph{Promuovere il migliore.} Quando bisogna scegliere un leader per un
team (programmazione, design, art), molto spesso si sceglie \textbf{il migliore
tecnico}: il miglior programmatore, il miglior designer. È spesso la scelta
sbagliata: è come prendere un idraulico per fare l'elettricista. Fare il
\emph{people manager} non è come fare il designer o il programmatore; si può
essere un ottimo programmatore e un pessimo manager, e viceversa. Le aziende
l'hanno capito e lo fanno meno di prima. A Blizzard c'era un grave problema di
\emph{middle management}: molte persone in ruoli importanti erano state
promosse cercando le skill sbagliate, erano frustrate dal non saper fare il
proprio lavoro e schiacciavano chi stava sotto per paura di fare brutta
figura.

Un partecipante cita una teoria simile: una persona viene promossa finché non
raggiunge \emph{il proprio livello di incompetenza}, dove si blocca\footnote{È
il \emph{principio di Peter} (Laurence J.\ Peter, 1969), citato in lezione
come ``la teoria di Peter''.}. Un altro cita Michael Scott di \emph{The
Office}, su cui il relatore tornerà (sezione~\ref{sec:z1-quando}).

% ---------------------------------------------------------------------
\section{Perché la cultura conta: retention, costi e fiducia}
\label{sec:z1-cultura}

\subsection{L'ambiente di lavoro e la retention}

Perché tutto questo è così importante? Perché l'ambiente di lavoro conta
moltissimo per la \textbf{retention}, cioè per trattenere le persone e farle
lavorare in modo efficiente. Se contasse solo lo stipendio, molti non
lascerebbero mai Riot o Blizzard; e più si va avanti con l'età, più si impara
a dare valore anche al resto.

Inoltre, secondo gli studi citati nelle slide, chi vive il proprio lavoro in
modo negativo è \textbf{meno efficiente nel problem solving}: l'energia
mentale va verso ``come faccio a non farmi dare la colpa'' o ``come faccio a
fare bella figura'', non verso la soluzione del problema. In ambienti negativi,
dove ci si deve preoccupare della propria ``sopravvivenza'', si lavora
peggio.

\subsection{Il costo di chi se ne va}

Ogni volta che qualcuno si dimette o viene licenziato, il costo di trovare
una nuova persona, fare i colloqui, formarla e farle attraversare forming,
storming, norming\dots\ è circa \textbf{un terzo del suo stipendio}. Per questo
le aziende tendono a fare più \emph{retention} che \emph{acquisition}, almeno
con le figure considerate utili alla produzione.

Allo stesso tempo si vedono sempre più spesso \emph{layoff} e \emph{hiring
wave}: si licenziano 3.000 persone e sei mesi dopo se ne riassumono 2.500 o
3.500, ma tutte più junior. Il relatore ne ha scritto in un articolo su
Eurogamer.it, in cui spiega anche alcuni ``trucchi del mestiere'': perché i
licenziamenti avvengono sempre in certi periodi e sempre ``a gruppi'' (prima
Bethesda, poi Blizzard, poi un'altra\dots).

\paragraph{L'offerta di Riot.} A Riot esisteva una pratica (di cui il relatore
non ricorda il nome) per cui, firmato il contratto, si aveva \textbf{un anno
di tempo per dire ``non mi piace, non me la sento, me ne vado''} e ricevere
subito come buonuscita un anno di stipendio. Avevano stimato che fosse molto
meglio pagare subito un anno a qualcuno e lasciarlo andare che tenerlo due,
tre, quattro anni come \emph{underperformer}: la responsabilità della scelta
di assumere era dell'azienda\footnote{La pratica, nota come ``The Offer'', è
descritta in fonti pubbliche con condizioni diverse (un periodo più breve e
una somma pari a una frazione dello stipendio annuo). Si riportano le
condizioni come raccontate in lezione; andrebbero verificate.}. L'hanno usata
pochissimi, ma il solo fatto che esistesse faceva lavorare meglio anche lui da
manager, perché sapeva che chi non si trovava bene poteva andarsene con un
anno di tempo per trovare un altro lavoro.

\paragraph{Il ricircolo.} C'è un forte ricircolo nell'industria: molti ex
colleghi passano da Ubisoft a Paradox, da Paradox a Funcom, da Funcom a
Blizzard. Il relatore ha scelto di uscirne lavorando solo per aziende
\textbf{indipendenti}: Stormind non è un indie di cinque persone, ma è
indipendente, e non deve accontentare azionisti. Quando vede un'azienda con
utili miliardari licenziare migliaia di persone si chiede ``dov'è la linea''
(\emph{when enough is enough}): con azionisti che vogliono un ritorno, se
quello è l'obiettivo stanno facendo la cosa giusta, ma lui non condivide
l'obiettivo\footnote{In lezione si parla di ``6,5 miliardi di utile'' e di
``3.000 persone licenziate''; le cifre sono riportate come dette e non sono
state verificate.}.

\paragraph{L'open space vuoto.} A Riot, durante un'ondata di licenziamenti,
del suo team di trenta persone è rimasto solo lui: ``teniamo te come
coordinatore, poi ricostruirai il team''. Ha rifiutato di ricostruire un team
cancellato ``con un colpo di spugna'' e ha deciso che non avrebbe più lavorato
per una grande azienda. Nessuno aveva pensato di spostarlo, e ha lavorato per
tre mesi in un open space di cento metri quadri con la sua scrivania e altre
ventinove vuote. Su una lavagna aveva scritto i nomi dei suoi ragazzi e li
cancellava man mano che, anche grazie ai suoi \emph{referral}, trovavano
lavoro. Cancellato l'ultimo nome, si è dimesso: per lui il lavoro era finito.

\paragraph{Team che si dissolvono.} Uno studente cita un articolo su Larian
Studios, il cui management è contrario a ricostruire ogni volta il team a fine
progetto. È lo stesso problema: si pensa ai ``giochi Blizzard'' o ai ``giochi
BioWare'', ma le persone che li hanno fatti non ci sono più; è rimasto solo il
brand, e non ha senso avere le stesse aspettative da un gioco fatto da persone
diverse. E anche chi sopravvive ai layoff vive pensando che il prossimo
potrebbe essere lui: alcune persone, licenziate nell'ultima ondata a
Blizzard, gli hanno detto ``erano due anni che lo sapevo''. Vivevano a Cork,
``un paesotto irlandese dove o vai per lavorare a Blizzard o non ci vai'',
avevano comprato casa: non ci si può aspettare chissà che morale.

\subsection{La cultura è forte quando si lavora per gli altri}

\begin{quote}
  \emph{La cultura è forte quando le persone lavorano per gli altri e a favore
  degli altri, e non per se stesse.}
\end{quote}

Se si dà feedback per far migliorare l'altro e non per imporre il proprio
modo, se si aiuta per il bene del progetto e non per fare bella figura, la
cultura si rafforza. È dai \textbf{piccoli gesti} che emerge la cultura
aziendale: ad esempio, dal non dover chiedere ``perché mi stai aiutando?''.

\paragraph{L'incognita.} \textbf{L'incognita è spesso dannosa quanto le
cattive intenzioni.} In un esperimento (``mi pare a Stoccolma, cinque o sei
anni fa'') si distribuivano per strada banconote da 500 euro, e la maggior
parte delle persone le rifiutava chiedendosi ``dov'è la fregatura?''
\footnote{L'esperimento è riportato in modo vago in lezione e non è stato
possibile identificarlo con certezza.}. Lo stesso accadeva al relatore,
arrivato a PG e a Stormind da Riot: a ``vedo che sei in difficoltà, posso
aiutarti?'' la risposta era sempre ``perché?''. Alla fine rispondeva: ``perché
quando sarò in difficoltà io vorrei che tu mi aiutassi''.

Quando si comincia a lavorare per gli altri e si smette di pensare ``così avrò
la promozione, così avrò l'aumento'', promozioni, aumenti e rispetto arrivano
da soli. A Blizzard aveva fallito tre volte i colloqui per essere promosso a
manager (inteso come coordinatore del team); quando ha deciso di smettere di
pensarci e di fare solo il bene del proprio lavoro, la volta dopo è stato
promosso. Funziona così: chi aiuta i colleghi e fa andare bene il progetto
viene notato (``questa persona fa migliorare cinque persone, non solo se
stessa: vale di più'').

\subsection{Il lavoro è valutato come opera d'insieme}

Il relatore riconosce che buona parte del suo successo è venuto dal fatto che
il suo team lavorava bene, non dalla sua bravura personale. Anche i progetti
degli studenti, presentati a Diego, a Enrico, ad Andrea, verranno valutati
\textbf{come opera d'insieme}. Nessuno dice ``il \emph{combat system} di questo
gioco l'ha fatto X, che bravo'': si dice ``BioWare ha fatto un bel gioco'',
``Blizzard ha fatto un bel gioco''.

Lo ha replicato anche lato HR. In molte aziende ci sono \textbf{bonus
personali} legati a obiettivi (se il gioco vende X, se la feature è pronta in
tempo). Il relatore li ha dimezzati: metà restano personali, metà diventano
\textbf{bonus di team}, che tutti ricevono se il team raggiunge gli obiettivi.
Un incentivo concreto a collaborare: non solo si dice che il lavoro di gruppo è
importante, lo si dimostra dandogli più valore economico. Aveva proposto il
100\% di bonus di team, ma gli è stato concesso il 50 e 50; ha funzionato
comunque, anche se sarebbe stato ancora più potente. Ovviamente, alla
domanda ``perché mi aiuti?'' si può anche rispondere ``perché voglio il
bonus'': un aiuto dato ``per il motivo sbagliato'' era un compromesso che era
disposto ad accettare.

\subsection{Il coraggio di chiedere aiuto}

Un partecipante osserva che aiutarsi dovrebbe essere la norma. Il modello
diffuso, invece, è ``devi essere forte e indipendente, non chiedere mai una
mano'', e vale in tutte le direzioni (dal maschilismo alla ``donna forte e
indipendente che non ha bisogno di nessuno''): \textbf{non ammettere di avere
bisogno di aiuto è sempre uno dei motivi per cui le aziende vanno male}, che si
tratti di uomini, donne o persone non binarie. Il problema è non avere la
sicurezza di poter alzare la mano.

Il relatore ha cercato di togliere ogni conseguenza dall'aver alzato la mano:
se si \emph{depenalizza} qualcosa, le persone non sono più disincentivate a
farla. Ha detto espressamente che non avrebbe mai valutato nessuno per il
numero di volte in cui chiedeva aiuto o per il numero di feedback negativi
ricevuti, e questo ha funzionato.

\subsection{Perché si dà un feedback}

\emph{Perché avete dato il vostro ultimo feedback? Che cosa volevate ottenere?}
Le risposte della classe:
\begin{itemize}
  \item per \textbf{riconoscere il valore} di un collega, anche verso il
        manager, perché qualcuno che non aveva visibilità vedesse quanto
        quella persona fosse efficace;
  \item per \textbf{spingere una persona a superarsi};
  \item per aiutare un collega a rendere il prodotto all'altezza di ciò che
        voleva rappresentare;
  \item per \textbf{genuinità}: c'era qualcosa che poteva essere fatto meglio.
\end{itemize}
Il relatore aggiunge che la forma cambia col tempo: all'inizio era ``devo
darti un feedback, hai dieci minuti?''; ora è un messaggio informale su
Discord, perché ha costruito un rapporto di fiducia e sa che non verrà letto
come ``l'HR manager mi scrive, sono nei guai?'', ma come ``Stefano è così, mi
dà feedback, resterà tra noi e lo risolveremo insieme''.

Un altro motivo è \textbf{risolvere incompatibilità caratteriali}: le
conversazioni più difficili, del tipo ``io e te non stiamo funzionando, cosa
possiamo fare diversamente?''. Sono importanti perché, se due responsabili non
vanno d'accordo, il team se ne accorge, comincia a raccontarsi storie (``si
odiano, ci saranno problemi'') e, per difendere il proprio manager, comincia a
giudicare male il lavoro dell'altro gruppo. È la dinamica da fermare:
\textbf{valutare gli altri solo per difendere la propria fazione}.

Dare e ricevere feedback sono \textbf{due muscoli da allenare} (il relatore è
sempre stato più bravo a darlo che a riceverlo), e si impara col tempo. È come
fare colloqui: fare colloqui è una skill a parte, che non indica quanto si è
bravi nel proprio lavoro. Essere scartati non significa non saper fare quel
lavoro: spesso significa aver fatto male il colloquio o che è stato preferito
qualcun altro.

\paragraph{Feedback pubblici e privati.} Uno studente racconta di aver dato un
feedback negativo a un superiore, senza un motivo particolare, per dimostrare
a sé e ai colleghi che non si viene licenziati per questo. Il relatore coglie
l'occasione per precisare che \textbf{i feedback negativi si danno
individualmente}: lui dà moltissimi feedback di cui il resto dell'azienda non
sa nulla. Reagirebbe molto male se, in un meeting, qualcuno alzasse la mano per
dire ``comunque questa cosa non l'hai fatta bene'': quello non è feedback, è
discussione, accettabile solo se è un brainstorming e se l'opinione è stata
chiesta. Alcune aziende hanno piattaforme pubbliche di feedback, ma spesso
servono solo ad applaudire; lo stesso accade nelle aziende che
\emph{obbligano} a lasciare un certo numero di feedback a settimana, che
diventano ``lui è bravo, lei è competente, anche questa settimana l'abbiamo
fatto''. Bisognerebbe chiedere a chi le adotta cosa pensa di ottenere (forse
trasparenza), ma difficilmente una persona ``attaccata'' pubblicamente reagisce
bene.

% ---------------------------------------------------------------------
\section{Among Us: la caccia al traditore}
\label{sec:z1-among-us}

Prima della pausa pranzo il relatore mostra tre film di fantascienza di
qualche decennio fa e chiede che cosa li accomuni\footnote{I titoli non sono
leggibili nella trascrizione: viene citato \emph{Event Horizon} da uno
studente, e il film ``in mezzo'' è indicato come ``l'originale'',
verosimilmente \emph{La cosa} (\emph{The Thing}).}. In ognuno c'è un gruppo di
persone con un obiettivo e la presenza, reale o percepita, di qualcuno che
rema contro: \textbf{l'energia del gruppo viene spesa per trovare il
responsabile}.

\paragraph{L'esperimento degli alieni.} Nei master in presenza il relatore
propone un esercizio. Distribuisce segretamente cartellini plastificati con
scritto ``essere umano'' o ``alieno''; le carte alieno contengono istruzioni di
comportamento con cui gli alieni, che cercano di fingersi umani, si
tradirebbero (dare sempre ragione a chi ha parlato prima, finire ogni frase
con una domanda, parlare in un certo modo). Lo scopo: in dieci minuti formare
gruppi, e \textbf{vince il gruppo più numeroso senza alieni}; basta un alieno
per perdere. Si formano sempre gruppi piccoli (due, cinque, sei persone),
perché la tendenza è ``mi fido al massimo di voi due''.

Ma vincono i gruppi più grandi, per un motivo semplice: il relatore
\textbf{non distribuisce nessuna carta alieno}. Mostra a tutti entrambe le
carte, ma consegna solo carte ``umano''. Vuole simulare ciò che accade in
ufficio: la tendenza a percepire sempre qualcuno che rema contro (i manager,
l'altro team, quel collega), quando in realtà non c'è nessuno, ``è tutto nella
nostra testa''. Quando chiede le motivazioni, emergono sospetti
\emph{fabbricati}: ``faceva sempre domande'', ``parlava solo all'infinito''
(``no, parla così e basta''). Ci si fabbrica da soli i motivi per non fidarsi
degli altri, ed è comunissimo.

\paragraph{Among Us.} L'autosabotaggio è, per tornare ad \emph{Among Us},
\textbf{concentrarsi sul cercare cadaveri e traditori invece di fare le
task}: non si vince perché non si sono completate le task, mentre si cercava il
colpevole. Succede tantissimo negli uffici, ed è il motivo per cui l'80\% dei
problemi che il relatore risolveva erano di comunicazione. Entrando nei team
sentiva immancabilmente: ``facciamo il possibile, ma questa lotta continua con
l'altro gruppo è estenuante''. ``Perché c'è questo problema?'' ``Non lo so,
c'era già quando sono arrivato.'' Un paradosso: c'è un problema, nessuno sa
dire quale sia, ma c'è sempre stato e quindi lo si accetta.

\paragraph{L'esperimento delle scimmie.} È l'esempio, ``finto o reale'', delle
scimmie in gabbia\footnote{È il noto racconto delle ``cinque scimmie'', spesso
presentato come esperimento scientifico ma probabilmente apocrifo (lo stesso
relatore lo definisce ``finto barra reale''); nella trascrizione il numero
delle scimmie varia.}. In una gabbia con alcune scimmie pende una banana
elettrificata: chi la tocca fa prendere la scossa a tutte. Le scimmie imparano
e, quando una si avvicina alla banana, le altre la picchiano. Poi si
sostituisce una scimmia alla volta con una nuova: la nuova va verso la banana,
le altre la picchiano, e smette. Si continua finché nella gabbia ci sono solo
scimmie nuove, poi si toglie la corrente. Nessuna scimmia tocca la banana:
nessuna ha mai preso la scossa, ma tutte hanno imparato che chi tocca la
banana prende le botte.

Nei team è lo stesso: il motivo per cui si ``odia'' l'altro team si
trasforma, ma un vero motivo non c'è. Ne derivano la sfiducia, il non alzare
mai la mano per chiedere aiuto, il non dire mai ``bel lavoro'' agli altri
perché non sia mai che qualcuno ti veda fare un complimento ``a quelli''.
Risolvere queste dinamiche è lungo e tedioso, e sempre pieno di ``non mi
ricordo'', ``io non c'ero'' e ``abbiamo sempre fatto così'', la cosa più
frustrante per chi cerca di risolverle.

% ---------------------------------------------------------------------
\section{Il principio del default to trust}
\label{sec:z1-default-trust}

\subsection{La definizione}

Il relatore indica questa come la slide più importante della presentazione.
\textbf{Default to trust} è uno dei valori aziendali di Riot Games, che lui ha
tradotto, adattato e ``esportato'':

\begin{quote}
  \emph{Approccio ogni situazione di conflitto partendo dal presupposto che una
  persona abbia agito in un certo modo perché, con le informazioni a sua
  disposizione in quel momento, ha pensato che fosse la soluzione migliore, e
  non perché sia un'incapace o volesse danneggiare il progetto
  intenzionalmente.}
\end{quote}

In altre parole: \textbf{di default, mi fido}.

\subsection{L'esempio del ritardo al daily stand-up}

A Riot il relatore aveva un \emph{daily stand-up} dal vivo alle 15 con il suo
team di venti persone. Una ragazza arrivava regolarmente alle 15:15, per una,
due, tre settimane. Alla quarta ci ha parlato. \emph{Come iniziereste la
conversazione?} Le proposte della classe:
\begin{itemize}
  \item sottolineare l'importanza di esserci dall'inizio, perché potrebbe
        perdersi qualcosa che le interessa, e per rispetto dei colleghi
        presenti;
  \item ``devi arrivare in orario'', spiegando perché è importante;
  \item \textbf{chiedere perché}.
\end{itemize}

Il relatore osserva, anche dalle sue consulenze fuori dai videogiochi, che
\textbf{più si sale con la seniority}, più la risposta è ``devi smettere di
arrivare in ritardo'', spesso con un \emph{or else} finale (``altrimenti
succede quest'altra cosa''). Più si scende con l'età, e quindi con meno
abitudini consolidate, più la prima domanda diventa ``dimmi perché arrivi in
ritardo''. Già spiegare perché è importante è ``miglia avanti'' rispetto al
semplice ``smetti''.

Partire dal perché gli ha permesso di scoprire la causa: il meeting precedente
della ragazza, in un'altra zona del campus a un quarto d'ora di distanza,
finiva sempre tardi, e lei doveva attraversare il campus di corsa. Non aveva
la confidenza per dire all'altro manager, che insisteva perché si restasse
fino alla fine, che doveva andare; sacrificava quindi il meeting del relatore,
sapendo che era il più comprensivo dei due.

Il relatore si è offerto di \textbf{mediare}: ha detto all'altro manager che
o anticipava l'inizio del suo meeting, così da finire comunque in tempo,
oppure lasciava uscire la ragazza, perché quel meeting gli serviva. C'è anche
una questione di \textbf{economia di scala}, spesso sottovalutata: se venti
persone aspettano qualcuno per dieci minuti, si sono sprecati dieci minuti per
venti. Per lo stesso motivo il workshop è iniziato alle 10:05: iniziare alle
10:10 avrebbe significato far aspettare cinquanta persone cinque minuti in
più. A Riot c'era una scrivania vuota con un pupazzo gonfiabile, \textbf{Mr.\
Late}: sommando tutti i minuti di ritardo a meeting, call e conversazioni, in
un anno veniva fuori praticamente \emph{una persona a tempo pieno}. Serviva a
ricordare che anche cinque minuti contano, e che bisogna rispettare gli altri
arrivando in orario o almeno avvisando.

\paragraph{Gli effetti.} La conversazione era avvenuta nel primissimo periodo a
Riot, quando il team non aveva ancora motivo di fidarsi di lui. Risolvendo il
problema ha eliminato anche tutto lo strato di micro-dinamiche spiacevoli che
un manager impara a vedere (due colleghi che si guardano con la risatina
quando lei entra in ritardo). La ragazza ha cominciato a entrare tranquilla,
con una certa leggerezza: non le faceva piacere arrivare in ritardo (``ci sono
anche persone che lo fanno apposta, e quello è un problema da affrontare
diversamente''). Soprattutto, con i colleghi non parlava di lui come del capo
che minacciava, ma come di chi le aveva risolto un problema: era diventato
un \textbf{alleato}. La settimana dopo altre tre persone sono andate da lui
con un problema da risolvere. Ripetuto cinque, sei, sette, dodici volte, in un
paio di mesi poteva dire quasi qualsiasi cosa nei meeting sapendo che
richieste e critiche sarebbero state percepite come ``me lo chiede perché vuole
darmi una mano'', non per fare bella figura o per presentare i successi del
team ai suoi superiori.

\subsection{Condividere i successi, schermare i fallimenti}

Un altro elemento fondamentale (ripreso nella seconda giornata): il manager
\textbf{condivide i successi e protegge il team dai fallimenti}. ``Se va male
qualcosa è colpa mia; se va bene qualcosa, voglio che ne parliate voi.'' Il
lavoro del manager non dovrebbe essere valutato su ciò che presenta del lavoro
del team: tutti sanno comunque chi è il manager. Tanto vale lasciar presentare
il team, che potrà dire con orgoglio di aver presentato a Mike
Morhaime\footnote{Cofondatore ed ex presidente di Blizzard; nella
trascrizione ``Mike Morayme''.} la propria campagna europea per
\emph{StarCraft~II}. Ci sono manager che dicono ``come mi prendo le colpe, mi
prendo i meriti'', ma i meriti li hanno già: è inutile ostentarli. Con i
fallimenti invece il manager è la \textbf{rete di sicurezza} del team. Il
relatore ne ha fatto una questione di principio: \textbf{non scaricare mai la
colpa su qualcun altro}.

\paragraph{Colpevole o causa.} Un esempio che il relatore attribuisce a Pirate
Software, creatore di contenuti ed ex dipendente Blizzard\footnote{Il relatore
aggiunge che si tratterebbe del figlio dello sviluppatore a cui si sarebbe
ispirato il personaggio del giocatore ``senza vita'' nell'episodio di
\emph{South Park} dedicato a \emph{World of Warcraft}. L'informazione è
riportata come detta in lezione e non è verificata.}. Contattando il servizio
clienti Blizzard per motivi personali, aveva scoperto che i colleghi gli
avevano scritto nei ticket informazioni confidenziali che non avrebbero dovuto
condividere in quel modo, nemmeno con un collega. Non ha segnalato le persone
da punire: ha detto ``c'è un problema, evidentemente le policy non sono
chiare''. Quando gli hanno chiesto i nomi, ha rifiutato: ``se vi dico chi è,
punite quelli, e gli altri cinquecento continueranno a farlo; non avremo
cambiato la causa alla fonte''.

\begin{quote}
  \emph{È inutile trovare il colpevole. Voglio trovare la causa. Perché il
  colpevole cambia.}
\end{quote}

Il colpevole a volte è vittima della situazione: per questo la definizione
dice \emph{``con le informazioni a sua disposizione in quel momento''}. Se il
relatore avesse rimproverato la ragazza in ritardo, non solo non avrebbe
risolto il problema, ma avrebbe creato un'\textbf{ingiustizia}: lei, che
attraversava il campus di corsa per una colpa non sua e di cui non poteva
parlare, sarebbe stata anche rimproverata.

Nessuno fa un brutto lavoro apposta: nessuno è contento quando sbaglia o
perde. Perché al lavoro diamo per scontato che gli altri sbaglino
intenzionalmente? Se qualcuno ha fatto una cosa male, è perché in quel momento
credeva fosse la scelta migliore. La domanda giusta non è ``perché hai fatto
questa cosa che è stata un fallimento?'', ma:

\begin{quote}
  \emph{Che cosa ti ha fatto pensare che questa fosse la scelta giusta?}
\end{quote}

Magari erano limiti di budget, di tempo, di scopo: si agisce su quelli, non
sulla persona. Se si cambia la persona ma le si danno gli stessi fattori, è
molto probabile che arrivi alle stesse conclusioni, e ci si troverà a dover
``disciplinare'' anche lei.

\paragraph{La reputazione costruita sui rumor.} Durante il giro di
presentazione in un ufficio (non in nessuna delle aziende citate), al relatore
è stato detto: ``questo è il dipartimento di design; se ti serve qualcosa
chiedi a lei, non a lui, perché lui consegna sempre in ritardo''. Oltre a
trovarlo poco professionale, una settimana dopo ha chiesto: ``cosa ti ha fatto
questa persona?'' ``A me niente, però ho sentito che\dots''. Si torna alle
scimmie: magari quella persona non consegna in ritardo da mesi, ma ormai lo
dicono tutti. È un circolo vizioso da fermare subito.

\paragraph{Niente escalation inutili.} Lo stesso vale per il feedback: quando
qualcuno va dal relatore dicendo ``ho un problema con questa persona'', lui
chiede: \textbf{``perché ne stai parlando a me e non a lei?''}. ``Perché sei il
mio capo'' non basta. Passare da lui significa dare per scontato che l'altra
persona non abbia interesse a risolvere, mentre magari basta parlarle. È la
strada più facile (``lo dico a Stefano, ci pensa Stefano''), e proprio per
questo bisogna mettere dei paletti: il manager interviene quando serve, ma le
cose che si possono gestire da soli vanno gestite da soli. (Quando invece il
confronto diretto fallisce, l'escalation è legittima: si veda la
sezione~\ref{sec:z1-qa-finale}.)

% ---------------------------------------------------------------------
\section{Traiettorie di carriera e valutazione delle performance}
\label{sec:z1-traiettorie}

\subsection{Quale traiettoria vorreste?}

Il relatore mostra un grafico con quattro curve colorate, contrassegnate
anche da lettere (A, B, C, D) per eventuali persone daltoniche (``a me è
capitato che qualcuno scoprisse di esserlo durante una lezione''). Le curve
rappresentano la performance lavorativa nel tempo. Due domande: \emph{quale
vorreste essere? E quale pensate realisticamente di essere?} Un partecipante
aspira a B e pensa di essere C; quasi tutti scelgono B.

Perché non A, che sale più rapidamente? Un partecipante risponde che una
crescita così repentina può essere più dannosa che utile, come una ``rush'' di
studio molto faticosa che poi non si riesce a sostenere. Il relatore propone
una lettura suggerita da un amico psicologo, che spiega la forma del grafico:
di B, C e D \textbf{vediamo dove finiscono}; di A no. Potrebbe continuare a
salire oltre la scala del grafico, o crollare sotto D. Non saperlo ci spinge a
preferire B. (A chi obietta che anche le altre proseguono oltre il bordo, il
relatore precisa: il bordo destro è ``oggi'', e di A non sappiamo come sia
arrivata fin lì.)

\subsection{La carriera non è una linea retta}

Tutti immaginano la carriera come una linea: si parte da X e, attraverso
stage e posizioni, si arriva a Y. In realtà carriera e performance avranno un
andamento irregolare, con alti e bassi. Come non si accetta che un team possa
temporaneamente \emph{underperformare}, tanto meno lo si accetta da se stessi.
Bisogna accettare che ci saranno momenti bassi.

Per questo il relatore \textbf{detesta le valutazioni annuali delle
performance}: prendono un punto casuale della traiettoria, che se si è
fortunati è il picco più alto e se si è sfortunati il più basso, e lo usano
per valutare un intero anno. Insegna ai manager a fare più punti di controllo
e a guardare altre cose:
\begin{enumerate}
  \item \textbf{quanto tempo passa tra un picco e l'altro}, cioè quanto
        \emph{cooldown} serve a una persona. È il motivo per cui, dopo la
        Games Week, tutti staccavano una o due settimane. Il relatore stesso,
        che si impegna moltissimo e tende al burnout, non potrebbe rifare lo
        stesso workshop il lunedì e il martedì successivi con la stessa
        qualità;
  \item \textbf{la differenza tra la miglior performance e la peggiore}: che
        cosa ha portato a lavorare così bene, e così male? Con chi si
        lavorava, su cosa ci si concentrava, quanto tempo e budget c'erano?
        L'obiettivo è capire come \emph{replicare} le condizioni che fanno
        lavorare bene, invece del solito ``hai lavorato male, è colpa tua''.
        Il relatore ha avuto molti manager per cui quando lavorava male era
        colpa sua e quando lavorava bene ``che bravo'', ma sapeva che raramente
        era l'unico fattore della propria performance. È la fallacia del ``se
        sai fare 100, perché non fai sempre 100?'', come nella scena di
        \emph{Schindler's List} sull'operaio che produce una vite in pochi
        secondi: non sempre un ritmo è sostenibile;
  \item \textbf{il fallimento di oggi confrontato con quello di sei mesi fa}.
        Il relatore non è contento della presentazione di oggi (l'avrebbe
        migliorata per altri due anni), ma rispetto a quelle di cinque anni fa
        vede un grande miglioramento di chiarezza: il suo peggior fallimento di
        oggi equivale al suo primo successo di inizio carriera.
\end{enumerate}
Quello che conta è la \textbf{traiettoria} dall'inizio dell'anno a oggi, non
i singoli punti. E se si riesce a capire insieme quando e perché si ha
successo e quando e perché si fallisce, si può lavorare per replicare solo le
situazioni che fanno lavorare bene.

\paragraph{Lo stesso vale per i giochi.} È il motivo principale per cui i
giochi delle grandi case vanno bene o male. Quando si riescono a replicare le
condizioni che hanno prodotto \emph{Warcraft~II}, esce \emph{Warcraft~III};
quando non ci si riesce, esce ``il nuovo \emph{Dragon Age}'': la gente non è
arrabbiata perché il gioco non è bello, ma perché non è come gli altri
\emph{Dragon Age}; non ne replica il \emph{feeling}, anche se potrebbe
essere un gioco favoloso. L'eccezione positiva è il passaggio da \emph{Dune} a
\emph{Dune~II}: il primo era di fatto un'avventura grafica, il secondo ha
generato il genere degli strategici in tempo reale. Gli sviluppatori avevano
capito che il successo non stava nel tipo di gioco, ma nell'ambientazione che
attirava le persone: un bel gioco in quell'ambientazione sarebbe piaciuto
comunque. Il relatore comprò \emph{Dune~II} aspettandosi un seguito
dell'avventura e, pur sorpreso, lo finì.

\subsection{Lying, hiding, faking}

Quando si è valutati su un singolo punto nel tempo, e da quel punto dipendono
bonus, stipendio e il rischio di licenziamento, si comincia a fare
\textbf{lying, hiding e faking}:
\begin{description}
  \item[Lying] dire che è sempre andato tutto bene, che si finirà sempre in
        orario, che non c'è nessun problema;
  \item[Hiding] nascondere: ``non era un mio task'', ``non c'ero'', ``non ho
        visto niente'';
  \item[Faking] fingere, anche in cose apparentemente innocue come ``sì, sì,
        ho capito tutto'' quando si ha solo paura di dire ``me lo rispieghi?
        Non ho capito niente''.
\end{description}

\subsection{``Non ho capito'' e la conferma delle istruzioni}

Dire ``non ho capito, ripetimelo'' è ancora visto come una figura da sciocco.
Il relatore non si arrabbierà mai se qualcuno gli chiede di rispiegare
un'intera slide: le cose ovvie per lui possono non esserlo per altri. Più si
conosce qualcosa, \textbf{più è difficile spiegarla}, perché si danno per
scontate molte cose. Ne parlava proprio con Diego, che gli diceva che fare
l'insegnante diventa sempre più difficile man mano che si diventa esperti. È
come spiegare \emph{EVE Online}, uno dei giochi più complessi in circolazione:
chi lo conosce bene dà per scontato che ``ovviamente non puoi fare questa
cosa'', mentre per un nuovo giocatore non è affatto ovvio. A scuola chi dice
``non ho capito'' viene visto come lo ``scemo''; nessuno si chiede ``non è che
l'ho spiegato male?''.

Il racconto più ``traumatico'' dei ragazzi che fanno stage nelle grandi
aziende è sempre lo stesso: ``Cinzia, devi fare questa cosa.'' Cinzia la fa.
``Non te l'avevo chiesta così.'' Il problema è che nessuno ha
\textbf{confermato che Cinzia avesse capito}. Il consiglio, valido sempre:
\begin{itemize}
  \item quando si \emph{ricevono} istruzioni, ripeterle: ``ho capito bene?
        Vuoi che faccia questa cosa in questo modo?'';
  \item quando si \emph{danno} istruzioni, chiedere a chi le riceve di
        ripeterle, non perché lo si ritenga incapace, ma per capire che cosa
        ha capito dalla propria spiegazione (``quindi vuoi un quadrato?''
        ``No, ho detto un triangolo\dots'' ``Però mi hai detto una forma con
        quattro angoli.'' ``Hai ragione: voglio un quadrato'').
\end{itemize}
Frasi come ``non te lo ripeto due volte'' o ``stammi dietro'' portano solo a
fare le cose come si pensa che vadano fatte e a sentirsi dire a fine progetto
``non la volevo così'': che si tratti di una feature, di un design o di un
modello 3D. (Un partecipante osserva che lo stesso accade insegnando le lingue:
chi insegna dà per scontato, ad esempio, che si sappia dell'esistenza del
genere neutro in tedesco.)

\paragraph{Come convincere il team a chiedere?} Il relatore mette sempre come
premessa che \textbf{l'obiettivo non è trovare la persona più intelligente
nella stanza, ma fare un buon lavoro}. Che Gianni capisca alla prima e un altro
alla quinta non gli interessa, se alla fine entrambi lavorano bene: ripetere
cinque volte è comunque meno lavoro che far rifare le cose sette volte. È un
investimento a lungo termine, e non dissimula alcun fastidio: davvero non gli
dà fastidio. Ma ripetere le stesse parole non aiuta: ``non ho sentito'' si
risolve ripetendo; ``non ho capito'' si risolve cambiando spiegazione. Meglio
ancora: ``dimmi che cosa hai capito tu, e vediamo cosa non torna''. Se
intravedesse nel team uno sguardo del tipo ``è la quinta volta che lo
chiede'', interverrebbe: spesso chi giudica è proprio chi dice ``ho capito'' e
poi sbaglia. Preferisce di gran lunga chi chiede cinque volte.

\subsection{Il feedback più utile che abbiate ricevuto}

\emph{Qual è il feedback più utile che abbiate mai ricevuto?} Il relatore ne ha
almeno dieci che gli hanno cambiato il modo di lavorare. Le risposte:
\begin{itemize}
  \item un semplice \textbf{``grazie''} per un lavoro fatto, più utile di un
        ``questa cosa fa schifo'';
  \item un feedback sulla propria \textbf{impulsività} nel rispondere: fermarsi
        un attimo e capire di cosa si sta parlando prima di intervenire;
  \item \textbf{essere se stessi}, anche al lavoro, e non scusarsi
        continuamente per ogni cosa.
\end{itemize}
I feedback che toccano di più sono quelli \textbf{veri}. Il ``grazie''
significa ``ho visto quello che hai fatto e l'ho apprezzato''. Quelli negativi
che si ricordano sono quelli che hanno fatto realizzare ``questa cosa che
pensavo di saper fare, non la so fare'', oppure ``la so fare, ma questa volta
non mi è venuta'' (non è che se una cosa una volta non riesce non la si sa più
fare: può capitare di sbagliare una porta). Arrivano sempre da un
\textbf{punto cieco}: qualcosa che non si sapeva fosse percepito dagli altri.
Conta più ciò che le persone ti dicono di ciò che pensi che pensino di te; e
ciò che il relatore pensa che gli altri pensino di lui è sempre terribile,
mentre il loro giudizio si rivela molto più positivo del previsto.

\paragraph{Il denominatore comune siete voi.} Si potrebbe pensare che questi
concetti non siano applicabili da chi arriva in un'azienda senza poterne
cambiare la cultura. Ma c'è sempre un denominatore comune a cui applicarli:
\emph{se stessi}. La cultura aziendale può non spingere a dare feedback, può
vedere male chi chiede aiuto; ma il relatore ha visto che quando lui lo fa, le
persone intorno tendono a farlo, e un gruppo che lo fa lavora meglio. Perfino
i CEO più ostici a un certo punto gli chiedevano: ``perché il tuo team ora
performa meglio degli altri?''. La differenza è il default to trust: una
pratica che col tempo fa operare il team \emph{come un team} e non come un
insieme di individui. Un videogioco è un progetto corale: c'è chi fa giochi da
solo, ma non si fa da soli \emph{Red Dead Redemption~2}; e se il gruppo non è
affiatato e allineato non esce \emph{Red Dead Redemption~2}, esce
\emph{Anthem} (``oggi la mia fissazione è \emph{Anthem}, domani sarà
un'altra'').

% ---------------------------------------------------------------------
\section{Come dare un feedback}
\label{sec:z1-dare}

\subsection{Concentrarsi sull'azione e sui suoi effetti}

Il primo principio: \textbf{il feedback non riguarda mai la persona}. Non si
dice ``tu hai fatto questa cosa che non va bene'', ma ``questa cosa non va
bene: come facciamo a non farla più?''. Che sia un comportamento, un'azione,
una decisione o una strategia, ci si concentra su quello, e si
\textbf{esplicita l'effetto} che ha. Nell'esempio del ritardo l'effetto era
``quindici minuti di attesa moltiplicati per venti persone''. In un progetto
può essere una catena:
\begin{quote}
  se non rispetti le deadline, l'animazione che viene dopo resta ferma; se
  loro sono in ritardo, la \emph{vertical slice} è in ritardo; se la vertical
  slice non è pronta per la Gamescom, non andiamo alla Gamescom e non troviamo
  un publisher; se non troviamo un publisher, chiudiamo.
\end{quote}
Più la conseguenza è grave, più dovrebbe sembrare piccolo il gesto di dire
``datemi una mano'': è il \emph{male minore}. Poi, a cose fatte, ci si dà
feedback: cosa ha funzionato, cosa no, cosa fare di più, di meno, cosa smettere
di fare.

\subsection{Il pericolo degli ambienti solo positivi}

Gli ambienti dove ci si fa sempre grandi complimenti e nessuno parla dei
problemi creano difficoltà vere alle aziende. Il problema non è mai la
mancanza di complimenti, ma un \textbf{eccesso di ottimismo}: si comincia a
pensare che le cose andranno come si spera, perché tutti intorno danno solo
feedback positivi.

A Blizzard i feedback determinavano gli \textbf{aumenti annuali}: ognuno
doveva indicare cinque persone che avrebbero dato feedback su di lui, e
ovviamente tutti indicavano gli amici, o colleghi di progetti andati bene. Un
sistema nato con le migliori intenzioni diventava uno \emph{scambio di voti}
(``se mi metti follow ti metto follow''), senza alcun miglioramento.

Al contrario, \textbf{le conversazioni più brutte sono state le più utili}:
affrontare i licenziamenti a Riot, o la cancellazione di \emph{Project Titan}
a Blizzard, il progetto che doveva essere il nuovo MMO dopo \emph{World of
Warcraft} e da cui poi è nato \emph{Overwatch}\footnote{Nella trascrizione si
legge anche ``doveva essere Warcraft 2'', probabilmente un lapsus o un errore
di trascrizione.}. Andando a vedere perché lo si cancellava, c'erano solo
motivazioni razionali e pragmatiche, di business, chiare a tutti, non un ``Mike
Morhaime ci odia''. Sono anche le conversazioni in cui ha lavorato su se
stesso: è molto \textbf{avverso ai conflitti}, e questo a volte gli ha
procurato conseguenze molto più gravi di quanto sarebbe costato un piccolo
disagio. Dieci anni fa faceva fatica ad accettare questo difetto; ha dovuto
subire le conseguenze di conflitti evitati e poi cresciuti.

\subsection{Che cosa fare}

Le linee guida della slide, che il relatore legge rapidamente perché saranno
disponibili nel materiale:
\begin{description}
  \item[Attenzione al linguaggio del corpo] dare un feedback a braccia
        incrociate o guardando altrove è diverso dal guardare l'altro negli
        occhi.
  \item[Sii specifico] con esempi tangibili.
  \item[Esplicita gli effetti] come visto sopra.
  \item[Discuti il prima possibile] non aspettare un mese per parlare di
        qualcosa successo ieri, altrimenti nessuno ricorderà più bene cosa sia
        accaduto.
  \item[Concentrati sull'azione e sul problema, non sulla persona] il nemico è
        l'azione, non la persona che si ha davanti. Non bisogna far cambiare
        l'altro: bisogna evitare che il problema si ripresenti.
\end{description}

\paragraph{Gli insulti sono rivolti al ruolo.} Un esempio solo apparentemente
scollegato. Il relatore ha lavorato su \emph{League of Legends}, ``non
esattamente noto per avere una community da orsetti del cuore'', in varie
\emph{task force} sulla tossicità. A chi smetteva di giocare veniva inviato un
sondaggio sul perché; il 33\% indicava la tossicità della community. Dopo varie
campagne di sensibilizzazione e cambiamenti si è scesi al 30\%: sembra solo
un 3\% in meno, ma valeva circa \textbf{120 milioni di revenue in più
all'anno}, perché quei giocatori restavano esposti a skin e contenuti e
continuavano a comprare. Anche qui, piccoli guadagni che si sommano.

Lavorando con degli psicologi, aveva notato che gli insulti \textbf{non sono
quasi mai rivolti alla persona, ma al ruolo}: ``il mio jungler fa schifo'',
``il mio DPS non è capace'', ``il nostro portiere è una pippa''. Chiamare
qualcuno per ruolo lo \emph{depersonalizza} e rende più facile insultarlo. Lo
stesso nell'esempio del giro dell'ufficio: non ``non chiedere a Giulio'', ma
``non chiedere a quel grafico, consegna sempre in ritardo''. Concentrarsi
sull'azione, senza attaccare la persona, fa sì che il feedback non sia (e non
venga percepito come) un attacco personale.

\paragraph{Considera il carattere.} Qui tornano utili i test di personalità.
Servono più agli studenti che al relatore (che sa già che usciranno 53 profili
diversi, senza un pattern): bisogna conoscere il proprio profilo e
\textbf{condividerlo con chi lavora con noi}, chiedendo anche agli altri di
fare lo stesso. In ufficio il relatore ha usato una versione più avanzata, a
pagamento, con un report di circa trenta pagine: il \textbf{profilo
DISC}\footnote{Nella trascrizione ``disk''. Si tratta del modello DISC
(Dominance, Influence, Steadiness, Conscientiousness).}. Ha mappato tutta
l'azienda e ha fatto stampare il profilo di ciascuno sulla sua scrivania, così
che tutti sapessero con chi stavano parlando. Capire come ragionano gli altri
(``è introverso, non gliene parlo davanti a tutti'') è una delle cose che gli
sono tornate più utili: \emph{adattare ciò che si vuole dire a chi si ha
davanti}.

\paragraph{Assicurati di non disturbare.} Sotto il profilo di ognuno c'era un
piccolo \textbf{semaforo}: verde, ``parlami pure, interrompimi''; giallo,
``interrompimi solo se è urgente''; rosso, ``mi sto concentrando, non mi
parlare''.

\paragraph{Lascia finire il lavoro.} Durante gli eventi dal vivo il relatore si
era reso conto di quanto fosse grave dare feedback mentre qualcuno sta ancora
lavorando. Se una cosa sta andando male, dirlo in quel momento non aiuta: tra
la pressione di finire e quella del feedback, peggiorano entrambe le cose, e
anche il rapporto. Meglio lasciarla finire, anche male, e parlarne dopo.

\subsection{Che cosa non fare}

\begin{description}
  \item[Non basarti su supposizioni] ``mi hanno detto che ieri in un meeting
        hai fatto\dots'': o c'eri e l'hai visto, o non esiste.
  \item[Non giudicare la persona.]
  \item[Non fornire feedback vago.]
  \item[Non accettare e non fornire feedback di terze parti] Se qualcuno vi
        dicesse ``mi hanno detto che hai fatto un brutto lavoro'', il primo
        pensiero non sarebbe ``cosa ho fatto male?'', ma \emph{``chi te l'ha
        detto? Chi è la vipera che ho in seno? Chi è l'impostore?''}. Siamo
        fatti così. Ogni conversazione che inizia con ``mi hanno detto che'' è
        inutile: l'altro non ascolterà, penserà solo a chi dare la caccia, e
        spesso penserà alla persona sbagliata, creando antipatie tra persone
        che non c'entravano nulla. Quando al relatore dicono ``mi hanno detto
        che lui ha fatto\dots'', risponde ``vai a parlarne con lui''; se gli
        dicono ``mi hanno detto che tu hai fatto\dots'', risponde ``mandami la
        persona che te l'ha detto, altrimenti non ne parliamo''. Non per
        cattiveria: se chiedesse ``che cosa ho detto di preciso a Marco?'' e la
        risposta fosse ``non lo so, non c'ero'', di cosa si starebbe parlando?
        \textbf{Parlate con la fonte}, altrimenti fate come se quel feedback non
        esistesse.
  \item[Non dare consigli non richiesti] come visto, non portare la soluzione:
        la soluzione deve venire da chi riceve il feedback.
  \item[``Non mi piace'' non è un feedback] ``Penso che questo design non
        funzioni perché, visto che il nostro obiettivo è\dots'' ha argomenti;
        ``non mi piace il combat system'' è un'opinione. Anche ``io l'avrei
        fatto diversamente'' non è un feedback.
  \item[Accettare non vuol dire fare come proponi] deve essere chiaro a chi dà
        il feedback: l'aspettativa è essere ascoltati e presi in
        considerazione, ma se la responsabilità è dell'altro, deciderà
        l'altro.
  \item[Non offenderti] se il feedback non viene seguito. Restarci male
        significa pensare di saperne più dell'altro, cioè \emph{venire meno al
        default to trust}. ``Decidi tu'' significa ``mi fido di te, perché sei
        lo specialista, e mi fido che mi ascolterai e farai la cosa migliore
        per tutti'', che potrebbe anche non essere ciò che ho detto io. Se
        voglio che l'altro faccia sempre come dico io, mi sto fidando solo di
        me stesso.
\end{description}

\subsection{Il primo feedback e il debriefing}

La cosa più difficile è dare il \textbf{primo} feedback: è un muscolo da
allenare. Chi non l'ha mai fatto (soprattutto con questa struttura) dovrebbe
provarci, anche su minuzie (``abbiamo concordato di usare WhatsApp nel gruppo:
non potremmo usare Telegram?''), anche solo per vedere come va. Altre
indicazioni:
\begin{itemize}
  \item feedback \textbf{solo uno a uno}, mai davanti ad altri, soprattutto se
        negativo;
  \item \textbf{prepararsi} prima;
  \item fare il \textbf{debriefing del feedback}: dopo, fermarsi a pensare a
        com'è andato. Il relatore, quando media un feedback tra due persone,
        si sente spesso rispondere ``benissimo'' da una e ``malissimo''
        dall'altra: la percezione è molto diversa. Non si fa problemi a tornare
        dalla stessa persona: ``ieri ti ho dato un feedback: com'è andata? Sei
        soddisfatto? Come poteva andare meglio?''. Con le persone con cui si ha
        meno a che fare, questo aiuta molto la volta successiva.
\end{itemize}

% ---------------------------------------------------------------------
\section{Come ricevere un feedback}
\label{sec:z1-ricevere}

Per il relatore è la parte più difficile. Si riassume in quattro passi:
\begin{enumerate}
  \item \textbf{Ascolta}: si parte sempre ascoltando.
  \item \textbf{Chiedi dettagli}, se serve, ma per capire meglio, non per
        difenderti.
  \item \textbf{Accetta, a prescindere}: non si sta decidendo chi ha ragione;
        si sta ascoltando un'opinione. Non c'è nulla da cui difendersi.
  \item \textbf{Agisci}: fatti un piano. Può essere banale (``mi hanno detto
        che interrompo sempre: cerco di non interrompere''), oppure ``devo
        migliorare questa cosa: come?''. Pensa anche a \emph{chi} potrebbe
        aiutarti e a \emph{che supporto} ti serve: il piano non deve essere
        per forza ``devo risolvere tutto da solo''.
\end{enumerate}
Una domanda che molti sottovalutano: \emph{``sono junior in un'azienda: sarò
mica il primo ad avere questo problema?''}. Probabilmente qualcun altro l'ha
già avuto: basta chiedere ``tu come hai fatto?''.

\subsection{Il feedback positivo}

Tutto questo vale per i feedback costruttivi negativi, a cui si pensa di
default; ma esistono anche i \textbf{feedback positivi}, e vanno usati. Quando
si riceve un ``grazie'' o un ``bravo'', la domanda è: \emph{come faccio a
ricreare le stesse condizioni?}

Il relatore ha lavorato anche in aziende dove il feedback era solo negativo.
Il rischio? Si immagini l'impiegato ideale: vuole sempre migliorare, è sempre
pronto a dire ``magari sbaglio io''. Se nessuno gli dice cosa sta facendo
\emph{bene}, spinto dal desiderio di migliorare rischia di cambiare cose che
erano già buone, e di farle peggio. Il rinforzo positivo \textbf{non è solo una
gentilezza}: è un atto pratico, che evita che le persone peggiorino ciò che
sanno già fare bene.

\subsection{Il feedback accidentale}

Esiste infine il \textbf{feedback accidentale}, molto frequente negli uffici:
\begin{itemize}
  \item il ``chiedi a lei, non a quel grafico'' del giro dell'ufficio è un
        feedback negativo su una persona, dato accidentalmente, né richiesto
        né voluto;
  \item a Riot li chiamavano \textbf{drive-by}, come nelle scene dei film in
        cui passa un'auto, si abbassa il finestrino, si spara e si riparte. In
        molte aziende con una ``cultura del feedback'' in cui tutti devono
        dire la loro, si riceve una valanga di commenti del tipo ``comunque il
        sito a me non piace, ciao''. Non si sa perché, non si sa cosa fare di
        diverso, non si sa che rischi ci siano: \emph{a che cosa è servita
        questa conversazione?} A niente.
\end{itemize}

% ---------------------------------------------------------------------
\section{Quando, quanto e con chi}
\label{sec:z1-quando}

\subsection{Frequentemente e recentemente}

\emph{Quando avete ricevuto il vostro ultimo feedback?} Per gli studenti,
qualche giorno fa (nel corso) o qualche mese fa (al lavoro, dopo una
consegna). Come aveva osservato Silvia, il feedback andrebbe dato
spesso.

Il relatore usa l'esempio di una nave da crociera (``la Costa Concordia'') che
deve andare da A a B. La rotta ideale (azzurra) è quasi impossibile: nessun
progetto va esattamente come previsto. C'è la rotta ``panoramica'' (verde), che
gira intorno all'isola d'Elba; ci sono poi rotte (arancione e rosa) che a un
certo punto si allontanano molto e poi correggono bruscamente: qualcuno si è
accorto che si stava sbagliando, ``ma non è quello che avevo chiesto'',
``non è dove dovevamo andare''. Sono i famosi \textbf{sanity check} o
\emph{checkpoint}.

L'ideale è avere \textbf{tanti piccoli check per fare tante micro-correzioni}:
correggere la rotta di un paio di gradi è molto più rapido che farlo dopo
essersi allontanati. Bastano trenta secondi di ``ehi, vado bene?''. Spesso non
lo si fa per paura di sembrare insicuri o poco autonomi, ma per il relatore è
l'opposto: quando qualcuno gli chiede ``vado bene così?'' e lui risponde
``sì'', si alleggerisce anche del proprio carico mentale. Un capo che assegna
un compito e lo vede solo alla fine vive sempre col dubbio (``avranno capito?
Lo staranno facendo giusto? Arriverà in tempo?''); con un check di dieci
secondi al giorno (``tutto ok, Stefano'') è tranquillo, e la persona va avanti
sapendo che quanto fatto finora va bene. Non serve fare ogni giorno il giro del
dipartimento: anche ogni uno o due giorni basta a evitare la rotta verde.

Per il feedback è lo stesso: \textbf{va dato subito}. Dire oggi ``ieri hai
fatto questa cosa'' è molto diverso da dire ``negli ultimi sei mesi hai fatto
questa cosa''. Il relatore si è sentito dire più volte, alla valutazione di
dicembre, ``da maggio hai fatto questa cosa che non andava''. La sua risposta:
``e aspetti dicembre per dirmelo? Mi hai lasciato sbagliare per sei mesi?''. È
una mancanza di rispetto, e la conversazione non può più essere costruttiva.
Per questo non ama le valutazioni annuali: i manager crescono abituandosi a
prendere appunti delle cose negative per tirarle fuori alla valutazione
(``a marzo hai fatto questa cosa'' ``sì, però poi ho smesso'' ``sì, però a
marzo l'hai fatta'').

In sintesi:
\begin{description}
  \item[Frequentemente] riguarda l'intervallo tra un check e l'altro.
  \item[Recentemente] significa non lasciar passare un anno prima di dare il
        feedback, altrimenti non serve più a niente.
\end{description}

\subsection{Quanto dura un feedback}

Un feedback dura \textbf{da trenta secondi a quindici minuti}; raramente di
più. Chi non ne ha mai dato uno dovrebbe sforzarsi di darne almeno uno, magari
positivo per cominciare. Ogni feedback dato o ricevuto è un
\textbf{deposito nel proprio ``fondo fiduciario''} (\emph{trust fund}), che
costruisce la reputazione futura. Il relatore si è costruito la reputazione di
chi dà feedback il più possibile oggettivi e analitici: giusti o sbagliati che
siano, costringono a pensare. Non impone mai il suo punto di vista (``si fa
come dico io''), ma chiede ``tu come faresti? Come posso aiutarti?''. Il
feedback, per lui, serve a \emph{far crescere qualcun altro}.

\subsection{I manager peggiori: ``a che numero sto pensando?''}

I peggiori manager che il relatore abbia avuto erano quelli che volevano che si
facesse come volevano loro, anche quando chiedevano ``tu come la faresti?''.
In realtà la domanda era ``\emph{a che numero sto pensando?}'': se dicevi 5,
``no, dovremmo fare 9''; se dicevi 9, ``no, 5''. Lo stress costante di dover
indovinare cosa vuole il capo nasce dalla mancanza di fiducia (``se mi chiedi
come fare, fammelo fare'').

Lo stesso gli succede con i CEO che lo chiamano per ``sistemare la cultura
aziendale'' e poi dicono ``non sono sicuro, io farei diversamente''. La sua
risposta: ``lo so che faresti diversamente, per questo mi hai chiamato''. Se si
chiama un esperto, bisogna lasciargli fare il suo lavoro: non si va dal
meccanico a cambiare le gomme per poi dirgli ``non dovresti girare un po' di
più quel bullone?''. Chiedere un parere a un esperto significa dargli anche il
potere di scegliere. Il problema è il \emph{micromanagement}, o in generale il
non accettare punti di vista diversi dal proprio (magari perché si è il
fondatore e si pensa di essere l'unico a saper fare le cose).

\subsection{Dare feedback a un superiore egocentrico}

\paragraph{Come si dà un feedback negativo a un superiore con un grande ego?}
Il trucco è \textbf{non puntare sull'ego, ma sull'azione e sui suoi effetti}.
Non ``tu fai così e io mi sento così'': anzi, si alimenta il suo ego
rendendolo il \emph{main character} dell'avventura, l'eroe che risolve il
problema (``tu, grande manager, aiutaci a risolvere questo problema che
affligge noi che facciamo l'operatività''). Non è l'ideale (``io cambierei
manager''), ma pochissimi possono scegliersi i colleghi: il relatore stesso ha
quasi sempre ricevuto team preassemblati (``ecco le chiavi, questa è la tua
macchina''). Bisogna giocare con le carte che si hanno in mano.

\paragraph{Un esempio pratico.} Con un manager che fa micromanagement e mette
sempre ``il becco'' nel proprio lavoro, si può dire: ``se ti fidi di me per
questo lavoro, con le tue capacità potresti dedicarti a questa cosa più
strategica e di lungo termine: pensa quanto sarebbe più utile per l'azienda''.
In sostanza: ``guarda cosa potresti raggiungere se mi lasciassi fare il mio
lavoro''. Non come un ``tu fai il tuo e io faccio il mio'', ma come un
``ti sollevo da questa parte del carico, che so che ritieni importante: lascia
che ci pensi io, così puoi dedicarti ad altro'', che magari gli dà anche più
visibilità.

Si gioca su ciò che quella persona pensa di saper fare, e che probabilmente sa
anche fare. Anche Michael Scott di \emph{The Office}, per quanto preso in giro,
è raccontato come un venditore bravissimo, che ha chiuso i contratti più
importanti per l'azienda, e rispetto ai manager che gli succedono è forse
``uno dei meno peggiori''. Tutti, anche i manager peggiori, hanno punti di
forza: \textbf{giocate su quelli}.

% ---------------------------------------------------------------------
\section{Il feedback come regalo}
\label{sec:z1-regalo}

Per chiudere, il relatore risponde alla domanda d'apertura, che aveva lasciato
apposta in sospeso: che cos'è il feedback per lui?

\begin{quote}
  \emph{Il feedback è un regalo che mi viene fatto da qualcuno che so che vuole
  il mio interesse.}
\end{quote}

Ha costruito un team di persone che si fidano di lui e che non hanno motivo di
sabotarlo o di ferirlo: per questo interpreta anche le cose peggiori con il
default to trust. È un regalo, e cerca di prenderlo in considerazione il più
possibile.

Ci è arrivato partendo da se stesso. Oltre a essere avverso ai conflitti, si
definisce ``avverso alla comunicazione umana'': paradossalmente, una persona
introversa e timida finita a fare comunicazione, marketing e ora HR (``succede
spesso''). Prima di ordinare una pizza al telefono ha dieci minuti di
preparazione mentale, e se l'ordine delle domande cambia va in confusione. Per
lui un feedback di dieci minuti sono spesso \textbf{tre ore di preparazione}:
costruisce un \textbf{albero dei dialoghi}, come in un'avventura grafica
personale, con tutte le opzioni. Poi spesso finisce con ``hai fatto questa
cosa, andrebbe migliorata'' ``hai ragione, grazie'', e basta.

Così ha pensato: chissà quanti si fanno gli stessi problemi con lui, e passano
ore a prepararsi per dirgli ``questa cosa l'hai fatta male'', sapendo che
riceveranno comunque quell'occhiata fugace del ``come ti permetti''. Eppure
vengono a dirglielo. Non può far altro che rispondere: ``grazie, hai ragione,
cercherò di fare meglio''. È così che ha metabolizzato il ricevere feedback:
capendo che non c'è nessuna cattiva intenzione e che, se non proprio felice
(``ci ho messo una pietra sopra''), può almeno essere \emph{grato} che ci sia
ancora qualcuno che ha così a cuore il suo interesse da dargli feedback.

\paragraph{Il silenzio fa più paura.} Non bisogna spaventarsi dei feedback
negativi che emergono dai sondaggi aziendali: bisogna spaventarsi del
\textbf{silenzio}, quando nessuno dice più nulla. Significa che nessuno si
fida più che si sappiano risolvere i problemi.

\paragraph{La sindrome dell'impostore.} Il suo motivatore personale, all'inizio,
era soprattutto ``non voglio fallire'': per questo cercava di migliorare, e
migliorare è diventato accettare feedback, perché più feedback riceve, più
migliora. Il feedback è un \textbf{acceleratore}, e lui cerca di incentivare
gli altri a darglielo. Resta però sempre la vocina del ``e se si accorgono che
non sei capace?''. Con la \emph{impostor syndrome} convive da anni e
continuerà a conviverci; se ce l'avete, conviveteci anche voi. Avrete sempre
la sensazione che ``se ne accorgeranno'': forse sì, ma sapete fare anche
altre cose. Non concentratevi solo su ciò che non sapete fare. Gli errori
capitano: la perfezione non è di questo mondo.

\paragraph{Esercitarsi in accademia.} Per fare pratica con il feedback,
l'ambiente migliore è quello dei colleghi di corso: persone nella stessa
situazione, con esperienza simile e lo stesso obiettivo, in un contesto in cui
non si può essere licenziati e l'unico effetto a lungo termine è ``ho parlato
con questa persona, ora le cose vanno meglio''. E bisogna avere il coraggio di
parlare anche con i docenti (``dire a Diego: secondo me questa lezione potevi
farla in questo modo''; ``so che è una persona che lo accetta''), senza farsi
bloccare dalla gerarchia. Nell'azienda giusta (come quella in cui lavora ora) il
relatore va sempre dal CEO a dare feedback, sapendo che non verrà licenziato:
il CEO gli ha detto ``ti ho preso per darmi feedback''.

La frase di chiusura:

\begin{quote}
  \emph{Per superare le vostre difficoltà, tutto quello che vi viene chiesto è
  il coraggio di chiedere aiuto.}
\end{quote}

Chiedere aiuto non è una debolezza: è un segno di consapevolezza di sé
(\emph{self-awareness}) e di forza personale. Se tutti lo facessero, i
progetti andrebbero molto meglio; spesso invece ci si fa frenare dal ``non
voglio fare brutta figura'', ``non voglio passare per quello scemo'', e sono
proprio queste paure a rallentarci e a impedirci di raggiungere gli obiettivi.

% ---------------------------------------------------------------------
\section{Domande e risposte finali}
\label{sec:z1-qa-finale}

L'ultima mezz'ora è dedicata a domande libere.

\paragraph{E se una persona peggiora per uno o due anni di fila?} Se nessuno
gliel'ha mai detto, c'è un grosso problema, non suo ma dell'azienda. Se invece
c'è consapevolezza, il percorso è: ``cosa vuoi fare per risolvere?''; ``ho
provato questo, questo e questo''; ``visto che non ha funzionato, vuoi un mio
consiglio? Proviamo con quello?''; ``che supporto ti serve?''. Se, con tutto il
supporto, il tempo e le soluzioni concordate per mesi, il miglioramento non
arriva, per il relatore la soluzione è: \textbf{``dimmi un prezzo per andare
via''}. Preferisce pagare una buonuscita (\emph{severance package}) che
avviare la strada dei richiami formali sulla performance fino al
licenziamento per giusta causa. A un certo punto andare avanti non è un favore
per nessuno.

\paragraph{È vero che nei videogiochi si viene sfruttati? Meglio un piccolo
studio italiano?} Per il relatore si è quasi più a rischio a Blizzard che in un
piccolo studio di Roma o Torino. Nel piccolo studio si sa dove ci si va a
cacciare: non ci sarà il team building a Santo Domingo, ma ci si concentra su
un gioco alla volta e, se va bene, si fa il successivo. A Blizzard si può
lavorare su \emph{Overwatch}, che vende milioni, e perdere comunque il posto,
in modo totalmente scollegato dal proprio lavoro. Gli stipendi più alti
all'estero compensano anche questo rischio, che spesso non dipende da come va
il gioco (anzi, ``il gioco non ha fatto come volevamo'' è spesso una scusa per
licenziamenti che non avevano motivo).

\paragraph{Cambiare settore a trent'anni, da web developer full remote?} Lo
stipendio indicato dallo studente non è fuori dalla portata di vari studi
italiani, anzi è superabile. Non va esclusa la possibilità di lavorare da
remoto per aziende piccole ma estere. Alcuni amici del relatore, pagati con
stipendi italiani, si sono trasferiti a Bali o in paesi con un costo della vita
più basso; il collega di \emph{PUBG} percepisce uno stipendio olandese vivendo
a Milano, che per quanto cara non costa come i Paesi Bassi.

\paragraph{Come si diventa manager?} Dipende dall'azienda: alcune promuovono i
migliori (processo tendenzialmente sbagliato), altre propongono un percorso di
\emph{people management}. Diventare manager significa \textbf{abbandonare ciò
che si sa fare per imparare a fare altro}. Il relatore ha dovuto accettare che
non sarebbe più stato ``il miglior community manager italiano del mondo''
(``eravamo in pochi, quindi era facile''): dedicandosi alle persone non poteva
più fare il suo lavoro bene come prima. È stata una negoziazione con se
stesso: fare il manager significa non essere più l'esperto di qualcosa, ma
saper gestire e far crescere persone. Significa anche \textbf{cedere il
controllo della propria carriera}: se il team decidesse di non fare più nulla,
il manager ne pagherebbe le conseguenze. Il suo lavoro diventa scegliere le
persone giuste e farle lavorare bene. Lui lo esplicita sempre: sa di affidare
al team parte della propria carriera, e il team sa che lui combatterà per
loro, anche prendendosi le colpe.

\paragraph{L'inglese.} \textbf{Va assolutamente imparato}, anche lavorando in
Italia. Non saperlo preclude circa il 95\% del mercato (tutto quello fuori
dall'Italia) e soprattutto il \emph{networking} con persone che potrebbero dare
occasioni future: l'inglese è spesso l'unica lingua in comune con chi lavora
all'estero ed è oggi la lingua universale dei videogiochi.
\begin{itemize}
  \item \emph{Che livello serve?} Non importa B2 o C2: bisogna saper sostenere
        una conversazione, guardare almeno Netflix con sottotitoli in inglese,
        comunicare senza dodici pause e senza ``come si dice\dots''. Il
        certificato non verrà mai chiesto: verrà chiesto se si sa parlare
        inglese. Il relatore ha sentito persone con un C2 che non sapevano
        parlarlo.
  \item \emph{Come allenarlo senza occasioni di conversazione?} Netflix, Prime
        Video, Now TV, Udemy, Coursera, YouTube, Duolingo, siti di scambio
        linguistico con madrelingua, Discord: non sono mai esistiti così tanti
        mezzi gratuiti. Il relatore ha imparato l'inglese giocando (con una
        pronuncia sbagliata, perché leggeva senza sentire, su cui ha poi
        dovuto lavorare); guardare un film in inglese con i sottotitoli in
        italiano insegna anche la traduzione.
  \item Oltre all'ascolto passivo è fondamentale \textbf{parlare}: in una
        flotta di \emph{EVE Online} con 150 persone quasi tutte madrelingua
        inglese si è esposti a decine di accenti e modi di dire (mille modi di
        dire ``come stai'' oltre a \emph{how are you}); lavorando all'estero si
        impara a distinguere gli accenti (Manchester, Londra) e ad apprezzarli
        anche nei media.
\end{itemize}

\paragraph{Quanto deve sapere un junior?} Dipende da quanto è tecnico il
ruolo. Un game designer junior che non sa qualcosa può essere affiancato da un
senior; un programmatore che non sa programmare, in quel momento, non è utile
all'azienda.

\paragraph{Annunci ``junior con cinque anni di esperienza''.} ``Cercasi
ventenne con trent'anni di esperienza.'' Non è l'annuncio per un junior: è un
errore dell'azienda, che in realtà cerca un \emph{mid}. Riceverà profili
junior, si stupirà che non abbiano cinque anni di esperienza, rifarà l'annuncio
e avrà perso sei mesi. Il consiglio: tenere d'occhio quella pagina, perché
dopo qualche mese l'annuncio verrà corretto.

\paragraph{Una struttura per il post-mortem.} A chi chiede un formato di base,
il relatore anticipa quanto vedrà nella seconda giornata a proposito di
brainstorming: una lavagna \textbf{divisa in quattro aree}, su cui ognuno
attacca tutti i post-it che vuole:
\begin{center}
  \begin{tabular}{@{}p{0.44\textwidth}p{0.44\textwidth}@{}}
    \toprule
    \textbf{Smettere di fare} & \textbf{Iniziare a fare} \\
    ciò che non ha funzionato & ciò che non abbiamo fatto e andava fatto \\
    \midrule
    \textbf{Fare di meno} & \textbf{Fare di più} \\
    ciò che ha fatto perdere tempo rispetto all'impatto ottenuto &
    ciò che ha funzionato ma, fatto con troppa leggerezza, non ha dato i
    frutti sperati \\
    \bottomrule
  \end{tabular}
\end{center}
Con i post-it ognuno ha voce e tutto risulta catalogato; per il report basta
una foto della lavagna (o uno screenshot, se la lavagna è digitale, come Miro).

\paragraph{Come sbloccare la paralisi del feedback verso il management?} Uno
studente ha visto manager dire ``non abbiate paura, la mia porta è sempre
aperta'' senza che cambiasse nulla. C'è una dissonanza tra ciò che dicono e
ciò che fanno: è ``la mamma che dice vieni qua che non ti faccio niente''.
Qualcosa deve aver fatto pensare alle persone che non è sicuro. Prima va
sistemato quello; poi si può provare con i \textbf{feedback anonimi}.
L'obiettivo del relatore resta il feedback aperto e firmato, ma a Stormind ha
iniziato con feedback anonimi verso il management per verificare se
l'anonimato fosse un fattore. Non lo era (evidentemente non c'erano problemi a
parlare con il CEO), e la volta successiva intende abbandonarli. Ma è un buon
\emph{check}: se senza anonimato non arriva nulla e con l'anonimato arrivano
cento feedback, il problema è chiaro.

\paragraph{Quando è giusto rivolgersi a un superiore?} La strada giusta resta
parlare direttamente con l'interessato. Ma se si prova e non si trova una
soluzione, la cosa più ovvia è andare dal proprio manager insieme: ``ci
aiuti?''. Quando il relatore parlava di mediare i dialoghi, era proprio questo:
fare da moderatore, aiutare a sbloccare i punti fermi o, se non si riesce a
scegliere tra A e B, scegliere lui. L'escalation, quando necessaria, va fatta.

\paragraph{Il tempo libero veniva rispettato?} No, e se l'è dovuto imporre
lui. Il lavoro gli piaceva così tanto che continuare a casa non gli sembrava
lavorare. Gli si è ritorto contro con il lancio di \emph{Diablo~III}, fatto in
orario americano (le due o l'una di notte in Europa), con l'ufficio chiuso
perché a nessuno era stato chiesto di restare. Il lancio va malissimo, i
giocatori protestano, e lui, sveglio, risponde sui forum fino alle tre, quattro,
cinque del mattino: ``non posso risolvere il problema, ma resto qui ad
aggiornarvi con quello che riesco a sapere dal team di sviluppo''. Il giorno
dopo trova sulla scrivania un \textbf{richiamo formale dell'HR}: aver lavorato
fuori orario senza autorizzazione, condotta non professionale. Per lui è stato
un interruttore: ``ora me ne vado''. Due settimane dopo faceva colloqui con
Riot. Ancora oggi è contrario a quel richiamo, ma gli ha fatto capire come
ragiona un'azienda: ``hai un contratto 9--18, non mi interessa se l'hai fatto
con le migliori intenzioni''. Eppure per anni ha goduto di un'ottima
reputazione in quelle community come ``quello che è rimasto con noi fino alle
cinque del mattino'', e anche l'azienda ne ha beneficiato.

Ci sono poi aziende in cui, soprattutto per i developer, viene richiesto il
famigerato \textbf{crunch}: lavoro straordinario continuativo, senza pause.
Dipende molto dall'azienda; sta venendo (``grazie al cielo'') deprecato, perché
molte aziende capiscono che non è sano ed è visto male anche dal pubblico, ma
esiste ancora.

\paragraph{E se il tuo capo demotiva qualcuno del tuo team?} Qui ``diventa
personale'': il team del relatore si fida di lui perché sa che in questi casi
``diventa una bestia''. ``Te la puoi prendere con me quanto vuoi, ma il mio
team non lo tocchi.'' Lo gestisce direttamente con l'altro manager, per quanto
brutta possa diventare la conversazione. Lo sforzo più grande è \textbf{non
andarne poi a parlare con il team} (``ragazzi, ho fatto questa cosa''), perché
minerebbe la fiducia verso l'alto: non vuole un rapporto ottimo fino a lui e
pessimo oltre. L'obiettivo è che l'altro manager ``si faccia perdonare''
cambiando il proprio comportamento, e che il team percepisca da solo il
cambiamento.

\paragraph{Chiusura della giornata.} Il relatore chiede un feedback collettivo
in chat (``se domani ci dai \emph{more of the same} siamo contenti, oppure
ditemi cosa limare''), si scusa per i ripetuti crash della connessione e
ringrazia per la proattività nelle domande, che gli servono anche come
\emph{check}: se nessuno chiede nulla su un argomento, significa che può
andare avanti. Racconta di aver passato una settimana a decidere che cosa
mettere nelle slide e di averle poi preparate in poco tempo (rifacendone
alcune fino alle tre di notte). Ricorda infine di usare il test di
personalità per sé e per il proprio gruppo, facendolo fare a chi non l'ha
ancora fatto: tornerà molto utile per cooperare meglio.

% ---------------------------------------------------------------------
\section{Anticipazioni sulla seconda giornata}
\label{sec:z1-domani}

La seconda giornata del workshop sarà più pratica: ``oggi vi ho spiegato
perché, domani vi insegno come lo fanno queste aziende''. Gli argomenti
previsti:
\begin{description}
  \item[Time management] come organizzarsi una giornata professionale in
        azienda: come strutturare una \emph{task list}, cosa fare e cosa non
        fare, come e cosa prioritizzare.
  \item[Brainstorming] i gruppi dovranno fare moltissimi meeting (come facciamo
        il gioco, come lo chiamiamo, che scopo gli diamo). Se non si sanno fare
        bene, i brainstorming sono \textbf{dannosi ed eterni}: un brainstorming
        dovrebbe durare \textbf{al massimo due ore}. Verranno mostrati metodi
        che accorciano la durata, massimizzano l'effetto e, soprattutto,
        \textbf{fanno partecipare tutti}: non c'è niente di peggio di un
        brainstorming in cui parlano in due per il 90\% del tempo. Il classico
        ``ragazzi, non esistono idee sbagliate, sparate'' è il modo peggiore di
        farlo; non è colpa di nessuno, ma si può fare diversamente.
  \item[Guidare da dentro] come essere leader nel proprio team anche senza
        esserne il leader: sviluppare responsabilità personale e verso gli
        altri. Non essere il manager di qualcuno non significa non esserne
        responsabili: la persona accanto è responsabilità di tutti, perché si
        è un team e non un insieme di individui. ``Se qualcuno cade, tutti lo
        facciamo rialzare: non c'è l'addetto al rialzamento.'' È la versione
        \emph{advanced} della giornata: come comportarsi da manager senza
        averne il titolo.
\end{description}
Torneranno inoltre il \emph{service level agreement} e il tema della
condivisione dei successi e della protezione del team dai fallimenti.

% ---------------------------------------------------------------------
\section*{Punti chiave del capitolo}
\addcontentsline{toc}{section}{Punti chiave del capitolo}

\begin{itemize}
  \item ``La capacità di un gruppo di persone di fare cose straordinarie
        dipende dal modo in cui queste persone operano insieme.'' Il 90\% dei
        problemi dei team sono \textbf{problemi di comunicazione}.
  \item Le aziende si liberano sempre più presto dello \textbf{useless
        genius}: le competenze si trasmettono, il carattere no. Il
        \textbf{fit culturale} va valutato anche dal candidato.
  \item Non esiste un unico modo di dare feedback: dipende dal contesto e
        soprattutto da \textbf{chi lo riceve}. Il feedback è una somma di
        piccoli miglioramenti periodici; tenersi le cose nel cassetto non
        serve mai.
  \item I team attraversano \textbf{forming, storming, norming, performing e
        adjourning}. Lo storming non si salta; si passa la maggior parte del
        tempo in forming; la bravura sta nell'accelerare i passaggi e nel
        fare bene il \textbf{post-mortem}, anche dei successi.
  \item Gli obiettivi di un progetto vanno fissati \textbf{all'inizio}: solo
        così si può valutare se è andato bene e scegliere tra A e B.
  \item Il circolo virtuoso è \textbf{feedback $\rightarrow$ miglioramento
        $\rightarrow$ fiducia}. Accettare un feedback non significa fare come
        dice l'altro: decide chi ha la responsabilità, ma deve argomentare.
  \item ``Aiutami a poterti aiutare'': meglio dire in anticipo ``potrei non
        farcela'' che ``non ce l'ho fatta''. Chiedere aiuto non è una
        debolezza.
  \item Quando si dà feedback \textbf{non si porta la soluzione}: la si
        costruisce insieme a chi lo riceve.
  \item La cultura è forte quando si lavora \textbf{per gli altri}; il lavoro
        è valutato come opera d'insieme (anche con bonus di team).
  \item Come in \emph{Among Us}, i team sprecano energie a cercare il
        traditore invece di fare le task. \textbf{Default to trust}: chi ha
        sbagliato ha agito con le informazioni che aveva. Cercare
        \textbf{la causa, non il colpevole}: ``che cosa ti ha fatto pensare
        che fosse la scelta giusta?''.
  \item Carriere e performance non sono lineari: vanno valutate le
        \textbf{traiettorie}, non un punto dell'anno. Le valutazioni puntuali
        generano \emph{lying, hiding, faking}.
  \item Confermare sempre le istruzioni facendosele ripetere; ``non ho
        capito'' non è una colpa.
  \item Dare feedback: sull'\textbf{azione} e sui suoi \textbf{effetti}, specifico,
        presto, uno a uno, adattato al carattere, senza terze parti, senza
        ``non mi piace'', senza offendersi se non viene seguito; poi fare il
        debriefing.
  \item Ricevere feedback: \textbf{ascolta, chiedi dettagli, accetta, agisci}.
        Il feedback positivo serve a non peggiorare ciò che si fa bene; il
        feedback accidentale e il \emph{drive-by} non servono a nulla.
  \item Feedback \textbf{frequenti e recenti}: tanti piccoli \emph{sanity
        check} costano meno di una grande correzione di rotta.
  \item Il feedback è \textbf{un regalo} di chi ha a cuore il nostro
        interesse; fa paura il silenzio, non le critiche.
\end{itemize}

%    % =====================================================================
%  Lezione zero 2 — Gestire il tempo, generare idee, guidare da dentro
%  Fonte: trascrizione della seconda "lezione zero" (GD0_2), seconda e
%  ultima giornata del workshop di due giorni tenuto prima dell'inizio del
%  corso di Game Design, dedicato alla figura professionale e al lavoro in
%  team. Prosegue direttamente il capitolo 0.1 (cultura del feedback).
%  Pacchetti richiesti nel preambolo del libro:
%    \usepackage[italian]{babel}, \usepackage[T1]{fontenc},
%    \usepackage[utf8]{inputenc} (solo pdfLaTeX), \usepackage{booktabs}
%
%  Collocazione nel libro: parte "Lezioni zero", subito dopo
%  GD0_1_cultura_del_feedback, con la stessa numerazione propria (0.2).
%  Nel file principale del libro, ad esempio:
%
%    \part{Lezioni zero: il game designer nell'industria}
%    \renewcommand{\thechapter}{0.\arabic{chapter}}
%    \setcounter{chapter}{0}
%    % =====================================================================
%  Lezione zero 1 — Lavorare in squadra: la cultura del feedback
%  Fonte: trascrizione della prima "lezione zero" (GD0), prima giornata di
%  un workshop di due giorni tenuto prima dell'inizio del corso di Game
%  Design, dedicato alla figura professionale e al lavoro in team.
%  Pacchetti richiesti nel preambolo del libro:
%    \usepackage[italian]{babel}, \usepackage[T1]{fontenc},
%    \usepackage[utf8]{inputenc} (solo pdfLaTeX), \usepackage{booktabs}
%
%  Collocazione consigliata nel libro: le lezioni zero precedono il corso,
%  quindi vanno in una parte separata PRIMA di GD1, con numerazione propria
%  (0.1, 0.2, ...) così che i capitoli del corso conservino i numeri 1, 2, 3.
%  Nel file principale del libro, ad esempio:
%
%    \part{Lezioni zero: il game designer nell'industria}
%    \renewcommand{\thechapter}{0.\arabic{chapter}}
%    \setcounter{chapter}{0}
%    \input{Seminari/GD0_1_cultura_del_feedback}
%    % ... altre lezioni zero ...
%    \part{Il corso di Game Design}
%    \renewcommand{\thechapter}{\arabic{chapter}}
%    \setcounter{chapter}{0}
%    \input{GD1/GD1_introduzione_al_game_design}
%    % ...
%
%  (Con hyperref aggiungere anche \renewcommand{\theHchapter}{z.\arabic{chapter}}
%  nella parte zero, per evitare ancore duplicate.)
%  Tutte le etichette di questo capitolo hanno il prefisso "z1:" per non
%  collidere con quelle dei capitoli del corso.
% =====================================================================

\chapter[La cultura del feedback]{Lavorare in squadra: la cultura del feedback}
\label{cap:z1-feedback}

Questo capitolo apre le \emph{lezioni zero}: incontri che precedono l'inizio
vero e proprio del corso di Game Design e che non riguardano una tecnica di
design in senso stretto, ma \textbf{la figura professionale} di chi lavora
nell'industria dei videogiochi. Prima di imparare a progettare un gioco
conviene infatti capire in che tipo di ambiente lo si progetterà: un ambiente
fatto di team, di gerarchie, di scadenze e, soprattutto, di persone che devono
riuscire a lavorare insieme.

La lezione è la prima giornata di un \textbf{workshop di due giorni} (circa
dieci ore in totale, pause pranzo comprese) tenuto da Stefano
Cozzi\footnote{Il cognome è riportato come compare nella trascrizione
automatica, che in più punti storpia i nomi propri; andrebbe verificato.
Nella stessa trascrizione lo studio per cui lavora appare come
``Stormine'', ``Stormline'' e ``Stormind'': si tratta di Stormind Games,
studio con sede a Catania, coerentemente con quanto detto in lezione.}, un
professionista con circa quindici anni di esperienza nel settore, tra
Blizzard Entertainment, Riot Games e, oggi, Stormind Games. Il workshop è
nato da un confronto con i coordinatori dell'accademia\footnote{Nella
trascrizione l'accademia è resa in modi diversi (``Hive'', ``a dive'',
``D.I.V.''); i coordinatori citati sono Enrico e Andrea, insieme al
direttore Diego Ricchiuti. Nel resto del capitolo si usa semplicemente
``l'accademia'', come negli altri capitoli.}: serviva un momento dedicato a
\emph{come ci si comporta in un ambiente professionale}, cosa che riguarda sia
il lavoro futuro sia i progetti di gruppo svolti durante il percorso di
studi, dove ci si deve già abituare a comportarsi ``come se si fosse in
un'azienda''.

La prima giornata ha un unico grande tema, sezionato in più parti:
\textbf{la cultura del feedback}. La seconda giornata (che sarà oggetto della
prossima lezione zero) sarà più pratica: \emph{time management},
\emph{brainstorming} e il ``guidare da dentro'' un team anche senza esserne
il leader (sezione~\ref{sec:z1-domani}).

Il percorso del capitolo segue quello della lezione:
\begin{enumerate}
  \item il formato del workshop e il percorso professionale del relatore
        (sezione~\ref{sec:z1-relatore});
  \item perché le aziende tengono tanto al lavoro di squadra e al
        \emph{fit} culturale (sezione~\ref{sec:z1-squadra});
  \item che cos'è il feedback e da che cosa dipende il modo di darlo
        (sezione~\ref{sec:z1-cose-feedback});
  \item le cinque fasi di vita di un team e il valore del
        \emph{post-mortem} (sezioni~\ref{sec:z1-fasi-team}
        e~\ref{sec:z1-qa-fasi});
  \item il confronto con i colleghi e alcune domande sul lavoro nei grandi
        studi (sezioni~\ref{sec:z1-rivalita} e~\ref{sec:z1-industria});
  \item il circolo virtuoso tra feedback, miglioramento e fiducia, e i suoi
        effetti sulla \emph{retention} e sulla cultura aziendale
        (sezioni~\ref{sec:z1-triangolo} e~\ref{sec:z1-cultura});
  \item il ``gioco del traditore'' e il principio del
        \textbf{default to trust} (sezioni~\ref{sec:z1-among-us}
        e~\ref{sec:z1-default-trust});
  \item le traiettorie di carriera e la valutazione delle performance
        (sezione~\ref{sec:z1-traiettorie});
  \item come si dà, come si riceve e quando si dà un feedback
        (sezioni~\ref{sec:z1-dare}, \ref{sec:z1-ricevere}
        e~\ref{sec:z1-quando});
  \item il feedback come regalo, le domande finali e le anticipazioni
        della seconda giornata (sezioni~\ref{sec:z1-regalo},
        \ref{sec:z1-qa-finale} e~\ref{sec:z1-domani}).
\end{enumerate}

Una premessa di metodo: il relatore ha alternato continuamente teoria,
domande alla classe e racconti della propria esperienza. I racconti non sono
digressioni: sono gli esempi concreti con cui i principi vengono resi
tangibili, e per questo sono riportati con cura.

% ---------------------------------------------------------------------
\section{Il workshop e il relatore}
\label{sec:z1-relatore}

\subsection{Formato e regole di partecipazione}

Il relatore ha cercato di ``riassumere quindici anni di esperienza in dieci
ore''. Ha materiale per circa il doppio delle ore disponibili: le cose
strettamente necessarie vengono trattate per prime, il resto è tenuto ``in
riserva'' come \emph{nice to have}. Le regole del workshop:
\begin{itemize}
  \item le domande si fanno \textbf{alla fine di ogni slide}, così da chiudere
        prima il discorso; non c'è un'unica sessione finale di domande e
        risposte (anche se ci sarà), perché il relatore preferisce fermarsi
        spesso e non fare ``uno spiegone di dieci ore'';
  \item le domande non devono per forza riguardare la slide appena vista: se
        una slide dà lo spunto per altro, se ne può parlare; al massimo sarà
        il relatore a rispondere in breve o a rimandare a una slide
        successiva;
  \item si può chiedere anche di esperienze passate, opinioni e sensazioni;
        se per NDA o per serietà professionale non può rispondere, lo dirà
        esplicitamente (``non lasciate che sia un \emph{blocker}'');
  \item su Zoom si usa l'alzata di mano o la chat, che il relatore tiene
        aperta sul secondo monitor;
  \item non ci sono pause ufficiali: la lezione è registrata e ognuno è
        libero di assentarsi quando serve, cercando però di non disturbare il
        flusso della lezione;
  \item le slide verranno fornite integralmente tramite la segreteria
        dell'accademia (Vanessa), insieme alle registrazioni: le slide senza
        il parlato servono poco, ma sono un buon promemoria.
\end{itemize}

\subsection{Il percorso professionale}

Il relatore si presenta (scherzosamente) come la ``guida escursionistica''
che accompagna la classe nella ``foresta incantata'' della cultura del
feedback, perché per molti è un territorio sconosciuto. Il suo percorso:
\begin{itemize}
  \item in Italia ha lavorato per aziende che fornivano servizi a terzi, tra
        cui il lavoro su \emph{Warhammer Online}, e per un canale televisivo
        sul digitale terrestre dedicato ai videogiochi\footnote{Nella
        trascrizione: ``Digital Game Network Channel''. Il nome esatto del
        canale andrebbe verificato.};
  \item a \textbf{Blizzard Entertainment} (circa sei anni, a Parigi) ha fatto
        il \emph{community manager} per l'Italia (\emph{StarCraft~II},
        \emph{World of Warcraft} da \emph{Mists of Pandaria} in poi,
        \emph{Hearthstone}, \emph{Diablo~III}, le espansioni di
        \emph{StarCraft}, il lancio di Battle.net), per poi diventare il
        manager dei community manager europei;
  \item a \textbf{Riot Games} (circa cinque anni, a Dublino) è stato
        brevemente community manager italiano e poi \emph{publishing
        manager} di un team di circa trenta persone per l'Europa. Il
        \emph{publishing} non è solo community: è l'unione di
        \textbf{community, marketing e PR};
  \item a \textbf{PG Esports} si è occupato di eventi dal vivo e online
        (ad esempio la Milano Games Week e tornei di \emph{Valorant} e
        \emph{League of Legends}), restando in contatto con Riot e Blizzard;
  \item da circa un anno è tornato nello sviluppo, a \textbf{Stormind
        Games}, dove oggi si occupa di risorse umane e organizzazione del
        lavoro; lo studio verrà presentato più ampiamente al \emph{Career
        Day} dell'accademia insieme alla sua HR manager;
  \item fa inoltre consulenze aziendali anche al di fuori dei videogiochi e
        insegna in master universitari (IULM di Milano).
\end{itemize}

Il suo ruolo ricorrente, sia a Blizzard sia a Riot, è stato quello di
\emph{Mr.~Wolf} (``per chi ha visto \emph{Pulp Fiction}''): quello che
risolve i problemi. Veniva inserito nei team più in difficoltà, ad esempio
dal punto di vista delle performance, per sistemarli e passare al successivo.
Era presente in alcuni dei momenti più critici, come il lancio di
\emph{Diablo~III} o il lancio di espansioni di \emph{World of
Warcraft}\footnote{La trascrizione riporta ``il lancio con i Gary's'' e
``World of Trenor'': si tratta verosimilmente di \emph{Mists of Pandaria} e
di \emph{Warlords of Draenor}.}, a fare \emph{damage control}.

La scoperta fondamentale di questi anni è che i problemi da risolvere
\textbf{quasi mai riguardavano il prodotto} (non era certo lui a sistemare i
server di \emph{Diablo~III}), ma i team che c'erano dietro: persone che
litigavano, non collaboravano, non riuscivano a stare nella stessa stanza
senza urlarsi contro. E nel 90\% dei casi tutto questo si riduce a
\textbf{problemi di comunicazione}: è raro che perfino le incompatibilità
caratteriali non siano risolvibili parlando.

\subsection{Struttura non vuol dire burocrazia}

A Stormind il relatore sta accompagnando un team passato da 30 a 95 persone
in meno di un anno, attraverso tutti i cambiamenti di cui parla il workshop.
Darsi una struttura viene spesso percepito male, come ``più processi, più
procedure'', ma bisogna distinguere:
\begin{itemize}
  \item un processo come ``timbrare il cartellino'' può essere davvero
        invasivo e noioso;
  \item definire \emph{come funziona un documento di game design, chi lo
        approva, chi lo visiona e chi no} è invece una struttura necessaria.
\end{itemize}
Uno studio di 30 persone nella stessa sede può farne a meno: quando c'è un
dubbio ci si alza e si va alla scrivania del collega. Uno studio di 95
persone sparse per l'Italia (e in parte per l'Europa) non può più farlo:
servono una struttura diversa, un modo di lavorare diverso e una
pianificazione del futuro diversa.

% ---------------------------------------------------------------------
\section{Perché il lavoro è sempre lavoro di squadra}
\label{sec:z1-squadra}

\subsection{Non si va in ufficio per fare amicizia}

Il primo messaggio è netto: \textbf{non si va in ufficio per fare amici} o
per trovare compagni di vita. Si va in ufficio perché c'è un lavoro da fare, e
lo scopo è riuscire a farlo \emph{anche con chi sta antipatico}. Con le
persone simpatiche è facile: diventeranno amici, e non sarà mai un problema.
È con chi sta antipatico che le cose si complicano, ed è lì che emergono i
\textbf{veri professionisti}: quelli che riescono a lavorare con chiunque
senza che l'antipatia pesi sul modo in cui lavorano e collaborano, mettendo il
prodotto e il team davanti al proprio ego e alle proprie preferenze personali.

La frase che fa da filo conduttore a tutto il workshop (e che il relatore
promette di riproporre ``fino alla nausea'') è:

\begin{quote}
  \emph{La capacità di un gruppo di persone di fare cose straordinarie
  dipende dal modo in cui queste persone operano insieme.}
\end{quote}

Le aziende raramente pensano che l'innesto di una singola persona possa
cambiare il destino di un progetto. Questo è tanto più vero quanto più grande
è il team, ma vale anche per i piccoli studi, anche per un gruppo di cinque
persone. Come negli sport di squadra: l'individuo conta fino a un certo punto,
e anche se si ha Dončić in squadra, se si gioca uno contro cinque si perde
comunque\footnote{Il riferimento è alle finali NBA in corso nei giorni del
workshop tra Dallas Mavericks e Boston Celtics (giugno 2024), il che
permette di datare approssimativamente la lezione.}.

\subsection{Lo \emph{useless genius} e il costo delle persone tossiche}

Molti dipartimenti HR usano la definizione di \textbf{useless
genius}\footnote{Il concetto è più noto in letteratura manageriale come
\emph{brilliant jerk}; in lezione è usata l'espressione \emph{useless
genius}.}: si può essere bravissimi, a volte i migliori nel proprio lavoro,
ma se non ci si sa comportare, relazionare e operare in un gruppo si è
fondamentalmente inutili ai fini del gruppo. Nel migliore dei casi inutili,
nel peggiore \emph{dannosi}.

In passato le aziende tendevano a ``chiudere un occhio'': la persona è
fondamentale, la teniamo. Oggi la domanda è diventata: \emph{quanto costa
all'azienda, oltre al suo stipendio, tenere una persona che crea un ambiente
tossico e peggiora chi le sta intorno?} Ci sono costi sommersi che le aziende
cominciano a vedere, e la tendenza è liberarsi di queste persone il prima
possibile. Il ragionamento è semplice:
\begin{itemize}
  \item le \textbf{competenze si possono trasmettere} o sostituire: si trova
        qualcuno che fa lo stesso lavoro magari nel doppio del tempo, oppure
        più persone che se lo dividono;
  \item il \textbf{carattere delle persone} invece cambia molto raramente.
\end{itemize}

\subsection{Il fit culturale}

Quando qualcuno \textbf{non è un \emph{fit} culturale} per un'azienda è
meglio separarsi presto. Ma non essere adatti a \emph{quell'}azienda non
significa non essere adatti a lavorare: è perfettamente legittimo che una
persona inadatta a un'azienda sia adattissima a un'altra.

Il relatore l'ha vissuto in prima persona. Dopo sei anni a Blizzard viene
contattato da Riot e affronta un processo di selezione lunghissimo (``quasi
undici colloqui''), al termine del quale viene scartato. Scrive allora una
mail che gli ha ``cambiato la traiettoria di carriera'': chiede
semplicemente perché non è stato ritenuto adatto. La risposta: in nessun
colloquio si era esposto abbastanza da far capire \emph{perché} volesse
lasciare Blizzard, e quindi pensavano che in realtà stesse bene dov'era. Il
motivo era che, per educazione professionale, non gli piace parlare male del
proprio posto di lavoro; portata all'eccesso, però, questa riservatezza era
diventata un ``è tutto bellissimo, ma vorrei provare Riot'', senza un perché.
Lo invitano comunque a un incontro dal vivo (``dal vivo è diverso''), e alla
fine viene assunto.

Una volta arrivato, i sei mesi di prova (\emph{probation}) sono serviti tutti
per smettere di ragionare e lavorare ``come a Blizzard''. La differenza
culturale era enorme, tanto che in Riot esisteva un gruppo di supporto
interno tra ex dipendenti Blizzard: un collega gliel'ha descritto come
``uscire da una relazione tossica e doversi riabituare a una relazione
normale''\footnote{Il relatore precisa che parla della Blizzard di allora,
che non esiste più: oggi è parte di Microsoft e nessuno dei suoi ex colleghi
vi lavora ancora, anche a causa dei licenziamenti di massa.}.
\begin{description}
  \item[Blizzard] era \emph{developer first}: il focus era chiaramente sugli
        sviluppatori, prima di tutti gli altri. Dopo sei anni c'erano ancora
        meeting a cui non poteva partecipare perché non abbastanza
        \emph{senior}, informazioni che non poteva sapere, piani che non
        poteva vedere.
  \item[Riot] era \emph{player first}: prima i giocatori, e i dipendenti come
        mezzo per trasferire il feedback dei giocatori dentro il gioco. Il
        secondo giorno lo portano in una stanza e gli mostrano il piano dei
        cinque anni successivi. Lui, sorpreso, pensa a un errore (``sono qui
        da oggi, non vorrei mettervi nei guai''); la risposta è: ``sei un
        dipendente Riot o no? Allora non ci sono livelli di accesso: devi
        sapere queste cose per fare meglio il tuo lavoro''.
\end{description}

Da questa esperienza ha imparato molto sulla \textbf{trasparenza}: condividere
le informazioni aiuta il lavoro. Non sempre, però: a volte non si condivide
qualcosa non per segretezza, ma perché non è ancora pronta per essere
presentata, esattamente come non si mostra un gioco quando è ancora solo
teoria (``te lo faccio vedere quando ho qualcosa di tangibile'').

La lezione per gli studenti è doppia: le aziende valuteranno il \emph{fit}
culturale sempre di più, ma capire se si è un \emph{fit} richiede anche
\textbf{lavoro da parte del candidato}. Bisogna avere la presenza di spirito
di dire: ``conosco i valori di questa azienda, non mando nemmeno il CV,
perché so che durerei un mese e mezzo''.

% ---------------------------------------------------------------------
\section{Che cos'è il feedback}
\label{sec:z1-cose-feedback}

\subsection{Le prime risposte della classe}

La prima domanda alla classe è: \emph{quando sentite la parola
``feedback'', a cosa pensate?} Le risposte raccolte, a voce e in chat:
\begin{itemize}
  \item un \textbf{riscontro} su ciò che si fa e su come lo si fa, sul lavoro
        o sul modo in cui ci si interfaccia con il gruppo; può essere
        lavorativo oppure comunicativo e comportamentale. Chi ha risposto lo
        vede come qualcosa che, se l'azienda non dà, va attivamente
        richiesto, perché è importante sapere come sta andando per
        migliorarsi;
  \item né positivo né negativo, ma \textbf{costruttivo}: in entrambi i casi
        serve a migliorarsi;
  \item uno \emph{scambio di opinioni}, un \emph{riscontro produttivo};
  \item un feedback ``di reazione'', non verbale: si fa una cosa e si
        osserva come cambia l'ambiente o il comportamento delle persone
        intorno (anche perché a voce la gente spesso si tiene le cose
        negative);
  \item per un designer, anche il \textbf{linguaggio del corpo}: osservando
        un giocatore giocare si può dedurre l'emozione che prova (frustrazione,
        incomprensione).
\end{itemize}

Sulla forma (scritto o orale?) le opinioni divergono: c'è chi lo preferisce
\emph{mille volte} scritto; chi pensa che dipenda dalle dimensioni
dell'azienda (orale nelle piccole, scritto in quelle con team di centinaia di
persone); chi pensa che dipenda dall'importanza del feedback (scritto se
tecnico o pieno di dati, da tenere sotto gli occhi; orale se è un consiglio
veloce).

Il relatore usa due esempi per mostrare che \emph{dipende}:
\begin{itemize}
  \item un feedback sulla programmazione ha senso darlo scritto, magari
        commentando direttamente il codice, più che ``andando a prendere un
        caffè'';
  \item ma se il feedback è ``ieri, durante il meeting, ho percepito male il
        fatto che mi interrompessi di continuo'', ce lo si aspetta
        \textbf{orale}.
\end{itemize}

\subsection{Da che cosa dipende il modo di dare feedback}

Lo \emph{spoiler} della giornata è che \textbf{non esiste una risposta univoca}
su come dare feedback. Il workshop non insegna ``il'' modo, ma fornisce
\emph{tutte le opzioni}, perché il modo giusto cambia in base a molti fattori:
\begin{itemize}
  \item a chi lo si dà: un collega, un superiore, il proprio capo, un proprio
        \emph{report} (un sottoposto);
  \item se è personale o professionale;
  \item se è la prima volta o la terza volta che si dà lo stesso feedback;
  \item se riguarda un progetto importante o secondario.
\end{itemize}

Manca però l'elemento fondamentale, che la classe individua subito:
\textbf{l'altro}. A parità di tutti gli altri fattori, il \emph{carattere
della persona che riceve il feedback} dovrebbe far cambiare il modo in cui lo
si dà. Per questo, prima della lezione, agli studenti era stato chiesto di
fare un test di personalità: non un test di psicologia avanzata, ma uno
strumento per rendersi conto di quanta varietà ci sia anche solo in una
classe di 53 persone (il che non vuol dire trovare 52 modi diversi di dare
feedback, ma lavorarci un po' sì).

Il relatore porta se stesso come esempio. Preferisce:
\begin{itemize}
  \item i feedback \textbf{positivi scritti}, perché è a disagio a ricevere
        complimenti in pubblico;
  \item i feedback sul proprio lavoro \textbf{diretti}: ``vai subito al
        punto e non girarmi intorno''.
\end{itemize}
È molto orgoglioso del suo lavoro e, quando qualcuno lo critica, la sua
primissima reazione istintiva è ``come ti permetti''. Negli anni ha imparato
a farla durare un secondo, ma chi gli dà feedback di persona dice di
vederla ancora negli occhi. Per mitigarla, dice esplicitamente a chi gli dà
feedback: \emph{``ignora il modo in cui reagisco nei primi cinque secondi,
vai avanti, poi mi passa''}. È considerato da tutti molto aperto al feedback
e non ha problemi a dire ``hai ragione, facciamo diversamente'': ma ci sono
reazioni istintive e caratteriali che è difficile controllare al cento per
cento.

\subsection{Con chi si lavora meglio}

Guardando le persone con cui si lavora abitualmente (anche nei progetti
dell'accademia), si noterà che si lavora meglio con due profili:
\begin{enumerate}
  \item chi è \textbf{esattamente come noi}, ma qui è un 50 e 50: alcuni
        caratteri vanno d'accordo con i propri simili, altri non vogliono
        nessuno con lo stesso carattere nel gruppo (e si potrebbe rientrare
        in questa categoria);
  \item chi è \textbf{affine}: ha qualcosa in comune con noi, ma non tutto,
        e proprio per questo può dare punti di vista differenti.
\end{enumerate}
Il lavoro della giornata è capire come valorizzare il secondo caso: ottenere
feedback da persone con caratteri, visioni e background diversi, per
completare una visione del proprio lavoro che è necessariamente limitata. È
come esplorare una mappa con la nebbia di guerra: un'altra persona, al di là
del muro, vede cose che noi non vediamo.

\subsection{Il feedback come somma di piccoli miglioramenti}

Per il relatore il feedback è questo: \textbf{applicare tanti piccoli
miglioramenti periodici}, che si sommano nel tempo. Migliorare dello 0,1\%
ogni giorno porta a fine anno a un miglioramento enorme, e questo vale per i
prodotti, per i giochi, per le persone e per i team.

\paragraph{L'esempio del ciclismo britannico.} La squadra di ciclismo
britannica, fino a una quindicina d'anni fa, non aveva vinto quasi nulla.
Cambiato il team manager, tutti si aspettavano grandi cambiamenti sportivi
(convocazioni, strategie di gara). Arrivarono invece \emph{minuscoli
cambiamenti non sportivi}: un altro chef in trasferta, materassi migliori nei
camper per riposare meglio, un sellino più comodo a fine tappa. Cambiamenti
marginali che da soli non avrebbero cambiato niente; eppure, dopo due o tre
anni di questo lavoro quotidiano, la squadra ha cominciato a vincere a
ripetizione, Olimpiadi comprese, cambiando il modo in cui il ciclismo è
visto in Gran Bretagna\footnote{È la nota strategia delle \emph{marginal
gains} associata a Dave Brailsford (British Cycling e Team Sky); il nome non
è citato in lezione.}. Certo, si è investito anche sulle persone giuste, ma
tutto è partito dai piccoli miglioramenti periodici.

Per gli studenti vale lo stesso: chiarire un conflitto con una persona con
cui si ha ancora poco a che fare, o una situazione che ha lasciato
``l'amaro in bocca'', sembra inutile oggi, ma darà i suoi frutti tra tre, sei,
nove mesi. Sono \emph{investimenti sul futuro}.

\paragraph{Peggiorare ogni giorno.} Esiste anche il percorso inverso:
lasciar correre, ``buttare sotto il tappeto'' le cose perché ``tanto è
inutile parlarne''. Queste cose tornano sempre fuori dopo otto mesi, come il
\emph{libro dei rancori} dei nani di Warhammer: ``ti ricordi quando, un anno
e mezzo fa, all'inizio del progetto, mi hai detto quella cosa? Mi sono
offeso''. A quel punto è quasi impossibile risolvere, perché nessuno ricorda
più con precisione cosa sia successo (``ma quale cena?'', ``forse era in
pizzeria'', ``no, eravamo da un'altra parte''). \textbf{Tenersi le cose nel
cassetto per tirarle fuori quando servono è sempre sconsigliato}: non servono
ed è impossibile usarle in modo costruttivo.

% ---------------------------------------------------------------------
\section{Le cinque fasi di vita di un team}
\label{sec:z1-fasi-team}

\subsection{Perché conoscerle}

Quasi tutti gli studenti lavorano già in gruppo (al primo anno su un gioco da
tavolo, negli anni successivi su un videogioco) e si saranno accorti che un
team attraversa delle fasi: il modo di lavorare insieme cambia nel tempo.
Spesso ci si fa prendere di sorpresa: ``eravamo così affiatati, cosa è
successo?''.

Un esempio da una consulenza fuori dai videogiochi: un'azienda farmaceutica
gli affida un team commerciale (rappresentanti che presentano medicinali a
ospedali, medici e privati). Il problema: ``prima andava tutto bene; ora che
siamo in diciotto dovremmo lavorare meno e meglio, e invece va peggio''. Va
peggio perché le persone nuove sono appena arrivate, mentre il management si
aspettava già di vedere i risultati dell'investimento.

Fino a pochi anni fa le fasi di un team venivano disegnate come una
\textbf{linea dritta verso l'alto}: le fasi erano le stesse, ma viste come un
miglioramento costante. Non si vedevano le parti \emph{fisiologiche} del
lavoro in team, cioè anche i peggioramenti. Accettarle è il primo passo per
accettare i fallimenti: nessun percorso personale o di team è lineare.

Il relatore ammette che il suo modo di lavorare è nato per il 70\% dal
proprio impegno a imparare, fare bene e rispettare le persone, e per il 30\%
dall'aver lavorato con ``incapaci colossali'' che gli hanno insegnato come
\emph{non} fare le cose. A Blizzard un manager gli fece capire che, pur non
sapendo ancora che tipo di \emph{people manager} sarebbe diventato, sapeva
benissimo quale tipo di manager \emph{non} voleva essere.

\subsection{Forming, storming, norming, performing, adjourning}

Le fasi sono cinque, un po' come le fasi del lutto. Molte aziende ne
considerano solo quattro (``e sono quelle che solitamente chiudono''); le
aziende al passo con i tempi ne considerano cinque\footnote{Il modello è
quello di Bruce Tuckman (1965), a cui la quinta fase, \emph{adjourning}, fu
aggiunta nel 1977. Il nome di Tuckman non è citato in lezione; la
trascrizione riporta la quinta fase come ``journing''.}.

\begin{description}
  \item[Forming] Si viene messi insieme e si comincia a lavorare. Non si sa
        nulla degli altri; spesso non si conoscono nemmeno i nomi, ed è
        normale.
  \item[Storming] Si cominciano a ``dare le testate come le capre in
        montagna per determinare il territorio'': questa cosa la faccio io,
        questa la vorrei fare così, chi decide? Si stabilisce chi fa cosa e
        come si divide il lavoro. Nel grafico è una \textbf{fase di
        peggioramento}: che il team litighi o discuta civilmente non cambia
        nulla, il team deve investire parte di tempo ed energie per risolvere
        queste questioni. È normale e \textbf{non è evitabile}.
  \item[Norming] Si comincia a capire come lavorano gli altri, cosa sanno
        fare e come far funzionare il gruppo; iniziano i miglioramenti. È una
        forma di \emph{accettazione} degli altri, che deriva dall'aver capito
        cosa devo fare io e cosa devono fare loro.
  \item[Performing] C'è uno sforzo \textbf{proattivo}: ``questa cosa
        toccherebbe a me, ma so che Alessio la fa molto meglio, quindi gli
        chiedo aiuto''. Si sfruttano i punti di forza degli altri, perché li
        si conosce e perché si sa che passare il testimone non avrà effetti
        negativi (il capo non dirà ``perché non l'hai fatta tu?'').
  \item[Adjourning] Il progetto finisce, si chiude o viene riorganizzato, il
        team viene sciolto. Quello che andrebbe fatto è analizzare
        \textbf{cosa ha funzionato, cosa no e cosa migliorare} per il
        progetto successivo.
\end{description}

\subsection{La fase che tutti saltano}

``Spoiler incredibile'': l'\emph{adjourning} è la fase che molte aziende
saltano, \textbf{perché c'è sempre il progetto dopo}.

A PG Esports, per la Milano Games Week (come per qualsiasi evento dal vivo a
cadenza annuale: Lucca Comics, Etna Comics, Comicon\dots), finita l'edizione
di un anno si cominciava dal giorno dopo a lavorare a quella successiva,
perché un anno sembra tanto ma non basta mai. Nei videogiochi è lo stesso:
finito un gioco, a parte le patch di mantenimento e del \emph{day one}, si
passa subito al successivo.

Il cambiamento introdotto dal relatore (in parte già praticato, ma da lui
strutturato) è stato: finita la Games Week, \textbf{tutti in ferie per una
settimana}; poi ci si ritrova per definire cosa ha funzionato e cosa no, con
un \textbf{post-mortem} (o \emph{debrief}); solo dopo si inizia a lavorare
all'edizione successiva. Senza questa fase manca la parte di
\emph{learning}, cioè di ciò che si è imparato dal progetto.

Il consiglio vale per gli studenti: quando si finisce il gioco da tavolo, e
poi il gioco digitale, bisogna trovare il tempo per il post-mortem,
altrimenti si butta via \textbf{il 90\% dell'utilità del progetto}. Una
struttura semplice per farlo è descritta nella sezione~\ref{sec:z1-qa-finale}.

% ---------------------------------------------------------------------
\section{Domande sulle fasi del team}
\label{sec:z1-qa-fasi}

\subsection{Lo storming che ritorna}

\paragraph{Capita che lo storming si ritagli spazio anche nelle altre fasi?}
Lo \emph{storming} è una fase precisa, che però molti manager, coordinatori e
team cercano di \textbf{saltare}: ``se io sono bravo come manager e loro sono
bravi come team, passeremo da forming a norming senza problemi''. È un errore,
perché senza una fase in cui si discute davvero non si sviluppa una
\emph{skill} fondamentale: \textbf{saper discutere}. Tutti conoscono persone
non abituate a farlo, per cui ogni minimo disaccordo diventa un affronto
personale (``avevamo deciso la pizza, perché andiamo al sushi? Perché mi fate
questo affronto?''): sono persone che non sono abituate ad avere conflitti
costruttivi.

Se si pensa di essere in \emph{norming} e continuano a emergere conflitti, il
punto è che \textbf{non si era davvero in norming}: lo storming non ``si
trascina'' nelle fasi successive, semplicemente non se n'è mai usciti.
Allo stesso modo, a chi chiede se ogni mini-fase di storming durante la vita
del gruppo sia frutto di uno storming fatto male o solo di carattere, la
risposta è già contenuta qui.

Un partecipante osserva che spesso, ancor prima di finire un progetto, si è
già messi su quello successivo: è il modo migliore per essere in costante
\emph{damage control}. Era l'esperienza del relatore a Blizzard: veniva
spostato fisicamente da un team all'altro, e finito il lancio di
\emph{Diablo~III} era già su \emph{World of Warcraft}, senza vivere il
``dopo''. Si è dovuto ritagliare da solo lo spazio per un debrief personale,
ed è lì che ha deciso che tipo di manager non voleva essere. Il consiglio:
\textbf{se non lo fa il team, fatelo voi}.

\subsection{Dove si passa la maggior parte del tempo}

\emph{In quale fase si passa più tempo nella propria carriera?} Gli studenti
rispondono \emph{storming} e \emph{norming}. Il relatore, in quindici anni,
calcola di aver passato in \emph{performing} circa \textbf{dieci mesi}. Non è
pessimismo: arrivare a performing deve essere sempre l'obiettivo, sia da
manager sia da membro del team; ma è difficile arrivarci, ed è normale che non
sempre ci si arrivi.

Lo storming può durare a lungo, ma con team molto esperti dura pochissimo,
perché c'è chi lo sa gestire. In realtà \textbf{la maggior parte del tempo si
passa in forming}, perché ogni equilibrio raggiunto vale per \emph{quelle}
persone, \emph{quel} progetto, \emph{quelle} tempistiche. Basta che cambi un
qualsiasi fattore (una deadline spostata, una persona nuova, qualcuno che si
dimette, un progetto cancellato) e si ricomincia, perché i bilanciamenti
precedenti non valgono più.

Esempio: un team di cinque persone affiatatissimo, in pieno performing. La
dirigenza, entusiasta, ``raddoppia le risorse'' e lo porta a dieci. O si
lascia tutto com'è (e i cinque nuovi non fanno niente), o bisogna ridiscutere
chi fa cosa e come i nuovi possono aiutare: che piaccia o no, si ricomincia
da forming.

Quindi la bravura di un manager, ma anche di un membro del team,
\textbf{non sta nel passare tanto tempo in performing, ma nell'accelerare il
passaggio da una fase all'altra}, e nel fare bene la fase di
\emph{adjourning}.

\subsection{Analizzare anche i successi}

C'è la tendenza a fare post-mortem solo dei progetti andati male: ``questa
cosa non ha funzionato, vediamo perché''. Va benissimo. Ma quando un progetto
va bene si dà per scontato che sia andato bene ``perché era buono''. Bisogna
invece chiedersi \emph{perché} è andato bene: era il prodotto? La
tempistica? Il modo in cui abbiamo fatto una cosa piuttosto che un'altra?
\textbf{Se non si capisce davvero il motivo del proprio successo, si
tenderanno a replicare le cose sbagliate.}

\paragraph{Blockbuster.} Quando arrivò la notizia che Netflix avrebbe
lanciato un servizio di streaming, Blockbuster (la catena di noleggio di
videocassette e DVD) aveva tutto il tempo per preparare un'offerta
concorrente. In una riunione del management raccontata in un libro (di cui il
relatore non ricorda il titolo), la direzione sembrava orientata a investire
nello streaming, finché un dirigente fece notare che il 12\% del fatturato
derivava dalle \emph{multe per riconsegna tardiva}, impossibili con lo
streaming, mentre secondo le stime lo streaming avrebbe prodotto il 10\% del
fatturato. ``Non possiamo rinunciare a questo 2\%'': niente streaming. Oggi
non c'è più un Blockbuster ``manco a cercarlo col lanternino''. Non avevano
capito qual era il vero motivo del loro successo, cioè il loro modello di
business\footnote{Le percentuali sono quelle riferite in lezione e non sono
state verificate.}.

\paragraph{The Founder.} Nel film sulla nascita di McDonald's c'è una scena in
cui il futuro socio del protagonista gli dice: ``tu pensi di essere nella
\emph{food industry}, e invece sei nella \emph{real estate industry}''. Non
dovrebbe investire sugli hamburger, ma sui terreni su cui costruisce i
ristoranti. Anche per un team vale lo stesso: bisogna capire che cosa
l'ha reso davvero di successo, oppure no.

\subsection{L'autoanalisi personale}

\paragraph{E quando si cambia azienda?} Uno studente chiede se il debrief sia
possibile anche quando si lascia un'azienda. Non si può tornare in ufficio a
discutere della propria esperienza, ma lo si può fare \textbf{per se stessi}.
Il relatore l'ha fatto, anche se all'epoca in modo poco consapevole, ed è
così che è passato dal \emph{community management} al \emph{people
management}. Analizzando cosa aveva fatto bene a Blizzard si era reso conto
che:
\begin{itemize}
  \item aveva costruito un ottimo rapporto di fiducia con i giocatori (con
        alcuni proprietari di fansite parla ancora oggi);
  \item pur senza la seniority, era diventato un punto di riferimento per il
        team, visto come imparziale ed equilibrato anche su temi che non lo
        riguardavano direttamente;
  \item sapeva \textbf{risolvere problemi}: durante il lancio di
        \emph{Diablo~III} era il più tranquillo di tutti, non perché
        sminuisse il problema, ma perché sapeva che si sarebbe risolto con i
        suoi tempi; il lavoro non era spegnere l'incendio ma ``far sfollare
        la gente'' e aspettare che tutto tornasse in ordine.
\end{itemize}
E aveva capito cosa sapeva fare meno bene, ad esempio \textbf{gestire il
proprio manager}. Avere un manager non significa che il manager gestisca te:
il manager gestisce il tuo lavoro, ma tu devi saper gestire i tuoi
\emph{stakeholder}. Spesso lo stakeholder è il proprio manager; da manager,
gli stakeholder saranno CEO, CFO e così via. Gestirli significa conoscerli,
sapere come portare loro le informazioni, come far presente che qualcosa non
va. Fa parte del lavoro, anche se a lui non piace, e ci ha lavorato.

Nel mese e mezzo di pausa prima di entrare in Riot si è concentrato su
questo: ``cosa porto come mio punto di forza?''. L'autoanalisi è culminata
nel chiedere espressamente alla persona che l'aveva assunto: ``prima mi avete
detto di no, ora mi avete preso: perché?''.

\paragraph{Come si capisce se un progetto è andato bene?} Dipende
dall'\textbf{obiettivo}. Se l'obiettivo era vendere milioni di copie e il
gioco le ha vendute, bisogna capire perché. Se, più ragionevolmente,
l'obiettivo era un gioco da tavolo che permettesse di mettere in pratica le
conoscenze del corso, va valutato per quello: ha permesso di farlo? Se sì,
perché? Se no, perché no? L'obiettivo va fissato \textbf{all'inizio del
progetto, non alla fine}, proprio per poterlo valutare dopo.

\subsection{Quando il team cambia di continuo}

\paragraph{Tutti questi cambiamenti non sono dannosi per il progetto?}
Assolutamente sì, ma quasi mai sono volontari. I progetti andati male sono
spesso progetti grandi.

Il relatore racconta di \emph{Cyberpunk 2077}. Conosceva persone in CD
Projekt Red (``il mondo è piccolissimo, le cose si sanno''). Alla Gamescom va
con Diego Ricchiuti, direttore dell'accademia, a vedere la presentazione del
primo gameplay, quella ``in cui regalavano il giubbottino da samurai''.
Finito il video, tra gli applausi, si guardano e capiscono entrambi che
\textbf{il gioco non è pronto}, a soli quattro mesi dall'uscita prevista:
Diego aveva notato il combattimento e la guida. Gli amici di CD Projekt Red
(uno aveva lavorato con lui a Blizzard per sei anni) gli confermano tra le
righe che ``ce la faremo, ma ci sono ancora cose da sistemare''. Il gioco fu
poi rimandato più volte\footnote{In lezione il relatore dice ``due, forse
tre volte''; l'uscita di \emph{Cyberpunk 2077} è stata effettivamente
rinviata tre volte prima del dicembre 2020.}.

Per deformazione professionale, chi ha sviluppato l'occhio critico non
riesce più a godersi un trailer: in un trailer di \emph{The Day Before} Diego
gli aveva fatto notare che l'effetto della fiammata di un fucile era stato
chiaramente aggiunto in postproduzione e non faceva parte del gameplay. ``È
il triste epilogo di questa storia'': non riescono più a godersi un trailer.

Il punto è che più il team è grande, più queste variabili vengono imposte, e
più è facile che il fallimento \textbf{non dipenda dal fatto che il prodotto
sia brutto} (\emph{Cyberpunk} è un bellissimo gioco sotto molti punti di
vista, e il relatore lo sta rigiocando ora che è stato corretto ed espanso).
Nei grandi gruppi qualcuno arriva, qualcuno si dimette, cambia il budget,
arriva il rinvio; nel caso di \emph{Anthem}, il team ha fatto così tante volte
forming--storming--forming--storming che ne è uscito ``un prodotto che
dovevamo espellere: l'abbiamo espulso, buttato fuori dal nostro organismo pur
di liberarcene''.

Tra i cambiamenti c'è anche la crescita verticale: se il capo viene promosso
perché lavora benissimo, bisogna scegliere un nuovo capo, ed è sempre un
caos (chi ci resta male, chi pensava di meritarlo, chi voleva il suo amico,
chi non sopporta il prescelto). Per questo il relatore è \textbf{contrario
alle promozioni durante i progetti}: si rischia di rovinare tutto il lavoro di
storming.

\subsection{Obiettivi chiari nei piccoli gruppi}

Nei gruppi piccoli, come quelli degli studenti, il problema deriva più
spesso dal \textbf{non avere chiaro che cosa si vuole ottenere}: un prodotto
puramente educativo e formativo, o un prodotto richiesto dall'accademia che,
se viene bene, si potrebbe anche vendere? Un gioco commerciale o un gioco
artistico? Va chiarito subito, perché ci saranno sempre momenti in cui
scegliere tra A e B, e a guidare la scelta sarà la domanda: \emph{quale delle
due opzioni ci porta più vicino all'obiettivo?} Se l'obiettivo non è chiaro,
non si sa cosa scegliere.

Questo porta molte aziende a fare \textbf{scelte di design che in realtà sono
scelte commerciali}. Le microtransazioni sono scelte commerciali: nessuno ha
mai pensato ``sai cosa piacerebbe tanto ai giocatori? Pagare le bustine''.
Prima le bustine si guadagnavano giocando; ora la tendenza è ``potresti
giocare, ma se paghi è meglio''. Si pensi alle polemiche su Ubisoft per
l'opzione di ``saltare il gioco'' pagando, o ai famosi \emph{boost} di livello
di \emph{World of Warcraft}: ``devo pagarti per saltare il gioco?''.

Una studentessa (Silvia) cita \emph{Final Fantasy XIV}, che a un certo punto fu
``resettato'' perché non funzionava: un esempio positivo. Il relatore aggiunge
quello negativo di \emph{Age of Conan}, che non è mai riuscito a fare quel
reset iniziale (``è una schifezza, ve lo ripresentiamo quando è pronto''). Ci
vuole consapevolezza, ma anche i mezzi: un conto è Square Enix, che può
permetterselo, un altro è un'azienda che non può chiudere tutto per un anno
continuando a pagare gli stipendi.

Riassumendo: nei piccoli gruppi i problemi sono più di
\textbf{comunicazione e chiarezza} (obiettivi non chiari o che si spostano);
nelle grandi aziende di più di \textbf{persone che vanno e vengono, obiettivi
e tempistiche che cambiano}. In accademia le scadenze sono fisse (il gioco da
tavolo si consegna a fine anno scolastico, a giugno, e giugno resta); in
azienda può succedere che una feature prevista per la \emph{vertical slice}
di luglio venga anticipata a maggio per portarla alla Gamescom, e un sistema
perfetto per consegnare a giugno può non esserlo più per consegnare a maggio.

\subsection{L'ambiente determina la qualità del prodotto}

\paragraph{È vero che chi passa per Electronic Arts finisce per fare brutti
giochi?} Uno studente cita BioWare, che faceva ottimi giochi finché era
autonoma. Il problema non è dei developer, che non sono ``improvvisamente
diventati incapaci'': è evidente che il problema è altrove, e che
\textbf{l'ambiente in cui si fa un gioco determina molto spesso la qualità del
prodotto}.

Un esempio locale positivo: Tommaso Bonanni, conosciuto dal relatore insieme a
Diego e già docente dell'accademia, ha fondato a Roma Caracal Games Studio,
applicando fin dall'inizio questi concetti di cultura aziendale. L'azienda è
rimasta piccola, ma ha molto chiaro cosa vuole fare e come. Quando si
predispone fin dall'inizio l'azienda e il team a lavorare in un certo modo è
fattibilissimo; è molto più difficile quando si è EA, si acquisisce
un'azienda e le si dice ``da oggi si fa così''.

Il relatore ha vissuto la transizione da Blizzard ad Activision Blizzard (uno
dei motivi per cui se n'è andato, volontariamente). Ha visto la cultura
interna cambiare per gradi:
\begin{enumerate}
  \item ``facciamo giochi fighissimi, perché i giochi devono essere
        fighissimi'';
  \item ``facciamo giochi fighi così vendono di più'';
  \item ``vendiamo di più facendo giochi fighi'' (sembra uguale, ma è molto
        diverso);
  \item ``facciamo più soldi''.
\end{enumerate}
Lo si capiva dalle scelte di design, osservando quale fosse il peso più
grande sulla bilancia. Lo stesso varrà per gli studenti, nei loro giochi e nei
colloqui di lavoro: capire subito che cosa vuole l'altra parte.

\paragraph{Leader e visione.} Nei gruppi dell'accademia, dal secondo anno, c'è
una persona che assume il ruolo di \emph{lead} o \emph{scrum master}. È
diverso dall'essere un \emph{game director} (``il gioco si fa come dico io,
aiutatemi a farlo con la mia visione''), come avviene nei grandi studi: da lì
nascono i nomi nei titoli, come \emph{Sid Meier's Civilization}, che indicano
la persona che ha fatto quel gioco, oppure un puro branding, come \emph{Tom
Clancy's Rainbow Six}.

\paragraph{Inclusion come jolly.} A margine, uno studente osserva che in molte
aziende ``la risposta giusta è sempre \emph{inclusion}'', usata come un
jolly. È comunissimo, conferma il relatore: \emph{diversity and inclusion}
viene spesso usata per dire ``così includiamo più clienti'', ma non è quello
il suo punto (``altro discorso'').

% ---------------------------------------------------------------------
\section{Confrontarsi con chi è più bravo}
\label{sec:z1-rivalita}

Arrivato a Riot dopo sei anni a Blizzard, il relatore pensava di avere un buon
bagaglio di esperienza. Si è trovato invece ``come un fante buttato in mezzo
ai Navy SEAL'': persone che avevano superato 11, 12, 13 colloqui, con un
livello molto alto.

Individua subito un collega francese che, nella sua testa, diventa ``la
versione migliore di me'': stesso ruolo (community manager), esperienza
simile (arrivava da Ubisoft, come quasi tutti i francesi di Riot), ma molto
più spigliato, sempre con la battuta pronta, ``l'anima della festa''. Il
relatore invece è introverso: un workshop come questo per lui è
\emph{emotivamente drenante}, non ama parlare in pubblico pur amando
condividere le proprie esperienze, e non riesce a sostenere conversazioni
come ``cosa hai fatto questo weekend?''. Lavorando raramente insieme, aveva
sviluppato una rivalità ``da allenatore di Pokémon'' (lui era Ash, il
relatore l'antagonista).

Dopo sei mesi, a un evento di \emph{team building} aziendale (a Tenerife), il
suo capo, molto attento a queste dinamiche, li mette apposta nello stesso
gruppo. A fine serata, complice l'open bar, il relatore gli confessa di averlo
mal sopportato dal primo giorno per invidia della sua estroversione. Il
collega ride e risponde che pochi giorni prima aveva detto al capo la stessa
cosa al contrario: invidiava a Stefano il suo essere sempre organizzato, sul
pezzo, presentarsi nella sala riunioni giusta (mentre lui finiva in quella
sbagliata), arrivare con gli appunti, con i dati a supporto e le analisi di
mercato, mentre lui spesso non aveva nemmeno letto la mail del meeting.
Un'invidia in due direzioni, che li ha portati a lavorare molto insieme e a
migliorarsi a vicenda.

Le lezioni:
\begin{itemize}
  \item non bisogna demoralizzarsi vedendo qualcuno che fa lo stesso lavoro
        ed è più avanti sotto alcuni aspetti: spesso c'è un bilanciamento. I
        ``punti caratteristica'' sono limitati: se li ha messi in carisma,
        non li avrà messi in saggezza;
  \item conviene \textbf{chiedere a queste persone come mai} sono così brave
        in qualcosa, come si fa con un progetto andato bene: magari non lo si
        farà mai bene quanto loro, ma si comincerà ad avere il
        \emph{mindset} giusto; e viceversa, loro saranno spinti a chiedere a
        noi;
  \item spesso \textbf{le proprie doti le scopriamo dagli altri}. Il relatore
        non si era mai accorto di essere ``quello che resta calmo durante le
        emergenze'' (il lancio di \emph{Diablo~III}, quello di
        \emph{Draenor}, i problemi dal vivo alla Games Week) finché non
        gliel'hanno fatto notare tutti: per lui era la normalità. A chi pensa
        ``non sono capace, forse non è il mio percorso'', può darsi che
        semplicemente nessuno gli abbia ancora detto ``questa cosa la sai
        fare davvero bene, concentrati su questo''.
\end{itemize}

% ---------------------------------------------------------------------
\section{Domande sul lavoro nell'industria}
\label{sec:z1-industria}

Diverse domande degli studenti hanno riguardato le condizioni di lavoro nei
grandi studi internazionali. Le risposte sono raccolte qui per argomento.

\subsection{Presenza, remoto e trasferimenti}

\paragraph{Grandi studi: quasi solo in presenza.} Sia a Blizzard sia a Riot si
lavorava \textbf{solo in presenza}, senza smart working: sei anni a Parigi e
cinque a Dublino. Si parlava solo inglese: anche nei meeting tra soli
italiani si tendeva a comunicare e a scrivere in inglese, per correttezza
verso chi si fosse aggiunto e per lasciare una traccia comprensibile a tutti.
Il francese serviva per fare la spesa e sopravvivere, non in ufficio.
L'ambiente era estremamente eterogeneo (francesi, polacchi, spagnoli e
molti altri): durante il colloquio dal vivo a Dublino, mentre si preoccupava
di essersi vestito in modo adeguato, vide passare dietro la parete a vetro un
collega svedese con i dreadlock, a piedi scalzi, con una pianta sotto un
braccio e un aquilone sotto l'altro, e capì di essere ``a posto così''.
Un'esposizione a molte culture e modi di pensare che il relatore considera
bellissima.

Anche oggi, tranne alcuni studi particolari, la maggior parte del lavoro si
svolge di persona: amici a Paradox lavorano a Stoccolma, i pochi rimasti a
Blizzard si sono dovuti trasferire a Irvine. Per posizioni \emph{junior} o
\emph{entry level} i grandi studi raramente concedono il lavoro da casa; per
profili senior può esserci qualche eccezione (ad esempio nove mesi da remoto
e tre in sede).

\paragraph{Relocation.} Spesso tocca spostarsi, ma il relatore non ha mai
pagato un trasloco di tasca propria: queste aziende offrono la
\emph{relocation} (aiuto a trovare casa, ad aprire il conto in banca e
supporto pratico), perché è un investimento sulla persona.

\paragraph{Colloqui mentre si lavora altrove.} Per i colloqui in presenza
prendeva ferie, senza dire dove andava. Il consiglio generale: per fare
colloqui \textbf{prendete ferie o permessi}, non datevi malati, ``che non è un
bel modo di comportarsi''.

\paragraph{Remoto ``doppio''.} Un collega lavora per l'azienda coreana che
gestisce \emph{PUBG}, con ufficio ad Amsterdam, ma lavora da Milano: un
``doppio remoto''. Si trova, ma in genere un ambiente anglofono con sedi in
Europa chiederà di trasferirsi in Europa.

\paragraph{Famiglia e trasferimenti.} Il relatore ha avuto la fortuna di
sposare una persona che fa lo stesso lavoro (conosciuta lavorando su
\emph{Warhammer Online}, poi con lui a Blizzard e a Riot, nella
localizzazione e nel doppiaggio, ad esempio di \emph{World of Warcraft} in
italiano): avendo competenze rare, spesso venivano assunti entrambi, al
massimo a pochi mesi di distanza. Per i developer la soluzione frequente è
che la famiglia segua la persona assunta, perché lo stipendio basta per
tutti. Ma non sempre c'è equilibrio, e a volte uno dei due deve sacrificarsi:
sua moglie ha lasciato Blizzard per seguirlo a Dublino (per poi essere
assunta da Riot); lui ha lasciato Riot per tornare in Italia quando lei ha
trovato lavoro qui. Ha dovuto rinunciare più volte a lavorare in Islanda per
CCP (\emph{EVE Online}), perché non era una scelta solo sua. Si fanno
compromessi, ma si riesce a far quadrare.

\paragraph{Ci si sposta per ogni progetto?} No: si sono spostati
esclusivamente per cambiare azienda. Cambiando progetto nella stessa azienda
è raro, quasi mai (``non l'ho mai visto succedere''), che venga chiesto di
trasferirsi. In Francia stavano benissimo (vivevano fuori Parigi; ``i
problemi con i francesi sono in realtà problemi con i parigini''), parlavano
francese ed erano vicini alla doppia cittadinanza: se ne sono andati perché
non vedevano un futuro a Blizzard, per il modo in cui l'azienda trattava i
colleghi.

\paragraph{Il mercato oggi.} Non esiste un \emph{silver bullet} per trovare
lavoro domani dove si vuole. Con tutti i \emph{layoff} recenti ci sono
migliaia di esperti del settore che cercano lavoro: serve pazienza, ma è
fattibile.

\paragraph{Studi indie e remoto.} Negli studi più piccoli il remoto è più
diffuso: a Stormind circa 60 persone su 95 lavorano da remoto, ma lo studio
si è strutturato fin dall'inizio per farlo. Serve consapevolezza. Bending
Spoons (che sviluppa app per mobile, non giochi, ed è l'azienda dell'app
Immuni), pur essendo grandissima e con sede a Milano, lavora da remoto e
lascia l'ufficio a chi vuole: lo smart working è posizionato quasi come un
\emph{benefit}, ma permette anche di avere un ufficio da 300 persone invece
che da 700. Le sedi fisiche costano decine di milioni all'anno (Blizzard ha
chiuso le sedi di Amsterdam e di Parigi). Tuttavia, per lavori molto
verticali e specifici gli studi preferiscono avere le persone dal vivo; i
ruoli più spesso remoti sono quelli che non fanno sviluppo (customer service,
PR\dots).

Anche gli artisti lavorano spesso in presenza: l'accademia può dare solo una
\emph{simulazione} di quanto stretto sarà il contatto tra programmazione,
design e art. ``Sarà ancora di più.''

\paragraph{La comunicazione da remoto.} È più difficile. Pur essendo
estremamente introverso, il relatore \textbf{obbliga tutti a usare la webcam}
nei meeting in cui conta: non in un aggiornamento con novanta persone
collegate, ma certamente in un colloquio a due sulla performance dell'ultimo
anno. La componente fisica è una parte enorme della
comunicazione\footnote{In lezione si parla di ``quasi l'ottanta per cento'';
il dato va preso come indicativo, poiché le percentuali di questo tipo
circolano in forme diverse e sono spesso citate fuori contesto.}: anche se si
è timidi, accendere la webcam è l'approssimazione più vicina all'essere dal
vivo. Negli ambienti dove si lavorava senza webcam la situazione era
``esponenzialmente peggiore''.

\subsection{Differenze culturali nel team}

\paragraph{Gestire sensibilità diverse.} A Riot i team erano divisi per
\emph{lingua}, non per nazione: il team di lingua spagnola includeva anche un
messicano, quello di lingua russa includeva russi e ucraini, proprio nel
periodo della guerra in Crimea. Tra i colleghi non c'era alcun problema (erano
amici), ma ci si è posti la domanda: \emph{come comunichiamo con una community
che sarà divisa su questo tema?} Nei giochi online la divisione era evidente:
in \emph{Rocket League} c'era chi, con un compagno di squadra con la bandiera
avversaria, si faceva otto autogol pur di farlo perdere. Allo stesso modo,
con molti americani mandati a Dublino a coprire posizioni scoperte, c'erano le
elezioni (Clinton--Trump, Biden--Trump) e varie guerre: serviva sempre una
certa sensibilità.

Imparare a navigare queste situazioni lo ha reso un professionista migliore:
gli ha insegnato a fermarsi e chiedersi ``io la percepisco così; gli altri
dieci, con background, storie e lingue diverse, come la penseranno?'', a
ragionare più \emph{long term} e più \emph{big picture}, con più empatia.

\paragraph{Stili di comunicazione.} Un partecipante che ha lavorato nel Regno
Unito chiede dello stile comunicativo italiano, a volte frainteso. Uno dei più
grandi complimenti ricevuti dal relatore è stato ``non sembri italiano''
(``gesticoli poco per essere italiano'', ``parli bene inglese per essere
italiano''). Ha imparato ad \textbf{adattare la comunicazione alla persona}:
con spagnoli e francesi la battuta sulla nazione funziona, perché c'è un
``odio--amore'' condiviso; con russi e giapponesi no. Doveva adattare anche
l'umorismo. Un collega metà russo e metà messicano diceva di aver preso ``il
lato peggiore di entrambe le nazionalità'', ed è stato un ottimo maestro per
imparare a leggere culture diverse.

Non bisogna generalizzare, ma gli stereotipi possono essere un punto di
partenza ragionevole, \emph{sapendo che troverete sempre le eccezioni}. E
vanno letti correttamente: gli americani tendono a essere più estroversi e ad
avere sempre qualcosa da dire non perché ``sono fatti così'', ma perché
vengono educati così (all'università tra i corsi c'è il \emph{debate}, saper
argomentare in pubblico); un collega giapponese può essere più introverso per
ragioni culturali. Nel dubbio si parte da lì, poi si conosce la persona.

\subsection{Negoziare, discriminazione e stipendi}

\paragraph{La famiglia come leva.} Avere un partner riduce le opzioni, ma le
aziende lo sanno: uno spostamento è per due. Si può usare come leva per
negoziare (un compenso diverso, un supporto per il trasloco): se un'azienda ha
deciso di volere quella persona, 2, 3, 4 o 10.000 euro in più all'anno, per
un'azienda che fattura centinaia di milioni o miliardi, non cambiano nulla.

\paragraph{Essere sposati è un motivo per non essere assunti?} No. In alcuni
paesi è proprio illegale discriminare su questo: in Francia era illegale
chiedere CV con data di nascita, foto, sesso o religione, e chiedere se si
vogliono figli. Purtroppo, dai racconti dei suoi studenti di master (in
settori come \emph{beauty and wellness} e \emph{food and wine}), in piccole
realtà italiane queste domande (``hai 25 anni, hai intenzione di
sposarti?'') vengono ancora fatte, e servono a filtrare i candidati. Nelle
grandi aziende ci sono sistemi per evitarlo; anzi, molte danno \emph{più}
benefit a chi ha famiglia (``porta pure la famiglia, è tutto pagato''),
perché le persone competenti hanno spesso una certa età e quindi,
ragionevolmente, una famiglia.

\paragraph{Gli stipendi.} Dipendono molto da dove si vive. Lo stipendio a Riot
non era paragonabile a quello di Blizzard perché in Francia era tassato al
52\% e in Irlanda al 17\%\footnote{Aliquote come riportate in lezione,
riferite alla situazione personale del relatore.}: nessuno stipendio
competitivo può compensare il 35\% di differenza di tasse. Negli Stati Uniti,
con pochissima previdenza pubblica, gli stipendi sono molto più alti perché
poi ci si paga tutto (``se hai un raffreddore paghi tu''). Bisogna valutare il
potere d'acquisto di quei soldi nel paese in cui si andrà a vivere.

\subsection{Freelance e \emph{work for hire}}

All'inizio della sessione pomeridiana uno studente che lavora nell'audio
chiede quali siano le condizioni per chi lavora come \emph{freelance}.
\begin{itemize}
  \item Dipende molto dal progetto: spesso il contratto è a progetto, per
        portare a termine una parte o un intero progetto.
  \item Non tutte le aziende usano freelance: Riot e Blizzard lo facevano
        pochissimo e preferivano assumere (magari trattando poi la persona
        ``come un freelance'', lasciandola a casa a fine progetto).
  \item Spesso i lavori offerti come freelance sono in realtà
        \textbf{sostitutivi di un full time} (``freelance, ma sei ore al
        giorno'').
  \item Si può fare, ma bisogna essere molto chiari nel definire l'impegno
        (quante ore, come vengono tracciate): non tutti specificano cosa
        significa davvero ``freelance''. Lavorando così per due o tre aziende
        si può arrivare a un tempo pieno.
  \item È più comune il \emph{work for hire}: aziende che forniscono personale
        ad altre aziende; si è assunti da X ma si lavora ``embeddati'' presso
        Y (una sorta di subappalto). Anche lo studio del relatore lo fa.
  \item Il puro freelance di quattro mesi esiste, ma di solito per quei
        quattro mesi diventa un lavoro a tempo pieno; è raro ``quattro mesi,
        due ore al giorno''. Nelle grandi aziende è improbabile lavorare nello
        stesso mese per Ubisoft, Riot e Blizzard tre ore ciascuna; è più
        realistico lavorare per più startup indie italiane, che non possono
        permettersi un'assunzione ma possono pagare due ore al giorno.
  \item Molto passa dai contatti: trovare le aziende giuste e rendere esplicite
        le condizioni.
\end{itemize}

% ---------------------------------------------------------------------
\section{Il circolo virtuoso: feedback, miglioramento, fiducia}
\label{sec:z1-triangolo}

\subsection{Il triangolo}

L'obiettivo a cui il workshop vuole arrivare è un \textbf{triangolo che si
autoalimenta}:
\begin{enumerate}
  \item si dà un \textbf{feedback} con l'obiettivo di migliorare le
        performance (inclusi elementi caratteriali);
  \item la persona \textbf{migliora} e vede di averne beneficiato;
  \item questo genera \textbf{fiducia}, che rende più facile dare e ricevere
        il feedback successivo.
\end{enumerate}
Togliendo uno qualsiasi dei tre elementi, il circolo non è altrettanto forte.

In qualsiasi fase si parta, il rispetto per i colleghi è inizialmente dettato
dall'\emph{educazione}, non dalla fiducia. La fiducia arriva quando i colleghi
cominciano a dire cose giuste e sensate e a \textbf{mantenere la parola}.
Dare e ricevere feedback è il modo in cui, in un certo senso,
``inganniamo il cervello'' per creare una \emph{case history} di fiducia: si
ha motivo di fidarsi di più, di ascoltare di più, di essere più ricettivi. E
si diventa sempre più bravi ad accettare il feedback e ad agire di
conseguenza.

\subsection{Accettare un feedback non vuol dire fare come dice l'altro}

Dieci anni fa il relatore non era ostile ai feedback, ma molto
\emph{difensivo}: ``sì, sì, va bene, ora ci guardo''; oppure, a chi gli diceva
``avresti dovuto farla così'', rispondeva spiegando perché l'aveva fatta in
quel modo. Oggi ha imparato a dire ``ok, grazie, ora ci ragiono'': non deve
difendere la propria posizione da chi viene solo a dire ``hai pensato a
questa cosa?''. Prima lo viveva come un attacco (``mi viene a dire che non so
fare il mio lavoro''), ora come un ``ottimo punto di vista, lo prendo in
considerazione''.

Con una postilla fondamentale: \textbf{accettare il feedback non significa
fare come dice l'altro}. ``Mi hai detto che preferisci questo mob verde. Ho
pensato anche a farlo verde? Sì. Penso comunque che giallo sia meglio? Sì, lo
faccio giallo.'' Il relatore accetta \emph{sempre} il feedback, valido o no:
garantisce di tenerlo a mente quando lavora. Che poi faccia come gli viene
detto è un altro discorso.

\subsection{Fiducia e scadenze}

La fiducia funziona allo stesso modo con gli impegni. Se si chiede a qualcuno
una cosa per venerdì e la consegna venerdì, la volta dopo gliela si chiederà
per venerdì. Se invece arriva il giovedì della settimana dopo, la volta dopo
gli si darà una \textbf{deadline anticipata}, come con chi arriva sempre in
ritardo agli appuntamenti (``se ti dico alle 7 arrivi alle 7 e mezza; la
prossima volta ti dico alle 6 e mezza'').

Al lavoro questo è dannosissimo, perché crea \textbf{deadline fittizie} che
non sono mai quelle vere, e nascono discussioni del tipo ``perché devo farlo
in due giorni se si fa in cinque?'' (``sì, se tu lo facessi in cinque''). La
conversazione giusta è un'altra: \emph{``che problemi hai avuto nel
consegnare venerdì?'' ``Questo, questo e questo.'' ``Cosa posso fare per
evitare che succeda?''}

Il relatore non vuole mai sentirsi dire ``non sono riuscito a farlo''.
Preferisce di gran lunga sentirsi dire, tre giorni prima, \textbf{``potrei non
riuscire a farlo, per questi motivi''}. Chi ha la visione d'insieme può dare
più budget, più persone, spostare la deadline, cambiare lo scopo del progetto:
può fare qualsiasi cosa per aiutare. Ma se lo si scopre il venerdì, cinque
minuti prima della deadline, non può fare più niente. Il messaggio è:
\textbf{``aiutami a poterti aiutare''}. Non penserà mai male di chi dice ``non
ce la faccio per questi motivi'', perché non sta dicendo ``non ho voglia'': sta
indicando i problemi, e lavorare su quei problemi è compito del manager. Non
può invece fidarsi di chi fino al giorno prima dice ``tranquillo, ce la
faccio'' e il giorno dopo ``non ce l'ho fatta''.

Questo è fondamentale per chi lavora in programmazione, design o art, perché
ci saranno sempre deadline: bisogna sempre \textbf{lavorare sulle aspettative},
dire al proprio manager (o al proprio report) cosa è lecito aspettarsi. Non
bisogna fare come i pinguini di \emph{Madagascar} (``sorridete e salutate, va
tutto bene''): se c'è un problema, va detto, altrimenti si peggiora solo la
propria situazione.

Da italiani si tende a non voler mai alzare la mano per chiedere aiuto
(``no, no, ce la facciamo, lo faccio io''). A volte non ce la si può fare, non
perché non si è capaci, ma perché vari fattori lo impediscono. Accettarlo è
importante: \textbf{è molto meglio avvisare e chiedere una mano che dire
``non ce l'ho fatta''}.

\subsection{Domande sul dare e ricevere feedback}

\paragraph{``Ne parliamo a fare, se poi fai come vuoi?''} Uno studente (Cristian)
riceve feedback che, dopo averli analizzati, non ritiene congrui; quando non
li segue gli viene detto ``ne parliamo a fare?''. C'è un problema di
\textbf{aspettative}: l'idea che accettare un feedback voglia dire fare come
dice l'altro. Accettarlo significa prendere in considerazione ciò che è stato
detto. Chi decide diversamente ha però un \textbf{onere}: argomentare perché
si fa comunque a modo proprio (``mi hai detto verde, la facciamo gialla
perché\dots''). Poi bisogna solo aspettare i risultati per vedere chi aveva
ragione.

La decisione spetta sempre \textbf{a chi ha la responsabilità di quella
decisione}. Se il game director deve scegliere tra A e B, anche se il
feedback dice B, la scelta è sua perché la responsabilità è sua. Non sarebbe
corretto essere ritenuti responsabili di una scelta delegata a chi dà
feedback. La risposta può essere semplicemente: ``l'ho tenuto in
considerazione, ma visto che ne sono responsabile io e le conseguenze sono
mie, ho scelto di fare X''.

\paragraph{E se l'altro fa ripicca?} Un partecipante chiede di chi, vedendo
il proprio feedback non accolto, si vendica ignorando a sua volta feedback
buoni. Per il relatore è un problema caratteriale: chi si comporta così ``è
già sulla strada per uscire''. Non c'è spazio per l'ego in team di questo
tipo. Il fenomeno nasce dall'idea che \emph{il feedback negativo faccia fare
brutta figura}, che il fatto che si potesse fare meglio tolga qualcosa alle
proprie capacità.

Non è così: \textbf{il feedback non è un giudizio}. Si può passare un anno a
Riot ricevendo solo feedback negativi e avere una valutazione annuale
stellare, perché ricevere molti feedback negativi significa essere una persona
approcciabile, che agisce sul feedback e invita gli altri a darlo. Il relatore
ha visto valutazioni molto peggiori per chi riceveva solo feedback positivi,
perché la gente aveva paura di dargliene di negativi: ma solo ``pacche sulle
spalle'' non fanno migliorare. Il problema va affrontato parlando con la
persona; a volte le conversazioni iniziano con \textbf{``ti do un feedback sul
feedback''}, cioè sul modo in cui l'altro ha dato il suo.

Esempi dalla sua esperienza:
\begin{itemize}
  \item lui chiede spesso di \textbf{non ricevere feedback senza esempi
        tangibili}, perché fa fatica ad astrarre: ``sei stato maleducato'' non
        significa nulla, ``ieri al meeting mi hai risposto male su questa
        cosa'' sì;
  \item ad altri serve l'opposto: ha ricevuto il feedback ``quando devi dirmi
        che ho sbagliato, puoi girarci un po' intorno e darmi del contesto?'',
        perché il suo stile diretto (``questa cosa non ha funzionato, dovremmo
        fare così'') feriva. Ha adattato il suo modo di dare feedback: con
        alcuni è diretto, con altri gira intorno. L'importante è saperlo.
\end{itemize}

\subsection{Non portare la soluzione}

Un consiglio centrale: \textbf{quando si dà un feedback, non si dà una
soluzione}. Non si dice ``io farei così'', ma ``secondo me questa cosa si
potrebbe fare diversamente, per questo e questo motivo: come potremmo farla
diversamente?''. È come chi smette di fumare perché vuole e chi smette perché
glielo dicono moglie e figli: se non si è convinti, spesso non si smette.

Quando il relatore diceva ``c'è questo problema, io farei così'', il meeting
successivo iniziava inevitabilmente con ``ci ho provato come hai detto tu, con
me non ha funzionato''. Sempre. Perché la soluzione era sua. La
responsabilità della soluzione deve essere invece \textbf{congiunta}: ogni
feedback deve creare un piccolo team che ha come obiettivo risolvere quel
problema. Questo toglie l'antagonismo (non ``ti dico che non sai fare il tuo
lavoro'', ma ``ti voglio aiutare a farlo meglio, lavoriamoci insieme'') e
sposta il focus: il problema non sei tu, è ciò che sta accadendo.

\paragraph{E se sono i piani alti a tenere all'oscuro dei problemi?} È una
questione di cultura aziendale. Il relatore fa sempre presente al proprio capo
il concetto di \textbf{service level agreement} (che verrà ripreso nella
seconda giornata): ``se mi dai questi input, ti garantisco questo output in
questi tempi''. Se il capo chiede delle cose senza dire budget, tempi e
perché, non può aspettarsi lo stesso risultato. È una comunicazione a due
vie: si dice al manager cosa serve per fare il lavoro, e lui deve fornirlo o
accettare che il risultato non sarà quello voluto.

\subsection{Amici, parenti e parole proibite}

\paragraph{``Non si lavora con amici e parenti''?} Per sua definizione
personale il relatore ha pochissimi amici (``stanno sulle dita di una mano'')
e una soglia molto alta per considerare amico un collega. Allo stesso tempo,
in quindici anni e con centinaia di colleghi, ha avuto solo \textbf{due}
esperienze in cui non è riuscito a lavorare con qualcuno: non serve essere
amici per far funzionare le cose. Con gli amici, anzi, gli viene più facile:
se dice a un amico che una cosa ``fa schifo'', l'amico non pensa ``come si
permette'', ma ``se perfino Stefano, che è mio amico, me lo dice, allora è
proprio così''. Diventa un rafforzativo, grazie alla fiducia costruita prima.
Il problema del lavorare con amici e parenti nasce dal \emph{non voler} dire
all'amico che sta sbagliando; se invece la confidenza permette di dirlo, si
può saltare la ``confezione'' del messaggio e passare direttamente al
contenuto.

\paragraph{Le frasi peggiori.} Alla domanda se esistano parole peggiori di
``ottimo lavoro'', il relatore indica la sua classifica: ``ottimo lavoro''
(detto in modo vuoto)\footnote{Il passaggio della trascrizione è
frammentario: si intende verosimilmente il complimento generico e di
circostanza, che non dice nulla di utile.}; poi la frase che lo ``fa uscire di
testa'', \textbf{``abbiamo sempre fatto così''} (``e infatti avete chiuso'',
come Blockbuster); infine \textbf{``siamo una famiglia''}, riferito al posto di
lavoro. Un'azienda non è mai una famiglia, perché non ci si può aspettare lo
stesso impegno da chi la possiede e da chi ci lavora: non c'è un
corrispettivo proporzionale tra impegno e ritorno. Per quanto inclusiva,
flessibile e razionale, non sarà mai una famiglia.

\subsection{Ego, muri e il peggior modo di scegliere un leader}

\paragraph{Chi mette davanti l'ego e alza muri.} Esistono, ma di solito durano
pochissimo in azienda. Non si tratta di avere un ego forte, ma di \textbf{non
saper ricevere feedback}. Spesso sono persone bravissime nel loro lavoro che
poi si chiedono perché tutti intorno diventino più bravi di loro: perché gli
altri accettano feedback e hanno più fonti da cui attingere, mentre loro si
sono rintanati nella loro torre.

Molti manager tendono ad \textbf{assumere persone un filo peggiori di loro},
per non fare brutta figura. Il relatore assume solo persone più brave di lui,
``perché non è una gara'': ha lavorato con molti community manager più bravi
di lui come community manager, ma lui era più bravo a fare il manager. Non
viene valutato su quanto sa fare il community manager, ma su quanto il suo
team sa farlo. Un partecipante chiede se la paura di avere qualcuno di
migliore esista anche in ambito accademico. Assolutamente sì: al relatore è
capitato di non essere preso come docente in un corso di formazione sui
videogiochi perché il responsabile aveva percepito che ne sapesse più di
lui\footnote{La trascrizione dice ``che non sapevo più di lui'', ma il senso
del passaggio richiede il contrario.}. ``Ma perché te ne importa? Tu sei quello che ha portato chi
sa fare il suo lavoro: ben venga.'' O si ha questa cultura, o si ha paura per
il proprio posto.

\paragraph{Promuovere il migliore.} Quando bisogna scegliere un leader per un
team (programmazione, design, art), molto spesso si sceglie \textbf{il migliore
tecnico}: il miglior programmatore, il miglior designer. È spesso la scelta
sbagliata: è come prendere un idraulico per fare l'elettricista. Fare il
\emph{people manager} non è come fare il designer o il programmatore; si può
essere un ottimo programmatore e un pessimo manager, e viceversa. Le aziende
l'hanno capito e lo fanno meno di prima. A Blizzard c'era un grave problema di
\emph{middle management}: molte persone in ruoli importanti erano state
promosse cercando le skill sbagliate, erano frustrate dal non saper fare il
proprio lavoro e schiacciavano chi stava sotto per paura di fare brutta
figura.

Un partecipante cita una teoria simile: una persona viene promossa finché non
raggiunge \emph{il proprio livello di incompetenza}, dove si blocca\footnote{È
il \emph{principio di Peter} (Laurence J.\ Peter, 1969), citato in lezione
come ``la teoria di Peter''.}. Un altro cita Michael Scott di \emph{The
Office}, su cui il relatore tornerà (sezione~\ref{sec:z1-quando}).

% ---------------------------------------------------------------------
\section{Perché la cultura conta: retention, costi e fiducia}
\label{sec:z1-cultura}

\subsection{L'ambiente di lavoro e la retention}

Perché tutto questo è così importante? Perché l'ambiente di lavoro conta
moltissimo per la \textbf{retention}, cioè per trattenere le persone e farle
lavorare in modo efficiente. Se contasse solo lo stipendio, molti non
lascerebbero mai Riot o Blizzard; e più si va avanti con l'età, più si impara
a dare valore anche al resto.

Inoltre, secondo gli studi citati nelle slide, chi vive il proprio lavoro in
modo negativo è \textbf{meno efficiente nel problem solving}: l'energia
mentale va verso ``come faccio a non farmi dare la colpa'' o ``come faccio a
fare bella figura'', non verso la soluzione del problema. In ambienti negativi,
dove ci si deve preoccupare della propria ``sopravvivenza'', si lavora
peggio.

\subsection{Il costo di chi se ne va}

Ogni volta che qualcuno si dimette o viene licenziato, il costo di trovare
una nuova persona, fare i colloqui, formarla e farle attraversare forming,
storming, norming\dots\ è circa \textbf{un terzo del suo stipendio}. Per questo
le aziende tendono a fare più \emph{retention} che \emph{acquisition}, almeno
con le figure considerate utili alla produzione.

Allo stesso tempo si vedono sempre più spesso \emph{layoff} e \emph{hiring
wave}: si licenziano 3.000 persone e sei mesi dopo se ne riassumono 2.500 o
3.500, ma tutte più junior. Il relatore ne ha scritto in un articolo su
Eurogamer.it, in cui spiega anche alcuni ``trucchi del mestiere'': perché i
licenziamenti avvengono sempre in certi periodi e sempre ``a gruppi'' (prima
Bethesda, poi Blizzard, poi un'altra\dots).

\paragraph{L'offerta di Riot.} A Riot esisteva una pratica (di cui il relatore
non ricorda il nome) per cui, firmato il contratto, si aveva \textbf{un anno
di tempo per dire ``non mi piace, non me la sento, me ne vado''} e ricevere
subito come buonuscita un anno di stipendio. Avevano stimato che fosse molto
meglio pagare subito un anno a qualcuno e lasciarlo andare che tenerlo due,
tre, quattro anni come \emph{underperformer}: la responsabilità della scelta
di assumere era dell'azienda\footnote{La pratica, nota come ``The Offer'', è
descritta in fonti pubbliche con condizioni diverse (un periodo più breve e
una somma pari a una frazione dello stipendio annuo). Si riportano le
condizioni come raccontate in lezione; andrebbero verificate.}. L'hanno usata
pochissimi, ma il solo fatto che esistesse faceva lavorare meglio anche lui da
manager, perché sapeva che chi non si trovava bene poteva andarsene con un
anno di tempo per trovare un altro lavoro.

\paragraph{Il ricircolo.} C'è un forte ricircolo nell'industria: molti ex
colleghi passano da Ubisoft a Paradox, da Paradox a Funcom, da Funcom a
Blizzard. Il relatore ha scelto di uscirne lavorando solo per aziende
\textbf{indipendenti}: Stormind non è un indie di cinque persone, ma è
indipendente, e non deve accontentare azionisti. Quando vede un'azienda con
utili miliardari licenziare migliaia di persone si chiede ``dov'è la linea''
(\emph{when enough is enough}): con azionisti che vogliono un ritorno, se
quello è l'obiettivo stanno facendo la cosa giusta, ma lui non condivide
l'obiettivo\footnote{In lezione si parla di ``6,5 miliardi di utile'' e di
``3.000 persone licenziate''; le cifre sono riportate come dette e non sono
state verificate.}.

\paragraph{L'open space vuoto.} A Riot, durante un'ondata di licenziamenti,
del suo team di trenta persone è rimasto solo lui: ``teniamo te come
coordinatore, poi ricostruirai il team''. Ha rifiutato di ricostruire un team
cancellato ``con un colpo di spugna'' e ha deciso che non avrebbe più lavorato
per una grande azienda. Nessuno aveva pensato di spostarlo, e ha lavorato per
tre mesi in un open space di cento metri quadri con la sua scrivania e altre
ventinove vuote. Su una lavagna aveva scritto i nomi dei suoi ragazzi e li
cancellava man mano che, anche grazie ai suoi \emph{referral}, trovavano
lavoro. Cancellato l'ultimo nome, si è dimesso: per lui il lavoro era finito.

\paragraph{Team che si dissolvono.} Uno studente cita un articolo su Larian
Studios, il cui management è contrario a ricostruire ogni volta il team a fine
progetto. È lo stesso problema: si pensa ai ``giochi Blizzard'' o ai ``giochi
BioWare'', ma le persone che li hanno fatti non ci sono più; è rimasto solo il
brand, e non ha senso avere le stesse aspettative da un gioco fatto da persone
diverse. E anche chi sopravvive ai layoff vive pensando che il prossimo
potrebbe essere lui: alcune persone, licenziate nell'ultima ondata a
Blizzard, gli hanno detto ``erano due anni che lo sapevo''. Vivevano a Cork,
``un paesotto irlandese dove o vai per lavorare a Blizzard o non ci vai'',
avevano comprato casa: non ci si può aspettare chissà che morale.

\subsection{La cultura è forte quando si lavora per gli altri}

\begin{quote}
  \emph{La cultura è forte quando le persone lavorano per gli altri e a favore
  degli altri, e non per se stesse.}
\end{quote}

Se si dà feedback per far migliorare l'altro e non per imporre il proprio
modo, se si aiuta per il bene del progetto e non per fare bella figura, la
cultura si rafforza. È dai \textbf{piccoli gesti} che emerge la cultura
aziendale: ad esempio, dal non dover chiedere ``perché mi stai aiutando?''.

\paragraph{L'incognita.} \textbf{L'incognita è spesso dannosa quanto le
cattive intenzioni.} In un esperimento (``mi pare a Stoccolma, cinque o sei
anni fa'') si distribuivano per strada banconote da 500 euro, e la maggior
parte delle persone le rifiutava chiedendosi ``dov'è la fregatura?''
\footnote{L'esperimento è riportato in modo vago in lezione e non è stato
possibile identificarlo con certezza.}. Lo stesso accadeva al relatore,
arrivato a PG e a Stormind da Riot: a ``vedo che sei in difficoltà, posso
aiutarti?'' la risposta era sempre ``perché?''. Alla fine rispondeva: ``perché
quando sarò in difficoltà io vorrei che tu mi aiutassi''.

Quando si comincia a lavorare per gli altri e si smette di pensare ``così avrò
la promozione, così avrò l'aumento'', promozioni, aumenti e rispetto arrivano
da soli. A Blizzard aveva fallito tre volte i colloqui per essere promosso a
manager (inteso come coordinatore del team); quando ha deciso di smettere di
pensarci e di fare solo il bene del proprio lavoro, la volta dopo è stato
promosso. Funziona così: chi aiuta i colleghi e fa andare bene il progetto
viene notato (``questa persona fa migliorare cinque persone, non solo se
stessa: vale di più'').

\subsection{Il lavoro è valutato come opera d'insieme}

Il relatore riconosce che buona parte del suo successo è venuto dal fatto che
il suo team lavorava bene, non dalla sua bravura personale. Anche i progetti
degli studenti, presentati a Diego, a Enrico, ad Andrea, verranno valutati
\textbf{come opera d'insieme}. Nessuno dice ``il \emph{combat system} di questo
gioco l'ha fatto X, che bravo'': si dice ``BioWare ha fatto un bel gioco'',
``Blizzard ha fatto un bel gioco''.

Lo ha replicato anche lato HR. In molte aziende ci sono \textbf{bonus
personali} legati a obiettivi (se il gioco vende X, se la feature è pronta in
tempo). Il relatore li ha dimezzati: metà restano personali, metà diventano
\textbf{bonus di team}, che tutti ricevono se il team raggiunge gli obiettivi.
Un incentivo concreto a collaborare: non solo si dice che il lavoro di gruppo è
importante, lo si dimostra dandogli più valore economico. Aveva proposto il
100\% di bonus di team, ma gli è stato concesso il 50 e 50; ha funzionato
comunque, anche se sarebbe stato ancora più potente. Ovviamente, alla
domanda ``perché mi aiuti?'' si può anche rispondere ``perché voglio il
bonus'': un aiuto dato ``per il motivo sbagliato'' era un compromesso che era
disposto ad accettare.

\subsection{Il coraggio di chiedere aiuto}

Un partecipante osserva che aiutarsi dovrebbe essere la norma. Il modello
diffuso, invece, è ``devi essere forte e indipendente, non chiedere mai una
mano'', e vale in tutte le direzioni (dal maschilismo alla ``donna forte e
indipendente che non ha bisogno di nessuno''): \textbf{non ammettere di avere
bisogno di aiuto è sempre uno dei motivi per cui le aziende vanno male}, che si
tratti di uomini, donne o persone non binarie. Il problema è non avere la
sicurezza di poter alzare la mano.

Il relatore ha cercato di togliere ogni conseguenza dall'aver alzato la mano:
se si \emph{depenalizza} qualcosa, le persone non sono più disincentivate a
farla. Ha detto espressamente che non avrebbe mai valutato nessuno per il
numero di volte in cui chiedeva aiuto o per il numero di feedback negativi
ricevuti, e questo ha funzionato.

\subsection{Perché si dà un feedback}

\emph{Perché avete dato il vostro ultimo feedback? Che cosa volevate ottenere?}
Le risposte della classe:
\begin{itemize}
  \item per \textbf{riconoscere il valore} di un collega, anche verso il
        manager, perché qualcuno che non aveva visibilità vedesse quanto
        quella persona fosse efficace;
  \item per \textbf{spingere una persona a superarsi};
  \item per aiutare un collega a rendere il prodotto all'altezza di ciò che
        voleva rappresentare;
  \item per \textbf{genuinità}: c'era qualcosa che poteva essere fatto meglio.
\end{itemize}
Il relatore aggiunge che la forma cambia col tempo: all'inizio era ``devo
darti un feedback, hai dieci minuti?''; ora è un messaggio informale su
Discord, perché ha costruito un rapporto di fiducia e sa che non verrà letto
come ``l'HR manager mi scrive, sono nei guai?'', ma come ``Stefano è così, mi
dà feedback, resterà tra noi e lo risolveremo insieme''.

Un altro motivo è \textbf{risolvere incompatibilità caratteriali}: le
conversazioni più difficili, del tipo ``io e te non stiamo funzionando, cosa
possiamo fare diversamente?''. Sono importanti perché, se due responsabili non
vanno d'accordo, il team se ne accorge, comincia a raccontarsi storie (``si
odiano, ci saranno problemi'') e, per difendere il proprio manager, comincia a
giudicare male il lavoro dell'altro gruppo. È la dinamica da fermare:
\textbf{valutare gli altri solo per difendere la propria fazione}.

Dare e ricevere feedback sono \textbf{due muscoli da allenare} (il relatore è
sempre stato più bravo a darlo che a riceverlo), e si impara col tempo. È come
fare colloqui: fare colloqui è una skill a parte, che non indica quanto si è
bravi nel proprio lavoro. Essere scartati non significa non saper fare quel
lavoro: spesso significa aver fatto male il colloquio o che è stato preferito
qualcun altro.

\paragraph{Feedback pubblici e privati.} Uno studente racconta di aver dato un
feedback negativo a un superiore, senza un motivo particolare, per dimostrare
a sé e ai colleghi che non si viene licenziati per questo. Il relatore coglie
l'occasione per precisare che \textbf{i feedback negativi si danno
individualmente}: lui dà moltissimi feedback di cui il resto dell'azienda non
sa nulla. Reagirebbe molto male se, in un meeting, qualcuno alzasse la mano per
dire ``comunque questa cosa non l'hai fatta bene'': quello non è feedback, è
discussione, accettabile solo se è un brainstorming e se l'opinione è stata
chiesta. Alcune aziende hanno piattaforme pubbliche di feedback, ma spesso
servono solo ad applaudire; lo stesso accade nelle aziende che
\emph{obbligano} a lasciare un certo numero di feedback a settimana, che
diventano ``lui è bravo, lei è competente, anche questa settimana l'abbiamo
fatto''. Bisognerebbe chiedere a chi le adotta cosa pensa di ottenere (forse
trasparenza), ma difficilmente una persona ``attaccata'' pubblicamente reagisce
bene.

% ---------------------------------------------------------------------
\section{Among Us: la caccia al traditore}
\label{sec:z1-among-us}

Prima della pausa pranzo il relatore mostra tre film di fantascienza di
qualche decennio fa e chiede che cosa li accomuni\footnote{I titoli non sono
leggibili nella trascrizione: viene citato \emph{Event Horizon} da uno
studente, e il film ``in mezzo'' è indicato come ``l'originale'',
verosimilmente \emph{La cosa} (\emph{The Thing}).}. In ognuno c'è un gruppo di
persone con un obiettivo e la presenza, reale o percepita, di qualcuno che
rema contro: \textbf{l'energia del gruppo viene spesa per trovare il
responsabile}.

\paragraph{L'esperimento degli alieni.} Nei master in presenza il relatore
propone un esercizio. Distribuisce segretamente cartellini plastificati con
scritto ``essere umano'' o ``alieno''; le carte alieno contengono istruzioni di
comportamento con cui gli alieni, che cercano di fingersi umani, si
tradirebbero (dare sempre ragione a chi ha parlato prima, finire ogni frase
con una domanda, parlare in un certo modo). Lo scopo: in dieci minuti formare
gruppi, e \textbf{vince il gruppo più numeroso senza alieni}; basta un alieno
per perdere. Si formano sempre gruppi piccoli (due, cinque, sei persone),
perché la tendenza è ``mi fido al massimo di voi due''.

Ma vincono i gruppi più grandi, per un motivo semplice: il relatore
\textbf{non distribuisce nessuna carta alieno}. Mostra a tutti entrambe le
carte, ma consegna solo carte ``umano''. Vuole simulare ciò che accade in
ufficio: la tendenza a percepire sempre qualcuno che rema contro (i manager,
l'altro team, quel collega), quando in realtà non c'è nessuno, ``è tutto nella
nostra testa''. Quando chiede le motivazioni, emergono sospetti
\emph{fabbricati}: ``faceva sempre domande'', ``parlava solo all'infinito''
(``no, parla così e basta''). Ci si fabbrica da soli i motivi per non fidarsi
degli altri, ed è comunissimo.

\paragraph{Among Us.} L'autosabotaggio è, per tornare ad \emph{Among Us},
\textbf{concentrarsi sul cercare cadaveri e traditori invece di fare le
task}: non si vince perché non si sono completate le task, mentre si cercava il
colpevole. Succede tantissimo negli uffici, ed è il motivo per cui l'80\% dei
problemi che il relatore risolveva erano di comunicazione. Entrando nei team
sentiva immancabilmente: ``facciamo il possibile, ma questa lotta continua con
l'altro gruppo è estenuante''. ``Perché c'è questo problema?'' ``Non lo so,
c'era già quando sono arrivato.'' Un paradosso: c'è un problema, nessuno sa
dire quale sia, ma c'è sempre stato e quindi lo si accetta.

\paragraph{L'esperimento delle scimmie.} È l'esempio, ``finto o reale'', delle
scimmie in gabbia\footnote{È il noto racconto delle ``cinque scimmie'', spesso
presentato come esperimento scientifico ma probabilmente apocrifo (lo stesso
relatore lo definisce ``finto barra reale''); nella trascrizione il numero
delle scimmie varia.}. In una gabbia con alcune scimmie pende una banana
elettrificata: chi la tocca fa prendere la scossa a tutte. Le scimmie imparano
e, quando una si avvicina alla banana, le altre la picchiano. Poi si
sostituisce una scimmia alla volta con una nuova: la nuova va verso la banana,
le altre la picchiano, e smette. Si continua finché nella gabbia ci sono solo
scimmie nuove, poi si toglie la corrente. Nessuna scimmia tocca la banana:
nessuna ha mai preso la scossa, ma tutte hanno imparato che chi tocca la
banana prende le botte.

Nei team è lo stesso: il motivo per cui si ``odia'' l'altro team si
trasforma, ma un vero motivo non c'è. Ne derivano la sfiducia, il non alzare
mai la mano per chiedere aiuto, il non dire mai ``bel lavoro'' agli altri
perché non sia mai che qualcuno ti veda fare un complimento ``a quelli''.
Risolvere queste dinamiche è lungo e tedioso, e sempre pieno di ``non mi
ricordo'', ``io non c'ero'' e ``abbiamo sempre fatto così'', la cosa più
frustrante per chi cerca di risolverle.

% ---------------------------------------------------------------------
\section{Il principio del default to trust}
\label{sec:z1-default-trust}

\subsection{La definizione}

Il relatore indica questa come la slide più importante della presentazione.
\textbf{Default to trust} è uno dei valori aziendali di Riot Games, che lui ha
tradotto, adattato e ``esportato'':

\begin{quote}
  \emph{Approccio ogni situazione di conflitto partendo dal presupposto che una
  persona abbia agito in un certo modo perché, con le informazioni a sua
  disposizione in quel momento, ha pensato che fosse la soluzione migliore, e
  non perché sia un'incapace o volesse danneggiare il progetto
  intenzionalmente.}
\end{quote}

In altre parole: \textbf{di default, mi fido}.

\subsection{L'esempio del ritardo al daily stand-up}

A Riot il relatore aveva un \emph{daily stand-up} dal vivo alle 15 con il suo
team di venti persone. Una ragazza arrivava regolarmente alle 15:15, per una,
due, tre settimane. Alla quarta ci ha parlato. \emph{Come iniziereste la
conversazione?} Le proposte della classe:
\begin{itemize}
  \item sottolineare l'importanza di esserci dall'inizio, perché potrebbe
        perdersi qualcosa che le interessa, e per rispetto dei colleghi
        presenti;
  \item ``devi arrivare in orario'', spiegando perché è importante;
  \item \textbf{chiedere perché}.
\end{itemize}

Il relatore osserva, anche dalle sue consulenze fuori dai videogiochi, che
\textbf{più si sale con la seniority}, più la risposta è ``devi smettere di
arrivare in ritardo'', spesso con un \emph{or else} finale (``altrimenti
succede quest'altra cosa''). Più si scende con l'età, e quindi con meno
abitudini consolidate, più la prima domanda diventa ``dimmi perché arrivi in
ritardo''. Già spiegare perché è importante è ``miglia avanti'' rispetto al
semplice ``smetti''.

Partire dal perché gli ha permesso di scoprire la causa: il meeting precedente
della ragazza, in un'altra zona del campus a un quarto d'ora di distanza,
finiva sempre tardi, e lei doveva attraversare il campus di corsa. Non aveva
la confidenza per dire all'altro manager, che insisteva perché si restasse
fino alla fine, che doveva andare; sacrificava quindi il meeting del relatore,
sapendo che era il più comprensivo dei due.

Il relatore si è offerto di \textbf{mediare}: ha detto all'altro manager che
o anticipava l'inizio del suo meeting, così da finire comunque in tempo,
oppure lasciava uscire la ragazza, perché quel meeting gli serviva. C'è anche
una questione di \textbf{economia di scala}, spesso sottovalutata: se venti
persone aspettano qualcuno per dieci minuti, si sono sprecati dieci minuti per
venti. Per lo stesso motivo il workshop è iniziato alle 10:05: iniziare alle
10:10 avrebbe significato far aspettare cinquanta persone cinque minuti in
più. A Riot c'era una scrivania vuota con un pupazzo gonfiabile, \textbf{Mr.\
Late}: sommando tutti i minuti di ritardo a meeting, call e conversazioni, in
un anno veniva fuori praticamente \emph{una persona a tempo pieno}. Serviva a
ricordare che anche cinque minuti contano, e che bisogna rispettare gli altri
arrivando in orario o almeno avvisando.

\paragraph{Gli effetti.} La conversazione era avvenuta nel primissimo periodo a
Riot, quando il team non aveva ancora motivo di fidarsi di lui. Risolvendo il
problema ha eliminato anche tutto lo strato di micro-dinamiche spiacevoli che
un manager impara a vedere (due colleghi che si guardano con la risatina
quando lei entra in ritardo). La ragazza ha cominciato a entrare tranquilla,
con una certa leggerezza: non le faceva piacere arrivare in ritardo (``ci sono
anche persone che lo fanno apposta, e quello è un problema da affrontare
diversamente''). Soprattutto, con i colleghi non parlava di lui come del capo
che minacciava, ma come di chi le aveva risolto un problema: era diventato
un \textbf{alleato}. La settimana dopo altre tre persone sono andate da lui
con un problema da risolvere. Ripetuto cinque, sei, sette, dodici volte, in un
paio di mesi poteva dire quasi qualsiasi cosa nei meeting sapendo che
richieste e critiche sarebbero state percepite come ``me lo chiede perché vuole
darmi una mano'', non per fare bella figura o per presentare i successi del
team ai suoi superiori.

\subsection{Condividere i successi, schermare i fallimenti}

Un altro elemento fondamentale (ripreso nella seconda giornata): il manager
\textbf{condivide i successi e protegge il team dai fallimenti}. ``Se va male
qualcosa è colpa mia; se va bene qualcosa, voglio che ne parliate voi.'' Il
lavoro del manager non dovrebbe essere valutato su ciò che presenta del lavoro
del team: tutti sanno comunque chi è il manager. Tanto vale lasciar presentare
il team, che potrà dire con orgoglio di aver presentato a Mike
Morhaime\footnote{Cofondatore ed ex presidente di Blizzard; nella
trascrizione ``Mike Morayme''.} la propria campagna europea per
\emph{StarCraft~II}. Ci sono manager che dicono ``come mi prendo le colpe, mi
prendo i meriti'', ma i meriti li hanno già: è inutile ostentarli. Con i
fallimenti invece il manager è la \textbf{rete di sicurezza} del team. Il
relatore ne ha fatto una questione di principio: \textbf{non scaricare mai la
colpa su qualcun altro}.

\paragraph{Colpevole o causa.} Un esempio che il relatore attribuisce a Pirate
Software, creatore di contenuti ed ex dipendente Blizzard\footnote{Il relatore
aggiunge che si tratterebbe del figlio dello sviluppatore a cui si sarebbe
ispirato il personaggio del giocatore ``senza vita'' nell'episodio di
\emph{South Park} dedicato a \emph{World of Warcraft}. L'informazione è
riportata come detta in lezione e non è verificata.}. Contattando il servizio
clienti Blizzard per motivi personali, aveva scoperto che i colleghi gli
avevano scritto nei ticket informazioni confidenziali che non avrebbero dovuto
condividere in quel modo, nemmeno con un collega. Non ha segnalato le persone
da punire: ha detto ``c'è un problema, evidentemente le policy non sono
chiare''. Quando gli hanno chiesto i nomi, ha rifiutato: ``se vi dico chi è,
punite quelli, e gli altri cinquecento continueranno a farlo; non avremo
cambiato la causa alla fonte''.

\begin{quote}
  \emph{È inutile trovare il colpevole. Voglio trovare la causa. Perché il
  colpevole cambia.}
\end{quote}

Il colpevole a volte è vittima della situazione: per questo la definizione
dice \emph{``con le informazioni a sua disposizione in quel momento''}. Se il
relatore avesse rimproverato la ragazza in ritardo, non solo non avrebbe
risolto il problema, ma avrebbe creato un'\textbf{ingiustizia}: lei, che
attraversava il campus di corsa per una colpa non sua e di cui non poteva
parlare, sarebbe stata anche rimproverata.

Nessuno fa un brutto lavoro apposta: nessuno è contento quando sbaglia o
perde. Perché al lavoro diamo per scontato che gli altri sbaglino
intenzionalmente? Se qualcuno ha fatto una cosa male, è perché in quel momento
credeva fosse la scelta migliore. La domanda giusta non è ``perché hai fatto
questa cosa che è stata un fallimento?'', ma:

\begin{quote}
  \emph{Che cosa ti ha fatto pensare che questa fosse la scelta giusta?}
\end{quote}

Magari erano limiti di budget, di tempo, di scopo: si agisce su quelli, non
sulla persona. Se si cambia la persona ma le si danno gli stessi fattori, è
molto probabile che arrivi alle stesse conclusioni, e ci si troverà a dover
``disciplinare'' anche lei.

\paragraph{La reputazione costruita sui rumor.} Durante il giro di
presentazione in un ufficio (non in nessuna delle aziende citate), al relatore
è stato detto: ``questo è il dipartimento di design; se ti serve qualcosa
chiedi a lei, non a lui, perché lui consegna sempre in ritardo''. Oltre a
trovarlo poco professionale, una settimana dopo ha chiesto: ``cosa ti ha fatto
questa persona?'' ``A me niente, però ho sentito che\dots''. Si torna alle
scimmie: magari quella persona non consegna in ritardo da mesi, ma ormai lo
dicono tutti. È un circolo vizioso da fermare subito.

\paragraph{Niente escalation inutili.} Lo stesso vale per il feedback: quando
qualcuno va dal relatore dicendo ``ho un problema con questa persona'', lui
chiede: \textbf{``perché ne stai parlando a me e non a lei?''}. ``Perché sei il
mio capo'' non basta. Passare da lui significa dare per scontato che l'altra
persona non abbia interesse a risolvere, mentre magari basta parlarle. È la
strada più facile (``lo dico a Stefano, ci pensa Stefano''), e proprio per
questo bisogna mettere dei paletti: il manager interviene quando serve, ma le
cose che si possono gestire da soli vanno gestite da soli. (Quando invece il
confronto diretto fallisce, l'escalation è legittima: si veda la
sezione~\ref{sec:z1-qa-finale}.)

% ---------------------------------------------------------------------
\section{Traiettorie di carriera e valutazione delle performance}
\label{sec:z1-traiettorie}

\subsection{Quale traiettoria vorreste?}

Il relatore mostra un grafico con quattro curve colorate, contrassegnate
anche da lettere (A, B, C, D) per eventuali persone daltoniche (``a me è
capitato che qualcuno scoprisse di esserlo durante una lezione''). Le curve
rappresentano la performance lavorativa nel tempo. Due domande: \emph{quale
vorreste essere? E quale pensate realisticamente di essere?} Un partecipante
aspira a B e pensa di essere C; quasi tutti scelgono B.

Perché non A, che sale più rapidamente? Un partecipante risponde che una
crescita così repentina può essere più dannosa che utile, come una ``rush'' di
studio molto faticosa che poi non si riesce a sostenere. Il relatore propone
una lettura suggerita da un amico psicologo, che spiega la forma del grafico:
di B, C e D \textbf{vediamo dove finiscono}; di A no. Potrebbe continuare a
salire oltre la scala del grafico, o crollare sotto D. Non saperlo ci spinge a
preferire B. (A chi obietta che anche le altre proseguono oltre il bordo, il
relatore precisa: il bordo destro è ``oggi'', e di A non sappiamo come sia
arrivata fin lì.)

\subsection{La carriera non è una linea retta}

Tutti immaginano la carriera come una linea: si parte da X e, attraverso
stage e posizioni, si arriva a Y. In realtà carriera e performance avranno un
andamento irregolare, con alti e bassi. Come non si accetta che un team possa
temporaneamente \emph{underperformare}, tanto meno lo si accetta da se stessi.
Bisogna accettare che ci saranno momenti bassi.

Per questo il relatore \textbf{detesta le valutazioni annuali delle
performance}: prendono un punto casuale della traiettoria, che se si è
fortunati è il picco più alto e se si è sfortunati il più basso, e lo usano
per valutare un intero anno. Insegna ai manager a fare più punti di controllo
e a guardare altre cose:
\begin{enumerate}
  \item \textbf{quanto tempo passa tra un picco e l'altro}, cioè quanto
        \emph{cooldown} serve a una persona. È il motivo per cui, dopo la
        Games Week, tutti staccavano una o due settimane. Il relatore stesso,
        che si impegna moltissimo e tende al burnout, non potrebbe rifare lo
        stesso workshop il lunedì e il martedì successivi con la stessa
        qualità;
  \item \textbf{la differenza tra la miglior performance e la peggiore}: che
        cosa ha portato a lavorare così bene, e così male? Con chi si
        lavorava, su cosa ci si concentrava, quanto tempo e budget c'erano?
        L'obiettivo è capire come \emph{replicare} le condizioni che fanno
        lavorare bene, invece del solito ``hai lavorato male, è colpa tua''.
        Il relatore ha avuto molti manager per cui quando lavorava male era
        colpa sua e quando lavorava bene ``che bravo'', ma sapeva che raramente
        era l'unico fattore della propria performance. È la fallacia del ``se
        sai fare 100, perché non fai sempre 100?'', come nella scena di
        \emph{Schindler's List} sull'operaio che produce una vite in pochi
        secondi: non sempre un ritmo è sostenibile;
  \item \textbf{il fallimento di oggi confrontato con quello di sei mesi fa}.
        Il relatore non è contento della presentazione di oggi (l'avrebbe
        migliorata per altri due anni), ma rispetto a quelle di cinque anni fa
        vede un grande miglioramento di chiarezza: il suo peggior fallimento di
        oggi equivale al suo primo successo di inizio carriera.
\end{enumerate}
Quello che conta è la \textbf{traiettoria} dall'inizio dell'anno a oggi, non
i singoli punti. E se si riesce a capire insieme quando e perché si ha
successo e quando e perché si fallisce, si può lavorare per replicare solo le
situazioni che fanno lavorare bene.

\paragraph{Lo stesso vale per i giochi.} È il motivo principale per cui i
giochi delle grandi case vanno bene o male. Quando si riescono a replicare le
condizioni che hanno prodotto \emph{Warcraft~II}, esce \emph{Warcraft~III};
quando non ci si riesce, esce ``il nuovo \emph{Dragon Age}'': la gente non è
arrabbiata perché il gioco non è bello, ma perché non è come gli altri
\emph{Dragon Age}; non ne replica il \emph{feeling}, anche se potrebbe
essere un gioco favoloso. L'eccezione positiva è il passaggio da \emph{Dune} a
\emph{Dune~II}: il primo era di fatto un'avventura grafica, il secondo ha
generato il genere degli strategici in tempo reale. Gli sviluppatori avevano
capito che il successo non stava nel tipo di gioco, ma nell'ambientazione che
attirava le persone: un bel gioco in quell'ambientazione sarebbe piaciuto
comunque. Il relatore comprò \emph{Dune~II} aspettandosi un seguito
dell'avventura e, pur sorpreso, lo finì.

\subsection{Lying, hiding, faking}

Quando si è valutati su un singolo punto nel tempo, e da quel punto dipendono
bonus, stipendio e il rischio di licenziamento, si comincia a fare
\textbf{lying, hiding e faking}:
\begin{description}
  \item[Lying] dire che è sempre andato tutto bene, che si finirà sempre in
        orario, che non c'è nessun problema;
  \item[Hiding] nascondere: ``non era un mio task'', ``non c'ero'', ``non ho
        visto niente'';
  \item[Faking] fingere, anche in cose apparentemente innocue come ``sì, sì,
        ho capito tutto'' quando si ha solo paura di dire ``me lo rispieghi?
        Non ho capito niente''.
\end{description}

\subsection{``Non ho capito'' e la conferma delle istruzioni}

Dire ``non ho capito, ripetimelo'' è ancora visto come una figura da sciocco.
Il relatore non si arrabbierà mai se qualcuno gli chiede di rispiegare
un'intera slide: le cose ovvie per lui possono non esserlo per altri. Più si
conosce qualcosa, \textbf{più è difficile spiegarla}, perché si danno per
scontate molte cose. Ne parlava proprio con Diego, che gli diceva che fare
l'insegnante diventa sempre più difficile man mano che si diventa esperti. È
come spiegare \emph{EVE Online}, uno dei giochi più complessi in circolazione:
chi lo conosce bene dà per scontato che ``ovviamente non puoi fare questa
cosa'', mentre per un nuovo giocatore non è affatto ovvio. A scuola chi dice
``non ho capito'' viene visto come lo ``scemo''; nessuno si chiede ``non è che
l'ho spiegato male?''.

Il racconto più ``traumatico'' dei ragazzi che fanno stage nelle grandi
aziende è sempre lo stesso: ``Cinzia, devi fare questa cosa.'' Cinzia la fa.
``Non te l'avevo chiesta così.'' Il problema è che nessuno ha
\textbf{confermato che Cinzia avesse capito}. Il consiglio, valido sempre:
\begin{itemize}
  \item quando si \emph{ricevono} istruzioni, ripeterle: ``ho capito bene?
        Vuoi che faccia questa cosa in questo modo?'';
  \item quando si \emph{danno} istruzioni, chiedere a chi le riceve di
        ripeterle, non perché lo si ritenga incapace, ma per capire che cosa
        ha capito dalla propria spiegazione (``quindi vuoi un quadrato?''
        ``No, ho detto un triangolo\dots'' ``Però mi hai detto una forma con
        quattro angoli.'' ``Hai ragione: voglio un quadrato'').
\end{itemize}
Frasi come ``non te lo ripeto due volte'' o ``stammi dietro'' portano solo a
fare le cose come si pensa che vadano fatte e a sentirsi dire a fine progetto
``non la volevo così'': che si tratti di una feature, di un design o di un
modello 3D. (Un partecipante osserva che lo stesso accade insegnando le lingue:
chi insegna dà per scontato, ad esempio, che si sappia dell'esistenza del
genere neutro in tedesco.)

\paragraph{Come convincere il team a chiedere?} Il relatore mette sempre come
premessa che \textbf{l'obiettivo non è trovare la persona più intelligente
nella stanza, ma fare un buon lavoro}. Che Gianni capisca alla prima e un altro
alla quinta non gli interessa, se alla fine entrambi lavorano bene: ripetere
cinque volte è comunque meno lavoro che far rifare le cose sette volte. È un
investimento a lungo termine, e non dissimula alcun fastidio: davvero non gli
dà fastidio. Ma ripetere le stesse parole non aiuta: ``non ho sentito'' si
risolve ripetendo; ``non ho capito'' si risolve cambiando spiegazione. Meglio
ancora: ``dimmi che cosa hai capito tu, e vediamo cosa non torna''. Se
intravedesse nel team uno sguardo del tipo ``è la quinta volta che lo
chiede'', interverrebbe: spesso chi giudica è proprio chi dice ``ho capito'' e
poi sbaglia. Preferisce di gran lunga chi chiede cinque volte.

\subsection{Il feedback più utile che abbiate ricevuto}

\emph{Qual è il feedback più utile che abbiate mai ricevuto?} Il relatore ne ha
almeno dieci che gli hanno cambiato il modo di lavorare. Le risposte:
\begin{itemize}
  \item un semplice \textbf{``grazie''} per un lavoro fatto, più utile di un
        ``questa cosa fa schifo'';
  \item un feedback sulla propria \textbf{impulsività} nel rispondere: fermarsi
        un attimo e capire di cosa si sta parlando prima di intervenire;
  \item \textbf{essere se stessi}, anche al lavoro, e non scusarsi
        continuamente per ogni cosa.
\end{itemize}
I feedback che toccano di più sono quelli \textbf{veri}. Il ``grazie''
significa ``ho visto quello che hai fatto e l'ho apprezzato''. Quelli negativi
che si ricordano sono quelli che hanno fatto realizzare ``questa cosa che
pensavo di saper fare, non la so fare'', oppure ``la so fare, ma questa volta
non mi è venuta'' (non è che se una cosa una volta non riesce non la si sa più
fare: può capitare di sbagliare una porta). Arrivano sempre da un
\textbf{punto cieco}: qualcosa che non si sapeva fosse percepito dagli altri.
Conta più ciò che le persone ti dicono di ciò che pensi che pensino di te; e
ciò che il relatore pensa che gli altri pensino di lui è sempre terribile,
mentre il loro giudizio si rivela molto più positivo del previsto.

\paragraph{Il denominatore comune siete voi.} Si potrebbe pensare che questi
concetti non siano applicabili da chi arriva in un'azienda senza poterne
cambiare la cultura. Ma c'è sempre un denominatore comune a cui applicarli:
\emph{se stessi}. La cultura aziendale può non spingere a dare feedback, può
vedere male chi chiede aiuto; ma il relatore ha visto che quando lui lo fa, le
persone intorno tendono a farlo, e un gruppo che lo fa lavora meglio. Perfino
i CEO più ostici a un certo punto gli chiedevano: ``perché il tuo team ora
performa meglio degli altri?''. La differenza è il default to trust: una
pratica che col tempo fa operare il team \emph{come un team} e non come un
insieme di individui. Un videogioco è un progetto corale: c'è chi fa giochi da
solo, ma non si fa da soli \emph{Red Dead Redemption~2}; e se il gruppo non è
affiatato e allineato non esce \emph{Red Dead Redemption~2}, esce
\emph{Anthem} (``oggi la mia fissazione è \emph{Anthem}, domani sarà
un'altra'').

% ---------------------------------------------------------------------
\section{Come dare un feedback}
\label{sec:z1-dare}

\subsection{Concentrarsi sull'azione e sui suoi effetti}

Il primo principio: \textbf{il feedback non riguarda mai la persona}. Non si
dice ``tu hai fatto questa cosa che non va bene'', ma ``questa cosa non va
bene: come facciamo a non farla più?''. Che sia un comportamento, un'azione,
una decisione o una strategia, ci si concentra su quello, e si
\textbf{esplicita l'effetto} che ha. Nell'esempio del ritardo l'effetto era
``quindici minuti di attesa moltiplicati per venti persone''. In un progetto
può essere una catena:
\begin{quote}
  se non rispetti le deadline, l'animazione che viene dopo resta ferma; se
  loro sono in ritardo, la \emph{vertical slice} è in ritardo; se la vertical
  slice non è pronta per la Gamescom, non andiamo alla Gamescom e non troviamo
  un publisher; se non troviamo un publisher, chiudiamo.
\end{quote}
Più la conseguenza è grave, più dovrebbe sembrare piccolo il gesto di dire
``datemi una mano'': è il \emph{male minore}. Poi, a cose fatte, ci si dà
feedback: cosa ha funzionato, cosa no, cosa fare di più, di meno, cosa smettere
di fare.

\subsection{Il pericolo degli ambienti solo positivi}

Gli ambienti dove ci si fa sempre grandi complimenti e nessuno parla dei
problemi creano difficoltà vere alle aziende. Il problema non è mai la
mancanza di complimenti, ma un \textbf{eccesso di ottimismo}: si comincia a
pensare che le cose andranno come si spera, perché tutti intorno danno solo
feedback positivi.

A Blizzard i feedback determinavano gli \textbf{aumenti annuali}: ognuno
doveva indicare cinque persone che avrebbero dato feedback su di lui, e
ovviamente tutti indicavano gli amici, o colleghi di progetti andati bene. Un
sistema nato con le migliori intenzioni diventava uno \emph{scambio di voti}
(``se mi metti follow ti metto follow''), senza alcun miglioramento.

Al contrario, \textbf{le conversazioni più brutte sono state le più utili}:
affrontare i licenziamenti a Riot, o la cancellazione di \emph{Project Titan}
a Blizzard, il progetto che doveva essere il nuovo MMO dopo \emph{World of
Warcraft} e da cui poi è nato \emph{Overwatch}\footnote{Nella trascrizione si
legge anche ``doveva essere Warcraft 2'', probabilmente un lapsus o un errore
di trascrizione.}. Andando a vedere perché lo si cancellava, c'erano solo
motivazioni razionali e pragmatiche, di business, chiare a tutti, non un ``Mike
Morhaime ci odia''. Sono anche le conversazioni in cui ha lavorato su se
stesso: è molto \textbf{avverso ai conflitti}, e questo a volte gli ha
procurato conseguenze molto più gravi di quanto sarebbe costato un piccolo
disagio. Dieci anni fa faceva fatica ad accettare questo difetto; ha dovuto
subire le conseguenze di conflitti evitati e poi cresciuti.

\subsection{Che cosa fare}

Le linee guida della slide, che il relatore legge rapidamente perché saranno
disponibili nel materiale:
\begin{description}
  \item[Attenzione al linguaggio del corpo] dare un feedback a braccia
        incrociate o guardando altrove è diverso dal guardare l'altro negli
        occhi.
  \item[Sii specifico] con esempi tangibili.
  \item[Esplicita gli effetti] come visto sopra.
  \item[Discuti il prima possibile] non aspettare un mese per parlare di
        qualcosa successo ieri, altrimenti nessuno ricorderà più bene cosa sia
        accaduto.
  \item[Concentrati sull'azione e sul problema, non sulla persona] il nemico è
        l'azione, non la persona che si ha davanti. Non bisogna far cambiare
        l'altro: bisogna evitare che il problema si ripresenti.
\end{description}

\paragraph{Gli insulti sono rivolti al ruolo.} Un esempio solo apparentemente
scollegato. Il relatore ha lavorato su \emph{League of Legends}, ``non
esattamente noto per avere una community da orsetti del cuore'', in varie
\emph{task force} sulla tossicità. A chi smetteva di giocare veniva inviato un
sondaggio sul perché; il 33\% indicava la tossicità della community. Dopo varie
campagne di sensibilizzazione e cambiamenti si è scesi al 30\%: sembra solo
un 3\% in meno, ma valeva circa \textbf{120 milioni di revenue in più
all'anno}, perché quei giocatori restavano esposti a skin e contenuti e
continuavano a comprare. Anche qui, piccoli guadagni che si sommano.

Lavorando con degli psicologi, aveva notato che gli insulti \textbf{non sono
quasi mai rivolti alla persona, ma al ruolo}: ``il mio jungler fa schifo'',
``il mio DPS non è capace'', ``il nostro portiere è una pippa''. Chiamare
qualcuno per ruolo lo \emph{depersonalizza} e rende più facile insultarlo. Lo
stesso nell'esempio del giro dell'ufficio: non ``non chiedere a Giulio'', ma
``non chiedere a quel grafico, consegna sempre in ritardo''. Concentrarsi
sull'azione, senza attaccare la persona, fa sì che il feedback non sia (e non
venga percepito come) un attacco personale.

\paragraph{Considera il carattere.} Qui tornano utili i test di personalità.
Servono più agli studenti che al relatore (che sa già che usciranno 53 profili
diversi, senza un pattern): bisogna conoscere il proprio profilo e
\textbf{condividerlo con chi lavora con noi}, chiedendo anche agli altri di
fare lo stesso. In ufficio il relatore ha usato una versione più avanzata, a
pagamento, con un report di circa trenta pagine: il \textbf{profilo
DISC}\footnote{Nella trascrizione ``disk''. Si tratta del modello DISC
(Dominance, Influence, Steadiness, Conscientiousness).}. Ha mappato tutta
l'azienda e ha fatto stampare il profilo di ciascuno sulla sua scrivania, così
che tutti sapessero con chi stavano parlando. Capire come ragionano gli altri
(``è introverso, non gliene parlo davanti a tutti'') è una delle cose che gli
sono tornate più utili: \emph{adattare ciò che si vuole dire a chi si ha
davanti}.

\paragraph{Assicurati di non disturbare.} Sotto il profilo di ognuno c'era un
piccolo \textbf{semaforo}: verde, ``parlami pure, interrompimi''; giallo,
``interrompimi solo se è urgente''; rosso, ``mi sto concentrando, non mi
parlare''.

\paragraph{Lascia finire il lavoro.} Durante gli eventi dal vivo il relatore si
era reso conto di quanto fosse grave dare feedback mentre qualcuno sta ancora
lavorando. Se una cosa sta andando male, dirlo in quel momento non aiuta: tra
la pressione di finire e quella del feedback, peggiorano entrambe le cose, e
anche il rapporto. Meglio lasciarla finire, anche male, e parlarne dopo.

\subsection{Che cosa non fare}

\begin{description}
  \item[Non basarti su supposizioni] ``mi hanno detto che ieri in un meeting
        hai fatto\dots'': o c'eri e l'hai visto, o non esiste.
  \item[Non giudicare la persona.]
  \item[Non fornire feedback vago.]
  \item[Non accettare e non fornire feedback di terze parti] Se qualcuno vi
        dicesse ``mi hanno detto che hai fatto un brutto lavoro'', il primo
        pensiero non sarebbe ``cosa ho fatto male?'', ma \emph{``chi te l'ha
        detto? Chi è la vipera che ho in seno? Chi è l'impostore?''}. Siamo
        fatti così. Ogni conversazione che inizia con ``mi hanno detto che'' è
        inutile: l'altro non ascolterà, penserà solo a chi dare la caccia, e
        spesso penserà alla persona sbagliata, creando antipatie tra persone
        che non c'entravano nulla. Quando al relatore dicono ``mi hanno detto
        che lui ha fatto\dots'', risponde ``vai a parlarne con lui''; se gli
        dicono ``mi hanno detto che tu hai fatto\dots'', risponde ``mandami la
        persona che te l'ha detto, altrimenti non ne parliamo''. Non per
        cattiveria: se chiedesse ``che cosa ho detto di preciso a Marco?'' e la
        risposta fosse ``non lo so, non c'ero'', di cosa si starebbe parlando?
        \textbf{Parlate con la fonte}, altrimenti fate come se quel feedback non
        esistesse.
  \item[Non dare consigli non richiesti] come visto, non portare la soluzione:
        la soluzione deve venire da chi riceve il feedback.
  \item[``Non mi piace'' non è un feedback] ``Penso che questo design non
        funzioni perché, visto che il nostro obiettivo è\dots'' ha argomenti;
        ``non mi piace il combat system'' è un'opinione. Anche ``io l'avrei
        fatto diversamente'' non è un feedback.
  \item[Accettare non vuol dire fare come proponi] deve essere chiaro a chi dà
        il feedback: l'aspettativa è essere ascoltati e presi in
        considerazione, ma se la responsabilità è dell'altro, deciderà
        l'altro.
  \item[Non offenderti] se il feedback non viene seguito. Restarci male
        significa pensare di saperne più dell'altro, cioè \emph{venire meno al
        default to trust}. ``Decidi tu'' significa ``mi fido di te, perché sei
        lo specialista, e mi fido che mi ascolterai e farai la cosa migliore
        per tutti'', che potrebbe anche non essere ciò che ho detto io. Se
        voglio che l'altro faccia sempre come dico io, mi sto fidando solo di
        me stesso.
\end{description}

\subsection{Il primo feedback e il debriefing}

La cosa più difficile è dare il \textbf{primo} feedback: è un muscolo da
allenare. Chi non l'ha mai fatto (soprattutto con questa struttura) dovrebbe
provarci, anche su minuzie (``abbiamo concordato di usare WhatsApp nel gruppo:
non potremmo usare Telegram?''), anche solo per vedere come va. Altre
indicazioni:
\begin{itemize}
  \item feedback \textbf{solo uno a uno}, mai davanti ad altri, soprattutto se
        negativo;
  \item \textbf{prepararsi} prima;
  \item fare il \textbf{debriefing del feedback}: dopo, fermarsi a pensare a
        com'è andato. Il relatore, quando media un feedback tra due persone,
        si sente spesso rispondere ``benissimo'' da una e ``malissimo''
        dall'altra: la percezione è molto diversa. Non si fa problemi a tornare
        dalla stessa persona: ``ieri ti ho dato un feedback: com'è andata? Sei
        soddisfatto? Come poteva andare meglio?''. Con le persone con cui si ha
        meno a che fare, questo aiuta molto la volta successiva.
\end{itemize}

% ---------------------------------------------------------------------
\section{Come ricevere un feedback}
\label{sec:z1-ricevere}

Per il relatore è la parte più difficile. Si riassume in quattro passi:
\begin{enumerate}
  \item \textbf{Ascolta}: si parte sempre ascoltando.
  \item \textbf{Chiedi dettagli}, se serve, ma per capire meglio, non per
        difenderti.
  \item \textbf{Accetta, a prescindere}: non si sta decidendo chi ha ragione;
        si sta ascoltando un'opinione. Non c'è nulla da cui difendersi.
  \item \textbf{Agisci}: fatti un piano. Può essere banale (``mi hanno detto
        che interrompo sempre: cerco di non interrompere''), oppure ``devo
        migliorare questa cosa: come?''. Pensa anche a \emph{chi} potrebbe
        aiutarti e a \emph{che supporto} ti serve: il piano non deve essere
        per forza ``devo risolvere tutto da solo''.
\end{enumerate}
Una domanda che molti sottovalutano: \emph{``sono junior in un'azienda: sarò
mica il primo ad avere questo problema?''}. Probabilmente qualcun altro l'ha
già avuto: basta chiedere ``tu come hai fatto?''.

\subsection{Il feedback positivo}

Tutto questo vale per i feedback costruttivi negativi, a cui si pensa di
default; ma esistono anche i \textbf{feedback positivi}, e vanno usati. Quando
si riceve un ``grazie'' o un ``bravo'', la domanda è: \emph{come faccio a
ricreare le stesse condizioni?}

Il relatore ha lavorato anche in aziende dove il feedback era solo negativo.
Il rischio? Si immagini l'impiegato ideale: vuole sempre migliorare, è sempre
pronto a dire ``magari sbaglio io''. Se nessuno gli dice cosa sta facendo
\emph{bene}, spinto dal desiderio di migliorare rischia di cambiare cose che
erano già buone, e di farle peggio. Il rinforzo positivo \textbf{non è solo una
gentilezza}: è un atto pratico, che evita che le persone peggiorino ciò che
sanno già fare bene.

\subsection{Il feedback accidentale}

Esiste infine il \textbf{feedback accidentale}, molto frequente negli uffici:
\begin{itemize}
  \item il ``chiedi a lei, non a quel grafico'' del giro dell'ufficio è un
        feedback negativo su una persona, dato accidentalmente, né richiesto
        né voluto;
  \item a Riot li chiamavano \textbf{drive-by}, come nelle scene dei film in
        cui passa un'auto, si abbassa il finestrino, si spara e si riparte. In
        molte aziende con una ``cultura del feedback'' in cui tutti devono
        dire la loro, si riceve una valanga di commenti del tipo ``comunque il
        sito a me non piace, ciao''. Non si sa perché, non si sa cosa fare di
        diverso, non si sa che rischi ci siano: \emph{a che cosa è servita
        questa conversazione?} A niente.
\end{itemize}

% ---------------------------------------------------------------------
\section{Quando, quanto e con chi}
\label{sec:z1-quando}

\subsection{Frequentemente e recentemente}

\emph{Quando avete ricevuto il vostro ultimo feedback?} Per gli studenti,
qualche giorno fa (nel corso) o qualche mese fa (al lavoro, dopo una
consegna). Come aveva osservato Silvia, il feedback andrebbe dato
spesso.

Il relatore usa l'esempio di una nave da crociera (``la Costa Concordia'') che
deve andare da A a B. La rotta ideale (azzurra) è quasi impossibile: nessun
progetto va esattamente come previsto. C'è la rotta ``panoramica'' (verde), che
gira intorno all'isola d'Elba; ci sono poi rotte (arancione e rosa) che a un
certo punto si allontanano molto e poi correggono bruscamente: qualcuno si è
accorto che si stava sbagliando, ``ma non è quello che avevo chiesto'',
``non è dove dovevamo andare''. Sono i famosi \textbf{sanity check} o
\emph{checkpoint}.

L'ideale è avere \textbf{tanti piccoli check per fare tante micro-correzioni}:
correggere la rotta di un paio di gradi è molto più rapido che farlo dopo
essersi allontanati. Bastano trenta secondi di ``ehi, vado bene?''. Spesso non
lo si fa per paura di sembrare insicuri o poco autonomi, ma per il relatore è
l'opposto: quando qualcuno gli chiede ``vado bene così?'' e lui risponde
``sì'', si alleggerisce anche del proprio carico mentale. Un capo che assegna
un compito e lo vede solo alla fine vive sempre col dubbio (``avranno capito?
Lo staranno facendo giusto? Arriverà in tempo?''); con un check di dieci
secondi al giorno (``tutto ok, Stefano'') è tranquillo, e la persona va avanti
sapendo che quanto fatto finora va bene. Non serve fare ogni giorno il giro del
dipartimento: anche ogni uno o due giorni basta a evitare la rotta verde.

Per il feedback è lo stesso: \textbf{va dato subito}. Dire oggi ``ieri hai
fatto questa cosa'' è molto diverso da dire ``negli ultimi sei mesi hai fatto
questa cosa''. Il relatore si è sentito dire più volte, alla valutazione di
dicembre, ``da maggio hai fatto questa cosa che non andava''. La sua risposta:
``e aspetti dicembre per dirmelo? Mi hai lasciato sbagliare per sei mesi?''. È
una mancanza di rispetto, e la conversazione non può più essere costruttiva.
Per questo non ama le valutazioni annuali: i manager crescono abituandosi a
prendere appunti delle cose negative per tirarle fuori alla valutazione
(``a marzo hai fatto questa cosa'' ``sì, però poi ho smesso'' ``sì, però a
marzo l'hai fatta'').

In sintesi:
\begin{description}
  \item[Frequentemente] riguarda l'intervallo tra un check e l'altro.
  \item[Recentemente] significa non lasciar passare un anno prima di dare il
        feedback, altrimenti non serve più a niente.
\end{description}

\subsection{Quanto dura un feedback}

Un feedback dura \textbf{da trenta secondi a quindici minuti}; raramente di
più. Chi non ne ha mai dato uno dovrebbe sforzarsi di darne almeno uno, magari
positivo per cominciare. Ogni feedback dato o ricevuto è un
\textbf{deposito nel proprio ``fondo fiduciario''} (\emph{trust fund}), che
costruisce la reputazione futura. Il relatore si è costruito la reputazione di
chi dà feedback il più possibile oggettivi e analitici: giusti o sbagliati che
siano, costringono a pensare. Non impone mai il suo punto di vista (``si fa
come dico io''), ma chiede ``tu come faresti? Come posso aiutarti?''. Il
feedback, per lui, serve a \emph{far crescere qualcun altro}.

\subsection{I manager peggiori: ``a che numero sto pensando?''}

I peggiori manager che il relatore abbia avuto erano quelli che volevano che si
facesse come volevano loro, anche quando chiedevano ``tu come la faresti?''.
In realtà la domanda era ``\emph{a che numero sto pensando?}'': se dicevi 5,
``no, dovremmo fare 9''; se dicevi 9, ``no, 5''. Lo stress costante di dover
indovinare cosa vuole il capo nasce dalla mancanza di fiducia (``se mi chiedi
come fare, fammelo fare'').

Lo stesso gli succede con i CEO che lo chiamano per ``sistemare la cultura
aziendale'' e poi dicono ``non sono sicuro, io farei diversamente''. La sua
risposta: ``lo so che faresti diversamente, per questo mi hai chiamato''. Se si
chiama un esperto, bisogna lasciargli fare il suo lavoro: non si va dal
meccanico a cambiare le gomme per poi dirgli ``non dovresti girare un po' di
più quel bullone?''. Chiedere un parere a un esperto significa dargli anche il
potere di scegliere. Il problema è il \emph{micromanagement}, o in generale il
non accettare punti di vista diversi dal proprio (magari perché si è il
fondatore e si pensa di essere l'unico a saper fare le cose).

\subsection{Dare feedback a un superiore egocentrico}

\paragraph{Come si dà un feedback negativo a un superiore con un grande ego?}
Il trucco è \textbf{non puntare sull'ego, ma sull'azione e sui suoi effetti}.
Non ``tu fai così e io mi sento così'': anzi, si alimenta il suo ego
rendendolo il \emph{main character} dell'avventura, l'eroe che risolve il
problema (``tu, grande manager, aiutaci a risolvere questo problema che
affligge noi che facciamo l'operatività''). Non è l'ideale (``io cambierei
manager''), ma pochissimi possono scegliersi i colleghi: il relatore stesso ha
quasi sempre ricevuto team preassemblati (``ecco le chiavi, questa è la tua
macchina''). Bisogna giocare con le carte che si hanno in mano.

\paragraph{Un esempio pratico.} Con un manager che fa micromanagement e mette
sempre ``il becco'' nel proprio lavoro, si può dire: ``se ti fidi di me per
questo lavoro, con le tue capacità potresti dedicarti a questa cosa più
strategica e di lungo termine: pensa quanto sarebbe più utile per l'azienda''.
In sostanza: ``guarda cosa potresti raggiungere se mi lasciassi fare il mio
lavoro''. Non come un ``tu fai il tuo e io faccio il mio'', ma come un
``ti sollevo da questa parte del carico, che so che ritieni importante: lascia
che ci pensi io, così puoi dedicarti ad altro'', che magari gli dà anche più
visibilità.

Si gioca su ciò che quella persona pensa di saper fare, e che probabilmente sa
anche fare. Anche Michael Scott di \emph{The Office}, per quanto preso in giro,
è raccontato come un venditore bravissimo, che ha chiuso i contratti più
importanti per l'azienda, e rispetto ai manager che gli succedono è forse
``uno dei meno peggiori''. Tutti, anche i manager peggiori, hanno punti di
forza: \textbf{giocate su quelli}.

% ---------------------------------------------------------------------
\section{Il feedback come regalo}
\label{sec:z1-regalo}

Per chiudere, il relatore risponde alla domanda d'apertura, che aveva lasciato
apposta in sospeso: che cos'è il feedback per lui?

\begin{quote}
  \emph{Il feedback è un regalo che mi viene fatto da qualcuno che so che vuole
  il mio interesse.}
\end{quote}

Ha costruito un team di persone che si fidano di lui e che non hanno motivo di
sabotarlo o di ferirlo: per questo interpreta anche le cose peggiori con il
default to trust. È un regalo, e cerca di prenderlo in considerazione il più
possibile.

Ci è arrivato partendo da se stesso. Oltre a essere avverso ai conflitti, si
definisce ``avverso alla comunicazione umana'': paradossalmente, una persona
introversa e timida finita a fare comunicazione, marketing e ora HR (``succede
spesso''). Prima di ordinare una pizza al telefono ha dieci minuti di
preparazione mentale, e se l'ordine delle domande cambia va in confusione. Per
lui un feedback di dieci minuti sono spesso \textbf{tre ore di preparazione}:
costruisce un \textbf{albero dei dialoghi}, come in un'avventura grafica
personale, con tutte le opzioni. Poi spesso finisce con ``hai fatto questa
cosa, andrebbe migliorata'' ``hai ragione, grazie'', e basta.

Così ha pensato: chissà quanti si fanno gli stessi problemi con lui, e passano
ore a prepararsi per dirgli ``questa cosa l'hai fatta male'', sapendo che
riceveranno comunque quell'occhiata fugace del ``come ti permetti''. Eppure
vengono a dirglielo. Non può far altro che rispondere: ``grazie, hai ragione,
cercherò di fare meglio''. È così che ha metabolizzato il ricevere feedback:
capendo che non c'è nessuna cattiva intenzione e che, se non proprio felice
(``ci ho messo una pietra sopra''), può almeno essere \emph{grato} che ci sia
ancora qualcuno che ha così a cuore il suo interesse da dargli feedback.

\paragraph{Il silenzio fa più paura.} Non bisogna spaventarsi dei feedback
negativi che emergono dai sondaggi aziendali: bisogna spaventarsi del
\textbf{silenzio}, quando nessuno dice più nulla. Significa che nessuno si
fida più che si sappiano risolvere i problemi.

\paragraph{La sindrome dell'impostore.} Il suo motivatore personale, all'inizio,
era soprattutto ``non voglio fallire'': per questo cercava di migliorare, e
migliorare è diventato accettare feedback, perché più feedback riceve, più
migliora. Il feedback è un \textbf{acceleratore}, e lui cerca di incentivare
gli altri a darglielo. Resta però sempre la vocina del ``e se si accorgono che
non sei capace?''. Con la \emph{impostor syndrome} convive da anni e
continuerà a conviverci; se ce l'avete, conviveteci anche voi. Avrete sempre
la sensazione che ``se ne accorgeranno'': forse sì, ma sapete fare anche
altre cose. Non concentratevi solo su ciò che non sapete fare. Gli errori
capitano: la perfezione non è di questo mondo.

\paragraph{Esercitarsi in accademia.} Per fare pratica con il feedback,
l'ambiente migliore è quello dei colleghi di corso: persone nella stessa
situazione, con esperienza simile e lo stesso obiettivo, in un contesto in cui
non si può essere licenziati e l'unico effetto a lungo termine è ``ho parlato
con questa persona, ora le cose vanno meglio''. E bisogna avere il coraggio di
parlare anche con i docenti (``dire a Diego: secondo me questa lezione potevi
farla in questo modo''; ``so che è una persona che lo accetta''), senza farsi
bloccare dalla gerarchia. Nell'azienda giusta (come quella in cui lavora ora) il
relatore va sempre dal CEO a dare feedback, sapendo che non verrà licenziato:
il CEO gli ha detto ``ti ho preso per darmi feedback''.

La frase di chiusura:

\begin{quote}
  \emph{Per superare le vostre difficoltà, tutto quello che vi viene chiesto è
  il coraggio di chiedere aiuto.}
\end{quote}

Chiedere aiuto non è una debolezza: è un segno di consapevolezza di sé
(\emph{self-awareness}) e di forza personale. Se tutti lo facessero, i
progetti andrebbero molto meglio; spesso invece ci si fa frenare dal ``non
voglio fare brutta figura'', ``non voglio passare per quello scemo'', e sono
proprio queste paure a rallentarci e a impedirci di raggiungere gli obiettivi.

% ---------------------------------------------------------------------
\section{Domande e risposte finali}
\label{sec:z1-qa-finale}

L'ultima mezz'ora è dedicata a domande libere.

\paragraph{E se una persona peggiora per uno o due anni di fila?} Se nessuno
gliel'ha mai detto, c'è un grosso problema, non suo ma dell'azienda. Se invece
c'è consapevolezza, il percorso è: ``cosa vuoi fare per risolvere?''; ``ho
provato questo, questo e questo''; ``visto che non ha funzionato, vuoi un mio
consiglio? Proviamo con quello?''; ``che supporto ti serve?''. Se, con tutto il
supporto, il tempo e le soluzioni concordate per mesi, il miglioramento non
arriva, per il relatore la soluzione è: \textbf{``dimmi un prezzo per andare
via''}. Preferisce pagare una buonuscita (\emph{severance package}) che
avviare la strada dei richiami formali sulla performance fino al
licenziamento per giusta causa. A un certo punto andare avanti non è un favore
per nessuno.

\paragraph{È vero che nei videogiochi si viene sfruttati? Meglio un piccolo
studio italiano?} Per il relatore si è quasi più a rischio a Blizzard che in un
piccolo studio di Roma o Torino. Nel piccolo studio si sa dove ci si va a
cacciare: non ci sarà il team building a Santo Domingo, ma ci si concentra su
un gioco alla volta e, se va bene, si fa il successivo. A Blizzard si può
lavorare su \emph{Overwatch}, che vende milioni, e perdere comunque il posto,
in modo totalmente scollegato dal proprio lavoro. Gli stipendi più alti
all'estero compensano anche questo rischio, che spesso non dipende da come va
il gioco (anzi, ``il gioco non ha fatto come volevamo'' è spesso una scusa per
licenziamenti che non avevano motivo).

\paragraph{Cambiare settore a trent'anni, da web developer full remote?} Lo
stipendio indicato dallo studente non è fuori dalla portata di vari studi
italiani, anzi è superabile. Non va esclusa la possibilità di lavorare da
remoto per aziende piccole ma estere. Alcuni amici del relatore, pagati con
stipendi italiani, si sono trasferiti a Bali o in paesi con un costo della vita
più basso; il collega di \emph{PUBG} percepisce uno stipendio olandese vivendo
a Milano, che per quanto cara non costa come i Paesi Bassi.

\paragraph{Come si diventa manager?} Dipende dall'azienda: alcune promuovono i
migliori (processo tendenzialmente sbagliato), altre propongono un percorso di
\emph{people management}. Diventare manager significa \textbf{abbandonare ciò
che si sa fare per imparare a fare altro}. Il relatore ha dovuto accettare che
non sarebbe più stato ``il miglior community manager italiano del mondo''
(``eravamo in pochi, quindi era facile''): dedicandosi alle persone non poteva
più fare il suo lavoro bene come prima. È stata una negoziazione con se
stesso: fare il manager significa non essere più l'esperto di qualcosa, ma
saper gestire e far crescere persone. Significa anche \textbf{cedere il
controllo della propria carriera}: se il team decidesse di non fare più nulla,
il manager ne pagherebbe le conseguenze. Il suo lavoro diventa scegliere le
persone giuste e farle lavorare bene. Lui lo esplicita sempre: sa di affidare
al team parte della propria carriera, e il team sa che lui combatterà per
loro, anche prendendosi le colpe.

\paragraph{L'inglese.} \textbf{Va assolutamente imparato}, anche lavorando in
Italia. Non saperlo preclude circa il 95\% del mercato (tutto quello fuori
dall'Italia) e soprattutto il \emph{networking} con persone che potrebbero dare
occasioni future: l'inglese è spesso l'unica lingua in comune con chi lavora
all'estero ed è oggi la lingua universale dei videogiochi.
\begin{itemize}
  \item \emph{Che livello serve?} Non importa B2 o C2: bisogna saper sostenere
        una conversazione, guardare almeno Netflix con sottotitoli in inglese,
        comunicare senza dodici pause e senza ``come si dice\dots''. Il
        certificato non verrà mai chiesto: verrà chiesto se si sa parlare
        inglese. Il relatore ha sentito persone con un C2 che non sapevano
        parlarlo.
  \item \emph{Come allenarlo senza occasioni di conversazione?} Netflix, Prime
        Video, Now TV, Udemy, Coursera, YouTube, Duolingo, siti di scambio
        linguistico con madrelingua, Discord: non sono mai esistiti così tanti
        mezzi gratuiti. Il relatore ha imparato l'inglese giocando (con una
        pronuncia sbagliata, perché leggeva senza sentire, su cui ha poi
        dovuto lavorare); guardare un film in inglese con i sottotitoli in
        italiano insegna anche la traduzione.
  \item Oltre all'ascolto passivo è fondamentale \textbf{parlare}: in una
        flotta di \emph{EVE Online} con 150 persone quasi tutte madrelingua
        inglese si è esposti a decine di accenti e modi di dire (mille modi di
        dire ``come stai'' oltre a \emph{how are you}); lavorando all'estero si
        impara a distinguere gli accenti (Manchester, Londra) e ad apprezzarli
        anche nei media.
\end{itemize}

\paragraph{Quanto deve sapere un junior?} Dipende da quanto è tecnico il
ruolo. Un game designer junior che non sa qualcosa può essere affiancato da un
senior; un programmatore che non sa programmare, in quel momento, non è utile
all'azienda.

\paragraph{Annunci ``junior con cinque anni di esperienza''.} ``Cercasi
ventenne con trent'anni di esperienza.'' Non è l'annuncio per un junior: è un
errore dell'azienda, che in realtà cerca un \emph{mid}. Riceverà profili
junior, si stupirà che non abbiano cinque anni di esperienza, rifarà l'annuncio
e avrà perso sei mesi. Il consiglio: tenere d'occhio quella pagina, perché
dopo qualche mese l'annuncio verrà corretto.

\paragraph{Una struttura per il post-mortem.} A chi chiede un formato di base,
il relatore anticipa quanto vedrà nella seconda giornata a proposito di
brainstorming: una lavagna \textbf{divisa in quattro aree}, su cui ognuno
attacca tutti i post-it che vuole:
\begin{center}
  \begin{tabular}{@{}p{0.44\textwidth}p{0.44\textwidth}@{}}
    \toprule
    \textbf{Smettere di fare} & \textbf{Iniziare a fare} \\
    ciò che non ha funzionato & ciò che non abbiamo fatto e andava fatto \\
    \midrule
    \textbf{Fare di meno} & \textbf{Fare di più} \\
    ciò che ha fatto perdere tempo rispetto all'impatto ottenuto &
    ciò che ha funzionato ma, fatto con troppa leggerezza, non ha dato i
    frutti sperati \\
    \bottomrule
  \end{tabular}
\end{center}
Con i post-it ognuno ha voce e tutto risulta catalogato; per il report basta
una foto della lavagna (o uno screenshot, se la lavagna è digitale, come Miro).

\paragraph{Come sbloccare la paralisi del feedback verso il management?} Uno
studente ha visto manager dire ``non abbiate paura, la mia porta è sempre
aperta'' senza che cambiasse nulla. C'è una dissonanza tra ciò che dicono e
ciò che fanno: è ``la mamma che dice vieni qua che non ti faccio niente''.
Qualcosa deve aver fatto pensare alle persone che non è sicuro. Prima va
sistemato quello; poi si può provare con i \textbf{feedback anonimi}.
L'obiettivo del relatore resta il feedback aperto e firmato, ma a Stormind ha
iniziato con feedback anonimi verso il management per verificare se
l'anonimato fosse un fattore. Non lo era (evidentemente non c'erano problemi a
parlare con il CEO), e la volta successiva intende abbandonarli. Ma è un buon
\emph{check}: se senza anonimato non arriva nulla e con l'anonimato arrivano
cento feedback, il problema è chiaro.

\paragraph{Quando è giusto rivolgersi a un superiore?} La strada giusta resta
parlare direttamente con l'interessato. Ma se si prova e non si trova una
soluzione, la cosa più ovvia è andare dal proprio manager insieme: ``ci
aiuti?''. Quando il relatore parlava di mediare i dialoghi, era proprio questo:
fare da moderatore, aiutare a sbloccare i punti fermi o, se non si riesce a
scegliere tra A e B, scegliere lui. L'escalation, quando necessaria, va fatta.

\paragraph{Il tempo libero veniva rispettato?} No, e se l'è dovuto imporre
lui. Il lavoro gli piaceva così tanto che continuare a casa non gli sembrava
lavorare. Gli si è ritorto contro con il lancio di \emph{Diablo~III}, fatto in
orario americano (le due o l'una di notte in Europa), con l'ufficio chiuso
perché a nessuno era stato chiesto di restare. Il lancio va malissimo, i
giocatori protestano, e lui, sveglio, risponde sui forum fino alle tre, quattro,
cinque del mattino: ``non posso risolvere il problema, ma resto qui ad
aggiornarvi con quello che riesco a sapere dal team di sviluppo''. Il giorno
dopo trova sulla scrivania un \textbf{richiamo formale dell'HR}: aver lavorato
fuori orario senza autorizzazione, condotta non professionale. Per lui è stato
un interruttore: ``ora me ne vado''. Due settimane dopo faceva colloqui con
Riot. Ancora oggi è contrario a quel richiamo, ma gli ha fatto capire come
ragiona un'azienda: ``hai un contratto 9--18, non mi interessa se l'hai fatto
con le migliori intenzioni''. Eppure per anni ha goduto di un'ottima
reputazione in quelle community come ``quello che è rimasto con noi fino alle
cinque del mattino'', e anche l'azienda ne ha beneficiato.

Ci sono poi aziende in cui, soprattutto per i developer, viene richiesto il
famigerato \textbf{crunch}: lavoro straordinario continuativo, senza pause.
Dipende molto dall'azienda; sta venendo (``grazie al cielo'') deprecato, perché
molte aziende capiscono che non è sano ed è visto male anche dal pubblico, ma
esiste ancora.

\paragraph{E se il tuo capo demotiva qualcuno del tuo team?} Qui ``diventa
personale'': il team del relatore si fida di lui perché sa che in questi casi
``diventa una bestia''. ``Te la puoi prendere con me quanto vuoi, ma il mio
team non lo tocchi.'' Lo gestisce direttamente con l'altro manager, per quanto
brutta possa diventare la conversazione. Lo sforzo più grande è \textbf{non
andarne poi a parlare con il team} (``ragazzi, ho fatto questa cosa''), perché
minerebbe la fiducia verso l'alto: non vuole un rapporto ottimo fino a lui e
pessimo oltre. L'obiettivo è che l'altro manager ``si faccia perdonare''
cambiando il proprio comportamento, e che il team percepisca da solo il
cambiamento.

\paragraph{Chiusura della giornata.} Il relatore chiede un feedback collettivo
in chat (``se domani ci dai \emph{more of the same} siamo contenti, oppure
ditemi cosa limare''), si scusa per i ripetuti crash della connessione e
ringrazia per la proattività nelle domande, che gli servono anche come
\emph{check}: se nessuno chiede nulla su un argomento, significa che può
andare avanti. Racconta di aver passato una settimana a decidere che cosa
mettere nelle slide e di averle poi preparate in poco tempo (rifacendone
alcune fino alle tre di notte). Ricorda infine di usare il test di
personalità per sé e per il proprio gruppo, facendolo fare a chi non l'ha
ancora fatto: tornerà molto utile per cooperare meglio.

% ---------------------------------------------------------------------
\section{Anticipazioni sulla seconda giornata}
\label{sec:z1-domani}

La seconda giornata del workshop sarà più pratica: ``oggi vi ho spiegato
perché, domani vi insegno come lo fanno queste aziende''. Gli argomenti
previsti:
\begin{description}
  \item[Time management] come organizzarsi una giornata professionale in
        azienda: come strutturare una \emph{task list}, cosa fare e cosa non
        fare, come e cosa prioritizzare.
  \item[Brainstorming] i gruppi dovranno fare moltissimi meeting (come facciamo
        il gioco, come lo chiamiamo, che scopo gli diamo). Se non si sanno fare
        bene, i brainstorming sono \textbf{dannosi ed eterni}: un brainstorming
        dovrebbe durare \textbf{al massimo due ore}. Verranno mostrati metodi
        che accorciano la durata, massimizzano l'effetto e, soprattutto,
        \textbf{fanno partecipare tutti}: non c'è niente di peggio di un
        brainstorming in cui parlano in due per il 90\% del tempo. Il classico
        ``ragazzi, non esistono idee sbagliate, sparate'' è il modo peggiore di
        farlo; non è colpa di nessuno, ma si può fare diversamente.
  \item[Guidare da dentro] come essere leader nel proprio team anche senza
        esserne il leader: sviluppare responsabilità personale e verso gli
        altri. Non essere il manager di qualcuno non significa non esserne
        responsabili: la persona accanto è responsabilità di tutti, perché si
        è un team e non un insieme di individui. ``Se qualcuno cade, tutti lo
        facciamo rialzare: non c'è l'addetto al rialzamento.'' È la versione
        \emph{advanced} della giornata: come comportarsi da manager senza
        averne il titolo.
\end{description}
Torneranno inoltre il \emph{service level agreement} e il tema della
condivisione dei successi e della protezione del team dai fallimenti.

% ---------------------------------------------------------------------
\section*{Punti chiave del capitolo}
\addcontentsline{toc}{section}{Punti chiave del capitolo}

\begin{itemize}
  \item ``La capacità di un gruppo di persone di fare cose straordinarie
        dipende dal modo in cui queste persone operano insieme.'' Il 90\% dei
        problemi dei team sono \textbf{problemi di comunicazione}.
  \item Le aziende si liberano sempre più presto dello \textbf{useless
        genius}: le competenze si trasmettono, il carattere no. Il
        \textbf{fit culturale} va valutato anche dal candidato.
  \item Non esiste un unico modo di dare feedback: dipende dal contesto e
        soprattutto da \textbf{chi lo riceve}. Il feedback è una somma di
        piccoli miglioramenti periodici; tenersi le cose nel cassetto non
        serve mai.
  \item I team attraversano \textbf{forming, storming, norming, performing e
        adjourning}. Lo storming non si salta; si passa la maggior parte del
        tempo in forming; la bravura sta nell'accelerare i passaggi e nel
        fare bene il \textbf{post-mortem}, anche dei successi.
  \item Gli obiettivi di un progetto vanno fissati \textbf{all'inizio}: solo
        così si può valutare se è andato bene e scegliere tra A e B.
  \item Il circolo virtuoso è \textbf{feedback $\rightarrow$ miglioramento
        $\rightarrow$ fiducia}. Accettare un feedback non significa fare come
        dice l'altro: decide chi ha la responsabilità, ma deve argomentare.
  \item ``Aiutami a poterti aiutare'': meglio dire in anticipo ``potrei non
        farcela'' che ``non ce l'ho fatta''. Chiedere aiuto non è una
        debolezza.
  \item Quando si dà feedback \textbf{non si porta la soluzione}: la si
        costruisce insieme a chi lo riceve.
  \item La cultura è forte quando si lavora \textbf{per gli altri}; il lavoro
        è valutato come opera d'insieme (anche con bonus di team).
  \item Come in \emph{Among Us}, i team sprecano energie a cercare il
        traditore invece di fare le task. \textbf{Default to trust}: chi ha
        sbagliato ha agito con le informazioni che aveva. Cercare
        \textbf{la causa, non il colpevole}: ``che cosa ti ha fatto pensare
        che fosse la scelta giusta?''.
  \item Carriere e performance non sono lineari: vanno valutate le
        \textbf{traiettorie}, non un punto dell'anno. Le valutazioni puntuali
        generano \emph{lying, hiding, faking}.
  \item Confermare sempre le istruzioni facendosele ripetere; ``non ho
        capito'' non è una colpa.
  \item Dare feedback: sull'\textbf{azione} e sui suoi \textbf{effetti}, specifico,
        presto, uno a uno, adattato al carattere, senza terze parti, senza
        ``non mi piace'', senza offendersi se non viene seguito; poi fare il
        debriefing.
  \item Ricevere feedback: \textbf{ascolta, chiedi dettagli, accetta, agisci}.
        Il feedback positivo serve a non peggiorare ciò che si fa bene; il
        feedback accidentale e il \emph{drive-by} non servono a nulla.
  \item Feedback \textbf{frequenti e recenti}: tanti piccoli \emph{sanity
        check} costano meno di una grande correzione di rotta.
  \item Il feedback è \textbf{un regalo} di chi ha a cuore il nostro
        interesse; fa paura il silenzio, non le critiche.
\end{itemize}

%    % =====================================================================
%  Lezione zero 2 — Gestire il tempo, generare idee, guidare da dentro
%  Fonte: trascrizione della seconda "lezione zero" (GD0_2), seconda e
%  ultima giornata del workshop di due giorni tenuto prima dell'inizio del
%  corso di Game Design, dedicato alla figura professionale e al lavoro in
%  team. Prosegue direttamente il capitolo 0.1 (cultura del feedback).
%  Pacchetti richiesti nel preambolo del libro:
%    \usepackage[italian]{babel}, \usepackage[T1]{fontenc},
%    \usepackage[utf8]{inputenc} (solo pdfLaTeX), \usepackage{booktabs}
%
%  Collocazione nel libro: parte "Lezioni zero", subito dopo
%  GD0_1_cultura_del_feedback, con la stessa numerazione propria (0.2).
%  Nel file principale del libro, ad esempio:
%
%    \part{Lezioni zero: il game designer nell'industria}
%    \renewcommand{\thechapter}{0.\arabic{chapter}}
%    \setcounter{chapter}{0}
%    \input{Seminari/GD0_1_cultura_del_feedback}
%    \input{Seminari/GD0_2_tempo_idee_e_leadership}
%    \part{Il corso di Game Design}
%    \renewcommand{\thechapter}{\arabic{chapter}}
%    \setcounter{chapter}{0}
%    \input{GD1/GD1_introduzione_al_game_design}
%    % ...
%
%  (Con hyperref aggiungere anche \renewcommand{\theHchapter}{z.\arabic{chapter}}
%  nella parte zero, per evitare ancore duplicate.)
%  Tutte le etichette di questo capitolo hanno il prefisso "z2:"; i
%  rimandi al capitolo precedente usano le etichette "z1:" già definite in
%  GD0_1_cultura_del_feedback.tex.
% =====================================================================

\chapter[Tempo, idee e leadership]{Gestire il tempo, generare idee, guidare da dentro}
\label{cap:z2-tempo-idee-leadership}

Questo capitolo raccoglie la \textbf{seconda giornata} del workshop di due
giorni con cui si aprono le lezioni zero, tenuto dallo stesso relatore del
capitolo~\ref{cap:z1-feedback}, Stefano Cozzi\footnote{Per il cognome e per
il nome dello studio (Stormind Games) valgono le avvertenze della nota
all'inizio del capitolo~\ref{cap:z1-feedback}. In questa giornata la
trascrizione riporta soltanto il nome di battesimo, ``Stefano'', con cui il
relatore si cita più volte nei suoi esempi.}. Se la prima giornata spiegava
\emph{perché} le aziende tengono tanto alla cultura del feedback, la seconda
è dichiaratamente più pratica: ``ieri vi ho spiegato perché, oggi vi insegno
come lo fanno queste aziende''. I concetti della prima giornata (il
\emph{default to trust}, il ``potrei non farcela'' detto in anticipo, la
condivisione dei successi, il feedback frequente) ritornano di continuo,
applicati alla gestione del tempo, alle riunioni creative e alla leadership.

La giornata si è svolta di domenica, a chiusura di un fine settimana di
workshop\footnote{Il relatore dice ``oggi chiaramente è domenica'' e saluta
la classe con ``buon resto del weekend''; la prima giornata si era quindi
svolta di sabato. Più avanti parla di ``dodici ore'' di workshop in totale,
mentre nel capitolo~\ref{cap:z1-feedback} si parlava di circa dieci ore: la
differenza dipende probabilmente dal conteggio delle pause.}, con il
relatore alle prese con la voce un po' calata dal giorno prima (``la sfida
principale di oggi sarà contro la mia voce'').

Il percorso del capitolo segue quello della lezione:
\begin{enumerate}
  \item il programma della giornata (sezione~\ref{sec:z2-programma});
  \item le domande sulla prima giornata, che si trasformano in una lunga
        sessione su \textbf{colloqui e selezione}
        (sezione~\ref{sec:z2-colloqui});
  \item il \textbf{time management}: pianificare la giornata, le liste, le
        checklist, le interruzioni (sezione~\ref{sec:z2-pianificare});
  \item la \textbf{matrice di Eisenhower} per prioritizzare
        (sezione~\ref{sec:z2-eisenhower});
  \item come proteggere il proprio tempo: sprechi, distrazioni, il saper dire
        di no (anche ai meeting), il mito del multitasking e la chiusura della
        giornata (sezione~\ref{sec:z2-proteggere});
  \item le domande sul time management e una parentesi su tempo libero, ferie
        e benessere in azienda (sezioni~\ref{sec:z2-qa-tempo}
        e~\ref{sec:z2-benessere});
  \item il \textbf{brainstorming}: obiettivi, problemi tipici, un metodo
        completo (\emph{rapid ideation} e \emph{dot voting}) e la
        \textbf{matrice impatto--sforzo} per decidere
        (sezioni~\ref{sec:z2-brainstorming}, \ref{sec:z2-metodo}
        e~\ref{sec:z2-impatto});
  \item i consigli pratici, le domande sul brainstorming e il principio del
        \textbf{disagree and commit} (sezioni~\ref{sec:z2-qa-brainstorming}
        e~\ref{sec:z2-disagree});
  \item il \textbf{leading from within}: perché la leadership non è un
        titolo, le carriere divise tra manager ed esperti
        (sezione~\ref{sec:z2-leadership});
  \item il ``momento Piero Angela'' sulla biologia del gruppo e le domande
        che ne sono nate (sezioni~\ref{sec:z2-biologia}
        e~\ref{sec:z2-qa-biologia});
  \item la \textbf{piramide della fiducia}, il prendersi cura del gruppo e
        che cosa cambia quando si diventa manager
        (sezioni~\ref{sec:z2-piramide}, \ref{sec:z2-gruppo}
        e~\ref{sec:z2-manager});
  \item le domande finali e la chiusura del workshop
        (sezioni~\ref{sec:z2-qa-finale} e~\ref{sec:z2-chiusura}).
\end{enumerate}

Come per il capitolo precedente, una premessa di metodo: buona parte della
giornata è fatta di domande della classe e di racconti dell'esperienza del
relatore. Sono riportati con cura perché sono il modo in cui i principi
vengono resi concreti, e perché spesso contengono le indicazioni più
direttamente utilizzabili (soprattutto sui colloqui di lavoro).

% ---------------------------------------------------------------------
\section{Il programma della giornata}
\label{sec:z2-programma}

Come annunciato alla fine della prima giornata
(sezione~\ref{sec:z1-domani}), gli argomenti sono due:
\begin{itemize}
  \item un ``doppio topic'', \textbf{time management e brainstorming};
  \item il \textbf{leading from within}, che il relatore definisce ``la
        versione avanzata di ieri'': come guidare un team stando dentro al
        team, verso i colleghi e verso il gruppo con cui si lavora.
\end{itemize}
Poiché l'ordine è indifferente, il relatore lascia scegliere alla classe: dalle
poche preferenze espresse, la leadership va al pomeriggio, così che le domande
finali si concentrino su quella, mentre la mattina è dedicata a time
management e brainstorming. Il materiale c'è, ma il relatore avverte che la
giornata potrebbe finire prima delle quattro: non vuole trattenere nessuno se
non c'è niente di specifico da fare. Per l'eventuale tempo avanzato prima della
pausa pranzo propone un \textbf{argomento bonus}: i colloqui di lavoro,
selezione e \emph{recruitment}, visto che il relatore ha fatto e fa tuttora
selezione del personale e che per molti studenti sarà la cosa più immediata da
affrontare finito il percorso di studi (domande comuni, come rispondere,
eccetera).

Prima di cominciare, però, il relatore apre a \textbf{domande sulla giornata
precedente}: la classe ha avuto modo di ``metabolizzare'' i contenuti e
potrebbero essere emersi dubbi a posteriori. È proprio da una di queste domande
che nasce, in anticipo, la parte sui colloqui.

% ---------------------------------------------------------------------
\section{Colloqui di lavoro: capire che cosa cerca l'azienda}
\label{sec:z2-colloqui}

\subsection{L'annuncio come prima fonte di informazione}

La domanda: il giorno prima si era detto che è importante capire che cosa
l'azienda cerca davvero; ma se il primo colloquio è con le risorse umane, cioè
con qualcuno che magari ``di tecnico non sa'', come si fa a capirlo? Il
relatore risponde subito, nel dubbio che più tardi manchi il tempo.

La prima fonte è \textbf{l'annuncio}. Non viene scritto a caso: è approvato
sia dall'HR sia dall'\emph{hiring manager}\footnote{Nella trascrizione
``Airing Manager'' e ``Iron Manager''; allo stesso modo ``e-char'' e ``char''
stanno per HR.}, cioè la persona che seguirà il candidato una volta assunto,
e spesso dall'intero dipartimento HR. È calcolato e scritto in un certo modo
per far capire che cosa serve. Molti annunci moderni, oltre alla parte su
\emph{che cosa si verrà a fare}, contengono una parte su \emph{che cosa ci si
aspetta da questa persona}, cioè il comportamento e il carattere che servono
in quella posizione: ed è quella che dice molto sulla \textbf{cultura
aziendale}.

L'esempio è quello di due annunci per la stessa posizione (programmatore
senior) descritta in modo diverso:
\begin{enumerate}
  \item ``cerchiamo un programmatore esperto che possa prendersi cura dei
        nostri programmatori junior, aiutarli a crescere, sviluppare un
        ambiente di lavoro inclusivo e stimolante, in modo che i nostri
        progetti futuri siano all'avanguardia'';
  \item ``cerchiamo un programmatore esperto che sappia rispettare le
        deadline, seguire le indicazioni dei dati, sia indipendente nel suo
        lavoro e sappia relazionarsi con gli \emph{stakeholder} per capire e
        proporre nuove direzioni per il progetto''.
\end{enumerate}
Nel primo caso ci si aspetta un team di junior, probabilmente con un'età media
più bassa, con cui ci si relazionerà soprattutto; per un senior è l'occasione
di sviluppare le doti da \emph{people manager}, e al colloquio arriveranno
domande come ``hai mai lavorato con dei team? Con che età media? Ti sei mai
trovato a guidare un gruppo?''. Nel secondo caso il senior deve eseguire un
piano deciso o proporre qualcosa, e si relazionerà soprattutto \emph{verso
l'alto}; al colloquio arriverà piuttosto ``fammi questo test di
programmazione''. Sono, di fatto, due colloqui completamente diversi.

\subsection{Quanti colloqui e con chi}

Ormai è molto raro essere assunti con un solo colloquio: ne servono almeno
due, realisticamente \textbf{tre}, che è diventata la media. L'iter tipico:
\begin{enumerate}
  \item il primo con l'HR;
  \item il secondo con l'HR o con il futuro manager;
  \item il terzo sicuramente con il futuro manager, perché l'azienda vuole far
        conoscere al candidato la persona con cui lavorerà, e viceversa.
\end{enumerate}
Già le domande diverse dei vari interlocutori aiutano a capire che cosa
cercano. Ma c'è una parte fondamentale, che secondo il relatore molti
sottovalutano: il momento finale delle \textbf{``hai domande?''}.

\subsection{Le domande da fare al selezionatore}

Un principio di base: \textbf{i colloqui non sono ``devo farmi selezionare
dall'azienda'', sono ``dobbiamo selezionarci a vicenda''}. Le domande vanno
quindi preparate: ci sarà sempre qualcosa che, dalla \emph{job description} o
dal colloquio, non è chiaro al cento per cento. Nel dubbio, il relatore
suggerisce tre domande ``che funzionano sempre'':

\paragraph{``Se iniziassi domani, quale sarebbe il mio \emph{day to day}?''}
Ottiene due effetti. Il primo è pratico: descrive la giornata tipo e fa capire
se piacerebbe o no. Il secondo è più psicologico, ma il relatore sa che
funziona: costringe l'interlocutore a pensare al candidato \emph{già dentro
l'azienda}, cioè a che cosa potrebbe fare da subito e a quale impatto potrebbe
avere con le sue competenze.

\paragraph{``Se aveste una bacchetta magica, che cosa cambiereste subito
dell'azienda?''} Sposta la conversazione su di loro e va bene sia con l'HR sia
con l'hiring manager. Li obbliga a indicare una sfida interna che vorrebbero
risolvere, ed è sempre un'informazione utile: se entrambi rispondono che
cambierebbero il modo in cui si prendono le decisioni, viene da pensare che non
ci sia un percorso chiaro e trasparente da parte di chi sta in alto.

\paragraph{Lo scambio: ``vi do un motivo per non prendermi, voi datemene uno
per non venire''.} È più avanzata e il relatore non la usa sempre, ma la trova
particolarmente interessante. Permette al candidato di esprimere i propri dubbi
in modo costruttivo (``mi dite che lavorerei a stretto contatto con un team di
junior: un motivo per non prendermi è che non ho anni di esperienza nella
gestione di giovani sviluppatori, sarebbe la mia prima volta; però penso di
poterlo fare perché ho questo interesse e queste competenze''). Dall'altra
parte, il selezionatore dirà che cosa secondo lui sarebbe più difficile per il
candidato: si scopre così che cosa sta già pensando, e l'informazione torna
utile anche per i colloqui successivi.

\paragraph{E se il selezionatore si innervosisce?} ``Se un recruiter si
innervosisce per le domande del candidato, ha sbagliato lavoro.'' Al massimo il
relatore si innervosisce se un senior gli chiede che cosa faccia di preciso un
programmatore senior (``lo fai da dieci anni, lo saprai tu''); ma se un junior
chiede che cosa dovrebbe fare ogni giorno, lo si spiega volentieri. E in realtà
vale anche per i senior, perché \textbf{in ogni azienda lo stesso ruolo fa cose
diverse}: conviene sempre farsi spiegare che cosa intendono per quel titolo.
L'esempio viene dal mondo delle community: \emph{community manager} a Riot
indicava la persona a capo del team di \emph{community representative} di una
nazionalità (``community manager italiano''); a Blizzard indicava la persona
sopra l'intero dipartimento che gestiva quei team. Candidarsi per lo stesso
titolo non garantisce di aver capito per che cosa ci si sta candidando, e i
recruiter lo sanno.

Se poi un recruiter si innervosisce comunque, sta dando più risposte di quelle
che pensa: vuol dire che l'azienda non vuole domande, o non vuole interazioni
del tipo ``ma perché facciamo questa cosa?''. Che abbia formato i suoi
recruiter a innervosirsi, o che accetti che lo facciano, è informativo.

\subsection{Le domande che dicono molto sull'azienda}

Allo stesso modo, come detto il giorno prima, è informativo se un recruiter
chiede a una donna se vuole avere figli: ``ah, quindi a voi interessa sapere se
voglio avere figli: questo già mi dice che probabilmente non voglio lavorare
qui''.

Un esempio \emph{gender neutral}: la richiesta di vedere la \textbf{busta paga
precedente}. ``Quanto è la tua RAL? Mandami una busta paga.'' La risposta del
relatore è no: ``ti ho detto qual è la mia RAL; se vogliamo instaurare un
rapporto in cui non ti fidi di quello che dico, abbiamo un altro problema''. È
di nuovo \emph{default to trust} (sezione~\ref{sec:z1-default-trust}). Il
relatore è quasi sicuro che legalmente non si possa chiedere la busta paga di
un candidato\footnote{È l'opinione del relatore, espressa come tale (``sono
quasi sicuro''); non si tratta di un parere legale e non è stata verificata.}:
a quel punto diventa lui antagonista del recruiter, e visto che la
conversazione è finita in una ``zona grigia'', fa lui le domande (``come mai è
legale questa cosa che mi chiedi?''). Ha già deciso che probabilmente non
lavorerà lì, ma così magari evita che i candidati successivi debbano subire le
stesse richieste. Cita l'esempio perché sa che in Italia è abbastanza comune,
mentre all'estero non gli è mai successo, anche se avrebbe potuto ``sparare
qualsiasi cifra''.

Il motivo è che \textbf{non c'è alcun onere di dimostrare quanto si
guadagna}: chi spara troppo alto si autoesclude. Recruiter e manager non
arrivano al colloquio con carta bianca né con l'indicazione di risparmiare il
più possibile, ma con un \emph{range}: ``per questa posizione possiamo
spendere dai 30 ai 40 mila euro l'anno; paga questa persona a seconda della
sua seniority''. Se un candidato che guadagna 30 dice di prenderne 55 per
``provare a fare i soldi'' e il budget arriva al massimo a 40, il relatore non
chiederebbe nessuna prova: direbbe ``grazie, siamo in un range diverso, possiamo
salutarci''. La domanda sulla busta paga è quindi \emph{formulata male}, perché
è poco furba. Quella giusta è: \textbf{``hai già una RAL in mente per questa
posizione?''}. Si ottengono le stesse informazioni senza passare per la zona
grigia.

\subsection{La cifra minima mentale}

E se la domanda la fanno a un junior che non ha mai lavorato? Il consiglio vale
per tutti: \textbf{preparatevi sempre la vostra cifra minima mentale}, per
qualsiasi offerta di lavoro, prima ancora di fare il colloquio. Non serve
conoscere la media italiana per quella posizione o quanto paga quell'azienda:
serve sapere \emph{quanto si è disposti ad accettare senza andare in negativo},
tenendo conto dei propri costi (la benzina per andare in ufficio ogni giorno, un
eventuale trasferimento, e così via). Quella cifra va sempre avuta in testa.

\subsection{La scrematura iniziale e il portfolio}

Uno studente chiede come funzioni la risposta a una candidatura: arriva un
rifiuto diretto, oppure c'è una scrematura iniziale e poi i colloqui? E vale
per tutti i ruoli? Sì: c'è una \textbf{scrematura iniziale} fatta dall'HR o, in
alcune aziende, direttamente dal software dell'HR. Si scarta chi è chiaramente
fuori da ciò che si cerca:
\begin{itemize}
  \item si cerca un \emph{graphic designer} e arriva il CV di un programmatore
        (``faccio questi esempi perché purtroppo succede'');
  \item si cerca un \emph{animator} tecnico bravo con MotionBuilder, lo si
        scrive nell'annuncio, e il candidato risponde ``so fare tutto tranne
        MotionBuilder'': c'è un motivo tecnico per non poterlo prendere;
  \item problemi di seniority: si cerca un senior e risponde un junior con un
        anno di esperienza.
\end{itemize}
Se la figura non è già chiaramente fuori target, si fa il primo colloquio.

\paragraph{Gli anni di esperienza negli annunci junior.} Spesso gli annunci
junior chiedono come minimo tre anni di esperienza. \textbf{Se l'annuncio dice
``junior'', ignorate i requisiti di esperienza}: o sono lì per sbaglio, o sono lì
per scremare. Il secondo caso è più \emph{borderline}, ma il relatore l'ha
visto fare in alcune aziende americane: si mettono requisiti molto alti,
chiaramente irraggiungibili, per vedere chi ha la ``faccia tosta'' di mandare
comunque il CV (``non me ne frega se chiedete un \emph{bachelor's degree}, non ce
l'ho, mando lo stesso''). Così però si filtra una grossa fetta di popolazione e
si seleziona un certo tipo di carattere: si stanno cercando ``i Top Gun''. Non è
quindi necessario iniziare con uno stage: si può iniziare direttamente da
junior.

\paragraph{Chi guarda il portfolio.} Il relatore non ha mai avuto un portfolio,
ma per le figure che lo richiedono parte comunque dal CV (per escludere, ad
esempio, il designer che si candida da programmatore) e guarda il portfolio
solo a parità con altri candidati: da HR può aprirlo e dire ``bello'', ma non
può dare un giudizio tecnico. Il portfolio lo apre chi sarà l'hiring manager,
che è sempre molto più tecnico. L'ordine quindi è: l'HR guarda il CV e si salva
il portfolio; se il candidato passa la prima scrematura, CV e portfolio vanno
a un tecnico che conferma (``sì, sa fare queste cose'') o boccia (``il
portfolio indica chiaramente che non sarebbe in grado, digli pure che non
facciamo il colloquio'').

\subsection{Test pratici, da casa e con qualsiasi strumento}

\paragraph{Si può chiedere un test pratico al posto del portfolio?} Uno
studente chiede se, con un portfolio che non rispecchia le proprie capacità, si
possa chiedere direttamente un test. No: se l'azienda ha bisogno di verificare
una competenza ha già pronto un test e lo fa fare a prescindere. Il relatore,
per esempio, fa fare un test ai programmatori e lo fa correggere dal suo senior
developer; per gli artisti si chiedono animazioni di prova, chiaramente veloci.
Non chiede mai test troppo lunghi, perché non vuole far lavorare nessuno gratis:
non chiede ``fammi un pitch di progetto'', chiede uno spunto per capire come si
ragiona e come si lavora. Per il resto, \textbf{se il portfolio non è
indicativo, va aggiornato il portfolio}.

\paragraph{Test dal vivo o a casa?} Il relatore li fa fare sempre a casa. Sa
benissimo che un programmatore a casa lavorerà con due monitor, il codice da
una parte e Google o ChatGPT dall'altra. Ne ha discusso con i suoi senior
programmer: interessa che qualcuno sappia le cose a memoria, o che sappia
trovarle e fare quel lavoro a prescindere? A lui interessa che il lavoro
\textbf{venga fatto}: che il candidato abbia imparato il codice a memoria o che
in venti minuti trovi le informazioni, le metta insieme, le faccia elaborare a
ChatGPT, le corregga e consegni il codice che serve, non ha preferenze. Quando
si lavora le indicazioni saranno ``fai questa cosa il prima possibile'': se la
si sa fare anche con strumenti aggiuntivi, ben venga, ``anzi, spiegalo anche
agli altri''. Farla a memoria va benissimo, ``ma perché mettere l'\emph{hard
mode}?''. E comunque ci sono cose che non si possono delegare a uno strumento:
se arriva un codice che non funziona e va fatto il \emph{debug}, bisogna
saperlo fare e saper progettare i test. La figura più completa, per il
relatore, è anche quella che sa trovarsi le informazioni e gli aiuti.

Le prove dal vivo sono soprattutto domande di \emph{problem solving}: il
relatore vuole sentire come si ragiona e chiede espressamente di \textbf{pensare
ad alta voce}, per seguire il percorso mentale che porta alla conclusione. Sono
più domande impegnative che veri test. I test veri e propri si fanno con i
propri tempi e i propri strumenti: ``scegli il momento, ti mandiamo il test, hai
due ore per rimandarlo indietro''. Questo è l'accordo; in quelle due ore si può
fare quello che si vuole.

\subsection{Junior, mid e senior}

\paragraph{Perché ci sono così pochi annunci per figure mid?} Uno studente ha
notato su LinkedIn e sulle piattaforme di recruiting un divario notevole tra le
richieste di junior e quelle di mid. Il motivo è che per molte aziende,
Stormind inclusa, \textbf{la differenza tra junior e mid è sparita}: o si è
senior o non lo si è. Ci si è resi conto che la differenza tra un junior e un
mid è che il mid si è abituato a lavorare in quell'azienda e il junior non
ancora; ma un mid che arriva con tre anni di esperienza in un altro studio (ad
esempio Milestone\footnote{Nella trascrizione ``se ha fatto tre anni in
milestone'': si tratta verosimilmente di Milestone, studio di sviluppo
milanese.}) va comunque ``ripreso come junior'', perché si lavora in modo
diverso e deve reimparare. Un senior, invece, ha molte più competenze e le sa
applicare in un certo modo, e quella distinzione resta netta.

\paragraph{Ci si può candidare anche se il livello non coincide?} Un'altra
studentessa dava per scontato che, se un annuncio non specifica ``junior'', sia
per mid. Se non è specificato ``senior'', secondo il relatore si può mandare il
CV chiedendo di valutare. Certo, se è specificato junior e si è junior le
possibilità di essere contattati sono maggiori. Ma vale lo stesso discorso del
livello di inglese fatto il giorno prima: si cerca qualcuno che \emph{sappia
fare quelle cose}, non una certificazione. Se serve un programmatore che sappia
usare Unreal Engine~5 e sistemare la cosa X o Y, che sia mid da due anni o
junior da uno non è un problema.

\subsection{Gli strumenti di intelligenza artificiale}

Una domanda in chat riguarda l'uso quotidiano dell'intelligenza artificiale in
azienda. Il relatore sta facendo fare corsi di IA al suo team ``perché voglio che
la gente lavori meno'': vuole che si sfrutti qualsiasi strumento disponibile,
perché gli altri lo stanno facendo e non vuole restare indietro. Se un lavoro
si può fare in tre ore di meno, non c'è alcun vantaggio a farlo in modo lungo e
tortuoso. Il rischio esiste: se si fa fare tutto a ChatGPT e non ci si tiene al
passo con il proprio \emph{know-how}, quando ci sarà da risolvere qualcosa ci si
potrebbe ``impallare''. Ma sta al singolo restare aggiornato e saper fare il
proprio lavoro: \textbf{gli strumenti devono essere strumenti, non un sostituto
delle competenze}. ``Non è una partita di calcio, con regole su come si devono
fare le cose o un \emph{dress code}'': il lavoro va fatto nel modo più efficiente
possibile.

A uno studente a cui era saltata la connessione proprio in questo momento il
relatore risponde che la lezione è registrata e si può recuperare dal video.
Poi chiude la parentesi: ``la parte difficile era finire ieri in tempo, ce
l'abbiamo fatta''.

% ---------------------------------------------------------------------
\section{Time management: pianificare la giornata}
\label{sec:z2-pianificare}

\subsection{Time management o prioritizzazione?}

La domanda di partenza alla classe è: \emph{chi di voi, in questo momento, fa
time management su se stesso}, cioè si organizza la giornata all'inizio e
organizza il proprio tempo? Qualcuno risponde di sì, qualcuno ha uno
\emph{schedule} settimanale con orari da rispettare, qualcuno ``ci prova''
(``hai fatto bene a specificare \emph{ci provo}: magari non la rispetto''). Il
relatore lo chiede perché è uno di quegli argomenti in cui molti, dopo aver
visto la spiegazione, si accorgono di non fare davvero time management: fanno
piuttosto \textbf{prioritizzazione}, ad esempio una lista di cose da fare, che
è un po' diverso.

Nessuna delle tecniche è complicatissima, ma è utile sapersi muovere. Il
relatore colloca volutamente questo modulo \emph{dopo} quello sul feedback: saper
gestire il proprio tempo permette anche di dare feedback agli altri sulla loro
gestione del tempo.

\subsection{La \emph{must do list}}

Il punto di partenza è il \textbf{planning iniziale}: non improvvisare la
giornata. Improvvisando si è in balia di una grande limitazione umana: ci si
dimentica le cose, oppure ci si lascia trascinare dal task più urgente o
dall'ultimo arrivato. Esistono modi di dire come ``se arriva una mail o rispondi
subito o diventa un task per domani'', perché il ``rispondo dopo'' non arriva
mai; più avanti si vedranno modi più raffinati di gestire le zone grigie, ma la
pianificazione resta fondamentale.

La prima cosa da fare \emph{non} è la lista delle cose da fare oggi, ma la lista
delle cose che \textbf{assolutamente} vanno fatte oggi. Le due liste non devono
coincidere: la seconda richiede un livello di ragionamento in più. Sono le cose
per cui, se alle sei si stacca e non sono fatte, si resta finché non sono fatte,
perché sono imprescindibili.

\paragraph{Il primo momento per dire ``potrei non farcela''.} La \emph{must do
list} è anche il primo indicatore della giornata. Se è troppo corposa (dipende
naturalmente dalla complessità dei task), quello è il momento, ``alle 9:01 se
iniziate alle 9'', per dire al proprio manager \emph{oggi sarò in difficoltà}.
Alle 9:01 il manager ha l'intera giornata per aiutare; se invece si pecca di
ottimismo o di presunzione (``ce la posso fare'') e glielo si dice alle 17:59,
non può fare nulla. È lo stesso discorso del giorno prima
(sezione~\ref{sec:z1-triangolo}): anche facendo la to-do list si applicano i
concetti della cultura del feedback, ed è il motivo per cui questo modulo viene
dopo. Ci saranno vari momenti di questo tipo nella giornata, ma questo è il
primo.

\paragraph{E se ci si accorge a metà giornata di non farcela?} Uno studente
chiede che cosa fare se, convinto di farcela, si accorge a metà giornata di non
poter completare bene la \emph{must do list}. Il relatore lo concederebbe
\emph{una volta}: se si vogliono testare i propri limiti, va bene, ma
comunque si dice al capo ``oggi potrei essere tirato, ci provo'', e sarà lui a
valutare se il rischio è accettabile. Dipende dal rischio: se il \emph{must do}
è consegnare una milestone a un cliente, ``nel dubbio facciamolo fare a
qualcun altro, non possiamo correre quel rischio''; se il rischio è che domani
un animator riceva il codice con tre ore di ritardo, lo si può correre. ``È un
po' questo il lavoro.''

\paragraph{Cinque minuti e una stima.} Tutto questo va fatto in circa cinque
minuti. Se non è chiaro che cosa si deve fare in una giornata, \textbf{si
chiede}: ogni giorno dovrebbe essere chiaro. Poi si cerca di \textbf{stimare il
tempo} di ogni cosa. Si diventa più bravi col tempo: alla quinta volta che si
fa il report mensile si sa che ci vuole un'ora e ventidue minuti; la prima volta
si può chiedere a un collega quanto ci mette lui, sapendo che per noi potrebbe
essere diverso.

\subsection{To-do list e checklist}

Esistono molti strumenti per tenere to-do list e checklist; il relatore non ne
consiglia uno in particolare e cita come esempi Any.do, i promemoria di Windows
o Samsung Notes\footnote{Nella trascrizione ``AnyDo, Windows Memo''. Any.do è
un'app di liste e promemoria; ``Windows Memo'' indica probabilmente le note o i
promemoria integrati in Windows (Sticky Notes, Microsoft To Do).}: ``potete usare
davvero la qualunque''. L'importante è tenere la lista \textbf{visibile} da
qualche parte sul computer. Alcuni accorgimenti:
\begin{itemize}
  \item i \textbf{must do} vanno sempre \textbf{in cima} alla lista;
  \item quando durante la giornata arriva un task nuovo, la cosa da fare subito
        non è \emph{eseguirlo}, ma \textbf{aggiungerlo alla lista}: ``se arriva
        una mail, aggiungete subito alla lista \emph{rispondi alla mail}'',
        decidendo se è diventato un \emph{must do} o se è comunque un task per
        oggi;
  \item non limitarsi alle cose operative: il relatore ha iniziato ad
        aggiungere anche task come ``parla con Tizio'' o ``parla con Caio'',
        perché è da un po' che non chiede loro un aggiornamento; se non li
        mette in lista si accorge di non farli;
  \item spuntare le voci completate: per quanto minimo, rilascia
        \textbf{dopamina} (se ne parlerà nel pomeriggio,
        sezione~\ref{sec:z2-biologia}), aiuta a motivarsi e ad andare avanti.
        Sono piccole cose, ma si sommano.
\end{itemize}

\paragraph{Priorità contro recenza.} Che cosa succede se si esegue sempre
l'ultimo task arrivato? Lo chiede il relatore alla classe. Si interrompe il
flusso di lavoro e si cambia continuamente attività, certo; ma soprattutto si
stanno \textbf{mescolando priorità e recenza}: si sta dicendo ``non importa come
ho prioritizzato i task, la mia priorità è sempre l'ultimo arrivato'', che non
vuol dire fare le cose più importanti. È il motivo per cui tanti dipartimenti IT,
quando si chiede loro qualcosa, rispondono ``mi apri un ticket?''. È diventata
quasi una battuta, ma il concetto è chiarissimo: devono poter prioritizzare
quella richiesta rispetto alle altre 112 arrivate in giornata, e senza ticket
non possono confrontarla con le altre. Il relatore è totalmente a favore: ``certo
che ti apro un ticket, altrimenti tu come fai a tracciare il tuo lavoro?''.

\subsection{La checklist: l'esempio dell'aviazione}

La \textbf{checklist} non è la task list. La task list contiene le cose da fare;
la checklist contiene \textbf{i passi da seguire quando si esegue un task
specifico}. L'esempio è quello delle \textbf{checklist pre-volo}, introdotte
nell'aviazione (secondo il relatore tra gli anni Sessanta e Settanta): chi vola
spesso nota che i piloti seguono sempre un ordine simile nei controlli, tanto
che il relatore capisce quanto manca al decollo da quando controllano i flap, i
motori, quando accendono e spengono le luci. L'introduzione delle checklist
avrebbe ridotto di circa il 75\% gli incidenti in fase di decollo\footnote{Date
e percentuale sono quelle riferite in lezione e non sono state verificate.
L'origine della checklist pre-volo viene di solito fatta risalire al 1935, dopo
l'incidente del prototipo del Boeing B-17, causato da un controllo dimenticato da
un equipaggio molto esperto; la diffusione sistematica delle checklist e dei
protocolli di cabina è invece successiva.}, e soprattutto per i \textbf{piloti
più senior}. Perché? Perché chi ha fatto un decollo centomila volte non va a
controllare l'alettone destro: ha funzionato centocinquantamila volte, funzionerà
anche stavolta. Quando ci si abitua a fare qualcosa, si sbaglia proprio quando
qualcosa non va. La checklist ha obbligato a fare quei controlli ogni volta e ha
ridotto in modo drastico gli incidenti.

Il relatore usa checklist anche per cose banali, ad esempio le \textbf{email}:
essendo una persona abbastanza ansiosa su queste cose, se prepara una mail alla
leggera finisce per passare troppo tempo a ricontrollare l'oggetto, il
destinatario, l'allegato, e comunque si dimentica qualcosa. Seguire una checklist
sembra più complesso, ma \textbf{fa risparmiare tempo}: se la checklist è
completa, si sa di aver fatto tutto e si può andare avanti senza continuare a
pensare ``ho chiuso il gas? ho spento la luce?''. È come la lista della spesa: se
tutto è segnato e preso, punto. È un investimento: trenta secondi in più per
fare la cosa, invece di cinque minuti a sperare di averla fatta giusta.

\paragraph{E se ci si dimentica qualcosa nella checklist?} Succede. A quel
punto \emph{preparare la checklist} è un task in sé: un giorno in cui si ha
tempo si fa quel task pensando a quali passi si seguono e a come assicurarsi che
funzioni. Nessuna checklist è perfetta: ci saranno sempre compromessi tra ``devo
davvero controllarlo ogni volta?'' e ``posso farne a meno?''. Però, in generale,
\textbf{se qualcosa anche solo una volta ha dato problemi, va in checklist},
perché vuol dire che potrebbe darne di nuovo. Nessun manager rimprovererà
``come hai fatto a non prevedere che questa cosa poteva andare male?'': si
diventa esperti anche nelle checklist di un task facendo quel task, ed è un
processo iterativo. Se qualcosa è stato dimenticato e ha creato problemi, lo si
aggiunge per la volta dopo. Sbagliare una o due volte non è un problema;
l'importante è che la situazione migliori la volta successiva.

\paragraph{L'inizio di un junior e l'\emph{onboarding}.} A uno studente che
osserva che l'inizio di un junior è quindi inevitabilmente un po' travagliato,
il relatore risponde che è giusto così: ``se volessi il dipendente perfetto da
subito, cercherei un senior, non un junior''. Allo stesso modo non capisce le
aziende che prendono un junior e non lo fanno seguire da nessuno: che cosa ci si
aspetta, che assorba passivamente tutto quello che ha intorno e impari ``per
osmosi''? Può succedere, ma imparerà a fare le cose come vuole lui, perché
nessuno gli sta dicendo ``meglio farle così''. Moltissime aziende, oltre al
manager (che magari ha quindici \emph{report} e non ha il tempo materiale di
seguire un junior come andrebbe seguito), assegnano un \textbf{referente}, di
solito un senior, con il compito di prendersi cura della persona nuova e
insegnarle a fare le cose. Altrimenti l'\textbf{onboarding}, cioè il tempo
necessario per arrivare al cento per cento nel proprio lavoro, dura un anno.
Più risorse l'azienda dà all'inizio, prima la persona diventa operativa: è
nell'interesse dell'azienda dare il massimo supporto a chi arriva.

Il relatore diffida delle aziende che ``vogliono vedere se uno se la cava''. È
il discorso dei cento metri e della corsa a ostacoli\footnote{Il relatore lo
introduce come un discorso già fatto (``si torna al discorso: vuoi fare i cento
metri o la corsa a ostacoli?''), ma il passaggio non compare nella trascrizione
della prima giornata riportata nel capitolo~\ref{cap:z1-feedback}.}: non si può
far correre uno sui cento metri e l'altro con gli ostacoli e aspettarsi tempi
simili, perché chi ha gli ostacoli ci metterà per forza di più. Il supporto ai junior è un
investimento sensato.

\subsection{\emph{First contact resolution}}

Un concetto che viene dai call center: \textbf{quante richieste vengono risolte
al primo contatto}. È importante quando un task richiede iterazioni o
comunicazione con altri: bisogna organizzarsi come se si dovesse risolvere la
cosa \emph{in una sola passata}. Se viene chiesto un modello 3D, lo si fa giusto
la prima volta, al meglio delle proprie possibilità. Se si deve parlare con
qualcuno, si raccolgono tutte le richieste in una.

L'esempio è quello dei ticket che il relatore mandava ai colleghi del supporto di
Blizzard: ``ciao, ho un problema, non riesco a loggarmi'' --- ``mi dici lo
username?'' --- ``è questo'' --- ``mi dici l'email?'' --- ``è questa'' --- e così
via, con un paio d'ore tra un messaggio e l'altro perché il ticket deve fare il
suo giro. Tutto poteva essere riassunto in ``ciao, mi servono username, email e
questo e quest'altro''. Unificare le richieste su una singola cosa aiuta a
risolverla il prima possibile; bisogna avere chiaro \emph{che cosa serve prima
di fare una richiesta}, in modo che ce ne sia una sola. In molti sistemi di
supporto, se si apre un secondo ticket per aggiungere informazioni al primo, il
primo, anche se era quasi in lavorazione, viene rimandato in fondo alla coda:
serve a disincentivare chi manda continuamente ticket del tipo ``mi sono
ricordato questa cosa''. O si manda tutto subito, o ci vorrà tempo.

\subsection{Non iniziare senza le informazioni}

Corollario: \textbf{se non si ha tutto quello che serve per completare un task,
non lo si inizia}. Prima ci si procura quello che manca. Al relatore capita con
la pianificazione dei due anni di lavoro successivi: se non ha tutte le
informazioni, prima lavora per averle (``mi dai questa informazione, mi dai
quell'altra''), poi pianifica; altrimenti pianificherebbe qualcosa di cui non ha
completa visibilità.

\subsection{Il costo delle interruzioni}

Alternare continuamente l'attenzione tra i task e, soprattutto, reagire di
continuo agli input esterni (la mail, il messaggio, il collega che parla) fa
\textbf{sparire il tempo dalla giornata}. Se mentre si prepara un report in
Excel si viene interrotti (dal gatto, da una mail, da un collega, da una
telefonata) e poi si torna al foglio, servono in media \textbf{un quarto d'ora}
per tornare mentalmente al cento per cento; chi è particolarmente bravo ci mette
dieci minuti, ma si parla comunque di dieci-quindici minuti\footnote{Ordini di
grandezza come riportati in lezione. Le ricerche più citate sul tema (ad esempio
quelle di Gloria Mark sulle interruzioni nel lavoro d'ufficio) indicano tempi
medi di ripresa anche superiori, intorno ai venti minuti.}. Se si viene
interrotti di continuo, si passa più tempo a cercare di tornare al cento per
cento che a lavorare. È una partita persa in partenza.

Per questo il giorno prima si era parlato dei \textbf{semafori} sulla scrivania
(sezione~\ref{sec:z1-dare}): il semaforo può essere anche semplicemente il
proprio stato su Slack, su Discord o su qualunque strumento si usi con i
colleghi. Se si è occupati, si è occupati: è una propria valutazione che,
accettando di essere disturbati di continuo, quella cosa non si riuscirà a
farla, o almeno non in tempo.

% ---------------------------------------------------------------------
\section{Prioritizzare: la matrice di Eisenhower}
\label{sec:z2-eisenhower}

Ora che si hanno task list, checklist, \emph{must do} e \emph{nice to do}, resta
la domanda di prima: se arriva un task nuovo e si decide di non fare per prime le
cose nuove, \emph{come si capisce se è davvero urgente}? Serve un metodo, ed è la
\textbf{matrice di Eisenhower}\footnote{Nella trascrizione anche ``Jason
Howard''. Dal nome del presidente statunitense Dwight D.~Eisenhower, a cui si
attribuisce la distinzione tra urgente e importante.} (la classe prova a
indovinarne il nome proponendo altre matrici, come una matrice dei rischi, ma è
un'altra cosa). Ogni task che arriva va posizionato su questa tabella:

\begin{center}
  \begin{tabular}{@{}p{0.2\textwidth}p{0.36\textwidth}p{0.36\textwidth}@{}}
    \toprule
    & \textbf{Urgente} & \textbf{Non urgente} \\
    \midrule
    \textbf{Importante} &
    \textbf{Fallo} (\emph{do}): diventa un task di oggi. &
    \textbf{Pianificalo} (\emph{schedule}): va fatto, ma non subito. \\
    \addlinespace
    \textbf{Non importante} &
    \textbf{Delegalo} (\emph{delegate}): va fatto, ma non serve che lo faccia
    tu. &
    \textbf{Eliminalo} (\emph{delete}): non lo fai. \\
    \bottomrule
  \end{tabular}
\end{center}

\paragraph{Urgente e importante: lo fai.} Sono di solito task con una
\textbf{deadline molto chiara} (entro oggi, entro un'ora, entro X) e con
\textbf{conseguenze} se non vengono fatti: da qui l'``importante''.

\paragraph{Importante ma non urgente: lo pianifichi.} Va fatto, ma non subito:
o lo si mette tra le cose di oggi ``ma dopo'', o anche alla settimana successiva.
L'esempio: un collaboratore scrive ``il mio rinnovo è tra tre mesi, ne possiamo
parlare?''. È importante, perché il rinnovo di una persona è importante; non è
urgente, se ne può parlare anche domani, così da potersi preparare e fare altre
cose nel frattempo.

\paragraph{Urgente ma non importante: lo deleghi.} Qui ``non importante''
\textbf{non} vuol dire che il task non sia importante, ma che \textbf{non è
importante che lo faccia tu}, che è molto diverso. Se arriva ``mi serve questo
cambio a questa animazione 3D'' e ciò che conta è il cambio, non chi lo fa, si può
dire ``sono d'accordo che è urgente, ma non serve che lo faccia io'' e chiedere a
qualcun altro. Sono task che vanno fatti ma non richiedono la propria particolare
\emph{expertise}.

Il relatore insiste molto su questo punto con manager e senior: quando arriva un
task urgente, i senior tendono a dire ``lo faccio io, perché è urgente''. Ma se un
junior può farlo, lo si fa fare al junior e si investono le competenze più alte
del senior in qualcos'altro: il senior può fare il lavoro del junior, ma il
junior non può fare il lavoro del senior mentre il senior fa quello del junior.
Le persone non sono intercambiabili. Si usa il \textbf{minimo sforzo
necessario}: se una cosa la può fare un mid, la fa il mid, non il senior.

\paragraph{Né urgente né importante: lo elimini.} È il quadrante che quasi
nessuno usa: tutti fanno, delegano o pianificano, ma non pensano mai a
\emph{delete}. Invece è fondamentale. A volte è corretto dire: ``questa cosa non
dovrei farla io, non è urgente, non la faccio; sappi che non la faccio, e se c'è
un problema ne parli con il manager''. Non si sta delegando a qualcun altro: si
sta dicendo che quella cosa così com'è \textbf{non andrebbe fatta}, perché non
ne sono chiari l'importanza e il perché. Nella propria lista di oggi non ha
spazio, e non è abbastanza importante per essere fatta da sé nemmeno nei giorni
successivi.

Un esempio di \emph{delete} che non è un vero \emph{delete}: arriva una richiesta
di meeting di tre ore sul futuro del progetto, con invitati tutti gli art
director, il game director e così via. Il relatore spesso risponde ``io non
vengo'': per lui è \emph{delete}, ma il meeting si fa lo stesso. Uno studente
osserva che a volte si tratta di una richiesta mal formulata o poco chiara: è
esattamente così. ``Nella sua forma per me è no; se poi mi spieghi meglio perché
dovrei farlo io o perché è urgente, allora cambia.''

% ---------------------------------------------------------------------
\section{Proteggere il proprio tempo}
\label{sec:z2-proteggere}

\subsection{Evitare gli sprechi: un limite che nessuno può insegnare}

Evitare gli sprechi di tempo è difficile, soprattutto da remoto: è una sfida
costante con se stessi. Ha molto a che fare con il carattere, il
\emph{self-control} e (``passatemi anche questa'') la professionalità. Il
relatore può dare tutti i consigli del mondo, ma se dentro di sé non si percepisce
l'urgenza di fare il proprio lavoro quando lo si deve fare, non lo si può
insegnare. Se si è a casa e invece di lavorare si gioca a qualcosa su Steam,
nessuno deve spiegare che non si dovrebbe. Ci sono limiti oltre i quali non si
può andare.

Quello che si può fare è sapere che \textbf{ogni interruzione ha un costo
sommerso}: ogni messaggio, domanda, telefonata, ogni cambio di focus. Sono i
quindici minuti di prima, che possono arrivare a trenta a seconda del tipo di
task.

\paragraph{È più difficile tornare su un task logico o creativo?} Il relatore lo
chiede alla classe: fa più fatica un artista a tornare sul suo modello 3D o un
\emph{data analyst} a tornare sul suo foglio Excel? Le risposte si dividono
(``creativo'', ``entrambi''), ed è proprio il punto: dipende a chi lo si chiede. I
creativi rispondono ``creativo'', gli analisti ``analitico'', perché ognuno ha una
chiara comprensione di che cosa serve per essere davvero concentrati su ciò che
fa. Il creativo sa che per essere pienamente creativo gli servono l'ambiente, la
concentrazione, gli spunti, gli esempi da tenere chiari in testa; l'analista deve
avere in testa i conti, i fogli di dati e così via. Sono \textbf{ugualmente
difficili}, ma solo chi fa quella cosa ne apprezza fino in fondo la difficoltà. Si
può avere un assaggio dell'altra parte, ma la piena comprensione arriva solo
facendo: è come guardare uno sport in tv e pensare ``potrei farlo anch'io'',
finché non ci si trova lì.

\subsection{Distrazioni volontarie e involontarie}

Uno studente, che sta leggendo manuali di \emph{design thinking}, chiede come
conciliare tutto questo con l'idea che distrarsi, cioè non restare focalizzati,
aiuti a trovare un nuovo punto di vista per risolvere un problema. Qual è il
giusto equilibrio? Il relatore è \textbf{totalmente d'accordo} con quel concetto:
sta parlando di \textbf{interruzioni non volontarie}. Sa benissimo che dopo cinque
ore su un foglio Excel ``esplode la testa''; ma se si pianifica un quarto d'ora di
pausa e lo si rispetta, si sta investendo sulla propria creatività ed
efficienza. Non è un investimento dover rispondere a una chiamata ogni tre
minuti.

Anzi, il consiglio è \textbf{crearsi interruzioni pianificate}, sotto il proprio
controllo: fare pause, fare altro, soprattutto nel \emph{problem solving}. Il
relatore stesso, se non trova una soluzione, si ferma, fa un altro task e poi ci
torna; ma è una sua scelta. La discriminante è che la pausa sia \emph{scelta} e
non imposta: anche se si fanno le pause giuste, se si viene interrotti proprio
nel momento di massima concentrazione (raggiunta magari grazie alla pausa) si
butta via tutto. Uno studente cita la \textbf{tecnica del pomodoro}: è
esattamente quello, piccole pause pianificate che aiutano a essere creativi,
purché restino pianificate. Se i quindici minuti diventano ``vabbè, facciamo
mezz'ora'' o ``lo faccio nel pomeriggio'', non è più la tecnica del pomodoro:
diventano interruzioni, o un pretesto per non lavorare.

\subsection{Il peggior nemico siamo noi}

E chi sono i peggiori nemici in questo? ``Noi stessi'': siamo noi ad
autoconvincerci che venti minuti vanno bene quanto quindici. Le
\textbf{distrazioni autoinflitte} sono spesso le peggiori. Il relatore porta
l'esempio di una sua giornata, che ha annotato perché lo aveva fatto ridere
quanto tempo era riuscito a buttare via senza che nessuno glielo imponesse:
\begin{itemize}
  \item gli piace avere qualcosa in sottofondo mentre lavora (``ha a che fare con
        l'ADHD sicuramente''\footnote{Nella trascrizione ``IDACD''. Più avanti il
        relatore dice di avere l'ADHD ``come dicevo ieri''; il capitolo 0.1 non
        lo riporta, ma riferisce il suo racconto sull'``overthinking'' anche
        nell'ordinare una pizza (sezione~\ref{sec:z1-regalo}).}),
        e in genere lo aiuta. Il problema è quando il sottofondo è troppo
        interessante: la musica va bene, ma se nella playlist di Spotify parte un
        podcast di Alessandro Barbero sulla guerra civile americana, finisce per
        ascoltare quello più del lavoro;
  \item un'abitudine che ha poi perso: se faceva straordinari alle due di notte,
        da solo, per qualcosa di urgentissimo, apriva comunque Slack a
        controllare i messaggi, e la mail. ``Come se alle due di notte qualcuno
        mi scrivesse qualcosa di urgente'';
  \item la pausa con un po' di TikTok, dopo la quale si accorgeva di aver
        dimenticato su quale slide della presentazione stava lavorando.
\end{itemize}
La regola: \textbf{inserire i break è meglio che subirli passivamente}. Un conto
è dire ``faccio una pausa intenzionalmente'', un conto è ``mi sono distratto e ho
perso due ore di lavoro''. Chi ne è particolarmente soggetto non deve
demoralizzarsi: è anche un allenamento.

\subsection{Spezzettare i task grandi}

Un altro concetto utile: se si fa molta fatica a iniziare i task grossi, perché
sembrano infiniti (``più tardi inizio, più tardi inizia questa sofferenza''),
\textbf{vanno spezzettati}. Un task grosso non è quasi mai un task grosso: è
spesso un insieme di task piccoli. Si comincia a farne dei pezzi e si va avanti
piano piano, spuntandoli. Aiuta tantissimo anche nel \emph{problem solving}: un
problema gigante va scomposto in problemi più piccoli, da risolvere per primi, per
poi ricomporre il puzzle.

\subsection{Accettare di dire di no}

Tornando alla matrice di Eisenhower: bisogna \textbf{accettare di dover dire di
no}, perché dire sempre sì porta a problemi ben peggiori di quanto si pensi.

\paragraph{Il collega che diceva sempre sì.} A Riot il relatore aveva un collega
estremamente disponibile, amichevole e collaborativo: qualsiasi cosa gli si
chiedesse, rispondeva ``certo, la faccio io'' (``mi porti fuori il cane?'' ---
``certo''). Il problema è che era umano: a un certo punto si supera la quantità di
cose umanamente possibili in una giornata. In più tendeva a offrirsi volontario
per far vedere che era intraprendente e aveva iniziativa, e nei progetti che
coordinava il relatore doveva spesso rifiutare la sua candidatura perché era
sicuro che non ce l'avrebbe fatta. Il suo dire sempre sì, a un certo punto, lo ha
danneggiato: era quello che dice sempre che lo fa lui, ma non è detto che lo
faccia. ``Ce la fai?'' --- ``Sì, sì, assolutamente'' --- e poi ``scusa, non ce l'ho
fatta''. ``Ok, però ora hai messo nei guai me e il resto del team.'' Quella
disponibilità era più dannosa di un ``no, non posso''. È lo stesso discorso del
giorno prima: dire all'ultimo momento che non ce la si fa è peggio che dirlo a
inizio settimana.

Uno studente lo paragona alla \emph{sindrome dello yes man}. Non proprio: lo
\emph{yes man} dà sempre ragione all'altro perché è avverso al conflitto e non vuole
la discussione del ``spiegami perché sei in disaccordo''. Il collega, invece,
pensava davvero di poter fare molto più di quanto potesse, e rifiutava l'idea di
non riuscirci.

\paragraph{``Ci hai pensato bene?''} Vale per i progetti e per le deadline. Con i
junior soprattutto, quando a ``vuoi fare questa cosa?'' arriva un sì troppo
convinto, il relatore chiede sempre: ``ci hai pensato bene? Me lo stai dicendo
per farmi contento? Sappi che se mi dici di no non penso peggio di te; penso
peggio di te se mi dici di sì e poi non ce la fai. Vuoi rivalutare la risposta, e
magari dirmi \emph{non lo so} invece di un sì categorico?''.

\paragraph{Che cos'è una sconfitta.} Uno studente chiede se davanti alle
``sconfitte'' non si dovrebbe rimediare. Il relatore chiarisce che per lui la
sconfitta non è il task non fatto: quella è una sconfitta \emph{a breve termine},
si troverà un modo di farlo, è un danno ``ma chi se ne frega''. La vera sconfitta,
un danno \emph{a lungo termine}, è aver pensato di riuscirci e non esserci
riusciti: se non si risolve il motivo per cui si è detto sì quando era no, lo si
ripeterà altre cinque, dieci, quindici volte. Nella testa del collega di Riot c'era
sempre una scusante (``c'erano questi imprevisti che mi hanno bloccato''), come
quella con cui ci si giustifica le interruzioni. Il relatore gli rispondeva: ``se
ogni volta c'è un imprevisto, forse c'è un \emph{pattern}: il problema non è che
non riesci a fare le cose, è che non prevedi mai gli imprevisti''. Dopo sette
ritardi su dieci dovrebbe scattare l'ingranaggio mentale (``forse sbaglio io a
dire sempre di sì''), ma non è sempre facile farlo scattare, e a volte serve una
spinta. È una delle situazioni più frequenti. E la cosa più difficile è che,
quando glielo si fa notare, queste persone lo prendono come un ``non sei capace''.
Il messaggio è invece: \textbf{``sei bravo, ma sei umano anche tu''}. Non si può
fare un numero infinito di cose in un'ora: bisogna accettare i propri limiti ``dati
dalla carne'' e non pensare di essere un cyborg.

\paragraph{Il cliente non ha sempre ragione.} Uno studente osserva che vale anche
per i project manager e per il principio ``il cliente ha sempre ragione''. Anzi:
il relatore ha lavorato con molti clienti e a volte \textbf{il cliente non sa che
cosa vuole}, che è molto diverso dall'avere sempre ragione. Soprattutto con i
clienti che chiedevano pitch creativi, la sua squadra presentava sì un'opzione
con la loro richiesta, ma anche le proprie (``se vuoi raggiungere questo, sarebbe
meglio fare così''), e molto spesso il cliente sceglieva quella proposta dal team.
Spesso il cliente ha in mente \emph{una soluzione}, ma bisogna capire \emph{quale
obiettivo} ha in mente e poi dargli una soluzione: non vuole la soluzione, vuole il
\emph{goal}, e se c'è un modo migliore per arrivarci, forse è meglio un'altra
soluzione.

\subsection{Dire di no ai meeting}

La prima cosa a cui il relatore ha imparato a dire di no sono i \textbf{meeting}.
Ha passato anni in ambienti in cui il suo calendario era letteralmente solo
meeting dalla mattina alla sera: estenuante per lui e inutile per il team. Tra i
fattori c'è il \textbf{FOMO}: tanti manager vogliono essere in tutti i meeting per
non perdersi niente ed essere sempre informati su quello che fa il team. In realtà
vanno in \emph{burnout}, passano i pomeriggi nei meeting e finiscono per essere i
più distratti, quelli che non ascoltano e che fanno domande inutili solo per far
sentire che fanno domande, senza aiutare davvero.

Prima di accettare un invito il relatore si fa una serie di domande, \textbf{come
una checklist} (e le fa anche agli altri):
\begin{enumerate}
  \item \textbf{Sono davvero necessario?} Servo davvero io, o se mancassi non
        succederebbe niente? Ai meeting ``di cortesia'' non va: se pensa di
        essere necessario lo dice lui, e il fatto che si faccia un meeting
        importante e urgente senza di lui non danneggia il suo ego. ``Non
        aggiungermi solo perché pensi che serva che io venga aggiunto.''
  \item \textbf{Uno di noi potrebbe partecipare e informare il resto del team?}
        C'è già qualcuno del suo team? Se sì, non può andare solo lui? Per un
        meeting informativo (una nuova policy, una nuova procedura da spiegare a
        tutti i dipartimenti), andarci tutti significa investire tre ore a
        testa; se ci va uno solo e poi spiega agli altri, si risparmia tempo.
  \item \textbf{Il riassunto sarà reso disponibile?} Per il relatore è
        obbligatorio. Per lo stesso motivo per cui la lezione viene
        registrata: se ci si perde un pezzo lo si può rivedere, se qualcosa non è
        chiaro lo si può risentire. Ma c'è anche un motivo di \emph{office
        politics}: se in un meeting si chiede ``mi fai questa cosa per venerdì?''
        e la risposta è sì, non si vuole sentire venerdì ``intendevo venerdì
        prossimo'' o ``non ero sicuro''. Nelle note del meeting c'è scritto che hai
        detto sì. ``Lo so che è un po' paraculo, ma in realtà è così che si
        lavora'': bisogna dare più peso alla parola data, e quindi tenerne
        traccia. Finché tutto resta orale, niente rimane scritto e tutto è
        interpretabile. Oggi è anche facile: Zoom, con l'assistente IA, produce
        in automatico trascrizione e sottotitoli, senza bisogno di qualcuno
        addetto alle \emph{meeting minutes}. In ogni caso, se in un meeting viene
        presa una decisione, il relatore vuole che sia \textbf{ripetuta per
        iscritto}: non deve ipotizzare, deve poter andare a rileggere.
  \item \textbf{Il meeting aggiunge valore rispetto al documento?} Se arriva un
        invito con allegato un documento di venti pagine, il relatore chiede: ``il
        meeting consiste nel leggermi il documento? Perché altrimenti non
        vengo''. Se non c'è un valore aggiunto, se lo legge da solo e manda
        conferma di averlo letto. È lo stesso motivo per cui si sconsigliano le
        presentazioni in cui si leggono le slide; e infatti, osserva, in quel
        momento lui stesso stava leggendo una slide, ma aggiungendo informazioni.
        ``Se tutta la lezione fosse io che vi leggo le slide, non servirei io: vi
        manderei le slide.''
  \item \textbf{Per che cos'altro potrei usare questo tempo?} È forse la più
        importante. Anche ammesso che serva, che serva proprio lui e che il
        meeting sia importante, in quel momento dovrebbe fare qualcosa di più
        importante? Se sì, fa quella. Non bisogna avere paura di dire ``no, non
        vengo'' anche se si serve, oppure ``ho un'altra cosa importante,
        spostiamo il meeting''. Il relatore si è sforzato di essere più
        \textbf{assertivo}: ``capisco la richiesta, ma ho qualcosa di più urgente
        da fare in questo momento''.
\end{enumerate}

\paragraph{E le richieste a voce?} Uno studente chiede se una richiesta fatta a
voce da un senior (``devi fare questo'') vada accettata così o fatta mandare per
mail o per messaggio. Accettarla a voce ci può stare. Ma se la richiesta è ``devi
farla in questo modo particolare e preciso'', è meglio farsi mandare il modo
preciso \textbf{per iscritto}: non si deve lasciare a interpretazione \emph{come}
fare quella cosa, per non ricadere nel ``ah, ma io non la volevo così''. Oppure si
conferma: ``sì, la faccio; va bene se la faccio così, così e così?''. Se dice sì,
si è a posto.

\paragraph{Non dovrebbero pensarci gli organizzatori?} Dalla chat: queste domande
non dovrebbero farsele già gli organizzatori dei meeting? Sì, sarebbe meglio che
tutti lo facessero, e il relatore stesso se le fa quando invita qualcuno: ``è
davvero necessario Marco? È davvero necessario Andrea?'', e solo se sì li invita.
Ma lavora in aziende anche di cinquecento dipendenti, dove non è detto che tutti
conoscano o applichino questa logica; per questo la applica per sé, sperando che
la applichino tutti. Ha sempre chiarito con tutti che da lui non arriveranno mai
inviti di cortesia o dettati dalla gerarchia (``deve esserci anche il mio capo''):
``devi esserci tu; se vuoi portarti il capo, portatelo, ma non lo invito solo
perché qui si vuole sempre il capo presente: sono sicuro che avrà di meglio da
fare''.

\paragraph{E i task?} Dire di no vale naturalmente anche per i task: si applica la
matrice di Eisenhower (sezione~\ref{sec:z2-eisenhower}) e si decide se fare,
delegare, pianificare o cancellare. Non c'è molto da aggiungere.

\subsection{Il multitasking non esiste}

``Questo potrà urtare tanti ego'', ma \textbf{il multitasking non esiste}. Si può
essere più bravi o più portati di altri a fare più cose contemporaneamente; la
realtà è che, concentrandosi su più cose, \textbf{la somma non farà mai cento}. Se
si lavora su due cose non si lavora al 50\% su una e al 50\% sull'altra, ma al 40\%
e al 40\%, e il 20\% va perso in ``costo di manutenzione'', cioè nella capacità di
seguirle entrambe. Il consiglio è dedicarsi al cento per cento a un task, poi
passare al successivo, poi al successivo, invece di farne cinque insieme: si può
avere la percezione di andare più veloci, ma non è così.

Uno studente cita una lettura di neuroscienze secondo cui il cervello è in grado
di fare una sola operazione alla volta. Il relatore conferma: smettiamola di
parlarne come di un pregio o di una \emph{skill}. È un modo che alcuni cervelli
trovano per svagarsi, l'equivalente di accendere un podcast: fare due cose
insieme vuol dire cercare di distrarsi alternandole, e farle entrambe peggio di
quanto si farebbero una alla volta. A chi preferisce fare tante cose insieme
suggerisce di provare a farne una alla volta: le farà più velocemente. Anche il
semplice passaggio da una cosa all'altra non è istantaneo: ci sono sempre una
decina di secondi in cui il cervello si chiede ``che cosa stavo facendo?''; per
cinquanta cambi, già questo spiega la perdita di velocità. È una questione di
\textbf{bandwidth mentale}, un concetto che tornerà nel pomeriggio: anche lo
\emph{switch} occupa banda.

\subsection{La fine della giornata}

A fine giornata si \textbf{rimette a fuoco la giornata}: si riguardano la task list
e la \emph{must do list} e si conferma di aver fatto tutto. È letteralmente il
controllo prima di uscire di casa: chiavi, porta chiusa, cellulare, portafoglio.
Al relatore è capitato di dimenticare un \emph{must do} rimasto in lista per tutto
il giorno, ed entra in gioco quella che chiama la \textbf{sindrome del
tirannosauro}\footnote{Il riferimento è alla celebre scena di \emph{Jurassic Park},
in cui si dice che il tirannosauro vede la preda solo se si muove.}: se una cosa è
lì, non è detto che la si veda finché non si muove. Averla in lista tutto il giorno
non vuol dire averla percepita come da fare.

Può anche capitare di non farcela perché quel giorno sono arrivate così tante cose
urgenti e importanti da rendere inutili le stime: la giornata no capita e non c'è
nessun problema. Ma se era davvero un \emph{must do}, \textbf{va fatto comunque}:
è lì che scattano la richiesta di straordinario o il ``devo farlo comunque''.
Attenzione: se si chiude la giornata con un \emph{must do} ancora aperto e si
decide ``vabbè, lo faccio domani'', può starci, ma allora \textbf{non doveva essere
in quella lista}, perché era rimandabile.

Al relatore questo controllo serve anche ``per quieto vivere'': se chiude ogni
giornata con la \emph{must do list} completa, sa che magari non ha fatto tutto
quello che c'era nella to-do list, ma i \emph{must do} sì. ``Il mio, oggi, l'ho
fatto'': può stare tranquillo senza continuare a chiedersi se ha fatto quello che
doveva.

% ---------------------------------------------------------------------
\section{Domande sul time management}
\label{sec:z2-qa-tempo}

\paragraph{C'è un numero medio ``sano'' di \emph{must do} quotidiani?} No, perché un
conto è avere dieci \emph{must do} come ``rispondere a dieci mail'', un conto è
averne uno solo come ``inventarmi un nuovo gioco''. Dipende dalla dimensione: se ne
gestiscono di più se sono minuscoli, di meno se sono giganteschi. Ci sono lavori
che per definizione hanno tantissimi \emph{must do} quotidiani e lavori che ne hanno
uno al giorno. Si imparerà da sé qual è il numero giusto.

\paragraph{C'è un modo per stimare quanto si riesce a fare in una giornata?} ``Lo so
che non è una bella risposta'', ma l'unico modo è la propria \textbf{storia}:
sapere quanto si riesce a fare nel quotidiano. È normale e fisiologico che le
prime stime siano sbagliate, e probabilmente eccessivamente ottimistiche;
l'importante è rendersene conto. È letteralmente la \emph{seniority}: abituarsi a
fare qualcosa insegna quanto ci si mette. Va benissimo metterci di più all'inizio,
ed è normale che il tempo cambi man mano che si migliora: ciò che oggi si stima in
un'ora, tra sei mesi lo si stimerà in mezz'ora. Chi ha lavorato con Scrum lo sa:
su un progetto nuovo tutti sbagliano la dimensione delle \emph{stories} all'inizio,
perché nessuno sa quanto ci vorrà, e al terzo \emph{sprint} consecutivo è più
chiaro. Le stime sono stime, altrimenti si chiamerebbero ``inserire quanto tempo ci
metterai''.

Se però si dà una stima e poi il lavoro si rivela il triplo, il problema non è la
stima: è che \textbf{bisognava chiedere più informazioni} prima di darla. La
risposta giusta non è la stima, ma ``per darti una stima accurata mi serve sapere
questo, questo e questo''. Il problema non è sbagliare la stima, è non darsi i mezzi
per azzeccarla.

\paragraph{Come si fa a far entrare questi metodi nella routine? È solo
determinazione?} È un misto. Il relatore risponde da persona con ADHD, con la
tendenza a distrarsi e a fare un mucchio di \emph{overthinking} (``come dicevo ieri,
anche ordinare la pizza è overthinking''). È un po' disciplina, ma soprattutto far
diventare qualcosa un'\textbf{abitudine}: lo sforzo per fare le checklist era molto
più alto un anno prima, oggi è un'abitudine. Di solito servono \textbf{circa tre
mesi} perché qualcosa diventi un'abitudine\footnote{Stima come riportata in lezione.
Gli studi più citati sulla formazione delle abitudini indicano tempi molto variabili
da persona a persona, con una mediana intorno ai due mesi.}: la teoria dice che chi
inizia oggi ad andare in palestra ogni giorno, dopo tre mesi non penserà più ``oggi
devo andare in palestra'', ma la palestra farà parte della sua giornata. Non è più
uno sforzo, ma ``è così che lavoro''.

\paragraph{Chi fa davvero time management?} Prima di passare al brainstorming il
relatore torna alla domanda iniziale: qualcuno, ora che sa che cosa intendeva,
cambierebbe risposta? Uno studente risponde che lo fa, ma senza checklist, perché
non ne sente il bisogno (non gli capita spesso di dimenticare le cose); la
\emph{must do list} invece la segue sempre, e le cose meno importanti capitano nelle
giornate in cui ne salta qualcuna. Va benissimo: non esiste un unico modo, ognuno
troverà il proprio. Il relatore porta il suo esempio perché sa che è servito a molte
persone con cui ha lavorato; con chi faceva fatica ad applicare un metodo (come
chiedeva prima uno studente) ne hanno trovati altri, perché ciò che funziona per uno
non è detto che funzioni per un altro. L'importante è che \textbf{ognuno trovi il
proprio metodo}, perché la soluzione non è non fare time management: ci si sentirà
più ``liberi'' improvvisando la giornata, ma non la si sta sfruttando a dovere.

E non si parla dei video sul \emph{mindset} in cui ci si alza alle cinque, si fa la
checklist e si fa colazione leggendo le mail: organizzarsi meglio la giornata
\textbf{dà più tempo libero}. E per il relatore il tempo libero delle persone è
fondamentale, perché è ciò che permette loro di lavorare meglio.

% ---------------------------------------------------------------------
\section{Una parentesi: tempo libero, ferie e benessere}
\label{sec:z2-benessere}

\subsection{Le tre condizioni del relatore}

Il relatore è un sostenitore della \textbf{settimana lavorativa di quattro giorni},
che ha applicato in varie aziende; dove non ci è riuscito, ha introdotto almeno
\textbf{venerdì liberi}, due volte al mese o una volta al mese, cioè giornate in
cui non si fa niente. Come le pause brevi sono importanti per essere creativi,
anche le pause lunghe, cioè giornate intere, lo sono per essere creativi ed
efficienti. È lo stesso motivo per cui, dopo la Games Week, dava una settimana
libera alle persone (sezione~\ref{sec:z1-fasi-team}): ha bisogno che siano
riposate, e quella settimana ``persa'' è in realtà un investimento per le
settimane successive. In chat qualcuno lo paragona al ``metodo finlandese'': è un
misto tra quello e il metodo del buonsenso.

Il relatore viene spesso chiamato a \textbf{cambiare le culture aziendali}, ed è un
po' il suo focus. Come un allenatore che accetta una squadra solo se può portare il
suo staff (il preparatore, il tecnico), lui chiede tre elementi:
\begin{enumerate}
  \item possibilmente \textbf{quattro giorni lavorativi}, o altrimenti giornate
        libere per dare spazio alle persone;
  \item \textbf{ferie illimitate}: nessun limite massimo, ma un \textbf{limite
        minimo garantito}. Ferie illimitate non devono voler dire poter scegliere
        di non prenderle: ``devi prendere almeno X giorni, fino a infinito'';
  \item la \textbf{feedback culture}.
\end{enumerate}
Senza questi tre elementi non garantisce il risultato: è un modo di gestire le
aspettative fin dall'inizio, come si diceva il giorno prima.

\paragraph{La feedback culture.} Il perché è stato spiegato nella prima giornata, ma
il relatore lo sottolinea: un ambiente in cui ci si può fidare del team vale
moltissimo. Quando va a fare lezione allo IULM, anche in settimana, non deve
preoccuparsi di che cosa stia facendo il team mentre lui non lavora (``oggi è
domenica e non ho questa preoccupazione a prescindere, ma anche se fosse mercoledì
non l'avrei''): sa che stanno facendo del loro meglio, e non gli interessa sapere
altro.

\paragraph{Ferie e permessi illimitati.} Nell'esperienza del relatore, sia a Riot sia
a Blizzard, avere ferie illimitate significava spesso \emph{prenderne meno} di quelle
obbligatorie: c'è una vocina che dice ``ti serve davvero un giorno di ferie oggi?
Tanto sono illimitate, te le prenderai quando ti serviranno'', e sapere di averne
illimitate spingeva a lavorare di più. Per lui era altrettanto dannoso, e per questo
ha introdotto un \textbf{minimo obbligatorio di 24 giorni} di ferie.

C'è stato però un effetto non previsto, e molto desiderato: in tutte le aziende in
cui ha introdotto ferie e permessi illimitati, \textbf{i giorni di malattia sono
calati del 70\%}\footnote{Dato riferito dal relatore sulla base della propria
esperienza, non verificabile.}. La gente si ammala di meno? O piuttosto il giorno in
cui ci si sveglia e si pensa ``oggi non ne ho voglia, mi do malato'' diventa il
giorno in cui si prende un permesso? Si smette di usare come malattia un giorno che
non era di malattia.

\paragraph{L'\emph{Hangover Day}.} Sia a Riot sia a PG Esports il relatore aveva
introdotto quello che a Riot, ``essendo Riot a Dublino'', si chiamava
\textbf{Hangover Day}\footnote{Nella trascrizione ``Angover Day''.}: una giornata
all'anno in cui si aveva il diritto di telefonare alle nove del mattino e dire ``oggi
non vengo perché non ne ho voglia'' (tradotto: perché ieri sera ho bevuto troppo).
Altrimenti quella giornata sarebbe diventata un giorno di malattia o di ferie. Il
fatto è che \textbf{quasi nessuno lo usava}: il relatore stesso a volte ci ha
pensato, ma poi la mattina si chiedeva ``voglio davvero giocarmelo oggi? Tutto
sommato potrei anche andare'', e ci andava. Ma anche solo sapere di averlo lo faceva
stare tranquillo. Allo stesso modo, con le ferie illimitate, chi voleva prendersi un
mese lo faceva (``buon per loro''), chi non lo faceva andava bene lo stesso:
l'importante era che staccassero. C'era un \textbf{benessere più diffuso} nel team,
perché sapevano che, volendo, avrebbero potuto.

\paragraph{Il \emph{Neko Day}.} Il primo punto, i giorni liberi, era chiamato
\textbf{Neko Day}, da \emph{neko}, ``gatto'' in giapponese: il giorno in cui si può
fare il gatto, rannicchiati davanti al camino a non fare niente. Perché
introdurlo? Chi lavora dal lunedì al venerdì passa di solito il sabato o la domenica
a fare le commissioni che non può fare nei giorni lavorativi: la spesa, le bollette,
il bollo dell'auto. In questo modo sfrutta meno il weekend. Con un venerdì libero al
mese (o due, o tre) le persone facevano quel giorno le commissioni del weekend, e
sabato e domenica si riposavano davvero (``anche non fare niente andava bene,
chiaramente''). Il lunedì arrivavano molto più cariche e riposate. Questo ha aiutato
tantissimo a mantenere un team sano, fisicamente e mentalmente.

\paragraph{Il supporto psicologico.} A PG Esports il relatore ha introdotto anche il
\textbf{supporto psicoterapeutico offerto dall'azienda}: sei sedute da un'ora con
uno psicoterapeuta, online. Era il periodo in cui si parlava del bonus psicologo
statale, ``che poi non è mai arrivato''\footnote{Così in lezione. In realtà il bonus
psicologo è stato istituito in Italia nel 2022, ma con fondi limitati rispetto alle
richieste e con forti ritardi nell'erogazione; il relatore si riferisce
probabilmente a questo.}, e lui ha pensato: l'idea è buona, perché non introdurla a
livello aziendale? Erano gli anni dei lockdown: l'azienda stava perdendo molte persone
che andavano in depressione o non avevano più voglia di lavorare, perché avevano
difficoltà a vivere quel periodo (``io stesso non è che l'abbia vissuto benissimo'').
Qualcuno ne ha usufruito e gli ha detto: ``mi stavo per dimettere, poi ho risolto le
cose e ho preferito restare''.

Secondo il relatore oggi c'è una consapevolezza maggiore da parte delle aziende
sulla \textbf{salute mentale}: vuol dire dare più ferie, dare questo tipo di
supporto, dare più \emph{benefit} (anche i buoni pasto). Il benessere non è solo
economico, è anche mentale, ed è una cosa positiva da sapere.

% ---------------------------------------------------------------------
\section{Brainstorming: obiettivo e problemi tipici}
\label{sec:z2-brainstorming}

\subsection{Perché il brainstorming riguarda tutti}

Secondo il relatore questa è la parte che servirà di più, sia nei gruppi e nei
progetti del percorso di studi (il gioco da tavolo, il gioco digitale) sia nelle
aziende. Rispetto al time management ha una pertinenza più immediata: tutti
hanno esempi di \textbf{brainstorming andati male}, da cui si è usciti esausti, con
la sensazione di non aver ottenuto niente, o di non aver parlato abbastanza (``ha
parlato solo quello''), o di non essere sicuri della scelta. Il brainstorming lascia
una sensazione molto chiara, a fine meeting, che qualcosa non torna; il time
management dà piuttosto la tristezza del ``non ce l'ho fatta a fare quello che
dovevo oggi, come mai?''.

\subsection{Generare idee, non decidere}

Il primo chiarimento riguarda \textbf{l'obiettivo}, che il relatore vede spesso
travisato. L'obiettivo del brainstorming è \textbf{generare idee}; invece diventa
spesso decisionale, con lo scopo di scegliere che cosa fare. Ma allora non sarebbe
un brainstorming, sarebbe un \emph{decision making meeting}. Generare molte idee e
scegliere la migliore nello stesso meeting non funziona né per generare idee né per
scegliere la migliore, perché sono \textbf{due metodi che si contrappongono}. Il
brainstorming va fatto quando si devono generare idee \emph{per poi} scegliere che
cosa fare: deve venire \textbf{prima} del \emph{decision making}, come primo passo
della fase di esplorazione di un nuovo progetto.

\subsection{``Non esistono cattive idee''?}

Bisogna evitare che i membri del team \textbf{autofiltrino} le proprie idee per paura
di essere rifiutati o presi alla leggera. Da qui nasce la frase ``non esistono
cattive idee, buttate fuori tutto'', che è vera solo fino a un certo punto. Le
cattive idee esistono, oggettivamente; quello che non si deve fare è
\emph{colpevolizzarle}. Né si deve per forza perseguire ogni opzione: parlare di
opzioni che sono effettivamente brutte è uno spreco di tempo, e il brainstorming ha
anche un problema di gestione del tempo. Il relatore vede spesso brainstorming
gonfiati a ore e ore in cui si parla anche di cose che non meritano nemmeno dieci
minuti, come i feedback buoni ma inutili del giorno prima. \textbf{Trattare tutte le
idee allo stesso modo è sbagliato}, perché non hanno la stessa qualità. Bisogna
trovare un modo di valutarle senza che chi le ha proposte si senta preso di mira,
bullizzato o sminuito.

Per questo a ogni brainstorming dovrebbe seguire un \textbf{follow-up}, il momento in
cui si decide che cosa fare: le due parti vanno separate, altrimenti ``si crea solo
casino''. Il vero lavoro arriva quando le idee vengono messe a terra: c'è la parte di
brainstorming e di \emph{theorycrafting}, e la parte esecutiva viene dopo.

\subsection{I problemi più comuni}

\begin{description}
  \item[Le prime idee monopolizzano la sessione] Le prime due idee ``sparate a
        caso'' sono quelle su cui ci si concentra per la maggior parte del tempo,
        sviscerandone i dettagli (``intendevi così o così?''), per poi chiudere con
        ``altre idee? Questa, questa e questa? Sì, no, ok, andiamo avanti''. Si
        parla troppo delle prime idee e poco delle ultime.
  \item[Poche persone parlano per la maggior parte del tempo] È un effetto
        collegato al precedente.
  \item[\emph{Anchoring}] Ci si lega alle prime idee, e spesso si votano o si
        scelgono non perché siano le migliori, ma perché sembra di conoscerle meglio
        di quelle menzionate in trenta secondi alla fine. Non c'è nessun motivo per
        pensare che un'idea detta alla fine sia peggiore: era solo arrivata dopo.
        C'è un problema di \emph{timing}. Anzi, un'idea venuta istantaneamente è
        forse quella meno ragionata, perché ci si è pensato meno tempo.
  \item[Silenzi imbarazzanti] Spesso le persone arrivano al meeting senza nemmeno
        sapere che è un brainstorming. Quando il relatore convoca un brainstorming,
        comunica prima a tutti il \emph{topic}, il tipo di idee cercate, i paletti e
        le restrizioni, così che anche chi è meno portato all'improvvisazione arrivi
        con qualche idea pronta. Se la gente arriva impreparata rischia di non dire
        niente per un'ora, e allora perché è stata invitata? Non è un difetto loro:
        tanti manager dicevano al relatore ``ci sono persone inadatte ai
        brainstorming perché non sanno improvvisare''. No: si sta facendo
        brainstorming con un metodo che premia solo chi sa improvvisare ed è
        abbastanza estroverso da parlare sul momento. Non significa che chi non
        parla non avrebbe potuto contribuire, ma che non gli è stato dato modo di
        farlo.
  \item[Estroversi e pensatori rapidi dominano] Le personalità estroverse e chi
        pensa velocemente tendono a dominare la conversazione, più o meno
        intenzionalmente, senza lasciare agli altri il tempo di contribuire. Non c'è
        uno spazio in cui un introverso possa dire ``veramente io volevo\dots'':
        ``no, no, stiamo parlando di questa bellissima idea, facciamo questa''.
\end{description}

% ---------------------------------------------------------------------
\section{Un metodo di brainstorming}
\label{sec:z2-metodo}

Il relatore propone una soluzione principale, declinata in vari modi, che tocca
tutti questi punti.

\subsection{Preparazione e ambiente inclusivo}

\paragraph{Arrivare con un obiettivo chiaro.} Bisogna arrivare al brainstorming con
chiaro e preciso che cosa si sta cercando, cioè \textbf{che cosa si sta cercando di
risolvere}, \textbf{perché} e soprattutto \textbf{con quali vincoli}. Si stanno solo
raccogliendo idee per un cliente? Si cercano idee fattibili, o idee innovative? Che
cosa si vuole produrre da questo meeting? Deve essere chiaro a tutti \emph{prima}
del meeting, non spiegato durante, perché altrimenti si ricade nel problema delle
persone non pronte a intervenire.

\paragraph{Non giudicare le idee degli altri.} Al relatore è capitato spessissimo di
sentire un'idea terribile da una persona e, mezz'ora dopo, l'idea perfetta dalla
stessa persona. Ma se all'idea terribile tutti rispondono ``che stupidaggine, per
favore stai zitto'', quella persona non dirà la successiva, perché non vuole
rivivere l'esperienza. Giudicando le idee degli altri, di fatto, li si zittisce.
Le idee migliori nascono quando si propone quella che si pensa sia l'idea perfetta,
magari non è buona e non viene scelta (``fa niente''), si ascoltano anche le altre e
poi la si \textbf{riformula} integrando i pareri sentiti nel frattempo. Sono queste
\textbf{idee combinate} a essere spesso le soluzioni migliori.

\subsection{Separare generazione e discussione: \emph{brainwriting} e
\emph{rapid ideation}}

Con il classico ``non ci sono cattive idee, sparate'', succede spesso che uno prenda
la parola, spieghi la sua idea e il suo concetto, e gli altri comincino subito a
fargli domande (``ma quindi la vedi così?''): si sta unendo la persona che
\emph{genera} l'idea con la persona che la \emph{discute}, nello stesso momento.
Invece bisogna \textbf{dividere la generazione dalla discussione}: un momento in cui
si generano idee e un momento in cui si discutono. È il principio del
\textbf{brainwriting}. Funziona molto quando nel team c'è una figura con una
posizione gerarchica più alta: se quella persona esprime un'idea e i due momenti
sono uniti, tutti tenderanno ad allinearsi. E si vuole che \emph{tutti} possano
avere idee.

Lo strumento è la \textbf{rapid ideation}. Il relatore di solito dice: ``per i primi
cinque minuti non parliamo; scrivete su dei post-it il pitch delle vostre idee''. Se
ne possono scrivere quanti se ne vuole, dieci o uno, e si mettono da parte. Ma il
tempo è di \textbf{cinque minuti}: non dieci, non due, e soprattutto non mezz'ora.
Il \textbf{limite di tempo breve} è importantissimo, perché costringe,
consapevolmente o meno, a filtrare le idee mentalmente e a scrivere solo quelle che
si ritengono migliori: magari non è chiaro perché si preferisca un'idea all'altra,
ma il cervello sta facendo il lavoro in background e fa venire in mente le cose più
importanti o più giuste. È ottimo anche per i team che tendono a distrarsi, perché
sono solo cinque minuti.

\subsection{\emph{Dot voting}}

Poi si passa al \textbf{dot voting}:
\begin{enumerate}
  \item si attaccano \textbf{tutte} le idee su una lavagna fisica, un muro o una
        lavagna digitale (come Miro\footnote{Nella trascrizione ``mirror board'':
        si tratta della lavagna digitale Miro, già citata nella prima
        giornata.});
  \item si \textbf{uniscono le idee che hanno lo stesso concetto di base}. Se si sta
        pensando a un evento dal vivo e la proposta del relatore è un parco di
        dinosauri, mentre quella di un collega è un parco di dinosauri ``del
        paleolitico'' con i tirannosauri, le si raggruppa sotto ``parco di
        dinosauri'': non sono due idee separate. Un altro gruppo potrebbe essere,
        ad esempio, un parco a tema Avengers;
  \item ognuno ha \textbf{da tre a cinque voti} e va alla lavagna a mettere i propri
        pallini sulle idee che gli piacciono e di cui vorrebbe sapere di più.
\end{enumerate}
Il risultato tipico: un'idea con molti voti, una con un paio, altre con un numero
discreto.

\paragraph{Che cosa si vede subito?} Il relatore lo chiede alla classe. Si vede una
\textbf{preferenza} per un tema o per un'idea. Non c'è \textbf{polarizzazione} su una
o due idee iniziali, perché c'è varietà. Il voto è \textbf{anonimo}: non ci sono nomi.
E soprattutto manca un'altra cosa importantissima: \textbf{il giudizio negativo
espresso}. Se c'è un'idea con sette voti, una con cinque e una con tre, è
ragionevole parlare prima di quelle, perché tutti hanno espresso interesse e vogliono
saperne di più; ma nessuno sta dicendo ``non voglio assolutamente parlare di questa
perché mi fa schifo''. È semplicemente oggettivo che quell'idea non ha ricevuto voti.

\paragraph{Autovotarsi.} Il sistema permette anche di votare la propria idea, e non c'è
nessun problema: se la propria idea ha un pallino e gli altri ne mettono sette su
un'altra, il proprio voto è ininfluente. Il relatore vedeva spesso che all'inizio la
gente tendeva ad autovotarsi, poi capiva che era inutile, perché non sarebbe mai stato
il voto decisivo, e tanto valeva votare le idee che interessavano davvero. La
tendenza spariva con il tempo.

\paragraph{Che cosa si è ottenuto.} Tre ore di meeting sono state filtrate in una
discussione concentrata solo sulle idee che tutti ritengono interessanti, con un
\emph{agreement} di fondo: non si parla dell'idea numero uno perché è uscita per prima,
non è una questione di \emph{timing} ma di qualità. Il relatore ha vissuto tanti
meeting con la frustrazione interna del ``stiamo sprecando tempo su questa idea,
quando vorrei dire la mia che secondo me è meglio''. Qui non si può sprecare tempo,
perché si parla di qualcosa che tutti vogliono sentire.

\paragraph{E se si formano sottogruppi che si votano a vicenda?} Uno studente chiede se
non possano esistere ``cordate''. Anche se esistessero, significherebbe che il gruppo di
quel meeting vuole parlare di quella cosa, e va bene: l'importante è non sprecare tempo
su idee che non interessano a nessuno. Se si investono due ore, quarantacinque minuti
su ciascuna di tre idee vanno benissimo; dieci minuti per ogni idea no. E comunque il
voto è anonimo: non si può sapere che un'idea è di Gianni per dare supporto a Gianni,
quindi un eventuale sostegno reciproco sarebbe involontario.

\paragraph{E chi è lento a elaborare?} Uno studente osserva che il problema degli
introversi potrebbe ripresentarsi con i post-it: con un tempo così ristretto, chi è più
lento potrebbe andare ancora più in panico. Per questo bisogna \textbf{arrivare
preparati}: si conosce il topic in anticipo, quindi almeno un paio di idee già
formalizzate si hanno in mente. Magari non se ne scriveranno quindici, ma un paio sì.

\subsection{Quanti voti e quali varianti}

Dare più pallini a testa sarebbe controproducente? Il relatore indicava da tre a
cinque voti. Chiede alla classe quanti siano i gruppi di lavoro: dai quattro agli otto,
alcuni gruppi del secondo anno sono nove. La sua prassi:
\begin{itemize}
  \item \textbf{dai cinque partecipanti in su}, \textbf{tre voti} a testa, più che
        sufficienti per popolare molto la lavagna;
  \item \textbf{fino a cinque partecipanti}, \textbf{cinque voti} a testa, che
        permettono anche di variare un po'.
\end{itemize}
Esistono varianti: ad esempio tre voti di colori diversi, il blu per l'idea più
innovativa, il verde per quella più bella, il giallo per quella più facile, così da
qualificare \emph{perché} un'idea piace. Se l'obiettivo è proporre a un cliente tre
pitch (uno facile, uno medio, uno avanzato), la votazione è già fatta. All'inizio,
però, meglio partire da tre o cinque voti generici; con il tempo si imparerà ad
accelerare.

\subsection{Il moderatore e gli invitati}

\paragraph{Serve un moderatore?} Può essere utile. Il relatore fa spessissimo il
moderatore: partecipa ai brainstorming senza fare altro che moderare, ``perché so che
le persone con le idee fighe sono loro e non io''. Ci si può anche autogestire. Nei team
piccoli non ha senso ``perdere'' una persona per la moderazione; nei team più grandi
(dai cinque in su, sette-otto persone), se qualcuno dice ``non penso di avere idee
particolari su questo, vi modero così vi aiuto'', è un ottimo modo di contribuire.

\paragraph{Chi invitare?} Dipende da che cosa si cerca. Per inventare un nuovo gioco il
relatore invita i game director, gli art director e magari il \emph{publishing manager}
per un input commerciale; per decidere il \emph{combat system} invita magari il senior
game designer e il game director. La domanda che si fa chi convoca il meeting è: dato
l'obiettivo, chi è più ragionevole che possa contribuire?

Un 3D artist della classe chiede se allora gli artisti 3D non vengano inclusi. Dipende:
se l'obiettivo è definire la linea creativa dello stile dei personaggi, il relatore
invita anche il 3D artist, e non solo i \emph{concept artist}. Il 3D artist può dare un
input diverso dal concept: vota un'idea perché per lui è \emph{più fattibile}. Un
concept può essere bellissimo ma ``un inferno'' da modellare, e chi sa di doverlo fare
entro giugno non vota un'idea che sa già infattibile.

In realtà si può invitare chiunque: per un pitch su un'area dedicata al basket alla
Games Week, il relatore ha invitato un tecnico di regia che gioca a basket, per avere
anche la prospettiva di chi ha familiarità con il tema. A uno studente che porta
l'esempio di un elemento di interfaccia sul personaggio, il relatore risponde che,
volendo, si trova un motivo per invitare chiunque: l'importante è avere chiaro
\textbf{chi è necessario}. Il gruppo di lavoro può essere di otto e il brainstorming di
dodici, con input esterni a chi poi farà il lavoro. Basta non esagerare: se la maggior
parte dei voti arriva da chi non farà quel lavoro, si vota su qualcosa che poi non
faranno loro, e non ha più senso.

% ---------------------------------------------------------------------
\section{Dalla votazione alla decisione: la matrice impatto--sforzo}
\label{sec:z2-impatto}

\subsection{Massimo impatto, minimo sforzo}

Fatto questo, non si è finito: ora si sa quali idee interessano di più, ma non si è
ancora deciso che cosa fare. Serve la parte decisionale. L'obiettivo di qualsiasi
azienda (ed è il motivo per cui al relatore non importa se il programmatore fa il test
da casa cercando su Google) è \textbf{massimizzare l'impatto e minimizzare lo sforzo}.
Vale per i progetti, i task, il recruitment, per qualsiasi cosa: ``non perché sono
pigro, ma perché sarebbe stupido fare diversamente''.

Le idee selezionate si mettono quindi in un'altra matrice, con due assi:
\begin{itemize}
  \item sull'asse verticale il \textbf{valore}, cioè quanto \textbf{impatto} avrebbe
        l'idea (in alto l'impatto alto);
  \item sull'asse orizzontale lo \textbf{sforzo} (\emph{effort}), da basso (a sinistra)
        ad alto (a destra).
\end{itemize}
In alto a sinistra ci sono quindi le idee ad alto impatto e basso sforzo, in alto a
destra quelle ad alto impatto e alto sforzo, e così via.

\paragraph{L'esempio Adidas.} Se il cliente è Adidas, con la linea di abbigliamento per
il basket, un'idea a basso valore potrebbe essere una campagna su TikTok in cui gli
utenti condividono i loro momenti con un pallone da basket; un'idea ad alto valore
potrebbe essere fondare un parco divertimenti dove i giocatori NBA vengono a sfidarti
uno contro uno. Questo non dice che una sia meglio dell'altra: dice che l'impatto è
diverso. Allo stesso modo si valuta lo sforzo: la campagna social è a basso sforzo, il
parco divertimenti ad alto sforzo per molti motivi (i permessi di costruzione, i
giocatori NBA).

\subsection{Posizionare le idee, insieme e con precisione}

Il gruppo che a quel punto è \textbf{il gruppo decisionale}, cioè chi farà il lavoro,
deve posizionare sulla matrice le idee ritenute migliori, e deve essere \textbf{d'accordo
tutto il gruppo}: ``siamo tutti d'accordo che il parco dei dinosauri di Stefano e Andrea è
ad alto valore ma anche ad alto sforzo? Sì? Allora lo mettiamo qui''. E bisogna indicare
\textbf{un punto preciso}, non solo il quadrante: piazzare un'idea in un punto non è come
piazzarla poco più in là. Il risultato è una sorta di \textbf{mosaico} delle idee sul
grafico.

\subsection{Leggere la matrice in diagonale}

Davanti a una situazione tipica, con varie idee sparse, il relatore chiede alla classe:
se foste i project manager, se doveste fare voi queste cose, \emph{quali} fareste? Al
plurale, perché spesso da un brainstorming si esce con più idee realizzabili e si deve
scegliere su quali concentrarsi. Le risposte variano: c'è chi sceglie le due con più
impatto, chi aggiunge che dipende da tempi, budget e risorse, chi preferisce quella più
innovativa ma anche più rischiosa.

Il metodo che secondo il relatore va applicato è \textbf{procedere in diagonale}, dalle
idee a basso sforzo e alto impatto verso quelle ad alto sforzo e basso
impatto\footnote{Nella slide le idee erano rappresentate da post-it colorati (rossa,
rosa o salmone, azzurra, verde acqua, viola, blu, gialla), e le indicazioni del
relatore si riferiscono a quei colori. Dalla sola trascrizione non è possibile
ricostruire con certezza la posizione di ciascun post-it, e alcune indicazioni sui
colori risultano contraddittorie; qui è riportato il principio generale, che il
relatore chiarisce nelle risposte alle domande.}:
\begin{enumerate}
  \item un'idea a \textbf{basso sforzo e alto impatto} (nella slide, quella rossa) va
        fatta \textbf{per prima, sempre}: ``non dovrebbe essere nemmeno in discussione'';
  \item poi, a seconda delle risorse e del tempo, si possono fare
        \textbf{contemporaneamente} le idee che stanno sulla stessa diagonale
        successiva, che hanno cioè un rapporto simile tra impatto e sforzo;
  \item poi le successive;
  \item le idee ad \textbf{alto sforzo e basso impatto} (nella slide, la gialla)
        \textbf{non vanno nemmeno considerate}. Nessuno, infatti, l'aveva scelta.
\end{enumerate}
Per le scelte intermedie valgono due criteri, chiariti rispondendo a uno studente che
non capiva perché alcune idee fossero preferibili ad altre: \textbf{a parità di sforzo si
sceglie l'idea con più impatto}; \textbf{a parità di impatto si sceglie quella con meno
sforzo}. Con risorse infinite si farebbero anche le idee più costose; ma di solito le
risorse sono limitate, e allora tra due idee con lo stesso sforzo e un impatto molto
diverso la scelta è evidente. Se ci piacerebbe fare un'idea più ambiziosa ma abbiamo
risorse solo per una meno costosa sulla stessa diagonale, va bene lo stesso: abbiamo
comunque stabilito che è meglio delle idee sulle diagonali successive. Quanto fare dipende
dalla \textbf{bandwidth} a disposizione.

\paragraph{``La gialla era la mia idea: devo offendermi?''} No. Qualcuno ha detto che
l'autore fa schifo? No. Si è concordato che la sua idea, per quanto votata (``qualcuno ha
votato per la mia idea''), è piuttosto difficile da realizzare rispetto al suo impatto.
Non si sta dicendo che non sa più fare il suo lavoro. Anzi, se lui stesso dovesse
scegliere tra la propria idea gialla e la rossa di un altro, sceglierebbe la rossa: non è
nemmeno una questione di discussione. La matrice rende \textbf{visivamente evidente} che
cosa fare, e permette di aggirare molte discussioni che non sono più soggettive:
oggettivamente, dovremmo fare così.

\subsection{Il ruolo del moderatore e dell'esperto}

È qui che lavora il moderatore: ``non è più una discussione; visto il posizionamento,
dobbiamo fare queste cose''. Non si decide a maggioranza dove mettere le idee: se uno
posiziona un'idea e un altro non è d'accordo, difficilmente dirà ``la rossa dovrebbe stare
al posto della gialla''; dirà ``secondo me è un po' più a destra, un po' più in basso'', e
si fanno gli aggiustamenti. Se invece nascono discussioni del tipo ``la rossa dovrebbe
stare al posto della gialla'', vuol dire che non si è definita bene l'idea e non è chiaro
quanto sia davvero difficile da realizzare.

Oppure interviene \textbf{l'esperto}: è l'esempio del concept artist e del 3D artist. Il
concept artist dice che un'idea è a basso sforzo e alto impatto; il 3D artist risponde ``ne
condivido l'impatto, ma dovendola fare io ti assicuro che lo sforzo non è basso, è molto
alto''. Allora è un altro discorso, perché l'esperto sta dicendo che lo sforzo è diverso
da quello stimato. Ma alla fine tutto deve essere concordato.

\subsection{Chi decide}

\paragraph{Il gruppo decisionale può essere diverso da quello del brainstorming?} Sì: a
questo punto escono quelli invitati come ``consulenti'', che hanno dato idee ma non devono
poi realizzarle. Farli votare anche qui significherebbe perdere il collegamento tra chi
deve fare le cose e che cosa si decide di fare. In quel momento non hanno responsabilità:
hanno dato un input, e se decidessero che cosa devono fare gli altri dovrebbero esserne
i manager. Non dovrebbero poter dire agli altri che cosa fare se non ne hanno la
responsabilità. Vale per tutti i team, piccoli e grandi; secondo il relatore chi decide
dovrebbe essere \textbf{al massimo una decina di persone}. In un team di dieci che fa il
brainstorming tutto insieme, va benissimo che chi vota coincida con chi ha partecipato:
l'importante è che \textbf{chi vota sia chi farà il lavoro}.

\paragraph{Lo sforzo è il costo o la fatica?} Uno studente chiede se l'alto impatto
includa anche il costo. No, il costo sta nello sforzo, che per il relatore va
\textbf{ragionato in risorse}: denaro, tempo e così via. Un'idea può avere un basso costo
monetario ma un alto costo in tempo, e allora va bilanciata e si troverà a metà. È lì che
si spostano le idee e che avviene la conversazione: i tecnici dicono che lo sforzo tecnico
è un certo, poi arriva l'\emph{accountant} che dice che a livello economico è un altro, e
la si sposta.

\paragraph{E il piacere di fare il progetto?} In un gruppo di due o tre persone, osserva
uno studente, si finisce per votare anche per preferenza personale: quando si fa un gioco
bisogna avere anche il piacere di farlo. Ma \textbf{anche quello è impatto}: l'impatto
sulla motivazione. Come lo sforzo sono le risorse, l'impatto è l'impatto \emph{su di te,
sul team, sul mercato}, su ciò che si vuole. Il relatore lo ha tenuto volutamente
generico: può essere quanti soldi si fanno o quanto è utile per il proprio CV, e
l'impatto cambia tra un progetto accademico e uno professionale.

% ---------------------------------------------------------------------
\section{Consigli pratici e domande sul brainstorming}
\label{sec:z2-qa-brainstorming}

\subsection{Musica e cambio di scenario}

``Sembrano tutte stupidaggini, ma vi assicuro che aiutano.''

\paragraph{Musica di sottofondo.} È dimostrato che aiuta\footnote{Le affermazioni sugli
effetti della musica (chiave maggiore, rumore ambientale) sono riportate come dette in
lezione e non sono state verificate.}. Meglio una musica con \textbf{tempo e volume
costanti}: niente \emph{bass drop}, niente variazioni di tema o di ritmo. Va bene la
musica fatta di rumore ambientale a basso volume; è stato rilevato anche che le canzoni
in \textbf{chiave maggiore} aiutano la parte creativa a produrre idee più originali.
Meglio la \textbf{musica strumentale}, senza testo: per quanto ci si impegni a non
ascoltarlo, il cervello cercherà di capire le parole. È un effetto marginale e non
succede niente se capita la canzone con il testo, ma si mette in conto che ``l'uno per
cento del cervello'' sarà impegnato a capirlo, come quando in treno si ascolta, che lo si
voglia o no, la telefonata di qualcun altro. È \emph{min-maxing}.

\paragraph{Cambiare scenario.} A Riot c'era una stanza dedicata al brainstorming, pensata
per essere creativa: le finestre sul fiume, le porte a vetri colorati, tutto fatto per
ispirare creatività. Ma fare brainstorming sempre nella stessa stanza generava sempre lo
stesso tipo di idee: anche solo cambiare ambiente aiutava ad averne di diverse. Il
relatore spesso spostava il meeting, ad esempio nello Starbucks sotto l'ufficio. Lo
spazio fisico gioca un ruolo importante, ma anche quello virtuale: se lo si fa sempre su
Discord, si può provare con ognuno uno sfondo diverso e simpatico. Oppure: ha sempre
moderato la stessa persona? Si prova con un altro moderatore. Si cambiano i fattori
perché i risultati non siano sempre gli stessi.

Uno studente chiede se il cambio di scenario possa essere anche qualcosa come passarsi una
palla per avere la parola, o un altro espediente per tenere viva l'attenzione. Con questo
modello, che non richiede di prendere la parola, l'esempio della palla perde senso, ma come
esempio va benissimo: anche il comportamento fisico e il \emph{modus operandi} sono un ottimo
cambio di scenario. Come lo sono cambiare il moderatore o \textbf{cambiare l'obiettivo}: se
il cliente chiede X, gli si darà un'idea X, ma si può provare anche a votare per un'idea Y,
facendo più votazioni con vari criteri di valutazione.

\subsection{Domande}

\paragraph{L'idea scelta deve essere approvata da qualcun altro?} Se un manager dice ``fate
questo brainstorming senza di me'', sta \textbf{delegando la scelta}: ``mi fido di quello che
sceglierete''. Avendo preparato prima limiti, paletti e \emph{constraint}, si aspetta che la
soluzione portata sia in linea con essi; ma se non partecipa deve fidarsi della scelta. Se
poi la scelta debba passare anche da un livello superiore dipende dall'azienda e dagli
accordi precedenti.

\paragraph{Si può fare un brainstorming ibrido, tra presenza e remoto?} Uno studente fa
l'esempio di un gruppo in cui uno è a New York e gli altri a Milano. Sì, senza problemi: si
fa una call, il portatile punta verso il muro con le idee, o le si mostra fisicamente con
il cellulare. Chi è da remoto scrive i suoi post-it, ad esempio mandandoli su WhatsApp al
moderatore o a qualcuno presente che li scrive per lui, e vota dicendo ``mettimi un pallino
qui, qui e qui'' (in quel caso il voto non è anonimo, ma si può fare tranquillamente); se
deve parlare, parla.

\paragraph{E se i colleghi non riconoscono o non condividono il metodo?} Ci deve essere un
\textbf{agreement}: non si può arrivare e dire ``facciamo così'' se gli altri non vogliono.
Questo metodo funziona solo se tutti sono d'accordo a usarlo. Il relatore ne ha proposto
uno tra i cento metodi di brainstorming esistenti, un buon compromesso tra semplicità ed
efficienza (ce ne sono di più semplici e di molto più efficienti ma più complessi). Se il
gruppo non ne vede il valore, al massimo insisterebbe per \textbf{fare almeno una prova},
altrimenti diventa chiusura mentale; se poi per quel team non funziona, se ne cercano
altri.

\paragraph{E nei gruppi molto piccoli, dove diventa un duello?} Uno studente ha lavorato in
gruppi di tre o anche di due persone, dove il brainstorming diventa quasi un duello in cui
ognuno difende la propria idea. Bisogna \textbf{aspettare la discussione}: invece di ``scegliamo
tra la mia idea e la tua'', si cerca di concordare quale idea abbia l'impatto maggiore e
quella lo sforzo minore. E se non si è d'accordo, bisogna \textbf{portare esempi}: se uno
dice ``questa idea costa di più'', l'altro chiede ``in che senso? Che cosa costa di più?''. Se
non sa rispondere, non si può fare qualcosa senza sapere come va a finire: meglio rimandare
la discussione a quando si avranno le informazioni. A volte la domanda giusta è \emph{perché
non abbiamo queste informazioni?}: forse è \textbf{prematuro} decidere, perché non si hanno
ancora tutti i fattori.

Nei team piccoli non dovrebbe esserci un ``si fa come dici tu o come dico io'', una scelta a
pelle: si dovrebbero conoscere così bene le proprie proposte da poterle posizionare con
assoluta certezza sulla matrice, e avere anche la presenza di spirito di dire ``la tua è
chiaramente migliore, si fa la tua''. Si fa la cosa migliore per entrambi, che dovrebbe
essere oggettiva e non soggettiva. Non è facile arrivarci, ma deve essere l'obiettivo del
team.

% ---------------------------------------------------------------------
\section{Disagree and commit}
\label{sec:z2-disagree}

\subsection{Quando la scelta non ci convince}

Uno studente teme che chi vede scartata la propria idea perda motivazione, in un gruppo
piccolo come in uno grande. Non è un deragliamento: è una conseguenza potenziale di
qualsiasi scelta. Il relatore ha visto gente demotivata perfino dalla settimana di quattro
giorni: ci sarà sempre qualcuno più o meno motivato dalle scelte. Ma qual era l'alternativa?
Un brainstorming classico in cui tutti parlano, chi più chi meno, e si sceglie un'idea che
tutti pensano sia la migliore senza esserne davvero sicuri, perché non si è votato, né
espresso una preferenza, né discusso davvero le idee migliori.

Qui entra il concetto di teamwork del \textbf{disagree and commit}: ``non sono d'accordo con
questa decisione, non mi piace, non penso sia la migliore, ma la faccio''. Per due motivi:
\begin{enumerate}
  \item perché \textbf{è il mio lavoro}, e dovrebbe essere il punto principale;
  \item perché \textbf{voi ritenete che sia la scelta migliore}. Se si è la voce fuori dal coro
        e tutti hanno scelto un'altra strada, non si rema contro: è \emph{default to trust}.
        Se gli altri vi hanno visto più potenzialità, e (grande premessa) \textbf{hanno
        ascoltato} il proprio disaccordo e i suoi motivi, e hanno comunque scelto così, ci si
        fida di loro.
\end{enumerate}
Con una postilla: se le cose vanno male, \textbf{non si dà la colpa agli altri} e non si dice
``ve l'avevo detto'', come se non si facesse parte del team. Disagree and commit vuol dire ``mi
fido del team e facciamo questa cosa''; se poi si rivela sbagliata, \emph{il team} ha
sbagliato, e la volta dopo si cambierà il modo di scegliere. Altrimenti non è disagree and
commit, è ``\emph{disagree and ostracize}'': disaccordo più ostilità. Non bisogna cercare il
colpevole, perché si è deciso tutti insieme (sezione~\ref{sec:z1-default-trust}).

Se si pensa davvero che si stia sbagliando clamorosamente il posizionamento di un'idea, si
dice piuttosto ``non andiamo avanti finché non definiamo bene lo sforzo di questa cosa''; ma
bisogna saperlo giustificare, altrimenti è chi dissente a dover fare un passo indietro.

\subsection{La demotivazione fa parte del lavoro}

La demotivazione ci sarà sempre, perché non si sceglierà mai solo ciò che piace: per quanto
piaccia il proprio lavoro, ci saranno task sgradevoli e scelte che non spettano a noi. E
ci saranno scelte (``ve lo dico per esperienza personale'') che si è convinti siano
sbagliate, al punto di pensare ``piuttosto lascio il team'', ma che si finisce per accettare,
per poi rendersi conto che gli altri avevano ragione. Come c'è un \emph{bias} verso l'amico e
il nemico, c'è anche un \textbf{bias verso noi stessi}: crediamo di conoscere bene una cosa, e
se altre dieci persone dicono ``secondo me è fattibile'', bisognerebbe aprirsi alla
possibilità che ne sappiano più di noi.

``Non lo accetto ma lo faccio'' dovrebbe essere alla base del lavoro insieme, perché è
\emph{default to trust}: rifiutarsi (``non mi fido di voi, non lo faccio'') significa non
fidarsi del giudizio del proprio team. E dire ad alta voce ``non sono d'accordo, ma ti darò
comunque il massimo'' è la cosa che \textbf{fa guadagnare di più il rispetto degli altri}.
Non c'è niente di meglio che vedere qualcuno che aveva un'opinione diversa impegnarsi per
far vedere di avere torto, facendo andare bene la cosa. Per il relatore dimostra solo che è
una persona matura e importante; ha avuto persone che, passate da questa situazione, gli
hanno fatto pensare ``è molto più matura di quanto credessi'', e che ha valorizzato molto di
più in seguito, perché sa di potersi fidare di chi fa disagree and commit: è stato aperto
nella comunicazione, ha detto che cosa secondo lui non andava, ma non ha fatto niente per
remare contro, anzi ha dato il suo supporto. ``Non posso chiedere di più a una persona.''

\subsection{Chi arriva a lavori iniziati}

\paragraph{Se il team cambia in corso d'opera, chi subentra come si rapporta alle decisioni
già prese?} (per dimissioni, promozioni e così via). Deve \textbf{accettare lo status quo}. Al
relatore capita spesso di entrare in team con progetti già in corso: si fa dire che cosa è
stato scelto finora e va \textbf{in commit di default}: ``non ho preso parte a questa decisione,
quindi non sono in \emph{agreement}, ma ora faccio parte del team, quindi vado in commit''.

\paragraph{E se l'esperienza dice che una scelta ha uno sforzo molto diverso?} Ad esempio
che un'idea ``rossa'' ha in realtà lo sforzo di una ``viola''. Lo si fa presente, ma dipende dal
\textbf{tempismo}: se la scelta è appena stata fatta e c'è ancora modo di cambiarla, la si
segnala; se si è al 90\% del progetto, è controproducente dire ``fermi tutti, passo
indietro'', e a un certo punto bisogna rispettare la scelta del team (``avete scelto così,
andiamo avanti e vediamo come va''). E se si sa che, calcolando ferie e periodo estivo,
un'animazione prevista per giugno arriverà a settembre? È molto \emph{case by case}: a
seconda dell'apporto si può modificare il risultato o no. È come riaprire un processo se
emergono nuove prove: se dimostrano che l'imputato non è l'assassino, è giusto riaprirlo; se
dimostrano che l'assassino ha sparato con la sinistra invece che con la destra, conta fino a
un certo punto. È giusto segnalarlo, ma non sempre quello che si porta giustifica uno
``\emph{scratch that}, rifacciamo tutto'': diventa \textbf{un feedback, non una risoluzione}
(``state facendo così, questo è il mio feedback, poi vediamo come va'').

Con cinque minuti di anticipo si va in pausa pranzo, con l'accordo per il pomeriggio: circa
un'ora di presentazione più ``standard'' sul leading from within, poi domande e risposte,
``e quando finiscono le domande abbiamo finito''.

% ---------------------------------------------------------------------
\section{Leading from within: perché parlare di leadership}
\label{sec:z2-leadership}

\subsection{La ripresa del pomeriggio}

Alla ripresa, mentre si aspettano gli ultimi arrivati, il relatore chiede se durante la
pausa siano venute in mente domande. L'unica è logistica: le slide verranno date a fine
workshop? Sì: il relatore le passerà all'accademia (a Vanessa, Andrea ed
Enrico\footnote{Nella trascrizione l'accademia è chiamata ``Ive'' e ``Hive'', come nella
prima giornata; per i nomi dei referenti vedi la nota all'inizio del
capitolo~\ref{cap:z1-feedback}.}), che le farà avere alla classe. Gli appunti vanno
comunque presi, meglio su quello che dice che su quello che mostra.

L'ultima parte, dice, dovrebbe essere abbastanza veloce da presentare, ma vale la pena
fermarsi a parlarne, perché è un discorso più avanzato. È la logica chiusura di quanto
detto tra la prima giornata e la mattina: tirare le fila e mostrare come tutto faccia parte
di un unico filone di ragionamento e progettualità, con un unico scopo, un unico obiettivo e
un unico modo di ragionare.

La voce sembra a posto, ma il relatore è stanco dopo due giorni, come immagina lo sia la
classe, nello \emph{stretch} finale di due giornate molto piene. Sono cose che in condizioni
normali spiega al suo team nell'arco di settimane, con il tempo di metterle in pratica,
discuterle e approfondirle. Qui è tutto concentrato: servirà tempo per metterlo in atto e
fare le proprie prove, quindi non bisogna spaventarsi se non verrà tutto giusto la prima
volta.

\subsection{Il leader non è il più bravo}

Perché parlare di leadership in un workshop non rivolto solo ai leader? La risposta
completa arriverà più avanti, ma il punto di partenza è che il relatore \textbf{non crede
nei leader di team o di progetto che detengono da soli il sapere e il modo di fare le
cose}. Come detto il giorno prima, essere un leader, un manager o un capo non dovrebbe
voler dire essere il migliore in quel lavoro, ma essere \textbf{la persona più adatta a
ricoprire quel ruolo}: guidare un gruppo significa \textbf{occuparsi delle persone} al suo
interno.

L'obiettivo che il relatore si dà, e che secondo lui ogni manager dovrebbe avere in mente,
è che quando si torna dal lavoro ci si senta \textbf{ispirati, sicuri e realizzati}, e che
questo sia \textbf{un diritto naturale di tutti}, non un lusso riservato a chi ricopre certe
posizioni, a pochi fortunati o ``a qualcuno sì e a qualcuno no''. Ha notato,
specialmente in Italia, la tendenza a vedere tutto ciò che viene prima del diventare
manager come la \emph{gavetta} necessaria per arrivarci. Invece chiunque faccia parte di
un'organizzazione, che sia un'azienda di videogiochi o altro, dovrebbe sentirsi soddisfatto,
ispirato e sicuro in quello che fa; e non è corretto che un'azienda tratti bene le persone
``da manager in su''. Anche quando non lo dicono, a volte i fatti sono quelli: da manager si
viene trattati in un certo modo e si ha accesso a certe cose, altrimenti no. Non si tratta di
\textbf{appiattire la gerarchia}: si tratta di assicurarsi che, pur dando cose diverse ai
diversi livelli, a tutti i livelli venga riconosciuto uno strato comune di
\textbf{rispetto, fiducia e professionalità}.

\subsection{Le carriere divise: manager o \emph{subject matter expert}}

Uno studente chiede conferma: quindi un \emph{lead artist} non è necessariamente più bravo
di un senior a fare l'artista? Esatto. E qui il relatore apre una parentesi. In molte
aziende, e soprattutto in quelle in cui tende a lavorare lui, la motivazione principale per
diventare manager è avere una promozione e quindi un aumento, più che il desiderio di
prendersi cura del team e delle persone. Come si ovvia al problema? \textbf{Dividendo le
carriere}. Arrivati al livello senior, il relatore fa con ognuno dei suoi senior
\textbf{``the Talk''}, per esporre le due opzioni che hanno davanti:
\begin{description}
  \item[Busta A: diventare manager] È il percorso standard: si valuta se si è adatti a fare
        il leader e a guidare un team, e si danno gli strumenti per farlo, perché
        ``nessuno nasce leader imparato'': prima di provarlo sulla propria pelle nessuno sa
        che cosa voglia dire. Per quanto si possa essere carismatici, fare il leader è
        un'altra cosa, e non vuol dire solo essere carismatici. Ci sono formazione,
        strumenti e un percorso di crescita.
  \item[Busta B: diventare \emph{subject matter expert}] Una persona che va oltre il senior
        nella stessa disciplina: diventa \textbf{l'esperto dell'intera azienda} su quella
        disciplina, perché la sua priorità non è aiutare gli altri ma essere il migliore.
        ``Non c'è assolutamente nulla di male.'' Il suo ruolo è essere il più specializzato
        in quel campo, il punto di riferimento da cui tutti vanno a farsi spiegare le cose.
\end{description}
Il concetto è abbastanza nuovo, e il punto chiave è che l'azienda \textbf{paga il
\emph{subject matter expert} e il leader esattamente allo stesso modo}: non c'è un guadagno
economico nello scegliere di fare il manager invece dello specialista, e il relatore valuta
quel tipo di conoscenza quanto il saper gestire persone. Così chi vuole essere il migliore,
o anche solo un aumento, può continuare a fare ciò che sa fare e avere l'aumento; chi vuole
l'aumento \emph{e} prendersi cura delle persone sceglie l'altra strada. Così si fa una scrematura sul
\emph{motivo} per cui si vuole fare il manager. Sembra una scelta piccola, ma è molto
importante capire se ci si vuole occupare di se stessi o degli altri. Non è un giudizio
(``non sto dicendo che sia meglio l'uno o l'altro''): le due cose \textbf{non convivono}, e
bisogna sapere che cosa si vuole fare, sapendo che l'azienda tratterà allo stesso modo le
due scelte.

\paragraph{Si può essere entrambi?} Si può scegliere ad esempio il percorso da lead artist
e voler comunque continuare a migliorare le proprie capacità. Ma, come nelle simulazioni di
conversazione del giorno prima, il relatore chiederebbe: ``se avessi le risorse per fare solo
una delle due cose, quale sceglieresti?''. Perché non le si avrà: non ci sarà il tempo e il
modo di fare bene entrambe. Il suo ``cinismo'' arriva a dire che bisogna scegliere, perché
tenersi buone entrambe le strade porta a farle male tutte e due. Preferisce un ottimo
\emph{people manager} o un ottimo specialista; le figure intermedie, in questo campo,
rischiano di diventare pericolose per sé e per il team.

\paragraph{Specializzazioni verticali.} La strada dello specialista permette
specializzazioni molto verticali: in chat qualcuno cita il \emph{mecha designer}, ma può
esserci anche lo specialista di animazioni delle espressioni facciali. Naturalmente l'azienda
deve avere bisogno di quella competenza (non ci si specializza in cose che non servono), e non
è detto che ne serva uno solo. È un percorso di carriera parallelo e alternativo a quello del
manager ``standard''.

% ---------------------------------------------------------------------
\section{Il ``momento Piero Angela'': la biologia del gruppo}
\label{sec:z2-biologia}

\subsection{Perché una digressione scientifica}

Il relatore annuncia un ``piccolo excursus biologico'', fondamentale per capire dove vuole
arrivare, e prova a mettere un sottofondo musicale adatto: la condivisione dell'audio non
funziona, quindi manda il link in chat perché ognuno lo faccia partire per conto suo. Un
studente lo battezza subito ``momento Piero Angela''. Il relatore cerca di far combaciare le
motivazioni aziendali del suo metodo di leadership e di progettualità con \textbf{basi
scientifiche}: in passato ha collaborato con varie università che si occupavano di
neuromarketing e neurobiologia, per avere una conferma o una smentita del fatto che ciò che
ipotizzava ``non fosse una cavolata colossale''. Semplificando moltissimo, si guardano solo
alcune sostanze che si traducono direttamente nel modo in cui si vive l'essere in un gruppo,
in un team, in un'azienda\footnote{L'impianto di questa parte (endorfine, dopamina,
serotonina, ossitocina e cortisolo, il villaggio protetto dai pericoli esterni, il ``cerchio
di sicurezza'') ricalca da vicino quello del libro di Simon Sinek \emph{Leaders Eat Last}
(2014), che il relatore cita poco più avanti. Le descrizioni biologiche sono molto
semplificate, come avverte lo stesso relatore, e vanno lette come tali.}.

\subsection{Endorfine}

Le \textbf{endorfine} sono rilasciate per \textbf{mascherare il dolore}. Lo sa chi corre, chi
fa pesi o ginnastica, ma anche chi ha riso così tanto per una battuta da avere male ai
fianchi: è quella sensazione di dolore in cui però si è contenti di provarlo. Il \emph{runner's
high} è l'esempio più calzante: essere distrutti dopo una corsa ma soddisfatti. È la risposta
biologica allo sforzarsi in cose che sono per il nostro bene, un modo con cui il corpo ci
``convince'' a farle.

\subsection{Dopamina}

La \textbf{dopamina} è più familiare: è la \textbf{ricompensa per un'azione}. Dalle cose più
banali (si mangia e si riceve dopamina) a quelle viste in mattinata: anche completare una
checklist rilascia dopamina (sezione~\ref{sec:z2-pianificare}). Viene rilasciata quando si
vede l'obiettivo, come il traguardo durante una corsa; è la stessa cosa che accadeva ai
cacciatori del villaggio in tempi antichi, che rilasciavano dopamina quando vedevano le tracce
e venivano così incentivati a seguirle. Sono concetti fondamentalmente ``animali'', ma si
applicano tuttora anche in azienda.

\paragraph{Una nota di colore: la dipendenza.} La dopamina \textbf{crea dipendenza}: molti casi
di ossessione per alcol, nicotina, gioco d'azzardo e social media sono problemi legati alla
dipendenza da dopamina. Secondo il relatore non è un caso che le diagnosi di ADD e ADHD siano
aumentate del 60\% negli ultimi dieci anni\footnote{Nella trascrizione ``EDD e HD''. Il dato
e il nesso causale con la maggiore disponibilità di alcol, nicotina, gioco d'azzardo e social
media sono riportati come detti in lezione e non sono stati verificati; l'aumento delle
diagnosi viene in genere attribuito anche ad altri fattori, come la maggiore consapevolezza e
il cambiamento dei criteri diagnostici.}, perché buona parte di quegli stimoli è diventata
ampiamente disponibile proprio in quel periodo. Anche la dipendenza dal lavoro, a volte, ha la
stessa origine: si vuole il rispetto che deriva dal lavoro, il \emph{rush} che deriva
dall'averlo.

\subsection{Il villaggio di diecimila anni fa}

Si fa poi un salto di diecimila anni, in un villaggio. Gli esseri umani, per quanto più o meno
introversi o estroversi, sono \textbf{animali sociali} che tendono a vivere esperienze di
gruppo; le aziende nascono da un gruppo di persone che si danno un obiettivo condiviso, come
in antichità un gruppo di persone che voleva sopravvivere insieme. Quello che non è cambiato è
che \textbf{fuori dal gruppo ci sono dei pericoli}: tigri, lupi, draghi (``do per scontato
che nell'antichità ci fossero i draghi, perché D\&D mi insegna così''). Per evitarli, i gruppi
tendono a riunirsi e a creare difese intorno a sé: da qui le prime difese dei primi villaggi,
per isolarsi dai pericoli esterni. ``Vi giuro che avrà tutto senso.''

\subsection{Serotonina}

La \textbf{serotonina} viene rilasciata anche \emph{indirettamente}: non deve per forza
riguardare noi. Quando si assiste a una cerimonia o celebrazione pubblica che include una
persona cara, si rilascia serotonina come se si fosse noi l'oggetto del riconoscimento. È uno
dei motivi per cui si fanno diplomi e lauree ``all'americana'', con la consegna sul palco
davanti a tutti: gli studi avrebbero mostrato che anche i genitori in platea rilasciano
serotonina, come chi riceve il premio. Lo stesso vale per parenti, amici, matrimoni: sono
esperienze condivise, e si tende a condividere in pubblico riconoscimenti e successi perché in
tutti i partecipanti viene rilasciata serotonina.

Purtroppo ha a che fare anche con i beni materiali: la tendenza all'\emph{hoarding}, cioè a
collezionare eccessivamente qualcosa, nasce dal fatto che quell'attività rilascia serotonina,
ingannando la mente nel credere di aver ottenuto un successo.

È chiamata anche \textbf{la sostanza della leadership}, perché aiuta a riconoscere in qualcun
altro chi ha valore per il gruppo. Nel villaggio, quando arriva un pericolo, c'è sempre un
gruppo di persone che dice ``ce ne occupiamo noi'': il gruppo sa che si stanno prendendo un
rischio, e la serotonina è il modo per riconoscere che meritano rispetto. Il relatore consiglia
il libro di Simon Sinek \emph{Leaders Eat Last}\footnote{Nella trascrizione ``Leadership
Lust''. Il titolo è \emph{Leaders Eat Last: Why Some Teams Pull Together and Others Don't}
(2014).}: i leader mangiano per ultimi,
all'opposto di quanto accadeva in antichità, quando al ritorno dalla caccia mangiavano per
primi i cacciatori, che dovevano essere in forze per la caccia successiva, poi donne e
bambini, poi gli altri. C'è un \textbf{riconoscimento sociale}: tu fai qualcosa per proteggere
il gruppo, noi ti riconosciamo qualcos'altro in cambio.

Le aziende sono strutturate in modo molto diverso? È una domanda retorica: quando si dice che
serve un capo, un manager, un rappresentante, si sta dicendo che quella persona dovrebbe avere
l'interesse di proteggere il gruppo e di prendersi cura di tutti. Per questo sapere che il
proprio capo guadagna di più non fa particolarmente arrabbiare: lo si vede come quasi naturale,
perché più responsabilità si ha verso il gruppo più si guadagna, dato che quello dovrebbe essere
il suo scopo. Ed è anche il motivo per cui il relatore si arrabbia molto quando legge di CEO o
\emph{executive} che, a fronte di licenziamenti di migliaia di persone, si assegnano dei bonus.
È una rabbia non del tutto razionale, che sente da dentro: è come se si stesse violando un
\textbf{contratto sociale}. Il contratto è ``tu ti prendi dei rischi per me e mi difendi, io ti
riconosco più cibo, più soldi''; quando si abusa di quella posizione per un beneficio personale
a scapito di chi si dovrebbe proteggere, tutti lo vedono come un'ingiustizia, perché si viene
meno al contratto stabilito diecimila anni fa con la nostra biologia.

\subsection{Ossitocina}

L'\textbf{ossitocina}\footnote{Nella trascrizione ``ossidocina''.} è la \textbf{sensazione di
sicurezza} che si genera quando si è vicini a qualcuno. Non è un caso che tra i legami di
affinità e collaborazione più forti ci siano quelli dei team che si uniscono contro un
\textbf{nemico comune}: i team che superano insieme avversità e difficoltà ne escono molto
legati. Tra le persone che il relatore sente più spesso in questo periodo ci sono quelle con cui
ha condiviso momenti difficili, un cattivo manager, una brutta esperienza lavorativa: sa che
erano lì con lui. Vale anche per i veterani di guerra o per i sopravvissuti a una stessa
tragedia, che sviluppano un rapporto molto forte (se ne è parlato molto negli Stati Uniti a
proposito dei veterani), perché hanno rilasciato immense dosi di ossitocina. È anche ciò che fa
accettare subito il figlio alla madre quando lo prende in braccio dopo il parto, ed è il motivo
per cui, a migliaia di anni di distanza, si stringono ancora accordi con una stretta di mano o
un abbraccio: c'è sempre qualcosa di fisico, generato dal bisogno di ossitocina.

Viene generata anche quando si compiono o si \textbf{assiste ad atti di generosità} senza
aspettarsi nulla in cambio. Se il relatore dicesse di aver donato cinquanta euro a
un'associazione di beneficenza, la reazione sarebbe ``ok, bene''; se dicesse che ogni sabato
pomeriggio fa volontariato, lo si vedrebbe come qualcosa di molto più prezioso, e la cifra
potrebbe crescere parecchio prima che qualcuno la consideri più preziosa dei suoi pomeriggi:
tutti sanno che la vera donazione è \textbf{il tempo} investito, che vale più del bene
materiale. E basta assistere: vedere qualcuno che aiuta un'anziana ad attraversare la strada fa
stare meglio anche noi. È percepita a livello di gruppo. Infine, l'ossitocina rafforza il
\textbf{sistema immunitario}: sentirsi sicuri e tranquilli fa stare meglio e, in un certo senso,
ammalare di meno.

\subsection{Cortisolo}

L'opposto è il \textbf{cortisolo}: la sensazione di stress, paranoia e paura, la risposta
\emph{fight or flight} (o scappo o combatto), con battito accelerato, muscoli in tensione e così
via. Tra i suoi effetti:
\begin{itemize}
  \item \textbf{blocca il sistema immunitario}: se si sta per morire, al corpo non importa di
        prendersi un raffreddore. È il motivo per cui le persone stressate si ammalano di più;
  \item \textbf{blocca l'ossitocina}: non ci si può sentire tranquilli quando si è in ansia;
  \item soprattutto, e questo è il punto più importante per le aziende, è
        \textbf{contagioso}. In qualsiasi documentario, quando una gazzella che bruca alza la
        testa e percepisce un movimento nella savana, scappano tutte, non solo lei. Succede,
        non sorprendentemente, anche in ufficio: quando qualcuno ha paura perché potrebbero
        esserci licenziamenti, o perché il manager è molto aggressivo, la paura si diffonde nel
        gruppo ed è molto dannosa.
\end{itemize}
È anche l'effetto alla base della \textbf{sindrome dell'impostore}: sentirsi costantemente
inadeguati perché si teme che succeda qualcosa significa vivere in uno stato di stress. Ed è il
motivo per cui chi è insoddisfatto del lavoro tende a stare peggio anche fisicamente e ad
ammalarsi di più.

\begin{center}
  \begin{tabular}{@{}p{0.17\textwidth}p{0.36\textwidth}p{0.37\textwidth}@{}}
    \toprule
    \textbf{Sostanza} & \textbf{Che cosa fa (in sintesi)} & \textbf{Esempi dalla lezione} \\
    \midrule
    Endorfine & mascherano il dolore di uno sforzo che ci fa bene &
      \emph{runner's high}, ridere fino ad avere male ai fianchi \\
    \addlinespace
    Dopamina & ricompensa un'azione, si attiva vedendo l'obiettivo; crea dipendenza &
      spuntare una checklist, le tracce per il cacciatore, gioco d'azzardo, social \\
    \addlinespace
    Serotonina & riconoscimento sociale, anche indiretto; ``sostanza della leadership'' &
      lauree e premi in pubblico, chi difende il villaggio, \emph{hoarding} \\
    \addlinespace
    Ossitocina & sicurezza e legame; rafforza il sistema immunitario &
      il nemico comune, i veterani, la stretta di mano, il volontariato \\
    \addlinespace
    Cortisolo & stress, \emph{fight or flight}; blocca difese immunitarie e ossitocina; è
      contagioso & la gazzella che scappa, la paura dei licenziamenti, la sindrome
      dell'impostore \\
    \bottomrule
  \end{tabular}
\end{center}

% ---------------------------------------------------------------------
\section{Domande sulla biologia del gruppo}
\label{sec:z2-qa-biologia}

\paragraph{Il gioco d'azzardo include i gacha e le loot box?} Assolutamente sì, anzi ne sono il
\textbf{driver principale}: i gacha, i FIFA Ultimate Team, lo ``sbustare'' da \emph{Magic} a
\emph{Hearthstone}, e in generale il ``vediamo se stavolta mi va bene'' sono tutti gioco
d'azzardo. Lo studente aggiunge che molti giochi mobile sono costruiti su questi meccanismi e che
c'è chi gioca più a quelli che ai giochi ``normali''. È il motivo per cui alcuni stati hanno
dichiarato illegali certi giochi\footnote{Il caso più noto è quello del Belgio, che dal 2018
considera gioco d'azzardo le \emph{loot box} a pagamento; anche i Paesi Bassi le hanno
contestate.}.

\paragraph{Come si genera l'ossitocina nei rapporti da remoto?} ``Bellissima domanda.'' È il motivo
per cui, come detto il giorno prima, il relatore insiste perché i meeting si facciano almeno in
\textbf{webcam}, perché i riconoscimenti siano dati e i successi celebrati \textbf{pubblicamente}
(quando non mette troppo a disagio la persona), e in generale cerca di creare \textbf{situazioni
condivise}. Anche lavorando da remoto, insiste per avere il \emph{team building} dal vivo e almeno
\textbf{un incontro all'anno} di persona in ufficio, perché crea situazioni di rapporto diretto.

\paragraph{Come si contrasta il cortisolo?} In chat qualcuno suggerisce la meditazione: sì,
assolutamente. Per il lavoro, la risposta arriva subito dopo nella presentazione.

\paragraph{Il nemico comune è una strategia che un manager può usare?} Spesso \textbf{il manager è
il nemico comune}, senza bisogno di sforzarsi. Ma il relatore chiede di capire bene il punto:
alcuni manager applicano consapevolmente il concetto ``mi ergo io a nemico del team, così sono
tutti uniti contro di me e lavorano bene'', e ne fanno uno stile di management. Decidono di essere
il \emph{bad cop}, il \emph{villain}, ed è così che ottengono risultati. È una scelta lecita e
funziona; il relatore preferisce fare diversamente, ma non per questo dice che non funzioni. Molte
volte si sottovaluta quanto sia consapevole: ``non si rende conto di quanto sia antipatico a
tutti''; forse se ne rende perfettamente conto e lo fa apposta per essere il bersaglio. Il relatore
l'ha visto fare spesso come soluzione disperata, quando il team era troppo impegnato a litigare al
suo interno: ``piuttosto che vederli litigare tra loro, gli do io un motivo per odiare me, così si
uniscono contro di me''. Non lo consiglia assolutamente, ma l'ha visto funzionare.

Uno studente racconta che nella sua azienda il manager fa vedere \emph{il lavoro stesso} come il
nemico (``prima finiamo, meglio è''). È un modello più basilare ma, come dice lo studente stesso,
funziona.

\paragraph{Quanto è etico puntare sulla dipendenza del giocatore per fare soldi?} ``Non è
assolutamente etico, ed è assolutamente profittevole, altrimenti le aziende non lo farebbero.'' Il
relatore racconta una cosa che lo ha sconvolto quando l'ha scoperta, a Blizzard. Lavorava su
\emph{Hearthstone}\footnote{Nella trascrizione ``Arsdol''.} e si stava discutendo di come
monetizzare il gioco; venne preso come consulente un esperto, israeliano se ricorda bene, che
lavorava a Malta per società di poker online e di \emph{gambling}. Raccontò che molte aziende che
monetizzano sulla dipendenza dal gioco d'azzardo avevano sviluppato metodi più o meno etici (``più
spesso meno''). L'esempio culminante: analizzavano gli utenti occasionali, quelli che giocavano
magari una settimana al mese, per capire \emph{quando arrivava loro lo stipendio} (di solito si
deposita lo stipendio e poi si gioca quello che resta), per colpirli il giorno dopo con email di
offerte per tornare a giocare. Li si stimolava proprio nel momento in cui avevano disponibilità, sapendo
che erano sensibili, per aumentare le probabilità di riaverli come giocatori. ``È orribile. È
terrificante e, come tutte le cose terrificanti, è anche altamente funzionale.''

Gli ha fatto riflettere molto su dove ci si possa spingere per fare soldi. La versione etica, per
lui, è ad esempio far uscire un gioco sotto Natale, quando si è più disposti a regalarlo ai figli:
si stanno ottimizzando le proprie possibilità, e lo accetta. Si scade nel \textbf{predatorio}
quando si sfrutta la debolezza di qualcuno in un momento preciso per estorcergli soldi che non può
permettersi di perdere. Poi ognuno tragga le proprie conclusioni: ``non è un workshop sull'etica del
lavoro''. Non pensa di dover imporre la propria opinione, né si arroga il diritto di dire che chi lo
fa è una bestia senza cuore, visto che è il loro lavoro: ``chiunque riesca poi a dormire in pace con
se stesso, buon percorso''.

% ---------------------------------------------------------------------
\section{La piramide della fiducia}
\label{sec:z2-piramide}

\subsection{La leadership non è un titolo}

Chiuso il ``momento Piero Angela'', il relatore arriva al punto. La frase della prima giornata (``la
capacità di un gruppo di persone di fare cose straordinarie dipende dal modo in cui queste persone
operano insieme'') porta a capire che \textbf{la leadership non è un titolo lavorativo}. Essere
manager da contratto, o avere ``manager'' nel titolo, rende un manager, non un leader. Il relatore
non ha mai visto un team rispettare qualcuno solo perché da contratto doveva farlo: si può avere
paura o timore del proprio manager, ma non lo si deve rispettare perché è il manager. \textbf{Il
rispetto si guadagna}, e si possono influenzare le persone intorno solo attraverso il rispetto
guadagnato, quando si sono accumulati abbastanza esempi e \emph{case history} da giustificare
rispetto e fiducia. Altrimenti, dice, sarebbe il primo a chiedere ``perché vi fidate di me, se non
vi ho ancora dato motivo di farlo?'': non sarebbe razionale.

Da qui la \textbf{piramide della fiducia}, che si costruisce partendo dalla base.

\subsection{Alla base: l'integrità morale}

Per essere un buon leader o un buon manager bisogna \textbf{voler genuinamente aiutare gli altri
membri del team}. È il motivo del discorso sulle carriere divise: è a questo livello che si
screma chi vuole la promozione per i soldi. Chi vuole essere leader non deve cercare il guadagno
personale, ma quello del gruppo. Come detto il giorno prima, il relatore è pronto a litigare per il
suo team e a prendersi le colpe per il suo team, perché vuole il bene del suo team e sa che poi loro
si prenderanno cura di lui. Chi ha in mente questo sta in questa categoria. Poi:
\begin{itemize}
  \item \textbf{essere onesti}. Il relatore può affermare di non aver mai mentito
        intenzionalmente a un collega, a un cliente o a un giocatore. Gli è capitato di dare
        informazioni errate perché in quel momento sapeva quelle: ad esempio annunciare che
        un'espansione sarebbe uscita in un certo mese, e poi l'uscita è slittata\footnote{Nella
        trascrizione l'esempio non è coerente: prima si parla di un'uscita a novembre rimandata
        di una settimana, poi di ``pensavo fosse settembre''. Il senso resta chiaro.}. Ha detto
        una cosa sbagliata, ma tutti hanno capito che non stava mentendo: in quel momento pensava
        che fosse vera;
  \item soprattutto, \textbf{accettare la responsabilità dei fallimenti}: quando si sbaglia,
        saper dire ``ho sbagliato''. Il relatore lo ha messo perché ``non avete idea di quanti
        facciano fatica'' ad ammettere di poter sbagliare: c'è sempre una scusante, una
        mitigazione, un ``sì, ma''. Lui vuole sentire se è colpa tua o no;
  \item \textbf{chiedere scusa} e \textbf{ammettere di non sapere}: anche dire ``non lo so'' viene
        spesso evitato, come accettare un fallimento.
\end{itemize}

\paragraph{``Raccontami il tuo più grande fallimento''.} Per lo stesso motivo, nei colloqui il
relatore chiede spesso, più che ``parlami di un tuo grande successo'' o di un progetto andato bene,
\textbf{``mi racconti il tuo più grande fallimento?''}, e osserva come lo raccontano. Chi non si
aspetta la domanda spesso racconta di quella volta in cui ha provato in tutti i modi a fare in un
certo modo, ma il capo gli ha detto di fare diversamente e il progetto è andato male. ``Non mi hai
risposto: mi hai raccontato un progetto andato male, ma il tuo messaggio è che volevi fare
diversamente e qualcuno ti ha detto di fare così. Quindi non è un tuo fallimento.'' È come
rispondere a ``qual è la tua debolezza?'' con ``sono così attento ai dettagli che a volte\dots'':
si fa passare un punto di forza per una debolezza. Il relatore vuole un esempio in cui \emph{tu},
senza altri fattori, hai sbagliato, e come ti sei comportato. Gli dice molto su come quella persona
reagirà ai suoi feedback: se ha davanti ``un muro'' che dirà sempre ``non è colpa mia'', non è adatta
a un ambiente in cui il feedback è alla base di ogni comunicazione. Chi non ammette di aver mai
sbagliato in vita sua, lo scarta.

\paragraph{È possibile colpevolizzarsi troppo per un fallimento?} Sì, assolutamente. Rientra molto
nella sindrome dell'impostore, o nell'aver sempre vissuto rapporti di lavoro tossici: ci si abitua ad
accettare che ci vengano urlate contro le cose e che sia sempre colpa propria, e ci si autoconvince
che lo sia.

\subsection{Al centro: l'affidabilità}

Superato questo livello si arriva all'\textbf{affidabilità}. Per essere un buon leader non bisogna
essere il migliore, ma affidabili: basta essere riconosciuti come bravi in qualcosa, non
necessariamente i migliori, e quell'affidabilità ce la si porterà anche nella leadership. Significa:
\begin{itemize}
  \item \textbf{rispettare la parola data}: se si promette ``parleremo di un aumento al prossimo
        slot disponibile'', se ne parlerà. Non ci si dimentica, e non si dice poi ``in realtà non si
        poteva fare, ma ho voluto provarci fino all'ultimo''. Una promessa va mantenuta;
  \item \textbf{impegnarsi per il tempo necessario} a raggiungere gli obiettivi;
  \item \textbf{saper rifiutare un lavoro} quando necessario: è il saper dire di no della mattina,
        ancora la matrice di Eisenhower (sezione~\ref{sec:z2-eisenhower}). ``Sono sempre gli
        stessi concetti.''
\end{itemize}

\subsection{In cima: la leadership, cioè dare l'esempio}

L'ultimo gradino è la \textbf{leadership}\footnote{Nella trascrizione ``dealership''.}, cioè
\textbf{dare l'esempio}. Non bisogna essere i più bravi; bisogna saper \textbf{offrirsi volontari
non solo quando il progetto è un successo sicuro, ma anche quando c'è il rischio di un
fallimento}. Sono tutti bravi a offrirsi per il gioco nuovo che andrà sicuramente bene.

\paragraph{Il rappresentante dei licenziati.} Al relatore è capitato di offrirsi per rappresentare
il team durante i \emph{layoff}, con il rischio di essere ostracizzato dall'azienda; con quelle
persone è ancora in contatto. Quando Blizzard fece i licenziamenti nella sede francese, vicino a
Versailles\footnote{Nella trascrizione ``i layoff di Versailles''. La sede europea di Blizzard si
trovava a Vélizy-Villacoublay, nei pressi di Versailles; il relatore non precisa l'anno.}, che
coinvolsero quasi tutto l'ufficio, lui non era tra i licenziati, ma venne \textbf{eletto
rappresentante sindacale} di chi lo era. Per lui è stato molto importante: pur non essendo parte in
causa, lo ritennero degno della loro fiducia. Per quanto ne sapevano, avrebbe potuto negoziare e
accettare tutto quello che proponeva Blizzard, perché non ne era toccato: che si negoziasse un mese
di liquidazione in più o in meno non gli portava un soldo. Avrebbe persino potuto usare la
situazione per negoziare un aumento per sé (``gli facciamo accettare tre mesi invece di sette e a te
diamo un aumento''): non era fuori discussione. Se si sono fidati è solo perché c'era una piramide di
questo tipo e lui aveva dato loro abbastanza motivi per farlo.

\paragraph{Autorevolezza, non autorità.} A uno studente che lo riassume così, il relatore risponde
che è proprio questo: un carisma diverso dall'essere simpatici, l'essere un punto di riferimento,
\textbf{più un pilastro di supporto che il motore} che fa andare avanti le cose. Non il motore davanti
a tutti, ma il \textbf{meccanico} che, quando si rompe qualcosa, sai che non ti frega: ti sistema le
cose senza ``farti la cresta'' sulle pastiglie dei freni.

\paragraph{Non abbandonare nessuno.} E soprattutto \textbf{non abbandonare chi è in difficoltà}.
Sembra una banalità, ma è alla base del motto dei Marines statunitensi, \emph{no one left behind}:
sapere che, se si viene feriti, verranno a prenderti anche sotto i missili è un forte elemento
motivazionale per correre quel rischio. In azienda è il motivo per cui il relatore ribadisce sempre
che, se qualcuno è in difficoltà, la sua prima opzione non è licenziarlo ma \textbf{dargli una mano},
e lo ha fatto a volte anche all'eccesso, per far vedere agli altri che quando saranno in difficoltà
non li lascerà in mezzo a una strada. A volte ce l'ha fatta: ha esempi di persone ``a tanto così''
dal licenziamento recuperate per i capelli e diventate ottimi \emph{performer}. Altre volte la
persona non era recuperabile e alla fine è stata licenziata. Ma chi attraversa un periodo difficile
(il \emph{burnout}, la depressione, una relazione finita) sa che la prima risposta non sarà ``ci sei
servito finora, ora puoi andartene'', ma ``che cosa possiamo fare?''. E questo fa lavorare tutti
tranquilli.

% ---------------------------------------------------------------------
\section{Prendersi cura del gruppo}
\label{sec:z2-gruppo}

\subsection{La soluzione: più leader possibile}

In un caso del genere manager e leader coincidono, ed è l'ideale: un manager dovrebbe essere anche
un leader. Ma la vera soluzione è \textbf{avere più leader possibile}. Se l'obiettivo del leader è
prendersi cura degli altri, e il manager dovrebbe essere ``il primo tra i leader'', ancora meglio è
che \textbf{tutti siano leader del proprio gruppo}, perché la leadership non deriva da un titolo né
dall'essere il manager. Il relatore ha visto tanti gruppi andare avanti perché qualcuno al loro
interno era il leader, non perché lo era il manager.

Vuol dire mettere in pratica tutto il discorso della prima giornata su reciprocità, gentilezza,
aiuto e collaborazione. Si sa che, nel peggiore dei casi, c'è il manager, ma in condizioni normali ci
sono i colleghi: e allora si va da chi è in difficoltà a dire ``posso darti una mano?''. Così, quando
saremo noi in difficoltà, magari senza nemmeno rendercene conto, sarà qualcun altro a dirci ``ti vedo
in difficoltà, come posso aiutarti?''.

\paragraph{L'esempio è contagioso, nel bene e nel male.} Quando si ha una posizione anche solo di
rispetto agli occhi degli altri, tutto ciò che si fa ha effetto sugli altri. In negativo vale la
\textbf{slippery slope}\footnote{L'esempio del cartone abbandonato ricorda la cosiddetta ``teoria
delle finestre rotte'' (\emph{broken windows theory}), secondo cui i piccoli segni di degrado
tollerati favoriscono comportamenti peggiori.}: se si accetta un cartone di pizza abbandonato per
strada, domani se ne accetteranno due, dopodomani tre. Ma vale anche in positivo: chi vede qualcuno
fare qualcosa di positivo pensa ``lo voglio fare anch'io''. ``Come fai una cosa, fai tutto'', dice un
proverbio che il relatore attribuisce alla tradizione cinese\footnote{La massima (``how you do
anything is how you do everything'') è attribuita a fonti diverse, spesso alla tradizione zen;
l'attribuzione è incerta.}.

\subsection{Il villaggio di oggi}

Si torna al villaggio del ``momento Piero Angela'', con i difensori e le bestie (reali o no) intorno.
Oggi fuori dal gruppo non ci sono più tigri, ma le \textbf{deadline}, l'\textbf{overwork}, i
\textbf{layoff}, il manager che urla, le aspettative, i \textbf{KPI}. Avere manager che si prendono
cura del gruppo è un buon inizio, ma è ancora meglio avere \textbf{il gruppo che si prende cura del
gruppo}. Non bisogna più limitarsi a pensare che sia responsabilità del manager prendersi cura di
noi: è responsabilità nostra prenderci cura del gruppo di cui facciamo parte.

Che cosa serve? La piramide vista prima. Serve essere percepiti come \textbf{autentici}: il
relatore ormai riconosce abbastanza facilmente chi si vuole ergere a leader perché gli piace
\emph{sentirsi} leader, non perché vuole \emph{fare} il leader, e quella mancanza di genuinità non
aiuta. In generale bisogna \textbf{generare sicurezza} nel resto del gruppo.

\subsection{Smettere di dire ``ci penserà qualcun altro''}

Un esempio concreto è smettere di dire \textbf{``ci penserà qualcun altro''}. Nello sviluppo è
evidente: se si trova un bug che non riguarda il proprio lavoro, tutti sono d'accordo che vada
comunque segnalato, perché altrimenti potrebbe arrivare in \emph{live}, nella \emph{release}.

\paragraph{Le lasagne di PG Esports.} Circa un anno e mezzo prima, quando lavorava ancora a PG
Esports (che organizza anche eventi dal vivo), l'azienda ospitava in ufficio dei potenziali sponsor
venuti a conoscere il team per decidere se sponsorizzare la Milano Games Week. C'era un catering, con
i classici buffet. In mezzo alla confusione (una ventina di persone di PG Esports e una trentina tra
manager, sponsor ed \emph{executive}), il relatore vede per terra, in mezzo alla stanza, una mezza
porzione di lasagne abbandonata. Senza pensarci la raccoglie per portarla via: non è sua, ma è un
piatto in mezzo alla stanza. Nel tragitto fino al bancone del buffet, almeno altri sei tra executive e
sponsor che non lo conoscevano gli consegnano il loro piatto non finito, come a un cameriere: nella
loro testa l'unico motivo per cui qualcuno potesse raccogliere un piatto da terra era che fosse il
suo lavoro. Sul momento non ci fa caso; arrivato al bancone si gira e vede tutto il suo team che ride
della scena.

Torna da loro e ne nasce una discussione: ``avete visto quanto è stato educativo?''. Avevano visto
esattamente come ragionano i gruppi che non sono abituati a pensare ``adesso ci penso io'', ma ``ci
penserà qualcun altro''. In quelle persone non c'era malizia: sicuramente era un cameriere, perché
altrimenti farlo? ``Non è il mio lavoro, lo avrà fatto perché è il suo.'' E diventa un problema quando
la stessa tendenza (``non c'entra con il mio lavoro, ci penserà qualcun altro'') si applica a
qualsiasi altra cosa. Il relatore si aspetta che \textbf{tutti raccolgano quel piattino}: che chi vede
un bug, se non può correggerlo lo segnali, e se può correggerlo lo corregga anche se non è il suo
lavoro, ``perché è un piattino in mezzo alla stanza''.

\paragraph{Da reattivi a proattivi.} Significa passare da un atteggiamento \textbf{reattivo}, cioè di
fatto passivo, a uno \textbf{proattivo}: prevedere i problemi e cercarli intorno a sé per risolverli,
invece di aspettare che esplodano (``finché non esplode non è un problema mio''). Ragionare nell'ottica
del ``e se succedesse, e io me ne fossi accorto senza fare niente?''. Uno studente lo chiama
\emph{mindset da tester}: il relatore si aspetta che, se si vede qualcosa, ci si assicuri di non essere
gli unici ad averlo visto. A volte basta farlo presente: il relatore non potrebbe mai correggere un
bug, quindi se lo trova lo segnala, e va bene così. Ma non bisogna dare per scontato che ``è un bug così
evidente che qualcuno l'avrà già segnalato''. È il paradosso dell'incidente stradale: \textbf{più
l'incidente è grosso, più è probabile che nessuno chiami l'ambulanza}, perché tutti pensano che
qualcuno l'abbia già fatto\footnote{È il cosiddetto \emph{effetto spettatore} (\emph{bystander
effect}), studiato in psicologia sociale a partire dagli anni Sessanta.}.

\subsection{Domande sul ``me ne occupo io''}

\paragraph{Se correggo un bug, devo avvisare il collega?} Dopo una breve pausa per problemi audio
(``facciamo \emph{double testing}''), uno studente chiede se, correggendo il lavoro di un collega, il
dovere sia poi andare a dargli il feedback. A prescindere bisogna informare qualcuno; e la
conversazione dovrebbe essere: ``ho trovato questa cosa, vuoi che te la sistemi io?''. Correggere
senza dirlo a nessuno potrebbe essere un danno: magari era \emph{working as intended}, o qualcuno ci
stava già lavorando.

\paragraph{E se divento ``quello che risolve quel problema''?} Uno studente porta la sua esperienza:
in un'azienda in cui ha lavorato tre anni, risolvendo un problema facile che non era di sua
competenza è diventato, di fatto, l'incaricato di quel problema, distolto dalle sue mansioni. È
un'arma a doppio taglio. Per il relatore lì il problema è del \textbf{management}. In alcune aziende si
applica il criterio ``chi trova il problema lo risolve'' (``se te ne sei accorto tu, pensaci tu''), ed
è un altro modo di scaricare la responsabilità sugli altri. Non è ciò che dice lui: deve essere chi
trova il problema a valutare se è giusto e positivo che lo risolva lui. Il cartone di pizza lo
raccoglie perché ci mette dieci secondi; ma se ha trovato un bug nel sistema di pagamento che richiede
tre settimane di lavoro, e chi l'ha introdotto fa finta di niente, non deve risolverlo lui: deve far
presente che c'è un problema, che è stato ignorato e che va risolto (``pensateci voi, ma io vi dico
che c'è''). Va fatto caso per caso: se è immediato, facile e nel proprio interesse, lo si risolve.
Ma il relatore sconsiglia alle aziende la regola ``se lo trovi lo risolvi'', perché \textbf{disincentiva
a trovare i problemi}: qualcuno potrebbe non segnalarli per non avere altro lavoro e aspettare che li
trovi qualcun altro, ed è negativo quanto la situazione descritta. Per chiudere, ne discuterebbe con il
management: ``possiamo rivedere questa policy? Incentiva l'azione sbagliata''.

\paragraph{E se ci dicono di non mettere mano?} Uno studente che lavora come tester racconta che spesso
potrebbero correggere certe cose, ma viene loro detto di non toccarle. Esatto: si comunica. ``Ho trovato
questa cosa, ne siete consapevoli? Sì? Ci state già lavorando? No? Posso farlo io? Sì? Allora lo
faccio.'' Non bisogna esagerare con l'iniziativa: anche qui vale il \emph{default to trust}. Si parte
dal presupposto che qualcuno potrebbe averlo già visto e starci già lavorando, quindi si chiede e si
offre il proprio supporto.

% ---------------------------------------------------------------------
\section{Se si diventa manager}
\label{sec:z2-manager}

\subsection{Meriti, colpe e \emph{big picture}}

Una breve digressione per chi, un giorno, diventerà manager: che cosa cambierebbe? Oltre a tutto
quanto detto, ci sono il discorso su \textbf{meriti e colpe} e quello su \textbf{progetto e task}:
\textbf{delegare i meriti e prendersi le colpe}, e pensare al progetto più che ai singoli task, cioè
alla \emph{big picture}, invece che al ``voglio tornare a fare le cose che sapevo fare''.

È un altro problema dei manager promossi dal gruppo di chi sapeva fare bene un certo lavoro: quando
hanno l'occasione di tornare a farlo, lo fanno molto volentieri. ``C'è un task da fare, lo faccio io'',
magari con la scusa di dare l'esempio. No: si sta andando nella propria \textbf{comfort zone}. Il
relatore lo capisce (anche lui tornerebbe volentieri a fare ciò che faceva prima), ma bisogna accettare
che \textbf{quelle cose non le si deve più fare}.

\subsection{Il ruolo è più grande della persona}

Anche da manager, il ruolo che si ricopre per il gruppo è \textbf{più importante di sé}. Non importa
chi sia il community manager di Blizzard, di Riot o di un'altra azienda: conta ciò che rappresenta per i
giocatori in quanto community manager, non il fatto che sia lui piuttosto che un altro.

Prendersi cura del gruppo vuol dire applicare concetti molto semplici:
\begin{center}
  \begin{tabular}{@{}p{0.42\textwidth}p{0.48\textwidth}@{}}
    \toprule
    \textbf{Invece di dire\dots} & \textbf{\dots dire} \\
    \midrule
    ``Qui non si accettano errori.'' &
      ``Spiegami perché abbiamo sbagliato'': fammi capire perché è successo
      (sezione~\ref{sec:z1-default-trust}). \\
    \addlinespace
    ``Qui si fa come dico io.'' & ``Tu come faresti?'' \\
    \bottomrule
  \end{tabular}
\end{center}

\subsection{Il tetto della bravura del capo}

Il secondo punto è cruciale soprattutto quando si promuove il più bravo in un lavoro: se il criterio
noto è che si promuove il più bravo, il promosso lo sa, e la prima cosa che penserà è ``faccio lavorare
tutti come me, visto che sono il più bravo''. Il relatore chiede alla classe: che cosa si nega l'azienda
con questa dinamica? Anche ammesso che sia davvero il più bravo, se tutti fanno come lui diventeranno
tutti più bravi, no?

Il problema è che \textbf{si fissa come tetto massimo la bravura di quella persona}: se tutti fanno solo
come dice lui, si arriverà al massimo al suo livello, mentre ognuno ha il proprio metodo. Da qui quanto
detto il giorno prima sulla tendenza dei manager ad assumere solo persone peggiori di sé: se il criterio
di promozione è la bravura, si prenderanno sempre persone un po' meno brave per evitare che diventino il
proprio capo. Se invece le carriere sono divise (``se sei bravo fai questo, se vuoi fare il \emph{people
manager} fai quest'altro''), il manager è incentivato a prendere persone \textbf{migliori di sé}; e se,
come detto il giorno prima, il bonus dipende anche dalle performance del gruppo, è nell'interesse di
tutti che il gruppo sia pieno di superstar, un ulteriore incentivo ad assumere i più bravi.

\paragraph{L'opzione C.} C'è poi un'altra cosa, forse minima, ma per il relatore molto importante: si
perdono le \textbf{occasioni di innovare}. Nelle aziende precedenti era visto come il punto di riferimento
per il \emph{problem solving}: in un'emergenza, se c'era da scegliere tra A e B, ``decide Stefano''. Andava
benissimo, ma la tendenza era che a ogni emergenza venivano da lui. Così ci si negava \textbf{l'opzione C}:
quella a cui lui non avrebbe pensato e qualcun altro sì. Magari la prima volta sbaglia, la seconda anche,
e la terza ha un'idea migliore della sua, proprio perché lui non è più l'esperto di quel settore. Fissare
tutto come ``al massimo decide Stefano'' significa non accettare altre idee: un limite molto grosso
all'innovazione, all'aggiornamento e al fare le cose in modo più moderno e flessibile.

\subsection{Celebrare i successi e i fallimenti}

\paragraph{Celebrare i successi del team.} Il relatore vuole assicurarsi che tutti capiscano quando si è
ottenuta una vittoria o un successo: come detto il giorno prima, se non si sa che cosa si fa bene si
rischia di smettere di farlo. ``Questa cosa l'abbiamo fatta bene, complimenti, bravi tutti.''

\paragraph{Celebrare i fallimenti.} Sui fallimenti si era spinto ben oltre il semplice accettarli: aveva
indetto \textbf{meeting in cui si celebravano i fallimenti}, con le persone che raccontavano i propri errori
per distribuire quella conoscenza al resto del team. ``Non mi fa arrabbiare l'errore fatto, mi fa arrabbiare
lo stesso errore fatto quindici volte.'' Condividere un fallimento con il gruppo non significa ``guardate
questa persona sciocca che ha fatto questo errore'', ma ``questa persona ci sta permettendo di evitare lo
stesso errore in quindici'': sta salvando l'azienda dal ripeterlo quindici volte, perché la gente ha paura di
dire agli altri che ha sbagliato. Mostrare dove si è sbagliato diventava quindi un incentivo, un servizio
all'azienda. E si torna al punto di prima: se il proprio ruolo per il gruppo è più importante di sé, non
importa dover alzare la mano e dire ``ho sbagliato, ho fatto così'', perché gli altri penseranno ``grazie,
così io non lo faccio''. Si dovrebbe essere lieti che qualcuno abbia fatto quell'errore prima di noi, e
andare a ringraziarlo per averlo condiviso.

\paragraph{Condividere i successi.} Infine bisogna \textbf{dare il merito a chi lo ha}. Al relatore è
capitato spesso di ricevere ringraziamenti come ``grazie a Stefano e al suo team per aver fatto questa
cosa'', e la sua risposta è stata ``grazie, ma ha fatto tutto Gianni''. Non bisogna avere paura di dire
``grazie del feedback, ma è stato lui'': è giusto ringraziare lui e il team (il relatore si sente parte del
team, chiaramente). Non si ottiene niente dal dire ``grazie, riferirò i complimenti al team'': il relatore
non ha mai creduto nel ``riferire'' che il management ha detto ``ottimo lavoro''. Chiede al management di
dirlo \textbf{direttamente} al team, perché sentirlo dire dal management ha un impatto diverso che sentirlo
riferire da lui. In questo senso, osserva uno studente, ogni gruppo è il gruppo di sostegno di se stesso:
``assolutamente sì''.

\subsection{Quanto ci vuole?}

Non c'è una scadenza. Il rispetto non si ottiene pagando per fare prima, ``con il \emph{boost} a livello
60''. Servono tempo ed energia, ed è difficilmente misurabile: il relatore non può dire quanto il team si
fidi di lui in questo momento. Lo percepisce da come lavorano, ma lo si vedrà davvero solo quando ci sarà un
problema grosso; fino ad allora è tutta una sua percezione. Il team lo sente ripetere queste cose ogni giorno
e lo vede metterle in pratica ogni giorno, e lui può solo sperare che il messaggio passi; quando ci saranno
davvero ``gatte da pelare'' capirà come è messo, perché non c'è momento migliore di una difficoltà per vedere
come reagisce un team.

\textbf{Essere leader è prendersi un impegno a lungo termine verso i propri colleghi}: non ci sono modi per
velocizzarlo. Bisogna fare il proprio lavoro ogni giorno ``a testa bassa'', aiutando gli altri, facendolo
vedere ed \textbf{esplicitando perché lo si fa}. Il relatore lo fa quotidianamente: ``ti sto aiutando; nel
caso avessi un dubbio, lo faccio perché il mio lavoro è aiutarti e io ne traggo vantaggio dal fatto che tu
lavori meglio. Sono qui per te, ma lo faccio anche per me: non pensare che sia solo gentilezza. È il mio
lavoro, ma sappi che lo faccio spontaneamente''.

\subsection{L'ultima slide}

L'obiettivo dovrebbe sempre essere lo stesso: tutti dovremmo poter finire la giornata, staccare e tornare a
casa sentendoci
\begin{itemize}
  \item \textbf{ispirati}, perché abbiamo fatto qualcosa che ha un effetto tangibile;
  \item \textbf{sicuri}, perché sappiamo di avere colleghi, un leader e dei manager che si prendono cura di
        noi;
  \item \textbf{realizzati}, per quello che facciamo e perché è ciò che ci piace fare.
\end{itemize}
E questo dovrebbe essere \textbf{un diritto naturale di tutti}, non un ``sì, ma solo il manager'', ``sì, ma
solo se sei un esperto'', ``sì, ma solo se sei un game director''. Ognuno ha le proprie capacità: bisogna
solo trovare ciò che piace e che si vuole fare.

% ---------------------------------------------------------------------
\section{Domande e risposte finali}
\label{sec:z2-qa-finale}

Finita la presentazione, si passa alle domande: ``chiedete tutto quello che volete, perché è
l'ultima possibilità''. Le domande sono riportate per temi.

\subsection{Leadership e organizzazione del lavoro}

\paragraph{La leadership riguarda tutti o solo chi vuole diventare leader?} Il discorso sui
manager vale solo per i manager; ma il consiglio è \textbf{essere tutti leader del proprio gruppo,
da dentro}. Essere un leader significa sapere che gli altri possono fare affidamento su di sé, e
di queste figure in un team ce ne possono essere infinite: non è una gara a chi è il leader.
``Siate tutti una di quelle figure a cui gli altri guardano quando sono in difficoltà.''

\paragraph{Da manager non si hanno più orari e si fanno straordinari di continuo?} No: il relatore
ha notato di essere molto più efficiente come manager quando si dà \textbf{paletti molto precisi}.
Se stacca alle sei e alle sei e cinque arriva un messaggio, risponde ``se non mi dici che è
urgentissimo, non lo guardo nemmeno, lo leggo domani''; se non è urgente, lo guarda il giorno dopo.
Ma gestisce le aspettative: ``se è urgente ci sono e ci sarò sempre, ma non aspettatevi che faccia
orari eccessivi solo perché sono manager, così come non me lo aspetto da voi''. Si aspetta che
ognuno faccia ciò che deve; se è urgente tutti faranno di più, ma se non lo è, no.

\paragraph{L'orario davvero flessibile.} Qui il relatore aggiunge una cosa che si era dimenticato di
dire parlando di ferie e permessi illimitati: introduce anche l'\textbf{orario davvero
flessibile}. Non il flessibile di molte aziende, in cui si può iniziare alle 9 e finire alle 18
oppure iniziare alle 10 e finire alle 19: se oggi si vuole iniziare alle 10 e finire alle 15, e si è
fatto tutto quello che si doveva, va bene così. Ci si prendono le ore libere in più per altro, per
staccare, per riposarsi, sapendo che se la settimana dopo bisognerà lavorare dalle 9 alle 21, lo si
farà (o si dirà ``ieri ho fatto dalle 9 alle 21, domani faccio dalle 10 alle 15''). L'obiettivo dato
è un lavoro e un progetto, non un orario: si lavora per completare i propri task e rispettare le
deadline. L'orario indicato è solo indicativo, ed è quello dei meeting: \textbf{nessun meeting prima
delle 10 o dopo le 17}. Se poi qualcuno vuole lavorare dalle 10 alle 11 e poi di notte, perché quel
giorno preferisce lavorare di notte, più rilassato, faccia pure: l'importante è che trovi il tempo per
riposare e staccare. Ci si autogestisce senza timbrare cartellini; l'unico accordo è ``hai un task, va
consegnato entro il giorno X e mi garantisci che venga fatto; come ci arrivi non mi interessa''.

\paragraph{Le è mai capitato di essere aiutato da qualcuno, in un momento di difficoltà?} Sì: come
detto, questo comportamento è \textbf{contagioso in senso positivo}. Se ogni volta che qualcuno è in
difficoltà qualcun altro lo aiuta, prima o poi capita anche a chi aiuta. Il relatore non lo dice
``tanto per dire'': ha avuto momenti in cui faceva davvero fatica ad alzarsi dal letto per lavorare,
perché esausto o demotivato per altri motivi, e il team gli ha detto ``stai a letto, ci pensiamo noi
oggi''. Questo lo aiuta ad aiutare, perché sa che la cosa è bidirezionale.

\paragraph{Tra le soft skill del workshop, da quale partire?} \textbf{Il feedback}: saperlo dare e
ricevere. È la più difficile, molto più del ``voglio aiutarti'', ma è anche quella su cui si costruisce
tutto il resto: attraverso il feedback si costruisce la \textbf{fiducia}, attraverso la fiducia il
\textbf{rispetto}, e il rispetto diventa la base portante di tutto il resto.

\paragraph{E se il leader se ne va?} A una domanda su che cosa fare quando viene meno la figura di
riferimento del gruppo\footnote{La domanda non è leggibile nella trascrizione; dalla risposta si
capisce che riguardava il caso in cui il leader che incarnava la visione del gruppo se ne va.}, il
relatore risponde che la domanda da farsi è un'altra: \textbf{è cambiato l'obiettivo?} La direzione è
sempre la stessa? Se sì, chi è la persona più adatta a portarci lì? Non conta tanto chi guida, ma
l'obiettivo che ci accomunava, di solito la visione di un leader in cui tutti si allineano: ci
crediamo ancora tutti, o è cambiato qualcosa? È il momento di fare autoanalisi: ``ero qui per quella
persona o per la visione?''. Si dice, ed è vero, che \textbf{le persone non si dimettono dai posti di
lavoro, ma da capi che non rispettano più}, perché spesso l'unica interfaccia con l'azienda è il
proprio capo: le dimissioni non sono un messaggio per l'azienda, ma per il capo (``non voglio più
lavorare con te''). Allo stesso modo, se non c'è più il leader della visione, bisogna prima capire se
si condivide ancora qualcosa.

\paragraph{Lavorare così è più faticoso?} Sì e no. ``Non voglio demotivarvi, perché non è il mio
lavoro'', ma non è tutto rose e fiori: lavorare in questo modo è anche più stancante
\textbf{emotivamente}, perché non si può più pensare ``faccio il mio, alle cinque ho finito e chi si è
visto si è visto''. Si sente un senso di responsabilità verso gli altri. Succede che si finisca il
proprio lavoro del pomeriggio, ci si giri e si veda un collega che finirà alle nove di sera, e
l'istinto sarà ``posso farti un pezzo del lavoro?'': così si finisce entrambi alle sei. In un certo
senso ci si accolla più lavoro, con la consapevolezza che quando si sarà in difficoltà ci sarà qualcuno
che si girerà a dire ``ti vedo stressato, posso darti una mano?''. È bellissimo, ma non dà solo benefici
diretti: richiede anche un investimento personale che non sempre le persone vogliono fare. Pro e
contro: è il modo che sceglierebbe il relatore, ma non è per tutti. Conosce persone che vogliono poter
staccare appena finito, e va benissimo: si cerca un'azienda in cui si possa lavorare così.

\subsection{Condividere fallimenti e successi in pratica}

\paragraph{Come si celebrano i fallimenti e si condividono i successi?} Condividere i successi sembra
più semplice (``bravo Mario che ha fatto questa cosa''). Il relatore dà uno spunto ancora migliore:
\textbf{``bravo Mario, e ora Mario ci racconta come l'ha fatto e com'è andata''}. Sono sottigliezze,
ma contano: raccontare in prima persona ciò che si è fatto dà un moto d'orgoglio che non dà sentirlo
raccontare da un altro.

Per i fallimenti, che sembrano dover partire dalla persona stessa, l'obiettivo era che tutti li
condividessero spontaneamente; ma all'inizio era difficile, perché non viene naturale alzare la mano e
dire ``ho sbagliato''. Per questo il relatore ha creato artificialmente dei \textbf{``festival del
fallimento''}, dal tono chiaramente ironico: durante una pausa pranzo, mentre tutti mangiavano insieme,
a turno si saliva ``sul palco'' a raccontare la stupidaggine fatta in settimana. Molto spesso iniziava
lui, raccontando i propri errori (``ho sbagliato questa cosa, sappiate che non si fa così''). Dopo aver
forzato l'interazione per una, due, tre settimane, alla quarta lo facevano da soli: era diventato
organico dire a tutti ``ragazzi, attenzione, ho fatto questa sciocchezza, non fatela anche voi'', e c'era
perfino un \textbf{tag apposito su Slack}. È stato il team a decidere come condividere i fallimenti: come
sempre, invece di imporre il modo, il relatore ha lasciato scegliere a loro, e funziona. È impegnativo
tanto quanto lo si vuole rendere tale.

\paragraph{Che tipo di fallimenti si condividono? Anche quelli personali?} Di solito si parte da quelli
\textbf{professionali}, cioè da ciò che si riconosce come utile anche agli altri: qualcosa che ha
rallentato, fatto perdere tempo o che rischia di ripetersi (``sappiate che, per evitare che lo facciate
anche voi\dots''). Il resto è più da conversazione privata, \emph{one to one}, con la persona con cui ci
si vuole sfogare, o con gli amici: se non ha un potenziale impatto sugli altri, non c'è un incentivo a
dirlo. E un errore che si è fatto per il proprio tipo di personalità? Qui tornano i test di personalità
della prima giornata: conoscendo il proprio profilo e quello dei colleghi, si sa che qualcuno ha lo
stesso carattere, e quindi lo si aiuta.

\subsection{Cultura aziendale e formazione}

\paragraph{Come si arriva a poter fare queste cose in azienda?} Conta molto la cultura aziendale. Il
relatore sa di lavorare bene perché oggi ha il \emph{background} e i successi alle spalle per dimostrare
che il suo metodo funziona, e quindi glielo lasciano fare; ma se dieci anni prima avesse proposto le
stesse cose, ``mi avrebbero mandato a quel paese''. C'è un momento e un posto per fare le cose, e non è
sempre possibile fare ciò che si ha in testa; ma ci si può lavorare. Per lui la domanda è diventata:
``sono convinto che così sarebbe meglio; come faccio ad avere carta bianca?''. \textbf{Portando
risultati}: si è concentrato su quello, e ci è arrivato progressivamente.

\paragraph{Come funziona la formazione con workshop come questo, in azienda?} Il relatore li fa in modo
un po' più rilassato ed esteso, con molte più interazioni ed esempi, perché in azienda si possono usare
esempi che si hanno davanti ogni giorno, e parlare di problemi e discussioni diventa molto più reale. Ma
in generale è sempre \textbf{formazione pratica}, ripetuta finché il messaggio non è recepito. Il segno
che qualcosa è stato recepito è quando arriva una persona nuova e \textbf{non è lui a spiegarle le cose}:
è il team stesso a spiegarle che cos'è la feedback culture, e lui non deve parlare. È la conferma che il
concetto viene perpetuato dal team.

\paragraph{Quanti manager dell'industria condividono questi principi?} Il relatore ha fatto molti
interventi all'IGDA di Roma e ne farà uno a Milano la settimana successiva (``il 24, mi pare''); interviene
spesso a workshop ed eventi, e il riscontro è sempre ``che bello, vorremmo farlo anche noi in Italia'', e ci
stanno lavorando. La moglie lavora in un'azienda che non ha niente a che fare con i videogiochi e lo stanno
già applicando. Magari in Italia è tra i primi a parlarne, perché il paese è ancora un po' indietro, ma non è
assolutamente l'unico: è più diffuso di quanto si pensi, non è una cosa di nicchia.

\paragraph{Questi spunti si applicano anche all'IT e alla consulenza?} ``Praticamente tutti.'' Si applicano
facilmente quando c'è un gruppo di persone su cui applicarli; se si cambia continuamente cliente è più
difficile portare questa mentalità da loro, ma questo non impedisce a chi lavora con noi di applicarla
direttamente con noi.

\paragraph{Ci si può tenere in contatto?} Su LinkedIn: il relatore accetta sempre tutti e risponde alle
domande quando riesce (ne riceve tante, ma cerca di stare al passo almeno su base settimanale). E, ``a costo
di sembrare quello che si fa pubblicità'', chi lavora in un'azienda che fa fatica ad applicare questi
concetti o non sa che esistono può metterla in contatto con lui: è spesso così che ha trovato nuovi modi di
applicarli. Vale anche per chi tra sei mesi lavorerà altrove. I suoi profili sono pubblici e non c'è nessun
problema a contattarlo.

\subsection{Performance e persone in difficoltà}

\paragraph{Se un'azienda è a corto di personale, conviene comunque eliminare chi performa poco?} Il giorno
prima si era detto che conviene separarsi presto da chi non performa (sezione~\ref{sec:z1-squadra}); ma non
si finisce per caricare di lavoro tutti gli altri? Per il relatore \textbf{eliminare non vuol dire non
sostituire}: vuol dire cambiare la persona con qualcuno magari più junior, ma più adatto come \emph{fit}
culturale. Il vecchio discorso del ``fare lo stesso lavoro con meno gente'' non è sostenibile: lo si può fare
fino a un certo punto, in caso di necessità, ma non è una soluzione. Non è ottimizzare, è far andare in
\emph{burnout} chi resta.

\paragraph{Ha più senso tenere una persona poco performante e non in sintonia con il team, piuttosto che
licenziarla senza poterla sostituire?} No: se una persona è poco in sintonia, performa poco ed è in generale
tossica e negativa per il team, \textbf{prima ci si separa, meglio è}. Questo non vuol dire non provare a
recuperarla, ma bisogna andare dritti al punto: ``abbiamo un problema: ti vedo poco performante, non ti vedo
allineato con il team, vedo che il team non si fida di te. Che cosa vogliamo fare?''. Se la risposta è ``non lo
so'', la strada porta dritta all'uscita. Bisogna \textbf{recidere i rami morti} il prima possibile, facendo
però attenzione a essere sicuri che si tratti di un ramo morto e non di un ramo recuperabile: tagliare un ramo
recuperabile mette negli altri la paura di essere i prossimi, ``alla prima foglia gialla''. Quello che si
vuole è che, quando un ramo è morto, la pianta smetta di alimentarlo, e ne crescerà un altro. Entrambi i
messaggi devono essere chiari: \textbf{non si subisce passivamente la situazione}, e allo stesso tempo
\textbf{non si abbandona nessuno} (se qualcuno è in difficoltà, la prima soluzione non è tagliare il ramo). È
un compromesso tra le due cose. (A questo punto arriva, puntuale, il crash della connessione: ``non era ancora
successo oggi''.)

\paragraph{Che cosa fare se qualcuno nel team sembra soffrire di depressione?} È un discorso complesso, perché
il relatore \textbf{non è uno psicologo}. Può dare la sua disponibilità a concedere permessi, coprire il lavoro
della persona e in generale darle modo di seguire un percorso con uno specialista: non può risolvere la
depressione, ma può alleviarne la parte lavorativa e permettere di ricevere supporto da chi è qualificato. È lo
stesso motivo per cui, a Blizzard e a Riot, quando arrivavano ticket con minacce di suicidio (``mi avete
bannato, quindi mi ammazzo'', ticket purtroppo abbastanza comuni), il \emph{customer service} era istruito a
non intervenire direttamente: non dovevano essere loro a fermare la persona, perché non qualificati, ma
indirizzarla a qualcuno di qualificato, prendendo tempo. Una persona non qualificata che prova a farlo rischia
di peggiorare la situazione; come davanti a un incidente si chiama l'ambulanza e non si fa una tracheotomia.
Per il relatore salute mentale e salute fisica sono la stessa cosa: come non si azzarderebbe a operare qualcuno
di appendicite, non si azzarda a fargli da terapeuta.

\paragraph{Come si danno feedback adatti a ogni tipo di personalità?} Bisognerebbe leggere i risultati del
profilo della persona sullo stesso sito su cui si è fatto il test di personalità della prima giornata, che
spiega anche come interagire con quel profilo. In generale: con chi è più \textbf{analitico}, portare dati a
supporto; con chi è più \textbf{emotivo}, metterlo in una condizione di comfort; con chi è molto
\textbf{pratico}, andare dritti al punto senza girarci intorno. Poi fare pratica. Oppure il modo ``perfetto''
per il relatore: provarci come viene e poi \textbf{chiedere} ``andava bene così o preferisci in un altro modo?
Potevo fare meglio?''. Bisogna creare anche quel rapporto: fidarsi di discutere anche di queste cose.

\paragraph{Se nel team fossimo quasi tutti leader, non ci sarebbero ``troppi galli''?} No, perché per leader si
intende chi si prende cura degli altri, e magari lo fossero tutti. Non significa che tutti scelgano che cosa
fare: spesso lo sceglierà comunque il manager, o guiderà la scelta. Ma avere tutte persone con iniziativa e
capaci di scegliere quando serve è qualcosa per cui il relatore ``metterebbe la firma''. Non ha mai corso il
rischio dei ``troppi galli'': se ci sono quei caratteri, il team è stato assemblato male, perché non tutti sono
compatibili con questo modo di lavorare e magari non si hanno le persone giuste per applicarlo.

\paragraph{Se la fiducia richiede tempo ed esperienza, quanto è difficile per un junior?} Meno di quanto si
pensi: proprio perché su un junior non ci sono aspettative, un junior che chiede ``posso darti una mano?'' dà un
segno chiaro che tiene più al team che a se stesso. Si può essere leader anche non dal punto di vista
professionale, ma come chi va a dare supporto agli altri: ``il medico del team, l'\emph{healer}'', e lo si può
fare anche da junior. È così che ha iniziato il relatore: andare ad aiutare a risolvere problemi non suoi gli ha
fatto accelerare molto il percorso per ottenere rispetto e stima, e lo ha esposto a cose a cui altrimenti non
sarebbe mai arrivato. È una questione di carattere: se ci si tiene agli altri, lo si fa vedere, e il resto
arriva come risultato. ``È quasi un reddito passivo'': un investimento nel far vedere che ci si è, e il resto
arriva (``Stefano si è preso cura di me, quindi ora mi prendo cura di lui'').

\paragraph{Aiutare gli altri solo dopo aver finito il proprio lavoro?} Dipende dall'urgenza. Il relatore spesso,
se qualcosa è davvero urgente, mette in pausa ciò che sta facendo, risolve e poi torna al suo lavoro; ma è troppo
\emph{case by case}. Se due colleghi litigano su chi debba avere le ferie di Natale, lo si può risolvere anche
domani; se litigano perché si è perso il \emph{repository} e domani c'è la consegna, lo si risolve subito.

\paragraph{Come mantenere motivato un team su progetti non retribuiti, come quelli accademici?} Essendo il più
\textbf{trasparenti} possibile nei processi decisionali: se si deve decidere se fare un action o un RPG, tutti
devono sapere perché si fa l'uno invece dell'altro (da qui il metodo di brainstorming della mattina). Deve essere
sempre chiaro a tutti perché si stanno facendo le cose in quel modo, perché \textbf{la demotivazione nasce dal
disallineamento} rispetto all'obiettivo del gruppo, e l'unico modo di essere allineati è essere sempre
consapevoli e partecipi del processo che decide dove sta andando il gruppo. E soprattutto: \textbf{non chiudere
mai un meeting senza aver concordato che cosa si vuole fare} e senza essere tutti d'accordo. Se non lo si è,
bisogna continuare a parlarne, per quanto scomodo, lungo e noioso: si deve sempre uscire da un meeting
sapendo che si è tutti d'accordo, o al massimo in \emph{disagree and commit} (sezione~\ref{sec:z2-disagree}).
``Quello è il massimo accettabile.''

\subsection{Carriera, colloqui e candidature}

\paragraph{Come ha iniziato a lavorare all'estero?} Negli annunci online si trovano spesso requisiti come un titolo
di studio conseguito nel paese o l'abilitazione a lavorarci. Questo vale per i paesi fuori dall'Unione europea: per
lavorare negli Stati Uniti o in Sud America serve un visto o una \emph{green card}. In molti casi, però, è proprio
ottenere il lavoro a permettere di lavorare lì: quando il relatore doveva andare a lavorare in California (a Irvine,
a Santa Monica, dove capitava), era l'azienda a occuparsi dei permessi necessari, perché ``per venire qui devi avere
un lavoro, e il lavoro te lo diamo noi''.

Il suo inizio è stato abbastanza semplice. In Italia lavorava per Digital Bros, occupandosi dei giochi online gestiti
dall'azienda e di quello che era il canale Game Network sul digitale terrestre\footnote{Questo passaggio chiarisce il
dubbio segnalato nel capitolo~\ref{cap:z1-feedback} sul ``Digital Game Network Channel'': si tratta di Game Network,
canale televisivo dedicato ai videogiochi legato al gruppo Digital Bros.}. Vide un annuncio a Dublino per
\emph{Warhammer Online}, mandò la candidatura, fece i colloqui e gli dissero ``ti paghiamo la \emph{relocation}, vieni
qui''. Come detto il giorno prima, sono le aziende a occuparsi del trasferimento: se devono creare a Dublino un'unità
che segue dodici lingue europee, sanno di dover far arrivare persone da fuori, ed è già a budget. Non è nemmeno una
richiesta da fare: è prevista. Per il relatore, da questo punto di vista, è indifferente lavorare per un'azienda
italiana o estera: se si è adatti, si è adatti. Per questo insisteva sull'inglese: se lo si parla, è indifferente di
che cosa ci si occupi e quale sia la propria lingua madre, a meno che la posizione non sia esplicitamente ``community
manager italiano''. Un programmatore polacco o italiano non cambia niente. Lavorare all'estero non è più difficile: c'è
più concorrenza solo se l'azienda cerca attivamente persone da tutto il mondo e non si pone problemi di \emph{location}.
Cambiano l'azienda, il fatto di lavorare da remoto o dal vivo, l'eventuale trasferimento; per il resto non cambia
molto.

\paragraph{Gli anni passati nell'IT contano se si vuole passare ai videogiochi?} Dipende da che cosa si vuole fare: se
da responsabile IT si vuole diventare programmatore, bisogna saper dimostrare di saper programmare. Vale anche per i
passaggi interni: se un game designer vuole diventare art director, al relatore va bene, deve solo dimostrare di saper
fare l'art director. Anzi, incentiva le \textbf{promozioni interne}, perché si conosce già il candidato; ma gli
spostamenti non sono liberi ``a sentimento'': bisogna saper fare quella cosa. I recruiter valuteranno non tanto gli anni
quanto ciò che si sa fare; gli anni servono comunque, perché dimostrano di saper lavorare in un ambiente professionale,
che è un altro insieme di competenze. Ma se si sa fare ciò che chiedono, aver lavorato nell'azienda X o Y non dovrebbe
cambiare molto.

\paragraph{Quanto vengono considerate le candidature spontanee?} Non tantissimo. Soprattutto per le aziende non grandi,
se non stanno cercando qualcuno è perché non se lo possono permettere: una candidatura spontanea non le convincerà ad
assumere. Più l'azienda è piccola, più conviene aspettare l'annuncio, perché se non c'è vuol dire che non c'è budget, e
potrebbero rispondere ``ti vorremmo, ma non possiamo prenderti''. C'è anche il rischio che finisca nel vuoto (l'HR magari
non controlla nemmeno la mail, perché non sta assumendo) o, peggio, che rispondano ``non ci servi'' e una settimana dopo
esca la posizione, con il candidato già nel cestino dei rifiutati. Si rischia di \textbf{bruciarsi} la possibilità.

\paragraph{Come ci si accorge di essere riconosciuti come leader, e come lo si ``vende'' a un colloquio?} Uno studente
osserva che la leadership, così descritta, non la si riconosce da sé ma la riconoscono gli altri. Internamente il
relatore non ha mai dovuto ``venderla'', perché parlavano i risultati. Finché non lo si vive è difficile capirlo, ma
quando si arriva alla posizione in cui il team rispetta e vede qualcuno come punto di riferimento, non ci sono grandi
dubbi: è evidente che vengono da lui. Si può riflettere su \textbf{quanto spesso gli altri vengono da noi}, cercano di
emulare il nostro modo di fare, vanno a rivedere le nostre cose per capire dove hanno sbagliato, chiedono un parere, ci
parlano di un problema: la somma di tutto questo significa molto spesso che si fidano. Per ``venderla'' all'esterno si
può dire di essersi occupati di un team in cui, attraverso il lavoro e il rispetto degli altri, si è arrivati a essere
visti come punto di riferimento; ma il suggerimento è \textbf{portare esempi tangibili} (``c'era questo progetto, non
sapevamo come risolverlo, il team mi ha chiesto una mano e l'ho risolto così''), invece di ``incensarsi'' dicendo ``sono
il punto di riferimento del team''. Se il relatore racconta di essere stato eletto rappresentante dei licenziati senza
essere tra loro, un'azienda capisce che si fidavano di lui, punto: non ci sono altre conclusioni da trarre.

\paragraph{Quali sono gli errori più gravi dei junior ai colloqui, a livello di soft skill?} In ordine di priorità:
\begin{enumerate}
  \item \textbf{non arrivare preparati}, cioè \textbf{senza esempi}. Il relatore ha sempre un esempio pronto per ogni
        cosa di cui vuole parlare: per la fiducia, per gli \emph{achievement}, per i risultati. Servono sempre esempi
        tangibili per far capire che cosa si intende;
  \item \textbf{non avere domande alla fine}, perché dimostra di non avere chiaro che cosa si andrebbe a fare. Anche solo
        chiedere che cosa si farebbe, che cosa si può migliorare, o ``che impressioni hai avuto finora dal mio
        colloquio?'': si possono fare domande molto aperte al recruiter per capire come sta andando;
  \item \textbf{non leggere bene la \emph{job description}}: presentarsi con vaghezza, senza aver capito che cosa cerca
        l'azienda. Lo si impara anche con la pratica; ma il designer che si candida per una posizione da programmatore,
        senza nulla nel CV che faccia pensare che sappia anche programmare, il relatore non solo lo scarta: lo mette in
        una \textbf{blacklist} di persone che forse non considererà mai. Se non si riesce a interpretare un testo come una
        \emph{job description}, è un problema troppo grosso da superare in una collaborazione;
  \item \textbf{errori di battitura e CV impaginati male}, che molti sottovalutano. Il relatore ha letto letteralmente,
        in una lettera di presentazione per una posizione da \emph{copywriter}, la frase ``sono una persona molto attenta
        ai dettaggi''. Come si può considerare quella persona per quel ruolo? Non ha nemmeno incollato la lettera in Word
        per il controllo automatico: non è un errore sfuggito al correttore perché la parola esiste con un altro senso,
        ``dettaggi'' viene segnato in rosso in qualsiasi lingua. Come affidarle la comunicazione aziendale? Vale per tutti
        gli errori: ciò che si manda è \textbf{il primo biglietto da visita} con l'azienda, e non si manda un biglietto da
        visita sgrammaticato. Un conto è non andare ai colloqui in giacca e cravatta (``non l'ho mai fatto e probabilmente
        non lo farò mai''), un altro è presentarsi con errori;
  \item un altro esempio comune: \textbf{l'email sbagliata nel CV}. Se la mail torna indietro, il relatore non telefona
        per farsi dare quella corretta: se non si è in grado di scrivere giusta la propria email, non vale la pena. Più
        l'errore è banale e facilmente evitabile, più lo fa arrabbiare (``scusate se mi sono inalberato''): evidentemente
        il candidato non teneva abbastanza a quel posto.
\end{enumerate}

\paragraph{Quanto è importante la lettera di presentazione?} Sempre di più: sempre più aziende la richiedono come
obbligatoria, quindi conviene prepararne una, sapendo di doverla adattare al destinatario. Deve essere \textbf{meno lunga
di una pagina}, in \textbf{tre paragrafi}:
\begin{enumerate}
  \item \textbf{chi siete}: un minimo di \emph{backstory}, che cosa avete fatto, che cosa volete fare, da dove arrivate;
  \item \textbf{perché avete mandato il CV}: che cosa vi ha interessato di questa azienda;
  \item \textbf{perché dovrebbero chiamare voi}: che cosa dovrebbe spingere l'azienda a scegliere proprio voi.
\end{enumerate}
Uno studente dice di riuscire bene con i primi due, ma di trovare difficile il terzo, che è poi la domanda che tutte le
aziende fanno. In realtà basta \textbf{collegarsi alla \emph{job description}}: se cercano un 3D artist in grado di fare
X, Y e Z, si conferma di averlo letto e si spiega che con la propria esperienza si è imparato a fare proprio quelle cose.
Non vogliono sentire niente di particolare, ``non devi cambiargli il mondo''.

\paragraph{Che cosa guarda un recruiter in un portfolio?} Il relatore li guarda poco: fa un \emph{sanity check} generale,
per vedere che non sia evidente perfino a lui, ``a digiuno di tante cose tecniche'', che c'è un problema (se se ne accorge
perfino lui, il problema c'è). Per il resto guarda che il portfolio sia \textbf{coerente con il CV}: se arriva un CV da
senior con dodici progetti completati e anni di lavoro, e il portfolio sono quattro immagini, non è coerente. Vuole una
concordanza tra ciò che si racconta e ciò che si vede.

\paragraph{Che strumenti per creare un CV?} \textbf{Canva}, che ha un catalogo di template gratuiti e a pagamento
amplissimo e personalizzabile: anche se non piace, lo si prende e lo si adatta, ma fa già il 90\% del lavoro. Insieme a
\textbf{ChatGPT}, perché ``è inutile far finta che non esista'': ci si fa scrivere la lettera di presentazione e poi la si
corregge e la si adatta, e lo stesso per il CV. L'importante è renderli adatti e coerenti, non lasciarli solo frutto di
ChatGPT; ma per fare il grosso del lavoro è un ottimo strumento.

\paragraph{Conviene lasciare nel CV esperienze non del settore?} Uno studente ha alle spalle il liceo, vari stage e
lavori che non c'entrano niente con la modellazione 3D. Si possono mettere finché c'è spazio: il relatore consiglia
\textbf{al massimo due pagine}, e quando ha cominciato a sforare nella terza ha iniziato a togliere ciò che non era più
rilevante. Ormai nel suo CV la formazione è un paragrafo: ha così tante esperienze lavorative che non interessa a nessuno
che abbia fatto il liceo scientifico o classico, e non lo mette nemmeno più. Si filtra in base a ciò che è più rilevante
per lo spazio disponibile.

\paragraph{Che cosa fare con le aziende che dicono ``ti faremo sapere'' e poi spariscono?} Non aspettare uno o due anni
per scrivere. Il suggerimento: tra le domande finali del colloquio, l'ultima del relatore è sempre \textbf{``tra quanto
mi fate sapere?''}. Se rispondono dieci giorni, all'undicesimo scrive: ``mi avete detto dieci''. Non si sente in difetto a
mettere pressione, perché sono stati loro a dare la scadenza. E se viene ignorato del tutto, quell'azienda per lui è
``morta'': non si ricandiderà mai, perché non vede perché il rispetto debba andare in una sola direzione.

\paragraph{Quanto essere trasparenti, nel CV e nella lettera, su problemi personali come l'ADHD?} Uno studente con
ADHD\footnote{Nella trascrizione ``essendo 10D''; nel resto della risposta ``di'' e ``dei 10D'' stanno per ADHD.} dice di
tendere ad avere performance altalenanti. Il relatore lo invita a \textbf{non usare l'ADHD come scusa} per essere
altalenante: anche lui ha l'ADHD, ma non è altalenante. Esserlo è senz'altro un \emph{malus}, ma esistono i mezzi per non
avere problemi di questo tipo: lui ha trovato i suoi metodi (la to-do list, la \emph{must do list}, il delegare il lavoro
al team, perché occuparsi delle cose per gli altri lo mantiene molto motivato, sveglio, all'erta, attento a come vanno le
cose). Sta anche a ciascuno trovare un lavoro in cui compensare ciò che viene richiesto: se domani lo rimettessero a fare
il programmatore, come all'inizio, probabilmente impazzirebbe, perché non riusciva a restare concentrato sul codice oltre
un certo tempo. Lì ha capito che non poteva fare il programmatore, ma ha trovato qualcos'altro che gli permette di essere
com'è e fare ciò che fa. Bisogna capire i propri punti di forza e di debolezza e giocare su quelli: si possono avere tutti
i difetti del mondo e voler comunque fare una cosa, ma bisogna accettare che sarà più difficile.

Quanto alla trasparenza, il consiglio è \textbf{non mentire mai nel CV}, ``perché tanto vi beccano''. E per il discorso della
\emph{bandwidth}: tutte le energie spese a mentire, e a tenere viva la bugia, sono energie tolte allo spiegare bene ciò che
si sa fare. È un investimento nel posto sbagliato. Bisogna concentrarsi su ciò che si sa fare, con esempi tangibili, e dire
``non lo so fare'' quando non lo si sa fare. Il peggio non è essere scoperti, ma \textbf{non esserlo}: se si finge di sapere
il C++, si viene assunti e messi a lavorare in C++, che cosa si fa il giorno dopo? Si rischia davvero di rovinarsi la
carriera. Se non si è un \emph{fit}, si sarà il \emph{fit} per la prossima. L'incubo del relatore è essere assunto in un
posto dove pensano che sappia fare qualcosa che non sa fare; anzi, tende a esplicitarlo: ``cercate questa cosa? Io non la so
fare, va bene lo stesso?'' --- ``Sì, te la insegniamo noi'' --- ``Voglio solo che sappiate che non la so fare''. Meglio
questo che vivere con l'ansia che prima o poi si accorgano che non sa programmare.

Gli capita spesso con chi mente sull'inglese. Al primo colloquio chiede ``sai parlare inglese? A che livello?'' --- ``C1,
\emph{fluent}, quasi madrelingua'' --- ``Perfetto'', e lascia cadere l'argomento. All'ultimo colloquio: ``ti va se lo facciamo
in inglese?'' --- ``No, ma io intendevo che\dots'' --- ``No, mi avevi detto \emph{fluent}: ora facciamo il colloquio in
inglese''. Prima o poi le bugie raggiungono chi le dice.

% ---------------------------------------------------------------------
\section{Chiusura del workshop}
\label{sec:z2-chiusura}

Il relatore chiede alla classe di far sapere com'è andata a Vanessa, il punto di contatto con l'accademia per il workshop:
sarà lei a riferirgli eventuali lamentele o suggerimenti, che gli permetteranno di fare meglio i prossimi. Ringrazia per aver
``retto per dodici ore'' e spera che sia stato utile, consapevole che tutti sono stanchi quanto lui.

Gli studenti rispondono che è stato chiarificatore su molte cose che non sapevano, soprattutto su come interfacciarsi nei
team, e che sarà materiale da rivedere per fissare bene i concetti, vista la mole in due giorni; qualcuno si aspettava che
sarebbe stato molto più pesante e l'ha trovato invece scorrevole. Il relatore ribadisce che \textbf{niente sostituirà
l'esperienza diretta}: il workshop dà un'infarinatura e un punto di vista, ma i concetti si metteranno in pratica solo nel
\emph{day to day}. Bisogna avere \textbf{pazienza con se stessi}, sapendo che non verrà tutto bene dal primo giorno.
Confessa di essere stato un po' preoccupato per la durata, ma il materiale c'era, e farlo in un giorno avrebbe rischiato di
trasformarlo in una corsa: ``rischiava di essere uno \emph{sprint}, e invece è una maratona''.

Infine ricorda che si può aggiungerlo su LinkedIn, che registrazione e slide (in forma integrale) verranno fatte avere
tramite Vanessa (``dite che vi ho autorizzato io a chiederle'') e che, per qualsiasi dubbio, resta a disposizione. ``Buon
resto del weekend.''

% ---------------------------------------------------------------------
\section*{Punti chiave del capitolo}
\addcontentsline{toc}{section}{Punti chiave del capitolo}

\begin{itemize}
  \item I colloqui sono una \textbf{selezione reciproca}. L'annuncio dice molto sulla cultura aziendale; alla fine del
        colloquio vanno sempre fatte domande (il \emph{day to day}, la ``bacchetta magica'', ``tra quanto mi fate
        sapere?''). Si arriva con una \textbf{cifra minima mentale} e con \textbf{esempi tangibili}; non si mente mai.
  \item Time management non è una lista di cose da fare: si parte dalla \textbf{must do list}, che è anche il primo
        momento per dire ``oggi potrei non farcela''; si stimano i tempi e si chiede quando qualcosa non è chiaro.
  \item Ogni nuovo task va \textbf{prima aggiunto alla lista}, non eseguito: priorità non è recenza. Le
        \textbf{checklist} fanno risparmiare tempo, anche (e soprattutto) ai più esperti.
  \item \textbf{First contact resolution}: raccogliere tutto ciò che serve in una sola richiesta, e non iniziare un task
        senza le informazioni necessarie. Ogni interruzione costa dieci-quindici minuti di concentrazione.
  \item \textbf{Matrice di Eisenhower}: fai, pianifica, delega, elimina. ``Non importante'' significa che non è importante
        che lo faccia tu; il quadrante dimenticato è \emph{delete}.
  \item Le pause vanno \textbf{scelte, non subite}; i task grossi vanno spezzettati; bisogna \textbf{saper dire di no},
        anche ai meeting (sono necessario? può andarci uno solo? ci sarà un resoconto scritto? per che cos'altro potrei usare
        questo tempo?). \textbf{Il multitasking non esiste}.
  \item A fine giornata si verifica la \emph{must do list}: un \emph{must do} rimandabile non era un \emph{must do}.
  \item Ferie illimitate con un minimo, giornate libere e feedback culture sono, per il relatore, le condizioni per un team
        sano; il benessere è anche mentale.
  \item Il \textbf{brainstorming genera idee, non decide}. Va preparato (topic, obiettivo, vincoli), separa la
        \textbf{generazione} dalla \textbf{discussione} (\emph{rapid ideation} di cinque minuti, \emph{dot voting}
        anonimo) e combatte l'\emph{anchoring} e il dominio degli estroversi.
  \item La decisione si prende con la \textbf{matrice impatto--sforzo}: prima le idee a basso sforzo e alto impatto, poi in
        diagonale. Lo sforzo sono le risorse, l'impatto comprende anche la motivazione; decide chi farà il lavoro.
  \item \textbf{Disagree and commit}: si dichiara il disaccordo, poi ci si impegna al massimo. Se va male ha sbagliato il
        team, non ``ve l'avevo detto''.
  \item \textbf{La leadership non è un titolo}: il rispetto si guadagna. Le carriere vanno divise tra \textbf{manager} e
        \textbf{subject matter expert}, pagati allo stesso modo, così che diventi manager solo chi vuole prendersi cura degli
        altri.
  \item Dopamina, serotonina e ossitocina legano il gruppo; il \textbf{cortisolo} è contagioso e lo disgrega. Il
        ``contratto sociale'' del villaggio vale ancora: chi protegge il gruppo riceve di più, chi ne abusa tradisce.
  \item La \textbf{piramide della fiducia}: integrità (aiutare davvero, onestà, ammettere gli errori), affidabilità
        (mantenere la parola, saper rifiutare), leadership (dare l'esempio, offrirsi anche quando si rischia, non lasciare
        indietro nessuno).
  \item Il gruppo deve prendersi cura del gruppo: \textbf{raccogliere il piattino}, smettere di dire ``ci penserà qualcun
        altro'', segnalare e chiedere prima di intervenire.
  \item Da manager: delegare i meriti e prendersi le colpe, lasciare la comfort zone, chiedere ``tu come faresti?'',
        celebrare successi e \textbf{fallimenti}, dare il merito a chi lo ha. Essere leader è un \textbf{impegno a lungo
        termine}.
  \item L'obiettivo: tornare a casa \textbf{ispirati, sicuri e realizzati}, come diritto di tutti e non privilegio di pochi.
\end{itemize}

%    \part{Il corso di Game Design}
%    \renewcommand{\thechapter}{\arabic{chapter}}
%    \setcounter{chapter}{0}
%    % =====================================================================
%  Capitolo 1 — Introduzione al corso e al mestiere del game designer
%  Fonte: trascrizione della lezione GD1 (Game Design, primo anno)
%  Pacchetti richiesti nel preambolo del libro:
%    \usepackage[italian]{babel}, \usepackage[T1]{fontenc},
%    \usepackage[utf8]{inputenc} (solo pdfLaTeX), \usepackage{booktabs}
% =====================================================================

\chapter[Il mestiere del game designer]{Il mestiere del game designer: introduzione al corso}
\label{cap:introduzione-game-design}

Questo primo capitolo fa da soglia. Prima di affrontare metodi, framework e
documenti, è necessario capire \emph{in che modo} si impara il game design,
\emph{con chi} lo si fa, \emph{con quali strumenti} e, soprattutto,
\emph{che cosa fa concretamente} un game designer nel suo lavoro quotidiano.
Il capitolo segue l'impostazione della lezione inaugurale del corso: si parte
dal contesto formativo (il corso come simulazione di lavoro), si passa al
metodo di apprendimento (l'\emph{inquiry-based learning}), si elencano
letture e risorse di riferimento, e si arriva infine a una prima definizione
operativa del ruolo del designer, dei limiti entro cui lavora e delle sue
possibili specializzazioni.

Una premessa terminologica: il game design è una disciplina giovane e quasi
tutta la sua letteratura è scritta in inglese. Per questo motivo molti termini
tecnici verranno lasciati in lingua originale (\emph{pipeline},
\emph{pitching}, \emph{tuning}, \emph{feature team}, \dots): sono i termini che
si useranno nel lavoro, ed è importante imparare a riconoscerli e a usarli
correttamente fin da subito.

% ---------------------------------------------------------------------
\section{Il contesto: un corso come simulazione di lavoro}
\label{sec:simulazione-lavoro}

\subsection{Il principio fondamentale}

Il primo concetto da fissare è che il corso \textbf{si presenta come una
simulazione del lavoro}. Dal momento in cui si entra in aula si viene trattati
come \emph{potenziali professionisti dell'industria}: non soltanto studenti,
ma lavoratori effettivi di uno studio (l'accademia), con una gerarchia, dei
referenti, dei compiti e delle scadenze.

Ne derivano alcune conseguenze pratiche:
\begin{itemize}
  \item il docente non assegna un ``compitino'' da risolvere per accumulare
        punti o presenze: ogni \emph{task} è pensato per far dedicare tempo
        all'apprendimento di qualcosa di concreto;
  \item non consegnare un task non danneggia il docente (che ``vive ancora
        domani e viene pagato ugualmente''), ma lo studente, che perde
        l'occasione di mettere le mani su un argomento e diventare più abile;
  \item docenti e studenti devono trattarsi \textbf{da colleghi}, perché è
        molto probabile che un domani lo diventino davvero.
\end{itemize}

\subsection{Perché l'impegno conta più che mai}

Il lavoro del game designer è descritto come \emph{molto romantico, ma anche
molto duro}, e sempre più duro col passare degli anni. Durante il periodo del
Covid l'industria ha visto grandi investimenti; quei ``rubinetti'' sono stati
chiusi da tempo, il rapporto con i publisher è diventato più complesso e le
opportunità lavorative non sono più quelle di tre o quattro anni fa. Questo
significa:
\begin{itemize}
  \item più competizione per ogni posizione;
  \item aspettative più alte: ci si aspetta che i candidati siano ancora più
        preparati e sicuri delle proprie competenze;
  \item un unico modo per diventarlo: \textbf{lavorare}. Più tempo si dedica
        oggi ad apprendere il mestiere, più facilmente si sarà assumibili
        domani e più agevole sarà il lavoro di tutti i giorni.
\end{itemize}

\subsection{Sbagliare qui per non sbagliare fuori}

Il mantra del corso è: \textbf{quanto più provate qui, anche in modo goffo e
sbagliando, tanto meno sbaglierete là fuori}. In accademia, per quanto la
simulazione sia realistica, nessuno toglie lo stipendio e nessuno licenzia da
un momento all'altro. Nell'industria sì, e un errore può costare molto denaro
all'azienda, ancora di più se l'azienda è la propria.

\subsection{Il valore del fallimento}

Il fallimento, se affrontato con la mentalità corretta, è probabilmente
\emph{la cosa più vantaggiosa} che si possa fare durante la formazione:
\begin{itemize}
  \item un successo (``avete superato questa deadline'') conferma le proprie
        abilità, ma insegna relativamente poco;
  \item dopo $n$ fallimenti si impara moltissimo, e i successi arrivano;
  \item è fondamentale non demoralizzarsi: il feedback negativo non è mai un
        giudizio sulla persona, ma su ciò che è stato consegnato.
\end{itemize}

Il docente riconosce apertamente che dedicare tempo ed energia a un progetto e
vederlo rifiutato, o vederlo non uscire mai, è demotivante: è successo a lui e
lo ha visto accadere anche su progetti molto grandi, con team numerosi. Il
punto però non è ciò che si è ottenuto materialmente, ma \textbf{quanto si è
cresciuti}. Da ogni esperienza ci si portano a casa due cose:
\begin{enumerate}
  \item \textbf{l'esperienza}: quanto si è diventati bravi, veloci,
        praticamente abili nell'applicare un framework, quanto si è cresciuti
        come designer;
  \item \textbf{il rapporto umano}: nel lavoro si ``suda sangue'' fianco a
        fianco con altre persone, e questo crea legami molto forti che
        portano avanti l'industria nel senso più positivo del termine: nascono
        progetti privati, ci si porta da un'azienda all'altra, si formano team
        con basi solide.
\end{enumerate}

Va evitato anche l'eccesso opposto: la cultura del \emph{grinding} si è
appropriata di questi discorsi in modo distorto. Non bisogna diventare robot.

% ---------------------------------------------------------------------
\section{Organizzazione del corso e gerarchia}
\label{sec:gerarchia-corso}

\subsection{La struttura gerarchica}

La struttura dei referenti del corso ricalca volutamente una
\textbf{gerarchia tipica dell'industria}:
\begin{itemize}
  \item al vertice c'è il \textbf{direttore del corso} (Diego Ricciuti), che
        corrisponde al \emph{creative director}: supervisiona la visione del
        progetto (il concetto di ``visione'' verrà approfondito più avanti);
  \item sotto di lui ci sono i \textbf{referenti di dipartimento}, uno per
        disciplina, con un corrispettivo per ciascuna sede (Milano e Roma);
  \item in fondo ci sono gli studenti, che occupano il primo gradino di
        un'azienda: l'\emph{intern} (stagista) o il \emph{junior}, il passo
        che precede il lavoratore ``regolare''.
\end{itemize}

\begin{table}[htbp]
  \centering
  \begin{tabular}{@{}p{0.34\textwidth}p{0.29\textwidth}p{0.29\textwidth}@{}}
    \toprule
    \textbf{Disciplina} & \textbf{Milano} & \textbf{Roma} \\
    \midrule
    Direzione (\emph{creative director}) & \multicolumn{2}{p{0.58\textwidth}}{Diego Ricciuti} \\
    Game design      & Mattia Podini        & Giacomo ``Jack'' Danna \\
    Level design     & Stefano Zocchi       & Andrea Dispoli \\
    Narrative design & \multicolumn{2}{p{0.58\textwidth}}{Daniele Fusetto (lezioni condivise), Giacomo Masi} \\
    Programmazione (Unreal) & Alessandro Pregnolato & Matteo Velardi \\
    \bottomrule
  \end{tabular}
  \caption{Referenti del corso per disciplina e sede (nomi come riportati in lezione).}
  \label{tab:referenti}
\end{table}

Per il primo anno, il referente diretto degli studenti di game design è il
docente di game design: è a lui che ci si rivolge per dubbi e domande. Altre
figure utili sono il direttore didattico (severo, ``giustamente severo'', da
contattare solo per questioni particolarmente importanti), l'amministrazione
(contratti e documentazione) e la segreteria.

\subsection{I canali di comunicazione}

\begin{description}
  \item[Discord] è il canale per le comunicazioni ufficiali del corso e per
        le domande che non richiedono urgenza immediata; può avere qualche
        ritardo.
  \item[E-mail] è il canale per tutto ciò che è formale o ``ancora più
        ufficiale'' e per le comunicazioni rapide e immediate.
  \item[Telegram] esistono residui di vecchi canali precedenti a Discord, da
        evitare.
\end{description}

Le domande che valgono per tutti vanno fatte in classe (anche da un portavoce
a nome del gruppo); quelle singole e personali via Discord o, se formali, via
e-mail. Se il docente non risponde a una domanda c'è un motivo, ad esempio
ragioni legali, e lo dirà esplicitamente.

\subsection{I tre moduli e la settimana tipo}

Le tre discipline principali del primo anno (e non solo) sono:
\begin{itemize}
  \item \textbf{Game Design} (tendenzialmente il lunedì);
  \item \textbf{Level Design} (tendenzialmente il mercoledì);
  \item \textbf{Programmazione} (tendenzialmente il venerdì).
\end{itemize}
Nel primo mese le lezioni sono tutte di game design (lunedì, mercoledì e
venerdì). Sono previste anche lezioni di narrativa, ma il piano formativo si
aggiorna ogni anno: le ore di narrativa sono state ridotte a favore di quelle
di programmazione. Le poche lezioni di narrativa daranno quindi dei
\emph{cenni}, non la preparazione di un narrative designer. Ciò che serve
davvero è imparare \textbf{come comunicano gli altri dipartimenti} e come
lavorare con persone che fanno un lavoro diverso ma complementare al proprio.

\subsection{Le due anime delle lezioni di game design}

\begin{description}
  \item[Design methods (metodologie)] lezioni prevalentemente teoriche, ma con
        una buona componente pratica, in cui si studiano metodologie di design
        e i cosiddetti \emph{framework} dell'industria.
  \item[Documentazione e pitching] il cuore del lavoro quotidiano del
        designer.
\end{description}

\paragraph{La documentazione.}
Il docente si definisce ``il più grande nerd della documentazione'' che si
possa incontrare. È una parte dura, perché a pochi piace documentare, ma il
lavoro del designer \emph{è principalmente questo}: scrivere documenti e, da
quei documenti, passare alla pratica e all'implementazione. Chi non sa
documentare bene parte male. Quando non si è presenti, il documento è il
proprio \textbf{biglietto da visita}: qualcun altro lo leggerà, quindi va
scritto nel modo più chiaro e comprensibile possibile.

Un esempio concreto del perché la documentazione conta: CD Projekt Red ha
sviluppato un proprio engine proprietario, ma buona parte della sua
documentazione è stata ``buttata all'acqua''. Chi deve usarlo o lo sa già
usare, o ha competenze di un altro engine da trasferire, oppure è in
difficoltà. Si documenta \textbf{perché altri possano leggere}.

\paragraph{Il pitching.}
Il \emph{pitching} è l'esposizione, la presentazione orale. La parte creativa
del team presenta spesso:
\begin{itemize}
  \item le proprie idee ai colleghi;
  \item il gioco a convention ed eventi;
  \item il pitch del gioco agli investitori, per ottenere i fondi necessari a
        realizzarlo.
\end{itemize}
Saper parlare e comunicare chiaramente le idee è fondamentale. In sintesi:
\emph{il game designer parla}, e bisogna farsene una ragione.

\subsection{Altre componenti del corso}

\begin{description}
  \item[Gioco critico e playtest] la classe viene divisa in due gruppi: uno
        gioca, l'altro osserva. Chi osserva non applaude ``come le
        scimmiette'': annota, segna e analizza il \emph{playtest}, esattamente
        come si fa nell'industria. Il playtest fa parte delle mansioni del
        designer, ma esiste anche come specializzazione autonoma.
  \item[Pipeline] termine che indica tutto ciò che riguarda i processi
        formalizzati dell'industria; di solito si intende la \emph{pipeline di
        produzione}: come si divide il team, qual è la sua gerarchia, come il
        \emph{producer} gestisce le deadline. In breve: \emph{come} e
        \emph{quando} si fanno le cose. Al primo anno è meno rilevante (si
        dovrà ``nuotare'' un po' da soli); dal secondo anno si vedranno le
        metodologie effettivamente applicate ai progetti digitali. Si
        scoprirà quanto sia difficile lavorare a un progetto digitale, e
        ancor più fisico, senza un piano di sviluppo comprovato, come quelli
        usati per qualsiasi software in qualsiasi parte del mondo.
  \item[Engine] il motore di riferimento è \textbf{Unreal Engine}
        (sezione~\ref{sec:engine}).
  \item[Programmazione] principalmente tramite \emph{Blueprint} in Unreal,
        più eventuali sistemi di versionamento; il corso parte da zero.
  \item[Progetti pratici] probabilmente la cosa più importante. Al primo
        anno i progetti di gruppo restano circoscritti al corso di design,
        perché mancano ancora alcune nozioni (come comunicare, affrontare
        problemi, dare e ricevere feedback). In seguito si allargheranno a
        progetti digitali, molto più difficili di quelli fisici, con membri
        di altri team: un rapporto non unilaterale, ma ``covalente''.
\end{description}

% ---------------------------------------------------------------------
\section{Il metodo: l'inquiry-based learning}
\label{sec:inquiry}

\subsection{Che cos'è}

Il corso è costruito secondo l'\textbf{inquiry-based learning}
(apprendimento basato sull'indagine), un approccio pedagogico nato intorno
agli anni Sessanta, che si richiama a pensatori come Dewey e Piaget. In
generale, le discipline da cui il design attinge sono quasi sempre la
sociologia, la psicologia e, talvolta, l'economia.

Il funzionamento è un ciclo:
\begin{enumerate}
  \item il docente fornisce metodologie e framework dell'industria durante
        la lezione;
  \item lo studente, fuori dalle ore di lezione, lavora sui \emph{task}:
        mette a terra le nozioni, le maneggia, le assorbe meglio;
  \item da questo lavoro nascono nuovi dubbi;
  \item i dubbi tornano in aula: lo studente chiede, il docente risponde
        (ammesso che dare la risposta sia utile, e ammesso che una risposta
        esista).
\end{enumerate}

\subsection{Non sempre esiste una risposta giusta}

Non tutto ciò che si studia ha ``la risposta migliore possibile in assoluto
nel vuoto''. Esistono metodologie, limiti e framework che aiutano il lavoro,
ma non sono fissi: cambiano da azienda ad azienda, da lavoro a lavoro, da
progetto a progetto, e a volte vengono stravolti. Sta alla
\textbf{sensibilità del designer} capire quale strumento o framework usare in
un dato momento, e questa sensibilità si impara solo con l'esperienza. Nel
lavoro nessuno dirà ``usa questa cosa'' o ``è meglio quest'altra'': sarà il
designer a scegliere.

Anche l'esperienza del docente ha confini: sei-sette anni di lavoro, un modo
personale di approcciare il design e un certo insieme di progetti (ad esempio
nessun MMO). Oltre quei confini si cercano risposte nella ricerca, studiando
come aziende con esperienza specifica (ad esempio Riot) hanno affrontato
problemi simili e hanno condiviso le loro soluzioni, sempre valutando quali
metodi siano adatti alla dimensione del team, al tipo di progetto e al tempo
disponibile.

\subsection{Sensibilità, non memoria}

Non interessa imparare a memoria un framework o la teoria di un autore. Certe
terminologie vanno conosciute, ma per ragioni di comunicazione e chiarezza.
Ciò che conta è sviluppare la \textbf{capacità di farsi domande}, di farle
agli altri, di fare ricerca e di trovare una propria risposta. È esattamente
ciò che si fa nel lavoro: davanti a qualcosa che non si sa fare, l'abilità
del designer sta nel cercare come qualcun altro ha risolto un problema simile.

\subsection{Le lezioni come corsi di aggiornamento}

Le lezioni frontali simulano i \emph{corsi di aggiornamento} aziendali: le
sessioni che un lead di design tiene al suo team di design (o il lead di
level design al suo team, e così via). Alla fine dell'aggiornamento c'è quasi
sempre un lavoro da fare, un obiettivo da portare a terra:
\textbf{l'obiettivo della settimana}. Si lavora di settimana in settimana, e
il perché verrà chiarito in seguito.

Struttura tipica di una lezione:
\begin{enumerate}
  \item breve sessione di domande e risposte sulla lezione precedente;
  \item analisi di eventuali esercizi particolarmente critici;
  \item lezione vera e propria;
  \item discussione di gruppo, se c'è tempo;
  \item assegnazione del task della settimana.
\end{enumerate}
Il docente tende a ``parlare e straparlare'', ad aprire parentesi e a saltare
da un argomento all'altro: l'unico modo per fermarlo è che siano gli studenti
a interromperlo con domande.

\subsection{La gerarchia della competenza}

L'inquiry-based learning può essere rappresentato come una scala di
competenza:
\begin{enumerate}
  \item la sola lezione frontale porta a \textbf{conoscenze incomplete},
        perché manca la parte pratica e di ricerca;
  \item il lavoro a casa (task, ricerca) porta a una competenza \textbf{più
        completa}, non sbagliata in sé, ma ancora parziale e complementare ad
        altre;
  \item la \textbf{condivisione} dei task e dei feedback con gli altri porta
        idealmente alla competenza più completa possibile;
  \item il picco della competenza è la sua \textbf{applicazione nei
        progetti} (fisici o digitali, di gruppo): ``con le cose che avete
        imparato e ricercato, riuscite a fare un gioco?''. Quando la risposta
        diventa positiva, o comunque si inizia a navigare nei processi di
        design per realizzarne uno, quella è la vittoria.
\end{enumerate}
Non si parte subito creando giochi: si parte dall'analisi, dal design e dalle
regole generali di design, prima ``nel vuoto'', poi contestualizzando i
framework, infine unendo i puntini nei progetti.

Rimarranno comunque delle lacune, per tutti, docente compreso, perché il game
design \textbf{non è una scienza esatta}, anzi, a rigore non è affatto una
scienza. Prende in prestito metodologie dalle discipline scientifiche, ma
contiene una componente imprevedibile e incalcolabile: \textbf{la
produzione}. Si possono applicare perfettamente le regole di un framework, e
quel framework può comunque non essere adatto al gioco, al proprio modo di
sviluppare o al tempo di produzione disponibile (alcuni framework richiedono
una ricerca e un'attenzione che il tempo non concede).

C'è poi la componente creativa, la parte ``romantica'' del lavoro: esiste, ma
è \textbf{metà del lavoro}. L'altra metà è ricerca, metodo scientifico al
100\%. Dagli stessi dati $n$ designer diversi, con sensibilità diverse,
possono ottenere $n$ risposte diverse. Non tutte saranno valide: vanno
applicate, messe a terra, e questa scrematura fa parte del lavoro.

\subsection{I livelli dell'apprendimento}

Un secondo modello, a gradini, descrive i livelli di apprendimento (è la
cosiddetta tassonomia SOLO):
\begin{description}
  \item[Pre-strutturale] il docente dà il compito e lo studente lo risolve,
        possibilmente il prima possibile. È esattamente ciò che \emph{non} si
        vuole.
  \item[Uni-strutturale] si affronta una sola disciplina: è il primo passo
        necessario, ed è il motivo per cui nel primo mese ci sono solo lezioni
        di game design.
  \item[Multi-strutturale] entrano in gioco altri dipartimenti: level design,
        game design, programmazione e così via.
  \item[Relazionale] si capisce che alcune cose valgono per più
        dipartimenti: alcuni framework di narrativa derivano dal design o
        possono essere usati per fare design; le regole di documentazione si
        applicano anche ad altri dipartimenti; le regole di level design si
        mescolano con quelle di game design. Dopotutto, per fare un gioco
        servono molti dipartimenti.
  \item[Astratto esteso] il punto finale: muoversi autonomamente e capire
        cosa usare e quando, sulla base della propria esperienza. L'unico modo
        per arrivarci è fare pratica e fare progetti.
\end{description}

% ---------------------------------------------------------------------
\section{Task, feedback e regole del corso}
\label{sec:task-feedback}

\subsection{I task come opportunità}

I task non sono strettamente obbligatori (esiste un numero minimo di
consegne), ma sono \textbf{opportunità}. Quanto e come ci si lavora dipende
dallo studente. Un task fatto in cinque minuti, o in cui si è stati ``furbi'',
non verrà letto né riceverà feedback.

Domande da porsi davanti a ogni task:
\begin{enumerate}
  \item \textbf{Perché} mi è stato assegnato? Qual è l'obiettivo?
  \item Come ci lavoro per ottenere un risultato \textbf{che abbia senso}? Se
        la ricerca e i dati indicano una direzione, è inutile andare in
        un'altra.
\end{enumerate}

\subsection{Quanto lavorare: un calcolo indicativo}

Se il corso è una simulazione di lavoro, vale il carico di un lavoro normale:
\begin{itemize}
  \item un lavoro a tempo pieno comprende da 30 a 45 ore settimanali; si
        considera una media di \textbf{40 ore};
  \item si sottraggono circa 9 ore di lezione: restano $\approx 31$, diciamo
        \textbf{30 ore};
  \item con circa 3 task a settimana: $30 / 3 = $ \textbf{circa 10 ore per
        task}.
\end{itemize}
Il numero può variare, ma rende l'idea: dedicare un'ora a un task e poi fare
altro non fa guadagnare nulla a nessuno. Lo studente ``butta soldi'', il
docente risparmia tempo leggendo il task in un minuto (o non leggendolo).
Dedicando invece attenzione, ricerca, cura nella formattazione e
sperimentazione, e facendolo con cognizione di causa, si impara.

\subsection{Niente voti, solo feedback}

Ad eccezione del corso di programmazione, non ci sono voti (né ``ottimo'', né
10, né 100). Ci sono \textbf{feedback}:
\begin{itemize}
  \item in game design, generalmente sulla documentazione;
  \item in level design, sui progetti pratici.
\end{itemize}
Il docente indica cosa va bene e cosa va modificato, concentrandosi in
particolare su ciò che va cambiato. Il feedback \textbf{non è mai un giudizio
sulla persona}, ma su ciò che è stato consegnato nel concreto.

\subsection{La condivisione dei task}

Tutti i task consegnati vengono \textbf{condivisi} con la classe, e possono
essere letti solo dopo aver ricevuto i feedback del docente. Si cresce così
non solo con i propri feedback, ma anche con quelli dati agli altri. Regola
importante: ciò che il docente dice a uno studente \textbf{vale per tutti}. Se
a una persona viene detto ``questa cosa non va mai più fatta'', da quel
momento vale per l'intera classe.

\subsection{Non copiare}

Copiare i task degli altri è una delle cose peggiori che si possano fare, non
per una questione morale, ma perché annulla il senso del percorso. Si può
dedicare a ogni task il tempo e l'attenzione che si ritengono opportuni, ed è
normale che col tempo ci si specializzi (preferendo, ad esempio, il level
design o la programmazione tecnica) e si investa meno dove si sente meno
necessità. In questi casi l'alternativa corretta è:
\begin{itemize}
  \item non consegnare;
  \item studiare ciò che hanno fatto gli altri;
  \item eventualmente riconsegnare in seguito.
\end{itemize}
Ci saranno momenti in cui si verrà sommersi dalle cose da fare: bisogna
imparare a scegliere.

\subsection{Il ``non leggo''}

Il docente ha piena facoltà di non dare alcun feedback a un task (il
cosiddetto \textbf{``non leggo''}) se:
\begin{itemize}
  \item non si è lavorato con attenzione;
  \item non si rispettano le indicazioni date;
  \item è stata violata una regola estremamente importante.
\end{itemize}
Un task ``non letto'' equivale a un task non consegnato e non entra nel
conteggio. Attenzione: un task \emph{perfetto} e un task che \emph{fallisce}
nel comprendere o applicare un framework sono \textbf{ugualmente validi},
purché ci si sia impegnati e si sia dedicato tempo, pazienza e ``sudore''.

\paragraph{Riconsegna.} Un task non letto, in ritardo o che si vuole
migliorare può essere riconsegnato e verrà conteggiato, a patto che:
\begin{itemize}
  \item i problemi segnalati siano stati risolti;
  \item sia stato fatto ulteriore lavoro, integrando anche tutto ciò che nel
        frattempo è stato detto agli altri colleghi nei loro feedback.
\end{itemize}
Il task va sviluppato con le competenze che si hanno nel momento della
riconsegna, non con quelle del momento della prima consegna.

\subsection{Requisiti per superare l'anno}

\begin{itemize}
  \item consegnare almeno l'\textbf{80\% dei task} assegnati (soglia
        malleabile);
  \item partecipare almeno all'\textbf{80\% delle lezioni}, in presenza o
        da remoto, preferibilmente in presenza; i ritardi vengono valutati
        caso per caso;
  \item comportarsi da persone corrette e rispettose (è stata firmata una
        carta dello studente): non urlare, non alzare le mani, non
        vandalizzare, non compiere atti illegali.
\end{itemize}
Impegnarsi di più porta inoltre a essere selezionati per varie opportunità.

\subsection{Scadenze}

\begin{itemize}
  \item la scadenza di un task è \textbf{sempre prima dell'inizio della
        lezione} successiva indicata;
  \item se la classe, come team, mantiene comportamenti scorretti, il docente
        può \textbf{anticipare le deadline} per tutti.
\end{itemize}

\subsection{Accessibilità e DSA}

Gli studenti con disturbi specifici dell'apprendimento o neurodivergenze
(dislessia, ADHD, autismo, ecc.) sono invitati a segnalarlo. Quando si
affronteranno le regole di formattazione e documentazione (errori
grammaticali, mancanze, incoerenze) se ne terrà conto e si potranno trovare
dei compromessi, pur ricordando che nel mondo del lavoro nessuno farà
``sconti''. La documentazione può essere inviata all'amministrazione o alla
segreteria, che la comunicherà a tutti i docenti, ma il consiglio è parlarne
anche direttamente con ogni singolo docente.

% ---------------------------------------------------------------------
\section{Le competenze trasversali}
\label{sec:soft-skill}

\subsection{I giochi si fanno in gruppo}

Si possono fare giochi da soli, ed è un'impresa ammirevole; ma lavorando da
soli \textbf{tutto ciò che si guadagna in libertà si perde in produttività}.

\paragraph{Il caso Stardew Valley.} \emph{Stardew Valley} è stato realizzato
da una sola persona, Eric Barone, in anni di sviluppo e di sacrifici: faceva
un altro lavoro per mantenersi (in un cinema) e di notte programmava, faceva
l'arte e componeva la musica. Il risultato è stato straordinario e lo ha reso
ricco, ma si tratta di un caso estremamente raro. Quanti giochi simili, nati
nelle stesse condizioni, sono falliti? Molti di più.

\paragraph{Il rischio d'impresa.} Lavorare da soli non è sbagliato di per sé
e ha aspetti molto positivi: si ha il privilegio di creare esattamente ciò che
si vuole. Ma ci si fa carico anche di tutto il \textbf{rischio d'impresa}: se
si lavora a un gioco per cinque anni e poi il gioco va male, resta il
vantaggio di averlo pubblicato, ma il dopo è sperare di vivere di quella
release o di finanziare con essa il gioco successivo. Le grandi aziende hanno
gli stessi problemi, ma li distribuiscono su più persone e più responsabilità.

L'obiettivo di un professionista è \textbf{saper lavorare in diverse
condizioni}, e la più comune è quella del dipendente (o del freelance /
libero professionista con partita IVA): si viene assunti per le proprie
capacità da un'azienda, piccola o grande che sia, e si lavora con altre
persone.

\subsection{Essere persone con cui si vuole lavorare}

Se le persone non vogliono lavorare con voi, quello è \textbf{il primo motivo
di licenziamento}. E se non si viene licenziati, il clima nell'azienda
diventa terribile: la produttività crolla, la motivazione crolla. Le
\emph{soft skill}, cioè l'insieme di abilità che rendono una persona
piacevole con cui lavorare, vanno coltivate e diventano sempre più importanti
col passare degli anni (relativamente meno ora, molto di più quando i team
saranno più grandi). Anche chi è un ``lupo solitario'' deve imparare a:
\begin{itemize}
  \item lavorare con altre persone;
  \item comunicare con altre persone;
  \item trovare un punto in comune da cui partire per fare le cose.
\end{itemize}

\subsection{Indipendenza e proattività}

Sono le altre due parole chiave del corso:
\begin{description}
  \item[Indipendenza] spesso si verrà lasciati da soli, non per cattiveria ma
        per sfida: quando si cerca una risposta da soli, ciò che conta non è
        tanto la risposta quanto \textbf{il processo}. Una risposta data dal
        docente non vale niente.
  \item[Proattività] chiedere, domandare, agire da soli senza aspettare
        direttive, perché spesso quelle direttive non ci sono, a maggior
        ragione nei team piccoli.
\end{description}
Tutto questo si concretizza nel parlare, nel presentare e nel fare
\emph{post-mortem} (l'analisi a posteriori di un progetto concluso, che verrà
approfondita più avanti), e in generale nell'attenzione alla crescita
personale.

\subsection{Fare domande}

Il corso insiste sulle domande: alzare la mano, intervenire anche ``a
bomba'', fermare il docente ogni volta che qualcosa non è chiaro. Se una
persona non conosce qualcosa, molto probabilmente c'è almeno un'altra persona
che non la conosce. Il docente può dare per scontata la conoscenza di termini
che non lo sono, ed è del tutto normale, soprattutto all'inizio.

% ---------------------------------------------------------------------
\section{Strumenti: engine, Unreal e materiali}
\label{sec:engine}

\subsection{Unreal Engine e gli altri motori}

Il motore di riferimento del corso, usato in level design, è \textbf{Unreal
Engine}: ``un gran mattone'', ma che dà a un designer molte possibilità. Il
vantaggio principale è che permette a un designer non pratico di
programmazione di \textbf{prototipare tramite i Blueprint} (il sistema di
scripting visuale di Unreal). Questo non impedisce di passare a Unity o ad
altri engine, e viceversa: il docente stesso ha studiato e lavorato a lungo
con Unreal, ma la sua ultima esperienza professionale (come level designer) è
stata su Unity, un engine che non conosceva, a dimostrazione che il salto da
un engine all'altro è possibile.

Panorama dei motori citati:
\begin{itemize}
  \item \textbf{Unreal Engine} e \textbf{Unity}: i due più diffusi, con
        modelli di sorgente ``più o meno'' aperto;
  \item \textbf{Godot}: engine open source;
  \item \textbf{CryEngine}: proprietario di Crytek;
  \item \textbf{Snowdrop}: engine di Ubisoft;
  \item Rockstar usa un proprio engine proprietario;
  \item alcune aziende usano versioni \textbf{personalizzate} di engine
        esistenti: ad esempio Milestone, nei giochi di corse, usa un Unreal
        customizzato;
  \item CD Projekt Red ha un engine proprietario (con i problemi di
        documentazione già citati).
\end{itemize}
In sintesi: ogni azienda può avere un engine proprietario oppure una versione
modificata di un engine commerciale, ma si parte sempre da una base.

\subsection{Programmazione: cosa ci si aspetta da un designer}

Non è richiesto alcun bagaglio pregresso di programmazione: il corso parte da
zero. Gli studenti sono stati selezionati (il corso è a numero chiuso) per
competenze analitiche e critiche sopra la media. Anche nel mercato del lavoro,
per una posizione di game designer generalmente \textbf{non ci si aspetta
competenza di scripting} (C++, Java, C\#, ecc.), a meno di voler fare il
\emph{technical designer}, e anche lì varia da azienda ad azienda.

\subsection{Versione e preparazione}

Indicazioni pratiche emerse in aula:
\begin{itemize}
  \item scaricare Unreal Engine; la versione consigliata per il corso è
        stata indicata nella \textbf{5.6} (con conferma successiva). La
        stabilità dipende da cosa si deve fare, da quali tool e plugin si
        usano: i problemi nascono spesso dai plugin. Nel passaggio da Unreal 4
        a Unreal 5, ad esempio, il sistema \emph{Landscape} era inizialmente
        inutilizzabile, mentre in Unreal 4 funzionava molto bene;
  \item è consigliato installare tramite \textbf{Epic Games Launcher}, più
        comodo: permette anche di riscattare giochi gratuiti periodici e di
        accedere ai pacchetti di \textbf{Fab} (il marketplace di asset
        acquisito da Epic), che ogni mese offre alcuni asset gratuiti utili
        per esercitarsi nel level design (ad esempio per fare \emph{speed
        level design});
  \item molti pacchetti di UI o di sistemi di interazione sono utili per
        prototipare al volo, ma spesso inutilmente complicati;
  \item un plugin di \textbf{materiali di prototipazione a griglia} (adattivi,
        con metriche, allineabili alla griglia, con materiali già pronti e
        colori modificabili, dal costo di circa 5 euro) è indicato come
        ``salvavita'' per il level designer; se necessario, i pacchetti
        indispensabili vengono forniti;
  \item portare \textbf{sempre il portatile} alle lezioni di level design
        (da novembre, a partire dall'interfaccia dell'engine); alle lezioni
        di game design basta che lo porti almeno metà della classe, per i
        lavori a coppie. Serve comunque qualcosa per prendere appunti
        (``tavolette d'argilla e cannucce, se volete fare i babilonesi'',
        oppure carta e penna).
\end{itemize}
Chi ha un portatile poco potente dovrebbe comunque provare a installare
l'engine e aprire un template per vedere come gira: all'inizio non si faranno
cose pesanti. Nelle prime lezioni di level design è importante mettere subito
le mani sull'interfaccia: ascoltare dove si trova uno strumento (ad esempio lo
\emph{snapping}) senza provarlo porta a dimenticarlo.

% ---------------------------------------------------------------------
\section{Letture e risorse di riferimento}
\label{sec:risorse}

\subsection{I libri}

\begin{description}
  \item[\emph{The Art of Game Design: A Book of Lenses} (Jesse Schell)] il
        testo \textbf{davvero importante}. Alcune lezioni finali dell'anno si
        baseranno sulle sue terminologie. È utile perché apre la mente sulla
        possibilità di applicare metodologie diverse a giochi diversi, ma
        soprattutto per la \textbf{terminologia} comune. Esiste anche in
        audiolibro (circa 25 ore).
  \item[\emph{Challenges for Game Designers} (Brenda Romero e Ian Schreiber)]
        libro \emph{entry level} con moltissimi esercizi, molto caro al
        direttore del corso.
  \item[\emph{21st Century Game Design} (Chris Bateman)] utile come
        retrospettiva: è un libro visionario che ha previsto metodi e
        soprattutto \emph{archetipizzazioni} dei giocatori tuttora valide nel
        design e nel marketing. Servirà per una delle lezioni più importanti
        dell'anno. È utile leggerlo per vedere quanto aveva previsto, ma è
        utile ``fino a un certo punto''.
  \item[\emph{An Architectural Approach to Level Design} (Christopher
        Totten)] non strettamente necessario, ma di lettura semplice. Ha due
        anime: parte dall'architettura (babilonese, greca, ecc.) e la trasla
        su esempi di level design (come \emph{Doom} e \emph{Team Fortress}).
        Un tempo ne esisteva un riassunto su Gamasutra (oggi
        \emph{Game Developer}).
  \item[\emph{A Theory of Fun for Game Design} (Raph Koster)] utile fino a un
        certo punto: contiene diverse idee superate, che lo stesso Koster ha
        affrontato in vari talk reperibili nel GDC Vault. Verrà ripreso nelle
        prime lezioni, poi messo da parte.
\end{description}

\subsection{Canali e conferenze}

\begin{description}
  \item[GDC (Game Developers Conference)] il suo canale raccoglie talk e
        panel estremamente importanti: è ``uno strumento misterioso che ci
        servirà più tardi''. \textbf{Attenzione}: per molti talk mancano ancora
        le competenze per interpretarli, e alcuni potrebbero fuorviare o
        mettere in testa idee strane. La GDC si concentra spesso su
        produttività e mobile.
  \item[Il \emph{Dragon Speech} di Chris Crawford] il talk del fondatore della
        GDC: appartiene a un passato lontano, ma è molto motivante. Da vedere
        subito.
  \item[Game Maker's Toolkit] canale particolarmente utile per nozioni
        pratiche, con video brevi e alcune interviste.
  \item[Nordic Game] altra conferenza importante, con talk e panel di esperti
        del settore.
\end{description}

\subsection{Associazioni e premi: come leggerli}

Le associazioni di settore (ad esempio IIDEA in Italia, o la IGDA) non sono
il male, anzi: sono un canale importante per affacciarsi all'estero o per
accedere a fondi. Va però ricordato che \textbf{tutelano i propri associati}:
non esiste solo il loro palinsesto, e bisogna valutare quanto spazio occupano
nel panorama della visibilità.

Lo stesso vale per i \textbf{premi} (Game Awards, ``Best Mobile Game'',
``Games for Change'', ecc.): al designer \emph{non interessano}. Che un gioco
abbia vinto ``più di 90 premi'' non dice nulla di utile. Servono a dare
visibilità e lustro ai giochi ed è giusto che esistano, perché ampliano il
panorama e il numero di giocatori. Ma ciò che interessa al designer sono i
\textbf{dati}:
\begin{itemize}
  \item come sono state implementate certe cose;
  \item come è stato fatto il marketing e dove il gioco è stato presentato;
  \item non quante wishlist ha fatto, ma \emph{come} le ha fatte;
  \item in che periodo storico è uscito e quali erano i competitor in quel
        momento.
\end{itemize}

% ---------------------------------------------------------------------
\section{Che cosa fa un game designer}
\label{sec:cosa-fa}

Per arrivare a una prima definizione, la classe ha condotto un brainstorming:
ognuno doveva indicare una cosa che, secondo lui, fa il game designer (senza
ripetere quelle già dette). Ne è emersa la lista seguente, commentata punto
per punto.

\subsection{Progettare meccaniche e struttura del gioco}

Il ruolo più generico è quello del \textbf{core mechanic designer}, detto
anche \emph{gameplay designer} e, un tempo, \emph{black box designer}: si
occupa delle \textbf{meccaniche di gioco}. Si va dal piccolo al grande (in
funzione anche della dimensione dell'azienda):
\begin{itemize}
  \item nel \textbf{piccolo} ci sono le singole meccaniche;
  \item nel \textbf{grande} ci sono i \textbf{sistemi}. Nei giochi
        \emph{sistemici}, con molti sistemi interconnessi (ad esempio
        \emph{GTA}), esistono specializzazioni dedicate a bilanciare e
        supervisionare i sistemi: è il \emph{system design}.
\end{itemize}

\paragraph{Quando un gioco è ``sistemico''?}
Non dipende solo da come le meccaniche interagiscono con il giocatore (le
\emph{dinamiche}, secondo i framework che passano dalle meccaniche alle
dinamiche, che verranno studiati più avanti), ma anche da \emph{quante}
meccaniche ci sono e da quanto interagiscono tra loro. Non diremmo che
\emph{Super Mario} sia un gioco sistemico, eppure contiene dei sistemi: la
linea è sottile. Quando i sistemi aumentano molto, sono complessi e si
influenzano a vicenda, il gioco diventa sistemico. Esempio: in \emph{The
Legend of Zelda: Breath of the Wild} il vento modifica quanto il personaggio
sale o scende; il fuoco può incendiare l'erba; l'erba prende fuoco se colpita
da una freccia infuocata; e così via. Nel mercato moderno i giochi sistemici
sono spesso, ma non sempre, legati a mappe \emph{open world}.

\subsection{Le 3C: camera, character, control}

La base del gioco, da cui generalmente si inizia a lavorare (anche nei
prototipi), è la ``sacra trinità videoludica'' delle \textbf{3C}:
\begin{description}
  \item[Camera] che tipo di inquadratura ha il gioco: prima persona, terza
        persona, isometrica, ecc.
  \item[Character] quanto è grande il personaggio, che tipo di personaggio è
        e, di conseguenza, quali azioni può compiere.
  \item[Control] i controlli: che cosa il personaggio può fare, quali sono i
        \emph{verbi} di gioco.
\end{description}
Le 3C sono definite dal designer, ed esiste anche una specializzazione per
questo lavoro. Esistono varianti: quella del direttore del corso è
\textbf{4C + J}, cioè \emph{camera, character, control, conflict} (conflitto)
e \emph{juice} (il ``succo'', il feedback che rende piacevole l'interazione).
Verranno approfondite più avanti.

\subsection{Il personaggio}

Il designer definisce il personaggio dal punto di vista delle meccaniche
(dimensioni, azioni, capacità), all'interno delle 3C. L'aspetto estetico è
invece competenza di altri dipartimenti.

\subsection{Documentare le informazioni per gli altri dipartimenti}

Assolutamente vero: la documentazione è una delle prime attività del
designer. Più l'azienda è grande, più questo lavoro diventa prominente; più è
piccola, meno è rilevante. In \textbf{pre-produzione}, però, è sempre
fondamentale.

\subsection{Mediatore e leader ``al servizio''}

Il designer è un \textbf{mediatore} nei conflitti interni al team, una sorta
di collante che si relaziona con dipartimenti diversi (non sempre gli stessi,
a seconda della fase di produzione). A volte fa anche da leader, a seconda di
come sono gestite la pipeline e la gerarchia.

Questa leadership è difficile da definire perché \textbf{non è verticale}: è
una leadership \emph{al servizio} (\emph{servant leadership}). Il designer
``è asservito al programmatore'': il suo compito non è arrivare con un'idea e
scaricarla addosso agli altri, né portare la soluzione e lasciare il problema
a qualcun altro, ma \textbf{seguire il processo a basso livello}, sporcandosi
le mani.

\paragraph{I feature team.} Nei team medio-piccoli o medio-grandi possono
esistere sotto-team (\emph{feature team}) dedicati a un obiettivo specifico:
ad esempio un designer incaricato del character o della camera, affiancato da
uno o due programmatori e da un artista. In questi casi il designer
supervisiona il lavoro e gestisce la produzione del sotto-team.

\subsection{Gestire l'avanzamento del prodotto}

Esiste una specializzazione che si occupa solo di questo, il
\textbf{producer}, a un livello gerarchico completamente diverso (e che
potrebbe essere una strada percorribile al termine del percorso di studi).
All'interno di un feature team anche il designer può farlo, ma non è detto.
Può certamente \emph{dare la direzione}: nella sua documentazione può
indicare che per arrivare a uno stato giocabile del prototipo serve un certo
lavoro di arte e di programmazione, per lo stato successivo un altro, e così
via. In ambito produttivo esistono metodologie strutturate come
\textbf{Agile}, che verranno viste in seguito.

\subsection{Presentare le idee}

Il pitching (sezione~\ref{sec:gerarchia-corso}) è richiesto al designer sia
in pre-produzione sia in produzione, in misura che dipende dall'azienda. Ad
esempio Ubisoft lascia ai designer, e in realtà a qualsiasi membro, la
possibilità di \emph{pitchare} un'idea prima di una nuova produzione, perché
possa iniziare a prendere vita.

\subsection{Bozze e prototipi}

Le bozze vanno bene, ma devono essere \textbf{chiare}. C'è un'eterna lotta tra
documentazione e gioco giocabile: se si fa solo documentazione il gioco non
avanza; se si fa solo il gioco manca la documentazione e si perde la visione.
Sono due cose \textbf{complementari}.

Il modo per comunicare un'idea in forma grezza ma giocabile è il
\textbf{prototipo}. Un prototipo può essere qualsiasi cosa: di carta, in
Unreal con Blueprint ``brutti''. \emph{Non deve essere bello, deve essere
giocabile}.

\paragraph{Esempio: il prototipo di carta di Mario + Rabbids.} Durante lo
sviluppo di \emph{Mario + Rabbids}, i programmatori di Ubisoft erano impegnati
su altri progetti (come \emph{Ghost Recon}). I designer realizzarono allora
un \textbf{gioco da tavolo} che riproduceva le meccaniche del gioco, quelle
di uno strategico a turni sul modello di \emph{XCOM}. E funzionava.
L'obiettivo era \textbf{validare un'idea}, validare un gameplay.

Altri modi per rendere chiara un'idea sono i \emph{mock-up} nella
documentazione e le immagini modificate o i \emph{mash-up} di qualsiasi
natura. Lato documentazione l'importante è che sia chiaro; lato prototipo,
che sia giocabile.

\subsection{Fare ricerca per l'idea, e poi smontarla}

Il designer ha una componente creativa importante e fa ricerca per arrivare a
un'idea. Ma ciò che deve poi fare necessariamente è \textbf{smontarla}. Il
suo lavoro è scientifico, esattamente come il \textbf{metodo scientifico}, il
cui obiettivo è \emph{confutare un'ipotesi}: si arriva a un'ipotesi tramite
sperimentazione e poi si cerca $n$ volte di confutarla; se resiste, la si
assume come valida \emph{in un certo contesto} (le regole della fisica
classica, ad esempio, non valgono nella meccanica quantistica).

Allo stesso modo il designer arriva a un'idea (tramite ricerca o chiedendo
agli altri: le idee le hanno tutti, non sono un privilegio del designer),
cerca di smontarla e poi cerca di \textbf{metterla a terra}. Se l'idea non è
applicabile, non è implementabile, finché non si trova un modo per
realizzarla va scartata. Può essere l'idea più bella dell'universo, ma
\textbf{se non può essere messa a terra, non esiste}.

\paragraph{L'originalità.} Serve un'idea originale? Dipende dall'obiettivo e
dalla visione del progetto:
\begin{itemize}
  \item se l'obiettivo è vendere molto, un'idea originale non serve
        necessariamente;
  \item se si vuole creare un nuovo filone, serve un'idea, se non originale,
        almeno interessante: può essere un discorso di marketing o una
        meccanica innovativa, con il rischio di non vendere.
\end{itemize}
Il ruolo del \emph{director} e il concetto di \emph{visione} verranno
approfonditi in seguito.

\subsection{Analizzare i competitor}

Arrivati a un'idea (``vogliamo un certo tipo di combattimento''), si va a
vedere se qualcuno l'ha già fatta e \emph{come}. Le meccaniche non sono
privilegio di nessuno, salvo quando vengono brevettate, come il
\textbf{Nemesis System} di \emph{La Terra di Mezzo: L'Ombra di Mordor}.

\paragraph{Esempio: il combattimento da ridimensionare.} In un progetto
(presso Stormind Games) il team si trovò con un sistema di combattimento
troppo costoso: mancavano le risorse e gli animatori per realizzarlo. Poco
output lato gameplay significava non poter fare qualcosa di complesso e
profondo come in \emph{Elden Ring}; al contrario, il team aveva molti
artisti e poteva puntare su qualcosa di più artistico. Si iniziò così a
orbitare intorno al corpo a corpo di \emph{The Last of Us} (parte I e II).
Sul canale della GDC esistono diversi talk che raccontano cosa è stato fatto
e come, con successi, insuccessi e costi (ad esempio del \emph{motion
capture}): da quei talk si ricavarono moltissime informazioni utili.

I competitor possono essere:
\begin{itemize}
  \item \textbf{diretti}: un gioco molto simile al proprio;
  \item \textbf{di mercato}: un gioco che domina il mercato in cui ci si vuole
        affacciare;
  \item \textbf{indiretti}: un gioco con un sistema simile al proprio, ma
        tutto il resto diverso.
\end{itemize}
Regola d'oro: \textbf{là fuori, nascosto nell'ombra, c'è sicuramente un gioco
molto simile al vostro}. Va solo trovato. E se per qualche ragione non ha
venduto, bisogna capire perché.

\subsection{Lavorare in engine e fare tuning}

Il level designer lavora in engine, così come il technical designer; ma
qualsiasi game designer, a seconda di ciò che deve fare, può finire in
engine, soprattutto per il \textbf{bilanciamento} e il \textbf{tuning}.

Il \emph{tuning} è l'azione di bilanciare con precisione i dati. Esempio: in
\emph{Super Mario} si modifica quanto in alto salta il personaggio, con che
velocità, in quanto tempo raggiunge l'apice, e così via. Questo lavoro rende
il gioco:
\begin{itemize}
  \item ad alto livello, dotato della \textbf{progressione} giusta;
  \item a basso livello, \textbf{piacevole da giocare}.
\end{itemize}
Il \emph{fine tuning} è un lavoro di dettaglio, paziente e minuzioso.

\subsection{Bilanciare}

Il bilanciamento, emerso tardi nel brainstorming, è una delle attività
centrali del designer e ha anche una specializzazione dedicata
(l'\emph{economy designer}, sezione~\ref{sec:tipi-designer}).

\subsection{Analizzare il mercato}

In pre-produzione, di solito, sì: capire le tendenze, i generi che cambiano
nel tempo, la compatibilità del gioco con il momento storico, in vista ad
esempio di un possibile publisher. Va capito quanto, come e con quanto tempo
farlo, perché si è sempre limitati da molte cose.

\subsection{Sintesi}

\begin{quote}
Il designer parte dall'idea ma non si ferma lì: la confuta, lavora sui
dettagli, bilancia, fa \emph{fine tuning}; lavora con il team e fa da
\emph{manager} per diversi dipartimenti; e naviga costantemente tra i limiti.
\end{quote}

% ---------------------------------------------------------------------
\section{I limiti del design}
\label{sec:limiti}

I limiti ci sono sempre, e nel mondo del design sono spesso numerosi. I più
importanti entro cui navigare sono:
\begin{description}
  \item[Tempo] ad alto livello, l'idea deve essere realizzabile nel tempo di
        produzione disponibile; a basso livello, il proprio lavoro
        settimanale deve essere fattibile nel tempo che si ha. Non ci si può
        dividere tra mille feature team se realisticamente se ne può seguire
        uno solo.
  \item[Tecnologia] si può o non si può fare qualcosa? Ci sono i server per un
        multiplayer con un certo numero di giocatori? Ci si può permettere,
        ad esempio, una fauna che reagisce alla fisica come quella di
        \emph{GTA VI}, che richiede il meglio dei \emph{technical artist}?
  \item[Background e competenze] che cosa sa fare il team? Realisticamente,
        chi ha lavorato dieci anni su giochi horror viene assunto per un
        horror, perché ha le competenze per risolvere problemi già visti.
  \item[Persone e motivazione] il team ha la motivazione e le competenze per
        fare un certo tipo di gioco, e per lavorare insieme?
  \item[Denaro] tutto ruota intorno ai soldi. Che si lavori da soli o in mille,
        le persone vanno pagate e il gioco deve avere un ritorno economico.
        Bisogna decidere dove tracciare la linea: raggiungere il
        \textbf{break-even} (guadagnare esattamente quanto speso) o puntare a
        grandi guadagni. Con un \textbf{publisher}, inoltre, una parte
        consistente dei ricavi va a lui (ripartizioni dell'ordine di 80/20 o
        70/30, che dipendono dal contratto firmato).
\end{description}

\subsection{Non siete ``gli uomini delle idee''}

Per qualcuno può essere un problema scoprirlo, ma il designer non è
``l'uomo (o la donna) delle idee'': accanto alla parte creativa c'è una parte
tecnica, e si è fortemente limitati da molti fattori. Questi limiti, però,
sono un aiuto:
\begin{itemize}
  \item riducono il ventaglio di possibilità potenzialmente distruttive per il
        team (cose che non si riescono a fare, che non piacciono o per cui
        mancano le risorse);
  \item aiutano a mirare bene, trovando soluzioni creative in uno spazio più
        piccolo e più attuabile.
\end{itemize}

\begin{quote}
\centering
\textbf{Le idee sono nulla, l'implementazione è tutto.}
\end{quote}

% ---------------------------------------------------------------------
\section{Ruoli e specializzazioni}
\label{sec:tipi-designer}

\subsection{La scala gerarchica}

\begin{description}
  \item[Creative Director (CD)] mantiene la \textbf{visione} e dà l'idea, in
        modo che tutto segua una certa direzione. In teoria può diventarlo
        chiunque, ma con un'esperienza notevole, nell'ordine dei decenni: si
        assume che in quel tempo abbia visto come lavorano molti dipartimenti.
        Naturalmente un CD con un background artistico ha meno competenze di
        design, e viceversa.
  \item[Director di disciplina] uno per dipartimento: \emph{game director},
        \emph{level director}, \emph{art director}, ecc.
  \item[Lead] il capo del dipartimento, che gestisce i gruppi di persone. Chi
        lavora in un dipartimento risponde al proprio lead (ad esempio un game
        designer risponde al \emph{lead designer}, non al game director).
  \item[Senior / Principal] i profili più esperti (a seconda dell'azienda).
  \item[Junior] il profilo di inizio carriera.
  \item[Intern] lo stagista.
\end{description}
Lo schema si replica uno a uno per ogni dipartimento: arte, suono, animazione
(a volte arte e animazione insieme, o \emph{rigging} e animazione insieme),
programmazione, UX/UI e così via. A fare la differenza, a parità di ruolo, è
il \textbf{livello di seniority}.

\subsection{UX e UI}

UX e UI stanno a metà tra game design e arte: a seconda dell'azienda il ruolo
si chiama \emph{UX/UI Artist} o \emph{UX/UI Designer}.
\begin{description}
  \item[UX (User Experience)] come una persona interagisce con le cose. Esempio:
        su uno smartphone, un pulsante accende e spegne lo schermo; per
        digitare la password si scorre verso l'alto e compaiono i punti da
        toccare.
  \item[UI (User Interface)] come sono ``vestite'' queste cose: ad esempio
        l'aspetto della barra dei punti vita (HP), realizzata con asset 2D o
        3D.
\end{description}

\subsection{Le specializzazioni del game design}

Più l'azienda è grande, più la specializzazione è puntuale.
\begin{description}
  \item[Gameplay / Core mechanics designer] (sono la stessa cosa) si occupa
        delle meccaniche e delle 3C.
  \item[System designer] si occupa dei sistemi e delle loro interazioni.
  \item[Economy designer] si occupa del bilanciamento; lavora soprattutto con
        fogli di calcolo (Excel, Google Sheets).
  \item[Combat designer] si occupa del combattimento; negli annunci di lavoro
        indica quasi sempre un combattimento \emph{action}, con gameplay
        reattivo (non strategico), anche se non è una regola.
  \item[Multiplayer designer] si occupa delle interazioni multigiocatore.
  \item[Technical designer] fa da ponte tra game design e programmazione.
\end{description}

\subsection{Le specializzazioni del level design}

\begin{description}
  \item[Encounter designer] posiziona i nemici nella mappa in modo che gli
        scontri siano bilanciati e la progressione equilibrata come
        desiderato.
  \item[World / map designer] popola il mondo di gioco.
  \item[Quest designer] popola il mondo con le missioni e pensa alla
        progressione al loro interno, cioè a dove il giocatore deve andare (si
        vedano ad esempio gli annunci di lavoro di CD Projekt Red o di
        Larian).
  \item[Technical level designer] fa da ponte tra level design e
        programmazione.
\end{description}

\subsection{Discipline vicine: suono, UI, produzione}

Il monte ore del corso non permette di uscire come \emph{narrative designer}
(o \emph{sound designer}) completi. Ciò non toglie che si possano coltivare
queste discipline:
\begin{itemize}
  \item nei progetti di gruppo si può occuparsi anche del suono, della UI/UX,
        o provare ad applicare metodologie come Agile per fare
        \emph{producing} e gestire il team;
  \item il tempo dedicato a queste attività va però sottratto al game design,
        e dipende da quante persone ci sono nel team, da quanto lavoro c'è e
        da quanto è grande il progetto;
  \item avere sviluppato e rilasciato un progetto, anche piccolo (anche una
        \emph{game jam}), documentando i propri processi, è già un ottimo
        argomento da presentare in futuro;
  \item la via migliore resta lavorare nell'industria e, da lì, chiedere di
        imparare anche altre discipline.
\end{itemize}
Un gioco si può fare su qualsiasi cosa: sugli scarafaggi, sui piedini delle
sedie, sui pennarelli. Si parte avvantaggiati da ciò che è vicino alla propria
sensibilità: chi ama la musica può pensare a un gioco con meccaniche legate
alla musica, con un gameplay innovativo ma contenuto e realizzabile. Il design
è un cappello che comprende molte discipline.

% ---------------------------------------------------------------------
\section{Approfondimenti emersi in lezione}
\label{sec:approfondimenti}

\subsection{Roguelike e roguelite}

Un esempio di come il marketing conia generi: Supergiant Games, con
\emph{Hades}, rese celebre il termine \textbf{roguelite} al posto di
\emph{roguelike} (pur non essendo il primo roguelite in assoluto). La
differenza:
\begin{itemize}
  \item il \textbf{roguelike} si fonda su una serie di punti fissati in una
        sorta di \emph{manifesto} che definisce che cosa rende un gioco
        roguelike; uno su tutti: a ogni morte si ricomincia la \emph{run} da
        capo;
  \item il \textbf{roguelite} rompe alcuni di questi punti, in particolare
        introducendo una \textbf{progressione esterna alla run}: il giocatore
        può potenziarsi fuori dalla singola partita per vincere più
        facilmente in quelle successive. \emph{Hades} ne è il caso principe.
\end{itemize}
Analogamente, un mobile game sviluppato dal direttore del corso fu
promosso come ``il primo \emph{donation game} al mondo'': un gioco che
convertiva il tempo di gioco in donazioni reali a enti benefici per animali.
Anche qui un neologismo creato lato marketing per definire un genere.

\subsection{Applicazioni fuori dall'industria videoludica}

Le competenze sviluppate non servono solo ai videogiochi:
\begin{itemize}
  \item i \textbf{render architettonici} oggi possono essere realizzati con
        Unreal, e diversi professionisti formati nel game design lavorano in
        questo campo;
  \item la \textbf{gamification}, cioè l'applicazione di regole videoludiche
        a contesti non ludici (mostre interattive, progetti culturali,
        consulenze come quella per applicare questi concetti alle macchine da
        allenamento di un'azienda di fitness), è in Italia uno dei modi più
        redditizi di spendere queste competenze, spesso con finanziamenti
        pubblici.
\end{itemize}

% ---------------------------------------------------------------------
\section{Il primo task}
\label{sec:primo-task}

\begin{enumerate}
  \item Creare un indirizzo e-mail \textbf{professionale} (consigliato: un indirizzo Gmail nella forma
        \texttt{nome.cognome@gmail.com}; l'anno di nascita è accettabile).
        Va bene anche un indirizzo professionale già esistente usato per
        lavoro. Da evitare indirizzi poco seri (``non \texttt{blackfire67}'').
  \item Una persona della classe raccoglie tutti gli indirizzi in un unico
        documento e lo invia via e-mail al docente, che aggiungerà tutti al
        Drive condiviso del corso (dove verranno caricate lezioni, slide e
        materiali, per avere i documenti tutti nello stesso posto).
  \item Installare Discord e iscriversi al server del corso, dove passano le
        comunicazioni ufficiali. È consuetudine, e utile, creare anche un
        gruppo tra studenti.
\end{enumerate}
\textbf{Scadenza}: prima della lezione del venerdì.

% ---------------------------------------------------------------------
\section*{Punti chiave del capitolo}
\addcontentsline{toc}{section}{Punti chiave del capitolo}

\begin{itemize}
  \item Il corso è una \textbf{simulazione di lavoro}: ci si tratta da
        colleghi, con una gerarchia da studio di sviluppo.
  \item Si impara con l'\textbf{inquiry-based learning}: lezione, task,
        dubbi, domande, condivisione, progetti.
  \item Il fallimento affrontato con la mentalità giusta è la principale fonte
        di crescita; i feedback riguardano il lavoro, mai la persona.
  \item Il lavoro del designer è per metà creativo e per metà
        \textbf{ricerca e metodo scientifico}; il game design non è una
        scienza esatta.
  \item \textbf{Documentazione} e \textbf{pitching} sono il cuore del
        mestiere.
  \item Il designer progetta meccaniche e sistemi (partendo dalle
        \textbf{3C}), documenta, media, prototipa, confuta le idee, studia i
        competitor, bilancia e fa tuning.
  \item La leadership del designer è \textbf{al servizio} del team, non
        verticale.
  \item I limiti (tempo, tecnologia, competenze, persone, denaro) restringono
        e indirizzano la creatività.
  \item \textbf{Le idee sono nulla, l'implementazione è tutto.}
\end{itemize}

%    % ...
%
%  (Con hyperref aggiungere anche \renewcommand{\theHchapter}{z.\arabic{chapter}}
%  nella parte zero, per evitare ancore duplicate.)
%  Tutte le etichette di questo capitolo hanno il prefisso "z2:"; i
%  rimandi al capitolo precedente usano le etichette "z1:" già definite in
%  GD0_1_cultura_del_feedback.tex.
% =====================================================================

\chapter[Tempo, idee e leadership]{Gestire il tempo, generare idee, guidare da dentro}
\label{cap:z2-tempo-idee-leadership}

Questo capitolo raccoglie la \textbf{seconda giornata} del workshop di due
giorni con cui si aprono le lezioni zero, tenuto dallo stesso relatore del
capitolo~\ref{cap:z1-feedback}, Stefano Cozzi\footnote{Per il cognome e per
il nome dello studio (Stormind Games) valgono le avvertenze della nota
all'inizio del capitolo~\ref{cap:z1-feedback}. In questa giornata la
trascrizione riporta soltanto il nome di battesimo, ``Stefano'', con cui il
relatore si cita più volte nei suoi esempi.}. Se la prima giornata spiegava
\emph{perché} le aziende tengono tanto alla cultura del feedback, la seconda
è dichiaratamente più pratica: ``ieri vi ho spiegato perché, oggi vi insegno
come lo fanno queste aziende''. I concetti della prima giornata (il
\emph{default to trust}, il ``potrei non farcela'' detto in anticipo, la
condivisione dei successi, il feedback frequente) ritornano di continuo,
applicati alla gestione del tempo, alle riunioni creative e alla leadership.

La giornata si è svolta di domenica, a chiusura di un fine settimana di
workshop\footnote{Il relatore dice ``oggi chiaramente è domenica'' e saluta
la classe con ``buon resto del weekend''; la prima giornata si era quindi
svolta di sabato. Più avanti parla di ``dodici ore'' di workshop in totale,
mentre nel capitolo~\ref{cap:z1-feedback} si parlava di circa dieci ore: la
differenza dipende probabilmente dal conteggio delle pause.}, con il
relatore alle prese con la voce un po' calata dal giorno prima (``la sfida
principale di oggi sarà contro la mia voce'').

Il percorso del capitolo segue quello della lezione:
\begin{enumerate}
  \item il programma della giornata (sezione~\ref{sec:z2-programma});
  \item le domande sulla prima giornata, che si trasformano in una lunga
        sessione su \textbf{colloqui e selezione}
        (sezione~\ref{sec:z2-colloqui});
  \item il \textbf{time management}: pianificare la giornata, le liste, le
        checklist, le interruzioni (sezione~\ref{sec:z2-pianificare});
  \item la \textbf{matrice di Eisenhower} per prioritizzare
        (sezione~\ref{sec:z2-eisenhower});
  \item come proteggere il proprio tempo: sprechi, distrazioni, il saper dire
        di no (anche ai meeting), il mito del multitasking e la chiusura della
        giornata (sezione~\ref{sec:z2-proteggere});
  \item le domande sul time management e una parentesi su tempo libero, ferie
        e benessere in azienda (sezioni~\ref{sec:z2-qa-tempo}
        e~\ref{sec:z2-benessere});
  \item il \textbf{brainstorming}: obiettivi, problemi tipici, un metodo
        completo (\emph{rapid ideation} e \emph{dot voting}) e la
        \textbf{matrice impatto--sforzo} per decidere
        (sezioni~\ref{sec:z2-brainstorming}, \ref{sec:z2-metodo}
        e~\ref{sec:z2-impatto});
  \item i consigli pratici, le domande sul brainstorming e il principio del
        \textbf{disagree and commit} (sezioni~\ref{sec:z2-qa-brainstorming}
        e~\ref{sec:z2-disagree});
  \item il \textbf{leading from within}: perché la leadership non è un
        titolo, le carriere divise tra manager ed esperti
        (sezione~\ref{sec:z2-leadership});
  \item il ``momento Piero Angela'' sulla biologia del gruppo e le domande
        che ne sono nate (sezioni~\ref{sec:z2-biologia}
        e~\ref{sec:z2-qa-biologia});
  \item la \textbf{piramide della fiducia}, il prendersi cura del gruppo e
        che cosa cambia quando si diventa manager
        (sezioni~\ref{sec:z2-piramide}, \ref{sec:z2-gruppo}
        e~\ref{sec:z2-manager});
  \item le domande finali e la chiusura del workshop
        (sezioni~\ref{sec:z2-qa-finale} e~\ref{sec:z2-chiusura}).
\end{enumerate}

Come per il capitolo precedente, una premessa di metodo: buona parte della
giornata è fatta di domande della classe e di racconti dell'esperienza del
relatore. Sono riportati con cura perché sono il modo in cui i principi
vengono resi concreti, e perché spesso contengono le indicazioni più
direttamente utilizzabili (soprattutto sui colloqui di lavoro).

% ---------------------------------------------------------------------
\section{Il programma della giornata}
\label{sec:z2-programma}

Come annunciato alla fine della prima giornata
(sezione~\ref{sec:z1-domani}), gli argomenti sono due:
\begin{itemize}
  \item un ``doppio topic'', \textbf{time management e brainstorming};
  \item il \textbf{leading from within}, che il relatore definisce ``la
        versione avanzata di ieri'': come guidare un team stando dentro al
        team, verso i colleghi e verso il gruppo con cui si lavora.
\end{itemize}
Poiché l'ordine è indifferente, il relatore lascia scegliere alla classe: dalle
poche preferenze espresse, la leadership va al pomeriggio, così che le domande
finali si concentrino su quella, mentre la mattina è dedicata a time
management e brainstorming. Il materiale c'è, ma il relatore avverte che la
giornata potrebbe finire prima delle quattro: non vuole trattenere nessuno se
non c'è niente di specifico da fare. Per l'eventuale tempo avanzato prima della
pausa pranzo propone un \textbf{argomento bonus}: i colloqui di lavoro,
selezione e \emph{recruitment}, visto che il relatore ha fatto e fa tuttora
selezione del personale e che per molti studenti sarà la cosa più immediata da
affrontare finito il percorso di studi (domande comuni, come rispondere,
eccetera).

Prima di cominciare, però, il relatore apre a \textbf{domande sulla giornata
precedente}: la classe ha avuto modo di ``metabolizzare'' i contenuti e
potrebbero essere emersi dubbi a posteriori. È proprio da una di queste domande
che nasce, in anticipo, la parte sui colloqui.

% ---------------------------------------------------------------------
\section{Colloqui di lavoro: capire che cosa cerca l'azienda}
\label{sec:z2-colloqui}

\subsection{L'annuncio come prima fonte di informazione}

La domanda: il giorno prima si era detto che è importante capire che cosa
l'azienda cerca davvero; ma se il primo colloquio è con le risorse umane, cioè
con qualcuno che magari ``di tecnico non sa'', come si fa a capirlo? Il
relatore risponde subito, nel dubbio che più tardi manchi il tempo.

La prima fonte è \textbf{l'annuncio}. Non viene scritto a caso: è approvato
sia dall'HR sia dall'\emph{hiring manager}\footnote{Nella trascrizione
``Airing Manager'' e ``Iron Manager''; allo stesso modo ``e-char'' e ``char''
stanno per HR.}, cioè la persona che seguirà il candidato una volta assunto,
e spesso dall'intero dipartimento HR. È calcolato e scritto in un certo modo
per far capire che cosa serve. Molti annunci moderni, oltre alla parte su
\emph{che cosa si verrà a fare}, contengono una parte su \emph{che cosa ci si
aspetta da questa persona}, cioè il comportamento e il carattere che servono
in quella posizione: ed è quella che dice molto sulla \textbf{cultura
aziendale}.

L'esempio è quello di due annunci per la stessa posizione (programmatore
senior) descritta in modo diverso:
\begin{enumerate}
  \item ``cerchiamo un programmatore esperto che possa prendersi cura dei
        nostri programmatori junior, aiutarli a crescere, sviluppare un
        ambiente di lavoro inclusivo e stimolante, in modo che i nostri
        progetti futuri siano all'avanguardia'';
  \item ``cerchiamo un programmatore esperto che sappia rispettare le
        deadline, seguire le indicazioni dei dati, sia indipendente nel suo
        lavoro e sappia relazionarsi con gli \emph{stakeholder} per capire e
        proporre nuove direzioni per il progetto''.
\end{enumerate}
Nel primo caso ci si aspetta un team di junior, probabilmente con un'età media
più bassa, con cui ci si relazionerà soprattutto; per un senior è l'occasione
di sviluppare le doti da \emph{people manager}, e al colloquio arriveranno
domande come ``hai mai lavorato con dei team? Con che età media? Ti sei mai
trovato a guidare un gruppo?''. Nel secondo caso il senior deve eseguire un
piano deciso o proporre qualcosa, e si relazionerà soprattutto \emph{verso
l'alto}; al colloquio arriverà piuttosto ``fammi questo test di
programmazione''. Sono, di fatto, due colloqui completamente diversi.

\subsection{Quanti colloqui e con chi}

Ormai è molto raro essere assunti con un solo colloquio: ne servono almeno
due, realisticamente \textbf{tre}, che è diventata la media. L'iter tipico:
\begin{enumerate}
  \item il primo con l'HR;
  \item il secondo con l'HR o con il futuro manager;
  \item il terzo sicuramente con il futuro manager, perché l'azienda vuole far
        conoscere al candidato la persona con cui lavorerà, e viceversa.
\end{enumerate}
Già le domande diverse dei vari interlocutori aiutano a capire che cosa
cercano. Ma c'è una parte fondamentale, che secondo il relatore molti
sottovalutano: il momento finale delle \textbf{``hai domande?''}.

\subsection{Le domande da fare al selezionatore}

Un principio di base: \textbf{i colloqui non sono ``devo farmi selezionare
dall'azienda'', sono ``dobbiamo selezionarci a vicenda''}. Le domande vanno
quindi preparate: ci sarà sempre qualcosa che, dalla \emph{job description} o
dal colloquio, non è chiaro al cento per cento. Nel dubbio, il relatore
suggerisce tre domande ``che funzionano sempre'':

\paragraph{``Se iniziassi domani, quale sarebbe il mio \emph{day to day}?''}
Ottiene due effetti. Il primo è pratico: descrive la giornata tipo e fa capire
se piacerebbe o no. Il secondo è più psicologico, ma il relatore sa che
funziona: costringe l'interlocutore a pensare al candidato \emph{già dentro
l'azienda}, cioè a che cosa potrebbe fare da subito e a quale impatto potrebbe
avere con le sue competenze.

\paragraph{``Se aveste una bacchetta magica, che cosa cambiereste subito
dell'azienda?''} Sposta la conversazione su di loro e va bene sia con l'HR sia
con l'hiring manager. Li obbliga a indicare una sfida interna che vorrebbero
risolvere, ed è sempre un'informazione utile: se entrambi rispondono che
cambierebbero il modo in cui si prendono le decisioni, viene da pensare che non
ci sia un percorso chiaro e trasparente da parte di chi sta in alto.

\paragraph{Lo scambio: ``vi do un motivo per non prendermi, voi datemene uno
per non venire''.} È più avanzata e il relatore non la usa sempre, ma la trova
particolarmente interessante. Permette al candidato di esprimere i propri dubbi
in modo costruttivo (``mi dite che lavorerei a stretto contatto con un team di
junior: un motivo per non prendermi è che non ho anni di esperienza nella
gestione di giovani sviluppatori, sarebbe la mia prima volta; però penso di
poterlo fare perché ho questo interesse e queste competenze''). Dall'altra
parte, il selezionatore dirà che cosa secondo lui sarebbe più difficile per il
candidato: si scopre così che cosa sta già pensando, e l'informazione torna
utile anche per i colloqui successivi.

\paragraph{E se il selezionatore si innervosisce?} ``Se un recruiter si
innervosisce per le domande del candidato, ha sbagliato lavoro.'' Al massimo il
relatore si innervosisce se un senior gli chiede che cosa faccia di preciso un
programmatore senior (``lo fai da dieci anni, lo saprai tu''); ma se un junior
chiede che cosa dovrebbe fare ogni giorno, lo si spiega volentieri. E in realtà
vale anche per i senior, perché \textbf{in ogni azienda lo stesso ruolo fa cose
diverse}: conviene sempre farsi spiegare che cosa intendono per quel titolo.
L'esempio viene dal mondo delle community: \emph{community manager} a Riot
indicava la persona a capo del team di \emph{community representative} di una
nazionalità (``community manager italiano''); a Blizzard indicava la persona
sopra l'intero dipartimento che gestiva quei team. Candidarsi per lo stesso
titolo non garantisce di aver capito per che cosa ci si sta candidando, e i
recruiter lo sanno.

Se poi un recruiter si innervosisce comunque, sta dando più risposte di quelle
che pensa: vuol dire che l'azienda non vuole domande, o non vuole interazioni
del tipo ``ma perché facciamo questa cosa?''. Che abbia formato i suoi
recruiter a innervosirsi, o che accetti che lo facciano, è informativo.

\subsection{Le domande che dicono molto sull'azienda}

Allo stesso modo, come detto il giorno prima, è informativo se un recruiter
chiede a una donna se vuole avere figli: ``ah, quindi a voi interessa sapere se
voglio avere figli: questo già mi dice che probabilmente non voglio lavorare
qui''.

Un esempio \emph{gender neutral}: la richiesta di vedere la \textbf{busta paga
precedente}. ``Quanto è la tua RAL? Mandami una busta paga.'' La risposta del
relatore è no: ``ti ho detto qual è la mia RAL; se vogliamo instaurare un
rapporto in cui non ti fidi di quello che dico, abbiamo un altro problema''. È
di nuovo \emph{default to trust} (sezione~\ref{sec:z1-default-trust}). Il
relatore è quasi sicuro che legalmente non si possa chiedere la busta paga di
un candidato\footnote{È l'opinione del relatore, espressa come tale (``sono
quasi sicuro''); non si tratta di un parere legale e non è stata verificata.}:
a quel punto diventa lui antagonista del recruiter, e visto che la
conversazione è finita in una ``zona grigia'', fa lui le domande (``come mai è
legale questa cosa che mi chiedi?''). Ha già deciso che probabilmente non
lavorerà lì, ma così magari evita che i candidati successivi debbano subire le
stesse richieste. Cita l'esempio perché sa che in Italia è abbastanza comune,
mentre all'estero non gli è mai successo, anche se avrebbe potuto ``sparare
qualsiasi cifra''.

Il motivo è che \textbf{non c'è alcun onere di dimostrare quanto si
guadagna}: chi spara troppo alto si autoesclude. Recruiter e manager non
arrivano al colloquio con carta bianca né con l'indicazione di risparmiare il
più possibile, ma con un \emph{range}: ``per questa posizione possiamo
spendere dai 30 ai 40 mila euro l'anno; paga questa persona a seconda della
sua seniority''. Se un candidato che guadagna 30 dice di prenderne 55 per
``provare a fare i soldi'' e il budget arriva al massimo a 40, il relatore non
chiederebbe nessuna prova: direbbe ``grazie, siamo in un range diverso, possiamo
salutarci''. La domanda sulla busta paga è quindi \emph{formulata male}, perché
è poco furba. Quella giusta è: \textbf{``hai già una RAL in mente per questa
posizione?''}. Si ottengono le stesse informazioni senza passare per la zona
grigia.

\subsection{La cifra minima mentale}

E se la domanda la fanno a un junior che non ha mai lavorato? Il consiglio vale
per tutti: \textbf{preparatevi sempre la vostra cifra minima mentale}, per
qualsiasi offerta di lavoro, prima ancora di fare il colloquio. Non serve
conoscere la media italiana per quella posizione o quanto paga quell'azienda:
serve sapere \emph{quanto si è disposti ad accettare senza andare in negativo},
tenendo conto dei propri costi (la benzina per andare in ufficio ogni giorno, un
eventuale trasferimento, e così via). Quella cifra va sempre avuta in testa.

\subsection{La scrematura iniziale e il portfolio}

Uno studente chiede come funzioni la risposta a una candidatura: arriva un
rifiuto diretto, oppure c'è una scrematura iniziale e poi i colloqui? E vale
per tutti i ruoli? Sì: c'è una \textbf{scrematura iniziale} fatta dall'HR o, in
alcune aziende, direttamente dal software dell'HR. Si scarta chi è chiaramente
fuori da ciò che si cerca:
\begin{itemize}
  \item si cerca un \emph{graphic designer} e arriva il CV di un programmatore
        (``faccio questi esempi perché purtroppo succede'');
  \item si cerca un \emph{animator} tecnico bravo con MotionBuilder, lo si
        scrive nell'annuncio, e il candidato risponde ``so fare tutto tranne
        MotionBuilder'': c'è un motivo tecnico per non poterlo prendere;
  \item problemi di seniority: si cerca un senior e risponde un junior con un
        anno di esperienza.
\end{itemize}
Se la figura non è già chiaramente fuori target, si fa il primo colloquio.

\paragraph{Gli anni di esperienza negli annunci junior.} Spesso gli annunci
junior chiedono come minimo tre anni di esperienza. \textbf{Se l'annuncio dice
``junior'', ignorate i requisiti di esperienza}: o sono lì per sbaglio, o sono lì
per scremare. Il secondo caso è più \emph{borderline}, ma il relatore l'ha
visto fare in alcune aziende americane: si mettono requisiti molto alti,
chiaramente irraggiungibili, per vedere chi ha la ``faccia tosta'' di mandare
comunque il CV (``non me ne frega se chiedete un \emph{bachelor's degree}, non ce
l'ho, mando lo stesso''). Così però si filtra una grossa fetta di popolazione e
si seleziona un certo tipo di carattere: si stanno cercando ``i Top Gun''. Non è
quindi necessario iniziare con uno stage: si può iniziare direttamente da
junior.

\paragraph{Chi guarda il portfolio.} Il relatore non ha mai avuto un portfolio,
ma per le figure che lo richiedono parte comunque dal CV (per escludere, ad
esempio, il designer che si candida da programmatore) e guarda il portfolio
solo a parità con altri candidati: da HR può aprirlo e dire ``bello'', ma non
può dare un giudizio tecnico. Il portfolio lo apre chi sarà l'hiring manager,
che è sempre molto più tecnico. L'ordine quindi è: l'HR guarda il CV e si salva
il portfolio; se il candidato passa la prima scrematura, CV e portfolio vanno
a un tecnico che conferma (``sì, sa fare queste cose'') o boccia (``il
portfolio indica chiaramente che non sarebbe in grado, digli pure che non
facciamo il colloquio'').

\subsection{Test pratici, da casa e con qualsiasi strumento}

\paragraph{Si può chiedere un test pratico al posto del portfolio?} Uno
studente chiede se, con un portfolio che non rispecchia le proprie capacità, si
possa chiedere direttamente un test. No: se l'azienda ha bisogno di verificare
una competenza ha già pronto un test e lo fa fare a prescindere. Il relatore,
per esempio, fa fare un test ai programmatori e lo fa correggere dal suo senior
developer; per gli artisti si chiedono animazioni di prova, chiaramente veloci.
Non chiede mai test troppo lunghi, perché non vuole far lavorare nessuno gratis:
non chiede ``fammi un pitch di progetto'', chiede uno spunto per capire come si
ragiona e come si lavora. Per il resto, \textbf{se il portfolio non è
indicativo, va aggiornato il portfolio}.

\paragraph{Test dal vivo o a casa?} Il relatore li fa fare sempre a casa. Sa
benissimo che un programmatore a casa lavorerà con due monitor, il codice da
una parte e Google o ChatGPT dall'altra. Ne ha discusso con i suoi senior
programmer: interessa che qualcuno sappia le cose a memoria, o che sappia
trovarle e fare quel lavoro a prescindere? A lui interessa che il lavoro
\textbf{venga fatto}: che il candidato abbia imparato il codice a memoria o che
in venti minuti trovi le informazioni, le metta insieme, le faccia elaborare a
ChatGPT, le corregga e consegni il codice che serve, non ha preferenze. Quando
si lavora le indicazioni saranno ``fai questa cosa il prima possibile'': se la
si sa fare anche con strumenti aggiuntivi, ben venga, ``anzi, spiegalo anche
agli altri''. Farla a memoria va benissimo, ``ma perché mettere l'\emph{hard
mode}?''. E comunque ci sono cose che non si possono delegare a uno strumento:
se arriva un codice che non funziona e va fatto il \emph{debug}, bisogna
saperlo fare e saper progettare i test. La figura più completa, per il
relatore, è anche quella che sa trovarsi le informazioni e gli aiuti.

Le prove dal vivo sono soprattutto domande di \emph{problem solving}: il
relatore vuole sentire come si ragiona e chiede espressamente di \textbf{pensare
ad alta voce}, per seguire il percorso mentale che porta alla conclusione. Sono
più domande impegnative che veri test. I test veri e propri si fanno con i
propri tempi e i propri strumenti: ``scegli il momento, ti mandiamo il test, hai
due ore per rimandarlo indietro''. Questo è l'accordo; in quelle due ore si può
fare quello che si vuole.

\subsection{Junior, mid e senior}

\paragraph{Perché ci sono così pochi annunci per figure mid?} Uno studente ha
notato su LinkedIn e sulle piattaforme di recruiting un divario notevole tra le
richieste di junior e quelle di mid. Il motivo è che per molte aziende,
Stormind inclusa, \textbf{la differenza tra junior e mid è sparita}: o si è
senior o non lo si è. Ci si è resi conto che la differenza tra un junior e un
mid è che il mid si è abituato a lavorare in quell'azienda e il junior non
ancora; ma un mid che arriva con tre anni di esperienza in un altro studio (ad
esempio Milestone\footnote{Nella trascrizione ``se ha fatto tre anni in
milestone'': si tratta verosimilmente di Milestone, studio di sviluppo
milanese.}) va comunque ``ripreso come junior'', perché si lavora in modo
diverso e deve reimparare. Un senior, invece, ha molte più competenze e le sa
applicare in un certo modo, e quella distinzione resta netta.

\paragraph{Ci si può candidare anche se il livello non coincide?} Un'altra
studentessa dava per scontato che, se un annuncio non specifica ``junior'', sia
per mid. Se non è specificato ``senior'', secondo il relatore si può mandare il
CV chiedendo di valutare. Certo, se è specificato junior e si è junior le
possibilità di essere contattati sono maggiori. Ma vale lo stesso discorso del
livello di inglese fatto il giorno prima: si cerca qualcuno che \emph{sappia
fare quelle cose}, non una certificazione. Se serve un programmatore che sappia
usare Unreal Engine~5 e sistemare la cosa X o Y, che sia mid da due anni o
junior da uno non è un problema.

\subsection{Gli strumenti di intelligenza artificiale}

Una domanda in chat riguarda l'uso quotidiano dell'intelligenza artificiale in
azienda. Il relatore sta facendo fare corsi di IA al suo team ``perché voglio che
la gente lavori meno'': vuole che si sfrutti qualsiasi strumento disponibile,
perché gli altri lo stanno facendo e non vuole restare indietro. Se un lavoro
si può fare in tre ore di meno, non c'è alcun vantaggio a farlo in modo lungo e
tortuoso. Il rischio esiste: se si fa fare tutto a ChatGPT e non ci si tiene al
passo con il proprio \emph{know-how}, quando ci sarà da risolvere qualcosa ci si
potrebbe ``impallare''. Ma sta al singolo restare aggiornato e saper fare il
proprio lavoro: \textbf{gli strumenti devono essere strumenti, non un sostituto
delle competenze}. ``Non è una partita di calcio, con regole su come si devono
fare le cose o un \emph{dress code}'': il lavoro va fatto nel modo più efficiente
possibile.

A uno studente a cui era saltata la connessione proprio in questo momento il
relatore risponde che la lezione è registrata e si può recuperare dal video.
Poi chiude la parentesi: ``la parte difficile era finire ieri in tempo, ce
l'abbiamo fatta''.

% ---------------------------------------------------------------------
\section{Time management: pianificare la giornata}
\label{sec:z2-pianificare}

\subsection{Time management o prioritizzazione?}

La domanda di partenza alla classe è: \emph{chi di voi, in questo momento, fa
time management su se stesso}, cioè si organizza la giornata all'inizio e
organizza il proprio tempo? Qualcuno risponde di sì, qualcuno ha uno
\emph{schedule} settimanale con orari da rispettare, qualcuno ``ci prova''
(``hai fatto bene a specificare \emph{ci provo}: magari non la rispetto''). Il
relatore lo chiede perché è uno di quegli argomenti in cui molti, dopo aver
visto la spiegazione, si accorgono di non fare davvero time management: fanno
piuttosto \textbf{prioritizzazione}, ad esempio una lista di cose da fare, che
è un po' diverso.

Nessuna delle tecniche è complicatissima, ma è utile sapersi muovere. Il
relatore colloca volutamente questo modulo \emph{dopo} quello sul feedback: saper
gestire il proprio tempo permette anche di dare feedback agli altri sulla loro
gestione del tempo.

\subsection{La \emph{must do list}}

Il punto di partenza è il \textbf{planning iniziale}: non improvvisare la
giornata. Improvvisando si è in balia di una grande limitazione umana: ci si
dimentica le cose, oppure ci si lascia trascinare dal task più urgente o
dall'ultimo arrivato. Esistono modi di dire come ``se arriva una mail o rispondi
subito o diventa un task per domani'', perché il ``rispondo dopo'' non arriva
mai; più avanti si vedranno modi più raffinati di gestire le zone grigie, ma la
pianificazione resta fondamentale.

La prima cosa da fare \emph{non} è la lista delle cose da fare oggi, ma la lista
delle cose che \textbf{assolutamente} vanno fatte oggi. Le due liste non devono
coincidere: la seconda richiede un livello di ragionamento in più. Sono le cose
per cui, se alle sei si stacca e non sono fatte, si resta finché non sono fatte,
perché sono imprescindibili.

\paragraph{Il primo momento per dire ``potrei non farcela''.} La \emph{must do
list} è anche il primo indicatore della giornata. Se è troppo corposa (dipende
naturalmente dalla complessità dei task), quello è il momento, ``alle 9:01 se
iniziate alle 9'', per dire al proprio manager \emph{oggi sarò in difficoltà}.
Alle 9:01 il manager ha l'intera giornata per aiutare; se invece si pecca di
ottimismo o di presunzione (``ce la posso fare'') e glielo si dice alle 17:59,
non può fare nulla. È lo stesso discorso del giorno prima
(sezione~\ref{sec:z1-triangolo}): anche facendo la to-do list si applicano i
concetti della cultura del feedback, ed è il motivo per cui questo modulo viene
dopo. Ci saranno vari momenti di questo tipo nella giornata, ma questo è il
primo.

\paragraph{E se ci si accorge a metà giornata di non farcela?} Uno studente
chiede che cosa fare se, convinto di farcela, si accorge a metà giornata di non
poter completare bene la \emph{must do list}. Il relatore lo concederebbe
\emph{una volta}: se si vogliono testare i propri limiti, va bene, ma
comunque si dice al capo ``oggi potrei essere tirato, ci provo'', e sarà lui a
valutare se il rischio è accettabile. Dipende dal rischio: se il \emph{must do}
è consegnare una milestone a un cliente, ``nel dubbio facciamolo fare a
qualcun altro, non possiamo correre quel rischio''; se il rischio è che domani
un animator riceva il codice con tre ore di ritardo, lo si può correre. ``È un
po' questo il lavoro.''

\paragraph{Cinque minuti e una stima.} Tutto questo va fatto in circa cinque
minuti. Se non è chiaro che cosa si deve fare in una giornata, \textbf{si
chiede}: ogni giorno dovrebbe essere chiaro. Poi si cerca di \textbf{stimare il
tempo} di ogni cosa. Si diventa più bravi col tempo: alla quinta volta che si
fa il report mensile si sa che ci vuole un'ora e ventidue minuti; la prima volta
si può chiedere a un collega quanto ci mette lui, sapendo che per noi potrebbe
essere diverso.

\subsection{To-do list e checklist}

Esistono molti strumenti per tenere to-do list e checklist; il relatore non ne
consiglia uno in particolare e cita come esempi Any.do, i promemoria di Windows
o Samsung Notes\footnote{Nella trascrizione ``AnyDo, Windows Memo''. Any.do è
un'app di liste e promemoria; ``Windows Memo'' indica probabilmente le note o i
promemoria integrati in Windows (Sticky Notes, Microsoft To Do).}: ``potete usare
davvero la qualunque''. L'importante è tenere la lista \textbf{visibile} da
qualche parte sul computer. Alcuni accorgimenti:
\begin{itemize}
  \item i \textbf{must do} vanno sempre \textbf{in cima} alla lista;
  \item quando durante la giornata arriva un task nuovo, la cosa da fare subito
        non è \emph{eseguirlo}, ma \textbf{aggiungerlo alla lista}: ``se arriva
        una mail, aggiungete subito alla lista \emph{rispondi alla mail}'',
        decidendo se è diventato un \emph{must do} o se è comunque un task per
        oggi;
  \item non limitarsi alle cose operative: il relatore ha iniziato ad
        aggiungere anche task come ``parla con Tizio'' o ``parla con Caio'',
        perché è da un po' che non chiede loro un aggiornamento; se non li
        mette in lista si accorge di non farli;
  \item spuntare le voci completate: per quanto minimo, rilascia
        \textbf{dopamina} (se ne parlerà nel pomeriggio,
        sezione~\ref{sec:z2-biologia}), aiuta a motivarsi e ad andare avanti.
        Sono piccole cose, ma si sommano.
\end{itemize}

\paragraph{Priorità contro recenza.} Che cosa succede se si esegue sempre
l'ultimo task arrivato? Lo chiede il relatore alla classe. Si interrompe il
flusso di lavoro e si cambia continuamente attività, certo; ma soprattutto si
stanno \textbf{mescolando priorità e recenza}: si sta dicendo ``non importa come
ho prioritizzato i task, la mia priorità è sempre l'ultimo arrivato'', che non
vuol dire fare le cose più importanti. È il motivo per cui tanti dipartimenti IT,
quando si chiede loro qualcosa, rispondono ``mi apri un ticket?''. È diventata
quasi una battuta, ma il concetto è chiarissimo: devono poter prioritizzare
quella richiesta rispetto alle altre 112 arrivate in giornata, e senza ticket
non possono confrontarla con le altre. Il relatore è totalmente a favore: ``certo
che ti apro un ticket, altrimenti tu come fai a tracciare il tuo lavoro?''.

\subsection{La checklist: l'esempio dell'aviazione}

La \textbf{checklist} non è la task list. La task list contiene le cose da fare;
la checklist contiene \textbf{i passi da seguire quando si esegue un task
specifico}. L'esempio è quello delle \textbf{checklist pre-volo}, introdotte
nell'aviazione (secondo il relatore tra gli anni Sessanta e Settanta): chi vola
spesso nota che i piloti seguono sempre un ordine simile nei controlli, tanto
che il relatore capisce quanto manca al decollo da quando controllano i flap, i
motori, quando accendono e spengono le luci. L'introduzione delle checklist
avrebbe ridotto di circa il 75\% gli incidenti in fase di decollo\footnote{Date
e percentuale sono quelle riferite in lezione e non sono state verificate.
L'origine della checklist pre-volo viene di solito fatta risalire al 1935, dopo
l'incidente del prototipo del Boeing B-17, causato da un controllo dimenticato da
un equipaggio molto esperto; la diffusione sistematica delle checklist e dei
protocolli di cabina è invece successiva.}, e soprattutto per i \textbf{piloti
più senior}. Perché? Perché chi ha fatto un decollo centomila volte non va a
controllare l'alettone destro: ha funzionato centocinquantamila volte, funzionerà
anche stavolta. Quando ci si abitua a fare qualcosa, si sbaglia proprio quando
qualcosa non va. La checklist ha obbligato a fare quei controlli ogni volta e ha
ridotto in modo drastico gli incidenti.

Il relatore usa checklist anche per cose banali, ad esempio le \textbf{email}:
essendo una persona abbastanza ansiosa su queste cose, se prepara una mail alla
leggera finisce per passare troppo tempo a ricontrollare l'oggetto, il
destinatario, l'allegato, e comunque si dimentica qualcosa. Seguire una checklist
sembra più complesso, ma \textbf{fa risparmiare tempo}: se la checklist è
completa, si sa di aver fatto tutto e si può andare avanti senza continuare a
pensare ``ho chiuso il gas? ho spento la luce?''. È come la lista della spesa: se
tutto è segnato e preso, punto. È un investimento: trenta secondi in più per
fare la cosa, invece di cinque minuti a sperare di averla fatta giusta.

\paragraph{E se ci si dimentica qualcosa nella checklist?} Succede. A quel
punto \emph{preparare la checklist} è un task in sé: un giorno in cui si ha
tempo si fa quel task pensando a quali passi si seguono e a come assicurarsi che
funzioni. Nessuna checklist è perfetta: ci saranno sempre compromessi tra ``devo
davvero controllarlo ogni volta?'' e ``posso farne a meno?''. Però, in generale,
\textbf{se qualcosa anche solo una volta ha dato problemi, va in checklist},
perché vuol dire che potrebbe darne di nuovo. Nessun manager rimprovererà
``come hai fatto a non prevedere che questa cosa poteva andare male?'': si
diventa esperti anche nelle checklist di un task facendo quel task, ed è un
processo iterativo. Se qualcosa è stato dimenticato e ha creato problemi, lo si
aggiunge per la volta dopo. Sbagliare una o due volte non è un problema;
l'importante è che la situazione migliori la volta successiva.

\paragraph{L'inizio di un junior e l'\emph{onboarding}.} A uno studente che
osserva che l'inizio di un junior è quindi inevitabilmente un po' travagliato,
il relatore risponde che è giusto così: ``se volessi il dipendente perfetto da
subito, cercherei un senior, non un junior''. Allo stesso modo non capisce le
aziende che prendono un junior e non lo fanno seguire da nessuno: che cosa ci si
aspetta, che assorba passivamente tutto quello che ha intorno e impari ``per
osmosi''? Può succedere, ma imparerà a fare le cose come vuole lui, perché
nessuno gli sta dicendo ``meglio farle così''. Moltissime aziende, oltre al
manager (che magari ha quindici \emph{report} e non ha il tempo materiale di
seguire un junior come andrebbe seguito), assegnano un \textbf{referente}, di
solito un senior, con il compito di prendersi cura della persona nuova e
insegnarle a fare le cose. Altrimenti l'\textbf{onboarding}, cioè il tempo
necessario per arrivare al cento per cento nel proprio lavoro, dura un anno.
Più risorse l'azienda dà all'inizio, prima la persona diventa operativa: è
nell'interesse dell'azienda dare il massimo supporto a chi arriva.

Il relatore diffida delle aziende che ``vogliono vedere se uno se la cava''. È
il discorso dei cento metri e della corsa a ostacoli\footnote{Il relatore lo
introduce come un discorso già fatto (``si torna al discorso: vuoi fare i cento
metri o la corsa a ostacoli?''), ma il passaggio non compare nella trascrizione
della prima giornata riportata nel capitolo~\ref{cap:z1-feedback}.}: non si può
far correre uno sui cento metri e l'altro con gli ostacoli e aspettarsi tempi
simili, perché chi ha gli ostacoli ci metterà per forza di più. Il supporto ai junior è un
investimento sensato.

\subsection{\emph{First contact resolution}}

Un concetto che viene dai call center: \textbf{quante richieste vengono risolte
al primo contatto}. È importante quando un task richiede iterazioni o
comunicazione con altri: bisogna organizzarsi come se si dovesse risolvere la
cosa \emph{in una sola passata}. Se viene chiesto un modello 3D, lo si fa giusto
la prima volta, al meglio delle proprie possibilità. Se si deve parlare con
qualcuno, si raccolgono tutte le richieste in una.

L'esempio è quello dei ticket che il relatore mandava ai colleghi del supporto di
Blizzard: ``ciao, ho un problema, non riesco a loggarmi'' --- ``mi dici lo
username?'' --- ``è questo'' --- ``mi dici l'email?'' --- ``è questa'' --- e così
via, con un paio d'ore tra un messaggio e l'altro perché il ticket deve fare il
suo giro. Tutto poteva essere riassunto in ``ciao, mi servono username, email e
questo e quest'altro''. Unificare le richieste su una singola cosa aiuta a
risolverla il prima possibile; bisogna avere chiaro \emph{che cosa serve prima
di fare una richiesta}, in modo che ce ne sia una sola. In molti sistemi di
supporto, se si apre un secondo ticket per aggiungere informazioni al primo, il
primo, anche se era quasi in lavorazione, viene rimandato in fondo alla coda:
serve a disincentivare chi manda continuamente ticket del tipo ``mi sono
ricordato questa cosa''. O si manda tutto subito, o ci vorrà tempo.

\subsection{Non iniziare senza le informazioni}

Corollario: \textbf{se non si ha tutto quello che serve per completare un task,
non lo si inizia}. Prima ci si procura quello che manca. Al relatore capita con
la pianificazione dei due anni di lavoro successivi: se non ha tutte le
informazioni, prima lavora per averle (``mi dai questa informazione, mi dai
quell'altra''), poi pianifica; altrimenti pianificherebbe qualcosa di cui non ha
completa visibilità.

\subsection{Il costo delle interruzioni}

Alternare continuamente l'attenzione tra i task e, soprattutto, reagire di
continuo agli input esterni (la mail, il messaggio, il collega che parla) fa
\textbf{sparire il tempo dalla giornata}. Se mentre si prepara un report in
Excel si viene interrotti (dal gatto, da una mail, da un collega, da una
telefonata) e poi si torna al foglio, servono in media \textbf{un quarto d'ora}
per tornare mentalmente al cento per cento; chi è particolarmente bravo ci mette
dieci minuti, ma si parla comunque di dieci-quindici minuti\footnote{Ordini di
grandezza come riportati in lezione. Le ricerche più citate sul tema (ad esempio
quelle di Gloria Mark sulle interruzioni nel lavoro d'ufficio) indicano tempi
medi di ripresa anche superiori, intorno ai venti minuti.}. Se si viene
interrotti di continuo, si passa più tempo a cercare di tornare al cento per
cento che a lavorare. È una partita persa in partenza.

Per questo il giorno prima si era parlato dei \textbf{semafori} sulla scrivania
(sezione~\ref{sec:z1-dare}): il semaforo può essere anche semplicemente il
proprio stato su Slack, su Discord o su qualunque strumento si usi con i
colleghi. Se si è occupati, si è occupati: è una propria valutazione che,
accettando di essere disturbati di continuo, quella cosa non si riuscirà a
farla, o almeno non in tempo.

% ---------------------------------------------------------------------
\section{Prioritizzare: la matrice di Eisenhower}
\label{sec:z2-eisenhower}

Ora che si hanno task list, checklist, \emph{must do} e \emph{nice to do}, resta
la domanda di prima: se arriva un task nuovo e si decide di non fare per prime le
cose nuove, \emph{come si capisce se è davvero urgente}? Serve un metodo, ed è la
\textbf{matrice di Eisenhower}\footnote{Nella trascrizione anche ``Jason
Howard''. Dal nome del presidente statunitense Dwight D.~Eisenhower, a cui si
attribuisce la distinzione tra urgente e importante.} (la classe prova a
indovinarne il nome proponendo altre matrici, come una matrice dei rischi, ma è
un'altra cosa). Ogni task che arriva va posizionato su questa tabella:

\begin{center}
  \begin{tabular}{@{}p{0.2\textwidth}p{0.36\textwidth}p{0.36\textwidth}@{}}
    \toprule
    & \textbf{Urgente} & \textbf{Non urgente} \\
    \midrule
    \textbf{Importante} &
    \textbf{Fallo} (\emph{do}): diventa un task di oggi. &
    \textbf{Pianificalo} (\emph{schedule}): va fatto, ma non subito. \\
    \addlinespace
    \textbf{Non importante} &
    \textbf{Delegalo} (\emph{delegate}): va fatto, ma non serve che lo faccia
    tu. &
    \textbf{Eliminalo} (\emph{delete}): non lo fai. \\
    \bottomrule
  \end{tabular}
\end{center}

\paragraph{Urgente e importante: lo fai.} Sono di solito task con una
\textbf{deadline molto chiara} (entro oggi, entro un'ora, entro X) e con
\textbf{conseguenze} se non vengono fatti: da qui l'``importante''.

\paragraph{Importante ma non urgente: lo pianifichi.} Va fatto, ma non subito:
o lo si mette tra le cose di oggi ``ma dopo'', o anche alla settimana successiva.
L'esempio: un collaboratore scrive ``il mio rinnovo è tra tre mesi, ne possiamo
parlare?''. È importante, perché il rinnovo di una persona è importante; non è
urgente, se ne può parlare anche domani, così da potersi preparare e fare altre
cose nel frattempo.

\paragraph{Urgente ma non importante: lo deleghi.} Qui ``non importante''
\textbf{non} vuol dire che il task non sia importante, ma che \textbf{non è
importante che lo faccia tu}, che è molto diverso. Se arriva ``mi serve questo
cambio a questa animazione 3D'' e ciò che conta è il cambio, non chi lo fa, si può
dire ``sono d'accordo che è urgente, ma non serve che lo faccia io'' e chiedere a
qualcun altro. Sono task che vanno fatti ma non richiedono la propria particolare
\emph{expertise}.

Il relatore insiste molto su questo punto con manager e senior: quando arriva un
task urgente, i senior tendono a dire ``lo faccio io, perché è urgente''. Ma se un
junior può farlo, lo si fa fare al junior e si investono le competenze più alte
del senior in qualcos'altro: il senior può fare il lavoro del junior, ma il
junior non può fare il lavoro del senior mentre il senior fa quello del junior.
Le persone non sono intercambiabili. Si usa il \textbf{minimo sforzo
necessario}: se una cosa la può fare un mid, la fa il mid, non il senior.

\paragraph{Né urgente né importante: lo elimini.} È il quadrante che quasi
nessuno usa: tutti fanno, delegano o pianificano, ma non pensano mai a
\emph{delete}. Invece è fondamentale. A volte è corretto dire: ``questa cosa non
dovrei farla io, non è urgente, non la faccio; sappi che non la faccio, e se c'è
un problema ne parli con il manager''. Non si sta delegando a qualcun altro: si
sta dicendo che quella cosa così com'è \textbf{non andrebbe fatta}, perché non
ne sono chiari l'importanza e il perché. Nella propria lista di oggi non ha
spazio, e non è abbastanza importante per essere fatta da sé nemmeno nei giorni
successivi.

Un esempio di \emph{delete} che non è un vero \emph{delete}: arriva una richiesta
di meeting di tre ore sul futuro del progetto, con invitati tutti gli art
director, il game director e così via. Il relatore spesso risponde ``io non
vengo'': per lui è \emph{delete}, ma il meeting si fa lo stesso. Uno studente
osserva che a volte si tratta di una richiesta mal formulata o poco chiara: è
esattamente così. ``Nella sua forma per me è no; se poi mi spieghi meglio perché
dovrei farlo io o perché è urgente, allora cambia.''

% ---------------------------------------------------------------------
\section{Proteggere il proprio tempo}
\label{sec:z2-proteggere}

\subsection{Evitare gli sprechi: un limite che nessuno può insegnare}

Evitare gli sprechi di tempo è difficile, soprattutto da remoto: è una sfida
costante con se stessi. Ha molto a che fare con il carattere, il
\emph{self-control} e (``passatemi anche questa'') la professionalità. Il
relatore può dare tutti i consigli del mondo, ma se dentro di sé non si percepisce
l'urgenza di fare il proprio lavoro quando lo si deve fare, non lo si può
insegnare. Se si è a casa e invece di lavorare si gioca a qualcosa su Steam,
nessuno deve spiegare che non si dovrebbe. Ci sono limiti oltre i quali non si
può andare.

Quello che si può fare è sapere che \textbf{ogni interruzione ha un costo
sommerso}: ogni messaggio, domanda, telefonata, ogni cambio di focus. Sono i
quindici minuti di prima, che possono arrivare a trenta a seconda del tipo di
task.

\paragraph{È più difficile tornare su un task logico o creativo?} Il relatore lo
chiede alla classe: fa più fatica un artista a tornare sul suo modello 3D o un
\emph{data analyst} a tornare sul suo foglio Excel? Le risposte si dividono
(``creativo'', ``entrambi''), ed è proprio il punto: dipende a chi lo si chiede. I
creativi rispondono ``creativo'', gli analisti ``analitico'', perché ognuno ha una
chiara comprensione di che cosa serve per essere davvero concentrati su ciò che
fa. Il creativo sa che per essere pienamente creativo gli servono l'ambiente, la
concentrazione, gli spunti, gli esempi da tenere chiari in testa; l'analista deve
avere in testa i conti, i fogli di dati e così via. Sono \textbf{ugualmente
difficili}, ma solo chi fa quella cosa ne apprezza fino in fondo la difficoltà. Si
può avere un assaggio dell'altra parte, ma la piena comprensione arriva solo
facendo: è come guardare uno sport in tv e pensare ``potrei farlo anch'io'',
finché non ci si trova lì.

\subsection{Distrazioni volontarie e involontarie}

Uno studente, che sta leggendo manuali di \emph{design thinking}, chiede come
conciliare tutto questo con l'idea che distrarsi, cioè non restare focalizzati,
aiuti a trovare un nuovo punto di vista per risolvere un problema. Qual è il
giusto equilibrio? Il relatore è \textbf{totalmente d'accordo} con quel concetto:
sta parlando di \textbf{interruzioni non volontarie}. Sa benissimo che dopo cinque
ore su un foglio Excel ``esplode la testa''; ma se si pianifica un quarto d'ora di
pausa e lo si rispetta, si sta investendo sulla propria creatività ed
efficienza. Non è un investimento dover rispondere a una chiamata ogni tre
minuti.

Anzi, il consiglio è \textbf{crearsi interruzioni pianificate}, sotto il proprio
controllo: fare pause, fare altro, soprattutto nel \emph{problem solving}. Il
relatore stesso, se non trova una soluzione, si ferma, fa un altro task e poi ci
torna; ma è una sua scelta. La discriminante è che la pausa sia \emph{scelta} e
non imposta: anche se si fanno le pause giuste, se si viene interrotti proprio
nel momento di massima concentrazione (raggiunta magari grazie alla pausa) si
butta via tutto. Uno studente cita la \textbf{tecnica del pomodoro}: è
esattamente quello, piccole pause pianificate che aiutano a essere creativi,
purché restino pianificate. Se i quindici minuti diventano ``vabbè, facciamo
mezz'ora'' o ``lo faccio nel pomeriggio'', non è più la tecnica del pomodoro:
diventano interruzioni, o un pretesto per non lavorare.

\subsection{Il peggior nemico siamo noi}

E chi sono i peggiori nemici in questo? ``Noi stessi'': siamo noi ad
autoconvincerci che venti minuti vanno bene quanto quindici. Le
\textbf{distrazioni autoinflitte} sono spesso le peggiori. Il relatore porta
l'esempio di una sua giornata, che ha annotato perché lo aveva fatto ridere
quanto tempo era riuscito a buttare via senza che nessuno glielo imponesse:
\begin{itemize}
  \item gli piace avere qualcosa in sottofondo mentre lavora (``ha a che fare con
        l'ADHD sicuramente''\footnote{Nella trascrizione ``IDACD''. Più avanti il
        relatore dice di avere l'ADHD ``come dicevo ieri''; il capitolo 0.1 non
        lo riporta, ma riferisce il suo racconto sull'``overthinking'' anche
        nell'ordinare una pizza (sezione~\ref{sec:z1-regalo}).}),
        e in genere lo aiuta. Il problema è quando il sottofondo è troppo
        interessante: la musica va bene, ma se nella playlist di Spotify parte un
        podcast di Alessandro Barbero sulla guerra civile americana, finisce per
        ascoltare quello più del lavoro;
  \item un'abitudine che ha poi perso: se faceva straordinari alle due di notte,
        da solo, per qualcosa di urgentissimo, apriva comunque Slack a
        controllare i messaggi, e la mail. ``Come se alle due di notte qualcuno
        mi scrivesse qualcosa di urgente'';
  \item la pausa con un po' di TikTok, dopo la quale si accorgeva di aver
        dimenticato su quale slide della presentazione stava lavorando.
\end{itemize}
La regola: \textbf{inserire i break è meglio che subirli passivamente}. Un conto
è dire ``faccio una pausa intenzionalmente'', un conto è ``mi sono distratto e ho
perso due ore di lavoro''. Chi ne è particolarmente soggetto non deve
demoralizzarsi: è anche un allenamento.

\subsection{Spezzettare i task grandi}

Un altro concetto utile: se si fa molta fatica a iniziare i task grossi, perché
sembrano infiniti (``più tardi inizio, più tardi inizia questa sofferenza''),
\textbf{vanno spezzettati}. Un task grosso non è quasi mai un task grosso: è
spesso un insieme di task piccoli. Si comincia a farne dei pezzi e si va avanti
piano piano, spuntandoli. Aiuta tantissimo anche nel \emph{problem solving}: un
problema gigante va scomposto in problemi più piccoli, da risolvere per primi, per
poi ricomporre il puzzle.

\subsection{Accettare di dire di no}

Tornando alla matrice di Eisenhower: bisogna \textbf{accettare di dover dire di
no}, perché dire sempre sì porta a problemi ben peggiori di quanto si pensi.

\paragraph{Il collega che diceva sempre sì.} A Riot il relatore aveva un collega
estremamente disponibile, amichevole e collaborativo: qualsiasi cosa gli si
chiedesse, rispondeva ``certo, la faccio io'' (``mi porti fuori il cane?'' ---
``certo''). Il problema è che era umano: a un certo punto si supera la quantità di
cose umanamente possibili in una giornata. In più tendeva a offrirsi volontario
per far vedere che era intraprendente e aveva iniziativa, e nei progetti che
coordinava il relatore doveva spesso rifiutare la sua candidatura perché era
sicuro che non ce l'avrebbe fatta. Il suo dire sempre sì, a un certo punto, lo ha
danneggiato: era quello che dice sempre che lo fa lui, ma non è detto che lo
faccia. ``Ce la fai?'' --- ``Sì, sì, assolutamente'' --- e poi ``scusa, non ce l'ho
fatta''. ``Ok, però ora hai messo nei guai me e il resto del team.'' Quella
disponibilità era più dannosa di un ``no, non posso''. È lo stesso discorso del
giorno prima: dire all'ultimo momento che non ce la si fa è peggio che dirlo a
inizio settimana.

Uno studente lo paragona alla \emph{sindrome dello yes man}. Non proprio: lo
\emph{yes man} dà sempre ragione all'altro perché è avverso al conflitto e non vuole
la discussione del ``spiegami perché sei in disaccordo''. Il collega, invece,
pensava davvero di poter fare molto più di quanto potesse, e rifiutava l'idea di
non riuscirci.

\paragraph{``Ci hai pensato bene?''} Vale per i progetti e per le deadline. Con i
junior soprattutto, quando a ``vuoi fare questa cosa?'' arriva un sì troppo
convinto, il relatore chiede sempre: ``ci hai pensato bene? Me lo stai dicendo
per farmi contento? Sappi che se mi dici di no non penso peggio di te; penso
peggio di te se mi dici di sì e poi non ce la fai. Vuoi rivalutare la risposta, e
magari dirmi \emph{non lo so} invece di un sì categorico?''.

\paragraph{Che cos'è una sconfitta.} Uno studente chiede se davanti alle
``sconfitte'' non si dovrebbe rimediare. Il relatore chiarisce che per lui la
sconfitta non è il task non fatto: quella è una sconfitta \emph{a breve termine},
si troverà un modo di farlo, è un danno ``ma chi se ne frega''. La vera sconfitta,
un danno \emph{a lungo termine}, è aver pensato di riuscirci e non esserci
riusciti: se non si risolve il motivo per cui si è detto sì quando era no, lo si
ripeterà altre cinque, dieci, quindici volte. Nella testa del collega di Riot c'era
sempre una scusante (``c'erano questi imprevisti che mi hanno bloccato''), come
quella con cui ci si giustifica le interruzioni. Il relatore gli rispondeva: ``se
ogni volta c'è un imprevisto, forse c'è un \emph{pattern}: il problema non è che
non riesci a fare le cose, è che non prevedi mai gli imprevisti''. Dopo sette
ritardi su dieci dovrebbe scattare l'ingranaggio mentale (``forse sbaglio io a
dire sempre di sì''), ma non è sempre facile farlo scattare, e a volte serve una
spinta. È una delle situazioni più frequenti. E la cosa più difficile è che,
quando glielo si fa notare, queste persone lo prendono come un ``non sei capace''.
Il messaggio è invece: \textbf{``sei bravo, ma sei umano anche tu''}. Non si può
fare un numero infinito di cose in un'ora: bisogna accettare i propri limiti ``dati
dalla carne'' e non pensare di essere un cyborg.

\paragraph{Il cliente non ha sempre ragione.} Uno studente osserva che vale anche
per i project manager e per il principio ``il cliente ha sempre ragione''. Anzi:
il relatore ha lavorato con molti clienti e a volte \textbf{il cliente non sa che
cosa vuole}, che è molto diverso dall'avere sempre ragione. Soprattutto con i
clienti che chiedevano pitch creativi, la sua squadra presentava sì un'opzione
con la loro richiesta, ma anche le proprie (``se vuoi raggiungere questo, sarebbe
meglio fare così''), e molto spesso il cliente sceglieva quella proposta dal team.
Spesso il cliente ha in mente \emph{una soluzione}, ma bisogna capire \emph{quale
obiettivo} ha in mente e poi dargli una soluzione: non vuole la soluzione, vuole il
\emph{goal}, e se c'è un modo migliore per arrivarci, forse è meglio un'altra
soluzione.

\subsection{Dire di no ai meeting}

La prima cosa a cui il relatore ha imparato a dire di no sono i \textbf{meeting}.
Ha passato anni in ambienti in cui il suo calendario era letteralmente solo
meeting dalla mattina alla sera: estenuante per lui e inutile per il team. Tra i
fattori c'è il \textbf{FOMO}: tanti manager vogliono essere in tutti i meeting per
non perdersi niente ed essere sempre informati su quello che fa il team. In realtà
vanno in \emph{burnout}, passano i pomeriggi nei meeting e finiscono per essere i
più distratti, quelli che non ascoltano e che fanno domande inutili solo per far
sentire che fanno domande, senza aiutare davvero.

Prima di accettare un invito il relatore si fa una serie di domande, \textbf{come
una checklist} (e le fa anche agli altri):
\begin{enumerate}
  \item \textbf{Sono davvero necessario?} Servo davvero io, o se mancassi non
        succederebbe niente? Ai meeting ``di cortesia'' non va: se pensa di
        essere necessario lo dice lui, e il fatto che si faccia un meeting
        importante e urgente senza di lui non danneggia il suo ego. ``Non
        aggiungermi solo perché pensi che serva che io venga aggiunto.''
  \item \textbf{Uno di noi potrebbe partecipare e informare il resto del team?}
        C'è già qualcuno del suo team? Se sì, non può andare solo lui? Per un
        meeting informativo (una nuova policy, una nuova procedura da spiegare a
        tutti i dipartimenti), andarci tutti significa investire tre ore a
        testa; se ci va uno solo e poi spiega agli altri, si risparmia tempo.
  \item \textbf{Il riassunto sarà reso disponibile?} Per il relatore è
        obbligatorio. Per lo stesso motivo per cui la lezione viene
        registrata: se ci si perde un pezzo lo si può rivedere, se qualcosa non è
        chiaro lo si può risentire. Ma c'è anche un motivo di \emph{office
        politics}: se in un meeting si chiede ``mi fai questa cosa per venerdì?''
        e la risposta è sì, non si vuole sentire venerdì ``intendevo venerdì
        prossimo'' o ``non ero sicuro''. Nelle note del meeting c'è scritto che hai
        detto sì. ``Lo so che è un po' paraculo, ma in realtà è così che si
        lavora'': bisogna dare più peso alla parola data, e quindi tenerne
        traccia. Finché tutto resta orale, niente rimane scritto e tutto è
        interpretabile. Oggi è anche facile: Zoom, con l'assistente IA, produce
        in automatico trascrizione e sottotitoli, senza bisogno di qualcuno
        addetto alle \emph{meeting minutes}. In ogni caso, se in un meeting viene
        presa una decisione, il relatore vuole che sia \textbf{ripetuta per
        iscritto}: non deve ipotizzare, deve poter andare a rileggere.
  \item \textbf{Il meeting aggiunge valore rispetto al documento?} Se arriva un
        invito con allegato un documento di venti pagine, il relatore chiede: ``il
        meeting consiste nel leggermi il documento? Perché altrimenti non
        vengo''. Se non c'è un valore aggiunto, se lo legge da solo e manda
        conferma di averlo letto. È lo stesso motivo per cui si sconsigliano le
        presentazioni in cui si leggono le slide; e infatti, osserva, in quel
        momento lui stesso stava leggendo una slide, ma aggiungendo informazioni.
        ``Se tutta la lezione fosse io che vi leggo le slide, non servirei io: vi
        manderei le slide.''
  \item \textbf{Per che cos'altro potrei usare questo tempo?} È forse la più
        importante. Anche ammesso che serva, che serva proprio lui e che il
        meeting sia importante, in quel momento dovrebbe fare qualcosa di più
        importante? Se sì, fa quella. Non bisogna avere paura di dire ``no, non
        vengo'' anche se si serve, oppure ``ho un'altra cosa importante,
        spostiamo il meeting''. Il relatore si è sforzato di essere più
        \textbf{assertivo}: ``capisco la richiesta, ma ho qualcosa di più urgente
        da fare in questo momento''.
\end{enumerate}

\paragraph{E le richieste a voce?} Uno studente chiede se una richiesta fatta a
voce da un senior (``devi fare questo'') vada accettata così o fatta mandare per
mail o per messaggio. Accettarla a voce ci può stare. Ma se la richiesta è ``devi
farla in questo modo particolare e preciso'', è meglio farsi mandare il modo
preciso \textbf{per iscritto}: non si deve lasciare a interpretazione \emph{come}
fare quella cosa, per non ricadere nel ``ah, ma io non la volevo così''. Oppure si
conferma: ``sì, la faccio; va bene se la faccio così, così e così?''. Se dice sì,
si è a posto.

\paragraph{Non dovrebbero pensarci gli organizzatori?} Dalla chat: queste domande
non dovrebbero farsele già gli organizzatori dei meeting? Sì, sarebbe meglio che
tutti lo facessero, e il relatore stesso se le fa quando invita qualcuno: ``è
davvero necessario Marco? È davvero necessario Andrea?'', e solo se sì li invita.
Ma lavora in aziende anche di cinquecento dipendenti, dove non è detto che tutti
conoscano o applichino questa logica; per questo la applica per sé, sperando che
la applichino tutti. Ha sempre chiarito con tutti che da lui non arriveranno mai
inviti di cortesia o dettati dalla gerarchia (``deve esserci anche il mio capo''):
``devi esserci tu; se vuoi portarti il capo, portatelo, ma non lo invito solo
perché qui si vuole sempre il capo presente: sono sicuro che avrà di meglio da
fare''.

\paragraph{E i task?} Dire di no vale naturalmente anche per i task: si applica la
matrice di Eisenhower (sezione~\ref{sec:z2-eisenhower}) e si decide se fare,
delegare, pianificare o cancellare. Non c'è molto da aggiungere.

\subsection{Il multitasking non esiste}

``Questo potrà urtare tanti ego'', ma \textbf{il multitasking non esiste}. Si può
essere più bravi o più portati di altri a fare più cose contemporaneamente; la
realtà è che, concentrandosi su più cose, \textbf{la somma non farà mai cento}. Se
si lavora su due cose non si lavora al 50\% su una e al 50\% sull'altra, ma al 40\%
e al 40\%, e il 20\% va perso in ``costo di manutenzione'', cioè nella capacità di
seguirle entrambe. Il consiglio è dedicarsi al cento per cento a un task, poi
passare al successivo, poi al successivo, invece di farne cinque insieme: si può
avere la percezione di andare più veloci, ma non è così.

Uno studente cita una lettura di neuroscienze secondo cui il cervello è in grado
di fare una sola operazione alla volta. Il relatore conferma: smettiamola di
parlarne come di un pregio o di una \emph{skill}. È un modo che alcuni cervelli
trovano per svagarsi, l'equivalente di accendere un podcast: fare due cose
insieme vuol dire cercare di distrarsi alternandole, e farle entrambe peggio di
quanto si farebbero una alla volta. A chi preferisce fare tante cose insieme
suggerisce di provare a farne una alla volta: le farà più velocemente. Anche il
semplice passaggio da una cosa all'altra non è istantaneo: ci sono sempre una
decina di secondi in cui il cervello si chiede ``che cosa stavo facendo?''; per
cinquanta cambi, già questo spiega la perdita di velocità. È una questione di
\textbf{bandwidth mentale}, un concetto che tornerà nel pomeriggio: anche lo
\emph{switch} occupa banda.

\subsection{La fine della giornata}

A fine giornata si \textbf{rimette a fuoco la giornata}: si riguardano la task list
e la \emph{must do list} e si conferma di aver fatto tutto. È letteralmente il
controllo prima di uscire di casa: chiavi, porta chiusa, cellulare, portafoglio.
Al relatore è capitato di dimenticare un \emph{must do} rimasto in lista per tutto
il giorno, ed entra in gioco quella che chiama la \textbf{sindrome del
tirannosauro}\footnote{Il riferimento è alla celebre scena di \emph{Jurassic Park},
in cui si dice che il tirannosauro vede la preda solo se si muove.}: se una cosa è
lì, non è detto che la si veda finché non si muove. Averla in lista tutto il giorno
non vuol dire averla percepita come da fare.

Può anche capitare di non farcela perché quel giorno sono arrivate così tante cose
urgenti e importanti da rendere inutili le stime: la giornata no capita e non c'è
nessun problema. Ma se era davvero un \emph{must do}, \textbf{va fatto comunque}:
è lì che scattano la richiesta di straordinario o il ``devo farlo comunque''.
Attenzione: se si chiude la giornata con un \emph{must do} ancora aperto e si
decide ``vabbè, lo faccio domani'', può starci, ma allora \textbf{non doveva essere
in quella lista}, perché era rimandabile.

Al relatore questo controllo serve anche ``per quieto vivere'': se chiude ogni
giornata con la \emph{must do list} completa, sa che magari non ha fatto tutto
quello che c'era nella to-do list, ma i \emph{must do} sì. ``Il mio, oggi, l'ho
fatto'': può stare tranquillo senza continuare a chiedersi se ha fatto quello che
doveva.

% ---------------------------------------------------------------------
\section{Domande sul time management}
\label{sec:z2-qa-tempo}

\paragraph{C'è un numero medio ``sano'' di \emph{must do} quotidiani?} No, perché un
conto è avere dieci \emph{must do} come ``rispondere a dieci mail'', un conto è
averne uno solo come ``inventarmi un nuovo gioco''. Dipende dalla dimensione: se ne
gestiscono di più se sono minuscoli, di meno se sono giganteschi. Ci sono lavori
che per definizione hanno tantissimi \emph{must do} quotidiani e lavori che ne hanno
uno al giorno. Si imparerà da sé qual è il numero giusto.

\paragraph{C'è un modo per stimare quanto si riesce a fare in una giornata?} ``Lo so
che non è una bella risposta'', ma l'unico modo è la propria \textbf{storia}:
sapere quanto si riesce a fare nel quotidiano. È normale e fisiologico che le
prime stime siano sbagliate, e probabilmente eccessivamente ottimistiche;
l'importante è rendersene conto. È letteralmente la \emph{seniority}: abituarsi a
fare qualcosa insegna quanto ci si mette. Va benissimo metterci di più all'inizio,
ed è normale che il tempo cambi man mano che si migliora: ciò che oggi si stima in
un'ora, tra sei mesi lo si stimerà in mezz'ora. Chi ha lavorato con Scrum lo sa:
su un progetto nuovo tutti sbagliano la dimensione delle \emph{stories} all'inizio,
perché nessuno sa quanto ci vorrà, e al terzo \emph{sprint} consecutivo è più
chiaro. Le stime sono stime, altrimenti si chiamerebbero ``inserire quanto tempo ci
metterai''.

Se però si dà una stima e poi il lavoro si rivela il triplo, il problema non è la
stima: è che \textbf{bisognava chiedere più informazioni} prima di darla. La
risposta giusta non è la stima, ma ``per darti una stima accurata mi serve sapere
questo, questo e questo''. Il problema non è sbagliare la stima, è non darsi i mezzi
per azzeccarla.

\paragraph{Come si fa a far entrare questi metodi nella routine? È solo
determinazione?} È un misto. Il relatore risponde da persona con ADHD, con la
tendenza a distrarsi e a fare un mucchio di \emph{overthinking} (``come dicevo ieri,
anche ordinare la pizza è overthinking''). È un po' disciplina, ma soprattutto far
diventare qualcosa un'\textbf{abitudine}: lo sforzo per fare le checklist era molto
più alto un anno prima, oggi è un'abitudine. Di solito servono \textbf{circa tre
mesi} perché qualcosa diventi un'abitudine\footnote{Stima come riportata in lezione.
Gli studi più citati sulla formazione delle abitudini indicano tempi molto variabili
da persona a persona, con una mediana intorno ai due mesi.}: la teoria dice che chi
inizia oggi ad andare in palestra ogni giorno, dopo tre mesi non penserà più ``oggi
devo andare in palestra'', ma la palestra farà parte della sua giornata. Non è più
uno sforzo, ma ``è così che lavoro''.

\paragraph{Chi fa davvero time management?} Prima di passare al brainstorming il
relatore torna alla domanda iniziale: qualcuno, ora che sa che cosa intendeva,
cambierebbe risposta? Uno studente risponde che lo fa, ma senza checklist, perché
non ne sente il bisogno (non gli capita spesso di dimenticare le cose); la
\emph{must do list} invece la segue sempre, e le cose meno importanti capitano nelle
giornate in cui ne salta qualcuna. Va benissimo: non esiste un unico modo, ognuno
troverà il proprio. Il relatore porta il suo esempio perché sa che è servito a molte
persone con cui ha lavorato; con chi faceva fatica ad applicare un metodo (come
chiedeva prima uno studente) ne hanno trovati altri, perché ciò che funziona per uno
non è detto che funzioni per un altro. L'importante è che \textbf{ognuno trovi il
proprio metodo}, perché la soluzione non è non fare time management: ci si sentirà
più ``liberi'' improvvisando la giornata, ma non la si sta sfruttando a dovere.

E non si parla dei video sul \emph{mindset} in cui ci si alza alle cinque, si fa la
checklist e si fa colazione leggendo le mail: organizzarsi meglio la giornata
\textbf{dà più tempo libero}. E per il relatore il tempo libero delle persone è
fondamentale, perché è ciò che permette loro di lavorare meglio.

% ---------------------------------------------------------------------
\section{Una parentesi: tempo libero, ferie e benessere}
\label{sec:z2-benessere}

\subsection{Le tre condizioni del relatore}

Il relatore è un sostenitore della \textbf{settimana lavorativa di quattro giorni},
che ha applicato in varie aziende; dove non ci è riuscito, ha introdotto almeno
\textbf{venerdì liberi}, due volte al mese o una volta al mese, cioè giornate in
cui non si fa niente. Come le pause brevi sono importanti per essere creativi,
anche le pause lunghe, cioè giornate intere, lo sono per essere creativi ed
efficienti. È lo stesso motivo per cui, dopo la Games Week, dava una settimana
libera alle persone (sezione~\ref{sec:z1-fasi-team}): ha bisogno che siano
riposate, e quella settimana ``persa'' è in realtà un investimento per le
settimane successive. In chat qualcuno lo paragona al ``metodo finlandese'': è un
misto tra quello e il metodo del buonsenso.

Il relatore viene spesso chiamato a \textbf{cambiare le culture aziendali}, ed è un
po' il suo focus. Come un allenatore che accetta una squadra solo se può portare il
suo staff (il preparatore, il tecnico), lui chiede tre elementi:
\begin{enumerate}
  \item possibilmente \textbf{quattro giorni lavorativi}, o altrimenti giornate
        libere per dare spazio alle persone;
  \item \textbf{ferie illimitate}: nessun limite massimo, ma un \textbf{limite
        minimo garantito}. Ferie illimitate non devono voler dire poter scegliere
        di non prenderle: ``devi prendere almeno X giorni, fino a infinito'';
  \item la \textbf{feedback culture}.
\end{enumerate}
Senza questi tre elementi non garantisce il risultato: è un modo di gestire le
aspettative fin dall'inizio, come si diceva il giorno prima.

\paragraph{La feedback culture.} Il perché è stato spiegato nella prima giornata, ma
il relatore lo sottolinea: un ambiente in cui ci si può fidare del team vale
moltissimo. Quando va a fare lezione allo IULM, anche in settimana, non deve
preoccuparsi di che cosa stia facendo il team mentre lui non lavora (``oggi è
domenica e non ho questa preoccupazione a prescindere, ma anche se fosse mercoledì
non l'avrei''): sa che stanno facendo del loro meglio, e non gli interessa sapere
altro.

\paragraph{Ferie e permessi illimitati.} Nell'esperienza del relatore, sia a Riot sia
a Blizzard, avere ferie illimitate significava spesso \emph{prenderne meno} di quelle
obbligatorie: c'è una vocina che dice ``ti serve davvero un giorno di ferie oggi?
Tanto sono illimitate, te le prenderai quando ti serviranno'', e sapere di averne
illimitate spingeva a lavorare di più. Per lui era altrettanto dannoso, e per questo
ha introdotto un \textbf{minimo obbligatorio di 24 giorni} di ferie.

C'è stato però un effetto non previsto, e molto desiderato: in tutte le aziende in
cui ha introdotto ferie e permessi illimitati, \textbf{i giorni di malattia sono
calati del 70\%}\footnote{Dato riferito dal relatore sulla base della propria
esperienza, non verificabile.}. La gente si ammala di meno? O piuttosto il giorno in
cui ci si sveglia e si pensa ``oggi non ne ho voglia, mi do malato'' diventa il
giorno in cui si prende un permesso? Si smette di usare come malattia un giorno che
non era di malattia.

\paragraph{L'\emph{Hangover Day}.} Sia a Riot sia a PG Esports il relatore aveva
introdotto quello che a Riot, ``essendo Riot a Dublino'', si chiamava
\textbf{Hangover Day}\footnote{Nella trascrizione ``Angover Day''.}: una giornata
all'anno in cui si aveva il diritto di telefonare alle nove del mattino e dire ``oggi
non vengo perché non ne ho voglia'' (tradotto: perché ieri sera ho bevuto troppo).
Altrimenti quella giornata sarebbe diventata un giorno di malattia o di ferie. Il
fatto è che \textbf{quasi nessuno lo usava}: il relatore stesso a volte ci ha
pensato, ma poi la mattina si chiedeva ``voglio davvero giocarmelo oggi? Tutto
sommato potrei anche andare'', e ci andava. Ma anche solo sapere di averlo lo faceva
stare tranquillo. Allo stesso modo, con le ferie illimitate, chi voleva prendersi un
mese lo faceva (``buon per loro''), chi non lo faceva andava bene lo stesso:
l'importante era che staccassero. C'era un \textbf{benessere più diffuso} nel team,
perché sapevano che, volendo, avrebbero potuto.

\paragraph{Il \emph{Neko Day}.} Il primo punto, i giorni liberi, era chiamato
\textbf{Neko Day}, da \emph{neko}, ``gatto'' in giapponese: il giorno in cui si può
fare il gatto, rannicchiati davanti al camino a non fare niente. Perché
introdurlo? Chi lavora dal lunedì al venerdì passa di solito il sabato o la domenica
a fare le commissioni che non può fare nei giorni lavorativi: la spesa, le bollette,
il bollo dell'auto. In questo modo sfrutta meno il weekend. Con un venerdì libero al
mese (o due, o tre) le persone facevano quel giorno le commissioni del weekend, e
sabato e domenica si riposavano davvero (``anche non fare niente andava bene,
chiaramente''). Il lunedì arrivavano molto più cariche e riposate. Questo ha aiutato
tantissimo a mantenere un team sano, fisicamente e mentalmente.

\paragraph{Il supporto psicologico.} A PG Esports il relatore ha introdotto anche il
\textbf{supporto psicoterapeutico offerto dall'azienda}: sei sedute da un'ora con
uno psicoterapeuta, online. Era il periodo in cui si parlava del bonus psicologo
statale, ``che poi non è mai arrivato''\footnote{Così in lezione. In realtà il bonus
psicologo è stato istituito in Italia nel 2022, ma con fondi limitati rispetto alle
richieste e con forti ritardi nell'erogazione; il relatore si riferisce
probabilmente a questo.}, e lui ha pensato: l'idea è buona, perché non introdurla a
livello aziendale? Erano gli anni dei lockdown: l'azienda stava perdendo molte persone
che andavano in depressione o non avevano più voglia di lavorare, perché avevano
difficoltà a vivere quel periodo (``io stesso non è che l'abbia vissuto benissimo'').
Qualcuno ne ha usufruito e gli ha detto: ``mi stavo per dimettere, poi ho risolto le
cose e ho preferito restare''.

Secondo il relatore oggi c'è una consapevolezza maggiore da parte delle aziende
sulla \textbf{salute mentale}: vuol dire dare più ferie, dare questo tipo di
supporto, dare più \emph{benefit} (anche i buoni pasto). Il benessere non è solo
economico, è anche mentale, ed è una cosa positiva da sapere.

% ---------------------------------------------------------------------
\section{Brainstorming: obiettivo e problemi tipici}
\label{sec:z2-brainstorming}

\subsection{Perché il brainstorming riguarda tutti}

Secondo il relatore questa è la parte che servirà di più, sia nei gruppi e nei
progetti del percorso di studi (il gioco da tavolo, il gioco digitale) sia nelle
aziende. Rispetto al time management ha una pertinenza più immediata: tutti
hanno esempi di \textbf{brainstorming andati male}, da cui si è usciti esausti, con
la sensazione di non aver ottenuto niente, o di non aver parlato abbastanza (``ha
parlato solo quello''), o di non essere sicuri della scelta. Il brainstorming lascia
una sensazione molto chiara, a fine meeting, che qualcosa non torna; il time
management dà piuttosto la tristezza del ``non ce l'ho fatta a fare quello che
dovevo oggi, come mai?''.

\subsection{Generare idee, non decidere}

Il primo chiarimento riguarda \textbf{l'obiettivo}, che il relatore vede spesso
travisato. L'obiettivo del brainstorming è \textbf{generare idee}; invece diventa
spesso decisionale, con lo scopo di scegliere che cosa fare. Ma allora non sarebbe
un brainstorming, sarebbe un \emph{decision making meeting}. Generare molte idee e
scegliere la migliore nello stesso meeting non funziona né per generare idee né per
scegliere la migliore, perché sono \textbf{due metodi che si contrappongono}. Il
brainstorming va fatto quando si devono generare idee \emph{per poi} scegliere che
cosa fare: deve venire \textbf{prima} del \emph{decision making}, come primo passo
della fase di esplorazione di un nuovo progetto.

\subsection{``Non esistono cattive idee''?}

Bisogna evitare che i membri del team \textbf{autofiltrino} le proprie idee per paura
di essere rifiutati o presi alla leggera. Da qui nasce la frase ``non esistono
cattive idee, buttate fuori tutto'', che è vera solo fino a un certo punto. Le
cattive idee esistono, oggettivamente; quello che non si deve fare è
\emph{colpevolizzarle}. Né si deve per forza perseguire ogni opzione: parlare di
opzioni che sono effettivamente brutte è uno spreco di tempo, e il brainstorming ha
anche un problema di gestione del tempo. Il relatore vede spesso brainstorming
gonfiati a ore e ore in cui si parla anche di cose che non meritano nemmeno dieci
minuti, come i feedback buoni ma inutili del giorno prima. \textbf{Trattare tutte le
idee allo stesso modo è sbagliato}, perché non hanno la stessa qualità. Bisogna
trovare un modo di valutarle senza che chi le ha proposte si senta preso di mira,
bullizzato o sminuito.

Per questo a ogni brainstorming dovrebbe seguire un \textbf{follow-up}, il momento in
cui si decide che cosa fare: le due parti vanno separate, altrimenti ``si crea solo
casino''. Il vero lavoro arriva quando le idee vengono messe a terra: c'è la parte di
brainstorming e di \emph{theorycrafting}, e la parte esecutiva viene dopo.

\subsection{I problemi più comuni}

\begin{description}
  \item[Le prime idee monopolizzano la sessione] Le prime due idee ``sparate a
        caso'' sono quelle su cui ci si concentra per la maggior parte del tempo,
        sviscerandone i dettagli (``intendevi così o così?''), per poi chiudere con
        ``altre idee? Questa, questa e questa? Sì, no, ok, andiamo avanti''. Si
        parla troppo delle prime idee e poco delle ultime.
  \item[Poche persone parlano per la maggior parte del tempo] È un effetto
        collegato al precedente.
  \item[\emph{Anchoring}] Ci si lega alle prime idee, e spesso si votano o si
        scelgono non perché siano le migliori, ma perché sembra di conoscerle meglio
        di quelle menzionate in trenta secondi alla fine. Non c'è nessun motivo per
        pensare che un'idea detta alla fine sia peggiore: era solo arrivata dopo.
        C'è un problema di \emph{timing}. Anzi, un'idea venuta istantaneamente è
        forse quella meno ragionata, perché ci si è pensato meno tempo.
  \item[Silenzi imbarazzanti] Spesso le persone arrivano al meeting senza nemmeno
        sapere che è un brainstorming. Quando il relatore convoca un brainstorming,
        comunica prima a tutti il \emph{topic}, il tipo di idee cercate, i paletti e
        le restrizioni, così che anche chi è meno portato all'improvvisazione arrivi
        con qualche idea pronta. Se la gente arriva impreparata rischia di non dire
        niente per un'ora, e allora perché è stata invitata? Non è un difetto loro:
        tanti manager dicevano al relatore ``ci sono persone inadatte ai
        brainstorming perché non sanno improvvisare''. No: si sta facendo
        brainstorming con un metodo che premia solo chi sa improvvisare ed è
        abbastanza estroverso da parlare sul momento. Non significa che chi non
        parla non avrebbe potuto contribuire, ma che non gli è stato dato modo di
        farlo.
  \item[Estroversi e pensatori rapidi dominano] Le personalità estroverse e chi
        pensa velocemente tendono a dominare la conversazione, più o meno
        intenzionalmente, senza lasciare agli altri il tempo di contribuire. Non c'è
        uno spazio in cui un introverso possa dire ``veramente io volevo\dots'':
        ``no, no, stiamo parlando di questa bellissima idea, facciamo questa''.
\end{description}

% ---------------------------------------------------------------------
\section{Un metodo di brainstorming}
\label{sec:z2-metodo}

Il relatore propone una soluzione principale, declinata in vari modi, che tocca
tutti questi punti.

\subsection{Preparazione e ambiente inclusivo}

\paragraph{Arrivare con un obiettivo chiaro.} Bisogna arrivare al brainstorming con
chiaro e preciso che cosa si sta cercando, cioè \textbf{che cosa si sta cercando di
risolvere}, \textbf{perché} e soprattutto \textbf{con quali vincoli}. Si stanno solo
raccogliendo idee per un cliente? Si cercano idee fattibili, o idee innovative? Che
cosa si vuole produrre da questo meeting? Deve essere chiaro a tutti \emph{prima}
del meeting, non spiegato durante, perché altrimenti si ricade nel problema delle
persone non pronte a intervenire.

\paragraph{Non giudicare le idee degli altri.} Al relatore è capitato spessissimo di
sentire un'idea terribile da una persona e, mezz'ora dopo, l'idea perfetta dalla
stessa persona. Ma se all'idea terribile tutti rispondono ``che stupidaggine, per
favore stai zitto'', quella persona non dirà la successiva, perché non vuole
rivivere l'esperienza. Giudicando le idee degli altri, di fatto, li si zittisce.
Le idee migliori nascono quando si propone quella che si pensa sia l'idea perfetta,
magari non è buona e non viene scelta (``fa niente''), si ascoltano anche le altre e
poi la si \textbf{riformula} integrando i pareri sentiti nel frattempo. Sono queste
\textbf{idee combinate} a essere spesso le soluzioni migliori.

\subsection{Separare generazione e discussione: \emph{brainwriting} e
\emph{rapid ideation}}

Con il classico ``non ci sono cattive idee, sparate'', succede spesso che uno prenda
la parola, spieghi la sua idea e il suo concetto, e gli altri comincino subito a
fargli domande (``ma quindi la vedi così?''): si sta unendo la persona che
\emph{genera} l'idea con la persona che la \emph{discute}, nello stesso momento.
Invece bisogna \textbf{dividere la generazione dalla discussione}: un momento in cui
si generano idee e un momento in cui si discutono. È il principio del
\textbf{brainwriting}. Funziona molto quando nel team c'è una figura con una
posizione gerarchica più alta: se quella persona esprime un'idea e i due momenti
sono uniti, tutti tenderanno ad allinearsi. E si vuole che \emph{tutti} possano
avere idee.

Lo strumento è la \textbf{rapid ideation}. Il relatore di solito dice: ``per i primi
cinque minuti non parliamo; scrivete su dei post-it il pitch delle vostre idee''. Se
ne possono scrivere quanti se ne vuole, dieci o uno, e si mettono da parte. Ma il
tempo è di \textbf{cinque minuti}: non dieci, non due, e soprattutto non mezz'ora.
Il \textbf{limite di tempo breve} è importantissimo, perché costringe,
consapevolmente o meno, a filtrare le idee mentalmente e a scrivere solo quelle che
si ritengono migliori: magari non è chiaro perché si preferisca un'idea all'altra,
ma il cervello sta facendo il lavoro in background e fa venire in mente le cose più
importanti o più giuste. È ottimo anche per i team che tendono a distrarsi, perché
sono solo cinque minuti.

\subsection{\emph{Dot voting}}

Poi si passa al \textbf{dot voting}:
\begin{enumerate}
  \item si attaccano \textbf{tutte} le idee su una lavagna fisica, un muro o una
        lavagna digitale (come Miro\footnote{Nella trascrizione ``mirror board'':
        si tratta della lavagna digitale Miro, già citata nella prima
        giornata.});
  \item si \textbf{uniscono le idee che hanno lo stesso concetto di base}. Se si sta
        pensando a un evento dal vivo e la proposta del relatore è un parco di
        dinosauri, mentre quella di un collega è un parco di dinosauri ``del
        paleolitico'' con i tirannosauri, le si raggruppa sotto ``parco di
        dinosauri'': non sono due idee separate. Un altro gruppo potrebbe essere,
        ad esempio, un parco a tema Avengers;
  \item ognuno ha \textbf{da tre a cinque voti} e va alla lavagna a mettere i propri
        pallini sulle idee che gli piacciono e di cui vorrebbe sapere di più.
\end{enumerate}
Il risultato tipico: un'idea con molti voti, una con un paio, altre con un numero
discreto.

\paragraph{Che cosa si vede subito?} Il relatore lo chiede alla classe. Si vede una
\textbf{preferenza} per un tema o per un'idea. Non c'è \textbf{polarizzazione} su una
o due idee iniziali, perché c'è varietà. Il voto è \textbf{anonimo}: non ci sono nomi.
E soprattutto manca un'altra cosa importantissima: \textbf{il giudizio negativo
espresso}. Se c'è un'idea con sette voti, una con cinque e una con tre, è
ragionevole parlare prima di quelle, perché tutti hanno espresso interesse e vogliono
saperne di più; ma nessuno sta dicendo ``non voglio assolutamente parlare di questa
perché mi fa schifo''. È semplicemente oggettivo che quell'idea non ha ricevuto voti.

\paragraph{Autovotarsi.} Il sistema permette anche di votare la propria idea, e non c'è
nessun problema: se la propria idea ha un pallino e gli altri ne mettono sette su
un'altra, il proprio voto è ininfluente. Il relatore vedeva spesso che all'inizio la
gente tendeva ad autovotarsi, poi capiva che era inutile, perché non sarebbe mai stato
il voto decisivo, e tanto valeva votare le idee che interessavano davvero. La
tendenza spariva con il tempo.

\paragraph{Che cosa si è ottenuto.} Tre ore di meeting sono state filtrate in una
discussione concentrata solo sulle idee che tutti ritengono interessanti, con un
\emph{agreement} di fondo: non si parla dell'idea numero uno perché è uscita per prima,
non è una questione di \emph{timing} ma di qualità. Il relatore ha vissuto tanti
meeting con la frustrazione interna del ``stiamo sprecando tempo su questa idea,
quando vorrei dire la mia che secondo me è meglio''. Qui non si può sprecare tempo,
perché si parla di qualcosa che tutti vogliono sentire.

\paragraph{E se si formano sottogruppi che si votano a vicenda?} Uno studente chiede se
non possano esistere ``cordate''. Anche se esistessero, significherebbe che il gruppo di
quel meeting vuole parlare di quella cosa, e va bene: l'importante è non sprecare tempo
su idee che non interessano a nessuno. Se si investono due ore, quarantacinque minuti
su ciascuna di tre idee vanno benissimo; dieci minuti per ogni idea no. E comunque il
voto è anonimo: non si può sapere che un'idea è di Gianni per dare supporto a Gianni,
quindi un eventuale sostegno reciproco sarebbe involontario.

\paragraph{E chi è lento a elaborare?} Uno studente osserva che il problema degli
introversi potrebbe ripresentarsi con i post-it: con un tempo così ristretto, chi è più
lento potrebbe andare ancora più in panico. Per questo bisogna \textbf{arrivare
preparati}: si conosce il topic in anticipo, quindi almeno un paio di idee già
formalizzate si hanno in mente. Magari non se ne scriveranno quindici, ma un paio sì.

\subsection{Quanti voti e quali varianti}

Dare più pallini a testa sarebbe controproducente? Il relatore indicava da tre a
cinque voti. Chiede alla classe quanti siano i gruppi di lavoro: dai quattro agli otto,
alcuni gruppi del secondo anno sono nove. La sua prassi:
\begin{itemize}
  \item \textbf{dai cinque partecipanti in su}, \textbf{tre voti} a testa, più che
        sufficienti per popolare molto la lavagna;
  \item \textbf{fino a cinque partecipanti}, \textbf{cinque voti} a testa, che
        permettono anche di variare un po'.
\end{itemize}
Esistono varianti: ad esempio tre voti di colori diversi, il blu per l'idea più
innovativa, il verde per quella più bella, il giallo per quella più facile, così da
qualificare \emph{perché} un'idea piace. Se l'obiettivo è proporre a un cliente tre
pitch (uno facile, uno medio, uno avanzato), la votazione è già fatta. All'inizio,
però, meglio partire da tre o cinque voti generici; con il tempo si imparerà ad
accelerare.

\subsection{Il moderatore e gli invitati}

\paragraph{Serve un moderatore?} Può essere utile. Il relatore fa spessissimo il
moderatore: partecipa ai brainstorming senza fare altro che moderare, ``perché so che
le persone con le idee fighe sono loro e non io''. Ci si può anche autogestire. Nei team
piccoli non ha senso ``perdere'' una persona per la moderazione; nei team più grandi
(dai cinque in su, sette-otto persone), se qualcuno dice ``non penso di avere idee
particolari su questo, vi modero così vi aiuto'', è un ottimo modo di contribuire.

\paragraph{Chi invitare?} Dipende da che cosa si cerca. Per inventare un nuovo gioco il
relatore invita i game director, gli art director e magari il \emph{publishing manager}
per un input commerciale; per decidere il \emph{combat system} invita magari il senior
game designer e il game director. La domanda che si fa chi convoca il meeting è: dato
l'obiettivo, chi è più ragionevole che possa contribuire?

Un 3D artist della classe chiede se allora gli artisti 3D non vengano inclusi. Dipende:
se l'obiettivo è definire la linea creativa dello stile dei personaggi, il relatore
invita anche il 3D artist, e non solo i \emph{concept artist}. Il 3D artist può dare un
input diverso dal concept: vota un'idea perché per lui è \emph{più fattibile}. Un
concept può essere bellissimo ma ``un inferno'' da modellare, e chi sa di doverlo fare
entro giugno non vota un'idea che sa già infattibile.

In realtà si può invitare chiunque: per un pitch su un'area dedicata al basket alla
Games Week, il relatore ha invitato un tecnico di regia che gioca a basket, per avere
anche la prospettiva di chi ha familiarità con il tema. A uno studente che porta
l'esempio di un elemento di interfaccia sul personaggio, il relatore risponde che,
volendo, si trova un motivo per invitare chiunque: l'importante è avere chiaro
\textbf{chi è necessario}. Il gruppo di lavoro può essere di otto e il brainstorming di
dodici, con input esterni a chi poi farà il lavoro. Basta non esagerare: se la maggior
parte dei voti arriva da chi non farà quel lavoro, si vota su qualcosa che poi non
faranno loro, e non ha più senso.

% ---------------------------------------------------------------------
\section{Dalla votazione alla decisione: la matrice impatto--sforzo}
\label{sec:z2-impatto}

\subsection{Massimo impatto, minimo sforzo}

Fatto questo, non si è finito: ora si sa quali idee interessano di più, ma non si è
ancora deciso che cosa fare. Serve la parte decisionale. L'obiettivo di qualsiasi
azienda (ed è il motivo per cui al relatore non importa se il programmatore fa il test
da casa cercando su Google) è \textbf{massimizzare l'impatto e minimizzare lo sforzo}.
Vale per i progetti, i task, il recruitment, per qualsiasi cosa: ``non perché sono
pigro, ma perché sarebbe stupido fare diversamente''.

Le idee selezionate si mettono quindi in un'altra matrice, con due assi:
\begin{itemize}
  \item sull'asse verticale il \textbf{valore}, cioè quanto \textbf{impatto} avrebbe
        l'idea (in alto l'impatto alto);
  \item sull'asse orizzontale lo \textbf{sforzo} (\emph{effort}), da basso (a sinistra)
        ad alto (a destra).
\end{itemize}
In alto a sinistra ci sono quindi le idee ad alto impatto e basso sforzo, in alto a
destra quelle ad alto impatto e alto sforzo, e così via.

\paragraph{L'esempio Adidas.} Se il cliente è Adidas, con la linea di abbigliamento per
il basket, un'idea a basso valore potrebbe essere una campagna su TikTok in cui gli
utenti condividono i loro momenti con un pallone da basket; un'idea ad alto valore
potrebbe essere fondare un parco divertimenti dove i giocatori NBA vengono a sfidarti
uno contro uno. Questo non dice che una sia meglio dell'altra: dice che l'impatto è
diverso. Allo stesso modo si valuta lo sforzo: la campagna social è a basso sforzo, il
parco divertimenti ad alto sforzo per molti motivi (i permessi di costruzione, i
giocatori NBA).

\subsection{Posizionare le idee, insieme e con precisione}

Il gruppo che a quel punto è \textbf{il gruppo decisionale}, cioè chi farà il lavoro,
deve posizionare sulla matrice le idee ritenute migliori, e deve essere \textbf{d'accordo
tutto il gruppo}: ``siamo tutti d'accordo che il parco dei dinosauri di Stefano e Andrea è
ad alto valore ma anche ad alto sforzo? Sì? Allora lo mettiamo qui''. E bisogna indicare
\textbf{un punto preciso}, non solo il quadrante: piazzare un'idea in un punto non è come
piazzarla poco più in là. Il risultato è una sorta di \textbf{mosaico} delle idee sul
grafico.

\subsection{Leggere la matrice in diagonale}

Davanti a una situazione tipica, con varie idee sparse, il relatore chiede alla classe:
se foste i project manager, se doveste fare voi queste cose, \emph{quali} fareste? Al
plurale, perché spesso da un brainstorming si esce con più idee realizzabili e si deve
scegliere su quali concentrarsi. Le risposte variano: c'è chi sceglie le due con più
impatto, chi aggiunge che dipende da tempi, budget e risorse, chi preferisce quella più
innovativa ma anche più rischiosa.

Il metodo che secondo il relatore va applicato è \textbf{procedere in diagonale}, dalle
idee a basso sforzo e alto impatto verso quelle ad alto sforzo e basso
impatto\footnote{Nella slide le idee erano rappresentate da post-it colorati (rossa,
rosa o salmone, azzurra, verde acqua, viola, blu, gialla), e le indicazioni del
relatore si riferiscono a quei colori. Dalla sola trascrizione non è possibile
ricostruire con certezza la posizione di ciascun post-it, e alcune indicazioni sui
colori risultano contraddittorie; qui è riportato il principio generale, che il
relatore chiarisce nelle risposte alle domande.}:
\begin{enumerate}
  \item un'idea a \textbf{basso sforzo e alto impatto} (nella slide, quella rossa) va
        fatta \textbf{per prima, sempre}: ``non dovrebbe essere nemmeno in discussione'';
  \item poi, a seconda delle risorse e del tempo, si possono fare
        \textbf{contemporaneamente} le idee che stanno sulla stessa diagonale
        successiva, che hanno cioè un rapporto simile tra impatto e sforzo;
  \item poi le successive;
  \item le idee ad \textbf{alto sforzo e basso impatto} (nella slide, la gialla)
        \textbf{non vanno nemmeno considerate}. Nessuno, infatti, l'aveva scelta.
\end{enumerate}
Per le scelte intermedie valgono due criteri, chiariti rispondendo a uno studente che
non capiva perché alcune idee fossero preferibili ad altre: \textbf{a parità di sforzo si
sceglie l'idea con più impatto}; \textbf{a parità di impatto si sceglie quella con meno
sforzo}. Con risorse infinite si farebbero anche le idee più costose; ma di solito le
risorse sono limitate, e allora tra due idee con lo stesso sforzo e un impatto molto
diverso la scelta è evidente. Se ci piacerebbe fare un'idea più ambiziosa ma abbiamo
risorse solo per una meno costosa sulla stessa diagonale, va bene lo stesso: abbiamo
comunque stabilito che è meglio delle idee sulle diagonali successive. Quanto fare dipende
dalla \textbf{bandwidth} a disposizione.

\paragraph{``La gialla era la mia idea: devo offendermi?''} No. Qualcuno ha detto che
l'autore fa schifo? No. Si è concordato che la sua idea, per quanto votata (``qualcuno ha
votato per la mia idea''), è piuttosto difficile da realizzare rispetto al suo impatto.
Non si sta dicendo che non sa più fare il suo lavoro. Anzi, se lui stesso dovesse
scegliere tra la propria idea gialla e la rossa di un altro, sceglierebbe la rossa: non è
nemmeno una questione di discussione. La matrice rende \textbf{visivamente evidente} che
cosa fare, e permette di aggirare molte discussioni che non sono più soggettive:
oggettivamente, dovremmo fare così.

\subsection{Il ruolo del moderatore e dell'esperto}

È qui che lavora il moderatore: ``non è più una discussione; visto il posizionamento,
dobbiamo fare queste cose''. Non si decide a maggioranza dove mettere le idee: se uno
posiziona un'idea e un altro non è d'accordo, difficilmente dirà ``la rossa dovrebbe stare
al posto della gialla''; dirà ``secondo me è un po' più a destra, un po' più in basso'', e
si fanno gli aggiustamenti. Se invece nascono discussioni del tipo ``la rossa dovrebbe
stare al posto della gialla'', vuol dire che non si è definita bene l'idea e non è chiaro
quanto sia davvero difficile da realizzare.

Oppure interviene \textbf{l'esperto}: è l'esempio del concept artist e del 3D artist. Il
concept artist dice che un'idea è a basso sforzo e alto impatto; il 3D artist risponde ``ne
condivido l'impatto, ma dovendola fare io ti assicuro che lo sforzo non è basso, è molto
alto''. Allora è un altro discorso, perché l'esperto sta dicendo che lo sforzo è diverso
da quello stimato. Ma alla fine tutto deve essere concordato.

\subsection{Chi decide}

\paragraph{Il gruppo decisionale può essere diverso da quello del brainstorming?} Sì: a
questo punto escono quelli invitati come ``consulenti'', che hanno dato idee ma non devono
poi realizzarle. Farli votare anche qui significherebbe perdere il collegamento tra chi
deve fare le cose e che cosa si decide di fare. In quel momento non hanno responsabilità:
hanno dato un input, e se decidessero che cosa devono fare gli altri dovrebbero esserne
i manager. Non dovrebbero poter dire agli altri che cosa fare se non ne hanno la
responsabilità. Vale per tutti i team, piccoli e grandi; secondo il relatore chi decide
dovrebbe essere \textbf{al massimo una decina di persone}. In un team di dieci che fa il
brainstorming tutto insieme, va benissimo che chi vota coincida con chi ha partecipato:
l'importante è che \textbf{chi vota sia chi farà il lavoro}.

\paragraph{Lo sforzo è il costo o la fatica?} Uno studente chiede se l'alto impatto
includa anche il costo. No, il costo sta nello sforzo, che per il relatore va
\textbf{ragionato in risorse}: denaro, tempo e così via. Un'idea può avere un basso costo
monetario ma un alto costo in tempo, e allora va bilanciata e si troverà a metà. È lì che
si spostano le idee e che avviene la conversazione: i tecnici dicono che lo sforzo tecnico
è un certo, poi arriva l'\emph{accountant} che dice che a livello economico è un altro, e
la si sposta.

\paragraph{E il piacere di fare il progetto?} In un gruppo di due o tre persone, osserva
uno studente, si finisce per votare anche per preferenza personale: quando si fa un gioco
bisogna avere anche il piacere di farlo. Ma \textbf{anche quello è impatto}: l'impatto
sulla motivazione. Come lo sforzo sono le risorse, l'impatto è l'impatto \emph{su di te,
sul team, sul mercato}, su ciò che si vuole. Il relatore lo ha tenuto volutamente
generico: può essere quanti soldi si fanno o quanto è utile per il proprio CV, e
l'impatto cambia tra un progetto accademico e uno professionale.

% ---------------------------------------------------------------------
\section{Consigli pratici e domande sul brainstorming}
\label{sec:z2-qa-brainstorming}

\subsection{Musica e cambio di scenario}

``Sembrano tutte stupidaggini, ma vi assicuro che aiutano.''

\paragraph{Musica di sottofondo.} È dimostrato che aiuta\footnote{Le affermazioni sugli
effetti della musica (chiave maggiore, rumore ambientale) sono riportate come dette in
lezione e non sono state verificate.}. Meglio una musica con \textbf{tempo e volume
costanti}: niente \emph{bass drop}, niente variazioni di tema o di ritmo. Va bene la
musica fatta di rumore ambientale a basso volume; è stato rilevato anche che le canzoni
in \textbf{chiave maggiore} aiutano la parte creativa a produrre idee più originali.
Meglio la \textbf{musica strumentale}, senza testo: per quanto ci si impegni a non
ascoltarlo, il cervello cercherà di capire le parole. È un effetto marginale e non
succede niente se capita la canzone con il testo, ma si mette in conto che ``l'uno per
cento del cervello'' sarà impegnato a capirlo, come quando in treno si ascolta, che lo si
voglia o no, la telefonata di qualcun altro. È \emph{min-maxing}.

\paragraph{Cambiare scenario.} A Riot c'era una stanza dedicata al brainstorming, pensata
per essere creativa: le finestre sul fiume, le porte a vetri colorati, tutto fatto per
ispirare creatività. Ma fare brainstorming sempre nella stessa stanza generava sempre lo
stesso tipo di idee: anche solo cambiare ambiente aiutava ad averne di diverse. Il
relatore spesso spostava il meeting, ad esempio nello Starbucks sotto l'ufficio. Lo
spazio fisico gioca un ruolo importante, ma anche quello virtuale: se lo si fa sempre su
Discord, si può provare con ognuno uno sfondo diverso e simpatico. Oppure: ha sempre
moderato la stessa persona? Si prova con un altro moderatore. Si cambiano i fattori
perché i risultati non siano sempre gli stessi.

Uno studente chiede se il cambio di scenario possa essere anche qualcosa come passarsi una
palla per avere la parola, o un altro espediente per tenere viva l'attenzione. Con questo
modello, che non richiede di prendere la parola, l'esempio della palla perde senso, ma come
esempio va benissimo: anche il comportamento fisico e il \emph{modus operandi} sono un ottimo
cambio di scenario. Come lo sono cambiare il moderatore o \textbf{cambiare l'obiettivo}: se
il cliente chiede X, gli si darà un'idea X, ma si può provare anche a votare per un'idea Y,
facendo più votazioni con vari criteri di valutazione.

\subsection{Domande}

\paragraph{L'idea scelta deve essere approvata da qualcun altro?} Se un manager dice ``fate
questo brainstorming senza di me'', sta \textbf{delegando la scelta}: ``mi fido di quello che
sceglierete''. Avendo preparato prima limiti, paletti e \emph{constraint}, si aspetta che la
soluzione portata sia in linea con essi; ma se non partecipa deve fidarsi della scelta. Se
poi la scelta debba passare anche da un livello superiore dipende dall'azienda e dagli
accordi precedenti.

\paragraph{Si può fare un brainstorming ibrido, tra presenza e remoto?} Uno studente fa
l'esempio di un gruppo in cui uno è a New York e gli altri a Milano. Sì, senza problemi: si
fa una call, il portatile punta verso il muro con le idee, o le si mostra fisicamente con
il cellulare. Chi è da remoto scrive i suoi post-it, ad esempio mandandoli su WhatsApp al
moderatore o a qualcuno presente che li scrive per lui, e vota dicendo ``mettimi un pallino
qui, qui e qui'' (in quel caso il voto non è anonimo, ma si può fare tranquillamente); se
deve parlare, parla.

\paragraph{E se i colleghi non riconoscono o non condividono il metodo?} Ci deve essere un
\textbf{agreement}: non si può arrivare e dire ``facciamo così'' se gli altri non vogliono.
Questo metodo funziona solo se tutti sono d'accordo a usarlo. Il relatore ne ha proposto
uno tra i cento metodi di brainstorming esistenti, un buon compromesso tra semplicità ed
efficienza (ce ne sono di più semplici e di molto più efficienti ma più complessi). Se il
gruppo non ne vede il valore, al massimo insisterebbe per \textbf{fare almeno una prova},
altrimenti diventa chiusura mentale; se poi per quel team non funziona, se ne cercano
altri.

\paragraph{E nei gruppi molto piccoli, dove diventa un duello?} Uno studente ha lavorato in
gruppi di tre o anche di due persone, dove il brainstorming diventa quasi un duello in cui
ognuno difende la propria idea. Bisogna \textbf{aspettare la discussione}: invece di ``scegliamo
tra la mia idea e la tua'', si cerca di concordare quale idea abbia l'impatto maggiore e
quella lo sforzo minore. E se non si è d'accordo, bisogna \textbf{portare esempi}: se uno
dice ``questa idea costa di più'', l'altro chiede ``in che senso? Che cosa costa di più?''. Se
non sa rispondere, non si può fare qualcosa senza sapere come va a finire: meglio rimandare
la discussione a quando si avranno le informazioni. A volte la domanda giusta è \emph{perché
non abbiamo queste informazioni?}: forse è \textbf{prematuro} decidere, perché non si hanno
ancora tutti i fattori.

Nei team piccoli non dovrebbe esserci un ``si fa come dici tu o come dico io'', una scelta a
pelle: si dovrebbero conoscere così bene le proprie proposte da poterle posizionare con
assoluta certezza sulla matrice, e avere anche la presenza di spirito di dire ``la tua è
chiaramente migliore, si fa la tua''. Si fa la cosa migliore per entrambi, che dovrebbe
essere oggettiva e non soggettiva. Non è facile arrivarci, ma deve essere l'obiettivo del
team.

% ---------------------------------------------------------------------
\section{Disagree and commit}
\label{sec:z2-disagree}

\subsection{Quando la scelta non ci convince}

Uno studente teme che chi vede scartata la propria idea perda motivazione, in un gruppo
piccolo come in uno grande. Non è un deragliamento: è una conseguenza potenziale di
qualsiasi scelta. Il relatore ha visto gente demotivata perfino dalla settimana di quattro
giorni: ci sarà sempre qualcuno più o meno motivato dalle scelte. Ma qual era l'alternativa?
Un brainstorming classico in cui tutti parlano, chi più chi meno, e si sceglie un'idea che
tutti pensano sia la migliore senza esserne davvero sicuri, perché non si è votato, né
espresso una preferenza, né discusso davvero le idee migliori.

Qui entra il concetto di teamwork del \textbf{disagree and commit}: ``non sono d'accordo con
questa decisione, non mi piace, non penso sia la migliore, ma la faccio''. Per due motivi:
\begin{enumerate}
  \item perché \textbf{è il mio lavoro}, e dovrebbe essere il punto principale;
  \item perché \textbf{voi ritenete che sia la scelta migliore}. Se si è la voce fuori dal coro
        e tutti hanno scelto un'altra strada, non si rema contro: è \emph{default to trust}.
        Se gli altri vi hanno visto più potenzialità, e (grande premessa) \textbf{hanno
        ascoltato} il proprio disaccordo e i suoi motivi, e hanno comunque scelto così, ci si
        fida di loro.
\end{enumerate}
Con una postilla: se le cose vanno male, \textbf{non si dà la colpa agli altri} e non si dice
``ve l'avevo detto'', come se non si facesse parte del team. Disagree and commit vuol dire ``mi
fido del team e facciamo questa cosa''; se poi si rivela sbagliata, \emph{il team} ha
sbagliato, e la volta dopo si cambierà il modo di scegliere. Altrimenti non è disagree and
commit, è ``\emph{disagree and ostracize}'': disaccordo più ostilità. Non bisogna cercare il
colpevole, perché si è deciso tutti insieme (sezione~\ref{sec:z1-default-trust}).

Se si pensa davvero che si stia sbagliando clamorosamente il posizionamento di un'idea, si
dice piuttosto ``non andiamo avanti finché non definiamo bene lo sforzo di questa cosa''; ma
bisogna saperlo giustificare, altrimenti è chi dissente a dover fare un passo indietro.

\subsection{La demotivazione fa parte del lavoro}

La demotivazione ci sarà sempre, perché non si sceglierà mai solo ciò che piace: per quanto
piaccia il proprio lavoro, ci saranno task sgradevoli e scelte che non spettano a noi. E
ci saranno scelte (``ve lo dico per esperienza personale'') che si è convinti siano
sbagliate, al punto di pensare ``piuttosto lascio il team'', ma che si finisce per accettare,
per poi rendersi conto che gli altri avevano ragione. Come c'è un \emph{bias} verso l'amico e
il nemico, c'è anche un \textbf{bias verso noi stessi}: crediamo di conoscere bene una cosa, e
se altre dieci persone dicono ``secondo me è fattibile'', bisognerebbe aprirsi alla
possibilità che ne sappiano più di noi.

``Non lo accetto ma lo faccio'' dovrebbe essere alla base del lavoro insieme, perché è
\emph{default to trust}: rifiutarsi (``non mi fido di voi, non lo faccio'') significa non
fidarsi del giudizio del proprio team. E dire ad alta voce ``non sono d'accordo, ma ti darò
comunque il massimo'' è la cosa che \textbf{fa guadagnare di più il rispetto degli altri}.
Non c'è niente di meglio che vedere qualcuno che aveva un'opinione diversa impegnarsi per
far vedere di avere torto, facendo andare bene la cosa. Per il relatore dimostra solo che è
una persona matura e importante; ha avuto persone che, passate da questa situazione, gli
hanno fatto pensare ``è molto più matura di quanto credessi'', e che ha valorizzato molto di
più in seguito, perché sa di potersi fidare di chi fa disagree and commit: è stato aperto
nella comunicazione, ha detto che cosa secondo lui non andava, ma non ha fatto niente per
remare contro, anzi ha dato il suo supporto. ``Non posso chiedere di più a una persona.''

\subsection{Chi arriva a lavori iniziati}

\paragraph{Se il team cambia in corso d'opera, chi subentra come si rapporta alle decisioni
già prese?} (per dimissioni, promozioni e così via). Deve \textbf{accettare lo status quo}. Al
relatore capita spesso di entrare in team con progetti già in corso: si fa dire che cosa è
stato scelto finora e va \textbf{in commit di default}: ``non ho preso parte a questa decisione,
quindi non sono in \emph{agreement}, ma ora faccio parte del team, quindi vado in commit''.

\paragraph{E se l'esperienza dice che una scelta ha uno sforzo molto diverso?} Ad esempio
che un'idea ``rossa'' ha in realtà lo sforzo di una ``viola''. Lo si fa presente, ma dipende dal
\textbf{tempismo}: se la scelta è appena stata fatta e c'è ancora modo di cambiarla, la si
segnala; se si è al 90\% del progetto, è controproducente dire ``fermi tutti, passo
indietro'', e a un certo punto bisogna rispettare la scelta del team (``avete scelto così,
andiamo avanti e vediamo come va''). E se si sa che, calcolando ferie e periodo estivo,
un'animazione prevista per giugno arriverà a settembre? È molto \emph{case by case}: a
seconda dell'apporto si può modificare il risultato o no. È come riaprire un processo se
emergono nuove prove: se dimostrano che l'imputato non è l'assassino, è giusto riaprirlo; se
dimostrano che l'assassino ha sparato con la sinistra invece che con la destra, conta fino a
un certo punto. È giusto segnalarlo, ma non sempre quello che si porta giustifica uno
``\emph{scratch that}, rifacciamo tutto'': diventa \textbf{un feedback, non una risoluzione}
(``state facendo così, questo è il mio feedback, poi vediamo come va'').

Con cinque minuti di anticipo si va in pausa pranzo, con l'accordo per il pomeriggio: circa
un'ora di presentazione più ``standard'' sul leading from within, poi domande e risposte,
``e quando finiscono le domande abbiamo finito''.

% ---------------------------------------------------------------------
\section{Leading from within: perché parlare di leadership}
\label{sec:z2-leadership}

\subsection{La ripresa del pomeriggio}

Alla ripresa, mentre si aspettano gli ultimi arrivati, il relatore chiede se durante la
pausa siano venute in mente domande. L'unica è logistica: le slide verranno date a fine
workshop? Sì: il relatore le passerà all'accademia (a Vanessa, Andrea ed
Enrico\footnote{Nella trascrizione l'accademia è chiamata ``Ive'' e ``Hive'', come nella
prima giornata; per i nomi dei referenti vedi la nota all'inizio del
capitolo~\ref{cap:z1-feedback}.}), che le farà avere alla classe. Gli appunti vanno
comunque presi, meglio su quello che dice che su quello che mostra.

L'ultima parte, dice, dovrebbe essere abbastanza veloce da presentare, ma vale la pena
fermarsi a parlarne, perché è un discorso più avanzato. È la logica chiusura di quanto
detto tra la prima giornata e la mattina: tirare le fila e mostrare come tutto faccia parte
di un unico filone di ragionamento e progettualità, con un unico scopo, un unico obiettivo e
un unico modo di ragionare.

La voce sembra a posto, ma il relatore è stanco dopo due giorni, come immagina lo sia la
classe, nello \emph{stretch} finale di due giornate molto piene. Sono cose che in condizioni
normali spiega al suo team nell'arco di settimane, con il tempo di metterle in pratica,
discuterle e approfondirle. Qui è tutto concentrato: servirà tempo per metterlo in atto e
fare le proprie prove, quindi non bisogna spaventarsi se non verrà tutto giusto la prima
volta.

\subsection{Il leader non è il più bravo}

Perché parlare di leadership in un workshop non rivolto solo ai leader? La risposta
completa arriverà più avanti, ma il punto di partenza è che il relatore \textbf{non crede
nei leader di team o di progetto che detengono da soli il sapere e il modo di fare le
cose}. Come detto il giorno prima, essere un leader, un manager o un capo non dovrebbe
voler dire essere il migliore in quel lavoro, ma essere \textbf{la persona più adatta a
ricoprire quel ruolo}: guidare un gruppo significa \textbf{occuparsi delle persone} al suo
interno.

L'obiettivo che il relatore si dà, e che secondo lui ogni manager dovrebbe avere in mente,
è che quando si torna dal lavoro ci si senta \textbf{ispirati, sicuri e realizzati}, e che
questo sia \textbf{un diritto naturale di tutti}, non un lusso riservato a chi ricopre certe
posizioni, a pochi fortunati o ``a qualcuno sì e a qualcuno no''. Ha notato,
specialmente in Italia, la tendenza a vedere tutto ciò che viene prima del diventare
manager come la \emph{gavetta} necessaria per arrivarci. Invece chiunque faccia parte di
un'organizzazione, che sia un'azienda di videogiochi o altro, dovrebbe sentirsi soddisfatto,
ispirato e sicuro in quello che fa; e non è corretto che un'azienda tratti bene le persone
``da manager in su''. Anche quando non lo dicono, a volte i fatti sono quelli: da manager si
viene trattati in un certo modo e si ha accesso a certe cose, altrimenti no. Non si tratta di
\textbf{appiattire la gerarchia}: si tratta di assicurarsi che, pur dando cose diverse ai
diversi livelli, a tutti i livelli venga riconosciuto uno strato comune di
\textbf{rispetto, fiducia e professionalità}.

\subsection{Le carriere divise: manager o \emph{subject matter expert}}

Uno studente chiede conferma: quindi un \emph{lead artist} non è necessariamente più bravo
di un senior a fare l'artista? Esatto. E qui il relatore apre una parentesi. In molte
aziende, e soprattutto in quelle in cui tende a lavorare lui, la motivazione principale per
diventare manager è avere una promozione e quindi un aumento, più che il desiderio di
prendersi cura del team e delle persone. Come si ovvia al problema? \textbf{Dividendo le
carriere}. Arrivati al livello senior, il relatore fa con ognuno dei suoi senior
\textbf{``the Talk''}, per esporre le due opzioni che hanno davanti:
\begin{description}
  \item[Busta A: diventare manager] È il percorso standard: si valuta se si è adatti a fare
        il leader e a guidare un team, e si danno gli strumenti per farlo, perché
        ``nessuno nasce leader imparato'': prima di provarlo sulla propria pelle nessuno sa
        che cosa voglia dire. Per quanto si possa essere carismatici, fare il leader è
        un'altra cosa, e non vuol dire solo essere carismatici. Ci sono formazione,
        strumenti e un percorso di crescita.
  \item[Busta B: diventare \emph{subject matter expert}] Una persona che va oltre il senior
        nella stessa disciplina: diventa \textbf{l'esperto dell'intera azienda} su quella
        disciplina, perché la sua priorità non è aiutare gli altri ma essere il migliore.
        ``Non c'è assolutamente nulla di male.'' Il suo ruolo è essere il più specializzato
        in quel campo, il punto di riferimento da cui tutti vanno a farsi spiegare le cose.
\end{description}
Il concetto è abbastanza nuovo, e il punto chiave è che l'azienda \textbf{paga il
\emph{subject matter expert} e il leader esattamente allo stesso modo}: non c'è un guadagno
economico nello scegliere di fare il manager invece dello specialista, e il relatore valuta
quel tipo di conoscenza quanto il saper gestire persone. Così chi vuole essere il migliore,
o anche solo un aumento, può continuare a fare ciò che sa fare e avere l'aumento; chi vuole
l'aumento \emph{e} prendersi cura delle persone sceglie l'altra strada. Così si fa una scrematura sul
\emph{motivo} per cui si vuole fare il manager. Sembra una scelta piccola, ma è molto
importante capire se ci si vuole occupare di se stessi o degli altri. Non è un giudizio
(``non sto dicendo che sia meglio l'uno o l'altro''): le due cose \textbf{non convivono}, e
bisogna sapere che cosa si vuole fare, sapendo che l'azienda tratterà allo stesso modo le
due scelte.

\paragraph{Si può essere entrambi?} Si può scegliere ad esempio il percorso da lead artist
e voler comunque continuare a migliorare le proprie capacità. Ma, come nelle simulazioni di
conversazione del giorno prima, il relatore chiederebbe: ``se avessi le risorse per fare solo
una delle due cose, quale sceglieresti?''. Perché non le si avrà: non ci sarà il tempo e il
modo di fare bene entrambe. Il suo ``cinismo'' arriva a dire che bisogna scegliere, perché
tenersi buone entrambe le strade porta a farle male tutte e due. Preferisce un ottimo
\emph{people manager} o un ottimo specialista; le figure intermedie, in questo campo,
rischiano di diventare pericolose per sé e per il team.

\paragraph{Specializzazioni verticali.} La strada dello specialista permette
specializzazioni molto verticali: in chat qualcuno cita il \emph{mecha designer}, ma può
esserci anche lo specialista di animazioni delle espressioni facciali. Naturalmente l'azienda
deve avere bisogno di quella competenza (non ci si specializza in cose che non servono), e non
è detto che ne serva uno solo. È un percorso di carriera parallelo e alternativo a quello del
manager ``standard''.

% ---------------------------------------------------------------------
\section{Il ``momento Piero Angela'': la biologia del gruppo}
\label{sec:z2-biologia}

\subsection{Perché una digressione scientifica}

Il relatore annuncia un ``piccolo excursus biologico'', fondamentale per capire dove vuole
arrivare, e prova a mettere un sottofondo musicale adatto: la condivisione dell'audio non
funziona, quindi manda il link in chat perché ognuno lo faccia partire per conto suo. Un
studente lo battezza subito ``momento Piero Angela''. Il relatore cerca di far combaciare le
motivazioni aziendali del suo metodo di leadership e di progettualità con \textbf{basi
scientifiche}: in passato ha collaborato con varie università che si occupavano di
neuromarketing e neurobiologia, per avere una conferma o una smentita del fatto che ciò che
ipotizzava ``non fosse una cavolata colossale''. Semplificando moltissimo, si guardano solo
alcune sostanze che si traducono direttamente nel modo in cui si vive l'essere in un gruppo,
in un team, in un'azienda\footnote{L'impianto di questa parte (endorfine, dopamina,
serotonina, ossitocina e cortisolo, il villaggio protetto dai pericoli esterni, il ``cerchio
di sicurezza'') ricalca da vicino quello del libro di Simon Sinek \emph{Leaders Eat Last}
(2014), che il relatore cita poco più avanti. Le descrizioni biologiche sono molto
semplificate, come avverte lo stesso relatore, e vanno lette come tali.}.

\subsection{Endorfine}

Le \textbf{endorfine} sono rilasciate per \textbf{mascherare il dolore}. Lo sa chi corre, chi
fa pesi o ginnastica, ma anche chi ha riso così tanto per una battuta da avere male ai
fianchi: è quella sensazione di dolore in cui però si è contenti di provarlo. Il \emph{runner's
high} è l'esempio più calzante: essere distrutti dopo una corsa ma soddisfatti. È la risposta
biologica allo sforzarsi in cose che sono per il nostro bene, un modo con cui il corpo ci
``convince'' a farle.

\subsection{Dopamina}

La \textbf{dopamina} è più familiare: è la \textbf{ricompensa per un'azione}. Dalle cose più
banali (si mangia e si riceve dopamina) a quelle viste in mattinata: anche completare una
checklist rilascia dopamina (sezione~\ref{sec:z2-pianificare}). Viene rilasciata quando si
vede l'obiettivo, come il traguardo durante una corsa; è la stessa cosa che accadeva ai
cacciatori del villaggio in tempi antichi, che rilasciavano dopamina quando vedevano le tracce
e venivano così incentivati a seguirle. Sono concetti fondamentalmente ``animali'', ma si
applicano tuttora anche in azienda.

\paragraph{Una nota di colore: la dipendenza.} La dopamina \textbf{crea dipendenza}: molti casi
di ossessione per alcol, nicotina, gioco d'azzardo e social media sono problemi legati alla
dipendenza da dopamina. Secondo il relatore non è un caso che le diagnosi di ADD e ADHD siano
aumentate del 60\% negli ultimi dieci anni\footnote{Nella trascrizione ``EDD e HD''. Il dato
e il nesso causale con la maggiore disponibilità di alcol, nicotina, gioco d'azzardo e social
media sono riportati come detti in lezione e non sono stati verificati; l'aumento delle
diagnosi viene in genere attribuito anche ad altri fattori, come la maggiore consapevolezza e
il cambiamento dei criteri diagnostici.}, perché buona parte di quegli stimoli è diventata
ampiamente disponibile proprio in quel periodo. Anche la dipendenza dal lavoro, a volte, ha la
stessa origine: si vuole il rispetto che deriva dal lavoro, il \emph{rush} che deriva
dall'averlo.

\subsection{Il villaggio di diecimila anni fa}

Si fa poi un salto di diecimila anni, in un villaggio. Gli esseri umani, per quanto più o meno
introversi o estroversi, sono \textbf{animali sociali} che tendono a vivere esperienze di
gruppo; le aziende nascono da un gruppo di persone che si danno un obiettivo condiviso, come
in antichità un gruppo di persone che voleva sopravvivere insieme. Quello che non è cambiato è
che \textbf{fuori dal gruppo ci sono dei pericoli}: tigri, lupi, draghi (``do per scontato
che nell'antichità ci fossero i draghi, perché D\&D mi insegna così''). Per evitarli, i gruppi
tendono a riunirsi e a creare difese intorno a sé: da qui le prime difese dei primi villaggi,
per isolarsi dai pericoli esterni. ``Vi giuro che avrà tutto senso.''

\subsection{Serotonina}

La \textbf{serotonina} viene rilasciata anche \emph{indirettamente}: non deve per forza
riguardare noi. Quando si assiste a una cerimonia o celebrazione pubblica che include una
persona cara, si rilascia serotonina come se si fosse noi l'oggetto del riconoscimento. È uno
dei motivi per cui si fanno diplomi e lauree ``all'americana'', con la consegna sul palco
davanti a tutti: gli studi avrebbero mostrato che anche i genitori in platea rilasciano
serotonina, come chi riceve il premio. Lo stesso vale per parenti, amici, matrimoni: sono
esperienze condivise, e si tende a condividere in pubblico riconoscimenti e successi perché in
tutti i partecipanti viene rilasciata serotonina.

Purtroppo ha a che fare anche con i beni materiali: la tendenza all'\emph{hoarding}, cioè a
collezionare eccessivamente qualcosa, nasce dal fatto che quell'attività rilascia serotonina,
ingannando la mente nel credere di aver ottenuto un successo.

È chiamata anche \textbf{la sostanza della leadership}, perché aiuta a riconoscere in qualcun
altro chi ha valore per il gruppo. Nel villaggio, quando arriva un pericolo, c'è sempre un
gruppo di persone che dice ``ce ne occupiamo noi'': il gruppo sa che si stanno prendendo un
rischio, e la serotonina è il modo per riconoscere che meritano rispetto. Il relatore consiglia
il libro di Simon Sinek \emph{Leaders Eat Last}\footnote{Nella trascrizione ``Leadership
Lust''. Il titolo è \emph{Leaders Eat Last: Why Some Teams Pull Together and Others Don't}
(2014).}: i leader mangiano per ultimi,
all'opposto di quanto accadeva in antichità, quando al ritorno dalla caccia mangiavano per
primi i cacciatori, che dovevano essere in forze per la caccia successiva, poi donne e
bambini, poi gli altri. C'è un \textbf{riconoscimento sociale}: tu fai qualcosa per proteggere
il gruppo, noi ti riconosciamo qualcos'altro in cambio.

Le aziende sono strutturate in modo molto diverso? È una domanda retorica: quando si dice che
serve un capo, un manager, un rappresentante, si sta dicendo che quella persona dovrebbe avere
l'interesse di proteggere il gruppo e di prendersi cura di tutti. Per questo sapere che il
proprio capo guadagna di più non fa particolarmente arrabbiare: lo si vede come quasi naturale,
perché più responsabilità si ha verso il gruppo più si guadagna, dato che quello dovrebbe essere
il suo scopo. Ed è anche il motivo per cui il relatore si arrabbia molto quando legge di CEO o
\emph{executive} che, a fronte di licenziamenti di migliaia di persone, si assegnano dei bonus.
È una rabbia non del tutto razionale, che sente da dentro: è come se si stesse violando un
\textbf{contratto sociale}. Il contratto è ``tu ti prendi dei rischi per me e mi difendi, io ti
riconosco più cibo, più soldi''; quando si abusa di quella posizione per un beneficio personale
a scapito di chi si dovrebbe proteggere, tutti lo vedono come un'ingiustizia, perché si viene
meno al contratto stabilito diecimila anni fa con la nostra biologia.

\subsection{Ossitocina}

L'\textbf{ossitocina}\footnote{Nella trascrizione ``ossidocina''.} è la \textbf{sensazione di
sicurezza} che si genera quando si è vicini a qualcuno. Non è un caso che tra i legami di
affinità e collaborazione più forti ci siano quelli dei team che si uniscono contro un
\textbf{nemico comune}: i team che superano insieme avversità e difficoltà ne escono molto
legati. Tra le persone che il relatore sente più spesso in questo periodo ci sono quelle con cui
ha condiviso momenti difficili, un cattivo manager, una brutta esperienza lavorativa: sa che
erano lì con lui. Vale anche per i veterani di guerra o per i sopravvissuti a una stessa
tragedia, che sviluppano un rapporto molto forte (se ne è parlato molto negli Stati Uniti a
proposito dei veterani), perché hanno rilasciato immense dosi di ossitocina. È anche ciò che fa
accettare subito il figlio alla madre quando lo prende in braccio dopo il parto, ed è il motivo
per cui, a migliaia di anni di distanza, si stringono ancora accordi con una stretta di mano o
un abbraccio: c'è sempre qualcosa di fisico, generato dal bisogno di ossitocina.

Viene generata anche quando si compiono o si \textbf{assiste ad atti di generosità} senza
aspettarsi nulla in cambio. Se il relatore dicesse di aver donato cinquanta euro a
un'associazione di beneficenza, la reazione sarebbe ``ok, bene''; se dicesse che ogni sabato
pomeriggio fa volontariato, lo si vedrebbe come qualcosa di molto più prezioso, e la cifra
potrebbe crescere parecchio prima che qualcuno la consideri più preziosa dei suoi pomeriggi:
tutti sanno che la vera donazione è \textbf{il tempo} investito, che vale più del bene
materiale. E basta assistere: vedere qualcuno che aiuta un'anziana ad attraversare la strada fa
stare meglio anche noi. È percepita a livello di gruppo. Infine, l'ossitocina rafforza il
\textbf{sistema immunitario}: sentirsi sicuri e tranquilli fa stare meglio e, in un certo senso,
ammalare di meno.

\subsection{Cortisolo}

L'opposto è il \textbf{cortisolo}: la sensazione di stress, paranoia e paura, la risposta
\emph{fight or flight} (o scappo o combatto), con battito accelerato, muscoli in tensione e così
via. Tra i suoi effetti:
\begin{itemize}
  \item \textbf{blocca il sistema immunitario}: se si sta per morire, al corpo non importa di
        prendersi un raffreddore. È il motivo per cui le persone stressate si ammalano di più;
  \item \textbf{blocca l'ossitocina}: non ci si può sentire tranquilli quando si è in ansia;
  \item soprattutto, e questo è il punto più importante per le aziende, è
        \textbf{contagioso}. In qualsiasi documentario, quando una gazzella che bruca alza la
        testa e percepisce un movimento nella savana, scappano tutte, non solo lei. Succede,
        non sorprendentemente, anche in ufficio: quando qualcuno ha paura perché potrebbero
        esserci licenziamenti, o perché il manager è molto aggressivo, la paura si diffonde nel
        gruppo ed è molto dannosa.
\end{itemize}
È anche l'effetto alla base della \textbf{sindrome dell'impostore}: sentirsi costantemente
inadeguati perché si teme che succeda qualcosa significa vivere in uno stato di stress. Ed è il
motivo per cui chi è insoddisfatto del lavoro tende a stare peggio anche fisicamente e ad
ammalarsi di più.

\begin{center}
  \begin{tabular}{@{}p{0.17\textwidth}p{0.36\textwidth}p{0.37\textwidth}@{}}
    \toprule
    \textbf{Sostanza} & \textbf{Che cosa fa (in sintesi)} & \textbf{Esempi dalla lezione} \\
    \midrule
    Endorfine & mascherano il dolore di uno sforzo che ci fa bene &
      \emph{runner's high}, ridere fino ad avere male ai fianchi \\
    \addlinespace
    Dopamina & ricompensa un'azione, si attiva vedendo l'obiettivo; crea dipendenza &
      spuntare una checklist, le tracce per il cacciatore, gioco d'azzardo, social \\
    \addlinespace
    Serotonina & riconoscimento sociale, anche indiretto; ``sostanza della leadership'' &
      lauree e premi in pubblico, chi difende il villaggio, \emph{hoarding} \\
    \addlinespace
    Ossitocina & sicurezza e legame; rafforza il sistema immunitario &
      il nemico comune, i veterani, la stretta di mano, il volontariato \\
    \addlinespace
    Cortisolo & stress, \emph{fight or flight}; blocca difese immunitarie e ossitocina; è
      contagioso & la gazzella che scappa, la paura dei licenziamenti, la sindrome
      dell'impostore \\
    \bottomrule
  \end{tabular}
\end{center}

% ---------------------------------------------------------------------
\section{Domande sulla biologia del gruppo}
\label{sec:z2-qa-biologia}

\paragraph{Il gioco d'azzardo include i gacha e le loot box?} Assolutamente sì, anzi ne sono il
\textbf{driver principale}: i gacha, i FIFA Ultimate Team, lo ``sbustare'' da \emph{Magic} a
\emph{Hearthstone}, e in generale il ``vediamo se stavolta mi va bene'' sono tutti gioco
d'azzardo. Lo studente aggiunge che molti giochi mobile sono costruiti su questi meccanismi e che
c'è chi gioca più a quelli che ai giochi ``normali''. È il motivo per cui alcuni stati hanno
dichiarato illegali certi giochi\footnote{Il caso più noto è quello del Belgio, che dal 2018
considera gioco d'azzardo le \emph{loot box} a pagamento; anche i Paesi Bassi le hanno
contestate.}.

\paragraph{Come si genera l'ossitocina nei rapporti da remoto?} ``Bellissima domanda.'' È il motivo
per cui, come detto il giorno prima, il relatore insiste perché i meeting si facciano almeno in
\textbf{webcam}, perché i riconoscimenti siano dati e i successi celebrati \textbf{pubblicamente}
(quando non mette troppo a disagio la persona), e in generale cerca di creare \textbf{situazioni
condivise}. Anche lavorando da remoto, insiste per avere il \emph{team building} dal vivo e almeno
\textbf{un incontro all'anno} di persona in ufficio, perché crea situazioni di rapporto diretto.

\paragraph{Come si contrasta il cortisolo?} In chat qualcuno suggerisce la meditazione: sì,
assolutamente. Per il lavoro, la risposta arriva subito dopo nella presentazione.

\paragraph{Il nemico comune è una strategia che un manager può usare?} Spesso \textbf{il manager è
il nemico comune}, senza bisogno di sforzarsi. Ma il relatore chiede di capire bene il punto:
alcuni manager applicano consapevolmente il concetto ``mi ergo io a nemico del team, così sono
tutti uniti contro di me e lavorano bene'', e ne fanno uno stile di management. Decidono di essere
il \emph{bad cop}, il \emph{villain}, ed è così che ottengono risultati. È una scelta lecita e
funziona; il relatore preferisce fare diversamente, ma non per questo dice che non funzioni. Molte
volte si sottovaluta quanto sia consapevole: ``non si rende conto di quanto sia antipatico a
tutti''; forse se ne rende perfettamente conto e lo fa apposta per essere il bersaglio. Il relatore
l'ha visto fare spesso come soluzione disperata, quando il team era troppo impegnato a litigare al
suo interno: ``piuttosto che vederli litigare tra loro, gli do io un motivo per odiare me, così si
uniscono contro di me''. Non lo consiglia assolutamente, ma l'ha visto funzionare.

Uno studente racconta che nella sua azienda il manager fa vedere \emph{il lavoro stesso} come il
nemico (``prima finiamo, meglio è''). È un modello più basilare ma, come dice lo studente stesso,
funziona.

\paragraph{Quanto è etico puntare sulla dipendenza del giocatore per fare soldi?} ``Non è
assolutamente etico, ed è assolutamente profittevole, altrimenti le aziende non lo farebbero.'' Il
relatore racconta una cosa che lo ha sconvolto quando l'ha scoperta, a Blizzard. Lavorava su
\emph{Hearthstone}\footnote{Nella trascrizione ``Arsdol''.} e si stava discutendo di come
monetizzare il gioco; venne preso come consulente un esperto, israeliano se ricorda bene, che
lavorava a Malta per società di poker online e di \emph{gambling}. Raccontò che molte aziende che
monetizzano sulla dipendenza dal gioco d'azzardo avevano sviluppato metodi più o meno etici (``più
spesso meno''). L'esempio culminante: analizzavano gli utenti occasionali, quelli che giocavano
magari una settimana al mese, per capire \emph{quando arrivava loro lo stipendio} (di solito si
deposita lo stipendio e poi si gioca quello che resta), per colpirli il giorno dopo con email di
offerte per tornare a giocare. Li si stimolava proprio nel momento in cui avevano disponibilità, sapendo
che erano sensibili, per aumentare le probabilità di riaverli come giocatori. ``È orribile. È
terrificante e, come tutte le cose terrificanti, è anche altamente funzionale.''

Gli ha fatto riflettere molto su dove ci si possa spingere per fare soldi. La versione etica, per
lui, è ad esempio far uscire un gioco sotto Natale, quando si è più disposti a regalarlo ai figli:
si stanno ottimizzando le proprie possibilità, e lo accetta. Si scade nel \textbf{predatorio}
quando si sfrutta la debolezza di qualcuno in un momento preciso per estorcergli soldi che non può
permettersi di perdere. Poi ognuno tragga le proprie conclusioni: ``non è un workshop sull'etica del
lavoro''. Non pensa di dover imporre la propria opinione, né si arroga il diritto di dire che chi lo
fa è una bestia senza cuore, visto che è il loro lavoro: ``chiunque riesca poi a dormire in pace con
se stesso, buon percorso''.

% ---------------------------------------------------------------------
\section{La piramide della fiducia}
\label{sec:z2-piramide}

\subsection{La leadership non è un titolo}

Chiuso il ``momento Piero Angela'', il relatore arriva al punto. La frase della prima giornata (``la
capacità di un gruppo di persone di fare cose straordinarie dipende dal modo in cui queste persone
operano insieme'') porta a capire che \textbf{la leadership non è un titolo lavorativo}. Essere
manager da contratto, o avere ``manager'' nel titolo, rende un manager, non un leader. Il relatore
non ha mai visto un team rispettare qualcuno solo perché da contratto doveva farlo: si può avere
paura o timore del proprio manager, ma non lo si deve rispettare perché è il manager. \textbf{Il
rispetto si guadagna}, e si possono influenzare le persone intorno solo attraverso il rispetto
guadagnato, quando si sono accumulati abbastanza esempi e \emph{case history} da giustificare
rispetto e fiducia. Altrimenti, dice, sarebbe il primo a chiedere ``perché vi fidate di me, se non
vi ho ancora dato motivo di farlo?'': non sarebbe razionale.

Da qui la \textbf{piramide della fiducia}, che si costruisce partendo dalla base.

\subsection{Alla base: l'integrità morale}

Per essere un buon leader o un buon manager bisogna \textbf{voler genuinamente aiutare gli altri
membri del team}. È il motivo del discorso sulle carriere divise: è a questo livello che si
screma chi vuole la promozione per i soldi. Chi vuole essere leader non deve cercare il guadagno
personale, ma quello del gruppo. Come detto il giorno prima, il relatore è pronto a litigare per il
suo team e a prendersi le colpe per il suo team, perché vuole il bene del suo team e sa che poi loro
si prenderanno cura di lui. Chi ha in mente questo sta in questa categoria. Poi:
\begin{itemize}
  \item \textbf{essere onesti}. Il relatore può affermare di non aver mai mentito
        intenzionalmente a un collega, a un cliente o a un giocatore. Gli è capitato di dare
        informazioni errate perché in quel momento sapeva quelle: ad esempio annunciare che
        un'espansione sarebbe uscita in un certo mese, e poi l'uscita è slittata\footnote{Nella
        trascrizione l'esempio non è coerente: prima si parla di un'uscita a novembre rimandata
        di una settimana, poi di ``pensavo fosse settembre''. Il senso resta chiaro.}. Ha detto
        una cosa sbagliata, ma tutti hanno capito che non stava mentendo: in quel momento pensava
        che fosse vera;
  \item soprattutto, \textbf{accettare la responsabilità dei fallimenti}: quando si sbaglia,
        saper dire ``ho sbagliato''. Il relatore lo ha messo perché ``non avete idea di quanti
        facciano fatica'' ad ammettere di poter sbagliare: c'è sempre una scusante, una
        mitigazione, un ``sì, ma''. Lui vuole sentire se è colpa tua o no;
  \item \textbf{chiedere scusa} e \textbf{ammettere di non sapere}: anche dire ``non lo so'' viene
        spesso evitato, come accettare un fallimento.
\end{itemize}

\paragraph{``Raccontami il tuo più grande fallimento''.} Per lo stesso motivo, nei colloqui il
relatore chiede spesso, più che ``parlami di un tuo grande successo'' o di un progetto andato bene,
\textbf{``mi racconti il tuo più grande fallimento?''}, e osserva come lo raccontano. Chi non si
aspetta la domanda spesso racconta di quella volta in cui ha provato in tutti i modi a fare in un
certo modo, ma il capo gli ha detto di fare diversamente e il progetto è andato male. ``Non mi hai
risposto: mi hai raccontato un progetto andato male, ma il tuo messaggio è che volevi fare
diversamente e qualcuno ti ha detto di fare così. Quindi non è un tuo fallimento.'' È come
rispondere a ``qual è la tua debolezza?'' con ``sono così attento ai dettagli che a volte\dots'':
si fa passare un punto di forza per una debolezza. Il relatore vuole un esempio in cui \emph{tu},
senza altri fattori, hai sbagliato, e come ti sei comportato. Gli dice molto su come quella persona
reagirà ai suoi feedback: se ha davanti ``un muro'' che dirà sempre ``non è colpa mia'', non è adatta
a un ambiente in cui il feedback è alla base di ogni comunicazione. Chi non ammette di aver mai
sbagliato in vita sua, lo scarta.

\paragraph{È possibile colpevolizzarsi troppo per un fallimento?} Sì, assolutamente. Rientra molto
nella sindrome dell'impostore, o nell'aver sempre vissuto rapporti di lavoro tossici: ci si abitua ad
accettare che ci vengano urlate contro le cose e che sia sempre colpa propria, e ci si autoconvince
che lo sia.

\subsection{Al centro: l'affidabilità}

Superato questo livello si arriva all'\textbf{affidabilità}. Per essere un buon leader non bisogna
essere il migliore, ma affidabili: basta essere riconosciuti come bravi in qualcosa, non
necessariamente i migliori, e quell'affidabilità ce la si porterà anche nella leadership. Significa:
\begin{itemize}
  \item \textbf{rispettare la parola data}: se si promette ``parleremo di un aumento al prossimo
        slot disponibile'', se ne parlerà. Non ci si dimentica, e non si dice poi ``in realtà non si
        poteva fare, ma ho voluto provarci fino all'ultimo''. Una promessa va mantenuta;
  \item \textbf{impegnarsi per il tempo necessario} a raggiungere gli obiettivi;
  \item \textbf{saper rifiutare un lavoro} quando necessario: è il saper dire di no della mattina,
        ancora la matrice di Eisenhower (sezione~\ref{sec:z2-eisenhower}). ``Sono sempre gli
        stessi concetti.''
\end{itemize}

\subsection{In cima: la leadership, cioè dare l'esempio}

L'ultimo gradino è la \textbf{leadership}\footnote{Nella trascrizione ``dealership''.}, cioè
\textbf{dare l'esempio}. Non bisogna essere i più bravi; bisogna saper \textbf{offrirsi volontari
non solo quando il progetto è un successo sicuro, ma anche quando c'è il rischio di un
fallimento}. Sono tutti bravi a offrirsi per il gioco nuovo che andrà sicuramente bene.

\paragraph{Il rappresentante dei licenziati.} Al relatore è capitato di offrirsi per rappresentare
il team durante i \emph{layoff}, con il rischio di essere ostracizzato dall'azienda; con quelle
persone è ancora in contatto. Quando Blizzard fece i licenziamenti nella sede francese, vicino a
Versailles\footnote{Nella trascrizione ``i layoff di Versailles''. La sede europea di Blizzard si
trovava a Vélizy-Villacoublay, nei pressi di Versailles; il relatore non precisa l'anno.}, che
coinvolsero quasi tutto l'ufficio, lui non era tra i licenziati, ma venne \textbf{eletto
rappresentante sindacale} di chi lo era. Per lui è stato molto importante: pur non essendo parte in
causa, lo ritennero degno della loro fiducia. Per quanto ne sapevano, avrebbe potuto negoziare e
accettare tutto quello che proponeva Blizzard, perché non ne era toccato: che si negoziasse un mese
di liquidazione in più o in meno non gli portava un soldo. Avrebbe persino potuto usare la
situazione per negoziare un aumento per sé (``gli facciamo accettare tre mesi invece di sette e a te
diamo un aumento''): non era fuori discussione. Se si sono fidati è solo perché c'era una piramide di
questo tipo e lui aveva dato loro abbastanza motivi per farlo.

\paragraph{Autorevolezza, non autorità.} A uno studente che lo riassume così, il relatore risponde
che è proprio questo: un carisma diverso dall'essere simpatici, l'essere un punto di riferimento,
\textbf{più un pilastro di supporto che il motore} che fa andare avanti le cose. Non il motore davanti
a tutti, ma il \textbf{meccanico} che, quando si rompe qualcosa, sai che non ti frega: ti sistema le
cose senza ``farti la cresta'' sulle pastiglie dei freni.

\paragraph{Non abbandonare nessuno.} E soprattutto \textbf{non abbandonare chi è in difficoltà}.
Sembra una banalità, ma è alla base del motto dei Marines statunitensi, \emph{no one left behind}:
sapere che, se si viene feriti, verranno a prenderti anche sotto i missili è un forte elemento
motivazionale per correre quel rischio. In azienda è il motivo per cui il relatore ribadisce sempre
che, se qualcuno è in difficoltà, la sua prima opzione non è licenziarlo ma \textbf{dargli una mano},
e lo ha fatto a volte anche all'eccesso, per far vedere agli altri che quando saranno in difficoltà
non li lascerà in mezzo a una strada. A volte ce l'ha fatta: ha esempi di persone ``a tanto così''
dal licenziamento recuperate per i capelli e diventate ottimi \emph{performer}. Altre volte la
persona non era recuperabile e alla fine è stata licenziata. Ma chi attraversa un periodo difficile
(il \emph{burnout}, la depressione, una relazione finita) sa che la prima risposta non sarà ``ci sei
servito finora, ora puoi andartene'', ma ``che cosa possiamo fare?''. E questo fa lavorare tutti
tranquilli.

% ---------------------------------------------------------------------
\section{Prendersi cura del gruppo}
\label{sec:z2-gruppo}

\subsection{La soluzione: più leader possibile}

In un caso del genere manager e leader coincidono, ed è l'ideale: un manager dovrebbe essere anche
un leader. Ma la vera soluzione è \textbf{avere più leader possibile}. Se l'obiettivo del leader è
prendersi cura degli altri, e il manager dovrebbe essere ``il primo tra i leader'', ancora meglio è
che \textbf{tutti siano leader del proprio gruppo}, perché la leadership non deriva da un titolo né
dall'essere il manager. Il relatore ha visto tanti gruppi andare avanti perché qualcuno al loro
interno era il leader, non perché lo era il manager.

Vuol dire mettere in pratica tutto il discorso della prima giornata su reciprocità, gentilezza,
aiuto e collaborazione. Si sa che, nel peggiore dei casi, c'è il manager, ma in condizioni normali ci
sono i colleghi: e allora si va da chi è in difficoltà a dire ``posso darti una mano?''. Così, quando
saremo noi in difficoltà, magari senza nemmeno rendercene conto, sarà qualcun altro a dirci ``ti vedo
in difficoltà, come posso aiutarti?''.

\paragraph{L'esempio è contagioso, nel bene e nel male.} Quando si ha una posizione anche solo di
rispetto agli occhi degli altri, tutto ciò che si fa ha effetto sugli altri. In negativo vale la
\textbf{slippery slope}\footnote{L'esempio del cartone abbandonato ricorda la cosiddetta ``teoria
delle finestre rotte'' (\emph{broken windows theory}), secondo cui i piccoli segni di degrado
tollerati favoriscono comportamenti peggiori.}: se si accetta un cartone di pizza abbandonato per
strada, domani se ne accetteranno due, dopodomani tre. Ma vale anche in positivo: chi vede qualcuno
fare qualcosa di positivo pensa ``lo voglio fare anch'io''. ``Come fai una cosa, fai tutto'', dice un
proverbio che il relatore attribuisce alla tradizione cinese\footnote{La massima (``how you do
anything is how you do everything'') è attribuita a fonti diverse, spesso alla tradizione zen;
l'attribuzione è incerta.}.

\subsection{Il villaggio di oggi}

Si torna al villaggio del ``momento Piero Angela'', con i difensori e le bestie (reali o no) intorno.
Oggi fuori dal gruppo non ci sono più tigri, ma le \textbf{deadline}, l'\textbf{overwork}, i
\textbf{layoff}, il manager che urla, le aspettative, i \textbf{KPI}. Avere manager che si prendono
cura del gruppo è un buon inizio, ma è ancora meglio avere \textbf{il gruppo che si prende cura del
gruppo}. Non bisogna più limitarsi a pensare che sia responsabilità del manager prendersi cura di
noi: è responsabilità nostra prenderci cura del gruppo di cui facciamo parte.

Che cosa serve? La piramide vista prima. Serve essere percepiti come \textbf{autentici}: il
relatore ormai riconosce abbastanza facilmente chi si vuole ergere a leader perché gli piace
\emph{sentirsi} leader, non perché vuole \emph{fare} il leader, e quella mancanza di genuinità non
aiuta. In generale bisogna \textbf{generare sicurezza} nel resto del gruppo.

\subsection{Smettere di dire ``ci penserà qualcun altro''}

Un esempio concreto è smettere di dire \textbf{``ci penserà qualcun altro''}. Nello sviluppo è
evidente: se si trova un bug che non riguarda il proprio lavoro, tutti sono d'accordo che vada
comunque segnalato, perché altrimenti potrebbe arrivare in \emph{live}, nella \emph{release}.

\paragraph{Le lasagne di PG Esports.} Circa un anno e mezzo prima, quando lavorava ancora a PG
Esports (che organizza anche eventi dal vivo), l'azienda ospitava in ufficio dei potenziali sponsor
venuti a conoscere il team per decidere se sponsorizzare la Milano Games Week. C'era un catering, con
i classici buffet. In mezzo alla confusione (una ventina di persone di PG Esports e una trentina tra
manager, sponsor ed \emph{executive}), il relatore vede per terra, in mezzo alla stanza, una mezza
porzione di lasagne abbandonata. Senza pensarci la raccoglie per portarla via: non è sua, ma è un
piatto in mezzo alla stanza. Nel tragitto fino al bancone del buffet, almeno altri sei tra executive e
sponsor che non lo conoscevano gli consegnano il loro piatto non finito, come a un cameriere: nella
loro testa l'unico motivo per cui qualcuno potesse raccogliere un piatto da terra era che fosse il
suo lavoro. Sul momento non ci fa caso; arrivato al bancone si gira e vede tutto il suo team che ride
della scena.

Torna da loro e ne nasce una discussione: ``avete visto quanto è stato educativo?''. Avevano visto
esattamente come ragionano i gruppi che non sono abituati a pensare ``adesso ci penso io'', ma ``ci
penserà qualcun altro''. In quelle persone non c'era malizia: sicuramente era un cameriere, perché
altrimenti farlo? ``Non è il mio lavoro, lo avrà fatto perché è il suo.'' E diventa un problema quando
la stessa tendenza (``non c'entra con il mio lavoro, ci penserà qualcun altro'') si applica a
qualsiasi altra cosa. Il relatore si aspetta che \textbf{tutti raccolgano quel piattino}: che chi vede
un bug, se non può correggerlo lo segnali, e se può correggerlo lo corregga anche se non è il suo
lavoro, ``perché è un piattino in mezzo alla stanza''.

\paragraph{Da reattivi a proattivi.} Significa passare da un atteggiamento \textbf{reattivo}, cioè di
fatto passivo, a uno \textbf{proattivo}: prevedere i problemi e cercarli intorno a sé per risolverli,
invece di aspettare che esplodano (``finché non esplode non è un problema mio''). Ragionare nell'ottica
del ``e se succedesse, e io me ne fossi accorto senza fare niente?''. Uno studente lo chiama
\emph{mindset da tester}: il relatore si aspetta che, se si vede qualcosa, ci si assicuri di non essere
gli unici ad averlo visto. A volte basta farlo presente: il relatore non potrebbe mai correggere un
bug, quindi se lo trova lo segnala, e va bene così. Ma non bisogna dare per scontato che ``è un bug così
evidente che qualcuno l'avrà già segnalato''. È il paradosso dell'incidente stradale: \textbf{più
l'incidente è grosso, più è probabile che nessuno chiami l'ambulanza}, perché tutti pensano che
qualcuno l'abbia già fatto\footnote{È il cosiddetto \emph{effetto spettatore} (\emph{bystander
effect}), studiato in psicologia sociale a partire dagli anni Sessanta.}.

\subsection{Domande sul ``me ne occupo io''}

\paragraph{Se correggo un bug, devo avvisare il collega?} Dopo una breve pausa per problemi audio
(``facciamo \emph{double testing}''), uno studente chiede se, correggendo il lavoro di un collega, il
dovere sia poi andare a dargli il feedback. A prescindere bisogna informare qualcuno; e la
conversazione dovrebbe essere: ``ho trovato questa cosa, vuoi che te la sistemi io?''. Correggere
senza dirlo a nessuno potrebbe essere un danno: magari era \emph{working as intended}, o qualcuno ci
stava già lavorando.

\paragraph{E se divento ``quello che risolve quel problema''?} Uno studente porta la sua esperienza:
in un'azienda in cui ha lavorato tre anni, risolvendo un problema facile che non era di sua
competenza è diventato, di fatto, l'incaricato di quel problema, distolto dalle sue mansioni. È
un'arma a doppio taglio. Per il relatore lì il problema è del \textbf{management}. In alcune aziende si
applica il criterio ``chi trova il problema lo risolve'' (``se te ne sei accorto tu, pensaci tu''), ed
è un altro modo di scaricare la responsabilità sugli altri. Non è ciò che dice lui: deve essere chi
trova il problema a valutare se è giusto e positivo che lo risolva lui. Il cartone di pizza lo
raccoglie perché ci mette dieci secondi; ma se ha trovato un bug nel sistema di pagamento che richiede
tre settimane di lavoro, e chi l'ha introdotto fa finta di niente, non deve risolverlo lui: deve far
presente che c'è un problema, che è stato ignorato e che va risolto (``pensateci voi, ma io vi dico
che c'è''). Va fatto caso per caso: se è immediato, facile e nel proprio interesse, lo si risolve.
Ma il relatore sconsiglia alle aziende la regola ``se lo trovi lo risolvi'', perché \textbf{disincentiva
a trovare i problemi}: qualcuno potrebbe non segnalarli per non avere altro lavoro e aspettare che li
trovi qualcun altro, ed è negativo quanto la situazione descritta. Per chiudere, ne discuterebbe con il
management: ``possiamo rivedere questa policy? Incentiva l'azione sbagliata''.

\paragraph{E se ci dicono di non mettere mano?} Uno studente che lavora come tester racconta che spesso
potrebbero correggere certe cose, ma viene loro detto di non toccarle. Esatto: si comunica. ``Ho trovato
questa cosa, ne siete consapevoli? Sì? Ci state già lavorando? No? Posso farlo io? Sì? Allora lo
faccio.'' Non bisogna esagerare con l'iniziativa: anche qui vale il \emph{default to trust}. Si parte
dal presupposto che qualcuno potrebbe averlo già visto e starci già lavorando, quindi si chiede e si
offre il proprio supporto.

% ---------------------------------------------------------------------
\section{Se si diventa manager}
\label{sec:z2-manager}

\subsection{Meriti, colpe e \emph{big picture}}

Una breve digressione per chi, un giorno, diventerà manager: che cosa cambierebbe? Oltre a tutto
quanto detto, ci sono il discorso su \textbf{meriti e colpe} e quello su \textbf{progetto e task}:
\textbf{delegare i meriti e prendersi le colpe}, e pensare al progetto più che ai singoli task, cioè
alla \emph{big picture}, invece che al ``voglio tornare a fare le cose che sapevo fare''.

È un altro problema dei manager promossi dal gruppo di chi sapeva fare bene un certo lavoro: quando
hanno l'occasione di tornare a farlo, lo fanno molto volentieri. ``C'è un task da fare, lo faccio io'',
magari con la scusa di dare l'esempio. No: si sta andando nella propria \textbf{comfort zone}. Il
relatore lo capisce (anche lui tornerebbe volentieri a fare ciò che faceva prima), ma bisogna accettare
che \textbf{quelle cose non le si deve più fare}.

\subsection{Il ruolo è più grande della persona}

Anche da manager, il ruolo che si ricopre per il gruppo è \textbf{più importante di sé}. Non importa
chi sia il community manager di Blizzard, di Riot o di un'altra azienda: conta ciò che rappresenta per i
giocatori in quanto community manager, non il fatto che sia lui piuttosto che un altro.

Prendersi cura del gruppo vuol dire applicare concetti molto semplici:
\begin{center}
  \begin{tabular}{@{}p{0.42\textwidth}p{0.48\textwidth}@{}}
    \toprule
    \textbf{Invece di dire\dots} & \textbf{\dots dire} \\
    \midrule
    ``Qui non si accettano errori.'' &
      ``Spiegami perché abbiamo sbagliato'': fammi capire perché è successo
      (sezione~\ref{sec:z1-default-trust}). \\
    \addlinespace
    ``Qui si fa come dico io.'' & ``Tu come faresti?'' \\
    \bottomrule
  \end{tabular}
\end{center}

\subsection{Il tetto della bravura del capo}

Il secondo punto è cruciale soprattutto quando si promuove il più bravo in un lavoro: se il criterio
noto è che si promuove il più bravo, il promosso lo sa, e la prima cosa che penserà è ``faccio lavorare
tutti come me, visto che sono il più bravo''. Il relatore chiede alla classe: che cosa si nega l'azienda
con questa dinamica? Anche ammesso che sia davvero il più bravo, se tutti fanno come lui diventeranno
tutti più bravi, no?

Il problema è che \textbf{si fissa come tetto massimo la bravura di quella persona}: se tutti fanno solo
come dice lui, si arriverà al massimo al suo livello, mentre ognuno ha il proprio metodo. Da qui quanto
detto il giorno prima sulla tendenza dei manager ad assumere solo persone peggiori di sé: se il criterio
di promozione è la bravura, si prenderanno sempre persone un po' meno brave per evitare che diventino il
proprio capo. Se invece le carriere sono divise (``se sei bravo fai questo, se vuoi fare il \emph{people
manager} fai quest'altro''), il manager è incentivato a prendere persone \textbf{migliori di sé}; e se,
come detto il giorno prima, il bonus dipende anche dalle performance del gruppo, è nell'interesse di
tutti che il gruppo sia pieno di superstar, un ulteriore incentivo ad assumere i più bravi.

\paragraph{L'opzione C.} C'è poi un'altra cosa, forse minima, ma per il relatore molto importante: si
perdono le \textbf{occasioni di innovare}. Nelle aziende precedenti era visto come il punto di riferimento
per il \emph{problem solving}: in un'emergenza, se c'era da scegliere tra A e B, ``decide Stefano''. Andava
benissimo, ma la tendenza era che a ogni emergenza venivano da lui. Così ci si negava \textbf{l'opzione C}:
quella a cui lui non avrebbe pensato e qualcun altro sì. Magari la prima volta sbaglia, la seconda anche,
e la terza ha un'idea migliore della sua, proprio perché lui non è più l'esperto di quel settore. Fissare
tutto come ``al massimo decide Stefano'' significa non accettare altre idee: un limite molto grosso
all'innovazione, all'aggiornamento e al fare le cose in modo più moderno e flessibile.

\subsection{Celebrare i successi e i fallimenti}

\paragraph{Celebrare i successi del team.} Il relatore vuole assicurarsi che tutti capiscano quando si è
ottenuta una vittoria o un successo: come detto il giorno prima, se non si sa che cosa si fa bene si
rischia di smettere di farlo. ``Questa cosa l'abbiamo fatta bene, complimenti, bravi tutti.''

\paragraph{Celebrare i fallimenti.} Sui fallimenti si era spinto ben oltre il semplice accettarli: aveva
indetto \textbf{meeting in cui si celebravano i fallimenti}, con le persone che raccontavano i propri errori
per distribuire quella conoscenza al resto del team. ``Non mi fa arrabbiare l'errore fatto, mi fa arrabbiare
lo stesso errore fatto quindici volte.'' Condividere un fallimento con il gruppo non significa ``guardate
questa persona sciocca che ha fatto questo errore'', ma ``questa persona ci sta permettendo di evitare lo
stesso errore in quindici'': sta salvando l'azienda dal ripeterlo quindici volte, perché la gente ha paura di
dire agli altri che ha sbagliato. Mostrare dove si è sbagliato diventava quindi un incentivo, un servizio
all'azienda. E si torna al punto di prima: se il proprio ruolo per il gruppo è più importante di sé, non
importa dover alzare la mano e dire ``ho sbagliato, ho fatto così'', perché gli altri penseranno ``grazie,
così io non lo faccio''. Si dovrebbe essere lieti che qualcuno abbia fatto quell'errore prima di noi, e
andare a ringraziarlo per averlo condiviso.

\paragraph{Condividere i successi.} Infine bisogna \textbf{dare il merito a chi lo ha}. Al relatore è
capitato spesso di ricevere ringraziamenti come ``grazie a Stefano e al suo team per aver fatto questa
cosa'', e la sua risposta è stata ``grazie, ma ha fatto tutto Gianni''. Non bisogna avere paura di dire
``grazie del feedback, ma è stato lui'': è giusto ringraziare lui e il team (il relatore si sente parte del
team, chiaramente). Non si ottiene niente dal dire ``grazie, riferirò i complimenti al team'': il relatore
non ha mai creduto nel ``riferire'' che il management ha detto ``ottimo lavoro''. Chiede al management di
dirlo \textbf{direttamente} al team, perché sentirlo dire dal management ha un impatto diverso che sentirlo
riferire da lui. In questo senso, osserva uno studente, ogni gruppo è il gruppo di sostegno di se stesso:
``assolutamente sì''.

\subsection{Quanto ci vuole?}

Non c'è una scadenza. Il rispetto non si ottiene pagando per fare prima, ``con il \emph{boost} a livello
60''. Servono tempo ed energia, ed è difficilmente misurabile: il relatore non può dire quanto il team si
fidi di lui in questo momento. Lo percepisce da come lavorano, ma lo si vedrà davvero solo quando ci sarà un
problema grosso; fino ad allora è tutta una sua percezione. Il team lo sente ripetere queste cose ogni giorno
e lo vede metterle in pratica ogni giorno, e lui può solo sperare che il messaggio passi; quando ci saranno
davvero ``gatte da pelare'' capirà come è messo, perché non c'è momento migliore di una difficoltà per vedere
come reagisce un team.

\textbf{Essere leader è prendersi un impegno a lungo termine verso i propri colleghi}: non ci sono modi per
velocizzarlo. Bisogna fare il proprio lavoro ogni giorno ``a testa bassa'', aiutando gli altri, facendolo
vedere ed \textbf{esplicitando perché lo si fa}. Il relatore lo fa quotidianamente: ``ti sto aiutando; nel
caso avessi un dubbio, lo faccio perché il mio lavoro è aiutarti e io ne traggo vantaggio dal fatto che tu
lavori meglio. Sono qui per te, ma lo faccio anche per me: non pensare che sia solo gentilezza. È il mio
lavoro, ma sappi che lo faccio spontaneamente''.

\subsection{L'ultima slide}

L'obiettivo dovrebbe sempre essere lo stesso: tutti dovremmo poter finire la giornata, staccare e tornare a
casa sentendoci
\begin{itemize}
  \item \textbf{ispirati}, perché abbiamo fatto qualcosa che ha un effetto tangibile;
  \item \textbf{sicuri}, perché sappiamo di avere colleghi, un leader e dei manager che si prendono cura di
        noi;
  \item \textbf{realizzati}, per quello che facciamo e perché è ciò che ci piace fare.
\end{itemize}
E questo dovrebbe essere \textbf{un diritto naturale di tutti}, non un ``sì, ma solo il manager'', ``sì, ma
solo se sei un esperto'', ``sì, ma solo se sei un game director''. Ognuno ha le proprie capacità: bisogna
solo trovare ciò che piace e che si vuole fare.

% ---------------------------------------------------------------------
\section{Domande e risposte finali}
\label{sec:z2-qa-finale}

Finita la presentazione, si passa alle domande: ``chiedete tutto quello che volete, perché è
l'ultima possibilità''. Le domande sono riportate per temi.

\subsection{Leadership e organizzazione del lavoro}

\paragraph{La leadership riguarda tutti o solo chi vuole diventare leader?} Il discorso sui
manager vale solo per i manager; ma il consiglio è \textbf{essere tutti leader del proprio gruppo,
da dentro}. Essere un leader significa sapere che gli altri possono fare affidamento su di sé, e
di queste figure in un team ce ne possono essere infinite: non è una gara a chi è il leader.
``Siate tutti una di quelle figure a cui gli altri guardano quando sono in difficoltà.''

\paragraph{Da manager non si hanno più orari e si fanno straordinari di continuo?} No: il relatore
ha notato di essere molto più efficiente come manager quando si dà \textbf{paletti molto precisi}.
Se stacca alle sei e alle sei e cinque arriva un messaggio, risponde ``se non mi dici che è
urgentissimo, non lo guardo nemmeno, lo leggo domani''; se non è urgente, lo guarda il giorno dopo.
Ma gestisce le aspettative: ``se è urgente ci sono e ci sarò sempre, ma non aspettatevi che faccia
orari eccessivi solo perché sono manager, così come non me lo aspetto da voi''. Si aspetta che
ognuno faccia ciò che deve; se è urgente tutti faranno di più, ma se non lo è, no.

\paragraph{L'orario davvero flessibile.} Qui il relatore aggiunge una cosa che si era dimenticato di
dire parlando di ferie e permessi illimitati: introduce anche l'\textbf{orario davvero
flessibile}. Non il flessibile di molte aziende, in cui si può iniziare alle 9 e finire alle 18
oppure iniziare alle 10 e finire alle 19: se oggi si vuole iniziare alle 10 e finire alle 15, e si è
fatto tutto quello che si doveva, va bene così. Ci si prendono le ore libere in più per altro, per
staccare, per riposarsi, sapendo che se la settimana dopo bisognerà lavorare dalle 9 alle 21, lo si
farà (o si dirà ``ieri ho fatto dalle 9 alle 21, domani faccio dalle 10 alle 15''). L'obiettivo dato
è un lavoro e un progetto, non un orario: si lavora per completare i propri task e rispettare le
deadline. L'orario indicato è solo indicativo, ed è quello dei meeting: \textbf{nessun meeting prima
delle 10 o dopo le 17}. Se poi qualcuno vuole lavorare dalle 10 alle 11 e poi di notte, perché quel
giorno preferisce lavorare di notte, più rilassato, faccia pure: l'importante è che trovi il tempo per
riposare e staccare. Ci si autogestisce senza timbrare cartellini; l'unico accordo è ``hai un task, va
consegnato entro il giorno X e mi garantisci che venga fatto; come ci arrivi non mi interessa''.

\paragraph{Le è mai capitato di essere aiutato da qualcuno, in un momento di difficoltà?} Sì: come
detto, questo comportamento è \textbf{contagioso in senso positivo}. Se ogni volta che qualcuno è in
difficoltà qualcun altro lo aiuta, prima o poi capita anche a chi aiuta. Il relatore non lo dice
``tanto per dire'': ha avuto momenti in cui faceva davvero fatica ad alzarsi dal letto per lavorare,
perché esausto o demotivato per altri motivi, e il team gli ha detto ``stai a letto, ci pensiamo noi
oggi''. Questo lo aiuta ad aiutare, perché sa che la cosa è bidirezionale.

\paragraph{Tra le soft skill del workshop, da quale partire?} \textbf{Il feedback}: saperlo dare e
ricevere. È la più difficile, molto più del ``voglio aiutarti'', ma è anche quella su cui si costruisce
tutto il resto: attraverso il feedback si costruisce la \textbf{fiducia}, attraverso la fiducia il
\textbf{rispetto}, e il rispetto diventa la base portante di tutto il resto.

\paragraph{E se il leader se ne va?} A una domanda su che cosa fare quando viene meno la figura di
riferimento del gruppo\footnote{La domanda non è leggibile nella trascrizione; dalla risposta si
capisce che riguardava il caso in cui il leader che incarnava la visione del gruppo se ne va.}, il
relatore risponde che la domanda da farsi è un'altra: \textbf{è cambiato l'obiettivo?} La direzione è
sempre la stessa? Se sì, chi è la persona più adatta a portarci lì? Non conta tanto chi guida, ma
l'obiettivo che ci accomunava, di solito la visione di un leader in cui tutti si allineano: ci
crediamo ancora tutti, o è cambiato qualcosa? È il momento di fare autoanalisi: ``ero qui per quella
persona o per la visione?''. Si dice, ed è vero, che \textbf{le persone non si dimettono dai posti di
lavoro, ma da capi che non rispettano più}, perché spesso l'unica interfaccia con l'azienda è il
proprio capo: le dimissioni non sono un messaggio per l'azienda, ma per il capo (``non voglio più
lavorare con te''). Allo stesso modo, se non c'è più il leader della visione, bisogna prima capire se
si condivide ancora qualcosa.

\paragraph{Lavorare così è più faticoso?} Sì e no. ``Non voglio demotivarvi, perché non è il mio
lavoro'', ma non è tutto rose e fiori: lavorare in questo modo è anche più stancante
\textbf{emotivamente}, perché non si può più pensare ``faccio il mio, alle cinque ho finito e chi si è
visto si è visto''. Si sente un senso di responsabilità verso gli altri. Succede che si finisca il
proprio lavoro del pomeriggio, ci si giri e si veda un collega che finirà alle nove di sera, e
l'istinto sarà ``posso farti un pezzo del lavoro?'': così si finisce entrambi alle sei. In un certo
senso ci si accolla più lavoro, con la consapevolezza che quando si sarà in difficoltà ci sarà qualcuno
che si girerà a dire ``ti vedo stressato, posso darti una mano?''. È bellissimo, ma non dà solo benefici
diretti: richiede anche un investimento personale che non sempre le persone vogliono fare. Pro e
contro: è il modo che sceglierebbe il relatore, ma non è per tutti. Conosce persone che vogliono poter
staccare appena finito, e va benissimo: si cerca un'azienda in cui si possa lavorare così.

\subsection{Condividere fallimenti e successi in pratica}

\paragraph{Come si celebrano i fallimenti e si condividono i successi?} Condividere i successi sembra
più semplice (``bravo Mario che ha fatto questa cosa''). Il relatore dà uno spunto ancora migliore:
\textbf{``bravo Mario, e ora Mario ci racconta come l'ha fatto e com'è andata''}. Sono sottigliezze,
ma contano: raccontare in prima persona ciò che si è fatto dà un moto d'orgoglio che non dà sentirlo
raccontare da un altro.

Per i fallimenti, che sembrano dover partire dalla persona stessa, l'obiettivo era che tutti li
condividessero spontaneamente; ma all'inizio era difficile, perché non viene naturale alzare la mano e
dire ``ho sbagliato''. Per questo il relatore ha creato artificialmente dei \textbf{``festival del
fallimento''}, dal tono chiaramente ironico: durante una pausa pranzo, mentre tutti mangiavano insieme,
a turno si saliva ``sul palco'' a raccontare la stupidaggine fatta in settimana. Molto spesso iniziava
lui, raccontando i propri errori (``ho sbagliato questa cosa, sappiate che non si fa così''). Dopo aver
forzato l'interazione per una, due, tre settimane, alla quarta lo facevano da soli: era diventato
organico dire a tutti ``ragazzi, attenzione, ho fatto questa sciocchezza, non fatela anche voi'', e c'era
perfino un \textbf{tag apposito su Slack}. È stato il team a decidere come condividere i fallimenti: come
sempre, invece di imporre il modo, il relatore ha lasciato scegliere a loro, e funziona. È impegnativo
tanto quanto lo si vuole rendere tale.

\paragraph{Che tipo di fallimenti si condividono? Anche quelli personali?} Di solito si parte da quelli
\textbf{professionali}, cioè da ciò che si riconosce come utile anche agli altri: qualcosa che ha
rallentato, fatto perdere tempo o che rischia di ripetersi (``sappiate che, per evitare che lo facciate
anche voi\dots''). Il resto è più da conversazione privata, \emph{one to one}, con la persona con cui ci
si vuole sfogare, o con gli amici: se non ha un potenziale impatto sugli altri, non c'è un incentivo a
dirlo. E un errore che si è fatto per il proprio tipo di personalità? Qui tornano i test di personalità
della prima giornata: conoscendo il proprio profilo e quello dei colleghi, si sa che qualcuno ha lo
stesso carattere, e quindi lo si aiuta.

\subsection{Cultura aziendale e formazione}

\paragraph{Come si arriva a poter fare queste cose in azienda?} Conta molto la cultura aziendale. Il
relatore sa di lavorare bene perché oggi ha il \emph{background} e i successi alle spalle per dimostrare
che il suo metodo funziona, e quindi glielo lasciano fare; ma se dieci anni prima avesse proposto le
stesse cose, ``mi avrebbero mandato a quel paese''. C'è un momento e un posto per fare le cose, e non è
sempre possibile fare ciò che si ha in testa; ma ci si può lavorare. Per lui la domanda è diventata:
``sono convinto che così sarebbe meglio; come faccio ad avere carta bianca?''. \textbf{Portando
risultati}: si è concentrato su quello, e ci è arrivato progressivamente.

\paragraph{Come funziona la formazione con workshop come questo, in azienda?} Il relatore li fa in modo
un po' più rilassato ed esteso, con molte più interazioni ed esempi, perché in azienda si possono usare
esempi che si hanno davanti ogni giorno, e parlare di problemi e discussioni diventa molto più reale. Ma
in generale è sempre \textbf{formazione pratica}, ripetuta finché il messaggio non è recepito. Il segno
che qualcosa è stato recepito è quando arriva una persona nuova e \textbf{non è lui a spiegarle le cose}:
è il team stesso a spiegarle che cos'è la feedback culture, e lui non deve parlare. È la conferma che il
concetto viene perpetuato dal team.

\paragraph{Quanti manager dell'industria condividono questi principi?} Il relatore ha fatto molti
interventi all'IGDA di Roma e ne farà uno a Milano la settimana successiva (``il 24, mi pare''); interviene
spesso a workshop ed eventi, e il riscontro è sempre ``che bello, vorremmo farlo anche noi in Italia'', e ci
stanno lavorando. La moglie lavora in un'azienda che non ha niente a che fare con i videogiochi e lo stanno
già applicando. Magari in Italia è tra i primi a parlarne, perché il paese è ancora un po' indietro, ma non è
assolutamente l'unico: è più diffuso di quanto si pensi, non è una cosa di nicchia.

\paragraph{Questi spunti si applicano anche all'IT e alla consulenza?} ``Praticamente tutti.'' Si applicano
facilmente quando c'è un gruppo di persone su cui applicarli; se si cambia continuamente cliente è più
difficile portare questa mentalità da loro, ma questo non impedisce a chi lavora con noi di applicarla
direttamente con noi.

\paragraph{Ci si può tenere in contatto?} Su LinkedIn: il relatore accetta sempre tutti e risponde alle
domande quando riesce (ne riceve tante, ma cerca di stare al passo almeno su base settimanale). E, ``a costo
di sembrare quello che si fa pubblicità'', chi lavora in un'azienda che fa fatica ad applicare questi
concetti o non sa che esistono può metterla in contatto con lui: è spesso così che ha trovato nuovi modi di
applicarli. Vale anche per chi tra sei mesi lavorerà altrove. I suoi profili sono pubblici e non c'è nessun
problema a contattarlo.

\subsection{Performance e persone in difficoltà}

\paragraph{Se un'azienda è a corto di personale, conviene comunque eliminare chi performa poco?} Il giorno
prima si era detto che conviene separarsi presto da chi non performa (sezione~\ref{sec:z1-squadra}); ma non
si finisce per caricare di lavoro tutti gli altri? Per il relatore \textbf{eliminare non vuol dire non
sostituire}: vuol dire cambiare la persona con qualcuno magari più junior, ma più adatto come \emph{fit}
culturale. Il vecchio discorso del ``fare lo stesso lavoro con meno gente'' non è sostenibile: lo si può fare
fino a un certo punto, in caso di necessità, ma non è una soluzione. Non è ottimizzare, è far andare in
\emph{burnout} chi resta.

\paragraph{Ha più senso tenere una persona poco performante e non in sintonia con il team, piuttosto che
licenziarla senza poterla sostituire?} No: se una persona è poco in sintonia, performa poco ed è in generale
tossica e negativa per il team, \textbf{prima ci si separa, meglio è}. Questo non vuol dire non provare a
recuperarla, ma bisogna andare dritti al punto: ``abbiamo un problema: ti vedo poco performante, non ti vedo
allineato con il team, vedo che il team non si fida di te. Che cosa vogliamo fare?''. Se la risposta è ``non lo
so'', la strada porta dritta all'uscita. Bisogna \textbf{recidere i rami morti} il prima possibile, facendo
però attenzione a essere sicuri che si tratti di un ramo morto e non di un ramo recuperabile: tagliare un ramo
recuperabile mette negli altri la paura di essere i prossimi, ``alla prima foglia gialla''. Quello che si
vuole è che, quando un ramo è morto, la pianta smetta di alimentarlo, e ne crescerà un altro. Entrambi i
messaggi devono essere chiari: \textbf{non si subisce passivamente la situazione}, e allo stesso tempo
\textbf{non si abbandona nessuno} (se qualcuno è in difficoltà, la prima soluzione non è tagliare il ramo). È
un compromesso tra le due cose. (A questo punto arriva, puntuale, il crash della connessione: ``non era ancora
successo oggi''.)

\paragraph{Che cosa fare se qualcuno nel team sembra soffrire di depressione?} È un discorso complesso, perché
il relatore \textbf{non è uno psicologo}. Può dare la sua disponibilità a concedere permessi, coprire il lavoro
della persona e in generale darle modo di seguire un percorso con uno specialista: non può risolvere la
depressione, ma può alleviarne la parte lavorativa e permettere di ricevere supporto da chi è qualificato. È lo
stesso motivo per cui, a Blizzard e a Riot, quando arrivavano ticket con minacce di suicidio (``mi avete
bannato, quindi mi ammazzo'', ticket purtroppo abbastanza comuni), il \emph{customer service} era istruito a
non intervenire direttamente: non dovevano essere loro a fermare la persona, perché non qualificati, ma
indirizzarla a qualcuno di qualificato, prendendo tempo. Una persona non qualificata che prova a farlo rischia
di peggiorare la situazione; come davanti a un incidente si chiama l'ambulanza e non si fa una tracheotomia.
Per il relatore salute mentale e salute fisica sono la stessa cosa: come non si azzarderebbe a operare qualcuno
di appendicite, non si azzarda a fargli da terapeuta.

\paragraph{Come si danno feedback adatti a ogni tipo di personalità?} Bisognerebbe leggere i risultati del
profilo della persona sullo stesso sito su cui si è fatto il test di personalità della prima giornata, che
spiega anche come interagire con quel profilo. In generale: con chi è più \textbf{analitico}, portare dati a
supporto; con chi è più \textbf{emotivo}, metterlo in una condizione di comfort; con chi è molto
\textbf{pratico}, andare dritti al punto senza girarci intorno. Poi fare pratica. Oppure il modo ``perfetto''
per il relatore: provarci come viene e poi \textbf{chiedere} ``andava bene così o preferisci in un altro modo?
Potevo fare meglio?''. Bisogna creare anche quel rapporto: fidarsi di discutere anche di queste cose.

\paragraph{Se nel team fossimo quasi tutti leader, non ci sarebbero ``troppi galli''?} No, perché per leader si
intende chi si prende cura degli altri, e magari lo fossero tutti. Non significa che tutti scelgano che cosa
fare: spesso lo sceglierà comunque il manager, o guiderà la scelta. Ma avere tutte persone con iniziativa e
capaci di scegliere quando serve è qualcosa per cui il relatore ``metterebbe la firma''. Non ha mai corso il
rischio dei ``troppi galli'': se ci sono quei caratteri, il team è stato assemblato male, perché non tutti sono
compatibili con questo modo di lavorare e magari non si hanno le persone giuste per applicarlo.

\paragraph{Se la fiducia richiede tempo ed esperienza, quanto è difficile per un junior?} Meno di quanto si
pensi: proprio perché su un junior non ci sono aspettative, un junior che chiede ``posso darti una mano?'' dà un
segno chiaro che tiene più al team che a se stesso. Si può essere leader anche non dal punto di vista
professionale, ma come chi va a dare supporto agli altri: ``il medico del team, l'\emph{healer}'', e lo si può
fare anche da junior. È così che ha iniziato il relatore: andare ad aiutare a risolvere problemi non suoi gli ha
fatto accelerare molto il percorso per ottenere rispetto e stima, e lo ha esposto a cose a cui altrimenti non
sarebbe mai arrivato. È una questione di carattere: se ci si tiene agli altri, lo si fa vedere, e il resto
arriva come risultato. ``È quasi un reddito passivo'': un investimento nel far vedere che ci si è, e il resto
arriva (``Stefano si è preso cura di me, quindi ora mi prendo cura di lui'').

\paragraph{Aiutare gli altri solo dopo aver finito il proprio lavoro?} Dipende dall'urgenza. Il relatore spesso,
se qualcosa è davvero urgente, mette in pausa ciò che sta facendo, risolve e poi torna al suo lavoro; ma è troppo
\emph{case by case}. Se due colleghi litigano su chi debba avere le ferie di Natale, lo si può risolvere anche
domani; se litigano perché si è perso il \emph{repository} e domani c'è la consegna, lo si risolve subito.

\paragraph{Come mantenere motivato un team su progetti non retribuiti, come quelli accademici?} Essendo il più
\textbf{trasparenti} possibile nei processi decisionali: se si deve decidere se fare un action o un RPG, tutti
devono sapere perché si fa l'uno invece dell'altro (da qui il metodo di brainstorming della mattina). Deve essere
sempre chiaro a tutti perché si stanno facendo le cose in quel modo, perché \textbf{la demotivazione nasce dal
disallineamento} rispetto all'obiettivo del gruppo, e l'unico modo di essere allineati è essere sempre
consapevoli e partecipi del processo che decide dove sta andando il gruppo. E soprattutto: \textbf{non chiudere
mai un meeting senza aver concordato che cosa si vuole fare} e senza essere tutti d'accordo. Se non lo si è,
bisogna continuare a parlarne, per quanto scomodo, lungo e noioso: si deve sempre uscire da un meeting
sapendo che si è tutti d'accordo, o al massimo in \emph{disagree and commit} (sezione~\ref{sec:z2-disagree}).
``Quello è il massimo accettabile.''

\subsection{Carriera, colloqui e candidature}

\paragraph{Come ha iniziato a lavorare all'estero?} Negli annunci online si trovano spesso requisiti come un titolo
di studio conseguito nel paese o l'abilitazione a lavorarci. Questo vale per i paesi fuori dall'Unione europea: per
lavorare negli Stati Uniti o in Sud America serve un visto o una \emph{green card}. In molti casi, però, è proprio
ottenere il lavoro a permettere di lavorare lì: quando il relatore doveva andare a lavorare in California (a Irvine,
a Santa Monica, dove capitava), era l'azienda a occuparsi dei permessi necessari, perché ``per venire qui devi avere
un lavoro, e il lavoro te lo diamo noi''.

Il suo inizio è stato abbastanza semplice. In Italia lavorava per Digital Bros, occupandosi dei giochi online gestiti
dall'azienda e di quello che era il canale Game Network sul digitale terrestre\footnote{Questo passaggio chiarisce il
dubbio segnalato nel capitolo~\ref{cap:z1-feedback} sul ``Digital Game Network Channel'': si tratta di Game Network,
canale televisivo dedicato ai videogiochi legato al gruppo Digital Bros.}. Vide un annuncio a Dublino per
\emph{Warhammer Online}, mandò la candidatura, fece i colloqui e gli dissero ``ti paghiamo la \emph{relocation}, vieni
qui''. Come detto il giorno prima, sono le aziende a occuparsi del trasferimento: se devono creare a Dublino un'unità
che segue dodici lingue europee, sanno di dover far arrivare persone da fuori, ed è già a budget. Non è nemmeno una
richiesta da fare: è prevista. Per il relatore, da questo punto di vista, è indifferente lavorare per un'azienda
italiana o estera: se si è adatti, si è adatti. Per questo insisteva sull'inglese: se lo si parla, è indifferente di
che cosa ci si occupi e quale sia la propria lingua madre, a meno che la posizione non sia esplicitamente ``community
manager italiano''. Un programmatore polacco o italiano non cambia niente. Lavorare all'estero non è più difficile: c'è
più concorrenza solo se l'azienda cerca attivamente persone da tutto il mondo e non si pone problemi di \emph{location}.
Cambiano l'azienda, il fatto di lavorare da remoto o dal vivo, l'eventuale trasferimento; per il resto non cambia
molto.

\paragraph{Gli anni passati nell'IT contano se si vuole passare ai videogiochi?} Dipende da che cosa si vuole fare: se
da responsabile IT si vuole diventare programmatore, bisogna saper dimostrare di saper programmare. Vale anche per i
passaggi interni: se un game designer vuole diventare art director, al relatore va bene, deve solo dimostrare di saper
fare l'art director. Anzi, incentiva le \textbf{promozioni interne}, perché si conosce già il candidato; ma gli
spostamenti non sono liberi ``a sentimento'': bisogna saper fare quella cosa. I recruiter valuteranno non tanto gli anni
quanto ciò che si sa fare; gli anni servono comunque, perché dimostrano di saper lavorare in un ambiente professionale,
che è un altro insieme di competenze. Ma se si sa fare ciò che chiedono, aver lavorato nell'azienda X o Y non dovrebbe
cambiare molto.

\paragraph{Quanto vengono considerate le candidature spontanee?} Non tantissimo. Soprattutto per le aziende non grandi,
se non stanno cercando qualcuno è perché non se lo possono permettere: una candidatura spontanea non le convincerà ad
assumere. Più l'azienda è piccola, più conviene aspettare l'annuncio, perché se non c'è vuol dire che non c'è budget, e
potrebbero rispondere ``ti vorremmo, ma non possiamo prenderti''. C'è anche il rischio che finisca nel vuoto (l'HR magari
non controlla nemmeno la mail, perché non sta assumendo) o, peggio, che rispondano ``non ci servi'' e una settimana dopo
esca la posizione, con il candidato già nel cestino dei rifiutati. Si rischia di \textbf{bruciarsi} la possibilità.

\paragraph{Come ci si accorge di essere riconosciuti come leader, e come lo si ``vende'' a un colloquio?} Uno studente
osserva che la leadership, così descritta, non la si riconosce da sé ma la riconoscono gli altri. Internamente il
relatore non ha mai dovuto ``venderla'', perché parlavano i risultati. Finché non lo si vive è difficile capirlo, ma
quando si arriva alla posizione in cui il team rispetta e vede qualcuno come punto di riferimento, non ci sono grandi
dubbi: è evidente che vengono da lui. Si può riflettere su \textbf{quanto spesso gli altri vengono da noi}, cercano di
emulare il nostro modo di fare, vanno a rivedere le nostre cose per capire dove hanno sbagliato, chiedono un parere, ci
parlano di un problema: la somma di tutto questo significa molto spesso che si fidano. Per ``venderla'' all'esterno si
può dire di essersi occupati di un team in cui, attraverso il lavoro e il rispetto degli altri, si è arrivati a essere
visti come punto di riferimento; ma il suggerimento è \textbf{portare esempi tangibili} (``c'era questo progetto, non
sapevamo come risolverlo, il team mi ha chiesto una mano e l'ho risolto così''), invece di ``incensarsi'' dicendo ``sono
il punto di riferimento del team''. Se il relatore racconta di essere stato eletto rappresentante dei licenziati senza
essere tra loro, un'azienda capisce che si fidavano di lui, punto: non ci sono altre conclusioni da trarre.

\paragraph{Quali sono gli errori più gravi dei junior ai colloqui, a livello di soft skill?} In ordine di priorità:
\begin{enumerate}
  \item \textbf{non arrivare preparati}, cioè \textbf{senza esempi}. Il relatore ha sempre un esempio pronto per ogni
        cosa di cui vuole parlare: per la fiducia, per gli \emph{achievement}, per i risultati. Servono sempre esempi
        tangibili per far capire che cosa si intende;
  \item \textbf{non avere domande alla fine}, perché dimostra di non avere chiaro che cosa si andrebbe a fare. Anche solo
        chiedere che cosa si farebbe, che cosa si può migliorare, o ``che impressioni hai avuto finora dal mio
        colloquio?'': si possono fare domande molto aperte al recruiter per capire come sta andando;
  \item \textbf{non leggere bene la \emph{job description}}: presentarsi con vaghezza, senza aver capito che cosa cerca
        l'azienda. Lo si impara anche con la pratica; ma il designer che si candida per una posizione da programmatore,
        senza nulla nel CV che faccia pensare che sappia anche programmare, il relatore non solo lo scarta: lo mette in
        una \textbf{blacklist} di persone che forse non considererà mai. Se non si riesce a interpretare un testo come una
        \emph{job description}, è un problema troppo grosso da superare in una collaborazione;
  \item \textbf{errori di battitura e CV impaginati male}, che molti sottovalutano. Il relatore ha letto letteralmente,
        in una lettera di presentazione per una posizione da \emph{copywriter}, la frase ``sono una persona molto attenta
        ai dettaggi''. Come si può considerare quella persona per quel ruolo? Non ha nemmeno incollato la lettera in Word
        per il controllo automatico: non è un errore sfuggito al correttore perché la parola esiste con un altro senso,
        ``dettaggi'' viene segnato in rosso in qualsiasi lingua. Come affidarle la comunicazione aziendale? Vale per tutti
        gli errori: ciò che si manda è \textbf{il primo biglietto da visita} con l'azienda, e non si manda un biglietto da
        visita sgrammaticato. Un conto è non andare ai colloqui in giacca e cravatta (``non l'ho mai fatto e probabilmente
        non lo farò mai''), un altro è presentarsi con errori;
  \item un altro esempio comune: \textbf{l'email sbagliata nel CV}. Se la mail torna indietro, il relatore non telefona
        per farsi dare quella corretta: se non si è in grado di scrivere giusta la propria email, non vale la pena. Più
        l'errore è banale e facilmente evitabile, più lo fa arrabbiare (``scusate se mi sono inalberato''): evidentemente
        il candidato non teneva abbastanza a quel posto.
\end{enumerate}

\paragraph{Quanto è importante la lettera di presentazione?} Sempre di più: sempre più aziende la richiedono come
obbligatoria, quindi conviene prepararne una, sapendo di doverla adattare al destinatario. Deve essere \textbf{meno lunga
di una pagina}, in \textbf{tre paragrafi}:
\begin{enumerate}
  \item \textbf{chi siete}: un minimo di \emph{backstory}, che cosa avete fatto, che cosa volete fare, da dove arrivate;
  \item \textbf{perché avete mandato il CV}: che cosa vi ha interessato di questa azienda;
  \item \textbf{perché dovrebbero chiamare voi}: che cosa dovrebbe spingere l'azienda a scegliere proprio voi.
\end{enumerate}
Uno studente dice di riuscire bene con i primi due, ma di trovare difficile il terzo, che è poi la domanda che tutte le
aziende fanno. In realtà basta \textbf{collegarsi alla \emph{job description}}: se cercano un 3D artist in grado di fare
X, Y e Z, si conferma di averlo letto e si spiega che con la propria esperienza si è imparato a fare proprio quelle cose.
Non vogliono sentire niente di particolare, ``non devi cambiargli il mondo''.

\paragraph{Che cosa guarda un recruiter in un portfolio?} Il relatore li guarda poco: fa un \emph{sanity check} generale,
per vedere che non sia evidente perfino a lui, ``a digiuno di tante cose tecniche'', che c'è un problema (se se ne accorge
perfino lui, il problema c'è). Per il resto guarda che il portfolio sia \textbf{coerente con il CV}: se arriva un CV da
senior con dodici progetti completati e anni di lavoro, e il portfolio sono quattro immagini, non è coerente. Vuole una
concordanza tra ciò che si racconta e ciò che si vede.

\paragraph{Che strumenti per creare un CV?} \textbf{Canva}, che ha un catalogo di template gratuiti e a pagamento
amplissimo e personalizzabile: anche se non piace, lo si prende e lo si adatta, ma fa già il 90\% del lavoro. Insieme a
\textbf{ChatGPT}, perché ``è inutile far finta che non esista'': ci si fa scrivere la lettera di presentazione e poi la si
corregge e la si adatta, e lo stesso per il CV. L'importante è renderli adatti e coerenti, non lasciarli solo frutto di
ChatGPT; ma per fare il grosso del lavoro è un ottimo strumento.

\paragraph{Conviene lasciare nel CV esperienze non del settore?} Uno studente ha alle spalle il liceo, vari stage e
lavori che non c'entrano niente con la modellazione 3D. Si possono mettere finché c'è spazio: il relatore consiglia
\textbf{al massimo due pagine}, e quando ha cominciato a sforare nella terza ha iniziato a togliere ciò che non era più
rilevante. Ormai nel suo CV la formazione è un paragrafo: ha così tante esperienze lavorative che non interessa a nessuno
che abbia fatto il liceo scientifico o classico, e non lo mette nemmeno più. Si filtra in base a ciò che è più rilevante
per lo spazio disponibile.

\paragraph{Che cosa fare con le aziende che dicono ``ti faremo sapere'' e poi spariscono?} Non aspettare uno o due anni
per scrivere. Il suggerimento: tra le domande finali del colloquio, l'ultima del relatore è sempre \textbf{``tra quanto
mi fate sapere?''}. Se rispondono dieci giorni, all'undicesimo scrive: ``mi avete detto dieci''. Non si sente in difetto a
mettere pressione, perché sono stati loro a dare la scadenza. E se viene ignorato del tutto, quell'azienda per lui è
``morta'': non si ricandiderà mai, perché non vede perché il rispetto debba andare in una sola direzione.

\paragraph{Quanto essere trasparenti, nel CV e nella lettera, su problemi personali come l'ADHD?} Uno studente con
ADHD\footnote{Nella trascrizione ``essendo 10D''; nel resto della risposta ``di'' e ``dei 10D'' stanno per ADHD.} dice di
tendere ad avere performance altalenanti. Il relatore lo invita a \textbf{non usare l'ADHD come scusa} per essere
altalenante: anche lui ha l'ADHD, ma non è altalenante. Esserlo è senz'altro un \emph{malus}, ma esistono i mezzi per non
avere problemi di questo tipo: lui ha trovato i suoi metodi (la to-do list, la \emph{must do list}, il delegare il lavoro
al team, perché occuparsi delle cose per gli altri lo mantiene molto motivato, sveglio, all'erta, attento a come vanno le
cose). Sta anche a ciascuno trovare un lavoro in cui compensare ciò che viene richiesto: se domani lo rimettessero a fare
il programmatore, come all'inizio, probabilmente impazzirebbe, perché non riusciva a restare concentrato sul codice oltre
un certo tempo. Lì ha capito che non poteva fare il programmatore, ma ha trovato qualcos'altro che gli permette di essere
com'è e fare ciò che fa. Bisogna capire i propri punti di forza e di debolezza e giocare su quelli: si possono avere tutti
i difetti del mondo e voler comunque fare una cosa, ma bisogna accettare che sarà più difficile.

Quanto alla trasparenza, il consiglio è \textbf{non mentire mai nel CV}, ``perché tanto vi beccano''. E per il discorso della
\emph{bandwidth}: tutte le energie spese a mentire, e a tenere viva la bugia, sono energie tolte allo spiegare bene ciò che
si sa fare. È un investimento nel posto sbagliato. Bisogna concentrarsi su ciò che si sa fare, con esempi tangibili, e dire
``non lo so fare'' quando non lo si sa fare. Il peggio non è essere scoperti, ma \textbf{non esserlo}: se si finge di sapere
il C++, si viene assunti e messi a lavorare in C++, che cosa si fa il giorno dopo? Si rischia davvero di rovinarsi la
carriera. Se non si è un \emph{fit}, si sarà il \emph{fit} per la prossima. L'incubo del relatore è essere assunto in un
posto dove pensano che sappia fare qualcosa che non sa fare; anzi, tende a esplicitarlo: ``cercate questa cosa? Io non la so
fare, va bene lo stesso?'' --- ``Sì, te la insegniamo noi'' --- ``Voglio solo che sappiate che non la so fare''. Meglio
questo che vivere con l'ansia che prima o poi si accorgano che non sa programmare.

Gli capita spesso con chi mente sull'inglese. Al primo colloquio chiede ``sai parlare inglese? A che livello?'' --- ``C1,
\emph{fluent}, quasi madrelingua'' --- ``Perfetto'', e lascia cadere l'argomento. All'ultimo colloquio: ``ti va se lo facciamo
in inglese?'' --- ``No, ma io intendevo che\dots'' --- ``No, mi avevi detto \emph{fluent}: ora facciamo il colloquio in
inglese''. Prima o poi le bugie raggiungono chi le dice.

% ---------------------------------------------------------------------
\section{Chiusura del workshop}
\label{sec:z2-chiusura}

Il relatore chiede alla classe di far sapere com'è andata a Vanessa, il punto di contatto con l'accademia per il workshop:
sarà lei a riferirgli eventuali lamentele o suggerimenti, che gli permetteranno di fare meglio i prossimi. Ringrazia per aver
``retto per dodici ore'' e spera che sia stato utile, consapevole che tutti sono stanchi quanto lui.

Gli studenti rispondono che è stato chiarificatore su molte cose che non sapevano, soprattutto su come interfacciarsi nei
team, e che sarà materiale da rivedere per fissare bene i concetti, vista la mole in due giorni; qualcuno si aspettava che
sarebbe stato molto più pesante e l'ha trovato invece scorrevole. Il relatore ribadisce che \textbf{niente sostituirà
l'esperienza diretta}: il workshop dà un'infarinatura e un punto di vista, ma i concetti si metteranno in pratica solo nel
\emph{day to day}. Bisogna avere \textbf{pazienza con se stessi}, sapendo che non verrà tutto bene dal primo giorno.
Confessa di essere stato un po' preoccupato per la durata, ma il materiale c'era, e farlo in un giorno avrebbe rischiato di
trasformarlo in una corsa: ``rischiava di essere uno \emph{sprint}, e invece è una maratona''.

Infine ricorda che si può aggiungerlo su LinkedIn, che registrazione e slide (in forma integrale) verranno fatte avere
tramite Vanessa (``dite che vi ho autorizzato io a chiederle'') e che, per qualsiasi dubbio, resta a disposizione. ``Buon
resto del weekend.''

% ---------------------------------------------------------------------
\section*{Punti chiave del capitolo}
\addcontentsline{toc}{section}{Punti chiave del capitolo}

\begin{itemize}
  \item I colloqui sono una \textbf{selezione reciproca}. L'annuncio dice molto sulla cultura aziendale; alla fine del
        colloquio vanno sempre fatte domande (il \emph{day to day}, la ``bacchetta magica'', ``tra quanto mi fate
        sapere?''). Si arriva con una \textbf{cifra minima mentale} e con \textbf{esempi tangibili}; non si mente mai.
  \item Time management non è una lista di cose da fare: si parte dalla \textbf{must do list}, che è anche il primo
        momento per dire ``oggi potrei non farcela''; si stimano i tempi e si chiede quando qualcosa non è chiaro.
  \item Ogni nuovo task va \textbf{prima aggiunto alla lista}, non eseguito: priorità non è recenza. Le
        \textbf{checklist} fanno risparmiare tempo, anche (e soprattutto) ai più esperti.
  \item \textbf{First contact resolution}: raccogliere tutto ciò che serve in una sola richiesta, e non iniziare un task
        senza le informazioni necessarie. Ogni interruzione costa dieci-quindici minuti di concentrazione.
  \item \textbf{Matrice di Eisenhower}: fai, pianifica, delega, elimina. ``Non importante'' significa che non è importante
        che lo faccia tu; il quadrante dimenticato è \emph{delete}.
  \item Le pause vanno \textbf{scelte, non subite}; i task grossi vanno spezzettati; bisogna \textbf{saper dire di no},
        anche ai meeting (sono necessario? può andarci uno solo? ci sarà un resoconto scritto? per che cos'altro potrei usare
        questo tempo?). \textbf{Il multitasking non esiste}.
  \item A fine giornata si verifica la \emph{must do list}: un \emph{must do} rimandabile non era un \emph{must do}.
  \item Ferie illimitate con un minimo, giornate libere e feedback culture sono, per il relatore, le condizioni per un team
        sano; il benessere è anche mentale.
  \item Il \textbf{brainstorming genera idee, non decide}. Va preparato (topic, obiettivo, vincoli), separa la
        \textbf{generazione} dalla \textbf{discussione} (\emph{rapid ideation} di cinque minuti, \emph{dot voting}
        anonimo) e combatte l'\emph{anchoring} e il dominio degli estroversi.
  \item La decisione si prende con la \textbf{matrice impatto--sforzo}: prima le idee a basso sforzo e alto impatto, poi in
        diagonale. Lo sforzo sono le risorse, l'impatto comprende anche la motivazione; decide chi farà il lavoro.
  \item \textbf{Disagree and commit}: si dichiara il disaccordo, poi ci si impegna al massimo. Se va male ha sbagliato il
        team, non ``ve l'avevo detto''.
  \item \textbf{La leadership non è un titolo}: il rispetto si guadagna. Le carriere vanno divise tra \textbf{manager} e
        \textbf{subject matter expert}, pagati allo stesso modo, così che diventi manager solo chi vuole prendersi cura degli
        altri.
  \item Dopamina, serotonina e ossitocina legano il gruppo; il \textbf{cortisolo} è contagioso e lo disgrega. Il
        ``contratto sociale'' del villaggio vale ancora: chi protegge il gruppo riceve di più, chi ne abusa tradisce.
  \item La \textbf{piramide della fiducia}: integrità (aiutare davvero, onestà, ammettere gli errori), affidabilità
        (mantenere la parola, saper rifiutare), leadership (dare l'esempio, offrirsi anche quando si rischia, non lasciare
        indietro nessuno).
  \item Il gruppo deve prendersi cura del gruppo: \textbf{raccogliere il piattino}, smettere di dire ``ci penserà qualcun
        altro'', segnalare e chiedere prima di intervenire.
  \item Da manager: delegare i meriti e prendersi le colpe, lasciare la comfort zone, chiedere ``tu come faresti?'',
        celebrare successi e \textbf{fallimenti}, dare il merito a chi lo ha. Essere leader è un \textbf{impegno a lungo
        termine}.
  \item L'obiettivo: tornare a casa \textbf{ispirati, sicuri e realizzati}, come diritto di tutti e non privilegio di pochi.
\end{itemize}

%    % =====================================================================
%  Lezione zero 3 — Agile e Scrum nella pratica
%  Fonte: trascrizione della terza "lezione zero" (GD0_3): una lezione
%  introdotta da Diego Ricchiuti e tenuta da due ospiti, Simone Cicchetti
%  (ex studente, oggi level designer e producer) e Cristian Liberatore
%  (studente del terzo anno, game designer e project manager), che
%  raccontano come hanno applicato davvero Agile e Scrum nei loro progetti.
%  Pacchetti richiesti nel preambolo del libro:
%    \usepackage[italian]{babel}, \usepackage[T1]{fontenc},
%    \usepackage[utf8]{inputenc} (solo pdfLaTeX), \usepackage{booktabs}
%
%  Collocazione nel libro: parte "Lezioni zero", subito dopo
%  GD0_2_tempo_idee_e_leadership, con la stessa numerazione propria (0.3).
%  Nel file principale del libro, ad esempio:
%
%    \part{Lezioni zero: il game designer nell'industria}
%    \renewcommand{\thechapter}{0.\arabic{chapter}}
%    \setcounter{chapter}{0}
%    % =====================================================================
%  Lezione zero 1 — Lavorare in squadra: la cultura del feedback
%  Fonte: trascrizione della prima "lezione zero" (GD0), prima giornata di
%  un workshop di due giorni tenuto prima dell'inizio del corso di Game
%  Design, dedicato alla figura professionale e al lavoro in team.
%  Pacchetti richiesti nel preambolo del libro:
%    \usepackage[italian]{babel}, \usepackage[T1]{fontenc},
%    \usepackage[utf8]{inputenc} (solo pdfLaTeX), \usepackage{booktabs}
%
%  Collocazione consigliata nel libro: le lezioni zero precedono il corso,
%  quindi vanno in una parte separata PRIMA di GD1, con numerazione propria
%  (0.1, 0.2, ...) così che i capitoli del corso conservino i numeri 1, 2, 3.
%  Nel file principale del libro, ad esempio:
%
%    \part{Lezioni zero: il game designer nell'industria}
%    \renewcommand{\thechapter}{0.\arabic{chapter}}
%    \setcounter{chapter}{0}
%    \input{Seminari/GD0_1_cultura_del_feedback}
%    % ... altre lezioni zero ...
%    \part{Il corso di Game Design}
%    \renewcommand{\thechapter}{\arabic{chapter}}
%    \setcounter{chapter}{0}
%    \input{GD1/GD1_introduzione_al_game_design}
%    % ...
%
%  (Con hyperref aggiungere anche \renewcommand{\theHchapter}{z.\arabic{chapter}}
%  nella parte zero, per evitare ancore duplicate.)
%  Tutte le etichette di questo capitolo hanno il prefisso "z1:" per non
%  collidere con quelle dei capitoli del corso.
% =====================================================================

\chapter[La cultura del feedback]{Lavorare in squadra: la cultura del feedback}
\label{cap:z1-feedback}

Questo capitolo apre le \emph{lezioni zero}: incontri che precedono l'inizio
vero e proprio del corso di Game Design e che non riguardano una tecnica di
design in senso stretto, ma \textbf{la figura professionale} di chi lavora
nell'industria dei videogiochi. Prima di imparare a progettare un gioco
conviene infatti capire in che tipo di ambiente lo si progetterà: un ambiente
fatto di team, di gerarchie, di scadenze e, soprattutto, di persone che devono
riuscire a lavorare insieme.

La lezione è la prima giornata di un \textbf{workshop di due giorni} (circa
dieci ore in totale, pause pranzo comprese) tenuto da Stefano
Cozzi\footnote{Il cognome è riportato come compare nella trascrizione
automatica, che in più punti storpia i nomi propri; andrebbe verificato.
Nella stessa trascrizione lo studio per cui lavora appare come
``Stormine'', ``Stormline'' e ``Stormind'': si tratta di Stormind Games,
studio con sede a Catania, coerentemente con quanto detto in lezione.}, un
professionista con circa quindici anni di esperienza nel settore, tra
Blizzard Entertainment, Riot Games e, oggi, Stormind Games. Il workshop è
nato da un confronto con i coordinatori dell'accademia\footnote{Nella
trascrizione l'accademia è resa in modi diversi (``Hive'', ``a dive'',
``D.I.V.''); i coordinatori citati sono Enrico e Andrea, insieme al
direttore Diego Ricchiuti. Nel resto del capitolo si usa semplicemente
``l'accademia'', come negli altri capitoli.}: serviva un momento dedicato a
\emph{come ci si comporta in un ambiente professionale}, cosa che riguarda sia
il lavoro futuro sia i progetti di gruppo svolti durante il percorso di
studi, dove ci si deve già abituare a comportarsi ``come se si fosse in
un'azienda''.

La prima giornata ha un unico grande tema, sezionato in più parti:
\textbf{la cultura del feedback}. La seconda giornata (che sarà oggetto della
prossima lezione zero) sarà più pratica: \emph{time management},
\emph{brainstorming} e il ``guidare da dentro'' un team anche senza esserne
il leader (sezione~\ref{sec:z1-domani}).

Il percorso del capitolo segue quello della lezione:
\begin{enumerate}
  \item il formato del workshop e il percorso professionale del relatore
        (sezione~\ref{sec:z1-relatore});
  \item perché le aziende tengono tanto al lavoro di squadra e al
        \emph{fit} culturale (sezione~\ref{sec:z1-squadra});
  \item che cos'è il feedback e da che cosa dipende il modo di darlo
        (sezione~\ref{sec:z1-cose-feedback});
  \item le cinque fasi di vita di un team e il valore del
        \emph{post-mortem} (sezioni~\ref{sec:z1-fasi-team}
        e~\ref{sec:z1-qa-fasi});
  \item il confronto con i colleghi e alcune domande sul lavoro nei grandi
        studi (sezioni~\ref{sec:z1-rivalita} e~\ref{sec:z1-industria});
  \item il circolo virtuoso tra feedback, miglioramento e fiducia, e i suoi
        effetti sulla \emph{retention} e sulla cultura aziendale
        (sezioni~\ref{sec:z1-triangolo} e~\ref{sec:z1-cultura});
  \item il ``gioco del traditore'' e il principio del
        \textbf{default to trust} (sezioni~\ref{sec:z1-among-us}
        e~\ref{sec:z1-default-trust});
  \item le traiettorie di carriera e la valutazione delle performance
        (sezione~\ref{sec:z1-traiettorie});
  \item come si dà, come si riceve e quando si dà un feedback
        (sezioni~\ref{sec:z1-dare}, \ref{sec:z1-ricevere}
        e~\ref{sec:z1-quando});
  \item il feedback come regalo, le domande finali e le anticipazioni
        della seconda giornata (sezioni~\ref{sec:z1-regalo},
        \ref{sec:z1-qa-finale} e~\ref{sec:z1-domani}).
\end{enumerate}

Una premessa di metodo: il relatore ha alternato continuamente teoria,
domande alla classe e racconti della propria esperienza. I racconti non sono
digressioni: sono gli esempi concreti con cui i principi vengono resi
tangibili, e per questo sono riportati con cura.

% ---------------------------------------------------------------------
\section{Il workshop e il relatore}
\label{sec:z1-relatore}

\subsection{Formato e regole di partecipazione}

Il relatore ha cercato di ``riassumere quindici anni di esperienza in dieci
ore''. Ha materiale per circa il doppio delle ore disponibili: le cose
strettamente necessarie vengono trattate per prime, il resto è tenuto ``in
riserva'' come \emph{nice to have}. Le regole del workshop:
\begin{itemize}
  \item le domande si fanno \textbf{alla fine di ogni slide}, così da chiudere
        prima il discorso; non c'è un'unica sessione finale di domande e
        risposte (anche se ci sarà), perché il relatore preferisce fermarsi
        spesso e non fare ``uno spiegone di dieci ore'';
  \item le domande non devono per forza riguardare la slide appena vista: se
        una slide dà lo spunto per altro, se ne può parlare; al massimo sarà
        il relatore a rispondere in breve o a rimandare a una slide
        successiva;
  \item si può chiedere anche di esperienze passate, opinioni e sensazioni;
        se per NDA o per serietà professionale non può rispondere, lo dirà
        esplicitamente (``non lasciate che sia un \emph{blocker}'');
  \item su Zoom si usa l'alzata di mano o la chat, che il relatore tiene
        aperta sul secondo monitor;
  \item non ci sono pause ufficiali: la lezione è registrata e ognuno è
        libero di assentarsi quando serve, cercando però di non disturbare il
        flusso della lezione;
  \item le slide verranno fornite integralmente tramite la segreteria
        dell'accademia (Vanessa), insieme alle registrazioni: le slide senza
        il parlato servono poco, ma sono un buon promemoria.
\end{itemize}

\subsection{Il percorso professionale}

Il relatore si presenta (scherzosamente) come la ``guida escursionistica''
che accompagna la classe nella ``foresta incantata'' della cultura del
feedback, perché per molti è un territorio sconosciuto. Il suo percorso:
\begin{itemize}
  \item in Italia ha lavorato per aziende che fornivano servizi a terzi, tra
        cui il lavoro su \emph{Warhammer Online}, e per un canale televisivo
        sul digitale terrestre dedicato ai videogiochi\footnote{Nella
        trascrizione: ``Digital Game Network Channel''. Il nome esatto del
        canale andrebbe verificato.};
  \item a \textbf{Blizzard Entertainment} (circa sei anni, a Parigi) ha fatto
        il \emph{community manager} per l'Italia (\emph{StarCraft~II},
        \emph{World of Warcraft} da \emph{Mists of Pandaria} in poi,
        \emph{Hearthstone}, \emph{Diablo~III}, le espansioni di
        \emph{StarCraft}, il lancio di Battle.net), per poi diventare il
        manager dei community manager europei;
  \item a \textbf{Riot Games} (circa cinque anni, a Dublino) è stato
        brevemente community manager italiano e poi \emph{publishing
        manager} di un team di circa trenta persone per l'Europa. Il
        \emph{publishing} non è solo community: è l'unione di
        \textbf{community, marketing e PR};
  \item a \textbf{PG Esports} si è occupato di eventi dal vivo e online
        (ad esempio la Milano Games Week e tornei di \emph{Valorant} e
        \emph{League of Legends}), restando in contatto con Riot e Blizzard;
  \item da circa un anno è tornato nello sviluppo, a \textbf{Stormind
        Games}, dove oggi si occupa di risorse umane e organizzazione del
        lavoro; lo studio verrà presentato più ampiamente al \emph{Career
        Day} dell'accademia insieme alla sua HR manager;
  \item fa inoltre consulenze aziendali anche al di fuori dei videogiochi e
        insegna in master universitari (IULM di Milano).
\end{itemize}

Il suo ruolo ricorrente, sia a Blizzard sia a Riot, è stato quello di
\emph{Mr.~Wolf} (``per chi ha visto \emph{Pulp Fiction}''): quello che
risolve i problemi. Veniva inserito nei team più in difficoltà, ad esempio
dal punto di vista delle performance, per sistemarli e passare al successivo.
Era presente in alcuni dei momenti più critici, come il lancio di
\emph{Diablo~III} o il lancio di espansioni di \emph{World of
Warcraft}\footnote{La trascrizione riporta ``il lancio con i Gary's'' e
``World of Trenor'': si tratta verosimilmente di \emph{Mists of Pandaria} e
di \emph{Warlords of Draenor}.}, a fare \emph{damage control}.

La scoperta fondamentale di questi anni è che i problemi da risolvere
\textbf{quasi mai riguardavano il prodotto} (non era certo lui a sistemare i
server di \emph{Diablo~III}), ma i team che c'erano dietro: persone che
litigavano, non collaboravano, non riuscivano a stare nella stessa stanza
senza urlarsi contro. E nel 90\% dei casi tutto questo si riduce a
\textbf{problemi di comunicazione}: è raro che perfino le incompatibilità
caratteriali non siano risolvibili parlando.

\subsection{Struttura non vuol dire burocrazia}

A Stormind il relatore sta accompagnando un team passato da 30 a 95 persone
in meno di un anno, attraverso tutti i cambiamenti di cui parla il workshop.
Darsi una struttura viene spesso percepito male, come ``più processi, più
procedure'', ma bisogna distinguere:
\begin{itemize}
  \item un processo come ``timbrare il cartellino'' può essere davvero
        invasivo e noioso;
  \item definire \emph{come funziona un documento di game design, chi lo
        approva, chi lo visiona e chi no} è invece una struttura necessaria.
\end{itemize}
Uno studio di 30 persone nella stessa sede può farne a meno: quando c'è un
dubbio ci si alza e si va alla scrivania del collega. Uno studio di 95
persone sparse per l'Italia (e in parte per l'Europa) non può più farlo:
servono una struttura diversa, un modo di lavorare diverso e una
pianificazione del futuro diversa.

% ---------------------------------------------------------------------
\section{Perché il lavoro è sempre lavoro di squadra}
\label{sec:z1-squadra}

\subsection{Non si va in ufficio per fare amicizia}

Il primo messaggio è netto: \textbf{non si va in ufficio per fare amici} o
per trovare compagni di vita. Si va in ufficio perché c'è un lavoro da fare, e
lo scopo è riuscire a farlo \emph{anche con chi sta antipatico}. Con le
persone simpatiche è facile: diventeranno amici, e non sarà mai un problema.
È con chi sta antipatico che le cose si complicano, ed è lì che emergono i
\textbf{veri professionisti}: quelli che riescono a lavorare con chiunque
senza che l'antipatia pesi sul modo in cui lavorano e collaborano, mettendo il
prodotto e il team davanti al proprio ego e alle proprie preferenze personali.

La frase che fa da filo conduttore a tutto il workshop (e che il relatore
promette di riproporre ``fino alla nausea'') è:

\begin{quote}
  \emph{La capacità di un gruppo di persone di fare cose straordinarie
  dipende dal modo in cui queste persone operano insieme.}
\end{quote}

Le aziende raramente pensano che l'innesto di una singola persona possa
cambiare il destino di un progetto. Questo è tanto più vero quanto più grande
è il team, ma vale anche per i piccoli studi, anche per un gruppo di cinque
persone. Come negli sport di squadra: l'individuo conta fino a un certo punto,
e anche se si ha Dončić in squadra, se si gioca uno contro cinque si perde
comunque\footnote{Il riferimento è alle finali NBA in corso nei giorni del
workshop tra Dallas Mavericks e Boston Celtics (giugno 2024), il che
permette di datare approssimativamente la lezione.}.

\subsection{Lo \emph{useless genius} e il costo delle persone tossiche}

Molti dipartimenti HR usano la definizione di \textbf{useless
genius}\footnote{Il concetto è più noto in letteratura manageriale come
\emph{brilliant jerk}; in lezione è usata l'espressione \emph{useless
genius}.}: si può essere bravissimi, a volte i migliori nel proprio lavoro,
ma se non ci si sa comportare, relazionare e operare in un gruppo si è
fondamentalmente inutili ai fini del gruppo. Nel migliore dei casi inutili,
nel peggiore \emph{dannosi}.

In passato le aziende tendevano a ``chiudere un occhio'': la persona è
fondamentale, la teniamo. Oggi la domanda è diventata: \emph{quanto costa
all'azienda, oltre al suo stipendio, tenere una persona che crea un ambiente
tossico e peggiora chi le sta intorno?} Ci sono costi sommersi che le aziende
cominciano a vedere, e la tendenza è liberarsi di queste persone il prima
possibile. Il ragionamento è semplice:
\begin{itemize}
  \item le \textbf{competenze si possono trasmettere} o sostituire: si trova
        qualcuno che fa lo stesso lavoro magari nel doppio del tempo, oppure
        più persone che se lo dividono;
  \item il \textbf{carattere delle persone} invece cambia molto raramente.
\end{itemize}

\subsection{Il fit culturale}

Quando qualcuno \textbf{non è un \emph{fit} culturale} per un'azienda è
meglio separarsi presto. Ma non essere adatti a \emph{quell'}azienda non
significa non essere adatti a lavorare: è perfettamente legittimo che una
persona inadatta a un'azienda sia adattissima a un'altra.

Il relatore l'ha vissuto in prima persona. Dopo sei anni a Blizzard viene
contattato da Riot e affronta un processo di selezione lunghissimo (``quasi
undici colloqui''), al termine del quale viene scartato. Scrive allora una
mail che gli ha ``cambiato la traiettoria di carriera'': chiede
semplicemente perché non è stato ritenuto adatto. La risposta: in nessun
colloquio si era esposto abbastanza da far capire \emph{perché} volesse
lasciare Blizzard, e quindi pensavano che in realtà stesse bene dov'era. Il
motivo era che, per educazione professionale, non gli piace parlare male del
proprio posto di lavoro; portata all'eccesso, però, questa riservatezza era
diventata un ``è tutto bellissimo, ma vorrei provare Riot'', senza un perché.
Lo invitano comunque a un incontro dal vivo (``dal vivo è diverso''), e alla
fine viene assunto.

Una volta arrivato, i sei mesi di prova (\emph{probation}) sono serviti tutti
per smettere di ragionare e lavorare ``come a Blizzard''. La differenza
culturale era enorme, tanto che in Riot esisteva un gruppo di supporto
interno tra ex dipendenti Blizzard: un collega gliel'ha descritto come
``uscire da una relazione tossica e doversi riabituare a una relazione
normale''\footnote{Il relatore precisa che parla della Blizzard di allora,
che non esiste più: oggi è parte di Microsoft e nessuno dei suoi ex colleghi
vi lavora ancora, anche a causa dei licenziamenti di massa.}.
\begin{description}
  \item[Blizzard] era \emph{developer first}: il focus era chiaramente sugli
        sviluppatori, prima di tutti gli altri. Dopo sei anni c'erano ancora
        meeting a cui non poteva partecipare perché non abbastanza
        \emph{senior}, informazioni che non poteva sapere, piani che non
        poteva vedere.
  \item[Riot] era \emph{player first}: prima i giocatori, e i dipendenti come
        mezzo per trasferire il feedback dei giocatori dentro il gioco. Il
        secondo giorno lo portano in una stanza e gli mostrano il piano dei
        cinque anni successivi. Lui, sorpreso, pensa a un errore (``sono qui
        da oggi, non vorrei mettervi nei guai''); la risposta è: ``sei un
        dipendente Riot o no? Allora non ci sono livelli di accesso: devi
        sapere queste cose per fare meglio il tuo lavoro''.
\end{description}

Da questa esperienza ha imparato molto sulla \textbf{trasparenza}: condividere
le informazioni aiuta il lavoro. Non sempre, però: a volte non si condivide
qualcosa non per segretezza, ma perché non è ancora pronta per essere
presentata, esattamente come non si mostra un gioco quando è ancora solo
teoria (``te lo faccio vedere quando ho qualcosa di tangibile'').

La lezione per gli studenti è doppia: le aziende valuteranno il \emph{fit}
culturale sempre di più, ma capire se si è un \emph{fit} richiede anche
\textbf{lavoro da parte del candidato}. Bisogna avere la presenza di spirito
di dire: ``conosco i valori di questa azienda, non mando nemmeno il CV,
perché so che durerei un mese e mezzo''.

% ---------------------------------------------------------------------
\section{Che cos'è il feedback}
\label{sec:z1-cose-feedback}

\subsection{Le prime risposte della classe}

La prima domanda alla classe è: \emph{quando sentite la parola
``feedback'', a cosa pensate?} Le risposte raccolte, a voce e in chat:
\begin{itemize}
  \item un \textbf{riscontro} su ciò che si fa e su come lo si fa, sul lavoro
        o sul modo in cui ci si interfaccia con il gruppo; può essere
        lavorativo oppure comunicativo e comportamentale. Chi ha risposto lo
        vede come qualcosa che, se l'azienda non dà, va attivamente
        richiesto, perché è importante sapere come sta andando per
        migliorarsi;
  \item né positivo né negativo, ma \textbf{costruttivo}: in entrambi i casi
        serve a migliorarsi;
  \item uno \emph{scambio di opinioni}, un \emph{riscontro produttivo};
  \item un feedback ``di reazione'', non verbale: si fa una cosa e si
        osserva come cambia l'ambiente o il comportamento delle persone
        intorno (anche perché a voce la gente spesso si tiene le cose
        negative);
  \item per un designer, anche il \textbf{linguaggio del corpo}: osservando
        un giocatore giocare si può dedurre l'emozione che prova (frustrazione,
        incomprensione).
\end{itemize}

Sulla forma (scritto o orale?) le opinioni divergono: c'è chi lo preferisce
\emph{mille volte} scritto; chi pensa che dipenda dalle dimensioni
dell'azienda (orale nelle piccole, scritto in quelle con team di centinaia di
persone); chi pensa che dipenda dall'importanza del feedback (scritto se
tecnico o pieno di dati, da tenere sotto gli occhi; orale se è un consiglio
veloce).

Il relatore usa due esempi per mostrare che \emph{dipende}:
\begin{itemize}
  \item un feedback sulla programmazione ha senso darlo scritto, magari
        commentando direttamente il codice, più che ``andando a prendere un
        caffè'';
  \item ma se il feedback è ``ieri, durante il meeting, ho percepito male il
        fatto che mi interrompessi di continuo'', ce lo si aspetta
        \textbf{orale}.
\end{itemize}

\subsection{Da che cosa dipende il modo di dare feedback}

Lo \emph{spoiler} della giornata è che \textbf{non esiste una risposta univoca}
su come dare feedback. Il workshop non insegna ``il'' modo, ma fornisce
\emph{tutte le opzioni}, perché il modo giusto cambia in base a molti fattori:
\begin{itemize}
  \item a chi lo si dà: un collega, un superiore, il proprio capo, un proprio
        \emph{report} (un sottoposto);
  \item se è personale o professionale;
  \item se è la prima volta o la terza volta che si dà lo stesso feedback;
  \item se riguarda un progetto importante o secondario.
\end{itemize}

Manca però l'elemento fondamentale, che la classe individua subito:
\textbf{l'altro}. A parità di tutti gli altri fattori, il \emph{carattere
della persona che riceve il feedback} dovrebbe far cambiare il modo in cui lo
si dà. Per questo, prima della lezione, agli studenti era stato chiesto di
fare un test di personalità: non un test di psicologia avanzata, ma uno
strumento per rendersi conto di quanta varietà ci sia anche solo in una
classe di 53 persone (il che non vuol dire trovare 52 modi diversi di dare
feedback, ma lavorarci un po' sì).

Il relatore porta se stesso come esempio. Preferisce:
\begin{itemize}
  \item i feedback \textbf{positivi scritti}, perché è a disagio a ricevere
        complimenti in pubblico;
  \item i feedback sul proprio lavoro \textbf{diretti}: ``vai subito al
        punto e non girarmi intorno''.
\end{itemize}
È molto orgoglioso del suo lavoro e, quando qualcuno lo critica, la sua
primissima reazione istintiva è ``come ti permetti''. Negli anni ha imparato
a farla durare un secondo, ma chi gli dà feedback di persona dice di
vederla ancora negli occhi. Per mitigarla, dice esplicitamente a chi gli dà
feedback: \emph{``ignora il modo in cui reagisco nei primi cinque secondi,
vai avanti, poi mi passa''}. È considerato da tutti molto aperto al feedback
e non ha problemi a dire ``hai ragione, facciamo diversamente'': ma ci sono
reazioni istintive e caratteriali che è difficile controllare al cento per
cento.

\subsection{Con chi si lavora meglio}

Guardando le persone con cui si lavora abitualmente (anche nei progetti
dell'accademia), si noterà che si lavora meglio con due profili:
\begin{enumerate}
  \item chi è \textbf{esattamente come noi}, ma qui è un 50 e 50: alcuni
        caratteri vanno d'accordo con i propri simili, altri non vogliono
        nessuno con lo stesso carattere nel gruppo (e si potrebbe rientrare
        in questa categoria);
  \item chi è \textbf{affine}: ha qualcosa in comune con noi, ma non tutto,
        e proprio per questo può dare punti di vista differenti.
\end{enumerate}
Il lavoro della giornata è capire come valorizzare il secondo caso: ottenere
feedback da persone con caratteri, visioni e background diversi, per
completare una visione del proprio lavoro che è necessariamente limitata. È
come esplorare una mappa con la nebbia di guerra: un'altra persona, al di là
del muro, vede cose che noi non vediamo.

\subsection{Il feedback come somma di piccoli miglioramenti}

Per il relatore il feedback è questo: \textbf{applicare tanti piccoli
miglioramenti periodici}, che si sommano nel tempo. Migliorare dello 0,1\%
ogni giorno porta a fine anno a un miglioramento enorme, e questo vale per i
prodotti, per i giochi, per le persone e per i team.

\paragraph{L'esempio del ciclismo britannico.} La squadra di ciclismo
britannica, fino a una quindicina d'anni fa, non aveva vinto quasi nulla.
Cambiato il team manager, tutti si aspettavano grandi cambiamenti sportivi
(convocazioni, strategie di gara). Arrivarono invece \emph{minuscoli
cambiamenti non sportivi}: un altro chef in trasferta, materassi migliori nei
camper per riposare meglio, un sellino più comodo a fine tappa. Cambiamenti
marginali che da soli non avrebbero cambiato niente; eppure, dopo due o tre
anni di questo lavoro quotidiano, la squadra ha cominciato a vincere a
ripetizione, Olimpiadi comprese, cambiando il modo in cui il ciclismo è
visto in Gran Bretagna\footnote{È la nota strategia delle \emph{marginal
gains} associata a Dave Brailsford (British Cycling e Team Sky); il nome non
è citato in lezione.}. Certo, si è investito anche sulle persone giuste, ma
tutto è partito dai piccoli miglioramenti periodici.

Per gli studenti vale lo stesso: chiarire un conflitto con una persona con
cui si ha ancora poco a che fare, o una situazione che ha lasciato
``l'amaro in bocca'', sembra inutile oggi, ma darà i suoi frutti tra tre, sei,
nove mesi. Sono \emph{investimenti sul futuro}.

\paragraph{Peggiorare ogni giorno.} Esiste anche il percorso inverso:
lasciar correre, ``buttare sotto il tappeto'' le cose perché ``tanto è
inutile parlarne''. Queste cose tornano sempre fuori dopo otto mesi, come il
\emph{libro dei rancori} dei nani di Warhammer: ``ti ricordi quando, un anno
e mezzo fa, all'inizio del progetto, mi hai detto quella cosa? Mi sono
offeso''. A quel punto è quasi impossibile risolvere, perché nessuno ricorda
più con precisione cosa sia successo (``ma quale cena?'', ``forse era in
pizzeria'', ``no, eravamo da un'altra parte''). \textbf{Tenersi le cose nel
cassetto per tirarle fuori quando servono è sempre sconsigliato}: non servono
ed è impossibile usarle in modo costruttivo.

% ---------------------------------------------------------------------
\section{Le cinque fasi di vita di un team}
\label{sec:z1-fasi-team}

\subsection{Perché conoscerle}

Quasi tutti gli studenti lavorano già in gruppo (al primo anno su un gioco da
tavolo, negli anni successivi su un videogioco) e si saranno accorti che un
team attraversa delle fasi: il modo di lavorare insieme cambia nel tempo.
Spesso ci si fa prendere di sorpresa: ``eravamo così affiatati, cosa è
successo?''.

Un esempio da una consulenza fuori dai videogiochi: un'azienda farmaceutica
gli affida un team commerciale (rappresentanti che presentano medicinali a
ospedali, medici e privati). Il problema: ``prima andava tutto bene; ora che
siamo in diciotto dovremmo lavorare meno e meglio, e invece va peggio''. Va
peggio perché le persone nuove sono appena arrivate, mentre il management si
aspettava già di vedere i risultati dell'investimento.

Fino a pochi anni fa le fasi di un team venivano disegnate come una
\textbf{linea dritta verso l'alto}: le fasi erano le stesse, ma viste come un
miglioramento costante. Non si vedevano le parti \emph{fisiologiche} del
lavoro in team, cioè anche i peggioramenti. Accettarle è il primo passo per
accettare i fallimenti: nessun percorso personale o di team è lineare.

Il relatore ammette che il suo modo di lavorare è nato per il 70\% dal
proprio impegno a imparare, fare bene e rispettare le persone, e per il 30\%
dall'aver lavorato con ``incapaci colossali'' che gli hanno insegnato come
\emph{non} fare le cose. A Blizzard un manager gli fece capire che, pur non
sapendo ancora che tipo di \emph{people manager} sarebbe diventato, sapeva
benissimo quale tipo di manager \emph{non} voleva essere.

\subsection{Forming, storming, norming, performing, adjourning}

Le fasi sono cinque, un po' come le fasi del lutto. Molte aziende ne
considerano solo quattro (``e sono quelle che solitamente chiudono''); le
aziende al passo con i tempi ne considerano cinque\footnote{Il modello è
quello di Bruce Tuckman (1965), a cui la quinta fase, \emph{adjourning}, fu
aggiunta nel 1977. Il nome di Tuckman non è citato in lezione; la
trascrizione riporta la quinta fase come ``journing''.}.

\begin{description}
  \item[Forming] Si viene messi insieme e si comincia a lavorare. Non si sa
        nulla degli altri; spesso non si conoscono nemmeno i nomi, ed è
        normale.
  \item[Storming] Si cominciano a ``dare le testate come le capre in
        montagna per determinare il territorio'': questa cosa la faccio io,
        questa la vorrei fare così, chi decide? Si stabilisce chi fa cosa e
        come si divide il lavoro. Nel grafico è una \textbf{fase di
        peggioramento}: che il team litighi o discuta civilmente non cambia
        nulla, il team deve investire parte di tempo ed energie per risolvere
        queste questioni. È normale e \textbf{non è evitabile}.
  \item[Norming] Si comincia a capire come lavorano gli altri, cosa sanno
        fare e come far funzionare il gruppo; iniziano i miglioramenti. È una
        forma di \emph{accettazione} degli altri, che deriva dall'aver capito
        cosa devo fare io e cosa devono fare loro.
  \item[Performing] C'è uno sforzo \textbf{proattivo}: ``questa cosa
        toccherebbe a me, ma so che Alessio la fa molto meglio, quindi gli
        chiedo aiuto''. Si sfruttano i punti di forza degli altri, perché li
        si conosce e perché si sa che passare il testimone non avrà effetti
        negativi (il capo non dirà ``perché non l'hai fatta tu?'').
  \item[Adjourning] Il progetto finisce, si chiude o viene riorganizzato, il
        team viene sciolto. Quello che andrebbe fatto è analizzare
        \textbf{cosa ha funzionato, cosa no e cosa migliorare} per il
        progetto successivo.
\end{description}

\subsection{La fase che tutti saltano}

``Spoiler incredibile'': l'\emph{adjourning} è la fase che molte aziende
saltano, \textbf{perché c'è sempre il progetto dopo}.

A PG Esports, per la Milano Games Week (come per qualsiasi evento dal vivo a
cadenza annuale: Lucca Comics, Etna Comics, Comicon\dots), finita l'edizione
di un anno si cominciava dal giorno dopo a lavorare a quella successiva,
perché un anno sembra tanto ma non basta mai. Nei videogiochi è lo stesso:
finito un gioco, a parte le patch di mantenimento e del \emph{day one}, si
passa subito al successivo.

Il cambiamento introdotto dal relatore (in parte già praticato, ma da lui
strutturato) è stato: finita la Games Week, \textbf{tutti in ferie per una
settimana}; poi ci si ritrova per definire cosa ha funzionato e cosa no, con
un \textbf{post-mortem} (o \emph{debrief}); solo dopo si inizia a lavorare
all'edizione successiva. Senza questa fase manca la parte di
\emph{learning}, cioè di ciò che si è imparato dal progetto.

Il consiglio vale per gli studenti: quando si finisce il gioco da tavolo, e
poi il gioco digitale, bisogna trovare il tempo per il post-mortem,
altrimenti si butta via \textbf{il 90\% dell'utilità del progetto}. Una
struttura semplice per farlo è descritta nella sezione~\ref{sec:z1-qa-finale}.

% ---------------------------------------------------------------------
\section{Domande sulle fasi del team}
\label{sec:z1-qa-fasi}

\subsection{Lo storming che ritorna}

\paragraph{Capita che lo storming si ritagli spazio anche nelle altre fasi?}
Lo \emph{storming} è una fase precisa, che però molti manager, coordinatori e
team cercano di \textbf{saltare}: ``se io sono bravo come manager e loro sono
bravi come team, passeremo da forming a norming senza problemi''. È un errore,
perché senza una fase in cui si discute davvero non si sviluppa una
\emph{skill} fondamentale: \textbf{saper discutere}. Tutti conoscono persone
non abituate a farlo, per cui ogni minimo disaccordo diventa un affronto
personale (``avevamo deciso la pizza, perché andiamo al sushi? Perché mi fate
questo affronto?''): sono persone che non sono abituate ad avere conflitti
costruttivi.

Se si pensa di essere in \emph{norming} e continuano a emergere conflitti, il
punto è che \textbf{non si era davvero in norming}: lo storming non ``si
trascina'' nelle fasi successive, semplicemente non se n'è mai usciti.
Allo stesso modo, a chi chiede se ogni mini-fase di storming durante la vita
del gruppo sia frutto di uno storming fatto male o solo di carattere, la
risposta è già contenuta qui.

Un partecipante osserva che spesso, ancor prima di finire un progetto, si è
già messi su quello successivo: è il modo migliore per essere in costante
\emph{damage control}. Era l'esperienza del relatore a Blizzard: veniva
spostato fisicamente da un team all'altro, e finito il lancio di
\emph{Diablo~III} era già su \emph{World of Warcraft}, senza vivere il
``dopo''. Si è dovuto ritagliare da solo lo spazio per un debrief personale,
ed è lì che ha deciso che tipo di manager non voleva essere. Il consiglio:
\textbf{se non lo fa il team, fatelo voi}.

\subsection{Dove si passa la maggior parte del tempo}

\emph{In quale fase si passa più tempo nella propria carriera?} Gli studenti
rispondono \emph{storming} e \emph{norming}. Il relatore, in quindici anni,
calcola di aver passato in \emph{performing} circa \textbf{dieci mesi}. Non è
pessimismo: arrivare a performing deve essere sempre l'obiettivo, sia da
manager sia da membro del team; ma è difficile arrivarci, ed è normale che non
sempre ci si arrivi.

Lo storming può durare a lungo, ma con team molto esperti dura pochissimo,
perché c'è chi lo sa gestire. In realtà \textbf{la maggior parte del tempo si
passa in forming}, perché ogni equilibrio raggiunto vale per \emph{quelle}
persone, \emph{quel} progetto, \emph{quelle} tempistiche. Basta che cambi un
qualsiasi fattore (una deadline spostata, una persona nuova, qualcuno che si
dimette, un progetto cancellato) e si ricomincia, perché i bilanciamenti
precedenti non valgono più.

Esempio: un team di cinque persone affiatatissimo, in pieno performing. La
dirigenza, entusiasta, ``raddoppia le risorse'' e lo porta a dieci. O si
lascia tutto com'è (e i cinque nuovi non fanno niente), o bisogna ridiscutere
chi fa cosa e come i nuovi possono aiutare: che piaccia o no, si ricomincia
da forming.

Quindi la bravura di un manager, ma anche di un membro del team,
\textbf{non sta nel passare tanto tempo in performing, ma nell'accelerare il
passaggio da una fase all'altra}, e nel fare bene la fase di
\emph{adjourning}.

\subsection{Analizzare anche i successi}

C'è la tendenza a fare post-mortem solo dei progetti andati male: ``questa
cosa non ha funzionato, vediamo perché''. Va benissimo. Ma quando un progetto
va bene si dà per scontato che sia andato bene ``perché era buono''. Bisogna
invece chiedersi \emph{perché} è andato bene: era il prodotto? La
tempistica? Il modo in cui abbiamo fatto una cosa piuttosto che un'altra?
\textbf{Se non si capisce davvero il motivo del proprio successo, si
tenderanno a replicare le cose sbagliate.}

\paragraph{Blockbuster.} Quando arrivò la notizia che Netflix avrebbe
lanciato un servizio di streaming, Blockbuster (la catena di noleggio di
videocassette e DVD) aveva tutto il tempo per preparare un'offerta
concorrente. In una riunione del management raccontata in un libro (di cui il
relatore non ricorda il titolo), la direzione sembrava orientata a investire
nello streaming, finché un dirigente fece notare che il 12\% del fatturato
derivava dalle \emph{multe per riconsegna tardiva}, impossibili con lo
streaming, mentre secondo le stime lo streaming avrebbe prodotto il 10\% del
fatturato. ``Non possiamo rinunciare a questo 2\%'': niente streaming. Oggi
non c'è più un Blockbuster ``manco a cercarlo col lanternino''. Non avevano
capito qual era il vero motivo del loro successo, cioè il loro modello di
business\footnote{Le percentuali sono quelle riferite in lezione e non sono
state verificate.}.

\paragraph{The Founder.} Nel film sulla nascita di McDonald's c'è una scena in
cui il futuro socio del protagonista gli dice: ``tu pensi di essere nella
\emph{food industry}, e invece sei nella \emph{real estate industry}''. Non
dovrebbe investire sugli hamburger, ma sui terreni su cui costruisce i
ristoranti. Anche per un team vale lo stesso: bisogna capire che cosa
l'ha reso davvero di successo, oppure no.

\subsection{L'autoanalisi personale}

\paragraph{E quando si cambia azienda?} Uno studente chiede se il debrief sia
possibile anche quando si lascia un'azienda. Non si può tornare in ufficio a
discutere della propria esperienza, ma lo si può fare \textbf{per se stessi}.
Il relatore l'ha fatto, anche se all'epoca in modo poco consapevole, ed è
così che è passato dal \emph{community management} al \emph{people
management}. Analizzando cosa aveva fatto bene a Blizzard si era reso conto
che:
\begin{itemize}
  \item aveva costruito un ottimo rapporto di fiducia con i giocatori (con
        alcuni proprietari di fansite parla ancora oggi);
  \item pur senza la seniority, era diventato un punto di riferimento per il
        team, visto come imparziale ed equilibrato anche su temi che non lo
        riguardavano direttamente;
  \item sapeva \textbf{risolvere problemi}: durante il lancio di
        \emph{Diablo~III} era il più tranquillo di tutti, non perché
        sminuisse il problema, ma perché sapeva che si sarebbe risolto con i
        suoi tempi; il lavoro non era spegnere l'incendio ma ``far sfollare
        la gente'' e aspettare che tutto tornasse in ordine.
\end{itemize}
E aveva capito cosa sapeva fare meno bene, ad esempio \textbf{gestire il
proprio manager}. Avere un manager non significa che il manager gestisca te:
il manager gestisce il tuo lavoro, ma tu devi saper gestire i tuoi
\emph{stakeholder}. Spesso lo stakeholder è il proprio manager; da manager,
gli stakeholder saranno CEO, CFO e così via. Gestirli significa conoscerli,
sapere come portare loro le informazioni, come far presente che qualcosa non
va. Fa parte del lavoro, anche se a lui non piace, e ci ha lavorato.

Nel mese e mezzo di pausa prima di entrare in Riot si è concentrato su
questo: ``cosa porto come mio punto di forza?''. L'autoanalisi è culminata
nel chiedere espressamente alla persona che l'aveva assunto: ``prima mi avete
detto di no, ora mi avete preso: perché?''.

\paragraph{Come si capisce se un progetto è andato bene?} Dipende
dall'\textbf{obiettivo}. Se l'obiettivo era vendere milioni di copie e il
gioco le ha vendute, bisogna capire perché. Se, più ragionevolmente,
l'obiettivo era un gioco da tavolo che permettesse di mettere in pratica le
conoscenze del corso, va valutato per quello: ha permesso di farlo? Se sì,
perché? Se no, perché no? L'obiettivo va fissato \textbf{all'inizio del
progetto, non alla fine}, proprio per poterlo valutare dopo.

\subsection{Quando il team cambia di continuo}

\paragraph{Tutti questi cambiamenti non sono dannosi per il progetto?}
Assolutamente sì, ma quasi mai sono volontari. I progetti andati male sono
spesso progetti grandi.

Il relatore racconta di \emph{Cyberpunk 2077}. Conosceva persone in CD
Projekt Red (``il mondo è piccolissimo, le cose si sanno''). Alla Gamescom va
con Diego Ricchiuti, direttore dell'accademia, a vedere la presentazione del
primo gameplay, quella ``in cui regalavano il giubbottino da samurai''.
Finito il video, tra gli applausi, si guardano e capiscono entrambi che
\textbf{il gioco non è pronto}, a soli quattro mesi dall'uscita prevista:
Diego aveva notato il combattimento e la guida. Gli amici di CD Projekt Red
(uno aveva lavorato con lui a Blizzard per sei anni) gli confermano tra le
righe che ``ce la faremo, ma ci sono ancora cose da sistemare''. Il gioco fu
poi rimandato più volte\footnote{In lezione il relatore dice ``due, forse
tre volte''; l'uscita di \emph{Cyberpunk 2077} è stata effettivamente
rinviata tre volte prima del dicembre 2020.}.

Per deformazione professionale, chi ha sviluppato l'occhio critico non
riesce più a godersi un trailer: in un trailer di \emph{The Day Before} Diego
gli aveva fatto notare che l'effetto della fiammata di un fucile era stato
chiaramente aggiunto in postproduzione e non faceva parte del gameplay. ``È
il triste epilogo di questa storia'': non riescono più a godersi un trailer.

Il punto è che più il team è grande, più queste variabili vengono imposte, e
più è facile che il fallimento \textbf{non dipenda dal fatto che il prodotto
sia brutto} (\emph{Cyberpunk} è un bellissimo gioco sotto molti punti di
vista, e il relatore lo sta rigiocando ora che è stato corretto ed espanso).
Nei grandi gruppi qualcuno arriva, qualcuno si dimette, cambia il budget,
arriva il rinvio; nel caso di \emph{Anthem}, il team ha fatto così tante volte
forming--storming--forming--storming che ne è uscito ``un prodotto che
dovevamo espellere: l'abbiamo espulso, buttato fuori dal nostro organismo pur
di liberarcene''.

Tra i cambiamenti c'è anche la crescita verticale: se il capo viene promosso
perché lavora benissimo, bisogna scegliere un nuovo capo, ed è sempre un
caos (chi ci resta male, chi pensava di meritarlo, chi voleva il suo amico,
chi non sopporta il prescelto). Per questo il relatore è \textbf{contrario
alle promozioni durante i progetti}: si rischia di rovinare tutto il lavoro di
storming.

\subsection{Obiettivi chiari nei piccoli gruppi}

Nei gruppi piccoli, come quelli degli studenti, il problema deriva più
spesso dal \textbf{non avere chiaro che cosa si vuole ottenere}: un prodotto
puramente educativo e formativo, o un prodotto richiesto dall'accademia che,
se viene bene, si potrebbe anche vendere? Un gioco commerciale o un gioco
artistico? Va chiarito subito, perché ci saranno sempre momenti in cui
scegliere tra A e B, e a guidare la scelta sarà la domanda: \emph{quale delle
due opzioni ci porta più vicino all'obiettivo?} Se l'obiettivo non è chiaro,
non si sa cosa scegliere.

Questo porta molte aziende a fare \textbf{scelte di design che in realtà sono
scelte commerciali}. Le microtransazioni sono scelte commerciali: nessuno ha
mai pensato ``sai cosa piacerebbe tanto ai giocatori? Pagare le bustine''.
Prima le bustine si guadagnavano giocando; ora la tendenza è ``potresti
giocare, ma se paghi è meglio''. Si pensi alle polemiche su Ubisoft per
l'opzione di ``saltare il gioco'' pagando, o ai famosi \emph{boost} di livello
di \emph{World of Warcraft}: ``devo pagarti per saltare il gioco?''.

Una studentessa (Silvia) cita \emph{Final Fantasy XIV}, che a un certo punto fu
``resettato'' perché non funzionava: un esempio positivo. Il relatore aggiunge
quello negativo di \emph{Age of Conan}, che non è mai riuscito a fare quel
reset iniziale (``è una schifezza, ve lo ripresentiamo quando è pronto''). Ci
vuole consapevolezza, ma anche i mezzi: un conto è Square Enix, che può
permetterselo, un altro è un'azienda che non può chiudere tutto per un anno
continuando a pagare gli stipendi.

Riassumendo: nei piccoli gruppi i problemi sono più di
\textbf{comunicazione e chiarezza} (obiettivi non chiari o che si spostano);
nelle grandi aziende di più di \textbf{persone che vanno e vengono, obiettivi
e tempistiche che cambiano}. In accademia le scadenze sono fisse (il gioco da
tavolo si consegna a fine anno scolastico, a giugno, e giugno resta); in
azienda può succedere che una feature prevista per la \emph{vertical slice}
di luglio venga anticipata a maggio per portarla alla Gamescom, e un sistema
perfetto per consegnare a giugno può non esserlo più per consegnare a maggio.

\subsection{L'ambiente determina la qualità del prodotto}

\paragraph{È vero che chi passa per Electronic Arts finisce per fare brutti
giochi?} Uno studente cita BioWare, che faceva ottimi giochi finché era
autonoma. Il problema non è dei developer, che non sono ``improvvisamente
diventati incapaci'': è evidente che il problema è altrove, e che
\textbf{l'ambiente in cui si fa un gioco determina molto spesso la qualità del
prodotto}.

Un esempio locale positivo: Tommaso Bonanni, conosciuto dal relatore insieme a
Diego e già docente dell'accademia, ha fondato a Roma Caracal Games Studio,
applicando fin dall'inizio questi concetti di cultura aziendale. L'azienda è
rimasta piccola, ma ha molto chiaro cosa vuole fare e come. Quando si
predispone fin dall'inizio l'azienda e il team a lavorare in un certo modo è
fattibilissimo; è molto più difficile quando si è EA, si acquisisce
un'azienda e le si dice ``da oggi si fa così''.

Il relatore ha vissuto la transizione da Blizzard ad Activision Blizzard (uno
dei motivi per cui se n'è andato, volontariamente). Ha visto la cultura
interna cambiare per gradi:
\begin{enumerate}
  \item ``facciamo giochi fighissimi, perché i giochi devono essere
        fighissimi'';
  \item ``facciamo giochi fighi così vendono di più'';
  \item ``vendiamo di più facendo giochi fighi'' (sembra uguale, ma è molto
        diverso);
  \item ``facciamo più soldi''.
\end{enumerate}
Lo si capiva dalle scelte di design, osservando quale fosse il peso più
grande sulla bilancia. Lo stesso varrà per gli studenti, nei loro giochi e nei
colloqui di lavoro: capire subito che cosa vuole l'altra parte.

\paragraph{Leader e visione.} Nei gruppi dell'accademia, dal secondo anno, c'è
una persona che assume il ruolo di \emph{lead} o \emph{scrum master}. È
diverso dall'essere un \emph{game director} (``il gioco si fa come dico io,
aiutatemi a farlo con la mia visione''), come avviene nei grandi studi: da lì
nascono i nomi nei titoli, come \emph{Sid Meier's Civilization}, che indicano
la persona che ha fatto quel gioco, oppure un puro branding, come \emph{Tom
Clancy's Rainbow Six}.

\paragraph{Inclusion come jolly.} A margine, uno studente osserva che in molte
aziende ``la risposta giusta è sempre \emph{inclusion}'', usata come un
jolly. È comunissimo, conferma il relatore: \emph{diversity and inclusion}
viene spesso usata per dire ``così includiamo più clienti'', ma non è quello
il suo punto (``altro discorso'').

% ---------------------------------------------------------------------
\section{Confrontarsi con chi è più bravo}
\label{sec:z1-rivalita}

Arrivato a Riot dopo sei anni a Blizzard, il relatore pensava di avere un buon
bagaglio di esperienza. Si è trovato invece ``come un fante buttato in mezzo
ai Navy SEAL'': persone che avevano superato 11, 12, 13 colloqui, con un
livello molto alto.

Individua subito un collega francese che, nella sua testa, diventa ``la
versione migliore di me'': stesso ruolo (community manager), esperienza
simile (arrivava da Ubisoft, come quasi tutti i francesi di Riot), ma molto
più spigliato, sempre con la battuta pronta, ``l'anima della festa''. Il
relatore invece è introverso: un workshop come questo per lui è
\emph{emotivamente drenante}, non ama parlare in pubblico pur amando
condividere le proprie esperienze, e non riesce a sostenere conversazioni
come ``cosa hai fatto questo weekend?''. Lavorando raramente insieme, aveva
sviluppato una rivalità ``da allenatore di Pokémon'' (lui era Ash, il
relatore l'antagonista).

Dopo sei mesi, a un evento di \emph{team building} aziendale (a Tenerife), il
suo capo, molto attento a queste dinamiche, li mette apposta nello stesso
gruppo. A fine serata, complice l'open bar, il relatore gli confessa di averlo
mal sopportato dal primo giorno per invidia della sua estroversione. Il
collega ride e risponde che pochi giorni prima aveva detto al capo la stessa
cosa al contrario: invidiava a Stefano il suo essere sempre organizzato, sul
pezzo, presentarsi nella sala riunioni giusta (mentre lui finiva in quella
sbagliata), arrivare con gli appunti, con i dati a supporto e le analisi di
mercato, mentre lui spesso non aveva nemmeno letto la mail del meeting.
Un'invidia in due direzioni, che li ha portati a lavorare molto insieme e a
migliorarsi a vicenda.

Le lezioni:
\begin{itemize}
  \item non bisogna demoralizzarsi vedendo qualcuno che fa lo stesso lavoro
        ed è più avanti sotto alcuni aspetti: spesso c'è un bilanciamento. I
        ``punti caratteristica'' sono limitati: se li ha messi in carisma,
        non li avrà messi in saggezza;
  \item conviene \textbf{chiedere a queste persone come mai} sono così brave
        in qualcosa, come si fa con un progetto andato bene: magari non lo si
        farà mai bene quanto loro, ma si comincerà ad avere il
        \emph{mindset} giusto; e viceversa, loro saranno spinti a chiedere a
        noi;
  \item spesso \textbf{le proprie doti le scopriamo dagli altri}. Il relatore
        non si era mai accorto di essere ``quello che resta calmo durante le
        emergenze'' (il lancio di \emph{Diablo~III}, quello di
        \emph{Draenor}, i problemi dal vivo alla Games Week) finché non
        gliel'hanno fatto notare tutti: per lui era la normalità. A chi pensa
        ``non sono capace, forse non è il mio percorso'', può darsi che
        semplicemente nessuno gli abbia ancora detto ``questa cosa la sai
        fare davvero bene, concentrati su questo''.
\end{itemize}

% ---------------------------------------------------------------------
\section{Domande sul lavoro nell'industria}
\label{sec:z1-industria}

Diverse domande degli studenti hanno riguardato le condizioni di lavoro nei
grandi studi internazionali. Le risposte sono raccolte qui per argomento.

\subsection{Presenza, remoto e trasferimenti}

\paragraph{Grandi studi: quasi solo in presenza.} Sia a Blizzard sia a Riot si
lavorava \textbf{solo in presenza}, senza smart working: sei anni a Parigi e
cinque a Dublino. Si parlava solo inglese: anche nei meeting tra soli
italiani si tendeva a comunicare e a scrivere in inglese, per correttezza
verso chi si fosse aggiunto e per lasciare una traccia comprensibile a tutti.
Il francese serviva per fare la spesa e sopravvivere, non in ufficio.
L'ambiente era estremamente eterogeneo (francesi, polacchi, spagnoli e
molti altri): durante il colloquio dal vivo a Dublino, mentre si preoccupava
di essersi vestito in modo adeguato, vide passare dietro la parete a vetro un
collega svedese con i dreadlock, a piedi scalzi, con una pianta sotto un
braccio e un aquilone sotto l'altro, e capì di essere ``a posto così''.
Un'esposizione a molte culture e modi di pensare che il relatore considera
bellissima.

Anche oggi, tranne alcuni studi particolari, la maggior parte del lavoro si
svolge di persona: amici a Paradox lavorano a Stoccolma, i pochi rimasti a
Blizzard si sono dovuti trasferire a Irvine. Per posizioni \emph{junior} o
\emph{entry level} i grandi studi raramente concedono il lavoro da casa; per
profili senior può esserci qualche eccezione (ad esempio nove mesi da remoto
e tre in sede).

\paragraph{Relocation.} Spesso tocca spostarsi, ma il relatore non ha mai
pagato un trasloco di tasca propria: queste aziende offrono la
\emph{relocation} (aiuto a trovare casa, ad aprire il conto in banca e
supporto pratico), perché è un investimento sulla persona.

\paragraph{Colloqui mentre si lavora altrove.} Per i colloqui in presenza
prendeva ferie, senza dire dove andava. Il consiglio generale: per fare
colloqui \textbf{prendete ferie o permessi}, non datevi malati, ``che non è un
bel modo di comportarsi''.

\paragraph{Remoto ``doppio''.} Un collega lavora per l'azienda coreana che
gestisce \emph{PUBG}, con ufficio ad Amsterdam, ma lavora da Milano: un
``doppio remoto''. Si trova, ma in genere un ambiente anglofono con sedi in
Europa chiederà di trasferirsi in Europa.

\paragraph{Famiglia e trasferimenti.} Il relatore ha avuto la fortuna di
sposare una persona che fa lo stesso lavoro (conosciuta lavorando su
\emph{Warhammer Online}, poi con lui a Blizzard e a Riot, nella
localizzazione e nel doppiaggio, ad esempio di \emph{World of Warcraft} in
italiano): avendo competenze rare, spesso venivano assunti entrambi, al
massimo a pochi mesi di distanza. Per i developer la soluzione frequente è
che la famiglia segua la persona assunta, perché lo stipendio basta per
tutti. Ma non sempre c'è equilibrio, e a volte uno dei due deve sacrificarsi:
sua moglie ha lasciato Blizzard per seguirlo a Dublino (per poi essere
assunta da Riot); lui ha lasciato Riot per tornare in Italia quando lei ha
trovato lavoro qui. Ha dovuto rinunciare più volte a lavorare in Islanda per
CCP (\emph{EVE Online}), perché non era una scelta solo sua. Si fanno
compromessi, ma si riesce a far quadrare.

\paragraph{Ci si sposta per ogni progetto?} No: si sono spostati
esclusivamente per cambiare azienda. Cambiando progetto nella stessa azienda
è raro, quasi mai (``non l'ho mai visto succedere''), che venga chiesto di
trasferirsi. In Francia stavano benissimo (vivevano fuori Parigi; ``i
problemi con i francesi sono in realtà problemi con i parigini''), parlavano
francese ed erano vicini alla doppia cittadinanza: se ne sono andati perché
non vedevano un futuro a Blizzard, per il modo in cui l'azienda trattava i
colleghi.

\paragraph{Il mercato oggi.} Non esiste un \emph{silver bullet} per trovare
lavoro domani dove si vuole. Con tutti i \emph{layoff} recenti ci sono
migliaia di esperti del settore che cercano lavoro: serve pazienza, ma è
fattibile.

\paragraph{Studi indie e remoto.} Negli studi più piccoli il remoto è più
diffuso: a Stormind circa 60 persone su 95 lavorano da remoto, ma lo studio
si è strutturato fin dall'inizio per farlo. Serve consapevolezza. Bending
Spoons (che sviluppa app per mobile, non giochi, ed è l'azienda dell'app
Immuni), pur essendo grandissima e con sede a Milano, lavora da remoto e
lascia l'ufficio a chi vuole: lo smart working è posizionato quasi come un
\emph{benefit}, ma permette anche di avere un ufficio da 300 persone invece
che da 700. Le sedi fisiche costano decine di milioni all'anno (Blizzard ha
chiuso le sedi di Amsterdam e di Parigi). Tuttavia, per lavori molto
verticali e specifici gli studi preferiscono avere le persone dal vivo; i
ruoli più spesso remoti sono quelli che non fanno sviluppo (customer service,
PR\dots).

Anche gli artisti lavorano spesso in presenza: l'accademia può dare solo una
\emph{simulazione} di quanto stretto sarà il contatto tra programmazione,
design e art. ``Sarà ancora di più.''

\paragraph{La comunicazione da remoto.} È più difficile. Pur essendo
estremamente introverso, il relatore \textbf{obbliga tutti a usare la webcam}
nei meeting in cui conta: non in un aggiornamento con novanta persone
collegate, ma certamente in un colloquio a due sulla performance dell'ultimo
anno. La componente fisica è una parte enorme della
comunicazione\footnote{In lezione si parla di ``quasi l'ottanta per cento'';
il dato va preso come indicativo, poiché le percentuali di questo tipo
circolano in forme diverse e sono spesso citate fuori contesto.}: anche se si
è timidi, accendere la webcam è l'approssimazione più vicina all'essere dal
vivo. Negli ambienti dove si lavorava senza webcam la situazione era
``esponenzialmente peggiore''.

\subsection{Differenze culturali nel team}

\paragraph{Gestire sensibilità diverse.} A Riot i team erano divisi per
\emph{lingua}, non per nazione: il team di lingua spagnola includeva anche un
messicano, quello di lingua russa includeva russi e ucraini, proprio nel
periodo della guerra in Crimea. Tra i colleghi non c'era alcun problema (erano
amici), ma ci si è posti la domanda: \emph{come comunichiamo con una community
che sarà divisa su questo tema?} Nei giochi online la divisione era evidente:
in \emph{Rocket League} c'era chi, con un compagno di squadra con la bandiera
avversaria, si faceva otto autogol pur di farlo perdere. Allo stesso modo,
con molti americani mandati a Dublino a coprire posizioni scoperte, c'erano le
elezioni (Clinton--Trump, Biden--Trump) e varie guerre: serviva sempre una
certa sensibilità.

Imparare a navigare queste situazioni lo ha reso un professionista migliore:
gli ha insegnato a fermarsi e chiedersi ``io la percepisco così; gli altri
dieci, con background, storie e lingue diverse, come la penseranno?'', a
ragionare più \emph{long term} e più \emph{big picture}, con più empatia.

\paragraph{Stili di comunicazione.} Un partecipante che ha lavorato nel Regno
Unito chiede dello stile comunicativo italiano, a volte frainteso. Uno dei più
grandi complimenti ricevuti dal relatore è stato ``non sembri italiano''
(``gesticoli poco per essere italiano'', ``parli bene inglese per essere
italiano''). Ha imparato ad \textbf{adattare la comunicazione alla persona}:
con spagnoli e francesi la battuta sulla nazione funziona, perché c'è un
``odio--amore'' condiviso; con russi e giapponesi no. Doveva adattare anche
l'umorismo. Un collega metà russo e metà messicano diceva di aver preso ``il
lato peggiore di entrambe le nazionalità'', ed è stato un ottimo maestro per
imparare a leggere culture diverse.

Non bisogna generalizzare, ma gli stereotipi possono essere un punto di
partenza ragionevole, \emph{sapendo che troverete sempre le eccezioni}. E
vanno letti correttamente: gli americani tendono a essere più estroversi e ad
avere sempre qualcosa da dire non perché ``sono fatti così'', ma perché
vengono educati così (all'università tra i corsi c'è il \emph{debate}, saper
argomentare in pubblico); un collega giapponese può essere più introverso per
ragioni culturali. Nel dubbio si parte da lì, poi si conosce la persona.

\subsection{Negoziare, discriminazione e stipendi}

\paragraph{La famiglia come leva.} Avere un partner riduce le opzioni, ma le
aziende lo sanno: uno spostamento è per due. Si può usare come leva per
negoziare (un compenso diverso, un supporto per il trasloco): se un'azienda ha
deciso di volere quella persona, 2, 3, 4 o 10.000 euro in più all'anno, per
un'azienda che fattura centinaia di milioni o miliardi, non cambiano nulla.

\paragraph{Essere sposati è un motivo per non essere assunti?} No. In alcuni
paesi è proprio illegale discriminare su questo: in Francia era illegale
chiedere CV con data di nascita, foto, sesso o religione, e chiedere se si
vogliono figli. Purtroppo, dai racconti dei suoi studenti di master (in
settori come \emph{beauty and wellness} e \emph{food and wine}), in piccole
realtà italiane queste domande (``hai 25 anni, hai intenzione di
sposarti?'') vengono ancora fatte, e servono a filtrare i candidati. Nelle
grandi aziende ci sono sistemi per evitarlo; anzi, molte danno \emph{più}
benefit a chi ha famiglia (``porta pure la famiglia, è tutto pagato''),
perché le persone competenti hanno spesso una certa età e quindi,
ragionevolmente, una famiglia.

\paragraph{Gli stipendi.} Dipendono molto da dove si vive. Lo stipendio a Riot
non era paragonabile a quello di Blizzard perché in Francia era tassato al
52\% e in Irlanda al 17\%\footnote{Aliquote come riportate in lezione,
riferite alla situazione personale del relatore.}: nessuno stipendio
competitivo può compensare il 35\% di differenza di tasse. Negli Stati Uniti,
con pochissima previdenza pubblica, gli stipendi sono molto più alti perché
poi ci si paga tutto (``se hai un raffreddore paghi tu''). Bisogna valutare il
potere d'acquisto di quei soldi nel paese in cui si andrà a vivere.

\subsection{Freelance e \emph{work for hire}}

All'inizio della sessione pomeridiana uno studente che lavora nell'audio
chiede quali siano le condizioni per chi lavora come \emph{freelance}.
\begin{itemize}
  \item Dipende molto dal progetto: spesso il contratto è a progetto, per
        portare a termine una parte o un intero progetto.
  \item Non tutte le aziende usano freelance: Riot e Blizzard lo facevano
        pochissimo e preferivano assumere (magari trattando poi la persona
        ``come un freelance'', lasciandola a casa a fine progetto).
  \item Spesso i lavori offerti come freelance sono in realtà
        \textbf{sostitutivi di un full time} (``freelance, ma sei ore al
        giorno'').
  \item Si può fare, ma bisogna essere molto chiari nel definire l'impegno
        (quante ore, come vengono tracciate): non tutti specificano cosa
        significa davvero ``freelance''. Lavorando così per due o tre aziende
        si può arrivare a un tempo pieno.
  \item È più comune il \emph{work for hire}: aziende che forniscono personale
        ad altre aziende; si è assunti da X ma si lavora ``embeddati'' presso
        Y (una sorta di subappalto). Anche lo studio del relatore lo fa.
  \item Il puro freelance di quattro mesi esiste, ma di solito per quei
        quattro mesi diventa un lavoro a tempo pieno; è raro ``quattro mesi,
        due ore al giorno''. Nelle grandi aziende è improbabile lavorare nello
        stesso mese per Ubisoft, Riot e Blizzard tre ore ciascuna; è più
        realistico lavorare per più startup indie italiane, che non possono
        permettersi un'assunzione ma possono pagare due ore al giorno.
  \item Molto passa dai contatti: trovare le aziende giuste e rendere esplicite
        le condizioni.
\end{itemize}

% ---------------------------------------------------------------------
\section{Il circolo virtuoso: feedback, miglioramento, fiducia}
\label{sec:z1-triangolo}

\subsection{Il triangolo}

L'obiettivo a cui il workshop vuole arrivare è un \textbf{triangolo che si
autoalimenta}:
\begin{enumerate}
  \item si dà un \textbf{feedback} con l'obiettivo di migliorare le
        performance (inclusi elementi caratteriali);
  \item la persona \textbf{migliora} e vede di averne beneficiato;
  \item questo genera \textbf{fiducia}, che rende più facile dare e ricevere
        il feedback successivo.
\end{enumerate}
Togliendo uno qualsiasi dei tre elementi, il circolo non è altrettanto forte.

In qualsiasi fase si parta, il rispetto per i colleghi è inizialmente dettato
dall'\emph{educazione}, non dalla fiducia. La fiducia arriva quando i colleghi
cominciano a dire cose giuste e sensate e a \textbf{mantenere la parola}.
Dare e ricevere feedback è il modo in cui, in un certo senso,
``inganniamo il cervello'' per creare una \emph{case history} di fiducia: si
ha motivo di fidarsi di più, di ascoltare di più, di essere più ricettivi. E
si diventa sempre più bravi ad accettare il feedback e ad agire di
conseguenza.

\subsection{Accettare un feedback non vuol dire fare come dice l'altro}

Dieci anni fa il relatore non era ostile ai feedback, ma molto
\emph{difensivo}: ``sì, sì, va bene, ora ci guardo''; oppure, a chi gli diceva
``avresti dovuto farla così'', rispondeva spiegando perché l'aveva fatta in
quel modo. Oggi ha imparato a dire ``ok, grazie, ora ci ragiono'': non deve
difendere la propria posizione da chi viene solo a dire ``hai pensato a
questa cosa?''. Prima lo viveva come un attacco (``mi viene a dire che non so
fare il mio lavoro''), ora come un ``ottimo punto di vista, lo prendo in
considerazione''.

Con una postilla fondamentale: \textbf{accettare il feedback non significa
fare come dice l'altro}. ``Mi hai detto che preferisci questo mob verde. Ho
pensato anche a farlo verde? Sì. Penso comunque che giallo sia meglio? Sì, lo
faccio giallo.'' Il relatore accetta \emph{sempre} il feedback, valido o no:
garantisce di tenerlo a mente quando lavora. Che poi faccia come gli viene
detto è un altro discorso.

\subsection{Fiducia e scadenze}

La fiducia funziona allo stesso modo con gli impegni. Se si chiede a qualcuno
una cosa per venerdì e la consegna venerdì, la volta dopo gliela si chiederà
per venerdì. Se invece arriva il giovedì della settimana dopo, la volta dopo
gli si darà una \textbf{deadline anticipata}, come con chi arriva sempre in
ritardo agli appuntamenti (``se ti dico alle 7 arrivi alle 7 e mezza; la
prossima volta ti dico alle 6 e mezza'').

Al lavoro questo è dannosissimo, perché crea \textbf{deadline fittizie} che
non sono mai quelle vere, e nascono discussioni del tipo ``perché devo farlo
in due giorni se si fa in cinque?'' (``sì, se tu lo facessi in cinque''). La
conversazione giusta è un'altra: \emph{``che problemi hai avuto nel
consegnare venerdì?'' ``Questo, questo e questo.'' ``Cosa posso fare per
evitare che succeda?''}

Il relatore non vuole mai sentirsi dire ``non sono riuscito a farlo''.
Preferisce di gran lunga sentirsi dire, tre giorni prima, \textbf{``potrei non
riuscire a farlo, per questi motivi''}. Chi ha la visione d'insieme può dare
più budget, più persone, spostare la deadline, cambiare lo scopo del progetto:
può fare qualsiasi cosa per aiutare. Ma se lo si scopre il venerdì, cinque
minuti prima della deadline, non può fare più niente. Il messaggio è:
\textbf{``aiutami a poterti aiutare''}. Non penserà mai male di chi dice ``non
ce la faccio per questi motivi'', perché non sta dicendo ``non ho voglia'': sta
indicando i problemi, e lavorare su quei problemi è compito del manager. Non
può invece fidarsi di chi fino al giorno prima dice ``tranquillo, ce la
faccio'' e il giorno dopo ``non ce l'ho fatta''.

Questo è fondamentale per chi lavora in programmazione, design o art, perché
ci saranno sempre deadline: bisogna sempre \textbf{lavorare sulle aspettative},
dire al proprio manager (o al proprio report) cosa è lecito aspettarsi. Non
bisogna fare come i pinguini di \emph{Madagascar} (``sorridete e salutate, va
tutto bene''): se c'è un problema, va detto, altrimenti si peggiora solo la
propria situazione.

Da italiani si tende a non voler mai alzare la mano per chiedere aiuto
(``no, no, ce la facciamo, lo faccio io''). A volte non ce la si può fare, non
perché non si è capaci, ma perché vari fattori lo impediscono. Accettarlo è
importante: \textbf{è molto meglio avvisare e chiedere una mano che dire
``non ce l'ho fatta''}.

\subsection{Domande sul dare e ricevere feedback}

\paragraph{``Ne parliamo a fare, se poi fai come vuoi?''} Uno studente (Cristian)
riceve feedback che, dopo averli analizzati, non ritiene congrui; quando non
li segue gli viene detto ``ne parliamo a fare?''. C'è un problema di
\textbf{aspettative}: l'idea che accettare un feedback voglia dire fare come
dice l'altro. Accettarlo significa prendere in considerazione ciò che è stato
detto. Chi decide diversamente ha però un \textbf{onere}: argomentare perché
si fa comunque a modo proprio (``mi hai detto verde, la facciamo gialla
perché\dots''). Poi bisogna solo aspettare i risultati per vedere chi aveva
ragione.

La decisione spetta sempre \textbf{a chi ha la responsabilità di quella
decisione}. Se il game director deve scegliere tra A e B, anche se il
feedback dice B, la scelta è sua perché la responsabilità è sua. Non sarebbe
corretto essere ritenuti responsabili di una scelta delegata a chi dà
feedback. La risposta può essere semplicemente: ``l'ho tenuto in
considerazione, ma visto che ne sono responsabile io e le conseguenze sono
mie, ho scelto di fare X''.

\paragraph{E se l'altro fa ripicca?} Un partecipante chiede di chi, vedendo
il proprio feedback non accolto, si vendica ignorando a sua volta feedback
buoni. Per il relatore è un problema caratteriale: chi si comporta così ``è
già sulla strada per uscire''. Non c'è spazio per l'ego in team di questo
tipo. Il fenomeno nasce dall'idea che \emph{il feedback negativo faccia fare
brutta figura}, che il fatto che si potesse fare meglio tolga qualcosa alle
proprie capacità.

Non è così: \textbf{il feedback non è un giudizio}. Si può passare un anno a
Riot ricevendo solo feedback negativi e avere una valutazione annuale
stellare, perché ricevere molti feedback negativi significa essere una persona
approcciabile, che agisce sul feedback e invita gli altri a darlo. Il relatore
ha visto valutazioni molto peggiori per chi riceveva solo feedback positivi,
perché la gente aveva paura di dargliene di negativi: ma solo ``pacche sulle
spalle'' non fanno migliorare. Il problema va affrontato parlando con la
persona; a volte le conversazioni iniziano con \textbf{``ti do un feedback sul
feedback''}, cioè sul modo in cui l'altro ha dato il suo.

Esempi dalla sua esperienza:
\begin{itemize}
  \item lui chiede spesso di \textbf{non ricevere feedback senza esempi
        tangibili}, perché fa fatica ad astrarre: ``sei stato maleducato'' non
        significa nulla, ``ieri al meeting mi hai risposto male su questa
        cosa'' sì;
  \item ad altri serve l'opposto: ha ricevuto il feedback ``quando devi dirmi
        che ho sbagliato, puoi girarci un po' intorno e darmi del contesto?'',
        perché il suo stile diretto (``questa cosa non ha funzionato, dovremmo
        fare così'') feriva. Ha adattato il suo modo di dare feedback: con
        alcuni è diretto, con altri gira intorno. L'importante è saperlo.
\end{itemize}

\subsection{Non portare la soluzione}

Un consiglio centrale: \textbf{quando si dà un feedback, non si dà una
soluzione}. Non si dice ``io farei così'', ma ``secondo me questa cosa si
potrebbe fare diversamente, per questo e questo motivo: come potremmo farla
diversamente?''. È come chi smette di fumare perché vuole e chi smette perché
glielo dicono moglie e figli: se non si è convinti, spesso non si smette.

Quando il relatore diceva ``c'è questo problema, io farei così'', il meeting
successivo iniziava inevitabilmente con ``ci ho provato come hai detto tu, con
me non ha funzionato''. Sempre. Perché la soluzione era sua. La
responsabilità della soluzione deve essere invece \textbf{congiunta}: ogni
feedback deve creare un piccolo team che ha come obiettivo risolvere quel
problema. Questo toglie l'antagonismo (non ``ti dico che non sai fare il tuo
lavoro'', ma ``ti voglio aiutare a farlo meglio, lavoriamoci insieme'') e
sposta il focus: il problema non sei tu, è ciò che sta accadendo.

\paragraph{E se sono i piani alti a tenere all'oscuro dei problemi?} È una
questione di cultura aziendale. Il relatore fa sempre presente al proprio capo
il concetto di \textbf{service level agreement} (che verrà ripreso nella
seconda giornata): ``se mi dai questi input, ti garantisco questo output in
questi tempi''. Se il capo chiede delle cose senza dire budget, tempi e
perché, non può aspettarsi lo stesso risultato. È una comunicazione a due
vie: si dice al manager cosa serve per fare il lavoro, e lui deve fornirlo o
accettare che il risultato non sarà quello voluto.

\subsection{Amici, parenti e parole proibite}

\paragraph{``Non si lavora con amici e parenti''?} Per sua definizione
personale il relatore ha pochissimi amici (``stanno sulle dita di una mano'')
e una soglia molto alta per considerare amico un collega. Allo stesso tempo,
in quindici anni e con centinaia di colleghi, ha avuto solo \textbf{due}
esperienze in cui non è riuscito a lavorare con qualcuno: non serve essere
amici per far funzionare le cose. Con gli amici, anzi, gli viene più facile:
se dice a un amico che una cosa ``fa schifo'', l'amico non pensa ``come si
permette'', ma ``se perfino Stefano, che è mio amico, me lo dice, allora è
proprio così''. Diventa un rafforzativo, grazie alla fiducia costruita prima.
Il problema del lavorare con amici e parenti nasce dal \emph{non voler} dire
all'amico che sta sbagliando; se invece la confidenza permette di dirlo, si
può saltare la ``confezione'' del messaggio e passare direttamente al
contenuto.

\paragraph{Le frasi peggiori.} Alla domanda se esistano parole peggiori di
``ottimo lavoro'', il relatore indica la sua classifica: ``ottimo lavoro''
(detto in modo vuoto)\footnote{Il passaggio della trascrizione è
frammentario: si intende verosimilmente il complimento generico e di
circostanza, che non dice nulla di utile.}; poi la frase che lo ``fa uscire di
testa'', \textbf{``abbiamo sempre fatto così''} (``e infatti avete chiuso'',
come Blockbuster); infine \textbf{``siamo una famiglia''}, riferito al posto di
lavoro. Un'azienda non è mai una famiglia, perché non ci si può aspettare lo
stesso impegno da chi la possiede e da chi ci lavora: non c'è un
corrispettivo proporzionale tra impegno e ritorno. Per quanto inclusiva,
flessibile e razionale, non sarà mai una famiglia.

\subsection{Ego, muri e il peggior modo di scegliere un leader}

\paragraph{Chi mette davanti l'ego e alza muri.} Esistono, ma di solito durano
pochissimo in azienda. Non si tratta di avere un ego forte, ma di \textbf{non
saper ricevere feedback}. Spesso sono persone bravissime nel loro lavoro che
poi si chiedono perché tutti intorno diventino più bravi di loro: perché gli
altri accettano feedback e hanno più fonti da cui attingere, mentre loro si
sono rintanati nella loro torre.

Molti manager tendono ad \textbf{assumere persone un filo peggiori di loro},
per non fare brutta figura. Il relatore assume solo persone più brave di lui,
``perché non è una gara'': ha lavorato con molti community manager più bravi
di lui come community manager, ma lui era più bravo a fare il manager. Non
viene valutato su quanto sa fare il community manager, ma su quanto il suo
team sa farlo. Un partecipante chiede se la paura di avere qualcuno di
migliore esista anche in ambito accademico. Assolutamente sì: al relatore è
capitato di non essere preso come docente in un corso di formazione sui
videogiochi perché il responsabile aveva percepito che ne sapesse più di
lui\footnote{La trascrizione dice ``che non sapevo più di lui'', ma il senso
del passaggio richiede il contrario.}. ``Ma perché te ne importa? Tu sei quello che ha portato chi
sa fare il suo lavoro: ben venga.'' O si ha questa cultura, o si ha paura per
il proprio posto.

\paragraph{Promuovere il migliore.} Quando bisogna scegliere un leader per un
team (programmazione, design, art), molto spesso si sceglie \textbf{il migliore
tecnico}: il miglior programmatore, il miglior designer. È spesso la scelta
sbagliata: è come prendere un idraulico per fare l'elettricista. Fare il
\emph{people manager} non è come fare il designer o il programmatore; si può
essere un ottimo programmatore e un pessimo manager, e viceversa. Le aziende
l'hanno capito e lo fanno meno di prima. A Blizzard c'era un grave problema di
\emph{middle management}: molte persone in ruoli importanti erano state
promosse cercando le skill sbagliate, erano frustrate dal non saper fare il
proprio lavoro e schiacciavano chi stava sotto per paura di fare brutta
figura.

Un partecipante cita una teoria simile: una persona viene promossa finché non
raggiunge \emph{il proprio livello di incompetenza}, dove si blocca\footnote{È
il \emph{principio di Peter} (Laurence J.\ Peter, 1969), citato in lezione
come ``la teoria di Peter''.}. Un altro cita Michael Scott di \emph{The
Office}, su cui il relatore tornerà (sezione~\ref{sec:z1-quando}).

% ---------------------------------------------------------------------
\section{Perché la cultura conta: retention, costi e fiducia}
\label{sec:z1-cultura}

\subsection{L'ambiente di lavoro e la retention}

Perché tutto questo è così importante? Perché l'ambiente di lavoro conta
moltissimo per la \textbf{retention}, cioè per trattenere le persone e farle
lavorare in modo efficiente. Se contasse solo lo stipendio, molti non
lascerebbero mai Riot o Blizzard; e più si va avanti con l'età, più si impara
a dare valore anche al resto.

Inoltre, secondo gli studi citati nelle slide, chi vive il proprio lavoro in
modo negativo è \textbf{meno efficiente nel problem solving}: l'energia
mentale va verso ``come faccio a non farmi dare la colpa'' o ``come faccio a
fare bella figura'', non verso la soluzione del problema. In ambienti negativi,
dove ci si deve preoccupare della propria ``sopravvivenza'', si lavora
peggio.

\subsection{Il costo di chi se ne va}

Ogni volta che qualcuno si dimette o viene licenziato, il costo di trovare
una nuova persona, fare i colloqui, formarla e farle attraversare forming,
storming, norming\dots\ è circa \textbf{un terzo del suo stipendio}. Per questo
le aziende tendono a fare più \emph{retention} che \emph{acquisition}, almeno
con le figure considerate utili alla produzione.

Allo stesso tempo si vedono sempre più spesso \emph{layoff} e \emph{hiring
wave}: si licenziano 3.000 persone e sei mesi dopo se ne riassumono 2.500 o
3.500, ma tutte più junior. Il relatore ne ha scritto in un articolo su
Eurogamer.it, in cui spiega anche alcuni ``trucchi del mestiere'': perché i
licenziamenti avvengono sempre in certi periodi e sempre ``a gruppi'' (prima
Bethesda, poi Blizzard, poi un'altra\dots).

\paragraph{L'offerta di Riot.} A Riot esisteva una pratica (di cui il relatore
non ricorda il nome) per cui, firmato il contratto, si aveva \textbf{un anno
di tempo per dire ``non mi piace, non me la sento, me ne vado''} e ricevere
subito come buonuscita un anno di stipendio. Avevano stimato che fosse molto
meglio pagare subito un anno a qualcuno e lasciarlo andare che tenerlo due,
tre, quattro anni come \emph{underperformer}: la responsabilità della scelta
di assumere era dell'azienda\footnote{La pratica, nota come ``The Offer'', è
descritta in fonti pubbliche con condizioni diverse (un periodo più breve e
una somma pari a una frazione dello stipendio annuo). Si riportano le
condizioni come raccontate in lezione; andrebbero verificate.}. L'hanno usata
pochissimi, ma il solo fatto che esistesse faceva lavorare meglio anche lui da
manager, perché sapeva che chi non si trovava bene poteva andarsene con un
anno di tempo per trovare un altro lavoro.

\paragraph{Il ricircolo.} C'è un forte ricircolo nell'industria: molti ex
colleghi passano da Ubisoft a Paradox, da Paradox a Funcom, da Funcom a
Blizzard. Il relatore ha scelto di uscirne lavorando solo per aziende
\textbf{indipendenti}: Stormind non è un indie di cinque persone, ma è
indipendente, e non deve accontentare azionisti. Quando vede un'azienda con
utili miliardari licenziare migliaia di persone si chiede ``dov'è la linea''
(\emph{when enough is enough}): con azionisti che vogliono un ritorno, se
quello è l'obiettivo stanno facendo la cosa giusta, ma lui non condivide
l'obiettivo\footnote{In lezione si parla di ``6,5 miliardi di utile'' e di
``3.000 persone licenziate''; le cifre sono riportate come dette e non sono
state verificate.}.

\paragraph{L'open space vuoto.} A Riot, durante un'ondata di licenziamenti,
del suo team di trenta persone è rimasto solo lui: ``teniamo te come
coordinatore, poi ricostruirai il team''. Ha rifiutato di ricostruire un team
cancellato ``con un colpo di spugna'' e ha deciso che non avrebbe più lavorato
per una grande azienda. Nessuno aveva pensato di spostarlo, e ha lavorato per
tre mesi in un open space di cento metri quadri con la sua scrivania e altre
ventinove vuote. Su una lavagna aveva scritto i nomi dei suoi ragazzi e li
cancellava man mano che, anche grazie ai suoi \emph{referral}, trovavano
lavoro. Cancellato l'ultimo nome, si è dimesso: per lui il lavoro era finito.

\paragraph{Team che si dissolvono.} Uno studente cita un articolo su Larian
Studios, il cui management è contrario a ricostruire ogni volta il team a fine
progetto. È lo stesso problema: si pensa ai ``giochi Blizzard'' o ai ``giochi
BioWare'', ma le persone che li hanno fatti non ci sono più; è rimasto solo il
brand, e non ha senso avere le stesse aspettative da un gioco fatto da persone
diverse. E anche chi sopravvive ai layoff vive pensando che il prossimo
potrebbe essere lui: alcune persone, licenziate nell'ultima ondata a
Blizzard, gli hanno detto ``erano due anni che lo sapevo''. Vivevano a Cork,
``un paesotto irlandese dove o vai per lavorare a Blizzard o non ci vai'',
avevano comprato casa: non ci si può aspettare chissà che morale.

\subsection{La cultura è forte quando si lavora per gli altri}

\begin{quote}
  \emph{La cultura è forte quando le persone lavorano per gli altri e a favore
  degli altri, e non per se stesse.}
\end{quote}

Se si dà feedback per far migliorare l'altro e non per imporre il proprio
modo, se si aiuta per il bene del progetto e non per fare bella figura, la
cultura si rafforza. È dai \textbf{piccoli gesti} che emerge la cultura
aziendale: ad esempio, dal non dover chiedere ``perché mi stai aiutando?''.

\paragraph{L'incognita.} \textbf{L'incognita è spesso dannosa quanto le
cattive intenzioni.} In un esperimento (``mi pare a Stoccolma, cinque o sei
anni fa'') si distribuivano per strada banconote da 500 euro, e la maggior
parte delle persone le rifiutava chiedendosi ``dov'è la fregatura?''
\footnote{L'esperimento è riportato in modo vago in lezione e non è stato
possibile identificarlo con certezza.}. Lo stesso accadeva al relatore,
arrivato a PG e a Stormind da Riot: a ``vedo che sei in difficoltà, posso
aiutarti?'' la risposta era sempre ``perché?''. Alla fine rispondeva: ``perché
quando sarò in difficoltà io vorrei che tu mi aiutassi''.

Quando si comincia a lavorare per gli altri e si smette di pensare ``così avrò
la promozione, così avrò l'aumento'', promozioni, aumenti e rispetto arrivano
da soli. A Blizzard aveva fallito tre volte i colloqui per essere promosso a
manager (inteso come coordinatore del team); quando ha deciso di smettere di
pensarci e di fare solo il bene del proprio lavoro, la volta dopo è stato
promosso. Funziona così: chi aiuta i colleghi e fa andare bene il progetto
viene notato (``questa persona fa migliorare cinque persone, non solo se
stessa: vale di più'').

\subsection{Il lavoro è valutato come opera d'insieme}

Il relatore riconosce che buona parte del suo successo è venuto dal fatto che
il suo team lavorava bene, non dalla sua bravura personale. Anche i progetti
degli studenti, presentati a Diego, a Enrico, ad Andrea, verranno valutati
\textbf{come opera d'insieme}. Nessuno dice ``il \emph{combat system} di questo
gioco l'ha fatto X, che bravo'': si dice ``BioWare ha fatto un bel gioco'',
``Blizzard ha fatto un bel gioco''.

Lo ha replicato anche lato HR. In molte aziende ci sono \textbf{bonus
personali} legati a obiettivi (se il gioco vende X, se la feature è pronta in
tempo). Il relatore li ha dimezzati: metà restano personali, metà diventano
\textbf{bonus di team}, che tutti ricevono se il team raggiunge gli obiettivi.
Un incentivo concreto a collaborare: non solo si dice che il lavoro di gruppo è
importante, lo si dimostra dandogli più valore economico. Aveva proposto il
100\% di bonus di team, ma gli è stato concesso il 50 e 50; ha funzionato
comunque, anche se sarebbe stato ancora più potente. Ovviamente, alla
domanda ``perché mi aiuti?'' si può anche rispondere ``perché voglio il
bonus'': un aiuto dato ``per il motivo sbagliato'' era un compromesso che era
disposto ad accettare.

\subsection{Il coraggio di chiedere aiuto}

Un partecipante osserva che aiutarsi dovrebbe essere la norma. Il modello
diffuso, invece, è ``devi essere forte e indipendente, non chiedere mai una
mano'', e vale in tutte le direzioni (dal maschilismo alla ``donna forte e
indipendente che non ha bisogno di nessuno''): \textbf{non ammettere di avere
bisogno di aiuto è sempre uno dei motivi per cui le aziende vanno male}, che si
tratti di uomini, donne o persone non binarie. Il problema è non avere la
sicurezza di poter alzare la mano.

Il relatore ha cercato di togliere ogni conseguenza dall'aver alzato la mano:
se si \emph{depenalizza} qualcosa, le persone non sono più disincentivate a
farla. Ha detto espressamente che non avrebbe mai valutato nessuno per il
numero di volte in cui chiedeva aiuto o per il numero di feedback negativi
ricevuti, e questo ha funzionato.

\subsection{Perché si dà un feedback}

\emph{Perché avete dato il vostro ultimo feedback? Che cosa volevate ottenere?}
Le risposte della classe:
\begin{itemize}
  \item per \textbf{riconoscere il valore} di un collega, anche verso il
        manager, perché qualcuno che non aveva visibilità vedesse quanto
        quella persona fosse efficace;
  \item per \textbf{spingere una persona a superarsi};
  \item per aiutare un collega a rendere il prodotto all'altezza di ciò che
        voleva rappresentare;
  \item per \textbf{genuinità}: c'era qualcosa che poteva essere fatto meglio.
\end{itemize}
Il relatore aggiunge che la forma cambia col tempo: all'inizio era ``devo
darti un feedback, hai dieci minuti?''; ora è un messaggio informale su
Discord, perché ha costruito un rapporto di fiducia e sa che non verrà letto
come ``l'HR manager mi scrive, sono nei guai?'', ma come ``Stefano è così, mi
dà feedback, resterà tra noi e lo risolveremo insieme''.

Un altro motivo è \textbf{risolvere incompatibilità caratteriali}: le
conversazioni più difficili, del tipo ``io e te non stiamo funzionando, cosa
possiamo fare diversamente?''. Sono importanti perché, se due responsabili non
vanno d'accordo, il team se ne accorge, comincia a raccontarsi storie (``si
odiano, ci saranno problemi'') e, per difendere il proprio manager, comincia a
giudicare male il lavoro dell'altro gruppo. È la dinamica da fermare:
\textbf{valutare gli altri solo per difendere la propria fazione}.

Dare e ricevere feedback sono \textbf{due muscoli da allenare} (il relatore è
sempre stato più bravo a darlo che a riceverlo), e si impara col tempo. È come
fare colloqui: fare colloqui è una skill a parte, che non indica quanto si è
bravi nel proprio lavoro. Essere scartati non significa non saper fare quel
lavoro: spesso significa aver fatto male il colloquio o che è stato preferito
qualcun altro.

\paragraph{Feedback pubblici e privati.} Uno studente racconta di aver dato un
feedback negativo a un superiore, senza un motivo particolare, per dimostrare
a sé e ai colleghi che non si viene licenziati per questo. Il relatore coglie
l'occasione per precisare che \textbf{i feedback negativi si danno
individualmente}: lui dà moltissimi feedback di cui il resto dell'azienda non
sa nulla. Reagirebbe molto male se, in un meeting, qualcuno alzasse la mano per
dire ``comunque questa cosa non l'hai fatta bene'': quello non è feedback, è
discussione, accettabile solo se è un brainstorming e se l'opinione è stata
chiesta. Alcune aziende hanno piattaforme pubbliche di feedback, ma spesso
servono solo ad applaudire; lo stesso accade nelle aziende che
\emph{obbligano} a lasciare un certo numero di feedback a settimana, che
diventano ``lui è bravo, lei è competente, anche questa settimana l'abbiamo
fatto''. Bisognerebbe chiedere a chi le adotta cosa pensa di ottenere (forse
trasparenza), ma difficilmente una persona ``attaccata'' pubblicamente reagisce
bene.

% ---------------------------------------------------------------------
\section{Among Us: la caccia al traditore}
\label{sec:z1-among-us}

Prima della pausa pranzo il relatore mostra tre film di fantascienza di
qualche decennio fa e chiede che cosa li accomuni\footnote{I titoli non sono
leggibili nella trascrizione: viene citato \emph{Event Horizon} da uno
studente, e il film ``in mezzo'' è indicato come ``l'originale'',
verosimilmente \emph{La cosa} (\emph{The Thing}).}. In ognuno c'è un gruppo di
persone con un obiettivo e la presenza, reale o percepita, di qualcuno che
rema contro: \textbf{l'energia del gruppo viene spesa per trovare il
responsabile}.

\paragraph{L'esperimento degli alieni.} Nei master in presenza il relatore
propone un esercizio. Distribuisce segretamente cartellini plastificati con
scritto ``essere umano'' o ``alieno''; le carte alieno contengono istruzioni di
comportamento con cui gli alieni, che cercano di fingersi umani, si
tradirebbero (dare sempre ragione a chi ha parlato prima, finire ogni frase
con una domanda, parlare in un certo modo). Lo scopo: in dieci minuti formare
gruppi, e \textbf{vince il gruppo più numeroso senza alieni}; basta un alieno
per perdere. Si formano sempre gruppi piccoli (due, cinque, sei persone),
perché la tendenza è ``mi fido al massimo di voi due''.

Ma vincono i gruppi più grandi, per un motivo semplice: il relatore
\textbf{non distribuisce nessuna carta alieno}. Mostra a tutti entrambe le
carte, ma consegna solo carte ``umano''. Vuole simulare ciò che accade in
ufficio: la tendenza a percepire sempre qualcuno che rema contro (i manager,
l'altro team, quel collega), quando in realtà non c'è nessuno, ``è tutto nella
nostra testa''. Quando chiede le motivazioni, emergono sospetti
\emph{fabbricati}: ``faceva sempre domande'', ``parlava solo all'infinito''
(``no, parla così e basta''). Ci si fabbrica da soli i motivi per non fidarsi
degli altri, ed è comunissimo.

\paragraph{Among Us.} L'autosabotaggio è, per tornare ad \emph{Among Us},
\textbf{concentrarsi sul cercare cadaveri e traditori invece di fare le
task}: non si vince perché non si sono completate le task, mentre si cercava il
colpevole. Succede tantissimo negli uffici, ed è il motivo per cui l'80\% dei
problemi che il relatore risolveva erano di comunicazione. Entrando nei team
sentiva immancabilmente: ``facciamo il possibile, ma questa lotta continua con
l'altro gruppo è estenuante''. ``Perché c'è questo problema?'' ``Non lo so,
c'era già quando sono arrivato.'' Un paradosso: c'è un problema, nessuno sa
dire quale sia, ma c'è sempre stato e quindi lo si accetta.

\paragraph{L'esperimento delle scimmie.} È l'esempio, ``finto o reale'', delle
scimmie in gabbia\footnote{È il noto racconto delle ``cinque scimmie'', spesso
presentato come esperimento scientifico ma probabilmente apocrifo (lo stesso
relatore lo definisce ``finto barra reale''); nella trascrizione il numero
delle scimmie varia.}. In una gabbia con alcune scimmie pende una banana
elettrificata: chi la tocca fa prendere la scossa a tutte. Le scimmie imparano
e, quando una si avvicina alla banana, le altre la picchiano. Poi si
sostituisce una scimmia alla volta con una nuova: la nuova va verso la banana,
le altre la picchiano, e smette. Si continua finché nella gabbia ci sono solo
scimmie nuove, poi si toglie la corrente. Nessuna scimmia tocca la banana:
nessuna ha mai preso la scossa, ma tutte hanno imparato che chi tocca la
banana prende le botte.

Nei team è lo stesso: il motivo per cui si ``odia'' l'altro team si
trasforma, ma un vero motivo non c'è. Ne derivano la sfiducia, il non alzare
mai la mano per chiedere aiuto, il non dire mai ``bel lavoro'' agli altri
perché non sia mai che qualcuno ti veda fare un complimento ``a quelli''.
Risolvere queste dinamiche è lungo e tedioso, e sempre pieno di ``non mi
ricordo'', ``io non c'ero'' e ``abbiamo sempre fatto così'', la cosa più
frustrante per chi cerca di risolverle.

% ---------------------------------------------------------------------
\section{Il principio del default to trust}
\label{sec:z1-default-trust}

\subsection{La definizione}

Il relatore indica questa come la slide più importante della presentazione.
\textbf{Default to trust} è uno dei valori aziendali di Riot Games, che lui ha
tradotto, adattato e ``esportato'':

\begin{quote}
  \emph{Approccio ogni situazione di conflitto partendo dal presupposto che una
  persona abbia agito in un certo modo perché, con le informazioni a sua
  disposizione in quel momento, ha pensato che fosse la soluzione migliore, e
  non perché sia un'incapace o volesse danneggiare il progetto
  intenzionalmente.}
\end{quote}

In altre parole: \textbf{di default, mi fido}.

\subsection{L'esempio del ritardo al daily stand-up}

A Riot il relatore aveva un \emph{daily stand-up} dal vivo alle 15 con il suo
team di venti persone. Una ragazza arrivava regolarmente alle 15:15, per una,
due, tre settimane. Alla quarta ci ha parlato. \emph{Come iniziereste la
conversazione?} Le proposte della classe:
\begin{itemize}
  \item sottolineare l'importanza di esserci dall'inizio, perché potrebbe
        perdersi qualcosa che le interessa, e per rispetto dei colleghi
        presenti;
  \item ``devi arrivare in orario'', spiegando perché è importante;
  \item \textbf{chiedere perché}.
\end{itemize}

Il relatore osserva, anche dalle sue consulenze fuori dai videogiochi, che
\textbf{più si sale con la seniority}, più la risposta è ``devi smettere di
arrivare in ritardo'', spesso con un \emph{or else} finale (``altrimenti
succede quest'altra cosa''). Più si scende con l'età, e quindi con meno
abitudini consolidate, più la prima domanda diventa ``dimmi perché arrivi in
ritardo''. Già spiegare perché è importante è ``miglia avanti'' rispetto al
semplice ``smetti''.

Partire dal perché gli ha permesso di scoprire la causa: il meeting precedente
della ragazza, in un'altra zona del campus a un quarto d'ora di distanza,
finiva sempre tardi, e lei doveva attraversare il campus di corsa. Non aveva
la confidenza per dire all'altro manager, che insisteva perché si restasse
fino alla fine, che doveva andare; sacrificava quindi il meeting del relatore,
sapendo che era il più comprensivo dei due.

Il relatore si è offerto di \textbf{mediare}: ha detto all'altro manager che
o anticipava l'inizio del suo meeting, così da finire comunque in tempo,
oppure lasciava uscire la ragazza, perché quel meeting gli serviva. C'è anche
una questione di \textbf{economia di scala}, spesso sottovalutata: se venti
persone aspettano qualcuno per dieci minuti, si sono sprecati dieci minuti per
venti. Per lo stesso motivo il workshop è iniziato alle 10:05: iniziare alle
10:10 avrebbe significato far aspettare cinquanta persone cinque minuti in
più. A Riot c'era una scrivania vuota con un pupazzo gonfiabile, \textbf{Mr.\
Late}: sommando tutti i minuti di ritardo a meeting, call e conversazioni, in
un anno veniva fuori praticamente \emph{una persona a tempo pieno}. Serviva a
ricordare che anche cinque minuti contano, e che bisogna rispettare gli altri
arrivando in orario o almeno avvisando.

\paragraph{Gli effetti.} La conversazione era avvenuta nel primissimo periodo a
Riot, quando il team non aveva ancora motivo di fidarsi di lui. Risolvendo il
problema ha eliminato anche tutto lo strato di micro-dinamiche spiacevoli che
un manager impara a vedere (due colleghi che si guardano con la risatina
quando lei entra in ritardo). La ragazza ha cominciato a entrare tranquilla,
con una certa leggerezza: non le faceva piacere arrivare in ritardo (``ci sono
anche persone che lo fanno apposta, e quello è un problema da affrontare
diversamente''). Soprattutto, con i colleghi non parlava di lui come del capo
che minacciava, ma come di chi le aveva risolto un problema: era diventato
un \textbf{alleato}. La settimana dopo altre tre persone sono andate da lui
con un problema da risolvere. Ripetuto cinque, sei, sette, dodici volte, in un
paio di mesi poteva dire quasi qualsiasi cosa nei meeting sapendo che
richieste e critiche sarebbero state percepite come ``me lo chiede perché vuole
darmi una mano'', non per fare bella figura o per presentare i successi del
team ai suoi superiori.

\subsection{Condividere i successi, schermare i fallimenti}

Un altro elemento fondamentale (ripreso nella seconda giornata): il manager
\textbf{condivide i successi e protegge il team dai fallimenti}. ``Se va male
qualcosa è colpa mia; se va bene qualcosa, voglio che ne parliate voi.'' Il
lavoro del manager non dovrebbe essere valutato su ciò che presenta del lavoro
del team: tutti sanno comunque chi è il manager. Tanto vale lasciar presentare
il team, che potrà dire con orgoglio di aver presentato a Mike
Morhaime\footnote{Cofondatore ed ex presidente di Blizzard; nella
trascrizione ``Mike Morayme''.} la propria campagna europea per
\emph{StarCraft~II}. Ci sono manager che dicono ``come mi prendo le colpe, mi
prendo i meriti'', ma i meriti li hanno già: è inutile ostentarli. Con i
fallimenti invece il manager è la \textbf{rete di sicurezza} del team. Il
relatore ne ha fatto una questione di principio: \textbf{non scaricare mai la
colpa su qualcun altro}.

\paragraph{Colpevole o causa.} Un esempio che il relatore attribuisce a Pirate
Software, creatore di contenuti ed ex dipendente Blizzard\footnote{Il relatore
aggiunge che si tratterebbe del figlio dello sviluppatore a cui si sarebbe
ispirato il personaggio del giocatore ``senza vita'' nell'episodio di
\emph{South Park} dedicato a \emph{World of Warcraft}. L'informazione è
riportata come detta in lezione e non è verificata.}. Contattando il servizio
clienti Blizzard per motivi personali, aveva scoperto che i colleghi gli
avevano scritto nei ticket informazioni confidenziali che non avrebbero dovuto
condividere in quel modo, nemmeno con un collega. Non ha segnalato le persone
da punire: ha detto ``c'è un problema, evidentemente le policy non sono
chiare''. Quando gli hanno chiesto i nomi, ha rifiutato: ``se vi dico chi è,
punite quelli, e gli altri cinquecento continueranno a farlo; non avremo
cambiato la causa alla fonte''.

\begin{quote}
  \emph{È inutile trovare il colpevole. Voglio trovare la causa. Perché il
  colpevole cambia.}
\end{quote}

Il colpevole a volte è vittima della situazione: per questo la definizione
dice \emph{``con le informazioni a sua disposizione in quel momento''}. Se il
relatore avesse rimproverato la ragazza in ritardo, non solo non avrebbe
risolto il problema, ma avrebbe creato un'\textbf{ingiustizia}: lei, che
attraversava il campus di corsa per una colpa non sua e di cui non poteva
parlare, sarebbe stata anche rimproverata.

Nessuno fa un brutto lavoro apposta: nessuno è contento quando sbaglia o
perde. Perché al lavoro diamo per scontato che gli altri sbaglino
intenzionalmente? Se qualcuno ha fatto una cosa male, è perché in quel momento
credeva fosse la scelta migliore. La domanda giusta non è ``perché hai fatto
questa cosa che è stata un fallimento?'', ma:

\begin{quote}
  \emph{Che cosa ti ha fatto pensare che questa fosse la scelta giusta?}
\end{quote}

Magari erano limiti di budget, di tempo, di scopo: si agisce su quelli, non
sulla persona. Se si cambia la persona ma le si danno gli stessi fattori, è
molto probabile che arrivi alle stesse conclusioni, e ci si troverà a dover
``disciplinare'' anche lei.

\paragraph{La reputazione costruita sui rumor.} Durante il giro di
presentazione in un ufficio (non in nessuna delle aziende citate), al relatore
è stato detto: ``questo è il dipartimento di design; se ti serve qualcosa
chiedi a lei, non a lui, perché lui consegna sempre in ritardo''. Oltre a
trovarlo poco professionale, una settimana dopo ha chiesto: ``cosa ti ha fatto
questa persona?'' ``A me niente, però ho sentito che\dots''. Si torna alle
scimmie: magari quella persona non consegna in ritardo da mesi, ma ormai lo
dicono tutti. È un circolo vizioso da fermare subito.

\paragraph{Niente escalation inutili.} Lo stesso vale per il feedback: quando
qualcuno va dal relatore dicendo ``ho un problema con questa persona'', lui
chiede: \textbf{``perché ne stai parlando a me e non a lei?''}. ``Perché sei il
mio capo'' non basta. Passare da lui significa dare per scontato che l'altra
persona non abbia interesse a risolvere, mentre magari basta parlarle. È la
strada più facile (``lo dico a Stefano, ci pensa Stefano''), e proprio per
questo bisogna mettere dei paletti: il manager interviene quando serve, ma le
cose che si possono gestire da soli vanno gestite da soli. (Quando invece il
confronto diretto fallisce, l'escalation è legittima: si veda la
sezione~\ref{sec:z1-qa-finale}.)

% ---------------------------------------------------------------------
\section{Traiettorie di carriera e valutazione delle performance}
\label{sec:z1-traiettorie}

\subsection{Quale traiettoria vorreste?}

Il relatore mostra un grafico con quattro curve colorate, contrassegnate
anche da lettere (A, B, C, D) per eventuali persone daltoniche (``a me è
capitato che qualcuno scoprisse di esserlo durante una lezione''). Le curve
rappresentano la performance lavorativa nel tempo. Due domande: \emph{quale
vorreste essere? E quale pensate realisticamente di essere?} Un partecipante
aspira a B e pensa di essere C; quasi tutti scelgono B.

Perché non A, che sale più rapidamente? Un partecipante risponde che una
crescita così repentina può essere più dannosa che utile, come una ``rush'' di
studio molto faticosa che poi non si riesce a sostenere. Il relatore propone
una lettura suggerita da un amico psicologo, che spiega la forma del grafico:
di B, C e D \textbf{vediamo dove finiscono}; di A no. Potrebbe continuare a
salire oltre la scala del grafico, o crollare sotto D. Non saperlo ci spinge a
preferire B. (A chi obietta che anche le altre proseguono oltre il bordo, il
relatore precisa: il bordo destro è ``oggi'', e di A non sappiamo come sia
arrivata fin lì.)

\subsection{La carriera non è una linea retta}

Tutti immaginano la carriera come una linea: si parte da X e, attraverso
stage e posizioni, si arriva a Y. In realtà carriera e performance avranno un
andamento irregolare, con alti e bassi. Come non si accetta che un team possa
temporaneamente \emph{underperformare}, tanto meno lo si accetta da se stessi.
Bisogna accettare che ci saranno momenti bassi.

Per questo il relatore \textbf{detesta le valutazioni annuali delle
performance}: prendono un punto casuale della traiettoria, che se si è
fortunati è il picco più alto e se si è sfortunati il più basso, e lo usano
per valutare un intero anno. Insegna ai manager a fare più punti di controllo
e a guardare altre cose:
\begin{enumerate}
  \item \textbf{quanto tempo passa tra un picco e l'altro}, cioè quanto
        \emph{cooldown} serve a una persona. È il motivo per cui, dopo la
        Games Week, tutti staccavano una o due settimane. Il relatore stesso,
        che si impegna moltissimo e tende al burnout, non potrebbe rifare lo
        stesso workshop il lunedì e il martedì successivi con la stessa
        qualità;
  \item \textbf{la differenza tra la miglior performance e la peggiore}: che
        cosa ha portato a lavorare così bene, e così male? Con chi si
        lavorava, su cosa ci si concentrava, quanto tempo e budget c'erano?
        L'obiettivo è capire come \emph{replicare} le condizioni che fanno
        lavorare bene, invece del solito ``hai lavorato male, è colpa tua''.
        Il relatore ha avuto molti manager per cui quando lavorava male era
        colpa sua e quando lavorava bene ``che bravo'', ma sapeva che raramente
        era l'unico fattore della propria performance. È la fallacia del ``se
        sai fare 100, perché non fai sempre 100?'', come nella scena di
        \emph{Schindler's List} sull'operaio che produce una vite in pochi
        secondi: non sempre un ritmo è sostenibile;
  \item \textbf{il fallimento di oggi confrontato con quello di sei mesi fa}.
        Il relatore non è contento della presentazione di oggi (l'avrebbe
        migliorata per altri due anni), ma rispetto a quelle di cinque anni fa
        vede un grande miglioramento di chiarezza: il suo peggior fallimento di
        oggi equivale al suo primo successo di inizio carriera.
\end{enumerate}
Quello che conta è la \textbf{traiettoria} dall'inizio dell'anno a oggi, non
i singoli punti. E se si riesce a capire insieme quando e perché si ha
successo e quando e perché si fallisce, si può lavorare per replicare solo le
situazioni che fanno lavorare bene.

\paragraph{Lo stesso vale per i giochi.} È il motivo principale per cui i
giochi delle grandi case vanno bene o male. Quando si riescono a replicare le
condizioni che hanno prodotto \emph{Warcraft~II}, esce \emph{Warcraft~III};
quando non ci si riesce, esce ``il nuovo \emph{Dragon Age}'': la gente non è
arrabbiata perché il gioco non è bello, ma perché non è come gli altri
\emph{Dragon Age}; non ne replica il \emph{feeling}, anche se potrebbe
essere un gioco favoloso. L'eccezione positiva è il passaggio da \emph{Dune} a
\emph{Dune~II}: il primo era di fatto un'avventura grafica, il secondo ha
generato il genere degli strategici in tempo reale. Gli sviluppatori avevano
capito che il successo non stava nel tipo di gioco, ma nell'ambientazione che
attirava le persone: un bel gioco in quell'ambientazione sarebbe piaciuto
comunque. Il relatore comprò \emph{Dune~II} aspettandosi un seguito
dell'avventura e, pur sorpreso, lo finì.

\subsection{Lying, hiding, faking}

Quando si è valutati su un singolo punto nel tempo, e da quel punto dipendono
bonus, stipendio e il rischio di licenziamento, si comincia a fare
\textbf{lying, hiding e faking}:
\begin{description}
  \item[Lying] dire che è sempre andato tutto bene, che si finirà sempre in
        orario, che non c'è nessun problema;
  \item[Hiding] nascondere: ``non era un mio task'', ``non c'ero'', ``non ho
        visto niente'';
  \item[Faking] fingere, anche in cose apparentemente innocue come ``sì, sì,
        ho capito tutto'' quando si ha solo paura di dire ``me lo rispieghi?
        Non ho capito niente''.
\end{description}

\subsection{``Non ho capito'' e la conferma delle istruzioni}

Dire ``non ho capito, ripetimelo'' è ancora visto come una figura da sciocco.
Il relatore non si arrabbierà mai se qualcuno gli chiede di rispiegare
un'intera slide: le cose ovvie per lui possono non esserlo per altri. Più si
conosce qualcosa, \textbf{più è difficile spiegarla}, perché si danno per
scontate molte cose. Ne parlava proprio con Diego, che gli diceva che fare
l'insegnante diventa sempre più difficile man mano che si diventa esperti. È
come spiegare \emph{EVE Online}, uno dei giochi più complessi in circolazione:
chi lo conosce bene dà per scontato che ``ovviamente non puoi fare questa
cosa'', mentre per un nuovo giocatore non è affatto ovvio. A scuola chi dice
``non ho capito'' viene visto come lo ``scemo''; nessuno si chiede ``non è che
l'ho spiegato male?''.

Il racconto più ``traumatico'' dei ragazzi che fanno stage nelle grandi
aziende è sempre lo stesso: ``Cinzia, devi fare questa cosa.'' Cinzia la fa.
``Non te l'avevo chiesta così.'' Il problema è che nessuno ha
\textbf{confermato che Cinzia avesse capito}. Il consiglio, valido sempre:
\begin{itemize}
  \item quando si \emph{ricevono} istruzioni, ripeterle: ``ho capito bene?
        Vuoi che faccia questa cosa in questo modo?'';
  \item quando si \emph{danno} istruzioni, chiedere a chi le riceve di
        ripeterle, non perché lo si ritenga incapace, ma per capire che cosa
        ha capito dalla propria spiegazione (``quindi vuoi un quadrato?''
        ``No, ho detto un triangolo\dots'' ``Però mi hai detto una forma con
        quattro angoli.'' ``Hai ragione: voglio un quadrato'').
\end{itemize}
Frasi come ``non te lo ripeto due volte'' o ``stammi dietro'' portano solo a
fare le cose come si pensa che vadano fatte e a sentirsi dire a fine progetto
``non la volevo così'': che si tratti di una feature, di un design o di un
modello 3D. (Un partecipante osserva che lo stesso accade insegnando le lingue:
chi insegna dà per scontato, ad esempio, che si sappia dell'esistenza del
genere neutro in tedesco.)

\paragraph{Come convincere il team a chiedere?} Il relatore mette sempre come
premessa che \textbf{l'obiettivo non è trovare la persona più intelligente
nella stanza, ma fare un buon lavoro}. Che Gianni capisca alla prima e un altro
alla quinta non gli interessa, se alla fine entrambi lavorano bene: ripetere
cinque volte è comunque meno lavoro che far rifare le cose sette volte. È un
investimento a lungo termine, e non dissimula alcun fastidio: davvero non gli
dà fastidio. Ma ripetere le stesse parole non aiuta: ``non ho sentito'' si
risolve ripetendo; ``non ho capito'' si risolve cambiando spiegazione. Meglio
ancora: ``dimmi che cosa hai capito tu, e vediamo cosa non torna''. Se
intravedesse nel team uno sguardo del tipo ``è la quinta volta che lo
chiede'', interverrebbe: spesso chi giudica è proprio chi dice ``ho capito'' e
poi sbaglia. Preferisce di gran lunga chi chiede cinque volte.

\subsection{Il feedback più utile che abbiate ricevuto}

\emph{Qual è il feedback più utile che abbiate mai ricevuto?} Il relatore ne ha
almeno dieci che gli hanno cambiato il modo di lavorare. Le risposte:
\begin{itemize}
  \item un semplice \textbf{``grazie''} per un lavoro fatto, più utile di un
        ``questa cosa fa schifo'';
  \item un feedback sulla propria \textbf{impulsività} nel rispondere: fermarsi
        un attimo e capire di cosa si sta parlando prima di intervenire;
  \item \textbf{essere se stessi}, anche al lavoro, e non scusarsi
        continuamente per ogni cosa.
\end{itemize}
I feedback che toccano di più sono quelli \textbf{veri}. Il ``grazie''
significa ``ho visto quello che hai fatto e l'ho apprezzato''. Quelli negativi
che si ricordano sono quelli che hanno fatto realizzare ``questa cosa che
pensavo di saper fare, non la so fare'', oppure ``la so fare, ma questa volta
non mi è venuta'' (non è che se una cosa una volta non riesce non la si sa più
fare: può capitare di sbagliare una porta). Arrivano sempre da un
\textbf{punto cieco}: qualcosa che non si sapeva fosse percepito dagli altri.
Conta più ciò che le persone ti dicono di ciò che pensi che pensino di te; e
ciò che il relatore pensa che gli altri pensino di lui è sempre terribile,
mentre il loro giudizio si rivela molto più positivo del previsto.

\paragraph{Il denominatore comune siete voi.} Si potrebbe pensare che questi
concetti non siano applicabili da chi arriva in un'azienda senza poterne
cambiare la cultura. Ma c'è sempre un denominatore comune a cui applicarli:
\emph{se stessi}. La cultura aziendale può non spingere a dare feedback, può
vedere male chi chiede aiuto; ma il relatore ha visto che quando lui lo fa, le
persone intorno tendono a farlo, e un gruppo che lo fa lavora meglio. Perfino
i CEO più ostici a un certo punto gli chiedevano: ``perché il tuo team ora
performa meglio degli altri?''. La differenza è il default to trust: una
pratica che col tempo fa operare il team \emph{come un team} e non come un
insieme di individui. Un videogioco è un progetto corale: c'è chi fa giochi da
solo, ma non si fa da soli \emph{Red Dead Redemption~2}; e se il gruppo non è
affiatato e allineato non esce \emph{Red Dead Redemption~2}, esce
\emph{Anthem} (``oggi la mia fissazione è \emph{Anthem}, domani sarà
un'altra'').

% ---------------------------------------------------------------------
\section{Come dare un feedback}
\label{sec:z1-dare}

\subsection{Concentrarsi sull'azione e sui suoi effetti}

Il primo principio: \textbf{il feedback non riguarda mai la persona}. Non si
dice ``tu hai fatto questa cosa che non va bene'', ma ``questa cosa non va
bene: come facciamo a non farla più?''. Che sia un comportamento, un'azione,
una decisione o una strategia, ci si concentra su quello, e si
\textbf{esplicita l'effetto} che ha. Nell'esempio del ritardo l'effetto era
``quindici minuti di attesa moltiplicati per venti persone''. In un progetto
può essere una catena:
\begin{quote}
  se non rispetti le deadline, l'animazione che viene dopo resta ferma; se
  loro sono in ritardo, la \emph{vertical slice} è in ritardo; se la vertical
  slice non è pronta per la Gamescom, non andiamo alla Gamescom e non troviamo
  un publisher; se non troviamo un publisher, chiudiamo.
\end{quote}
Più la conseguenza è grave, più dovrebbe sembrare piccolo il gesto di dire
``datemi una mano'': è il \emph{male minore}. Poi, a cose fatte, ci si dà
feedback: cosa ha funzionato, cosa no, cosa fare di più, di meno, cosa smettere
di fare.

\subsection{Il pericolo degli ambienti solo positivi}

Gli ambienti dove ci si fa sempre grandi complimenti e nessuno parla dei
problemi creano difficoltà vere alle aziende. Il problema non è mai la
mancanza di complimenti, ma un \textbf{eccesso di ottimismo}: si comincia a
pensare che le cose andranno come si spera, perché tutti intorno danno solo
feedback positivi.

A Blizzard i feedback determinavano gli \textbf{aumenti annuali}: ognuno
doveva indicare cinque persone che avrebbero dato feedback su di lui, e
ovviamente tutti indicavano gli amici, o colleghi di progetti andati bene. Un
sistema nato con le migliori intenzioni diventava uno \emph{scambio di voti}
(``se mi metti follow ti metto follow''), senza alcun miglioramento.

Al contrario, \textbf{le conversazioni più brutte sono state le più utili}:
affrontare i licenziamenti a Riot, o la cancellazione di \emph{Project Titan}
a Blizzard, il progetto che doveva essere il nuovo MMO dopo \emph{World of
Warcraft} e da cui poi è nato \emph{Overwatch}\footnote{Nella trascrizione si
legge anche ``doveva essere Warcraft 2'', probabilmente un lapsus o un errore
di trascrizione.}. Andando a vedere perché lo si cancellava, c'erano solo
motivazioni razionali e pragmatiche, di business, chiare a tutti, non un ``Mike
Morhaime ci odia''. Sono anche le conversazioni in cui ha lavorato su se
stesso: è molto \textbf{avverso ai conflitti}, e questo a volte gli ha
procurato conseguenze molto più gravi di quanto sarebbe costato un piccolo
disagio. Dieci anni fa faceva fatica ad accettare questo difetto; ha dovuto
subire le conseguenze di conflitti evitati e poi cresciuti.

\subsection{Che cosa fare}

Le linee guida della slide, che il relatore legge rapidamente perché saranno
disponibili nel materiale:
\begin{description}
  \item[Attenzione al linguaggio del corpo] dare un feedback a braccia
        incrociate o guardando altrove è diverso dal guardare l'altro negli
        occhi.
  \item[Sii specifico] con esempi tangibili.
  \item[Esplicita gli effetti] come visto sopra.
  \item[Discuti il prima possibile] non aspettare un mese per parlare di
        qualcosa successo ieri, altrimenti nessuno ricorderà più bene cosa sia
        accaduto.
  \item[Concentrati sull'azione e sul problema, non sulla persona] il nemico è
        l'azione, non la persona che si ha davanti. Non bisogna far cambiare
        l'altro: bisogna evitare che il problema si ripresenti.
\end{description}

\paragraph{Gli insulti sono rivolti al ruolo.} Un esempio solo apparentemente
scollegato. Il relatore ha lavorato su \emph{League of Legends}, ``non
esattamente noto per avere una community da orsetti del cuore'', in varie
\emph{task force} sulla tossicità. A chi smetteva di giocare veniva inviato un
sondaggio sul perché; il 33\% indicava la tossicità della community. Dopo varie
campagne di sensibilizzazione e cambiamenti si è scesi al 30\%: sembra solo
un 3\% in meno, ma valeva circa \textbf{120 milioni di revenue in più
all'anno}, perché quei giocatori restavano esposti a skin e contenuti e
continuavano a comprare. Anche qui, piccoli guadagni che si sommano.

Lavorando con degli psicologi, aveva notato che gli insulti \textbf{non sono
quasi mai rivolti alla persona, ma al ruolo}: ``il mio jungler fa schifo'',
``il mio DPS non è capace'', ``il nostro portiere è una pippa''. Chiamare
qualcuno per ruolo lo \emph{depersonalizza} e rende più facile insultarlo. Lo
stesso nell'esempio del giro dell'ufficio: non ``non chiedere a Giulio'', ma
``non chiedere a quel grafico, consegna sempre in ritardo''. Concentrarsi
sull'azione, senza attaccare la persona, fa sì che il feedback non sia (e non
venga percepito come) un attacco personale.

\paragraph{Considera il carattere.} Qui tornano utili i test di personalità.
Servono più agli studenti che al relatore (che sa già che usciranno 53 profili
diversi, senza un pattern): bisogna conoscere il proprio profilo e
\textbf{condividerlo con chi lavora con noi}, chiedendo anche agli altri di
fare lo stesso. In ufficio il relatore ha usato una versione più avanzata, a
pagamento, con un report di circa trenta pagine: il \textbf{profilo
DISC}\footnote{Nella trascrizione ``disk''. Si tratta del modello DISC
(Dominance, Influence, Steadiness, Conscientiousness).}. Ha mappato tutta
l'azienda e ha fatto stampare il profilo di ciascuno sulla sua scrivania, così
che tutti sapessero con chi stavano parlando. Capire come ragionano gli altri
(``è introverso, non gliene parlo davanti a tutti'') è una delle cose che gli
sono tornate più utili: \emph{adattare ciò che si vuole dire a chi si ha
davanti}.

\paragraph{Assicurati di non disturbare.} Sotto il profilo di ognuno c'era un
piccolo \textbf{semaforo}: verde, ``parlami pure, interrompimi''; giallo,
``interrompimi solo se è urgente''; rosso, ``mi sto concentrando, non mi
parlare''.

\paragraph{Lascia finire il lavoro.} Durante gli eventi dal vivo il relatore si
era reso conto di quanto fosse grave dare feedback mentre qualcuno sta ancora
lavorando. Se una cosa sta andando male, dirlo in quel momento non aiuta: tra
la pressione di finire e quella del feedback, peggiorano entrambe le cose, e
anche il rapporto. Meglio lasciarla finire, anche male, e parlarne dopo.

\subsection{Che cosa non fare}

\begin{description}
  \item[Non basarti su supposizioni] ``mi hanno detto che ieri in un meeting
        hai fatto\dots'': o c'eri e l'hai visto, o non esiste.
  \item[Non giudicare la persona.]
  \item[Non fornire feedback vago.]
  \item[Non accettare e non fornire feedback di terze parti] Se qualcuno vi
        dicesse ``mi hanno detto che hai fatto un brutto lavoro'', il primo
        pensiero non sarebbe ``cosa ho fatto male?'', ma \emph{``chi te l'ha
        detto? Chi è la vipera che ho in seno? Chi è l'impostore?''}. Siamo
        fatti così. Ogni conversazione che inizia con ``mi hanno detto che'' è
        inutile: l'altro non ascolterà, penserà solo a chi dare la caccia, e
        spesso penserà alla persona sbagliata, creando antipatie tra persone
        che non c'entravano nulla. Quando al relatore dicono ``mi hanno detto
        che lui ha fatto\dots'', risponde ``vai a parlarne con lui''; se gli
        dicono ``mi hanno detto che tu hai fatto\dots'', risponde ``mandami la
        persona che te l'ha detto, altrimenti non ne parliamo''. Non per
        cattiveria: se chiedesse ``che cosa ho detto di preciso a Marco?'' e la
        risposta fosse ``non lo so, non c'ero'', di cosa si starebbe parlando?
        \textbf{Parlate con la fonte}, altrimenti fate come se quel feedback non
        esistesse.
  \item[Non dare consigli non richiesti] come visto, non portare la soluzione:
        la soluzione deve venire da chi riceve il feedback.
  \item[``Non mi piace'' non è un feedback] ``Penso che questo design non
        funzioni perché, visto che il nostro obiettivo è\dots'' ha argomenti;
        ``non mi piace il combat system'' è un'opinione. Anche ``io l'avrei
        fatto diversamente'' non è un feedback.
  \item[Accettare non vuol dire fare come proponi] deve essere chiaro a chi dà
        il feedback: l'aspettativa è essere ascoltati e presi in
        considerazione, ma se la responsabilità è dell'altro, deciderà
        l'altro.
  \item[Non offenderti] se il feedback non viene seguito. Restarci male
        significa pensare di saperne più dell'altro, cioè \emph{venire meno al
        default to trust}. ``Decidi tu'' significa ``mi fido di te, perché sei
        lo specialista, e mi fido che mi ascolterai e farai la cosa migliore
        per tutti'', che potrebbe anche non essere ciò che ho detto io. Se
        voglio che l'altro faccia sempre come dico io, mi sto fidando solo di
        me stesso.
\end{description}

\subsection{Il primo feedback e il debriefing}

La cosa più difficile è dare il \textbf{primo} feedback: è un muscolo da
allenare. Chi non l'ha mai fatto (soprattutto con questa struttura) dovrebbe
provarci, anche su minuzie (``abbiamo concordato di usare WhatsApp nel gruppo:
non potremmo usare Telegram?''), anche solo per vedere come va. Altre
indicazioni:
\begin{itemize}
  \item feedback \textbf{solo uno a uno}, mai davanti ad altri, soprattutto se
        negativo;
  \item \textbf{prepararsi} prima;
  \item fare il \textbf{debriefing del feedback}: dopo, fermarsi a pensare a
        com'è andato. Il relatore, quando media un feedback tra due persone,
        si sente spesso rispondere ``benissimo'' da una e ``malissimo''
        dall'altra: la percezione è molto diversa. Non si fa problemi a tornare
        dalla stessa persona: ``ieri ti ho dato un feedback: com'è andata? Sei
        soddisfatto? Come poteva andare meglio?''. Con le persone con cui si ha
        meno a che fare, questo aiuta molto la volta successiva.
\end{itemize}

% ---------------------------------------------------------------------
\section{Come ricevere un feedback}
\label{sec:z1-ricevere}

Per il relatore è la parte più difficile. Si riassume in quattro passi:
\begin{enumerate}
  \item \textbf{Ascolta}: si parte sempre ascoltando.
  \item \textbf{Chiedi dettagli}, se serve, ma per capire meglio, non per
        difenderti.
  \item \textbf{Accetta, a prescindere}: non si sta decidendo chi ha ragione;
        si sta ascoltando un'opinione. Non c'è nulla da cui difendersi.
  \item \textbf{Agisci}: fatti un piano. Può essere banale (``mi hanno detto
        che interrompo sempre: cerco di non interrompere''), oppure ``devo
        migliorare questa cosa: come?''. Pensa anche a \emph{chi} potrebbe
        aiutarti e a \emph{che supporto} ti serve: il piano non deve essere
        per forza ``devo risolvere tutto da solo''.
\end{enumerate}
Una domanda che molti sottovalutano: \emph{``sono junior in un'azienda: sarò
mica il primo ad avere questo problema?''}. Probabilmente qualcun altro l'ha
già avuto: basta chiedere ``tu come hai fatto?''.

\subsection{Il feedback positivo}

Tutto questo vale per i feedback costruttivi negativi, a cui si pensa di
default; ma esistono anche i \textbf{feedback positivi}, e vanno usati. Quando
si riceve un ``grazie'' o un ``bravo'', la domanda è: \emph{come faccio a
ricreare le stesse condizioni?}

Il relatore ha lavorato anche in aziende dove il feedback era solo negativo.
Il rischio? Si immagini l'impiegato ideale: vuole sempre migliorare, è sempre
pronto a dire ``magari sbaglio io''. Se nessuno gli dice cosa sta facendo
\emph{bene}, spinto dal desiderio di migliorare rischia di cambiare cose che
erano già buone, e di farle peggio. Il rinforzo positivo \textbf{non è solo una
gentilezza}: è un atto pratico, che evita che le persone peggiorino ciò che
sanno già fare bene.

\subsection{Il feedback accidentale}

Esiste infine il \textbf{feedback accidentale}, molto frequente negli uffici:
\begin{itemize}
  \item il ``chiedi a lei, non a quel grafico'' del giro dell'ufficio è un
        feedback negativo su una persona, dato accidentalmente, né richiesto
        né voluto;
  \item a Riot li chiamavano \textbf{drive-by}, come nelle scene dei film in
        cui passa un'auto, si abbassa il finestrino, si spara e si riparte. In
        molte aziende con una ``cultura del feedback'' in cui tutti devono
        dire la loro, si riceve una valanga di commenti del tipo ``comunque il
        sito a me non piace, ciao''. Non si sa perché, non si sa cosa fare di
        diverso, non si sa che rischi ci siano: \emph{a che cosa è servita
        questa conversazione?} A niente.
\end{itemize}

% ---------------------------------------------------------------------
\section{Quando, quanto e con chi}
\label{sec:z1-quando}

\subsection{Frequentemente e recentemente}

\emph{Quando avete ricevuto il vostro ultimo feedback?} Per gli studenti,
qualche giorno fa (nel corso) o qualche mese fa (al lavoro, dopo una
consegna). Come aveva osservato Silvia, il feedback andrebbe dato
spesso.

Il relatore usa l'esempio di una nave da crociera (``la Costa Concordia'') che
deve andare da A a B. La rotta ideale (azzurra) è quasi impossibile: nessun
progetto va esattamente come previsto. C'è la rotta ``panoramica'' (verde), che
gira intorno all'isola d'Elba; ci sono poi rotte (arancione e rosa) che a un
certo punto si allontanano molto e poi correggono bruscamente: qualcuno si è
accorto che si stava sbagliando, ``ma non è quello che avevo chiesto'',
``non è dove dovevamo andare''. Sono i famosi \textbf{sanity check} o
\emph{checkpoint}.

L'ideale è avere \textbf{tanti piccoli check per fare tante micro-correzioni}:
correggere la rotta di un paio di gradi è molto più rapido che farlo dopo
essersi allontanati. Bastano trenta secondi di ``ehi, vado bene?''. Spesso non
lo si fa per paura di sembrare insicuri o poco autonomi, ma per il relatore è
l'opposto: quando qualcuno gli chiede ``vado bene così?'' e lui risponde
``sì'', si alleggerisce anche del proprio carico mentale. Un capo che assegna
un compito e lo vede solo alla fine vive sempre col dubbio (``avranno capito?
Lo staranno facendo giusto? Arriverà in tempo?''); con un check di dieci
secondi al giorno (``tutto ok, Stefano'') è tranquillo, e la persona va avanti
sapendo che quanto fatto finora va bene. Non serve fare ogni giorno il giro del
dipartimento: anche ogni uno o due giorni basta a evitare la rotta verde.

Per il feedback è lo stesso: \textbf{va dato subito}. Dire oggi ``ieri hai
fatto questa cosa'' è molto diverso da dire ``negli ultimi sei mesi hai fatto
questa cosa''. Il relatore si è sentito dire più volte, alla valutazione di
dicembre, ``da maggio hai fatto questa cosa che non andava''. La sua risposta:
``e aspetti dicembre per dirmelo? Mi hai lasciato sbagliare per sei mesi?''. È
una mancanza di rispetto, e la conversazione non può più essere costruttiva.
Per questo non ama le valutazioni annuali: i manager crescono abituandosi a
prendere appunti delle cose negative per tirarle fuori alla valutazione
(``a marzo hai fatto questa cosa'' ``sì, però poi ho smesso'' ``sì, però a
marzo l'hai fatta'').

In sintesi:
\begin{description}
  \item[Frequentemente] riguarda l'intervallo tra un check e l'altro.
  \item[Recentemente] significa non lasciar passare un anno prima di dare il
        feedback, altrimenti non serve più a niente.
\end{description}

\subsection{Quanto dura un feedback}

Un feedback dura \textbf{da trenta secondi a quindici minuti}; raramente di
più. Chi non ne ha mai dato uno dovrebbe sforzarsi di darne almeno uno, magari
positivo per cominciare. Ogni feedback dato o ricevuto è un
\textbf{deposito nel proprio ``fondo fiduciario''} (\emph{trust fund}), che
costruisce la reputazione futura. Il relatore si è costruito la reputazione di
chi dà feedback il più possibile oggettivi e analitici: giusti o sbagliati che
siano, costringono a pensare. Non impone mai il suo punto di vista (``si fa
come dico io''), ma chiede ``tu come faresti? Come posso aiutarti?''. Il
feedback, per lui, serve a \emph{far crescere qualcun altro}.

\subsection{I manager peggiori: ``a che numero sto pensando?''}

I peggiori manager che il relatore abbia avuto erano quelli che volevano che si
facesse come volevano loro, anche quando chiedevano ``tu come la faresti?''.
In realtà la domanda era ``\emph{a che numero sto pensando?}'': se dicevi 5,
``no, dovremmo fare 9''; se dicevi 9, ``no, 5''. Lo stress costante di dover
indovinare cosa vuole il capo nasce dalla mancanza di fiducia (``se mi chiedi
come fare, fammelo fare'').

Lo stesso gli succede con i CEO che lo chiamano per ``sistemare la cultura
aziendale'' e poi dicono ``non sono sicuro, io farei diversamente''. La sua
risposta: ``lo so che faresti diversamente, per questo mi hai chiamato''. Se si
chiama un esperto, bisogna lasciargli fare il suo lavoro: non si va dal
meccanico a cambiare le gomme per poi dirgli ``non dovresti girare un po' di
più quel bullone?''. Chiedere un parere a un esperto significa dargli anche il
potere di scegliere. Il problema è il \emph{micromanagement}, o in generale il
non accettare punti di vista diversi dal proprio (magari perché si è il
fondatore e si pensa di essere l'unico a saper fare le cose).

\subsection{Dare feedback a un superiore egocentrico}

\paragraph{Come si dà un feedback negativo a un superiore con un grande ego?}
Il trucco è \textbf{non puntare sull'ego, ma sull'azione e sui suoi effetti}.
Non ``tu fai così e io mi sento così'': anzi, si alimenta il suo ego
rendendolo il \emph{main character} dell'avventura, l'eroe che risolve il
problema (``tu, grande manager, aiutaci a risolvere questo problema che
affligge noi che facciamo l'operatività''). Non è l'ideale (``io cambierei
manager''), ma pochissimi possono scegliersi i colleghi: il relatore stesso ha
quasi sempre ricevuto team preassemblati (``ecco le chiavi, questa è la tua
macchina''). Bisogna giocare con le carte che si hanno in mano.

\paragraph{Un esempio pratico.} Con un manager che fa micromanagement e mette
sempre ``il becco'' nel proprio lavoro, si può dire: ``se ti fidi di me per
questo lavoro, con le tue capacità potresti dedicarti a questa cosa più
strategica e di lungo termine: pensa quanto sarebbe più utile per l'azienda''.
In sostanza: ``guarda cosa potresti raggiungere se mi lasciassi fare il mio
lavoro''. Non come un ``tu fai il tuo e io faccio il mio'', ma come un
``ti sollevo da questa parte del carico, che so che ritieni importante: lascia
che ci pensi io, così puoi dedicarti ad altro'', che magari gli dà anche più
visibilità.

Si gioca su ciò che quella persona pensa di saper fare, e che probabilmente sa
anche fare. Anche Michael Scott di \emph{The Office}, per quanto preso in giro,
è raccontato come un venditore bravissimo, che ha chiuso i contratti più
importanti per l'azienda, e rispetto ai manager che gli succedono è forse
``uno dei meno peggiori''. Tutti, anche i manager peggiori, hanno punti di
forza: \textbf{giocate su quelli}.

% ---------------------------------------------------------------------
\section{Il feedback come regalo}
\label{sec:z1-regalo}

Per chiudere, il relatore risponde alla domanda d'apertura, che aveva lasciato
apposta in sospeso: che cos'è il feedback per lui?

\begin{quote}
  \emph{Il feedback è un regalo che mi viene fatto da qualcuno che so che vuole
  il mio interesse.}
\end{quote}

Ha costruito un team di persone che si fidano di lui e che non hanno motivo di
sabotarlo o di ferirlo: per questo interpreta anche le cose peggiori con il
default to trust. È un regalo, e cerca di prenderlo in considerazione il più
possibile.

Ci è arrivato partendo da se stesso. Oltre a essere avverso ai conflitti, si
definisce ``avverso alla comunicazione umana'': paradossalmente, una persona
introversa e timida finita a fare comunicazione, marketing e ora HR (``succede
spesso''). Prima di ordinare una pizza al telefono ha dieci minuti di
preparazione mentale, e se l'ordine delle domande cambia va in confusione. Per
lui un feedback di dieci minuti sono spesso \textbf{tre ore di preparazione}:
costruisce un \textbf{albero dei dialoghi}, come in un'avventura grafica
personale, con tutte le opzioni. Poi spesso finisce con ``hai fatto questa
cosa, andrebbe migliorata'' ``hai ragione, grazie'', e basta.

Così ha pensato: chissà quanti si fanno gli stessi problemi con lui, e passano
ore a prepararsi per dirgli ``questa cosa l'hai fatta male'', sapendo che
riceveranno comunque quell'occhiata fugace del ``come ti permetti''. Eppure
vengono a dirglielo. Non può far altro che rispondere: ``grazie, hai ragione,
cercherò di fare meglio''. È così che ha metabolizzato il ricevere feedback:
capendo che non c'è nessuna cattiva intenzione e che, se non proprio felice
(``ci ho messo una pietra sopra''), può almeno essere \emph{grato} che ci sia
ancora qualcuno che ha così a cuore il suo interesse da dargli feedback.

\paragraph{Il silenzio fa più paura.} Non bisogna spaventarsi dei feedback
negativi che emergono dai sondaggi aziendali: bisogna spaventarsi del
\textbf{silenzio}, quando nessuno dice più nulla. Significa che nessuno si
fida più che si sappiano risolvere i problemi.

\paragraph{La sindrome dell'impostore.} Il suo motivatore personale, all'inizio,
era soprattutto ``non voglio fallire'': per questo cercava di migliorare, e
migliorare è diventato accettare feedback, perché più feedback riceve, più
migliora. Il feedback è un \textbf{acceleratore}, e lui cerca di incentivare
gli altri a darglielo. Resta però sempre la vocina del ``e se si accorgono che
non sei capace?''. Con la \emph{impostor syndrome} convive da anni e
continuerà a conviverci; se ce l'avete, conviveteci anche voi. Avrete sempre
la sensazione che ``se ne accorgeranno'': forse sì, ma sapete fare anche
altre cose. Non concentratevi solo su ciò che non sapete fare. Gli errori
capitano: la perfezione non è di questo mondo.

\paragraph{Esercitarsi in accademia.} Per fare pratica con il feedback,
l'ambiente migliore è quello dei colleghi di corso: persone nella stessa
situazione, con esperienza simile e lo stesso obiettivo, in un contesto in cui
non si può essere licenziati e l'unico effetto a lungo termine è ``ho parlato
con questa persona, ora le cose vanno meglio''. E bisogna avere il coraggio di
parlare anche con i docenti (``dire a Diego: secondo me questa lezione potevi
farla in questo modo''; ``so che è una persona che lo accetta''), senza farsi
bloccare dalla gerarchia. Nell'azienda giusta (come quella in cui lavora ora) il
relatore va sempre dal CEO a dare feedback, sapendo che non verrà licenziato:
il CEO gli ha detto ``ti ho preso per darmi feedback''.

La frase di chiusura:

\begin{quote}
  \emph{Per superare le vostre difficoltà, tutto quello che vi viene chiesto è
  il coraggio di chiedere aiuto.}
\end{quote}

Chiedere aiuto non è una debolezza: è un segno di consapevolezza di sé
(\emph{self-awareness}) e di forza personale. Se tutti lo facessero, i
progetti andrebbero molto meglio; spesso invece ci si fa frenare dal ``non
voglio fare brutta figura'', ``non voglio passare per quello scemo'', e sono
proprio queste paure a rallentarci e a impedirci di raggiungere gli obiettivi.

% ---------------------------------------------------------------------
\section{Domande e risposte finali}
\label{sec:z1-qa-finale}

L'ultima mezz'ora è dedicata a domande libere.

\paragraph{E se una persona peggiora per uno o due anni di fila?} Se nessuno
gliel'ha mai detto, c'è un grosso problema, non suo ma dell'azienda. Se invece
c'è consapevolezza, il percorso è: ``cosa vuoi fare per risolvere?''; ``ho
provato questo, questo e questo''; ``visto che non ha funzionato, vuoi un mio
consiglio? Proviamo con quello?''; ``che supporto ti serve?''. Se, con tutto il
supporto, il tempo e le soluzioni concordate per mesi, il miglioramento non
arriva, per il relatore la soluzione è: \textbf{``dimmi un prezzo per andare
via''}. Preferisce pagare una buonuscita (\emph{severance package}) che
avviare la strada dei richiami formali sulla performance fino al
licenziamento per giusta causa. A un certo punto andare avanti non è un favore
per nessuno.

\paragraph{È vero che nei videogiochi si viene sfruttati? Meglio un piccolo
studio italiano?} Per il relatore si è quasi più a rischio a Blizzard che in un
piccolo studio di Roma o Torino. Nel piccolo studio si sa dove ci si va a
cacciare: non ci sarà il team building a Santo Domingo, ma ci si concentra su
un gioco alla volta e, se va bene, si fa il successivo. A Blizzard si può
lavorare su \emph{Overwatch}, che vende milioni, e perdere comunque il posto,
in modo totalmente scollegato dal proprio lavoro. Gli stipendi più alti
all'estero compensano anche questo rischio, che spesso non dipende da come va
il gioco (anzi, ``il gioco non ha fatto come volevamo'' è spesso una scusa per
licenziamenti che non avevano motivo).

\paragraph{Cambiare settore a trent'anni, da web developer full remote?} Lo
stipendio indicato dallo studente non è fuori dalla portata di vari studi
italiani, anzi è superabile. Non va esclusa la possibilità di lavorare da
remoto per aziende piccole ma estere. Alcuni amici del relatore, pagati con
stipendi italiani, si sono trasferiti a Bali o in paesi con un costo della vita
più basso; il collega di \emph{PUBG} percepisce uno stipendio olandese vivendo
a Milano, che per quanto cara non costa come i Paesi Bassi.

\paragraph{Come si diventa manager?} Dipende dall'azienda: alcune promuovono i
migliori (processo tendenzialmente sbagliato), altre propongono un percorso di
\emph{people management}. Diventare manager significa \textbf{abbandonare ciò
che si sa fare per imparare a fare altro}. Il relatore ha dovuto accettare che
non sarebbe più stato ``il miglior community manager italiano del mondo''
(``eravamo in pochi, quindi era facile''): dedicandosi alle persone non poteva
più fare il suo lavoro bene come prima. È stata una negoziazione con se
stesso: fare il manager significa non essere più l'esperto di qualcosa, ma
saper gestire e far crescere persone. Significa anche \textbf{cedere il
controllo della propria carriera}: se il team decidesse di non fare più nulla,
il manager ne pagherebbe le conseguenze. Il suo lavoro diventa scegliere le
persone giuste e farle lavorare bene. Lui lo esplicita sempre: sa di affidare
al team parte della propria carriera, e il team sa che lui combatterà per
loro, anche prendendosi le colpe.

\paragraph{L'inglese.} \textbf{Va assolutamente imparato}, anche lavorando in
Italia. Non saperlo preclude circa il 95\% del mercato (tutto quello fuori
dall'Italia) e soprattutto il \emph{networking} con persone che potrebbero dare
occasioni future: l'inglese è spesso l'unica lingua in comune con chi lavora
all'estero ed è oggi la lingua universale dei videogiochi.
\begin{itemize}
  \item \emph{Che livello serve?} Non importa B2 o C2: bisogna saper sostenere
        una conversazione, guardare almeno Netflix con sottotitoli in inglese,
        comunicare senza dodici pause e senza ``come si dice\dots''. Il
        certificato non verrà mai chiesto: verrà chiesto se si sa parlare
        inglese. Il relatore ha sentito persone con un C2 che non sapevano
        parlarlo.
  \item \emph{Come allenarlo senza occasioni di conversazione?} Netflix, Prime
        Video, Now TV, Udemy, Coursera, YouTube, Duolingo, siti di scambio
        linguistico con madrelingua, Discord: non sono mai esistiti così tanti
        mezzi gratuiti. Il relatore ha imparato l'inglese giocando (con una
        pronuncia sbagliata, perché leggeva senza sentire, su cui ha poi
        dovuto lavorare); guardare un film in inglese con i sottotitoli in
        italiano insegna anche la traduzione.
  \item Oltre all'ascolto passivo è fondamentale \textbf{parlare}: in una
        flotta di \emph{EVE Online} con 150 persone quasi tutte madrelingua
        inglese si è esposti a decine di accenti e modi di dire (mille modi di
        dire ``come stai'' oltre a \emph{how are you}); lavorando all'estero si
        impara a distinguere gli accenti (Manchester, Londra) e ad apprezzarli
        anche nei media.
\end{itemize}

\paragraph{Quanto deve sapere un junior?} Dipende da quanto è tecnico il
ruolo. Un game designer junior che non sa qualcosa può essere affiancato da un
senior; un programmatore che non sa programmare, in quel momento, non è utile
all'azienda.

\paragraph{Annunci ``junior con cinque anni di esperienza''.} ``Cercasi
ventenne con trent'anni di esperienza.'' Non è l'annuncio per un junior: è un
errore dell'azienda, che in realtà cerca un \emph{mid}. Riceverà profili
junior, si stupirà che non abbiano cinque anni di esperienza, rifarà l'annuncio
e avrà perso sei mesi. Il consiglio: tenere d'occhio quella pagina, perché
dopo qualche mese l'annuncio verrà corretto.

\paragraph{Una struttura per il post-mortem.} A chi chiede un formato di base,
il relatore anticipa quanto vedrà nella seconda giornata a proposito di
brainstorming: una lavagna \textbf{divisa in quattro aree}, su cui ognuno
attacca tutti i post-it che vuole:
\begin{center}
  \begin{tabular}{@{}p{0.44\textwidth}p{0.44\textwidth}@{}}
    \toprule
    \textbf{Smettere di fare} & \textbf{Iniziare a fare} \\
    ciò che non ha funzionato & ciò che non abbiamo fatto e andava fatto \\
    \midrule
    \textbf{Fare di meno} & \textbf{Fare di più} \\
    ciò che ha fatto perdere tempo rispetto all'impatto ottenuto &
    ciò che ha funzionato ma, fatto con troppa leggerezza, non ha dato i
    frutti sperati \\
    \bottomrule
  \end{tabular}
\end{center}
Con i post-it ognuno ha voce e tutto risulta catalogato; per il report basta
una foto della lavagna (o uno screenshot, se la lavagna è digitale, come Miro).

\paragraph{Come sbloccare la paralisi del feedback verso il management?} Uno
studente ha visto manager dire ``non abbiate paura, la mia porta è sempre
aperta'' senza che cambiasse nulla. C'è una dissonanza tra ciò che dicono e
ciò che fanno: è ``la mamma che dice vieni qua che non ti faccio niente''.
Qualcosa deve aver fatto pensare alle persone che non è sicuro. Prima va
sistemato quello; poi si può provare con i \textbf{feedback anonimi}.
L'obiettivo del relatore resta il feedback aperto e firmato, ma a Stormind ha
iniziato con feedback anonimi verso il management per verificare se
l'anonimato fosse un fattore. Non lo era (evidentemente non c'erano problemi a
parlare con il CEO), e la volta successiva intende abbandonarli. Ma è un buon
\emph{check}: se senza anonimato non arriva nulla e con l'anonimato arrivano
cento feedback, il problema è chiaro.

\paragraph{Quando è giusto rivolgersi a un superiore?} La strada giusta resta
parlare direttamente con l'interessato. Ma se si prova e non si trova una
soluzione, la cosa più ovvia è andare dal proprio manager insieme: ``ci
aiuti?''. Quando il relatore parlava di mediare i dialoghi, era proprio questo:
fare da moderatore, aiutare a sbloccare i punti fermi o, se non si riesce a
scegliere tra A e B, scegliere lui. L'escalation, quando necessaria, va fatta.

\paragraph{Il tempo libero veniva rispettato?} No, e se l'è dovuto imporre
lui. Il lavoro gli piaceva così tanto che continuare a casa non gli sembrava
lavorare. Gli si è ritorto contro con il lancio di \emph{Diablo~III}, fatto in
orario americano (le due o l'una di notte in Europa), con l'ufficio chiuso
perché a nessuno era stato chiesto di restare. Il lancio va malissimo, i
giocatori protestano, e lui, sveglio, risponde sui forum fino alle tre, quattro,
cinque del mattino: ``non posso risolvere il problema, ma resto qui ad
aggiornarvi con quello che riesco a sapere dal team di sviluppo''. Il giorno
dopo trova sulla scrivania un \textbf{richiamo formale dell'HR}: aver lavorato
fuori orario senza autorizzazione, condotta non professionale. Per lui è stato
un interruttore: ``ora me ne vado''. Due settimane dopo faceva colloqui con
Riot. Ancora oggi è contrario a quel richiamo, ma gli ha fatto capire come
ragiona un'azienda: ``hai un contratto 9--18, non mi interessa se l'hai fatto
con le migliori intenzioni''. Eppure per anni ha goduto di un'ottima
reputazione in quelle community come ``quello che è rimasto con noi fino alle
cinque del mattino'', e anche l'azienda ne ha beneficiato.

Ci sono poi aziende in cui, soprattutto per i developer, viene richiesto il
famigerato \textbf{crunch}: lavoro straordinario continuativo, senza pause.
Dipende molto dall'azienda; sta venendo (``grazie al cielo'') deprecato, perché
molte aziende capiscono che non è sano ed è visto male anche dal pubblico, ma
esiste ancora.

\paragraph{E se il tuo capo demotiva qualcuno del tuo team?} Qui ``diventa
personale'': il team del relatore si fida di lui perché sa che in questi casi
``diventa una bestia''. ``Te la puoi prendere con me quanto vuoi, ma il mio
team non lo tocchi.'' Lo gestisce direttamente con l'altro manager, per quanto
brutta possa diventare la conversazione. Lo sforzo più grande è \textbf{non
andarne poi a parlare con il team} (``ragazzi, ho fatto questa cosa''), perché
minerebbe la fiducia verso l'alto: non vuole un rapporto ottimo fino a lui e
pessimo oltre. L'obiettivo è che l'altro manager ``si faccia perdonare''
cambiando il proprio comportamento, e che il team percepisca da solo il
cambiamento.

\paragraph{Chiusura della giornata.} Il relatore chiede un feedback collettivo
in chat (``se domani ci dai \emph{more of the same} siamo contenti, oppure
ditemi cosa limare''), si scusa per i ripetuti crash della connessione e
ringrazia per la proattività nelle domande, che gli servono anche come
\emph{check}: se nessuno chiede nulla su un argomento, significa che può
andare avanti. Racconta di aver passato una settimana a decidere che cosa
mettere nelle slide e di averle poi preparate in poco tempo (rifacendone
alcune fino alle tre di notte). Ricorda infine di usare il test di
personalità per sé e per il proprio gruppo, facendolo fare a chi non l'ha
ancora fatto: tornerà molto utile per cooperare meglio.

% ---------------------------------------------------------------------
\section{Anticipazioni sulla seconda giornata}
\label{sec:z1-domani}

La seconda giornata del workshop sarà più pratica: ``oggi vi ho spiegato
perché, domani vi insegno come lo fanno queste aziende''. Gli argomenti
previsti:
\begin{description}
  \item[Time management] come organizzarsi una giornata professionale in
        azienda: come strutturare una \emph{task list}, cosa fare e cosa non
        fare, come e cosa prioritizzare.
  \item[Brainstorming] i gruppi dovranno fare moltissimi meeting (come facciamo
        il gioco, come lo chiamiamo, che scopo gli diamo). Se non si sanno fare
        bene, i brainstorming sono \textbf{dannosi ed eterni}: un brainstorming
        dovrebbe durare \textbf{al massimo due ore}. Verranno mostrati metodi
        che accorciano la durata, massimizzano l'effetto e, soprattutto,
        \textbf{fanno partecipare tutti}: non c'è niente di peggio di un
        brainstorming in cui parlano in due per il 90\% del tempo. Il classico
        ``ragazzi, non esistono idee sbagliate, sparate'' è il modo peggiore di
        farlo; non è colpa di nessuno, ma si può fare diversamente.
  \item[Guidare da dentro] come essere leader nel proprio team anche senza
        esserne il leader: sviluppare responsabilità personale e verso gli
        altri. Non essere il manager di qualcuno non significa non esserne
        responsabili: la persona accanto è responsabilità di tutti, perché si
        è un team e non un insieme di individui. ``Se qualcuno cade, tutti lo
        facciamo rialzare: non c'è l'addetto al rialzamento.'' È la versione
        \emph{advanced} della giornata: come comportarsi da manager senza
        averne il titolo.
\end{description}
Torneranno inoltre il \emph{service level agreement} e il tema della
condivisione dei successi e della protezione del team dai fallimenti.

% ---------------------------------------------------------------------
\section*{Punti chiave del capitolo}
\addcontentsline{toc}{section}{Punti chiave del capitolo}

\begin{itemize}
  \item ``La capacità di un gruppo di persone di fare cose straordinarie
        dipende dal modo in cui queste persone operano insieme.'' Il 90\% dei
        problemi dei team sono \textbf{problemi di comunicazione}.
  \item Le aziende si liberano sempre più presto dello \textbf{useless
        genius}: le competenze si trasmettono, il carattere no. Il
        \textbf{fit culturale} va valutato anche dal candidato.
  \item Non esiste un unico modo di dare feedback: dipende dal contesto e
        soprattutto da \textbf{chi lo riceve}. Il feedback è una somma di
        piccoli miglioramenti periodici; tenersi le cose nel cassetto non
        serve mai.
  \item I team attraversano \textbf{forming, storming, norming, performing e
        adjourning}. Lo storming non si salta; si passa la maggior parte del
        tempo in forming; la bravura sta nell'accelerare i passaggi e nel
        fare bene il \textbf{post-mortem}, anche dei successi.
  \item Gli obiettivi di un progetto vanno fissati \textbf{all'inizio}: solo
        così si può valutare se è andato bene e scegliere tra A e B.
  \item Il circolo virtuoso è \textbf{feedback $\rightarrow$ miglioramento
        $\rightarrow$ fiducia}. Accettare un feedback non significa fare come
        dice l'altro: decide chi ha la responsabilità, ma deve argomentare.
  \item ``Aiutami a poterti aiutare'': meglio dire in anticipo ``potrei non
        farcela'' che ``non ce l'ho fatta''. Chiedere aiuto non è una
        debolezza.
  \item Quando si dà feedback \textbf{non si porta la soluzione}: la si
        costruisce insieme a chi lo riceve.
  \item La cultura è forte quando si lavora \textbf{per gli altri}; il lavoro
        è valutato come opera d'insieme (anche con bonus di team).
  \item Come in \emph{Among Us}, i team sprecano energie a cercare il
        traditore invece di fare le task. \textbf{Default to trust}: chi ha
        sbagliato ha agito con le informazioni che aveva. Cercare
        \textbf{la causa, non il colpevole}: ``che cosa ti ha fatto pensare
        che fosse la scelta giusta?''.
  \item Carriere e performance non sono lineari: vanno valutate le
        \textbf{traiettorie}, non un punto dell'anno. Le valutazioni puntuali
        generano \emph{lying, hiding, faking}.
  \item Confermare sempre le istruzioni facendosele ripetere; ``non ho
        capito'' non è una colpa.
  \item Dare feedback: sull'\textbf{azione} e sui suoi \textbf{effetti}, specifico,
        presto, uno a uno, adattato al carattere, senza terze parti, senza
        ``non mi piace'', senza offendersi se non viene seguito; poi fare il
        debriefing.
  \item Ricevere feedback: \textbf{ascolta, chiedi dettagli, accetta, agisci}.
        Il feedback positivo serve a non peggiorare ciò che si fa bene; il
        feedback accidentale e il \emph{drive-by} non servono a nulla.
  \item Feedback \textbf{frequenti e recenti}: tanti piccoli \emph{sanity
        check} costano meno di una grande correzione di rotta.
  \item Il feedback è \textbf{un regalo} di chi ha a cuore il nostro
        interesse; fa paura il silenzio, non le critiche.
\end{itemize}

%    % =====================================================================
%  Lezione zero 2 — Gestire il tempo, generare idee, guidare da dentro
%  Fonte: trascrizione della seconda "lezione zero" (GD0_2), seconda e
%  ultima giornata del workshop di due giorni tenuto prima dell'inizio del
%  corso di Game Design, dedicato alla figura professionale e al lavoro in
%  team. Prosegue direttamente il capitolo 0.1 (cultura del feedback).
%  Pacchetti richiesti nel preambolo del libro:
%    \usepackage[italian]{babel}, \usepackage[T1]{fontenc},
%    \usepackage[utf8]{inputenc} (solo pdfLaTeX), \usepackage{booktabs}
%
%  Collocazione nel libro: parte "Lezioni zero", subito dopo
%  GD0_1_cultura_del_feedback, con la stessa numerazione propria (0.2).
%  Nel file principale del libro, ad esempio:
%
%    \part{Lezioni zero: il game designer nell'industria}
%    \renewcommand{\thechapter}{0.\arabic{chapter}}
%    \setcounter{chapter}{0}
%    \input{Seminari/GD0_1_cultura_del_feedback}
%    \input{Seminari/GD0_2_tempo_idee_e_leadership}
%    \part{Il corso di Game Design}
%    \renewcommand{\thechapter}{\arabic{chapter}}
%    \setcounter{chapter}{0}
%    \input{GD1/GD1_introduzione_al_game_design}
%    % ...
%
%  (Con hyperref aggiungere anche \renewcommand{\theHchapter}{z.\arabic{chapter}}
%  nella parte zero, per evitare ancore duplicate.)
%  Tutte le etichette di questo capitolo hanno il prefisso "z2:"; i
%  rimandi al capitolo precedente usano le etichette "z1:" già definite in
%  GD0_1_cultura_del_feedback.tex.
% =====================================================================

\chapter[Tempo, idee e leadership]{Gestire il tempo, generare idee, guidare da dentro}
\label{cap:z2-tempo-idee-leadership}

Questo capitolo raccoglie la \textbf{seconda giornata} del workshop di due
giorni con cui si aprono le lezioni zero, tenuto dallo stesso relatore del
capitolo~\ref{cap:z1-feedback}, Stefano Cozzi\footnote{Per il cognome e per
il nome dello studio (Stormind Games) valgono le avvertenze della nota
all'inizio del capitolo~\ref{cap:z1-feedback}. In questa giornata la
trascrizione riporta soltanto il nome di battesimo, ``Stefano'', con cui il
relatore si cita più volte nei suoi esempi.}. Se la prima giornata spiegava
\emph{perché} le aziende tengono tanto alla cultura del feedback, la seconda
è dichiaratamente più pratica: ``ieri vi ho spiegato perché, oggi vi insegno
come lo fanno queste aziende''. I concetti della prima giornata (il
\emph{default to trust}, il ``potrei non farcela'' detto in anticipo, la
condivisione dei successi, il feedback frequente) ritornano di continuo,
applicati alla gestione del tempo, alle riunioni creative e alla leadership.

La giornata si è svolta di domenica, a chiusura di un fine settimana di
workshop\footnote{Il relatore dice ``oggi chiaramente è domenica'' e saluta
la classe con ``buon resto del weekend''; la prima giornata si era quindi
svolta di sabato. Più avanti parla di ``dodici ore'' di workshop in totale,
mentre nel capitolo~\ref{cap:z1-feedback} si parlava di circa dieci ore: la
differenza dipende probabilmente dal conteggio delle pause.}, con il
relatore alle prese con la voce un po' calata dal giorno prima (``la sfida
principale di oggi sarà contro la mia voce'').

Il percorso del capitolo segue quello della lezione:
\begin{enumerate}
  \item il programma della giornata (sezione~\ref{sec:z2-programma});
  \item le domande sulla prima giornata, che si trasformano in una lunga
        sessione su \textbf{colloqui e selezione}
        (sezione~\ref{sec:z2-colloqui});
  \item il \textbf{time management}: pianificare la giornata, le liste, le
        checklist, le interruzioni (sezione~\ref{sec:z2-pianificare});
  \item la \textbf{matrice di Eisenhower} per prioritizzare
        (sezione~\ref{sec:z2-eisenhower});
  \item come proteggere il proprio tempo: sprechi, distrazioni, il saper dire
        di no (anche ai meeting), il mito del multitasking e la chiusura della
        giornata (sezione~\ref{sec:z2-proteggere});
  \item le domande sul time management e una parentesi su tempo libero, ferie
        e benessere in azienda (sezioni~\ref{sec:z2-qa-tempo}
        e~\ref{sec:z2-benessere});
  \item il \textbf{brainstorming}: obiettivi, problemi tipici, un metodo
        completo (\emph{rapid ideation} e \emph{dot voting}) e la
        \textbf{matrice impatto--sforzo} per decidere
        (sezioni~\ref{sec:z2-brainstorming}, \ref{sec:z2-metodo}
        e~\ref{sec:z2-impatto});
  \item i consigli pratici, le domande sul brainstorming e il principio del
        \textbf{disagree and commit} (sezioni~\ref{sec:z2-qa-brainstorming}
        e~\ref{sec:z2-disagree});
  \item il \textbf{leading from within}: perché la leadership non è un
        titolo, le carriere divise tra manager ed esperti
        (sezione~\ref{sec:z2-leadership});
  \item il ``momento Piero Angela'' sulla biologia del gruppo e le domande
        che ne sono nate (sezioni~\ref{sec:z2-biologia}
        e~\ref{sec:z2-qa-biologia});
  \item la \textbf{piramide della fiducia}, il prendersi cura del gruppo e
        che cosa cambia quando si diventa manager
        (sezioni~\ref{sec:z2-piramide}, \ref{sec:z2-gruppo}
        e~\ref{sec:z2-manager});
  \item le domande finali e la chiusura del workshop
        (sezioni~\ref{sec:z2-qa-finale} e~\ref{sec:z2-chiusura}).
\end{enumerate}

Come per il capitolo precedente, una premessa di metodo: buona parte della
giornata è fatta di domande della classe e di racconti dell'esperienza del
relatore. Sono riportati con cura perché sono il modo in cui i principi
vengono resi concreti, e perché spesso contengono le indicazioni più
direttamente utilizzabili (soprattutto sui colloqui di lavoro).

% ---------------------------------------------------------------------
\section{Il programma della giornata}
\label{sec:z2-programma}

Come annunciato alla fine della prima giornata
(sezione~\ref{sec:z1-domani}), gli argomenti sono due:
\begin{itemize}
  \item un ``doppio topic'', \textbf{time management e brainstorming};
  \item il \textbf{leading from within}, che il relatore definisce ``la
        versione avanzata di ieri'': come guidare un team stando dentro al
        team, verso i colleghi e verso il gruppo con cui si lavora.
\end{itemize}
Poiché l'ordine è indifferente, il relatore lascia scegliere alla classe: dalle
poche preferenze espresse, la leadership va al pomeriggio, così che le domande
finali si concentrino su quella, mentre la mattina è dedicata a time
management e brainstorming. Il materiale c'è, ma il relatore avverte che la
giornata potrebbe finire prima delle quattro: non vuole trattenere nessuno se
non c'è niente di specifico da fare. Per l'eventuale tempo avanzato prima della
pausa pranzo propone un \textbf{argomento bonus}: i colloqui di lavoro,
selezione e \emph{recruitment}, visto che il relatore ha fatto e fa tuttora
selezione del personale e che per molti studenti sarà la cosa più immediata da
affrontare finito il percorso di studi (domande comuni, come rispondere,
eccetera).

Prima di cominciare, però, il relatore apre a \textbf{domande sulla giornata
precedente}: la classe ha avuto modo di ``metabolizzare'' i contenuti e
potrebbero essere emersi dubbi a posteriori. È proprio da una di queste domande
che nasce, in anticipo, la parte sui colloqui.

% ---------------------------------------------------------------------
\section{Colloqui di lavoro: capire che cosa cerca l'azienda}
\label{sec:z2-colloqui}

\subsection{L'annuncio come prima fonte di informazione}

La domanda: il giorno prima si era detto che è importante capire che cosa
l'azienda cerca davvero; ma se il primo colloquio è con le risorse umane, cioè
con qualcuno che magari ``di tecnico non sa'', come si fa a capirlo? Il
relatore risponde subito, nel dubbio che più tardi manchi il tempo.

La prima fonte è \textbf{l'annuncio}. Non viene scritto a caso: è approvato
sia dall'HR sia dall'\emph{hiring manager}\footnote{Nella trascrizione
``Airing Manager'' e ``Iron Manager''; allo stesso modo ``e-char'' e ``char''
stanno per HR.}, cioè la persona che seguirà il candidato una volta assunto,
e spesso dall'intero dipartimento HR. È calcolato e scritto in un certo modo
per far capire che cosa serve. Molti annunci moderni, oltre alla parte su
\emph{che cosa si verrà a fare}, contengono una parte su \emph{che cosa ci si
aspetta da questa persona}, cioè il comportamento e il carattere che servono
in quella posizione: ed è quella che dice molto sulla \textbf{cultura
aziendale}.

L'esempio è quello di due annunci per la stessa posizione (programmatore
senior) descritta in modo diverso:
\begin{enumerate}
  \item ``cerchiamo un programmatore esperto che possa prendersi cura dei
        nostri programmatori junior, aiutarli a crescere, sviluppare un
        ambiente di lavoro inclusivo e stimolante, in modo che i nostri
        progetti futuri siano all'avanguardia'';
  \item ``cerchiamo un programmatore esperto che sappia rispettare le
        deadline, seguire le indicazioni dei dati, sia indipendente nel suo
        lavoro e sappia relazionarsi con gli \emph{stakeholder} per capire e
        proporre nuove direzioni per il progetto''.
\end{enumerate}
Nel primo caso ci si aspetta un team di junior, probabilmente con un'età media
più bassa, con cui ci si relazionerà soprattutto; per un senior è l'occasione
di sviluppare le doti da \emph{people manager}, e al colloquio arriveranno
domande come ``hai mai lavorato con dei team? Con che età media? Ti sei mai
trovato a guidare un gruppo?''. Nel secondo caso il senior deve eseguire un
piano deciso o proporre qualcosa, e si relazionerà soprattutto \emph{verso
l'alto}; al colloquio arriverà piuttosto ``fammi questo test di
programmazione''. Sono, di fatto, due colloqui completamente diversi.

\subsection{Quanti colloqui e con chi}

Ormai è molto raro essere assunti con un solo colloquio: ne servono almeno
due, realisticamente \textbf{tre}, che è diventata la media. L'iter tipico:
\begin{enumerate}
  \item il primo con l'HR;
  \item il secondo con l'HR o con il futuro manager;
  \item il terzo sicuramente con il futuro manager, perché l'azienda vuole far
        conoscere al candidato la persona con cui lavorerà, e viceversa.
\end{enumerate}
Già le domande diverse dei vari interlocutori aiutano a capire che cosa
cercano. Ma c'è una parte fondamentale, che secondo il relatore molti
sottovalutano: il momento finale delle \textbf{``hai domande?''}.

\subsection{Le domande da fare al selezionatore}

Un principio di base: \textbf{i colloqui non sono ``devo farmi selezionare
dall'azienda'', sono ``dobbiamo selezionarci a vicenda''}. Le domande vanno
quindi preparate: ci sarà sempre qualcosa che, dalla \emph{job description} o
dal colloquio, non è chiaro al cento per cento. Nel dubbio, il relatore
suggerisce tre domande ``che funzionano sempre'':

\paragraph{``Se iniziassi domani, quale sarebbe il mio \emph{day to day}?''}
Ottiene due effetti. Il primo è pratico: descrive la giornata tipo e fa capire
se piacerebbe o no. Il secondo è più psicologico, ma il relatore sa che
funziona: costringe l'interlocutore a pensare al candidato \emph{già dentro
l'azienda}, cioè a che cosa potrebbe fare da subito e a quale impatto potrebbe
avere con le sue competenze.

\paragraph{``Se aveste una bacchetta magica, che cosa cambiereste subito
dell'azienda?''} Sposta la conversazione su di loro e va bene sia con l'HR sia
con l'hiring manager. Li obbliga a indicare una sfida interna che vorrebbero
risolvere, ed è sempre un'informazione utile: se entrambi rispondono che
cambierebbero il modo in cui si prendono le decisioni, viene da pensare che non
ci sia un percorso chiaro e trasparente da parte di chi sta in alto.

\paragraph{Lo scambio: ``vi do un motivo per non prendermi, voi datemene uno
per non venire''.} È più avanzata e il relatore non la usa sempre, ma la trova
particolarmente interessante. Permette al candidato di esprimere i propri dubbi
in modo costruttivo (``mi dite che lavorerei a stretto contatto con un team di
junior: un motivo per non prendermi è che non ho anni di esperienza nella
gestione di giovani sviluppatori, sarebbe la mia prima volta; però penso di
poterlo fare perché ho questo interesse e queste competenze''). Dall'altra
parte, il selezionatore dirà che cosa secondo lui sarebbe più difficile per il
candidato: si scopre così che cosa sta già pensando, e l'informazione torna
utile anche per i colloqui successivi.

\paragraph{E se il selezionatore si innervosisce?} ``Se un recruiter si
innervosisce per le domande del candidato, ha sbagliato lavoro.'' Al massimo il
relatore si innervosisce se un senior gli chiede che cosa faccia di preciso un
programmatore senior (``lo fai da dieci anni, lo saprai tu''); ma se un junior
chiede che cosa dovrebbe fare ogni giorno, lo si spiega volentieri. E in realtà
vale anche per i senior, perché \textbf{in ogni azienda lo stesso ruolo fa cose
diverse}: conviene sempre farsi spiegare che cosa intendono per quel titolo.
L'esempio viene dal mondo delle community: \emph{community manager} a Riot
indicava la persona a capo del team di \emph{community representative} di una
nazionalità (``community manager italiano''); a Blizzard indicava la persona
sopra l'intero dipartimento che gestiva quei team. Candidarsi per lo stesso
titolo non garantisce di aver capito per che cosa ci si sta candidando, e i
recruiter lo sanno.

Se poi un recruiter si innervosisce comunque, sta dando più risposte di quelle
che pensa: vuol dire che l'azienda non vuole domande, o non vuole interazioni
del tipo ``ma perché facciamo questa cosa?''. Che abbia formato i suoi
recruiter a innervosirsi, o che accetti che lo facciano, è informativo.

\subsection{Le domande che dicono molto sull'azienda}

Allo stesso modo, come detto il giorno prima, è informativo se un recruiter
chiede a una donna se vuole avere figli: ``ah, quindi a voi interessa sapere se
voglio avere figli: questo già mi dice che probabilmente non voglio lavorare
qui''.

Un esempio \emph{gender neutral}: la richiesta di vedere la \textbf{busta paga
precedente}. ``Quanto è la tua RAL? Mandami una busta paga.'' La risposta del
relatore è no: ``ti ho detto qual è la mia RAL; se vogliamo instaurare un
rapporto in cui non ti fidi di quello che dico, abbiamo un altro problema''. È
di nuovo \emph{default to trust} (sezione~\ref{sec:z1-default-trust}). Il
relatore è quasi sicuro che legalmente non si possa chiedere la busta paga di
un candidato\footnote{È l'opinione del relatore, espressa come tale (``sono
quasi sicuro''); non si tratta di un parere legale e non è stata verificata.}:
a quel punto diventa lui antagonista del recruiter, e visto che la
conversazione è finita in una ``zona grigia'', fa lui le domande (``come mai è
legale questa cosa che mi chiedi?''). Ha già deciso che probabilmente non
lavorerà lì, ma così magari evita che i candidati successivi debbano subire le
stesse richieste. Cita l'esempio perché sa che in Italia è abbastanza comune,
mentre all'estero non gli è mai successo, anche se avrebbe potuto ``sparare
qualsiasi cifra''.

Il motivo è che \textbf{non c'è alcun onere di dimostrare quanto si
guadagna}: chi spara troppo alto si autoesclude. Recruiter e manager non
arrivano al colloquio con carta bianca né con l'indicazione di risparmiare il
più possibile, ma con un \emph{range}: ``per questa posizione possiamo
spendere dai 30 ai 40 mila euro l'anno; paga questa persona a seconda della
sua seniority''. Se un candidato che guadagna 30 dice di prenderne 55 per
``provare a fare i soldi'' e il budget arriva al massimo a 40, il relatore non
chiederebbe nessuna prova: direbbe ``grazie, siamo in un range diverso, possiamo
salutarci''. La domanda sulla busta paga è quindi \emph{formulata male}, perché
è poco furba. Quella giusta è: \textbf{``hai già una RAL in mente per questa
posizione?''}. Si ottengono le stesse informazioni senza passare per la zona
grigia.

\subsection{La cifra minima mentale}

E se la domanda la fanno a un junior che non ha mai lavorato? Il consiglio vale
per tutti: \textbf{preparatevi sempre la vostra cifra minima mentale}, per
qualsiasi offerta di lavoro, prima ancora di fare il colloquio. Non serve
conoscere la media italiana per quella posizione o quanto paga quell'azienda:
serve sapere \emph{quanto si è disposti ad accettare senza andare in negativo},
tenendo conto dei propri costi (la benzina per andare in ufficio ogni giorno, un
eventuale trasferimento, e così via). Quella cifra va sempre avuta in testa.

\subsection{La scrematura iniziale e il portfolio}

Uno studente chiede come funzioni la risposta a una candidatura: arriva un
rifiuto diretto, oppure c'è una scrematura iniziale e poi i colloqui? E vale
per tutti i ruoli? Sì: c'è una \textbf{scrematura iniziale} fatta dall'HR o, in
alcune aziende, direttamente dal software dell'HR. Si scarta chi è chiaramente
fuori da ciò che si cerca:
\begin{itemize}
  \item si cerca un \emph{graphic designer} e arriva il CV di un programmatore
        (``faccio questi esempi perché purtroppo succede'');
  \item si cerca un \emph{animator} tecnico bravo con MotionBuilder, lo si
        scrive nell'annuncio, e il candidato risponde ``so fare tutto tranne
        MotionBuilder'': c'è un motivo tecnico per non poterlo prendere;
  \item problemi di seniority: si cerca un senior e risponde un junior con un
        anno di esperienza.
\end{itemize}
Se la figura non è già chiaramente fuori target, si fa il primo colloquio.

\paragraph{Gli anni di esperienza negli annunci junior.} Spesso gli annunci
junior chiedono come minimo tre anni di esperienza. \textbf{Se l'annuncio dice
``junior'', ignorate i requisiti di esperienza}: o sono lì per sbaglio, o sono lì
per scremare. Il secondo caso è più \emph{borderline}, ma il relatore l'ha
visto fare in alcune aziende americane: si mettono requisiti molto alti,
chiaramente irraggiungibili, per vedere chi ha la ``faccia tosta'' di mandare
comunque il CV (``non me ne frega se chiedete un \emph{bachelor's degree}, non ce
l'ho, mando lo stesso''). Così però si filtra una grossa fetta di popolazione e
si seleziona un certo tipo di carattere: si stanno cercando ``i Top Gun''. Non è
quindi necessario iniziare con uno stage: si può iniziare direttamente da
junior.

\paragraph{Chi guarda il portfolio.} Il relatore non ha mai avuto un portfolio,
ma per le figure che lo richiedono parte comunque dal CV (per escludere, ad
esempio, il designer che si candida da programmatore) e guarda il portfolio
solo a parità con altri candidati: da HR può aprirlo e dire ``bello'', ma non
può dare un giudizio tecnico. Il portfolio lo apre chi sarà l'hiring manager,
che è sempre molto più tecnico. L'ordine quindi è: l'HR guarda il CV e si salva
il portfolio; se il candidato passa la prima scrematura, CV e portfolio vanno
a un tecnico che conferma (``sì, sa fare queste cose'') o boccia (``il
portfolio indica chiaramente che non sarebbe in grado, digli pure che non
facciamo il colloquio'').

\subsection{Test pratici, da casa e con qualsiasi strumento}

\paragraph{Si può chiedere un test pratico al posto del portfolio?} Uno
studente chiede se, con un portfolio che non rispecchia le proprie capacità, si
possa chiedere direttamente un test. No: se l'azienda ha bisogno di verificare
una competenza ha già pronto un test e lo fa fare a prescindere. Il relatore,
per esempio, fa fare un test ai programmatori e lo fa correggere dal suo senior
developer; per gli artisti si chiedono animazioni di prova, chiaramente veloci.
Non chiede mai test troppo lunghi, perché non vuole far lavorare nessuno gratis:
non chiede ``fammi un pitch di progetto'', chiede uno spunto per capire come si
ragiona e come si lavora. Per il resto, \textbf{se il portfolio non è
indicativo, va aggiornato il portfolio}.

\paragraph{Test dal vivo o a casa?} Il relatore li fa fare sempre a casa. Sa
benissimo che un programmatore a casa lavorerà con due monitor, il codice da
una parte e Google o ChatGPT dall'altra. Ne ha discusso con i suoi senior
programmer: interessa che qualcuno sappia le cose a memoria, o che sappia
trovarle e fare quel lavoro a prescindere? A lui interessa che il lavoro
\textbf{venga fatto}: che il candidato abbia imparato il codice a memoria o che
in venti minuti trovi le informazioni, le metta insieme, le faccia elaborare a
ChatGPT, le corregga e consegni il codice che serve, non ha preferenze. Quando
si lavora le indicazioni saranno ``fai questa cosa il prima possibile'': se la
si sa fare anche con strumenti aggiuntivi, ben venga, ``anzi, spiegalo anche
agli altri''. Farla a memoria va benissimo, ``ma perché mettere l'\emph{hard
mode}?''. E comunque ci sono cose che non si possono delegare a uno strumento:
se arriva un codice che non funziona e va fatto il \emph{debug}, bisogna
saperlo fare e saper progettare i test. La figura più completa, per il
relatore, è anche quella che sa trovarsi le informazioni e gli aiuti.

Le prove dal vivo sono soprattutto domande di \emph{problem solving}: il
relatore vuole sentire come si ragiona e chiede espressamente di \textbf{pensare
ad alta voce}, per seguire il percorso mentale che porta alla conclusione. Sono
più domande impegnative che veri test. I test veri e propri si fanno con i
propri tempi e i propri strumenti: ``scegli il momento, ti mandiamo il test, hai
due ore per rimandarlo indietro''. Questo è l'accordo; in quelle due ore si può
fare quello che si vuole.

\subsection{Junior, mid e senior}

\paragraph{Perché ci sono così pochi annunci per figure mid?} Uno studente ha
notato su LinkedIn e sulle piattaforme di recruiting un divario notevole tra le
richieste di junior e quelle di mid. Il motivo è che per molte aziende,
Stormind inclusa, \textbf{la differenza tra junior e mid è sparita}: o si è
senior o non lo si è. Ci si è resi conto che la differenza tra un junior e un
mid è che il mid si è abituato a lavorare in quell'azienda e il junior non
ancora; ma un mid che arriva con tre anni di esperienza in un altro studio (ad
esempio Milestone\footnote{Nella trascrizione ``se ha fatto tre anni in
milestone'': si tratta verosimilmente di Milestone, studio di sviluppo
milanese.}) va comunque ``ripreso come junior'', perché si lavora in modo
diverso e deve reimparare. Un senior, invece, ha molte più competenze e le sa
applicare in un certo modo, e quella distinzione resta netta.

\paragraph{Ci si può candidare anche se il livello non coincide?} Un'altra
studentessa dava per scontato che, se un annuncio non specifica ``junior'', sia
per mid. Se non è specificato ``senior'', secondo il relatore si può mandare il
CV chiedendo di valutare. Certo, se è specificato junior e si è junior le
possibilità di essere contattati sono maggiori. Ma vale lo stesso discorso del
livello di inglese fatto il giorno prima: si cerca qualcuno che \emph{sappia
fare quelle cose}, non una certificazione. Se serve un programmatore che sappia
usare Unreal Engine~5 e sistemare la cosa X o Y, che sia mid da due anni o
junior da uno non è un problema.

\subsection{Gli strumenti di intelligenza artificiale}

Una domanda in chat riguarda l'uso quotidiano dell'intelligenza artificiale in
azienda. Il relatore sta facendo fare corsi di IA al suo team ``perché voglio che
la gente lavori meno'': vuole che si sfrutti qualsiasi strumento disponibile,
perché gli altri lo stanno facendo e non vuole restare indietro. Se un lavoro
si può fare in tre ore di meno, non c'è alcun vantaggio a farlo in modo lungo e
tortuoso. Il rischio esiste: se si fa fare tutto a ChatGPT e non ci si tiene al
passo con il proprio \emph{know-how}, quando ci sarà da risolvere qualcosa ci si
potrebbe ``impallare''. Ma sta al singolo restare aggiornato e saper fare il
proprio lavoro: \textbf{gli strumenti devono essere strumenti, non un sostituto
delle competenze}. ``Non è una partita di calcio, con regole su come si devono
fare le cose o un \emph{dress code}'': il lavoro va fatto nel modo più efficiente
possibile.

A uno studente a cui era saltata la connessione proprio in questo momento il
relatore risponde che la lezione è registrata e si può recuperare dal video.
Poi chiude la parentesi: ``la parte difficile era finire ieri in tempo, ce
l'abbiamo fatta''.

% ---------------------------------------------------------------------
\section{Time management: pianificare la giornata}
\label{sec:z2-pianificare}

\subsection{Time management o prioritizzazione?}

La domanda di partenza alla classe è: \emph{chi di voi, in questo momento, fa
time management su se stesso}, cioè si organizza la giornata all'inizio e
organizza il proprio tempo? Qualcuno risponde di sì, qualcuno ha uno
\emph{schedule} settimanale con orari da rispettare, qualcuno ``ci prova''
(``hai fatto bene a specificare \emph{ci provo}: magari non la rispetto''). Il
relatore lo chiede perché è uno di quegli argomenti in cui molti, dopo aver
visto la spiegazione, si accorgono di non fare davvero time management: fanno
piuttosto \textbf{prioritizzazione}, ad esempio una lista di cose da fare, che
è un po' diverso.

Nessuna delle tecniche è complicatissima, ma è utile sapersi muovere. Il
relatore colloca volutamente questo modulo \emph{dopo} quello sul feedback: saper
gestire il proprio tempo permette anche di dare feedback agli altri sulla loro
gestione del tempo.

\subsection{La \emph{must do list}}

Il punto di partenza è il \textbf{planning iniziale}: non improvvisare la
giornata. Improvvisando si è in balia di una grande limitazione umana: ci si
dimentica le cose, oppure ci si lascia trascinare dal task più urgente o
dall'ultimo arrivato. Esistono modi di dire come ``se arriva una mail o rispondi
subito o diventa un task per domani'', perché il ``rispondo dopo'' non arriva
mai; più avanti si vedranno modi più raffinati di gestire le zone grigie, ma la
pianificazione resta fondamentale.

La prima cosa da fare \emph{non} è la lista delle cose da fare oggi, ma la lista
delle cose che \textbf{assolutamente} vanno fatte oggi. Le due liste non devono
coincidere: la seconda richiede un livello di ragionamento in più. Sono le cose
per cui, se alle sei si stacca e non sono fatte, si resta finché non sono fatte,
perché sono imprescindibili.

\paragraph{Il primo momento per dire ``potrei non farcela''.} La \emph{must do
list} è anche il primo indicatore della giornata. Se è troppo corposa (dipende
naturalmente dalla complessità dei task), quello è il momento, ``alle 9:01 se
iniziate alle 9'', per dire al proprio manager \emph{oggi sarò in difficoltà}.
Alle 9:01 il manager ha l'intera giornata per aiutare; se invece si pecca di
ottimismo o di presunzione (``ce la posso fare'') e glielo si dice alle 17:59,
non può fare nulla. È lo stesso discorso del giorno prima
(sezione~\ref{sec:z1-triangolo}): anche facendo la to-do list si applicano i
concetti della cultura del feedback, ed è il motivo per cui questo modulo viene
dopo. Ci saranno vari momenti di questo tipo nella giornata, ma questo è il
primo.

\paragraph{E se ci si accorge a metà giornata di non farcela?} Uno studente
chiede che cosa fare se, convinto di farcela, si accorge a metà giornata di non
poter completare bene la \emph{must do list}. Il relatore lo concederebbe
\emph{una volta}: se si vogliono testare i propri limiti, va bene, ma
comunque si dice al capo ``oggi potrei essere tirato, ci provo'', e sarà lui a
valutare se il rischio è accettabile. Dipende dal rischio: se il \emph{must do}
è consegnare una milestone a un cliente, ``nel dubbio facciamolo fare a
qualcun altro, non possiamo correre quel rischio''; se il rischio è che domani
un animator riceva il codice con tre ore di ritardo, lo si può correre. ``È un
po' questo il lavoro.''

\paragraph{Cinque minuti e una stima.} Tutto questo va fatto in circa cinque
minuti. Se non è chiaro che cosa si deve fare in una giornata, \textbf{si
chiede}: ogni giorno dovrebbe essere chiaro. Poi si cerca di \textbf{stimare il
tempo} di ogni cosa. Si diventa più bravi col tempo: alla quinta volta che si
fa il report mensile si sa che ci vuole un'ora e ventidue minuti; la prima volta
si può chiedere a un collega quanto ci mette lui, sapendo che per noi potrebbe
essere diverso.

\subsection{To-do list e checklist}

Esistono molti strumenti per tenere to-do list e checklist; il relatore non ne
consiglia uno in particolare e cita come esempi Any.do, i promemoria di Windows
o Samsung Notes\footnote{Nella trascrizione ``AnyDo, Windows Memo''. Any.do è
un'app di liste e promemoria; ``Windows Memo'' indica probabilmente le note o i
promemoria integrati in Windows (Sticky Notes, Microsoft To Do).}: ``potete usare
davvero la qualunque''. L'importante è tenere la lista \textbf{visibile} da
qualche parte sul computer. Alcuni accorgimenti:
\begin{itemize}
  \item i \textbf{must do} vanno sempre \textbf{in cima} alla lista;
  \item quando durante la giornata arriva un task nuovo, la cosa da fare subito
        non è \emph{eseguirlo}, ma \textbf{aggiungerlo alla lista}: ``se arriva
        una mail, aggiungete subito alla lista \emph{rispondi alla mail}'',
        decidendo se è diventato un \emph{must do} o se è comunque un task per
        oggi;
  \item non limitarsi alle cose operative: il relatore ha iniziato ad
        aggiungere anche task come ``parla con Tizio'' o ``parla con Caio'',
        perché è da un po' che non chiede loro un aggiornamento; se non li
        mette in lista si accorge di non farli;
  \item spuntare le voci completate: per quanto minimo, rilascia
        \textbf{dopamina} (se ne parlerà nel pomeriggio,
        sezione~\ref{sec:z2-biologia}), aiuta a motivarsi e ad andare avanti.
        Sono piccole cose, ma si sommano.
\end{itemize}

\paragraph{Priorità contro recenza.} Che cosa succede se si esegue sempre
l'ultimo task arrivato? Lo chiede il relatore alla classe. Si interrompe il
flusso di lavoro e si cambia continuamente attività, certo; ma soprattutto si
stanno \textbf{mescolando priorità e recenza}: si sta dicendo ``non importa come
ho prioritizzato i task, la mia priorità è sempre l'ultimo arrivato'', che non
vuol dire fare le cose più importanti. È il motivo per cui tanti dipartimenti IT,
quando si chiede loro qualcosa, rispondono ``mi apri un ticket?''. È diventata
quasi una battuta, ma il concetto è chiarissimo: devono poter prioritizzare
quella richiesta rispetto alle altre 112 arrivate in giornata, e senza ticket
non possono confrontarla con le altre. Il relatore è totalmente a favore: ``certo
che ti apro un ticket, altrimenti tu come fai a tracciare il tuo lavoro?''.

\subsection{La checklist: l'esempio dell'aviazione}

La \textbf{checklist} non è la task list. La task list contiene le cose da fare;
la checklist contiene \textbf{i passi da seguire quando si esegue un task
specifico}. L'esempio è quello delle \textbf{checklist pre-volo}, introdotte
nell'aviazione (secondo il relatore tra gli anni Sessanta e Settanta): chi vola
spesso nota che i piloti seguono sempre un ordine simile nei controlli, tanto
che il relatore capisce quanto manca al decollo da quando controllano i flap, i
motori, quando accendono e spengono le luci. L'introduzione delle checklist
avrebbe ridotto di circa il 75\% gli incidenti in fase di decollo\footnote{Date
e percentuale sono quelle riferite in lezione e non sono state verificate.
L'origine della checklist pre-volo viene di solito fatta risalire al 1935, dopo
l'incidente del prototipo del Boeing B-17, causato da un controllo dimenticato da
un equipaggio molto esperto; la diffusione sistematica delle checklist e dei
protocolli di cabina è invece successiva.}, e soprattutto per i \textbf{piloti
più senior}. Perché? Perché chi ha fatto un decollo centomila volte non va a
controllare l'alettone destro: ha funzionato centocinquantamila volte, funzionerà
anche stavolta. Quando ci si abitua a fare qualcosa, si sbaglia proprio quando
qualcosa non va. La checklist ha obbligato a fare quei controlli ogni volta e ha
ridotto in modo drastico gli incidenti.

Il relatore usa checklist anche per cose banali, ad esempio le \textbf{email}:
essendo una persona abbastanza ansiosa su queste cose, se prepara una mail alla
leggera finisce per passare troppo tempo a ricontrollare l'oggetto, il
destinatario, l'allegato, e comunque si dimentica qualcosa. Seguire una checklist
sembra più complesso, ma \textbf{fa risparmiare tempo}: se la checklist è
completa, si sa di aver fatto tutto e si può andare avanti senza continuare a
pensare ``ho chiuso il gas? ho spento la luce?''. È come la lista della spesa: se
tutto è segnato e preso, punto. È un investimento: trenta secondi in più per
fare la cosa, invece di cinque minuti a sperare di averla fatta giusta.

\paragraph{E se ci si dimentica qualcosa nella checklist?} Succede. A quel
punto \emph{preparare la checklist} è un task in sé: un giorno in cui si ha
tempo si fa quel task pensando a quali passi si seguono e a come assicurarsi che
funzioni. Nessuna checklist è perfetta: ci saranno sempre compromessi tra ``devo
davvero controllarlo ogni volta?'' e ``posso farne a meno?''. Però, in generale,
\textbf{se qualcosa anche solo una volta ha dato problemi, va in checklist},
perché vuol dire che potrebbe darne di nuovo. Nessun manager rimprovererà
``come hai fatto a non prevedere che questa cosa poteva andare male?'': si
diventa esperti anche nelle checklist di un task facendo quel task, ed è un
processo iterativo. Se qualcosa è stato dimenticato e ha creato problemi, lo si
aggiunge per la volta dopo. Sbagliare una o due volte non è un problema;
l'importante è che la situazione migliori la volta successiva.

\paragraph{L'inizio di un junior e l'\emph{onboarding}.} A uno studente che
osserva che l'inizio di un junior è quindi inevitabilmente un po' travagliato,
il relatore risponde che è giusto così: ``se volessi il dipendente perfetto da
subito, cercherei un senior, non un junior''. Allo stesso modo non capisce le
aziende che prendono un junior e non lo fanno seguire da nessuno: che cosa ci si
aspetta, che assorba passivamente tutto quello che ha intorno e impari ``per
osmosi''? Può succedere, ma imparerà a fare le cose come vuole lui, perché
nessuno gli sta dicendo ``meglio farle così''. Moltissime aziende, oltre al
manager (che magari ha quindici \emph{report} e non ha il tempo materiale di
seguire un junior come andrebbe seguito), assegnano un \textbf{referente}, di
solito un senior, con il compito di prendersi cura della persona nuova e
insegnarle a fare le cose. Altrimenti l'\textbf{onboarding}, cioè il tempo
necessario per arrivare al cento per cento nel proprio lavoro, dura un anno.
Più risorse l'azienda dà all'inizio, prima la persona diventa operativa: è
nell'interesse dell'azienda dare il massimo supporto a chi arriva.

Il relatore diffida delle aziende che ``vogliono vedere se uno se la cava''. È
il discorso dei cento metri e della corsa a ostacoli\footnote{Il relatore lo
introduce come un discorso già fatto (``si torna al discorso: vuoi fare i cento
metri o la corsa a ostacoli?''), ma il passaggio non compare nella trascrizione
della prima giornata riportata nel capitolo~\ref{cap:z1-feedback}.}: non si può
far correre uno sui cento metri e l'altro con gli ostacoli e aspettarsi tempi
simili, perché chi ha gli ostacoli ci metterà per forza di più. Il supporto ai junior è un
investimento sensato.

\subsection{\emph{First contact resolution}}

Un concetto che viene dai call center: \textbf{quante richieste vengono risolte
al primo contatto}. È importante quando un task richiede iterazioni o
comunicazione con altri: bisogna organizzarsi come se si dovesse risolvere la
cosa \emph{in una sola passata}. Se viene chiesto un modello 3D, lo si fa giusto
la prima volta, al meglio delle proprie possibilità. Se si deve parlare con
qualcuno, si raccolgono tutte le richieste in una.

L'esempio è quello dei ticket che il relatore mandava ai colleghi del supporto di
Blizzard: ``ciao, ho un problema, non riesco a loggarmi'' --- ``mi dici lo
username?'' --- ``è questo'' --- ``mi dici l'email?'' --- ``è questa'' --- e così
via, con un paio d'ore tra un messaggio e l'altro perché il ticket deve fare il
suo giro. Tutto poteva essere riassunto in ``ciao, mi servono username, email e
questo e quest'altro''. Unificare le richieste su una singola cosa aiuta a
risolverla il prima possibile; bisogna avere chiaro \emph{che cosa serve prima
di fare una richiesta}, in modo che ce ne sia una sola. In molti sistemi di
supporto, se si apre un secondo ticket per aggiungere informazioni al primo, il
primo, anche se era quasi in lavorazione, viene rimandato in fondo alla coda:
serve a disincentivare chi manda continuamente ticket del tipo ``mi sono
ricordato questa cosa''. O si manda tutto subito, o ci vorrà tempo.

\subsection{Non iniziare senza le informazioni}

Corollario: \textbf{se non si ha tutto quello che serve per completare un task,
non lo si inizia}. Prima ci si procura quello che manca. Al relatore capita con
la pianificazione dei due anni di lavoro successivi: se non ha tutte le
informazioni, prima lavora per averle (``mi dai questa informazione, mi dai
quell'altra''), poi pianifica; altrimenti pianificherebbe qualcosa di cui non ha
completa visibilità.

\subsection{Il costo delle interruzioni}

Alternare continuamente l'attenzione tra i task e, soprattutto, reagire di
continuo agli input esterni (la mail, il messaggio, il collega che parla) fa
\textbf{sparire il tempo dalla giornata}. Se mentre si prepara un report in
Excel si viene interrotti (dal gatto, da una mail, da un collega, da una
telefonata) e poi si torna al foglio, servono in media \textbf{un quarto d'ora}
per tornare mentalmente al cento per cento; chi è particolarmente bravo ci mette
dieci minuti, ma si parla comunque di dieci-quindici minuti\footnote{Ordini di
grandezza come riportati in lezione. Le ricerche più citate sul tema (ad esempio
quelle di Gloria Mark sulle interruzioni nel lavoro d'ufficio) indicano tempi
medi di ripresa anche superiori, intorno ai venti minuti.}. Se si viene
interrotti di continuo, si passa più tempo a cercare di tornare al cento per
cento che a lavorare. È una partita persa in partenza.

Per questo il giorno prima si era parlato dei \textbf{semafori} sulla scrivania
(sezione~\ref{sec:z1-dare}): il semaforo può essere anche semplicemente il
proprio stato su Slack, su Discord o su qualunque strumento si usi con i
colleghi. Se si è occupati, si è occupati: è una propria valutazione che,
accettando di essere disturbati di continuo, quella cosa non si riuscirà a
farla, o almeno non in tempo.

% ---------------------------------------------------------------------
\section{Prioritizzare: la matrice di Eisenhower}
\label{sec:z2-eisenhower}

Ora che si hanno task list, checklist, \emph{must do} e \emph{nice to do}, resta
la domanda di prima: se arriva un task nuovo e si decide di non fare per prime le
cose nuove, \emph{come si capisce se è davvero urgente}? Serve un metodo, ed è la
\textbf{matrice di Eisenhower}\footnote{Nella trascrizione anche ``Jason
Howard''. Dal nome del presidente statunitense Dwight D.~Eisenhower, a cui si
attribuisce la distinzione tra urgente e importante.} (la classe prova a
indovinarne il nome proponendo altre matrici, come una matrice dei rischi, ma è
un'altra cosa). Ogni task che arriva va posizionato su questa tabella:

\begin{center}
  \begin{tabular}{@{}p{0.2\textwidth}p{0.36\textwidth}p{0.36\textwidth}@{}}
    \toprule
    & \textbf{Urgente} & \textbf{Non urgente} \\
    \midrule
    \textbf{Importante} &
    \textbf{Fallo} (\emph{do}): diventa un task di oggi. &
    \textbf{Pianificalo} (\emph{schedule}): va fatto, ma non subito. \\
    \addlinespace
    \textbf{Non importante} &
    \textbf{Delegalo} (\emph{delegate}): va fatto, ma non serve che lo faccia
    tu. &
    \textbf{Eliminalo} (\emph{delete}): non lo fai. \\
    \bottomrule
  \end{tabular}
\end{center}

\paragraph{Urgente e importante: lo fai.} Sono di solito task con una
\textbf{deadline molto chiara} (entro oggi, entro un'ora, entro X) e con
\textbf{conseguenze} se non vengono fatti: da qui l'``importante''.

\paragraph{Importante ma non urgente: lo pianifichi.} Va fatto, ma non subito:
o lo si mette tra le cose di oggi ``ma dopo'', o anche alla settimana successiva.
L'esempio: un collaboratore scrive ``il mio rinnovo è tra tre mesi, ne possiamo
parlare?''. È importante, perché il rinnovo di una persona è importante; non è
urgente, se ne può parlare anche domani, così da potersi preparare e fare altre
cose nel frattempo.

\paragraph{Urgente ma non importante: lo deleghi.} Qui ``non importante''
\textbf{non} vuol dire che il task non sia importante, ma che \textbf{non è
importante che lo faccia tu}, che è molto diverso. Se arriva ``mi serve questo
cambio a questa animazione 3D'' e ciò che conta è il cambio, non chi lo fa, si può
dire ``sono d'accordo che è urgente, ma non serve che lo faccia io'' e chiedere a
qualcun altro. Sono task che vanno fatti ma non richiedono la propria particolare
\emph{expertise}.

Il relatore insiste molto su questo punto con manager e senior: quando arriva un
task urgente, i senior tendono a dire ``lo faccio io, perché è urgente''. Ma se un
junior può farlo, lo si fa fare al junior e si investono le competenze più alte
del senior in qualcos'altro: il senior può fare il lavoro del junior, ma il
junior non può fare il lavoro del senior mentre il senior fa quello del junior.
Le persone non sono intercambiabili. Si usa il \textbf{minimo sforzo
necessario}: se una cosa la può fare un mid, la fa il mid, non il senior.

\paragraph{Né urgente né importante: lo elimini.} È il quadrante che quasi
nessuno usa: tutti fanno, delegano o pianificano, ma non pensano mai a
\emph{delete}. Invece è fondamentale. A volte è corretto dire: ``questa cosa non
dovrei farla io, non è urgente, non la faccio; sappi che non la faccio, e se c'è
un problema ne parli con il manager''. Non si sta delegando a qualcun altro: si
sta dicendo che quella cosa così com'è \textbf{non andrebbe fatta}, perché non
ne sono chiari l'importanza e il perché. Nella propria lista di oggi non ha
spazio, e non è abbastanza importante per essere fatta da sé nemmeno nei giorni
successivi.

Un esempio di \emph{delete} che non è un vero \emph{delete}: arriva una richiesta
di meeting di tre ore sul futuro del progetto, con invitati tutti gli art
director, il game director e così via. Il relatore spesso risponde ``io non
vengo'': per lui è \emph{delete}, ma il meeting si fa lo stesso. Uno studente
osserva che a volte si tratta di una richiesta mal formulata o poco chiara: è
esattamente così. ``Nella sua forma per me è no; se poi mi spieghi meglio perché
dovrei farlo io o perché è urgente, allora cambia.''

% ---------------------------------------------------------------------
\section{Proteggere il proprio tempo}
\label{sec:z2-proteggere}

\subsection{Evitare gli sprechi: un limite che nessuno può insegnare}

Evitare gli sprechi di tempo è difficile, soprattutto da remoto: è una sfida
costante con se stessi. Ha molto a che fare con il carattere, il
\emph{self-control} e (``passatemi anche questa'') la professionalità. Il
relatore può dare tutti i consigli del mondo, ma se dentro di sé non si percepisce
l'urgenza di fare il proprio lavoro quando lo si deve fare, non lo si può
insegnare. Se si è a casa e invece di lavorare si gioca a qualcosa su Steam,
nessuno deve spiegare che non si dovrebbe. Ci sono limiti oltre i quali non si
può andare.

Quello che si può fare è sapere che \textbf{ogni interruzione ha un costo
sommerso}: ogni messaggio, domanda, telefonata, ogni cambio di focus. Sono i
quindici minuti di prima, che possono arrivare a trenta a seconda del tipo di
task.

\paragraph{È più difficile tornare su un task logico o creativo?} Il relatore lo
chiede alla classe: fa più fatica un artista a tornare sul suo modello 3D o un
\emph{data analyst} a tornare sul suo foglio Excel? Le risposte si dividono
(``creativo'', ``entrambi''), ed è proprio il punto: dipende a chi lo si chiede. I
creativi rispondono ``creativo'', gli analisti ``analitico'', perché ognuno ha una
chiara comprensione di che cosa serve per essere davvero concentrati su ciò che
fa. Il creativo sa che per essere pienamente creativo gli servono l'ambiente, la
concentrazione, gli spunti, gli esempi da tenere chiari in testa; l'analista deve
avere in testa i conti, i fogli di dati e così via. Sono \textbf{ugualmente
difficili}, ma solo chi fa quella cosa ne apprezza fino in fondo la difficoltà. Si
può avere un assaggio dell'altra parte, ma la piena comprensione arriva solo
facendo: è come guardare uno sport in tv e pensare ``potrei farlo anch'io'',
finché non ci si trova lì.

\subsection{Distrazioni volontarie e involontarie}

Uno studente, che sta leggendo manuali di \emph{design thinking}, chiede come
conciliare tutto questo con l'idea che distrarsi, cioè non restare focalizzati,
aiuti a trovare un nuovo punto di vista per risolvere un problema. Qual è il
giusto equilibrio? Il relatore è \textbf{totalmente d'accordo} con quel concetto:
sta parlando di \textbf{interruzioni non volontarie}. Sa benissimo che dopo cinque
ore su un foglio Excel ``esplode la testa''; ma se si pianifica un quarto d'ora di
pausa e lo si rispetta, si sta investendo sulla propria creatività ed
efficienza. Non è un investimento dover rispondere a una chiamata ogni tre
minuti.

Anzi, il consiglio è \textbf{crearsi interruzioni pianificate}, sotto il proprio
controllo: fare pause, fare altro, soprattutto nel \emph{problem solving}. Il
relatore stesso, se non trova una soluzione, si ferma, fa un altro task e poi ci
torna; ma è una sua scelta. La discriminante è che la pausa sia \emph{scelta} e
non imposta: anche se si fanno le pause giuste, se si viene interrotti proprio
nel momento di massima concentrazione (raggiunta magari grazie alla pausa) si
butta via tutto. Uno studente cita la \textbf{tecnica del pomodoro}: è
esattamente quello, piccole pause pianificate che aiutano a essere creativi,
purché restino pianificate. Se i quindici minuti diventano ``vabbè, facciamo
mezz'ora'' o ``lo faccio nel pomeriggio'', non è più la tecnica del pomodoro:
diventano interruzioni, o un pretesto per non lavorare.

\subsection{Il peggior nemico siamo noi}

E chi sono i peggiori nemici in questo? ``Noi stessi'': siamo noi ad
autoconvincerci che venti minuti vanno bene quanto quindici. Le
\textbf{distrazioni autoinflitte} sono spesso le peggiori. Il relatore porta
l'esempio di una sua giornata, che ha annotato perché lo aveva fatto ridere
quanto tempo era riuscito a buttare via senza che nessuno glielo imponesse:
\begin{itemize}
  \item gli piace avere qualcosa in sottofondo mentre lavora (``ha a che fare con
        l'ADHD sicuramente''\footnote{Nella trascrizione ``IDACD''. Più avanti il
        relatore dice di avere l'ADHD ``come dicevo ieri''; il capitolo 0.1 non
        lo riporta, ma riferisce il suo racconto sull'``overthinking'' anche
        nell'ordinare una pizza (sezione~\ref{sec:z1-regalo}).}),
        e in genere lo aiuta. Il problema è quando il sottofondo è troppo
        interessante: la musica va bene, ma se nella playlist di Spotify parte un
        podcast di Alessandro Barbero sulla guerra civile americana, finisce per
        ascoltare quello più del lavoro;
  \item un'abitudine che ha poi perso: se faceva straordinari alle due di notte,
        da solo, per qualcosa di urgentissimo, apriva comunque Slack a
        controllare i messaggi, e la mail. ``Come se alle due di notte qualcuno
        mi scrivesse qualcosa di urgente'';
  \item la pausa con un po' di TikTok, dopo la quale si accorgeva di aver
        dimenticato su quale slide della presentazione stava lavorando.
\end{itemize}
La regola: \textbf{inserire i break è meglio che subirli passivamente}. Un conto
è dire ``faccio una pausa intenzionalmente'', un conto è ``mi sono distratto e ho
perso due ore di lavoro''. Chi ne è particolarmente soggetto non deve
demoralizzarsi: è anche un allenamento.

\subsection{Spezzettare i task grandi}

Un altro concetto utile: se si fa molta fatica a iniziare i task grossi, perché
sembrano infiniti (``più tardi inizio, più tardi inizia questa sofferenza''),
\textbf{vanno spezzettati}. Un task grosso non è quasi mai un task grosso: è
spesso un insieme di task piccoli. Si comincia a farne dei pezzi e si va avanti
piano piano, spuntandoli. Aiuta tantissimo anche nel \emph{problem solving}: un
problema gigante va scomposto in problemi più piccoli, da risolvere per primi, per
poi ricomporre il puzzle.

\subsection{Accettare di dire di no}

Tornando alla matrice di Eisenhower: bisogna \textbf{accettare di dover dire di
no}, perché dire sempre sì porta a problemi ben peggiori di quanto si pensi.

\paragraph{Il collega che diceva sempre sì.} A Riot il relatore aveva un collega
estremamente disponibile, amichevole e collaborativo: qualsiasi cosa gli si
chiedesse, rispondeva ``certo, la faccio io'' (``mi porti fuori il cane?'' ---
``certo''). Il problema è che era umano: a un certo punto si supera la quantità di
cose umanamente possibili in una giornata. In più tendeva a offrirsi volontario
per far vedere che era intraprendente e aveva iniziativa, e nei progetti che
coordinava il relatore doveva spesso rifiutare la sua candidatura perché era
sicuro che non ce l'avrebbe fatta. Il suo dire sempre sì, a un certo punto, lo ha
danneggiato: era quello che dice sempre che lo fa lui, ma non è detto che lo
faccia. ``Ce la fai?'' --- ``Sì, sì, assolutamente'' --- e poi ``scusa, non ce l'ho
fatta''. ``Ok, però ora hai messo nei guai me e il resto del team.'' Quella
disponibilità era più dannosa di un ``no, non posso''. È lo stesso discorso del
giorno prima: dire all'ultimo momento che non ce la si fa è peggio che dirlo a
inizio settimana.

Uno studente lo paragona alla \emph{sindrome dello yes man}. Non proprio: lo
\emph{yes man} dà sempre ragione all'altro perché è avverso al conflitto e non vuole
la discussione del ``spiegami perché sei in disaccordo''. Il collega, invece,
pensava davvero di poter fare molto più di quanto potesse, e rifiutava l'idea di
non riuscirci.

\paragraph{``Ci hai pensato bene?''} Vale per i progetti e per le deadline. Con i
junior soprattutto, quando a ``vuoi fare questa cosa?'' arriva un sì troppo
convinto, il relatore chiede sempre: ``ci hai pensato bene? Me lo stai dicendo
per farmi contento? Sappi che se mi dici di no non penso peggio di te; penso
peggio di te se mi dici di sì e poi non ce la fai. Vuoi rivalutare la risposta, e
magari dirmi \emph{non lo so} invece di un sì categorico?''.

\paragraph{Che cos'è una sconfitta.} Uno studente chiede se davanti alle
``sconfitte'' non si dovrebbe rimediare. Il relatore chiarisce che per lui la
sconfitta non è il task non fatto: quella è una sconfitta \emph{a breve termine},
si troverà un modo di farlo, è un danno ``ma chi se ne frega''. La vera sconfitta,
un danno \emph{a lungo termine}, è aver pensato di riuscirci e non esserci
riusciti: se non si risolve il motivo per cui si è detto sì quando era no, lo si
ripeterà altre cinque, dieci, quindici volte. Nella testa del collega di Riot c'era
sempre una scusante (``c'erano questi imprevisti che mi hanno bloccato''), come
quella con cui ci si giustifica le interruzioni. Il relatore gli rispondeva: ``se
ogni volta c'è un imprevisto, forse c'è un \emph{pattern}: il problema non è che
non riesci a fare le cose, è che non prevedi mai gli imprevisti''. Dopo sette
ritardi su dieci dovrebbe scattare l'ingranaggio mentale (``forse sbaglio io a
dire sempre di sì''), ma non è sempre facile farlo scattare, e a volte serve una
spinta. È una delle situazioni più frequenti. E la cosa più difficile è che,
quando glielo si fa notare, queste persone lo prendono come un ``non sei capace''.
Il messaggio è invece: \textbf{``sei bravo, ma sei umano anche tu''}. Non si può
fare un numero infinito di cose in un'ora: bisogna accettare i propri limiti ``dati
dalla carne'' e non pensare di essere un cyborg.

\paragraph{Il cliente non ha sempre ragione.} Uno studente osserva che vale anche
per i project manager e per il principio ``il cliente ha sempre ragione''. Anzi:
il relatore ha lavorato con molti clienti e a volte \textbf{il cliente non sa che
cosa vuole}, che è molto diverso dall'avere sempre ragione. Soprattutto con i
clienti che chiedevano pitch creativi, la sua squadra presentava sì un'opzione
con la loro richiesta, ma anche le proprie (``se vuoi raggiungere questo, sarebbe
meglio fare così''), e molto spesso il cliente sceglieva quella proposta dal team.
Spesso il cliente ha in mente \emph{una soluzione}, ma bisogna capire \emph{quale
obiettivo} ha in mente e poi dargli una soluzione: non vuole la soluzione, vuole il
\emph{goal}, e se c'è un modo migliore per arrivarci, forse è meglio un'altra
soluzione.

\subsection{Dire di no ai meeting}

La prima cosa a cui il relatore ha imparato a dire di no sono i \textbf{meeting}.
Ha passato anni in ambienti in cui il suo calendario era letteralmente solo
meeting dalla mattina alla sera: estenuante per lui e inutile per il team. Tra i
fattori c'è il \textbf{FOMO}: tanti manager vogliono essere in tutti i meeting per
non perdersi niente ed essere sempre informati su quello che fa il team. In realtà
vanno in \emph{burnout}, passano i pomeriggi nei meeting e finiscono per essere i
più distratti, quelli che non ascoltano e che fanno domande inutili solo per far
sentire che fanno domande, senza aiutare davvero.

Prima di accettare un invito il relatore si fa una serie di domande, \textbf{come
una checklist} (e le fa anche agli altri):
\begin{enumerate}
  \item \textbf{Sono davvero necessario?} Servo davvero io, o se mancassi non
        succederebbe niente? Ai meeting ``di cortesia'' non va: se pensa di
        essere necessario lo dice lui, e il fatto che si faccia un meeting
        importante e urgente senza di lui non danneggia il suo ego. ``Non
        aggiungermi solo perché pensi che serva che io venga aggiunto.''
  \item \textbf{Uno di noi potrebbe partecipare e informare il resto del team?}
        C'è già qualcuno del suo team? Se sì, non può andare solo lui? Per un
        meeting informativo (una nuova policy, una nuova procedura da spiegare a
        tutti i dipartimenti), andarci tutti significa investire tre ore a
        testa; se ci va uno solo e poi spiega agli altri, si risparmia tempo.
  \item \textbf{Il riassunto sarà reso disponibile?} Per il relatore è
        obbligatorio. Per lo stesso motivo per cui la lezione viene
        registrata: se ci si perde un pezzo lo si può rivedere, se qualcosa non è
        chiaro lo si può risentire. Ma c'è anche un motivo di \emph{office
        politics}: se in un meeting si chiede ``mi fai questa cosa per venerdì?''
        e la risposta è sì, non si vuole sentire venerdì ``intendevo venerdì
        prossimo'' o ``non ero sicuro''. Nelle note del meeting c'è scritto che hai
        detto sì. ``Lo so che è un po' paraculo, ma in realtà è così che si
        lavora'': bisogna dare più peso alla parola data, e quindi tenerne
        traccia. Finché tutto resta orale, niente rimane scritto e tutto è
        interpretabile. Oggi è anche facile: Zoom, con l'assistente IA, produce
        in automatico trascrizione e sottotitoli, senza bisogno di qualcuno
        addetto alle \emph{meeting minutes}. In ogni caso, se in un meeting viene
        presa una decisione, il relatore vuole che sia \textbf{ripetuta per
        iscritto}: non deve ipotizzare, deve poter andare a rileggere.
  \item \textbf{Il meeting aggiunge valore rispetto al documento?} Se arriva un
        invito con allegato un documento di venti pagine, il relatore chiede: ``il
        meeting consiste nel leggermi il documento? Perché altrimenti non
        vengo''. Se non c'è un valore aggiunto, se lo legge da solo e manda
        conferma di averlo letto. È lo stesso motivo per cui si sconsigliano le
        presentazioni in cui si leggono le slide; e infatti, osserva, in quel
        momento lui stesso stava leggendo una slide, ma aggiungendo informazioni.
        ``Se tutta la lezione fosse io che vi leggo le slide, non servirei io: vi
        manderei le slide.''
  \item \textbf{Per che cos'altro potrei usare questo tempo?} È forse la più
        importante. Anche ammesso che serva, che serva proprio lui e che il
        meeting sia importante, in quel momento dovrebbe fare qualcosa di più
        importante? Se sì, fa quella. Non bisogna avere paura di dire ``no, non
        vengo'' anche se si serve, oppure ``ho un'altra cosa importante,
        spostiamo il meeting''. Il relatore si è sforzato di essere più
        \textbf{assertivo}: ``capisco la richiesta, ma ho qualcosa di più urgente
        da fare in questo momento''.
\end{enumerate}

\paragraph{E le richieste a voce?} Uno studente chiede se una richiesta fatta a
voce da un senior (``devi fare questo'') vada accettata così o fatta mandare per
mail o per messaggio. Accettarla a voce ci può stare. Ma se la richiesta è ``devi
farla in questo modo particolare e preciso'', è meglio farsi mandare il modo
preciso \textbf{per iscritto}: non si deve lasciare a interpretazione \emph{come}
fare quella cosa, per non ricadere nel ``ah, ma io non la volevo così''. Oppure si
conferma: ``sì, la faccio; va bene se la faccio così, così e così?''. Se dice sì,
si è a posto.

\paragraph{Non dovrebbero pensarci gli organizzatori?} Dalla chat: queste domande
non dovrebbero farsele già gli organizzatori dei meeting? Sì, sarebbe meglio che
tutti lo facessero, e il relatore stesso se le fa quando invita qualcuno: ``è
davvero necessario Marco? È davvero necessario Andrea?'', e solo se sì li invita.
Ma lavora in aziende anche di cinquecento dipendenti, dove non è detto che tutti
conoscano o applichino questa logica; per questo la applica per sé, sperando che
la applichino tutti. Ha sempre chiarito con tutti che da lui non arriveranno mai
inviti di cortesia o dettati dalla gerarchia (``deve esserci anche il mio capo''):
``devi esserci tu; se vuoi portarti il capo, portatelo, ma non lo invito solo
perché qui si vuole sempre il capo presente: sono sicuro che avrà di meglio da
fare''.

\paragraph{E i task?} Dire di no vale naturalmente anche per i task: si applica la
matrice di Eisenhower (sezione~\ref{sec:z2-eisenhower}) e si decide se fare,
delegare, pianificare o cancellare. Non c'è molto da aggiungere.

\subsection{Il multitasking non esiste}

``Questo potrà urtare tanti ego'', ma \textbf{il multitasking non esiste}. Si può
essere più bravi o più portati di altri a fare più cose contemporaneamente; la
realtà è che, concentrandosi su più cose, \textbf{la somma non farà mai cento}. Se
si lavora su due cose non si lavora al 50\% su una e al 50\% sull'altra, ma al 40\%
e al 40\%, e il 20\% va perso in ``costo di manutenzione'', cioè nella capacità di
seguirle entrambe. Il consiglio è dedicarsi al cento per cento a un task, poi
passare al successivo, poi al successivo, invece di farne cinque insieme: si può
avere la percezione di andare più veloci, ma non è così.

Uno studente cita una lettura di neuroscienze secondo cui il cervello è in grado
di fare una sola operazione alla volta. Il relatore conferma: smettiamola di
parlarne come di un pregio o di una \emph{skill}. È un modo che alcuni cervelli
trovano per svagarsi, l'equivalente di accendere un podcast: fare due cose
insieme vuol dire cercare di distrarsi alternandole, e farle entrambe peggio di
quanto si farebbero una alla volta. A chi preferisce fare tante cose insieme
suggerisce di provare a farne una alla volta: le farà più velocemente. Anche il
semplice passaggio da una cosa all'altra non è istantaneo: ci sono sempre una
decina di secondi in cui il cervello si chiede ``che cosa stavo facendo?''; per
cinquanta cambi, già questo spiega la perdita di velocità. È una questione di
\textbf{bandwidth mentale}, un concetto che tornerà nel pomeriggio: anche lo
\emph{switch} occupa banda.

\subsection{La fine della giornata}

A fine giornata si \textbf{rimette a fuoco la giornata}: si riguardano la task list
e la \emph{must do list} e si conferma di aver fatto tutto. È letteralmente il
controllo prima di uscire di casa: chiavi, porta chiusa, cellulare, portafoglio.
Al relatore è capitato di dimenticare un \emph{must do} rimasto in lista per tutto
il giorno, ed entra in gioco quella che chiama la \textbf{sindrome del
tirannosauro}\footnote{Il riferimento è alla celebre scena di \emph{Jurassic Park},
in cui si dice che il tirannosauro vede la preda solo se si muove.}: se una cosa è
lì, non è detto che la si veda finché non si muove. Averla in lista tutto il giorno
non vuol dire averla percepita come da fare.

Può anche capitare di non farcela perché quel giorno sono arrivate così tante cose
urgenti e importanti da rendere inutili le stime: la giornata no capita e non c'è
nessun problema. Ma se era davvero un \emph{must do}, \textbf{va fatto comunque}:
è lì che scattano la richiesta di straordinario o il ``devo farlo comunque''.
Attenzione: se si chiude la giornata con un \emph{must do} ancora aperto e si
decide ``vabbè, lo faccio domani'', può starci, ma allora \textbf{non doveva essere
in quella lista}, perché era rimandabile.

Al relatore questo controllo serve anche ``per quieto vivere'': se chiude ogni
giornata con la \emph{must do list} completa, sa che magari non ha fatto tutto
quello che c'era nella to-do list, ma i \emph{must do} sì. ``Il mio, oggi, l'ho
fatto'': può stare tranquillo senza continuare a chiedersi se ha fatto quello che
doveva.

% ---------------------------------------------------------------------
\section{Domande sul time management}
\label{sec:z2-qa-tempo}

\paragraph{C'è un numero medio ``sano'' di \emph{must do} quotidiani?} No, perché un
conto è avere dieci \emph{must do} come ``rispondere a dieci mail'', un conto è
averne uno solo come ``inventarmi un nuovo gioco''. Dipende dalla dimensione: se ne
gestiscono di più se sono minuscoli, di meno se sono giganteschi. Ci sono lavori
che per definizione hanno tantissimi \emph{must do} quotidiani e lavori che ne hanno
uno al giorno. Si imparerà da sé qual è il numero giusto.

\paragraph{C'è un modo per stimare quanto si riesce a fare in una giornata?} ``Lo so
che non è una bella risposta'', ma l'unico modo è la propria \textbf{storia}:
sapere quanto si riesce a fare nel quotidiano. È normale e fisiologico che le
prime stime siano sbagliate, e probabilmente eccessivamente ottimistiche;
l'importante è rendersene conto. È letteralmente la \emph{seniority}: abituarsi a
fare qualcosa insegna quanto ci si mette. Va benissimo metterci di più all'inizio,
ed è normale che il tempo cambi man mano che si migliora: ciò che oggi si stima in
un'ora, tra sei mesi lo si stimerà in mezz'ora. Chi ha lavorato con Scrum lo sa:
su un progetto nuovo tutti sbagliano la dimensione delle \emph{stories} all'inizio,
perché nessuno sa quanto ci vorrà, e al terzo \emph{sprint} consecutivo è più
chiaro. Le stime sono stime, altrimenti si chiamerebbero ``inserire quanto tempo ci
metterai''.

Se però si dà una stima e poi il lavoro si rivela il triplo, il problema non è la
stima: è che \textbf{bisognava chiedere più informazioni} prima di darla. La
risposta giusta non è la stima, ma ``per darti una stima accurata mi serve sapere
questo, questo e questo''. Il problema non è sbagliare la stima, è non darsi i mezzi
per azzeccarla.

\paragraph{Come si fa a far entrare questi metodi nella routine? È solo
determinazione?} È un misto. Il relatore risponde da persona con ADHD, con la
tendenza a distrarsi e a fare un mucchio di \emph{overthinking} (``come dicevo ieri,
anche ordinare la pizza è overthinking''). È un po' disciplina, ma soprattutto far
diventare qualcosa un'\textbf{abitudine}: lo sforzo per fare le checklist era molto
più alto un anno prima, oggi è un'abitudine. Di solito servono \textbf{circa tre
mesi} perché qualcosa diventi un'abitudine\footnote{Stima come riportata in lezione.
Gli studi più citati sulla formazione delle abitudini indicano tempi molto variabili
da persona a persona, con una mediana intorno ai due mesi.}: la teoria dice che chi
inizia oggi ad andare in palestra ogni giorno, dopo tre mesi non penserà più ``oggi
devo andare in palestra'', ma la palestra farà parte della sua giornata. Non è più
uno sforzo, ma ``è così che lavoro''.

\paragraph{Chi fa davvero time management?} Prima di passare al brainstorming il
relatore torna alla domanda iniziale: qualcuno, ora che sa che cosa intendeva,
cambierebbe risposta? Uno studente risponde che lo fa, ma senza checklist, perché
non ne sente il bisogno (non gli capita spesso di dimenticare le cose); la
\emph{must do list} invece la segue sempre, e le cose meno importanti capitano nelle
giornate in cui ne salta qualcuna. Va benissimo: non esiste un unico modo, ognuno
troverà il proprio. Il relatore porta il suo esempio perché sa che è servito a molte
persone con cui ha lavorato; con chi faceva fatica ad applicare un metodo (come
chiedeva prima uno studente) ne hanno trovati altri, perché ciò che funziona per uno
non è detto che funzioni per un altro. L'importante è che \textbf{ognuno trovi il
proprio metodo}, perché la soluzione non è non fare time management: ci si sentirà
più ``liberi'' improvvisando la giornata, ma non la si sta sfruttando a dovere.

E non si parla dei video sul \emph{mindset} in cui ci si alza alle cinque, si fa la
checklist e si fa colazione leggendo le mail: organizzarsi meglio la giornata
\textbf{dà più tempo libero}. E per il relatore il tempo libero delle persone è
fondamentale, perché è ciò che permette loro di lavorare meglio.

% ---------------------------------------------------------------------
\section{Una parentesi: tempo libero, ferie e benessere}
\label{sec:z2-benessere}

\subsection{Le tre condizioni del relatore}

Il relatore è un sostenitore della \textbf{settimana lavorativa di quattro giorni},
che ha applicato in varie aziende; dove non ci è riuscito, ha introdotto almeno
\textbf{venerdì liberi}, due volte al mese o una volta al mese, cioè giornate in
cui non si fa niente. Come le pause brevi sono importanti per essere creativi,
anche le pause lunghe, cioè giornate intere, lo sono per essere creativi ed
efficienti. È lo stesso motivo per cui, dopo la Games Week, dava una settimana
libera alle persone (sezione~\ref{sec:z1-fasi-team}): ha bisogno che siano
riposate, e quella settimana ``persa'' è in realtà un investimento per le
settimane successive. In chat qualcuno lo paragona al ``metodo finlandese'': è un
misto tra quello e il metodo del buonsenso.

Il relatore viene spesso chiamato a \textbf{cambiare le culture aziendali}, ed è un
po' il suo focus. Come un allenatore che accetta una squadra solo se può portare il
suo staff (il preparatore, il tecnico), lui chiede tre elementi:
\begin{enumerate}
  \item possibilmente \textbf{quattro giorni lavorativi}, o altrimenti giornate
        libere per dare spazio alle persone;
  \item \textbf{ferie illimitate}: nessun limite massimo, ma un \textbf{limite
        minimo garantito}. Ferie illimitate non devono voler dire poter scegliere
        di non prenderle: ``devi prendere almeno X giorni, fino a infinito'';
  \item la \textbf{feedback culture}.
\end{enumerate}
Senza questi tre elementi non garantisce il risultato: è un modo di gestire le
aspettative fin dall'inizio, come si diceva il giorno prima.

\paragraph{La feedback culture.} Il perché è stato spiegato nella prima giornata, ma
il relatore lo sottolinea: un ambiente in cui ci si può fidare del team vale
moltissimo. Quando va a fare lezione allo IULM, anche in settimana, non deve
preoccuparsi di che cosa stia facendo il team mentre lui non lavora (``oggi è
domenica e non ho questa preoccupazione a prescindere, ma anche se fosse mercoledì
non l'avrei''): sa che stanno facendo del loro meglio, e non gli interessa sapere
altro.

\paragraph{Ferie e permessi illimitati.} Nell'esperienza del relatore, sia a Riot sia
a Blizzard, avere ferie illimitate significava spesso \emph{prenderne meno} di quelle
obbligatorie: c'è una vocina che dice ``ti serve davvero un giorno di ferie oggi?
Tanto sono illimitate, te le prenderai quando ti serviranno'', e sapere di averne
illimitate spingeva a lavorare di più. Per lui era altrettanto dannoso, e per questo
ha introdotto un \textbf{minimo obbligatorio di 24 giorni} di ferie.

C'è stato però un effetto non previsto, e molto desiderato: in tutte le aziende in
cui ha introdotto ferie e permessi illimitati, \textbf{i giorni di malattia sono
calati del 70\%}\footnote{Dato riferito dal relatore sulla base della propria
esperienza, non verificabile.}. La gente si ammala di meno? O piuttosto il giorno in
cui ci si sveglia e si pensa ``oggi non ne ho voglia, mi do malato'' diventa il
giorno in cui si prende un permesso? Si smette di usare come malattia un giorno che
non era di malattia.

\paragraph{L'\emph{Hangover Day}.} Sia a Riot sia a PG Esports il relatore aveva
introdotto quello che a Riot, ``essendo Riot a Dublino'', si chiamava
\textbf{Hangover Day}\footnote{Nella trascrizione ``Angover Day''.}: una giornata
all'anno in cui si aveva il diritto di telefonare alle nove del mattino e dire ``oggi
non vengo perché non ne ho voglia'' (tradotto: perché ieri sera ho bevuto troppo).
Altrimenti quella giornata sarebbe diventata un giorno di malattia o di ferie. Il
fatto è che \textbf{quasi nessuno lo usava}: il relatore stesso a volte ci ha
pensato, ma poi la mattina si chiedeva ``voglio davvero giocarmelo oggi? Tutto
sommato potrei anche andare'', e ci andava. Ma anche solo sapere di averlo lo faceva
stare tranquillo. Allo stesso modo, con le ferie illimitate, chi voleva prendersi un
mese lo faceva (``buon per loro''), chi non lo faceva andava bene lo stesso:
l'importante era che staccassero. C'era un \textbf{benessere più diffuso} nel team,
perché sapevano che, volendo, avrebbero potuto.

\paragraph{Il \emph{Neko Day}.} Il primo punto, i giorni liberi, era chiamato
\textbf{Neko Day}, da \emph{neko}, ``gatto'' in giapponese: il giorno in cui si può
fare il gatto, rannicchiati davanti al camino a non fare niente. Perché
introdurlo? Chi lavora dal lunedì al venerdì passa di solito il sabato o la domenica
a fare le commissioni che non può fare nei giorni lavorativi: la spesa, le bollette,
il bollo dell'auto. In questo modo sfrutta meno il weekend. Con un venerdì libero al
mese (o due, o tre) le persone facevano quel giorno le commissioni del weekend, e
sabato e domenica si riposavano davvero (``anche non fare niente andava bene,
chiaramente''). Il lunedì arrivavano molto più cariche e riposate. Questo ha aiutato
tantissimo a mantenere un team sano, fisicamente e mentalmente.

\paragraph{Il supporto psicologico.} A PG Esports il relatore ha introdotto anche il
\textbf{supporto psicoterapeutico offerto dall'azienda}: sei sedute da un'ora con
uno psicoterapeuta, online. Era il periodo in cui si parlava del bonus psicologo
statale, ``che poi non è mai arrivato''\footnote{Così in lezione. In realtà il bonus
psicologo è stato istituito in Italia nel 2022, ma con fondi limitati rispetto alle
richieste e con forti ritardi nell'erogazione; il relatore si riferisce
probabilmente a questo.}, e lui ha pensato: l'idea è buona, perché non introdurla a
livello aziendale? Erano gli anni dei lockdown: l'azienda stava perdendo molte persone
che andavano in depressione o non avevano più voglia di lavorare, perché avevano
difficoltà a vivere quel periodo (``io stesso non è che l'abbia vissuto benissimo'').
Qualcuno ne ha usufruito e gli ha detto: ``mi stavo per dimettere, poi ho risolto le
cose e ho preferito restare''.

Secondo il relatore oggi c'è una consapevolezza maggiore da parte delle aziende
sulla \textbf{salute mentale}: vuol dire dare più ferie, dare questo tipo di
supporto, dare più \emph{benefit} (anche i buoni pasto). Il benessere non è solo
economico, è anche mentale, ed è una cosa positiva da sapere.

% ---------------------------------------------------------------------
\section{Brainstorming: obiettivo e problemi tipici}
\label{sec:z2-brainstorming}

\subsection{Perché il brainstorming riguarda tutti}

Secondo il relatore questa è la parte che servirà di più, sia nei gruppi e nei
progetti del percorso di studi (il gioco da tavolo, il gioco digitale) sia nelle
aziende. Rispetto al time management ha una pertinenza più immediata: tutti
hanno esempi di \textbf{brainstorming andati male}, da cui si è usciti esausti, con
la sensazione di non aver ottenuto niente, o di non aver parlato abbastanza (``ha
parlato solo quello''), o di non essere sicuri della scelta. Il brainstorming lascia
una sensazione molto chiara, a fine meeting, che qualcosa non torna; il time
management dà piuttosto la tristezza del ``non ce l'ho fatta a fare quello che
dovevo oggi, come mai?''.

\subsection{Generare idee, non decidere}

Il primo chiarimento riguarda \textbf{l'obiettivo}, che il relatore vede spesso
travisato. L'obiettivo del brainstorming è \textbf{generare idee}; invece diventa
spesso decisionale, con lo scopo di scegliere che cosa fare. Ma allora non sarebbe
un brainstorming, sarebbe un \emph{decision making meeting}. Generare molte idee e
scegliere la migliore nello stesso meeting non funziona né per generare idee né per
scegliere la migliore, perché sono \textbf{due metodi che si contrappongono}. Il
brainstorming va fatto quando si devono generare idee \emph{per poi} scegliere che
cosa fare: deve venire \textbf{prima} del \emph{decision making}, come primo passo
della fase di esplorazione di un nuovo progetto.

\subsection{``Non esistono cattive idee''?}

Bisogna evitare che i membri del team \textbf{autofiltrino} le proprie idee per paura
di essere rifiutati o presi alla leggera. Da qui nasce la frase ``non esistono
cattive idee, buttate fuori tutto'', che è vera solo fino a un certo punto. Le
cattive idee esistono, oggettivamente; quello che non si deve fare è
\emph{colpevolizzarle}. Né si deve per forza perseguire ogni opzione: parlare di
opzioni che sono effettivamente brutte è uno spreco di tempo, e il brainstorming ha
anche un problema di gestione del tempo. Il relatore vede spesso brainstorming
gonfiati a ore e ore in cui si parla anche di cose che non meritano nemmeno dieci
minuti, come i feedback buoni ma inutili del giorno prima. \textbf{Trattare tutte le
idee allo stesso modo è sbagliato}, perché non hanno la stessa qualità. Bisogna
trovare un modo di valutarle senza che chi le ha proposte si senta preso di mira,
bullizzato o sminuito.

Per questo a ogni brainstorming dovrebbe seguire un \textbf{follow-up}, il momento in
cui si decide che cosa fare: le due parti vanno separate, altrimenti ``si crea solo
casino''. Il vero lavoro arriva quando le idee vengono messe a terra: c'è la parte di
brainstorming e di \emph{theorycrafting}, e la parte esecutiva viene dopo.

\subsection{I problemi più comuni}

\begin{description}
  \item[Le prime idee monopolizzano la sessione] Le prime due idee ``sparate a
        caso'' sono quelle su cui ci si concentra per la maggior parte del tempo,
        sviscerandone i dettagli (``intendevi così o così?''), per poi chiudere con
        ``altre idee? Questa, questa e questa? Sì, no, ok, andiamo avanti''. Si
        parla troppo delle prime idee e poco delle ultime.
  \item[Poche persone parlano per la maggior parte del tempo] È un effetto
        collegato al precedente.
  \item[\emph{Anchoring}] Ci si lega alle prime idee, e spesso si votano o si
        scelgono non perché siano le migliori, ma perché sembra di conoscerle meglio
        di quelle menzionate in trenta secondi alla fine. Non c'è nessun motivo per
        pensare che un'idea detta alla fine sia peggiore: era solo arrivata dopo.
        C'è un problema di \emph{timing}. Anzi, un'idea venuta istantaneamente è
        forse quella meno ragionata, perché ci si è pensato meno tempo.
  \item[Silenzi imbarazzanti] Spesso le persone arrivano al meeting senza nemmeno
        sapere che è un brainstorming. Quando il relatore convoca un brainstorming,
        comunica prima a tutti il \emph{topic}, il tipo di idee cercate, i paletti e
        le restrizioni, così che anche chi è meno portato all'improvvisazione arrivi
        con qualche idea pronta. Se la gente arriva impreparata rischia di non dire
        niente per un'ora, e allora perché è stata invitata? Non è un difetto loro:
        tanti manager dicevano al relatore ``ci sono persone inadatte ai
        brainstorming perché non sanno improvvisare''. No: si sta facendo
        brainstorming con un metodo che premia solo chi sa improvvisare ed è
        abbastanza estroverso da parlare sul momento. Non significa che chi non
        parla non avrebbe potuto contribuire, ma che non gli è stato dato modo di
        farlo.
  \item[Estroversi e pensatori rapidi dominano] Le personalità estroverse e chi
        pensa velocemente tendono a dominare la conversazione, più o meno
        intenzionalmente, senza lasciare agli altri il tempo di contribuire. Non c'è
        uno spazio in cui un introverso possa dire ``veramente io volevo\dots'':
        ``no, no, stiamo parlando di questa bellissima idea, facciamo questa''.
\end{description}

% ---------------------------------------------------------------------
\section{Un metodo di brainstorming}
\label{sec:z2-metodo}

Il relatore propone una soluzione principale, declinata in vari modi, che tocca
tutti questi punti.

\subsection{Preparazione e ambiente inclusivo}

\paragraph{Arrivare con un obiettivo chiaro.} Bisogna arrivare al brainstorming con
chiaro e preciso che cosa si sta cercando, cioè \textbf{che cosa si sta cercando di
risolvere}, \textbf{perché} e soprattutto \textbf{con quali vincoli}. Si stanno solo
raccogliendo idee per un cliente? Si cercano idee fattibili, o idee innovative? Che
cosa si vuole produrre da questo meeting? Deve essere chiaro a tutti \emph{prima}
del meeting, non spiegato durante, perché altrimenti si ricade nel problema delle
persone non pronte a intervenire.

\paragraph{Non giudicare le idee degli altri.} Al relatore è capitato spessissimo di
sentire un'idea terribile da una persona e, mezz'ora dopo, l'idea perfetta dalla
stessa persona. Ma se all'idea terribile tutti rispondono ``che stupidaggine, per
favore stai zitto'', quella persona non dirà la successiva, perché non vuole
rivivere l'esperienza. Giudicando le idee degli altri, di fatto, li si zittisce.
Le idee migliori nascono quando si propone quella che si pensa sia l'idea perfetta,
magari non è buona e non viene scelta (``fa niente''), si ascoltano anche le altre e
poi la si \textbf{riformula} integrando i pareri sentiti nel frattempo. Sono queste
\textbf{idee combinate} a essere spesso le soluzioni migliori.

\subsection{Separare generazione e discussione: \emph{brainwriting} e
\emph{rapid ideation}}

Con il classico ``non ci sono cattive idee, sparate'', succede spesso che uno prenda
la parola, spieghi la sua idea e il suo concetto, e gli altri comincino subito a
fargli domande (``ma quindi la vedi così?''): si sta unendo la persona che
\emph{genera} l'idea con la persona che la \emph{discute}, nello stesso momento.
Invece bisogna \textbf{dividere la generazione dalla discussione}: un momento in cui
si generano idee e un momento in cui si discutono. È il principio del
\textbf{brainwriting}. Funziona molto quando nel team c'è una figura con una
posizione gerarchica più alta: se quella persona esprime un'idea e i due momenti
sono uniti, tutti tenderanno ad allinearsi. E si vuole che \emph{tutti} possano
avere idee.

Lo strumento è la \textbf{rapid ideation}. Il relatore di solito dice: ``per i primi
cinque minuti non parliamo; scrivete su dei post-it il pitch delle vostre idee''. Se
ne possono scrivere quanti se ne vuole, dieci o uno, e si mettono da parte. Ma il
tempo è di \textbf{cinque minuti}: non dieci, non due, e soprattutto non mezz'ora.
Il \textbf{limite di tempo breve} è importantissimo, perché costringe,
consapevolmente o meno, a filtrare le idee mentalmente e a scrivere solo quelle che
si ritengono migliori: magari non è chiaro perché si preferisca un'idea all'altra,
ma il cervello sta facendo il lavoro in background e fa venire in mente le cose più
importanti o più giuste. È ottimo anche per i team che tendono a distrarsi, perché
sono solo cinque minuti.

\subsection{\emph{Dot voting}}

Poi si passa al \textbf{dot voting}:
\begin{enumerate}
  \item si attaccano \textbf{tutte} le idee su una lavagna fisica, un muro o una
        lavagna digitale (come Miro\footnote{Nella trascrizione ``mirror board'':
        si tratta della lavagna digitale Miro, già citata nella prima
        giornata.});
  \item si \textbf{uniscono le idee che hanno lo stesso concetto di base}. Se si sta
        pensando a un evento dal vivo e la proposta del relatore è un parco di
        dinosauri, mentre quella di un collega è un parco di dinosauri ``del
        paleolitico'' con i tirannosauri, le si raggruppa sotto ``parco di
        dinosauri'': non sono due idee separate. Un altro gruppo potrebbe essere,
        ad esempio, un parco a tema Avengers;
  \item ognuno ha \textbf{da tre a cinque voti} e va alla lavagna a mettere i propri
        pallini sulle idee che gli piacciono e di cui vorrebbe sapere di più.
\end{enumerate}
Il risultato tipico: un'idea con molti voti, una con un paio, altre con un numero
discreto.

\paragraph{Che cosa si vede subito?} Il relatore lo chiede alla classe. Si vede una
\textbf{preferenza} per un tema o per un'idea. Non c'è \textbf{polarizzazione} su una
o due idee iniziali, perché c'è varietà. Il voto è \textbf{anonimo}: non ci sono nomi.
E soprattutto manca un'altra cosa importantissima: \textbf{il giudizio negativo
espresso}. Se c'è un'idea con sette voti, una con cinque e una con tre, è
ragionevole parlare prima di quelle, perché tutti hanno espresso interesse e vogliono
saperne di più; ma nessuno sta dicendo ``non voglio assolutamente parlare di questa
perché mi fa schifo''. È semplicemente oggettivo che quell'idea non ha ricevuto voti.

\paragraph{Autovotarsi.} Il sistema permette anche di votare la propria idea, e non c'è
nessun problema: se la propria idea ha un pallino e gli altri ne mettono sette su
un'altra, il proprio voto è ininfluente. Il relatore vedeva spesso che all'inizio la
gente tendeva ad autovotarsi, poi capiva che era inutile, perché non sarebbe mai stato
il voto decisivo, e tanto valeva votare le idee che interessavano davvero. La
tendenza spariva con il tempo.

\paragraph{Che cosa si è ottenuto.} Tre ore di meeting sono state filtrate in una
discussione concentrata solo sulle idee che tutti ritengono interessanti, con un
\emph{agreement} di fondo: non si parla dell'idea numero uno perché è uscita per prima,
non è una questione di \emph{timing} ma di qualità. Il relatore ha vissuto tanti
meeting con la frustrazione interna del ``stiamo sprecando tempo su questa idea,
quando vorrei dire la mia che secondo me è meglio''. Qui non si può sprecare tempo,
perché si parla di qualcosa che tutti vogliono sentire.

\paragraph{E se si formano sottogruppi che si votano a vicenda?} Uno studente chiede se
non possano esistere ``cordate''. Anche se esistessero, significherebbe che il gruppo di
quel meeting vuole parlare di quella cosa, e va bene: l'importante è non sprecare tempo
su idee che non interessano a nessuno. Se si investono due ore, quarantacinque minuti
su ciascuna di tre idee vanno benissimo; dieci minuti per ogni idea no. E comunque il
voto è anonimo: non si può sapere che un'idea è di Gianni per dare supporto a Gianni,
quindi un eventuale sostegno reciproco sarebbe involontario.

\paragraph{E chi è lento a elaborare?} Uno studente osserva che il problema degli
introversi potrebbe ripresentarsi con i post-it: con un tempo così ristretto, chi è più
lento potrebbe andare ancora più in panico. Per questo bisogna \textbf{arrivare
preparati}: si conosce il topic in anticipo, quindi almeno un paio di idee già
formalizzate si hanno in mente. Magari non se ne scriveranno quindici, ma un paio sì.

\subsection{Quanti voti e quali varianti}

Dare più pallini a testa sarebbe controproducente? Il relatore indicava da tre a
cinque voti. Chiede alla classe quanti siano i gruppi di lavoro: dai quattro agli otto,
alcuni gruppi del secondo anno sono nove. La sua prassi:
\begin{itemize}
  \item \textbf{dai cinque partecipanti in su}, \textbf{tre voti} a testa, più che
        sufficienti per popolare molto la lavagna;
  \item \textbf{fino a cinque partecipanti}, \textbf{cinque voti} a testa, che
        permettono anche di variare un po'.
\end{itemize}
Esistono varianti: ad esempio tre voti di colori diversi, il blu per l'idea più
innovativa, il verde per quella più bella, il giallo per quella più facile, così da
qualificare \emph{perché} un'idea piace. Se l'obiettivo è proporre a un cliente tre
pitch (uno facile, uno medio, uno avanzato), la votazione è già fatta. All'inizio,
però, meglio partire da tre o cinque voti generici; con il tempo si imparerà ad
accelerare.

\subsection{Il moderatore e gli invitati}

\paragraph{Serve un moderatore?} Può essere utile. Il relatore fa spessissimo il
moderatore: partecipa ai brainstorming senza fare altro che moderare, ``perché so che
le persone con le idee fighe sono loro e non io''. Ci si può anche autogestire. Nei team
piccoli non ha senso ``perdere'' una persona per la moderazione; nei team più grandi
(dai cinque in su, sette-otto persone), se qualcuno dice ``non penso di avere idee
particolari su questo, vi modero così vi aiuto'', è un ottimo modo di contribuire.

\paragraph{Chi invitare?} Dipende da che cosa si cerca. Per inventare un nuovo gioco il
relatore invita i game director, gli art director e magari il \emph{publishing manager}
per un input commerciale; per decidere il \emph{combat system} invita magari il senior
game designer e il game director. La domanda che si fa chi convoca il meeting è: dato
l'obiettivo, chi è più ragionevole che possa contribuire?

Un 3D artist della classe chiede se allora gli artisti 3D non vengano inclusi. Dipende:
se l'obiettivo è definire la linea creativa dello stile dei personaggi, il relatore
invita anche il 3D artist, e non solo i \emph{concept artist}. Il 3D artist può dare un
input diverso dal concept: vota un'idea perché per lui è \emph{più fattibile}. Un
concept può essere bellissimo ma ``un inferno'' da modellare, e chi sa di doverlo fare
entro giugno non vota un'idea che sa già infattibile.

In realtà si può invitare chiunque: per un pitch su un'area dedicata al basket alla
Games Week, il relatore ha invitato un tecnico di regia che gioca a basket, per avere
anche la prospettiva di chi ha familiarità con il tema. A uno studente che porta
l'esempio di un elemento di interfaccia sul personaggio, il relatore risponde che,
volendo, si trova un motivo per invitare chiunque: l'importante è avere chiaro
\textbf{chi è necessario}. Il gruppo di lavoro può essere di otto e il brainstorming di
dodici, con input esterni a chi poi farà il lavoro. Basta non esagerare: se la maggior
parte dei voti arriva da chi non farà quel lavoro, si vota su qualcosa che poi non
faranno loro, e non ha più senso.

% ---------------------------------------------------------------------
\section{Dalla votazione alla decisione: la matrice impatto--sforzo}
\label{sec:z2-impatto}

\subsection{Massimo impatto, minimo sforzo}

Fatto questo, non si è finito: ora si sa quali idee interessano di più, ma non si è
ancora deciso che cosa fare. Serve la parte decisionale. L'obiettivo di qualsiasi
azienda (ed è il motivo per cui al relatore non importa se il programmatore fa il test
da casa cercando su Google) è \textbf{massimizzare l'impatto e minimizzare lo sforzo}.
Vale per i progetti, i task, il recruitment, per qualsiasi cosa: ``non perché sono
pigro, ma perché sarebbe stupido fare diversamente''.

Le idee selezionate si mettono quindi in un'altra matrice, con due assi:
\begin{itemize}
  \item sull'asse verticale il \textbf{valore}, cioè quanto \textbf{impatto} avrebbe
        l'idea (in alto l'impatto alto);
  \item sull'asse orizzontale lo \textbf{sforzo} (\emph{effort}), da basso (a sinistra)
        ad alto (a destra).
\end{itemize}
In alto a sinistra ci sono quindi le idee ad alto impatto e basso sforzo, in alto a
destra quelle ad alto impatto e alto sforzo, e così via.

\paragraph{L'esempio Adidas.} Se il cliente è Adidas, con la linea di abbigliamento per
il basket, un'idea a basso valore potrebbe essere una campagna su TikTok in cui gli
utenti condividono i loro momenti con un pallone da basket; un'idea ad alto valore
potrebbe essere fondare un parco divertimenti dove i giocatori NBA vengono a sfidarti
uno contro uno. Questo non dice che una sia meglio dell'altra: dice che l'impatto è
diverso. Allo stesso modo si valuta lo sforzo: la campagna social è a basso sforzo, il
parco divertimenti ad alto sforzo per molti motivi (i permessi di costruzione, i
giocatori NBA).

\subsection{Posizionare le idee, insieme e con precisione}

Il gruppo che a quel punto è \textbf{il gruppo decisionale}, cioè chi farà il lavoro,
deve posizionare sulla matrice le idee ritenute migliori, e deve essere \textbf{d'accordo
tutto il gruppo}: ``siamo tutti d'accordo che il parco dei dinosauri di Stefano e Andrea è
ad alto valore ma anche ad alto sforzo? Sì? Allora lo mettiamo qui''. E bisogna indicare
\textbf{un punto preciso}, non solo il quadrante: piazzare un'idea in un punto non è come
piazzarla poco più in là. Il risultato è una sorta di \textbf{mosaico} delle idee sul
grafico.

\subsection{Leggere la matrice in diagonale}

Davanti a una situazione tipica, con varie idee sparse, il relatore chiede alla classe:
se foste i project manager, se doveste fare voi queste cose, \emph{quali} fareste? Al
plurale, perché spesso da un brainstorming si esce con più idee realizzabili e si deve
scegliere su quali concentrarsi. Le risposte variano: c'è chi sceglie le due con più
impatto, chi aggiunge che dipende da tempi, budget e risorse, chi preferisce quella più
innovativa ma anche più rischiosa.

Il metodo che secondo il relatore va applicato è \textbf{procedere in diagonale}, dalle
idee a basso sforzo e alto impatto verso quelle ad alto sforzo e basso
impatto\footnote{Nella slide le idee erano rappresentate da post-it colorati (rossa,
rosa o salmone, azzurra, verde acqua, viola, blu, gialla), e le indicazioni del
relatore si riferiscono a quei colori. Dalla sola trascrizione non è possibile
ricostruire con certezza la posizione di ciascun post-it, e alcune indicazioni sui
colori risultano contraddittorie; qui è riportato il principio generale, che il
relatore chiarisce nelle risposte alle domande.}:
\begin{enumerate}
  \item un'idea a \textbf{basso sforzo e alto impatto} (nella slide, quella rossa) va
        fatta \textbf{per prima, sempre}: ``non dovrebbe essere nemmeno in discussione'';
  \item poi, a seconda delle risorse e del tempo, si possono fare
        \textbf{contemporaneamente} le idee che stanno sulla stessa diagonale
        successiva, che hanno cioè un rapporto simile tra impatto e sforzo;
  \item poi le successive;
  \item le idee ad \textbf{alto sforzo e basso impatto} (nella slide, la gialla)
        \textbf{non vanno nemmeno considerate}. Nessuno, infatti, l'aveva scelta.
\end{enumerate}
Per le scelte intermedie valgono due criteri, chiariti rispondendo a uno studente che
non capiva perché alcune idee fossero preferibili ad altre: \textbf{a parità di sforzo si
sceglie l'idea con più impatto}; \textbf{a parità di impatto si sceglie quella con meno
sforzo}. Con risorse infinite si farebbero anche le idee più costose; ma di solito le
risorse sono limitate, e allora tra due idee con lo stesso sforzo e un impatto molto
diverso la scelta è evidente. Se ci piacerebbe fare un'idea più ambiziosa ma abbiamo
risorse solo per una meno costosa sulla stessa diagonale, va bene lo stesso: abbiamo
comunque stabilito che è meglio delle idee sulle diagonali successive. Quanto fare dipende
dalla \textbf{bandwidth} a disposizione.

\paragraph{``La gialla era la mia idea: devo offendermi?''} No. Qualcuno ha detto che
l'autore fa schifo? No. Si è concordato che la sua idea, per quanto votata (``qualcuno ha
votato per la mia idea''), è piuttosto difficile da realizzare rispetto al suo impatto.
Non si sta dicendo che non sa più fare il suo lavoro. Anzi, se lui stesso dovesse
scegliere tra la propria idea gialla e la rossa di un altro, sceglierebbe la rossa: non è
nemmeno una questione di discussione. La matrice rende \textbf{visivamente evidente} che
cosa fare, e permette di aggirare molte discussioni che non sono più soggettive:
oggettivamente, dovremmo fare così.

\subsection{Il ruolo del moderatore e dell'esperto}

È qui che lavora il moderatore: ``non è più una discussione; visto il posizionamento,
dobbiamo fare queste cose''. Non si decide a maggioranza dove mettere le idee: se uno
posiziona un'idea e un altro non è d'accordo, difficilmente dirà ``la rossa dovrebbe stare
al posto della gialla''; dirà ``secondo me è un po' più a destra, un po' più in basso'', e
si fanno gli aggiustamenti. Se invece nascono discussioni del tipo ``la rossa dovrebbe
stare al posto della gialla'', vuol dire che non si è definita bene l'idea e non è chiaro
quanto sia davvero difficile da realizzare.

Oppure interviene \textbf{l'esperto}: è l'esempio del concept artist e del 3D artist. Il
concept artist dice che un'idea è a basso sforzo e alto impatto; il 3D artist risponde ``ne
condivido l'impatto, ma dovendola fare io ti assicuro che lo sforzo non è basso, è molto
alto''. Allora è un altro discorso, perché l'esperto sta dicendo che lo sforzo è diverso
da quello stimato. Ma alla fine tutto deve essere concordato.

\subsection{Chi decide}

\paragraph{Il gruppo decisionale può essere diverso da quello del brainstorming?} Sì: a
questo punto escono quelli invitati come ``consulenti'', che hanno dato idee ma non devono
poi realizzarle. Farli votare anche qui significherebbe perdere il collegamento tra chi
deve fare le cose e che cosa si decide di fare. In quel momento non hanno responsabilità:
hanno dato un input, e se decidessero che cosa devono fare gli altri dovrebbero esserne
i manager. Non dovrebbero poter dire agli altri che cosa fare se non ne hanno la
responsabilità. Vale per tutti i team, piccoli e grandi; secondo il relatore chi decide
dovrebbe essere \textbf{al massimo una decina di persone}. In un team di dieci che fa il
brainstorming tutto insieme, va benissimo che chi vota coincida con chi ha partecipato:
l'importante è che \textbf{chi vota sia chi farà il lavoro}.

\paragraph{Lo sforzo è il costo o la fatica?} Uno studente chiede se l'alto impatto
includa anche il costo. No, il costo sta nello sforzo, che per il relatore va
\textbf{ragionato in risorse}: denaro, tempo e così via. Un'idea può avere un basso costo
monetario ma un alto costo in tempo, e allora va bilanciata e si troverà a metà. È lì che
si spostano le idee e che avviene la conversazione: i tecnici dicono che lo sforzo tecnico
è un certo, poi arriva l'\emph{accountant} che dice che a livello economico è un altro, e
la si sposta.

\paragraph{E il piacere di fare il progetto?} In un gruppo di due o tre persone, osserva
uno studente, si finisce per votare anche per preferenza personale: quando si fa un gioco
bisogna avere anche il piacere di farlo. Ma \textbf{anche quello è impatto}: l'impatto
sulla motivazione. Come lo sforzo sono le risorse, l'impatto è l'impatto \emph{su di te,
sul team, sul mercato}, su ciò che si vuole. Il relatore lo ha tenuto volutamente
generico: può essere quanti soldi si fanno o quanto è utile per il proprio CV, e
l'impatto cambia tra un progetto accademico e uno professionale.

% ---------------------------------------------------------------------
\section{Consigli pratici e domande sul brainstorming}
\label{sec:z2-qa-brainstorming}

\subsection{Musica e cambio di scenario}

``Sembrano tutte stupidaggini, ma vi assicuro che aiutano.''

\paragraph{Musica di sottofondo.} È dimostrato che aiuta\footnote{Le affermazioni sugli
effetti della musica (chiave maggiore, rumore ambientale) sono riportate come dette in
lezione e non sono state verificate.}. Meglio una musica con \textbf{tempo e volume
costanti}: niente \emph{bass drop}, niente variazioni di tema o di ritmo. Va bene la
musica fatta di rumore ambientale a basso volume; è stato rilevato anche che le canzoni
in \textbf{chiave maggiore} aiutano la parte creativa a produrre idee più originali.
Meglio la \textbf{musica strumentale}, senza testo: per quanto ci si impegni a non
ascoltarlo, il cervello cercherà di capire le parole. È un effetto marginale e non
succede niente se capita la canzone con il testo, ma si mette in conto che ``l'uno per
cento del cervello'' sarà impegnato a capirlo, come quando in treno si ascolta, che lo si
voglia o no, la telefonata di qualcun altro. È \emph{min-maxing}.

\paragraph{Cambiare scenario.} A Riot c'era una stanza dedicata al brainstorming, pensata
per essere creativa: le finestre sul fiume, le porte a vetri colorati, tutto fatto per
ispirare creatività. Ma fare brainstorming sempre nella stessa stanza generava sempre lo
stesso tipo di idee: anche solo cambiare ambiente aiutava ad averne di diverse. Il
relatore spesso spostava il meeting, ad esempio nello Starbucks sotto l'ufficio. Lo
spazio fisico gioca un ruolo importante, ma anche quello virtuale: se lo si fa sempre su
Discord, si può provare con ognuno uno sfondo diverso e simpatico. Oppure: ha sempre
moderato la stessa persona? Si prova con un altro moderatore. Si cambiano i fattori
perché i risultati non siano sempre gli stessi.

Uno studente chiede se il cambio di scenario possa essere anche qualcosa come passarsi una
palla per avere la parola, o un altro espediente per tenere viva l'attenzione. Con questo
modello, che non richiede di prendere la parola, l'esempio della palla perde senso, ma come
esempio va benissimo: anche il comportamento fisico e il \emph{modus operandi} sono un ottimo
cambio di scenario. Come lo sono cambiare il moderatore o \textbf{cambiare l'obiettivo}: se
il cliente chiede X, gli si darà un'idea X, ma si può provare anche a votare per un'idea Y,
facendo più votazioni con vari criteri di valutazione.

\subsection{Domande}

\paragraph{L'idea scelta deve essere approvata da qualcun altro?} Se un manager dice ``fate
questo brainstorming senza di me'', sta \textbf{delegando la scelta}: ``mi fido di quello che
sceglierete''. Avendo preparato prima limiti, paletti e \emph{constraint}, si aspetta che la
soluzione portata sia in linea con essi; ma se non partecipa deve fidarsi della scelta. Se
poi la scelta debba passare anche da un livello superiore dipende dall'azienda e dagli
accordi precedenti.

\paragraph{Si può fare un brainstorming ibrido, tra presenza e remoto?} Uno studente fa
l'esempio di un gruppo in cui uno è a New York e gli altri a Milano. Sì, senza problemi: si
fa una call, il portatile punta verso il muro con le idee, o le si mostra fisicamente con
il cellulare. Chi è da remoto scrive i suoi post-it, ad esempio mandandoli su WhatsApp al
moderatore o a qualcuno presente che li scrive per lui, e vota dicendo ``mettimi un pallino
qui, qui e qui'' (in quel caso il voto non è anonimo, ma si può fare tranquillamente); se
deve parlare, parla.

\paragraph{E se i colleghi non riconoscono o non condividono il metodo?} Ci deve essere un
\textbf{agreement}: non si può arrivare e dire ``facciamo così'' se gli altri non vogliono.
Questo metodo funziona solo se tutti sono d'accordo a usarlo. Il relatore ne ha proposto
uno tra i cento metodi di brainstorming esistenti, un buon compromesso tra semplicità ed
efficienza (ce ne sono di più semplici e di molto più efficienti ma più complessi). Se il
gruppo non ne vede il valore, al massimo insisterebbe per \textbf{fare almeno una prova},
altrimenti diventa chiusura mentale; se poi per quel team non funziona, se ne cercano
altri.

\paragraph{E nei gruppi molto piccoli, dove diventa un duello?} Uno studente ha lavorato in
gruppi di tre o anche di due persone, dove il brainstorming diventa quasi un duello in cui
ognuno difende la propria idea. Bisogna \textbf{aspettare la discussione}: invece di ``scegliamo
tra la mia idea e la tua'', si cerca di concordare quale idea abbia l'impatto maggiore e
quella lo sforzo minore. E se non si è d'accordo, bisogna \textbf{portare esempi}: se uno
dice ``questa idea costa di più'', l'altro chiede ``in che senso? Che cosa costa di più?''. Se
non sa rispondere, non si può fare qualcosa senza sapere come va a finire: meglio rimandare
la discussione a quando si avranno le informazioni. A volte la domanda giusta è \emph{perché
non abbiamo queste informazioni?}: forse è \textbf{prematuro} decidere, perché non si hanno
ancora tutti i fattori.

Nei team piccoli non dovrebbe esserci un ``si fa come dici tu o come dico io'', una scelta a
pelle: si dovrebbero conoscere così bene le proprie proposte da poterle posizionare con
assoluta certezza sulla matrice, e avere anche la presenza di spirito di dire ``la tua è
chiaramente migliore, si fa la tua''. Si fa la cosa migliore per entrambi, che dovrebbe
essere oggettiva e non soggettiva. Non è facile arrivarci, ma deve essere l'obiettivo del
team.

% ---------------------------------------------------------------------
\section{Disagree and commit}
\label{sec:z2-disagree}

\subsection{Quando la scelta non ci convince}

Uno studente teme che chi vede scartata la propria idea perda motivazione, in un gruppo
piccolo come in uno grande. Non è un deragliamento: è una conseguenza potenziale di
qualsiasi scelta. Il relatore ha visto gente demotivata perfino dalla settimana di quattro
giorni: ci sarà sempre qualcuno più o meno motivato dalle scelte. Ma qual era l'alternativa?
Un brainstorming classico in cui tutti parlano, chi più chi meno, e si sceglie un'idea che
tutti pensano sia la migliore senza esserne davvero sicuri, perché non si è votato, né
espresso una preferenza, né discusso davvero le idee migliori.

Qui entra il concetto di teamwork del \textbf{disagree and commit}: ``non sono d'accordo con
questa decisione, non mi piace, non penso sia la migliore, ma la faccio''. Per due motivi:
\begin{enumerate}
  \item perché \textbf{è il mio lavoro}, e dovrebbe essere il punto principale;
  \item perché \textbf{voi ritenete che sia la scelta migliore}. Se si è la voce fuori dal coro
        e tutti hanno scelto un'altra strada, non si rema contro: è \emph{default to trust}.
        Se gli altri vi hanno visto più potenzialità, e (grande premessa) \textbf{hanno
        ascoltato} il proprio disaccordo e i suoi motivi, e hanno comunque scelto così, ci si
        fida di loro.
\end{enumerate}
Con una postilla: se le cose vanno male, \textbf{non si dà la colpa agli altri} e non si dice
``ve l'avevo detto'', come se non si facesse parte del team. Disagree and commit vuol dire ``mi
fido del team e facciamo questa cosa''; se poi si rivela sbagliata, \emph{il team} ha
sbagliato, e la volta dopo si cambierà il modo di scegliere. Altrimenti non è disagree and
commit, è ``\emph{disagree and ostracize}'': disaccordo più ostilità. Non bisogna cercare il
colpevole, perché si è deciso tutti insieme (sezione~\ref{sec:z1-default-trust}).

Se si pensa davvero che si stia sbagliando clamorosamente il posizionamento di un'idea, si
dice piuttosto ``non andiamo avanti finché non definiamo bene lo sforzo di questa cosa''; ma
bisogna saperlo giustificare, altrimenti è chi dissente a dover fare un passo indietro.

\subsection{La demotivazione fa parte del lavoro}

La demotivazione ci sarà sempre, perché non si sceglierà mai solo ciò che piace: per quanto
piaccia il proprio lavoro, ci saranno task sgradevoli e scelte che non spettano a noi. E
ci saranno scelte (``ve lo dico per esperienza personale'') che si è convinti siano
sbagliate, al punto di pensare ``piuttosto lascio il team'', ma che si finisce per accettare,
per poi rendersi conto che gli altri avevano ragione. Come c'è un \emph{bias} verso l'amico e
il nemico, c'è anche un \textbf{bias verso noi stessi}: crediamo di conoscere bene una cosa, e
se altre dieci persone dicono ``secondo me è fattibile'', bisognerebbe aprirsi alla
possibilità che ne sappiano più di noi.

``Non lo accetto ma lo faccio'' dovrebbe essere alla base del lavoro insieme, perché è
\emph{default to trust}: rifiutarsi (``non mi fido di voi, non lo faccio'') significa non
fidarsi del giudizio del proprio team. E dire ad alta voce ``non sono d'accordo, ma ti darò
comunque il massimo'' è la cosa che \textbf{fa guadagnare di più il rispetto degli altri}.
Non c'è niente di meglio che vedere qualcuno che aveva un'opinione diversa impegnarsi per
far vedere di avere torto, facendo andare bene la cosa. Per il relatore dimostra solo che è
una persona matura e importante; ha avuto persone che, passate da questa situazione, gli
hanno fatto pensare ``è molto più matura di quanto credessi'', e che ha valorizzato molto di
più in seguito, perché sa di potersi fidare di chi fa disagree and commit: è stato aperto
nella comunicazione, ha detto che cosa secondo lui non andava, ma non ha fatto niente per
remare contro, anzi ha dato il suo supporto. ``Non posso chiedere di più a una persona.''

\subsection{Chi arriva a lavori iniziati}

\paragraph{Se il team cambia in corso d'opera, chi subentra come si rapporta alle decisioni
già prese?} (per dimissioni, promozioni e così via). Deve \textbf{accettare lo status quo}. Al
relatore capita spesso di entrare in team con progetti già in corso: si fa dire che cosa è
stato scelto finora e va \textbf{in commit di default}: ``non ho preso parte a questa decisione,
quindi non sono in \emph{agreement}, ma ora faccio parte del team, quindi vado in commit''.

\paragraph{E se l'esperienza dice che una scelta ha uno sforzo molto diverso?} Ad esempio
che un'idea ``rossa'' ha in realtà lo sforzo di una ``viola''. Lo si fa presente, ma dipende dal
\textbf{tempismo}: se la scelta è appena stata fatta e c'è ancora modo di cambiarla, la si
segnala; se si è al 90\% del progetto, è controproducente dire ``fermi tutti, passo
indietro'', e a un certo punto bisogna rispettare la scelta del team (``avete scelto così,
andiamo avanti e vediamo come va''). E se si sa che, calcolando ferie e periodo estivo,
un'animazione prevista per giugno arriverà a settembre? È molto \emph{case by case}: a
seconda dell'apporto si può modificare il risultato o no. È come riaprire un processo se
emergono nuove prove: se dimostrano che l'imputato non è l'assassino, è giusto riaprirlo; se
dimostrano che l'assassino ha sparato con la sinistra invece che con la destra, conta fino a
un certo punto. È giusto segnalarlo, ma non sempre quello che si porta giustifica uno
``\emph{scratch that}, rifacciamo tutto'': diventa \textbf{un feedback, non una risoluzione}
(``state facendo così, questo è il mio feedback, poi vediamo come va'').

Con cinque minuti di anticipo si va in pausa pranzo, con l'accordo per il pomeriggio: circa
un'ora di presentazione più ``standard'' sul leading from within, poi domande e risposte,
``e quando finiscono le domande abbiamo finito''.

% ---------------------------------------------------------------------
\section{Leading from within: perché parlare di leadership}
\label{sec:z2-leadership}

\subsection{La ripresa del pomeriggio}

Alla ripresa, mentre si aspettano gli ultimi arrivati, il relatore chiede se durante la
pausa siano venute in mente domande. L'unica è logistica: le slide verranno date a fine
workshop? Sì: il relatore le passerà all'accademia (a Vanessa, Andrea ed
Enrico\footnote{Nella trascrizione l'accademia è chiamata ``Ive'' e ``Hive'', come nella
prima giornata; per i nomi dei referenti vedi la nota all'inizio del
capitolo~\ref{cap:z1-feedback}.}), che le farà avere alla classe. Gli appunti vanno
comunque presi, meglio su quello che dice che su quello che mostra.

L'ultima parte, dice, dovrebbe essere abbastanza veloce da presentare, ma vale la pena
fermarsi a parlarne, perché è un discorso più avanzato. È la logica chiusura di quanto
detto tra la prima giornata e la mattina: tirare le fila e mostrare come tutto faccia parte
di un unico filone di ragionamento e progettualità, con un unico scopo, un unico obiettivo e
un unico modo di ragionare.

La voce sembra a posto, ma il relatore è stanco dopo due giorni, come immagina lo sia la
classe, nello \emph{stretch} finale di due giornate molto piene. Sono cose che in condizioni
normali spiega al suo team nell'arco di settimane, con il tempo di metterle in pratica,
discuterle e approfondirle. Qui è tutto concentrato: servirà tempo per metterlo in atto e
fare le proprie prove, quindi non bisogna spaventarsi se non verrà tutto giusto la prima
volta.

\subsection{Il leader non è il più bravo}

Perché parlare di leadership in un workshop non rivolto solo ai leader? La risposta
completa arriverà più avanti, ma il punto di partenza è che il relatore \textbf{non crede
nei leader di team o di progetto che detengono da soli il sapere e il modo di fare le
cose}. Come detto il giorno prima, essere un leader, un manager o un capo non dovrebbe
voler dire essere il migliore in quel lavoro, ma essere \textbf{la persona più adatta a
ricoprire quel ruolo}: guidare un gruppo significa \textbf{occuparsi delle persone} al suo
interno.

L'obiettivo che il relatore si dà, e che secondo lui ogni manager dovrebbe avere in mente,
è che quando si torna dal lavoro ci si senta \textbf{ispirati, sicuri e realizzati}, e che
questo sia \textbf{un diritto naturale di tutti}, non un lusso riservato a chi ricopre certe
posizioni, a pochi fortunati o ``a qualcuno sì e a qualcuno no''. Ha notato,
specialmente in Italia, la tendenza a vedere tutto ciò che viene prima del diventare
manager come la \emph{gavetta} necessaria per arrivarci. Invece chiunque faccia parte di
un'organizzazione, che sia un'azienda di videogiochi o altro, dovrebbe sentirsi soddisfatto,
ispirato e sicuro in quello che fa; e non è corretto che un'azienda tratti bene le persone
``da manager in su''. Anche quando non lo dicono, a volte i fatti sono quelli: da manager si
viene trattati in un certo modo e si ha accesso a certe cose, altrimenti no. Non si tratta di
\textbf{appiattire la gerarchia}: si tratta di assicurarsi che, pur dando cose diverse ai
diversi livelli, a tutti i livelli venga riconosciuto uno strato comune di
\textbf{rispetto, fiducia e professionalità}.

\subsection{Le carriere divise: manager o \emph{subject matter expert}}

Uno studente chiede conferma: quindi un \emph{lead artist} non è necessariamente più bravo
di un senior a fare l'artista? Esatto. E qui il relatore apre una parentesi. In molte
aziende, e soprattutto in quelle in cui tende a lavorare lui, la motivazione principale per
diventare manager è avere una promozione e quindi un aumento, più che il desiderio di
prendersi cura del team e delle persone. Come si ovvia al problema? \textbf{Dividendo le
carriere}. Arrivati al livello senior, il relatore fa con ognuno dei suoi senior
\textbf{``the Talk''}, per esporre le due opzioni che hanno davanti:
\begin{description}
  \item[Busta A: diventare manager] È il percorso standard: si valuta se si è adatti a fare
        il leader e a guidare un team, e si danno gli strumenti per farlo, perché
        ``nessuno nasce leader imparato'': prima di provarlo sulla propria pelle nessuno sa
        che cosa voglia dire. Per quanto si possa essere carismatici, fare il leader è
        un'altra cosa, e non vuol dire solo essere carismatici. Ci sono formazione,
        strumenti e un percorso di crescita.
  \item[Busta B: diventare \emph{subject matter expert}] Una persona che va oltre il senior
        nella stessa disciplina: diventa \textbf{l'esperto dell'intera azienda} su quella
        disciplina, perché la sua priorità non è aiutare gli altri ma essere il migliore.
        ``Non c'è assolutamente nulla di male.'' Il suo ruolo è essere il più specializzato
        in quel campo, il punto di riferimento da cui tutti vanno a farsi spiegare le cose.
\end{description}
Il concetto è abbastanza nuovo, e il punto chiave è che l'azienda \textbf{paga il
\emph{subject matter expert} e il leader esattamente allo stesso modo}: non c'è un guadagno
economico nello scegliere di fare il manager invece dello specialista, e il relatore valuta
quel tipo di conoscenza quanto il saper gestire persone. Così chi vuole essere il migliore,
o anche solo un aumento, può continuare a fare ciò che sa fare e avere l'aumento; chi vuole
l'aumento \emph{e} prendersi cura delle persone sceglie l'altra strada. Così si fa una scrematura sul
\emph{motivo} per cui si vuole fare il manager. Sembra una scelta piccola, ma è molto
importante capire se ci si vuole occupare di se stessi o degli altri. Non è un giudizio
(``non sto dicendo che sia meglio l'uno o l'altro''): le due cose \textbf{non convivono}, e
bisogna sapere che cosa si vuole fare, sapendo che l'azienda tratterà allo stesso modo le
due scelte.

\paragraph{Si può essere entrambi?} Si può scegliere ad esempio il percorso da lead artist
e voler comunque continuare a migliorare le proprie capacità. Ma, come nelle simulazioni di
conversazione del giorno prima, il relatore chiederebbe: ``se avessi le risorse per fare solo
una delle due cose, quale sceglieresti?''. Perché non le si avrà: non ci sarà il tempo e il
modo di fare bene entrambe. Il suo ``cinismo'' arriva a dire che bisogna scegliere, perché
tenersi buone entrambe le strade porta a farle male tutte e due. Preferisce un ottimo
\emph{people manager} o un ottimo specialista; le figure intermedie, in questo campo,
rischiano di diventare pericolose per sé e per il team.

\paragraph{Specializzazioni verticali.} La strada dello specialista permette
specializzazioni molto verticali: in chat qualcuno cita il \emph{mecha designer}, ma può
esserci anche lo specialista di animazioni delle espressioni facciali. Naturalmente l'azienda
deve avere bisogno di quella competenza (non ci si specializza in cose che non servono), e non
è detto che ne serva uno solo. È un percorso di carriera parallelo e alternativo a quello del
manager ``standard''.

% ---------------------------------------------------------------------
\section{Il ``momento Piero Angela'': la biologia del gruppo}
\label{sec:z2-biologia}

\subsection{Perché una digressione scientifica}

Il relatore annuncia un ``piccolo excursus biologico'', fondamentale per capire dove vuole
arrivare, e prova a mettere un sottofondo musicale adatto: la condivisione dell'audio non
funziona, quindi manda il link in chat perché ognuno lo faccia partire per conto suo. Un
studente lo battezza subito ``momento Piero Angela''. Il relatore cerca di far combaciare le
motivazioni aziendali del suo metodo di leadership e di progettualità con \textbf{basi
scientifiche}: in passato ha collaborato con varie università che si occupavano di
neuromarketing e neurobiologia, per avere una conferma o una smentita del fatto che ciò che
ipotizzava ``non fosse una cavolata colossale''. Semplificando moltissimo, si guardano solo
alcune sostanze che si traducono direttamente nel modo in cui si vive l'essere in un gruppo,
in un team, in un'azienda\footnote{L'impianto di questa parte (endorfine, dopamina,
serotonina, ossitocina e cortisolo, il villaggio protetto dai pericoli esterni, il ``cerchio
di sicurezza'') ricalca da vicino quello del libro di Simon Sinek \emph{Leaders Eat Last}
(2014), che il relatore cita poco più avanti. Le descrizioni biologiche sono molto
semplificate, come avverte lo stesso relatore, e vanno lette come tali.}.

\subsection{Endorfine}

Le \textbf{endorfine} sono rilasciate per \textbf{mascherare il dolore}. Lo sa chi corre, chi
fa pesi o ginnastica, ma anche chi ha riso così tanto per una battuta da avere male ai
fianchi: è quella sensazione di dolore in cui però si è contenti di provarlo. Il \emph{runner's
high} è l'esempio più calzante: essere distrutti dopo una corsa ma soddisfatti. È la risposta
biologica allo sforzarsi in cose che sono per il nostro bene, un modo con cui il corpo ci
``convince'' a farle.

\subsection{Dopamina}

La \textbf{dopamina} è più familiare: è la \textbf{ricompensa per un'azione}. Dalle cose più
banali (si mangia e si riceve dopamina) a quelle viste in mattinata: anche completare una
checklist rilascia dopamina (sezione~\ref{sec:z2-pianificare}). Viene rilasciata quando si
vede l'obiettivo, come il traguardo durante una corsa; è la stessa cosa che accadeva ai
cacciatori del villaggio in tempi antichi, che rilasciavano dopamina quando vedevano le tracce
e venivano così incentivati a seguirle. Sono concetti fondamentalmente ``animali'', ma si
applicano tuttora anche in azienda.

\paragraph{Una nota di colore: la dipendenza.} La dopamina \textbf{crea dipendenza}: molti casi
di ossessione per alcol, nicotina, gioco d'azzardo e social media sono problemi legati alla
dipendenza da dopamina. Secondo il relatore non è un caso che le diagnosi di ADD e ADHD siano
aumentate del 60\% negli ultimi dieci anni\footnote{Nella trascrizione ``EDD e HD''. Il dato
e il nesso causale con la maggiore disponibilità di alcol, nicotina, gioco d'azzardo e social
media sono riportati come detti in lezione e non sono stati verificati; l'aumento delle
diagnosi viene in genere attribuito anche ad altri fattori, come la maggiore consapevolezza e
il cambiamento dei criteri diagnostici.}, perché buona parte di quegli stimoli è diventata
ampiamente disponibile proprio in quel periodo. Anche la dipendenza dal lavoro, a volte, ha la
stessa origine: si vuole il rispetto che deriva dal lavoro, il \emph{rush} che deriva
dall'averlo.

\subsection{Il villaggio di diecimila anni fa}

Si fa poi un salto di diecimila anni, in un villaggio. Gli esseri umani, per quanto più o meno
introversi o estroversi, sono \textbf{animali sociali} che tendono a vivere esperienze di
gruppo; le aziende nascono da un gruppo di persone che si danno un obiettivo condiviso, come
in antichità un gruppo di persone che voleva sopravvivere insieme. Quello che non è cambiato è
che \textbf{fuori dal gruppo ci sono dei pericoli}: tigri, lupi, draghi (``do per scontato
che nell'antichità ci fossero i draghi, perché D\&D mi insegna così''). Per evitarli, i gruppi
tendono a riunirsi e a creare difese intorno a sé: da qui le prime difese dei primi villaggi,
per isolarsi dai pericoli esterni. ``Vi giuro che avrà tutto senso.''

\subsection{Serotonina}

La \textbf{serotonina} viene rilasciata anche \emph{indirettamente}: non deve per forza
riguardare noi. Quando si assiste a una cerimonia o celebrazione pubblica che include una
persona cara, si rilascia serotonina come se si fosse noi l'oggetto del riconoscimento. È uno
dei motivi per cui si fanno diplomi e lauree ``all'americana'', con la consegna sul palco
davanti a tutti: gli studi avrebbero mostrato che anche i genitori in platea rilasciano
serotonina, come chi riceve il premio. Lo stesso vale per parenti, amici, matrimoni: sono
esperienze condivise, e si tende a condividere in pubblico riconoscimenti e successi perché in
tutti i partecipanti viene rilasciata serotonina.

Purtroppo ha a che fare anche con i beni materiali: la tendenza all'\emph{hoarding}, cioè a
collezionare eccessivamente qualcosa, nasce dal fatto che quell'attività rilascia serotonina,
ingannando la mente nel credere di aver ottenuto un successo.

È chiamata anche \textbf{la sostanza della leadership}, perché aiuta a riconoscere in qualcun
altro chi ha valore per il gruppo. Nel villaggio, quando arriva un pericolo, c'è sempre un
gruppo di persone che dice ``ce ne occupiamo noi'': il gruppo sa che si stanno prendendo un
rischio, e la serotonina è il modo per riconoscere che meritano rispetto. Il relatore consiglia
il libro di Simon Sinek \emph{Leaders Eat Last}\footnote{Nella trascrizione ``Leadership
Lust''. Il titolo è \emph{Leaders Eat Last: Why Some Teams Pull Together and Others Don't}
(2014).}: i leader mangiano per ultimi,
all'opposto di quanto accadeva in antichità, quando al ritorno dalla caccia mangiavano per
primi i cacciatori, che dovevano essere in forze per la caccia successiva, poi donne e
bambini, poi gli altri. C'è un \textbf{riconoscimento sociale}: tu fai qualcosa per proteggere
il gruppo, noi ti riconosciamo qualcos'altro in cambio.

Le aziende sono strutturate in modo molto diverso? È una domanda retorica: quando si dice che
serve un capo, un manager, un rappresentante, si sta dicendo che quella persona dovrebbe avere
l'interesse di proteggere il gruppo e di prendersi cura di tutti. Per questo sapere che il
proprio capo guadagna di più non fa particolarmente arrabbiare: lo si vede come quasi naturale,
perché più responsabilità si ha verso il gruppo più si guadagna, dato che quello dovrebbe essere
il suo scopo. Ed è anche il motivo per cui il relatore si arrabbia molto quando legge di CEO o
\emph{executive} che, a fronte di licenziamenti di migliaia di persone, si assegnano dei bonus.
È una rabbia non del tutto razionale, che sente da dentro: è come se si stesse violando un
\textbf{contratto sociale}. Il contratto è ``tu ti prendi dei rischi per me e mi difendi, io ti
riconosco più cibo, più soldi''; quando si abusa di quella posizione per un beneficio personale
a scapito di chi si dovrebbe proteggere, tutti lo vedono come un'ingiustizia, perché si viene
meno al contratto stabilito diecimila anni fa con la nostra biologia.

\subsection{Ossitocina}

L'\textbf{ossitocina}\footnote{Nella trascrizione ``ossidocina''.} è la \textbf{sensazione di
sicurezza} che si genera quando si è vicini a qualcuno. Non è un caso che tra i legami di
affinità e collaborazione più forti ci siano quelli dei team che si uniscono contro un
\textbf{nemico comune}: i team che superano insieme avversità e difficoltà ne escono molto
legati. Tra le persone che il relatore sente più spesso in questo periodo ci sono quelle con cui
ha condiviso momenti difficili, un cattivo manager, una brutta esperienza lavorativa: sa che
erano lì con lui. Vale anche per i veterani di guerra o per i sopravvissuti a una stessa
tragedia, che sviluppano un rapporto molto forte (se ne è parlato molto negli Stati Uniti a
proposito dei veterani), perché hanno rilasciato immense dosi di ossitocina. È anche ciò che fa
accettare subito il figlio alla madre quando lo prende in braccio dopo il parto, ed è il motivo
per cui, a migliaia di anni di distanza, si stringono ancora accordi con una stretta di mano o
un abbraccio: c'è sempre qualcosa di fisico, generato dal bisogno di ossitocina.

Viene generata anche quando si compiono o si \textbf{assiste ad atti di generosità} senza
aspettarsi nulla in cambio. Se il relatore dicesse di aver donato cinquanta euro a
un'associazione di beneficenza, la reazione sarebbe ``ok, bene''; se dicesse che ogni sabato
pomeriggio fa volontariato, lo si vedrebbe come qualcosa di molto più prezioso, e la cifra
potrebbe crescere parecchio prima che qualcuno la consideri più preziosa dei suoi pomeriggi:
tutti sanno che la vera donazione è \textbf{il tempo} investito, che vale più del bene
materiale. E basta assistere: vedere qualcuno che aiuta un'anziana ad attraversare la strada fa
stare meglio anche noi. È percepita a livello di gruppo. Infine, l'ossitocina rafforza il
\textbf{sistema immunitario}: sentirsi sicuri e tranquilli fa stare meglio e, in un certo senso,
ammalare di meno.

\subsection{Cortisolo}

L'opposto è il \textbf{cortisolo}: la sensazione di stress, paranoia e paura, la risposta
\emph{fight or flight} (o scappo o combatto), con battito accelerato, muscoli in tensione e così
via. Tra i suoi effetti:
\begin{itemize}
  \item \textbf{blocca il sistema immunitario}: se si sta per morire, al corpo non importa di
        prendersi un raffreddore. È il motivo per cui le persone stressate si ammalano di più;
  \item \textbf{blocca l'ossitocina}: non ci si può sentire tranquilli quando si è in ansia;
  \item soprattutto, e questo è il punto più importante per le aziende, è
        \textbf{contagioso}. In qualsiasi documentario, quando una gazzella che bruca alza la
        testa e percepisce un movimento nella savana, scappano tutte, non solo lei. Succede,
        non sorprendentemente, anche in ufficio: quando qualcuno ha paura perché potrebbero
        esserci licenziamenti, o perché il manager è molto aggressivo, la paura si diffonde nel
        gruppo ed è molto dannosa.
\end{itemize}
È anche l'effetto alla base della \textbf{sindrome dell'impostore}: sentirsi costantemente
inadeguati perché si teme che succeda qualcosa significa vivere in uno stato di stress. Ed è il
motivo per cui chi è insoddisfatto del lavoro tende a stare peggio anche fisicamente e ad
ammalarsi di più.

\begin{center}
  \begin{tabular}{@{}p{0.17\textwidth}p{0.36\textwidth}p{0.37\textwidth}@{}}
    \toprule
    \textbf{Sostanza} & \textbf{Che cosa fa (in sintesi)} & \textbf{Esempi dalla lezione} \\
    \midrule
    Endorfine & mascherano il dolore di uno sforzo che ci fa bene &
      \emph{runner's high}, ridere fino ad avere male ai fianchi \\
    \addlinespace
    Dopamina & ricompensa un'azione, si attiva vedendo l'obiettivo; crea dipendenza &
      spuntare una checklist, le tracce per il cacciatore, gioco d'azzardo, social \\
    \addlinespace
    Serotonina & riconoscimento sociale, anche indiretto; ``sostanza della leadership'' &
      lauree e premi in pubblico, chi difende il villaggio, \emph{hoarding} \\
    \addlinespace
    Ossitocina & sicurezza e legame; rafforza il sistema immunitario &
      il nemico comune, i veterani, la stretta di mano, il volontariato \\
    \addlinespace
    Cortisolo & stress, \emph{fight or flight}; blocca difese immunitarie e ossitocina; è
      contagioso & la gazzella che scappa, la paura dei licenziamenti, la sindrome
      dell'impostore \\
    \bottomrule
  \end{tabular}
\end{center}

% ---------------------------------------------------------------------
\section{Domande sulla biologia del gruppo}
\label{sec:z2-qa-biologia}

\paragraph{Il gioco d'azzardo include i gacha e le loot box?} Assolutamente sì, anzi ne sono il
\textbf{driver principale}: i gacha, i FIFA Ultimate Team, lo ``sbustare'' da \emph{Magic} a
\emph{Hearthstone}, e in generale il ``vediamo se stavolta mi va bene'' sono tutti gioco
d'azzardo. Lo studente aggiunge che molti giochi mobile sono costruiti su questi meccanismi e che
c'è chi gioca più a quelli che ai giochi ``normali''. È il motivo per cui alcuni stati hanno
dichiarato illegali certi giochi\footnote{Il caso più noto è quello del Belgio, che dal 2018
considera gioco d'azzardo le \emph{loot box} a pagamento; anche i Paesi Bassi le hanno
contestate.}.

\paragraph{Come si genera l'ossitocina nei rapporti da remoto?} ``Bellissima domanda.'' È il motivo
per cui, come detto il giorno prima, il relatore insiste perché i meeting si facciano almeno in
\textbf{webcam}, perché i riconoscimenti siano dati e i successi celebrati \textbf{pubblicamente}
(quando non mette troppo a disagio la persona), e in generale cerca di creare \textbf{situazioni
condivise}. Anche lavorando da remoto, insiste per avere il \emph{team building} dal vivo e almeno
\textbf{un incontro all'anno} di persona in ufficio, perché crea situazioni di rapporto diretto.

\paragraph{Come si contrasta il cortisolo?} In chat qualcuno suggerisce la meditazione: sì,
assolutamente. Per il lavoro, la risposta arriva subito dopo nella presentazione.

\paragraph{Il nemico comune è una strategia che un manager può usare?} Spesso \textbf{il manager è
il nemico comune}, senza bisogno di sforzarsi. Ma il relatore chiede di capire bene il punto:
alcuni manager applicano consapevolmente il concetto ``mi ergo io a nemico del team, così sono
tutti uniti contro di me e lavorano bene'', e ne fanno uno stile di management. Decidono di essere
il \emph{bad cop}, il \emph{villain}, ed è così che ottengono risultati. È una scelta lecita e
funziona; il relatore preferisce fare diversamente, ma non per questo dice che non funzioni. Molte
volte si sottovaluta quanto sia consapevole: ``non si rende conto di quanto sia antipatico a
tutti''; forse se ne rende perfettamente conto e lo fa apposta per essere il bersaglio. Il relatore
l'ha visto fare spesso come soluzione disperata, quando il team era troppo impegnato a litigare al
suo interno: ``piuttosto che vederli litigare tra loro, gli do io un motivo per odiare me, così si
uniscono contro di me''. Non lo consiglia assolutamente, ma l'ha visto funzionare.

Uno studente racconta che nella sua azienda il manager fa vedere \emph{il lavoro stesso} come il
nemico (``prima finiamo, meglio è''). È un modello più basilare ma, come dice lo studente stesso,
funziona.

\paragraph{Quanto è etico puntare sulla dipendenza del giocatore per fare soldi?} ``Non è
assolutamente etico, ed è assolutamente profittevole, altrimenti le aziende non lo farebbero.'' Il
relatore racconta una cosa che lo ha sconvolto quando l'ha scoperta, a Blizzard. Lavorava su
\emph{Hearthstone}\footnote{Nella trascrizione ``Arsdol''.} e si stava discutendo di come
monetizzare il gioco; venne preso come consulente un esperto, israeliano se ricorda bene, che
lavorava a Malta per società di poker online e di \emph{gambling}. Raccontò che molte aziende che
monetizzano sulla dipendenza dal gioco d'azzardo avevano sviluppato metodi più o meno etici (``più
spesso meno''). L'esempio culminante: analizzavano gli utenti occasionali, quelli che giocavano
magari una settimana al mese, per capire \emph{quando arrivava loro lo stipendio} (di solito si
deposita lo stipendio e poi si gioca quello che resta), per colpirli il giorno dopo con email di
offerte per tornare a giocare. Li si stimolava proprio nel momento in cui avevano disponibilità, sapendo
che erano sensibili, per aumentare le probabilità di riaverli come giocatori. ``È orribile. È
terrificante e, come tutte le cose terrificanti, è anche altamente funzionale.''

Gli ha fatto riflettere molto su dove ci si possa spingere per fare soldi. La versione etica, per
lui, è ad esempio far uscire un gioco sotto Natale, quando si è più disposti a regalarlo ai figli:
si stanno ottimizzando le proprie possibilità, e lo accetta. Si scade nel \textbf{predatorio}
quando si sfrutta la debolezza di qualcuno in un momento preciso per estorcergli soldi che non può
permettersi di perdere. Poi ognuno tragga le proprie conclusioni: ``non è un workshop sull'etica del
lavoro''. Non pensa di dover imporre la propria opinione, né si arroga il diritto di dire che chi lo
fa è una bestia senza cuore, visto che è il loro lavoro: ``chiunque riesca poi a dormire in pace con
se stesso, buon percorso''.

% ---------------------------------------------------------------------
\section{La piramide della fiducia}
\label{sec:z2-piramide}

\subsection{La leadership non è un titolo}

Chiuso il ``momento Piero Angela'', il relatore arriva al punto. La frase della prima giornata (``la
capacità di un gruppo di persone di fare cose straordinarie dipende dal modo in cui queste persone
operano insieme'') porta a capire che \textbf{la leadership non è un titolo lavorativo}. Essere
manager da contratto, o avere ``manager'' nel titolo, rende un manager, non un leader. Il relatore
non ha mai visto un team rispettare qualcuno solo perché da contratto doveva farlo: si può avere
paura o timore del proprio manager, ma non lo si deve rispettare perché è il manager. \textbf{Il
rispetto si guadagna}, e si possono influenzare le persone intorno solo attraverso il rispetto
guadagnato, quando si sono accumulati abbastanza esempi e \emph{case history} da giustificare
rispetto e fiducia. Altrimenti, dice, sarebbe il primo a chiedere ``perché vi fidate di me, se non
vi ho ancora dato motivo di farlo?'': non sarebbe razionale.

Da qui la \textbf{piramide della fiducia}, che si costruisce partendo dalla base.

\subsection{Alla base: l'integrità morale}

Per essere un buon leader o un buon manager bisogna \textbf{voler genuinamente aiutare gli altri
membri del team}. È il motivo del discorso sulle carriere divise: è a questo livello che si
screma chi vuole la promozione per i soldi. Chi vuole essere leader non deve cercare il guadagno
personale, ma quello del gruppo. Come detto il giorno prima, il relatore è pronto a litigare per il
suo team e a prendersi le colpe per il suo team, perché vuole il bene del suo team e sa che poi loro
si prenderanno cura di lui. Chi ha in mente questo sta in questa categoria. Poi:
\begin{itemize}
  \item \textbf{essere onesti}. Il relatore può affermare di non aver mai mentito
        intenzionalmente a un collega, a un cliente o a un giocatore. Gli è capitato di dare
        informazioni errate perché in quel momento sapeva quelle: ad esempio annunciare che
        un'espansione sarebbe uscita in un certo mese, e poi l'uscita è slittata\footnote{Nella
        trascrizione l'esempio non è coerente: prima si parla di un'uscita a novembre rimandata
        di una settimana, poi di ``pensavo fosse settembre''. Il senso resta chiaro.}. Ha detto
        una cosa sbagliata, ma tutti hanno capito che non stava mentendo: in quel momento pensava
        che fosse vera;
  \item soprattutto, \textbf{accettare la responsabilità dei fallimenti}: quando si sbaglia,
        saper dire ``ho sbagliato''. Il relatore lo ha messo perché ``non avete idea di quanti
        facciano fatica'' ad ammettere di poter sbagliare: c'è sempre una scusante, una
        mitigazione, un ``sì, ma''. Lui vuole sentire se è colpa tua o no;
  \item \textbf{chiedere scusa} e \textbf{ammettere di non sapere}: anche dire ``non lo so'' viene
        spesso evitato, come accettare un fallimento.
\end{itemize}

\paragraph{``Raccontami il tuo più grande fallimento''.} Per lo stesso motivo, nei colloqui il
relatore chiede spesso, più che ``parlami di un tuo grande successo'' o di un progetto andato bene,
\textbf{``mi racconti il tuo più grande fallimento?''}, e osserva come lo raccontano. Chi non si
aspetta la domanda spesso racconta di quella volta in cui ha provato in tutti i modi a fare in un
certo modo, ma il capo gli ha detto di fare diversamente e il progetto è andato male. ``Non mi hai
risposto: mi hai raccontato un progetto andato male, ma il tuo messaggio è che volevi fare
diversamente e qualcuno ti ha detto di fare così. Quindi non è un tuo fallimento.'' È come
rispondere a ``qual è la tua debolezza?'' con ``sono così attento ai dettagli che a volte\dots'':
si fa passare un punto di forza per una debolezza. Il relatore vuole un esempio in cui \emph{tu},
senza altri fattori, hai sbagliato, e come ti sei comportato. Gli dice molto su come quella persona
reagirà ai suoi feedback: se ha davanti ``un muro'' che dirà sempre ``non è colpa mia'', non è adatta
a un ambiente in cui il feedback è alla base di ogni comunicazione. Chi non ammette di aver mai
sbagliato in vita sua, lo scarta.

\paragraph{È possibile colpevolizzarsi troppo per un fallimento?} Sì, assolutamente. Rientra molto
nella sindrome dell'impostore, o nell'aver sempre vissuto rapporti di lavoro tossici: ci si abitua ad
accettare che ci vengano urlate contro le cose e che sia sempre colpa propria, e ci si autoconvince
che lo sia.

\subsection{Al centro: l'affidabilità}

Superato questo livello si arriva all'\textbf{affidabilità}. Per essere un buon leader non bisogna
essere il migliore, ma affidabili: basta essere riconosciuti come bravi in qualcosa, non
necessariamente i migliori, e quell'affidabilità ce la si porterà anche nella leadership. Significa:
\begin{itemize}
  \item \textbf{rispettare la parola data}: se si promette ``parleremo di un aumento al prossimo
        slot disponibile'', se ne parlerà. Non ci si dimentica, e non si dice poi ``in realtà non si
        poteva fare, ma ho voluto provarci fino all'ultimo''. Una promessa va mantenuta;
  \item \textbf{impegnarsi per il tempo necessario} a raggiungere gli obiettivi;
  \item \textbf{saper rifiutare un lavoro} quando necessario: è il saper dire di no della mattina,
        ancora la matrice di Eisenhower (sezione~\ref{sec:z2-eisenhower}). ``Sono sempre gli
        stessi concetti.''
\end{itemize}

\subsection{In cima: la leadership, cioè dare l'esempio}

L'ultimo gradino è la \textbf{leadership}\footnote{Nella trascrizione ``dealership''.}, cioè
\textbf{dare l'esempio}. Non bisogna essere i più bravi; bisogna saper \textbf{offrirsi volontari
non solo quando il progetto è un successo sicuro, ma anche quando c'è il rischio di un
fallimento}. Sono tutti bravi a offrirsi per il gioco nuovo che andrà sicuramente bene.

\paragraph{Il rappresentante dei licenziati.} Al relatore è capitato di offrirsi per rappresentare
il team durante i \emph{layoff}, con il rischio di essere ostracizzato dall'azienda; con quelle
persone è ancora in contatto. Quando Blizzard fece i licenziamenti nella sede francese, vicino a
Versailles\footnote{Nella trascrizione ``i layoff di Versailles''. La sede europea di Blizzard si
trovava a Vélizy-Villacoublay, nei pressi di Versailles; il relatore non precisa l'anno.}, che
coinvolsero quasi tutto l'ufficio, lui non era tra i licenziati, ma venne \textbf{eletto
rappresentante sindacale} di chi lo era. Per lui è stato molto importante: pur non essendo parte in
causa, lo ritennero degno della loro fiducia. Per quanto ne sapevano, avrebbe potuto negoziare e
accettare tutto quello che proponeva Blizzard, perché non ne era toccato: che si negoziasse un mese
di liquidazione in più o in meno non gli portava un soldo. Avrebbe persino potuto usare la
situazione per negoziare un aumento per sé (``gli facciamo accettare tre mesi invece di sette e a te
diamo un aumento''): non era fuori discussione. Se si sono fidati è solo perché c'era una piramide di
questo tipo e lui aveva dato loro abbastanza motivi per farlo.

\paragraph{Autorevolezza, non autorità.} A uno studente che lo riassume così, il relatore risponde
che è proprio questo: un carisma diverso dall'essere simpatici, l'essere un punto di riferimento,
\textbf{più un pilastro di supporto che il motore} che fa andare avanti le cose. Non il motore davanti
a tutti, ma il \textbf{meccanico} che, quando si rompe qualcosa, sai che non ti frega: ti sistema le
cose senza ``farti la cresta'' sulle pastiglie dei freni.

\paragraph{Non abbandonare nessuno.} E soprattutto \textbf{non abbandonare chi è in difficoltà}.
Sembra una banalità, ma è alla base del motto dei Marines statunitensi, \emph{no one left behind}:
sapere che, se si viene feriti, verranno a prenderti anche sotto i missili è un forte elemento
motivazionale per correre quel rischio. In azienda è il motivo per cui il relatore ribadisce sempre
che, se qualcuno è in difficoltà, la sua prima opzione non è licenziarlo ma \textbf{dargli una mano},
e lo ha fatto a volte anche all'eccesso, per far vedere agli altri che quando saranno in difficoltà
non li lascerà in mezzo a una strada. A volte ce l'ha fatta: ha esempi di persone ``a tanto così''
dal licenziamento recuperate per i capelli e diventate ottimi \emph{performer}. Altre volte la
persona non era recuperabile e alla fine è stata licenziata. Ma chi attraversa un periodo difficile
(il \emph{burnout}, la depressione, una relazione finita) sa che la prima risposta non sarà ``ci sei
servito finora, ora puoi andartene'', ma ``che cosa possiamo fare?''. E questo fa lavorare tutti
tranquilli.

% ---------------------------------------------------------------------
\section{Prendersi cura del gruppo}
\label{sec:z2-gruppo}

\subsection{La soluzione: più leader possibile}

In un caso del genere manager e leader coincidono, ed è l'ideale: un manager dovrebbe essere anche
un leader. Ma la vera soluzione è \textbf{avere più leader possibile}. Se l'obiettivo del leader è
prendersi cura degli altri, e il manager dovrebbe essere ``il primo tra i leader'', ancora meglio è
che \textbf{tutti siano leader del proprio gruppo}, perché la leadership non deriva da un titolo né
dall'essere il manager. Il relatore ha visto tanti gruppi andare avanti perché qualcuno al loro
interno era il leader, non perché lo era il manager.

Vuol dire mettere in pratica tutto il discorso della prima giornata su reciprocità, gentilezza,
aiuto e collaborazione. Si sa che, nel peggiore dei casi, c'è il manager, ma in condizioni normali ci
sono i colleghi: e allora si va da chi è in difficoltà a dire ``posso darti una mano?''. Così, quando
saremo noi in difficoltà, magari senza nemmeno rendercene conto, sarà qualcun altro a dirci ``ti vedo
in difficoltà, come posso aiutarti?''.

\paragraph{L'esempio è contagioso, nel bene e nel male.} Quando si ha una posizione anche solo di
rispetto agli occhi degli altri, tutto ciò che si fa ha effetto sugli altri. In negativo vale la
\textbf{slippery slope}\footnote{L'esempio del cartone abbandonato ricorda la cosiddetta ``teoria
delle finestre rotte'' (\emph{broken windows theory}), secondo cui i piccoli segni di degrado
tollerati favoriscono comportamenti peggiori.}: se si accetta un cartone di pizza abbandonato per
strada, domani se ne accetteranno due, dopodomani tre. Ma vale anche in positivo: chi vede qualcuno
fare qualcosa di positivo pensa ``lo voglio fare anch'io''. ``Come fai una cosa, fai tutto'', dice un
proverbio che il relatore attribuisce alla tradizione cinese\footnote{La massima (``how you do
anything is how you do everything'') è attribuita a fonti diverse, spesso alla tradizione zen;
l'attribuzione è incerta.}.

\subsection{Il villaggio di oggi}

Si torna al villaggio del ``momento Piero Angela'', con i difensori e le bestie (reali o no) intorno.
Oggi fuori dal gruppo non ci sono più tigri, ma le \textbf{deadline}, l'\textbf{overwork}, i
\textbf{layoff}, il manager che urla, le aspettative, i \textbf{KPI}. Avere manager che si prendono
cura del gruppo è un buon inizio, ma è ancora meglio avere \textbf{il gruppo che si prende cura del
gruppo}. Non bisogna più limitarsi a pensare che sia responsabilità del manager prendersi cura di
noi: è responsabilità nostra prenderci cura del gruppo di cui facciamo parte.

Che cosa serve? La piramide vista prima. Serve essere percepiti come \textbf{autentici}: il
relatore ormai riconosce abbastanza facilmente chi si vuole ergere a leader perché gli piace
\emph{sentirsi} leader, non perché vuole \emph{fare} il leader, e quella mancanza di genuinità non
aiuta. In generale bisogna \textbf{generare sicurezza} nel resto del gruppo.

\subsection{Smettere di dire ``ci penserà qualcun altro''}

Un esempio concreto è smettere di dire \textbf{``ci penserà qualcun altro''}. Nello sviluppo è
evidente: se si trova un bug che non riguarda il proprio lavoro, tutti sono d'accordo che vada
comunque segnalato, perché altrimenti potrebbe arrivare in \emph{live}, nella \emph{release}.

\paragraph{Le lasagne di PG Esports.} Circa un anno e mezzo prima, quando lavorava ancora a PG
Esports (che organizza anche eventi dal vivo), l'azienda ospitava in ufficio dei potenziali sponsor
venuti a conoscere il team per decidere se sponsorizzare la Milano Games Week. C'era un catering, con
i classici buffet. In mezzo alla confusione (una ventina di persone di PG Esports e una trentina tra
manager, sponsor ed \emph{executive}), il relatore vede per terra, in mezzo alla stanza, una mezza
porzione di lasagne abbandonata. Senza pensarci la raccoglie per portarla via: non è sua, ma è un
piatto in mezzo alla stanza. Nel tragitto fino al bancone del buffet, almeno altri sei tra executive e
sponsor che non lo conoscevano gli consegnano il loro piatto non finito, come a un cameriere: nella
loro testa l'unico motivo per cui qualcuno potesse raccogliere un piatto da terra era che fosse il
suo lavoro. Sul momento non ci fa caso; arrivato al bancone si gira e vede tutto il suo team che ride
della scena.

Torna da loro e ne nasce una discussione: ``avete visto quanto è stato educativo?''. Avevano visto
esattamente come ragionano i gruppi che non sono abituati a pensare ``adesso ci penso io'', ma ``ci
penserà qualcun altro''. In quelle persone non c'era malizia: sicuramente era un cameriere, perché
altrimenti farlo? ``Non è il mio lavoro, lo avrà fatto perché è il suo.'' E diventa un problema quando
la stessa tendenza (``non c'entra con il mio lavoro, ci penserà qualcun altro'') si applica a
qualsiasi altra cosa. Il relatore si aspetta che \textbf{tutti raccolgano quel piattino}: che chi vede
un bug, se non può correggerlo lo segnali, e se può correggerlo lo corregga anche se non è il suo
lavoro, ``perché è un piattino in mezzo alla stanza''.

\paragraph{Da reattivi a proattivi.} Significa passare da un atteggiamento \textbf{reattivo}, cioè di
fatto passivo, a uno \textbf{proattivo}: prevedere i problemi e cercarli intorno a sé per risolverli,
invece di aspettare che esplodano (``finché non esplode non è un problema mio''). Ragionare nell'ottica
del ``e se succedesse, e io me ne fossi accorto senza fare niente?''. Uno studente lo chiama
\emph{mindset da tester}: il relatore si aspetta che, se si vede qualcosa, ci si assicuri di non essere
gli unici ad averlo visto. A volte basta farlo presente: il relatore non potrebbe mai correggere un
bug, quindi se lo trova lo segnala, e va bene così. Ma non bisogna dare per scontato che ``è un bug così
evidente che qualcuno l'avrà già segnalato''. È il paradosso dell'incidente stradale: \textbf{più
l'incidente è grosso, più è probabile che nessuno chiami l'ambulanza}, perché tutti pensano che
qualcuno l'abbia già fatto\footnote{È il cosiddetto \emph{effetto spettatore} (\emph{bystander
effect}), studiato in psicologia sociale a partire dagli anni Sessanta.}.

\subsection{Domande sul ``me ne occupo io''}

\paragraph{Se correggo un bug, devo avvisare il collega?} Dopo una breve pausa per problemi audio
(``facciamo \emph{double testing}''), uno studente chiede se, correggendo il lavoro di un collega, il
dovere sia poi andare a dargli il feedback. A prescindere bisogna informare qualcuno; e la
conversazione dovrebbe essere: ``ho trovato questa cosa, vuoi che te la sistemi io?''. Correggere
senza dirlo a nessuno potrebbe essere un danno: magari era \emph{working as intended}, o qualcuno ci
stava già lavorando.

\paragraph{E se divento ``quello che risolve quel problema''?} Uno studente porta la sua esperienza:
in un'azienda in cui ha lavorato tre anni, risolvendo un problema facile che non era di sua
competenza è diventato, di fatto, l'incaricato di quel problema, distolto dalle sue mansioni. È
un'arma a doppio taglio. Per il relatore lì il problema è del \textbf{management}. In alcune aziende si
applica il criterio ``chi trova il problema lo risolve'' (``se te ne sei accorto tu, pensaci tu''), ed
è un altro modo di scaricare la responsabilità sugli altri. Non è ciò che dice lui: deve essere chi
trova il problema a valutare se è giusto e positivo che lo risolva lui. Il cartone di pizza lo
raccoglie perché ci mette dieci secondi; ma se ha trovato un bug nel sistema di pagamento che richiede
tre settimane di lavoro, e chi l'ha introdotto fa finta di niente, non deve risolverlo lui: deve far
presente che c'è un problema, che è stato ignorato e che va risolto (``pensateci voi, ma io vi dico
che c'è''). Va fatto caso per caso: se è immediato, facile e nel proprio interesse, lo si risolve.
Ma il relatore sconsiglia alle aziende la regola ``se lo trovi lo risolvi'', perché \textbf{disincentiva
a trovare i problemi}: qualcuno potrebbe non segnalarli per non avere altro lavoro e aspettare che li
trovi qualcun altro, ed è negativo quanto la situazione descritta. Per chiudere, ne discuterebbe con il
management: ``possiamo rivedere questa policy? Incentiva l'azione sbagliata''.

\paragraph{E se ci dicono di non mettere mano?} Uno studente che lavora come tester racconta che spesso
potrebbero correggere certe cose, ma viene loro detto di non toccarle. Esatto: si comunica. ``Ho trovato
questa cosa, ne siete consapevoli? Sì? Ci state già lavorando? No? Posso farlo io? Sì? Allora lo
faccio.'' Non bisogna esagerare con l'iniziativa: anche qui vale il \emph{default to trust}. Si parte
dal presupposto che qualcuno potrebbe averlo già visto e starci già lavorando, quindi si chiede e si
offre il proprio supporto.

% ---------------------------------------------------------------------
\section{Se si diventa manager}
\label{sec:z2-manager}

\subsection{Meriti, colpe e \emph{big picture}}

Una breve digressione per chi, un giorno, diventerà manager: che cosa cambierebbe? Oltre a tutto
quanto detto, ci sono il discorso su \textbf{meriti e colpe} e quello su \textbf{progetto e task}:
\textbf{delegare i meriti e prendersi le colpe}, e pensare al progetto più che ai singoli task, cioè
alla \emph{big picture}, invece che al ``voglio tornare a fare le cose che sapevo fare''.

È un altro problema dei manager promossi dal gruppo di chi sapeva fare bene un certo lavoro: quando
hanno l'occasione di tornare a farlo, lo fanno molto volentieri. ``C'è un task da fare, lo faccio io'',
magari con la scusa di dare l'esempio. No: si sta andando nella propria \textbf{comfort zone}. Il
relatore lo capisce (anche lui tornerebbe volentieri a fare ciò che faceva prima), ma bisogna accettare
che \textbf{quelle cose non le si deve più fare}.

\subsection{Il ruolo è più grande della persona}

Anche da manager, il ruolo che si ricopre per il gruppo è \textbf{più importante di sé}. Non importa
chi sia il community manager di Blizzard, di Riot o di un'altra azienda: conta ciò che rappresenta per i
giocatori in quanto community manager, non il fatto che sia lui piuttosto che un altro.

Prendersi cura del gruppo vuol dire applicare concetti molto semplici:
\begin{center}
  \begin{tabular}{@{}p{0.42\textwidth}p{0.48\textwidth}@{}}
    \toprule
    \textbf{Invece di dire\dots} & \textbf{\dots dire} \\
    \midrule
    ``Qui non si accettano errori.'' &
      ``Spiegami perché abbiamo sbagliato'': fammi capire perché è successo
      (sezione~\ref{sec:z1-default-trust}). \\
    \addlinespace
    ``Qui si fa come dico io.'' & ``Tu come faresti?'' \\
    \bottomrule
  \end{tabular}
\end{center}

\subsection{Il tetto della bravura del capo}

Il secondo punto è cruciale soprattutto quando si promuove il più bravo in un lavoro: se il criterio
noto è che si promuove il più bravo, il promosso lo sa, e la prima cosa che penserà è ``faccio lavorare
tutti come me, visto che sono il più bravo''. Il relatore chiede alla classe: che cosa si nega l'azienda
con questa dinamica? Anche ammesso che sia davvero il più bravo, se tutti fanno come lui diventeranno
tutti più bravi, no?

Il problema è che \textbf{si fissa come tetto massimo la bravura di quella persona}: se tutti fanno solo
come dice lui, si arriverà al massimo al suo livello, mentre ognuno ha il proprio metodo. Da qui quanto
detto il giorno prima sulla tendenza dei manager ad assumere solo persone peggiori di sé: se il criterio
di promozione è la bravura, si prenderanno sempre persone un po' meno brave per evitare che diventino il
proprio capo. Se invece le carriere sono divise (``se sei bravo fai questo, se vuoi fare il \emph{people
manager} fai quest'altro''), il manager è incentivato a prendere persone \textbf{migliori di sé}; e se,
come detto il giorno prima, il bonus dipende anche dalle performance del gruppo, è nell'interesse di
tutti che il gruppo sia pieno di superstar, un ulteriore incentivo ad assumere i più bravi.

\paragraph{L'opzione C.} C'è poi un'altra cosa, forse minima, ma per il relatore molto importante: si
perdono le \textbf{occasioni di innovare}. Nelle aziende precedenti era visto come il punto di riferimento
per il \emph{problem solving}: in un'emergenza, se c'era da scegliere tra A e B, ``decide Stefano''. Andava
benissimo, ma la tendenza era che a ogni emergenza venivano da lui. Così ci si negava \textbf{l'opzione C}:
quella a cui lui non avrebbe pensato e qualcun altro sì. Magari la prima volta sbaglia, la seconda anche,
e la terza ha un'idea migliore della sua, proprio perché lui non è più l'esperto di quel settore. Fissare
tutto come ``al massimo decide Stefano'' significa non accettare altre idee: un limite molto grosso
all'innovazione, all'aggiornamento e al fare le cose in modo più moderno e flessibile.

\subsection{Celebrare i successi e i fallimenti}

\paragraph{Celebrare i successi del team.} Il relatore vuole assicurarsi che tutti capiscano quando si è
ottenuta una vittoria o un successo: come detto il giorno prima, se non si sa che cosa si fa bene si
rischia di smettere di farlo. ``Questa cosa l'abbiamo fatta bene, complimenti, bravi tutti.''

\paragraph{Celebrare i fallimenti.} Sui fallimenti si era spinto ben oltre il semplice accettarli: aveva
indetto \textbf{meeting in cui si celebravano i fallimenti}, con le persone che raccontavano i propri errori
per distribuire quella conoscenza al resto del team. ``Non mi fa arrabbiare l'errore fatto, mi fa arrabbiare
lo stesso errore fatto quindici volte.'' Condividere un fallimento con il gruppo non significa ``guardate
questa persona sciocca che ha fatto questo errore'', ma ``questa persona ci sta permettendo di evitare lo
stesso errore in quindici'': sta salvando l'azienda dal ripeterlo quindici volte, perché la gente ha paura di
dire agli altri che ha sbagliato. Mostrare dove si è sbagliato diventava quindi un incentivo, un servizio
all'azienda. E si torna al punto di prima: se il proprio ruolo per il gruppo è più importante di sé, non
importa dover alzare la mano e dire ``ho sbagliato, ho fatto così'', perché gli altri penseranno ``grazie,
così io non lo faccio''. Si dovrebbe essere lieti che qualcuno abbia fatto quell'errore prima di noi, e
andare a ringraziarlo per averlo condiviso.

\paragraph{Condividere i successi.} Infine bisogna \textbf{dare il merito a chi lo ha}. Al relatore è
capitato spesso di ricevere ringraziamenti come ``grazie a Stefano e al suo team per aver fatto questa
cosa'', e la sua risposta è stata ``grazie, ma ha fatto tutto Gianni''. Non bisogna avere paura di dire
``grazie del feedback, ma è stato lui'': è giusto ringraziare lui e il team (il relatore si sente parte del
team, chiaramente). Non si ottiene niente dal dire ``grazie, riferirò i complimenti al team'': il relatore
non ha mai creduto nel ``riferire'' che il management ha detto ``ottimo lavoro''. Chiede al management di
dirlo \textbf{direttamente} al team, perché sentirlo dire dal management ha un impatto diverso che sentirlo
riferire da lui. In questo senso, osserva uno studente, ogni gruppo è il gruppo di sostegno di se stesso:
``assolutamente sì''.

\subsection{Quanto ci vuole?}

Non c'è una scadenza. Il rispetto non si ottiene pagando per fare prima, ``con il \emph{boost} a livello
60''. Servono tempo ed energia, ed è difficilmente misurabile: il relatore non può dire quanto il team si
fidi di lui in questo momento. Lo percepisce da come lavorano, ma lo si vedrà davvero solo quando ci sarà un
problema grosso; fino ad allora è tutta una sua percezione. Il team lo sente ripetere queste cose ogni giorno
e lo vede metterle in pratica ogni giorno, e lui può solo sperare che il messaggio passi; quando ci saranno
davvero ``gatte da pelare'' capirà come è messo, perché non c'è momento migliore di una difficoltà per vedere
come reagisce un team.

\textbf{Essere leader è prendersi un impegno a lungo termine verso i propri colleghi}: non ci sono modi per
velocizzarlo. Bisogna fare il proprio lavoro ogni giorno ``a testa bassa'', aiutando gli altri, facendolo
vedere ed \textbf{esplicitando perché lo si fa}. Il relatore lo fa quotidianamente: ``ti sto aiutando; nel
caso avessi un dubbio, lo faccio perché il mio lavoro è aiutarti e io ne traggo vantaggio dal fatto che tu
lavori meglio. Sono qui per te, ma lo faccio anche per me: non pensare che sia solo gentilezza. È il mio
lavoro, ma sappi che lo faccio spontaneamente''.

\subsection{L'ultima slide}

L'obiettivo dovrebbe sempre essere lo stesso: tutti dovremmo poter finire la giornata, staccare e tornare a
casa sentendoci
\begin{itemize}
  \item \textbf{ispirati}, perché abbiamo fatto qualcosa che ha un effetto tangibile;
  \item \textbf{sicuri}, perché sappiamo di avere colleghi, un leader e dei manager che si prendono cura di
        noi;
  \item \textbf{realizzati}, per quello che facciamo e perché è ciò che ci piace fare.
\end{itemize}
E questo dovrebbe essere \textbf{un diritto naturale di tutti}, non un ``sì, ma solo il manager'', ``sì, ma
solo se sei un esperto'', ``sì, ma solo se sei un game director''. Ognuno ha le proprie capacità: bisogna
solo trovare ciò che piace e che si vuole fare.

% ---------------------------------------------------------------------
\section{Domande e risposte finali}
\label{sec:z2-qa-finale}

Finita la presentazione, si passa alle domande: ``chiedete tutto quello che volete, perché è
l'ultima possibilità''. Le domande sono riportate per temi.

\subsection{Leadership e organizzazione del lavoro}

\paragraph{La leadership riguarda tutti o solo chi vuole diventare leader?} Il discorso sui
manager vale solo per i manager; ma il consiglio è \textbf{essere tutti leader del proprio gruppo,
da dentro}. Essere un leader significa sapere che gli altri possono fare affidamento su di sé, e
di queste figure in un team ce ne possono essere infinite: non è una gara a chi è il leader.
``Siate tutti una di quelle figure a cui gli altri guardano quando sono in difficoltà.''

\paragraph{Da manager non si hanno più orari e si fanno straordinari di continuo?} No: il relatore
ha notato di essere molto più efficiente come manager quando si dà \textbf{paletti molto precisi}.
Se stacca alle sei e alle sei e cinque arriva un messaggio, risponde ``se non mi dici che è
urgentissimo, non lo guardo nemmeno, lo leggo domani''; se non è urgente, lo guarda il giorno dopo.
Ma gestisce le aspettative: ``se è urgente ci sono e ci sarò sempre, ma non aspettatevi che faccia
orari eccessivi solo perché sono manager, così come non me lo aspetto da voi''. Si aspetta che
ognuno faccia ciò che deve; se è urgente tutti faranno di più, ma se non lo è, no.

\paragraph{L'orario davvero flessibile.} Qui il relatore aggiunge una cosa che si era dimenticato di
dire parlando di ferie e permessi illimitati: introduce anche l'\textbf{orario davvero
flessibile}. Non il flessibile di molte aziende, in cui si può iniziare alle 9 e finire alle 18
oppure iniziare alle 10 e finire alle 19: se oggi si vuole iniziare alle 10 e finire alle 15, e si è
fatto tutto quello che si doveva, va bene così. Ci si prendono le ore libere in più per altro, per
staccare, per riposarsi, sapendo che se la settimana dopo bisognerà lavorare dalle 9 alle 21, lo si
farà (o si dirà ``ieri ho fatto dalle 9 alle 21, domani faccio dalle 10 alle 15''). L'obiettivo dato
è un lavoro e un progetto, non un orario: si lavora per completare i propri task e rispettare le
deadline. L'orario indicato è solo indicativo, ed è quello dei meeting: \textbf{nessun meeting prima
delle 10 o dopo le 17}. Se poi qualcuno vuole lavorare dalle 10 alle 11 e poi di notte, perché quel
giorno preferisce lavorare di notte, più rilassato, faccia pure: l'importante è che trovi il tempo per
riposare e staccare. Ci si autogestisce senza timbrare cartellini; l'unico accordo è ``hai un task, va
consegnato entro il giorno X e mi garantisci che venga fatto; come ci arrivi non mi interessa''.

\paragraph{Le è mai capitato di essere aiutato da qualcuno, in un momento di difficoltà?} Sì: come
detto, questo comportamento è \textbf{contagioso in senso positivo}. Se ogni volta che qualcuno è in
difficoltà qualcun altro lo aiuta, prima o poi capita anche a chi aiuta. Il relatore non lo dice
``tanto per dire'': ha avuto momenti in cui faceva davvero fatica ad alzarsi dal letto per lavorare,
perché esausto o demotivato per altri motivi, e il team gli ha detto ``stai a letto, ci pensiamo noi
oggi''. Questo lo aiuta ad aiutare, perché sa che la cosa è bidirezionale.

\paragraph{Tra le soft skill del workshop, da quale partire?} \textbf{Il feedback}: saperlo dare e
ricevere. È la più difficile, molto più del ``voglio aiutarti'', ma è anche quella su cui si costruisce
tutto il resto: attraverso il feedback si costruisce la \textbf{fiducia}, attraverso la fiducia il
\textbf{rispetto}, e il rispetto diventa la base portante di tutto il resto.

\paragraph{E se il leader se ne va?} A una domanda su che cosa fare quando viene meno la figura di
riferimento del gruppo\footnote{La domanda non è leggibile nella trascrizione; dalla risposta si
capisce che riguardava il caso in cui il leader che incarnava la visione del gruppo se ne va.}, il
relatore risponde che la domanda da farsi è un'altra: \textbf{è cambiato l'obiettivo?} La direzione è
sempre la stessa? Se sì, chi è la persona più adatta a portarci lì? Non conta tanto chi guida, ma
l'obiettivo che ci accomunava, di solito la visione di un leader in cui tutti si allineano: ci
crediamo ancora tutti, o è cambiato qualcosa? È il momento di fare autoanalisi: ``ero qui per quella
persona o per la visione?''. Si dice, ed è vero, che \textbf{le persone non si dimettono dai posti di
lavoro, ma da capi che non rispettano più}, perché spesso l'unica interfaccia con l'azienda è il
proprio capo: le dimissioni non sono un messaggio per l'azienda, ma per il capo (``non voglio più
lavorare con te''). Allo stesso modo, se non c'è più il leader della visione, bisogna prima capire se
si condivide ancora qualcosa.

\paragraph{Lavorare così è più faticoso?} Sì e no. ``Non voglio demotivarvi, perché non è il mio
lavoro'', ma non è tutto rose e fiori: lavorare in questo modo è anche più stancante
\textbf{emotivamente}, perché non si può più pensare ``faccio il mio, alle cinque ho finito e chi si è
visto si è visto''. Si sente un senso di responsabilità verso gli altri. Succede che si finisca il
proprio lavoro del pomeriggio, ci si giri e si veda un collega che finirà alle nove di sera, e
l'istinto sarà ``posso farti un pezzo del lavoro?'': così si finisce entrambi alle sei. In un certo
senso ci si accolla più lavoro, con la consapevolezza che quando si sarà in difficoltà ci sarà qualcuno
che si girerà a dire ``ti vedo stressato, posso darti una mano?''. È bellissimo, ma non dà solo benefici
diretti: richiede anche un investimento personale che non sempre le persone vogliono fare. Pro e
contro: è il modo che sceglierebbe il relatore, ma non è per tutti. Conosce persone che vogliono poter
staccare appena finito, e va benissimo: si cerca un'azienda in cui si possa lavorare così.

\subsection{Condividere fallimenti e successi in pratica}

\paragraph{Come si celebrano i fallimenti e si condividono i successi?} Condividere i successi sembra
più semplice (``bravo Mario che ha fatto questa cosa''). Il relatore dà uno spunto ancora migliore:
\textbf{``bravo Mario, e ora Mario ci racconta come l'ha fatto e com'è andata''}. Sono sottigliezze,
ma contano: raccontare in prima persona ciò che si è fatto dà un moto d'orgoglio che non dà sentirlo
raccontare da un altro.

Per i fallimenti, che sembrano dover partire dalla persona stessa, l'obiettivo era che tutti li
condividessero spontaneamente; ma all'inizio era difficile, perché non viene naturale alzare la mano e
dire ``ho sbagliato''. Per questo il relatore ha creato artificialmente dei \textbf{``festival del
fallimento''}, dal tono chiaramente ironico: durante una pausa pranzo, mentre tutti mangiavano insieme,
a turno si saliva ``sul palco'' a raccontare la stupidaggine fatta in settimana. Molto spesso iniziava
lui, raccontando i propri errori (``ho sbagliato questa cosa, sappiate che non si fa così''). Dopo aver
forzato l'interazione per una, due, tre settimane, alla quarta lo facevano da soli: era diventato
organico dire a tutti ``ragazzi, attenzione, ho fatto questa sciocchezza, non fatela anche voi'', e c'era
perfino un \textbf{tag apposito su Slack}. È stato il team a decidere come condividere i fallimenti: come
sempre, invece di imporre il modo, il relatore ha lasciato scegliere a loro, e funziona. È impegnativo
tanto quanto lo si vuole rendere tale.

\paragraph{Che tipo di fallimenti si condividono? Anche quelli personali?} Di solito si parte da quelli
\textbf{professionali}, cioè da ciò che si riconosce come utile anche agli altri: qualcosa che ha
rallentato, fatto perdere tempo o che rischia di ripetersi (``sappiate che, per evitare che lo facciate
anche voi\dots''). Il resto è più da conversazione privata, \emph{one to one}, con la persona con cui ci
si vuole sfogare, o con gli amici: se non ha un potenziale impatto sugli altri, non c'è un incentivo a
dirlo. E un errore che si è fatto per il proprio tipo di personalità? Qui tornano i test di personalità
della prima giornata: conoscendo il proprio profilo e quello dei colleghi, si sa che qualcuno ha lo
stesso carattere, e quindi lo si aiuta.

\subsection{Cultura aziendale e formazione}

\paragraph{Come si arriva a poter fare queste cose in azienda?} Conta molto la cultura aziendale. Il
relatore sa di lavorare bene perché oggi ha il \emph{background} e i successi alle spalle per dimostrare
che il suo metodo funziona, e quindi glielo lasciano fare; ma se dieci anni prima avesse proposto le
stesse cose, ``mi avrebbero mandato a quel paese''. C'è un momento e un posto per fare le cose, e non è
sempre possibile fare ciò che si ha in testa; ma ci si può lavorare. Per lui la domanda è diventata:
``sono convinto che così sarebbe meglio; come faccio ad avere carta bianca?''. \textbf{Portando
risultati}: si è concentrato su quello, e ci è arrivato progressivamente.

\paragraph{Come funziona la formazione con workshop come questo, in azienda?} Il relatore li fa in modo
un po' più rilassato ed esteso, con molte più interazioni ed esempi, perché in azienda si possono usare
esempi che si hanno davanti ogni giorno, e parlare di problemi e discussioni diventa molto più reale. Ma
in generale è sempre \textbf{formazione pratica}, ripetuta finché il messaggio non è recepito. Il segno
che qualcosa è stato recepito è quando arriva una persona nuova e \textbf{non è lui a spiegarle le cose}:
è il team stesso a spiegarle che cos'è la feedback culture, e lui non deve parlare. È la conferma che il
concetto viene perpetuato dal team.

\paragraph{Quanti manager dell'industria condividono questi principi?} Il relatore ha fatto molti
interventi all'IGDA di Roma e ne farà uno a Milano la settimana successiva (``il 24, mi pare''); interviene
spesso a workshop ed eventi, e il riscontro è sempre ``che bello, vorremmo farlo anche noi in Italia'', e ci
stanno lavorando. La moglie lavora in un'azienda che non ha niente a che fare con i videogiochi e lo stanno
già applicando. Magari in Italia è tra i primi a parlarne, perché il paese è ancora un po' indietro, ma non è
assolutamente l'unico: è più diffuso di quanto si pensi, non è una cosa di nicchia.

\paragraph{Questi spunti si applicano anche all'IT e alla consulenza?} ``Praticamente tutti.'' Si applicano
facilmente quando c'è un gruppo di persone su cui applicarli; se si cambia continuamente cliente è più
difficile portare questa mentalità da loro, ma questo non impedisce a chi lavora con noi di applicarla
direttamente con noi.

\paragraph{Ci si può tenere in contatto?} Su LinkedIn: il relatore accetta sempre tutti e risponde alle
domande quando riesce (ne riceve tante, ma cerca di stare al passo almeno su base settimanale). E, ``a costo
di sembrare quello che si fa pubblicità'', chi lavora in un'azienda che fa fatica ad applicare questi
concetti o non sa che esistono può metterla in contatto con lui: è spesso così che ha trovato nuovi modi di
applicarli. Vale anche per chi tra sei mesi lavorerà altrove. I suoi profili sono pubblici e non c'è nessun
problema a contattarlo.

\subsection{Performance e persone in difficoltà}

\paragraph{Se un'azienda è a corto di personale, conviene comunque eliminare chi performa poco?} Il giorno
prima si era detto che conviene separarsi presto da chi non performa (sezione~\ref{sec:z1-squadra}); ma non
si finisce per caricare di lavoro tutti gli altri? Per il relatore \textbf{eliminare non vuol dire non
sostituire}: vuol dire cambiare la persona con qualcuno magari più junior, ma più adatto come \emph{fit}
culturale. Il vecchio discorso del ``fare lo stesso lavoro con meno gente'' non è sostenibile: lo si può fare
fino a un certo punto, in caso di necessità, ma non è una soluzione. Non è ottimizzare, è far andare in
\emph{burnout} chi resta.

\paragraph{Ha più senso tenere una persona poco performante e non in sintonia con il team, piuttosto che
licenziarla senza poterla sostituire?} No: se una persona è poco in sintonia, performa poco ed è in generale
tossica e negativa per il team, \textbf{prima ci si separa, meglio è}. Questo non vuol dire non provare a
recuperarla, ma bisogna andare dritti al punto: ``abbiamo un problema: ti vedo poco performante, non ti vedo
allineato con il team, vedo che il team non si fida di te. Che cosa vogliamo fare?''. Se la risposta è ``non lo
so'', la strada porta dritta all'uscita. Bisogna \textbf{recidere i rami morti} il prima possibile, facendo
però attenzione a essere sicuri che si tratti di un ramo morto e non di un ramo recuperabile: tagliare un ramo
recuperabile mette negli altri la paura di essere i prossimi, ``alla prima foglia gialla''. Quello che si
vuole è che, quando un ramo è morto, la pianta smetta di alimentarlo, e ne crescerà un altro. Entrambi i
messaggi devono essere chiari: \textbf{non si subisce passivamente la situazione}, e allo stesso tempo
\textbf{non si abbandona nessuno} (se qualcuno è in difficoltà, la prima soluzione non è tagliare il ramo). È
un compromesso tra le due cose. (A questo punto arriva, puntuale, il crash della connessione: ``non era ancora
successo oggi''.)

\paragraph{Che cosa fare se qualcuno nel team sembra soffrire di depressione?} È un discorso complesso, perché
il relatore \textbf{non è uno psicologo}. Può dare la sua disponibilità a concedere permessi, coprire il lavoro
della persona e in generale darle modo di seguire un percorso con uno specialista: non può risolvere la
depressione, ma può alleviarne la parte lavorativa e permettere di ricevere supporto da chi è qualificato. È lo
stesso motivo per cui, a Blizzard e a Riot, quando arrivavano ticket con minacce di suicidio (``mi avete
bannato, quindi mi ammazzo'', ticket purtroppo abbastanza comuni), il \emph{customer service} era istruito a
non intervenire direttamente: non dovevano essere loro a fermare la persona, perché non qualificati, ma
indirizzarla a qualcuno di qualificato, prendendo tempo. Una persona non qualificata che prova a farlo rischia
di peggiorare la situazione; come davanti a un incidente si chiama l'ambulanza e non si fa una tracheotomia.
Per il relatore salute mentale e salute fisica sono la stessa cosa: come non si azzarderebbe a operare qualcuno
di appendicite, non si azzarda a fargli da terapeuta.

\paragraph{Come si danno feedback adatti a ogni tipo di personalità?} Bisognerebbe leggere i risultati del
profilo della persona sullo stesso sito su cui si è fatto il test di personalità della prima giornata, che
spiega anche come interagire con quel profilo. In generale: con chi è più \textbf{analitico}, portare dati a
supporto; con chi è più \textbf{emotivo}, metterlo in una condizione di comfort; con chi è molto
\textbf{pratico}, andare dritti al punto senza girarci intorno. Poi fare pratica. Oppure il modo ``perfetto''
per il relatore: provarci come viene e poi \textbf{chiedere} ``andava bene così o preferisci in un altro modo?
Potevo fare meglio?''. Bisogna creare anche quel rapporto: fidarsi di discutere anche di queste cose.

\paragraph{Se nel team fossimo quasi tutti leader, non ci sarebbero ``troppi galli''?} No, perché per leader si
intende chi si prende cura degli altri, e magari lo fossero tutti. Non significa che tutti scelgano che cosa
fare: spesso lo sceglierà comunque il manager, o guiderà la scelta. Ma avere tutte persone con iniziativa e
capaci di scegliere quando serve è qualcosa per cui il relatore ``metterebbe la firma''. Non ha mai corso il
rischio dei ``troppi galli'': se ci sono quei caratteri, il team è stato assemblato male, perché non tutti sono
compatibili con questo modo di lavorare e magari non si hanno le persone giuste per applicarlo.

\paragraph{Se la fiducia richiede tempo ed esperienza, quanto è difficile per un junior?} Meno di quanto si
pensi: proprio perché su un junior non ci sono aspettative, un junior che chiede ``posso darti una mano?'' dà un
segno chiaro che tiene più al team che a se stesso. Si può essere leader anche non dal punto di vista
professionale, ma come chi va a dare supporto agli altri: ``il medico del team, l'\emph{healer}'', e lo si può
fare anche da junior. È così che ha iniziato il relatore: andare ad aiutare a risolvere problemi non suoi gli ha
fatto accelerare molto il percorso per ottenere rispetto e stima, e lo ha esposto a cose a cui altrimenti non
sarebbe mai arrivato. È una questione di carattere: se ci si tiene agli altri, lo si fa vedere, e il resto
arriva come risultato. ``È quasi un reddito passivo'': un investimento nel far vedere che ci si è, e il resto
arriva (``Stefano si è preso cura di me, quindi ora mi prendo cura di lui'').

\paragraph{Aiutare gli altri solo dopo aver finito il proprio lavoro?} Dipende dall'urgenza. Il relatore spesso,
se qualcosa è davvero urgente, mette in pausa ciò che sta facendo, risolve e poi torna al suo lavoro; ma è troppo
\emph{case by case}. Se due colleghi litigano su chi debba avere le ferie di Natale, lo si può risolvere anche
domani; se litigano perché si è perso il \emph{repository} e domani c'è la consegna, lo si risolve subito.

\paragraph{Come mantenere motivato un team su progetti non retribuiti, come quelli accademici?} Essendo il più
\textbf{trasparenti} possibile nei processi decisionali: se si deve decidere se fare un action o un RPG, tutti
devono sapere perché si fa l'uno invece dell'altro (da qui il metodo di brainstorming della mattina). Deve essere
sempre chiaro a tutti perché si stanno facendo le cose in quel modo, perché \textbf{la demotivazione nasce dal
disallineamento} rispetto all'obiettivo del gruppo, e l'unico modo di essere allineati è essere sempre
consapevoli e partecipi del processo che decide dove sta andando il gruppo. E soprattutto: \textbf{non chiudere
mai un meeting senza aver concordato che cosa si vuole fare} e senza essere tutti d'accordo. Se non lo si è,
bisogna continuare a parlarne, per quanto scomodo, lungo e noioso: si deve sempre uscire da un meeting
sapendo che si è tutti d'accordo, o al massimo in \emph{disagree and commit} (sezione~\ref{sec:z2-disagree}).
``Quello è il massimo accettabile.''

\subsection{Carriera, colloqui e candidature}

\paragraph{Come ha iniziato a lavorare all'estero?} Negli annunci online si trovano spesso requisiti come un titolo
di studio conseguito nel paese o l'abilitazione a lavorarci. Questo vale per i paesi fuori dall'Unione europea: per
lavorare negli Stati Uniti o in Sud America serve un visto o una \emph{green card}. In molti casi, però, è proprio
ottenere il lavoro a permettere di lavorare lì: quando il relatore doveva andare a lavorare in California (a Irvine,
a Santa Monica, dove capitava), era l'azienda a occuparsi dei permessi necessari, perché ``per venire qui devi avere
un lavoro, e il lavoro te lo diamo noi''.

Il suo inizio è stato abbastanza semplice. In Italia lavorava per Digital Bros, occupandosi dei giochi online gestiti
dall'azienda e di quello che era il canale Game Network sul digitale terrestre\footnote{Questo passaggio chiarisce il
dubbio segnalato nel capitolo~\ref{cap:z1-feedback} sul ``Digital Game Network Channel'': si tratta di Game Network,
canale televisivo dedicato ai videogiochi legato al gruppo Digital Bros.}. Vide un annuncio a Dublino per
\emph{Warhammer Online}, mandò la candidatura, fece i colloqui e gli dissero ``ti paghiamo la \emph{relocation}, vieni
qui''. Come detto il giorno prima, sono le aziende a occuparsi del trasferimento: se devono creare a Dublino un'unità
che segue dodici lingue europee, sanno di dover far arrivare persone da fuori, ed è già a budget. Non è nemmeno una
richiesta da fare: è prevista. Per il relatore, da questo punto di vista, è indifferente lavorare per un'azienda
italiana o estera: se si è adatti, si è adatti. Per questo insisteva sull'inglese: se lo si parla, è indifferente di
che cosa ci si occupi e quale sia la propria lingua madre, a meno che la posizione non sia esplicitamente ``community
manager italiano''. Un programmatore polacco o italiano non cambia niente. Lavorare all'estero non è più difficile: c'è
più concorrenza solo se l'azienda cerca attivamente persone da tutto il mondo e non si pone problemi di \emph{location}.
Cambiano l'azienda, il fatto di lavorare da remoto o dal vivo, l'eventuale trasferimento; per il resto non cambia
molto.

\paragraph{Gli anni passati nell'IT contano se si vuole passare ai videogiochi?} Dipende da che cosa si vuole fare: se
da responsabile IT si vuole diventare programmatore, bisogna saper dimostrare di saper programmare. Vale anche per i
passaggi interni: se un game designer vuole diventare art director, al relatore va bene, deve solo dimostrare di saper
fare l'art director. Anzi, incentiva le \textbf{promozioni interne}, perché si conosce già il candidato; ma gli
spostamenti non sono liberi ``a sentimento'': bisogna saper fare quella cosa. I recruiter valuteranno non tanto gli anni
quanto ciò che si sa fare; gli anni servono comunque, perché dimostrano di saper lavorare in un ambiente professionale,
che è un altro insieme di competenze. Ma se si sa fare ciò che chiedono, aver lavorato nell'azienda X o Y non dovrebbe
cambiare molto.

\paragraph{Quanto vengono considerate le candidature spontanee?} Non tantissimo. Soprattutto per le aziende non grandi,
se non stanno cercando qualcuno è perché non se lo possono permettere: una candidatura spontanea non le convincerà ad
assumere. Più l'azienda è piccola, più conviene aspettare l'annuncio, perché se non c'è vuol dire che non c'è budget, e
potrebbero rispondere ``ti vorremmo, ma non possiamo prenderti''. C'è anche il rischio che finisca nel vuoto (l'HR magari
non controlla nemmeno la mail, perché non sta assumendo) o, peggio, che rispondano ``non ci servi'' e una settimana dopo
esca la posizione, con il candidato già nel cestino dei rifiutati. Si rischia di \textbf{bruciarsi} la possibilità.

\paragraph{Come ci si accorge di essere riconosciuti come leader, e come lo si ``vende'' a un colloquio?} Uno studente
osserva che la leadership, così descritta, non la si riconosce da sé ma la riconoscono gli altri. Internamente il
relatore non ha mai dovuto ``venderla'', perché parlavano i risultati. Finché non lo si vive è difficile capirlo, ma
quando si arriva alla posizione in cui il team rispetta e vede qualcuno come punto di riferimento, non ci sono grandi
dubbi: è evidente che vengono da lui. Si può riflettere su \textbf{quanto spesso gli altri vengono da noi}, cercano di
emulare il nostro modo di fare, vanno a rivedere le nostre cose per capire dove hanno sbagliato, chiedono un parere, ci
parlano di un problema: la somma di tutto questo significa molto spesso che si fidano. Per ``venderla'' all'esterno si
può dire di essersi occupati di un team in cui, attraverso il lavoro e il rispetto degli altri, si è arrivati a essere
visti come punto di riferimento; ma il suggerimento è \textbf{portare esempi tangibili} (``c'era questo progetto, non
sapevamo come risolverlo, il team mi ha chiesto una mano e l'ho risolto così''), invece di ``incensarsi'' dicendo ``sono
il punto di riferimento del team''. Se il relatore racconta di essere stato eletto rappresentante dei licenziati senza
essere tra loro, un'azienda capisce che si fidavano di lui, punto: non ci sono altre conclusioni da trarre.

\paragraph{Quali sono gli errori più gravi dei junior ai colloqui, a livello di soft skill?} In ordine di priorità:
\begin{enumerate}
  \item \textbf{non arrivare preparati}, cioè \textbf{senza esempi}. Il relatore ha sempre un esempio pronto per ogni
        cosa di cui vuole parlare: per la fiducia, per gli \emph{achievement}, per i risultati. Servono sempre esempi
        tangibili per far capire che cosa si intende;
  \item \textbf{non avere domande alla fine}, perché dimostra di non avere chiaro che cosa si andrebbe a fare. Anche solo
        chiedere che cosa si farebbe, che cosa si può migliorare, o ``che impressioni hai avuto finora dal mio
        colloquio?'': si possono fare domande molto aperte al recruiter per capire come sta andando;
  \item \textbf{non leggere bene la \emph{job description}}: presentarsi con vaghezza, senza aver capito che cosa cerca
        l'azienda. Lo si impara anche con la pratica; ma il designer che si candida per una posizione da programmatore,
        senza nulla nel CV che faccia pensare che sappia anche programmare, il relatore non solo lo scarta: lo mette in
        una \textbf{blacklist} di persone che forse non considererà mai. Se non si riesce a interpretare un testo come una
        \emph{job description}, è un problema troppo grosso da superare in una collaborazione;
  \item \textbf{errori di battitura e CV impaginati male}, che molti sottovalutano. Il relatore ha letto letteralmente,
        in una lettera di presentazione per una posizione da \emph{copywriter}, la frase ``sono una persona molto attenta
        ai dettaggi''. Come si può considerare quella persona per quel ruolo? Non ha nemmeno incollato la lettera in Word
        per il controllo automatico: non è un errore sfuggito al correttore perché la parola esiste con un altro senso,
        ``dettaggi'' viene segnato in rosso in qualsiasi lingua. Come affidarle la comunicazione aziendale? Vale per tutti
        gli errori: ciò che si manda è \textbf{il primo biglietto da visita} con l'azienda, e non si manda un biglietto da
        visita sgrammaticato. Un conto è non andare ai colloqui in giacca e cravatta (``non l'ho mai fatto e probabilmente
        non lo farò mai''), un altro è presentarsi con errori;
  \item un altro esempio comune: \textbf{l'email sbagliata nel CV}. Se la mail torna indietro, il relatore non telefona
        per farsi dare quella corretta: se non si è in grado di scrivere giusta la propria email, non vale la pena. Più
        l'errore è banale e facilmente evitabile, più lo fa arrabbiare (``scusate se mi sono inalberato''): evidentemente
        il candidato non teneva abbastanza a quel posto.
\end{enumerate}

\paragraph{Quanto è importante la lettera di presentazione?} Sempre di più: sempre più aziende la richiedono come
obbligatoria, quindi conviene prepararne una, sapendo di doverla adattare al destinatario. Deve essere \textbf{meno lunga
di una pagina}, in \textbf{tre paragrafi}:
\begin{enumerate}
  \item \textbf{chi siete}: un minimo di \emph{backstory}, che cosa avete fatto, che cosa volete fare, da dove arrivate;
  \item \textbf{perché avete mandato il CV}: che cosa vi ha interessato di questa azienda;
  \item \textbf{perché dovrebbero chiamare voi}: che cosa dovrebbe spingere l'azienda a scegliere proprio voi.
\end{enumerate}
Uno studente dice di riuscire bene con i primi due, ma di trovare difficile il terzo, che è poi la domanda che tutte le
aziende fanno. In realtà basta \textbf{collegarsi alla \emph{job description}}: se cercano un 3D artist in grado di fare
X, Y e Z, si conferma di averlo letto e si spiega che con la propria esperienza si è imparato a fare proprio quelle cose.
Non vogliono sentire niente di particolare, ``non devi cambiargli il mondo''.

\paragraph{Che cosa guarda un recruiter in un portfolio?} Il relatore li guarda poco: fa un \emph{sanity check} generale,
per vedere che non sia evidente perfino a lui, ``a digiuno di tante cose tecniche'', che c'è un problema (se se ne accorge
perfino lui, il problema c'è). Per il resto guarda che il portfolio sia \textbf{coerente con il CV}: se arriva un CV da
senior con dodici progetti completati e anni di lavoro, e il portfolio sono quattro immagini, non è coerente. Vuole una
concordanza tra ciò che si racconta e ciò che si vede.

\paragraph{Che strumenti per creare un CV?} \textbf{Canva}, che ha un catalogo di template gratuiti e a pagamento
amplissimo e personalizzabile: anche se non piace, lo si prende e lo si adatta, ma fa già il 90\% del lavoro. Insieme a
\textbf{ChatGPT}, perché ``è inutile far finta che non esista'': ci si fa scrivere la lettera di presentazione e poi la si
corregge e la si adatta, e lo stesso per il CV. L'importante è renderli adatti e coerenti, non lasciarli solo frutto di
ChatGPT; ma per fare il grosso del lavoro è un ottimo strumento.

\paragraph{Conviene lasciare nel CV esperienze non del settore?} Uno studente ha alle spalle il liceo, vari stage e
lavori che non c'entrano niente con la modellazione 3D. Si possono mettere finché c'è spazio: il relatore consiglia
\textbf{al massimo due pagine}, e quando ha cominciato a sforare nella terza ha iniziato a togliere ciò che non era più
rilevante. Ormai nel suo CV la formazione è un paragrafo: ha così tante esperienze lavorative che non interessa a nessuno
che abbia fatto il liceo scientifico o classico, e non lo mette nemmeno più. Si filtra in base a ciò che è più rilevante
per lo spazio disponibile.

\paragraph{Che cosa fare con le aziende che dicono ``ti faremo sapere'' e poi spariscono?} Non aspettare uno o due anni
per scrivere. Il suggerimento: tra le domande finali del colloquio, l'ultima del relatore è sempre \textbf{``tra quanto
mi fate sapere?''}. Se rispondono dieci giorni, all'undicesimo scrive: ``mi avete detto dieci''. Non si sente in difetto a
mettere pressione, perché sono stati loro a dare la scadenza. E se viene ignorato del tutto, quell'azienda per lui è
``morta'': non si ricandiderà mai, perché non vede perché il rispetto debba andare in una sola direzione.

\paragraph{Quanto essere trasparenti, nel CV e nella lettera, su problemi personali come l'ADHD?} Uno studente con
ADHD\footnote{Nella trascrizione ``essendo 10D''; nel resto della risposta ``di'' e ``dei 10D'' stanno per ADHD.} dice di
tendere ad avere performance altalenanti. Il relatore lo invita a \textbf{non usare l'ADHD come scusa} per essere
altalenante: anche lui ha l'ADHD, ma non è altalenante. Esserlo è senz'altro un \emph{malus}, ma esistono i mezzi per non
avere problemi di questo tipo: lui ha trovato i suoi metodi (la to-do list, la \emph{must do list}, il delegare il lavoro
al team, perché occuparsi delle cose per gli altri lo mantiene molto motivato, sveglio, all'erta, attento a come vanno le
cose). Sta anche a ciascuno trovare un lavoro in cui compensare ciò che viene richiesto: se domani lo rimettessero a fare
il programmatore, come all'inizio, probabilmente impazzirebbe, perché non riusciva a restare concentrato sul codice oltre
un certo tempo. Lì ha capito che non poteva fare il programmatore, ma ha trovato qualcos'altro che gli permette di essere
com'è e fare ciò che fa. Bisogna capire i propri punti di forza e di debolezza e giocare su quelli: si possono avere tutti
i difetti del mondo e voler comunque fare una cosa, ma bisogna accettare che sarà più difficile.

Quanto alla trasparenza, il consiglio è \textbf{non mentire mai nel CV}, ``perché tanto vi beccano''. E per il discorso della
\emph{bandwidth}: tutte le energie spese a mentire, e a tenere viva la bugia, sono energie tolte allo spiegare bene ciò che
si sa fare. È un investimento nel posto sbagliato. Bisogna concentrarsi su ciò che si sa fare, con esempi tangibili, e dire
``non lo so fare'' quando non lo si sa fare. Il peggio non è essere scoperti, ma \textbf{non esserlo}: se si finge di sapere
il C++, si viene assunti e messi a lavorare in C++, che cosa si fa il giorno dopo? Si rischia davvero di rovinarsi la
carriera. Se non si è un \emph{fit}, si sarà il \emph{fit} per la prossima. L'incubo del relatore è essere assunto in un
posto dove pensano che sappia fare qualcosa che non sa fare; anzi, tende a esplicitarlo: ``cercate questa cosa? Io non la so
fare, va bene lo stesso?'' --- ``Sì, te la insegniamo noi'' --- ``Voglio solo che sappiate che non la so fare''. Meglio
questo che vivere con l'ansia che prima o poi si accorgano che non sa programmare.

Gli capita spesso con chi mente sull'inglese. Al primo colloquio chiede ``sai parlare inglese? A che livello?'' --- ``C1,
\emph{fluent}, quasi madrelingua'' --- ``Perfetto'', e lascia cadere l'argomento. All'ultimo colloquio: ``ti va se lo facciamo
in inglese?'' --- ``No, ma io intendevo che\dots'' --- ``No, mi avevi detto \emph{fluent}: ora facciamo il colloquio in
inglese''. Prima o poi le bugie raggiungono chi le dice.

% ---------------------------------------------------------------------
\section{Chiusura del workshop}
\label{sec:z2-chiusura}

Il relatore chiede alla classe di far sapere com'è andata a Vanessa, il punto di contatto con l'accademia per il workshop:
sarà lei a riferirgli eventuali lamentele o suggerimenti, che gli permetteranno di fare meglio i prossimi. Ringrazia per aver
``retto per dodici ore'' e spera che sia stato utile, consapevole che tutti sono stanchi quanto lui.

Gli studenti rispondono che è stato chiarificatore su molte cose che non sapevano, soprattutto su come interfacciarsi nei
team, e che sarà materiale da rivedere per fissare bene i concetti, vista la mole in due giorni; qualcuno si aspettava che
sarebbe stato molto più pesante e l'ha trovato invece scorrevole. Il relatore ribadisce che \textbf{niente sostituirà
l'esperienza diretta}: il workshop dà un'infarinatura e un punto di vista, ma i concetti si metteranno in pratica solo nel
\emph{day to day}. Bisogna avere \textbf{pazienza con se stessi}, sapendo che non verrà tutto bene dal primo giorno.
Confessa di essere stato un po' preoccupato per la durata, ma il materiale c'era, e farlo in un giorno avrebbe rischiato di
trasformarlo in una corsa: ``rischiava di essere uno \emph{sprint}, e invece è una maratona''.

Infine ricorda che si può aggiungerlo su LinkedIn, che registrazione e slide (in forma integrale) verranno fatte avere
tramite Vanessa (``dite che vi ho autorizzato io a chiederle'') e che, per qualsiasi dubbio, resta a disposizione. ``Buon
resto del weekend.''

% ---------------------------------------------------------------------
\section*{Punti chiave del capitolo}
\addcontentsline{toc}{section}{Punti chiave del capitolo}

\begin{itemize}
  \item I colloqui sono una \textbf{selezione reciproca}. L'annuncio dice molto sulla cultura aziendale; alla fine del
        colloquio vanno sempre fatte domande (il \emph{day to day}, la ``bacchetta magica'', ``tra quanto mi fate
        sapere?''). Si arriva con una \textbf{cifra minima mentale} e con \textbf{esempi tangibili}; non si mente mai.
  \item Time management non è una lista di cose da fare: si parte dalla \textbf{must do list}, che è anche il primo
        momento per dire ``oggi potrei non farcela''; si stimano i tempi e si chiede quando qualcosa non è chiaro.
  \item Ogni nuovo task va \textbf{prima aggiunto alla lista}, non eseguito: priorità non è recenza. Le
        \textbf{checklist} fanno risparmiare tempo, anche (e soprattutto) ai più esperti.
  \item \textbf{First contact resolution}: raccogliere tutto ciò che serve in una sola richiesta, e non iniziare un task
        senza le informazioni necessarie. Ogni interruzione costa dieci-quindici minuti di concentrazione.
  \item \textbf{Matrice di Eisenhower}: fai, pianifica, delega, elimina. ``Non importante'' significa che non è importante
        che lo faccia tu; il quadrante dimenticato è \emph{delete}.
  \item Le pause vanno \textbf{scelte, non subite}; i task grossi vanno spezzettati; bisogna \textbf{saper dire di no},
        anche ai meeting (sono necessario? può andarci uno solo? ci sarà un resoconto scritto? per che cos'altro potrei usare
        questo tempo?). \textbf{Il multitasking non esiste}.
  \item A fine giornata si verifica la \emph{must do list}: un \emph{must do} rimandabile non era un \emph{must do}.
  \item Ferie illimitate con un minimo, giornate libere e feedback culture sono, per il relatore, le condizioni per un team
        sano; il benessere è anche mentale.
  \item Il \textbf{brainstorming genera idee, non decide}. Va preparato (topic, obiettivo, vincoli), separa la
        \textbf{generazione} dalla \textbf{discussione} (\emph{rapid ideation} di cinque minuti, \emph{dot voting}
        anonimo) e combatte l'\emph{anchoring} e il dominio degli estroversi.
  \item La decisione si prende con la \textbf{matrice impatto--sforzo}: prima le idee a basso sforzo e alto impatto, poi in
        diagonale. Lo sforzo sono le risorse, l'impatto comprende anche la motivazione; decide chi farà il lavoro.
  \item \textbf{Disagree and commit}: si dichiara il disaccordo, poi ci si impegna al massimo. Se va male ha sbagliato il
        team, non ``ve l'avevo detto''.
  \item \textbf{La leadership non è un titolo}: il rispetto si guadagna. Le carriere vanno divise tra \textbf{manager} e
        \textbf{subject matter expert}, pagati allo stesso modo, così che diventi manager solo chi vuole prendersi cura degli
        altri.
  \item Dopamina, serotonina e ossitocina legano il gruppo; il \textbf{cortisolo} è contagioso e lo disgrega. Il
        ``contratto sociale'' del villaggio vale ancora: chi protegge il gruppo riceve di più, chi ne abusa tradisce.
  \item La \textbf{piramide della fiducia}: integrità (aiutare davvero, onestà, ammettere gli errori), affidabilità
        (mantenere la parola, saper rifiutare), leadership (dare l'esempio, offrirsi anche quando si rischia, non lasciare
        indietro nessuno).
  \item Il gruppo deve prendersi cura del gruppo: \textbf{raccogliere il piattino}, smettere di dire ``ci penserà qualcun
        altro'', segnalare e chiedere prima di intervenire.
  \item Da manager: delegare i meriti e prendersi le colpe, lasciare la comfort zone, chiedere ``tu come faresti?'',
        celebrare successi e \textbf{fallimenti}, dare il merito a chi lo ha. Essere leader è un \textbf{impegno a lungo
        termine}.
  \item L'obiettivo: tornare a casa \textbf{ispirati, sicuri e realizzati}, come diritto di tutti e non privilegio di pochi.
\end{itemize}

%    % =====================================================================
%  Lezione zero 3 — Agile e Scrum nella pratica
%  Fonte: trascrizione della terza "lezione zero" (GD0_3): una lezione
%  introdotta da Diego Ricchiuti e tenuta da due ospiti, Simone Cicchetti
%  (ex studente, oggi level designer e producer) e Cristian Liberatore
%  (studente del terzo anno, game designer e project manager), che
%  raccontano come hanno applicato davvero Agile e Scrum nei loro progetti.
%  Pacchetti richiesti nel preambolo del libro:
%    \usepackage[italian]{babel}, \usepackage[T1]{fontenc},
%    \usepackage[utf8]{inputenc} (solo pdfLaTeX), \usepackage{booktabs}
%
%  Collocazione nel libro: parte "Lezioni zero", subito dopo
%  GD0_2_tempo_idee_e_leadership, con la stessa numerazione propria (0.3).
%  Nel file principale del libro, ad esempio:
%
%    \part{Lezioni zero: il game designer nell'industria}
%    \renewcommand{\thechapter}{0.\arabic{chapter}}
%    \setcounter{chapter}{0}
%    \input{Seminari/GD0_1_cultura_del_feedback}
%    \input{Seminari/GD0_2_tempo_idee_e_leadership}
%    \input{Seminari/GD0_3_agile_e_scrum_nella_pratica}
%    \part{Il corso di Game Design}
%    \renewcommand{\thechapter}{\arabic{chapter}}
%    \setcounter{chapter}{0}
%    \input{GD1/GD1_introduzione_al_game_design}
%    % ...
%
%  (Con hyperref aggiungere anche \renewcommand{\theHchapter}{z.\arabic{chapter}}
%  nella parte zero, per evitare ancore duplicate.)
%  Tutte le etichette di questo capitolo hanno il prefisso "z3:"; i
%  rimandi ai capitoli precedenti usano le etichette "z1:" e "z2:" già
%  definite in GD0_1 e GD0_2, e alcune etichette dei capitoli del corso
%  (GD1, GD3).
% =====================================================================

\chapter[Agile e Scrum nella pratica]{Agile e Scrum nella pratica: l'esperienza di chi ci è passato}
\label{cap:z3-agile-scrum}

Le prime due lezioni zero (capitoli~\ref{cap:z1-feedback}
e~\ref{cap:z2-tempo-idee-leadership}) guardavano al lavoro in team dal punto
di vista di un professionista con quindici anni di industria alle spalle.
Questa terza lezione cambia prospettiva: a parlare non è un veterano, ma
\textbf{due persone che hanno fatto lo stesso percorso di studi} di chi
ascolta e che, meglio di chiunque altro fino a quel momento, sono riuscite a
\emph{internalizzare} Agile e Scrum nei progetti dell'accademia.

La lezione è introdotta e chiusa da Diego Ricchiuti\footnote{Nella
trascrizione compare soltanto il nome di battesimo, ``Diego'', con cui i due
ospiti si rivolgono più volte a lui; il cognome è quello indicato per il
direttore dell'accademia nel capitolo~\ref{cap:z1-feedback}. Dal contesto è
anche il docente che ha già spiegato Scrum e Agile alla classe nelle lezioni
precedenti (``Diego vi avrà già ampiamente indirizzato'', ``come avete fatto
con Diego nelle lezioni scorse'').}, che ha già spiegato alla classe la teoria
di Agile e Scrum e che interviene di tanto in tanto per sottolineare o
completare un passaggio. Gli ospiti sono due:
\begin{itemize}
  \item \textbf{Simone Cicchetti}, ex studente dell'accademia, oggi level
        designer e producer, che parla per circa un'ora e mezza della teoria
        ``vista dalla vita reale'' e degli errori più comuni;
  \item \textbf{Cristian Liberatore}\footnote{Nella trascrizione il nome
        compare sia come ``Cristian'' sia come ``Christian''; qui si usa la
        prima forma, quella con cui viene presentato all'inizio della
        lezione. Anche i nomi dell'accademia sono resi in modi diversi
        (``Hive'', ``a dive'', ``IVE Online'', ``Live Online''): come negli
        altri capitoli si usa semplicemente ``l'accademia''; ``IVE Online''
        sembra indicare il corso a distanza dell'accademia.}, studente del
        terzo anno, game designer e project manager di due progetti nati
        come progetti d'anno, che segue passo per passo il ciclo di Scrum
        con gli esempi del suo gioco.
\end{itemize}
Diego li ringrazia per il tempo dedicato: Simone lavora, quindi sta
sacrificando tempo personale, e Cristian è uno studente con lo stesso carico
di lavoro della classe. Sono stati scelti perché, tra gli studenti avuti
fino a quel momento, sono quelli che hanno davvero \textbf{internalizzato}
Scrum e Agile, ottenendo grandi risultati produttivi nel progetto scolastico
e, nel caso di Simone, anche dopo l'accademia\footnote{La trascrizione dice
``nel caso di Christian post-accademia'', ma è Simone l'ex studente che oggi
lavora, mentre Cristian è ancora iscritto: il riferimento al dopo-accademia
è stato quindi attribuito a Simone.}.

Perché insistere sulla parola \emph{internalizzare}? Perché, dice Diego,
``utilizzare Scrum e Agile sembra molto facile, ma poi in realtà quando li
provi ad applicare veramente ti rendi conto che facili come sembravano quando
facciamo lezione non lo sono''. Se invece li si internalizza davvero, si
ottengono grandi risultati. L'idea della lezione è che possa essere utile
vedere l'esperienza di qualcuno \emph{molto più vicino} agli studenti di
quanto lo sia il docente, che ormai usa questi strumenti da anni. Il formato
annunciato è di circa un'ora e mezza per ciascun ospite, domande comprese.

Due avvertenze di lettura. La prima: i due interventi sono fatti di teoria
richiamata in breve e di moltissimi esempi concreti, spesso con numeri,
errori e ripensamenti; sono riportati con cura perché sono proprio la parte
che la teoria non dà. La seconda: i due relatori usano in modo
intercambiabile alcuni termini (\emph{framework} e metodologia,
\emph{post-mortem} e revisione della sprint, \emph{score}, peso e punteggio
dei task), e il capitolo segue il loro uso, segnalando in nota dove si
discosta dalla terminologia ufficiale di Scrum.

Il percorso del capitolo segue quello della lezione:
\begin{enumerate}
  \item un'apertura organizzativa e la presentazione degli ospiti
        (sezione~\ref{sec:z3-apertura});
  \item \textbf{Simone}: chi è, un sondaggio sulla classe e le prime
        esperienze di Scrum degli studenti (sezione~\ref{sec:z3-simone});
  \item il \textbf{mindset Agile}: perché non è un framework e perché va
        adattato (sezione~\ref{sec:z3-mindset});
  \item le \textbf{sprint} e le riunioni di pianificazione e
        \emph{post-mortem} (sezione~\ref{sec:z3-sprint});
  \item le \textbf{cerimonie} e la regola della \emph{consistency}, con le
        critiche al \emph{daily scrum} e al \emph{planning poker}
        (sezione~\ref{sec:z3-cerimonie});
  \item il \textbf{backlog} e gli obiettivi dei task
        (sezione~\ref{sec:z3-backlog-simone});
  \item la \textbf{velocity}: perché non usare le ore, la sequenza di
        Fibonacci e come la velocity permette di ``prevedere il futuro''
        (sezione~\ref{sec:z3-velocity});
  \item il riepilogo della creazione di una sprint e la conclusione sulla
        \textbf{fiducia} (sezione~\ref{sec:z3-ricapitolo});
  \item le domande a Simone, con la costruzione in diretta di un foglio per
        il calcolo della velocity (sezione~\ref{sec:z3-qa-simone});
  \item \textbf{Cristian}: chi è e come imposta la sua parte
        (sezione~\ref{sec:z3-cristian});
  \item il caso di studio, il gioco \emph{L'Inf}, e il ``trauma'' del cambio
        di team (sezione~\ref{sec:z3-linf});
  \item il \textbf{loop di Scrum} in cinque fasi
        (sezione~\ref{sec:z3-loop}): la stesura del backlog
        (sezione~\ref{sec:z3-backlog}), la selezione dei task
        (sezione~\ref{sec:z3-selezione}), l'implementazione con i suoi
        controlli, i \emph{blocker} e i feedback
        (sezione~\ref{sec:z3-implementazione}), il calcolo della velocity
        (sezione~\ref{sec:z3-calcolo}) e il \emph{tuning} della produzione,
        con il problema del perfezionismo (sezione~\ref{sec:z3-tuning});
  \item una domanda su come comunicare un taglio al team
        (sezione~\ref{sec:z3-qa-taglio});
  \item il \textbf{team management}: il team prima di Scrum, la produzione
        senza stress, il mindset condiviso, il \emph{post-mortem} e la
        \textbf{sincerità} (sezione~\ref{sec:z3-team});
  \item la chiusura della lezione (sezione~\ref{sec:z3-chiusura}).
\end{enumerate}

% ---------------------------------------------------------------------
\section{Apertura e presentazione degli ospiti}
\label{sec:z3-apertura}

\subsection{Una registrazione andata persa}

La lezione si apre con una comunicazione organizzativa\footnote{La
trascrizione comincia a frase già iniziata (``in una modalità non video fin
tanto che ce la chiamate poi la converte alla fine''): l'inizio della
spiegazione non è stato registrato e non è chiaro a quale lezione precedente
si riferisca.}. Il programma usato per registrare una lezione precedente
registra in una modalità che va convertita alla fine; durante la conversione
il computer di Diego, probabilmente per colpa degli aggiornamenti di Windows,
si è bloccato e la conversione è ``andata a pallino''. Poiché si tratta di
una lezione importante, Diego si rende disponibile a \textbf{rifarla
interamente} per registrarla di nuovo oppure, come la classe preferisce, a
fornire un audio, un video o qualunque altro formato sia più comodo, che poi
caricherà. Chiede agli studenti di pensarci e di fargli sapere.

\subsection{Chi sono gli ospiti e perché sono qui}

Poi Diego presenta gli ospiti speciali della giornata: Cristian Liberatore,
studente del terzo anno, e Simone Cicchetti, ex studente dell'accademia (che
in un primo momento ``sparisce'' dalla griglia della videochiamata, finché
Diego non lo ritrova in alto a destra). Le ragioni della scelta, il valore
della parola \emph{internalizzare} e il formato (un'ora e mezza ciascuno, con
domande) sono quelli già riassunti all'inizio del capitolo. Detto questo,
Diego si silenzia e lascia la parola a Simone.

% ---------------------------------------------------------------------
\section{Simone: il punto di vista del producer}
\label{sec:z3-simone}

\subsection{Chi è il relatore}

Simone si presenta mentre condivide lo schermo (``adesso tutto è imploso,
perfetto''). È uno degli ex studenti del terzo anno e oggi lavora come
\textbf{level designer}; collabora con Diego su alcuni progetti legati alla
gestione dei team, al \emph{game design}, al \emph{framework design} e ad
altre ``cose molto fighe'' che alcuni studenti forse conoscono\footnote{Il
passaggio è confuso nella trascrizione (``lavoro con diego alcune alcuni
progetti diciamo sul gestione di pro dei team game design framework
design''): la ricostruzione è approssimativa. Più avanti Simone dice anche di
lavorare per una società australiana.}. Sta inoltre continuando a lavorare al
progetto del suo terzo anno (il \emph{Biggest One}\footnote{Dalla lezione
si capisce che i progetti d'anno dell'accademia si chiamano \emph{Big One}
(primo anno), \emph{Bigger One} (secondo anno) e \emph{Biggest One} (terzo
anno); la trascrizione li scrive in vari modi (``biggest one'', ``bigger
one'', ``BiggerOne'').}), un gioco chiamato \emph{ISO}\footnote{Il nome del
gioco compare sempre come ``ISO''; non è possibile stabilire dalla sola
trascrizione se sia la grafia corretta (potrebbe essere, ad esempio,
\emph{Iso} o una sigla).}, per il quale è level designer, game designer,
social media manager e soprattutto \textbf{producer}: si occupa di gestire il
lavoro delle persone, di far rispettare le cerimonie di Scrum e, ``punto
centrale della questione'', di ``rompere le scatole alla gente''.

\subsection{Una lezione partecipativa}

Il tempo a disposizione è molto, quindi Simone vuole una lezione il più
possibile \textbf{partecipativa e cooperativa}: a differenza del solito, le
domande vanno fatte in qualsiasi momento, scrivendo in chat per attirare la
sua attenzione o alzando la ``manina'' con le emoji, che controlla su un
secondo monitor con la griglia dei partecipanti. Al massimo aspetterà la fine
della slide per non perdere il filo.

L'obiettivo dichiarato non è spiegare la teoria, che la classe ha già fatto
con Diego, ma darne un breve \emph{recap} e poi spiegare \textbf{a che cosa
serve nel pratico e com'è nella vita reale}: ``come spesso accade, la teoria
è una cosa, la vita reale è un'altra''.

\subsection{Un piccolo sondaggio}

Prima di partire, Simone chiede chi abbia ``un amore incommensurabile e
infinito'' per il producing, lo \emph{scrum mastering} e il management in
generale: nessuno, e uno studente, Gavi\footnote{Così nella trascrizione;
potrebbe trattarsi di un soprannome o di un nome storpiato.}, risponde
spammando il ``no''. Poi la domanda opposta: chi lo odia così tanto da
preferire ``affrontare un alligatore in una lotta greco-romana''? Ancora
nessuno. La classe è quindi, in generale, ``tutta intermedia''.

Chi si è già dovuto occupare di queste cose nei progetti d'anno o nei lavori
di gruppo? Rispondono due studenti, Mattia e Lorenzo\footnote{Nella
trascrizione ``da noi Mattia, Podini'' e ``da noi Lorenzoaldasso'': i cognomi
non sono ricostruibili con sicurezza (``Podini'' potrebbe essere il cognome
di Mattia o il nome di un altro studente).}, ognuno per il proprio team.
Simone chiede solo i punti salienti: che cosa è andato bene, che cosa no e
perché.

\paragraph{L'esperienza di Mattia.} Il suo gruppo usa Scrum da pochissimo,
circa una settimana. È lui a gestire la \emph{Scrum board}, con l'aiuto degli
altri due membri, ma il metodo non è ancora applicato a pieno. Mattia è
interessato a Scrum, ma la \textbf{responsabilità di guidare altre persone}
oltre a se stesso gli pesa parecchio, e usare Scrum in modo serrato tutti i
giorni gli fa ``una bella paura'': se qualche meccanismo si rompe,
bisognerebbe correre ai ripari. Che cosa intende per ``rompere un
meccanismo''? Che se una cerimonia non viene rispettata, tornare allo
\emph{status quo} non sarebbe semplice, anche per la sua mancanza di
esperienza nell'organizzare i team. Simone riassume: la preoccupazione è più
sul piano umano che su quello delle ore di lavoro necessarie. Mattia
riconosce che Scrum avrà le sue ``quasi magiche proprietà'', ma non avendolo
ancora visto all'opera nella vita reale non può dirlo con certezza.

\paragraph{L'esperienza di Lorenzo.} Niente di diverso: anche il suo gruppo
ha iniziato da pochissimo e proprio il giorno prima ha fatto la prima
riunione di Scrum, decidendo insieme i pesi dei task. Alla domanda se il
ruolo sia arrivato per voto popolare, Lorenzo risponde che si è offerto lui:
conoscere un \emph{framework} di lavoro, e soprattutto capire come
applicarlo, gli interessa parecchio.

Simone trova molto interessante quanto detto da Mattia sulle cerimonie e sul
doverle rispettare, e promette di tornarci più avanti
(sezione~\ref{sec:z3-cerimonie}).

% ---------------------------------------------------------------------
\section{Il mindset Agile}
\label{sec:z3-mindset}

\subsection{Agile non è un framework}

La prima slide parte da una frase che tutti conoscono: \textbf{Agile non è un
framework, è un mindset} (Simone dice di aver ``percepito la furia negli
occhi di Diego'' quando la classe esitava a rispondere). La frase sembra una
di quelle che servono ``per fare i fighi con i colleghi'', ma ha un
significato preciso: un \emph{framework} presuppone \textbf{regole da
seguire}, mentre Agile \textbf{non ha regole}. È una serie di \textbf{linee
guida da internalizzare}: la parola chiave è che devono diventare
\emph{interne}, al punto da non dover più pensare attivamente a quali siano
le regole, i punti o la mentalità di Agile. Bisogna capire la logica che c'è
sotto.

L'analogia è quella della \textbf{lingua straniera}: all'inizio si ragiona in
italiano, si costruisce la frase in italiano, la si traduce e la si dice; man
mano questi passaggi si annullano e si comincia a ragionare direttamente
nella nuova lingua. Con Agile funziona allo stesso modo.

\subsection{I quattro valori, e perché non vanno presi alla lettera}

Simone rilegge rapidamente i quattro valori del manifesto Agile:
\begin{enumerate}
  \item gli \textbf{individui e le interazioni} sono più importanti dei
        processi e degli strumenti;
  \item un \textbf{prodotto funzionante} è più importante di una
        documentazione esaustiva\footnote{Nella lezione ``documentazione
        comprensibile''; nella formulazione corrente del manifesto Agile il
        termine è \emph{comprehensive documentation}, cioè documentazione
        esaustiva.};
  \item la \textbf{collaborazione con il cliente} è più importante della
        negoziazione del contratto;
  \item \textbf{rispondere al cambiamento} è più importante di seguire un
        piano.
\end{enumerate}
Come è scritto sotto al manifesto stesso, questo non significa che le cose a
destra, quelle ``meno importanti'', non contino o non vadano seguite. Ma
Simone va oltre: \textbf{non significa nemmeno che queste regole debbano
essere seguite davvero}. Non è detto che individui e interazioni debbano
essere per forza più importanti di processi e strumenti. Il punto chiave di
Agile è \emph{sentire} quale sia il procedimento e la mentalità giusti e
\textbf{adattare Agile al proprio team}, perché ogni team e ogni progetto è
diverso.

Se un team e un progetto, per vari motivi, hanno bisogno di pagine e pagine
di documentazione, allora per loro la documentazione è più importante del
prodotto funzionante. Gli esempi sono infiniti; uno è il \textbf{gioco di
carte}: lì il prodotto funzionante è relativamente meno importante della
documentazione, perché le logiche si possono scrivere in qualsiasi momento,
ma senza una documentazione solida a supporto ``cadrà tutto quanto
immediatamente''. Quindi non bisogna seguire ``a menadito'' i punti del
manifesto, ma capirne la logica di fondo, capire \textbf{a che cosa serve} e
capire \textbf{se serve}: non è detto che Agile serva.

\subsection{Altri mindset e i framework}

Esistono molti altri \emph{mindset}. Di recente va molto di moda la
mentalità \textbf{lean}, che Simone definisce ``già il 2.0''\footnote{Così
nella lezione. Storicamente il \emph{lean} nasce nella produzione
industriale (il sistema Toyota) e precede Agile; nello sviluppo software il
\emph{lean software development} si è diffuso negli anni Duemila. Il ``2.0''
va letto come una battuta sul fatto che oggi è di moda, non come una
derivazione da Agile.}. L'importante, aggiunge con Diego, è \textbf{non usare
Waterfall}, ``perché Waterfall nel gaming\dots{} non dirò nient'altro''.

Sulla base di Agile sono state create delle \textbf{metodologie}, che
Simone si corregge chiamando \emph{framework}: e queste sì che hanno delle
regole. L'esempio è Scrum, ma ce ne sono infiniti, come Kanban. Si basano
tutti sul principio Agile dell'\textbf{inspect and adapt} (ispeziona e
adattati, analizza e adattati), che si può riassumere in una parola chiave:
\textbf{iterazioni}. Scrum è completamente basato sulle iterazioni, e anche lo
sviluppo dei videogiochi lo è, ``molto più dello sviluppo software''.

Il resto della sua parte è dedicato ai punti di Scrum, il \emph{framework} da
cui presumibilmente partono tutti gli studenti, e agli \textbf{errori
classici} in cui si incappa spesso. L'ideale sarebbe trasmettere la logica
alla base di tutto, ma Simone avverte che sarà difficile, perché ``sono cose
che si iniziano a capire mettendole in pratica''. Per questo invita di nuovo a
fare domande sulle proprie esperienze.

% ---------------------------------------------------------------------
\section{Le sprint e le riunioni}
\label{sec:z3-sprint}

\subsection{Il libro di Scrum}

Simone mostra il libro in cui viene raccontata la genesi del
\emph{framework} Scrum, scritto da Jeff Sutherland\footnote{Si tratta
verosimilmente di J.~Sutherland, \emph{Scrum: The Art of Doing Twice the Work
in Half the Time} (2014), edito in italiano come \emph{Fare il doppio in metà
tempo}. Il titolo non viene pronunciato: Simone mostra la copertina e più
avanti cita ``Scrum, di Jeff Sutherland''. La lunghezza indicata (``circa 300
pagine'') è approssimativa.}: un ``libricino'' di circa trecento pagine.
Avverte che userà le parole ``metodologia'' e ``framework'' in modo
intercambiabile. Lo consiglia assolutamente a chi è interessato al tema,
perché è scorrevole e contiene \emph{insight} molto interessanti; chi non è
particolarmente interessato può evitarlo del tutto, perché non avrà niente da
offrire che non venga detto in questa lezione.

\subsection{Che cos'è una sprint}

Scrum, a differenza di Agile, è un \emph{framework} a tutti gli effetti: ci
sono regole da seguire in modo preciso e da rispettare. Il modo in cui
abbraccia il principio dell'\emph{inspect and adapt} sono le
\textbf{sprint}: periodi di tempo tipicamente di \textbf{una o due
settimane}. Possono essere anche di quattro, come si studia, ma secondo Simone
non ha senso farle così lunghe a meno di non lavorare a un grande progetto
``con duecento miliardi di persone che ci stanno lavorando
costantemente''\footnote{Nella trascrizione ``a meno che non state facendo un
sesto plà con 200 miliardi di persone'': ``sesto plà'' è incomprensibile e
potrebbe essere una storpiatura di ``tripla A''; il senso, in ogni caso, è
quello di una produzione enorme.}.

Prima di ogni sprint c'è una riunione di \textbf{sprint planning}, dopo ogni
sprint una riunione di \textbf{post-mortem}\footnote{Nella guida ufficiale di
Scrum la riunione di fine sprint è divisa in \emph{sprint review} (che cosa è
stato prodotto) e \emph{sprint retrospective} (come ha lavorato il team);
qui, come nel resto della lezione, entrambe confluiscono nel termine
\emph{post-mortem}, che in senso stretto indica l'analisi a progetto
concluso (sezione~\ref{sec:z1-qa-fasi}).}.

\subsection{La regola d'oro: meno riunioni possibili}

``Le riunioni non piacciono a nessuno'': sono, dice Simone, una ``rottura di
scatole infinita'', e chiunque, soprattutto chi non è designer, non aspetta
altro che poter dire ``questa riunione è inutile, dobbiamo andarcene tutti a
lavorare perché abbiamo cose più importanti da fare''. In chat qualcuno
risponde con il meme di Bernie Sanders a braccia incrociate; Simone avrebbe
voluto mettere quello del cane che urla sulla scrivania perché ``tutta la
riunione poteva essere una mail'', e promette che dalla prossima volta ogni
slide avrà un meme alla fine.

Da qui la sua \textbf{regola d'oro}: \textbf{fare meno riunioni possibile}.
Sottolinea però che è una regola ``per quanto mi riguarda'': magari un team
non vede l'ora di fare riunioni, magari si portano le pizzette a ogni
incontro (``non so se con le pizzette le riunioni siano effettivamente
efficienti, però se funziona, perché no?'').

Una mossa molto semplice per ridurle è \textbf{unire sprint planning e
post-mortem} in un'unica riunione: si fa letteralmente la metà delle
riunioni. Ciò che si fa in quella riunione è:
\begin{itemize}
  \item \textbf{pianificare} la prossima sprint (\emph{sprint planning});
  \item \textbf{analizzare} la sprint precedente (\emph{post-mortem}),
        indicando in particolare i problemi incontrati e le cose rimaste
        bloccate.
\end{itemize}

\subsection{Una riunione deve essere rapida}

La riunione unica deve essere \textbf{rapida}. Durante l'ultimo anno, per il
\emph{Biggest One}, il team di Simone contava circa quindici persone, e alle
riunioni di planning e post-mortem ne partecipavano in media una dozzina.
Quando sono ``diventati bravi'', le riunioni duravano \textbf{venticinque
minuti, con dodici persone che parlavano tutte}. Quello è l'obiettivo, ``anzi,
il vostro obiettivo è fare anche meno''.

Il motivo per cui si sfora è una \textbf{tentazione} costante: quella di
\textbf{risolvere i problemi}. Quando qualcuno descrive un problema, un paio
di persone cominciano a discuterne (``hai provato a fare questo?'', ``no,
dovresti fare quest'altro''), e in un attimo sono passati dieci minuti in cui
hanno parlato solo in due senza che la riunione sia andata avanti di un
passo. In queste riunioni non si fa brainstorming e non si ragiona su come
risolvere i problemi: ``bloccatelo sul nascere, picchiate quelle persone, fate
qualsiasi cosa per impedirglielo''.

Quando c'è un problema, si \textbf{prende un appuntamento per risolverlo}.
Simone mostra come dovrebbe suonare un report:
\begin{quote}
``La settimana scorsa mi sono occupato della logica di movimento del
personaggio, ho fatto un po' di tuning e di testing e ho preparato il block
out del primo livello. L'unico problema è che il blueprint che faceva volare
via gli uccelli non mi ha funzionato bene: visto che tu, Cristian, l'hai
utilizzato in un altro livello, più tardi ci sentiamo su Discord quindici
minuti e mi spieghi come l'hai settato. Questa settimana mi occuperò di
completare il block out e di risolvere il problema del blueprint.''
\end{quote}
Fine del report: \textbf{circa un minuto}, e c'è tutto. Tutti devono fare
così: trovare una persona con cui risolvere il problema, identificare
\emph{come} risolverlo e andare oltre. ``C'è tempo e luogo per ogni cosa'', ma
la riunione di sprint planning non è il momento giusto.

\subsection{La sprint come unità di tempo}

Ultimo punto: le sprint sono \textbf{l'unità di tempo} fondamentale per ogni
cosa. Le \emph{milestone} non si misurano in mesi, giorni o settimane, ma
\textbf{in sprint}. Bisogna prendere l'abitudine di dire ``completeremo la
preproduzione in otto sprint'' e non ``completeremo la preproduzione tra due
mesi''. È importantissimo per il calcolo della velocity
(sezione~\ref{sec:z3-velocity}).

% ---------------------------------------------------------------------
\section{Le cerimonie e la regola della consistency}
\label{sec:z3-cerimonie}

\subsection{Le cerimonie di Scrum secondo Simone}

Le regole di Scrum si chiamano \textbf{cerimonie}. Simone le divide in due
gruppi\footnote{La classificazione è quella proposta da Simone. Nella guida
ufficiale di Scrum gli ``eventi'' sono la sprint, lo \emph{sprint planning},
il \emph{daily scrum}, la \emph{sprint review} e la \emph{sprint
retrospective}; il backlog è un ``artefatto'' e il calcolo della velocity
e il \emph{planning poker} sono pratiche diffuse ma non obbligatorie.}:
\begin{center}
  \begin{tabular}{@{}p{0.42\textwidth}p{0.48\textwidth}@{}}
    \toprule
    \textbf{Cerimonie principali} & \textbf{Altre cerimonie} \\
    \midrule
    planning; post-mortem; calcolo della velocity; creazione del backlog &
    \emph{daily scrum}; \emph{poker time estimation} (\emph{planning poker})\\
    \bottomrule
  \end{tabular}
\end{center}
Le quattro principali sono quelle che ``non potete evitare di fare,
altrimenti non state facendo Scrum'': il che, precisa, non è necessariamente
un problema (ci torna alla fine, sezione~\ref{sec:z3-ricapitolo}). Le altre
due le tiene separate perché a lui non piacciono granché, e prima di spiegare
perché beve un sorso d'acqua ``per lasciarvi la suspense''.

\subsection{Consistency: una cerimonia saltata fa cadere le altre}

In Scrum c'è una parola chiave da ricordare: \textbf{consistency}. Bisogna
evitare a tutti i costi di creare \textbf{precedenti di
\emph{inconsistency}}, cioè di incostanza: con Scrum bisogna rimanere
costanti. \textbf{Se si salta una sola cerimonia si crea un precedente che,
a effetto domino, butta giù tutto il resto.} Questo è il punto già sollevato
da Mattia. Ma, avverte Simone, questo \emph{non} vuol dire che ci si debba
forzare: ed è questo il punto chiave.

\subsection{Perché il daily scrum non funziona più}

Il \emph{daily scrum} era pensato ``per il millennio scorso''. Scrum è nato
molto tempo fa\footnote{Simone non ricorda la data esatta (``fine degli anni
Novanta, forse fine degli anni Ottanta''). Il termine \emph{scrum} applicato
allo sviluppo di prodotti risale a un articolo di Takeuchi e Nonaka del 1986;
il \emph{framework} Scrum fu presentato da Jeff Sutherland e Ken Schwaber nel
1995.}, per team che si incontravano tutti i giorni in ufficio alla stessa
ora, con un orario d'ufficio standard: team di sviluppo software, quindi per
lo più programmatori (anche i più ``artistici'', come gli UX designer, erano
alla fine programmatori), con una mentalità piuttosto simile e, soprattutto,
con lo stesso orario.

Oggi, Covid a parte, questo non è più né fattibile né sensato. Gli studenti
non hanno nemmeno un ufficio: si incontravano a volte in sede quando si
poteva, ma per lo più lavorano online, e \textbf{non tutti lavorano allo
stesso orario}. ``E va bene così'': c'è chi lavora di notte (la salute
potrebbe risentirne, ma se si sentono bene così, va bene), e quelle persone
non scriveranno il proprio \emph{daily scrum} all'orario in cui lo fanno
tutti. Senza arrivare a casi del genere, Simone lavora per una
\textbf{società australiana}, dal lato opposto del globo: quando lui si
sveglia, loro stanno finendo di lavorare, e il suo \emph{daily scrum} cadrebbe
nella loro notte.

In queste condizioni è difficile tenere traccia del \emph{daily scrum} e
rispettarlo, e prima o poi qualcuno lo salterà. Quello è un
\textbf{precedente di incostanza}: qualcuno comincia a dire ``ok, non lo faccio
neanche io'', poi ``ma visto che non faccio più gli scrum, a che servono le
riunioni? Questa volta non vengo'', ``e io neanche'': ``e da lì è il caos''.

La conclusione è netta: \textbf{quando qualcosa non funziona, tagliatelo}. Il
senso di ``individui e interazioni sono più importanti di processi e
strumenti'' è proprio questo: se un processo o uno strumento non funziona
\textbf{non è colpa vostra, è il team che non ne ha bisogno}. Non bisogna
fissarsi su ciò che si vuole fare, ma sentire di che cosa ha bisogno il team e
adattarsi a quel bisogno.

\subsection{Perché il planning poker rallenta tutto}

Il \emph{planning poker} (che Simone chiama \emph{poker time estimation})
non gli piace perché lo trova \textbf{estremamente lungo e tedioso}: rallenta
un processo già lento di suo. Creare il backlog e valutare i task è già
lunghissimo (``ci vuole almeno un'intera sprint di riunioni per fare un
backlog''); aggiungerci il planning poker significa moltiplicare il tempo. Si
informa con Diego su come sia stato spiegato alla classe: con il metodo di
Fibonacci, ``\emph{power-uppato}''\footnote{Nella trascrizione ``è rimasta
nelle lezioni di trial'' e ``è stata power upata'': il senso sembra essere che
il planning poker è stato spiegato nelle lezioni precedenti in una versione
potenziata con la sequenza di Fibonacci. ``Trial'' potrebbe essere una
storpiatura.}, il che farà risparmiare un po' di spiegazioni più avanti.

Simone capisce la necessità di sentire le opinioni di tutti
(``\emph{empowering the team}'', eccetera), ma l'effetto è rendere il processo
ancora più lento. E poiché le riunioni non piacciono a nessuno, se si
quintuplica il tempo necessario ``si parte subito male'', perché il team si
annoia già. Anche questo, quindi, lui personalmente non lo usa.

\subsection{Il punto chiave: crearsi il proprio framework}

La cosa più importante è che \textbf{non bisogna eseguire tutto a menadito}.
I \emph{framework} sono stati pensati per essere generici, ma anche per
essere \textbf{adattati}. La procedura che Simone suggerisce:
\begin{enumerate}
  \item si parte \textbf{rispettando il framework al 100\%};
  \item quando qualcosa non funziona, si capisce \emph{perché} non funziona;
  \item si smette di farlo, si risolve, si fa altro;
  \item si capisce qual è il framework più adatto e \textbf{ci si crea il
        proprio}.
\end{enumerate}
``Questo è il punto chiave di tutto questo.''

% ---------------------------------------------------------------------
\section{Il backlog}
\label{sec:z3-backlog-simone}

\subsection{Una ``rottura di scatole'' inevitabile}

Sulla creazione del backlog Simone spende ``giusto due parole'', perché la
parte è molto tecnica e i consigli sono pochi. Creare il backlog è faticoso,
``non ci sono mezze misure per dirlo'': richiede moltissimo tempo e lavoro ed
è intenso a livello mentale, perché bisogna ragionare tantissimo.

\subsection{Il backlog a scaglioni}

In teoria i backlog si creano all'inizio dello sviluppo: in un mondo perfetto
si conoscerebbero fin da subito tutti i task necessari per completare il
progetto. Nel mondo reale questo non accade. Ciò che accadrà, più
probabilmente, è che si andrà avanti \textbf{per milestone}. È ciò che Simone
ha cominciato a fare con \emph{ISO} e con cui si trova bene:
\begin{itemize}
  \item ha fatto il backlog della \textbf{preproduzione}, con tutto ciò che
        serviva per la preproduzione;
  \item finita la preproduzione, ha fatto quello della \textbf{produzione};
  \item finita la produzione farà quello della fase successiva
        (``espansione, build finale, non lo so''), a seconda di come si
        svilupperanno le cose.
\end{itemize}
Fare i backlog \textbf{a scaglioni} li rende più semplici e il processo meno
lungo; e tanto, comunque, bisognerà aggiungere cose in corso d'opera: ``non si
scappa''.

\subsection{Le epic}

L'organizzazione in \textbf{epic} è molto utile, soprattutto per la stesura:
si comincia ragionando per \emph{epic} e poi si scende ai task. Un esempio di
epic è ``creazione del personaggio''. Che cosa contiene il personaggio?
Movimento, tuning, modello, animazioni, materiale, il design delle feature
che dovrà avere, la creazione di tutte queste feature, e così via. Il
processo diventa più \emph{streamlined}, meno pesante e più semplice da
seguire.

\subsection{Ogni task deve avere un obiettivo chiaro}

Fondamentale: \textbf{i task devono avere un obiettivo chiaro}. Un consiglio
comune è scriverli come \textbf{user story}: ``come utente voglio avere un
personaggio che si senta potente, che io riesca a percepire come potente''.
Aiuta moltissimo nello sviluppo software e parecchio nei videogiochi per
alcune cose, come appunto il personaggio. In generale è buona pratica
\textbf{spiegare l'obiettivo e non la soluzione}: dire che il personaggio deve
muoversi, saltare e fare le capriole, invece di dire che ``quando premo il
tasto W deve subire una forza che lo spinge in avanti''. Però non è
fondamentale.

Ciò che è davvero importante è che il task abbia un \textbf{obiettivo
verificabile}: qualcosa di cui si possa dire ``è completo'' o ``non è
completo''. Altrimenti le persone \textbf{non lo completeranno}, o lo
completeranno in modo diverso da come lo si immaginava. Un esempio di task
senza obiettivo chiaro è \textbf{``completare il livello~5''}:
\begin{itemize}
  \item per Simone vuol dire aver creato il block out con qualche elemento
        estetico per sapere dove vanno le cose e come funzionano con i
        moduli già disponibili, e basta;
  \item per qualcun altro potrebbe voler dire inserire illuminazione,
        \emph{post-processing}, \emph{particles}, tutti gli shader fatti
        bene, ottimizzazione e \emph{level streaming};
  \item per altri ancora potrebbe voler dire arrivare all'illuminazione,
        lasciando \emph{level streaming} e ottimizzazione a una seconda fase.
\end{itemize}
È la cosa più importante da ricordare quando si crea il backlog e si
scrivono i task.

% ---------------------------------------------------------------------
\section{La velocity}
\label{sec:z3-velocity}

``Bellissimo, bellissimo'': il calcolo della velocity è per Simone ``la cosa
più figa e allo stesso tempo più odiata''.

\subsection{Il peso dei task}

Come la classe sa, i task sono valutati in base al loro \textbf{peso} (o
punteggio, o \emph{score}). Il peso dipende da quanto un task è difficile e/o
da quanto tempo richiede: è un po' una zona grigia, e Simone consiglia di
valutarlo come \textbf{combinazione delle due cose}.

\subsection{Mai le ore}

Il consiglio fondamentale, che Simone non capisce perché si continui a
contraddire, è \textbf{non usare le ore per valutare i task}, ``per niente,
neanche alla lontana''. Nel libro di Sutherland c'è un paragrafo intero che
spiega perché. Diego commenta che chi lo consiglia di solito fa ``Agile
Waterfall'' e non ha nemmeno letto i libri su Agile.

Il motivo: \textbf{il cervello umano è completamente incapace di stimare le
ore}, cioè di prevedere il tempo necessario per fare qualcosa: ``non siamo
pensati per questo''. Siamo invece \textbf{molto bravi con le misure
relative}: fissato un task che vale 10, sappiamo dire bene quale vale 5,
quale 7, quale 2. Sappiamo dire che un task è molto più grande di un altro.

\subsection{La sequenza di Fibonacci}

Per aiutarsi ulteriormente, Simone usa la \textbf{sequenza di Fibonacci} (1,
2, 3, 5, 8, 13, 21, 34, \dots), anche questo un consiglio di Sutherland che
``nessuno segue''. Sutherland racconta addirittura di un team che usava le
\textbf{taglie dei cani} (``questo task è un pinscher, questo è un golden
retriever'')\footnote{Nel libro di Sutherland le stime per razze canine sono
citate come esempio di scala relativa alternativa ai numeri; nella
trascrizione ``pincher''.}; a Simone sembra eccessivo, e ``i buoni vecchi
numeri'' bastano. Fibonacci ha due scopi molto utili:
\begin{enumerate}
  \item \textbf{Gli intervalli crescenti assorbono l'errore di stima.} Man
        mano che un task diventa più grande e complesso, cresce anche
        l'errore associato. È facile dire che aprire la mail vale 2 e
        rispondere a tutti sul server Slack vale 3: che cosa può andare
        storto? La mail non si apre e bisogna reinserire la password; è
        esploso qualcosa e ci sono cinquanta messaggi invece dei soliti
        venti. Niente che cambi granché. È molto più difficile, invece,
        valutare la differenza tra creare il \emph{behavior} di movimento del
        personaggio e costruire un intero livello: possono succedere
        infinite cose, soprattutto quando c'è di mezzo la programmazione.
        Il behavior del personaggio è un 20 o un 21? ``È un 21, punto'':
        non è un 34, perché non è così difficile, e non è un 13, perché
        sarebbe troppo poco. Magari la cosa più corretta sarebbe stata 23,
        ma non importa: lo scopo è \textbf{appianare} i risultati.
  \item \textbf{Meno scelta, meno discussione.} I valori sono
        ``infinitamente arbitrari'', ``completamente, al 100\%'', e nessun
        altro al mondo li condividerà; ma Fibonacci li rende un po' più
        condivisibili, perché con meno possibilità di scelta ci sono meno
        occasioni di discutere.
\end{enumerate}
Che i valori siano arbitrari non è un problema: finché \textbf{è una persona
sola a darli}, sono tutti arbitrari allo stesso modo, quindi
\textbf{consistenti}; e se sono consistenti hanno una logica e si possono
usare per fare calcoli.

\subsection{Che cos'è la velocity e perché serve}

Sommando il peso dei task completati in una sprint si ottiene la
\textbf{velocity}. Ed è ``fighissima'' perché permette, in sostanza, di
\textbf{prevedere il futuro}.

Simone mostra il grafico della velocity della preproduzione di \emph{ISO}
durante l'anno scolastico. La consegna era a maggio (``10 maggio, o poco
dopo''). Già a inizio marzo, guardando l'andamento, aveva capito che \textbf{non
sarebbero riusciti a finire in tempo}. Il ragionamento era semplice: mancano
dodici sprint; quanto facciamo ogni sprint? Circa 89; quanto manca nel
backlog? Circa 2000. ``Non ce la facciamo''\footnote{I numeri sono quelli
detti a voce: 12 sprint per circa 89 punti danno poco più di 1000 punti,
circa la metà dei 2000 rimasti nel backlog.}.

La velocity aiuta a prevedere il futuro e quindi a \textbf{prendere
provvedimenti prima}. A marzo Simone ha cambiato i piani: prepariamo meno
cose, smettiamo di perdere tempo su ciò che ci blocca, consegniamo qualcosa di
diverso. In più, la velocity del team è \textbf{aumentata} sempre di più, e
alla fine sono arrivati in fondo alla preproduzione \textbf{senza crunch}:
``zero, completamente zero crunch''. Nessuno ha dovuto lavorare di più;
l'ultima settimana non hanno nemmeno toccato il progetto: hanno fatto la
build, l'hanno spedita e per una settimana non ci hanno lavorato. E sono
riusciti a consegnare anche di più del previsto, perché Simone ha
\textbf{analizzato i trend precedenti}: perché in un certo punto c'era un
picco? Che cosa li bloccava e faceva diminuire la produttività? L'ha
scoperto, l'ha capito e hanno agito.

\subsection{Come calcolarla in pratica}

Il consiglio è di \textbf{trovare un modo per fare il calcolo}. Simone usava
Monday, che non aveva uno strumento costruito apposta, ma ci è riuscito lo
stesso; si può fare tranquillamente con \textbf{Google Sheets}, ``col
sorriso''. Per farlo davvero bene bisognerebbe scrivere un po' di script
dentro il foglio, ``niente di improponibile''. Quando comincia a girare
``funziona da dio'': ``assolutamente fatelo, non posso dirlo abbastanza''. (Un
esempio costruito in diretta è nella sezione~\ref{sec:z3-qa-simone}.)

\subsection{Come assegnare i punteggi}

Per assegnare gli \emph{score} conviene partire dagli
\textbf{estremi}: quali task sono un 34? Quali sono un 1? Da lì si mettono i
valori intermedi. Gli estremi fanno da \textbf{ancore} e facilitano il
processo di valutazione \emph{relativa} descritto sopra.

Lo stesso vale per le \textbf{priorità}, se si usa un sistema di priorità:
regole simili, ma con \textbf{intervalli più piccoli} e senza Fibonacci, che
lì non serve. Simone usa una scala da 1 a 5; ma non la considera così
fondamentale.

\subsection{Che cosa blocca che cosa}

Durante il lavoro sul backlog bisogna capire \textbf{che cosa blocca che
cosa}. Ad esempio: ci sono 29 task, per un valore complessivo di 400, tutti
relativi alla creazione dei livelli. Ma se non c'è il movimento del
personaggio, non ha senso fare i livelli, perché vanno costruiti attorno a
quel movimento: la creazione del personaggio \textbf{blocca} i livelli e
quindi ha priorità più alta. Da qui si comincia anche a capire quale reparto
si occupa di che cosa: probabilmente all'inizio i programmatori faranno i
\emph{behavior} fondamentali (come il personaggio), mentre i designer
lavoreranno su documentazione e ricerca. Ma dipende dal gioco: alcuni non
costruiscono i livelli attorno al personaggio, bensì il personaggio attorno ai
livelli, quindi anche l'esempio ``è opinabile''.

% ---------------------------------------------------------------------
\section{Ricapitolando: come nasce una sprint}
\label{sec:z3-ricapitolo}

Prima di passare alle domande, perché ``quello che avevo da dirvi di pratico è
quasi finito'', Simone ricorda che il grosso non sono le regole, ma
\textbf{capire la logica alla base e integrarla nel proprio team e nel
proprio progetto}. Poi riassume la creazione di una sprint in quattro passi.

\begin{enumerate}
  \item \textbf{Post-mortem della sprint precedente.} Si vede che cosa ha
        fatto ciascuno, si capisce che cosa è bloccato e si
        \emph{programma un incontro} per risolvere i problemi.
  \item \textbf{Pianificazione.} Dal backlog si guarda che cosa blocca che
        cosa, quali sono le priorità e che cosa si può fare in questa sprint,
        e si prendono i task.
  \item \textbf{Analizzare i problemi e agire.} Si guarda la velocity: è in
        aumento? In calo? Si è in un picco? L'andamento a dente di sega,
        ``su, giù, su, giù'', è \textbf{assolutamente normale}, ed è ciò che
        ci si deve aspettare; e man mano che il progetto va avanti la
        velocity \textbf{tenderà a diminuire}. Le domande da farsi: come sta
        andando il progetto? Chi sta facendo di più e chi di meno, e perché?
        Qualcuno ha avuto un calo di rendimento? C'è un problema che si
        ripresenta spesso, un collo di bottiglia? E come risolverlo,
        soprattutto \textbf{ascoltando il resto del team}.
  \item \textbf{Rimanere consistenti}, in particolare con Scrum.
\end{enumerate}

\subsection{Ascoltare il team}

Come scrum master o producer, il lavoro più importante è e rimane
\textbf{ascoltare il team}. Tutto il resto non conta se poi, quando il team
dice una cosa, non lo si ascolta. È lo stesso principio dell'adattare il
\emph{framework} al team:
\begin{itemize}
  \item se il team odia il \emph{daily scrum}, non si fa: non serve a niente;
  \item se il team fa fatica a prendere i task dal backlog, non glielo si fa
        fare: si assegnano i task in modo diverso, ``non fategli manco vedere
        il backlog'';
  \item se il team vuole sprint di due settimane invece che di una, si fanno
        di due settimane.
\end{itemize}
Magari le persone si organizzano personalmente in modo diverso: Simone, per
esempio, organizza tutto il suo lavoro, su tutti i fronti, \textbf{su base
settimanale}, e una sprint di due settimane gli sconvolgerebbe i piani. Se il
team lavora per lo più così, si fa così. ``Adattatevi'' non vuol dire solo
adattarsi ai bisogni del software, ma \textbf{soprattutto ai bisogni del
team}.

\subsection{Nessun framework va bene al 100\%}

Esistono infiniti \emph{framework}, e \textbf{nessuno andrà bene al 100\% per
un team}. Bisogna provarli, il più possibile, e adattarli: cambiare qualche
regola, togliere qualche cerimonia. Il \emph{framework} più usato oggi, ad
esempio, è \textbf{Kanban}, che non ha nemmeno le sprint\footnote{Simone dice
prima ``Kanban neanche ce l'ha il backlog'' e poi si corregge: ``scusatemi,
neanche ce l'ha le sprint, ha soltanto un gigantesco backlog''.}: ha soltanto
un gigantesco backlog. ``Ogni team è diverso, ogni team ha bisogni diversi,
quindi adattatevi'' (una frase che Simone calcola di aver ripetuto ``almeno 49
volte'').

\subsection{La nota ``smielata'': la fiducia}

La conclusione: alla fine di tutto, \textbf{Agile non si basa sulle
iterazioni}, né sull'efficienza, né sulla creazione del miglior videogioco o
software del mondo, né sul controllo. L'unica cosa su cui si basa davvero è
\textbf{la fiducia}:
\begin{itemize}
  \item la fiducia che il team ripone in chi lo guida, e che chi lo guida
        ripone nel team: ``queste persone lavoreranno nel miglior modo
        possibile, sempre e comunque'';
  \item la fiducia del team nel dire ``se qualcosa va male, so che lo scrum
        master ha un piano di riserva che ci risolverà il problema, e il mio
        lavoro non sarà sprecato'';
  \item la fiducia di chi guida nel dire ``non posso controllarvi
        costantemente, non posso assicurarmi che facciate sempre tutto il
        lavoro o che seguiate le cerimonie come dovreste, però mi fido'': mi
        fido che lo farete, che se ci saranno problemi ve lo direte, che
        farete sempre del vostro meglio.
\end{itemize}
Bisogna \textbf{costruire fiducia nel team e trasmettere fiducia al team}. È
lo stesso principio del \emph{default to trust} della prima lezione zero
(sezione~\ref{sec:z1-default-trust}).

Simone ha finito con mezz'ora d'anticipo e apre alle domande (``go crazy'').
Diego apprezza molto il finale, ``smielato ma in realtà estremamente
veritiero'': lui è più orientato all'ambito professionale, ma crede che in
Agile ci debba essere davvero questa fiducia, e che sia anche una delle sue
forze. Simone pensava che non l'avrebbe apprezzato, tanto era smielato
(``mentre lo scrivevo ho percepito il diabete che aumentava''); Diego gli
suggerisce, la prossima volta, di aggiungere immagini con gli occhi a cuore e
qualche gattino.

% ---------------------------------------------------------------------
\section{Domande a Simone}
\label{sec:z3-qa-simone}

\subsection{La dimensione del team allunga le riunioni?}

Uno studente, che lavora in un team di tre persone (``essendo in tre le cose
vengono un po' più facili''), chiede quanto la \textbf{grandezza del team}
influisca sui tempi delle riunioni e delle cerimonie, e se ci sia un modo per
mantenerle rapide comunque.

Influisce assolutamente, ``per una mera questione matematica''. Ma in genere
è possibile far sì che un team di qualsiasi dimensione completi una riunione
in \textbf{una trentina di minuti al massimo}: quello è l'obiettivo, ``anche a
costo di dire meno cose, anche a costo di essere meno precisi''. A meno di
avere un team di una cinquantina di persone, la riunione durerà sempre trenta
minuti; e con cinquanta persone non si fa comunque una riunione tutti insieme,
perché è impossibile: si fanno le \textbf{riunioni di reparto}, e poi i
\textbf{capi reparto} fanno la loro riunione di sprint planning. I numeri
tornano piccoli: a una riunione ci saranno al massimo una quindicina di
persone, venti ``se proprio i pianeti si sono allineati'', e venti persone
possono stare in trenta minuti.

\subsection{Quanto deve imporsi lo scrum master? Il timekeeper}

Lo stesso studente chiede quanto lo scrum master debba \textbf{imporsi} sugli
altri, sul piano umano, per far rispettare i tempi. Moltissimo, risponde
Simone, ma \textbf{non deve per forza essere lo scrum master a farlo}. Nel
team australiano in cui lavora, all'inizio di ogni riunione si sceglie, in
modo pseudo-casuale, tra chi è disponibile o lanciando un dado, un
\textbf{timekeeper}: la persona che ha ``lo scettro del comando'' e tiene in
mano un cronometro, perché nessuno dedichi ai propri report più tempo del
dovuto. È \textbf{comunemente accettato da tutti} che quella persona abbia il
diritto di ``rompere le scatole'': è socialmente accettato.

Se le riunioni si dilungano troppo, si può proporre di introdurre la regola
del timekeeper, scelto ogni volta a caso o sempre lo stesso: ha il diritto di
richiamare gli altri perché, oggettivamente, il tempo scorre.

\subsection{E se il problema blocca tutti?}

Ultima domanda dello stesso studente: si è detto che i problemi si risolvono
dopo la riunione; ma se un problema è molto grave, se il fatto di non averlo
risolto blocca la produzione di più persone (``praticamente un disastro''),
vale la pena fermarsi e parlarne tutti insieme, o si dedica comunque quel
tempo solo ai report e poi si fa una riunione in più?

La risposta semplice è \textbf{no}: non si risolve nella riunione di sprint
planning e post-mortem. ``L'unica cosa che mi viene in mente che possa essere
così tanto grave è se la stanza in cui state facendo il meeting va a fuoco.''
La riunione non è pensata per quello. Subito dopo, piuttosto, si avvia
un'altra riunione con tutte le persone necessarie (``che durerà 25 ore'') per
risolvere il problema. Ma finché si fa planning e post-mortem, si fa
\textbf{quello e basta}, anche se il problema riguarda tutti.

\subsection{E se mancano delle persone?}

Si fa comunque la riunione? La si rimanda? Simone la fa \textbf{sempre}, a
meno di essere da solo, per \textbf{rispetto delle persone che sono venute}.
Le assenze possono avere cause diverse. Chi non è venuto deve mandargli un
messaggio su Telegram, anche vocale, in cui fa di fatto il proprio report.

Se una persona manca spesso, Simone ci parla e cerca di capire perché.
Magari l'orario della riunione non va bene perché ha altro da fare: nel
gruppo di \emph{ISO} c'è un ragazzo che lavora a settimane alterne di mattina
e di pomeriggio, e quando lavora di pomeriggio, tornando a casa, trova spesso
traffico; quindi a settimane alterne non può venire, o arriva molto in
ritardo. ``Lo so, me lo posso aspettare, non è un problema.'' Se invece
qualcuno manca in modo pseudo-casuale, bisogna capire il motivo: non gli va
bene l'orario? Ne preferirebbe un altro? Pensa che le riunioni non servano?
Non ha voglia di farle? Da lì si capisce il problema e si cerca di
risolverlo, spiegando che \textbf{le riunioni sono importanti} e che è
importante esserci tutti, e chiedendo ``dimmi come possiamo aiutarti a
rendertele meno pesanti''.

\subsection{Quando è la persona a non adattarsi}

Se però, nonostante tutto, la persona ``continua a fare la \emph{drama
queen}'', a un certo punto bisogna ammettere che esistono anche persone non
adatte a lavorare in un certo contesto, e si può valutare se lo sia.

Diego interviene per completare: in Agile tutti vogliono ottenere la cosa
migliore per il team e per il metodo di lavoro di ciascuno. Ma se un individuo
non si adatta al team, o il team non riesce ad adattarsi a lui, che sia o
meno una \emph{drama queen}, \textbf{la presenza di quella persona nel team va
valutata}. ``Nessuno, a parte casi eccezionali, tipo Miyamoto, vale il lavoro
degli altri.'' Le \emph{soft skill} restano alla base. E aggiunge che non è
solo questione di soft skill: ci sarà la \emph{drama queen}, quello che non si
adatta, ``lo stronzo o la stronza'' (il genere, precisa, è irrilevante),
quello che non ha voglia\footnote{Nella trascrizione ``stronzo
stronza,umbaquin'': la parte finale è incomprensibile.}; ma c'è anche la
persona che ha voglia di fare ma non si adatta al team. Se non ci riesce, o il
team non riesce ad adattarsi a lei, potenzialmente deve andare via.

Simone conferma che è un punto molto importante, che rientra nel discorso
dell'\textbf{ascoltare di che cosa ha bisogno il team}: se il team non riesce
a lavorare con un individuo, ``non c'è niente di male, è così e basta''. Per
fortuna non è troppo comune: a un certo livello professionale è difficile che
succeda, soprattutto in questa industria, dove le persone, se non stessero
lavorando insieme, potrebbero diventare amiche, perché hanno un interesse in
comune e sono molto appassionate. Ma può succedere.

\subsection{Come si calcola la velocity, in pratica}

Diego invita gli studenti ad approfittare di una persona ``che ha fatto il
vostro percorso''. Una domanda sul calcolo: la velocity è semplicemente la
\textbf{somma dei pesi dei task completati} nel tempo trascorso, cioè nella
sprint? Esatto. Si fa la sommatoria dei task \textbf{taggati come
\emph{done}} che hanno la stessa sprint (o la stessa data di consegna, a
seconda di come la si indica): si prendono tutti i task con ``sprint 5'', di
questi quelli in stato \emph{done}, e se ne sommano i pesi. Allo stesso modo
si possono sommare i task non completati, quelli rimandati, e così via.

\paragraph{Uno sguardo ai fogli personali.} Simone vorrebbe mostrare il foglio
di \emph{ISO}, ma è difficile da leggere per la quantità di dati; si chiede se
mostrare i propri fogli rompa qualche NDA, poi nota che sta già condividendo
lo schermo, ``quindi gli NDA sono già stati rotti''. Mostra allora uno dei
suoi fogli personali: organizza il proprio lavoro \textbf{settimana per
settimana}, con il nome della settimana, l'elenco dei task, la somma di ciò
che era da fare, quanto è stato effettivamente fatto, quanto è stato tagliato
e quanto rimandato per cause diverse. Il risultato è un grafico con
l'andamento settimanale di questi valori.

\paragraph{La ``speedrun velocity''.} Poi decide di fare ``una cosa più
figa'' e di costruire un foglio in diretta, che chiama \emph{speedrun
velocity}. La struttura minima\footnote{I riferimenti alle celle detti a voce
durante la costruzione (``B2B'', ``G1'', ``F1'', ``C2, F2'') sono confusi
nella trascrizione; la tabella e la formula sono una ricostruzione dello
schema, non una copia esatta del foglio mostrato.}:
\begin{center}
  \begin{tabular}{@{}llll@{}}
    \toprule
    \textbf{Task name} & \textbf{Score} & \textbf{Sprint} & \textbf{Status} \\
    \midrule
    Task 1 & 5  & Sprint 1 & done \\
    Task 2 & 13 & Sprint 1 & WIP \\
    Task 3 & 3  & Sprint 2 & cut \\
    \dots  & \dots & \dots & \dots \\
    \bottomrule
  \end{tabular}
\end{center}
Nella colonna \emph{Status} inserisce una convalida dei dati con i tre valori
\emph{cut}, \emph{WIP} e \emph{done}, ``per semplificarsi la vita''. Poi
scrive una serie di task con punteggi a caso (copiati da un altro foglio) e
assegna le sprint, avvertendo che ``ovviamente non dovrebbe essere fatto
esattamente così, ma \emph{for the sake of showcase}''; aggiunge qualche
colore, perché ``sono pazzo'', anche se non è la cosa fondamentale.

Il calcolo è una somma condizionale su più criteri (\texttt{SUMIFS}, in
italiano \texttt{SOMMA.PIÙ.SE}): si somma la colonna dei punteggi per le
righe in cui la sprint è quella cercata \emph{e} lo stato è \emph{done}. Ad
esempio:
\begin{center}
  \texttt{=SUMIFS(B:B; C:C; "Sprint 1"; D:D; "done")}
\end{center}
Lo stesso si può fare con \emph{cut} per sapere quanto è stato tagliato. Si
ripete per la sprint~2, e così via, e con i risultati si inserisce un
\textbf{grafico a linee} (Sheets propone il tipo di grafico sbagliato, e
Simone lo corregge). Con più sprint si ottiene un andamento come quello dei
suoi fogli personali o di \emph{ISO}. Chi vuole fare le cose per bene può
aggiungere altre analisi: chi ha fatto più task, chi ha più peso e carico di
lavoro, lo stato generale del progetto (``magari il backlog fosse così
vuoto'').

Era, dice, per essere sicuro di riuscirci; è disponibile a rifare qualche
passaggio, ma alla classe è tutto chiaro. Uno studente chiede se il foglio si
possa avere: Simone lo mette in una cartella e sarà Diego a farlo avere alla
classe.

\subsection{I problemi più gravi incontrati con Agile}

Diego osserva che Simone ha capito come applicare i principi di Agile con una
velocità sorprendente, più di quanto sia successo a lui quando gli fu
spiegato, e che all'inizio si fa spesso ``Waterfall Agile'', cioè Waterfall
con le sprint, credendo di fare Agile\footnote{Il passaggio è frammentato e non
è del tutto chiaro chi abbia fatto ``Waterfall Agile'' all'inizio: dalla
risposta di Simone (``ho fatto anch'io waterfall agile per un po''') sembra
trattarsi di un errore comune, che Diego attribuisce anche a se stesso o agli
studenti in generale.}. Poi la domanda: quali sono stati i problemi più gravi?

\paragraph{I task ``tenuti in caldo''.} Simone si è accorto che le persone non
chiudevano i task al momento giusto. Se diceva che un task richiedeva, per
esempio, quattro ore, qualcuno lo finiva in un'ora e poi, invece di segnarlo
come completato, aspettava (``me lo rivedo più tardi, tanto ho ancora altre
tre ore''). Non funziona così: \textbf{quando un task è completato, è
completato}, punto. Non c'è ``me lo rivedo tra un po''' né ``devo farlo in
queste quattro ore, quindi ci metto quattro ore''. Questo falsava i risultati,
e Simone ha cambiato subito approccio\footnote{L'esempio parla di stime in
ore, in apparente contrasto con il consiglio di non usarle
(sezione~\ref{sec:z3-velocity}); è plausibile che si riferisca proprio a una
fase precedente in cui le ore erano ancora usate, e che il ``cambiamento
immediato'' sia stato anche questo.}.

\paragraph{Le riunioni che nessuno voleva fare.} Durante il secondo anno le
persone non volevano fare le riunioni settimanali. All'inizio le facevano solo
tra designer; poi le hanno estese a tutti, perché non aveva senso limitarle,
ma la risposta era ``questa riunione è inutile, non ha senso''. In quel caso
Simone non è riuscito a risolvere il problema: non sapeva come farlo. C'erano
problemi anche sul piano umano, non tra i designer ma con membri di altre
classi: problemi che, col senno di poi, si sarebbero potuti risolvere, ma
all'epoca non ci sono riusciti. La soluzione adottata è stata non fare le
riunioni con tutto il team: gli altri membri mandavano un \textbf{recap
settimanale per messaggio}. Non era un metodo del tutto fallimentare, ma si
perdono molti benefici delle riunioni tutti insieme: il \emph{team building},
il confronto, il vedersi, il parlare. È stato, dice, ``un grosso problema e
anche un grosso fallimento''.

Per il resto, niente di così grosso da valere la pena di raccontarlo:
valutazioni di task sbagliate, sprint fatte male, e un periodo di
``Waterfall Agile'' anche per lui, finché non ha capito di sbagliare.

\subsection{Che cosa rispondere a chi dice che le riunioni sono inutili}

E oggi, se qualcuno dicesse ``per me queste riunioni sono inutili'', che cosa
gli risponderebbe? \textbf{``Perché?''}. E da lì si capisce qual è il
problema: non vuole dedicarci tempo? Durano troppo? Gli costa fatica? Ha
problemi con le riunioni perché soffre di qualche forma d'ansia e non gli
piace stare con tutti i membri del team, anche fuori dalle riunioni (``tipo
ci becchiamo a giocare online a qualcosa'')? Bisogna capire il perché e
\textbf{tirarlo fuori dalle persone}, perché spesso le persone stesse non sanno
il perché delle cose che dicono e sentono.

\subsection{I saluti a Simone}

Diego ringrazia Simone e spera che gli studenti abbiano appreso il più
possibile: un'esperienza più diretta di quella di Simone e Cristian non si può
avere, perché hanno fatto letteralmente lo stesso percorso (la lezione di
Diego è ``fondamentalmente la stessa'' da un anno all'altro, con qualche
aggiunta). Ricorda che Simone ha concesso anche il foglio della
\emph{speedrun velocity}, che Diego passerà alla classe, e lo saluta per non
rubargli altro tempo.

% ---------------------------------------------------------------------
\section{Cristian: il punto di vista del game designer e project manager}
\label{sec:z3-cristian}

\subsection{Chi è il relatore}

Il secondo intervento è di Cristian, che attiva la webcam (``mi vedete super
zoomato, spero non sia una brutta vista'') e condivide le slide, mentre Diego
cerca di farle andare in modalità presentazione perché vengano meglio nella
registrazione. Cristian ringrazia Diego per la presentazione e si presenta:
\begin{itemize}
  \item è uno \textbf{studente del terzo anno} di Game Design, quindi ha fatto
        lo stesso percorso della classe;
  \item è anche \textbf{tutor} del corso di Game Design per gli studenti
        dell'accademia online;
  \item è \textbf{game designer e project manager} di due progetti:
        \begin{itemize}
          \item \emph{Gamers in Time}\footnote{Così nella trascrizione;
                il titolo non è stato verificato.}, un \textbf{gioco da
                tavolo}, vincitore del \emph{Big One} del primo anno;
          \item \emph{L'Inf}\footnote{Il titolo compare in molte forme
                (``l'Inf'', ``Linf'', ``Leanf'', ``Linfa''); qui si usa
                \emph{L'Inf}, la più frequente. Non è chiaro se la grafia
                corretta sia \emph{Linf}, \emph{L'Inf} o \emph{Linfa}.}, un
                \textbf{videogioco} che ha vinto, ``tra virgolette'', il
                \emph{Bigger One} del secondo anno e che il team ha avuto
                l'opportunità di continuare come \emph{Biggest One} del terzo
                anno, opportunità che ha accettato.
        \end{itemize}
\end{itemize}
In entrambi si è occupato di tutta la parte di design, insieme ad altri
colleghi, ed è per lo più anche \textbf{project manager, scrum master,
producer}, ``come volete definirmi''.

\subsection{Come è strutturata la sua parte}

Oltre all'esperienza personale, la lezione ha due parti principali:
\begin{enumerate}
  \item la \textbf{produzione} effettiva tramite Scrum;
  \item la \textbf{gestione del team} in un ambiente di lavoro Scrum.
\end{enumerate}
Possono sembrare la stessa cosa, ma ``il progetto tendenzialmente è una cosa
e il team è un'altra'': ci sono ovvi rapporti tra le due parti, ma per
ciascuna ci sono precauzioni specifiche.

\subsection{Le premesse}

\begin{itemize}
  \item Quello che descrive è \textbf{un} percorso, un modo di applicare
        Scrum, \textbf{non l'unico}: se non si è d'accordo, o se il team non
        ``sopporta'' questo metodo, si possono fare tutti i cambiamenti che si
        vuole (e qualche consiglio lo darà strada facendo).
  \item Non è una lezione descrittiva come quelle di Diego (che cos'è Scrum,
        le definizioni, le cerimonie, come si calcolano la velocity e il peso
        dei task), ma un \textbf{confronto} sull'esperienza dei due progetti,
        più su \emph{L'Inf} che su \emph{Gamers in Time}, perché Scrum l'ha
        applicato al videogioco e non al gioco da tavolo.
  \item \textbf{Tutto è in continua evoluzione}. Per questo vuole le domande
        della classe: il primo obiettivo è chiarire i dubbi e far uscire tutti
        dalla lezione avendo capito un po' di più la mentalità agile, ``con
        tutti voi stessi''; il secondo è ricevere \textbf{feedback} per
        migliorare non solo le lezioni, ma lo stesso modo in cui applica
        Scrum.
  \item Ha un po' di mal di gola e si fermerà spesso a bere: si scusa in
        anticipo.
\end{itemize}

% ---------------------------------------------------------------------
\section{Il caso di studio: \emph{L'Inf}}
\label{sec:z3-linf}

\subsection{Le caratteristiche del gioco}

Qualcuno in classe forse conosce Cristian o ha provato il gioco a Level Up,
l'evento in cui è stato presentato. Le caratteristiche principali:
\begin{itemize}
  \item è \textbf{single player} e offline: niente co-op, niente
        multiplayer, niente server;
  \item è un'\textbf{avventura fantasy}: un gioco di \textbf{durata fissa}, che
        tendenzialmente si gioca per una \emph{run}; chi l'ha apprezzato di
        più e vuole vedere tutti gli aspetti del design e delle meccaniche, o
        esplorarne i segreti, lo rigiocherà;
  \item il giocatore ha delle \textbf{magie colorate} con cui affronta
        \textbf{orde di nemici colorati}: spara le magie e, a seconda di
        quanti nemici ``scoppiano'' con un solo colpo, ottiene vari bonus,
        come succede in \emph{Candy Crush}.
\end{itemize}

\paragraph{Dati e passione.} Diego interrompe perché ritiene questa parte molto
importante: com'è nato l'elemento davvero innovativo di \emph{L'Inf}? Dallo
\textbf{studio dei dati}, ma anche dalla \textbf{passione del team}. I dati
danno una direzione: le ricerche di Cristian indicavano il successo delle
meccaniche \emph{match-3}\footnote{Nella trascrizione ``match free'' e
``match tree'': si tratta delle meccaniche \emph{match-3} (allineare tre o
più elementi dello stesso colore), tipiche di \emph{Candy Crush}.}; e
Cristian è anche, principalmente, un giocatore mobile. La combinazione tra
passione e ricerca ha prodotto un progetto molto valido sul piano della
fattibilità (single player, fantasy, espandibile).

Cristian aggiunge un'altra ricerca fondamentale, da cui è stato poi delineato
il \textbf{documento di visione}\footnote{Nella trascrizione ``iVision'' e
``l'iVision'': si tratta quasi certamente della \emph{vision}, il documento
di visione (sezione~\ref{sec:perche-documentare}).} del gioco: quella su
che cosa chiedesse Epic Games per finanziare i progetti con l'\textbf{Epic
MegaGrant}. Il team ha individuato tre elementi fondamentali\footnote{È la
lettura dei requisiti fatta dal team, come riferita da Cristian; non è stata
verificata sui criteri ufficiali del programma.}:
\begin{enumerate}
  \item \textbf{innovazione};
  \item \textbf{temi ambientali}: attraverso i colori e l'ambientazione, il
        gioco vuole criticare l'inquinamento globale;
  \item un \textbf{elemento narrativo}.
\end{enumerate}
Questi tre elementi, le ricerche specifiche sui vari tipi di gioco e la
passione (di Cristian in primis, come lead, ma anche degli altri membri,
incuriositi dalla proposta) hanno portato a un progetto che al momento sembra
funzionare: ``stiamo andando nella direzione giusta, almeno secondo quanto
previsto; non si può mai prevedere al 100\% ciò che accadrà''.

\paragraph{Il contenuto.} Oltre alla meccanica \emph{match-3}, il gioco
prevede:
\begin{itemize}
  \item \textbf{quattro modalità di sparo};
  \item \textbf{cinque livelli}, con esplorazione e fasi di combattimento
        contro le orde;
  \item \textbf{puzzle};
  \item \textbf{cinematic} per lo più ambientali: niente attori, doppiatori o
        mimica facciale, ma scene di presentazione dell'ambiente di gioco;
  \item una \textbf{boss fight} finale.
\end{itemize}
Il perché di questo elenco sarà chiaro tra poco.

\subsection{Un ``trauma'' di produzione: il cambio di team}

La produzione si è sviluppata in due sezioni, e \emph{L'Inf}, ``per fortuna o
purtroppo'', ha avuto un grave problema: è la risposta, anticipata, alla
domanda fatta a Simone sui problemi più gravi. Come tutte le produzioni di
videogiochi, ha subito dei ``traumi''. Il più grande: \textbf{il 100\% del
team originario è stato sostituito}. Non per scelta, ma perché i membri non si
sono iscritti al terzo anno: avevano altri obiettivi; c'è chi cerca lavoro e
chi ha aperto un'azienda. Cristian, volendo continuare il terzo anno e
soprattutto \emph{L'Inf}, un progetto in cui crede e che ha dimostrato il suo
potenziale a Level Up e con i feedback costanti dei giocatori, ha dovuto fare
dei \textbf{colloqui} per formare un nuovo team che si adattasse meglio al tipo
di sviluppo del gioco\footnote{La frase è in parte persa nella trascrizione
(``ho dovuto fare dei colloqui Grazie a tutti. Cioè, servisse una parola
brutta, ma semplicemente si adattasse di più al tipo di sviluppo''): il
senso ricostruito è che ``colloqui'' è una parola forte, ma che si trattava
di scegliere persone adatte a questa fase.}.

La composizione è cambiata così:
\begin{center}
  \begin{tabular}{@{}lp{0.55\textwidth}@{}}
    \toprule
    \textbf{Team} & \textbf{Composizione} \\
    \midrule
    originario & 2 game designer, 2 level designer, 1 technical designer \\
    attuale    & 3 game e level designer, 2 technical designer \\
    \bottomrule
  \end{tabular}
\end{center}
Il motivo: la preproduzione, cioè la parte più importante di game design, era
finita, e servivano \textbf{più tecnici}. Il gioco si basa molto sulle
meccaniche: un sistema di combattimento innovativo che richiede di scriptare
le meccaniche, diversi tipi di sparo, puzzle, una boss fight. È molto più
basato sulle meccaniche che su narrativa o ambientazione. Allo stesso tempo
servivano level designer, perché il gioco doveva prendere forma e i livelli
andavano costruiti.

\subsection{Perché raccontare tutto questo}

Non per vendere il gioco (``anche se magari è un effetto collaterale ben
voluto''), ma per mostrare che \textbf{ogni progetto, a seconda di come è
composto} sul piano del design, della produzione e soprattutto del team,
\textbf{ha un modo diverso di applicare Scrum}. Per questo Cristian è sceso nel
dettaglio delle caratteristiche: ognuna comporta ambiti di produzione diversi.
\begin{itemize}
  \item I livelli richiedono level designer.
  \item Il \emph{match-3} e i diversi tipi di combattimento richiedono
        tecnici.
  \item Il single player è più un vantaggio che un limite: un gioco
        multiplayer non è il doppio, ma ``tipo dieci volte più difficile'' da
        gestire a livello produttivo.
  \item Il genere \emph{fantasy adventure} deriva soprattutto dalle ricerche.
\end{itemize}
Quindi \textbf{ogni progetto ha il suo framework di Scrum}. E il fatto di
aver cambiato team a metà produzione, finita la \emph{vertical slice}, con
persone completamente nuove (ex compagni di studio, che però non avevano mai
visto come funzionava il gioco), ha avuto diverse conseguenze, che Cristian
rimanda alle domande di fine lezione.

% ---------------------------------------------------------------------
\section{Il loop di Scrum}
\label{sec:z3-loop}

L'applicazione di Scrum è un \textbf{loop} molto semplice, simile a quello
mostrato da Simone:
\begin{enumerate}
  \item \textbf{creazione del backlog};
  \item \textbf{selezione dei task} da svolgere, a ogni sprint;
  \item \textbf{implementazione dei task} durante la sprint, con controlli e
        \emph{scrum} (con la frequenza che si decide);
  \item \textbf{calcolo della velocity}, per vedere com'è andata e come
        migliorare la produzione nelle sprint successive;
  \item \textbf{interpretazione} della velocity per fare il \emph{tuning}
        della produzione; e si ricomincia.
\end{enumerate}
Cristian le affronta una per una, con esempi della sua esperienza e consigli.
È contento che la classe stia cominciando proprio adesso, perché questa parte
potrebbe tornare utile già dalla settimana successiva, quando gli studenti
applicheranno il \emph{framework}.

% ---------------------------------------------------------------------
\section{La stesura del backlog}
\label{sec:z3-backlog}

\subsection{Un backlog per milestone}

Anche Cristian, all'inizio, cercava di prevedere il più possibile tutto il
gioco, dal prototipo fino al rilascio. Oggi invece lo struttura
\textbf{per milestone}: il backlog dell'\emph{alpha}, poi quello della
\emph{beta}, poi quello della \emph{gold}. Questo non impedisce di cominciare
a prevedere che cosa ci sarà in beta e in gold, ma dà un'idea molto chiara di
dove si sta andando, soprattutto per capire che cosa tagliare o aggiungere
quando ci si accorge che qualcosa serve prima o dopo, e adattarsi di
conseguenza: ``la cosa più importante del mindset agile è adattarsi ed essere
flessibili''.

\subsection{Da dove vengono i task}

Il gioco deve pur nascere da qualche parte: i task vengono dai
\textbf{documenti chiave} della produzione, già visti nel corso
(capitolo~\ref{cap:produzione-documentazione}):
\begin{itemize}
  \item il \textbf{documento di visione};
  \item i \textbf{GDD} delle varie feature, che dettagliano la visione;
  \item il \textbf{versioning} delle feature, a cui si dovrebbe pensare già
        mentre si scrive il GDD;
  \item gli \textbf{LDD} (\emph{level design document}) e tutti gli altri.
\end{itemize}
L'ordine dipende dal gioco: in un gioco fortemente basato sulle meccaniche è
probabile che gli LDD arrivino dopo lo sviluppo del \emph{core}; in un gioco
basato sull'esplorazione non succede proprio l'inverso, perché le meccaniche
servono comunque, ma le due cose avvengono quasi in contemporanea.

Una volta che si ha la documentazione con tutte le descrizioni e le
caratteristiche delle feature, \textbf{dalla versione 0 alla versione X}, si
ha già una serie di implementazioni da realizzare. Un esempio:
\begin{itemize}
  \item \textbf{versione 1}: il giocatore deve camminare; si usa il
        manichino base di Unreal;
  \item \textbf{versione 2}: si implementa la \emph{skeletal mesh} fatta dal
        3D artist e si adatta la camera al movimento, con gli effetti di
        vibrazione a seconda che il personaggio corra, cammini o sia
        abbassato;
  \item \textbf{versioni successive}: quando il personaggio si muove su
        terreni diversi solleva tipi diversi di polvere, di più su un terreno
        molto polveroso. Non è fondamentale e potrebbe essere una delle
        ultime versioni, verso la fine.
\end{itemize}
Il consiglio è che, quando si fa il \emph{versioning}, si inseriscano
\textbf{tutti gli ambiti necessari} per quella feature: programmazione, arte
(se serve un asset), animazione, UI, VFX, suoni. Se tutti i \emph{field}
compaiono già nella documentazione, la creazione del backlog sarà più chiara
e più spedita.

\subsection{Come organizzare i task}

I giochi hanno moltissimi task. Come organizzarli perché siano leggibili non
solo da chi li ha scritti (scrum master o project manager), ma anche da chi
deve navigarli per capirci qualcosa? Cristian usa tre tipi principali di
organizzazione.

\paragraph{1. Per feature (o epic).} Si suddividono i task per feature,
collegandosi alle \emph{epic} di Simone. Ad esempio, l'epic è il movimento
del personaggio e i task sono camminata, corsa, salto, volo, ``dipende dal
gioco''.

\paragraph{Tutto in un solo posto.} La suddivisione non va fatta per forza in
documenti diversi, anzi Cristian lo sconsiglia: se tutto è nello stesso
strumento (Google Sheets, Monday, ClickUp, Trello) si usano i
\textbf{filtri} e si trova tutto lì, senza navigare tra mille siti o mille
documenti.

\paragraph{Per ambito, con un color code.} Quando si suddivide per ambito
(programmazione, arte, UI, eccetera), conviene usare un \textbf{codice
colore}. Cristian lo usa ``per qualsiasi cosa'': uno per le milestone, uno
per gli ambiti di produzione. Diego gli chiede di mostrare il suo:
\begin{center}
  \begin{tabular}{@{}lll@{}}
    \toprule
    \textbf{Ambito o milestone} & \textbf{Colore} & \textbf{Perché} \\
    \midrule
    Suoni        & viola        & --- \\
    VFX          & rosso        & sta imparando a usare Niagara \\
    Level design & verde        & ``erano finiti i colori'' \\
    Game design  & arancione    & è il suo colore preferito \\
    \midrule
    Alpha        & azzurro      & --- \\
    Beta         & viola        & --- \\
    Gold         & giallo scuro & il colore richiama il nome \\
    \bottomrule
  \end{tabular}
\end{center}
Come si vede, è un misto di funzionalità e preferenze: la cosa più importante
è che \textbf{tutto il team sia d'accordo}.

\paragraph{2. Per milestone, con le sotto-milestone.} Un secondo metodo è
suddividere il backlog per milestone. Ma se il backlog dell'alpha è separato
da quelli di beta e gold, come si fa a dividere in alpha, beta e gold un
backlog in cui tutto è alpha? Cristian suddivide le milestone in
\textbf{sotto-milestone}: l'alpha dell'alpha, la beta dell'alpha, la gold
dell'alpha. Serve ad avere idee più chiare, a tenere sempre in mente lo
\emph{scope} senza perderlo e soprattutto a capire meglio \textbf{se il gioco
è fattibile}. Una suddivisione più fine permette anche di calcolare con più
precisione (una precisione che ``lascia il tempo che trova'') quante sprint
servono per ogni sotto-milestone. La dissonanza tra ``dividete tutto in alpha,
beta e gold'' e ``dentro l'alpha trovate alpha, beta e gold'' è quindi voluta:
si parla delle sotto-versioni.

\paragraph{3. Per priorità.} Chi non vuole usare le sotto-milestone può usare
il sistema di \textbf{priorità} di cui parlava Simone: una scala numerica, da
1 a 5 o da 1 a 10, a seconda di come si riesce meglio a distinguere le
priorità. Serve soprattutto a capire quali task selezionare
(sezione~\ref{sec:z3-selezione}). In realtà è più o meno la stessa cosa con un
approccio diverso: indicare alpha, beta e gold dell'alpha crea già di per sé
delle priorità. ``Dipende tutto da voi.''

\subsection{Come scrivere i task e dove}

Il consiglio è scriverli \textbf{dal punto di vista dell'utente}. Ma, come
diceva Simone, non sempre ``il gioco vale la candela''. Lo è soprattutto in
un caso: se il programmatore legge un task dell'artista 3D scritto in modo
troppo tecnico e non lo capisce, si crea dissonanza tra i reparti e possono
nascere attriti. Quindi è sempre consigliabile, ma non obbligatorio: spesso
basta una chiamata (``prima semplicemente ti giravi verso il tuo compagno e
gli chiedevi che cosa significa''). ``Però è giusto che vi dica le cose come
dovrebbero essere fatte.''

Dove scriverli? In qualsiasi strumento: Trello, Monday, ClickUp, gratuiti o a
pagamento, oppure un foglio Google costruito da sé. Dall'esperienza di
Cristian emerge che \textbf{Google Sheets sta antipatico alle persone}: già il
nome ``foglio di calcolo'' spaventa. Strumenti come Monday o ClickUp hanno più
colori e soprattutto \textbf{mandano le notifiche}, quindi sono più accessibili
a tutti. ``Decidete voi.''

\subsection{Non prevedere tutto: la flessibilità}

È fondamentale capire che \textbf{prevedere tutti i task passo per passo}, in
un ordine preciso (``dopo la camminata faremo il salto, dopo il salto il
sistema di sparo, e così via all'infinito''), va contro i principi di Agile,
in particolare contro il \textbf{quarto valore}: rispondere ai cambiamenti più
che seguire un piano.

Da qui il consiglio centrale: la \textbf{flessibilità} deve essere il cuore
della propria filosofia di lavoro e del proprio modo di applicare Scrum.
Bisogna essere sempre pronti a \textbf{tagliare} nel modo più agile possibile:
se si è capita davvero la mentalità agile, si riesce a tagliare una feature
\textbf{senza bloccare la produzione delle altre}.

L'esempio più eclatante viene dall'inizio di \emph{L'Inf}: Cristian faceva il
backlog di tutto il gioco, arrivando a circa \textbf{400 task} di ogni ambito,
nel dettaglio: per la camminata arrivava a indicare come si muoveva la camera
quando il personaggio girava a destra, e come cambiava quando saltava o cadeva
da un dirupo. ``Vi viene subito da pensare: ma come fai a prevedere tutto?''
Infatti la soluzione è \textbf{non prevedere tutto}: capire a grandi linee ciò
che si vuole, qual è l'obiettivo della feature, senza scendere troppo nel
dettaglio, perché il dettaglio crea di per sé un \textbf{lavoro Waterfall}, che
è ciò che bisogna assolutamente evitare. E scrivere un backlog così dettagliato
richiede tantissimo tempo, che è limitato: la produzione ne risente in tutti i
modi possibili.

\paragraph{Bonus: gli elementi innovativi.} Se il gioco (il proprio
\emph{Bigger One}) ha elementi innovativi, la flessibilità va portata
\textbf{ai massimi livelli}. \emph{L'Inf} era partito con una meccanica di
\textbf{assorbimento e sparo}: si assorbe un nemico colorato, ad esempio
giallo, e lo si spara contro altri nemici gialli, prendendo ispirazione da
altri giochi. Nel loro caso, però, creava più problemi che altro: per due o
tre ore di gioco, ``assorbire, sparare, assorbire, sparare'' risultava
\textbf{troppo ripetitivo}, e hanno deciso di tagliare la meccanica. Il
problema è che l'hanno tagliata dopo che Cristian aveva già scritto
\textbf{600 task nel backlog e 12 GDD}: tempo perso. Per questo la
flessibilità, negli elementi innovativi ma in generale, ``è la prima cosa''.

\subsection{Quanto dettaglio?}

Quando un backlog è abbastanza dettagliato e quando lo è troppo poco?
``La risposta purtroppo non c'è.'' Viene da due fattori:
\begin{enumerate}
  \item l'\textbf{esperienza}: man mano che si scrivono GDD e si va avanti
        nello sviluppo, ci si ``tuna'';
  \item il \textbf{trial and error}: si sbaglia e si torna indietro, il che
        riporta, di nuovo, all'esperienza.
\end{enumerate}

% ---------------------------------------------------------------------
\section{La selezione dei task}
\label{sec:z3-selezione}

Con un bellissimo backlog ordinato per colori, ambiti e priorità, quali task
si selezionano? ``Da qualche parte bisogna pur cominciare.'' Le domande
principali sono due: \textbf{quali} selezionare e \textbf{quanti}.

\subsection{Quali task}

\paragraph{La priorità.} Il primo criterio è la priorità: i task da fare
prima sono le scelte più urgenti, e quindi si scelgono quelli\footnote{L'inizio
di questo passaggio manca nella trascrizione (``quanti selezionare. cui
scegliere prima, sarà la scelta più urgente da fare''): il primo criterio è
stato ricostruito dal contesto, anche perché poco dopo Cristian dice che,
risolto il problema delle competenze, ``si ritorna a selezionarli per
priorità''.}.

\paragraph{Le capacità attuali degli sviluppatori.} Il secondo criterio sono
le \textbf{abilità dei developer}. Cristian precisa che non intende in alcun
modo giudicarne la professionalità: intende letteralmente \textbf{che cosa sa
fare lo sviluppatore in quel momento}. Se un programmatore non ha mai usato
Unreal, non è il momento di dargli task di implementazione in Unreal: bisogna
prevedere anche il \textbf{periodo di apprendimento}, eventualmente come task
vero e proprio (``hai questa sprint per imparare a implementare l'IA'').

È successo due mesi prima con uno dei suoi technical designer: gli ha dato
\textbf{due settimane} dicendogli ``in queste due settimane non ti preoccupare
di niente, il resto lo produciamo noi''. Doveva essere d'accordo, e lo era,
anzi era contento di approfondire l'argomento: studiare come funzionano le IA
(anche grazie a vari link passati da Diego) e soprattutto come implementarle
nel motore. Per due settimane ha dovuto ``solamente'' pensare
all'apprendimento. Il consiglio: \textbf{prevedere questi ``tempi morti''},
perché ne ha bisogno il team, e ne ha bisogno chi pianifica per non falsare la
previsione dello \emph{scope}. Se in questo momento (``sottolineo, in questo
momento'') un programmatore non sa fare una cosa, la si rimanda e si mette
altro; prima o poi imparerà, e a quel punto si torna a selezionare per
priorità.

\subsection{Quanti task}

Il primo criterio sulla quantità è la \textbf{velocity}, ma Cristian l'ha
messo nella slide solo per completezza del loop: se non si ha ancora una
velocity, non ha senso. Per ora la si trascura e si fa finta che sia la
\textbf{prima sprint}, la prima applicazione di Scrum.

Nelle prime sprint la quantità si basa sulle \textbf{capacità} e sulla
\textbf{disponibilità} degli sviluppatori:
\begin{itemize}
  \item \textbf{disponibilità}: per quanta parte della sprint il technical
        designer (per riprendere l'esempio) è disponibile? Non in termini di
        ore, perché Cristian concorda con Simone che le ore sono sempre molto
        relative;
  \item \textbf{performance}: gli studenti si conoscono ormai da quasi un
        anno e mezzo. Se uno sviluppatore non ha mai fatto un livello in
        Unreal (per motivi suoi, magari problemi al PC) e dice ``voglio fare
        un livello e ci metto due settimane'', sta a chi pianifica cominciare
        a capire \textbf{chi può fare che cosa}. Tutti possono fare tutto, se
        si prevede il tempo di apprendimento; ma se Cristian arrivasse a dire
        ``in due settimane vi faccio l'IA del vostro gioco implementandola in
        Unreal'', ``vi sto prendendo in giro''.
\end{itemize}
Queste sono le analisi preliminari quando non si ha una velocity. Se invece
non si conosce affatto il team, l'unico punto di partenza è la
\textbf{disponibilità} degli sviluppatori, da cui si stima quanto si riesce a
fare.

Tanto, \textbf{nelle prime sprint si sbaglieranno completamente tutte le
previsioni}, ``ed è giusto così, e ben venga'': se non si sbaglia,
probabilmente si sta applicando Waterfall, o si sta applicando male Scrum.
All'inizio si tratta solo di buttarsi: ``questa settimana abbiamo previsto 15
task; ne abbiamo fatti 40. Bene, la prossima ne proviamo 40''; poi magari se ne
fanno 25, e così via, \textbf{tarando} man mano la quantità.

\subsection{Definire bene il goal}

Per ogni sprint, e anche per ogni singolo task, bisogna \textbf{definire bene
l'obiettivo}. Serve a mantenere concentrati gli sviluppatori: se so che entro
questa settimana devo finire il block out del livello, so dove sto andando,
lavoro per quel risultato e soprattutto \textbf{non perdo tempo}, come succede
spessissimo. Se durante il block out qualcuno si mette a inserire effetti
visivi ambientali o il \emph{post-process}, viene il dubbio che non stia
lavorando per l'obiettivo: è un \textbf{campanello d'allarme}, perché potrebbe
far mancare l'obiettivo per cui la sprint era stata pensata.

\subsection{Bonus: progettare il ritmo dello sviluppo}

Un consiglio che Cristian ha sentito raramente, ma che ha deciso di
sperimentare: poiché gli studenti sono designer, e tendenzialmente anche lo
scrum master del loro progetto lo sarà, possono \textbf{progettare il
\emph{pace} dello sviluppo}.

Di solito si pensa (e potrebbe anche essere giusto) che l'obiettivo di Scrum
sia trovare una \emph{velocity target} e mantenerla costante per tutta la
produzione. Può essere una soluzione; ma con entrambi i suoi team Cristian ha
notato che, per come sono fatti gli esseri umani, ci sono \textbf{picchi di
attenzione e di produzione}. Con l'esperienza si impara a percepirli in
anticipo e a progettare di conseguenza; oppure sono gli sviluppatori stessi a
dirlo: se la settimana prima hanno implementato 12 nemici, 15 modalità di sparo
e 13 cutscene, quella dopo diranno ``andiamoci piano''. Di solito non succede,
perché le sprint sono programmate apposta per adattarsi alle capacità degli
sviluppatori, ma è un metodo possibile.

Che cosa intende, in concreto? Se per un periodo si sono implementate molte
feature o sistemi che però, visivamente, \textbf{non hanno portato un aumento
di qualità percepibile nel gioco}, gli sviluppatori non si sentono
riconosciuti e ripagati dello sforzo. Allora, la settimana dopo, ci si
concentra su qualcosa che dà \textbf{soddisfazione}. Questo lo permettono le
sprint, lo permette Scrum e non Waterfall: è una differenza fondamentale.

\paragraph{L'esempio di \emph{L'Inf}.} Per più di otto mesi il gioco è rimasto
con \textbf{quattro tipi di nemici}. Nonostante l'impegno (quattro nuove
modalità di sparo, la boss fight, tantissima programmazione ``di
background'', cioè non visibile al giocatore ma necessaria al funzionamento
del gioco), il team cominciava a percepire che \textbf{non c'era corrispondenza
tra impegno e qualità percepita}. Cristian ha allora deciso di limitare per un
periodo la produzione degli elementi di background (programmazione come l'IA,
tutto ciò che faceva proseguire il gioco ``sotto''), di inserire altri tipi di
nemici, che tutti non vedevano l'ora di avere (``sapevo che era uno di quei
jolly da sfruttare per far risalire il morale del team''), e di dedicare tempo
a effetti visivi e sonori, a migliorare tutto ciò che era estetico. Diego
suggerisce il termine: la \textbf{juiciness}. Ciò che vedeva il giocatore lo
vedevano anche gli sviluppatori, che si sentivano più appagati e sentivano i
propri sforzi ricompensati. ``Adesso loro non lo sanno, ma li ributterò sotto
un treno di lavoro'' (``in realtà lo sanno, però non me lo dicono'').

Per Cristian è un punto fondamentale, che gli è costato ``mesi di
sofferenze'' di esperienza per capirlo. Diego invita la classe a fare domande
su un \emph{insight} ottenuto per tentativi ed errori.

% ---------------------------------------------------------------------
\section{L'implementazione: controlli, blocker e feedback}
\label{sec:z3-implementazione}

\subsection{Il check dei progressi}

Selezionati i task, viene la parte ``più tranquilla'' dal punto di vista di
Scrum (i problemi di implementazione, ovviamente, sono altri): parte la
settimana di lavoro, ognuno ha i suoi task, i \emph{task owner} hanno scelto
che cosa fare, e questo prosegue per tutta la durata della sprint.

Durante l'implementazione si fa un \textbf{check dei progressi}. Non serve a
controllare che cosa stanno facendo le persone: di base bisogna
\textbf{fidarsi} che, se il programmatore dice ``farò questo'', lo farà entro i
limiti della sprint. Cristian parla qui di \emph{daily scrum} perché è ciò che
si dovrebbe fare in generale, anche se il suo team non li fa ogni giorno
(sezione~\ref{sec:z3-team}). Oltre alle solite domande (che cosa hai fatto
ieri, che cosa farai oggi), \textbf{la cosa più importante è trovare i
blocker}.

\subsection{Agire subito sui blocker}

Se il programmatore dice ``non ho potuto implementare il personaggio perché mi
manca l'asset, la \emph{skeletal mesh} del protagonista'', bisogna
\textbf{agire subito}, in due modi:
\begin{enumerate}
  \item \textbf{dare altri task} al programmatore, perché altrimenti resta lì
        ad aspettare, e questo non lo si vuole assolutamente;
  \item \textbf{capire che cosa è successo} alla skeletal mesh, parlando con
        la 3D artist. Se è solo questione di tempo, non è davvero un problema:
        è normalissimo aver assegnato l'implementazione prima che l'asset
        fosse pronto (``è normale sbagliare, l'importante è rimediare''). Se
        invece anche lì c'è stato un blocco, bisogna \textbf{risalire a
        cascata} fino a trovarlo.
\end{enumerate}
Magari si scopre che il problema sono proprio i game designer: quando hanno
fatto la \emph{mood board} non sono stati abbastanza chiari, e la 3D artist ha
dovuto spendere una settimana a cercare reference per conto suo, mentre il
designer pensava che quelle fornite bastassero e non si è posto il problema
della fase di ricerca delle reference da parte di chi fa la grafica o il 3D
(Cristian parla al femminile perché la 3D artist del suo team è una donna, ma
il discorso è generale). Anche questo, come l'apprendimento, è uno di quei
task che non bisogna trascurare.

\subsection{Analizzare i blocker prima di risolverli}

Il consiglio è \textbf{analizzare i blocker}, non risolverli appena si
presentano, perché si rischia di trascurare il problema alla radice. Cristian
richiama una delle prime lezioni del primo anno, quella su \textbf{bisogni e
problemi}: bisogna capire se il blocker è un problema o un bisogno, e agire di
conseguenza. Il guaio nasce quando lo si trascura o, peggio, quando lo si
risolve in fretta senza averlo considerato: ``poi vi si ritorcerà contro''.

\paragraph{L'esempio dei muri invisibili.} \emph{L'Inf} è un gioco di
esplorazione ambientale, per lo più all'aperto: su cinque livelli, circa
l'80\% del gioco è all'esterno. Per dare \emph{meaning} al livello finale,
quello della boss fight e della rivelazione narrativa su dove siano finiti
tutti gli esseri umani del pianeta, il team ha deciso di sottolinearne la
diversità rendendolo \textbf{al chiuso}.

La camera è in terza persona e abbastanza lontana dal giocatore, non classica.
Dare le collisioni alla camera sembrava fastidioso, perché a ogni rotazione
``sbatte'' contro i muri. La prima idea: \textbf{rendere invisibili i muri}
quando la camera ci passa dietro. Non si potevano far sparire tutti insieme
quando il giocatore era solo dietro il muro destro, quindi ogni muro aveva il
suo \emph{trigger}, che stabiliva quante e quali pareti dovessero sparire. Il
livello dura mezz'ora, con una buona decina di minuti di esplorazione al
chiuso, completamente circondata da pareti: \textbf{un mese di lavoro} solo
per assegnare i trigger, cambiare i materiali e gestire il tutto. ``Era l'idea
geniale che ci era venuta in quel momento.''

Quando hanno consegnato la beta e l'hanno fatta provare ai professori e ad altri
giocatori, tutti hanno detto più o meno la stessa cosa: i muri che spariscono
\textbf{non sono eleganti e rompono l'immersione}. Hanno tagliato tutto e
rimesso le collisioni alla camera. Se avessero analizzato bene il problema
prima, avrebbero capito che bastava \textbf{gestire meglio gli spazi}: non fare
curve strette ma allargarle, così che anche ruotando la camera il giocatore
tendenzialmente non sbatta contro il muro. Era ``assolutamente una questione di
riallargare più o meno tutti gli spazi''; ed è ciò che hanno fatto alla fine,
insieme alle collisioni. ``Ad averlo fatto prima avremmo risparmiato due
settimane di tempo''\footnote{Poco prima Cristian parla di ``un mese buono'' di
lavoro sui trigger (correggendosi da ``due mesi''); le ``due settimane'' sono
forse il risparmio netto stimato, tenendo conto del lavoro di allargamento
degli spazi che è stato comunque necessario. La trascrizione non permette di
chiarirlo.}. Ecco perché è importante analizzare i blocker.

\subsection{Bonus: i feedback costanti}

Cristian nota che, soprattutto ultimamente, in una fase di
\emph{lockdown}\footnote{Il riferimento al lockdown colloca la lezione nel
periodo della pandemia di Covid-19, coerentemente con l'accenno di Simone al
Covid e alle sedi chiuse (``quando si poteva andare'').} in cui tutti sono
costantemente in contatto su Discord, ``anche più di quanto lo saremmo nella
vita vera'', c'è un \textbf{bisogno costante di dare feedback} sul lavoro.

Lui è pro feedback in qualsiasi momento, anzi ``ossessionato'': all'inizio ha
chiesto alla classe di criticarlo e di dargli consigli sulle slide, perché è
l'unico modo per migliorare. Il problema è che \textbf{i feedback costanti
producono un calo di produttività enorme}. L'esempio è personale: Cristian è
anche il VFX artist del team, e sugli elementi estetici c'è sempre modo di
migliorare, quindi riceve di continuo feedback (``io lo farei più grande, rosso,
verde, giallo, a forma di triangolo, a forma di quadrato''). Sono feedback
correttissimi: chi li dà, soprattutto in produzione, non lo fa perché pensa
che tu sia scarso, ma per il bene del gioco, che vive e vuole ben prodotto. I
migliori mandano anche video di riferimento (``ho visto questo VFX in questo
gioco, sapresti rifarlo?''), e questo ben venga. Il problema è che siano
\textbf{costanti}.

Il consiglio: \textbf{limitare la richiesta e l'ascolto dei feedback a fasi
specifiche} (alla fine della sprint, oppure in un controllo intermedio). Dipende
anche da che cosa si sta facendo: con i VFX molte cose dette oggi sono
sistemate il giorno dopo; ma se chi costruisce un livello, in block out o in
\emph{environment}, riceve mille feedback per ogni metro quadro (``un po' più a
destra, un po' più a sinistra''), il lavoro si ferma lì. Sia perché chi lo fa
si \textbf{demoralizza}, sentendosi costantemente criticato, sia per una
questione di fiducia, tornando al discorso di Simone: ``fidati di me come io mi
fido di te. Fammi completare il task; poi, quando ti chiedo il feedback, o più
avanti, è ben accetto, anzi ben venga''.

Cristian si scusa per l'apparente presunzione di definire ``fondamentali'' i
propri consigli: lo fa perché, se li avesse avuti prima, si sarebbe risparmiato
tantissimi problemi. Diego scherza che andranno nelle slide ``con la mia faccia
super inquietante''; Cristian dice che può passare anche le sue slide.

% ---------------------------------------------------------------------
\section{Il calcolo della velocity}
\label{sec:z3-calcolo}

\subsection{Confrontare il fatto con lo stimato}

Sulla parte pratica Simone è stato più esaustivo, quindi Cristian si concentra
sull'interpretazione\footnote{Un breve tratto della spiegazione iniziale è
perso nella trascrizione, che riprende a metà frase (``nel comparare quello
che è stato fatto da quello che è stato stimato'').}. La velocity consiste nel
\textbf{confrontare ciò che è stato fatto con ciò che era stato stimato}, in
numero di task o in pesi, a seconda di come la si gestisce. Da lì si capisce
che cosa è successo e che cosa si può prevedere:
\begin{itemize}
  \item \textbf{fatto $>$ stimato}: si può prevedere di più; nella sprint
        successiva si aumentano i pesi o il numero di task;
  \item \textbf{fatto $<$ stimato}: o si è esagerato con la stima (succede
        sempre, così come succede sempre il caso opposto), oppure ci sono stati
        problemi gravi o \emph{blocker} nella produzione, ed è fondamentale
        capire quale delle due cose sia successa;
  \item \textbf{fatto $=$ stimato}: succede quasi mai, ma è l'obiettivo. Le
        possibilità sono due: si è trovato il proprio \textbf{bilanciamento}, e
        si può continuare così; oppure \textbf{non si sta applicando bene
        Scrum}, ad esempio perché si calcola la velocity su ciò che si è fatto
        invece che sulla stima.
\end{itemize}

\subsection{Più di un numero}

L'obiettivo non è arrivare a una velocity pari a 1, intesa come rapporto tra
fatto e stimato (100 punti stimati, 100 fatti). La velocity è \textbf{molto più
di un numero}: dà almeno tre informazioni.

\paragraph{1. La bontà degli strumenti di stima.} Planning poker, Fibonacci e
tutti gli strumenti di stima sono utilissimi, ma \textbf{vanno testati}: la
velocity mostra se ci sono metodi sbagliati di stima dei pesi o del numero dei
task.

\paragraph{2. La raggiungibilità dello scope.} Se per tutta la produzione la
velocity resta, ad esempio, a 0,8, sempre al di sotto di quanto dovrebbe essere
(in teoria non dovrebbe succedere, perché ci si dovrebbe tarare), e si sono
previsti dieci livelli, si capisce già da subito che \textbf{non ci si
arriverà mai}, a meno di aggiungere uno o due anni di sviluppo.

\paragraph{L'esempio degli otto livelli.} Finita la \emph{vertical slice}, il
team ha ricontrollato tutta la preproduzione e impostato il gioco come un gioco
completo, non come la creazione di una parte: c'è tutta la storia, ci sono
tutti i livelli, tutto ciò che fa il giocatore dall'inizio alla fine, chi
incontra e quando; tutto scritto, ma tarato e lasciato \textbf{flessibile},
``altrimenti ricadiamo nella trappola del Waterfall''.

A settembre, all'inizio dell'anno (``per farvi capire quanto il processo è
rapido''), avevano previsto \textbf{otto livelli}. Sembrava una stima
credibile: il vecchio team, con due level designer che volevano fare
esattamente quel lavoro, produceva un livello in \textbf{due o tre settimane}, a
una qualità abbastanza elevata: magari mancavano l'illuminazione, il
\emph{foliage} e alcuni dettagli, ma il livello era \textbf{giocabile}, che è
la parte fondamentale. Cristian, unico rimasto del vecchio team, si è basato su
quella stima: ``da qualche parte dovevo cominciare''.

Ma i level designer del nuovo team \textbf{non sono level designer puri}: sono
soprattutto game designer, interessati al game design, che hanno deciso di
sperimentare con il level design per necessità (il game design serve sempre, ma
se manca un level designer, visto che il percorso di studi insegna sia level
sia game design, ci si adatta). C'è chi si è adattato a fare il level designer,
e chi ha deciso proprio quell'anno di specializzarsi nel level design: ``perché
avrei dovuto dirgli no, tu non sei un level designer? Anzi, ben venga''.
L'unica cosa è stata prevedere che avesse più bisogno di apprendimento: magari
sbaglierà di più, dovrà informarsi su come si fanno i livelli, avrà bisogno di
più feedback. L'importante è capirlo all'inizio.

Guardando l'andamento del primo livello, Cristian vedeva che \textbf{in un mese
avevano a malapena raggiunto il block out}. Ha convocato una \textbf{riunione
di emergenza}, lunghissima (non una revisione della sprint da venticinque
minuti come quelle di Simone, ma un altro tipo di riunione). Il suo obiettivo
non era tanto convincerli che lo \emph{scope} era troppo alto (ci è riuscito),
quanto \textbf{assicurarsi che non la prendessero come una sconfitta
personale}. È stato difficile ma interessantissimo: non se la sono presa né con
se stessi, per non riuscire a performare, né con lui, che stava di fatto
dicendo ``non ce la farete''. Anzi, nonostante il taglio, sono stati
\textbf{loro stessi a proporre la soluzione}: ``ok, abbiamo capito; facciamo
cinque livelli, e se vediamo che arriviamo (speriamo a fine marzo) con tutti i
livelli, poi possiamo decidere di aumentarli o di cambiare parti del gioco''.
Cristian l'aveva già in mente, ma il fatto che l'abbiano proposta loro l'ha reso
``veramente un successo''.

Tutto questo grazie a un numero. Se Cristian fosse arrivato tra un mese e
mezzo, dopo che avevano già fatto quattro livelli ``con l'acqua alla gola'',
dicendo ``ne dobbiamo fare altri quattro entro marzo'', e poi a fine marzo
``basta, non ce la facciamo'', sarebbe stato molto più grave\footnote{Il
passaggio è confuso nella trascrizione (``mi ha portato ad affrontare queste
situazioni che affrontate a fine marzo'', ``non potne quattro''): il senso
ricostruito è che la velocity ha permesso di anticipare a inizio produzione una
decisione che, presa a fine marzo, sarebbe stata molto più dolorosa.}.
``Ringraziamo in questo caso la velocity.''

\paragraph{3. Lo stato del team.} La velocity permette di capire lo
\textbf{status del team e le sue performance}, perché si basa su una stima fatta
a partire dalla settimana precedente. Si ricollega al ``progettare il \emph{pace}
dello sviluppo'' (sezione~\ref{sec:z3-selezione}): ci si può basare proprio sulla
velocity per capire qual è il ritmo attuale e prevedere i prossimi. ``Tutto
torna.''

\subsection{Bonus: mai tornare indietro}

Un bonus che ``più che un bonus è un'altra lezione'': \textbf{prevedete sempre
task che portano a un incremento del prodotto}. La produzione può rallentare,
perché possono esserci difficoltà, ma \textbf{mai, mai, mai andare indietro}.

È un'altra lezione insegnata da \emph{L'Inf}, legata alla stesura del backlog.
Quando Cristian aveva scritto tutto il backlog (``400 task: come vedete ogni
tanto dico 400, 600, 300, giusto per dire che erano veramente tanti'') e si
accorgeva di aver sbagliato, invece di adattarsi e appoggiarsi alla produzione
in corso \textbf{fermava la produzione} (``ditemi pure che ero folle: lo ero'')
e riscriveva da capo tutti i task, arrivando magari a mille o riducendoli.

Lo stesso vale per i \textbf{documenti}. ``La documentazione è bellissima'':
Cristian era un grande fan della documentazione (Diego può confermarlo), ma
lavorando con Scrum, per ragioni di produttività, ha dovuto \textbf{scendere a
compromessi}: darsi un limite su quanto un documento è bello e ricco di
dettagli. Il motivo è sempre l'esempio della prima meccanica, l'assorbimento e
sparo: tutti i GDD del gioco si basavano su di essa, con tutte le interazioni,
anche secondarie (che cosa si poteva assorbire? Nemici, oggetti, chiavi, cose
per recuperare vita, bonus, malus, qualsiasi cosa). Tagliata la feature, tutti i
GDD sono stati ``letteralmente presi e buttati'', e Cristian ha deciso di
rifarli tutti da capo.

Questo approccio potrebbe funzionare in Waterfall (``finiti i documenti,
iniziamo la produzione''), ma l'obiettivo era \textbf{consegnare un videogioco,
non fare documenti}; e soprattutto il gioco era innovativo. È successo altre due
volte: in tutto Cristian ha riscritto \textbf{backlog e GDD da capo più
volte}\footnote{Nella trascrizione: ``è successo altre due volte quindi per ben
quattro volte io ho riscritto tutto il backlog e tutti i GDD''. Il conteggio non
torna del tutto (una volta più altre due fa tre); forse Cristian conta anche la
prima stesura.}, riportando ogni volta la produzione all'inizio (``ok, possiamo
iniziare il prototipo; no, è sbagliato, ricominciamo''). Lo racconta per far
capire quanto ci ha messo a entrare nel mindset agile, e per tranquillizzare la
classe: non indovinarci subito è normale, ``succederà il 90\% delle volte''.

Quando poi hanno fatto il prototipo, si sono accorti che \textbf{tutta la
documentazione era sbagliata}: con un gioco innovativo non si poteva sapere che
cosa avrebbe funzionato senza prima vederlo, perché non c'erano reference:
nessuno aveva fatto un gioco d'avventura \emph{match-3}. ``Il \emph{match-3}
esiste, i giochi d'avventura esistono, mischiarli insieme è una follia'', e lo
dice ``per criticarmi, non per vantarmi, perché non lo farò più''.

% ---------------------------------------------------------------------
\section{Il tuning della produzione}
\label{sec:z3-tuning}

\subsection{Che cosa e quando}

Il \emph{tuning} della produzione si basa sull'analisi fatta: dove ci si trova,
qual è lo stato delle cose, che cosa si potrebbe fare meglio e stimare
meglio\footnote{Anche qui l'inizio della spiegazione è parzialmente perso nella
trascrizione (``che in parte si basa su ci si trovano, qual è il loro
status'').}; per poi tornare all'inizio del loop con più dati su cui basare una
stima corretta.

Si può fare \textbf{a ogni sprint}, a seconda della durata, oppure \textbf{ogni
serie di sprint}: ``vedetelo voi''. Cristian suggerisce la seconda opzione
perché ritoccare continuamente il backlog per tagliare o aggiungere feature va
contro il principio del lavorare sempre per incrementare il prodotto. Va da sé
che, se si tratta di un taglio molto importante, \textbf{prima lo si prende e
meglio è}: ``ragionateci bene''.

\subsection{Su che cosa basarsi}

\paragraph{Sulle performance, cioè sulla velocity.} Se la velocity è alta, si
può dire: ``a questo punto possiamo reintrodurre qualche feature che era nei
\emph{nice to have}\footnote{Nella trascrizione ``nice to the epoch'': quasi
certamente \emph{nice to have}.} o che era addirittura stata tagliata''. Se è
bassa, si sistemano già milestone e backlog in base al risultato.

\paragraph{Sui feedback.} A Level Up \emph{L'Inf} ha ricevuto tantissimi
feedback, per lo più da giocatori. Trattandosi di una \emph{vertical slice},
spesso dicevano cose che il team già sapeva e che avrebbe implementato nella
sprint successiva (ma, per questioni di tempo, da qualche parte bisogna
tagliare). Altri, invece, hanno parlato di \textbf{accessibilità}.

\paragraph{Il cambio di camera.} A questi si sono aggiunti i feedback costanti
di Diego sulla camera: più che un ``dovete cambiarla'', un ``perché non provate a
cambiare camera?''. Diego interviene: ``c'è voluto del tempo a convincere
Cristian: otto mesi''. Alla fine l'ha fatto e sta funzionando, anche se ha
pesato molto sul team. Il motivo: il \textbf{twin-stick shooter}\footnote{Nella
trascrizione ``twin steep shooter''.} richiede una grande
\textbf{dimestichezza con i controlli} (muoversi e mirare contemporaneamente con
le due levette) e di per sé esclude gran parte del pubblico casual. I giocatori
più casual dicevano proprio questo: ``non riesco a camminare e mirare
contemporaneamente''. Il team se l'è segnato, ne ha discusso e alla fine ha
accettato il cambio di camera, con tutte le conseguenze.

\subsection{Flessibilità anche nella struttura dei documenti}

Un promemoria sulla \textbf{flessibilità}: più si è stati flessibili all'inizio,
più si riesce a rimettere mano a backlog e produzione in seguito. E questo vale
anche, banalmente, per \textbf{come è strutturato un documento}: se ogni casella
fa riferimento a un'altra casella, che fa riferimento a un'altra due righe
sopra, modificare quel documento (o un'implementazione costruita allo stesso
modo) richiederà molto più tempo. Flessibilità in tutto: nella scrittura dei
documenti e nella produzione (``non lo ripeto, altrimenti ricomincio'').

\subsection{Evitare il perfezionismo}

``A proposito di perfetto'': il consiglio sul \emph{tuning} è \textbf{evitare
in tutti i modi il perfezionismo}. Cristian ci sta ancora lavorando
``tantissimo'' (``anche qui Diego lo sa meglio di me''): è migliorato molto e nei
progetti ci riesce, ma in altre cose pretende ancora troppo. Il perfezionismo
\textbf{va contro la flessibilità e crea inevitabilmente dei blocker}. Lo sta
vivendo in prima persona, non per sé ma con altri membri del team.

\paragraph{Il block out con il foliage.} Il task è ``fai il block out di questa
parte del livello'' (dare l'intero block out di livelli da mezz'ora richiede
comunque tempo, quindi il task riguarda la parte iniziale). E nel block out
Cristian si ritrova il \emph{foliage}, il \emph{landscape} scolpito bene, le
rocce. ``Va benissimo perché ti piace il lavoro\dots{} anzi no, sto proprio
dicendo che non va bene.'' L'approccio di voler consegnare il task nel miglior
modo possibile è comprensibile; il problema è che questo perfezionismo
\textbf{blocca tutti gli altri}, per moltissimi motivi. Se si passano due ore a
mettere il foliage e poi, quando si va in gold a rifinire il livello, si trova un
foliage migliore per il \emph{mood} estetico del gioco, quel lavoro va buttato.
Il perfezionismo crea \textbf{buchi di produzione} che alla lunga diventano
letteralmente mesi di lavoro buttato o da buttare: un peccato per chi lo fa e
per chi lavora con lui.

\paragraph{La pressione sugli altri.} Un'altra conseguenza: il perfezionismo
\textbf{mette pressione a tutti gli altri}. Chi lavora alla seconda parte del
block out, con il task di renderla giocabile, vede che la prima parte ha
foliage, rocce e tutti i props (non l'illuminazione, che negli ambienti aperti
è tendenzialmente globale, salvo situazioni specifiche) e si sente in dovere di
raggiungere quel livello di rifinitura\footnote{Nella trascrizione ``si sentirà
in dovere di raggiliere la pozione della vita'': passaggio incomprensibile, il
senso è ricostruito dal contesto.}, trascurando la parte di gameplay, che invece
nella prima parte era stata fatta. Il perfezionismo di uno ha generato
\textbf{ansia} nell'altro, ed è da evitare in ogni modo.

\paragraph{I criteri di accettazione.} È un altro motivo per cui sprint e task
devono avere un \textbf{goal} e un \textbf{criterio di accettabilità}, per
usare i termini di Scrum: il task si considera fatto quando è \emph{done}, e va
bene così, perché lo si è valutato prima. Tutto ciò che si fa oltre il
\emph{done} è \textbf{tempo perso}. Può andare bene che quello che fai adesso
non dovrai farlo dopo, ma è sbagliato come approccio: se il lavoro è stato fatto
bene prima, si può andare avanti tranquilli, e al passo successivo ci si penserà
dopo.

% ---------------------------------------------------------------------
\section{Domanda: come comunicare un taglio al team}
\label{sec:z3-qa-taglio}

Prima della seconda parte, che Cristian promette più breve, una domanda: sul
piano relazionale ed emotivo, come ha gestito il dover convincere, o comunque
spiegare, al team la necessità del taglio da otto a cinque livelli?

\paragraph{Ricordare come funzionano i feedback.} ``Ottima domanda.''
Emotivamente è stato pesante, perché non tutti reagiscono ai feedback allo
stesso modo. Per questo, a ogni riunione, Cristian ricorda i principi dei
feedback (``so che sembra banale, so che non serve''), in particolare una
clausola che tutti dimenticano: \textbf{il feedback si può accettare o
rifiutare}. Se non si è d'accordo, lo si può ignorare; invece tutti partono dal
presupposto che un feedback vada seguito al 100\%\footnote{Nella prima lezione
zero il relatore invitava ad accettare il feedback ``a prescindere''
(sezione~\ref{sec:z1-ricevere}); lì però ``accettare'' significava ascoltarlo
senza difendersi, non necessariamente applicarlo. Le due indicazioni sono
quindi compatibili: si ascolta tutto, poi si decide che cosa farne.}.

\paragraph{Partire dai dati.} Sembra brutto affrontare una questione emotiva
partendo dai dati, ma l'obiettivo era \textbf{far arrivare loro alla soluzione},
non dargliela. Ed è andata così: sono stati loro a dire ``per adesso facciamo
cinque livelli, e poi, \textbf{se ci riusciamo}, implementiamo gli altri''.
Quel ``se ci riusciamo'' era ciò che Cristian voleva sentire: se fosse rimasto
solo ``ne facciamo cinque'', l'avrebbero vissuto come un limite su se stessi
(``non ho fatto in tempo a farne otto''). Il ``se ci riusciamo'' gli ha
comunicato che avevano capito che \textbf{non era una critica al loro modo di
lavorare}, né un modo per dire ``non siete capaci'': lo aveva sottolineato più
volte, magari non con parole così dirette, ma ``un po' \emph{soft-skillato},
come diciamo noi designer''.

Hanno colto che non c'erano alternative (``otto non li facciamo, punto''), ma
anche le possibilità: facciamo cinque livelli, abbiamo altro tempo,
perfezioniamo ciò che abbiamo fatto; vogliamo aggiungere una nuova modalità di
gioco? Benissimo. Altri due livelli? Vediamo, in base al tempo. Esserci
riuscito lo ha reso ``super fiero''.

% ---------------------------------------------------------------------
\section{Il team management}
\label{sec:z3-team}

\subsection{Senza team non esiste Scrum}

Cristian accelera, per stare nei tempi. Il punto da capire è che
\textbf{Scrum si adatta al team}: quindi il team è la cosa più importante di
tutte, più di Scrum stesso. ``Se io mi adatto al team, il team deve esserci.
Senza team non esiste Scrum.''

\subsection{Una produzione senza stress}

Di conseguenza, una produzione in Scrum \textbf{dovrebbe rendere il team
rilassato}. Come il team di Simone, anche il primo team di Cristian (quello fino
alla \emph{vertical slice}) \textbf{non ha mai fatto crunch}. Il nuovo team sì,
per i grossi cambiamenti già raccontati: la camera è cambiata completamente, e
``cambiare la camera in un gioco è fare un altro gioco''. Cristian aveva già
spostato la scadenza (``invece di finire a marzo finiamo ad aprile''), ma il team
si è sentito comunque in dovere di risolvere prima lo scompenso creato da quei
cambiamenti.

In generale ci si dovrebbe sentire meno stressati per un motivo semplice: chi
guida sta dando i task che lo sviluppatore ha chiesto, verbalmente o dimostrando
di saperli fare. ``So che in due settimane devo fare cinque cose, faccio queste
cinque cose'': non c'è la pressione di dover fare il più possibile. Se poi uno è
perfezionista e passa da cinque cose a quindici perché ``ha tempo'', il problema
(``sembra brutto dirlo'') è suo, nasce nel singolo individuo; e chi guida deve
basarsi sul team, sulle performance del team.

\subsection{Mai la velocity sui singoli}

Da qui un consiglio sulla velocity: \textbf{non calcolate mai, mai, mai la
velocity delle singole persone}, e nemmeno i task fatti dai singoli. Lo si può
fare per sé, ma tenendolo ``chiuso in un'agenda scritta a penna'': se un membro
del team scopre che si stanno valutando le sue prestazioni individuali,
\textbf{si sente giudicato} e soprattutto non sa più come dire che non vuole
fare una cosa.

A Cristian è successo: ``mi hanno sgamato''. Si è dovuto scusare, anche se non
erano scuse: ha detto la verità, cioè che lo faceva più per capire lui la
produzione che per valutare le loro performance. Ma spiegarlo è complicato,
quindi è meglio non farlo: tanto la velocity si basa sulle performance del team,
e i dati individuali non servono. Se li si vuole per sé, per capire a chi dare
un task piuttosto che a un altro (se una persona in un'ora fa dieci
implementazioni e un'altra una, ci si regola quando si selezionano i task),
non bisogna assolutamente farli ricadere sulle persone.

\subsection{Primo aspetto: adattare Scrum al team}

Per una produzione senza stress servono due aspetti chiave. Il primo è che
\textbf{Scrum va adattato al team}: si prende la base di Scrum (quella di Diego,
dei libri, di Simone e di Cristian), la si \textbf{testa}, si \textbf{analizzano}
la velocity e tutto il resto, la si \textbf{tara} (``\emph{daily scrum}? No,
facciamolo ogni tre giorni, ogni settimana, o non facciamolo''), la si
\textbf{riprova}, e si ripete il ciclo per tutta la produzione. Alcuni esempi da
\emph{L'Inf}:
\begin{itemize}
  \item i \textbf{daily scrum} li fanno \textbf{ogni tre giorni}, ``se va
        bene'', e dipende anche dal tipo di produzione;
  \item la \textbf{lunghezza delle sprint} è flessibile (una, due o tre
        settimane) e dipende:
        \begin{itemize}
          \item dalla \textbf{disponibilità del team}: se tutto il team non
                c'è, è inutile fare una sprint di una settimana;
          \item soprattutto dalla \textbf{fase di produzione}: quando si
                implementano script e in una settimana si riescono a
                implementare tantissime meccaniche, la sprint dura una
                settimana, al termine della quale si controlla ciò che è stato
                fatto e si sceglie che cos'altro implementare; quando si passa al
                level design, con lavori molto dispendiosi in termini di tempo,
                le sprint si diluiscono, e infatti in quel momento sono a sprint
                di due settimane, concentrati sul completamento dei
                livelli\footnote{Nella trascrizione ``sull'ultima mento,
                sull'ultima azione, non mi ricordo, scusate. Per completare il
                livello'': forse ``l'ultima milestone''.};
        \end{itemize}
  \item la velocity si calcola \textbf{sui task e sul loro peso, non sulle ore
        lavorative}, come diceva Simone. Soprattutto per degli studenti, che
        non hanno ore fisse. In un ambiente di lavoro in cui tutti sono
        disponibili dalle 9 alle 18 potrebbe cominciare ad avere senso:
        ``secondo me no, ma ve lo dico per correttezza''.
\end{itemize}

\subsection{Secondo aspetto: tutti devono conoscere il mindset}

Il secondo aspetto è che \textbf{tutti i membri del team devono conoscere il
mindset di Scrum}\footnote{A questo punto Cristian chiede a Diego cinque minuti
in più per chiudere bene, e Diego acconsente.}. Per due motivi:
\begin{itemize}
  \item se tutti lo conoscono, il rischio di cadere nel Waterfall o nel
        perfezionismo si riduce, perché ognuno sa che cosa è giusto e che cosa
        è sbagliato;
  \item tutto il team ha \textbf{lo stesso goal}: se la sprint chiede il block
        out, tutti sanno che, secondo la mentalità agile, va bene farlo ma non
        deve essere perfetto. Bisogna fare i \emph{done}, non i
        \emph{perfect}.
\end{itemize}
Come far capire a tutti la mentalità di Scrum? Cristian usa due modi.
\begin{enumerate}
  \item \textbf{Parlarne costantemente}, a ogni riunione, come sta facendo con
        la classe: sia per migliorare lui stesso e farsi entrare il mindset,
        sia per trasmetterlo agli altri.
  \item Il più efficace: \textbf{lo shipping}. Il team riesce a consegnare
        quasi costantemente. Ci sono momenti in cui un problema cade
        all'ultimo, come in quella settimana, con alcuni strumenti
        particolari che li hanno rallentati, ma si risolve tutto. Dimostrando
        di riuscire a \textbf{mantenere le milestone}, il team ha acquisito
        fiducia in lui e ha capito quanto è potente Agile. Nella
        \emph{vertical slice} consegnavano le build addirittura prima di
        quanto richiesto; e la consegna prima di Natale, nota da settembre, è
        stata rispettata nonostante il cambio di camera e tutto il resto. Il
        messaggio: \textbf{non è importante che ciò che si consegna sia
        perfetto, ma che rispetti le condizioni di accettabilità} che ci si è
        dati.
\end{enumerate}
Sono i due metodi che funzionano meglio finora; se ne troverà altri, li
condividerà: ``tutto è in continua evoluzione''.

\subsection{Il post-mortem}

Il \emph{post-mortem} è \textbf{la garanzia di capire se ciò che si sta facendo
è corretto}: serve a vedere come stanno andando il team e la produzione, se il
mindset di Scrum è stato compreso o no, e ad agire di conseguenza.

\textbf{Quando farlo} lo si decide liberamente: a ogni sprint, a ogni gruppo di
sprint, a ogni milestone.

Ciò che Cristian sottolinea è che \textbf{nel post-mortem devono parlare tutti},
anche lo sviluppatore che non vuole mai parlare: bisogna dargli il tempo per
farlo. Un modo è \textbf{organizzarlo in anticipo}: se è tra una settimana, lo si
comunica la settimana prima. Cristian prepara sempre un \textbf{documento della
riunione} prima e lo aggiorna dopo, con una sorta di elenco: sembra burocratico,
ma aiuta a mantenere il focus. Bisogna far capire a ciascuno che tra una
settimana parlerà, e non può sottrarsi (se si assenta ci perde lui). E bisogna
dirgli di \textbf{prepararsi} qualcosa da dire, magari scrivendolo o facendo una
piccola revisione: se ti viene in mente una cosa da dire la settimana dopo e non
la scrivi, finirai per pensare ``sai che c'è, non la dico''; se l'hai scritta,
non puoi dimenticartela, e hai qualcosa di cui parlare. È fondamentale perché
bisogna tenere conto di \textbf{tutti i tipi di personalità}: non solo i
giocatori hanno personalità diverse, anche gli sviluppatori hanno personalità
con bisogni diversi.

\subsection{La sincerità}

Cristian propone di tagliare la parte ``sdolcinata'', ma Diego vuole sentirla
per intero (``non si sfugge al miele''; ``poi mi manderete gli IBAN e vi pagherò
il tempo perso''). Come Diego sa, Cristian cerca sempre di metterci ``qualcosa di
filosofico''. Per \textbf{sincerità} intende che, in qualsiasi momento della
produzione, \textbf{tutti i membri del team devono sentirsi protetti in ciò che
dicono} e liberi di esprimere opinioni e feedback. È importantissima sia nella
produzione sia nel team.

\paragraph{Nella produzione.} Senza sincerità la produzione va a picco; si può
comunque consegnare, non è obbligatoria, ma con la sincerità si realizza un
prodotto migliore. Se a inizio produzione si pensa che una meccanica ``faccia
schifo'' (prende un \emph{player type} completamente fuori target e richiede
troppo lavoro per essere implementata), \textbf{è quello il momento di dirlo}.
Altrimenti ``il karma vuole che tu stesso che pensi questa cosa sarai colui che la
andrà a implementare'', e la farai male, perché non vuoi farla. Se non l'hai
detto, o l'hai detto troppo tardi, ``ti mangerai le mani''.

E il \textbf{``ve l'avevo detto''} va risparmiato: anche se il giorno dopo arriva
la conferma di ciò che si era detto, anche se è umano, va evitato, perché crea
attriti nel team che non si risolvono più (Cristian parla per esperienza
diretta). È lo stesso spirito del \emph{disagree and commit}
(sezione~\ref{sec:z2-disagree}).

\paragraph{In Scrum.} Se uno sviluppatore dice ``questa settimana so fare 15
task'' e ne fa 7, la prima volta passa. La settimana dopo dice di nuovo ``ne
faccio 15, tranquillo'', magari inventando scuse (``quel problema non ce l'ho
più''). Una volta, due volte; \textbf{alla terza bisogna chiedere sincerità}: qual
è il problema? Non hai tempo?

Perché sul tempo, ``ragazzi'', \textbf{è importante dire quanto tempo si ha, non
garantirlo}. Se il team (cinque persone, nel caso di Cristian) dice di avere due
ore al giorno, lui pianifica una produzione da dieci ore al giorno; se il team
dice di avere otto ore al giorno, pianifica di conseguenza (``80
ore''\footnote{Così nella trascrizione; con cinque persone a otto ore al giorno
il totale sarebbe 40 ore, oppure 80 su due giorni. Il senso dell'esempio non
cambia.}), e se poi in realtà le ore sono dieci, gli diranno ``ci dai troppi
task''.

Bisogna essere sinceri \textbf{prima di tutto con se stessi}: se un task richiede
due giorni, si dice due giorni; magari tre, perché un ``tempo cuscinetto'' ben
venga, ma lo si dice per com'è. E non bisogna proteggersi dal timore che ``se
dico la verità nessuno si fiderà più di me'' o ``mi daranno più task'': è compito
dello scrum master, ma di tutto il team, far capire che \textbf{si è tutti sulla
stessa barca}.

\paragraph{Nel team.} Gli attriti generati dalla mancanza di sincerità si
trasformano in \textbf{relazioni migliori} quando c'è sincerità. Le
incomprensioni vanno \textbf{risolte subito}, affrontandole; magari porteranno a
un litigio, se proprio quello è il finale, ma tra persone mature si dovrebbe
riuscire a risolvere in modo pacifico. Per questo, sulla slide, la sincerità è
definita \textbf{``rischiosa ma necessaria''}.

% ---------------------------------------------------------------------
\section{Chiusura della lezione}
\label{sec:z3-chiusura}

Cristian salta la slide di riepilogo (la lezione si potrà rivedere online) e
lascia i suoi contatti: \textbf{Telegram}, dove molti studenti lo trovano già nei
gruppi (in particolare quelli dell'accademia online), oppure l'\textbf{email},
che Diego può fornire. Si rende disponibile per qualsiasi cosa, sulla lezione o
su un'idea da sottoporgli (``ho pensato così, secondo te può funzionare?''),
precisando di non volersi sostituire in alcun modo a Diego. Si scusa per i dieci
minuti finali ``di corsa''.

Diego lo ringrazia per la disponibilità, ricordando che anche lui, da studente,
ha un carico di lavoro notevole. L'obiettivo della lezione era far vivere alla
classe un'esperienza e dare informazioni simili a quelle di chi ha fatto lo
stesso percorso, ma ha \textbf{internalizzato Agile} meglio di altri: tutti hanno
avuto le stesse lezioni, ma Cristian e Simone sono quelli che le hanno capite
meglio. Ci sono molti \emph{insight} utili in entrambi gli interventi, e i punti
chiave di Agile sono stati ripetuti più volte, ``perché alla fine Agile è
comprendere i punti chiave''. ``Grazie ancora, e buon pranzo.''\footnote{La
trascrizione si chiude con la parola ``grazie'' ripetuta diverse centinaia di
volte: è un artefatto tipico della trascrizione automatica su un tratto di
audio silenzioso o di saluti sovrapposti, non un contenuto della lezione.}

% ---------------------------------------------------------------------
\section*{Punti chiave del capitolo}
\addcontentsline{toc}{section}{Punti chiave del capitolo}

\begin{itemize}
  \item \textbf{Agile è un mindset, non un framework}: non ha regole, ha una
        logica da \textbf{internalizzare}, come una lingua che a un certo punto
        si smette di tradurre. I quattro valori non vanno seguiti alla lettera:
        vanno adattati al team e al progetto, e non è detto che Agile serva.
  \item I \emph{framework} (Scrum, Kanban, \dots) hanno regole e si basano
        sull'\textbf{inspect and adapt}, cioè sulle iterazioni. Si parte
        rispettandoli al 100\%, si capisce che cosa non funziona e perché, e
        \textbf{ci si costruisce il proprio}.
  \item \textbf{Meno riunioni possibili}: unire planning e post-mortem, stare in
        25--30 minuti, report da un minuto. \textbf{Nelle riunioni non si
        risolvono i problemi}: si fissa un appuntamento con chi può aiutare. Un
        \textbf{timekeeper} può far rispettare i tempi.
  \item Le \textbf{sprint} (una o due settimane) sono l'unità di tempo: le
        milestone si misurano in sprint, non in mesi.
  \item \textbf{Consistency}: una cerimonia saltata fa crollare le altre a
        effetto domino. Per questo ciò che non funziona per il team (spesso il
        \emph{daily scrum} o il \emph{planning poker}) va \textbf{tagliato},
        non trascinato.
  \item Il \textbf{backlog} si fa \textbf{a scaglioni}, per milestone, partendo
        dalle \emph{epic} e dai documenti (visione, GDD, \emph{versioning} con
        tutti gli ambiti, LDD). Si organizza per feature, per ambito con un
        \textbf{color code} condiviso, per milestone e sotto-milestone o per
        priorità, tutto in un solo strumento.
  \item Ogni task (e ogni sprint) ha un \textbf{obiettivo verificabile} e un
        \textbf{criterio di accettazione}: ``completare il livello 5'' non lo è.
        Le \emph{user story} aiutano, ma non sono obbligatorie.
  \item \textbf{Mai stimare in ore}: si stima in modo \textbf{relativo}, con la
        sequenza di \textbf{Fibonacci}, partendo dagli estremi. I valori sono
        arbitrari ma \textbf{consistenti}, e quindi utilizzabili.
  \item La \textbf{velocity} (somma dei pesi dei task \emph{done} per sprint,
        confrontata con lo stimato) \textbf{prevede il futuro}: dice se gli
        strumenti di stima funzionano, se lo \emph{scope} è raggiungibile e come
        sta il team. Permette di tagliare per tempo e di arrivare in fondo
        \textbf{senza crunch}. L'andamento a dente di sega è normale.
  \item Le prime sprint \textbf{si sbagliano tutte}, ed è giusto: la quantità si
        tara sulle disponibilità e sulle capacità attuali, prevedendo i
        \textbf{tempi di apprendimento}.
  \item Si può \textbf{progettare il ritmo dello sviluppo}: dopo periodi di
        lavoro invisibile, sprint che danno soddisfazione (nuovi nemici, VFX,
        \emph{juiciness}) per il morale del team.
  \item I \textbf{blocker} vanno trovati subito e \textbf{analizzati alla
        radice} prima di essere risolti: la soluzione più rapida (i muri
        invisibili) può costare settimane.
  \item I \textbf{feedback costanti} abbattono la produttività: vanno raccolti in
        momenti specifici. Il feedback si può accettare o rifiutare.
  \item \textbf{Flessibilità} in tutto, soprattutto negli elementi innovativi:
        non prevedere tutto nel dettaglio (è Waterfall), non fermare mai la
        produzione per riscrivere backlog e documenti, lavorare sempre per
        \textbf{incrementare il prodotto}.
  \item Il \textbf{perfezionismo} crea blocker, lavoro buttato e ansia negli
        altri: si fanno i \emph{done}, non i \emph{perfect}.
  \item \textbf{Il team viene prima di Scrum}. La velocity non si calcola mai sui
        singoli. Tutti devono conoscere il mindset: parlarne a ogni riunione e,
        soprattutto, \textbf{consegnare} rispettando le milestone.
  \item Nel \textbf{post-mortem} parlano tutti: va annunciato prima, con un
        documento, chiedendo a ciascuno di prepararsi.
  \item Alla base di tutto ci sono la \textbf{fiducia} reciproca e la
        \textbf{sincerità}: dire subito ciò che non va, dire quanto tempo si ha
        (non garantirlo), evitare il ``ve l'avevo detto''. La sincerità è
        \textbf{rischiosa ma necessaria}.
\end{itemize}

%    \part{Il corso di Game Design}
%    \renewcommand{\thechapter}{\arabic{chapter}}
%    \setcounter{chapter}{0}
%    % =====================================================================
%  Capitolo 1 — Introduzione al corso e al mestiere del game designer
%  Fonte: trascrizione della lezione GD1 (Game Design, primo anno)
%  Pacchetti richiesti nel preambolo del libro:
%    \usepackage[italian]{babel}, \usepackage[T1]{fontenc},
%    \usepackage[utf8]{inputenc} (solo pdfLaTeX), \usepackage{booktabs}
% =====================================================================

\chapter[Il mestiere del game designer]{Il mestiere del game designer: introduzione al corso}
\label{cap:introduzione-game-design}

Questo primo capitolo fa da soglia. Prima di affrontare metodi, framework e
documenti, è necessario capire \emph{in che modo} si impara il game design,
\emph{con chi} lo si fa, \emph{con quali strumenti} e, soprattutto,
\emph{che cosa fa concretamente} un game designer nel suo lavoro quotidiano.
Il capitolo segue l'impostazione della lezione inaugurale del corso: si parte
dal contesto formativo (il corso come simulazione di lavoro), si passa al
metodo di apprendimento (l'\emph{inquiry-based learning}), si elencano
letture e risorse di riferimento, e si arriva infine a una prima definizione
operativa del ruolo del designer, dei limiti entro cui lavora e delle sue
possibili specializzazioni.

Una premessa terminologica: il game design è una disciplina giovane e quasi
tutta la sua letteratura è scritta in inglese. Per questo motivo molti termini
tecnici verranno lasciati in lingua originale (\emph{pipeline},
\emph{pitching}, \emph{tuning}, \emph{feature team}, \dots): sono i termini che
si useranno nel lavoro, ed è importante imparare a riconoscerli e a usarli
correttamente fin da subito.

% ---------------------------------------------------------------------
\section{Il contesto: un corso come simulazione di lavoro}
\label{sec:simulazione-lavoro}

\subsection{Il principio fondamentale}

Il primo concetto da fissare è che il corso \textbf{si presenta come una
simulazione del lavoro}. Dal momento in cui si entra in aula si viene trattati
come \emph{potenziali professionisti dell'industria}: non soltanto studenti,
ma lavoratori effettivi di uno studio (l'accademia), con una gerarchia, dei
referenti, dei compiti e delle scadenze.

Ne derivano alcune conseguenze pratiche:
\begin{itemize}
  \item il docente non assegna un ``compitino'' da risolvere per accumulare
        punti o presenze: ogni \emph{task} è pensato per far dedicare tempo
        all'apprendimento di qualcosa di concreto;
  \item non consegnare un task non danneggia il docente (che ``vive ancora
        domani e viene pagato ugualmente''), ma lo studente, che perde
        l'occasione di mettere le mani su un argomento e diventare più abile;
  \item docenti e studenti devono trattarsi \textbf{da colleghi}, perché è
        molto probabile che un domani lo diventino davvero.
\end{itemize}

\subsection{Perché l'impegno conta più che mai}

Il lavoro del game designer è descritto come \emph{molto romantico, ma anche
molto duro}, e sempre più duro col passare degli anni. Durante il periodo del
Covid l'industria ha visto grandi investimenti; quei ``rubinetti'' sono stati
chiusi da tempo, il rapporto con i publisher è diventato più complesso e le
opportunità lavorative non sono più quelle di tre o quattro anni fa. Questo
significa:
\begin{itemize}
  \item più competizione per ogni posizione;
  \item aspettative più alte: ci si aspetta che i candidati siano ancora più
        preparati e sicuri delle proprie competenze;
  \item un unico modo per diventarlo: \textbf{lavorare}. Più tempo si dedica
        oggi ad apprendere il mestiere, più facilmente si sarà assumibili
        domani e più agevole sarà il lavoro di tutti i giorni.
\end{itemize}

\subsection{Sbagliare qui per non sbagliare fuori}

Il mantra del corso è: \textbf{quanto più provate qui, anche in modo goffo e
sbagliando, tanto meno sbaglierete là fuori}. In accademia, per quanto la
simulazione sia realistica, nessuno toglie lo stipendio e nessuno licenzia da
un momento all'altro. Nell'industria sì, e un errore può costare molto denaro
all'azienda, ancora di più se l'azienda è la propria.

\subsection{Il valore del fallimento}

Il fallimento, se affrontato con la mentalità corretta, è probabilmente
\emph{la cosa più vantaggiosa} che si possa fare durante la formazione:
\begin{itemize}
  \item un successo (``avete superato questa deadline'') conferma le proprie
        abilità, ma insegna relativamente poco;
  \item dopo $n$ fallimenti si impara moltissimo, e i successi arrivano;
  \item è fondamentale non demoralizzarsi: il feedback negativo non è mai un
        giudizio sulla persona, ma su ciò che è stato consegnato.
\end{itemize}

Il docente riconosce apertamente che dedicare tempo ed energia a un progetto e
vederlo rifiutato, o vederlo non uscire mai, è demotivante: è successo a lui e
lo ha visto accadere anche su progetti molto grandi, con team numerosi. Il
punto però non è ciò che si è ottenuto materialmente, ma \textbf{quanto si è
cresciuti}. Da ogni esperienza ci si portano a casa due cose:
\begin{enumerate}
  \item \textbf{l'esperienza}: quanto si è diventati bravi, veloci,
        praticamente abili nell'applicare un framework, quanto si è cresciuti
        come designer;
  \item \textbf{il rapporto umano}: nel lavoro si ``suda sangue'' fianco a
        fianco con altre persone, e questo crea legami molto forti che
        portano avanti l'industria nel senso più positivo del termine: nascono
        progetti privati, ci si porta da un'azienda all'altra, si formano team
        con basi solide.
\end{enumerate}

Va evitato anche l'eccesso opposto: la cultura del \emph{grinding} si è
appropriata di questi discorsi in modo distorto. Non bisogna diventare robot.

% ---------------------------------------------------------------------
\section{Organizzazione del corso e gerarchia}
\label{sec:gerarchia-corso}

\subsection{La struttura gerarchica}

La struttura dei referenti del corso ricalca volutamente una
\textbf{gerarchia tipica dell'industria}:
\begin{itemize}
  \item al vertice c'è il \textbf{direttore del corso} (Diego Ricciuti), che
        corrisponde al \emph{creative director}: supervisiona la visione del
        progetto (il concetto di ``visione'' verrà approfondito più avanti);
  \item sotto di lui ci sono i \textbf{referenti di dipartimento}, uno per
        disciplina, con un corrispettivo per ciascuna sede (Milano e Roma);
  \item in fondo ci sono gli studenti, che occupano il primo gradino di
        un'azienda: l'\emph{intern} (stagista) o il \emph{junior}, il passo
        che precede il lavoratore ``regolare''.
\end{itemize}

\begin{table}[htbp]
  \centering
  \begin{tabular}{@{}p{0.34\textwidth}p{0.29\textwidth}p{0.29\textwidth}@{}}
    \toprule
    \textbf{Disciplina} & \textbf{Milano} & \textbf{Roma} \\
    \midrule
    Direzione (\emph{creative director}) & \multicolumn{2}{p{0.58\textwidth}}{Diego Ricciuti} \\
    Game design      & Mattia Podini        & Giacomo ``Jack'' Danna \\
    Level design     & Stefano Zocchi       & Andrea Dispoli \\
    Narrative design & \multicolumn{2}{p{0.58\textwidth}}{Daniele Fusetto (lezioni condivise), Giacomo Masi} \\
    Programmazione (Unreal) & Alessandro Pregnolato & Matteo Velardi \\
    \bottomrule
  \end{tabular}
  \caption{Referenti del corso per disciplina e sede (nomi come riportati in lezione).}
  \label{tab:referenti}
\end{table}

Per il primo anno, il referente diretto degli studenti di game design è il
docente di game design: è a lui che ci si rivolge per dubbi e domande. Altre
figure utili sono il direttore didattico (severo, ``giustamente severo'', da
contattare solo per questioni particolarmente importanti), l'amministrazione
(contratti e documentazione) e la segreteria.

\subsection{I canali di comunicazione}

\begin{description}
  \item[Discord] è il canale per le comunicazioni ufficiali del corso e per
        le domande che non richiedono urgenza immediata; può avere qualche
        ritardo.
  \item[E-mail] è il canale per tutto ciò che è formale o ``ancora più
        ufficiale'' e per le comunicazioni rapide e immediate.
  \item[Telegram] esistono residui di vecchi canali precedenti a Discord, da
        evitare.
\end{description}

Le domande che valgono per tutti vanno fatte in classe (anche da un portavoce
a nome del gruppo); quelle singole e personali via Discord o, se formali, via
e-mail. Se il docente non risponde a una domanda c'è un motivo, ad esempio
ragioni legali, e lo dirà esplicitamente.

\subsection{I tre moduli e la settimana tipo}

Le tre discipline principali del primo anno (e non solo) sono:
\begin{itemize}
  \item \textbf{Game Design} (tendenzialmente il lunedì);
  \item \textbf{Level Design} (tendenzialmente il mercoledì);
  \item \textbf{Programmazione} (tendenzialmente il venerdì).
\end{itemize}
Nel primo mese le lezioni sono tutte di game design (lunedì, mercoledì e
venerdì). Sono previste anche lezioni di narrativa, ma il piano formativo si
aggiorna ogni anno: le ore di narrativa sono state ridotte a favore di quelle
di programmazione. Le poche lezioni di narrativa daranno quindi dei
\emph{cenni}, non la preparazione di un narrative designer. Ciò che serve
davvero è imparare \textbf{come comunicano gli altri dipartimenti} e come
lavorare con persone che fanno un lavoro diverso ma complementare al proprio.

\subsection{Le due anime delle lezioni di game design}

\begin{description}
  \item[Design methods (metodologie)] lezioni prevalentemente teoriche, ma con
        una buona componente pratica, in cui si studiano metodologie di design
        e i cosiddetti \emph{framework} dell'industria.
  \item[Documentazione e pitching] il cuore del lavoro quotidiano del
        designer.
\end{description}

\paragraph{La documentazione.}
Il docente si definisce ``il più grande nerd della documentazione'' che si
possa incontrare. È una parte dura, perché a pochi piace documentare, ma il
lavoro del designer \emph{è principalmente questo}: scrivere documenti e, da
quei documenti, passare alla pratica e all'implementazione. Chi non sa
documentare bene parte male. Quando non si è presenti, il documento è il
proprio \textbf{biglietto da visita}: qualcun altro lo leggerà, quindi va
scritto nel modo più chiaro e comprensibile possibile.

Un esempio concreto del perché la documentazione conta: CD Projekt Red ha
sviluppato un proprio engine proprietario, ma buona parte della sua
documentazione è stata ``buttata all'acqua''. Chi deve usarlo o lo sa già
usare, o ha competenze di un altro engine da trasferire, oppure è in
difficoltà. Si documenta \textbf{perché altri possano leggere}.

\paragraph{Il pitching.}
Il \emph{pitching} è l'esposizione, la presentazione orale. La parte creativa
del team presenta spesso:
\begin{itemize}
  \item le proprie idee ai colleghi;
  \item il gioco a convention ed eventi;
  \item il pitch del gioco agli investitori, per ottenere i fondi necessari a
        realizzarlo.
\end{itemize}
Saper parlare e comunicare chiaramente le idee è fondamentale. In sintesi:
\emph{il game designer parla}, e bisogna farsene una ragione.

\subsection{Altre componenti del corso}

\begin{description}
  \item[Gioco critico e playtest] la classe viene divisa in due gruppi: uno
        gioca, l'altro osserva. Chi osserva non applaude ``come le
        scimmiette'': annota, segna e analizza il \emph{playtest}, esattamente
        come si fa nell'industria. Il playtest fa parte delle mansioni del
        designer, ma esiste anche come specializzazione autonoma.
  \item[Pipeline] termine che indica tutto ciò che riguarda i processi
        formalizzati dell'industria; di solito si intende la \emph{pipeline di
        produzione}: come si divide il team, qual è la sua gerarchia, come il
        \emph{producer} gestisce le deadline. In breve: \emph{come} e
        \emph{quando} si fanno le cose. Al primo anno è meno rilevante (si
        dovrà ``nuotare'' un po' da soli); dal secondo anno si vedranno le
        metodologie effettivamente applicate ai progetti digitali. Si
        scoprirà quanto sia difficile lavorare a un progetto digitale, e
        ancor più fisico, senza un piano di sviluppo comprovato, come quelli
        usati per qualsiasi software in qualsiasi parte del mondo.
  \item[Engine] il motore di riferimento è \textbf{Unreal Engine}
        (sezione~\ref{sec:engine}).
  \item[Programmazione] principalmente tramite \emph{Blueprint} in Unreal,
        più eventuali sistemi di versionamento; il corso parte da zero.
  \item[Progetti pratici] probabilmente la cosa più importante. Al primo
        anno i progetti di gruppo restano circoscritti al corso di design,
        perché mancano ancora alcune nozioni (come comunicare, affrontare
        problemi, dare e ricevere feedback). In seguito si allargheranno a
        progetti digitali, molto più difficili di quelli fisici, con membri
        di altri team: un rapporto non unilaterale, ma ``covalente''.
\end{description}

% ---------------------------------------------------------------------
\section{Il metodo: l'inquiry-based learning}
\label{sec:inquiry}

\subsection{Che cos'è}

Il corso è costruito secondo l'\textbf{inquiry-based learning}
(apprendimento basato sull'indagine), un approccio pedagogico nato intorno
agli anni Sessanta, che si richiama a pensatori come Dewey e Piaget. In
generale, le discipline da cui il design attinge sono quasi sempre la
sociologia, la psicologia e, talvolta, l'economia.

Il funzionamento è un ciclo:
\begin{enumerate}
  \item il docente fornisce metodologie e framework dell'industria durante
        la lezione;
  \item lo studente, fuori dalle ore di lezione, lavora sui \emph{task}:
        mette a terra le nozioni, le maneggia, le assorbe meglio;
  \item da questo lavoro nascono nuovi dubbi;
  \item i dubbi tornano in aula: lo studente chiede, il docente risponde
        (ammesso che dare la risposta sia utile, e ammesso che una risposta
        esista).
\end{enumerate}

\subsection{Non sempre esiste una risposta giusta}

Non tutto ciò che si studia ha ``la risposta migliore possibile in assoluto
nel vuoto''. Esistono metodologie, limiti e framework che aiutano il lavoro,
ma non sono fissi: cambiano da azienda ad azienda, da lavoro a lavoro, da
progetto a progetto, e a volte vengono stravolti. Sta alla
\textbf{sensibilità del designer} capire quale strumento o framework usare in
un dato momento, e questa sensibilità si impara solo con l'esperienza. Nel
lavoro nessuno dirà ``usa questa cosa'' o ``è meglio quest'altra'': sarà il
designer a scegliere.

Anche l'esperienza del docente ha confini: sei-sette anni di lavoro, un modo
personale di approcciare il design e un certo insieme di progetti (ad esempio
nessun MMO). Oltre quei confini si cercano risposte nella ricerca, studiando
come aziende con esperienza specifica (ad esempio Riot) hanno affrontato
problemi simili e hanno condiviso le loro soluzioni, sempre valutando quali
metodi siano adatti alla dimensione del team, al tipo di progetto e al tempo
disponibile.

\subsection{Sensibilità, non memoria}

Non interessa imparare a memoria un framework o la teoria di un autore. Certe
terminologie vanno conosciute, ma per ragioni di comunicazione e chiarezza.
Ciò che conta è sviluppare la \textbf{capacità di farsi domande}, di farle
agli altri, di fare ricerca e di trovare una propria risposta. È esattamente
ciò che si fa nel lavoro: davanti a qualcosa che non si sa fare, l'abilità
del designer sta nel cercare come qualcun altro ha risolto un problema simile.

\subsection{Le lezioni come corsi di aggiornamento}

Le lezioni frontali simulano i \emph{corsi di aggiornamento} aziendali: le
sessioni che un lead di design tiene al suo team di design (o il lead di
level design al suo team, e così via). Alla fine dell'aggiornamento c'è quasi
sempre un lavoro da fare, un obiettivo da portare a terra:
\textbf{l'obiettivo della settimana}. Si lavora di settimana in settimana, e
il perché verrà chiarito in seguito.

Struttura tipica di una lezione:
\begin{enumerate}
  \item breve sessione di domande e risposte sulla lezione precedente;
  \item analisi di eventuali esercizi particolarmente critici;
  \item lezione vera e propria;
  \item discussione di gruppo, se c'è tempo;
  \item assegnazione del task della settimana.
\end{enumerate}
Il docente tende a ``parlare e straparlare'', ad aprire parentesi e a saltare
da un argomento all'altro: l'unico modo per fermarlo è che siano gli studenti
a interromperlo con domande.

\subsection{La gerarchia della competenza}

L'inquiry-based learning può essere rappresentato come una scala di
competenza:
\begin{enumerate}
  \item la sola lezione frontale porta a \textbf{conoscenze incomplete},
        perché manca la parte pratica e di ricerca;
  \item il lavoro a casa (task, ricerca) porta a una competenza \textbf{più
        completa}, non sbagliata in sé, ma ancora parziale e complementare ad
        altre;
  \item la \textbf{condivisione} dei task e dei feedback con gli altri porta
        idealmente alla competenza più completa possibile;
  \item il picco della competenza è la sua \textbf{applicazione nei
        progetti} (fisici o digitali, di gruppo): ``con le cose che avete
        imparato e ricercato, riuscite a fare un gioco?''. Quando la risposta
        diventa positiva, o comunque si inizia a navigare nei processi di
        design per realizzarne uno, quella è la vittoria.
\end{enumerate}
Non si parte subito creando giochi: si parte dall'analisi, dal design e dalle
regole generali di design, prima ``nel vuoto'', poi contestualizzando i
framework, infine unendo i puntini nei progetti.

Rimarranno comunque delle lacune, per tutti, docente compreso, perché il game
design \textbf{non è una scienza esatta}, anzi, a rigore non è affatto una
scienza. Prende in prestito metodologie dalle discipline scientifiche, ma
contiene una componente imprevedibile e incalcolabile: \textbf{la
produzione}. Si possono applicare perfettamente le regole di un framework, e
quel framework può comunque non essere adatto al gioco, al proprio modo di
sviluppare o al tempo di produzione disponibile (alcuni framework richiedono
una ricerca e un'attenzione che il tempo non concede).

C'è poi la componente creativa, la parte ``romantica'' del lavoro: esiste, ma
è \textbf{metà del lavoro}. L'altra metà è ricerca, metodo scientifico al
100\%. Dagli stessi dati $n$ designer diversi, con sensibilità diverse,
possono ottenere $n$ risposte diverse. Non tutte saranno valide: vanno
applicate, messe a terra, e questa scrematura fa parte del lavoro.

\subsection{I livelli dell'apprendimento}

Un secondo modello, a gradini, descrive i livelli di apprendimento (è la
cosiddetta tassonomia SOLO):
\begin{description}
  \item[Pre-strutturale] il docente dà il compito e lo studente lo risolve,
        possibilmente il prima possibile. È esattamente ciò che \emph{non} si
        vuole.
  \item[Uni-strutturale] si affronta una sola disciplina: è il primo passo
        necessario, ed è il motivo per cui nel primo mese ci sono solo lezioni
        di game design.
  \item[Multi-strutturale] entrano in gioco altri dipartimenti: level design,
        game design, programmazione e così via.
  \item[Relazionale] si capisce che alcune cose valgono per più
        dipartimenti: alcuni framework di narrativa derivano dal design o
        possono essere usati per fare design; le regole di documentazione si
        applicano anche ad altri dipartimenti; le regole di level design si
        mescolano con quelle di game design. Dopotutto, per fare un gioco
        servono molti dipartimenti.
  \item[Astratto esteso] il punto finale: muoversi autonomamente e capire
        cosa usare e quando, sulla base della propria esperienza. L'unico modo
        per arrivarci è fare pratica e fare progetti.
\end{description}

% ---------------------------------------------------------------------
\section{Task, feedback e regole del corso}
\label{sec:task-feedback}

\subsection{I task come opportunità}

I task non sono strettamente obbligatori (esiste un numero minimo di
consegne), ma sono \textbf{opportunità}. Quanto e come ci si lavora dipende
dallo studente. Un task fatto in cinque minuti, o in cui si è stati ``furbi'',
non verrà letto né riceverà feedback.

Domande da porsi davanti a ogni task:
\begin{enumerate}
  \item \textbf{Perché} mi è stato assegnato? Qual è l'obiettivo?
  \item Come ci lavoro per ottenere un risultato \textbf{che abbia senso}? Se
        la ricerca e i dati indicano una direzione, è inutile andare in
        un'altra.
\end{enumerate}

\subsection{Quanto lavorare: un calcolo indicativo}

Se il corso è una simulazione di lavoro, vale il carico di un lavoro normale:
\begin{itemize}
  \item un lavoro a tempo pieno comprende da 30 a 45 ore settimanali; si
        considera una media di \textbf{40 ore};
  \item si sottraggono circa 9 ore di lezione: restano $\approx 31$, diciamo
        \textbf{30 ore};
  \item con circa 3 task a settimana: $30 / 3 = $ \textbf{circa 10 ore per
        task}.
\end{itemize}
Il numero può variare, ma rende l'idea: dedicare un'ora a un task e poi fare
altro non fa guadagnare nulla a nessuno. Lo studente ``butta soldi'', il
docente risparmia tempo leggendo il task in un minuto (o non leggendolo).
Dedicando invece attenzione, ricerca, cura nella formattazione e
sperimentazione, e facendolo con cognizione di causa, si impara.

\subsection{Niente voti, solo feedback}

Ad eccezione del corso di programmazione, non ci sono voti (né ``ottimo'', né
10, né 100). Ci sono \textbf{feedback}:
\begin{itemize}
  \item in game design, generalmente sulla documentazione;
  \item in level design, sui progetti pratici.
\end{itemize}
Il docente indica cosa va bene e cosa va modificato, concentrandosi in
particolare su ciò che va cambiato. Il feedback \textbf{non è mai un giudizio
sulla persona}, ma su ciò che è stato consegnato nel concreto.

\subsection{La condivisione dei task}

Tutti i task consegnati vengono \textbf{condivisi} con la classe, e possono
essere letti solo dopo aver ricevuto i feedback del docente. Si cresce così
non solo con i propri feedback, ma anche con quelli dati agli altri. Regola
importante: ciò che il docente dice a uno studente \textbf{vale per tutti}. Se
a una persona viene detto ``questa cosa non va mai più fatta'', da quel
momento vale per l'intera classe.

\subsection{Non copiare}

Copiare i task degli altri è una delle cose peggiori che si possano fare, non
per una questione morale, ma perché annulla il senso del percorso. Si può
dedicare a ogni task il tempo e l'attenzione che si ritengono opportuni, ed è
normale che col tempo ci si specializzi (preferendo, ad esempio, il level
design o la programmazione tecnica) e si investa meno dove si sente meno
necessità. In questi casi l'alternativa corretta è:
\begin{itemize}
  \item non consegnare;
  \item studiare ciò che hanno fatto gli altri;
  \item eventualmente riconsegnare in seguito.
\end{itemize}
Ci saranno momenti in cui si verrà sommersi dalle cose da fare: bisogna
imparare a scegliere.

\subsection{Il ``non leggo''}

Il docente ha piena facoltà di non dare alcun feedback a un task (il
cosiddetto \textbf{``non leggo''}) se:
\begin{itemize}
  \item non si è lavorato con attenzione;
  \item non si rispettano le indicazioni date;
  \item è stata violata una regola estremamente importante.
\end{itemize}
Un task ``non letto'' equivale a un task non consegnato e non entra nel
conteggio. Attenzione: un task \emph{perfetto} e un task che \emph{fallisce}
nel comprendere o applicare un framework sono \textbf{ugualmente validi},
purché ci si sia impegnati e si sia dedicato tempo, pazienza e ``sudore''.

\paragraph{Riconsegna.} Un task non letto, in ritardo o che si vuole
migliorare può essere riconsegnato e verrà conteggiato, a patto che:
\begin{itemize}
  \item i problemi segnalati siano stati risolti;
  \item sia stato fatto ulteriore lavoro, integrando anche tutto ciò che nel
        frattempo è stato detto agli altri colleghi nei loro feedback.
\end{itemize}
Il task va sviluppato con le competenze che si hanno nel momento della
riconsegna, non con quelle del momento della prima consegna.

\subsection{Requisiti per superare l'anno}

\begin{itemize}
  \item consegnare almeno l'\textbf{80\% dei task} assegnati (soglia
        malleabile);
  \item partecipare almeno all'\textbf{80\% delle lezioni}, in presenza o
        da remoto, preferibilmente in presenza; i ritardi vengono valutati
        caso per caso;
  \item comportarsi da persone corrette e rispettose (è stata firmata una
        carta dello studente): non urlare, non alzare le mani, non
        vandalizzare, non compiere atti illegali.
\end{itemize}
Impegnarsi di più porta inoltre a essere selezionati per varie opportunità.

\subsection{Scadenze}

\begin{itemize}
  \item la scadenza di un task è \textbf{sempre prima dell'inizio della
        lezione} successiva indicata;
  \item se la classe, come team, mantiene comportamenti scorretti, il docente
        può \textbf{anticipare le deadline} per tutti.
\end{itemize}

\subsection{Accessibilità e DSA}

Gli studenti con disturbi specifici dell'apprendimento o neurodivergenze
(dislessia, ADHD, autismo, ecc.) sono invitati a segnalarlo. Quando si
affronteranno le regole di formattazione e documentazione (errori
grammaticali, mancanze, incoerenze) se ne terrà conto e si potranno trovare
dei compromessi, pur ricordando che nel mondo del lavoro nessuno farà
``sconti''. La documentazione può essere inviata all'amministrazione o alla
segreteria, che la comunicherà a tutti i docenti, ma il consiglio è parlarne
anche direttamente con ogni singolo docente.

% ---------------------------------------------------------------------
\section{Le competenze trasversali}
\label{sec:soft-skill}

\subsection{I giochi si fanno in gruppo}

Si possono fare giochi da soli, ed è un'impresa ammirevole; ma lavorando da
soli \textbf{tutto ciò che si guadagna in libertà si perde in produttività}.

\paragraph{Il caso Stardew Valley.} \emph{Stardew Valley} è stato realizzato
da una sola persona, Eric Barone, in anni di sviluppo e di sacrifici: faceva
un altro lavoro per mantenersi (in un cinema) e di notte programmava, faceva
l'arte e componeva la musica. Il risultato è stato straordinario e lo ha reso
ricco, ma si tratta di un caso estremamente raro. Quanti giochi simili, nati
nelle stesse condizioni, sono falliti? Molti di più.

\paragraph{Il rischio d'impresa.} Lavorare da soli non è sbagliato di per sé
e ha aspetti molto positivi: si ha il privilegio di creare esattamente ciò che
si vuole. Ma ci si fa carico anche di tutto il \textbf{rischio d'impresa}: se
si lavora a un gioco per cinque anni e poi il gioco va male, resta il
vantaggio di averlo pubblicato, ma il dopo è sperare di vivere di quella
release o di finanziare con essa il gioco successivo. Le grandi aziende hanno
gli stessi problemi, ma li distribuiscono su più persone e più responsabilità.

L'obiettivo di un professionista è \textbf{saper lavorare in diverse
condizioni}, e la più comune è quella del dipendente (o del freelance /
libero professionista con partita IVA): si viene assunti per le proprie
capacità da un'azienda, piccola o grande che sia, e si lavora con altre
persone.

\subsection{Essere persone con cui si vuole lavorare}

Se le persone non vogliono lavorare con voi, quello è \textbf{il primo motivo
di licenziamento}. E se non si viene licenziati, il clima nell'azienda
diventa terribile: la produttività crolla, la motivazione crolla. Le
\emph{soft skill}, cioè l'insieme di abilità che rendono una persona
piacevole con cui lavorare, vanno coltivate e diventano sempre più importanti
col passare degli anni (relativamente meno ora, molto di più quando i team
saranno più grandi). Anche chi è un ``lupo solitario'' deve imparare a:
\begin{itemize}
  \item lavorare con altre persone;
  \item comunicare con altre persone;
  \item trovare un punto in comune da cui partire per fare le cose.
\end{itemize}

\subsection{Indipendenza e proattività}

Sono le altre due parole chiave del corso:
\begin{description}
  \item[Indipendenza] spesso si verrà lasciati da soli, non per cattiveria ma
        per sfida: quando si cerca una risposta da soli, ciò che conta non è
        tanto la risposta quanto \textbf{il processo}. Una risposta data dal
        docente non vale niente.
  \item[Proattività] chiedere, domandare, agire da soli senza aspettare
        direttive, perché spesso quelle direttive non ci sono, a maggior
        ragione nei team piccoli.
\end{description}
Tutto questo si concretizza nel parlare, nel presentare e nel fare
\emph{post-mortem} (l'analisi a posteriori di un progetto concluso, che verrà
approfondita più avanti), e in generale nell'attenzione alla crescita
personale.

\subsection{Fare domande}

Il corso insiste sulle domande: alzare la mano, intervenire anche ``a
bomba'', fermare il docente ogni volta che qualcosa non è chiaro. Se una
persona non conosce qualcosa, molto probabilmente c'è almeno un'altra persona
che non la conosce. Il docente può dare per scontata la conoscenza di termini
che non lo sono, ed è del tutto normale, soprattutto all'inizio.

% ---------------------------------------------------------------------
\section{Strumenti: engine, Unreal e materiali}
\label{sec:engine}

\subsection{Unreal Engine e gli altri motori}

Il motore di riferimento del corso, usato in level design, è \textbf{Unreal
Engine}: ``un gran mattone'', ma che dà a un designer molte possibilità. Il
vantaggio principale è che permette a un designer non pratico di
programmazione di \textbf{prototipare tramite i Blueprint} (il sistema di
scripting visuale di Unreal). Questo non impedisce di passare a Unity o ad
altri engine, e viceversa: il docente stesso ha studiato e lavorato a lungo
con Unreal, ma la sua ultima esperienza professionale (come level designer) è
stata su Unity, un engine che non conosceva, a dimostrazione che il salto da
un engine all'altro è possibile.

Panorama dei motori citati:
\begin{itemize}
  \item \textbf{Unreal Engine} e \textbf{Unity}: i due più diffusi, con
        modelli di sorgente ``più o meno'' aperto;
  \item \textbf{Godot}: engine open source;
  \item \textbf{CryEngine}: proprietario di Crytek;
  \item \textbf{Snowdrop}: engine di Ubisoft;
  \item Rockstar usa un proprio engine proprietario;
  \item alcune aziende usano versioni \textbf{personalizzate} di engine
        esistenti: ad esempio Milestone, nei giochi di corse, usa un Unreal
        customizzato;
  \item CD Projekt Red ha un engine proprietario (con i problemi di
        documentazione già citati).
\end{itemize}
In sintesi: ogni azienda può avere un engine proprietario oppure una versione
modificata di un engine commerciale, ma si parte sempre da una base.

\subsection{Programmazione: cosa ci si aspetta da un designer}

Non è richiesto alcun bagaglio pregresso di programmazione: il corso parte da
zero. Gli studenti sono stati selezionati (il corso è a numero chiuso) per
competenze analitiche e critiche sopra la media. Anche nel mercato del lavoro,
per una posizione di game designer generalmente \textbf{non ci si aspetta
competenza di scripting} (C++, Java, C\#, ecc.), a meno di voler fare il
\emph{technical designer}, e anche lì varia da azienda ad azienda.

\subsection{Versione e preparazione}

Indicazioni pratiche emerse in aula:
\begin{itemize}
  \item scaricare Unreal Engine; la versione consigliata per il corso è
        stata indicata nella \textbf{5.6} (con conferma successiva). La
        stabilità dipende da cosa si deve fare, da quali tool e plugin si
        usano: i problemi nascono spesso dai plugin. Nel passaggio da Unreal 4
        a Unreal 5, ad esempio, il sistema \emph{Landscape} era inizialmente
        inutilizzabile, mentre in Unreal 4 funzionava molto bene;
  \item è consigliato installare tramite \textbf{Epic Games Launcher}, più
        comodo: permette anche di riscattare giochi gratuiti periodici e di
        accedere ai pacchetti di \textbf{Fab} (il marketplace di asset
        acquisito da Epic), che ogni mese offre alcuni asset gratuiti utili
        per esercitarsi nel level design (ad esempio per fare \emph{speed
        level design});
  \item molti pacchetti di UI o di sistemi di interazione sono utili per
        prototipare al volo, ma spesso inutilmente complicati;
  \item un plugin di \textbf{materiali di prototipazione a griglia} (adattivi,
        con metriche, allineabili alla griglia, con materiali già pronti e
        colori modificabili, dal costo di circa 5 euro) è indicato come
        ``salvavita'' per il level designer; se necessario, i pacchetti
        indispensabili vengono forniti;
  \item portare \textbf{sempre il portatile} alle lezioni di level design
        (da novembre, a partire dall'interfaccia dell'engine); alle lezioni
        di game design basta che lo porti almeno metà della classe, per i
        lavori a coppie. Serve comunque qualcosa per prendere appunti
        (``tavolette d'argilla e cannucce, se volete fare i babilonesi'',
        oppure carta e penna).
\end{itemize}
Chi ha un portatile poco potente dovrebbe comunque provare a installare
l'engine e aprire un template per vedere come gira: all'inizio non si faranno
cose pesanti. Nelle prime lezioni di level design è importante mettere subito
le mani sull'interfaccia: ascoltare dove si trova uno strumento (ad esempio lo
\emph{snapping}) senza provarlo porta a dimenticarlo.

% ---------------------------------------------------------------------
\section{Letture e risorse di riferimento}
\label{sec:risorse}

\subsection{I libri}

\begin{description}
  \item[\emph{The Art of Game Design: A Book of Lenses} (Jesse Schell)] il
        testo \textbf{davvero importante}. Alcune lezioni finali dell'anno si
        baseranno sulle sue terminologie. È utile perché apre la mente sulla
        possibilità di applicare metodologie diverse a giochi diversi, ma
        soprattutto per la \textbf{terminologia} comune. Esiste anche in
        audiolibro (circa 25 ore).
  \item[\emph{Challenges for Game Designers} (Brenda Romero e Ian Schreiber)]
        libro \emph{entry level} con moltissimi esercizi, molto caro al
        direttore del corso.
  \item[\emph{21st Century Game Design} (Chris Bateman)] utile come
        retrospettiva: è un libro visionario che ha previsto metodi e
        soprattutto \emph{archetipizzazioni} dei giocatori tuttora valide nel
        design e nel marketing. Servirà per una delle lezioni più importanti
        dell'anno. È utile leggerlo per vedere quanto aveva previsto, ma è
        utile ``fino a un certo punto''.
  \item[\emph{An Architectural Approach to Level Design} (Christopher
        Totten)] non strettamente necessario, ma di lettura semplice. Ha due
        anime: parte dall'architettura (babilonese, greca, ecc.) e la trasla
        su esempi di level design (come \emph{Doom} e \emph{Team Fortress}).
        Un tempo ne esisteva un riassunto su Gamasutra (oggi
        \emph{Game Developer}).
  \item[\emph{A Theory of Fun for Game Design} (Raph Koster)] utile fino a un
        certo punto: contiene diverse idee superate, che lo stesso Koster ha
        affrontato in vari talk reperibili nel GDC Vault. Verrà ripreso nelle
        prime lezioni, poi messo da parte.
\end{description}

\subsection{Canali e conferenze}

\begin{description}
  \item[GDC (Game Developers Conference)] il suo canale raccoglie talk e
        panel estremamente importanti: è ``uno strumento misterioso che ci
        servirà più tardi''. \textbf{Attenzione}: per molti talk mancano ancora
        le competenze per interpretarli, e alcuni potrebbero fuorviare o
        mettere in testa idee strane. La GDC si concentra spesso su
        produttività e mobile.
  \item[Il \emph{Dragon Speech} di Chris Crawford] il talk del fondatore della
        GDC: appartiene a un passato lontano, ma è molto motivante. Da vedere
        subito.
  \item[Game Maker's Toolkit] canale particolarmente utile per nozioni
        pratiche, con video brevi e alcune interviste.
  \item[Nordic Game] altra conferenza importante, con talk e panel di esperti
        del settore.
\end{description}

\subsection{Associazioni e premi: come leggerli}

Le associazioni di settore (ad esempio IIDEA in Italia, o la IGDA) non sono
il male, anzi: sono un canale importante per affacciarsi all'estero o per
accedere a fondi. Va però ricordato che \textbf{tutelano i propri associati}:
non esiste solo il loro palinsesto, e bisogna valutare quanto spazio occupano
nel panorama della visibilità.

Lo stesso vale per i \textbf{premi} (Game Awards, ``Best Mobile Game'',
``Games for Change'', ecc.): al designer \emph{non interessano}. Che un gioco
abbia vinto ``più di 90 premi'' non dice nulla di utile. Servono a dare
visibilità e lustro ai giochi ed è giusto che esistano, perché ampliano il
panorama e il numero di giocatori. Ma ciò che interessa al designer sono i
\textbf{dati}:
\begin{itemize}
  \item come sono state implementate certe cose;
  \item come è stato fatto il marketing e dove il gioco è stato presentato;
  \item non quante wishlist ha fatto, ma \emph{come} le ha fatte;
  \item in che periodo storico è uscito e quali erano i competitor in quel
        momento.
\end{itemize}

% ---------------------------------------------------------------------
\section{Che cosa fa un game designer}
\label{sec:cosa-fa}

Per arrivare a una prima definizione, la classe ha condotto un brainstorming:
ognuno doveva indicare una cosa che, secondo lui, fa il game designer (senza
ripetere quelle già dette). Ne è emersa la lista seguente, commentata punto
per punto.

\subsection{Progettare meccaniche e struttura del gioco}

Il ruolo più generico è quello del \textbf{core mechanic designer}, detto
anche \emph{gameplay designer} e, un tempo, \emph{black box designer}: si
occupa delle \textbf{meccaniche di gioco}. Si va dal piccolo al grande (in
funzione anche della dimensione dell'azienda):
\begin{itemize}
  \item nel \textbf{piccolo} ci sono le singole meccaniche;
  \item nel \textbf{grande} ci sono i \textbf{sistemi}. Nei giochi
        \emph{sistemici}, con molti sistemi interconnessi (ad esempio
        \emph{GTA}), esistono specializzazioni dedicate a bilanciare e
        supervisionare i sistemi: è il \emph{system design}.
\end{itemize}

\paragraph{Quando un gioco è ``sistemico''?}
Non dipende solo da come le meccaniche interagiscono con il giocatore (le
\emph{dinamiche}, secondo i framework che passano dalle meccaniche alle
dinamiche, che verranno studiati più avanti), ma anche da \emph{quante}
meccaniche ci sono e da quanto interagiscono tra loro. Non diremmo che
\emph{Super Mario} sia un gioco sistemico, eppure contiene dei sistemi: la
linea è sottile. Quando i sistemi aumentano molto, sono complessi e si
influenzano a vicenda, il gioco diventa sistemico. Esempio: in \emph{The
Legend of Zelda: Breath of the Wild} il vento modifica quanto il personaggio
sale o scende; il fuoco può incendiare l'erba; l'erba prende fuoco se colpita
da una freccia infuocata; e così via. Nel mercato moderno i giochi sistemici
sono spesso, ma non sempre, legati a mappe \emph{open world}.

\subsection{Le 3C: camera, character, control}

La base del gioco, da cui generalmente si inizia a lavorare (anche nei
prototipi), è la ``sacra trinità videoludica'' delle \textbf{3C}:
\begin{description}
  \item[Camera] che tipo di inquadratura ha il gioco: prima persona, terza
        persona, isometrica, ecc.
  \item[Character] quanto è grande il personaggio, che tipo di personaggio è
        e, di conseguenza, quali azioni può compiere.
  \item[Control] i controlli: che cosa il personaggio può fare, quali sono i
        \emph{verbi} di gioco.
\end{description}
Le 3C sono definite dal designer, ed esiste anche una specializzazione per
questo lavoro. Esistono varianti: quella del direttore del corso è
\textbf{4C + J}, cioè \emph{camera, character, control, conflict} (conflitto)
e \emph{juice} (il ``succo'', il feedback che rende piacevole l'interazione).
Verranno approfondite più avanti.

\subsection{Il personaggio}

Il designer definisce il personaggio dal punto di vista delle meccaniche
(dimensioni, azioni, capacità), all'interno delle 3C. L'aspetto estetico è
invece competenza di altri dipartimenti.

\subsection{Documentare le informazioni per gli altri dipartimenti}

Assolutamente vero: la documentazione è una delle prime attività del
designer. Più l'azienda è grande, più questo lavoro diventa prominente; più è
piccola, meno è rilevante. In \textbf{pre-produzione}, però, è sempre
fondamentale.

\subsection{Mediatore e leader ``al servizio''}

Il designer è un \textbf{mediatore} nei conflitti interni al team, una sorta
di collante che si relaziona con dipartimenti diversi (non sempre gli stessi,
a seconda della fase di produzione). A volte fa anche da leader, a seconda di
come sono gestite la pipeline e la gerarchia.

Questa leadership è difficile da definire perché \textbf{non è verticale}: è
una leadership \emph{al servizio} (\emph{servant leadership}). Il designer
``è asservito al programmatore'': il suo compito non è arrivare con un'idea e
scaricarla addosso agli altri, né portare la soluzione e lasciare il problema
a qualcun altro, ma \textbf{seguire il processo a basso livello}, sporcandosi
le mani.

\paragraph{I feature team.} Nei team medio-piccoli o medio-grandi possono
esistere sotto-team (\emph{feature team}) dedicati a un obiettivo specifico:
ad esempio un designer incaricato del character o della camera, affiancato da
uno o due programmatori e da un artista. In questi casi il designer
supervisiona il lavoro e gestisce la produzione del sotto-team.

\subsection{Gestire l'avanzamento del prodotto}

Esiste una specializzazione che si occupa solo di questo, il
\textbf{producer}, a un livello gerarchico completamente diverso (e che
potrebbe essere una strada percorribile al termine del percorso di studi).
All'interno di un feature team anche il designer può farlo, ma non è detto.
Può certamente \emph{dare la direzione}: nella sua documentazione può
indicare che per arrivare a uno stato giocabile del prototipo serve un certo
lavoro di arte e di programmazione, per lo stato successivo un altro, e così
via. In ambito produttivo esistono metodologie strutturate come
\textbf{Agile}, che verranno viste in seguito.

\subsection{Presentare le idee}

Il pitching (sezione~\ref{sec:gerarchia-corso}) è richiesto al designer sia
in pre-produzione sia in produzione, in misura che dipende dall'azienda. Ad
esempio Ubisoft lascia ai designer, e in realtà a qualsiasi membro, la
possibilità di \emph{pitchare} un'idea prima di una nuova produzione, perché
possa iniziare a prendere vita.

\subsection{Bozze e prototipi}

Le bozze vanno bene, ma devono essere \textbf{chiare}. C'è un'eterna lotta tra
documentazione e gioco giocabile: se si fa solo documentazione il gioco non
avanza; se si fa solo il gioco manca la documentazione e si perde la visione.
Sono due cose \textbf{complementari}.

Il modo per comunicare un'idea in forma grezza ma giocabile è il
\textbf{prototipo}. Un prototipo può essere qualsiasi cosa: di carta, in
Unreal con Blueprint ``brutti''. \emph{Non deve essere bello, deve essere
giocabile}.

\paragraph{Esempio: il prototipo di carta di Mario + Rabbids.} Durante lo
sviluppo di \emph{Mario + Rabbids}, i programmatori di Ubisoft erano impegnati
su altri progetti (come \emph{Ghost Recon}). I designer realizzarono allora
un \textbf{gioco da tavolo} che riproduceva le meccaniche del gioco, quelle
di uno strategico a turni sul modello di \emph{XCOM}. E funzionava.
L'obiettivo era \textbf{validare un'idea}, validare un gameplay.

Altri modi per rendere chiara un'idea sono i \emph{mock-up} nella
documentazione e le immagini modificate o i \emph{mash-up} di qualsiasi
natura. Lato documentazione l'importante è che sia chiaro; lato prototipo,
che sia giocabile.

\subsection{Fare ricerca per l'idea, e poi smontarla}

Il designer ha una componente creativa importante e fa ricerca per arrivare a
un'idea. Ma ciò che deve poi fare necessariamente è \textbf{smontarla}. Il
suo lavoro è scientifico, esattamente come il \textbf{metodo scientifico}, il
cui obiettivo è \emph{confutare un'ipotesi}: si arriva a un'ipotesi tramite
sperimentazione e poi si cerca $n$ volte di confutarla; se resiste, la si
assume come valida \emph{in un certo contesto} (le regole della fisica
classica, ad esempio, non valgono nella meccanica quantistica).

Allo stesso modo il designer arriva a un'idea (tramite ricerca o chiedendo
agli altri: le idee le hanno tutti, non sono un privilegio del designer),
cerca di smontarla e poi cerca di \textbf{metterla a terra}. Se l'idea non è
applicabile, non è implementabile, finché non si trova un modo per
realizzarla va scartata. Può essere l'idea più bella dell'universo, ma
\textbf{se non può essere messa a terra, non esiste}.

\paragraph{L'originalità.} Serve un'idea originale? Dipende dall'obiettivo e
dalla visione del progetto:
\begin{itemize}
  \item se l'obiettivo è vendere molto, un'idea originale non serve
        necessariamente;
  \item se si vuole creare un nuovo filone, serve un'idea, se non originale,
        almeno interessante: può essere un discorso di marketing o una
        meccanica innovativa, con il rischio di non vendere.
\end{itemize}
Il ruolo del \emph{director} e il concetto di \emph{visione} verranno
approfonditi in seguito.

\subsection{Analizzare i competitor}

Arrivati a un'idea (``vogliamo un certo tipo di combattimento''), si va a
vedere se qualcuno l'ha già fatta e \emph{come}. Le meccaniche non sono
privilegio di nessuno, salvo quando vengono brevettate, come il
\textbf{Nemesis System} di \emph{La Terra di Mezzo: L'Ombra di Mordor}.

\paragraph{Esempio: il combattimento da ridimensionare.} In un progetto
(presso Stormind Games) il team si trovò con un sistema di combattimento
troppo costoso: mancavano le risorse e gli animatori per realizzarlo. Poco
output lato gameplay significava non poter fare qualcosa di complesso e
profondo come in \emph{Elden Ring}; al contrario, il team aveva molti
artisti e poteva puntare su qualcosa di più artistico. Si iniziò così a
orbitare intorno al corpo a corpo di \emph{The Last of Us} (parte I e II).
Sul canale della GDC esistono diversi talk che raccontano cosa è stato fatto
e come, con successi, insuccessi e costi (ad esempio del \emph{motion
capture}): da quei talk si ricavarono moltissime informazioni utili.

I competitor possono essere:
\begin{itemize}
  \item \textbf{diretti}: un gioco molto simile al proprio;
  \item \textbf{di mercato}: un gioco che domina il mercato in cui ci si vuole
        affacciare;
  \item \textbf{indiretti}: un gioco con un sistema simile al proprio, ma
        tutto il resto diverso.
\end{itemize}
Regola d'oro: \textbf{là fuori, nascosto nell'ombra, c'è sicuramente un gioco
molto simile al vostro}. Va solo trovato. E se per qualche ragione non ha
venduto, bisogna capire perché.

\subsection{Lavorare in engine e fare tuning}

Il level designer lavora in engine, così come il technical designer; ma
qualsiasi game designer, a seconda di ciò che deve fare, può finire in
engine, soprattutto per il \textbf{bilanciamento} e il \textbf{tuning}.

Il \emph{tuning} è l'azione di bilanciare con precisione i dati. Esempio: in
\emph{Super Mario} si modifica quanto in alto salta il personaggio, con che
velocità, in quanto tempo raggiunge l'apice, e così via. Questo lavoro rende
il gioco:
\begin{itemize}
  \item ad alto livello, dotato della \textbf{progressione} giusta;
  \item a basso livello, \textbf{piacevole da giocare}.
\end{itemize}
Il \emph{fine tuning} è un lavoro di dettaglio, paziente e minuzioso.

\subsection{Bilanciare}

Il bilanciamento, emerso tardi nel brainstorming, è una delle attività
centrali del designer e ha anche una specializzazione dedicata
(l'\emph{economy designer}, sezione~\ref{sec:tipi-designer}).

\subsection{Analizzare il mercato}

In pre-produzione, di solito, sì: capire le tendenze, i generi che cambiano
nel tempo, la compatibilità del gioco con il momento storico, in vista ad
esempio di un possibile publisher. Va capito quanto, come e con quanto tempo
farlo, perché si è sempre limitati da molte cose.

\subsection{Sintesi}

\begin{quote}
Il designer parte dall'idea ma non si ferma lì: la confuta, lavora sui
dettagli, bilancia, fa \emph{fine tuning}; lavora con il team e fa da
\emph{manager} per diversi dipartimenti; e naviga costantemente tra i limiti.
\end{quote}

% ---------------------------------------------------------------------
\section{I limiti del design}
\label{sec:limiti}

I limiti ci sono sempre, e nel mondo del design sono spesso numerosi. I più
importanti entro cui navigare sono:
\begin{description}
  \item[Tempo] ad alto livello, l'idea deve essere realizzabile nel tempo di
        produzione disponibile; a basso livello, il proprio lavoro
        settimanale deve essere fattibile nel tempo che si ha. Non ci si può
        dividere tra mille feature team se realisticamente se ne può seguire
        uno solo.
  \item[Tecnologia] si può o non si può fare qualcosa? Ci sono i server per un
        multiplayer con un certo numero di giocatori? Ci si può permettere,
        ad esempio, una fauna che reagisce alla fisica come quella di
        \emph{GTA VI}, che richiede il meglio dei \emph{technical artist}?
  \item[Background e competenze] che cosa sa fare il team? Realisticamente,
        chi ha lavorato dieci anni su giochi horror viene assunto per un
        horror, perché ha le competenze per risolvere problemi già visti.
  \item[Persone e motivazione] il team ha la motivazione e le competenze per
        fare un certo tipo di gioco, e per lavorare insieme?
  \item[Denaro] tutto ruota intorno ai soldi. Che si lavori da soli o in mille,
        le persone vanno pagate e il gioco deve avere un ritorno economico.
        Bisogna decidere dove tracciare la linea: raggiungere il
        \textbf{break-even} (guadagnare esattamente quanto speso) o puntare a
        grandi guadagni. Con un \textbf{publisher}, inoltre, una parte
        consistente dei ricavi va a lui (ripartizioni dell'ordine di 80/20 o
        70/30, che dipendono dal contratto firmato).
\end{description}

\subsection{Non siete ``gli uomini delle idee''}

Per qualcuno può essere un problema scoprirlo, ma il designer non è
``l'uomo (o la donna) delle idee'': accanto alla parte creativa c'è una parte
tecnica, e si è fortemente limitati da molti fattori. Questi limiti, però,
sono un aiuto:
\begin{itemize}
  \item riducono il ventaglio di possibilità potenzialmente distruttive per il
        team (cose che non si riescono a fare, che non piacciono o per cui
        mancano le risorse);
  \item aiutano a mirare bene, trovando soluzioni creative in uno spazio più
        piccolo e più attuabile.
\end{itemize}

\begin{quote}
\centering
\textbf{Le idee sono nulla, l'implementazione è tutto.}
\end{quote}

% ---------------------------------------------------------------------
\section{Ruoli e specializzazioni}
\label{sec:tipi-designer}

\subsection{La scala gerarchica}

\begin{description}
  \item[Creative Director (CD)] mantiene la \textbf{visione} e dà l'idea, in
        modo che tutto segua una certa direzione. In teoria può diventarlo
        chiunque, ma con un'esperienza notevole, nell'ordine dei decenni: si
        assume che in quel tempo abbia visto come lavorano molti dipartimenti.
        Naturalmente un CD con un background artistico ha meno competenze di
        design, e viceversa.
  \item[Director di disciplina] uno per dipartimento: \emph{game director},
        \emph{level director}, \emph{art director}, ecc.
  \item[Lead] il capo del dipartimento, che gestisce i gruppi di persone. Chi
        lavora in un dipartimento risponde al proprio lead (ad esempio un game
        designer risponde al \emph{lead designer}, non al game director).
  \item[Senior / Principal] i profili più esperti (a seconda dell'azienda).
  \item[Junior] il profilo di inizio carriera.
  \item[Intern] lo stagista.
\end{description}
Lo schema si replica uno a uno per ogni dipartimento: arte, suono, animazione
(a volte arte e animazione insieme, o \emph{rigging} e animazione insieme),
programmazione, UX/UI e così via. A fare la differenza, a parità di ruolo, è
il \textbf{livello di seniority}.

\subsection{UX e UI}

UX e UI stanno a metà tra game design e arte: a seconda dell'azienda il ruolo
si chiama \emph{UX/UI Artist} o \emph{UX/UI Designer}.
\begin{description}
  \item[UX (User Experience)] come una persona interagisce con le cose. Esempio:
        su uno smartphone, un pulsante accende e spegne lo schermo; per
        digitare la password si scorre verso l'alto e compaiono i punti da
        toccare.
  \item[UI (User Interface)] come sono ``vestite'' queste cose: ad esempio
        l'aspetto della barra dei punti vita (HP), realizzata con asset 2D o
        3D.
\end{description}

\subsection{Le specializzazioni del game design}

Più l'azienda è grande, più la specializzazione è puntuale.
\begin{description}
  \item[Gameplay / Core mechanics designer] (sono la stessa cosa) si occupa
        delle meccaniche e delle 3C.
  \item[System designer] si occupa dei sistemi e delle loro interazioni.
  \item[Economy designer] si occupa del bilanciamento; lavora soprattutto con
        fogli di calcolo (Excel, Google Sheets).
  \item[Combat designer] si occupa del combattimento; negli annunci di lavoro
        indica quasi sempre un combattimento \emph{action}, con gameplay
        reattivo (non strategico), anche se non è una regola.
  \item[Multiplayer designer] si occupa delle interazioni multigiocatore.
  \item[Technical designer] fa da ponte tra game design e programmazione.
\end{description}

\subsection{Le specializzazioni del level design}

\begin{description}
  \item[Encounter designer] posiziona i nemici nella mappa in modo che gli
        scontri siano bilanciati e la progressione equilibrata come
        desiderato.
  \item[World / map designer] popola il mondo di gioco.
  \item[Quest designer] popola il mondo con le missioni e pensa alla
        progressione al loro interno, cioè a dove il giocatore deve andare (si
        vedano ad esempio gli annunci di lavoro di CD Projekt Red o di
        Larian).
  \item[Technical level designer] fa da ponte tra level design e
        programmazione.
\end{description}

\subsection{Discipline vicine: suono, UI, produzione}

Il monte ore del corso non permette di uscire come \emph{narrative designer}
(o \emph{sound designer}) completi. Ciò non toglie che si possano coltivare
queste discipline:
\begin{itemize}
  \item nei progetti di gruppo si può occuparsi anche del suono, della UI/UX,
        o provare ad applicare metodologie come Agile per fare
        \emph{producing} e gestire il team;
  \item il tempo dedicato a queste attività va però sottratto al game design,
        e dipende da quante persone ci sono nel team, da quanto lavoro c'è e
        da quanto è grande il progetto;
  \item avere sviluppato e rilasciato un progetto, anche piccolo (anche una
        \emph{game jam}), documentando i propri processi, è già un ottimo
        argomento da presentare in futuro;
  \item la via migliore resta lavorare nell'industria e, da lì, chiedere di
        imparare anche altre discipline.
\end{itemize}
Un gioco si può fare su qualsiasi cosa: sugli scarafaggi, sui piedini delle
sedie, sui pennarelli. Si parte avvantaggiati da ciò che è vicino alla propria
sensibilità: chi ama la musica può pensare a un gioco con meccaniche legate
alla musica, con un gameplay innovativo ma contenuto e realizzabile. Il design
è un cappello che comprende molte discipline.

% ---------------------------------------------------------------------
\section{Approfondimenti emersi in lezione}
\label{sec:approfondimenti}

\subsection{Roguelike e roguelite}

Un esempio di come il marketing conia generi: Supergiant Games, con
\emph{Hades}, rese celebre il termine \textbf{roguelite} al posto di
\emph{roguelike} (pur non essendo il primo roguelite in assoluto). La
differenza:
\begin{itemize}
  \item il \textbf{roguelike} si fonda su una serie di punti fissati in una
        sorta di \emph{manifesto} che definisce che cosa rende un gioco
        roguelike; uno su tutti: a ogni morte si ricomincia la \emph{run} da
        capo;
  \item il \textbf{roguelite} rompe alcuni di questi punti, in particolare
        introducendo una \textbf{progressione esterna alla run}: il giocatore
        può potenziarsi fuori dalla singola partita per vincere più
        facilmente in quelle successive. \emph{Hades} ne è il caso principe.
\end{itemize}
Analogamente, un mobile game sviluppato dal direttore del corso fu
promosso come ``il primo \emph{donation game} al mondo'': un gioco che
convertiva il tempo di gioco in donazioni reali a enti benefici per animali.
Anche qui un neologismo creato lato marketing per definire un genere.

\subsection{Applicazioni fuori dall'industria videoludica}

Le competenze sviluppate non servono solo ai videogiochi:
\begin{itemize}
  \item i \textbf{render architettonici} oggi possono essere realizzati con
        Unreal, e diversi professionisti formati nel game design lavorano in
        questo campo;
  \item la \textbf{gamification}, cioè l'applicazione di regole videoludiche
        a contesti non ludici (mostre interattive, progetti culturali,
        consulenze come quella per applicare questi concetti alle macchine da
        allenamento di un'azienda di fitness), è in Italia uno dei modi più
        redditizi di spendere queste competenze, spesso con finanziamenti
        pubblici.
\end{itemize}

% ---------------------------------------------------------------------
\section{Il primo task}
\label{sec:primo-task}

\begin{enumerate}
  \item Creare un indirizzo e-mail \textbf{professionale} (consigliato: un indirizzo Gmail nella forma
        \texttt{nome.cognome@gmail.com}; l'anno di nascita è accettabile).
        Va bene anche un indirizzo professionale già esistente usato per
        lavoro. Da evitare indirizzi poco seri (``non \texttt{blackfire67}'').
  \item Una persona della classe raccoglie tutti gli indirizzi in un unico
        documento e lo invia via e-mail al docente, che aggiungerà tutti al
        Drive condiviso del corso (dove verranno caricate lezioni, slide e
        materiali, per avere i documenti tutti nello stesso posto).
  \item Installare Discord e iscriversi al server del corso, dove passano le
        comunicazioni ufficiali. È consuetudine, e utile, creare anche un
        gruppo tra studenti.
\end{enumerate}
\textbf{Scadenza}: prima della lezione del venerdì.

% ---------------------------------------------------------------------
\section*{Punti chiave del capitolo}
\addcontentsline{toc}{section}{Punti chiave del capitolo}

\begin{itemize}
  \item Il corso è una \textbf{simulazione di lavoro}: ci si tratta da
        colleghi, con una gerarchia da studio di sviluppo.
  \item Si impara con l'\textbf{inquiry-based learning}: lezione, task,
        dubbi, domande, condivisione, progetti.
  \item Il fallimento affrontato con la mentalità giusta è la principale fonte
        di crescita; i feedback riguardano il lavoro, mai la persona.
  \item Il lavoro del designer è per metà creativo e per metà
        \textbf{ricerca e metodo scientifico}; il game design non è una
        scienza esatta.
  \item \textbf{Documentazione} e \textbf{pitching} sono il cuore del
        mestiere.
  \item Il designer progetta meccaniche e sistemi (partendo dalle
        \textbf{3C}), documenta, media, prototipa, confuta le idee, studia i
        competitor, bilancia e fa tuning.
  \item La leadership del designer è \textbf{al servizio} del team, non
        verticale.
  \item I limiti (tempo, tecnologia, competenze, persone, denaro) restringono
        e indirizzano la creatività.
  \item \textbf{Le idee sono nulla, l'implementazione è tutto.}
\end{itemize}

%    % ...
%
%  (Con hyperref aggiungere anche \renewcommand{\theHchapter}{z.\arabic{chapter}}
%  nella parte zero, per evitare ancore duplicate.)
%  Tutte le etichette di questo capitolo hanno il prefisso "z3:"; i
%  rimandi ai capitoli precedenti usano le etichette "z1:" e "z2:" già
%  definite in GD0_1 e GD0_2, e alcune etichette dei capitoli del corso
%  (GD1, GD3).
% =====================================================================

\chapter[Agile e Scrum nella pratica]{Agile e Scrum nella pratica: l'esperienza di chi ci è passato}
\label{cap:z3-agile-scrum}

Le prime due lezioni zero (capitoli~\ref{cap:z1-feedback}
e~\ref{cap:z2-tempo-idee-leadership}) guardavano al lavoro in team dal punto
di vista di un professionista con quindici anni di industria alle spalle.
Questa terza lezione cambia prospettiva: a parlare non è un veterano, ma
\textbf{due persone che hanno fatto lo stesso percorso di studi} di chi
ascolta e che, meglio di chiunque altro fino a quel momento, sono riuscite a
\emph{internalizzare} Agile e Scrum nei progetti dell'accademia.

La lezione è introdotta e chiusa da Diego Ricchiuti\footnote{Nella
trascrizione compare soltanto il nome di battesimo, ``Diego'', con cui i due
ospiti si rivolgono più volte a lui; il cognome è quello indicato per il
direttore dell'accademia nel capitolo~\ref{cap:z1-feedback}. Dal contesto è
anche il docente che ha già spiegato Scrum e Agile alla classe nelle lezioni
precedenti (``Diego vi avrà già ampiamente indirizzato'', ``come avete fatto
con Diego nelle lezioni scorse'').}, che ha già spiegato alla classe la teoria
di Agile e Scrum e che interviene di tanto in tanto per sottolineare o
completare un passaggio. Gli ospiti sono due:
\begin{itemize}
  \item \textbf{Simone Cicchetti}, ex studente dell'accademia, oggi level
        designer e producer, che parla per circa un'ora e mezza della teoria
        ``vista dalla vita reale'' e degli errori più comuni;
  \item \textbf{Cristian Liberatore}\footnote{Nella trascrizione il nome
        compare sia come ``Cristian'' sia come ``Christian''; qui si usa la
        prima forma, quella con cui viene presentato all'inizio della
        lezione. Anche i nomi dell'accademia sono resi in modi diversi
        (``Hive'', ``a dive'', ``IVE Online'', ``Live Online''): come negli
        altri capitoli si usa semplicemente ``l'accademia''; ``IVE Online''
        sembra indicare il corso a distanza dell'accademia.}, studente del
        terzo anno, game designer e project manager di due progetti nati
        come progetti d'anno, che segue passo per passo il ciclo di Scrum
        con gli esempi del suo gioco.
\end{itemize}
Diego li ringrazia per il tempo dedicato: Simone lavora, quindi sta
sacrificando tempo personale, e Cristian è uno studente con lo stesso carico
di lavoro della classe. Sono stati scelti perché, tra gli studenti avuti
fino a quel momento, sono quelli che hanno davvero \textbf{internalizzato}
Scrum e Agile, ottenendo grandi risultati produttivi nel progetto scolastico
e, nel caso di Simone, anche dopo l'accademia\footnote{La trascrizione dice
``nel caso di Christian post-accademia'', ma è Simone l'ex studente che oggi
lavora, mentre Cristian è ancora iscritto: il riferimento al dopo-accademia
è stato quindi attribuito a Simone.}.

Perché insistere sulla parola \emph{internalizzare}? Perché, dice Diego,
``utilizzare Scrum e Agile sembra molto facile, ma poi in realtà quando li
provi ad applicare veramente ti rendi conto che facili come sembravano quando
facciamo lezione non lo sono''. Se invece li si internalizza davvero, si
ottengono grandi risultati. L'idea della lezione è che possa essere utile
vedere l'esperienza di qualcuno \emph{molto più vicino} agli studenti di
quanto lo sia il docente, che ormai usa questi strumenti da anni. Il formato
annunciato è di circa un'ora e mezza per ciascun ospite, domande comprese.

Due avvertenze di lettura. La prima: i due interventi sono fatti di teoria
richiamata in breve e di moltissimi esempi concreti, spesso con numeri,
errori e ripensamenti; sono riportati con cura perché sono proprio la parte
che la teoria non dà. La seconda: i due relatori usano in modo
intercambiabile alcuni termini (\emph{framework} e metodologia,
\emph{post-mortem} e revisione della sprint, \emph{score}, peso e punteggio
dei task), e il capitolo segue il loro uso, segnalando in nota dove si
discosta dalla terminologia ufficiale di Scrum.

Il percorso del capitolo segue quello della lezione:
\begin{enumerate}
  \item un'apertura organizzativa e la presentazione degli ospiti
        (sezione~\ref{sec:z3-apertura});
  \item \textbf{Simone}: chi è, un sondaggio sulla classe e le prime
        esperienze di Scrum degli studenti (sezione~\ref{sec:z3-simone});
  \item il \textbf{mindset Agile}: perché non è un framework e perché va
        adattato (sezione~\ref{sec:z3-mindset});
  \item le \textbf{sprint} e le riunioni di pianificazione e
        \emph{post-mortem} (sezione~\ref{sec:z3-sprint});
  \item le \textbf{cerimonie} e la regola della \emph{consistency}, con le
        critiche al \emph{daily scrum} e al \emph{planning poker}
        (sezione~\ref{sec:z3-cerimonie});
  \item il \textbf{backlog} e gli obiettivi dei task
        (sezione~\ref{sec:z3-backlog-simone});
  \item la \textbf{velocity}: perché non usare le ore, la sequenza di
        Fibonacci e come la velocity permette di ``prevedere il futuro''
        (sezione~\ref{sec:z3-velocity});
  \item il riepilogo della creazione di una sprint e la conclusione sulla
        \textbf{fiducia} (sezione~\ref{sec:z3-ricapitolo});
  \item le domande a Simone, con la costruzione in diretta di un foglio per
        il calcolo della velocity (sezione~\ref{sec:z3-qa-simone});
  \item \textbf{Cristian}: chi è e come imposta la sua parte
        (sezione~\ref{sec:z3-cristian});
  \item il caso di studio, il gioco \emph{L'Inf}, e il ``trauma'' del cambio
        di team (sezione~\ref{sec:z3-linf});
  \item il \textbf{loop di Scrum} in cinque fasi
        (sezione~\ref{sec:z3-loop}): la stesura del backlog
        (sezione~\ref{sec:z3-backlog}), la selezione dei task
        (sezione~\ref{sec:z3-selezione}), l'implementazione con i suoi
        controlli, i \emph{blocker} e i feedback
        (sezione~\ref{sec:z3-implementazione}), il calcolo della velocity
        (sezione~\ref{sec:z3-calcolo}) e il \emph{tuning} della produzione,
        con il problema del perfezionismo (sezione~\ref{sec:z3-tuning});
  \item una domanda su come comunicare un taglio al team
        (sezione~\ref{sec:z3-qa-taglio});
  \item il \textbf{team management}: il team prima di Scrum, la produzione
        senza stress, il mindset condiviso, il \emph{post-mortem} e la
        \textbf{sincerità} (sezione~\ref{sec:z3-team});
  \item la chiusura della lezione (sezione~\ref{sec:z3-chiusura}).
\end{enumerate}

% ---------------------------------------------------------------------
\section{Apertura e presentazione degli ospiti}
\label{sec:z3-apertura}

\subsection{Una registrazione andata persa}

La lezione si apre con una comunicazione organizzativa\footnote{La
trascrizione comincia a frase già iniziata (``in una modalità non video fin
tanto che ce la chiamate poi la converte alla fine''): l'inizio della
spiegazione non è stato registrato e non è chiaro a quale lezione precedente
si riferisca.}. Il programma usato per registrare una lezione precedente
registra in una modalità che va convertita alla fine; durante la conversione
il computer di Diego, probabilmente per colpa degli aggiornamenti di Windows,
si è bloccato e la conversione è ``andata a pallino''. Poiché si tratta di
una lezione importante, Diego si rende disponibile a \textbf{rifarla
interamente} per registrarla di nuovo oppure, come la classe preferisce, a
fornire un audio, un video o qualunque altro formato sia più comodo, che poi
caricherà. Chiede agli studenti di pensarci e di fargli sapere.

\subsection{Chi sono gli ospiti e perché sono qui}

Poi Diego presenta gli ospiti speciali della giornata: Cristian Liberatore,
studente del terzo anno, e Simone Cicchetti, ex studente dell'accademia (che
in un primo momento ``sparisce'' dalla griglia della videochiamata, finché
Diego non lo ritrova in alto a destra). Le ragioni della scelta, il valore
della parola \emph{internalizzare} e il formato (un'ora e mezza ciascuno, con
domande) sono quelli già riassunti all'inizio del capitolo. Detto questo,
Diego si silenzia e lascia la parola a Simone.

% ---------------------------------------------------------------------
\section{Simone: il punto di vista del producer}
\label{sec:z3-simone}

\subsection{Chi è il relatore}

Simone si presenta mentre condivide lo schermo (``adesso tutto è imploso,
perfetto''). È uno degli ex studenti del terzo anno e oggi lavora come
\textbf{level designer}; collabora con Diego su alcuni progetti legati alla
gestione dei team, al \emph{game design}, al \emph{framework design} e ad
altre ``cose molto fighe'' che alcuni studenti forse conoscono\footnote{Il
passaggio è confuso nella trascrizione (``lavoro con diego alcune alcuni
progetti diciamo sul gestione di pro dei team game design framework
design''): la ricostruzione è approssimativa. Più avanti Simone dice anche di
lavorare per una società australiana.}. Sta inoltre continuando a lavorare al
progetto del suo terzo anno (il \emph{Biggest One}\footnote{Dalla lezione
si capisce che i progetti d'anno dell'accademia si chiamano \emph{Big One}
(primo anno), \emph{Bigger One} (secondo anno) e \emph{Biggest One} (terzo
anno); la trascrizione li scrive in vari modi (``biggest one'', ``bigger
one'', ``BiggerOne'').}), un gioco chiamato \emph{ISO}\footnote{Il nome del
gioco compare sempre come ``ISO''; non è possibile stabilire dalla sola
trascrizione se sia la grafia corretta (potrebbe essere, ad esempio,
\emph{Iso} o una sigla).}, per il quale è level designer, game designer,
social media manager e soprattutto \textbf{producer}: si occupa di gestire il
lavoro delle persone, di far rispettare le cerimonie di Scrum e, ``punto
centrale della questione'', di ``rompere le scatole alla gente''.

\subsection{Una lezione partecipativa}

Il tempo a disposizione è molto, quindi Simone vuole una lezione il più
possibile \textbf{partecipativa e cooperativa}: a differenza del solito, le
domande vanno fatte in qualsiasi momento, scrivendo in chat per attirare la
sua attenzione o alzando la ``manina'' con le emoji, che controlla su un
secondo monitor con la griglia dei partecipanti. Al massimo aspetterà la fine
della slide per non perdere il filo.

L'obiettivo dichiarato non è spiegare la teoria, che la classe ha già fatto
con Diego, ma darne un breve \emph{recap} e poi spiegare \textbf{a che cosa
serve nel pratico e com'è nella vita reale}: ``come spesso accade, la teoria
è una cosa, la vita reale è un'altra''.

\subsection{Un piccolo sondaggio}

Prima di partire, Simone chiede chi abbia ``un amore incommensurabile e
infinito'' per il producing, lo \emph{scrum mastering} e il management in
generale: nessuno, e uno studente, Gavi\footnote{Così nella trascrizione;
potrebbe trattarsi di un soprannome o di un nome storpiato.}, risponde
spammando il ``no''. Poi la domanda opposta: chi lo odia così tanto da
preferire ``affrontare un alligatore in una lotta greco-romana''? Ancora
nessuno. La classe è quindi, in generale, ``tutta intermedia''.

Chi si è già dovuto occupare di queste cose nei progetti d'anno o nei lavori
di gruppo? Rispondono due studenti, Mattia e Lorenzo\footnote{Nella
trascrizione ``da noi Mattia, Podini'' e ``da noi Lorenzoaldasso'': i cognomi
non sono ricostruibili con sicurezza (``Podini'' potrebbe essere il cognome
di Mattia o il nome di un altro studente).}, ognuno per il proprio team.
Simone chiede solo i punti salienti: che cosa è andato bene, che cosa no e
perché.

\paragraph{L'esperienza di Mattia.} Il suo gruppo usa Scrum da pochissimo,
circa una settimana. È lui a gestire la \emph{Scrum board}, con l'aiuto degli
altri due membri, ma il metodo non è ancora applicato a pieno. Mattia è
interessato a Scrum, ma la \textbf{responsabilità di guidare altre persone}
oltre a se stesso gli pesa parecchio, e usare Scrum in modo serrato tutti i
giorni gli fa ``una bella paura'': se qualche meccanismo si rompe,
bisognerebbe correre ai ripari. Che cosa intende per ``rompere un
meccanismo''? Che se una cerimonia non viene rispettata, tornare allo
\emph{status quo} non sarebbe semplice, anche per la sua mancanza di
esperienza nell'organizzare i team. Simone riassume: la preoccupazione è più
sul piano umano che su quello delle ore di lavoro necessarie. Mattia
riconosce che Scrum avrà le sue ``quasi magiche proprietà'', ma non avendolo
ancora visto all'opera nella vita reale non può dirlo con certezza.

\paragraph{L'esperienza di Lorenzo.} Niente di diverso: anche il suo gruppo
ha iniziato da pochissimo e proprio il giorno prima ha fatto la prima
riunione di Scrum, decidendo insieme i pesi dei task. Alla domanda se il
ruolo sia arrivato per voto popolare, Lorenzo risponde che si è offerto lui:
conoscere un \emph{framework} di lavoro, e soprattutto capire come
applicarlo, gli interessa parecchio.

Simone trova molto interessante quanto detto da Mattia sulle cerimonie e sul
doverle rispettare, e promette di tornarci più avanti
(sezione~\ref{sec:z3-cerimonie}).

% ---------------------------------------------------------------------
\section{Il mindset Agile}
\label{sec:z3-mindset}

\subsection{Agile non è un framework}

La prima slide parte da una frase che tutti conoscono: \textbf{Agile non è un
framework, è un mindset} (Simone dice di aver ``percepito la furia negli
occhi di Diego'' quando la classe esitava a rispondere). La frase sembra una
di quelle che servono ``per fare i fighi con i colleghi'', ma ha un
significato preciso: un \emph{framework} presuppone \textbf{regole da
seguire}, mentre Agile \textbf{non ha regole}. È una serie di \textbf{linee
guida da internalizzare}: la parola chiave è che devono diventare
\emph{interne}, al punto da non dover più pensare attivamente a quali siano
le regole, i punti o la mentalità di Agile. Bisogna capire la logica che c'è
sotto.

L'analogia è quella della \textbf{lingua straniera}: all'inizio si ragiona in
italiano, si costruisce la frase in italiano, la si traduce e la si dice; man
mano questi passaggi si annullano e si comincia a ragionare direttamente
nella nuova lingua. Con Agile funziona allo stesso modo.

\subsection{I quattro valori, e perché non vanno presi alla lettera}

Simone rilegge rapidamente i quattro valori del manifesto Agile:
\begin{enumerate}
  \item gli \textbf{individui e le interazioni} sono più importanti dei
        processi e degli strumenti;
  \item un \textbf{prodotto funzionante} è più importante di una
        documentazione esaustiva\footnote{Nella lezione ``documentazione
        comprensibile''; nella formulazione corrente del manifesto Agile il
        termine è \emph{comprehensive documentation}, cioè documentazione
        esaustiva.};
  \item la \textbf{collaborazione con il cliente} è più importante della
        negoziazione del contratto;
  \item \textbf{rispondere al cambiamento} è più importante di seguire un
        piano.
\end{enumerate}
Come è scritto sotto al manifesto stesso, questo non significa che le cose a
destra, quelle ``meno importanti'', non contino o non vadano seguite. Ma
Simone va oltre: \textbf{non significa nemmeno che queste regole debbano
essere seguite davvero}. Non è detto che individui e interazioni debbano
essere per forza più importanti di processi e strumenti. Il punto chiave di
Agile è \emph{sentire} quale sia il procedimento e la mentalità giusti e
\textbf{adattare Agile al proprio team}, perché ogni team e ogni progetto è
diverso.

Se un team e un progetto, per vari motivi, hanno bisogno di pagine e pagine
di documentazione, allora per loro la documentazione è più importante del
prodotto funzionante. Gli esempi sono infiniti; uno è il \textbf{gioco di
carte}: lì il prodotto funzionante è relativamente meno importante della
documentazione, perché le logiche si possono scrivere in qualsiasi momento,
ma senza una documentazione solida a supporto ``cadrà tutto quanto
immediatamente''. Quindi non bisogna seguire ``a menadito'' i punti del
manifesto, ma capirne la logica di fondo, capire \textbf{a che cosa serve} e
capire \textbf{se serve}: non è detto che Agile serva.

\subsection{Altri mindset e i framework}

Esistono molti altri \emph{mindset}. Di recente va molto di moda la
mentalità \textbf{lean}, che Simone definisce ``già il 2.0''\footnote{Così
nella lezione. Storicamente il \emph{lean} nasce nella produzione
industriale (il sistema Toyota) e precede Agile; nello sviluppo software il
\emph{lean software development} si è diffuso negli anni Duemila. Il ``2.0''
va letto come una battuta sul fatto che oggi è di moda, non come una
derivazione da Agile.}. L'importante, aggiunge con Diego, è \textbf{non usare
Waterfall}, ``perché Waterfall nel gaming\dots{} non dirò nient'altro''.

Sulla base di Agile sono state create delle \textbf{metodologie}, che
Simone si corregge chiamando \emph{framework}: e queste sì che hanno delle
regole. L'esempio è Scrum, ma ce ne sono infiniti, come Kanban. Si basano
tutti sul principio Agile dell'\textbf{inspect and adapt} (ispeziona e
adattati, analizza e adattati), che si può riassumere in una parola chiave:
\textbf{iterazioni}. Scrum è completamente basato sulle iterazioni, e anche lo
sviluppo dei videogiochi lo è, ``molto più dello sviluppo software''.

Il resto della sua parte è dedicato ai punti di Scrum, il \emph{framework} da
cui presumibilmente partono tutti gli studenti, e agli \textbf{errori
classici} in cui si incappa spesso. L'ideale sarebbe trasmettere la logica
alla base di tutto, ma Simone avverte che sarà difficile, perché ``sono cose
che si iniziano a capire mettendole in pratica''. Per questo invita di nuovo a
fare domande sulle proprie esperienze.

% ---------------------------------------------------------------------
\section{Le sprint e le riunioni}
\label{sec:z3-sprint}

\subsection{Il libro di Scrum}

Simone mostra il libro in cui viene raccontata la genesi del
\emph{framework} Scrum, scritto da Jeff Sutherland\footnote{Si tratta
verosimilmente di J.~Sutherland, \emph{Scrum: The Art of Doing Twice the Work
in Half the Time} (2014), edito in italiano come \emph{Fare il doppio in metà
tempo}. Il titolo non viene pronunciato: Simone mostra la copertina e più
avanti cita ``Scrum, di Jeff Sutherland''. La lunghezza indicata (``circa 300
pagine'') è approssimativa.}: un ``libricino'' di circa trecento pagine.
Avverte che userà le parole ``metodologia'' e ``framework'' in modo
intercambiabile. Lo consiglia assolutamente a chi è interessato al tema,
perché è scorrevole e contiene \emph{insight} molto interessanti; chi non è
particolarmente interessato può evitarlo del tutto, perché non avrà niente da
offrire che non venga detto in questa lezione.

\subsection{Che cos'è una sprint}

Scrum, a differenza di Agile, è un \emph{framework} a tutti gli effetti: ci
sono regole da seguire in modo preciso e da rispettare. Il modo in cui
abbraccia il principio dell'\emph{inspect and adapt} sono le
\textbf{sprint}: periodi di tempo tipicamente di \textbf{una o due
settimane}. Possono essere anche di quattro, come si studia, ma secondo Simone
non ha senso farle così lunghe a meno di non lavorare a un grande progetto
``con duecento miliardi di persone che ci stanno lavorando
costantemente''\footnote{Nella trascrizione ``a meno che non state facendo un
sesto plà con 200 miliardi di persone'': ``sesto plà'' è incomprensibile e
potrebbe essere una storpiatura di ``tripla A''; il senso, in ogni caso, è
quello di una produzione enorme.}.

Prima di ogni sprint c'è una riunione di \textbf{sprint planning}, dopo ogni
sprint una riunione di \textbf{post-mortem}\footnote{Nella guida ufficiale di
Scrum la riunione di fine sprint è divisa in \emph{sprint review} (che cosa è
stato prodotto) e \emph{sprint retrospective} (come ha lavorato il team);
qui, come nel resto della lezione, entrambe confluiscono nel termine
\emph{post-mortem}, che in senso stretto indica l'analisi a progetto
concluso (sezione~\ref{sec:z1-qa-fasi}).}.

\subsection{La regola d'oro: meno riunioni possibili}

``Le riunioni non piacciono a nessuno'': sono, dice Simone, una ``rottura di
scatole infinita'', e chiunque, soprattutto chi non è designer, non aspetta
altro che poter dire ``questa riunione è inutile, dobbiamo andarcene tutti a
lavorare perché abbiamo cose più importanti da fare''. In chat qualcuno
risponde con il meme di Bernie Sanders a braccia incrociate; Simone avrebbe
voluto mettere quello del cane che urla sulla scrivania perché ``tutta la
riunione poteva essere una mail'', e promette che dalla prossima volta ogni
slide avrà un meme alla fine.

Da qui la sua \textbf{regola d'oro}: \textbf{fare meno riunioni possibile}.
Sottolinea però che è una regola ``per quanto mi riguarda'': magari un team
non vede l'ora di fare riunioni, magari si portano le pizzette a ogni
incontro (``non so se con le pizzette le riunioni siano effettivamente
efficienti, però se funziona, perché no?'').

Una mossa molto semplice per ridurle è \textbf{unire sprint planning e
post-mortem} in un'unica riunione: si fa letteralmente la metà delle
riunioni. Ciò che si fa in quella riunione è:
\begin{itemize}
  \item \textbf{pianificare} la prossima sprint (\emph{sprint planning});
  \item \textbf{analizzare} la sprint precedente (\emph{post-mortem}),
        indicando in particolare i problemi incontrati e le cose rimaste
        bloccate.
\end{itemize}

\subsection{Una riunione deve essere rapida}

La riunione unica deve essere \textbf{rapida}. Durante l'ultimo anno, per il
\emph{Biggest One}, il team di Simone contava circa quindici persone, e alle
riunioni di planning e post-mortem ne partecipavano in media una dozzina.
Quando sono ``diventati bravi'', le riunioni duravano \textbf{venticinque
minuti, con dodici persone che parlavano tutte}. Quello è l'obiettivo, ``anzi,
il vostro obiettivo è fare anche meno''.

Il motivo per cui si sfora è una \textbf{tentazione} costante: quella di
\textbf{risolvere i problemi}. Quando qualcuno descrive un problema, un paio
di persone cominciano a discuterne (``hai provato a fare questo?'', ``no,
dovresti fare quest'altro''), e in un attimo sono passati dieci minuti in cui
hanno parlato solo in due senza che la riunione sia andata avanti di un
passo. In queste riunioni non si fa brainstorming e non si ragiona su come
risolvere i problemi: ``bloccatelo sul nascere, picchiate quelle persone, fate
qualsiasi cosa per impedirglielo''.

Quando c'è un problema, si \textbf{prende un appuntamento per risolverlo}.
Simone mostra come dovrebbe suonare un report:
\begin{quote}
``La settimana scorsa mi sono occupato della logica di movimento del
personaggio, ho fatto un po' di tuning e di testing e ho preparato il block
out del primo livello. L'unico problema è che il blueprint che faceva volare
via gli uccelli non mi ha funzionato bene: visto che tu, Cristian, l'hai
utilizzato in un altro livello, più tardi ci sentiamo su Discord quindici
minuti e mi spieghi come l'hai settato. Questa settimana mi occuperò di
completare il block out e di risolvere il problema del blueprint.''
\end{quote}
Fine del report: \textbf{circa un minuto}, e c'è tutto. Tutti devono fare
così: trovare una persona con cui risolvere il problema, identificare
\emph{come} risolverlo e andare oltre. ``C'è tempo e luogo per ogni cosa'', ma
la riunione di sprint planning non è il momento giusto.

\subsection{La sprint come unità di tempo}

Ultimo punto: le sprint sono \textbf{l'unità di tempo} fondamentale per ogni
cosa. Le \emph{milestone} non si misurano in mesi, giorni o settimane, ma
\textbf{in sprint}. Bisogna prendere l'abitudine di dire ``completeremo la
preproduzione in otto sprint'' e non ``completeremo la preproduzione tra due
mesi''. È importantissimo per il calcolo della velocity
(sezione~\ref{sec:z3-velocity}).

% ---------------------------------------------------------------------
\section{Le cerimonie e la regola della consistency}
\label{sec:z3-cerimonie}

\subsection{Le cerimonie di Scrum secondo Simone}

Le regole di Scrum si chiamano \textbf{cerimonie}. Simone le divide in due
gruppi\footnote{La classificazione è quella proposta da Simone. Nella guida
ufficiale di Scrum gli ``eventi'' sono la sprint, lo \emph{sprint planning},
il \emph{daily scrum}, la \emph{sprint review} e la \emph{sprint
retrospective}; il backlog è un ``artefatto'' e il calcolo della velocity
e il \emph{planning poker} sono pratiche diffuse ma non obbligatorie.}:
\begin{center}
  \begin{tabular}{@{}p{0.42\textwidth}p{0.48\textwidth}@{}}
    \toprule
    \textbf{Cerimonie principali} & \textbf{Altre cerimonie} \\
    \midrule
    planning; post-mortem; calcolo della velocity; creazione del backlog &
    \emph{daily scrum}; \emph{poker time estimation} (\emph{planning poker})\\
    \bottomrule
  \end{tabular}
\end{center}
Le quattro principali sono quelle che ``non potete evitare di fare,
altrimenti non state facendo Scrum'': il che, precisa, non è necessariamente
un problema (ci torna alla fine, sezione~\ref{sec:z3-ricapitolo}). Le altre
due le tiene separate perché a lui non piacciono granché, e prima di spiegare
perché beve un sorso d'acqua ``per lasciarvi la suspense''.

\subsection{Consistency: una cerimonia saltata fa cadere le altre}

In Scrum c'è una parola chiave da ricordare: \textbf{consistency}. Bisogna
evitare a tutti i costi di creare \textbf{precedenti di
\emph{inconsistency}}, cioè di incostanza: con Scrum bisogna rimanere
costanti. \textbf{Se si salta una sola cerimonia si crea un precedente che,
a effetto domino, butta giù tutto il resto.} Questo è il punto già sollevato
da Mattia. Ma, avverte Simone, questo \emph{non} vuol dire che ci si debba
forzare: ed è questo il punto chiave.

\subsection{Perché il daily scrum non funziona più}

Il \emph{daily scrum} era pensato ``per il millennio scorso''. Scrum è nato
molto tempo fa\footnote{Simone non ricorda la data esatta (``fine degli anni
Novanta, forse fine degli anni Ottanta''). Il termine \emph{scrum} applicato
allo sviluppo di prodotti risale a un articolo di Takeuchi e Nonaka del 1986;
il \emph{framework} Scrum fu presentato da Jeff Sutherland e Ken Schwaber nel
1995.}, per team che si incontravano tutti i giorni in ufficio alla stessa
ora, con un orario d'ufficio standard: team di sviluppo software, quindi per
lo più programmatori (anche i più ``artistici'', come gli UX designer, erano
alla fine programmatori), con una mentalità piuttosto simile e, soprattutto,
con lo stesso orario.

Oggi, Covid a parte, questo non è più né fattibile né sensato. Gli studenti
non hanno nemmeno un ufficio: si incontravano a volte in sede quando si
poteva, ma per lo più lavorano online, e \textbf{non tutti lavorano allo
stesso orario}. ``E va bene così'': c'è chi lavora di notte (la salute
potrebbe risentirne, ma se si sentono bene così, va bene), e quelle persone
non scriveranno il proprio \emph{daily scrum} all'orario in cui lo fanno
tutti. Senza arrivare a casi del genere, Simone lavora per una
\textbf{società australiana}, dal lato opposto del globo: quando lui si
sveglia, loro stanno finendo di lavorare, e il suo \emph{daily scrum} cadrebbe
nella loro notte.

In queste condizioni è difficile tenere traccia del \emph{daily scrum} e
rispettarlo, e prima o poi qualcuno lo salterà. Quello è un
\textbf{precedente di incostanza}: qualcuno comincia a dire ``ok, non lo faccio
neanche io'', poi ``ma visto che non faccio più gli scrum, a che servono le
riunioni? Questa volta non vengo'', ``e io neanche'': ``e da lì è il caos''.

La conclusione è netta: \textbf{quando qualcosa non funziona, tagliatelo}. Il
senso di ``individui e interazioni sono più importanti di processi e
strumenti'' è proprio questo: se un processo o uno strumento non funziona
\textbf{non è colpa vostra, è il team che non ne ha bisogno}. Non bisogna
fissarsi su ciò che si vuole fare, ma sentire di che cosa ha bisogno il team e
adattarsi a quel bisogno.

\subsection{Perché il planning poker rallenta tutto}

Il \emph{planning poker} (che Simone chiama \emph{poker time estimation})
non gli piace perché lo trova \textbf{estremamente lungo e tedioso}: rallenta
un processo già lento di suo. Creare il backlog e valutare i task è già
lunghissimo (``ci vuole almeno un'intera sprint di riunioni per fare un
backlog''); aggiungerci il planning poker significa moltiplicare il tempo. Si
informa con Diego su come sia stato spiegato alla classe: con il metodo di
Fibonacci, ``\emph{power-uppato}''\footnote{Nella trascrizione ``è rimasta
nelle lezioni di trial'' e ``è stata power upata'': il senso sembra essere che
il planning poker è stato spiegato nelle lezioni precedenti in una versione
potenziata con la sequenza di Fibonacci. ``Trial'' potrebbe essere una
storpiatura.}, il che farà risparmiare un po' di spiegazioni più avanti.

Simone capisce la necessità di sentire le opinioni di tutti
(``\emph{empowering the team}'', eccetera), ma l'effetto è rendere il processo
ancora più lento. E poiché le riunioni non piacciono a nessuno, se si
quintuplica il tempo necessario ``si parte subito male'', perché il team si
annoia già. Anche questo, quindi, lui personalmente non lo usa.

\subsection{Il punto chiave: crearsi il proprio framework}

La cosa più importante è che \textbf{non bisogna eseguire tutto a menadito}.
I \emph{framework} sono stati pensati per essere generici, ma anche per
essere \textbf{adattati}. La procedura che Simone suggerisce:
\begin{enumerate}
  \item si parte \textbf{rispettando il framework al 100\%};
  \item quando qualcosa non funziona, si capisce \emph{perché} non funziona;
  \item si smette di farlo, si risolve, si fa altro;
  \item si capisce qual è il framework più adatto e \textbf{ci si crea il
        proprio}.
\end{enumerate}
``Questo è il punto chiave di tutto questo.''

% ---------------------------------------------------------------------
\section{Il backlog}
\label{sec:z3-backlog-simone}

\subsection{Una ``rottura di scatole'' inevitabile}

Sulla creazione del backlog Simone spende ``giusto due parole'', perché la
parte è molto tecnica e i consigli sono pochi. Creare il backlog è faticoso,
``non ci sono mezze misure per dirlo'': richiede moltissimo tempo e lavoro ed
è intenso a livello mentale, perché bisogna ragionare tantissimo.

\subsection{Il backlog a scaglioni}

In teoria i backlog si creano all'inizio dello sviluppo: in un mondo perfetto
si conoscerebbero fin da subito tutti i task necessari per completare il
progetto. Nel mondo reale questo non accade. Ciò che accadrà, più
probabilmente, è che si andrà avanti \textbf{per milestone}. È ciò che Simone
ha cominciato a fare con \emph{ISO} e con cui si trova bene:
\begin{itemize}
  \item ha fatto il backlog della \textbf{preproduzione}, con tutto ciò che
        serviva per la preproduzione;
  \item finita la preproduzione, ha fatto quello della \textbf{produzione};
  \item finita la produzione farà quello della fase successiva
        (``espansione, build finale, non lo so''), a seconda di come si
        svilupperanno le cose.
\end{itemize}
Fare i backlog \textbf{a scaglioni} li rende più semplici e il processo meno
lungo; e tanto, comunque, bisognerà aggiungere cose in corso d'opera: ``non si
scappa''.

\subsection{Le epic}

L'organizzazione in \textbf{epic} è molto utile, soprattutto per la stesura:
si comincia ragionando per \emph{epic} e poi si scende ai task. Un esempio di
epic è ``creazione del personaggio''. Che cosa contiene il personaggio?
Movimento, tuning, modello, animazioni, materiale, il design delle feature
che dovrà avere, la creazione di tutte queste feature, e così via. Il
processo diventa più \emph{streamlined}, meno pesante e più semplice da
seguire.

\subsection{Ogni task deve avere un obiettivo chiaro}

Fondamentale: \textbf{i task devono avere un obiettivo chiaro}. Un consiglio
comune è scriverli come \textbf{user story}: ``come utente voglio avere un
personaggio che si senta potente, che io riesca a percepire come potente''.
Aiuta moltissimo nello sviluppo software e parecchio nei videogiochi per
alcune cose, come appunto il personaggio. In generale è buona pratica
\textbf{spiegare l'obiettivo e non la soluzione}: dire che il personaggio deve
muoversi, saltare e fare le capriole, invece di dire che ``quando premo il
tasto W deve subire una forza che lo spinge in avanti''. Però non è
fondamentale.

Ciò che è davvero importante è che il task abbia un \textbf{obiettivo
verificabile}: qualcosa di cui si possa dire ``è completo'' o ``non è
completo''. Altrimenti le persone \textbf{non lo completeranno}, o lo
completeranno in modo diverso da come lo si immaginava. Un esempio di task
senza obiettivo chiaro è \textbf{``completare il livello~5''}:
\begin{itemize}
  \item per Simone vuol dire aver creato il block out con qualche elemento
        estetico per sapere dove vanno le cose e come funzionano con i
        moduli già disponibili, e basta;
  \item per qualcun altro potrebbe voler dire inserire illuminazione,
        \emph{post-processing}, \emph{particles}, tutti gli shader fatti
        bene, ottimizzazione e \emph{level streaming};
  \item per altri ancora potrebbe voler dire arrivare all'illuminazione,
        lasciando \emph{level streaming} e ottimizzazione a una seconda fase.
\end{itemize}
È la cosa più importante da ricordare quando si crea il backlog e si
scrivono i task.

% ---------------------------------------------------------------------
\section{La velocity}
\label{sec:z3-velocity}

``Bellissimo, bellissimo'': il calcolo della velocity è per Simone ``la cosa
più figa e allo stesso tempo più odiata''.

\subsection{Il peso dei task}

Come la classe sa, i task sono valutati in base al loro \textbf{peso} (o
punteggio, o \emph{score}). Il peso dipende da quanto un task è difficile e/o
da quanto tempo richiede: è un po' una zona grigia, e Simone consiglia di
valutarlo come \textbf{combinazione delle due cose}.

\subsection{Mai le ore}

Il consiglio fondamentale, che Simone non capisce perché si continui a
contraddire, è \textbf{non usare le ore per valutare i task}, ``per niente,
neanche alla lontana''. Nel libro di Sutherland c'è un paragrafo intero che
spiega perché. Diego commenta che chi lo consiglia di solito fa ``Agile
Waterfall'' e non ha nemmeno letto i libri su Agile.

Il motivo: \textbf{il cervello umano è completamente incapace di stimare le
ore}, cioè di prevedere il tempo necessario per fare qualcosa: ``non siamo
pensati per questo''. Siamo invece \textbf{molto bravi con le misure
relative}: fissato un task che vale 10, sappiamo dire bene quale vale 5,
quale 7, quale 2. Sappiamo dire che un task è molto più grande di un altro.

\subsection{La sequenza di Fibonacci}

Per aiutarsi ulteriormente, Simone usa la \textbf{sequenza di Fibonacci} (1,
2, 3, 5, 8, 13, 21, 34, \dots), anche questo un consiglio di Sutherland che
``nessuno segue''. Sutherland racconta addirittura di un team che usava le
\textbf{taglie dei cani} (``questo task è un pinscher, questo è un golden
retriever'')\footnote{Nel libro di Sutherland le stime per razze canine sono
citate come esempio di scala relativa alternativa ai numeri; nella
trascrizione ``pincher''.}; a Simone sembra eccessivo, e ``i buoni vecchi
numeri'' bastano. Fibonacci ha due scopi molto utili:
\begin{enumerate}
  \item \textbf{Gli intervalli crescenti assorbono l'errore di stima.} Man
        mano che un task diventa più grande e complesso, cresce anche
        l'errore associato. È facile dire che aprire la mail vale 2 e
        rispondere a tutti sul server Slack vale 3: che cosa può andare
        storto? La mail non si apre e bisogna reinserire la password; è
        esploso qualcosa e ci sono cinquanta messaggi invece dei soliti
        venti. Niente che cambi granché. È molto più difficile, invece,
        valutare la differenza tra creare il \emph{behavior} di movimento del
        personaggio e costruire un intero livello: possono succedere
        infinite cose, soprattutto quando c'è di mezzo la programmazione.
        Il behavior del personaggio è un 20 o un 21? ``È un 21, punto'':
        non è un 34, perché non è così difficile, e non è un 13, perché
        sarebbe troppo poco. Magari la cosa più corretta sarebbe stata 23,
        ma non importa: lo scopo è \textbf{appianare} i risultati.
  \item \textbf{Meno scelta, meno discussione.} I valori sono
        ``infinitamente arbitrari'', ``completamente, al 100\%'', e nessun
        altro al mondo li condividerà; ma Fibonacci li rende un po' più
        condivisibili, perché con meno possibilità di scelta ci sono meno
        occasioni di discutere.
\end{enumerate}
Che i valori siano arbitrari non è un problema: finché \textbf{è una persona
sola a darli}, sono tutti arbitrari allo stesso modo, quindi
\textbf{consistenti}; e se sono consistenti hanno una logica e si possono
usare per fare calcoli.

\subsection{Che cos'è la velocity e perché serve}

Sommando il peso dei task completati in una sprint si ottiene la
\textbf{velocity}. Ed è ``fighissima'' perché permette, in sostanza, di
\textbf{prevedere il futuro}.

Simone mostra il grafico della velocity della preproduzione di \emph{ISO}
durante l'anno scolastico. La consegna era a maggio (``10 maggio, o poco
dopo''). Già a inizio marzo, guardando l'andamento, aveva capito che \textbf{non
sarebbero riusciti a finire in tempo}. Il ragionamento era semplice: mancano
dodici sprint; quanto facciamo ogni sprint? Circa 89; quanto manca nel
backlog? Circa 2000. ``Non ce la facciamo''\footnote{I numeri sono quelli
detti a voce: 12 sprint per circa 89 punti danno poco più di 1000 punti,
circa la metà dei 2000 rimasti nel backlog.}.

La velocity aiuta a prevedere il futuro e quindi a \textbf{prendere
provvedimenti prima}. A marzo Simone ha cambiato i piani: prepariamo meno
cose, smettiamo di perdere tempo su ciò che ci blocca, consegniamo qualcosa di
diverso. In più, la velocity del team è \textbf{aumentata} sempre di più, e
alla fine sono arrivati in fondo alla preproduzione \textbf{senza crunch}:
``zero, completamente zero crunch''. Nessuno ha dovuto lavorare di più;
l'ultima settimana non hanno nemmeno toccato il progetto: hanno fatto la
build, l'hanno spedita e per una settimana non ci hanno lavorato. E sono
riusciti a consegnare anche di più del previsto, perché Simone ha
\textbf{analizzato i trend precedenti}: perché in un certo punto c'era un
picco? Che cosa li bloccava e faceva diminuire la produttività? L'ha
scoperto, l'ha capito e hanno agito.

\subsection{Come calcolarla in pratica}

Il consiglio è di \textbf{trovare un modo per fare il calcolo}. Simone usava
Monday, che non aveva uno strumento costruito apposta, ma ci è riuscito lo
stesso; si può fare tranquillamente con \textbf{Google Sheets}, ``col
sorriso''. Per farlo davvero bene bisognerebbe scrivere un po' di script
dentro il foglio, ``niente di improponibile''. Quando comincia a girare
``funziona da dio'': ``assolutamente fatelo, non posso dirlo abbastanza''. (Un
esempio costruito in diretta è nella sezione~\ref{sec:z3-qa-simone}.)

\subsection{Come assegnare i punteggi}

Per assegnare gli \emph{score} conviene partire dagli
\textbf{estremi}: quali task sono un 34? Quali sono un 1? Da lì si mettono i
valori intermedi. Gli estremi fanno da \textbf{ancore} e facilitano il
processo di valutazione \emph{relativa} descritto sopra.

Lo stesso vale per le \textbf{priorità}, se si usa un sistema di priorità:
regole simili, ma con \textbf{intervalli più piccoli} e senza Fibonacci, che
lì non serve. Simone usa una scala da 1 a 5; ma non la considera così
fondamentale.

\subsection{Che cosa blocca che cosa}

Durante il lavoro sul backlog bisogna capire \textbf{che cosa blocca che
cosa}. Ad esempio: ci sono 29 task, per un valore complessivo di 400, tutti
relativi alla creazione dei livelli. Ma se non c'è il movimento del
personaggio, non ha senso fare i livelli, perché vanno costruiti attorno a
quel movimento: la creazione del personaggio \textbf{blocca} i livelli e
quindi ha priorità più alta. Da qui si comincia anche a capire quale reparto
si occupa di che cosa: probabilmente all'inizio i programmatori faranno i
\emph{behavior} fondamentali (come il personaggio), mentre i designer
lavoreranno su documentazione e ricerca. Ma dipende dal gioco: alcuni non
costruiscono i livelli attorno al personaggio, bensì il personaggio attorno ai
livelli, quindi anche l'esempio ``è opinabile''.

% ---------------------------------------------------------------------
\section{Ricapitolando: come nasce una sprint}
\label{sec:z3-ricapitolo}

Prima di passare alle domande, perché ``quello che avevo da dirvi di pratico è
quasi finito'', Simone ricorda che il grosso non sono le regole, ma
\textbf{capire la logica alla base e integrarla nel proprio team e nel
proprio progetto}. Poi riassume la creazione di una sprint in quattro passi.

\begin{enumerate}
  \item \textbf{Post-mortem della sprint precedente.} Si vede che cosa ha
        fatto ciascuno, si capisce che cosa è bloccato e si
        \emph{programma un incontro} per risolvere i problemi.
  \item \textbf{Pianificazione.} Dal backlog si guarda che cosa blocca che
        cosa, quali sono le priorità e che cosa si può fare in questa sprint,
        e si prendono i task.
  \item \textbf{Analizzare i problemi e agire.} Si guarda la velocity: è in
        aumento? In calo? Si è in un picco? L'andamento a dente di sega,
        ``su, giù, su, giù'', è \textbf{assolutamente normale}, ed è ciò che
        ci si deve aspettare; e man mano che il progetto va avanti la
        velocity \textbf{tenderà a diminuire}. Le domande da farsi: come sta
        andando il progetto? Chi sta facendo di più e chi di meno, e perché?
        Qualcuno ha avuto un calo di rendimento? C'è un problema che si
        ripresenta spesso, un collo di bottiglia? E come risolverlo,
        soprattutto \textbf{ascoltando il resto del team}.
  \item \textbf{Rimanere consistenti}, in particolare con Scrum.
\end{enumerate}

\subsection{Ascoltare il team}

Come scrum master o producer, il lavoro più importante è e rimane
\textbf{ascoltare il team}. Tutto il resto non conta se poi, quando il team
dice una cosa, non lo si ascolta. È lo stesso principio dell'adattare il
\emph{framework} al team:
\begin{itemize}
  \item se il team odia il \emph{daily scrum}, non si fa: non serve a niente;
  \item se il team fa fatica a prendere i task dal backlog, non glielo si fa
        fare: si assegnano i task in modo diverso, ``non fategli manco vedere
        il backlog'';
  \item se il team vuole sprint di due settimane invece che di una, si fanno
        di due settimane.
\end{itemize}
Magari le persone si organizzano personalmente in modo diverso: Simone, per
esempio, organizza tutto il suo lavoro, su tutti i fronti, \textbf{su base
settimanale}, e una sprint di due settimane gli sconvolgerebbe i piani. Se il
team lavora per lo più così, si fa così. ``Adattatevi'' non vuol dire solo
adattarsi ai bisogni del software, ma \textbf{soprattutto ai bisogni del
team}.

\subsection{Nessun framework va bene al 100\%}

Esistono infiniti \emph{framework}, e \textbf{nessuno andrà bene al 100\% per
un team}. Bisogna provarli, il più possibile, e adattarli: cambiare qualche
regola, togliere qualche cerimonia. Il \emph{framework} più usato oggi, ad
esempio, è \textbf{Kanban}, che non ha nemmeno le sprint\footnote{Simone dice
prima ``Kanban neanche ce l'ha il backlog'' e poi si corregge: ``scusatemi,
neanche ce l'ha le sprint, ha soltanto un gigantesco backlog''.}: ha soltanto
un gigantesco backlog. ``Ogni team è diverso, ogni team ha bisogni diversi,
quindi adattatevi'' (una frase che Simone calcola di aver ripetuto ``almeno 49
volte'').

\subsection{La nota ``smielata'': la fiducia}

La conclusione: alla fine di tutto, \textbf{Agile non si basa sulle
iterazioni}, né sull'efficienza, né sulla creazione del miglior videogioco o
software del mondo, né sul controllo. L'unica cosa su cui si basa davvero è
\textbf{la fiducia}:
\begin{itemize}
  \item la fiducia che il team ripone in chi lo guida, e che chi lo guida
        ripone nel team: ``queste persone lavoreranno nel miglior modo
        possibile, sempre e comunque'';
  \item la fiducia del team nel dire ``se qualcosa va male, so che lo scrum
        master ha un piano di riserva che ci risolverà il problema, e il mio
        lavoro non sarà sprecato'';
  \item la fiducia di chi guida nel dire ``non posso controllarvi
        costantemente, non posso assicurarmi che facciate sempre tutto il
        lavoro o che seguiate le cerimonie come dovreste, però mi fido'': mi
        fido che lo farete, che se ci saranno problemi ve lo direte, che
        farete sempre del vostro meglio.
\end{itemize}
Bisogna \textbf{costruire fiducia nel team e trasmettere fiducia al team}. È
lo stesso principio del \emph{default to trust} della prima lezione zero
(sezione~\ref{sec:z1-default-trust}).

Simone ha finito con mezz'ora d'anticipo e apre alle domande (``go crazy'').
Diego apprezza molto il finale, ``smielato ma in realtà estremamente
veritiero'': lui è più orientato all'ambito professionale, ma crede che in
Agile ci debba essere davvero questa fiducia, e che sia anche una delle sue
forze. Simone pensava che non l'avrebbe apprezzato, tanto era smielato
(``mentre lo scrivevo ho percepito il diabete che aumentava''); Diego gli
suggerisce, la prossima volta, di aggiungere immagini con gli occhi a cuore e
qualche gattino.

% ---------------------------------------------------------------------
\section{Domande a Simone}
\label{sec:z3-qa-simone}

\subsection{La dimensione del team allunga le riunioni?}

Uno studente, che lavora in un team di tre persone (``essendo in tre le cose
vengono un po' più facili''), chiede quanto la \textbf{grandezza del team}
influisca sui tempi delle riunioni e delle cerimonie, e se ci sia un modo per
mantenerle rapide comunque.

Influisce assolutamente, ``per una mera questione matematica''. Ma in genere
è possibile far sì che un team di qualsiasi dimensione completi una riunione
in \textbf{una trentina di minuti al massimo}: quello è l'obiettivo, ``anche a
costo di dire meno cose, anche a costo di essere meno precisi''. A meno di
avere un team di una cinquantina di persone, la riunione durerà sempre trenta
minuti; e con cinquanta persone non si fa comunque una riunione tutti insieme,
perché è impossibile: si fanno le \textbf{riunioni di reparto}, e poi i
\textbf{capi reparto} fanno la loro riunione di sprint planning. I numeri
tornano piccoli: a una riunione ci saranno al massimo una quindicina di
persone, venti ``se proprio i pianeti si sono allineati'', e venti persone
possono stare in trenta minuti.

\subsection{Quanto deve imporsi lo scrum master? Il timekeeper}

Lo stesso studente chiede quanto lo scrum master debba \textbf{imporsi} sugli
altri, sul piano umano, per far rispettare i tempi. Moltissimo, risponde
Simone, ma \textbf{non deve per forza essere lo scrum master a farlo}. Nel
team australiano in cui lavora, all'inizio di ogni riunione si sceglie, in
modo pseudo-casuale, tra chi è disponibile o lanciando un dado, un
\textbf{timekeeper}: la persona che ha ``lo scettro del comando'' e tiene in
mano un cronometro, perché nessuno dedichi ai propri report più tempo del
dovuto. È \textbf{comunemente accettato da tutti} che quella persona abbia il
diritto di ``rompere le scatole'': è socialmente accettato.

Se le riunioni si dilungano troppo, si può proporre di introdurre la regola
del timekeeper, scelto ogni volta a caso o sempre lo stesso: ha il diritto di
richiamare gli altri perché, oggettivamente, il tempo scorre.

\subsection{E se il problema blocca tutti?}

Ultima domanda dello stesso studente: si è detto che i problemi si risolvono
dopo la riunione; ma se un problema è molto grave, se il fatto di non averlo
risolto blocca la produzione di più persone (``praticamente un disastro''),
vale la pena fermarsi e parlarne tutti insieme, o si dedica comunque quel
tempo solo ai report e poi si fa una riunione in più?

La risposta semplice è \textbf{no}: non si risolve nella riunione di sprint
planning e post-mortem. ``L'unica cosa che mi viene in mente che possa essere
così tanto grave è se la stanza in cui state facendo il meeting va a fuoco.''
La riunione non è pensata per quello. Subito dopo, piuttosto, si avvia
un'altra riunione con tutte le persone necessarie (``che durerà 25 ore'') per
risolvere il problema. Ma finché si fa planning e post-mortem, si fa
\textbf{quello e basta}, anche se il problema riguarda tutti.

\subsection{E se mancano delle persone?}

Si fa comunque la riunione? La si rimanda? Simone la fa \textbf{sempre}, a
meno di essere da solo, per \textbf{rispetto delle persone che sono venute}.
Le assenze possono avere cause diverse. Chi non è venuto deve mandargli un
messaggio su Telegram, anche vocale, in cui fa di fatto il proprio report.

Se una persona manca spesso, Simone ci parla e cerca di capire perché.
Magari l'orario della riunione non va bene perché ha altro da fare: nel
gruppo di \emph{ISO} c'è un ragazzo che lavora a settimane alterne di mattina
e di pomeriggio, e quando lavora di pomeriggio, tornando a casa, trova spesso
traffico; quindi a settimane alterne non può venire, o arriva molto in
ritardo. ``Lo so, me lo posso aspettare, non è un problema.'' Se invece
qualcuno manca in modo pseudo-casuale, bisogna capire il motivo: non gli va
bene l'orario? Ne preferirebbe un altro? Pensa che le riunioni non servano?
Non ha voglia di farle? Da lì si capisce il problema e si cerca di
risolverlo, spiegando che \textbf{le riunioni sono importanti} e che è
importante esserci tutti, e chiedendo ``dimmi come possiamo aiutarti a
rendertele meno pesanti''.

\subsection{Quando è la persona a non adattarsi}

Se però, nonostante tutto, la persona ``continua a fare la \emph{drama
queen}'', a un certo punto bisogna ammettere che esistono anche persone non
adatte a lavorare in un certo contesto, e si può valutare se lo sia.

Diego interviene per completare: in Agile tutti vogliono ottenere la cosa
migliore per il team e per il metodo di lavoro di ciascuno. Ma se un individuo
non si adatta al team, o il team non riesce ad adattarsi a lui, che sia o
meno una \emph{drama queen}, \textbf{la presenza di quella persona nel team va
valutata}. ``Nessuno, a parte casi eccezionali, tipo Miyamoto, vale il lavoro
degli altri.'' Le \emph{soft skill} restano alla base. E aggiunge che non è
solo questione di soft skill: ci sarà la \emph{drama queen}, quello che non si
adatta, ``lo stronzo o la stronza'' (il genere, precisa, è irrilevante),
quello che non ha voglia\footnote{Nella trascrizione ``stronzo
stronza,umbaquin'': la parte finale è incomprensibile.}; ma c'è anche la
persona che ha voglia di fare ma non si adatta al team. Se non ci riesce, o il
team non riesce ad adattarsi a lei, potenzialmente deve andare via.

Simone conferma che è un punto molto importante, che rientra nel discorso
dell'\textbf{ascoltare di che cosa ha bisogno il team}: se il team non riesce
a lavorare con un individuo, ``non c'è niente di male, è così e basta''. Per
fortuna non è troppo comune: a un certo livello professionale è difficile che
succeda, soprattutto in questa industria, dove le persone, se non stessero
lavorando insieme, potrebbero diventare amiche, perché hanno un interesse in
comune e sono molto appassionate. Ma può succedere.

\subsection{Come si calcola la velocity, in pratica}

Diego invita gli studenti ad approfittare di una persona ``che ha fatto il
vostro percorso''. Una domanda sul calcolo: la velocity è semplicemente la
\textbf{somma dei pesi dei task completati} nel tempo trascorso, cioè nella
sprint? Esatto. Si fa la sommatoria dei task \textbf{taggati come
\emph{done}} che hanno la stessa sprint (o la stessa data di consegna, a
seconda di come la si indica): si prendono tutti i task con ``sprint 5'', di
questi quelli in stato \emph{done}, e se ne sommano i pesi. Allo stesso modo
si possono sommare i task non completati, quelli rimandati, e così via.

\paragraph{Uno sguardo ai fogli personali.} Simone vorrebbe mostrare il foglio
di \emph{ISO}, ma è difficile da leggere per la quantità di dati; si chiede se
mostrare i propri fogli rompa qualche NDA, poi nota che sta già condividendo
lo schermo, ``quindi gli NDA sono già stati rotti''. Mostra allora uno dei
suoi fogli personali: organizza il proprio lavoro \textbf{settimana per
settimana}, con il nome della settimana, l'elenco dei task, la somma di ciò
che era da fare, quanto è stato effettivamente fatto, quanto è stato tagliato
e quanto rimandato per cause diverse. Il risultato è un grafico con
l'andamento settimanale di questi valori.

\paragraph{La ``speedrun velocity''.} Poi decide di fare ``una cosa più
figa'' e di costruire un foglio in diretta, che chiama \emph{speedrun
velocity}. La struttura minima\footnote{I riferimenti alle celle detti a voce
durante la costruzione (``B2B'', ``G1'', ``F1'', ``C2, F2'') sono confusi
nella trascrizione; la tabella e la formula sono una ricostruzione dello
schema, non una copia esatta del foglio mostrato.}:
\begin{center}
  \begin{tabular}{@{}llll@{}}
    \toprule
    \textbf{Task name} & \textbf{Score} & \textbf{Sprint} & \textbf{Status} \\
    \midrule
    Task 1 & 5  & Sprint 1 & done \\
    Task 2 & 13 & Sprint 1 & WIP \\
    Task 3 & 3  & Sprint 2 & cut \\
    \dots  & \dots & \dots & \dots \\
    \bottomrule
  \end{tabular}
\end{center}
Nella colonna \emph{Status} inserisce una convalida dei dati con i tre valori
\emph{cut}, \emph{WIP} e \emph{done}, ``per semplificarsi la vita''. Poi
scrive una serie di task con punteggi a caso (copiati da un altro foglio) e
assegna le sprint, avvertendo che ``ovviamente non dovrebbe essere fatto
esattamente così, ma \emph{for the sake of showcase}''; aggiunge qualche
colore, perché ``sono pazzo'', anche se non è la cosa fondamentale.

Il calcolo è una somma condizionale su più criteri (\texttt{SUMIFS}, in
italiano \texttt{SOMMA.PIÙ.SE}): si somma la colonna dei punteggi per le
righe in cui la sprint è quella cercata \emph{e} lo stato è \emph{done}. Ad
esempio:
\begin{center}
  \texttt{=SUMIFS(B:B; C:C; "Sprint 1"; D:D; "done")}
\end{center}
Lo stesso si può fare con \emph{cut} per sapere quanto è stato tagliato. Si
ripete per la sprint~2, e così via, e con i risultati si inserisce un
\textbf{grafico a linee} (Sheets propone il tipo di grafico sbagliato, e
Simone lo corregge). Con più sprint si ottiene un andamento come quello dei
suoi fogli personali o di \emph{ISO}. Chi vuole fare le cose per bene può
aggiungere altre analisi: chi ha fatto più task, chi ha più peso e carico di
lavoro, lo stato generale del progetto (``magari il backlog fosse così
vuoto'').

Era, dice, per essere sicuro di riuscirci; è disponibile a rifare qualche
passaggio, ma alla classe è tutto chiaro. Uno studente chiede se il foglio si
possa avere: Simone lo mette in una cartella e sarà Diego a farlo avere alla
classe.

\subsection{I problemi più gravi incontrati con Agile}

Diego osserva che Simone ha capito come applicare i principi di Agile con una
velocità sorprendente, più di quanto sia successo a lui quando gli fu
spiegato, e che all'inizio si fa spesso ``Waterfall Agile'', cioè Waterfall
con le sprint, credendo di fare Agile\footnote{Il passaggio è frammentato e non
è del tutto chiaro chi abbia fatto ``Waterfall Agile'' all'inizio: dalla
risposta di Simone (``ho fatto anch'io waterfall agile per un po''') sembra
trattarsi di un errore comune, che Diego attribuisce anche a se stesso o agli
studenti in generale.}. Poi la domanda: quali sono stati i problemi più gravi?

\paragraph{I task ``tenuti in caldo''.} Simone si è accorto che le persone non
chiudevano i task al momento giusto. Se diceva che un task richiedeva, per
esempio, quattro ore, qualcuno lo finiva in un'ora e poi, invece di segnarlo
come completato, aspettava (``me lo rivedo più tardi, tanto ho ancora altre
tre ore''). Non funziona così: \textbf{quando un task è completato, è
completato}, punto. Non c'è ``me lo rivedo tra un po''' né ``devo farlo in
queste quattro ore, quindi ci metto quattro ore''. Questo falsava i risultati,
e Simone ha cambiato subito approccio\footnote{L'esempio parla di stime in
ore, in apparente contrasto con il consiglio di non usarle
(sezione~\ref{sec:z3-velocity}); è plausibile che si riferisca proprio a una
fase precedente in cui le ore erano ancora usate, e che il ``cambiamento
immediato'' sia stato anche questo.}.

\paragraph{Le riunioni che nessuno voleva fare.} Durante il secondo anno le
persone non volevano fare le riunioni settimanali. All'inizio le facevano solo
tra designer; poi le hanno estese a tutti, perché non aveva senso limitarle,
ma la risposta era ``questa riunione è inutile, non ha senso''. In quel caso
Simone non è riuscito a risolvere il problema: non sapeva come farlo. C'erano
problemi anche sul piano umano, non tra i designer ma con membri di altre
classi: problemi che, col senno di poi, si sarebbero potuti risolvere, ma
all'epoca non ci sono riusciti. La soluzione adottata è stata non fare le
riunioni con tutto il team: gli altri membri mandavano un \textbf{recap
settimanale per messaggio}. Non era un metodo del tutto fallimentare, ma si
perdono molti benefici delle riunioni tutti insieme: il \emph{team building},
il confronto, il vedersi, il parlare. È stato, dice, ``un grosso problema e
anche un grosso fallimento''.

Per il resto, niente di così grosso da valere la pena di raccontarlo:
valutazioni di task sbagliate, sprint fatte male, e un periodo di
``Waterfall Agile'' anche per lui, finché non ha capito di sbagliare.

\subsection{Che cosa rispondere a chi dice che le riunioni sono inutili}

E oggi, se qualcuno dicesse ``per me queste riunioni sono inutili'', che cosa
gli risponderebbe? \textbf{``Perché?''}. E da lì si capisce qual è il
problema: non vuole dedicarci tempo? Durano troppo? Gli costa fatica? Ha
problemi con le riunioni perché soffre di qualche forma d'ansia e non gli
piace stare con tutti i membri del team, anche fuori dalle riunioni (``tipo
ci becchiamo a giocare online a qualcosa'')? Bisogna capire il perché e
\textbf{tirarlo fuori dalle persone}, perché spesso le persone stesse non sanno
il perché delle cose che dicono e sentono.

\subsection{I saluti a Simone}

Diego ringrazia Simone e spera che gli studenti abbiano appreso il più
possibile: un'esperienza più diretta di quella di Simone e Cristian non si può
avere, perché hanno fatto letteralmente lo stesso percorso (la lezione di
Diego è ``fondamentalmente la stessa'' da un anno all'altro, con qualche
aggiunta). Ricorda che Simone ha concesso anche il foglio della
\emph{speedrun velocity}, che Diego passerà alla classe, e lo saluta per non
rubargli altro tempo.

% ---------------------------------------------------------------------
\section{Cristian: il punto di vista del game designer e project manager}
\label{sec:z3-cristian}

\subsection{Chi è il relatore}

Il secondo intervento è di Cristian, che attiva la webcam (``mi vedete super
zoomato, spero non sia una brutta vista'') e condivide le slide, mentre Diego
cerca di farle andare in modalità presentazione perché vengano meglio nella
registrazione. Cristian ringrazia Diego per la presentazione e si presenta:
\begin{itemize}
  \item è uno \textbf{studente del terzo anno} di Game Design, quindi ha fatto
        lo stesso percorso della classe;
  \item è anche \textbf{tutor} del corso di Game Design per gli studenti
        dell'accademia online;
  \item è \textbf{game designer e project manager} di due progetti:
        \begin{itemize}
          \item \emph{Gamers in Time}\footnote{Così nella trascrizione;
                il titolo non è stato verificato.}, un \textbf{gioco da
                tavolo}, vincitore del \emph{Big One} del primo anno;
          \item \emph{L'Inf}\footnote{Il titolo compare in molte forme
                (``l'Inf'', ``Linf'', ``Leanf'', ``Linfa''); qui si usa
                \emph{L'Inf}, la più frequente. Non è chiaro se la grafia
                corretta sia \emph{Linf}, \emph{L'Inf} o \emph{Linfa}.}, un
                \textbf{videogioco} che ha vinto, ``tra virgolette'', il
                \emph{Bigger One} del secondo anno e che il team ha avuto
                l'opportunità di continuare come \emph{Biggest One} del terzo
                anno, opportunità che ha accettato.
        \end{itemize}
\end{itemize}
In entrambi si è occupato di tutta la parte di design, insieme ad altri
colleghi, ed è per lo più anche \textbf{project manager, scrum master,
producer}, ``come volete definirmi''.

\subsection{Come è strutturata la sua parte}

Oltre all'esperienza personale, la lezione ha due parti principali:
\begin{enumerate}
  \item la \textbf{produzione} effettiva tramite Scrum;
  \item la \textbf{gestione del team} in un ambiente di lavoro Scrum.
\end{enumerate}
Possono sembrare la stessa cosa, ma ``il progetto tendenzialmente è una cosa
e il team è un'altra'': ci sono ovvi rapporti tra le due parti, ma per
ciascuna ci sono precauzioni specifiche.

\subsection{Le premesse}

\begin{itemize}
  \item Quello che descrive è \textbf{un} percorso, un modo di applicare
        Scrum, \textbf{non l'unico}: se non si è d'accordo, o se il team non
        ``sopporta'' questo metodo, si possono fare tutti i cambiamenti che si
        vuole (e qualche consiglio lo darà strada facendo).
  \item Non è una lezione descrittiva come quelle di Diego (che cos'è Scrum,
        le definizioni, le cerimonie, come si calcolano la velocity e il peso
        dei task), ma un \textbf{confronto} sull'esperienza dei due progetti,
        più su \emph{L'Inf} che su \emph{Gamers in Time}, perché Scrum l'ha
        applicato al videogioco e non al gioco da tavolo.
  \item \textbf{Tutto è in continua evoluzione}. Per questo vuole le domande
        della classe: il primo obiettivo è chiarire i dubbi e far uscire tutti
        dalla lezione avendo capito un po' di più la mentalità agile, ``con
        tutti voi stessi''; il secondo è ricevere \textbf{feedback} per
        migliorare non solo le lezioni, ma lo stesso modo in cui applica
        Scrum.
  \item Ha un po' di mal di gola e si fermerà spesso a bere: si scusa in
        anticipo.
\end{itemize}

% ---------------------------------------------------------------------
\section{Il caso di studio: \emph{L'Inf}}
\label{sec:z3-linf}

\subsection{Le caratteristiche del gioco}

Qualcuno in classe forse conosce Cristian o ha provato il gioco a Level Up,
l'evento in cui è stato presentato. Le caratteristiche principali:
\begin{itemize}
  \item è \textbf{single player} e offline: niente co-op, niente
        multiplayer, niente server;
  \item è un'\textbf{avventura fantasy}: un gioco di \textbf{durata fissa}, che
        tendenzialmente si gioca per una \emph{run}; chi l'ha apprezzato di
        più e vuole vedere tutti gli aspetti del design e delle meccaniche, o
        esplorarne i segreti, lo rigiocherà;
  \item il giocatore ha delle \textbf{magie colorate} con cui affronta
        \textbf{orde di nemici colorati}: spara le magie e, a seconda di
        quanti nemici ``scoppiano'' con un solo colpo, ottiene vari bonus,
        come succede in \emph{Candy Crush}.
\end{itemize}

\paragraph{Dati e passione.} Diego interrompe perché ritiene questa parte molto
importante: com'è nato l'elemento davvero innovativo di \emph{L'Inf}? Dallo
\textbf{studio dei dati}, ma anche dalla \textbf{passione del team}. I dati
danno una direzione: le ricerche di Cristian indicavano il successo delle
meccaniche \emph{match-3}\footnote{Nella trascrizione ``match free'' e
``match tree'': si tratta delle meccaniche \emph{match-3} (allineare tre o
più elementi dello stesso colore), tipiche di \emph{Candy Crush}.}; e
Cristian è anche, principalmente, un giocatore mobile. La combinazione tra
passione e ricerca ha prodotto un progetto molto valido sul piano della
fattibilità (single player, fantasy, espandibile).

Cristian aggiunge un'altra ricerca fondamentale, da cui è stato poi delineato
il \textbf{documento di visione}\footnote{Nella trascrizione ``iVision'' e
``l'iVision'': si tratta quasi certamente della \emph{vision}, il documento
di visione (sezione~\ref{sec:perche-documentare}).} del gioco: quella su
che cosa chiedesse Epic Games per finanziare i progetti con l'\textbf{Epic
MegaGrant}. Il team ha individuato tre elementi fondamentali\footnote{È la
lettura dei requisiti fatta dal team, come riferita da Cristian; non è stata
verificata sui criteri ufficiali del programma.}:
\begin{enumerate}
  \item \textbf{innovazione};
  \item \textbf{temi ambientali}: attraverso i colori e l'ambientazione, il
        gioco vuole criticare l'inquinamento globale;
  \item un \textbf{elemento narrativo}.
\end{enumerate}
Questi tre elementi, le ricerche specifiche sui vari tipi di gioco e la
passione (di Cristian in primis, come lead, ma anche degli altri membri,
incuriositi dalla proposta) hanno portato a un progetto che al momento sembra
funzionare: ``stiamo andando nella direzione giusta, almeno secondo quanto
previsto; non si può mai prevedere al 100\% ciò che accadrà''.

\paragraph{Il contenuto.} Oltre alla meccanica \emph{match-3}, il gioco
prevede:
\begin{itemize}
  \item \textbf{quattro modalità di sparo};
  \item \textbf{cinque livelli}, con esplorazione e fasi di combattimento
        contro le orde;
  \item \textbf{puzzle};
  \item \textbf{cinematic} per lo più ambientali: niente attori, doppiatori o
        mimica facciale, ma scene di presentazione dell'ambiente di gioco;
  \item una \textbf{boss fight} finale.
\end{itemize}
Il perché di questo elenco sarà chiaro tra poco.

\subsection{Un ``trauma'' di produzione: il cambio di team}

La produzione si è sviluppata in due sezioni, e \emph{L'Inf}, ``per fortuna o
purtroppo'', ha avuto un grave problema: è la risposta, anticipata, alla
domanda fatta a Simone sui problemi più gravi. Come tutte le produzioni di
videogiochi, ha subito dei ``traumi''. Il più grande: \textbf{il 100\% del
team originario è stato sostituito}. Non per scelta, ma perché i membri non si
sono iscritti al terzo anno: avevano altri obiettivi; c'è chi cerca lavoro e
chi ha aperto un'azienda. Cristian, volendo continuare il terzo anno e
soprattutto \emph{L'Inf}, un progetto in cui crede e che ha dimostrato il suo
potenziale a Level Up e con i feedback costanti dei giocatori, ha dovuto fare
dei \textbf{colloqui} per formare un nuovo team che si adattasse meglio al tipo
di sviluppo del gioco\footnote{La frase è in parte persa nella trascrizione
(``ho dovuto fare dei colloqui Grazie a tutti. Cioè, servisse una parola
brutta, ma semplicemente si adattasse di più al tipo di sviluppo''): il
senso ricostruito è che ``colloqui'' è una parola forte, ma che si trattava
di scegliere persone adatte a questa fase.}.

La composizione è cambiata così:
\begin{center}
  \begin{tabular}{@{}lp{0.55\textwidth}@{}}
    \toprule
    \textbf{Team} & \textbf{Composizione} \\
    \midrule
    originario & 2 game designer, 2 level designer, 1 technical designer \\
    attuale    & 3 game e level designer, 2 technical designer \\
    \bottomrule
  \end{tabular}
\end{center}
Il motivo: la preproduzione, cioè la parte più importante di game design, era
finita, e servivano \textbf{più tecnici}. Il gioco si basa molto sulle
meccaniche: un sistema di combattimento innovativo che richiede di scriptare
le meccaniche, diversi tipi di sparo, puzzle, una boss fight. È molto più
basato sulle meccaniche che su narrativa o ambientazione. Allo stesso tempo
servivano level designer, perché il gioco doveva prendere forma e i livelli
andavano costruiti.

\subsection{Perché raccontare tutto questo}

Non per vendere il gioco (``anche se magari è un effetto collaterale ben
voluto''), ma per mostrare che \textbf{ogni progetto, a seconda di come è
composto} sul piano del design, della produzione e soprattutto del team,
\textbf{ha un modo diverso di applicare Scrum}. Per questo Cristian è sceso nel
dettaglio delle caratteristiche: ognuna comporta ambiti di produzione diversi.
\begin{itemize}
  \item I livelli richiedono level designer.
  \item Il \emph{match-3} e i diversi tipi di combattimento richiedono
        tecnici.
  \item Il single player è più un vantaggio che un limite: un gioco
        multiplayer non è il doppio, ma ``tipo dieci volte più difficile'' da
        gestire a livello produttivo.
  \item Il genere \emph{fantasy adventure} deriva soprattutto dalle ricerche.
\end{itemize}
Quindi \textbf{ogni progetto ha il suo framework di Scrum}. E il fatto di
aver cambiato team a metà produzione, finita la \emph{vertical slice}, con
persone completamente nuove (ex compagni di studio, che però non avevano mai
visto come funzionava il gioco), ha avuto diverse conseguenze, che Cristian
rimanda alle domande di fine lezione.

% ---------------------------------------------------------------------
\section{Il loop di Scrum}
\label{sec:z3-loop}

L'applicazione di Scrum è un \textbf{loop} molto semplice, simile a quello
mostrato da Simone:
\begin{enumerate}
  \item \textbf{creazione del backlog};
  \item \textbf{selezione dei task} da svolgere, a ogni sprint;
  \item \textbf{implementazione dei task} durante la sprint, con controlli e
        \emph{scrum} (con la frequenza che si decide);
  \item \textbf{calcolo della velocity}, per vedere com'è andata e come
        migliorare la produzione nelle sprint successive;
  \item \textbf{interpretazione} della velocity per fare il \emph{tuning}
        della produzione; e si ricomincia.
\end{enumerate}
Cristian le affronta una per una, con esempi della sua esperienza e consigli.
È contento che la classe stia cominciando proprio adesso, perché questa parte
potrebbe tornare utile già dalla settimana successiva, quando gli studenti
applicheranno il \emph{framework}.

% ---------------------------------------------------------------------
\section{La stesura del backlog}
\label{sec:z3-backlog}

\subsection{Un backlog per milestone}

Anche Cristian, all'inizio, cercava di prevedere il più possibile tutto il
gioco, dal prototipo fino al rilascio. Oggi invece lo struttura
\textbf{per milestone}: il backlog dell'\emph{alpha}, poi quello della
\emph{beta}, poi quello della \emph{gold}. Questo non impedisce di cominciare
a prevedere che cosa ci sarà in beta e in gold, ma dà un'idea molto chiara di
dove si sta andando, soprattutto per capire che cosa tagliare o aggiungere
quando ci si accorge che qualcosa serve prima o dopo, e adattarsi di
conseguenza: ``la cosa più importante del mindset agile è adattarsi ed essere
flessibili''.

\subsection{Da dove vengono i task}

Il gioco deve pur nascere da qualche parte: i task vengono dai
\textbf{documenti chiave} della produzione, già visti nel corso
(capitolo~\ref{cap:produzione-documentazione}):
\begin{itemize}
  \item il \textbf{documento di visione};
  \item i \textbf{GDD} delle varie feature, che dettagliano la visione;
  \item il \textbf{versioning} delle feature, a cui si dovrebbe pensare già
        mentre si scrive il GDD;
  \item gli \textbf{LDD} (\emph{level design document}) e tutti gli altri.
\end{itemize}
L'ordine dipende dal gioco: in un gioco fortemente basato sulle meccaniche è
probabile che gli LDD arrivino dopo lo sviluppo del \emph{core}; in un gioco
basato sull'esplorazione non succede proprio l'inverso, perché le meccaniche
servono comunque, ma le due cose avvengono quasi in contemporanea.

Una volta che si ha la documentazione con tutte le descrizioni e le
caratteristiche delle feature, \textbf{dalla versione 0 alla versione X}, si
ha già una serie di implementazioni da realizzare. Un esempio:
\begin{itemize}
  \item \textbf{versione 1}: il giocatore deve camminare; si usa il
        manichino base di Unreal;
  \item \textbf{versione 2}: si implementa la \emph{skeletal mesh} fatta dal
        3D artist e si adatta la camera al movimento, con gli effetti di
        vibrazione a seconda che il personaggio corra, cammini o sia
        abbassato;
  \item \textbf{versioni successive}: quando il personaggio si muove su
        terreni diversi solleva tipi diversi di polvere, di più su un terreno
        molto polveroso. Non è fondamentale e potrebbe essere una delle
        ultime versioni, verso la fine.
\end{itemize}
Il consiglio è che, quando si fa il \emph{versioning}, si inseriscano
\textbf{tutti gli ambiti necessari} per quella feature: programmazione, arte
(se serve un asset), animazione, UI, VFX, suoni. Se tutti i \emph{field}
compaiono già nella documentazione, la creazione del backlog sarà più chiara
e più spedita.

\subsection{Come organizzare i task}

I giochi hanno moltissimi task. Come organizzarli perché siano leggibili non
solo da chi li ha scritti (scrum master o project manager), ma anche da chi
deve navigarli per capirci qualcosa? Cristian usa tre tipi principali di
organizzazione.

\paragraph{1. Per feature (o epic).} Si suddividono i task per feature,
collegandosi alle \emph{epic} di Simone. Ad esempio, l'epic è il movimento
del personaggio e i task sono camminata, corsa, salto, volo, ``dipende dal
gioco''.

\paragraph{Tutto in un solo posto.} La suddivisione non va fatta per forza in
documenti diversi, anzi Cristian lo sconsiglia: se tutto è nello stesso
strumento (Google Sheets, Monday, ClickUp, Trello) si usano i
\textbf{filtri} e si trova tutto lì, senza navigare tra mille siti o mille
documenti.

\paragraph{Per ambito, con un color code.} Quando si suddivide per ambito
(programmazione, arte, UI, eccetera), conviene usare un \textbf{codice
colore}. Cristian lo usa ``per qualsiasi cosa'': uno per le milestone, uno
per gli ambiti di produzione. Diego gli chiede di mostrare il suo:
\begin{center}
  \begin{tabular}{@{}lll@{}}
    \toprule
    \textbf{Ambito o milestone} & \textbf{Colore} & \textbf{Perché} \\
    \midrule
    Suoni        & viola        & --- \\
    VFX          & rosso        & sta imparando a usare Niagara \\
    Level design & verde        & ``erano finiti i colori'' \\
    Game design  & arancione    & è il suo colore preferito \\
    \midrule
    Alpha        & azzurro      & --- \\
    Beta         & viola        & --- \\
    Gold         & giallo scuro & il colore richiama il nome \\
    \bottomrule
  \end{tabular}
\end{center}
Come si vede, è un misto di funzionalità e preferenze: la cosa più importante
è che \textbf{tutto il team sia d'accordo}.

\paragraph{2. Per milestone, con le sotto-milestone.} Un secondo metodo è
suddividere il backlog per milestone. Ma se il backlog dell'alpha è separato
da quelli di beta e gold, come si fa a dividere in alpha, beta e gold un
backlog in cui tutto è alpha? Cristian suddivide le milestone in
\textbf{sotto-milestone}: l'alpha dell'alpha, la beta dell'alpha, la gold
dell'alpha. Serve ad avere idee più chiare, a tenere sempre in mente lo
\emph{scope} senza perderlo e soprattutto a capire meglio \textbf{se il gioco
è fattibile}. Una suddivisione più fine permette anche di calcolare con più
precisione (una precisione che ``lascia il tempo che trova'') quante sprint
servono per ogni sotto-milestone. La dissonanza tra ``dividete tutto in alpha,
beta e gold'' e ``dentro l'alpha trovate alpha, beta e gold'' è quindi voluta:
si parla delle sotto-versioni.

\paragraph{3. Per priorità.} Chi non vuole usare le sotto-milestone può usare
il sistema di \textbf{priorità} di cui parlava Simone: una scala numerica, da
1 a 5 o da 1 a 10, a seconda di come si riesce meglio a distinguere le
priorità. Serve soprattutto a capire quali task selezionare
(sezione~\ref{sec:z3-selezione}). In realtà è più o meno la stessa cosa con un
approccio diverso: indicare alpha, beta e gold dell'alpha crea già di per sé
delle priorità. ``Dipende tutto da voi.''

\subsection{Come scrivere i task e dove}

Il consiglio è scriverli \textbf{dal punto di vista dell'utente}. Ma, come
diceva Simone, non sempre ``il gioco vale la candela''. Lo è soprattutto in
un caso: se il programmatore legge un task dell'artista 3D scritto in modo
troppo tecnico e non lo capisce, si crea dissonanza tra i reparti e possono
nascere attriti. Quindi è sempre consigliabile, ma non obbligatorio: spesso
basta una chiamata (``prima semplicemente ti giravi verso il tuo compagno e
gli chiedevi che cosa significa''). ``Però è giusto che vi dica le cose come
dovrebbero essere fatte.''

Dove scriverli? In qualsiasi strumento: Trello, Monday, ClickUp, gratuiti o a
pagamento, oppure un foglio Google costruito da sé. Dall'esperienza di
Cristian emerge che \textbf{Google Sheets sta antipatico alle persone}: già il
nome ``foglio di calcolo'' spaventa. Strumenti come Monday o ClickUp hanno più
colori e soprattutto \textbf{mandano le notifiche}, quindi sono più accessibili
a tutti. ``Decidete voi.''

\subsection{Non prevedere tutto: la flessibilità}

È fondamentale capire che \textbf{prevedere tutti i task passo per passo}, in
un ordine preciso (``dopo la camminata faremo il salto, dopo il salto il
sistema di sparo, e così via all'infinito''), va contro i principi di Agile,
in particolare contro il \textbf{quarto valore}: rispondere ai cambiamenti più
che seguire un piano.

Da qui il consiglio centrale: la \textbf{flessibilità} deve essere il cuore
della propria filosofia di lavoro e del proprio modo di applicare Scrum.
Bisogna essere sempre pronti a \textbf{tagliare} nel modo più agile possibile:
se si è capita davvero la mentalità agile, si riesce a tagliare una feature
\textbf{senza bloccare la produzione delle altre}.

L'esempio più eclatante viene dall'inizio di \emph{L'Inf}: Cristian faceva il
backlog di tutto il gioco, arrivando a circa \textbf{400 task} di ogni ambito,
nel dettaglio: per la camminata arrivava a indicare come si muoveva la camera
quando il personaggio girava a destra, e come cambiava quando saltava o cadeva
da un dirupo. ``Vi viene subito da pensare: ma come fai a prevedere tutto?''
Infatti la soluzione è \textbf{non prevedere tutto}: capire a grandi linee ciò
che si vuole, qual è l'obiettivo della feature, senza scendere troppo nel
dettaglio, perché il dettaglio crea di per sé un \textbf{lavoro Waterfall}, che
è ciò che bisogna assolutamente evitare. E scrivere un backlog così dettagliato
richiede tantissimo tempo, che è limitato: la produzione ne risente in tutti i
modi possibili.

\paragraph{Bonus: gli elementi innovativi.} Se il gioco (il proprio
\emph{Bigger One}) ha elementi innovativi, la flessibilità va portata
\textbf{ai massimi livelli}. \emph{L'Inf} era partito con una meccanica di
\textbf{assorbimento e sparo}: si assorbe un nemico colorato, ad esempio
giallo, e lo si spara contro altri nemici gialli, prendendo ispirazione da
altri giochi. Nel loro caso, però, creava più problemi che altro: per due o
tre ore di gioco, ``assorbire, sparare, assorbire, sparare'' risultava
\textbf{troppo ripetitivo}, e hanno deciso di tagliare la meccanica. Il
problema è che l'hanno tagliata dopo che Cristian aveva già scritto
\textbf{600 task nel backlog e 12 GDD}: tempo perso. Per questo la
flessibilità, negli elementi innovativi ma in generale, ``è la prima cosa''.

\subsection{Quanto dettaglio?}

Quando un backlog è abbastanza dettagliato e quando lo è troppo poco?
``La risposta purtroppo non c'è.'' Viene da due fattori:
\begin{enumerate}
  \item l'\textbf{esperienza}: man mano che si scrivono GDD e si va avanti
        nello sviluppo, ci si ``tuna'';
  \item il \textbf{trial and error}: si sbaglia e si torna indietro, il che
        riporta, di nuovo, all'esperienza.
\end{enumerate}

% ---------------------------------------------------------------------
\section{La selezione dei task}
\label{sec:z3-selezione}

Con un bellissimo backlog ordinato per colori, ambiti e priorità, quali task
si selezionano? ``Da qualche parte bisogna pur cominciare.'' Le domande
principali sono due: \textbf{quali} selezionare e \textbf{quanti}.

\subsection{Quali task}

\paragraph{La priorità.} Il primo criterio è la priorità: i task da fare
prima sono le scelte più urgenti, e quindi si scelgono quelli\footnote{L'inizio
di questo passaggio manca nella trascrizione (``quanti selezionare. cui
scegliere prima, sarà la scelta più urgente da fare''): il primo criterio è
stato ricostruito dal contesto, anche perché poco dopo Cristian dice che,
risolto il problema delle competenze, ``si ritorna a selezionarli per
priorità''.}.

\paragraph{Le capacità attuali degli sviluppatori.} Il secondo criterio sono
le \textbf{abilità dei developer}. Cristian precisa che non intende in alcun
modo giudicarne la professionalità: intende letteralmente \textbf{che cosa sa
fare lo sviluppatore in quel momento}. Se un programmatore non ha mai usato
Unreal, non è il momento di dargli task di implementazione in Unreal: bisogna
prevedere anche il \textbf{periodo di apprendimento}, eventualmente come task
vero e proprio (``hai questa sprint per imparare a implementare l'IA'').

È successo due mesi prima con uno dei suoi technical designer: gli ha dato
\textbf{due settimane} dicendogli ``in queste due settimane non ti preoccupare
di niente, il resto lo produciamo noi''. Doveva essere d'accordo, e lo era,
anzi era contento di approfondire l'argomento: studiare come funzionano le IA
(anche grazie a vari link passati da Diego) e soprattutto come implementarle
nel motore. Per due settimane ha dovuto ``solamente'' pensare
all'apprendimento. Il consiglio: \textbf{prevedere questi ``tempi morti''},
perché ne ha bisogno il team, e ne ha bisogno chi pianifica per non falsare la
previsione dello \emph{scope}. Se in questo momento (``sottolineo, in questo
momento'') un programmatore non sa fare una cosa, la si rimanda e si mette
altro; prima o poi imparerà, e a quel punto si torna a selezionare per
priorità.

\subsection{Quanti task}

Il primo criterio sulla quantità è la \textbf{velocity}, ma Cristian l'ha
messo nella slide solo per completezza del loop: se non si ha ancora una
velocity, non ha senso. Per ora la si trascura e si fa finta che sia la
\textbf{prima sprint}, la prima applicazione di Scrum.

Nelle prime sprint la quantità si basa sulle \textbf{capacità} e sulla
\textbf{disponibilità} degli sviluppatori:
\begin{itemize}
  \item \textbf{disponibilità}: per quanta parte della sprint il technical
        designer (per riprendere l'esempio) è disponibile? Non in termini di
        ore, perché Cristian concorda con Simone che le ore sono sempre molto
        relative;
  \item \textbf{performance}: gli studenti si conoscono ormai da quasi un
        anno e mezzo. Se uno sviluppatore non ha mai fatto un livello in
        Unreal (per motivi suoi, magari problemi al PC) e dice ``voglio fare
        un livello e ci metto due settimane'', sta a chi pianifica cominciare
        a capire \textbf{chi può fare che cosa}. Tutti possono fare tutto, se
        si prevede il tempo di apprendimento; ma se Cristian arrivasse a dire
        ``in due settimane vi faccio l'IA del vostro gioco implementandola in
        Unreal'', ``vi sto prendendo in giro''.
\end{itemize}
Queste sono le analisi preliminari quando non si ha una velocity. Se invece
non si conosce affatto il team, l'unico punto di partenza è la
\textbf{disponibilità} degli sviluppatori, da cui si stima quanto si riesce a
fare.

Tanto, \textbf{nelle prime sprint si sbaglieranno completamente tutte le
previsioni}, ``ed è giusto così, e ben venga'': se non si sbaglia,
probabilmente si sta applicando Waterfall, o si sta applicando male Scrum.
All'inizio si tratta solo di buttarsi: ``questa settimana abbiamo previsto 15
task; ne abbiamo fatti 40. Bene, la prossima ne proviamo 40''; poi magari se ne
fanno 25, e così via, \textbf{tarando} man mano la quantità.

\subsection{Definire bene il goal}

Per ogni sprint, e anche per ogni singolo task, bisogna \textbf{definire bene
l'obiettivo}. Serve a mantenere concentrati gli sviluppatori: se so che entro
questa settimana devo finire il block out del livello, so dove sto andando,
lavoro per quel risultato e soprattutto \textbf{non perdo tempo}, come succede
spessissimo. Se durante il block out qualcuno si mette a inserire effetti
visivi ambientali o il \emph{post-process}, viene il dubbio che non stia
lavorando per l'obiettivo: è un \textbf{campanello d'allarme}, perché potrebbe
far mancare l'obiettivo per cui la sprint era stata pensata.

\subsection{Bonus: progettare il ritmo dello sviluppo}

Un consiglio che Cristian ha sentito raramente, ma che ha deciso di
sperimentare: poiché gli studenti sono designer, e tendenzialmente anche lo
scrum master del loro progetto lo sarà, possono \textbf{progettare il
\emph{pace} dello sviluppo}.

Di solito si pensa (e potrebbe anche essere giusto) che l'obiettivo di Scrum
sia trovare una \emph{velocity target} e mantenerla costante per tutta la
produzione. Può essere una soluzione; ma con entrambi i suoi team Cristian ha
notato che, per come sono fatti gli esseri umani, ci sono \textbf{picchi di
attenzione e di produzione}. Con l'esperienza si impara a percepirli in
anticipo e a progettare di conseguenza; oppure sono gli sviluppatori stessi a
dirlo: se la settimana prima hanno implementato 12 nemici, 15 modalità di sparo
e 13 cutscene, quella dopo diranno ``andiamoci piano''. Di solito non succede,
perché le sprint sono programmate apposta per adattarsi alle capacità degli
sviluppatori, ma è un metodo possibile.

Che cosa intende, in concreto? Se per un periodo si sono implementate molte
feature o sistemi che però, visivamente, \textbf{non hanno portato un aumento
di qualità percepibile nel gioco}, gli sviluppatori non si sentono
riconosciuti e ripagati dello sforzo. Allora, la settimana dopo, ci si
concentra su qualcosa che dà \textbf{soddisfazione}. Questo lo permettono le
sprint, lo permette Scrum e non Waterfall: è una differenza fondamentale.

\paragraph{L'esempio di \emph{L'Inf}.} Per più di otto mesi il gioco è rimasto
con \textbf{quattro tipi di nemici}. Nonostante l'impegno (quattro nuove
modalità di sparo, la boss fight, tantissima programmazione ``di
background'', cioè non visibile al giocatore ma necessaria al funzionamento
del gioco), il team cominciava a percepire che \textbf{non c'era corrispondenza
tra impegno e qualità percepita}. Cristian ha allora deciso di limitare per un
periodo la produzione degli elementi di background (programmazione come l'IA,
tutto ciò che faceva proseguire il gioco ``sotto''), di inserire altri tipi di
nemici, che tutti non vedevano l'ora di avere (``sapevo che era uno di quei
jolly da sfruttare per far risalire il morale del team''), e di dedicare tempo
a effetti visivi e sonori, a migliorare tutto ciò che era estetico. Diego
suggerisce il termine: la \textbf{juiciness}. Ciò che vedeva il giocatore lo
vedevano anche gli sviluppatori, che si sentivano più appagati e sentivano i
propri sforzi ricompensati. ``Adesso loro non lo sanno, ma li ributterò sotto
un treno di lavoro'' (``in realtà lo sanno, però non me lo dicono'').

Per Cristian è un punto fondamentale, che gli è costato ``mesi di
sofferenze'' di esperienza per capirlo. Diego invita la classe a fare domande
su un \emph{insight} ottenuto per tentativi ed errori.

% ---------------------------------------------------------------------
\section{L'implementazione: controlli, blocker e feedback}
\label{sec:z3-implementazione}

\subsection{Il check dei progressi}

Selezionati i task, viene la parte ``più tranquilla'' dal punto di vista di
Scrum (i problemi di implementazione, ovviamente, sono altri): parte la
settimana di lavoro, ognuno ha i suoi task, i \emph{task owner} hanno scelto
che cosa fare, e questo prosegue per tutta la durata della sprint.

Durante l'implementazione si fa un \textbf{check dei progressi}. Non serve a
controllare che cosa stanno facendo le persone: di base bisogna
\textbf{fidarsi} che, se il programmatore dice ``farò questo'', lo farà entro i
limiti della sprint. Cristian parla qui di \emph{daily scrum} perché è ciò che
si dovrebbe fare in generale, anche se il suo team non li fa ogni giorno
(sezione~\ref{sec:z3-team}). Oltre alle solite domande (che cosa hai fatto
ieri, che cosa farai oggi), \textbf{la cosa più importante è trovare i
blocker}.

\subsection{Agire subito sui blocker}

Se il programmatore dice ``non ho potuto implementare il personaggio perché mi
manca l'asset, la \emph{skeletal mesh} del protagonista'', bisogna
\textbf{agire subito}, in due modi:
\begin{enumerate}
  \item \textbf{dare altri task} al programmatore, perché altrimenti resta lì
        ad aspettare, e questo non lo si vuole assolutamente;
  \item \textbf{capire che cosa è successo} alla skeletal mesh, parlando con
        la 3D artist. Se è solo questione di tempo, non è davvero un problema:
        è normalissimo aver assegnato l'implementazione prima che l'asset
        fosse pronto (``è normale sbagliare, l'importante è rimediare''). Se
        invece anche lì c'è stato un blocco, bisogna \textbf{risalire a
        cascata} fino a trovarlo.
\end{enumerate}
Magari si scopre che il problema sono proprio i game designer: quando hanno
fatto la \emph{mood board} non sono stati abbastanza chiari, e la 3D artist ha
dovuto spendere una settimana a cercare reference per conto suo, mentre il
designer pensava che quelle fornite bastassero e non si è posto il problema
della fase di ricerca delle reference da parte di chi fa la grafica o il 3D
(Cristian parla al femminile perché la 3D artist del suo team è una donna, ma
il discorso è generale). Anche questo, come l'apprendimento, è uno di quei
task che non bisogna trascurare.

\subsection{Analizzare i blocker prima di risolverli}

Il consiglio è \textbf{analizzare i blocker}, non risolverli appena si
presentano, perché si rischia di trascurare il problema alla radice. Cristian
richiama una delle prime lezioni del primo anno, quella su \textbf{bisogni e
problemi}: bisogna capire se il blocker è un problema o un bisogno, e agire di
conseguenza. Il guaio nasce quando lo si trascura o, peggio, quando lo si
risolve in fretta senza averlo considerato: ``poi vi si ritorcerà contro''.

\paragraph{L'esempio dei muri invisibili.} \emph{L'Inf} è un gioco di
esplorazione ambientale, per lo più all'aperto: su cinque livelli, circa
l'80\% del gioco è all'esterno. Per dare \emph{meaning} al livello finale,
quello della boss fight e della rivelazione narrativa su dove siano finiti
tutti gli esseri umani del pianeta, il team ha deciso di sottolinearne la
diversità rendendolo \textbf{al chiuso}.

La camera è in terza persona e abbastanza lontana dal giocatore, non classica.
Dare le collisioni alla camera sembrava fastidioso, perché a ogni rotazione
``sbatte'' contro i muri. La prima idea: \textbf{rendere invisibili i muri}
quando la camera ci passa dietro. Non si potevano far sparire tutti insieme
quando il giocatore era solo dietro il muro destro, quindi ogni muro aveva il
suo \emph{trigger}, che stabiliva quante e quali pareti dovessero sparire. Il
livello dura mezz'ora, con una buona decina di minuti di esplorazione al
chiuso, completamente circondata da pareti: \textbf{un mese di lavoro} solo
per assegnare i trigger, cambiare i materiali e gestire il tutto. ``Era l'idea
geniale che ci era venuta in quel momento.''

Quando hanno consegnato la beta e l'hanno fatta provare ai professori e ad altri
giocatori, tutti hanno detto più o meno la stessa cosa: i muri che spariscono
\textbf{non sono eleganti e rompono l'immersione}. Hanno tagliato tutto e
rimesso le collisioni alla camera. Se avessero analizzato bene il problema
prima, avrebbero capito che bastava \textbf{gestire meglio gli spazi}: non fare
curve strette ma allargarle, così che anche ruotando la camera il giocatore
tendenzialmente non sbatta contro il muro. Era ``assolutamente una questione di
riallargare più o meno tutti gli spazi''; ed è ciò che hanno fatto alla fine,
insieme alle collisioni. ``Ad averlo fatto prima avremmo risparmiato due
settimane di tempo''\footnote{Poco prima Cristian parla di ``un mese buono'' di
lavoro sui trigger (correggendosi da ``due mesi''); le ``due settimane'' sono
forse il risparmio netto stimato, tenendo conto del lavoro di allargamento
degli spazi che è stato comunque necessario. La trascrizione non permette di
chiarirlo.}. Ecco perché è importante analizzare i blocker.

\subsection{Bonus: i feedback costanti}

Cristian nota che, soprattutto ultimamente, in una fase di
\emph{lockdown}\footnote{Il riferimento al lockdown colloca la lezione nel
periodo della pandemia di Covid-19, coerentemente con l'accenno di Simone al
Covid e alle sedi chiuse (``quando si poteva andare'').} in cui tutti sono
costantemente in contatto su Discord, ``anche più di quanto lo saremmo nella
vita vera'', c'è un \textbf{bisogno costante di dare feedback} sul lavoro.

Lui è pro feedback in qualsiasi momento, anzi ``ossessionato'': all'inizio ha
chiesto alla classe di criticarlo e di dargli consigli sulle slide, perché è
l'unico modo per migliorare. Il problema è che \textbf{i feedback costanti
producono un calo di produttività enorme}. L'esempio è personale: Cristian è
anche il VFX artist del team, e sugli elementi estetici c'è sempre modo di
migliorare, quindi riceve di continuo feedback (``io lo farei più grande, rosso,
verde, giallo, a forma di triangolo, a forma di quadrato''). Sono feedback
correttissimi: chi li dà, soprattutto in produzione, non lo fa perché pensa
che tu sia scarso, ma per il bene del gioco, che vive e vuole ben prodotto. I
migliori mandano anche video di riferimento (``ho visto questo VFX in questo
gioco, sapresti rifarlo?''), e questo ben venga. Il problema è che siano
\textbf{costanti}.

Il consiglio: \textbf{limitare la richiesta e l'ascolto dei feedback a fasi
specifiche} (alla fine della sprint, oppure in un controllo intermedio). Dipende
anche da che cosa si sta facendo: con i VFX molte cose dette oggi sono
sistemate il giorno dopo; ma se chi costruisce un livello, in block out o in
\emph{environment}, riceve mille feedback per ogni metro quadro (``un po' più a
destra, un po' più a sinistra''), il lavoro si ferma lì. Sia perché chi lo fa
si \textbf{demoralizza}, sentendosi costantemente criticato, sia per una
questione di fiducia, tornando al discorso di Simone: ``fidati di me come io mi
fido di te. Fammi completare il task; poi, quando ti chiedo il feedback, o più
avanti, è ben accetto, anzi ben venga''.

Cristian si scusa per l'apparente presunzione di definire ``fondamentali'' i
propri consigli: lo fa perché, se li avesse avuti prima, si sarebbe risparmiato
tantissimi problemi. Diego scherza che andranno nelle slide ``con la mia faccia
super inquietante''; Cristian dice che può passare anche le sue slide.

% ---------------------------------------------------------------------
\section{Il calcolo della velocity}
\label{sec:z3-calcolo}

\subsection{Confrontare il fatto con lo stimato}

Sulla parte pratica Simone è stato più esaustivo, quindi Cristian si concentra
sull'interpretazione\footnote{Un breve tratto della spiegazione iniziale è
perso nella trascrizione, che riprende a metà frase (``nel comparare quello
che è stato fatto da quello che è stato stimato'').}. La velocity consiste nel
\textbf{confrontare ciò che è stato fatto con ciò che era stato stimato}, in
numero di task o in pesi, a seconda di come la si gestisce. Da lì si capisce
che cosa è successo e che cosa si può prevedere:
\begin{itemize}
  \item \textbf{fatto $>$ stimato}: si può prevedere di più; nella sprint
        successiva si aumentano i pesi o il numero di task;
  \item \textbf{fatto $<$ stimato}: o si è esagerato con la stima (succede
        sempre, così come succede sempre il caso opposto), oppure ci sono stati
        problemi gravi o \emph{blocker} nella produzione, ed è fondamentale
        capire quale delle due cose sia successa;
  \item \textbf{fatto $=$ stimato}: succede quasi mai, ma è l'obiettivo. Le
        possibilità sono due: si è trovato il proprio \textbf{bilanciamento}, e
        si può continuare così; oppure \textbf{non si sta applicando bene
        Scrum}, ad esempio perché si calcola la velocity su ciò che si è fatto
        invece che sulla stima.
\end{itemize}

\subsection{Più di un numero}

L'obiettivo non è arrivare a una velocity pari a 1, intesa come rapporto tra
fatto e stimato (100 punti stimati, 100 fatti). La velocity è \textbf{molto più
di un numero}: dà almeno tre informazioni.

\paragraph{1. La bontà degli strumenti di stima.} Planning poker, Fibonacci e
tutti gli strumenti di stima sono utilissimi, ma \textbf{vanno testati}: la
velocity mostra se ci sono metodi sbagliati di stima dei pesi o del numero dei
task.

\paragraph{2. La raggiungibilità dello scope.} Se per tutta la produzione la
velocity resta, ad esempio, a 0,8, sempre al di sotto di quanto dovrebbe essere
(in teoria non dovrebbe succedere, perché ci si dovrebbe tarare), e si sono
previsti dieci livelli, si capisce già da subito che \textbf{non ci si
arriverà mai}, a meno di aggiungere uno o due anni di sviluppo.

\paragraph{L'esempio degli otto livelli.} Finita la \emph{vertical slice}, il
team ha ricontrollato tutta la preproduzione e impostato il gioco come un gioco
completo, non come la creazione di una parte: c'è tutta la storia, ci sono
tutti i livelli, tutto ciò che fa il giocatore dall'inizio alla fine, chi
incontra e quando; tutto scritto, ma tarato e lasciato \textbf{flessibile},
``altrimenti ricadiamo nella trappola del Waterfall''.

A settembre, all'inizio dell'anno (``per farvi capire quanto il processo è
rapido''), avevano previsto \textbf{otto livelli}. Sembrava una stima
credibile: il vecchio team, con due level designer che volevano fare
esattamente quel lavoro, produceva un livello in \textbf{due o tre settimane}, a
una qualità abbastanza elevata: magari mancavano l'illuminazione, il
\emph{foliage} e alcuni dettagli, ma il livello era \textbf{giocabile}, che è
la parte fondamentale. Cristian, unico rimasto del vecchio team, si è basato su
quella stima: ``da qualche parte dovevo cominciare''.

Ma i level designer del nuovo team \textbf{non sono level designer puri}: sono
soprattutto game designer, interessati al game design, che hanno deciso di
sperimentare con il level design per necessità (il game design serve sempre, ma
se manca un level designer, visto che il percorso di studi insegna sia level
sia game design, ci si adatta). C'è chi si è adattato a fare il level designer,
e chi ha deciso proprio quell'anno di specializzarsi nel level design: ``perché
avrei dovuto dirgli no, tu non sei un level designer? Anzi, ben venga''.
L'unica cosa è stata prevedere che avesse più bisogno di apprendimento: magari
sbaglierà di più, dovrà informarsi su come si fanno i livelli, avrà bisogno di
più feedback. L'importante è capirlo all'inizio.

Guardando l'andamento del primo livello, Cristian vedeva che \textbf{in un mese
avevano a malapena raggiunto il block out}. Ha convocato una \textbf{riunione
di emergenza}, lunghissima (non una revisione della sprint da venticinque
minuti come quelle di Simone, ma un altro tipo di riunione). Il suo obiettivo
non era tanto convincerli che lo \emph{scope} era troppo alto (ci è riuscito),
quanto \textbf{assicurarsi che non la prendessero come una sconfitta
personale}. È stato difficile ma interessantissimo: non se la sono presa né con
se stessi, per non riuscire a performare, né con lui, che stava di fatto
dicendo ``non ce la farete''. Anzi, nonostante il taglio, sono stati
\textbf{loro stessi a proporre la soluzione}: ``ok, abbiamo capito; facciamo
cinque livelli, e se vediamo che arriviamo (speriamo a fine marzo) con tutti i
livelli, poi possiamo decidere di aumentarli o di cambiare parti del gioco''.
Cristian l'aveva già in mente, ma il fatto che l'abbiano proposta loro l'ha reso
``veramente un successo''.

Tutto questo grazie a un numero. Se Cristian fosse arrivato tra un mese e
mezzo, dopo che avevano già fatto quattro livelli ``con l'acqua alla gola'',
dicendo ``ne dobbiamo fare altri quattro entro marzo'', e poi a fine marzo
``basta, non ce la facciamo'', sarebbe stato molto più grave\footnote{Il
passaggio è confuso nella trascrizione (``mi ha portato ad affrontare queste
situazioni che affrontate a fine marzo'', ``non potne quattro''): il senso
ricostruito è che la velocity ha permesso di anticipare a inizio produzione una
decisione che, presa a fine marzo, sarebbe stata molto più dolorosa.}.
``Ringraziamo in questo caso la velocity.''

\paragraph{3. Lo stato del team.} La velocity permette di capire lo
\textbf{status del team e le sue performance}, perché si basa su una stima fatta
a partire dalla settimana precedente. Si ricollega al ``progettare il \emph{pace}
dello sviluppo'' (sezione~\ref{sec:z3-selezione}): ci si può basare proprio sulla
velocity per capire qual è il ritmo attuale e prevedere i prossimi. ``Tutto
torna.''

\subsection{Bonus: mai tornare indietro}

Un bonus che ``più che un bonus è un'altra lezione'': \textbf{prevedete sempre
task che portano a un incremento del prodotto}. La produzione può rallentare,
perché possono esserci difficoltà, ma \textbf{mai, mai, mai andare indietro}.

È un'altra lezione insegnata da \emph{L'Inf}, legata alla stesura del backlog.
Quando Cristian aveva scritto tutto il backlog (``400 task: come vedete ogni
tanto dico 400, 600, 300, giusto per dire che erano veramente tanti'') e si
accorgeva di aver sbagliato, invece di adattarsi e appoggiarsi alla produzione
in corso \textbf{fermava la produzione} (``ditemi pure che ero folle: lo ero'')
e riscriveva da capo tutti i task, arrivando magari a mille o riducendoli.

Lo stesso vale per i \textbf{documenti}. ``La documentazione è bellissima'':
Cristian era un grande fan della documentazione (Diego può confermarlo), ma
lavorando con Scrum, per ragioni di produttività, ha dovuto \textbf{scendere a
compromessi}: darsi un limite su quanto un documento è bello e ricco di
dettagli. Il motivo è sempre l'esempio della prima meccanica, l'assorbimento e
sparo: tutti i GDD del gioco si basavano su di essa, con tutte le interazioni,
anche secondarie (che cosa si poteva assorbire? Nemici, oggetti, chiavi, cose
per recuperare vita, bonus, malus, qualsiasi cosa). Tagliata la feature, tutti i
GDD sono stati ``letteralmente presi e buttati'', e Cristian ha deciso di
rifarli tutti da capo.

Questo approccio potrebbe funzionare in Waterfall (``finiti i documenti,
iniziamo la produzione''), ma l'obiettivo era \textbf{consegnare un videogioco,
non fare documenti}; e soprattutto il gioco era innovativo. È successo altre due
volte: in tutto Cristian ha riscritto \textbf{backlog e GDD da capo più
volte}\footnote{Nella trascrizione: ``è successo altre due volte quindi per ben
quattro volte io ho riscritto tutto il backlog e tutti i GDD''. Il conteggio non
torna del tutto (una volta più altre due fa tre); forse Cristian conta anche la
prima stesura.}, riportando ogni volta la produzione all'inizio (``ok, possiamo
iniziare il prototipo; no, è sbagliato, ricominciamo''). Lo racconta per far
capire quanto ci ha messo a entrare nel mindset agile, e per tranquillizzare la
classe: non indovinarci subito è normale, ``succederà il 90\% delle volte''.

Quando poi hanno fatto il prototipo, si sono accorti che \textbf{tutta la
documentazione era sbagliata}: con un gioco innovativo non si poteva sapere che
cosa avrebbe funzionato senza prima vederlo, perché non c'erano reference:
nessuno aveva fatto un gioco d'avventura \emph{match-3}. ``Il \emph{match-3}
esiste, i giochi d'avventura esistono, mischiarli insieme è una follia'', e lo
dice ``per criticarmi, non per vantarmi, perché non lo farò più''.

% ---------------------------------------------------------------------
\section{Il tuning della produzione}
\label{sec:z3-tuning}

\subsection{Che cosa e quando}

Il \emph{tuning} della produzione si basa sull'analisi fatta: dove ci si trova,
qual è lo stato delle cose, che cosa si potrebbe fare meglio e stimare
meglio\footnote{Anche qui l'inizio della spiegazione è parzialmente perso nella
trascrizione (``che in parte si basa su ci si trovano, qual è il loro
status'').}; per poi tornare all'inizio del loop con più dati su cui basare una
stima corretta.

Si può fare \textbf{a ogni sprint}, a seconda della durata, oppure \textbf{ogni
serie di sprint}: ``vedetelo voi''. Cristian suggerisce la seconda opzione
perché ritoccare continuamente il backlog per tagliare o aggiungere feature va
contro il principio del lavorare sempre per incrementare il prodotto. Va da sé
che, se si tratta di un taglio molto importante, \textbf{prima lo si prende e
meglio è}: ``ragionateci bene''.

\subsection{Su che cosa basarsi}

\paragraph{Sulle performance, cioè sulla velocity.} Se la velocity è alta, si
può dire: ``a questo punto possiamo reintrodurre qualche feature che era nei
\emph{nice to have}\footnote{Nella trascrizione ``nice to the epoch'': quasi
certamente \emph{nice to have}.} o che era addirittura stata tagliata''. Se è
bassa, si sistemano già milestone e backlog in base al risultato.

\paragraph{Sui feedback.} A Level Up \emph{L'Inf} ha ricevuto tantissimi
feedback, per lo più da giocatori. Trattandosi di una \emph{vertical slice},
spesso dicevano cose che il team già sapeva e che avrebbe implementato nella
sprint successiva (ma, per questioni di tempo, da qualche parte bisogna
tagliare). Altri, invece, hanno parlato di \textbf{accessibilità}.

\paragraph{Il cambio di camera.} A questi si sono aggiunti i feedback costanti
di Diego sulla camera: più che un ``dovete cambiarla'', un ``perché non provate a
cambiare camera?''. Diego interviene: ``c'è voluto del tempo a convincere
Cristian: otto mesi''. Alla fine l'ha fatto e sta funzionando, anche se ha
pesato molto sul team. Il motivo: il \textbf{twin-stick shooter}\footnote{Nella
trascrizione ``twin steep shooter''.} richiede una grande
\textbf{dimestichezza con i controlli} (muoversi e mirare contemporaneamente con
le due levette) e di per sé esclude gran parte del pubblico casual. I giocatori
più casual dicevano proprio questo: ``non riesco a camminare e mirare
contemporaneamente''. Il team se l'è segnato, ne ha discusso e alla fine ha
accettato il cambio di camera, con tutte le conseguenze.

\subsection{Flessibilità anche nella struttura dei documenti}

Un promemoria sulla \textbf{flessibilità}: più si è stati flessibili all'inizio,
più si riesce a rimettere mano a backlog e produzione in seguito. E questo vale
anche, banalmente, per \textbf{come è strutturato un documento}: se ogni casella
fa riferimento a un'altra casella, che fa riferimento a un'altra due righe
sopra, modificare quel documento (o un'implementazione costruita allo stesso
modo) richiederà molto più tempo. Flessibilità in tutto: nella scrittura dei
documenti e nella produzione (``non lo ripeto, altrimenti ricomincio'').

\subsection{Evitare il perfezionismo}

``A proposito di perfetto'': il consiglio sul \emph{tuning} è \textbf{evitare
in tutti i modi il perfezionismo}. Cristian ci sta ancora lavorando
``tantissimo'' (``anche qui Diego lo sa meglio di me''): è migliorato molto e nei
progetti ci riesce, ma in altre cose pretende ancora troppo. Il perfezionismo
\textbf{va contro la flessibilità e crea inevitabilmente dei blocker}. Lo sta
vivendo in prima persona, non per sé ma con altri membri del team.

\paragraph{Il block out con il foliage.} Il task è ``fai il block out di questa
parte del livello'' (dare l'intero block out di livelli da mezz'ora richiede
comunque tempo, quindi il task riguarda la parte iniziale). E nel block out
Cristian si ritrova il \emph{foliage}, il \emph{landscape} scolpito bene, le
rocce. ``Va benissimo perché ti piace il lavoro\dots{} anzi no, sto proprio
dicendo che non va bene.'' L'approccio di voler consegnare il task nel miglior
modo possibile è comprensibile; il problema è che questo perfezionismo
\textbf{blocca tutti gli altri}, per moltissimi motivi. Se si passano due ore a
mettere il foliage e poi, quando si va in gold a rifinire il livello, si trova un
foliage migliore per il \emph{mood} estetico del gioco, quel lavoro va buttato.
Il perfezionismo crea \textbf{buchi di produzione} che alla lunga diventano
letteralmente mesi di lavoro buttato o da buttare: un peccato per chi lo fa e
per chi lavora con lui.

\paragraph{La pressione sugli altri.} Un'altra conseguenza: il perfezionismo
\textbf{mette pressione a tutti gli altri}. Chi lavora alla seconda parte del
block out, con il task di renderla giocabile, vede che la prima parte ha
foliage, rocce e tutti i props (non l'illuminazione, che negli ambienti aperti
è tendenzialmente globale, salvo situazioni specifiche) e si sente in dovere di
raggiungere quel livello di rifinitura\footnote{Nella trascrizione ``si sentirà
in dovere di raggiliere la pozione della vita'': passaggio incomprensibile, il
senso è ricostruito dal contesto.}, trascurando la parte di gameplay, che invece
nella prima parte era stata fatta. Il perfezionismo di uno ha generato
\textbf{ansia} nell'altro, ed è da evitare in ogni modo.

\paragraph{I criteri di accettazione.} È un altro motivo per cui sprint e task
devono avere un \textbf{goal} e un \textbf{criterio di accettabilità}, per
usare i termini di Scrum: il task si considera fatto quando è \emph{done}, e va
bene così, perché lo si è valutato prima. Tutto ciò che si fa oltre il
\emph{done} è \textbf{tempo perso}. Può andare bene che quello che fai adesso
non dovrai farlo dopo, ma è sbagliato come approccio: se il lavoro è stato fatto
bene prima, si può andare avanti tranquilli, e al passo successivo ci si penserà
dopo.

% ---------------------------------------------------------------------
\section{Domanda: come comunicare un taglio al team}
\label{sec:z3-qa-taglio}

Prima della seconda parte, che Cristian promette più breve, una domanda: sul
piano relazionale ed emotivo, come ha gestito il dover convincere, o comunque
spiegare, al team la necessità del taglio da otto a cinque livelli?

\paragraph{Ricordare come funzionano i feedback.} ``Ottima domanda.''
Emotivamente è stato pesante, perché non tutti reagiscono ai feedback allo
stesso modo. Per questo, a ogni riunione, Cristian ricorda i principi dei
feedback (``so che sembra banale, so che non serve''), in particolare una
clausola che tutti dimenticano: \textbf{il feedback si può accettare o
rifiutare}. Se non si è d'accordo, lo si può ignorare; invece tutti partono dal
presupposto che un feedback vada seguito al 100\%\footnote{Nella prima lezione
zero il relatore invitava ad accettare il feedback ``a prescindere''
(sezione~\ref{sec:z1-ricevere}); lì però ``accettare'' significava ascoltarlo
senza difendersi, non necessariamente applicarlo. Le due indicazioni sono
quindi compatibili: si ascolta tutto, poi si decide che cosa farne.}.

\paragraph{Partire dai dati.} Sembra brutto affrontare una questione emotiva
partendo dai dati, ma l'obiettivo era \textbf{far arrivare loro alla soluzione},
non dargliela. Ed è andata così: sono stati loro a dire ``per adesso facciamo
cinque livelli, e poi, \textbf{se ci riusciamo}, implementiamo gli altri''.
Quel ``se ci riusciamo'' era ciò che Cristian voleva sentire: se fosse rimasto
solo ``ne facciamo cinque'', l'avrebbero vissuto come un limite su se stessi
(``non ho fatto in tempo a farne otto''). Il ``se ci riusciamo'' gli ha
comunicato che avevano capito che \textbf{non era una critica al loro modo di
lavorare}, né un modo per dire ``non siete capaci'': lo aveva sottolineato più
volte, magari non con parole così dirette, ma ``un po' \emph{soft-skillato},
come diciamo noi designer''.

Hanno colto che non c'erano alternative (``otto non li facciamo, punto''), ma
anche le possibilità: facciamo cinque livelli, abbiamo altro tempo,
perfezioniamo ciò che abbiamo fatto; vogliamo aggiungere una nuova modalità di
gioco? Benissimo. Altri due livelli? Vediamo, in base al tempo. Esserci
riuscito lo ha reso ``super fiero''.

% ---------------------------------------------------------------------
\section{Il team management}
\label{sec:z3-team}

\subsection{Senza team non esiste Scrum}

Cristian accelera, per stare nei tempi. Il punto da capire è che
\textbf{Scrum si adatta al team}: quindi il team è la cosa più importante di
tutte, più di Scrum stesso. ``Se io mi adatto al team, il team deve esserci.
Senza team non esiste Scrum.''

\subsection{Una produzione senza stress}

Di conseguenza, una produzione in Scrum \textbf{dovrebbe rendere il team
rilassato}. Come il team di Simone, anche il primo team di Cristian (quello fino
alla \emph{vertical slice}) \textbf{non ha mai fatto crunch}. Il nuovo team sì,
per i grossi cambiamenti già raccontati: la camera è cambiata completamente, e
``cambiare la camera in un gioco è fare un altro gioco''. Cristian aveva già
spostato la scadenza (``invece di finire a marzo finiamo ad aprile''), ma il team
si è sentito comunque in dovere di risolvere prima lo scompenso creato da quei
cambiamenti.

In generale ci si dovrebbe sentire meno stressati per un motivo semplice: chi
guida sta dando i task che lo sviluppatore ha chiesto, verbalmente o dimostrando
di saperli fare. ``So che in due settimane devo fare cinque cose, faccio queste
cinque cose'': non c'è la pressione di dover fare il più possibile. Se poi uno è
perfezionista e passa da cinque cose a quindici perché ``ha tempo'', il problema
(``sembra brutto dirlo'') è suo, nasce nel singolo individuo; e chi guida deve
basarsi sul team, sulle performance del team.

\subsection{Mai la velocity sui singoli}

Da qui un consiglio sulla velocity: \textbf{non calcolate mai, mai, mai la
velocity delle singole persone}, e nemmeno i task fatti dai singoli. Lo si può
fare per sé, ma tenendolo ``chiuso in un'agenda scritta a penna'': se un membro
del team scopre che si stanno valutando le sue prestazioni individuali,
\textbf{si sente giudicato} e soprattutto non sa più come dire che non vuole
fare una cosa.

A Cristian è successo: ``mi hanno sgamato''. Si è dovuto scusare, anche se non
erano scuse: ha detto la verità, cioè che lo faceva più per capire lui la
produzione che per valutare le loro performance. Ma spiegarlo è complicato,
quindi è meglio non farlo: tanto la velocity si basa sulle performance del team,
e i dati individuali non servono. Se li si vuole per sé, per capire a chi dare
un task piuttosto che a un altro (se una persona in un'ora fa dieci
implementazioni e un'altra una, ci si regola quando si selezionano i task),
non bisogna assolutamente farli ricadere sulle persone.

\subsection{Primo aspetto: adattare Scrum al team}

Per una produzione senza stress servono due aspetti chiave. Il primo è che
\textbf{Scrum va adattato al team}: si prende la base di Scrum (quella di Diego,
dei libri, di Simone e di Cristian), la si \textbf{testa}, si \textbf{analizzano}
la velocity e tutto il resto, la si \textbf{tara} (``\emph{daily scrum}? No,
facciamolo ogni tre giorni, ogni settimana, o non facciamolo''), la si
\textbf{riprova}, e si ripete il ciclo per tutta la produzione. Alcuni esempi da
\emph{L'Inf}:
\begin{itemize}
  \item i \textbf{daily scrum} li fanno \textbf{ogni tre giorni}, ``se va
        bene'', e dipende anche dal tipo di produzione;
  \item la \textbf{lunghezza delle sprint} è flessibile (una, due o tre
        settimane) e dipende:
        \begin{itemize}
          \item dalla \textbf{disponibilità del team}: se tutto il team non
                c'è, è inutile fare una sprint di una settimana;
          \item soprattutto dalla \textbf{fase di produzione}: quando si
                implementano script e in una settimana si riescono a
                implementare tantissime meccaniche, la sprint dura una
                settimana, al termine della quale si controlla ciò che è stato
                fatto e si sceglie che cos'altro implementare; quando si passa al
                level design, con lavori molto dispendiosi in termini di tempo,
                le sprint si diluiscono, e infatti in quel momento sono a sprint
                di due settimane, concentrati sul completamento dei
                livelli\footnote{Nella trascrizione ``sull'ultima mento,
                sull'ultima azione, non mi ricordo, scusate. Per completare il
                livello'': forse ``l'ultima milestone''.};
        \end{itemize}
  \item la velocity si calcola \textbf{sui task e sul loro peso, non sulle ore
        lavorative}, come diceva Simone. Soprattutto per degli studenti, che
        non hanno ore fisse. In un ambiente di lavoro in cui tutti sono
        disponibili dalle 9 alle 18 potrebbe cominciare ad avere senso:
        ``secondo me no, ma ve lo dico per correttezza''.
\end{itemize}

\subsection{Secondo aspetto: tutti devono conoscere il mindset}

Il secondo aspetto è che \textbf{tutti i membri del team devono conoscere il
mindset di Scrum}\footnote{A questo punto Cristian chiede a Diego cinque minuti
in più per chiudere bene, e Diego acconsente.}. Per due motivi:
\begin{itemize}
  \item se tutti lo conoscono, il rischio di cadere nel Waterfall o nel
        perfezionismo si riduce, perché ognuno sa che cosa è giusto e che cosa
        è sbagliato;
  \item tutto il team ha \textbf{lo stesso goal}: se la sprint chiede il block
        out, tutti sanno che, secondo la mentalità agile, va bene farlo ma non
        deve essere perfetto. Bisogna fare i \emph{done}, non i
        \emph{perfect}.
\end{itemize}
Come far capire a tutti la mentalità di Scrum? Cristian usa due modi.
\begin{enumerate}
  \item \textbf{Parlarne costantemente}, a ogni riunione, come sta facendo con
        la classe: sia per migliorare lui stesso e farsi entrare il mindset,
        sia per trasmetterlo agli altri.
  \item Il più efficace: \textbf{lo shipping}. Il team riesce a consegnare
        quasi costantemente. Ci sono momenti in cui un problema cade
        all'ultimo, come in quella settimana, con alcuni strumenti
        particolari che li hanno rallentati, ma si risolve tutto. Dimostrando
        di riuscire a \textbf{mantenere le milestone}, il team ha acquisito
        fiducia in lui e ha capito quanto è potente Agile. Nella
        \emph{vertical slice} consegnavano le build addirittura prima di
        quanto richiesto; e la consegna prima di Natale, nota da settembre, è
        stata rispettata nonostante il cambio di camera e tutto il resto. Il
        messaggio: \textbf{non è importante che ciò che si consegna sia
        perfetto, ma che rispetti le condizioni di accettabilità} che ci si è
        dati.
\end{enumerate}
Sono i due metodi che funzionano meglio finora; se ne troverà altri, li
condividerà: ``tutto è in continua evoluzione''.

\subsection{Il post-mortem}

Il \emph{post-mortem} è \textbf{la garanzia di capire se ciò che si sta facendo
è corretto}: serve a vedere come stanno andando il team e la produzione, se il
mindset di Scrum è stato compreso o no, e ad agire di conseguenza.

\textbf{Quando farlo} lo si decide liberamente: a ogni sprint, a ogni gruppo di
sprint, a ogni milestone.

Ciò che Cristian sottolinea è che \textbf{nel post-mortem devono parlare tutti},
anche lo sviluppatore che non vuole mai parlare: bisogna dargli il tempo per
farlo. Un modo è \textbf{organizzarlo in anticipo}: se è tra una settimana, lo si
comunica la settimana prima. Cristian prepara sempre un \textbf{documento della
riunione} prima e lo aggiorna dopo, con una sorta di elenco: sembra burocratico,
ma aiuta a mantenere il focus. Bisogna far capire a ciascuno che tra una
settimana parlerà, e non può sottrarsi (se si assenta ci perde lui). E bisogna
dirgli di \textbf{prepararsi} qualcosa da dire, magari scrivendolo o facendo una
piccola revisione: se ti viene in mente una cosa da dire la settimana dopo e non
la scrivi, finirai per pensare ``sai che c'è, non la dico''; se l'hai scritta,
non puoi dimenticartela, e hai qualcosa di cui parlare. È fondamentale perché
bisogna tenere conto di \textbf{tutti i tipi di personalità}: non solo i
giocatori hanno personalità diverse, anche gli sviluppatori hanno personalità
con bisogni diversi.

\subsection{La sincerità}

Cristian propone di tagliare la parte ``sdolcinata'', ma Diego vuole sentirla
per intero (``non si sfugge al miele''; ``poi mi manderete gli IBAN e vi pagherò
il tempo perso''). Come Diego sa, Cristian cerca sempre di metterci ``qualcosa di
filosofico''. Per \textbf{sincerità} intende che, in qualsiasi momento della
produzione, \textbf{tutti i membri del team devono sentirsi protetti in ciò che
dicono} e liberi di esprimere opinioni e feedback. È importantissima sia nella
produzione sia nel team.

\paragraph{Nella produzione.} Senza sincerità la produzione va a picco; si può
comunque consegnare, non è obbligatoria, ma con la sincerità si realizza un
prodotto migliore. Se a inizio produzione si pensa che una meccanica ``faccia
schifo'' (prende un \emph{player type} completamente fuori target e richiede
troppo lavoro per essere implementata), \textbf{è quello il momento di dirlo}.
Altrimenti ``il karma vuole che tu stesso che pensi questa cosa sarai colui che la
andrà a implementare'', e la farai male, perché non vuoi farla. Se non l'hai
detto, o l'hai detto troppo tardi, ``ti mangerai le mani''.

E il \textbf{``ve l'avevo detto''} va risparmiato: anche se il giorno dopo arriva
la conferma di ciò che si era detto, anche se è umano, va evitato, perché crea
attriti nel team che non si risolvono più (Cristian parla per esperienza
diretta). È lo stesso spirito del \emph{disagree and commit}
(sezione~\ref{sec:z2-disagree}).

\paragraph{In Scrum.} Se uno sviluppatore dice ``questa settimana so fare 15
task'' e ne fa 7, la prima volta passa. La settimana dopo dice di nuovo ``ne
faccio 15, tranquillo'', magari inventando scuse (``quel problema non ce l'ho
più''). Una volta, due volte; \textbf{alla terza bisogna chiedere sincerità}: qual
è il problema? Non hai tempo?

Perché sul tempo, ``ragazzi'', \textbf{è importante dire quanto tempo si ha, non
garantirlo}. Se il team (cinque persone, nel caso di Cristian) dice di avere due
ore al giorno, lui pianifica una produzione da dieci ore al giorno; se il team
dice di avere otto ore al giorno, pianifica di conseguenza (``80
ore''\footnote{Così nella trascrizione; con cinque persone a otto ore al giorno
il totale sarebbe 40 ore, oppure 80 su due giorni. Il senso dell'esempio non
cambia.}), e se poi in realtà le ore sono dieci, gli diranno ``ci dai troppi
task''.

Bisogna essere sinceri \textbf{prima di tutto con se stessi}: se un task richiede
due giorni, si dice due giorni; magari tre, perché un ``tempo cuscinetto'' ben
venga, ma lo si dice per com'è. E non bisogna proteggersi dal timore che ``se
dico la verità nessuno si fiderà più di me'' o ``mi daranno più task'': è compito
dello scrum master, ma di tutto il team, far capire che \textbf{si è tutti sulla
stessa barca}.

\paragraph{Nel team.} Gli attriti generati dalla mancanza di sincerità si
trasformano in \textbf{relazioni migliori} quando c'è sincerità. Le
incomprensioni vanno \textbf{risolte subito}, affrontandole; magari porteranno a
un litigio, se proprio quello è il finale, ma tra persone mature si dovrebbe
riuscire a risolvere in modo pacifico. Per questo, sulla slide, la sincerità è
definita \textbf{``rischiosa ma necessaria''}.

% ---------------------------------------------------------------------
\section{Chiusura della lezione}
\label{sec:z3-chiusura}

Cristian salta la slide di riepilogo (la lezione si potrà rivedere online) e
lascia i suoi contatti: \textbf{Telegram}, dove molti studenti lo trovano già nei
gruppi (in particolare quelli dell'accademia online), oppure l'\textbf{email},
che Diego può fornire. Si rende disponibile per qualsiasi cosa, sulla lezione o
su un'idea da sottoporgli (``ho pensato così, secondo te può funzionare?''),
precisando di non volersi sostituire in alcun modo a Diego. Si scusa per i dieci
minuti finali ``di corsa''.

Diego lo ringrazia per la disponibilità, ricordando che anche lui, da studente,
ha un carico di lavoro notevole. L'obiettivo della lezione era far vivere alla
classe un'esperienza e dare informazioni simili a quelle di chi ha fatto lo
stesso percorso, ma ha \textbf{internalizzato Agile} meglio di altri: tutti hanno
avuto le stesse lezioni, ma Cristian e Simone sono quelli che le hanno capite
meglio. Ci sono molti \emph{insight} utili in entrambi gli interventi, e i punti
chiave di Agile sono stati ripetuti più volte, ``perché alla fine Agile è
comprendere i punti chiave''. ``Grazie ancora, e buon pranzo.''\footnote{La
trascrizione si chiude con la parola ``grazie'' ripetuta diverse centinaia di
volte: è un artefatto tipico della trascrizione automatica su un tratto di
audio silenzioso o di saluti sovrapposti, non un contenuto della lezione.}

% ---------------------------------------------------------------------
\section*{Punti chiave del capitolo}
\addcontentsline{toc}{section}{Punti chiave del capitolo}

\begin{itemize}
  \item \textbf{Agile è un mindset, non un framework}: non ha regole, ha una
        logica da \textbf{internalizzare}, come una lingua che a un certo punto
        si smette di tradurre. I quattro valori non vanno seguiti alla lettera:
        vanno adattati al team e al progetto, e non è detto che Agile serva.
  \item I \emph{framework} (Scrum, Kanban, \dots) hanno regole e si basano
        sull'\textbf{inspect and adapt}, cioè sulle iterazioni. Si parte
        rispettandoli al 100\%, si capisce che cosa non funziona e perché, e
        \textbf{ci si costruisce il proprio}.
  \item \textbf{Meno riunioni possibili}: unire planning e post-mortem, stare in
        25--30 minuti, report da un minuto. \textbf{Nelle riunioni non si
        risolvono i problemi}: si fissa un appuntamento con chi può aiutare. Un
        \textbf{timekeeper} può far rispettare i tempi.
  \item Le \textbf{sprint} (una o due settimane) sono l'unità di tempo: le
        milestone si misurano in sprint, non in mesi.
  \item \textbf{Consistency}: una cerimonia saltata fa crollare le altre a
        effetto domino. Per questo ciò che non funziona per il team (spesso il
        \emph{daily scrum} o il \emph{planning poker}) va \textbf{tagliato},
        non trascinato.
  \item Il \textbf{backlog} si fa \textbf{a scaglioni}, per milestone, partendo
        dalle \emph{epic} e dai documenti (visione, GDD, \emph{versioning} con
        tutti gli ambiti, LDD). Si organizza per feature, per ambito con un
        \textbf{color code} condiviso, per milestone e sotto-milestone o per
        priorità, tutto in un solo strumento.
  \item Ogni task (e ogni sprint) ha un \textbf{obiettivo verificabile} e un
        \textbf{criterio di accettazione}: ``completare il livello 5'' non lo è.
        Le \emph{user story} aiutano, ma non sono obbligatorie.
  \item \textbf{Mai stimare in ore}: si stima in modo \textbf{relativo}, con la
        sequenza di \textbf{Fibonacci}, partendo dagli estremi. I valori sono
        arbitrari ma \textbf{consistenti}, e quindi utilizzabili.
  \item La \textbf{velocity} (somma dei pesi dei task \emph{done} per sprint,
        confrontata con lo stimato) \textbf{prevede il futuro}: dice se gli
        strumenti di stima funzionano, se lo \emph{scope} è raggiungibile e come
        sta il team. Permette di tagliare per tempo e di arrivare in fondo
        \textbf{senza crunch}. L'andamento a dente di sega è normale.
  \item Le prime sprint \textbf{si sbagliano tutte}, ed è giusto: la quantità si
        tara sulle disponibilità e sulle capacità attuali, prevedendo i
        \textbf{tempi di apprendimento}.
  \item Si può \textbf{progettare il ritmo dello sviluppo}: dopo periodi di
        lavoro invisibile, sprint che danno soddisfazione (nuovi nemici, VFX,
        \emph{juiciness}) per il morale del team.
  \item I \textbf{blocker} vanno trovati subito e \textbf{analizzati alla
        radice} prima di essere risolti: la soluzione più rapida (i muri
        invisibili) può costare settimane.
  \item I \textbf{feedback costanti} abbattono la produttività: vanno raccolti in
        momenti specifici. Il feedback si può accettare o rifiutare.
  \item \textbf{Flessibilità} in tutto, soprattutto negli elementi innovativi:
        non prevedere tutto nel dettaglio (è Waterfall), non fermare mai la
        produzione per riscrivere backlog e documenti, lavorare sempre per
        \textbf{incrementare il prodotto}.
  \item Il \textbf{perfezionismo} crea blocker, lavoro buttato e ansia negli
        altri: si fanno i \emph{done}, non i \emph{perfect}.
  \item \textbf{Il team viene prima di Scrum}. La velocity non si calcola mai sui
        singoli. Tutti devono conoscere il mindset: parlarne a ogni riunione e,
        soprattutto, \textbf{consegnare} rispettando le milestone.
  \item Nel \textbf{post-mortem} parlano tutti: va annunciato prima, con un
        documento, chiedendo a ciascuno di prepararsi.
  \item Alla base di tutto ci sono la \textbf{fiducia} reciproca e la
        \textbf{sincerità}: dire subito ciò che non va, dire quanto tempo si ha
        (non garantirlo), evitare il ``ve l'avevo detto''. La sincerità è
        \textbf{rischiosa ma necessaria}.
\end{itemize}

%    \part{Il corso di Game Design}
%    \renewcommand{\thechapter}{\arabic{chapter}}
%    \setcounter{chapter}{0}
%    % =====================================================================
%  Capitolo 1 — Introduzione al corso e al mestiere del game designer
%  Fonte: trascrizione della lezione GD1 (Game Design, primo anno)
%  Pacchetti richiesti nel preambolo del libro:
%    \usepackage[italian]{babel}, \usepackage[T1]{fontenc},
%    \usepackage[utf8]{inputenc} (solo pdfLaTeX), \usepackage{booktabs}
% =====================================================================

\chapter[Il mestiere del game designer]{Il mestiere del game designer: introduzione al corso}
\label{cap:introduzione-game-design}

Questo primo capitolo fa da soglia. Prima di affrontare metodi, framework e
documenti, è necessario capire \emph{in che modo} si impara il game design,
\emph{con chi} lo si fa, \emph{con quali strumenti} e, soprattutto,
\emph{che cosa fa concretamente} un game designer nel suo lavoro quotidiano.
Il capitolo segue l'impostazione della lezione inaugurale del corso: si parte
dal contesto formativo (il corso come simulazione di lavoro), si passa al
metodo di apprendimento (l'\emph{inquiry-based learning}), si elencano
letture e risorse di riferimento, e si arriva infine a una prima definizione
operativa del ruolo del designer, dei limiti entro cui lavora e delle sue
possibili specializzazioni.

Una premessa terminologica: il game design è una disciplina giovane e quasi
tutta la sua letteratura è scritta in inglese. Per questo motivo molti termini
tecnici verranno lasciati in lingua originale (\emph{pipeline},
\emph{pitching}, \emph{tuning}, \emph{feature team}, \dots): sono i termini che
si useranno nel lavoro, ed è importante imparare a riconoscerli e a usarli
correttamente fin da subito.

% ---------------------------------------------------------------------
\section{Il contesto: un corso come simulazione di lavoro}
\label{sec:simulazione-lavoro}

\subsection{Il principio fondamentale}

Il primo concetto da fissare è che il corso \textbf{si presenta come una
simulazione del lavoro}. Dal momento in cui si entra in aula si viene trattati
come \emph{potenziali professionisti dell'industria}: non soltanto studenti,
ma lavoratori effettivi di uno studio (l'accademia), con una gerarchia, dei
referenti, dei compiti e delle scadenze.

Ne derivano alcune conseguenze pratiche:
\begin{itemize}
  \item il docente non assegna un ``compitino'' da risolvere per accumulare
        punti o presenze: ogni \emph{task} è pensato per far dedicare tempo
        all'apprendimento di qualcosa di concreto;
  \item non consegnare un task non danneggia il docente (che ``vive ancora
        domani e viene pagato ugualmente''), ma lo studente, che perde
        l'occasione di mettere le mani su un argomento e diventare più abile;
  \item docenti e studenti devono trattarsi \textbf{da colleghi}, perché è
        molto probabile che un domani lo diventino davvero.
\end{itemize}

\subsection{Perché l'impegno conta più che mai}

Il lavoro del game designer è descritto come \emph{molto romantico, ma anche
molto duro}, e sempre più duro col passare degli anni. Durante il periodo del
Covid l'industria ha visto grandi investimenti; quei ``rubinetti'' sono stati
chiusi da tempo, il rapporto con i publisher è diventato più complesso e le
opportunità lavorative non sono più quelle di tre o quattro anni fa. Questo
significa:
\begin{itemize}
  \item più competizione per ogni posizione;
  \item aspettative più alte: ci si aspetta che i candidati siano ancora più
        preparati e sicuri delle proprie competenze;
  \item un unico modo per diventarlo: \textbf{lavorare}. Più tempo si dedica
        oggi ad apprendere il mestiere, più facilmente si sarà assumibili
        domani e più agevole sarà il lavoro di tutti i giorni.
\end{itemize}

\subsection{Sbagliare qui per non sbagliare fuori}

Il mantra del corso è: \textbf{quanto più provate qui, anche in modo goffo e
sbagliando, tanto meno sbaglierete là fuori}. In accademia, per quanto la
simulazione sia realistica, nessuno toglie lo stipendio e nessuno licenzia da
un momento all'altro. Nell'industria sì, e un errore può costare molto denaro
all'azienda, ancora di più se l'azienda è la propria.

\subsection{Il valore del fallimento}

Il fallimento, se affrontato con la mentalità corretta, è probabilmente
\emph{la cosa più vantaggiosa} che si possa fare durante la formazione:
\begin{itemize}
  \item un successo (``avete superato questa deadline'') conferma le proprie
        abilità, ma insegna relativamente poco;
  \item dopo $n$ fallimenti si impara moltissimo, e i successi arrivano;
  \item è fondamentale non demoralizzarsi: il feedback negativo non è mai un
        giudizio sulla persona, ma su ciò che è stato consegnato.
\end{itemize}

Il docente riconosce apertamente che dedicare tempo ed energia a un progetto e
vederlo rifiutato, o vederlo non uscire mai, è demotivante: è successo a lui e
lo ha visto accadere anche su progetti molto grandi, con team numerosi. Il
punto però non è ciò che si è ottenuto materialmente, ma \textbf{quanto si è
cresciuti}. Da ogni esperienza ci si portano a casa due cose:
\begin{enumerate}
  \item \textbf{l'esperienza}: quanto si è diventati bravi, veloci,
        praticamente abili nell'applicare un framework, quanto si è cresciuti
        come designer;
  \item \textbf{il rapporto umano}: nel lavoro si ``suda sangue'' fianco a
        fianco con altre persone, e questo crea legami molto forti che
        portano avanti l'industria nel senso più positivo del termine: nascono
        progetti privati, ci si porta da un'azienda all'altra, si formano team
        con basi solide.
\end{enumerate}

Va evitato anche l'eccesso opposto: la cultura del \emph{grinding} si è
appropriata di questi discorsi in modo distorto. Non bisogna diventare robot.

% ---------------------------------------------------------------------
\section{Organizzazione del corso e gerarchia}
\label{sec:gerarchia-corso}

\subsection{La struttura gerarchica}

La struttura dei referenti del corso ricalca volutamente una
\textbf{gerarchia tipica dell'industria}:
\begin{itemize}
  \item al vertice c'è il \textbf{direttore del corso} (Diego Ricciuti), che
        corrisponde al \emph{creative director}: supervisiona la visione del
        progetto (il concetto di ``visione'' verrà approfondito più avanti);
  \item sotto di lui ci sono i \textbf{referenti di dipartimento}, uno per
        disciplina, con un corrispettivo per ciascuna sede (Milano e Roma);
  \item in fondo ci sono gli studenti, che occupano il primo gradino di
        un'azienda: l'\emph{intern} (stagista) o il \emph{junior}, il passo
        che precede il lavoratore ``regolare''.
\end{itemize}

\begin{table}[htbp]
  \centering
  \begin{tabular}{@{}p{0.34\textwidth}p{0.29\textwidth}p{0.29\textwidth}@{}}
    \toprule
    \textbf{Disciplina} & \textbf{Milano} & \textbf{Roma} \\
    \midrule
    Direzione (\emph{creative director}) & \multicolumn{2}{p{0.58\textwidth}}{Diego Ricciuti} \\
    Game design      & Mattia Podini        & Giacomo ``Jack'' Danna \\
    Level design     & Stefano Zocchi       & Andrea Dispoli \\
    Narrative design & \multicolumn{2}{p{0.58\textwidth}}{Daniele Fusetto (lezioni condivise), Giacomo Masi} \\
    Programmazione (Unreal) & Alessandro Pregnolato & Matteo Velardi \\
    \bottomrule
  \end{tabular}
  \caption{Referenti del corso per disciplina e sede (nomi come riportati in lezione).}
  \label{tab:referenti}
\end{table}

Per il primo anno, il referente diretto degli studenti di game design è il
docente di game design: è a lui che ci si rivolge per dubbi e domande. Altre
figure utili sono il direttore didattico (severo, ``giustamente severo'', da
contattare solo per questioni particolarmente importanti), l'amministrazione
(contratti e documentazione) e la segreteria.

\subsection{I canali di comunicazione}

\begin{description}
  \item[Discord] è il canale per le comunicazioni ufficiali del corso e per
        le domande che non richiedono urgenza immediata; può avere qualche
        ritardo.
  \item[E-mail] è il canale per tutto ciò che è formale o ``ancora più
        ufficiale'' e per le comunicazioni rapide e immediate.
  \item[Telegram] esistono residui di vecchi canali precedenti a Discord, da
        evitare.
\end{description}

Le domande che valgono per tutti vanno fatte in classe (anche da un portavoce
a nome del gruppo); quelle singole e personali via Discord o, se formali, via
e-mail. Se il docente non risponde a una domanda c'è un motivo, ad esempio
ragioni legali, e lo dirà esplicitamente.

\subsection{I tre moduli e la settimana tipo}

Le tre discipline principali del primo anno (e non solo) sono:
\begin{itemize}
  \item \textbf{Game Design} (tendenzialmente il lunedì);
  \item \textbf{Level Design} (tendenzialmente il mercoledì);
  \item \textbf{Programmazione} (tendenzialmente il venerdì).
\end{itemize}
Nel primo mese le lezioni sono tutte di game design (lunedì, mercoledì e
venerdì). Sono previste anche lezioni di narrativa, ma il piano formativo si
aggiorna ogni anno: le ore di narrativa sono state ridotte a favore di quelle
di programmazione. Le poche lezioni di narrativa daranno quindi dei
\emph{cenni}, non la preparazione di un narrative designer. Ciò che serve
davvero è imparare \textbf{come comunicano gli altri dipartimenti} e come
lavorare con persone che fanno un lavoro diverso ma complementare al proprio.

\subsection{Le due anime delle lezioni di game design}

\begin{description}
  \item[Design methods (metodologie)] lezioni prevalentemente teoriche, ma con
        una buona componente pratica, in cui si studiano metodologie di design
        e i cosiddetti \emph{framework} dell'industria.
  \item[Documentazione e pitching] il cuore del lavoro quotidiano del
        designer.
\end{description}

\paragraph{La documentazione.}
Il docente si definisce ``il più grande nerd della documentazione'' che si
possa incontrare. È una parte dura, perché a pochi piace documentare, ma il
lavoro del designer \emph{è principalmente questo}: scrivere documenti e, da
quei documenti, passare alla pratica e all'implementazione. Chi non sa
documentare bene parte male. Quando non si è presenti, il documento è il
proprio \textbf{biglietto da visita}: qualcun altro lo leggerà, quindi va
scritto nel modo più chiaro e comprensibile possibile.

Un esempio concreto del perché la documentazione conta: CD Projekt Red ha
sviluppato un proprio engine proprietario, ma buona parte della sua
documentazione è stata ``buttata all'acqua''. Chi deve usarlo o lo sa già
usare, o ha competenze di un altro engine da trasferire, oppure è in
difficoltà. Si documenta \textbf{perché altri possano leggere}.

\paragraph{Il pitching.}
Il \emph{pitching} è l'esposizione, la presentazione orale. La parte creativa
del team presenta spesso:
\begin{itemize}
  \item le proprie idee ai colleghi;
  \item il gioco a convention ed eventi;
  \item il pitch del gioco agli investitori, per ottenere i fondi necessari a
        realizzarlo.
\end{itemize}
Saper parlare e comunicare chiaramente le idee è fondamentale. In sintesi:
\emph{il game designer parla}, e bisogna farsene una ragione.

\subsection{Altre componenti del corso}

\begin{description}
  \item[Gioco critico e playtest] la classe viene divisa in due gruppi: uno
        gioca, l'altro osserva. Chi osserva non applaude ``come le
        scimmiette'': annota, segna e analizza il \emph{playtest}, esattamente
        come si fa nell'industria. Il playtest fa parte delle mansioni del
        designer, ma esiste anche come specializzazione autonoma.
  \item[Pipeline] termine che indica tutto ciò che riguarda i processi
        formalizzati dell'industria; di solito si intende la \emph{pipeline di
        produzione}: come si divide il team, qual è la sua gerarchia, come il
        \emph{producer} gestisce le deadline. In breve: \emph{come} e
        \emph{quando} si fanno le cose. Al primo anno è meno rilevante (si
        dovrà ``nuotare'' un po' da soli); dal secondo anno si vedranno le
        metodologie effettivamente applicate ai progetti digitali. Si
        scoprirà quanto sia difficile lavorare a un progetto digitale, e
        ancor più fisico, senza un piano di sviluppo comprovato, come quelli
        usati per qualsiasi software in qualsiasi parte del mondo.
  \item[Engine] il motore di riferimento è \textbf{Unreal Engine}
        (sezione~\ref{sec:engine}).
  \item[Programmazione] principalmente tramite \emph{Blueprint} in Unreal,
        più eventuali sistemi di versionamento; il corso parte da zero.
  \item[Progetti pratici] probabilmente la cosa più importante. Al primo
        anno i progetti di gruppo restano circoscritti al corso di design,
        perché mancano ancora alcune nozioni (come comunicare, affrontare
        problemi, dare e ricevere feedback). In seguito si allargheranno a
        progetti digitali, molto più difficili di quelli fisici, con membri
        di altri team: un rapporto non unilaterale, ma ``covalente''.
\end{description}

% ---------------------------------------------------------------------
\section{Il metodo: l'inquiry-based learning}
\label{sec:inquiry}

\subsection{Che cos'è}

Il corso è costruito secondo l'\textbf{inquiry-based learning}
(apprendimento basato sull'indagine), un approccio pedagogico nato intorno
agli anni Sessanta, che si richiama a pensatori come Dewey e Piaget. In
generale, le discipline da cui il design attinge sono quasi sempre la
sociologia, la psicologia e, talvolta, l'economia.

Il funzionamento è un ciclo:
\begin{enumerate}
  \item il docente fornisce metodologie e framework dell'industria durante
        la lezione;
  \item lo studente, fuori dalle ore di lezione, lavora sui \emph{task}:
        mette a terra le nozioni, le maneggia, le assorbe meglio;
  \item da questo lavoro nascono nuovi dubbi;
  \item i dubbi tornano in aula: lo studente chiede, il docente risponde
        (ammesso che dare la risposta sia utile, e ammesso che una risposta
        esista).
\end{enumerate}

\subsection{Non sempre esiste una risposta giusta}

Non tutto ciò che si studia ha ``la risposta migliore possibile in assoluto
nel vuoto''. Esistono metodologie, limiti e framework che aiutano il lavoro,
ma non sono fissi: cambiano da azienda ad azienda, da lavoro a lavoro, da
progetto a progetto, e a volte vengono stravolti. Sta alla
\textbf{sensibilità del designer} capire quale strumento o framework usare in
un dato momento, e questa sensibilità si impara solo con l'esperienza. Nel
lavoro nessuno dirà ``usa questa cosa'' o ``è meglio quest'altra'': sarà il
designer a scegliere.

Anche l'esperienza del docente ha confini: sei-sette anni di lavoro, un modo
personale di approcciare il design e un certo insieme di progetti (ad esempio
nessun MMO). Oltre quei confini si cercano risposte nella ricerca, studiando
come aziende con esperienza specifica (ad esempio Riot) hanno affrontato
problemi simili e hanno condiviso le loro soluzioni, sempre valutando quali
metodi siano adatti alla dimensione del team, al tipo di progetto e al tempo
disponibile.

\subsection{Sensibilità, non memoria}

Non interessa imparare a memoria un framework o la teoria di un autore. Certe
terminologie vanno conosciute, ma per ragioni di comunicazione e chiarezza.
Ciò che conta è sviluppare la \textbf{capacità di farsi domande}, di farle
agli altri, di fare ricerca e di trovare una propria risposta. È esattamente
ciò che si fa nel lavoro: davanti a qualcosa che non si sa fare, l'abilità
del designer sta nel cercare come qualcun altro ha risolto un problema simile.

\subsection{Le lezioni come corsi di aggiornamento}

Le lezioni frontali simulano i \emph{corsi di aggiornamento} aziendali: le
sessioni che un lead di design tiene al suo team di design (o il lead di
level design al suo team, e così via). Alla fine dell'aggiornamento c'è quasi
sempre un lavoro da fare, un obiettivo da portare a terra:
\textbf{l'obiettivo della settimana}. Si lavora di settimana in settimana, e
il perché verrà chiarito in seguito.

Struttura tipica di una lezione:
\begin{enumerate}
  \item breve sessione di domande e risposte sulla lezione precedente;
  \item analisi di eventuali esercizi particolarmente critici;
  \item lezione vera e propria;
  \item discussione di gruppo, se c'è tempo;
  \item assegnazione del task della settimana.
\end{enumerate}
Il docente tende a ``parlare e straparlare'', ad aprire parentesi e a saltare
da un argomento all'altro: l'unico modo per fermarlo è che siano gli studenti
a interromperlo con domande.

\subsection{La gerarchia della competenza}

L'inquiry-based learning può essere rappresentato come una scala di
competenza:
\begin{enumerate}
  \item la sola lezione frontale porta a \textbf{conoscenze incomplete},
        perché manca la parte pratica e di ricerca;
  \item il lavoro a casa (task, ricerca) porta a una competenza \textbf{più
        completa}, non sbagliata in sé, ma ancora parziale e complementare ad
        altre;
  \item la \textbf{condivisione} dei task e dei feedback con gli altri porta
        idealmente alla competenza più completa possibile;
  \item il picco della competenza è la sua \textbf{applicazione nei
        progetti} (fisici o digitali, di gruppo): ``con le cose che avete
        imparato e ricercato, riuscite a fare un gioco?''. Quando la risposta
        diventa positiva, o comunque si inizia a navigare nei processi di
        design per realizzarne uno, quella è la vittoria.
\end{enumerate}
Non si parte subito creando giochi: si parte dall'analisi, dal design e dalle
regole generali di design, prima ``nel vuoto'', poi contestualizzando i
framework, infine unendo i puntini nei progetti.

Rimarranno comunque delle lacune, per tutti, docente compreso, perché il game
design \textbf{non è una scienza esatta}, anzi, a rigore non è affatto una
scienza. Prende in prestito metodologie dalle discipline scientifiche, ma
contiene una componente imprevedibile e incalcolabile: \textbf{la
produzione}. Si possono applicare perfettamente le regole di un framework, e
quel framework può comunque non essere adatto al gioco, al proprio modo di
sviluppare o al tempo di produzione disponibile (alcuni framework richiedono
una ricerca e un'attenzione che il tempo non concede).

C'è poi la componente creativa, la parte ``romantica'' del lavoro: esiste, ma
è \textbf{metà del lavoro}. L'altra metà è ricerca, metodo scientifico al
100\%. Dagli stessi dati $n$ designer diversi, con sensibilità diverse,
possono ottenere $n$ risposte diverse. Non tutte saranno valide: vanno
applicate, messe a terra, e questa scrematura fa parte del lavoro.

\subsection{I livelli dell'apprendimento}

Un secondo modello, a gradini, descrive i livelli di apprendimento (è la
cosiddetta tassonomia SOLO):
\begin{description}
  \item[Pre-strutturale] il docente dà il compito e lo studente lo risolve,
        possibilmente il prima possibile. È esattamente ciò che \emph{non} si
        vuole.
  \item[Uni-strutturale] si affronta una sola disciplina: è il primo passo
        necessario, ed è il motivo per cui nel primo mese ci sono solo lezioni
        di game design.
  \item[Multi-strutturale] entrano in gioco altri dipartimenti: level design,
        game design, programmazione e così via.
  \item[Relazionale] si capisce che alcune cose valgono per più
        dipartimenti: alcuni framework di narrativa derivano dal design o
        possono essere usati per fare design; le regole di documentazione si
        applicano anche ad altri dipartimenti; le regole di level design si
        mescolano con quelle di game design. Dopotutto, per fare un gioco
        servono molti dipartimenti.
  \item[Astratto esteso] il punto finale: muoversi autonomamente e capire
        cosa usare e quando, sulla base della propria esperienza. L'unico modo
        per arrivarci è fare pratica e fare progetti.
\end{description}

% ---------------------------------------------------------------------
\section{Task, feedback e regole del corso}
\label{sec:task-feedback}

\subsection{I task come opportunità}

I task non sono strettamente obbligatori (esiste un numero minimo di
consegne), ma sono \textbf{opportunità}. Quanto e come ci si lavora dipende
dallo studente. Un task fatto in cinque minuti, o in cui si è stati ``furbi'',
non verrà letto né riceverà feedback.

Domande da porsi davanti a ogni task:
\begin{enumerate}
  \item \textbf{Perché} mi è stato assegnato? Qual è l'obiettivo?
  \item Come ci lavoro per ottenere un risultato \textbf{che abbia senso}? Se
        la ricerca e i dati indicano una direzione, è inutile andare in
        un'altra.
\end{enumerate}

\subsection{Quanto lavorare: un calcolo indicativo}

Se il corso è una simulazione di lavoro, vale il carico di un lavoro normale:
\begin{itemize}
  \item un lavoro a tempo pieno comprende da 30 a 45 ore settimanali; si
        considera una media di \textbf{40 ore};
  \item si sottraggono circa 9 ore di lezione: restano $\approx 31$, diciamo
        \textbf{30 ore};
  \item con circa 3 task a settimana: $30 / 3 = $ \textbf{circa 10 ore per
        task}.
\end{itemize}
Il numero può variare, ma rende l'idea: dedicare un'ora a un task e poi fare
altro non fa guadagnare nulla a nessuno. Lo studente ``butta soldi'', il
docente risparmia tempo leggendo il task in un minuto (o non leggendolo).
Dedicando invece attenzione, ricerca, cura nella formattazione e
sperimentazione, e facendolo con cognizione di causa, si impara.

\subsection{Niente voti, solo feedback}

Ad eccezione del corso di programmazione, non ci sono voti (né ``ottimo'', né
10, né 100). Ci sono \textbf{feedback}:
\begin{itemize}
  \item in game design, generalmente sulla documentazione;
  \item in level design, sui progetti pratici.
\end{itemize}
Il docente indica cosa va bene e cosa va modificato, concentrandosi in
particolare su ciò che va cambiato. Il feedback \textbf{non è mai un giudizio
sulla persona}, ma su ciò che è stato consegnato nel concreto.

\subsection{La condivisione dei task}

Tutti i task consegnati vengono \textbf{condivisi} con la classe, e possono
essere letti solo dopo aver ricevuto i feedback del docente. Si cresce così
non solo con i propri feedback, ma anche con quelli dati agli altri. Regola
importante: ciò che il docente dice a uno studente \textbf{vale per tutti}. Se
a una persona viene detto ``questa cosa non va mai più fatta'', da quel
momento vale per l'intera classe.

\subsection{Non copiare}

Copiare i task degli altri è una delle cose peggiori che si possano fare, non
per una questione morale, ma perché annulla il senso del percorso. Si può
dedicare a ogni task il tempo e l'attenzione che si ritengono opportuni, ed è
normale che col tempo ci si specializzi (preferendo, ad esempio, il level
design o la programmazione tecnica) e si investa meno dove si sente meno
necessità. In questi casi l'alternativa corretta è:
\begin{itemize}
  \item non consegnare;
  \item studiare ciò che hanno fatto gli altri;
  \item eventualmente riconsegnare in seguito.
\end{itemize}
Ci saranno momenti in cui si verrà sommersi dalle cose da fare: bisogna
imparare a scegliere.

\subsection{Il ``non leggo''}

Il docente ha piena facoltà di non dare alcun feedback a un task (il
cosiddetto \textbf{``non leggo''}) se:
\begin{itemize}
  \item non si è lavorato con attenzione;
  \item non si rispettano le indicazioni date;
  \item è stata violata una regola estremamente importante.
\end{itemize}
Un task ``non letto'' equivale a un task non consegnato e non entra nel
conteggio. Attenzione: un task \emph{perfetto} e un task che \emph{fallisce}
nel comprendere o applicare un framework sono \textbf{ugualmente validi},
purché ci si sia impegnati e si sia dedicato tempo, pazienza e ``sudore''.

\paragraph{Riconsegna.} Un task non letto, in ritardo o che si vuole
migliorare può essere riconsegnato e verrà conteggiato, a patto che:
\begin{itemize}
  \item i problemi segnalati siano stati risolti;
  \item sia stato fatto ulteriore lavoro, integrando anche tutto ciò che nel
        frattempo è stato detto agli altri colleghi nei loro feedback.
\end{itemize}
Il task va sviluppato con le competenze che si hanno nel momento della
riconsegna, non con quelle del momento della prima consegna.

\subsection{Requisiti per superare l'anno}

\begin{itemize}
  \item consegnare almeno l'\textbf{80\% dei task} assegnati (soglia
        malleabile);
  \item partecipare almeno all'\textbf{80\% delle lezioni}, in presenza o
        da remoto, preferibilmente in presenza; i ritardi vengono valutati
        caso per caso;
  \item comportarsi da persone corrette e rispettose (è stata firmata una
        carta dello studente): non urlare, non alzare le mani, non
        vandalizzare, non compiere atti illegali.
\end{itemize}
Impegnarsi di più porta inoltre a essere selezionati per varie opportunità.

\subsection{Scadenze}

\begin{itemize}
  \item la scadenza di un task è \textbf{sempre prima dell'inizio della
        lezione} successiva indicata;
  \item se la classe, come team, mantiene comportamenti scorretti, il docente
        può \textbf{anticipare le deadline} per tutti.
\end{itemize}

\subsection{Accessibilità e DSA}

Gli studenti con disturbi specifici dell'apprendimento o neurodivergenze
(dislessia, ADHD, autismo, ecc.) sono invitati a segnalarlo. Quando si
affronteranno le regole di formattazione e documentazione (errori
grammaticali, mancanze, incoerenze) se ne terrà conto e si potranno trovare
dei compromessi, pur ricordando che nel mondo del lavoro nessuno farà
``sconti''. La documentazione può essere inviata all'amministrazione o alla
segreteria, che la comunicherà a tutti i docenti, ma il consiglio è parlarne
anche direttamente con ogni singolo docente.

% ---------------------------------------------------------------------
\section{Le competenze trasversali}
\label{sec:soft-skill}

\subsection{I giochi si fanno in gruppo}

Si possono fare giochi da soli, ed è un'impresa ammirevole; ma lavorando da
soli \textbf{tutto ciò che si guadagna in libertà si perde in produttività}.

\paragraph{Il caso Stardew Valley.} \emph{Stardew Valley} è stato realizzato
da una sola persona, Eric Barone, in anni di sviluppo e di sacrifici: faceva
un altro lavoro per mantenersi (in un cinema) e di notte programmava, faceva
l'arte e componeva la musica. Il risultato è stato straordinario e lo ha reso
ricco, ma si tratta di un caso estremamente raro. Quanti giochi simili, nati
nelle stesse condizioni, sono falliti? Molti di più.

\paragraph{Il rischio d'impresa.} Lavorare da soli non è sbagliato di per sé
e ha aspetti molto positivi: si ha il privilegio di creare esattamente ciò che
si vuole. Ma ci si fa carico anche di tutto il \textbf{rischio d'impresa}: se
si lavora a un gioco per cinque anni e poi il gioco va male, resta il
vantaggio di averlo pubblicato, ma il dopo è sperare di vivere di quella
release o di finanziare con essa il gioco successivo. Le grandi aziende hanno
gli stessi problemi, ma li distribuiscono su più persone e più responsabilità.

L'obiettivo di un professionista è \textbf{saper lavorare in diverse
condizioni}, e la più comune è quella del dipendente (o del freelance /
libero professionista con partita IVA): si viene assunti per le proprie
capacità da un'azienda, piccola o grande che sia, e si lavora con altre
persone.

\subsection{Essere persone con cui si vuole lavorare}

Se le persone non vogliono lavorare con voi, quello è \textbf{il primo motivo
di licenziamento}. E se non si viene licenziati, il clima nell'azienda
diventa terribile: la produttività crolla, la motivazione crolla. Le
\emph{soft skill}, cioè l'insieme di abilità che rendono una persona
piacevole con cui lavorare, vanno coltivate e diventano sempre più importanti
col passare degli anni (relativamente meno ora, molto di più quando i team
saranno più grandi). Anche chi è un ``lupo solitario'' deve imparare a:
\begin{itemize}
  \item lavorare con altre persone;
  \item comunicare con altre persone;
  \item trovare un punto in comune da cui partire per fare le cose.
\end{itemize}

\subsection{Indipendenza e proattività}

Sono le altre due parole chiave del corso:
\begin{description}
  \item[Indipendenza] spesso si verrà lasciati da soli, non per cattiveria ma
        per sfida: quando si cerca una risposta da soli, ciò che conta non è
        tanto la risposta quanto \textbf{il processo}. Una risposta data dal
        docente non vale niente.
  \item[Proattività] chiedere, domandare, agire da soli senza aspettare
        direttive, perché spesso quelle direttive non ci sono, a maggior
        ragione nei team piccoli.
\end{description}
Tutto questo si concretizza nel parlare, nel presentare e nel fare
\emph{post-mortem} (l'analisi a posteriori di un progetto concluso, che verrà
approfondita più avanti), e in generale nell'attenzione alla crescita
personale.

\subsection{Fare domande}

Il corso insiste sulle domande: alzare la mano, intervenire anche ``a
bomba'', fermare il docente ogni volta che qualcosa non è chiaro. Se una
persona non conosce qualcosa, molto probabilmente c'è almeno un'altra persona
che non la conosce. Il docente può dare per scontata la conoscenza di termini
che non lo sono, ed è del tutto normale, soprattutto all'inizio.

% ---------------------------------------------------------------------
\section{Strumenti: engine, Unreal e materiali}
\label{sec:engine}

\subsection{Unreal Engine e gli altri motori}

Il motore di riferimento del corso, usato in level design, è \textbf{Unreal
Engine}: ``un gran mattone'', ma che dà a un designer molte possibilità. Il
vantaggio principale è che permette a un designer non pratico di
programmazione di \textbf{prototipare tramite i Blueprint} (il sistema di
scripting visuale di Unreal). Questo non impedisce di passare a Unity o ad
altri engine, e viceversa: il docente stesso ha studiato e lavorato a lungo
con Unreal, ma la sua ultima esperienza professionale (come level designer) è
stata su Unity, un engine che non conosceva, a dimostrazione che il salto da
un engine all'altro è possibile.

Panorama dei motori citati:
\begin{itemize}
  \item \textbf{Unreal Engine} e \textbf{Unity}: i due più diffusi, con
        modelli di sorgente ``più o meno'' aperto;
  \item \textbf{Godot}: engine open source;
  \item \textbf{CryEngine}: proprietario di Crytek;
  \item \textbf{Snowdrop}: engine di Ubisoft;
  \item Rockstar usa un proprio engine proprietario;
  \item alcune aziende usano versioni \textbf{personalizzate} di engine
        esistenti: ad esempio Milestone, nei giochi di corse, usa un Unreal
        customizzato;
  \item CD Projekt Red ha un engine proprietario (con i problemi di
        documentazione già citati).
\end{itemize}
In sintesi: ogni azienda può avere un engine proprietario oppure una versione
modificata di un engine commerciale, ma si parte sempre da una base.

\subsection{Programmazione: cosa ci si aspetta da un designer}

Non è richiesto alcun bagaglio pregresso di programmazione: il corso parte da
zero. Gli studenti sono stati selezionati (il corso è a numero chiuso) per
competenze analitiche e critiche sopra la media. Anche nel mercato del lavoro,
per una posizione di game designer generalmente \textbf{non ci si aspetta
competenza di scripting} (C++, Java, C\#, ecc.), a meno di voler fare il
\emph{technical designer}, e anche lì varia da azienda ad azienda.

\subsection{Versione e preparazione}

Indicazioni pratiche emerse in aula:
\begin{itemize}
  \item scaricare Unreal Engine; la versione consigliata per il corso è
        stata indicata nella \textbf{5.6} (con conferma successiva). La
        stabilità dipende da cosa si deve fare, da quali tool e plugin si
        usano: i problemi nascono spesso dai plugin. Nel passaggio da Unreal 4
        a Unreal 5, ad esempio, il sistema \emph{Landscape} era inizialmente
        inutilizzabile, mentre in Unreal 4 funzionava molto bene;
  \item è consigliato installare tramite \textbf{Epic Games Launcher}, più
        comodo: permette anche di riscattare giochi gratuiti periodici e di
        accedere ai pacchetti di \textbf{Fab} (il marketplace di asset
        acquisito da Epic), che ogni mese offre alcuni asset gratuiti utili
        per esercitarsi nel level design (ad esempio per fare \emph{speed
        level design});
  \item molti pacchetti di UI o di sistemi di interazione sono utili per
        prototipare al volo, ma spesso inutilmente complicati;
  \item un plugin di \textbf{materiali di prototipazione a griglia} (adattivi,
        con metriche, allineabili alla griglia, con materiali già pronti e
        colori modificabili, dal costo di circa 5 euro) è indicato come
        ``salvavita'' per il level designer; se necessario, i pacchetti
        indispensabili vengono forniti;
  \item portare \textbf{sempre il portatile} alle lezioni di level design
        (da novembre, a partire dall'interfaccia dell'engine); alle lezioni
        di game design basta che lo porti almeno metà della classe, per i
        lavori a coppie. Serve comunque qualcosa per prendere appunti
        (``tavolette d'argilla e cannucce, se volete fare i babilonesi'',
        oppure carta e penna).
\end{itemize}
Chi ha un portatile poco potente dovrebbe comunque provare a installare
l'engine e aprire un template per vedere come gira: all'inizio non si faranno
cose pesanti. Nelle prime lezioni di level design è importante mettere subito
le mani sull'interfaccia: ascoltare dove si trova uno strumento (ad esempio lo
\emph{snapping}) senza provarlo porta a dimenticarlo.

% ---------------------------------------------------------------------
\section{Letture e risorse di riferimento}
\label{sec:risorse}

\subsection{I libri}

\begin{description}
  \item[\emph{The Art of Game Design: A Book of Lenses} (Jesse Schell)] il
        testo \textbf{davvero importante}. Alcune lezioni finali dell'anno si
        baseranno sulle sue terminologie. È utile perché apre la mente sulla
        possibilità di applicare metodologie diverse a giochi diversi, ma
        soprattutto per la \textbf{terminologia} comune. Esiste anche in
        audiolibro (circa 25 ore).
  \item[\emph{Challenges for Game Designers} (Brenda Romero e Ian Schreiber)]
        libro \emph{entry level} con moltissimi esercizi, molto caro al
        direttore del corso.
  \item[\emph{21st Century Game Design} (Chris Bateman)] utile come
        retrospettiva: è un libro visionario che ha previsto metodi e
        soprattutto \emph{archetipizzazioni} dei giocatori tuttora valide nel
        design e nel marketing. Servirà per una delle lezioni più importanti
        dell'anno. È utile leggerlo per vedere quanto aveva previsto, ma è
        utile ``fino a un certo punto''.
  \item[\emph{An Architectural Approach to Level Design} (Christopher
        Totten)] non strettamente necessario, ma di lettura semplice. Ha due
        anime: parte dall'architettura (babilonese, greca, ecc.) e la trasla
        su esempi di level design (come \emph{Doom} e \emph{Team Fortress}).
        Un tempo ne esisteva un riassunto su Gamasutra (oggi
        \emph{Game Developer}).
  \item[\emph{A Theory of Fun for Game Design} (Raph Koster)] utile fino a un
        certo punto: contiene diverse idee superate, che lo stesso Koster ha
        affrontato in vari talk reperibili nel GDC Vault. Verrà ripreso nelle
        prime lezioni, poi messo da parte.
\end{description}

\subsection{Canali e conferenze}

\begin{description}
  \item[GDC (Game Developers Conference)] il suo canale raccoglie talk e
        panel estremamente importanti: è ``uno strumento misterioso che ci
        servirà più tardi''. \textbf{Attenzione}: per molti talk mancano ancora
        le competenze per interpretarli, e alcuni potrebbero fuorviare o
        mettere in testa idee strane. La GDC si concentra spesso su
        produttività e mobile.
  \item[Il \emph{Dragon Speech} di Chris Crawford] il talk del fondatore della
        GDC: appartiene a un passato lontano, ma è molto motivante. Da vedere
        subito.
  \item[Game Maker's Toolkit] canale particolarmente utile per nozioni
        pratiche, con video brevi e alcune interviste.
  \item[Nordic Game] altra conferenza importante, con talk e panel di esperti
        del settore.
\end{description}

\subsection{Associazioni e premi: come leggerli}

Le associazioni di settore (ad esempio IIDEA in Italia, o la IGDA) non sono
il male, anzi: sono un canale importante per affacciarsi all'estero o per
accedere a fondi. Va però ricordato che \textbf{tutelano i propri associati}:
non esiste solo il loro palinsesto, e bisogna valutare quanto spazio occupano
nel panorama della visibilità.

Lo stesso vale per i \textbf{premi} (Game Awards, ``Best Mobile Game'',
``Games for Change'', ecc.): al designer \emph{non interessano}. Che un gioco
abbia vinto ``più di 90 premi'' non dice nulla di utile. Servono a dare
visibilità e lustro ai giochi ed è giusto che esistano, perché ampliano il
panorama e il numero di giocatori. Ma ciò che interessa al designer sono i
\textbf{dati}:
\begin{itemize}
  \item come sono state implementate certe cose;
  \item come è stato fatto il marketing e dove il gioco è stato presentato;
  \item non quante wishlist ha fatto, ma \emph{come} le ha fatte;
  \item in che periodo storico è uscito e quali erano i competitor in quel
        momento.
\end{itemize}

% ---------------------------------------------------------------------
\section{Che cosa fa un game designer}
\label{sec:cosa-fa}

Per arrivare a una prima definizione, la classe ha condotto un brainstorming:
ognuno doveva indicare una cosa che, secondo lui, fa il game designer (senza
ripetere quelle già dette). Ne è emersa la lista seguente, commentata punto
per punto.

\subsection{Progettare meccaniche e struttura del gioco}

Il ruolo più generico è quello del \textbf{core mechanic designer}, detto
anche \emph{gameplay designer} e, un tempo, \emph{black box designer}: si
occupa delle \textbf{meccaniche di gioco}. Si va dal piccolo al grande (in
funzione anche della dimensione dell'azienda):
\begin{itemize}
  \item nel \textbf{piccolo} ci sono le singole meccaniche;
  \item nel \textbf{grande} ci sono i \textbf{sistemi}. Nei giochi
        \emph{sistemici}, con molti sistemi interconnessi (ad esempio
        \emph{GTA}), esistono specializzazioni dedicate a bilanciare e
        supervisionare i sistemi: è il \emph{system design}.
\end{itemize}

\paragraph{Quando un gioco è ``sistemico''?}
Non dipende solo da come le meccaniche interagiscono con il giocatore (le
\emph{dinamiche}, secondo i framework che passano dalle meccaniche alle
dinamiche, che verranno studiati più avanti), ma anche da \emph{quante}
meccaniche ci sono e da quanto interagiscono tra loro. Non diremmo che
\emph{Super Mario} sia un gioco sistemico, eppure contiene dei sistemi: la
linea è sottile. Quando i sistemi aumentano molto, sono complessi e si
influenzano a vicenda, il gioco diventa sistemico. Esempio: in \emph{The
Legend of Zelda: Breath of the Wild} il vento modifica quanto il personaggio
sale o scende; il fuoco può incendiare l'erba; l'erba prende fuoco se colpita
da una freccia infuocata; e così via. Nel mercato moderno i giochi sistemici
sono spesso, ma non sempre, legati a mappe \emph{open world}.

\subsection{Le 3C: camera, character, control}

La base del gioco, da cui generalmente si inizia a lavorare (anche nei
prototipi), è la ``sacra trinità videoludica'' delle \textbf{3C}:
\begin{description}
  \item[Camera] che tipo di inquadratura ha il gioco: prima persona, terza
        persona, isometrica, ecc.
  \item[Character] quanto è grande il personaggio, che tipo di personaggio è
        e, di conseguenza, quali azioni può compiere.
  \item[Control] i controlli: che cosa il personaggio può fare, quali sono i
        \emph{verbi} di gioco.
\end{description}
Le 3C sono definite dal designer, ed esiste anche una specializzazione per
questo lavoro. Esistono varianti: quella del direttore del corso è
\textbf{4C + J}, cioè \emph{camera, character, control, conflict} (conflitto)
e \emph{juice} (il ``succo'', il feedback che rende piacevole l'interazione).
Verranno approfondite più avanti.

\subsection{Il personaggio}

Il designer definisce il personaggio dal punto di vista delle meccaniche
(dimensioni, azioni, capacità), all'interno delle 3C. L'aspetto estetico è
invece competenza di altri dipartimenti.

\subsection{Documentare le informazioni per gli altri dipartimenti}

Assolutamente vero: la documentazione è una delle prime attività del
designer. Più l'azienda è grande, più questo lavoro diventa prominente; più è
piccola, meno è rilevante. In \textbf{pre-produzione}, però, è sempre
fondamentale.

\subsection{Mediatore e leader ``al servizio''}

Il designer è un \textbf{mediatore} nei conflitti interni al team, una sorta
di collante che si relaziona con dipartimenti diversi (non sempre gli stessi,
a seconda della fase di produzione). A volte fa anche da leader, a seconda di
come sono gestite la pipeline e la gerarchia.

Questa leadership è difficile da definire perché \textbf{non è verticale}: è
una leadership \emph{al servizio} (\emph{servant leadership}). Il designer
``è asservito al programmatore'': il suo compito non è arrivare con un'idea e
scaricarla addosso agli altri, né portare la soluzione e lasciare il problema
a qualcun altro, ma \textbf{seguire il processo a basso livello}, sporcandosi
le mani.

\paragraph{I feature team.} Nei team medio-piccoli o medio-grandi possono
esistere sotto-team (\emph{feature team}) dedicati a un obiettivo specifico:
ad esempio un designer incaricato del character o della camera, affiancato da
uno o due programmatori e da un artista. In questi casi il designer
supervisiona il lavoro e gestisce la produzione del sotto-team.

\subsection{Gestire l'avanzamento del prodotto}

Esiste una specializzazione che si occupa solo di questo, il
\textbf{producer}, a un livello gerarchico completamente diverso (e che
potrebbe essere una strada percorribile al termine del percorso di studi).
All'interno di un feature team anche il designer può farlo, ma non è detto.
Può certamente \emph{dare la direzione}: nella sua documentazione può
indicare che per arrivare a uno stato giocabile del prototipo serve un certo
lavoro di arte e di programmazione, per lo stato successivo un altro, e così
via. In ambito produttivo esistono metodologie strutturate come
\textbf{Agile}, che verranno viste in seguito.

\subsection{Presentare le idee}

Il pitching (sezione~\ref{sec:gerarchia-corso}) è richiesto al designer sia
in pre-produzione sia in produzione, in misura che dipende dall'azienda. Ad
esempio Ubisoft lascia ai designer, e in realtà a qualsiasi membro, la
possibilità di \emph{pitchare} un'idea prima di una nuova produzione, perché
possa iniziare a prendere vita.

\subsection{Bozze e prototipi}

Le bozze vanno bene, ma devono essere \textbf{chiare}. C'è un'eterna lotta tra
documentazione e gioco giocabile: se si fa solo documentazione il gioco non
avanza; se si fa solo il gioco manca la documentazione e si perde la visione.
Sono due cose \textbf{complementari}.

Il modo per comunicare un'idea in forma grezza ma giocabile è il
\textbf{prototipo}. Un prototipo può essere qualsiasi cosa: di carta, in
Unreal con Blueprint ``brutti''. \emph{Non deve essere bello, deve essere
giocabile}.

\paragraph{Esempio: il prototipo di carta di Mario + Rabbids.} Durante lo
sviluppo di \emph{Mario + Rabbids}, i programmatori di Ubisoft erano impegnati
su altri progetti (come \emph{Ghost Recon}). I designer realizzarono allora
un \textbf{gioco da tavolo} che riproduceva le meccaniche del gioco, quelle
di uno strategico a turni sul modello di \emph{XCOM}. E funzionava.
L'obiettivo era \textbf{validare un'idea}, validare un gameplay.

Altri modi per rendere chiara un'idea sono i \emph{mock-up} nella
documentazione e le immagini modificate o i \emph{mash-up} di qualsiasi
natura. Lato documentazione l'importante è che sia chiaro; lato prototipo,
che sia giocabile.

\subsection{Fare ricerca per l'idea, e poi smontarla}

Il designer ha una componente creativa importante e fa ricerca per arrivare a
un'idea. Ma ciò che deve poi fare necessariamente è \textbf{smontarla}. Il
suo lavoro è scientifico, esattamente come il \textbf{metodo scientifico}, il
cui obiettivo è \emph{confutare un'ipotesi}: si arriva a un'ipotesi tramite
sperimentazione e poi si cerca $n$ volte di confutarla; se resiste, la si
assume come valida \emph{in un certo contesto} (le regole della fisica
classica, ad esempio, non valgono nella meccanica quantistica).

Allo stesso modo il designer arriva a un'idea (tramite ricerca o chiedendo
agli altri: le idee le hanno tutti, non sono un privilegio del designer),
cerca di smontarla e poi cerca di \textbf{metterla a terra}. Se l'idea non è
applicabile, non è implementabile, finché non si trova un modo per
realizzarla va scartata. Può essere l'idea più bella dell'universo, ma
\textbf{se non può essere messa a terra, non esiste}.

\paragraph{L'originalità.} Serve un'idea originale? Dipende dall'obiettivo e
dalla visione del progetto:
\begin{itemize}
  \item se l'obiettivo è vendere molto, un'idea originale non serve
        necessariamente;
  \item se si vuole creare un nuovo filone, serve un'idea, se non originale,
        almeno interessante: può essere un discorso di marketing o una
        meccanica innovativa, con il rischio di non vendere.
\end{itemize}
Il ruolo del \emph{director} e il concetto di \emph{visione} verranno
approfonditi in seguito.

\subsection{Analizzare i competitor}

Arrivati a un'idea (``vogliamo un certo tipo di combattimento''), si va a
vedere se qualcuno l'ha già fatta e \emph{come}. Le meccaniche non sono
privilegio di nessuno, salvo quando vengono brevettate, come il
\textbf{Nemesis System} di \emph{La Terra di Mezzo: L'Ombra di Mordor}.

\paragraph{Esempio: il combattimento da ridimensionare.} In un progetto
(presso Stormind Games) il team si trovò con un sistema di combattimento
troppo costoso: mancavano le risorse e gli animatori per realizzarlo. Poco
output lato gameplay significava non poter fare qualcosa di complesso e
profondo come in \emph{Elden Ring}; al contrario, il team aveva molti
artisti e poteva puntare su qualcosa di più artistico. Si iniziò così a
orbitare intorno al corpo a corpo di \emph{The Last of Us} (parte I e II).
Sul canale della GDC esistono diversi talk che raccontano cosa è stato fatto
e come, con successi, insuccessi e costi (ad esempio del \emph{motion
capture}): da quei talk si ricavarono moltissime informazioni utili.

I competitor possono essere:
\begin{itemize}
  \item \textbf{diretti}: un gioco molto simile al proprio;
  \item \textbf{di mercato}: un gioco che domina il mercato in cui ci si vuole
        affacciare;
  \item \textbf{indiretti}: un gioco con un sistema simile al proprio, ma
        tutto il resto diverso.
\end{itemize}
Regola d'oro: \textbf{là fuori, nascosto nell'ombra, c'è sicuramente un gioco
molto simile al vostro}. Va solo trovato. E se per qualche ragione non ha
venduto, bisogna capire perché.

\subsection{Lavorare in engine e fare tuning}

Il level designer lavora in engine, così come il technical designer; ma
qualsiasi game designer, a seconda di ciò che deve fare, può finire in
engine, soprattutto per il \textbf{bilanciamento} e il \textbf{tuning}.

Il \emph{tuning} è l'azione di bilanciare con precisione i dati. Esempio: in
\emph{Super Mario} si modifica quanto in alto salta il personaggio, con che
velocità, in quanto tempo raggiunge l'apice, e così via. Questo lavoro rende
il gioco:
\begin{itemize}
  \item ad alto livello, dotato della \textbf{progressione} giusta;
  \item a basso livello, \textbf{piacevole da giocare}.
\end{itemize}
Il \emph{fine tuning} è un lavoro di dettaglio, paziente e minuzioso.

\subsection{Bilanciare}

Il bilanciamento, emerso tardi nel brainstorming, è una delle attività
centrali del designer e ha anche una specializzazione dedicata
(l'\emph{economy designer}, sezione~\ref{sec:tipi-designer}).

\subsection{Analizzare il mercato}

In pre-produzione, di solito, sì: capire le tendenze, i generi che cambiano
nel tempo, la compatibilità del gioco con il momento storico, in vista ad
esempio di un possibile publisher. Va capito quanto, come e con quanto tempo
farlo, perché si è sempre limitati da molte cose.

\subsection{Sintesi}

\begin{quote}
Il designer parte dall'idea ma non si ferma lì: la confuta, lavora sui
dettagli, bilancia, fa \emph{fine tuning}; lavora con il team e fa da
\emph{manager} per diversi dipartimenti; e naviga costantemente tra i limiti.
\end{quote}

% ---------------------------------------------------------------------
\section{I limiti del design}
\label{sec:limiti}

I limiti ci sono sempre, e nel mondo del design sono spesso numerosi. I più
importanti entro cui navigare sono:
\begin{description}
  \item[Tempo] ad alto livello, l'idea deve essere realizzabile nel tempo di
        produzione disponibile; a basso livello, il proprio lavoro
        settimanale deve essere fattibile nel tempo che si ha. Non ci si può
        dividere tra mille feature team se realisticamente se ne può seguire
        uno solo.
  \item[Tecnologia] si può o non si può fare qualcosa? Ci sono i server per un
        multiplayer con un certo numero di giocatori? Ci si può permettere,
        ad esempio, una fauna che reagisce alla fisica come quella di
        \emph{GTA VI}, che richiede il meglio dei \emph{technical artist}?
  \item[Background e competenze] che cosa sa fare il team? Realisticamente,
        chi ha lavorato dieci anni su giochi horror viene assunto per un
        horror, perché ha le competenze per risolvere problemi già visti.
  \item[Persone e motivazione] il team ha la motivazione e le competenze per
        fare un certo tipo di gioco, e per lavorare insieme?
  \item[Denaro] tutto ruota intorno ai soldi. Che si lavori da soli o in mille,
        le persone vanno pagate e il gioco deve avere un ritorno economico.
        Bisogna decidere dove tracciare la linea: raggiungere il
        \textbf{break-even} (guadagnare esattamente quanto speso) o puntare a
        grandi guadagni. Con un \textbf{publisher}, inoltre, una parte
        consistente dei ricavi va a lui (ripartizioni dell'ordine di 80/20 o
        70/30, che dipendono dal contratto firmato).
\end{description}

\subsection{Non siete ``gli uomini delle idee''}

Per qualcuno può essere un problema scoprirlo, ma il designer non è
``l'uomo (o la donna) delle idee'': accanto alla parte creativa c'è una parte
tecnica, e si è fortemente limitati da molti fattori. Questi limiti, però,
sono un aiuto:
\begin{itemize}
  \item riducono il ventaglio di possibilità potenzialmente distruttive per il
        team (cose che non si riescono a fare, che non piacciono o per cui
        mancano le risorse);
  \item aiutano a mirare bene, trovando soluzioni creative in uno spazio più
        piccolo e più attuabile.
\end{itemize}

\begin{quote}
\centering
\textbf{Le idee sono nulla, l'implementazione è tutto.}
\end{quote}

% ---------------------------------------------------------------------
\section{Ruoli e specializzazioni}
\label{sec:tipi-designer}

\subsection{La scala gerarchica}

\begin{description}
  \item[Creative Director (CD)] mantiene la \textbf{visione} e dà l'idea, in
        modo che tutto segua una certa direzione. In teoria può diventarlo
        chiunque, ma con un'esperienza notevole, nell'ordine dei decenni: si
        assume che in quel tempo abbia visto come lavorano molti dipartimenti.
        Naturalmente un CD con un background artistico ha meno competenze di
        design, e viceversa.
  \item[Director di disciplina] uno per dipartimento: \emph{game director},
        \emph{level director}, \emph{art director}, ecc.
  \item[Lead] il capo del dipartimento, che gestisce i gruppi di persone. Chi
        lavora in un dipartimento risponde al proprio lead (ad esempio un game
        designer risponde al \emph{lead designer}, non al game director).
  \item[Senior / Principal] i profili più esperti (a seconda dell'azienda).
  \item[Junior] il profilo di inizio carriera.
  \item[Intern] lo stagista.
\end{description}
Lo schema si replica uno a uno per ogni dipartimento: arte, suono, animazione
(a volte arte e animazione insieme, o \emph{rigging} e animazione insieme),
programmazione, UX/UI e così via. A fare la differenza, a parità di ruolo, è
il \textbf{livello di seniority}.

\subsection{UX e UI}

UX e UI stanno a metà tra game design e arte: a seconda dell'azienda il ruolo
si chiama \emph{UX/UI Artist} o \emph{UX/UI Designer}.
\begin{description}
  \item[UX (User Experience)] come una persona interagisce con le cose. Esempio:
        su uno smartphone, un pulsante accende e spegne lo schermo; per
        digitare la password si scorre verso l'alto e compaiono i punti da
        toccare.
  \item[UI (User Interface)] come sono ``vestite'' queste cose: ad esempio
        l'aspetto della barra dei punti vita (HP), realizzata con asset 2D o
        3D.
\end{description}

\subsection{Le specializzazioni del game design}

Più l'azienda è grande, più la specializzazione è puntuale.
\begin{description}
  \item[Gameplay / Core mechanics designer] (sono la stessa cosa) si occupa
        delle meccaniche e delle 3C.
  \item[System designer] si occupa dei sistemi e delle loro interazioni.
  \item[Economy designer] si occupa del bilanciamento; lavora soprattutto con
        fogli di calcolo (Excel, Google Sheets).
  \item[Combat designer] si occupa del combattimento; negli annunci di lavoro
        indica quasi sempre un combattimento \emph{action}, con gameplay
        reattivo (non strategico), anche se non è una regola.
  \item[Multiplayer designer] si occupa delle interazioni multigiocatore.
  \item[Technical designer] fa da ponte tra game design e programmazione.
\end{description}

\subsection{Le specializzazioni del level design}

\begin{description}
  \item[Encounter designer] posiziona i nemici nella mappa in modo che gli
        scontri siano bilanciati e la progressione equilibrata come
        desiderato.
  \item[World / map designer] popola il mondo di gioco.
  \item[Quest designer] popola il mondo con le missioni e pensa alla
        progressione al loro interno, cioè a dove il giocatore deve andare (si
        vedano ad esempio gli annunci di lavoro di CD Projekt Red o di
        Larian).
  \item[Technical level designer] fa da ponte tra level design e
        programmazione.
\end{description}

\subsection{Discipline vicine: suono, UI, produzione}

Il monte ore del corso non permette di uscire come \emph{narrative designer}
(o \emph{sound designer}) completi. Ciò non toglie che si possano coltivare
queste discipline:
\begin{itemize}
  \item nei progetti di gruppo si può occuparsi anche del suono, della UI/UX,
        o provare ad applicare metodologie come Agile per fare
        \emph{producing} e gestire il team;
  \item il tempo dedicato a queste attività va però sottratto al game design,
        e dipende da quante persone ci sono nel team, da quanto lavoro c'è e
        da quanto è grande il progetto;
  \item avere sviluppato e rilasciato un progetto, anche piccolo (anche una
        \emph{game jam}), documentando i propri processi, è già un ottimo
        argomento da presentare in futuro;
  \item la via migliore resta lavorare nell'industria e, da lì, chiedere di
        imparare anche altre discipline.
\end{itemize}
Un gioco si può fare su qualsiasi cosa: sugli scarafaggi, sui piedini delle
sedie, sui pennarelli. Si parte avvantaggiati da ciò che è vicino alla propria
sensibilità: chi ama la musica può pensare a un gioco con meccaniche legate
alla musica, con un gameplay innovativo ma contenuto e realizzabile. Il design
è un cappello che comprende molte discipline.

% ---------------------------------------------------------------------
\section{Approfondimenti emersi in lezione}
\label{sec:approfondimenti}

\subsection{Roguelike e roguelite}

Un esempio di come il marketing conia generi: Supergiant Games, con
\emph{Hades}, rese celebre il termine \textbf{roguelite} al posto di
\emph{roguelike} (pur non essendo il primo roguelite in assoluto). La
differenza:
\begin{itemize}
  \item il \textbf{roguelike} si fonda su una serie di punti fissati in una
        sorta di \emph{manifesto} che definisce che cosa rende un gioco
        roguelike; uno su tutti: a ogni morte si ricomincia la \emph{run} da
        capo;
  \item il \textbf{roguelite} rompe alcuni di questi punti, in particolare
        introducendo una \textbf{progressione esterna alla run}: il giocatore
        può potenziarsi fuori dalla singola partita per vincere più
        facilmente in quelle successive. \emph{Hades} ne è il caso principe.
\end{itemize}
Analogamente, un mobile game sviluppato dal direttore del corso fu
promosso come ``il primo \emph{donation game} al mondo'': un gioco che
convertiva il tempo di gioco in donazioni reali a enti benefici per animali.
Anche qui un neologismo creato lato marketing per definire un genere.

\subsection{Applicazioni fuori dall'industria videoludica}

Le competenze sviluppate non servono solo ai videogiochi:
\begin{itemize}
  \item i \textbf{render architettonici} oggi possono essere realizzati con
        Unreal, e diversi professionisti formati nel game design lavorano in
        questo campo;
  \item la \textbf{gamification}, cioè l'applicazione di regole videoludiche
        a contesti non ludici (mostre interattive, progetti culturali,
        consulenze come quella per applicare questi concetti alle macchine da
        allenamento di un'azienda di fitness), è in Italia uno dei modi più
        redditizi di spendere queste competenze, spesso con finanziamenti
        pubblici.
\end{itemize}

% ---------------------------------------------------------------------
\section{Il primo task}
\label{sec:primo-task}

\begin{enumerate}
  \item Creare un indirizzo e-mail \textbf{professionale} (consigliato: un indirizzo Gmail nella forma
        \texttt{nome.cognome@gmail.com}; l'anno di nascita è accettabile).
        Va bene anche un indirizzo professionale già esistente usato per
        lavoro. Da evitare indirizzi poco seri (``non \texttt{blackfire67}'').
  \item Una persona della classe raccoglie tutti gli indirizzi in un unico
        documento e lo invia via e-mail al docente, che aggiungerà tutti al
        Drive condiviso del corso (dove verranno caricate lezioni, slide e
        materiali, per avere i documenti tutti nello stesso posto).
  \item Installare Discord e iscriversi al server del corso, dove passano le
        comunicazioni ufficiali. È consuetudine, e utile, creare anche un
        gruppo tra studenti.
\end{enumerate}
\textbf{Scadenza}: prima della lezione del venerdì.

% ---------------------------------------------------------------------
\section*{Punti chiave del capitolo}
\addcontentsline{toc}{section}{Punti chiave del capitolo}

\begin{itemize}
  \item Il corso è una \textbf{simulazione di lavoro}: ci si tratta da
        colleghi, con una gerarchia da studio di sviluppo.
  \item Si impara con l'\textbf{inquiry-based learning}: lezione, task,
        dubbi, domande, condivisione, progetti.
  \item Il fallimento affrontato con la mentalità giusta è la principale fonte
        di crescita; i feedback riguardano il lavoro, mai la persona.
  \item Il lavoro del designer è per metà creativo e per metà
        \textbf{ricerca e metodo scientifico}; il game design non è una
        scienza esatta.
  \item \textbf{Documentazione} e \textbf{pitching} sono il cuore del
        mestiere.
  \item Il designer progetta meccaniche e sistemi (partendo dalle
        \textbf{3C}), documenta, media, prototipa, confuta le idee, studia i
        competitor, bilancia e fa tuning.
  \item La leadership del designer è \textbf{al servizio} del team, non
        verticale.
  \item I limiti (tempo, tecnologia, competenze, persone, denaro) restringono
        e indirizzano la creatività.
  \item \textbf{Le idee sono nulla, l'implementazione è tutto.}
\end{itemize}

%    % ...
%
%  (Con hyperref aggiungere anche \renewcommand{\theHchapter}{z.\arabic{chapter}}
%  nella parte zero, per evitare ancore duplicate.)
%  Tutte le etichette di questo capitolo hanno il prefisso "z3:"; i
%  rimandi ai capitoli precedenti usano le etichette "z1:" e "z2:" già
%  definite in GD0_1 e GD0_2, e alcune etichette dei capitoli del corso
%  (GD1, GD3).
% =====================================================================

\chapter[Agile e Scrum nella pratica]{Agile e Scrum nella pratica: l'esperienza di chi ci è passato}
\label{cap:z3-agile-scrum}

Le prime due lezioni zero (capitoli~\ref{cap:z1-feedback}
e~\ref{cap:z2-tempo-idee-leadership}) guardavano al lavoro in team dal punto
di vista di un professionista con quindici anni di industria alle spalle.
Questa terza lezione cambia prospettiva: a parlare non è un veterano, ma
\textbf{due persone che hanno fatto lo stesso percorso di studi} di chi
ascolta e che, meglio di chiunque altro fino a quel momento, sono riuscite a
\emph{internalizzare} Agile e Scrum nei progetti dell'accademia.

La lezione è introdotta e chiusa da Diego Ricchiuti\footnote{Nella
trascrizione compare soltanto il nome di battesimo, ``Diego'', con cui i due
ospiti si rivolgono più volte a lui; il cognome è quello indicato per il
direttore dell'accademia nel capitolo~\ref{cap:z1-feedback}. Dal contesto è
anche il docente che ha già spiegato Scrum e Agile alla classe nelle lezioni
precedenti (``Diego vi avrà già ampiamente indirizzato'', ``come avete fatto
con Diego nelle lezioni scorse'').}, che ha già spiegato alla classe la teoria
di Agile e Scrum e che interviene di tanto in tanto per sottolineare o
completare un passaggio. Gli ospiti sono due:
\begin{itemize}
  \item \textbf{Simone Cicchetti}, ex studente dell'accademia, oggi level
        designer e producer, che parla per circa un'ora e mezza della teoria
        ``vista dalla vita reale'' e degli errori più comuni;
  \item \textbf{Cristian Liberatore}\footnote{Nella trascrizione il nome
        compare sia come ``Cristian'' sia come ``Christian''; qui si usa la
        prima forma, quella con cui viene presentato all'inizio della
        lezione. Anche i nomi dell'accademia sono resi in modi diversi
        (``Hive'', ``a dive'', ``IVE Online'', ``Live Online''): come negli
        altri capitoli si usa semplicemente ``l'accademia''; ``IVE Online''
        sembra indicare il corso a distanza dell'accademia.}, studente del
        terzo anno, game designer e project manager di due progetti nati
        come progetti d'anno, che segue passo per passo il ciclo di Scrum
        con gli esempi del suo gioco.
\end{itemize}
Diego li ringrazia per il tempo dedicato: Simone lavora, quindi sta
sacrificando tempo personale, e Cristian è uno studente con lo stesso carico
di lavoro della classe. Sono stati scelti perché, tra gli studenti avuti
fino a quel momento, sono quelli che hanno davvero \textbf{internalizzato}
Scrum e Agile, ottenendo grandi risultati produttivi nel progetto scolastico
e, nel caso di Simone, anche dopo l'accademia\footnote{La trascrizione dice
``nel caso di Christian post-accademia'', ma è Simone l'ex studente che oggi
lavora, mentre Cristian è ancora iscritto: il riferimento al dopo-accademia
è stato quindi attribuito a Simone.}.

Perché insistere sulla parola \emph{internalizzare}? Perché, dice Diego,
``utilizzare Scrum e Agile sembra molto facile, ma poi in realtà quando li
provi ad applicare veramente ti rendi conto che facili come sembravano quando
facciamo lezione non lo sono''. Se invece li si internalizza davvero, si
ottengono grandi risultati. L'idea della lezione è che possa essere utile
vedere l'esperienza di qualcuno \emph{molto più vicino} agli studenti di
quanto lo sia il docente, che ormai usa questi strumenti da anni. Il formato
annunciato è di circa un'ora e mezza per ciascun ospite, domande comprese.

Due avvertenze di lettura. La prima: i due interventi sono fatti di teoria
richiamata in breve e di moltissimi esempi concreti, spesso con numeri,
errori e ripensamenti; sono riportati con cura perché sono proprio la parte
che la teoria non dà. La seconda: i due relatori usano in modo
intercambiabile alcuni termini (\emph{framework} e metodologia,
\emph{post-mortem} e revisione della sprint, \emph{score}, peso e punteggio
dei task), e il capitolo segue il loro uso, segnalando in nota dove si
discosta dalla terminologia ufficiale di Scrum.

Il percorso del capitolo segue quello della lezione:
\begin{enumerate}
  \item un'apertura organizzativa e la presentazione degli ospiti
        (sezione~\ref{sec:z3-apertura});
  \item \textbf{Simone}: chi è, un sondaggio sulla classe e le prime
        esperienze di Scrum degli studenti (sezione~\ref{sec:z3-simone});
  \item il \textbf{mindset Agile}: perché non è un framework e perché va
        adattato (sezione~\ref{sec:z3-mindset});
  \item le \textbf{sprint} e le riunioni di pianificazione e
        \emph{post-mortem} (sezione~\ref{sec:z3-sprint});
  \item le \textbf{cerimonie} e la regola della \emph{consistency}, con le
        critiche al \emph{daily scrum} e al \emph{planning poker}
        (sezione~\ref{sec:z3-cerimonie});
  \item il \textbf{backlog} e gli obiettivi dei task
        (sezione~\ref{sec:z3-backlog-simone});
  \item la \textbf{velocity}: perché non usare le ore, la sequenza di
        Fibonacci e come la velocity permette di ``prevedere il futuro''
        (sezione~\ref{sec:z3-velocity});
  \item il riepilogo della creazione di una sprint e la conclusione sulla
        \textbf{fiducia} (sezione~\ref{sec:z3-ricapitolo});
  \item le domande a Simone, con la costruzione in diretta di un foglio per
        il calcolo della velocity (sezione~\ref{sec:z3-qa-simone});
  \item \textbf{Cristian}: chi è e come imposta la sua parte
        (sezione~\ref{sec:z3-cristian});
  \item il caso di studio, il gioco \emph{L'Inf}, e il ``trauma'' del cambio
        di team (sezione~\ref{sec:z3-linf});
  \item il \textbf{loop di Scrum} in cinque fasi
        (sezione~\ref{sec:z3-loop}): la stesura del backlog
        (sezione~\ref{sec:z3-backlog}), la selezione dei task
        (sezione~\ref{sec:z3-selezione}), l'implementazione con i suoi
        controlli, i \emph{blocker} e i feedback
        (sezione~\ref{sec:z3-implementazione}), il calcolo della velocity
        (sezione~\ref{sec:z3-calcolo}) e il \emph{tuning} della produzione,
        con il problema del perfezionismo (sezione~\ref{sec:z3-tuning});
  \item una domanda su come comunicare un taglio al team
        (sezione~\ref{sec:z3-qa-taglio});
  \item il \textbf{team management}: il team prima di Scrum, la produzione
        senza stress, il mindset condiviso, il \emph{post-mortem} e la
        \textbf{sincerità} (sezione~\ref{sec:z3-team});
  \item la chiusura della lezione (sezione~\ref{sec:z3-chiusura}).
\end{enumerate}

% ---------------------------------------------------------------------
\section{Apertura e presentazione degli ospiti}
\label{sec:z3-apertura}

\subsection{Una registrazione andata persa}

La lezione si apre con una comunicazione organizzativa\footnote{La
trascrizione comincia a frase già iniziata (``in una modalità non video fin
tanto che ce la chiamate poi la converte alla fine''): l'inizio della
spiegazione non è stato registrato e non è chiaro a quale lezione precedente
si riferisca.}. Il programma usato per registrare una lezione precedente
registra in una modalità che va convertita alla fine; durante la conversione
il computer di Diego, probabilmente per colpa degli aggiornamenti di Windows,
si è bloccato e la conversione è ``andata a pallino''. Poiché si tratta di
una lezione importante, Diego si rende disponibile a \textbf{rifarla
interamente} per registrarla di nuovo oppure, come la classe preferisce, a
fornire un audio, un video o qualunque altro formato sia più comodo, che poi
caricherà. Chiede agli studenti di pensarci e di fargli sapere.

\subsection{Chi sono gli ospiti e perché sono qui}

Poi Diego presenta gli ospiti speciali della giornata: Cristian Liberatore,
studente del terzo anno, e Simone Cicchetti, ex studente dell'accademia (che
in un primo momento ``sparisce'' dalla griglia della videochiamata, finché
Diego non lo ritrova in alto a destra). Le ragioni della scelta, il valore
della parola \emph{internalizzare} e il formato (un'ora e mezza ciascuno, con
domande) sono quelli già riassunti all'inizio del capitolo. Detto questo,
Diego si silenzia e lascia la parola a Simone.

% ---------------------------------------------------------------------
\section{Simone: il punto di vista del producer}
\label{sec:z3-simone}

\subsection{Chi è il relatore}

Simone si presenta mentre condivide lo schermo (``adesso tutto è imploso,
perfetto''). È uno degli ex studenti del terzo anno e oggi lavora come
\textbf{level designer}; collabora con Diego su alcuni progetti legati alla
gestione dei team, al \emph{game design}, al \emph{framework design} e ad
altre ``cose molto fighe'' che alcuni studenti forse conoscono\footnote{Il
passaggio è confuso nella trascrizione (``lavoro con diego alcune alcuni
progetti diciamo sul gestione di pro dei team game design framework
design''): la ricostruzione è approssimativa. Più avanti Simone dice anche di
lavorare per una società australiana.}. Sta inoltre continuando a lavorare al
progetto del suo terzo anno (il \emph{Biggest One}\footnote{Dalla lezione
si capisce che i progetti d'anno dell'accademia si chiamano \emph{Big One}
(primo anno), \emph{Bigger One} (secondo anno) e \emph{Biggest One} (terzo
anno); la trascrizione li scrive in vari modi (``biggest one'', ``bigger
one'', ``BiggerOne'').}), un gioco chiamato \emph{ISO}\footnote{Il nome del
gioco compare sempre come ``ISO''; non è possibile stabilire dalla sola
trascrizione se sia la grafia corretta (potrebbe essere, ad esempio,
\emph{Iso} o una sigla).}, per il quale è level designer, game designer,
social media manager e soprattutto \textbf{producer}: si occupa di gestire il
lavoro delle persone, di far rispettare le cerimonie di Scrum e, ``punto
centrale della questione'', di ``rompere le scatole alla gente''.

\subsection{Una lezione partecipativa}

Il tempo a disposizione è molto, quindi Simone vuole una lezione il più
possibile \textbf{partecipativa e cooperativa}: a differenza del solito, le
domande vanno fatte in qualsiasi momento, scrivendo in chat per attirare la
sua attenzione o alzando la ``manina'' con le emoji, che controlla su un
secondo monitor con la griglia dei partecipanti. Al massimo aspetterà la fine
della slide per non perdere il filo.

L'obiettivo dichiarato non è spiegare la teoria, che la classe ha già fatto
con Diego, ma darne un breve \emph{recap} e poi spiegare \textbf{a che cosa
serve nel pratico e com'è nella vita reale}: ``come spesso accade, la teoria
è una cosa, la vita reale è un'altra''.

\subsection{Un piccolo sondaggio}

Prima di partire, Simone chiede chi abbia ``un amore incommensurabile e
infinito'' per il producing, lo \emph{scrum mastering} e il management in
generale: nessuno, e uno studente, Gavi\footnote{Così nella trascrizione;
potrebbe trattarsi di un soprannome o di un nome storpiato.}, risponde
spammando il ``no''. Poi la domanda opposta: chi lo odia così tanto da
preferire ``affrontare un alligatore in una lotta greco-romana''? Ancora
nessuno. La classe è quindi, in generale, ``tutta intermedia''.

Chi si è già dovuto occupare di queste cose nei progetti d'anno o nei lavori
di gruppo? Rispondono due studenti, Mattia e Lorenzo\footnote{Nella
trascrizione ``da noi Mattia, Podini'' e ``da noi Lorenzoaldasso'': i cognomi
non sono ricostruibili con sicurezza (``Podini'' potrebbe essere il cognome
di Mattia o il nome di un altro studente).}, ognuno per il proprio team.
Simone chiede solo i punti salienti: che cosa è andato bene, che cosa no e
perché.

\paragraph{L'esperienza di Mattia.} Il suo gruppo usa Scrum da pochissimo,
circa una settimana. È lui a gestire la \emph{Scrum board}, con l'aiuto degli
altri due membri, ma il metodo non è ancora applicato a pieno. Mattia è
interessato a Scrum, ma la \textbf{responsabilità di guidare altre persone}
oltre a se stesso gli pesa parecchio, e usare Scrum in modo serrato tutti i
giorni gli fa ``una bella paura'': se qualche meccanismo si rompe,
bisognerebbe correre ai ripari. Che cosa intende per ``rompere un
meccanismo''? Che se una cerimonia non viene rispettata, tornare allo
\emph{status quo} non sarebbe semplice, anche per la sua mancanza di
esperienza nell'organizzare i team. Simone riassume: la preoccupazione è più
sul piano umano che su quello delle ore di lavoro necessarie. Mattia
riconosce che Scrum avrà le sue ``quasi magiche proprietà'', ma non avendolo
ancora visto all'opera nella vita reale non può dirlo con certezza.

\paragraph{L'esperienza di Lorenzo.} Niente di diverso: anche il suo gruppo
ha iniziato da pochissimo e proprio il giorno prima ha fatto la prima
riunione di Scrum, decidendo insieme i pesi dei task. Alla domanda se il
ruolo sia arrivato per voto popolare, Lorenzo risponde che si è offerto lui:
conoscere un \emph{framework} di lavoro, e soprattutto capire come
applicarlo, gli interessa parecchio.

Simone trova molto interessante quanto detto da Mattia sulle cerimonie e sul
doverle rispettare, e promette di tornarci più avanti
(sezione~\ref{sec:z3-cerimonie}).

% ---------------------------------------------------------------------
\section{Il mindset Agile}
\label{sec:z3-mindset}

\subsection{Agile non è un framework}

La prima slide parte da una frase che tutti conoscono: \textbf{Agile non è un
framework, è un mindset} (Simone dice di aver ``percepito la furia negli
occhi di Diego'' quando la classe esitava a rispondere). La frase sembra una
di quelle che servono ``per fare i fighi con i colleghi'', ma ha un
significato preciso: un \emph{framework} presuppone \textbf{regole da
seguire}, mentre Agile \textbf{non ha regole}. È una serie di \textbf{linee
guida da internalizzare}: la parola chiave è che devono diventare
\emph{interne}, al punto da non dover più pensare attivamente a quali siano
le regole, i punti o la mentalità di Agile. Bisogna capire la logica che c'è
sotto.

L'analogia è quella della \textbf{lingua straniera}: all'inizio si ragiona in
italiano, si costruisce la frase in italiano, la si traduce e la si dice; man
mano questi passaggi si annullano e si comincia a ragionare direttamente
nella nuova lingua. Con Agile funziona allo stesso modo.

\subsection{I quattro valori, e perché non vanno presi alla lettera}

Simone rilegge rapidamente i quattro valori del manifesto Agile:
\begin{enumerate}
  \item gli \textbf{individui e le interazioni} sono più importanti dei
        processi e degli strumenti;
  \item un \textbf{prodotto funzionante} è più importante di una
        documentazione esaustiva\footnote{Nella lezione ``documentazione
        comprensibile''; nella formulazione corrente del manifesto Agile il
        termine è \emph{comprehensive documentation}, cioè documentazione
        esaustiva.};
  \item la \textbf{collaborazione con il cliente} è più importante della
        negoziazione del contratto;
  \item \textbf{rispondere al cambiamento} è più importante di seguire un
        piano.
\end{enumerate}
Come è scritto sotto al manifesto stesso, questo non significa che le cose a
destra, quelle ``meno importanti'', non contino o non vadano seguite. Ma
Simone va oltre: \textbf{non significa nemmeno che queste regole debbano
essere seguite davvero}. Non è detto che individui e interazioni debbano
essere per forza più importanti di processi e strumenti. Il punto chiave di
Agile è \emph{sentire} quale sia il procedimento e la mentalità giusti e
\textbf{adattare Agile al proprio team}, perché ogni team e ogni progetto è
diverso.

Se un team e un progetto, per vari motivi, hanno bisogno di pagine e pagine
di documentazione, allora per loro la documentazione è più importante del
prodotto funzionante. Gli esempi sono infiniti; uno è il \textbf{gioco di
carte}: lì il prodotto funzionante è relativamente meno importante della
documentazione, perché le logiche si possono scrivere in qualsiasi momento,
ma senza una documentazione solida a supporto ``cadrà tutto quanto
immediatamente''. Quindi non bisogna seguire ``a menadito'' i punti del
manifesto, ma capirne la logica di fondo, capire \textbf{a che cosa serve} e
capire \textbf{se serve}: non è detto che Agile serva.

\subsection{Altri mindset e i framework}

Esistono molti altri \emph{mindset}. Di recente va molto di moda la
mentalità \textbf{lean}, che Simone definisce ``già il 2.0''\footnote{Così
nella lezione. Storicamente il \emph{lean} nasce nella produzione
industriale (il sistema Toyota) e precede Agile; nello sviluppo software il
\emph{lean software development} si è diffuso negli anni Duemila. Il ``2.0''
va letto come una battuta sul fatto che oggi è di moda, non come una
derivazione da Agile.}. L'importante, aggiunge con Diego, è \textbf{non usare
Waterfall}, ``perché Waterfall nel gaming\dots{} non dirò nient'altro''.

Sulla base di Agile sono state create delle \textbf{metodologie}, che
Simone si corregge chiamando \emph{framework}: e queste sì che hanno delle
regole. L'esempio è Scrum, ma ce ne sono infiniti, come Kanban. Si basano
tutti sul principio Agile dell'\textbf{inspect and adapt} (ispeziona e
adattati, analizza e adattati), che si può riassumere in una parola chiave:
\textbf{iterazioni}. Scrum è completamente basato sulle iterazioni, e anche lo
sviluppo dei videogiochi lo è, ``molto più dello sviluppo software''.

Il resto della sua parte è dedicato ai punti di Scrum, il \emph{framework} da
cui presumibilmente partono tutti gli studenti, e agli \textbf{errori
classici} in cui si incappa spesso. L'ideale sarebbe trasmettere la logica
alla base di tutto, ma Simone avverte che sarà difficile, perché ``sono cose
che si iniziano a capire mettendole in pratica''. Per questo invita di nuovo a
fare domande sulle proprie esperienze.

% ---------------------------------------------------------------------
\section{Le sprint e le riunioni}
\label{sec:z3-sprint}

\subsection{Il libro di Scrum}

Simone mostra il libro in cui viene raccontata la genesi del
\emph{framework} Scrum, scritto da Jeff Sutherland\footnote{Si tratta
verosimilmente di J.~Sutherland, \emph{Scrum: The Art of Doing Twice the Work
in Half the Time} (2014), edito in italiano come \emph{Fare il doppio in metà
tempo}. Il titolo non viene pronunciato: Simone mostra la copertina e più
avanti cita ``Scrum, di Jeff Sutherland''. La lunghezza indicata (``circa 300
pagine'') è approssimativa.}: un ``libricino'' di circa trecento pagine.
Avverte che userà le parole ``metodologia'' e ``framework'' in modo
intercambiabile. Lo consiglia assolutamente a chi è interessato al tema,
perché è scorrevole e contiene \emph{insight} molto interessanti; chi non è
particolarmente interessato può evitarlo del tutto, perché non avrà niente da
offrire che non venga detto in questa lezione.

\subsection{Che cos'è una sprint}

Scrum, a differenza di Agile, è un \emph{framework} a tutti gli effetti: ci
sono regole da seguire in modo preciso e da rispettare. Il modo in cui
abbraccia il principio dell'\emph{inspect and adapt} sono le
\textbf{sprint}: periodi di tempo tipicamente di \textbf{una o due
settimane}. Possono essere anche di quattro, come si studia, ma secondo Simone
non ha senso farle così lunghe a meno di non lavorare a un grande progetto
``con duecento miliardi di persone che ci stanno lavorando
costantemente''\footnote{Nella trascrizione ``a meno che non state facendo un
sesto plà con 200 miliardi di persone'': ``sesto plà'' è incomprensibile e
potrebbe essere una storpiatura di ``tripla A''; il senso, in ogni caso, è
quello di una produzione enorme.}.

Prima di ogni sprint c'è una riunione di \textbf{sprint planning}, dopo ogni
sprint una riunione di \textbf{post-mortem}\footnote{Nella guida ufficiale di
Scrum la riunione di fine sprint è divisa in \emph{sprint review} (che cosa è
stato prodotto) e \emph{sprint retrospective} (come ha lavorato il team);
qui, come nel resto della lezione, entrambe confluiscono nel termine
\emph{post-mortem}, che in senso stretto indica l'analisi a progetto
concluso (sezione~\ref{sec:z1-qa-fasi}).}.

\subsection{La regola d'oro: meno riunioni possibili}

``Le riunioni non piacciono a nessuno'': sono, dice Simone, una ``rottura di
scatole infinita'', e chiunque, soprattutto chi non è designer, non aspetta
altro che poter dire ``questa riunione è inutile, dobbiamo andarcene tutti a
lavorare perché abbiamo cose più importanti da fare''. In chat qualcuno
risponde con il meme di Bernie Sanders a braccia incrociate; Simone avrebbe
voluto mettere quello del cane che urla sulla scrivania perché ``tutta la
riunione poteva essere una mail'', e promette che dalla prossima volta ogni
slide avrà un meme alla fine.

Da qui la sua \textbf{regola d'oro}: \textbf{fare meno riunioni possibile}.
Sottolinea però che è una regola ``per quanto mi riguarda'': magari un team
non vede l'ora di fare riunioni, magari si portano le pizzette a ogni
incontro (``non so se con le pizzette le riunioni siano effettivamente
efficienti, però se funziona, perché no?'').

Una mossa molto semplice per ridurle è \textbf{unire sprint planning e
post-mortem} in un'unica riunione: si fa letteralmente la metà delle
riunioni. Ciò che si fa in quella riunione è:
\begin{itemize}
  \item \textbf{pianificare} la prossima sprint (\emph{sprint planning});
  \item \textbf{analizzare} la sprint precedente (\emph{post-mortem}),
        indicando in particolare i problemi incontrati e le cose rimaste
        bloccate.
\end{itemize}

\subsection{Una riunione deve essere rapida}

La riunione unica deve essere \textbf{rapida}. Durante l'ultimo anno, per il
\emph{Biggest One}, il team di Simone contava circa quindici persone, e alle
riunioni di planning e post-mortem ne partecipavano in media una dozzina.
Quando sono ``diventati bravi'', le riunioni duravano \textbf{venticinque
minuti, con dodici persone che parlavano tutte}. Quello è l'obiettivo, ``anzi,
il vostro obiettivo è fare anche meno''.

Il motivo per cui si sfora è una \textbf{tentazione} costante: quella di
\textbf{risolvere i problemi}. Quando qualcuno descrive un problema, un paio
di persone cominciano a discuterne (``hai provato a fare questo?'', ``no,
dovresti fare quest'altro''), e in un attimo sono passati dieci minuti in cui
hanno parlato solo in due senza che la riunione sia andata avanti di un
passo. In queste riunioni non si fa brainstorming e non si ragiona su come
risolvere i problemi: ``bloccatelo sul nascere, picchiate quelle persone, fate
qualsiasi cosa per impedirglielo''.

Quando c'è un problema, si \textbf{prende un appuntamento per risolverlo}.
Simone mostra come dovrebbe suonare un report:
\begin{quote}
``La settimana scorsa mi sono occupato della logica di movimento del
personaggio, ho fatto un po' di tuning e di testing e ho preparato il block
out del primo livello. L'unico problema è che il blueprint che faceva volare
via gli uccelli non mi ha funzionato bene: visto che tu, Cristian, l'hai
utilizzato in un altro livello, più tardi ci sentiamo su Discord quindici
minuti e mi spieghi come l'hai settato. Questa settimana mi occuperò di
completare il block out e di risolvere il problema del blueprint.''
\end{quote}
Fine del report: \textbf{circa un minuto}, e c'è tutto. Tutti devono fare
così: trovare una persona con cui risolvere il problema, identificare
\emph{come} risolverlo e andare oltre. ``C'è tempo e luogo per ogni cosa'', ma
la riunione di sprint planning non è il momento giusto.

\subsection{La sprint come unità di tempo}

Ultimo punto: le sprint sono \textbf{l'unità di tempo} fondamentale per ogni
cosa. Le \emph{milestone} non si misurano in mesi, giorni o settimane, ma
\textbf{in sprint}. Bisogna prendere l'abitudine di dire ``completeremo la
preproduzione in otto sprint'' e non ``completeremo la preproduzione tra due
mesi''. È importantissimo per il calcolo della velocity
(sezione~\ref{sec:z3-velocity}).

% ---------------------------------------------------------------------
\section{Le cerimonie e la regola della consistency}
\label{sec:z3-cerimonie}

\subsection{Le cerimonie di Scrum secondo Simone}

Le regole di Scrum si chiamano \textbf{cerimonie}. Simone le divide in due
gruppi\footnote{La classificazione è quella proposta da Simone. Nella guida
ufficiale di Scrum gli ``eventi'' sono la sprint, lo \emph{sprint planning},
il \emph{daily scrum}, la \emph{sprint review} e la \emph{sprint
retrospective}; il backlog è un ``artefatto'' e il calcolo della velocity
e il \emph{planning poker} sono pratiche diffuse ma non obbligatorie.}:
\begin{center}
  \begin{tabular}{@{}p{0.42\textwidth}p{0.48\textwidth}@{}}
    \toprule
    \textbf{Cerimonie principali} & \textbf{Altre cerimonie} \\
    \midrule
    planning; post-mortem; calcolo della velocity; creazione del backlog &
    \emph{daily scrum}; \emph{poker time estimation} (\emph{planning poker})\\
    \bottomrule
  \end{tabular}
\end{center}
Le quattro principali sono quelle che ``non potete evitare di fare,
altrimenti non state facendo Scrum'': il che, precisa, non è necessariamente
un problema (ci torna alla fine, sezione~\ref{sec:z3-ricapitolo}). Le altre
due le tiene separate perché a lui non piacciono granché, e prima di spiegare
perché beve un sorso d'acqua ``per lasciarvi la suspense''.

\subsection{Consistency: una cerimonia saltata fa cadere le altre}

In Scrum c'è una parola chiave da ricordare: \textbf{consistency}. Bisogna
evitare a tutti i costi di creare \textbf{precedenti di
\emph{inconsistency}}, cioè di incostanza: con Scrum bisogna rimanere
costanti. \textbf{Se si salta una sola cerimonia si crea un precedente che,
a effetto domino, butta giù tutto il resto.} Questo è il punto già sollevato
da Mattia. Ma, avverte Simone, questo \emph{non} vuol dire che ci si debba
forzare: ed è questo il punto chiave.

\subsection{Perché il daily scrum non funziona più}

Il \emph{daily scrum} era pensato ``per il millennio scorso''. Scrum è nato
molto tempo fa\footnote{Simone non ricorda la data esatta (``fine degli anni
Novanta, forse fine degli anni Ottanta''). Il termine \emph{scrum} applicato
allo sviluppo di prodotti risale a un articolo di Takeuchi e Nonaka del 1986;
il \emph{framework} Scrum fu presentato da Jeff Sutherland e Ken Schwaber nel
1995.}, per team che si incontravano tutti i giorni in ufficio alla stessa
ora, con un orario d'ufficio standard: team di sviluppo software, quindi per
lo più programmatori (anche i più ``artistici'', come gli UX designer, erano
alla fine programmatori), con una mentalità piuttosto simile e, soprattutto,
con lo stesso orario.

Oggi, Covid a parte, questo non è più né fattibile né sensato. Gli studenti
non hanno nemmeno un ufficio: si incontravano a volte in sede quando si
poteva, ma per lo più lavorano online, e \textbf{non tutti lavorano allo
stesso orario}. ``E va bene così'': c'è chi lavora di notte (la salute
potrebbe risentirne, ma se si sentono bene così, va bene), e quelle persone
non scriveranno il proprio \emph{daily scrum} all'orario in cui lo fanno
tutti. Senza arrivare a casi del genere, Simone lavora per una
\textbf{società australiana}, dal lato opposto del globo: quando lui si
sveglia, loro stanno finendo di lavorare, e il suo \emph{daily scrum} cadrebbe
nella loro notte.

In queste condizioni è difficile tenere traccia del \emph{daily scrum} e
rispettarlo, e prima o poi qualcuno lo salterà. Quello è un
\textbf{precedente di incostanza}: qualcuno comincia a dire ``ok, non lo faccio
neanche io'', poi ``ma visto che non faccio più gli scrum, a che servono le
riunioni? Questa volta non vengo'', ``e io neanche'': ``e da lì è il caos''.

La conclusione è netta: \textbf{quando qualcosa non funziona, tagliatelo}. Il
senso di ``individui e interazioni sono più importanti di processi e
strumenti'' è proprio questo: se un processo o uno strumento non funziona
\textbf{non è colpa vostra, è il team che non ne ha bisogno}. Non bisogna
fissarsi su ciò che si vuole fare, ma sentire di che cosa ha bisogno il team e
adattarsi a quel bisogno.

\subsection{Perché il planning poker rallenta tutto}

Il \emph{planning poker} (che Simone chiama \emph{poker time estimation})
non gli piace perché lo trova \textbf{estremamente lungo e tedioso}: rallenta
un processo già lento di suo. Creare il backlog e valutare i task è già
lunghissimo (``ci vuole almeno un'intera sprint di riunioni per fare un
backlog''); aggiungerci il planning poker significa moltiplicare il tempo. Si
informa con Diego su come sia stato spiegato alla classe: con il metodo di
Fibonacci, ``\emph{power-uppato}''\footnote{Nella trascrizione ``è rimasta
nelle lezioni di trial'' e ``è stata power upata'': il senso sembra essere che
il planning poker è stato spiegato nelle lezioni precedenti in una versione
potenziata con la sequenza di Fibonacci. ``Trial'' potrebbe essere una
storpiatura.}, il che farà risparmiare un po' di spiegazioni più avanti.

Simone capisce la necessità di sentire le opinioni di tutti
(``\emph{empowering the team}'', eccetera), ma l'effetto è rendere il processo
ancora più lento. E poiché le riunioni non piacciono a nessuno, se si
quintuplica il tempo necessario ``si parte subito male'', perché il team si
annoia già. Anche questo, quindi, lui personalmente non lo usa.

\subsection{Il punto chiave: crearsi il proprio framework}

La cosa più importante è che \textbf{non bisogna eseguire tutto a menadito}.
I \emph{framework} sono stati pensati per essere generici, ma anche per
essere \textbf{adattati}. La procedura che Simone suggerisce:
\begin{enumerate}
  \item si parte \textbf{rispettando il framework al 100\%};
  \item quando qualcosa non funziona, si capisce \emph{perché} non funziona;
  \item si smette di farlo, si risolve, si fa altro;
  \item si capisce qual è il framework più adatto e \textbf{ci si crea il
        proprio}.
\end{enumerate}
``Questo è il punto chiave di tutto questo.''

% ---------------------------------------------------------------------
\section{Il backlog}
\label{sec:z3-backlog-simone}

\subsection{Una ``rottura di scatole'' inevitabile}

Sulla creazione del backlog Simone spende ``giusto due parole'', perché la
parte è molto tecnica e i consigli sono pochi. Creare il backlog è faticoso,
``non ci sono mezze misure per dirlo'': richiede moltissimo tempo e lavoro ed
è intenso a livello mentale, perché bisogna ragionare tantissimo.

\subsection{Il backlog a scaglioni}

In teoria i backlog si creano all'inizio dello sviluppo: in un mondo perfetto
si conoscerebbero fin da subito tutti i task necessari per completare il
progetto. Nel mondo reale questo non accade. Ciò che accadrà, più
probabilmente, è che si andrà avanti \textbf{per milestone}. È ciò che Simone
ha cominciato a fare con \emph{ISO} e con cui si trova bene:
\begin{itemize}
  \item ha fatto il backlog della \textbf{preproduzione}, con tutto ciò che
        serviva per la preproduzione;
  \item finita la preproduzione, ha fatto quello della \textbf{produzione};
  \item finita la produzione farà quello della fase successiva
        (``espansione, build finale, non lo so''), a seconda di come si
        svilupperanno le cose.
\end{itemize}
Fare i backlog \textbf{a scaglioni} li rende più semplici e il processo meno
lungo; e tanto, comunque, bisognerà aggiungere cose in corso d'opera: ``non si
scappa''.

\subsection{Le epic}

L'organizzazione in \textbf{epic} è molto utile, soprattutto per la stesura:
si comincia ragionando per \emph{epic} e poi si scende ai task. Un esempio di
epic è ``creazione del personaggio''. Che cosa contiene il personaggio?
Movimento, tuning, modello, animazioni, materiale, il design delle feature
che dovrà avere, la creazione di tutte queste feature, e così via. Il
processo diventa più \emph{streamlined}, meno pesante e più semplice da
seguire.

\subsection{Ogni task deve avere un obiettivo chiaro}

Fondamentale: \textbf{i task devono avere un obiettivo chiaro}. Un consiglio
comune è scriverli come \textbf{user story}: ``come utente voglio avere un
personaggio che si senta potente, che io riesca a percepire come potente''.
Aiuta moltissimo nello sviluppo software e parecchio nei videogiochi per
alcune cose, come appunto il personaggio. In generale è buona pratica
\textbf{spiegare l'obiettivo e non la soluzione}: dire che il personaggio deve
muoversi, saltare e fare le capriole, invece di dire che ``quando premo il
tasto W deve subire una forza che lo spinge in avanti''. Però non è
fondamentale.

Ciò che è davvero importante è che il task abbia un \textbf{obiettivo
verificabile}: qualcosa di cui si possa dire ``è completo'' o ``non è
completo''. Altrimenti le persone \textbf{non lo completeranno}, o lo
completeranno in modo diverso da come lo si immaginava. Un esempio di task
senza obiettivo chiaro è \textbf{``completare il livello~5''}:
\begin{itemize}
  \item per Simone vuol dire aver creato il block out con qualche elemento
        estetico per sapere dove vanno le cose e come funzionano con i
        moduli già disponibili, e basta;
  \item per qualcun altro potrebbe voler dire inserire illuminazione,
        \emph{post-processing}, \emph{particles}, tutti gli shader fatti
        bene, ottimizzazione e \emph{level streaming};
  \item per altri ancora potrebbe voler dire arrivare all'illuminazione,
        lasciando \emph{level streaming} e ottimizzazione a una seconda fase.
\end{itemize}
È la cosa più importante da ricordare quando si crea il backlog e si
scrivono i task.

% ---------------------------------------------------------------------
\section{La velocity}
\label{sec:z3-velocity}

``Bellissimo, bellissimo'': il calcolo della velocity è per Simone ``la cosa
più figa e allo stesso tempo più odiata''.

\subsection{Il peso dei task}

Come la classe sa, i task sono valutati in base al loro \textbf{peso} (o
punteggio, o \emph{score}). Il peso dipende da quanto un task è difficile e/o
da quanto tempo richiede: è un po' una zona grigia, e Simone consiglia di
valutarlo come \textbf{combinazione delle due cose}.

\subsection{Mai le ore}

Il consiglio fondamentale, che Simone non capisce perché si continui a
contraddire, è \textbf{non usare le ore per valutare i task}, ``per niente,
neanche alla lontana''. Nel libro di Sutherland c'è un paragrafo intero che
spiega perché. Diego commenta che chi lo consiglia di solito fa ``Agile
Waterfall'' e non ha nemmeno letto i libri su Agile.

Il motivo: \textbf{il cervello umano è completamente incapace di stimare le
ore}, cioè di prevedere il tempo necessario per fare qualcosa: ``non siamo
pensati per questo''. Siamo invece \textbf{molto bravi con le misure
relative}: fissato un task che vale 10, sappiamo dire bene quale vale 5,
quale 7, quale 2. Sappiamo dire che un task è molto più grande di un altro.

\subsection{La sequenza di Fibonacci}

Per aiutarsi ulteriormente, Simone usa la \textbf{sequenza di Fibonacci} (1,
2, 3, 5, 8, 13, 21, 34, \dots), anche questo un consiglio di Sutherland che
``nessuno segue''. Sutherland racconta addirittura di un team che usava le
\textbf{taglie dei cani} (``questo task è un pinscher, questo è un golden
retriever'')\footnote{Nel libro di Sutherland le stime per razze canine sono
citate come esempio di scala relativa alternativa ai numeri; nella
trascrizione ``pincher''.}; a Simone sembra eccessivo, e ``i buoni vecchi
numeri'' bastano. Fibonacci ha due scopi molto utili:
\begin{enumerate}
  \item \textbf{Gli intervalli crescenti assorbono l'errore di stima.} Man
        mano che un task diventa più grande e complesso, cresce anche
        l'errore associato. È facile dire che aprire la mail vale 2 e
        rispondere a tutti sul server Slack vale 3: che cosa può andare
        storto? La mail non si apre e bisogna reinserire la password; è
        esploso qualcosa e ci sono cinquanta messaggi invece dei soliti
        venti. Niente che cambi granché. È molto più difficile, invece,
        valutare la differenza tra creare il \emph{behavior} di movimento del
        personaggio e costruire un intero livello: possono succedere
        infinite cose, soprattutto quando c'è di mezzo la programmazione.
        Il behavior del personaggio è un 20 o un 21? ``È un 21, punto'':
        non è un 34, perché non è così difficile, e non è un 13, perché
        sarebbe troppo poco. Magari la cosa più corretta sarebbe stata 23,
        ma non importa: lo scopo è \textbf{appianare} i risultati.
  \item \textbf{Meno scelta, meno discussione.} I valori sono
        ``infinitamente arbitrari'', ``completamente, al 100\%'', e nessun
        altro al mondo li condividerà; ma Fibonacci li rende un po' più
        condivisibili, perché con meno possibilità di scelta ci sono meno
        occasioni di discutere.
\end{enumerate}
Che i valori siano arbitrari non è un problema: finché \textbf{è una persona
sola a darli}, sono tutti arbitrari allo stesso modo, quindi
\textbf{consistenti}; e se sono consistenti hanno una logica e si possono
usare per fare calcoli.

\subsection{Che cos'è la velocity e perché serve}

Sommando il peso dei task completati in una sprint si ottiene la
\textbf{velocity}. Ed è ``fighissima'' perché permette, in sostanza, di
\textbf{prevedere il futuro}.

Simone mostra il grafico della velocity della preproduzione di \emph{ISO}
durante l'anno scolastico. La consegna era a maggio (``10 maggio, o poco
dopo''). Già a inizio marzo, guardando l'andamento, aveva capito che \textbf{non
sarebbero riusciti a finire in tempo}. Il ragionamento era semplice: mancano
dodici sprint; quanto facciamo ogni sprint? Circa 89; quanto manca nel
backlog? Circa 2000. ``Non ce la facciamo''\footnote{I numeri sono quelli
detti a voce: 12 sprint per circa 89 punti danno poco più di 1000 punti,
circa la metà dei 2000 rimasti nel backlog.}.

La velocity aiuta a prevedere il futuro e quindi a \textbf{prendere
provvedimenti prima}. A marzo Simone ha cambiato i piani: prepariamo meno
cose, smettiamo di perdere tempo su ciò che ci blocca, consegniamo qualcosa di
diverso. In più, la velocity del team è \textbf{aumentata} sempre di più, e
alla fine sono arrivati in fondo alla preproduzione \textbf{senza crunch}:
``zero, completamente zero crunch''. Nessuno ha dovuto lavorare di più;
l'ultima settimana non hanno nemmeno toccato il progetto: hanno fatto la
build, l'hanno spedita e per una settimana non ci hanno lavorato. E sono
riusciti a consegnare anche di più del previsto, perché Simone ha
\textbf{analizzato i trend precedenti}: perché in un certo punto c'era un
picco? Che cosa li bloccava e faceva diminuire la produttività? L'ha
scoperto, l'ha capito e hanno agito.

\subsection{Come calcolarla in pratica}

Il consiglio è di \textbf{trovare un modo per fare il calcolo}. Simone usava
Monday, che non aveva uno strumento costruito apposta, ma ci è riuscito lo
stesso; si può fare tranquillamente con \textbf{Google Sheets}, ``col
sorriso''. Per farlo davvero bene bisognerebbe scrivere un po' di script
dentro il foglio, ``niente di improponibile''. Quando comincia a girare
``funziona da dio'': ``assolutamente fatelo, non posso dirlo abbastanza''. (Un
esempio costruito in diretta è nella sezione~\ref{sec:z3-qa-simone}.)

\subsection{Come assegnare i punteggi}

Per assegnare gli \emph{score} conviene partire dagli
\textbf{estremi}: quali task sono un 34? Quali sono un 1? Da lì si mettono i
valori intermedi. Gli estremi fanno da \textbf{ancore} e facilitano il
processo di valutazione \emph{relativa} descritto sopra.

Lo stesso vale per le \textbf{priorità}, se si usa un sistema di priorità:
regole simili, ma con \textbf{intervalli più piccoli} e senza Fibonacci, che
lì non serve. Simone usa una scala da 1 a 5; ma non la considera così
fondamentale.

\subsection{Che cosa blocca che cosa}

Durante il lavoro sul backlog bisogna capire \textbf{che cosa blocca che
cosa}. Ad esempio: ci sono 29 task, per un valore complessivo di 400, tutti
relativi alla creazione dei livelli. Ma se non c'è il movimento del
personaggio, non ha senso fare i livelli, perché vanno costruiti attorno a
quel movimento: la creazione del personaggio \textbf{blocca} i livelli e
quindi ha priorità più alta. Da qui si comincia anche a capire quale reparto
si occupa di che cosa: probabilmente all'inizio i programmatori faranno i
\emph{behavior} fondamentali (come il personaggio), mentre i designer
lavoreranno su documentazione e ricerca. Ma dipende dal gioco: alcuni non
costruiscono i livelli attorno al personaggio, bensì il personaggio attorno ai
livelli, quindi anche l'esempio ``è opinabile''.

% ---------------------------------------------------------------------
\section{Ricapitolando: come nasce una sprint}
\label{sec:z3-ricapitolo}

Prima di passare alle domande, perché ``quello che avevo da dirvi di pratico è
quasi finito'', Simone ricorda che il grosso non sono le regole, ma
\textbf{capire la logica alla base e integrarla nel proprio team e nel
proprio progetto}. Poi riassume la creazione di una sprint in quattro passi.

\begin{enumerate}
  \item \textbf{Post-mortem della sprint precedente.} Si vede che cosa ha
        fatto ciascuno, si capisce che cosa è bloccato e si
        \emph{programma un incontro} per risolvere i problemi.
  \item \textbf{Pianificazione.} Dal backlog si guarda che cosa blocca che
        cosa, quali sono le priorità e che cosa si può fare in questa sprint,
        e si prendono i task.
  \item \textbf{Analizzare i problemi e agire.} Si guarda la velocity: è in
        aumento? In calo? Si è in un picco? L'andamento a dente di sega,
        ``su, giù, su, giù'', è \textbf{assolutamente normale}, ed è ciò che
        ci si deve aspettare; e man mano che il progetto va avanti la
        velocity \textbf{tenderà a diminuire}. Le domande da farsi: come sta
        andando il progetto? Chi sta facendo di più e chi di meno, e perché?
        Qualcuno ha avuto un calo di rendimento? C'è un problema che si
        ripresenta spesso, un collo di bottiglia? E come risolverlo,
        soprattutto \textbf{ascoltando il resto del team}.
  \item \textbf{Rimanere consistenti}, in particolare con Scrum.
\end{enumerate}

\subsection{Ascoltare il team}

Come scrum master o producer, il lavoro più importante è e rimane
\textbf{ascoltare il team}. Tutto il resto non conta se poi, quando il team
dice una cosa, non lo si ascolta. È lo stesso principio dell'adattare il
\emph{framework} al team:
\begin{itemize}
  \item se il team odia il \emph{daily scrum}, non si fa: non serve a niente;
  \item se il team fa fatica a prendere i task dal backlog, non glielo si fa
        fare: si assegnano i task in modo diverso, ``non fategli manco vedere
        il backlog'';
  \item se il team vuole sprint di due settimane invece che di una, si fanno
        di due settimane.
\end{itemize}
Magari le persone si organizzano personalmente in modo diverso: Simone, per
esempio, organizza tutto il suo lavoro, su tutti i fronti, \textbf{su base
settimanale}, e una sprint di due settimane gli sconvolgerebbe i piani. Se il
team lavora per lo più così, si fa così. ``Adattatevi'' non vuol dire solo
adattarsi ai bisogni del software, ma \textbf{soprattutto ai bisogni del
team}.

\subsection{Nessun framework va bene al 100\%}

Esistono infiniti \emph{framework}, e \textbf{nessuno andrà bene al 100\% per
un team}. Bisogna provarli, il più possibile, e adattarli: cambiare qualche
regola, togliere qualche cerimonia. Il \emph{framework} più usato oggi, ad
esempio, è \textbf{Kanban}, che non ha nemmeno le sprint\footnote{Simone dice
prima ``Kanban neanche ce l'ha il backlog'' e poi si corregge: ``scusatemi,
neanche ce l'ha le sprint, ha soltanto un gigantesco backlog''.}: ha soltanto
un gigantesco backlog. ``Ogni team è diverso, ogni team ha bisogni diversi,
quindi adattatevi'' (una frase che Simone calcola di aver ripetuto ``almeno 49
volte'').

\subsection{La nota ``smielata'': la fiducia}

La conclusione: alla fine di tutto, \textbf{Agile non si basa sulle
iterazioni}, né sull'efficienza, né sulla creazione del miglior videogioco o
software del mondo, né sul controllo. L'unica cosa su cui si basa davvero è
\textbf{la fiducia}:
\begin{itemize}
  \item la fiducia che il team ripone in chi lo guida, e che chi lo guida
        ripone nel team: ``queste persone lavoreranno nel miglior modo
        possibile, sempre e comunque'';
  \item la fiducia del team nel dire ``se qualcosa va male, so che lo scrum
        master ha un piano di riserva che ci risolverà il problema, e il mio
        lavoro non sarà sprecato'';
  \item la fiducia di chi guida nel dire ``non posso controllarvi
        costantemente, non posso assicurarmi che facciate sempre tutto il
        lavoro o che seguiate le cerimonie come dovreste, però mi fido'': mi
        fido che lo farete, che se ci saranno problemi ve lo direte, che
        farete sempre del vostro meglio.
\end{itemize}
Bisogna \textbf{costruire fiducia nel team e trasmettere fiducia al team}. È
lo stesso principio del \emph{default to trust} della prima lezione zero
(sezione~\ref{sec:z1-default-trust}).

Simone ha finito con mezz'ora d'anticipo e apre alle domande (``go crazy'').
Diego apprezza molto il finale, ``smielato ma in realtà estremamente
veritiero'': lui è più orientato all'ambito professionale, ma crede che in
Agile ci debba essere davvero questa fiducia, e che sia anche una delle sue
forze. Simone pensava che non l'avrebbe apprezzato, tanto era smielato
(``mentre lo scrivevo ho percepito il diabete che aumentava''); Diego gli
suggerisce, la prossima volta, di aggiungere immagini con gli occhi a cuore e
qualche gattino.

% ---------------------------------------------------------------------
\section{Domande a Simone}
\label{sec:z3-qa-simone}

\subsection{La dimensione del team allunga le riunioni?}

Uno studente, che lavora in un team di tre persone (``essendo in tre le cose
vengono un po' più facili''), chiede quanto la \textbf{grandezza del team}
influisca sui tempi delle riunioni e delle cerimonie, e se ci sia un modo per
mantenerle rapide comunque.

Influisce assolutamente, ``per una mera questione matematica''. Ma in genere
è possibile far sì che un team di qualsiasi dimensione completi una riunione
in \textbf{una trentina di minuti al massimo}: quello è l'obiettivo, ``anche a
costo di dire meno cose, anche a costo di essere meno precisi''. A meno di
avere un team di una cinquantina di persone, la riunione durerà sempre trenta
minuti; e con cinquanta persone non si fa comunque una riunione tutti insieme,
perché è impossibile: si fanno le \textbf{riunioni di reparto}, e poi i
\textbf{capi reparto} fanno la loro riunione di sprint planning. I numeri
tornano piccoli: a una riunione ci saranno al massimo una quindicina di
persone, venti ``se proprio i pianeti si sono allineati'', e venti persone
possono stare in trenta minuti.

\subsection{Quanto deve imporsi lo scrum master? Il timekeeper}

Lo stesso studente chiede quanto lo scrum master debba \textbf{imporsi} sugli
altri, sul piano umano, per far rispettare i tempi. Moltissimo, risponde
Simone, ma \textbf{non deve per forza essere lo scrum master a farlo}. Nel
team australiano in cui lavora, all'inizio di ogni riunione si sceglie, in
modo pseudo-casuale, tra chi è disponibile o lanciando un dado, un
\textbf{timekeeper}: la persona che ha ``lo scettro del comando'' e tiene in
mano un cronometro, perché nessuno dedichi ai propri report più tempo del
dovuto. È \textbf{comunemente accettato da tutti} che quella persona abbia il
diritto di ``rompere le scatole'': è socialmente accettato.

Se le riunioni si dilungano troppo, si può proporre di introdurre la regola
del timekeeper, scelto ogni volta a caso o sempre lo stesso: ha il diritto di
richiamare gli altri perché, oggettivamente, il tempo scorre.

\subsection{E se il problema blocca tutti?}

Ultima domanda dello stesso studente: si è detto che i problemi si risolvono
dopo la riunione; ma se un problema è molto grave, se il fatto di non averlo
risolto blocca la produzione di più persone (``praticamente un disastro''),
vale la pena fermarsi e parlarne tutti insieme, o si dedica comunque quel
tempo solo ai report e poi si fa una riunione in più?

La risposta semplice è \textbf{no}: non si risolve nella riunione di sprint
planning e post-mortem. ``L'unica cosa che mi viene in mente che possa essere
così tanto grave è se la stanza in cui state facendo il meeting va a fuoco.''
La riunione non è pensata per quello. Subito dopo, piuttosto, si avvia
un'altra riunione con tutte le persone necessarie (``che durerà 25 ore'') per
risolvere il problema. Ma finché si fa planning e post-mortem, si fa
\textbf{quello e basta}, anche se il problema riguarda tutti.

\subsection{E se mancano delle persone?}

Si fa comunque la riunione? La si rimanda? Simone la fa \textbf{sempre}, a
meno di essere da solo, per \textbf{rispetto delle persone che sono venute}.
Le assenze possono avere cause diverse. Chi non è venuto deve mandargli un
messaggio su Telegram, anche vocale, in cui fa di fatto il proprio report.

Se una persona manca spesso, Simone ci parla e cerca di capire perché.
Magari l'orario della riunione non va bene perché ha altro da fare: nel
gruppo di \emph{ISO} c'è un ragazzo che lavora a settimane alterne di mattina
e di pomeriggio, e quando lavora di pomeriggio, tornando a casa, trova spesso
traffico; quindi a settimane alterne non può venire, o arriva molto in
ritardo. ``Lo so, me lo posso aspettare, non è un problema.'' Se invece
qualcuno manca in modo pseudo-casuale, bisogna capire il motivo: non gli va
bene l'orario? Ne preferirebbe un altro? Pensa che le riunioni non servano?
Non ha voglia di farle? Da lì si capisce il problema e si cerca di
risolverlo, spiegando che \textbf{le riunioni sono importanti} e che è
importante esserci tutti, e chiedendo ``dimmi come possiamo aiutarti a
rendertele meno pesanti''.

\subsection{Quando è la persona a non adattarsi}

Se però, nonostante tutto, la persona ``continua a fare la \emph{drama
queen}'', a un certo punto bisogna ammettere che esistono anche persone non
adatte a lavorare in un certo contesto, e si può valutare se lo sia.

Diego interviene per completare: in Agile tutti vogliono ottenere la cosa
migliore per il team e per il metodo di lavoro di ciascuno. Ma se un individuo
non si adatta al team, o il team non riesce ad adattarsi a lui, che sia o
meno una \emph{drama queen}, \textbf{la presenza di quella persona nel team va
valutata}. ``Nessuno, a parte casi eccezionali, tipo Miyamoto, vale il lavoro
degli altri.'' Le \emph{soft skill} restano alla base. E aggiunge che non è
solo questione di soft skill: ci sarà la \emph{drama queen}, quello che non si
adatta, ``lo stronzo o la stronza'' (il genere, precisa, è irrilevante),
quello che non ha voglia\footnote{Nella trascrizione ``stronzo
stronza,umbaquin'': la parte finale è incomprensibile.}; ma c'è anche la
persona che ha voglia di fare ma non si adatta al team. Se non ci riesce, o il
team non riesce ad adattarsi a lei, potenzialmente deve andare via.

Simone conferma che è un punto molto importante, che rientra nel discorso
dell'\textbf{ascoltare di che cosa ha bisogno il team}: se il team non riesce
a lavorare con un individuo, ``non c'è niente di male, è così e basta''. Per
fortuna non è troppo comune: a un certo livello professionale è difficile che
succeda, soprattutto in questa industria, dove le persone, se non stessero
lavorando insieme, potrebbero diventare amiche, perché hanno un interesse in
comune e sono molto appassionate. Ma può succedere.

\subsection{Come si calcola la velocity, in pratica}

Diego invita gli studenti ad approfittare di una persona ``che ha fatto il
vostro percorso''. Una domanda sul calcolo: la velocity è semplicemente la
\textbf{somma dei pesi dei task completati} nel tempo trascorso, cioè nella
sprint? Esatto. Si fa la sommatoria dei task \textbf{taggati come
\emph{done}} che hanno la stessa sprint (o la stessa data di consegna, a
seconda di come la si indica): si prendono tutti i task con ``sprint 5'', di
questi quelli in stato \emph{done}, e se ne sommano i pesi. Allo stesso modo
si possono sommare i task non completati, quelli rimandati, e così via.

\paragraph{Uno sguardo ai fogli personali.} Simone vorrebbe mostrare il foglio
di \emph{ISO}, ma è difficile da leggere per la quantità di dati; si chiede se
mostrare i propri fogli rompa qualche NDA, poi nota che sta già condividendo
lo schermo, ``quindi gli NDA sono già stati rotti''. Mostra allora uno dei
suoi fogli personali: organizza il proprio lavoro \textbf{settimana per
settimana}, con il nome della settimana, l'elenco dei task, la somma di ciò
che era da fare, quanto è stato effettivamente fatto, quanto è stato tagliato
e quanto rimandato per cause diverse. Il risultato è un grafico con
l'andamento settimanale di questi valori.

\paragraph{La ``speedrun velocity''.} Poi decide di fare ``una cosa più
figa'' e di costruire un foglio in diretta, che chiama \emph{speedrun
velocity}. La struttura minima\footnote{I riferimenti alle celle detti a voce
durante la costruzione (``B2B'', ``G1'', ``F1'', ``C2, F2'') sono confusi
nella trascrizione; la tabella e la formula sono una ricostruzione dello
schema, non una copia esatta del foglio mostrato.}:
\begin{center}
  \begin{tabular}{@{}llll@{}}
    \toprule
    \textbf{Task name} & \textbf{Score} & \textbf{Sprint} & \textbf{Status} \\
    \midrule
    Task 1 & 5  & Sprint 1 & done \\
    Task 2 & 13 & Sprint 1 & WIP \\
    Task 3 & 3  & Sprint 2 & cut \\
    \dots  & \dots & \dots & \dots \\
    \bottomrule
  \end{tabular}
\end{center}
Nella colonna \emph{Status} inserisce una convalida dei dati con i tre valori
\emph{cut}, \emph{WIP} e \emph{done}, ``per semplificarsi la vita''. Poi
scrive una serie di task con punteggi a caso (copiati da un altro foglio) e
assegna le sprint, avvertendo che ``ovviamente non dovrebbe essere fatto
esattamente così, ma \emph{for the sake of showcase}''; aggiunge qualche
colore, perché ``sono pazzo'', anche se non è la cosa fondamentale.

Il calcolo è una somma condizionale su più criteri (\texttt{SUMIFS}, in
italiano \texttt{SOMMA.PIÙ.SE}): si somma la colonna dei punteggi per le
righe in cui la sprint è quella cercata \emph{e} lo stato è \emph{done}. Ad
esempio:
\begin{center}
  \texttt{=SUMIFS(B:B; C:C; "Sprint 1"; D:D; "done")}
\end{center}
Lo stesso si può fare con \emph{cut} per sapere quanto è stato tagliato. Si
ripete per la sprint~2, e così via, e con i risultati si inserisce un
\textbf{grafico a linee} (Sheets propone il tipo di grafico sbagliato, e
Simone lo corregge). Con più sprint si ottiene un andamento come quello dei
suoi fogli personali o di \emph{ISO}. Chi vuole fare le cose per bene può
aggiungere altre analisi: chi ha fatto più task, chi ha più peso e carico di
lavoro, lo stato generale del progetto (``magari il backlog fosse così
vuoto'').

Era, dice, per essere sicuro di riuscirci; è disponibile a rifare qualche
passaggio, ma alla classe è tutto chiaro. Uno studente chiede se il foglio si
possa avere: Simone lo mette in una cartella e sarà Diego a farlo avere alla
classe.

\subsection{I problemi più gravi incontrati con Agile}

Diego osserva che Simone ha capito come applicare i principi di Agile con una
velocità sorprendente, più di quanto sia successo a lui quando gli fu
spiegato, e che all'inizio si fa spesso ``Waterfall Agile'', cioè Waterfall
con le sprint, credendo di fare Agile\footnote{Il passaggio è frammentato e non
è del tutto chiaro chi abbia fatto ``Waterfall Agile'' all'inizio: dalla
risposta di Simone (``ho fatto anch'io waterfall agile per un po''') sembra
trattarsi di un errore comune, che Diego attribuisce anche a se stesso o agli
studenti in generale.}. Poi la domanda: quali sono stati i problemi più gravi?

\paragraph{I task ``tenuti in caldo''.} Simone si è accorto che le persone non
chiudevano i task al momento giusto. Se diceva che un task richiedeva, per
esempio, quattro ore, qualcuno lo finiva in un'ora e poi, invece di segnarlo
come completato, aspettava (``me lo rivedo più tardi, tanto ho ancora altre
tre ore''). Non funziona così: \textbf{quando un task è completato, è
completato}, punto. Non c'è ``me lo rivedo tra un po''' né ``devo farlo in
queste quattro ore, quindi ci metto quattro ore''. Questo falsava i risultati,
e Simone ha cambiato subito approccio\footnote{L'esempio parla di stime in
ore, in apparente contrasto con il consiglio di non usarle
(sezione~\ref{sec:z3-velocity}); è plausibile che si riferisca proprio a una
fase precedente in cui le ore erano ancora usate, e che il ``cambiamento
immediato'' sia stato anche questo.}.

\paragraph{Le riunioni che nessuno voleva fare.} Durante il secondo anno le
persone non volevano fare le riunioni settimanali. All'inizio le facevano solo
tra designer; poi le hanno estese a tutti, perché non aveva senso limitarle,
ma la risposta era ``questa riunione è inutile, non ha senso''. In quel caso
Simone non è riuscito a risolvere il problema: non sapeva come farlo. C'erano
problemi anche sul piano umano, non tra i designer ma con membri di altre
classi: problemi che, col senno di poi, si sarebbero potuti risolvere, ma
all'epoca non ci sono riusciti. La soluzione adottata è stata non fare le
riunioni con tutto il team: gli altri membri mandavano un \textbf{recap
settimanale per messaggio}. Non era un metodo del tutto fallimentare, ma si
perdono molti benefici delle riunioni tutti insieme: il \emph{team building},
il confronto, il vedersi, il parlare. È stato, dice, ``un grosso problema e
anche un grosso fallimento''.

Per il resto, niente di così grosso da valere la pena di raccontarlo:
valutazioni di task sbagliate, sprint fatte male, e un periodo di
``Waterfall Agile'' anche per lui, finché non ha capito di sbagliare.

\subsection{Che cosa rispondere a chi dice che le riunioni sono inutili}

E oggi, se qualcuno dicesse ``per me queste riunioni sono inutili'', che cosa
gli risponderebbe? \textbf{``Perché?''}. E da lì si capisce qual è il
problema: non vuole dedicarci tempo? Durano troppo? Gli costa fatica? Ha
problemi con le riunioni perché soffre di qualche forma d'ansia e non gli
piace stare con tutti i membri del team, anche fuori dalle riunioni (``tipo
ci becchiamo a giocare online a qualcosa'')? Bisogna capire il perché e
\textbf{tirarlo fuori dalle persone}, perché spesso le persone stesse non sanno
il perché delle cose che dicono e sentono.

\subsection{I saluti a Simone}

Diego ringrazia Simone e spera che gli studenti abbiano appreso il più
possibile: un'esperienza più diretta di quella di Simone e Cristian non si può
avere, perché hanno fatto letteralmente lo stesso percorso (la lezione di
Diego è ``fondamentalmente la stessa'' da un anno all'altro, con qualche
aggiunta). Ricorda che Simone ha concesso anche il foglio della
\emph{speedrun velocity}, che Diego passerà alla classe, e lo saluta per non
rubargli altro tempo.

% ---------------------------------------------------------------------
\section{Cristian: il punto di vista del game designer e project manager}
\label{sec:z3-cristian}

\subsection{Chi è il relatore}

Il secondo intervento è di Cristian, che attiva la webcam (``mi vedete super
zoomato, spero non sia una brutta vista'') e condivide le slide, mentre Diego
cerca di farle andare in modalità presentazione perché vengano meglio nella
registrazione. Cristian ringrazia Diego per la presentazione e si presenta:
\begin{itemize}
  \item è uno \textbf{studente del terzo anno} di Game Design, quindi ha fatto
        lo stesso percorso della classe;
  \item è anche \textbf{tutor} del corso di Game Design per gli studenti
        dell'accademia online;
  \item è \textbf{game designer e project manager} di due progetti:
        \begin{itemize}
          \item \emph{Gamers in Time}\footnote{Così nella trascrizione;
                il titolo non è stato verificato.}, un \textbf{gioco da
                tavolo}, vincitore del \emph{Big One} del primo anno;
          \item \emph{L'Inf}\footnote{Il titolo compare in molte forme
                (``l'Inf'', ``Linf'', ``Leanf'', ``Linfa''); qui si usa
                \emph{L'Inf}, la più frequente. Non è chiaro se la grafia
                corretta sia \emph{Linf}, \emph{L'Inf} o \emph{Linfa}.}, un
                \textbf{videogioco} che ha vinto, ``tra virgolette'', il
                \emph{Bigger One} del secondo anno e che il team ha avuto
                l'opportunità di continuare come \emph{Biggest One} del terzo
                anno, opportunità che ha accettato.
        \end{itemize}
\end{itemize}
In entrambi si è occupato di tutta la parte di design, insieme ad altri
colleghi, ed è per lo più anche \textbf{project manager, scrum master,
producer}, ``come volete definirmi''.

\subsection{Come è strutturata la sua parte}

Oltre all'esperienza personale, la lezione ha due parti principali:
\begin{enumerate}
  \item la \textbf{produzione} effettiva tramite Scrum;
  \item la \textbf{gestione del team} in un ambiente di lavoro Scrum.
\end{enumerate}
Possono sembrare la stessa cosa, ma ``il progetto tendenzialmente è una cosa
e il team è un'altra'': ci sono ovvi rapporti tra le due parti, ma per
ciascuna ci sono precauzioni specifiche.

\subsection{Le premesse}

\begin{itemize}
  \item Quello che descrive è \textbf{un} percorso, un modo di applicare
        Scrum, \textbf{non l'unico}: se non si è d'accordo, o se il team non
        ``sopporta'' questo metodo, si possono fare tutti i cambiamenti che si
        vuole (e qualche consiglio lo darà strada facendo).
  \item Non è una lezione descrittiva come quelle di Diego (che cos'è Scrum,
        le definizioni, le cerimonie, come si calcolano la velocity e il peso
        dei task), ma un \textbf{confronto} sull'esperienza dei due progetti,
        più su \emph{L'Inf} che su \emph{Gamers in Time}, perché Scrum l'ha
        applicato al videogioco e non al gioco da tavolo.
  \item \textbf{Tutto è in continua evoluzione}. Per questo vuole le domande
        della classe: il primo obiettivo è chiarire i dubbi e far uscire tutti
        dalla lezione avendo capito un po' di più la mentalità agile, ``con
        tutti voi stessi''; il secondo è ricevere \textbf{feedback} per
        migliorare non solo le lezioni, ma lo stesso modo in cui applica
        Scrum.
  \item Ha un po' di mal di gola e si fermerà spesso a bere: si scusa in
        anticipo.
\end{itemize}

% ---------------------------------------------------------------------
\section{Il caso di studio: \emph{L'Inf}}
\label{sec:z3-linf}

\subsection{Le caratteristiche del gioco}

Qualcuno in classe forse conosce Cristian o ha provato il gioco a Level Up,
l'evento in cui è stato presentato. Le caratteristiche principali:
\begin{itemize}
  \item è \textbf{single player} e offline: niente co-op, niente
        multiplayer, niente server;
  \item è un'\textbf{avventura fantasy}: un gioco di \textbf{durata fissa}, che
        tendenzialmente si gioca per una \emph{run}; chi l'ha apprezzato di
        più e vuole vedere tutti gli aspetti del design e delle meccaniche, o
        esplorarne i segreti, lo rigiocherà;
  \item il giocatore ha delle \textbf{magie colorate} con cui affronta
        \textbf{orde di nemici colorati}: spara le magie e, a seconda di
        quanti nemici ``scoppiano'' con un solo colpo, ottiene vari bonus,
        come succede in \emph{Candy Crush}.
\end{itemize}

\paragraph{Dati e passione.} Diego interrompe perché ritiene questa parte molto
importante: com'è nato l'elemento davvero innovativo di \emph{L'Inf}? Dallo
\textbf{studio dei dati}, ma anche dalla \textbf{passione del team}. I dati
danno una direzione: le ricerche di Cristian indicavano il successo delle
meccaniche \emph{match-3}\footnote{Nella trascrizione ``match free'' e
``match tree'': si tratta delle meccaniche \emph{match-3} (allineare tre o
più elementi dello stesso colore), tipiche di \emph{Candy Crush}.}; e
Cristian è anche, principalmente, un giocatore mobile. La combinazione tra
passione e ricerca ha prodotto un progetto molto valido sul piano della
fattibilità (single player, fantasy, espandibile).

Cristian aggiunge un'altra ricerca fondamentale, da cui è stato poi delineato
il \textbf{documento di visione}\footnote{Nella trascrizione ``iVision'' e
``l'iVision'': si tratta quasi certamente della \emph{vision}, il documento
di visione (sezione~\ref{sec:perche-documentare}).} del gioco: quella su
che cosa chiedesse Epic Games per finanziare i progetti con l'\textbf{Epic
MegaGrant}. Il team ha individuato tre elementi fondamentali\footnote{È la
lettura dei requisiti fatta dal team, come riferita da Cristian; non è stata
verificata sui criteri ufficiali del programma.}:
\begin{enumerate}
  \item \textbf{innovazione};
  \item \textbf{temi ambientali}: attraverso i colori e l'ambientazione, il
        gioco vuole criticare l'inquinamento globale;
  \item un \textbf{elemento narrativo}.
\end{enumerate}
Questi tre elementi, le ricerche specifiche sui vari tipi di gioco e la
passione (di Cristian in primis, come lead, ma anche degli altri membri,
incuriositi dalla proposta) hanno portato a un progetto che al momento sembra
funzionare: ``stiamo andando nella direzione giusta, almeno secondo quanto
previsto; non si può mai prevedere al 100\% ciò che accadrà''.

\paragraph{Il contenuto.} Oltre alla meccanica \emph{match-3}, il gioco
prevede:
\begin{itemize}
  \item \textbf{quattro modalità di sparo};
  \item \textbf{cinque livelli}, con esplorazione e fasi di combattimento
        contro le orde;
  \item \textbf{puzzle};
  \item \textbf{cinematic} per lo più ambientali: niente attori, doppiatori o
        mimica facciale, ma scene di presentazione dell'ambiente di gioco;
  \item una \textbf{boss fight} finale.
\end{itemize}
Il perché di questo elenco sarà chiaro tra poco.

\subsection{Un ``trauma'' di produzione: il cambio di team}

La produzione si è sviluppata in due sezioni, e \emph{L'Inf}, ``per fortuna o
purtroppo'', ha avuto un grave problema: è la risposta, anticipata, alla
domanda fatta a Simone sui problemi più gravi. Come tutte le produzioni di
videogiochi, ha subito dei ``traumi''. Il più grande: \textbf{il 100\% del
team originario è stato sostituito}. Non per scelta, ma perché i membri non si
sono iscritti al terzo anno: avevano altri obiettivi; c'è chi cerca lavoro e
chi ha aperto un'azienda. Cristian, volendo continuare il terzo anno e
soprattutto \emph{L'Inf}, un progetto in cui crede e che ha dimostrato il suo
potenziale a Level Up e con i feedback costanti dei giocatori, ha dovuto fare
dei \textbf{colloqui} per formare un nuovo team che si adattasse meglio al tipo
di sviluppo del gioco\footnote{La frase è in parte persa nella trascrizione
(``ho dovuto fare dei colloqui Grazie a tutti. Cioè, servisse una parola
brutta, ma semplicemente si adattasse di più al tipo di sviluppo''): il
senso ricostruito è che ``colloqui'' è una parola forte, ma che si trattava
di scegliere persone adatte a questa fase.}.

La composizione è cambiata così:
\begin{center}
  \begin{tabular}{@{}lp{0.55\textwidth}@{}}
    \toprule
    \textbf{Team} & \textbf{Composizione} \\
    \midrule
    originario & 2 game designer, 2 level designer, 1 technical designer \\
    attuale    & 3 game e level designer, 2 technical designer \\
    \bottomrule
  \end{tabular}
\end{center}
Il motivo: la preproduzione, cioè la parte più importante di game design, era
finita, e servivano \textbf{più tecnici}. Il gioco si basa molto sulle
meccaniche: un sistema di combattimento innovativo che richiede di scriptare
le meccaniche, diversi tipi di sparo, puzzle, una boss fight. È molto più
basato sulle meccaniche che su narrativa o ambientazione. Allo stesso tempo
servivano level designer, perché il gioco doveva prendere forma e i livelli
andavano costruiti.

\subsection{Perché raccontare tutto questo}

Non per vendere il gioco (``anche se magari è un effetto collaterale ben
voluto''), ma per mostrare che \textbf{ogni progetto, a seconda di come è
composto} sul piano del design, della produzione e soprattutto del team,
\textbf{ha un modo diverso di applicare Scrum}. Per questo Cristian è sceso nel
dettaglio delle caratteristiche: ognuna comporta ambiti di produzione diversi.
\begin{itemize}
  \item I livelli richiedono level designer.
  \item Il \emph{match-3} e i diversi tipi di combattimento richiedono
        tecnici.
  \item Il single player è più un vantaggio che un limite: un gioco
        multiplayer non è il doppio, ma ``tipo dieci volte più difficile'' da
        gestire a livello produttivo.
  \item Il genere \emph{fantasy adventure} deriva soprattutto dalle ricerche.
\end{itemize}
Quindi \textbf{ogni progetto ha il suo framework di Scrum}. E il fatto di
aver cambiato team a metà produzione, finita la \emph{vertical slice}, con
persone completamente nuove (ex compagni di studio, che però non avevano mai
visto come funzionava il gioco), ha avuto diverse conseguenze, che Cristian
rimanda alle domande di fine lezione.

% ---------------------------------------------------------------------
\section{Il loop di Scrum}
\label{sec:z3-loop}

L'applicazione di Scrum è un \textbf{loop} molto semplice, simile a quello
mostrato da Simone:
\begin{enumerate}
  \item \textbf{creazione del backlog};
  \item \textbf{selezione dei task} da svolgere, a ogni sprint;
  \item \textbf{implementazione dei task} durante la sprint, con controlli e
        \emph{scrum} (con la frequenza che si decide);
  \item \textbf{calcolo della velocity}, per vedere com'è andata e come
        migliorare la produzione nelle sprint successive;
  \item \textbf{interpretazione} della velocity per fare il \emph{tuning}
        della produzione; e si ricomincia.
\end{enumerate}
Cristian le affronta una per una, con esempi della sua esperienza e consigli.
È contento che la classe stia cominciando proprio adesso, perché questa parte
potrebbe tornare utile già dalla settimana successiva, quando gli studenti
applicheranno il \emph{framework}.

% ---------------------------------------------------------------------
\section{La stesura del backlog}
\label{sec:z3-backlog}

\subsection{Un backlog per milestone}

Anche Cristian, all'inizio, cercava di prevedere il più possibile tutto il
gioco, dal prototipo fino al rilascio. Oggi invece lo struttura
\textbf{per milestone}: il backlog dell'\emph{alpha}, poi quello della
\emph{beta}, poi quello della \emph{gold}. Questo non impedisce di cominciare
a prevedere che cosa ci sarà in beta e in gold, ma dà un'idea molto chiara di
dove si sta andando, soprattutto per capire che cosa tagliare o aggiungere
quando ci si accorge che qualcosa serve prima o dopo, e adattarsi di
conseguenza: ``la cosa più importante del mindset agile è adattarsi ed essere
flessibili''.

\subsection{Da dove vengono i task}

Il gioco deve pur nascere da qualche parte: i task vengono dai
\textbf{documenti chiave} della produzione, già visti nel corso
(capitolo~\ref{cap:produzione-documentazione}):
\begin{itemize}
  \item il \textbf{documento di visione};
  \item i \textbf{GDD} delle varie feature, che dettagliano la visione;
  \item il \textbf{versioning} delle feature, a cui si dovrebbe pensare già
        mentre si scrive il GDD;
  \item gli \textbf{LDD} (\emph{level design document}) e tutti gli altri.
\end{itemize}
L'ordine dipende dal gioco: in un gioco fortemente basato sulle meccaniche è
probabile che gli LDD arrivino dopo lo sviluppo del \emph{core}; in un gioco
basato sull'esplorazione non succede proprio l'inverso, perché le meccaniche
servono comunque, ma le due cose avvengono quasi in contemporanea.

Una volta che si ha la documentazione con tutte le descrizioni e le
caratteristiche delle feature, \textbf{dalla versione 0 alla versione X}, si
ha già una serie di implementazioni da realizzare. Un esempio:
\begin{itemize}
  \item \textbf{versione 1}: il giocatore deve camminare; si usa il
        manichino base di Unreal;
  \item \textbf{versione 2}: si implementa la \emph{skeletal mesh} fatta dal
        3D artist e si adatta la camera al movimento, con gli effetti di
        vibrazione a seconda che il personaggio corra, cammini o sia
        abbassato;
  \item \textbf{versioni successive}: quando il personaggio si muove su
        terreni diversi solleva tipi diversi di polvere, di più su un terreno
        molto polveroso. Non è fondamentale e potrebbe essere una delle
        ultime versioni, verso la fine.
\end{itemize}
Il consiglio è che, quando si fa il \emph{versioning}, si inseriscano
\textbf{tutti gli ambiti necessari} per quella feature: programmazione, arte
(se serve un asset), animazione, UI, VFX, suoni. Se tutti i \emph{field}
compaiono già nella documentazione, la creazione del backlog sarà più chiara
e più spedita.

\subsection{Come organizzare i task}

I giochi hanno moltissimi task. Come organizzarli perché siano leggibili non
solo da chi li ha scritti (scrum master o project manager), ma anche da chi
deve navigarli per capirci qualcosa? Cristian usa tre tipi principali di
organizzazione.

\paragraph{1. Per feature (o epic).} Si suddividono i task per feature,
collegandosi alle \emph{epic} di Simone. Ad esempio, l'epic è il movimento
del personaggio e i task sono camminata, corsa, salto, volo, ``dipende dal
gioco''.

\paragraph{Tutto in un solo posto.} La suddivisione non va fatta per forza in
documenti diversi, anzi Cristian lo sconsiglia: se tutto è nello stesso
strumento (Google Sheets, Monday, ClickUp, Trello) si usano i
\textbf{filtri} e si trova tutto lì, senza navigare tra mille siti o mille
documenti.

\paragraph{Per ambito, con un color code.} Quando si suddivide per ambito
(programmazione, arte, UI, eccetera), conviene usare un \textbf{codice
colore}. Cristian lo usa ``per qualsiasi cosa'': uno per le milestone, uno
per gli ambiti di produzione. Diego gli chiede di mostrare il suo:
\begin{center}
  \begin{tabular}{@{}lll@{}}
    \toprule
    \textbf{Ambito o milestone} & \textbf{Colore} & \textbf{Perché} \\
    \midrule
    Suoni        & viola        & --- \\
    VFX          & rosso        & sta imparando a usare Niagara \\
    Level design & verde        & ``erano finiti i colori'' \\
    Game design  & arancione    & è il suo colore preferito \\
    \midrule
    Alpha        & azzurro      & --- \\
    Beta         & viola        & --- \\
    Gold         & giallo scuro & il colore richiama il nome \\
    \bottomrule
  \end{tabular}
\end{center}
Come si vede, è un misto di funzionalità e preferenze: la cosa più importante
è che \textbf{tutto il team sia d'accordo}.

\paragraph{2. Per milestone, con le sotto-milestone.} Un secondo metodo è
suddividere il backlog per milestone. Ma se il backlog dell'alpha è separato
da quelli di beta e gold, come si fa a dividere in alpha, beta e gold un
backlog in cui tutto è alpha? Cristian suddivide le milestone in
\textbf{sotto-milestone}: l'alpha dell'alpha, la beta dell'alpha, la gold
dell'alpha. Serve ad avere idee più chiare, a tenere sempre in mente lo
\emph{scope} senza perderlo e soprattutto a capire meglio \textbf{se il gioco
è fattibile}. Una suddivisione più fine permette anche di calcolare con più
precisione (una precisione che ``lascia il tempo che trova'') quante sprint
servono per ogni sotto-milestone. La dissonanza tra ``dividete tutto in alpha,
beta e gold'' e ``dentro l'alpha trovate alpha, beta e gold'' è quindi voluta:
si parla delle sotto-versioni.

\paragraph{3. Per priorità.} Chi non vuole usare le sotto-milestone può usare
il sistema di \textbf{priorità} di cui parlava Simone: una scala numerica, da
1 a 5 o da 1 a 10, a seconda di come si riesce meglio a distinguere le
priorità. Serve soprattutto a capire quali task selezionare
(sezione~\ref{sec:z3-selezione}). In realtà è più o meno la stessa cosa con un
approccio diverso: indicare alpha, beta e gold dell'alpha crea già di per sé
delle priorità. ``Dipende tutto da voi.''

\subsection{Come scrivere i task e dove}

Il consiglio è scriverli \textbf{dal punto di vista dell'utente}. Ma, come
diceva Simone, non sempre ``il gioco vale la candela''. Lo è soprattutto in
un caso: se il programmatore legge un task dell'artista 3D scritto in modo
troppo tecnico e non lo capisce, si crea dissonanza tra i reparti e possono
nascere attriti. Quindi è sempre consigliabile, ma non obbligatorio: spesso
basta una chiamata (``prima semplicemente ti giravi verso il tuo compagno e
gli chiedevi che cosa significa''). ``Però è giusto che vi dica le cose come
dovrebbero essere fatte.''

Dove scriverli? In qualsiasi strumento: Trello, Monday, ClickUp, gratuiti o a
pagamento, oppure un foglio Google costruito da sé. Dall'esperienza di
Cristian emerge che \textbf{Google Sheets sta antipatico alle persone}: già il
nome ``foglio di calcolo'' spaventa. Strumenti come Monday o ClickUp hanno più
colori e soprattutto \textbf{mandano le notifiche}, quindi sono più accessibili
a tutti. ``Decidete voi.''

\subsection{Non prevedere tutto: la flessibilità}

È fondamentale capire che \textbf{prevedere tutti i task passo per passo}, in
un ordine preciso (``dopo la camminata faremo il salto, dopo il salto il
sistema di sparo, e così via all'infinito''), va contro i principi di Agile,
in particolare contro il \textbf{quarto valore}: rispondere ai cambiamenti più
che seguire un piano.

Da qui il consiglio centrale: la \textbf{flessibilità} deve essere il cuore
della propria filosofia di lavoro e del proprio modo di applicare Scrum.
Bisogna essere sempre pronti a \textbf{tagliare} nel modo più agile possibile:
se si è capita davvero la mentalità agile, si riesce a tagliare una feature
\textbf{senza bloccare la produzione delle altre}.

L'esempio più eclatante viene dall'inizio di \emph{L'Inf}: Cristian faceva il
backlog di tutto il gioco, arrivando a circa \textbf{400 task} di ogni ambito,
nel dettaglio: per la camminata arrivava a indicare come si muoveva la camera
quando il personaggio girava a destra, e come cambiava quando saltava o cadeva
da un dirupo. ``Vi viene subito da pensare: ma come fai a prevedere tutto?''
Infatti la soluzione è \textbf{non prevedere tutto}: capire a grandi linee ciò
che si vuole, qual è l'obiettivo della feature, senza scendere troppo nel
dettaglio, perché il dettaglio crea di per sé un \textbf{lavoro Waterfall}, che
è ciò che bisogna assolutamente evitare. E scrivere un backlog così dettagliato
richiede tantissimo tempo, che è limitato: la produzione ne risente in tutti i
modi possibili.

\paragraph{Bonus: gli elementi innovativi.} Se il gioco (il proprio
\emph{Bigger One}) ha elementi innovativi, la flessibilità va portata
\textbf{ai massimi livelli}. \emph{L'Inf} era partito con una meccanica di
\textbf{assorbimento e sparo}: si assorbe un nemico colorato, ad esempio
giallo, e lo si spara contro altri nemici gialli, prendendo ispirazione da
altri giochi. Nel loro caso, però, creava più problemi che altro: per due o
tre ore di gioco, ``assorbire, sparare, assorbire, sparare'' risultava
\textbf{troppo ripetitivo}, e hanno deciso di tagliare la meccanica. Il
problema è che l'hanno tagliata dopo che Cristian aveva già scritto
\textbf{600 task nel backlog e 12 GDD}: tempo perso. Per questo la
flessibilità, negli elementi innovativi ma in generale, ``è la prima cosa''.

\subsection{Quanto dettaglio?}

Quando un backlog è abbastanza dettagliato e quando lo è troppo poco?
``La risposta purtroppo non c'è.'' Viene da due fattori:
\begin{enumerate}
  \item l'\textbf{esperienza}: man mano che si scrivono GDD e si va avanti
        nello sviluppo, ci si ``tuna'';
  \item il \textbf{trial and error}: si sbaglia e si torna indietro, il che
        riporta, di nuovo, all'esperienza.
\end{enumerate}

% ---------------------------------------------------------------------
\section{La selezione dei task}
\label{sec:z3-selezione}

Con un bellissimo backlog ordinato per colori, ambiti e priorità, quali task
si selezionano? ``Da qualche parte bisogna pur cominciare.'' Le domande
principali sono due: \textbf{quali} selezionare e \textbf{quanti}.

\subsection{Quali task}

\paragraph{La priorità.} Il primo criterio è la priorità: i task da fare
prima sono le scelte più urgenti, e quindi si scelgono quelli\footnote{L'inizio
di questo passaggio manca nella trascrizione (``quanti selezionare. cui
scegliere prima, sarà la scelta più urgente da fare''): il primo criterio è
stato ricostruito dal contesto, anche perché poco dopo Cristian dice che,
risolto il problema delle competenze, ``si ritorna a selezionarli per
priorità''.}.

\paragraph{Le capacità attuali degli sviluppatori.} Il secondo criterio sono
le \textbf{abilità dei developer}. Cristian precisa che non intende in alcun
modo giudicarne la professionalità: intende letteralmente \textbf{che cosa sa
fare lo sviluppatore in quel momento}. Se un programmatore non ha mai usato
Unreal, non è il momento di dargli task di implementazione in Unreal: bisogna
prevedere anche il \textbf{periodo di apprendimento}, eventualmente come task
vero e proprio (``hai questa sprint per imparare a implementare l'IA'').

È successo due mesi prima con uno dei suoi technical designer: gli ha dato
\textbf{due settimane} dicendogli ``in queste due settimane non ti preoccupare
di niente, il resto lo produciamo noi''. Doveva essere d'accordo, e lo era,
anzi era contento di approfondire l'argomento: studiare come funzionano le IA
(anche grazie a vari link passati da Diego) e soprattutto come implementarle
nel motore. Per due settimane ha dovuto ``solamente'' pensare
all'apprendimento. Il consiglio: \textbf{prevedere questi ``tempi morti''},
perché ne ha bisogno il team, e ne ha bisogno chi pianifica per non falsare la
previsione dello \emph{scope}. Se in questo momento (``sottolineo, in questo
momento'') un programmatore non sa fare una cosa, la si rimanda e si mette
altro; prima o poi imparerà, e a quel punto si torna a selezionare per
priorità.

\subsection{Quanti task}

Il primo criterio sulla quantità è la \textbf{velocity}, ma Cristian l'ha
messo nella slide solo per completezza del loop: se non si ha ancora una
velocity, non ha senso. Per ora la si trascura e si fa finta che sia la
\textbf{prima sprint}, la prima applicazione di Scrum.

Nelle prime sprint la quantità si basa sulle \textbf{capacità} e sulla
\textbf{disponibilità} degli sviluppatori:
\begin{itemize}
  \item \textbf{disponibilità}: per quanta parte della sprint il technical
        designer (per riprendere l'esempio) è disponibile? Non in termini di
        ore, perché Cristian concorda con Simone che le ore sono sempre molto
        relative;
  \item \textbf{performance}: gli studenti si conoscono ormai da quasi un
        anno e mezzo. Se uno sviluppatore non ha mai fatto un livello in
        Unreal (per motivi suoi, magari problemi al PC) e dice ``voglio fare
        un livello e ci metto due settimane'', sta a chi pianifica cominciare
        a capire \textbf{chi può fare che cosa}. Tutti possono fare tutto, se
        si prevede il tempo di apprendimento; ma se Cristian arrivasse a dire
        ``in due settimane vi faccio l'IA del vostro gioco implementandola in
        Unreal'', ``vi sto prendendo in giro''.
\end{itemize}
Queste sono le analisi preliminari quando non si ha una velocity. Se invece
non si conosce affatto il team, l'unico punto di partenza è la
\textbf{disponibilità} degli sviluppatori, da cui si stima quanto si riesce a
fare.

Tanto, \textbf{nelle prime sprint si sbaglieranno completamente tutte le
previsioni}, ``ed è giusto così, e ben venga'': se non si sbaglia,
probabilmente si sta applicando Waterfall, o si sta applicando male Scrum.
All'inizio si tratta solo di buttarsi: ``questa settimana abbiamo previsto 15
task; ne abbiamo fatti 40. Bene, la prossima ne proviamo 40''; poi magari se ne
fanno 25, e così via, \textbf{tarando} man mano la quantità.

\subsection{Definire bene il goal}

Per ogni sprint, e anche per ogni singolo task, bisogna \textbf{definire bene
l'obiettivo}. Serve a mantenere concentrati gli sviluppatori: se so che entro
questa settimana devo finire il block out del livello, so dove sto andando,
lavoro per quel risultato e soprattutto \textbf{non perdo tempo}, come succede
spessissimo. Se durante il block out qualcuno si mette a inserire effetti
visivi ambientali o il \emph{post-process}, viene il dubbio che non stia
lavorando per l'obiettivo: è un \textbf{campanello d'allarme}, perché potrebbe
far mancare l'obiettivo per cui la sprint era stata pensata.

\subsection{Bonus: progettare il ritmo dello sviluppo}

Un consiglio che Cristian ha sentito raramente, ma che ha deciso di
sperimentare: poiché gli studenti sono designer, e tendenzialmente anche lo
scrum master del loro progetto lo sarà, possono \textbf{progettare il
\emph{pace} dello sviluppo}.

Di solito si pensa (e potrebbe anche essere giusto) che l'obiettivo di Scrum
sia trovare una \emph{velocity target} e mantenerla costante per tutta la
produzione. Può essere una soluzione; ma con entrambi i suoi team Cristian ha
notato che, per come sono fatti gli esseri umani, ci sono \textbf{picchi di
attenzione e di produzione}. Con l'esperienza si impara a percepirli in
anticipo e a progettare di conseguenza; oppure sono gli sviluppatori stessi a
dirlo: se la settimana prima hanno implementato 12 nemici, 15 modalità di sparo
e 13 cutscene, quella dopo diranno ``andiamoci piano''. Di solito non succede,
perché le sprint sono programmate apposta per adattarsi alle capacità degli
sviluppatori, ma è un metodo possibile.

Che cosa intende, in concreto? Se per un periodo si sono implementate molte
feature o sistemi che però, visivamente, \textbf{non hanno portato un aumento
di qualità percepibile nel gioco}, gli sviluppatori non si sentono
riconosciuti e ripagati dello sforzo. Allora, la settimana dopo, ci si
concentra su qualcosa che dà \textbf{soddisfazione}. Questo lo permettono le
sprint, lo permette Scrum e non Waterfall: è una differenza fondamentale.

\paragraph{L'esempio di \emph{L'Inf}.} Per più di otto mesi il gioco è rimasto
con \textbf{quattro tipi di nemici}. Nonostante l'impegno (quattro nuove
modalità di sparo, la boss fight, tantissima programmazione ``di
background'', cioè non visibile al giocatore ma necessaria al funzionamento
del gioco), il team cominciava a percepire che \textbf{non c'era corrispondenza
tra impegno e qualità percepita}. Cristian ha allora deciso di limitare per un
periodo la produzione degli elementi di background (programmazione come l'IA,
tutto ciò che faceva proseguire il gioco ``sotto''), di inserire altri tipi di
nemici, che tutti non vedevano l'ora di avere (``sapevo che era uno di quei
jolly da sfruttare per far risalire il morale del team''), e di dedicare tempo
a effetti visivi e sonori, a migliorare tutto ciò che era estetico. Diego
suggerisce il termine: la \textbf{juiciness}. Ciò che vedeva il giocatore lo
vedevano anche gli sviluppatori, che si sentivano più appagati e sentivano i
propri sforzi ricompensati. ``Adesso loro non lo sanno, ma li ributterò sotto
un treno di lavoro'' (``in realtà lo sanno, però non me lo dicono'').

Per Cristian è un punto fondamentale, che gli è costato ``mesi di
sofferenze'' di esperienza per capirlo. Diego invita la classe a fare domande
su un \emph{insight} ottenuto per tentativi ed errori.

% ---------------------------------------------------------------------
\section{L'implementazione: controlli, blocker e feedback}
\label{sec:z3-implementazione}

\subsection{Il check dei progressi}

Selezionati i task, viene la parte ``più tranquilla'' dal punto di vista di
Scrum (i problemi di implementazione, ovviamente, sono altri): parte la
settimana di lavoro, ognuno ha i suoi task, i \emph{task owner} hanno scelto
che cosa fare, e questo prosegue per tutta la durata della sprint.

Durante l'implementazione si fa un \textbf{check dei progressi}. Non serve a
controllare che cosa stanno facendo le persone: di base bisogna
\textbf{fidarsi} che, se il programmatore dice ``farò questo'', lo farà entro i
limiti della sprint. Cristian parla qui di \emph{daily scrum} perché è ciò che
si dovrebbe fare in generale, anche se il suo team non li fa ogni giorno
(sezione~\ref{sec:z3-team}). Oltre alle solite domande (che cosa hai fatto
ieri, che cosa farai oggi), \textbf{la cosa più importante è trovare i
blocker}.

\subsection{Agire subito sui blocker}

Se il programmatore dice ``non ho potuto implementare il personaggio perché mi
manca l'asset, la \emph{skeletal mesh} del protagonista'', bisogna
\textbf{agire subito}, in due modi:
\begin{enumerate}
  \item \textbf{dare altri task} al programmatore, perché altrimenti resta lì
        ad aspettare, e questo non lo si vuole assolutamente;
  \item \textbf{capire che cosa è successo} alla skeletal mesh, parlando con
        la 3D artist. Se è solo questione di tempo, non è davvero un problema:
        è normalissimo aver assegnato l'implementazione prima che l'asset
        fosse pronto (``è normale sbagliare, l'importante è rimediare''). Se
        invece anche lì c'è stato un blocco, bisogna \textbf{risalire a
        cascata} fino a trovarlo.
\end{enumerate}
Magari si scopre che il problema sono proprio i game designer: quando hanno
fatto la \emph{mood board} non sono stati abbastanza chiari, e la 3D artist ha
dovuto spendere una settimana a cercare reference per conto suo, mentre il
designer pensava che quelle fornite bastassero e non si è posto il problema
della fase di ricerca delle reference da parte di chi fa la grafica o il 3D
(Cristian parla al femminile perché la 3D artist del suo team è una donna, ma
il discorso è generale). Anche questo, come l'apprendimento, è uno di quei
task che non bisogna trascurare.

\subsection{Analizzare i blocker prima di risolverli}

Il consiglio è \textbf{analizzare i blocker}, non risolverli appena si
presentano, perché si rischia di trascurare il problema alla radice. Cristian
richiama una delle prime lezioni del primo anno, quella su \textbf{bisogni e
problemi}: bisogna capire se il blocker è un problema o un bisogno, e agire di
conseguenza. Il guaio nasce quando lo si trascura o, peggio, quando lo si
risolve in fretta senza averlo considerato: ``poi vi si ritorcerà contro''.

\paragraph{L'esempio dei muri invisibili.} \emph{L'Inf} è un gioco di
esplorazione ambientale, per lo più all'aperto: su cinque livelli, circa
l'80\% del gioco è all'esterno. Per dare \emph{meaning} al livello finale,
quello della boss fight e della rivelazione narrativa su dove siano finiti
tutti gli esseri umani del pianeta, il team ha deciso di sottolinearne la
diversità rendendolo \textbf{al chiuso}.

La camera è in terza persona e abbastanza lontana dal giocatore, non classica.
Dare le collisioni alla camera sembrava fastidioso, perché a ogni rotazione
``sbatte'' contro i muri. La prima idea: \textbf{rendere invisibili i muri}
quando la camera ci passa dietro. Non si potevano far sparire tutti insieme
quando il giocatore era solo dietro il muro destro, quindi ogni muro aveva il
suo \emph{trigger}, che stabiliva quante e quali pareti dovessero sparire. Il
livello dura mezz'ora, con una buona decina di minuti di esplorazione al
chiuso, completamente circondata da pareti: \textbf{un mese di lavoro} solo
per assegnare i trigger, cambiare i materiali e gestire il tutto. ``Era l'idea
geniale che ci era venuta in quel momento.''

Quando hanno consegnato la beta e l'hanno fatta provare ai professori e ad altri
giocatori, tutti hanno detto più o meno la stessa cosa: i muri che spariscono
\textbf{non sono eleganti e rompono l'immersione}. Hanno tagliato tutto e
rimesso le collisioni alla camera. Se avessero analizzato bene il problema
prima, avrebbero capito che bastava \textbf{gestire meglio gli spazi}: non fare
curve strette ma allargarle, così che anche ruotando la camera il giocatore
tendenzialmente non sbatta contro il muro. Era ``assolutamente una questione di
riallargare più o meno tutti gli spazi''; ed è ciò che hanno fatto alla fine,
insieme alle collisioni. ``Ad averlo fatto prima avremmo risparmiato due
settimane di tempo''\footnote{Poco prima Cristian parla di ``un mese buono'' di
lavoro sui trigger (correggendosi da ``due mesi''); le ``due settimane'' sono
forse il risparmio netto stimato, tenendo conto del lavoro di allargamento
degli spazi che è stato comunque necessario. La trascrizione non permette di
chiarirlo.}. Ecco perché è importante analizzare i blocker.

\subsection{Bonus: i feedback costanti}

Cristian nota che, soprattutto ultimamente, in una fase di
\emph{lockdown}\footnote{Il riferimento al lockdown colloca la lezione nel
periodo della pandemia di Covid-19, coerentemente con l'accenno di Simone al
Covid e alle sedi chiuse (``quando si poteva andare'').} in cui tutti sono
costantemente in contatto su Discord, ``anche più di quanto lo saremmo nella
vita vera'', c'è un \textbf{bisogno costante di dare feedback} sul lavoro.

Lui è pro feedback in qualsiasi momento, anzi ``ossessionato'': all'inizio ha
chiesto alla classe di criticarlo e di dargli consigli sulle slide, perché è
l'unico modo per migliorare. Il problema è che \textbf{i feedback costanti
producono un calo di produttività enorme}. L'esempio è personale: Cristian è
anche il VFX artist del team, e sugli elementi estetici c'è sempre modo di
migliorare, quindi riceve di continuo feedback (``io lo farei più grande, rosso,
verde, giallo, a forma di triangolo, a forma di quadrato''). Sono feedback
correttissimi: chi li dà, soprattutto in produzione, non lo fa perché pensa
che tu sia scarso, ma per il bene del gioco, che vive e vuole ben prodotto. I
migliori mandano anche video di riferimento (``ho visto questo VFX in questo
gioco, sapresti rifarlo?''), e questo ben venga. Il problema è che siano
\textbf{costanti}.

Il consiglio: \textbf{limitare la richiesta e l'ascolto dei feedback a fasi
specifiche} (alla fine della sprint, oppure in un controllo intermedio). Dipende
anche da che cosa si sta facendo: con i VFX molte cose dette oggi sono
sistemate il giorno dopo; ma se chi costruisce un livello, in block out o in
\emph{environment}, riceve mille feedback per ogni metro quadro (``un po' più a
destra, un po' più a sinistra''), il lavoro si ferma lì. Sia perché chi lo fa
si \textbf{demoralizza}, sentendosi costantemente criticato, sia per una
questione di fiducia, tornando al discorso di Simone: ``fidati di me come io mi
fido di te. Fammi completare il task; poi, quando ti chiedo il feedback, o più
avanti, è ben accetto, anzi ben venga''.

Cristian si scusa per l'apparente presunzione di definire ``fondamentali'' i
propri consigli: lo fa perché, se li avesse avuti prima, si sarebbe risparmiato
tantissimi problemi. Diego scherza che andranno nelle slide ``con la mia faccia
super inquietante''; Cristian dice che può passare anche le sue slide.

% ---------------------------------------------------------------------
\section{Il calcolo della velocity}
\label{sec:z3-calcolo}

\subsection{Confrontare il fatto con lo stimato}

Sulla parte pratica Simone è stato più esaustivo, quindi Cristian si concentra
sull'interpretazione\footnote{Un breve tratto della spiegazione iniziale è
perso nella trascrizione, che riprende a metà frase (``nel comparare quello
che è stato fatto da quello che è stato stimato'').}. La velocity consiste nel
\textbf{confrontare ciò che è stato fatto con ciò che era stato stimato}, in
numero di task o in pesi, a seconda di come la si gestisce. Da lì si capisce
che cosa è successo e che cosa si può prevedere:
\begin{itemize}
  \item \textbf{fatto $>$ stimato}: si può prevedere di più; nella sprint
        successiva si aumentano i pesi o il numero di task;
  \item \textbf{fatto $<$ stimato}: o si è esagerato con la stima (succede
        sempre, così come succede sempre il caso opposto), oppure ci sono stati
        problemi gravi o \emph{blocker} nella produzione, ed è fondamentale
        capire quale delle due cose sia successa;
  \item \textbf{fatto $=$ stimato}: succede quasi mai, ma è l'obiettivo. Le
        possibilità sono due: si è trovato il proprio \textbf{bilanciamento}, e
        si può continuare così; oppure \textbf{non si sta applicando bene
        Scrum}, ad esempio perché si calcola la velocity su ciò che si è fatto
        invece che sulla stima.
\end{itemize}

\subsection{Più di un numero}

L'obiettivo non è arrivare a una velocity pari a 1, intesa come rapporto tra
fatto e stimato (100 punti stimati, 100 fatti). La velocity è \textbf{molto più
di un numero}: dà almeno tre informazioni.

\paragraph{1. La bontà degli strumenti di stima.} Planning poker, Fibonacci e
tutti gli strumenti di stima sono utilissimi, ma \textbf{vanno testati}: la
velocity mostra se ci sono metodi sbagliati di stima dei pesi o del numero dei
task.

\paragraph{2. La raggiungibilità dello scope.} Se per tutta la produzione la
velocity resta, ad esempio, a 0,8, sempre al di sotto di quanto dovrebbe essere
(in teoria non dovrebbe succedere, perché ci si dovrebbe tarare), e si sono
previsti dieci livelli, si capisce già da subito che \textbf{non ci si
arriverà mai}, a meno di aggiungere uno o due anni di sviluppo.

\paragraph{L'esempio degli otto livelli.} Finita la \emph{vertical slice}, il
team ha ricontrollato tutta la preproduzione e impostato il gioco come un gioco
completo, non come la creazione di una parte: c'è tutta la storia, ci sono
tutti i livelli, tutto ciò che fa il giocatore dall'inizio alla fine, chi
incontra e quando; tutto scritto, ma tarato e lasciato \textbf{flessibile},
``altrimenti ricadiamo nella trappola del Waterfall''.

A settembre, all'inizio dell'anno (``per farvi capire quanto il processo è
rapido''), avevano previsto \textbf{otto livelli}. Sembrava una stima
credibile: il vecchio team, con due level designer che volevano fare
esattamente quel lavoro, produceva un livello in \textbf{due o tre settimane}, a
una qualità abbastanza elevata: magari mancavano l'illuminazione, il
\emph{foliage} e alcuni dettagli, ma il livello era \textbf{giocabile}, che è
la parte fondamentale. Cristian, unico rimasto del vecchio team, si è basato su
quella stima: ``da qualche parte dovevo cominciare''.

Ma i level designer del nuovo team \textbf{non sono level designer puri}: sono
soprattutto game designer, interessati al game design, che hanno deciso di
sperimentare con il level design per necessità (il game design serve sempre, ma
se manca un level designer, visto che il percorso di studi insegna sia level
sia game design, ci si adatta). C'è chi si è adattato a fare il level designer,
e chi ha deciso proprio quell'anno di specializzarsi nel level design: ``perché
avrei dovuto dirgli no, tu non sei un level designer? Anzi, ben venga''.
L'unica cosa è stata prevedere che avesse più bisogno di apprendimento: magari
sbaglierà di più, dovrà informarsi su come si fanno i livelli, avrà bisogno di
più feedback. L'importante è capirlo all'inizio.

Guardando l'andamento del primo livello, Cristian vedeva che \textbf{in un mese
avevano a malapena raggiunto il block out}. Ha convocato una \textbf{riunione
di emergenza}, lunghissima (non una revisione della sprint da venticinque
minuti come quelle di Simone, ma un altro tipo di riunione). Il suo obiettivo
non era tanto convincerli che lo \emph{scope} era troppo alto (ci è riuscito),
quanto \textbf{assicurarsi che non la prendessero come una sconfitta
personale}. È stato difficile ma interessantissimo: non se la sono presa né con
se stessi, per non riuscire a performare, né con lui, che stava di fatto
dicendo ``non ce la farete''. Anzi, nonostante il taglio, sono stati
\textbf{loro stessi a proporre la soluzione}: ``ok, abbiamo capito; facciamo
cinque livelli, e se vediamo che arriviamo (speriamo a fine marzo) con tutti i
livelli, poi possiamo decidere di aumentarli o di cambiare parti del gioco''.
Cristian l'aveva già in mente, ma il fatto che l'abbiano proposta loro l'ha reso
``veramente un successo''.

Tutto questo grazie a un numero. Se Cristian fosse arrivato tra un mese e
mezzo, dopo che avevano già fatto quattro livelli ``con l'acqua alla gola'',
dicendo ``ne dobbiamo fare altri quattro entro marzo'', e poi a fine marzo
``basta, non ce la facciamo'', sarebbe stato molto più grave\footnote{Il
passaggio è confuso nella trascrizione (``mi ha portato ad affrontare queste
situazioni che affrontate a fine marzo'', ``non potne quattro''): il senso
ricostruito è che la velocity ha permesso di anticipare a inizio produzione una
decisione che, presa a fine marzo, sarebbe stata molto più dolorosa.}.
``Ringraziamo in questo caso la velocity.''

\paragraph{3. Lo stato del team.} La velocity permette di capire lo
\textbf{status del team e le sue performance}, perché si basa su una stima fatta
a partire dalla settimana precedente. Si ricollega al ``progettare il \emph{pace}
dello sviluppo'' (sezione~\ref{sec:z3-selezione}): ci si può basare proprio sulla
velocity per capire qual è il ritmo attuale e prevedere i prossimi. ``Tutto
torna.''

\subsection{Bonus: mai tornare indietro}

Un bonus che ``più che un bonus è un'altra lezione'': \textbf{prevedete sempre
task che portano a un incremento del prodotto}. La produzione può rallentare,
perché possono esserci difficoltà, ma \textbf{mai, mai, mai andare indietro}.

È un'altra lezione insegnata da \emph{L'Inf}, legata alla stesura del backlog.
Quando Cristian aveva scritto tutto il backlog (``400 task: come vedete ogni
tanto dico 400, 600, 300, giusto per dire che erano veramente tanti'') e si
accorgeva di aver sbagliato, invece di adattarsi e appoggiarsi alla produzione
in corso \textbf{fermava la produzione} (``ditemi pure che ero folle: lo ero'')
e riscriveva da capo tutti i task, arrivando magari a mille o riducendoli.

Lo stesso vale per i \textbf{documenti}. ``La documentazione è bellissima'':
Cristian era un grande fan della documentazione (Diego può confermarlo), ma
lavorando con Scrum, per ragioni di produttività, ha dovuto \textbf{scendere a
compromessi}: darsi un limite su quanto un documento è bello e ricco di
dettagli. Il motivo è sempre l'esempio della prima meccanica, l'assorbimento e
sparo: tutti i GDD del gioco si basavano su di essa, con tutte le interazioni,
anche secondarie (che cosa si poteva assorbire? Nemici, oggetti, chiavi, cose
per recuperare vita, bonus, malus, qualsiasi cosa). Tagliata la feature, tutti i
GDD sono stati ``letteralmente presi e buttati'', e Cristian ha deciso di
rifarli tutti da capo.

Questo approccio potrebbe funzionare in Waterfall (``finiti i documenti,
iniziamo la produzione''), ma l'obiettivo era \textbf{consegnare un videogioco,
non fare documenti}; e soprattutto il gioco era innovativo. È successo altre due
volte: in tutto Cristian ha riscritto \textbf{backlog e GDD da capo più
volte}\footnote{Nella trascrizione: ``è successo altre due volte quindi per ben
quattro volte io ho riscritto tutto il backlog e tutti i GDD''. Il conteggio non
torna del tutto (una volta più altre due fa tre); forse Cristian conta anche la
prima stesura.}, riportando ogni volta la produzione all'inizio (``ok, possiamo
iniziare il prototipo; no, è sbagliato, ricominciamo''). Lo racconta per far
capire quanto ci ha messo a entrare nel mindset agile, e per tranquillizzare la
classe: non indovinarci subito è normale, ``succederà il 90\% delle volte''.

Quando poi hanno fatto il prototipo, si sono accorti che \textbf{tutta la
documentazione era sbagliata}: con un gioco innovativo non si poteva sapere che
cosa avrebbe funzionato senza prima vederlo, perché non c'erano reference:
nessuno aveva fatto un gioco d'avventura \emph{match-3}. ``Il \emph{match-3}
esiste, i giochi d'avventura esistono, mischiarli insieme è una follia'', e lo
dice ``per criticarmi, non per vantarmi, perché non lo farò più''.

% ---------------------------------------------------------------------
\section{Il tuning della produzione}
\label{sec:z3-tuning}

\subsection{Che cosa e quando}

Il \emph{tuning} della produzione si basa sull'analisi fatta: dove ci si trova,
qual è lo stato delle cose, che cosa si potrebbe fare meglio e stimare
meglio\footnote{Anche qui l'inizio della spiegazione è parzialmente perso nella
trascrizione (``che in parte si basa su ci si trovano, qual è il loro
status'').}; per poi tornare all'inizio del loop con più dati su cui basare una
stima corretta.

Si può fare \textbf{a ogni sprint}, a seconda della durata, oppure \textbf{ogni
serie di sprint}: ``vedetelo voi''. Cristian suggerisce la seconda opzione
perché ritoccare continuamente il backlog per tagliare o aggiungere feature va
contro il principio del lavorare sempre per incrementare il prodotto. Va da sé
che, se si tratta di un taglio molto importante, \textbf{prima lo si prende e
meglio è}: ``ragionateci bene''.

\subsection{Su che cosa basarsi}

\paragraph{Sulle performance, cioè sulla velocity.} Se la velocity è alta, si
può dire: ``a questo punto possiamo reintrodurre qualche feature che era nei
\emph{nice to have}\footnote{Nella trascrizione ``nice to the epoch'': quasi
certamente \emph{nice to have}.} o che era addirittura stata tagliata''. Se è
bassa, si sistemano già milestone e backlog in base al risultato.

\paragraph{Sui feedback.} A Level Up \emph{L'Inf} ha ricevuto tantissimi
feedback, per lo più da giocatori. Trattandosi di una \emph{vertical slice},
spesso dicevano cose che il team già sapeva e che avrebbe implementato nella
sprint successiva (ma, per questioni di tempo, da qualche parte bisogna
tagliare). Altri, invece, hanno parlato di \textbf{accessibilità}.

\paragraph{Il cambio di camera.} A questi si sono aggiunti i feedback costanti
di Diego sulla camera: più che un ``dovete cambiarla'', un ``perché non provate a
cambiare camera?''. Diego interviene: ``c'è voluto del tempo a convincere
Cristian: otto mesi''. Alla fine l'ha fatto e sta funzionando, anche se ha
pesato molto sul team. Il motivo: il \textbf{twin-stick shooter}\footnote{Nella
trascrizione ``twin steep shooter''.} richiede una grande
\textbf{dimestichezza con i controlli} (muoversi e mirare contemporaneamente con
le due levette) e di per sé esclude gran parte del pubblico casual. I giocatori
più casual dicevano proprio questo: ``non riesco a camminare e mirare
contemporaneamente''. Il team se l'è segnato, ne ha discusso e alla fine ha
accettato il cambio di camera, con tutte le conseguenze.

\subsection{Flessibilità anche nella struttura dei documenti}

Un promemoria sulla \textbf{flessibilità}: più si è stati flessibili all'inizio,
più si riesce a rimettere mano a backlog e produzione in seguito. E questo vale
anche, banalmente, per \textbf{come è strutturato un documento}: se ogni casella
fa riferimento a un'altra casella, che fa riferimento a un'altra due righe
sopra, modificare quel documento (o un'implementazione costruita allo stesso
modo) richiederà molto più tempo. Flessibilità in tutto: nella scrittura dei
documenti e nella produzione (``non lo ripeto, altrimenti ricomincio'').

\subsection{Evitare il perfezionismo}

``A proposito di perfetto'': il consiglio sul \emph{tuning} è \textbf{evitare
in tutti i modi il perfezionismo}. Cristian ci sta ancora lavorando
``tantissimo'' (``anche qui Diego lo sa meglio di me''): è migliorato molto e nei
progetti ci riesce, ma in altre cose pretende ancora troppo. Il perfezionismo
\textbf{va contro la flessibilità e crea inevitabilmente dei blocker}. Lo sta
vivendo in prima persona, non per sé ma con altri membri del team.

\paragraph{Il block out con il foliage.} Il task è ``fai il block out di questa
parte del livello'' (dare l'intero block out di livelli da mezz'ora richiede
comunque tempo, quindi il task riguarda la parte iniziale). E nel block out
Cristian si ritrova il \emph{foliage}, il \emph{landscape} scolpito bene, le
rocce. ``Va benissimo perché ti piace il lavoro\dots{} anzi no, sto proprio
dicendo che non va bene.'' L'approccio di voler consegnare il task nel miglior
modo possibile è comprensibile; il problema è che questo perfezionismo
\textbf{blocca tutti gli altri}, per moltissimi motivi. Se si passano due ore a
mettere il foliage e poi, quando si va in gold a rifinire il livello, si trova un
foliage migliore per il \emph{mood} estetico del gioco, quel lavoro va buttato.
Il perfezionismo crea \textbf{buchi di produzione} che alla lunga diventano
letteralmente mesi di lavoro buttato o da buttare: un peccato per chi lo fa e
per chi lavora con lui.

\paragraph{La pressione sugli altri.} Un'altra conseguenza: il perfezionismo
\textbf{mette pressione a tutti gli altri}. Chi lavora alla seconda parte del
block out, con il task di renderla giocabile, vede che la prima parte ha
foliage, rocce e tutti i props (non l'illuminazione, che negli ambienti aperti
è tendenzialmente globale, salvo situazioni specifiche) e si sente in dovere di
raggiungere quel livello di rifinitura\footnote{Nella trascrizione ``si sentirà
in dovere di raggiliere la pozione della vita'': passaggio incomprensibile, il
senso è ricostruito dal contesto.}, trascurando la parte di gameplay, che invece
nella prima parte era stata fatta. Il perfezionismo di uno ha generato
\textbf{ansia} nell'altro, ed è da evitare in ogni modo.

\paragraph{I criteri di accettazione.} È un altro motivo per cui sprint e task
devono avere un \textbf{goal} e un \textbf{criterio di accettabilità}, per
usare i termini di Scrum: il task si considera fatto quando è \emph{done}, e va
bene così, perché lo si è valutato prima. Tutto ciò che si fa oltre il
\emph{done} è \textbf{tempo perso}. Può andare bene che quello che fai adesso
non dovrai farlo dopo, ma è sbagliato come approccio: se il lavoro è stato fatto
bene prima, si può andare avanti tranquilli, e al passo successivo ci si penserà
dopo.

% ---------------------------------------------------------------------
\section{Domanda: come comunicare un taglio al team}
\label{sec:z3-qa-taglio}

Prima della seconda parte, che Cristian promette più breve, una domanda: sul
piano relazionale ed emotivo, come ha gestito il dover convincere, o comunque
spiegare, al team la necessità del taglio da otto a cinque livelli?

\paragraph{Ricordare come funzionano i feedback.} ``Ottima domanda.''
Emotivamente è stato pesante, perché non tutti reagiscono ai feedback allo
stesso modo. Per questo, a ogni riunione, Cristian ricorda i principi dei
feedback (``so che sembra banale, so che non serve''), in particolare una
clausola che tutti dimenticano: \textbf{il feedback si può accettare o
rifiutare}. Se non si è d'accordo, lo si può ignorare; invece tutti partono dal
presupposto che un feedback vada seguito al 100\%\footnote{Nella prima lezione
zero il relatore invitava ad accettare il feedback ``a prescindere''
(sezione~\ref{sec:z1-ricevere}); lì però ``accettare'' significava ascoltarlo
senza difendersi, non necessariamente applicarlo. Le due indicazioni sono
quindi compatibili: si ascolta tutto, poi si decide che cosa farne.}.

\paragraph{Partire dai dati.} Sembra brutto affrontare una questione emotiva
partendo dai dati, ma l'obiettivo era \textbf{far arrivare loro alla soluzione},
non dargliela. Ed è andata così: sono stati loro a dire ``per adesso facciamo
cinque livelli, e poi, \textbf{se ci riusciamo}, implementiamo gli altri''.
Quel ``se ci riusciamo'' era ciò che Cristian voleva sentire: se fosse rimasto
solo ``ne facciamo cinque'', l'avrebbero vissuto come un limite su se stessi
(``non ho fatto in tempo a farne otto''). Il ``se ci riusciamo'' gli ha
comunicato che avevano capito che \textbf{non era una critica al loro modo di
lavorare}, né un modo per dire ``non siete capaci'': lo aveva sottolineato più
volte, magari non con parole così dirette, ma ``un po' \emph{soft-skillato},
come diciamo noi designer''.

Hanno colto che non c'erano alternative (``otto non li facciamo, punto''), ma
anche le possibilità: facciamo cinque livelli, abbiamo altro tempo,
perfezioniamo ciò che abbiamo fatto; vogliamo aggiungere una nuova modalità di
gioco? Benissimo. Altri due livelli? Vediamo, in base al tempo. Esserci
riuscito lo ha reso ``super fiero''.

% ---------------------------------------------------------------------
\section{Il team management}
\label{sec:z3-team}

\subsection{Senza team non esiste Scrum}

Cristian accelera, per stare nei tempi. Il punto da capire è che
\textbf{Scrum si adatta al team}: quindi il team è la cosa più importante di
tutte, più di Scrum stesso. ``Se io mi adatto al team, il team deve esserci.
Senza team non esiste Scrum.''

\subsection{Una produzione senza stress}

Di conseguenza, una produzione in Scrum \textbf{dovrebbe rendere il team
rilassato}. Come il team di Simone, anche il primo team di Cristian (quello fino
alla \emph{vertical slice}) \textbf{non ha mai fatto crunch}. Il nuovo team sì,
per i grossi cambiamenti già raccontati: la camera è cambiata completamente, e
``cambiare la camera in un gioco è fare un altro gioco''. Cristian aveva già
spostato la scadenza (``invece di finire a marzo finiamo ad aprile''), ma il team
si è sentito comunque in dovere di risolvere prima lo scompenso creato da quei
cambiamenti.

In generale ci si dovrebbe sentire meno stressati per un motivo semplice: chi
guida sta dando i task che lo sviluppatore ha chiesto, verbalmente o dimostrando
di saperli fare. ``So che in due settimane devo fare cinque cose, faccio queste
cinque cose'': non c'è la pressione di dover fare il più possibile. Se poi uno è
perfezionista e passa da cinque cose a quindici perché ``ha tempo'', il problema
(``sembra brutto dirlo'') è suo, nasce nel singolo individuo; e chi guida deve
basarsi sul team, sulle performance del team.

\subsection{Mai la velocity sui singoli}

Da qui un consiglio sulla velocity: \textbf{non calcolate mai, mai, mai la
velocity delle singole persone}, e nemmeno i task fatti dai singoli. Lo si può
fare per sé, ma tenendolo ``chiuso in un'agenda scritta a penna'': se un membro
del team scopre che si stanno valutando le sue prestazioni individuali,
\textbf{si sente giudicato} e soprattutto non sa più come dire che non vuole
fare una cosa.

A Cristian è successo: ``mi hanno sgamato''. Si è dovuto scusare, anche se non
erano scuse: ha detto la verità, cioè che lo faceva più per capire lui la
produzione che per valutare le loro performance. Ma spiegarlo è complicato,
quindi è meglio non farlo: tanto la velocity si basa sulle performance del team,
e i dati individuali non servono. Se li si vuole per sé, per capire a chi dare
un task piuttosto che a un altro (se una persona in un'ora fa dieci
implementazioni e un'altra una, ci si regola quando si selezionano i task),
non bisogna assolutamente farli ricadere sulle persone.

\subsection{Primo aspetto: adattare Scrum al team}

Per una produzione senza stress servono due aspetti chiave. Il primo è che
\textbf{Scrum va adattato al team}: si prende la base di Scrum (quella di Diego,
dei libri, di Simone e di Cristian), la si \textbf{testa}, si \textbf{analizzano}
la velocity e tutto il resto, la si \textbf{tara} (``\emph{daily scrum}? No,
facciamolo ogni tre giorni, ogni settimana, o non facciamolo''), la si
\textbf{riprova}, e si ripete il ciclo per tutta la produzione. Alcuni esempi da
\emph{L'Inf}:
\begin{itemize}
  \item i \textbf{daily scrum} li fanno \textbf{ogni tre giorni}, ``se va
        bene'', e dipende anche dal tipo di produzione;
  \item la \textbf{lunghezza delle sprint} è flessibile (una, due o tre
        settimane) e dipende:
        \begin{itemize}
          \item dalla \textbf{disponibilità del team}: se tutto il team non
                c'è, è inutile fare una sprint di una settimana;
          \item soprattutto dalla \textbf{fase di produzione}: quando si
                implementano script e in una settimana si riescono a
                implementare tantissime meccaniche, la sprint dura una
                settimana, al termine della quale si controlla ciò che è stato
                fatto e si sceglie che cos'altro implementare; quando si passa al
                level design, con lavori molto dispendiosi in termini di tempo,
                le sprint si diluiscono, e infatti in quel momento sono a sprint
                di due settimane, concentrati sul completamento dei
                livelli\footnote{Nella trascrizione ``sull'ultima mento,
                sull'ultima azione, non mi ricordo, scusate. Per completare il
                livello'': forse ``l'ultima milestone''.};
        \end{itemize}
  \item la velocity si calcola \textbf{sui task e sul loro peso, non sulle ore
        lavorative}, come diceva Simone. Soprattutto per degli studenti, che
        non hanno ore fisse. In un ambiente di lavoro in cui tutti sono
        disponibili dalle 9 alle 18 potrebbe cominciare ad avere senso:
        ``secondo me no, ma ve lo dico per correttezza''.
\end{itemize}

\subsection{Secondo aspetto: tutti devono conoscere il mindset}

Il secondo aspetto è che \textbf{tutti i membri del team devono conoscere il
mindset di Scrum}\footnote{A questo punto Cristian chiede a Diego cinque minuti
in più per chiudere bene, e Diego acconsente.}. Per due motivi:
\begin{itemize}
  \item se tutti lo conoscono, il rischio di cadere nel Waterfall o nel
        perfezionismo si riduce, perché ognuno sa che cosa è giusto e che cosa
        è sbagliato;
  \item tutto il team ha \textbf{lo stesso goal}: se la sprint chiede il block
        out, tutti sanno che, secondo la mentalità agile, va bene farlo ma non
        deve essere perfetto. Bisogna fare i \emph{done}, non i
        \emph{perfect}.
\end{itemize}
Come far capire a tutti la mentalità di Scrum? Cristian usa due modi.
\begin{enumerate}
  \item \textbf{Parlarne costantemente}, a ogni riunione, come sta facendo con
        la classe: sia per migliorare lui stesso e farsi entrare il mindset,
        sia per trasmetterlo agli altri.
  \item Il più efficace: \textbf{lo shipping}. Il team riesce a consegnare
        quasi costantemente. Ci sono momenti in cui un problema cade
        all'ultimo, come in quella settimana, con alcuni strumenti
        particolari che li hanno rallentati, ma si risolve tutto. Dimostrando
        di riuscire a \textbf{mantenere le milestone}, il team ha acquisito
        fiducia in lui e ha capito quanto è potente Agile. Nella
        \emph{vertical slice} consegnavano le build addirittura prima di
        quanto richiesto; e la consegna prima di Natale, nota da settembre, è
        stata rispettata nonostante il cambio di camera e tutto il resto. Il
        messaggio: \textbf{non è importante che ciò che si consegna sia
        perfetto, ma che rispetti le condizioni di accettabilità} che ci si è
        dati.
\end{enumerate}
Sono i due metodi che funzionano meglio finora; se ne troverà altri, li
condividerà: ``tutto è in continua evoluzione''.

\subsection{Il post-mortem}

Il \emph{post-mortem} è \textbf{la garanzia di capire se ciò che si sta facendo
è corretto}: serve a vedere come stanno andando il team e la produzione, se il
mindset di Scrum è stato compreso o no, e ad agire di conseguenza.

\textbf{Quando farlo} lo si decide liberamente: a ogni sprint, a ogni gruppo di
sprint, a ogni milestone.

Ciò che Cristian sottolinea è che \textbf{nel post-mortem devono parlare tutti},
anche lo sviluppatore che non vuole mai parlare: bisogna dargli il tempo per
farlo. Un modo è \textbf{organizzarlo in anticipo}: se è tra una settimana, lo si
comunica la settimana prima. Cristian prepara sempre un \textbf{documento della
riunione} prima e lo aggiorna dopo, con una sorta di elenco: sembra burocratico,
ma aiuta a mantenere il focus. Bisogna far capire a ciascuno che tra una
settimana parlerà, e non può sottrarsi (se si assenta ci perde lui). E bisogna
dirgli di \textbf{prepararsi} qualcosa da dire, magari scrivendolo o facendo una
piccola revisione: se ti viene in mente una cosa da dire la settimana dopo e non
la scrivi, finirai per pensare ``sai che c'è, non la dico''; se l'hai scritta,
non puoi dimenticartela, e hai qualcosa di cui parlare. È fondamentale perché
bisogna tenere conto di \textbf{tutti i tipi di personalità}: non solo i
giocatori hanno personalità diverse, anche gli sviluppatori hanno personalità
con bisogni diversi.

\subsection{La sincerità}

Cristian propone di tagliare la parte ``sdolcinata'', ma Diego vuole sentirla
per intero (``non si sfugge al miele''; ``poi mi manderete gli IBAN e vi pagherò
il tempo perso''). Come Diego sa, Cristian cerca sempre di metterci ``qualcosa di
filosofico''. Per \textbf{sincerità} intende che, in qualsiasi momento della
produzione, \textbf{tutti i membri del team devono sentirsi protetti in ciò che
dicono} e liberi di esprimere opinioni e feedback. È importantissima sia nella
produzione sia nel team.

\paragraph{Nella produzione.} Senza sincerità la produzione va a picco; si può
comunque consegnare, non è obbligatoria, ma con la sincerità si realizza un
prodotto migliore. Se a inizio produzione si pensa che una meccanica ``faccia
schifo'' (prende un \emph{player type} completamente fuori target e richiede
troppo lavoro per essere implementata), \textbf{è quello il momento di dirlo}.
Altrimenti ``il karma vuole che tu stesso che pensi questa cosa sarai colui che la
andrà a implementare'', e la farai male, perché non vuoi farla. Se non l'hai
detto, o l'hai detto troppo tardi, ``ti mangerai le mani''.

E il \textbf{``ve l'avevo detto''} va risparmiato: anche se il giorno dopo arriva
la conferma di ciò che si era detto, anche se è umano, va evitato, perché crea
attriti nel team che non si risolvono più (Cristian parla per esperienza
diretta). È lo stesso spirito del \emph{disagree and commit}
(sezione~\ref{sec:z2-disagree}).

\paragraph{In Scrum.} Se uno sviluppatore dice ``questa settimana so fare 15
task'' e ne fa 7, la prima volta passa. La settimana dopo dice di nuovo ``ne
faccio 15, tranquillo'', magari inventando scuse (``quel problema non ce l'ho
più''). Una volta, due volte; \textbf{alla terza bisogna chiedere sincerità}: qual
è il problema? Non hai tempo?

Perché sul tempo, ``ragazzi'', \textbf{è importante dire quanto tempo si ha, non
garantirlo}. Se il team (cinque persone, nel caso di Cristian) dice di avere due
ore al giorno, lui pianifica una produzione da dieci ore al giorno; se il team
dice di avere otto ore al giorno, pianifica di conseguenza (``80
ore''\footnote{Così nella trascrizione; con cinque persone a otto ore al giorno
il totale sarebbe 40 ore, oppure 80 su due giorni. Il senso dell'esempio non
cambia.}), e se poi in realtà le ore sono dieci, gli diranno ``ci dai troppi
task''.

Bisogna essere sinceri \textbf{prima di tutto con se stessi}: se un task richiede
due giorni, si dice due giorni; magari tre, perché un ``tempo cuscinetto'' ben
venga, ma lo si dice per com'è. E non bisogna proteggersi dal timore che ``se
dico la verità nessuno si fiderà più di me'' o ``mi daranno più task'': è compito
dello scrum master, ma di tutto il team, far capire che \textbf{si è tutti sulla
stessa barca}.

\paragraph{Nel team.} Gli attriti generati dalla mancanza di sincerità si
trasformano in \textbf{relazioni migliori} quando c'è sincerità. Le
incomprensioni vanno \textbf{risolte subito}, affrontandole; magari porteranno a
un litigio, se proprio quello è il finale, ma tra persone mature si dovrebbe
riuscire a risolvere in modo pacifico. Per questo, sulla slide, la sincerità è
definita \textbf{``rischiosa ma necessaria''}.

% ---------------------------------------------------------------------
\section{Chiusura della lezione}
\label{sec:z3-chiusura}

Cristian salta la slide di riepilogo (la lezione si potrà rivedere online) e
lascia i suoi contatti: \textbf{Telegram}, dove molti studenti lo trovano già nei
gruppi (in particolare quelli dell'accademia online), oppure l'\textbf{email},
che Diego può fornire. Si rende disponibile per qualsiasi cosa, sulla lezione o
su un'idea da sottoporgli (``ho pensato così, secondo te può funzionare?''),
precisando di non volersi sostituire in alcun modo a Diego. Si scusa per i dieci
minuti finali ``di corsa''.

Diego lo ringrazia per la disponibilità, ricordando che anche lui, da studente,
ha un carico di lavoro notevole. L'obiettivo della lezione era far vivere alla
classe un'esperienza e dare informazioni simili a quelle di chi ha fatto lo
stesso percorso, ma ha \textbf{internalizzato Agile} meglio di altri: tutti hanno
avuto le stesse lezioni, ma Cristian e Simone sono quelli che le hanno capite
meglio. Ci sono molti \emph{insight} utili in entrambi gli interventi, e i punti
chiave di Agile sono stati ripetuti più volte, ``perché alla fine Agile è
comprendere i punti chiave''. ``Grazie ancora, e buon pranzo.''\footnote{La
trascrizione si chiude con la parola ``grazie'' ripetuta diverse centinaia di
volte: è un artefatto tipico della trascrizione automatica su un tratto di
audio silenzioso o di saluti sovrapposti, non un contenuto della lezione.}

% ---------------------------------------------------------------------
\section*{Punti chiave del capitolo}
\addcontentsline{toc}{section}{Punti chiave del capitolo}

\begin{itemize}
  \item \textbf{Agile è un mindset, non un framework}: non ha regole, ha una
        logica da \textbf{internalizzare}, come una lingua che a un certo punto
        si smette di tradurre. I quattro valori non vanno seguiti alla lettera:
        vanno adattati al team e al progetto, e non è detto che Agile serva.
  \item I \emph{framework} (Scrum, Kanban, \dots) hanno regole e si basano
        sull'\textbf{inspect and adapt}, cioè sulle iterazioni. Si parte
        rispettandoli al 100\%, si capisce che cosa non funziona e perché, e
        \textbf{ci si costruisce il proprio}.
  \item \textbf{Meno riunioni possibili}: unire planning e post-mortem, stare in
        25--30 minuti, report da un minuto. \textbf{Nelle riunioni non si
        risolvono i problemi}: si fissa un appuntamento con chi può aiutare. Un
        \textbf{timekeeper} può far rispettare i tempi.
  \item Le \textbf{sprint} (una o due settimane) sono l'unità di tempo: le
        milestone si misurano in sprint, non in mesi.
  \item \textbf{Consistency}: una cerimonia saltata fa crollare le altre a
        effetto domino. Per questo ciò che non funziona per il team (spesso il
        \emph{daily scrum} o il \emph{planning poker}) va \textbf{tagliato},
        non trascinato.
  \item Il \textbf{backlog} si fa \textbf{a scaglioni}, per milestone, partendo
        dalle \emph{epic} e dai documenti (visione, GDD, \emph{versioning} con
        tutti gli ambiti, LDD). Si organizza per feature, per ambito con un
        \textbf{color code} condiviso, per milestone e sotto-milestone o per
        priorità, tutto in un solo strumento.
  \item Ogni task (e ogni sprint) ha un \textbf{obiettivo verificabile} e un
        \textbf{criterio di accettazione}: ``completare il livello 5'' non lo è.
        Le \emph{user story} aiutano, ma non sono obbligatorie.
  \item \textbf{Mai stimare in ore}: si stima in modo \textbf{relativo}, con la
        sequenza di \textbf{Fibonacci}, partendo dagli estremi. I valori sono
        arbitrari ma \textbf{consistenti}, e quindi utilizzabili.
  \item La \textbf{velocity} (somma dei pesi dei task \emph{done} per sprint,
        confrontata con lo stimato) \textbf{prevede il futuro}: dice se gli
        strumenti di stima funzionano, se lo \emph{scope} è raggiungibile e come
        sta il team. Permette di tagliare per tempo e di arrivare in fondo
        \textbf{senza crunch}. L'andamento a dente di sega è normale.
  \item Le prime sprint \textbf{si sbagliano tutte}, ed è giusto: la quantità si
        tara sulle disponibilità e sulle capacità attuali, prevedendo i
        \textbf{tempi di apprendimento}.
  \item Si può \textbf{progettare il ritmo dello sviluppo}: dopo periodi di
        lavoro invisibile, sprint che danno soddisfazione (nuovi nemici, VFX,
        \emph{juiciness}) per il morale del team.
  \item I \textbf{blocker} vanno trovati subito e \textbf{analizzati alla
        radice} prima di essere risolti: la soluzione più rapida (i muri
        invisibili) può costare settimane.
  \item I \textbf{feedback costanti} abbattono la produttività: vanno raccolti in
        momenti specifici. Il feedback si può accettare o rifiutare.
  \item \textbf{Flessibilità} in tutto, soprattutto negli elementi innovativi:
        non prevedere tutto nel dettaglio (è Waterfall), non fermare mai la
        produzione per riscrivere backlog e documenti, lavorare sempre per
        \textbf{incrementare il prodotto}.
  \item Il \textbf{perfezionismo} crea blocker, lavoro buttato e ansia negli
        altri: si fanno i \emph{done}, non i \emph{perfect}.
  \item \textbf{Il team viene prima di Scrum}. La velocity non si calcola mai sui
        singoli. Tutti devono conoscere il mindset: parlarne a ogni riunione e,
        soprattutto, \textbf{consegnare} rispettando le milestone.
  \item Nel \textbf{post-mortem} parlano tutti: va annunciato prima, con un
        documento, chiedendo a ciascuno di prepararsi.
  \item Alla base di tutto ci sono la \textbf{fiducia} reciproca e la
        \textbf{sincerità}: dire subito ciò che non va, dire quanto tempo si ha
        (non garantirlo), evitare il ``ve l'avevo detto''. La sincerità è
        \textbf{rischiosa ma necessaria}.
\end{itemize}

%    \part{Il corso di Game Design}
%    \renewcommand{\thechapter}{\arabic{chapter}}
%    \setcounter{chapter}{0}
%    % =====================================================================
%  Capitolo 1 — Introduzione al corso e al mestiere del game designer
%  Fonte: trascrizione della lezione GD1 (Game Design, primo anno)
%  Pacchetti richiesti nel preambolo del libro:
%    \usepackage[italian]{babel}, \usepackage[T1]{fontenc},
%    \usepackage[utf8]{inputenc} (solo pdfLaTeX), \usepackage{booktabs}
% =====================================================================

\chapter[Il mestiere del game designer]{Il mestiere del game designer: introduzione al corso}
\label{cap:introduzione-game-design}

Questo primo capitolo fa da soglia. Prima di affrontare metodi, framework e
documenti, è necessario capire \emph{in che modo} si impara il game design,
\emph{con chi} lo si fa, \emph{con quali strumenti} e, soprattutto,
\emph{che cosa fa concretamente} un game designer nel suo lavoro quotidiano.
Il capitolo segue l'impostazione della lezione inaugurale del corso: si parte
dal contesto formativo (il corso come simulazione di lavoro), si passa al
metodo di apprendimento (l'\emph{inquiry-based learning}), si elencano
letture e risorse di riferimento, e si arriva infine a una prima definizione
operativa del ruolo del designer, dei limiti entro cui lavora e delle sue
possibili specializzazioni.

Una premessa terminologica: il game design è una disciplina giovane e quasi
tutta la sua letteratura è scritta in inglese. Per questo motivo molti termini
tecnici verranno lasciati in lingua originale (\emph{pipeline},
\emph{pitching}, \emph{tuning}, \emph{feature team}, \dots): sono i termini che
si useranno nel lavoro, ed è importante imparare a riconoscerli e a usarli
correttamente fin da subito.

% ---------------------------------------------------------------------
\section{Il contesto: un corso come simulazione di lavoro}
\label{sec:simulazione-lavoro}

\subsection{Il principio fondamentale}

Il primo concetto da fissare è che il corso \textbf{si presenta come una
simulazione del lavoro}. Dal momento in cui si entra in aula si viene trattati
come \emph{potenziali professionisti dell'industria}: non soltanto studenti,
ma lavoratori effettivi di uno studio (l'accademia), con una gerarchia, dei
referenti, dei compiti e delle scadenze.

Ne derivano alcune conseguenze pratiche:
\begin{itemize}
  \item il docente non assegna un ``compitino'' da risolvere per accumulare
        punti o presenze: ogni \emph{task} è pensato per far dedicare tempo
        all'apprendimento di qualcosa di concreto;
  \item non consegnare un task non danneggia il docente (che ``vive ancora
        domani e viene pagato ugualmente''), ma lo studente, che perde
        l'occasione di mettere le mani su un argomento e diventare più abile;
  \item docenti e studenti devono trattarsi \textbf{da colleghi}, perché è
        molto probabile che un domani lo diventino davvero.
\end{itemize}

\subsection{Perché l'impegno conta più che mai}

Il lavoro del game designer è descritto come \emph{molto romantico, ma anche
molto duro}, e sempre più duro col passare degli anni. Durante il periodo del
Covid l'industria ha visto grandi investimenti; quei ``rubinetti'' sono stati
chiusi da tempo, il rapporto con i publisher è diventato più complesso e le
opportunità lavorative non sono più quelle di tre o quattro anni fa. Questo
significa:
\begin{itemize}
  \item più competizione per ogni posizione;
  \item aspettative più alte: ci si aspetta che i candidati siano ancora più
        preparati e sicuri delle proprie competenze;
  \item un unico modo per diventarlo: \textbf{lavorare}. Più tempo si dedica
        oggi ad apprendere il mestiere, più facilmente si sarà assumibili
        domani e più agevole sarà il lavoro di tutti i giorni.
\end{itemize}

\subsection{Sbagliare qui per non sbagliare fuori}

Il mantra del corso è: \textbf{quanto più provate qui, anche in modo goffo e
sbagliando, tanto meno sbaglierete là fuori}. In accademia, per quanto la
simulazione sia realistica, nessuno toglie lo stipendio e nessuno licenzia da
un momento all'altro. Nell'industria sì, e un errore può costare molto denaro
all'azienda, ancora di più se l'azienda è la propria.

\subsection{Il valore del fallimento}

Il fallimento, se affrontato con la mentalità corretta, è probabilmente
\emph{la cosa più vantaggiosa} che si possa fare durante la formazione:
\begin{itemize}
  \item un successo (``avete superato questa deadline'') conferma le proprie
        abilità, ma insegna relativamente poco;
  \item dopo $n$ fallimenti si impara moltissimo, e i successi arrivano;
  \item è fondamentale non demoralizzarsi: il feedback negativo non è mai un
        giudizio sulla persona, ma su ciò che è stato consegnato.
\end{itemize}

Il docente riconosce apertamente che dedicare tempo ed energia a un progetto e
vederlo rifiutato, o vederlo non uscire mai, è demotivante: è successo a lui e
lo ha visto accadere anche su progetti molto grandi, con team numerosi. Il
punto però non è ciò che si è ottenuto materialmente, ma \textbf{quanto si è
cresciuti}. Da ogni esperienza ci si portano a casa due cose:
\begin{enumerate}
  \item \textbf{l'esperienza}: quanto si è diventati bravi, veloci,
        praticamente abili nell'applicare un framework, quanto si è cresciuti
        come designer;
  \item \textbf{il rapporto umano}: nel lavoro si ``suda sangue'' fianco a
        fianco con altre persone, e questo crea legami molto forti che
        portano avanti l'industria nel senso più positivo del termine: nascono
        progetti privati, ci si porta da un'azienda all'altra, si formano team
        con basi solide.
\end{enumerate}

Va evitato anche l'eccesso opposto: la cultura del \emph{grinding} si è
appropriata di questi discorsi in modo distorto. Non bisogna diventare robot.

% ---------------------------------------------------------------------
\section{Organizzazione del corso e gerarchia}
\label{sec:gerarchia-corso}

\subsection{La struttura gerarchica}

La struttura dei referenti del corso ricalca volutamente una
\textbf{gerarchia tipica dell'industria}:
\begin{itemize}
  \item al vertice c'è il \textbf{direttore del corso} (Diego Ricciuti), che
        corrisponde al \emph{creative director}: supervisiona la visione del
        progetto (il concetto di ``visione'' verrà approfondito più avanti);
  \item sotto di lui ci sono i \textbf{referenti di dipartimento}, uno per
        disciplina, con un corrispettivo per ciascuna sede (Milano e Roma);
  \item in fondo ci sono gli studenti, che occupano il primo gradino di
        un'azienda: l'\emph{intern} (stagista) o il \emph{junior}, il passo
        che precede il lavoratore ``regolare''.
\end{itemize}

\begin{table}[htbp]
  \centering
  \begin{tabular}{@{}p{0.34\textwidth}p{0.29\textwidth}p{0.29\textwidth}@{}}
    \toprule
    \textbf{Disciplina} & \textbf{Milano} & \textbf{Roma} \\
    \midrule
    Direzione (\emph{creative director}) & \multicolumn{2}{p{0.58\textwidth}}{Diego Ricciuti} \\
    Game design      & Mattia Podini        & Giacomo ``Jack'' Danna \\
    Level design     & Stefano Zocchi       & Andrea Dispoli \\
    Narrative design & \multicolumn{2}{p{0.58\textwidth}}{Daniele Fusetto (lezioni condivise), Giacomo Masi} \\
    Programmazione (Unreal) & Alessandro Pregnolato & Matteo Velardi \\
    \bottomrule
  \end{tabular}
  \caption{Referenti del corso per disciplina e sede (nomi come riportati in lezione).}
  \label{tab:referenti}
\end{table}

Per il primo anno, il referente diretto degli studenti di game design è il
docente di game design: è a lui che ci si rivolge per dubbi e domande. Altre
figure utili sono il direttore didattico (severo, ``giustamente severo'', da
contattare solo per questioni particolarmente importanti), l'amministrazione
(contratti e documentazione) e la segreteria.

\subsection{I canali di comunicazione}

\begin{description}
  \item[Discord] è il canale per le comunicazioni ufficiali del corso e per
        le domande che non richiedono urgenza immediata; può avere qualche
        ritardo.
  \item[E-mail] è il canale per tutto ciò che è formale o ``ancora più
        ufficiale'' e per le comunicazioni rapide e immediate.
  \item[Telegram] esistono residui di vecchi canali precedenti a Discord, da
        evitare.
\end{description}

Le domande che valgono per tutti vanno fatte in classe (anche da un portavoce
a nome del gruppo); quelle singole e personali via Discord o, se formali, via
e-mail. Se il docente non risponde a una domanda c'è un motivo, ad esempio
ragioni legali, e lo dirà esplicitamente.

\subsection{I tre moduli e la settimana tipo}

Le tre discipline principali del primo anno (e non solo) sono:
\begin{itemize}
  \item \textbf{Game Design} (tendenzialmente il lunedì);
  \item \textbf{Level Design} (tendenzialmente il mercoledì);
  \item \textbf{Programmazione} (tendenzialmente il venerdì).
\end{itemize}
Nel primo mese le lezioni sono tutte di game design (lunedì, mercoledì e
venerdì). Sono previste anche lezioni di narrativa, ma il piano formativo si
aggiorna ogni anno: le ore di narrativa sono state ridotte a favore di quelle
di programmazione. Le poche lezioni di narrativa daranno quindi dei
\emph{cenni}, non la preparazione di un narrative designer. Ciò che serve
davvero è imparare \textbf{come comunicano gli altri dipartimenti} e come
lavorare con persone che fanno un lavoro diverso ma complementare al proprio.

\subsection{Le due anime delle lezioni di game design}

\begin{description}
  \item[Design methods (metodologie)] lezioni prevalentemente teoriche, ma con
        una buona componente pratica, in cui si studiano metodologie di design
        e i cosiddetti \emph{framework} dell'industria.
  \item[Documentazione e pitching] il cuore del lavoro quotidiano del
        designer.
\end{description}

\paragraph{La documentazione.}
Il docente si definisce ``il più grande nerd della documentazione'' che si
possa incontrare. È una parte dura, perché a pochi piace documentare, ma il
lavoro del designer \emph{è principalmente questo}: scrivere documenti e, da
quei documenti, passare alla pratica e all'implementazione. Chi non sa
documentare bene parte male. Quando non si è presenti, il documento è il
proprio \textbf{biglietto da visita}: qualcun altro lo leggerà, quindi va
scritto nel modo più chiaro e comprensibile possibile.

Un esempio concreto del perché la documentazione conta: CD Projekt Red ha
sviluppato un proprio engine proprietario, ma buona parte della sua
documentazione è stata ``buttata all'acqua''. Chi deve usarlo o lo sa già
usare, o ha competenze di un altro engine da trasferire, oppure è in
difficoltà. Si documenta \textbf{perché altri possano leggere}.

\paragraph{Il pitching.}
Il \emph{pitching} è l'esposizione, la presentazione orale. La parte creativa
del team presenta spesso:
\begin{itemize}
  \item le proprie idee ai colleghi;
  \item il gioco a convention ed eventi;
  \item il pitch del gioco agli investitori, per ottenere i fondi necessari a
        realizzarlo.
\end{itemize}
Saper parlare e comunicare chiaramente le idee è fondamentale. In sintesi:
\emph{il game designer parla}, e bisogna farsene una ragione.

\subsection{Altre componenti del corso}

\begin{description}
  \item[Gioco critico e playtest] la classe viene divisa in due gruppi: uno
        gioca, l'altro osserva. Chi osserva non applaude ``come le
        scimmiette'': annota, segna e analizza il \emph{playtest}, esattamente
        come si fa nell'industria. Il playtest fa parte delle mansioni del
        designer, ma esiste anche come specializzazione autonoma.
  \item[Pipeline] termine che indica tutto ciò che riguarda i processi
        formalizzati dell'industria; di solito si intende la \emph{pipeline di
        produzione}: come si divide il team, qual è la sua gerarchia, come il
        \emph{producer} gestisce le deadline. In breve: \emph{come} e
        \emph{quando} si fanno le cose. Al primo anno è meno rilevante (si
        dovrà ``nuotare'' un po' da soli); dal secondo anno si vedranno le
        metodologie effettivamente applicate ai progetti digitali. Si
        scoprirà quanto sia difficile lavorare a un progetto digitale, e
        ancor più fisico, senza un piano di sviluppo comprovato, come quelli
        usati per qualsiasi software in qualsiasi parte del mondo.
  \item[Engine] il motore di riferimento è \textbf{Unreal Engine}
        (sezione~\ref{sec:engine}).
  \item[Programmazione] principalmente tramite \emph{Blueprint} in Unreal,
        più eventuali sistemi di versionamento; il corso parte da zero.
  \item[Progetti pratici] probabilmente la cosa più importante. Al primo
        anno i progetti di gruppo restano circoscritti al corso di design,
        perché mancano ancora alcune nozioni (come comunicare, affrontare
        problemi, dare e ricevere feedback). In seguito si allargheranno a
        progetti digitali, molto più difficili di quelli fisici, con membri
        di altri team: un rapporto non unilaterale, ma ``covalente''.
\end{description}

% ---------------------------------------------------------------------
\section{Il metodo: l'inquiry-based learning}
\label{sec:inquiry}

\subsection{Che cos'è}

Il corso è costruito secondo l'\textbf{inquiry-based learning}
(apprendimento basato sull'indagine), un approccio pedagogico nato intorno
agli anni Sessanta, che si richiama a pensatori come Dewey e Piaget. In
generale, le discipline da cui il design attinge sono quasi sempre la
sociologia, la psicologia e, talvolta, l'economia.

Il funzionamento è un ciclo:
\begin{enumerate}
  \item il docente fornisce metodologie e framework dell'industria durante
        la lezione;
  \item lo studente, fuori dalle ore di lezione, lavora sui \emph{task}:
        mette a terra le nozioni, le maneggia, le assorbe meglio;
  \item da questo lavoro nascono nuovi dubbi;
  \item i dubbi tornano in aula: lo studente chiede, il docente risponde
        (ammesso che dare la risposta sia utile, e ammesso che una risposta
        esista).
\end{enumerate}

\subsection{Non sempre esiste una risposta giusta}

Non tutto ciò che si studia ha ``la risposta migliore possibile in assoluto
nel vuoto''. Esistono metodologie, limiti e framework che aiutano il lavoro,
ma non sono fissi: cambiano da azienda ad azienda, da lavoro a lavoro, da
progetto a progetto, e a volte vengono stravolti. Sta alla
\textbf{sensibilità del designer} capire quale strumento o framework usare in
un dato momento, e questa sensibilità si impara solo con l'esperienza. Nel
lavoro nessuno dirà ``usa questa cosa'' o ``è meglio quest'altra'': sarà il
designer a scegliere.

Anche l'esperienza del docente ha confini: sei-sette anni di lavoro, un modo
personale di approcciare il design e un certo insieme di progetti (ad esempio
nessun MMO). Oltre quei confini si cercano risposte nella ricerca, studiando
come aziende con esperienza specifica (ad esempio Riot) hanno affrontato
problemi simili e hanno condiviso le loro soluzioni, sempre valutando quali
metodi siano adatti alla dimensione del team, al tipo di progetto e al tempo
disponibile.

\subsection{Sensibilità, non memoria}

Non interessa imparare a memoria un framework o la teoria di un autore. Certe
terminologie vanno conosciute, ma per ragioni di comunicazione e chiarezza.
Ciò che conta è sviluppare la \textbf{capacità di farsi domande}, di farle
agli altri, di fare ricerca e di trovare una propria risposta. È esattamente
ciò che si fa nel lavoro: davanti a qualcosa che non si sa fare, l'abilità
del designer sta nel cercare come qualcun altro ha risolto un problema simile.

\subsection{Le lezioni come corsi di aggiornamento}

Le lezioni frontali simulano i \emph{corsi di aggiornamento} aziendali: le
sessioni che un lead di design tiene al suo team di design (o il lead di
level design al suo team, e così via). Alla fine dell'aggiornamento c'è quasi
sempre un lavoro da fare, un obiettivo da portare a terra:
\textbf{l'obiettivo della settimana}. Si lavora di settimana in settimana, e
il perché verrà chiarito in seguito.

Struttura tipica di una lezione:
\begin{enumerate}
  \item breve sessione di domande e risposte sulla lezione precedente;
  \item analisi di eventuali esercizi particolarmente critici;
  \item lezione vera e propria;
  \item discussione di gruppo, se c'è tempo;
  \item assegnazione del task della settimana.
\end{enumerate}
Il docente tende a ``parlare e straparlare'', ad aprire parentesi e a saltare
da un argomento all'altro: l'unico modo per fermarlo è che siano gli studenti
a interromperlo con domande.

\subsection{La gerarchia della competenza}

L'inquiry-based learning può essere rappresentato come una scala di
competenza:
\begin{enumerate}
  \item la sola lezione frontale porta a \textbf{conoscenze incomplete},
        perché manca la parte pratica e di ricerca;
  \item il lavoro a casa (task, ricerca) porta a una competenza \textbf{più
        completa}, non sbagliata in sé, ma ancora parziale e complementare ad
        altre;
  \item la \textbf{condivisione} dei task e dei feedback con gli altri porta
        idealmente alla competenza più completa possibile;
  \item il picco della competenza è la sua \textbf{applicazione nei
        progetti} (fisici o digitali, di gruppo): ``con le cose che avete
        imparato e ricercato, riuscite a fare un gioco?''. Quando la risposta
        diventa positiva, o comunque si inizia a navigare nei processi di
        design per realizzarne uno, quella è la vittoria.
\end{enumerate}
Non si parte subito creando giochi: si parte dall'analisi, dal design e dalle
regole generali di design, prima ``nel vuoto'', poi contestualizzando i
framework, infine unendo i puntini nei progetti.

Rimarranno comunque delle lacune, per tutti, docente compreso, perché il game
design \textbf{non è una scienza esatta}, anzi, a rigore non è affatto una
scienza. Prende in prestito metodologie dalle discipline scientifiche, ma
contiene una componente imprevedibile e incalcolabile: \textbf{la
produzione}. Si possono applicare perfettamente le regole di un framework, e
quel framework può comunque non essere adatto al gioco, al proprio modo di
sviluppare o al tempo di produzione disponibile (alcuni framework richiedono
una ricerca e un'attenzione che il tempo non concede).

C'è poi la componente creativa, la parte ``romantica'' del lavoro: esiste, ma
è \textbf{metà del lavoro}. L'altra metà è ricerca, metodo scientifico al
100\%. Dagli stessi dati $n$ designer diversi, con sensibilità diverse,
possono ottenere $n$ risposte diverse. Non tutte saranno valide: vanno
applicate, messe a terra, e questa scrematura fa parte del lavoro.

\subsection{I livelli dell'apprendimento}

Un secondo modello, a gradini, descrive i livelli di apprendimento (è la
cosiddetta tassonomia SOLO):
\begin{description}
  \item[Pre-strutturale] il docente dà il compito e lo studente lo risolve,
        possibilmente il prima possibile. È esattamente ciò che \emph{non} si
        vuole.
  \item[Uni-strutturale] si affronta una sola disciplina: è il primo passo
        necessario, ed è il motivo per cui nel primo mese ci sono solo lezioni
        di game design.
  \item[Multi-strutturale] entrano in gioco altri dipartimenti: level design,
        game design, programmazione e così via.
  \item[Relazionale] si capisce che alcune cose valgono per più
        dipartimenti: alcuni framework di narrativa derivano dal design o
        possono essere usati per fare design; le regole di documentazione si
        applicano anche ad altri dipartimenti; le regole di level design si
        mescolano con quelle di game design. Dopotutto, per fare un gioco
        servono molti dipartimenti.
  \item[Astratto esteso] il punto finale: muoversi autonomamente e capire
        cosa usare e quando, sulla base della propria esperienza. L'unico modo
        per arrivarci è fare pratica e fare progetti.
\end{description}

% ---------------------------------------------------------------------
\section{Task, feedback e regole del corso}
\label{sec:task-feedback}

\subsection{I task come opportunità}

I task non sono strettamente obbligatori (esiste un numero minimo di
consegne), ma sono \textbf{opportunità}. Quanto e come ci si lavora dipende
dallo studente. Un task fatto in cinque minuti, o in cui si è stati ``furbi'',
non verrà letto né riceverà feedback.

Domande da porsi davanti a ogni task:
\begin{enumerate}
  \item \textbf{Perché} mi è stato assegnato? Qual è l'obiettivo?
  \item Come ci lavoro per ottenere un risultato \textbf{che abbia senso}? Se
        la ricerca e i dati indicano una direzione, è inutile andare in
        un'altra.
\end{enumerate}

\subsection{Quanto lavorare: un calcolo indicativo}

Se il corso è una simulazione di lavoro, vale il carico di un lavoro normale:
\begin{itemize}
  \item un lavoro a tempo pieno comprende da 30 a 45 ore settimanali; si
        considera una media di \textbf{40 ore};
  \item si sottraggono circa 9 ore di lezione: restano $\approx 31$, diciamo
        \textbf{30 ore};
  \item con circa 3 task a settimana: $30 / 3 = $ \textbf{circa 10 ore per
        task}.
\end{itemize}
Il numero può variare, ma rende l'idea: dedicare un'ora a un task e poi fare
altro non fa guadagnare nulla a nessuno. Lo studente ``butta soldi'', il
docente risparmia tempo leggendo il task in un minuto (o non leggendolo).
Dedicando invece attenzione, ricerca, cura nella formattazione e
sperimentazione, e facendolo con cognizione di causa, si impara.

\subsection{Niente voti, solo feedback}

Ad eccezione del corso di programmazione, non ci sono voti (né ``ottimo'', né
10, né 100). Ci sono \textbf{feedback}:
\begin{itemize}
  \item in game design, generalmente sulla documentazione;
  \item in level design, sui progetti pratici.
\end{itemize}
Il docente indica cosa va bene e cosa va modificato, concentrandosi in
particolare su ciò che va cambiato. Il feedback \textbf{non è mai un giudizio
sulla persona}, ma su ciò che è stato consegnato nel concreto.

\subsection{La condivisione dei task}

Tutti i task consegnati vengono \textbf{condivisi} con la classe, e possono
essere letti solo dopo aver ricevuto i feedback del docente. Si cresce così
non solo con i propri feedback, ma anche con quelli dati agli altri. Regola
importante: ciò che il docente dice a uno studente \textbf{vale per tutti}. Se
a una persona viene detto ``questa cosa non va mai più fatta'', da quel
momento vale per l'intera classe.

\subsection{Non copiare}

Copiare i task degli altri è una delle cose peggiori che si possano fare, non
per una questione morale, ma perché annulla il senso del percorso. Si può
dedicare a ogni task il tempo e l'attenzione che si ritengono opportuni, ed è
normale che col tempo ci si specializzi (preferendo, ad esempio, il level
design o la programmazione tecnica) e si investa meno dove si sente meno
necessità. In questi casi l'alternativa corretta è:
\begin{itemize}
  \item non consegnare;
  \item studiare ciò che hanno fatto gli altri;
  \item eventualmente riconsegnare in seguito.
\end{itemize}
Ci saranno momenti in cui si verrà sommersi dalle cose da fare: bisogna
imparare a scegliere.

\subsection{Il ``non leggo''}

Il docente ha piena facoltà di non dare alcun feedback a un task (il
cosiddetto \textbf{``non leggo''}) se:
\begin{itemize}
  \item non si è lavorato con attenzione;
  \item non si rispettano le indicazioni date;
  \item è stata violata una regola estremamente importante.
\end{itemize}
Un task ``non letto'' equivale a un task non consegnato e non entra nel
conteggio. Attenzione: un task \emph{perfetto} e un task che \emph{fallisce}
nel comprendere o applicare un framework sono \textbf{ugualmente validi},
purché ci si sia impegnati e si sia dedicato tempo, pazienza e ``sudore''.

\paragraph{Riconsegna.} Un task non letto, in ritardo o che si vuole
migliorare può essere riconsegnato e verrà conteggiato, a patto che:
\begin{itemize}
  \item i problemi segnalati siano stati risolti;
  \item sia stato fatto ulteriore lavoro, integrando anche tutto ciò che nel
        frattempo è stato detto agli altri colleghi nei loro feedback.
\end{itemize}
Il task va sviluppato con le competenze che si hanno nel momento della
riconsegna, non con quelle del momento della prima consegna.

\subsection{Requisiti per superare l'anno}

\begin{itemize}
  \item consegnare almeno l'\textbf{80\% dei task} assegnati (soglia
        malleabile);
  \item partecipare almeno all'\textbf{80\% delle lezioni}, in presenza o
        da remoto, preferibilmente in presenza; i ritardi vengono valutati
        caso per caso;
  \item comportarsi da persone corrette e rispettose (è stata firmata una
        carta dello studente): non urlare, non alzare le mani, non
        vandalizzare, non compiere atti illegali.
\end{itemize}
Impegnarsi di più porta inoltre a essere selezionati per varie opportunità.

\subsection{Scadenze}

\begin{itemize}
  \item la scadenza di un task è \textbf{sempre prima dell'inizio della
        lezione} successiva indicata;
  \item se la classe, come team, mantiene comportamenti scorretti, il docente
        può \textbf{anticipare le deadline} per tutti.
\end{itemize}

\subsection{Accessibilità e DSA}

Gli studenti con disturbi specifici dell'apprendimento o neurodivergenze
(dislessia, ADHD, autismo, ecc.) sono invitati a segnalarlo. Quando si
affronteranno le regole di formattazione e documentazione (errori
grammaticali, mancanze, incoerenze) se ne terrà conto e si potranno trovare
dei compromessi, pur ricordando che nel mondo del lavoro nessuno farà
``sconti''. La documentazione può essere inviata all'amministrazione o alla
segreteria, che la comunicherà a tutti i docenti, ma il consiglio è parlarne
anche direttamente con ogni singolo docente.

% ---------------------------------------------------------------------
\section{Le competenze trasversali}
\label{sec:soft-skill}

\subsection{I giochi si fanno in gruppo}

Si possono fare giochi da soli, ed è un'impresa ammirevole; ma lavorando da
soli \textbf{tutto ciò che si guadagna in libertà si perde in produttività}.

\paragraph{Il caso Stardew Valley.} \emph{Stardew Valley} è stato realizzato
da una sola persona, Eric Barone, in anni di sviluppo e di sacrifici: faceva
un altro lavoro per mantenersi (in un cinema) e di notte programmava, faceva
l'arte e componeva la musica. Il risultato è stato straordinario e lo ha reso
ricco, ma si tratta di un caso estremamente raro. Quanti giochi simili, nati
nelle stesse condizioni, sono falliti? Molti di più.

\paragraph{Il rischio d'impresa.} Lavorare da soli non è sbagliato di per sé
e ha aspetti molto positivi: si ha il privilegio di creare esattamente ciò che
si vuole. Ma ci si fa carico anche di tutto il \textbf{rischio d'impresa}: se
si lavora a un gioco per cinque anni e poi il gioco va male, resta il
vantaggio di averlo pubblicato, ma il dopo è sperare di vivere di quella
release o di finanziare con essa il gioco successivo. Le grandi aziende hanno
gli stessi problemi, ma li distribuiscono su più persone e più responsabilità.

L'obiettivo di un professionista è \textbf{saper lavorare in diverse
condizioni}, e la più comune è quella del dipendente (o del freelance /
libero professionista con partita IVA): si viene assunti per le proprie
capacità da un'azienda, piccola o grande che sia, e si lavora con altre
persone.

\subsection{Essere persone con cui si vuole lavorare}

Se le persone non vogliono lavorare con voi, quello è \textbf{il primo motivo
di licenziamento}. E se non si viene licenziati, il clima nell'azienda
diventa terribile: la produttività crolla, la motivazione crolla. Le
\emph{soft skill}, cioè l'insieme di abilità che rendono una persona
piacevole con cui lavorare, vanno coltivate e diventano sempre più importanti
col passare degli anni (relativamente meno ora, molto di più quando i team
saranno più grandi). Anche chi è un ``lupo solitario'' deve imparare a:
\begin{itemize}
  \item lavorare con altre persone;
  \item comunicare con altre persone;
  \item trovare un punto in comune da cui partire per fare le cose.
\end{itemize}

\subsection{Indipendenza e proattività}

Sono le altre due parole chiave del corso:
\begin{description}
  \item[Indipendenza] spesso si verrà lasciati da soli, non per cattiveria ma
        per sfida: quando si cerca una risposta da soli, ciò che conta non è
        tanto la risposta quanto \textbf{il processo}. Una risposta data dal
        docente non vale niente.
  \item[Proattività] chiedere, domandare, agire da soli senza aspettare
        direttive, perché spesso quelle direttive non ci sono, a maggior
        ragione nei team piccoli.
\end{description}
Tutto questo si concretizza nel parlare, nel presentare e nel fare
\emph{post-mortem} (l'analisi a posteriori di un progetto concluso, che verrà
approfondita più avanti), e in generale nell'attenzione alla crescita
personale.

\subsection{Fare domande}

Il corso insiste sulle domande: alzare la mano, intervenire anche ``a
bomba'', fermare il docente ogni volta che qualcosa non è chiaro. Se una
persona non conosce qualcosa, molto probabilmente c'è almeno un'altra persona
che non la conosce. Il docente può dare per scontata la conoscenza di termini
che non lo sono, ed è del tutto normale, soprattutto all'inizio.

% ---------------------------------------------------------------------
\section{Strumenti: engine, Unreal e materiali}
\label{sec:engine}

\subsection{Unreal Engine e gli altri motori}

Il motore di riferimento del corso, usato in level design, è \textbf{Unreal
Engine}: ``un gran mattone'', ma che dà a un designer molte possibilità. Il
vantaggio principale è che permette a un designer non pratico di
programmazione di \textbf{prototipare tramite i Blueprint} (il sistema di
scripting visuale di Unreal). Questo non impedisce di passare a Unity o ad
altri engine, e viceversa: il docente stesso ha studiato e lavorato a lungo
con Unreal, ma la sua ultima esperienza professionale (come level designer) è
stata su Unity, un engine che non conosceva, a dimostrazione che il salto da
un engine all'altro è possibile.

Panorama dei motori citati:
\begin{itemize}
  \item \textbf{Unreal Engine} e \textbf{Unity}: i due più diffusi, con
        modelli di sorgente ``più o meno'' aperto;
  \item \textbf{Godot}: engine open source;
  \item \textbf{CryEngine}: proprietario di Crytek;
  \item \textbf{Snowdrop}: engine di Ubisoft;
  \item Rockstar usa un proprio engine proprietario;
  \item alcune aziende usano versioni \textbf{personalizzate} di engine
        esistenti: ad esempio Milestone, nei giochi di corse, usa un Unreal
        customizzato;
  \item CD Projekt Red ha un engine proprietario (con i problemi di
        documentazione già citati).
\end{itemize}
In sintesi: ogni azienda può avere un engine proprietario oppure una versione
modificata di un engine commerciale, ma si parte sempre da una base.

\subsection{Programmazione: cosa ci si aspetta da un designer}

Non è richiesto alcun bagaglio pregresso di programmazione: il corso parte da
zero. Gli studenti sono stati selezionati (il corso è a numero chiuso) per
competenze analitiche e critiche sopra la media. Anche nel mercato del lavoro,
per una posizione di game designer generalmente \textbf{non ci si aspetta
competenza di scripting} (C++, Java, C\#, ecc.), a meno di voler fare il
\emph{technical designer}, e anche lì varia da azienda ad azienda.

\subsection{Versione e preparazione}

Indicazioni pratiche emerse in aula:
\begin{itemize}
  \item scaricare Unreal Engine; la versione consigliata per il corso è
        stata indicata nella \textbf{5.6} (con conferma successiva). La
        stabilità dipende da cosa si deve fare, da quali tool e plugin si
        usano: i problemi nascono spesso dai plugin. Nel passaggio da Unreal 4
        a Unreal 5, ad esempio, il sistema \emph{Landscape} era inizialmente
        inutilizzabile, mentre in Unreal 4 funzionava molto bene;
  \item è consigliato installare tramite \textbf{Epic Games Launcher}, più
        comodo: permette anche di riscattare giochi gratuiti periodici e di
        accedere ai pacchetti di \textbf{Fab} (il marketplace di asset
        acquisito da Epic), che ogni mese offre alcuni asset gratuiti utili
        per esercitarsi nel level design (ad esempio per fare \emph{speed
        level design});
  \item molti pacchetti di UI o di sistemi di interazione sono utili per
        prototipare al volo, ma spesso inutilmente complicati;
  \item un plugin di \textbf{materiali di prototipazione a griglia} (adattivi,
        con metriche, allineabili alla griglia, con materiali già pronti e
        colori modificabili, dal costo di circa 5 euro) è indicato come
        ``salvavita'' per il level designer; se necessario, i pacchetti
        indispensabili vengono forniti;
  \item portare \textbf{sempre il portatile} alle lezioni di level design
        (da novembre, a partire dall'interfaccia dell'engine); alle lezioni
        di game design basta che lo porti almeno metà della classe, per i
        lavori a coppie. Serve comunque qualcosa per prendere appunti
        (``tavolette d'argilla e cannucce, se volete fare i babilonesi'',
        oppure carta e penna).
\end{itemize}
Chi ha un portatile poco potente dovrebbe comunque provare a installare
l'engine e aprire un template per vedere come gira: all'inizio non si faranno
cose pesanti. Nelle prime lezioni di level design è importante mettere subito
le mani sull'interfaccia: ascoltare dove si trova uno strumento (ad esempio lo
\emph{snapping}) senza provarlo porta a dimenticarlo.

% ---------------------------------------------------------------------
\section{Letture e risorse di riferimento}
\label{sec:risorse}

\subsection{I libri}

\begin{description}
  \item[\emph{The Art of Game Design: A Book of Lenses} (Jesse Schell)] il
        testo \textbf{davvero importante}. Alcune lezioni finali dell'anno si
        baseranno sulle sue terminologie. È utile perché apre la mente sulla
        possibilità di applicare metodologie diverse a giochi diversi, ma
        soprattutto per la \textbf{terminologia} comune. Esiste anche in
        audiolibro (circa 25 ore).
  \item[\emph{Challenges for Game Designers} (Brenda Romero e Ian Schreiber)]
        libro \emph{entry level} con moltissimi esercizi, molto caro al
        direttore del corso.
  \item[\emph{21st Century Game Design} (Chris Bateman)] utile come
        retrospettiva: è un libro visionario che ha previsto metodi e
        soprattutto \emph{archetipizzazioni} dei giocatori tuttora valide nel
        design e nel marketing. Servirà per una delle lezioni più importanti
        dell'anno. È utile leggerlo per vedere quanto aveva previsto, ma è
        utile ``fino a un certo punto''.
  \item[\emph{An Architectural Approach to Level Design} (Christopher
        Totten)] non strettamente necessario, ma di lettura semplice. Ha due
        anime: parte dall'architettura (babilonese, greca, ecc.) e la trasla
        su esempi di level design (come \emph{Doom} e \emph{Team Fortress}).
        Un tempo ne esisteva un riassunto su Gamasutra (oggi
        \emph{Game Developer}).
  \item[\emph{A Theory of Fun for Game Design} (Raph Koster)] utile fino a un
        certo punto: contiene diverse idee superate, che lo stesso Koster ha
        affrontato in vari talk reperibili nel GDC Vault. Verrà ripreso nelle
        prime lezioni, poi messo da parte.
\end{description}

\subsection{Canali e conferenze}

\begin{description}
  \item[GDC (Game Developers Conference)] il suo canale raccoglie talk e
        panel estremamente importanti: è ``uno strumento misterioso che ci
        servirà più tardi''. \textbf{Attenzione}: per molti talk mancano ancora
        le competenze per interpretarli, e alcuni potrebbero fuorviare o
        mettere in testa idee strane. La GDC si concentra spesso su
        produttività e mobile.
  \item[Il \emph{Dragon Speech} di Chris Crawford] il talk del fondatore della
        GDC: appartiene a un passato lontano, ma è molto motivante. Da vedere
        subito.
  \item[Game Maker's Toolkit] canale particolarmente utile per nozioni
        pratiche, con video brevi e alcune interviste.
  \item[Nordic Game] altra conferenza importante, con talk e panel di esperti
        del settore.
\end{description}

\subsection{Associazioni e premi: come leggerli}

Le associazioni di settore (ad esempio IIDEA in Italia, o la IGDA) non sono
il male, anzi: sono un canale importante per affacciarsi all'estero o per
accedere a fondi. Va però ricordato che \textbf{tutelano i propri associati}:
non esiste solo il loro palinsesto, e bisogna valutare quanto spazio occupano
nel panorama della visibilità.

Lo stesso vale per i \textbf{premi} (Game Awards, ``Best Mobile Game'',
``Games for Change'', ecc.): al designer \emph{non interessano}. Che un gioco
abbia vinto ``più di 90 premi'' non dice nulla di utile. Servono a dare
visibilità e lustro ai giochi ed è giusto che esistano, perché ampliano il
panorama e il numero di giocatori. Ma ciò che interessa al designer sono i
\textbf{dati}:
\begin{itemize}
  \item come sono state implementate certe cose;
  \item come è stato fatto il marketing e dove il gioco è stato presentato;
  \item non quante wishlist ha fatto, ma \emph{come} le ha fatte;
  \item in che periodo storico è uscito e quali erano i competitor in quel
        momento.
\end{itemize}

% ---------------------------------------------------------------------
\section{Che cosa fa un game designer}
\label{sec:cosa-fa}

Per arrivare a una prima definizione, la classe ha condotto un brainstorming:
ognuno doveva indicare una cosa che, secondo lui, fa il game designer (senza
ripetere quelle già dette). Ne è emersa la lista seguente, commentata punto
per punto.

\subsection{Progettare meccaniche e struttura del gioco}

Il ruolo più generico è quello del \textbf{core mechanic designer}, detto
anche \emph{gameplay designer} e, un tempo, \emph{black box designer}: si
occupa delle \textbf{meccaniche di gioco}. Si va dal piccolo al grande (in
funzione anche della dimensione dell'azienda):
\begin{itemize}
  \item nel \textbf{piccolo} ci sono le singole meccaniche;
  \item nel \textbf{grande} ci sono i \textbf{sistemi}. Nei giochi
        \emph{sistemici}, con molti sistemi interconnessi (ad esempio
        \emph{GTA}), esistono specializzazioni dedicate a bilanciare e
        supervisionare i sistemi: è il \emph{system design}.
\end{itemize}

\paragraph{Quando un gioco è ``sistemico''?}
Non dipende solo da come le meccaniche interagiscono con il giocatore (le
\emph{dinamiche}, secondo i framework che passano dalle meccaniche alle
dinamiche, che verranno studiati più avanti), ma anche da \emph{quante}
meccaniche ci sono e da quanto interagiscono tra loro. Non diremmo che
\emph{Super Mario} sia un gioco sistemico, eppure contiene dei sistemi: la
linea è sottile. Quando i sistemi aumentano molto, sono complessi e si
influenzano a vicenda, il gioco diventa sistemico. Esempio: in \emph{The
Legend of Zelda: Breath of the Wild} il vento modifica quanto il personaggio
sale o scende; il fuoco può incendiare l'erba; l'erba prende fuoco se colpita
da una freccia infuocata; e così via. Nel mercato moderno i giochi sistemici
sono spesso, ma non sempre, legati a mappe \emph{open world}.

\subsection{Le 3C: camera, character, control}

La base del gioco, da cui generalmente si inizia a lavorare (anche nei
prototipi), è la ``sacra trinità videoludica'' delle \textbf{3C}:
\begin{description}
  \item[Camera] che tipo di inquadratura ha il gioco: prima persona, terza
        persona, isometrica, ecc.
  \item[Character] quanto è grande il personaggio, che tipo di personaggio è
        e, di conseguenza, quali azioni può compiere.
  \item[Control] i controlli: che cosa il personaggio può fare, quali sono i
        \emph{verbi} di gioco.
\end{description}
Le 3C sono definite dal designer, ed esiste anche una specializzazione per
questo lavoro. Esistono varianti: quella del direttore del corso è
\textbf{4C + J}, cioè \emph{camera, character, control, conflict} (conflitto)
e \emph{juice} (il ``succo'', il feedback che rende piacevole l'interazione).
Verranno approfondite più avanti.

\subsection{Il personaggio}

Il designer definisce il personaggio dal punto di vista delle meccaniche
(dimensioni, azioni, capacità), all'interno delle 3C. L'aspetto estetico è
invece competenza di altri dipartimenti.

\subsection{Documentare le informazioni per gli altri dipartimenti}

Assolutamente vero: la documentazione è una delle prime attività del
designer. Più l'azienda è grande, più questo lavoro diventa prominente; più è
piccola, meno è rilevante. In \textbf{pre-produzione}, però, è sempre
fondamentale.

\subsection{Mediatore e leader ``al servizio''}

Il designer è un \textbf{mediatore} nei conflitti interni al team, una sorta
di collante che si relaziona con dipartimenti diversi (non sempre gli stessi,
a seconda della fase di produzione). A volte fa anche da leader, a seconda di
come sono gestite la pipeline e la gerarchia.

Questa leadership è difficile da definire perché \textbf{non è verticale}: è
una leadership \emph{al servizio} (\emph{servant leadership}). Il designer
``è asservito al programmatore'': il suo compito non è arrivare con un'idea e
scaricarla addosso agli altri, né portare la soluzione e lasciare il problema
a qualcun altro, ma \textbf{seguire il processo a basso livello}, sporcandosi
le mani.

\paragraph{I feature team.} Nei team medio-piccoli o medio-grandi possono
esistere sotto-team (\emph{feature team}) dedicati a un obiettivo specifico:
ad esempio un designer incaricato del character o della camera, affiancato da
uno o due programmatori e da un artista. In questi casi il designer
supervisiona il lavoro e gestisce la produzione del sotto-team.

\subsection{Gestire l'avanzamento del prodotto}

Esiste una specializzazione che si occupa solo di questo, il
\textbf{producer}, a un livello gerarchico completamente diverso (e che
potrebbe essere una strada percorribile al termine del percorso di studi).
All'interno di un feature team anche il designer può farlo, ma non è detto.
Può certamente \emph{dare la direzione}: nella sua documentazione può
indicare che per arrivare a uno stato giocabile del prototipo serve un certo
lavoro di arte e di programmazione, per lo stato successivo un altro, e così
via. In ambito produttivo esistono metodologie strutturate come
\textbf{Agile}, che verranno viste in seguito.

\subsection{Presentare le idee}

Il pitching (sezione~\ref{sec:gerarchia-corso}) è richiesto al designer sia
in pre-produzione sia in produzione, in misura che dipende dall'azienda. Ad
esempio Ubisoft lascia ai designer, e in realtà a qualsiasi membro, la
possibilità di \emph{pitchare} un'idea prima di una nuova produzione, perché
possa iniziare a prendere vita.

\subsection{Bozze e prototipi}

Le bozze vanno bene, ma devono essere \textbf{chiare}. C'è un'eterna lotta tra
documentazione e gioco giocabile: se si fa solo documentazione il gioco non
avanza; se si fa solo il gioco manca la documentazione e si perde la visione.
Sono due cose \textbf{complementari}.

Il modo per comunicare un'idea in forma grezza ma giocabile è il
\textbf{prototipo}. Un prototipo può essere qualsiasi cosa: di carta, in
Unreal con Blueprint ``brutti''. \emph{Non deve essere bello, deve essere
giocabile}.

\paragraph{Esempio: il prototipo di carta di Mario + Rabbids.} Durante lo
sviluppo di \emph{Mario + Rabbids}, i programmatori di Ubisoft erano impegnati
su altri progetti (come \emph{Ghost Recon}). I designer realizzarono allora
un \textbf{gioco da tavolo} che riproduceva le meccaniche del gioco, quelle
di uno strategico a turni sul modello di \emph{XCOM}. E funzionava.
L'obiettivo era \textbf{validare un'idea}, validare un gameplay.

Altri modi per rendere chiara un'idea sono i \emph{mock-up} nella
documentazione e le immagini modificate o i \emph{mash-up} di qualsiasi
natura. Lato documentazione l'importante è che sia chiaro; lato prototipo,
che sia giocabile.

\subsection{Fare ricerca per l'idea, e poi smontarla}

Il designer ha una componente creativa importante e fa ricerca per arrivare a
un'idea. Ma ciò che deve poi fare necessariamente è \textbf{smontarla}. Il
suo lavoro è scientifico, esattamente come il \textbf{metodo scientifico}, il
cui obiettivo è \emph{confutare un'ipotesi}: si arriva a un'ipotesi tramite
sperimentazione e poi si cerca $n$ volte di confutarla; se resiste, la si
assume come valida \emph{in un certo contesto} (le regole della fisica
classica, ad esempio, non valgono nella meccanica quantistica).

Allo stesso modo il designer arriva a un'idea (tramite ricerca o chiedendo
agli altri: le idee le hanno tutti, non sono un privilegio del designer),
cerca di smontarla e poi cerca di \textbf{metterla a terra}. Se l'idea non è
applicabile, non è implementabile, finché non si trova un modo per
realizzarla va scartata. Può essere l'idea più bella dell'universo, ma
\textbf{se non può essere messa a terra, non esiste}.

\paragraph{L'originalità.} Serve un'idea originale? Dipende dall'obiettivo e
dalla visione del progetto:
\begin{itemize}
  \item se l'obiettivo è vendere molto, un'idea originale non serve
        necessariamente;
  \item se si vuole creare un nuovo filone, serve un'idea, se non originale,
        almeno interessante: può essere un discorso di marketing o una
        meccanica innovativa, con il rischio di non vendere.
\end{itemize}
Il ruolo del \emph{director} e il concetto di \emph{visione} verranno
approfonditi in seguito.

\subsection{Analizzare i competitor}

Arrivati a un'idea (``vogliamo un certo tipo di combattimento''), si va a
vedere se qualcuno l'ha già fatta e \emph{come}. Le meccaniche non sono
privilegio di nessuno, salvo quando vengono brevettate, come il
\textbf{Nemesis System} di \emph{La Terra di Mezzo: L'Ombra di Mordor}.

\paragraph{Esempio: il combattimento da ridimensionare.} In un progetto
(presso Stormind Games) il team si trovò con un sistema di combattimento
troppo costoso: mancavano le risorse e gli animatori per realizzarlo. Poco
output lato gameplay significava non poter fare qualcosa di complesso e
profondo come in \emph{Elden Ring}; al contrario, il team aveva molti
artisti e poteva puntare su qualcosa di più artistico. Si iniziò così a
orbitare intorno al corpo a corpo di \emph{The Last of Us} (parte I e II).
Sul canale della GDC esistono diversi talk che raccontano cosa è stato fatto
e come, con successi, insuccessi e costi (ad esempio del \emph{motion
capture}): da quei talk si ricavarono moltissime informazioni utili.

I competitor possono essere:
\begin{itemize}
  \item \textbf{diretti}: un gioco molto simile al proprio;
  \item \textbf{di mercato}: un gioco che domina il mercato in cui ci si vuole
        affacciare;
  \item \textbf{indiretti}: un gioco con un sistema simile al proprio, ma
        tutto il resto diverso.
\end{itemize}
Regola d'oro: \textbf{là fuori, nascosto nell'ombra, c'è sicuramente un gioco
molto simile al vostro}. Va solo trovato. E se per qualche ragione non ha
venduto, bisogna capire perché.

\subsection{Lavorare in engine e fare tuning}

Il level designer lavora in engine, così come il technical designer; ma
qualsiasi game designer, a seconda di ciò che deve fare, può finire in
engine, soprattutto per il \textbf{bilanciamento} e il \textbf{tuning}.

Il \emph{tuning} è l'azione di bilanciare con precisione i dati. Esempio: in
\emph{Super Mario} si modifica quanto in alto salta il personaggio, con che
velocità, in quanto tempo raggiunge l'apice, e così via. Questo lavoro rende
il gioco:
\begin{itemize}
  \item ad alto livello, dotato della \textbf{progressione} giusta;
  \item a basso livello, \textbf{piacevole da giocare}.
\end{itemize}
Il \emph{fine tuning} è un lavoro di dettaglio, paziente e minuzioso.

\subsection{Bilanciare}

Il bilanciamento, emerso tardi nel brainstorming, è una delle attività
centrali del designer e ha anche una specializzazione dedicata
(l'\emph{economy designer}, sezione~\ref{sec:tipi-designer}).

\subsection{Analizzare il mercato}

In pre-produzione, di solito, sì: capire le tendenze, i generi che cambiano
nel tempo, la compatibilità del gioco con il momento storico, in vista ad
esempio di un possibile publisher. Va capito quanto, come e con quanto tempo
farlo, perché si è sempre limitati da molte cose.

\subsection{Sintesi}

\begin{quote}
Il designer parte dall'idea ma non si ferma lì: la confuta, lavora sui
dettagli, bilancia, fa \emph{fine tuning}; lavora con il team e fa da
\emph{manager} per diversi dipartimenti; e naviga costantemente tra i limiti.
\end{quote}

% ---------------------------------------------------------------------
\section{I limiti del design}
\label{sec:limiti}

I limiti ci sono sempre, e nel mondo del design sono spesso numerosi. I più
importanti entro cui navigare sono:
\begin{description}
  \item[Tempo] ad alto livello, l'idea deve essere realizzabile nel tempo di
        produzione disponibile; a basso livello, il proprio lavoro
        settimanale deve essere fattibile nel tempo che si ha. Non ci si può
        dividere tra mille feature team se realisticamente se ne può seguire
        uno solo.
  \item[Tecnologia] si può o non si può fare qualcosa? Ci sono i server per un
        multiplayer con un certo numero di giocatori? Ci si può permettere,
        ad esempio, una fauna che reagisce alla fisica come quella di
        \emph{GTA VI}, che richiede il meglio dei \emph{technical artist}?
  \item[Background e competenze] che cosa sa fare il team? Realisticamente,
        chi ha lavorato dieci anni su giochi horror viene assunto per un
        horror, perché ha le competenze per risolvere problemi già visti.
  \item[Persone e motivazione] il team ha la motivazione e le competenze per
        fare un certo tipo di gioco, e per lavorare insieme?
  \item[Denaro] tutto ruota intorno ai soldi. Che si lavori da soli o in mille,
        le persone vanno pagate e il gioco deve avere un ritorno economico.
        Bisogna decidere dove tracciare la linea: raggiungere il
        \textbf{break-even} (guadagnare esattamente quanto speso) o puntare a
        grandi guadagni. Con un \textbf{publisher}, inoltre, una parte
        consistente dei ricavi va a lui (ripartizioni dell'ordine di 80/20 o
        70/30, che dipendono dal contratto firmato).
\end{description}

\subsection{Non siete ``gli uomini delle idee''}

Per qualcuno può essere un problema scoprirlo, ma il designer non è
``l'uomo (o la donna) delle idee'': accanto alla parte creativa c'è una parte
tecnica, e si è fortemente limitati da molti fattori. Questi limiti, però,
sono un aiuto:
\begin{itemize}
  \item riducono il ventaglio di possibilità potenzialmente distruttive per il
        team (cose che non si riescono a fare, che non piacciono o per cui
        mancano le risorse);
  \item aiutano a mirare bene, trovando soluzioni creative in uno spazio più
        piccolo e più attuabile.
\end{itemize}

\begin{quote}
\centering
\textbf{Le idee sono nulla, l'implementazione è tutto.}
\end{quote}

% ---------------------------------------------------------------------
\section{Ruoli e specializzazioni}
\label{sec:tipi-designer}

\subsection{La scala gerarchica}

\begin{description}
  \item[Creative Director (CD)] mantiene la \textbf{visione} e dà l'idea, in
        modo che tutto segua una certa direzione. In teoria può diventarlo
        chiunque, ma con un'esperienza notevole, nell'ordine dei decenni: si
        assume che in quel tempo abbia visto come lavorano molti dipartimenti.
        Naturalmente un CD con un background artistico ha meno competenze di
        design, e viceversa.
  \item[Director di disciplina] uno per dipartimento: \emph{game director},
        \emph{level director}, \emph{art director}, ecc.
  \item[Lead] il capo del dipartimento, che gestisce i gruppi di persone. Chi
        lavora in un dipartimento risponde al proprio lead (ad esempio un game
        designer risponde al \emph{lead designer}, non al game director).
  \item[Senior / Principal] i profili più esperti (a seconda dell'azienda).
  \item[Junior] il profilo di inizio carriera.
  \item[Intern] lo stagista.
\end{description}
Lo schema si replica uno a uno per ogni dipartimento: arte, suono, animazione
(a volte arte e animazione insieme, o \emph{rigging} e animazione insieme),
programmazione, UX/UI e così via. A fare la differenza, a parità di ruolo, è
il \textbf{livello di seniority}.

\subsection{UX e UI}

UX e UI stanno a metà tra game design e arte: a seconda dell'azienda il ruolo
si chiama \emph{UX/UI Artist} o \emph{UX/UI Designer}.
\begin{description}
  \item[UX (User Experience)] come una persona interagisce con le cose. Esempio:
        su uno smartphone, un pulsante accende e spegne lo schermo; per
        digitare la password si scorre verso l'alto e compaiono i punti da
        toccare.
  \item[UI (User Interface)] come sono ``vestite'' queste cose: ad esempio
        l'aspetto della barra dei punti vita (HP), realizzata con asset 2D o
        3D.
\end{description}

\subsection{Le specializzazioni del game design}

Più l'azienda è grande, più la specializzazione è puntuale.
\begin{description}
  \item[Gameplay / Core mechanics designer] (sono la stessa cosa) si occupa
        delle meccaniche e delle 3C.
  \item[System designer] si occupa dei sistemi e delle loro interazioni.
  \item[Economy designer] si occupa del bilanciamento; lavora soprattutto con
        fogli di calcolo (Excel, Google Sheets).
  \item[Combat designer] si occupa del combattimento; negli annunci di lavoro
        indica quasi sempre un combattimento \emph{action}, con gameplay
        reattivo (non strategico), anche se non è una regola.
  \item[Multiplayer designer] si occupa delle interazioni multigiocatore.
  \item[Technical designer] fa da ponte tra game design e programmazione.
\end{description}

\subsection{Le specializzazioni del level design}

\begin{description}
  \item[Encounter designer] posiziona i nemici nella mappa in modo che gli
        scontri siano bilanciati e la progressione equilibrata come
        desiderato.
  \item[World / map designer] popola il mondo di gioco.
  \item[Quest designer] popola il mondo con le missioni e pensa alla
        progressione al loro interno, cioè a dove il giocatore deve andare (si
        vedano ad esempio gli annunci di lavoro di CD Projekt Red o di
        Larian).
  \item[Technical level designer] fa da ponte tra level design e
        programmazione.
\end{description}

\subsection{Discipline vicine: suono, UI, produzione}

Il monte ore del corso non permette di uscire come \emph{narrative designer}
(o \emph{sound designer}) completi. Ciò non toglie che si possano coltivare
queste discipline:
\begin{itemize}
  \item nei progetti di gruppo si può occuparsi anche del suono, della UI/UX,
        o provare ad applicare metodologie come Agile per fare
        \emph{producing} e gestire il team;
  \item il tempo dedicato a queste attività va però sottratto al game design,
        e dipende da quante persone ci sono nel team, da quanto lavoro c'è e
        da quanto è grande il progetto;
  \item avere sviluppato e rilasciato un progetto, anche piccolo (anche una
        \emph{game jam}), documentando i propri processi, è già un ottimo
        argomento da presentare in futuro;
  \item la via migliore resta lavorare nell'industria e, da lì, chiedere di
        imparare anche altre discipline.
\end{itemize}
Un gioco si può fare su qualsiasi cosa: sugli scarafaggi, sui piedini delle
sedie, sui pennarelli. Si parte avvantaggiati da ciò che è vicino alla propria
sensibilità: chi ama la musica può pensare a un gioco con meccaniche legate
alla musica, con un gameplay innovativo ma contenuto e realizzabile. Il design
è un cappello che comprende molte discipline.

% ---------------------------------------------------------------------
\section{Approfondimenti emersi in lezione}
\label{sec:approfondimenti}

\subsection{Roguelike e roguelite}

Un esempio di come il marketing conia generi: Supergiant Games, con
\emph{Hades}, rese celebre il termine \textbf{roguelite} al posto di
\emph{roguelike} (pur non essendo il primo roguelite in assoluto). La
differenza:
\begin{itemize}
  \item il \textbf{roguelike} si fonda su una serie di punti fissati in una
        sorta di \emph{manifesto} che definisce che cosa rende un gioco
        roguelike; uno su tutti: a ogni morte si ricomincia la \emph{run} da
        capo;
  \item il \textbf{roguelite} rompe alcuni di questi punti, in particolare
        introducendo una \textbf{progressione esterna alla run}: il giocatore
        può potenziarsi fuori dalla singola partita per vincere più
        facilmente in quelle successive. \emph{Hades} ne è il caso principe.
\end{itemize}
Analogamente, un mobile game sviluppato dal direttore del corso fu
promosso come ``il primo \emph{donation game} al mondo'': un gioco che
convertiva il tempo di gioco in donazioni reali a enti benefici per animali.
Anche qui un neologismo creato lato marketing per definire un genere.

\subsection{Applicazioni fuori dall'industria videoludica}

Le competenze sviluppate non servono solo ai videogiochi:
\begin{itemize}
  \item i \textbf{render architettonici} oggi possono essere realizzati con
        Unreal, e diversi professionisti formati nel game design lavorano in
        questo campo;
  \item la \textbf{gamification}, cioè l'applicazione di regole videoludiche
        a contesti non ludici (mostre interattive, progetti culturali,
        consulenze come quella per applicare questi concetti alle macchine da
        allenamento di un'azienda di fitness), è in Italia uno dei modi più
        redditizi di spendere queste competenze, spesso con finanziamenti
        pubblici.
\end{itemize}

% ---------------------------------------------------------------------
\section{Il primo task}
\label{sec:primo-task}

\begin{enumerate}
  \item Creare un indirizzo e-mail \textbf{professionale} (consigliato: un indirizzo Gmail nella forma
        \texttt{nome.cognome@gmail.com}; l'anno di nascita è accettabile).
        Va bene anche un indirizzo professionale già esistente usato per
        lavoro. Da evitare indirizzi poco seri (``non \texttt{blackfire67}'').
  \item Una persona della classe raccoglie tutti gli indirizzi in un unico
        documento e lo invia via e-mail al docente, che aggiungerà tutti al
        Drive condiviso del corso (dove verranno caricate lezioni, slide e
        materiali, per avere i documenti tutti nello stesso posto).
  \item Installare Discord e iscriversi al server del corso, dove passano le
        comunicazioni ufficiali. È consuetudine, e utile, creare anche un
        gruppo tra studenti.
\end{enumerate}
\textbf{Scadenza}: prima della lezione del venerdì.

% ---------------------------------------------------------------------
\section*{Punti chiave del capitolo}
\addcontentsline{toc}{section}{Punti chiave del capitolo}

\begin{itemize}
  \item Il corso è una \textbf{simulazione di lavoro}: ci si tratta da
        colleghi, con una gerarchia da studio di sviluppo.
  \item Si impara con l'\textbf{inquiry-based learning}: lezione, task,
        dubbi, domande, condivisione, progetti.
  \item Il fallimento affrontato con la mentalità giusta è la principale fonte
        di crescita; i feedback riguardano il lavoro, mai la persona.
  \item Il lavoro del designer è per metà creativo e per metà
        \textbf{ricerca e metodo scientifico}; il game design non è una
        scienza esatta.
  \item \textbf{Documentazione} e \textbf{pitching} sono il cuore del
        mestiere.
  \item Il designer progetta meccaniche e sistemi (partendo dalle
        \textbf{3C}), documenta, media, prototipa, confuta le idee, studia i
        competitor, bilancia e fa tuning.
  \item La leadership del designer è \textbf{al servizio} del team, non
        verticale.
  \item I limiti (tempo, tecnologia, competenze, persone, denaro) restringono
        e indirizzano la creatività.
  \item \textbf{Le idee sono nulla, l'implementazione è tutto.}
\end{itemize}

%    % ...
%
%  (Con hyperref aggiungere anche \renewcommand{\theHchapter}{z.\arabic{chapter}}
%  nella parte zero, per evitare ancore duplicate.)
%  Tutte le etichette di questo capitolo hanno il prefisso "z3:"; i
%  rimandi ai capitoli precedenti usano le etichette "z1:" e "z2:" già
%  definite in GD0_1 e GD0_2, e alcune etichette dei capitoli del corso
%  (GD1, GD3).
% =====================================================================

\chapter[Agile e Scrum nella pratica]{Agile e Scrum nella pratica: l'esperienza di chi ci è passato}
\label{cap:z3-agile-scrum}

Le prime due lezioni zero (capitoli~\ref{cap:z1-feedback}
e~\ref{cap:z2-tempo-idee-leadership}) guardavano al lavoro in team dal punto
di vista di un professionista con quindici anni di industria alle spalle.
Questa terza lezione cambia prospettiva: a parlare non è un veterano, ma
\textbf{due persone che hanno fatto lo stesso percorso di studi} di chi
ascolta e che, meglio di chiunque altro fino a quel momento, sono riuscite a
\emph{internalizzare} Agile e Scrum nei progetti dell'accademia.

La lezione è introdotta e chiusa da Diego Ricchiuti\footnote{Nella
trascrizione compare soltanto il nome di battesimo, ``Diego'', con cui i due
ospiti si rivolgono più volte a lui; il cognome è quello indicato per il
direttore dell'accademia nel capitolo~\ref{cap:z1-feedback}. Dal contesto è
anche il docente che ha già spiegato Scrum e Agile alla classe nelle lezioni
precedenti (``Diego vi avrà già ampiamente indirizzato'', ``come avete fatto
con Diego nelle lezioni scorse'').}, che ha già spiegato alla classe la teoria
di Agile e Scrum e che interviene di tanto in tanto per sottolineare o
completare un passaggio. Gli ospiti sono due:
\begin{itemize}
  \item \textbf{Simone Cicchetti}, ex studente dell'accademia, oggi level
        designer e producer, che parla per circa un'ora e mezza della teoria
        ``vista dalla vita reale'' e degli errori più comuni;
  \item \textbf{Cristian Liberatore}\footnote{Nella trascrizione il nome
        compare sia come ``Cristian'' sia come ``Christian''; qui si usa la
        prima forma, quella con cui viene presentato all'inizio della
        lezione. Anche i nomi dell'accademia sono resi in modi diversi
        (``Hive'', ``a dive'', ``IVE Online'', ``Live Online''): come negli
        altri capitoli si usa semplicemente ``l'accademia''; ``IVE Online''
        sembra indicare il corso a distanza dell'accademia.}, studente del
        terzo anno, game designer e project manager di due progetti nati
        come progetti d'anno, che segue passo per passo il ciclo di Scrum
        con gli esempi del suo gioco.
\end{itemize}
Diego li ringrazia per il tempo dedicato: Simone lavora, quindi sta
sacrificando tempo personale, e Cristian è uno studente con lo stesso carico
di lavoro della classe. Sono stati scelti perché, tra gli studenti avuti
fino a quel momento, sono quelli che hanno davvero \textbf{internalizzato}
Scrum e Agile, ottenendo grandi risultati produttivi nel progetto scolastico
e, nel caso di Simone, anche dopo l'accademia\footnote{La trascrizione dice
``nel caso di Christian post-accademia'', ma è Simone l'ex studente che oggi
lavora, mentre Cristian è ancora iscritto: il riferimento al dopo-accademia
è stato quindi attribuito a Simone.}.

Perché insistere sulla parola \emph{internalizzare}? Perché, dice Diego,
``utilizzare Scrum e Agile sembra molto facile, ma poi in realtà quando li
provi ad applicare veramente ti rendi conto che facili come sembravano quando
facciamo lezione non lo sono''. Se invece li si internalizza davvero, si
ottengono grandi risultati. L'idea della lezione è che possa essere utile
vedere l'esperienza di qualcuno \emph{molto più vicino} agli studenti di
quanto lo sia il docente, che ormai usa questi strumenti da anni. Il formato
annunciato è di circa un'ora e mezza per ciascun ospite, domande comprese.

Due avvertenze di lettura. La prima: i due interventi sono fatti di teoria
richiamata in breve e di moltissimi esempi concreti, spesso con numeri,
errori e ripensamenti; sono riportati con cura perché sono proprio la parte
che la teoria non dà. La seconda: i due relatori usano in modo
intercambiabile alcuni termini (\emph{framework} e metodologia,
\emph{post-mortem} e revisione della sprint, \emph{score}, peso e punteggio
dei task), e il capitolo segue il loro uso, segnalando in nota dove si
discosta dalla terminologia ufficiale di Scrum.

Il percorso del capitolo segue quello della lezione:
\begin{enumerate}
  \item un'apertura organizzativa e la presentazione degli ospiti
        (sezione~\ref{sec:z3-apertura});
  \item \textbf{Simone}: chi è, un sondaggio sulla classe e le prime
        esperienze di Scrum degli studenti (sezione~\ref{sec:z3-simone});
  \item il \textbf{mindset Agile}: perché non è un framework e perché va
        adattato (sezione~\ref{sec:z3-mindset});
  \item le \textbf{sprint} e le riunioni di pianificazione e
        \emph{post-mortem} (sezione~\ref{sec:z3-sprint});
  \item le \textbf{cerimonie} e la regola della \emph{consistency}, con le
        critiche al \emph{daily scrum} e al \emph{planning poker}
        (sezione~\ref{sec:z3-cerimonie});
  \item il \textbf{backlog} e gli obiettivi dei task
        (sezione~\ref{sec:z3-backlog-simone});
  \item la \textbf{velocity}: perché non usare le ore, la sequenza di
        Fibonacci e come la velocity permette di ``prevedere il futuro''
        (sezione~\ref{sec:z3-velocity});
  \item il riepilogo della creazione di una sprint e la conclusione sulla
        \textbf{fiducia} (sezione~\ref{sec:z3-ricapitolo});
  \item le domande a Simone, con la costruzione in diretta di un foglio per
        il calcolo della velocity (sezione~\ref{sec:z3-qa-simone});
  \item \textbf{Cristian}: chi è e come imposta la sua parte
        (sezione~\ref{sec:z3-cristian});
  \item il caso di studio, il gioco \emph{L'Inf}, e il ``trauma'' del cambio
        di team (sezione~\ref{sec:z3-linf});
  \item il \textbf{loop di Scrum} in cinque fasi
        (sezione~\ref{sec:z3-loop}): la stesura del backlog
        (sezione~\ref{sec:z3-backlog}), la selezione dei task
        (sezione~\ref{sec:z3-selezione}), l'implementazione con i suoi
        controlli, i \emph{blocker} e i feedback
        (sezione~\ref{sec:z3-implementazione}), il calcolo della velocity
        (sezione~\ref{sec:z3-calcolo}) e il \emph{tuning} della produzione,
        con il problema del perfezionismo (sezione~\ref{sec:z3-tuning});
  \item una domanda su come comunicare un taglio al team
        (sezione~\ref{sec:z3-qa-taglio});
  \item il \textbf{team management}: il team prima di Scrum, la produzione
        senza stress, il mindset condiviso, il \emph{post-mortem} e la
        \textbf{sincerità} (sezione~\ref{sec:z3-team});
  \item la chiusura della lezione (sezione~\ref{sec:z3-chiusura}).
\end{enumerate}

% ---------------------------------------------------------------------
\section{Apertura e presentazione degli ospiti}
\label{sec:z3-apertura}

\subsection{Una registrazione andata persa}

La lezione si apre con una comunicazione organizzativa\footnote{La
trascrizione comincia a frase già iniziata (``in una modalità non video fin
tanto che ce la chiamate poi la converte alla fine''): l'inizio della
spiegazione non è stato registrato e non è chiaro a quale lezione precedente
si riferisca.}. Il programma usato per registrare una lezione precedente
registra in una modalità che va convertita alla fine; durante la conversione
il computer di Diego, probabilmente per colpa degli aggiornamenti di Windows,
si è bloccato e la conversione è ``andata a pallino''. Poiché si tratta di
una lezione importante, Diego si rende disponibile a \textbf{rifarla
interamente} per registrarla di nuovo oppure, come la classe preferisce, a
fornire un audio, un video o qualunque altro formato sia più comodo, che poi
caricherà. Chiede agli studenti di pensarci e di fargli sapere.

\subsection{Chi sono gli ospiti e perché sono qui}

Poi Diego presenta gli ospiti speciali della giornata: Cristian Liberatore,
studente del terzo anno, e Simone Cicchetti, ex studente dell'accademia (che
in un primo momento ``sparisce'' dalla griglia della videochiamata, finché
Diego non lo ritrova in alto a destra). Le ragioni della scelta, il valore
della parola \emph{internalizzare} e il formato (un'ora e mezza ciascuno, con
domande) sono quelli già riassunti all'inizio del capitolo. Detto questo,
Diego si silenzia e lascia la parola a Simone.

% ---------------------------------------------------------------------
\section{Simone: il punto di vista del producer}
\label{sec:z3-simone}

\subsection{Chi è il relatore}

Simone si presenta mentre condivide lo schermo (``adesso tutto è imploso,
perfetto''). È uno degli ex studenti del terzo anno e oggi lavora come
\textbf{level designer}; collabora con Diego su alcuni progetti legati alla
gestione dei team, al \emph{game design}, al \emph{framework design} e ad
altre ``cose molto fighe'' che alcuni studenti forse conoscono\footnote{Il
passaggio è confuso nella trascrizione (``lavoro con diego alcune alcuni
progetti diciamo sul gestione di pro dei team game design framework
design''): la ricostruzione è approssimativa. Più avanti Simone dice anche di
lavorare per una società australiana.}. Sta inoltre continuando a lavorare al
progetto del suo terzo anno (il \emph{Biggest One}\footnote{Dalla lezione
si capisce che i progetti d'anno dell'accademia si chiamano \emph{Big One}
(primo anno), \emph{Bigger One} (secondo anno) e \emph{Biggest One} (terzo
anno); la trascrizione li scrive in vari modi (``biggest one'', ``bigger
one'', ``BiggerOne'').}), un gioco chiamato \emph{ISO}\footnote{Il nome del
gioco compare sempre come ``ISO''; non è possibile stabilire dalla sola
trascrizione se sia la grafia corretta (potrebbe essere, ad esempio,
\emph{Iso} o una sigla).}, per il quale è level designer, game designer,
social media manager e soprattutto \textbf{producer}: si occupa di gestire il
lavoro delle persone, di far rispettare le cerimonie di Scrum e, ``punto
centrale della questione'', di ``rompere le scatole alla gente''.

\subsection{Una lezione partecipativa}

Il tempo a disposizione è molto, quindi Simone vuole una lezione il più
possibile \textbf{partecipativa e cooperativa}: a differenza del solito, le
domande vanno fatte in qualsiasi momento, scrivendo in chat per attirare la
sua attenzione o alzando la ``manina'' con le emoji, che controlla su un
secondo monitor con la griglia dei partecipanti. Al massimo aspetterà la fine
della slide per non perdere il filo.

L'obiettivo dichiarato non è spiegare la teoria, che la classe ha già fatto
con Diego, ma darne un breve \emph{recap} e poi spiegare \textbf{a che cosa
serve nel pratico e com'è nella vita reale}: ``come spesso accade, la teoria
è una cosa, la vita reale è un'altra''.

\subsection{Un piccolo sondaggio}

Prima di partire, Simone chiede chi abbia ``un amore incommensurabile e
infinito'' per il producing, lo \emph{scrum mastering} e il management in
generale: nessuno, e uno studente, Gavi\footnote{Così nella trascrizione;
potrebbe trattarsi di un soprannome o di un nome storpiato.}, risponde
spammando il ``no''. Poi la domanda opposta: chi lo odia così tanto da
preferire ``affrontare un alligatore in una lotta greco-romana''? Ancora
nessuno. La classe è quindi, in generale, ``tutta intermedia''.

Chi si è già dovuto occupare di queste cose nei progetti d'anno o nei lavori
di gruppo? Rispondono due studenti, Mattia e Lorenzo\footnote{Nella
trascrizione ``da noi Mattia, Podini'' e ``da noi Lorenzoaldasso'': i cognomi
non sono ricostruibili con sicurezza (``Podini'' potrebbe essere il cognome
di Mattia o il nome di un altro studente).}, ognuno per il proprio team.
Simone chiede solo i punti salienti: che cosa è andato bene, che cosa no e
perché.

\paragraph{L'esperienza di Mattia.} Il suo gruppo usa Scrum da pochissimo,
circa una settimana. È lui a gestire la \emph{Scrum board}, con l'aiuto degli
altri due membri, ma il metodo non è ancora applicato a pieno. Mattia è
interessato a Scrum, ma la \textbf{responsabilità di guidare altre persone}
oltre a se stesso gli pesa parecchio, e usare Scrum in modo serrato tutti i
giorni gli fa ``una bella paura'': se qualche meccanismo si rompe,
bisognerebbe correre ai ripari. Che cosa intende per ``rompere un
meccanismo''? Che se una cerimonia non viene rispettata, tornare allo
\emph{status quo} non sarebbe semplice, anche per la sua mancanza di
esperienza nell'organizzare i team. Simone riassume: la preoccupazione è più
sul piano umano che su quello delle ore di lavoro necessarie. Mattia
riconosce che Scrum avrà le sue ``quasi magiche proprietà'', ma non avendolo
ancora visto all'opera nella vita reale non può dirlo con certezza.

\paragraph{L'esperienza di Lorenzo.} Niente di diverso: anche il suo gruppo
ha iniziato da pochissimo e proprio il giorno prima ha fatto la prima
riunione di Scrum, decidendo insieme i pesi dei task. Alla domanda se il
ruolo sia arrivato per voto popolare, Lorenzo risponde che si è offerto lui:
conoscere un \emph{framework} di lavoro, e soprattutto capire come
applicarlo, gli interessa parecchio.

Simone trova molto interessante quanto detto da Mattia sulle cerimonie e sul
doverle rispettare, e promette di tornarci più avanti
(sezione~\ref{sec:z3-cerimonie}).

% ---------------------------------------------------------------------
\section{Il mindset Agile}
\label{sec:z3-mindset}

\subsection{Agile non è un framework}

La prima slide parte da una frase che tutti conoscono: \textbf{Agile non è un
framework, è un mindset} (Simone dice di aver ``percepito la furia negli
occhi di Diego'' quando la classe esitava a rispondere). La frase sembra una
di quelle che servono ``per fare i fighi con i colleghi'', ma ha un
significato preciso: un \emph{framework} presuppone \textbf{regole da
seguire}, mentre Agile \textbf{non ha regole}. È una serie di \textbf{linee
guida da internalizzare}: la parola chiave è che devono diventare
\emph{interne}, al punto da non dover più pensare attivamente a quali siano
le regole, i punti o la mentalità di Agile. Bisogna capire la logica che c'è
sotto.

L'analogia è quella della \textbf{lingua straniera}: all'inizio si ragiona in
italiano, si costruisce la frase in italiano, la si traduce e la si dice; man
mano questi passaggi si annullano e si comincia a ragionare direttamente
nella nuova lingua. Con Agile funziona allo stesso modo.

\subsection{I quattro valori, e perché non vanno presi alla lettera}

Simone rilegge rapidamente i quattro valori del manifesto Agile:
\begin{enumerate}
  \item gli \textbf{individui e le interazioni} sono più importanti dei
        processi e degli strumenti;
  \item un \textbf{prodotto funzionante} è più importante di una
        documentazione esaustiva\footnote{Nella lezione ``documentazione
        comprensibile''; nella formulazione corrente del manifesto Agile il
        termine è \emph{comprehensive documentation}, cioè documentazione
        esaustiva.};
  \item la \textbf{collaborazione con il cliente} è più importante della
        negoziazione del contratto;
  \item \textbf{rispondere al cambiamento} è più importante di seguire un
        piano.
\end{enumerate}
Come è scritto sotto al manifesto stesso, questo non significa che le cose a
destra, quelle ``meno importanti'', non contino o non vadano seguite. Ma
Simone va oltre: \textbf{non significa nemmeno che queste regole debbano
essere seguite davvero}. Non è detto che individui e interazioni debbano
essere per forza più importanti di processi e strumenti. Il punto chiave di
Agile è \emph{sentire} quale sia il procedimento e la mentalità giusti e
\textbf{adattare Agile al proprio team}, perché ogni team e ogni progetto è
diverso.

Se un team e un progetto, per vari motivi, hanno bisogno di pagine e pagine
di documentazione, allora per loro la documentazione è più importante del
prodotto funzionante. Gli esempi sono infiniti; uno è il \textbf{gioco di
carte}: lì il prodotto funzionante è relativamente meno importante della
documentazione, perché le logiche si possono scrivere in qualsiasi momento,
ma senza una documentazione solida a supporto ``cadrà tutto quanto
immediatamente''. Quindi non bisogna seguire ``a menadito'' i punti del
manifesto, ma capirne la logica di fondo, capire \textbf{a che cosa serve} e
capire \textbf{se serve}: non è detto che Agile serva.

\subsection{Altri mindset e i framework}

Esistono molti altri \emph{mindset}. Di recente va molto di moda la
mentalità \textbf{lean}, che Simone definisce ``già il 2.0''\footnote{Così
nella lezione. Storicamente il \emph{lean} nasce nella produzione
industriale (il sistema Toyota) e precede Agile; nello sviluppo software il
\emph{lean software development} si è diffuso negli anni Duemila. Il ``2.0''
va letto come una battuta sul fatto che oggi è di moda, non come una
derivazione da Agile.}. L'importante, aggiunge con Diego, è \textbf{non usare
Waterfall}, ``perché Waterfall nel gaming\dots{} non dirò nient'altro''.

Sulla base di Agile sono state create delle \textbf{metodologie}, che
Simone si corregge chiamando \emph{framework}: e queste sì che hanno delle
regole. L'esempio è Scrum, ma ce ne sono infiniti, come Kanban. Si basano
tutti sul principio Agile dell'\textbf{inspect and adapt} (ispeziona e
adattati, analizza e adattati), che si può riassumere in una parola chiave:
\textbf{iterazioni}. Scrum è completamente basato sulle iterazioni, e anche lo
sviluppo dei videogiochi lo è, ``molto più dello sviluppo software''.

Il resto della sua parte è dedicato ai punti di Scrum, il \emph{framework} da
cui presumibilmente partono tutti gli studenti, e agli \textbf{errori
classici} in cui si incappa spesso. L'ideale sarebbe trasmettere la logica
alla base di tutto, ma Simone avverte che sarà difficile, perché ``sono cose
che si iniziano a capire mettendole in pratica''. Per questo invita di nuovo a
fare domande sulle proprie esperienze.

% ---------------------------------------------------------------------
\section{Le sprint e le riunioni}
\label{sec:z3-sprint}

\subsection{Il libro di Scrum}

Simone mostra il libro in cui viene raccontata la genesi del
\emph{framework} Scrum, scritto da Jeff Sutherland\footnote{Si tratta
verosimilmente di J.~Sutherland, \emph{Scrum: The Art of Doing Twice the Work
in Half the Time} (2014), edito in italiano come \emph{Fare il doppio in metà
tempo}. Il titolo non viene pronunciato: Simone mostra la copertina e più
avanti cita ``Scrum, di Jeff Sutherland''. La lunghezza indicata (``circa 300
pagine'') è approssimativa.}: un ``libricino'' di circa trecento pagine.
Avverte che userà le parole ``metodologia'' e ``framework'' in modo
intercambiabile. Lo consiglia assolutamente a chi è interessato al tema,
perché è scorrevole e contiene \emph{insight} molto interessanti; chi non è
particolarmente interessato può evitarlo del tutto, perché non avrà niente da
offrire che non venga detto in questa lezione.

\subsection{Che cos'è una sprint}

Scrum, a differenza di Agile, è un \emph{framework} a tutti gli effetti: ci
sono regole da seguire in modo preciso e da rispettare. Il modo in cui
abbraccia il principio dell'\emph{inspect and adapt} sono le
\textbf{sprint}: periodi di tempo tipicamente di \textbf{una o due
settimane}. Possono essere anche di quattro, come si studia, ma secondo Simone
non ha senso farle così lunghe a meno di non lavorare a un grande progetto
``con duecento miliardi di persone che ci stanno lavorando
costantemente''\footnote{Nella trascrizione ``a meno che non state facendo un
sesto plà con 200 miliardi di persone'': ``sesto plà'' è incomprensibile e
potrebbe essere una storpiatura di ``tripla A''; il senso, in ogni caso, è
quello di una produzione enorme.}.

Prima di ogni sprint c'è una riunione di \textbf{sprint planning}, dopo ogni
sprint una riunione di \textbf{post-mortem}\footnote{Nella guida ufficiale di
Scrum la riunione di fine sprint è divisa in \emph{sprint review} (che cosa è
stato prodotto) e \emph{sprint retrospective} (come ha lavorato il team);
qui, come nel resto della lezione, entrambe confluiscono nel termine
\emph{post-mortem}, che in senso stretto indica l'analisi a progetto
concluso (sezione~\ref{sec:z1-qa-fasi}).}.

\subsection{La regola d'oro: meno riunioni possibili}

``Le riunioni non piacciono a nessuno'': sono, dice Simone, una ``rottura di
scatole infinita'', e chiunque, soprattutto chi non è designer, non aspetta
altro che poter dire ``questa riunione è inutile, dobbiamo andarcene tutti a
lavorare perché abbiamo cose più importanti da fare''. In chat qualcuno
risponde con il meme di Bernie Sanders a braccia incrociate; Simone avrebbe
voluto mettere quello del cane che urla sulla scrivania perché ``tutta la
riunione poteva essere una mail'', e promette che dalla prossima volta ogni
slide avrà un meme alla fine.

Da qui la sua \textbf{regola d'oro}: \textbf{fare meno riunioni possibile}.
Sottolinea però che è una regola ``per quanto mi riguarda'': magari un team
non vede l'ora di fare riunioni, magari si portano le pizzette a ogni
incontro (``non so se con le pizzette le riunioni siano effettivamente
efficienti, però se funziona, perché no?'').

Una mossa molto semplice per ridurle è \textbf{unire sprint planning e
post-mortem} in un'unica riunione: si fa letteralmente la metà delle
riunioni. Ciò che si fa in quella riunione è:
\begin{itemize}
  \item \textbf{pianificare} la prossima sprint (\emph{sprint planning});
  \item \textbf{analizzare} la sprint precedente (\emph{post-mortem}),
        indicando in particolare i problemi incontrati e le cose rimaste
        bloccate.
\end{itemize}

\subsection{Una riunione deve essere rapida}

La riunione unica deve essere \textbf{rapida}. Durante l'ultimo anno, per il
\emph{Biggest One}, il team di Simone contava circa quindici persone, e alle
riunioni di planning e post-mortem ne partecipavano in media una dozzina.
Quando sono ``diventati bravi'', le riunioni duravano \textbf{venticinque
minuti, con dodici persone che parlavano tutte}. Quello è l'obiettivo, ``anzi,
il vostro obiettivo è fare anche meno''.

Il motivo per cui si sfora è una \textbf{tentazione} costante: quella di
\textbf{risolvere i problemi}. Quando qualcuno descrive un problema, un paio
di persone cominciano a discuterne (``hai provato a fare questo?'', ``no,
dovresti fare quest'altro''), e in un attimo sono passati dieci minuti in cui
hanno parlato solo in due senza che la riunione sia andata avanti di un
passo. In queste riunioni non si fa brainstorming e non si ragiona su come
risolvere i problemi: ``bloccatelo sul nascere, picchiate quelle persone, fate
qualsiasi cosa per impedirglielo''.

Quando c'è un problema, si \textbf{prende un appuntamento per risolverlo}.
Simone mostra come dovrebbe suonare un report:
\begin{quote}
``La settimana scorsa mi sono occupato della logica di movimento del
personaggio, ho fatto un po' di tuning e di testing e ho preparato il block
out del primo livello. L'unico problema è che il blueprint che faceva volare
via gli uccelli non mi ha funzionato bene: visto che tu, Cristian, l'hai
utilizzato in un altro livello, più tardi ci sentiamo su Discord quindici
minuti e mi spieghi come l'hai settato. Questa settimana mi occuperò di
completare il block out e di risolvere il problema del blueprint.''
\end{quote}
Fine del report: \textbf{circa un minuto}, e c'è tutto. Tutti devono fare
così: trovare una persona con cui risolvere il problema, identificare
\emph{come} risolverlo e andare oltre. ``C'è tempo e luogo per ogni cosa'', ma
la riunione di sprint planning non è il momento giusto.

\subsection{La sprint come unità di tempo}

Ultimo punto: le sprint sono \textbf{l'unità di tempo} fondamentale per ogni
cosa. Le \emph{milestone} non si misurano in mesi, giorni o settimane, ma
\textbf{in sprint}. Bisogna prendere l'abitudine di dire ``completeremo la
preproduzione in otto sprint'' e non ``completeremo la preproduzione tra due
mesi''. È importantissimo per il calcolo della velocity
(sezione~\ref{sec:z3-velocity}).

% ---------------------------------------------------------------------
\section{Le cerimonie e la regola della consistency}
\label{sec:z3-cerimonie}

\subsection{Le cerimonie di Scrum secondo Simone}

Le regole di Scrum si chiamano \textbf{cerimonie}. Simone le divide in due
gruppi\footnote{La classificazione è quella proposta da Simone. Nella guida
ufficiale di Scrum gli ``eventi'' sono la sprint, lo \emph{sprint planning},
il \emph{daily scrum}, la \emph{sprint review} e la \emph{sprint
retrospective}; il backlog è un ``artefatto'' e il calcolo della velocity
e il \emph{planning poker} sono pratiche diffuse ma non obbligatorie.}:
\begin{center}
  \begin{tabular}{@{}p{0.42\textwidth}p{0.48\textwidth}@{}}
    \toprule
    \textbf{Cerimonie principali} & \textbf{Altre cerimonie} \\
    \midrule
    planning; post-mortem; calcolo della velocity; creazione del backlog &
    \emph{daily scrum}; \emph{poker time estimation} (\emph{planning poker})\\
    \bottomrule
  \end{tabular}
\end{center}
Le quattro principali sono quelle che ``non potete evitare di fare,
altrimenti non state facendo Scrum'': il che, precisa, non è necessariamente
un problema (ci torna alla fine, sezione~\ref{sec:z3-ricapitolo}). Le altre
due le tiene separate perché a lui non piacciono granché, e prima di spiegare
perché beve un sorso d'acqua ``per lasciarvi la suspense''.

\subsection{Consistency: una cerimonia saltata fa cadere le altre}

In Scrum c'è una parola chiave da ricordare: \textbf{consistency}. Bisogna
evitare a tutti i costi di creare \textbf{precedenti di
\emph{inconsistency}}, cioè di incostanza: con Scrum bisogna rimanere
costanti. \textbf{Se si salta una sola cerimonia si crea un precedente che,
a effetto domino, butta giù tutto il resto.} Questo è il punto già sollevato
da Mattia. Ma, avverte Simone, questo \emph{non} vuol dire che ci si debba
forzare: ed è questo il punto chiave.

\subsection{Perché il daily scrum non funziona più}

Il \emph{daily scrum} era pensato ``per il millennio scorso''. Scrum è nato
molto tempo fa\footnote{Simone non ricorda la data esatta (``fine degli anni
Novanta, forse fine degli anni Ottanta''). Il termine \emph{scrum} applicato
allo sviluppo di prodotti risale a un articolo di Takeuchi e Nonaka del 1986;
il \emph{framework} Scrum fu presentato da Jeff Sutherland e Ken Schwaber nel
1995.}, per team che si incontravano tutti i giorni in ufficio alla stessa
ora, con un orario d'ufficio standard: team di sviluppo software, quindi per
lo più programmatori (anche i più ``artistici'', come gli UX designer, erano
alla fine programmatori), con una mentalità piuttosto simile e, soprattutto,
con lo stesso orario.

Oggi, Covid a parte, questo non è più né fattibile né sensato. Gli studenti
non hanno nemmeno un ufficio: si incontravano a volte in sede quando si
poteva, ma per lo più lavorano online, e \textbf{non tutti lavorano allo
stesso orario}. ``E va bene così'': c'è chi lavora di notte (la salute
potrebbe risentirne, ma se si sentono bene così, va bene), e quelle persone
non scriveranno il proprio \emph{daily scrum} all'orario in cui lo fanno
tutti. Senza arrivare a casi del genere, Simone lavora per una
\textbf{società australiana}, dal lato opposto del globo: quando lui si
sveglia, loro stanno finendo di lavorare, e il suo \emph{daily scrum} cadrebbe
nella loro notte.

In queste condizioni è difficile tenere traccia del \emph{daily scrum} e
rispettarlo, e prima o poi qualcuno lo salterà. Quello è un
\textbf{precedente di incostanza}: qualcuno comincia a dire ``ok, non lo faccio
neanche io'', poi ``ma visto che non faccio più gli scrum, a che servono le
riunioni? Questa volta non vengo'', ``e io neanche'': ``e da lì è il caos''.

La conclusione è netta: \textbf{quando qualcosa non funziona, tagliatelo}. Il
senso di ``individui e interazioni sono più importanti di processi e
strumenti'' è proprio questo: se un processo o uno strumento non funziona
\textbf{non è colpa vostra, è il team che non ne ha bisogno}. Non bisogna
fissarsi su ciò che si vuole fare, ma sentire di che cosa ha bisogno il team e
adattarsi a quel bisogno.

\subsection{Perché il planning poker rallenta tutto}

Il \emph{planning poker} (che Simone chiama \emph{poker time estimation})
non gli piace perché lo trova \textbf{estremamente lungo e tedioso}: rallenta
un processo già lento di suo. Creare il backlog e valutare i task è già
lunghissimo (``ci vuole almeno un'intera sprint di riunioni per fare un
backlog''); aggiungerci il planning poker significa moltiplicare il tempo. Si
informa con Diego su come sia stato spiegato alla classe: con il metodo di
Fibonacci, ``\emph{power-uppato}''\footnote{Nella trascrizione ``è rimasta
nelle lezioni di trial'' e ``è stata power upata'': il senso sembra essere che
il planning poker è stato spiegato nelle lezioni precedenti in una versione
potenziata con la sequenza di Fibonacci. ``Trial'' potrebbe essere una
storpiatura.}, il che farà risparmiare un po' di spiegazioni più avanti.

Simone capisce la necessità di sentire le opinioni di tutti
(``\emph{empowering the team}'', eccetera), ma l'effetto è rendere il processo
ancora più lento. E poiché le riunioni non piacciono a nessuno, se si
quintuplica il tempo necessario ``si parte subito male'', perché il team si
annoia già. Anche questo, quindi, lui personalmente non lo usa.

\subsection{Il punto chiave: crearsi il proprio framework}

La cosa più importante è che \textbf{non bisogna eseguire tutto a menadito}.
I \emph{framework} sono stati pensati per essere generici, ma anche per
essere \textbf{adattati}. La procedura che Simone suggerisce:
\begin{enumerate}
  \item si parte \textbf{rispettando il framework al 100\%};
  \item quando qualcosa non funziona, si capisce \emph{perché} non funziona;
  \item si smette di farlo, si risolve, si fa altro;
  \item si capisce qual è il framework più adatto e \textbf{ci si crea il
        proprio}.
\end{enumerate}
``Questo è il punto chiave di tutto questo.''

% ---------------------------------------------------------------------
\section{Il backlog}
\label{sec:z3-backlog-simone}

\subsection{Una ``rottura di scatole'' inevitabile}

Sulla creazione del backlog Simone spende ``giusto due parole'', perché la
parte è molto tecnica e i consigli sono pochi. Creare il backlog è faticoso,
``non ci sono mezze misure per dirlo'': richiede moltissimo tempo e lavoro ed
è intenso a livello mentale, perché bisogna ragionare tantissimo.

\subsection{Il backlog a scaglioni}

In teoria i backlog si creano all'inizio dello sviluppo: in un mondo perfetto
si conoscerebbero fin da subito tutti i task necessari per completare il
progetto. Nel mondo reale questo non accade. Ciò che accadrà, più
probabilmente, è che si andrà avanti \textbf{per milestone}. È ciò che Simone
ha cominciato a fare con \emph{ISO} e con cui si trova bene:
\begin{itemize}
  \item ha fatto il backlog della \textbf{preproduzione}, con tutto ciò che
        serviva per la preproduzione;
  \item finita la preproduzione, ha fatto quello della \textbf{produzione};
  \item finita la produzione farà quello della fase successiva
        (``espansione, build finale, non lo so''), a seconda di come si
        svilupperanno le cose.
\end{itemize}
Fare i backlog \textbf{a scaglioni} li rende più semplici e il processo meno
lungo; e tanto, comunque, bisognerà aggiungere cose in corso d'opera: ``non si
scappa''.

\subsection{Le epic}

L'organizzazione in \textbf{epic} è molto utile, soprattutto per la stesura:
si comincia ragionando per \emph{epic} e poi si scende ai task. Un esempio di
epic è ``creazione del personaggio''. Che cosa contiene il personaggio?
Movimento, tuning, modello, animazioni, materiale, il design delle feature
che dovrà avere, la creazione di tutte queste feature, e così via. Il
processo diventa più \emph{streamlined}, meno pesante e più semplice da
seguire.

\subsection{Ogni task deve avere un obiettivo chiaro}

Fondamentale: \textbf{i task devono avere un obiettivo chiaro}. Un consiglio
comune è scriverli come \textbf{user story}: ``come utente voglio avere un
personaggio che si senta potente, che io riesca a percepire come potente''.
Aiuta moltissimo nello sviluppo software e parecchio nei videogiochi per
alcune cose, come appunto il personaggio. In generale è buona pratica
\textbf{spiegare l'obiettivo e non la soluzione}: dire che il personaggio deve
muoversi, saltare e fare le capriole, invece di dire che ``quando premo il
tasto W deve subire una forza che lo spinge in avanti''. Però non è
fondamentale.

Ciò che è davvero importante è che il task abbia un \textbf{obiettivo
verificabile}: qualcosa di cui si possa dire ``è completo'' o ``non è
completo''. Altrimenti le persone \textbf{non lo completeranno}, o lo
completeranno in modo diverso da come lo si immaginava. Un esempio di task
senza obiettivo chiaro è \textbf{``completare il livello~5''}:
\begin{itemize}
  \item per Simone vuol dire aver creato il block out con qualche elemento
        estetico per sapere dove vanno le cose e come funzionano con i
        moduli già disponibili, e basta;
  \item per qualcun altro potrebbe voler dire inserire illuminazione,
        \emph{post-processing}, \emph{particles}, tutti gli shader fatti
        bene, ottimizzazione e \emph{level streaming};
  \item per altri ancora potrebbe voler dire arrivare all'illuminazione,
        lasciando \emph{level streaming} e ottimizzazione a una seconda fase.
\end{itemize}
È la cosa più importante da ricordare quando si crea il backlog e si
scrivono i task.

% ---------------------------------------------------------------------
\section{La velocity}
\label{sec:z3-velocity}

``Bellissimo, bellissimo'': il calcolo della velocity è per Simone ``la cosa
più figa e allo stesso tempo più odiata''.

\subsection{Il peso dei task}

Come la classe sa, i task sono valutati in base al loro \textbf{peso} (o
punteggio, o \emph{score}). Il peso dipende da quanto un task è difficile e/o
da quanto tempo richiede: è un po' una zona grigia, e Simone consiglia di
valutarlo come \textbf{combinazione delle due cose}.

\subsection{Mai le ore}

Il consiglio fondamentale, che Simone non capisce perché si continui a
contraddire, è \textbf{non usare le ore per valutare i task}, ``per niente,
neanche alla lontana''. Nel libro di Sutherland c'è un paragrafo intero che
spiega perché. Diego commenta che chi lo consiglia di solito fa ``Agile
Waterfall'' e non ha nemmeno letto i libri su Agile.

Il motivo: \textbf{il cervello umano è completamente incapace di stimare le
ore}, cioè di prevedere il tempo necessario per fare qualcosa: ``non siamo
pensati per questo''. Siamo invece \textbf{molto bravi con le misure
relative}: fissato un task che vale 10, sappiamo dire bene quale vale 5,
quale 7, quale 2. Sappiamo dire che un task è molto più grande di un altro.

\subsection{La sequenza di Fibonacci}

Per aiutarsi ulteriormente, Simone usa la \textbf{sequenza di Fibonacci} (1,
2, 3, 5, 8, 13, 21, 34, \dots), anche questo un consiglio di Sutherland che
``nessuno segue''. Sutherland racconta addirittura di un team che usava le
\textbf{taglie dei cani} (``questo task è un pinscher, questo è un golden
retriever'')\footnote{Nel libro di Sutherland le stime per razze canine sono
citate come esempio di scala relativa alternativa ai numeri; nella
trascrizione ``pincher''.}; a Simone sembra eccessivo, e ``i buoni vecchi
numeri'' bastano. Fibonacci ha due scopi molto utili:
\begin{enumerate}
  \item \textbf{Gli intervalli crescenti assorbono l'errore di stima.} Man
        mano che un task diventa più grande e complesso, cresce anche
        l'errore associato. È facile dire che aprire la mail vale 2 e
        rispondere a tutti sul server Slack vale 3: che cosa può andare
        storto? La mail non si apre e bisogna reinserire la password; è
        esploso qualcosa e ci sono cinquanta messaggi invece dei soliti
        venti. Niente che cambi granché. È molto più difficile, invece,
        valutare la differenza tra creare il \emph{behavior} di movimento del
        personaggio e costruire un intero livello: possono succedere
        infinite cose, soprattutto quando c'è di mezzo la programmazione.
        Il behavior del personaggio è un 20 o un 21? ``È un 21, punto'':
        non è un 34, perché non è così difficile, e non è un 13, perché
        sarebbe troppo poco. Magari la cosa più corretta sarebbe stata 23,
        ma non importa: lo scopo è \textbf{appianare} i risultati.
  \item \textbf{Meno scelta, meno discussione.} I valori sono
        ``infinitamente arbitrari'', ``completamente, al 100\%'', e nessun
        altro al mondo li condividerà; ma Fibonacci li rende un po' più
        condivisibili, perché con meno possibilità di scelta ci sono meno
        occasioni di discutere.
\end{enumerate}
Che i valori siano arbitrari non è un problema: finché \textbf{è una persona
sola a darli}, sono tutti arbitrari allo stesso modo, quindi
\textbf{consistenti}; e se sono consistenti hanno una logica e si possono
usare per fare calcoli.

\subsection{Che cos'è la velocity e perché serve}

Sommando il peso dei task completati in una sprint si ottiene la
\textbf{velocity}. Ed è ``fighissima'' perché permette, in sostanza, di
\textbf{prevedere il futuro}.

Simone mostra il grafico della velocity della preproduzione di \emph{ISO}
durante l'anno scolastico. La consegna era a maggio (``10 maggio, o poco
dopo''). Già a inizio marzo, guardando l'andamento, aveva capito che \textbf{non
sarebbero riusciti a finire in tempo}. Il ragionamento era semplice: mancano
dodici sprint; quanto facciamo ogni sprint? Circa 89; quanto manca nel
backlog? Circa 2000. ``Non ce la facciamo''\footnote{I numeri sono quelli
detti a voce: 12 sprint per circa 89 punti danno poco più di 1000 punti,
circa la metà dei 2000 rimasti nel backlog.}.

La velocity aiuta a prevedere il futuro e quindi a \textbf{prendere
provvedimenti prima}. A marzo Simone ha cambiato i piani: prepariamo meno
cose, smettiamo di perdere tempo su ciò che ci blocca, consegniamo qualcosa di
diverso. In più, la velocity del team è \textbf{aumentata} sempre di più, e
alla fine sono arrivati in fondo alla preproduzione \textbf{senza crunch}:
``zero, completamente zero crunch''. Nessuno ha dovuto lavorare di più;
l'ultima settimana non hanno nemmeno toccato il progetto: hanno fatto la
build, l'hanno spedita e per una settimana non ci hanno lavorato. E sono
riusciti a consegnare anche di più del previsto, perché Simone ha
\textbf{analizzato i trend precedenti}: perché in un certo punto c'era un
picco? Che cosa li bloccava e faceva diminuire la produttività? L'ha
scoperto, l'ha capito e hanno agito.

\subsection{Come calcolarla in pratica}

Il consiglio è di \textbf{trovare un modo per fare il calcolo}. Simone usava
Monday, che non aveva uno strumento costruito apposta, ma ci è riuscito lo
stesso; si può fare tranquillamente con \textbf{Google Sheets}, ``col
sorriso''. Per farlo davvero bene bisognerebbe scrivere un po' di script
dentro il foglio, ``niente di improponibile''. Quando comincia a girare
``funziona da dio'': ``assolutamente fatelo, non posso dirlo abbastanza''. (Un
esempio costruito in diretta è nella sezione~\ref{sec:z3-qa-simone}.)

\subsection{Come assegnare i punteggi}

Per assegnare gli \emph{score} conviene partire dagli
\textbf{estremi}: quali task sono un 34? Quali sono un 1? Da lì si mettono i
valori intermedi. Gli estremi fanno da \textbf{ancore} e facilitano il
processo di valutazione \emph{relativa} descritto sopra.

Lo stesso vale per le \textbf{priorità}, se si usa un sistema di priorità:
regole simili, ma con \textbf{intervalli più piccoli} e senza Fibonacci, che
lì non serve. Simone usa una scala da 1 a 5; ma non la considera così
fondamentale.

\subsection{Che cosa blocca che cosa}

Durante il lavoro sul backlog bisogna capire \textbf{che cosa blocca che
cosa}. Ad esempio: ci sono 29 task, per un valore complessivo di 400, tutti
relativi alla creazione dei livelli. Ma se non c'è il movimento del
personaggio, non ha senso fare i livelli, perché vanno costruiti attorno a
quel movimento: la creazione del personaggio \textbf{blocca} i livelli e
quindi ha priorità più alta. Da qui si comincia anche a capire quale reparto
si occupa di che cosa: probabilmente all'inizio i programmatori faranno i
\emph{behavior} fondamentali (come il personaggio), mentre i designer
lavoreranno su documentazione e ricerca. Ma dipende dal gioco: alcuni non
costruiscono i livelli attorno al personaggio, bensì il personaggio attorno ai
livelli, quindi anche l'esempio ``è opinabile''.

% ---------------------------------------------------------------------
\section{Ricapitolando: come nasce una sprint}
\label{sec:z3-ricapitolo}

Prima di passare alle domande, perché ``quello che avevo da dirvi di pratico è
quasi finito'', Simone ricorda che il grosso non sono le regole, ma
\textbf{capire la logica alla base e integrarla nel proprio team e nel
proprio progetto}. Poi riassume la creazione di una sprint in quattro passi.

\begin{enumerate}
  \item \textbf{Post-mortem della sprint precedente.} Si vede che cosa ha
        fatto ciascuno, si capisce che cosa è bloccato e si
        \emph{programma un incontro} per risolvere i problemi.
  \item \textbf{Pianificazione.} Dal backlog si guarda che cosa blocca che
        cosa, quali sono le priorità e che cosa si può fare in questa sprint,
        e si prendono i task.
  \item \textbf{Analizzare i problemi e agire.} Si guarda la velocity: è in
        aumento? In calo? Si è in un picco? L'andamento a dente di sega,
        ``su, giù, su, giù'', è \textbf{assolutamente normale}, ed è ciò che
        ci si deve aspettare; e man mano che il progetto va avanti la
        velocity \textbf{tenderà a diminuire}. Le domande da farsi: come sta
        andando il progetto? Chi sta facendo di più e chi di meno, e perché?
        Qualcuno ha avuto un calo di rendimento? C'è un problema che si
        ripresenta spesso, un collo di bottiglia? E come risolverlo,
        soprattutto \textbf{ascoltando il resto del team}.
  \item \textbf{Rimanere consistenti}, in particolare con Scrum.
\end{enumerate}

\subsection{Ascoltare il team}

Come scrum master o producer, il lavoro più importante è e rimane
\textbf{ascoltare il team}. Tutto il resto non conta se poi, quando il team
dice una cosa, non lo si ascolta. È lo stesso principio dell'adattare il
\emph{framework} al team:
\begin{itemize}
  \item se il team odia il \emph{daily scrum}, non si fa: non serve a niente;
  \item se il team fa fatica a prendere i task dal backlog, non glielo si fa
        fare: si assegnano i task in modo diverso, ``non fategli manco vedere
        il backlog'';
  \item se il team vuole sprint di due settimane invece che di una, si fanno
        di due settimane.
\end{itemize}
Magari le persone si organizzano personalmente in modo diverso: Simone, per
esempio, organizza tutto il suo lavoro, su tutti i fronti, \textbf{su base
settimanale}, e una sprint di due settimane gli sconvolgerebbe i piani. Se il
team lavora per lo più così, si fa così. ``Adattatevi'' non vuol dire solo
adattarsi ai bisogni del software, ma \textbf{soprattutto ai bisogni del
team}.

\subsection{Nessun framework va bene al 100\%}

Esistono infiniti \emph{framework}, e \textbf{nessuno andrà bene al 100\% per
un team}. Bisogna provarli, il più possibile, e adattarli: cambiare qualche
regola, togliere qualche cerimonia. Il \emph{framework} più usato oggi, ad
esempio, è \textbf{Kanban}, che non ha nemmeno le sprint\footnote{Simone dice
prima ``Kanban neanche ce l'ha il backlog'' e poi si corregge: ``scusatemi,
neanche ce l'ha le sprint, ha soltanto un gigantesco backlog''.}: ha soltanto
un gigantesco backlog. ``Ogni team è diverso, ogni team ha bisogni diversi,
quindi adattatevi'' (una frase che Simone calcola di aver ripetuto ``almeno 49
volte'').

\subsection{La nota ``smielata'': la fiducia}

La conclusione: alla fine di tutto, \textbf{Agile non si basa sulle
iterazioni}, né sull'efficienza, né sulla creazione del miglior videogioco o
software del mondo, né sul controllo. L'unica cosa su cui si basa davvero è
\textbf{la fiducia}:
\begin{itemize}
  \item la fiducia che il team ripone in chi lo guida, e che chi lo guida
        ripone nel team: ``queste persone lavoreranno nel miglior modo
        possibile, sempre e comunque'';
  \item la fiducia del team nel dire ``se qualcosa va male, so che lo scrum
        master ha un piano di riserva che ci risolverà il problema, e il mio
        lavoro non sarà sprecato'';
  \item la fiducia di chi guida nel dire ``non posso controllarvi
        costantemente, non posso assicurarmi che facciate sempre tutto il
        lavoro o che seguiate le cerimonie come dovreste, però mi fido'': mi
        fido che lo farete, che se ci saranno problemi ve lo direte, che
        farete sempre del vostro meglio.
\end{itemize}
Bisogna \textbf{costruire fiducia nel team e trasmettere fiducia al team}. È
lo stesso principio del \emph{default to trust} della prima lezione zero
(sezione~\ref{sec:z1-default-trust}).

Simone ha finito con mezz'ora d'anticipo e apre alle domande (``go crazy'').
Diego apprezza molto il finale, ``smielato ma in realtà estremamente
veritiero'': lui è più orientato all'ambito professionale, ma crede che in
Agile ci debba essere davvero questa fiducia, e che sia anche una delle sue
forze. Simone pensava che non l'avrebbe apprezzato, tanto era smielato
(``mentre lo scrivevo ho percepito il diabete che aumentava''); Diego gli
suggerisce, la prossima volta, di aggiungere immagini con gli occhi a cuore e
qualche gattino.

% ---------------------------------------------------------------------
\section{Domande a Simone}
\label{sec:z3-qa-simone}

\subsection{La dimensione del team allunga le riunioni?}

Uno studente, che lavora in un team di tre persone (``essendo in tre le cose
vengono un po' più facili''), chiede quanto la \textbf{grandezza del team}
influisca sui tempi delle riunioni e delle cerimonie, e se ci sia un modo per
mantenerle rapide comunque.

Influisce assolutamente, ``per una mera questione matematica''. Ma in genere
è possibile far sì che un team di qualsiasi dimensione completi una riunione
in \textbf{una trentina di minuti al massimo}: quello è l'obiettivo, ``anche a
costo di dire meno cose, anche a costo di essere meno precisi''. A meno di
avere un team di una cinquantina di persone, la riunione durerà sempre trenta
minuti; e con cinquanta persone non si fa comunque una riunione tutti insieme,
perché è impossibile: si fanno le \textbf{riunioni di reparto}, e poi i
\textbf{capi reparto} fanno la loro riunione di sprint planning. I numeri
tornano piccoli: a una riunione ci saranno al massimo una quindicina di
persone, venti ``se proprio i pianeti si sono allineati'', e venti persone
possono stare in trenta minuti.

\subsection{Quanto deve imporsi lo scrum master? Il timekeeper}

Lo stesso studente chiede quanto lo scrum master debba \textbf{imporsi} sugli
altri, sul piano umano, per far rispettare i tempi. Moltissimo, risponde
Simone, ma \textbf{non deve per forza essere lo scrum master a farlo}. Nel
team australiano in cui lavora, all'inizio di ogni riunione si sceglie, in
modo pseudo-casuale, tra chi è disponibile o lanciando un dado, un
\textbf{timekeeper}: la persona che ha ``lo scettro del comando'' e tiene in
mano un cronometro, perché nessuno dedichi ai propri report più tempo del
dovuto. È \textbf{comunemente accettato da tutti} che quella persona abbia il
diritto di ``rompere le scatole'': è socialmente accettato.

Se le riunioni si dilungano troppo, si può proporre di introdurre la regola
del timekeeper, scelto ogni volta a caso o sempre lo stesso: ha il diritto di
richiamare gli altri perché, oggettivamente, il tempo scorre.

\subsection{E se il problema blocca tutti?}

Ultima domanda dello stesso studente: si è detto che i problemi si risolvono
dopo la riunione; ma se un problema è molto grave, se il fatto di non averlo
risolto blocca la produzione di più persone (``praticamente un disastro''),
vale la pena fermarsi e parlarne tutti insieme, o si dedica comunque quel
tempo solo ai report e poi si fa una riunione in più?

La risposta semplice è \textbf{no}: non si risolve nella riunione di sprint
planning e post-mortem. ``L'unica cosa che mi viene in mente che possa essere
così tanto grave è se la stanza in cui state facendo il meeting va a fuoco.''
La riunione non è pensata per quello. Subito dopo, piuttosto, si avvia
un'altra riunione con tutte le persone necessarie (``che durerà 25 ore'') per
risolvere il problema. Ma finché si fa planning e post-mortem, si fa
\textbf{quello e basta}, anche se il problema riguarda tutti.

\subsection{E se mancano delle persone?}

Si fa comunque la riunione? La si rimanda? Simone la fa \textbf{sempre}, a
meno di essere da solo, per \textbf{rispetto delle persone che sono venute}.
Le assenze possono avere cause diverse. Chi non è venuto deve mandargli un
messaggio su Telegram, anche vocale, in cui fa di fatto il proprio report.

Se una persona manca spesso, Simone ci parla e cerca di capire perché.
Magari l'orario della riunione non va bene perché ha altro da fare: nel
gruppo di \emph{ISO} c'è un ragazzo che lavora a settimane alterne di mattina
e di pomeriggio, e quando lavora di pomeriggio, tornando a casa, trova spesso
traffico; quindi a settimane alterne non può venire, o arriva molto in
ritardo. ``Lo so, me lo posso aspettare, non è un problema.'' Se invece
qualcuno manca in modo pseudo-casuale, bisogna capire il motivo: non gli va
bene l'orario? Ne preferirebbe un altro? Pensa che le riunioni non servano?
Non ha voglia di farle? Da lì si capisce il problema e si cerca di
risolverlo, spiegando che \textbf{le riunioni sono importanti} e che è
importante esserci tutti, e chiedendo ``dimmi come possiamo aiutarti a
rendertele meno pesanti''.

\subsection{Quando è la persona a non adattarsi}

Se però, nonostante tutto, la persona ``continua a fare la \emph{drama
queen}'', a un certo punto bisogna ammettere che esistono anche persone non
adatte a lavorare in un certo contesto, e si può valutare se lo sia.

Diego interviene per completare: in Agile tutti vogliono ottenere la cosa
migliore per il team e per il metodo di lavoro di ciascuno. Ma se un individuo
non si adatta al team, o il team non riesce ad adattarsi a lui, che sia o
meno una \emph{drama queen}, \textbf{la presenza di quella persona nel team va
valutata}. ``Nessuno, a parte casi eccezionali, tipo Miyamoto, vale il lavoro
degli altri.'' Le \emph{soft skill} restano alla base. E aggiunge che non è
solo questione di soft skill: ci sarà la \emph{drama queen}, quello che non si
adatta, ``lo stronzo o la stronza'' (il genere, precisa, è irrilevante),
quello che non ha voglia\footnote{Nella trascrizione ``stronzo
stronza,umbaquin'': la parte finale è incomprensibile.}; ma c'è anche la
persona che ha voglia di fare ma non si adatta al team. Se non ci riesce, o il
team non riesce ad adattarsi a lei, potenzialmente deve andare via.

Simone conferma che è un punto molto importante, che rientra nel discorso
dell'\textbf{ascoltare di che cosa ha bisogno il team}: se il team non riesce
a lavorare con un individuo, ``non c'è niente di male, è così e basta''. Per
fortuna non è troppo comune: a un certo livello professionale è difficile che
succeda, soprattutto in questa industria, dove le persone, se non stessero
lavorando insieme, potrebbero diventare amiche, perché hanno un interesse in
comune e sono molto appassionate. Ma può succedere.

\subsection{Come si calcola la velocity, in pratica}

Diego invita gli studenti ad approfittare di una persona ``che ha fatto il
vostro percorso''. Una domanda sul calcolo: la velocity è semplicemente la
\textbf{somma dei pesi dei task completati} nel tempo trascorso, cioè nella
sprint? Esatto. Si fa la sommatoria dei task \textbf{taggati come
\emph{done}} che hanno la stessa sprint (o la stessa data di consegna, a
seconda di come la si indica): si prendono tutti i task con ``sprint 5'', di
questi quelli in stato \emph{done}, e se ne sommano i pesi. Allo stesso modo
si possono sommare i task non completati, quelli rimandati, e così via.

\paragraph{Uno sguardo ai fogli personali.} Simone vorrebbe mostrare il foglio
di \emph{ISO}, ma è difficile da leggere per la quantità di dati; si chiede se
mostrare i propri fogli rompa qualche NDA, poi nota che sta già condividendo
lo schermo, ``quindi gli NDA sono già stati rotti''. Mostra allora uno dei
suoi fogli personali: organizza il proprio lavoro \textbf{settimana per
settimana}, con il nome della settimana, l'elenco dei task, la somma di ciò
che era da fare, quanto è stato effettivamente fatto, quanto è stato tagliato
e quanto rimandato per cause diverse. Il risultato è un grafico con
l'andamento settimanale di questi valori.

\paragraph{La ``speedrun velocity''.} Poi decide di fare ``una cosa più
figa'' e di costruire un foglio in diretta, che chiama \emph{speedrun
velocity}. La struttura minima\footnote{I riferimenti alle celle detti a voce
durante la costruzione (``B2B'', ``G1'', ``F1'', ``C2, F2'') sono confusi
nella trascrizione; la tabella e la formula sono una ricostruzione dello
schema, non una copia esatta del foglio mostrato.}:
\begin{center}
  \begin{tabular}{@{}llll@{}}
    \toprule
    \textbf{Task name} & \textbf{Score} & \textbf{Sprint} & \textbf{Status} \\
    \midrule
    Task 1 & 5  & Sprint 1 & done \\
    Task 2 & 13 & Sprint 1 & WIP \\
    Task 3 & 3  & Sprint 2 & cut \\
    \dots  & \dots & \dots & \dots \\
    \bottomrule
  \end{tabular}
\end{center}
Nella colonna \emph{Status} inserisce una convalida dei dati con i tre valori
\emph{cut}, \emph{WIP} e \emph{done}, ``per semplificarsi la vita''. Poi
scrive una serie di task con punteggi a caso (copiati da un altro foglio) e
assegna le sprint, avvertendo che ``ovviamente non dovrebbe essere fatto
esattamente così, ma \emph{for the sake of showcase}''; aggiunge qualche
colore, perché ``sono pazzo'', anche se non è la cosa fondamentale.

Il calcolo è una somma condizionale su più criteri (\texttt{SUMIFS}, in
italiano \texttt{SOMMA.PIÙ.SE}): si somma la colonna dei punteggi per le
righe in cui la sprint è quella cercata \emph{e} lo stato è \emph{done}. Ad
esempio:
\begin{center}
  \texttt{=SUMIFS(B:B; C:C; "Sprint 1"; D:D; "done")}
\end{center}
Lo stesso si può fare con \emph{cut} per sapere quanto è stato tagliato. Si
ripete per la sprint~2, e così via, e con i risultati si inserisce un
\textbf{grafico a linee} (Sheets propone il tipo di grafico sbagliato, e
Simone lo corregge). Con più sprint si ottiene un andamento come quello dei
suoi fogli personali o di \emph{ISO}. Chi vuole fare le cose per bene può
aggiungere altre analisi: chi ha fatto più task, chi ha più peso e carico di
lavoro, lo stato generale del progetto (``magari il backlog fosse così
vuoto'').

Era, dice, per essere sicuro di riuscirci; è disponibile a rifare qualche
passaggio, ma alla classe è tutto chiaro. Uno studente chiede se il foglio si
possa avere: Simone lo mette in una cartella e sarà Diego a farlo avere alla
classe.

\subsection{I problemi più gravi incontrati con Agile}

Diego osserva che Simone ha capito come applicare i principi di Agile con una
velocità sorprendente, più di quanto sia successo a lui quando gli fu
spiegato, e che all'inizio si fa spesso ``Waterfall Agile'', cioè Waterfall
con le sprint, credendo di fare Agile\footnote{Il passaggio è frammentato e non
è del tutto chiaro chi abbia fatto ``Waterfall Agile'' all'inizio: dalla
risposta di Simone (``ho fatto anch'io waterfall agile per un po''') sembra
trattarsi di un errore comune, che Diego attribuisce anche a se stesso o agli
studenti in generale.}. Poi la domanda: quali sono stati i problemi più gravi?

\paragraph{I task ``tenuti in caldo''.} Simone si è accorto che le persone non
chiudevano i task al momento giusto. Se diceva che un task richiedeva, per
esempio, quattro ore, qualcuno lo finiva in un'ora e poi, invece di segnarlo
come completato, aspettava (``me lo rivedo più tardi, tanto ho ancora altre
tre ore''). Non funziona così: \textbf{quando un task è completato, è
completato}, punto. Non c'è ``me lo rivedo tra un po''' né ``devo farlo in
queste quattro ore, quindi ci metto quattro ore''. Questo falsava i risultati,
e Simone ha cambiato subito approccio\footnote{L'esempio parla di stime in
ore, in apparente contrasto con il consiglio di non usarle
(sezione~\ref{sec:z3-velocity}); è plausibile che si riferisca proprio a una
fase precedente in cui le ore erano ancora usate, e che il ``cambiamento
immediato'' sia stato anche questo.}.

\paragraph{Le riunioni che nessuno voleva fare.} Durante il secondo anno le
persone non volevano fare le riunioni settimanali. All'inizio le facevano solo
tra designer; poi le hanno estese a tutti, perché non aveva senso limitarle,
ma la risposta era ``questa riunione è inutile, non ha senso''. In quel caso
Simone non è riuscito a risolvere il problema: non sapeva come farlo. C'erano
problemi anche sul piano umano, non tra i designer ma con membri di altre
classi: problemi che, col senno di poi, si sarebbero potuti risolvere, ma
all'epoca non ci sono riusciti. La soluzione adottata è stata non fare le
riunioni con tutto il team: gli altri membri mandavano un \textbf{recap
settimanale per messaggio}. Non era un metodo del tutto fallimentare, ma si
perdono molti benefici delle riunioni tutti insieme: il \emph{team building},
il confronto, il vedersi, il parlare. È stato, dice, ``un grosso problema e
anche un grosso fallimento''.

Per il resto, niente di così grosso da valere la pena di raccontarlo:
valutazioni di task sbagliate, sprint fatte male, e un periodo di
``Waterfall Agile'' anche per lui, finché non ha capito di sbagliare.

\subsection{Che cosa rispondere a chi dice che le riunioni sono inutili}

E oggi, se qualcuno dicesse ``per me queste riunioni sono inutili'', che cosa
gli risponderebbe? \textbf{``Perché?''}. E da lì si capisce qual è il
problema: non vuole dedicarci tempo? Durano troppo? Gli costa fatica? Ha
problemi con le riunioni perché soffre di qualche forma d'ansia e non gli
piace stare con tutti i membri del team, anche fuori dalle riunioni (``tipo
ci becchiamo a giocare online a qualcosa'')? Bisogna capire il perché e
\textbf{tirarlo fuori dalle persone}, perché spesso le persone stesse non sanno
il perché delle cose che dicono e sentono.

\subsection{I saluti a Simone}

Diego ringrazia Simone e spera che gli studenti abbiano appreso il più
possibile: un'esperienza più diretta di quella di Simone e Cristian non si può
avere, perché hanno fatto letteralmente lo stesso percorso (la lezione di
Diego è ``fondamentalmente la stessa'' da un anno all'altro, con qualche
aggiunta). Ricorda che Simone ha concesso anche il foglio della
\emph{speedrun velocity}, che Diego passerà alla classe, e lo saluta per non
rubargli altro tempo.

% ---------------------------------------------------------------------
\section{Cristian: il punto di vista del game designer e project manager}
\label{sec:z3-cristian}

\subsection{Chi è il relatore}

Il secondo intervento è di Cristian, che attiva la webcam (``mi vedete super
zoomato, spero non sia una brutta vista'') e condivide le slide, mentre Diego
cerca di farle andare in modalità presentazione perché vengano meglio nella
registrazione. Cristian ringrazia Diego per la presentazione e si presenta:
\begin{itemize}
  \item è uno \textbf{studente del terzo anno} di Game Design, quindi ha fatto
        lo stesso percorso della classe;
  \item è anche \textbf{tutor} del corso di Game Design per gli studenti
        dell'accademia online;
  \item è \textbf{game designer e project manager} di due progetti:
        \begin{itemize}
          \item \emph{Gamers in Time}\footnote{Così nella trascrizione;
                il titolo non è stato verificato.}, un \textbf{gioco da
                tavolo}, vincitore del \emph{Big One} del primo anno;
          \item \emph{L'Inf}\footnote{Il titolo compare in molte forme
                (``l'Inf'', ``Linf'', ``Leanf'', ``Linfa''); qui si usa
                \emph{L'Inf}, la più frequente. Non è chiaro se la grafia
                corretta sia \emph{Linf}, \emph{L'Inf} o \emph{Linfa}.}, un
                \textbf{videogioco} che ha vinto, ``tra virgolette'', il
                \emph{Bigger One} del secondo anno e che il team ha avuto
                l'opportunità di continuare come \emph{Biggest One} del terzo
                anno, opportunità che ha accettato.
        \end{itemize}
\end{itemize}
In entrambi si è occupato di tutta la parte di design, insieme ad altri
colleghi, ed è per lo più anche \textbf{project manager, scrum master,
producer}, ``come volete definirmi''.

\subsection{Come è strutturata la sua parte}

Oltre all'esperienza personale, la lezione ha due parti principali:
\begin{enumerate}
  \item la \textbf{produzione} effettiva tramite Scrum;
  \item la \textbf{gestione del team} in un ambiente di lavoro Scrum.
\end{enumerate}
Possono sembrare la stessa cosa, ma ``il progetto tendenzialmente è una cosa
e il team è un'altra'': ci sono ovvi rapporti tra le due parti, ma per
ciascuna ci sono precauzioni specifiche.

\subsection{Le premesse}

\begin{itemize}
  \item Quello che descrive è \textbf{un} percorso, un modo di applicare
        Scrum, \textbf{non l'unico}: se non si è d'accordo, o se il team non
        ``sopporta'' questo metodo, si possono fare tutti i cambiamenti che si
        vuole (e qualche consiglio lo darà strada facendo).
  \item Non è una lezione descrittiva come quelle di Diego (che cos'è Scrum,
        le definizioni, le cerimonie, come si calcolano la velocity e il peso
        dei task), ma un \textbf{confronto} sull'esperienza dei due progetti,
        più su \emph{L'Inf} che su \emph{Gamers in Time}, perché Scrum l'ha
        applicato al videogioco e non al gioco da tavolo.
  \item \textbf{Tutto è in continua evoluzione}. Per questo vuole le domande
        della classe: il primo obiettivo è chiarire i dubbi e far uscire tutti
        dalla lezione avendo capito un po' di più la mentalità agile, ``con
        tutti voi stessi''; il secondo è ricevere \textbf{feedback} per
        migliorare non solo le lezioni, ma lo stesso modo in cui applica
        Scrum.
  \item Ha un po' di mal di gola e si fermerà spesso a bere: si scusa in
        anticipo.
\end{itemize}

% ---------------------------------------------------------------------
\section{Il caso di studio: \emph{L'Inf}}
\label{sec:z3-linf}

\subsection{Le caratteristiche del gioco}

Qualcuno in classe forse conosce Cristian o ha provato il gioco a Level Up,
l'evento in cui è stato presentato. Le caratteristiche principali:
\begin{itemize}
  \item è \textbf{single player} e offline: niente co-op, niente
        multiplayer, niente server;
  \item è un'\textbf{avventura fantasy}: un gioco di \textbf{durata fissa}, che
        tendenzialmente si gioca per una \emph{run}; chi l'ha apprezzato di
        più e vuole vedere tutti gli aspetti del design e delle meccaniche, o
        esplorarne i segreti, lo rigiocherà;
  \item il giocatore ha delle \textbf{magie colorate} con cui affronta
        \textbf{orde di nemici colorati}: spara le magie e, a seconda di
        quanti nemici ``scoppiano'' con un solo colpo, ottiene vari bonus,
        come succede in \emph{Candy Crush}.
\end{itemize}

\paragraph{Dati e passione.} Diego interrompe perché ritiene questa parte molto
importante: com'è nato l'elemento davvero innovativo di \emph{L'Inf}? Dallo
\textbf{studio dei dati}, ma anche dalla \textbf{passione del team}. I dati
danno una direzione: le ricerche di Cristian indicavano il successo delle
meccaniche \emph{match-3}\footnote{Nella trascrizione ``match free'' e
``match tree'': si tratta delle meccaniche \emph{match-3} (allineare tre o
più elementi dello stesso colore), tipiche di \emph{Candy Crush}.}; e
Cristian è anche, principalmente, un giocatore mobile. La combinazione tra
passione e ricerca ha prodotto un progetto molto valido sul piano della
fattibilità (single player, fantasy, espandibile).

Cristian aggiunge un'altra ricerca fondamentale, da cui è stato poi delineato
il \textbf{documento di visione}\footnote{Nella trascrizione ``iVision'' e
``l'iVision'': si tratta quasi certamente della \emph{vision}, il documento
di visione (sezione~\ref{sec:perche-documentare}).} del gioco: quella su
che cosa chiedesse Epic Games per finanziare i progetti con l'\textbf{Epic
MegaGrant}. Il team ha individuato tre elementi fondamentali\footnote{È la
lettura dei requisiti fatta dal team, come riferita da Cristian; non è stata
verificata sui criteri ufficiali del programma.}:
\begin{enumerate}
  \item \textbf{innovazione};
  \item \textbf{temi ambientali}: attraverso i colori e l'ambientazione, il
        gioco vuole criticare l'inquinamento globale;
  \item un \textbf{elemento narrativo}.
\end{enumerate}
Questi tre elementi, le ricerche specifiche sui vari tipi di gioco e la
passione (di Cristian in primis, come lead, ma anche degli altri membri,
incuriositi dalla proposta) hanno portato a un progetto che al momento sembra
funzionare: ``stiamo andando nella direzione giusta, almeno secondo quanto
previsto; non si può mai prevedere al 100\% ciò che accadrà''.

\paragraph{Il contenuto.} Oltre alla meccanica \emph{match-3}, il gioco
prevede:
\begin{itemize}
  \item \textbf{quattro modalità di sparo};
  \item \textbf{cinque livelli}, con esplorazione e fasi di combattimento
        contro le orde;
  \item \textbf{puzzle};
  \item \textbf{cinematic} per lo più ambientali: niente attori, doppiatori o
        mimica facciale, ma scene di presentazione dell'ambiente di gioco;
  \item una \textbf{boss fight} finale.
\end{itemize}
Il perché di questo elenco sarà chiaro tra poco.

\subsection{Un ``trauma'' di produzione: il cambio di team}

La produzione si è sviluppata in due sezioni, e \emph{L'Inf}, ``per fortuna o
purtroppo'', ha avuto un grave problema: è la risposta, anticipata, alla
domanda fatta a Simone sui problemi più gravi. Come tutte le produzioni di
videogiochi, ha subito dei ``traumi''. Il più grande: \textbf{il 100\% del
team originario è stato sostituito}. Non per scelta, ma perché i membri non si
sono iscritti al terzo anno: avevano altri obiettivi; c'è chi cerca lavoro e
chi ha aperto un'azienda. Cristian, volendo continuare il terzo anno e
soprattutto \emph{L'Inf}, un progetto in cui crede e che ha dimostrato il suo
potenziale a Level Up e con i feedback costanti dei giocatori, ha dovuto fare
dei \textbf{colloqui} per formare un nuovo team che si adattasse meglio al tipo
di sviluppo del gioco\footnote{La frase è in parte persa nella trascrizione
(``ho dovuto fare dei colloqui Grazie a tutti. Cioè, servisse una parola
brutta, ma semplicemente si adattasse di più al tipo di sviluppo''): il
senso ricostruito è che ``colloqui'' è una parola forte, ma che si trattava
di scegliere persone adatte a questa fase.}.

La composizione è cambiata così:
\begin{center}
  \begin{tabular}{@{}lp{0.55\textwidth}@{}}
    \toprule
    \textbf{Team} & \textbf{Composizione} \\
    \midrule
    originario & 2 game designer, 2 level designer, 1 technical designer \\
    attuale    & 3 game e level designer, 2 technical designer \\
    \bottomrule
  \end{tabular}
\end{center}
Il motivo: la preproduzione, cioè la parte più importante di game design, era
finita, e servivano \textbf{più tecnici}. Il gioco si basa molto sulle
meccaniche: un sistema di combattimento innovativo che richiede di scriptare
le meccaniche, diversi tipi di sparo, puzzle, una boss fight. È molto più
basato sulle meccaniche che su narrativa o ambientazione. Allo stesso tempo
servivano level designer, perché il gioco doveva prendere forma e i livelli
andavano costruiti.

\subsection{Perché raccontare tutto questo}

Non per vendere il gioco (``anche se magari è un effetto collaterale ben
voluto''), ma per mostrare che \textbf{ogni progetto, a seconda di come è
composto} sul piano del design, della produzione e soprattutto del team,
\textbf{ha un modo diverso di applicare Scrum}. Per questo Cristian è sceso nel
dettaglio delle caratteristiche: ognuna comporta ambiti di produzione diversi.
\begin{itemize}
  \item I livelli richiedono level designer.
  \item Il \emph{match-3} e i diversi tipi di combattimento richiedono
        tecnici.
  \item Il single player è più un vantaggio che un limite: un gioco
        multiplayer non è il doppio, ma ``tipo dieci volte più difficile'' da
        gestire a livello produttivo.
  \item Il genere \emph{fantasy adventure} deriva soprattutto dalle ricerche.
\end{itemize}
Quindi \textbf{ogni progetto ha il suo framework di Scrum}. E il fatto di
aver cambiato team a metà produzione, finita la \emph{vertical slice}, con
persone completamente nuove (ex compagni di studio, che però non avevano mai
visto come funzionava il gioco), ha avuto diverse conseguenze, che Cristian
rimanda alle domande di fine lezione.

% ---------------------------------------------------------------------
\section{Il loop di Scrum}
\label{sec:z3-loop}

L'applicazione di Scrum è un \textbf{loop} molto semplice, simile a quello
mostrato da Simone:
\begin{enumerate}
  \item \textbf{creazione del backlog};
  \item \textbf{selezione dei task} da svolgere, a ogni sprint;
  \item \textbf{implementazione dei task} durante la sprint, con controlli e
        \emph{scrum} (con la frequenza che si decide);
  \item \textbf{calcolo della velocity}, per vedere com'è andata e come
        migliorare la produzione nelle sprint successive;
  \item \textbf{interpretazione} della velocity per fare il \emph{tuning}
        della produzione; e si ricomincia.
\end{enumerate}
Cristian le affronta una per una, con esempi della sua esperienza e consigli.
È contento che la classe stia cominciando proprio adesso, perché questa parte
potrebbe tornare utile già dalla settimana successiva, quando gli studenti
applicheranno il \emph{framework}.

% ---------------------------------------------------------------------
\section{La stesura del backlog}
\label{sec:z3-backlog}

\subsection{Un backlog per milestone}

Anche Cristian, all'inizio, cercava di prevedere il più possibile tutto il
gioco, dal prototipo fino al rilascio. Oggi invece lo struttura
\textbf{per milestone}: il backlog dell'\emph{alpha}, poi quello della
\emph{beta}, poi quello della \emph{gold}. Questo non impedisce di cominciare
a prevedere che cosa ci sarà in beta e in gold, ma dà un'idea molto chiara di
dove si sta andando, soprattutto per capire che cosa tagliare o aggiungere
quando ci si accorge che qualcosa serve prima o dopo, e adattarsi di
conseguenza: ``la cosa più importante del mindset agile è adattarsi ed essere
flessibili''.

\subsection{Da dove vengono i task}

Il gioco deve pur nascere da qualche parte: i task vengono dai
\textbf{documenti chiave} della produzione, già visti nel corso
(capitolo~\ref{cap:produzione-documentazione}):
\begin{itemize}
  \item il \textbf{documento di visione};
  \item i \textbf{GDD} delle varie feature, che dettagliano la visione;
  \item il \textbf{versioning} delle feature, a cui si dovrebbe pensare già
        mentre si scrive il GDD;
  \item gli \textbf{LDD} (\emph{level design document}) e tutti gli altri.
\end{itemize}
L'ordine dipende dal gioco: in un gioco fortemente basato sulle meccaniche è
probabile che gli LDD arrivino dopo lo sviluppo del \emph{core}; in un gioco
basato sull'esplorazione non succede proprio l'inverso, perché le meccaniche
servono comunque, ma le due cose avvengono quasi in contemporanea.

Una volta che si ha la documentazione con tutte le descrizioni e le
caratteristiche delle feature, \textbf{dalla versione 0 alla versione X}, si
ha già una serie di implementazioni da realizzare. Un esempio:
\begin{itemize}
  \item \textbf{versione 1}: il giocatore deve camminare; si usa il
        manichino base di Unreal;
  \item \textbf{versione 2}: si implementa la \emph{skeletal mesh} fatta dal
        3D artist e si adatta la camera al movimento, con gli effetti di
        vibrazione a seconda che il personaggio corra, cammini o sia
        abbassato;
  \item \textbf{versioni successive}: quando il personaggio si muove su
        terreni diversi solleva tipi diversi di polvere, di più su un terreno
        molto polveroso. Non è fondamentale e potrebbe essere una delle
        ultime versioni, verso la fine.
\end{itemize}
Il consiglio è che, quando si fa il \emph{versioning}, si inseriscano
\textbf{tutti gli ambiti necessari} per quella feature: programmazione, arte
(se serve un asset), animazione, UI, VFX, suoni. Se tutti i \emph{field}
compaiono già nella documentazione, la creazione del backlog sarà più chiara
e più spedita.

\subsection{Come organizzare i task}

I giochi hanno moltissimi task. Come organizzarli perché siano leggibili non
solo da chi li ha scritti (scrum master o project manager), ma anche da chi
deve navigarli per capirci qualcosa? Cristian usa tre tipi principali di
organizzazione.

\paragraph{1. Per feature (o epic).} Si suddividono i task per feature,
collegandosi alle \emph{epic} di Simone. Ad esempio, l'epic è il movimento
del personaggio e i task sono camminata, corsa, salto, volo, ``dipende dal
gioco''.

\paragraph{Tutto in un solo posto.} La suddivisione non va fatta per forza in
documenti diversi, anzi Cristian lo sconsiglia: se tutto è nello stesso
strumento (Google Sheets, Monday, ClickUp, Trello) si usano i
\textbf{filtri} e si trova tutto lì, senza navigare tra mille siti o mille
documenti.

\paragraph{Per ambito, con un color code.} Quando si suddivide per ambito
(programmazione, arte, UI, eccetera), conviene usare un \textbf{codice
colore}. Cristian lo usa ``per qualsiasi cosa'': uno per le milestone, uno
per gli ambiti di produzione. Diego gli chiede di mostrare il suo:
\begin{center}
  \begin{tabular}{@{}lll@{}}
    \toprule
    \textbf{Ambito o milestone} & \textbf{Colore} & \textbf{Perché} \\
    \midrule
    Suoni        & viola        & --- \\
    VFX          & rosso        & sta imparando a usare Niagara \\
    Level design & verde        & ``erano finiti i colori'' \\
    Game design  & arancione    & è il suo colore preferito \\
    \midrule
    Alpha        & azzurro      & --- \\
    Beta         & viola        & --- \\
    Gold         & giallo scuro & il colore richiama il nome \\
    \bottomrule
  \end{tabular}
\end{center}
Come si vede, è un misto di funzionalità e preferenze: la cosa più importante
è che \textbf{tutto il team sia d'accordo}.

\paragraph{2. Per milestone, con le sotto-milestone.} Un secondo metodo è
suddividere il backlog per milestone. Ma se il backlog dell'alpha è separato
da quelli di beta e gold, come si fa a dividere in alpha, beta e gold un
backlog in cui tutto è alpha? Cristian suddivide le milestone in
\textbf{sotto-milestone}: l'alpha dell'alpha, la beta dell'alpha, la gold
dell'alpha. Serve ad avere idee più chiare, a tenere sempre in mente lo
\emph{scope} senza perderlo e soprattutto a capire meglio \textbf{se il gioco
è fattibile}. Una suddivisione più fine permette anche di calcolare con più
precisione (una precisione che ``lascia il tempo che trova'') quante sprint
servono per ogni sotto-milestone. La dissonanza tra ``dividete tutto in alpha,
beta e gold'' e ``dentro l'alpha trovate alpha, beta e gold'' è quindi voluta:
si parla delle sotto-versioni.

\paragraph{3. Per priorità.} Chi non vuole usare le sotto-milestone può usare
il sistema di \textbf{priorità} di cui parlava Simone: una scala numerica, da
1 a 5 o da 1 a 10, a seconda di come si riesce meglio a distinguere le
priorità. Serve soprattutto a capire quali task selezionare
(sezione~\ref{sec:z3-selezione}). In realtà è più o meno la stessa cosa con un
approccio diverso: indicare alpha, beta e gold dell'alpha crea già di per sé
delle priorità. ``Dipende tutto da voi.''

\subsection{Come scrivere i task e dove}

Il consiglio è scriverli \textbf{dal punto di vista dell'utente}. Ma, come
diceva Simone, non sempre ``il gioco vale la candela''. Lo è soprattutto in
un caso: se il programmatore legge un task dell'artista 3D scritto in modo
troppo tecnico e non lo capisce, si crea dissonanza tra i reparti e possono
nascere attriti. Quindi è sempre consigliabile, ma non obbligatorio: spesso
basta una chiamata (``prima semplicemente ti giravi verso il tuo compagno e
gli chiedevi che cosa significa''). ``Però è giusto che vi dica le cose come
dovrebbero essere fatte.''

Dove scriverli? In qualsiasi strumento: Trello, Monday, ClickUp, gratuiti o a
pagamento, oppure un foglio Google costruito da sé. Dall'esperienza di
Cristian emerge che \textbf{Google Sheets sta antipatico alle persone}: già il
nome ``foglio di calcolo'' spaventa. Strumenti come Monday o ClickUp hanno più
colori e soprattutto \textbf{mandano le notifiche}, quindi sono più accessibili
a tutti. ``Decidete voi.''

\subsection{Non prevedere tutto: la flessibilità}

È fondamentale capire che \textbf{prevedere tutti i task passo per passo}, in
un ordine preciso (``dopo la camminata faremo il salto, dopo il salto il
sistema di sparo, e così via all'infinito''), va contro i principi di Agile,
in particolare contro il \textbf{quarto valore}: rispondere ai cambiamenti più
che seguire un piano.

Da qui il consiglio centrale: la \textbf{flessibilità} deve essere il cuore
della propria filosofia di lavoro e del proprio modo di applicare Scrum.
Bisogna essere sempre pronti a \textbf{tagliare} nel modo più agile possibile:
se si è capita davvero la mentalità agile, si riesce a tagliare una feature
\textbf{senza bloccare la produzione delle altre}.

L'esempio più eclatante viene dall'inizio di \emph{L'Inf}: Cristian faceva il
backlog di tutto il gioco, arrivando a circa \textbf{400 task} di ogni ambito,
nel dettaglio: per la camminata arrivava a indicare come si muoveva la camera
quando il personaggio girava a destra, e come cambiava quando saltava o cadeva
da un dirupo. ``Vi viene subito da pensare: ma come fai a prevedere tutto?''
Infatti la soluzione è \textbf{non prevedere tutto}: capire a grandi linee ciò
che si vuole, qual è l'obiettivo della feature, senza scendere troppo nel
dettaglio, perché il dettaglio crea di per sé un \textbf{lavoro Waterfall}, che
è ciò che bisogna assolutamente evitare. E scrivere un backlog così dettagliato
richiede tantissimo tempo, che è limitato: la produzione ne risente in tutti i
modi possibili.

\paragraph{Bonus: gli elementi innovativi.} Se il gioco (il proprio
\emph{Bigger One}) ha elementi innovativi, la flessibilità va portata
\textbf{ai massimi livelli}. \emph{L'Inf} era partito con una meccanica di
\textbf{assorbimento e sparo}: si assorbe un nemico colorato, ad esempio
giallo, e lo si spara contro altri nemici gialli, prendendo ispirazione da
altri giochi. Nel loro caso, però, creava più problemi che altro: per due o
tre ore di gioco, ``assorbire, sparare, assorbire, sparare'' risultava
\textbf{troppo ripetitivo}, e hanno deciso di tagliare la meccanica. Il
problema è che l'hanno tagliata dopo che Cristian aveva già scritto
\textbf{600 task nel backlog e 12 GDD}: tempo perso. Per questo la
flessibilità, negli elementi innovativi ma in generale, ``è la prima cosa''.

\subsection{Quanto dettaglio?}

Quando un backlog è abbastanza dettagliato e quando lo è troppo poco?
``La risposta purtroppo non c'è.'' Viene da due fattori:
\begin{enumerate}
  \item l'\textbf{esperienza}: man mano che si scrivono GDD e si va avanti
        nello sviluppo, ci si ``tuna'';
  \item il \textbf{trial and error}: si sbaglia e si torna indietro, il che
        riporta, di nuovo, all'esperienza.
\end{enumerate}

% ---------------------------------------------------------------------
\section{La selezione dei task}
\label{sec:z3-selezione}

Con un bellissimo backlog ordinato per colori, ambiti e priorità, quali task
si selezionano? ``Da qualche parte bisogna pur cominciare.'' Le domande
principali sono due: \textbf{quali} selezionare e \textbf{quanti}.

\subsection{Quali task}

\paragraph{La priorità.} Il primo criterio è la priorità: i task da fare
prima sono le scelte più urgenti, e quindi si scelgono quelli\footnote{L'inizio
di questo passaggio manca nella trascrizione (``quanti selezionare. cui
scegliere prima, sarà la scelta più urgente da fare''): il primo criterio è
stato ricostruito dal contesto, anche perché poco dopo Cristian dice che,
risolto il problema delle competenze, ``si ritorna a selezionarli per
priorità''.}.

\paragraph{Le capacità attuali degli sviluppatori.} Il secondo criterio sono
le \textbf{abilità dei developer}. Cristian precisa che non intende in alcun
modo giudicarne la professionalità: intende letteralmente \textbf{che cosa sa
fare lo sviluppatore in quel momento}. Se un programmatore non ha mai usato
Unreal, non è il momento di dargli task di implementazione in Unreal: bisogna
prevedere anche il \textbf{periodo di apprendimento}, eventualmente come task
vero e proprio (``hai questa sprint per imparare a implementare l'IA'').

È successo due mesi prima con uno dei suoi technical designer: gli ha dato
\textbf{due settimane} dicendogli ``in queste due settimane non ti preoccupare
di niente, il resto lo produciamo noi''. Doveva essere d'accordo, e lo era,
anzi era contento di approfondire l'argomento: studiare come funzionano le IA
(anche grazie a vari link passati da Diego) e soprattutto come implementarle
nel motore. Per due settimane ha dovuto ``solamente'' pensare
all'apprendimento. Il consiglio: \textbf{prevedere questi ``tempi morti''},
perché ne ha bisogno il team, e ne ha bisogno chi pianifica per non falsare la
previsione dello \emph{scope}. Se in questo momento (``sottolineo, in questo
momento'') un programmatore non sa fare una cosa, la si rimanda e si mette
altro; prima o poi imparerà, e a quel punto si torna a selezionare per
priorità.

\subsection{Quanti task}

Il primo criterio sulla quantità è la \textbf{velocity}, ma Cristian l'ha
messo nella slide solo per completezza del loop: se non si ha ancora una
velocity, non ha senso. Per ora la si trascura e si fa finta che sia la
\textbf{prima sprint}, la prima applicazione di Scrum.

Nelle prime sprint la quantità si basa sulle \textbf{capacità} e sulla
\textbf{disponibilità} degli sviluppatori:
\begin{itemize}
  \item \textbf{disponibilità}: per quanta parte della sprint il technical
        designer (per riprendere l'esempio) è disponibile? Non in termini di
        ore, perché Cristian concorda con Simone che le ore sono sempre molto
        relative;
  \item \textbf{performance}: gli studenti si conoscono ormai da quasi un
        anno e mezzo. Se uno sviluppatore non ha mai fatto un livello in
        Unreal (per motivi suoi, magari problemi al PC) e dice ``voglio fare
        un livello e ci metto due settimane'', sta a chi pianifica cominciare
        a capire \textbf{chi può fare che cosa}. Tutti possono fare tutto, se
        si prevede il tempo di apprendimento; ma se Cristian arrivasse a dire
        ``in due settimane vi faccio l'IA del vostro gioco implementandola in
        Unreal'', ``vi sto prendendo in giro''.
\end{itemize}
Queste sono le analisi preliminari quando non si ha una velocity. Se invece
non si conosce affatto il team, l'unico punto di partenza è la
\textbf{disponibilità} degli sviluppatori, da cui si stima quanto si riesce a
fare.

Tanto, \textbf{nelle prime sprint si sbaglieranno completamente tutte le
previsioni}, ``ed è giusto così, e ben venga'': se non si sbaglia,
probabilmente si sta applicando Waterfall, o si sta applicando male Scrum.
All'inizio si tratta solo di buttarsi: ``questa settimana abbiamo previsto 15
task; ne abbiamo fatti 40. Bene, la prossima ne proviamo 40''; poi magari se ne
fanno 25, e così via, \textbf{tarando} man mano la quantità.

\subsection{Definire bene il goal}

Per ogni sprint, e anche per ogni singolo task, bisogna \textbf{definire bene
l'obiettivo}. Serve a mantenere concentrati gli sviluppatori: se so che entro
questa settimana devo finire il block out del livello, so dove sto andando,
lavoro per quel risultato e soprattutto \textbf{non perdo tempo}, come succede
spessissimo. Se durante il block out qualcuno si mette a inserire effetti
visivi ambientali o il \emph{post-process}, viene il dubbio che non stia
lavorando per l'obiettivo: è un \textbf{campanello d'allarme}, perché potrebbe
far mancare l'obiettivo per cui la sprint era stata pensata.

\subsection{Bonus: progettare il ritmo dello sviluppo}

Un consiglio che Cristian ha sentito raramente, ma che ha deciso di
sperimentare: poiché gli studenti sono designer, e tendenzialmente anche lo
scrum master del loro progetto lo sarà, possono \textbf{progettare il
\emph{pace} dello sviluppo}.

Di solito si pensa (e potrebbe anche essere giusto) che l'obiettivo di Scrum
sia trovare una \emph{velocity target} e mantenerla costante per tutta la
produzione. Può essere una soluzione; ma con entrambi i suoi team Cristian ha
notato che, per come sono fatti gli esseri umani, ci sono \textbf{picchi di
attenzione e di produzione}. Con l'esperienza si impara a percepirli in
anticipo e a progettare di conseguenza; oppure sono gli sviluppatori stessi a
dirlo: se la settimana prima hanno implementato 12 nemici, 15 modalità di sparo
e 13 cutscene, quella dopo diranno ``andiamoci piano''. Di solito non succede,
perché le sprint sono programmate apposta per adattarsi alle capacità degli
sviluppatori, ma è un metodo possibile.

Che cosa intende, in concreto? Se per un periodo si sono implementate molte
feature o sistemi che però, visivamente, \textbf{non hanno portato un aumento
di qualità percepibile nel gioco}, gli sviluppatori non si sentono
riconosciuti e ripagati dello sforzo. Allora, la settimana dopo, ci si
concentra su qualcosa che dà \textbf{soddisfazione}. Questo lo permettono le
sprint, lo permette Scrum e non Waterfall: è una differenza fondamentale.

\paragraph{L'esempio di \emph{L'Inf}.} Per più di otto mesi il gioco è rimasto
con \textbf{quattro tipi di nemici}. Nonostante l'impegno (quattro nuove
modalità di sparo, la boss fight, tantissima programmazione ``di
background'', cioè non visibile al giocatore ma necessaria al funzionamento
del gioco), il team cominciava a percepire che \textbf{non c'era corrispondenza
tra impegno e qualità percepita}. Cristian ha allora deciso di limitare per un
periodo la produzione degli elementi di background (programmazione come l'IA,
tutto ciò che faceva proseguire il gioco ``sotto''), di inserire altri tipi di
nemici, che tutti non vedevano l'ora di avere (``sapevo che era uno di quei
jolly da sfruttare per far risalire il morale del team''), e di dedicare tempo
a effetti visivi e sonori, a migliorare tutto ciò che era estetico. Diego
suggerisce il termine: la \textbf{juiciness}. Ciò che vedeva il giocatore lo
vedevano anche gli sviluppatori, che si sentivano più appagati e sentivano i
propri sforzi ricompensati. ``Adesso loro non lo sanno, ma li ributterò sotto
un treno di lavoro'' (``in realtà lo sanno, però non me lo dicono'').

Per Cristian è un punto fondamentale, che gli è costato ``mesi di
sofferenze'' di esperienza per capirlo. Diego invita la classe a fare domande
su un \emph{insight} ottenuto per tentativi ed errori.

% ---------------------------------------------------------------------
\section{L'implementazione: controlli, blocker e feedback}
\label{sec:z3-implementazione}

\subsection{Il check dei progressi}

Selezionati i task, viene la parte ``più tranquilla'' dal punto di vista di
Scrum (i problemi di implementazione, ovviamente, sono altri): parte la
settimana di lavoro, ognuno ha i suoi task, i \emph{task owner} hanno scelto
che cosa fare, e questo prosegue per tutta la durata della sprint.

Durante l'implementazione si fa un \textbf{check dei progressi}. Non serve a
controllare che cosa stanno facendo le persone: di base bisogna
\textbf{fidarsi} che, se il programmatore dice ``farò questo'', lo farà entro i
limiti della sprint. Cristian parla qui di \emph{daily scrum} perché è ciò che
si dovrebbe fare in generale, anche se il suo team non li fa ogni giorno
(sezione~\ref{sec:z3-team}). Oltre alle solite domande (che cosa hai fatto
ieri, che cosa farai oggi), \textbf{la cosa più importante è trovare i
blocker}.

\subsection{Agire subito sui blocker}

Se il programmatore dice ``non ho potuto implementare il personaggio perché mi
manca l'asset, la \emph{skeletal mesh} del protagonista'', bisogna
\textbf{agire subito}, in due modi:
\begin{enumerate}
  \item \textbf{dare altri task} al programmatore, perché altrimenti resta lì
        ad aspettare, e questo non lo si vuole assolutamente;
  \item \textbf{capire che cosa è successo} alla skeletal mesh, parlando con
        la 3D artist. Se è solo questione di tempo, non è davvero un problema:
        è normalissimo aver assegnato l'implementazione prima che l'asset
        fosse pronto (``è normale sbagliare, l'importante è rimediare''). Se
        invece anche lì c'è stato un blocco, bisogna \textbf{risalire a
        cascata} fino a trovarlo.
\end{enumerate}
Magari si scopre che il problema sono proprio i game designer: quando hanno
fatto la \emph{mood board} non sono stati abbastanza chiari, e la 3D artist ha
dovuto spendere una settimana a cercare reference per conto suo, mentre il
designer pensava che quelle fornite bastassero e non si è posto il problema
della fase di ricerca delle reference da parte di chi fa la grafica o il 3D
(Cristian parla al femminile perché la 3D artist del suo team è una donna, ma
il discorso è generale). Anche questo, come l'apprendimento, è uno di quei
task che non bisogna trascurare.

\subsection{Analizzare i blocker prima di risolverli}

Il consiglio è \textbf{analizzare i blocker}, non risolverli appena si
presentano, perché si rischia di trascurare il problema alla radice. Cristian
richiama una delle prime lezioni del primo anno, quella su \textbf{bisogni e
problemi}: bisogna capire se il blocker è un problema o un bisogno, e agire di
conseguenza. Il guaio nasce quando lo si trascura o, peggio, quando lo si
risolve in fretta senza averlo considerato: ``poi vi si ritorcerà contro''.

\paragraph{L'esempio dei muri invisibili.} \emph{L'Inf} è un gioco di
esplorazione ambientale, per lo più all'aperto: su cinque livelli, circa
l'80\% del gioco è all'esterno. Per dare \emph{meaning} al livello finale,
quello della boss fight e della rivelazione narrativa su dove siano finiti
tutti gli esseri umani del pianeta, il team ha deciso di sottolinearne la
diversità rendendolo \textbf{al chiuso}.

La camera è in terza persona e abbastanza lontana dal giocatore, non classica.
Dare le collisioni alla camera sembrava fastidioso, perché a ogni rotazione
``sbatte'' contro i muri. La prima idea: \textbf{rendere invisibili i muri}
quando la camera ci passa dietro. Non si potevano far sparire tutti insieme
quando il giocatore era solo dietro il muro destro, quindi ogni muro aveva il
suo \emph{trigger}, che stabiliva quante e quali pareti dovessero sparire. Il
livello dura mezz'ora, con una buona decina di minuti di esplorazione al
chiuso, completamente circondata da pareti: \textbf{un mese di lavoro} solo
per assegnare i trigger, cambiare i materiali e gestire il tutto. ``Era l'idea
geniale che ci era venuta in quel momento.''

Quando hanno consegnato la beta e l'hanno fatta provare ai professori e ad altri
giocatori, tutti hanno detto più o meno la stessa cosa: i muri che spariscono
\textbf{non sono eleganti e rompono l'immersione}. Hanno tagliato tutto e
rimesso le collisioni alla camera. Se avessero analizzato bene il problema
prima, avrebbero capito che bastava \textbf{gestire meglio gli spazi}: non fare
curve strette ma allargarle, così che anche ruotando la camera il giocatore
tendenzialmente non sbatta contro il muro. Era ``assolutamente una questione di
riallargare più o meno tutti gli spazi''; ed è ciò che hanno fatto alla fine,
insieme alle collisioni. ``Ad averlo fatto prima avremmo risparmiato due
settimane di tempo''\footnote{Poco prima Cristian parla di ``un mese buono'' di
lavoro sui trigger (correggendosi da ``due mesi''); le ``due settimane'' sono
forse il risparmio netto stimato, tenendo conto del lavoro di allargamento
degli spazi che è stato comunque necessario. La trascrizione non permette di
chiarirlo.}. Ecco perché è importante analizzare i blocker.

\subsection{Bonus: i feedback costanti}

Cristian nota che, soprattutto ultimamente, in una fase di
\emph{lockdown}\footnote{Il riferimento al lockdown colloca la lezione nel
periodo della pandemia di Covid-19, coerentemente con l'accenno di Simone al
Covid e alle sedi chiuse (``quando si poteva andare'').} in cui tutti sono
costantemente in contatto su Discord, ``anche più di quanto lo saremmo nella
vita vera'', c'è un \textbf{bisogno costante di dare feedback} sul lavoro.

Lui è pro feedback in qualsiasi momento, anzi ``ossessionato'': all'inizio ha
chiesto alla classe di criticarlo e di dargli consigli sulle slide, perché è
l'unico modo per migliorare. Il problema è che \textbf{i feedback costanti
producono un calo di produttività enorme}. L'esempio è personale: Cristian è
anche il VFX artist del team, e sugli elementi estetici c'è sempre modo di
migliorare, quindi riceve di continuo feedback (``io lo farei più grande, rosso,
verde, giallo, a forma di triangolo, a forma di quadrato''). Sono feedback
correttissimi: chi li dà, soprattutto in produzione, non lo fa perché pensa
che tu sia scarso, ma per il bene del gioco, che vive e vuole ben prodotto. I
migliori mandano anche video di riferimento (``ho visto questo VFX in questo
gioco, sapresti rifarlo?''), e questo ben venga. Il problema è che siano
\textbf{costanti}.

Il consiglio: \textbf{limitare la richiesta e l'ascolto dei feedback a fasi
specifiche} (alla fine della sprint, oppure in un controllo intermedio). Dipende
anche da che cosa si sta facendo: con i VFX molte cose dette oggi sono
sistemate il giorno dopo; ma se chi costruisce un livello, in block out o in
\emph{environment}, riceve mille feedback per ogni metro quadro (``un po' più a
destra, un po' più a sinistra''), il lavoro si ferma lì. Sia perché chi lo fa
si \textbf{demoralizza}, sentendosi costantemente criticato, sia per una
questione di fiducia, tornando al discorso di Simone: ``fidati di me come io mi
fido di te. Fammi completare il task; poi, quando ti chiedo il feedback, o più
avanti, è ben accetto, anzi ben venga''.

Cristian si scusa per l'apparente presunzione di definire ``fondamentali'' i
propri consigli: lo fa perché, se li avesse avuti prima, si sarebbe risparmiato
tantissimi problemi. Diego scherza che andranno nelle slide ``con la mia faccia
super inquietante''; Cristian dice che può passare anche le sue slide.

% ---------------------------------------------------------------------
\section{Il calcolo della velocity}
\label{sec:z3-calcolo}

\subsection{Confrontare il fatto con lo stimato}

Sulla parte pratica Simone è stato più esaustivo, quindi Cristian si concentra
sull'interpretazione\footnote{Un breve tratto della spiegazione iniziale è
perso nella trascrizione, che riprende a metà frase (``nel comparare quello
che è stato fatto da quello che è stato stimato'').}. La velocity consiste nel
\textbf{confrontare ciò che è stato fatto con ciò che era stato stimato}, in
numero di task o in pesi, a seconda di come la si gestisce. Da lì si capisce
che cosa è successo e che cosa si può prevedere:
\begin{itemize}
  \item \textbf{fatto $>$ stimato}: si può prevedere di più; nella sprint
        successiva si aumentano i pesi o il numero di task;
  \item \textbf{fatto $<$ stimato}: o si è esagerato con la stima (succede
        sempre, così come succede sempre il caso opposto), oppure ci sono stati
        problemi gravi o \emph{blocker} nella produzione, ed è fondamentale
        capire quale delle due cose sia successa;
  \item \textbf{fatto $=$ stimato}: succede quasi mai, ma è l'obiettivo. Le
        possibilità sono due: si è trovato il proprio \textbf{bilanciamento}, e
        si può continuare così; oppure \textbf{non si sta applicando bene
        Scrum}, ad esempio perché si calcola la velocity su ciò che si è fatto
        invece che sulla stima.
\end{itemize}

\subsection{Più di un numero}

L'obiettivo non è arrivare a una velocity pari a 1, intesa come rapporto tra
fatto e stimato (100 punti stimati, 100 fatti). La velocity è \textbf{molto più
di un numero}: dà almeno tre informazioni.

\paragraph{1. La bontà degli strumenti di stima.} Planning poker, Fibonacci e
tutti gli strumenti di stima sono utilissimi, ma \textbf{vanno testati}: la
velocity mostra se ci sono metodi sbagliati di stima dei pesi o del numero dei
task.

\paragraph{2. La raggiungibilità dello scope.} Se per tutta la produzione la
velocity resta, ad esempio, a 0,8, sempre al di sotto di quanto dovrebbe essere
(in teoria non dovrebbe succedere, perché ci si dovrebbe tarare), e si sono
previsti dieci livelli, si capisce già da subito che \textbf{non ci si
arriverà mai}, a meno di aggiungere uno o due anni di sviluppo.

\paragraph{L'esempio degli otto livelli.} Finita la \emph{vertical slice}, il
team ha ricontrollato tutta la preproduzione e impostato il gioco come un gioco
completo, non come la creazione di una parte: c'è tutta la storia, ci sono
tutti i livelli, tutto ciò che fa il giocatore dall'inizio alla fine, chi
incontra e quando; tutto scritto, ma tarato e lasciato \textbf{flessibile},
``altrimenti ricadiamo nella trappola del Waterfall''.

A settembre, all'inizio dell'anno (``per farvi capire quanto il processo è
rapido''), avevano previsto \textbf{otto livelli}. Sembrava una stima
credibile: il vecchio team, con due level designer che volevano fare
esattamente quel lavoro, produceva un livello in \textbf{due o tre settimane}, a
una qualità abbastanza elevata: magari mancavano l'illuminazione, il
\emph{foliage} e alcuni dettagli, ma il livello era \textbf{giocabile}, che è
la parte fondamentale. Cristian, unico rimasto del vecchio team, si è basato su
quella stima: ``da qualche parte dovevo cominciare''.

Ma i level designer del nuovo team \textbf{non sono level designer puri}: sono
soprattutto game designer, interessati al game design, che hanno deciso di
sperimentare con il level design per necessità (il game design serve sempre, ma
se manca un level designer, visto che il percorso di studi insegna sia level
sia game design, ci si adatta). C'è chi si è adattato a fare il level designer,
e chi ha deciso proprio quell'anno di specializzarsi nel level design: ``perché
avrei dovuto dirgli no, tu non sei un level designer? Anzi, ben venga''.
L'unica cosa è stata prevedere che avesse più bisogno di apprendimento: magari
sbaglierà di più, dovrà informarsi su come si fanno i livelli, avrà bisogno di
più feedback. L'importante è capirlo all'inizio.

Guardando l'andamento del primo livello, Cristian vedeva che \textbf{in un mese
avevano a malapena raggiunto il block out}. Ha convocato una \textbf{riunione
di emergenza}, lunghissima (non una revisione della sprint da venticinque
minuti come quelle di Simone, ma un altro tipo di riunione). Il suo obiettivo
non era tanto convincerli che lo \emph{scope} era troppo alto (ci è riuscito),
quanto \textbf{assicurarsi che non la prendessero come una sconfitta
personale}. È stato difficile ma interessantissimo: non se la sono presa né con
se stessi, per non riuscire a performare, né con lui, che stava di fatto
dicendo ``non ce la farete''. Anzi, nonostante il taglio, sono stati
\textbf{loro stessi a proporre la soluzione}: ``ok, abbiamo capito; facciamo
cinque livelli, e se vediamo che arriviamo (speriamo a fine marzo) con tutti i
livelli, poi possiamo decidere di aumentarli o di cambiare parti del gioco''.
Cristian l'aveva già in mente, ma il fatto che l'abbiano proposta loro l'ha reso
``veramente un successo''.

Tutto questo grazie a un numero. Se Cristian fosse arrivato tra un mese e
mezzo, dopo che avevano già fatto quattro livelli ``con l'acqua alla gola'',
dicendo ``ne dobbiamo fare altri quattro entro marzo'', e poi a fine marzo
``basta, non ce la facciamo'', sarebbe stato molto più grave\footnote{Il
passaggio è confuso nella trascrizione (``mi ha portato ad affrontare queste
situazioni che affrontate a fine marzo'', ``non potne quattro''): il senso
ricostruito è che la velocity ha permesso di anticipare a inizio produzione una
decisione che, presa a fine marzo, sarebbe stata molto più dolorosa.}.
``Ringraziamo in questo caso la velocity.''

\paragraph{3. Lo stato del team.} La velocity permette di capire lo
\textbf{status del team e le sue performance}, perché si basa su una stima fatta
a partire dalla settimana precedente. Si ricollega al ``progettare il \emph{pace}
dello sviluppo'' (sezione~\ref{sec:z3-selezione}): ci si può basare proprio sulla
velocity per capire qual è il ritmo attuale e prevedere i prossimi. ``Tutto
torna.''

\subsection{Bonus: mai tornare indietro}

Un bonus che ``più che un bonus è un'altra lezione'': \textbf{prevedete sempre
task che portano a un incremento del prodotto}. La produzione può rallentare,
perché possono esserci difficoltà, ma \textbf{mai, mai, mai andare indietro}.

È un'altra lezione insegnata da \emph{L'Inf}, legata alla stesura del backlog.
Quando Cristian aveva scritto tutto il backlog (``400 task: come vedete ogni
tanto dico 400, 600, 300, giusto per dire che erano veramente tanti'') e si
accorgeva di aver sbagliato, invece di adattarsi e appoggiarsi alla produzione
in corso \textbf{fermava la produzione} (``ditemi pure che ero folle: lo ero'')
e riscriveva da capo tutti i task, arrivando magari a mille o riducendoli.

Lo stesso vale per i \textbf{documenti}. ``La documentazione è bellissima'':
Cristian era un grande fan della documentazione (Diego può confermarlo), ma
lavorando con Scrum, per ragioni di produttività, ha dovuto \textbf{scendere a
compromessi}: darsi un limite su quanto un documento è bello e ricco di
dettagli. Il motivo è sempre l'esempio della prima meccanica, l'assorbimento e
sparo: tutti i GDD del gioco si basavano su di essa, con tutte le interazioni,
anche secondarie (che cosa si poteva assorbire? Nemici, oggetti, chiavi, cose
per recuperare vita, bonus, malus, qualsiasi cosa). Tagliata la feature, tutti i
GDD sono stati ``letteralmente presi e buttati'', e Cristian ha deciso di
rifarli tutti da capo.

Questo approccio potrebbe funzionare in Waterfall (``finiti i documenti,
iniziamo la produzione''), ma l'obiettivo era \textbf{consegnare un videogioco,
non fare documenti}; e soprattutto il gioco era innovativo. È successo altre due
volte: in tutto Cristian ha riscritto \textbf{backlog e GDD da capo più
volte}\footnote{Nella trascrizione: ``è successo altre due volte quindi per ben
quattro volte io ho riscritto tutto il backlog e tutti i GDD''. Il conteggio non
torna del tutto (una volta più altre due fa tre); forse Cristian conta anche la
prima stesura.}, riportando ogni volta la produzione all'inizio (``ok, possiamo
iniziare il prototipo; no, è sbagliato, ricominciamo''). Lo racconta per far
capire quanto ci ha messo a entrare nel mindset agile, e per tranquillizzare la
classe: non indovinarci subito è normale, ``succederà il 90\% delle volte''.

Quando poi hanno fatto il prototipo, si sono accorti che \textbf{tutta la
documentazione era sbagliata}: con un gioco innovativo non si poteva sapere che
cosa avrebbe funzionato senza prima vederlo, perché non c'erano reference:
nessuno aveva fatto un gioco d'avventura \emph{match-3}. ``Il \emph{match-3}
esiste, i giochi d'avventura esistono, mischiarli insieme è una follia'', e lo
dice ``per criticarmi, non per vantarmi, perché non lo farò più''.

% ---------------------------------------------------------------------
\section{Il tuning della produzione}
\label{sec:z3-tuning}

\subsection{Che cosa e quando}

Il \emph{tuning} della produzione si basa sull'analisi fatta: dove ci si trova,
qual è lo stato delle cose, che cosa si potrebbe fare meglio e stimare
meglio\footnote{Anche qui l'inizio della spiegazione è parzialmente perso nella
trascrizione (``che in parte si basa su ci si trovano, qual è il loro
status'').}; per poi tornare all'inizio del loop con più dati su cui basare una
stima corretta.

Si può fare \textbf{a ogni sprint}, a seconda della durata, oppure \textbf{ogni
serie di sprint}: ``vedetelo voi''. Cristian suggerisce la seconda opzione
perché ritoccare continuamente il backlog per tagliare o aggiungere feature va
contro il principio del lavorare sempre per incrementare il prodotto. Va da sé
che, se si tratta di un taglio molto importante, \textbf{prima lo si prende e
meglio è}: ``ragionateci bene''.

\subsection{Su che cosa basarsi}

\paragraph{Sulle performance, cioè sulla velocity.} Se la velocity è alta, si
può dire: ``a questo punto possiamo reintrodurre qualche feature che era nei
\emph{nice to have}\footnote{Nella trascrizione ``nice to the epoch'': quasi
certamente \emph{nice to have}.} o che era addirittura stata tagliata''. Se è
bassa, si sistemano già milestone e backlog in base al risultato.

\paragraph{Sui feedback.} A Level Up \emph{L'Inf} ha ricevuto tantissimi
feedback, per lo più da giocatori. Trattandosi di una \emph{vertical slice},
spesso dicevano cose che il team già sapeva e che avrebbe implementato nella
sprint successiva (ma, per questioni di tempo, da qualche parte bisogna
tagliare). Altri, invece, hanno parlato di \textbf{accessibilità}.

\paragraph{Il cambio di camera.} A questi si sono aggiunti i feedback costanti
di Diego sulla camera: più che un ``dovete cambiarla'', un ``perché non provate a
cambiare camera?''. Diego interviene: ``c'è voluto del tempo a convincere
Cristian: otto mesi''. Alla fine l'ha fatto e sta funzionando, anche se ha
pesato molto sul team. Il motivo: il \textbf{twin-stick shooter}\footnote{Nella
trascrizione ``twin steep shooter''.} richiede una grande
\textbf{dimestichezza con i controlli} (muoversi e mirare contemporaneamente con
le due levette) e di per sé esclude gran parte del pubblico casual. I giocatori
più casual dicevano proprio questo: ``non riesco a camminare e mirare
contemporaneamente''. Il team se l'è segnato, ne ha discusso e alla fine ha
accettato il cambio di camera, con tutte le conseguenze.

\subsection{Flessibilità anche nella struttura dei documenti}

Un promemoria sulla \textbf{flessibilità}: più si è stati flessibili all'inizio,
più si riesce a rimettere mano a backlog e produzione in seguito. E questo vale
anche, banalmente, per \textbf{come è strutturato un documento}: se ogni casella
fa riferimento a un'altra casella, che fa riferimento a un'altra due righe
sopra, modificare quel documento (o un'implementazione costruita allo stesso
modo) richiederà molto più tempo. Flessibilità in tutto: nella scrittura dei
documenti e nella produzione (``non lo ripeto, altrimenti ricomincio'').

\subsection{Evitare il perfezionismo}

``A proposito di perfetto'': il consiglio sul \emph{tuning} è \textbf{evitare
in tutti i modi il perfezionismo}. Cristian ci sta ancora lavorando
``tantissimo'' (``anche qui Diego lo sa meglio di me''): è migliorato molto e nei
progetti ci riesce, ma in altre cose pretende ancora troppo. Il perfezionismo
\textbf{va contro la flessibilità e crea inevitabilmente dei blocker}. Lo sta
vivendo in prima persona, non per sé ma con altri membri del team.

\paragraph{Il block out con il foliage.} Il task è ``fai il block out di questa
parte del livello'' (dare l'intero block out di livelli da mezz'ora richiede
comunque tempo, quindi il task riguarda la parte iniziale). E nel block out
Cristian si ritrova il \emph{foliage}, il \emph{landscape} scolpito bene, le
rocce. ``Va benissimo perché ti piace il lavoro\dots{} anzi no, sto proprio
dicendo che non va bene.'' L'approccio di voler consegnare il task nel miglior
modo possibile è comprensibile; il problema è che questo perfezionismo
\textbf{blocca tutti gli altri}, per moltissimi motivi. Se si passano due ore a
mettere il foliage e poi, quando si va in gold a rifinire il livello, si trova un
foliage migliore per il \emph{mood} estetico del gioco, quel lavoro va buttato.
Il perfezionismo crea \textbf{buchi di produzione} che alla lunga diventano
letteralmente mesi di lavoro buttato o da buttare: un peccato per chi lo fa e
per chi lavora con lui.

\paragraph{La pressione sugli altri.} Un'altra conseguenza: il perfezionismo
\textbf{mette pressione a tutti gli altri}. Chi lavora alla seconda parte del
block out, con il task di renderla giocabile, vede che la prima parte ha
foliage, rocce e tutti i props (non l'illuminazione, che negli ambienti aperti
è tendenzialmente globale, salvo situazioni specifiche) e si sente in dovere di
raggiungere quel livello di rifinitura\footnote{Nella trascrizione ``si sentirà
in dovere di raggiliere la pozione della vita'': passaggio incomprensibile, il
senso è ricostruito dal contesto.}, trascurando la parte di gameplay, che invece
nella prima parte era stata fatta. Il perfezionismo di uno ha generato
\textbf{ansia} nell'altro, ed è da evitare in ogni modo.

\paragraph{I criteri di accettazione.} È un altro motivo per cui sprint e task
devono avere un \textbf{goal} e un \textbf{criterio di accettabilità}, per
usare i termini di Scrum: il task si considera fatto quando è \emph{done}, e va
bene così, perché lo si è valutato prima. Tutto ciò che si fa oltre il
\emph{done} è \textbf{tempo perso}. Può andare bene che quello che fai adesso
non dovrai farlo dopo, ma è sbagliato come approccio: se il lavoro è stato fatto
bene prima, si può andare avanti tranquilli, e al passo successivo ci si penserà
dopo.

% ---------------------------------------------------------------------
\section{Domanda: come comunicare un taglio al team}
\label{sec:z3-qa-taglio}

Prima della seconda parte, che Cristian promette più breve, una domanda: sul
piano relazionale ed emotivo, come ha gestito il dover convincere, o comunque
spiegare, al team la necessità del taglio da otto a cinque livelli?

\paragraph{Ricordare come funzionano i feedback.} ``Ottima domanda.''
Emotivamente è stato pesante, perché non tutti reagiscono ai feedback allo
stesso modo. Per questo, a ogni riunione, Cristian ricorda i principi dei
feedback (``so che sembra banale, so che non serve''), in particolare una
clausola che tutti dimenticano: \textbf{il feedback si può accettare o
rifiutare}. Se non si è d'accordo, lo si può ignorare; invece tutti partono dal
presupposto che un feedback vada seguito al 100\%\footnote{Nella prima lezione
zero il relatore invitava ad accettare il feedback ``a prescindere''
(sezione~\ref{sec:z1-ricevere}); lì però ``accettare'' significava ascoltarlo
senza difendersi, non necessariamente applicarlo. Le due indicazioni sono
quindi compatibili: si ascolta tutto, poi si decide che cosa farne.}.

\paragraph{Partire dai dati.} Sembra brutto affrontare una questione emotiva
partendo dai dati, ma l'obiettivo era \textbf{far arrivare loro alla soluzione},
non dargliela. Ed è andata così: sono stati loro a dire ``per adesso facciamo
cinque livelli, e poi, \textbf{se ci riusciamo}, implementiamo gli altri''.
Quel ``se ci riusciamo'' era ciò che Cristian voleva sentire: se fosse rimasto
solo ``ne facciamo cinque'', l'avrebbero vissuto come un limite su se stessi
(``non ho fatto in tempo a farne otto''). Il ``se ci riusciamo'' gli ha
comunicato che avevano capito che \textbf{non era una critica al loro modo di
lavorare}, né un modo per dire ``non siete capaci'': lo aveva sottolineato più
volte, magari non con parole così dirette, ma ``un po' \emph{soft-skillato},
come diciamo noi designer''.

Hanno colto che non c'erano alternative (``otto non li facciamo, punto''), ma
anche le possibilità: facciamo cinque livelli, abbiamo altro tempo,
perfezioniamo ciò che abbiamo fatto; vogliamo aggiungere una nuova modalità di
gioco? Benissimo. Altri due livelli? Vediamo, in base al tempo. Esserci
riuscito lo ha reso ``super fiero''.

% ---------------------------------------------------------------------
\section{Il team management}
\label{sec:z3-team}

\subsection{Senza team non esiste Scrum}

Cristian accelera, per stare nei tempi. Il punto da capire è che
\textbf{Scrum si adatta al team}: quindi il team è la cosa più importante di
tutte, più di Scrum stesso. ``Se io mi adatto al team, il team deve esserci.
Senza team non esiste Scrum.''

\subsection{Una produzione senza stress}

Di conseguenza, una produzione in Scrum \textbf{dovrebbe rendere il team
rilassato}. Come il team di Simone, anche il primo team di Cristian (quello fino
alla \emph{vertical slice}) \textbf{non ha mai fatto crunch}. Il nuovo team sì,
per i grossi cambiamenti già raccontati: la camera è cambiata completamente, e
``cambiare la camera in un gioco è fare un altro gioco''. Cristian aveva già
spostato la scadenza (``invece di finire a marzo finiamo ad aprile''), ma il team
si è sentito comunque in dovere di risolvere prima lo scompenso creato da quei
cambiamenti.

In generale ci si dovrebbe sentire meno stressati per un motivo semplice: chi
guida sta dando i task che lo sviluppatore ha chiesto, verbalmente o dimostrando
di saperli fare. ``So che in due settimane devo fare cinque cose, faccio queste
cinque cose'': non c'è la pressione di dover fare il più possibile. Se poi uno è
perfezionista e passa da cinque cose a quindici perché ``ha tempo'', il problema
(``sembra brutto dirlo'') è suo, nasce nel singolo individuo; e chi guida deve
basarsi sul team, sulle performance del team.

\subsection{Mai la velocity sui singoli}

Da qui un consiglio sulla velocity: \textbf{non calcolate mai, mai, mai la
velocity delle singole persone}, e nemmeno i task fatti dai singoli. Lo si può
fare per sé, ma tenendolo ``chiuso in un'agenda scritta a penna'': se un membro
del team scopre che si stanno valutando le sue prestazioni individuali,
\textbf{si sente giudicato} e soprattutto non sa più come dire che non vuole
fare una cosa.

A Cristian è successo: ``mi hanno sgamato''. Si è dovuto scusare, anche se non
erano scuse: ha detto la verità, cioè che lo faceva più per capire lui la
produzione che per valutare le loro performance. Ma spiegarlo è complicato,
quindi è meglio non farlo: tanto la velocity si basa sulle performance del team,
e i dati individuali non servono. Se li si vuole per sé, per capire a chi dare
un task piuttosto che a un altro (se una persona in un'ora fa dieci
implementazioni e un'altra una, ci si regola quando si selezionano i task),
non bisogna assolutamente farli ricadere sulle persone.

\subsection{Primo aspetto: adattare Scrum al team}

Per una produzione senza stress servono due aspetti chiave. Il primo è che
\textbf{Scrum va adattato al team}: si prende la base di Scrum (quella di Diego,
dei libri, di Simone e di Cristian), la si \textbf{testa}, si \textbf{analizzano}
la velocity e tutto il resto, la si \textbf{tara} (``\emph{daily scrum}? No,
facciamolo ogni tre giorni, ogni settimana, o non facciamolo''), la si
\textbf{riprova}, e si ripete il ciclo per tutta la produzione. Alcuni esempi da
\emph{L'Inf}:
\begin{itemize}
  \item i \textbf{daily scrum} li fanno \textbf{ogni tre giorni}, ``se va
        bene'', e dipende anche dal tipo di produzione;
  \item la \textbf{lunghezza delle sprint} è flessibile (una, due o tre
        settimane) e dipende:
        \begin{itemize}
          \item dalla \textbf{disponibilità del team}: se tutto il team non
                c'è, è inutile fare una sprint di una settimana;
          \item soprattutto dalla \textbf{fase di produzione}: quando si
                implementano script e in una settimana si riescono a
                implementare tantissime meccaniche, la sprint dura una
                settimana, al termine della quale si controlla ciò che è stato
                fatto e si sceglie che cos'altro implementare; quando si passa al
                level design, con lavori molto dispendiosi in termini di tempo,
                le sprint si diluiscono, e infatti in quel momento sono a sprint
                di due settimane, concentrati sul completamento dei
                livelli\footnote{Nella trascrizione ``sull'ultima mento,
                sull'ultima azione, non mi ricordo, scusate. Per completare il
                livello'': forse ``l'ultima milestone''.};
        \end{itemize}
  \item la velocity si calcola \textbf{sui task e sul loro peso, non sulle ore
        lavorative}, come diceva Simone. Soprattutto per degli studenti, che
        non hanno ore fisse. In un ambiente di lavoro in cui tutti sono
        disponibili dalle 9 alle 18 potrebbe cominciare ad avere senso:
        ``secondo me no, ma ve lo dico per correttezza''.
\end{itemize}

\subsection{Secondo aspetto: tutti devono conoscere il mindset}

Il secondo aspetto è che \textbf{tutti i membri del team devono conoscere il
mindset di Scrum}\footnote{A questo punto Cristian chiede a Diego cinque minuti
in più per chiudere bene, e Diego acconsente.}. Per due motivi:
\begin{itemize}
  \item se tutti lo conoscono, il rischio di cadere nel Waterfall o nel
        perfezionismo si riduce, perché ognuno sa che cosa è giusto e che cosa
        è sbagliato;
  \item tutto il team ha \textbf{lo stesso goal}: se la sprint chiede il block
        out, tutti sanno che, secondo la mentalità agile, va bene farlo ma non
        deve essere perfetto. Bisogna fare i \emph{done}, non i
        \emph{perfect}.
\end{itemize}
Come far capire a tutti la mentalità di Scrum? Cristian usa due modi.
\begin{enumerate}
  \item \textbf{Parlarne costantemente}, a ogni riunione, come sta facendo con
        la classe: sia per migliorare lui stesso e farsi entrare il mindset,
        sia per trasmetterlo agli altri.
  \item Il più efficace: \textbf{lo shipping}. Il team riesce a consegnare
        quasi costantemente. Ci sono momenti in cui un problema cade
        all'ultimo, come in quella settimana, con alcuni strumenti
        particolari che li hanno rallentati, ma si risolve tutto. Dimostrando
        di riuscire a \textbf{mantenere le milestone}, il team ha acquisito
        fiducia in lui e ha capito quanto è potente Agile. Nella
        \emph{vertical slice} consegnavano le build addirittura prima di
        quanto richiesto; e la consegna prima di Natale, nota da settembre, è
        stata rispettata nonostante il cambio di camera e tutto il resto. Il
        messaggio: \textbf{non è importante che ciò che si consegna sia
        perfetto, ma che rispetti le condizioni di accettabilità} che ci si è
        dati.
\end{enumerate}
Sono i due metodi che funzionano meglio finora; se ne troverà altri, li
condividerà: ``tutto è in continua evoluzione''.

\subsection{Il post-mortem}

Il \emph{post-mortem} è \textbf{la garanzia di capire se ciò che si sta facendo
è corretto}: serve a vedere come stanno andando il team e la produzione, se il
mindset di Scrum è stato compreso o no, e ad agire di conseguenza.

\textbf{Quando farlo} lo si decide liberamente: a ogni sprint, a ogni gruppo di
sprint, a ogni milestone.

Ciò che Cristian sottolinea è che \textbf{nel post-mortem devono parlare tutti},
anche lo sviluppatore che non vuole mai parlare: bisogna dargli il tempo per
farlo. Un modo è \textbf{organizzarlo in anticipo}: se è tra una settimana, lo si
comunica la settimana prima. Cristian prepara sempre un \textbf{documento della
riunione} prima e lo aggiorna dopo, con una sorta di elenco: sembra burocratico,
ma aiuta a mantenere il focus. Bisogna far capire a ciascuno che tra una
settimana parlerà, e non può sottrarsi (se si assenta ci perde lui). E bisogna
dirgli di \textbf{prepararsi} qualcosa da dire, magari scrivendolo o facendo una
piccola revisione: se ti viene in mente una cosa da dire la settimana dopo e non
la scrivi, finirai per pensare ``sai che c'è, non la dico''; se l'hai scritta,
non puoi dimenticartela, e hai qualcosa di cui parlare. È fondamentale perché
bisogna tenere conto di \textbf{tutti i tipi di personalità}: non solo i
giocatori hanno personalità diverse, anche gli sviluppatori hanno personalità
con bisogni diversi.

\subsection{La sincerità}

Cristian propone di tagliare la parte ``sdolcinata'', ma Diego vuole sentirla
per intero (``non si sfugge al miele''; ``poi mi manderete gli IBAN e vi pagherò
il tempo perso''). Come Diego sa, Cristian cerca sempre di metterci ``qualcosa di
filosofico''. Per \textbf{sincerità} intende che, in qualsiasi momento della
produzione, \textbf{tutti i membri del team devono sentirsi protetti in ciò che
dicono} e liberi di esprimere opinioni e feedback. È importantissima sia nella
produzione sia nel team.

\paragraph{Nella produzione.} Senza sincerità la produzione va a picco; si può
comunque consegnare, non è obbligatoria, ma con la sincerità si realizza un
prodotto migliore. Se a inizio produzione si pensa che una meccanica ``faccia
schifo'' (prende un \emph{player type} completamente fuori target e richiede
troppo lavoro per essere implementata), \textbf{è quello il momento di dirlo}.
Altrimenti ``il karma vuole che tu stesso che pensi questa cosa sarai colui che la
andrà a implementare'', e la farai male, perché non vuoi farla. Se non l'hai
detto, o l'hai detto troppo tardi, ``ti mangerai le mani''.

E il \textbf{``ve l'avevo detto''} va risparmiato: anche se il giorno dopo arriva
la conferma di ciò che si era detto, anche se è umano, va evitato, perché crea
attriti nel team che non si risolvono più (Cristian parla per esperienza
diretta). È lo stesso spirito del \emph{disagree and commit}
(sezione~\ref{sec:z2-disagree}).

\paragraph{In Scrum.} Se uno sviluppatore dice ``questa settimana so fare 15
task'' e ne fa 7, la prima volta passa. La settimana dopo dice di nuovo ``ne
faccio 15, tranquillo'', magari inventando scuse (``quel problema non ce l'ho
più''). Una volta, due volte; \textbf{alla terza bisogna chiedere sincerità}: qual
è il problema? Non hai tempo?

Perché sul tempo, ``ragazzi'', \textbf{è importante dire quanto tempo si ha, non
garantirlo}. Se il team (cinque persone, nel caso di Cristian) dice di avere due
ore al giorno, lui pianifica una produzione da dieci ore al giorno; se il team
dice di avere otto ore al giorno, pianifica di conseguenza (``80
ore''\footnote{Così nella trascrizione; con cinque persone a otto ore al giorno
il totale sarebbe 40 ore, oppure 80 su due giorni. Il senso dell'esempio non
cambia.}), e se poi in realtà le ore sono dieci, gli diranno ``ci dai troppi
task''.

Bisogna essere sinceri \textbf{prima di tutto con se stessi}: se un task richiede
due giorni, si dice due giorni; magari tre, perché un ``tempo cuscinetto'' ben
venga, ma lo si dice per com'è. E non bisogna proteggersi dal timore che ``se
dico la verità nessuno si fiderà più di me'' o ``mi daranno più task'': è compito
dello scrum master, ma di tutto il team, far capire che \textbf{si è tutti sulla
stessa barca}.

\paragraph{Nel team.} Gli attriti generati dalla mancanza di sincerità si
trasformano in \textbf{relazioni migliori} quando c'è sincerità. Le
incomprensioni vanno \textbf{risolte subito}, affrontandole; magari porteranno a
un litigio, se proprio quello è il finale, ma tra persone mature si dovrebbe
riuscire a risolvere in modo pacifico. Per questo, sulla slide, la sincerità è
definita \textbf{``rischiosa ma necessaria''}.

% ---------------------------------------------------------------------
\section{Chiusura della lezione}
\label{sec:z3-chiusura}

Cristian salta la slide di riepilogo (la lezione si potrà rivedere online) e
lascia i suoi contatti: \textbf{Telegram}, dove molti studenti lo trovano già nei
gruppi (in particolare quelli dell'accademia online), oppure l'\textbf{email},
che Diego può fornire. Si rende disponibile per qualsiasi cosa, sulla lezione o
su un'idea da sottoporgli (``ho pensato così, secondo te può funzionare?''),
precisando di non volersi sostituire in alcun modo a Diego. Si scusa per i dieci
minuti finali ``di corsa''.

Diego lo ringrazia per la disponibilità, ricordando che anche lui, da studente,
ha un carico di lavoro notevole. L'obiettivo della lezione era far vivere alla
classe un'esperienza e dare informazioni simili a quelle di chi ha fatto lo
stesso percorso, ma ha \textbf{internalizzato Agile} meglio di altri: tutti hanno
avuto le stesse lezioni, ma Cristian e Simone sono quelli che le hanno capite
meglio. Ci sono molti \emph{insight} utili in entrambi gli interventi, e i punti
chiave di Agile sono stati ripetuti più volte, ``perché alla fine Agile è
comprendere i punti chiave''. ``Grazie ancora, e buon pranzo.''\footnote{La
trascrizione si chiude con la parola ``grazie'' ripetuta diverse centinaia di
volte: è un artefatto tipico della trascrizione automatica su un tratto di
audio silenzioso o di saluti sovrapposti, non un contenuto della lezione.}

% ---------------------------------------------------------------------
\section*{Punti chiave del capitolo}
\addcontentsline{toc}{section}{Punti chiave del capitolo}

\begin{itemize}
  \item \textbf{Agile è un mindset, non un framework}: non ha regole, ha una
        logica da \textbf{internalizzare}, come una lingua che a un certo punto
        si smette di tradurre. I quattro valori non vanno seguiti alla lettera:
        vanno adattati al team e al progetto, e non è detto che Agile serva.
  \item I \emph{framework} (Scrum, Kanban, \dots) hanno regole e si basano
        sull'\textbf{inspect and adapt}, cioè sulle iterazioni. Si parte
        rispettandoli al 100\%, si capisce che cosa non funziona e perché, e
        \textbf{ci si costruisce il proprio}.
  \item \textbf{Meno riunioni possibili}: unire planning e post-mortem, stare in
        25--30 minuti, report da un minuto. \textbf{Nelle riunioni non si
        risolvono i problemi}: si fissa un appuntamento con chi può aiutare. Un
        \textbf{timekeeper} può far rispettare i tempi.
  \item Le \textbf{sprint} (una o due settimane) sono l'unità di tempo: le
        milestone si misurano in sprint, non in mesi.
  \item \textbf{Consistency}: una cerimonia saltata fa crollare le altre a
        effetto domino. Per questo ciò che non funziona per il team (spesso il
        \emph{daily scrum} o il \emph{planning poker}) va \textbf{tagliato},
        non trascinato.
  \item Il \textbf{backlog} si fa \textbf{a scaglioni}, per milestone, partendo
        dalle \emph{epic} e dai documenti (visione, GDD, \emph{versioning} con
        tutti gli ambiti, LDD). Si organizza per feature, per ambito con un
        \textbf{color code} condiviso, per milestone e sotto-milestone o per
        priorità, tutto in un solo strumento.
  \item Ogni task (e ogni sprint) ha un \textbf{obiettivo verificabile} e un
        \textbf{criterio di accettazione}: ``completare il livello 5'' non lo è.
        Le \emph{user story} aiutano, ma non sono obbligatorie.
  \item \textbf{Mai stimare in ore}: si stima in modo \textbf{relativo}, con la
        sequenza di \textbf{Fibonacci}, partendo dagli estremi. I valori sono
        arbitrari ma \textbf{consistenti}, e quindi utilizzabili.
  \item La \textbf{velocity} (somma dei pesi dei task \emph{done} per sprint,
        confrontata con lo stimato) \textbf{prevede il futuro}: dice se gli
        strumenti di stima funzionano, se lo \emph{scope} è raggiungibile e come
        sta il team. Permette di tagliare per tempo e di arrivare in fondo
        \textbf{senza crunch}. L'andamento a dente di sega è normale.
  \item Le prime sprint \textbf{si sbagliano tutte}, ed è giusto: la quantità si
        tara sulle disponibilità e sulle capacità attuali, prevedendo i
        \textbf{tempi di apprendimento}.
  \item Si può \textbf{progettare il ritmo dello sviluppo}: dopo periodi di
        lavoro invisibile, sprint che danno soddisfazione (nuovi nemici, VFX,
        \emph{juiciness}) per il morale del team.
  \item I \textbf{blocker} vanno trovati subito e \textbf{analizzati alla
        radice} prima di essere risolti: la soluzione più rapida (i muri
        invisibili) può costare settimane.
  \item I \textbf{feedback costanti} abbattono la produttività: vanno raccolti in
        momenti specifici. Il feedback si può accettare o rifiutare.
  \item \textbf{Flessibilità} in tutto, soprattutto negli elementi innovativi:
        non prevedere tutto nel dettaglio (è Waterfall), non fermare mai la
        produzione per riscrivere backlog e documenti, lavorare sempre per
        \textbf{incrementare il prodotto}.
  \item Il \textbf{perfezionismo} crea blocker, lavoro buttato e ansia negli
        altri: si fanno i \emph{done}, non i \emph{perfect}.
  \item \textbf{Il team viene prima di Scrum}. La velocity non si calcola mai sui
        singoli. Tutti devono conoscere il mindset: parlarne a ogni riunione e,
        soprattutto, \textbf{consegnare} rispettando le milestone.
  \item Nel \textbf{post-mortem} parlano tutti: va annunciato prima, con un
        documento, chiedendo a ciascuno di prepararsi.
  \item Alla base di tutto ci sono la \textbf{fiducia} reciproca e la
        \textbf{sincerità}: dire subito ciò che non va, dire quanto tempo si ha
        (non garantirlo), evitare il ``ve l'avevo detto''. La sincerità è
        \textbf{rischiosa ma necessaria}.
\end{itemize}
